\chapter{Norms of lifts and the glued integral structure}\label{chapter-integral-structure}

This chapter bounds the auxiliary lattice of Chapter~\ref{chapter-inert-prime-balance} from
above and proves the three properties of its glued integral structure that
Chapter~\ref{chapter-assembly} needs. Each theta lift of Chapter~\ref{chapter-unitary-theta}
multiplies Petersson metrics by explicit scalars, and its image lattice has bounded index at the
places where the Hecke characters are twisted; with the covolume bound of
Chapter~\ref{chapter-covolumes} these give the division inequality
$\mu(CC)+h_u\le\mathrm{MAIN}(u)+\text{error}$ for the lattice $CC$ glued across tuples of Hecke
characters. The columns of the glued structure are nearly orthogonal, by bounds for Poincaré
series, and its row space has nearly full dimension, by Green's description of the cuspidal
characters of $GL_r(\mathbf F_g)$.

\section{Indefinite-side data}\label{section-is-indefinite-data}

\begin{definition}\label{is-indefinite-datum}
Let $K$ be a number field with an involution $a\mapsto\bar a$, and let $\mathcal O_K$ be its ring
of integers. For a finite place $w$ of $K$ write $\bar w$ for the conjugate place, $p_w$ for its
residue characteristic, $e_w=v_w(p_w)$ for its ramification index and
$\ell_w=\log Nw/[K:\mathbf Q]$. A \emph{metric} on $K^{d}$ is a family $X=(X_v)$ of invertible
complex $d\times d$ matrices indexed by the infinite places $v$ of $K$; it gives $K^{d}$ the
Hermitian norm $|y|_{X,v}=|X_v^{-1}\sigma_v(y)|$ at $v$, where $\sigma_v$ is the embedding attached
to $v$. An \emph{indefinite-side datum} over $K$ consists of the following.
\begin{enumerate}
\item Base data: a subfield $K_0\subset K$, an infinite place $v_0$, a set $\mathcal T$ of
  \emph{tuples} with a base tuple $\mu_0\in\mathcal T$, a finite set $\Pi$ of \emph{parts} and a
  finite set $J$ of \emph{lifts}.
\item A definite side: for each tuple $t$, a Hermitian matrix $H_V(t)$ of size $n_V(t)$, a lattice
  $M_t\subset K^{n_V(t)}$, a matrix $P_t\in M_{d_Q(t)\times n_V(t)}(K)$, a metric $X_Q(t)$ on
  $K^{d_Q(t)}$ and subspaces $Q_p(t)\subset K^{d_Q(t)}$ for $p\in\Pi$; for each lift $j$ and tuple
  $t$, a Hermitian matrix $H_T(j,t)$ of size $d_Z(j,t)$, a metric $X_T(j,t)$, a lattice
  $T_{j,t}\subset K^{d_Z(j,t)}$ and subspaces $Z_p(j,t)$ for $p\in\Pi$.
\item Indefinite data: a map $w\mapsto\mu_w$ from the finite places of $K$ to $\mathcal T$; a
  finite set $\mathrm{Tw}$ of finite places, the \emph{twisted places}; an integer $D\ge0$, a metric
  $X_E$ on $K^{D}$, subspaces $E_p\subset K^{D}$ for $p\in\Pi$ and matrices
  $\iota_t\in M_{D\times d_Q(t)}(K)$ for $t\in\mathcal T$; for each lift $j$, an integer $d_W(j)$,
  a metric $X_W(j)$ on $K^{d_W(j)}$, subspaces $W_p(j)$, a lattice $M_j\subset K^{d_W(j)}$ and
  matrices $\theta_{j,t}\in M_{d_W(j)\times d_Z(j,t)}(K)$ for $t\in\mathcal T$; and real numbers
  $\kappa_v$, $\lambda_{j,t,p,v}$ and $\mathrm{err}_0(w,j)$, for infinite places $v$, lifts $j$,
  tuples $t$, parts $p$ and finite places $w$.
\end{enumerate}
For a lift $j$, a tuple $t$ and a finite place $w$, the \emph{index} of the $j$-th image lattice at
$w$ is
\[
 I_w(j,t)=\operatorname{len}_w\big(M_j/(M_j\cap\theta_{j,t}T_{j,t})\big)
 -\operatorname{len}_w\big(\theta_{j,t}T_{j,t}/(M_j\cap\theta_{j,t}T_{j,t})\big),
\]
where $\operatorname{len}_w$ is the length of the localization at $w$. The \emph{glued lattice}
$RR\subset K^{D}$ is the $\mathcal O_K$-module of vectors that lie, for every finite place $w$, in
the localization at $w$ of $\iota_{\mu_w}P_{\mu_w}M_{\mu_w}$; it carries the metric $X_E$.
\end{definition}
\leanref{Formalism/LocalToGlobal/IndefiniteForms/IndefSide.lean}{Formalism.L2.Indef.DefBase}
\leanref{Formalism/LocalToGlobal/IndefiniteForms/IndefSide.lean}{Formalism.L2.Indef.DefSide}
\leanref{Formalism/LocalToGlobal/IndefiniteForms/IndefSide.lean}{Formalism.L2.Indef.IndefData}
\leanref{Formalism/LocalToGlobal/IndefiniteForms/IndefSide.lean}{Formalism.L2.Indef.IndefData.Iw}
\leanref{Formalism/LocalToGlobal/IndefiniteForms/IndefSide.lean}{Formalism.L2.Indef.IndefData.RR}

In the data used below the matrices $\iota_t$ and $\theta_{j,t}$ are bijective, $\kappa_v>0$ and
$\lambda_{j,t,p,v}>0$, and the scalars have the following meaning. Whenever $t=\mu_0$ or $v=v_0$,
the matrix $\theta_{j,t}$ multiplies the Hermitian form of $X_T(j,t)$ at $v$ by $\lambda_{j,t,p,v}$
on $\theta_{j,t}^{-1}(W_p(j))$, and $\iota_t$ multiplies the Hermitian form of $X_Q(t)$ at $v$ by
$\kappa_v\prod_j\lambda_{j,t,p,v}$ on $\iota_t^{-1}(E_p)$. So $\lambda_{j,t,p,v}$ is the factor by
which the $j$-th lift changes the Petersson metric on the $p$-th part at $v$, and
$\mathrm{err}_0(w,j)$ is an error term at the twisted place $w$.
\leanref{Formalism/LocalToGlobal/IndefiniteForms/IndefSide.lean}{Formalism.L2.Indef.IndefData.RowLatticeBoundForBasis}

\begin{definition}\label{is-numerical-bounds}
Let an indefinite-side datum over $K$ be given, with tuples $\mathcal T$ and lifts $J$
(Definition~\ref{is-indefinite-datum}). \emph{Numerical bounds} for it consist of real functions
$F_0(w,j)$, $J_0(t,w,j)$, $E_1(j)$, $\ell_3(j)$, $\ell_6(j)$ of finite places $w$, tuples $t$ and
lifts $j$, a real number $h_N$, and two real numbers $\ell_a,\delta_e$ that bound the passage from
rows to columns. \emph{Scale data} consist of real functions $F_0(w,j)$, $J_0(t,w,j)$,
$\ell_3^{m}(j)$, $\ell_6^{m}(j)$ and real numbers $E_1^{u},\ell_3^{u},\ell_6^{u},x,b,\ell_a,\delta_e$.
For real constants $C$ and $C_H$, scale data give the numerical bounds $N(C,C_H)$ with the same
$F_0,J_0,\ell_a,\delta_e$ and
\[
 E_1(j)=CE_1^{u},\qquad \ell_3(j)=\ell_3^{m}(j)+C\ell_3^{u},\qquad
 \ell_6(j)=\ell_6^{m}(j)-C\ell_6^{u},\qquad h_N=x^{C_H}b.
\]
The datum satisfies the \emph{covolume bound} for numerical bounds $N$ if
$\operatorname{lcv}(M_j)/D\le\ell_3(j)$ for every lift $j$, where
$\operatorname{lcv}(M)=\frac{2}{[K:\mathbf Q]}\log\operatorname{covol}(M)$ is the normalized
logarithmic covolume of $M_j$ in the metric $X_W(j)$.
\end{definition}
\leanref{Formalism/LocalToGlobal/IndefiniteForms/IndefSide.lean}{Formalism.L2.Indef.IndefNum}
\leanref{Formalism/LocalToGlobal/IndefiniteForms/IndefSide.lean}{Formalism.L2.Indef.IndefScale.num}
\leanref{Formalism/LocalToGlobal/IndefiniteForms/IndefSide.lean}{Formalism.L2.Indef.IndefData.L3}

Chapter~\ref{chapter-covolumes} proves the covolume bound for the lattices of lifted forms
(Theorem~\ref{cv-covolume-bound}).

\begin{definition}\label{is-lower-norm-bound}
Let an indefinite-side datum over $K$ be given, with base tuple $\mu_0$, parts $\Pi$, lifts $J$ and
scalars $\lambda_{j,t,p,v}$ (Definition~\ref{is-indefinite-datum}), and let $N$ be numerical bounds
for it with functions $\ell_6(j)$ (Definition~\ref{is-numerical-bounds}). The datum satisfies the
\emph{lower norm bound} for $N$ if
\[
 \log\lambda_{j,\mu_0,p,v}\ge\ell_6(j)
\]
for every lift $j$, every part $p$ and every infinite place $v$ of $K$.
\end{definition}
\leanref{Formalism/LocalToGlobal/IndefiniteForms/IndefSide.lean}{Formalism.L2.Indef.IndefData.L6}

\begin{definition}\label{is-surjectivity-bound}
Let an indefinite-side datum over $K$ be given, with twisted places $\mathrm{Tw}$, tuple map
$w\mapsto\mu_w$, integer $D$, error terms $\mathrm{err}_0(w,j)$ and indices $I_w(j,t)$
(Definition~\ref{is-indefinite-datum}), and let $N$ be numerical bounds for it with functions
$F_0,J_0,E_1$ (Definition~\ref{is-numerical-bounds}). The datum satisfies the \emph{surjectivity
bound} for $N$ if the following hold.
\begin{enumerate}
\item For every twisted place $w$,
\[
 \sum_{j\in J}\big(I_w(j,\mu_w)+I_{\bar w}(j,\mu_w)\big)
 \le\Big(\sum_{j\in J}\big(2e_wF_0(w,j)+e_w\,\mathrm{err}_0(w,j)\big)+\frac{2e_w}{p_w^{2}}\Big)D.
\]
\item $\mathrm{err}_0(w,j)\ge0$ for every twisted place $w$ and every lift $j$.
\item For every lift $j$,
\[
 \sum_{w\in\mathrm{Tw}}\ell_w\big(e_w\,\mathrm{err}_0(w,j)+e_wJ_0(\mu_w,w,j)\big)\le E_1(j).
\]
\end{enumerate}
\end{definition}
\leanref{Formalism/LocalToGlobal/IndefiniteForms/IndefSide.lean}{Formalism.L2.Indef.IndefData.ON}

\begin{definition}\label{is-functional-bound}
Let an indefinite-side datum over $K$ be given, with glued lattice $RR$ in the metric $X_E$ and
scalars $\kappa_v$ (Definition~\ref{is-indefinite-datum}), and let $N$ be numerical bounds for it
with real number $h_N$ (Definition~\ref{is-numerical-bounds}). Let $r_\infty$ be the number of
infinite places of $K$. The datum satisfies the \emph{functional bound} for $N$ if
\[
 \mu_{\max}(RR)\le h_N-\frac{1}{r_\infty}\sum_v\log\kappa_v,
\]
where $\mu_{\max}$ is the largest normalized slope of Chapter~\ref{chapter-hermitian-slopes}.
\end{definition}
\leanref{Formalism/LocalToGlobal/IndefiniteForms/IndefSide.lean}{Formalism.L2.Indef.IndefData.FUN}

\section{Families and points}\label{section-is-points}

Families of definite-side data, their size data, classes of families and the phrase
\emph{eventually along} a class are those of Definitions~\ref{pd-family} and~\ref{pd-size-data}.
For a member at scale $x$ and exponent $\varepsilon$ we use the
quantities $b=\exp(x^{\varepsilon})$, $a=xb$ and $X''=X''(n)$ of its scale, and for a place $w$ of
the coefficient field $L$ the numbers $p_w$, $e_w$ and $\ell_w$ of
Definition~\ref{is-indefinite-datum}. For a class $\mathcal A$ of families,
$\mathcal A^{\mathrm{des}}=\operatorname{Des}(\mathcal A_N,3)$ is its class of designed families with
constant $3$ (Definition~\ref{asm-designed}).

The \emph{standard families} form the class $\mathcal H$ of families $\mathcal F$ such that
$\mathcal F$ is designed with constant $3$ in its own class, that is, lies in
$\operatorname{Des}(\mathcal C_N,3)$ for the class $\mathcal C$ of all families
(Definition~\ref{asm-designed}), and in addition satisfies the page and window conditions and the
bound on the first modulus of its duty data (Definition~\ref{asm-duty}) and the condition on the
chosen walls. In particular the stabilizer denominators of $\mathcal F$ divide Minkowski's number,
$C_{YL}\le9$, and eventually $b^{m/2}\le T_u$ and $\log T_u\le C_{T_u}mx^{\varepsilon}$ for every word
$u$.
\leanref{Development/AbelianVarietyTheory/Definitions/Divisors/TupleReadBaseTupleEqFamily.lean}{GlobalTheory.BigTop.FamHyp}
\leanref{Formalism/SharpRestriction/ParamsHeightRankConstant.lean}{Formalism.MainNotions.TuSize}

\begin{definition}\label{is-twisted-division-bound}
Let $\mathcal F$ be a family of exponent $\varepsilon$ (Definition~\ref{pd-family}). For a member
$i$ at scale $x$, with definite-side datum $D$ of ranks $k,m,n$, $f=2n$, write $\eta_j(w)$ for the
local defect of $D$ at the place $w$ for the lift $j$, $D_0$ for the dimension of the base rows of
$D$, $T_2$ for the set of integers that gives the thick places of the host
(Definition~\ref{hf-host}), and
\[
 B_0=f\,m\,b\,x^{\varepsilon},\qquad N_c=\frac{|T_2|\,k\,b\,x^{2\varepsilon}}{n^{2}}.
\]
The family satisfies the \emph{division bound at twisted places} if there are $C$ and $x_0$ such that
for every $x\ge x_0$, every member at $x$ and every lift $1\le j\le m$,
\[
 \sum_{w\in\operatorname{supp},\ \mu_w\ne\mu_0}\ell_w\,e_w\,\eta_j(w)\le C\,D_0\,x^{\varepsilon}(B_0+N_c).
\]
\end{definition}
\leanref{Formalism/SharpRestriction/GlobalTheoremInputs.lean}{Formalism.MainNotions.DIVinN}
\leanref{Formalism/SharpRestriction/GlobalTheoremInputs.lean}{Formalism.MainNotions.etaSum}

\begin{definition}\label{is-point}
Let $\mathcal F$ be a family (Definition~\ref{pd-family}), let $x>x_1$ be a scale and $i$ a member
at $x$, with definite-side datum $D$ and coefficient field $L$. A \emph{point} $p$ of $\mathcal F$ at
$(x,i)$ is a realization of $D$ by indefinite Hermitian data. It consists of two indefinite-side
data $S_0(p)$, the \emph{rows}, and $S_1(p)$, the \emph{columns}, over $L$
(Definition~\ref{is-indefinite-datum}), with the same base data and the same map $w\mapsto\mu_w$,
where $\mu_0$ is the trivial tuple of Hecke characters and $\mu_w$ the tuple attached to the place
$w$; scale data $Z(p)$ (Definition~\ref{is-numerical-bounds}), and so numerical bounds
$N_p(C,C_H)$ for all $C,C_H$; the theta kernels, weights, phases and level data of the realization;
and compatibilities with $D$: the twisted places of $S_0(p)$ are the places $w$ with
$\mu_w\ne\mu_0$, the column pairings of $S_1(p)$ factor through the periods of $D$, the error terms
$\mathrm{err}_0$ of $S_0(p)$ are the canonical ones defined by $F_0$, and the orders of the
stabilizers divide the stabilizer denominator of $D$. The \emph{point conditions} are a further list
of compatibilities of the realization with $D$, its theta kernels and its local data; the
\emph{further conditions} require that the recursion, core and series data of $p$ are tied to
those of $D$, and that a finite place $w$ is twisted whenever, for some lift $j$, one of
$I_w(j,\mu_0)$, $I_{\bar w}(j,\mu_0)$, $F_0(w,j)$, $F_0(\bar w,j)$, or the relative index of
$T_{j,\mu_0}$ and its dual for $H_T(j,\mu_0)$ at $w$ or at $\bar w$, is nonzero; here the indices are
those of the rows $S_0(p)$ and $F_0$ is that of $N_p(0,0)$.
\end{definition}
\leanref{Development/AbelianVarietyTheory/Definitions/Divisors/TupleReadBaseTupleEqFamily.lean}{GlobalTheory.BigTop.PtAt}
\leanref{Development/AutomorphicArithmetic/ThetaCorrespondence/ThetaPairLifts/ThirdPlaceSeries/Realised.lean}{GlobalTheory.BigTop.PtAt.PointConditions'}
\leanref{Development/AbelianVarietyTheory/Statements/PlaceShiftStatements/SingleIndexPlacesAreTwisted.lean}{GlobalTheory.BigOn.More}

\section{Realization data and the bounds for the lifts}\label{section-is-realisation}

\subsection{Complex data and realizations}\label{subsection-is-complex-data}

\begin{definition}\label{is-complex-data}
Let $K$ be a number field with an involution, and let base data over $K$ be given
(Definition~\ref{is-indefinite-datum}(1)), with tuples $\mathcal T$, base tuple $\mu_0$, parts
$\Pi$, lifts $J$ and infinite place $v_0$; write $\sigma_0$ for the embedding attached to $v_0$. A
\emph{definite ambient} consists of $K$-vector spaces $R_t$ of \emph{row forms} and $L_{j,t}$ of
\emph{lift forms}, for tuples $t$ and lifts $j$, each with a family of $K$-linear operators, the
\emph{Hecke operators}, and a subspace, its \emph{Hecke block}. \emph{Complex data} on a definite ambient consist of the following.
\begin{enumerate}
\item Finite-dimensional complex Hilbert spaces $F_t$ with $\sigma_0$-semilinear maps
  $j_t\colon R_t\to F_t$; a finite-dimensional Hilbert space $E$, the \emph{top space}; and linear
  maps $\Theta_t\colon F_t\to E$, the \emph{theta lifts}.
\item Subspaces $E_p\subset E$ for $p\in\Pi$, the \emph{parts}, and subspaces $E^{b}\subset E$ for
  $b\in\{0,1\}$, the \emph{levels}.
\item For each lift $j$: finite-dimensional Hilbert spaces $F_{j,t}$ with $\sigma_0$-semilinear maps
  $j_{j,t}\colon L_{j,t}\to F_{j,t}$, a finite-dimensional Hilbert space $W_j$, linear maps
  $\theta'_{j,t}\colon F_{j,t}\to W_j$, and subspaces $W_p(j)\subset W_j$ for $p\in\Pi$.
\item For each level $b$ and lift $j$: an integer $d_W(b,j)$, for each infinite place $v$ a Hilbert
  space $W_{b,j,v}$ with vectors $w_{b,j,v}(1),\dots,w_{b,j,v}(d_W(b,j))$, a linear map
  $e_{b,j}\colon W_{b,j,v_0}\to W_j$, and an $\mathcal O_K$-submodule $M_{b,j}\subset K^{d_W(b,j)}$;
  and a positive number $n_b$.
\item A map $w\mapsto\mu_w$ from the finite places of $K$ to $\mathcal T$, a finite set
  $\mathrm{Tw}$ of finite places, and real numbers $\kappa_v$, $\lambda_{j,t,p,v}$ and
  $\mathrm{err}_0(w,j)$.
\end{enumerate}
The \emph{definite subspaces} of the complex data are $Z_t=j_t^{-1}(\Theta_t^*E)\subset R_t$ and
$Z_{j,t}=j_{j,t}^{-1}(\theta'^*_{j,t}W_j)\subset L_{j,t}$.

A \emph{face} for the complex data assigns to the definite subspaces, at each level $b$, a definite
side (Definition~\ref{is-indefinite-datum}(2)) with integers $d_Q(t)$ and $d_Z(j,t)$, together with
coordinate maps $c_{b,t}\colon K^{d_Q(t)}\to R_t$ and $c_{b,j,t}\colon K^{d_Z(j,t)}\to L_{j,t}$.
The \emph{indefinite-side datum at level $b$} of the complex data and the face has these base data and
this definite side, the map $w\mapsto\mu_w$, the twisted places $\mathrm{Tw}$, the scalars
$\kappa_v$, $\lambda_{j,t,p,v}$ and the error terms $\mathrm{err}_0(w,j)$ of (5), and the following
objects.
\begin{enumerate}
\item[(a)] Let $i_b(t)\colon\mathbf C^{d_Q(t)}\to E$ send the $k$-th basis vector to
  $\Theta_tj_tc_{b,t}(e_k)$. Then $D=d_Q(\mu_0)$, and $\iota_t\in M_{D\times d_Q(t)}(K)$ is the matrix
  whose image under $\sigma_0$ expresses $i_b(t)$ through $i_b(\mu_0)$.
\item[(b)] The metric $X_E$ is, at $v_0$, a metric for the Gram matrix
  $n_b^{-1}\big(\langle i_b(\mu_0)e_k,i_b(\mu_0)e_l\rangle\big)_{k,l}$, and at every other infinite
  place $v$ it is the metric $X_Q(\mu_0)$ scaled by $\kappa_v\prod_j\lambda_{j,\mu_0,p,v}$ on the
  coordinates of the $p$-th part. The parts are $E_p^K=\{x\in K^{D}:i_b(\mu_0)\sigma_0(x)\in E_p\}$.
\item[(c)] Let $t_b(j,t)\colon\mathbf C^{d_Z(j,t)}\to W_j$ send $e_k$ to
  $\theta'_{j,t}j_{j,t}c_{b,j,t}(e_k)$, and $w_b(j)\colon\mathbf C^{d_W(b,j)}\to W_j$ send $e_k$ to
  $e_{b,j}(w_{b,j,v_0}(k))$. Then $d_W(j)=d_W(b,j)$, $M_j=M_{b,j}$, the metric $X_W(j)$ is at each $v$
  a metric for the Gram matrix of the vectors $w_{b,j,v}(k)$, the parts of $K^{d_W(j)}$ are the
  preimages of the $W_p(j)$ under $w_b(j)\circ\sigma_0$, and $\theta_{j,t}\in M_{d_W(j)\times
  d_Z(j,t)}(K)$ is the matrix whose image under $\sigma_0$ expresses $t_b(j,t)$ through $w_b(j)$.
\end{enumerate}
The datum at level $0$ is the \emph{rows} of the complex data, and the datum at level $1$ its
\emph{columns}.
\end{definition}
\leanref{Formalism/LocalToGlobal/IndefiniteForms/IndefSide.lean}{Formalism.L2.Indef.DefAmb}
\leanref{Formalism/LocalToGlobal/IndefiniteReal/PartsOrth.lean}{Formalism.IndefR.Cx.CxData}
\leanref{Formalism/LocalToGlobal/IndefiniteReal/PartsOrth.lean}{Formalism.IndefR.Cx.CxData.zs}
\leanref{Formalism/LocalToGlobal/IndefiniteReal/PartsOrth.lean}{Formalism.IndefR.Cx.CxData.data}

\begin{definition}\label{is-realisation}
Let $D$ be a definite-side datum (Definition~\ref{th-definite-datum}) with host field $E_h/F_h$,
coefficient field $L$ with the embedding $\sigma_0$ of its base place $v_0$, ranks $m$ and $r$, and
the Hermitian spaces $Y_j$, $X'$, $V_i$ and $W_i$ of Definition~\ref{th-spaces}. Let a family of
definite forms of $D$ be given: a finite group $G$, base data over $L$ with tuples $\mathcal T$ and
the lifts $1,\dots,m$, a definite ambient over these base data whose row forms at $t$ are algebraic
modular forms on $U(X')$ and whose lift forms at $(j,t)$ are algebraic modular forms on $U(V_j)$,
for each tuple, sets of \emph{inert indices} with Hecke operators on the row and lift forms, and for
each lift $j$ a map on types, the \emph{transfer} at $j$. Here a \emph{type} is a system of eigenvalues
of these operators: a Hecke polynomial at each place of the family and a value at each inert index. The
\emph{types of the base} are the types whose Hecke polynomial is $1$ at every place of the family that
is not a place of the row forms at the base tuple $\mu_0$, whose value is $0$ at every inert index that
is not an inert index of those row forms, and whose eigenspace in the row forms at $\mu_0$ contains a nonzero
vector of the base space at the full level, the level of the trivial subgroup of $G$ with the trivial
character. The family has \emph{lift-distinct types} if, for every lift $j$, the transfer at $j$ sends
any two distinct types of the base to types that differ at some place or at some inert index of the
lift forms at $(j,\mu_0)$.

A \emph{realization} of $D$ for this family over a number field $K_s\subset L$ consists of:
\begin{enumerate}
\item level and character data: Hecke characters $\mu^0_j$, $\eta$ and $\xi_t$ with ramified sets
  $Q_t$ as in a level system (Definition~\ref{th-grid}), characters of the residue fields, and lattices
  with adelic matrices in the spaces $Y_j$, $X'$, $V_i$ and $W_i$; a finite set $P_g$ of places of
  $E_h$ and an isomorphism of $G$ with $\prod_{P\in P_g}GL_r(k_P)$, $k_P$ the residue field; maps from
  places of the family to places of $E_h$; and a finite set $S_b$ of places;
\item an archimedean family $\theta$ and a standard character $\psi$ of $\mathbf A_{F_h}$
  (Definition~\ref{th-weil-datum}(1));
\item models of discrete series for $U(Y_j)$, $U(X')$ and, at each embedding $\sigma$ of $L$,
  $U(W_i)$; coefficient spaces for $X'$ and the $V_i$; and a realization of the definite family by
  algebraic modular forms with these coefficients;
\item Hecke modules $H_X(t)$, $H_V(i,t)$, $H_E$ and $H_W(i,\sigma)$, representations of $G$ on $H_E$
  and the $H_W(i,\sigma)$, and $\sigma_0$-semilinear maps from the row forms at $t$ to $H_X(t)$ and from
  the lift forms at $(i,t)$ to $H_V(i,t)$;
\item Haar measures and fundamental domains on the adelic unitary groups of $Y_j$, $X'$, $V_i$ and
  $W_i$;
\item Weil data (Definition~\ref{th-weil-datum}(3)) for the pairs $(Y_k,Y_{k+1})$,
  $(Y_{m-1},X')$ and $(V_i,W_i)$ at each tuple, with vacuum vectors and test vectors for each pair;
\item the iterated theta maps between the models of $Y_{k+1}$ and $Y_k$ and from $X'$ to $Y_{m-1}$,
  their restrictions to spaces of theta multiplicities, the theta lifts
  $\mathrm{Th}_t\colon H_X(t)\to H_E$, matrices $B_V(i)$ at the places of $\mathcal T$, the maps
  $\mathrm{io}_P(i)$ from the lowest $K$-types of the model of $W_i$ and the coefficients of $V_i$ to
  polynomials, and theta maps $\mathrm{th}_F(i,t)$ and $\mathrm{th}_G(i,t)\colon H_V(i,t)\to
  H_W(i,\sigma_0)$;
\item for each lift $i$ an integer $n_W(i)$ and a $K_s$-rational structure $FW_i$ on the spaces
  $H_W(i,\sigma)$, with a character of $G$;
\item a map $w\mapsto\mu_w$ from the finite places of $L$ to $\mathcal T$, a finite set
  $\mathrm{Tw}$ of finite places, a real number $\kappa$, real numbers $\lambda_{i,t,p,v}$, and real
  numbers $\mathrm{err}_0(w,i)$;
\item maps from the inert indices to places of $E_h$, a finite set $S_F$ of finite places of $F_h$,
  \emph{phases} $\phi_i\in\mathbf C$ for the lifts, and real numbers $\lambda^{(3)}_{i,t,p,v}$.
\end{enumerate}
These data are subject to one condition: the base place $v_0$ is the place of the embedding
$\sigma_D\colon L\to\mathbf C$ of the coefficient field of $D$. Of the embeddings of $L$ at $v_0$,
$\sigma_0$ denotes $\sigma_D$, so $\sigma_0=\sigma_D$.
A \emph{face} of the family consists of subgroups $P_0,P_1\subset G$ with characters
$\psi_0\colon P_0\to L^\times$ and $\psi_1\colon P_1\to L^\times$, and, for every choice of definite
subspaces, definite columns, their periods, and scale data (Definition~\ref{is-numerical-bounds}).
The realization and a face give complex data (Definition~\ref{is-complex-data}), whose embedding attached
to $v_0$ is $\sigma_D$ or its complex conjugate; in the second case the complex spaces below carry the
conjugate complex structure. $F_t$ and $F_{i,t}$
are the completions of $H_X(t)$ and $H_V(i,t)$, $E=H_E$ with $\Theta_t=\mathrm{Th}_t$, $W_i=H_W(i,\sigma_0)$
with $\theta'_{i,t}=\mathrm{th}_G(i,t)$, the parts and levels come from the face, and the scalars are
those of (9) and (10). The rows $S_0(R,F)$ and columns $S_1(R,F)$ of the realization at the face are
the rows and columns of these complex data, and $N(R,F;C,C_H)$ are the numerical bounds of the scale
data of the face.
\end{definition}
\leanref{Formalism/LocalToGlobal/Divisibility/HeckeRI.lean}{Formalism.L2.DIV.FamIn3}
\leanref{Formalism/LocalToGlobal/IndefiniteReal/LinksDefs.lean}{Formalism.IndefR.Links.RealDat}
\leanref{Formalism/LocalToGlobal/IndefiniteReal/LinksDefs.lean}{Formalism.IndefR.Links.RealDat.hv0}
\leanref{Formalism/LocalToGlobal/IndefiniteReal/LinksDefs.lean}{Formalism.IndefR.Links.sg0}
\leanref{Formalism/LocalToGlobal/IndefiniteReal/ZeroPartEigenvalueConvention.lean}{Formalism.IndefR.Links.RealDat3}
\leanref{Formalism/LocalToGlobal/IndefiniteReal/ZeroPartEigenvalueConvention.lean}{Formalism.IndefR.Links.Face3}

\begin{definition}\label{is-one-map-data}
Let $R$ be a realization of a definite-side datum $D$ (Definition~\ref{is-realisation}), with the
models of discrete series of the $W_i$ at $\sigma_0$. Let a \emph{core} of $D$ be given: the
spaces of algebraic modular forms on the $V_i$ at the levels of $D$, with their coefficient field.
\emph{One-map data} for $R$ consist of: a \emph{weight family} $\kappa_i$ at each level $i+1$; for each
lift $i$ and embedding $\sigma$ of $L$ a vector $v_0(i,\sigma)$ in the model of $W_i$ at $\sigma$; a
\emph{weight tie} between the core, the realization, the $\kappa_i$ and the $v_0$, and a
\emph{conjugation tie} for it; and for each lift $i$ \emph{one-map data at $i$}: a presentation of the
$i$-th lift of the core at the base tuple as a single theta map, consisting of a coefficient field
containing that of the core, an identification of the space of the core at level $i+1$ with $W_i$,
classes of lattices in $V_i$ with their stabilizer orders and volumes, local constants at the bad
places, a polynomial realization of the coefficients of $V_i$ with a highest weight vector, and the
coefficient functions and adelic theta data of the presentation. One-map data at $i$ give the
$i$-th lift \emph{as one map with constants $c$ and phase $\phi$} if the presentation is equivalent to
the lift of the realization for the weight $\kappa_i$, agrees with the theta lift of the core with the
constants $c$ and the phase $\phi$, and has rational phase. The data are \emph{pinned} at $i$ if
their coefficient field is the image of $L$ and they are pinned to the core at the vector
$v_0(i,\sigma_0)$.
\end{definition}
\leanref{Formalism/LocalToGlobal/IndefiniteReal/OneMapT.lean}{Formalism.IndefR.FWT.OMDat}
\leanref{Formalism/LocalToGlobal/IndefiniteReal/OneMapT.lean}{Formalism.IndefR.FWT.OMDat.OneMapOf}
\leanref{Formalism/LocalToGlobal/IndefiniteReal/OneMapT.lean}{Formalism.IndefR.FWT.OMPins}

\begin{definition}\label{is-bounded-rows}
Let $R$ be a realization of a definite-side datum $D$ with one-map data (Definitions~\ref{is-realisation}
and~\ref{is-one-map-data}), and let $\phi_i$ be phases for the lifts. The realization has
\emph{bounded rows} with these phases if:
\begin{enumerate}
\item for every lift $i$ there is a vector $u$ of the model of $W_i$ at $\sigma_0$ that is
  \emph{extremal}, that is, of norm $1$, in the lowest $K$-type and of extremal weight at every
  place, and that \emph{pairs to the delta}: there are lowest-$K$-type vectors $u_\rho$ with
  components those of $u$ and a coefficient vector $w$ of norm $1$ for $V_i$ such that
  $\mathrm{io}_P(i)(\otimes_\rho u_\rho,w)$ is the delta polynomial $\Delta_i$ of $R$;
\item for every lift $i$ the $i$-th lift is given by the one-map data at $i$ as one map for some
  constants and phase; and, if the ramified primes $Q_t$ of the tuples (Definition~\ref{th-grid})
  satisfy the \emph{guard}, that two primes in one $Q_t$ with the same residue characteristic lie over
  the same prime of the subfield $K$ of the host and that no prime in a $Q_t$ contains the integer
  $Q_Sq\,N(\mathfrak c_K)$ of the host, then there are one-map data at $i$, pinned, over a field
  containing $L$, that give the $i$-th lift as one map with the canonical constants of the core and
  the phase $\phi_i$.
\end{enumerate}
\end{definition}
\leanref{Development/AutomorphicArithmetic/ThetaCorrespondence/ThetaPairLifts/ThirdPlaceSeries/IndefRRealTB.lean}{Formalism.IndefR.IndefiniteRealisationBoundedRows}
\leanref{Formalism/LocalToGlobal/IndefiniteReal/ExistsChosenExtremalVector.lean}{Formalism.IndefR.FWT.IoExt}
\leanref{Formalism/LocalToGlobal/IndefiniteReal/OneMapT.lean}{Formalism.IndefR.FWT.OneMapRowR}

\begin{definition}\label{is-primitive-rows}
Let $R$ be a realization of a definite-side datum $D$ for a family of definite forms, and $F$ a face
(Definition~\ref{is-realisation}), with complex data $X$ (Definition~\ref{is-complex-data}). The
realization has \emph{primitive rows} at $F$ if the following hold.
\begin{enumerate}
\item (Rational Hecke operators.) For every tuple $t$ and lift $j$, the Hecke operators and the inert
  Hecke operators on the row forms at $t$ and the lift forms at $(j,t)$ have matrices with entries in
  the subfield of $L$ on which the involution is complex conjugation at every embedding, and each has
  an adjoint among them for the class-sum pairing.
\item (Intertwining.) There are operators on $E$ and on each $W_j$ that the theta lifts
  $\Theta_t$ and $\theta'_{j,t}$ intertwine with the Hecke operators; at the base tuple they commute
  in pairs, and the parts $E_p$ and $W_p(j)$ are the sums of their joint generalized eigenspaces for
  the systems of eigenvalues of the part $p$.
\item (The group.) $G$ acts unitarily on $E$ and on each $W_j$, and the theta lifts are
  $G$-equivariant.
\item (Blocks.) The definite subspaces of $X$ lie in the Hecke blocks of the ambient.
\item (Rationality tests.) There are finitely many functionals separating the points of $E$, and of
  each $W_j$, whose values on the theta lifts of row and lift forms, and on the coordinate vectors
  $w_0(j)$ of the rows, lie in $\sigma_0(L)$.
\item (Surjectivity of the lifts.) For any operators intertwining the lifts, the adjoints of the
  $\theta'_{j,\mu_0}$ map the good parts of the $W_j$ onto the good parts of the lift forms.
\item (Positivity.) $d_Q(\mu_0)>0$ at both levels, $\kappa>0$, and $\lambda^{(3)}_{i,t,p,v}>0$.
\item (Level types.) Let $i^{\mathrm{pr}}(t)$ and $t^{\mathrm{pr}}(j,t)$ be the coordinate maps
  $i_0(t)$ and $t_0(j,t)$ of Definition~\ref{is-complex-data} for the face with the trivial subgroup
  and character. At level $b$ the level $E^{b}$ is the intersection of the image of
  $i^{\mathrm{pr}}(\mu_0)$ with the subspace on which $P_b$ acts through $\sigma_0\circ\psi_b$; the
  maps $i^{\mathrm{pr}}(t)$ are injective and have the same image for all $t$; and the image of
  $w_0(j)$ is the intersection of the sum over $t$ of the images of $t^{\mathrm{pr}}(j,t)$ with the
  subspace on which $P_0$ acts through $\sigma_0\circ\psi_0^{-1}$.
\item (Scalars.) For every lift $i$, tuple $t$ and part $p$,
  $\mathrm{th}_G(i,t)\,\mathrm{th}_G(i,t)^*=\lambda^{(3)}_{i,t,p,v_0}$ on $W_p(i)$, and
  $\mathrm{Th}_t\mathrm{Th}_t^*=\kappa\prod_i\lambda^{(3)}_{i,t,p,v_0}$ on $E_p$.
\item (Parts.) At both levels $b$ the spaces $E_p\cap E^{b}$ are orthogonal and span $E^{b}$; the
  spaces $W_p(j)\cap w_0(j)(\mathbf C^{d_W(0,j)})$ are orthogonal and span the image of $w_0(j)$; and
  $\dim\big(W_p(j)\cap w_0(j)(\mathbf C^{d_W(0,j)})\big)=\dim(E_p\cap E^{0})$ for every $j$ and $p$.
\item (Unitary structure.) The rational structures $FW_i$ carry $L$-independent vectors to
  $\mathbf C$-independent ones, span, and are $G$-equivariant up to the character of $FW_i$.
\item (Lattices and conjugates.) The modules $M_{0,j}$ are lattices, and at every infinite place
  $v\ne v_0$ the matrix $\theta_{j,\mu_0}$ of the rows multiplies the Hermitian form of
  $X_T(j,\mu_0)$ at $v$ by $\lambda_{j,\mu_0,p,v}$ on the vectors whose image under
  $t_0(j,\mu_0)\circ\sigma_0$ lies in $W_p(j)$.
\end{enumerate}
\end{definition}
\leanref{Formalism/LocalToGlobal/IndefiniteReal/BaseLiftsPackageOfPrimitiveLinks.lean}{Formalism.IndefR.Links.PrimRows3}
\leanref{Formalism/LocalToGlobal/IndefiniteReal/PrimitiveStatementsWithPassage.lean}{Formalism.IndefR.Cx.HeckeIntertwiningDefinite}
\leanref{Formalism/LocalToGlobal/IndefiniteReal/PrimitiveStatementsWithPassage.lean}{Formalism.IndefR.Cx.LevType3}
\leanref{Formalism/LocalToGlobal/IndefiniteReal/PartsOrth.lean}{Formalism.IndefR.Cx.CxData.GalD}

\begin{definition}\label{is-row-side-family}
Let $D$ be a definite-side datum (Definition~\ref{th-definite-datum}) with coefficient field $L$,
division data and base point $p_0$ (Definition~\ref{pd-period-data}). A \emph{row-side family} for
$D$ consists of: a finite group $G$; a family of definite forms (Definition~\ref{is-realisation})
over $L$, with tuples $\mathcal T$, base data with parts $\Pi$ (the \emph{types} of the base), and
for each lift $j$ and part $p$ the \emph{isotypic part} of type $p$ of the $j$-th lift; subgroups
$P_0,P_s\subset G$ with characters $\psi_0,\psi_s$; definite subspaces; rows $S_0$ and columns $S_1$
(Definition~\ref{is-indefinite-datum}) on the definite sides of $(P_0,\psi_0)$ and $(P_s,\psi_s)$;
definite columns $D_c$ with periods and column data (Definition~\ref{is-division-package}(3)) for
$S_1$; for each word $u$ an $\mathcal O_L$-module $N^{C}_u$ in the column of $u$; and identifications
$I$ of the tuples, words, lattices and spaces of $D$ with those of the family: a map $t\mapsto\mu_t$
to tuples of Hecke characters, isomorphisms $e_{RR}$ and $e_{RR,0}$ from $L^{D}$ to the spaces of the
glued lattices of $D$ and of $p_0$, isomorphisms $e_{N,u}$ from the columns to the spaces of the
$N_u$, a bijection $u\mapsto w_u$ from words to the members of $D_c$, isomorphisms $e_{M}(i,t)$ to
the twisted spaces of $D$, and maps $j_W(i)$ to the lifted spaces of $D$. These are tied to the
ambient and lattice data of $D$ and of $p_0$.
\end{definition}
\leanref{Formalism/LocalToGlobal/HilbertRowConditions/ColumnsFunctionAtSubspaces.lean}{Formalism.L2.Rows.RowSideFamily}

\begin{definition}\label{is-division-rows}
Let $P$ be a row-side family for a definite-side datum $D$ (Definition~\ref{is-row-side-family}).
Call a tuple $t$ \emph{used} if $\mu_t$ is used by $D$ (Definition~\ref{pd-period-data}(6)). The
\emph{division rows} hold for $P$ if:
\begin{enumerate}
\item $e_{RR}$ identifies the glued lattice $RR$ of $S_1$ with its metric $X_E$ with the glued
  lattice of the sharp row side of $D$ with its metric, and $e_{RR,0}$ identifies the glued lattice of
  $S_0$ with that of the base point $p_0$, as lattices and at every infinite place as metrics;
\item for every lift $i$ and used tuple $t$, the space of forms of $D$ for $\mu_t$ at level $i+1$ is
  spanned over $\mathbf C$ by the images under $e_M(i,t)$ of the columns of the coordinate matrix of
  the $i$-th lift at $t$;
\item for every lift $i$, used tuple $t$ and $x$, the theta lift of $D$ at level $i+1$ of the form with
  coordinates $x$ is $j_W(i)$ applied to $\theta_{i,t}$ times the projected coordinates of $x$;
\item for every lift $i$ and used tuple $t$ and every finite place $w$, the index of $D$ at $(\mu_t,i+1,w)$
  is the index $I_w(i,t)$ of $S_0$, and the relative index of $D$ is $j_w(H_T(i,t),T_{i,t})$; and the
  local defect $\eta_{i+1}(w)$ of $D$ is that of the family at $w$ if $w$ is twisted for $S_0$, and $0$
  otherwise;
\item for every word $u$ and all $x,y$, the division pairing $B^{\mathrm{div}}_u$ of $D$
  (Definition~\ref{th-bessel-identity}) at $(e_{N,u}x,e_{RR}y)$ is the pairing $B_{w_u}(x)(y)$ of the
  column data.
\end{enumerate}
\end{definition}
\leanref{Formalism/LocalToGlobal/Divisibility/TupUsed.lean}{Formalism.L2.DIV.DivRows3G}
\leanref{Formalism/LocalToGlobal/Divisibility/RowSideTie.lean}{Formalism.L2.DIV.HLatTie}
\leanref{Formalism/LocalToGlobal/Divisibility/TupUsed.lean}{Formalism.L2.DIV.NumTieG}

\begin{definition}\label{is-realisation-characterization}
Let $R$ be a realization of a definite-side datum $D$ for a family of definite forms, with a face $F$
(Definition~\ref{is-realisation}), and use the notation of that definition. The realization satisfies
its \emph{characterization} at $F$ if it satisfies the following conditions.
\begin{enumerate}
\item (Level characters.) The characters $\mu^0_j$, $\eta$ and $\xi_t$ are those of the level system of
  $D$ (Definition~\ref{th-grid}); $t\mapsto\xi_t$ is injective, $\xi_{\mu_0}=1$, and each $\xi_t$ has
  finite order, is trivial on the ideles of $F_h$ and is tamely ramified; $P_g$ is the set of places of
  $E_h$ over the level primes $g_0,g_1$ that the columns of $D$ choose; each $Q_t$ is a set of twisting
  places prime to the levels of $D$ and of its words, whose places over $p_{de}$ or $p'_{de}$ lie over a
  single prime of the subfield $K$ of the host, and $\xi_t$ is ramified at the places of $Q_t$ and their
  conjugates and at no others; the characters of the residue fields are the restrictions of $(\mu^0_j)^{-1}$ at the
  places of $P_g$ and of $\xi_t$ at those of $Q_t$; and the adelic matrices of $Y_j$, $X'$, $V_i$ and
  $W_i$ are self-dual frames of their lattices.
\item (Group.) The reductions of the level actions on $X'$ and the $V_i$ are given by the isomorphism
  $G\cong\prod_{P\in P_g}GL_r(k_P)$; under it, for a generic family of additive characters of the $k_P$,
  $P_0$ is the upper unitriangular subgroup with the character built from them, and $P_1$, after a
  fixed reordering of the coordinates, is the sharp symmetric subgroup with its character; $G$ is
  identified with the group of $D$, the places matching $P_g$ and the idempotent of the family having
  the coefficients of the type idempotent of $D$; and the character of $FW_i$ is the product over
  $P\in P_g$ of the characters $\chi_{i+1}$ of the residue fields applied to the determinants.
\item (Family.) The family tie of Definition~\ref{is-ties-lattice-systems}(4) holds; $P_g\subset S_b$;
  $L$ contains a primitive root of unity of order $|k_P|$ for every $P$ in $P_g$ or in some $Q_t$;
  every place split off $S_b$ comes from a place of the family; the Hecke polynomials of the row and
  lift forms have constant term $1$ and degree at most the rank of $X'$ or $V_i$; and their values at
  the places of the family are given by the characters $\mu^0_{i+1}\xi_t$ and $\eta$.
\item (Models.) $\theta$ is an archimedean family and the $V_i$ are positive definite; the coefficient
  spaces of $X'$ and the $V_i$ are coefficient spaces, with a dual pairing for $X'$; the Hecke modules
  $H_X(t)$, $H_V(i,t)$, $H_E$ and $H_W(i,\sigma)$ are the spaces of adelic forms of their levels and
  characters, with their Petersson products; $G$ acts on $H_E$ and $H_W(i,\sigma)$ through the level
  action; the semilinear maps from the row and lift forms are those of the realization by algebraic
  modular forms; the measures and fundamental domains are those of the levels; the matrices $B_V(i)$
  are orthonormal bases; and the model of $W_i$ at $\sigma_0$ puts its positive signs first.
\item (Chain of theta maps.) The Weil data of $(Y_k,Y_{k+1})$ and $(Y_{m-1},X')$ have the characters
  of the chain; their test vectors are the step vectors of the level data, and their pure tensors are
  theta functions; the iterated theta maps are nonzero, bounded and equivariant for the Weil
  representations; the maps $\mathrm{io}$ of the chain are those of these data; and each theta lift
  $\mathrm{Th}_t$ is the composite of the chain.
\item (Theta lifts.) The Weil data of $(V_i,W_i)$ have the characters of the theta lift; their test
  vectors are the level vectors and their pure tensors are theta functions; $\mathrm{io}_P(i)$
  intertwines the actions and takes the value $\Delta_i$ at some pair of unit vectors; and
  $\mathrm{th}_F(i,t)$ takes values in the forms on $U(W_i)$, is the theta integral of the lift forms,
  and gives $\mathrm{th}_G(i,t)$ in $H_W(i,\sigma_0)$.
\item (Scalars.) The scalars $\lambda_{i,t,p,v}$ and $\kappa$ satisfy the conventions for the scalars at
  $F$ of Lemma~\ref{is-onto-step-links-scalars}.
\item (Tuples.) A finite place $w$ is twisted if, for $w'=w$ or $w'=\bar w$ and some lift $j$, one of
  $I_{w'}(j,\mu_0)$, $I_{\bar w'}(j,\mu_0)$, the relative index of $T_{j,\mu_0}$ and its dual for
  $H_T(j,\mu_0)$ at $w'$, $F_0(w',j)$ and $F_0(\bar w',j)$ is nonzero, for the rows of $R$ at $F$ and the
  scale data of $F$; $\mu_w=\mu_0$ if $w$ is not twisted and $\mu_w\ne\mu_0$ if it is; every tuple other
  than $\mu_0$ is some $\mu_w$; $\mu_{\bar w}=\mu_w$; and for a twisted place $w$,
  $\xi_{\mu_w}=\delta_w\psi_w$ for the character $\delta_w$ of the columns at the place of $E_h$ below $w$
  and a twisting character $\psi_w$ of that place with ramified set $Q_w$.
\end{enumerate}
The \emph{core conditions} are (1) to (6). The \emph{weak characterization} at $F$ consists of (1) to
(4); conditions (5) and (6) for those iterated theta maps, maps $\mathrm{io}$, theta lifts and theta maps
that exist; the \emph{conditions on the inert indices}, that the maps of Definition~\ref{is-realisation}(10)
from the inert indices to places of $E_h$ are tied to the family, reach every place off $S_b$ at which it
is inert, carry the inert Hecke operators to the local Hecke operators of their labels, and that $S_F$ is
the set of places of $F_h$ below $S_b$; the \emph{fullness} of the Weil data of the chain and of the
theta lifts; the conventions of (7) for the scalars $\lambda^{(3)}_{i,t,p,v}$ and $\kappa$ at the level
$0$ of $F$; the conditions of (8) other than the first; that the phases $\phi_i$ have absolute value $1$
and that the $i$-th lift at $\mu_0$, if it is a constant of absolute value $1$ times a map into $FW_i$,
is $\phi_i$ times one (conjugated if $\sigma_0$ is not the embedding of $v_0$); and $S_b=P_g$. When all
the maps of (5) and (6) exist, the weak characterization gives (1) to (6).
\end{definition}
\leanref{Formalism/LocalToGlobal/IndefiniteReal/LinksChar.lean}{Formalism.IndefR.Links.RealDat.Char}
\leanref{Formalism/LocalToGlobal/IndefiniteReal/LinksDefs.lean}{Formalism.IndefR.Links.LevCh.Is}
\leanref{Formalism/LocalToGlobal/IndefiniteReal/LinksChar.lean}{Formalism.IndefR.Links.RealDat.IsGrp}
\leanref{Formalism/LocalToGlobal/IndefiniteReal/LinksChar.lean}{Formalism.IndefR.Links.RealDat.IsFam}
\leanref{Formalism/LocalToGlobal/IndefiniteReal/LinksChar.lean}{Formalism.IndefR.Links.RealDat.IsModels}
\leanref{Formalism/LocalToGlobal/IndefiniteReal/LinksChar.lean}{Formalism.IndefR.Links.RealDat.IsChain}
\leanref{Formalism/LocalToGlobal/IndefiniteReal/LinksChar.lean}{Formalism.IndefR.Links.RealDat.IsThP}
\leanref{Formalism/LocalToGlobal/IndefiniteReal/LinksChar.lean}{Formalism.IndefR.Links.RealDat.IsTup}
\leanref{Formalism/LocalToGlobal/IndefiniteReal/ZeroPartEigenvalueConvention.lean}{Formalism.IndefR.Links.RealDat3.WeakCharacterisation}
\leanref{Formalism/LocalToGlobal/IndefiniteReal/ZeroPartEigenvalueConvention.lean}{Formalism.IndefR.Links.RealDat3.IsInert}
\leanref{Formalism/LocalToGlobal/IndefiniteReal/LinksW.lean}{Formalism.IndefR.Links.RealDat.BaseW}
\leanref{Formalism/LocalToGlobal/IndefiniteReal/ZeroPartEigenvalueConvention.lean}{Formalism.IndefR.Links.RealDat3.characterisation_of_weak}
% Lean: the characterization is RealDat.Char (lev, grp, fam, models, chain, thP, scal, tup); the core
% conditions are GlobalTheory.Hk3.Core (lev, grp, fam, models, chain, thP). (1) is LevCh.Is (mu0, eta,
% xi, xi_inj, xi_base, xi_fin, xi_triv, xi_tame, Pg, Q_pl, Q_off, Q_de, Q_ram, chi, xiR, betaY,
% betaX, betaV, betaW); (2) is IsGrp (red_R, red_L, levels with IsGenericPsi, P0, PsharpSym and
% IsSymOrd, record, chiW); (3) is IsFam (tie, Sb, roots, pl_onto, degR, degL, muVal, etaVal); (4) is
% IsModels (21 fields, models_inf to models_order); (5) is IsChain (omegaY, omegaX, vY, vX, ptY, ptX,
% ItY, ItX, ioY, ioX, Th); (6) is IsThP (omegaP, vP, pt_theta, ioP, thF_mem, thF, thG); (7) is
% RealDat.IsScal (Lemma is-onto-step-links-scalars); (8) is RealDat.IsTup (tw, mu_off, mu_used, mu_ne,
% mu_conj, mu_on). The weak characterization is RealDat3.WeakCharacterisation: base = BaseW (lev, grp,
% fam, models, IsChainW, IsThPW, whose clauses hold when the map exists), inert = IsInert (tie, onto,
% hR, hL, SF), fullY, fullX, fullP (IsWeilFull), scal = RealDat3.IsScal (with lam3), tup =
% RealDat3.IsTup (without tw), phase and allPl (Sb = Pg). characterisation_of_weak, with ExAt,
% gives RealDat3.Char3, which contains (1) to (6). The items summarize the Lean fields; their full
% content is in the Lean.

\begin{definition}\label{is-realisation-join}
Let $P$ be a row-side family for a definite-side datum $D$ (Definition~\ref{is-row-side-family}), and
let an \emph{indefinite realization} of $P$ be given: a realization $R$ of $D$ for the family of $P$
(Definition~\ref{is-realisation}) that satisfies the weak characterization at the face of $P$
(Definition~\ref{is-realisation-characterization}), has the
standard character $\psi$ with sign $-1$, and whose complex data at that face have the definite
subspaces, rows, columns, definite columns, periods and column data of $P$, together with an automorphic setting
with Poincaré series for the words and an isometry $s_E$ from $E$ to the space $\mathcal H$ of that
setting. Let $\sigma_D$ be the embedding of $L$ into $\mathbf C$ of $D$. The base place $v_0$ is the
place of $\sigma_D$ (Definition~\ref{is-realisation}), so the embedding attached to $v_0$ in the complex
data (Definition~\ref{is-complex-data}) is $c\circ\sigma_D$, where $c$ is the identity of $\mathbf C$ if
that embedding is $\sigma_D$ and complex conjugation otherwise; $\sigma_D$ and $c$ are applied to vectors
entry by entry. The realization is \emph{joined} if there are: a bijection $u\mapsto m(u)$ from words to the
members of $D_c$; isomorphisms $c_\varphi(u)$ from the space $M_u$ of $D$ to the column of $m(u)$ and a
linear map $c_F$ from the coordinates of the base rows to the ambient space of $D$ at $\mu_0$; and a
$c$-semilinear map $e_R$ from the top forms of $D$ to $E$, such that:
\begin{enumerate}
\item (Chain.) For every word $u$ and all $x,y$ with entries in $L$, the column period of $D$ of
  $c_\varphi(u)^{-1}\sigma_D(x)$ against $c_F(\sigma_D(y))$ is $\sigma_D$ of the period pairing
  $\mathrm{Pt}_{m(u)}(\mu_0)(x,y)$.
\item (Base vectors.) The base space $V_0(\mu_0)$ of $D$ lies in the image of $c_F$.
\item (Rows.) $e_R$ maps onto the level $E^{1}$, multiplies norms by a constant $c_0>0$, and
  $\Theta_{\mu_0}j_{\mu_0}$ of the form with coordinates $y\in L^{n_V(\mu_0)}$ is a nonzero multiple of
  $e_R(\iota_{\mu_0}c_F(\sigma_D(y)))$, the multiple independent of $y$.
\item (Columns.) For every word $u$ the norm of $c(c_\varphi(u)\varphi)$ in the metric of the column
  of $m(u)$ at $v_0$ is a constant multiple of the norm of $\varphi$.
\item (Pairings.) For every word $u$ there is $\kappa_u\ne0$ such that, for $F$ a top form, $\varphi\in
  M_u$ and $T$ a Poincaré datum of $u$ whose column dictionary is the division vector of $\varphi$, the
  inner product of $s_E(e_RF)$ with the Poincaré vector of $c(c_\varphi(u)\varphi)$ is $c$ of
  $\kappa_u$ times the inner product of the conjugated automorphic form of $F$ with the Poincaré series
  of $T$.
\end{enumerate}
\end{definition}
\leanref{Development/AutomorphicArithmetic/ThetaCorrespondence/ThetaPairLifts/ThirdPlaceSeries/IndefRRealT.lean}{Formalism.IndefR.IndefiniteRealisation}
\leanref{Development/AutomorphicArithmetic/ThetaCorrespondence/ThetaPairLifts/ThirdPlaceSeries/ColJoinSOfSTiesR.lean}{Formalism.L2.SG.JoinSI.IndefiniteRealisationJoin}
\leanref{Formalism/LocalToGlobal/SiegelLatticeTheory/FlagStabilizerTensors/JoinSI.lean}{Formalism.L2.SG.JoinSI.RowTieT}
\leanref{Formalism/LocalToGlobal/SiegelLatticeTheory/FlagStabilizerTensors/JoinSI.lean}{Formalism.L2.SG.JoinSI.PairTieT}

\subsection{Record objects and their estimates}\label{subsection-is-record-objects}

\begin{definition}\label{is-record-object}
Let $S$ be an indefinite-side datum over $K$ with lifts $J$, rank $D$ and twisted places
$\mathrm{Tw}$ (Definition~\ref{is-indefinite-datum}), let $P$ be a finite group, and let
$c=(C,C_{b,1},C_{b,2})$ be real numbers with $14\le C_{b,1}$ and $66C+4\le C_{b,2}$. A \emph{record
object} for $S$, $P$ and $c$ consists of: for each lift $j$ an integer $n_Z(j)$, a representation of
$P$ on $K^{n_Z(j)}$, a subspace $V_0(j)$, an endomorphism $e_j$, an $\mathcal O_K$-lattice $M_s(j)$
and a linear map $\varphi_j\colon K^{d_W(j)}\to K^{n_Z(j)}$; a subgroup $U\subset P$ and a function
$P\to\mathcal O_K$; for each finite place $w$ and lift $j$, $\mathcal O_K$-submodules
$\Theta(w,j)$, $\Theta_0(w,j)$, $\Theta_1(w,j)$, $\Theta_K(w,j)$, $\Theta_J(w,j)$ and $X_3(w,j)$ of
$K^{n_Z(j)}$, elements $\sigma(w,j),a(w,j)\in K$, integers $E(w,j)$ and $n_X(w,j)$, real numbers
$r_1(w,j)$, $r_2(w,j)$ and $\varepsilon_0(w,j)$, maps $\theta'_{w,j}\colon K^{n_X(w,j)}\to K^{n_Z(j)}$
with endomorphisms $e_X$, $\mathrm{pr}$ and lattices $M_t(w,j)$ of $K^{n_X(w,j)}$; numbers $d_0(j)$,
$J_p(j,\cdot)$ and weights $\mathrm{wt}(j)$; and real numbers $k,m,f,b,b',S_c,\eta,T,x,\varepsilon,y$
and further numerical parameters, the sizes of the member. For a point $p$ (Definition~\ref{is-point})
and a pair of \emph{runs} $r$ at $p$, which choose at every twisted place laws of the local
correspondence (Definition~\ref{gal-llc-laws}) and levels for the lifts, pinned to the realization
and to the own place of the twisted place, the record object is \emph{recorded} at $(p,r)$ if its
spaces, lattices and numbers are those that $p$ and $r$ define.
\end{definition}
\leanref{Formalism/LocalToGlobal/IndefiniteForms/IndefiniteTree/ExampleSatisfiesConditionsAtFive.lean}{Formalism.L2.Indef.Tree.OnObj}
\leanref{Formalism/LocalToGlobal/IndefiniteForms/IndefiniteTree/ExampleSatisfiesConditionsAtFive.lean}{Formalism.L2.Indef.Tree.OnConst}
\leanref{Development/AbelianVarietyTheory/Statements/PlaceShiftStatements/SingleIndexPlacesAreTwisted.lean}{GlobalTheory.BigOn.RecAt}
\leanref{Development/AbelianVarietyTheory/Definitions/Divisors/RealisedChainTorusEq.lean}{GlobalTheory.BigOn.RunsOk}

\begin{lemma}\label{is-canonical-error}
Let $\mathrm{Tw}$ be a finite set and $J$ a nonempty finite set. For $w\in\mathrm{Tw}$ and $j\in J$
let $I(w,j)$, $F(w,j)$, $\mathrm{err}(w,j)$ and $J_0(w,j)\geq0$ be real numbers, and let $e_w>0$,
$c_w$ and $\ell_w\geq0$ be real numbers. Let $D>0$ and $E,E'$ be real with $2E\le E'$. Put
\[
 \mathrm{err}_0(w,j)=\frac{\max\big(0,\ \frac1D\sum_{j'}I(w,j')-\sum_{j'}2e_wF(w,j')-c_w\big)}{e_w\,|J|}.
\]
Suppose that for all $w\in\mathrm{Tw}$ and $j\in J$:
\begin{enumerate}
\item $\sum_{j'}I(w,j')\le\big(\sum_{j'}(2e_wF(w,j')+e_w\,\mathrm{err}(w,j'))+c_w\big)D$;
\item $\mathrm{err}(w,j)\geq0$;
\item $\sum_{w\in\mathrm{Tw}}\ell_w\big(e_w\,\mathrm{err}(w,j)+e_wJ_0(w,j)\big)\le E$.
\end{enumerate}
Then (1) and (2) hold with $\mathrm{err}_0$ in place of $\mathrm{err}$, and
$\sum_{w\in\mathrm{Tw}}\ell_w\big(e_w\,\mathrm{err}_0(w,j)+e_wJ_0(w,j)\big)\le E'$ for every $j$.
\end{lemma}
\leanref{Development/AbelianVarietyTheory/Proofs/DivisorProofs/OnOfOnWithCanonical.lean}{GlobalTheory.BigOn.canonical_error_real}
\begin{proof}
Write $Q_w=\frac1D\sum_{j'}I(w,j')-\sum_{j'}2e_wF(w,j')-c_w$, so that
$\sum_je_w\,\mathrm{err}_0(w,j)=\max(0,Q_w)$. Then (1) for $\mathrm{err}_0$ reads
$\frac1D\sum_{j'}I(w,j')\le\sum_{j'}2e_wF(w,j')+\max(0,Q_w)+c_w$, which holds because $Q_w\le\max(0,Q_w)$.
Condition (2) for $\mathrm{err}_0$ holds because the numerator is nonnegative and $e_w|J|>0$.

For the sum, fix $j$ and $w\in\mathrm{Tw}$. By (1), $Q_w\le\sum_{j'}e_w\,\mathrm{err}(w,j')$, and by (2)
this sum is nonnegative, so $|J|\,\ell_we_w\,\mathrm{err}_0(w,j)=\ell_w\max(0,Q_w)\le\sum_{j'}\ell_we_w\,
\mathrm{err}(w,j')$. Summing over $w$ and using (3) together with $J_0\geq0$ gives
$|J|\sum_w\ell_we_w\,\mathrm{err}_0(w,j)\le\sum_{j'}E=|J|E$. By (3) and $\mathrm{err}\geq0$,
$\sum_w\ell_we_wJ_0(w,j)\le E$. Adding the two bounds gives $E+E\le E'$.
\end{proof}

\begin{lemma}\label{is-hook-lower-bound}
Let $D$ be a definite-side datum with size data at scale $x$ (Definition~\ref{pd-size-data}), with
the quantity $b$, and let $j$ be a lift at which the lifts have the sizes required by the size data.
Let $\log\mathrm{hook}_j(\rho)$ be the logarithmic hook sum of the weights of the $j$-th lift at the
type $\rho$, $L_p\geq0$ the logarithmic size of the size data, $f=2n$ and $n_V$ the rank of the
columns, and put $\mathrm{low}(j)=\sum_\rho\log\mathrm{hook}_j(\rho)-8fn_V^2L_p$. If $C_{sz}\geq8$,
then
\[
 \sum_\rho\log\mathrm{hook}_j(\rho)-C_{sz}fn_V^2L_p\le\mathrm{low}(j).
\]
\end{lemma}
\leanref{Development/AbelianVarietyTheory/Proofs/DivisorProofs/HookLogSumSubSizeTermLeLowerBoundStatement.lean}{GlobalTheory.T.hookLogSum_sub_sizeTerm_le_lowerBound}
\leanref{Development/AbelianVarietyTheory/Definitions/Divisors/RealisedChainTorusEq.lean}{GlobalTheory.BigOn.lowB}
\begin{proof}
The sizes of the lifts give $L_p\geq0$. The hook sum of the scale data at $j$ is the sum over the
types of the logarithmic hook sums of the positive parts of the weights of the $j$-th lift, which is
the first term of $\mathrm{low}(j)$. Since $fn_V^2L_p\geq0$ and $C_{sz}\geq8$,
$C_{sz}fn_V^2L_p\geq8fn_V^2L_p$.
\end{proof}

\begin{definition}\label{is-onto-record-numbers}
Let $O$ be a record object for an indefinite-side datum $S$ over $K$, a finite group $P$ and
constants $c=(C,C_{b,1},C_{b,2})$ (Definition~\ref{is-record-object}), with the notation there;
write $\rho_j$ for its representation of $P$ on $K^{n_Z(j)}$ and $\psi\colon P\to\mathcal O_K$ for
its function. The \emph{further parameters} of $O$ are real numbers $D_F$, $X_2$, $v_d$, $r_0$,
$\delta_2$ and $N_c$, real functions $\vartheta(l)$ and $\alpha(l)$ of integers $l\geq0$, integers
$T_2$, $p_{de}$, $g_0$, $g_1$, $q$ and $d$, a finite set $S_m$ of integers, and elements
$\nu(j)\in K$ for the lifts $j$. For a lift $j$, a finite place $v$ of $K$ and an integer $l\geq1$,
the following quantities are attached to $O$.
\begin{enumerate}
\item The eigenspace $V_\psi(j)=\{z\in K^{n_Z(j)}:\rho_j(u)z=\psi(u)z\text{ for all }u\in U\}$, its
  dimension $d_0(j)=\dim_KV_\psi(j)$, and $Z_g(j)=V_\psi(j)\cap V_0(j)$.
\item The numbers $b=\exp(x^\varepsilon)$, so that $\log b=x^\varepsilon$, and $y=1/(4(fk)^3)$,
  $S_e=S_c+\eta$, $Y=CD_F$,
  \[
   E_1^{u}(O)=x^\varepsilon\big(fmb\log b+N_c\big),\qquad E_J=30\,T_2\,mb\,x^{2\varepsilon}+2mb.
  \]
\item The numbers $s_B(j)=\mathrm{wt}(j)+fS_e$, $s_C(j)=Cm(2+v_d)s_B(j)$,
  $A_1(j)=C\big(m\,\mathrm{wt}(j)\log k+fkS_e\big)$, $A_w(j)=Ck\,\mathrm{wt}(j)$ and
  \[
   a_E(j,l)=\frac{A_1(j)}{l}+A_w(j)\vartheta(l)+T\alpha(l)+\begin{cases}A_w(j)/l^{2}&\text{if }l>X_2,\\
   0&\text{otherwise.}\end{cases}
  \]
\item The local numbers $e_0(v,j)=\big(E(v,j)-v_v(\sigma(v,j))\big)/e_v$ and
  $\nu_v(j)=v_v(\nu(j))/e_v$, where $v_v$ is the valuation at $v$ and $e_v$ the ramification index of
  Definition~\ref{is-indefinite-datum}; and $r(v,j)=r_2(v,j)$ if $p_v=p_{de}$, $r(v,j)=r_1(v,j)$
  otherwise.
\item The \emph{bad primes} $\{g_0,g_1,q,d\}\cup S_m$.
\end{enumerate}
For $\mathcal O_K$-submodules $A,B$ of a $K$-vector space and a finite place $w$, write
$\operatorname{idx}_w(A,B)=\operatorname{len}_w\big(A/(A\cap B)\big)-\operatorname{len}_w\big(B/(A\cap
B)\big)$, with $\operatorname{len}_w$ as in Definition~\ref{is-indefinite-datum}, and $A_{(w)}$ for
the localization of $A$ at $w$. An $\mathcal O_K$-submodule is a \emph{lattice in} a subspace $W$ if
it is finitely generated and spans $W$ over $K$.
\end{definition}
\leanref{Formalism/LocalToGlobal/IndefiniteForms/IndefiniteTree/ExampleSatisfiesConditionsAtFive.lean}{Formalism.L2.Indef.Tree.OnObj.Zg}
\leanref{Formalism/LocalToGlobal/IndefiniteForms/IndefiniteTree/ExampleSatisfiesConditionsAtFive.lean}{Formalism.L2.Indef.Tree.OnObj.d0}
\leanref{Formalism/LocalToGlobal/IndefiniteForms/IndefiniteTree/ExampleSatisfiesConditionsAtFive.lean}{Formalism.L2.Indef.Tree.OnObj.E1u}
\leanref{Formalism/LocalToGlobal/IndefiniteForms/IndefiniteTree/ExampleSatisfiesConditionsAtFive.lean}{Formalism.L2.Indef.Tree.OnObj.EJ}
\leanref{Formalism/LocalToGlobal/IndefiniteForms/IndefiniteTree/ExampleSatisfiesConditionsAtFive.lean}{Formalism.L2.Indef.Tree.OnObj.aE}
\leanref{Formalism/LocalToGlobal/IndefiniteForms/IndefiniteTree/ExampleSatisfiesConditionsAtFive.lean}{Formalism.L2.Indef.Tree.OnObj.e0}
\leanref{Formalism/LocalToGlobal/IndefiniteForms/IndefiniteTree/ExampleSatisfiesConditionsAtFive.lean}{Formalism.L2.Indef.Tree.OnObj.Bad}
\leanref{Formalism/MirabolicMackeyTheory/StatementL.lean}{Formalism.EigenvalueLevelCeiling.eig}

\begin{definition}\label{is-onto-page-estimates}
Let $O$ be a record object with constants $c=(C,C_{b,1},C_{b,2})$, with the notation of
Definitions~\ref{is-record-object} and~\ref{is-onto-record-numbers}. Let $(k_i)_{i\in I}$ be a finite
family of finite fields and $f_N$ an integer. The record object satisfies the \emph{page estimates}
for $(k_i)$ and $f_N$ if the following hold.
\begin{enumerate}
\item (Page.) $C\geq1$, $k\geq1$, $m\geq1$, $f\geq8$, $x\geq1$, and $T$, $S_c$, $\eta$, $b'$, $v_d$,
  $N_c$ are nonnegative; $g_0,g_1,q,d$ are primes with $g_0\ne g_1$, $q\ne d$ and
  $q,d\notin\{g_0,g_1\}$, and none of them lies in $S_m$.
\item (Sizes.) $0\le\mathrm{wt}(j)\le2fb$ for every lift $j$; $1<Y\le mb$ and $\log Y\le2\log b$;
  $kS_e\log b\le mb$; $k\le X_2$; $\log k\le\log b$; and $1\le\log b$.
\item (Scale.) $\log x\le x^\varepsilon$ and $T_2\le f$.
\item (Bad primes.) $C_{b,1}\geq0$, $C_{b,2}\geq0$, and for every lift $j$
  \[
  \begin{aligned}
   \big(2(s_B(j)+T+1)+4fkb'\big)\log(qd)&\le C_{b,1}\,fmb,\\
   \big(2s_C(j)+2T+2\big)\textstyle\sum_{l\in S_m}\log l&\le C_{b,2}\,fmb\log b,\\
   \big(4T+4fr_0^{2}\big)\log(g_0g_1)&\le4b.
  \end{aligned}
  \]
\item (Levels.) Every $k_i$ has $g_0$ or $g_1$ elements, and the number of $k_i$ with $g_0$
  elements and the number with $g_1$ elements are both $f_N$.
\item (The $y$-bound.) $2f/(\min(g_0,g_1)-1)\le y$.
\end{enumerate}
\end{definition}
\leanref{Formalism/LocalToGlobal/IndefiniteForms/IndefiniteTree/ExampleSatisfiesConditionsAtFive.lean}{Formalism.L2.Indef.Tree.LeafPage}
\leanref{Formalism/LocalToGlobal/IndefiniteForms/IndefiniteTree/ExampleSatisfiesConditionsAtFive.lean}{Formalism.L2.Indef.Tree.LeafSizes}
\leanref{Formalism/LocalToGlobal/IndefiniteForms/IndefiniteTree/ExampleSatisfiesConditionsAtFive.lean}{Formalism.L2.Indef.Tree.LeafX}
\leanref{Formalism/LocalToGlobal/IndefiniteForms/IndefiniteTree/ExampleSatisfiesConditionsAtFive.lean}{Formalism.L2.Indef.Tree.LeafBad}
\leanref{Formalism/LocalToGlobal/IndefiniteForms/IndefiniteTree/ONLeaves.lean}{Formalism.L2.Indef.Tree.LeafLevelSq}
\leanref{Formalism/LocalToGlobal/IndefiniteForms/IndefiniteTree/ONLeaves.lean}{Formalism.L2.Indef.Tree.LeafYBound}

\begin{definition}\label{is-onto-geometric-estimates}
Let $O$ be a record object for an indefinite-side datum $S$ over $K$, with the notation of
Definitions~\ref{is-record-object} and~\ref{is-onto-record-numbers}. In this definition $j$ runs
over the lifts and $v$ over the twisted places $w$ of $S$ and their conjugates $\bar w$. The record
object satisfies the \emph{geometric estimates} if the following hold.
\begin{enumerate}
\item (Supremum.) $M_s(j)$ is a lattice in $K^{n_Z(j)}$ stable under $\rho_j(P)$.
\item (Span.) $\Theta(v,j)$ and $\Theta_K(v,j)$ are lattices in $K^{n_Z(j)}$.
\item (Containment.) $\Theta(v,j)\subset M_s(j)$, $\Theta(v,j)$ is stable under $\rho_j(P)$,
  $\sigma(v,j)\ne0$ and $\Theta(v,j)\subset\sigma(v,j)\Theta_K(v,j)$.
\item (Data.) $\Theta_0(v,j)$ is finitely generated and contained in $V_\psi(j)$, and
  $|U|\cdot\big(\Theta_K(v,j)\cap V_\psi(j)\big)\subset\Theta_0(v,j)$.
\item (Index bound.) There are $\mathcal O_K$-submodules $\Theta_1'(v,j)$ of $K^{n_Z(j)}$, elements
  $a'(v,j)\in K$, real numbers $r_1'(u,j),r_2'(u,j)$ for all finite places $u$, and a real number
  $\delta_2'$, such that, with $r'(u,j)=r_2'(u,j)$ if $p_u=p_{de}$ and $r'(u,j)=r_1'(u,j)$
  otherwise:
  \begin{enumerate}
  \item $\Theta_1'(v,j)$ is a lattice in $K^{n_Z(j)}$ contained in $\Theta(v,j)$, $a'(v,j)\ne0$,
    $v_v(a'(v,j))=E(v,j)$, and for every subspace $W\subset K^{n_Z(j)}$
    \[
     \operatorname{idx}_v\big(M_s(j)\cap W,\ \Theta_1'(v,j)\cap W\big)\le E(v,j)\dim W+r'(v,j);
    \]
  \item $0\le\delta_2'\le1$, $r_1'(u,j)\geq0$ and $r_2'(u,j)\geq0$ for all finite places $u$, and for
    every lift $j$ and every finite set $S'$ of finite places,
    $\sum_{u\in S'}\ell_ur_i'(u,j)\le n_Z(j)\,fmb\,\delta_2'$ for $i=1,2$.
  \end{enumerate}
\end{enumerate}
\end{definition}
\leanref{Formalism/LocalToGlobal/IndefiniteForms/IndefiniteTree/ExampleSatisfiesConditionsAtFive.lean}{Formalism.L2.Indef.Tree.LeafSup}
\leanref{Formalism/LocalToGlobal/IndefiniteForms/IndefiniteTree/ExampleSatisfiesConditionsAtFive.lean}{Formalism.L2.Indef.Tree.LeafSpan}
\leanref{Formalism/LocalToGlobal/IndefiniteForms/IndefiniteTree/ExampleSatisfiesConditionsAtFive.lean}{Formalism.L2.Indef.Tree.LeafLemmaS}
\leanref{Formalism/LocalToGlobal/IndefiniteForms/IndefiniteTree/ExampleSatisfiesConditionsAtFive.lean}{Formalism.L2.Indef.Tree.LeafData}
\leanref{Formalism/LocalToGlobal/IndefiniteForms/IndefiniteTree/ExampleSatisfiesConditionsAtFive.lean}{Formalism.L2.Indef.Tree.LatticeIndexBoundAtPairs}
\leanref{Development/AbelianVarietyTheory/Definitions/Divisors/RealisedChainTorusEq.lean}{GlobalTheory.BigOn.HasIndexBound}

\begin{definition}\label{is-onto-depth-estimates}
Let $O$ be a record object for an indefinite-side datum $S$ over $K$, with the notation of
Definitions~\ref{is-record-object} and~\ref{is-onto-record-numbers}, and let $N$ be numerical bounds
for $S$ with function $F_0$ (Definition~\ref{is-numerical-bounds}). In this definition $j$ runs over
the lifts and $w$ over the twisted places of $S$, with residue characteristic $p_w$. The record
object satisfies the \emph{depth estimates} for $N$ if the following hold.
\begin{enumerate}
\item (Good places.) $\vartheta(l)\geq0$ and $\alpha(l)\geq0$ for every $l$; for every finite set
  $Q$ of primes, $\sum_{l\in Q}\vartheta(l)\log l\le C/k$ and $T\sum_{l\in Q}\alpha(l)\log l\le
  Cmb$; and if $p_w$ is not a bad prime, then
  \[
   e_0(w,j)\le F_0(w,j)+a_E(j,p_w)+1\quad\text{and}\quad e_0(\bar w,j)\le F_0(w,j)+a_E(j,p_w)+1
  \]
  when $p_w\le Y$, while $e_0(w,j)=e_0(\bar w,j)=0$ and $F_0(w,j)\geq0$ when $p_w>Y$.
\item (Places over $q$ and $d$.) If $p_w\in\{q,d\}$, then $e_0(w,j)$ and $e_0(\bar w,j)$ are at
  most $s_B(j)+T+1$; $F_0(w,j)=0$ if $p_w=q$, and $-2fkb'\le F_0(w,j)\le0$ if $p_w=d$.
\item (Small primes.) If $p_w\in S_m$, then $e_0(w,j)$ and $e_0(\bar w,j)$ are at most
  $s_C(j)+T+1$, and $F_0(w,j)=0$.
\item (Level primes.) If $p_w\in\{g_0,g_1\}$, then $e_0(w,j)\le2T-\nu_w(j)$,
  $e_0(\bar w,j)\le2T-\nu_{\bar w}(j)$ and $\nu_w(j)+\nu_{\bar w}(j)=-2F_0(w,j)$.
\end{enumerate}
\end{definition}
\leanref{Formalism/LocalToGlobal/IndefiniteForms/IndefiniteTree/ExampleSatisfiesConditionsAtFive.lean}{Formalism.L2.Indef.Tree.LeafDepthGood}
\leanref{Formalism/LocalToGlobal/IndefiniteForms/IndefiniteTree/ExampleSatisfiesConditionsAtFive.lean}{Formalism.L2.Indef.Tree.LeafDepthRam}
\leanref{Formalism/LocalToGlobal/IndefiniteForms/IndefiniteTree/ExampleSatisfiesConditionsAtFive.lean}{Formalism.L2.Indef.Tree.LeafDepth2}
\leanref{Formalism/LocalToGlobal/IndefiniteForms/IndefiniteTree/ExampleSatisfiesConditionsAtFive.lean}{Formalism.L2.Indef.Tree.LeafDepthLevel}

\begin{definition}\label{is-onto-local-estimates}
Let $O$ be a record object for an indefinite-side datum $S$ over $K$ with rank $D$, with the notation
of Definitions~\ref{is-record-object} and~\ref{is-onto-record-numbers}, and let $N$ be numerical
bounds for $S$ with function $J_0$ (Definition~\ref{is-numerical-bounds}). In this definition $j$
runs over the lifts, $w$ over the twisted places of $S$, and $v$ over the twisted places and their
conjugates.
\begin{enumerate}
\item (Rank.) $O$ has the \emph{rank at the twisted places} if $e_j(\Theta_0(v,j))$ and
  $e_j(\Theta_J(v,j))$ are lattices in $Z_g(j)$ and $\Theta_0(v,j)_{(v)}\subset\Theta_J(v,j)_{(v)}$.
\item (Twist.) $O$ has a \emph{twist bound} if there are real numbers $\varepsilon'(v,j)\geq0$ with
  \[
   \operatorname{idx}_v\big(M_s(j)\cap Z_g(j),\ X_3(v,j)\big)\le\operatorname{idx}_v\big(M_s(j)\cap
   Z_g(j),\ e_j\Theta_J(v,j)\big)+e_v\,\varepsilon'(v,j)
  \]
  for all $v$ and $j$, and $\sum_j\big(\varepsilon'(w,j)+\varepsilon'(\bar w,j)\big)\le2/p_w^{2}$ for
  every $w$.
\item (Share.) $O$ has the \emph{share bound} if for every $j$ there is $\epsilon\le y$ with
  $(1-\epsilon)\,n_Z(j)\le|U|\,d_0(j)$.
\item (Bad places.) $O$ has the \emph{bad-place bound} if for every $j$ there is
  $\delta\in[0,\frac14]$ with $(1-\delta)\,d_0(j)\le D$.
\item (Exponents.) $O$ has the \emph{exponent bound} for $N$ if $J_0(\mu_w,w,j)\geq0$ for all $w$
  and $j$, and there are real numbers $J'(j,l)\geq0$, for lifts $j$ and integers $l\geq0$, such that
  $\sum_{l\in Q}J'(j,l)\log l\le E_J$ for every $j$ and every finite set $Q$ of primes, and
  $J_0(\mu_w,w,j)\le J'(j,p_w)$ for all $w$ and $j$.
\end{enumerate}
\end{definition}
\leanref{Development/AbelianVarietyTheory/Proofs/FoundationProofs/TraceLatticeLeThetaLatticeAtPlace.lean}{GlobalTheory.TraceClassRank.RankAtPlaces}
\leanref{Development/AbelianVarietyTheory/Definitions/Divisors/RealisedChainTorusEq.lean}{GlobalTheory.BigOn.HasTwistBound}
\leanref{Formalism/LocalToGlobal/IndefiniteForms/IndefiniteTree/ExampleSatisfiesConditionsAtFive.lean}{Formalism.L2.Indef.Tree.LeafTW1}
\leanref{Formalism/LocalToGlobal/IndefiniteForms/IndefiniteTree/ExampleSatisfiesConditionsAtFive.lean}{Formalism.L2.Indef.Tree.LeafShareI}
\leanref{Formalism/LocalToGlobal/IndefiniteForms/IndefiniteTree/ExampleSatisfiesConditionsAtFive.lean}{Formalism.L2.Indef.Tree.LeafBadShare}
\leanref{Development/AbelianVarietyTheory/Definitions/Divisors/RealisedChainTorusEq.lean}{GlobalTheory.BigOn.HasExpBound}
\leanref{Formalism/LocalToGlobal/IndefiniteForms/IndefiniteTree/ExampleSatisfiesConditionsAtFive.lean}{Formalism.L2.Indef.Tree.LeafJSum}
\leanref{Formalism/LocalToGlobal/IndefiniteForms/IndefiniteTree/ExampleSatisfiesConditionsAtFive.lean}{Formalism.L2.Indef.Tree.LeafJExp}

\begin{lemma}\label{is-onto-recorded-numbers}
Let $p$ be a point at a member of a family (Definition~\ref{is-point}), whose definite-side datum
has ranks $k,m,n$ and $f=2n$, scale $(x,\varepsilon)$, size data and host. Let a pair of runs at $p$
be given, and let $O$ be a record object recorded at $p$ and these runs
(Definition~\ref{is-record-object}), with the notation of Definition~\ref{is-onto-record-numbers}.
\begin{enumerate}
\item The numbers $k,m,f,x,\varepsilon$ of $O$ are those of the member, and $r_0=k+m-1$; $b$,
  $b'$, $X_2=X''(n)$ and $S_c$ are the quantities $b,b',X''(n)$ and the cloud size of the size data
  (Definition~\ref{pd-size-data}), and $\eta=\eta_{sz}$ is its size $\eta_{sz}$
  (Definition~\ref{is-kernel-ties}); $\mathrm{wt}(j)=2fb$ for every lift $j$;
  $\vartheta(l)$ is the supremum of the numbers $l^{-i}$ over the odd $i\geq3$ with $l^{i}>k^{3/2}$,
  and $\alpha(l)=0$; $v_d=\log_2Q_S$; and $T$ is the integer $T_O$ of the parameter choices of the
  member and $D_F$ is their top degree.
\item The integers $g_0,g_1,q,d,p_{de}$ are the primes of the same names of the host
  (Definitions~\ref{hf-host} and~\ref{pd-family}), $S_m$ is the set of primes dividing $Q_S$, $T_2$ is
  the number of elements of the set $T_2$ of the host, and $N_c=|T_2|\,k\,b\,x^{2\varepsilon}/n^{2}$.
\item Let $FW_j$ be the rational structure of the $j$-th lift of the realization of $p$, with its
  coordinates $L^{n_W(j)}$ and its character $\chi_j$ of $G$ (Definition~\ref{is-realisation}(8)), and
  let $P_0,\psi_0$ be the subgroup and character of the face of $p$. Then $n_Z(j)=n_W(j)$,
  $|U|=|P_0|$, and $d_0(j)$ is the dimension of the subspace $FW_j^{0}$ of $L^{n_W(j)}$ on which $P_0$
  acts through $\chi_j\psi_0^{-1}$.
\end{enumerate}
\end{lemma}
\leanref{Development/AbelianVarietyTheory/Definitions/Jets/CmQuadraticFormEqAlternatingForm.lean}{GlobalTheory.RunDepth.RunTie}
\leanref{Development/AbelianVarietyTheory/Definitions/Divisors/RealisedChainTorusEq.lean}{GlobalTheory.BigOn.IsRec}
\leanref{Development/AbelianVarietyTheory/Proofs/DivisorProofs/LeafShareIOfShareW.lean}{GlobalTheory.ShareOn.Comparison₂.tied_dims}
\begin{proof}
(1) and (2) are among the conditions that make $O$ recorded: the depth tie of $O$ with the core of
$p$ gives the values of $k,m,f,x,\varepsilon,S_c,\eta,b',X_2,v_d,\mathrm{wt},\vartheta,\alpha,T$ and
of $g_0,g_1,q,d$ and $S_m$, and the remaining conditions give those of $r_0$, $T_2$, $N_c$, $D_F$
and $p_{de}$. Then $b=\exp(x^\varepsilon)$ is the quantity $b$ of the size data, and $N_c$ is the
number $N_c$ of the size data, $|T_2|kbx^{\varepsilon}/(n\cdot nx^{-\varepsilon})$.

(3) Fix a place $v$. The level tie of $O$ at $v$ gives a linear isomorphism $K^{n_Z(j)}\to
L^{n_W(j)}$, so $n_Z(j)=n_W(j)$. The tie identifies $U$ with the preimage of $P_0$ under the
inclusion of $P$ in $G$. By the links of the realization, $P_0$ consists of the elements of $G$ that
are upper unitriangular at every place of $P_g$; these lie in the mirabolic subgroups, so $P_0$ lies
in the image of $P$, and $|U|=|P_0|$. The tie also carries $\rho_j$ to the action of $P$ on $FW_j$
through this inclusion and $\psi$ to $\chi_j\psi_0^{-1}$ on $U$, so the isomorphism carries
$V_\psi(j)$ onto $FW_j^{0}$.
\end{proof}

\begin{lemma}\label{is-onto-page-numerics}
Let $\mathcal H$ be the class of standard families. Eventually along $\mathcal H$, every family has a
scale $x_0$ with the following property. Let $p$ be a point at a member at a scale $x\geq x_0$, let a
pair of runs at $p$ be given, and let $O$ be a record object recorded at $p$ and these runs, with
constants $c=(C,C_{b,1},C_{b,2})$ such that $C\geq1$, $14\le C_{b,1}$, $66C+4\le C_{b,2}$ and
$3C\le x$. Then $O$ satisfies (1) to (4) of Definition~\ref{is-onto-page-estimates}.
\end{lemma}
\leanref{Development/AbelianVarietyTheory/Proofs/FoundationProofs/ThreeLeScaleOfPowBound.lean}{GlobalTheory.IndexAssembly.numericLeaves}
\leanref{Development/AbelianVarietyTheory/Proofs/FoundationProofs/ThreeLeScaleOfPowBound.lean}{GlobalTheory.IndexAssembly.pageLeaves}
\leanref{Development/AbelianVarietyTheory/Proofs/FoundationProofs/ThreeLeScaleOfPowBound.lean}{GlobalTheory.IndexAssembly.SizeRows}
\leanref{Development/AbelianVarietyTheory/Proofs/DivisorProofs/IndexBoundParameterRequirementsEventuallyFamilyStatement.lean}{GlobalTheory.T.indexBoundParameterRequirements_eventually_family}
\begin{proof}
By Lemma~\ref{is-onto-recorded-numbers}, (1) to (4) are inequalities among the ranks, the size data,
the host and the parameter choices of the member and the constants $c$; the parameter $\delta_2$ of
$O$ does not enter them. Eventually along $\mathcal H$ every family has a scale beyond which every
member satisfies the \emph{size rows}: $\eta_{sz}b^{4/5}\le w''$, $11\le x^{\varepsilon_z/5}$ for the
exponent $\varepsilon_z$ of the scale, $x^{-\varepsilon}\le\delta_2\le1/(4L_w)$ for the number
$\delta_2$ of the size data and the integer $L_w$ of the host, $g_i\le b^{9}$ for $i=0,1$,
$\log(qd)\le m$, $kb'\le b$, $N(m)\le X'$, $m\geq44$, and two bounds on the top degree $D_F$ and
the constant $C_{nW}$ of the parameter choices. Together with $n\geq4$ and $k\geq m+3$, the
conditions on $q,d,g_0,g_1$ and $Q_S$ of the host, the bounds $kS_c\le w''$ and
$T\le b^{1/2}$ of the parameter choices, $C\geq1$, $14\le C_{b,1}$, $66C+4\le C_{b,2}$ and $3C\le
x$, these give (1) to (4) by a computation.
\end{proof}

\begin{lemma}\label{is-onto-twist-transport}
Let $K$ be a number field, $v$ a finite place of $K$ with residue characteristic $p_v$ and
ramification index $e_v=v_v(p_v)$, and $V$ a finite-dimensional $K$-vector space. Let $\mathcal O'$
be a discrete valuation ring with fraction field $F'$ and normalized valuation $v'$, and let
$\sigma\colon K\to F'$ be a field embedding such that, for $x\in K$, $v_v(x)\geq0$ if and only if
$\sigma(x)\in\mathcal O'$. Let $W$ be an $F'$-vector space and $R\colon V\to W$ a $\sigma$-semilinear
map that carries some basis of $V$ to an $F'$-linearly independent family. For an
$\mathcal O_K$-submodule $A$ of $V$ write $\mathcal O'R(A)$ for the $\mathcal O'$-submodule of $W$
spanned by $R(A)$. For $\mathcal O'$-submodules $A',B'$ of $W$ such that $A'/(A'\cap B')$ and
$B'/(A'\cap B')$ have finite length put
\[
 [A':B']=\frac{\operatorname{len}_{\mathcal O'}\big(A'/(A'\cap B')\big)
 -\operatorname{len}_{\mathcal O'}\big(B'/(A'\cap B')\big)}{v'(p_v)}.
\]
Let $Z\subset V$ be a subspace, $M$ and $\Theta$ lattices in $Z$ in the sense of
Definition~\ref{is-onto-record-numbers}, $X$ an $\mathcal O_K$-submodule of $V$ and $\varepsilon$ a
real number. Suppose that $[\mathcal O'R(\Theta):\mathcal O'R(X)]$ is defined and at most
$\varepsilon$. Then, with $\operatorname{idx}_v$ as in Definition~\ref{is-onto-record-numbers},
\[
 \operatorname{idx}_v(M,X)\le\operatorname{idx}_v(M,\Theta)+e_v\,\varepsilon.
\]
\end{lemma}
\leanref{Development/AbelianVarietyTheory/Proofs/FoundationProofs/LeafTW1OfAllConjugatesGood.lean}{GlobalTheory.TwistIndex.lix_le_of_read}
\leanref{Development/AbelianVarietyTheory/Proofs/FoundationProofs/LeafTW1OfAllConjugatesGood.lean}{GlobalTheory.TwistIndex.relLength_read}
\begin{proof}
Let $\mathcal O_{(v)}$ be the localization of $\mathcal O_K$ at $v$, and write $A_{(v)}$ for the
localization of an $\mathcal O_K$-module $A$, so that $\operatorname{len}_v$ is the length over
$\mathcal O_{(v)}$. The hypothesis on $\sigma$ says that $\sigma$ maps $\mathcal O_{(v)}$ into
$\mathcal O'$, and that an element of $\mathcal O_{(v)}$ is a unit if and only if its image is one. So
$\sigma$ is an injective local homomorphism of discrete valuation rings. Let $e\geq1$ be its
ramification index, so that $v'(\sigma(x))=e\,v_v(x)$ for $x\in\mathcal O_{(v)}$; in particular
$v'(p_v)=e\,e_v$.

Let $g\colon\mathcal O'\otimes_{\mathcal O_{(v)}}V\to W$ be the $\mathcal O'$-linear map with
$g(s\otimes x)=sR(x)$. It is injective. Indeed, every element $t$ has a multiple $ct$, with
$0\ne c\in\mathcal O_{(v)}$, of the form $\sum_as_a\otimes b_a$ for the basis $(b_a)$ of the
hypothesis. If $g(t)=0$, then $\sum_as_aR(b_a)=0$, so all $s_a$ vanish by the independence of the
$R(b_a)$. Then $ct=0$, and $t=0$ because $c$ acts invertibly on $V$. For every
$\mathcal O_K$-submodule $A$ of $V$, the map $g$ carries $\mathcal O'\otimes A_{(v)}$ onto
$\mathcal O'R(A)$.

The ring $\mathcal O'$ is torsion-free, hence flat, over $\mathcal O_{(v)}$, and faithfully flat since
$\sigma$ is local. So $\mathcal O'\otimes-$ preserves intersections, quotients and strict inclusions,
and $\operatorname{len}_{\mathcal O'}(\mathcal O'\otimes N)=e\operatorname{len}_v(N)$ for every
$\mathcal O_{(v)}$-module $N$, both sides being finite or infinite together: for the residue field
$N$ of $\mathcal O_{(v)}$ the left side is the length $e$ of $\mathcal O'/\pi\mathcal O'$, with $\pi$
a uniformizer of $\mathcal O_{(v)}$, and both sides are additive in short exact sequences. Hence,
for $\mathcal O_K$-submodules $A,B$ of $V$,
\[
 \operatorname{len}_{\mathcal O'}\Big(\mathcal O'R(A)\big/\big(\mathcal O'R(A)\cap\mathcal O'R(B)\big)\Big)
 =e\operatorname{len}_v\big(A_{(v)}/(A_{(v)}\cap B_{(v)})\big).
\]

By hypothesis the left side is finite for $(A,B)=(\Theta,X)$ and $(X,\Theta)$, so $\Theta_{(v)}$ and
$X_{(v)}$ are commensurable, and
$[\mathcal O'R(\Theta):\mathcal O'R(X)]=e\operatorname{idx}_v(\Theta,X)/(e\,e_v)$. Hence
$\operatorname{idx}_v(\Theta,X)\le e_v\varepsilon$. As $M$ and $\Theta$ are lattices in $Z$, $M_{(v)}$
and $\Theta_{(v)}$ are commensurable as well. For three pairwise commensurable
$\mathcal O_{(v)}$-modules the signed index is additive,
$\operatorname{idx}_v(A,C)=\operatorname{idx}_v(A,B)+\operatorname{idx}_v(B,C)$, since each term is
$\operatorname{len}_v(\cdot/N)-\operatorname{len}_v(\cdot/N)$ for $N=A_{(v)}\cap B_{(v)}\cap C_{(v)}$.
So $\operatorname{idx}_v(M,X)=\operatorname{idx}_v(M,\Theta)+\operatorname{idx}_v(\Theta,X)\le
\operatorname{idx}_v(M,\Theta)+e_v\varepsilon$.
\end{proof}

\subsection{Two-level data and the definite case}\label{subsection-is-two-level-data}

\begin{definition}\label{is-onto-share-flag-data}
Let $E$ be a CM field with complex conjugation $a\mapsto\bar a$, and let $r\geq0$ be an integer. For a
matrix $A\in M_r(E)$ with $\bar A_{ab}=A_{ba}$ put $h(u,v)=\sum_{a,b}u_aA_{ab}\bar v_b$ for $u,v\in E^r$,
and $h_\sigma(u,v)=\sum_{a,b}u_a\sigma(A_{ab})\overline{v_b}$ for an embedding $\sigma\colon E\to\mathbf C$
and $u,v\in\mathbf C^r$; write $\sigma(v)$ and $M^\sigma$ for $\sigma$ applied to the entries of a vector
$v$ or a matrix $M$ over $E$. \emph{Flag data} $X$ of rank $r$ over $E$ consist of:
\begin{enumerate}
\item a set $\Psi$ of embeddings $E\to\mathbf C$ that contains $\tau$ if and only if it does not contain
  $\bar\tau$;
\item such a matrix $A$ with $\det A\ne0$, and an integer $0\le p_0\le r$ such that, for every
  $\tau\in\Psi$, some basis of $\mathbf C^r$ has Gram matrix $J=\operatorname{diag}(1,\dots,1,-1,\dots,-1)$,
  with $p_0$ entries $1$, for $h_\tau$;
\item a finitely generated $\mathcal O_E$-submodule $\Lambda$ of $E^r$ that spans $E^r$, with
  $d\,h(x,y)\in\mathcal O_E$ for $x,y\in\Lambda$ and some $d\ne0$ in $\mathcal O_E$;
\item a subgroup $\Gamma$ of finite index in the group of elements of $GL_r(E)$ that preserve $h$ and
  $\Lambda$, such that $\gamma=1$ whenever $\gamma\in\Gamma$, $x\in\mathcal D_X$ and $b\in\mathcal B^\Psi$
  satisfy $\gamma x=xb$, in the notation below;
\item for every \emph{weight} $\kappa\in\mathbf Z^{\Psi\times\{1,\dots,r\}}$ a function $\chi_\kappa$ on
  $GL_r(E)$ whose restriction to $\Gamma$ is a character of finite order, with
  $\chi_{\kappa+\kappa'}=\chi_\kappa\chi_{\kappa'}$ on $\Gamma$;
\item an element $\xi_0\ne0$ with $\bar\xi_0=-\xi_0$, an element $\kappa_0\in\{2\xi_0,-2\xi_0\}$ with
  $\operatorname{Im}\tau(\kappa_0)<0$ for $\tau\in\Psi$, and an integer $m>0$ such that, for every
  $\xi\in\Lambda$ with $h(\xi,\xi)=0$ and every algebraic integer $s\in E$ with $\bar s=s$, the
  transvection $v\mapsto v+ms\kappa_0h(v,\xi)\xi$ lies in $\Gamma$;
\item a real number $0<k_0\le2$ with $|\tau(\xi_0)|=k_0/2$ for $\tau\in\Psi$, such that
  $\operatorname{Tr}_{E/\mathbf Q}(-\xi_0h(x,y))\in\mathbf Z$ for $x,y\in\Lambda$ and every homomorphism
  $\Lambda\to\mathbf Z$ is $y\mapsto\operatorname{Tr}_{E/\mathbf Q}(-\xi_0h(x,y))$ for some $x\in\Lambda$;
\item a finitely generated $\mathcal O_E$-submodule $L'$ of $E^r$ that is \emph{self-dual} for $h$, that
  is, $h(x,y)$ is an algebraic integer for $x,y\in L'$, and every $x\in E^r$ with $h(x,L')$ integral lies
  in $L'$; integers $N_A,N_C>0$ with $N_A\Lambda\subset L'$ and $N_CL'\subset\Lambda$; and an integer
  $g_X>0$ such that $\Gamma$ contains every element of $GL_r(E)$ that preserves $h$ and $\Lambda$ and
  satisfies $\gamma y-y\in g_X\Lambda$ for $y\in\Lambda$.
\end{enumerate}
For $1\le a\le r$ let $\operatorname{sgn}(a)=1$ if $a\le p_0$ and $\operatorname{sgn}(a)=-1$ otherwise,
and $\epsilon_i=\prod_{a=1}^{i}\operatorname{sgn}(a)$. Let $\mathcal B\subset GL_r(\mathbf C)$ be the group
of invertible upper triangular matrices. For $x=(x_\tau)_{\tau\in\Psi}$ in $M_r(\mathbf C)^\Psi$ let
$d_i(x,\tau)$ be the leading principal $i\times i$ minor of the Gram matrix
$\big(h_\tau(x_\tau e_a,x_\tau e_b)\big)_{a,b}$ of the columns of $x_\tau$, so that $d_0(x,\tau)=1$. The
\emph{flag domain} $\mathcal D_X$ is the set of $x$ with $\det x_\tau\ne0$ and
$\epsilon_i\operatorname{Re}d_i(x,\tau)>0$ for all $\tau\in\Psi$ and $0\le i\le r$; it depends only on
$\Psi$, $A$ and $p_0$. The group $\mathcal B^\Psi$ acts by $(xb)_\tau=x_\tau b_\tau$, and $GL_r(E)$ by
$(\gamma x)_\tau=\gamma^\tau x_\tau$. The point $x$ is \emph{orthonormal} if the Gram matrix of the columns
of every $x_\tau$ is $J$.

A \emph{section} of weight $\kappa$ is a function $P$ on $M_r(\mathbf C)^\Psi$, holomorphic on
$\mathcal D_X$, with $P(xb)=\prod_{\tau,i}(b_\tau)_{ii}^{\kappa_{\tau,i}}P(x)$ and $P(\gamma
x)=\chi_\kappa(\gamma)P(x)$ for $x\in\mathcal D_X$, $b\in\mathcal B^\Psi$ and $\gamma\in\Gamma$. It is
\emph{bounded} if the function
\[
 \|P\|_\kappa(x)=|P(x)|\prod_{\tau\in\Psi}\prod_{i=1}^{r}\Big|\frac{d_i(x,\tau)}{d_{i-1}(x,\tau)}\Big|^{-\kappa_{\tau,i}/2}
\]
is bounded on $\mathcal D_X$. A vector $\xi\in\Lambda$ is \emph{isotropic} if $\xi\ne0$ and $h(\xi,\xi)=0$.
For such $\xi$ and $s\in\mathbf R^\Psi$ let $\phi_\xi(s)x$ have the components $f_\tau x_\tau$, where
$f_\tau(v)=v-ic_\tau s_\tau h_\tau(v,\tau(\xi))\tau(\xi)$ and $c_\tau=2|\tau(\xi_0)|$. Let $o_\kappa>0$ be a
fixed integer with $\chi_\kappa^{o_\kappa}=1$ on $\Gamma$, and $\Pi_\kappa\subset\mathbf R^\Psi$ the subgroup
generated by the vectors $(mo_\kappa\operatorname{Re}\tau(a))_\tau$ for the algebraic integers $a\in E$ with
$\bar a=a$. The \emph{constant term} of $P$ along $\xi$ at $x$ is
$\operatorname{vol}(F_\kappa)^{-1}\int_{F_\kappa}P(\phi_\xi(s)x)\,ds$ for a fixed measurable fundamental
domain $F_\kappa$ of $\Pi_\kappa$ in $\mathbf R^\Psi$ of finite positive volume, and $0$ if there is none.
A \emph{cusp form} of weight $\kappa$ is a bounded section of weight $\kappa$ whose constant terms along
all isotropic $\xi\in\Lambda$ vanish at every $x\in\mathcal D_X$; $S_\kappa(X)$ is the set of cusp forms.

For $x\in\mathcal D_X$ and $\tau\in\Psi$, let $x_1,\dots,x_r$ be the columns of $x_\tau$ and $g$ the Gram
matrix of $x_1,\dots,x_{p_0}$ for $h_\tau$. The \emph{majorant} of $x$ at $\tau$ is
\[
 m_{x,\tau}(v)=2\operatorname{Re}\sum_{a,b\le p_0}h_\tau(v,x_a)(g^{-1})_{ab}\overline{h_\tau(v,x_b)}
 -\operatorname{Re}h_\tau(v,v)\qquad(v\in\mathbf C^r).
\]
There is \emph{no contraction} at $x$ if every $N\in M_r(E)$ with $NL'\subset L'$ for which, at every
$\tau\in\Psi$, some $0\le c_\tau<1$ satisfies $m_{x,\tau}(N^\tau u)\le c_\tau^2m_{x,\tau}(u)$ for all
$u\in\mathbf C^r$, is zero.
\end{definition}
\leanref{ReciprocityTheory/ReductiveGroups/ArithmeticGeometry/FundamentalDomains/Signature.lean}{ReciprocityTheory.ChartRestriction.G.FD.Data}
\leanref{ReciprocityTheory/ReductiveGroups/ArithmeticGeometry/FundamentalDomains/FundamentalDomainPointBound.lean}{ReciprocityTheory.ChartRestriction.G.FD.DataX}
\leanref{ReciprocityTheory/ReductiveGroups/ArithmeticGeometry/FundamentalDomains/Signature.lean}{ReciprocityTheory.ChartRestriction.G.FD.GG}
\leanref{ReciprocityTheory/ReductiveGroups/ArithmeticGeometry/FundamentalDomains/Signature.lean}{ReciprocityTheory.ChartRestriction.G.FD.Data.IsSec}
\leanref{ReciprocityTheory/ReductiveGroups/ArithmeticGeometry/FundamentalDomains/Signature.lean}{ReciprocityTheory.ChartRestriction.G.FD.Data.pnorm}
\leanref{ReciprocityTheory/ReductiveGroups/ArithmeticGeometry/FundamentalDomains/ModerateCoefficientPositivity.lean}{ReciprocityTheory.ChartRestriction.G.FD.Data.fj}
\leanref{ReciprocityTheory/ReductiveGroups/ArithmeticGeometry/FundamentalDomains/ModerateCoefficientPositivity.lean}{ReciprocityTheory.ChartRestriction.G.FD.Data.SS}
\leanref{ReciprocityTheory/ReductiveGroups/ArithmeticGeometry/FundamentalDomains/Signature.lean}{ReciprocityTheory.ChartRestriction.G.FD.Data.maj}
\leanref{ReciprocityTheory/ReductiveGroups/ArithmeticGeometry/ArakelovVolumes/UpperBounds/IsOrth.lean}{ReciprocityTheory.ChartRestriction.G.AR.Up.IsOrth}
\leanref{Development/AbelianVarietyTheory/BlockTheory/PoincareCuspForms/SmallUnipotentElements.lean}{GlobalTheory.IndefiniteRowSideFullHecke.SmallUnipotentElements.NoContraction}

\begin{definition}\label{is-onto-share-two-level}
Let $E$ be a CM field, $r\geq0$ an integer and $U$ a group acting on a set $C$. \emph{Two-level data}
$M$ of rank $r$ over $E$ for $U$ and $C$ consist of $\Psi$, $A$, $p_0$ and $\xi_0$ as in
Definition~\ref{is-onto-share-flag-data}; flag data $X_c$ with these $\Psi$, $A$, $p_0$, $\xi_0$ for $c\in C$,
with groups $\Gamma_c$, characters $\chi_{c,\kappa}$ and modules $\Lambda_c$, $L'_c$, such that every
$\chi_{c,\kappa}$ is trivial on $\Gamma_c$ and $\Gamma_c$ preserves $L'_c$; and elements $t(u,c)\in GL_r(E)$
for $u\in U$ and $c\in C$ that preserve $h$, with $t(u,c)\Lambda_c=\Lambda_{uc}$, $t(u,c)L'_c=L'_{uc}$,
$\gamma\in\Gamma_c$ if and only if $t(u,c)\gamma t(u,c)^{-1}\in\Gamma_{uc}$,
$t(uv,c)\big(t(u,vc)t(v,c)\big)^{-1}\in\Gamma_{uvc}$, and $u=1$ whenever $uc=c$ and $t(u,c)\in\Gamma_c$.
All $X_c$ have the same flag domain $\mathcal D_M$.

For a weight $\kappa$, a \emph{family} of weight $\kappa$ is $F=(F_c)_{c\in C}$ with $F_c\in S_\kappa(X_c)$,
and $(u^*F)_c(x)=F_{uc}(t(u,c)x)$. The space $S_\kappa(M)$ of \emph{cusp forms} of $M$ of weight $\kappa$
is the $\mathbf C$-span of the functions $(c,x)\mapsto F_c(x)$ on $C\times\mathcal D_M$, for the families
$F$ of weight $\kappa$. A representation $\rho$ of $U$ on $S_\kappa(M)$ is \emph{pulled back} from the
action on $C$ if there are $a(u,c)\in\mathbf C^\times$ with $(\rho(u^{-1})F)(c,x)=a(u,c)(u^*F)_c(x)$ for
every family $F$ of weight $\kappa$, every $c$ and every $x\in\mathcal D_M$. For a subgroup $H\subset U$
and a character $\lambda\colon H\to\mathbf C^\times$, $S_\kappa(M)^{H,\lambda}$ is the subspace on which
$\rho(h)$ acts as $\lambda(h)$ for $h\in H$.

Put $\dim X(M)=|\Psi|p_0(r-p_0)$, $\dim Y(M)=|\Psi|r(r-1)/2$ and $s(M)=2r+50r^2\dim X(M)+2$. For an
integer $s$ let $\delta_s$ be the weight with $(\delta_s)_{\tau,i}=-\lfloor s/2\rfloor$ for $i\le p_0$ and
$\lfloor s/2\rfloor$ for $i>p_0$. The data are \emph{definite} if $p_0\in\{0,r\}$. For real numbers
$\varepsilon,c_0$, the cusp forms of $M$ have \emph{ranks} $(\varepsilon,c_0)$ at $\kappa$ and $s$ if
$S_\kappa(M)$, $S_{\kappa+2\delta_s}(M)$ and $S_{\kappa-\delta_s}(M)$ are finite-dimensional,
$(1-\varepsilon)\dim S_\kappa(M)\le\dim S_{\kappa-\delta_s}(M)$ and $\dim S_{\kappa+2\delta_s}(M)\le
c_0\dim S_{\kappa-\delta_s}(M)$.

The \emph{level hypothesis} holds for $M$ if there is a prime $g>3r$ that splits completely in $E$ such
that $(\gamma t(u,c)-1)^rv\in PL'_c$ for all $c\in C$, all $u\in U$ with $uc=c$, all $\gamma\in\Gamma_c$,
all prime ideals $P\ni g$ of $\mathcal O_E$ and all $v\in L'_c$. For a subgroup $H\subset U$, $M|_H$ is $M$
with $U$ replaced by $H$, and the level hypothesis holds for $H$ if it holds for $M|_H$.

Let $n=\dim Y(M)$ and $k_{\tau,i}=2i-r-1$. Weights $L,F$ and an integer $T$ are \emph{admissible} for
$M$ if, with $F'=F-k-(n+2)L$, for every $\tau\in\Psi$: $L_{\tau,i}-L_{\tau,i+1}\geq1$ and
$F'_{\tau,i}-F'_{\tau,i+1}\geq0$ for $1\le i<r$ with $i\ne p_0$, and $L_{\tau,r}-L_{\tau,1}\geq1.9T$ and
$F'_{\tau,r}-F'_{\tau,1}\geq1$. The data satisfy the \emph{archimedean condition} at $T$ for a weight
$\kappa$ if $\kappa=F+TL$ for weights $L,F$ admissible with $T$, such that every $X_c$ carries an
\emph{archimedean model} for $(L,F,T)$: Alpöge's archimedean data of $X_c$, consisting of sections
of $X_c$ of weights built from $L$ and $F$, invariant functions on $\mathcal D_M$ with bounds for their
Levi forms, and further conditions, which this book does not unfold.
\end{definition}
\leanref{Development/AbelianVarietyTheory/BlockTheory/PoincareCuspForms/TwoLevelLogarithm.lean}{GlobalTheory.IndefiniteRowSideFullHecke.FrameWeightGapBounds.TwoLevels}
\leanref{Development/AbelianVarietyTheory/BlockTheory/PoincareCuspForms/TwoLevelLogarithm.lean}{GlobalTheory.IndefiniteRowSideFullHecke.FrameWeightGapBounds.TwoLevels.cusp}
\leanref{Development/AbelianVarietyTheory/BlockTheory/PoincareCuspForms/TwoLevelLogarithm.lean}{GlobalTheory.IndefiniteRowSideFullHecke.FrameWeightGapBounds.TwoLevels.pull}
\leanref{Development/AbelianVarietyTheory/BlockTheory/PoincareCuspForms/TwoLevelLogarithm.lean}{GlobalTheory.IndefiniteRowSideFullHecke.FrameWeightGapBounds.deMul}
\leanref{Development/AbelianVarietyTheory/BlockTheory/PoincareCuspForms/TwoLevelLogarithm.lean}{GlobalTheory.IndefiniteRowSideFullHecke.FrameWeightGapBounds.LevelHyp}
\leanref{Development/AbelianVarietyTheory/Definitions/Divisors/ExistsArchOne.lean}{GlobalTheory.ShareOn.IsPullRep}
\leanref{Development/AbelianVarietyTheory/Definitions/Divisors/ExistsArchOne.lean}{GlobalTheory.ShareOn.Ranks}
\leanref{Development/AbelianVarietyTheory/Definitions/Divisors/ExistsArchOne.lean}{GlobalTheory.ShareOn.ArchOn}
\leanref{Development/AbelianVarietyTheory/Definitions/Divisors/ExistsArchOne.lean}{GlobalTheory.ShareOn.sOf}
\leanref{Development/AbelianVarietyTheory/Definitions/Divisors/BadElt.lean}{GlobalTheory.BadCount.eig}

\begin{lemma}\label{is-onto-share-majorant}
Let $X$ be flag data of rank $r$ over $E$ (Definition~\ref{is-onto-share-flag-data}), let
$x\in\mathcal D_X$ and $\tau\in\Psi$, and write $|u|^2=\sum_a|u_a|^2$.
\begin{enumerate}
\item $m_{x,\tau}(v)\geq0$ for $v\in\mathbf C^r$, and $m_{x,\tau}(v)>0$ for $v\ne0$.
\item For $b\in\mathcal B^\Psi$, $xb\in\mathcal D_X$ and $m_{xb,\tau}=m_{x,\tau}$.
\item Every orthonormal point lies in $\mathcal D_X$, and $x=x'b$ for an orthonormal $x'$ and some
  $b\in\mathcal B^\Psi$.
\item If $x$ is orthonormal, then $\det x_\tau\ne0$ and $m_{x,\tau}(v)=|x_\tau^{-1}v|^2$.
\item If $p_0\in\{0,r\}$, then $\mathcal D_X$ is the set of points with $\det x_\tau\ne0$ for all $\tau$, and
  $m_{x,\tau}(\beta^\tau v)=m_{x,\tau}(v)$ for every $\beta\in GL_r(E)$ that preserves $h$.
\end{enumerate}
\end{lemma}
\leanref{ReciprocityTheory/ReductiveGroups/ArithmeticGeometry/FundamentalDomains/MajorantConjugateTransfer.lean}{ReciprocityTheory.ChartRestriction.G.FD.MajorantConjugateTransfer.maj_nonneg}
\leanref{ReciprocityTheory/ReductiveGroups/ArithmeticGeometry/FundamentalDomains/MajPos.lean}{ReciprocityTheory.ChartRestriction.G.FD.Data.maj_pos}
\leanref{ReciprocityTheory/ReductiveGroups/ArithmeticGeometry/FundamentalDomains/OrthogonalFrameDomain.lean}{ReciprocityTheory.ChartRestriction.G.FD.OrthogonalFrameDomain.rmul_mem_dom}
\leanref{Development/AbelianVarietyTheory/Proofs/TorusHaarMeasure/MajLactRmul.lean}{GlobalTheory.MajInv.maj_rmul}
\leanref{ReciprocityTheory/ReductiveGroups/ArithmeticGeometry/ArakelovVolumes/UpperBounds/OrthoStatementPf.lean}{ReciprocityTheory.ChartRestriction.G.AR.Up.ortho_holds}
\leanref{ReciprocityTheory/ReductiveGroups/ArithmeticGeometry/ArakelovVolumes/UpperBounds/InterpStatementPfH.lean}{ReciprocityTheory.ChartRestriction.G.AR.Up.IW.orth_mem_dom}
\leanref{Development/AbelianVarietyTheory/Proofs/TorusHaarMeasure/MajorantEqSumNormSqOrthonormalFrame.lean}{GlobalTheory.MajorantOrthFrame.maj_eq_sum_norm_sq_of_isOrth}
\leanref{Development/AbelianVarietyTheory/Proofs/DivisorProofs/DefiniteFrameLeadingMinorsSignStatement.lean}{GlobalTheory.T.re_dmin_sign_pos_of_definite}
\leanref{ReciprocityTheory/ReductiveGroups/ArithmeticGeometry/FundamentalDomains/MajorantConjugateTransfer.lean}{ReciprocityTheory.ChartRestriction.G.FD.MajorantConjugateTransfer.maj_transfer}
\begin{proof}
Write $h=h_\tau$, let $x_1,\dots,x_r$ be the columns of $x_\tau$, $G=(h(x_a,x_b))_{a,b}$ their Gram
matrix and $W$ the span of $x_1,\dots,x_{p_0}$. Since $\bar A_{ab}=A_{ba}$, $h(v,u)=\overline{h(u,v)}$, so
$G$ is Hermitian and its leading minors $d_i(x,\tau)$ are real.

(1) For $i\le p_0$, $\epsilon_i=1$, so $d_i(x,\tau)>0$. These are the leading minors of the upper
left $p_0\times p_0$ block $g$ of $G$, so $g$ is positive definite by Sylvester's criterion
\cite[Theorem~7.2.5]{HJ13}. For $v\in\mathbf C^r$ put $c_a=\sum_{b\le p_0}h(v,x_b)(g^{-1})_{ba}$ and
$y=\sum_{a\le p_0}c_ax_a\in W$. Then $h(y,x_{a'})=h(v,x_{a'})$ for $a'\le p_0$, so $w=v-y$ is orthogonal
to $W$, and the sum in the definition of $m_{x,\tau}(v)$ is $h(y,v)$. As $h(y,w)=h(w,y)=0$,
\[
 m_{x,\tau}(v)=2h(y,y)-\big(h(y,y)+h(w,w)\big)=h(y,y)-h(w,w).
\]
Here $h(y,y)\geq0$, with equality only if $y=0$, since $g$ is positive definite. Suppose $w\ne0$ and
$h(w,w)\geq0$. By (2) of Definition~\ref{is-onto-share-flag-data} there are vectors
$n_1,\dots,n_{r-p_0}$ with $h(n_b,n_{b'})=-\delta_{bb'}$. The $r+1$ vectors $w,x_1,\dots,x_{p_0},n_1,
\dots,n_{r-p_0}$ satisfy a nontrivial relation $\alpha w+y'=\sum_b\gamma_bn_b$ with $y'=\sum_a\beta_ax_a$.
Call this vector $u$. As $w$ is orthogonal to $y'$, $h(u,u)=|\alpha|^2h(w,w)+h(y',y')\geq0$, while
$h(u,u)=-\sum_b|\gamma_b|^2$. So all $\gamma_b$ vanish and $u=0$. Then $0=h(u,y')=h(y',y')$, so all
$\beta_a$ vanish, and then $\alpha=0$, against the nontriviality. Hence $h(w,w)<0$ when $w\ne0$, and
(1) follows.

(2) Let $b\in\mathcal B^\Psi$. The Gram matrix of the columns of $x_\tau b_\tau$ is $b_\tau^tG\bar b_\tau$.
As $b_\tau$ is upper triangular, the leading $i\times i$ block of this matrix is $\beta_i^tG_i\bar\beta_i$,
where $G_i$ and $\beta_i$ are the leading $i\times i$ blocks of $G$ and $b_\tau$. So
$d_i(xb,\tau)=|\det\beta_i|^2d_i(x,\tau)$ with $\det\beta_i\ne0$, and $\det(x_\tau b_\tau)\ne0$; hence
$xb\in\mathcal D_X$. The first $p_0$ columns of $x_\tau b_\tau$ are combinations of $x_1,\dots,x_{p_0}$,
and conversely, as $b_\tau^{-1}$ is upper triangular; so they span $W$. By the proof of (1),
$m_{x,\tau}(v)=2\operatorname{Re}h(y,v)-\operatorname{Re}h(v,v)$, where $y\in W$ and $v-y$ is orthogonal
to $W$. Such a $y$ is unique: the difference of two of them lies in $W$ and is orthogonal to $W$, so it
vanishes as $h$ is positive definite on $W$. The same description holds for $xb$, with the same $W$,
so $m_{xb,\tau}=m_{x,\tau}$.

(3) Let $x$ be orthonormal. Then $d_i(x,\tau)=\prod_{a\le i}\operatorname{sgn}(a)=\epsilon_i$, so
$\epsilon_id_i(x,\tau)=1>0$. The Gram matrix is $x_\tau^tA^\tau\bar x_\tau$, so its determinant
$\det(x_\tau)\tau(\det A)\overline{\det x_\tau}=\det J$ is nonzero, and $\det x_\tau\ne0$. So
$x\in\mathcal D_X$.

Now let $x\in\mathcal D_X$ be arbitrary. The leading blocks $G_c$ are invertible, as $d_c(x,\tau)\ne0$. Let
$u$ be the upper unitriangular matrix whose $c$-th column has the entries
$u_{jc}=-(\rho_cG_{c-1}^{-1})_j$ for $j<c$, where $\rho_c=(G_{c1},\dots,G_{c,c-1})$. Then
$(u^tG)_{cm}=G_{cm}-(\rho_cG_{c-1}^{-1}G_{c-1})_m=0$ for $m<c$, so the entries of $u^tG\bar u$ below
the diagonal vanish. It is Hermitian, hence diagonal, equal to $\operatorname{diag}(\delta_1,\dots,\delta_r)$ with real
$\delta_k$. Since $u$ is unitriangular, $u^tG\bar u$ has the leading minors of $G$, so
$d_i(x,\tau)=\delta_1\cdots\delta_i$. Then
$\epsilon_kd_k(x,\tau)=\epsilon_{k-1}d_{k-1}(x,\tau)\operatorname{sgn}(k)\delta_k$ with both $\epsilon$-terms
positive, so $\operatorname{sgn}(k)\delta_k>0$. Put
$b_\tau=u\operatorname{diag}(|\delta_1|^{-1/2},\dots,|\delta_r|^{-1/2})$, an element of $\mathcal B$. The
Gram matrix of the columns of $x_\tau b_\tau$ is $\operatorname{diag}(\delta_k/|\delta_k|)=J$. So
$x'=xb$ is orthonormal, and $x=x'b^{-1}$ with $b^{-1}\in\mathcal B^\Psi$.

(4) Let $x$ be orthonormal. Then $\det x_\tau\ne0$ by (3). Put $c=x_\tau^{-1}v$, so that $v=\sum_ac_ax_a$.
Then $h(v,x_a)=\operatorname{sgn}(a)c_a$, $h(v,v)=\sum_a\operatorname{sgn}(a)|c_a|^2$, and $g$ is the
identity matrix. So
$m_{x,\tau}(v)=2\sum_{a\le p_0}|c_a|^2-\sum_a\operatorname{sgn}(a)|c_a|^2=\sum_a|c_a|^2=|x_\tau^{-1}v|^2$.

(5) The Gram matrix of a basis for the definite form $h_\tau$ has leading minors of the signs
$\epsilon_i$, and $W$ is $0$ or $\mathbf C^r$. By the proof of (2), $m_{x,\tau}$ depends only on $h_\tau$
and $W$, and $\beta^\tau$ preserves both.
\end{proof}

\begin{lemma}\label{is-onto-share-free-summands}
Let $K$ be a field, $\Gamma$ a finite group and $\rho$ a representation of $\Gamma$ on a finite-dimensional
$K$-vector space $V$. Let $H\subset\Gamma$ be a subgroup acting on a finite set $C'$ such that $hc=c$ implies
$h=1$, and let $V=\bigoplus_{c\in C'}A_c$ with $\rho(h)A_c\subset A_{hc}$ for $h\in H$.
\begin{enumerate}
\item There are an integer $n$ with $n|H|=\dim V$ and an injective linear map $j\colon V\to K[H]^n$ with
  $j(\rho(h)v)=[h]\,j(v)$ for $h\in H$, the product taken in each coordinate.
\item If $K$ has characteristic $0$, then $|H'|\dim V^{H',\lambda}=\dim V$ for every subgroup $H'\subset H$
  and every character $\lambda\colon H'\to K^\times$, where $V^{H',\lambda}$ is the subspace on which
  $\rho(h)$ acts as $\lambda(h)$ for $h\in H'$.
\end{enumerate}
\end{lemma}
\leanref{Development/AbelianVarietyTheory/Proofs/DivisorProofs/EmbedsInRegularRepresentationCopiesOfFreePermuteStatement.lean}{GlobalTheory.T.embedsInRegularRepresentationCopies_of_free_permuted_summands}
\leanref{Development/AbelianVarietyTheory/Proofs/DivisorProofs/EmbedsInRegularRepresentationCopiesOfFreePermute.lean}{GlobalTheory.BadCount.embReg_of_free_permuted_summands}
\leanref{Development/AbelianVarietyTheory/Proofs/DivisorProofs/CardMulFinrankCharacterSpaceOfFreePermutedSummanStatement.lean}{GlobalTheory.T.card_mul_finrank_characterSpace_of_free_permuted_summands}
\leanref{Development/AbelianVarietyTheory/Proofs/DivisorProofs/CardMulFinrankCharacterSpaceOfFreePermutedSumman.lean}{GlobalTheory.BadCount.card_mul_finrank_eig_of_free}
\begin{proof}
Write $p_c\colon V\to A_c$ for the projections; then $p_{hc}(\rho(h)v)=\rho(h)p_c(v)$. (1) Since $\rho(h)$
maps $A_c$ injectively into $A_{hc}$ for every $h$, $\dim A_{hc}=\dim A_c$. Choose a point $c_o$ in each
orbit $o$ of $H$. As the action is free, $(h,o)\mapsto hc_o$ is a bijection from $H\times C'/H$ onto $C'$,
so $n=\sum_o\dim A_{c_o}$ satisfies $n|H|=\dim V$. For each $o$ and each coordinate function $f$ of a
basis of $A_{c_o}$ put $j_f(v)=\sum_{h\in H}f\big(p_{c_o}(\rho(h)^{-1}v)\big)[h]$; then
$j_f(\rho(g)v)=[g]j_f(v)$, by the substitution $h\mapsto gh$. If all $j_f(v)$ vanish, then
$p_{c_o}(\rho(h)^{-1}v)=0$ for all $h$ and $o$, so $p_{hc_o}(v)=\rho(h)p_{c_o}(\rho(h)^{-1}v)=0$ and $v=0$.

(2) The operator $\frac1{|H'|}\sum_{h\in H'}\lambda(h)^{-1}\rho(h)$ is a projection onto $V^{H',\lambda}$,
so its trace is $\dim V^{H',\lambda}$. For $h\ne1$, $\rho(h)$ maps every $A_c$ into $A_{hc}$ with $hc\ne c$,
so in a basis adapted to the decomposition its diagonal entries vanish and its trace is $0$. Hence the
trace of the projection is $\dim V/|H'|$.
\end{proof}

\begin{lemma}\label{is-onto-share-components}
Let $M$ be two-level data for a group $U$ and a finite set $C$ (Definition~\ref{is-onto-share-two-level}),
let $\kappa$ be a weight, and for $c\in C$ let $A_c\subset S_\kappa(M)$ be the subspace of cusp forms that
vanish on $\{c'\}\times\mathcal D_M$ for $c'\ne c$.
\begin{enumerate}
\item $S_\kappa(M)=\bigoplus_cA_c$, and if $\rho$ is a representation of $U$ on $S_\kappa(M)$ pulled back
  from the action on $C$, then $\rho(u)A_c\subset A_{uc}$.
\item If, for every $c$, the restrictions to $\mathcal D_M$ of the elements of $S_\kappa(X_c)$ span a
  finite-dimensional space, then $S_\kappa(M)$ is finite-dimensional.
\end{enumerate}
\end{lemma}
\leanref{Development/AbelianVarietyTheory/Proofs/DivisorProofs/TwoLevelCuspFormsDirectSumComponentsStatement.lean}{GlobalTheory.T.twoLevelCuspForms_directSum_components}
\leanref{Development/AbelianVarietyTheory/Proofs/DivisorProofs/TwoLevelCuspFormsDirectSumComponents.lean}{GlobalTheory.ShareOn.Decomp.pullRep_onComp}
\leanref{Development/AbelianVarietyTheory/Proofs/DivisorProofs/TwoLevelCuspFiniteOfComponentsStatement.lean}{GlobalTheory.T.twoLevelCusp_finite_of_components}
\begin{proof}
(1) The zero function is a cusp form. So for a family $F$ and $c\in C$, the family $F^{(c)}$ equal to $F_c$
at $c$ and to $0$ elsewhere is again a family; the map $f\mapsto f\cdot1_{\{c\}\times\mathcal D_M}$ therefore
carries $S_\kappa(M)$ into $A_c$, and into the span of the families that vanish off $c$. Every $f$ is the
sum of these components, and the $A_c$ are independent since their supports are disjoint. If $F$ vanishes
off $c$, then $(\rho(u^{-1})F)(c',x)=a(u,c')F_{uc'}(t(u,c')x)$ vanishes unless $uc'=c$; by linearity
$\rho(u^{-1})A_c\subset A_{u^{-1}c}$.

(2) Restriction to $\{c\}\times\mathcal D_M$ maps $S_\kappa(M)$ into the span for $X_c$, and the product of
these maps over the finitely many $c$ is injective.
\end{proof}

\begin{lemma}[Borel-Weil \cite{Ser54,Jan03}]\label{is-onto-share-borel-weil}
Let $r\geq0$ and $a_1,\dots,a_r\geq0$ be integers, and let $P$ be a holomorphic function on
$GL_r(\mathbf C)$ with $P(yb)=\prod_ib_{ii}^{a_i}P(y)$ for $y\in GL_r(\mathbf C)$ and $b\in\mathcal B$. Then
there is a homogeneous polynomial $F$ of degree $\sum_ia_i$ in the entries with $P(y)=F(y)$ on
$GL_r(\mathbf C)$.
\end{lemma}
\leanref{ReciprocityTheory/SemistableReduction/LatAll.lean}{ReciprocityTheory.MultiplicativeHeights.Closed.bwPoly}
\leanref{ReciprocityTheory/SemistableReduction/TorsionNonvanishStatement.lean}{ReciprocityTheory.MultiplicativeHeights.BWPolyStatement}

\begin{lemma}\label{is-onto-share-definite-finite}
Let $X$ be flag data of rank $r$ over $E$ with $p_0\in\{0,r\}$, and let $\kappa$ be a weight. Then the
restrictions to $\mathcal D_X$ of the cusp forms of weight $\kappa$ span a finite-dimensional space.
\end{lemma}
\leanref{Development/AbelianVarietyTheory/Proofs/DivisorProofs/FiniteDimensionalCuspFormsDefiniteUnitaryFlagDomStatement.lean}{GlobalTheory.T.finiteDimensional_cuspForms_definiteUnitaryFlagDomain}
\leanref{Development/AbelianVarietyTheory/Proofs/DivisorProofs/FiniteDimensionalCuspFormsDefiniteUnitaryFlagDom.lean}{GlobalTheory.DefiniteFin.finiteDimensional_cusp_definite}
\leanref{Development/AbelianVarietyTheory/Proofs/DivisorProofs/FiniteDimensionalCuspFormsDefiniteUnitaryFlagDom.lean}{GlobalTheory.DefiniteFin.finiteDimensional_span_of_slices}
\leanref{Development/AbelianVarietyTheory/Proofs/DivisorProofs/FrameSectionSliceEqPolyDivDetPow.lean}{GlobalTheory.section_slice_eq_poly_div_det_pow}
\begin{proof}
By Lemma~\ref{is-onto-share-majorant}(5), $\mathcal D_X=GL_r(\mathbf C)^\Psi$. Put
$N=\sum_{\tau,i}|\kappa_{\tau,i}|$, so that $\kappa_{\tau,i}+N\geq0$. Let $P$ be a section of weight
$\kappa$, $x\in\mathcal D_X$ and $\tau\in\Psi$, and let $x[y]$ be $x$ with its $\tau$-component replaced by
$y\in GL_r(\mathbf C)$. Applying the equivariance of $P$ to the element of $\mathcal B^\Psi$ equal to $b$ at
$\tau$ and to $1$ elsewhere shows that $y\mapsto P(x[y])\det(y)^N$ satisfies the hypothesis of
Lemma~\ref{is-onto-share-borel-weil} with $a_i=\kappa_{\tau,i}+N$. So $P(x[y])=F(y)\det(y)^{-N}$ with $F$
homogeneous of degree $d_\tau=\sum_i(\kappa_{\tau,i}+N)$: every slice of $P$ in the $\tau$-th variable lies
in the finite-dimensional space $W_\tau$ of the functions $F\det^{-N}$ with $\deg F\le d_\tau$.

Let $V$ be the space of functions on $GL_r(\mathbf C)^\Psi$ all of whose slices in the $\tau$-th variable
lie in $W_\tau$, for every $\tau$. Choose finite sets $T_\tau\subset GL_r(\mathbf C)$ such that an element of
$W_\tau$ that vanishes on $T_\tau$ vanishes identically. If $f\in V$ vanishes on $\prod_\tau T_\tau$, then,
changing one variable at a time, $f$ vanishes on every product of sets $T_\tau$ and $GL_r(\mathbf C)$, and
finally everywhere. So restriction to the finite set $\prod_\tau T_\tau$ is injective on $V$, $V$ is
finite-dimensional, and it contains the restrictions of all cusp forms of weight $\kappa$.
\end{proof}

\begin{lemma}\label{is-onto-share-unipotent}
Let $K$ be a number field, $r,g\geq0$ integers and $\gamma\in M_r(K)$, and for every embedding
$\sigma\colon K\to\mathbf C$ let $Q_\sigma\in GL_r(\mathbf C)$. Write $|u|^2=\sum_a|u_a|^2$. Suppose:
\begin{enumerate}
\item $\operatorname{tr}(\gamma^n)-r\in g\mathcal O_K$ for $1\le n\le r$;
\item $(e^{1/3}+1)r<g$;
\item $|Q_\sigma\gamma^\sigma Q_\sigma^{-1}u|^2\le e^{2/(3r)}|u|^2$ for every $\sigma$ and every
  $u\in\mathbf C^r$.
\end{enumerate}
Then $(\gamma-1)^r=0$, and $\sum_{a,b}\big|\big(Q_\sigma(\gamma^\sigma-1)Q_\sigma^{-1}\big)_{ab}\big|^2<1$ for
every $\sigma$.
\end{lemma}
\leanref{ReciprocityTheory/SemistableReduction/Unipotent.lean}{ReciprocityTheory.MultiplicativeHeights.unipotent_holds}
\leanref{ReciprocityTheory/SemistableReduction/LatSemiSt.lean}{ReciprocityTheory.MultiplicativeHeights.UnipotentStatement}
\begin{proof}
Put $A_\sigma=Q_\sigma\gamma^\sigma Q_\sigma^{-1}$, so that $\operatorname{tr}(A_\sigma^n)=\sigma(\operatorname{tr}\gamma^n)$.
By (3), $|A_\sigma^nu|^2\le e^{2n/(3r)}|u|^2$; for a basis vector $u=e_i$ this gives
$|(A_\sigma^n)_{ii}|^2\le e^{2n/(3r)}\le e^{2/3}$ when $n\le r$, and so $|\operatorname{tr}A_\sigma^n|\le
re^{1/3}$.

Let $1\le n\le r$ and write $\operatorname{tr}\gamma^n-r=gz$ with $z\in\mathcal O_K$. If $z\ne0$, then
$|\sigma(z)|\geq1$ for some $\sigma$, since the product of the $|\sigma(z)|$ is $|N_{K/\mathbf Q}(z)|\geq1$;
then $g\le g|\sigma(z)|=|\operatorname{tr}A_\sigma^n-r|\le re^{1/3}+r$, against (2). So
$\operatorname{tr}\gamma^n=r$ for $0\le n\le r$. Hence $\operatorname{tr}P(\gamma)=rP(1)$ for every
polynomial $P$ of degree at most $r$, and $\operatorname{tr}(\gamma-1)^k=0$ for $1\le k\le r$. In
characteristic $0$ a matrix $N$ of size $r$ with $\operatorname{tr}N^k=0$ for $1\le k\le r$ satisfies
$N^r=0$, by Newton's identities; so $(\gamma-1)^r=0$.

For any $B\in M_r(\mathbf C)$, $\sum_{a,b}|(B-1)_{ab}|^2=\sum_{a,b}|B_{ab}|^2-2\operatorname{Re}\operatorname{tr}B+r$.
For $B=A_\sigma$, $\sum_{a,b}|B_{ab}|^2=\sum_b|A_\sigma e_b|^2\le re^{2/(3r)}$ by (3), and
$\operatorname{tr}A_\sigma=\sigma(\operatorname{tr}\gamma)=r$. As $Q_\sigma(\gamma^\sigma-1)Q_\sigma^{-1}=A_\sigma-1$,
the sum in question is at most $r(e^{2/(3r)}-1)$, which is less than $1$ for $r\geq1$ because
$e^y-1<\frac32y$ for $0<y\le\frac23$; for $r=0$ it is $0$.
\end{proof}

\begin{lemma}\label{is-onto-share-trace}
Let $X$ be flag data of rank $r$ over $E$, and let $g$ be a prime that splits completely in $E$. Let
$\delta\in GL_r(E)$ preserve $h$ and $L'$, and suppose that for every prime $P$ of $\mathcal O_E$ over $g$
there is $k\geq0$ with $(\delta-1)^kL'\subset PL'$. Then $\operatorname{tr}\delta-r\in g\mathcal O_E$.
\end{lemma}
\leanref{Development/AbelianVarietyTheory/Proofs/DivisorProofs/SmallUnipotentUnitaryElementIsOneProof.lean}{GlobalTheory.SmallElement.trace_sub_mem}
\leanref{Development/AbelianVarietyTheory/Proofs/TorusHaarMeasure/TraceMemPrimeOfPowMulVecMemSmul.lean}{GlobalTheory.trace_mem_prime_of_pow_mulVec_mem_smul}
\begin{proof}
The matrix $N=\delta-1$ maps $L'$ into $L'$, and $\operatorname{tr}N=\operatorname{tr}\delta-r$. Let $P$ be a
prime over $g$ and $k$ as in the hypothesis, and $m=k+1$, so that $N^mL'\subset PL'$. As $L'$ is finitely
generated, there is a monic $q\in\mathcal O_E[X]$ whose coefficient of $X^i$ lies in $P^{\deg q-i}$ for
$i<\deg q$, with $q(N^m)=0$ on $L'$, hence on $E^r$, which $L'$ spans. The characteristic polynomial
$\chi_N$ divides $q(X^m)^r$; both are monic, and $\chi_N$ has coefficients in $\mathcal O_E$, so the
divisibility holds in $\mathcal O_E[X]$. Modulo $P$, $q(X^m)^r$ is a power of $X$, so $\chi_N\equiv X^r$, and
its coefficient $-\operatorname{tr}N$ of $X^{r-1}$ lies in $P$. Since $g$ splits completely, $g\mathcal O_E$ is
the product of the distinct primes over $g$, so $\operatorname{tr}N\in g\mathcal O_E$.
\end{proof}

\begin{lemma}\label{is-onto-share-small}
Let $X$ be flag data of rank $r$ over $E$, and $g>3r$ a prime that splits completely in $E$. Let $\Gamma^\circ$
be a subgroup of $GL_r(E)$ whose elements preserve $h$ and $L'$ and satisfy, for every prime $P$ of
$\mathcal O_E$ over $g$, $(\gamma-1)^kL'\subset PL'$ for some $k\geq0$. Let $x\in\mathcal D_X$ be a point at
which there is no contraction, and let $\gamma\in\Gamma^\circ$ satisfy $m_{x,\tau}(\gamma^\tau u)\le
e^{2/(3r)}m_{x,\tau}(u)$ for all $\tau\in\Psi$ and $u\in\mathbf C^r$. Then $\gamma=1$.
\end{lemma}
\leanref{Development/AbelianVarietyTheory/Proofs/DivisorProofs/SmallUnipotentUnitaryElementIsOneProofStatement.lean}{GlobalTheory.T.smallUnipotentUnitaryElementIsOne_proof}
\leanref{Development/AbelianVarietyTheory/Proofs/DivisorProofs/SmallUnipotentUnitaryElementIsOneProof.lean}{GlobalTheory.SmallElement.smallElement}
\begin{proof}
For $r=0$ there is nothing to prove. By Lemma~\ref{is-onto-share-majorant}(2) to (4), $x=x'b$ with $x'$
orthonormal, $b\in\mathcal B^\Psi$ and $m_{x,\tau}(v)=|x'^{-1}_\tau v|^2$. For an embedding $\sigma$ of $E$
put $Q_\sigma=x'^{-1}_\sigma$ if $\sigma\in\Psi$, and $Q_\sigma=\overline{x'^{-1}_{\bar\sigma}}$ otherwise, in
which case $\bar\sigma\in\Psi$. For $\tau\in\Psi$, the hypothesis at $v=x'_\tau u$ reads
$|Q_\tau\gamma^\tau Q_\tau^{-1}u|^2\le e^{2/(3r)}|u|^2$. For $\sigma\notin\Psi$,
$Q_\sigma\gamma^\sigma Q_\sigma^{-1}$ is the complex conjugate of $Q_{\bar\sigma}\gamma^{\bar\sigma}Q_{\bar\sigma}^{-1}$,
so the same bound holds. Every power $\gamma^n$ lies in $\Gamma^\circ$, so $\operatorname{tr}\gamma^n-r\in
g\mathcal O_E$ by Lemma~\ref{is-onto-share-trace}. Since $e^{1/3}<2$ and $g>3r$, $(e^{1/3}+1)r<g$.
Lemma~\ref{is-onto-share-unipotent} gives $S_\sigma=\sum_{a,b}|(Q_\sigma(\gamma^\sigma-1)Q_\sigma^{-1})_{ab}|^2<1$
for every $\sigma$.

Put $N=\gamma-1$; it maps $L'$ into $L'$. For $\tau\in\Psi$ and $u\in\mathbf C^r$,
\[
 m_{x,\tau}(N^\tau u)=\big|\big(Q_\tau(\gamma^\tau-1)Q_\tau^{-1}\big)Q_\tau u\big|^2\le S_\tau|Q_\tau u|^2=S_\tau\,m_{x,\tau}(u),
\]
since the operator norm is at most the Hilbert-Schmidt norm. So $N$ is contracting at every $\tau$ with
$c_\tau=S_\tau^{1/2}<1$, and $N=0$ as there is no contraction at $x$.
\end{proof}

\begin{lemma}\label{is-onto-share-degree}
Let $K$ be a number field and $r\geq0$ an integer, and for every embedding $\sigma\colon K\to\mathbf C$ let
$H_\sigma$ be a positive definite Hermitian $r\times r$ matrix. For families $b=(b_1,\dots,b_k)$ and
$c=(c_1,\dots,c_m)$ in $K^r$ put $G_\sigma(b,c)=\big(\sigma(b_i)^*H_\sigma\sigma(c_j)\big)_{i,j}$ and
$G_\sigma(b)=G_\sigma(b,b)$. For an $\mathcal O_K$-submodule $L_1\subset K^r$ and elements $b_i\in L_1$
that form a basis of $KL_1$ put
\[
 \deg(L_1,b)=\log\Big[L_1:\textstyle\sum_i\mathcal O_Kb_i\Big]-\frac12\sum_\sigma\log\det G_\sigma(b).
\]
Let $L\subset K^r$ be a finitely generated $\mathcal O_K$-submodule, and $b_1,\dots,b_k$,
$c_1,\dots,c_m$ elements of $L$ that together form a basis of $KL$. Put $L_b=L\cap\sum_iKb_i$,
$Q_\sigma=G_\sigma(c)-G_\sigma(c,b)G_\sigma(b)^{-1}G_\sigma(b,c)$ and
\[
 \deg(L;b,c)=\log\Big[L:L_b+\textstyle\sum_j\mathcal O_Kc_j\Big]-\frac12\sum_\sigma\log\det Q_\sigma.
\]
Then the indices are finite, $\det G_\sigma(b)$ and $\det Q_\sigma$ are positive, and
$\deg\big(L,(b,c)\big)=\deg(L_b,b)+\deg(L;b,c)$.
\end{lemma}
\leanref{ReciprocityTheory/SemistableReduction/DegAddHolds.lean}{ReciprocityTheory.MultiplicativeHeights.degAdd_holds}
\leanref{ReciprocityTheory/SemistableReduction/LatSemiSt.lean}{ReciprocityTheory.MultiplicativeHeights.DegAddStatement}
\begin{proof}
Each generator of $L$ has a multiple by a nonzero integer in $B+C$, where $B=\sum_i\mathcal O_Kb_i$ and
$C=\sum_j\mathcal O_Kc_j$; so $[L:B+C]$ is finite. As $b$ and $c$ are independent,
$L_b\cap(B+C)=B$, so $[L_b+C:B+C]=[L_b:B]$ and
\[
 [L:B+C]=[L_b:B]\cdot[L:L_b+C].
\]
For each $\sigma$ the vectors $\sigma(b_i),\sigma(c_j)$ are linearly independent over $\mathbf C$, so the Gram
matrices $G_\sigma(b)$ and $G_\sigma((b,c))$ are positive definite. The latter has the blocks
$G_\sigma(b)$, $G_\sigma(b,c)$, $G_\sigma(c,b)$, $G_\sigma(c)$, so the Schur complement formula gives
$\det G_\sigma((b,c))=\det G_\sigma(b)\det Q_\sigma$, and $\det Q_\sigma>0$. Taking logarithms gives the
claim.
\end{proof}

\begin{lemma}\label{is-onto-semistable-basis}
Let $K$ be a number field and $r\geq0$ an integer, and for every embedding $\sigma\colon K\to\mathbf C$ let
$H_\sigma$ be a positive definite Hermitian $r\times r$ matrix; write $G_\sigma$ and $\deg$ as in
Lemma~\ref{is-onto-share-degree}. Let $L_1\subset K^r$ be a finitely generated $\mathcal O_K$-submodule, and
let $b=(b_1,\dots,b_k)$ and $b'=(b'_1,\dots,b'_k)$ be families in $L_1$ that are both bases of $KL_1$. Then
$\deg(L_1,b)=\deg(L_1,b')$.
\end{lemma}
\leanref{ReciprocityTheory/SemistableReduction/SemiDegB.lean}{ReciprocityTheory.MultiplicativeHeights.latDegBasis_holds}
\leanref{ReciprocityTheory/SemistableReduction/SemiDegBB.lean}{ReciprocityTheory.MultiplicativeHeights.relIndex_span_comb}
\begin{proof}
Two facts are used. First, if $a_1,\dots,a_k\in K^r$ are linearly independent over $K$ and $N\in
M_k(\mathcal O_K)$ has $\det N\ne0$, then the elements $\sum_iN_{ij}a_i$ span a submodule of index
$|N_{K/\mathbf Q}(\det N)|$ in $\sum_i\mathcal O_Ka_i$: in the basis $a$, the inclusion is $N$ acting on
$\mathcal O_K^k$, a free $\mathbf Z$-module of rank $k[K:\mathbf Q]$, on which its determinant is
$N_{K/\mathbf Q}(\det N)$. Second, $\sum_\sigma\log|\sigma(a)|^2=2\log|N_{K/\mathbf Q}(a)|$ for $a\in K^\times$,
the sum running over all embeddings.

Put $B=\sum_i\mathcal O_Kb_i$ and $B'=\sum_i\mathcal O_Kb'_i$. As $L_1$ is finitely generated and lies in
$KB$, there is $D\ne0$ in $\mathcal O_K$ with $DL_1\subset B$. Write $Db'_j=\sum_iM_{ij}b_i$ with $M\in
M_k(\mathcal O_K)$; then $\det M\ne0$, as $Db'$ is a basis of $KL_1$. Let $C=\sum_j\mathcal O_KDb'_j$. By the
first fact, $[B:C]=|N_{K/\mathbf Q}(\det M)|$ and $[B':C]=|N_{K/\mathbf Q}(D)|^k$. For every $\sigma$,
$G_\sigma(Db')=\sigma(M)^*G_\sigma(b)\sigma(M)=|\sigma(D)|^2G_\sigma(b')$, so
\[
 \det G_\sigma(Db')=|\sigma(\det M)|^2\det G_\sigma(b)=|\sigma(D)|^{2k}\det G_\sigma(b').
\]
By the second fact and $[L_1:C]=[L_1:B][B:C]$,
$\deg(L_1,b)=\log[L_1:C]-\frac12\sum_\sigma\log\det G_\sigma(Db')$. In the same way, with $[L_1:C]=[L_1:B'][B':C]$,
$\deg(L_1,b')$ is the same number.
\end{proof}

\begin{lemma}\label{is-onto-semistable-det}
Let $A,C\in M_n(\mathbf C)$ with $A$ positive definite and $C$ Hermitian.
\begin{enumerate}
\item If $|v^*Cv|\le v^*Av$ for all $v\in\mathbf C^n$, then $|\det C|\le\det A$.
\item If $v^*Av\le v^*Cv$ for all $v\in\mathbf C^n$, then $0<\det A\le\det C$.
\end{enumerate}
\end{lemma}
\leanref{ReciprocityTheory/SemistableReduction/SemiDegDA.lean}{ReciprocityTheory.MultiplicativeHeights.sdd_absDet}
\leanref{ReciprocityTheory/SemistableReduction/SemiDegDA.lean}{ReciprocityTheory.MultiplicativeHeights.sdd_det_mono}
\leanref{ReciprocityTheory/SemistableReduction/SemiDegD.lean}{ReciprocityTheory.MultiplicativeHeights.absDet_holds}
\begin{proof}
Write $A=Q^*Q$ with $Q$ invertible, and put $D=(Q^{-1})^*CQ^{-1}$. It is Hermitian, $\det C=|\det Q|^2\det D$
and $\det A=|\det Q|^2>0$. Let $w$ be a unit eigenvector of $D$ with eigenvalue $\lambda$, and $v=Q^{-1}w$.
Then $v^*Av=1$ and $v^*Cv=\lambda$. Under (1) every eigenvalue of $D$ has $|\lambda|\le1$, so $|\det
D|\le1$; under (2) every eigenvalue is at least $1$, so $\det D\geq1$.
\end{proof}

\begin{lemma}\label{is-onto-semistable-volume}
Let $K$ be a number field with an automorphism $a\mapsto a^c$ of order at most $2$ such that
$\sigma(a^c)=\overline{\sigma(a)}$ for every embedding $\sigma\colon K\to\mathbf C$, let $B\in M_r(K)$ with
$B_{ab}^c=B_{ba}$, and put $S(x,y)=\sum_{a,b}x_a^cB_{ab}y_b$. Let $L_1\subset K^r$ be a finitely generated
$\mathcal O_K$-submodule, $b_1,\dots,b_k\in L_1$ a basis of $KL_1$, and $\Gamma=(S(b_i,b_j))_{i,j}$.
\begin{enumerate}
\item If $S(x,y)$ is an algebraic integer for $x,y\in L_1$ and $\det\Gamma\ne0$, then
  $[L_1:\sum_i\mathcal O_Kb_i]^2\le|N_{K/\mathbf Q}(\det\Gamma)|$.
\item If $L_1$ spans $K^r$ and is self-dual for $S$, that is, $S(x,y)$ is an algebraic integer for
  $x,y\in L_1$ and every $x\in K^r$ with $S(x,L_1)$ integral lies in $L_1$, then
  $|N_{K/\mathbf Q}(\det\Gamma)|\le[L_1:\sum_i\mathcal O_Kb_i]^2$.
\end{enumerate}
\end{lemma}
\leanref{ReciprocityTheory/SemistableReduction/SemiDegC.lean}{ReciprocityTheory.MultiplicativeHeights.integralIndex_holds}
\leanref{ReciprocityTheory/SemistableReduction/SemiDegC.lean}{ReciprocityTheory.MultiplicativeHeights.selfDualIndex_holds}
\leanref{ReciprocityTheory/SemistableReduction/SemiDegCE.lean}{ReciprocityTheory.MultiplicativeHeights.SemiDegIdx.split_norm}
\leanref{ReciprocityTheory/SemistableReduction/SemiDegCB.lean}{ReciprocityTheory.MultiplicativeHeights.SemiDegIdx.lattice_split}
\begin{proof}
Identify $K^k$ with $KL_1$ by $t\mapsto\sum_it_ib_i$. This carries $\mathcal O_K^k$ onto
$\sum_i\mathcal O_Kb_i$, $L_1$ onto a finitely generated submodule $X\supset\mathcal O_K^k$, and $S$ to
$S'(s,t)=\sum_{i,j}s_i^c\Gamma_{ij}t_j$; so $[L_1:\sum_i\mathcal O_Kb_i]=[X:\mathcal O_K^k]$. By the Hermite normal
form over the Dedekind domain $\mathcal O_K$ \cite[\S1.4]{Coh00}, there are finitely generated fractional
ideals $C_1,\dots,C_k\supset\mathcal O_K$ and vectors $z_1,\dots,z_k\in K^k$ with $(z_i)_l=0$ for $l<i$ and
$(z_i)_i=1$, such that $X=\bigoplus_iC_iz_i$; projecting away one coordinate at a time gives $[X:\mathcal
O_K^k]=\prod_i[C_i:\mathcal O_K]$. Let $Z$ be the matrix with rows $z_i$, so that $\det Z=1$, put $C=\prod_iC_i$
and $C^c=\{a^c:a\in C\}$, and let $J=C^c\cdot\det\Gamma\cdot C$. The norm of $C_i$ is $[C_i:\mathcal O_K]^{-1}$,
and $a\mapsto a^c$ preserves norms, so $N(J)\,[X:\mathcal O_K^k]^2=|N_{K/\mathbf Q}(\det\Gamma)|$.

(1) Let $s_i^c\in C_i$ and $t_i\in C_i$ for every $i$. The matrix $\big(S'(s_i^cz_i,t_jz_j)\big)_{i,j}$ has
integral entries, as $s_i^cz_i$ and $t_jz_j$ lie in $X$, and it equals
$\operatorname{diag}(s)\,Z^c\Gamma Z^t\operatorname{diag}(t)$. So its determinant $\prod_is_i\cdot\det\Gamma\cdot
\prod_it_i$ is integral. These products generate $J$, so $J\subset\mathcal O_K$, and $N(J)\geq1$ as $J\ne0$.

(2) If $\det\Gamma=0$ there is nothing to prove. Otherwise, $X$ is self-dual for $S'$. Let $u_j\in K$ be such
that $u_jw$ is integral whenever $w^c\in C_j$. Put $Q=Z^c\Gamma$, of determinant $\det\Gamma$, and
$y_j=Q^{-1}(u_je_j)$. For $t_i\in C_i$, $S'(\sum_it_iz_i,y_j)=t_j^cu_j$ is integral, and so is its conjugate
$S'(y_j,\sum_it_iz_i)$; hence $y_j\in X$, and $y_j=\sum_iT_{ij}z_i$ with $T_{ij}\in C_i$. The matrix with columns
$y_j$ is $Z^tT=Q^{-1}\operatorname{diag}(u)$, so $\prod_ju_j=\det\Gamma\cdot\det T\in\det\Gamma\cdot C$. Hence
$(C^c)^{-1}\subset\det\Gamma\cdot C$, that is, $\mathcal O_K\subset J$. With the proof of (1), $J=\mathcal O_K$,
and $|N_{K/\mathbf Q}(\det\Gamma)|=[X:\mathcal O_K^k]^2$.
\end{proof}

\begin{lemma}[Witt decomposition \cite{Sch85}]\label{is-onto-semistable-witt}
Let $K$ be a number field with an automorphism $a\mapsto a^c$ of order at most $2$ such that
$\sigma(a^c)=\overline{\sigma(a)}$ for every embedding $\sigma\colon K\to\mathbf C$, and let $B\in M_r(K)$ with
$B_{ab}^c=B_{ba}$ and $\det B\ne0$. Then there are $P\in GL_r(K)$, an integer $m$ with $2m\le r$, and elements
$g_i\in K$ for $m<i\le r-m$, such that $(P^c)^tBP=W$, where $W_{a,r+1-a}=1$ for $a\le m$ and for $a>r-m$,
$W_{aa}=g_a$ for $m<a\le r-m$, and all other entries of $W$ vanish; and such that
$\sum_{m<i\le r-m}t_i^cg_it_i\ne0$ for every nonzero family $(t_i)_{m<i\le r-m}$ in $K$. Moreover
$g_i^c=g_i\ne0$ for $m<i\le r-m$.
\end{lemma}
\leanref{ReciprocityTheory/SemistableReduction/SemiDegA.lean}{ReciprocityTheory.MultiplicativeHeights.wittBasis_holds}
\leanref{ReciprocityTheory/SemistableReduction/SemiDegSt.lean}{ReciprocityTheory.MultiplicativeHeights.WittBasisStatement}
\leanref{ReciprocityTheory/SemistableReduction/SemiDegG.lean}{ReciprocityTheory.MultiplicativeHeights.SemiPt.g_fix}
\leanref{ReciprocityTheory/SemistableReduction/SemiDegZA.lean}{ReciprocityTheory.MultiplicativeHeights.SemiPt.g_ne_zero}

\begin{lemma}\label{is-onto-semistable-graded}
Let $K$ be a number field and $r\geq0$ an integer, and let $\delta_{\sigma,i}>0$ be real numbers, for the
embeddings $\sigma\colon K\to\mathbf C$ and $1\le i\le r$. Write $\deg_\delta(B_1)$ for the degree of
Lemma~\ref{is-onto-share-degree} of a finitely generated $\mathcal O_K$-submodule $B_1\subset K^r$, with
respect to the matrices $\operatorname{diag}(\delta_{\sigma,1},\dots,\delta_{\sigma,r})$ and any family in
$B_1$ that is a basis of $KB_1$; it does not depend on the family by
Lemma~\ref{is-onto-semistable-basis}, and $\deg_\delta(0)=0$. For $0\le e\le r$ let $V_e$ be the set of
$x\in K^r$ with $x_i=0$ for $i>e$, and $\pi_e\colon K^r\to K^r$ the map that keeps the coordinates $x_i$ with
$i>e$ and replaces the others by $0$.
\begin{enumerate}
\item For every finitely generated submodule $B\subset K^r$ and $0\le e\le r$,
  $\deg_\delta(B)\le\deg_\delta(B\cap V_e)+\deg_\delta(\pi_eB)$.
\item Let $L'\subset K^r$ be a finitely generated submodule and $0\le e'\le e\le r$. If all submodules of
  $L'\cap V_{e'}$ and all submodules of $\pi_{e'}(L'\cap V_e)$ have $\deg_\delta\le0$, then so do all
  submodules of $L'\cap V_e$.
\end{enumerate}
\end{lemma}
\leanref{ReciprocityTheory/SemistableReduction/SemiDegD.lean}{ReciprocityTheory.MultiplicativeHeights.gradedExt_holds}
\leanref{ReciprocityTheory/SemistableReduction/SemiDegD.lean}{ReciprocityTheory.MultiplicativeHeights.sdd_gram_step}
\leanref{ReciprocityTheory/SemistableReduction/SemiDegDB.lean}{ReciprocityTheory.MultiplicativeHeights.sdd_index}
\leanref{ReciprocityTheory/SemistableReduction/SemiDegW.lean}{ReciprocityTheory.MultiplicativeHeights.SemiPt.flag_step}
\begin{proof}
(1) Let $b_0$ be a family in $B\cap V_e$ that is a basis of $K(B\cap V_e)$, $b_1$ a family in $\pi_eB$ that is a
basis of $K\pi_eB$, and $c_j\in B$ with $\pi_ec_j=b_{1,j}$. The family $(b_0,c)$ is a basis of $KB$ in $B$.
Indeed, applying $\pi_e$ to a linear relation kills $b_0$, so the coefficients of $c$ vanish, and then those of
$b_0$. If $x\in B$, then $\pi_ex=\sum_j\lambda_jb_{1,j}$, and $x-\sum_j\lambda_jc_j$ lies in $V_e\cap KB$, whose
elements have nonzero multiples in $B\cap V_e$. Moreover $B\cap Kb_0=B\cap V_e$. By
Lemmas~\ref{is-onto-semistable-basis} and~\ref{is-onto-share-degree},
\[
 \deg_\delta(B)=\deg_\delta(B\cap V_e)+\deg(B;b_0,c),
\]
with the quotient degree of Lemma~\ref{is-onto-share-degree} for $D_\sigma=\operatorname{diag}(\delta_{\sigma,1},
\dots,\delta_{\sigma,r})$. It remains to show $\deg(B;b_0,c)\le\deg_\delta(\pi_eB)$, computed with $b_1$.

\emph{Indices.} $[B:(B\cap V_e)+\sum_j\mathcal O_Kc_j]=[\pi_eB:\sum_j\mathcal O_Kb_{1,j}]$, since $\pi_e$ maps $B$
onto $\pi_eB$ and $(B\cap V_e)+\sum_j\mathcal O_Kc_j$ onto $\sum_j\mathcal O_Kb_{1,j}$, and its kernel on $B$ is
$B\cap V_e$.

\emph{Gram determinants.} Fix $\sigma$, let $M_0$, $M_c$ and $E_1$ be the matrices with columns
$\sigma(b_{0,i})$, $\sigma(c_j)$ and $\sigma(b_{1,j})$, and $E=M_c-M_0G_\sigma(b_0)^{-1}G_\sigma(b_0,c)$. A direct
computation gives $Q_\sigma=E^*D_\sigma E$ for the matrix $Q_\sigma$ of Lemma~\ref{is-onto-share-degree}. The
rows of $M_0$ with index $i>e$ vanish, so the rows of $E$ and $E_1$ with index $i>e$ agree, while those of $E_1$
with $i\le e$ vanish. Hence, for $v\in\mathbf C^{k_1}$, where $k_1$ is the length of $b_1$,
\[
 v^*G_\sigma(b_1)v=\sum_{i>e}\delta_{\sigma,i}|(Ev)_i|^2\le\sum_i\delta_{\sigma,i}|(Ev)_i|^2=v^*Q_\sigma v,
\]
and Lemma~\ref{is-onto-semistable-det}(2) gives $0<\det G_\sigma(b_1)\le\det Q_\sigma$. With the indices, this
gives $\deg(B;b_0,c)\le\deg_\delta(\pi_eB)$.

(2) Let $B\subset L'\cap V_e$ be a submodule; it is finitely generated. Then $B\cap V_{e'}\subset L'\cap V_{e'}$
and $\pi_{e'}B\subset\pi_{e'}(L'\cap V_e)$, so (1) gives $\deg_\delta(B)\le0$.
\end{proof}

\begin{lemma}\label{is-onto-semistable-line}
Let $K$ be a number field, $\delta_{\sigma,i}>0$ and $\deg_\delta$ as in
Lemma~\ref{is-onto-semistable-graded}, and $n=[K:\mathbf Q]$. Let $1\le i\le r$, $e_i$ the $i$-th standard
basis vector, and $G\subset Ke_i$ a finitely generated $\mathcal O_K$-submodule that contains $x=ae_i$ with
$a\ne0$.
\begin{enumerate}
\item $\deg_\delta(G)=\log[G:\mathcal O_Kx]-\log|N_{K/\mathbf Q}(a)|-\frac12\sum_\sigma\log\delta_{\sigma,i}$; if
  $\delta_{\sigma,i}=d$ for every $\sigma$, then $\frac12\sum_\sigma\log\delta_{\sigma,i}=\frac n2\log d$.
\item If $\deg_\delta(G)\le0$, then $\deg_\delta(G_1)\le0$ for every submodule $G_1\subset G$.
\end{enumerate}
\end{lemma}
\leanref{ReciprocityTheory/SemistableReduction/SemiDegR.lean}{ReciprocityTheory.MultiplicativeHeights.SemiPt.latDeg_single}
\leanref{ReciprocityTheory/SemistableReduction/SemiDegR.lean}{ReciprocityTheory.MultiplicativeHeights.SemiPt.sum_log_single}
\leanref{ReciprocityTheory/SemistableReduction/SemiDegR.lean}{ReciprocityTheory.MultiplicativeHeights.SemiPt.rankOne_le_zero}
\leanref{ReciprocityTheory/SemistableReduction/SemiDegE.lean}{ReciprocityTheory.MultiplicativeHeights.SemiPt.latDeg_le_of_le}
\begin{proof}
(1) The Gram matrix of $x$ at $\sigma$ is $|\sigma(a)|^2\delta_{\sigma,i}$, and
$\sum_\sigma\log|\sigma(a)|^2=2\log|N_{K/\mathbf Q}(a)|$.

(2) A submodule $G_1$ has rank at most $1$ over $K$. If $G_1=0$, its degree is $0$. Otherwise it has a basis
$y$ of $KG_1=KG$, and $[G:\mathcal O_Ky]=[G:G_1][G_1:\mathcal O_Ky]\geq[G_1:\mathcal O_Ky]$, so
$\deg_\delta(G_1)\le\deg_\delta(G)\le0$, both computed with $y$.
\end{proof}

\begin{lemma}\label{is-onto-semistable-pair}
Let $K$ be a number field with an automorphism $a\mapsto a^c$ of order at most $2$ such that
$\sigma(a^c)=\overline{\sigma(a)}$ for every embedding $\sigma$, and let $\delta_{\sigma,i}>0$ and
$\deg_\delta$ be as in Lemma~\ref{is-onto-semistable-graded}. Let $1\le j<p\le r$ and $d>0$ with
$\delta_{\sigma,j}=d$ and $\delta_{\sigma,p}=d^{-1}$ for every $\sigma$. Let $G_j\subset Ke_j$ and $G_p\subset
Ke_p$ be finitely generated $\mathcal O_K$-submodules, containing $x$ with $x_j\ne0$ and $y$ with $y_p\ne0$
respectively, such that $u_j^cv_p$ is an algebraic integer for all $u\in G_j$ and $v\in G_p$. If
$\deg_\delta(G_j)=0$, then $\deg_\delta(G_p)\le0$.
\end{lemma}
\leanref{ReciprocityTheory/SemistableReduction/SemiDegI.lean}{ReciprocityTheory.MultiplicativeHeights.SemiPt.top_le_zero}
\leanref{ReciprocityTheory/SemistableReduction/SemiDegH.lean}{ReciprocityTheory.MultiplicativeHeights.SemiPt.spanIdx_pair}
\begin{proof}
Let $n=[K:\mathbf Q]$, and put $A=[G_j:\mathcal O_Kx]$, $B=[G_p:\mathcal O_Ky]$, $X=|N_{K/\mathbf Q}(x_j)|$ and
$Y=|N_{K/\mathbf Q}(y_p)|$. By
Lemma~\ref{is-onto-semistable-line}(1), $0=\deg_\delta(G_j)=\log A-\log X-\frac n2\log d$ and
$\deg_\delta(G_p)=\log B-\log Y+\frac n2\log d$, so it suffices to show $AB\le XY$.

Let $T=G_j+G_p$, with the basis $x,y$ of $KT$, and $S_0(u,v)=u_j^cv_p+u_p^cv_j$, the form of
Lemma~\ref{is-onto-semistable-volume} for the matrix with entries $1$ at $(j,p)$ and $(p,j)$ and $0$ elsewhere.
It is integral on $T$: for $u=u'+u''$ and $v=v'+v''$ with $u',v'\in G_j$ and $u'',v''\in G_p$,
$S_0(u,v)=u_j'^cv''_p+(v_j'^cu''_p)^c$. Its Gram matrix on $x,y$ has the entries $x_j^cy_p$ and $y_p^cx_j$ off
the diagonal and $0$ on it, with determinant of absolute norm $X^2Y^2\ne0$. As the coordinates $j$ and $p$
separate $G_j$ from $G_p$, $[T:\mathcal O_Kx+\mathcal O_Ky]=AB$. Lemma~\ref{is-onto-semistable-volume}(1) gives
$(AB)^2\le X^2Y^2$.
\end{proof}

\begin{lemma}\label{is-onto-semistable-middle}
Let $K$ be a number field with an automorphism $a\mapsto a^c$ of order at most $2$ such that
$\sigma(a^c)=\overline{\sigma(a)}$ for every embedding $\sigma$, and let $\delta_{\sigma,i}>0$ and
$\deg_\delta$ be as in Lemma~\ref{is-onto-semistable-graded}. Let $m\geq0$ with $2m\le r$, $I=\{i:m<i\le
r-m\}$, and $g_i\in K$ for $i\in I$ with $g_i^c=g_i$ and $\sum_{i\in I}t_i^cg_it_i\ne0$ for every nonzero
$(t_i)_{i\in I}$, and suppose $\delta_{\sigma,i}=|\sigma(g_i)|$ for $i\in I$ and every $\sigma$. Let $G\subset
K^r$ be a finitely generated $\mathcal O_K$-submodule with $x_i=0$ for $x\in G$ and $i\notin I$, such that
$\sum_{i\in I}x_i^cg_iy_i$ is an algebraic integer for $x,y\in G$. Then $\deg_\delta(G_1)\le0$ for every
submodule $G_1\subset G$.
\end{lemma}
\leanref{ReciprocityTheory/SemistableReduction/SemiDegN.lean}{ReciprocityTheory.MultiplicativeHeights.SemiPt.mid_le_zero}
\begin{proof}
Let $c_1,\dots,c_k\in G_1$ be a basis of $KG_1$ and $\Gamma$ the Gram matrix of the $c_a$ for the form
$S_{\mathrm{mid}}(x,y)=\sum_{i\in I}x_i^cg_iy_i$. If $\Gamma t=0$, then $z=\sum_at_ac_a$ has
$S_{\mathrm{mid}}(z,z)=0$, so $z=0$ and $t=0$; hence $\det\Gamma\ne0$. Since $S_{\mathrm{mid}}$ is integral on
$G_1$, Lemma~\ref{is-onto-semistable-volume}(1) gives $2\log[G_1:\sum_a\mathcal O_Kc_a]\le\log|N_{K/\mathbf
Q}(\det\Gamma)|=\sum_\sigma\log|\sigma(\det\Gamma)|$.

Fix $\sigma$, let $E$ be the matrix with columns $\sigma(c_a)$ and $D_\sigma=\operatorname{diag}(\delta_{\sigma,i})$.
The matrix $\sigma(\Gamma)$ is Hermitian, and for $v\in\mathbf C^k$, with $w=Ev$,
\[
 |v^*\sigma(\Gamma)v|=\Big|\sum_{i\in I}\sigma(g_i)|w_i|^2\Big|\le\sum_i\delta_{\sigma,i}|w_i|^2=v^*E^*D_\sigma Ev.
\]
By Lemma~\ref{is-onto-semistable-det}(1), $|\sigma(\det\Gamma)|\le\det G_\sigma(c)$. Summing the logarithms,
$2\log[G_1:\sum_a\mathcal O_Kc_a]\le\sum_\sigma\log\det G_\sigma(c)$, that is, $\deg_\delta(G_1)\le0$.
\end{proof}

\begin{lemma}\label{is-onto-semistable-majorant}
Let $K$ be a number field with an automorphism $a\mapsto a^c$ of order at most $2$ such that
$\sigma(a^c)=\overline{\sigma(a)}$ for every embedding $\sigma$, and $B\in M_r(K)$ with $B_{ab}^c=B_{ba}$. Let
$P\in GL_r(K)$, $m\geq0$ with $2m\le r$, and $g_i\in K$ with $g_i^c=g_i\ne0$ for $m<i\le r-m$, such that
$(P^c)^tBP=W$, where $W$ has the entries $1$ at $(a,r+1-a)$ for $a\le m$ and for $a>r-m$, $g_a$ at $(a,a)$ for
$m<a\le r-m$, and $0$ elsewhere; put $Q=P^{-1}$. Let $d_1,\dots,d_m>0$, and for every embedding $\sigma$ put $\delta_{\sigma,i}=d_i$ for $i\le m$,
$\delta_{\sigma,i}=d_{r+1-i}^{-1}$ for $i>r-m$ and $\delta_{\sigma,i}=|\sigma(g_i)|$ otherwise,
$D_\sigma=\operatorname{diag}(\delta_{\sigma,1},\dots,\delta_{\sigma,r})$ and
$H_\sigma=(Q^\sigma)^*D_\sigma Q^\sigma$. Then $H_\sigma$ is positive definite Hermitian,
$H_\sigma(B^\sigma)^{-1}H_\sigma=B^\sigma$, $H_{\bar\sigma}=\overline{H_\sigma}$, and for every finitely generated
submodule $L_1\subset K^r$ and every family $b$ in $L_1$ that is a basis of $KL_1$, the degree of $L_1$ for
the $H_\sigma$ computed with $b$ equals $\deg_\delta(QL_1)$, in the notation of
Lemma~\ref{is-onto-semistable-graded}.
\end{lemma}
\leanref{ReciprocityTheory/SemistableReduction/SemiDegG.lean}{ReciprocityTheory.MultiplicativeHeights.SemiPt.exists_adapted_majorant}
\leanref{ReciprocityTheory/SemistableReduction/SemiDegF.lean}{ReciprocityTheory.MultiplicativeHeights.SemiPt.wVal_adMet}
\leanref{ReciprocityTheory/SemistableReduction/SemiDegT.lean}{ReciprocityTheory.MultiplicativeHeights.SemiPt.latDeg_transport}
\begin{proof}
From $(P^c)^tBP=W$ and $\sigma(a^c)=\overline{\sigma(a)}$, $B^\sigma=(Q^\sigma)^*W^\sigma Q^\sigma$. The matrix
$W^\sigma D_\sigma^{-1}W^\sigma$ is $D_\sigma$. Indeed, the nonzero entries of $W$ are at the positions
$(a,a^*)$, where $a^*=r+1-a$ for $a\le m$ or $a>r-m$ and $a^*=a$ otherwise; as $a\mapsto a^*$ is an
involution, $W^\sigma D_\sigma^{-1}W^\sigma$ is diagonal with entries
$\sigma(W_{aa^*})\delta_{\sigma,a^*}^{-1}\sigma(W_{a^*a})$. For $a\le m$ this entry is
$\delta_{\sigma,r+1-a}^{-1}=d_a$, for $a>r-m$ it is $\delta_{\sigma,r+1-a}^{-1}=d_{r+1-a}^{-1}$, and otherwise
it is $\sigma(g_a)^2/|\sigma(g_a)|=|\sigma(g_a)|$, as $\sigma(g_a)$ is real because $g_a^c=g_a$. Hence
\[
 B^\sigma H_\sigma^{-1}B^\sigma=(Q^\sigma)^*W^\sigma D_\sigma^{-1}W^\sigma Q^\sigma=H_\sigma,
\]
which is equivalent to $H_\sigma(B^\sigma)^{-1}H_\sigma=B^\sigma$. The matrix $H_\sigma$ is positive definite, as
$D_\sigma$ is and $Q^\sigma$ is invertible. Since $|\overline{\sigma(g_i)}|=|\sigma(g_i)|$,
$\delta_{\bar\sigma,i}=\delta_{\sigma,i}$, and $Q^{\bar\sigma}=\overline{Q^\sigma}$; so $H_{\bar\sigma}=\overline{H_\sigma}$.
Finally $\sigma(b_i)^*H_\sigma\sigma(b_j)=\sigma(Qb_i)^*D_\sigma\sigma(Qb_j)$, and $Q$ is injective, so
$[L_1:\sum_i\mathcal O_Kb_i]=[QL_1:\sum_i\mathcal O_KQb_i]$.
\end{proof}

\begin{lemma}\label{is-onto-share-semistable}
Let $K$ be a number field with an automorphism $a\mapsto a^c$ of order at most $2$ such that
$\sigma(a^c)=\overline{\sigma(a)}$ for every embedding $\sigma\colon K\to\mathbf C$, let $B\in M_r(K)$ with
$B_{ab}^c=B_{ba}$, and put $S(x,y)=\sum_{a,b}x_a^cB_{ab}y_b$. Let $L\subset K^r$ be a finitely generated
$\mathcal O_K$-submodule that spans $K^r$ and is self-dual for $S$: $S(x,y)$ is an algebraic integer for
$x,y\in L$, and every $x\in K^r$ with $S(x,L)$ integral lies in $L$. Then there are positive definite
Hermitian matrices $H_\sigma$, one for each embedding $\sigma$, with $H_\sigma(B^\sigma)^{-1}H_\sigma=B^\sigma$
and $H_{\bar\sigma}=\overline{H_\sigma}$, such that:
\begin{enumerate}
\item $L$ is semistable of degree $0$: for an $\mathcal O_K$-submodule $L_1\subset K^r$ and elements
  $b_1,\dots,b_k\in L_1$ that form a basis of $KL_1$, put
  \[
   \deg(L_1)=\log\Big[L_1:\textstyle\sum_i\mathcal O_Kb_i\Big]-\frac12\sum_\sigma\log\det\big(\sigma(b_i)^*H_\sigma\sigma(b_j)\big)_{i,j};
  \]
  then $\deg(L)=0$ for every such basis of $L$, and $\deg(L_1)\le0$ for every submodule $L_1\subset L$
  and every such basis of $L_1$;
\item every $N\in M_r(K)$ with $NL\subset L$ and $(N^\sigma u)^*H_\sigma(N^\sigma u)<u^*H_\sigma u$ for every
  $\sigma$ and every $u\ne0$ is zero.
\end{enumerate}
\end{lemma}
\leanref{ReciprocityTheory/SemistableReduction/SemiAll.lean}{ReciprocityTheory.MultiplicativeHeights.Closed.semistable}
\leanref{ReciprocityTheory/SemistableReduction/LatSemiSt.lean}{ReciprocityTheory.MultiplicativeHeights.SemistableStatement}
\leanref{ReciprocityTheory/SemistableReduction/SemiContract.lean}{ReciprocityTheory.MultiplicativeHeights.semiContract_holds}
\leanref{ReciprocityTheory/SemistableReduction/SemiContractB.lean}{ReciprocityTheory.MultiplicativeHeights.sc_gram_lt}
\leanref{ReciprocityTheory/SemistableReduction/SemiAll.lean}{ReciprocityTheory.MultiplicativeHeights.Closed.semiPoint}
\leanref{ReciprocityTheory/SemistableReduction/SemiDegZA.lean}{ReciprocityTheory.MultiplicativeHeights.SemiPt.semiPoint_of}
\leanref{ReciprocityTheory/SemistableReduction/SemiDegW.lean}{ReciprocityTheory.MultiplicativeHeights.SemiPt.adapted_semistable}
\leanref{ReciprocityTheory/SemistableReduction/SemiDegZA.lean}{ReciprocityTheory.MultiplicativeHeights.SemiPt.latDeg_L_nonneg}
\leanref{ReciprocityTheory/SemistableReduction/SemiDegZ.lean}{ReciprocityTheory.MultiplicativeHeights.SemiPt.detG_ne_zero}
\begin{proof}
Let $n=[K:\mathbf Q]$, and write $\deg_\delta$, $V_e$ and $\pi_e$ as in Lemma~\ref{is-onto-semistable-graded}.

\emph{The form is nondegenerate.} Suppose $\det B=0$, and let $v\ne0$ with $B^tv=0$. For $w=(v_a^c)_a$,
$S(w,y)=\sum_{a,b}v_aB_{ab}y_b=0$ for all $y$, so $S(tw,L)=0$ for all $t\in K$, and $Kw\subset L$ by
self-duality. But the coordinates of the elements of the finitely generated module $L$ lie in a fractional
ideal of $K$, which is not $K$. So $\det B\ne0$.

\emph{Witt coordinates.} Let $P$, $m$, $g_i$ and $W$ be as in Lemma~\ref{is-onto-semistable-witt}, $Q=P^{-1}$
and $L'=QL$, a finitely generated submodule that spans $K^r$. For $x,y\in K^r$,
\[
 S(Px,Py)=\sum_{a\le m\text{ or }a>r-m}x_a^cy_{r+1-a}+\sum_{m<a\le r-m}x_a^cg_ay_a,
\]
which is integral for $x,y\in L'$. For $1\le j\le r$ let $G_j=\pi_{j-1}(L'\cap V_j)$, the set of the
$x_je_j$ for $x\in L'\cap V_j$; it is a finitely generated submodule of $Ke_j$. As $L'$ spans $K^r$, it contains
$a_je_j$ for some nonzero $a_j\in\mathcal O_K$, which lies in $G_j$. For $j\le m$ put
\[
 d_j=\exp\Big(\frac2n\big(\log[G_j:\mathcal O_Ka_je_j]-\log|N_{K/\mathbf Q}(a_j)|\big)\Big)>0,
\]
and let $\delta$ and $H_\sigma$ be as in Lemma~\ref{is-onto-semistable-majorant} for these $d_j$. By that lemma,
$H_\sigma$ is positive definite Hermitian with $H_\sigma(B^\sigma)^{-1}H_\sigma=B^\sigma$ and
$H_{\bar\sigma}=\overline{H_\sigma}$.

\emph{All submodules of $L'$ have $\deg_\delta\le0$.} Apply Lemma~\ref{is-onto-semistable-graded}(2) along the
chain $V_0\subset V_1\subset\dots\subset V_m\subset V_{r-m}\subset V_{r-m+1}\subset\dots\subset V_r=K^r$,
starting from $L'\cap V_0=0$. Three inputs are needed.
\begin{itemize}
\item From $V_{j-1}$ to $V_j$, for $j\le m$, the module is $G_j$. Here $\delta_{\sigma,j}=d_j$ for every $\sigma$,
  so Lemma~\ref{is-onto-semistable-line}(1) and the choice of $d_j$ give $\deg_\delta(G_j)=0$, and
  Lemma~\ref{is-onto-semistable-line}(2) gives the bound for its submodules.
\item From $V_m$ to $V_{r-m}$, the module $\pi_m(L'\cap V_{r-m})$ consists of vectors supported on
  $m<i\le r-m$. For $x,y\in L'\cap V_{r-m}$ the first sum in $S(Px,Py)$ vanishes: for $a\le m$,
  $y_{r+1-a}=0$, and for $a>r-m$, $x_a=0$. So $\sum_{m<i\le r-m}x_i^cg_iy_i$ is integral on this module, and
  Lemma~\ref{is-onto-semistable-middle} applies, as $\delta_{\sigma,i}=|\sigma(g_i)|$ for $m<i\le r-m$.
\item From $V_{p-1}$ to $V_p$, for $p>r-m$, the module is $G_p$. Let $j=r+1-p$, so that $j\le m<p$. For
  $x\in L'\cap V_j$ and $y\in L'\cap V_p$, $S(Px,Py)=x_j^cy_p$: the terms with $a<j$ have $y_{r+1-a}=0$, as
  $r+1-a>p$, and those with $a>j$ have $x_a=0$. So $u_j^cv_p$ is integral for $u\in G_j$ and $v\in G_p$. As
  $\delta_{\sigma,j}=d_j$, $\delta_{\sigma,p}=d_j^{-1}$ and $\deg_\delta(G_j)=0$,
  Lemma~\ref{is-onto-semistable-pair} gives $\deg_\delta(G_p)\le0$, and Lemma~\ref{is-onto-semistable-line}(2)
  the bound for its submodules.
\end{itemize}

\emph{The bound $\deg\le0$.} Let $L_1\subset L$ be a submodule and $b$ a family in $L_1$ that is a basis of
$KL_1$. By Lemma~\ref{is-onto-semistable-majorant}, $\deg(L_1)=\deg_\delta(QL_1)$, which is at most $0$ since
$QL_1\subset L'$.

\emph{The degree of $L$.} Let $b_1,\dots,b_r\in L$ be a basis of $K^r$, $M$ the matrix with columns $b_i$, and
$\Gamma=(S(b_i,b_j))_{i,j}=(M^c)^tBM$. Since $H_\sigma(B^\sigma)^{-1}H_\sigma=B^\sigma$, $(\det H_\sigma)^2=(\det
B^\sigma)^2$, and $\det H_\sigma>0$; so $\det H_\sigma=|\sigma(\det B)|$. Then
$\det\big(\sigma(b_i)^*H_\sigma\sigma(b_j)\big)=|\sigma(\det M)|^2\det H_\sigma=|\sigma(\det\Gamma)|$, and
\[
 \deg(L)=\log\Big[L:\textstyle\sum_i\mathcal O_Kb_i\Big]-\frac12\log|N_{K/\mathbf Q}(\det\Gamma)|\geq0
\]
by Lemma~\ref{is-onto-semistable-volume}(2). With the bound $\deg(L)\le0$, $\deg(L)=0$. This proves (1).

We deduce (2) from (1). Write $\deg$ as in Lemma~\ref{is-onto-share-degree} for the matrices $H_\sigma$.
Let $N\in M_r(K)$ satisfy $NL\subset L$ and $(N^\sigma u)^*H_\sigma N^\sigma u<u^*H_\sigma u$ for all
$\sigma$ and $u\ne0$, and suppose $N\ne0$. Let $k$ and $m$ be the dimensions of the kernel and the image
of $N$ on $K^r$, so that $k+m=r$ and $m>0$. As $L$ spans $K^r$, every vector has a nonzero multiple in
$L$. So there are $b_1,\dots,b_k\in L$ that form a basis of $\ker N$, and $c_1,\dots,c_m\in L$ such that
$Nc_1,\dots,Nc_m$ form a basis of the image. Applying $N$ to a linear relation among the $b_i$ and $c_j$
kills the $b_i$, so the coefficients of the $c_j$ vanish, and then those of the $b_i$. So $b,c$ is a
basis of $K^r$ contained in $L$. Put $L_b=L\cap\ker N$ and $C=\sum_j\mathcal O_Kc_j$.

\emph{Degrees.} By (1), $\deg(L,(b,c))=0$, $\deg(L_b,b)\le0$, and $\deg(NL,Nc)\le0$; here $NL\subset L$
has the basis $Nc_1,\dots,Nc_m$ of its span.

\emph{Indices.} $[NL:\sum_j\mathcal O_KNc_j]=[L:L_b+C]$. Indeed $N$ maps $L$ onto $NL$ and $L_b+C$ onto
$\sum_j\mathcal O_KNc_j$, and if $x\in L$ has $Nx=Nz$ with $z\in C$, then $x-z\in L\cap\ker N=L_b$.

\emph{Gram determinants.} For every $\sigma$, $0<\det G_\sigma(Nc)<\det Q_\sigma$, with $Q_\sigma$ as in
Lemma~\ref{is-onto-share-degree}. Let $B_\sigma$ and $C_\sigma$ be the matrices with columns
$\sigma(b_i)$ and $\sigma(c_j)$, and $C'=C_\sigma-B_\sigma G_\sigma(b)^{-1}G_\sigma(b,c)$. Then
$Q_\sigma=C'^*H_\sigma C'$. As $N^\sigma B_\sigma=0$, $N^\sigma C'$ has the columns $\sigma(Nc_j)$, so
$G_\sigma(Nc)=C'^*(N^\sigma)^*H_\sigma N^\sigma C'$. The columns of $C'$ are independent, since the
$\sigma(b_i)$ and $\sigma(c_j)$ are, and $H_\sigma-(N^\sigma)^*H_\sigma N^\sigma$ is positive definite by
hypothesis. So $Q_\sigma=P+E$ with $P=G_\sigma(Nc)$ and
$E=C'^*(H_\sigma-(N^\sigma)^*H_\sigma N^\sigma)C'$ both positive definite of size $m\geq1$. Write
$P=TT^*$. Then $\det(P+E)=\det P\cdot\det(1+T^{-1}E(T^{-1})^*)$, and the eigenvalues of
$T^{-1}E(T^{-1})^*$ are positive, so $\det(P+E)>\det P>0$.

By the indices and the Gram determinants, summing the logarithms over $\sigma$ gives
$\deg(L;b,c)<\deg(NL,Nc)$. By Lemma~\ref{is-onto-share-degree} and the degrees,
\[
 0=\deg\big(L,(b,c)\big)=\deg(L_b,b)+\deg(L;b,c)<\deg(L_b,b)+\deg(NL,Nc)\le0,
\]
a contradiction. So $N=0$.
\end{proof}

\begin{lemma}\label{is-onto-share-balanced}
Let $X$ be flag data of rank $r$ over $E$ (Definition~\ref{is-onto-share-flag-data}). There is a point
$x\in\mathcal D_X$ with the following two properties.
\begin{enumerate}
\item There is no contraction at $x$.
\item For every nonzero $v\in L'$,
  \[
   \sum_{\tau\in\Psi}m_{x,\tau}\big(\tau(v)\big)\geq\tfrac12[E:\mathbf Q].
  \]
\end{enumerate}
\end{lemma}
\leanref{Development/AbelianVarietyTheory/Proofs/DivisorProofs/ExistsBalancedFrameCMUnitaryProofStatement.lean}{GlobalTheory.T.existsBalancedFrameCMUnitary_proof}
\leanref{Development/AbelianVarietyTheory/Proofs/DivisorProofs/ExistsBalancedFrameCMUnitaryProof.lean}{GlobalTheory.BalancedPointPf.balancedPoint}
\begin{proof}
Apply Lemma~\ref{is-onto-share-semistable} to $K=E$ with complex conjugation, $B=A^{t}$, so that
$S(x,y)=h(y,x)$, and $L=L'$, which is finitely generated, spans $E^r$ as it contains $N_A\Lambda$, and is
self-dual. Let $H$ be the resulting family of matrices. For $\tau\in\Psi$, condition (2) of
Definition~\ref{is-onto-share-flag-data} gives an invertible $P$ whose columns have Gram matrix $J$, that
is, $P^*(A^\tau)^tP=J$. As $H_\tau$ is positive definite and $H_\tau((A^{\tau})^t)^{-1}H_\tau=(A^\tau)^t$,
linear algebra gives $x_\tau$ with $x_\tau^*(A^\tau)^tx_\tau=J$ and $x_\tau^*H_\tau x_\tau=1$. The first
identity says that $x=(x_\tau)$ is orthonormal, so $x\in\mathcal D_X$, and by
Lemma~\ref{is-onto-share-majorant}(4) $m_{x,\tau}(u)=|x_\tau^{-1}u|^2=u^*H_\tau u$.

(1) Let $N$ be contracting at $x$, with constants $c_\tau<1$. For $\tau\in\Psi$ and $u\ne0$,
$(N^\tau u)^*H_\tau(N^\tau u)\le c_\tau^2u^*H_\tau u<u^*H_\tau u$, as $H_\tau$ is positive definite. For
$\sigma\notin\Psi$, $\bar\sigma\in\Psi$, $H_\sigma=\overline{H_{\bar\sigma}}$ and
$N^\sigma u=\overline{N^{\bar\sigma}\bar u}$, so the inequality at $\sigma$ and $u$ is the one at $\bar\sigma$
and $\bar u$. Lemma~\ref{is-onto-share-semistable}(2) gives $N=0$.

(2) Let $v\in L'$ be nonzero, and put $q(\sigma)=\sigma(v)^*H_\sigma\sigma(v)$ for every embedding
$\sigma$ of $E$. Then $q(\sigma)>0$, as $H_\sigma$ is positive definite, and $q(\bar\sigma)=q(\sigma)$, as
$H_{\bar\sigma}=\overline{H_\sigma}$. The submodule $\mathcal O_Ev$ of $L'$ has the basis $v$, and its index
over $\mathcal O_Ev$ is $1$, so Lemma~\ref{is-onto-share-semistable}(1) gives
$\deg(\mathcal O_Ev)=-\frac12\sum_\sigma\log q(\sigma)\le0$. The embeddings of $E$ are the $\tau$ and
$\bar\tau$ for $\tau\in\Psi$, each once, by (1) of Definition~\ref{is-onto-share-flag-data}. So
$|\Psi|=\frac12[E:\mathbf Q]$ and $\sum_{\tau\in\Psi}\log q(\tau)\geq0$. Since $\log t\le t-1$ for $t>0$,
$\sum_{\tau\in\Psi}q(\tau)\geq|\Psi|$. Finally $m_{x,\tau}(\tau(v))=q(\tau)$ for $\tau\in\Psi$.

Part (2) is not used in this chapter.
\end{proof}

\begin{lemma}\label{is-onto-share-free}
Let $M$ be definite two-level data for $U$ and $C$ (Definition~\ref{is-onto-share-two-level}) for which the
level hypothesis holds. If $u\in U$ and $c\in C$ satisfy $uc=c$, then $u=1$.
\end{lemma}
\leanref{Development/AbelianVarietyTheory/Proofs/DivisorProofs/DefiniteCuspFormsFiniteRegularRepresentation.lean}{GlobalTheory.ShareOn.Definite.free_of_definite}
\leanref{Development/AbelianVarietyTheory/Proofs/DivisorProofs/ExistsUpperArithmeticGroupUnipotentUnitary.lean}{GlobalTheory.UpperArith.upper_subgroup}
\begin{proof}
Let $g$ be the prime of the level hypothesis. Let $\Gamma^\circ$ be the set of products $\gamma t(v,c)$ with
$vc=c$ and $\gamma\in\Gamma_c$. It is a subgroup: the cocycle condition at $(1,1)$ puts $t(1,c)$ in
$\Gamma_c$, so $1=t(1,c)^{-1}t(1,c)\in\Gamma^\circ$; for $vc=v'c=c$,
\[
 \gamma t(v,c)\,\gamma't(v',c)=\gamma\big(t(v,c)\gamma't(v,c)^{-1}\big)\big(t(vv',c)(t(v,c)t(v',c))^{-1}\big)^{-1}t(vv',c),
\]
and the inverse of $\gamma t(v,c)$ is a product of the same kind with $v^{-1}$, by the conjugation and
cocycle conditions. Its elements preserve $h$ and $L'_c$, since $t(v,c)$ and $\Gamma_c$ do, and the level
hypothesis gives $(\beta-1)^rL'_c\subset PL'_c$ for $\beta\in\Gamma^\circ$ and every prime $P$ over $g$.

By Lemma~\ref{is-onto-share-balanced} there is $x\in\mathcal D_M$ at which there is no contraction for
$X_c$. The element $t(u,c)$ lies in $\Gamma^\circ$ and preserves $h$, so by
Lemma~\ref{is-onto-share-majorant}(5) and (1),
$m_{x,\tau}(t(u,c)^\tau v)=m_{x,\tau}(v)\le e^{2/(3r)}m_{x,\tau}(v)$. Lemma~\ref{is-onto-share-small} gives
$t(u,c)=1\in\Gamma_c$, and the last condition on $t$ gives $u=1$.
\end{proof}

\begin{lemma}\label{is-onto-share-definite}
Let $M$ be definite two-level data of rank $r$ over $E$ for a finite group $U$ acting on a finite set $C$
(Definition~\ref{is-onto-share-two-level}), $\kappa$ a weight, and $\rho$ a representation of $U$ on
$S_\kappa(M)$ pulled back from the action on $C$. Then $S_\kappa(M)$ is finite-dimensional. Let $H\subset U$
be a subgroup for which the level hypothesis holds. Then there is an integer $n$ with $n|H|=\dim
S_\kappa(M)$ such that $S_\kappa(M)$ embeds $H$-equivariantly into $n$ copies of the left regular
representation of $H$, as in Lemma~\ref{is-onto-share-free-summands}(1), and
\[
 |H|\dim S_\kappa(M)^{H,\lambda}=\dim S_\kappa(M)
\]
for every character $\lambda\colon H\to\mathbf C^\times$.
\end{lemma}
\leanref{Development/AbelianVarietyTheory/Proofs/DivisorProofs/DefiniteCuspFormsFiniteRegularRepresentationStatement.lean}{GlobalTheory.T.definite_cuspForms_finite_regularRepresentation}
\leanref{Development/AbelianVarietyTheory/Proofs/DivisorProofs/DefiniteCuspFormsFiniteRegularRepresentation.lean}{GlobalTheory.ShareOn.Definite.share_definite}
\begin{proof}
Every $X_c$ has $p_0\in\{0,r\}$, so Lemmas~\ref{is-onto-share-definite-finite}
and~\ref{is-onto-share-components}(2) give the finiteness. The data $M|_H$ are definite and satisfy the
level hypothesis, so $H$ acts freely on $C$ by Lemma~\ref{is-onto-share-free}. By
Lemma~\ref{is-onto-share-components}(1), $S_\kappa(M)=\bigoplus_cA_c$ with $\rho(h)A_c\subset A_{hc}$.
Lemma~\ref{is-onto-share-free-summands}, with $\Gamma=U$, $C'=C$ and $H'=H$, gives both claims.
\end{proof}

\subsection{Bounded sections and free balls}\label{subsection-is-free-balls}

\begin{lemma}\label{is-onto-share-pullback}
Let $M$ be two-level data for $U$ and $C$, $F$ a family of weight $\kappa$, $u\in U$ and $c\in C$. Then
$x\mapsto(u^*F)_c(x)=F_{uc}(t(u,c)x)$ is a bounded section of weight $\kappa$ for $X_c$.
\end{lemma}
\leanref{Development/AbelianVarietyTheory/Proofs/DivisorProofs/CuspFormPullbackIsBoundedSectionStatement.lean}{GlobalTheory.T.cuspForm_pullback_isBoundedSection}
\leanref{Development/AbelianVarietyTheory/Proofs/DivisorProofs/CuspFormPullbackIsBoundedSection.lean}{GlobalTheory.PullBdd.isBddSec_transport}
\begin{proof}
Put $\beta=t(u,c)$. Since $\beta^\tau$ preserves $h_\tau$, the Gram matrices of $\beta x$ and $x$ agree, so
$d_i(\beta x,\tau)=d_i(x,\tau)$, $\beta x\in\mathcal D_M$ for $x\in\mathcal D_M$, and
$\|F_{uc}\|_\kappa(\beta x)=\|(u^*F)_c\|_\kappa(x)$. The function is holomorphic on $\mathcal D_M$, and
$\beta(xb)=(\beta x)b$ gives the equivariance under $\mathcal B^\Psi$. For $\gamma\in\Gamma_c$,
$\beta\gamma=(\beta\gamma\beta^{-1})\beta$ with $\beta\gamma\beta^{-1}\in\Gamma_{uc}$, so
$F_{uc}(\beta\gamma x)=F_{uc}(\beta x)$, as $\chi_{uc,\kappa}$ is trivial on $\Gamma_{uc}$; and
$\chi_{c,\kappa}(\gamma)=1$.
\end{proof}

\begin{lemma}[Fourier-Jacobi expansions \cite{MS02,Lan13}]\label{is-onto-sections-fourier}
Let $X$ be flag data of rank $r$ over $E$ (Definition~\ref{is-onto-share-flag-data}), with the notation
there. For an isotropic $\xi\in\Lambda$, a weight $\kappa$, $\nu\in\mathbf R^\Psi$, a function $P$ on
$M_r(\mathbf C)^\Psi$ and $x\in\mathcal D_X$, the \emph{Fourier-Jacobi coefficient} $P^{\kappa}_{\xi,\nu}(x)$
is
\[
 \operatorname{vol}(F_\kappa)^{-1}\int_{F_\kappa}P(\phi_\xi(s)x)\,e^{-4\pi i\sum_\tau\nu_\tau s_\tau}\,ds
\]
if $2\sum_\tau\nu_\tau p_\tau\in\mathbf Z$ for every $p\in\Pi_\kappa$ and $F_\kappa$ exists, and $0$ otherwise;
for $\nu=0$ it is the constant term of $P$ along $\xi$ at $x$. Let $\kappa,\kappa'$ be weights, $P$ a bounded
section of weight $\kappa$, $Q$ a bounded section of weight $\kappa'$, $\xi\in\Lambda$ isotropic and
$x\in\mathcal D_X$.
\begin{enumerate}
\item If $P^\kappa_{\xi,\nu}(x)\ne0$, then $\nu_\tau\geq0$ for every $\tau\in\Psi$.
\item For every $\nu\in\mathbf R^\Psi$, the family $\big(P^\kappa_{\xi,\nu'}(x)\,Q^{\kappa'}_{\xi,\nu-\nu'}(x)\big)_{\nu'\in
  \mathbf R^\Psi}$ is summable, with sum $(PQ)^{\kappa+\kappa'}_{\xi,\nu}(x)$.
\end{enumerate}
\end{lemma}
\leanref{ReciprocityTheory/ReductiveGroups/ArithmeticGeometry/FundamentalDomains/ModerateCoefficientPositivity.lean}{ReciprocityTheory.ChartRestriction.G.FD.ModerateCoefficientPositivity}
\leanref{ReciprocityTheory/ReductiveGroups/ArithmeticGeometry/FundamentalDomains/ModerateCoefficientPositivity.lean}{ReciprocityTheory.ChartRestriction.G.FD.CoefficientConvolution}
\leanref{ReciprocityTheory/ReductiveGroups/ArithmeticGeometry/FundamentalDomains/ModerateCoefficientPositivityHolds.lean}{ReciprocityTheory.ChartRestriction.G.FD.moderateCoefficientPositivity_holds}
\leanref{ReciprocityTheory/ReductiveGroups/ArithmeticGeometry/FundamentalDomains/CoefficientConvolutionHolds.lean}{ReciprocityTheory.ChartRestriction.G.FD.coefficientConvolution_holds}
\leanref{ReciprocityTheory/ReductiveGroups/ArithmeticGeometry/FundamentalDomains/BoundedSectionAlgebra.lean}{ReciprocityTheory.ChartRestriction.G.FD.BoundedSectionAlgebra.mod_of_bdd}
\leanref{ReciprocityTheory/ReductiveGroups/ArithmeticGeometry/FundamentalDomains/BoundedSectionAlgebra.lean}{ReciprocityTheory.ChartRestriction.G.FD.BoundedSectionAlgebra.sign_of_bdd}

\begin{lemma}\label{is-onto-share-sections}
Let $X$ be flag data and $\kappa,\kappa'$ weights.
\begin{enumerate}
\item $S_\kappa(X)$ is a complex vector space, and $PQ\in S_{\kappa+\kappa'}(X)$ for $P\in S_\kappa(X)$ and
  every bounded section $Q$ of weight $\kappa'$.
\item $\mathcal D_X$ is open and connected.
\end{enumerate}
\end{lemma}
\leanref{ReciprocityTheory/ReductiveGroups/ArithmeticGeometry/FundamentalDomains/SiegelSpaceAlgebraHolds.lean}{ReciprocityTheory.ChartRestriction.G.FD.siegelSpaceAlgebra_holds}
\leanref{ReciprocityTheory/ReductiveGroups/ArithmeticGeometry/FundamentalDomains/BoundedSectionAlgebra.lean}{ReciprocityTheory.ChartRestriction.G.FD.BoundedSectionAlgebra.siegelSpaceAlgebra_of}
\leanref{ReciprocityTheory/ReductiveGroups/ArithmeticGeometry/FundamentalDomains/FourierCoefficientDomain.lean}{ReciprocityTheory.ChartRestriction.G.FD.FourierCoefficientDomain.isOpen_dom}
\leanref{ReciprocityTheory/ReductiveGroups/ArithmeticGeometry/TaylorDiscriminants/WickVanishing/GGConn.lean}{ReciprocityTheory.ChartRestriction.G.TB.WickVanishing.flag_domain_frames_connected}
\leanref{ReciprocityTheory/ReductiveGroups/ArithmeticGeometry/TaylorDiscriminants/WickVanishing/GGConnA.lean}{ReciprocityTheory.ChartRestriction.G.TB.WickVanishing.GGC.isPreconnected_ok}
\begin{proof}
(1) The sum and the scalar multiples of sections of weight $\kappa$ are sections of weight $\kappa$, and
$\|P+Q\|_\kappa\le\|P\|_\kappa+\|Q\|_\kappa$, so bounded sections form a vector space. For real $s$, the Gram
matrices of $\phi_\xi(s)x$ and $x$ have the same leading minors, because $\xi$ is isotropic; so the factor
$\|P\|_\kappa(x)/|P(x)|$ is constant along $s\mapsto\phi_\xi(s)x$, and a bounded $P$ is bounded on these
curves. Hence the integrals defining constant terms converge, and constant terms are linear in $P$. So
$S_\kappa(X)$ is a vector space; it contains $0$.

Let $P\in S_\kappa(X)$ and $Q$ a bounded section of weight $\kappa'$. Then $PQ$ is holomorphic on
$\mathcal D_X$; it is equivariant of weight $\kappa+\kappa'$ under $\mathcal B^\Psi$, and under $\Gamma$ through
$\chi_\kappa\chi_{\kappa'}=\chi_{\kappa+\kappa'}$; and $\|PQ\|_{\kappa+\kappa'}=\|P\|_\kappa\|Q\|_{\kappa'}$ is
bounded. Let $\xi$ be isotropic and $x\in\mathcal D_X$. By Lemma~\ref{is-onto-sections-fourier}(2), the
constant term of $PQ$ along $\xi$ at $x$ is $\sum_{\nu'}P^\kappa_{\xi,\nu'}(x)Q^{\kappa'}_{\xi,-\nu'}(x)$. If a term
is nonzero, Lemma~\ref{is-onto-sections-fourier}(1) gives $\nu'_\tau\geq0$ and $-\nu'_\tau\geq0$ for all $\tau$,
so $\nu'=0$; but $P^\kappa_{\xi,0}(x)=0$, as $P$ is a cusp form. So every term vanishes, and $PQ$ is a cusp form.

(2) The set $\mathcal D_X$ is defined by finitely many conditions $\det x_\tau\ne0$ and
$\epsilon_i\operatorname{Re}d_i(x,\tau)>0$ on continuous functions of $x$, so it is open. For connectedness,
write $J$ also for the diagonal matrix of the signs $\operatorname{sgn}(a)$, and fix for every $\tau\in\Psi$ an
invertible $P_\tau$ whose columns have Gram matrix $J$ for $h_\tau$ (Definition~\ref{is-onto-share-flag-data}(2)).
For $y\in GL_r(\mathbf C)$ put $G(y)=y^tJ\bar y$, the Gram matrix of the columns of $P_\tau y$. Call $y$
\emph{positive} if $JG(y)$ is positive definite.

If every $y_\tau$ is positive, then $x=(P_\tau y_\tau)_\tau$ lies in $\mathcal D_X$. Indeed the leading
$i\times i$ block of $JG(y_\tau)$ is positive definite, and it is the product of the leading block of $J$,
of determinant $\epsilon_i$, with that of $G(y_\tau)$, of determinant $d_i(x,\tau)$; so
$\epsilon_id_i(x,\tau)>0$, and $\det y_\tau\ne0$.

Call $y$ \emph{block diagonal} if $y_{ka}=0$ whenever exactly one of $k,a$ is at most $p_0$. An invertible
block diagonal $y$ is positive, as then $JG(y)=\bar y^*\bar y$. The tuples of invertible block diagonal
matrices form a connected set $V$: for two of them $a,b$ the determinants of the components of $(1-z)a+zb$
are polynomials in $z\in\mathbf C$, nonzero at $z=0$, so they vanish at only finitely many $z$, and the
complement of a finite set in $\mathbf C$ is connected. In the same way the tuples of invertible upper
triangular matrices, $\mathcal B^\Psi$, form a connected set.

Let $y$ be positive. As $JG(y)$ and $G(y)$ are Hermitian, $JG(y)=G(y)J$, so $G(y)_{ab}=0$ when exactly one of
$a,b$ is at most $p_0$. For $0\le t\le1$ let $y(t)$ be $y$ with the entries $y_{ka}$, for $k,a$ on different
sides of $p_0$, multiplied by $t$. A direct computation with this block form gives
$JG(y(t))=JG(y)+W^*W$, where $W_{ka}=\sqrt{1-t^2}\,\bar y_{ka}$ for $k,a$ on different sides of $p_0$ and
$W_{ka}=0$ otherwise. So $y(t)$ is positive, and the path $t\mapsto y(t)$ joins $y=y(1)$ to the block
diagonal $y(0)$ through positive matrices. Hence the set $K$ of tuples of positive matrices is connected,
and so is its image $\{(P_\tau y_\tau)\}$ in $\mathcal D_X$, which contains $(P_\tau)_\tau$.

Let $x\in\mathcal D_X$. By Lemma~\ref{is-onto-share-majorant}(2) and (3), the connected set
$\{xb:b\in\mathcal B^\Psi\}$ lies in $\mathcal D_X$, contains $x$ and contains an orthonormal point $x'$. For
$y_\tau=P_\tau^{-1}x'_\tau$, $G(y_\tau)=J$, so $y_\tau$ is positive and $x'$ lies in the image of $K$. So
every point of $\mathcal D_X$ is joined to $(P_\tau)_\tau$ inside $\mathcal D_X$ by a connected set, and
$\mathcal D_X$ is connected.
\end{proof}

\begin{lemma}\label{is-onto-step-free-ball-majorant}
Let $0\le p\le r$ be integers and $J=\operatorname{diag}(1,\dots,1,-1,\dots,-1)$, with $p$ entries $1$. For
$u\in\mathbf C^r$ write $|u|^2=\sum_a|u_a|^2$, and $u_+$ and $u_-$ for the vectors of its first $p$ and of
its last $r-p$ coordinates. Let $H$ be a positive definite Hermitian $r\times r$ matrix with $HJH=J$, and
$0\le s<1$ a real number such that $|u_-|^2\le s^2|u_+|^2$ for every $u$ with $JHu=u$. Then
\[
 u^*Hu\le\frac{1+s}{1-s}|u|^2\quad\text{and}\quad|u|^2\le\frac{1+s}{1-s}\,u^*Hu\qquad(u\in\mathbf C^r).
\]
\end{lemma}
\leanref{ReciprocityTheory/SemistableReduction/ShareBall.lean}{ReciprocityTheory.MultiplicativeHeights.boundedBall_holds}
\leanref{ReciprocityTheory/SemistableReduction/ShareSt.lean}{ReciprocityTheory.MultiplicativeHeights.BoundedBallStatement}
\leanref{ReciprocityTheory/SemistableReduction/ShareBallC.lean}{ReciprocityTheory.MultiplicativeHeights.sb_setup}
\leanref{ReciprocityTheory/SemistableReduction/ShareBallB.lean}{ReciprocityTheory.MultiplicativeHeights.sb_real_bounds}
\begin{proof}
Let $V_+$ and $V_-$ be the eigenspaces of $JH$ for the eigenvalues $1$ and $-1$. As
$(JH)^2=J(HJH)=J^2=1$, every $u$ is $x+y$ with $x=\frac12(u+JHu)\in V_+$ and $y=\frac12(u-JHu)\in V_-$.
For $x\in V_+$ and $y\in V_-$, $Hx=J(JHx)=Jx$ and $Hy=-Jy$. As $H$ and $J$ are Hermitian,
$x^*Jy=(Hx)^*y=x^*Hy=-x^*Jy$, so $x^*Jy=0$, that is, $x_+^*y_+=x_-^*y_-$.

\emph{The map $x\mapsto x_+$ is a bijection from $V_+$ onto $\mathbf C^p$.} If $x\in V_+$ and $x_+=0$, then
$x^*Hx=x^*Jx=-|x_-|^2\le0$, so $x=0$, as $H$ is positive definite. In the same way $y\in V_-$ and $y_-=0$
give $y^*Hy=-y^*Jy=-|y_+|^2\le0$ and $y=0$. So $\dim V_+\le p$ and $\dim V_-\le r-p$, and both are
equalities since $V_++V_-=\mathbf C^r$.

\emph{The space $V_-$.} Let $y\in V_-$, and choose $x\in V_+$ with $x_+=y_+$. Then
$|y_+|^2=x_+^*y_+=x_-^*y_-$, so $|y_+|^2\le|x_-||y_-|\le s|x_+||y_-|=s|y_+||y_-|$, and $|y_+|\le s|y_-|$.

\emph{The bounds.} Let $u=x+y$ as above, and put $a=|x_+|$, $b=|x_-|$, $c=|y_+|$, $d=|y_-|$ and
$t=\operatorname{Re}(x_+^*y_+)=\operatorname{Re}(x_-^*y_-)$. Then $b\le sa$ and $c\le sd$, and $|t|\le ac$ and
$|t|\le bd$ by the Cauchy-Schwarz inequality. As $x^*Jy=y^*Jx=0$,
\[
 u^*Hu=x^*Jx-y^*Jy=a^2-b^2-c^2+d^2,\qquad|u|^2=a^2+b^2+c^2+d^2+4t.
\]
So the two claims are $s(a^2+d^2)+b^2+c^2+2(1+s)t\geq0$ and $s(a^2+d^2)-(b^2+c^2)-2(1-s)t\geq0$.

For the second, $t\le ac\le sad$, $b^2\le s^2a^2$ and $c^2\le s^2d^2$ bound the left side below by
$s(1-s)(a^2+d^2)-2s(1-s)ad=s(1-s)(a-d)^2\geq0$.

For the first it suffices that $2(1+s)|t|\le s(a^2+d^2)+b^2+c^2$. The substitution
$(a,b,c,d)\mapsto(d,c,b,a)$ preserves the hypotheses and this inequality, so it suffices to treat the
case $ac\le bd$, in which $|t|\le ac$. If $d=0$, then $c=0$ and the inequality is clear. If $d>0$, then
$a^2c^2\le b^2d^2$ and $c\le sd\le d$ give
\[
 d^2\big(s(a^2+d^2)+b^2+c^2-2(1+s)ac\big)\geq(sd^2+c^2)(a-d)^2+2ad(d-c)(sd-c)\geq0.
\]
% Lean: boundedBall_holds: hdec (x and y), hHx, hHy, horth and horth' (x^*Jy = 0), sb_setup (the
% bijection, through the injective map F = D_p P + D_n N), hNprop (the space V_-), hq_eq, hn_eq, hts1,
% hts2 and sb_real_bounds, whose first inequality is two_mul_one_add_mul_le_of_le_mul applied to
% (a, b, c, d) or to (d, c, b, a), with the identity above (e1, e2).
\end{proof}

\begin{lemma}\label{is-onto-step-free-ball-minor}
Let $X$ consist of the data (1) to (6) of flag data of rank $r$ over $E$
(Definition~\ref{is-onto-share-flag-data}), and let $J=\operatorname{diag}(1,\dots,1,-1,\dots,-1)$, with
$p_0$ entries $1$. For a square matrix $N$ write $\ell_i(N)$ for its leading principal $i\times i$ minor.
\begin{enumerate}
\item If $g\in M_r(\mathbf C)$ satisfies $g^*Jg=J$, then $1+\sum_{a>p_0\geq b}|g_{ab}|^2\le|\ell_{p_0}(g)|^2$.
\item If $x_0$ and $x_1$ are orthonormal points and $\tau\in\Psi$, then $g=x_{0,\tau}^{-1}x_{1,\tau}$
  satisfies $g^*Jg=J$.
\item Let $x_0$ be an orthonormal point, $x\in\mathcal D_X$ and $\tau\in\Psi$. Then
  $\ell_{p_0}(x_{0,\tau}^{-1}x_\tau)\ne0$, and the quotient
  \[
   q_\tau(x_0,x)=\frac{|d_{p_0}(x,\tau)|}{|\ell_{p_0}(x_{0,\tau}^{-1}x_\tau)|^2}
  \]
  satisfies $0<q_\tau(x_0,x)\le1$ and $q_\tau(x_0,xb)=q_\tau(x_0,x)$ for $b\in\mathcal B^\Psi$.
\end{enumerate}
\end{lemma}
\leanref{Development/AbelianVarietyTheory/Proofs/TorusHaarMeasure/OneAddLeLminorSq.lean}{GlobalTheory.Up.BallExPf.one_add_le_lminor_sq}
\leanref{Development/AbelianVarietyTheory/Proofs/TorusHaarMeasure/OneLeMinorOfIsOrth.lean}{GlobalTheory.FramePotential.one_le_minor_of_isOrth}
\leanref{Development/AbelianVarietyTheory/Proofs/TorusHaarMeasure/OneLeMinorOfIsOrth.lean}{GlobalTheory.FramePotential.factor_mem}
\leanref{Development/AbelianVarietyTheory/Proofs/TorusHaarMeasure/OneLeMinorOfIsOrth.lean}{GlobalTheory.FramePotential.factor_rmul}
\begin{proof}
(1) Let $g_{11}$ be the leading $p_0\times p_0$ block of $g$ and $g_{21}$ the block of the entries $g_{ab}$
with $a>p_0\geq b$. The leading $p_0\times p_0$ block of the identity $g^*Jg=J$ reads
$g_{11}^*g_{11}-g_{21}^*g_{21}=1$. The Hermitian matrix $g_{21}^*g_{21}$ has eigenvalues $\lambda_k\geq0$, so
\[
 |\ell_{p_0}(g)|^2=\det(g_{11}^*g_{11})=\prod_k(1+\lambda_k)\geq1+\sum_k\lambda_k=1+\sum_{a>p_0\geq b}|g_{ab}|^2.
\]

(2) Orthonormal points have invertible components (Lemma~\ref{is-onto-share-majorant}(4)), and
$x_{1,\tau}=x_{0,\tau}g$. The Gram matrices of the columns of $x_{0,\tau}$ and $x_{1,\tau}$ are both $J$,
so $J=x_{1,\tau}^tA^\tau\bar x_{1,\tau}=g^t\big(x_{0,\tau}^tA^\tau\bar x_{0,\tau}\big)\bar g=g^tJ\bar g$.
Complex conjugation gives $g^*Jg=J$, as $J$ is real.

(3) For $b\in\mathcal B^\Psi$ let $\beta_i$ be the leading $i\times i$ block of $b_\tau$. As $b_\tau$ is upper
triangular, the leading $i\times i$ block of $Nb_\tau$ is that of $N$ times $\beta_i$, so
$\ell_i(Nb_\tau)=\ell_i(N)\det\beta_i$, with $\det\beta_i\ne0$; and
$d_i(xb,\tau)=|\det\beta_i|^2d_i(x,\tau)$, as in the proof of Lemma~\ref{is-onto-share-majorant}(2). By
Lemma~\ref{is-onto-share-majorant}(3), $x=x_1b$ with $x_1$ orthonormal and $b\in\mathcal B^\Psi$. By (2)
and (1), $g=x_{0,\tau}^{-1}x_{1,\tau}$ has $|\ell_{p_0}(g)|\geq1$, so
$\ell_{p_0}(x_{0,\tau}^{-1}x_\tau)=\ell_{p_0}(g)\det\beta_{p_0}\ne0$. The Gram matrix of the columns of
$x_{1,\tau}$ is $J$, so $|d_{p_0}(x_1,\tau)|=1$, and
\[
 q_\tau(x_0,x)=\frac{|\det\beta_{p_0}|^2\,|d_{p_0}(x_1,\tau)|}{|\det\beta_{p_0}|^2\,|\ell_{p_0}(g)|^2}
 =|\ell_{p_0}(g)|^{-2}\in(0,1].
\]
For $b'\in\mathcal B^\Psi$, $xb'\in\mathcal D_X$ by Lemma~\ref{is-onto-share-majorant}(2), and the factors
$|\det\beta'_{p_0}|^2$ of the numerator and of the denominator cancel in the same way, so
$q_\tau(x_0,xb')=q_\tau(x_0,x)$.
% Lean: (1) is one_add_le_lminor_sq (key: A^H A = 1 + C^H C, and sumBd_det_one_add_ge); the book's
% (2) is the Gram computation of one_le_minor_of_isOrth (FrameDegreeData.gram_mul) and of frame_bound
% (gram_iff); (3) is factor_mem (ortho_holds, factor_rmul, norm_dmin_of_isOrth and
% one_le_minor_of_isOrth through one_le_norm_lminor_sq) and factor_rmul (FrameDegreeData.dmin_rmul and
% lminor_mul_upper).
\end{proof}

\begin{definition}\label{is-onto-step-free-ball}
Let $X$ consist of the data (1) to (6) of flag data of rank $r$ over $E$
(Definition~\ref{is-onto-share-flag-data}), and let $x_0$ be an orthonormal point. For $x\in\mathcal D_X$ put
\[
 \rho(x_0,x)=-\log\prod_{\tau\in\Psi}q_\tau(x_0,x),
\]
with the quotients $q_\tau(x_0,x)\in(0,1]$ of Lemma~\ref{is-onto-step-free-ball-minor}(3). Let
$G\subset GL_r(E)$ be a subgroup and $\rho_0$ a real number. Then $G$ has a \emph{free ball} of radius
$\rho_0$ at $x_0$ if $\Gamma\subset G$, every element of $G$ preserves $h$ and $\Lambda$, and
$\rho(x_0,\gamma x)\geq\rho_0$ for every $\gamma\ne1$ in $G$ and every $x\in\mathcal D_X$ with
$\rho(x_0,x)<\rho_0$; here $\gamma x\in\mathcal D_X$, as $\gamma$ preserves $h$.
\end{definition}
\leanref{Development/AbelianVarietyTheory/Definitions/UpperBounds/FreeBallIffOfDimXEqZero.lean}{GlobalTheory.Up.rhX}
\leanref{Development/AbelianVarietyTheory/Definitions/UpperBounds/FreeBallIffOfDimXEqZero.lean}{GlobalTheory.Up.FreeBall}

\begin{lemma}\label{is-onto-step-free-ball-bound}
Let $X$ consist of the data (1) to (6) of flag data of rank $r$ over $E$
(Definition~\ref{is-onto-share-flag-data}), let $x_0$ be an orthonormal point, and write
$|u|^2=\sum_a|u_a|^2$ for $u\in\mathbf C^r$.
\begin{enumerate}
\item For $x\in\mathcal D_X$ and $b\in\mathcal B^\Psi$,
  $\rho(x_0,x)=-\sum_{\tau\in\Psi}\log q_\tau(x_0,x)\geq0$ and $\rho(x_0,xb)=\rho(x_0,x)$
  (Definition~\ref{is-onto-step-free-ball}).
\item Let $r\geq1$, put $\beta_r=(1+1/(7r))/(1-1/(7r))$, and let $x_1$ be an orthonormal point with
  $\rho(x_0,x_1)<1/(49r^2)$. Then for every $\tau\in\Psi$ the matrix $g=x_{0,\tau}^{-1}x_{1,\tau}$
  satisfies $|g^{-1}u|^2\le\beta_r|u|^2$ and $|gu|^2\le\beta_r|u|^2$ for every $u\in\mathbf C^r$.
  Moreover $\beta_r\le e^{1/(3r)}$.
\end{enumerate}
\end{lemma}
\leanref{Development/AbelianVarietyTheory/Definitions/UpperBounds/FreeBallIffOfDimXEqZero.lean}{GlobalTheory.Up.rhX_eq_sum_log}
\leanref{Development/AbelianVarietyTheory/Definitions/UpperBounds/FreeBallIffOfDimXEqZero.lean}{GlobalTheory.Up.rhX_nonneg}
\leanref{Development/AbelianVarietyTheory/Definitions/UpperBounds/FreeBallIffOfDimXEqZero.lean}{GlobalTheory.Up.rhX_rmul}
\leanref{Development/AbelianVarietyTheory/Proofs/TorusHaarMeasure/OneAddLeLminorSq.lean}{GlobalTheory.Up.BallExPf.frame_bound}
\leanref{Development/AbelianVarietyTheory/Proofs/TorusHaarMeasure/OneAddLeLminorSq.lean}{GlobalTheory.Up.BallExPf.norm_bound}
\leanref{ReciprocityTheory/SemistableReduction/ShareBallA.lean}{ReciprocityTheory.MultiplicativeHeights.sb_ballRatio_le}
\begin{proof}
(1) By Lemma~\ref{is-onto-step-free-ball-minor}(3) the quotients lie in $(0,1]$, so the logarithm of their
product is the sum of their logarithms, each at most $0$; and each quotient is unchanged when $x$ is
replaced by $xb$.

(2) Put $s=1/(7r)$, so that $0<s<1$, $s^2=1/(49r^2)$ and $\beta_r=(1+s)/(1-s)$. Let $J$ be as in
Lemma~\ref{is-onto-step-free-ball-minor}; by part (2) of that lemma, $g^*Jg=J$.

\emph{The minor.} The Gram matrix of the columns of $x_{1,\tau}$ is $J$, so $|d_{p_0}(x_1,\tau)|=1$ and
$q_\tau(x_0,x_1)=|\ell_{p_0}(g)|^{-2}$. By (1) every term $-\log q_{\tau'}(x_0,x_1)$ is nonnegative, so
$\log|\ell_{p_0}(g)|^2=-\log q_\tau(x_0,x_1)\le\rho(x_0,x_1)<s^2$. Let
$c=\sum_{a>p_0\geq b}|g_{ab}|^2$. By Lemma~\ref{is-onto-step-free-ball-minor}(1),
$\log(1+c)\le\log|\ell_{p_0}(g)|^2<s^2$; as $c/(1+c)\le\log(1+c)$, $c\le s^2(1+c)$.

\emph{The majorant.} Let $H=(g^{-1})^*g^{-1}$, a positive definite Hermitian matrix. As $g^*Jg=J$ and
$J^2=1$, $g^{-1}=Jg^*J$, so $gJg^*=J$ and $H=Jgg^*J$; hence $HJH=Jg(g^*Jg)g^*J=JgJg^*J=J$. Let $v$
satisfy $JHv=v$, that is, $gg^*Jv=v$, and put $w=g^{-1}v=Jg^*Jv$. Then $Jw=g^*Jv=g^{-1}v=w$, so the last
$r-p_0$ coordinates of $w$ vanish. With the notation of Lemma~\ref{is-onto-step-free-ball-majorant} for
$p=p_0$, $v_-$ is the product of the block $(g_{ab})_{a>p_0\geq b}$ with $w_+$, as $v=gw$, so
$|v_-|^2\le c|w|^2$; and $|v_+|^2-|v_-|^2=v^*Jv=w^*g^*Jgw=w^*Jw=|w|^2$. Hence
$(1+c)|v_-|^2\le c|v_+|^2\le s^2(1+c)|v_+|^2$, and $|v_-|^2\le s^2|v_+|^2$. So
Lemma~\ref{is-onto-step-free-ball-majorant} applies to $H$: it gives $|g^{-1}u|^2=u^*Hu\le\beta_r|u|^2$, and,
applied to $gu$, $|gu|^2\le\beta_r(gu)^*H(gu)=\beta_r|u|^2$.

\emph{The constant.} With $y=1/(7r)\le1/7$,
$(1+\frac73y)(1-y)-(1+y)=\frac13y(1-7y)\geq0$, so $\beta_r\le1+\frac73y=1+\frac1{3r}\le e^{1/(3r)}$.
% Lean: (1) is rhX_eq_sum_log, rhX_nonneg and rhX_rmul, all through factor_mem and factor_rmul. (2) is
% frame_bound, which applies norm_bound with s = 1/(7r) (ballRatio r = (1 + s)/(1 - s)), gram_iff for
% g^H J g = J, and neg_log_factor_le_rhX with norm_dmin_of_isOrth for the minor. norm_bound: hcs
% (log_le_log, one_sub_inv_le_log_of_pos), H = (g^-1)^H g^-1 with hinv and hinvH, the hypothesis of
% boundedBall_holds through hJw, hDw, b1 (sb_normSq_mulVec_le), b2 and b5, and hq_ct_mul. The constant
% is sb_ballRatio_le.
\end{proof}

\begin{lemma}\label{is-onto-step-free-ball-exists}
Let $X$ be flag data of rank $r\geq1$ over $E$ (Definition~\ref{is-onto-share-flag-data}), and $g>3r$ a
prime that splits completely in $E$. Let $G\subset GL_r(E)$ be a subgroup that contains $\Gamma$, whose
elements preserve $h$, $L'$ and $\Lambda$ and satisfy, for every prime $P$ of $\mathcal O_E$ over $g$,
$(\gamma-1)^kL'\subset PL'$ for some $k\geq0$. Let $x\in\mathcal D_X$ be a point at which there is no
contraction. Then there are an orthonormal point $x_0$ and $b\in\mathcal B^\Psi$ with $x=x_0b$ such that
$G$ has a free ball of radius $1/(49r^2)$ at $x_0$ (Definition~\ref{is-onto-step-free-ball}).
\end{lemma}
\leanref{Development/AbelianVarietyTheory/Proofs/DivisorProofs/FreeBallNoContractionUnitaryUpperLevelStatement.lean}{GlobalTheory.T.freeBall_noContraction_unitaryUpperLevel}
\leanref{Development/AbelianVarietyTheory/Proofs/DivisorProofs/FreeBallNoContractionUnitaryUpperLevel.lean}{GlobalTheory.FreeBallAt.freeBall_at}
\leanref{Development/AbelianVarietyTheory/Proofs/DivisorProofs/FreeBallNoContractionUnitaryUpperLevel.lean}{GlobalTheory.FreeBallAt.opNormLe_of_ball}
\begin{proof}
By Lemma~\ref{is-onto-share-majorant}(3), $x=x_0b$ with $x_0$ orthonormal and $b\in\mathcal B^\Psi$. By
Lemma~\ref{is-onto-share-majorant}(2), $m_{x_0,\tau}=m_{x,\tau}$ for every $\tau$, so there is no contraction
at $x_0$ either. The group $G$ contains $\Gamma$, and its elements preserve $h$ and $\Lambda$.

Let $\gamma\ne1$ in $G$ and $z\in\mathcal D_X$ with $\rho(x_0,z)<1/(49r^2)$, and suppose that
$\rho(x_0,\gamma z)<1/(49r^2)$. By Lemma~\ref{is-onto-share-majorant}(3), $z=z_1b'$ with $z_1$ orthonormal
and $b'\in\mathcal B^\Psi$. Then $\gamma z=(\gamma z_1)b'$, so by Lemma~\ref{is-onto-step-free-ball-bound}(1)
$\rho(x_0,z_1)=\rho(x_0,z)$ and $\rho(x_0,\gamma z_1)=\rho(x_0,\gamma z)$; and $\gamma z_1$ is orthonormal, as
$\gamma$ preserves $h$. Let $\tau\in\Psi$ and $u\in\mathbf C^r$. Lemma~\ref{is-onto-step-free-ball-bound}(2),
for $x_0$ and $z_1$ with the vector $x_{0,\tau}^{-1}u$, and for $x_0$ and $\gamma z_1$ with the vector
$z_{1,\tau}^{-1}u$, gives
\[
 |z_{1,\tau}^{-1}u|^2\le\beta_r|x_{0,\tau}^{-1}u|^2,\qquad
 |x_{0,\tau}^{-1}\gamma^\tau u|^2\le\beta_r|z_{1,\tau}^{-1}u|^2.
\]
By Lemma~\ref{is-onto-share-majorant}(4), $m_{x_0,\tau}(v)=|x_{0,\tau}^{-1}v|^2$, so
$m_{x_0,\tau}(\gamma^\tau u)\le\beta_r^2m_{x_0,\tau}(u)\le e^{2/(3r)}m_{x_0,\tau}(u)$.
Lemma~\ref{is-onto-share-small}, with $\Gamma^\circ=G$ and the point $x_0$, gives $\gamma=1$, a
contradiction. So $\rho(x_0,\gamma z)\geq1/(49r^2)$.
% Lean: FreeBallAt.freeBall_at: ortho_holds for x1 and b, noContraction_of_rmul (MajInv.maj_rmul),
% ortho_holds for z1 and b', FundSetPf.lact_rmul and rhX_rmul, isOrth_lact_of_isUStab, opNormLe_of_ball
% (ModeHalfA.maj_orth, frame_bound twice, sb_ballRatio_le), and smallUnipotentUnitaryElementIsOne_proof
% (Lemma is-onto-share-small). The Lean also allows r = 0, where its radius 1 / (49 r^2) is 0 and
% rhX_nonneg makes the claim empty.
\end{proof}

\subsection{Interpolation when $p_0(r-p_0)\geq4$}\label{subsection-is-interpolation-large}

\begin{situation}\label{is-onto-step-orbit-situation}
Let $X$ consist of the data (1) to (6) of flag data of rank $r$ over $E$
(Definition~\ref{is-onto-share-flag-data}), with $0<p_0<r$, such that $\chi_\kappa$ is trivial on $\Gamma$
for every weight $\kappa$. Let $G\subset GL_r(E)$ be a subgroup and $x_0$ an orthonormal point such that $G$
has a free ball of radius $1/(49r^2)$ at $x_0$ (Definition~\ref{is-onto-step-free-ball}). Let $s$ be an even
integer with $s\geq2r+50r^2|\Psi|p_0(r-p_0)+1$, and let $\delta_s$ be the weight with
$(\delta_s)_{\tau,i}=-s/2$ for $i\le p_0$ and $(\delta_s)_{\tau,i}=s/2$ for $i>p_0$.
\end{situation}
\leanref{ReciprocityTheory/ReductiveGroups/ArithmeticGeometry/FundamentalDomains/Signature.lean}{ReciprocityTheory.ChartRestriction.G.FD.Data}
\leanref{Development/AbelianVarietyTheory/BlockTheory/PoincareCuspForms/TwoLevelLogarithm.lean}{GlobalTheory.IndefiniteRowSideFullHecke.FrameWeightGapBounds.deMul}

\begin{lemma}\label{is-onto-step-orbit-interpolation-large-span}
Let $A$ be a set, $S$ a complex vector space of functions $A\to\mathbf C$, $a\in A$, and $(p_i)_{i\in I}$ a finite
family of points of $A$. Suppose that for every real number $\eta>0$ there is $u\in S$ with $|u(a)-1|\le\eta$ and
$|u(p_i)|\le\eta$ for every $i\in I$. Then there is $u\in S$ with $u(a)=1$ and $u(p_i)=0$ for every $i\in I$.
\end{lemma}
\leanref{Development/AbelianVarietyTheory/Proofs/TorusHaarMeasure/SumRangeBlocks.lean}{GlobalTheory.Up.OrbitInterpolation.exists_of_approx}
\begin{proof}
The map $u\mapsto\big(u(a),(u(p_i))_{i\in I}\big)$ from $S$ to $\mathbf C\times\mathbf C^I$ is linear. Its image is a
linear subspace of a finite-dimensional space, so it is closed. By hypothesis the point $(1,0)$ is a limit of points
of the image, for the largest absolute value of the coordinates. So $(1,0)$ lies in the image.
% Lean: exists_of_approx: L = (proj a).prod (pi (proj (p i))), Submodule.closed_of_finiteDimensional for the
% image, and Metric.mem_closure_iff with the u given for eta / 2.
\end{proof}

\begin{lemma}\label{is-onto-step-orbit-interpolation-large-weight-near}
Let $X$ consist of the data (1) to (6) of flag data of rank $r$ over $E$ (Definition~\ref{is-onto-share-flag-data}),
use the charts $c_y$ of Definition~\ref{is-onto-share-counts-model}, write $|z|$ for the Euclidean norm of
$z\in\mathbf C^n$ and $\ell_i(N)$ for the leading principal $i\times i$ minor of a square matrix $N$, and let $x_0$ be
an orthonormal point. For $y\in M_r(\mathbf C)^\Psi$ put
\[
 q_{x_0}(y)=\prod_{\tau\in\Psi}\frac{|d_{p_0}(y,\tau)|}{|\ell_{p_0}(x_{0,\tau}^{-1}y_\tau)|^2},
\]
a quotient by $0$ being read as $0$. Then there are real numbers $\varrho>0$ and $C$ such that
$q_{x_0}(c_x(z))\le C\,q_{x_0}(x)$ for every orthonormal point $x$ and every $z\in\mathbf C^n$ with $|z|<\varrho$.
\end{lemma}
\leanref{Development/AbelianVarietyTheory/Proofs/TorusHaarMeasure/MulSgnMatMulConjTranspose.lean}{GlobalTheory.MinorRatioNear.minorRatio_near}
\leanref{Development/AbelianVarietyTheory/Proofs/TorusHaarMeasure/MulSgnMatMulConjTranspose.lean}{GlobalTheory.MinorRatioNear.exists_block_det_ge}
\leanref{Development/AbelianVarietyTheory/Proofs/TorusHaarMeasure/MulSgnMatMulConjTranspose.lean}{GlobalTheory.MinorRatioNear.norm_entry_le_one}
\begin{proof}
Let $J$ be the diagonal matrix with $p_0$ entries $1$ followed by $r-p_0$ entries $-1$, and write $M^*$ for the
conjugate transpose of a complex matrix $M$.

\emph{A block of $g$.} Let $g\in M_r(\mathbf C)$ with $gJg^*=J$, let $g_0$ be its leading $p_0\times p_0$ block and $R$
the $p_0\times r$ matrix of its first $p_0$ rows. The leading $p_0\times p_0$ block of $gJg^*=J$ is $RJR^*=1$, so every
row vector $v\in\mathbf C^{p_0}$ satisfies $\sum_a\operatorname{sgn}(a)|(vR)_a|^2=|v|^2$. As $(vR)_a=(vg_0)_a$ for
$a\le p_0$, this reads
\[
 |v|^2+\sum_{a>p_0}|(vR)_a|^2=|vg_0|^2 .
\]
So $vg_0=0$ forces $v=0$, and $g_0$ is invertible. For $v$ a row of $g_0^{-1}$, $vg_0$ is a standard basis vector, and
the identity gives $\sum_{a>p_0}|(vR)_a|^2\le1$. So every entry of $g_0^{-1}R$ has absolute value at most $1$; its
first $p_0$ columns form the identity matrix.

\emph{A lower bound for the minor.} As the determinant is continuous, there is $\eta>0$ such that
$|\det N|\geq\frac12$ for every complex $p_0\times p_0$ matrix $N$ whose entries differ from those of the identity
matrix by less than $\eta$. Put $\varrho=\eta/(r+1)$. Let $g$ be as above and $N$ an $r\times r$ matrix whose entries
differ from those of the identity matrix by at most $\varrho$, and write $M'$ for the matrix of the first $p_0$
columns of a matrix $M$. As $R1'=g_0$, the leading $p_0\times p_0$ block of $gN$ is
$RN'=g_0\big(1+g_0^{-1}R(N-1)'\big)$. Each entry of $g_0^{-1}R(N-1)'$ is a sum of $r$ products of an entry of $g_0^{-1}R$
with an entry of $N-1$, so its absolute value is at most $r\varrho<\eta$. Hence
$|\ell_{p_0}(gN)|\geq\frac12|\ell_{p_0}(g)|$.

\emph{The bound.} For an orthonormal point $x$ the Gram matrix of the columns of $c_x(z)_\tau=x_\tau\ell_\tau(z)$ is
$\ell_\tau(z)^{\mathsf T}J\overline{\ell_\tau(z)}$, which does not depend on $x$; so
$d_{p_0}(c_x(z),\tau)=d_{p_0}(c_{x_0}(z),\tau)$. Let $B\geq0$ bound the continuous functions
$z\mapsto|d_{p_0}(c_{x_0}(z),\tau)|$ on the closed ball $|z|\le\varrho$, and put $C=(4B)^{|\Psi|}$. Let $x$ be
orthonormal, $|z|<\varrho$ and $\tau\in\Psi$, and put $g=x_{0,\tau}^{-1}x_\tau$. By
Lemma~\ref{is-onto-step-free-ball-minor}(2), $g^*Jg=J$; as $J^2=1$, $g^{-1}=Jg^*J$ and $gJg^*=J$. The entries of
$\ell_\tau(z)-1$ are coordinates of $z$ or $0$, so their absolute values are at most $|z|<\varrho$; and
$x_{0,\tau}^{-1}c_x(z)_\tau=g\ell_\tau(z)$. By the previous step,
$|\ell_{p_0}(x_{0,\tau}^{-1}c_x(z)_\tau)|\geq\frac12|\ell_{p_0}(g)|$, and $\ell_{p_0}(g)\ne0$ by the first step. As $x$ is
orthonormal, $|d_{p_0}(x,\tau)|=1$. So the factor of $q_{x_0}(c_x(z))$ at $\tau$ is at most
$B/(\frac12|\ell_{p_0}(g)|)^2$, which is $4B$ times the factor of $q_{x_0}(x)$ at $\tau$. Multiplying over $\tau$ gives
the claim.
% Lean: minorRatio_near: exists_block_det_ge (exists_det_near_one for eta, det_block_ne_zero and
% norm_entry_le_one through signed_sum_rows and neg_sum_rows, with delta / (r + 1) for rho), isometry_of_orth
% (mul_sgnMat_mul_conjTranspose), dmin_cpt_eq, the bound B on the closed ball (exists_bound_of_continuousOn),
% norm_nmat_sub_one_le and IC.norm_dmin_orth; the constant is the product over tau of 4 B.
\end{proof}

\begin{lemma}\label{is-onto-step-orbit-interpolation-large-weight-power}
Let $X$ consist of the data (1) to (6) of flag data of rank $r$ over $E$, use the notation of
Definition~\ref{is-onto-share-counts-section-levi}, write $|z|$ for the Euclidean norm of $z\in\mathbf C^n$, and let
$x_0$ be an orthonormal point and $\nu>0$ a real number. Let $q\colon M_r(\mathbf C)^\Psi\to[0,\infty)$ be a measurable
function such that $0<q\le1$ on $\mathcal D_X$, $q(xb)=q(x)$ for $x\in\mathcal D_X$ and $b\in\mathcal B^\Psi$, and
$q(c_{x_0}(z))=e^{-\rho_{x_0}(z)}$ for $z\in\Omega(x_0)$. Suppose that there are real numbers $\varrho>0$, $C$ and
$C_s$ such that $q(c_x(z))\le C\,q(x)$ for every orthonormal point $x$ and every $z$ with $|z|<\varrho$, and
$\sum_{\gamma\in\Gamma}q(\gamma x)^{\nu/2}\le C_s$ for every orthonormal point $x$. Then the function $\kappa=q^\nu$ is
measurable, and:
\begin{enumerate}
\item $0<\kappa\le1$ on $\mathcal D_X$, and $\kappa(xb)=\kappa(x)$ for $x\in\mathcal D_X$ and $b\in\mathcal B^\Psi$;
\item $\kappa(c_{x_0}(z))=e^{-\nu\rho_{x_0}(z)}$ for $z\in\Omega(x_0)$;
\item $\sum_{\gamma\in\Gamma}\kappa(\gamma x)^{1/2}\le C_s$ for every $x\in\mathcal D_X$;
\item $\kappa(c_x(z))\le\max(C,0)^\nu\kappa(x)$ for every orthonormal point $x$ and every $z$ with $|z|<\varrho$.
\end{enumerate}
\end{lemma}
\leanref{Development/AbelianVarietyTheory/Proofs/TorusHaarMeasure/KaFrameOfDelta.lean}{GlobalTheory.kaFrame_of_delta}
\begin{proof}
The function $\kappa$ is measurable as $q$ is, and (1) and (2) follow from the properties of $q$.

(3) Let $x\in\mathcal D_X$. By Lemma~\ref{is-onto-share-majorant}(3), $x=x_1b$ with $x_1$ orthonormal and
$b\in\mathcal B^\Psi$. For $\gamma\in\Gamma$, $\gamma x=(\gamma x_1)b$, and $\gamma x_1\in\mathcal D_X$, as $\gamma$ preserves
$h$ and so $\gamma x_1$ has the Gram matrices of $x_1$. So $\kappa(\gamma x)^{1/2}=q(\gamma x_1)^{\nu/2}$, and summing over
$\gamma$ gives (3).

(4) $q(c_x(z))\le\max(C,0)\,q(x)$, and $t\mapsto t^\nu$ is increasing on $[0,\infty)$.
% Lean: kaFrame_of_delta, with delta for q and bb for nu: ortho_holds, lact_rmul_torus_haar_measure
% (lact gamma (rmul y b) = rmul (lact gamma y) b), OntoS.dom_lact, Real.rpow_mul for sqrt(q^bb) = q^(bb / 2),
% and Real.rpow_le_rpow with Real.mul_rpow for (4).
\end{proof}

\begin{lemma}\label{is-onto-step-orbit-interpolation-large-weight}
Let $X$ consist of the data (1) to (6) of flag data of rank $r$ over $E$, such that $\Gamma$ acts properly
discontinuously on $\mathcal D_X$ modulo $\mathcal B^\Psi$ (Lemma~\ref{is-onto-share-counts-section-discontinuous}), use
the notation of Definition~\ref{is-onto-share-counts-section-levi}, write $|z|$ for the Euclidean norm of
$z\in\mathbf C^n$, and let $x_0$ be an orthonormal point and $\nu>4|\Psi|r^2$ a real number. Then there is a measurable
function $\kappa\colon M_r(\mathbf C)^\Psi\to\mathbf R$ such that:
\begin{enumerate}
\item $0<\kappa\le1$ on $\mathcal D_X$, and $\kappa(xb)=\kappa(x)$ for $x\in\mathcal D_X$ and $b\in\mathcal B^\Psi$;
\item $\kappa(c_{x_0}(z))=e^{-\nu\rho_{x_0}(z)}$ for $z\in\Omega(x_0)$;
\item there is a real number $C_s$ with $\sum_{\gamma\in\Gamma}\kappa(\gamma x)^{1/2}\le C_s$ for every $x\in\mathcal D_X$;
\item there are real numbers $\varrho>0$ and $C$ with $\kappa(c_x(z))\le C\kappa(x)$ for every orthonormal point $x$
  and every $z$ with $|z|<\varrho$.
\end{enumerate}
\end{lemma}
\leanref{Development/AbelianVarietyTheory/Proofs/TorusHaarMeasure/KaFrameOfDelta.lean}{GlobalTheory.kaFrame}
\leanref{Development/AbelianVarietyTheory/Proofs/TorusHaarMeasure/OneLeMinorOfIsOrth.lean}{GlobalTheory.FramePotential.frame_potential_basic}
\leanref{Development/AbelianVarietyTheory/Proofs/TorusHaarMeasure/IsUpqSubmatrix.lean}{GlobalTheory.FrameSumBound.frame_sum_bound}
\leanref{Development/AbelianVarietyTheory/Proofs/TorusHaarMeasure/ExistsLeNormDev.lean}{GlobalTheory.grpSep}
\begin{proof}
Let $q=q_{x_0}$ be the function of Lemma~\ref{is-onto-step-orbit-interpolation-large-weight-near}. It satisfies the
hypotheses of Lemma~\ref{is-onto-step-orbit-interpolation-large-weight-power} for $\nu$, which then gives
$\kappa=q^\nu$.

\emph{The values of $q$.} The function $q$ is a product of quotients of continuous functions, so it is measurable.
For $x\in\mathcal D_X$, $q(x)=\prod_\tau q_\tau(x_0,x)$ in the notation of
Lemma~\ref{is-onto-step-free-ball-minor}(3); so $0<q(x)\le1$ and $q(xb)=q(x)$ for $b\in\mathcal B^\Psi$, by that lemma.
For $z\in\Omega(x_0)$, the matrix $x_{0,\tau}^{-1}c_{x_0}(z)_\tau=\ell_\tau(z)$ has leading minors $1$, and
$d_{p_0}(c_{x_0}(z),\tau)\ne0$; so $q(c_{x_0}(z))=\prod_\tau|d_{p_0}(c_{x_0}(z),\tau)|=e^{-\rho_{x_0}(z)}$. The local bound is
Lemma~\ref{is-onto-step-orbit-interpolation-large-weight-near}.

\emph{The sum.} Let $c>0$ be the number of Lemma~\ref{is-onto-share-counts-section-poincare-count-separation} for
$x_0$, and let $J$ be as in Lemma~\ref{is-onto-step-free-ball-minor}. Let $x$ be orthonormal. For $\gamma\in\Gamma$ the
point $\gamma x$ is orthonormal, as $\gamma$ preserves $h$. Put $g_{\gamma,\tau}=x_{0,\tau}^{-1}(\gamma x)_\tau$, and let
$M_{\gamma,\tau}$ be $g_{\gamma,\tau}$ with its rows and columns put in the order $p_0+1,\dots,r,1,\dots,p_0$. By
Lemma~\ref{is-onto-step-free-ball-minor}(2), $g_{\gamma,\tau}^*Jg_{\gamma,\tau}=J$. The reordering takes $J$ to
$-J_{r-p_0,p_0}$, in the notation of Lemma~\ref{is-onto-share-counts-section-poincare-count-sum}, so
$M_{\gamma,\tau}^*J_{r-p_0,p_0}M_{\gamma,\tau}=J_{r-p_0,p_0}$. For $\gamma\ne\gamma'$ in $\Gamma$,
\[
 g_{\gamma',\tau}\,g_{\gamma,\tau}^{-1}=x_{0,\tau}^{-1}(\gamma'\gamma^{-1})^\tau x_{0,\tau},
\]
and $M_{\gamma',\tau}M_{\gamma,\tau}^{-1}$ is this matrix reordered; by
Lemma~\ref{is-onto-share-counts-section-poincare-count-separation} for $\gamma'\gamma^{-1}\ne1$, for some $\tau$ it differs
from $1$ by at least $c$ in some entry. The lower right $p_0\times p_0$ block $D(M_{\gamma,\tau})$ of $M_{\gamma,\tau}$ is the
leading block of $g_{\gamma,\tau}$, and $|d_{p_0}(\gamma x,\tau)|=1$; so
\[
 q(\gamma x)^{\nu/2}=\prod_{\tau\in\Psi}|\ell_{p_0}(g_{\gamma,\tau})|^{-\nu}=\prod_{\tau\in\Psi}|\det D(M_{\gamma,\tau})|^{-\nu}.
\]
Lemma~\ref{is-onto-share-counts-section-poincare-count-sum}, for $p=r-p_0$, $q=p_0$, $f=|\Psi|$, $s=\nu/2>2|\Psi|r^2$,
the number $c$, $I=\Gamma$ and the family $(M_{\gamma,\tau})$, gives a number $C_s$, which does not depend on $x$, with
$\sum_\gamma q(\gamma x)^{\nu/2}\le C_s$.
% Lean: kaFrame = kaFrame_of_delta applied to frame_potential_basic (measurability, factor_mem, factor_rmul,
% and lminor_cpt for the chart), minorRatio_near (Lemma is-onto-step-orbit-interpolation-large-weight-near), and
% frame_sum_bound grpSep with s = bb / 2. frame_sum_bound: hsep gives c (grpSep, whose statement is that of Lemma
% is-onto-share-counts-section-poincare-count-separation), sumBd_sum (Lemma
% is-onto-share-counts-section-poincare-count-sum) with p = r - p0, q = p0, f = |Psi|, applied to the matrices
% (mat x0 tau)^-1 * mat (lact gamma x) tau reordered by swapIdx (isUpq_frame through isUpq_submatrix and
% sgn_submatrix for M^H Jf M = Jf; quot_eq for the separation), and IC.norm_dmin_orth for |d_p0| = 1, with the
% lower right block of the reordered matrix equal to the leading block (hm). Unlike PoinPf.Mq in the proof of Lemma
% is-onto-share-counts-section-poincare-count(3), the matrices are not conjugated.
\end{proof}

\begin{definition}\label{is-onto-step-orbit-interpolation-large-pole}
Let $X$ consist of the data (1) to (6) of flag data of rank $r$ over $E$, use the notation of
Definitions~\ref{is-onto-share-counts-section-levi} and~\ref{is-onto-step-free-ball}, and put
$\dim X=|\Psi|p_0(r-p_0)$. Let $G\subset GL_r(E)$ be a subgroup, $x_0$ an orthonormal point, and
$\kappa,\varepsilon,C,\delta$ real numbers with $\delta>0$. A \emph{pole function} for $(G,x_0,\kappa,\varepsilon,C,\delta)$
is a function $\psi\colon M_r(\mathbf C)^\Psi\to\mathbf R$ such that:
\begin{enumerate}
\item $\psi(\gamma xb)=\psi(x)$ for $\gamma\in G$, $x\in\mathcal D_X$ and $b\in\mathcal B^\Psi$;
\item $\psi\le0$ on $\mathcal D_X$;
\item for every orthonormal point $x_1$, the function $\psi\circ c_{x_1}$ is smooth on $\Omega(x_1)$, and
  $\mathcal L_{\psi\circ c_{x_1}}(z;\xi)\geq-\kappa\,\mathcal L_{\rho_{x_1}}(z;\xi)$ for $z\in\Omega(x_1)$ and
  $\xi\in\mathbf C^n$;
\item $\psi(x)\le\dim X\cdot\log\big(\rho(x_0,x)+\delta\big)+C$ for $x\in\mathcal D_X$ with $\rho(x_0,x)\le\varepsilon$;
\item $\psi(x)\geq-C$ for $x\in\mathcal D_X$ with $\varepsilon/2\le\rho(x_0,x)\le\varepsilon$.
\end{enumerate}
\end{definition}
\leanref{Development/AbelianVarietyTheory/Definitions/UpperBounds/FreeBallIffOfDimXEqZero.lean}{GlobalTheory.Up.IsPole}
\leanref{Development/AbelianVarietyTheory/Definitions/UpperBounds/FreeBallIffOfDimXEqZero.lean}{GlobalTheory.Up.dimX}

\begin{lemma}\label{is-onto-step-orbit-interpolation-large-pole-exists-cutoff}
Let $\rho_0>0$ and $\theta>1$ be real numbers. Then there are real numbers $\varepsilon$ and $C$ with
$0<\varepsilon<\rho_0$ such that for every real number $0<\delta\le1$ there is a smooth function $g$ on $(-\delta,\infty)$
with:
\begin{enumerate}
\item $g(u)=0$ for $u\geq\rho_0-\varepsilon$;
\item $g(u)\le0$, $g'(u)\geq0$ and $g'(u)+u\,g''(u)\geq-\theta/\rho_0$ for $u\geq0$;
\item $g(u)\le\log(u+\delta)+C$ for $0\le u\le\varepsilon$;
\item $g(u)\geq-C$ for $u\geq\varepsilon/2$.
\end{enumerate}
\end{lemma}
\leanref{Development/AbelianVarietyTheory/Proofs/TorusHaarMeasure/ExistsSmoothStepSharpSlope.lean}{GlobalTheory.exists_cutoff_log_sharp}
\leanref{Development/AbelianVarietyTheory/Proofs/TorusHaarMeasure/ExistsSmoothStepSharpSlope.lean}{GlobalTheory.exists_smooth_step_sharp_slope}
\begin{proof}
Put $\varepsilon=\rho_0(\theta-1)/(8\theta)$ and $C=2\rho_0/\varepsilon+|\log\varepsilon|$. Then $0<4\varepsilon<\rho_0$, and as
$4\theta\varepsilon\le\rho_0(\theta-1)$, $\rho_0\le\theta(\rho_0-4\varepsilon)$, that is, $1/(\rho_0-4\varepsilon)\le\theta/\rho_0$.

\emph{A step.} Let $\sigma$ be a smooth function on $\mathbf R$ with values in $[0,1]$ that vanishes for $t\le\varepsilon$
and for $t\geq\rho_0-\varepsilon$ and equals $1$ on $[2\varepsilon,\rho_0-2\varepsilon]$. Put $F(u)=\int_\varepsilon^u\sigma(t)\,dt$,
$I=F(\rho_0-\varepsilon)$ and $\eta=1-F/I$. As $0\le\sigma\le1$ and $\sigma=1$ on $[2\varepsilon,\rho_0-2\varepsilon]$,
$I\geq\rho_0-4\varepsilon>0$; $F(u)=0$ for $u\le\varepsilon$, $F(u)=I$ for $u\geq\rho_0-\varepsilon$, and $0\le F\le I$. So
$\eta$ is smooth with values in $[0,1]$, $\eta(u)=1$ for $u\le\varepsilon$, $\eta(u)=0$ for $u\geq\rho_0-\varepsilon$, and
$\eta'=-\sigma/I\geq-1/(\rho_0-4\varepsilon)$.

\emph{The function.} Let $0<\delta\le1$. On $(-\delta,\infty)$ put $f(t)=\eta(t)/(t+\delta)$, a smooth function with
$f\geq0$, and $g(u)=\int_{\rho_0}^uf(t)\,dt$. Then $g$ is smooth, $g'=f$ and $g''=f'$.

(1) As $f=0$ on $[\rho_0-\varepsilon,\infty)$, $g(u)=0$ for $u\geq\rho_0-\varepsilon$.

(2) Let $u\geq0$. If $u\le\rho_0$, then $g(u)=-\int_u^{\rho_0}f\le0$, and otherwise $g(u)=0$ by (1). Next $g'(u)=f(u)\geq0$, and
\[
 g'(u)+u\,g''(u)=\frac{\eta(u)}{u+\delta}+u\,\frac{\eta'(u)(u+\delta)-\eta(u)}{(u+\delta)^2}
 =\frac{\eta(u)\,\delta}{(u+\delta)^2}+\frac{u\,\eta'(u)}{u+\delta}\geq-\frac u{(\rho_0-4\varepsilon)(u+\delta)}
 \geq-\frac\theta{\rho_0}.
\]

(3) Let $0\le u\le\varepsilon$. As $\eta=1$ on $[u,\varepsilon]$,
$g(u)=g(\varepsilon)-\int_u^\varepsilon\frac{dt}{t+\delta}=g(\varepsilon)-\log(\varepsilon+\delta)+\log(u+\delta)$. By (2),
$g(\varepsilon)\le0$, and $\log\varepsilon\le\log(\varepsilon+\delta)$; so $g(u)\le\log(u+\delta)-\log\varepsilon\le\log(u+\delta)+C$.

(4) Let $u\geq\varepsilon/2$. If $u\geq\rho_0$, then $g(u)=0\geq-C$. Otherwise $f(t)\le1/(t+\delta)\le2/\varepsilon$ for
$t\geq\varepsilon/2$, so $g(u)=-\int_u^{\rho_0}f\geq-2(\rho_0-u)/\varepsilon\geq-2\rho_0/\varepsilon\geq-C$.
% Lean: exists_cutoff_log_sharp: eps from 8 theta eps = rho0 (theta - 1), hslope, the step eta of
% exists_smooth_step_sharp_slope (sigma is the ContDiffBump at rho0 / 2 with radii rho0 / 2 - 2 eps and
% rho0 / 2 - eps; hIge for I >= rho0 - 4 eps), f = eta / (t + delta), g = the integral from rho0 (hd, hzero, hle0),
% the derivative computation key, hsplit and hlog for (3), hb for (4), and C = 2 rho0 / eps + |log eps|. The Lean's
% g is defined on R, smooth on Ioi (-delta); only its values on [0, infinity) are used.
\end{proof}

\begin{lemma}\label{is-onto-step-orbit-interpolation-large-pole-exists-block}
Let $0\le p\le r$ be integers, $J=\operatorname{diag}(1,\dots,1,-1,\dots,-1)$ with $p$ entries $1$, and write $M^*$ for the
conjugate transpose of a complex matrix $M$ and $\ell_p(N)$ for the leading principal $p\times p$ minor of a square
matrix $N$. Let $c<1$ be a real number. Then there is a real number $\eta>0$ such that
\[
 c\,|\ell_p(g)|\le|\ell_p(gn)|\le2\,|\ell_p(g)|
\]
for all $g,n\in M_r(\mathbf C)$ with $gJg^*=J$ and $|n_{ab}-[a=b]|\le\eta$ for all $a,b$, where $[a=b]$ is $1$ if $a=b$
and $0$ otherwise.
\end{lemma}
\leanref{Development/AbelianVarietyTheory/Proofs/DivisorProofs/ExistsBlockDetGe.lean}{GlobalTheory.PoleExistence.exists_block_det_ge'}
\leanref{Development/AbelianVarietyTheory/Proofs/DivisorProofs/OrbitSumExpPotentialUniformBound.lean}{GlobalTheory.OrbitSum.exists_block_det_le}
\begin{proof}
Let $g$ satisfy $gJg^*=J$, let $g_0$ be its leading $p\times p$ block and $R$ the $p\times r$ matrix of its first $p$
rows. The first step of the proof of Lemma~\ref{is-onto-step-orbit-interpolation-large-weight-near}, which uses only
$gJg^*=J$, shows that $g_0$ is invertible and that every entry of $g_0^{-1}R$ has absolute value at most $1$.

As the determinant is continuous, there is $\eta'>0$ such that $c\le|\det N|\le2$ for every complex $p\times p$ matrix $N$
whose entries differ from those of the identity matrix by less than $\eta'$. Put $\eta=\eta'/(r+1)$. Let $n$ be as in the
statement, and write $M'$ for the matrix of the first $p$ columns of a matrix $M$. The leading $p\times p$ block of $gn$ is
$Rn'=g_0\big(1+g_0^{-1}R(n-1)'\big)$. Each entry of $g_0^{-1}R(n-1)'$ is a sum of $r$ products of an entry of $g_0^{-1}R$
with an entry of $n-1$, so its absolute value is at most $r\eta<\eta'$. Hence
$\ell_p(gn)=\ell_p(g)\det\big(1+g_0^{-1}R(n-1)'\big)$ with $c\le|\det(1+g_0^{-1}R(n-1)')|\le2$.
% Lean: exists_block_det_ge' (exists_det_ge for c) and OrbitSum.exists_block_det_le (exists_det_le_two), two lemmas
% with the same proof: det_block_ne_zero and norm_entry_le_one (MinorRatioNear), e1 to e3, and the bound by
% delta / (r + 1). The book takes the smaller of the two numbers.
\end{proof}

\begin{lemma}\label{is-onto-step-orbit-interpolation-large-pole-exists-translate}
Let $X$ consist of the data (1) to (6) of flag data of rank $r$ over $E$, and use the notation of
Definition~\ref{is-onto-step-free-ball}.
\begin{enumerate}
\item Let $\gamma\in GL_r(E)$ preserve $h$, let $y_0$ be an orthonormal point and $y\in\mathcal D_X$. Then $\gamma^{-1}y_0$ is
  orthonormal, $\gamma y\in\mathcal D_X$, and $\rho(y_0,\gamma y)=\rho(\gamma^{-1}y_0,y)$.
\item Let $G\subset GL_r(E)$ be a subgroup with a free ball of radius $\rho_0$ at an orthonormal point $x_0$, let
  $x\in\mathcal D_X$, and let $\gamma,\gamma'\in G$ with $\rho(x_0,\gamma x)<\rho_0$ and $\rho(x_0,\gamma'x)<\rho_0$. Then
  $\gamma=\gamma'$.
\end{enumerate}
\end{lemma}
\leanref{Development/AbelianVarietyTheory/Definitions/UpperBounds/FreeBallIffOfDimXEqZero.lean}{GlobalTheory.Up.rhX_lact}
\leanref{Development/AbelianVarietyTheory/Definitions/UpperBounds/FreeBallIffOfDimXEqZero.lean}{GlobalTheory.Up.rhX_lact_lact}
\leanref{Development/AbelianVarietyTheory/Definitions/UpperBounds/FreeBallIffOfDimXEqZero.lean}{GlobalTheory.Up.FreeBall.unique}
\begin{proof}
(1) As $\gamma$ preserves $h$, the Gram matrix of the columns of $(\gamma w)_\tau=\gamma^\tau w_\tau$ is that of $w_\tau$, for
every $w\in M_r(\mathbf C)^\Psi$ and $\tau\in\Psi$. So $\gamma^{-1}y_0$ is orthonormal, $\gamma y\in\mathcal D_X$ and
$d_{p_0}(\gamma y,\tau)=d_{p_0}(y,\tau)$. Moreover $(\gamma^{-1}y_0)_\tau^{-1}y_\tau=y_{0,\tau}^{-1}\gamma^\tau y_\tau=y_{0,\tau}^{-1}(\gamma y)_\tau$.
So the quotients of Lemma~\ref{is-onto-step-free-ball-minor}(3) satisfy $q_\tau(y_0,\gamma y)=q_\tau(\gamma^{-1}y_0,y)$, and
$\rho(y_0,\gamma y)=\rho(\gamma^{-1}y_0,y)$.

(2) The point $\gamma x$ lies in $\mathcal D_X$ by (1), as the elements of $G$ preserve $h$. If $\gamma'\gamma^{-1}\ne1$,
then, as $\rho(x_0,\gamma x)<\rho_0$, the free ball gives $\rho(x_0,\gamma'x)=\rho(x_0,\gamma'\gamma^{-1}\cdot\gamma x)\geq\rho_0$,
a contradiction. So $\gamma'=\gamma$.
% Lean: (1) is rhX_lact (through rhX_lact_lact: MajInv.gram_lact for the minors and the cancellation of the
% matrices of gamma) with isOrth_lact_of_isUStab and MajInv.lact_mem_dom. (2) is FreeBall.unique, through
% FreeBall.eq_one applied to gamma' gamma^-1 at the point gamma x.
\end{proof}

\begin{lemma}\label{is-onto-step-orbit-interpolation-large-pole-exists-lower}
Let $X$ consist of the data (1) to (6) of flag data of rank $r$ over $E$, use the notation of
Definition~\ref{is-onto-step-free-ball}, and let $x\in\mathcal D_X$ and $t>0$. Then $x$ has a neighbourhood $V$ in
$M_r(\mathbf C)^\Psi$ such that $\rho(x_2,y)\geq\rho(x_2,x)-t$ for every $y\in V\cap\mathcal D_X$ and every orthonormal
point $x_2$.
\end{lemma}
\leanref{Development/AbelianVarietyTheory/Proofs/DivisorProofs/ExistsBlockDetGe.lean}{GlobalTheory.PoleExistence.rhX_lower}
\leanref{Development/AbelianVarietyTheory/Proofs/DivisorProofs/ExistsBlockDetGe.lean}{GlobalTheory.PoleExistence.rhX_lower_orth}
\begin{proof}
Write $\ell_{p_0}(N)$ for the leading principal $p_0\times p_0$ minor of a square matrix $N$, and let $J$ be as in
Lemma~\ref{is-onto-step-free-ball-minor}.

\emph{An orthonormal point.} Suppose first that $x$ is orthonormal. Put $k=|\Psi|$ and $u=t/(3k+1)$, and let $\eta$ be the
number of Lemma~\ref{is-onto-step-orbit-interpolation-large-pole-exists-block} for $p=p_0$ and $c=e^{-u}$. As $x$ is
orthonormal, $|d_{p_0}(x,\tau)|=1$ for every $\tau$. By continuity there is a neighbourhood $V$ of $x$ such that every
$y\in V$ satisfies $|d_{p_0}(y,\tau)|<e^u$ and $|(x_\tau^{-1}y_\tau)_{ab}-[a=b]|<\eta$ for all $\tau$, $a$ and $b$. Let
$y\in V\cap\mathcal D_X$, let $x_2$ be orthonormal, and for $\tau\in\Psi$ put $g_\tau=x_{2,\tau}^{-1}x_\tau$. By
Lemma~\ref{is-onto-step-free-ball-minor}(2), $g_\tau^*Jg_\tau=J$, so $g_\tau Jg_\tau^*=J$ as $J^2=1$. As
$x_{2,\tau}^{-1}y_\tau=g_\tau\,(x_\tau^{-1}y_\tau)$, Lemma~\ref{is-onto-step-orbit-interpolation-large-pole-exists-block} gives
$|\ell_{p_0}(x_{2,\tau}^{-1}y_\tau)|\geq e^{-u}|\ell_{p_0}(g_\tau)|$. By Definition~\ref{is-onto-step-free-ball} and
Lemma~\ref{is-onto-step-free-ball-minor}(3),
\[
 \rho(x_2,y)=-\sum_\tau\log|d_{p_0}(y,\tau)|+\sum_\tau\log|\ell_{p_0}(x_{2,\tau}^{-1}y_\tau)|^2
 \geq-ku+\sum_\tau\big(\log|\ell_{p_0}(g_\tau)|^2-2u\big),
\]
and the same formula for $x$, where $\log|d_{p_0}(x,\tau)|=0$, gives $\rho(x_2,x)=\sum_\tau\log|\ell_{p_0}(g_\tau)|^2$. So
$\rho(x_2,y)\geq\rho(x_2,x)-3ku\geq\rho(x_2,x)-t$.

\emph{The general case.} By Lemma~\ref{is-onto-share-majorant}(3), $x=x_1b$ with $x_1$ orthonormal and
$b\in\mathcal B^\Psi$. Let $V_1$ be the neighbourhood of $x_1$ for $t$ of the first step, and $V$ the set of the $y$ with
$yb^{-1}\in V_1$, a neighbourhood of $x$. Let $y\in V\cap\mathcal D_X$. Then $yb^{-1}\in\mathcal D_X$ by
Lemma~\ref{is-onto-share-majorant}(2), and by Lemma~\ref{is-onto-step-free-ball-bound}(1)
$\rho(x_2,y)=\rho(x_2,yb^{-1})\geq\rho(x_2,x_1)-t=\rho(x_2,x)-t$.
% Lean: rhX_lower_orth: u = t / (3 k + 1), exists_block_det_ge' with c = exp(-u), ha (norm_dmin_of_isOrth and
% the continuity of dmin) and hb (the entries of (mat x1)^-1 mat y - 1), rhX_eq_sum, s1, s2, s3 with
% isometry_of_orth. rhX_lower: ortho_holds, the continuity of y -> rmul y b^-1, OntoS.dom_rmul and rhX_rmul.
\end{proof}

\begin{definition}\label{is-onto-step-orbit-interpolation-large-pole-exists-sharp}
Write $\partial_j$ for $\partial/\partial z_j$ on $\mathbf C^n$ and $\mathcal L_\phi$ for the Levi form of a real function
$\phi$ of class $C^2$ (Definition~\ref{is-onto-share-counts-model}). Let $z,\xi\in\mathbf C^n$. A real function $\phi$ of
class $C^2$ near $z$ has a \emph{sharp gradient} at $z$ in the direction $\xi$ if $\phi(z)\geq0$,
$\mathcal L_\phi(z;\xi)\geq0$ and
\[
 \Big|\sum_j\partial_j\phi(z)\,\xi_j\Big|^2\le\phi(z)\,\mathcal L_\phi(z;\xi).
\]
\end{definition}
\leanref{Development/AbelianVarietyTheory/Proofs/TorusHaarMeasure/ContinuousInvLeadingBlockEntry.lean}{GlobalTheory.SharpGradient.SharpAt}

\begin{lemma}\label{is-onto-step-orbit-interpolation-large-pole-exists-sharp-rules}
Let $z,\xi\in\mathbf C^n$, and use the notation of
Definition~\ref{is-onto-step-orbit-interpolation-large-pole-exists-sharp}.
\begin{enumerate}
\item Let $v$ be a real function of class $C^2$ near $z$ and $k$ a real function of class $C^2$ near $v(z)$. Then
  \[
   \mathcal L_{k\circ v}(z;\xi)=k'(v(z))\,\mathcal L_v(z;\xi)+k''(v(z))\,\Big|\sum_j\partial_jv(z)\,\xi_j\Big|^2.
  \]
\item Let $g_1,\dots,g_m$ be holomorphic functions near $z$ with $\sum_i|g_i(z)|^2<1$. Then
  $-\log\big(1-\sum_i|g_i|^2\big)$ has a sharp gradient at $z$ in the direction $\xi$.
\item A finite sum of functions with sharp gradients at $z$ in the direction $\xi$ has a sharp gradient at $z$ in the
  direction $\xi$.
\item Let $\psi$ have a sharp gradient at $z$ in the direction $\xi$, and let $\phi$ be a real function of class $C^2$
  near $z$ with $\psi\le\phi$ near $z$ and $\psi(z)=\phi(z)$. Then $\phi$ has a sharp gradient at $z$ in the direction
  $\xi$.
\end{enumerate}
\end{lemma}
\leanref{ReciprocityTheory/L2DbarEstimates/SkodaWeights/LeviVA.lean}{ReciprocityTheory.L2DbarEstimates.H.leviChain}
\leanref{Development/AbelianVarietyTheory/Proofs/TorusHaarMeasure/ContinuousInvLeadingBlockEntry.lean}{GlobalTheory.SharpGradient.sharpAt_ball}
\leanref{Development/AbelianVarietyTheory/Proofs/TorusHaarMeasure/ContinuousInvLeadingBlockEntry.lean}{GlobalTheory.SharpGradient.sharpAt_sum}
\leanref{Development/AbelianVarietyTheory/Proofs/TorusHaarMeasure/ContinuousInvLeadingBlockEntry.lean}{GlobalTheory.SharpGradient.SharpAt.of_minorant}
\begin{proof}
Write $\bar\partial_j=\partial/\partial\bar z_j$.

(1) As $v$ is real, $\bar\partial_kv=\overline{\partial_kv}$. By the chain rule
$\partial_j\bar\partial_k(k\circ v)=k'(v)\,\partial_j\bar\partial_kv+k''(v)\,\partial_jv\,\overline{\partial_kv}$; multiply by
$\xi_j\bar\xi_k$ and sum over $j$ and $k$.

(2) Put $S=\sum_i|g_i|^2$, $t=S(z)\in[0,1)$, $\eta_i=\sum_j\partial_jg_i(z)\xi_j$ and $\phi=-\log(1-S)$. As the $g_i$ are
holomorphic, $\partial_jS=\sum_i\partial_jg_i\,\bar g_i$ and $\partial_j\bar\partial_kS=\sum_i\partial_jg_i\,\overline{\partial_kg_i}$.
So $\sum_j\partial_jS(z)\xi_j=a$ with $a=\sum_i\eta_i\overline{g_i(z)}$, and $\mathcal L_S(z;\xi)=M$ with $M=\sum_i|\eta_i|^2$; by
the Cauchy-Schwarz inequality, $N=|a|^2\le Mt$. With $k(u)=-\log(1-u)$, $k'(u)=(1-u)^{-1}$ and $k''(u)=(1-u)^{-2}$, (1)
gives $\mathcal L_\phi(z;\xi)=M/(1-t)+N/(1-t)^2\geq0$, and $\sum_j\partial_j\phi(z)\xi_j=a/(1-t)$. Also
$\phi(z)=-\log(1-t)\geq t\geq0$. Hence
\[
 \phi(z)\,\mathcal L_\phi(z;\xi)\geq\frac{tM}{1-t}+\frac{tN}{(1-t)^2}\geq\frac N{1-t}+\frac{tN}{(1-t)^2}=\frac N{(1-t)^2}
 =\Big|\sum_j\partial_j\phi(z)\,\xi_j\Big|^2.
\]

(3) Let $\phi_1,\dots,\phi_m$ have sharp gradients at $z$ in the direction $\xi$, and put $s_i=\phi_i(z)$,
$L_i=\mathcal L_{\phi_i}(z;\xi)$ and $a_i=\sum_j\partial_j\phi_i(z)\xi_j$, so that $|a_i|^2\le s_iL_i$. The sum $\sum_i\phi_i$ has
the value $\sum_is_i\geq0$ and the Levi form $\sum_iL_i\geq0$ at $z$, and by the Cauchy-Schwarz inequality
$|\sum_ia_i|^2\le(\sum_i(s_iL_i)^{1/2})^2\le(\sum_is_i)(\sum_iL_i)$.

(4) The function $\phi-\psi$ has a local minimum at $z$. So the first derivatives of $\phi$ and $\psi$ at $z$ agree, and
$\psi-\phi$, which has a local maximum at $z$, has a Levi form $\le0$ there, as in the proof of
Lemma~\ref{is-onto-share-counts-section-kronecker-solve-chart}; that is,
$\mathcal L_\psi(z;\xi)\le\mathcal L_\phi(z;\xi)$. Hence $\phi(z)=\psi(z)\geq0$, $\mathcal L_\phi(z;\xi)\geq0$, and
\[
 \Big|\sum_j\partial_j\phi(z)\,\xi_j\Big|^2=\Big|\sum_j\partial_j\psi(z)\,\xi_j\Big|^2\le\psi(z)\,\mathcal L_\psi(z;\xi)
 \le\phi(z)\,\mathcal L_\phi(z;\xi).
\]
% Lean: (1) is H.leviChain (LeviChainStatement), with wd and wdb the Wirtinger derivatives. (2) is sharpAt_ball:
% k t = -log (1 - t), hd1, hd2, hw, hNM (normSq_sum_le), hL, hG and hlog. (3) is sharpAt_sum and sharpAt_sum'
% (normSq_sum_le through Finset.sum_sq_le_sum_mul_sum_of_sq_le_mul). (4) is SharpAt.of_minorant (hfd from
% IsLocalMin.fderiv_eq_zero, and leviV_re_nonpos_of_isLocalMax for psi + (-1) phi).
\end{proof}

\begin{lemma}\label{is-onto-step-orbit-interpolation-large-pole-exists-transfer}
Let $X$ consist of the data (1) to (6) of flag data of rank $r$ over $E$, use the notation of
Definitions~\ref{is-onto-share-counts-section-levi} and~\ref{is-onto-step-free-ball}, write $\partial_j$ for
$\partial/\partial z_j$, and let $x_1$ and $x_2$ be orthonormal points. Put $f(z)=\rho(x_2,c_{x_1}(z))$ for
$z\in\Omega(x_1)$. Then $f$ is smooth on $\Omega(x_1)$, and for every $z\in\Omega(x_1)$ and every $\xi\in\mathbf C^n$:
$f(z)\geq0$, $\mathcal L_f(z;\xi)=\mathcal L_{\rho_{x_1}}(z;\xi)$, and
\[
 \Big|\sum_j\partial_jf(z)\,\xi_j\Big|^2\le f(z)\,\mathcal L_{\rho_{x_1}}(z;\xi).
\]
\end{lemma}
\leanref{Development/AbelianVarietyTheory/Definitions/UpperBounds/FreeBallIffOfDimXEqZero.lean}{GlobalTheory.Up.RhTr}
\leanref{Development/AbelianVarietyTheory/Proofs/TorusHaarMeasure/ContDiffAtLogMinorProd.lean}{GlobalTheory.PotentialTransfer.potential_transfer}
\leanref{Development/AbelianVarietyTheory/Proofs/TorusHaarMeasure/ContDiffAtLogMinorProd.lean}{GlobalTheory.PotentialTransfer.rhX_eq_sum_graph}
\leanref{Development/AbelianVarietyTheory/Proofs/TorusHaarMeasure/ContDiffAtLogMinorProd.lean}{GlobalTheory.PotentialTransfer.leviO_rhX}
\leanref{Development/AbelianVarietyTheory/Proofs/TorusHaarMeasure/ContDiffAtLogMinorProd.lean}{GlobalTheory.PotentialTransfer.sharpAt_rhX_cpt}
\leanref{Development/AbelianVarietyTheory/Proofs/TorusHaarMeasure/ContinuousInvLeadingBlockEntry.lean}{GlobalTheory.SharpGradient.exists_unitary_rows}
\leanref{Development/AbelianVarietyTheory/Proofs/TorusHaarMeasure/ContinuousInvLeadingBlockEntry.lean}{GlobalTheory.SharpGradient.rows_le_neg_log_det}
\begin{proof}
Write $\ell_{p_0}(N)$ for the leading principal $p_0\times p_0$ minor of a square matrix $N$, $M^*$ for the conjugate
transpose and $M^{\mathsf T}$ for the transpose of a matrix $M$, and let $J$ be as in
Lemma~\ref{is-onto-step-free-ball-minor}.

\emph{The function.} Let $y\in\Omega(x_1)$. Then $c_{x_1}(y)\in\mathcal D_X$, and by Definition~\ref{is-onto-step-free-ball}
and Lemma~\ref{is-onto-step-free-ball-minor}(3)
\[
 f(y)=-\sum_\tau\log|d_{p_0}(c_{x_1}(y),\tau)|+\log|m(y)|^2=\rho_{x_1}(y)+\log|m(y)|^2,\qquad
 m(y)=\prod_{\tau\in\Psi}\ell_{p_0}\big(x_{2,\tau}^{-1}c_{x_1}(y)_\tau\big),
\]
with $m(y)\ne0$. The function $\rho_{x_1}$ is smooth on $\Omega(x_1)$
(Lemma~\ref{is-onto-share-counts-section-kronecker-solve-chart}), and $m$ is a polynomial in $y$. So $f$ is smooth on
$\Omega(x_1)$. Near $z\in\Omega(x_1)$, $\log|m|^2$ is $\mu+\bar\mu$ for a holomorphic logarithm $\mu$ of $m$, so its Levi form
vanishes, and $\mathcal L_f(z;\xi)=\mathcal L_{\rho_{x_1}}(z;\xi)$ by linearity. And $f(z)\geq0$ by
Lemma~\ref{is-onto-step-free-ball-bound}(1).

\emph{Graph matrices.} Let $x\in\mathcal D_X$ and $\tau\in\Psi$. Put $g=x_{2,\tau}^{-1}x_\tau$, let $t$ be the transpose of its
leading $p_0\times p_0$ block, and $l$ the $p_0\times r$ matrix with $l_{ak}=g_{ka}$ for $k>p_0$ and $l_{ak}=0$ for $k\le p_0$.
Then $\det t=\ell_{p_0}(g)\ne0$, by Lemma~\ref{is-onto-step-free-ball-minor}(3); put $W_\tau(x)=t^{-1}l$. As $x_2$ is
orthonormal, the Gram matrix of the columns of $x_\tau=x_{2,\tau}g$ is $g^{\mathsf T}J\bar g$, and its leading
$p_0\times p_0$ block, the Gram matrix of the first $p_0$ columns of $x_\tau$, is $tt^*-ll^*$. This block is positive
definite, as in the proof of Lemma~\ref{is-onto-share-majorant}(1). Hence
$1-W_\tau(x)W_\tau(x)^*=t^{-1}(tt^*-ll^*)(t^{-1})^*$ is positive definite, with determinant
$|d_{p_0}(x,\tau)|/|\ell_{p_0}(g)|^2$, and
\[
 \rho(x_2,x)=-\sum_{\tau\in\Psi}\log\det\big(1-W_\tau(x)W_\tau(x)^*\big).
\]
For $y$ near a point of $\Omega(x_1)$, the entries of $W_\tau(c_{x_1}(y))$ are polynomials in $y$ divided by the nonzero
polynomial $\ell_{p_0}(x_{2,\tau}^{-1}c_{x_1}(y)_\tau)$, so they are holomorphic in $y$.

\emph{The gradient.} Let $z\in\Omega(x_1)$ and $\xi\in\mathbf C^n$. For $\tau\in\Psi$ the positive definite Hermitian
matrix $1-W_\tau(c_{x_1}(z))W_\tau(c_{x_1}(z))^*$ is $U_\tau\Lambda_\tau U_\tau^*$ with $U_\tau$ unitary and $\Lambda_\tau$ diagonal
with positive entries. For $y$ near $z$ put $G_\tau(y)=U_\tau^*W_\tau(c_{x_1}(y))$, a matrix of holomorphic functions. As
$U_\tau^*(1-W_\tau W_\tau^*)U_\tau=1-G_\tau G_\tau^*$ is positive definite, its diagonal entries
$1-\sum_b|G_\tau(y)_{ab}|^2$ are positive. Put
\[
 \psi(y)=\sum_{\tau\in\Psi}\sum_a-\log\Big(1-\sum_b|G_\tau(y)_{ab}|^2\Big).
\]
As $1-G_\tau(z)G_\tau(z)^*=\Lambda_\tau$ is diagonal, $\psi(z)=f(z)$. For $y$ near $z$, Hadamard's inequality
\cite[Theorem~7.8.1]{HJ13} for the positive definite matrices $1-G_\tau(y)G_\tau(y)^*$, whose determinants are those of
$1-W_\tau(c_{x_1}(y))W_\tau(c_{x_1}(y))^*$, gives $\psi(y)\le f(y)$. By
Lemma~\ref{is-onto-step-orbit-interpolation-large-pole-exists-sharp-rules}(2) and (3), $\psi$ has a sharp gradient at
$z$ in the direction $\xi$, and by part (4) of that lemma so has $f$. As $\mathcal L_f(z;\xi)=\mathcal L_{\rho_{x_1}}(z;\xi)$,
this is the claim.
% Lean: potential_transfer: contDiffAt_rhX_cpt and rhX_cpt_eq_add (minorProd, minorProd_ne_zero through
% lminor_inv_mul_ne_zero, contDiffAt_log_minorProd), rhX_nonneg, leviO_rhX (ExF.levi_logNormSq_zero), and
% sharpAt_rhX_cpt. The graph matrices are topT, lowT and graphT, with leading_gram, one_sub_graph, graph_posDef
% (gramPos_eq, Data.gramPos_posDef) and rhX_eq_sum_graph; contDiffAt_graph for holomorphy (adjugate). The gradient:
% exists_unitary_rows (the eigenvectorUnitary, conj_diag), sharpAt_sum', sharpAt_sum and sharpAt_ball (row_lt_one
% for the diagonal entries), and SharpAt.of_minorant with rows_le_neg_log_det, which applies det_le_prod_diag,
% Hadamard's inequality, proved in the Lean through the AM-GM inequality RgB.amgm.
% Size: about 1,400 lines of Lean below potential_transfer (text count through namespaced names, stopping at
% statements the book cites).
\end{proof}

\begin{definition}\label{is-onto-step-orbit-interpolation-large-pole-exists-function}
Let $X$ consist of the data (1) to (6) of flag data of rank $r$ over $E$, use the notation of
Definition~\ref{is-onto-step-free-ball}, let $G\subset GL_r(E)$ be a subgroup with a free ball of radius $\rho_0$ at an
orthonormal point $x_0$, and let $k\colon[0,\infty)\to\mathbf R$ be a function. For $x\in\mathcal D_X$ put
$\psi_k(x)=k(\rho(x_0,\gamma x))$ if some $\gamma\in G$ satisfies $\rho(x_0,\gamma x)<\rho_0$, which determines $\gamma$ by
Lemma~\ref{is-onto-step-orbit-interpolation-large-pole-exists-translate}(2), and $\psi_k(x)=0$ otherwise; for
$x\notin\mathcal D_X$ put $\psi_k(x)=0$.
\end{definition}
\leanref{Development/AbelianVarietyTheory/Proofs/DivisorProofs/ExistsBlockDetGe.lean}{GlobalTheory.PoleExistence.poleFn}

\begin{lemma}\label{is-onto-step-orbit-interpolation-large-pole-exists-function-local}
Let $X$, $G$, $x_0$, $\rho_0$ and $k$ be as in
Definition~\ref{is-onto-step-orbit-interpolation-large-pole-exists-function}, and use the notation of
Definition~\ref{is-onto-share-counts-section-levi}.
\begin{enumerate}
\item $\psi_k(\gamma xb)=\psi_k(x)$ for $\gamma\in G$, $x\in\mathcal D_X$ and $b\in\mathcal B^\Psi$.
\item $\psi_k(x)=k(\rho(x_0,x))$ for $x\in\mathcal D_X$ with $\rho(x_0,x)<\rho_0$.
\item Let $a<\rho_0$ be a real number with $k(u)=0$ for $u\geq a$, let $x_1$ be an orthonormal point and
  $z\in\Omega(x_1)$. Then either there is an orthonormal point $x_2$ with $\psi_k(c_{x_1}(w))=k(\rho(x_2,c_{x_1}(w)))$ for
  all $w$ near $z$, or $\psi_k(c_{x_1}(w))=0$ for all $w$ near $z$.
\end{enumerate}
\end{lemma}
\leanref{Development/AbelianVarietyTheory/Proofs/DivisorProofs/ExistsBlockDetGe.lean}{GlobalTheory.PoleExistence.poleFn_inv}
\leanref{Development/AbelianVarietyTheory/Proofs/DivisorProofs/ExistsBlockDetGe.lean}{GlobalTheory.PoleExistence.poleFn_eq}
\leanref{Development/AbelianVarietyTheory/Proofs/DivisorProofs/ExistsBlockDetGe.lean}{GlobalTheory.PoleExistence.poleFn_eq_zero}
\leanref{Development/AbelianVarietyTheory/Proofs/DivisorProofs/ExistsBlockDetGe.lean}{GlobalTheory.PoleExistence.poleFn_local}
\begin{proof}
For $\gamma\in G$ and $y\in\mathcal D_X$, $\gamma y\in\mathcal D_X$ and $\gamma^{-1}x_0$ is orthonormal, and
$\rho(x_0,\gamma y)=\rho(\gamma^{-1}x_0,y)$, by Lemma~\ref{is-onto-step-orbit-interpolation-large-pole-exists-translate}(1),
as $\gamma$ preserves $h$.

(1) Put $y=\gamma xb$, which lies in $\mathcal D_X$ by Lemma~\ref{is-onto-share-majorant}(2). For $g\in GL_r(E)$,
$gy=(g\gamma x)b$, so $\rho(x_0,gy)=\rho(x_0,g\gamma x)$ by Lemma~\ref{is-onto-step-free-ball-bound}(1). If some $g\in G$ satisfies
$\rho(x_0,gx)<\rho_0$, then $g'=g\gamma^{-1}\in G$ satisfies $\rho(x_0,g'y)=\rho(x_0,gx)<\rho_0$, and
$\psi_k(y)=k(\rho(x_0,gx))=\psi_k(x)$. Otherwise no $g'\in G$ satisfies $\rho(x_0,g'y)<\rho_0$, as $g'\gamma$ would satisfy
$\rho(x_0,g'\gamma x)<\rho_0$, and $\psi_k(y)=0=\psi_k(x)$.

(2) Take $\gamma=1$ in the definition.

(3) Put $y=c_{x_1}(z)\in\mathcal D_X$. Suppose first that some $\gamma\in G$ satisfies $\rho(x_0,\gamma y)<\rho_0$, and put
$x_2=\gamma^{-1}x_0$. The function $y'\mapsto\rho(x_2,y')$ is continuous on $\mathcal D_X$, as minus the logarithm of the
product of the continuous positive functions $q_\tau(x_2,\cdot)$ of Lemma~\ref{is-onto-step-free-ball-minor}(3). As
$\Omega(x_1)$ is open, every $w$ near $z$ satisfies $c_{x_1}(w)\in\mathcal D_X$ and
$\rho(x_0,\gamma c_{x_1}(w))=\rho(x_2,c_{x_1}(w))<\rho_0$, and then $\psi_k(c_{x_1}(w))=k(\rho(x_2,c_{x_1}(w)))$.

Suppose now that $\rho(x_0,\gamma y)\geq\rho_0$ for every $\gamma\in G$. Let $V$ be the neighbourhood of $y$ of
Lemma~\ref{is-onto-step-orbit-interpolation-large-pole-exists-lower} for $t=\rho_0-a>0$. Every $w$ near $z$ satisfies
$c_{x_1}(w)\in V\cap\mathcal D_X$, and then, for every $\gamma\in G$,
\[
 \rho(x_0,\gamma c_{x_1}(w))=\rho(\gamma^{-1}x_0,c_{x_1}(w))\geq\rho(\gamma^{-1}x_0,y)-(\rho_0-a)=\rho(x_0,\gamma y)-\rho_0+a\geq a,
\]
so $k(\rho(x_0,\gamma c_{x_1}(w)))=0$ for every $\gamma\in G$, and $\psi_k(c_{x_1}(w))=0$.
% Lean: (1) is poleFn_inv (e, poleFn_eq for g gamma^-1, MajInv.lact_mem_dom, OntoS.dom_rmul); (2) is poleFn_eq with
% gamma = 1 (as in hone in poleEx_of_rhTr), through FreeBall.unique; (3) is poleFn_local: the first case with
% rhX_lact, continuousOn_rhX and poleFn_eq, the second with rhX_lower for t = rho0 - a and poleFn_eq_zero. The
% Lean's poleFn takes the value k (rhX x0 (lact gamma x)) for a chosen gamma at every x; the book sets
% psi_k = 0 off the flag domain, where nothing is asked of it.
\end{proof}

\begin{lemma}\label{is-onto-step-orbit-interpolation-large-pole-exists}
Let $X$ consist of the data (1) to (6) of flag data of rank $r$ over $E$, put $\dim X=|\Psi|p_0(r-p_0)$, and let
$G\subset GL_r(E)$ be a subgroup, $x_0$ an orthonormal point, and $\rho_0>0$ and $\theta>1$ real numbers such that
$G$ has a free ball of radius $\rho_0$ at $x_0$ (Definition~\ref{is-onto-step-free-ball}). Then there are real
numbers $\varepsilon$ and $C$ with $0<\varepsilon<\rho_0$ such that for every real number $0<\delta\le1$ there is a pole
function for $(G,x_0,\theta\dim X/\rho_0,\varepsilon,C,\delta)$
(Definition~\ref{is-onto-step-orbit-interpolation-large-pole}).
\end{lemma}
\leanref{Development/AbelianVarietyTheory/Definitions/UpperBounds/FreeBallIffOfDimXEqZero.lean}{GlobalTheory.Up.PoleEx}
\leanref{Development/AbelianVarietyTheory/Proofs/DivisorProofs/ExistsBlockDetGe.lean}{GlobalTheory.PoleExistence.poleEx_holds}
\begin{proof}
Put $d=\dim X$ and $\kappa=\theta d/\rho_0$, and write $\partial_j$ for $\partial/\partial z_j$.
Lemma~\ref{is-onto-step-orbit-interpolation-large-pole-exists-cutoff} for $\rho_0$ and $\theta$ gives $\varepsilon$ and $C_0$
with $0<\varepsilon<\rho_0$; put $C=dC_0$. Let $0<\delta\le1$, let $g$ be the function of that lemma for $\delta$, put
$k=d\,g$ on $[0,\infty)$, and let $\psi=\psi_k$ be the function of
Definition~\ref{is-onto-step-orbit-interpolation-large-pole-exists-function}. It is a pole function for
$(G,x_0,\kappa,\varepsilon,C,\delta)$; the five conditions of
Definition~\ref{is-onto-step-orbit-interpolation-large-pole} are checked in turn.

(1) This is Lemma~\ref{is-onto-step-orbit-interpolation-large-pole-exists-function-local}(1).

(2) For $x\in\mathcal D_X$, $\psi(x)$ is $0$ or $d\,g(u)$ with $u=\rho(x_0,\gamma x)\geq0$
(Lemma~\ref{is-onto-step-free-ball-bound}(1)), and $g(u)\le0$.

(3) Let $x_1$ be an orthonormal point, $z\in\Omega(x_1)$ and $\xi\in\mathbf C^n$, and put
$L=\mathcal L_{\rho_{x_1}}(z;\xi)$, which is $\geq0$ by Lemma~\ref{is-onto-step-orbit-interpolation-large-solve-levi}. As
$k(u)=0$ for $u\geq\rho_0-\varepsilon$, Lemma~\ref{is-onto-step-orbit-interpolation-large-pole-exists-function-local}(3)
applies with $a=\rho_0-\varepsilon$. If $\psi\circ c_{x_1}$ vanishes near $z$, it is smooth near $z$ with Levi form $0\geq-\kappa L$
at $z$. Otherwise $\psi\circ c_{x_1}=d\,g\circ f$ near $z$, where $f(w)=\rho(x_2,c_{x_1}(w))$ for an orthonormal point $x_2$. By
Lemma~\ref{is-onto-step-orbit-interpolation-large-pole-exists-transfer}, $f$ is smooth with values $\geq0$, so
$\psi\circ c_{x_1}$ is smooth near $z$; and with $u=f(z)$ and $N=|\sum_j\partial_jf(z)\xi_j|^2$, the same lemma gives
$\mathcal L_f(z;\xi)=L$ and $N\le uL$. By Lemma~\ref{is-onto-step-orbit-interpolation-large-pole-exists-sharp-rules}(1),
\[
 \mathcal L_{\psi\circ c_{x_1}}(z;\xi)=d\big(g'(u)L+g''(u)N\big).
\]
If $g''(u)\geq0$, this is $\geq0\geq-\kappa L$, as $g'(u)\geq0$. Otherwise $g''(u)N\geq g''(u)uL$, and it is at least
$d\big(g'(u)+ug''(u)\big)L\geq-d(\theta/\rho_0)L=-\kappa L$. So $\psi\circ c_{x_1}$ is smooth on $\Omega(x_1)$, with the
bound for its Levi form.

(4) Let $x\in\mathcal D_X$ with $\rho(x_0,x)\le\varepsilon$. As $\varepsilon<\rho_0$,
Lemma~\ref{is-onto-step-orbit-interpolation-large-pole-exists-function-local}(2) gives
$\psi(x)=d\,g(\rho(x_0,x))\le d\log(\rho(x_0,x)+\delta)+dC_0$, by (3) of
Lemma~\ref{is-onto-step-orbit-interpolation-large-pole-exists-cutoff}.

(5) Let $x\in\mathcal D_X$ with $\varepsilon/2\le\rho(x_0,x)\le\varepsilon$. As in (4), $\psi(x)=d\,g(\rho(x_0,x))\geq-dC_0=-C$.
% Lean: poleEx_holds = poleEx_of_rhTr GlobalTheory.PotentialTransfer.potential_transfer. poleEx_of_rhTr:
% exists_cutoff_log_sharp (Lemma is-onto-step-orbit-interpolation-large-pole-exists-cutoff), D = dimX, the constant
% D * C, psi = poleFn X G x0 rho0 (D * g); the five fields: poleFn_inv; poleFn with rhX_nonneg and the first bound
% of hgd; poleFn_local with a = rho0 - eps, RhTr (potential_transfer, Lemma
% is-onto-step-orbit-interpolation-large-pole-exists-transfer), H.leviChain, leviRh_nonneg (Lemma
% is-onto-step-orbit-interpolation-large-solve-levi) and the case split on the sign of g''; hone (poleFn_eq with
% gamma = 1) with hgu and hgl.
% Size: about 3,400 lines of Lean below poleEx_holds (text count through namespaced names, stopping at statements
% the book cites): about 2,000 in poleEx_of_rhTr with exists_cutoff_log_sharp and the poleFn lemmas, and about
% 1,400 in potential_transfer. Nothing below this proof remains open.
\end{proof}

\begin{lemma}\label{is-onto-step-orbit-interpolation-large-cutoff}
Let $\varepsilon>0$ be a real number. Then there are a smooth function $\chi\colon\mathbf R\to\mathbf R$ and a real
number $D$ such that $\chi(u)=1$ for $u\le\varepsilon/2$, $\chi(u)=0$ for $u\geq\varepsilon$, $|\chi|\le1$,
$\chi'(u)^2\le D$ for every $u$, and $\chi'(u)=0$ for $u$ outside $[\varepsilon/2,\varepsilon]$.
\end{lemma}
\leanref{Development/AbelianVarietyTheory/Proofs/TorusHaarMeasure/ExistsPhlaBound.lean}{GlobalTheory.PoleDbar.exists_cutoff}
\begin{proof}
Let $\sigma$ be a smooth function on $\mathbf R$ with values in $[0,1]$, equal to $0$ on $(-\infty,0]$ and to $1$ on
$[1,\infty)$, and put $\chi(u)=\sigma(2-2u/\varepsilon)$. For $u\le\varepsilon/2$, $2-2u/\varepsilon\geq1$, and for
$u\geq\varepsilon$, $2-2u/\varepsilon\le0$; this gives the values of $\chi$. On the open sets $u<\varepsilon/2$ and
$u>\varepsilon$ the function $\chi$ is constant, so $\chi'$ vanishes outside $[\varepsilon/2,\varepsilon]$. The
continuous function $\chi'^2$ is bounded on the compact interval $[\varepsilon/2,\varepsilon]$ by some $D\geq0$.
% Lean: exists_cutoff, with Real.smoothTransition for sigma and max D 0 for D.
\end{proof}

\begin{lemma}\label{is-onto-step-orbit-interpolation-large-block}
Let $X$ consist of the data (1) to (6) of flag data of rank $r$ over $E$, use the notation of
Definitions~\ref{is-onto-share-counts-section-levi} and~\ref{is-onto-step-free-ball}, and let $x_0$ be an orthonormal
point.
\begin{enumerate}
\item For $z\in\Omega(x_0)$ and $\tau\in\Psi$, $0<|d_{p_0}(c_{x_0}(z),\tau)|\le1$ and
  $\rho(x_0,c_{x_0}(z))=\rho_{x_0}(z)\geq0$.
\item Let $\varepsilon$ be a real number and $\lambda$ a weight with $\lambda_{\tau,i}=\lambda_{\tau,j}$ whenever
  $i,j\le p_0$ or $i,j>p_0$. Then there is a real number $A$ such that $-\varphi_\lambda(c_{x_0}(z))\le A$ for every
  $z\in\Omega(x_0)$ with $\rho_{x_0}(z)\le\varepsilon$.
\end{enumerate}
\end{lemma}
\leanref{Development/AbelianVarietyTheory/Proofs/TorusHaarMeasure/ExistsPhlaBound.lean}{GlobalTheory.PoleDbar.exists_phla_bound}
\leanref{Development/AbelianVarietyTheory/Proofs/TorusHaarMeasure/ExistsPhlaBound.lean}{GlobalTheory.PoleDbar.phla_block}
\leanref{Development/AbelianVarietyTheory/Proofs/TorusHaarMeasure/ExistsPhlaBound.lean}{GlobalTheory.PoleDbar.abs_Lm_le}
\leanref{Development/AbelianVarietyTheory/Definitions/UpperBounds/FreeBallIffOfDimXEqZero.lean}{GlobalTheory.Up.rhX_cpt}
\begin{proof}
(1) Let $z\in\Omega(x_0)$ and $x=c_{x_0}(z)\in\mathcal D_X$. The matrix $x_{0,\tau}^{-1}x_\tau=\ell_\tau(z)$ has leading
minors $1$, so in the notation of Lemma~\ref{is-onto-step-free-ball-minor}(3),
$q_\tau(x_0,x)=|d_{p_0}(x,\tau)|$, which lies in $(0,1]$ by that lemma. Hence
$\rho(x_0,x)=-\sum_\tau\log|d_{p_0}(x,\tau)|=\rho_{x_0}(z)$, a sum of terms $\geq0$.

(2) Let $a_\tau$ and $a'_\tau$ be the values of $\lambda_{\tau,i}$ for $i\le p_0$ and for $i>p_0$. Let
$z\in\Omega(x_0)$, $x=c_{x_0}(z)$ and $L_{\tau,i}=\log|d_i(x,\tau)|$. Then $L_{\tau,0}=0$, and $L_{\tau,r}=0$ as in the
proof of Lemma~\ref{is-onto-share-counts-section-curvature}. Summing the formula of
Definition~\ref{is-onto-share-counts-section-levi} over the two blocks,
\[
 \varphi_\lambda(x)=\sum_\tau\big(a_\tau(L_{\tau,p_0}-L_{\tau,0})+a'_\tau(L_{\tau,r}-L_{\tau,p_0})\big)
 =\sum_\tau(a_\tau-a'_\tau)L_{\tau,p_0}.
\]
By (1), every $-L_{\sigma,p_0}$ is $\geq0$, so $|L_{\tau,p_0}|\le\sum_\sigma(-L_{\sigma,p_0})=\rho_{x_0}(z)$. Hence
$-\varphi_\lambda(x)\le\sum_\tau|a_\tau-a'_\tau|\,\rho_{x_0}(z)\le A$ when $\rho_{x_0}(z)\le\varepsilon$, with
$A=\varepsilon\sum_\tau|a_\tau-a'_\tau|$.
% Lean: (1) is FramePotential.lminor_cpt, factor_mem and rhX_cpt (through frame_potential_basic). (2) is
% exists_phla_bound: blockConst_iff, phla_block (IW.phla_eq, the summation by blocks sum_block, IW.Lm_zero and
% IW.Lm_last_cpt), and abs_Lm_le (neg_log_factor_le_rhX, factor_mem, lminor_cpt, rhX_cpt). The same bound,
% exponentiated, is OrbitInterpolation.exp_neg_phla_le, proved the same way through sum_fin_blocks.
\end{proof}

\begin{lemma}\label{is-onto-step-orbit-interpolation-large-sublevel}
Let $X$ consist of the data (1) to (6) of flag data of rank $r$ over $E$, use the notation of
Definition~\ref{is-onto-share-counts-section-levi}, and let $x_0$ be an orthonormal point.
\begin{enumerate}
\item Let $K\subset\mathcal D_X$ be compact. Then the integral of $D_{x_0}(z)$ over the set of the $z\in\Omega(x_0)$ with
  $c_{x_0}(z)\in K\mathcal B^\Psi$ is finite.
\item For every real number $a$, the integral of $D_{x_0}(z)$ over the set of the $z\in\Omega(x_0)$ with
  $\rho_{x_0}(z)\le a$ is finite.
\end{enumerate}
\end{lemma}
\leanref{Development/AbelianVarietyTheory/Proofs/TorusHaarMeasure/MuFLtTopOfIsCompact.lean}{GlobalTheory.SatFin.muF_sat_lt_top}
\leanref{Development/AbelianVarietyTheory/Proofs/TorusHaarMeasure/MuFSublevelLtTop.lean}{GlobalTheory.SublevelFin.muF_sublevel_lt_top}
\leanref{Development/AbelianVarietyTheory/Proofs/TorusHaarMeasure/MuFSublevelLtTop.lean}{GlobalTheory.SublevelFin.entry_sq_le}
\leanref{Development/AbelianVarietyTheory/Proofs/TorusHaarMeasure/MuFSublevelLtTop.lean}{GlobalTheory.SublevelFin.minor_sq_le}
\begin{proof}
For an orthonormal point $y$ write $\mu_y$ for the measure $D_y(z)\,dz$ on $\Omega(y)$, and for a set
$T\subset M_r(\mathbf C)^\Psi$ let $\operatorname{sat}_y(T)$ be the set of the $z\in\Omega(y)$ with
$c_y(z)\in T\mathcal B^\Psi$.

(1) The sets $T\subset M_r(\mathbf C)^\Psi$ with $\mu_{x_0}(\operatorname{sat}_{x_0}(T))<\infty$ are closed under taking
subsets and finite unions. As $K$ is compact, it suffices to show that every $x\in K$ has a neighbourhood $U$ in $M_r(\mathbf C)^\Psi$ with
$\mu_{x_0}(\operatorname{sat}_{x_0}(U))<\infty$. By Lemma~\ref{is-onto-share-majorant}(3), $x=x_1b_1$ with $x_1$
orthonormal and $b_1\in\mathcal B^\Psi$. Let $\Delta$, $z(\cdot)$ and $b(\cdot)$ be as in
Lemma~\ref{is-onto-share-counts-section-transfer-lu} for $x_1$; by part (3) of that lemma, $\Delta(x)\ne0$ and
$z(x)=0$. Choose $\varrho>0$ such that the closed ball $\bar B$ of radius $\varrho$ around $0$ lies in the open set
$\Omega(x_1)$, and let $U$ be the set of the $y\in\mathcal D_X$ with $\Delta(y)\ne0$ and $|z(y)|\le\varrho$; it is a
neighbourhood of $x$, as $z(\cdot)$ is continuous. By part (2), each $y\in U$ is $c_{x_1}(z(y))b(y)$, so
$U\mathcal B^\Psi\subset K_1\mathcal B^\Psi$ for the compact set $K_1=c_{x_1}(\bar B)\subset\mathcal D_X$.

The set $K_1\mathcal B^\Psi$ consists of the $y$ with $\det y_\tau\ne0$ for every $\tau$ such that, for some $x'\in K_1$,
$\operatorname{adj}(x'_\tau)\,y_\tau$ is upper triangular for every $\tau$. It is the intersection of an open set with the
projection of a closed subset of $M_r(\mathbf C)^\Psi\times K_1$, which is closed as $K_1$ is compact; so it is
measurable, and it is invariant under $\mathcal B^\Psi$. Lemma~\ref{is-onto-share-counts-section-poincare-measure}(1),
for its indicator function, gives $\mu_{x_0}(\operatorname{sat}_{x_0}(K_1))=\mu_{x_1}(\operatorname{sat}_{x_1}(K_1))$. If
$z\in\Omega(x_1)$ and $c_{x_1}(z)=c_{x_1}(z')b$ with $z'\in\bar B$ and $b\in\mathcal B^\Psi$, then $z=z'$ by part (3) of
Lemma~\ref{is-onto-share-counts-section-transfer-lu}; so $\operatorname{sat}_{x_1}(K_1)\subset\bar B$. Finally
$\mu_{x_1}(\bar B)<\infty$, as $D_{x_1}$ is continuous on the compact set $\bar B\subset\Omega(x_1)$.

(2) Let $J$ be as in Lemma~\ref{is-onto-step-free-ball-minor}, and put $M=((r+1)e^a)^{1/2}$. Let $z\in\Omega(x_0)$ with
$\rho_{x_0}(z)\le a$. By Lemma~\ref{is-onto-share-majorant}(3), $c_{x_0}(z)=xb$ with $x$ orthonormal and
$b\in\mathcal B^\Psi$. Put $g_\tau=x_{0,\tau}^{-1}x_\tau$. By Lemma~\ref{is-onto-step-orbit-interpolation-large-block}(1),
Lemma~\ref{is-onto-step-free-ball-minor}(3) and $|d_{p_0}(x,\tau)|=1$,
\[
 e^{-\rho_{x_0}(z)}=\prod_\tau q_\tau(x_0,c_{x_0}(z))=\prod_\tau q_\tau(x_0,x)=\prod_\tau|\ell_{p_0}(g_\tau)|^{-2}.
\]
By parts (1) and (2) of Lemma~\ref{is-onto-step-free-ball-minor}, every $|\ell_{p_0}(g_\sigma)|$ is at least $1$, so
$|\ell_{p_0}(g_\tau)|^2\le e^{\rho_{x_0}(z)}\le e^a$. As $g_\tau^*Jg_\tau=J$, the matrix obtained from $g_\tau$ by putting its
rows and columns in the order $p_0+1,\dots,r,1,\dots,p_0$ satisfies the hypothesis of the bound $(*)$ in the proof of
Lemma~\ref{is-onto-share-counts-section-poincare-count-sum}, for $p=r-p_0$ and $q=p_0$, and its lower right
$p_0\times p_0$ block is the leading block of $g_\tau$. So every entry of $g_\tau$ satisfies
$|(g_\tau)_{ij}|^2\le(r+1)|\ell_{p_0}(g_\tau)|^2\le M^2$. Hence $x$ lies in the set $K$ of the orthonormal points
$x_0g$, with $g\in M_r(\mathbf C)^\Psi$ having entries of absolute value at most $M$. This set is compact, as the image of a
compact set under a continuous map intersected with the closed set of orthonormal points, and it lies in
$\mathcal D_X$ by Lemma~\ref{is-onto-share-majorant}(3). So the set in (2) lies in $\operatorname{sat}_{x_0}(K)$, and (1)
applies.
% Lean: (1) is SatFin.muF_sat_lt_top: sat_eq_preimage, IsCompact.induction_on, ortho_holds, lu_holds, the
% neighbourhood {y in dom | Delta y != 0} cap zOf^-1 (closedBall 0 rho), muF_borSat_eq (measInv_holds for the
% indicator of borSat, measurable by FrameMeas.measurableSet_borSat through isClosedMap_fst_of_compactSpace),
% FundSetPf.cpt_uniq, and muF_lt_top_of_isCompact. (2) is muF_sublevel_lt_top: minor_sq_le (frame_potential_basic,
% factor_rmul, norm_dmin_of_isOrth, one_le_minor_of_isOrth), entry_sq_le (sumBd_entry_bound for
% FrameSumBound.isUpq_frame, the matrix reordered by swapIdx), isClosed_isOrth, and the null set Omega(x0)^c.
\end{proof}

\begin{lemma}\label{is-onto-step-orbit-interpolation-large-gradient}
Let $X$ consist of the data (1) to (6) of flag data of rank $r$ over $E$, use the notation of
Definition~\ref{is-onto-share-counts-section-levi}, write $\partial_j$ for $\partial/\partial z_j$, and let $x_0$ be an
orthonormal point. Then for every $z\in\Omega(x_0)$ and every $\xi\in\mathbf C^n$
\[
 \Big|\sum_j\partial_j\rho_{x_0}(z)\,\xi_j\Big|^2\le|\Psi|\min(p_0,r-p_0)\,\mathcal L_{\rho_{x_0}}(z;\xi).
\]
\end{lemma}
\leanref{ReciprocityTheory/ReductiveGroups/ArithmeticGeometry/ArakelovVolumes/UpperBounds/FundSetMeasStatement.lean}{ReciprocityTheory.ChartRestriction.G.AR.Up.RhGradFStatement}
\leanref{ReciprocityTheory/ReductiveGroups/ArithmeticGeometry/ArakelovVolumes/UpperBounds/RhGradFPf.lean}{ReciprocityTheory.ChartRestriction.G.AR.Up.RhGradFStatement_holds}
\leanref{ReciprocityTheory/ReductiveGroups/ArithmeticGeometry/ArakelovVolumes/UpperBounds/RhGradStatementPf.lean}{ReciprocityTheory.ChartRestriction.G.AR.Up.RhGradStatement_holds}
\leanref{ReciprocityTheory/ReductiveGroups/ArithmeticGeometry/ArakelovVolumes/UpperBounds/RhGradStatementPfB.lean}{ReciprocityTheory.ChartRestriction.G.AR.Up.RgD.block}
\leanref{ReciprocityTheory/ReductiveGroups/ArithmeticGeometry/ArakelovVolumes/UpperBounds/RhGradStatementPfA.lean}{ReciprocityTheory.ChartRestriction.G.AR.Up.RgA.grad_le_of_minorant}
\begin{proof}
Put $p=p_0$, $q=r-p_0$, $m=\min(p,q)$ and $C=|\Psi|m$, and use the matrices $Z_\tau(z)$ of
Lemma~\ref{is-onto-share-counts-section-kronecker-solve-chart-block}, of size $p\times q$. If $m=0$ or $\Psi$ is empty, the
matrices $1-Z_\tau(z)Z_\tau(z)^*$ have determinant $1$, so $\rho_{x_0}=0$ on $\Omega(x_0)$ by part (4) of that lemma, and
both sides vanish. Suppose that $m\geq1$ and $\Psi\ne\emptyset$.

\emph{Matrix balls.} Let $\mathcal U$ be the open set of the families $Z=(Z_\tau)_{\tau\in\Psi}$ of complex $p\times q$
matrices such that every $1-Z_\tau Z_\tau^*$ is positive definite, and on it put
$\phi(Z)=-\sum_\tau\log\det(1-Z_\tau Z_\tau^*)$. By parts (1), (3) and (4) of
Lemma~\ref{is-onto-share-counts-section-kronecker-solve-chart-block}, the map $z\mapsto Z(z)=(Z_\tau(z))_\tau$ is
holomorphic, it maps $\Omega(x_0)$ into $\mathcal U$, and $\rho_{x_0}=\phi\circ Z$ on $\Omega(x_0)$. For a holomorphic map the
chain rule gives $\sum_j\partial_j(\phi\circ Z)(z)\,\xi_j=\sum_a\partial_a\phi(Z(z))\,\zeta_a$ and
$\mathcal L_{\phi\circ Z}(z;\xi)=\mathcal L_\phi(Z(z);\zeta)$, with $\zeta=DZ(z)\xi$ and $a$ running over the entries of the
matrices, as in the proof of Lemma~\ref{is-onto-share-counts-section-transfer}. So it suffices to show
$|\sum_a\partial_a\phi(Z_0)\zeta_a|^2\le C\,\mathcal L_\phi(Z_0;\zeta)$ for $Z_0\in\mathcal U$ and every $\zeta$.

\emph{Comparison functions.} Fix $Z_0\in\mathcal U$ and $\tau\in\Psi$. Suppose first that $p\le q$, so $m=p$. The matrix
$A_0=1-Z_{0,\tau}Z_{0,\tau}^*$ is positive definite; write $A_0^{-1}=y^*y$ with a $p\times p$ matrix $y$. Put
\[
 Q_\tau(Z)=\operatorname{Re}\operatorname{tr}\big((1-Z_\tau Z_\tau^*)A_0^{-1}\big)=\operatorname{tr}(y^*y)-\sum_{a,b}|(yZ_\tau)_{ab}|^2 .
\]
It is a constant minus a sum of squares of absolute values of linear functions of $Z$, so its Levi form is $\le0$
everywhere, and $Q_\tau(Z_0)=\operatorname{tr}(1)=m$. For $Z\in\mathcal U$ the matrix $y(1-Z_\tau Z_\tau^*)y^*$ is positive
definite, with trace $Q_\tau(Z)$ and determinant $\det(1-Z_\tau Z_\tau^*)/\det A_0$; the inequality between the geometric
and the arithmetic mean of its eigenvalues gives
\[
 \det(1-Z_\tau Z_\tau^*)^{1/m}\le(\det A_0)^{1/m}\,Q_\tau(Z)/m .
\]
If $q<p$, the same holds with $1-Z_\tau^*Z_\tau$ and $1-Z_{0,\tau}^*Z_{0,\tau}$ in place of $1-Z_\tau Z_\tau^*$ and $A_0$:
these $q\times q$ matrices are positive definite, as $(1-Z^*Z)^{-1}=1+Z^*(1-ZZ^*)^{-1}Z$, and
$\det(1-Z^*Z)=\det(1-ZZ^*)$.

\emph{A minorant.} Taking logarithms, summing over $\tau$, and using the concavity of $\log$ for the mean of the
numbers $Q_\tau(Z)/m$, with $\Theta=\sum_\tau Q_\tau$,
\[
 \phi(Z)\geq\phi(Z_0)-m\sum_\tau\log\frac{Q_\tau(Z)}m\geq\phi(Z_0)-C\log\frac{\Theta(Z)}C\qquad(Z\in\mathcal U),
\]
with equality at $Z_0$, where $\Theta(Z_0)=C$. So $\phi-\psi$, with $\psi=-C\log\Theta$, has a local minimum at $Z_0$.

\emph{The bound.} At this minimum the first derivatives of $\phi-\psi$ vanish, and its Levi form is $\geq0$, as in the
proof of Lemma~\ref{is-onto-share-counts-section-kronecker-solve-chart}. So
$\partial_a\phi(Z_0)=-C\,\partial_a\Theta(Z_0)/\Theta(Z_0)$, and, as $\mathcal L_\Theta\le0$,
\[
 \mathcal L_\phi(Z_0;\zeta)\geq\mathcal L_\psi(Z_0;\zeta)=-\frac C{\Theta(Z_0)}\mathcal L_\Theta(Z_0;\zeta)
 +\frac C{\Theta(Z_0)^2}\Big|\sum_a\partial_a\Theta(Z_0)\zeta_a\Big|^2\geq\frac C{\Theta(Z_0)^2}\Big|\sum_a\partial_a\Theta(Z_0)\zeta_a\Big|^2 .
\]
Hence $|\sum_a\partial_a\phi(Z_0)\zeta_a|^2=C^2|\sum_a\partial_a\Theta(Z_0)\zeta_a|^2/\Theta(Z_0)^2\le C\,\mathcal L_\phi(Z_0;\zeta)$.
% Lean: RhGradFStatement_holds = RGF.rhGradF_holds: Zmap (the entries of the ExF.Zt X z tau), Zmap_mem, rh_Zmap
% (rhF = rh o Zmap on OmF x0), Zmap_contDiff, and the chain rules gc_sum_wd_comp and gc_leviV_comp with
% fc_leviV_congr; then RhGradStatement_holds on Dset f p q with f = |Psi|: the case f = 0 or min p q = 0 (nD = 0),
% and RgE.rhGrad_main, with RgD.block (Q through exists_factor, Q_formula and levi_quad_nonpos; posDef_swap and
% det_one_sub_mul_comm when q < p), RgB.amgm (amgm_psd, the weighted AM-GM inequality for the eigenvalues), hkey
% with the concavity of log (strictConcaveOn_log_Ioi), and RgA.grad_le_of_minorant (leviChain and
% leviV_re_nonpos_of_isLocalMax for the local maximum of psi - rh).
% Size: about 1,200 lines of Lean below RhGradFStatement_holds (text count through namespaced names, stopping at
% statements the book cites), much of it the calculus of Wirtinger derivatives.
\end{proof}

\begin{lemma}\label{is-onto-step-orbit-interpolation-large-solve-levi}
Let $X$ consist of the data (1) to (6) of flag data of rank $r$ over $E$, use the notation of
Definition~\ref{is-onto-share-counts-section-levi}, and let $x_0$ be an orthonormal point. Then
$\mathcal L_{\rho_{x_0}}(z;\xi)\geq0$ for every $z\in\Omega(x_0)$ and every $\xi\in\mathbf C^n$.
\end{lemma}
\leanref{Development/AbelianVarietyTheory/Definitions/UpperBounds/FreeBallIffOfDimXEqZero.lean}{GlobalTheory.Up.leviRh_nonneg}
\leanref{ReciprocityTheory/ReductiveGroups/ArithmeticGeometry/ArakelovVolumes/UpperBounds/OntoStatementPfB.lean}{ReciprocityTheory.ChartRestriction.G.AR.Up.OntoB.levi_rh_nonneg}
\begin{proof}
Write $\rho=\rho_{x_0}$, $\partial_j=\partial/\partial z_j$ and $\bar\partial_j=\partial/\partial\bar z_j$, and put
$R=|\Psi|\min(p_0,r-p_0)$. Let $z\in\Omega(x_0)$ and $\xi\in\mathbf C^n$. If $R>0$, then by
Lemma~\ref{is-onto-step-orbit-interpolation-large-gradient}
\[
 R\,\mathcal L_\rho(z;\xi)\geq\Big|\sum_j\partial_j\rho(z)\,\xi_j\Big|^2\geq0,
\]
and $\mathcal L_\rho(z;\xi)\geq0$. If $R=0$, the same lemma for the standard basis vectors $\xi$ gives
$\partial_j\rho=0$ on $\Omega(x_0)$ for every $j$. As $\rho$ is real, $\bar\partial_j\rho=\overline{\partial_j\rho}=0$ on the
open set $\Omega(x_0)$, so every $\partial_k\bar\partial_j\rho$ vanishes there, and $\mathcal L_\rho(z;\xi)=0$.
% Lean: Up.leviRh_nonneg = OntoB.levi_rh_nonneg RhGradFStatement_holds. The case R = 0 (hC): hwd for the Wirtinger
% derivatives wd j through Pi.single j 1, then wdb j = 0 on OmF x0 through the real and imaginary parts p and q of
% the real derivative, and leviV = 0 through hev and fderiv of a function that vanishes near z; the case R > 0 is
% le_of_mul_le_mul_left.
\end{proof}

\begin{lemma}\label{is-onto-step-orbit-interpolation-large-solve}
Let $X$ consist of the data (1) to (6) of flag data of rank $r$ over $E$, and use the notation of
Definitions~\ref{is-onto-share-counts-section-levi} and~\ref{is-onto-step-orbit-interpolation-large-pole}. Let
$G\subset GL_r(E)$ be a subgroup, $x_0$ an orthonormal point, $\lambda$ a weight with
$\lambda_{\tau,i}=\lambda_{\tau,j}$ whenever $i,j\le p_0$ or $i,j>p_0$, and $\kappa,\varepsilon,C,\nu,m$ real numbers with
$\varepsilon>0$, $\nu\geq0$ and $\nu+\kappa<m$, such that
\[
 m\,\mathcal L_{\rho_{x_0}}(z;\xi)\le\mathcal L^{x_0}_\lambda(z;\xi)-\mathcal L_{\log D_{x_0}}(z;\xi)
\]
for every $z\in\Omega(x_0)$ and every $\xi\in\mathbf C^n$. Then there are a continuous function
$\chi\colon\mathbf R\to\mathbf R$ with $\chi(u)=1$ for $u\le\varepsilon/2$, $\chi(u)=0$ for $u\geq\varepsilon$ and
$|\chi|\le1$, and a real number $M\geq0$, such that for every real number $\delta>0$ and every pole function $\psi$ for
$(G,x_0,\kappa,\varepsilon,C,\delta)$ there is a holomorphic function $F$ on $\Omega(x_0)$ with
\[
 \int_{\Omega(x_0)}\big|\chi(\rho_{x_0}(z))-F(z)\big|^2e^{-\varphi_\lambda(c_{x_0}(z))-\psi(c_{x_0}(z))}\,
 e^{\nu\rho_{x_0}(z)}\,D_{x_0}(z)\,dz\le M.
\]
\end{lemma}
\leanref{Development/AbelianVarietyTheory/Definitions/UpperBounds/FreeBallIffOfDimXEqZero.lean}{GlobalTheory.Up.PoleSolve}
\leanref{Development/AbelianVarietyTheory/Proofs/TorusHaarMeasure/ExistsPhlaBound.lean}{GlobalTheory.PoleDbar.poleSolve}
\begin{proof}
Write $c=c_{x_0}$, $\Omega=\Omega(x_0)$ and $\rho=\rho_{x_0}$. Let $\chi$ and $D_\chi$ be the function and the number of
Lemma~\ref{is-onto-step-orbit-interpolation-large-cutoff} for $\varepsilon$, and $A$ the number of
Lemma~\ref{is-onto-step-orbit-interpolation-large-block}(2) for $\varepsilon$ and $\lambda$. Put
$R=|\Psi|\min(p_0,r-p_0)$ and $K=D_\chi R/(m-\kappa-\nu)$, which is $\geq0$ as $D_\chi\geq\chi'(0)^2$. Let $S$ be the set
of the $z\in\Omega$ with $\varepsilon/2\le\rho(z)\le\varepsilon$. By
Lemma~\ref{is-onto-step-orbit-interpolation-large-sublevel}(2), $I_S=\int_SD_{x_0}(z)\,dz$ is finite. Put
$M=Ke^{A+C+\nu\varepsilon}I_S$.

Let $\delta>0$ and let $\psi$ be a pole function for $(G,x_0,\kappa,\varepsilon,C,\delta)$.

\emph{The weight.} On $\Omega$ put $\Phi=\varphi_\lambda\circ c+\psi\circ c-\log D_{x_0}-\nu\rho$. It is smooth: the minors
$d_i(c(z),\tau)$ are nonzero polynomials in the real and imaginary parts of $z$ on $\Omega$, and $\psi\circ c$ is smooth
by (3) of Definition~\ref{is-onto-step-orbit-interpolation-large-pole}. By the linearity of the Levi form, the
hypothesis and (3) of that definition,
\[
 \mathcal L_\Phi(z;\xi)=\mathcal L^{x_0}_\lambda(z;\xi)-\mathcal L_{\log D_{x_0}}(z;\xi)+\mathcal L_{\psi\circ c}(z;\xi)
 -\nu\,\mathcal L_\rho(z;\xi)\geq(m-\kappa-\nu)\,\mathcal L_\rho(z;\xi),
\]
which is $\geq0$, as $m-\kappa-\nu>0$ and $\mathcal L_\rho(z;\xi)\geq0$ by
Lemma~\ref{is-onto-step-orbit-interpolation-large-solve-levi}. Moreover
$e^{-\Phi}=e^{-\varphi_\lambda\circ c-\psi\circ c}e^{\nu\rho}D_{x_0}$.

\emph{The solution.} Apply Theorem~\ref{is-onto-share-counts-section-kronecker-solve-hormander} to $U=\Omega$, the
function $e(z)=\rho(z)+\log(1+|z|^2)$, the weight $\Phi$, $F_0=\chi\circ\rho$, and $h$ equal to $K$ on $S$ and to $0$
elsewhere. The function $e$ is smooth, $\mathcal L_e\geq0$, and its sublevel sets in $\Omega$ are compact, as in the
proof of Lemma~\ref{is-onto-share-counts-section-kronecker-solve}, by
Lemma~\ref{is-onto-share-counts-section-kronecker-solve-chart}. The function $F_0$ is smooth on $\Omega$. As $\rho$ is
real, $\bar\partial_j(\chi\circ\rho)=\chi'(\rho)\,\overline{\partial_j\rho}$, so
\[
 \Big|\sum_j\bar\partial_jF_0(z)\,\bar\xi_j\Big|^2=\chi'(\rho(z))^2\,\Big|\sum_j\partial_j\rho(z)\,\xi_j\Big|^2 .
\]
For $z\in S$ this is at most $D_\chi R\,\mathcal L_\rho(z;\xi)\le K\,\mathcal L_\Phi(z;\xi)$, by
Lemma~\ref{is-onto-step-orbit-interpolation-large-gradient}; for $z\notin S$ it is $0$, as
$\rho(z)\notin[\varepsilon/2,\varepsilon]$. The theorem gives a holomorphic function $F$ on $\Omega$ with
\[
 \int_\Omega|F_0-F|^2e^{-\Phi}\,dz\le\int_\Omega h\,e^{-\Phi}\,dz
 =K\int_Se^{-\varphi_\lambda(c(z))-\psi(c(z))}e^{\nu\rho(z)}D_{x_0}(z)\,dz .
\]
The left side is the integral of the statement.

\emph{The bound.} Let $z\in S$. Then $-\varphi_\lambda(c(z))\le A$, as $\rho(z)\le\varepsilon$. By
Lemma~\ref{is-onto-step-orbit-interpolation-large-block}(1), $\rho(x_0,c(z))=\rho(z)$ lies in
$[\varepsilon/2,\varepsilon]$, so $-\psi(c(z))\le C$ by (5) of Definition~\ref{is-onto-step-orbit-interpolation-large-pole}.
And $\nu\rho(z)\le\nu\varepsilon$. So the right side is at most $Ke^{A+C+\nu\varepsilon}I_S=M$.
% Lean: poleSolve: exists_cutoff, exists_phla_bound, R = |Psi| min p0 (r - p0), K = D R / (m - kappa - bb), S =
% rhF^-1 [eps / 2, eps] with SublevelFin.muF_sublevel_lt_top, M = K exp(A + C + bb eps) muF(S); the weight phi
% (phla o cpt + psi o cpt - log densF - bb rhF) with IW.phla_cpt_cdAt, OntoB.logDens_contDiffAt,
% OntoB.rhF_contDiffAt, the Levi bound hphil and hexp; HE.hormander_open with eF (ExF.eF_contDiffAt,
% levi_eF_nonneg, through ExF.levi_rhF_nonneg of Lemma is-onto-share-counts-section-kronecker-solve-chart,
% ExF.isCompact_subF), F0 = chi o rhF with wdb_ofReal_comp, h = K on S, RhGradFStatement_holds and
% Up.leviRh_nonneg (Lemma is-onto-step-orbit-interpolation-large-solve-levi, for the weight phi); the last bound uses
% hA, the fifth field of IsPole, and lintegral_muF. The Lean's PoleSolve
% quantifies over every real delta; the book's pole functions have delta > 0.
\end{proof}

\begin{lemma}\label{is-onto-step-orbit-interpolation-large-near-zero}
Let $X$ consist of the data (1) to (6) of flag data of rank $r$ over $E$, use the notation of
Definition~\ref{is-onto-share-counts-section-levi}, write $|z|$ for the Euclidean norm of $z\in\mathbf C^n$, and let
$x_0$ be an orthonormal point. For $z\in\mathbf C^n$ let $S(z)$ be the sum of the $|z_{\tau,a,b}|^2$ over the coordinates
with $b\le p_0<a$. Then there are real numbers $\varrho>0$ and $C$ such that every $z$ with $|z|<\varrho$ lies in
$\Omega(x_0)$ and satisfies $\rho_{x_0}(z)\le C\,S(z)$.
\end{lemma}
\leanref{Development/AbelianVarietyTheory/Proofs/TorusHaarMeasure/ContinuousInvLeadingBlockEntry.lean}{GlobalTheory.SharpGradient.rhF_le_near_zero}
\leanref{Development/AbelianVarietyTheory/Proofs/TorusHaarMeasure/ContinuousInvLeadingBlockEntry.lean}{GlobalTheory.SharpGradient.neg_log_det_le}
\begin{proof}
\emph{A matrix bound.} Let $Z$ be a complex $p\times q$ matrix such that $1-ZZ^*$ is positive definite and
$\sum_{a,b}|Z_{ab}|^2\le\frac12$. Then $-\log\det(1-ZZ^*)\le2\sum_{a,b}|Z_{ab}|^2$. Indeed, by the spectral theorem there
is a unitary matrix $U$ such that $U^*(1-ZZ^*)U$ is diagonal. Its diagonal entries are $1-t_a$ with
$t_a=\sum_b|(U^*Z)_{ab}|^2$, and they are positive. As $U$ is unitary, $\sum_at_a=\sum_{a,b}|Z_{ab}|^2\le\frac12$, so
$0\le t_a\le\frac12$; and $\det(1-ZZ^*)=\prod_a(1-t_a)$. For $0\le t\le\frac12$,
$-\log(1-t)\le(1-t)^{-1}-1\le2t$, the last inequality as $(1+2t)(1-t)=1+t(1-2t)\geq1$.

\emph{Near $0$.} Use the matrices $N_\tau(z)$, $Y_\tau(z)$, $Z_\tau(z)$ and $R_\tau(z)$ of
Lemma~\ref{is-onto-share-counts-section-kronecker-solve-chart-block}. As $Y_\tau(0)=0$, $Z_\tau(0)=0$ and $R_\tau(0)=1$,
so $0\in\Omega(x_0)$ by part (3) of that lemma. The set $\Omega(x_0)$ is open, the entries of $Z_\tau(z)$ are
polynomials in $z$ by part (1), and so are those of $N_\tau(z)^{-1}$, the adjugate of the unitriangular matrix
$N_\tau(z)$; and $N_\tau(0)=1$. So there is $\varrho>0$ such that for $|z|<\varrho$: $z\in\Omega(x_0)$,
$\sum_{a,b}|Z_\tau(z)_{ab}|^2<\frac12$ for every $\tau$, and $|(N_\tau(z)^{-1})_{ba}|<|(N_\tau(0)^{-1})_{ba}|+1$ for all
$\tau$, $a$ and $b$.

Let $|z|<\varrho$ and $\tau\in\Psi$. The entries of $Y_\tau(z)$ are coordinates $z_{\tau,a,b}$ with $b\le p_0<a$, so their
absolute values are at most $S(z)^{1/2}$. Each entry of $Z_\tau(z)=(Y_\tau(z)N_\tau(z)^{-1})^{\mathsf T}$ is a sum of $p_0$
products of such an entry with an entry of $N_\tau(z)^{-1}$, so its absolute value is at most $(p_0+1)S(z)^{1/2}$. By
part (4) of Lemma~\ref{is-onto-share-counts-section-kronecker-solve-chart-block} and the matrix bound,
\[
 \rho_{x_0}(z)=-\sum_\tau\log\det R_\tau(z)\le2\sum_\tau\sum_{a,b}|Z_\tau(z)_{ab}|^2\le2|\Psi|p_0(r-p_0)(p_0+1)^2S(z).
\]
% Lean: rhF_le_near_zero: hZ0 and h0 (ExF.mem_OmF_iff), the eventual conditions e1 (ExF.isOpen_OmF), e2 (ExF.cd_Zt)
% and e3 (continuous_inv_leadingBlock_entry, through the adjugate), norm_Wm_sq_le for the entries of Y,
% rhF_eq_sum_re and neg_log_det_le (exists_unitary_rows with the eigenvectorUnitary of 1 - Z Z^H, conj_diag,
% Real.one_sub_inv_le_log_of_pos). The Lean's constant is the sum over tau of 2 sum_a sum_m (sum_b
% (|(N_tau(0)^-1)_ba| + 1))^2, which is the one above.
\end{proof}

\begin{lemma}\label{is-onto-step-orbit-interpolation-large-blowup-product}
Let $d\geq1$ and $k\geq0$ be integers and $\varepsilon>0$ and $M$ real numbers, write $|x|$ for the Euclidean norm,
and let $B_j$ be the open ball of radius $\varepsilon$ around $0$ in $\mathbf C^j$. Then there is a real number
$t_0>0$ such that for every real number $0<t\le t_0$ and every holomorphic function $u$ on $B_d\times B_k$
\[
 M|u(0)|^2\le\int_{B_d\times B_k}|u(x,\zeta)|^2\max(|x|^2,t)^{-d}\,dx\,d\zeta .
\]
\end{lemma}
\leanref{ReciprocityTheory/L2DbarEstimates/PlurisubharmonicCharts/APrimeA.lean}{ReciprocityTheory.L2DbarEstimates.Y.ap_prod_blowup}
\leanref{ReciprocityTheory/L2DbarEstimates/PlurisubharmonicCharts/APrimeA.lean}{ReciprocityTheory.L2DbarEstimates.Y.apra_blow}
\begin{proof}
\emph{Mean values.} Let $u$ be holomorphic on the open ball of radius $R$ around $0$ in $\mathbf C^j$, and $w\geq0$ a
measurable function on $\mathbf R$ that vanishes on $[R,\infty)$. Then
$|u(0)|^2\int_{\mathbf C^j}w(|y|)\,dy\le\int_{\mathbf C^j}|u(y)|^2w(|y|)\,dy$. Indeed, for $0<s<R$, $u(0)$ is the mean of
$u$ over the sphere $|y|=s$, so by the Cauchy-Schwarz inequality the mean of $|u|^2$ over it is at least $|u(0)|^2$;
integrate against $w(s)$ in polar coordinates.

\emph{The weight.} For $t>0$ let $w_t(s)=\max(s^2,t)^{-d}$ for $s<\varepsilon$ and $w_t(s)=0$ otherwise, and put
$I(t)=\int_{\mathbf C^d}w_t(|x|)\,dx$. As $t$ decreases, $w_t$ increases, and for $0<|x|<\varepsilon$, $w_t(|x|)$ tends to
$|x|^{-2d}$. By monotone convergence, $I(1/(j+1))$ tends to $\int_{|x|<\varepsilon}|x|^{-2d}\,dx$ as $j\to\infty$, and this
integral is infinite: in polar coordinates on $\mathbf C^d=\mathbf R^{2d}$ it is a positive multiple of
$\int_0^\varepsilon s^{-2d}s^{2d-1}\,ds$. Choose $j$ with $\operatorname{vol}(B_k)\,I(1/(j+1))>M$, and put $t_0=1/(j+1)$.

\emph{The bound.} Let $0<t\le t_0$ and $u$ holomorphic on $B_d\times B_k$. For $\zeta\in B_k$, the mean values for
$x\mapsto u(x,\zeta)$ and $w_t$ give
$|u(0,\zeta)|^2I(t)\le\int_{B_d}|u(x,\zeta)|^2\max(|x|^2,t)^{-d}\,dx$. For $\zeta\mapsto u(0,\zeta)$ and the indicator
function of $(-\infty,\varepsilon)$ they give $|u(0)|^2\operatorname{vol}(B_k)\le\int_{B_k}|u(0,\zeta)|^2\,d\zeta$. As
$I(t)\geq I(t_0)$,
\[
 M|u(0)|^2\le\operatorname{vol}(B_k)\,I(t)\,|u(0)|^2\le I(t)\int_{B_k}|u(0,\zeta)|^2\,d\zeta\le
 \int_{B_k}\int_{B_d}|u(x,\zeta)|^2\max(|x|^2,t)^{-d}\,dx\,d\zeta,
\]
which is the integral of the statement by Fubini's theorem.
% Lean: ap_prod_blowup: apraW d eps t for w_t, apraW_anti, apra_slice and apra_base (the mean value statement
% MVStatement, proved as L2DbarEstimates.mv), apra_blow (lintegral_iSup, and MVIntStatement d, the divergence of the
% integral of |y|^(-2d) over the unit ball, proved as L2DbarEstimates.mvInt), and lintegral_prod_symm.
\end{proof}

\begin{lemma}\label{is-onto-step-orbit-interpolation-large-blowup}
Let $n$ and $d\geq1$ be integers, write $|v|$ for the Euclidean norm and $B(v_0,\delta)$ for the open ball of radius
$\delta$ around $v_0$, and let $U\subset\mathbf C^n$ be an open set, $p\colon U\to\mathbf C^d$ a holomorphic map and
$v_0\in U$ a point with $p(v_0)=0$ at which the derivative of $p$ is surjective. Then there is a real number
$\delta>0$ such that the closed ball of radius $\delta$ around $v_0$ lies in $U$, and for every real number $M$ there is
a real number $t_0>0$ with
\[
 M|u(v_0)|^2\le\int_{B(v_0,\delta)}|u(v)|^2\max(|p(v)|^2,t)^{-d}\,dv
\]
for every $0<t\le t_0$ and every holomorphic function $u$ on $B(v_0,\delta)$.
\end{lemma}
\leanref{ReciprocityTheory/L2DbarEstimates/PlurisubharmonicCharts/APrimeLoc.lean}{ReciprocityTheory.L2DbarEstimates.Y.ap_local_blowup}
\begin{proof}
\emph{Coordinates.} Let $\Lambda$ be the derivative of $p$ at $v_0$. As $\Lambda$ is surjective, $d\le n$, and there is a
linear isomorphism $A\colon\mathbf C^n\to\mathbf C^d\times\mathbf C^{n-d}$ whose first component is $\Lambda$; write $L_2$ for
its second component. For $v\in U$ put $\Psi(v)=A^{-1}\big(p(v),L_2(v-v_0)\big)$. It is holomorphic, $\Psi(v_0)=0$, and its
derivative at $v_0$ is $A^{-1}A=1$. By the holomorphic inverse function theorem there is a holomorphic map $G$ on a
neighbourhood of $0$ with $\Psi(G(y))=y$ there and $G(\Psi(v))=v$ for $v$ near $v_0$. Let $C'$ be $1$ plus the absolute
value of the real Jacobian determinant of $\Psi$ at $v_0$. Choose $\delta>0$ such that the closed ball of radius
$\delta$ around $v_0$ lies in $U$, and $G(\Psi(v))=v$ and the real Jacobian determinant of $\Psi$ at $v$ has absolute
value less than $C'$ for $v\in B(v_0,\delta)$. Choose $r_1>0$ such that for $|y|<r_1$, $G$ is defined and holomorphic
near $y$, $\Psi(G(y))=y$ and $G(y)\in B(v_0,\delta)$. Choose $\varepsilon>0$ such that $A^{-1}$ maps
$\mathrm{Box}=B_d\times B_{n-d}$ into $B(0,r_1)$, where $B_j$ is the open ball of radius $\varepsilon$ around $0$ in
$\mathbf C^j$. There is a real number $K>0$ with
$\int_{\mathrm{Box}}h(A^{-1}y)\,dy=K\int_{A^{-1}(\mathrm{Box})}h(x)\,dx$ for every measurable $h\geq0$.

\emph{The bound.} Let $M$ be real, and let $t_0$ be the number of
Lemma~\ref{is-onto-step-orbit-interpolation-large-blowup-product} for $d$, $n-d$, $\varepsilon$ and $MKC'$. Let
$0<t\le t_0$, let $u$ be holomorphic on $B(v_0,\delta)$, and put $g(v)=|u(v)|^2\max(|p(v)|^2,t)^{-d}$. The function
$\tilde u(y)=u(G(A^{-1}y))$ is holomorphic on $\mathrm{Box}$, and $\tilde u(0)=u(v_0)$. For $y\in\mathrm{Box}$,
$A\Psi(G(A^{-1}y))=y$, and the first component of the left side is $p(G(A^{-1}y))$; so $p(G(A^{-1}y))=y_1$. Hence
\[
 MKC'|u(v_0)|^2\le\int_{\mathrm{Box}}|\tilde u(y)|^2\max(|y_1|^2,t)^{-d}\,dy=\int_{\mathrm{Box}}g(G(A^{-1}y))\,dy
 =K\int_Tg(G(x))\,dx
\]
with $T=A^{-1}(\mathrm{Box})$. Let $s$ be the set of the $v\in B(v_0,\delta)$ with $\Psi(v)\in T$. It is measurable,
$\Psi$ is injective on $s$, and $\Psi(s)=T$: for $x\in T$, $|x|<r_1$, so $G(x)\in B(v_0,\delta)$ and $\Psi(G(x))=x$. The
change of variables $x=\Psi(v)$ gives
\[
 \int_Tg(G(x))\,dx=\int_s|\det D_{\mathbf R}\Psi(v)|\,g(G(\Psi(v)))\,dv\le C'\int_{B(v_0,\delta)}g(v)\,dv .
\]
So $MKC'|u(v_0)|^2\le KC'\int_{B(v_0,\delta)}g(v)\,dv$, and dividing by $KC'>0$ gives the claim.
% Lean: ap_local_blowup: aprb_complement for A, Psi = A.symm (p v, Lsnd v - Lsnd v0),
% HasStrictFDerivAt.localInverse for G, the bound C on |det fderiv R Psi| near v0, the radius r of hrb and
% eps = r / (||A.symm|| + 1), exists_const_lintegral_comp_continuousLinearEquiv for K, ap_prod_blowup with
% M K C, and lintegral_image_eq_lintegral_abs_det_fderiv_mul for the change of variables.
\end{proof}

\begin{lemma}\label{is-onto-step-orbit-interpolation-large-point}
Let $X$ consist of the data (1) to (6) of flag data of rank $r$ over $E$, use the notation of
Definition~\ref{is-onto-share-counts-section-levi}, write $|z|$ for the Euclidean norm of $z\in\mathbf C^n$, and suppose
that $\dim X=|\Psi|p_0(r-p_0)\geq1$. Let $\varrho>0$ and $c>0$ be real numbers. Then there is a real number
$\delta_0>0$ such that for every orthonormal point $x_1$, every real number $0<\delta\le\delta_0$ and every holomorphic
function $u$ on the set $B_1$ of the $z\in\Omega(x_1)$ with $|z|<\varrho$
\[
 |u(0)|^2\le c\int_{B_1}|u(z)|^2\big(\rho_{x_1}(z)+\delta\big)^{-\dim X}D_{x_1}(z)\,dz .
\]
\end{lemma}
\leanref{Development/AbelianVarietyTheory/Definitions/UpperBounds/FreeBallIffOfDimXEqZero.lean}{GlobalTheory.Up.PolePoint}
\leanref{Development/AbelianVarietyTheory/Proofs/TorusHaarMeasure/ExistsMixedProj.lean}{GlobalTheory.Up.polePoint_holds}
\leanref{Development/AbelianVarietyTheory/Proofs/TorusHaarMeasure/ExistsMixedProj.lean}{GlobalTheory.Up.PoleValue.muF_eq}
\leanref{Development/AbelianVarietyTheory/Definitions/UpperBounds/FreeBallIffOfDimXEqZero.lean}{GlobalTheory.Up.card_mixedEnt}
\begin{proof}
If there is no orthonormal point, there is nothing to prove; let $x_0$ be one. For an orthonormal point $x_1$, the
Gram matrix of the columns of $c_{x_1}(z)_\tau=x_{1,\tau}\ell_\tau(z)$ is $\ell_\tau(z)^{\mathsf T}J\overline{\ell_\tau(z)}$, with
$J$ as in Lemma~\ref{is-onto-step-free-ball-minor}, and $\det c_{x_1}(z)_\tau=\det x_{1,\tau}\ne0$
(Lemma~\ref{is-onto-share-majorant}(4)). So $\Omega(x_1)=\Omega(x_0)$, $\rho_{x_1}=\rho_{x_0}$ and $D_{x_1}=D_{x_0}$, and it
suffices to find $\delta_0$ for $x_1=x_0$. Write $\Omega$, $\rho$ and $D$ for $\Omega(x_0)$, $\rho_{x_0}$ and $D_{x_0}$,
and $d=\dim X$.

\emph{The density.} As $c_{x_0}(0)=x_0$ lies in $\mathcal D_X$, $D$ is continuous and positive at $0$; so there are
$\varepsilon_0>0$ and $m>0$ with $D(z)\geq m$ for $|z|<\varepsilon_0$.

\emph{The potential.} Lemma~\ref{is-onto-step-orbit-interpolation-large-near-zero} gives $\varrho_1>0$ and $C$ such that
every $z$ with $|z|<\varrho_1$ lies in $\Omega$ and satisfies $\rho(z)\le CS(z)$, where $S(z)$ is the sum of the
$|z_{\tau,a,b}|^2$ over the coordinates with $b\le p_0<a$. And $\rho\geq0$ on $\Omega$
(Lemma~\ref{is-onto-share-counts-section-kronecker-solve-chart}). There are $d$ such coordinates: $|\Psi|$ choices of
$\tau$, $p_0$ of $b$ and $r-p_0$ of $a$. Let $L\colon\mathbf C^n\to\mathbf C^d$ be the projection to them; it is linear and
surjective, and $|Lz|^2=S(z)$.

\emph{The blowup.} Put $\varrho'=\min(\varrho,\varrho_1,\varepsilon_0)$ and $K=\max(C,0)+1$.
Lemma~\ref{is-onto-step-orbit-interpolation-large-blowup}, for the ball $|z|<\varrho'$, the map $L$ and $v_0=0$, gives
$\delta'>0$ such that the closed ball of radius $\delta'$ lies in the ball of radius $\varrho'$, and, for
$M=K^d/(cm)$, a number $t_0>0$. Put $\delta_0=t_0$.

Let $0<\delta\le\delta_0$ and $u$ holomorphic on $B_1$. The ball $|v|<\delta'$ lies in $B_1$, as $\varrho'\le\varrho$ and
$\varrho'\le\varrho_1$. So
\[
 \frac{K^d}{cm}|u(0)|^2\le\int_{|v|<\delta'}|u(v)|^2\max(|Lv|^2,\delta)^{-d}\,dv .
\]
For $|v|<\delta'$: $0<\rho(v)+\delta\le C|Lv|^2+\delta\le K\max(|Lv|^2,\delta)$, so
$K^{-d}\max(|Lv|^2,\delta)^{-d}\le(\rho(v)+\delta)^{-d}$; and $D(v)\geq m$. Hence
\[
 |u(0)|^2\le cm\int_{|v|<\delta'}|u(v)|^2(\rho(v)+\delta)^{-d}\,dv\le c\int_{B_1}|u(v)|^2(\rho(v)+\delta)^{-d}D(v)\,dv .
\]
% Lean: polePoint_holds: PoleValue.OmF_eq, rhF_eq and muF_eq for the reduction to x0; m = D(0) / 2 near 0 (hdens);
% SharpGradient.rhF_sharp_gradient, of which only rhF >= 0 and the part rhF_le_near_zero are used; L from
% card_mixedEnt (the number of mixed entries is dimX); Y.ap_local_blowup with L2DbarEstimates.mv and mvInt d,
% M = K^d / (c m) and K = max C 0 + 1; hpt and the comparison of measures hmu.
\end{proof}

\begin{lemma}\label{is-onto-step-orbit-interpolation-large-ball}
Let $X$ consist of the data (1) to (6) of flag data of rank $r$ over $E$, use the notation of
Definition~\ref{is-onto-share-counts-section-levi}, write $|z|$ for the Euclidean norm of $z\in\mathbf C^n$, and let
$\lambda$ be a weight, $\varepsilon>0$ a real number, and $\kappa\colon M_r(\mathbf C)^\Psi\to\mathbf R$ a function with
$\kappa>0$ on $\mathcal D_X$. Suppose that there are real numbers $\varrho_2>0$ and $C_\kappa$ with
$\kappa(c_x(z))\le C_\kappa\kappa(x)$ for every orthonormal point $x$ and every $z$ with $|z|<\varrho_2$. Then there are
real numbers $\varrho>0$ and $A>0$ such that for every orthonormal point $x$ and every $z$ with $|z|<\varrho$:
$z\in\Omega(x)$, $\rho_x(z)\le\varepsilon$, and
\[
 1\le A\,\kappa(x)\,\frac{N_\lambda(c_x(z))^2}{\kappa(c_x(z))}.
\]
\end{lemma}
\leanref{Development/AbelianVarietyTheory/Proofs/TorusHaarMeasure/SumRangeBlocks.lean}{GlobalTheory.Up.OrbitInterpolation.exists_ball}
\leanref{Development/AbelianVarietyTheory/Proofs/TorusHaarMeasure/DminCptEq.lean}{GlobalTheory.OrthNear.orthNear}
\begin{proof}
By (2) of Definition~\ref{is-onto-share-flag-data}, for every $\tau$ some basis of $\mathbf C^r$ has the Gram matrix
$J$ of Lemma~\ref{is-onto-step-free-ball-minor} for $h_\tau$; so there is an orthonormal point $x_*$. For an orthonormal
point $x$, the Gram matrix of the columns of $c_x(z)_\tau=x_\tau\ell_\tau(z)$ is $\ell_\tau(z)^{\mathsf T}J\overline{\ell_\tau(z)}$,
and $\det c_x(z)_\tau=\det x_\tau\ne0$ (Lemma~\ref{is-onto-share-majorant}(4)). So $z\in\Omega(x)$ if and only if
$z\in\Omega(x_*)$, and then $N_\lambda(c_x(z))=N_\lambda(c_{x_*}(z))$ and $\rho_x(z)=\rho_{x_*}(z)$.

The set $\Omega(x_*)$ is open and contains $0$ (Lemma~\ref{is-onto-share-majorant}(3)); choose $\varrho_0>0$ such that the
closed ball $\bar B$ of radius $\varrho_0$ around $0$ lies in it. The function $z\mapsto N_\lambda(c_{x_*}(z))$ is
continuous and positive on $\bar B$, so it is at least some $c>0$ there. The function $\rho_{x_*}$ is continuous at $0$,
with $\rho_{x_*}(0)=0$ as $|d_{p_0}(x_*,\tau)|=1$; choose $\varrho_3>0$ with $\rho_{x_*}(z)<\varepsilon$ for
$|z|<\varrho_3$. Put $\varrho=\min(\varrho_0,\varrho_2,\varrho_3)$ and $A=(|C_\kappa|+1)/c^2$.

Let $x$ be orthonormal and $|z|<\varrho$. Then $z\in\Omega(x)$ and $\rho_x(z)=\rho_{x_*}(z)\le\varepsilon$. As
$0<\kappa(c_x(z))\le C_\kappa\kappa(x)\le(|C_\kappa|+1)\kappa(x)$, the number $(|C_\kappa|+1)\kappa(x)/\kappa(c_x(z))$ is at
least $1$, and so is $N_\lambda(c_x(z))^2/c^2$. Their product is the right side of the claim.
% Lean: exists_ball with OrthNear.orthNear for the orthonormal point x0F built from X.signature
% (densF_eq_of_isOrth, nfac_cpt_eq_of_isOrth, cpt_mem_dom_iff_of_isOrth, the minima on the closed ball of radius
% rho0 / 2; its bounds on the density are not used here), OrthNear.dmin_cpt_eq and rhX_self for rho, and
% KronMetric.exp_neg_phla for exp(-phla) = N_lambda^2. A = (|C_kappa| + 1) / c^2.
\end{proof}

\begin{lemma}\label{is-onto-step-orbit-interpolation-large-approx}
Let $X$ consist of the data (1) to (6) of flag data of rank $r$ over $E$, use the notation of
Definitions~\ref{is-onto-share-counts-section-levi}, \ref{is-onto-share-counts-section-gauge}
and~\ref{is-onto-step-free-ball} and the gauge $W_\lambda$ of
Lemma~\ref{is-onto-share-counts-section-kronecker-gauge}, write $|z|$ for the Euclidean norm of $z\in\mathbf C^n$, and
suppose that $\dim X=|\Psi|p_0(r-p_0)\geq1$. Let $G\subset GL_r(E)$ be a subgroup, $x_0$ an orthonormal point,
$\lambda$ a weight with $\lambda_{\tau,i}=\lambda_{\tau,j}$ whenever $i,j\le p_0$ or $i,j>p_0$, and $\rho_0>0$, $\nu\geq0$ and
$m$ real numbers, such that $G$ has a free ball of radius $\rho_0$ at $x_0$, $\nu+\dim X/\rho_0<m$, and
$m\,\mathcal L_{\rho_{x_0}}(z;\xi)\le\mathcal L^{x_0}_\lambda(z;\xi)-\mathcal L_{\log D_{x_0}}(z;\xi)$ for every
$z\in\Omega(x_0)$ and $\xi\in\mathbf C^n$. Let $\kappa\colon M_r(\mathbf C)^\Psi\to\mathbf R$ be a measurable function with
$0<\kappa\le1$ on $\mathcal D_X$, $\kappa(xb)=\kappa(x)$ for $x\in\mathcal D_X$ and $b\in\mathcal B^\Psi$, and
$\kappa(c_{x_0}(z))=e^{-\nu\rho_{x_0}(z)}$ for $z\in\Omega(x_0)$, such that there are real numbers $\varrho_2>0$ and
$C_\kappa$ with $\kappa(c_x(z))\le C_\kappa\kappa(x)$ for every orthonormal point $x$ and every $z$ with $|z|<\varrho_2$.
Let $\eta>0$. Then there are a holomorphic function $F$ on $\Omega(x_0)$ that is square integrable for $W_\lambda$ and
$\nu$, and a function $P$ that is $\mathcal B$-equivariant of weight $\lambda$ with $P\circ c_{x_0}=F$ on $\Omega(x_0)$,
such that $|P(x_0)-1|\le\eta$ and $|P(gx_0)|\le\eta\,\kappa(gx_0)^{1/2}$ for every $g\ne1$ in $G$.
\end{lemma}
\leanref{Development/AbelianVarietyTheory/Proofs/TorusHaarMeasure/SumRangeBlocks.lean}{GlobalTheory.Up.OrbitInterpolation.exists_approx}
\leanref{Development/AbelianVarietyTheory/Proofs/TorusHaarMeasure/SumRangeBlocks.lean}{GlobalTheory.Up.OrbitInterpolation.lintegral_lt_top}
\leanref{Development/AbelianVarietyTheory/Proofs/TorusHaarMeasure/SumRangeBlocks.lean}{GlobalTheory.Up.OrbitInterpolation.exists_density}
\leanref{Development/AbelianVarietyTheory/Proofs/TorusHaarMeasure/SumRangeBlocks.lean}{GlobalTheory.Up.OrbitInterpolation.rpow_le_weight}
\leanref{Development/AbelianVarietyTheory/Proofs/TorusHaarMeasure/ContinuousOnOfCharts.lean}{GlobalTheory.KronMetric.continuousOn_of_charts}
\begin{proof}
Write $c=c_{x_0}$, $\Omega=\Omega(x_0)$, $\rho=\rho_{x_0}$ and $d=\dim X$. For an orthonormal point $x_1$,
$\Omega(x_1)=\Omega$, $\rho_{x_1}=\rho$ and $D_{x_1}=D_{x_0}$, as in the proof of
Lemma~\ref{is-onto-step-orbit-interpolation-large-point}; and for $g\in G$, $gx_0$ is orthonormal, as $g$ preserves $h$,
and $c_{gx_0}(z)=g\,c(z)$.

\emph{The choices.} As $m-\nu>d/\rho_0>0$, there is $\theta>1$ with $\theta d/\rho_0<m-\nu$.
Lemma~\ref{is-onto-step-orbit-interpolation-large-pole-exists}, for $\rho_0$ and $\theta$, gives $\varepsilon$ and $C$ with
$0<\varepsilon<\rho_0$, and for $0<\delta\le1$ a pole function for $(G,x_0,\theta d/\rho_0,\varepsilon,C,\delta)$.
Lemma~\ref{is-onto-step-orbit-interpolation-large-solve}, for $\theta d/\rho_0$, $\varepsilon$, $C$, $\nu$ and $m$, gives
$\chi$ and $M\geq0$. Lemma~\ref{is-onto-step-orbit-interpolation-large-ball}, for $\lambda$, $\varepsilon/2$ and $\kappa$,
gives $\varrho$ and $A$. Put $D_0=e^CAM$ and $c_0=\eta^2/(D_0+1)$, so that $c_0>0$ and $c_0D_0\le\eta^2$.
Lemma~\ref{is-onto-step-orbit-interpolation-large-point}, for $\varrho$ and $c_0$, gives $\delta_0$. Put
$\delta=\min(\delta_0,1)$, let $\psi$ be the pole function for $\delta$, and let $F$ be the holomorphic function on
$\Omega$ of Lemma~\ref{is-onto-step-orbit-interpolation-large-solve} for $\delta$ and $\psi$:
\[
 \int_\Omega\big|\chi(\rho(z))-F(z)\big|^2e^{-\varphi_\lambda(c(z))-\psi(c(z))}e^{\nu\rho(z)}D_{x_0}(z)\,dz\le M.\tag{I}
\]

\emph{Square integrability.} Let $A_1$ be the number of Lemma~\ref{is-onto-step-orbit-interpolation-large-block}(2)
for $\varepsilon$ and $\lambda$. For $z\in\Omega$, $|F|^2\le2|\chi(\rho)|^2+2|\chi(\rho)-F|^2$; $|\chi|\le1$, and
$\chi(\rho(z))=0$ if $\rho(z)\geq\varepsilon$; and $e^{-\psi(c(z))}\geq1$ by (2) of
Definition~\ref{is-onto-step-orbit-interpolation-large-pole}. So
\[
 |F(z)|^2e^{-\varphi_\lambda(c(z))}e^{\nu\rho(z)}\le2e^{A_1+\nu\varepsilon}[\rho(z)\le\varepsilon]
 +2\big|\chi(\rho(z))-F(z)\big|^2e^{-\varphi_\lambda(c(z))-\psi(c(z))}e^{\nu\rho(z)},
\]
where $[\rho(z)\le\varepsilon]$ is $1$ or $0$. Integrating against $D_{x_0}(z)\,dz$, with
Lemma~\ref{is-onto-step-orbit-interpolation-large-sublevel}(2) for $a=\varepsilon$ and (I), and as
$e^{-\varphi_\lambda}=N_\lambda^2=W_\lambda$ on $\mathcal D_X$, $F$ is square integrable for $W_\lambda$ and $\nu$.

\emph{The function $P$.} As $\nu\geq0$ and $\rho\geq0$ on $\Omega$, $e^{\nu\rho}\geq1$. So for every compact set
$C'\subset\mathcal D_X$ the integral of $|F(z)|^2N_\lambda(c(z))^2D_{x_0}(z)$ over the $z\in\Omega$ with
$c(z)\in C'\mathcal B^\Psi$ is finite, and Lemma~\ref{is-onto-share-counts-section-kronecker-extend} gives a function $P$
that is $\mathcal B$-equivariant of weight $\lambda$ with $P\circ c=F$ on $\Omega$.

\emph{Comparison with the pole.} Let $g\in G$, $x_1=gx_0$, $z$ with $|z|<\varrho$, and $s\geq0$. By
Lemma~\ref{is-onto-step-orbit-interpolation-large-ball}, $z\in\Omega$, $\rho(z)\le\varepsilon/2$, and
$1\le A\kappa(x_1)N_\lambda(c_{x_1}(z))^2/\kappa(c_{x_1}(z))$. By
Lemma~\ref{is-onto-step-orbit-interpolation-large-block}(1), $\rho(x_0,c(z))=\rho(z)\le\varepsilon$; so (1) and (4) of
Definition~\ref{is-onto-step-orbit-interpolation-large-pole} give
$\psi(c_{x_1}(z))=\psi(c(z))\le d\log(\rho(z)+\delta)+C$, that is, $(\rho(z)+\delta)^{-d}\le e^Ce^{-\psi(c_{x_1}(z))}$. Hence
\[
 s\,(\rho(z)+\delta)^{-d}\le e^CA\kappa(x_1)\cdot s\,\frac{e^{-\varphi_\lambda(c_{x_1}(z))-\psi(c_{x_1}(z))}}{\kappa(c_{x_1}(z))}.
 \tag{II}
\]

\emph{The value at $x_0$.} Lemma~\ref{is-onto-step-orbit-interpolation-large-point}, for $x_1=x_0$ and $u=1-F$, gives
$|1-F(0)|^2\le c_0\int|1-F(z)|^2(\rho(z)+\delta)^{-d}D_{x_0}(z)\,dz$, over the $z\in\Omega$ with $|z|<\varrho$. For these
$z$, $\chi(\rho(z))=1$ and $\kappa(c(z))=e^{-\nu\rho(z)}$, so (II) for $g=1$ and $s=|1-F(z)|^2$ bounds the integrand by
$e^CA\kappa(x_0)$ times the integrand of (I). Hence $|1-F(0)|^2\le c_0e^CA\kappa(x_0)M\le c_0D_0\le\eta^2$, as
$\kappa(x_0)\le1$. As $P(x_0)=P(c(0))=F(0)$, $|P(x_0)-1|\le\eta$.

\emph{The values at $gx_0$.} Let $g\ne1$ in $G$ and $x_1=gx_0$. The function $\psi$ is continuous on $\mathcal D_X$: near
$x\in\mathcal D_X$ write $x=x'b'$ with $x'$ orthonormal (Lemma~\ref{is-onto-share-majorant}(3)), and let $\Delta$,
$z(\cdot)$ and $b(\cdot)$ be as in Lemma~\ref{is-onto-share-counts-section-transfer-lu} for $x'$; on the open set of the
$y\in\mathcal D_X$ with $\Delta(y)\ne0$, which contains $x$, $y=c_{x'}(z(y))b(y)$, so $\psi(y)=\psi(c_{x'}(z(y)))$ by (1)
of Definition~\ref{is-onto-step-orbit-interpolation-large-pole}, a continuous function by (3). Let $T$ be the set of the
$y\in\mathcal D_X$ with $\rho(x_0,y)\geq\varepsilon$, and put
\[
 \Phi(y)=|P(y)|^2N_\lambda(y)^2e^{-\psi(y)}/\kappa(y)\quad(y\in T),\qquad\Phi(y)=0\quad(y\notin T).
\]
It is measurable, as $P$, $N_\lambda$, $\psi$ and $\rho(x_0,\cdot)$ are continuous on the open set $\mathcal D_X$ and
$\kappa$ is measurable; and $\Phi(yb)=\Phi(y)$ for $y\in\mathcal D_X$ and $b\in\mathcal B^\Psi$: $\rho(x_0,yb)=\rho(x_0,y)$
(Lemma~\ref{is-onto-step-free-ball-bound}(1)), $|P(yb)|=|b^\lambda||P(y)|$, $N_\lambda(yb)=|b^\lambda|^{-1}N_\lambda(y)$ (proof
of Lemma~\ref{is-onto-share-counts-section-transfer}), and $\psi$ and $\kappa$ are invariant. Let $z\in\Omega$. If
$c(z)\in T$, then $\rho(z)=\rho(x_0,c(z))\geq\varepsilon$, so $\chi(\rho(z))=0$, $P(c(z))=F(z)$ and
$\kappa(c(z))=e^{-\nu\rho(z)}$, and $\Phi(c(z))D_{x_0}(z)$ is the integrand of (I); otherwise $\Phi(c(z))=0$. So
$\int_\Omega\Phi(c(z))D_{x_0}(z)\,dz\le M$, and by Lemma~\ref{is-onto-share-counts-section-poincare-measure}(1) the same
bound holds for $\int_\Omega\Phi(c_{x_1}(z))D_{x_0}(z)\,dz$.

Lemma~\ref{is-onto-step-orbit-interpolation-large-point}, for $x_1$ and the function $u=P\circ c_{x_1}$, holomorphic on
$\Omega$ as $P$ is holomorphic on $\mathcal D_X$, gives
$|P(x_1)|^2\le c_0\int|P(c_{x_1}(z))|^2(\rho(z)+\delta)^{-d}D_{x_0}(z)\,dz$, over the $z\in\Omega$ with $|z|<\varrho$. For
these $z$, $c(z)=g^{-1}c_{x_1}(z)$ lies in $\mathcal D_X$ and $\rho(x_0,c(z))=\rho(z)\le\varepsilon/2<\rho_0$; as $g\ne1$ and
$G$ has a free ball of radius $\rho_0$ at $x_0$, $\rho(x_0,c_{x_1}(z))=\rho(x_0,g\,c(z))\geq\rho_0>\varepsilon$. So
$c_{x_1}(z)\in T$, and (II) for $s=|P(c_{x_1}(z))|^2$ bounds the integrand by $e^CA\kappa(x_1)\Phi(c_{x_1}(z))D_{x_0}(z)$.
Hence $|P(gx_0)|^2\le c_0e^CA\kappa(x_1)M\le\eta^2\kappa(gx_0)$.
% Lean: exists_approx: theta from exists_between; PoleEx (Lemma is-onto-step-orbit-interpolation-large-pole-exists),
% PoleSolve (Lemma is-onto-step-orbit-interpolation-large-solve) with kappa = dimX theta / rho0, exists_ball (Lemma
% is-onto-step-orbit-interpolation-large-ball) for eps / 2, D = exp C A M, c = eta^2 / (D + 1), PolePoint (Lemma
% is-onto-step-orbit-interpolation-large-point) for rho and c, delta = min delta0 1; lintegral_lt_top (with
% exp_neg_phla_le and SublevelFin.muF_sublevel_lt_top) for the square integrability; ExtFStatement (Lemma
% is-onto-share-counts-section-kronecker-extend) for P; hw with rpow_le_weight for (II), where
% IsPole.le_log_orbit reads the pole bound at lact g x0 through rhX_lact_lact (the book reads it at x0 through
% the invariance of psi, the same step); value_le with PolePoint at x0 for the value at x0; and, for the values at
% g x0, KronMetric.kronMetric with psi (continuousOn_of_charts for the continuity of psi, and the gauge
% exp(-(phla + psi)) on X.dom), exists_density (OntoE.phiF_meas, OntoE.phiF_B, measInv_holds), the third field of
% FreeBall with rhX_lact, and value_le with PolePoint at g x0.
\end{proof}

\begin{lemma}\label{is-onto-step-orbit-interpolation-large-interpolation}
Let $X$ consist of the data (1) to (6) of flag data of rank $r$ over $E$, such that $\Gamma$ acts properly
discontinuously on $\mathcal D_X$ modulo $\mathcal B^\Psi$ (Lemma~\ref{is-onto-share-counts-section-discontinuous}), use
the notation of Definitions~\ref{is-onto-share-counts-section-levi} and~\ref{is-onto-step-free-ball}, and suppose that
$\dim X=|\Psi|p_0(r-p_0)\geq1$. Let $G\subset GL_r(E)$ be a subgroup, $x_0$ an orthonormal point, $\lambda$ a weight with
$\lambda_{\tau,i}=\lambda_{\tau,j}$ whenever $i,j\le p_0$ or $i,j>p_0$, and $\rho_0>0$, $\nu$ and $m$ real numbers, such that
$G$ has a free ball of radius $\rho_0$ at $x_0$, $4|\Psi|r^2<\nu$, $\nu+\dim X/\rho_0<m$,
$m\,\mathcal L_{\rho_{x_0}}(z;\xi)\le\mathcal L^{x_0}_\lambda(z;\xi)-\mathcal L_{\log D_{x_0}}(z;\xi)$ for every
$z\in\Omega(x_0)$ and $\xi\in\mathbf C^n$, and every bounded section of weight $\lambda$ is a cusp form. Then there is a
cusp form $P\in S_\lambda(X)$ with $P(x_0)\ne0$ and $P(\gamma x_0)=0$ for every $\gamma\in G$ outside $\Gamma$.
\end{lemma}
\leanref{Development/AbelianVarietyTheory/Definitions/UpperBounds/FreeBallIffOfDimXEqZero.lean}{GlobalTheory.Up.OrbitInterp}
\leanref{Development/AbelianVarietyTheory/Proofs/TorusHaarMeasure/SumRangeBlocks.lean}{GlobalTheory.Up.OrbitInterpolation.orbitInterp}
\begin{proof}
Let $\kappa$ and $C_s$ be as in Lemma~\ref{is-onto-step-orbit-interpolation-large-weight} for $x_0$ and $\nu$, and let
$W_\lambda$ be the gauge of Lemma~\ref{is-onto-share-counts-section-kronecker-gauge}, which is bounded below. The set
$S_\lambda(X)$ is a complex vector space (Lemma~\ref{is-onto-share-sections}(1)). As $\Gamma$ has finite index in the
group of the elements of $GL_r(E)$ that preserve $h$ and $\Lambda$, there is a finite set $T$ of such elements such that
each of them is $s\gamma$ for some $s\in T$ and $\gamma\in\Gamma$. Let $T'$ be the set of the $s\in T$ with $s^{-1}\in G$ and
$s^{-1}\notin\Gamma$. The elements of $G$ preserve $h$ and $\Lambda$, so the points $s^{-1}x_0$ lie in $\mathcal D_X$.

\emph{Approximation.} Let $\eta>0$ and $\eta'=\eta/(C_s+1)$. Lemma~\ref{is-onto-step-orbit-interpolation-large-approx},
for $\kappa$ and $\eta'$, gives $F$ and $P$; it applies as $\nu\geq0$. By
Lemma~\ref{is-onto-share-counts-section-poincare}, for $W_\lambda$, $\nu$, $F$ and $P$, the families
$(\chi_\lambda(\gamma)^{-1}P(\gamma x))_{\gamma\in\Gamma}$ are absolutely summable for $x\in\mathcal D_X$, and $Q=\mathcal P_\lambda P$
is a bounded section of weight $\lambda$, hence a cusp form. As $\chi_\lambda$ has finite order on $\Gamma$,
$|\chi_\lambda(\gamma)|=1$, and $\chi_\lambda(1)=1$. Every $\gamma\ne1$ in $\Gamma$ lies in $G$, so
\[
 |Q(x_0)-1|\le|P(x_0)-1|+\sum_{\gamma\ne1}|P(\gamma x_0)|\le\eta'+\eta'\sum_\gamma\kappa(\gamma x_0)^{1/2}\le\eta'(1+C_s)=\eta .
\]
Let $s\in T'$. For $\gamma\in\Gamma$, $\gamma s^{-1}\in G$, and $\gamma s^{-1}\ne1$, as $s^{-1}\notin\Gamma$. So
$|Q(s^{-1}x_0)|\le\sum_\gamma|P(\gamma s^{-1}x_0)|\le\eta'\sum_\gamma\kappa(\gamma s^{-1}x_0)^{1/2}\le\eta'C_s\le\eta$.

\emph{Interpolation.} Lemma~\ref{is-onto-step-orbit-interpolation-large-span}, for $S_\lambda(X)$, the point $x_0$ and
the points $s^{-1}x_0$ with $s\in T'$, gives $u\in S_\lambda(X)$ with $u(x_0)=1$ and $u(s^{-1}x_0)=0$ for $s\in T'$. Let
$\gamma\in G$ outside $\Gamma$. As $\gamma^{-1}$ preserves $h$ and $\Lambda$, $\gamma^{-1}=s\gamma_1$ with $s\in T$ and
$\gamma_1=s^{-1}\gamma^{-1}\in\Gamma$. Then $s^{-1}=\gamma_1\gamma\in G$, and $s^{-1}\notin\Gamma$, since otherwise
$\gamma=\gamma_1^{-1}s^{-1}\in\Gamma$; so $s\in T'$. As $u$ is a section of weight $\lambda$,
$u(\gamma x_0)=u(\gamma_1^{-1}s^{-1}x_0)=\chi_\lambda(\gamma_1^{-1})u(s^{-1}x_0)=0$. So $P=u$ works.
% Lean: orbitInterp: kaFrame (Lemma is-onto-step-orbit-interpolation-large-weight), KronMetric.kronMetric with
% psi = 0 (Lemma is-onto-share-counts-section-kronecker-gauge), siegelSpaceAlgebra_holds for the submodule S,
% X.Gamma_index for T, exists_of_approx (Lemma is-onto-step-orbit-interpolation-large-span) over the subtype of the
% s in T with s^-1 in G and not in Gamma, exists_approx (Lemma is-onto-step-orbit-interpolation-large-approx) with
% eta / (Cs + 1), PoinStatement (poin_holds, Lemma is-onto-share-counts-section-poincare) and the hypothesis hSS,
% tsum_eq_add_tsum_ite, chi_one, OntoE.chi_norm and tsum_le_of_le; at the end the section law of u
% (hu'.1.1.2.2) at the element s^-1 gamma^-1 of Gamma.
\end{proof}

\begin{lemma}\label{is-onto-step-orbit-interpolation-large-curvature}
Let $X$ consist of the data (1) to (6) of flag data of rank $r$ over $E$ with $0<p_0<r$, use the notation of
Definition~\ref{is-onto-share-counts-section-levi}, and let $s\geq0$ be an even integer and $\delta_s$ the weight with
$(\delta_s)_{\tau,i}=-s/2$ for $i\le p_0$ and $(\delta_s)_{\tau,i}=s/2$ for $i>p_0$. Then for every orthonormal point
$x_0$, every $z\in\Omega(x_0)$ and every $\xi\in\mathbf C^n$
\[
 (s-2r+2)\,\mathcal L_{\rho_{x_0}}(z;\xi)\le\mathcal L^{x_0}_{\delta_s}(z;\xi)-\mathcal L_{\log D_{x_0}}(z;\xi).
\]
\end{lemma}
\leanref{Development/AbelianVarietyTheory/Proofs/TorusHaarMeasure/DivisionWithBasisOfDivision.lean}{GlobalTheory.DivisionVolume.margin_iff}
\leanref{Development/AbelianVarietyTheory/Proofs/DivisorProofs/OrbitInterpolationDetWeightFreeBallIndefinite.lean}{GlobalTheory.OrbitInterpWeight.gapsPairing_deMul}
\leanref{Development/AbelianVarietyTheory/Proofs/TorusHaarMeasure/BasePsiLeviLowerBound.lean}{GlobalTheory.FrameThetaCompare.margin_of_gaps}
\leanref{Development/AbelianVarietyTheory/Proofs/TorusHaarMeasure/BasePsiLeviLowerBound.lean}{GlobalTheory.FrameThetaCompare.curv_lin}
\begin{proof}
Let $k$ be the weight with $k_{\tau,i}=2i-r-1$, and $\lambda=\delta_s-k$. Let $z\in\Omega(x_0)$ and $\xi\in\mathbf C^n$.

\emph{The gaps.} Let $\tau\in\Psi$. For $1\le i<r$ with $i\ne p_0$, $\delta_s$ takes the same value at $i$ and $i+1$, so
$\lambda_{\tau,i+1}=\lambda_{\tau,i}-2\le\lambda_{\tau,i}$. As $1\le p_0<r$,
$\lambda_{\tau,r}-\lambda_{\tau,1}=(s/2-(r-1))-(-s/2-(1-r))=s-2r+2$.

\emph{The weight $k$.} Lemma~\ref{is-onto-share-counts-section-transfer}, for $y=x_0$, gives an orthonormal point $x_1$
and a linear map $A$ with $\mathcal L^{x_0}_\kappa(z;\xi)=\mathcal L^{x_1}_\kappa(0;A\xi)$ for every weight $\kappa$, and by
Lemma~\ref{is-onto-share-counts-section-diagonal}(2) the right side is linear in $\kappa$. So
$\mathcal L^{x_0}_\lambda(z;\xi)=\mathcal L^{x_0}_{\delta_s}(z;\xi)-\mathcal L^{x_0}_k(z;\xi)$. Let $w\in\Omega(x_0)$ and
$\tau\in\Psi$, and write $L_i=\log|d_i(c_{x_0}(w),\tau)|$. Then $L_0=0$, and $L_r=0$ as in the proof of
Lemma~\ref{is-onto-share-counts-section-curvature}. As $k_{\tau,i+1}-k_{\tau,i}=2$, summation by parts gives
\[
 \sum_{i=1}^rk_{\tau,i}(L_i-L_{i-1})=k_{\tau,r}L_r-k_{\tau,1}L_0-\sum_{i=1}^{r-1}(k_{\tau,i+1}-k_{\tau,i})L_i=-2\sum_{i=0}^{r-1}L_i.
\]
Summing over $\tau$, $\varphi_k(c_{x_0}(w))=\log D_{x_0}(w)$. As $\Omega(x_0)$ is open, $\varphi_k\circ c_{x_0}$ and
$\log D_{x_0}$ have the same Levi form at $z$. Hence
$\mathcal L^{x_0}_\lambda(z;\xi)=\mathcal L^{x_0}_{\delta_s}(z;\xi)-\mathcal L_{\log D_{x_0}}(z;\xi)$.

\emph{The bound.} By the gaps, induction on $j-i$ gives $\lambda_{\tau,j}\le\lambda_{\tau,i}$ whenever $i\le j\le p_0$ or
$p_0<i\le j$. For $i\le p_0<j$ it follows that $\lambda_{\tau,i}\le\lambda_{\tau,1}$ and $\lambda_{\tau,r}\le\lambda_{\tau,j}$, so
$\lambda_{\tau,j}-\lambda_{\tau,i}\geq\lambda_{\tau,r}-\lambda_{\tau,1}=s-2r+2$. So
Lemma~\ref{is-onto-share-counts-section-curvature}, for $\lambda$ and $m=s-2r+2$, gives
$(s-2r+2)\mathcal L_{\rho_{x_0}}(z;\xi)\le\mathcal L^{x_0}_\lambda(z;\xi)$, which is the claim by the previous step.
% Lean: (margin_iff h0 h1 (deMul s) (s - 2 r + 2)).2 (gapsPairing_deMul X h0 h1 hs). gapsPairing_deMul is the
% step "the gaps": GapsPairing (deMul s - kapW) (s - 2 r + 2). The second direction of margin_iff first rewrites the
% right side with leviWeight_sub_leviDensity_eq_lemma (the step "the weight k": FrameThetaCompare.curv_lin, through
% transfer, levi_transfer, LeviZero.levi_phla_zero and diagForm_lin, and phla_kapW_cpt_eq_log_densF, the summation
% by parts sum_range_centered_mul_sub_eq with norm_dmin_cpt_last_eq_one, with leviV_congr on the open set OmF x0),
% then applies FrameThetaCompare.margin_of_gaps (the induction mono, then CurvMargin.curv_and_margin, Lemma
% is-onto-share-counts-section-curvature). The step "the weight k" is the case c = 0 of Lemma
% is-onto-share-counts-section-density, which margin_iff proves again this way.
\end{proof}

\begin{lemma}\label{is-onto-step-orbit-interpolation-large}
In Situation~\ref{is-onto-step-orbit-situation}, suppose that $p_0(r-p_0)\geq4$. Then there is a cusp form
$P\in S_{\delta_s}(X)$ with $P(x_0)\ne0$ and $P(\gamma x_0)=0$ for every $\gamma\in G$ outside $\Gamma$.
\end{lemma}
\leanref{Development/AbelianVarietyTheory/Proofs/DivisorProofs/OrbitInterpolationDetWeightFreeBallIndefinite.lean}{GlobalTheory.OrbitInterpWeight.large}
\begin{proof}
Put $d=\dim X=|\Psi|p_0(r-p_0)$. There is an embedding $E\to\mathbf C$, and by (1) of
Definition~\ref{is-onto-share-flag-data} it or its conjugate lies in $\Psi$; so $\Psi$ is not empty. As $0<p_0<r$ and
$p_0(r-p_0)\geq4$, $d\geq4|\Psi|\geq1$. Put $\rho_0=1/(49r^2)$, $\nu=4|\Psi|r^2+1$ and $m=s-2r+2$. Then $d/\rho_0=49r^2d$,
$4|\Psi|r^2\le r^2d$, and $m\geq50r^2d+3$ as $s\geq2r+50r^2d+1$; so
\[
 \nu+\frac d{\rho_0}\le r^2d+1+49r^2d<m .
\]
The hypotheses of Lemma~\ref{is-onto-step-orbit-interpolation-large-interpolation} hold for $G$, $x_0$, $\lambda=\delta_s$,
$\rho_0$, $\nu$ and $m$. The group $\Gamma$ acts properly discontinuously on $\mathcal D_X$ modulo $\mathcal B^\Psi$, by
Lemma~\ref{is-onto-share-counts-section-discontinuous}. The weight $\delta_s$ is constant on $i\le p_0$ and on $i>p_0$.
The group $G$ has a free ball of radius $\rho_0$ at $x_0$, and $4|\Psi|r^2<\nu$. The curvature inequality is
Lemma~\ref{is-onto-step-orbit-interpolation-large-curvature}. Every bounded section of weight $\delta_s$ is a cusp form, by
Lemma~\ref{is-onto-share-counts-section-cusp} for any $\tau_0\in\Psi$: $(\delta_s)_{\tau_0,i+1}=(\delta_s)_{\tau_0,i}$ for
$1\le i<r$ with $i\ne p_0$, and $(\delta_s)_{\tau_0,1}=-s/2<s/2=(\delta_s)_{\tau_0,r}$, as $s$ is even and positive. So
Lemma~\ref{is-onto-step-orbit-interpolation-large-interpolation} gives $P$.
% Lean: OrbitInterpWeight.large: hd (one_le_dimX_iff, through dimX_eq_zero_iff, which shows Psi nonempty from
% X.cmType), hDq (dimX >= 4 |Psi|), and Up.OrbitInterpolation.orbitInterp (Lemma
% is-onto-step-orbit-interpolation-large-interpolation) with bb = 4 |Psi| r^2 + 1, m = s - 2 r + 2 and
% rho0 = 1 / (49 r^2), ClosedDiscontinuity.propDisc_closed, blockConst_deMul, the curvature hcurv from margin_iff and
% gapsPairing_deMul (Lemma is-onto-step-orbit-interpolation-large-curvature), and cuspStatement_holds at tau in Psi.
% The inputs of orbitInterp are the proved statements normFactor_rightBorel_holds, normFactor_leftAction_holds,
% pointwiseNorm_invariance_holds (inside Lemmas is-onto-share-counts-section-kronecker-gauge and -transfer),
% poin_holds, extFStatement_holds, measInv_holds, PoleExistence.poleEx_holds, PoleDbar.poleSolve and
% Up.polePoint_holds. The hypothesis of the situation that chi_kappa is trivial on Gamma is not used.
% Size: nothing below this proof remains open; the last open step was Lemma
% is-onto-step-orbit-interpolation-large-pole-exists (about 3,400 lines of Lean). The whole of OrbitInterpWeight.large has about 6,500 lines in 450 declarations below it, by a text count through namespaced
% names that skips common short names and stops at statements the book cites; an earlier, looser text count gave
% 28,000 lines in 1,760 declarations.
\end{proof}

\subsection{Small signatures and the indefinite case}\label{subsection-is-small-signatures}

\begin{lemma}\label{is-onto-step-orbit-interpolation-small-integral-coordinates}
Let $X$ consist of the data (1) to (6) of flag data of rank $r$ over $E$, use the notation of
Definition~\ref{is-onto-share-counts-section-levi} and the blocks $N_\tau(z)$, $Y_\tau(z)$ and $N'_\tau(z)$ of $\ell_\tau(z)$ of
Lemma~\ref{is-onto-share-counts-section-kronecker-solve-chart-block}, write $M^*$ for the conjugate transpose of a
complex matrix $M$, and let $x_0$ be an orthonormal point. For $z\in\mathbf C^n$ let $T(z)\in\mathbf C^n$ have the
coordinates of $z$, except that $Y_\tau(T(z))=Y_\tau(z)N_\tau(z)$ for every $\tau\in\Psi$. Let $U$ be the set of the
$z\in\mathbf C^n$ such that every $Y_\tau(z)$ has operator norm less than $1$ for the Euclidean norms.
\begin{enumerate}
\item $T$ is a bijection of $\mathbf C^n$ onto itself that preserves Lebesgue measure, and
  \[
   c_{x_0}(T(z))_\tau=x_{0,\tau}\begin{pmatrix}N_\tau(z)&0\\ Y_\tau(z)N_\tau(z)&N'_\tau(z)\end{pmatrix}
   \qquad(z\in\mathbf C^n,\ \tau\in\Psi).
  \]
\item $T(z)\in\Omega(x_0)$ if and only if $z\in U$. For $z\in U$, every $1-Y_\tau(z)^*Y_\tau(z)$ is positive definite, and
  $\rho_{x_0}(T(z))=-\sum_\tau\log\det\big(1-Y_\tau(z)^*Y_\tau(z)\big)$.
\item For every real number $s$,
  \[
   \int_{\Omega(x_0)}e^{-s\rho_{x_0}(z)}D_{x_0}(z)\,dz=\int_Ue^{-s\rho_{x_0}(T(z))}D_{x_0}(T(z))\,dz .
  \]
\end{enumerate}
\end{lemma}
\leanref{Development/AbelianVarietyTheory/Proofs/TorusHaarMeasure/LeadingBlockLowerEntriesToBlockCoords.lean}{GlobalTheory.UpCoordCompare.toEntries}
\leanref{Development/AbelianVarietyTheory/Proofs/TorusHaarMeasure/LeadingBlockLowerEntriesToBlockCoords.lean}{GlobalTheory.UpCoordCompare.measurePreserving_toEntries}
\leanref{Development/AbelianVarietyTheory/Proofs/TorusHaarMeasure/LeadingBlockLowerEntriesToBlockCoords.lean}{GlobalTheory.UpCoordCompare.toEntries_mem_OmF_iff}
\leanref{Development/AbelianVarietyTheory/Proofs/TorusHaarMeasure/LeadingBlockLowerEntriesToBlockCoords.lean}{GlobalTheory.UpCoordCompare.rhF_toEntries}
\leanref{Development/AbelianVarietyTheory/Proofs/TorusHaarMeasure/LeadingBlockLowerEntriesToBlockCoords.lean}{GlobalTheory.UpCoordCompare.map_mu}
\leanref{Development/AbelianVarietyTheory/Proofs/DivisorProofs/LintegralExpRhoEqSetLIntegralUnitBallDomainStatement.lean}{GlobalTheory.T.lintegral_expRho_eq_setLIntegral_unitBallDomain}
\begin{proof}
Write $N=N_\tau(z)$ and $Y=Y_\tau(z)$.

(1) The coordinates of $N_\tau$ and $N'_\tau$ are not changed by $T$, so $\ell_\tau(T(z))$ is the matrix displayed, and
$c_{x_0}(T(z))_\tau=x_{0,\tau}\ell_\tau(T(z))$. As $N$ is lower unitriangular, $N^{-1}$ is its adjugate, a polynomial in $z$;
so the map that replaces every $Y_\tau(z)$ by $Y_\tau(z)N_\tau(z)^{-1}$ is a polynomial inverse of $T$. The coordinate
$(\tau,a,b)$ of $T(z)-z$ is $0$ unless $b\le p_0<a$, and then it is $\sum_{b<c\le p_0}z_{\tau,a,c}z_{\tau,c,b}$, which involves
only coordinates $(\tau,a',b')$ with $a'-b'<a-b$. Ordering the coordinates by $a-b$, the real derivative of $T$ at every
point is $1+K$ with $K$ strictly triangular, so its determinant is $1$, and $T$ preserves Lebesgue measure by the change of
variables formula.

(2) By (1), $N_\tau(T(z))=N$, so $Z_\tau(T(z))=(YNN^{-1})^{\mathsf T}=Y^{\mathsf T}$ and $R_\tau(T(z))=1-Y^{\mathsf T}\bar Y$, the
transpose of $1-Y^*Y$; it is positive definite if and only if $1-Y^*Y$ is, and has the same determinant. As
$v^*(1-Y^*Y)v=|v|^2-|Yv|^2$, the matrix $1-Y^*Y$ is positive definite if and only if $|Yv|<|v|$ for every $v\ne0$, which holds
if and only if the operator norm of $Y$ is less than $1$, by the compactness of the unit sphere. Parts (3) and (4) of
Lemma~\ref{is-onto-share-counts-section-kronecker-solve-chart-block} give the claims.

(3) By (2), $T^{-1}(\Omega(x_0))=U$, and by (1) the substitution $z\mapsto T(z)$ preserves Lebesgue measure.
% Lean: T is UpCoordCompare.toEntries = coordOf (nFrame z), nFrame = zFrame * leviFrame (nFrame_apply,
% cpt_toEntries: cpt x0 (toEntries z) = point x0 z); the inverse is ofEntries (coordEquiv, ext_blocks, Wm_toEntries,
% Wm_ofEntries); measurePreserving_toEntries is measurePreserving_of_triangular with the level a - b
% (det_one_add_of_isNilpotent, map_withDensity_abs_det_fderiv_eq_addHaar). (2) is toEntries_mem_OmF_iff (Am_toEntries,
% opNorm_lt_one_iff, ExF.mem_OmF_iff) and rhF_toEntries; U is ExhaustionFrames.Om, defined by the operator norm of
% zBlock. (3) is lintegral_exp_muF_eq_setLIntegral_Om through map_mu and densF_toEntries.
\end{proof}

\begin{lemma}\label{is-onto-step-orbit-interpolation-small-integral-minors}
Let $m,k\geq0$ be integers, write $M^*$ for the conjugate transpose of a complex matrix $M$ and $|\cdot|$ for the
Euclidean norm, and let $M$ be a complex $m\times k$ matrix.
\begin{enumerate}
\item Let $y\in\mathbf C^m$ be a column with $|y|<1$, put $A=\begin{pmatrix}1&0\\ y&M\end{pmatrix}$, an $(m+1)\times(k+1)$ matrix,
  and let $J'$ be the diagonal matrix with entries $1,-1,\dots,-1$, of size $m+1$. Then
  $(1-|y|^2)\det(M^*M)\le|\det(A^*J'A)|$.
\item Let $q\in\mathbf C^m$ be a row. Then $(1-|q|^2)\det(M^*M)\le\big|\det\big(M^*M-(qM)^*(qM)\big)\big|$.
\end{enumerate}
\end{lemma}
\leanref{Development/AbelianVarietyTheory/Proofs/DivisorProofs/UnitBallLeadingGramMinorLowerBoundFirst.lean}{GlobalTheory.IntBallMinorOne.schur_bound}
\leanref{Development/AbelianVarietyTheory/Proofs/DivisorProofs/UnitBallLeadingGramMinorLowerBoundFirst.lean}{GlobalTheory.IntBallMinorOne.norm_det_le_add_rankOne}
\leanref{Development/AbelianVarietyTheory/Proofs/DivisorProofs/UnitBallLeadingGramMinorLowerBoundLast.lean}{GlobalTheory.IntBall.MinorLast.det_sub_bound}
\begin{proof}
The matrix $Q=M^*M$ is positive semidefinite, so $\det Q\geq0$.

\emph{A rank one bound.} Let $v\in\mathbf C^k$ and $c\geq0$. Then $\det Q\le|\det(Q+cvv^*)|$. Indeed, if $\det Q\ne0$, then
$Q$ is positive definite, and $\det(Q+cvv^*)=\det Q\,(1+c\,v^*Q^{-1}v)$ with $v^*Q^{-1}v\geq0$; otherwise $\det Q=0$.

(1) Put $t=1-|y|^2>0$. Then
\[
 A^*J'A=\begin{pmatrix}t&-y^*M\\ -M^*y&-Q\end{pmatrix},
\]
and the Schur complement of the entry $t$ gives $\det(A^*J'A)=t\det\big(-Q-t^{-1}M^*yy^*M\big)=(-1)^kt\det(Q+t^{-1}vv^*)$ with
$v=M^*y$. The rank one bound gives the claim.

(2) If $\det Q=0$, the claim holds. Otherwise $Q$ is positive definite. Put $u=qM$ and $\Pi=MQ^{-1}M^*$. Then
$\det(Q-u^*u)=\det Q\,(1-uQ^{-1}u^*)=\det Q\,(1-q\Pi q^*)$. As $\Pi$ is Hermitian with $\Pi^2=\Pi$, the matrix
$1-\Pi=(1-\Pi)^*(1-\Pi)$ is positive semidefinite, so $q\Pi q^*\le qq^*=|q|^2$, and
$\det(Q-u^*u)\geq(1-|q|^2)\det Q$.
% Lean: (1) is IntBallMinorOne.schur_bound (det_fromBlocks11 for the Schur complement, hrk for the rank one term) with
% norm_det_le_add_rankOne (det_add_replicateCol_mul_replicateRow, the matrix determinant lemma). (2) is
% IntBall.MinorLast.det_sub_bound (Matrix.det_add_mul, hPP, hQ and the projection P).
\end{proof}

\begin{lemma}\label{is-onto-step-orbit-interpolation-small-integral-density}
Let $X$ consist of the data (1) to (6) of flag data of rank $r$ over $E$ with $0<p_0<r$ and $p_0\in\{1,r-1\}$, use the
notation of Lemma~\ref{is-onto-step-orbit-interpolation-small-integral-coordinates}, and let $x_0$ be an orthonormal point
and $s\geq2(r-1)$ a real number. For a complex square matrix $B$ write $\gamma_j(B)$ for the leading principal $j\times j$
minor of $B^*B$. For $z\in\mathbf C^n$ and $\tau\in\Psi$ let $B_\tau(z)$ be $N'_\tau(z)$ if $p_0=1$ and $N_\tau(z)$ otherwise, a
lower unitriangular $(r-1)\times(r-1)$ matrix. Then $\gamma_j(B_\tau(z))>0$ for $0\le j\le r-1$, and for $z\in U$
\[
 D_{x_0}(T(z))\,e^{-s\rho_{x_0}(T(z))}\le\prod_{\tau\in\Psi}\prod_{j=0}^{r-2}\gamma_j(B_\tau(z))^{-2}.
\]
\end{lemma}
\leanref{Development/AbelianVarietyTheory/Proofs/DivisorProofs/ExistsUniformBoundLintegralFlagDomainBallCase.lean}{GlobalTheory.IntBall.density_mul_exp_le}
\leanref{Development/AbelianVarietyTheory/Proofs/DivisorProofs/UnitBallLeadingGramMinorLowerBoundFirstStatement.lean}{GlobalTheory.T.unitBall_leadingGramMinor_lower_bound_first}
\leanref{Development/AbelianVarietyTheory/Proofs/DivisorProofs/UnitBallLeadingGramMinorLowerBoundLastStatement.lean}{GlobalTheory.T.unitBall_leadingGramMinor_lower_bound_last}
\leanref{Development/AbelianVarietyTheory/Proofs/DivisorProofs/DetOneSubOffDiagonalBlockGramLeOneBallStatement.lean}{GlobalTheory.T.detOneSubOffDiagonalBlockGram_le_one_ball}
\leanref{Development/AbelianVarietyTheory/Proofs/DivisorProofs/ColumnGramDeterminantDiagonalBlockUnitLowerPosStatement.lean}{GlobalTheory.T.columnGramDeterminant_diagonalBlockUnitLower_pos}
\leanref{Development/AbelianVarietyTheory/Proofs/DivisorProofs/ColumnGramDeterminantUnitLowerBlockLastEqOneStatement.lean}{GlobalTheory.T.columnGramDeterminant_unitLowerBlock_last_eq_one}
\begin{proof}
\emph{The minors $\gamma_j$.} Let $B=B_\tau(z)$. As $B$ is lower unitriangular, its first $j$ columns are linearly
independent, and $\gamma_j(B)$ is the determinant of their Gram matrix, which is positive definite; so $\gamma_j(B)>0$.
Also $\gamma_0(B)=1$ and $\gamma_{r-1}(B)=|\det B|^2=1$.

\emph{The minors of the Gram matrix.} Let $z\in U$ and $\tau\in\Psi$, write $N=N_\tau(z)$, $Y=Y_\tau(z)$ and $N'=N'_\tau(z)$,
put $t=\det(1-Y^*Y)$ and $d_i=d_i(c_{x_0}(T(z)),\tau)$, and let $J$ be as in Lemma~\ref{is-onto-step-free-ball-minor}. By
Lemma~\ref{is-onto-step-orbit-interpolation-small-integral-coordinates}(1),
$c_{x_0}(T(z))_\tau=x_{0,\tau}P$ with $P=\begin{pmatrix}N&0\\ YN&N'\end{pmatrix}$. As $x_0$ is orthonormal, the Gram matrix of
the columns of $c_{x_0}(T(z))_\tau$ is $P^{\mathsf T}J\bar P$, so $|d_i|=|\det(P_i^*JP_i)|$, where $P_i$ is the matrix of the first
$i$ columns of $P$.

If $p_0=1$, then $N=(1)$, $B=N'$, and $Y$ is a column $y$ of length $r-1$ with $|y|<1$, as $z\in U$; so $t=1-|y|^2$. For
$0\le k\le r-2$, $P_{k+1}=\begin{pmatrix}1&0\\ y&M\end{pmatrix}$ with $M$ the matrix of the first $k$ columns of $B$, and
Lemma~\ref{is-onto-step-orbit-interpolation-small-integral-minors}(1) gives $t\,\gamma_k(B)\le|d_{k+1}|$.

If $p_0=r-1$, then $N=B$, $N'=(1)$, and $Y$ is a row $q$ of length $r-1$, with $t=\det(1-q^*q)=1-|q|^2$. For
$0\le k\le r-2$, $P_{k+1}=\begin{pmatrix}M\\ qM\end{pmatrix}$ with $M$ the matrix of the first $k+1$ columns of $B$, and
$P_{k+1}^*JP_{k+1}=M^*M-(qM)^*(qM)$; Lemma~\ref{is-onto-step-orbit-interpolation-small-integral-minors}(2) gives
$t\,\gamma_{k+1}(B)\le|d_{k+1}|$.

\emph{The bound.} By Lemma~\ref{is-onto-step-orbit-interpolation-small-integral-coordinates}(2), $1-Y^*Y$ is positive
definite, so $t>0$; and $t\le1$, as $t$ is $1-|y|^2$ or $1-|q|^2$. Let $g_k$ be $\gamma_k(B)$ if $p_0=1$ and $\gamma_{k+1}(B)$
otherwise, so that $t\,g_k\le|d_{k+1}|$ for $0\le k\le r-2$. As $d_0=1$, $t\le1$ and $s\geq2(r-1)$,
\[
 \prod_{i=0}^{r-1}|d_i|^{-2}\cdot t^s=\prod_{k=0}^{r-2}|d_{k+1}|^{-2}\cdot t^s\le t^{-2(r-1)}t^s\prod_{k=0}^{r-2}g_k^{-2}
 \le\prod_{k=0}^{r-2}g_k^{-2}=\prod_{j=0}^{r-2}\gamma_j(B)^{-2},
\]
the last equality as $\gamma_0(B)=\gamma_{r-1}(B)=1$ when $p_0\ne1$. By
Lemma~\ref{is-onto-step-orbit-interpolation-small-integral-coordinates}(2), $e^{-s\rho_{x_0}(T(z))}=\prod_\tau t_\tau^s$ with
$t_\tau=\det(1-Y_\tau(z)^*Y_\tau(z))$, and $D_{x_0}(T(z))=\prod_\tau\prod_{i=0}^{r-1}|d_i(c_{x_0}(T(z)),\tau)|^{-2}$; the product
over $\tau$ of the bounds gives the claim.
% Lean: density_mul_exp_le (exp_neg_rho, prod_fin_split, norm_dmin_zero, prod_inv_sq_mul_rpow_le), with
% T.detOneSubOffDiagonalBlockGram_le_one_ball (IntBall.zDet_le_one, re_det_one_sub_le_of_unique), the positivity
% T.columnGramDeterminant_diagonalBlockUnitLower_pos (ColGramPos.colGram_blkW_pos), the last minor
% T.columnGramDeterminant_unitLowerBlock_last_eq_one (colGram_blkW_last), and the minor bounds
% T.unitBall_leadingGramMinor_lower_bound_first (IntBallMinorOne.zDet_mul_colGram_le_norm_dmin: mat_nFrame for p0 = 1,
% gram_mul_right, the blocks Zm and Nm, schur_bound) and T.unitBall_leadingGramMinor_lower_bound_last
% (MinorLast.zDet_mul_colGram_le_dmin: n_top, n_last, det_sub_bound). B_tau(z) is blkW p0 (zOf z tau), gamma_j is
% colGram, t is zDet; the functions density and rho of ExhaustionFrames are D_x0 o T and rho_x0 o T.
\end{proof}

\begin{lemma}\label{is-onto-step-orbit-interpolation-small-integral-blocks}
Let $2\le r\le4$ and $0<p_0<r$ be integers with $p_0\in\{1,r-1\}$, use the minors $\gamma_j$ of
Lemma~\ref{is-onto-step-orbit-interpolation-small-integral-density}, and put $m=r(r-1)/2$. For
$w=(w_{ab})_{1\le b<a\le r}$ in $\mathbf C^m$ let $L(w)$ be the lower unitriangular $r\times r$ matrix with the entries $w_{ab}$
below the diagonal, and $B(w)$ its lower right $(r-1)\times(r-1)$ block if $p_0=1$ and its upper left $(r-1)\times(r-1)$ block
otherwise. Put $h(w)=\prod_{j=0}^{r-2}\gamma_j(B(w))^{-2}$ if $|w_{ab}|\le1$ for every $b\le p_0<a$, and $h(w)=0$ otherwise.
Then $\int_{\mathbf C^m}h(w)\,dw<\infty$ for Lebesgue measure.
\end{lemma}
\leanref{Development/AbelianVarietyTheory/Proofs/DivisorProofs/ExistsUniformBoundLintegralFlagDomainBallCase.lean}{GlobalTheory.IntBall.block_fin}
\leanref{Development/AbelianVarietyTheory/Proofs/DivisorProofs/LintegralCutoffGramBoundMixedBlockTwoLtTopStatement.lean}{GlobalTheory.T.lintegral_cutoffGramBoundMixedBlock_two_lt_top}
\leanref{Development/AbelianVarietyTheory/Proofs/DivisorProofs/LintegralCutoffGramBoundMixedBlockRankThreeLtTopStatement.lean}{GlobalTheory.T.lintegral_cutoffGramBoundMixedBlock_rankThree_lt_top}
\leanref{Development/AbelianVarietyTheory/Proofs/DivisorProofs/LintegralCutoffGramBoundMixedBlockRankFourLtTopStatement.lean}{GlobalTheory.T.lintegral_cutoffGramBoundMixedBlock_rankFour_lt_top}
\leanref{Development/AbelianVarietyTheory/Proofs/DivisorProofs/FlagGramDensityInequalityComplexStatement.lean}{GlobalTheory.T.flagGram_density_inequality_complex}
\leanref{Development/AbelianVarietyTheory/Proofs/DivisorProofs/ColumnGramDeterminantUnitLowerThreeStatement.lean}{GlobalTheory.T.columnGramDeterminant_unitLower_three}
\leanref{Development/AbelianVarietyTheory/Proofs/DivisorProofs/LintegralUpdateSubShearStatement.lean}{GlobalTheory.T.lintegral_update_sub_shear'}
\begin{proof}
The coordinates $w_{ab}$ with $b\le p_0<a$ are the \emph{mixed} ones. For a real number $\sigma>1$,
$\int_{\mathbf C}(1+|t|^2)^{-\sigma}\,dt=\pi/(\sigma-1)$ is finite, by polar coordinates. Put $\gamma_j=\gamma_j(B(w))$.

\emph{$r=2$.} Then $B(w)=(1)$, $h$ is the indicator function of the disc $|w_{21}|\le1$, and its integral is $\pi$.

\emph{$r=3$.} Then $B(w)=\begin{pmatrix}1&0\\ a&1\end{pmatrix}$ with $a=w_{32}$ if $p_0=1$ and $a=w_{21}$ if $p_0=2$; the other two
coordinates are mixed. So $\gamma_0=1$, $\gamma_1=1+|a|^2$, and $h(w)\le\prod_{b<a'}\phi(w_{a'b})$ with
$\phi(t)=[|t|\le1]+(1+|t|^2)^{-2}$, where $[|t|\le1]$ is $1$ if $|t|\le1$ and $0$ otherwise. As $\int_{\mathbf C}\phi=2\pi$,
Fubini's theorem gives $\int h\le(2\pi)^3$.

\emph{$r=4$.} Then $B(w)=\begin{pmatrix}1&0&0\\ a&1&0\\ b&c&1\end{pmatrix}$, with $(a,b,c)=(w_{32},w_{42},w_{43})$ if $p_0=1$ and
$(a,b,c)=(w_{21},w_{31},w_{32})$ if $p_0=3$; the other three coordinates are mixed. The Gram matrix of the first two
columns $(1,a,b)$ and $(0,1,c)$ of $B(w)$ gives $\gamma_0=1$, $\gamma_1=1+|a|^2+|b|^2$ and $\gamma_2=1+|c|^2+|ac-b|^2$. Put
$\omega(t)=(1+|t|^2)^{-4/3}$. Then
\[
 (\gamma_1\gamma_2)^{-2}\le\omega(a)\,\omega(c)\big(\omega(b)+\omega(b-ac)\big).\tag{$*$}
\]
Indeed, for real $\alpha,\beta\geq1$, $\alpha^2\beta\le\alpha^3+\beta^3-1$: if $\alpha\le\beta$, then $\beta^3\geq\alpha^2\beta$ and
$\alpha^3\geq1$, and if $\beta\le\alpha$, then $\alpha^3\geq\alpha^2\beta$ and $\beta^3\geq1$. Put $A=|a|^2$, $B=|b|^2$, $C=|c|^2$ and
$D=|ac-b|^2$, and suppose $B\le D$. With $\alpha=(1+A)^{1/3}$, $\beta=(1+B)^{1/3}$ and $\alpha'=(1+C)^{1/3}$,
$\alpha^2\beta\le1+A+B=\gamma_1$ and $\alpha'^2\beta\le1+C+B\le\gamma_2$, so
$(\gamma_1\gamma_2)^{-2}\le\alpha^{-4}\beta^{-4}\alpha'^{-4}=\omega(a)\omega(b)\omega(c)$. If $D\le B$, the same argument with $(C,D,A,B)$ in
place of $(A,B,C,D)$ gives $(\gamma_1\gamma_2)^{-2}\le\omega(c)\,\omega(b-ac)\,\omega(a)$.

Let $V(w)$ be $\omega(a)\omega(b)\omega(c)$ if every mixed coordinate has absolute value at most $1$, and $0$ otherwise, and
let $\sigma(w)$ be $w$ with the coordinate $b$ replaced by $b-ac$. By $(*)$, $h\le V+V\circ\sigma$. For fixed values of
the other coordinates, $\sigma$ translates the coordinate $b$, which is not $a$ or $c$; so $\sigma$ preserves Lebesgue
measure by Fubini's theorem, and $\int V\circ\sigma=\int V=\pi^3\big(\int_{\mathbf C}\omega\big)^3=\pi^3(3\pi)^3$. Hence
$\int h<\infty$.
% Lean: block_fin, by cases on p0 and r (interval_cases). r = 2 is lintegral_hB_one_two_lt_top (gB 1 = 1, the closed
% ball). r = 3 is Fin3.lintegral_hB_lt_top (gB_one, gB_two, hB_le with dom = 1_ball + (1 + |x|^2)^(-2),
% lintegral_dom_lt_top through integrable_rpow_neg_one_add_norm_sq, RestrMeas.lintegral_fintype_prod_eq_prod).
% r = 4 is Fin4.lintegral_hB_four_lt_top (lintegral_hB_lt_top_of_pos with pa, pb, pc; gB_eq_entries through
% T.columnGramDeterminant_unitLower_three; pw = (1 + |t|^2)^(-4/3); hpt through T.flagGram_density_inequality_complex,
% proved as flagGram_density_ineq with FlagGram.key and sq_mul_le_cubes; T.lintegral_update_sub_shear', proved by a
% skew product). h is hB p0 and prod_j gamma_j^-2 is gB p0.
\end{proof}

\begin{lemma}\label{is-onto-step-orbit-interpolation-small-integral}
Let $X$ consist of the data (1) to (6) of flag data of rank $r$ over $E$ with $0<p_0<r$ and $p_0(r-p_0)\le3$, use the
notation of Definition~\ref{is-onto-share-counts-section-levi}, and let $s\geq2(r-1)$ be a real number. Then there is a
real number $I$ such that $\int_{\Omega(x_1)}e^{-s\rho_{x_1}(z)}D_{x_1}(z)\,dz\le I$ for every orthonormal point $x_1$.
\end{lemma}
\leanref{Development/AbelianVarietyTheory/Proofs/DivisorProofs/ExistsUniformBoundLintegralFlagDomainBallCaseStatement.lean}{GlobalTheory.T.exists_uniform_bound_lintegral_flagDomain_ball_case}
\leanref{Development/AbelianVarietyTheory/Proofs/DivisorProofs/ExistsUniformBoundLintegralFlagDomainBallCase.lean}{GlobalTheory.IntBall.intBall}
\begin{proof}
By (2) of Definition~\ref{is-onto-share-flag-data} there is an orthonormal point $x_0$. For every orthonormal point $x_1$,
$\Omega(x_1)=\Omega(x_0)$, $\rho_{x_1}=\rho_{x_0}$ and $D_{x_1}=D_{x_0}$, as in the proof of
Lemma~\ref{is-onto-step-orbit-interpolation-large-point}; so it suffices to bound the integral for $x_1=x_0$. As
$0<p_0<r$ and $p_0(r-p_0)\le3$, $p_0=1$ or $p_0=r-1$, since otherwise $p_0\geq2$ and $r-p_0\geq2$; and then $r\le4$.

Use the notation of Lemma~\ref{is-onto-step-orbit-interpolation-small-integral-coordinates}. For $z\in\mathbf C^n$ and
$\tau\in\Psi$ let $z_\tau=(z_{\tau,a,b})_{b<a}$, so that $B_\tau(z)=B(z_\tau)$ in the notation of
Lemmas~\ref{is-onto-step-orbit-interpolation-small-integral-density}
and~\ref{is-onto-step-orbit-interpolation-small-integral-blocks}, and let $h$ be the function of
Lemma~\ref{is-onto-step-orbit-interpolation-small-integral-blocks}. For $z\in U$ the entries of $Y_\tau(z)$, which are the
coordinates $z_{\tau,a,b}$ with $b\le p_0<a$, have absolute value at most the operator norm of $Y_\tau(z)$, which is less
than $1$. So, by Lemmas~\ref{is-onto-step-orbit-interpolation-small-integral-coordinates}(3)
and~\ref{is-onto-step-orbit-interpolation-small-integral-density},
\[
 \int_{\Omega(x_0)}e^{-s\rho_{x_0}(z)}D_{x_0}(z)\,dz\le\int_U\prod_{\tau\in\Psi}\prod_{j=0}^{r-2}\gamma_j(B_\tau(z))^{-2}\,dz
 \le\int_{\mathbf C^n}\prod_{\tau\in\Psi}h(z_\tau)\,dz=\Big(\int h(w)\,dw\Big)^{|\Psi|},
\]
the last equality by Fubini's theorem. By Lemma~\ref{is-onto-step-orbit-interpolation-small-integral-blocks} this is a
finite number $I$.
% Lean: T.exists_uniform_bound_lintegral_flagDomain_ball_case = IntBall.intBall: PoleValue.rhF_eq and muF_eq for the
% reduction to x0F (OrthNear.isOrth_x0F), and IntBall.lintegral_lt_top: ball_cases, T.lintegral_expRho_eq_setLIntegral_unitBallDomain
% (Lemma is-onto-step-orbit-interpolation-small-integral-coordinates(3)), density_mul_exp_le (Lemma
% is-onto-step-orbit-interpolation-small-integral-density), T.ofReal_gramDeterminantBound_le_prod_cutoff (Cutoff.ofReal_bnd_le,
% with norm_entry_le_opNorm for the entries of zBlock), T.lintegral_prod_chartCoordinates_eq_prod (lintegral_prod_zOf),
% T.measurable_cutoffGramBoundMixedBlock, and block_fin (Lemma is-onto-step-orbit-interpolation-small-integral-blocks).
% Size: about 1,750 lines of Lean below exists_uniform_bound_lintegral_flagDomain_ball_case (text count through
% namespaced names, stopping at statements the book cites). Nothing below this proof remains open.
\end{proof}

\begin{definition}\label{is-onto-step-orbit-interpolation-small-orbit-sum-fundamental}
Let $X$ consist of the data (1) to (6) of flag data of rank $r$ over $E$, use the notation of
Definition~\ref{is-onto-share-counts-section-levi}, and let $x_0$ be an orthonormal point. A \emph{fundamental set} for
$\Gamma$ at $x_0$ is a measurable set $F\subset\Omega(x_0)$ such that $z=z'$ whenever $z,z'\in F$, $\gamma\in\Gamma$ and
$b\in\mathcal B^\Psi$ satisfy $c_{x_0}(z')=\gamma c_{x_0}(z)b$, and such that for almost every $z\in\Omega(x_0)$ there are
$z'\in F$, $\gamma\in\Gamma$ and $b\in\mathcal B^\Psi$ with $c_{x_0}(z)=\gamma c_{x_0}(z')b$.
\end{definition}
\leanref{ReciprocityTheory/ReductiveGroups/ArithmeticGeometry/ArakelovVolumes/UpperBounds/FundSetMeasStatement.lean}{ReciprocityTheory.ChartRestriction.G.AR.Up.IsFundSet}

\begin{lemma}\label{is-onto-step-orbit-interpolation-small-orbit-sum-fundamental-exists}
Let $X$ consist of the data (1) to (6) of flag data of rank $r$ over $E$, such that $\Gamma$ acts properly
discontinuously on $\mathcal D_X$ modulo $\mathcal B^\Psi$ (Lemma~\ref{is-onto-share-counts-section-discontinuous}), and
let $x_0$ be an orthonormal point. Then there is a fundamental set for $\Gamma$ at $x_0$
(Definition~\ref{is-onto-step-orbit-interpolation-small-orbit-sum-fundamental}).
\end{lemma}
\leanref{ReciprocityTheory/ReductiveGroups/ArithmeticGeometry/ArakelovVolumes/UpperBounds/FundSetMeasStatement.lean}{ReciprocityTheory.ChartRestriction.G.AR.Up.FundSetStatement}
\leanref{ReciprocityTheory/ReductiveGroups/ArithmeticGeometry/ArakelovVolumes/UpperBounds/FundSetPf.lean}{ReciprocityTheory.ChartRestriction.G.AR.Up.fundSet_holds}
\leanref{ReciprocityTheory/ReductiveGroups/ArithmeticGeometry/ArakelovVolumes/UpperBounds/FundSetPf.lean}{ReciprocityTheory.ChartRestriction.G.AR.Up.FundSetPf.fundSet_main}
\leanref{ReciprocityTheory/ReductiveGroups/ArithmeticGeometry/ArakelovVolumes/UpperBounds/FundSetPf.lean}{ReciprocityTheory.ChartRestriction.G.AR.Up.FundSetPf.small}
\leanref{ReciprocityTheory/ReductiveGroups/ArithmeticGeometry/ArakelovVolumes/UpperBounds/FundSetPf.lean}{ReciprocityTheory.ChartRestriction.G.AR.Up.FundSetPf.charB}
\leanref{ReciprocityTheory/ReductiveGroups/ArithmeticGeometry/ArakelovVolumes/UpperBounds/FundSetPf.lean}{ReciprocityTheory.ChartRestriction.G.AR.Up.FundSetPf.meas_Sat}
\begin{proof}
Write $c=c_{x_0}$ and $\Omega=\Omega(x_0)$. For $z,z'\in\Omega$ write $z\sim z'$ if $\gamma c(z)=c(z')b$ for some $\gamma\in\Gamma$
and $b\in\mathcal B^\Psi$. As $\Gamma$ and $\mathcal B^\Psi$ are groups and their actions commute, $\sim$ is an equivalence
relation.

\emph{Triangular quotients.} Let $y,y'\in M_r(\mathbf C)^\Psi$ have invertible components, and let $\gamma\in GL_r(E)$. Then
$\gamma y=y'b$ for some $b\in\mathcal B^\Psi$ if and only if every matrix $\operatorname{adj}(y'_\tau)\gamma^\tau y_\tau$ is upper
triangular: if $\gamma y=y'b$, this matrix is $\det(y'_\tau)\,b_\tau$; conversely, $b_\tau=y'^{-1}_\tau\gamma^\tau y_\tau$ is then
invertible and upper triangular.

\emph{Small sets.} Every $z_0\in\Omega$ has an open neighbourhood $O$ such that $z=z'$ whenever $z,z'\in O\cap\Omega$ and
$z\sim z'$. Indeed, Lemma~\ref{is-onto-share-counts-section-discontinuous} at $c(z_0)$ and $c(z_0)$ gives neighbourhoods
$U$ and $U'$ of $c(z_0)$ and a finite set $\Phi\subset\Gamma$ that contains every $\gamma\in\Gamma$ with $\gamma u=u'b$ for some
$u\in U$, $u'\in U'$ and $b\in\mathcal B^\Psi$. For $\gamma\in\Phi$ with $\gamma\ne1$, $\gamma c(z_0)\ne c(z_0)b$ for every $b$, by (4)
of Definition~\ref{is-onto-share-flag-data}; so by the previous step some entry below the diagonal of some
$\operatorname{adj}(c(z_0)_\tau)\gamma^\tau c(z_0)_\tau$ is nonzero. As $c$ is continuous and $\Phi$ is finite, there is an open
set $O\ni z_0$ with $c(O)\subset U\cap U'$ such that, for all $z,z'\in O$ and $\gamma\in\Phi$ with $\gamma\ne1$, the same entry of
$\operatorname{adj}(c(z')_\tau)\gamma^\tau c(z)_\tau$ is nonzero. Let $z,z'\in O\cap\Omega$ with $\gamma c(z)=c(z')b$. If $\gamma=1$,
then $z=z'$ by Lemma~\ref{is-onto-share-counts-section-transfer-lu}(3) for $x_0$. Otherwise $\gamma\in\Phi$, and the previous
step contradicts the choice of $O$.

\emph{Saturations.} For a measurable set $A\subset\Omega$, the set $\operatorname{Sat}(A)$ of the $z\in\Omega$ with $z\sim z'$
for some $z'\in A$ is measurable. Indeed, $\Gamma$ is countable, as $E$ is. For $\gamma\in\Gamma$ let $T_\gamma$ be the set of the
pairs $(z,z')\in\Omega\times A$ such that every $\operatorname{adj}(c(z')_\tau)\gamma^\tau c(z)_\tau$ is upper triangular; it is
measurable, and by the first step $\operatorname{Sat}(A)$ is the union over $\gamma$ of the projections of the $T_\gamma$ to the
first factor. The projection is injective on $T_\gamma$: if $(z,z_1)$ and $(z,z_2)$ lie in $T_\gamma$, then
$c(z_1)b_1=\gamma c(z)=c(z_2)b_2$ for some $b_1,b_2\in\mathcal B^\Psi$, and $z_1=z_2$ by
Lemma~\ref{is-onto-share-counts-section-transfer-lu}(3). By the Lusin-Souslin theorem \cite[Theorem~15.1]{Kec95}, the
image of a Borel set under a continuous map that is injective on it is Borel; so each projection is measurable.

\emph{The set.} The open set $\Omega$ is covered by countably many of the sets $O$ of the second step, say $O_1,O_2,\dots$,
as $\mathbf C^n$ has a countable basis of open sets. Put $W_k=O_k\cap\Omega$, and let $F$ be the union over $k$ of the sets
$W_k\setminus\bigcup_{j<k}\operatorname{Sat}(W_j)$. It is measurable and lies in $\Omega$. Let $z,z'\in F$, with $z$ and $z'$ in
the sets of indices $k$ and $k'$ of this union, and let $c(z')=\gamma c(z)b$ with $\gamma\in\Gamma$ and $b\in\mathcal B^\Psi$. Then
$z'\sim z$. If $k<k'$, then $z'\in\operatorname{Sat}(W_k)$, and if $k'<k$, then $z\in\operatorname{Sat}(W_{k'})$; both are
excluded. So $k=k'$, and $z=z'$ by the choice of $O_k$. Let $z\in\Omega$. It lies in some $W_k$, so there is a least $k$ such
that $z\sim z'$ for some $z'\in W_k$. For $j<k$, $z'\notin\operatorname{Sat}(W_j)$, as otherwise $z$ would be equivalent to a
point of $W_j$. So $z'\in F$, and $\delta c(z)=c(z')b$ with $\delta\in\Gamma$ and $b\in\mathcal B^\Psi$ gives
$c(z)=\delta^{-1}c(z')b$.
% Lean: fundSet_holds = FundSetPf.fundSet_main: the relation R (R_refl, R_symm, R_trans), charB, small (PropDisc at
% cpt z0, Gamma_free, cont_ent, eventually_all_finite, cpt_uniq), countable_GL, meas_Sat (T as in the proof,
% MeasurableSet.image_of_continuousOn_injOn, cpt_uniq), TopologicalSpace.countable_cover_nhdsWithin for the cover, the
% set FF = union of W k minus the Sat (W j) for j < k, and Nat.find for the covering. The covering holds at every
% point of OmF x0, so muF-almost everywhere.
\end{proof}

\begin{lemma}\label{is-onto-step-orbit-interpolation-small-orbit-sum-unfold}
Let $X$ consist of the data (1) to (6) of flag data of rank $r$ over $E$, use the notation of
Definition~\ref{is-onto-share-counts-section-levi}, and let $x_0$ be an orthonormal point and $F$ a fundamental set for
$\Gamma$ at $x_0$ (Definition~\ref{is-onto-step-orbit-interpolation-small-orbit-sum-fundamental}).
\begin{enumerate}
\item Let $\phi\colon M_r(\mathbf C)^\Psi\to[0,\infty]$ be measurable with $\phi(yb)=\phi(y)$ for $y\in\mathcal D_X$ and
  $b\in\mathcal B^\Psi$. Then
  \[
   \int_{\Omega(x_0)}\phi(c_{x_0}(z))\,D_{x_0}(z)\,dz=\int_F\sum_{\gamma\in\Gamma}\phi(\gamma c_{x_0}(z))\,D_{x_0}(z)\,dz .
  \]
\item Let $x$ be an orthonormal point, $A\subset\Omega(x)$ a measurable set, and $\Phi\colon M_r(\mathbf C)^\Psi\to[0,\infty]$ a
  measurable function with $\Phi(\gamma y)=\Phi(y)=\Phi(yb)$ for $y\in\mathcal D_X$, $\gamma\in\Gamma$ and $b\in\mathcal B^\Psi$. Let
  $m\in[0,\infty]$ be at least the number of the $\gamma\in\Gamma$ with $\gamma c_x(z)=c_x(z')b$ for some $z,z'\in A$ and
  $b\in\mathcal B^\Psi$. Then
  \[
   \int_A\Phi(c_x(z))\,D_x(z)\,dz\le m\int_F\Phi(c_{x_0}(z))\,D_{x_0}(z)\,dz .
  \]
\end{enumerate}
\end{lemma}
\leanref{Development/AbelianVarietyTheory/Proofs/TorusHaarMeasure/MeasurableSetSatFrames.lean}{GlobalTheory.Unfold.lintegral_unfold}
\leanref{Development/AbelianVarietyTheory/Proofs/DivisorProofs/LintegralInvariantChartSetLeMulFundamentalSetStatement.lean}{GlobalTheory.T.lintegral_invariant_chartSet_le_mul_fundamentalSet}
\leanref{Development/AbelianVarietyTheory/Proofs/DivisorProofs/LintegralInvariantChartSetLeMulFundamentalSet.lean}{GlobalTheory.Unfold.lintegral_chart_le_mul_fundSet}
\leanref{Development/AbelianVarietyTheory/Proofs/TorusHaarMeasure/MeasurableSetSatFrames.lean}{GlobalTheory.Unfold.measurableSet_satFrames}
\begin{proof}
(1) Lemma~\ref{is-onto-share-counts-section-poincare-measure}(2) applies to $F$ and $\phi$, and a series of nonnegative
measurable functions can be integrated term by term.

(2) Let $S$ be the set of the points $c_x(z)b$ with $z\in A$ and $b\in\mathcal B^\Psi$. It is the image of the measurable set
$A\times\mathcal B^\Psi$ under the continuous map $(z,b)\mapsto c_x(z)b$, which is injective on it: $c_x(z)b=c_x(z')b'$ forces
$z=z'$ by Lemma~\ref{is-onto-share-counts-section-transfer-lu}(3) for $x$, and then $b=b'$, as $c_x(z)$ has invertible
components. So $S$ is measurable by the Lusin-Souslin theorem \cite[Theorem~15.1]{Kec95}, and $Sb=S$ for
$b\in\mathcal B^\Psi$. The function $\psi=\mathbf 1_S\Phi$ is measurable, with $\psi(yb)=\psi(y)$ for $y\in\mathcal D_X$ and
$b\in\mathcal B^\Psi$, and $c_x(z)\in S$ only for $z\in A$, by Lemma~\ref{is-onto-share-counts-section-transfer-lu}(3). So
Lemma~\ref{is-onto-share-counts-section-poincare-measure}(1) and (1) give
\[
 \int_A\Phi(c_x(z))D_x(z)\,dz=\int_{\Omega(x)}\psi(c_x(z))D_x(z)\,dz=\int_{\Omega(x_0)}\psi(c_{x_0}(z))D_{x_0}(z)\,dz
 =\int_F\sum_{\gamma\in\Gamma}\psi(\gamma c_{x_0}(z))D_{x_0}(z)\,dz .
\]
Let $z\in F$, put $y=c_{x_0}(z)\in\mathcal D_X$, and let $T_y$ be the set of the $\gamma\in\Gamma$ with $\gamma y\in S$. As
$\Phi(\gamma y)=\Phi(y)$, the sum $\sum_\gamma\psi(\gamma y)$ is $|T_y|\,\Phi(y)$. If $\gamma_0\in T_y$, the map $\gamma\mapsto\gamma\gamma_0^{-1}$
is injective on $T_y$, with values in the set counted by $m$: from $\gamma_0y=c_x(z_0)b_0$ and $\gamma y=c_x(z_1)b_1$ with
$z_0,z_1\in A$ follows $\gamma\gamma_0^{-1}c_x(z_0)=c_x(z_1)b_1b_0^{-1}$. So $|T_y|\le m$, and integrating over $F$ gives the
claim.
% Lean: (1) is Unfold.lintegral_unfold, proved again from measInv_holds by the argument of PairUnfoldPf.lunfold
% (Lemma is-onto-share-counts-section-poincare-measure(2)), with the sum inside the integral: the indicators of the
% translates of satFrames FF (measurableSet_satFrames), unique_translate, lintegral_tsum and muF_lact. (2) is
% T.lintegral_invariant_chartSet_le_mul_fundamentalSet = Unfold.lintegral_chart_le_mul_fundSet: psi is the indicator
% of satFrames x A times Phi (measurableSet_satFrames through MeasurableSet.image_of_continuousOn_injOn, cpt_uniq and
% rmul_right_cancel), cpt_mem_satFrames_iff, measInv_holds, lintegral_unfold, hpt and hTS for the count, and
% ENNReal.tsum_set_const.
\end{proof}

\begin{lemma}\label{is-onto-step-orbit-interpolation-small-orbit-sum-harnack}
Let $X$ consist of the data (1) to (6) of flag data of rank $r$ over $E$, use the notation of
Definitions~\ref{is-onto-share-counts-section-levi} and~\ref{is-onto-step-free-ball}, write $|z|$ for the Euclidean norm of
$z\in\mathbf C^n$, and let $x_0$ be an orthonormal point.
\begin{enumerate}
\item $\rho(x,y)=\rho(y,x)$ for all orthonormal points $x$ and $y$.
\item There are real numbers $\varrho>0$ and $K$ such that for all orthonormal points $x$ and $x_1$ and every $z\in\mathbf C^n$
  with $|z|<\varrho$: $z\in\Omega(x_1)$ and $\rho(x,c_{x_1}(z))\le\rho(x,x_1)+K$.
\end{enumerate}
\end{lemma}
\leanref{Development/AbelianVarietyTheory/Proofs/DivisorProofs/OrbitSumExpPotentialUniformBound.lean}{GlobalTheory.OrbitSum.rhX_symm}
\leanref{Development/AbelianVarietyTheory/Proofs/DivisorProofs/OrbitSumExpPotentialUniformBound.lean}{GlobalTheory.OrbitSum.norm_lminor_inv}
\leanref{Development/AbelianVarietyTheory/Proofs/DivisorProofs/OrbitSumExpPotentialUniformBound.lean}{GlobalTheory.OrbitSum.harnack}
\begin{proof}
Write $\ell_{p_0}(N)$ for the leading principal $p_0\times p_0$ minor of a square matrix $N$ and $M^*$ for the conjugate
transpose, and let $J$ be as in Lemma~\ref{is-onto-step-free-ball-minor}. For orthonormal points $x$ and $y$ and
$\tau\in\Psi$, $g=x_\tau^{-1}y_\tau$ satisfies $g^*Jg=J$ by Lemma~\ref{is-onto-step-free-ball-minor}(2), so $gJg^*=J$ and
$g^{-1}=Jg^*J$, as $J^2=1$; and $|d_{p_0}(y,\tau)|=1$. So by Definition~\ref{is-onto-step-free-ball} and
Lemma~\ref{is-onto-step-free-ball-minor}(3), $\rho(x,y)=\sum_\tau\log|\ell_{p_0}(x_\tau^{-1}y_\tau)|^2$.

(1) The leading $p_0\times p_0$ block of $Jg^*J$ is the conjugate transpose of that of $g$, so
$|\ell_{p_0}(y_\tau^{-1}x_\tau)|=|\ell_{p_0}(g^{-1})|=|\ell_{p_0}(g)|$, and the formula above gives $\rho(y,x)=\rho(x,y)$.

(2) Let $\eta$ be the number of Lemma~\ref{is-onto-step-orbit-interpolation-large-pole-exists-block} for $p=p_0$ and
$c=0$. The functions $z\mapsto d_{p_0}(c_{x_0}(z),\tau)$ are continuous, with absolute value $1$ at $z=0$, and $\Omega(x_0)$ is
an open set that contains $0$. So there is $\varrho$ with $0<\varrho\le\eta$ such that the ball $|z|<\varrho$ lies in $\Omega(x_0)$
and $|d_{p_0}(c_{x_0}(z),\tau)|\geq\frac12$ on it for every $\tau$. Put $K=|\Psi|\log8$. Let $x$ and $x_1$ be orthonormal and
$|z|<\varrho$. As in the proof of Lemma~\ref{is-onto-step-orbit-interpolation-large-point}, $\Omega(x_1)=\Omega(x_0)$ and
$d_{p_0}(c_{x_1}(z),\tau)=d_{p_0}(c_{x_0}(z),\tau)$; so $z\in\Omega(x_1)$ and $|d_{p_0}(c_{x_1}(z),\tau)|\geq\frac12$. Let $\tau\in\Psi$
and $g=x_\tau^{-1}x_{1,\tau}$, so that $gJg^*=J$. Then $x_\tau^{-1}c_{x_1}(z)_\tau=g\,\ell_\tau(z)$, and the entries of
$\ell_\tau(z)-1$ have absolute value at most $|z|<\eta$; so $|\ell_{p_0}(g\ell_\tau(z))|\le2|\ell_{p_0}(g)|$ by
Lemma~\ref{is-onto-step-orbit-interpolation-large-pole-exists-block}. As $|d_{p_0}(x_1,\tau)|=1$, the quotients of
Lemma~\ref{is-onto-step-free-ball-minor}(3) satisfy
\[
 q_\tau(x,c_{x_1}(z))=\frac{|d_{p_0}(c_{x_1}(z),\tau)|}{|\ell_{p_0}(g\ell_\tau(z))|^2}\geq\frac{1/2}{4|\ell_{p_0}(g)|^2}
 =\frac18\,q_\tau(x,x_1).
\]
Taking minus the logarithm of the product over $\tau$ gives $\rho(x,c_{x_1}(z))\le\rho(x,x_1)+|\Psi|\log8$.
% Lean: (1) is OrbitSum.rhX_symm, through norm_lminor_inv (hinv: g^-1 = J g^H J, hsub) and IC.norm_dmin_orth. (2) is
% OrbitSum.harnack: exists_block_det_le (Lemma is-onto-step-orbit-interpolation-large-pole-exists-block, its upper
% bound), rho1 from the continuity of the minors dmin at cpt x0 0 = x0 within 1/2, rho2 from IsOpen OmF x0, rho = the
% smallest of delta, rho1 and rho2, PoleValue.OmF_eq and dmin_cpt_eq for x1, hd, hab (norm_nmat_sub_one_le,
% isometry_of_orth, lminor_inv_mul_ne_zero), and hprod with K = |Psi| log 8.
\end{proof}

\begin{lemma}\label{is-onto-step-orbit-interpolation-small-orbit-sum}
Let $X$ consist of the data (1) to (6) of flag data of rank $r$ over $E$, such that $\Gamma$ acts properly
discontinuously on $\mathcal D_X$ modulo $\mathcal B^\Psi$ (Lemma~\ref{is-onto-share-counts-section-discontinuous}), use
the notation of Definitions~\ref{is-onto-share-counts-section-levi} and~\ref{is-onto-step-free-ball}, and let $x_0$ be
an orthonormal point and $s>0$ and $I$ real numbers such that $\int_{\Omega(x_1)}e^{-s\rho_{x_1}(z)}D_{x_1}(z)\,dz\le I$
for every orthonormal point $x_1$. Then there is a real number $C$ such that
$\sum_{\gamma\in\Gamma}e^{-s\rho(x_0,\gamma x)}\le C$ for every orthonormal point $x$.
\end{lemma}
\leanref{Development/AbelianVarietyTheory/Proofs/DivisorProofs/OrbitSumExpPotentialUniformBoundStatement.lean}{GlobalTheory.T.orbitSum_exp_potential_uniform_bound}
\leanref{Development/AbelianVarietyTheory/Proofs/DivisorProofs/OrbitSumExpPotentialUniformBound.lean}{GlobalTheory.OrbitSum.orbitSum_le}
\begin{proof}
Write $c=c_{x_0}$, $\Omega=\Omega(x_0)$ and $D=D_{x_0}$, and $|z|$ for the Euclidean norm of $z\in\mathbf C^n$.

\emph{The choices.} Let $\varrho_1$ and $K$ be the numbers of
Lemma~\ref{is-onto-step-orbit-interpolation-small-orbit-sum-harnack}(2) for $x_0$. By
Lemma~\ref{is-onto-share-counts-section-discontinuous} at $x_0$ and $x_0$, and as $c$ is continuous with $c(0)=x_0$, there is
$\varrho_2>0$ such that only finitely many $\gamma\in\Gamma$ satisfy $\gamma c(z)=c(z')b$ for some $z,z'$ with $|z|,|z'|<\varrho_2$
and some $b\in\mathcal B^\Psi$. Let $F$ be a fundamental set for $\Gamma$ at $x_0$
(Lemma~\ref{is-onto-step-orbit-interpolation-small-orbit-sum-fundamental-exists}). Put $\varrho=\min(\varrho_1,\varrho_2)$, and let
$A$ be the open ball $|z|<\varrho$; it lies in $\Omega$, by
Lemma~\ref{is-onto-step-orbit-interpolation-small-orbit-sum-harnack}(2) for $x=x_1=x_0$. Let $m$ be the number of the
$\gamma\in\Gamma$ with $\gamma c(z)=c(z')b$ for some $z,z'\in A$ and $b\in\mathcal B^\Psi$, which is finite. As $D>0$ on $\Omega$ and
$A$ is a nonempty open subset of $\Omega$, $v=\int_AD(z)\,dz>0$. Put $C=e^{sK}mI/v$.

Let $x$ be an orthonormal point. Put $\phi(y)=e^{-s\rho(x,y)}$ for $y\in\mathcal D_X$ and $\phi(y)=0$ otherwise, and
$G(y)=\sum_{\delta\in\Gamma}\phi(\delta y)$. The function $\rho(x,\cdot)$ is continuous on the open set $\mathcal D_X$, as minus the
logarithm of the product of the continuous positive functions $q_\tau(x,\cdot)$ of
Lemma~\ref{is-onto-step-free-ball-minor}(3); so $\phi$ and $G$ are measurable. By
Lemma~\ref{is-onto-step-free-ball-bound}(1), $\phi(yb)=\phi(y)$ and $G(yb)=G(y)$ for $y\in\mathcal D_X$ and
$b\in\mathcal B^\Psi$, and the substitution $\delta\mapsto\delta\gamma$ gives $G(\gamma y)=G(y)$ for $\gamma\in\Gamma$.

\emph{The sum is $G(x_0)$.} Let $\delta\in\Gamma$. The point $\delta x_0$ is orthonormal, and by
Lemmas~\ref{is-onto-step-orbit-interpolation-large-pole-exists-translate}(1)
and~\ref{is-onto-step-orbit-interpolation-small-orbit-sum-harnack}(1), $\rho(x_0,\delta^{-1}x)=\rho(\delta x_0,x)=\rho(x,\delta x_0)$.
So the substitution $\gamma=\delta^{-1}$ gives $\sum_{\gamma\in\Gamma}e^{-s\rho(x_0,\gamma x)}=\sum_\delta\phi(\delta x_0)=G(x_0)$.

\emph{Averaging.} Let $z\in A$ and $\delta\in\Gamma$.
Lemma~\ref{is-onto-step-orbit-interpolation-small-orbit-sum-harnack}(2), for $x$ and the orthonormal point $x_1=\delta x_0$,
gives $\rho(x,c_{\delta x_0}(z))\le\rho(x,\delta x_0)+K$, and $c_{\delta x_0}(z)=\delta c(z)$. As $s>0$,
$\phi(\delta x_0)\le e^{sK}\phi(\delta c(z))$, and summing over $\delta$, $G(x_0)\le e^{sK}G(c(z))$. Integrating over $A$ against
$D(z)\,dz$, and by Lemma~\ref{is-onto-step-orbit-interpolation-small-orbit-sum-unfold}(2) for $x_0$, $x_0$, $A$ and
$\Phi=G$,
\[
 G(x_0)\,v\le e^{sK}\int_AG(c(z))\,D(z)\,dz\le e^{sK}m\int_FG(c(z))\,D(z)\,dz .
\]

\emph{Unfolding.} By Lemma~\ref{is-onto-step-orbit-interpolation-small-orbit-sum-unfold}(1) for $\phi$,
Lemma~\ref{is-onto-share-counts-section-poincare-measure}(1) for $x_0$ and $x$, and
Lemma~\ref{is-onto-step-orbit-interpolation-large-block}(1) for $x$,
\[
 \int_FG(c(z))\,D(z)\,dz=\int_\Omega\phi(c(z))\,D(z)\,dz=\int_{\Omega(x)}\phi(c_x(z))\,D_x(z)\,dz
 =\int_{\Omega(x)}e^{-s\rho_x(z)}D_x(z)\,dz\le I.
\]
Hence $G(x_0)\le e^{sK}mI/v=C$.
% Lean: T.orbitSum_exp_potential_uniform_bound = OrbitSum.orbitSum_le: harnack (Lemma
% is-onto-step-orbit-interpolation-small-orbit-sum-harnack(2)), finite_mult (PropDisc at x0, x0 and the continuity of
% cpt x0 at 0), fundSet_holds (Lemma is-onto-step-orbit-interpolation-small-orbit-sum-fundamental-exists), A = ball 0 rho
% with hAO, m = encard S, v = muF x0 A (muF_ball_ne_zero), k = exp(s K), C = k m I / v; phi through
% Up.measurable_rhX and rhX_rmul, G with hGm, hGB and hGΓ; hsum (rhX_lact, rhX_symm, Lemmas
% is-onto-step-orbit-interpolation-large-pole-exists-translate(1) and -small-orbit-sum-harnack(1)), hHar
% (Unfold.cpt_lact), h1, h2 (Unfold.lintegral_chart_le_mul_fundSet), h3 (Unfold.lintegral_unfold), h4 (measInv_holds,
% Up.rhX_cpt). The Lean's phi is exp(-s rhX x y) at every y, and rhX is defined everywhere; the book sets phi = 0 off
% the flag domain, which changes no integral or sum used.
% Size: about 1,400 lines of Lean below orbitSum_exp_potential_uniform_bound (text count through namespaced names,
% stopping at statements the book cites). Nothing below this proof remains open.
\end{proof}

\begin{lemma}\label{is-onto-step-orbit-interpolation-small-weight}
Let $X$ consist of the data (1) to (6) of flag data of rank $r$ over $E$ with $0<p_0<r$ and $p_0(r-p_0)\le3$, such that
$\Gamma$ acts properly discontinuously on $\mathcal D_X$ modulo $\mathcal B^\Psi$
(Lemma~\ref{is-onto-share-counts-section-discontinuous}), use the notation of
Definitions~\ref{is-onto-share-counts-section-levi} and~\ref{is-onto-step-free-ball}, write $|z|$ for the Euclidean
norm of $z\in\mathbf C^n$, and let $x_0$ be an orthonormal point and $\nu>0$ a real number with $\nu\geq4(r-1)$.
\begin{enumerate}
\item There is a real number $C$ with $\sum_{\gamma\in\Gamma}e^{-(\nu/2)\rho(x_0,\gamma x)}\le C$ for every orthonormal point
  $x$.
\item There is a measurable function $\kappa\colon M_r(\mathbf C)^\Psi\to\mathbf R$ with $0<\kappa\le1$ on $\mathcal D_X$,
  $\kappa(xb)=\kappa(x)$ for $x\in\mathcal D_X$ and $b\in\mathcal B^\Psi$, and $\kappa(c_{x_0}(z))=e^{-\nu\rho_{x_0}(z)}$ for
  $z\in\Omega(x_0)$, such that there are real numbers $C_s$, $\varrho>0$ and $C'$ with
  $\sum_{\gamma\in\Gamma}\kappa(\gamma x)^{1/2}\le C_s$ for every $x\in\mathcal D_X$, and $\kappa(c_x(z))\le C'\kappa(x)$ for every
  orthonormal point $x$ and every $z$ with $|z|<\varrho$.
\end{enumerate}
\end{lemma}
\leanref{Development/AbelianVarietyTheory/Proofs/DivisorProofs/FreeBallCentreInterpolationDetWeightSmallSignatu.lean}{GlobalTheory.Up.OrbitInterpSmall.sum_bound}
\leanref{Development/AbelianVarietyTheory/Proofs/DivisorProofs/FreeBallCentreInterpolationDetWeightSmallSignatu.lean}{GlobalTheory.Up.OrbitInterpSmall.kappa_exists}
\begin{proof}
(1) As $\nu/2\geq2(r-1)$, Lemma~\ref{is-onto-step-orbit-interpolation-small-integral} for $s=\nu/2$ gives $I$, and
Lemma~\ref{is-onto-step-orbit-interpolation-small-orbit-sum} for $s=\nu/2>0$ and $I$ gives $C$.

(2) Let $q=q_{x_0}$ be the function of Lemma~\ref{is-onto-step-orbit-interpolation-large-weight-near}. As in the proof
of Lemma~\ref{is-onto-step-orbit-interpolation-large-weight}, $q$ is measurable, $0<q\le1$ on $\mathcal D_X$,
$q(xb)=q(x)$ for $x\in\mathcal D_X$ and $b\in\mathcal B^\Psi$, $q(c_{x_0}(z))=e^{-\rho_{x_0}(z)}$ for $z\in\Omega(x_0)$, and
the local bound of Lemma~\ref{is-onto-step-orbit-interpolation-large-weight-near} holds. For $x$ orthonormal and
$\gamma\in\Gamma$, $\gamma x\in\mathcal D_X$, so $q(\gamma x)=\prod_\tau q_\tau(x_0,\gamma x)=e^{-\rho(x_0,\gamma x)}$ by
Definition~\ref{is-onto-step-free-ball}, and $\sum_\gamma q(\gamma x)^{\nu/2}\le C$ by (1).
Lemma~\ref{is-onto-step-orbit-interpolation-large-weight-power} for $q$ and $\nu$ gives $\kappa=q^\nu$.
% Lean: (1) is sum_bound = T.orbitSum_exp_potential_uniform_bound with
% T.exists_uniform_bound_lintegral_flagDomain_ball_case (r >= 2). (2) is kappa_exists = kaFrame_of_delta with
% frame_potential_basic, minorRatio_near and sum_bound for s = bb / 2, through exp_neg_rhX.
\end{proof}

\begin{lemma}\label{is-onto-step-orbit-interpolation-small-poincare}
Let $X$ consist of the data (1) to (6) of flag data of rank $r$ over $E$, such that $\Gamma$ acts properly
discontinuously on $\mathcal D_X$ modulo $\mathcal B^\Psi$ (Lemma~\ref{is-onto-share-counts-section-discontinuous}), and
use the notation of Definitions~\ref{is-onto-share-counts-section-levi}, \ref{is-onto-share-counts-section-gauge}
and~\ref{is-onto-step-free-ball}. Let $x_0$ be an orthonormal point, $\lambda$ a weight, $W$ a gauge of weight $\lambda$
that is bounded below, $\nu>0$ a real number such that there is a real number $C$ with
$\sum_{\gamma\in\Gamma}e^{-(\nu/2)\rho(x_0,\gamma x)}\le C$ for every orthonormal point $x$, $f$ a holomorphic function on
$\Omega(x_0)$ that is square integrable for $W$ and $\nu$, and $P$ a function that is $\mathcal B$-equivariant of weight
$\lambda$ with $P\circ c_{x_0}=f$ on $\Omega(x_0)$. Then:
\begin{enumerate}
\item for every $x\in\mathcal D_X$ the family $\big(\chi_\lambda(\gamma)^{-1}P(\gamma x)\big)_{\gamma\in\Gamma}$ is absolutely
  summable;
\item $\mathcal P_\lambda P$ is a bounded section of weight $\lambda$.
\end{enumerate}
\end{lemma}
\leanref{Development/AbelianVarietyTheory/Proofs/DivisorProofs/PoinConvBdd.lean}{GlobalTheory.PoinSumBound.poin_conv_bdd}
\leanref{Development/AbelianVarietyTheory/Proofs/DivisorProofs/PoinConvBdd.lean}{GlobalTheory.PoinSumBound.sumBall}
\begin{proof}
Let $\rho_Y$ be the function of Lemma~\ref{is-onto-share-counts-section-poincare-count} for $x_0$. For $y\in\mathcal D_X$,
$\rho_Y(y)=\rho(x_0,y)$: both are $-\sum_\tau\big(\log|d_{p_0}(y,\tau)|-2\log|\ell_{p_0}(x_{0,\tau}^{-1}y_\tau)|\big)$, with nonzero
minors by Lemma~\ref{is-onto-step-free-ball-minor}(3). So $\sum_\gamma e^{-(\nu/2)\rho_Y(\gamma x)}\le C$ for every orthonormal
point $x$, as $\gamma x\in\mathcal D_X$.

The proof of Lemma~\ref{is-onto-share-counts-section-poincare-ball} uses $\nu>4|\Psi|r^2$ only at its last step, where
Lemma~\ref{is-onto-share-counts-section-poincare-count}(3) bounds $\sum_\gamma e^{-(\nu/2)\rho_Y(\gamma x)}$; its other steps
use $\nu\geq0$. With the bound above in place of that step, the conclusion of
Lemma~\ref{is-onto-share-counts-section-poincare-ball} holds for $W$, $\nu$, $f$ and $P$: there are $0<\varrho_0\le1$ and
$C_S$ with
\[
 \sum_{\gamma\in\Gamma}\Big(\int_{B_{\varrho_0}(x)}|P(\gamma c_x(w))|^2\,W(\gamma c_x(w))\,dw\Big)^{1/2}\le C_S
\]
for every orthonormal point $x$. The proofs of the local bound and of parts (1), (4) and (2) of
Lemma~\ref{is-onto-share-counts-section-poincare} use only this conclusion, Lemma~\ref{is-onto-share-counts-section-poincare-mean}
and the lower bound for $W$; so they apply here, and give (1) and (2).
% Lean: poin_conv_bdd: kaY_eq_exp (kaY x0 s = exp(-s rhX x0) on X.dom, through rhX_eq_sum) and kaSum_of; sumBall is
% the proof of PoinPf.sumBall_holds (Lemma is-onto-share-counts-section-poincare-ball) with kaSum_of in place of
% PoinPf.kaSum; ballB and locD are ball_of and loc_of of the proof of poin_holds (Lemma
% is-onto-share-counts-section-poincare), with subMean_holds and subMeanLoc_of; and the assembly
% (differentiableOn_tsum_loc, poin_rmul, poin_lact, FrameDegreeData.pnorm_rmul, pnorm_orth') is that of
% PoinPf.poin_of.
\end{proof}

\begin{lemma}\label{is-onto-step-orbit-interpolation-small-interpolation}
Let $X$ consist of the data (1) to (6) of flag data of rank $r$ over $E$ with $0<p_0<r$ and $p_0(r-p_0)\le3$, such that
$\Gamma$ acts properly discontinuously on $\mathcal D_X$ modulo $\mathcal B^\Psi$
(Lemma~\ref{is-onto-share-counts-section-discontinuous}), use the notation of
Definitions~\ref{is-onto-share-counts-section-levi} and~\ref{is-onto-step-free-ball}, and suppose that
$\dim X=|\Psi|p_0(r-p_0)\geq1$. Let $G\subset GL_r(E)$ be a subgroup, $x_0$ an orthonormal point, $\lambda$ a weight with
$\lambda_{\tau,i}=\lambda_{\tau,j}$ whenever $i,j\le p_0$ or $i,j>p_0$, and $\rho_0>0$, $\nu$ and $m$ real numbers, such that
$G$ has a free ball of radius $\rho_0$ at $x_0$, $\nu>0$, $\nu\geq4(r-1)$, $\nu+\dim X/\rho_0<m$,
$m\,\mathcal L_{\rho_{x_0}}(z;\xi)\le\mathcal L^{x_0}_\lambda(z;\xi)-\mathcal L_{\log D_{x_0}}(z;\xi)$ for every
$z\in\Omega(x_0)$ and $\xi\in\mathbf C^n$, and every bounded section of weight $\lambda$ is a cusp form. Then there is a
cusp form $P\in S_\lambda(X)$ with $P(x_0)\ne0$ and $P(\gamma x_0)=0$ for every $\gamma\in G$ outside $\Gamma$.
\end{lemma}
\leanref{Development/AbelianVarietyTheory/Proofs/DivisorProofs/FreeBallCentreInterpolationDetWeightSmallSignatu.lean}{GlobalTheory.Up.OrbitInterpSmall.orbitInterp_small}
\begin{proof}
Let $C$, $\kappa$ and $C_s$ be as in Lemma~\ref{is-onto-step-orbit-interpolation-small-weight}(1) and (2) for $x_0$ and
$\nu$, and let $W_\lambda$ be the gauge of Lemma~\ref{is-onto-share-counts-section-kronecker-gauge}, which is bounded
below. The set $S_\lambda(X)$ is a complex vector space (Lemma~\ref{is-onto-share-sections}(1)). As $\Gamma$ has finite
index in the group of the elements of $GL_r(E)$ that preserve $h$ and $\Lambda$, there is a finite set $T$ of such elements
such that each of them is $s\gamma$ for some $s\in T$ and $\gamma\in\Gamma$. Let $T'$ be the set of the $s\in T$ with
$s^{-1}\in G$ and $s^{-1}\notin\Gamma$. The elements of $G$ preserve $h$ and $\Lambda$, so the points $s^{-1}x_0$ lie in
$\mathcal D_X$.

\emph{Approximation.} Let $\eta>0$ and $\eta'=\eta/(C_s+1)$. Lemma~\ref{is-onto-step-orbit-interpolation-large-approx},
for $\kappa$ and $\eta'$, gives $F$ and $P$; it applies as $\nu\geq0$. By
Lemma~\ref{is-onto-step-orbit-interpolation-small-poincare}, for $W_\lambda$, $\nu$, $C$, $F$ and $P$, the families
$(\chi_\lambda(\gamma)^{-1}P(\gamma x))_{\gamma\in\Gamma}$ are absolutely summable for $x\in\mathcal D_X$, and $Q=\mathcal P_\lambda P$
is a bounded section of weight $\lambda$, hence a cusp form. As $\chi_\lambda$ has finite order on $\Gamma$,
$|\chi_\lambda(\gamma)|=1$, and $\chi_\lambda(1)=1$. Every $\gamma\ne1$ in $\Gamma$ lies in $G$, so
\[
 |Q(x_0)-1|\le|P(x_0)-1|+\sum_{\gamma\ne1}|P(\gamma x_0)|\le\eta'+\eta'\sum_\gamma\kappa(\gamma x_0)^{1/2}\le\eta'(1+C_s)=\eta .
\]
Let $s\in T'$. For $\gamma\in\Gamma$, $\gamma s^{-1}\in G$, and $\gamma s^{-1}\ne1$, as $s^{-1}\notin\Gamma$. So
$|Q(s^{-1}x_0)|\le\sum_\gamma|P(\gamma s^{-1}x_0)|\le\eta'\sum_\gamma\kappa(\gamma s^{-1}x_0)^{1/2}\le\eta'C_s\le\eta$.

\emph{Interpolation.} Lemma~\ref{is-onto-step-orbit-interpolation-large-span}, for $S_\lambda(X)$, the point $x_0$ and
the points $s^{-1}x_0$ with $s\in T'$, gives $u\in S_\lambda(X)$ with $u(x_0)=1$ and $u(s^{-1}x_0)=0$ for $s\in T'$. Let
$\gamma\in G$ outside $\Gamma$. As $\gamma^{-1}$ preserves $h$ and $\Lambda$, $\gamma^{-1}=s\gamma_1$ with $s\in T$ and
$\gamma_1=s^{-1}\gamma^{-1}\in\Gamma$. Then $s^{-1}=\gamma_1\gamma\in G$, and $s^{-1}\notin\Gamma$, since otherwise
$\gamma=\gamma_1^{-1}s^{-1}\in\Gamma$; so $s\in T'$. As $u$ is a section of weight $\lambda$,
$u(\gamma x_0)=u(\gamma_1^{-1}s^{-1}x_0)=\chi_\lambda(\gamma_1^{-1})u(s^{-1}x_0)=0$. So $P=u$ works.
% Lean: orbitInterp_small: kappa_exists (Lemma is-onto-step-orbit-interpolation-small-weight(2)), hsumE = sum_bound
% for bb / 2 (part (1)), KronMetric.kronMetric with psi = 0, siegelSpaceAlgebra_holds, X.Gamma_index,
% exists_of_approx, exists_approx (with the proved statements extFStatement_holds, measInv_holds, poleEx_holds,
% poleSolve measInv_holds and polePoint_holds) for eta / (Cs + 1), PoinSumBound.poin_conv_bdd (Lemma
% is-onto-step-orbit-interpolation-small-poincare) in place of PoinStatement, tsum_eq_add_tsum_ite, chi_one,
% OntoE.chi_norm, tsum_le_of_le, and the section law of u. It repeats the proof of orbitInterp (Lemma
% is-onto-step-orbit-interpolation-large-interpolation) line by line, and so does the book.
\end{proof}

\begin{lemma}\label{is-onto-step-orbit-interpolation-small}
In Situation~\ref{is-onto-step-orbit-situation}, suppose that $p_0(r-p_0)\le3$. Then there is a cusp form
$P\in S_{\delta_s}(X)$ with $P(x_0)\ne0$ and $P(\gamma x_0)=0$ for every $\gamma\in G$ outside $\Gamma$.
\end{lemma}
\leanref{Development/AbelianVarietyTheory/Proofs/DivisorProofs/FreeBallCentreInterpolationDetWeightSmallSignatuStatement.lean}{GlobalTheory.T.freeBallCentre_interpolation_detWeight_smallSignature}
\leanref{Development/AbelianVarietyTheory/Proofs/DivisorProofs/FreeBallCentreInterpolationDetWeightSmallSignatu.lean}{GlobalTheory.Up.OrbitInterpSmall.orbit_interp_small}
\begin{proof}
Put $d=\dim X=|\Psi|p_0(r-p_0)$. As in the proof of Lemma~\ref{is-onto-step-orbit-interpolation-large}, $\Psi$ is not
empty, so $d\geq1$; and $r\geq2$. Put $\rho_0=1/(49r^2)$, $\nu=4(r-1)+\frac12$ and $m=s-2r+2$. Then $\nu>0$,
$\nu\geq4(r-1)$, $m\geq50r^2d+3$ as $s\geq2r+50r^2d+1$, and $r^2d\geq r^2\geq4r-4$; so
\[
 \nu+\frac d{\rho_0}=4r-\frac72+49r^2d<50r^2d+3\le m .
\]
The hypotheses of Lemma~\ref{is-onto-step-orbit-interpolation-small-interpolation} hold for $G$, $x_0$,
$\lambda=\delta_s$, $\rho_0$, $\nu$ and $m$. The group $\Gamma$ acts properly discontinuously on $\mathcal D_X$ modulo
$\mathcal B^\Psi$, by Lemma~\ref{is-onto-share-counts-section-discontinuous}. The weight $\delta_s$ is constant on $i\le p_0$
and on $i>p_0$. The group $G$ has a free ball of radius $\rho_0$ at $x_0$, and $\nu>0$, $\nu\geq4(r-1)$. The curvature
inequality is Lemma~\ref{is-onto-step-orbit-interpolation-large-curvature}. Every bounded section of weight $\delta_s$ is a
cusp form, by Lemma~\ref{is-onto-share-counts-section-cusp} for any $\tau_0\in\Psi$, as in the proof of
Lemma~\ref{is-onto-step-orbit-interpolation-large}. So Lemma~\ref{is-onto-step-orbit-interpolation-small-interpolation}
gives $P$.
% Lean: T.freeBallCentre_interpolation_detWeight_smallSignature = Up.OrbitInterpSmall.orbit_interp_small,
% which applies orbitInterp_small (Lemma is-onto-step-orbit-interpolation-small-interpolation) with
% bb = 4 (r - 1) + 1/2, m = s - 2 r + 2 and rho0 = 1 / (49 r^2), one_le_dimX_iff, propDisc_closed, the block
% constancy of deMul, DivisionVolume.margin_iff with gapsPairing_deMul' (the same computation as gapsPairing_deMul,
% Lemma is-onto-step-orbit-interpolation-large-curvature) and cuspStatement_holds at tau in Psi.
% Size: nothing below this proof remains open; the last open steps were Lemmas
% is-onto-step-orbit-interpolation-large-pole-exists (about 3,400 lines of Lean),
% is-onto-step-orbit-interpolation-small-integral (about 1,750) and is-onto-step-orbit-interpolation-small-orbit-sum
% (about 1,400). The whole of orbit_interp_small has about 9,700 lines in 710 declarations below it (text count as
% for Lemma is-onto-step-orbit-interpolation-large); an earlier, looser count gave 29,000 lines.
\end{proof}

\begin{lemma}\label{is-onto-step-orbit-interpolation}
In Situation~\ref{is-onto-step-orbit-situation}, there is a cusp form $P\in S_{\delta_s}(X)$ with
$P(x_0)\ne0$ and $P(\gamma x_0)=0$ for every $\gamma\in G$ outside $\Gamma$.
\end{lemma}
\leanref{Development/AbelianVarietyTheory/Proofs/DivisorProofs/OrbitInterpolationDetWeightFreeBallIndefiniteStatement.lean}{GlobalTheory.T.orbitInterpolation_detWeight_freeBall_indefinite}
\leanref{Development/AbelianVarietyTheory/Proofs/DivisorProofs/OrbitInterpolationDetWeightFreeBallIndefinite.lean}{GlobalTheory.OrbitInterpWeight.orbit_interp_deMul}
\begin{proof}
If $p_0(r-p_0)\geq4$, this is Lemma~\ref{is-onto-step-orbit-interpolation-large}. Otherwise
$p_0(r-p_0)\le3$, and it is Lemma~\ref{is-onto-step-orbit-interpolation-small}.
% Lean: T.orbitInterpolation_detWeight_freeBall_indefinite = OrbitInterpWeight.orbit_interp_deMul, by
% cases on 4 <= p0 (r - p0).
\end{proof}

\begin{lemma}\label{is-onto-share-interpolation}
Let $M$ be two-level data of rank $r$ over $E$ for a finite group $U$ acting on a finite set $C$
(Definition~\ref{is-onto-share-two-level}), for which the level hypothesis holds, and let
$s\geq2r+50r^2\dim X(M)+1$ be even. Let $c\in C$, let $\Gamma^\circ_c$ be the set of products $\gamma t(v,c)$
with $vc=c$ and $\gamma\in\Gamma_c$, and let $x\in\mathcal D_M$ be a point at which there is no contraction
for $X_c$. Then there is a cusp form $P\in S_{\delta_s}(X_c)$ with $P(x)\ne0$ and $P(\beta x)=0$ for every
$\beta\in\Gamma^\circ_c$ outside $\Gamma_c$.
\end{lemma}
\leanref{Development/AbelianVarietyTheory/Proofs/DivisorProofs/ExistsCuspFormNonzeroAtOneLift.lean}{GlobalTheory.LiftInterp.comp_interp}
\leanref{Development/AbelianVarietyTheory/Proofs/DivisorProofs/ExistsCuspFormNonzeroAtOneLift.lean}{GlobalTheory.LiftInterp.exists_form_one_lift}
\leanref{Development/AbelianVarietyTheory/Proofs/DivisorProofs/UpperLevelTrivialDefiniteNoContraction.lean}{GlobalTheory.DefiniteLevel.eq_one_of_definite}
\leanref{Development/AbelianVarietyTheory/Proofs/DivisorProofs/ExistsDetPowerSectionNeZeroDefinite.lean}{GlobalTheory.DefConst.exists_ne_zero}
\leanref{Development/AbelianVarietyTheory/Proofs/DivisorProofs/ExistsDetPowerSectionNeZeroDefinite.lean}{GlobalTheory.DefConst.pnorm_detPow}
\leanref{Development/AbelianVarietyTheory/Proofs/DivisorProofs/ExistsDetPowerSectionNeZeroDefinite.lean}{GlobalTheory.DefConst.not_iso}
\begin{proof}
Write $X=X_c$, with $\Gamma=\Gamma_c$, $\Lambda=\Lambda_c$ and $L'=L'_c$, and let $g$ be the prime of the level
hypothesis. As in the proof of Lemma~\ref{is-onto-share-free}, $\Gamma^\circ_c$ is a subgroup of $GL_r(E)$, its
elements preserve $h$ and $L'$ and satisfy $(\beta-1)^rL'\subset PL'$ for every prime $P$ of $\mathcal O_E$
over $g$. It contains $\Gamma$, since $t(1,c)\in\Gamma$ by the cocycle condition, and its elements preserve
$\Lambda$, as $t(v,c)\Lambda_c=\Lambda_{vc}$ and $\Gamma$ preserves $\Lambda$.

\emph{The definite case $p_0\in\{0,r\}$.} Let $\beta\in\Gamma^\circ_c$. By
Lemma~\ref{is-onto-share-majorant}(5) and (1), $m_{x,\tau}(\beta^\tau u)=m_{x,\tau}(u)\le e^{2/(3r)}m_{x,\tau}(u)$
for all $\tau$ and $u$, so $\beta=1$ by Lemma~\ref{is-onto-share-small}. Hence $\Gamma^\circ_c=\Gamma=\{1\}$, and
the vanishing condition is empty. Let $a=\lfloor s/2\rfloor$ if $p_0=0$ and $a=-\lfloor s/2\rfloor$ if
$p_0=r$, so that $(\delta_s)_{\tau,i}=a$ for all $\tau$ and $i$, and put $P(y)=\prod_{\tau\in\Psi}\det(y_\tau)^a$.
It is holomorphic on $\mathcal D_X$, which is the set of $y$ with all $\det y_\tau\ne0$
(Lemma~\ref{is-onto-share-majorant}(5)), and $P(yb)=\prod_{\tau,i}(b_\tau)_{ii}^aP(y)$; the condition on $\Gamma$
holds as $\Gamma=\{1\}$. In the product defining $\|P\|_{\delta_s}(y)$ the ratios of leading minors telescope
to $d_r(y,\tau)=|\det y_\tau|^2\tau(\det A)$, so $\|P\|_{\delta_s}(y)=\prod_\tau|\tau(\det A)|^{-a/2}$ is
constant, and $P$ is bounded. There is no isotropic vector: by (2) of
Definition~\ref{is-onto-share-flag-data}, $h_\tau$ has Gram matrix $\pm1$ in some basis, so it is definite. So
$P$ is a cusp form, and $P(x)\ne0$.

\emph{The indefinite case $0<p_0<r$.} Here $r\geq2$, and $\Gamma^\circ_c$ and the prime $g$ satisfy the
hypotheses of Lemma~\ref{is-onto-step-free-ball-exists}. That lemma gives an orthonormal point $x_0$ and
$b\in\mathcal B^\Psi$ with $x=x_0b$ such that $\Gamma^\circ_c$ has a free ball of radius $1/(49r^2)$ at $x_0$
(Definition~\ref{is-onto-step-free-ball}). The characters $\chi_{c,\kappa}$ are trivial on $\Gamma$, and
$|\Psi|p_0(r-p_0)=\dim X(M)$. So Lemma~\ref{is-onto-step-orbit-interpolation}, with $G=\Gamma^\circ_c$, gives
a cusp form $P\in S_{\delta_s}(X)$ with $P(x_0)\ne0$ and $P(\beta x_0)=0$ for every $\beta\in\Gamma^\circ_c$
outside $\Gamma$.
% Lean: T.freeBall_noContraction_unitaryUpperLevel (Lemma is-onto-step-free-ball-exists) and
% T.orbitInterpolation_detWeight_freeBall_indefinite (Lemma is-onto-step-orbit-interpolation), with
% M.chi_one c for the characters.
Then $P$ also works at $x$. Indeed $P(x)=P(x_0b)=\prod_{\tau,i}(b_\tau)_{ii}^{(\delta_s)_{\tau,i}}P(x_0)\ne0$.
For $\beta\in\Gamma^\circ_c$ outside $\Gamma$, $\beta x_0\in\mathcal D_X$ as $\beta$ preserves $h$, and
$P(\beta x)=P((\beta x_0)b)=\prod_{\tau,i}(b_\tau)_{ii}^{(\delta_s)_{\tau,i}}P(\beta x_0)=0$.
\end{proof}

\begin{lemma}\label{is-onto-share-translates}
Let $M$ be two-level data of rank $r$ over $E$ for a finite group $U$ acting on a finite set $C$, for which
the level hypothesis holds, and let $s\geq2r+50r^2\dim X(M)+1$ be even. Then there is a family $u_0$ of
weight $\delta_s$ such that, for every $c\in C$, the determinant of the matrix
$\big(((vv')^*u_0)_c(x)\big)_{v,v'\in U}$ is nonzero at some $x\in\mathcal D_M$.
\end{lemma}
\leanref{Development/AbelianVarietyTheory/Proofs/DivisorProofs/ExistsCuspFormIndependentTranslatesStatement.lean}{GlobalTheory.T.exists_cuspForm_independent_translates}
\leanref{Development/AbelianVarietyTheory/BlockTheory/PoincareCuspForms/TwoLevelLogarithm.lean}{GlobalTheory.IndefiniteRowSideFullHecke.FrameWeightGapBounds.IndependentTranslates}
\leanref{Development/AbelianVarietyTheory/Proofs/DivisorProofs/ExistsCuspFormIndependentTranslates.lean}{GlobalTheory.TranslatesDet.independentTranslates}
\leanref{Development/AbelianVarietyTheory/Proofs/DivisorProofs/ExistsCuspFormIndependentTranslates.lean}{GlobalTheory.TranslatesDet.detNeAt}
\begin{proof}
Two facts on two-level data are used. First, $t(1,c)\in\Gamma_c$, by the cocycle condition for $u=v=1$.
Second, every component $F_c$ of a family satisfies $F_c(\gamma x)=F_c(x)$ for $\gamma\in\Gamma_c$, since
$\chi_{c,\kappa}$ is trivial on $\Gamma_c$.

For every $c\in C$, Lemma~\ref{is-onto-share-balanced}(1) gives a point $x_c\in\mathcal D_M$ at which there
is no contraction for $X_c$, and Lemma~\ref{is-onto-share-interpolation} a cusp form $P^{(c)}$ of $X_c$ of
weight $\delta_s$. Let $u^{(c)}$ be the family with $u^{(c)}_c=P^{(c)}$ and $u^{(c)}_{c'}=0$ for $c'\ne c$.
Then $u^{(c)}_c(x_c)\ne0$, and $(h^*u^{(c)})_c(x_c)=u^{(c)}_{hc}(t(h,c)x_c)=0$ for $h\ne1$: this is clear if
$hc\ne c$, and if $hc=c$, then $t(h,c)\in\Gamma^\circ_c$ and $t(h,c)\notin\Gamma_c$, by the last condition of
Definition~\ref{is-onto-share-two-level}.

Choose a point $c_q$ in every orbit $q$ of $U$ on $C$, and put $u_0=\sum_qu^{(c_q)}$, a family of weight
$\delta_s$. As $u^{(c_q)}$ vanishes off $c_q$, $u_0$ agrees with $u^{(c_q)}$ on the orbit $q$. Let $c\in C$,
write $c=kc_0$ with $c_0=c_q$ the chosen point of its orbit and $k\in U$, and put $u=u^{(c_0)}$,
$x_0=x_{c_0}$ and $x=t(k,c_0)x_0\in\mathcal D_M$. For $g\in U$, the cocycle condition gives
$t(gk,c_0)=\gamma\,t(g,kc_0)t(k,c_0)$ with $\gamma\in\Gamma_{gkc_0}$, so
\[
 (g^*u_0)_{kc_0}(x)=u_{gkc_0}\big(t(g,kc_0)t(k,c_0)x_0\big)=u_{gkc_0}\big(t(gk,c_0)x_0\big)=((gk)^*u)_{c_0}(x_0).
\]
So the entry of the matrix at $(h,h')$ is $((hh'k)^*u)_{c_0}(x_0)$. It vanishes unless $h'=h^{-1}k^{-1}$, and
then it is $(1^*u)_{c_0}(x_0)=u_{c_0}(t(1,c_0)x_0)=u_{c_0}(x_0)\ne0$. The matrix is therefore a permutation
matrix times $u_{c_0}(x_0)$, and its determinant is $\pm u_{c_0}(x_0)^{|U|}\ne0$.
\end{proof}

\begin{lemma}\label{is-onto-share-indefinite}
Let $M$ be two-level data of rank $r$ over $E$ for a finite group $U$ acting on a finite set $C$
(Definition~\ref{is-onto-share-two-level}), $\kappa$ a weight, $\rho$ a representation of $U$ on
$S_\kappa(M)$ pulled back from the action on $C$, and $s=s(M)$. Let $\varepsilon,c_0$ be real numbers such
that the cusp forms of $M$ have ranks $(\varepsilon,c_0)$ at $\kappa$ and $s$. Let $H\subset U$ be a
subgroup for which the level hypothesis holds, and $\lambda\colon H\to\mathbf C^\times$ a character. Then
\[
 \dim S_{\kappa-\delta_s}(M)\le|H|\dim S_\kappa(M)^{H,\lambda}\quad\text{and}\quad
 (1-\varepsilon)\dim S_\kappa(M)\le|H|\dim S_\kappa(M)^{H,\lambda}.
\]
\end{lemma}
\leanref{Development/AbelianVarietyTheory/Proofs/DivisorProofs/TwoLevelFlagDomainEigenspaceShareStatement.lean}{GlobalTheory.T.twoLevelFlagDomain_eigenspace_share}
\leanref{Development/AbelianVarietyTheory/Proofs/DivisorProofs/TwoLevelFlagDomainEigenspaceShare.lean}{GlobalTheory.ShareEig.share_eig}
\leanref{Development/AbelianVarietyTheory/Proofs/DivisorProofs/TwoLevelFlagDomainEigenspaceShare.lean}{GlobalTheory.ShareEig.eq_zero_of_det}
\begin{proof}
Write $S_\kappa=S_\kappa(M)$ and $S'=S_{\kappa-\delta_s}(M)$. Since $s$ is even and
$s\geq2r+50r^2\dim X(M)+1$, Lemma~\ref{is-onto-share-translates} for $M|_H$ gives a family $u_0$ of weight
$\delta_s$ whose matrix $\big(((hh')^*u_0)_c(x)\big)_{h,h'\in H}$ has a determinant $D_c(x)$ that is nonzero
somewhere on $\mathcal D_M$, for every $c$. By Lemma~\ref{is-onto-share-pullback}, every $(h^*u_0)_c$ is a
bounded section of weight $\delta_s$, so by Lemma~\ref{is-onto-share-sections}(1) the product of a family of
weight $\kappa-\delta_s$ with $h^*u_0$ is a family of weight $\kappa$. Every element of $S'$ is $(c,x)\mapsto
X_c(x)$ for a family $X$, by Lemma~\ref{is-onto-share-sections}(1); so multiplication by $h^*u_0$ is a linear
map $\mu_h\colon S'\to S_\kappa$.

The operator $\Sigma=\sum_{g\in H}\lambda(g)\rho(g^{-1})$ maps $S_\kappa$ into $S_\kappa^{H,\lambda}$:
$\rho(k)\Sigma=\sum_g\lambda(g)\rho(kg^{-1})=\lambda(k)\Sigma$, by the substitution $g\mapsto gk$. Let
$\Phi\colon S'\to(S_\kappa^{H,\lambda})^H$ be $\Phi(f)=(\Sigma\mu_hf)_{h\in H}$. We show that $\Phi$ is injective;
then $\dim S'\le|H|\dim S_\kappa^{H,\lambda}$, and the ranks give $(1-\varepsilon)\dim S_\kappa\le\dim S'$.

Let $\Phi(f)=0$ with $f$ given by a family $X$, and fix $c$ and $z\in\mathcal D_M$. With $a(g,c)$ the
constants of $\rho$, and since $g^*(h^*u_0)=(hg)^*u_0$ on $\mathcal D_M$ by the cocycle condition and the
triviality of the characters $\chi_{c',\kappa}$ on the groups $\Gamma_{c'}$,
\[
 0=(\Sigma\mu_hf)(c,z)=\sum_{g\in H}((hg)^*u_0)_c(z)\cdot\lambda(g)a(g,c)(g^*X)_c(z)\qquad(h\in H).
\]
So if $D_c(z)\ne0$, then $\lambda(g)a(g,c)(g^*X)_c(z)=0$ for all $g$; at $g=1$ this is $X_c(z)=0$, as
$1^*X=X$ on $\mathcal D_M$. The function $D_c$ is holomorphic on $\mathcal D_M$, a polynomial in the
holomorphic functions $((hh')^*u_0)_c$, and is nonzero somewhere, while $X_c$ is continuous. If
$X_c(z)\ne0$, then $X_c\ne0$ near $z$, so $D_c$ vanishes near $z$; as $\mathcal D_M$ is open and connected
(Lemma~\ref{is-onto-share-sections}(2)), $D_c$ vanishes identically by the identity theorem, a
contradiction. So $X_c=0$ on $\mathcal D_M$ for every $c$, and $f=0$.
\end{proof}

\subsection{The flat model and the upper count}\label{subsection-is-upper-count}

\begin{definition}\label{is-onto-share-counts-model}
Let $X$ be flag data of rank $r$ over $E$ (Definition~\ref{is-onto-share-flag-data}), and put
$n=|\Psi|r(r-1)/2$. Index the coordinates of $z\in\mathbf C^n$ by the triples $(\tau,a,b)$ with $\tau\in\Psi$
and $1\le b<a\le r$. For $y\in M_r(\mathbf C)^\Psi$ the \emph{chart} at $y$ is $z\mapsto c_y(z)$, where
$c_y(z)_\tau=y_\tau\ell_\tau(z)$ and $\ell_\tau(z)$ is the lower unitriangular matrix with the entries
$z_{\tau,a,b}$ below the diagonal. For a weight $\kappa$ and $x\in\mathcal D_X$ put
$N_\kappa(x)=\prod_{\tau,i}|d_i(x,\tau)/d_{i-1}(x,\tau)|^{-\kappa_{\tau,i}/2}$, so that $\|P\|_\kappa=|P|N_\kappa$. The
\emph{Levi form} at $z$ of a real function $g$ of class $C^2$ near $z\in\mathbf C^n$ is
$\xi\mapsto\sum_{j,k}\partial_j\bar\partial_kg(z)\,\xi_j\bar\xi_k$ on $\mathbf C^n$. For $x\in\mathcal D_X$ let
$\lambda_1(x)\geq0$ be the square root of the infimum of $\sum_\tau c_\tau m_{x,\tau}(\tau(v))$ over the nonzero
$v\in\Lambda$.

Let $L,F$ be weights and $T$ an integer. An archimedean model of $X$ for $(L,F,T)$
(Definition~\ref{is-onto-share-two-level}) consists of \emph{archimedean data} for $(L,F,T)$, \emph{margins}
for these archimedean data and a \emph{flat model} for them. These notions use only the data (1) to (6) of
$X$. Among their data and conditions are the following: (1) to (4) belong to the archimedean data, and (5)
to the flat model.
\begin{enumerate}
\item Integers $i,g\geq1$ with $g$ even, a real number $c>0$ and a weight $\Delta$, with $n\geq2$,
  $11n^2\le T$ and $Tcg\le1$.
\item Cusp forms $w_1,\dots,w_N$ of weight $L$ without common zero in $\mathcal D_X$, and a function
  $\phi\geq0$ on $\mathcal D_X$ with $\phi(\gamma xb)=\phi(x)$ for $\gamma\in\Gamma$ and $b\in\mathcal B^\Psi$, such
  that for some $c_a,\varepsilon_0,K>0$: $\phi(x)\geq c_a\lambda_1(x)^{-2}$ whenever $\lambda_1(x)<\varepsilon_0$,
  and $e^{-2c\phi}\le K\sum_k\|w_k\|_L^2$ on $\mathcal D_X$.
\item A section $\eta$ of weight $i(g/2)\Delta$ with $\|\eta\|^2_{i(g/2)\Delta}\le Ke^{-2i\phi}$ on $\mathcal D_X$,
  and cusp forms $\theta_1$ of weight $2L-\Delta$, $\theta_2$ of weight $(n+4)L-F$ and $\theta_0$ of weight
  $F$. None of $\eta,\theta_0,\theta_1,\theta_2$ vanishes identically on $\mathcal D_X$.
\item For all $y\in\mathcal D_X$ and all $z$ with $c_y(z)\in\mathcal D_X$, the Levi form at $z$ of
  $z'\mapsto-2\log N_L(c_y(z'))$ is positive definite.
\item The \emph{flat model}: cusp forms $w_0,s_1,\dots,s_n$ of weight $L$ and \emph{base points}
  $y_1,\dots,y_\delta\in\mathcal D_X$ such that the common zeros of the $s_k$ in $\mathcal D_X$ are the points
  $\gamma y_jb$ with $\gamma\in\Gamma$ and $b\in\mathcal B^\Psi$; $y_{j'}=\gamma y_jb$ only if $j'=j$;
  $w_0(y_j)\ne0$; the derivative at $0$ of $z\mapsto\big(s_k(c_{y_j}(z))/w_0(c_{y_j}(z))\big)_k$ is invertible;
  and for every open set $U$ that contains the $y_j$, and the points $\gamma xb$ with each of its points $x$,
  there is $K'$ with $\sum_k\|w_k\|_L^2\le K'\sum_k\|s_k\|_L^2$ on $\mathcal D_X\setminus U$.
\end{enumerate}
The other conditions bound the Levi forms of further functions; this book does not unfold them. Put
$S=\sum_k\|s_k\|_L^2$ on $\mathcal D_X$. For an integer $T'\geq0$ and a function $f$ on $\mathcal D_X$, the
\emph{quotient} $Q_{T',f}$ is $0$ where $S=0$, and elsewhere
\[
 Q_{T',f}=\|f\|^2_{F+T'L}\,\|\theta_2\|^2_{(n+4)L-F}\,\|\theta_1\|^2_{2L-\Delta}\,
 \big(\|\eta\|^2_{i(g/2)\Delta}\big)^{2/(ig)}\,S^{-(T'+n+13/2)}.
\]
\end{definition}
\leanref{Development/AbelianVarietyTheory/Definitions/Divisors/ExistsArchOne.lean}{GlobalTheory.ShareOn.ArchOne}
\leanref{ReciprocityTheory/ReductiveGroups/ArithmeticGeometry/ArakelovVolumes/UpperBounds/BasePsiLeviLowerBound.lean}{ReciprocityTheory.ChartRestriction.G.AR.Up.ArchInM}
\leanref{ReciprocityTheory/ReductiveGroups/ArithmeticGeometry/ArakelovVolumes/UpperBounds/BasePsiLeviLowerBound.lean}{ReciprocityTheory.ChartRestriction.G.AR.Up.FlatM}
\leanref{Development/AbelianVarietyTheory/Definitions/ThetaValues/Generic.lean}{GlobalTheory.Shares.Margins}
\leanref{ReciprocityTheory/ReductiveGroups/ArithmeticGeometry/ArakelovVolumes/ArakelovInstances/IsStrictDom.lean}{ReciprocityTheory.ChartRestriction.G.AR.Inst.cpt}
\leanref{ReciprocityTheory/ReductiveGroups/ArithmeticGeometry/FundamentalDomains/Signature.lean}{ReciprocityTheory.ChartRestriction.G.FD.Data.firstMinimumSq}
\leanref{Development/AbelianVarietyTheory/Proofs/DivisorProofs/CardIndependentBoundedSectionsShiftedWeightLe.lean}{GlobalTheory.UpperShift.quotT}

\begin{lemma}\label{is-onto-share-counts-jets}
Let $X$ consist of the data (1) to (6) of flag data of rank $r$ over $E$, with archimedean data for $(L,F,T)$,
a flat model for them and the charts of Definition~\ref{is-onto-share-counts-model}, and let $e\geq0$ be an
integer. For
$\alpha\in\mathbf N^n$ write $|\alpha|=\sum_k\alpha_k$ and $s^\alpha=\prod_ks_k^{\alpha_k}$.
\begin{enumerate}
\item Let $1\le j\le\delta$ and let $f$ be a holomorphic function on $\mathcal D_X$. There are $a_\alpha\in\mathbf C$,
  for $|\alpha|\le e$, such that as $z\to0$
  \[
   f(c_{y_j}(z))-\sum_{|\alpha|\le e}a_\alpha\,s^\alpha(c_{y_j}(z))\,w_0(c_{y_j}(z))^{e-|\alpha|}=O(|z|^{e+1}).
  \]
\item Let $f_1,\dots,f_m$ be holomorphic functions on $\mathcal D_X$ with $m>\delta\binom{e+n}{n}$. Then there is
  a nonzero $a\in\mathbf C^m$ with $\sum_qa_qf_q(c_{y_j}(z))=O(|z|^{e+1})$ as $z\to0$, for every $j$.
\end{enumerate}
\end{lemma}
\leanref{Development/AbelianVarietyTheory/Proofs/TorusHaarMeasure/NmatZeroMemBor.lean}{GlobalTheory.LeadJets.jets}
\leanref{Development/AbelianVarietyTheory/Proofs/CuspFormEvaluations/ExistsMulNe.lean}{GlobalTheory.UpperModel.exists_comb}
\begin{proof}
(1) The point $y_j$ is a common zero of the $s_k$, and $c_{y_j}(0)=y_j$. As $\mathcal D_X$ is open
(Lemma~\ref{is-onto-share-sections}(2)), $c_{y_j}(z)\in\mathcal D_X$ near $0$, the functions $f$, $s_k$ and
$w_0$ composed with $c_{y_j}$ are holomorphic near $0$, and $w_0(c_{y_j}(z))\ne0$ near $0$. Put
$\zeta_k(z)=s_k(c_{y_j}(z))/w_0(c_{y_j}(z))$. Then $\zeta(0)=0$, the derivative of $\zeta$ at $0$ is invertible,
and $s^\alpha w_0^{e-|\alpha|}=\zeta^\alpha w_0^e$ at $c_{y_j}(z)$. Induction on $a\le e+1$ gives coefficients for which
the difference $R$ in (1) is $O(|z|^a)$. For $a=0$ take all $a_\alpha=0$; then $R$ is continuous at $0$. Let
$a\le e$ and $R=O(|z|^a)$. As $\zeta$ is a coordinate system near $0$, the Taylor expansion of $R$ in the
coordinates $\zeta$ gives $b_\alpha\in\mathbf C$, for $|\alpha|=a$, with $R-\sum_{|\alpha|=a}b_\alpha\zeta^\alpha=O(|z|^{a+1})$.
Replace $a_\alpha$ by $a_\alpha+b_\alpha w_0(y_j)^{-e}$ for $|\alpha|=a$. The new difference is
$R-\sum_{|\alpha|=a}b_\alpha\big(w_0(c_{y_j}(z))/w_0(y_j)\big)^e\zeta^\alpha$, and it differs from
$R-\sum_{|\alpha|=a}b_\alpha\zeta^\alpha$ by $\sum_{|\alpha|=a}b_\alpha\big(1-(w_0(c_{y_j}(z))/w_0(y_j))^e\big)\zeta^\alpha$, a sum
of terms $O(|z|)\cdot O(|z|^a)$. So the new difference is $O(|z|^{a+1})$.

(2) For each $q$ and $j$ choose coefficients $a_{q,j,\alpha}$ for $f_q$ at $y_j$ by (1). The vectors
$(a_{q,j,\alpha})_{j,\alpha}$, for $1\le q\le m$, lie in a space of dimension $\delta\binom{e+n}{n}$, the number of
pairs $(j,\alpha)$ with $|\alpha|\le e$. So there is a nonzero $a\in\mathbf C^m$ with $\sum_qa_qa_{q,j,\alpha}=0$ for all
$j$ and $\alpha$, and then
\[
 \sum_qa_qf_q(c_{y_j}(z))=\sum_qa_q\Big(f_q(c_{y_j}(z))-\sum_{|\alpha|\le e}a_{q,j,\alpha}\,s^\alpha(c_{y_j}(z))\,
 w_0(c_{y_j}(z))^{e-|\alpha|}\Big)=O(|z|^{e+1}).
\]
% Lean: (1) is LeadJets.jets with T = e, a0 = e + 1 and t_j(z)_k = 1 if k = j and 0 otherwise, as
% exists_comb uses it; its induction is LeadJets.step, with taylor_expansion_local_coordinates for the
% b_alpha and leadVal_add. (2) is exists_comb, with ManyForms.card_multiLt_succ for the count.
\end{proof}

\begin{lemma}\label{is-onto-share-counts-reduction-lll}
Let $N\geq0$ be an integer, $A$ a positive definite real symmetric $N\times N$ matrix, and write
$A(x)=\sum_{i,j}x_iA_{ij}x_j$ for $x\in\mathbf R^N$. Let $m>0$ and $D$ be real numbers such that $A(x)\geq m$ for
every nonzero $x\in\mathbf Z^N$ and $\det A\le D$. Then there is a basis $e_1,\dots,e_N$ of $\mathbf Z^N$ such that
$A(e_i)\le2^{N(N-1)}D/m^{N-1}$ for every $i$, and such that every $x=\sum_ic_ie_i$ in $\mathbf Z^N$ satisfies
$c_i^2\le9^NA(x)/m$ for every $i$.
\end{lemma}
\leanref{ReciprocityTheory/ReductiveGroups/ArithmeticGeometry/ArakelovVolumes/IntegralModuleBases/IntegralModuleBasisStatement.lean}{ReciprocityTheory.ChartRestriction.G.AR.IntegralModuleBases.RedGramStatement}
\leanref{ReciprocityTheory/ReductiveGroups/ArithmeticGeometry/ArakelovVolumes/IntegralModuleBases/DetToMatrixSq.lean}{ReciprocityTheory.ChartRestriction.G.AR.IntegralModuleBases.redGram_holds}
\leanref{ReciprocityTheory/ReducedBases/LatticeReduction.lean}{ReciprocityTheory.ReducedBases.T2.IsRed}
\leanref{ReciprocityTheory/ReducedBases/ExistsRedOf.lean}{ReciprocityTheory.ReducedBases.T2.exists_red_of}
\leanref{ReciprocityTheory/ReductiveGroups/ArithmeticGeometry/ArakelovVolumes/IntegralModuleBases/NormGsoSqLeTwoPowMulOfValEqAdd.lean}{ReciprocityTheory.ChartRestriction.G.AR.IntegralModuleBases.norm_sq_le_pow_gs}
\leanref{ReciprocityTheory/ReductiveGroups/ArithmeticGeometry/ArakelovVolumes/IntegralModuleBases/NormGsoSqLeTwoPowMulOfValEqAdd.lean}{ReciprocityTheory.ChartRestriction.G.AR.IntegralModuleBases.red_norm_sq_le}
\leanref{ReciprocityTheory/ReductiveGroups/ArithmeticGeometry/ArakelovVolumes/IntegralModuleBases/NormGsoSqLeTwoPowMulOfValEqAdd.lean}{ReciprocityTheory.ChartRestriction.G.AR.IntegralModuleBases.red_coeff_sq_le}
\leanref{ReciprocityTheory/ReducedBases/LatticeReduction.lean}{ReciprocityTheory.ReducedBases.T2.coeff_bound}
\begin{proof}
Let $B$ be the positive square root of $A$, so that $A(x)=|Bx|^2$ for the Euclidean norm. For a basis
$e_1,\dots,e_N$ of $\mathbf Z^N$ put $b_i=Be_i$, let $b^*_1,\dots,b^*_N$ be the Gram--Schmidt vectors of
$b_1,\dots,b_N$, and let $\mu_{ij}=\langle b_i,b^*_j\rangle/|b^*_j|^2$, so that $b_i=\sum_{j\le i}\mu_{ij}b^*_j$ with
$\mu_{ii}=1$. The lattice $B\mathbf Z^N$ of $\mathbf R^N$ has a reduced basis in the sense of Lenstra, Lenstra and
Lovász, which their reduction algorithm produces from any basis \cite[Section~1]{LLL82}. Choose $e$ such that
$b_1,\dots,b_N$ is reduced. Then $|\mu_{ij}|\le1/2$ for $j<i$, $|b^*_j|^2\le2|b^*_{j+1}|^2$ for $j<N$, and
$|b_i|^2\le2^{i-1}|b^*_i|^2$ for every $i$ \cite[Proposition~1.6 and its proof]{LLL82}.

For every $j$, $e_j$ is a nonzero integral vector, so $m\le A(e_j)=|b_j|^2\le2^{j-1}|b^*_j|^2$, and
$|b^*_j|^2\geq m/2^{N-1}$. The Gram matrix of $b_1,\dots,b_N$ is $E^{\mathsf T}AE$, where the integer matrix $E$
with columns $e_1,\dots,e_N$ has determinant $\pm1$; so $\prod_j|b^*_j|^2=\det A\le D$. Hence
$|b^*_i|^2\le D\,(2^{N-1}/m)^{N-1}$, and
\[
 A(e_i)=|b_i|^2\le2^{N-1}|b^*_i|^2\le2^{N(N-1)}D/m^{N-1}.
\]

Let $x=\sum_lc_le_l\in\mathbf Z^N$. Then $Bx=\sum_lc_lb_l=\sum_j\nu_jb^*_j$ with $\nu_j=\sum_{l\geq j}c_l\mu_{lj}$, and
$A(x)=\sum_j\nu_j^2|b^*_j|^2\geq\nu_j^2m/2^{N-1}$. So $|\nu_j|\le M$ for $M=(2^{N-1}A(x)/m)^{1/2}$. As
$\nu_j=c_j+\sum_{l>j}c_l\mu_{lj}$, $|c_j|\le M+\frac12\sum_{l>j}|c_l|$. Put $S_k=M+\frac12\sum_{l\geq k}|c_l|$. Then
$S_{N+1}=M$ and $S_k=S_{k+1}+\frac12|c_k|\le\frac32S_{k+1}$, so $S_k\le(3/2)^{N+1-k}M$. Hence
$|c_j|\le S_{j+1}\le(3/2)^NM$, and $c_j^2\le(9/4)^N2^{N-1}A(x)/m\le9^NA(x)/m$.
% Lean: redGram_holds, with B = CFC.sqrt A and the vectors fB B (e i). The reduced basis is exists_red_of
% (IsRed: |mu| <= 1/2 and |b*_j|^2 <= 2 |b*_{j+1}|^2), from li_basis and hlow_of (the potential, the
% product of the leading Gram determinants, is bounded below); its proof takes a basis whose potential is
% within a factor 4/3 of the infimum, size-reduces it (size_reduce) and swaps two neighbours where the second
% condition fails (pot_swap, norm_swap_j), the termination argument of LLL82. The bound
% |b_i|^2 <= 2^i |b*_i|^2 (indices from 0) is norm_sq_le_pow_gs, through sum_gs_le and sum_lt_le; then
% gs_sq_ge, det_gram_eq with det_gram_fB, red_norm_sq_le, red_coeff_sq_le (norm_sq_comb gives A(x)) and
% coeff_bound for the descending induction.
\end{proof}

\begin{lemma}\label{is-onto-share-counts-reduction-lattice}
Let $X$ consist of the data (1) to (6) of flag data of rank $r$ over $E$
(Definition~\ref{is-onto-share-flag-data}). Then $[E:\mathbf Q]=2|\Psi|$, and $\Lambda$ is a free $\mathbf Z$-module
of rank $N=r[E:\mathbf Q]$. Every $\mathbf Z$-basis $g_1,\dots,g_N$ of $\Lambda$ is a $\mathbf Q$-basis of $E^r$, and
the vectors $(\tau(g_i))_{\tau\in\Psi}$, for $1\le i\le N$, form an $\mathbf R$-basis of $(\mathbf C^r)^\Psi$.
\end{lemma}
\leanref{ReciprocityTheory/ReductiveGroups/ArithmeticGeometry/FundamentalDomains/RealPlaceEmbeddings.lean}{ReciprocityTheory.ChartRestriction.G.FD.RealPlaceEmbeddings.two_card_Psi}
\leanref{ReciprocityTheory/ReductiveGroups/ArithmeticGeometry/FundamentalDomains/RealPlaceEmbeddings.lean}{ReciprocityTheory.ChartRestriction.G.FD.RealPlaceEmbeddings.linIndQ}
\leanref{ReciprocityTheory/ReductiveGroups/ArithmeticGeometry/FundamentalDomains/RealPlaceEmbeddings.lean}{ReciprocityTheory.ChartRestriction.G.FD.RealPlaceEmbeddings.spanQ_eq_top}
\leanref{ReciprocityTheory/ReductiveGroups/ArithmeticGeometry/FundamentalDomains/RealPlaceEmbeddings.lean}{ReciprocityTheory.ChartRestriction.G.FD.RealPlaceEmbeddings.card_qb}
\leanref{ReciprocityTheory/ReductiveGroups/ArithmeticGeometry/FundamentalDomains/RealPlaceEmbeddings.lean}{ReciprocityTheory.ChartRestriction.G.FD.RealPlaceEmbeddings.span_real_embeddings_eq_top}
\leanref{ReciprocityTheory/ReductiveGroups/ArithmeticGeometry/FundamentalDomains/RealPlaceEmbeddings.lean}{ReciprocityTheory.ChartRestriction.G.FD.RealPlaceEmbeddings.span_embV}
\leanref{ReciprocityTheory/ReductiveGroups/ArithmeticGeometry/FundamentalDomains/RealPlaceEmbeddings.lean}{ReciprocityTheory.ChartRestriction.G.FD.RealPlaceEmbeddings.linIndR}
\begin{proof}
By (1) of Definition~\ref{is-onto-share-flag-data}, $\tau\mapsto\bar\tau$ is a bijection from $\Psi$ onto the
embeddings of $E$ that are not in $\Psi$, so $[E:\mathbf Q]=2|\Psi|$. The module $\Lambda$ is finitely generated
over $\mathcal O_E$, which is finitely generated over $\mathbf Z$, and it is torsion free; so it is a free
$\mathbf Z$-module of finite rank. Let $g$ be a $\mathbf Z$-basis. It is linearly independent over $\mathbf Q$, by
clearing denominators. Its $\mathbf Q$-span contains $\Lambda$ and is stable under multiplication by $E$, since
every element of $E$ is a $\mathbf Q$-combination of elements of $\mathcal O_E$ and $\Lambda$ is an
$\mathcal O_E$-module. As $\Lambda$ spans $E^r$ over $E$, the $\mathbf Q$-span of $g$ is $E^r$. So $g$ is a
$\mathbf Q$-basis of $E^r$, and $N=r[E:\mathbf Q]$.

The CM field $E$ has only complex places. Fix an embedding $\sigma_w$ at each place $w$. The vectors
$(\sigma_w(a))_w$ for $a\in\mathcal O_E$ form a lattice in $\prod_w\mathbf C$ that spans it over $\mathbf R$
\cite[Chapter~I, \S5]{Neu99}. Let $u\colon\prod_w\mathbf C\to\mathbf C^\Psi$ take at $\tau\in\Psi$ the coordinate $z_w$
at the place $w$ of $\tau$ if $\sigma_w=\tau$, and $\bar z_w$ otherwise. It is $\mathbf R$-linear, injective
because every place has an embedding in $\Psi$, and so bijective, both spaces having real dimension
$[E:\mathbf Q]$; and it carries $(\sigma_w(a))_w$ to $(\tau(a))_{\tau\in\Psi}$. Hence the vectors
$(\tau(v))_{\tau\in\Psi}$ for $v\in E^r$ span $(\mathbf C^r)^\Psi$ over $\mathbf R$. Each of them is a
$\mathbf Q$-combination of the $N$ vectors $(\tau(g_i))_\tau$, which therefore span this space of real dimension
$2r|\Psi|=N$ and form a basis of it.
% Lean: two_card_Psi; RealPlaceEmbeddings.finite_int and free_int; linIndQ and spanQ_eq_top (with
% smul_mem_spanQ) give the Q-basis qb, and card_qb the rank. The map u is Tm on TraceIntegralBasis.VS, the
% product over the complex places; Tm_injective, finrank_Psi and TraceIntegralBasis.finrank_VS make it
% bijective, and TraceIntegralBasis.bas, the image of the integral basis (Mathlib's
% mixedEmbedding.latticeBasis), spans VS; then span_real_embeddings_eq_top, span_embV and linIndR.
\end{proof}

\begin{lemma}\label{is-onto-share-counts-reduction-normal}
Let $X$ consist of the data (1) to (6) of flag data of rank $r$ over $E$
(Definition~\ref{is-onto-share-flag-data}), and let $x\in\mathcal D_X$. Then there is $b\in\mathcal B^\Psi$ such that
$y=xb$ is orthonormal and $m_{x,\tau}(v)=|y_\tau^{-1}v|^2$ for every $\tau\in\Psi$ and $v\in\mathbf C^r$, where
$|u|^2=\sum_a|u_a|^2$.
\end{lemma}
\leanref{ReciprocityTheory/ReductiveGroups/ArithmeticGeometry/FundamentalDomains/MajorantNormalForms.lean}{ReciprocityTheory.ChartRestriction.G.FD.MajorantNormalForms.exists_normal_frame}
\leanref{ReciprocityTheory/ReductiveGroups/ArithmeticGeometry/FundamentalDomains/MajorantNormalForms.lean}{ReciprocityTheory.ChartRestriction.G.FD.MajorantNormalForms.hgram_posDef}
\leanref{ReciprocityTheory/ReductiveGroups/ArithmeticGeometry/FundamentalDomains/MajorantNormalForms.lean}{ReciprocityTheory.ChartRestriction.G.FD.MajorantNormalForms.exists_normal}
\leanref{ReciprocityTheory/ReductiveGroups/ArithmeticGeometry/FundamentalDomains/MajorantNormalForms.lean}{ReciprocityTheory.ChartRestriction.G.FD.MajorantNormalForms.gram_normal}
\leanref{ReciprocityTheory/ReductiveGroups/ArithmeticGeometry/FundamentalDomains/MajorantNormalForms.lean}{ReciprocityTheory.ChartRestriction.G.FD.MajorantNormalForms.maj_normal}
\leanref{ReciprocityTheory/ReductiveGroups/ArithmeticGeometry/FundamentalDomains/MajPos.lean}{ReciprocityTheory.ChartRestriction.G.FD.Data.Hf_self}
\begin{proof}
Fix $\tau$ and write $h=h_\tau$. Let $x_1,\dots,x_r$ be the columns of $x_\tau$, $U$ the span of
$x_1,\dots,x_{p_0}$, and $g$ their Gram matrix, which is invertible since $d_{p_0}(x,\tau)\ne0$. Put
$Pu=\sum_{a,b\le p_0}h(u,x_b)(g^{-1})_{ba}x_a$. Then $h(Pu,x_c)=h(u,x_c)$ for $c\le p_0$, so $P$ is the
$h$-orthogonal projection onto $U$: it fixes $U$, it vanishes on the vectors $h$-orthogonal to $U$, and
$h(Pu,w)=h(u,Pw)$. Put $H(u,w)=h(2Pu-u,w)$. It is a Hermitian form, and $H(v,v)=m_{x,\tau}(v)$ by the definition
of the majorant; so $H$ is positive definite by Lemma~\ref{is-onto-share-majorant}(1).

The Gram--Schmidt process for $H$, applied to $x_1,\dots,x_r$, gives an invertible upper triangular matrix
$\beta_\tau$ such that the columns $y_1,\dots,y_r$ of $x_\tau\beta_\tau$ satisfy $H(y_a,y_c)=\delta_{ac}$. For
$a\le p_0$, $y_a$ is a combination of $x_1,\dots,x_a$, so $Py_a=y_a$ and $H(y_a,w)=h(y_a,w)$ for every $w$. For
$a>p_0$ and $c\le p_0$, $0=H(y_a,y_c)=2h(y_a,Py_c)-h(y_a,y_c)=h(y_a,y_c)$. As $\beta_\tau^{-1}$ is upper
triangular, $x_1,\dots,x_{p_0}$ are combinations of $y_1,\dots,y_{p_0}$; so for $a>p_0$, $y_a$ is $h$-orthogonal
to $U$, $Py_a=0$ and $H(y_a,w)=-h(y_a,w)$. Hence $h(y_a,y_c)=\operatorname{sgn}(a)\delta_{ac}$ for all $a$ and $c$:
the Gram matrix of $x_\tau\beta_\tau$ is $J$. Let $b=(\beta_\tau)_\tau$ and $y=xb$. For $v\in\mathbf C^r$ write
$v=\sum_au_ay_a$ with $u=y_\tau^{-1}v$; then $m_{x,\tau}(v)=H(v,v)=\sum_a|u_a|^2$.
% Lean: exists_normal_frame. H is Data.Hf, with P = Data.ppos (ppos_colP, hermC_ppos_lemma_5) and
% Hf_conj_lemma_7, maj_eq_lemma_4 for H(v, v) = m(v); hgram_posDef uses Data.maj_pos. exists_normal takes
% the LDL decomposition (Mathlib's LDL.lowerInv, built by Gram-Schmidt) of the Gram matrix of the columns
% for H and B = lowerInv^T diag(D^(-1/2)); gram_normal and maj_normal are the last two steps.
\end{proof}

\begin{lemma}\label{is-onto-share-counts-reduction-forms}
Let $X$ consist of the data (1) to (6) of flag data of rank $r$ over $E$
(Definition~\ref{is-onto-share-flag-data}), and let $y\in M_r(\mathbf C)^\Psi$ have invertible components. For
$v\in E^r$ put $q_y(v)=\sum_{\tau\in\Psi}c_\tau|y_\tau^{-1}\tau(v)|^2$, where $|u|^2=\sum_a|u_a|^2$, and put
$C_X=\sum_{\tau\in\Psi}c_\tau^{-1}$.
\begin{enumerate}
\item Let $g_1,\dots,g_N$ be a $\mathbf Z$-basis of $\Lambda$. There is a positive definite real symmetric
  $N\times N$ matrix $G_y$ with $\sum_{i,j}x_i(G_y)_{ij}x_j=q_y\big(\sum_ix_ig_i\big)$ for $x\in\mathbf Z^N$. If $y$
  is orthonormal, then $\det G_y\le|\det B_g|$ for the matrix with entries
  $(B_g)_{ij}=\sum_{\tau\in\Psi}c_\tau\operatorname{Re}\tau\big(h(g_i,g_j)\big)$.
\item If $y$ is orthonormal, then $|\sigma(h(p,q))|\le\big(q_y(p)+q_y(q)\big)C_X$ for every embedding
  $\sigma\colon E\to\mathbf C$ and all $p,q\in E^r$.
\item $q_y(\alpha p)\le\big(\sum_{\tau\in\Psi}|\tau(\alpha)|^2\big)q_y(p)$ for $\alpha\in E$ and $p\in E^r$.
\item Let $\tau\in\Psi$, $1\le b\le r$, $a_1,\dots,a_k\in\mathbf Q$ and $f_1,\dots,f_k\in E^r$. Then
  $\big|\big(y_\tau^{-1}\tau(\sum_ia_if_i)\big)_b\big|\le\sum_i|a_i|\big(1+q_y(f_i)/c_\tau\big)$.
\end{enumerate}
\end{lemma}
\leanref{ReciprocityTheory/ReductiveGroups/ArithmeticGeometry/FundamentalDomains/QuadraticLatticeForms.lean}{ReciprocityTheory.ChartRestriction.G.FD.QuadraticLatticeForms.Qy}
\leanref{ReciprocityTheory/ReductiveGroups/ArithmeticGeometry/FundamentalDomains/QuadraticLatticeForms.lean}{ReciprocityTheory.ChartRestriction.G.FD.QuadraticLatticeForms.Am}
\leanref{ReciprocityTheory/ReductiveGroups/ArithmeticGeometry/FundamentalDomains/QuadraticLatticeForms.lean}{ReciprocityTheory.ChartRestriction.G.FD.QuadraticLatticeForms.Bm}
\leanref{ReciprocityTheory/ReductiveGroups/ArithmeticGeometry/FundamentalDomains/QuadraticLatticeForms.lean}{ReciprocityTheory.ChartRestriction.G.FD.QuadraticLatticeForms.Am_posDef}
\leanref{ReciprocityTheory/ReductiveGroups/ArithmeticGeometry/FundamentalDomains/QuadraticLatticeForms.lean}{ReciprocityTheory.ChartRestriction.G.FD.QuadraticLatticeForms.quad_int}
\leanref{ReciprocityTheory/ReductiveGroups/ArithmeticGeometry/FundamentalDomains/IntegralBasisReduction.lean}{ReciprocityTheory.ChartRestriction.G.FD.IntegralBasisReduction.ip_Sg}
\leanref{ReciprocityTheory/ReductiveGroups/ArithmeticGeometry/FundamentalDomains/IntegralBasisReduction.lean}{ReciprocityTheory.ChartRestriction.G.FD.IntegralBasisReduction.Bm_eq}
\leanref{ReciprocityTheory/ReductiveGroups/ArithmeticGeometry/FundamentalDomains/IntegralBasisReduction.lean}{ReciprocityTheory.ChartRestriction.G.FD.IntegralBasisReduction.det_Am_le}
\leanref{ReciprocityTheory/ReductiveGroups/ArithmeticGeometry/FundamentalDomains/HermitianPsiBounds.lean}{ReciprocityTheory.ChartRestriction.G.FD.HermitianPsiBounds.Cinv}
\leanref{ReciprocityTheory/ReductiveGroups/ArithmeticGeometry/FundamentalDomains/HermitianPsiBounds.lean}{ReciprocityTheory.ChartRestriction.G.FD.HermitianPsiBounds.herm_bound}
\leanref{ReciprocityTheory/ReductiveGroups/ArithmeticGeometry/FundamentalDomains/HermitianPsiBounds.lean}{ReciprocityTheory.ChartRestriction.G.FD.HermitianPsiBounds.Qy_smul}
\leanref{ReciprocityTheory/ReductiveGroups/ArithmeticGeometry/FundamentalDomains/CompactEntryBounds.lean}{ReciprocityTheory.ChartRestriction.G.FD.CompactEntryBounds.comb_bound}
\begin{proof}
Each $c_\tau=2|\tau(\xi_0)|$ is positive, as $\xi_0\ne0$. For $p,q\in E^r$, $\tau(h(p,q))=h_\tau(\tau(p),\tau(q))$,
as $\tau(\bar a)=\overline{\tau(a)}$ on the CM field $E$. If the Gram matrix of $y_\tau$ for $h_\tau$ is $J$, then
for $p',q'\in\mathbf C^r$, with $u=y_\tau^{-1}p'$ and $w=y_\tau^{-1}q'$ their coordinates in the columns of $y_\tau$,
\[
 h_\tau(p',q')=\sum_a\operatorname{sgn}(a)\,u_a\bar w_a.
\]

(1) On $V=(\mathbf C^r)^\Psi$ let $\langle u,w\rangle=\sum_\tau c_\tau\operatorname{Re}\sum_au_{\tau,a}\bar w_{\tau,a}$, a real
inner product, and $L_y(u)=(y_\tau^{-1}u_\tau)_\tau$, an $\mathbf R$-linear automorphism of $V$. By
Lemma~\ref{is-onto-share-counts-reduction-lattice} the vectors $z_i=L_y\big((\tau(g_i))_\tau\big)$ form an
$\mathbf R$-basis of $V$, so their Gram matrix $G_y=(\langle z_i,z_j\rangle)_{i,j}$ is positive definite. For
$x\in\mathbf Z^N$, $\sum_ix_iz_i=L_y\big((\tau(\sum_ix_ig_i))_\tau\big)$, whose squared norm is
$q_y(\sum_ix_ig_i)$. Now let $y$ be orthonormal, and let $S$ be the $\mathbf R$-linear map of $V$ that multiplies the
coordinate $(\tau,a)$ by $\operatorname{sgn}(a)$. By the formula above, $(B_g)_{ij}=\langle z_i,Sz_j\rangle$. With
$M$ the matrix of $S$ in the basis $(z_i)$, this says $B_g=G_yM$. As $S^2=1$, $\det M=\pm1$, so
$\det G_y=|\det B_g|$.

(2) Let $\sigma\in\Psi$, $u=y_\sigma^{-1}\sigma(p)$ and $w=y_\sigma^{-1}\sigma(q)$. By the formula above,
$|\sigma(h(p,q))|\le\sum_a|u_a||w_a|\le|u|^2+|w|^2$, and $|u|^2\le q_y(p)/c_\sigma\le q_y(p)C_X$; likewise
$|w|^2\le q_y(q)C_X$. If $\sigma\notin\Psi$, then $\bar\sigma\in\Psi$ by (1) of
Definition~\ref{is-onto-share-flag-data}, and $|\sigma(h(p,q))|=|\bar\sigma(h(p,q))|$.

(3) The term of $q_y(\alpha p)$ at $\tau$ is $|\tau(\alpha)|^2$ times that of $q_y(p)$, and
$|\tau(\alpha)|^2\le\sum_{\tau'\in\Psi}|\tau'(\alpha)|^2$.

(4) As $\tau(\sum_ia_if_i)=\sum_ia_i\tau(f_i)$, the left side is at most $\sum_i|a_i|t_i$ with
$t_i=|(y_\tau^{-1}\tau(f_i))_b|$. Here $c_\tau t_i^2\le q_y(f_i)$, and $t_i\le1+t_i^2\le1+q_y(f_i)/c_\tau$.
% Lean: QuadraticLatticeForms: c_pos; Qy, ip (the inner product) and Ly; G_y is Am, positive definite by
% Am_posDef through RealPlaceEmbeddings.linIndR, and quad_int is the identity for integral x. B_g is Bm;
% IntegralBasisReduction.ip_Sg (with hermC_normal for the formula above), Bm_eq (B_g = G_y M, M the matrix of
% Sg in zbas) and det_Am_le, which states det G_y <= |det B_g|. (2) is HermitianPsiBounds.herm_bound, through
% herm_bound_psi and sum_sq_le; (3) is Qy_smul; (4) is CompactEntryBounds.comb_bound, through comp_le_Qy.
\end{proof}

\begin{lemma}\label{is-onto-share-counts-reduction-basis}
Let $X$ consist of the data (1) to (6) of flag data of rank $r$ over $E$
(Definition~\ref{is-onto-share-flag-data}), fix a $\mathbf Z$-basis $g_1,\dots,g_N$ of $\Lambda$, and let $q_y$ and
$B_g$ be as in Lemma~\ref{is-onto-share-counts-reduction-forms}. Let $y$ be an orthonormal point and $m>0$ a real
number with $q_y(v)\geq m$ for every nonzero $v\in\Lambda$. Then $\Lambda$ has a $\mathbf Z$-basis $f_1,\dots,f_N$
such that $q_y(f_i)\le2^{N(N-1)}|\det B_g|/m^{N-1}$ for every $i$, and such that every $v=\sum_ic_if_i$ in
$\Lambda$ satisfies $c_i^2\le9^Nq_y(v)/m$ for every $i$.
\end{lemma}
\leanref{ReciprocityTheory/ReductiveGroups/ArithmeticGeometry/FundamentalDomains/IntegralBasisReduction.lean}{ReciprocityTheory.ChartRestriction.G.FD.IntegralBasisReduction.IsNormal}
\leanref{ReciprocityTheory/ReductiveGroups/ArithmeticGeometry/FundamentalDomains/IntegralBasisReduction.lean}{ReciprocityTheory.ChartRestriction.G.FD.IntegralBasisReduction.reduce}
\begin{proof}
The components of $y$ are invertible (Lemma~\ref{is-onto-share-majorant}(4)). Let $G_y$ be the matrix of
Lemma~\ref{is-onto-share-counts-reduction-forms}(1), and write $G_y(x)=\sum_{i,j}x_i(G_y)_{ij}x_j$. For nonzero
$x\in\mathbf Z^N$, $\sum_ix_ig_i$ is a nonzero element of $\Lambda$, so $G_y(x)=q_y(\sum_ix_ig_i)\geq m$; and
$\det G_y\le|\det B_g|$. Lemma~\ref{is-onto-share-counts-reduction-lll} for $G_y$, $m$ and $D=|\det B_g|$ gives a
basis $e_1,\dots,e_N$ of $\mathbf Z^N$. Put $f_i=\sum_l(e_i)_lg_l$. These form a $\mathbf Z$-basis of $\Lambda$, and
$q_y(f_i)=G_y(e_i)\le2^{N(N-1)}|\det B_g|/m^{N-1}$. If $v=\sum_ic_if_i\in\Lambda$, then $v=\sum_lx_lg_l$ with
$x=\sum_ic_ie_i$, so $c_i^2\le9^NG_y(x)/m=9^Nq_y(v)/m$.
% Lean: reduce, for g = RealPlaceEmbeddings.zb X (Module.finBasis); quad_int, Am_posDef, det_Am_le and
% redGram_holds, with f = eb.map g.equivFun.symm.
\end{proof}

\begin{lemma}\label{is-onto-share-counts-reduction-types}
Let $X$ consist of the data (1) to (6) of flag data of rank $r$ over $E$
(Definition~\ref{is-onto-share-flag-data}), and let $\omega_1,\dots,\omega_n$ be a $\mathbf Z$-basis of $\mathcal O_E$.
Let $f$ and $f'$ be $\mathbf Z$-bases of $\Lambda$ such that $h(f_i,f_j)=h(f'_i,f'_j)$ for all $i$ and $j$, and such
that for all $k$ and $i$ the coordinates of $\omega_kf_i$ in the basis $f$ are those of $\omega_kf'_i$ in the basis
$f'$. Then there is $\gamma\in GL_r(E)$ with $\gamma\Lambda=\Lambda$ and $h(\gamma u,\gamma w)=h(u,w)$ for
$u,w\in E^r$, such that $\gamma f_i=f'_i$ for every $i$, and $\gamma^{-1}v=\sum_ia_if_i$ for every
$v=\sum_ia_if'_i$ with $a_i\in\mathbf Q$.
\end{lemma}
\leanref{ReciprocityTheory/ReductiveGroups/ArithmeticGeometry/FundamentalDomains/HermitianLatticeMaps.lean}{ReciprocityTheory.ChartRestriction.G.FD.HermitianLatticeMaps.ty}
\leanref{ReciprocityTheory/ReductiveGroups/ArithmeticGeometry/FundamentalDomains/HermitianLatticeMaps.lean}{ReciprocityTheory.ChartRestriction.G.FD.HermitianLatticeMaps.exists_gl}
\begin{proof}
By Lemma~\ref{is-onto-share-counts-reduction-lattice}, $f$ and $f'$ are $\mathbf Q$-bases of $E^r$. Let $\varphi$ be
the $\mathbf Q$-linear map of $E^r$ with $\varphi(f_i)=f'_i$; its inverse sends $f'_i$ to $f_i$. For all $k$ and $i$,
$\varphi(\omega_kf_i)=\omega_kf'_i=\omega_k\varphi(f_i)$, by the hypothesis on the coordinates. Every
$\alpha\in E$ is a $\mathbf Q$-combination of the $\omega_k$, so $\varphi(\alpha f_i)=\alpha\varphi(f_i)$; as both sides
are $\mathbf Q$-linear in the vector, $\varphi(\alpha v)=\alpha\varphi(v)$ for every $v$, and $\varphi$ is $E$-linear.
The forms $(u,w)\mapsto h(\varphi u,\varphi w)$ and $h$ are $\mathbf Q$-bilinear and agree on the pairs $(f_i,f_j)$,
so they agree. And $\varphi$ maps $\Lambda=\sum_i\mathbf Zf_i$ onto $\sum_i\mathbf Zf'_i=\Lambda$. Let $\gamma$ be the
matrix of $\varphi$. Then $\gamma^{-1}(\sum_ia_if'_i)=\sum_ia_if_i$.
% Lean: exists_gl, with phi (the Q-linear map of the Q-bases qb), phi_omega and phi_smul for E-linearity,
% phi_herm (through hL and hR), phi_mem and phi_phi; ty is the pair of the Gram matrix and the
% coordinates of (RingOfIntegers.basis E k) * f i.
\end{proof}

\begin{lemma}\label{is-onto-share-counts-reduction-compact}
Let $X$ consist of the data (1) to (6) of flag data of rank $r$ over $E$
(Definition~\ref{is-onto-share-flag-data}), let $J$ be the diagonal matrix with entries
$\operatorname{sgn}(1),\dots,\operatorname{sgn}(r)$, and let $\tau\in\Psi$.
\begin{enumerate}
\item If $Y\in M_r(\mathbf C)$ has Gram matrix $J$ for $h_\tau$, then $Y$ is invertible and
  $|Y_{ab}|\le\sum_k|((A^\tau)^{-1})_{ka}|\,|(Y^{-1})_{bk}|$ for all $a$ and $b$.
\item Let $\gamma\in GL_r(E)$ with $h(\gamma u,\gamma w)=h(u,w)$ for $u,w\in E^r$, and let $Y\in M_r(\mathbf C)$.
  Then $\gamma^\tau Y$ and $Y$ have the same Gram matrix for $h_\tau$; and if $Y$ is invertible, then
  $(\gamma^\tau Y)^{-1}\tau(v)=Y^{-1}\tau(\gamma^{-1}v)$ for $v\in E^r$.
\item For $R\geq0$, the set $K_R$ of the orthonormal points $y\in M_r(\mathbf C)^\Psi$ with $|(y_\tau)_{ab}|\le R$ for
  all $\tau$, $a$ and $b$ is compact and contained in $\mathcal D_X$.
\end{enumerate}
\end{lemma}
\leanref{ReciprocityTheory/ReductiveGroups/ArithmeticGeometry/FundamentalDomains/CompactEntryBounds.lean}{ReciprocityTheory.ChartRestriction.G.FD.CompactEntryBounds.det_ne_of_gram}
\leanref{ReciprocityTheory/ReductiveGroups/ArithmeticGeometry/FundamentalDomains/CompactEntryBounds.lean}{ReciprocityTheory.ChartRestriction.G.FD.CompactEntryBounds.entry_bound}
\leanref{ReciprocityTheory/ReductiveGroups/ArithmeticGeometry/FundamentalDomains/CompactEntryBounds.lean}{ReciprocityTheory.ChartRestriction.G.FD.CompactEntryBounds.gram_stab}
\leanref{ReciprocityTheory/ReductiveGroups/ArithmeticGeometry/FundamentalDomains/CompactEntryBounds.lean}{ReciprocityTheory.ChartRestriction.G.FD.CompactEntryBounds.inv_lact}
\leanref{ReciprocityTheory/ReductiveGroups/ArithmeticGeometry/FundamentalDomains/CompactEntryBounds.lean}{ReciprocityTheory.ChartRestriction.G.FD.CompactEntryBounds.Kset}
\leanref{ReciprocityTheory/ReductiveGroups/ArithmeticGeometry/FundamentalDomains/CompactEntryBounds.lean}{ReciprocityTheory.ChartRestriction.G.FD.CompactEntryBounds.Kset_compact}
\leanref{ReciprocityTheory/ReductiveGroups/ArithmeticGeometry/FundamentalDomains/CompactEntryBounds.lean}{ReciprocityTheory.ChartRestriction.G.FD.CompactEntryBounds.Kset_sub}
\begin{proof}
(1) The Gram matrix of $Y$ for $h_\tau$ is $Y^{\mathsf T}A^\tau\bar Y$, and $\det A\ne0$. If it is $J$, then
$\det Y\ne0$, and transposing gives $Y^*(A^\tau)^{\mathsf T}Y=J$, so $Y=((A^\tau)^{\mathsf T})^{-1}(Y^*)^{-1}J$. Hence
$Y_{ab}=\operatorname{sgn}(b)\sum_k((A^\tau)^{-1})_{ka}\overline{(Y^{-1})_{bk}}$.

(2) For $u,w\in\mathbf C^r$,
$h_\tau(\gamma^\tau u,\gamma^\tau w)=\sum_{k,l}u_k\bar w_l\,\tau\big(h(\gamma e_k,\gamma e_l)\big)
=\sum_{k,l}u_k\bar w_l\,\tau\big(h(e_k,e_l)\big)=h_\tau(u,w)$; applied to the columns of $Y$, this gives the
first claim. The second follows from $(\gamma^\tau)^{-1}=(\gamma^{-1})^\tau$ and
$(\gamma^{-1})^\tau\tau(v)=\tau(\gamma^{-1}v)$.

(3) The Gram matrices of the $y_\tau$ depend continuously on $y$, so $K_R$ is closed; it is bounded, so it is
compact. Let $y\in K_R$. By (1) every $y_\tau$ is invertible, and the leading principal minors of $J$ give
$d_i(y,\tau)=\epsilon_i$, so $\epsilon_i\operatorname{Re}d_i(y,\tau)=\epsilon_i^2=1>0$. So $y\in\mathcal D_X$.
% Lean: (1) is det_ne_of_gram and entry_bound (gram_eq_transpose_mul_map_mul_fundamental_domains);
% (2) is gram_stab (hermC_stab, hermC_expand) and inv_lact; (3) is Kset with Kset_compact (closed by
% cont_gram', bounded) and Kset_sub.
\end{proof}

\begin{lemma}\label{is-onto-share-counts-reduction}
Let $X$ consist of the data (1) to (6) of flag data of rank $r$ over $E$
(Definition~\ref{is-onto-share-flag-data}), with $\lambda_1$ as in Definition~\ref{is-onto-share-counts-model}, and
let $\varepsilon>0$. Then there is a compact set $K\subset\mathcal D_X$ such that for every $x\in\mathcal D_X$ with
$\lambda_1(x)\geq\varepsilon$ there are $\gamma\in\Gamma$ and $b\in\mathcal B^\Psi$ with $\gamma xb\in K$.
\end{lemma}
\leanref{ReciprocityTheory/ReductiveGroups/ArithmeticGeometry/FundamentalDomains/ModelNefness.lean}{ReciprocityTheory.ChartRestriction.G.FD.reductionTheory_holds}
\leanref{ReciprocityTheory/ReductiveGroups/ArithmeticGeometry/FundamentalDomains/CuspidalExponentialDecay.lean}{ReciprocityTheory.ChartRestriction.G.FD.MahlerCompactnessCriterion}
\leanref{ReciprocityTheory/ReductiveGroups/ArithmeticGeometry/FundamentalDomains/HermitianPsiBounds.lean}{ReciprocityTheory.ChartRestriction.G.FD.HermitianPsiBounds.sq_le_Qy_of_le_firstMinimum}
\leanref{ReciprocityTheory/ReductiveGroups/ArithmeticGeometry/FundamentalDomains/FiniteTranslateSets.lean}{ReciprocityTheory.ChartRestriction.G.FD.FiniteTranslateSets.TS}
\leanref{ReciprocityTheory/ReductiveGroups/ArithmeticGeometry/FundamentalDomains/FiniteTranslateSets.lean}{ReciprocityTheory.ChartRestriction.G.FD.FiniteTranslateSets.TS_finite}
\leanref{ReciprocityTheory/ReductiveGroups/ArithmeticGeometry/FundamentalDomains/FiniteTranslateSets.lean}{ReciprocityTheory.ChartRestriction.G.FD.FiniteTranslateSets.Sh_finite}
\begin{proof}
Let $N=r[E:\mathbf Q]$ be the rank of $\Lambda$ (Lemma~\ref{is-onto-share-counts-reduction-lattice}). Fix a
$\mathbf Z$-basis $g$ of $\Lambda$, a $\mathbf Z$-basis $\omega_1,\dots,\omega_n$ of $\mathcal O_E$, and the element
$d\ne0$ of $\mathcal O_E$ of (3) of Definition~\ref{is-onto-share-flag-data}; let $q_y$, $B_g$ and $C_X$ be as in
Lemma~\ref{is-onto-share-counts-reduction-forms}. Put
\[
 B_0=\frac{2^{N(N-1)}|\det B_g|}{\varepsilon^{2(N-1)}},\qquad
 R_0=2B_0C_X\sum_\sigma|\sigma(d)|,\qquad
 T_0=\frac{9^NB_0}{\varepsilon^2}\sum_{k=1}^n\sum_{\tau\in\Psi}|\tau(\omega_k)|^2,
\]
where $\sigma$ runs over the embeddings of $E$ into $\mathbf C$.

\emph{Finitely many types.} The \emph{type} of a $\mathbf Z$-basis $f$ of $\Lambda$ is the pair of the matrix
$(h(f_i,f_j))_{i,j}$ and the integers $n_{kij}$ with $\omega_kf_i=\sum_jn_{kij}f_j$. Let $\mathcal F$ be the set of
the types of the $\mathbf Z$-bases $f$ of $\Lambda$ with $|\sigma(dh(f_i,f_j))|\le R_0$ for all $\sigma$, $i$ and $j$,
and $n_{kij}^2\le T_0$ for all $k$, $i$ and $j$. Each $dh(f_i,f_j)$ is an algebraic integer, by (3) of
Definition~\ref{is-onto-share-flag-data}. Only finitely many elements of $\mathcal O_E$ have all their conjugates
bounded by $R_0$, as $\mathcal O_E$ is a lattice in the Minkowski space of $E$ \cite[Chapter~I, \S5]{Neu99}; and only
finitely many integers $n$ satisfy $n^2\le T_0$. So $\mathcal F$ is finite. For each $t\in\mathcal F$ choose a
$\mathbf Z$-basis $f^t$ of $\Lambda$ of type $t$. By (4) of Definition~\ref{is-onto-share-flag-data} there is a
finite set $S\subset GL_r(E)$ such that every element of $GL_r(E)$ that preserves $h$ and $\Lambda$ lies in
$s\Gamma$ for some $s\in S$. For $t\in\mathcal F$, $s\in S$ and $1\le k\le r$, let $a_{tski}\in\mathbf Q$ be the
coordinates of $se_k$ in the $\mathbf Q$-basis $f^t$ of $E^r$ (Lemma~\ref{is-onto-share-counts-reduction-lattice}).
Put
\[
 Q_0=\sum_{t\in\mathcal F}\sum_{s\in S}\sum_{k=1}^r\sum_{i=1}^N|a_{tski}|,\qquad
 G_0=\sum_{\tau\in\Psi}\sum_{k,a}|((A^\tau)^{-1})_{ka}|,
\]
and let $K=K_{R_1}$ for $R_1=rG_0Q_0(1+B_0C_X)$, a compact subset of $\mathcal D_X$ by
Lemma~\ref{is-onto-share-counts-reduction-compact}(3).

\emph{A reduced basis.} Let $x\in\mathcal D_X$ with $\lambda_1(x)\geq\varepsilon$.
Lemma~\ref{is-onto-share-counts-reduction-normal} gives $b\in\mathcal B^\Psi$ such that $y=xb$ is orthonormal and
$m_{x,\tau}(u)=|y_\tau^{-1}u|^2$ for all $\tau$ and $u$. For nonzero $v\in\Lambda$, then
$q_y(v)=\sum_\tau c_\tau m_{x,\tau}(\tau(v))\geq\lambda_1(x)^2\geq\varepsilon^2$. So
Lemma~\ref{is-onto-share-counts-reduction-basis}, with $m=\varepsilon^2$, gives a $\mathbf Z$-basis $f$ of $\Lambda$
with $q_y(f_i)\le B_0$ for every $i$, such that the coordinates $c_i$ of every $v\in\Lambda$ in this basis satisfy
$c_i^2\le9^Nq_y(v)/\varepsilon^2$.

\emph{Its type lies in $\mathcal F$.} By Lemma~\ref{is-onto-share-counts-reduction-forms}(2), for every $\sigma$,
\[
 |\sigma(dh(f_i,f_j))|=|\sigma(d)|\,|\sigma(h(f_i,f_j))|\le\Big(\sum_{\sigma'}|\sigma'(d)|\Big)
 \big(q_y(f_i)+q_y(f_j)\big)C_X\le R_0.
\]
The integer $n_{kij}$ is the $j$-th coordinate of $\omega_kf_i\in\Lambda$, so by
Lemma~\ref{is-onto-share-counts-reduction-forms}(3)
\[
 n_{kij}^2\le\frac{9^N}{\varepsilon^2}\,q_y(\omega_kf_i)\le\frac{9^N}{\varepsilon^2}\Big(\sum_{\tau\in\Psi}|\tau(\omega_k)|^2\Big)B_0
 \le T_0.
\]

\emph{The translate.} Let $t\in\mathcal F$ be the type of $f$. Lemma~\ref{is-onto-share-counts-reduction-types}, for
$f$ and $f^t$, gives $\gamma\in GL_r(E)$ that preserves $h$ and $\Lambda$, with $\gamma^{-1}v=\sum_ia_if_i$ whenever
$v=\sum_ia_if^t_i$ with $a_i\in\mathbf Q$. Choose $s\in S$ with $\gamma_1=s^{-1}\gamma\in\Gamma$. It remains to see
that $\gamma_1xb=\gamma_1y$ lies in $K$. Let $\tau\in\Psi$ and $Y=\gamma_1^\tau y_\tau$. By
Lemma~\ref{is-onto-share-counts-reduction-compact}(2), $Y$ has the Gram matrix $J$ of $y_\tau$. For $1\le k\le r$,
$\gamma_1^{-1}e_k=\gamma^{-1}(se_k)=\sum_ia_{tski}f_i$. So, by
Lemma~\ref{is-onto-share-counts-reduction-compact}(2) and Lemma~\ref{is-onto-share-counts-reduction-forms}(4),
\[
 |(Y^{-1})_{bk}|=\big|\big(y_\tau^{-1}\tau(\gamma_1^{-1}e_k)\big)_b\big|
 \le\sum_i|a_{tski}|\big(1+q_y(f_i)/c_\tau\big)\le Q_0(1+B_0C_X),
\]
as $1/c_\tau\le C_X$. By Lemma~\ref{is-onto-share-counts-reduction-compact}(1),
$|Y_{ab}|\le\sum_k|((A^\tau)^{-1})_{ka}|\,|(Y^{-1})_{bk}|\le rG_0Q_0(1+B_0C_X)=R_1$. So $\gamma_1y\in K$.
% Lean: reductionTheory_holds proves MahlerCompactnessCriterion. B_0 is B (the bound of reduce for
% m = eps^2, with Bm X (RealPlaceEmbeddings.zb X)); R_0 is R = D ((B + B) Cinv), D the sum of |sigma(den)|
% over all embeddings; T_0 is T = 9^N (W B) / eps^2. The types are HermitianLatticeMaps.ty and F is
% FiniteTranslateSets.TS X R T, finite by TS_finite through Sh_finite (Mathlib's
% NumberField.Embeddings.finite_of_norm_le) and Sz_finite; rep chooses f^t, S comes from Gamma_index, the
% a_{tski} are q t s k i (the coordinates in qb (rep t)), Q_0 and G_0 are Q and G0, and K is
% CompactEntryBounds.Kset for the sum over k of G0 (Q (1 + B Cinv)). The steps are exists_normal_frame,
% sq_le_Qy_of_le_firstMinimum, reduce, herm_bound with Lambda_int, Qy_smul, exists_gl, gram_stab,
% entry_bound, inv_lact and comb_bound. About 2,800 lines of Lean in 190 declarations lie below
% reductionTheory_holds.
\end{proof}

\begin{lemma}\label{is-onto-share-counts-maximum}
Let $X$ consist of the data (1) to (6) of flag data of rank $r$ over $E$, with archimedean data for
$(L,F,T)$ and a flat model for them, and the notation of Definition~\ref{is-onto-share-counts-model}. Let
$T'\geq0$ be an integer with $cg(2T'+2n+13)<4$, and $f$ a
bounded section of weight $F+T'L$ with $f(c_{y_j}(z))=O(|z|^{T'+n+7})$ as $z\to0$, for every $j$. If
$Q_{T',f}(x)>0$ for some $x\in\mathcal D_X$, then $Q_{T',f}$ attains its maximum on $\mathcal D_X$ at a point $x_0$
with $Q_{T',f}(x_0)>0$ at which some $s_k$ does not vanish.
\end{lemma}
\leanref{Development/AbelianVarietyTheory/Proofs/DivisorProofs/CardIndependentBoundedSectionsShiftedWeightLe.lean}{GlobalTheory.UpperShift.mainT}
\leanref{Development/AbelianVarietyTheory/Proofs/DivisorProofs/CardIndependentBoundedSectionsShiftedWeightLe.lean}{GlobalTheory.UpperShift.exists_maxT}
\leanref{Development/AbelianVarietyTheory/Proofs/CuspFormEvaluations/NormQuotientThetaFlatAttainsSupOffCommonZerosStatement.lean}{GlobalTheory.T.normQuotientThetaFlat_attains_sup_off_common_zeros}
\leanref{Development/AbelianVarietyTheory/Proofs/CuspFormEvaluations/NormQuotientThetaFlatAttainsSupOffCommonZeros.lean}{GlobalTheory.QuotientMax.main}
\begin{proof}
Write $Q=Q_{T',f}$, $p=T'+n+13/2$ and
$\Theta=\|\theta_2\|^2_{(n+4)L-F}\|\theta_1\|^2_{2L-\Delta}(\|\eta\|^2_{i(g/2)\Delta})^{2/(ig)}$, so that
$Q=\|f\|^2_{F+T'L}\,\Theta\,S^{-p}$ where $S>0$.

\emph{Invariance and continuity.} Let $P$ be a section of weight $\kappa$, $\gamma\in\Gamma$, $b\in\mathcal B^\Psi$ and
$x\in\mathcal D_X$. Then $\|P\|_\kappa(\gamma xb)=\|P\|_\kappa(x)$. Indeed $|\chi_\kappa(\gamma)|=1$, as $\chi_\kappa$ has finite
order on $\Gamma$; $d_i(\gamma x,\tau)=d_i(x,\tau)$, as $\gamma^\tau$ preserves $h_\tau$; and
$d_i(xb,\tau)=|\det\beta_i|^2d_i(x,\tau)$ for the leading $i\times i$ block $\beta_i$ of $b_\tau$, as in the proof of
Lemma~\ref{is-onto-share-majorant}(2), so that $N_\kappa(xb)=\prod_{\tau,i}|(b_\tau)_{ii}|^{-\kappa_{\tau,i}}N_\kappa(x)$. Hence
$S$ and $Q$ are invariant under $x\mapsto\gamma xb$, which preserves $\mathcal D_X$. The functions $N_\kappa$ are
continuous on $\mathcal D_X$, so $Q$ is continuous on every subset of $\mathcal D_X$ on which $S>0$.

\emph{The theta factor.} The cusp forms $\theta_1,\theta_2$ are bounded, and
$(\|\eta\|^2_{i(g/2)\Delta})^{2/(ig)}\le K^{2/(ig)}e^{-(4/g)\phi}$ by (3) of
Definition~\ref{is-onto-share-counts-model}. So $\Theta\le Be^{-(4/g)\phi}\le B$ on $\mathcal D_X$ for some $B>0$,
as $\phi\geq0$.

Fix $x_1\in\mathcal D_X$ with $\epsilon=Q(x_1)>0$.

\emph{Near the base points.} Let $1\le j\le\delta$, and let $\zeta(z)=\big(s_k(c_{y_j}(z))/w_0(c_{y_j}(z))\big)_k$ near
$0$. As $\zeta(0)=0$ and its derivative at $0$ is invertible, $|\zeta(z)|\geq c'|z|$ near $0$ for some $c'>0$. As
$\|w_0\|_L$ is continuous and positive at $y_j$, $S(c_{y_j}(z))=\|w_0\|_L^2(c_{y_j}(z))\,|\zeta(z)|^2\geq d|z|^2$ near
$0$ for some $d>0$. As $N_{F+T'L}$ is bounded near $y_j$, $\|f\|_{F+T'L}(c_{y_j}(z))\le C|z|^{T'+n+7}$ near $0$. So
near $0$
\[
 Q(c_{y_j}(z))\le C^2Bd^{-p}|z|^{2(T'+n+7)-2p}=C^2Bd^{-p}|z|,
\]
also at $z=0$, where $S$ vanishes. Choose $\rho_j>0$ with $c_{y_j}(z)\in\mathcal D_X$ and $Q(c_{y_j}(z))<\epsilon$ for
$|z|<\rho_j$.

\emph{The set $U$.} Let $U$ be the set of the points $\gamma c_{y_j}(z)b$ with $1\le j\le\delta$, $\gamma\in\Gamma$,
$b\in\mathcal B^\Psi$ and $|z|<\rho_j$. It lies in $\mathcal D_X$, contains the base points, contains $\gamma xb$
with each of its points $x$, and $Q<\epsilon$ on $U$ by the invariance. It is open. Indeed, by the LU
decomposition \cite[Section~3.5]{HJ13}, $(l,b)\mapsto lb$, for $l$ lower unitriangular and $b$ invertible upper
triangular, is a homeomorphism onto the open set of matrices whose leading principal minors are nonzero.
As the components of $y_j$ are invertible, the points $c_{y_j}(z)b$ with $|z|<\rho_j$ and $b\in\mathcal B^\Psi$
form an open set, and so do its images under the homeomorphisms $x\mapsto\gamma x$.

\emph{Away from $U$.} By (5) and (2) of Definition~\ref{is-onto-share-counts-model},
$e^{-2c\phi}\le K\sum_k\|w_k\|_L^2\le KK'S$ on $\mathcal D_X\setminus U$, so $S\geq d'e^{-2c\phi}$ there for some $d'>0$.
As $f$ is bounded, $\|f\|_{F+T'L}\le C_f$, and on $\mathcal D_X\setminus U$
\[
 Q\le C_f^2Be^{-(4/g)\phi}\,d'^{-p}e^{2cp\phi}=C_f^2Bd'^{-p}e^{-\beta\phi},\qquad\beta=4/g-2cp.
\]
Here $\beta>0$, which is the hypothesis $cg(2T'+2n+13)<4$. By Lemma~\ref{is-onto-share-majorant}(1) the form
$v\mapsto\sum_\tau c_\tau m_{x,\tau}(\tau(v))$ is positive definite on $\Lambda\otimes\mathbf R$, in which $\Lambda\ne0$ is a
lattice; so $\lambda_1>0$ on $\mathcal D_X$. Since $\phi\geq c_a\lambda_1^{-2}$ where $\lambda_1<\varepsilon_0$, there is
$\varepsilon'>0$ with $Q<\epsilon$ at the points of $\mathcal D_X\setminus U$ where $\lambda_1<\varepsilon'$.

\emph{The maximum.} Let $K_c$ be the compact set of Lemma~\ref{is-onto-share-counts-reduction} for $\varepsilon'$.
Let $x\in\mathcal D_X$ with $Q(x)\geq\epsilon$. Then $x\notin U$ and $\lambda_1(x)\geq\varepsilon'$, so $\gamma xb\in K_c$
for some $\gamma$ and $b$; moreover $\gamma xb\notin U$, as otherwise $x\in U$, and $Q(\gamma xb)=Q(x)$. Now $S>0$ on
$K_c\setminus U$, so $Q$ is continuous on this compact set and attains its maximum on it at some $x_0$. For
$x=x_1$ this gives $Q(x_0)\geq\epsilon>0$. For every $x\in\mathcal D_X$, either $Q(x)<\epsilon\le Q(x_0)$, or
$Q(x)=Q(\gamma xb)\le Q(x_0)$. Finally $S(x_0)>0$, so some $s_k(x_0)\ne0$.
% Lean: exists_maxT, through mainT; for T' = T the Lean has the same argument as QuotientMax.exists_max
% (with rate_pos). quotT_inv (pointwiseNorm_invariance_holds), quotT_contOn (BigB.pnorm_contOn),
% thetaFac_le and thetaFac_bdd (p4), quotT_near (sumSq_lower with VanishPf.lower_of_det), sat_open
% (ExtFPf.luData_of lu_holds after ortho_holds, then LUB.isOpen_sat), sumSq_far (family_le_coords, p3,
% BigB.S_pos), quotT_far (p1, ModeHalfC.first_minimum_pos with BigB.lam_ne, and e^(beta ph) >= 1 + beta ph),
% reductionTheory_holds (Lemma is-onto-share-counts-reduction) and IsCompact.exists_isMaxOn.
\end{proof}

\begin{lemma}\label{is-onto-share-counts-no-maximum}
Let $X$ consist of the data (1) to (6) of flag data of rank $r$ over $E$, with archimedean data for
$(L,F,T)$ and a flat model for them, and the notation of Definition~\ref{is-onto-share-counts-model}. Let
$T'\geq0$ be an integer and $f$ a section of weight
$F+T'L$. Then $Q_{T',f}$ does not attain its maximum on $\mathcal D_X$ at a point $x_0$ with $Q_{T',f}(x_0)>0$ at
which some $s_k$ does not vanish.
\end{lemma}
\leanref{Development/AbelianVarietyTheory/Proofs/DivisorProofs/CardIndependentBoundedSectionsShiftedWeightLe.lean}{GlobalTheory.UpperShift.no_pos_maxT}
\leanref{Development/AbelianVarietyTheory/Proofs/CuspFormEvaluations/QuotientNoPositiveMaxOffCommonZerosStatement.lean}{GlobalTheory.T.quotient_no_positive_max_off_common_zeros}
\leanref{Development/AbelianVarietyTheory/Proofs/CuspFormEvaluations/QuotientNoPositiveMaxOffCommonZeros.lean}{GlobalTheory.QuotNoMax.no_localMax}
\leanref{Development/AbelianVarietyTheory/Proofs/CuspFormEvaluations/LeviLogSumSqFibreZeroAndNonposAtMaxStatement.lean}{GlobalTheory.T.levi_log_sum_sq_fibre_zero_and_nonpos_at_max}
\begin{proof}
Suppose that $Q=Q_{T',f}$ attains its maximum at such a point $x_0$, with $s_{k_0}(x_0)\ne0$. Since
$Q(x_0)>0$, none of $f,\theta_2,\theta_1,\eta$ vanishes at $x_0$. Put $p=T'+n+13/2$ and
$\psi_\kappa=-2\log N_\kappa$, where $N_\kappa$ is given by the same formula for real $\kappa$; then $\psi_\kappa$ is linear in $\kappa$,
and $\log\|P\|_\kappa^2=\log|P|^2-\psi_\kappa$. The weights $F+T'L$, $(n+4)L-F$, $2L-\Delta$, $\frac2{ig}\cdot i\frac g2\Delta$ and
$-pL$ add up to $-\frac12L$. So, near $x_0$,
\[
 \log Q=\log|f|^2+\log|\theta_2|^2+\log|\theta_1|^2+\tfrac2{ig}\log|\eta|^2-p\log\sum_k|s_k|^2+\tfrac12\psi_L.
\]
Let $c=c_{x_0}$. The functions $f\circ c$, $\theta_2\circ c$, $\theta_1\circ c$, $\eta\circ c$ and $s_k\circ c$ are
holomorphic near $0$, and $\log Q\circ c$ has a local maximum at $0$.

The derivative at $0$ of $z\mapsto\big(s_k(c(z))/s_{k_0}(c(z))\big)_k$ is not surjective, as its $k_0$-th
component is $0$; so its kernel contains some $\xi\ne0$. At $0$ and in the direction $\xi$, the Levi form of
each term of $\log Q\circ c$ but the last vanishes. For $\log|h\circ c|^2$, with $h$ holomorphic and
$h(x_0)\ne0$, this holds as $\log|h\circ c|^2$ is pluriharmonic near $0$. Write
$\log\sum_k|s_k\circ c|^2=\log|s_{k_0}\circ c|^2+\log\sum_k|v_k|^2$ with $v_k=(s_k/s_{k_0})\circ c$. The first term is
pluriharmonic, and the Levi form of the second at $0$ in the direction $\xi$ is
\[
 \frac{\sum_k|\partial_\xi v_k|^2}{\sum_k|v_k|^2}-\frac{\big|\sum_k\bar v_k\,\partial_\xi v_k\big|^2}{\big(\sum_k|v_k|^2\big)^2}=0,
\]
where $\partial_\xi=\sum_j\xi_j\partial_j$, as $\partial_\xi v_k(0)=0$ for every $k$. At a local maximum of a real
function of class $C^2$ its Levi form is $\le0$ in every direction: the restriction to the complex line
$t\mapsto t\xi$ has a local maximum at $0$, so its Laplacian, $4$ times the Levi form at $\xi$, is $\le0$. Hence
the Levi form of $\frac12\psi_L\circ c$ at $0$ in the direction $\xi$ is $\le0$, against (4) of
Definition~\ref{is-onto-share-counts-model}.
% Lean: no_pos_maxT; for T' = T the Lean has the same argument as QuotNoMax.no_pos_max. The logarithm is
% hlog, through BigC.log_pnorm_sq, OW.sum_pnorm_sq and the identities for phla (phla_add_nsmul, phla_sub,
% phla_nsmul); then no_localMax with exists_dir and T.levi_log_sum_sq_fibre_zero_and_nonpos_at_max
% (LeviFibre.main), and posL at z = 0.
\end{proof}

\begin{lemma}\label{is-onto-share-counts-upper}
Let $X$ consist of the data (1) to (6) of flag data of rank $r$ over $E$, and $n=|\Psi|r(r-1)/2$. Let $L,F$ be
weights and $T$ an integer, and let $X$ carry archimedean data for $(L,F,T)$ and a flat model for them, with
$\delta$ base points (Definition~\ref{is-onto-share-counts-model}). Then every family of bounded sections of weight $F+TL$ whose restrictions to $\mathcal D_X$ are
linearly independent has at most $\delta\binom{T+2n+6}{n}$ members, and every such family of weight
$F+(T+n+4)L$ has at most $\delta\binom{T+3n+10}{n}$ members.
\end{lemma}
\leanref{Development/AbelianVarietyTheory/Proofs/CuspFormEvaluations/ExistsMulNe.lean}{GlobalTheory.UpperModel.upper}
\leanref{Development/AbelianVarietyTheory/Proofs/DivisorProofs/CardIndependentBoundedSectionsShiftedWeightLeStatement.lean}{GlobalTheory.T.card_independent_boundedSections_shiftedWeight_le}
\leanref{Development/AbelianVarietyTheory/Proofs/DivisorProofs/CardIndependentBoundedSectionsShiftedWeightLe.lean}{GlobalTheory.UpperShift.upper_shift}
\begin{proof}
Let $T'$ be $T$ or $T+n+4$, and let $f_1,\dots,f_m$ be bounded sections of weight $F+T'L$ with independent
restrictions to $\mathcal D_X$. Suppose $m>\delta\binom{T'+2n+6}{n}$. Lemma~\ref{is-onto-share-counts-jets}(2), with
$e=T'+n+6$, gives a nonzero $a\in\mathbf C^m$ such that $f=\sum_qa_qf_q$ satisfies $f(c_{y_j}(z))=O(|z|^{T'+n+7})$
for every $j$. It is a bounded section of weight $F+T'L$, and it does not vanish identically on $\mathcal D_X$, by
the independence.

Some $s_k$ does not vanish identically on $\mathcal D_X$. Otherwise every point of $\mathcal D_X$, which is
nonempty as $\eta$ does not vanish identically on it, is a common zero of the $s_k$; so there is a base point
$y_j$. Then $c_{y_j}(z)\in\mathcal D_X$ near $0$, so $z\mapsto(s_k/w_0)(c_{y_j}(z))$ vanishes near $0$, and its derivative
at $0$ is $0$, which is not invertible as $n\geq2$.

The domain $\mathcal D_X$ is open and connected (Lemma~\ref{is-onto-share-sections}(2)), so by the identity
theorem the product of two holomorphic functions on $\mathcal D_X$ that do not vanish identically does not
vanish identically. Hence $f\theta_2\theta_1\eta s_k$ is nonzero at some $x\in\mathcal D_X$. There $S(x)>0$ and
$Q_{T',f}(x)>0$, as $N_\kappa>0$ on $\mathcal D_X$.

As $Tcg\le1$ and $T\geq11n^2\geq2n+11$, $cg(2T'+2n+13)\le(2T+4n+21)/T<4$. So
Lemma~\ref{is-onto-share-counts-maximum} gives a point $x_0$ at which some $s_k$ does not vanish and at which
$Q_{T',f}$ attains a positive maximum, and Lemma~\ref{is-onto-share-counts-no-maximum} excludes it. Hence
$m\le\delta\binom{T'+2n+6}{n}$, which is $\delta\binom{T+2n+6}{n}$ for $T'=T$ and $\delta\binom{T+3n+10}{n}$ for
$T'=T+n+4$.
% Lean: UpperModel.upper (T' = T, with QuotientMax.main and QuotNoMax.no_pos_max) and
% UpperShift.upper_shift (T' = T + n + 4, with mainT, no_pos_maxT and rate_shift), both through
% exists_comb, isBddSec_sum, exists_s_ne, exists_mul_ne and pnorm_pos.
\end{proof}

\subsection{Levi forms and sections at a base point}\label{subsection-is-base-point-sections}

\begin{definition}\label{is-onto-share-counts-section-levi}
Let $X$ consist of the data (1) to (6) of flag data of rank $r$ over $E$
(Definition~\ref{is-onto-share-flag-data}), put $n=|\Psi|r(r-1)/2$, and use the charts $c_y$, the functions
$N_\kappa$ and the Levi forms of Definition~\ref{is-onto-share-counts-model}. For a weight $\kappa$ put
$\varphi_\kappa=-2\log N_\kappa$ on $\mathcal D_X$, so that
\[
 \varphi_\kappa(x)=\sum_{\tau\in\Psi}\sum_{i=1}^r\kappa_{\tau,i}\big(\log|d_i(x,\tau)|-\log|d_{i-1}(x,\tau)|\big).
\]
For $y\in M_r(\mathbf C)^\Psi$ let $\Omega(y)$ be the open set of the $z\in\mathbf C^n$ with $c_y(z)\in\mathcal D_X$. On
$\Omega(y)$ the minors $d_i(c_y(z),\tau)$ are nonzero polynomials in the real and imaginary parts of $z$. Write
$\mathcal L_g(z;\xi)$ for the Levi form at $z$ of a smooth real function $g$ on $\Omega(y)$, in the direction
$\xi\in\mathbf C^n$, and $\mathcal L^y_\kappa(z;\xi)=\mathcal L_{\varphi_\kappa\circ c_y}(z;\xi)$. For an orthonormal point
$x_0$ and $z\in\Omega(x_0)$ put
\[
 \rho_{x_0}(z)=-\sum_{\tau\in\Psi}\log|d_{p_0}(c_{x_0}(z),\tau)|,\qquad
 D_{x_0}(z)=\prod_{\tau\in\Psi}\prod_{i=0}^{r-1}|d_i(c_{x_0}(z),\tau)|^{-2}.
\]

Let $L$ and $\lambda$ be weights. Real numbers $\nu,m,\eta$ are \emph{curvature constants} for $\lambda$ and $L$ if
$4|\Psi|r^2<\nu<m$, $\eta>0$, and for every orthonormal point $x_0$, every $z\in\Omega(x_0)$ and every
$\xi\in\mathbf C^n$ the number
$\mathcal C=\mathcal L^{x_0}_\lambda(z;\xi)-\mathcal L_{\log D_{x_0}}(z;\xi)-n\,\mathcal L^{x_0}_L(z;\xi)$ satisfies
\[
 m\,\mathcal L_{\rho_{x_0}}(z;\xi)\le\mathcal C\quad\text{and}\quad2\eta\,\mathcal L^{x_0}_L(z;\xi)\le\mathcal C.
\]
\end{definition}
\leanref{ReciprocityTheory/ReductiveGroups/ArithmeticGeometry/ArakelovVolumes/UpperBounds/FundSetMeasStatement.lean}{ReciprocityTheory.ChartRestriction.G.AR.Up.phla}
\leanref{ReciprocityTheory/ReductiveGroups/ArithmeticGeometry/ArakelovVolumes/UpperBounds/IsOrth.lean}{ReciprocityTheory.ChartRestriction.G.AR.Up.OmF}
\leanref{ReciprocityTheory/ReductiveGroups/ArithmeticGeometry/ArakelovVolumes/UpperBounds/IsOrth.lean}{ReciprocityTheory.ChartRestriction.G.AR.Up.rhF}
\leanref{ReciprocityTheory/ReductiveGroups/ArithmeticGeometry/ArakelovVolumes/UpperBounds/IsOrth.lean}{ReciprocityTheory.ChartRestriction.G.AR.Up.densF}
\leanref{ReciprocityTheory/L2DbarEstimates/IsLocalHomeomorph.lean}{ReciprocityTheory.L2DbarEstimates.leviV}
\leanref{ReciprocityTheory/ReductiveGroups/ArithmeticGeometry/ArakelovVolumes/UpperBounds/BasePsiLeviLowerBound.lean}{ReciprocityTheory.ChartRestriction.G.AR.Up.leviM}
\leanref{ReciprocityTheory/ReductiveGroups/ArithmeticGeometry/ArakelovVolumes/UpperBounds/WeightFunctionOnFrames.lean}{ReciprocityTheory.ChartRestriction.G.AR.Up.leviO}
\leanref{ReciprocityTheory/ReductiveGroups/ArithmeticGeometry/ArakelovVolumes/UpperBounds/WeightFunctionOnFrames.lean}{ReciprocityTheory.ChartRestriction.G.AR.Up.leviRh}
\leanref{ReciprocityTheory/ReductiveGroups/ArithmeticGeometry/ArakelovVolumes/UpperBounds/WeightFunctionOnFrames.lean}{ReciprocityTheory.ChartRestriction.G.AR.Up.leviDens}

\begin{lemma}\label{is-onto-share-counts-section-diagonal}
Let $X$ consist of the data (1) to (6) of flag data of rank $r$ over $E$, use the notation of
Definition~\ref{is-onto-share-counts-section-levi} and the signs $\operatorname{sgn}(a)$ of
Definition~\ref{is-onto-share-flag-data}, and let $x$ be an orthonormal point and $\xi\in\mathbf C^n$, with
coordinates $\xi_{\tau,a,b}$.
\begin{enumerate}
\item For $\tau\in\Psi$ and $0\le i\le r$, the Levi form at $0$ of $z\mapsto\log|d_i(c_x(z),\tau)|$ in the
  direction $\xi$ is the sum of $\operatorname{sgn}(a)\operatorname{sgn}(b)|\xi_{\tau,a,b}|^2$ over the pairs
  $1\le b\le i<a\le r$.
\item For every weight $\lambda$,
  \[
   \mathcal L^x_\lambda(0;\xi)=\sum_{\tau\in\Psi}\sum_{1\le b<a\le r}\operatorname{sgn}(a)\operatorname{sgn}(b)
   (\lambda_{\tau,b}-\lambda_{\tau,a})|\xi_{\tau,a,b}|^2.
  \]
\end{enumerate}
\end{lemma}
\leanref{Development/AbelianVarietyTheory/Proofs/TorusHaarMeasure/LeviFormWeightMetricOrthonormalDiagonalStatement.lean}{GlobalTheory.T.leviForm_weightMetric_orthonormal_diagonal}
\leanref{Development/AbelianVarietyTheory/Proofs/TorusHaarMeasure/LeviFormWeightMetricOrthonormalDiagonal.lean}{GlobalTheory.LeviZero.leviZero}
\leanref{Development/AbelianVarietyTheory/Proofs/TorusHaarMeasure/LeviFormWeightMetricOrthonormalDiagonal.lean}{GlobalTheory.LeviZero.levi_phla_zero}
\leanref{Development/AbelianVarietyTheory/Proofs/TorusHaarMeasure/LeviFormWeightMetricOrthonormalDiagonal.lean}{GlobalTheory.LeviZero.leviV_Lm}
\begin{proof}
(1) is a self-contained computation. The Gram matrix of the columns of $c_x(z)_\tau=x_\tau\ell_\tau(z)$ is
$\ell_\tau(z)^tJ\overline{\ell_\tau(z)}$, as that of $x_\tau$ is $J$, and (1) comes from the terms of degree two in
$z$ and $\bar z$ of its leading minors.

(2) The point $c_x(0)=x$ lies in $\mathcal D_X$ (Lemma~\ref{is-onto-share-majorant}(3)), so $0\in\Omega(x)$, and
on $\Omega(x)$
\[
 \varphi_\lambda(c_x(z))=\sum_{\tau\in\Psi}\sum_{i=1}^r\lambda_{\tau,i}\big(\log|d_i(c_x(z),\tau)|-\log|d_{i-1}(c_x(z),\tau)|\big).
\]
The Levi form is linear on smooth functions. Write $[P]$ for $1$ if the condition $P$ holds and $0$
otherwise. By (1), for $b<a$ the coefficient of $\operatorname{sgn}(a)\operatorname{sgn}(b)|\xi_{\tau,a,b}|^2$ in
$\mathcal L^x_\lambda(0;\xi)$ is
\[
 \sum_{i=1}^r\lambda_{\tau,i}\big([b\le i<a]-[b\le i-1<a]\big)=\sum_{i=1}^r\lambda_{\tau,i}\big([b\le i<a]-[b<i\le a]\big)
 =\lambda_{\tau,b}-\lambda_{\tau,a}.
\]
% Lean: (1) is LeviZero.leviV_Lm (the condition is b < i <= a for the Lean's indices a, b from 0), a computation
% with jets of the minors (HasJet.det, det_jet_sq, gram_cpt_apply). (2) is levi_phla_zero, through
% IW.phla_cpt_eventually, leviV_sum', leviV_add' and leviV_const_mul', with the coefficient comparison hi.
\end{proof}

\begin{lemma}\label{is-onto-share-counts-section-transfer-lu}
Let $X$ consist of the data (1) to (6) of flag data of rank $r$ over $E$, use the charts $c_y$ and the sets
$\Omega(y)$ of Definitions~\ref{is-onto-share-counts-model} and~\ref{is-onto-share-counts-section-levi}, and let
$x_0$ be an orthonormal point. For $x\in M_r(\mathbf C)^\Psi$ let $\Delta(x)$ be the product over $\tau\in\Psi$ and
$1\le k\le r$ of the leading principal $k\times k$ minors of $x_{0,\tau}^{-1}x_\tau$.
\begin{enumerate}
\item $\Delta$ is a polynomial in the entries of $x$ with $\Delta(x_0)=1$, and every neighbourhood of a point of
  $\mathcal D_X$ contains points $x$ with $\Delta(x)\ne0$.
\item On the open set of the $x\in\mathcal D_X$ with $\Delta(x)\ne0$ there are holomorphic maps $x\mapsto z(x)$ to
  $\Omega(x_0)$ and $x\mapsto b(x)$ to $\mathcal B^\Psi$ with $x=c_{x_0}(z(x))\,b(x)$.
\item For $z\in\Omega(x_0)$ and $b\in\mathcal B^\Psi$, the point $x=c_{x_0}(z)b$ satisfies $\Delta(x)\ne0$,
  $z(x)=z$ and $b(x)=b$.
\end{enumerate}
\end{lemma}
\leanref{ReciprocityTheory/ReductiveGroups/ArithmeticGeometry/ArakelovVolumes/UpperBounds/FundSetMeasStatement.lean}{ReciprocityTheory.ChartRestriction.G.AR.Up.LUStatement}
\leanref{ReciprocityTheory/ReductiveGroups/ArithmeticGeometry/ArakelovVolumes/UpperBounds/LUStatementPf.lean}{ReciprocityTheory.ChartRestriction.G.AR.Up.lu_holds}
\leanref{ReciprocityTheory/ReductiveGroups/ArithmeticGeometry/ArakelovVolumes/UpperBounds/LUStatementPfA.lean}{ReciprocityTheory.ChartRestriction.G.AR.Up.LUPf.Lm_unique}
\begin{proof}
\emph{Triangular factorization.} Let $M$ be an $r\times r$ matrix whose leading principal minors $\delta_1,\dots,\delta_r$
are nonzero, and put $\delta_0=1$. For $1\le i\le r$ let $M_{i-1}$ be the leading $(i-1)\times(i-1)$ block of $M$,
which is invertible, and let $w_i$ be the row vector with $w_iM_{i-1}=(M_{i1},\dots,M_{i,i-1})$. Let $L$ be the lower
unitriangular matrix whose $i$-th row is $(-w_i,1,0,\dots,0)$. Then $(LM)_{ij}=M_{ij}-\sum_{k<i}w_{ik}M_{kj}=0$ for
$j<i$, so $U=LM$ is upper triangular. If $L'$ is lower unitriangular and $L'M$ is upper triangular, the $i$-th row
of $L'$ satisfies the same linear system as that of $L$, so $L'=L$. By Cramer's rule the entries of $L$ and $U$ are
polynomials in the entries of $M$ divided by products of the $\delta_k$. The leading principal minors of a product
$NB$ of a lower unitriangular $N$ and an upper triangular $B$ are the products of the leading diagonal entries of
$B$; so the diagonal entries of $U$ are $\delta_i/\delta_{i-1}\ne0$. Hence $M=L^{-1}U$ with $L^{-1}$ lower unitriangular
and $U\in\mathcal B$, and this factorization is unique.

(1) As $\det x_{0,\tau}\ne0$ (Lemma~\ref{is-onto-share-majorant}(4)), $\Delta$ is a polynomial, and
$\Delta(x_0)=1$. Let $x\in\mathcal D_X$ and $\varepsilon>0$. The function $t\mapsto\Delta(x+t(x_0-x))$ is a
polynomial with the value $1$ at $t=1$, so it has finitely many zeros, and some real $t$ with $0<t|x_0-x|<\varepsilon$
is not one of them.

(2) Let $x\in\mathcal D_X$ with $\Delta(x)\ne0$, and apply the factorization to $M=x_{0,\tau}^{-1}x_\tau$ for each
$\tau$. Let $z(x)$ be the point whose coordinates $z_{\tau,a,b}$ are the entries of $L^{-1}$ below the diagonal, so
that $L^{-1}=\ell_\tau(z(x))$, and let $b(x)_\tau=U$. Then $x_\tau=x_{0,\tau}\ell_\tau(z(x))b(x)_\tau$, that is,
$x=c_{x_0}(z(x))b(x)$. Both maps are holomorphic on the open set where $\Delta\ne0$, as their entries are quotients
of polynomials by products of the minors. By Lemma~\ref{is-onto-share-majorant}(2), $c_{x_0}(z(x))=xb(x)^{-1}$ lies
in $\mathcal D_X$, so $z(x)\in\Omega(x_0)$.

(3) Here $x_{0,\tau}^{-1}x_\tau=\ell_\tau(z)b_\tau$, whose leading principal minors are products of diagonal entries of
$b_\tau$, so they are nonzero and $\Delta(x)\ne0$. By the uniqueness of the factorization, $z(x)=z$ and $b(x)=b$.
% Lean: lu_holds: DeltaF (the product includes k = 0, whose minor is 1), zF, bF and nF = (Lm M)^-1 with
% M = Mx x0 x; Lm (rows -wv M i, wv = row * (lb M i)^-1), Lm_mul_up, Lm_unique (row_sys), lminor_nb for the
% minors of N B; Delta_diff, zF_diffOn and bF_diffOn (diffOn_det, diffOn_inv); split_eq, zF_mem, bF_mem and
% uniq_all; orthogonalFrame_delta_eq_one and dense_delta (linePoly). Size: about 570 lines of Lean.
\end{proof}

\begin{lemma}\label{is-onto-share-counts-section-transfer-jacobian}
Let $X$ consist of the data (1) to (6) of flag data of rank $r$ over $E$, use the charts of
Definition~\ref{is-onto-share-counts-model} and the notation $b^\lambda$ of
Definition~\ref{is-onto-share-counts-section-gauge}, and let $k$ be the weight with $k_{\tau,i}=2i-r-1$. Let
$x_0,x_1\in M_r(\mathbf C)^\Psi$, $\gamma\in GL_r(E)$, $U\subset\mathbf C^n$ open, $\Phi\colon U\to\mathbf C^n$
holomorphic and $\beta\colon U\to M_r(\mathbf C)^\Psi$ a map such that, for every $w\in U$, $c_{x_0}(w)\in\mathcal D_X$,
$\beta(w)\in\mathcal B^\Psi$ and $c_{x_0}(w)=\gamma\,c_{x_1}(\Phi(w))\,\beta(w)$. Then
$\det D\Phi(w)=\beta(w)^k$ for every $w\in U$.
\end{lemma}
\leanref{ReciprocityTheory/ReductiveGroups/ArithmeticGeometry/ArakelovVolumes/ArakelovInstances/PnormInvStatement.lean}{ReciprocityTheory.ChartRestriction.G.AR.Inst.KapGGStatement}
\leanref{Development/AbelianVarietyTheory/Proofs/TorusHaarMeasure/MeasureStatements.lean}{GlobalTheory.FrameMeasureClosed.kapGG_holds}
\leanref{Development/AbelianVarietyTheory/Proofs/TorusHaarMeasure/DifferentiableAtAdjugate.lean}{GlobalTheory.UnitriJac.kapGG_of_det}
\leanref{Development/AbelianVarietyTheory/Proofs/TorusHaarMeasure/DetOfBlockTriangularInjective.lean}{GlobalTheory.det_conj_strictLower}
\leanref{Development/AbelianVarietyTheory/Proofs/TorusHaarMeasure/DetMulStrictLowerEqOne.lean}{GlobalTheory.det_mul_strictLower_eq_one}
\begin{proof}
This is a self-contained computation. Identify $\mathbf C^n$ with the families $(u_\tau)_{\tau\in\Psi}$ of strictly
lower triangular $r\times r$ matrices, so that $\ell_\tau(w)=1+w_\tau$. Fix $w\in U$ and $\tau$, write
$\ell=\ell_\tau(w)$, $\ell'=\ell_\tau(\Phi(w))$, $B=\beta(w)_\tau$, and let $v\in\mathbf C^n$. The map $\beta_\tau$ is
holomorphic, as $\beta_\tau=(\gamma^\tau x_{1,\tau}\ell_\tau\circ\Phi)^{-1}x_{0,\tau}\ell_\tau$ on $U$. Differentiating
$x_{0,\tau}\ell_\tau(w)=\gamma^\tau x_{1,\tau}\ell_\tau(\Phi(w))\beta(w)_\tau$ in the direction $v$, and using the relation
itself to eliminate $\gamma^\tau x_{1,\tau}$, gives
\[
 \ell'^{-1}\,(D\Phi(w)v)_\tau=B\,\ell^{-1}v_\tau\,B^{-1}-B'B^{-1},
\]
where $B'$ is the derivative of $\beta_\tau$ in the direction $v$. The matrix $B'B^{-1}$ is upper triangular, and the
left side is strictly lower triangular. So, with $\mathrm{sl}$ the strictly lower triangular part,
$\ell'^{-1}(D\Phi(w)v)_\tau=\mathrm{sl}(B\ell^{-1}v_\tau B^{-1})$. On the space of strictly lower triangular matrices,
multiplication on the left by a lower unitriangular matrix has determinant $1$, and $u\mapsto\mathrm{sl}(BuB^{-1})$
has determinant $\prod_iB_{ii}^{2i-r-1}$, as it is triangular for a suitable order of the entries with diagonal
entries $B_{aa}B_{bb}^{-1}$ at the entry $(a,b)$. Taking determinants over all $\tau$ gives
$\det D\Phi(w)=\prod_{\tau,i}(\beta(w)_\tau)_{ii}^{k_{\tau,i}}=\beta(w)^k$.
% Lean: FrameMeasureClosed.kapGG_holds = UnitriJac.kapGG_of_det det_conj_strictLower det_mul_strictLower_eq_one;
% fderiv_rel (deriv_rel) is the displayed relation, read through leftM and adM on the strictly lower entries
% (Low r), and the determinant is taken blockwise in tau (Matrix.det_blockDiagonal). The weight k is kapW X.
% Size: about 360 lines of Lean.
\end{proof}

\begin{lemma}\label{is-onto-share-counts-section-transfer}
Let $X$ consist of the data (1) to (6) of flag data of rank $r$ over $E$, use the notation of
Definition~\ref{is-onto-share-counts-section-levi}, and let $y\in M_r(\mathbf C)^\Psi$ and $z\in\Omega(y)$. Then
there are an orthonormal point $x_1$ and an injective $\mathbf C$-linear map $D\colon\mathbf C^n\to\mathbf C^n$ such
that $\mathcal L^y_\lambda(z;\xi)=\mathcal L^{x_1}_\lambda(0;D\xi)$ for every weight $\lambda$ and every
$\xi\in\mathbf C^n$.
\end{lemma}
\leanref{Development/AbelianVarietyTheory/Proofs/TorusHaarMeasure/BorelDiagNeZero.lean}{GlobalTheory.LeviTransfer.levi_transfer}
\begin{proof}
\emph{The change of charts.} Put $x=c_y(z)\in\mathcal D_X$. By Lemma~\ref{is-onto-share-majorant}(3), $x=x_1b_0$ for
an orthonormal point $x_1$ and some $b_0\in\mathcal B^\Psi$. Let $\Delta$, $z(\cdot)$ and $b(\cdot)$ be as in
Lemma~\ref{is-onto-share-counts-section-transfer-lu} for $x_1$. As $x=c_{x_1}(0)b_0$ and $0\in\Omega(x_1)$, part (3)
of that lemma gives $\Delta(x)\ne0$, $z(x)=0$ and $b(x)=b_0$. Let $U$ be the set of $w\in\Omega(y)$ with
$\Delta(c_y(w))\ne0$; it is open, as $c_y$ is polynomial, and it contains $z$. On $U$ put $\Phi(w)=z(c_y(w))$ and
$\beta(w)=b(c_y(w))$. By part (2), $\Phi$ and $\beta$ are holomorphic, $c_y(w)=c_{x_1}(\Phi(w))\beta(w)$ and
$\beta(w)\in\mathcal B^\Psi$ for $w\in U$, and $\Phi(z)=0$. Let $D=D\Phi(z)$. By
Lemma~\ref{is-onto-share-counts-section-transfer-jacobian}, for $x_0=y$ and $\gamma=1$, $\det D=\beta(z)^k\ne0$,
where $k_{\tau,i}=2i-r-1$; so $D$ is injective.

\emph{The potentials.} Let $x'\in\mathcal D_X$ and $b\in\mathcal B^\Psi$. By the proof of
Lemma~\ref{is-onto-share-majorant}(2), $d_i(x'b,\tau)=|\prod_{a\le i}(b_\tau)_{aa}|^2d_i(x',\tau)$, so
$N_\lambda(x'b)=|b^\lambda|^{-1}N_\lambda(x')$ and $\varphi_\lambda(x'b)=\varphi_\lambda(x')+\log|b^\lambda|^2$. Hence, for
$w\in U$,
\[
 \varphi_\lambda(c_y(w))=\varphi_\lambda(c_{x_1}(\Phi(w)))+\log|\beta(w)^\lambda|^2 .
\]
The function $w\mapsto\beta(w)^\lambda$ is holomorphic and nowhere zero on $U$, so near $z$ the last term is
$g+\bar g$ for a holomorphic logarithm $g$ of it, and its Levi form at $z$ vanishes. The function
$\varphi_\lambda\circ c_{x_1}$ is smooth near $\Phi(z)=0$, as the minors $d_i(c_{x_1}(\cdot),\tau)$ are nonzero
polynomials in the real and imaginary parts of the coordinates there. For a function $g'$ of class $C^2$ near
$\Phi(z)$ and the holomorphic map $\Phi$, the chain rule gives
$\partial_j\bar\partial_k(g'\circ\Phi)(z)=\sum_{a,b}(\partial_a\bar\partial_bg')(\Phi(z))\,\partial_j\Phi_a(z)\,
\overline{\partial_k\Phi_b(z)}$, that is, $\mathcal L_{g'\circ\Phi}(z;\xi)=\mathcal L_{g'}(\Phi(z);D\xi)$. With
$g'=\varphi_\lambda\circ c_{x_1}$ and the linearity of the Levi form,
$\mathcal L^y_\lambda(z;\xi)=\mathcal L^{x_1}_\lambda(0;D\xi)$ for every weight $\lambda$ and every $\xi$.
% Lean: levi_transfer: ortho_holds gives c_y(z) = x1 b; lu_holds at x1 gives zOf and bOf, and
% Phi = zOf o c_y, b = bOf o c_y on U = {w | cpt y w in dom, DeltaF (cpt y w) != 0} (ExF.isOpen_OmF);
% metric_rmul (from NfacRmulStatement, proved by NormFactor.nfac_transform) gives
% phla(c_y(w)) = phla(c_x1(Phi w)) + log |bchar la (b w)|^2; the smoothness is IW.phla_cpt_cdAt and
% HE.contDiffAt_of_punct (holomorphic functions are smooth); leviV_congr, leviV_ofReal_add,
% HE.levi_logNormSq_zero_at and RGF.gc_leviV_comp (Levi forms under holomorphic maps) give the identity; and
% KapGGStatement with gamma = 1 (FrameMeasureClosed.kapGG_holds) and borel_char_ne_zero give
% det D != 0. Size: about 2,000 lines of Lean below levi_transfer, kapGG_holds and nfac_transform (text count
% with tactic names and local names filtered out, stopping at statements the book cites); the earlier count of
% 9,800 lines followed local names such as b and hy to unrelated declarations.
\end{proof}

\begin{lemma}\label{is-onto-share-counts-section-curvature}
Let $X$ consist of the data (1) to (6) of flag data of rank $r$ over $E$, and use the notation of
Definition~\ref{is-onto-share-counts-section-levi}. Let $\lambda$ be a weight and $m$ a real number such that, for
every $\tau\in\Psi$, $\lambda_{\tau,j}\le\lambda_{\tau,i}$ whenever $i\le j\le p_0$ or $p_0<i\le j$, and
$m\le\lambda_{\tau,j}-\lambda_{\tau,i}$ whenever $i\le p_0<j$. Then
$m\,\mathcal L_{\rho_{x_0}}(z;\xi)\le\mathcal L^{x_0}_\lambda(z;\xi)$ for every orthonormal point $x_0$, every
$z\in\Omega(x_0)$ and every $\xi\in\mathbf C^n$.
\end{lemma}
\leanref{Development/AbelianVarietyTheory/Proofs/TorusHaarMeasure/CurvAndMarginClosed.lean}{GlobalTheory.CurvMarginClosed.curv_and_margin_closed}
\leanref{Development/AbelianVarietyTheory/Proofs/TorusHaarMeasure/PhlaBlockWtCpt.lean}{GlobalTheory.CurvMargin.curv_and_margin}
\leanref{Development/AbelianVarietyTheory/Proofs/TorusHaarMeasure/PhlaBlockWtCpt.lean}{GlobalTheory.CurvMargin.phla_blockWt_cpt}
\leanref{Development/AbelianVarietyTheory/Proofs/TorusHaarMeasure/PhlaBlockWtCpt.lean}{GlobalTheory.CurvMargin.leviRh_eq}
\begin{proof}
Let $\beta$ be the weight with $\beta_{\tau,i}=0$ for $i\le p_0$ and $\beta_{\tau,i}=1$ for $i>p_0$.

\emph{The function $\rho_{x_0}$.} Let $z\in\Omega(x_0)$ and $x=c_{x_0}(z)$. The Gram matrix of the columns of
$x_\tau=x_{0,\tau}\ell_\tau(z)$ is $\ell_\tau(z)^tJ\overline{\ell_\tau(z)}$, as that of $x_{0,\tau}$ is $J$. As
$\ell_\tau(z)$ is unitriangular, $|d_r(x,\tau)|=|\det J|=1$. So the sum defining $\varphi_\beta(x)$ telescopes to
$\sum_\tau\big(\log|d_r(x,\tau)|-\log|d_{p_0}(x,\tau)|\big)=\rho_{x_0}(z)$. Hence $\rho_{x_0}=\varphi_\beta\circ c_{x_0}$
on the open set $\Omega(x_0)$, and $\mathcal L_{\rho_{x_0}}(z;\xi)=\mathcal L^{x_0}_\beta(z;\xi)$.

\emph{Reduction to the centre.} Lemma~\ref{is-onto-share-counts-section-transfer}, for $y=x_0$, gives an
orthonormal point $x_1$ and a linear map $D$ with $\mathcal L^{x_0}_\kappa(z;\xi)=\mathcal L^{x_1}_\kappa(0;D\xi)$ for
$\kappa=\lambda$ and $\kappa=\beta$. Put $\zeta=D\xi$. By Lemma~\ref{is-onto-share-counts-section-diagonal}(2) for
$\lambda$ and for $\beta$, it suffices to show that for every $\tau$ and all $b<a$
\[
 m\operatorname{sgn}(a)\operatorname{sgn}(b)(\beta_{\tau,b}-\beta_{\tau,a})\le
 \operatorname{sgn}(a)\operatorname{sgn}(b)(\lambda_{\tau,b}-\lambda_{\tau,a}),
\]
and then to multiply by $|\zeta_{\tau,a,b}|^2$ and sum.

\emph{The coefficients.} If $a\le p_0$ or $b>p_0$, then $\operatorname{sgn}(a)\operatorname{sgn}(b)=1$ and
$\beta_{\tau,b}=\beta_{\tau,a}$, and the claim is $0\le\lambda_{\tau,b}-\lambda_{\tau,a}$, which holds by hypothesis
as $b<a$ lie on the same side of $p_0$. If $b\le p_0<a$, then $\operatorname{sgn}(a)\operatorname{sgn}(b)=-1$ and
$\beta_{\tau,b}-\beta_{\tau,a}=-1$, and the claim is $m\le\lambda_{\tau,a}-\lambda_{\tau,b}$, the last hypothesis.
% Lean: the second part of curv_and_margin_closed is that of CurvMargin.curv_and_margin, applied to
% T.leviForm_weightMetric_orthonormal_diagonal and levi_transfer. beta is blockWt; leviRh_eq and
% phla_blockWt_cpt (IW.phla_eq, tele_tail, IW.Lm_last_cpt) give the first step, and coeff_cases the three
% cases. The first part of curv_and_margin_closed, CurvGGStatement, is not used here.
\end{proof}

\begin{lemma}\label{is-onto-share-counts-section-density}
Let $X$ consist of the data (1) to (6) of flag data of rank $r$ over $E$, use the notation of
Definition~\ref{is-onto-share-counts-section-levi}, and let $k$ be the weight with $k_{\tau,i}=2i-r-1$. Let $x_0$ be
an orthonormal point, $z\in\Omega(x_0)$, $\lambda_1,\lambda_2$ weights and $c\geq0$ an integer. Then for every
$\xi\in\mathbf C^n$
\[
 \mathcal L^{x_0}_{\lambda_1-k-c\lambda_2}(z;\xi)=\mathcal L^{x_0}_{\lambda_1}(z;\xi)-\mathcal L_{\log D_{x_0}}(z;\xi)
 -c\,\mathcal L^{x_0}_{\lambda_2}(z;\xi).
\]
\end{lemma}
\leanref{ReciprocityTheory/ReductiveGroups/ArithmeticGeometry/ArakelovVolumes/UpperBounds/InterpStatementPfH.lean}{ReciprocityTheory.ChartRestriction.G.AR.Up.IW.leviM_comb}
\leanref{ReciprocityTheory/ReductiveGroups/ArithmeticGeometry/ArakelovVolumes/UpperBounds/InterpStatementPfH.lean}{ReciprocityTheory.ChartRestriction.G.AR.Up.IW.phla_comb}
\leanref{ReciprocityTheory/ReductiveGroups/ArithmeticGeometry/ArakelovVolumes/UpperBounds/InterpStatementPfH.lean}{ReciprocityTheory.ChartRestriction.G.AR.Up.IW.phla_kapW_cpt}
\begin{proof}
By its formula in Definition~\ref{is-onto-share-counts-section-levi}, $\varphi_\kappa$ is linear in $\kappa$, so
$\varphi_{\lambda_1-k-c\lambda_2}=\varphi_{\lambda_1}-\varphi_k-c\,\varphi_{\lambda_2}$ on $\mathcal D_X$.

Let $w\in\Omega(x_0)$, $x=c_{x_0}(w)$, $\tau\in\Psi$, and write $L_i=\log|d_i(x,\tau)|$. Then $L_0=0$, and $L_r=0$ as
in the proof of Lemma~\ref{is-onto-share-counts-section-curvature}. As $k_{\tau,i+1}-k_{\tau,i}=2$, summation by parts
gives
\[
 \sum_{i=1}^rk_{\tau,i}(L_i-L_{i-1})=k_{\tau,r}L_r-k_{\tau,1}L_0-\sum_{i=1}^{r-1}(k_{\tau,i+1}-k_{\tau,i})L_i
 =-2\sum_{i=0}^{r-1}L_i.
\]
Summing over $\tau$, $\varphi_k(c_{x_0}(w))=\log D_{x_0}(w)$. Hence
$\varphi_{\lambda_1-k-c\lambda_2}\circ c_{x_0}=\varphi_{\lambda_1}\circ c_{x_0}-\log D_{x_0}-c\,\varphi_{\lambda_2}\circ c_{x_0}$
on the open set $\Omega(x_0)$. The three functions on the right are smooth there, and the Levi form is linear.
% Lean: leviM_comb, through phla_comb (with IW.phla_eq), phla_kapW_cpt (log_densF_eq, the summation by parts
% sbp, Lm_last_cpt and Lm_zero) and leviV_sub_sub_const_mul_of_eventuallyEq, with phla_cpt_cdAt and
% logDens_cdAt for the smoothness.
\end{proof}

\begin{definition}\label{is-onto-share-counts-section-gauge}
Let $X$ consist of the data (1) to (6) of flag data of rank $r$ over $E$, use the notation of
Definition~\ref{is-onto-share-counts-section-levi}, and let $\lambda$ be a weight. For $b\in\mathcal B^\Psi$ put
$b^\lambda=\prod_{\tau,i}(b_\tau)_{ii}^{\lambda_{\tau,i}}$.
\begin{enumerate}
\item A \emph{gauge} of weight $\lambda$ is a measurable function $W$ on $M_r(\mathbf C)^\Psi$ with values in
  $[0,\infty]$ such that $W(xb)=|b^\lambda|^{-2}W(x)$ and $W(\gamma x)=W(x)$ for $x\in\mathcal D_X$,
  $b\in\mathcal B^\Psi$ and $\gamma\in\Gamma$. It is \emph{bounded below} if $W\geq cN_\lambda^2$ on $\mathcal D_X$
  for some real $c>0$.
\item A function $P$ on $M_r(\mathbf C)^\Psi$ is \emph{$\mathcal B$-equivariant} of weight $\lambda$ if it is
  holomorphic on $\mathcal D_X$ and $P(xb)=b^\lambda P(x)$ for $x\in\mathcal D_X$ and $b\in\mathcal B^\Psi$.
\item Let $x_0$ be an orthonormal point, $W$ a gauge of weight $\lambda$ and $\nu$ a real number. A function $f$ on
  $\Omega(x_0)$ is \emph{square integrable} for $W$ and $\nu$ if
  \[
   \int_{\Omega(x_0)}|f(z)|^2e^{\nu\rho_{x_0}(z)}\,W(c_{x_0}(z))\,D_{x_0}(z)\,dz<\infty,
  \]
  for the Lebesgue measure $dz$ on $\mathbf C^n$.
\item The \emph{Poincaré series} of a function $P$ on $M_r(\mathbf C)^\Psi$ is the function $\mathcal P_\lambda P$
  whose value at $x$ is $\sum_{\gamma\in\Gamma}\chi_\lambda(\gamma)^{-1}P(\gamma x)$ if this family is summable, and
  $0$ otherwise.
\end{enumerate}
\end{definition}
\leanref{ReciprocityTheory/ReductiveGroups/ArithmeticGeometry/FundamentalDomains/Signature.lean}{ReciprocityTheory.ChartRestriction.G.FD.bchar}
\leanref{ReciprocityTheory/ReductiveGroups/ArithmeticGeometry/ArakelovVolumes/UpperBounds/FundSetMeasStatement.lean}{ReciprocityTheory.ChartRestriction.G.AR.Up.IsMetricW}
\leanref{ReciprocityTheory/ReductiveGroups/ArithmeticGeometry/ArakelovVolumes/UpperBounds/FundSetMeasStatement.lean}{ReciprocityTheory.ChartRestriction.G.AR.Up.WLower}
\leanref{ReciprocityTheory/ReductiveGroups/ArithmeticGeometry/ArakelovVolumes/UpperBounds/FundSetMeasStatement.lean}{ReciprocityTheory.ChartRestriction.G.AR.Up.IsBEq}
\leanref{ReciprocityTheory/ReductiveGroups/ArithmeticGeometry/ArakelovVolumes/UpperBounds/FundSetMeasStatement.lean}{ReciprocityTheory.ChartRestriction.G.AR.Up.InAt}
\leanref{ReciprocityTheory/ReductiveGroups/ArithmeticGeometry/ArakelovVolumes/UpperBounds/IsOrth.lean}{ReciprocityTheory.ChartRestriction.G.AR.Up.kaF}
\leanref{ReciprocityTheory/ReductiveGroups/ArithmeticGeometry/ArakelovVolumes/UpperBounds/IsOrth.lean}{ReciprocityTheory.ChartRestriction.G.AR.Up.muF}
\leanref{ReciprocityTheory/ReductiveGroups/ArithmeticGeometry/ArakelovVolumes/UpperBounds/IsOrth.lean}{ReciprocityTheory.ChartRestriction.G.AR.Up.poin}

\begin{lemma}\label{is-onto-share-counts-section-kronecker-points}
Let $X$ consist of the data (1) to (6) of flag data of rank $r$ over $E$, with archimedean data for $(L,F,T)$ and a
flat model for them, with base points $y_1,\dots,y_\delta$ (Definition~\ref{is-onto-share-counts-model}). Let
$1\le j\le\delta$, $\gamma\in\Gamma$, $b\in\mathcal B^\Psi$ and $x_1=\gamma y_jb$. Then $w_0(x_1)\ne0$, $s_k(x_1)=0$ for
every $k$, and the derivative at $0$ of $z\mapsto\big(s_k(c_{x_1}(z))/w_0(c_{x_1}(z))\big)_k$ is invertible.
\end{lemma}
\leanref{Development/AbelianVarietyTheory/Proofs/TorusHaarMeasure/KronHolPsiOf.lean}{GlobalTheory.KronAssembly.sec_ne_zero}
\leanref{Development/AbelianVarietyTheory/Proofs/TorusHaarMeasure/DetFderivNeZeroOrbit.lean}{GlobalTheory.OrbitCoord.det_fderiv_ne_zero_orbit}
\leanref{Development/AbelianVarietyTheory/Proofs/TorusHaarMeasure/DetFderivNeZeroOrbit.lean}{GlobalTheory.OrbitCoord.exists_swap}
\begin{proof}
The point $x_1$ lies in $\mathcal D_X$, as $\gamma$ preserves $h$ and by Lemma~\ref{is-onto-share-majorant}(2). As $w_0$
is a section of weight $L$, $w_0(x_1)=\chi_L(\gamma)\,b^Lw_0(y_j)$, with $\chi_L(\gamma)\ne0$ as $\chi_L$ has finite
order on $\Gamma$, $b^L\ne0$ and $w_0(y_j)\ne0$. The $s_k$ vanish at $x_1$, as the common zeros of the $s_k$ in
$\mathcal D_X$ are the points $\gamma'y_{j'}b'$.

For $z$ near $0$ the leading principal minors of $b_\tau\ell_\tau(z)$ are nonzero, as at $z=0$ they are products of
diagonal entries of $b_\tau$. The triangular factorization in the proof of
Lemma~\ref{is-onto-share-counts-section-transfer-lu} writes $b_\tau\ell_\tau(z)=\ell_\tau(\Phi(z))\beta(z)_\tau$ near $0$, with
$\Phi$ and $\beta$ holomorphic, $\beta(z)\in\mathcal B^\Psi$, $\Phi(0)=0$ and $\beta(0)=b$, by uniqueness. So
$c_{x_1}(z)=\gamma\,c_{y_j}(\Phi(z))\,\beta(z)$ near $0$, and
Lemma~\ref{is-onto-share-counts-section-transfer-jacobian} gives $\det D\Phi(0)=b^k\ne0$, where $k_{\tau,i}=2i-r-1$.
The quotients $s_k/w_0$ of sections of the same weight $L$ do not change under $x\mapsto\gamma xb'$; so
$(s_k/w_0)(c_{x_1}(z))=(s_k/w_0)(c_{y_j}(\Phi(z)))$ near $0$. The derivative at $0$ is therefore the composite of the
derivative at $0$ of $z\mapsto((s_k/w_0)(c_{y_j}(z)))_k$, invertible by (5) of
Definition~\ref{is-onto-share-counts-model}, with $D\Phi(0)$.
% Lean: KronAssembly.sec_ne_zero, Fl.S_eq, and OrbitCoord.det_fderiv_ne_zero_orbit (exists_swap for the
% factorization of rmul b (nmat z), KapGGStatement for det D Phi(0) = bchar (kapW X) (bf 0), and
% fderiv_comp), as InterpWeight.core assembles them in key. Size: about 420 lines of Lean.
\end{proof}

\begin{lemma}\label{is-onto-share-counts-section-kronecker-gauge}
Let $X$ consist of the data (1) to (6) of flag data of rank $r$ over $E$, use the notation of
Definition~\ref{is-onto-share-counts-section-gauge}, and let $\lambda$ be a weight. Let $W_\lambda(x)=N_\lambda(x)^2$ for
$x\in\mathcal D_X$ and $W_\lambda(x)=0$ otherwise. Then $W_\lambda$ is a gauge of weight $\lambda$ that is bounded below.
\end{lemma}
\leanref{Development/AbelianVarietyTheory/Proofs/TorusHaarMeasure/ContinuousOnOfCharts.lean}{GlobalTheory.KronMetric.kronMetric}
\begin{proof}
The function $N_\lambda$ is continuous on the open set $\mathcal D_X$, so $W_\lambda$ is measurable. For
$b\in\mathcal B^\Psi$ and $x\in\mathcal D_X$, $N_\lambda(xb)=|b^\lambda|^{-1}N_\lambda(x)$, as in the proof of
Lemma~\ref{is-onto-share-counts-section-transfer}; and for $\gamma\in\Gamma$, $\gamma x$ has the same Gram matrices as $x$,
as $\gamma$ preserves $h$, so $N_\lambda(\gamma x)=N_\lambda(x)$. So $W_\lambda$ is a gauge, and $W_\lambda\geq1\cdot
N_\lambda^2$ on $\mathcal D_X$.
% Lean: KronMetric.kronMetric with psi = 0, for the indicator of X.dom of exp(-(phla la + psi)); the
% constant is exp(-C) with C = 0.
\end{proof}

\begin{lemma}\label{is-onto-share-counts-section-kronecker-solve-simple-zero}
Let $X$ consist of the data (1) to (6) of flag data of rank $r$ over $E$, with $n=|\Psi|r(r-1)/2>0$, use the charts
$c_y$ of Definition~\ref{is-onto-share-counts-model}, and write $|z|$ for the Euclidean norm of $z\in\mathbf C^n$. Let
$x_1$ be an orthonormal point, $\varrho>0$ a real number, and $u,w_0,s_1,\dots,s_n$ functions on $M_r(\mathbf C)^\Psi$
whose composites with $c_{x_1}$ are holomorphic on the ball $|z|<\varrho$. Suppose that $w_0(x_1)\ne0$, that
$s_k(x_1)=0$ for every $k$, that the derivative at $0$ of $z\mapsto\big(s_k(c_{x_1}(z))/w_0(c_{x_1}(z))\big)_k$ is
invertible, and that
\[
 \int_{|z|<\varrho}|u(c_{x_1}(z))|^2\Big(\sum_k|s_k(c_{x_1}(z))|^2\Big)^{-n}dz<\infty,
\]
the integrand being read as $0$ where the sum vanishes. Then $u(x_1)=0$.
\end{lemma}
\leanref{ReciprocityTheory/ReductiveGroups/ArithmeticGeometry/ArakelovVolumes/UpperBounds/FundSetMeasStatement.lean}{ReciprocityTheory.ChartRestriction.G.AR.Up.VanishAtSStatement}
\leanref{Development/AbelianVarietyTheory/Proofs/TorusHaarMeasure/VanishOfSimpleZero.lean}{GlobalTheory.vanishAtS_holds}
\leanref{Development/AbelianVarietyTheory/Proofs/TorusHaarMeasure/VanishOfSimpleZero.lean}{GlobalTheory.VanishAtSimpleZero.vanish_of_simple_zero}
\leanref{Development/AbelianVarietyTheory/Proofs/TorusHaarMeasure/VanishOfSimpleZero.lean}{GlobalTheory.VanishAtSimpleZero.lintegral_eq_top}
\begin{proof}
Write $\sigma_k(z)=s_k(c_{x_1}(z))$ and $\omega(z)=w_0(c_{x_1}(z))$, and suppose that $u(x_1)\ne0$. As $c_{x_1}(0)=x_1$, the
map $f=(\sigma_k/\omega)_k$ is holomorphic near $0$, $f(0)=0$, and its derivative $\Lambda$ at $0$ is invertible. So
$|\Lambda z|\geq\kappa|z|$ for some $\kappa>0$, and as $f(z)=\Lambda z+o(|z|)$ there is $\delta>0$ with $f(z)\ne0$ for
$0<|z|\le\delta$. Hence $S(z)=\sum_k|\sigma_k(z)|^2>0$ for $0<|z|\le\delta$. Each $\sigma_k$ is differentiable at $0$ with
$\sigma_k(0)=0$, so $|\sigma_k(z)|\le C_k|z|$ near $0$, and $S(z)\le B|z|^2$ with $B=\sum_kC_k^2+1$. By continuity,
$|u(c_{x_1}(z))|^2\geq|u(x_1)|^2/2$ near $0$. So there is $0<\varepsilon\le\varrho$ such that for $0<|z|<\varepsilon$ the
integrand is at least $\frac12|u(x_1)|^2B^{-n}|z|^{-2n}$. But $\int_{|z|<\varepsilon}|z|^{-2n}dz=\infty$ on
$\mathbf C^n=\mathbf R^{2n}$: in polar coordinates it is a positive multiple of $\int_0^\varepsilon t^{-2n}t^{2n-1}\,dt$.
This contradicts the hypothesis.
% Lean: VanishAtSStatement (which also assumes x1 in X.dom, true for orthonormal x1), proved by
% GlobalTheory.vanishAtS_holds = VanishAtSimpleZero.vanish_of_simple_zero for u o c, w0 o c and s_k o c
% (IW.cpt_zero): HE.iso_zero for the isolated zero, isBigO_sub for |s_k| <= C_k |z|, continuity of |u|^2 with
% the bound |u(x1)|^2 / 2, and lintegral_eq_top through L2DbarEstimates.I.vanish_div and C.mvDiv (the integral of
% (c |z|^2)^-n over a ball is infinite). The second proof in UpperBounds/VanishAtSStatementPf.lean
% (lower_of_det, crit_div through mvInt) is the same argument.
\end{proof}

\begin{lemma}\label{is-onto-share-counts-section-kronecker-solve-chart-matrices}
Let $p,q\geq0$ be integers, and write $M^*$ for the conjugate transpose of a complex matrix $M$.
\begin{enumerate}
\item Let $Z$ be a complex $p\times q$ matrix such that $1-ZZ^*$ is positive semidefinite. Then $\det(1-ZZ^*)$ is a
  real number with $0\le\det(1-ZZ^*)\le1$.
\item Let $Z$ and $W$ be complex $p\times q$ matrices such that $1-ZZ^*$ and $1-WW^*$ are positive definite. Then
  \[
   \det(1-ZZ^*)\det(1-WW^*)\le|\det(1-ZW^*)|^2.
  \]
\end{enumerate}
\end{lemma}
\leanref{ReciprocityTheory/ReductiveGroups/ArithmeticGeometry/ArakelovVolumes/UpperBounds/ExhaustFStatementPfE.lean}{ReciprocityTheory.ChartRestriction.G.AR.Up.ExF.det_one_sub_mul_conjTranspose_mem_Icc}
\leanref{ReciprocityTheory/ReductiveGroups/ArithmeticGeometry/ArakelovVolumes/UpperBounds/ExhaustStatementPfA.lean}{ReciprocityTheory.ChartRestriction.G.AR.Up.ExhA.det_mul_det_le_normSq_det}
\leanref{ReciprocityTheory/ReductiveGroups/ArithmeticGeometry/ArakelovVolumes/UpperBounds/ExhaustStatementPfA.lean}{ReciprocityTheory.ChartRestriction.G.AR.Up.ExhA.cross_psd}
\leanref{ReciprocityTheory/L2DbarEstimates/LeviForms/DetMonoDetOneAdd.lean}{ReciprocityTheory.L2DbarEstimates.I.maDet}
\begin{proof}
(1) The matrix $1-ZZ^*$ is Hermitian and positive semidefinite, so its eigenvalues are real and $\geq0$. For a unit
eigenvector $v$ with eigenvalue $\lambda$, $\lambda=v^*(1-ZZ^*)v=1-|Z^*v|^2\le1$. The determinant is the product of
the eigenvalues.

(2) Put $B=1-WW^*$ and $C=1-ZW^*$. As $B$ is positive definite, $\det B>0$ and $B^{-1}$ is Hermitian. Since
$(CB^{-1}-1)B=C-B$ and $C+C^*-B=(1-ZZ^*)+(W-Z)(W-Z)^*$, expanding gives
\[
 CB^{-1}C^*-(1-ZZ^*)=(CB^{-1}-1)B(CB^{-1}-1)^*+(W-Z)(W-Z)^*,
\]
and both terms on the right are positive semidefinite. So $v^*(1-ZZ^*)v\le v^*CB^{-1}C^*v$ for every
$v\in\mathbf C^p$, and Lemma~\ref{is-onto-semistable-det}(2) gives
$\det(1-ZZ^*)\le\det(CB^{-1}C^*)=|\det C|^2/\det B$. Multiply by $\det B$.
% Lean: (1) is ExF.det_one_sub_mul_conjTranspose_mem_Icc (private; the eigenvalues through
% IsHermitian.det_eq_prod_eigenvalues). (2) is ExhA.det_mul_det_le_normSq_det: cross_psd is the displayed
% identity, and the determinant comparison is I.maDet (DetMonoStatement, proved by
% Bl.DetMonotoneStarMulSelf.det_mono: B = X^H X, A = X^H (1 + D) X and det(1 + D) >= 1 by eigenvalues), the
% argument of Lemma is-onto-semistable-det(2), which the book cites as 1 - Z Z^H is positive definite here.
% Size: about 170 lines of Lean (the ExhA part of ExhaustStatementPfA, det_one_sub_mul_conjTranspose_mem_Icc
% and DetMonotoneStarMulSelf).
\end{proof}

\begin{lemma}\label{is-onto-share-counts-section-kronecker-solve-chart-domain}
Let $X$ consist of the data (1) to (6) of flag data of rank $r$ over $E$
(Definition~\ref{is-onto-share-flag-data}), and let $x\in M_r(\mathbf C)^\Psi$. For $\tau\in\Psi$ let $g_\tau(x)$ be
the Gram matrix $\big(h_\tau(x_\tau e_a,x_\tau e_b)\big)_{a,b\le p_0}$ of the first $p_0$ columns of $x_\tau$. Then
$x\in\mathcal D_X$ if and only if, for every $\tau\in\Psi$, $\det x_\tau\ne0$ and $g_\tau(x)$ is positive definite.
\end{lemma}
\leanref{ReciprocityTheory/ReductiveGroups/ArithmeticGeometry/ArakelovVolumes/UpperBounds/ExhaustFStatementPfB.lean}{ReciprocityTheory.ChartRestriction.G.AR.Up.ExF.mem_dom_of_gramPos}
\leanref{ReciprocityTheory/ReductiveGroups/ArithmeticGeometry/ArakelovVolumes/UpperBounds/ExhaustFStatementPfB.lean}{ReciprocityTheory.ChartRestriction.G.AR.Up.ExF.neg_of_perp'}
\leanref{ReciprocityTheory/ReductiveGroups/ArithmeticGeometry/ArakelovVolumes/UpperBounds/ExhaustFStatementPfB.lean}{ReciprocityTheory.ChartRestriction.G.AR.Up.ExF.psgn_dmin_le}
\leanref{ReciprocityTheory/ReductiveGroups/ArithmeticGeometry/ArakelovVolumes/UpperBounds/ExhaustFStatementPfB.lean}{ReciprocityTheory.ChartRestriction.G.AR.Up.ExF.psgn_dmin_gt}
\leanref{ReciprocityTheory/ReductiveGroups/ArithmeticGeometry/FundamentalDomains/LminorSubmatrixCastSucc.lean}{ReciprocityTheory.ChartRestriction.G.FD.Data.gramPos_posDef}
\begin{proof}
If $x\in\mathcal D_X$, then $\det x_\tau\ne0$, and $g_\tau(x)$ is positive definite as in the proof of
Lemma~\ref{is-onto-share-majorant}(1): its leading principal minors are the $d_i(x,\tau)$ with $i\le p_0$, which are
real and positive as $\epsilon_i=1$, and Sylvester's criterion \cite[Theorem~7.2.5]{HJ13} applies.

Conversely, suppose that the conditions hold. Fix $\tau$, write $h=h_\tau$, let $x_1,\dots,x_r$ be the columns of
$x_\tau$, and put $g=g_\tau(x)$. As $h(v,u)=\overline{h(u,v)}$, Gram matrices for $h$ are Hermitian.

\emph{Vectors orthogonal to the first columns.} Let $e\ne0$ in $\mathbf C^r$ with $h(e,x_b)=0$ for $b\le p_0$. Then
$h(e,e)<0$. Suppose instead that $h(e,e)\geq0$. By (2) of Definition~\ref{is-onto-share-flag-data} there is an
invertible matrix $P$ whose columns $P_1,\dots,P_r$ have the Gram matrix $J$ for $h$. The first $p_0$ coordinates
of $P^{-1}v$ are $p_0$ linear forms in $v$, so there is $(t,c_1,\dots,c_{p_0})\ne0$ such that $v=te+w$, with
$w=\sum_{b\le p_0}c_bx_b$, has $(P^{-1}v)_k=0$ for $k\le p_0$. Writing $v=\sum_ky_kP_k$ with $y=P^{-1}v$,
$h(v,v)=-\sum_{k>p_0}|y_k|^2\le0$. As $h(e,w)=h(w,e)=0$, also $h(v,v)=h(w,w)+|t|^2h(e,e)$, and $h(w,w)\geq0$ as $g$ is
positive definite. So $h(v,v)=0$, hence $y=0$ and $v=0$. Then $h(w,w)=0$, so all $c_b$ vanish, and $te=0$ gives
$t=0$, a contradiction.

\emph{The minors up to $p_0$.} For $i\le p_0$, $d_i(x,\tau)$ is the determinant of the leading $i\times i$ block of
$g$, which is positive definite; so $\epsilon_id_i(x,\tau)=d_i(x,\tau)>0$.

\emph{The minors beyond $p_0$.} Let $i=p_0+k$ with $1\le k\le r-p_0$. The Gram matrix of $x_1,\dots,x_i$ has the
blocks $\left(\begin{smallmatrix}g&B\\ C&D\end{smallmatrix}\right)$, where $D$ is the Gram matrix of
$x_{p_0+1},\dots,x_i$. As $g$ is invertible, $d_i(x,\tau)=\det g\cdot\det(D-Cg^{-1}B)$. Put $\alpha=Cg^{-1}$ and
$w_j=x_{p_0+j}-\sum_{a\le p_0}\alpha_{ja}x_a$ for $1\le j\le k$. Then $h(w_j,x_b)=C_{jb}-(\alpha g)_{jb}=0$ for
$b\le p_0$, so $h(w_j,w_l)=h(w_j,x_{p_0+l})=(D-\alpha B)_{jl}$: the matrix $D-Cg^{-1}B$ is the Gram matrix of
$w_1,\dots,w_k$. Let $\xi\ne0$ in $\mathbf C^k$ and $w=\sum_j\bar\xi_jw_j$. In $w$ the vector $x_{p_0+j}$ has the
coefficient $\bar\xi_j$, and $x_1,\dots,x_r$ are linearly independent as $\det x_\tau\ne0$; so $w\ne0$, and
$h(w,x_b)=0$ for $b\le p_0$. By the first step, $\xi^*(D-Cg^{-1}B)\xi=h(w,w)<0$. So $-(D-Cg^{-1}B)$ is positive
definite, and $\det(D-Cg^{-1}B)$ is $(-1)^k$ times a positive number. As $\det g>0$ and $\epsilon_i=(-1)^k$,
$\epsilon_id_i(x,\tau)>0$.

With $\det x_\tau\ne0$ for every $\tau$, this shows that $x\in\mathcal D_X$.
% Lean: ExF.mem_dom_of_gramPos (through Data.mem_dom_iff) from ExF.psgn_dmin_le (i <= p0, the principal
% submatrix of gramPos) and ExF.psgn_dmin_gt (i = p0 + k, with k >= 0: Matrix.det_fromBlocks for the Schur
% complement, wv for the w_j, hSw for their Gram matrix, hli for the independence, and psgn_add for
% epsilon_i = (-1)^k); the first step is ExF.neg_of_perp', where P comes from X.signature and the kernel vector
% of the p0 x (p0 + 1) matrix R gives (t, c). The converse is Data.gramPos_posDef (posDef_of_lminor). g_tau(x) is
% X.gramPos x tau. Size: about 300 lines of Lean (ExhaustFStatementPfB).
\end{proof}

\begin{lemma}\label{is-onto-share-counts-section-kronecker-solve-chart-block}
Let $X$ consist of the data (1) to (6) of flag data of rank $r$ over $E$, use the notation of
Definition~\ref{is-onto-share-counts-section-levi}, write $M^t$ and $M^*$ for the transpose and the conjugate
transpose of a complex matrix $M$, and let $x_0$ be an orthonormal point. For $z\in\mathbf C^n$ and $\tau\in\Psi$ write
the lower unitriangular matrix $\ell_\tau(z)$ of Definition~\ref{is-onto-share-counts-model} in blocks
\[
 \ell_\tau(z)=\begin{pmatrix}N_\tau(z)&0\\ Y_\tau(z)&N'_\tau(z)\end{pmatrix},
\]
with $N_\tau(z)$ of size $p_0\times p_0$ and $Y_\tau(z)$ of size $(r-p_0)\times p_0$, and put
\[
 Z_\tau(z)=\big(Y_\tau(z)N_\tau(z)^{-1}\big)^t,\qquad R_\tau(z)=1-Z_\tau(z)Z_\tau(z)^*,
\]
a $p_0\times(r-p_0)$ matrix and a Hermitian $p_0\times p_0$ matrix.
\begin{enumerate}
\item The entries of $Z_\tau(z)$ are polynomials in the coordinates of $z$.
\item The Gram matrix of the first $p_0$ columns of $c_{x_0}(z)_\tau$ is $N_\tau(z)^tR_\tau(z)\overline{N_\tau(z)}$, and
  $d_{p_0}(c_{x_0}(z),\tau)=\det R_\tau(z)$.
\item $z\in\Omega(x_0)$ if and only if $R_\tau(z)$ is positive definite for every $\tau\in\Psi$.
\item For $z\in\Omega(x_0)$, $\det R_\tau(z)>0$ for every $\tau\in\Psi$, and
  $\rho_{x_0}(z)=-\sum_{\tau\in\Psi}\log\det R_\tau(z)$.
\end{enumerate}
\end{lemma}
\leanref{ReciprocityTheory/ReductiveGroups/ArithmeticGeometry/ArakelovVolumes/UpperBounds/ExhaustFStatementPfC.lean}{ReciprocityTheory.ChartRestriction.G.AR.Up.ExF.Zt}
\leanref{ReciprocityTheory/ReductiveGroups/ArithmeticGeometry/ArakelovVolumes/UpperBounds/ExhaustFStatementPfC.lean}{ReciprocityTheory.ChartRestriction.G.AR.Up.ExF.Am}
\leanref{ReciprocityTheory/ReductiveGroups/ArithmeticGeometry/ArakelovVolumes/UpperBounds/ExhaustFStatementPfD.lean}{ReciprocityTheory.ChartRestriction.G.AR.Up.ExF.cd_Zt}
\leanref{ReciprocityTheory/ReductiveGroups/ArithmeticGeometry/ArakelovVolumes/UpperBounds/ExhaustFStatementPfC.lean}{ReciprocityTheory.ChartRestriction.G.AR.Up.ExF.gramPos_eq}
\leanref{ReciprocityTheory/ReductiveGroups/ArithmeticGeometry/ArakelovVolumes/UpperBounds/ExhaustFStatementPfC.lean}{ReciprocityTheory.ChartRestriction.G.AR.Up.ExF.det_gramPos}
\leanref{ReciprocityTheory/ReductiveGroups/ArithmeticGeometry/ArakelovVolumes/UpperBounds/ExhaustFStatementPfC.lean}{ReciprocityTheory.ChartRestriction.G.AR.Up.ExF.mem_OmF_iff}
\leanref{ReciprocityTheory/ReductiveGroups/ArithmeticGeometry/ArakelovVolumes/UpperBounds/ExhaustFStatementPfC.lean}{ReciprocityTheory.ChartRestriction.G.AR.Up.ExF.rhF_eq_sum}
\begin{proof}
Write $N=N_\tau(z)$, $Y=Y_\tau(z)$ and $Z=Z_\tau(z)$.

(1) The matrix $N$ is lower unitriangular, so $\det N=1$ and $N^{-1}$ is the adjugate of $N$, whose entries are
polynomials in those of $N$. The entries of $N$ and $Y$ are coordinates of $z$ or constants.

(2) As $x_0$ is orthonormal, the Gram matrix of the columns of $c_{x_0}(z)_\tau=x_{0,\tau}\ell_\tau(z)$ is
$\ell_\tau(z)^tJ\overline{\ell_\tau(z)}$. The first $p_0$ columns of $\ell_\tau(z)$ form the blocks $N$ and $Y$, and $J$
has the diagonal blocks $1$ of size $p_0$ and $-1$ of size $r-p_0$; so the Gram matrix of the first $p_0$ columns of
$c_{x_0}(z)_\tau$ is $N^t\bar N-Y^t\bar Y$. As $Y^t=N^tZ$, $Y^t\bar Y=N^tZ(N^tZ)^*=N^tZZ^*\bar N$, and the Gram matrix
is $N^tR_\tau(z)\bar N$. Its determinant is $d_{p_0}(c_{x_0}(z),\tau)$, and it equals $\det R_\tau(z)$, as $\det N=1$.

(3) As $\det\ell_\tau(z)=1$, $\det c_{x_0}(z)_\tau=\det x_{0,\tau}$, which is nonzero by
Lemma~\ref{is-onto-share-majorant}(4). By Lemma~\ref{is-onto-share-counts-section-kronecker-solve-chart-domain},
$z\in\Omega(x_0)$ if and only if the Gram matrices of (2) are positive definite for every $\tau$. As $N$ is invertible
and $\bar N=(N^t)^*$, $N^tR_\tau(z)\bar N$ is positive definite if and only if $R_\tau(z)$ is.

(4) A positive definite matrix has a positive determinant. So by (3) and (2),
$|d_{p_0}(c_{x_0}(z),\tau)|=\det R_\tau(z)>0$, and the formula is the definition of $\rho_{x_0}$.
% Lean: N_tau(z) is ExF.n1, Y_tau(z) is ExF.Wm, Z_tau(z) is ExF.Zt and R_tau(z) is ExF.Am, built from
% ExF.nM X z tau = mat (nmat X z) tau. (1) is ExF.cd_Zt (holomorphy, through unipotentFrame_inv_eq_adjugate and
% cd_nM); (2) is gramPos_cpt, unipotentFrame_transpose_mul, gramPos_eq, det_gramPos and leadingMinor_eq_gramDet;
% (3) is ExF.mem_OmF_iff, with det_nM and det_ne_of_orth for the determinant, Data.gramPos_posDef and
% ExF.mem_dom_of_gramPos (Lemma is-onto-share-counts-section-kronecker-solve-chart-domain); (4) is
% ExF.rhF_eq_sum (with log |det Am|) and rhF_eq_log. Size: about 230 lines of Lean (ExhaustFStatementPfC and the
% first part of ExhaustFStatementPfD).
\end{proof}

\begin{lemma}\label{is-onto-share-counts-section-kronecker-solve-chart}
Let $X$ consist of the data (1) to (6) of flag data of rank $r$ over $E$, use the notation of
Definition~\ref{is-onto-share-counts-section-levi}, write $|z|$ for the Euclidean norm of $z\in\mathbf C^n$, and let
$x_0$ be an orthonormal point. Then $\Omega(x_0)$ is an open connected set that contains $0$, and $\rho_{x_0}$ is
smooth on $\Omega(x_0)$. For every $z\in\Omega(x_0)$ and every $\xi\in\mathbf C^n$, $\rho_{x_0}(z)\geq0$ and
$\mathcal L_{\rho_{x_0}}(z;\xi)\geq0$. For every real number $a$, the set of the $z\in\Omega(x_0)$ with
$\rho_{x_0}(z)+\log(1+|z|^2)\le a$ is compact.
\end{lemma}
\leanref{ReciprocityTheory/ReductiveGroups/ArithmeticGeometry/ArakelovVolumes/UpperBounds/OntoStatementPfB.lean}{ReciprocityTheory.ChartRestriction.G.AR.Up.OntoB.isOpen_OmF}
\leanref{ReciprocityTheory/ReductiveGroups/ArithmeticGeometry/ArakelovVolumes/UpperBounds/OntoStatementPfB.lean}{ReciprocityTheory.ChartRestriction.G.AR.Up.OntoB.rhF_contDiffAt}
\leanref{ReciprocityTheory/ReductiveGroups/ArithmeticGeometry/ArakelovVolumes/UpperBounds/OmFConnPf.lean}{ReciprocityTheory.ChartRestriction.G.AR.Up.OmFConnStatement_holds}
\leanref{ReciprocityTheory/ReductiveGroups/ArithmeticGeometry/ArakelovVolumes/UpperBounds/ExhaustFStatementPfE.lean}{ReciprocityTheory.ChartRestriction.G.AR.Up.ExF.rhF_nonneg}
\leanref{ReciprocityTheory/ReductiveGroups/ArithmeticGeometry/ArakelovVolumes/UpperBounds/ExhaustFStatementPfE.lean}{ReciprocityTheory.ChartRestriction.G.AR.Up.ExF.levi_rhF_nonneg}
\leanref{ReciprocityTheory/ReductiveGroups/ArithmeticGeometry/ArakelovVolumes/UpperBounds/ExhaustFStatementPfE.lean}{ReciprocityTheory.ChartRestriction.G.AR.Up.ExF.isCompact_subF}
\begin{proof}
The set $\Omega(x_0)$ is the inverse image of the open set $\mathcal D_X$ under the polynomial map $c_{x_0}$, so it is
open, and it contains $0$, as $c_{x_0}(0)=x_0$ lies in $\mathcal D_X$ by Lemma~\ref{is-onto-share-majorant}(3). On
$\Omega(x_0)$ the minors $d_{p_0}(c_{x_0}(z),\tau)$ are nonzero polynomials in the real and imaginary parts of $z$, so
$\rho_{x_0}$ is smooth there.

Use the matrices $N_\tau(z)$, $Y_\tau(z)$, $Z_\tau(z)$ and $R_\tau(z)$ of
Lemma~\ref{is-onto-share-counts-section-kronecker-solve-chart-block}, and for $y\in\mathbf C^n$ put
$q(y)=\prod_{\tau\in\Psi}\det R_\tau(y)$. By parts (3) and (4) of that lemma, for $y\in\Omega(x_0)$ every $R_\tau(y)$ is
positive definite, $q(y)>0$ and $\rho_{x_0}(y)=-\log q(y)$.

\emph{Connectedness.} For a real number $t$ and $z\in\mathbf C^n$ let $u_t(z)$ be the point obtained from $z$ by
multiplying the coordinates $z_{\tau,a,b}$ with $b\le p_0<a$ by $t$, and $v_t(z)$ the point obtained by replacing
these coordinates by $0$ and multiplying the others by $t$. These coordinates are the entries of the blocks $Y_\tau$,
and the entries of $N_\tau$ below the diagonal are other coordinates. So $N_\tau(u_t(z))=N_\tau(z)$ and
$Y_\tau(u_t(z))=tY_\tau(z)$, whence $Z_\tau(u_t(z))=tZ_\tau(z)$ and
\[
 R_\tau(u_t(z))=(1-t^2)\cdot1+t^2R_\tau(z);
\]
and $Y_\tau(v_t(z))=0$, whence $R_\tau(v_t(z))=1$. Let $y\in\Omega(x_0)$. Every $R_\tau(y)$ is positive definite, hence
so is every $R_\tau(u_t(y))$ for $0\le t\le1$, and every $R_\tau(v_t(y))=1$ is positive definite. By
Lemma~\ref{is-onto-share-counts-section-kronecker-solve-chart-block}(3), $u_t(y)$ and $v_t(y)$ lie in $\Omega(x_0)$ for
$0\le t\le1$. The continuous path $t\mapsto v_t(y)$ on $[0,1]$ joins $v_0(y)=0$ to $v_1(y)=u_0(y)$, and the path
$t\mapsto u_t(y)$ joins $u_0(y)$ to $u_1(y)=y$. The union of their images is a connected subset of $\Omega(x_0)$ that
contains $0$ and $y$. So $\Omega(x_0)$ is connected.

\emph{Positivity.} Let $z\in\Omega(x_0)$. Every $\det R_\tau(z)$ is positive, as $R_\tau(z)$ is positive definite, and
at most $1$ by Lemma~\ref{is-onto-share-counts-section-kronecker-solve-chart-matrices}(1). So every term of $\rho_{x_0}(z)=-\sum_\tau\log\det R_\tau(z)$ is $\geq0$, and $\rho_{x_0}(z)\geq0$.

\emph{The Levi form.} Let $z_0\in\Omega(x_0)$ and $\xi\in\mathbf C^n$. For $y\in\mathbf C^n$ put
\[
 H(y)=\prod_{\tau\in\Psi}\det\big(1-Z_\tau(y)Z_\tau(z_0)^*\big).
\]
By Lemma~\ref{is-onto-share-counts-section-kronecker-solve-chart-block}(1), $H$ is a polynomial in the coordinates of
$y$, so it is holomorphic. For $y\in\Omega(x_0)$,
Lemma~\ref{is-onto-share-counts-section-kronecker-solve-chart-matrices}(2), for $Z=Z_\tau(y)$ and $W=Z_\tau(z_0)$, and
the product over $\tau$ give $q(y)q(z_0)\le|H(y)|^2$; in particular $H(y)\ne0$. At $y=z_0$ the factors of $H$ are the
positive numbers $\det R_\tau(z_0)$, so $|H(z_0)|^2=q(z_0)^2$. Hence the function $g=\rho_{x_0}+\log|H|^2$ on
$\Omega(x_0)$ satisfies
\[
 g(y)=\log\frac{|H(y)|^2}{q(y)}\geq\log q(z_0)=g(z_0)\qquad(y\in\Omega(x_0)),
\]
and as $\Omega(x_0)$ is open, $-g$ has a local maximum at $z_0$. The function $g$ is smooth on $\Omega(x_0)$, as $H$
does not vanish there. At a local
maximum of a real function of class $C^2$ its Levi form is $\le0$ in every direction: the restriction to the complex
line $t\mapsto z_0+t\xi$ has a local maximum at $0$, so its Laplacian, $4$ times the Levi form at $\xi$, is $\le0$.
Hence $\mathcal L_g(z_0;\xi)\geq0$. The function $\log|H|^2$ is pluriharmonic near $z_0$, as $H$ is holomorphic and
$H(z_0)\ne0$, so its Levi form at $z_0$ vanishes. As the Levi form is linear,
$\mathcal L_{\rho_{x_0}}(z_0;\xi)=\mathcal L_g(z_0;\xi)\geq0$.

\emph{The sublevel sets.} Let $a$ be a real number, $K_a$ the set of the $z\in\Omega(x_0)$ with
$\rho_{x_0}(z)+\log(1+|z|^2)\le a$, and $K'_a$ the set of the $z\in\mathbf C^n$ such that every $R_\tau(z)$ is
positive semidefinite and $(1+|z|^2)e^{-a}\le q(z)$. Then $K_a=K'_a$. If $z\in K_a$, then every $R_\tau(z)$ is
positive definite, and the inequality follows from $\rho_{x_0}(z)=-\log q(z)$ by taking exponentials. If $z\in K'_a$,
then $q(z)>0$, and every $\det R_\tau(z)$ is $\geq0$ by
Lemma~\ref{is-onto-share-counts-section-kronecker-solve-chart-matrices}(1); so every $\det R_\tau(z)$ is nonzero, and the
positive semidefinite matrices $R_\tau(z)$ are positive definite. By
Lemma~\ref{is-onto-share-counts-section-kronecker-solve-chart-block}(3), $z\in\Omega(x_0)$, and taking logarithms
gives $z\in K_a$.

The set $K'_a$ is closed. The entries of $R_\tau(z)$ are continuous in $z$, so for each $v\in\mathbf C^{p_0}$ the set
of the $z$ with $v^*R_\tau(z)v\geq0$ is closed; and $K'_a$ is the intersection of these sets, over all $\tau$ and $v$,
with the set where the continuous function $q(z)-(1+|z|^2)e^{-a}$ is $\geq0$. It is bounded: for $z\in K'_a$,
$q(z)\le1$ by Lemma~\ref{is-onto-share-counts-section-kronecker-solve-chart-matrices}(1), so $1+|z|^2\le e^a$. Hence
$K_a=K'_a$ is compact.
% Lean: openness and smoothness are OntoB.isOpen_OmF and OntoB.rhF_contDiffAt (kronSolve uses
% ExF.rhF_contDiffAt and ExF.eF_contDiffAt, through rhF_eq_sum and cd_detAm). Connectedness is
% OmFConnStatement_holds = OFC.omFConn_of_cm, IsPreconnected through isPreconnected_of_forall at 0: u_t is S1 and
% v_t is S2 (isW marks the coordinates of the blocks Y_tau), with scaledWm_linear,
% unipotentFrame_scaling_invariant, scaledZt_linear, retractedZt_eq_zero, S1_mem_OmF,
% radialRetraction_mem_chart, linearScaling_zero_eq_retraction, radialRetraction_at_zero and
% linearScaling_at_one. Positivity is ExF.rhF_nonneg (rhF_eq_sum, det_one_sub_mul_conjTranspose_mem_Icc). The
% Levi form is ExF.levi_rhF_nonneg: Hf is H and P is q, hle from ExhA.det_mul_det_le_normSq_det, hself and hmin;
% Y.LocalMaximumSecondDerivative.leviV_re_nonpos_of_isLocalMax for -g, leviV_ofReal_const_mul and
% leviV_ofReal_add for the linearity, and levi_logNormSq_zero (H.leviLogNormSq with one function) for
% log |H|^2. The sublevel sets are ExF.isCompact_subF: sub_eqF (K_a = K'_a, with exp(log(1 + |z|^2) - a) for
% (1 + |z|^2) e^-a), isClosed_psdF, and the bound |z| <= exp a. Size: about 1,200 lines of Lean in all
% (ExhaustFStatementPfB to PfE, OmFConnPf, the ExhA part of ExhaustStatementPfA and DetMonotoneStarMulSelf),
% of which about 500 are below this proof and not below the three lemmas before it.
\end{proof}

\begin{theorem}[Hörmander's $L^2$ estimate {\cite{Hor65}, \cite[Chapter~VIII, Theorems~4.5, 5.2 and~6.1]{Dem12}}]\label{is-onto-share-counts-section-kronecker-solve-hormander}
Let $U\subset\mathbf C^n$ be open. For a smooth real function $g$ on $U$ write $\mathcal L_g(z;\xi)$ for its Levi form,
and for a smooth complex function $g$ write $\bar\partial_jg=\partial g/\partial\bar z_j$. Let $e$ be a smooth real
function on $U$ with $\mathcal L_e\geq0$ on $U$ such that $\{z\in U:e(z)\le a\}$ is compact for every real number
$a$. Let $\varphi$ be a smooth real function on $U$ with $\mathcal L_\varphi\geq0$ on $U$, $F_0$ a smooth complex
function on $U$, and $h\colon U\to[0,\infty]$ a measurable function such that
\[
 \Big|\sum_j\bar\partial_jF_0(z)\,\bar\xi_j\Big|^2\le h(z)\,\mathcal L_\varphi(z;\xi)
\]
for all $z\in U$ and $\xi\in\mathbf C^n$, with $\infty\cdot0=0$. Then there is a holomorphic function $F$ on $U$ with
\[
 \int_U|F_0(z)-F(z)|^2e^{-\varphi(z)}\,dz\le\int_Uh(z)\,e^{-\varphi(z)}\,dz .
\]
\end{theorem}
\leanref{ReciprocityTheory/ReductiveGroups/ArithmeticGeometry/ArakelovVolumes/UpperBounds/InterpStatementPfC.lean}{ReciprocityTheory.ChartRestriction.G.AR.Up.HE.hormander_open}
\leanref{ReciprocityTheory/L2DbarEstimates/JacobianInterpolationStatement.lean}{ReciprocityTheory.L2DbarEstimates.HStatement}
\leanref{ReciprocityTheory/L2DbarEstimates/NormFibrations/Hormander.lean}{ReciprocityTheory.L2DbarEstimates.hormander}

\begin{theorem}[Extension of $L^2$ holomorphic functions {\cite[Chapter~VIII, Lemma~7.3, and Chapter~I, Corollary~3.30]{Dem12}}]\label{is-onto-share-counts-section-kronecker-solve-riemann}
Let $B\subset\mathbf C^N$ be an open ball, $g$ a holomorphic function on $B$, and $u$ a holomorphic function on
$B\setminus\{g=0\}$ with $\int_{B\setminus\{g=0\}}|u(z)|^2\,dz<\infty$. Then there is a holomorphic function $U$ on $B$
with $U=u$ on $B\setminus\{g=0\}$.
\end{theorem}
\leanref{ReciprocityTheory/L2DbarEstimates/JacobianInterpolationStatement.lean}{ReciprocityTheory.L2DbarEstimates.EXTStatement}
\leanref{ReciprocityTheory/L2DbarEstimates/NormFibrations/MvBlowup.lean}{ReciprocityTheory.L2DbarEstimates.ext}

\begin{lemma}\label{is-onto-share-counts-section-kronecker-solve-interpolation}
Let $U\subset\mathbf C^n$ be open, write $|z|$ for the Euclidean norm and $\mathcal L_g$ for the Levi form of a smooth
real function $g$, and for $\varepsilon>0$ fix a smooth function $\chi_\varepsilon$ on $\mathbf C^n$ with values in
$[0,1]$, equal to $1$ for $|z|\le\varepsilon$ and to $0$ for $|z|\geq2\varepsilon$. Let $e\geq0$ be a smooth function
on $U$ with $\mathcal L_e\geq0$ on $U$, such that $\{z\in U:e(z)\le a\}$ is compact for every real number $a$. Let
$s_1,\dots,s_n$ be holomorphic functions on $U$ such that, for some $k_0$, $s_{k_0}$ does not vanish identically on
any nonempty open subset of $U$, and put $S=\sum_k|s_k|^2$. Let $\psi$ be a smooth real function on $U$ with
$\mathcal L_\psi\geq0$ on $U$. Let $\varepsilon>0$ be such that the closed ball $\bar B$ of radius $2\varepsilon$
around $0$ lies in $U$, $0$ is the only point of $\bar B$ at which all the $s_k$ vanish, and for some $c>0$,
$\mathcal L_\psi(z;\xi)\geq c|\xi|^2$ for all $z\in\bar B$ and $\xi\in\mathbf C^n$. Then there is a holomorphic
function $F$ on $U$ with
\[
 \int_U|\chi_\varepsilon(z)-F(z)|^2e^{-\psi(z)}S(z)^{-n}\,dz<\infty,
\]
the integrand being read as $0$ where $S$ vanishes.
\end{lemma}
\leanref{ReciprocityTheory/ReductiveGroups/ArithmeticGeometry/ArakelovVolumes/UpperBounds/InterpStatementPfG.lean}{ReciprocityTheory.ChartRestriction.G.AR.Up.HE.interp_euclid}
\leanref{ReciprocityTheory/ReductiveGroups/ArithmeticGeometry/ArakelovVolumes/UpperBounds/InterpStatementPfD.lean}{ReciprocityTheory.ChartRestriction.G.AR.Up.HE.ext_glue}
\leanref{ReciprocityTheory/ReductiveGroups/ArithmeticGeometry/ArakelovVolumes/UpperBounds/InterpStatementPfE.lean}{ReciprocityTheory.ChartRestriction.G.AR.Up.HE.epunct_compact}
\leanref{ReciprocityTheory/ReductiveGroups/ArithmeticGeometry/ArakelovVolumes/UpperBounds/InterpStatementPfG.lean}{ReciprocityTheory.ChartRestriction.G.AR.Up.HE.wgt_levi}
\begin{proof}
Let $U'$ be the open set of the $z\in U$ with $s_{k_0}(z)\ne0$, and write $\chi=\chi_\varepsilon$.

\emph{The weight.} On $U'$, $S>0$, and $\Phi=\psi+n\log S$ is smooth. The function $\log S$ is plurisubharmonic on
$U'$, as the logarithm of a sum of squares of absolute values of holomorphic functions without common zero; so
$\mathcal L_\Phi\geq\mathcal L_\psi\geq0$ on $U'$.

\emph{The exhaustion.} On $U'$ put $e'=e+\log(1+|s_{k_0}|^2)-\log|s_{k_0}|^2$. It is smooth; $\log|s_{k_0}|^2$ is
pluriharmonic on $U'$ and $\log(1+|s_{k_0}|^2)$ is plurisubharmonic, so $\mathcal L_{e'}\geq0$. Let $a$ be real. If
$z\in U'$ and $e'(z)\le a$, then $e(z)\le e'(z)\le a$, and $\log(1+|s_{k_0}(z)|^{-2})\le a-e(z)\le a$, so
$e^a|s_{k_0}(z)|^2\geq1$. The set of $z\in U$ with $e(z)\le a$ and $e^a|s_{k_0}(z)|^2\geq1$ is compact and lies in $U'$,
and $e'$ is continuous on it; so $\{z\in U':e'(z)\le a\}$, a closed subset of it, is compact.

\emph{The datum.} Let $A$ be the set of $z$ with $\varepsilon\le|z|\le2\varepsilon$. The derivatives $\bar\partial_j\chi$
vanish off $A$, and $\sum_j|\bar\partial_j\chi|^2\le M$ on $\mathbf C^n$ for some $M$, as this function is continuous
with compact support. As $0\notin A\subset\bar B$, $S>0$ on $A$, so $A\subset U'$, and $e^{-\Phi}\le B_1$ on the compact
set $A$ for some $B_1$. Let $h=M/c$ on $A$ and $h=0$ elsewhere. For $z\in U'$ and $\xi\in\mathbf C^n$, if $z\in A$, then
by the Cauchy--Schwarz inequality
\[
 \Big|\sum_j\bar\partial_j\chi(z)\,\bar\xi_j\Big|^2\le M|\xi|^2\le\frac Mc\,\mathcal L_\psi(z;\xi)\le h(z)\,\mathcal L_\Phi(z;\xi),
\]
and if $z\notin A$ the left side is $0$.

\emph{The solution on $U'$.} Theorem~\ref{is-onto-share-counts-section-kronecker-solve-hormander}, for $U'$, $e'$,
$\varphi=\Phi$, $F_0=\chi$ and $h$, gives a holomorphic function $G$ on $U'$ with
\[
 \int_{U'}|\chi-G|^2e^{-\Phi}\,dz\le\int_{U'}h\,e^{-\Phi}\,dz\le\frac McB_1\operatorname{vol}(A)<\infty .
\]

\emph{Extension across the zeros of $s_{k_0}$.} Let $p\in U$, and let $B_p$ be an open ball around $p$ whose closure
lies in $U$. Let $K$ bound the continuous function $e^\psi S^n$ on the closure of $B_p$. On $B_p\cap U'$,
$e^\Phi=e^\psi S^n\le K$ and $|\chi|\le1$, so
\[
 |G|^2\le2|\chi|^2+2|\chi-G|^2\le2+2K|\chi-G|^2e^{-\Phi}.
\]
So $G$ is square integrable on $B_p\setminus\{s_{k_0}=0\}$, and
Theorem~\ref{is-onto-share-counts-section-kronecker-solve-riemann} gives a holomorphic function $F_p$ on $B_p$ equal to
$G$ there. Two such functions $F_p$ and $F_{p'}$ agree on $B_p\cap B_{p'}$: they agree off the zero set of $s_{k_0}$,
which contains no nonempty open subset of $U$, and they are continuous. So they define a holomorphic function $F$ on
$U$ with $F=G$ on $U'$.

\emph{The integral.} The zero set of $s_{k_0}$ in $U$ has measure zero, as $s_{k_0}$ does not vanish identically near
any point of $U$. On $U'$, $|\chi-F|^2e^{-\psi}S^{-n}=|\chi-G|^2e^{-\Phi}$. So the integral in the statement equals
$\int_{U'}|\chi-G|^2e^{-\Phi}\,dz$, which is finite.
% Lean: HE.interp_euclid with U' = {y in U | s k0 y != 0} (isOpen_punct), wgt psi s = psi + n log (nsq s)
% (wgt_smooth, wgt_levi through L2DbarEstimates.lseV), epunct e (s k0) = e + log(1 + |s k0|^2) - log |s k0|^2
% (epunct_smooth, epunct_levi with levi_logNormSq_zero_at, epunct_compact), ann eps for A, wdb_cut_bdd for M,
% h = indicator of ann of M / c, hormander_open, the local bound hL with K from exists_bound_of_continuousOn,
% ext_glue (L2DbarEstimates.ext on balls and the agreement through hgn), and zeros_null (extNull). The cutoff
% is HE.cut eps = the smooth bump around 0 with radii eps and 2 eps.
\end{proof}

\begin{lemma}\label{is-onto-share-counts-section-kronecker-solve}
Let $X$ consist of the data (1) to (6) of flag data of rank $r$ over $E$, with $n=|\Psi|r(r-1)/2>0$, use the
notation of Definitions~\ref{is-onto-share-counts-section-levi} and~\ref{is-onto-share-counts-section-gauge} and the
gauge $W_\lambda$ of Lemma~\ref{is-onto-share-counts-section-kronecker-gauge}, and write $|z|$ for the Euclidean norm of
$z\in\mathbf C^n$. Let $\lambda$ and $L$ be weights, $w_0,s_1,\dots,s_n$ bounded sections of weight $L$, and
$S=\sum_k\|s_k\|_L^2$. Let $x_0$ be an orthonormal point with $w_0(x_0)\ne0$, $s_k(x_0)=0$ for every $k$, and such
that the derivative at $0$ of $z\mapsto\big(s_k(c_{x_0}(z))/w_0(c_{x_0}(z))\big)_k$ is invertible. Let $\nu,m,\eta$ be
real numbers with $0\le\nu<m$ and $\eta>0$ such that, for every $z\in\Omega(x_0)$ and every $\xi\in\mathbf C^n$,
$\mathcal L^{x_0}_L(z;\xi)>0$ if $\xi\ne0$, and the number
$\mathcal C=\mathcal L^{x_0}_\lambda(z;\xi)-\mathcal L_{\log D_{x_0}}(z;\xi)-n\,\mathcal L^{x_0}_L(z;\xi)$ satisfies
\[
 m\,\mathcal L_{\rho_{x_0}}(z;\xi)\le\mathcal C\quad\text{and}\quad2\eta\,\mathcal L^{x_0}_L(z;\xi)\le\mathcal C.
\]
Then there are a real number $\varepsilon>0$, a continuous function $\chi$ on $\mathbf C^n$ with
$\chi(z)=1$ for $|z|\le\varepsilon$ and $\chi(z)=0$ for $|z|\geq2\varepsilon$, and a holomorphic function $F$ on
$\Omega(x_0)$, such that:
\begin{enumerate}
\item the closed ball of radius $2\varepsilon$ around $0$ lies in $\Omega(x_0)$, and $0$ is the only point $z$ of it with
  $s_k(c_{x_0}(z))=0$ for every $k$;
\item $\displaystyle\int_{\Omega(x_0)}|\chi(z)-F(z)|^2\,N_\lambda(c_{x_0}(z))^2\,e^{\nu\rho_{x_0}(z)}\,S(c_{x_0}(z))^{-n}\,
  D_{x_0}(z)\,dz<\infty$, the integrand being read as $0$ where $S$ vanishes;
\item $F(0)=1$; $F$ is square integrable for $W_\lambda$ and $\nu$; and for every compact set $C\subset\mathcal D_X$ the
  integral of $|F(z)|^2N_\lambda(c_{x_0}(z))^2D_{x_0}(z)$ over the $z\in\Omega(x_0)$ with $c_{x_0}(z)\in C\mathcal B^\Psi$ is
  finite.
\end{enumerate}
\end{lemma}
\leanref{Development/AbelianVarietyTheory/Proofs/TorusHaarMeasure/WeightContDiffOn.lean}{GlobalTheory.kronSolve}
\leanref{Development/AbelianVarietyTheory/Proofs/TorusHaarMeasure/LintegralLtTopOfCutoff.lean}{GlobalTheory.KronValue.value_and_integrability}
\begin{proof}
Write $c=c_{x_0}$, $\Omega=\Omega(x_0)$ and $\rho=\rho_{x_0}$, and for $z\in\Omega$ put $\sigma_k(z)=s_k(c(z))$ and
$\omega(z)=w_0(c(z))$. These functions are holomorphic on $\Omega$, as the sections are holomorphic on $\mathcal D_X$
and $c$ is polynomial. By Lemma~\ref{is-onto-share-counts-section-kronecker-solve-chart}, $\Omega$ is open and
connected and contains $0$, and $c(0)=x_0$.

\emph{The isolated zero.} The map $f=(\sigma_k/\omega)_k$ is holomorphic near $0$, as $\omega(0)\ne0$; $f(0)=0$, and its
derivative $\Lambda$ at $0$ is invertible. So $|\Lambda z|\geq\kappa|z|$ for some $\kappa>0$, and as
$f(z)=\Lambda z+o(|z|)$ there is $\delta>0$ with $f(z)\ne0$ for $0<|z|\le\delta$. Some coordinate $\sigma_{k_0}/\omega$
of $f$ does not vanish identically near $0$: otherwise $\Lambda=0$, which is not invertible as $n>0$. So $\sigma_{k_0}$
does not vanish identically near $0$, and by the identity theorem on the connected set $\Omega$ it does not vanish
identically on any nonempty open subset of $\Omega$. Let $\varrho>0$ be such that the ball of radius $\varrho$ around $0$
lies in $\Omega$, and put $\varepsilon=\min(\delta,\varrho)/4$. The closed ball $\bar B$ of radius $2\varepsilon$ lies in
$\Omega$, and if all $\sigma_k$ vanish at $z\in\bar B$, then $f(z)=0$ with $|z|\le\delta$, so $z=0$. This is (1).

\emph{The weight.} On $\Omega$ put $\Phi=\varphi_\lambda\circ c-\log D_{x_0}-n\,\varphi_L\circ c-\nu\rho$. It is smooth, as
the minors $d_i(c(z),\tau)$ are nonzero polynomials in the real and imaginary parts of $z$ there, and the Levi form
is linear: $\mathcal L_\Phi(z;\xi)=\mathcal C-\nu\,\mathcal L_\rho(z;\xi)$. Let $\xi\ne0$. Then
$\mathcal C\geq2\eta\,\mathcal L^{x_0}_L(z;\xi)>0$, and as $m>\nu\geq0$,
\[
 m\,\mathcal L_\Phi(z;\xi)=(m-\nu)\,\mathcal C+\nu\big(\mathcal C-m\,\mathcal L_\rho(z;\xi)\big)>0,
\]
so $\mathcal L_\Phi(z;\xi)>0$. The Levi form is continuous in $(z,\xi)$; its minimum over the compact set of the pairs
with $z\in\bar B$ and $|\xi|=1$ is a number $c_0>0$, and $\mathcal L_\Phi(z;\xi)\geq c_0|\xi|^2$ for $z\in\bar B$.

\emph{The exhaustion.} Put $e(z)=\rho(z)+\log(1+|z|^2)$ on $\Omega$. It is smooth, and $e\geq0$ as $\rho\geq0$ there
(Lemma~\ref{is-onto-share-counts-section-kronecker-solve-chart}). Its Levi form is that of $\rho$, which is $\geq0$
by the same lemma, plus that of $\log(1+|z|^2)=\log\big(|1|^2+\sum_j|z_j|^2\big)$, the logarithm of a sum of squares
of absolute values of holomorphic functions without common zero, which is $\geq0$. The sets $\{z\in\Omega:e(z)\le a\}$
are compact by the same lemma.

\emph{The solution.} Let $\chi=\chi_\varepsilon$ be the cutoff function of
Lemma~\ref{is-onto-share-counts-section-kronecker-solve-interpolation}; it is continuous, equal to $1$ for
$|z|\le\varepsilon$ and to $0$ for $|z|\geq2\varepsilon$. That lemma, for $U=\Omega$, $e$, the $\sigma_k$, $k_0$,
$\psi=\Phi$, $\varepsilon$ and $c_0$, gives a holomorphic function $F$ on $\Omega$ with
\[
 \int_\Omega|\chi(z)-F(z)|^2e^{-\Phi(z)}\Big(\sum_k|\sigma_k(z)|^2\Big)^{-n}dz<\infty .
\]
For $z\in\Omega$, $\|s_k\|_L(c(z))^2=|\sigma_k(z)|^2e^{-\varphi_L(c(z))}$, as $N_L^2=e^{-\varphi_L}$, so
$S(c(z))=e^{-\varphi_L(c(z))}\sum_k|\sigma_k(z)|^2$, and
\[
 N_\lambda(c(z))^2\,e^{\nu\rho(z)}\,S(c(z))^{-n}\,D_{x_0}(z)=e^{-\Phi(z)}\Big(\sum_k|\sigma_k(z)|^2\Big)^{-n}.
\]
So the integral in (2) is the integral above, and (2) holds.

\emph{The value at $0$.} Let $\Delta$, $z(\cdot)$ and $b(\cdot)$ be as in
Lemma~\ref{is-onto-share-counts-section-transfer-lu} for $x_0$; by part (3) of that lemma for $b=1$, $\Delta(c(z))\ne0$
and $z(c(z))=z$ for $z\in\Omega$. Let $\varrho'\le\varepsilon$ be such that the closed ball of radius $\varrho'$ lies in
$\Omega$, and put $u(y)=1-F(z(y))$ where $\Delta(y)\ne0$, and $u(y)=0$ otherwise; so $u(c(z))=1-F(z)$ for $|z|<\varrho'$.
There $\chi(z)=1$, so the integrand in (2) is $|u(c(z))|^2(\sum_k|\sigma_k(z)|^2)^{-n}$ times
$N_\lambda(c(z))^2e^{\nu\rho(z)}N_L(c(z))^{-2n}D_{x_0}(z)$. The last factor is continuous and positive on the closed ball,
so it is bounded below there by a positive number, and
$\int_{|z|<\varrho'}|u(c(z))|^2(\sum_k|\sigma_k(z)|^2)^{-n}dz<\infty$. By
Lemma~\ref{is-onto-share-counts-section-kronecker-solve-simple-zero} for $x_1=x_0$, $u(x_0)=0$. As $z(x_0)=0$,
$F(0)=1$.

\emph{Square integrability.} As the $s_k$ are bounded sections, $S\le M$ on $\mathcal D_X$ for some $M$. The zero set
of $\sigma_{k_0}$ in $\Omega$ has measure zero, as $\sigma_{k_0}$ does not vanish identically near any point of
$\Omega$; so $S(c(z))>0$, and $S(c(z))^{-n}\geq M^{-n}$, for almost every $z\in\Omega$. Put
$v(z)=N_\lambda(c(z))^2e^{\nu\rho(z)}$. As $|F|^2\le2|\chi|^2+2|\chi-F|^2$, for almost every $z\in\Omega$
\[
 |F(z)|^2v(z)\le2|\chi(z)|^2v(z)+2M^n|\chi(z)-F(z)|^2v(z)\,S(c(z))^{-n}.
\]
The first term is continuous on $\bar B$ and vanishes off $\bar B$, and $\bar B$ has finite measure for
$D_{x_0}(z)\,dz$; the second is integrable against $D_{x_0}(z)\,dz$ by (2). So
$\int_\Omega|F(z)|^2e^{\nu\rho(z)}N_\lambda(c(z))^2D_{x_0}(z)\,dz<\infty$. As $W_\lambda(c(z))=N_\lambda(c(z))^2$ for
$z\in\Omega$, $F$ is square integrable for $W_\lambda$ and $\nu$. Finally $N_\lambda(c(z))^2\le v(z)$, as $\nu\rho\geq0$;
so the integral over the $z\in\Omega$ with $c(z)\in C\mathcal B^\Psi$, for a compact set $C\subset\mathcal D_X$, is
finite too. This is (3).
% Lean: (1) and (2) are kronSolve with psi = 0: HE.iso_zero (isolated zero of f), HE.coord_ne_of_det (k0),
% HE.hgn_of with OmFConnStatement_holds, eps = min delta rho / 4; KronSolve.weight (Phi with the term psi),
% weight_contDiffOn and levi_weight; hWpos is the displayed inequality; HE.levi_unif for c0; eF with
% ExF.eF_contDiffAt, ExF.levi_rhF_nonneg, ExF.leviT_nonneg (H.leviLogNormSq), ExF.rhF_nonneg and
% ExF.isCompact_subF; HE.interp_euclid; HE.cut is the cutoff (cut_contDiff, cut_one, cut_zero); and
% KronSolve.integrand_eq (pnorm_sq) for the integrand. (3) is KronValue.value_and_integrability with delta =
% 2 eps: F 0 = 1 from VanishAtSStatement (GlobalTheory.vanishAtS_holds) for u = 1 - F o zOf (lu_holds),
% through lintegral_ball_lt_top; InAt from lintegral_lt_top_of_cutoff (with OntoD.null_zero); the compact-set
% clause over sat X x0 C with psi <= 0 and rhF_nonneg. Real.rpow gives 0^(-n) = 0, whence the reading of the
% integrand where S = 0. Nothing below this proof remains open: the chart facts are Lemma
% is-onto-share-counts-section-kronecker-solve-chart, now proved.
% Size: kronSolve is about 90 lines of Lean, value_and_integrability about 110 with 90 more in
% LintegralLtTopOfCutoff, and HE.interp_euclid with its helpers about 1,400 (InterpStatementPfC to PfJ), above
% the cited L2DbarEstimates.hormander and L2DbarEstimates.ext; the earlier count of 13,200 lines followed
% tactic and local names to unrelated declarations.
\end{proof}

\begin{lemma}\label{is-onto-share-counts-section-kronecker-extend-measure}
Let $X$ consist of the data (1) to (6) of flag data of rank $r$ over $E$, use the notation of
Definition~\ref{is-onto-share-counts-section-levi}, let $x_0$ be an orthonormal point and $C\subset\mathcal D_X$ a
compact set, and let $dx$ be the Lebesgue measure on $M_r(\mathbf C)^\Psi$. Then there is a real number $K$ such that
\[
 \int_CG(x)\,dx\le K\int_{\{z\in\Omega(x_0):\,c_{x_0}(z)\in C\mathcal B^\Psi\}}G(c_{x_0}(z))\,D_{x_0}(z)\,dz
\]
for every measurable function $G\colon M_r(\mathbf C)^\Psi\to[0,\infty]$ with $G(xb)=G(x)$ for $x\in\mathcal D_X$ and
$b\in\mathcal B^\Psi$.
\end{lemma}
\leanref{ReciprocityTheory/ReductiveGroups/ArithmeticGeometry/ArakelovVolumes/UpperBounds/FundSetMeasStatement.lean}{ReciprocityTheory.ChartRestriction.G.AR.Up.FrameMeasStatement}
\leanref{ReciprocityTheory/ReductiveGroups/ArithmeticGeometry/ArakelovVolumes/UpperBounds/FrameMeasStatementPf.lean}{ReciprocityTheory.ChartRestriction.G.AR.Up.FrameMeasStatement_holds}
\begin{proof}
This is a self-contained computation, through the change of variables $x=c_{x_0}(z)b$ of
Lemma~\ref{is-onto-share-counts-section-transfer-lu}.
% Lean: FrameMeasStatement_holds: the measurable equivalence ufC and its Haar factor, left multiplication by x0
% (lintegral_comp_LA, factor |prod det x0^r|^-2), MI.lintegral_flat (the factorization into lower unitriangular
% and upper triangular parts, Jacobian jacU), inner_le with the bound KR for the triangular parts over C
% (beta0_bound), and MI.lintegral_muX for the measure D_x0(z) dz; the points outside sat X x0 C contribute 0
% (Jp_eq_zero, mem_sat). Size: about 450 lines of Lean (FrameMeasStatementPf, PfA to PfC).
\end{proof}

\begin{lemma}\label{is-onto-share-counts-section-kronecker-extend}
Let $X$ consist of the data (1) to (6) of flag data of rank $r$ over $E$, use the notation of
Definition~\ref{is-onto-share-counts-section-levi}, and let $x_0$ be an orthonormal point and $\lambda$ a weight. Let $F$
be a holomorphic function on $\Omega(x_0)$ such that for every compact set $C\subset\mathcal D_X$ the integral of
$|F(z)|^2N_\lambda(c_{x_0}(z))^2D_{x_0}(z)$ over the $z\in\Omega(x_0)$ with $c_{x_0}(z)\in C\mathcal B^\Psi$ is finite. Then
there is a function $P$ that is $\mathcal B$-equivariant of weight $\lambda$
(Definition~\ref{is-onto-share-counts-section-gauge}) with $P\circ c_{x_0}=F$ on $\Omega(x_0)$.
\end{lemma}
\leanref{ReciprocityTheory/ReductiveGroups/ArithmeticGeometry/ArakelovVolumes/UpperBounds/FundSetMeasStatement.lean}{ReciprocityTheory.ChartRestriction.G.AR.Up.ExtFStatement}
\leanref{ReciprocityTheory/ReductiveGroups/ArithmeticGeometry/ArakelovVolumes/UpperBounds/KronHolPsiStatementHolds.lean}{ReciprocityTheory.ChartRestriction.G.AR.Up.extFStatement_holds}
\begin{proof}
Let $\Delta$, $z(\cdot)$ and $b(\cdot)$ be as in Lemma~\ref{is-onto-share-counts-section-transfer-lu} for $x_0$, and let
$\mathcal D'$ be the open set of the $x\in\mathcal D_X$ with $\Delta(x)\ne0$. By part (1) of that lemma and as
$\mathcal D_X$ is open, $\mathcal D'$ is dense in $\mathcal D_X$.

\emph{The function on $\mathcal D'$.} Put $P_0(x)=b(x)^\lambda F(z(x))$ for $x\in\mathcal D'$ and $P_0(x)=0$ otherwise.
By part (2), $P_0$ is holomorphic on $\mathcal D'$, as the diagonal entries of $b(x)$ do not vanish. For
$z\in\Omega(x_0)$, part (3) with $b=1$ gives $\Delta(c_{x_0}(z))\ne0$, $z(c_{x_0}(z))=z$ and $b(c_{x_0}(z))=1$, so
$P_0(c_{x_0}(z))=F(z)$. Let $x\in\mathcal D_X$ and $b\in\mathcal B^\Psi$. If $\Delta(x)\ne0$, then
$xb=c_{x_0}(z(x))\,b(x)b$, so by part (3) $\Delta(xb)\ne0$, $z(xb)=z(x)$ and $b(xb)=b(x)b$, and
$P_0(xb)=(b(x)b)^\lambda F(z(x))=b^\lambda P_0(x)$. If $\Delta(x)=0$, then $\Delta(xb)=0$, as otherwise the same
argument for $xb$ and $b^{-1}$ would give $\Delta(x)\ne0$; and both sides vanish. So $P_0(xb)=b^\lambda P_0(x)$.

\emph{Local square integrability.} Let $C\subset\mathcal D_X$ be compact, and $K$ the number of
Lemma~\ref{is-onto-share-counts-section-kronecker-extend-measure} for $C$. The function $G=|P_0|^2N_\lambda^2$ on
$\mathcal D_X$, and $0$ elsewhere, is measurable and satisfies $G(xb)=G(x)$, by the equivariance of $P_0$ and
$N_\lambda(xb)=|b^\lambda|^{-1}N_\lambda(x)$; and $G(c_{x_0}(z))=|F(z)|^2N_\lambda(c_{x_0}(z))^2$ for $z\in\Omega(x_0)$. As
$N_\lambda$ is continuous and positive on $\mathcal D_X$, it has a positive minimum $m_C$ on $C$. So
\[
 \int_C|P_0(x)|^2\,dx\le m_C^{-2}\int_CG(x)\,dx\le m_C^{-2}K\int_{\{z\in\Omega(x_0):\,c_{x_0}(z)\in C\mathcal B^\Psi\}}
 |F(z)|^2N_\lambda(c_{x_0}(z))^2D_{x_0}(z)\,dz<\infty .
\]

\emph{Local extensions.} Identify $M_r(\mathbf C)^\Psi$ with $\mathbf C^N$ by a linear isomorphism, which changes the
Lebesgue measure by a constant factor. For $x\in\mathcal D_X$ let $B_x$ be an open ball around $x$ whose closure lies
in $\mathcal D_X$. By the previous step for $C=\bar B_x$, $P_0$ is square integrable on $B_x\setminus\{\Delta=0\}$, where
it is holomorphic. Theorem~\ref{is-onto-share-counts-section-kronecker-solve-riemann}, for the polynomial $\Delta$,
gives a holomorphic function $U_x$ on $B_x$ equal to $P_0$ on $B_x\setminus\{\Delta=0\}$.

\emph{Gluing.} For $x,y\in\mathcal D_X$, $U_x$ and $U_y$ are continuous on $B_x\cap B_y$ and agree on its dense subset
$B_x\cap B_y\cap\mathcal D'$, so they agree on $B_x\cap B_y$. Put $P(y)=U_y(y)$ for $y\in\mathcal D_X$ and $P(y)=0$
otherwise. Near each $x\in\mathcal D_X$, $P=U_x$, so $P$ is holomorphic on $\mathcal D_X$, and $P=P_0$ on $\mathcal D'$.
In particular $P(c_{x_0}(z))=P_0(c_{x_0}(z))=F(z)$ for $z\in\Omega(x_0)$. For $b\in\mathcal B^\Psi$ the functions
$x\mapsto P(xb)$ and $x\mapsto b^\lambda P(x)$ are continuous on $\mathcal D_X$ and agree on the dense set $\mathcal D'$,
which is stable under $x\mapsto xb$; so they agree on $\mathcal D_X$, and $P$ is $\mathcal B$-equivariant of weight
$\lambda$.
% Lean: extFStatement_holds = extF_of lu_holds FrameMeasStatement_holds. luData_of gives Delta, zOf, bOf; P0 is
% ExtFPf.P0 (poincareLift_differentiableOn, poincareLift_on_chart through cpt_cell, poincareLift_borel_equivariant
% through cell_rmul and cell_of_rmul). local_ext: chartE identifies the frames with V N (lint_transfer for the
% constant), poincareLift_locally_squareIntegrable (FrameMeasStatement with G = Gf = ofReal (pnorm la P0 ^ 2),
% Gf_inv, Gf_cpt, the minimum m of nfac la on C), and L2DbarEstimates.ext. hagree is the gluing through
% Set.EqOn.of_subset_closure and D.dense; dense_cell for the equivariance. Size: about 450 lines of Lean
% (ExtFStatementPf, PfA, PfB), with the L^2 extension theorem below it.
\end{proof}

\begin{lemma}\label{is-onto-share-counts-section-kronecker-vanish}
Let $X$ consist of the data (1) to (6) of flag data of rank $r$ over $E$, with $n=|\Psi|r(r-1)/2>0$, use the
notation of Definitions~\ref{is-onto-share-counts-section-levi} and~\ref{is-onto-share-counts-section-gauge}, and write
$|z|$ for the Euclidean norm of $z\in\mathbf C^n$. Let $\lambda$ and $L$ be weights, $w_0,s_1,\dots,s_n$ sections of
weight $L$, and $S=\sum_k\|s_k\|_L^2$. Let $x_0$ and $x_1$ be orthonormal points, $\nu\geq0$ and $\delta$ real numbers
such that the closed ball of radius $\delta$ around $0$ lies in $\Omega(x_0)$, $\chi$ a function on $\mathbf C^n$ with
$\chi(z)=0$ for $|z|\geq\delta$, $P$ a function that is $\mathcal B$-equivariant of weight $\lambda$, and $F$ a
function on $\Omega(x_0)$ with $P\circ c_{x_0}=F$ on $\Omega(x_0)$, such that
\[
 \int_{\Omega(x_0)}|\chi(z)-F(z)|^2\,N_\lambda(c_{x_0}(z))^2\,e^{\nu\rho_{x_0}(z)}\,S(c_{x_0}(z))^{-n}\,D_{x_0}(z)\,dz<\infty,
\]
the integrand being read as $0$ where $S$ vanishes. Suppose that $x_1\ne c_{x_0}(z)b$ for every $z$ with $|z|\le\delta$
and every $b\in\mathcal B^\Psi$, that $w_0(x_1)\ne0$, that $s_k(x_1)=0$ for every $k$, and that the derivative at $0$
of $z\mapsto\big(s_k(c_{x_1}(z))/w_0(c_{x_1}(z))\big)_k$ is invertible. Then $P(x_1)=0$.
\end{lemma}
\leanref{Development/AbelianVarietyTheory/Proofs/TorusHaarMeasure/MeasurableIndicatorOfContOn.lean}{GlobalTheory.KronOff.vanish_off}
\leanref{Development/AbelianVarietyTheory/Proofs/TorusHaarMeasure/MeasurableIndicatorOfContOn.lean}{GlobalTheory.KronOff.isClosed_satBall}
\leanref{Development/AbelianVarietyTheory/Proofs/TorusHaarMeasure/MeasurableIndicatorOfContOn.lean}{GlobalTheory.KronOff.eq_of_mem_satBall}
\begin{proof}
\emph{A neighbourhood of $x_1$.} Let $T$ be the set of the $y\in M_r(\mathbf C)^\Psi$ for which some $z$ with $|z|\le\delta$
makes $\operatorname{adj}(c_{x_0}(z)_\tau)\,y_\tau$ upper triangular for every $\tau$, where $\operatorname{adj}$ is the
adjugate matrix. It is closed: it is the image under the projection to $M_r(\mathbf C)^\Psi$ of a closed set of pairs
$(z,y)$ with $|z|\le\delta$, and the projection along the compact ball is a closed map. As upper triangular matrices
form a group, $yb\in T$ if and only if $y\in T$, for $b\in\mathcal B^\Psi$. Each $c_{x_0}(z)$ with $|z|\le\delta$ lies in
$T$, as $\operatorname{adj}(c)\,c=\det(c)\cdot1$. If $y\in\mathcal D_X\cap T$, with $z$ as in the definition, then
$c_{x_0}(z)\in\mathcal D_X$, so $\det c_{x_0}(z)_\tau\ne0$, and $b_\tau=c_{x_0}(z)_\tau^{-1}y_\tau$ is upper triangular and
invertible; so $y=c_{x_0}(z)b$ with $b\in\mathcal B^\Psi$. Hence the open set $V=\mathcal D_X\setminus T$ contains
$x_1$, which lies in $\mathcal D_X$ by Lemma~\ref{is-onto-share-majorant}(3) and is not of the form $c_{x_0}(z)b$ with
$|z|\le\delta$; $yb\in V$ if and only if $y\in V$, for
$y\in\mathcal D_X$; and $c_{x_0}(z)\notin V$ for $|z|\le\delta$.

\emph{An invariant function.} Put $\phi(y)=\|P\|_\lambda(y)^2\big(\sum_k\|s_k\|_L(y)^2\big)^{-n}$ for $y\in V$, read as $0$
where the sum vanishes, and $\phi(y)=0$ for $y\notin V$. It is measurable, and $\phi(yb)=\phi(y)$ for $y\in\mathcal D_X$
and $b\in\mathcal B^\Psi$: $|P(yb)|=|b^\lambda||P(y)|$ and $N_\lambda(yb)=|b^\lambda|^{-1}N_\lambda(y)$, so
$\|P\|_\lambda(yb)=\|P\|_\lambda(y)$, and in the same way $\|s_k\|_L(yb)=\|s_k\|_L(y)$.

\emph{The integral at $x_0$.} Let $z\in\Omega(x_0)$ with $c_{x_0}(z)\in V$. Then $|z|>\delta$, so $\chi(z)=0$, and
$|F(z)|N_\lambda(c_{x_0}(z))=\|P\|_\lambda(c_{x_0}(z))$. As $\nu\geq0$ and $\rho_{x_0}\geq0$ on $\Omega(x_0)$
(Lemma~\ref{is-onto-share-counts-section-kronecker-solve-chart}),
\[
 \phi(c_{x_0}(z))\le|\chi(z)-F(z)|^2N_\lambda(c_{x_0}(z))^2e^{\nu\rho_{x_0}(z)}S(c_{x_0}(z))^{-n}.
\]
So $\int_{\Omega(x_0)}\phi(c_{x_0}(z))D_{x_0}(z)\,dz<\infty$, and by
Lemma~\ref{is-onto-share-counts-section-poincare-measure}(1), $\int_{\Omega(x_1)}\phi(c_{x_1}(z))D_{x_1}(z)\,dz<\infty$.

\emph{The integral at $x_1$.} The set of $z$ with $c_{x_1}(z)\in V$ is open and contains $0$; let $\varrho>0$ be such
that it contains the closed ball $\bar B$ of radius $\varrho$. On $\bar B$ the function
$N_\lambda(c_{x_1}(z))^2N_L(c_{x_1}(z))^{-2n}D_{x_1}(z)$ is continuous and positive, so it is at least some $m>0$. As
$\|s_k\|_L^2=|s_k|^2N_L^2$, for $|z|<\varrho$
\[
 \phi(c_{x_1}(z))\,D_{x_1}(z)\geq m\,|P(c_{x_1}(z))|^2\Big(\sum_k|s_k(c_{x_1}(z))|^2\Big)^{-n}.
\]
Hence $\int_{|z|<\varrho}|P(c_{x_1}(z))|^2(\sum_k|s_k(c_{x_1}(z))|^2)^{-n}dz<\infty$. The composites of $P$, $w_0$ and the
$s_k$ with $c_{x_1}$ are holomorphic on the ball $|z|<\varrho$, which $c_{x_1}$ maps into $\mathcal D_X$. So
Lemma~\ref{is-onto-share-counts-section-kronecker-solve-simple-zero} applies to $u=P$, and $P(x_1)=0$.
% Lean: KronOff.vanish_off with psi = 0: U = X.dom minus satBall X x0 delta (isClosed_satBall through
% isClosedMap_snd_of_compactSpace, rmul_mem_satBall_iff, cpt_mem_satBall, eq_of_mem_satBall for x1 in U), phi
% from the indicators of pnorm (measurable_indicator_of_contOn) with hphiinv (NfacRmulStatement and
% PnormInvStatement with gamma = 1), hL through rhF_nonneg and kaF <= 1, hR through MeasInvStatement
% (measInv_holds), the minimum m on the closed ball (exists_isMinOn), key, and VanishAtSStatement
% (GlobalTheory.vanishAtS_holds). The Lean's hypotheses are those of the statement (with psi continuous on
% X.dom and invariant); it does not use that F is holomorphic. Size: about 230 lines of Lean in
% MeasurableIndicatorOfContOn above the statements it cites.
\end{proof}

\begin{lemma}\label{is-onto-share-counts-section-kronecker}
Let $X$ consist of the data (1) to (6) of flag data of rank $r$ over $E$, with archimedean data for $(L,F,T)$ and a
flat model for them, with base points $y_1,\dots,y_\delta$ (Definition~\ref{is-onto-share-counts-model}), and use
the notation of Definitions~\ref{is-onto-share-counts-section-levi} and~\ref{is-onto-share-counts-section-gauge}.
Let $\lambda$ be a weight, $\nu,m,\eta$ curvature constants for $\lambda$ and $L$, and $1\le j\le\delta$. Then there are
an orthonormal point $x_0$, a gauge $W$ of weight $\lambda$ that is bounded below, a holomorphic function $f$ on
$\Omega(x_0)$ that is square integrable for $W$ and $\nu$, and a function $P$ that is $\mathcal B$-equivariant of
weight $\lambda$, such that $P\circ c_{x_0}=f$ on $\Omega(x_0)$, $P(y_j)\ne0$, and $P(\gamma y_{j'})=0$ for all
$\gamma\in\Gamma$ and $1\le j'\le\delta$ with $\gamma\ne1$ or $j'\ne j$.
\end{lemma}
\leanref{Development/AbelianVarietyTheory/Proofs/TorusHaarMeasure/InterpWeight.lean}{GlobalTheory.InterpWeight.core}
\leanref{Development/AbelianVarietyTheory/Proofs/DivisorProofs/LoweredWeightInterpolationAtPoints.lean}{GlobalTheory.LowerInterp.kron_at}
\begin{proof}
\emph{The point $x_0$.} By Lemma~\ref{is-onto-share-majorant}(3), $y_j=x_0b_0$ for an orthonormal point $x_0$ and some
$b_0\in\mathcal B^\Psi$. As $x_0=y_jb_0^{-1}$, Lemma~\ref{is-onto-share-counts-section-kronecker-points} gives
$w_0(x_0)\ne0$, $s_k(x_0)=0$ for every $k$, and the invertibility of the derivative at $0$ of
$z\mapsto\big(s_k(c_{x_0}(z))/w_0(c_{x_0}(z))\big)_k$.

\emph{The construction.} Let $W=W_\lambda$, a gauge of weight $\lambda$ that is bounded below
(Lemma~\ref{is-onto-share-counts-section-kronecker-gauge}). Lemma~\ref{is-onto-share-counts-section-kronecker-solve}
applies to $\lambda$, $L$, the cusp forms $w_0,s_1,\dots,s_n$ of the flat model, which are bounded sections of weight
$L$, and $x_0$: $n\geq2$ by (1) of Definition~\ref{is-onto-share-counts-model}; $0\le\nu<m$ and $\eta>0$, and the two
inequalities hold at $x_0$, as $\nu,m,\eta$ are curvature constants; and $\mathcal L^{x_0}_L(z;\xi)>0$ for
$z\in\Omega(x_0)$ and $\xi\ne0$ by (4) of Definition~\ref{is-onto-share-counts-model} for $y=x_0$, as
$\varphi_L=-2\log N_L$. It gives $\varepsilon$, $\chi$ and a holomorphic $f=F$ on $\Omega(x_0)$ with $f(0)=1$, square
integrable for $W$ and $\nu$, and with the integrability on the sets $C\mathcal B^\Psi$.
Lemma~\ref{is-onto-share-counts-section-kronecker-extend}
gives $P$, $\mathcal B$-equivariant of weight $\lambda$, with $P\circ c_{x_0}=f$ on $\Omega(x_0)$. As $c_{x_0}(0)=x_0$,
$P(x_0)=f(0)=1$, and $P(y_j)=P(x_0b_0)=b_0^\lambda\ne0$.

\emph{The other translates.} Let $\gamma\in\Gamma$ and $1\le j'\le\delta$ with $\gamma\ne1$ or $j'\ne j$. By
Lemma~\ref{is-onto-share-majorant}(3), $\gamma y_{j'}=x_1b_1$ for an orthonormal point $x_1$ and some
$b_1\in\mathcal B^\Psi$, and Lemma~\ref{is-onto-share-counts-section-kronecker-points} applies to
$x_1=\gamma y_{j'}b_1^{-1}$. Suppose that $x_1=c_{x_0}(z)b$ with $|z|\le2\varepsilon$ and $b\in\mathcal B^\Psi$. Then
$0=s_k(x_1)=b^Ls_k(c_{x_0}(z))$, so $s_k(c_{x_0}(z))=0$ for every $k$, and $z=0$ by
Lemma~\ref{is-onto-share-counts-section-kronecker-solve}(1). So $x_1=x_0b=y_jb_0^{-1}b$, and $\gamma y_{j'}=y_jb'$ with
$b'=b_0^{-1}bb_1$. By the flat model, $y_{j'}=\gamma^{-1}y_jb'$ forces $j'=j$; then $\gamma y_j=y_jb'$ forces $\gamma=1$, by
(4) of Definition~\ref{is-onto-share-flag-data}, a contradiction. Hence
Lemma~\ref{is-onto-share-counts-section-kronecker-vanish} applies, for $\delta=2\varepsilon$, $\chi$, $F=f$, $P$ and
$x_1$, with the closed ball of Lemma~\ref{is-onto-share-counts-section-kronecker-solve}(1) and the integral of its part
(2). So $P(x_1)=0$, and $P(\gamma y_{j'})=P(x_1b_1)=b_1^\lambda P(x_1)=0$.
% Lean: core (with bb, m, eta given) and kron_at (the same construction for a weight with an invariant
% function psi added, used with psi = 0; it returns gamma0 = 1). Both: KronAssembly.exists_orth for x0, the
% key facts at every translate of a base point (KronAssembly.sec_ne_zero, Fl.S_eq and
% OrbitCoord.det_fderiv_ne_zero_orbit from Fl.coord), KronMetric.kronMetric for W = N_lambda^2 on D_X (and 0
% off it), kronSolve for the cutoff chi and the solution F (with the weight
% e^{-phla} e^{bb rho} (sum_k ||s_k||_L^2)^{-n}), KronValue.value_and_integrability for F(0) = 1 and the square
% integrability, ExtFStatement for P, and KronOff.vanish_off (with Fl.xs_inj and X.Gamma_free) for the
% vanishing. The book's x0 and b0 are those of exists_orth for gamma = 1 (hx0, hx0'). kronSolve is called with
% hmar x0 ho and hstr x0 ho (the curvature conditions at x0), Xa.posL x0 for the positivity, Xa.n_ge for n >= 2
% and bb >= 0 from 4 |Psi| r^2 < bb. Nothing below this proof remains open; Lemma
% is-onto-share-counts-section-kronecker-solve-chart (about 1,200 lines of Lean) is now proved.
\end{proof}

\subsection{Poincaré series and the ranks of cusp forms}\label{subsection-is-poincare-ranks}

\begin{lemma}\label{is-onto-share-counts-section-discontinuous}
Let $X$ consist of the data (1) to (6) of flag data of rank $r$ over $E$. Then $\Gamma$ acts \emph{properly
discontinuously} on $\mathcal D_X$ modulo $\mathcal B^\Psi$: for $x,x'\in\mathcal D_X$ there are neighbourhoods $U$
of $x$ and $U'$ of $x'$ in $M_r(\mathbf C)^\Psi$ such that only finitely many $\gamma\in\Gamma$ satisfy $\gamma u=u'b$
for some $u\in U$, $u'\in U'$ and $b\in\mathcal B^\Psi$.
\end{lemma}
\leanref{ReciprocityTheory/ReductiveGroups/ArithmeticGeometry/ArakelovVolumes/ArakelovInstances/QuotientIsHausdorff.lean}{ReciprocityTheory.ChartRestriction.G.AR.Inst.PropDisc}
\leanref{Development/AbelianVarietyTheory/Proofs/TorusHaarMeasure/FirstMinimumInvariantContinuousPosOnFrames.lean}{GlobalTheory.ClosedDiscontinuity.propDisc_closed}
\begin{proof}
For $w\in\mathbf C^r$ write $|w|=\max_a|w_a|$.

\emph{Bounds near a point.} Let $y\in\mathcal D_X$. By Lemma~\ref{is-onto-share-majorant}(1) each $m_{y,\tau}$ is a
positive definite Hermitian form, so there are $0<c\le C$ with $c|w|^2\le m_{y,\tau}(w)\le C|w|^2$ for all $\tau$ and
$w$. The coefficients of $m_{y',\tau}$ depend continuously on $y'\in\mathcal D_X$, by the formula of
Definition~\ref{is-onto-share-flag-data}, and $\mathcal D_X$ is open; so there is a neighbourhood of $y$ in
$\mathcal D_X$ on which $m_{y,\tau}\le2m_{y',\tau}$ and $m_{y',\tau}\le2m_{y,\tau}$ for every $\tau$. On it
$\frac c2|w|^2\le m_{y',\tau}(w)\le2C|w|^2$. Choose such neighbourhoods $U$ of $x$ and $U'$ of $x'$, with the
constants $c,C$ taken at $x$ for $U$ and $c',C'$ taken at $x'$ for $U'$.

\emph{Invariance.} Let $\gamma\in\Gamma$, $u\in U$, $u'\in U'$ and $b\in\mathcal B^\Psi$ with $\gamma u=u'b$. As
$\gamma$ preserves $h$, the Gram matrix of the columns of $\gamma^\tau u_\tau$ is that of $u_\tau$, and
$h_\tau(\gamma^\tau w,\gamma^\tau v)=h_\tau(w,v)$; so the formula for the majorant gives
$m_{\gamma u,\tau}(\gamma^\tau w)=m_{u,\tau}(w)$. With Lemma~\ref{is-onto-share-majorant}(2),
$m_{u',\tau}(\gamma^\tau w)=m_{u'b,\tau}(\gamma^\tau w)=m_{u,\tau}(w)$. Hence
$\frac{c'}2|\gamma^\tau w|^2\le2C|w|^2$, that is, $|\gamma^\tau w|\le K|w|$ for every $\tau$ and $w$, where
$K=2(C/c')^{1/2}$ depends only on $U$ and $U'$.

\emph{Finiteness.} Let $S$ be a finite set that generates $\Lambda$ as an $\mathcal O_E$-module, and put
$R=K\sum_{s\in S}\sum_{\tau\in\Psi}|\tau(s)|$. For $\gamma$ as above and $s\in S$, $\gamma s\in\Lambda$, as $\gamma$
preserves $\Lambda$, and $|\tau(\gamma s)|=|\gamma^\tau\tau(s)|\le K|\tau(s)|\le R$ for every $\tau\in\Psi$. Let $F$ be
the set of $v\in\Lambda$ with $|\tau(v)|\le R$ for every $\tau\in\Psi$. It is finite. Indeed, there is $y\ne0$ in
$\mathcal O_E$ with $yv_a\in\mathcal O_E$ for every $v\in\Lambda$ and every $a$, namely a common denominator of the
coordinates of the elements of $S$. Every embedding $\sigma\colon E\to\mathbf C$ is $\tau$ or $\bar\tau$ for some
$\tau\in\Psi$, by (1) of Definition~\ref{is-onto-share-flag-data}. So for $v\in F$ each $yv_a$ is an algebraic
integer of $E$ with $|\sigma(yv_a)|\le R\sum_\psi|\psi(y)|$ for every embedding $\sigma$, the sum running over all
embeddings $\psi$. There are only finitely many such algebraic integers, as the coefficients of their characteristic
polynomials are integers bounded in terms of $R\sum_\psi|\psi(y)|$ and $[E:\mathbf Q]$.

As $\Lambda$ spans $E^r$, so does $S$, and $\gamma\in GL_r(E)$ is determined by the family $(\gamma s)_{s\in S}$, which
lies in $F^S$. So only finitely many $\gamma\in\Gamma$ satisfy $\gamma u=u'b$ for some $u\in U$, $u'\in U'$ and
$b\in\mathcal B^\Psi$.
% Lean: ClosedDiscontinuity.propDisc_closed = propDisc_arithmetic_group_frames: exists_nhds_maj_two_sided_norm_sq_bound
% (MajorantContinuity.compact_bounds at {x} and eventually_comparable with eps = 1), MajInv.maj_rmul and
% maj_lact for the invariance, K = sqrt(|C| / c') with le_sqrt_mul_of_sq_le_mul_sq (the Lean's constants are
% c / 2 and 2 C), finite_lattice_norm_emb_le (exists_common_denominator_lattice, finite_denominator_norm_le
% with NumberField.Embeddings.finite_of_norm_le, mem_Psi_or_conjugate_mem), and
% matrix_eq_of_mulVec_eq_on_span_generators for the injectivity of gamma -> (gamma s)_s. Size: about 200 lines
% of Lean; the earlier count of 8,700 lines followed the tactic name ext to an unrelated declaration.
\end{proof}

\begin{lemma}\label{is-onto-share-counts-section-poincare-measure}
Let $X$ consist of the data (1) to (6) of flag data of rank $r$ over $E$, and use the notation of
Definition~\ref{is-onto-share-counts-section-levi}. Let $x_0$ be an orthonormal point and $\phi\colon M_r(\mathbf
C)^\Psi\to[0,\infty]$ a measurable function with $\phi(xb)=\phi(x)$ for $x\in\mathcal D_X$ and $b\in\mathcal B^\Psi$.
\begin{enumerate}
\item For every orthonormal point $x_0'$,
  \[
   \int_{\Omega(x_0)}\phi(c_{x_0}(z))\,D_{x_0}(z)\,dz=\int_{\Omega(x_0')}\phi(c_{x_0'}(z))\,D_{x_0'}(z)\,dz .
  \]
\item Let $\Phi\subset\Omega(x_0)$ be a measurable set such that $z=z'$ whenever $z,z'\in\Phi$ satisfy
  $c_{x_0}(z')=\gamma c_{x_0}(z)b$ for some $\gamma\in\Gamma$ and $b\in\mathcal B^\Psi$, and such that for almost every
  $z\in\Omega(x_0)$ there are $z'\in\Phi$, $\gamma\in\Gamma$ and $b\in\mathcal B^\Psi$ with $c_{x_0}(z)=\gamma
  c_{x_0}(z')b$. Then
  \[
   \sum_{\gamma\in\Gamma}\int_\Phi\phi(\gamma c_{x_0}(z))\,D_{x_0}(z)\,dz=\int_{\Omega(x_0)}\phi(c_{x_0}(z))\,D_{x_0}(z)\,dz .
  \]
\end{enumerate}
\end{lemma}
\leanref{ReciprocityTheory/ReductiveGroups/ArithmeticGeometry/ArakelovVolumes/UpperBounds/FundSetMeasStatement.lean}{ReciprocityTheory.ChartRestriction.G.AR.Up.MeasInvStatement}
\leanref{ReciprocityTheory/ReductiveGroups/ArithmeticGeometry/ArakelovVolumes/UpperBounds/MeasInvStatementPf.lean}{ReciprocityTheory.ChartRestriction.G.AR.Up.measInv_holds}
\leanref{ReciprocityTheory/ReductiveGroups/ArithmeticGeometry/ArakelovVolumes/UpperBounds/PairUnfoldPfB.lean}{ReciprocityTheory.ChartRestriction.G.AR.Up.PairUnfoldPf.lunfold}
\begin{proof}
(1) is a self-contained computation: the measures $D_{x_0}(z)\,dz$ on the charts $\Omega(x_0)$ at orthonormal points
are the readings of one measure on $\mathcal D_X$ modulo $\mathcal B^\Psi$.

(2) For $\gamma\in\Gamma$ the point $\gamma x_0$ is orthonormal, as $\gamma$ preserves $h$, and
$\gamma c_{x_0}(z)=c_{\gamma x_0}(z)$; so $\Omega(\gamma x_0)=\Omega(x_0)$ and $D_{\gamma x_0}=D_{x_0}$. Let $T$ be the set of
points $c_{x_0}(z)b$ with $z\in\Phi$ and $b\in\mathcal B^\Psi$, which is measurable by
Lemma~\ref{is-onto-share-counts-section-transfer-lu}, and $\mathbf 1_T$ its indicator; it is invariant under
$\mathcal B^\Psi$, and $c_{x_0}(z)\in T$ only for $z\in\Phi$, by part (3) of that lemma. Applying (1) to $x_0$, $\gamma x_0$
and the function $y\mapsto\mathbf 1_T(\gamma^{-1}y)\phi(y)$ gives
\[
 \int_\Phi\phi(\gamma c_{x_0}(z))\,D_{x_0}(z)\,dz=\int_{\Omega(x_0)}\mathbf 1_T(\gamma^{-1}c_{x_0}(z))\,\phi(c_{x_0}(z))\,
 D_{x_0}(z)\,dz .
\]
For almost every $z$ there is exactly one $\gamma$ with $\gamma^{-1}c_{x_0}(z)\in T$: one by the covering condition, and
at most one by the injectivity condition. Summing over $\gamma$ gives (2).
% Lean: (1) is MeasInvStatement, proved by measInv_holds (MI.lintegral_cpt_muX_indep_of_isOrtho); the measure
% D_x0(z) dz on Omega(x0) is muF X x0. (2) is PairUnfoldPf.lunfold (lunfold_one with T D FF, ind_rmul, ind_cpt,
% muF_lact_upper_bounds, cpt_lact_upper_bounds, and the count of eq_of_lact_inv_cpt_mem_T). Size: about 250
% lines of Lean for (1) and 600 for (2).
\end{proof}

\begin{lemma}\label{is-onto-share-counts-section-poincare-mean}
Let $X$ consist of the data (1) to (6) of flag data of rank $r$ over $E$, and use the notation of
Definition~\ref{is-onto-share-counts-section-levi}. Let $\lambda$ be a weight and $0<\varrho\le1$, write $|w|$ for the
Euclidean norm of $w\in\mathbf C^n$, and for $y\in M_r(\mathbf C)^\Psi$ let $B_\varrho(y)$ be the set of the $w$ with
$|w|<\varrho$ and $c_y(w)\in\mathcal D_X$.
\begin{enumerate}
\item There is a real number $C_A$ such that for every holomorphic function $g$ on $\mathcal D_X$ and every orthonormal
  point $x$
  \[
   |g(x)|^2\le C_A\int_{B_\varrho(x)}|g(c_x(w))|^2\,N_\lambda(c_x(w))^2\,dw .
  \]
\item For every orthonormal point $x$ there are real numbers $\varepsilon>0$ and $C_A$ such that for every function
  $P$ that is $\mathcal B$-equivariant of weight $\lambda$ (Definition~\ref{is-onto-share-counts-section-gauge}) and
  every $y$ at distance less than $\varepsilon$ from $x$
  \[
   |P(y)|^2\le C_A\int_{B_\varrho(x)}|P(c_x(w))|^2\,N_\lambda(c_x(w))^2\,dw .
  \]
\end{enumerate}
\end{lemma}
\leanref{ReciprocityTheory/ReductiveGroups/ArithmeticGeometry/ArakelovVolumes/UpperBounds/PoinStatementPfD.lean}{ReciprocityTheory.ChartRestriction.G.AR.Up.PoinPf.subMean_holds}
\leanref{ReciprocityTheory/ReductiveGroups/ArithmeticGeometry/ArakelovVolumes/UpperBounds/PoinStatementPfE.lean}{ReciprocityTheory.ChartRestriction.G.AR.Up.PoinPf.subMeanLoc_of}
\leanref{ReciprocityTheory/L2DbarEstimates/WeakHolomorphy/MeasurableEmbeddingRotAbout.lean}{ReciprocityTheory.L2DbarEstimates.C.mv}
\begin{proof}
\emph{Mean values.} If $u$ is holomorphic on a ball $B$ in $\mathbf C^n$ with centre $c$, then $u(c)$ is the mean of $u$
over $B$, and by the Cauchy-Schwarz inequality $\operatorname{vol}(B)\,|u(c)|^2\le\int_B|u|^2$.

\emph{A uniform ball.} Let $x$ be orthonormal. The Gram matrix of the columns of $c_x(w)_\tau=x_\tau\ell_\tau(w)$ is
$\ell_\tau(w)^{\mathsf T}J\overline{\ell_\tau(w)}$, which does not depend on $x$; at $w=0$ its leading minors are the
$\epsilon_i$, and $\det\ell_\tau(w)=1$. So the conditions $c_x(w)\in\mathcal D_X$ and $N_\lambda(c_x(w))>\frac12$ do not
depend on the orthonormal point $x$, and they hold at $w=0$. By continuity there is $R_0>0$ such that they hold for
$|w|<R_0$. Put $R=\min(R_0,\varrho)$; then the ball $B(0,R)$ lies in $B_\varrho(x)$, and $4N_\lambda(c_x(w))^2\geq1$ on it.

(1) Apply the mean values to $u=g\circ c_x$ on $B(0,R)$:
\[
 |g(x)|^2\le\frac1{\operatorname{vol}B(0,R)}\int_{B(0,R)}|g(c_x(w))|^2\,dw\le
 \frac4{\operatorname{vol}B(0,R)}\int_{B_\varrho(x)}|g(c_x(w))|^2N_\lambda(c_x(w))^2\,dw .
\]

(2) Let $\Delta$, $z(\cdot)$ and $b(\cdot)$ be as in Lemma~\ref{is-onto-share-counts-section-transfer-lu} for $x$. By part
(3) of that lemma, $\Delta(x)\ne0$, $z(x)=0$ and $b(x)=1$. By continuity there is $\varepsilon>0$ such that every $y$ at
distance less than $\varepsilon$ from $x$ lies in $\mathcal D_X$ and satisfies $\Delta(y)\ne0$, $|z(y)|<R/2$ and
$|b(y)^\lambda|<2$. For such $y$, $y=c_x(z(y))b(y)$, so $P(y)=b(y)^\lambda P(c_x(z(y)))$, and the ball $B(z(y),R/2)$ lies in
$B(0,R)$. The mean values on $B(z(y),R/2)$ give, with $v=\operatorname{vol}B(0,R/2)$,
\[
 |P(y)|^2\le4\,|P(c_x(z(y)))|^2\le\frac4v\int_{B(0,R)}|P(c_x(w))|^2\,dw\le\frac{16}v\int_{B_\varrho(x)}|P(c_x(w))|^2
 N_\lambda(c_x(w))^2\,dw .
\]
% Lean: (1) is PoinPf.subMean_holds: good_ball (nfac_cpt_orth, dmin_cpt_orth, nfw), the mean value inequality
% L2DbarEstimates.C.mv with weight 1, and C_A = 4 / vol B(0, R). (2) is PoinPf.subMeanLoc_of lu_holds
% (ball_int_bound, bchar_contAt), with C_A = 4 v^-1 4. exp(-phla) is N_lambda^2.
\end{proof}

\begin{lemma}\label{is-onto-share-counts-section-poincare-count-separation}
Let $X$ consist of the data (1) to (6) of flag data of rank $r$ over $E$, such that $\Gamma$ acts properly
discontinuously on $\mathcal D_X$ modulo $\mathcal B^\Psi$ (Lemma~\ref{is-onto-share-counts-section-discontinuous}),
and let $x_0$ be an orthonormal point. Then there is a real number $c>0$ such that for every $\gamma\in\Gamma$ with
$\gamma\ne1$ there are $\tau\in\Psi$ and $1\le a,b\le r$ with
$\big|\big(x_{0,\tau}^{-1}(\gamma x_0)_\tau\big)_{ab}-\delta_{ab}\big|\geq c$.
\end{lemma}
\leanref{ReciprocityTheory/ReductiveGroups/ArithmeticGeometry/ArakelovVolumes/UpperBounds/PoinStatementPfF.lean}{ReciprocityTheory.ChartRestriction.G.AR.Up.PoinPf.sep}
\begin{proof}
For $y,y'\in M_r(\mathbf C)^\Psi$ let $\operatorname{dist}(y,y')$ be the largest absolute value of an entry of
$y-y'$. The point $x_0$ lies in $\mathcal D_X$ (Lemma~\ref{is-onto-share-majorant}(3)). Proper discontinuity at $x_0$
and $x_0$ gives neighbourhoods $U$ and $U'$ of $x_0$ such that only finitely many $\gamma\in\Gamma$ satisfy
$\gamma u=u'b$ for some $u\in U$, $u'\in U'$ and $b\in\mathcal B^\Psi$. Let $\varepsilon>0$ be such that $U'$ contains
every $y$ with $\operatorname{dist}(y,x_0)<\varepsilon$.

Let $\gamma\in\Gamma$ with $\gamma\ne1$. Then $\gamma x_0\ne x_0$, by (4) of Definition~\ref{is-onto-share-flag-data}
with $b=1$, so $\operatorname{dist}(\gamma x_0,x_0)>0$. If $\operatorname{dist}(\gamma x_0,x_0)<\varepsilon$, then
$\gamma x_0=u'\cdot1$ with $u=x_0\in U$ and $u'=\gamma x_0\in U'$, so $\gamma$ is one of the finitely many elements
above. Let $\varepsilon'>0$ be the minimum of $\varepsilon$ and of the numbers $\operatorname{dist}(\gamma x_0,x_0)$ for
these finitely many $\gamma\ne1$. Then $\operatorname{dist}(\gamma x_0,x_0)\geq\varepsilon'$ for every $\gamma\ne1$
in $\Gamma$.

Let $K_0>0$ bound the absolute values of the entries of $x_0$, and put $c=\varepsilon'/(rK_0+1)$. Let $\gamma\ne1$ and
$g_\tau=x_{0,\tau}^{-1}(\gamma x_0)_\tau$. If every entry of every $g_\tau-1$ had absolute value less than $c$, then
every entry of $(\gamma x_0)_\tau-x_{0,\tau}=x_{0,\tau}(g_\tau-1)$ would have absolute value less than
$rK_0c<\varepsilon'$, a contradiction.
% Lean: PoinPf.sep: PropDisc X at x0, x0 (hfin), Gamma_free with idF for gamma x0 != x0, the finite set T and
% its minimum eps1 (exists_min_image), eps' = min eps eps1, K0 = (sum of the norms of the entries of x0) + 1, and
% c = eps' / (r K0 + 1), with gM_mul for x0 gM = gamma x0. The distance is the sup distance on Frame.
\end{proof}

\begin{lemma}\label{is-onto-share-counts-section-poincare-count-sum}
Let $p,q,f\geq0$ be integers, let $J_{p,q}$ be the diagonal matrix with $p$ entries $1$ followed by $q$ entries $-1$,
and let $s>2f(p+q)^2$ and $c>0$ be real numbers. For a complex $(p+q)\times(p+q)$ matrix $g$ write $g^*=\bar
g^{\mathsf T}$ and $D(g)$ for its lower right $q\times q$ block. There is a real number $C$ with the following
property. Let $I$ be a set and $(g_{a,k})$, for $a\in I$ and $1\le k\le f$, complex $(p+q)\times(p+q)$ matrices with
$g_{a,k}^*J_{p,q}g_{a,k}=J_{p,q}$, such that for all $a\ne a'$ in $I$ there are $k$, $i$ and $j$ with
$\big|(g_{a',k}g_{a,k}^{-1})_{ij}-\delta_{ij}\big|\geq c$. Then
\[
 \sum_{a\in I}\prod_{k=1}^f|\det D(g_{a,k})|^{-2s}\le C .
\]
\end{lemma}
\leanref{ReciprocityTheory/ReductiveGroups/ArithmeticGeometry/ArakelovVolumes/UpperBounds/SumBoundStatementPfC.lean}{ReciprocityTheory.ChartRestriction.G.AR.Up.sumBd_sum}
\leanref{ReciprocityTheory/ReductiveGroups/ArithmeticGeometry/ArakelovVolumes/UpperBounds/SumBoundStatementPfC.lean}{ReciprocityTheory.ChartRestriction.G.AR.Up.sumBd_count_real}
\leanref{ReciprocityTheory/ReductiveGroups/ArithmeticGeometry/ArakelovVolumes/UpperBounds/SumBoundStatementPfB.lean}{ReciprocityTheory.ChartRestriction.G.AR.Up.sumBd_sep_diff}
\leanref{ReciprocityTheory/ReductiveGroups/ArithmeticGeometry/ArakelovVolumes/UpperBounds/SumBoundStatementPfA.lean}{ReciprocityTheory.ChartRestriction.G.AR.Up.sumBd_entry_bound}
\leanref{ReciprocityTheory/ReductiveGroups/ArithmeticGeometry/ArakelovVolumes/UpperBounds/SumBoundStatementPfA.lean}{ReciprocityTheory.ChartRestriction.G.AR.Up.sumBd_detD_ge}
\begin{proof}
Write $J=J_{p,q}$, and $\|m\|^2$ for the sum of the squares of the absolute values of the entries of a matrix $m$,
so that $\|m\|^2=\operatorname{tr}(m^*m)=\operatorname{tr}(mm^*)$.

\emph{Entries and the block $D$.} Let $g$ satisfy $g^*Jg=J$, with blocks $A$ ($p\times p$, upper left), $B$ (upper
right), $C'$ (lower left) and $D=D(g)$. As $J^2=1$, $g^{-1}=Jg^*J$, so $gJg^*=J$ and $|(g^{-1})_{ij}|=|g_{ji}|$.
Comparing blocks, $D^*D=1+B^*B$, $AA^*=1+BB^*$ and $A^*A=1+C'^*C'$. For a positive semidefinite Hermitian matrix
$H$ with eigenvalues $h_i\geq0$, $\det(1+H)=\prod_i(1+h_i)\geq1+\operatorname{tr}H$. Hence
$|\det D|^2=\det(1+B^*B)\geq1+\|B\|^2$; in particular $|\det D|\geq1$. Taking traces, $\|D\|^2=q+\|B\|^2$ and
$\|A\|^2=p+\|B\|^2=p+\|C'\|^2$. So every entry of $g$ satisfies
\[
 |g_{ij}|^2\le p+q+\|B\|^2\le p+q-1+|\det D|^2\le(p+q+1)|\det D|^2 .\tag{$*$}
\]

\emph{Separation.} Let $g$ satisfy $g^*Jg=J$ and have entries of absolute value at most $R>0$, let $h$ be a matrix,
and let $|(hg^{-1})_{ij}-\delta_{ij}|\geq c$. Then $|h_{il}-g_{il}|\geq c/((p+q+1)R)$ for some $l$. Indeed
$hg^{-1}-1=(h-g)g^{-1}$, so its $(i,j)$ entry is $\sum_l(h_{il}-g_{il})(g^{-1})_{lj}$ with $|(g^{-1})_{lj}|=|g_{jl}|\le R$;
if every $|h_{il}-g_{il}|$ were smaller, this entry would have absolute value less than $(p+q)c/(p+q+1)<c$.

\emph{Counting.} For $R\geq1$ let $I_R$ be the set of $a\in I$ such that every entry of every $g_{a,k}$ has absolute
value at most $R$. Put $\eta=c/(2(p+q+1)R)$ and $N=2f(p+q)^2$. For distinct $a,a'\in I_R$, the hypothesis and the
separation step, for $g=g_{a,k}$ and $h=g_{a',k}$, give an entry in which $g_{a',k}$ and $g_{a,k}$ differ by at least
$2\eta$. Send $a\in I_R$ to the $N$ integers $\lfloor\operatorname{Re}(g_{a,k})_{ij}/\eta\rfloor$ and
$\lfloor\operatorname{Im}(g_{a,k})_{ij}/\eta\rfloor$, which lie between $-\lceil R/\eta\rceil$ and $\lceil R/\eta\rceil$.
This map is injective: if $a$ and $a'$ had the same image, the real and imaginary parts of each entry would differ
by less than $\eta$, so the entries would differ by less than $2\eta$. As $R/\eta=2(p+q+1)R^2/c$ and $R\geq1$,
\[
 |I_R|\le\big(2\lceil R/\eta\rceil+1\big)^N\le(2R/\eta+3)^N\le(KR^2)^N,\qquad K=4(p+q+1)/c+3 .
\]

\emph{The sum.} By $(*)$ and $s\geq0$, $w_{a,k}=|\det D(g_{a,k})|^{-2s}$ lies in $(0,1]$. Put $B_0=4(p+q+1)$. Every
$a\in I$ lies in some $I_{2^m}$; let $m(a)$ be the least such $m$. If $m(a)=0$, then $\prod_kw_{a,k}\le1\le B_0^s$.
If $m(a)=m+1$, some entry of some $g_{a,k}$ has absolute value greater than $2^m$, so
$|\det D(g_{a,k})|^2>4^m/(p+q+1)$ by $(*)$, and $\prod_kw_{a,k}\le w_{a,k}\le\big((p+q+1)/4^m\big)^s=(B_0/4^{m(a)})^s$.
Hence
\[
 \sum_{a\in I}\prod_kw_{a,k}\le\sum_{m\geq0}|I_{2^m}|\,\Big(\frac{B_0}{4^m}\Big)^s\le\sum_{m\geq0}(K4^m)^N
 \Big(\frac{B_0}{4^m}\Big)^s=K^NB_0^s\sum_{m\geq0}\Big(\frac{4^N}{4^s}\Big)^m,
\]
which is finite, as $s>N$, and depends only on $p$, $q$, $f$, $s$ and $c$.
% Lean: sumBd_sum (the dyadic sets Bs m, hpt with Nat.find, sumBd_w_le_one, sumBd_w_le_big, sumBd_prod_le_big,
% sumBd_id and the geometric series with ratio rho = 4^D / 4^s, D = f (p + q) (p + q) 2); sumBd_count_real
% (eta = c / (p + q + 1) / R / 2, sumBd_count with the floors of the real and imaginary parts, card of Icc);
% sumBd_sep_diff (sumBd_mul_inv, sumBd_inv_entry'); sumBd_entry_bound and sumBd_detD_ge (sumBdIsUpq_blocks,
% sumBd_det_one_add_ge through the eigenvalues, sumBd_trace_ctm, sumBd_trace_mct). Here IsUpq g is
% g^H Jf g = Jf. The sum is a tsum in [0, infinity].
\end{proof}

\begin{lemma}\label{is-onto-share-counts-section-poincare-count}
Let $X$ consist of the data (1) to (6) of flag data of rank $r$ over $E$, such that $\Gamma$ acts properly
discontinuously on $\mathcal D_X$ modulo $\mathcal B^\Psi$ (Lemma~\ref{is-onto-share-counts-section-discontinuous}),
use the notation of Definition~\ref{is-onto-share-counts-section-levi}, and let $x_0$ be an orthonormal point. For
$y\in M_r(\mathbf C)^\Psi$ and $\tau\in\Psi$ let $\delta_\tau(y)$ be the leading principal $p_0\times p_0$ minor of
$x_{0,\tau}^{-1}y_\tau$, and put
\[
 \rho_Y(y)=-\sum_{\tau\in\Psi}\big(\log|d_{p_0}(y,\tau)|-2\log|\delta_\tau(y)|\big),
\]
where $\log0$ is read as $0$. Write $|w|$ for the Euclidean norm of $w\in\mathbf C^n$.
\begin{enumerate}
\item For $y\in\mathcal D_X$, $\delta_\tau(y)\ne0$ and $\rho_Y(yb)=\rho_Y(y)$ for $b\in\mathcal B^\Psi$; and
  $\rho_Y(c_{x_0}(z))=\rho_{x_0}(z)$ for $z\in\Omega(x_0)$.
\item For every real number $\nu\geq0$ there is $\varrho_1>0$ such that
  $e^{-\nu\rho_Y(c_x(w))}\le e^{3\nu|\Psi|\log2}\,e^{-\nu\rho_Y(x)}$ for every orthonormal point $x$ and every
  $w\in\mathbf C^n$ with $|w|<\varrho_1$.
\item For every real number $s>2|\Psi|r^2$ there is a real number $C$ such that
  $\sum_{\gamma\in\Gamma}e^{-s\rho_Y(\gamma x)}\le C$ for every orthonormal point $x$.
\end{enumerate}
\end{lemma}
\leanref{ReciprocityTheory/ReductiveGroups/ArithmeticGeometry/ArakelovVolumes/UpperBounds/PoinStatementPfF.lean}{ReciprocityTheory.ChartRestriction.G.AR.Up.PoinPf.rhY}
\leanref{ReciprocityTheory/ReductiveGroups/ArithmeticGeometry/ArakelovVolumes/UpperBounds/PoinStatementPfF.lean}{ReciprocityTheory.ChartRestriction.G.AR.Up.PoinPf.lminor_gM_ne}
\leanref{ReciprocityTheory/ReductiveGroups/ArithmeticGeometry/ArakelovVolumes/UpperBounds/PoinStatementPfF.lean}{ReciprocityTheory.ChartRestriction.G.AR.Up.PoinPf.kaY_rmul}
\leanref{ReciprocityTheory/ReductiveGroups/ArithmeticGeometry/ArakelovVolumes/UpperBounds/PoinStatementPfF.lean}{ReciprocityTheory.ChartRestriction.G.AR.Up.PoinPf.kaY_cpt}
\leanref{ReciprocityTheory/ReductiveGroups/ArithmeticGeometry/ArakelovVolumes/UpperBounds/PoinStatementPfH.lean}{ReciprocityTheory.ChartRestriction.G.AR.Up.PoinPf.ka_comp}
\leanref{ReciprocityTheory/ReductiveGroups/ArithmeticGeometry/ArakelovVolumes/UpperBounds/PoinStatementPfH.lean}{ReciprocityTheory.ChartRestriction.G.AR.Up.PoinPf.det_near_one}
\leanref{ReciprocityTheory/ReductiveGroups/ArithmeticGeometry/ArakelovVolumes/UpperBounds/PoinStatementPfH.lean}{ReciprocityTheory.ChartRestriction.G.AR.Up.PoinPf.Kc_entry}
\leanref{ReciprocityTheory/ReductiveGroups/ArithmeticGeometry/ArakelovVolumes/UpperBounds/PoinStatementPfH.lean}{ReciprocityTheory.ChartRestriction.G.AR.Up.PoinPf.pert_entry}
\leanref{ReciprocityTheory/ReductiveGroups/ArithmeticGeometry/ArakelovVolumes/UpperBounds/PoinStatementPfG.lean}{ReciprocityTheory.ChartRestriction.G.AR.Up.PoinPf.kaSum}
\leanref{ReciprocityTheory/ReductiveGroups/ArithmeticGeometry/ArakelovVolumes/UpperBounds/PoinStatementPfG.lean}{ReciprocityTheory.ChartRestriction.G.AR.Up.PoinPf.Mq_div}
\leanref{ReciprocityTheory/ReductiveGroups/ArithmeticGeometry/ArakelovVolumes/UpperBounds/PoinStatementPfF.lean}{ReciprocityTheory.ChartRestriction.G.AR.Up.PoinPf.kaY_orth}
\begin{proof}
(1) Let $y\in\mathcal D_X$ and $b\in\mathcal B^\Psi$. By the proof of Lemma~\ref{is-onto-share-majorant}(2),
$d_{p_0}(yb,\tau)=|\prod_{a\le p_0}(b_\tau)_{aa}|^2d_{p_0}(y,\tau)$, and as $b_\tau$ is upper triangular,
$\delta_\tau(yb)=\delta_\tau(y)\prod_{a\le p_0}(b_\tau)_{aa}$. By Lemma~\ref{is-onto-share-majorant}(3) it therefore
suffices to show $\delta_\tau(y)\ne0$ for orthonormal $y$. Then $g=x_{0,\tau}^{-1}y_\tau$ satisfies
$g^{\mathsf T}J\bar g=J$, as both points are orthonormal. Its leading $p_0\times p_0$ block $a$ and its lower left block
$c$ satisfy $a^{\mathsf T}\bar a=1+c^{\mathsf T}\bar c$, which is positive definite; so $\det a=\delta_\tau(y)\ne0$. The
invariance of $\rho_Y$ follows from the two formulas above. For $z\in\Omega(x_0)$,
$x_{0,\tau}^{-1}c_{x_0}(z)_\tau=\ell_\tau(z)$ has leading minors $1$, so $\rho_Y(c_{x_0}(z))=\rho_{x_0}(z)$.

\emph{The blocks of $g$.} Let $x$ be orthonormal and $\tau\in\Psi$. As in (1), $g=x_{0,\tau}^{-1}x_\tau$ satisfies
$g^{\mathsf T}J\bar g=J$; conjugating, $g^*Jg=J$ with $g^*=\bar g^{\mathsf T}$, so $g^{-1}=Jg^*J$ and $gJg^*=J$.
Write $a$, $b$ for the leading $p_0\times p_0$ block of $g$ and the block to its right. The upper left block of
$gJg^*=J$ is $aa^*-bb^*=1$, and $a$ is invertible by (1). So $K=a^{-1}b$ satisfies
$KK^*=a^{-1}(aa^*-1)(a^*)^{-1}=1-a^{-1}(a^{-1})^*$, whose diagonal entries are at most $1$. Hence every entry of
$K$ has absolute value at most $1$.

(2) Let $\nu\geq0$. As the determinant is continuous, there is $\eta>0$ such that every complex $p_0\times p_0$
matrix whose entries differ from those of the identity matrix by less than $\eta$ has determinant of absolute
value at least $\frac12$. For an orthonormal point $x$, the Gram matrix of the columns of $c_x(w)_\tau=x_\tau\ell_\tau(w)$
is $\ell_\tau(w)^{\mathsf T}J\overline{\ell_\tau(w)}$, which does not depend on $x$, and its leading minors are $\pm1$
at $w=0$. By continuity there is $R_1>0$ with $|d_{p_0}(c_x(w),\tau)|<2$ for $|w|<R_1$, every $\tau$ and every
orthonormal $x$. Put $\varrho_1=\min(R_1,\eta/(r+1))$.

Let $x$ be orthonormal, $|w|<\varrho_1$ and $\tau\in\Psi$. Then $\log|d_{p_0}(c_x(w),\tau)|\le\log2$, also when this
minor vanishes, as $\log0$ is read as $0$. Let $g$, $a$, $b$ and $K$ be as above, and write $\ell=\ell_\tau(w)$ in blocks,
with $\ell_{11}$ its leading $p_0\times p_0$ block and $\ell_{21}$ the block below it; the block of $\ell$ to the
right of $\ell_{11}$ is $0$. The leading $p_0\times p_0$ block of $x_{0,\tau}^{-1}c_x(w)_\tau=g\ell$ is
$a\ell_{11}+b\ell_{21}=a(\ell_{11}+K\ell_{21})$, so
\[
 \delta_\tau(c_x(w))=\delta_\tau(x)\det(\ell_{11}+K\ell_{21}).
\]
The entries of $\ell-1$ are coordinates of $w$, of absolute value at most $|w|$. So every entry of
$\ell_{11}+K\ell_{21}-1$ has absolute value at most $|w|+(r-p_0)|w|\le(r+1)|w|<\eta$, and
$|\det(\ell_{11}+K\ell_{21})|\geq\frac12$. Hence $\log|\delta_\tau(c_x(w))|\geq\log|\delta_\tau(x)|-\log2$. As
$|d_{p_0}(x,\tau)|=1$,
\[
 \log|d_{p_0}(c_x(w),\tau)|-2\log|\delta_\tau(c_x(w))|\le3\log2+\big(\log|d_{p_0}(x,\tau)|-2\log|\delta_\tau(x)|\big).
\]
Summing over $\tau$ gives $\rho_Y(c_x(w))\geq\rho_Y(x)-3|\Psi|\log2$. Multiplying by $-\nu\le0$ and exponentiating
gives (2).

(3) Let $s>2|\Psi|r^2$, and let $c>0$ be the number of
Lemma~\ref{is-onto-share-counts-section-poincare-count-separation}. Let $x$ be orthonormal. For $\gamma\in\Gamma$ the
point $\gamma x$ is orthonormal, as $\gamma$ preserves $h$; put $g_{\gamma,\tau}=x_{0,\tau}^{-1}(\gamma x)_\tau$, so that
$g_{\gamma,\tau}^{\mathsf T}J\bar g_{\gamma,\tau}=J$ as above. For $\gamma\ne\gamma'$,
\[
 g_{\gamma',\tau}\,g_{\gamma,\tau}^{-1}=x_{0,\tau}^{-1}(\gamma'\gamma^{-1})^\tau x_{0,\tau}
 =x_{0,\tau}^{-1}\big((\gamma'\gamma^{-1})x_0\big)_\tau,
\]
and by Lemma~\ref{is-onto-share-counts-section-poincare-count-separation} for $\gamma'\gamma^{-1}\ne1$ this matrix
differs from $1$ by at least $c$ in some entry, for some $\tau$. Let $M_{\gamma,\tau}$ be the complex conjugate of
$g_{\gamma,\tau}$ with its rows and columns put in the order $p_0+1,\dots,r,1,\dots,p_0$. The reordering takes $J$ to
$-J_{r-p_0,p_0}$, in the notation of Lemma~\ref{is-onto-share-counts-section-poincare-count-sum}, so
$M_{\gamma,\tau}^*J_{r-p_0,p_0}M_{\gamma,\tau}=J_{r-p_0,p_0}$; the matrices $M_{\gamma',\tau}M_{\gamma,\tau}^{-1}$ have
the same separation; and the lower right $p_0\times p_0$ block of $M_{\gamma,\tau}$ is the conjugate of the leading
block of $g_{\gamma,\tau}$, with determinant $\overline{\delta_\tau(\gamma x)}$. As $|d_{p_0}(\gamma x,\tau)|=1$,
\[
 e^{-s\rho_Y(\gamma x)}=\prod_{\tau\in\Psi}|\delta_\tau(\gamma x)|^{-2s}.
\]
Lemma~\ref{is-onto-share-counts-section-poincare-count-sum}, for $p=r-p_0$, $q=p_0$, $f=|\Psi|$, $I=\Gamma$ and the
family $(M_{\gamma,\tau})$, gives a bound $C$ for $\sum_\gamma e^{-s\rho_Y(\gamma x)}$ that does not depend on $x$.
% Lean: (1) is lminor_gM_ne (lminor_orth_ne with Mq_isUpq and det_lowerBlock_eq_conj_minor), rhY_rmul /
% kaY_rmul and rhY_cpt / kaY_cpt; kaY = exp(-bb rhY), and Real.log reads log 0 as 0. (2) is PoinPf.ka_comp
% (bb >= 0, no condition cpt x w in X.dom): det_near_one, e1 for R1 (cont_lminor, norm_lminor_zero,
% dmin_cpt_orth), hterm and hrh. The Lean works with Mq = the conjugate of g reordered by eS (the last r - p0
% indices first), so its block (Mq)_22^-1 (Mq)_21 is the conjugate of K = a^-1 b; Kc_entry is |K_ij| <= 1 from
% sumBdIsUpq_blocks, pert_entry the entry bound (r + 1)|w| (N_entries, nmat_entry), Mq_cpt' and det_split the
% factorization of the minor. (3) is PoinPf.kaSum: PoinPf.sep, sumBd_sum with p = r - p0, q = p0,
% f = |Psi|, Mq_isUpq, Mq_div, and kaY_orth for e^{-s rhY} at orthonormal points. Size: about 1,000 lines of Lean
% for (2) and (3) together.
\end{proof}

\begin{lemma}\label{is-onto-share-counts-section-poincare-ball}
Let $X$ consist of the data (1) to (6) of flag data of rank $r$ over $E$, such that $\Gamma$ acts properly
discontinuously on $\mathcal D_X$ modulo $\mathcal B^\Psi$, and use the notation of
Definitions~\ref{is-onto-share-counts-section-levi} and~\ref{is-onto-share-counts-section-gauge} and the sets
$B_\varrho(y)$ of Lemma~\ref{is-onto-share-counts-section-poincare-mean}. Let $x_0$ be an orthonormal point,
$\lambda$ a weight, $W$ a gauge of weight $\lambda$ that is bounded below, $\nu>4|\Psi|r^2$ a real number, $f$ a
holomorphic function on $\Omega(x_0)$ that is square integrable for $W$ and $\nu$, and $P$ a function that is
$\mathcal B$-equivariant of weight $\lambda$ with $P\circ c_{x_0}=f$ on $\Omega(x_0)$. That $W$ is bounded below and
that $f$ is holomorphic are not used in the proof. Then there are real numbers $0<\varrho_0\le1$ and $C$ such that
for every orthonormal point $x$
\[
 \sum_{\gamma\in\Gamma}\Big(\int_{B_{\varrho_0}(x)}|P(\gamma c_x(w))|^2\,W(\gamma c_x(w))\,dw\Big)^{1/2}\le C .
\]
\end{lemma}
\leanref{ReciprocityTheory/ReductiveGroups/ArithmeticGeometry/ArakelovVolumes/UpperBounds/PoinStatementPfI.lean}{ReciprocityTheory.ChartRestriction.G.AR.Up.PoinPf.sumBall_holds}
\begin{proof}
Let $\rho_Y$ be the function of Lemma~\ref{is-onto-share-counts-section-poincare-count} for $x_0$, and put
$\kappa(y)=e^{-\nu\rho_Y(y)}$ and $\psi(y)=|P(y)|^2W(y)/\kappa(y)$ for $y\in\mathcal D_X$, and $\psi=0$ outside
$\mathcal D_X$. By Lemma~\ref{is-onto-share-counts-section-poincare-count}(1) and the transformation laws of $P$ and
$W$, $\psi(yb)=\psi(y)$ for $b\in\mathcal B^\Psi$, and $\kappa(c_{x_0}(z))=e^{-\nu\rho_{x_0}(z)}$. So
\[
 N_f=\int_{\Omega(x_0)}\psi(c_{x_0}(z))\,D_{x_0}(z)\,dz=\int_{\Omega(x_0)}|f(z)|^2e^{\nu\rho_{x_0}(z)}W(c_{x_0}(z))\,
 D_{x_0}(z)\,dz<\infty .
\]

Let $\varrho_1$ be as in Lemma~\ref{is-onto-share-counts-section-poincare-count}(2) for $\nu$, which is positive, and
$K=e^{3\nu|\Psi|\log2}$. At an
orthonormal point $x'$ the minors $d_i(c_{x'}(w),\tau)$ are those of $\ell_\tau(w)^{\mathsf T}J\overline{\ell_\tau(w)}$,
which equal $\pm1$ at $w=0$; so there is $R_2>0$, independent of $x'$, with $|d_i(c_{x'}(w),\tau)|<2$ for $|w|<R_2$, and
then $D_{x'}(w)\geq c_d=4^{-r|\Psi|}$. Put $\varrho_0=\min(\varrho_1,R_2,1)$.

Let $x$ be orthonormal and $\gamma\in\Gamma$. Then $x'=\gamma x$ is orthonormal and $\gamma c_x(w)=c_{x'}(w)$, so
$B_{\varrho_0}(x)=B_{\varrho_0}(x')$. For $w$ in this set, Lemma~\ref{is-onto-share-counts-section-poincare-count}(2)
gives $\kappa(c_{x'}(w))\le K\kappa(x')$, and so
\[
 |P(c_{x'}(w))|^2W(c_{x'}(w))=\psi(c_{x'}(w))\,\kappa(c_{x'}(w))\le\frac{K\kappa(x')}{c_d}\,\psi(c_{x'}(w))\,D_{x'}(w).
\]
Integrating over $B_{\varrho_0}(x')\subset\Omega(x')$ and using
Lemma~\ref{is-onto-share-counts-section-poincare-measure}(1) for $x'$ and $x_0$,
\[
 \int_{B_{\varrho_0}(x)}|P(\gamma c_x(w))|^2W(\gamma c_x(w))\,dw\le\frac{K\kappa(\gamma x)}{c_d}\int_{\Omega(x')}
 \psi(c_{x'}(w))\,D_{x'}(w)\,dw=c_0\,e^{-\nu\rho_Y(\gamma x)},
\]
with $c_0=KN_f/c_d$. Taking square roots and summing over $\gamma$, the sum in the statement is at most
$c_0^{1/2}\sum_\gamma e^{-(\nu/2)\rho_Y(\gamma x)}$, which is bounded independently of $x$ by
Lemma~\ref{is-onto-share-counts-section-poincare-count}(3) for $s=\nu/2>2|\Psi|r^2$.
% Lean: PoinPf.sumBall_holds takes Hyp X x0 la W bb F P, whose fields WLower X la W and
% DifferentiableOn C F (OmF X x0) it discards (the two dashes in obtain); the statement keeps them, unused.
% Proof: ka_comp, e2 and cd for the density, psiF with psiF_rmul, psiF_mul and psiF_int,
% PairUnfoldPf.cpt_lact and isOrth_lact, measInv_holds, kaY_half and kaSum with s = bb / 2; Nn is N_f and
% c0 = K / cd * Nn.
\end{proof}

\begin{lemma}\label{is-onto-share-counts-section-poincare}
Let $X$ consist of the data (1) to (6) of flag data of rank $r$ over $E$, such that $\Gamma$ acts properly
discontinuously on $\mathcal D_X$ modulo $\mathcal B^\Psi$ (Lemma~\ref{is-onto-share-counts-section-discontinuous}),
and use the notation of Definitions~\ref{is-onto-share-counts-section-levi}
and~\ref{is-onto-share-counts-section-gauge}. Let $x_0$ be an orthonormal point, $\lambda$ a weight, $W$ a gauge of
weight $\lambda$ that is bounded below, $\nu>4|\Psi|r^2$ a real number, $f$ a holomorphic function on $\Omega(x_0)$
that is square integrable for $W$ and $\nu$, and $P$ a function that is $\mathcal B$-equivariant of weight $\lambda$
with $P\circ c_{x_0}=f$ on $\Omega(x_0)$. Then:
\begin{enumerate}
\item for every $x\in\mathcal D_X$ the family $\big(\chi_\lambda(\gamma)^{-1}P(\gamma x)\big)_{\gamma\in\Gamma}$ is
  absolutely summable;
\item $\mathcal P_\lambda P$ is a bounded section of weight $\lambda$;
\item let $\Phi\subset\Omega(x_0)$ be a measurable set such that $z=z'$ whenever $z,z'\in\Phi$ satisfy
  $c_{x_0}(z')=\gamma c_{x_0}(z)b$ for some $\gamma\in\Gamma$ and $b\in\mathcal B^\Psi$, and such that for almost every
  $z\in\Omega(x_0)$ there are $z'\in\Phi$, $\gamma\in\Gamma$ and $b\in\mathcal B^\Psi$ with
  $c_{x_0}(z)=\gamma c_{x_0}(z')b$; then
  \[
   \int_\Phi|\mathcal P_\lambda P(c_{x_0}(z))|^2\,W(c_{x_0}(z))\,D_{x_0}(z)\,dz<\infty;
  \]
\item for every real number $0<\varrho\le1$ there is a real number $C$ such that, for every orthonormal point $x$,
  $|\mathcal P_\lambda P(x)|$ is at most $C$ times the supremum in $[0,\infty]$ of
  $\big(N_\lambda(c_x(z))^2/W(c_x(z))\big)^{1/2}$ over the $z$ with $|z|<\varrho$ and $c_x(z)\in\mathcal D_X$;
\item for every weight $\lambda'$, every section $g$ of weight $\lambda'$ and every $x\in\mathcal D_X$,
  $\mathcal P_{\lambda+\lambda'}(gP)(x)=g(x)\,\mathcal P_\lambda P(x)$.
\end{enumerate}
\end{lemma}
\leanref{ReciprocityTheory/ReductiveGroups/ArithmeticGeometry/ArakelovVolumes/UpperBounds/FundSetMeasStatement.lean}{ReciprocityTheory.ChartRestriction.G.AR.Up.PoinStatement}
\leanref{ReciprocityTheory/ReductiveGroups/ArithmeticGeometry/ArakelovVolumes/UpperBounds/PoinStatementPf.lean}{ReciprocityTheory.ChartRestriction.G.AR.Up.poin_holds}
\begin{proof}
Let $c>0$ with $W\geq cN_\lambda^2$ on $\mathcal D_X$, use the sets $B_\varrho(y)$ of
Lemma~\ref{is-onto-share-counts-section-poincare-mean}, and let $\varrho_0$ and $C_S$ be the numbers of
Lemma~\ref{is-onto-share-counts-section-poincare-ball}. For an orthonormal point $x$ and $\gamma\in\Gamma$ put
\[
 I_\gamma(x)=\int_{B_{\varrho_0}(x)}|P(\gamma c_x(w))|^2\,W(\gamma c_x(w))\,dw,
\]
so that $\sum_\gamma I_\gamma(x)^{1/2}\le C_S$. As $\chi_\lambda$ has finite order on $\Gamma$, $|\chi_\lambda(\gamma)|=1$, so
the terms of the Poincaré series have the absolute values $|P(\gamma x)|$. For $\gamma\in\Gamma$ the function
$y\mapsto P(\gamma y)$ is $\mathcal B$-equivariant of weight $\lambda$, as $\gamma(yb)=(\gamma y)b$ and $\gamma$ preserves
$\mathcal D_X$; and $W(\gamma y)=W(y)$.

\emph{A local bound.} Let $x\in\mathcal D_X$. By Lemma~\ref{is-onto-share-majorant}(3) there is $b_0\in\mathcal B^\Psi$
with $x_1=xb_0$ orthonormal. Lemma~\ref{is-onto-share-counts-section-poincare-mean}(2) for $x_1$ and $\varrho_0$ gives
$\varepsilon$ and $C_A$; as $y\mapsto yb_0$ is continuous, there is $\varepsilon'>0$ such that $yb_0$ is within
$\varepsilon$ of $x_1$ and $y\in\mathcal D_X$ whenever $y$ is within $\varepsilon'$ of $x$. For such $y$ and every
$\gamma$, as $N_\lambda^2\le W/c$ on $\mathcal D_X$,
\[
 |b_0^\lambda|^2\,|P(\gamma y)|^2=|P(\gamma yb_0)|^2\le C_A\int_{B_{\varrho_0}(x_1)}|P(\gamma c_{x_1}(w))|^2
 N_\lambda(c_{x_1}(w))^2\,dw\le\frac{C_A}c\,I_\gamma(x_1).
\]
So $|P(\gamma y)|\le u_\gamma$ with $u_\gamma=|b_0^\lambda|^{-1}(C_A/c)^{1/2}I_\gamma(x_1)^{1/2}$, and $\sum_\gamma u_\gamma<\infty$.

(1) This is the local bound at $y=x$.

(4) Let $0<\varrho\le1$, $\varrho'=\min(\varrho,\varrho_0)$, and let $C_A$ be the constant of
Lemma~\ref{is-onto-share-counts-section-poincare-mean}(1) for $\varrho'$. Let $x$ be orthonormal and $S$ the supremum
in (4). For $w\in B_\varrho(x)$, $N_\lambda(c_x(w))^2\le S^2W(c_x(w))=S^2W(\gamma c_x(w))$. So that lemma, for the
holomorphic function $y\mapsto P(\gamma y)$, gives
\[
 |P(\gamma x)|^2\le C_A\int_{B_{\varrho'}(x)}|P(\gamma c_x(w))|^2N_\lambda(c_x(w))^2\,dw\le C_AS^2I_\gamma(x),
\]
and $|\mathcal P_\lambda P(x)|\le\sum_\gamma|P(\gamma x)|\le C_A^{1/2}C_S\,S$.

(2) Each term $y\mapsto\chi_\lambda(\gamma)^{-1}P(\gamma y)$ is holomorphic on $\mathcal D_X$, and by the local bound the
series is dominated near each point by a convergent series of constants; so its sum $\mathcal P_\lambda P$ is
holomorphic on $\mathcal D_X$. For $b\in\mathcal B^\Psi$ each term is multiplied by $b^\lambda$ when $y$ is replaced by
$yb$, so $\mathcal P_\lambda P(yb)=b^\lambda\mathcal P_\lambda P(y)$. For $\gamma'\in\Gamma$, the substitution
$\gamma\mapsto\gamma\gamma'$ and $\chi_\lambda(\gamma\gamma')=\chi_\lambda(\gamma)\chi_\lambda(\gamma')$ give
$\mathcal P_\lambda P(\gamma'y)=\chi_\lambda(\gamma')\mathcal P_\lambda P(y)$. So $\mathcal P_\lambda P$ is a section of
weight $\lambda$. By (4) for $\varrho=1$ and $S\le c^{-1/2}$, $|\mathcal P_\lambda P|\le C_1=C_A^{1/2}C_Sc^{-1/2}$ at
orthonormal points. For $x\in\mathcal D_X$ and $xb$ orthonormal, $\|\mathcal P_\lambda P\|_\lambda(x)=
\|\mathcal P_\lambda P\|_\lambda(xb)=|\mathcal P_\lambda P(xb)|\le C_1$, as $\|\cdot\|_\lambda$ is invariant under
$\mathcal B^\Psi$ and the minors of an orthonormal point are $\pm1$. So $\mathcal P_\lambda P$ is bounded.

(3) Let $\kappa$ and $\psi$ be as in the proof of Lemma~\ref{is-onto-share-counts-section-poincare-ball}, and let
$C_2$ bound $\sum_\gamma\kappa(\gamma x)$ for orthonormal $x$, by
Lemma~\ref{is-onto-share-counts-section-poincare-count}(3) for $s=\nu$. As $\kappa$ is invariant under
$\mathcal B^\Psi$ and $\gamma(xb)=(\gamma x)b$, the same bound holds for every $y\in\mathcal D_X$. For $z\in\Phi$ and
$y=c_{x_0}(z)$, the Cauchy-Schwarz inequality and $W(\gamma y)=W(y)$ give
\[
 |\mathcal P_\lambda P(y)|^2W(y)\le\Big(\sum_\gamma\frac{|P(\gamma y)|^2}{\kappa(\gamma y)}\Big)\Big(\sum_\gamma\kappa(\gamma y)\Big)
 W(y)\le C_2\sum_\gamma\psi(\gamma y).
\]
Integrating over $\Phi$ against $D_{x_0}(z)\,dz$ and using
Lemma~\ref{is-onto-share-counts-section-poincare-measure}(2) for $\psi$, the integral in (3) is at most
$C_2\int_{\Omega(x_0)}\psi(c_{x_0}(z))D_{x_0}(z)\,dz$, which is finite as in the proof of
Lemma~\ref{is-onto-share-counts-section-poincare-ball}.

(5) Let $x\in\mathcal D_X$. As $g(\gamma x)=\chi_{\lambda'}(\gamma)g(x)$ and $\chi_{\lambda+\lambda'}=\chi_\lambda\chi_{\lambda'}$
on $\Gamma$, the family defining $\mathcal P_{\lambda+\lambda'}(gP)(x)$ is $g(x)$ times the family defining
$\mathcal P_\lambda P(x)$. If $g(x)=0$ both sides vanish; otherwise the two families are summable together, and the
sums agree.
% Lean: poin_holds = poin_of_inputs lu_holds PoinPf.sumBall_holds PoinPf.ff_holds = PoinPf.poin_of (loc_of
% (subMeanLoc_of lu_holds) sumBall_holds) (ball_of subMean_holds sumBall_holds) ff_holds. The local bound is
% loc_of (isBEq_lact, the gauge bound hcW); (4) is ball_of with key_int and enorm_poin_le; (2) is poin_of:
% differentiableOn_tsum_loc, poin_rmul, poin_lact, and the bound horth (supR_le) with
% FrameDegreeData.pnorm_rmul and pnorm_orth'; (3) is ff_holds (cs_sq, kaSum, kaY_rmul, PairUnfoldPf.lunfold
% with measInv_holds); (5) is poin_mul. Nothing below this proof remains open.
\end{proof}

\begin{lemma}\label{is-onto-share-counts-section-separate}
Let $X$ consist of the data (1) to (6) of flag data of rank $r$ over $E$, with archimedean data for $(L,F,T)$ and a
flat model for them, with base points $y_1,\dots,y_\delta$ (Definition~\ref{is-onto-share-counts-model}), such
that $\Gamma$ acts properly discontinuously on $\mathcal D_X$ modulo $\mathcal B^\Psi$
(Lemma~\ref{is-onto-share-counts-section-discontinuous}). Let $\lambda$ be a weight with curvature constants
$\nu,m,\eta$ for $\lambda$ and $L$ (Definition~\ref{is-onto-share-counts-section-levi}), and $1\le j\le\delta$. Then
there is a bounded section $Q$ of weight $\lambda$ with $Q(y_j)\ne0$ and $Q(y_{j'})=0$ for $j'\ne j$.
\end{lemma}
\leanref{ReciprocityTheory/ReductiveGroups/ArithmeticGeometry/ArakelovVolumes/UpperBounds/InterpStatementPfB.lean}{ReciprocityTheory.ChartRestriction.G.AR.Up.poin_single}
\leanref{ReciprocityTheory/ReductiveGroups/ArithmeticGeometry/ArakelovVolumes/UpperBounds/InterpStatementPfB.lean}{ReciprocityTheory.ChartRestriction.G.AR.Up.poin_zero}
\leanref{ReciprocityTheory/ReductiveGroups/ArithmeticGeometry/ArakelovVolumes/UpperBounds/InterpStatementPfB.lean}{ReciprocityTheory.ChartRestriction.G.AR.Up.chi_ne_zero}
\begin{proof}
Lemma~\ref{is-onto-share-counts-section-kronecker} gives an orthonormal point $x_0$, a gauge $W$ of weight
$\lambda$ that is bounded below, a holomorphic function $f$ on $\Omega(x_0)$ that is square integrable for $W$ and
$\nu$, and a $\mathcal B$-equivariant function $P$ of weight $\lambda$ with $P\circ c_{x_0}=f$, $P(y_j)\ne0$ and
$P(\gamma y_{j'})=0$ whenever $\gamma\ne1$ or $j'\ne j$. As $\nu>4|\Psi|r^2$,
Lemma~\ref{is-onto-share-counts-section-poincare}(2) shows that $Q=\mathcal P_\lambda P$ is a bounded section of
weight $\lambda$.

At $y_j$, every term $\chi_\lambda(\gamma)^{-1}P(\gamma y_j)$ with $\gamma\ne1$ vanishes, so the family has at most one
nonzero term and $Q(y_j)=\chi_\lambda(1)^{-1}P(y_j)$. Here $\chi_\lambda(1)\ne0$, as $\chi_\lambda(1)^{o_\lambda}=1$; so
$Q(y_j)\ne0$. For $j'\ne j$, every term at $y_{j'}$ vanishes, so $Q(y_{j'})=0$.
% Lean: not one declaration. It is the common part of InterpWeight.interp_weight (after core, with the
% hypothesis PoinStatement) and LowerInterp.lowered_interp (after kron_at, with poin_holds): poin_single
% with gamma0 = 1 and chi_ne_zero, and poin_zero.
\end{proof}

\begin{lemma}\label{is-onto-share-counts-section-cusp-minors}
Let $X$ consist of the data (1) to (6) of flag data of rank $r$ over $E$
(Definition~\ref{is-onto-share-flag-data}), let $x\in\mathcal D_X$ and $\tau\in\Psi$, let $x_1,\dots,x_r$ be the
columns of $x_\tau$, and let $e\in\mathbf C^r$ be nonzero with $h_\tau(e,e)=0$. For $0\le i\le r$ let $W_i$ be the
span of $x_1,\dots,x_i$.
\begin{enumerate}
\item For each $i$ there is a unique $u_i\in W_i$ with $h_\tau(u_i,w)=h_\tau(e,w)$ for every $w\in W_i$. The number
  $\mu_i=h_\tau(u_i,u_i)$ is real, $\mu_0=\mu_r=0$, $\mu_i\le\mu_{i+1}$ for $0\le i<p_0$ and $\mu_{i+1}\le\mu_i$ for
  $p_0\le i<r$; so $0\le\mu_i\le\mu_{p_0}$ for every $i$. Moreover $m_{x,\tau}(e)=2\mu_{p_0}>0$.
\item Let $\alpha\in\mathbf C$ and let $f$ be the linear map $v\mapsto v+\alpha\,h_\tau(v,e)\,e$ of $\mathbf C^r$.
  Then $\det f=1$, and for $0\le i\le r$ the leading principal $i\times i$ minor of the Gram matrix of the columns
  of $fx_\tau$ for $h_\tau$ is $d_i(x,\tau)\,(1+2\operatorname{Re}(\alpha)\,\mu_i)$.
\end{enumerate}
\end{lemma}
\leanref{ReciprocityTheory/ReductiveGroups/ArithmeticGeometry/FundamentalDomains/PositiveMinorUnits.lean}{ReciprocityTheory.ChartRestriction.G.FD.PositiveMinorUnits.mN_bounds}
\leanref{ReciprocityTheory/ReductiveGroups/ArithmeticGeometry/FundamentalDomains/PositiveMinorUnits.lean}{ReciprocityTheory.ChartRestriction.G.FD.PositiveMinorUnits.mN_pos}
\leanref{ReciprocityTheory/ReductiveGroups/ArithmeticGeometry/FundamentalDomains/GramDeterminantReal.lean}{ReciprocityTheory.ChartRestriction.G.FD.GramDeterminantReal.det_step}
\leanref{ReciprocityTheory/ReductiveGroups/ArithmeticGeometry/FundamentalDomains/MajorantFamilyNorms.lean}{ReciprocityTheory.ChartRestriction.G.FD.MajorantFamilyNorms.maj_eq}
\leanref{ReciprocityTheory/ReductiveGroups/ArithmeticGeometry/FundamentalDomains/MajorantFamilyNorms.lean}{ReciprocityTheory.ChartRestriction.G.FD.MajorantFamilyNorms.det_flowMat}
\leanref{ReciprocityTheory/ReductiveGroups/ArithmeticGeometry/FundamentalDomains/GramFlowTransform.lean}{ReciprocityTheory.ChartRestriction.G.FD.GramFlowTransform.det_flw}
\begin{proof}
Write $h=h_\tau$, $d_i=d_i(x,\tau)$ and $g_i$ for the Gram matrix of $x_1,\dots,x_i$, so that $d_i=\det g_i$. As $h$
is Hermitian, $g_i$ is Hermitian and $d_i$ is real; as $x\in\mathcal D_X$, $\epsilon_id_i>0$. Put $c_a=h(e,x_a)$.

(1) A vector $\sum_{b\le i}y_bx_b$ satisfies the condition on $u_i$ if and only if $\sum_by_b(g_i)_{bd}=c_d$ for
$d\le i$. As $g_i$ is invertible, this system has exactly one solution, $y=cg_i^{-1}$ for the row vector
$c=(c_1,\dots,c_i)$. Then $\mu_i=\sum_b\bar y_bh(u_i,x_b)=\sum_b\bar y_bc_b$, and $\mu_i$ is real as $h$ is
Hermitian. Since $W_0=0$, $u_0=0$ and $\mu_0=0$. As $\det x_\tau\ne0$, $W_r=\mathbf C^r$; the vector $e-u_r$ is
orthogonal to $x_1,\dots,x_r$, so its coordinates $w$ in this basis satisfy $wg_r=0$, and $e=u_r$. Hence
$\mu_r=h(e,e)=0$.

Let $0\le i<r$, put $\Delta=u_{i+1}-u_i$, and let $y$ be the coefficient of $x_{i+1}$ in $u_{i+1}$. Both $u_{i+1}$ and
$u_i$ satisfy $h(\cdot,w)=h(e,w)$ for $w\in W_i$, so $\Delta$ is orthogonal to $W_i$, which contains $u_i$. So
$\mu_{i+1}=\mu_i+h(\Delta,\Delta)$. The vectors $x_1,\dots,x_i,\Delta$ arise from $x_1,\dots,x_{i+1}$ by a triangular
matrix with diagonal entries $1,\dots,1,y$, and their Gram matrix is block diagonal with blocks $g_i$ and
$h(\Delta,\Delta)$. Taking determinants, $|y|^2d_{i+1}=d_ih(\Delta,\Delta)$. As $\epsilon_{i+1}=\operatorname{sgn}(i+1)\epsilon_i$,
\[
 \operatorname{sgn}(i+1)\,(\mu_{i+1}-\mu_i)\,\epsilon_id_i=|y|^2\,\epsilon_{i+1}d_{i+1}\geq0 .
\]
So $\mu_{i+1}-\mu_i\geq0$ if $i+1\le p_0$ and $\le0$ otherwise. With $\mu_0=\mu_r=0$ this gives
$0\le\mu_i\le\mu_{p_0}$.

Suppose $\mu_{p_0}=0$. Then every $\mu_i$ vanishes, so at every step $h(\Delta,\Delta)=0$, and then $y=0$ as
$d_{i+1}\ne0$. So $\Delta\in W_i$ is orthogonal to $W_i$, and $\Delta=0$ as $g_i$ is invertible. Hence
$e=u_r=u_0=0$, a contradiction; so $\mu_{p_0}>0$. Finally, with $g=g_{p_0}$, the sum in the definition of
$m_{x,\tau}(e)$ is $\sum_{a,b\le p_0}c_a(g^{-1})_{ab}\bar c_b=\sum_by_b\bar c_b=\bar\mu_{p_0}=\mu_{p_0}$, and
$h(e,e)=0$; so $m_{x,\tau}(e)=2\mu_{p_0}$.

(2) The map $f$ is $1+\alpha ew^{\mathsf T}$ with $w^{\mathsf T}v=h(v,e)$, so $\det f=1+\alpha h(e,e)=1$. Using
$h(e,v)=\overline{h(v,e)}$ and $h(e,e)=0$,
\[
 h(fx_a,fx_b)=h(x_a,x_b)+(\alpha+\bar\alpha)\,\bar c_a\,c_b .
\]
So the leading $i\times i$ block of the Gram matrix of the columns of $fx_\tau$ is $g_i+2\operatorname{Re}(\alpha)\,
\bar c^{\mathsf T}c$, and by the matrix determinant lemma its determinant is
$d_i(1+2\operatorname{Re}(\alpha)\,cg_i^{-1}\bar c^{\mathsf T})=d_i(1+2\operatorname{Re}(\alpha)\,\mu_i)$, by the
computation of $\mu_i$ in (1).
% Lean: u_i is PositiveMinorUnits.pN (GramDeterminantReal.pv with yv = cv (gramF)^-1) and mu_i is mN. (1):
% mN_zero, pN_top and mN_top; mN_succ (pv_pyth) and mN_step (det_step, with the triangular change of basis
% M of determinant yv (Fin.last i)); mN_up, mN_down, mN_bounds; mN_pos (dv_zero, comb_orth_zero); and
% MajorantFamilyNorms.maj_eq with mt_famN and GramDeterminantReal.mt_eq. (2): det_flowMat and
% GramFlowTransform.det_flw (gramF_flw, Matrix.det_add_replicateCol_mul_replicateRow), with mt_famN, as in
% MajorantFamilyNorms.dmin_flow.
\end{proof}

\begin{lemma}\label{is-onto-share-counts-section-cusp-flow}
Let $X$ consist of the data (1) to (6) of flag data of rank $r$ over $E$, let $\xi\in\Lambda$ be isotropic and
$x\in\mathcal D_X$, and use the functions $N_\kappa$ of Definition~\ref{is-onto-share-counts-model}. For
$s\in\mathbf C^\Psi$ let $\phi_\xi(s)x$ be given by the formula of Definition~\ref{is-onto-share-flag-data}, with
complex $s$, and put $q_\tau(s)=1+c_\tau\,m_{x,\tau}(\tau(\xi))\operatorname{Im}s_\tau$ for $\tau\in\Psi$.
\begin{enumerate}
\item $\phi_\xi(s)x\in\mathcal D_X$ if and only if $q_\tau(s)>0$ for every $\tau\in\Psi$.
\item Let $0<p_0<r$, let $\lambda$ be a weight and let $\tau_0\in\Psi$ be such that
  $\lambda_{\tau_0,i+1}\le\lambda_{\tau_0,i}$ for $1\le i<r$ with $i\ne p_0$, and $\lambda_{\tau_0,1}<\lambda_{\tau_0,r}$.
  Let $s\in\mathbf C^\Psi$ with $s_\tau=0$ for $\tau\ne\tau_0$, $\operatorname{Im}s_{\tau_0}\geq0$ and
  $\phi_\xi(s)x\in\mathcal D_X$. Then $N_\lambda(x)\le q_{\tau_0}(s)^{-1/2}\,N_\lambda(\phi_\xi(s)x)$.
\end{enumerate}
\end{lemma}
\leanref{ReciprocityTheory/ReductiveGroups/ArithmeticGeometry/FundamentalDomains/FlowDomainCriterionHolds.lean}{ReciprocityTheory.ChartRestriction.G.FD.flowDomainCriterion_holds}
\leanref{ReciprocityTheory/ReductiveGroups/ArithmeticGeometry/ArakelovVolumes/UpperBounds/CuspStatementPf.lean}{ReciprocityTheory.ChartRestriction.G.AR.Up.FactorialDivisibilityBounds.rfac_le}
\leanref{ReciprocityTheory/ReductiveGroups/ArithmeticGeometry/ArakelovVolumes/UpperBounds/CuspStatementPf.lean}{ReciprocityTheory.ChartRestriction.G.AR.Up.FactorialDivisibilityBounds.abel_bound}
\leanref{ReciprocityTheory/ReductiveGroups/ArithmeticGeometry/ArakelovVolumes/UpperBounds/CuspStatementPf.lean}{ReciprocityTheory.ChartRestriction.G.AR.Up.FactorialDivisibilityBounds.nfac_div}
\begin{proof}
Let $\tau\in\Psi$ and $e=\tau(\xi)$. Then $e\ne0$, and $h_\tau(e,e)=\tau(h(\xi,\xi))=0$. The component $f_\tau$ of
$\phi_\xi(s)$ is the map of Lemma~\ref{is-onto-share-counts-section-cusp-minors}(2) for $\alpha=-ic_\tau s_\tau$, with
$\operatorname{Re}\alpha=c_\tau\operatorname{Im}s_\tau$. Let $\mu_{\tau,i}$ be the numbers of
Lemma~\ref{is-onto-share-counts-section-cusp-minors}(1) for $x$, $\tau$ and $e$. By that lemma, $\det f_\tau=1$ and
\[
 d_i(\phi_\xi(s)x,\tau)=d_i(x,\tau)\,a_{\tau,i},\qquad a_{\tau,i}=1+2c_\tau\operatorname{Im}(s_\tau)\,\mu_{\tau,i},
\]
and $a_{\tau,p_0}=q_\tau(s)$, as $m_{x,\tau}(e)=2\mu_{\tau,p_0}$.

(1) The matrix $f_\tau x_\tau$ is invertible. As $d_i(x,\tau)$ is real with $\epsilon_id_i(x,\tau)>0$ and $a_{\tau,i}$ is
real, $\epsilon_i\operatorname{Re}d_i(\phi_\xi(s)x,\tau)>0$ if and only if $a_{\tau,i}>0$. If $\phi_\xi(s)x\in\mathcal D_X$,
then $q_\tau(s)=a_{\tau,p_0}>0$. Conversely, if $q_\tau(s)>0$, then each $a_{\tau,i}$ lies between $a_{\tau,0}=1$ and
$a_{\tau,p_0}=q_\tau(s)$, as $0\le\mu_{\tau,i}\le\mu_{\tau,p_0}$; so $a_{\tau,i}>0$.

(2) By (1) all $a_{\tau,i}$ are positive, and $a_{\tau,i}=1$ for $\tau\ne\tau_0$. By the definition of $N_\lambda$,
\[
 \log\frac{N_\lambda(x)}{N_\lambda(\phi_\xi(s)x)}=\sum_{i=1}^r\frac{\lambda_i}2\,(L_i-L_{i-1}),
\]
where $\lambda_i=\lambda_{\tau_0,i}$ and $L_i=\log a_{\tau_0,i}$. Here $L_0=L_r=0$, as $\mu_{\tau_0,0}=\mu_{\tau_0,r}=0$.
Summation by parts and $1\le p_0\le r-1$ give
\[
 \sum_{i=1}^r\frac{\lambda_i}2\,(L_i-L_{i-1})=\sum_{i=1}^{r-1}\frac{\lambda_i-\lambda_{i+1}}2\,L_i
 =\frac{\lambda_1-\lambda_r}2\,L_{p_0}+\sum_{i=1}^{r-1}\frac{\lambda_i-\lambda_{i+1}}2\,(L_i-L_{p_0}).
\]
As $\operatorname{Im}s_{\tau_0}\geq0$ and $\mu_{\tau_0,i}\le\mu_{\tau_0,p_0}$, $L_i\le L_{p_0}$; for $i\ne p_0$,
$\lambda_i-\lambda_{i+1}\geq0$, and the term $i=p_0$ vanishes. So the last sum is $\le0$. As the $\lambda_i$ are
integers with $\lambda_1<\lambda_r$, $\lambda_1-\lambda_r\le-1$, and $L_{p_0}=\log q_{\tau_0}(s)\geq0$. Hence the
logarithm is at most $-\frac12\log q_{\tau_0}(s)$.
% Lean: (1) is flowDomainCriterion_holds (minors_iff with MajorantFamilyNorms.dmin_flow, mN_bounds, mN_pos,
% det_flowMat), with ttV = 1/(c maj), so that 1 + Im s / ttV = q. (2) is nfac_div, rfac_of_zero for tau
% different from tau0, and rfac_le (rfac_eq, abel_bound with a_k = la_(k+1)/2, hpa for la_1 - la_r <= -1,
% and hq for L_(p0) = log q).
\end{proof}

\begin{lemma}\label{is-onto-share-counts-section-cusp-constant}
Let $X$ consist of the data (1) to (6) of flag data of rank $r$ over $E$, and use the functions $N_\kappa$ of
Definition~\ref{is-onto-share-counts-model}. Let $\lambda$ be a weight, $P$ a bounded section of weight $\lambda$ and
$C$ a real number with $\|P\|_\lambda\le C$ on $\mathcal D_X$. Let $\xi\in\Lambda$ be isotropic, and let $G(x)$ be
the constant term of $P$ along $\xi$ at $x$. Then:
\begin{enumerate}
\item $|G(x)|\,N_\lambda(x)\le C$ for every $x\in\mathcal D_X$;
\item $G$ is holomorphic on $\mathcal D_X$;
\item $G(\phi_\xi(s)x)=G(x)$ for every $x\in\mathcal D_X$ and every $s\in\mathbf C^\Psi$ with
  $\phi_\xi(s)x\in\mathcal D_X$, where $\phi_\xi(s)x$ is given by the formula of
  Definition~\ref{is-onto-share-flag-data} with complex $s$.
\end{enumerate}
\end{lemma}
\leanref{ReciprocityTheory/ReductiveGroups/ArithmeticGeometry/FundamentalDomains/FlowCoefficientBounds.lean}{ReciprocityTheory.ChartRestriction.G.FD.FlowCoefficientBounds.fj_bound}
\leanref{ReciprocityTheory/ReductiveGroups/ArithmeticGeometry/FundamentalDomains/FlowCoefficientBounds.lean}{ReciprocityTheory.ChartRestriction.G.FD.FlowCoefficientBounds.flow_real}
\leanref{ReciprocityTheory/ReductiveGroups/ArithmeticGeometry/FundamentalDomains/CuspCoefficientFlow.lean}{ReciprocityTheory.ChartRestriction.G.FD.CuspCoefficientFlow.fj_diffOn}
\leanref{ReciprocityTheory/ReductiveGroups/ArithmeticGeometry/FundamentalDomains/CuspCoefficientFlow.lean}{ReciprocityTheory.ChartRestriction.G.FD.CuspCoefficientFlow.diffOn_param_bdd}
\leanref{ReciprocityTheory/ReductiveGroups/ArithmeticGeometry/FundamentalDomains/FjFlow.lean}{ReciprocityTheory.ChartRestriction.G.FD.Data.fj_flow_real}
\leanref{ReciprocityTheory/ReductiveGroups/ArithmeticGeometry/FundamentalDomains/TransvectionPeriodicity.lean}{ReciprocityTheory.ChartRestriction.G.FD.TransvectionPeriodicity.sec_periodic}
\leanref{ReciprocityTheory/ReductiveGroups/ArithmeticGeometry/FundamentalDomains/HolomorphicPeriodicityExtension.lean}{ReciprocityTheory.ChartRestriction.G.FD.holomorphic_periodicity_extension}
\begin{proof}
\emph{Real flows.} Let $s\in\mathbf R^\Psi$. In Lemma~\ref{is-onto-share-counts-section-cusp-minors}(2) for the
components of $\phi_\xi(s)$, $\operatorname{Re}\alpha=0$; so $\phi_\xi(s)x$ has the same leading minors as $x$, it
lies in $\mathcal D_X$ when $x$ does, and $N_\lambda(\phi_\xi(s)x)=N_\lambda(x)$. By
Lemma~\ref{is-eigenvalue-conjugate-dictionary-facts}(4), $\phi_\xi(t)\phi_\xi(s)=\phi_\xi(t+s)$ for
$s,t\in\mathbf R^\Psi$. If $\Pi_\lambda$ has no measurable fundamental domain of finite positive volume, then $G=0$,
and (1) to (3) hold, as $C\geq\|P\|_\lambda\geq0$. Otherwise let $F_\lambda$ be the fixed one, so that
$G(x)=\operatorname{vol}(F_\lambda)^{-1}\int_{F_\lambda}P(\phi_\xi(s)x)\,ds$.

(1) For $s\in\mathbf R^\Psi$, $|P(\phi_\xi(s)x)|\,N_\lambda(x)=\|P\|_\lambda(\phi_\xi(s)x)\le C$. Averaging over
$F_\lambda$ gives $|G(x)|\le C/N_\lambda(x)$.

(2) For each $s\in\mathbf R^\Psi$ the function $y\mapsto P(\phi_\xi(s)y)$ is holomorphic on $\mathcal D_X$, as
$\phi_\xi(s)$ is linear and preserves $\mathcal D_X$, and $(s,y)\mapsto P(\phi_\xi(s)y)$ is continuous. Let
$y_0\in\mathcal D_X$. As $\mathcal D_X$ is open and $N_\lambda$ is continuous and positive on it, there is a closed
ball $B$ around $y_0$ in $\mathcal D_X$ on which $N_\lambda>N_\lambda(y_0)/2$. On $B$,
$|P(\phi_\xi(s)y)|\le C/N_\lambda(y)\le2C/N_\lambda(y_0)$ for every $s$. By the Cauchy estimates on complex lines,
the derivatives in $y$ of these functions are bounded on the ball of half the radius, uniformly in $s$. As
$F_\lambda$ has finite measure, differentiation under the integral sign shows that $G$ is complex differentiable
at $y_0$.

(3) First let $s$ be real. Let $a\in E$ be an algebraic integer with $\bar a=a$, and $p\in\mathbf R^\Psi$ with
$p_\tau=mo_\lambda\tau(a)$, which is real. By (6) of Definition~\ref{is-onto-share-flag-data}, $\Gamma$ contains the
transvection $\gamma\colon v\mapsto v+ma\kappa_0h(v,\xi)\xi$. As $h(\xi,\xi)=0$, $\gamma^{o_\lambda}$ is
$v\mapsto v+o_\lambda ma\kappa_0h(v,\xi)\xi$. As $\bar\xi_0=-\xi_0$, $\tau(\kappa_0)$ is purely imaginary, with
negative imaginary part and absolute value $c_\tau$; so $\tau(\kappa_0)=-ic_\tau$, and $\gamma^{o_\lambda}$ acts at
$\tau$ as the component $f_\tau$ of $\phi_\xi(p)$. Hence for $y\in\mathcal D_X$
\[
 P(\phi_\xi(p)y)=P(\gamma^{o_\lambda}y)=\chi_\lambda(\gamma)^{o_\lambda}P(y)=P(y).
\]
The $p\in\mathbf R^\Psi$ with $P(\phi_\xi(p)y)=P(y)$ for every $y\in\mathcal D_X$ form a group, as the real flows
preserve $\mathcal D_X$ and compose additively; so this group contains $\Pi_\lambda$. Thus $t\mapsto P(\phi_\xi(t)x)$
is periodic under $\Pi_\lambda$, and as the integral of a periodic function over a fundamental domain does not
change under translation,
\[
 G(\phi_\xi(s)x)=\frac1{\operatorname{vol}(F_\lambda)}\int_{F_\lambda}P(\phi_\xi(t+s)x)\,dt
 =\frac1{\operatorname{vol}(F_\lambda)}\int_{F_\lambda}P(\phi_\xi(t)x)\,dt=G(x).
\]

Now let $s$ be complex. For a subset $A\subset\Psi$ consider the claim: $G(\phi_\xi(s)x)=G(x)$ whenever
$\phi_\xi(s)x\in\mathcal D_X$ and $\operatorname{Im}s_\tau=0$ for $\tau\notin A$. For $A=\emptyset$ this is the real
case. Suppose it holds for $A$, let $\tau_0\notin A$, and let $s$ satisfy $\phi_\xi(s)x\in\mathcal D_X$ and
$\operatorname{Im}s_\tau=0$ for $\tau\notin A\cup\{\tau_0\}$. For $z\in\mathbf C$ let $s(z)$ be $s$ with $s_{\tau_0}$
replaced by $z$, and let $V$ be the half plane of the $z$ with $1+c_{\tau_0}m_{x,\tau_0}(\tau_0(\xi))\operatorname{Im}z>0$.
It is open and convex, and it contains $\mathbf R$ and $s_{\tau_0}$. By
Lemma~\ref{is-onto-share-counts-section-cusp-flow}(1), $\phi_\xi(s(z))x\in\mathcal D_X$ for $z\in V$. The entries of
$\phi_\xi(s(z))x$ are polynomials in $z$, so by (2) the function $z\mapsto G(\phi_\xi(s(z))x)-G(x)$ is holomorphic on
$V$; by the claim for $A$ it vanishes on $\mathbf R$. By the identity theorem it vanishes on $V$, in particular at
$z=s_{\tau_0}$. So the claim holds for $A\cup\{\tau_0\}$, and for $A=\Psi$ it is (3).
% Lean: (1) is FlowCoefficientBounds.fj_bound (flow_real, norm_integral_le_of_norm_le_const). (2) is
% CuspCoefficientFlow.fj_diffOn: diffOn_param_bdd (hasFDerivAt_integral_of_dominated_of_fderiv_le, with the
% Cauchy estimate norm_fderiv_le_of_differentiableOn_ball_of_norm_le) and fj_hbd (continuousAt_nfac). (3) for
% real s is Data.fj_flow_real with nu = 0 (flow_flow, TransvectionPeriodicity.sec_periodic through
% sec_periodic_gen, transvMat_pow, emb_kappa and transvMat_map, and fd_translate); for complex s it is
% holomorphic_periodicity_extension with m = 0 (Finset.induction_on, d4l4_diff_flow and
% AnalyticOnNhd.eqOn_zero_of_preconnected_of_frequently_eq_zero), with flowDomainCriterion_holds.
\end{proof}

\begin{lemma}\label{is-onto-share-counts-section-cusp}
Let $X$ consist of the data (1) to (6) of flag data of rank $r$ over $E$ with $0<p_0<r$, let $\lambda$ be a weight,
and let $\tau_0\in\Psi$ be such that $\lambda_{\tau_0,i+1}\le\lambda_{\tau_0,i}$ for $1\le i<r$ with $i\ne p_0$, and
$\lambda_{\tau_0,1}<\lambda_{\tau_0,r}$. Then every bounded section of weight $\lambda$ is a cusp form.
\end{lemma}
\leanref{ReciprocityTheory/ReductiveGroups/ArithmeticGeometry/ArakelovVolumes/UpperBounds/FundSetMeasStatement.lean}{ReciprocityTheory.ChartRestriction.G.AR.Up.CuspStatement}
\leanref{ReciprocityTheory/ReductiveGroups/ArithmeticGeometry/ArakelovVolumes/UpperBounds/CuspStatementPf.lean}{ReciprocityTheory.ChartRestriction.G.AR.Up.cuspStatement_holds}
\begin{proof}
Use the functions $N_\kappa$ of Definition~\ref{is-onto-share-counts-model}. Let $P$ be a bounded section of
weight $\lambda$ and $C$ a real number with $\|P\|_\lambda\le C$ on $\mathcal D_X$. Let $\xi\in\Lambda$ be isotropic,
$x\in\mathcal D_X$, and $G$ the constant term of $P$ along $\xi$, as in
Lemma~\ref{is-onto-share-counts-section-cusp-constant}. Suppose that $G(x)\ne0$, and put $A=|G(x)|\,N_\lambda(x)$,
which is positive. Put $M=\big|\log|C/A|\big|+1$, where $\log0$ is read as $0$; so $M\geq1$.

By Lemma~\ref{is-onto-share-counts-section-cusp-minors}(1) for $e=\tau_0(\xi)$, $m_{x,\tau_0}(\tau_0(\xi))>0$; and
$c_{\tau_0}>0$, as $\xi_0\ne0$. Put $t_0=1/(c_{\tau_0}m_{x,\tau_0}(\tau_0(\xi)))$ and $t=t_0(e^{2M}-1)>0$, and let
$s\in\mathbf C^\Psi$ have $s_{\tau_0}=it$ and $s_\tau=0$ for $\tau\ne\tau_0$. In the notation of
Lemma~\ref{is-onto-share-counts-section-cusp-flow}, $q_{\tau_0}(s)=1+t/t_0=e^{2M}$ and $q_\tau(s)=1$ for
$\tau\ne\tau_0$. By that lemma, $y=\phi_\xi(s)x$ lies in $\mathcal D_X$ and $N_\lambda(x)\le e^{-M}N_\lambda(y)$. By
Lemma~\ref{is-onto-share-counts-section-cusp-constant}(3) and (1), $G(y)=G(x)$ and $|G(y)|\,N_\lambda(y)\le C$.
Hence
\[
 A=|G(y)|\,N_\lambda(x)\le e^{-M}|G(y)|\,N_\lambda(y)\le Ce^{-M}.
\]
If $C\le0$, then $Ce^{-M}\le0<A$. If $C>0$, then $\log(C/A)\le\big|\log(C/A)\big|<M$, so $C/A<e^M$ and
$Ce^{-M}<A$. Either case contradicts $A\le Ce^{-M}$. So the constant terms of $P$ along every isotropic $\xi\in\Lambda$ vanish at every point of
$\mathcal D_X$, and $P$ is a cusp form.
% Lean: cuspStatement_holds: G = fj la xi 0 P; hGb (FlowCoefficientBounds.fj_bound), hGd
% (CuspCoefficientFlow.fj_diffOn), hreal (Data.fj_flow_real) and hcx (holomorphic_periodicity_extension) are
% Lemma is-onto-share-counts-section-cusp-constant; hmaj and hmpos (MajorantFamilyNorms.maj_eq, mN_pos) give
% ttV > 0; hdom (flowDomainCriterion_holds) and hratio (nfac_div, rfac_of_zero, rfac_le) are Lemma
% is-onto-share-counts-section-cusp-flow; M = |log (C / A)| + 1 with Real.log (log |t|, log 0 = 0), h2 is
% A <= C e^-M, and the case split on C <= 0 (le_or_gt C 0) is at the end. Size: about 990 lines of
% Lean in 70 declarations below cuspStatement_holds (text count with the deps tool, which misses the
% dot-notation call fj_flow_real and its periodicity lemmas).
\end{proof}

\begin{lemma}\label{is-onto-share-counts-section-values}
Let $X$ consist of the data (1) to (6) of flag data of rank $r$ over $E$, with archimedean data for $(L,F,T)$ and a
flat model for them, with base points $y_1,\dots,y_\delta$ (Definition~\ref{is-onto-share-counts-model}), such
that $\Gamma$ acts properly discontinuously on $\mathcal D_X$ modulo $\mathcal B^\Psi$
(Lemma~\ref{is-onto-share-counts-section-discontinuous}). Let $G$ be a weight with curvature constants $\nu,m,\eta$
for $G$ and $L$ (Definition~\ref{is-onto-share-counts-section-levi}), such that every bounded section of weight $G$
is a cusp form. Then for every $v\in\mathbf C^\delta$ there is a cusp form $w$ of weight $G$ with $w(y_j)=v_j$ for
$1\le j\le\delta$.
\end{lemma}
\leanref{Development/AbelianVarietyTheory/Proofs/TorusHaarMeasure/InterpWeight.lean}{GlobalTheory.InterpWeight.interp_weight}
\leanref{ReciprocityTheory/ReductiveGroups/ArithmeticGeometry/ArakelovVolumes/UpperBounds/InterpStatementPfA.lean}{ReciprocityTheory.ChartRestriction.G.AR.Up.interp_of_kron}
\leanref{ReciprocityTheory/ReductiveGroups/ArithmeticGeometry/ArakelovVolumes/UpperBounds/InterpStatementPfA.lean}{ReciprocityTheory.ChartRestriction.G.AR.Up.ss_sub_ssgt0}
\begin{proof}
Let $1\le j\le\delta$. Lemma~\ref{is-onto-share-counts-section-separate} gives a bounded section $Q_j$ of weight $G$
with $Q_j(y_j)\ne0$ and $Q_j(y_{j'})=0$ for $j'\ne j$. It is a cusp form by hypothesis. As $S_G(X)$ is a vector space
(Lemma~\ref{is-onto-share-sections}(1)), $w_j=Q_j(y_j)^{-1}Q_j$ is a cusp form of weight $G$, with $w_j(y_{j'})=1$ if
$j'=j$ and $w_j(y_{j'})=0$ otherwise. Then $w=\sum_jv_jw_j$ is a cusp form of weight $G$, and $w(y_j)=v_j$ for every
$j$.
% Lean: interp_weight: interp_of_kron, for each j core (Lemma is-onto-share-counts-section-kronecker), the
% hypothesis PoinStatement, poin_single and poin_zero (Lemma is-onto-share-counts-section-separate), hcusp,
% ss_sub_ssgt0 (SS G is contained in SSgt G 0) and the scalar multiples of siegelSpaceAlgebra_holds.
\end{proof}

\begin{lemma}\label{is-onto-share-counts-section}
Let $X$ be flag data of rank $r$ over $E$ with $0<p_0<r$, and $n=|\Psi|r(r-1)/2$. Let $L,F$ be weights and $T$
an integer that are admissible in the sense of Definition~\ref{is-onto-share-two-level} for this $n$, such that
$X$ carries an archimedean model for $(L,F,T)$. Let $s\le T$ be even and $u\in S_{\delta_s}(X)$ not identically
zero on $\mathcal D_X$. Then there is a bounded section $w$ of weight $(n+4)L-2\delta_s$ that is not identically
zero on $\mathcal D_X$. The proof uses neither the margins of the model, nor the conditions on $F$ in the
admissibility, nor the parity of $s$, nor $u$.
\end{lemma}
\leanref{Development/AbelianVarietyTheory/Proofs/DivisorProofs/ExistsNonzeroBoundedSectionShiftedWeightStatement.lean}{GlobalTheory.T.exists_nonzero_boundedSection_shiftedWeight}
\leanref{Development/AbelianVarietyTheory/Proofs/DivisorProofs/ExistsNonzeroBoundedSectionShiftedWeight.lean}{GlobalTheory.ShareSection.exists_bddSec_shift}
\begin{proof}
Use the notation of Definition~\ref{is-onto-share-counts-model}. First, $\delta\geq1$. Otherwise, choose
$x_1\in\mathcal D_X$ with $\theta_0(x_1)\ne0$ and $k$ with $w_k(x_1)\ne0$. The product $\theta_0w_k^T$ of bounded
sections is a bounded section of weight $F+TL$, nonzero at $x_1$. By Lemma~\ref{is-onto-share-counts-upper} the
family consisting of it alone has at most $0\cdot\binom{T+2n+6}{n}=0$ members, a contradiction.

It remains to find a cusp form $w$ of weight $(n+4)L-2\delta_s$ with $w(y_j)=1$ for every base point $y_j$;
then $w$ is a bounded section and $w(y_1)\ne0$.
% Lean: the proof of delta >= 1 is UpperModel.upper for thp * w k ^ T (BoundedSectionAlgebra.bdd_mul).

Put $G=(n+4)L-2\delta_s$, let $k$ be the weight with $k_{\tau,i}=2i-r-1$, and for an integer $a\geq0$ let
$\mu_a=aL-2\delta_s-k$; then $G-k-nL=\mu_4$ and $G-k-(n+1)L=\mu_3$. As $0<p_0<r$, $r\geq2$; as $n\geq2$ and
$11n^2\le T$ (Definition~\ref{is-onto-share-counts-model}(1)), $\Psi$ is not empty and $T\geq1$. Write
$q=\lfloor s/2\rfloor$, so that $(\delta_s)_{\tau,i}=-q$ for $i\le p_0$ and $q$ for $i>p_0$, and $2q\le s\le T$.

\emph{The weights $\mu_a$.} Let $\tau\in\Psi$. By the admissibility, $L_{\tau,i}-L_{\tau,i+1}\geq1$ for $1\le i<r$
with $i\ne p_0$; so $L_{\tau,i}-L_{\tau,j}\geq j-i$ whenever $i\le j\le p_0$ or $p_0<i\le j$. For such $i\le j$,
$\delta_s$ takes the same value at $i$ and $j$, and $k_{\tau,j}-k_{\tau,i}=2(j-i)$, so
\[
 (\mu_a)_{\tau,i}-(\mu_a)_{\tau,j}=a(L_{\tau,i}-L_{\tau,j})+2(j-i)\geq0.
\]
For $i\le p_0<j$ the same bounds give $L_{\tau,i}\le L_{\tau,1}$ and $L_{\tau,r}\le L_{\tau,j}$, so
$L_{\tau,j}-L_{\tau,i}\geq L_{\tau,r}-L_{\tau,1}\geq1.9T$ by the admissibility, and, as $j-i\le r-1$,
\[
 (\mu_a)_{\tau,j}-(\mu_a)_{\tau,i}=a(L_{\tau,j}-L_{\tau,i})-4q-2(j-i)\geq1.9aT-2T-2(r-1).
\]

\emph{A size bound: $4|\Psi|r^2+2r\le2T$.} Put $e=r-1\geq1$. Then $|\Psi|e\le n$, as $r/2\geq1$. So
$|\Psi|r^2\le4|\Psi|e^2\le4(|\Psi|e)^2\le4n^2$, and $2r=2(e+1)\le4e\le4n\le2n^2$. Hence $4|\Psi|r^2+2r\le18n^2\le2T$.

\emph{Curvature constants.} Use the notation of Definition~\ref{is-onto-share-counts-section-levi}. Let $x_0$ be an
orthonormal point, $z\in\Omega(x_0)$ and $\xi\in\mathbf C^n$, and put
$\mathcal C=\mathcal L^{x_0}_G(z;\xi)-\mathcal L_{\log D_{x_0}}(z;\xi)-n\,\mathcal L^{x_0}_L(z;\xi)$. By
Lemma~\ref{is-onto-share-counts-section-density}, with $c=n$ and with $c=n+1$,
$\mathcal C=\mathcal L^{x_0}_{\mu_4}(z;\xi)$ and $\mathcal C-\mathcal L^{x_0}_L(z;\xi)=\mathcal L^{x_0}_{\mu_3}(z;\xi)$. Put
$m=4\cdot1.9T-2T-2(r-1)$. The bounds for $\mu_4$ and Lemma~\ref{is-onto-share-counts-section-curvature} give
$m\,\mathcal L_{\rho_{x_0}}(z;\xi)\le\mathcal C$. By the size bound, $3\cdot1.9T-2T-2(r-1)\geq0$, so the bounds for
$\mu_3$ and Lemma~\ref{is-onto-share-counts-section-curvature} with $m=0$ give $\mathcal L^{x_0}_{\mu_3}(z;\xi)\geq0$,
that is, $2\cdot\frac12\mathcal L^{x_0}_L(z;\xi)\le\mathcal C$. Finally $\nu=4|\Psi|r^2+1$ satisfies
$4|\Psi|r^2<\nu<m$, by the size bound. So $\nu$, $m$ and $\eta=\frac12$ are curvature constants for $G$ and $L$.

\emph{Bounded sections of weight $G$ are cusp forms.} Choose $\tau_0\in\Psi$. For $1\le i<r$ with $i\ne p_0$,
$\delta_s$ takes the same value at $i$ and $i+1$, so $G_{\tau_0,i}-G_{\tau_0,i+1}=(n+4)(L_{\tau_0,i}-L_{\tau_0,i+1})\geq0$.
And $G_{\tau_0,r}-G_{\tau_0,1}=(n+4)(L_{\tau_0,r}-L_{\tau_0,1})-4q\geq6\cdot1.9T-2T\geq1$, as $n\geq2$ and $T\geq1$. So
Lemma~\ref{is-onto-share-counts-section-cusp} applies to $G$ and $\tau_0$.

\emph{The section.} By Lemma~\ref{is-onto-share-counts-section-discontinuous}, $\Gamma$ acts properly
discontinuously on $\mathcal D_X$ modulo $\mathcal B^\Psi$. So Lemma~\ref{is-onto-share-counts-section-values}, for
$G$ and $v=(1,\dots,1)$, gives the cusp form $w$.
% Lean: exists_bddSec_shift: hR, hc1 (ThetaRecordCompare.two_le_r, card_mul_le_nd), size_le with
% card_mul_le_nd, e4 and e3; hmar with bb = 4 |Psi| r^2 + 1 and m = 4 (1.9 T) - 2 T - 2 (r - 1) from
% curv_and_margin_closed with block_fst, block_snd and cross for a = 4 (drop_of_gaps, spread_of_gaps,
% wt_apply) and IW.leviM_comb with k = n; hstr with eta = 1/2 from curv_and_margin_closed for a = 3 and m = 0
% (cross and the size bound) and IW.leviM_comb with k = n + 1; hcusp from ThetaRecordCompare.cusp_of_gaps with
% gaps_G and pair_G (n + 4 >= 6, T >= 1); then InterpWeight.interp_weight for G and v = 1, with
% ClosedDiscontinuity.propDisc_closed and the proved statements normFactor_rightBorel_holds,
% normFactor_leftAction_holds, pointwiseNorm_invariance_holds, poin_holds, extFStatement_holds,
% measInv_holds, vanishAtS_holds and KapPf.kappa_statement_holds. It takes Fl.xs 0 for y_1.
% Size: nothing below this proof remains open; the last open input, Lemma
% is-onto-share-counts-section-kronecker-solve-chart (the chart Omega(x0) is connected and pseudoconvex), is
% now proved; see the size comment there. The earlier counts of
% 16,700 and 27,000 lines followed tactic and local names (core, ext, b, hy) to unrelated declarations.
\end{proof}

\begin{lemma}\label{is-onto-share-counts-interpolation}
Let $X$ be flag data of rank $r$ over $E$ with $0<p_0<r$, and $n=|\Psi|r(r-1)/2$. Let $L,F$ be weights and $T$
an integer that are admissible in the sense of Definition~\ref{is-onto-share-two-level} for this $n$, such that
$X$ carries an archimedean model for $(L,F,T)$, with the flat model of
Definition~\ref{is-onto-share-counts-model}. Let $s\le T$ be even and $u\in S_{\delta_s}(X)$ not identically zero
on $\mathcal D_X$. Then for every $1\le j\le\delta$ there is a holomorphic function $h_j$ on $\mathcal D_X$ with
$h_j(y_j)=1$ and $h_j(y_{j'})=0$ for $j'\ne j$, such that $h_js^\alpha w_0^{T-|\alpha|}$ is a cusp form of weight
$F+TL-\delta_s$ for every $\alpha\in\mathbf N^n$ with $|\alpha|\le T$, where $s^\alpha=\prod_ks_k^{\alpha_k}$. The proof uses
neither the margins of the model, nor the parity of $s$, nor $u$.
\end{lemma}
\leanref{Development/AbelianVarietyTheory/Proofs/DivisorProofs/LoweredWeightInterpolationAtPointsStatement.lean}{GlobalTheory.T.loweredWeight_interpolation_at_points}
\leanref{Development/AbelianVarietyTheory/Proofs/DivisorProofs/LoweredWeightInterpolationAtPoints.lean}{GlobalTheory.LowerInterp.lowered_interp}
\leanref{Development/AbelianVarietyTheory/Proofs/DivisorProofs/LoweredWeightInterpolationAtPoints.lean}{GlobalTheory.LowerInterp.prod_mem_SS}
\begin{proof}
Fix $j$. There is a bounded section $P$ of weight $F-\delta_s$ with $P(y_j)\ne0$ and $P(y_{j'})=0$ for $j'\ne j$.
To see this, put $\lambda=F-\delta_s$, let $k$ be the weight with $k_{\tau,i}=2i-r-1$, let $F'=F-k-(n+2)L$ as in
Definition~\ref{is-onto-share-two-level}, and for an integer $c\geq0$ let $\mu_c=F'+cL-\delta_s$; then
$\lambda-k-nL=\mu_2$ and $\lambda-k-(n+1)L=\mu_1$. As $0<p_0<r$, $r\geq2$. Write $q=\lfloor s/2\rfloor$, so that
$2q\le s\le T$.

\emph{The weights $\mu_c$.} Let $\tau\in\Psi$. By the admissibility, $F'_{\tau,i}-F'_{\tau,i+1}\geq0$ and
$L_{\tau,i}-L_{\tau,i+1}\geq1$ for $1\le i<r$ with $i\ne p_0$; so $F'_{\tau,j}\le F'_{\tau,i}$ and $L_{\tau,j}\le L_{\tau,i}$
whenever $i\le j\le p_0$ or $p_0<i\le j$. For such $i\le j$, $\delta_s$ takes the same value at $i$ and $j$, so
$(\mu_c)_{\tau,j}\le(\mu_c)_{\tau,i}$. For $i\le p_0<j$ the same bounds give
$F'_{\tau,j}-F'_{\tau,i}\geq F'_{\tau,r}-F'_{\tau,1}\geq1$ and $L_{\tau,j}-L_{\tau,i}\geq L_{\tau,r}-L_{\tau,1}\geq1.9T$, by the
admissibility, so $(\mu_c)_{\tau,j}-(\mu_c)_{\tau,i}\geq1+1.9cT-2q\geq1+1.9cT-s$.

\emph{Curvature constants.} Use the notation of Definition~\ref{is-onto-share-counts-section-levi}. As $r\geq2$,
$|\Psi|(r-1)\le n$ and $r\le2(r-1)$, so $|\Psi|r^2\le4(|\Psi|(r-1))^2\le4n^2$ when $\Psi$ is not empty, and trivially
otherwise. Put $m=1+2\cdot1.9T-s$. As $n\geq2$, $11n^2\le T$ (Definition~\ref{is-onto-share-counts-model}(1)) and
$s\le T$, $4|\Psi|r^2\le16n^2<2T<m$. Put
$\nu=(4|\Psi|r^2+m)/2$, so that $4|\Psi|r^2<\nu<m$. Let $x_0$ be an orthonormal point, $z\in\Omega(x_0)$ and
$\xi\in\mathbf C^n$, and put $\mathcal C=\mathcal L^{x_0}_\lambda(z;\xi)-\mathcal L_{\log D_{x_0}}(z;\xi)-n\,\mathcal L^{x_0}_L(z;\xi)$.
By Lemma~\ref{is-onto-share-counts-section-density}, with $c=n$ and with $c=n+1$,
$\mathcal C=\mathcal L^{x_0}_{\mu_2}(z;\xi)$ and $\mathcal C-\mathcal L^{x_0}_L(z;\xi)=\mathcal L^{x_0}_{\mu_1}(z;\xi)$. The
bounds for $\mu_2$ and Lemma~\ref{is-onto-share-counts-section-curvature} give $m\,\mathcal L_{\rho_{x_0}}(z;\xi)\le\mathcal C$.
As $1+1.9T-s\geq0$, the bounds for $\mu_1$ and Lemma~\ref{is-onto-share-counts-section-curvature} with $m=0$ give
$\mathcal L^{x_0}_{\mu_1}(z;\xi)\geq0$, that is, $2\cdot\frac12\mathcal L^{x_0}_L(z;\xi)\le\mathcal C$. So $\nu$, $m$ and
$\eta=\frac12$ are curvature constants for $\lambda$ and $L$.

\emph{The section.} By Lemma~\ref{is-onto-share-counts-section-discontinuous}, $\Gamma$ acts properly
discontinuously on $\mathcal D_X$ modulo $\mathcal B^\Psi$, and Lemma~\ref{is-onto-share-counts-section-separate}
gives $P$.
% Lean: lowered_interp: e1 and e2; comb_facts s 2 and comb_facts s 1 (mono_of_gaps, pair_of) for mu_2 and
% mu_1; card_sq_le; hmar with bb = (4 |Psi| r^2 + m) / 2 and m = 1 + 2 (1.9 T) - s from
% curv_and_margin_closed and IW.leviM_comb with k = n; hstr with eta = 1/2 from curv_and_margin_closed with
% m = 0 and IW.leviM_comb with k = n + 1; then kron_at for la = F - deMul s and psi = 0 (Lemma
% is-onto-share-counts-section-kronecker), poin_holds with ClosedDiscontinuity.propDisc_closed (Lemmas
% is-onto-share-counts-section-poincare and -discontinuous), and poin_single, chi_ne_zero and poin_zero
% (Lemma is-onto-share-counts-section-separate). Size: nothing below this proof remains open; Lemma
% is-onto-share-counts-section-kronecker-solve-chart, its last open input, is now proved; see the size
% comment there.

Put $h_j=P/P(y_j)$. Let $|\alpha|\le T$. As $T\geq1$, the product $s^\alpha w_0^{T-|\alpha|}$ has at least one factor,
and each factor is a cusp form of weight $L$; so by Lemma~\ref{is-onto-share-sections}(1), applied one factor at a
time, $h_js^\alpha w_0^{T-|\alpha|}$ is a cusp form of weight $F-\delta_s+TL$.
% Lean: prod_mem_SS, with T >= 1 from T_ge and n_ge.
\end{proof}

\begin{lemma}\label{is-onto-share-counts-lower}
Let $X$ be flag data of rank $r$ over $E$ with $0<p_0<r$, and $n=|\Psi|r(r-1)/2$. Let $L,F$ be weights and $T$
an integer that are admissible in the sense of Definition~\ref{is-onto-share-two-level} for this $n$, such that
$X$ carries an archimedean model for $(L,F,T)$, with $\delta$ base points. Let $s\le T$ be even and
$u\in S_{\delta_s}(X)$ not identically zero on $\mathcal D_X$. Then there are $\delta\binom{T+n}{n}$ cusp forms of
weight $F+TL-\delta_s$ whose restrictions to $\mathcal D_X$ are linearly independent.
\end{lemma}
\leanref{Development/AbelianVarietyTheory/Proofs/DivisorProofs/LowerRankCuspFormsDetWeightOneComponentStatement.lean}{GlobalTheory.T.lowerRank_cuspForms_detWeight_oneComponent}
\leanref{Development/AbelianVarietyTheory/Proofs/DivisorProofs/LowerRankCuspFormsDetWeightOneComponent.lean}{GlobalTheory.ShareOn.lowerRank_oneComponent}
\leanref{Development/AbelianVarietyTheory/Proofs/DivisorProofs/LeadingMonomialProductsLinearIndependentStatement.lean}{GlobalTheory.T.leadingMonomialProducts_linearIndependent}
\leanref{Development/AbelianVarietyTheory/Proofs/DivisorProofs/LeadingMonomialProductsLinearIndependent.lean}{GlobalTheory.JetIndepAtPts.leadM_coeff_zero}
\begin{proof}
Use the notation of Definition~\ref{is-onto-share-counts-model}, and let $h_1,\dots,h_\delta$ be the functions of
Lemma~\ref{is-onto-share-counts-interpolation}. For $1\le j\le\delta$ and $\alpha\in\mathbf N^n$ with $|\alpha|\le T$, the
function $g_{j,\alpha}=h_js^\alpha w_0^{T-|\alpha|}$ is a cusp form of weight $F+TL-\delta_s$. There are
$\delta\binom{T+n}{n}$ of them. Let $c_{j,\alpha}\in\mathbf C$ with $\sum_{j,\alpha}c_{j,\alpha}g_{j,\alpha}=0$ on $\mathcal D_X$. By
induction on $d\le T$, $c_{j,\alpha}=0$ for all $j$ and all $\alpha$ with $|\alpha|=d$.

Let $d\le T$ and $1\le j\le\delta$, and suppose $c_{j',\alpha}=0$ for all $j'$ and all $\alpha$ with $|\alpha|<d$. Near $0$,
$c_{y_j}(z)\in\mathcal D_X$ and $w_0(c_{y_j}(z))\ne0$. Put $\zeta(z)=\big(s_k(c_{y_j}(z))/w_0(c_{y_j}(z))\big)_k$; then
$\zeta(z)=Az+O(|z|^2)$ with $A$ invertible, and
$g_{j',\alpha}(c_{y_j}(z))=h_{j'}(c_{y_j}(z))\,w_0(c_{y_j}(z))^T\zeta(z)^\alpha$. For $|\alpha|>d$ this is $O(|z|^{d+1})$. For
$|\alpha|=d$ it is $h_{j'}(y_j)w_0(y_j)^T(Az)^\alpha+O(|z|^{d+1})$, and $h_{j'}(y_j)$ is $1$ if $j'=j$ and $0$ otherwise. So
\[
 w_0(y_j)^T\sum_{|\alpha|=d}c_{j,\alpha}(Az)^\alpha=O(|z|^{d+1}).
\]
A homogeneous polynomial of degree $d$ that is $O(|z|^{d+1})$ at $0$ is zero, and the monomials $(Az)^\alpha$ with
$|\alpha|=d$ are linearly independent, as $A$ is invertible. Since $w_0(y_j)\ne0$, $c_{j,\alpha}=0$ for $|\alpha|=d$.
% Lean: lowerRank_oneComponent: T.loweredWeight_interpolation_at_points (Lemma
% is-onto-share-counts-interpolation) for the h_j, the index set Fin delta x multiLt n (T + 1) counted by
% Y.ind_count, and T.leadingMonomialProducts_linearIndependent = JetIndepAtPts.leadM_coeff_zero (strong
% induction on d, with DivLC.coordL for A, expand and high_O for the expansions, homog_L at the end).
\end{proof}

\begin{lemma}\label{is-onto-share-counts}
Let $X$ be flag data of rank $r$ over $E$ with $0<p_0<r$, and $n=|\Psi|r(r-1)/2$. Let $L,F$ be weights and
$T$ an integer that are admissible in the sense of Definition~\ref{is-onto-share-two-level} for this $n$,
such that $X$ carries an archimedean model for $(L,F,T)$. Let $s\le T$ be even and $u\in S_{\delta_s}(X)$
not identically zero on $\mathcal D_X$, and put $\kappa=F+TL$. Then $n\geq2$, $11n^2\le T$, and there is an
integer $\delta\geq0$ such that:
\begin{enumerate}
\item every family of cusp forms of weight $\kappa$, or of weight $\kappa-\delta_s$, whose restrictions to
  $\mathcal D_X$ are linearly independent has at most $\delta\binom{T+2n+6}{n}$ members;
\item every such family of weight $\kappa+2\delta_s$ has at most $\delta\binom{T+3n+10}{n}$ members;
\item there are $\delta\binom{T+n}{n}$ cusp forms of weight $\kappa-\delta_s$ whose restrictions to
  $\mathcal D_X$ are linearly independent.
\end{enumerate}
\end{lemma}
\leanref{Development/AbelianVarietyTheory/Proofs/DivisorProofs/TwoLevelFlagDomainCuspFormRanksShiftedWeights.lean}{GlobalTheory.ShareRanks.compT}
\leanref{Development/AbelianVarietyTheory/Proofs/DivisorProofs/TwoLevelFlagDomainCuspFormRanksShiftedWeights.lean}{GlobalTheory.ShareRanks.compX}
\leanref{Development/AbelianVarietyTheory/Proofs/DivisorProofs/TwoLevelFlagDomainCuspFormRanksShiftedWeights.lean}{GlobalTheory.ShareRanks.indep_mul}
\begin{proof}
The bounds $n\geq2$ and $11n^2\le T$ are among the conditions of the archimedean model
(Definition~\ref{is-onto-share-counts-model}(1)), and $\delta$ is the number of base points of its flat model.
% Lean: n_ge and T_ge are fields of ArchInM, and delta is Fl.delta of the flat model FlatM. The counts are
% UpperModel.upper and T.card_independent_boundedSections_shiftedWeight_le (Lemma
% is-onto-share-counts-upper), T.exists_nonzero_boundedSection_shiftedWeight (Lemma
% is-onto-share-counts-section) and T.lowerRank_cuspForms_detWeight_oneComponent (Lemma
% is-onto-share-counts-lower).

If $f_1,\dots,f_m$ are bounded sections with independent
restrictions to $\mathcal D_X$ and $w$ is a section not identically zero on $\mathcal D_X$, then the products
$f_qw$ have independent restrictions. Indeed, if $\sum_qa_qf_qw=0$ on $\mathcal D_X$, the holomorphic
function $\sum_qa_qf_q$ vanishes on the nonempty open set where $w\ne0$, hence on $\mathcal D_X$ by the
identity theorem, as $\mathcal D_X$ is open and connected (Lemma~\ref{is-onto-share-sections}(2)); so all
$a_q$ vanish.

(1) For weight $\kappa$ this is the first bound of Lemma~\ref{is-onto-share-counts-upper}, as cusp forms
are bounded sections. Let
$f_1,\dots,f_m$ be cusp forms of weight $\kappa-\delta_s$ with independent restrictions. As $u$ is a bounded
section of weight $\delta_s$, the $f_qu$ are cusp forms of weight $\kappa$ by
Lemma~\ref{is-onto-share-sections}(1), with independent restrictions. So $m\le\delta\binom{T+2n+6}{n}$.

(2) Lemma~\ref{is-onto-share-counts-section} gives a bounded section $w$ of weight $(n+4)L-2\delta_s$ that is
not identically zero on $\mathcal D_X$. Let $f_1,\dots,f_m$ be cusp forms of weight $\kappa+2\delta_s$ with
independent restrictions. The
products $f_qw$ are bounded sections of weight $\kappa+2\delta_s+(n+4)L-2\delta_s=F+(T+n+4)L$, as in the
proof of Lemma~\ref{is-onto-share-sections}(1), with independent restrictions, so
$m\le\delta\binom{T+3n+10}{n}$ by the second bound of Lemma~\ref{is-onto-share-counts-upper}.

(3) This is Lemma~\ref{is-onto-share-counts-lower}.
\end{proof}

\begin{lemma}\label{is-onto-share-binomial}
Let $T$ and $n\geq2$ be integers with $11n^2\le T$. Then
\[
 \binom{T+3n+10}{n}\le2\binom{T+n}{n}\quad\text{and}\quad
 \Big(1-\frac{n(n+6)}T\Big)\binom{T+2n+6}{n}\le\binom{T+n}{n}.
\]
\end{lemma}
\leanref{Development/AbelianVarietyTheory/Proofs/DivisorProofs/TwoBinomialRatioBoundsStatement.lean}{GlobalTheory.T.two_binomial_ratio_bounds}

\begin{lemma}\label{is-onto-share-ranks}
Let $M$ be two-level data of rank $r$ over $E$ for a finite group $U$ acting on a finite set $C$
(Definition~\ref{is-onto-share-two-level}) with $0<p_0<r$, for which the level hypothesis holds for the
trivial subgroup. Let $\kappa$ be a weight and $T$ an integer such that $M$ satisfies the archimedean
condition at $T$ for $\kappa$ and $s(M)\le T$, and let $n=\dim Y(M)$. Then the cusp forms of $M$ have ranks
$(n(n+6)/T,\,2)$ at $\kappa$ and $s(M)$.
\end{lemma}
\leanref{Development/AbelianVarietyTheory/Proofs/DivisorProofs/TwoLevelFlagDomainCuspFormRanksShiftedWeightsStatement.lean}{GlobalTheory.T.twoLevelFlagDomain_cuspForm_ranks_shifted_weights}
\leanref{Development/AbelianVarietyTheory/Proofs/DivisorProofs/TwoLevelFlagDomainCuspFormRanksShiftedWeights.lean}{GlobalTheory.ShareRanks.ranks}
\begin{proof}
Let $s=s(M)$ and $\kappa=F+TL$ as in the archimedean condition. Lemma~\ref{is-onto-share-translates} for
$M|_{\{1\}}$ gives a family $u_0$ of weight $\delta_s$ with $(1^*u_0)_c(x_c)=u_{0,c}(t(1,c)x_c)\ne0$ for
some $x_c\in\mathcal D_M$, for every $c$. Since $t(1,c)\in\Gamma_c$ and $\chi_{c,\delta_s}$ is trivial on
$\Gamma_c$, $u_{0,c}(x_c)\ne0$. Lemma~\ref{is-onto-share-counts} for $X_c$, with this $u=u_{0,c}$, gives
$n\geq2$, $11n^2\le T$ and an integer $\delta_c$ with the three counts.

For a weight $\kappa'$, an element of $S_{\kappa'}(M)$ restricts on each $\{c\}\times\mathcal D_M$ to an
element of the span of the restrictions of $S_{\kappa'}(X_c)$. If every family of elements of
$S_{\kappa'}(X_c)$ with independent restrictions has at most $N_c$ members, then $S_{\kappa'}(M)$ is
finite-dimensional with $\dim S_{\kappa'}(M)\le\sum_cN_c$, as restriction to the components is injective.
Conversely, cusp forms of $X_c$ extended by $0$ off $c$ give, by the third count,
$\sum_c\delta_c\binom{T+n}{n}\le\dim S_{\kappa-\delta_s}(M)$. So the three spaces are finite-dimensional,
and
\[
 \dim S_\kappa(M)\le\sum_c\delta_c\binom{T+2n+6}{n},\qquad
 \dim S_{\kappa+2\delta_s}(M)\le\sum_c\delta_c\binom{T+3n+10}{n}.
\]
By Lemma~\ref{is-onto-share-binomial}, $\dim S_{\kappa+2\delta_s}(M)\le2\sum_c\delta_c\binom{T+n}{n}\le
2\dim S_{\kappa-\delta_s}(M)$. If $1-n(n+6)/T\le0$, the remaining inequality is clear. Otherwise
\[
 \Big(1-\frac{n(n+6)}T\Big)\dim S_\kappa(M)\le\sum_c\delta_c\Big(1-\frac{n(n+6)}T\Big)\binom{T+2n+6}{n}
 \le\sum_c\delta_c\binom{T+n}{n}\le\dim S_{\kappa-\delta_s}(M).
\]
\end{proof}

\subsection{Runs and level structures}\label{subsection-is-runs}

\begin{definition}\label{is-onto-step-runs}
Let $p$ be a point of a family at $(x,i)$ (Definition~\ref{is-point}), with definite-side datum $D$
of host field $E_h/F_h$, coefficient field $L$ and lifts $1,\dots,m$, and with realization $R$, group
$G$, set $P_g$ of places of $E_h$, one-map data with weight families $\kappa_i$
(Definition~\ref{is-one-map-data}), and the primes $p_{de}\ne p'_{de}$ of the host.
\begin{enumerate}
\item The \emph{own prime} of a finite place $v$ of $L$ is $p'_{de}$ if $v$ lies over $p_{de}$, and
  $p_{de}$ otherwise. For a twisted place $v$, $\delta_v$ is the character $\delta'$ of $R$ if $v$
  lies over $p_{de}$, and $\delta$ otherwise.
\item A \emph{run} at $p$ chooses, for the core of $D$, at every level $l$ data for its
  $\nabla$-lattices and, for every integer $n\geq0$, delta data $\delta_l(n)$ with a prime
  $p_d(\delta_l(n))$; a character $\xi$ of $(\mathbf A_{F_h}\otimes_{F_h}E_h)^\times$ with values in
  $\mathbf C^\times$; and an ideal $\mathfrak h$ of $\mathcal O_{E_h}$. A \emph{pair of runs}
  $r=(r_1,r_2,\ell)$ consists of two runs and, for every lift $i$, a map $\ell_i$ from $G$ to the
  $\mathcal O_{E_h}$-linear automorphisms of the level module of the core at level $i+1$. The
  \emph{run used at} a finite place $v$ is $r_2$ if $v$ lies over $p_{de}$, and $r_1$ otherwise.
\item A run $s$ satisfies the \emph{laws} if its $\nabla$-data, its lattice data and its short
  generators satisfy the laws of the core, its quotient data are canonical, and at every level $l$ of
  a lift and every $n$ the level character of $\delta_l(n)$ is trivial, the ideal of $\delta_l(n)$ is
  $\mathfrak h$, the short generators of $\delta_l(n)$ do not depend on the choice of its
  $\nabla$-data, and $\xi$ takes at $\delta_l(n)$ the values that $\delta_l(n)$ prescribes. It is
  \emph{covered} at the lift $i$ if every $\delta_{i+1}(n)$ is $\delta_{i+1}(0)$ with other
  coefficients, and \emph{rational} at $i$ if $\xi(\det\gamma)=1$ for every rational point $\gamma$
  of the unitary group of its space at level $i+1$ and $\sigma_0$.
\item A run $s$ is \emph{pinned} at a twisted place $v$ and a lift $i$ if, with $q$ the own prime of
  $v$ and $n$ the integer of the weight family $\kappa_i$: $p_d(\delta_{i+1}(n))=q$; $\mathfrak h$ is
  contained in every place of $E_h$ over $q$ in the ramified set $Q_{\mu_v}$ of the tuple $\mu_v$; and
  there is a unitary character $\chi$ of $\mathbf A_{E_h}^\times$ equal to $\delta_v$ with
  $\xi(\det g)=\chi(y)$ whenever $y$ is a determinant root of an element $g$ of the adelic unitary
  group of $W_i$.
\item A pair of runs $r$ is \emph{pinned to the own places} if the following hold for every finite
  place $v$ and every lift $i$, with $s$ the run used at $v$:
  \begin{enumerate}
  \item $s$ satisfies the laws, and is covered and rational at $i$;
  \item if $v$ is twisted, $s$ is pinned at $v$ and $i$;
  \item (Level.) the level of $R$ at the lift $i$ is the level module of the core at level $i+1$,
    on which $G$ acts through $\ell_i$;
  \item (Action.) there is an action of $G$ on the cusp forms of $s$ at level $i+1$ and weight
    $\kappa_i$ that is the level action of $\ell_i$ and that the map of the weight tie carries to the
    action of $G$ on the rational structure $FW_i$;
  \item the level ideal of the core is $\prod_{P\in P_g}P$;
  \item (Own places.) $p_d(\delta_l(n))$ is the own prime of $v$ for every level $l$ of a lift and
    every $n$.
  \end{enumerate}
\end{enumerate}
\end{definition}
\leanref{Development/AbelianVarietyTheory/Definitions/Divisors/RealisedChainTorusEq.lean}{GlobalTheory.BigOn.ownAt}
\leanref{Development/AbelianVarietyTheory/Proofs/DivisorProofs/NotOntoOfCharAlone.lean}{GlobalTheory.RealIndexBound.RunData}
\leanref{Development/AbelianVarietyTheory/Definitions/Divisors/RealisedChainTorusEq.lean}{GlobalTheory.BigOn.TwoRuns}
\leanref{Development/AbelianVarietyTheory/Definitions/Divisors/RealisedChainTorusEq.lean}{GlobalTheory.BigOn.TwoRuns.used}
\leanref{Formalism/LocalToGlobal/ProjectiveLattices/FullExists.lean}{Formalism.L2.Pj.LawsFull}
\leanref{Development/AbelianVarietyTheory/Definitions/Jets/EmbShortThLeThetaPlain.lean}{GlobalTheory.FormTies.OneCover}
\leanref{Development/AbelianVarietyTheory/Definitions/Foundations/ThetaSeriesTies.lean}{GlobalTheory.OnTies.RunPin}
\leanref{Development/AbelianVarietyTheory/Definitions/Foundations/ThetaSeriesTies.lean}{GlobalTheory.OnTies.LevelPin}
\leanref{Development/AbelianVarietyTheory/Definitions/Foundations/ThetaSeriesTies.lean}{GlobalTheory.OnTies.ActionPin}
\leanref{Development/AbelianVarietyTheory/Definitions/Divisors/IsDeltaDelAtOwnAt.lean}{GlobalTheory.BigOn.OwnPin}

\begin{lemma}\label{is-onto-step-links-scalars}
Let $R$ be a realization of a definite-side datum with a face $F$ (Definition~\ref{is-realisation}),
with complex data as in Definition~\ref{is-complex-data}: the theta lifts $\mathrm{Th}_t$, the parts
$E_p\subset E$ and $W_p(i)\subset W_i$, the maps $\mathrm{th}_G(i,t)$, the maps $t_b(i,t)$ and the
vectors $w_{b,i,v}(k)$, the matrices $\theta_{i,t}$ of the datum at level $b$ and its metrics
$X_T(i,t)$; and let $\ell_6^m$ be the function of the scale data of $F$. Write $w_{b,v}(y)$ for
$\sum_ky_kw_{b,i,v}(k)$, and say that the part $p$ is \emph{visible} at $(b,i,v)$ if some $x$ with
$t_b(i,\mu_0)(\sigma_0x)\in W_p(i)$ has $|x|_{X_T(i,\mu_0),v}\ne0$. Then there are real numbers
$\lambda_{i,t,p,v}$ and $\kappa$ such that $R$ with these scalars satisfies the following
\emph{conventions for the scalars} at $F$.
\begin{enumerate}
\item If $W_p(i)\ne0$ and $\mathrm{th}_G(i,t)\,\mathrm{th}_G(i,t)^*=l$ on $W_p(i)$, then
  $\lambda_{i,t,p,v_0}=l$; if $W_p(i)=0$, then $\lambda_{i,t,p,v_0}=e^{\ell_6^m(i)}$.
\item Let $v\ne v_0$. Then $\lambda_{i,t,p,v}=\lambda_{i,\mu_0,p,v}$ for every tuple $t$. If $p$ is
  visible at $(b,i,v)$, either $b=0$ or $p$ is not visible at $(0,i,v)$, and $\langle
  w_{b,v}(\theta_{i,\mu_0}x),w_{b,v}(\theta_{i,\mu_0}y)\rangle=l\,\langle x,y\rangle_{X_T(i,\mu_0),v}$ for
  all $y$ and all $x$ with $t_b(i,\mu_0)(\sigma_0x)\in W_p(i)$, then $\lambda_{i,\mu_0,p,v}=l$. If $p$ is
  visible at neither level, then $\lambda_{i,\mu_0,p,v}=e^{\ell_6^m(i)}$.
\item If $\kappa'$ is a real number with $\mathrm{Th}_t\mathrm{Th}_t^*=\kappa'\prod_j\lambda_{j,t,p,v_0}$
  on $E_p$ for all $t$ and $p$, the product running over the lifts, and $E_p\ne0$ and
  $\prod_j\lambda_{j,t,p,v_0}\ne0$ for some $t$ and $p$, then $\kappa=\kappa'$.
\end{enumerate}
\end{lemma}
\leanref{Development/AbelianVarietyTheory/Proofs/IwasawaDecomposition/ExistsThetaLiftAdjointScalarOnPart.lean}{GlobalTheory.ScalChoice.exists_isScal}
\leanref{Formalism/LocalToGlobal/IndefiniteReal/LinksChar.lean}{Formalism.IndefR.Links.RealDat.IsScal}
\begin{proof}
Each convention fixes one scalar under a condition that allows at most one value, and the
conventions for different scalars are independent, except that (3) uses the scalars at $v_0$; so
choose the scalars at $v_0$, then those at the other places, then $\kappa$. In (1), if $W_p(i)\ne0$, a
nonzero vector of $W_p(i)$ determines $l$; take this $l$ if it exists, any value otherwise, and
$e^{\ell_6^m(i)}$ if $W_p(i)=0$. In (2), if $p$ is visible at $(0,i,v)$, a vector $x$ with
$|x|_{X_T(i,\mu_0),v}\ne0$, taken with $y=x$, determines $l$ at $b=0$; take this $l$ if it exists. If
$p$ is visible at $(1,i,v)$ only, do the same at $b=1$; if at neither, take $e^{\ell_6^m(i)}$; and
put $\lambda_{i,t,p,v}=\lambda_{i,\mu_0,p,v}$. In (3), with the scalars at $v_0$ chosen, a pair $(t,p)$
with $E_p\ne0$ and $\prod_j\lambda_{j,t,p,v_0}\ne0$ and a nonzero vector of $E_p$ determine $\kappa'$;
take it if it exists, and any value otherwise.
% Lean: exists_isScal (exists_thetaLift_adjoint_scalar_on_part and exists_eig for (1), exists_lamOff
% with OffEq.unique for (2), exists_pin with Lift.IsKap.unique for (3)). The six clauses are the fields
% lam0, zero_part_eigenvalue_abs, lam_tup, lam_off, lam_abs and kap of IsScal; PartAt is visibility.
\end{proof}

\begin{lemma}\label{is-onto-step-links}
Let $p$ be a point of a family (Definition~\ref{is-point}) that satisfies the point conditions, and
suppose that every finite place $w$ at which, for some lift $j$, one of the indices or local terms
listed in the further conditions is nonzero is twisted. Let $R$ be the realization of $p$ with the
scalars of its characterization, and $F$ its face. Then $R$ and $F$ have the following \emph{links}.
\begin{enumerate}
\item $R$ satisfies its characterization at $F$ (Definition~\ref{is-realisation-characterization}), and
  its family of definite forms satisfies the family conditions.
\item $R$ satisfies condition (11) of Definition~\ref{is-primitive-rows}, and its rows $S_0(R,F)$
  have the base links.
\item (Rank.) The rank $D$ of the rows is positive, and $d_W(j)=D$ for every lift $j$.
\item The one-map data of $R$ have a conjugation tie and are pinned at every lift
  (Definition~\ref{is-one-map-data}).
\item (Phases.) For every lift $i$, if the $i$-th lift is given by the one-map data at $i$ as one map
  for some constants and phase, then it is given by them as one map with the canonical constants of
  the core and the phase $1$.
\end{enumerate}
\end{lemma}
\leanref{Development/AutomorphicArithmetic/ThetaCorrespondence/ThetaPairLifts/ThirdPlaceSeries/Links2OfRowsStatement.lean}{GlobalTheory.T.links2_of_single_index_twisted}
\leanref{Development/AbelianVarietyTheory/Statements/PlaceShiftStatements/SingleIndexPlacesAreTwisted.lean}{GlobalTheory.BigOn.Links2}
\begin{proof}
Write $R_p$ for the realization of $p$ with its own scalars and $F_p$ for the face of $p$. The face $F$
is $F_p$ with the scale data $Z(p)$ of the point, and $R$ is $R_p$ with its scalars
$\lambda_{i,t,q,v}$ and $\kappa$ replaced by scalars for which it satisfies its characterization at
$F$, if there are such. The core conditions (Definition~\ref{is-eigenvalue-lift-data}), the rational
structures $FW_i$, the one-map data with their ties and pins, the dimensions and the indices of the rows
at the base level, and the base links involve neither the scalars nor the scale data; so they are the
same for $(R_p,F_p)$, $(R_p,F)$ and $(R,F)$.

Among the data of the point are the weak characterization of $R_p$ at $F_p$, the phases $1$ of $R_p$,
the conjugation tie and the pins of its one-map data, and a \emph{phase tie}: for every lift $i$, if
$R_p$ satisfies (11) of Definition~\ref{is-primitive-rows} and the one-map data at $i$ give the $i$-th
lift as one map for some constants and phase, they give it with the canonical constants of the core
and the phase $\phi_i$ of $R_p$, replaced by $\bar\phi_i$ if $\sigma_0$ is not the embedding of the
base place $v_0$. The point conditions give the existence conditions for the theta maps of $R_p$, the
family conditions with the conditions on the inert indices, lift-distinct types, and primitive rows
of $R_p$ at $F_p$. So Lemma~\ref{is-eigenvalue-base-inputs}(1) applies to $R_p$ and $F_p$: $R_p$
satisfies the core conditions at $F_p$, and the rows $S_0(R_p,F_p)$ have the base links.

(1) The characterization at $F$ (Definition~\ref{is-realisation-characterization}) consists of the core
conditions, the conventions for the scalars at $F$ (Lemma~\ref{is-onto-step-links-scalars}) and
conditions on the tuples. The core conditions hold.
Lemma~\ref{is-onto-step-links-scalars} gives scalars for which $R_p$ satisfies the conventions at $F$,
and neither the core conditions nor the conditions on the tuples involve the scalars. So once the
conditions on the tuples are checked for $R_p$ at $F$, $R_p$ with these scalars satisfies the
characterization at $F$, and therefore so does $R$. The conditions on the tuples other than the one on
the twisted places are conditions of the weak characterization. The remaining one asks that a finite
place $w$ be twisted if, for $w'=w$ or $w'=\bar w$ and some lift $j$, one of $I_{w'}(j,\mu_0)$,
$I_{\bar w'}(j,\mu_0)$, the relative index of $T_{j,\mu_0}$ and its dual for $H_T(j,\mu_0)$ at $w'$,
$F_0(w',j)$ and $F_0(\bar w',j)$ is nonzero, for the rows $S_0(R_p,F)$ and the scale data of $F$.
These indices are those of the rows $S_0(p)$, since the phases of $R_p$ are $1$, and $F_0$ is that of
$N_p(0,0)$; so each of these numbers is one of those listed in the further conditions, at $w$ or at
$\bar w$, and the hypothesis makes $w$ twisted. The family conditions with the conditions on the inert
indices contain the family conditions.

(2) Condition (11) is one of the primitive rows of $R_p$ and concerns only the $FW_i$. The base links of
$S_0(R_p,F_p)$ are those of $S_0(R,F)$.

(3) Lemma~\ref{asm-primitive-rows-base} applies to $R_p$ and $F_p$, by the family conditions with the
conditions on the inert indices, the lift-distinct types, the weak characterization and the primitive
rows. In its notation, the matrices $\theta_{j,\mu_0}$ and $\iota_{\mu_0}=1$ are bijective,
$\iota_{\mu_0}Q_p(\mu_0)=E^L_p$, $\theta_{j,\mu_0}Z_p(j,\mu_0)=W^L_p(j)$ and $\dim
W^L_p(j)=\dim E^L_p$, and the $Q_p(\mu_0)$, like the $Z_p(j,\mu_0)$, are independent and span. So the
$E^L_p$ are independent and span $L^D$, the $W^L_p(j)$ are independent and span $L^{d_W(j)}$, and
\[
 d_W(j)=\sum_p\dim W^L_p(j)=\sum_p\dim E^L_p=D.
\]
By Definition~\ref{is-complex-data}(a), $D=d_Q(\mu_0)$, which is positive by the positivity (7) of the
primitive rows.

(4) The conjugation tie and the pins of the one-map data are data of the point.

(5) Since the phases of $R_p$ are $1$, the phase tie gives the phase $1$, and (11) holds by (2).
% Lean: T.links2_of_single_index_twisted = BigOn.links2_of_rows_alt, from
% T.l_unif_base_links_rank_of_pt_split (SplitRows.lUnif_of_rows_alt, base_links_of_split_data,
% rank_conditions_of_split_data), T.bar_tie_om_pins_phase_tie_of_pt_split (ompins_onto_alt:
% conjugateReadingsVanishing_splitPart, thetaRowData_pins_splitPart,
% phaseCompatible_splitPart_of_liftUniformisation) and T.char_R2_of_exAt_of_singleIndexPlacesTwisted
% (splitPartRealisationWithScalars_char_of_weak, with ScalChoice.exists_isScal, char_charOf,
% characterisation_of_weak). R is R2 p = charOf p.R (F2 p), the Classical.epsilon choice of scalars
% with Char (F2 p) (scalOf); F2 p is p.F.lv2 with Zsc the numbers of p.Z. In (1), char_charOf takes
% lev, grp, fam, models, chain and thP of RealDat3.characterisation_of_weak p.charW hr.exAt, the
% tuple fields of the weak characterization, and tup_tw (F2 p) from
% T.splitPart_nonzeroIndexPlaces_twisted (tupTw_of_w1Row, w1_of_inW1_split_part with
% splitPartFace_baseLevel_indices_eq and _jw_eq); reScal_nonzeroIndexPlace_iff and the isX_reScal
% lemmas say the scalars do not enter. fam.Ok 1 is hr.ok.ok2. (2) is hr.prim.unif and, for the base
% links, complex_data_conditions_base_level_of_primitive_links (lifts0, ratTh, proviso0) with
% CharRealise.toCx3_eq_one p.R p.F p.phsOne, as in Lemma is-eigenvalue-base-inputs(1). (3) is
% indefiniteDataAtBaseLevel: IndefRows.D_pos and dW_eq. (4) is p.bar and p.pins; (5) is p.phase with
% p.phEff_one.
\end{proof}

\begin{lemma}\label{is-onto-step-runs-at-prime}
Let $p$ be a point of a family whose realization and face have the links of
Lemma~\ref{is-onto-step-links} and whose core ties hold (the further conditions,
Definition~\ref{is-point}), and let $l\in\{p_{de},p'_{de}\}$. Then there is a run $s$ at $p$
(Definition~\ref{is-onto-step-runs}) that satisfies the laws, is covered and rational at every lift,
has $p_d(\delta_j(n))=l$ for every level $j$ of a lift and every $n$, and is pinned at $v$ and $i$ for
every twisted place $v$ whose own prime is $l$ and every lift $i$.
\end{lemma}
\leanref{Development/AbelianVarietyTheory/Proofs/DivisorProofs/ExistsParameterRunWithLawsAndPinsAtSpecialPrimeStatement.lean}{GlobalTheory.T.exists_parameterRun_with_laws_and_pins_at_special_prime}
\begin{proof}
\openstep{The construction of a run at the special prime $l$ that satisfies the laws and is pinned at
the twisted places with own prime $l$ is not written out.}
% Lean: T.exists_parameterRun_with_laws_and_pins_at_special_prime = RunOnto.exists_run_pins.
% Size: about 394,000 lines of Lean in 22,000 theorems below it, by a text count that stops at
% statements the book already cites; about 9,000 of these lines lie below no other open step of
% Lemma is-onto-estimates.
\end{proof}

\begin{lemma}\label{is-onto-step-level-structure-congruence}
Let $p$ be a point of a family, with realization $R$, group $G$, host field $E_h$ and rank $r$, let
$i$ be a lift, and use the notation of Definition~\ref{is-onto-step-model-level-condition}: the level
$M$, $M_b$ of the core, the lattice $\Lambda=\Lambda_{i+1}$ with its map $\iota=\iota_{i+1}$, the frame
$\beta_i$, the closure $\hat L_0$ of the lattice of $W_i$, the group $\hat L_W$ and level lifts. Put
$\bar M_b$ for the complex conjugate of $M_b$. Suppose that:
\begin{enumerate}
\item $\hat L_0$ consists of the $z$ with $\beta_i^{-1}z\in\hat{\mathcal O}_{E_h}^{r}$;
\item $\iota$ is injective, and its image is the set of $x\in E_h^{r}$ whose image in
  $(E_h\otimes_{F_h}\mathbf A_{F_h,f})^{r}$ lies in $\hat L_W$;
\item $M_b=\prod_{P\in P_g}P$.
\end{enumerate}
Let $\lambda_g$, for $g\in G$, be $\mathcal O_{E_h}$-linear endomorphisms of $\Lambda$ such that every
$g\in G$ has a level lift $k$ with $\iota(\lambda_g(a))-k\,\iota(a)\in M\hat L_W$ for every
$a\in\bar M_b\Lambda$. Then for all $g,h\in G$ and $a\in\Lambda$,
\[
 \lambda_1(a)-a\in M_b\Lambda\quad\text{and}\quad\lambda_{gh}(a)-\lambda_g(\lambda_h(a))\in M_b\Lambda .
\]
\end{lemma}
\leanref{Development/AbelianVarietyTheory/Proofs/FoundationProofs/MulVecMemOfFrame.lean}{GlobalTheory.LevelActionPinned.congr_of_levIs}
\leanref{Development/AbelianVarietyTheory/Proofs/FoundationProofs/MulVecMemOfFrame.lean}{GlobalTheory.LevelActionPinned.congr_of_lifts}
\leanref{Development/AbelianVarietyTheory/Proofs/FoundationProofs/MulVecMemOfFrame.lean}{GlobalTheory.LevelActionPinned.congr_of_matrices}
\begin{proof}
Write $\hat L=\hat L_W$, $\hat{\mathcal O}=\hat{\mathcal O}_{E_h}$, and read $\iota(a)$ in
$(E_h\otimes_{F_h}\mathbf A_{F_h,f})^{r}$. By (2), $\iota(\Lambda)\subset\hat L$. For each $g$ choose a
level lift $k_g$ as in the hypothesis. The ideals satisfy $M_b\bar M_b=M\mathcal O_{E_h}$ and
$M_b+\bar M_b=\mathcal O_{E_h}$.

(a) \emph{A level lift $k$ preserves $\hat L$.} The matrix $y=\beta_i^{-1}k\beta_i$ lies in
$GL_{r}(\hat{\mathcal O})$, so by (1) $k\hat L_0\subset\hat L_0$, and $k$ commutes with the scalars
$a\in\mathfrak A$. The same holds for every matrix $\beta_iT\beta_i^{-1}$ with $T$ integral.

(b) \emph{If $k$ is a level lift of $1$, then $k\iota(a)-\iota(a)\in M\hat L$ for $a\in\bar M_b\Lambda$.}
The reduction of $y=\beta_i^{-1}k\beta_i$ modulo every $P\in P_g$ is $1$, so by (3) the entries of
$y-1$ lie in $M_b\hat{\mathcal O}$. For $b\in\bar M_b$ the entries of $b(y-1)$ lie in
$M_b\bar M_b\hat{\mathcal O}=M\hat{\mathcal O}$, so $b(y-1)=MT$ with $T$ integral, and $b(k-1)=Mm$ with
$m=\beta_iT\beta_i^{-1}$, which preserves $\hat L$ by (a). For $a\in\Lambda$, $\iota(ba)=b\iota(a)$ with
$\iota(a)\in\hat L$, so $k\iota(ba)-\iota(ba)=Mm\iota(a)\in M\hat L$. The elements $ba$ generate
$\bar M_b\Lambda$, and the claim follows by additivity.

(c) \emph{If $x\in\Lambda$ and $\iota(x)=Mw$ with $w\in\hat L$, then $x=Mx'$ with $x'\in\Lambda$.} The
vector $M^{-1}\iota(x)\in E_h^{r}$ has adelic image $w\in\hat L$, so by (2) it is $\iota(x')$ for some
$x'\in\Lambda$; as $\iota$ is injective, $x=Mx'$.

(d) \emph{If $f$ is an endomorphism of $\Lambda$ with $\iota(f(a))\in M\hat L$ for every
$a\in\bar M_b\Lambda$, then $f(\Lambda)\subset M_b\Lambda$.} Write $1=u+v$ with $u\in M_b$ and
$v\in\bar M_b$. For $a\in\Lambda$, $va\in\bar M_b\Lambda$, so $f(va)=Mx$ with $x\in\Lambda$ by (c), and
$f(a)=uf(a)+Mx\in M_b\Lambda$, as $M\in M_b$.

(e) Let $a\in\bar M_b\Lambda$. For $f=\lambda_1-1$,
$\iota(\lambda_1a)-\iota(a)=\big(\iota(\lambda_1a)-k_1\iota(a)\big)+\big(k_1\iota(a)-\iota(a)\big)\in M\hat L$
by the hypothesis and (b). For $f=\lambda_{gh}-\lambda_g\lambda_h$: the reductions of level lifts
multiply, so $k_gk_h$ is a level lift of $gh$ and $u=(k_gk_h)^{-1}k_{gh}$ is a level lift of $1$; and
$\lambda_h(a)\in\bar M_b\Lambda$. Then
\[
\begin{aligned}
 \iota(\lambda_{gh}a)-\iota(\lambda_g\lambda_ha)={}&\big(\iota(\lambda_{gh}a)-k_{gh}\iota(a)\big)
 +k_gk_h\big(u\iota(a)-\iota(a)\big)\\
 &-k_g\big(\iota(\lambda_ha)-k_h\iota(a)\big)-\big(\iota(\lambda_g\lambda_ha)-k_g\iota(\lambda_ha)\big),
\end{aligned}
\]
and each term lies in $M\hat L$, by the hypothesis, (b) and (a). By (d), both $f$ map $\Lambda$ into
$M_b\Lambda$.
% Lean: congr_of_levIs = congr_of_lifts (L = R.datWp i R.triv, l.prod for M_b Mbar_b = M, l.coprime,
% 3 <= M), which applies congr_of_matrices (the step (d) is its key, with hcop and exists_eq_nsmul for
% (c)) to the matrices finHom (kappa g); (a) is mulVec_mem_of_St (parLev_St, mem_latHat_iff from
% IsSDFrame, mulVec_mem_of_frame), (b) is lift_one_congr (mem_idS_prod, mulVec_smul_congr), and the
% product is Lev.lift_div_mem with rTup_mul and mulVec_mul_congr.
\end{proof}

\begin{situation}\label{is-onto-step-move-laws-situation}
Let $E$, the prime $d$, the elements $\sqrt{-d}$ and $w_0$, the square root $\theta$, the lattice $\Lambda$
of a PEL datum of signature $(p,q)$ over $E$ with its form $\langle\,,\rangle$, and the level
$l=(M,M_b,A)$ of this datum be as in Lemma~\ref{is-onto-step-leaves-move-isometry}. Let $X$ be the
moduli scheme at level $l$, over the field $K_l$, with its universal object $(\mathcal A,\eta,\eta')$
(Definition~\ref{cv-toroidal-moduli-object}(4) and~(5)), and let $F\subset\mathbf C$ be a number field
that contains $\tau(E)$ for every embedding $\tau$ of $E$, the $\operatorname{lcm}(M,A)$-th roots of unity
and $K_l$. Framed points of $X$ over $F$-algebras, and the image of a framed point under a homomorphism of
$F$-algebras, are as in Definitions~\ref{cv-volume-twisting-section}(3)
and~\ref{cv-uniformization-theta-rule-values}. For a framed point $y$ over an $F$-algebra $S$ and a family
$t$ of units of $S$, one for each vector of the bases of $y$ and one more, $y\cdot t$ is a framed point
over $S$ with the point of $y$, whose basis vector of index $a$ in each eigenspace is a combination of the
vectors of $y$ in that eigenspace of index at most $a$ with the entry of $t$ at that vector as the
coefficient of index $a$, and whose generator is the last entry of $t$ times that of $y$. For
$\mathcal O_E$-linear endomorphisms $g$ and $g'$ of $\Lambda$, a \emph{level automorphism} of $X$ for
$(g,g')$ is a pair $(\rho,\nu)$ of an automorphism $\rho$ of $X$ over $K_l$ and a morphism
$\nu\colon\mathcal A\to\mathcal A$ of $\mathcal O_E$-semi-abelian schemes over $\rho$
(Definition~\ref{cv-model-toroidal-semiabelian}(2)) with
\[
 \nu\circ\eta(g(a))=\eta(a)\circ\rho\quad\text{and}\quad\nu\circ\eta'(g'(a))=\eta'(a)\circ\rho
 \qquad(a\in\Lambda).
\]
For $g'=1$ this is a level automorphism with matrix $g$ (Definition~\ref{cv-realization-level-move}(1))
that fixes $\eta'$. The \emph{move} of a framed point by $(\rho,\nu)$ is that of
Definition~\ref{cv-realization-level-move}(2): it has the point $\rho\circ x$, where $x$ is the point of
the framed point, and the images under $\nu$ of its bases, and its generator is carried to the new point
by an isomorphism of the line bundle of $\mathcal A$ with its pullback by $\nu$ that respects the
rigidifications and the cubical structures.
\end{situation}

\begin{lemma}\label{is-onto-step-move-laws-inverse}
Let $\Lambda$, with its form $\langle\,,\rangle$, be as in Lemma~\ref{is-onto-step-leaves-move-isometry},
let $a_0,b_0\in\Lambda$ with $\langle a_0,b_0\rangle=1$, and let $(\omega_m,\nu_m)$ be a family of
similitudes of $\Lambda$ (Definition~\ref{cv-toroidal-moduli-object}(1)). Then there is a family of
similitudes $(\omega'_m,\nu'_m)$ of $\Lambda$ such that for every $m\geq1$ and $a\in\Lambda$ the elements
$\omega'_m(\omega_m(a))-a$ and $\omega_m(\omega'_m(a))-a$ lie in $m\Lambda$, and
$\nu_m\nu'_m\equiv1\pmod m$.
\end{lemma}
\leanref{Development/AbelianVarietyTheory/Proofs/ModuliProofs/ExistsFormEqOne.lean}{GlobalTheory.SimAlg.exists_inv}
\begin{proof}
Write $a\equiv b\pmod m$ for $a-b\in m\Lambda$. For each $m$, a family of similitudes gives an
$\mathcal O_E$-linear $h_m$ with $h_m(\omega_m(a))\equiv a$ and $\omega_m(h_m(a))\equiv a\pmod m$ for
every $a$; put $\omega'_m=h_m$ and $\nu'_m=\langle h_m(a_0),h_m(b_0)\rangle$. Then
$\nu_m\nu'_m\equiv\langle\omega_m h_m(a_0),\omega_m h_m(b_0)\rangle\equiv\langle a_0,b_0\rangle=1\pmod m$.
For $m,m'\geq1$ and $a\in\Lambda$,
\[
 h_m(a)\equiv h_m\big(\omega_{mm'}(h_{mm'}(a))\big)\equiv h_m\big(\omega_m(h_{mm'}(a))\big)\equiv
 h_{mm'}(a)\pmod m,
\]
using $\omega_{mm'}h_{mm'}\equiv1\pmod{mm'}$, $\omega_{mm'}\equiv\omega_m\pmod m$ and $h_m\omega_m\equiv1
\pmod m$. The map $\omega_m$ inverts $h_m$ modulo $m$. Finally, for $x,y\in\Lambda$, modulo $m$,
\[
 \langle h_mx,h_my\rangle\equiv\nu_m\nu'_m\langle h_mx,h_my\rangle\equiv\nu'_m\langle\omega_mh_mx,\omega_mh_my\rangle
 \equiv\nu'_m\langle x,y\rangle.
\]
So $(h_m,\nu'_m)$ is a family of similitudes with the stated properties.
% Lean: SimAlg.exists_inv: h n from g.inv n; the tower from h1, h2, h3 (modN_of_mul), nu' n =
% EF (h n a) (h n b) with dvd_mult_inv, inv from (hh n hn x).2 and .1, and mult from h1, h2, h3.
% exists_form_eq_one gives a_0, b_0 when p + q >= 1. About 60 lines.
\end{proof}

\begin{lemma}[Alp\"oge]\label{is-onto-step-move-laws-automorphism}
In Situation~\ref{is-onto-step-move-laws-situation}, suppose that $p+q\geq1$, and let
$(\omega_m,\nu_m)$ be a family of similitudes of $\Lambda$ with $\nu_A\equiv1\pmod A$. Then $X$ has a level
automorphism for $(\omega_M,\omega_A)$.
\end{lemma}
\leanref{Development/AbelianVarietyTheory/Proofs/UnitaryExistence/FrMoveTransfer.lean}{GlobalTheory.LevelActionLaws.exists_levAut}
\leanref{Development/AbelianVarietyTheory/Proofs/ModuliProofs/ExistsObjTwistBySimilitude.lean}{GlobalTheory.exists_obj_twist_by_similitude}
\leanref{Development/AbelianVarietyTheory/Proofs/ModuliProofs/CompEqIdOfTwists.lean}{GlobalTheory.Moduli.LevelAction.exists_twists}
\leanref{Development/AbelianVarietyTheory/Proofs/ModuliProofs/CompEqIdOfTwists.lean}{GlobalTheory.Moduli.LevelAction.comp_eq_id_of_twists}
\begin{proof}
\emph{Step 1: the twist.} Let $(G,\eta_G,\eta'_G)$ be an object at level $l$ over a scheme $S$ over $K_l$
(Definition~\ref{cv-toroidal-moduli-object}(4)). Then $(G,\eta_G\circ\omega_M,\eta'_G\circ\omega_A)$ is an
object at level $l$ over $S$. Indeed, both maps are additive; $\eta_G\circ\omega_M$ vanishes on
$M_b\Lambda$, as $\omega_M$ is $\mathcal O_E$-linear, and $\eta'_G\circ\omega_A$ on $A\Lambda$. At a point
$x$ of $S$ with values in an algebraically closed field, let $((\zeta'_m),(\iota_m))$ be the Tate system of
$G$ at $x$ given by the object. By Lemma~\ref{cv-toroidal-moduli-similitude}(1),
$(((\zeta'_m)^{\nu_m}),(\iota_m\circ\omega_m))$ is a Tate system of $G$ at $x$. Its root of unity of
index $A$ is $(\zeta'_A)^{\nu_A}=\zeta'_A$, as $\nu_A\equiv1\pmod A$ and $\zeta'_A$ has order $A$; and
$\iota_A(\omega_A(a))=\eta'_G(\omega_A(a))(x)$ for every $a$, and $\iota_M(\omega_M(a))=\eta_G(\omega_M(a))(x)$
for $a\in\bar M_b\Lambda$, as $\omega_M(a)\in\bar M_b\Lambda$. The conditions on $G$ alone, among them
condition (a) of Definition~\ref{cv-toroidal-family}(1), are unchanged. So the twisted triple satisfies
the conditions at $x$.

\emph{Step 2: the classifying map.} The scheme $X$ is smooth over the field $K_l$, so it is reduced. By
Step 1, $(\mathcal A,\eta\circ\omega_M,\eta'\circ\omega_A)$ is an object at level $l$ over $X$, and by the
universal property of $X$ (Definition~\ref{cv-toroidal-moduli-object}(5)) there are a morphism
$f\colon X\to X$ over $K_l$ and a morphism $\nu$ of objects from it to $(\mathcal A,\eta,\eta')$ over $f$.
That $\nu$ is a morphism of objects says $\nu\circ\eta(\omega_M(a))=\eta(a)\circ f$ and
$\nu\circ\eta'(\omega_A(a))=\eta'(a)\circ f$ for every $a$.

\emph{Step 3: $f$ is an automorphism.} As $p+q\geq1$ and the form is perfect, there are $a_0,b_0$ with
$\langle a_0,b_0\rangle=1$. Let $(\omega'_m,\nu'_m)$ be the family of
Lemma~\ref{is-onto-step-move-laws-inverse}. Then $\nu'_A\equiv\nu_A\nu'_A\equiv1\pmod A$, so Step 2 gives
$f'$ and $\nu'$ for it. The composite $\nu'\circ\nu$ is a morphism of $\mathcal O_E$-semi-abelian schemes
over $f'\circ f$, and for $a\in\Lambda$,
\[
 \nu'\circ\nu\circ\eta(a)=\nu'\circ\nu\circ\eta\big(\omega_M(\omega'_M(a))\big)
 =\nu'\circ\eta(\omega'_M(a))\circ f=\eta(a)\circ f'\circ f,
\]
since $\omega_M\omega'_M(a)-a\in M\Lambda\subset M_b\Lambda$, which $\eta$ kills; in the same way
$\nu'\circ\nu\circ\eta'(a)=\eta'(a)\circ f'\circ f$, as $\eta'$ kills $A\Lambda$. So $\nu'\circ\nu$ is a
morphism of objects from $(\mathcal A,\eta,\eta')$ to itself over $f'\circ f$, and so is the identity over
the identity of $X$. By the uniqueness in the universal property, $f'\circ f$ is the identity. Exchanging
the two families, $f\circ f'$ is the identity. So $(f,\nu)$ is a level automorphism for
$(\omega_M,\omega_A)$.
% Lean: exists_levAut: hsim = exists_obj_twist_by_similitude SimAlg.sim_group (LevelTwist.objTwist,
% whose geom field is LevelTwist.twist with zeta_zpow_congr; Step 1), hred from
% isReduced_isNormalScheme_isLocallyNoetherian_of_smooth_over_field, i = the identity of the universal
% object (ObjMaps.id_comp_cancel), g' from the fourth clause of SimAlg.sim_group (SimAlg.exists_inv,
% Lemma is-onto-step-move-laws-inverse), Moduli.LevelAction.exists_twists (Step 2: Xo.rep, with the
% inverse k of the identity map m and exists_comp_sabMap) and comp_eq_id_of_twists (Step 3:
% exists_comp_sabMap and the uniqueness of Xo.rep). The composites and inverses of morphisms of
% O_E-semi-abelian schemes (PullCub, SabComp, about 500 lines) are taken as routine. Size of
% exists_levAut, following imports and stopping at statements the book cites: 1,227 lines.
\end{proof}

\begin{lemma}[Alp\"oge]\label{is-onto-step-move-laws-unique}
In Situation~\ref{is-onto-step-move-laws-situation}, suppose that $p+q\geq1$. Let $g_M,g'_M,g_A,g'_A$ be
$\mathcal O_E$-linear endomorphisms of $\Lambda$ with $g_M(a)-g'_M(a)\in M_b\Lambda$ and
$g_A(a)-g'_A(a)\in A\Lambda$ for every $a$, let $(\rho,\nu)$ be a level automorphism of $X$ for
$(g_M,g_A)$ and $(\rho',\nu')$ one for $(g'_M,g'_A)$, and suppose that there is a family of similitudes
$(\omega_m,\nu_m)$ of $\Lambda$ with $\omega_M(a)-g_M(a)\in M_b\Lambda$ and $\omega_A(a)-g_A(a)\in
A\Lambda$ for every $a$, and $\nu_A\equiv1\pmod A$. Then $\rho=\rho'$, and $\nu=\nu'$ as morphisms of
schemes.
\end{lemma}
\leanref{Development/AbelianVarietyTheory/Proofs/UnitaryExistence/LevelMoveModuliUniqueUniversalTripleStatement.lean}{GlobalTheory.T.levelMove_moduli_unique_universalTriple}
\leanref{Development/AbelianVarietyTheory/Proofs/UnitaryExistence/LevelMoveModuliUniqueUniversalTriple.lean}{GlobalTheory.levMove_unique}
\begin{proof}
By Step 1 of the proof of Lemma~\ref{is-onto-step-move-laws-automorphism},
$o=(\mathcal A,\eta\circ\omega_M,\eta'\circ\omega_A)$ is an object at level $l$ over $X$. For $a\in\Lambda$,
as $\eta$ kills $M_b\Lambda$ and $\eta'$ kills $A\Lambda$,
\[
 \nu\circ\eta(\omega_M(a))=\nu\circ\eta(g_M(a))=\eta(a)\circ\rho,\qquad
 \nu'\circ\eta(\omega_M(a))=\nu'\circ\eta(g'_M(a))=\eta(a)\circ\rho',
\]
and in the same way $\nu\circ\eta'(\omega_A(a))=\eta'(a)\circ\rho$ and
$\nu'\circ\eta'(\omega_A(a))=\eta'(a)\circ\rho'$. So $\nu$ and $\nu'$ are morphisms of objects from $o$ to
$(\mathcal A,\eta,\eta')$ over $\rho$ and over $\rho'$. As $X$ is smooth over $K_l$, it is reduced, and the
uniqueness in the universal property of $X$ (Definition~\ref{cv-toroidal-moduli-object}(5)) gives
$\rho=\rho'$. Then $\nu$ and $\nu'$ are two morphisms of $\mathcal O_E$-semi-abelian schemes from $\mathcal A$ to
$\mathcal A$ over the same $\rho$ that agree on the sections $\eta(\omega_M(a))$ of the level structure of
$o$. By the rigidity of morphisms of objects at a level, they are equal.
\openstep{The rigidity of morphisms of objects at a level, that two morphisms of
$\mathcal O_E$-semi-abelian schemes over the same map of a reduced or locally noetherian base that agree
on the level structure are equal, is not written out.}
% Lean: levMove_unique: exists_obj_twist_by_similitude gives o with the identity map m; n and n' are
% m composed with nu and nu' (exists_comp_sabMap); c1, c2, d1, d2 with eta_congr and etaA_congr;
% rho = rho' from (Xo.rep _ Xo.s hred o).unique; then obj_rigid_closed firstFiveLaws.1 (the open step)
% and cancel_epi m.phi, m.phi an isomorphism. Sizes, following imports and stopping at statements the
% book cites: levMove_unique 2,928 lines, of which 683 below levMove_unique without
% obj_rigid_closed; obj_rigid_closed about 2,140 lines, 1,733 of them in theorems.
\end{proof}

\begin{lemma}[Alp\"oge]\label{is-onto-step-move-laws-moves}
In Situation~\ref{is-onto-step-move-laws-situation}, let $(\rho,\nu)$ and $(\rho',\nu')$ be pairs of a
morphism $\rho$ of $X$ to itself over $K_l$ and a morphism $\nu$ of $\mathcal O_E$-semi-abelian schemes
from $\mathcal A$ to itself over $\rho$, and let $S$ be an $F$-algebra.
\begin{enumerate}
\item Every framed point $y$ of $X$ over $S$ has a move by $(\rho,\nu)$; and if $\nu''$ is a morphism of
  $\mathcal O_E$-semi-abelian schemes over $\rho$ equal to $\nu$ as a morphism of schemes, $y'$ is a move of
  $y$ by $(\rho,\nu)$ and $y''$ one by $(\rho,\nu'')$, then $y'=y''$.
\item There is a morphism $\iota$ of $\mathcal O_E$-semi-abelian schemes over the identity of $X$ that is
  the identity of $\mathcal A$, and every framed point is its own move by $(1,\iota)$.
\item There is a morphism $\mu$ of $\mathcal O_E$-semi-abelian schemes over $\rho'\circ\rho$ that is
  $\nu'\circ\nu$ as a morphism of schemes, such that if $y'$ is a move of $y$ by $(\rho,\nu)$ and $y''$ a move
  of $y'$ by $(\rho',\nu')$, then $y''$ is a move of $y$ by $(\rho'\circ\rho,\mu)$.
\item Let $\varphi\colon S\to T$ be a homomorphism of $F$-algebras, $y$ and $y'$ framed points over $S$, and
  $z$ and $z'$ over $T$. If $z$ is the image of $y$ under $\varphi$, $y'$ is a move of $y$ and $z'$ a move of
  $z$ by $(\rho,\nu)$, then $z'$ is the image of $y'$ under $\varphi$.
\item Let $t$ be a family of units of $S$ as in Situation~\ref{is-onto-step-move-laws-situation}, $y,y',z,z'$
  framed points over $S$ with $z=y\cdot t$, and $y'$ and $z'$ moves of $y$ and $z$ by $(\rho,\nu)$. Then
  $z'=y'\cdot t$.
\end{enumerate}
\end{lemma}
\leanref{Development/AbelianVarietyTheory/Proofs/UnitaryExistence/ExistsUniqueFramedPointMoveUniversalTripleStatement.lean}{GlobalTheory.T.existsUnique_framedPoint_move_universalTriple}
\leanref{Development/AbelianVarietyTheory/Proofs/UnitaryExistence/FramedPointMoveCalculusStatement.lean}{GlobalTheory.T.framedPointMove_calculus}
\begin{proof}
(1) Let $x$ be the point of $y$. The move has the point $\rho\circ x$, the images under $\nu$ of the bases of
$y$, which are again bases of the eigenspaces since $\nu$ identifies $\mathcal A$ with the pullback of
itself by $\rho$ compatibly with $\mathcal O_E$, and as generator the section carried from that of $y$ by the
isomorphism of the line bundle with its pullback by $\nu$ that comes with $\nu$. Two moves through the same
morphism of schemes differ by an automorphism of the rigidified cubical line bundle on the abelian scheme
over $S$, which is $1$, as the global functions on an abelian scheme are constant along the fibres.
\openstep{The construction of the move of a framed point by $(\rho,\nu)$, and its uniqueness through the
normalized isomorphisms of the line bundle, are not written out.}
% Lean: T.existsUnique_framedPoint_move_universalTriple = GlobalTheory.exists_unique_frMove: existence
% with q (pullback.lift), LevelMoveTest.pinned_isIso, isGen_iff, exists_carried_forth for the
% generator, isTanAt_push, isEigAt_iff, exists_isTanAt, isCombAt_iff and exists_carried_back; uniqueness
% with SemiAbelianTripleChecks.sabMap_ext, LevelMoveTest.frMove_unique and
% iso_eq_of_maps_along_section with firstFiveLaws.1.globalFns. Size: about 613 lines below it, 461 in
% theorems.

(2) to (5)
\openstep{The identity, composition, base change and triangular change laws of the moves of framed
points are not written out.}
% Lean: T.framedPointMove_calculus = MoveCalc.frMove_calculus = move_id, move_comp, move_frMap and
% move_frBor. Sizes, following imports and stopping at statements the book cites: frMove_calculus about
% 1,919 lines (move_id 444, move_comp 463, move_frMap 1,382, move_frBor 1,458, overlapping); with
% exists_unique_frMove, 2,124 lines.
\end{proof}

\begin{lemma}[Alp\"oge]\label{is-onto-step-move-laws-framed}
In Situation~\ref{is-onto-step-move-laws-situation}, suppose that $p+q\geq1$ and that $A$ is prime to $M$.
For every $\mathcal O_E$-linear automorphism $\gamma$ of $\Lambda/M_b\Lambda$ fix an $\mathcal O_E$-linear
endomorphism $\tilde\gamma$ of $\Lambda$ that induces it. For framed points $y$ and $y'$ of $X$ over an
$F$-algebra $S$, call $y'$ a \emph{move of $y$ by $\gamma$} if it is the move of $y$ by a level automorphism
of $X$ for $(\tilde\gamma,1)$. Then for all automorphisms $\gamma,\delta$ of $\Lambda/M_b\Lambda$ and all
$F$-algebras $S$ and $T$:
\begin{enumerate}
\item every framed point $y$ over $S$ has a move by $\gamma$;
\item if $y'$ and $y''$ are moves of $y$ by $\gamma$, then $y''$ is the image of $y'$ under the identity of
  $S$;
\item $y$ is a move of itself by the identity;
\item if $y'$ is a move of $y$ by $\gamma$ and $y''$ a move of $y'$ by $\delta$, then $y''$ is a move of $y$
  by $\gamma\delta$;
\item if $\varphi\colon S\to T$ is a homomorphism of $F$-algebras, $z$ the image of $y$ under $\varphi$, and
  $y'$ and $z'$ moves of $y$ and $z$ by $\gamma$, then $z'$ is the image of $y'$ under $\varphi$;
\item if $z=y\cdot t$ for a family $t$ of units of $S$ as in
  Situation~\ref{is-onto-step-move-laws-situation}, and $y'$ and $z'$ are moves of $y$ and $z$ by $\gamma$,
  then $z'=y'\cdot t$.
\end{enumerate}
\end{lemma}
\leanref{Development/AbelianVarietyTheory/Proofs/JetHarmonics/ActLevFullLevelGroupFramedPoints.lean}{GlobalTheory.FullLevelAction.actLaws_levFull}
\leanref{Development/AbelianVarietyTheory/Proofs/JetHarmonics/ActLevFullLevelGroupFramedPointsStatement.lean}{GlobalTheory.T.actLev_full_level_group_framedPoints}
\leanref{Development/AbelianVarietyTheory/Proofs/UnitaryExistence/FrMoveTransfer.lean}{GlobalTheory.LevelActionLaws.exists_move}
\leanref{Development/AbelianVarietyTheory/Proofs/UnitaryExistence/FrMoveTransfer.lean}{GlobalTheory.LevelActionLaws.frMove_transfer}
\begin{proof}
\emph{Step 1: level automorphisms exist.} Let $\gamma$ be given. If $\tilde\gamma(a)\in M_b\Lambda$, then
$\gamma$ kills the class of $a$, so $a\in M_b\Lambda$. As $A$ is prime to $M$,
Lemma~\ref{cv-toroidal-moduli-similitude}(2), with $\mathfrak M=M_b$, the integer $A$ and
$f=\tilde\gamma$, gives a family of similitudes $(\omega_m,\nu_m)$ with
$\omega_M(a)-\tilde\gamma(a)\in M_b\Lambda$, $\omega_A(a)-a\in A\Lambda$ and $\nu_A\equiv1\pmod A$. By
Lemma~\ref{is-onto-step-move-laws-automorphism}, $X$ has a level automorphism $(\rho,\nu)$ for
$(\omega_M,\omega_A)$. As $\eta$ kills $M_b\Lambda$ and $\eta'$ kills $A\Lambda$,
$\nu\circ\eta(\tilde\gamma(a))=\nu\circ\eta(\omega_M(a))=\eta(a)\circ\rho$ and
$\nu\circ\eta'(a)=\nu\circ\eta'(\omega_A(a))=\eta'(a)\circ\rho$: $(\rho,\nu)$ is a level automorphism for
$(\tilde\gamma,1)$.

\emph{Step 2: the move does not depend on the level automorphism.} Let $(\rho,\nu)$ and $(\rho',\nu')$ be
level automorphisms for $(\tilde\gamma,1)$. The family of Step 1 satisfies the hypothesis of
Lemma~\ref{is-onto-step-move-laws-unique} with $g_M=g'_M=\tilde\gamma$ and $g_A=g'_A=1$, so $\rho=\rho'$ and
$\nu=\nu'$ as morphisms of schemes. If $y'$ is the move of $y$ by $(\rho',\nu')$, then by
Lemma~\ref{is-onto-step-move-laws-moves}(1) there is a move $y''$ of $y$ by $(\rho,\nu)$, and $y''=y'$; so
$y'$ is also the move of $y$ by $(\rho,\nu)$.

(1) By Step 1 and Lemma~\ref{is-onto-step-move-laws-moves}(1).

(2) Let $y'$ be the move of $y$ by a level automorphism $(\rho,\nu)$ for $(\tilde\gamma,1)$. By Step 2,
$y''$ is also a move of $y$ by $(\rho,\nu)$, and by Lemma~\ref{is-onto-step-move-laws-moves}(1),
$y''=y'$. Every framed point is its own image under the identity of $S$.

(3) Let $\iota$ be the morphism of Lemma~\ref{is-onto-step-move-laws-moves}(2). As $\tilde1(a)-a\in
M_b\Lambda$, which $\eta$ kills, $\iota\circ\eta(\tilde1(a))=\eta(a)$ and $\iota\circ\eta'(a)=\eta'(a)$, so
$(1,\iota)$ is a level automorphism for $(\tilde1,1)$, and $y$ is its own move by it.

(4) Let $y'$ be the move of $y$ by a level automorphism $(\rho,\nu)$ for $(\tilde\gamma,1)$, and $y''$ that of
$y'$ by one $(\rho',\nu')$ for $(\tilde\delta,1)$. Let $\mu$ be the morphism over $\rho'\circ\rho$ of
Lemma~\ref{is-onto-step-move-laws-moves}(3), so that $y''$ is the move of $y$ by $(\rho'\circ\rho,\mu)$. The
endomorphism $\widetilde{\gamma\delta}-\tilde\gamma\tilde\delta$ maps $\Lambda$ into $M_b\Lambda$, as both
induce $\gamma\delta$. So for $a\in\Lambda$,
\[
 \mu\circ\eta(\widetilde{\gamma\delta}(a))=\nu'\circ\nu\circ\eta(\tilde\gamma(\tilde\delta(a)))
 =\nu'\circ\eta(\tilde\delta(a))\circ\rho=\eta(a)\circ\rho'\circ\rho,
\]
and $\mu\circ\eta'(a)=\nu'\circ\nu\circ\eta'(a)=\eta'(a)\circ\rho'\circ\rho$. So $(\rho'\circ\rho,\mu)$ is a
level automorphism for $(\widetilde{\gamma\delta},1)$, and $y''$ is a move of $y$ by $\gamma\delta$.

(5) and (6) Let $y'$ be the move of $y$ by a level automorphism $(\rho,\nu)$ for $(\tilde\gamma,1)$. By
Step 2, $z'$ is the move of $z$ by $(\rho,\nu)$. Lemma~\ref{is-onto-step-move-laws-moves}(4) and~(5) give the
claims.
% Lean: actLaws_levFull (and actLev_levFull with A = 1): the six fields total, func, one, mul, nat and bor
% of ActLaws for levActR (levFull Dw l) = levR (liftLa gamma, id). Step 1 is exists_move (exists_sim =
% exists_sim_of_aut_mod, Lemma cv-toroidal-moduli-similitude(2), with the hypothesis from liftLa_spec;
% exists_levAut; eta_congr and etaA_congr). Step 2 is frMove_transfer (move_unique =
% T.levelMove_moduli_unique_universalTriple with exists_sim, and the uniqueness of
% T.existsUnique_framedPoint_move_universalTriple). (2) ends with Frames.frMap_unique_refl_iff (the
% identity image); (3) uses liftLa_one_sub and the identity move of T.framedPointMove_calculus; (4)
% liftLa_mul_sub and its composite; (5) and (6) its base change and triangular laws. The open parts are
% in Lemmas is-onto-step-move-laws-unique and is-onto-step-move-laws-moves.
\end{proof}

\begin{lemma}\label{is-onto-step-level-structure-cusp}
Let $s$ be a run at a point of a family that satisfies the laws (Definition~\ref{is-onto-step-runs}),
$l$ a level of a lift and $w$ a weight, and let $\gamma$ be an $\mathcal O_{E_h}$-linear automorphism of
the level module $\Lambda_l/M_b\Lambda_l$ of the core. Then the level action of $\gamma$, the change
of the level structures of the framed points through an endomorphism of $\Lambda_l$ that induces
$\gamma$, carries the cusp forms of $s$ at level $l$ and weight $w$ to cusp forms of $s$ at level $l$ and
weight $w$.
\end{lemma}
\leanref{Development/AbelianVarietyTheory/Proofs/JetHarmonics/LevelActionMemCuspFormsStatement.lean}{GlobalTheory.T.levelAction_mem_cuspForms}
\leanref{Development/AbelianVarietyTheory/Proofs/JetHarmonics/ExistsLevelChangeActionOnCuspForms.lean}{GlobalTheory.FormTies.exists_isLevAct}
\begin{proof}
The moves of framed points of the moduli scheme of the core at level $l$ by the automorphisms of the level
module satisfy the laws of Lemma~\ref{is-onto-step-move-laws-framed}: the datum of the core at level $l$ is
a PEL datum over $E_h$ with $p+q\geq1$, and the full level of the core is $1$, by the pins of the laws.
\openstep{The proof that the level action of an automorphism of the level module carries sections of a
weight to sections, and preserves their cuspidality at the good embeddings, is not written out.}
% Lean: T.levelAction_mem_cuspForms = LevelCusp.actLev_mem_cuspA: T.actLev_full_level_group_framedPoints
% (FullLevelAction.actLev_levFull, Lemma is-onto-step-move-laws-framed, with hpq from the pins
% (sigW) and A = 1 from hp.2.1.2) for the laws, actL_mem for sections, and LevelCusp.rdU_mem_SS with
% exists_match and laU_match for the cusp condition at each good embedding. The lift of gamma is liftLa
% gamma (levFull). Size: about 32,600 lines of Lean in 1,900 theorems below actLev_mem_cuspA (text
% count, stopping at statements the book cites), most of them shared with other open steps; the largest
% of the four open steps of this lemma group. GeoLevel.actLev_mem_cuspA, for the action of G through the
% one-map data (Lemma is-eigenvalue-conjugate-frame-marking-laws(2), "g f is a core form"), has the
% same proof with ActLev in place of T.actLev_full_level_group_framedPoints.
\end{proof}

\begin{lemma}\label{is-onto-step-level-structure}
Let $p$ be a point of a family that satisfies the point conditions, whose realization and face have
the links of Lemma~\ref{is-onto-step-links}, and whose core ties hold. Then there are maps $\ell_i$, for
the lifts $i$, from the group $G$ of the realization to the automorphisms of the level module of the
core at level $i+1$, such that every run at $p$ that satisfies the laws satisfies the level and
action conditions (c) and (d) of Definition~\ref{is-onto-step-runs}(5) with $\ell_i$ at every
lift $i$.
\end{lemma}
\leanref{Development/AbelianVarietyTheory/Proofs/DivisorProofs/FiniteGroupLevelStructureActionFixedRunsStatement.lean}{GlobalTheory.T.finiteGroup_levelStructure_action_fixedRuns}
\begin{proof}
Fix a lift $i$, and use the notation of Definition~\ref{is-onto-step-model-level-condition} at $i$:
$\Lambda=\Lambda_{i+1}$, $\iota=\iota_{i+1}$, the level $M$, $M_b$ of the core and $\hat L_W$. Let
$\lambda_g$, for $g\in G$, be the level maps of the one-map data at $i$, endomorphisms of $\Lambda$. For
$g\in G$ let $\ell_i(g)$ be the endomorphism of $\Lambda/M_b\Lambda$ induced by $\lambda_g$ if it is
bijective, and the identity otherwise.

The one-map data are pinned at $i$ (Lemma~\ref{is-onto-step-links}(4)). Among their pins is that the
$\lambda_g$ give the level of $R$ at $i$ in the form of Definition~\ref{is-onto-step-model-level-condition}
with $\lambda_g$ in place of $s_g$: every $g\in G$ has a level lift $k$ with
$\iota(\lambda_g(a))-k\,\iota(a)\in M\hat L_W$ for every $a\in\bar M_b\Lambda$.

Let $s$ be a run at $p$ that satisfies the laws. The laws of $s$ give that $\iota$ is injective with
image the set of vectors whose adelic image lies in $\hat L_W$; the core conditions of $R$, part of its
characterization (Lemma~\ref{is-onto-step-links}(1)), give that the frame $\beta_i$ describes $\hat L_0$
as in (1) of Lemma~\ref{is-onto-step-level-structure-congruence}; and the core ties give
$M_b=\prod_{P\in P_g}P$. So that lemma applies: $\lambda_1\equiv1$ and
$\lambda_{gh}\equiv\lambda_g\lambda_h$ modulo $M_b\Lambda$. Hence $\lambda_g\lambda_{g^{-1}}$ and
$\lambda_{g^{-1}}\lambda_g$ induce the identity on $\Lambda/M_b\Lambda$, the map induced by $\lambda_g$ is
bijective, and $\ell_i(g)$ is this map. Let $s_g$ be the lift of $\ell_i(g)$ to $\Lambda$ of
Definition~\ref{is-onto-step-model-level-condition}. Then $s_g(a)-\lambda_g(a)\in M_b\Lambda$ for every
$a\in\Lambda$.

\emph{The level condition (c).} Let $g\in G$, with the level lift $k$ of the pins. As $s_g-\lambda_g$ maps
$\Lambda$ into $M_b\Lambda$, it maps $\bar M_b\Lambda$ into $M_b\bar M_b\Lambda=M\Lambda$. So for
$a\in\bar M_b\Lambda$, $s_g(a)=\lambda_g(a)+My$ with $y\in\Lambda$, and $\iota(y)\in\hat L_W$. Hence
\[
 \iota(s_g(a))-k\,\iota(a)=\big(\iota(\lambda_g(a))-k\,\iota(a)\big)+M\iota(y)\in M\hat L_W,
\]
and $\ell_i$ gives the level of $R$ at $i$.

\emph{The action condition (d).} The point conditions include that the one-map data at $i$ give the
$i$-th lift as one map for some constants and phase. Part of this is that, for $g\in G$ with a level
lift, the level action of $G$ through the maps $\lambda_g$ carries the cusp forms of the core at level
$i+1$ and weight $\kappa_i$ to cusp forms; every $g$ has a level lift by the pins. Another point
condition, on the action of $G$, then says that the map $e_i$ of the weight tie carries this level
action to the action $\rho$ of $G$ on $FW_i$: $e_i(g\cdot F)=\rho(g)\,e_i(F)$ for every such cusp form
$F$. The cusp forms of $s$ at level $i+1$ are those of the core. By
Lemma~\ref{is-onto-step-level-structure-cusp}, the level action of $\ell_i(g)$, through $s_g$, carries
the cusp forms of $s$ at level $i+1$ and weight $\kappa_i$ to cusp forms, and so gives maps of these cusp
forms that are the level action of $\ell_i$. The level action through an endomorphism $f$ of $\Lambda$
depends on $f$ only modulo $M_b\Lambda$: the change of level structure through $f$ is a condition on the
values $\eta(f(a))$ of the universal level structure $\eta$ at the elements $a\in\Lambda$, and $\eta$ is a
homomorphism that is trivial on $M_b\Lambda$. As $s_g\equiv\lambda_g$ modulo $M_b\Lambda$, the level action of
$\ell_i$ is the level action through the $\lambda_g$. So the map of the weight tie carries it to the
action of $G$ on $FW_i$.
% Lean: T.finiteGroup_levelStructure_action_fixedRuns = LevelActionOnto.levelPin_actionPin, used in
% exists_runs_at_alt with the core Hk3.Core of the characterization, the one-map pins of the links,
% p.rhoGeo, the one-map rows (PointConditions'.oneMap) and the level ideal CoreTies.level. The proof:
% ell i g = quotEquiv Mb ((O i).lev g) (ofBijective if quotMap is bijective, refl otherwise); hlev is
% (hpins i).2.lift.lev (LevIs for (O i).lev); hι = TraceLattice.latEmb_iotaW h i (the laws of the run),
% hβ = hc.lev.βW i, and congr_of_levIs gives h1 and hm; hbij is quotMap_bijective; hspec is liftLa_spec.
% The level pin: Mfr * Mb = (M) (J.Q.l.prod), Submodule.mem_smul_pointwise_iff_exists for y, and
% MarkCoset.latG_spec for finVec (iota y) in latGHat. The action pin: hkeep from the first clause
% (EquivAt) of the OneMapOf of hrow i, hrho' = (hrho i).1 hkeep (RhoGeo: ActRho), hact_eq =
% OneMapAllCompare.actLev_congr with hA : l.A = 1 from the laws of the run (h.j.pins) and hspec,
% FormTies.exists_isLevAct for act; cuspTo is fw.e i. Size: about 33,000 lines of Lean below it (text
% count), nearly all below Lemma is-onto-step-level-structure-cusp.
\end{proof}

\begin{lemma}\label{is-onto-step-rational}
Let $R$ be a realization of a definite-side datum $D$ and $F$ a face (Definition~\ref{is-realisation}),
with one-map data and the subfield $K$ of the host. Suppose that $R$ satisfies its characterization at
$F$, condition (11) of Definition~\ref{is-primitive-rows} and the rank condition of
Lemma~\ref{is-onto-step-links}(3); that its rows have the base links; that its one-map data are pinned
at every lift; that $R$ has bounded rows with the phases $1$ (Definition~\ref{is-bounded-rows}) and its
ramified sets satisfy the guard of (2) of that definition; and that for every tuple $t$, with
$M_t=g_0g_1\prod_qq$ where $q$ runs over the orders of the residue fields of the places in $Q_t$, the
field $L$ contains the ray class field of $K$ of modulus $M_t$, in the following form: for every prime
$\mathfrak P$ of $L$ not dividing $M_t$, with $\mathfrak p=\mathfrak P\cap K$, there are $f\geq1$ and
$\alpha\in\mathcal O_K$ with $N\mathfrak P=N\mathfrak p^{f}$, $\mathfrak p^{f}=(\alpha)$ and
$\alpha\equiv1\pmod{M_t}$. Then $R$ is rational at the twisted places
(Definition~\ref{is-ties-class-sums}(4)).
\end{lemma}
\leanref{Development/AbelianVarietyTheory/Proofs/DivisorProofs/TwistedClassSumsRationalOfRayClassFieldsStatement.lean}{GlobalTheory.T.twistedClassSums_rational_of_rayClassFields}
\leanref{Development/AbelianVarietyTheory/BlockTheory/TwistedDiscreteSeries/RayFieldTheta.lean}{GlobalTheory.IndefiniteRowSideFullHecke.RayFieldTheta.ContainsRayField}
\begin{proof}
\openstep{The construction of the class sums at the twisted places from the ray class fields is not
written out.}
% Lean: T.twistedClassSums_rational_of_rayClassFields = TwistedRatBase.twistedRational_of_rayField.
% Size: about 196,000 lines of Lean in 11,700 theorems below it (text count, stopping at statements
% the book cites), about 1,000 of them below no other open step of Lemma is-onto-estimates.
\end{proof}

\subsection{Theta kernels and the base theta lattice}\label{subsection-is-base-theta-lattice}

\begin{definition}\label{is-onto-step-kernels}
Let $R$ be a realization of a definite-side datum with a face $F$, subgroup $P_0$ and character
$\psi_0$, let $j$ be a lift, and use the notation of Definitions~\ref{is-ties-class-sums}
and~\ref{is-ties-lattice-systems}. Let $v$ be a twisted place, with own prime $q$ and character
$\delta_v$ (Definition~\ref{is-onto-step-runs}(1)), and let $Q_v$ be the set of places of $E_h$
over $q$ in the ramified set $Q_{\mu_v}$ of the tuple $\mu_v$.
\begin{enumerate}
\item A \emph{base theta kernel} at $(j,v)$ is a pair $(\omega,\varphi)$ for which there are a unitary
  character $\chi$ of $\mathbf A_{E_h}^\times$ equal to $\delta_v$ and characters $\xi_P$ of the
  residue fields at the places $P$, restricting $\delta_v$ at the places of $Q_v$, such that $\omega$
  is a Weil representation of the pair $(V_j,W_j)$ for the theta characters of the tuple of $\mu_0$
  and $\chi$ at level $j+1$, and $\varphi$ is a theta test vector of $R$ with level at $Q_v$ given by
  the $\xi_P$.
\item For a base theta kernel $\kappa=(\omega,\varphi)$ and an element $h$ of the finite adelic group
  of $U(V_j)$, $\Theta_\kappa(h)$ is the set of $z\in L^{n_W(j)}$ for which there is $w$ integral for
  the lattice of $h$ such that $\Phi_j(z)$ is the trace $\sum_{g\in P_0}\sigma_D(\psi_0(g))\,g\vartheta$
  of the theta series $\vartheta$ of $\omega$ and $\varphi$ at $h$ for the polynomial of $w$, where
  $\sigma_D$ is the embedding of $L$ into $\mathbf C$ of the definite-side datum. Put
  $\Theta_\kappa=\bigcup_h\Theta_\kappa(h)$.
\item The \emph{base theta lattice} $\Theta^{\mathrm b}_v(j)$ is the $\mathcal O_L$-span of the sets
  $\Theta_\kappa$ over all base theta kernels $\kappa$ at $(j,v)$.
\item A family $(h_c)_{c\in C}$ of elements of the finite adelic group of $U(V_j)$ \emph{meets the
  classes} of the level $J_v$ of the data at $Q_v$ if every element of that group is $\gamma h_ck$ with
  $c\in C$, $\gamma\in U(V_j)(F_h)$ and $k\in J_v$.
\end{enumerate}
\end{definition}
\leanref{Development/AbelianVarietyTheory/Definitions/Foundations/ThetaSeriesTies.lean}{GlobalTheory.OnTies.IsKernel}
\leanref{Development/AbelianVarietyTheory/Proofs/FoundationProofs/KernelSpanOverClasses.lean}{GlobalTheory.TraceClassRank.IsBaseKernel}
\leanref{Development/AbelianVarietyTheory/Definitions/Foundations/ThetaSeriesTies.lean}{GlobalTheory.OnTies.seriesSet}
\leanref{Development/AbelianVarietyTheory/Proofs/FoundationProofs/KernelSpanOverClasses.lean}{GlobalTheory.TraceClassRank.seriesAt}
\leanref{Development/AbelianVarietyTheory/Definitions/Foundations/ThetaSeriesTies.lean}{GlobalTheory.OnTies.thetaBase}
\leanref{Development/AbelianVarietyTheory/Proofs/FoundationProofs/KernelSpanOverClasses.lean}{GlobalTheory.TraceClassRank.MeetsClasses}

\begin{lemma}\label{is-onto-step-kernel-unique-character}
Let $E$ be the host field of a definite-side datum and $\xi$ a Hecke character of $E$, a continuous
character of $\mathbf A_E^\times$ trivial on $E^\times$, with local components $\xi_P$ at the finite
places $P$. For a finite place $P$ and $k\in E^\times$ let $x_P(k)$ be the idele with component $k$ at
$P$ and $1$ at the other places. Let $c$ and $c'$ be continuous unitary characters of
$\mathbf A_E^\times$, trivial on $E^\times$, with $c(x_P(k))=c'(x_P(k))=\xi_P(k)$ for every finite
place $P$ and every $k\in E^\times$. Then $c=c'$.
\end{lemma}
\leanref{Development/AbelianVarietyTheory/Proofs/IwasawaDecomposition/HeckeCharacterIdeleCharacterMatchUnique.lean}{GlobalTheory.ChrEqUnique.idc_unique}
\leanref{Formalism/LocalToGlobal/IndefiniteReal/LinksDefs.lean}{Formalism.IndefR.Links.ChrEq}
\begin{proof}
Put $u=c\,c'^{-1}$. It is a continuous unitary character of $\mathbf A_E^\times$, trivial on $E^\times$
and on every $x_P(k)$.

First, $u(y)=1$ for every idele $y$ with archimedean component $1$. As $u$ is continuous, there is an
integer $N>0$ such that $u$ is trivial on the ideles with archimedean component $1$ whose finite
components $y_P$ are units with $y_P-1\in N\mathcal O_{E,P}$. Given $y$, let $S$ be the finite set of
the finite places at which $y$ is not a unit or which divide $N$, and choose for each $P\in S$ an
element $k_P\in E^\times$ with $|k_P|_P=|y_P|_P$ and $|y_P-k_P|_P\le|N|_P|k_P|_P$. Then
$y\prod_{P\in S}x_P(k_P)^{-1}$ is such an idele, so $u(y)=\prod_{P\in S}u(x_P(k_P))=1$.

Second, $u$ is trivial. Every idele is $y_\infty y_f$, with $y_\infty$ of finite components $1$ and
$y_f$ of archimedean component $1$, and $u(y_f)=1$ by the first step. For $\gamma\in E^\times$ with
archimedean part $\gamma_\infty$ and finite part $\gamma_f$,
$u(\gamma_\infty)=u(\gamma)u(\gamma_f)^{-1}=1$. As $E^\times$ is dense in $(E\otimes\mathbf R)^\times$
and $u$ is continuous, $u$ is trivial on the ideles with finite components $1$. Hence $u=1$, and
$c=c'$.
% Lean: idc_unique: c c'^{-1} satisfies ChrEq with the trivial Hecke character, and eq_one_of_chrEq
% concludes: HeckeFinTrivial.eq_one_of_trivial_on_finite (the second step, through a real chart of
% the archimedean adeles and the density of E) after CharLocProd.char_eq_prod_locK with
% CharFin.char_finUnit_finite_order (the integer N) and finiteIdele_exists_local_approx (the k_P).
\end{proof}

\begin{lemma}\label{is-onto-step-kernel-unique}
Let $R$ be a realization of a definite-side datum with host field $E_h/F_h$ and a face $F$, let $j$ be
a lift and $v$ a twisted place. Suppose that $R$ satisfies the core conditions at $F$
(Definition~\ref{is-eigenvalue-lift-data}), and that for every twisted place $v'$, with residue
characteristic $\ell$ and character $\delta_{v'}$ (Definition~\ref{is-onto-step-runs}(1)), there are a
Hecke character $\psi$ of $E_h$ and a finite set $Q$ of finite places of $E_h$ with
$\xi_{\mu_{v'}}=\delta_{v'}\psi$ such that: $\psi^{\ell^a}=1$ for some $a\geq0$; for every place $u$ of
$F_h$ the product of the components of $\psi$ at the places over $u$ is trivial on $F_h^\times$;
$\psi$ is tamely ramified at every finite place, and unramified at a finite place if and only if the
place is not in $Q$; its components at the places over $\ell$ are trivial; and no place of $Q$ is its
own complex conjugate. These hypotheses, and the hypothesis that $v$ is twisted, are not used in the
proof. Then any two base theta kernels at $(j,v)$ (Definition~\ref{is-onto-step-kernels}) have the
same theta series: for all $h$, $w$ and lowest $K$-type tensors, the theta series of the two kernels at
$h$ for the polynomial of $w$ agree.
\end{lemma}
\leanref{Development/AbelianVarietyTheory/Proofs/DivisorProofs/BaseThetaKernelSeriesUniqueStatement.lean}{GlobalTheory.T.baseThetaKernel_series_unique}
\leanref{Development/AbelianVarietyTheory/Proofs/FoundationProofs/KernelSpanOverClasses.lean}{GlobalTheory.TraceClassRank.BaseKernelUnique}
\begin{proof}
Let $\kappa=(\omega,\varphi)$ and $\kappa'=(\omega',\varphi')$ be base theta kernels at $(j,v)$, with
characters $\chi$, $\xi_P$ and $\chi'$, $\xi'_P$ as in Definition~\ref{is-onto-step-kernels}(1).

\emph{The characters.} Both $\chi$ and $\chi'$ are equal to $\delta_v$, so $\chi=\chi'$ by
Lemma~\ref{is-onto-step-kernel-unique-character}. Let $P\in Q_v$ and $u\in k_P^\times$. Then $u$ is the
reduction of some $a\in\mathcal O_{E_h}$, which is nonzero as $u\ne0$, and $\xi_P(u)$ and $\xi'_P(u)$
are both the value at $a$ of the local component of $\delta_v$ at $P$. So $\xi_P=\xi'_P$ for every
$P\in Q_v$.

\emph{The test vectors.} The vectors $\varphi$ and $\varphi'$ are level vectors, in the sense of
Lemma~\ref{is-eigenvalue-conjugate-theta-facts}, with the same vacuum, the vacuum vector of $R$ for the
pair $(V_j,W_j)$, for the same lattices and frames and the set $P_g\cup Q_v$, with the characters
$\xi_P$ and $\xi'_P$ respectively. Their place models are taken at the idempotent of
$\mathbf A_{E_h}$ attached to $P_g\cup Q_v$, the sum of the idempotents attached to its places; each
of these has component $1$ at its place and $0$ at the others, so it is unique. Hence a place model
$T$ of the common vacuum serves for both, and $\varphi=T(\Phi)$ and $\varphi'=T(\Phi')$ for the two
level functions. A level function is the product of factors that involve no character, an indicator
function and a factor at the places of $P_g$, with the factor $\prod_{P\in Q_v}\xi_P(d_P)^{-1}$, where
$d_P$ is the determinant of a block of the reduction of the argument at $P$ (the factor is $0$ if one
of these blocks is not invertible or the block below it is not $0$). So $\Phi=\Phi'$, and
$\varphi=\varphi'$.

\emph{The theta series.} Let $h$, $w$ and a tensor $u$ of lowest $K$-type vectors be given, and let
$P=\mathrm{io}_P(j)(u\otimes w)$. The theta series of $\kappa$ at $h$ for the polynomial of $w$ is
$g\mapsto\big(\omega(h,g)(\varphi\otimes P)\big)(0)$, where $\varphi\otimes P$ is formed with the
Cartan involutions $1$ of $V_j$ and those of the model of $W_j$ at $\sigma_0$ and the bases of $R$;
that of $\kappa'$ is the same expression with $\omega'$, as $\varphi'=\varphi$. By
Lemma~\ref{is-eigenvalue-conjugate-theta-facts}(3) $\varphi$ is a theta function, and by part (4)
of that lemma so is $\varphi\otimes P$. The Weil data $\omega$ and $\omega'$ have the same characters,
the theta characters of the tuple of $\mu_0$ and $\chi=\chi'$ at level $j+1$. So part (2) of that
lemma gives $\omega(h,g)(\varphi\otimes P)=\omega'(h,g)(\varphi\otimes P)$ for every $g$.
% Lean: T.baseThetaKernel_series_unique = BaseKernelUniq.baseKernelUnique: chrEq_unique (idc_unique,
% Lemma is-onto-step-kernel-unique-character) for chi; res_unique for the xi_P; thPVec_unique
% (idemSum_unique with HeckeParts.idem_unique, the uniqueness of the prime idempotents of
% IdemSupport.supported_group_and_prime_idempotents, and omegaQ_congr) for phi; then
% VecTheta.isThetaFunO_thPVec, InnerSteps.theta_pureTensorO and WeilUniq.weilDatum_unique, with the
% standard character from the Weil datum and IsInfFam theta from the vacuum. singleSeries is
% thetaKO omega h g (pureTensorO ... phi (ioP i (tprod u) w)). The hypotheses Hk3.Core R F, v in
% R.twPl and R.tup_mu_on are not used.
\end{proof}

\begin{lemma}\label{is-onto-step-kernel-span-move}
Let $R$ be a realization of a definite-side datum with host field $E_h/F_h$ and a face $F$, and suppose
that $R$ satisfies the core conditions at $F$. Let $j$ be a lift, $v$ a twisted place and $t$ a tuple,
with the notation of Definitions~\ref{is-onto-step-kernels} and~\ref{is-ties-lattice-systems}:
$\Gamma=U(V_j)(F_h)$, $G_f=U(V_j)(\mathbf A_{F_h,f})$, the level $J_v$, the coefficient space $W_c$ with
its action $\rho$, and vectors integral for the lattice of an element of $G_f$. Let $h,h'\in G_f$,
$\gamma\in\Gamma$ and $k\in J_v$ with $h=\gamma h'k$, and let $w\in W_c$ be integral for the lattice of
$h$. Then $\rho(\gamma)^{-1}w$ is integral for the lattice of $h'$.
\end{lemma}
\leanref{Development/AbelianVarietyTheory/Proofs/DivisorProofs/KernelThetaSeriesSpanOverClassRepresentatives.lean}{GlobalTheory.OneKernelSpan.inCoefLat_move}
\begin{proof}
By the core conditions the frame of $V_j$ is a frame of its standard lattice $L_0$, and $J_v$ lies in
the parahoric level of this frame, so $k\hat L_0=\hat L_0$. Let $\gamma_E\in GL_n(E_h)$ be the matrix of
$\gamma$. Then
\[
 L_h=E_h^n\cap\gamma h'k\hat L_0=E_h^n\cap\gamma h'\hat L_0=\gamma_E\big(E_h^n\cap h'\hat L_0\big)
 =\gamma_EL_{h'} .
\]
The construction $M\mapsto\nabla_{\lambda'}(M^\vee)$ of Definition~\ref{pd-weight-lattices} is
compatible with this change of lattice: $\nabla_{\lambda'}((\gamma_EM)^\vee)
=\rho_{\lambda'}(\check\gamma)\,\nabla_{\lambda'}(M^\vee)$. Now $w=\iota_{\lambda'}(P)$ for the map
$\iota_{\lambda'}$ of a family tie of $R$ and some $P\in\nabla_{\lambda'}(L_h^\vee)$, so
$P=\rho_{\lambda'}(\check\gamma)P'$ with $P'\in\nabla_{\lambda'}(L_{h'}^\vee)$. By the equivariance of
$\iota_{\lambda'}$ (Definition~\ref{is-ties-lattice-systems}(4)),
$w=\iota_{\lambda'}(\rho_{\lambda'}(\check\gamma)P')=\rho(\gamma)\iota_{\lambda'}(P')$, and
$\rho(\gamma)^{-1}w=\iota_{\lambda'}(P')$ is integral for the lattice of $h'$.
% Lean: inCoefLat_move (stated for the level dataLevel R i Q of any finite set Q): parLev_St for k,
% HLatMove.hLat_rat_mul_stab with the frame hR.lev.betaV i, Col.nabAtr_gmap, and the equivariance
% (hT.mapL i).1 of the family tie.
\end{proof}

\begin{definition}\label{is-onto-step-kernel-span-twisted}
Let $R$ be a realization of a definite-side datum $D$ with host field $E_h/F_h$, with the notation of
Definition~\ref{is-eigenvalue-lift-data}: the level primes $g_0,g_1$ and the integers $Q_S$, $q$ and $d$ of the host,
the prime $T$ of the system of spaces of $D$ (Definition~\ref{th-spaces}), the Hecke characters $\mu^0_j$ of $R$ and
the Hecke character $\tilde\eta$ of $R$. Let $j$ be a lift, $\nu$ a Hecke character of $E_h$ with local components
$\nu_P$, and $Q$ a finite set of finite places of $E_h$. For a finite place $P$ and $k\in E_h^\times$ let $x_P(k)$ be
the idele with component $k$ at $P$ and $1$ at the other places. A \emph{twisted kernel} of $R$ at $j$ for $(\nu,Q)$
is a tuple $(\chi,(\xi_P)_P,\omega,\varphi)$ of a unitary character $\chi$ of $\mathbf A_{E_h}^\times$, characters
$\xi_P$ of the residue fields $k_P^\times$ at the finite places $P$ of $E_h$, a Weil datum $\omega$ for $(V_j,W_j)$
(Definition~\ref{th-weil-datum}(3)) for the standard character of $R$, and a function $\varphi$ on the space
$\mathbb W$ of the pair $(V_j,W_j)$, such that:
\begin{enumerate}
\item $\chi$ is equal to $\nu$: it is continuous and trivial on $E_h^\times$, and $\chi(x_P(k))=\nu_P(k)$ for every
  finite place $P$ and every $k\in E_h^\times$;
\item $\nu^n=1$ for some integer $n\geq1$;
\item for every place $u$ of $F_h$ the product of the components of $\nu$ at the places of $E_h$ over $u$ is
  trivial on $F_h^\times$;
\item $\nu$ is tamely ramified at every finite place;
\item the residue field at every place $P$ of $Q$ has prime order $p_P$, with $p_P\ne g_0,g_1$ and $p_P\notin
  P^2$; and no place of $Q$ is the complex conjugate of a place of $Q$, itself included;
\item no place of $Q$ contains the integer $Q_Sqd\,T$;
\item $\nu$ is ramified at a finite place $P$ if and only if $P\in Q$ or $P$ is the complex conjugate of a place
  of $Q$;
\item for every $P\in Q$, $\xi_P$ restricts $\nu$ at $P$: $\xi_P(a\bmod P)=\nu_P(a)$ for every $a\in\mathcal
  O_{E_h}$ not in $P$;
\item $\omega$ has the theta characters of the tuple of $\mu_0$ and $\chi$ at the level $j+1$, that is, the
  characters $(\mu^0_j\chi,\mu^0_j\chi\tilde\eta^{-1})$;
\item $\varphi$ is a level vector for the Cartan involutions $1$ of $V_j$ and those of the model of $W_j$ at
  $\sigma_0$, the lattices and adelic matrices of $R$ for $(V_j,W_j)$, $P_g$, $Q$ and the $\xi_P$, whose vacuum is
  the vacuum vector of $R$ for $(V_j,W_j)$ (Lemma~\ref{is-eigenvalue-conjugate-theta-facts}).
\end{enumerate}
The \emph{level $J_Q$ of the data at $Q$} is the level
(Definition~\ref{asm-dependence-test-assembly-realization-equivariant-levels}(2)) of the level datum on $V_j$ with
the lattice and adelic matrix of $R$ at $V_j$, the integers $r=\operatorname{rk}W_j$, $t=\operatorname{rk}V_j-r$ and
$s=1$, the sets $P_g$ and $Q$, the trivial subgroup of $G_0$ and trivial characters, read in the finite adelic group
of $U(V_j)$. The level $J_v$ of Definition~\ref{is-onto-step-kernels}(4) is $J_{Q_v}$.
\end{definition}
\leanref{Development/AbelianVarietyTheory/Definitions/Divisors/KerDatTuple.lean}{GlobalTheory.OnCmp.KerDat}
\leanref{Development/AbelianVarietyTheory/Definitions/Divisors/KerDatTuple.lean}{GlobalTheory.OnCmp.IsWeilTw}
\leanref{Development/AbelianVarietyTheory/Proofs/FoundationProofs/KernelSpanOverClasses.lean}{GlobalTheory.TraceClassRank.dataLevel}

\begin{lemma}\label{is-onto-step-kernel-span-twisted-base}
Let $R$, $F$, $j$ and $v$ satisfy the hypotheses (1) and (2) of Lemma~\ref{is-onto-step-kernel-span} below,
with the notation of Definition~\ref{is-onto-step-kernels}, and let $\kappa=(\omega,\varphi)$ be a base theta kernel
at $(j,v)$. Then there are a unitary character $\chi$ of $\mathbf A_{E_h}^\times$ and characters $\xi_P$ of the
residue fields such that $(\chi,(\xi_P),\omega,\varphi)$ is a twisted kernel of $R$ at $j$ for $(\delta_v,Q_v)$
(Definition~\ref{is-onto-step-kernel-span-twisted}).
\end{lemma}
\leanref{Development/AbelianVarietyTheory/Proofs/DivisorProofs/BaseTupleThetaKernelTameCharacterStatement.lean}{GlobalTheory.T.baseTuple_thetaKernel_tameCharacter}
\leanref{Development/AbelianVarietyTheory/Proofs/DivisorProofs/BaseTupleThetaKernelTameCharacter.lean}{GlobalTheory.OnBaseKer.kerDat_of_isBaseKernel}
\begin{proof}
Let $\chi$ and $(\xi_P)$ be characters for $\kappa$ as in Definition~\ref{is-onto-step-kernels}(1). They give
conditions (1), (8), (9) and (10) of Definition~\ref{is-onto-step-kernel-span-twisted} for $\nu=\delta_v$ and
$Q=Q_v$: $\chi$ is equal to $\delta_v$, the $\xi_P$ restrict $\delta_v$ at the places of $Q_v$, $\omega$ has the theta
characters of the tuple of $\mu_0$ and $\chi$ at the level $j+1$, and $\varphi$ is a level vector for $P_g$, $Q_v$
and the $\xi_P$ whose vacuum is the vacuum vector of $R$. Let $q$ be the own prime of $v$ and $\ell$ its residue
characteristic.

\emph{Conditions (2) to (4).} The character $\delta_v$ is $\delta'$ if $\ell=p_{de}$, when $q=p'_{de}$, and $\delta$
otherwise, when $q=p_{de}$. By the core conditions on the level characters, the level data of $R$ form a level
system for $D$ (Definition~\ref{th-grid}), and by its condition (L3) $\delta_v$ has finite order, is trivial on
$\mathbf A_{F_h}^\times$, so that the product of its components at the places over a place of $F_h$ is trivial
on $F_h^\times$, and is tamely ramified at every finite place. Moreover it is unramified at every finite place
not over $q$, and at a place $P$ over $q$ its restriction to the units of $E_{h,P}$ has order $15$: an integer
$n\geq0$ satisfies $\delta_{v,P}(x)^n=1$ for every unit $x$ of $E_{h,P}$ if and only if $15\mid n$.

\emph{Conditions (5) and (6).} The set $Q_v$ lies in $Q_{\mu_v}$, and by (L5) of Definition~\ref{th-grid}
every ramified set $Q_t$ satisfies (5) and (6).

\emph{Condition (7).} Let $P$ be a finite place of $E_h$. Suppose first that $\delta_v$ is ramified at $P$. Then
$P$ lies over $q$. As $v$ is twisted, hypothesis (2) gives a Hecke character $\eta'$ with
$\xi_{\mu_v}=\delta_v\eta'$ and $\eta'^{\ell^a}=1$ for some $a\geq0$. If $\xi_{\mu_v}$ were unramified at $P$, then
$\delta_{v,P}(x)=\eta'_P(x)^{-1}$ for every unit $x$ at $P$, so $\delta_{v,P}(x)^{\ell^a}=1$; then $15\mid\ell^a$,
and $3$ and $5$ would both be equal to the prime $\ell$. So $\xi_{\mu_v}$ is ramified at $P$. By (L5) of
Definition~\ref{th-grid}, the ramified places of $\xi_t$ are the places of $Q_t$ and their complex
conjugates; so $P\in Q_{\mu_v}$, or $P$ is the conjugate of some $P'\in Q_{\mu_v}$. In the first case $P\in Q_v$,
as $P$ lies over $q$. In the second, $P$ contains the residue characteristic of $P'$, which is its own conjugate,
so $P$ and $P'$ lie over the same rational prime; hence $P'$ lies over $q$, and $P'\in Q_v$. Conversely, let $P$
lie in $Q_v$ or be the conjugate of a place of $Q_v$. In both cases $P$ lies over $q$, as $q$ is its own
conjugate. If $\delta_v$ were unramified at $P$, then $\delta_{v,P}(x)^1=1$ for every unit $x$ at $P$, and
$15\mid1$. So $\delta_v$ is ramified at $P$.
\end{proof}
% Lean: T.baseTuple_thetaKernel_tameCharacter = OnBaseKer.kerDat_of_isBaseKernel (68 lines), followed: chr, res,
% weil and vec from IsBaseKernel; fin, triv and tame from BigOn.isDelta_delAt_ownAt (Host.IsDelta (ownAt D v) of
% CharFrame.delAt (plc v): the fields isDel and isDel' of the CharFrame R.C); pl and off from hL.Q_pl and
% hL.Q_off at R.mu v; ram from hypothesis (2) (hmu v hv: psi with IsPsi.ord, psi^(l^a) = 1 and xi_mu = delta
% psi), hL.Q_ram, IsDelta.ine (15 | l^a, and 15 | 1 for the converse) and IsDelta.unr; the conjugate place P' is
% over q by hmem and hstar.

\begin{lemma}\label{is-onto-step-kernel-span-level-vector}
Let $R$ be a realization of a definite-side datum $D$ with host field $E_h/F_h$ and a face $F$ that satisfies
the core conditions at $F$ (Definition~\ref{is-eigenvalue-lift-data}), and let $\sigma_D$ be the embedding of
$L$ into $\mathbf C$ of $D$. Let $j$ be a lift, $\nu$ a Hecke character of $E_h$ whose values lie in
$\sigma_D(L)$, $Q$ a finite set of finite places of $E_h$, $(\chi,(\xi_P),\omega,\varphi)$ a twisted kernel of
$R$ at $j$ for $(\nu,Q)$, and $J_Q$ the level of the data at $Q$
(Definition~\ref{is-onto-step-kernel-span-twisted}). Conditions (3) and (6) of that definition are not used in
the proof. Then:
\begin{enumerate}
\item for every $k\in U(V_j)(\mathbf A_{F_h})$ with archimedean component $1$ and finite part in $J_Q$ there
  is a root of unity $z\in L$ with $\omega(k,1)\varphi=\sigma_D(z)\varphi$;
\item $\omega(1,k')\varphi=\varphi$ for every level lift $k'$ of $1$ in $U(W_j)(\mathbf A_{F_h,f})$
  (Lemma~\ref{is-eigenvalue-core-group}), read with archimedean component $1$.
\end{enumerate}
\end{lemma}
\leanref{Development/AbelianVarietyTheory/Proofs/DivisorProofs/ThetaKernelVectorMovedByRootsOfUnityStatement.lean}{GlobalTheory.T.thetaKernel_vector_moved_by_rootsOfUnity}
\leanref{Development/AbelianVarietyTheory/Proofs/DivisorProofs/ThetaKernelVectorMovedByRootsOfUnity.lean}{GlobalTheory.KerMove.kerDat_move}
\leanref{Development/AbelianVarietyTheory/Proofs/DivisorProofs/ThetaKernelVectorMovedByRootsOfUnity.lean}{GlobalTheory.KerMove.ramif_ker}
\leanref{Development/AbelianVarietyTheory/Proofs/DivisorProofs/ThetaKernelVectorMovedByRootsOfUnity.lean}{GlobalTheory.KerMove.unrOff_ker}
\leanref{Development/AbelianVarietyTheory/Proofs/DivisorProofs/ThetaKernelVectorMovedByRootsOfUnity.lean}{GlobalTheory.KerMove.val_ker}
\begin{proof}
Let $r=\operatorname{rk}W_j$ and $\tilde\eta$ the Hecke character of $R$. By condition (9), the Weil datum
$\omega$ has the characters $\chi_S=\mu^0_j\chi$ and $\chi_T=\mu^0_j\chi\tilde\eta^{-1}$, where $\mu^0_j$ is the
Hecke character of $R$ for the lift $j$, at the level $j+1$ (Definition~\ref{is-realisation}); by condition (10),
$\varphi$ is a level vector for $P_g$, $Q$ and the $\xi_P$ (Lemma~\ref{is-eigenvalue-conjugate-theta-facts}). By
the core conditions on the level characters, the level data of $R$ form a level system for $D$
(Definition~\ref{th-grid}). That definition numbers the characters $\mu^0$ and the spaces $V_i$ and $W_i$ from
$1$, so the character $\mu^0_j$ and the spaces $V_j$ and $W_j$ of the lift $j$ here are $\mu^0_{j+1}$, $V_{j+1}$
and $W_{j+1}$ there.

\emph{The places of $Q$.} Let $P$ be a finite place that contains $g_0$ or $g_1$. By condition (5), every place
$P'$ of $Q$ contains a prime $p_{P'}\ne g_0,g_1$, and so does its complex conjugate; as a place contains only one
rational prime, $P$ is neither in $Q$ nor conjugate to a place of $Q$. By (L1) of Definition~\ref{th-grid}, $P_g$
is the set of the places over $\mathfrak p_0$ or $\mathfrak p_1$, primes of $K$ with $g_l\in\mathfrak p_l$, so
every place of $P_g$ contains $g_0$ or $g_1$. Hence $P_g$ and $Q$ are disjoint, the places of $Q$ contain neither
$g_0$ nor $g_1$, and by conditions (1) and (7) $\chi$ is unramified at the places of $P_g$.

\emph{The level vector.} By (L7) of Definition~\ref{th-grid}, for $\Lambda_{V_{j+1}}$ and $\Lambda_{W_{j+1}}$
there, the adelic matrices of $V_j$ and $W_j$ are frames of self-dual lattices. By the core conditions on the
level characters, $\tilde\eta$ is equal to the character $\eta_L$ of (L3) of Definition~\ref{th-grid}; so it is
a Hecke character unramified at every finite place with archimedean exponents $\frac12$, and its restriction
to $\mathbf A_{F_h}^\times$ is $\eta$, as (L3) states for $\eta_L$. The Weil datum also gives this restriction: by
Definition~\ref{th-weil-datum}(3) the restrictions of $\chi_S$ and $\chi_T$ to $\mathbf A_{F_h}^\times$ are
$\eta^{\operatorname{rk}V_j}$ and $\eta^{\operatorname{rk}W_j}$, so $\tilde\eta=\chi_S\chi_T^{-1}$ restricts to
$\eta^{\operatorname{rk}V_j-\operatorname{rk}W_j}=\eta$, as $\operatorname{rk}V_j+\operatorname{rk}W_j=2r+t$ is
odd. So Lemma~\ref{is-eigenvalue-conjugate-frame-level-vector} applies to $S=V_j$, $T=W_j$, the frames of
$R$, $\omega$ and $\varphi$, with $P_g$ and $Q$; let $e$ be its idempotent, the sum of the idempotents
at the places of $P_g\cup Q$.

\emph{The characters off $P_g\cup Q$.} By (L2) of Definition~\ref{th-grid}, for $\mu^0_{j+1}$ there, the
character $\mu^0_j$ is unramified at every finite place that contains neither $g_0$ nor $g_1$. The places that
contain $g_0$ or $g_1$ lie over $\mathfrak p_l$ or $\bar{\mathfrak p}_l$, so they are in $P_g$ or conjugate to one
of its places. So $\mu^0_j$ is unramified at every finite place that is not in $P_g$ and not conjugate to a place
of $P_g$. By conditions (1) and (7), $\chi$ is unramified outside $Q$ and its conjugates, and $\tilde\eta$ is
unramified everywhere. A Hecke character that is unramified at every finite place outside a set of places and
their conjugates is trivial on every idele that has archimedean component $1$, is a unit at every finite place
and has component $1$ at the places of the set and their conjugates, as it is the product of its local
components. So $\chi_S$ and $\chi_T$ are trivial on these ideles for the set $P_g\cup Q$, as
Lemma~\ref{is-eigenvalue-conjugate-frame-level-vector} requires.

\emph{The values at $\det_e$.} Let $x\in U(V_j)(\mathbf A_{F_h})$ or $x\in U(W_j)(\mathbf A_{F_h})$ have
archimedean component $1$ and finite part stabilizing the lattice of the frame. Then $\det x$ is an
integral idele whose residue at every finite place $P$ is $\det\bar x_P$, and $\det_ex=ea+1-e$ with
$a=\det x$. A Hecke character that is tamely ramified at the places of $P_g\cup Q$, with the characters
$\rho_P$ of $k_P^\times$ there, takes at $\det_ex$ the value $\prod_{P\in P_g\cup Q}\rho_P(\det\bar x_P)$. By
(L6) of Definition~\ref{th-grid}, for $\mu^0_{j+1}$ there, $\mu^0_j$ has the characters $\chi_{j+1,P}^{-1}$ at
$P\in P_g$; and by (L2) it is unramified at the places of $Q$, which contain neither $g_0$ nor $g_1$. By
conditions (1), (4) and (8), $\chi$ is tamely ramified with the characters $\xi_P$ at $P\in Q$, and it is
unramified at the places of $P_g$. So
\[
 \mu^0_j(\det_ex)=\prod_{P\in P_g}\chi_{j+1,P}(\det\bar x_P)^{-1},\qquad
 \chi(\det_ex)=\prod_{P\in Q}\xi_P(\det\bar x_P),\qquad\tilde\eta(\det_ex)=1 .
\]

\emph{Roots of unity.} For $P\in Q$ and $u\in k_P^\times$, by condition (8) $\xi_P(u)$ is the value of $\nu$ at
the idele $x_P(a)$ for an element $a\in\mathcal O_{E_h}$ that reduces to $u$; it lies in $\sigma_D(L)$ by
hypothesis and has finite order by condition (2). For $g\in G$, $\prod_{P\in P_g}\chi_{j+1,P}(\det g_P)$ lies in $\sigma_D(L)$ by the
core conditions on the group, and it has finite order; applying this to the
element of $G$ with components $\operatorname{diag}(u_P,1,\dots,1)$ shows that
$\prod_{P\in P_g}\chi_{j+1,P}(u_P)$ is $\sigma_D$ of a root of unity of $L$ for all $u_P\in k_P^\times$.

(1) Let $k$ be as in (1). As its finite part lies in $J_Q$, $k$ stabilizes the lattice of $V_j$, its
reductions at the places of $P_g\cup Q$ are block upper triangular for the blocks
$(r,\operatorname{rk}V_j-r)$, and at $P\in P_g$ its upper left block $\bar k^{\mathrm{ul}}_P$ is $1$. So $(k,1)$
satisfies the conditions of Lemma~\ref{is-eigenvalue-conjugate-frame-vacuum}(2), with
$\bar k^{\mathrm{ul}}_P=1=\bar1_P$ at $P\in P_g$, and Lemma~\ref{is-eigenvalue-conjugate-frame-level-vector}
gives
\[
 \prod_{P\in Q}\xi_P(\det\bar k^{\mathrm{ul}}_P)^{-1}\cdot\omega(k,1)\varphi=\chi_T(\det_ek)^{-1}\varphi .
\]
The factor on the left is $\sigma_D$ of a root of unity of $L$. By the values at $\det_e$,
$\chi_T(\det_ek)^{-1}=\prod_{P\in P_g}\chi_{j+1,P}(\det\bar k_P)\prod_{P\in Q}\xi_P(\det\bar k_P)^{-1}$, which is
also $\sigma_D$ of a root of unity of $L$. Dividing gives $z$.

(2) Let $k'$ be a level lift of $1$. It stabilizes the lattice of $W_j$, and at $P\in P_g$ its reduction
$\bar k'_P$, an $r\times r$ matrix, is $1$. So $(1,k')$ satisfies the conditions of
Lemma~\ref{is-eigenvalue-conjugate-frame-vacuum}(2), with $\bar1^{\mathrm{ul}}_P=1=\bar k'_P$, and
Lemma~\ref{is-eigenvalue-conjugate-frame-level-vector} gives
\[
 \prod_{P\in Q}\xi_P(\det\bar k'_P)\cdot\omega(1,k')\varphi=\chi_S(\det_ek')\,\varphi .
\]
By the values at $\det_e$, as $\det\bar k'_P=1$ at $P\in P_g$,
$\chi_S(\det_ek')=\mu^0_j(\det_ek')\chi(\det_ek')=\prod_{P\in Q}\xi_P(\det\bar k'_P)$, which is nonzero. So
$\omega(1,k')\varphi=\varphi$.
% Lean: kerDat_move, for any KerDat R i xi chi xiR Q omega phi; of its fields it uses chr, fin, tame, pl, ram,
% res, weil and vec, and not triv or off. The places of Q are ramif_ker (exists_prime_mem from pl,
% not_two_primes, natCast_mem_conj, typePrime_mem). The level vector: hL := hc.lev, hW := hK.weil, hη from
% T.isUnramifiedEtaTilde_of_odd_exponentSum (hL.eta, R.C.isEta, hW.resS, hW.resT, D.R.t_odd), the
% self-dual lattices from T.isSelfDualCoordLattice_of_selfDualFrame (hL.betaV i, hL.betaW i), and
% VecInv.om_vec (Lemma is-eigenvalue-conjugate-frame-level-vector) for hK.vec. The characters off P_g u Q
% are unrOff_ker (ramif_ker; T.heckeChar_trivial_on_idempotentIntegralUnits for mu0, chi and eta); the
% values are val_ker (T.tameHeckeChar_splitSupported_eq_prod_residue three times, with hL.chi, hK.res and
% IsEtaT.unr) at detAt e (VecInv.exists_det_int). Roots of unity: rootLL_res (hv, hK.fin, hK.res) and
% rootLL_chi (hc.grp.chiW at R.eG.symm of diagU). (1) is hom (k, 1) with parLev_St, parLev_g, parLev_q,
% the trivial top blocks (hlev.2) and smul_root; (2) is hom (1, k') with hk'.2, det_twoDiagBlocks_properties
% (rW) and smul_fix.
\end{proof}

\begin{lemma}\label{is-onto-step-kernel-span-level}
Let $R$ be a realization of a definite-side datum $D$ with host field $E_h/F_h$ and a face $F$ that satisfies
the core conditions at $F$ (Definition~\ref{is-eigenvalue-lift-data}), and let $\sigma_D$ be the embedding of
$L$ into $\mathbf C$ of $D$. Let $j$ be a lift, $\nu$ a Hecke character of $E_h$ whose values lie in
$\sigma_D(L)$, $Q$ a finite set of finite places of $E_h$, $(\chi,(\xi_P),\omega,\varphi)$ a twisted kernel of
$R$ at $j$ for $(\nu,Q)$, and $J_Q$ the level of the data at $Q$
(Definition~\ref{is-onto-step-kernel-span-twisted}). Conditions (3) and (6) of that definition are not used in
the proof. Let $W_c$ be the coefficient space of $V_j$ (Definition~\ref{is-ties-lattice-systems}). For
$h\in U(V_j)(\mathbf A_{F_h})$ and $w\in W_c$ let $\vartheta(h,w)$ be the theta series of $\omega$ and $\varphi$
at $h$ for the polynomial of $w$, a function of the tensors of lowest $K$-type vectors and of $g\in
U(W_j)(\mathbf A_{F_h})$. Let $k\in U(V_j)(\mathbf A_{F_h})$ have archimedean component $1$ and finite part in
$J_Q$. Then there is a root of unity $z\in L$ with $\vartheta(hk,w)=\sigma_D(z)\vartheta(h,w)$ for all $h$ and
$w$.
\end{lemma}
\leanref{Development/AbelianVarietyTheory/Proofs/DivisorProofs/KernelThetaSeriesSpanOverClassRepresentatives.lean}{GlobalTheory.OneKernelSpan.series_level}
\begin{proof}
By Lemma~\ref{is-onto-step-kernel-span-level-vector}(1) there is a root of unity $z\in L$ with
$\omega(k,1)\varphi=\sigma_D(z)\varphi$.
% Lean: the first clause of T.thetaKernel_vector_moved_by_rootsOfUnity (KerMove.kerDat_move) for the KerDat hK,
% with hval for the values (Lemma is-onto-step-kernel-span-level-vector(1)).

Let $P$ be the polynomial of $w$ for a tensor of lowest $K$-type vectors. As $\varphi$ is a level vector by
condition (10), Lemma~\ref{is-eigenvalue-conjugate-theta-facts}(3) and (4) show that $\varphi$ and
$\varphi\otimes P$ are theta functions; by (4), applied to $(k,1)$, whose archimedean components are $1$,
$\omega(k,1)(\varphi\otimes P)=(\omega(k,1)\varphi)\otimes P=\sigma_D(z)(\varphi\otimes P)$. As $\omega$
is a homomorphism, for every $g$
\[
 \theta_\omega(hk,g;\varphi\otimes P)=\big(\omega(h,g)\,\omega(k,1)(\varphi\otimes P)\big)(0)
 =\sigma_D(z)\,\theta_\omega(h,g;\varphi\otimes P),
\]
which is $\vartheta(hk,w)=\sigma_D(z)\vartheta(h,w)$.
% Lean: series_level, for any KerDat R i xi chi xiR Q omega phi with the values of xi in the range of embLL:
% FiniteFock.om_pureTensorO_of_finite (pt_theta_of_kerDat, from the fields vec and weil, for the theta
% functions) and the second clause of thetaKO_rat_left_and_eigen_right.
\end{proof}

\begin{lemma}\label{is-onto-step-kernel-span-real-compact}
Let $D$, $\psi$, $\theta$, $S$ and $T$ be as in Lemma~\ref{is-eigenvalue-conjugate-theta-facts}, with Gram
matrices $G_S(\tau)$ and $G_T(\tau)$ at $\tau\in\mathcal T$, and put $n=\operatorname{rk}S$. Let $M$ be a model of
discrete series for $U(T)$ with frames $f_\tau$, signature matrices $J_\tau$, Cartan involutions
$s_\tau=f_\tau^{-1}J_\tau f_\tau$ and groups $K_\tau$ of unitary matrices commuting with $J_\tau$, whose signs at
$\tau$ are $1$ at the first $p_\tau$ indices and $-1$ at the other $q_\tau$, with $p_\tau+q_\tau=\operatorname{rk}T$.
Let $B_S$ be bases of $S$ with $B_{S,\tau}^*G_S(\tau)B_{S,\tau}=1$. For bases $B=(B_\tau)$ of $S$ and
$B'=(B'_\tau)$ of $T$ write $f\otimes_{B,B'}P$ for $f\otimes P$ formed with the Cartan involutions $1$ of $S$ and
$s_\tau$ of $T$ and these bases. For $k=(k_\tau)$ with $k_\tau\in K_\tau$ and unitary matrices $m=(m_\tau)$ of size
$n$, let $s_{k,m}$ be the substitution of the creation variables
\[
 (\tau,a,j)\longmapsto\sum_{a',b}c_{\tau,a'b}\,(\tau,a',b),\qquad
 c_{\tau,a'b}=\begin{cases}(k_\tau)_{a'a}(m_\tau^{-1})_{jb}&\text{if }a\le p_\tau,\\
 (k_\tau^{-1})_{aa'}(m_\tau)_{bj}&\text{if }a>p_\tau.\end{cases}
\]
\begin{enumerate}
\item (Change of bases.) For such $k$ and $m$, every theta function $f$ for $(S,T)$ and every polynomial $P$ in
  the creation variables, $f\otimes_{B_Sm,f^{-1}k}P=f\otimes_{B_S,f^{-1}}s_{k,m}(P)$, where
  $(B_Sm)_\tau=B_{S,\tau}m_\tau$ and $(f^{-1}k)_\tau=f_\tau^{-1}k_\tau$.
\item (Compact elements.) Let $\omega$ be a Weil datum for $(S,T)$ with characters $(\chi_S,\chi_T)$, let
  $a_\tau$ and $b_\tau$ be archimedean exponents of $\chi_S$ and $\chi_T$ at $\tau$, and suppose that
  $x_{1,\tau}=a_\tau+\frac n2$, $x_{2,\tau}=a_\tau-\frac n2$ and $x_{3,\tau}=-\big(\frac{p_\tau-q_\tau}2+b_\tau\big)$
  are integers. Let $f$ be a theta function annihilated by the annihilation operators for $1$ and $s_\tau$ at
  every $\tau$. Let $g\in U(S)(\mathbf A)$ and $h\in U(T)(\mathbf A)$ have finite components $1$, with images
  $g_\tau=\theta_\tau(g)$ and $h_\tau=\theta_\tau(h)$, such that $k_\tau=f_\tau h_\tau f_\tau^{-1}$ lies in
  $K_\tau$ for every $\tau$. Then $m_\tau=B_{S,\tau}^{-1}g_\tau B_{S,\tau}$ is unitary for every $\tau$, and for
  every polynomial $P$
  \[
   \omega(g,h)\big(f\otimes_{B_S,f^{-1}}P\big)=\prod_{\tau\in\mathcal T}\det(k^+_\tau)^{x_{1,\tau}}
   \det(k^-_\tau)^{x_{2,\tau}}\det(m_\tau)^{x_{3,\tau}}\cdot f\otimes_{B_S,f^{-1}}s_{k,m}(P),
  \]
  where $k^+_\tau$ and $k^-_\tau$ are the diagonal blocks of $k_\tau$ of sizes $p_\tau$ and $q_\tau$.
\end{enumerate}
\end{lemma}
\leanref{Development/AbelianVarietyTheory/Proofs/IwasawaDecomposition/FockPolynomialVectorChangeOrthonormalBasesSubstStatement.lean}{GlobalTheory.T.fockPolynomialVector_change_orthonormalBases_subst}
\leanref{Development/AbelianVarietyTheory/Proofs/IwasawaDecomposition/FockPolynomialVectorChangeOrthonormalBasesSubst.lean}{GlobalTheory.FockBasis.pureTensorO_basis}
\leanref{Development/AbelianVarietyTheory/Proofs/IwasawaDecomposition/WeilMaximalCompactFockSubstitutionCharacterStatement.lean}{GlobalTheory.T.weil_maximalCompact_fock_substitution_character}
\leanref{Development/AbelianVarietyTheory/Proofs/IwasawaDecomposition/WeilMaximalCompactFockSubstitutionCharacter.lean}{GlobalTheory.FockCompactAction.om_pureTensorO_eq_fact}
\begin{proof}
(1) This is a computation in the Fock model at the archimedean places. By the proof of
Lemma~\ref{is-eigenvalue-conjugate-theta-facts}, the creation operator of the variable $(\tau,a,j)$ for the
bases $B$ and $B'$ is $r_\tau\big(\partial_\tau(z)-i\,\partial_\tau(is_\tau z)\big)$ with
$z=B'_\tau E_{aj}B_\tau^{-1}$, for a number $r_\tau>0$ that depends only on $\tau$. Put
$z_{a'b}=f_\tau^{-1}E_{a'b}B_{S,\tau}^{-1}$, so that $is_\tau z_{a'b}=\pm iz_{a'b}$, with the sign of $a'$. For the
bases $B_Sm$ and $f^{-1}k$, $z=f_\tau^{-1}k_\tau E_{aj}m_\tau^{-1}B_{S,\tau}^{-1}=\sum_{a',b}(k_\tau)_{a'a}
(m_\tau^{-1})_{jb}z_{a'b}$, and $(k_\tau)_{a'a}=0$ unless $a$ and $a'$ have the same sign, as $k_\tau$ commutes with
$J_\tau$. The operator $\partial_\tau$ is real linear, so $z\mapsto\partial_\tau(z)-i\,\partial_\tau(is_\tau z)$
is complex linear on the span of the $z_{a'b}$ with $a'\le p_\tau$ and conjugate linear on the span of those
with $a'>p_\tau$. As $k_\tau$ and $m_\tau$ are unitary, $\overline{(k_\tau)_{a'a}(m_\tau^{-1})_{jb}}=
(k_\tau^{-1})_{aa'}(m_\tau)_{bj}$. So the creation operator of $(\tau,a,j)$ for $B_Sm$ and $f^{-1}k$ is
$\sum_{a',b}c_{\tau,a'b}$ times the creation operator of $(\tau,a',b)$ for $B_S$ and $f^{-1}$. The creation
operators commute on theta functions, by Lemma~\ref{is-eigenvalue-conjugate-frame-extremal-fock}(2) as in the
proof of its part (3), and $f\otimes P$ is $P$ evaluated at the creation operators applied to $f$; this gives
(1).

(2) As $g_\tau$ preserves $G_S(\tau)$ and $B_{S,\tau}^*G_S(\tau)B_{S,\tau}=1$, $m_\tau$ is unitary. As $k_\tau$
commutes with $J_\tau$, $h_\tau$ commutes with $s_\tau$, and $g_\tau$ commutes with $1$. So
Lemma~\ref{is-eigenvalue-conjugate-theta-facts}(4), with the same Cartan involutions and the bases
$g_\tau B_{S,\tau}=B_{S,\tau}m_\tau$ and $h_\tau f_\tau^{-1}=f_\tau^{-1}k_\tau$, gives
\[
 \omega(g,h)\big(f\otimes_{B_S,f^{-1}}P\big)=\big(\omega(g,h)f\big)\otimes_{B_Sm,f^{-1}k}P .
\]
The trace of $s_\tau$ is that of $J_\tau$, which is $p_\tau-q_\tau$, so $x_{3,\tau}=-(\frac12\operatorname{tr}s_\tau+
b_\tau)$, and $1$ is a Cartan involution of $G_S(\tau)=(B_{S,\tau}^{-1})^*B_{S,\tau}^{-1}$. By
Lemma~\ref{is-eigenvalue-conjugate-frame-phase-vacuum}, $\omega(g,h)f=cf$ with
$c=\prod_\tau(d^+_\tau)^{x_{1,\tau}}(d^-_\tau)^{x_{2,\tau}}\det(g_\tau)^{x_{3,\tau}}$. Conjugation by $f_\tau$
carries $\frac12(1\pm s_\tau)h_\tau+\frac12(1\mp s_\tau)$ to $\frac12(1\pm J_\tau)k_\tau+\frac12(1\mp J_\tau)$, and
$k_\tau$ is block diagonal, so $d^+_\tau=\det k^+_\tau$ and $d^-_\tau=\det k^-_\tau$; and $\det g_\tau=\det
m_\tau$. As $f\otimes P$ is linear in $f$, the right side above is $c\cdot f\otimes_{B_Sm,f^{-1}k}P$, and (1)
gives the claim.
\end{proof}
% Lean: (1) is T.fockPolynomialVector_change_orthonormalBases_subst = FockBasis.pureTensorO_basis: aPlus_unitMap_mul
% (the creation operator for the new bases as the combination with the coefficients of fsub, through
% cplxO_unitMap) and pure_subst with FockIndep.aPlus_comm' (the creation operators commute); fsub M.eps k u
% is s_{k,m}. (2) is T.weil_maximalCompact_fock_substitution_character = FockCompactAction.om_pureTensorO_eq_fact:
% conj_unitary for m, hc (h commutes with cart M), trace_Jm, T.vacuum_character_positive_definite_source (Lemma
% is-eigenvalue-conjugate-frame-phase-vacuum), OrientedHeis.om_pureTensorO (Lemma
% is-eigenvalue-conjugate-theta-facts(4)), (1), pureTensorO_smul, and det_conj for detP and detM, which are
% det k^+ and det k^-. The factor and the substitution are fact = fchar . fsub. Lemma
% is-eigenvalue-conjugate-frame-extremal-compact(1) is the case g = 1, h diagonal at one place.

\begin{lemma}\label{is-onto-step-kernel-span-real-fock-weights}
Let $D$ be a definite-side datum with ranks $k$, $m$, $t$ and $r=k+m-1$ (Definition~\ref{th-ranks}) and set
$\mathcal T$ of embeddings of its host field, and let $j$ be a lift. Put $n_V=\operatorname{rk}V_j$, and for
$\tau\in\mathcal T$ let $(p_\tau,q_\tau)$ be the signature of $W_j$ at $\tau$, $e_j(\tau)$ the $j$-th number $e$ of the
Harish-Chandra parameter of $D$ at $\tau$ (Definition~\ref{th-hc-parameter}), $\lambda_\tau$ the sequence of
Lemma~\ref{is-onto-step-exp-dominant} at $\tau$, $a^\pm_\tau=a^\pm_\tau(\sigma_0)$ the positive and negative parts of
the weights of the $j$-th lift at $\tau$ (Definition~\ref{is-eigenvalue-lift-data}), of lengths $p_\tau$ and $q_\tau$,
and $\lambda'_\tau$ the dominant weight of length $n_V$ that the Harish-Chandra parameters of $D$ give at the level
$j+1$ (Definition~\ref{is-ties-lattice-systems}(4)). Let $M$ be a model of discrete series for $U(W_j)$
(Definition~\ref{asm-dependence-rows-models}(1)) with the parameter of the models of $W_j$ at $\sigma_0$, whose signs
at every $\tau$ are $1$ at the first $p_\tau$ indices and $-1$ at the others. Then there are integers $x_{1,\tau}$,
$x_{2,\tau}$ and $x_{3,\tau}$ for $\tau\in\mathcal T$ with
\[
 x_{1,\tau}=\frac{n_V}2-e_j(\tau),\qquad x_{2,\tau}=-\frac{n_V}2-e_j(\tau),\qquad
 x_{3,\tau}=-\Big(\frac{p_\tau-q_\tau}2-e_j(\tau)-\frac12\Big),
\]
such that for every $\tau$:
\begin{enumerate}
\item the weight $\Lambda_\tau$ of the lowest $K$-type of $M$ at $\tau$
  (Definition~\ref{is-eigenvalue-conjugate-k-types}(3)) is
  $(a^+_{\tau,1}+x_{1,\tau},\dots,a^+_{\tau,p_\tau}+x_{1,\tau},-a^-_{\tau,q_\tau}+x_{2,\tau},\dots,
  -a^-_{\tau,1}+x_{2,\tau})$;
\item for $1\le k'\le n_V$, $\lambda'_{\tau,k'}=c^+_{k'}-c^-_{k'}-x_{3,\tau}$, where $c^+_{k'}=a^+_{\tau,k'}$ if
  $k'\le p_\tau$ and $c^+_{k'}=0$ otherwise, and $c^-_{k'}=a^-_{\tau,n_V+1-k'}$ if $n_V+1-k'\le q_\tau$ and
  $c^-_{k'}=0$ otherwise.
\end{enumerate}
\end{lemma}
\leanref{Development/AbelianVarietyTheory/Proofs/IwasawaDecomposition/LowestKTypeCoefficientWeightEqMinorWeightsStatement.lean}{GlobalTheory.T.lowestKType_coefficientWeight_eq_minorWeights}
\leanref{Development/AbelianVarietyTheory/Proofs/IwasawaDecomposition/LowestKTypeCoefficientWeightEqMinorWeights.lean}{GlobalTheory.JointHarm.weights_of_minors}
\leanref{Development/AbelianVarietyTheory/Proofs/IwasawaDecomposition/ThetaWeightArithmeticAtPlaceProofStatement.lean}{GlobalTheory.T.thetaWeightArithmeticAtPlace_proof}
\begin{proof}
Fix $\tau$, and write $p=p_\tau$, $q=q_\tau$, $e=e_j(\tau)$ and $\lambda=\lambda_\tau$. The signature of $W_j$ gives
$p+q=r$, and $n_V=r+t$. The positive and negative parts of the weights of the lift are
\[
 a^+_{\tau,a}=\lambda_a+e-(p-a)-\frac{t+1}2\quad(1\le a\le p),\qquad
 a^-_{\tau,c}=-\Big(e+\lambda_{p+q+1-c}+(q+1-c)+\frac{t-1}2\Big)\quad(1\le c\le q).
\]

\emph{The integers.} As $k$ and $m$ are even, $r=k+m-1$ is odd, and $t$ is odd, so $n_V=r+t=2s$ for an integer
$s$; and $e$ is an integer (Definition~\ref{th-hc-parameter}(1)). So $x_{1,\tau}=s-e$ and $x_{2,\tau}=-s-e$ are
integers. As $p+q=r=2s'+1$ for an integer $s'$, $\frac{p-q}2-\frac12=p-s'-1$, and $x_{3,\tau}=-(p-s'-1-e)$ is an
integer.

(1) The parameter of $M$ at $\tau$ is the holomorphic parameter of the Harish-Chandra parameter of $D$ at $\tau$:
the sequence $\lambda$ with the signs $1$ at the first $p$ indices and $-1$ at the other $q$. By
Definition~\ref{is-eigenvalue-conjugate-k-types}(3), $2\rho_{n,a}=q$ and $2\rho_{c,a}=p+1-2a$ for $a\le p$, and
$2\rho_{n,p+c}=-p$ and $2\rho_{c,p+c}=q+1-2c$ for $1\le c\le q$. So
$\Lambda_{\tau,a}=\lambda_a+\frac q2-\frac{p+1}2+a$ and $\Lambda_{\tau,p+c}=\lambda_{p+c}-\frac p2-\frac{q+1}2+c$.
With the formulas above, $\Lambda_{\tau,a}-a^+_{\tau,a}=\frac{p+q+t}2-e=x_{1,\tau}$, and, with $c'=q+1-c$,
$\Lambda_{\tau,p+c}+a^-_{\tau,c'}=-\frac{p+q+t}2-e=x_{2,\tau}$; as $c$ runs from $1$ to $q$, $c'$ runs from $q$ down
to $1$.

(2) By Lemma~\ref{is-onto-step-exp-dominant}, with $J$ the level $j+1$, $\lambda'_{\tau,k'}=\lambda_{k'}+k'-
\frac{n_V}2-1$ for $k'\le p$, $\lambda'_{\tau,k'}=\frac{p-q-1}2-e$ for $p<k'\le p+t$, and $\lambda'_{\tau,p+t+c}=
\lambda_{p+c}+c+\frac{t+p-q}2-1$ for $1\le c\le q$, where $\lambda_{p+c}=-b_{J+c}$. For $k'\le p$,
$n_V+1-k'>q$, and $a^+_{\tau,k'}-x_{3,\tau}=\lambda_{k'}+k'-\frac{p+q+t}2-1$. For $p<k'\le p+t$ both $c^\pm_{k'}$
are $0$, and $-x_{3,\tau}=\frac{p-q-1}2-e$. For $k'=p+t+c$, $n_V+1-k'=q+1-c\le q$, and
$-a^-_{\tau,q+1-c}-x_{3,\tau}=e+\lambda_{p+c}+c+\frac{t-1}2+\frac{p-q-1}2-e$, which is the value above.
\end{proof}
% Lean: T.lowestKType_coefficientWeight_eq_minorWeights = JointHarm.weights_of_minors, for any M : DSModel (Wc D i)
% (parAt (Wc D i) sigma) with sigma = embLL and the signs M.eps tau a = (a < p). The integers: hnr (odd_r and
% t_odd: r + t = 2 s), e_int, and d_int of ArithAt (T.thetaWeightArithmeticAtPlace_proof = PairWeightArith.arithAt,
% from p + q = r odd and e_int), x1 = -e + s, x2 = -e - s, x3 = -dShift. (1) is HolDS.lowK_holPar (lowK of the
% holomorphic parameter is kapW, through ThetaFiresProbe.parAt_fixed) with kap_plus and kap_minus of ArithAt; (2) is
% coef_eq of ArithAt (lamLdef = thetaWt + dShift, which is lamPrime_eq, Lemma is-onto-step-exp-dominant) with the
% case split of wtU. The formulas for a^+ and a^- are aPl and aMi, equal to aPlusQ and -dRowQ by cast_aPlus and
% aM_cast. 0-based indices in the Lean.

\begin{lemma}\label{is-onto-step-kernel-span-real-fock}
Let $R$ be a realization of a definite-side datum $D$ with a face $F$ that satisfies the core conditions at
$F$, with the notation of Definition~\ref{is-eigenvalue-lift-data}. Let $j$ be a lift, $\omega_0$ the Weil
datum of $R$ for $(V_j,W_j)$ at the base tuple $\mu_0$ (Definition~\ref{is-realisation}(6)), $W_c$ the
coefficient space of $V_j$ with its action $\rho$ of $U(V_j)(\mathbf A_{F_h})$, and $\mathrm{io}_P(j)$ the map of
$R$ from the tensors of lowest $K$-type vectors of the model of $W_j$ at $\sigma_0$ and the vectors of $W_c$ to
the polynomials (Definition~\ref{is-realisation}(7)). Write $f\otimes P$ for $f\otimes P$ formed with the Cartan
involutions $1$ of $V_j$ and those of the model of $W_j$ at $\sigma_0$ and the bases of $R$
(Lemma~\ref{is-eigenvalue-conjugate-theta-facts}). Let $Q$ be a finite set of places of $E_h$, $\xi_P$ a
character of $k_P^\times$ for every finite place $P$, and $\varphi$ a level vector for these Cartan
involutions, the lattices and adelic matrices of $R$ for $(V_j,W_j)$, $P_g$, $Q$ and the $\xi_P$, whose vacuum
is the vacuum vector of $R$ for $(V_j,W_j)$. Then
\[
 \omega_0(u,1)\big(\varphi\otimes\mathrm{io}_P(j)(t\otimes w)\big)=\varphi\otimes\mathrm{io}_P(j)\big(t\otimes\rho(u)w\big)
\]
for every $u\in U(V_j)(\mathbf A_{F_h})$ with finite components $1$, every tensor $t$ of lowest $K$-type vectors
and every $w\in W_c$.
\end{lemma}
\leanref{Development/AbelianVarietyTheory/Proofs/DivisorProofs/RealPointsActFockVectorsAnyDatumStatement.lean}{GlobalTheory.T.realPoints_act_fock_vectors_any_datum}
\leanref{Development/AbelianVarietyTheory/Proofs/DivisorProofs/RealPointsActFockVectorsAnyDatum.lean}{GlobalTheory.FockAnyDatum.om_ioP_of_datum}
\begin{proof}
Let $M$ be the model of $W_j$ at $\sigma_0$, with frames $f_\tau$, signature matrices $J_\tau$, Cartan
involutions $s_\tau=f_\tau^{-1}J_\tau f_\tau$ and groups $K_\tau$ of unitary matrices commuting with $J_\tau$,
and let $B_\tau$ be the bases of $V_j$ of $R$, so that $f\otimes P$ is $f\otimes_{B,f^{-1}}P$ in the notation of
Lemma~\ref{is-onto-step-kernel-span-real-compact}. Let $v_0$ be the test vector of $R$ for $(V_j,W_j)$ at
$\mu_0$ (Definition~\ref{is-realisation}(6)), and $r=\operatorname{rk}W_j$.

\emph{The setting.} By the core conditions on the models, $\theta$ is an archimedean family whose member at
each $\tau$ extends $\tau$; the bases $B_\tau$ are orthonormal, $B_\tau^*G_{V_j}(\tau)B_\tau=1$; and the signs of
$M$ at $\tau$ are $1$ at the first $p_\tau$ indices and $-1$ at the other $q_\tau$, where $(p_\tau,q_\tau)$ is the
signature of $W_j$ at $\tau$, with $p_\tau+q_\tau=r$. By the core conditions on the theta lifts, $\omega_0$ is a
Weil datum for $(V_j,W_j)$ with the characters $(\mu^0_j\xi_{\mu_0},\mu^0_j\xi_{\mu_0}\tilde\eta^{-1})$ at the
level $j+1$, $\tilde\eta$ the Hecke character of $R$, and $v_0$ is a level vector for $P_g$, $Q_{\mu_0}$ and the
residue characters of $\xi_{\mu_0}$. By Lemma~\ref{is-eigenvalue-conjugate-theta-facts}(3), $\varphi$ and $v_0$
are theta functions annihilated by the annihilation operators for $1$ and the $s_\tau$.

\emph{The exponents.} By the core conditions on the level characters, $\mu^0_j$ has archimedean exponent
$-e_j(\tau)$ at $\tau$, where $e_j(\tau)$ is the $j$-th number $e$ of the Harish-Chandra parameter of $D$ at
$\tau$ (Definition~\ref{th-hc-parameter}); $\tilde\eta$ has archimedean exponent $\frac12$; and $\xi_{\mu_0}$,
of finite order, has archimedean exponent $0$. So the characters of $\omega_0$ have the archimedean exponents
$-e_j(\tau)$ and $-e_j(\tau)-\frac12$ at $\tau$. With $n=\operatorname{rk}V_j$, the numbers
$x_{1,\tau}=\frac n2-e_j(\tau)$, $x_{2,\tau}=-\frac n2-e_j(\tau)$ and
$x_{3,\tau}=-\big(\frac{p_\tau-q_\tau}2-e_j(\tau)-\frac12\big)$ are integers, by
Lemma~\ref{is-onto-step-kernel-span-real-fock-weights} applied to $M$, which has the parameter of the models of $W_j$
at $\sigma_0$ and the signs above.

\emph{The substitution.} Let $u$ be as in the statement, and $m_\tau=B_\tau^{-1}\theta_\tau(u)B_\tau$.
Lemma~\ref{is-onto-step-kernel-span-real-compact}(2), applied to $\omega_0$ with $g=u$ and $h=1$, for which
$k_\tau=1\in K_\tau$, shows that every $m_\tau$ is unitary and that, for $f=\varphi$ and for $f=v_0$ and every
polynomial $P$,
\[
 \omega_0(u,1)(f\otimes P)=c\cdot f\otimes s(P),\qquad c=\prod_\tau\det(m_\tau)^{x_{3,\tau}},\quad s=s_{1,m}.
\]
The number $c$ and the substitution $s$ depend on $u$, and not on $f$.

\emph{The test vector.} Put $P_0=\mathrm{io}_P(j)(t\otimes w)$ and $P_1=\mathrm{io}_P(j)(t\otimes\rho(u)w)$. By
the core conditions on the theta lifts, $\mathrm{io}_P(j)$ intertwines $\omega_0$ with $\rho$ on the test vector:
$\omega_0(u,1)(v_0\otimes P_0)=v_0\otimes P_1$. With the substitution, $v_0\otimes(c\,s(P_0))=v_0\otimes P_1$. The
vector $v_0$ is not $0$, by the argument in the proof of
Lemma~\ref{is-eigenvalue-conjugate-frame-extremal-line}(2): the lattices of the frames of $V_j$ and $W_j$ are
self-dual by the core conditions on the level characters, $r\le n$, and the level function of $v_0$ takes the
value $1$ at the point chosen there. The bases $B_\tau$ are orthonormal and adapted to the Cartan involution
$1$, and the bases $f_\tau^{-1}$ are orthonormal and adapted to $s_\tau$, as in that proof. So
Lemma~\ref{is-eigenvalue-conjugate-frame-extremal-fock}(3) makes $P\mapsto v_0\otimes P$ injective, and
$c\,s(P_0)=P_1$.

\emph{Conclusion.} For $f=\varphi$,
$\omega_0(u,1)(\varphi\otimes P_0)=c\cdot\varphi\otimes s(P_0)=\varphi\otimes(c\,s(P_0))=\varphi\otimes P_1$.
\end{proof}
% Lean: FockAnyDatum.om_ioP_of_datum: hthphi and hthv (JointHarm.isThetaFunO_of_isThetaAt), hvphi and hvv
% (DatumFock.isFockVacAt_thPVec), hBS (JointHarm.onb_eq_one from hR.models.BV), x1, x2, x3 from
% T.lowestKType_coefficientWeight_eq_minorWeights (Lemma is-onto-step-kernel-span-real-fock-weights) with the exponents of
% T.thetaLiftInfinityTypes_atRealPlace (PairInfTypes.infTypesAt: chiV and chiW), act =
% T.weil_maximalCompact_fock_substitution_character (Lemma is-onto-step-kernel-span-real-compact(2)) at
% (u, 1), for which the frame conjugate of 1 is 1 (e); hcV (the Cartan involution 1 of V_j), hne
% (AdTh.VecNe.thPVec_ne_zero with the self-dual lattices of hR.lev.betaV and betaW), hinj
% (FockVecInj.pureTensorO_injective, Lemma is-eigenvalue-conjugate-frame-extremal-fock(3)), and key, the
% first clause of IsIoW for the core condition hR.thP.ioP i mu0 at the test vector R.vP i mu0 (ptP).

\begin{lemma}\label{is-onto-step-kernel-span-real-weil}
Let $R$ be a realization of a definite-side datum with host field $E_h/F_h$ and a face $F$ that satisfies the core
conditions at $F$ (Definition~\ref{is-eigenvalue-lift-data}). Let $j$ be a lift, $\nu$ a Hecke character of $E_h$,
$Q$ a finite set of finite places of $E_h$, $(\chi,(\xi_P),\omega,\varphi)$ a twisted kernel of $R$ at $j$ for
$(\nu,Q)$ (Definition~\ref{is-onto-step-kernel-span-twisted}), and $\omega_0$ the Weil datum of $R$ for
$(V_j,W_j)$ at the base tuple $\mu_0$ (Definition~\ref{is-realisation}(6)). Of the conditions on the twisted kernel,
only (1), (2) and (9) are used in the proof. Then $\omega(u,1)f=\omega_0(u,1)f$ for every $u\in
U(V_j)(\mathbf A_{F_h})$ with finite components $1$ and every theta function $f$ for $(V_j,W_j)$.
\end{lemma}
\leanref{Development/AbelianVarietyTheory/Proofs/DivisorProofs/WeilActionKernelEqBaseOnRealUnitaryPointsStatement.lean}{GlobalTheory.T.weilAction_kernel_eq_base_on_realUnitaryPoints}
\leanref{Development/AbelianVarietyTheory/Proofs/DivisorProofs/WeilActionKernelEqBaseOnRealUnitaryPoints.lean}{GlobalTheory.RealTwistKer.weil_kernel_real}
\begin{proof}
Let $\tilde\eta$ be the Hecke character of $R$. By (9), the Weil datum $\omega$ has the characters
$(\chi_S,\chi_T)=(\mu^0_j\chi,\mu^0_j\chi\tilde\eta^{-1})$, and by the core conditions on the theta lifts
$\omega_0$ has the characters $(\chi^0_S,\chi^0_T)=(\mu^0_j\xi_{\mu_0},\mu^0_j\xi_{\mu_0}\tilde\eta^{-1})$, at the
level $j+1$.

\emph{Two characters.} By (2) there is $n\geq1$ with $\nu^n=1$. Then $\chi^n$ is a continuous unitary character of
$\mathbf A_{E_h}^\times$, trivial on $E_h^\times$, and by (1) it takes the value $\nu_P(k)^n=1$ at the ideles $x_P(k)$
of Lemma~\ref{is-onto-step-kernel-unique-character}; by that lemma, applied to $\chi^n$ and the trivial character,
$\chi^n=1$. By the core conditions on the level characters $\xi_{\mu_0}=1$, and the unitary character
$\xi_{\mu_0}$ of $R$, which is equal to it, is $1$ by the same lemma.

\emph{A determinant root.} Let $u$ be as in the statement and $d=\det u$; its finite components are $1$. By
Lemma~\ref{is-eigenvalue-conjugate-frame-root-values}(1), $u$ has a determinant root $x_0$, so
$x_0=d\,\bar x_0$, and $|\theta_\tau(d)|=1$ for every $\tau\in\mathcal T$. Choose $s_\tau\in\mathbf C$ with
$s_\tau^{2n}=\theta_\tau(d)$, so that $|s_\tau|=1$, and let $y\in\mathbf A_{E_h}^\times$ have finite components
$1$ and $\theta_\tau(y)=s_\tau$ for every $\tau$. Then $y^n=d\,\bar y^n$: both sides have finite components $1$,
and $s_\tau^{2n}\bar s_\tau^n=s_\tau^n$ as $s_\tau\bar s_\tau=1$. So $y^n$ is a determinant root of $u$; and $1$ is
a determinant root of $1\in U(W_j)(\mathbf A_{F_h})$.

\emph{The twist.} Put $\upsilon=\chi_S(\chi^0_S)^{-1}=\chi_T(\chi^0_T)^{-1}=\chi\xi_{\mu_0}^{-1}$. It is continuous,
trivial on $E_h^\times$, and trivial on $\mathbf A_{F_h}^\times$, since $\chi_S$ and $\chi^0_S$ both restrict there
to $\eta^{\operatorname{rk}V_j}$ (Definition~\ref{th-weil-datum}(3)). By
Lemma~\ref{asm-dependence-test-assembly-first-layer-ratio-inputs-weil-twist}, with $\nu_S=\nu_T=\upsilon$, there is
a Weil datum $\omega'$ for $(V_j,W_j)$ with the characters $(\chi_S,\chi_T)$ such that
$\omega'(g,h)=\upsilon(x)^{-1}\upsilon(y')\,\omega_0(g,h)$ whenever $x$ and $y'$ are determinant roots of $g$ and
$h$. By Lemma~\ref{is-eigenvalue-conjugate-theta-facts}(2), $\omega$ and $\omega'$ agree on theta functions. So,
with the determinant roots $y^n$ of $u$ and $1$ of $1$,
\[
 \omega(u,1)f=\omega'(u,1)f=\upsilon(y^n)^{-1}\omega_0(u,1)f
\]
for every theta function $f$, and $\upsilon(y^n)=\chi^n(y)\,\xi_{\mu_0}(y^n)^{-1}=1$.
\end{proof}
% Lean: T.weilAction_kernel_eq_base_on_realUnitaryPoints = RealTwistKer.weil_kernel_real, for any KerDat R i xi
% chi xiR Q omega phi (Definition is-onto-step-kernel-span-twisted), of which it uses only chr, fin and weil. hc: ConjTwist.BaseInhab.chrEq_pow and
% ChrEqUnique.eq_one_of_chrEq for chi^n = 1; e2: TwistFactorCompare.eq_one_of_chrEq_one with hR.lev.xi_base and
% hR.lev.xi; exists_isDetRoot_pow (Member.isDetRoot_laws for x0, IsAlgClosed.exists_pow_nat_eq for s, and
% exists_idele_all for y); TwistFactorProof.weil_ratio (WeilChange.exists_weilDatum_twist, Lemma
% asm-dependence-test-assembly-first-layer-ratio-inputs-weil-twist, with isHeckeChar_mul_inv and resEta_parity, and
% T.weilDatum_unique_on_thetaFunctions, Lemma is-eigenvalue-conjugate-theta-facts(2)) with isDetRoot_one.

\begin{lemma}\label{is-onto-step-kernel-span-real}
Let $R$ be a realization of a definite-side datum with host field $E_h/F_h$ and a face $F$ that satisfies the
core conditions at $F$ (Definition~\ref{is-eigenvalue-lift-data}). Let $j$ be a lift, $\nu$ a Hecke character of
$E_h$, $Q$ a finite set of finite places of $E_h$, and $(\chi,(\xi_P),\omega,\varphi)$ a twisted kernel of $R$ at
$j$ for $(\nu,Q)$ (Definition~\ref{is-onto-step-kernel-span-twisted}). Of the conditions on the twisted kernel,
only (1), (2), (9) and (10) are used in the proof. Let $W_c$ be the coefficient space of $V_j$ with its action
$\rho$ of $U(V_j)(\mathbf A_{F_h})$ (Definition~\ref{is-ties-lattice-systems}), and for $h\in
U(V_j)(\mathbf A_{F_h})$ and $w\in W_c$ let $\vartheta(h,w)$ be the theta series of $\omega$ and $\varphi$ at $h$
for the polynomial of $w$, a function of the tensors of lowest $K$-type vectors and of $g\in
U(W_j)(\mathbf A_{F_h})$. Let $u\in U(V_j)(\mathbf A_{F_h})$ have finite components $1$. Then
$\vartheta(hu,w)=\vartheta(h,\rho(u)w)$ for all $h$ and $w$.
\end{lemma}
\leanref{Development/AbelianVarietyTheory/Proofs/DivisorProofs/KernelThetaSeriesSpanOverClassRepresentatives.lean}{GlobalTheory.OneKernelSpan.series_real}
\leanref{Development/AbelianVarietyTheory/Proofs/DivisorProofs/WeilActionKernelEqBaseOnRealUnitaryPointsStatement.lean}{GlobalTheory.T.weilAction_kernel_eq_base_on_realUnitaryPoints}
\leanref{Development/AbelianVarietyTheory/Proofs/DivisorProofs/RealPointsActFockVectorsAnyDatumStatement.lean}{GlobalTheory.T.realPoints_act_fock_vectors_any_datum}
\begin{proof}
Let $\omega_0$ be the Weil datum of $R$ for $(V_j,W_j)$ at the tuple $\mu_0$ (Definition~\ref{is-realisation}(6)).
The proof rests on two facts about $u$:
\begin{enumerate}
\item[(a)] $\omega(u,1)f=\omega_0(u,1)f$ for every theta function $f$;
\item[(b)] $\omega_0(u,1)(\varphi\otimes P)=\varphi\otimes P'$, where $P$ is the polynomial of $w$ and a tensor
  $t$ of lowest $K$-type vectors and $P'$ that of $\rho(u)w$ and $t$.
\end{enumerate}
Fact (a) is Lemma~\ref{is-onto-step-kernel-span-real-weil} for the twisted kernel, which uses its conditions (1),
(2) and (9). Fact (b) is Lemma~\ref{is-onto-step-kernel-span-real-fock}, with $P=\mathrm{io}_P(j)(t\otimes w)$ and
$P'=\mathrm{io}_P(j)(t\otimes\rho(u)w)$, as $\varphi$ is a level vector for $P_g$, $Q$ and the $\xi_P$ whose
vacuum is the vacuum vector of $R$ for $(V_j,W_j)$, by condition (10).
% Lean: (a) is T.weilAction_kernel_eq_base_on_realUnitaryPoints = RealTwistKer.weil_kernel_real (Lemma
% is-onto-step-kernel-span-real-weil) for hK; (b) is T.realPoints_act_fock_vectors_any_datum =
% FockAnyDatum.om_ioP_of_datum (Lemma is-onto-step-kernel-span-real-fock), applied to hK.vec.

Let $h$ and $w$ be given, $t$ a tensor of lowest $K$-type vectors with $P$ and $P'$ as in (b), and $g\in
U(W_j)(\mathbf A_{F_h})$. As $\omega$ is a homomorphism and $(hu,g)=(h,g)(u,1)$,
\[
 \vartheta(hu,w)(t)(g)=\big(\omega(h,g)\,\omega(u,1)(\varphi\otimes P)\big)(0).
\]
By conditions (9) and (10) and Lemma~\ref{is-eigenvalue-conjugate-theta-facts}(3) and (4), $\varphi\otimes P$ is a
theta function, so by (a) and (b), $\omega(u,1)(\varphi\otimes P)=\varphi\otimes P'$, and the right side is
$\vartheta(h,\rho(u)w)(t)(g)$.
% Lean: series_real, for any KerDat R i xi chi xiR Q omega phi: Prod.ext for (h u, g) = (h, g) (u, 1), map_mul,
% (a) with pt_theta_of_kerDat (from the fields vec and weil) for the theta function, and (b).
\end{proof}

\begin{lemma}\label{is-onto-step-kernel-span}
Let $R$ be a realization of a definite-side datum with host field $E_h/F_h$, group $G$ and a face $F$,
with subgroup $P_0$ and character $\psi_0$, let $j$ be a lift and $v$ a twisted place. Suppose that:
\begin{enumerate}
\item $R$ satisfies the core conditions at $F$ (Definition~\ref{is-eigenvalue-lift-data});
\item for every twisted place $v'$, with residue characteristic $\ell$ and character $\delta_{v'}$
  (Definition~\ref{is-onto-step-runs}(1)), there are a Hecke character $\eta$ of $E_h$ and a finite
  set $Q$ of finite places of $E_h$ with $\xi_{\mu_{v'}}=\delta_{v'}\eta$ such that: $\eta^{\ell^a}=1$
  for some $a\geq0$; for every place $u$ of $F_h$ the product of the components of $\eta$ at the places
  over $u$ is trivial on $F_h^\times$; $\eta$ is tamely ramified at every finite place, and unramified
  at a finite place if and only if the place is not in $Q$; its components at the places over $\ell$
  are trivial; and no place of $Q$ is its own complex conjugate;
\item the values of the character $\delta_v$ lie in $\sigma_D(L)$, for the embedding $\sigma_D$ of $L$
  into $\mathbf C$ of the definite-side datum.
\end{enumerate}
Let $U\subset G$ be a subgroup, $\psi\colon U\to L^\times$ a character and $\lambda\in\mathbf C$. For a
base theta kernel $\kappa$ at $(j,v)$ (Definition~\ref{is-onto-step-kernels}) define
$\Theta^{U,\psi,\lambda}_\kappa(h)$ and $\Theta^{U,\psi,\lambda}_\kappa$ as $\Theta_\kappa(h)$ and
$\Theta_\kappa$ in Definition~\ref{is-onto-step-kernels}(2), with the trace
$\lambda\sum_{g\in U}\sigma_D(\psi(g))^{-1}g\vartheta$ in place of
$\sum_{g\in P_0}\sigma_D(\psi_0(g))g\vartheta$, so that
$\Theta_\kappa(h)=\Theta^{P_0,\psi_0^{-1},1}_\kappa(h)$. Let $\kappa$ be a base theta kernel at $(j,v)$
and $(h_c)_{c\in C}$ a family that meets the classes of the level $J_v$. Then
$\Theta^{U,\psi,\lambda}_\kappa$ and $\bigcup_c\Theta^{U,\psi,\lambda}_\kappa(h_c)$ span the same
$\mathcal O_L$-module. In particular, $\Theta_\kappa$ and $\bigcup_c\Theta_\kappa(h_c)$ span the same
$\mathcal O_L$-module.
\end{lemma}
\leanref{Development/AbelianVarietyTheory/Proofs/DivisorProofs/KernelThetaSeriesSpanOverClassRepresentativesStatement.lean}{GlobalTheory.T.kernel_thetaSeries_span_over_class_representatives}
\leanref{Development/AbelianVarietyTheory/Proofs/FoundationProofs/KernelSpanOverClasses.lean}{GlobalTheory.TraceClassRank.KernelSpanOverClasses}
\begin{proof}
Write $\kappa=(\omega,\varphi)$. By Lemma~\ref{is-onto-step-kernel-span-twisted-base}, which uses hypotheses (1)
and (2), there are $\chi$ and $(\xi_P)$ such that $(\chi,(\xi_P),\omega,\varphi)$ is a twisted kernel of $R$ at $j$
for $(\delta_v,Q_v)$ (Definition~\ref{is-onto-step-kernel-span-twisted}); its level of the data at $Q_v$ is $J_v$.
Write $\vartheta(h,w)$ for the theta series of $\omega$ and $\varphi$ as in
Lemma~\ref{is-onto-step-kernel-span-level}, $\Gamma=U(V_j)(F_h)$, $G_f$ for the finite adelic group of
$U(V_j)$, read as the elements of $U(V_j)(\mathbf A_{F_h})$ with archimedean component $1$, and
$\operatorname{tr}$ for the trace $\lambda\sum_{g\in U}\sigma_D(\psi(g))^{-1}g$, a $\mathbf C$-linear
map. Each $h_c$ lies in $G_f$, so $\Theta^{U,\psi,\lambda}_\kappa(h_c)\subset\Theta^{U,\psi,\lambda}_\kappa$,
and the span of the union lies in the span of $\Theta^{U,\psi,\lambda}_\kappa$.

Conversely, let $g\in G_f$ and $x\in\Theta^{U,\psi,\lambda}_\kappa(g)$: there is $w$ integral for the
lattice of $g$ with $\Phi_j(x)=\operatorname{tr}\vartheta(g,w)$. As $(h_c)$ meets the classes of $J_v$,
$g=\gamma h_ck$ in $G_f$ for some $c$, $\gamma\in\Gamma$ and $k\in J_v$. Read $k$ in $G_f$ and put
$u=k^{-1}h_c^{-1}\gamma^{-1}g$. Its finite components are $1$, and $g=\gamma h_cku$.

\emph{The archimedean part.} By Lemma~\ref{is-onto-step-kernel-span-real}, for this twisted kernel,
$\vartheta(g,w)=\vartheta(\gamma h_ck,\rho(u)w)$. The core conditions on the models make $\rho$ trivial
on the elements with archimedean component $1$, so
$\rho(u)=\rho(k)^{-1}\rho(h_c)^{-1}\rho(\gamma)^{-1}\rho(g)=\rho(\gamma)^{-1}$.

\emph{The rational part.} Let $w'\in W_c$ and $P$ the polynomial of $w'$ for a tensor of lowest
$K$-type vectors. By Lemma~\ref{is-eigenvalue-conjugate-theta-facts}(3) and (4), $\varphi\otimes P$ is a
theta function, and so is $\omega(h',g')(\varphi\otimes P)$ for every $(h',g')$, by
Definition~\ref{th-weil-datum}(3). That definition also gives $(\omega(\gamma,1)f)(0)=f(0)$ for every
theta function $f$, so $\theta_\omega(\gamma h',g';\varphi\otimes P)=\theta_\omega(h',g';\varphi\otimes
P)$. Hence $\vartheta(\gamma h_ck,w')=\vartheta(h_ck,w')$.

\emph{The level.} By Lemma~\ref{is-onto-step-kernel-span-level}, for this twisted kernel, whose character
$\delta_v$ has its values in $\sigma_D(L)$ by hypothesis (3), there is a root of unity $z\in L$ with
$\vartheta(h_ck,w')=\sigma_D(z)\vartheta(h_c,w')$.

Together, with $w'=\rho(\gamma)^{-1}w$,
\[
 \vartheta(g,w)=\sigma_D(z)\,\vartheta\big(h_c,\rho(\gamma)^{-1}w\big),
\]
and $\rho(\gamma)^{-1}w$ is integral for the lattice of $h_c$ by
Lemma~\ref{is-onto-step-kernel-span-move}. The map $\Phi_j$ is semilinear for $\sigma_0=\sigma_D$ and
$\operatorname{tr}$ is linear, so
$\Phi_j(z^{-1}x)=\sigma_D(z)^{-1}\operatorname{tr}\vartheta(g,w)=\operatorname{tr}\vartheta(h_c,\rho(\gamma)^{-1}w)$,
and $z^{-1}x\in\Theta^{U,\psi,\lambda}_\kappa(h_c)$. A root of unity is integral, so $z\in\mathcal O_L$,
and $x=z\cdot z^{-1}x$ lies in the $\mathcal O_L$-span of $\bigcup_c\Theta^{U,\psi,\lambda}_\kappa(h_c)$.
% Lean: T.kernel_thetaSeries_span_over_class_representatives = OneKernelSpan.kernelSpanOverClasses_of.
% The Lean states it, as the book does, for every subgroup U, character psi and phase; it is used with
% (P0, psi0^-1, 1). The proof: the KerDat of the kernel (T.baseTuple_thetaKernel_tameCharacter), the
% decomposition hh.2 (fin g) with TraceIncl.KernLaws.exists_finite_lift for k, u = k^-1 h_c^-1 gamma^-1 g
% in GinfC, series_real, series_rat (the first clause of thetaKO_rat_left_and_eigen_right, through the
% om_aut of the Weil datum), series_level, hrho from (hR.models.cV i).1 (rho is 1 on GfC), inCoefLat_move,
% adForm_smul, and z integral (IsIntegral.of_pow). Size: about 25,000 lines of Lean below it (text count),
% nearly all below Lemmas is-onto-step-kernel-span-level and is-onto-step-kernel-span-real.
\end{proof}

\begin{lemma}\label{is-onto-step-span}
Let $R$, $F$, $j$ and $v$ be as in Lemma~\ref{is-onto-step-kernel-span}. Let $\kappa$ be a base theta
kernel at $(j,v)$ and $(h_c)_{c\in C}$ a family that meets the classes of the level $J_v$. Then the
base theta lattice $\Theta^{\mathrm b}_v(j)$ is the $\mathcal O_L$-span of $\bigcup_c\Theta_\kappa(h_c)$.
\end{lemma}
\leanref{Development/AbelianVarietyTheory/Proofs/FoundationProofs/KernelSpanOverClasses.lean}{GlobalTheory.TraceClassRank.SpanOverClasses}
\leanref{Development/AbelianVarietyTheory/Proofs/FoundationProofs/KernelSpanOverClasses.lean}{GlobalTheory.TraceClassRank.spanOverClasses_of}
\begin{proof}
By Lemma~\ref{is-onto-step-kernel-span}, the right side is the $\mathcal O_L$-span of $\Theta_\kappa$.
The set $\Theta_{\kappa'}(h)$ depends on the kernel $\kappa'$ only through its theta series, so by
Lemma~\ref{is-onto-step-kernel-unique} $\Theta_{\kappa'}=\Theta_\kappa$ for every base theta kernel
$\kappa'$ at $(j,v)$. Hence the union of the $\Theta_{\kappa'}$ over the base theta kernels is
$\Theta_\kappa$, and its span, the base theta lattice, is the span of $\Theta_\kappa$.
\end{proof}

\begin{lemma}\label{is-onto-step-export}
Assume the rapid decay of cusp forms on unitary groups, global Artin reciprocity, the existence of ray
class fields and the covering of the discrete series by representations with prescribed lowest
$K$-type (Theorem~\ref{lit-bundle-unitary}(1), (4) and (5)), and the Casselman-Wallach theorem
(Theorem~\ref{lit-casselman-wallach}). Let $\mathcal H$ be the class of standard families.
Eventually along $\mathcal H$, every family has a scale beyond which every point $p$ that satisfies
the point conditions and the further conditions has the following property. Let $R$ be its
realization, with the scalars of its characterization, and $F$ its face. For every twisted place $v$
and lift $j$, let $(i,b,Q,k',\psi)$ be the twist choice of the member at $v$
(Definition~\ref{asm-twist-condition}), with its discrete valuation ring $\mathcal O'$, its fraction
field $F'$, its embedding $\sigma_v$ of $L$ in $F'$ and the set $\Sigma$ of theta data that it gives
for $(i,j,b)$, and put $\varphi=\psi^{(b)}\circ\mathrm{art}$. Then there is a $\sigma_v$-semilinear map
$\mathrm{rd}$ from the coordinates $L^{n_W(j)}$ of $FW_j$ to $F'\otimes_{\mathcal O_i}F_{i,j,b}$ that
carries some basis to a linearly independent family, satisfies
$\mathrm{rd}\circ e_{\mathrm{good}}=e_{i,b}\circ\mathrm{rd}$ for the good operator of
Definition~\ref{is-ties-class-sums}(2), and carries the base theta lattice $\Theta^{\mathrm b}_v(j)$
and the theta lattice $\Xi_{j,\mu_v}$ to sets that span $\mathcal O'\otimes\Theta(\Sigma)$ and
$\mathcal O'\Theta_{\varphi^{(b)},Q}(\Sigma)$ over $\mathcal O'$.
\end{lemma}
\leanref{Development/AutomorphicArithmetic/ThetaCorrespondence/ThetaPairLifts/ThirdPlaceSeries/ExportStOfReadingStatement.lean}{GlobalTheory.T.exportSt_of_reciprocity_and_casselman_wallach}
\leanref{Development/AutomorphicArithmetic/ThetaCorrespondence/ThetaPairLifts/ThirdPlaceSeries/ExportSt.lean}{GlobalTheory.BigOn.ExportedStatement}
\begin{proof}
\openstep{The export of the base theta lattice and the theta lattice of the tuple to the local model
of the twist choice is not written out.}
% Lean: T.exportSt_of_reciprocity_and_casselman_wallach = exportSt_of_reading_alt =
% T.exportSt_of_discrete_series_cover = Exp.exportSt_of_hypotheses_alt, which feeds
% Exp.InductionConditions.exportSt with thirteen inputs: T.bigOn_canEx, T.bigOn_canProj,
% T.twist_admissible', T.at_statement_of_low_weight_twAdm (with T.lowAt), T.tieRat_of_twAdm_and_at,
% T.reciprocity_side_of_twAdm, T.adaptEx', T.conjBase_of_ray_class_field_and_seriesRat,
% T.conjTw_of_discrete_series_and_casselman_wallach (with T.twist_character_of_artin'),
% T.baseLe_of_cuspRapidDecay_artin_rayClassField, T.baseGe_of_cuspRapidDecay_artin_rayClassField_seriesRat,
% T.twist_upper_bound_of_prints and T.twist_lower_bound_of_prints. Each could be its own lemma. The
% rationality of theta series SerStd.SeriesRatS, a hypothesis of the Lean statement, comes from Artin
% reciprocity by seriesRatS_and_serInt_of_globalArtin (T.thetaSeriesRationalOnGoodCovers_of_untwisted
% with T.untwisted_theta_rational_on_good_covers); it belongs to this step. Size: about 789,000 lines
% of Lean in 46,000 theorems below the two (text count, stopping at statements the book cites); the
% export alone has 662,000, of which about 184,000 lie below no other open step of
% Lemma is-onto-estimates.
\end{proof}

\subsection{Theta modules and the geometric estimates}\label{subsection-is-geometric-estimates}

\begin{lemma}\label{is-onto-step-numerics}
Let $\mathcal H$ be the class of standard families, and $x_1$ a real number. Eventually along
$\mathcal H$, every family has a scale $x_0$ such that every member at a scale $x\geq x_0$, with
definite-side datum of ranks $k,m,n$, $f=2n$, rank $r$ and column rank $n_V$, size data with
quantities $b,b'$ and exponent $\varepsilon_z$ of its scale, and parameter choices with constants
$C_{nW}$ and $C_P$, satisfies: the size conditions of its size data; $C_{nW}=10$; $\log g_i\le9\log
b+1$ for $i=0,1$; $1000\le C_P$ and $C_P^{20}\le fr\log b$; $x_1\le x$; $3\le x$; $kb'\le b$;
$3\le x^{\varepsilon_z/5}$; and $8fr+4\le n_V$.
\end{lemma}
\leanref{Development/AbelianVarietyTheory/Proofs/DivisorProofs/RunNumericalConditionsHoldAlongFamilyStatement.lean}{GlobalTheory.T.runNumericalConditions_hold_along_family}
\leanref{Development/AbelianVarietyTheory/Definitions/Divisors/IsDeltaDelAtOwnAt.lean}{GlobalTheory.BigOn.PageNum}
\begin{proof}
By the proof of Lemma~\ref{is-onto-page-numerics}, there is $\varepsilon_1>0$ such that for every
exponent $\varepsilon\le\varepsilon_1$ there is a largeness threshold $\kappa$ for which every standard
family of exponent $\varepsilon$ that is large for $(\kappa,C_H)$ has a scale $x_a$ beyond which every
member satisfies the size rows. For such $\varepsilon$ let $\kappa'$ be the threshold with parts
$\max(\kappa_{\mathrm{pre}},1000)$, $\max(\kappa_{\mathrm{mid}},0)$ and $\max(\kappa_{\mathrm{last}},0)$.
Let $\mathcal F$ be a
standard family of exponent $\varepsilon$ that is large for $(\kappa',C_H)$. Then it is large for
$(\kappa,C_H)$. The list $c_{\mathrm{pre}}(\mathcal F)$ consists of the constant $C_P$ that the
parameter choices of all members share, so $C_P\geq1000$. Put $x_0=\max(x_a,10^{60},C_P^{20},x_1)$,
and let a member at a scale $x\geq x_0$ above the threshold of $\mathcal F$ be given.

The size conditions hold for every member above the threshold, the scale of its size data is $x$
(Definition~\ref{pd-family}), and $C_{nW}=10$ is one of the fixed constants of every family. By the
page conditions of the duty data (Definition~\ref{asm-duty}(2)), $g_i\le e^{9\log b}$, so $\log
g_i\le9\log b\le9\log b+1$. The size rows give $kb'\le b$ and $11\le x^{\varepsilon_z/5}$, so $3\le
x^{\varepsilon_z/5}$. Since $\varepsilon_z=\varepsilon/J\le\varepsilon<1$ and $x>1$, also
$x\geq x^{\varepsilon_z/5}\geq3$. By the choice of $x_0$, $x_1\le x$.

The size conditions give $x^{9/2+10\varepsilon-\varepsilon/100}\le k$, so $x\le k\le k+m-1=r$. They
also give $x^{\varepsilon/100}\log b\le f\log k$, whose left side is positive, so $f\geq1$. The
parameter choices give $(Q_Sqd)^{36}t^{150}\le b^{1/2}$, and the ranks satisfy $r+5\le t$ and
$r\geq6$. So $\log2\le\log r\le\log t\le(\log b)/300$, and $\log b\geq300\log2>4$. Hence
\[
 C_P^{20}\le x\le r\le fr\le fr\log b.
\]
Finally, the parameter choices write $n_V=r+t$ as $n_P+T_On_W$ with $n_P=8fr+4$, so
$8fr+4\le n_V$.
% Lean: OntoNum.pageNum_along, from T.indexBoundParameterRequirements_eventually_family (the size rows,
% with the threshold K.sup <1000, 0, 0, 0>), FinalRows.rows_of_pageFix (PageFix and 1000 <= CP give
% SizesAt, CnW = 10 from Fm.fixed, log g <= 9 L + 1, CP^20 <= zBase through scale_le_rank and
% hostSz_of_sizes (len: 4 <= L from Pm.log_r_le and six_le_r), at x >= max(10^60, CP^20)),
% IndexAssembly.three_le_scale_of_pow_bound and ParamsT.nV_ge (nV_eq: nV = TO nWn + 8 f r + 4).
\end{proof}

\begin{lemma}\label{is-onto-step-index-conditions}
There are constants $C_K$ and $x_1$ with the following property. Let $p$ be a point of a family whose
core ties hold, with a member whose definite-side datum, size data and parameter choices satisfy the
conditions of Lemma~\ref{is-onto-step-numerics} that involve $x_1$ (the size conditions, $C_{nW}=10$,
the bound on $\log g_i$, the bounds on $C_P$ and $x_1\le x$), and let $r$ be a pair of runs at $p$
pinned to the own places (Definition~\ref{is-onto-step-runs}). Then for every finite place $v$ and
every level $j$ of a lift, the run used at $v$ at level $j$, with the scalars of its own prime, the
weights of the parameter choices and the bound $\operatorname{rk}(M_j)\,fmb\,\delta_2$, satisfies the
hypotheses and the conclusion of the index-parameter conditions with the constants
$(C_K,20,C_P,e_P,C_P,e_P,7)$.
\end{lemma}
\leanref{Development/AbelianVarietyTheory/Proofs/DivisorProofs/RunsHoldOfTies.lean}{GlobalTheory.BigOn.runsHold_of_ties}
\leanref{Development/AbelianVarietyTheory/Definitions/Divisors/IsDeltaDelAtOwnAt.lean}{GlobalTheory.BigOn.RunsHold}
\begin{proof}
\openstep{The index-parameter conditions of the runs are not written out.}
% Lean: runsHold_of_ties applies T.exists_constants_indexBound_hypotheses_liftWeights (the constants
% C_K and x_1) to FinalJoin2.liftIndexParameterConditions_of_coreAndShortLaws, with the field
% conditions CoreTies.field1/field2 at the own prime, the canonical quotient data and halfEX of the
% core ties, the laws lvA of the run and OwnPin. Its own proof is short; the size is that of
% T.exists_constants_indexBound_hypotheses_liftWeights. Size: about 683,000 lines of Lean in 38,000
% theorems below it (text count, stopping at statements the book cites), nearly all shared with
% the other open steps of Lemma is-onto-estimates.
\end{proof}

\begin{definition}\label{is-onto-step-leaves-kappa}
Let $D$ be a definite-side datum with ranks $k,m,t$ and column rank $n_V=k+m-1+t$, the set $\mathcal T$
of its embeddings, and Harish-Chandra parameters with entries $u_i(\tau)$, $v_i(\tau)$, $b_j(\tau)$ and
$e_j(\tau)$ at $\tau\in\mathcal T$ (Definition~\ref{th-hc-parameter}). For $1\le l\le m$, $\tau\in\mathcal
T$, $1\le i\le k+l-1$ and $1\le c\le m-l$ put
\[
 a^+_{l,\tau,i}=\sigma_i(\tau)+e_l(\tau)-\frac{n_V+1-2i}2-\frac{(k+l-1)-(m-l)}2,\qquad
 a^-_{l,\tau,c}=b_{l+c}(\tau)-e_l(\tau)-\frac{t-1+2c}2,
\]
where $\sigma_i=v_i$ for $i\le k-1$ and $\sigma_i=-b_{i-k+1}$ for $k\le i\le k+m-1$; these are integers,
as $t$ is odd. Among the data of $D$ are integers $\kappa^{\mathrm{big}}_{l,\tau,i}$ and
$\kappa^{\mathrm{small}}_{l,\tau,c}$ for all levels $l$ and the same $\tau$, $i$ and $c$, its \emph{weight
blocks}. The datum \emph{has Harish-Chandra weight blocks} if $\kappa^{\mathrm{big}}_{l,\tau,i}=
a^+_{l,\tau,i}$ and $\kappa^{\mathrm{small}}_{l,\tau,c}=-a^-_{l,\tau,c}$ for all $1\le l\le m$, $\tau$, $i$
and $c$. In that case, if $n_V$ is even, the weight $\kappa'_l$ at a level $l$ with $1\le l\le m$
(Definition~\ref{is-onto-step-leaves-theta-module}(1)) has the entries
\[
 \kappa'^{+}_{l,\tau,i}=\max\big(a^+_{l,\tau,i},0\big),\qquad
 \kappa'^{-}_{l,\tau,c}=\max\big(a^-_{l,\tau,m-l+1-c},0\big),
\]
and the integer $n_V$; at the other levels $\kappa'_l$ has all entries $0$ and the integer $n_V$.
The entries $\kappa'^{\pm}_{l,\tau,i}$ do not increase with $i$: $a^+_{l,\tau,i}-a^+_{l,\tau,i+1}=
\sigma_i-\sigma_{i+1}-1\geq0$ and $a^-_{l,\tau,c+1}-a^-_{l,\tau,c}=b_{l+c+1}-b_{l+c}-1>0$, by the
inequalities between the Harish-Chandra parameters.
\end{definition}
\leanref{Formalism/LocalRingCardinality/DataKap.lean}{Formalism.BasicNotions.DV.KapDef}
\leanref{Formalism/LocalRingCardinality/DataKap.lean}{Formalism.BasicNotions.DV.aPlusQ}
\leanref{Formalism/LocalRingCardinality/DataKap.lean}{Formalism.BasicNotions.DV.aMinusQ}
\leanref{Formalism/SharpRestriction/Divisibility/AbelianOpenPart/KaSmallNonpos.lean}{Formalism.MainNotions.DIV.Cons.kaWt}
\leanref{Formalism/SharpRestriction/Divisibility/AbelianOpenPart/OEOsInjective.lean}{Formalism.MainNotions.DIV.Cons.kaR}
% Lean: the weight blocks are D.what2.kaBig j tau i and D.what2.kaSmall j tau c (Rows.Data); KapDef D is
% the condition kaBig = aPlusQ and kaSmall = -aMinusQ (0-indexed i and c: sigPar, eQ = e_j, and b at
% j + c); kaWt D hK hn j has aP = (kaBig j).toNat, aM tau c = (-kaSmall j tau (rev c)).toNat and n = nV
% for 1 <= j <= m, and WeightAlgebra.wscalar (all entries 0, n = nV) otherwise; the weight kappa'_l of
% the parameter choices is kaR Tp hK l = kaWt D hK (even_nV Tp) l. Integrality is aPlusQ_int and
% aMinusQ_int, monotonicity kaBig_anti and kaSmall_anti. The datum's weight blocks are among the
% further data of Definition th-definite-datum(9), which the book does not list.

\begin{definition}\label{is-onto-step-leaves-theta-module}
Let $s$ be a run at a point of a family that satisfies the laws (Definition~\ref{is-onto-step-runs}),
with definite-side datum of ranks $k,m,f$ and $r$, host field $E_h$, level primes $g_0,g_1$, primes $q$
and $d$ and integer $Q_S$ of the host, and elements $\pi_0,\pi_1\in\mathcal O_{E_h}$ of the error data
of the member, with $\pi_k\bar\pi_k=g_k$ (Definition~\ref{pd-family}(4)). Let $l$ be a level of a lift,
and use the notation of Definition~\ref{is-onto-step-index-bound-conclusion}: the coefficient field
$F$ of the core, into whose ring of integers $\mathcal O_{E_h}$ embeds, the spaces $S_l(w)$ with the
modules $M_l(w)$ of integral cusp forms, and the ring of forms $A_l$ with the maps $\iota_w$. Put
$g=g_0g_1$.
\begin{enumerate}
\item (Weights.) A weight $w$ at level $l$ has, at each of the $f$ embeddings $\tau$ of the datum,
  nonnegative integer entries $w^+_{\tau,1}\geq\dots\geq w^+_{\tau,k+l-1}$ and
  $w^-_{\tau,1}\geq\dots\geq w^-_{\tau,m-l}$, and an even integer $n_w\geq0$. Weights add entry by
  entry, $n$ included, and $cw$ is
  $w$ with every entry multiplied by the integer $c\geq0$. Put $D^\pm(w)=\sum_{\tau,i}w^\pm_{\tau,i}$,
  $\deg w=D^+(w)+D^-(w)$ and $c(w)=\lceil\deg w/2\rceil$, and let $d_Y$ be the dimension $\dim Y$ at
  level $l$.
\item (Products and integrality.) For $\phi\in S_l(w)$ and $\phi'\in S_l(w')$ there is
  $\phi''\in S_l(w+w')$ with $\iota_w(\phi)\iota_{w'}(\phi')=\iota_{w+w'}(\phi'')$. Each finite place
  $\lambda$ of $F$ has a set of points, and each point $y$ a function $v_{\lambda,y}$ on $A_l$ with
  values in $\mathbf R\cup\{+\infty\}$, such that
  $v_{\lambda,y}(\iota_w(\phi)\iota_{w'}(\phi'))=v_{\lambda,y}(\iota_w\phi)+v_{\lambda,y}(\iota_{w'}\phi')$,
  and $\phi\in M_l(w)$ if and only if $v_{\lambda,y}(\iota_w\phi)\geq0$ for every $\lambda$ and $y$.
\item (Theta series.) For each weight $w$ the run has two $\mathcal O_F$-submodules of $S_l(w)$: the
  module $\mathrm{Th}_l(w)$ of its theta series, and the module $\mathrm{Sh}_l(w)$ of its short theta
  series, the set of elements of $S_l(w)$ in the $\mathcal O_F$-span of the short series of weight $w$
  (item (7)). For $a\in\mathcal O_F$ put $\mathrm{Th}^a_l(w)=a\,\mathrm{Th}_l(w)$ and
  $\mathrm{Sh}^a_l(w)=a\,g^{frn_w}\,\mathrm{Sh}_l(w)$.
\item (Theta modules.) For $\mathcal O_F$-submodules $X,Y$ of $A_l$, $XY$ is the $\mathcal O_F$-span of
  the products $xy$, and $X^0=\mathcal O_F\cdot1$. Let $\kappa,w_L,w_F$ be weights and $T\geq0$ an
  integer, put $w_{F,0}=w_F+d_Yw_L$, and let $S=(s_W,s_F,s_0,s_K)$ be four elements of $\mathcal O_F$.
  The \emph{theta module} of $S$ is
  \[
  \begin{aligned}
   \Theta(S)={}&\iota_{w_L}\big(\mathrm{Th}^{s_W}_l(w_L)\big)^T\,\iota_{w_F}\big(\mathrm{Sh}^{s_F}_l(w_F)\big)
   +\iota_{w_L}\big(\mathrm{Th}^{s_W}_l(w_L)\big)^{T-d_Y}\,
   \iota_{w_{F,0}}\big(\mathrm{Sh}^{s_0}_l(w_{F,0})\big)\\
   &+\iota_\kappa\big(\mathrm{Sh}^{s_K}_l(\kappa)\big),
  \end{aligned}
  \]
  with $T-d_Y$ read as $0$ when $T<d_Y$. The \emph{plain theta module} $\Theta^{\mathrm p}$ is the same
  expression with $\mathrm{Th}_l$ and $\mathrm{Sh}_l$ in place of $\mathrm{Th}^{s_W}_l$ and of the
  three modules $\mathrm{Sh}^{s}_l$. The run satisfies \emph{(H2)} for $S$ if
  $\mathrm{Th}^{s_W}_l(w_L)\subset M_l(w_L)$, $\mathrm{Sh}^{s_F}_l(w_F)\subset M_l(w_F)$,
  $\mathrm{Sh}^{s_0}_l(w_{F,0})\subset M_l(w_{F,0})$ and $\mathrm{Sh}^{s_K}_l(\kappa)\subset
  M_l(\kappa)$.
\item (Scalars.) Put $N(w)=\prod_{k=0,1}\pi_k^{D^+(w)}\bar\pi_k^{D^-(w)}$, read in $\mathcal O_F$, and
  $P(a,b)=Q_S^aq^bd^b$; for a prime $p$ put $z_p(w)=p^{\deg w+2frn_w}$. With
  $e=T\deg w_L+\deg w_F$ and $c=Tc(w_L)+c(w_F)$, the \emph{scalars} of $p$ are
  \[
  \begin{aligned}
   s_W&=N(w_L)P(\deg w_L,c(w_L)),\qquad s_F=z_p(w_F)N(w_F)P(\deg w_F,c(w_F)),\\
   s_0&=z_p(w_{F,0})N(w_L)^{d_Y}N(w_F)P\big(d_Y\deg w_L+\deg w_F,\,d_Yc(w_L)+c(w_F)\big),\\
   s_K&=z_p(\kappa)N(w_L)^TN(w_F)P(e,c),
  \end{aligned}
  \]
  and the \emph{record scalar} is $\sigma=N(\kappa)\,g^{frn_{w_F}}P(e,c)$.
\item The theta module $\Theta_l(p)$ of Definition~\ref{is-onto-step-index-bound-conclusion} is
  $\Theta(S)$ for the scalars $S$ of $p$ and the weights $\kappa'_l$
  (Definition~\ref{is-onto-step-leaves-kappa}), $w_{L,l}$, $w_{F,l}$ and the integer $T=T_O$ of the
  parameter choices. For these weights write $\Theta^{\mathrm p}_l$ for the
  plain theta module and $\sigma_l$ for the record scalar.
\item (Short series.) A \emph{series datum} of weight $w$ is a triple $(\Lambda',N,P)$: a unimodular
  lattice $\Lambda'$ that is trivial at the bad primes and at the level of the delta data
  $\delta_l(n_w)$ of the run, an integer $N$ for which $\Lambda'$ has level $N$ relative to that of
  $\delta_l(n_w)$, and a harmonic polynomial $P$ for $\Lambda'$ and $w$. Let $\Lambda_l$ be the lattice
  of the core at level $l$ (Definition~\ref{is-onto-step-model-level-condition}), with its level $M$
  and level ideal $M_b$, and let $p+q$ be the rank of $\Lambda_l$ (the signature of the datum of the
  core at level $l$). A vector $x=(x_1,\dots,x_{n_w})\in\Lambda_l^{n_w}$ is \emph{allowed} if
  $\Lambda_l$ is spanned by $\bar M_b\Lambda_l$ and the $x_b$ with $b\le p+q$, or with any $b$ when
  $n_w=n_V$. For an
  allowed vector $x$, a \emph{short series} of $(\Lambda',N,P)$ at $x$
  is a section of weight $w$ that satisfies the short rule of $\delta_l(n_w)$ for $\Lambda'$, $N$, $x$
  and $P$.
\item (Right multiples.) Let $\gamma$ be an automorphism of the level module $\Lambda_l/M_b\Lambda_l$,
  with a lift $\tilde\gamma$ to $\Lambda_l$. A vector $x'$ is a \emph{right multiple} of $x$ by
  $\gamma$ if $\langle x'_b,ca\rangle\equiv\langle x_b,c\tilde\gamma(a)\rangle\pmod M$ for every index
  $b$ and every $a\in\Lambda_l$, where $\langle\,,\rangle$ is the pairing of the datum of the core at
  level $l$, $M$ the level and $c$ its separating element. The automorphism $\gamma$ has
  \emph{determinant one} if $\det(\tilde\gamma\otimes\mathcal O_{E_h}/M_b)=1$. For a subgroup $U$ of a
  group and a map $\ell$ from $U$ to the automorphisms of the level module, an \emph{orbit family} of
  weight $w$ for $\ell$ is a family $(\phi_u)_{u\in U}$ for which there are a series datum
  $(\Lambda',N,P)$ of weight $w$ and allowed vectors $x_0$ and $x_u$, $u\in U$, such that each $x_u$ is a
  right multiple of $x_0$ by $\ell(u)$ and each $\phi_u$ is a short series of $(\Lambda',N,P)$ at $x_u$.
\end{enumerate}
\end{definition}
\leanref{Formalism/SharpRestriction/Divisibility/WeightAlgebra.lean}{Formalism.MainNotions.WeightAlgebra.PJ}
\leanref{Formalism/SharpRestriction/Divisibility/WeightAlgebra.lean}{Formalism.MainNotions.WeightAlgebra.PJ.Theta}
\leanref{Formalism/SharpRestriction/Divisibility/WeightAlgebra.lean}{Formalism.MainNotions.WeightAlgebra.H2}
\leanref{Development/AbelianVarietyTheory/Definitions/Jets/EmbShortThLeThetaPlain.lean}{GlobalTheory.FormTies.thetaPlain}
\leanref{Formalism/SharpRestriction/Divisibility/AbelianOpenPart/OEOsInjective.lean}{Formalism.MainNotions.DIV.Cons.scalOfP}
\leanref{Development/AbelianVarietyTheory/Definitions/Jets/EmbShortThLeThetaPlain.lean}{GlobalTheory.FormTies.sigRec}
\leanref{Development/AbelianVarietyTheory/Proofs/FoundationProofs/TwistTrivialOnRational.lean}{GlobalTheory.TraceLattice.GoodDatum}
\leanref{Development/AbelianVarietyTheory/Proofs/FoundationProofs/TwistTrivialOnRational.lean}{GlobalTheory.TraceLattice.IsSeries}
\leanref{Development/AbelianVarietyTheory/Proofs/FoundationProofs/TwistTrivialOnRational.lean}{GlobalTheory.TraceLattice.DetOne}
\leanref{Development/AbelianVarietyTheory/Definitions/Foundations/ThetaSeriesTies.lean}{GlobalTheory.OnTies.IsRightMul}
\leanref{Development/AbelianVarietyTheory/Definitions/Foundations/ThetaSeriesTies.lean}{GlobalTheory.OnTies.IsOrbitFamily}
% Lean: the ring of forms at level l is recPJ h l hl : WeightAlgebra.PJ D (recB h.toLawsCore) P6 l, with
% Forms, inc, inc_mul, vAt, vAt_mul and Mlat_eq for (2); Th_l(w) and Sh_l(w) are its fields thAll w and
% shortTh w, which recPJ takes to be spanIn T.C l w (T.th l w) and spanIn T.C l w (T.sh l w). The weights
% are WeightFamilies.Wt (aP, aM, n); wadd and wsmul add and scale them; deg is DF = degP, D^+ and D^-
% are Dplus and Dminus, c(w) is cHalf and d_Y is WeightFamilies.dimY D l. Th^a and Sh^a are WWc a w and
% Thc a w (scaled a, with the factor (g0 g1)^(f r n_w)); X^T is the power of submodules, and T - d_Y
% is truncated subtraction. The scalars are scalOfP (sW, sF, s0, sK), with N(w) = nnW (EX.piG through
% iE), P(a, b) = pPow a b b and z_p = zOwnP; sigma is sigRec. A series datum is GoodDatum (GenLat with
% Unimod, TrivAt (badN D), TrivAt of the level of T.dd l n, LevOK N, and P in harmL); an allowed
% vector is AdmX; a short series is IsSeries (sec and IsShortRule); a right multiple is IsRightMul
% (Dw.EF, cSep, liftLa); an orbit family is IsOrbitFamily (T.Q.Xo l) T.Q.c (badN D) (wtA w)
% (T.dd l w.n).

\begin{lemma}\label{is-onto-step-leaves-weights}
Let $p$ be a point of a family, with definite-side datum of ranks $f$ and $r$ and column rank $n_V$,
and parameter choices with integers $T_O$, $n_P$, $n_W$ and $D_p$. Let $l$ be a level of a lift, with the
weights $\kappa=\kappa'_l$, $w_L=w_{L,l}$ and $w_F=w_{F,l}$ of the parameter choices
(Definition~\ref{is-onto-step-leaves-theta-module}). Then:
\begin{enumerate}
\item $w^\pm_{F,\tau,i}=\kappa^\pm_{\tau,i}-T_Ow^\pm_{L,\tau,i}$ for all $\tau$ and $i$, $n_{w_F}=n_P$ and
  $n_{w_L}=n_W$;
\item $n_\kappa=n_V=n_{w_F}+T_On_{w_L}$;
\item $d_Y\le T_O$;
\item $r\le n_{w_F}$.
\end{enumerate}
\end{lemma}
\leanref{Development/AbelianVarietyTheory/Proofs/JetHarmonics/RunWeightsSum.lean}{GlobalTheory.RunDepth.run_weights_sum}
\begin{proof}
(1) The weight $w_F$ is defined through the gaps between its consecutive entries, which are those of
$\kappa$ less $T_O$ times those of $w_L$; they are nonnegative by a computation with the floor
divisions by $T_O+d_Y+2\geq T_O+2$ that define the gaps of $w_L$. Summing the gaps gives the
entries. The integers
of $w_F$ and $w_L$ are $n_P$ and $n_W$ by definition.
(2) The weight $\kappa'_l$ has the integer $n_V$ (Definition~\ref{is-onto-step-leaves-kappa}), and the
parameter choices give
$n_V=r+t=n_P+T_On_W$.
(3) At a level of a lift $d_Y=f\binom{r}{2}$, which is $D_p$, and the parameter choices give
$40(D_p+2)\le T_O$.
(4) The parameter choices give $n_P=8fr+4$, and $f\geq1$.
% Lean: run_weights_sum: (1) is RunWeightTest.isF'_fWt for w_L = lWt and w_F = fWt (fGapP and fGapM
% with Int.toNat_of_nonneg, ofGaps_eq_sub_mul, mul_toNat_fdiv_add_two_le_of_add_two_le with
% T + 2 <= Tp); (2) is kaWt_n with Tp.split; (3) is DepthBoundTest.dimY_eq_flag with Tp.Dp_eq and
% Tp.TO_40; (4) is Tp.nP_eq.
\end{proof}

\begin{lemma}\label{is-onto-step-leaves-integral}
Let $s$ and $l$ be as in Definition~\ref{is-onto-step-leaves-theta-module}. Let $\kappa,w_L,w_F$ be
weights and $T\geq d_Y$ an integer with $w^\pm_{F,\tau,i}=\kappa^\pm_{\tau,i}-Tw^\pm_{L,\tau,i}$ for all
$\tau$ and $i$ and $n_\kappa=n_{w_F}+Tn_{w_L}$, and let $S$ be four elements of $\mathcal O_F$ for which
the run satisfies (H2). Then $\Theta(S)\subset\iota_\kappa(M_l(\kappa))$.
\end{lemma}
\leanref{Formalism/SharpRestriction/GenusPolyhedralBounds/RelativeIndexClassSets/ThetaLeMsupOfSetC.lean}{Formalism.MainNotions.Gaps.T0.theta_le_msup}
\begin{proof}
For weights $w$ and $w'$, $\iota_w(M_l(w))\,\iota_{w'}(M_l(w'))\subset\iota_{w+w'}(M_l(w+w'))$: for
$\phi\in M_l(w)$ and $\phi'\in M_l(w')$ there is $\phi''\in S_l(w+w')$ with
$\iota_w(\phi)\iota_{w'}(\phi')=\iota_{w+w'}(\phi'')$, and
$v_{\lambda,y}(\iota_{w+w'}\phi'')=v_{\lambda,y}(\iota_w\phi)+v_{\lambda,y}(\iota_{w'}\phi')\geq0$ for
all $\lambda$ and $y$, so $\phi''\in M_l(w+w')$. By induction on $n$,
$\iota_{w_L}(M_l(w_L))^n\,\iota_w(M_l(w))\subset\iota_{nw_L+w}(M_l(nw_L+w))$ for every weight $w$ and
every $n\geq0$. The hypotheses give $Tw_L+w_F=\kappa$, entry by entry and in $n$; and as $T\geq d_Y$,
$(T-d_Y)w_L+w_{F,0}=Tw_L+w_F=\kappa$. By (H2), $\mathrm{Th}^{s_W}_l(w_L)\subset M_l(w_L)$ and each of
the three modules $\mathrm{Sh}^{s}_l$ lies in the module $M_l$ of its weight. So the first two summands
of $\Theta(S)$ lie in $\iota_\kappa(M_l(\kappa))$ by the product rule with $n=T$ and $n=T-d_Y$, and the
third does by (H2).
% Lean: theta_le_msup: theta_first_le and theta_second_le (emb_pow_mul_le from msup_pow_mul_le and
% mlat_mul, with inc_mul, vAt_mul and Mlat_eq; wt_first and wt_second for the weights) and
% theta_third_le. About 110 lines.
\end{proof}

\begin{lemma}\label{is-onto-step-leaves-scalar}
Let $s$, $l$, $\kappa,w_L,w_F$ and $T\geq d_Y$ be as in Lemma~\ref{is-onto-step-leaves-integral}, let $p$
be a prime, and let $S=(s_W,s_F,s_0,s_K)$ be the scalars of $p$ and $\sigma$ the record scalar
(Definition~\ref{is-onto-step-leaves-theta-module}(5)). Then $\sigma\ne0$, $s_K\ne0$, and every element of
$\Theta(S)$ is $\sigma x$ for some $x\in\Theta^{\mathrm p}$.
\end{lemma}
\leanref{Development/AbelianVarietyTheory/Proofs/JetHarmonics/ThetaScaledSpanLeSigmaSmulUnscaledSpanStatement.lean}{GlobalTheory.T.theta_scaled_span_le_sigma_smul_unscaled_span}
\leanref{Development/AbelianVarietyTheory/Definitions/Jets/EmbShortThLeThetaPlain.lean}{GlobalTheory.FormTies.runScalar_kappa_ne_zero}
\begin{proof}
Since $\pi_k\bar\pi_k=g_k\ne0$, $N(w)\ne0$ for every weight $w$; $P(a,b)\ne0$, as $Q_S$, $q$ and $d$ are
positive; and $z_p(w)\ne0$. So $\sigma\ne0$ and $s_K\ne0$. The hypotheses give
$D^\pm(\kappa)=TD^\pm(w_L)+D^\pm(w_F)$, so $N(\kappa)=N(w_L)^TN(w_F)$; and
$P(a,b)P(a',b')=P(a+a',b+b')$.

For $a,b\in\mathcal O_F$ and $\mathcal O_F$-submodules $X,Y$ of $A_l$, $(aX)(bY)=ab\,XY$ and
$(aX)^n=a^nX^n$, and $aX\subset bX$ when $b$ divides $a$ in $\mathcal O_F$. Hence the first summand of
$\Theta(S)$ lies in $s_W^Ts_Fg^{frn_{w_F}}$ times the first summand of $\Theta^{\mathrm p}$, the second
in $s_W^{T-d_Y}s_0g^{frn_{w_{F,0}}}$ times the second, and the third in $s_Kg^{frn_\kappa}$ times the
third. With $n_{w_{F,0}}=n_{w_F}+d_Yn_{w_L}$, $n_\kappa=n_{w_F}+Tn_{w_L}$ and the identities above,
\[
 s_W^Ts_Fg^{frn_{w_F}}=z_p(w_F)\,\sigma,\qquad
 s_W^{T-d_Y}s_0g^{frn_{w_{F,0}}}=z_p(w_{F,0})g^{frd_Yn_{w_L}}\,\sigma,\qquad
 s_Kg^{frn_\kappa}=z_p(\kappa)g^{frTn_{w_L}}\,\sigma.
\]
So each summand of $\Theta(S)$ lies in $\sigma$ times the corresponding summand of $\Theta^{\mathrm p}$.
% Lean: T.theta_scaled_span_le_sigma_smul_unscaled_span = ScaledSpan.theta_le_sigma_smul, with
% nnW_ne_zero, nnW_eq (N(kappa) = N(w_L)^T N(w_F)), emb_WWc_le, emb_Thc_le, sc_mul_le, sc_pow_le and
% sc_le_of_dvd with the cofactors zOwnP; s_K is runScalar_kappa_ne_zero. About 130 lines.
\end{proof}

\begin{lemma}\label{is-onto-step-leaves-hypotheses-discrete}
Let $F$ be a number field and $V$ a finite-dimensional $F$-vector space, and for every embedding
$\sigma\colon F\to\mathbf C$ let $N_\sigma\colon V\to\mathbf R_{\geq0}$ be a function with
$N_\sigma(av)=|\sigma(a)|N_\sigma(v)$ and $N_\sigma(v+v')\le N_\sigma(v)+N_\sigma(v')$, vanishing only at
$0$. Let $M\subset V$ be an $\mathcal O_F$-submodule and $B$ a real number such that
$\sum_\sigma\log N_\sigma(v)\geq-B$ for every nonzero $v\in M$. Then $M$ is finitely generated.
\end{lemma}
\leanref{ReciprocityTheory/ReductiveGroups/ArithmeticGeometry/TaylorDiscriminants/DiscriminantEstimates/DiscriminantTheorems/TBDisc.lean}{discriminant_of_product_lower}
\leanref{ReciprocityTheory/ReductiveGroups/ArithmeticGeometry/TaylorDiscriminants/DiscriminantEstimates/DiscriminantTheorems/TBDisc.lean}{formalGroup_of_discriminant}
\begin{proof}
\emph{Uniform discreteness.} Put $\varepsilon=e^{-|B|}$. If $v\in M$ has $N_\sigma(v)<\varepsilon$ for every
$\sigma$, then $v=0$: otherwise every $N_\sigma(v)$ is positive and
$\sum_\sigma\log N_\sigma(v)<-[F:\mathbf Q]\,|B|\le-B$.

\emph{Finite generation.} Fix a basis $e_1,\dots,e_n$ of $V$, and for $v=\sum_ia_ie_i$ let
$c(v)=(\sigma(a_i))_{\sigma,i}\in\mathbf C^{[F:\mathbf Q]n}$, with the maximum norm. Then
$N_\sigma(v)\le\|c(v)\|\sum_iN_\sigma(e_i)\le C\|c(v)\|$ with $C=1+\sum_{\sigma,i}N_\sigma(e_i)$. So $v\in M$ with
$\|c(v)\|<\varepsilon/C$ is $0$: the additive map $c$ is injective on $M$, and $c(M)$ is a discrete subgroup
of $\mathbf C^{[F:\mathbf Q]n}$. A discrete subgroup of a finite-dimensional real vector space is a
finitely generated abelian group, so $M$ is finitely generated over $\mathbf Z$, hence over $\mathcal O_F$.
% Lean: discriminant_of_product_lower = TBDisc.discOfProdLower (UnifDiscrete with exp(-|A|)) and
% formalGroup_of_discriminant = fg_of_unifDiscrete (coordMap, norm_le_coord, DiscreteTopology of the
% range and Module.Finite over Z). The Lean's NormFam is the book's family N_sigma. About 170 lines.
\end{proof}

\begin{lemma}\label{is-onto-step-leaves-hypotheses-sup}
There is $C_0>0$ with the following property. Let $s$ be a run at a point of a family that satisfies the
laws, with definite-side datum of ranks $f$ and $r$, host field $E_h$ with discriminant $d_{E_h}$ and
level primes $g_0,g_1$, and size data with logarithmic size $L$ and the number $\log\Omega=\gamma+
\log(2\pi/\operatorname{rd}_{E_h})-\ell'$, where $\gamma$ is Euler's constant, $\operatorname{rd}_{E_h}$ the
root discriminant and $\ell'$ the logarithmic derivative at $1$ of the size data; suppose that $\log
g_i\le9L+1$ for $i=0,1$. Let $l$ be a level of a lift, $w$ a weight at level $l$ with $q_l$ negative
entries at each embedding, of size $|w|=\sum_{\tau,i}(w^+_{\tau,i}+w^-_{\tau,i})+f\max(q_l,1)rn_w+1$, and
let $C\geq C_0$ and $A\geq2$. Then every nonzero $\phi\in M_l(w)$ satisfies
\[
 \frac1{[F:\mathbf Q]}\sum_{\sigma\colon F\to\mathbf C}\log\|\phi\|_{\sigma,\infty}\geq
 -C\big(fr+L+\log|d_{E_h}|+|\log\Omega|\big)^A\,(1+\log|w|)\,|w|,
\]
where $F$ is the coefficient field of the core and $\|\phi\|_{\sigma,\infty}$ the supremum norm of $\phi$ at
$\sigma$.
\end{lemma}
\leanref{Development/AbelianVarietyTheory/Proofs/CubicalLineBundles/MeanLogSupNormLowerBoundOfCoreLawsStatement.lean}{GlobalTheory.T.meanLogSupNorm_lower_bound_of_coreLaws}
\leanref{Development/AbelianVarietyTheory/Proofs/CubicalLineBundles/MeanLogSupNormLowerBoundOfCoreLaws.lean}{GlobalTheory.lamwAt_of_lawsCore}
\begin{proof}
The proof follows the route of Lemma~\ref{cv-cusp-inputs}(3) for the moduli scheme of the core at level
$l$, at the good embedding $\sigma_0$ of $F$ that the core fixes, with the level $M=g_0g_1$ of the core.
Let $\phi\in M_l(w)$ be nonzero, let $p_l+q_l$ be the signature of the datum of the core at level $l$,
and put $s=\sum_{\tau,i}(w^+_{\tau,i}+w^-_{\tau,i})$ and $t=frn_w$.

\emph{Step 1: sizes.} The laws of the core and the datum give: $r\geq1$ and $p_l+q_l=r$;
$[E_h:\mathbf Q]=2f$; $\sqrt d\le|d_{E_h}|$ and $1\le|d_{E_h}|$; the spread of the weight that $w$ gives on
every component at $\sigma_0$ is at most $|w|$; and $s+t\le|w|$.

\emph{Step 2: near a maximum.} The reading of $\phi$ at $\sigma_0$ is a holomorphic cusp form on the
components at $\sigma_0$; its supremum is finite, and it is positive, as $\phi\ne0$ and, by the laws of
the core, a cusp form with supremum $0$ vanishes. So there are a component $Z$ and a frame $z_0$ of $Z$
with $|\phi|(y)\le2|\phi|(z_0)$ for every frame $y$ of $Z$, and $|\phi|(z_0)>0$.

\emph{Step 3: lattice data.} The uniformization of the core at $\sigma_0$ gives at $Z$ a lattice
$L'\subset E_h^r$, self-dual for the Hermitian matrix of $Z$, whose multiples by the ideal $\mathfrak A$
of the host span the lattice of $Z$, and a real number $k_0>0$ such that, with
$M'=g_0g_1\,\mathrm N(\mathfrak c_K)\,d$, $M'k_0=2g_0g_1\sqrt d$, and for every nonzero isotropic $v\in L'$
and every integral $b\in E_h$ fixed by complex conjugation the arithmetic group of $Z$ contains the
transvection of $v$ by $M'b$, on which the character of the weight is trivial. By Step 1 and
$\log g_i\le9L+1$, $\log(M'k_0)=\log(2g_0g_1\sqrt d)\le19L+3+\log|d_{E_h}|$.

\emph{Step 4: a CM frame.} Lemma~\ref{cv-cusp-frame}, in the form for this uniformization, applied to
$\phi$ at $z_0$ with the spread bound of Step 1 and the transvections of Step 3, gives a frame $x$ of $Z$
with $\phi(x)\ne0$ and a basis $v_1,\dots,v_r$ of $E_h^r$ in $L'$ with properties (1) to (4) of that
lemma for a number $\ell$ determined by $f$, $r$, an integer $p_0\le r$ of the lattice data,
$|d_{E_h}|$, $M'$, $k_0$ and $|w|$.

\emph{Step 5: the constant.} There is $C_0>0$, independent of all the data, such that this $\ell$ is at
least $1$ and, for all $s',t'\geq0$ with $s'+t'\le|w|$, $C\geq C_0$ and $A\geq2$,
\[
 \Big(\frac{1+|\log\Omega|}2+\frac12\log\ell\Big)(s'+t')+s'\log(r\ell)
 \le C\big(fr+L+\log|d_{E_h}|+|\log\Omega|\big)^A(1+\log|w|)\,|w|,
\]
given $\log(M'k_0)\le19L+3+\log|d_{E_h}|$.

\emph{Step 6: the CM point.} There is a Kronecker character $\chi$ of $\mathbf Q(\sqrt{-d})$ such that
the real part of $\log\Omega_{E_h}(\chi)$ (Definition~\ref{cv-log-omega}) is $\log\Omega$. The supremum norms of $\phi$ are
positive at every embedding of $F$. As in Lemma~\ref{cv-cusp-cm-point} and
Proposition~\ref{cv-cusp-cm-product}, the CM point of $x$ and $v_1,\dots,v_r$ gives
\[
 -\Big(\frac{1+|\log\Omega|}2+\frac12\log\ell\Big)(s+t)-s\log(r\ell)
 \le\frac1{[F:\mathbf Q]}\sum_{\sigma\colon F\to\mathbf C}\log\|\phi\|_{\sigma,\infty}.
\]
With Step 5 for $s'=s$ and $t'=t$, which is allowed by Step 1, this is the lemma.
\openstep{The sizes of Step 1, the lattice data and transvections of Step 3, the constant of Step 5, the
Kronecker character of Step 6, and the lower bound at the CM point for the uniformizations of the core,
which the book proves only in Situation~\ref{cv-cusp-setting} (Lemma~\ref{cv-cusp-cm-point} and
Proposition~\ref{cv-cusp-cm-product}), are not written out.}
% Lean, the top level: lamwAt_of_lawsCore. hj from cutR_zero; Step 1 is WtSizes.sizes (hr, hpq, hf,
% hsd, hd1, hlogd, hwf, hsz); Step 2 is mem_cusp at T.good_sigma0, holSub_SS with
% siegelSpaceAlgebra_holds, supE_cusp_ne_top, h.supDef and abelian2.half_max; Step 3 is
% LatticeData.dataXOfModel_of_gamAt (Xx with g = M and Lambda = AA L') and TransvModel.transv_model
% (hMk: M' k0 = 2 g0 g1 sqrt d; htr with h.chiTriv), and log_two_mul_mul_le; Step 4 is
% abelian2.cm_frame_of_eq (x, lam, LamBasis hLB, hPx) with La the frameBasisRadius; Step 5 is
% T.exists_constant_cmEvaluation_fits_room (= room_of_constants: hLa1 and hr'); Step 6 is
% LamT.exists_kronecker_logOm_core (Re logOm chi = large2.logOm), hsup, hSpos and
% T.meanLogSupNorm_lowerBound_cmFrame_nonvanishing (with h.unf, h.xiAbs, AA_ne_zero), then h3
% (meanLogSup) and the final linarith with hP (PP0v) and hOm. In Steps 5 and 6 the book writes
% (1+|log Omega|)/2 for the Lean's pDeg (Re logOm chi), as in Proposition cv-cusp-cm-product. Sizes,
% following imports and stopping at statements the book cites: lamwAt_of_lawsCore about 3,930 lines
% (2,858 in theorems); WtSizes.sizes 2,825 (2,022 in theorems, mostly sqrt d <= |d_E| through the
% ramification of E_h); dataXOfModel_of_gamAt with transv_model 1,032; exists_kronecker_logOm_core 823;
% exists_constant_cmEvaluation_fits_room 147; half_max, holSub_SS and supE_cusp_ne_top together 438.
% The CM-point bound T.meanLogSupNorm_lowerBound_cmFrame_nonvanishing is cited in Lemma cv-cusp-inputs,
% so lean_size stops at it.
% Earlier comment: lamwAt_of_lawsCore: half_max (a frame z0 where the norm is at least half its supremum),
% LatticeData.dataXOfModel_of_gamAt and TransvModel.transv_model (the data of the uniformization),
% Formalism.L2.Pj.abelian2.cm_frame_of_eq (Lemma cv-cusp-frame) for a CM frame, WtSizes.sizes (the sizes
% of the weight), T.exists_constant_cmEvaluation_fits_room (the constant C0), LamT.exists_kronecker_logOm_core
% and T.meanLogSupNorm_lowerBound_cmFrame_nonvanishing (cited in Lemma cv-cusp-inputs). GamAt comes from
% the laws (RunIn.gamAt_of_lawsCore'). Size, following imports and stopping at statements the book
% cites: about 3,900 lines below lamwAt_of_lawsCore, 2,900 of them in theorems; 2,800 of these lines are
% below WtSizes.sizes.
\end{proof}

\begin{lemma}\label{is-onto-step-leaves-hypotheses}
Let $s$ be a run at a point of a family that satisfies the laws, $l$ a level of a lift and $p$ a prime,
and suppose that $s$ satisfies the hypotheses of the index-parameter conditions at $(l,p)$
(Definition~\ref{is-onto-step-index-bound-conclusion}). Let $\kappa=\kappa'_l$, and let $S$ be the
scalars of $p$ for the weights and the integer of the parameter choices at $l$
(Definition~\ref{is-onto-step-leaves-theta-module}). Then the run satisfies (H2) for $S$, the module
$\mathrm{Sh}_l(\kappa)$ spans $S_l(\kappa)$ over $F$, and $M_l(\kappa)$ is a lattice in $S_l(\kappa)$.
\end{lemma}
\leanref{Development/AbelianVarietyTheory/Proofs/DivisorProofs/LeavesSupSpanData.lean}{GlobalTheory.BigOn.lat_span_of_hypC}
\leanref{Development/AbelianVarietyTheory/Proofs/CubicalLineBundles/WeightModuleIsLatticeOfShortSeriesSpan.lean}{GlobalTheory.latAt_record}
\begin{proof}
Condition (H2) for $S$ is one of the hypotheses of the index-parameter conditions. Another is a list of
conditions on the scalars, which include $s_K\ne0$ and a finite subset of
$\mathrm{Sh}^{s_K}_l(\kappa)$ that spans $S_l(\kappa)$ over $F$. As
$\mathrm{Sh}^{s_K}_l(\kappa)=s_Kg^{frn_\kappa}\,\mathrm{Sh}_l(\kappa)\subset\mathrm{Sh}_l(\kappa)$, the
module $\mathrm{Sh}_l(\kappa)$ spans $S_l(\kappa)$.

The module $M_l(\kappa)$ is finitely generated. Among the hypotheses are parameter conditions, which
include $\log g_i\le9L+1$ for $i=0,1$. Let $C_0$ be the constant of
Lemma~\ref{is-onto-step-leaves-hypotheses-sup}. That lemma, with $C=C_0$ and $A=2$, gives
$\sum_\sigma\log\|\phi\|_{\sigma,\infty}\geq-B$ for every nonzero $\phi\in M_l(\kappa)$, with
$B=[F:\mathbf Q]\,C_0\big(fr+L+\log|d_{E_h}|+|\log\Omega|\big)^2(1+\log|\kappa|)|\kappa|$. The supremum
norms vanish only at $0$, satisfy $\|a\phi\|_{\sigma,\infty}=|\sigma(a)|\,\|\phi\|_{\sigma,\infty}$ and the
triangle inequality, and $S_l(\kappa)$ is finite-dimensional; so $M_l(\kappa)$ is finitely generated by
Lemma~\ref{is-onto-step-leaves-hypotheses-discrete}.
% Lean: latAt_record = PullbackSections.latOfLam with GlobalTheory.T.meanLogSupNorm_lower_bound_of_coreLaws
% (C = C0, A = 2) and the parameter conditions ParA of the hypotheses (hA.2.2.2.2.1); hPL is the product
% lower bound (ProdLower) for the sup norms supN (supN_nonneg, supN_eq_zero, supN_smul, supN_add) on
% Mlat, then discriminant_of_product_lower and formalGroup_of_discriminant give FG. The open part is
% Lemma is-onto-step-leaves-hypotheses-sup.
It spans $S_l(\kappa)$: with $c=s_Kg^{frn_\kappa}\ne0$, for $\phi\in\mathrm{Sh}_l(\kappa)$ the element
$c\phi$ lies in $\mathrm{Sh}^{s_K}_l(\kappa)\subset M_l(\kappa)$ by (H2), so $\phi=c^{-1}(c\phi)$ lies in
the $F$-span of $M_l(\kappa)$, and $\mathrm{Sh}_l(\kappa)$ spans $S_l(\kappa)$.
% Lean: lat_span_of_hypC: hH2 is the H2 clause of HypC; hK (s_K /= 0) and hka (ShortG for kappa, the
% last clause of S4) come from the scalar conditions ScalB; RunLat.spanAt_of_shortG with thc_le_shortTh
% gives SpanAt; latAt_record gives LatAt, whose span part is the end of latOfLam.
\end{proof}

\begin{lemma}\label{is-onto-step-leaves-short}
Let $s$ and $l$ be as in Definition~\ref{is-onto-step-leaves-theta-module}, let $\kappa,w_L,w_F$ be
weights, $T\geq0$ an integer and $S$ the scalars of a prime $p$. If $M_l(\kappa)$ is a lattice in
$S_l(\kappa)$ and the run satisfies (H2) for $S$, then $\mathrm{Sh}_l(\kappa)$ is finitely generated.
Of these hypotheses the proof uses only that $M_l(\kappa)$ is finitely generated and that
$\mathrm{Sh}^{s_K}_l(\kappa)\subset M_l(\kappa)$.
\end{lemma}
\leanref{Development/AbelianVarietyTheory/Definitions/Jets/EmbShortThLeThetaPlain.lean}{GlobalTheory.FormTies.shortTh_fg}
\begin{proof}
The element $c=s_Kg^{frn_\kappa}$ of $\mathcal O_F$ is nonzero, as in the proof of
Lemma~\ref{is-onto-step-leaves-scalar}. Multiplication by $c$ is injective on the $F$-vector space
$S_l(\kappa)$ and maps $\mathrm{Sh}_l(\kappa)$ onto $\mathrm{Sh}^{s_K}_l(\kappa)$, a submodule of the
finitely generated module $M_l(\kappa)$ over the noetherian ring $\mathcal O_F$. So
$\mathrm{Sh}_l(\kappa)$ is finitely generated.
% Lean: shortTh_fg = shortTh_fg_of_Thc_le with runScalar_kappa_ne_zero and the fourth clause of H2.
\end{proof}

\begin{lemma}\label{is-onto-step-leaves-products}
Let $s$ be a run at a point of a family that satisfies the laws and is covered at a lift $j$, and put
$l=j+1$. Let $\kappa,w_L,w_F$ be weights and $T\geq d_Y$ an integer with
$w^\pm_{F,\tau,i}=\kappa^\pm_{\tau,i}-Tw^\pm_{L,\tau,i}$ for all $\tau$ and $i$,
$n_\kappa=n_{w_F}+Tn_{w_L}=n_V$ and $r\le n_{w_F}$
(Definition~\ref{is-onto-step-leaves-theta-module}). Then $\Theta^{\mathrm p}\subset
\iota_\kappa(\mathrm{Sh}_l(\kappa))$.
\end{lemma}
\leanref{Development/AbelianVarietyTheory/Proofs/JetHarmonics/ProductsShortRec.lean}{GlobalTheory.FormTies.productsShort_rec}
\begin{proof}
\openstep{The proof that the products of theta series of weight $w_L$ with short theta series are
short theta series of the sum weight, through the independence of the Taylor values of germs from the
formal chart, is not written out.}
% Lean: productsShort_rec = productsShort_of_oneFactor ShortTimesTheta.short_mul_theta_mem_span, with
% LawsSh, OneCover (covered), the trivial level character (T.dd l 0).l2.A = 1 of the laws (lvA) and
% chart independence ChartIndepSt2 (GlobalTheory.chartIndependence_unconditional, cited in Lemma
% cv-uniformization-theta-rule-unique). Size, following imports and stopping at statements the book
% cites (chart independence among them): about 35,100 lines, 29,800 of them in theorems.
\end{proof}

\begin{lemma}\label{is-onto-step-leaves-move-laws}
Let $s$ be a run at a point of a family that satisfies the laws, and $l$ a level of a lift. For an
automorphism $\gamma$ of the level module at level $l$, write $\gamma y$ for the move of a framed point
$y$ of the moduli scheme of the core at level $l$ by a level automorphism whose matrix induces $\gamma$
(Definition~\ref{cv-realization-level-move}), and $\gamma\phi$ for the image of a section $\phi$ at level
$l$ under the level action of $\gamma$ (Lemma~\ref{is-onto-step-level-structure-cusp}). Then:
\begin{enumerate}
\item if $y$ is the base change of a framed point over a test ring, a reduced finitely generated algebra
  over the coefficient field $F$ of the core, then so is $\gamma y$, and the value of $\gamma\phi$ at $y$
  is the value of $\phi$ at $\gamma y$;
\item $(\gamma\delta)\phi=\gamma(\delta\phi)$ for all automorphisms $\gamma,\delta$ of the level module.
\end{enumerate}
\end{lemma}
\leanref{Formalism/LocalToGlobal/ProjectiveLattices/SpMapAlgHomId.lean}{Formalism.L2.Pj.evalAt_actL}
\leanref{Development/AbelianVarietyTheory/Proofs/JetHarmonics/LevActMemMlat.lean}{GlobalTheory.FormTies.fromTest_mvPt}
\leanref{Formalism/LocalToGlobal/ProjectiveLattices/RhoWApply.lean}{Formalism.L2.Pj.actL_mul}
\leanref{Development/AbelianVarietyTheory/Proofs/JetHarmonics/ActLevFullLevelGroupFramedPointsStatement.lean}{GlobalTheory.T.actLev_full_level_group_framedPoints}
\begin{proof}
The level action of $\gamma$ maps $\phi$ to the section whose value at a framed point $y$ over a test
ring is the value of $\phi$ at $\gamma y$; at a framed point over any $F$-algebra, the value of a section
is read through a framed point over a test ring of which it is the base change. The datum of the core at
level $l$ is a PEL datum over $E_h$ with $p+q\geq1$, and the full level of the core is $1$, by the pins of
the laws; so Lemma~\ref{is-onto-step-move-laws-framed} applies. By its parts (1), (2), (4) and (5), every
framed point over an $F$-algebra has a move by $\gamma$; two moves of $y$ by $\gamma$ are images of each
other under the identity, at which sections take the same value; moves commute with images under
homomorphisms of $F$-algebras; and a move by $\delta$ of a move of $y$ by $\gamma$ is a move of $y$ by
$\gamma\delta$. Given these laws: (1) if $y=\varphi_*z$ with $z$
over a test ring, then $\gamma y=\varphi_*(\gamma z)$, and the value of $\gamma\phi$ at $y$ is the image
under $\varphi$ of the value of $\phi$ at $\gamma z$, which is the value of $\phi$ at $\gamma y$. (2) The
value of $(\gamma\delta)\phi$ at $y$ is the value of $\phi$ at $(\gamma\delta)y=\delta(\gamma y)$, which
is the value of $\delta\phi$ at $\gamma y$, that is, the value of $\gamma(\delta\phi)$ at $y$.
% Lean: evalAt_actL (ActLaws with EvalWD of the laws and frMap_trans_id), fromTest_mvPt (the field nat
% of ActLaws) and actL_mul (ActTotal, ActFun, ActMul and the naturality IsNatural of sections), with the
% laws T.actLev_full_level_group_framedPoints = FullLevelAction.actLev_levFull = actLaws_levFull
% (Lemma is-onto-step-move-laws-framed), which needs 1 <= p + q and A = 1 from the pins of the laws. The
% Lean's evalAt is 0 at points that are not base changes from test rings; the book does not use such
% points. Lemma is-onto-step-level-structure-cusp uses the same laws. The open parts of the laws are in
% Lemmas is-onto-step-move-laws-unique and is-onto-step-move-laws-moves.
\end{proof}

\begin{lemma}\label{is-onto-step-leaves-level-integral}
Let $X$ be the moduli scheme of a PEL datum over a CM field at a level, with a coefficient field $F$
(Definition~\ref{cv-realization-level-move}), let $(\rho,\nu)$ be a level automorphism of $X$, let $T$
be a complete discrete valuation ring with an embedding of $F$ into its fraction field, and let $y$ be a
framed point of $X$ over the fraction field of $T$ that is integral
(Situation~\ref{cv-definite-situation}(3)). Then the move of $y$ by $(\rho,\nu)$ is integral.
\end{lemma}
\leanref{Development/AbelianVarietyTheory/Proofs/JetHarmonics/IntegralFramedPointOfLevelMoveStatement.lean}{GlobalTheory.T.integral_framedPoint_of_level_move}
\leanref{Development/AbelianVarietyTheory/Proofs/JetHarmonics/IntegralFramedPointOfLevelMove.lean}{GlobalTheory.FormTies.LevelMove.integral_of_levR}
\begin{proof}
Let $y'$ be the move. That $y$ is integral gives an abelian scheme $G$ over $T$, with bases of the
eigenspaces of its Lie algebra and a generator $\theta$ of the sections of its line bundle, and an
identification $m$ of its generic fibre with the fibre $A_y$ of the universal abelian scheme at the point
of $y$, compatible with the actions of $\mathcal O_E$ and with the rigidified line bundles, that carries
the bases of $G$ to those of $y$ and $\theta$ to the generator of $y$. By
Definition~\ref{cv-realization-level-move}, $\nu$ restricts to an isomorphism $\nu_y\colon A_y\to A_{y'}$
compatible with the actions of $\mathcal O_E$, it comes with an isomorphism of the line bundle with its
pullback that respects the rigidifications, and the bases and the generator of $y'$ are the images of
those of $y$. So $\nu_y\circ m$ identifies the generic fibre of $G$ with $A_{y'}$ in the same way, and
$G$, with the same bases and generator, extends $y'$ over $T$.
% Lean: integral_of_levR: from levR, the move (rho, nu) with hx, hu, pm, w, hw, hs; from Integral, G, t',
% m1, m2, ub, theta, w1, w2, sigma, s1, theta2. The new map n = m1 composed with nu and its rigidified
% isomorphism w' come from exists_comp_rig, the section s' from thf_secUp; G, t', m2, ub (with the lifts
% moved by hn and hu), theta, w2, sigma and theta2 are kept, and the last clause is checked chart by
% chart (exists_chart, gs_appLE, gs_of_isoMaps). The Lean's integral framed point passes through an
% abelian variety t' over the fraction field with maps to both; the book's m is the composite. About
% 800 lines with the definitions below it.
\end{proof}

\begin{lemma}\label{is-onto-step-leaves-level-theta}
Let $s$ be a run at a point of a family that satisfies the laws, $l$ a level of a lift, $w$ a weight and
$\gamma$ an automorphism of the level module at level $l$. Then every theta series of $s$ of weight $w$
at level $l$ is a section of weight $w$, and the level action of $\gamma$ maps it to a theta series of
weight $w$.
\end{lemma}
\leanref{Development/AbelianVarietyTheory/Proofs/CuspFormEvaluations/ThetaRuleMemSec.lean}{GlobalTheory.thetaRule_mem_sec}
\leanref{Development/AbelianVarietyTheory/Proofs/JetHarmonics/LevelActionMemThetaSeriesStatement.lean}{GlobalTheory.T.levelAction_mem_thetaSeries}
\begin{proof}
\openstep{The proof that theta series are sections of their weight and that the level action carries
theta series to theta series, through the move of framed points and the independence of the theta
values from the formal chart, is not written out.}
% Lean: thetaRule_mem_sec (with GlobalTheory.abelianSchemeLaws_all, chart independence and the nabla pin
% h.nabPin of the laws) and T.levelAction_mem_thetaSeries = ThetaLevelMove.actLev_mem_th (with chart
% independence, GlobalTheory.chartIndependence_unconditional, cited in Lemma
% cv-uniformization-theta-rule-unique). Size, following imports and stopping at statements the book
% cites: about 6,300 lines below actLev_mem_th, 3,500 of them in theorems; about 700 below
% thetaRule_mem_sec.
\end{proof}

\begin{lemma}\label{is-onto-step-leaves-level-short}
Let $s$ be a run at a point of a family that satisfies the laws and is rational at a lift $j$, put
$l=j+1$, and let $w$ be a weight with $r\le n_w$ and $\gamma$ an automorphism of the level module at
level $l$ (Definition~\ref{is-onto-step-leaves-theta-module}). Then for every short series $\phi$ of
weight $w$ there are $u\in\mathcal O_F$ and a short series $\phi'$ of weight $w$ such that the level
action of $\gamma$ maps $\phi$ to $u\phi'$.
\end{lemma}
\leanref{Development/AbelianVarietyTheory/Proofs/JetHarmonics/LevelActionShortThetaSeriesRootOfUnityMulStatement.lean}{GlobalTheory.T.levelAction_shortThetaSeries_rootOfUnity_mul}
\begin{proof}
\openstep{The comparison of the level action on a short series with a short series at a moved vector,
with an integral scalar, is not written out.}
% Lean: T.levelAction_shortThetaSeries_rootOfUnity_mul = ShortMove.actLev_short: the same construction
% as exists_const (Lemma is-onto-step-leaves-move-const), from a completed isometry g inducing gamma,
% LevelActionLaws.exists_levAut, exists_root (u a root of unity in O_F, with the rationality and r <=
% n_w), levelCoverDatum_similitude_adjoint_transport, datum_move and the comparison of values with chart
% independence; f' is the short series at the moved vector. Size, following imports and stopping at
% statements the book cites: about 22,900 lines below actLev_short, 18,600 of them in theorems, nearly
% all shared with exists_const.
\end{proof}

\begin{lemma}\label{is-onto-step-leaves-level}
Let $s$ be a run at a point of a family that satisfies the laws and is rational at a lift $j$, and put
$l=j+1$. Let $\kappa,w_L,w_F$ be weights and $T\geq d_Y$ an integer with
$w^\pm_{F,\tau,i}=\kappa^\pm_{\tau,i}-Tw^\pm_{L,\tau,i}$ for all $\tau$ and $i$,
$n_\kappa=n_{w_F}+Tn_{w_L}$ and $r\le n_{w_F}$, and let $S$ be four elements of $\mathcal O_F$. Let
$\gamma$ be an automorphism of the level module of the core at level $l$ whose level action maps
$S_l(\kappa)$ into itself, and write $\gamma\phi$ for the image of $\phi\in S_l(\kappa)$. Then:
\begin{enumerate}
\item $\gamma\phi\in M_l(\kappa)$ for $\phi\in M_l(\kappa)$;
\item if $\iota_\kappa(\phi)\in\Theta(S)$, then $\iota_\kappa(\gamma\phi)\in\Theta(S)$.
\end{enumerate}
\end{lemma}
\leanref{Development/AbelianVarietyTheory/Proofs/JetHarmonics/LevActMemMlat.lean}{GlobalTheory.FormTies.levAct_mem_Mlat}
\leanref{Development/AbelianVarietyTheory/Proofs/JetHarmonics/ThetaStableOfAlgHom.lean}{GlobalTheory.FormTies.levAct_theta}
\leanref{Development/AbelianVarietyTheory/Proofs/JetHarmonics/ThetaStableOfAlgHom.lean}{GlobalTheory.FormTies.theta_stable_of_algHom}
\begin{proof}
(1) By the construction of the integral cusp forms, the points at a finite place $\lambda$ of $F$ in
Definition~\ref{is-onto-step-leaves-theta-module}(2) are the integral framed points $y$
(Situation~\ref{cv-definite-situation}(3)) over the fraction fields of complete discrete valuation rings
$T$ above $\lambda$ that are base changes of framed points over test rings, and
$v_{\lambda,y}(\iota_\kappa\phi)$ is the valuation of the value $\phi(y)$, normalized by the ramification
index of $T$ over $\lambda$. Let $\phi\in M_l(\kappa)$, and let $\lambda$, $T$ and $y$ be such. By
Lemma~\ref{is-onto-step-leaves-move-laws}(1), the move $\gamma y$ is again a base change of a framed point
over a test ring, and the value of $\gamma\phi$ at $y$ is the value of $\phi$ at $\gamma y$; by
Lemma~\ref{is-onto-step-leaves-level-integral}, $\gamma y$ is integral. Hence
$v_{\lambda,y}(\iota_\kappa(\gamma\phi))=v_{\lambda,\gamma y}(\iota_\kappa\phi)\geq0$ for all $\lambda$ and
$y$, and $\gamma\phi\in M_l(\kappa)$.

(2) The level action of $\gamma$ on the sections at level $l$ is an $F$-algebra homomorphism that
respects the weights. Applied to the components of the elements of $A_l$, it gives an $F$-algebra
endomorphism $\Phi$ of $A_l$ with $\Phi(\iota_w\phi)=\iota_w(\gamma\phi)$ for every weight $w$ and
$\phi\in S_l(w)$, where $\gamma\phi\in S_l(w)$ by Lemma~\ref{is-onto-step-level-structure-cusp}. The level
action of $\gamma$ (a) maps each theta series of weight $w_L$ to a theta series, by
Lemma~\ref{is-onto-step-leaves-level-theta}, and (b) maps each short series of weight
$w\in\{w_F,w_{F,0},\kappa\}$ to $u$ times a short series of weight $w$, for some $u\in\mathcal O_F$, by
Lemma~\ref{is-onto-step-leaves-level-short}, which applies as $r\le n_{w_F}\le n_{w_{F,0}}$ and
$r\le n_{w_F}\le n_\kappa$. Hence it maps the $\mathcal O_F$-span of the theta series of weight $w_L$, and that
of the short series of each of these weights, into itself, and so $\Phi$ maps
$\iota_{w_L}(\mathrm{Th}_l(w_L))$ and the three modules $\iota_w(\mathrm{Sh}_l(w))$ into themselves. As
$\Phi$ is a homomorphism of $\mathcal O_F$-algebras, $\Phi(aX)\subset aX$ when $\Phi(X)\subset X$, and
$\Phi$ carries products and powers of submodules to the products and powers of their images. So $\Phi$
maps each summand of $\Theta(S)$ into itself, and
$\iota_\kappa(\gamma\phi)=\Phi(\iota_\kappa\phi)\in\Theta(S)$.
% Lean: (1) is levAct_mem_Mlat: Mlat is latOf of the valuations recVals (PtV: a place la, a complete
% discrete valuation ring S over la with ramification e (OverLa), an IntFr y, that is FrPt with FromTest
% and Integral, and nval S.V e of evI y); for such a point, y' = mvPt (l g) y with fromTest_mvPt and
% T.integral_framedPoint_of_level_move (with mvPt_spec), and evalAt_actL with hact. (2) is levAct_theta:
% Phi = AddMonoidAlgebra.mapAlgHom of actLevAlg, hspan with T.exists_levelChangeAction_on_cuspForms and
% map_mem_spanIn, (a) is T.levelAction_mem_thetaSeries with thetaRule_mem_sec and chart independence,
% (b) is T.levelAction_shortThetaSeries_rootOfUnity_mul (it uses r <= n_w and the rationality xi), and
% theta_stable_of_algHom (about 45 lines) is the last step. The open parts are in Lemmas
% is-onto-step-leaves-move-laws, -level-theta and -level-short.
\end{proof}

\begin{lemma}\label{is-onto-step-leaves-form-tie}
Let $O$ be a record object for an indefinite-side datum over a number field $L$, a finite group $P$
and constants $c$ (Definition~\ref{is-record-object}), with the notation there; let $v$ be a finite place
of $L$ and $j$ a lift of the datum. Let $s$ be a run at a point of a family that satisfies the laws and
is covered and rational at $j$, put $l=j+1$, let $\kappa,w_L,w_F$ and $T=T_O$ be the weights and the
integer of the parameter choices at $l$, $p$ a prime and $S$ the scalars of $p$
(Definition~\ref{is-onto-step-leaves-theta-module}), and suppose that the coefficient field $F$ of the
core is a subfield of $L$. Let $z\colon L^{n_Z(j)}\to L\otimes_FS_l(\kappa)$ be an $L$-linear
isomorphism, write $X^z$ for the $\mathcal O_L$-span of the vectors $z^{-1}(1\otimes\phi)$ with $\phi$
in an $\mathcal O_F$-submodule $X$ of $S_l(\kappa)$, and let $\ell$ be a map from $P$ to the
automorphisms of the level module at level $l$ whose level actions map $S_l(\kappa)$ into itself.
Suppose that:
\begin{enumerate}
\item $M_s(j)=M_l(\kappa)^z$, $\Theta(v,j)=\{\phi\in S_l(\kappa):\iota_\kappa(\phi)\in\Theta_l(p)\}^z$,
  $\Theta_K(v,j)=\mathrm{Sh}_l(\kappa)^z$ and $\sigma(v,j)=\sigma_l$;
\item $\rho_j(g)z^{-1}(1\otimes\phi)=z^{-1}(1\otimes\ell(g)\phi)$ for $g\in P$ and $\phi\in S_l(\kappa)$;
\item $M_l(\kappa)$ is a lattice in $S_l(\kappa)$, $\mathrm{Sh}_l(\kappa)$ spans $S_l(\kappa)$ over $F$, and
  the run satisfies (H2) for $S$.
\end{enumerate}
Then $M_s(j)$ is a lattice in $L^{n_Z(j)}$ stable under $\rho_j(P)$; $\Theta(v,j)$ and $\Theta_K(v,j)$
are lattices in $L^{n_Z(j)}$; and $\Theta(v,j)\subset M_s(j)$, $\Theta(v,j)$ is stable under
$\rho_j(P)$, $\sigma(v,j)\ne0$ and $\Theta(v,j)\subset\sigma(v,j)\Theta_K(v,j)$.
\end{lemma}
\leanref{Development/AbelianVarietyTheory/Proofs/JetHarmonics/LeavesSupSpanLemmaSOfFormTieK.lean}{GlobalTheory.FormTies.leaves_sup_span_lemmaS_of_formTieK}
\leanref{Development/AbelianVarietyTheory/Proofs/JetHarmonics/SupLatticeIsLatInStable.lean}{GlobalTheory.FormTies.supLattice_isLatIn_stable}
\leanref{Development/AbelianVarietyTheory/Proofs/JetHarmonics/ThetaSpansIsLatIn.lean}{GlobalTheory.FormTies.theta_spans_isLatIn}
\leanref{Development/AbelianVarietyTheory/Proofs/JetHarmonics/ScaledSpanClauses.lean}{GlobalTheory.scaledSpan_clauses}
\begin{proof}
Put $\Theta_0=\{\phi\in S_l(\kappa):\iota_\kappa(\phi)\in\Theta_l(p)\}$ and $c_0=s_Kg^{frn_\kappa}$. By
Lemma~\ref{is-onto-step-leaves-weights}, $\kappa,w_L,w_F$ and $T$ satisfy the hypotheses of
Lemmas~\ref{is-onto-step-leaves-integral}, \ref{is-onto-step-leaves-scalar},
\ref{is-onto-step-leaves-products} and~\ref{is-onto-step-leaves-level}. By
Lemma~\ref{is-onto-step-leaves-short}, $\mathrm{Sh}_l(\kappa)$ is finitely generated; by
Lemma~\ref{is-onto-step-leaves-products}, $\Theta^{\mathrm p}_l\subset\iota_\kappa(\mathrm{Sh}_l(\kappa))$;
and by Lemma~\ref{is-onto-step-leaves-integral} and the injectivity of $\iota_\kappa$,
$\Theta_0\subset M_l(\kappa)$.

For $\mathcal O_F$-submodules $X\subset Y$ of $S_l(\kappa)$, $X^z\subset Y^z$; $X^z$ is finitely
generated if $X$ is, being generated by the $z^{-1}(1\otimes\phi)$ for $\phi$ in a set of generators of
$X$; and $X^z$ spans $L^{n_Z(j)}$ if $X$ spans $S_l(\kappa)$, since then the $1\otimes\phi$ span
$L\otimes_FS_l(\kappa)$.

\emph{Supremum.} $M_l(\kappa)$ is finitely generated and spans, so $M_s(j)=M_l(\kappa)^z$ is a lattice.
For $g\in P$, $\rho_j(g)$ maps each generator $z^{-1}(1\otimes\phi)$, $\phi\in M_l(\kappa)$, to
$z^{-1}(1\otimes\ell(g)\phi)$, and $\ell(g)\phi\in M_l(\kappa)$ by
Lemma~\ref{is-onto-step-leaves-level}(1). So $\rho_j(g)M_s(j)\subset M_s(j)$.

\emph{Span.} $\Theta(v,j)=\Theta_0^z\subset M_s(j)$ is finitely generated, as $\mathcal O_L$ is
noetherian. It spans: for $\phi\in\mathrm{Sh}_l(\kappa)$, $c_0\phi\in\mathrm{Sh}^{s_K}_l(\kappa)$, so
$\iota_\kappa(c_0\phi)$ lies in the third summand of $\Theta_l(p)$ and $c_0\phi\in\Theta_0$; as
$c_0\ne0$ (Lemma~\ref{is-onto-step-leaves-scalar}), $\phi$ lies in the $F$-span of $\Theta_0$, which is
therefore $S_l(\kappa)$. And $\Theta_K(v,j)=\mathrm{Sh}_l(\kappa)^z$ is finitely generated and spans, as
$\mathrm{Sh}_l(\kappa)$ is finitely generated and spans.

\emph{Containment.} $\Theta(v,j)=\Theta_0^z\subset M_l(\kappa)^z=M_s(j)$. For $g\in P$ and
$\phi\in\Theta_0$, $\ell(g)\phi\in\Theta_0$ by Lemma~\ref{is-onto-step-leaves-level}(2), so $\rho_j(g)$
maps the generators of $\Theta(v,j)$ into $\Theta(v,j)$. The element $\sigma(v,j)$ is the image in $L$
of $\sigma_l\ne0$ (Lemma~\ref{is-onto-step-leaves-scalar}), so it is nonzero. Let $\phi\in\Theta_0$. By
Lemma~\ref{is-onto-step-leaves-scalar}, $\iota_\kappa(\phi)=\sigma_lx$ with $x\in\Theta^{\mathrm p}_l$, so
$\iota_\kappa(\sigma_l^{-1}\phi)=x\in\iota_\kappa(\mathrm{Sh}_l(\kappa))$ and
$\sigma_l^{-1}\phi\in\mathrm{Sh}_l(\kappa)$. Hence
$z^{-1}(1\otimes\phi)=\sigma(v,j)\,z^{-1}(1\otimes\sigma_l^{-1}\phi)\in\sigma(v,j)\Theta_K(v,j)$.
% Lean: leaves_sup_span_lemmaS_of_formTieK, key at (v, i): RunDepth.run_weights_sum (h1-h4), the
% FormTie from FormTieK.toFormTie with shortTh_fg and productsShort_rec (ThK_fg, and ThK_ge: the z-span
% of the preimage of thetaPlain lies in ThK), then supLattice_isLatIn_stable (spanK_isLatIn, rho_eq and
% levAct_mem_Mlat), theta_spans_isLatIn (theta_le_msup for FG, SpanAt with c = sK (g0 g1)^(f r n) for
% the span of Th, ThK_fg and ThK_ge for ThK) and scaledSpan_clauses (theta_le_msup, rho_eq with
% levAct_theta, sigma_eq, and theta_scaled_span_le_sigma_smul_unscaled_span with ThK_ge). The Lean
% reaches the span of ThK through ThK_ge; with ThK_eq the two are the same.
\end{proof}

\begin{lemma}\label{is-onto-step-leaves-ties}
Let $p$ be a point of a family, $r$ a pair of runs at $p$ and $O$ a record object recorded at $(p,r)$,
with the notation of Definitions~\ref{is-record-object} and~\ref{is-onto-record-numbers} and the map
$\mathrm{inc}\colon P\to G$ of Lemma~\ref{is-ties-level-tie}. Let $v$ be a finite place of $L$, $s$ the
run used at $v$, $j$ a lift and $l=j+1$; let $z_j$ be the isomorphism of
Lemma~\ref{is-onto-step-index-bound-ties}(1), with the notation $X^z$ there, and put $\kappa=\kappa'_l$.
Then:
\begin{enumerate}
\item for every $g\in P$ the level action of $\ell_j(\mathrm{inc}(g))$ maps $S_l(\kappa)$ into itself,
  and $\rho_j(g)z_j^{-1}(1\otimes\phi)=z_j^{-1}\big(1\otimes\ell_j(\mathrm{inc}(g))\phi\big)$ for
  $\phi\in S_l(\kappa)$;
\item $\Theta_K(v,j)=\mathrm{Sh}_l(\kappa)^z$ and $\sigma(v,j)=\sigma_l$;
\item $\Theta_0(v,j)$ is the $\mathcal O_L$-span of the vectors
  $\sum_{u\in U}\psi(u)\,z_j^{-1}(1\otimes\phi_u)$, for the orbit families $(\phi_u)_{u\in U}$ of weight
  $\kappa$ for $\ell_j\circ\mathrm{inc}$ whose members lie in $S_l(\kappa)$
  (Definition~\ref{is-onto-step-leaves-theta-module}(8)).
\end{enumerate}
\end{lemma}
\leanref{Development/AbelianVarietyTheory/Definitions/Foundations/ThetaSeriesTies.lean}{GlobalTheory.OnTies.thetaEig}
\leanref{Development/AbelianVarietyTheory/Definitions/Jets/EmbShortThLeThetaPlain.lean}{GlobalTheory.FormTies.IsLevAct}
\begin{proof}
These are conditions of the level tie of $O$ at $(v,j)$, which is part of being recorded, as in the
proof of Lemma~\ref{is-onto-step-index-bound-ties}: (1) is its action condition with the condition on
$\rho_j$ of its form tie, (2) is the rest of the form tie, and (3) is its condition on $\Theta_0$.
% Lean: IsRec.tied gives OnObjLevelCompatibility, whose OnTie at (v, j) has l = r.lev j and
% inc = mirIncl R; (1) is OnTie.act_eq (IsLevAct for fun g => l (inc g)) with FormTieK.rho_eq, (2) is
% FormTieK.ThK_eq and sigma_eq (sigRec), and (3) is OnTie.Th0_eq (thetaEig, the traceSpan over
% IsOrbitFamily with T.dd (j + 1) (kappa j).n).
\end{proof}

\begin{lemma}\label{is-onto-step-leaves-move-isometry}
Let $E$ be a CM field, with complex conjugation $a\mapsto\bar a$. Let $d\equiv7\pmod{16}$ be a prime,
$\sqrt{-d}$ and $w_0$ elements of $\mathcal O_E$ with $(\sqrt{-d})^2=-d$ and $2w_0=1+\sqrt{-d}$, and
$\theta\in\mathbf C$ with $\theta^2=-d$, and suppose that no finite prime of the totally real subfield of
$E$ ramifies in $E$. Let $\Lambda$ be a finitely generated projective $\mathcal O_E$-module of rank $p+q$
with an alternating $\mathbf Z$-bilinear pairing $\langle\,,\rangle\colon\Lambda\times\Lambda\to\mathbf Z$
that identifies $\Lambda$ with $\operatorname{Hom}(\Lambda,\mathbf Z)$, satisfies $\langle
ax,y\rangle=\langle x,\bar ay\rangle$ for $a\in\mathcal O_E$, and has a complex structure of signature
$(p,q)$ for $\theta$ (Definition~\ref{cv-cm-point-periods}(5)). Let $(M,M_b,A)$ be a level of this
datum (Definition~\ref{cv-model-toroidal-semiabelian}(1), with $\mathfrak M=M_b$ and full level $A$):
an odd integer $M\geq3$, an ideal $M_b$ of $\mathcal O_E$ generated by elements of $\mathbf Z[w_0]$ with
$M_b\bar M_b=M\mathcal O_E$ and $M_b+\bar M_b=\mathcal O_E$, and an integer $A\geq1$. Let $\gamma$ be an
$\mathcal O_E$-linear automorphism of $\Lambda/M_b\Lambda$. Write $a\equiv b\pmod n$ for $a-b\in
n\Lambda$. Then there are $\mathcal O_E$-linear endomorphisms $F_n$ of $\Lambda$ and integers $\nu_n$,
one of each for every integer $n\geq1$, such that:
\begin{enumerate}
\item $F_{nm}(a)\equiv F_n(a)\pmod n$ for all $n,m\geq1$ and $a\in\Lambda$;
\item for every $n$ there is an endomorphism $G$ of $\Lambda$ with $G(F_n(a))\equiv a$ and
  $F_n(G(a))\equiv a\pmod n$ for every $a$;
\item $\langle F_na,F_nb\rangle\equiv\nu_n\langle a,b\rangle\pmod n$ for all $n$ and $a,b\in\Lambda$;
\item $\nu_n=1$ for every $n$;
\item $F_M$ induces $\gamma$ on $\Lambda/M_b\Lambda$;
\item $F_N(a)\equiv a\pmod N$ for every $N\geq1$ prime to $M$ and every $a$.
\end{enumerate}
By (1) to (3), $(F_n,\nu_n)_{n\geq1}$ is a family of similitudes of $\Lambda$
(Definition~\ref{cv-toroidal-moduli-object}(1)). Of the hypotheses on $d$, $\theta$, $E$, $\Lambda$ and
the level, the proof uses only that $d>0$, $(\sqrt{-d})^2=-d$, $2w_0=1+\sqrt{-d}$, that $\Lambda$ is
finitely generated, the three properties of $\langle\,,\rangle$, $M>0$, $M_b\bar M_b=M\mathcal O_E$ and
$M_b+\bar M_b=\mathcal O_E$; the others are not used.
\end{lemma}
\leanref{Development/AbelianVarietyTheory/Proofs/JetHarmonics/ExistsCompletedIsometryInducingHalfLevelAutomorpStatement.lean}{GlobalTheory.T.exists_completed_isometry_inducing_half_level_automorphism}
\leanref{Development/AbelianVarietyTheory/Proofs/JetHarmonics/ExistsCompletedIsometryInducingHalfLevelAutomorp.lean}{GlobalTheory.LevelIsometry.exists_isometry_mod}
\leanref{Development/AbelianVarietyTheory/Proofs/JetHarmonics/ExistsCompletedIsometryInducingHalfLevelAutomorp.lean}{GlobalTheory.LevelIsometry.exists_sim_trivial_away}
\leanref{Development/AbelianVarietyTheory/Proofs/ModuliProofs/ModNOfDvdModulus.lean}{GlobalTheory.SimLift.step_sq}
\leanref{Development/AbelianVarietyTheory/Proofs/ModuliProofs/ModNOfDvdModulus.lean}{GlobalTheory.SimLift.exists_inv}
\begin{proof}
Call an endomorphism $F$ of $\Lambda$ an \emph{isometry modulo} $n$ if $\langle
Fa,Fb\rangle\equiv\langle a,b\rangle\pmod n$ for all $a,b$. As the pairing identifies $\Lambda$ with
$\operatorname{Hom}(\Lambda,\mathbf Z)$, an element $x$ with $\langle x,b\rangle\in n\mathbf Z$ for every
$b$ lies in $n\Lambda$; and for every biadditive $B\colon\Lambda\times\Lambda\to\mathbf Z$ with
$B(ca,b)=B(a,\bar cb)$ for $c\in\mathcal O_E$ there is an $\mathcal O_E$-linear $u$ with $\langle
ua,b\rangle=B(a,b)$. Every complex embedding $\tau$ of $E$ has $\tau(\sqrt{-d})^2=-d<0$, so
$\tau(\sqrt{-d})$ is purely imaginary and complex conjugation sends $\sqrt{-d}$ to $-\sqrt{-d}$. With
$2w_0=1+\sqrt{-d}$ this gives $2\bar w_0=1-\sqrt{-d}=2-2w_0$, so $w_0+\bar w_0=1$.

\emph{Step 1: an isometry modulo $M$ inducing $\gamma$.} Write $1=e'+e$ with $e'\in M_b$ and
$e\in\bar M_b$, and let $f$ and $g$ be endomorphisms of $\Lambda$ inducing $\gamma$ and $\gamma^{-1}$, so
that $g(f(b))-b\in M_b\Lambda$. Let $h$ be the endomorphism with $\langle ha,b\rangle=\langle a,gb\rangle$,
and put $F=ef+e'h$. Then $F-f=e'(h-f)$ maps $\Lambda$ into $M_b\Lambda$, so $F$ induces $\gamma$. As $\bar
ee$ and $\bar e'e'$ lie in $M_b\bar M_b=M\mathcal O_E$, the numbers $\langle efa,efb\rangle=\langle
fa,\bar eefb\rangle$ and $\langle e'ha,e'hb\rangle$ are divisible by $M$. Next
\[
 \langle e'ha,efb\rangle=\langle ha,\bar e'efb\rangle=\langle a,\bar e'e\,g(fb)\rangle
 \equiv\langle a,\bar e'eb\rangle\pmod M,
\]
since $e(g(fb)-b)\in\bar M_bM_b\Lambda=M\Lambda$. Exchanging $a$ and $b$ and using that the pairing is
alternating, $\langle efa,e'hb\rangle\equiv\langle a,e'\bar eb\rangle\pmod M$. So $\langle
Fa,Fb\rangle\equiv\langle a,(\bar e'e+e'\bar e)b\rangle\pmod M$, and $\bar e'e+e'\bar e=1-\bar ee-\bar
e'e'$, as $(e+e')(\bar e+\bar e')=1$. Hence $F$ is an isometry modulo $M$.

\emph{Step 2: lifting.} Let $F$ be an isometry modulo $n$, and write $\langle Fa,Fb\rangle=\langle
a,b\rangle+nB(a,b)$. Then $B$ is alternating and $B(ca,b)=B(a,\bar cb)$, so $B(a,b)=\langle ua,b\rangle$ for
an endomorphism $u$, and $\langle a,ub\rangle=-B(b,a)=\langle ua,b\rangle$. Put $v=w_0u$. Then $\langle
a,vb\rangle=\langle\bar w_0ua,b\rangle$, so $\langle va,b\rangle+\langle a,vb\rangle=B(a,b)$. The endomorphism
$F'=F-nFv$ satisfies $F'\equiv F\pmod n$ and, since $\langle Fx,Fy\rangle\equiv\langle x,y\rangle\pmod n$,
\[
 \langle F'a,F'b\rangle\equiv\langle Fa,Fb\rangle-n\big(\langle va,b\rangle+\langle a,vb\rangle\big)
 =\langle a,b\rangle\pmod{n^2}.
\]

\emph{Step 3: the powers of $M$.} Starting from the map $F$ of Step 1 and applying Step 2 repeatedly,
there are isometries $s_k$ modulo $M^{k+1}$, $k\geq0$, with $s_0=F$ and $s_{k+1}\equiv s_k\pmod{M^{k+1}}$.
Hence $s_i\equiv s_j$ modulo every common divisor of $M^{i+1}$ and $M^{j+1}$.

\emph{Step 4: all moduli.} Every $n\geq1$ is a product $n=ab$ with $a\mid M^{k+1}$ for some $k$ and $b$
prime to $M$. Choose integers $\alpha,\beta$ with $\alpha a+\beta b=1$ and put $F_n=\beta b\,s_k+\alpha a$.
Then $F_n\equiv s_k\pmod a$ and $F_n\equiv1\pmod b$, so $F_n$ is an isometry modulo $a$ and modulo $b$,
hence modulo $n$; this is (3) with $\nu_n=1$, and (4). For a divisor $t$ of $n$: if $t$ divides some
$M^{i+1}$, then $F_n\equiv s_i\pmod t$; if $t$ is prime to $M$, then $F_n\equiv1\pmod t$. This gives (6),
and (5), since $F_M\equiv s_0=F\pmod M$ and $M\Lambda\subset M_b\Lambda$. For (1), $F_{nm}$ and $F_n$ are both congruent to $s_k$
modulo $a$ and to $1$ modulo $b$, and $a$ and $b$ are coprime. For (2), let $G$ be the endomorphism with
$\langle Ga,b\rangle=\langle a,F_nb\rangle$. Then $\langle GF_na,b\rangle=\langle F_na,F_nb\rangle\equiv\langle
a,b\rangle\pmod n$ for every $b$, so $GF_n\equiv1\pmod n$. Hence $G$ induces a surjective endomorphism of
the finitely generated $\mathcal O_E$-module $\Lambda/n\Lambda$, which is therefore bijective, with inverse
induced by $F_n$; so also $F_nG\equiv1\pmod n$.
% Lean: T.exists_completed_isometry_inducing_half_level_automorphism = LevelIsometry.exists_isometry_inducing.
% Step 1 is exists_isometry_mod (e' + e = 1 from l.coprime, f = liftLa gamma, g = liftLa gamma^-1, h from
% SimLift.exists_rep, key for the cross terms, m1 and m2); Steps 2 to 4 are exists_sim_trivial_away, with
% exists_tower (SimLift.step_sq, where v = H.w0 u, H.w0 the half-integer generator of the host),
% exists_split, exists_glue and modN_mul_of_coprime; (2) is SimLift.exists_inv (the adjoint, and
% OrzechProperty.injective_of_surjective_endomorphism). The Sim has nu n = 1 for every n (fun _ => 1);
% step_sq is stated for any multiplier nu prime to n, and is used here with nu = 1. w0 + bar w0 = 1 is
% HostModel.cmHost_bar_halfIntegerGenerator (bar_rt from conj_embedding_rt, which uses d_prime.pos and
% rt_sq). About 340 lines.
\end{proof}

\begin{lemma}\label{is-onto-step-leaves-move-character}
Let $R$ be a finite commutative ring whose residue fields all have at least $4$ elements, $r\geq0$ an
integer, and $\chi$ a homomorphism from $GL_r(R)$ to a commutative group. Then $\chi(g)=1$ for every
$g\in GL_r(R)$ with $\det g=1$.
\end{lemma}
\leanref{Development/AbelianVarietyTheory/Proofs/UnitaryExistence/GeneralLinearGroupCharEqOneOfDetEqOne.lean}{GlobalTheory.generalLinearGroup_char_eq_one_of_det_eq_one}
\leanref{Development/AbelianVarietyTheory/Proofs/UnitaryExistence/GeneralLinearGroupCharEqOneOfDetEqOne.lean}{GlobalTheory.ddChar_induction}
\begin{proof}
Write $e_a$ for the standard basis vectors of $R^r$, viewed as columns, and $y^{\mathrm t}$ for the
transpose.

\emph{Step 1: a unit $u$ with $u-1$ a unit.} Each residue field has an element other than $0$ and $1$.
By the Chinese remainder theorem there is $u\in R$ that is congruent modulo each maximal ideal to such
an element. Then $u$ and $u-1$ lie in no maximal ideal, so they are units.

\emph{Step 2: unipotents.} Let $N$ be a matrix with $N^2=0$, and let $d\in GL_r(R)$ and $p,q\in R$ satisfy
$dN=pN$, $Nd=qN$, with $q-p$ a unit. Then $\chi(1+N)=1$. Indeed, put $U(a)=1+a(q-p)^{-1}N$ for $a\in R$.
Then $U(a)U(a')=U(a+a')$ and $dU(q)=U(p)d$. As the target of $\chi$ is commutative,
$\chi(U(q))=\chi(U(p))$; and $U(q)=U(p)U(q-p)$, so $\chi(1+N)=\chi(U(q-p))=1$.

\emph{Step 3: elementary matrices.} Let $a$ be an index and $d$ the diagonal matrix with $u$ at $a$ and
$1$ elsewhere. For a row vector $y$ with $y_a=0$, $N=e_ay^{\mathrm t}$ has $N^2=0$, $dN=uN$ and $Nd=N$;
for a column vector $x$ with $x_a=0$, $N=xe_a^{\mathrm t}$ has $N^2=0$, $dN=N$ and $Nd=uN$. In both cases
$\pm(1-u)$ is a unit, so $\chi(1+N)=1$ by Step 2. These matrices have determinant one.

\emph{Step 4: diagonal matrices.} Conjugation by the permutation matrix of a transposition shows that
$\chi$ takes the same value on the diagonal matrices with a unit $v$ at the place $i$, and $1$ elsewhere,
for all $i$. A diagonal matrix of determinant one is the product of such matrices for the units $v_i$ on
its diagonal, so its value under $\chi$ is $\chi$ of the matrix with $\prod_iv_i=1$ at one place, which is
$1$.

\emph{Step 5: induction.} We show by induction on a set $s$ of indices: if $\det g=1$ and $g_{ij}=0$ for
$i\ne j$ with $i\notin s$ or $j\notin s$, then $\chi(g)=1$. For $s=\emptyset$ this is Step 4. Let $g$
satisfy the condition for $s\cup\{a\}$, $a\notin s$. The entry $(a,a)$ of $g^{-1}g=1$ gives $wv+b=1$ with
$v=g_{aa}$, $w=(g^{-1})_{aa}$ and $b=\sum_{k\in s}(g^{-1})_{ak}g_{ka}$. Choosing $t$ modulo each maximal
ideal and using the Chinese remainder theorem, there is $t\in R$ with $v+tb$ a unit. Multiplying $g$ on the
left by $1+e_ay^{\mathrm t}$, with $y_k=t(g^{-1})_{ak}$ for $k\in s$ and $y_k=0$ otherwise, makes the entry
$(a,a)$ equal to $C=v+tb$ and keeps the vanishing condition. A left multiplication by $1+xe_a^{\mathrm t}$
with $x_i=-g_{ia}C^{-1}$ for $i\ne a$ then clears the column $a$ off the diagonal, and a right
multiplication by $1+e_ay^{\mathrm t}$ with $y_j=-C^{-1}g_{aj}$ for $j\ne a$ clears the row $a$. By
Step 3 these factors do not change $\chi(g)$ or $\det g$, and the result satisfies the condition for $s$.
Taking $s$ to be all indices gives the lemma.
% Lean: generalLinearGroup_char_eq_one_of_det_eq_one: Step 1 is ddChar_exists_unit (ddChar_crt,
% ddChar_isUnit), Step 2 ddChar_unip_eq_one, Step 3 ddChar_row and ddChar_col (with ddChar_det_unip),
% Step 4 ddChar_diag_single and ddChar_diag, Step 5 ddChar_induction with ddChar_stable for t. About
% 400 lines.
\end{proof}

\begin{lemma}\label{is-onto-step-leaves-move-det-character}
Let $E$, $\Lambda$ and the level $(M,M_b,A)$ be as in Lemma~\ref{is-onto-step-leaves-move-isometry}, with
$M$ prime to $3$, and let $\chi$ be a homomorphism from the group of $\mathcal O_E$-linear automorphisms of
$\Lambda/M_b\Lambda$ to a commutative group. Let $\gamma$ be such an automorphism, induced by an
endomorphism $\tilde\gamma$ of $\Lambda$ with $\det(\tilde\gamma\otimes\mathcal O_E/M_b)=1$. Then
$\chi(\gamma)=1$.
\end{lemma}
\leanref{Development/AbelianVarietyTheory/Proofs/FoundationProofs/TwistTrivialOnRational.lean}{GlobalTheory.TraceLattice.char_eq_one_of_detOne}
\leanref{Development/AbelianVarietyTheory/Proofs/UnitaryExistence/DetCharacterReductionInvariant.lean}{GlobalTheory.char_eq_one_of_inverse_reduction_det_one}
\begin{proof}
The ring $R=\mathcal O_E/M_b$ is finite, as $M\in M_b$. A residue field of $R$ is a residue field of
$\mathcal O_E$ at a prime containing $M$, so its characteristic divides $M$; as $M$ is odd and prime to
$3$, it is neither $\mathbf F_2$ nor $\mathbf F_3$, and it has at least $4$ elements. The $R$-module
$\Lambda/M_b\Lambda=R\otimes\Lambda$ is projective and has the same rank $r$ at every prime of $R$, so it
is free. In a basis, the automorphisms of $\Lambda/M_b\Lambda$ become $GL_r(R)$, and $\gamma$ becomes the
matrix of $\tilde\gamma\otimes R$, of determinant one. Lemma~\ref{is-onto-step-leaves-move-character}
gives $\chi(\gamma)=1$.
% Lean: char_eq_one_of_detOne (h4: a residue field with 2 or 3 elements would have order prime to M,
% though M = 0 in it) applies char_eq_one_of_inverse_reduction_det_one with t = id to gamma^-1
% (free_quotient_tensor_of_projective, char_eq_one_of_inverse_det_one_of_basis), and inverts. The Lean
% asks M >= 3 odd and not divisible by 3, as the level does. About 550 lines with the previous lemma.
\end{proof}

\begin{lemma}\label{is-onto-step-leaves-move-const}
Let $s$ be a run at a point of a family that satisfies the laws and is rational at a lift $j$, put
$l=j+1$, let $F$ be the coefficient field of the core, and let $w$ be a weight with $r\le n_w$
(Definition~\ref{is-onto-step-leaves-theta-module}). For every automorphism $\gamma$ of the level module at
level $l$ there is $u(\gamma)\in F^\times$ with the following property. For every series datum
$(\Lambda',N,P)$ of weight $w$, all allowed vectors $x$ and $x'$ such that $x$ is a right multiple of $x'$
by $\gamma$, and short series $\phi$ and $\phi'$ of $(\Lambda',N,P)$ at $x$ and $x'$, the level action of
$\gamma$ maps $\phi$ to $u(\gamma)\phi'$.
\end{lemma}
\leanref{Development/AbelianVarietyTheory/Proofs/FoundationProofs/TwistTrivialOnRational.lean}{GlobalTheory.TraceLattice.exists_const}
\begin{proof}
\openstep{The comparison of the level action of an automorphism on short series at right multiples,
through the move of the framed points, the transport of the delta data and the independence of the
Taylor values from the formal chart, is not written out.}
% Lean: exists_const: a completed isometry g inducing gamma (Lemma is-onto-step-leaves-move-isometry),
% LevelActionLaws.exists_levAut for a level automorphism rho2 of the delta data, exists_root (with the
% rationality hxi and r <= n_w) for u /= 0 and the coefficient relation hcd, then
% T.levelCoverDatum_similitude_adjoint_transport (sigma_q), datum_move (the delta weights at x and x'),
% mvPt_spec, exists_unique_frMove, T.exists_conj_levelMap_by_levelMoves and
% T.liesOverFramed_of_square_levelMoves for the framed points, and
% ShortLiftInvariance.cover_product_indep_of_lift and T.shortThetaValue_under_levelMove with chart
% independence (GlobalTheory.chartIndependence_unconditional, cited in Lemma
% cv-uniformization-theta-rule-unique). Size, following imports and stopping at statements the book
% cites: about 23,000 lines below exists_const, 19,300 of them in theorems.
\end{proof}

\begin{lemma}\label{is-onto-step-leaves-move-last}
Let $p$ be a point of a family, with definite-side datum of ranks $k,m,t$ and Harish-Chandra parameters
$e_j(\tau),b_j(\tau)$ at the embeddings $\tau$ (Definition~\ref{th-hc-parameter}) that has
Harish-Chandra weight blocks (Definition~\ref{is-onto-step-leaves-kappa}), and parameter choices with
integer $T_O$, and let $l$ be a level of a lift. Then for every $\tau$ the last of the $k+l-1$ entries
$\kappa'^{+}_{l,\tau,i}$ of the weight $\kappa'_l$ is at least $2T_O^2$.
\end{lemma}
\leanref{Development/AbelianVarietyTheory/Counterexamples/JetSurvey/OrientationExclusive.lean}{GlobalTheory.DepthBoundProbe.kappa_last}
\leanref{Development/AbelianVarietyTheory/Proofs/FoundationProofs/LeafDataOfOnTieOfIntegral.lean}{GlobalTheory.TraceLattice.kaR_lastPos}
\begin{proof}
By Definition~\ref{is-onto-step-leaves-kappa}, the $i$-th positive entry of $\kappa'_l$ at $\tau$ is
$\max(a^+_{l,\tau,i},0)$ with $a^+_{l,\tau,i}=\sigma_i+e_l-\frac12(n_V+1-2i)-\frac12\big((k+l-1)-(m-l)\big)$,
where $\sigma_i=v_i$ for $i<k$ and $\sigma_i=-b_{i-k+1}$ for $k\le i<k+m$. At the last index $i=k+l-1$,
$a^+_{l,\tau,i}=e_l-b_l-(t+1)/2$, as $n_V=k+m-1+t$. The parameter choices satisfy $|e_l-b_l-h_0/2|\le2$ and $4T_O^2+2t\le h_0$ (as in (10)
of Lemma~\ref{asm-dependence-functional-witness-numbers-parameters-computations} and in
Lemma~\ref{asm-dependence-functional-witness-numbers-blocks-computations}), and $t\geq5$, as
$r+t=8fr+4+T_On_W$. So $a^+_{l,\tau,k+l-1}\geq2T_O^2+t-2-(t+1)/2\geq2T_O^2>0$, and the entry, which is
$a^+_{l,\tau,k+l-1}$, is at least $2T_O^2$.
% Lean: kappa_last: kaWt_big and the KapDef (aPlusQ with sigPar = -b and eQ = e at the last index),
% P.e_centred and Tp.heightBound_ge_order_sq, ht from Tp.split, nP_eq, TO_nW and nW_40; kaR_lastPos
% deduces 1 <= the entry from T_O >= 40. About 320 lines.
\end{proof}

\begin{lemma}\label{is-onto-step-leaves-move}
Let $s$ be a run at a point of a family that satisfies the laws and is rational at a lift $j$, put
$l=j+1$, and let $w$ be a weight with $r\le n_w$ and $(\Lambda',N,P)$ a series datum of weight $w$
(Definition~\ref{is-onto-step-leaves-theta-module}). Then:
\begin{enumerate}
\item for every automorphism $\gamma$ of the level module and every allowed vector $x$, some allowed
  vector is a right multiple of $x$ by $\gamma$;
\item for every allowed vector $x$ there is a short series of $(\Lambda',N,P)$ at $x$;
\item if $\gamma$ has determinant one, $x$ is a right multiple of $x'$ by $\gamma$, and $\phi$ and
  $\phi'$ are short series of $(\Lambda',N,P)$ at $x$ and $x'$, then the level action of $\gamma$ maps
  $\phi$ to $\phi'$;
\item if at every embedding $\tau$ the last of the $k+l-1$ positive entries $w^+_{\tau,i}$ of $w$ is at
  least $1$, where $k$ is the first rank of the definite-side datum, every short series of weight $w$ lies
  in $S_l(w)$.
\end{enumerate}
\end{lemma}
\leanref{Development/AbelianVarietyTheory/Proofs/FoundationProofs/TwistTrivialOnRational.lean}{GlobalTheory.TraceLattice.exists_rightMul}
\leanref{Development/AbelianVarietyTheory/Proofs/FoundationProofs/TwistTrivialOnRational.lean}{GlobalTheory.TraceLattice.exists_isSeries}
\leanref{Development/AbelianVarietyTheory/Proofs/FoundationProofs/TwistTrivialOnRational.lean}{GlobalTheory.TraceLattice.mv_series}
\begin{proof}
Write $\Lambda=\Lambda_l$, with the pairing $\langle\,,\rangle$ of the datum of the core at level $l$, the
level $M$, the level ideal $M_b$ and the separating element $c$, an element of $\bar M_b$ with $c-1\in
M_b$. The lattice $\Lambda$ with this pairing is the lattice of the PEL datum of the core at level $l$,
over the CM field $E_h$ with the prime $d$, the elements $\sqrt{-d}$ and $w_0$ and the square root $\theta$
of its host, and the level of the core, with its level ideal $M_b$, is a level of this datum; so
Lemma~\ref{is-onto-step-leaves-move-isometry} applies to $\Lambda$, this level and $\gamma$. For
$q\in M_b\Lambda$ and $y\in\Lambda$, $cq\in\bar M_bM_b\Lambda=M\Lambda$, so $\langle y,cq\rangle\in
M\mathbf Z$. Write $\tilde\gamma$ for the lift of $\gamma$ used in
Definition~\ref{is-onto-step-leaves-theta-module}(8).

(1) Let $F=F_M$ be the map of Lemma~\ref{is-onto-step-leaves-move-isometry} for $\gamma$, and $G$ an
endomorphism with $G(F(a))\equiv a\pmod{M\Lambda}$. Let $F^*$ and $G^*$ be the endomorphisms with $\langle
F^*a,b\rangle=\langle a,Fb\rangle$ and $\langle G^*a,b\rangle=\langle a,Gb\rangle$, and put $x'_b=F^*(x_b)$.
Then $\langle x'_b,ca\rangle=\langle x_b,cFa\rangle$, and $Fa-\tilde\gamma(a)\in M_b\Lambda$, as $F$ induces
$\gamma$; so $\langle x'_b,ca\rangle\equiv\langle x_b,c\tilde\gamma(a)\rangle\pmod M$, and $x'$ is a right
multiple of $x$ by $\gamma$. For $a\in\Lambda$, $F^*(G^*a)-a$ pairs with every $b$ to $\langle
a,G(Fb)-b\rangle\in M\mathbf Z$, so it lies in $M\Lambda\subset\bar M_b\Lambda$: $F^*$ is surjective modulo
$\bar M_b\Lambda$. As $x$ is allowed, every $a\in\Lambda$ is congruent modulo $\bar M_b\Lambda$ to $F^*(a')$
for some $a'$ in the span of $\bar M_b\Lambda$ and the relevant $x_b$, and $F^*(a')$ lies in the span of
$\bar M_b\Lambda$ and the corresponding $x'_b$. So $x'$ is allowed.

(2) Among the laws of the run is that for every series datum $(\Lambda',N,P)$ of a weight $w$ and every
allowed vector $x$ there is one and only one section of weight $w$ that satisfies the short rule of
$\delta_l(n_w)$ for $\Lambda'$, $N$, $x$ and $P$. This gives (2).

(3) Let $u$ be the function of Lemma~\ref{is-onto-step-leaves-move-const} for the weight $w$, so that
$\gamma\phi=u(\gamma)\phi'$ by that lemma. If $\phi'=0$, then $\gamma\phi=0=\phi'$. Suppose $\phi'\ne0$.
We show that $u(\alpha\beta)=u(\alpha)u(\beta)$ for all automorphisms $\alpha,\beta$ of the level module.
By (1), choose an allowed right multiple $x_1$ of $x'$ by $\alpha$ and an allowed right multiple $x_2$ of
$x_1$ by $\beta$, and by (2) short series $\phi_1$ and $\phi_2$ of $(\Lambda',N,P)$ at $x_1$ and $x_2$. Then
$x_2$ is a right multiple of $x'$ by $\alpha\beta$:
\[
 \langle x_{2,b},ca\rangle\equiv\langle x_{1,b},c\tilde\beta a\rangle\equiv\langle
 x'_b,c\tilde\alpha\tilde\beta a\rangle\equiv\langle x'_b,c\widetilde{\alpha\beta}a\rangle\pmod M,
\]
the last because $\tilde\alpha\tilde\beta-\widetilde{\alpha\beta}$ maps $\Lambda$ into $M_b\Lambda$. By
Lemma~\ref{is-onto-step-leaves-move-const}, $(\alpha\beta)\phi_2=u(\alpha\beta)\phi'$,
$\beta\phi_2=u(\beta)\phi_1$ and $\alpha\phi_1=u(\alpha)\phi'$; by
Lemma~\ref{is-onto-step-leaves-move-laws}, $(\alpha\beta)\phi_2=\alpha(\beta\phi_2)=u(\beta)u(\alpha)\phi'$.
As $\phi'\ne0$, $u(\alpha\beta)=u(\alpha)u(\beta)$. So $u$ is a homomorphism to the commutative group
$F^\times$. By the laws of the core, $M=g_0g_1$; the host gives $g_k\equiv1\pmod{Q_Sqd}$ with $4\mid Q_S$,
so the primes $g_k$ are congruent to $1$ modulo $4$ and differ from $3$, and $M$ is odd and prime to $3$.
As $\gamma$ has determinant one, Lemma~\ref{is-onto-step-leaves-move-det-character} gives $u(\gamma)=1$.

(4) The weight $w$ gives a weight of the datum of the core at level $l$ whose positive entries at an
embedding of the CM type are those of $w$, and the multiplicity $p$ of the datum at such an embedding is
$k+l-1$, by the pins of the laws; so the last positive entry of this weight is positive at every
embedding of the CM type. Among the laws of the run is that every short series of a weight of the datum
whose last positive entries at the embeddings of the CM type are positive is a cusp form. So every short
series of weight $w$ lies in $S_l(w)$.
% Lean: (1) is exists_rightMul: g from T.exists_completed_isometry_inducing_half_level_automorphism, g'
% = g.inv l.M, the adjoints fa and ga (exists_adjoint), hsurj with modN_of_dvd and natCast_mem_half (M
% in Mfr), IsRightMul from hfa, dvd_pair_of_mem and liftLa_spec, and AdmX from
% DatumCoverAdjoint.span_image_sup_eq_top. (2) is exists_isSeries (the law shTot of LawsSh; IsSeries.unique
% for the uniqueness). (3) is mv_series: choose u by exists_const; the case f' = 0; hmul with
% exists_rightMul, exists_isSeries, IsRightMul.mul and actL_mul (with
% T.actLev_full_level_group_framedPoints); then char_eq_one_of_detOne with not_three_dvd_level (M = g0 g1
% from h.pins, g_k mod 4 = 1 from D.ntypes.g_cond and D.host.tower_dvd). (4) is the hc of
% trace_mem_thetaEig, for any weight w with w.LastPos: WeightOnDatum.weight_on_datum (with the pins
% p = (sigW j).1 and ext0) turns w.LastPos into the LastPos of the datum's weight wtA w, and the law
% shCusp gives sh <= cutR (cutR_eq). Sizes, following imports and
% stopping at statements the book cites: exists_rightMul about 500 lines, exists_isSeries about 1,200
% (nearly all definitions), mv_series about 24,200, of which 23,000 below exists_const (Lemma
% is-onto-step-leaves-move-const) and the laws of the level action (Lemma is-onto-step-leaves-move-laws).
\end{proof}

\begin{lemma}\label{is-onto-step-leaves-det-congruence}
Let $E/F$ be an extension of number fields, $\mathbf A_{F,f}$ the finite adeles of $F$ and
$\hat{\mathcal O}_F$ its integral elements. For a subset $X$ of $E$ or of $E^n$, write $\hat X$ for the
subgroup of $E\otimes_F\mathbf A_{F,f}$, or of $(E\otimes_F\mathbf A_{F,f})^n$, generated by the $a\otimes x$
with $a\in\hat{\mathcal O}_F$ and $x\in X$. Let $L\subset E^n$ be a finitely generated $\mathcal
O_E$-submodule that spans $E^n$, $I\ne0$ an ideal of $\mathcal O_E$, $\Lambda$ an $\mathcal O_E$-module
with an injective $\mathcal O_E$-linear map $\iota\colon\Lambda\to E^n$ with image $L$, $f$ an
endomorphism of $\Lambda$ that induces an automorphism $\gamma$ of $\Lambda/I\Lambda$, and $u\in\mathcal
O_E$ whose class modulo $I$ is the determinant of $\gamma$ over $\mathcal O_E/I$. Let $G\in M_n(E\otimes_F\mathbf
A_{F,f})$ satisfy $G\iota(a)-\iota(f(a))\in I\hat L$ for every $a\in\Lambda$, where $I\hat L$ is the
subgroup generated by the $bz$ with $b\in I$ and $z\in\hat L$. Then $\det G-u\in\hat I$.
\end{lemma}
\leanref{Development/AbelianVarietyTheory/Proofs/DivisorProofs/LevelGroupLatticeCongruence.lean}{GlobalTheory.LevDet.det_sub_mem_completedIdeal_of_level}
\leanref{Development/AbelianVarietyTheory/Proofs/DivisorProofs/LevelGroupLatticeCongruence.lean}{GlobalTheory.LevDet.det_quot_eq}
\begin{proof}
\emph{Step 1: coordinates modulo $I$.} The module $L$ has a pseudo-basis: a basis $e_1,\dots,e_n$ of $E^n$
and nonzero fractional ideals $\mathfrak c_i$ with $L=\bigoplus_i\mathfrak c_ie_i$. For each $i$ choose
$x_i\in\mathfrak c_i$ and $t_i\in\mathfrak c_i^{-1}$ with $t_ix_i=1-m_i$ and $m_i\in I$. Let $a_j\in\Lambda$
with $\iota(a_j)=x_je_j$, and for $v\in\Lambda$ let $\varphi(v)\in\mathcal O_E^n$ have the entries $t_i$
times the $i$-th coordinate of $\iota(v)$ in the basis $e$; they are integral, as $t_i\mathfrak c_i\subset
\mathcal O_E$. Then $\varphi(a_j)_i\equiv\delta_{ij}\pmod I$, and $\sum_i\varphi(v)_ia_i-v\in I\Lambda$, since
its image under $\iota$ is $-\sum_im_iy_ie_i$ with $y_ie_i\in L$ the components of $\iota(v)$. So the classes
of the $a_j$ form a basis of $\Lambda/I\Lambda$ over $\mathcal O_E/I$, in which $\gamma$ has the matrix
$C=\big(\varphi(f(a_j))_i\big)_{i,j}$, and $\det C\equiv u\pmod I$.

\emph{Step 2: the matrix in the basis $e$.} Let $k=P^{-1}GP$, with $P$ the matrix whose columns are the
$e_j$; then $\det k=\det G$. Applied to the $v\in\Lambda$ with $\iota(v)=ye_j$, $y\in\mathfrak c_j$, the
hypothesis gives $k_{ij}y-y'_i\in\widehat{I\mathfrak c_i}$, where $y'_i$ is the $i$-th coordinate of
$\iota(f(v))$, an element of $\mathfrak c_i$. Hence $k_{ij}\in\widehat{\mathfrak c_i\mathfrak c_j^{-1}}$, and
$\det G=\det k$ lies in $\hat{\mathcal O}_E$, as the products of the ideals
$\mathfrak c_{i}\mathfrak c_{\pi(i)}^{-1}$ over a permutation $\pi$ are $\mathcal O_E$.

\emph{Step 3: comparison.} Put $k'_{ij}=t_ix_jk_{ij}$. Then $\det k'=\prod_i(1-m_i)\det G$. For $v=a_j$,
Step 2 gives $t_i\big(k_{ij}x_j-y'_i\big)\in\hat I$, and $t_iy'_i=C_{ij}$; so $k'_{ij}-C_{ij}\in\hat I$. All
entries of $k'$ and $C$ lie in $\hat{\mathcal O}_E$, and $\hat{\mathcal O}_E\hat I\subset\hat I$, so expanding
the determinants gives $\det k'-\det C\in\hat I$. As $\prod_i(1-m_i)\equiv1\pmod I$ and $\det
G\in\hat{\mathcal O}_E$, $\det G-\det C\in\hat I$, and with Step 1, $\det G-u\in\hat I$.
% Lean: det_sub_mem_completedIdeal_of_level: Step 1 is exists_pseudoBasis_adapted, exists_gen_mod,
% hsingle, phi with hphi, h1, h2 and det_quot_eq; Step 2 is P, Q, hQP, k = Q G P, hkdet, hk1, hcong, hk
% and hd (det_mem_adelicCompletionOfElementSet); Step 3 is k', hk'det, hk'C, det_sub_det_mem, h7 and h8.
% About 800 lines.
\end{proof}

\begin{lemma}\label{is-onto-step-leaves-det}
Let $p$ be a point of a family whose realization $R$ satisfies the conditions of its characterization
on its group at its face (Definition~\ref{is-realisation-characterization}(2)), with
group $G$ and subgroup $P_0$, and let a pair of runs at $p$ pinned to the own places be given, with maps
$\ell_i$ (Definition~\ref{is-onto-step-runs}). Then $\ell_i(u)$ has determinant one
(Definition~\ref{is-onto-step-leaves-theta-module}(8)) for every lift $i$ and every $u\in P_0$.
\end{lemma}
\leanref{Development/AbelianVarietyTheory/Proofs/FoundationProofs/LeafDataOfOnTieOfIntegral.lean}{GlobalTheory.TraceLattice.detOneOn_of_pins}
\leanref{Development/AbelianVarietyTheory/Proofs/FoundationProofs/ExistsDetOfMemSt.lean}{GlobalTheory.TraceLattice.detOne_of_levelPin}
\begin{proof}
Use the notation of Definition~\ref{is-onto-step-model-level-condition} at the lift $i$, with
$\Lambda=\Lambda_{i+1}$ and $\iota=\iota_{i+1}$. By the conditions of the characterization on the
group, $P_0$ consists of the elements of $G$ that are upper unitriangular at every place of $P_g$. Let
$u\in P_0$; then $\det u_P=1$ for every $P\in P_g$. By the level condition (c) of
Definition~\ref{is-onto-step-runs}(5), there are a level lift $k$ of $u$ and an $\mathcal O_{E_h}$-linear
map $s_u$ of $\Lambda$ inducing $\ell_i(u)$ with $\iota(s_u(a))-k\,\iota(a)\in M\hat L_W$ for $a\in\bar
M_b\Lambda$. The matrix $y=\beta_i^{-1}k\beta_i$ lies in $GL_{r}(\hat{\mathcal O}_{E_h})$, so
$\sigma=\det y=\det k$ is integral, and for $P\in P_g$ its reduction modulo $P$ is the determinant of the
reduction $u_P$ of $y$, which is $1$.

By (a) of Definition~\ref{is-onto-step-runs}(5) the run used at a finite place satisfies the laws, which
make $\iota$ injective with image the lattice $\mathfrak AL_0$ of $E_h^{r}$, whose closure in the sense
of Lemma~\ref{is-onto-step-leaves-det-congruence} is $\hat L_W$. Write $1=m+c$ with $m\in M_b$ and
$c\in\bar M_b$, and let $\tilde u\in\mathcal O_{E_h}$ represent the determinant of $\ell_i(u)$ over
$\mathcal O_{E_h}/M_b$. For $a\in\Lambda$, $ca\in\bar M_b\Lambda$, so
\[
 (ck)\,\iota(a)-\iota(s_u(a))=-m\,\iota(s_u(a))-\big(\iota(s_u(ca))-k\,\iota(ca)\big)\in M_b\hat L_W,
\]
as $m\in M_b$ and $M\in M_b$. Lemma~\ref{is-onto-step-leaves-det-congruence}, with $I=M_b$, $f=s_u$ and
$G=ck$, gives $c^{r}\sigma-\tilde u=\det(ck)-\tilde u\in\hat M_b$. Let $P\in P_g$; reducing modulo $P$,
where $c\equiv1$ as $m\in M_b\subset P$, and $\sigma\equiv1$, gives $\tilde u\equiv1\pmod P$. By (e) of
Definition~\ref{is-onto-step-runs}(5), $M_b=\prod_{P\in P_g}P$, a product of distinct primes, so
$\tilde u\equiv1\pmod{M_b}$. As $s_u$ induces $\ell_i(u)$, the determinant of $s_u\otimes\mathcal
O_{E_h}/M_b$ is that of $\ell_i(u)$, which is $1$.
% Lean: detOneOn_of_pins: hU (UnipPin) from OntoChain.mem_base_iff with the group levels of the
% characterization, Matrix.det_of_upperTriangular; detOne_of_levelPin (LevelPin gives kappa, St and the
% block of redFn, upper_block_val_submatrix, rank_Wc) and detOne_of_lift (latEmb_iotaW for the laws,
% GlobalTheory.Member.latG_lattice_and_completion for the lattice and its closure,
% exists_det_of_mem_St for sigma), which is det_eq_one_of_lift (l.coprime for 1 = m + c, the level
% ideal HalfPin as hPg, Lemma is-onto-step-leaves-det-congruence, exists_res_eq_zero for the
% reduction, and det_baseChange_lift). About 4,000 lines below detOne_of_levelPin.
\end{proof}

\begin{lemma}\label{is-onto-step-leaves-data}
Let $p$ be a point of a family whose realization satisfies the conditions of its characterization on
its group at its face (Definition~\ref{is-realisation-characterization}(2)), with subgroup $P_0$ and
character $\psi_0$, let a pair of runs at $p$ pinned to the
own places be given, and let $O$ be a record object recorded at $p$ and this pair, with the notation of
Definitions~\ref{is-record-object}
and~\ref{is-onto-record-numbers}. Let $v$ be a finite place of $L$ and $j$ a lift, and suppose that
$\Theta_K(v,j)$ is finitely generated. Then $\Theta_0(v,j)$ is finitely generated and contained in
$V_\psi(j)$, and $|U|\cdot\big(\Theta_K(v,j)\cap V_\psi(j)\big)\subset\Theta_0(v,j)$.
\end{lemma}
\leanref{Development/AbelianVarietyTheory/Proofs/FoundationProofs/LeafDataOfOnTieOfIntegral.lean}{GlobalTheory.TraceLattice.clauses_of_onTie}
\leanref{Development/AbelianVarietyTheory/Proofs/FoundationProofs/LeafDataOfOnTieOfIntegral.lean}{GlobalTheory.TraceLattice.trace_mem_thetaEig}
\begin{proof}
Let $s$ be the run used at $v$, $l=j+1$, $\kappa=\kappa'_l$ and $z=z_j$, with $\mathrm{inc}$ as in
Lemma~\ref{is-onto-step-leaves-ties}, and put $\gamma_u=\ell_j(\mathrm{inc}(u))$ for $u\in U$. As the pair
of runs is pinned to the own places, $s$ satisfies the laws and is rational at $j$. By
Lemma~\ref{is-ties-level-tie}(2), $\mathrm{inc}(U)\subset P_0$ and
$\psi(u)=\psi_0(\mathrm{inc}(u))^{-1}$; so $\psi(1)=1$ and $\psi(ab)=\psi(a)\psi(b)$ for $a,b\in U$, and
every $\gamma_u$ has determinant one by Lemma~\ref{is-onto-step-leaves-det}. By
Lemma~\ref{is-onto-step-leaves-weights}, $n_\kappa=n_V\geq n_{w_F}\geq r$, where $r$ is the rank of the
definite-side datum, so Lemma~\ref{is-onto-step-leaves-move} applies to the weight $\kappa$. By
Lemma~\ref{is-onto-step-leaves-move-last}, the last positive entry of $\kappa$ at every embedding is at
least $2T_O^2\geq1$, so by Lemma~\ref{is-onto-step-leaves-move}(4) every short series of weight $\kappa$
lies in $S_l(\kappa)$.

\emph{The trace.} For $y\in L^{n_Z(j)}$ put $\operatorname{tr}(y)=\sum_{u\in U}\psi(u)\rho_j(u)^{-1}y$.
For $a\in U$, substituting $u=u'a$,
\[
 \rho_j(a)\operatorname{tr}(y)=\sum_{u\in U}\psi(u)\rho_j(au^{-1})y=\sum_{u'\in U}\psi(u'a)\rho_j(u')^{-1}y
 =\psi(a)\operatorname{tr}(y),
\]
so $\operatorname{tr}(y)\in V_\psi(j)$. If $y\in V_\psi(j)$, then
$\rho_j(u)^{-1}y=\psi(u^{-1})y$ and $\psi(u)\psi(u^{-1})=1$, so $\operatorname{tr}(y)=|U|\,y$. If
$\phi_u,\phi\in S_l(\kappa)$ and $\gamma_u\phi_u=\phi$ for all $u\in U$, then by
Lemma~\ref{is-onto-step-leaves-ties}(1)
\[
 \sum_{u\in U}\psi(u)\,z^{-1}(1\otimes\phi_u)=\sum_{u\in U}\psi(u)\rho_j(u)^{-1}z^{-1}(1\otimes\gamma_u\phi_u)
 =\operatorname{tr}\big(z^{-1}(1\otimes\phi)\big).
\]

\emph{Finiteness.} Each member of an orbit family of weight $\kappa$ that lies in $S_l(\kappa)$ is a
short series, so it lies in $\mathrm{Sh}_l(\kappa)$. By Lemma~\ref{is-onto-step-leaves-ties}(2)
and~(3), the generators of $\Theta_0(v,j)$ lie in $\mathrm{Sh}_l(\kappa)^z=\Theta_K(v,j)$. So
$\Theta_0(v,j)\subset\Theta_K(v,j)$, and $\Theta_0(v,j)$ is finitely generated, as $\mathcal O_L$ is
noetherian.

\emph{The eigenspace.} Let $(\phi_u)$ be an orbit family of weight $\kappa$ for
$\ell_j\circ\mathrm{inc}$ with members in $S_l(\kappa)$, with series datum $(\Lambda',N,P)$ and allowed
vectors $x_0$ and $x_u$. By Lemma~\ref{is-onto-step-leaves-move}(2) and~(4) there is a short series
$\phi\in S_l(\kappa)$ of $(\Lambda',N,P)$ at $x_0$. Each $x_u$ is a right multiple of $x_0$ by
$\gamma_u$, which has determinant one, so by Lemma~\ref{is-onto-step-leaves-move}(3) $\gamma_u\phi_u=\phi$.
So the generator $\sum_u\psi(u)z^{-1}(1\otimes\phi_u)$ is $\operatorname{tr}(z^{-1}(1\otimes\phi))$,
which lies in $V_\psi(j)$. Hence $\Theta_0(v,j)\subset V_\psi(j)$.

\emph{The last inclusion.} Let $\phi$ be a short series of weight $\kappa$, of a series datum
$(\Lambda',N,P)$ at an allowed vector $x_0$; it lies in $S_l(\kappa)$. By
Lemma~\ref{is-onto-step-leaves-move}(1), (2) and~(4), choose for every $u\in U$ an allowed right multiple
$x_u$ of $x_0$ by $\gamma_u$ and a short series $\phi_u\in S_l(\kappa)$ of $(\Lambda',N,P)$ at $x_u$.
Then $(\phi_u)$ is an orbit family, $\gamma_u\phi_u=\phi$ by Lemma~\ref{is-onto-step-leaves-move}(3), and
$\operatorname{tr}(z^{-1}(1\otimes\phi))=\sum_u\psi(u)z^{-1}(1\otimes\phi_u)\in\Theta_0(v,j)$. The set of
$\phi'\in S_l(\kappa)$ with $\operatorname{tr}(z^{-1}(1\otimes\phi'))\in\Theta_0(v,j)$ is an
$\mathcal O_F$-submodule, and it contains the short series of weight $\kappa$, so it contains
$\mathrm{Sh}_l(\kappa)$. As $\operatorname{tr}$ is $\mathcal O_L$-linear and $\Theta_K(v,j)$ is spanned by
the $z^{-1}(1\otimes\phi')$ with $\phi'\in\mathrm{Sh}_l(\kappa)$, $\operatorname{tr}$ maps $\Theta_K(v,j)$
into $\Theta_0(v,j)$. For $y\in\Theta_K(v,j)\cap V_\psi(j)$, $|U|\,y=\operatorname{tr}(y)\in\Theta_0(v,j)$.
% Lean: clauses_of_onTie with t the OnTie: h1 and hpsi from OnTie.U_eq and psi_eq; hd from DetOneOn;
% OnTie.Th0_le_ThK (IsOrbitFamily.mem_shortGen with n_kappa = nV) for finiteness; thetaEig_le_eig
% (exists_isSeries, mv_series, sum_eq_trace_cusp, trace_mem_eig) for the eigenspace; and
% trace_mem_thetaEig (spanK_spanIn_le, trace_gen_mem with exists_rightMul, exists_isSeries and
% mv_series) with trace_of_mem_eig for the last inclusion.
\end{proof}

\begin{lemma}\label{is-onto-step-leaves}
For every constant $C_K$ there is $x_1$ with the following property. Let $p$ be a point of a family,
with realization $R$, face $F$, base tuple $\mu_0$ and tuple map $w\mapsto\mu_w$, such that $R$
satisfies the conditions of its characterization on its group at $F$
(Definition~\ref{is-realisation-characterization}(2)); that $\mu_{\bar w}=\mu_w$ for
every finite place $w$, $\mu_w=\mu_0$ for every finite place $w$ that is not twisted, and
$\mu_w\ne\mu_0$ for every twisted place $w$; that the core ties of $p$ hold; and that its member
satisfies the conditions of Lemma~\ref{is-onto-step-numerics} that involve $x_1$ (the size
conditions, $C_{nW}=10$, the bound on $\log g_i$, the bounds on $C_P$ and $x_1\le x$). Let $r$ be a
pair of runs at $p$ pinned to the own places whose runs satisfy the index-parameter conditions of
Lemma~\ref{is-onto-step-index-conditions} with $C_K$, and let $O$ be a record object recorded at
$(p,r)$ (Definition~\ref{is-record-object}). Then $O$ satisfies the supremum, span, containment and
data estimates (1) to (4) of Definition~\ref{is-onto-geometric-estimates}. The three conditions on
the tuples, the core ties, the conditions that involve $x_1$ and the own-places condition (f) of
Definition~\ref{is-onto-step-runs}(5) on $r$ are not used in the proof.
\end{lemma}
\leanref{Development/AbelianVarietyTheory/Proofs/DivisorProofs/LeavesSupSpanData.lean}{GlobalTheory.BigOn.leaves_sup_span_data}
\begin{proof}
Put $x_1=0$. Let $v$ be a finite
place of $L$ and $j$ a lift, let $s$ be the run used at $v$ and $q$ the own prime of $v$, put $l=j+1$,
and let $\kappa,w_L,w_F$ and $T=T_O$ be the weights and the integer of the parameter choices at $l$.
As $r$ is pinned to the own places, $s$ satisfies the laws and is covered and rational at $j$. As the
runs of $r$ satisfy the index-parameter conditions with $C_K$, $s$ satisfies their hypotheses at
$(l,q)$ (Definition~\ref{is-onto-step-index-bound-conclusion}(2)). So by
Lemma~\ref{is-onto-step-leaves-hypotheses}, $M_l(\kappa)$ is a lattice in $S_l(\kappa)$,
$\mathrm{Sh}_l(\kappa)$ spans $S_l(\kappa)$, and the run satisfies (H2) for the scalars of $q$. By
Lemma~\ref{is-onto-step-index-bound-ties}(1), where $\kappa_j=\kappa$, and
Lemma~\ref{is-onto-step-leaves-ties}(1) and~(2), the isomorphism $z_j$ and the map
$g\mapsto\ell_j(\mathrm{inc}(g))$ satisfy hypotheses (1) and (2) of
Lemma~\ref{is-onto-step-leaves-form-tie} with $p=q$. That lemma gives at $(v,j)$ the statements of the
supremum, span and containment estimates.

The supremum estimate (1) does not involve $v$, and $L$ has a finite place, so (1) holds. Applied at
every twisted place and at its conjugate, the lemma gives the span and containment estimates (2)
and~(3).

For the data estimate (4), let $v$ be a twisted place or the conjugate of one, and $j$ a lift. By the
span estimate, $\Theta_K(v,j)$ is a lattice, so it is finitely generated, and
Lemma~\ref{is-onto-step-leaves-data} gives (4) at $(v,j)$.
% Lean: leaves_sup_span_data takes x1 = 0, and uses neither PageNum, nor CoreTies, nor OwnPin, nor the
% tuple conditions tup_mu_conj, tup_mu_off and tup_mu_ne; the characterization enters only through
% IsGrp (its group levels, in mem_base_iff for detOneOn_of_pins). kappa = kaR from IsRec.kap; IsRec.tied
% gives the ties ht, and IsRec.ok (RunsOk) the laws, xiRat, cover, level and half pins of the run used
% at v. (1)-(3): leaves_sup_span_lemmaS_of_formTieK with GlobalTheory.chartIndependence_unconditional,
% the FormTieK and act_eq of the tie at (v, i), ownAt_pos, and lat_span_of_hypC applied to the HypC
% part of RunsHold at v; it takes v0 from exists_place3 for (1). (4): leafData_of_onTie_of_leafSpan
% with LeafSpan and TiedAt at w and conj w, built from the OnTie, the laws, xiRat, kaWt_n,
% kaR_lastPos and detOneOn_of_pins. Size, following imports and stopping at statements the book cites:
% the open parts are in Lemmas is-onto-step-leaves-hypotheses-sup (about 3,900 lines), -products (about
% 34,600), -move-laws (about 5,000), -level-theta (about 6,300), -level-short (about 22,200) and
% -move-const (about 21,600, nearly all shared with -level-short, and 6,300 of them not below -products).
\end{proof}

\subsection{The index bound and the lattice $\Theta_J$}\label{subsection-is-index-bound}

\begin{definition}\label{is-onto-step-index-bound-conclusion}
Let $s$ be a run at a point of a family that satisfies the laws (Definition~\ref{is-onto-step-runs}),
let $l$ be a level of a lift and $q$ a prime, and let $F$ be the coefficient field of the core, with
ring of integers $\mathcal O_F$. At the level $l$ the run has: for every weight $w$ an $F$-vector space
$S_l(w)$ of cusp forms of weight $w$, with the $\mathcal O_F$-module $M_l(w)$ of its integral cusp
forms; a commutative $F$-algebra $A_l$, its \emph{ring of forms}, with injective $F$-linear maps
$\iota_w\colon S_l(w)\to A_l$; for the weights $\kappa'_l$, $w_{L,l}$ and $w_{F,l}$ of the parameter
choices at level $l$, their integer $T=T_O$ and the scalars of $q$, an $\mathcal O_F$-submodule
$\Theta_l(q)$ of $A_l$, its \emph{theta module}; and a real function $\nu_l$ on the finite places of
$F$, its \emph{exponent function} for $w_{L,l}$. For a finite place $\lambda$ of $F$ let $F_\lambda$
be the completion, $\mathcal O_\lambda$ its ring of integers and $\ell_\lambda=\log
N\lambda/[F:\mathbf Q]$, and for an $\mathcal O_F$-submodule $X$ of $A_l$ let $\hat X_\lambda$ be the
$\mathcal O_\lambda$-span of the image of $X$ in $F_\lambda\otimes_FA_l$. Put
$M=\iota_{\kappa'_l}(M_l(\kappa'_l))$.
\begin{enumerate}
\item Let $\beta$ be a real number. The run satisfies the \emph{conclusion of the index-parameter
  conditions} at $(l,q)$ with bound $\beta$ if there is a function $\rho\geq0$ on the finite places of
  $F$ such that $\sum_\lambda\rho(\lambda)\ell_\lambda$ converges to a number at most $\beta$, and
  such that for every finite place $\lambda$, every integer $e\geq T\nu_l(\lambda)$ and every
  $F_\lambda$-subspace $V$ of $F_\lambda\otimes_FA_l$ contained in the $F_\lambda$-span of $\hat
  M_\lambda$, the $\mathcal O_\lambda$-module
  \[
   (\hat M_\lambda\cap V)\big/\big(\hat M_\lambda\cap\hat\Theta_l(q)_\lambda\cap V\big)
  \]
  has finite length, at most $e\dim_{F_\lambda}V+\rho(\lambda)$.
\item Let $r$ be a pair of runs at a point, with coefficient field $L$ and member of ranks $f$ and
  $m$, and let $b$ and $\delta_2$ be the numbers of the size data of the member. The runs of $r$
  satisfy the \emph{index-parameter conditions} with a constant $C_K$, as in
  Lemma~\ref{is-onto-step-index-conditions}, if for every finite place $v$ of $L$ and every lift $j$
  the run $s$ used at $v$, with the own prime $q$ of $v$, satisfies the hypotheses of these conditions
  with the constants $(C_K,20,C_P,e_P,C_P,e_P,7)$, and their conclusion at $(j+1,q)$ with bound
  $\operatorname{rk}(M_{j+1}(\kappa'_{j+1}))\,fmb\,\delta_2$, where $\operatorname{rk}$ is the
  dimension of the $F$-span.
\end{enumerate}
\end{definition}
\leanref{Formalism/SharpRestriction/Divisibility/WeightAlgebra.lean}{Formalism.MainNotions.WeightAlgebra.Concl}
\leanref{Formalism/SharpRestriction/Divisibility/WeightAlgebra.lean}{Formalism.MainNotions.WeightAlgebra.PJ}
\leanref{Formalism/SharpRestriction/Divisibility/WeightAlgebra.lean}{Formalism.MainNotions.WeightAlgebra.PJ.Theta}
\leanref{Development/AbelianVarietyTheory/Definitions/HeightFunctions/GrowthConstantBoundGcR.lean}{GlobalTheory.RunJoin.nusR}
% Lean: the run's data at level l are recPJ h l hl : WeightAlgebra.PJ D (recB h.toLawsCore) P6 l, with
% Forms (the ring of forms), inc w (the maps iota_w) and Theta S ka wL wF T (the theta module, for the
% scalars scalOfP at q); S_l(w) and M_l(w) are (recB h.toLawsCore).Cusp l w and .Mlat l w, which depend
% only on the core, and F is (recB _).KO = J.Q.c.F. The weights are kaR, lR and fR of the parameter
% choices, T = TO, and nu_l is nusR h l (lR l). In Concl, qlen la (loc Msup inf ofK V) (loc Theta inf
% ofK V) is the length of the quotient displayed in (1), and lN la = lw la is ell_lambda. Part (2) is
% RunsHold Tp hK EX CK r = HypC and Concl; the book does not unfold HypC.

\begin{lemma}\label{is-onto-step-index-bound-localize}
Let $F$ be a number field, $\lambda$ a finite place of $F$ with completion $F_\lambda$ and ring of
integers $\mathcal O_\lambda$, $\mathcal O_{(\lambda)}$ the localization of $\mathcal O_F$ at
$\lambda$, and $A$ an $F$-vector space. For an $\mathcal O_F$-submodule $X$ of $A$ write $X_{(\lambda)}$
for its localization, the $\mathcal O_{(\lambda)}$-span of $X$, and $\hat X$ for the
$\mathcal O_\lambda$-span of the image of $X$ in $F_\lambda\otimes_FA$; for an $F$-subspace $W$ of $A$
write $\hat W$ for the $F_\lambda$-span of its image.
\begin{enumerate}
\item Let $Y$ be an $\mathcal O_F$-submodule of $A$ and $y\in A$ with $1\otimes y\in\hat Y$. Then
  $y\in Y_{(\lambda)}$.
\item Let $X$ and $Y$ be $\mathcal O_F$-submodules of $A$ and $W$ an $F$-subspace, and suppose that
  $Q=(\hat X\cap\hat W)/(\hat X\cap\hat Y\cap\hat W)$ has finite length over $\mathcal O_\lambda$.
  Then
  \[
   \operatorname{len}_{\mathcal O_{(\lambda)}}\big((X\cap W)_{(\lambda)}/(X\cap W\cap Y)_{(\lambda)}\big)
   \le\operatorname{len}_{\mathcal O_\lambda}Q.
  \]
\end{enumerate}
\end{lemma}
\leanref{Formalism/SharpRestriction/GenusPolyhedralBounds/RelativeIndexClassSets/WeightParameterMemUnitInterval.lean}{Formalism.MainNotions.Gaps.LocalGlobalPassage.descent}
\leanref{Formalism/SharpRestriction/GenusPolyhedralBounds/RelativeIndexClassSets/WeightParameterMemUnitInterval.lean}{Formalism.MainNotions.Gaps.LocalGlobalPassage.l1}
\leanref{Formalism/SharpRestriction/GenusPolyhedralBounds/RelativeIndexClassSets/WeightParameterMemUnitInterval.lean}{Formalism.MainNotions.Gaps.LocalGlobalPassage.LengthLocalizationLeCompletion}
\begin{proof}
(1) The element $1\otimes y$ lies in the $\mathcal O_\lambda$-span of $1\otimes t_1,\dots,1\otimes t_k$
for some $t_a\in Y$. Let $N$ be the $\mathcal O_{(\lambda)}$-span of the $t_a$. It is finitely generated
and torsion-free over the discrete valuation ring $\mathcal O_{(\lambda)}$, so it has a basis
$u_1,\dots,u_m$, which is linearly independent over $F$. Extend it to a basis $B$ of $A$ over $F$. The
coordinates in $B$ of an element of $N$ lie in $\mathcal O_{(\lambda)}$ and vanish off the $u_i$. In
the basis $1\otimes B$ of $F_\lambda\otimes_FA$ the coordinates of $1\otimes x$ are those of $x$; so
the coordinates of the elements of the $\mathcal O_\lambda$-span of the $1\otimes t_a$ lie in
$\mathcal O_\lambda$ and vanish off the $1\otimes u_i$. For $1\otimes y$ this says that the coordinates
of $y$ lie in $F\cap\mathcal O_\lambda=\mathcal O_{(\lambda)}$ and vanish off the $u_i$. So $y\in
N\subset Y_{(\lambda)}$.

(2) Put $A'=(X\cap W)_{(\lambda)}$ and $B'=(X\cap W\cap Y)_{(\lambda)}$. Let $x\in A'$. Then $x\in W$,
and $sx\in X$ for some $s\in\mathcal O_F\setminus\lambda$, which is a unit of $\mathcal O_\lambda$; so
$1\otimes x\in\hat X\cap\hat W$. This gives an $\mathcal O_{(\lambda)}$-linear map $g\colon A'\to Q$,
$x\mapsto1\otimes x$. If $g(x)=g(y)$, then $1\otimes(x-y)\in\hat Y$, so $x-y\in Y_{(\lambda)}$ by (1),
and $x-y\in A'\cap Y_{(\lambda)}=B'$, as localization commutes with finite intersections. Since $Q$ has
finite length, $\pi^nQ=0$ for a uniformizer $\pi\in\mathcal O_F$ of $\lambda$ and some $n$. As $F$ is
dense in $F_\lambda$, every $c\in\mathcal O_\lambda$ is congruent modulo $\pi^n\mathcal O_\lambda$ to an
element of $\mathcal O_F$, so $c$ acts on $Q$ as an element of $\mathcal O_F$. Hence for an
$\mathcal O_{(\lambda)}$-submodule $N$ of $A'/B'$ the image under $g$ of its preimage in $A'$ is an
$\mathcal O_\lambda$-submodule $\Phi(N)$ of $Q$. Since $g(x)=g(y)$ only if $x-y\in B'$,
$\Phi(N)\subset\Phi(N')$ implies $N\subset N'$, so $\Phi$ is strictly increasing. A chain of submodules
of $A'/B'$ therefore gives a chain of the same length in $Q$, and the length of $A'/B'$ is at most that
of $Q$.
% Lean: descent (a basis u of the span L of the t over the localization OOv, extended by
% Basis.extend; hPF and hPH for the coordinates) and l1 (g with g_ker, exists_pow_smul_eq_zero for n,
% approx and exists_factor for the action, the strictly monotone Phi and Order.height_le_height_apply_of_strictMono).
\end{proof}

\begin{lemma}\label{is-onto-step-index-bound-passage}
Let $s$, $l$, $q$ and $F$ be as in Definition~\ref{is-onto-step-index-bound-conclusion}, with its
notation, and suppose that $s$ satisfies the conclusion of the index-parameter conditions at $(l,q)$
with a bound $\beta$ and a function $\rho$. Put $S=S_l(\kappa'_l)$, $M_0=M_l(\kappa'_l)$ and
$\Theta_0=\{f\in S:\iota_{\kappa'_l}(f)\in\Theta_l(q)\}$, and suppose that $M_0$ is a lattice in $S$:
finitely generated, and spanning $S$ over $F$. Let $\lambda$ be a finite place of $F$, with
$\operatorname{idx}_\lambda$ as in Definition~\ref{is-onto-record-numbers}. Then:
\begin{enumerate}
\item $M_0\cap\Theta_0$ spans $S$ over $F$;
\item for every integer $e\geq T\nu_l(\lambda)$ and every subspace $W_0\subset S$,
  $\operatorname{idx}_\lambda(M_0\cap W_0,M_0\cap\Theta_0\cap W_0)\le e\dim W_0+\rho(\lambda)$.
\end{enumerate}
\end{lemma}
\leanref{Formalism/SharpRestriction/GenusPolyhedralBounds/RelativeIndexClassSets/WeightParameterMemUnitInterval.lean}{Formalism.MainNotions.Gaps.LocalGlobalPassage.l2}
\leanref{Formalism/SharpRestriction/GenusPolyhedralBounds/RelativeIndexClassSets/WeightParameterMemUnitInterval.lean}{Formalism.MainNotions.Gaps.LocalGlobalPassage.span_of_fin}
\begin{proof}
Write $\iota=\iota_{\kappa'_l}$, $\Theta=\Theta_l(q)$, $M=\iota(M_0)$, and use the notation of
Lemma~\ref{is-onto-step-index-bound-localize} at $\lambda$. Let $e\geq T\nu_l(\lambda)$ be an integer
and $W_0\subset S$ a subspace. Since $M_0$ spans $S$, $\iota(W_0)$ lies in the $F$-span of $M$, so the
$F_\lambda$-span $V$ of the image of $\iota(W_0)$ lies in the $F_\lambda$-span of $\hat M$. As
$F_\lambda$ is flat over $F$ and $\iota$ is injective, $\dim_{F_\lambda}V=\dim\iota(W_0)=\dim W_0$. So
the conclusion of the index-parameter conditions applies to $V$: the module
$(\hat M\cap V)/(\hat M\cap\hat\Theta\cap V)$ has finite length, at most $e\dim W_0+\rho(\lambda)$.
Lemma~\ref{is-onto-step-index-bound-localize}(2), for $X=M$, $Y=\Theta$ and $W=\iota(W_0)$, whose
$\hat W$ is $V$, bounds the length of $(M\cap\iota(W_0))_{(\lambda)}/(M\cap\iota(W_0)\cap\Theta)_{(\lambda)}$
by the same number. As $\iota$ is injective, $M\cap\iota(W_0)=\iota(M_0\cap W_0)$ and
$M\cap\iota(W_0)\cap\Theta=\iota(M_0\cap W_0\cap\Theta_0)$, and $\iota$ carries localizations to
localizations. So
\[
 \operatorname{len}_{\mathcal O_{(\lambda)}}\big((M_0\cap W_0)_{(\lambda)}/(M_0\cap W_0\cap\Theta_0)_{(\lambda)}\big)
 \le e\dim W_0+\rho(\lambda),\tag{$*$}
\]
and the left side is finite.

(2) As $M_0\cap\Theta_0\cap W_0\subset M_0\cap W_0$, the index
$\operatorname{idx}_\lambda(M_0\cap W_0,M_0\cap\Theta_0\cap W_0)$ is the length in $(*)$.

(1) Apply $(*)$ with $W_0=S$ and an integer $e\geq T\nu_l(\lambda)$: the module
$M_{0,(\lambda)}/(M_0\cap\Theta_0)_{(\lambda)}$ has finite length. Let $x\in M_0$. The
$\mathcal O_{(\lambda)}$-linear map $c\mapsto cx$ from $\mathcal O_{(\lambda)}$ to this quotient is not
injective, since $\mathcal O_{(\lambda)}$ is not a field and so has infinite length. So $dx\in(M_0\cap
\Theta_0)_{(\lambda)}$ for some nonzero $d\in\mathcal O_{(\lambda)}$, and $x$ lies in the $F$-span of
$M_0\cap\Theta_0$. As $M_0$ spans $S$, so does $M_0\cap\Theta_0$.
% Lean: l2, with H1c = l1c (V inside locK la Msup), H1b = l1b (the dimension of V), l1' = l1, the
% transport idxRules_span_map, localization_map₄ and idxRules_len_map along inc, and
% idxRules_lidx_of_le for (2); (1) is span_of_fin at e0 = Q.eLa la (Run.eLa_ge) and W0 = top.
\end{proof}

\begin{lemma}\label{is-onto-step-index-bound-base-change}
Let $F\subset L$ be number fields, $Z_0$ a finite-dimensional $F$-vector space, $n\geq0$ an integer and
$\iota\colon Z_0\to L^n$ an $F$-linear map that carries $F$-linearly independent families to
$L$-linearly independent ones and whose image spans $L^n$ over $L$. For an $\mathcal O_F$-submodule $X$
of $Z_0$ write $X^\iota$ for the $\mathcal O_L$-span of $\iota(X)$. Let $v$ be a finite place of $L$,
$\lambda$ the place of $F$ below it and $e(v|\lambda)$ the ramification index, and let $A_0$ and $C_0$
be lattices in $Z_0$. Then $\operatorname{idx}_v(A_0^\iota,C_0^\iota)=e(v|\lambda)\operatorname{idx}_\lambda(A_0,C_0)$.
\end{lemma}
\leanref{Formalism/SharpRestriction/GenusPolyhedralBounds/RelativeIndexClassSets/WeightParameterMemUnitInterval.lean}{Formalism.MainNotions.Gaps.LocalGlobalPassage.l3}
\leanref{Formalism/SharpRestriction/GenusPolyhedralBounds/RelativeIndexClassSets/WeightParameterMemUnitInterval.lean}{Formalism.MainNotions.Gaps.LocalGlobalPassage.IndexMulRamificationIndex}
\begin{proof}
The localizations of $A_0$ and $C_0$ at $\lambda$ are free of rank $\dim Z_0$ over the discrete
valuation ring $\mathcal O_{(\lambda)}$: there are bases $a$ and $c$ of $Z_0$ over $F$ whose
$\mathcal O_{(\lambda)}$-spans are $A_{0,(\lambda)}$ and $C_{0,(\lambda)}$. Their images $\iota(a)$ and
$\iota(c)$ are bases of $L^n$ over $L$. An element of $\mathcal O_F\setminus\lambda$ lies outside $v$,
so the localizations of $A_0^\iota$ and $C_0^\iota$ at $v$ are the $\mathcal O_{(v)}$-spans of
$\iota(a)$ and $\iota(c)$. Let $\gamma_0$ be the automorphism of $Z_0$ with $\gamma_0(a_i)=c_i$, and
$\gamma$ that of $L^n$ with $\gamma(\iota(a_i))=\iota(c_i)$. If $x=\sum_ix_ia_i$ with $x_i\in F$, then
$\iota(x)=\sum_ix_i\iota(a_i)$; so the matrix of $\gamma$ in the basis $\iota(a)$ is the matrix of
$\gamma_0$ in the basis $a$, and $\det\gamma=\det\gamma_0\in F^\times$. For a basis $b$ of a vector
space over the fraction field of a discrete valuation ring and an automorphism $\delta$, the index of
the span of $b$ against the span of $\delta(b)$ is the valuation of $\det\delta$, by the elementary
divisor theorem. Hence $\operatorname{idx}_\lambda(A_0,C_0)=v_\lambda(\det\gamma_0)$ and
$\operatorname{idx}_v(A_0^\iota,C_0^\iota)=v_v(\det\gamma_0)=e(v|\lambda)v_\lambda(\det\gamma_0)$.
% Lean: l3 with l3a (the valuation of an element of K0 at v is e times that at lambda), exists_loc_basis,
% localization_up_congr and up_span_range, iota_eq_sum for the matrices, and PI.LatIdxDetW.lidx_free
% (the index of a free lattice against its image is the count of det) with lidx_congr.
\end{proof}

\begin{lemma}\label{is-onto-step-index-bound-extend}
Let $F\subset L$, $Z_0$, $\iota$, $v$, $\lambda$ and $e(v|\lambda)$ be as in
Lemma~\ref{is-onto-step-index-bound-base-change}, with its notation $X^\iota$. Let $N_0\subset M_0$ be
lattices in $Z_0$, $e_0$ an integer, $E\geq0$ an integer with $E=e(v|\lambda)e_0$, and $\rho$ a real
number, and suppose that
$\operatorname{idx}_\lambda(M_0\cap W_0,N_0\cap W_0)\le e_0\dim W_0+\rho$ for every subspace
$W_0\subset Z_0$. Let $\Theta$ be an $\mathcal O_L$-submodule of $L^n$ that contains $N_0^\iota$. Then
for every subspace $W\subset L^n$,
\[
 \operatorname{idx}_v(M_0^\iota\cap W,\Theta\cap W)\le E\dim W+e(v|\lambda)\rho.
\]
\end{lemma}
\leanref{Formalism/SharpRestriction/GenusPolyhedralBounds/RelativeIndexClassSets/WeightParameterMemUnitInterval.lean}{Formalism.MainNotions.Gaps.LocalGlobalPassage.l5}
\leanref{Formalism/SharpRestriction/GenusPolyhedralBounds/RelativeIndexClassSets/WeightParameterMemUnitInterval.lean}{Formalism.MainNotions.Gaps.LocalGlobalPassage.lidx_anti}
\begin{proof}
As $e(v|\lambda)>0$ and $E\geq0$, $e_0\geq0$. Let $a_0\in F^\times$ with $v_\lambda(a_0)=e_0$, a power
of a uniformizer. By Lemma~\ref{is-places-index-subspace}(2), for $F$, $\lambda$, the lattices
$N_0\subset M_0$ and $a_0$, the lattice $C_0=M_0\cap a_0^{-1}N_0$ (part (1) of that lemma) has
$\operatorname{idx}_\lambda(M_0,C_0)\le\rho$. By Lemma~\ref{is-onto-step-index-bound-base-change},
$\operatorname{idx}_v(M_0^\iota,C_0^\iota)=e(v|\lambda)\operatorname{idx}_\lambda(M_0,C_0)\le
e(v|\lambda)\rho$.

The modules $M_0^\iota$, $N_0^\iota$ and $C_0^\iota$ are lattices in $L^n$: each is spanned by the
image of a finite generating set, and spans $L^n$, since its $F$-preimage spans $Z_0$ and the image
of $\iota$ spans $L^n$. Let $a\in L$ be the image of $a_0$, so that $v_v(a)=e(v|\lambda)e_0=E$. For
$x\in C_0$, $a_0x\in N_0$ and $\iota(a_0x)=a\iota(x)$, so $C_0^\iota\subset M_0^\iota\cap
a^{-1}N_0^\iota$, which is a lattice by Lemma~\ref{is-places-index-subspace}(1). For lattices $A$ and
$B\subset B'$ in one space, the index is additive, $\operatorname{idx}_v(A,B)=\operatorname{idx}_v(A,B')+
\operatorname{len}_v(B'/B)$, so $\operatorname{idx}_v(A,B')\le\operatorname{idx}_v(A,B)$. Hence
$\operatorname{idx}_v(M_0^\iota,M_0^\iota\cap a^{-1}N_0^\iota)\le e(v|\lambda)\rho$, and
Lemma~\ref{is-places-index-subspace}(2), for $L$, $v$, the lattices $N_0^\iota\subset M_0^\iota$ and
$a$, gives
\[
 \operatorname{idx}_v(M_0^\iota\cap W,N_0^\iota\cap W)\le E\dim W+e(v|\lambda)\rho
\]
for every subspace $W\subset L^n$.

Let $W\subset L^n$ be a subspace, and put $A=M_0^\iota\cap W$ and $B'=N_0^\iota\cap W$, lattices in
$W$ with $B'\subset A$ and $B'\subset\Theta\cap W$. Then
\[
 \operatorname{idx}_v(A,\Theta\cap W)\le\operatorname{len}_v\big(A/(A\cap\Theta\cap W)\big)
 \le\operatorname{len}_v(A/B')=\operatorname{idx}_v(A,B'),
\]
and the bound follows.
% Lean: l5: he0nn, exists_vw_eq at lambda, three_forms (mp) and lattice_intersection over K0, l3,
% index_monotone (IdxMono, through idxRules_add), three_forms (mpr) over K, and lidx_anti.
\end{proof}

\begin{lemma}\label{is-onto-step-index-bound-sum}
Let $F\subset L$ be number fields. For finite places $\lambda$ of $F$ and $v$ of $L$ put
$\ell_\lambda=\log N\lambda/[F:\mathbf Q]$ and $\ell_v=\log Nv/[L:\mathbf Q]$, and let $\lambda_v$ be
the place of $F$ below $v$ and $e(v|\lambda_v)$ the ramification index. Let $r\geq0$ be a function on
the finite places of $F$ such that $\sum_\lambda r(\lambda)\ell_\lambda$ converges. Then for every
finite set $S'$ of finite places of $L$,
\[
 \sum_{v\in S'}\ell_v\,e(v|\lambda_v)\,r(\lambda_v)\le\sum_\lambda r(\lambda)\ell_\lambda.
\]
\end{lemma}
\leanref{Formalism/SharpRestriction/GenusPolyhedralBounds/RelativeIndexClassSets/WeightParameterMemUnitInterval.lean}{Formalism.MainNotions.Gaps.LocalGlobalPassage.l4}
\leanref{Formalism/SharpRestriction/GenusPolyhedralBounds/RelativeIndexClassSets/WeightParameterMemUnitInterval.lean}{Formalism.MainNotions.Gaps.LocalGlobalPassage.SumOverPlacesLeBaseSum}
\begin{proof}
For $v$ over $\lambda$ with residue degree $f(v|\lambda)$, $Nv=N\lambda^{f(v|\lambda)}$, so
$\ell_v=f(v|\lambda)\log N\lambda/[L:\mathbf Q]$. Group the terms of the left side by $\lambda=\lambda_v$.
For each $\lambda$, $\sum e(v|\lambda)f(v|\lambda)\le[L:F]$, the sum running over the $v\in S'$ above
$\lambda$; so these terms add up to at most $r(\lambda)\log N\lambda\,[L:F]/[L:\mathbf Q]
=r(\lambda)\ell_\lambda$, as $[L:\mathbf Q]=[L:F][F:\mathbf Q]$. The left side is therefore at most a
finite subsum of $\sum_\lambda r(\lambda)\ell_\lambda$, whose terms are nonnegative.
% Lean: l4, with lw_eq, sum_ef_le (Ideal.sum_ramification_inertia_eq_finrank) and Summable.sum_le_tsum.
\end{proof}

\begin{lemma}\label{is-onto-step-index-bound-exponent}
Let $s$ be a run at a point of a family that satisfies the laws, $l$ a level of a lift, $F$ the
coefficient field of the core, $T=T_O$ the integer of the parameter choices and $\nu_l$ the exponent
function of $s$ (Definition~\ref{is-onto-step-index-bound-conclusion}). Let $\lambda$ be a finite
place of $F$, with residue characteristic $p_\lambda$ and ramification index $e_\lambda$ over
$p_\lambda$, and let $\nu^*_l(\lambda)$ be the normalized exponent at $\lambda$ of the depth data of
the core at level $l$. Put $b_\lambda=1$ if $p_\lambda$ is one of the primes $g_0,g_1,q,d$ of the host
or divides its integer $Q_S$, and $b_\lambda=0$ otherwise, and
$E_l(\lambda)=\max\big(0,\lceil e_\lambda T(\nu^*_l(\lambda)+b_\lambda)\rceil\big)$. Then
$\nu_l(\lambda)=e_\lambda\nu^*_l(\lambda)$, and $T\nu_l(\lambda)\le E_l(\lambda)$.
\end{lemma}
\leanref{Development/AbelianVarietyTheory/Proofs/JetHarmonics/NusOfEqRamMulNuStarAt.lean}{GlobalTheory.RunDepthScale.nusOf_eq_ram_mul_nuStarAt}
\leanref{Development/AbelianVarietyTheory/Proofs/DivisorProofs/HasIndexBoundOfRuns.lean}{GlobalTheory.BigOn.nusR_le_eLam}
\leanref{Formalism/LocalToGlobal/ProjectiveLattices/PlaceBelow.lean}{Formalism.L2.Pj.DepPage.eLam}
\begin{proof}
The two functions are defined by the same case distinction on $p_\lambda$: $p_\lambda\in\{g_0,g_1\}$;
$p_\lambda$ a divisor of $Q_S$ or one of $q,d$; or neither. In each case they are computed from
the weight valuation of $w_{L,l}$ at $\lambda$, divided by $e_\lambda$, and the top degree of
$w_{L,l}$, and $\nu_l(\lambda)$ carries the factor $e_\lambda$ that $\nu^*_l(\lambda)$ does not; the
identity is a direct computation. Then $T\nu_l(\lambda)=e_\lambda T\nu^*_l(\lambda)\le e_\lambda
T(\nu^*_l(\lambda)+b_\lambda)\le E_l(\lambda)$.
% Lean: nusOf_eq_ram_mul_nuStarAt (nusOf_eq_nuN, with H1algLink.eWL_eq, prime_mem_iff and
% cmDatum_natCast_mem_heightOneSpectrum_iff, and WeightExponents.DF_wtA) and DepPage.le_eLam; nu*_l is
% nuStarAt (depOf T.toCore EX Tp hK) la l, b_lambda is IsBad, and E_l(lambda) is eLam (ramOverZ la)
% nu* (resCharOf la).
\end{proof}

\begin{lemma}\label{is-onto-step-index-bound-ties}
Let $p$ be a point of a family, with coefficient field $L$, $F$ the coefficient field of the core, $r$
a pair of runs at $p$, and $O$ a record object recorded at $(p,r)$ (Definition~\ref{is-record-object}).
Let $v$ be a finite place of $L$, $s$ the run used at $v$, $q$ the own prime of $v$, $\lambda$ the place
of $F$ below $v$ and $j$ a lift, and use the notation of
Definition~\ref{is-onto-step-index-bound-conclusion} at the level $l=j+1$. Then:
\begin{enumerate}
\item (Form tie.) The weight family $\kappa_j$ of the one-map data is $\kappa'_{j+1}$, and there is an
  $L$-linear isomorphism $z_j\colon L^{n_Z(j)}\to L\otimes_FS_{j+1}(\kappa'_{j+1})$, independent of
  $v$, with the following property. For an $\mathcal O_F$-submodule $X$ of $S_{j+1}(\kappa'_{j+1})$
  write $X^z$ for the $\mathcal O_L$-span of the vectors $z_j^{-1}(1\otimes f)$ with $f\in X$. Then
  $M_s(j)=M_{j+1}(\kappa'_{j+1})^z$, and $\Theta(v,j)=\Theta_0^z$ for the set $\Theta_0$ of those $f$
  with $\iota_{\kappa'_{j+1}}(f)\in\Theta_{j+1}(q)$, formed with the run $s$.
\item (Depth tie.) $E(v,j)=e(v|\lambda)E_{j+1}(\lambda)$, with $E_{j+1}(\lambda)$ as in
  Lemma~\ref{is-onto-step-index-bound-exponent}.
\item The numbers $f$, $m$ and $p_{de}$ of $O$ are the ranks $f$ and $m$ of the member and the prime
  $p_{de}$ of the host.
\end{enumerate}
\end{lemma}
\leanref{Development/AbelianVarietyTheory/Definitions/Jets/EmbShortThLeThetaPlain.lean}{GlobalTheory.FormTies.FormTieK}
\leanref{Development/AbelianVarietyTheory/Proofs/DivisorProofs/NotOntoOfCharAlone.lean}{GlobalTheory.RealIndexBound.OnObjLevelCompatibility}
\leanref{Formalism/LocalToGlobal/ProjectiveLattices/PlaceBelow.lean}{Formalism.L2.Pj.EAt}
\begin{proof}
These are among the conditions that make $O$ recorded at $(p,r)$. Part (1) is the weight condition of
the record and its form tie at $v$, which comes with the isomorphisms $z_j$ of its level tie; part (2)
is the depth tie of $O$ with the core of $p$; and part (3) is part of the depth tie and of the
remaining conditions, as in the proof of Lemma~\ref{is-onto-recorded-numbers}.
% Lean: IsRec.kap (kappa i = kaR (i + 1)); IsRec.tied gives OnObjLevelCompatibility, whose OnTie at v
% and i contains a FormTie (FormTieK: Ms_eq, Th_eq) for recPJ (laws v) (i + 1) with p = ownAt D v and
% the isomorphism z i; RunTie.E_eq gives E v i = EAt ... v (i + 1) = expE = ramOver v * eLam (ramOverZ
% la) (nuStarAt la (i + 1)) (resCharOf la); RunTie.f_eq, m_eq and IsRec.pde_eq give (3).
\end{proof}

\begin{lemma}\label{is-onto-step-index-bound-remainder}
Let $s$, $l$, $q$ and $F$ be as in Definition~\ref{is-onto-step-index-bound-conclusion}, put
$S=S_l(\kappa'_l)$, $M_0=M_l(\kappa'_l)$ and $\Theta_0=\{f\in S:\iota_{\kappa'_l}(f)\in\Theta_l(q)\}$,
and let $c$ be a real number. Let $L\supset F$ be a number field, $n\geq0$ an integer and
$z\colon L^n\to L\otimes_FS$ an $L$-linear isomorphism, and for an $\mathcal O_F$-submodule $X$ of $S$
write $X^z$ for the $\mathcal O_L$-span of the vectors $z^{-1}(1\otimes f)$ with $f\in X$. Suppose
that $M_0^z$ is a lattice in $L^n$, and that $s$ satisfies the conclusion of the index-parameter
conditions at $(l,q)$ with bound $\operatorname{rk}(M_0)c$, where $\operatorname{rk}(M_0)$ is the
dimension of the $F$-span of $M_0$. Then there is a function $r_L\geq0$ on the finite places of $L$
such that:
\begin{enumerate}
\item $\sum_{v\in S'}\ell_vr_L(v)\le nc$ for every finite set $S'$ of finite places of $L$, with
  $\ell_v$ as in Definition~\ref{is-indefinite-datum};
\item for every finite place $v$ of $L$, with $\lambda$ the place of $F$ below it, every integer
  $e\geq0$ with $T\nu_l(\lambda)\le e$ and every subspace $W\subset L^n$,
  \[
   \operatorname{idx}_v(M_0^z\cap W,\Theta_0^z\cap W)\le e(v|\lambda)\,e\dim W+r_L(v).
  \]
\end{enumerate}
\end{lemma}
\leanref{Development/AbelianVarietyTheory/Proofs/DivisorProofs/HasIndexBoundOfRuns.lean}{GlobalTheory.BigOn.exists_rem_of_concl}
\leanref{Development/AbelianVarietyTheory/Proofs/DivisorProofs/HasIndexBoundOfRuns.lean}{GlobalTheory.BigOn.isLattice_of_spanK}
\begin{proof}
Let $\iota\colon S\to L^n$ be the $F$-linear map $f\mapsto z^{-1}(1\otimes f)$, so that $X^z=X^\iota$
in the notation of Lemma~\ref{is-onto-step-index-bound-base-change}. It carries $F$-linearly
independent families to $L$-linearly independent ones, since $L$ is flat over $F$, and its image spans
$L^n$, since the vectors $1\otimes f$ span $L\otimes_FS$. In particular $\iota$ is injective.

\emph{$M_0$ is a lattice in $S$ of rank $n$.} The module $\iota(M_0)$ lies in $M_0^z$, which is
finitely generated over $\mathcal O_L$ and hence over $\mathcal O_F$, as $\mathcal O_L$ is a finite
$\mathcal O_F$-module. So $M_0$, isomorphic to a submodule of a finitely generated module over the
noetherian ring $\mathcal O_F$, is finitely generated. The $L$-span of $\iota(M_0)$ is $L^n$, and it is
the image under $z^{-1}$ of $L\otimes_F\operatorname{span}_F(M_0)$. So $L\otimes_F\operatorname{span}_F(M_0)
=L\otimes_FS$, and $\operatorname{span}_F(M_0)=S$, as base change from $F$ to $L$ is injective on
subspaces. Hence $\operatorname{rk}(M_0)=\dim_FS=\dim_L(L\otimes_FS)=n$.

Write $r$ for the function $\rho$ of the conclusion of the index-parameter conditions, so that
$\sum_\lambda r(\lambda)\ell_\lambda\le nc$, and put $r_L(v)=e(v|\lambda)r(\lambda)\geq0$ for $v$
above $\lambda$. Part (1) is Lemma~\ref{is-onto-step-index-bound-sum}.

(2) Let $v$, $e$ and $W$ be as in (2). By Lemma~\ref{is-onto-step-index-bound-passage}, applied at
$\lambda$, $N_0=M_0\cap\Theta_0$ spans $S$, so $N_0\subset M_0$ are lattices in $S$, and
$\operatorname{idx}_\lambda(M_0\cap W_0,N_0\cap W_0)\le e\dim W_0+r(\lambda)$ for every subspace
$W_0\subset S$. As $N_0\subset\Theta_0$, $N_0^z\subset\Theta_0^z$.
Lemma~\ref{is-onto-step-index-bound-extend}, for $Z_0=S$, $e_0=e$, $E=e(v|\lambda)e$,
$\rho=r(\lambda)$ and $\Theta=\Theta_0^z$, gives the bound.
% Lean: exists_rem_of_concl: isLattice_of_spanK (coordMap z injective by FaithfullyFlat.one_tmul_eq_zero_iff,
% fg_of_fg_map_injective, Submodule.baseChange_injective) and hn (X.rk = n); l2 at the Run Q with beta =
% rk * a * b * c * d; r_L(v) = ramificationIdx * r (v.under); l4 for (1); l5 with coordMap_linearIndependent,
% coordMap_span, N0 = M0 inf Th0 and ram_eq for (2).
\end{proof}

\begin{lemma}\label{is-onto-step-index-bound}
Let $C_K$ be a constant, $p$ a point of a family whose core ties hold, $r$ a pair of runs at $p$
pinned to the own places whose runs satisfy the index-parameter conditions of
Lemma~\ref{is-onto-step-index-conditions} with $C_K$, and $O$ a record object recorded at $(p,r)$
(Definition~\ref{is-record-object}) whose number $b$ is the quantity $b$ of the size data of the
member. Suppose that $O$ satisfies the supremum and span estimates (1) and (2) of
Definition~\ref{is-onto-geometric-estimates}. Then $O$ satisfies the index bound (5) of that
definition. The hypothesis that the core ties hold, and the own-places condition (f) of
Definition~\ref{is-onto-step-runs}(5) on $r$, are not used in the proof.
\end{lemma}
\leanref{Development/AbelianVarietyTheory/Proofs/DivisorProofs/HasIndexBoundOfRuns.lean}{GlobalTheory.BigOn.hasIndexBound_of_runs}
\leanref{Development/AbelianVarietyTheory/Proofs/DivisorProofs/HasIndexBoundOfRuns.lean}{GlobalTheory.BigOn.exists_rem_of_tie}
\begin{proof}
Let $L$ be the coefficient field of $p$ and $F$ that of the core, and use the notation of
Definition~\ref{is-onto-step-index-bound-conclusion}. Put $\Theta_1'(v,j)=\Theta(v,j)$ and
$\delta_2'=\delta_2$, the number of the size data, which satisfies $0<\delta_2\le1$. For a twisted
place $v$ or its conjugate and a lift $j$, let $a'(v,j)=\pi_v^{E(v,j)}$ for some $\pi_v\in\mathcal O_L$
with $v_v(\pi_v)=1$; then $a'(v,j)\ne0$ and $v_v(a'(v,j))=E(v,j)$. By the span estimate (2),
$\Theta_1'(v,j)$ is a lattice in $L^{n_Z(j)}$, and it is contained in $\Theta(v,j)$.

\emph{The remainders.} Fix a lift $j$. Let $C_2$ be the set of finite places of $L$ over $p_{de}$ and
$C_1$ the set of the other finite places. For $i=1,2$ there is a function $\rho_i\geq0$ on the finite
places of $L$ with $\sum_{u\in S'}\ell_u\rho_i(u)\le n_Z(j)\,fmb\,\delta_2$ for every finite set $S'$
of finite places, and
\[
 \operatorname{idx}_v\big(M_s(j)\cap W,\Theta(v,j)\cap W\big)\le E(v,j)\dim W+\rho_i(v)
\]
for every $v\in C_i$ and every subspace $W\subset L^{n_Z(j)}$. If $C_i$ is empty, $\rho_i=0$ has these
properties, as $n_Z(j)fmb\delta_2\geq0$.

Otherwise choose $v_1\in C_i$. The run used at a finite place and its own prime depend only on whether
the place lies over $p_{de}$ (Definition~\ref{is-onto-step-runs}(1) and (2)), so every place of $C_i$
has the run $s$ used at $v_1$ and the own prime $q$ of $v_1$. By
Lemma~\ref{is-onto-step-index-bound-ties}(1) at each $v\in C_i$, with $M_0=M_{j+1}(\kappa'_{j+1})$,
$\Theta_0$ formed with $s$ and $q$, and one isomorphism $z=z_j$,
\[
 M_s(j)=M_0^z\quad\text{and}\quad\Theta(v,j)=\Theta_0^z .
\]
By the supremum estimate (1), $M_s(j)$ is a lattice in $L^{n_Z(j)}$. Since the runs of $r$ satisfy the
index-parameter conditions with $C_K$, $s$ satisfies their conclusion at $(j+1,q)$ with bound
$\operatorname{rk}(M_0)fmb\delta_2$ (Definition~\ref{is-onto-step-index-bound-conclusion}(2)), where
$f$ and $m$ are the ranks of the member, which are those of $O$ by
Lemma~\ref{is-onto-step-index-bound-ties}(3), and $b$ and $\delta_2$ are the numbers of the size data;
$b$ is that of $O$ by hypothesis. So Lemma~\ref{is-onto-step-index-bound-remainder} applies, with
$c=fmb\delta_2$ and $n=n_Z(j)$; let $\rho_i$ be its function $r_L$. Its part (1) is the bound on the
sums. Let $v\in C_i$ and $\lambda$ the place of $F$ below $v$. By
Lemma~\ref{is-onto-step-index-bound-exponent}, $E_{j+1}(\lambda)$ is an integer $\geq0$ with
$T\nu_{j+1}(\lambda)\le E_{j+1}(\lambda)$, and by Lemma~\ref{is-onto-step-index-bound-ties}(2),
$E(v,j)=e(v|\lambda)E_{j+1}(\lambda)$. So part (2) of Lemma~\ref{is-onto-step-index-bound-remainder},
with $e=E_{j+1}(\lambda)$, gives the bound on the index.

\emph{Conclusion.} Put $r_1'(u,j)=\rho_1(u)$ and $r_2'(u,j)=\rho_2(u)$. The prime $p_{de}$ of $O$ is
that of the host (Lemma~\ref{is-onto-step-index-bound-ties}(3)), so $r'(v,j)=\rho_i(v)$ for the class
$C_i$ of $v$, and (a) and (b) of the index bound hold.
% Lean: hasIndexBound_of_runs: Theta_1' = O.Th, a from LocalGlobalPassage.exists_vw_eq (a power of a
% uniformizer, idxChange_unif), de2' = L.de2 (thickDensity_pos, thickDensity_le_one), and key for the
% two classes (plc v = pde iff q): for q realized by v_1, exists_rem_of_tie at the FormTieK (ht v_1 i)
% with LeafSup (hS i) and the Concl of RunsHold at v_1; e1 and e2 (TwoRuns.used and ownAt agree within
% the class) and h5 transport the form tie at v to the run used at v_1; hd.E_eq with nusR_le_eLam for
% the exponent; hd.f_eq, hd.m_eq and O.b = L.sc.b for the sums; and hO.pde_eq for r'. The core ties and
% OwnPin are not used. Size: about 2,700 lines of Lean in 150 theorems below hasIndexBound_of_runs.
\end{proof}

\begin{lemma}\label{is-onto-step-depth}
There is a constant $C_0$ with the following property. Let $p$ be a point of a family, $r$ a pair of
runs at $p$ pinned to the own places, and $O$ a record object recorded at $(p,r)$ with constants
$c=(C,C_{b,1},C_{b,2})$ (Definition~\ref{is-record-object}). Let $N$ be numerical bounds for the rows
whose function $F_0$ is, at the lift $j$, the local term of the definite-side datum of the member at
level $j+1$. Suppose that $O$ satisfies (1) and (2) of the page estimates
(Definition~\ref{is-onto-page-estimates}), that the member has $3\le x$, $kb'\le b$, $C_{nW}\geq10$
and $3\le x^{\varepsilon_z/5}$, and that $C\geq C_0$. Then $O$ satisfies the depth estimates for $N$
(Definition~\ref{is-onto-depth-estimates}).
\end{lemma}
\leanref{Development/AbelianVarietyTheory/Proofs/DivisorProofs/LeafDepthAllRecord.lean}{GlobalTheory.BigOn.leafDepth_all_record}
\begin{proof}
\openstep{The depth estimates at the good places, the places over $q$ and $d$, the small primes and
the level primes, from a colength bound at the split places, the inclusion of Hodge-Tate lattices and
the characters of the unramified companions, are not written out.}
% Lean: leafDepth_all_record = RunDepth.leafDepth_all_run applied to the three chart statements
% ChF.colength_bound_split_place_chart, ChF.hodgeTate_inclusion_chart and
% ChF.unramified_companion_chart_characters, with the depth tie of the record (IsRec.depth) and the
% laws LawsM of its tie. The three charts are natural separate pieces. Size: about 298,000 lines of
% Lean in 17,000 theorems below it (text count, stopping at statements the book cites), about 55,000
% of them below no other open step of Lemma is-onto-estimates.
\end{proof}

\begin{lemma}\label{is-onto-step-compare}
Let $\mathcal H$ be the class of standard families. Eventually along $\mathcal H$, every family has a
scale beyond which the following holds for every point $p$ that satisfies the point conditions and the
further conditions, every pair of runs $r$ at $p$ pinned to the own places, every twisted place $v$
and every lift $j$. Let $R$ be the realization of $p$, with the scalars of its characterization, and
$F$ its face, with subgroup $P_0$ and character $\psi_0$, and let $q$ be the own prime of $v$. Let
the \emph{trace lattice} of the run $s$ used at $v$ be the $\mathcal O_L$-span of the vectors
$\sum_{u\in P_0}\psi_0(u)^{-1}c(F_u)$, where $(F_u)_{u\in P_0}$ runs over the orbit families of cusp
forms of $s$ at level $j+1$ and weight $\kappa_j$ for its delta data $\delta_{j+1}(n)$ and the action
of $\ell_j$, and $c$ is the map of the weight tie from these cusp forms to the coordinates of $FW_j$.
Then for every finite place $\lambda$ of $L$ whose residue characteristic is not $q$, the localization
at $\lambda$ of the trace lattice is contained in the localization at $\lambda$ of the base theta
lattice $\Theta^{\mathrm b}_v(j)$.
\end{lemma}
\leanref{Development/AutomorphicArithmetic/ThetaCorrespondence/ThetaPairLifts/ThirdPlaceSeries/CompareOffOwnAtRunsStatement.lean}{GlobalTheory.T.bigOn_comparison_at_twisted_places}
\leanref{Development/AutomorphicArithmetic/ThetaCorrespondence/ThetaPairLifts/ThirdPlaceSeries/CompareOffOwnAtRuns.lean}{GlobalTheory.BigOn.compareOffOwn_atRuns_alt}
\leanref{Development/AbelianVarietyTheory/Proofs/FoundationProofs/TraceLatticeLeThetaLatticeAtPlace.lean}{GlobalTheory.TraceClassRank.CompareOffOwn}
\begin{proof}
\openstep{The comparison of the trace lattice of a run with the base theta lattice away from the own
prime is not written out.}
% Lean: compareOffOwn_atRuns_alt: the scale is that of T.runNumericalConditions_hold_along_family 0
% (for 8fr + 4 <= n_V); ClassSumSpanCore.classSumsSpan_core gives the span of the class sums, and
% T.traceLattice_le_span_thetaSeries_awayFromOwnPrime the comparison, with the links, IoExt, the
% one-map rows, betaStd, qGuard, CoreTies.fieldAt and delAt_mem, RecTies.small, and the pins of r.
% The two are natural separate pieces. Size: about 271,000 lines of Lean in 16,000 theorems below them
% (text count, stopping at statements the book cites), about 20,000 of them below no other open step
% of Lemma is-onto-estimates.
\end{proof}

\begin{lemma}\label{is-onto-step-theta-j-ties}
Let $p$ be a point of a family that satisfies the point conditions and the further conditions, $r$ a
pair of runs at $p$ pinned to the own places, and $O$ a record object recorded at $(p,r)$. Let $R$ be
the realization of $p$, with the scalars of its characterization, and $F$ its face, with subgroup $P_0$
and character $\psi_0$, and let $\alpha_j$ and $\mathrm{inc}$ be as in Lemma~\ref{is-ties-level-tie}. Let
$v$ be a twisted place and $j$ a lift. Then
\begin{enumerate}
\item $\Theta_J(v,j)=\alpha_j^{-1}\big(\Theta^{\mathrm b}_v(j)\big)$;
\item $\Theta_0(v,j)=\alpha_j^{-1}(T)$, where $T$ is the trace lattice of the run used at $v$
  (Lemma~\ref{is-onto-step-compare}), formed with the delta data $\delta_{j+1}(n)$ for the integer $n$ of
  the weight family $\kappa_j$.
\end{enumerate}
\end{lemma}
\leanref{Development/AbelianVarietyTheory/Definitions/Foundations/ThetaSeriesTies.lean}{GlobalTheory.OnTies.OnTie}
\leanref{Development/AbelianVarietyTheory/Proofs/FoundationProofs/TraceLatticeLeThetaLatticeAtPlace.lean}{GlobalTheory.TraceClassRank.Th0_eq_comap}
\begin{proof}
Let $s$ be the run used at $v$ and $q$ the own prime of $v$. Part (1), and the description of
$\Theta_0(v,j)$ as the $\mathcal O_L$-span of the vectors $\sum_{u\in U}\psi(u)\,z_j^{-1}(1\otimes F_u)$,
where $(F_u)_{u\in U}$ runs over the orbit families of cusp forms of $s$ at level $j+1$ and weight
$\kappa_j$ for $\delta_{j+1}(n)$ and the action of $\ell_j\circ\mathrm{inc}$, are conditions of the level
tie of $O$ at $(v,j)$, which is part of being recorded; $z_j$ is the isomorphism of the tie
(Lemma~\ref{is-onto-step-index-bound-ties}(1)). The map $\alpha_j$ is the composite of $z_j$ with the
map $c$ of the weight tie, so $\alpha_j(z_j^{-1}(1\otimes F))=c(F)$ for every cusp form $F$. By the links
(Lemma~\ref{is-onto-recorded-numbers}(3) and its proof), $\mathrm{inc}$ is injective and its image
contains $P_0$, so by Lemma~\ref{is-ties-level-tie}(2) it maps $U$ bijectively onto $P_0$, with
$\psi(u)=\psi_0(\mathrm{inc}(u))^{-1}$. Reindexing an orbit family along this bijection gives an orbit
family for the action of $\ell_j$, and conversely. Hence
\[
 \alpha_j\Big(\sum_{u\in U}\psi(u)\,z_j^{-1}(1\otimes F_u)\Big)
 =\sum_{u'\in P_0}\psi_0(u')^{-1}c(F_{\mathrm{inc}^{-1}(u')}),
\]
and $\alpha_j$ maps the generators of $\Theta_0(v,j)$ onto those of $T$. As $\alpha_j$ is bijective,
$\Theta_0(v,j)=\alpha_j^{-1}(T)$.
% Lean: OnTie.ThJ_eq and OnTie.Th0_eq (thetaEig, a traceSpan over IsOrbitFamily with T.d delta (j + 1)
% (kappa j).n), with Th0_eq_comap (toForms_tmul, comapEquiv for inc, isOrbitFamily_comp) and
% toForms_bijective; the hypotheses hi and hU are mirIncl_injective and mem_base_iff (grp levels).
\end{proof}

\begin{definition}\label{is-onto-step-theta-j-low-types-vectors}
Let $\epsilon\in\{\pm1\}^n$, let $J_\epsilon$ be the diagonal matrix with entries $\epsilon$, and let
$U(J_\epsilon)=\{g\in GL_n(\mathbf C):g^*J_\epsilon g=J_\epsilon\}$, with a Haar measure $dg$. Let $\sigma$ be
a unitary representation of $U(J_\epsilon)$ on a Hilbert space $H$: $\sigma(1)=1$,
$\sigma(ab)=\sigma(a)\sigma(b)$, every $\sigma(g)$ preserves the norm, and $g\mapsto\sigma(g)w$ is
continuous for every $w\in H$. Let $P\in GL_n(\mathbf C)$, and put $G=P^*J_\epsilon P$,
$U(G)=\{g\in GL_n(\mathbf C):g^*Gg=G\}$, $\mathfrak u(G)=\{X\in M_n(\mathbf C):X^*G+GX=0\}$, a real vector
space, and $\rho(g)=\sigma(PgP^{-1})$ for $g\in U(G)$.
\begin{enumerate}
\item A vector $w\in H$ is \emph{smooth} if for all $X_1,\dots,X_k\in\mathfrak u(G)$ the map
  $t\mapsto\rho\big(\exp(t_1X_1)\cdots\exp(t_kX_k)\big)w$ from $\mathbf R^k$ to $H$ is smooth. It is
  \emph{strongly smooth} if $t\mapsto\rho(x(t))w$ is smooth for every $k$ and every smooth map
  $x\colon\mathbf R^k\to M_n(\mathbf C)$ with values in $U(G)$. Write $H^\infty$ for the set of smooth
  vectors.
\item For $X\in\mathfrak u(G)$ and $w\in H$ put $Xw=\frac{d}{dt}\rho(\exp(tX))w\big|_{t=0}$. For a
  \emph{word} $l=(X_1,\dots,X_k)$ in $\mathfrak u(G)$ put $lw=X_1(X_2(\cdots(X_kw)))$, with $lw=w$ for the
  empty word, and for $h\in U(G)$ put $l^h=(h^{-1}X_1h,\dots,h^{-1}X_kh)$.
\item A linear functional $\ell$ on $H$ is \emph{Sobolev-continuous} if there are finitely many words
  $l_1,\dots,l_m$ in $\mathfrak u(G)$ and $C\in\mathbf R$ with $|\ell(w)|\le C\sum_i\|l_iw\|$ for every
  $w\in H^\infty$.
\item For a continuous function $\varphi\colon M_n(\mathbf C)\to\mathbf R$ with compact support, the
  \emph{Gårding operator} is the linear map
  $\gamma_\varphi(w)=\int_{U(J_\epsilon)}\varphi(g)\,\sigma(g)w\,dg$ on $H$.
\end{enumerate}
\end{definition}
\leanref{Formalism/LocalToGlobal/BesselHermitianTheta/TiedType.lean}{Formalism.L2.Theta.SmoothVec}
\leanref{Formalism/LocalToGlobal/BesselHermitianTheta/TiedType.lean}{Formalism.L2.Theta.dRep}
\leanref{Formalism/LocalToGlobal/BesselHermitianTheta/TiedType.lean}{Formalism.L2.Theta.Sobolev₂}
\leanref{Development/AbelianVarietyTheory/Proofs/DivisorProofs/MultiplicitySpaceEqZeroOfLowestKTypes.lean}{GlobalTheory.LowTypeDetermined.smooth}
\leanref{Development/AbelianVarietyTheory/Proofs/DivisorProofs/MultiplicitySpaceEqZeroOfLowestKTypes.lean}{GlobalTheory.LowTypeDetermined.smoothAll}
\leanref{Development/AbelianVarietyTheory/Proofs/DivisorProofs/ExistsSmoothVecLowKSub.lean}{GlobalTheory.LowTypeRestriction.garding}
\leanref{Formalism/LocalToGlobal/IndefiniteReal/FockTheta/DiscreteSeriesUnitary.lean}{Formalism.IndefR.Fock.SquareIntegrableRepresentations.IsURepOn}

\begin{lemma}\label{is-onto-step-theta-j-low-types-smooth}
Let $\epsilon$, $\sigma$, $H$, $P$, $G$ and $\rho$ be as in
Definition~\ref{is-onto-step-theta-j-low-types-vectors}, and suppose that $H\ne0$ and that $\sigma$ is
irreducible: the only closed subspaces of $H$ stable under every $\sigma(g)$ are $0$ and $H$. Then:
\begin{enumerate}
\item $H^\infty$ is a subspace of $H$. For $w\in H^\infty$, $X\in\mathfrak u(G)$, $h\in U(G)$ and a word
  $l$, the vectors $Xw$ and $\rho(h)w$ are smooth, and $l\rho(h)w=\rho(h)\,l^hw$.
\item $\langle y,Xw\rangle=-\langle Xy,w\rangle$ for $y,w\in H^\infty$ and $X\in\mathfrak u(G)$. For every
  $y\in H^\infty$ and every word $l$ there is $y'\in H$ with $\langle y,lw\rangle=\langle y',w\rangle$ for
  all $w\in H^\infty$.
\item Let $\varphi\colon M_n(\mathbf C)\to\mathbf R$ be smooth with compact support. Then
  $\|\gamma_\varphi(w)\|\le\|w\|\int|\varphi(g)|\,dg$, and $t\mapsto\sigma(x(t))\gamma_\varphi(w)$ is smooth
  for every $w\in H$ and every smooth map $x\colon\mathbf R^k\to M_n(\mathbf C)$ with values in
  $U(J_\epsilon)$. If $w\ne0$, $\varphi\geq0$, $\varphi(1)=1$ and $\|\sigma(g)w-w\|<\|w\|/2$ for every
  $g\in U(J_\epsilon)$ with $\varphi(g)\ne0$, then $\gamma_\varphi(w)\ne0$.
\item The strongly smooth vectors are smooth and form a dense subspace of $H$; in particular $H^\infty$
  is dense. For $w\in H^\infty$ the map $X\mapsto Xw$ is $\mathbf R$-linear on $\mathfrak u(G)$.
\item Let $T$ be a topological space, $a\colon T\to M_n(\mathbf C)$ a continuous map with values in
  $U(G)$, $l$ a word and $w\in H^\infty$. Then $t\mapsto l\rho(a(t))w$ is continuous.
\item For every Sobolev-continuous linear functional $\ell$ there are words $l_0,\dots,l_m$ and continuous
  linear functionals $\Lambda_0,\dots,\Lambda_m$ on $H$ with $\ell(w)=\sum_i\Lambda_i(l_iw)$ for every
  $w\in H^\infty$.
\end{enumerate}
\end{lemma}
\leanref{Development/AbelianVarietyTheory/Proofs/DivisorProofs/MultiplicitySpaceEqZeroOfLowestKTypes.lean}{GlobalTheory.LowTypeDetermined.dRep_mem}
\leanref{Development/AbelianVarietyTheory/Proofs/DivisorProofs/MultiplicitySpaceEqZeroOfLowestKTypes.lean}{GlobalTheory.LowTypeDetermined.translate_mem}
\leanref{Development/AbelianVarietyTheory/Proofs/DivisorProofs/MultiplicitySpaceEqZeroOfLowestKTypes.lean}{GlobalTheory.LowTypeDetermined.word_translate}
\leanref{Development/AbelianVarietyTheory/Proofs/DivisorProofs/MultiplicitySpaceEqZeroOfLowestKTypes.lean}{GlobalTheory.LowTypeDetermined.inner_dRep}
\leanref{Development/AbelianVarietyTheory/Proofs/DivisorProofs/MultiplicitySpaceEqZeroOfLowestKTypes.lean}{GlobalTheory.LowTypeDetermined.exists_adjoint}
\leanref{Development/AbelianVarietyTheory/Proofs/DivisorProofs/MultiplicitySpaceEqZeroOfLowestKTypes.lean}{GlobalTheory.LowTypeDetermined.continuous_garding}
\leanref{Development/AbelianVarietyTheory/Proofs/DivisorProofs/ExistsSmoothVecLowKSub.lean}{GlobalTheory.LowTypeRestriction.contDiff_garding}
\leanref{Development/AbelianVarietyTheory/Proofs/DivisorProofs/ExistsSmoothVecLowKSub.lean}{GlobalTheory.LowTypeRestriction.garding_ne_zero}
\leanref{Development/AbelianVarietyTheory/Proofs/IwasawaDecomposition/ContDiffSetIntegralNat.lean}{GlobalTheory.GardingVec.exists_garding}
\leanref{Development/AbelianVarietyTheory/Proofs/DivisorProofs/MultiplicitySpaceEqZeroOfLowestKTypes.lean}{GlobalTheory.LowTypeDetermined.dRep_sum}
\leanref{Development/AbelianVarietyTheory/Proofs/DivisorProofs/MultiplicitySpaceEqZeroOfLowestKTypes.lean}{GlobalTheory.LowTypeDetermined.continuous_word_orbit}
\leanref{Development/AbelianVarietyTheory/Proofs/DivisorProofs/MultiplicitySpaceEqZeroOfLowestKTypes.lean}{GlobalTheory.LowTypeDetermined.exists_factor}
\begin{proof}
For $X\in\mathfrak u(G)$ and $t\in\mathbf R$, $\exp(tX)\in U(G)$, and $h^{-1}Xh\in\mathfrak u(G)$ for
$h\in U(G)$.

(1) Sums and multiples of smooth maps are smooth, so $H^\infty$ is a subspace. Let $w\in H^\infty$ and
$X,X_1,\dots,X_k\in\mathfrak u(G)$, and write $e(t)=\exp(t_1X_1)\cdots\exp(t_kX_k)$. The map
$(t,s)\mapsto\rho(e(t)\exp(sX))w$ on $\mathbf R^{k+1}$ is smooth, and $\rho(e(t))Xw$ is its derivative in $s$
at $s=0$, because $\rho(e(t))$ is continuous and linear and $\rho$ is multiplicative on $U(G)$. So it is
smooth in $t$, and $Xw$ is smooth. As $\exp(t_iX_i)h=h\exp(t_ih^{-1}X_ih)$,
\[
 \rho(e(t))\rho(h)w=\rho(h)\,\rho\big(\exp(t_1h^{-1}X_1h)\cdots\exp(t_kh^{-1}X_kh)\big)w,
\]
which is smooth in $t$; so $\rho(h)w$ is smooth. For $k=1$ the derivative at $t=0$ gives
$X\rho(h)w=\rho(h)(h^{-1}Xh)w$, and induction on the length of $l$ gives $l\rho(h)w=\rho(h)l^hw$.

(2) As $\rho(\exp(tX))$ preserves the inner product, $\langle\rho(\exp(tX))y,\rho(\exp(tX))w\rangle=\langle
y,w\rangle$ for all $t$, and the derivative at $t=0$ gives $\langle Xy,w\rangle+\langle y,Xw\rangle=0$. For
a word $(X,l)$, with $X$ first, $\langle y,X(lw)\rangle=-\langle Xy,lw\rangle$ with $Xy$ smooth by (1), so
induction on the length gives $y'$.

(3) The integrand of $\gamma_\varphi(w)$ has norm $|\varphi(g)|\,\|w\|$, which gives the bound. Let $x$ be
as in (3). By the left invariance of $dg$,
\[
 \sigma(x(t))\gamma_\varphi(w)=\int_{U(J_\epsilon)}\varphi\big(x(t)^{-1}g\big)\,\sigma(g)w\,dg ,
\]
and $x(t)^{-1}=J_\epsilon x(t)^*J_\epsilon$ is smooth in $t$. For $t$ near a point $t_0$ the integrand
vanishes outside a fixed compact set, so the integral is smooth in $t$ by differentiation under the
integral sign. This is Gårding's argument \cite{Gar47}. For the last claim, put
$I=\int\varphi(g)\,dg$; it is positive, as $\varphi\geq0$ is continuous with $\varphi(1)=1$. Then
$\|\gamma_\varphi(w)-Iw\|=\|\int\varphi(g)(\sigma(g)w-w)\,dg\|\le I\|w\|/2<I\|w\|$, so
$\gamma_\varphi(w)\ne0$.

(4) The map $t\mapsto\exp(t_1X_1)\cdots\exp(t_kX_k)$ is smooth with values in $U(G)$, so strongly smooth
vectors are smooth. They form a subspace, and it is stable under every $\sigma(g)$, $g\in U(J_\epsilon)$:
$\rho(x(t))\sigma(g)w=\rho(x(t)P^{-1}gP)w$, and $t\mapsto x(t)P^{-1}gP$ is smooth with values in $U(G)$. It
is nonzero. Indeed, let $w\ne0$; by the continuity of $g\mapsto\sigma(g)w$ at $1$ there is a smooth
$\varphi\geq0$ with compact support and $\varphi(1)=1$ such that $\|\sigma(g)w-w\|<\|w\|/2$ wherever
$\varphi(g)\ne0$. By (3), $\gamma_\varphi(w)\ne0$, and it is strongly smooth, since
$\rho(x(t))\gamma_\varphi(w)=\sigma(Px(t)P^{-1})\gamma_\varphi(w)$ and $Px(t)P^{-1}$ is smooth with values
in $U(J_\epsilon)$. The closure of the strongly smooth vectors is then a nonzero closed subspace stable
under $\sigma$, and it is $H$ by irreducibility.

For the linearity let $B_1,\dots,B_k\in\mathfrak u(G)$ and $c\in\mathbf R^k$. For a strongly smooth $y$, the
map $F(t)=\rho(\exp(\sum_it_iB_i))y$ on $\mathbf R^k$ is smooth, and $(\sum_ic_iB_i)y$ is the derivative of
$s\mapsto F(sc)$ at $0$, which is $dF(0)(c)$; so $(\sum_ic_iB_i)y=\sum_ic_iB_iy$. For $w\in H^\infty$ and
strongly smooth $y$, (2) gives
$\langle y,(\sum_ic_iB_i)w\rangle=-\langle(\sum_ic_iB_i)y,w\rangle=-\sum_ic_i\langle B_iy,w\rangle
=\sum_ic_i\langle y,B_iw\rangle$. As the strongly smooth vectors are dense,
$(\sum_ic_iB_i)w=\sum_ic_iB_iw$.

(5) By (1), $l\rho(a(t))w=\rho(a(t))\,v(t)$ with $v(t)=l^{a(t)}w$. First, $t\mapsto f(l^{a(t)}w)$ is
continuous for every word $l$ and every linear map $f\colon H^\infty\to H$, by induction on the length
of $l$. For the empty word it is constant. For a word $(X,l')$, choose a basis $B_1,\dots,B_d$ of
$\mathfrak u(G)$ and write $a(t)^{-1}Xa(t)=\sum_ic_i(t)B_i$; the $c_i$ are continuous, as
$a(t)^{-1}=G^{-1}a(t)^*G$. By (4), $f\big((a(t)^{-1}Xa(t))\,l'^{a(t)}w\big)=\sum_ic_i(t)\,f\big(B_i\,l'^{a(t)}w\big)$,
and $f\circ B_i$ is a linear map $H^\infty\to H$ by (1), so the induction applies. Taking for $f$ the
inclusion, $v$ is continuous. Then
\[
 \rho(a(t))v(t)-\rho(a(t_0))v(t_0)=\rho(a(t))\big(v(t)-v(t_0)\big)+\big(\rho(a(t))-\rho(a(t_0))\big)v(t_0);
\]
the first term has norm $\|v(t)-v(t_0)\|$, and the second tends to $0$ as $t\to t_0$, because
$\rho(a(t))v(t_0)=\sigma(Pa(t)P^{-1})v(t_0)$ is continuous in $t$.

(6) Let $l_1,\dots,l_m$ and $C$ be as in the definition, let $l_0$ be the empty word, and give $H^{m+1}$
the norm $\max_i\|z_i\|$. Let $\Phi\colon H^\infty\to H^{m+1}$ be $w\mapsto(l_0w,\dots,l_mw)$. On the image of
$\Phi$ the map $f(z)=\ell(z_0)$ is linear, $f(\Phi(w))=\ell(w)$, and
$|f(\Phi(w))|\le|C|\sum_{i\geq1}\|l_iw\|\le|C|\,m\,\|\Phi(w)\|$. By the Hahn-Banach theorem $f$ extends to a
continuous linear functional $g$ on $H^{m+1}$. Let $\Lambda_i(x)=g(0,\dots,0,x,0,\dots,0)$, with $x$ in
the place $i$. Then $\ell(w)=g(\Phi(w))=\sum_i\Lambda_i(l_iw)$.
% Lean: (1) is smooth (a submodule), dRep_mem (contDiff_deriv_snoc), translate_mem (conj_mem_lie,
% exp_mul_eq, prod_mul_eq) and word_translate; (2) is inner_dRep and exists_adjoint; (3) is
% continuous_garding, contDiff_garding (garding_translate, contDiff_setIntegral_nat) and garding_ne_zero;
% (4) is smoothAll_le (contDiff_expProd, expProd_mem), the density step of eq_zero_of_translates
% (dense_of_irred with GardingVec.exists_garding, which repeats the argument of (3) with
% exists_contDiff_tsupport_subset), dRep_eq_fderiv, dRep_sum_all and dRep_sum; (5) is continuous_word
% (Module.finBasis of lie G) and continuous_word_orbit; (6) is exists_factor
% (exists_extension_norm_eq for the Hahn-Banach extension).
\end{proof}

\begin{lemma}\label{is-onto-step-theta-j-low-types-translates}
Let $\epsilon$, $\sigma$, $H$, $P$, $G$ and $\rho$ be as in
Definition~\ref{is-onto-step-theta-j-low-types-vectors}, with $\sigma$ irreducible. Let $\ell$ be a
Sobolev-continuous linear functional on $H$ and $u\ne0$ a smooth vector with $\ell(\sigma(g)u)=0$ for
every $g\in U(J_\epsilon)$. Then $\ell(w)=0$ for every smooth vector $w$.
\end{lemma}
\leanref{Development/AbelianVarietyTheory/Proofs/DivisorProofs/MultiplicitySpaceEqZeroOfLowestKTypes.lean}{GlobalTheory.LowTypeDetermined.eq_zero_of_translates}
\leanref{Development/AbelianVarietyTheory/Proofs/DivisorProofs/MultiplicitySpaceEqZeroOfLowestKTypes.lean}{GlobalTheory.LowTypeDetermined.eq_zero_of_bumps}
\begin{proof}
As $u\ne0$, $H\ne0$, so Lemma~\ref{is-onto-step-theta-j-low-types-smooth} applies; write (1) to (6) for
its parts. For $g\in U(J_\epsilon)$, $P^{-1}gP\in U(G)$ and $\sigma(g)=\rho(P^{-1}gP)$, so $\sigma(g)$ maps
$H^\infty$ into itself by (1). Let $w\in H^\infty$, and put $F(g)=\ell(\sigma(g)w)$ for $g\in U(J_\epsilon)$.

\emph{Continuity of $F$.} By (6) there are words $l_i$ and continuous functionals $\Lambda_i$ with
$F(g)=\sum_i\Lambda_i\big(l_i\rho(P^{-1}gP)w\big)$, which is continuous by (5).

\emph{Smoothing.} Let $\varphi\colon M_n(\mathbf C)\to\mathbf R$ be smooth with compact support. For a word
$l$ and $v\in H^\infty$, $g\mapsto\varphi(g)\,l\sigma(g)v$ is continuous with compact support by (5), and
\[
 l\gamma_\varphi(v)=\int_{U(J_\epsilon)}\varphi(g)\,l\sigma(g)v\,dg .
\]
Indeed, $\gamma_\varphi(v)$ is smooth by (3) and (4). For $y\in H^\infty$ and $y'$ as in (2),
$\langle y,l\gamma_\varphi(v)\rangle=\langle y',\gamma_\varphi(v)\rangle=\int\varphi(g)\langle
y',\sigma(g)v\rangle\,dg=\int\varphi(g)\langle y,l\sigma(g)v\rangle\,dg$, and $H^\infty$ is dense by (4). With
(6) this gives
\[
 \ell(\gamma_\varphi(v))=\int_{U(J_\epsilon)}\varphi(g)\,\ell(\sigma(g)v)\,dg\qquad(v\in H^\infty).
\]

\emph{Continuity of $\ell\circ\gamma_\varphi$.} On $H$, $\ell(\gamma_\varphi(v))=\sum_i\Lambda_i(A_iv)$ with
the linear maps $A_iv=l_i\gamma_\varphi(v)$ from $H$ to $H$. Each $A_i$ has a closed graph: if $v_j\to v$
and $A_iv_j\to z$, then for $y\in H^\infty$, with $y'$ as in (2), $\langle y,A_iv_j\rangle=\langle
y',\gamma_\varphi(v_j)\rangle$ tends to $\langle y',\gamma_\varphi(v)\rangle=\langle y,A_iv\rangle$, as
$\gamma_\varphi$ is continuous by (3); so $\langle y,z\rangle=\langle y,A_iv\rangle$ for a dense set of $y$,
and $z=A_iv$. By the closed graph theorem the $A_i$ are continuous, and so is $\ell\circ\gamma_\varphi$.

\emph{Vanishing.} The span of the vectors $\sigma(b)u$, $b\in U(J_\epsilon)$, is nonzero and stable under
$\sigma$, so it is dense by irreducibility. Each $\sigma(b)u$ is smooth, and by the smoothing formula
$\ell(\gamma_\varphi(\sigma(b)u))=\int\varphi(g)\,\ell(\sigma(gb)u)\,dg=0$. So the continuous functional
$\ell\circ\gamma_\varphi$ vanishes on a dense subspace, hence on $H$, and the smoothing formula for $v=w$
gives $\int\varphi(g)F(g)\,dg=0$.

\emph{Conclusion.} Suppose that $F(1)\ne0$. As $F$ is continuous, there is a smooth $\varphi\geq0$ with
compact support and $\varphi(1)=1$ such that $|F(g)-F(1)|<|F(1)|/2$ wherever $\varphi(g)\ne0$. With
$I=\int\varphi(g)\,dg>0$,
\[
 I\,|F(1)|=\Big|\int\varphi(g)\big(F(g)-F(1)\big)\,dg\Big|\le I\,|F(1)|/2,
\]
a contradiction. So $\ell(w)=F(1)=0$.
% Lean: eq_zero_of_translates: hd (dense_of_irred, exists_garding), htr (translate_mem), hcont
% (continuous_word_orbit), exists_factor and hF; hex and hform for the smoothing formula (exists_adjoint,
% integral_inner, ContinuousLinearMap.integral_comp_comm); hbd by continuous_of_seq_closed_graph and
% continuous_garding; hdu (dense_of_irred for the span of the translates of u) and hz by
% Continuous.ext_on; the conclusion is eq_zero_of_bumps (exists_bump of LowTypeRestriction).
\end{proof}

\begin{lemma}\label{is-onto-step-theta-j-low-types}
Let $D$ be a definite-side datum, $\theta$ an archimedean family that satisfies the core conditions on
the models, $V$ a Hermitian space of $D$, and $M$ a model of discrete series of $U(V)$ whose lowest
$K$-types are nonzero (Definition~\ref{asm-dependence-rows-models}(1)). Let $U$ be an element of a
multiplicity space of $M$ (Definition~\ref{asm-dependence-rows-models}(2)), a map from the tensors of
smooth vectors of $M$ to functions on $U(V)(\mathbf A_{F_h})$ with a level and a character. If $U(v)=0$
for every tensor $v$ of smooth vectors in the lowest $K$-types, then $U=0$. The hypotheses on $\theta$
and on the lowest $K$-types are not used in the proof.
\end{lemma}
\leanref{Development/AbelianVarietyTheory/Proofs/DivisorProofs/MultiplicitySpaceEqZeroOfLowestKTypesStatement.lean}{GlobalTheory.T.multiplicitySpace_eq_zero_of_lowestKTypes}
\leanref{Development/AbelianVarietyTheory/BlockTheory/DivisorBase/ResidueInjectivityAt.lean}{GlobalTheory.IndefiniteRowSideFullHecke.ResidueInjectivityAt.ResInjAt}
\leanref{Development/AbelianVarietyTheory/Proofs/DivisorProofs/MultiplicitySpaceEqZeroOfLowestKTypes.lean}{GlobalTheory.LowTypeDetermined.eq_zero_of_lowT}
\leanref{Development/AbelianVarietyTheory/Counterexamples/RotationExamples/SobolevContinuityLowGeneration.lean}{GlobalTheory.Survey.SobolevContinuityLowGeneration.mult_equiv_cont}
\leanref{Development/AbelianVarietyTheory/Counterexamples/RotationExamples/SobolevContinuityLowGeneration.lean}{GlobalTheory.Survey.SobolevContinuityLowGeneration.eq_zero_of_lowGenerates}
\leanref{Development/AbelianVarietyTheory/Proofs/DivisorProofs/MultiplicitySpaceEqZeroOfLowestKTypes.lean}{GlobalTheory.LowTypeDetermined.lowGenerates}
\leanref{Development/AbelianVarietyTheory/Proofs/DivisorProofs/ExistsSmoothVecLowKSub.lean}{GlobalTheory.LowTypeRestriction.exists_smoothVec_lowKSub}
\begin{proof}
Let $\mathcal T$ be the set of embeddings of $D$. For $\tau\in\mathcal T$ let $H_\tau$, $\sigma_\tau$,
$\epsilon_\tau$, $\lambda_\tau$ and $f_\tau$ be the data of $M$ at $\tau$, with $G_\tau=G_V(\tau)
=f_\tau^*J_{\epsilon_\tau}f_\tau$ and $\rho_\tau(g)=\sigma_\tau(f_\tau gf_\tau^{-1})$ for $g\in U(G_\tau)$.
Then $\epsilon_\tau$, $\sigma_\tau$, $H_\tau$ and $P=f_\tau$ are as in
Definition~\ref{is-onto-step-theta-j-low-types-vectors}, with $G=G_\tau$ and $\rho=\rho_\tau$; $\sigma_\tau$ is
irreducible; and the smooth vectors of $M$ at $\tau$ are the smooth vectors of that definition. Let
$K_\tau$ be the group of unitary matrices that commute with $J_{\epsilon_\tau}$, contained in
$U(J_{\epsilon_\tau})$, and $L_\tau\subset H_\tau$ the lowest $K$-type of $\sigma_\tau$: the sum of its
$K_\tau$-types (finite-dimensional $K_\tau$-stable subspaces, irreducible under $K_\tau$) with the highest
weight attached to $(\lambda_\tau,\epsilon_\tau)$. For a tuple $v=(v_\tau)$ of smooth vectors, $\tau\in\mathcal
T$ and a smooth $w\in H_\tau$, write $v[\tau\mapsto w]$ for $v$ with $v_\tau$ replaced by $w$.

\emph{Step 1: the multiplicity space.} $U$ is a linear combination of multiplicity maps. The following
two conditions hold for multiplicity maps by definition, and they are preserved by sums and scalar
multiples; for (b), a sum of Sobolev-continuous functionals is Sobolev-continuous, with the union of the
words and the larger constant. So $U$ satisfies them.
\begin{enumerate}
\item[(a)] For $\tau\in\mathcal T$, $g\in U(G_\tau)$ and a tuple $v$ of smooth vectors,
  $U(v[\tau\mapsto\rho_\tau(g)v_\tau])$ is the right translate of $U(v)$ by the image of $g$ under $\theta$.
\item[(b)] For $\tau\in\mathcal T$, a tuple $v$ of smooth vectors and a point $x$ of $U(V)(\mathbf A_{F_h})$,
  there is a Sobolev-continuous linear functional $\ell$ on $H_\tau$ with $\ell(w)=U(v[\tau\mapsto w])(x)$ for
  every smooth $w\in H_\tau$.
\end{enumerate}

\emph{Step 2: a smooth vector in the lowest $K$-type.} Fix $\tau$, and drop it from the notation. The
model gives a $K$-type $S$ with the highest weight attached to $(\lambda,\epsilon)$; let $w\in S$ be
nonzero. By the continuity of $g\mapsto\sigma(g)w$ at $1$ there is a smooth function $\varphi\geq0$ with
compact support, a function of $\operatorname{tr}((x-1)^*(x-1))$, with $\varphi(1)=1$ and
$\|\sigma(g)w-w\|<\|w\|/2$ wherever $\varphi(g)\ne0$. It satisfies $\varphi(k^*xk)=\varphi(x)$ for unitary
$k$. By Lemma~\ref{is-onto-step-theta-j-low-types-smooth}(3) and (4), $\gamma_\varphi(w)$ is a nonzero smooth
vector. For $k\in K$, $k^*=k^{-1}$, and the invariance of $\varphi$ and of $dg$ under left and right
translations give $\sigma(k)\gamma_\varphi(v)=\gamma_\varphi(\sigma(k)v)$ for every $v$. So $\gamma_\varphi$
restricts to a nonzero $K$-map on $S$; its kernel is a $K$-stable subspace other than $S$, hence $0$ by
irreducibility, and $\gamma_\varphi(S)$ is a $K$-type isomorphic to $S$, with the same highest weight. Hence
$\gamma_\varphi(w)$ is a nonzero smooth vector in $L$.

\emph{Step 3: the translates of the lowest $K$-type.} Let $\tau\in\mathcal T$ and let $\ell$ be a
Sobolev-continuous linear functional on $H_\tau$ with $\ell(\rho_\tau(g)u)=0$ for every smooth $u\in L_\tau$
and every $g\in U(G_\tau)$. Then $\ell(w)=0$ for every smooth $w$. Indeed, let $u\in L_\tau$ be smooth and
nonzero (Step 2). For $g\in U(J_{\epsilon_\tau})$, $h=f_\tau^{-1}gf_\tau$ lies in $U(G_\tau)$ and
$\sigma_\tau(g)u=\rho_\tau(h)u$, so $\ell(\sigma_\tau(g)u)=0$, and
Lemma~\ref{is-onto-step-theta-j-low-types-translates} applies.

\emph{Step 4: induction.} For $S\subset\mathcal T$ call a tuple $v$ of smooth vectors \emph{low off} $S$ if
$v_\tau\in L_\tau$ for every $\tau\notin S$. We show by induction on $|S|$ that $U(v)=0$ for every $v$ low
off $S$. For $S=\emptyset$ this is the hypothesis. Let $S=S'\cup\{\tau\}$ with $\tau\notin S'$, let $v$ be
low off $S$, and let $x$ be a point. Let $\ell$ be the functional of (b) for $\tau$, $v$ and $x$. For a
smooth $u\in L_\tau$ and $g\in U(G_\tau)$, the tuple $v'=v[\tau\mapsto u]$ is low off $S'$, so $U(v')=0$ by
induction, and by (a) $U(v[\tau\mapsto\rho_\tau(g)u])=U(v'[\tau\mapsto\rho_\tau(g)v'_\tau])$ is a translate of
$U(v')=0$. So $\ell(\rho_\tau(g)u)=0$, and by Step 3, $U(v)(x)=\ell(v_\tau)=0$. For $S=\mathcal T$ this
gives $U(v)=0$ for every tuple of smooth vectors, and $U=0$, as $U$ is multilinear.
% Lean: T.multiplicitySpace_eq_zero_of_lowestKTypes = LowTypeDetermined.eq_zero_of_lowT
% = SobolevContinuityLowGeneration.eq_zero_of_lowGenerates (lowGenerates M): Step 1 is mult_equiv_cont
% (span_induction, sobCont_add, sobCont_smul), Step 4 the Finset.induction_on in eq_zero_of_lowGenerates.
% Step 3 is lowGenerates (conj_inv_mem, conj_conj, eq_zero_of_translates); its u comes from
% LowTypeRestriction.lowT_ne_bot, which is Step 2 (exists_smoothVec_lowKSub: exists_bump with distOne,
% garding_ne_zero, garding_comm with indefinite_unitary_unimodular, isKTypeOf_map, contDiff_garding). The
% hypotheses IsInfFam theta and lowT M tau != bot of the Lean statement are not used. The Lean reads U
% on the smooth vectors (M.sm tau, the span of SmoothVec) and the lowest K-type vectors as lowT M tau,
% the smooth vectors of lowKSub. ResInj, the form used in Lemma is-onto-step-theta-j-determined, is
% LowTypeRestriction.resInj_of_determined, and ResInj.of_char for a realization with IsInfFam
% (hC.models.inf).
\end{proof}

\begin{lemma}\label{is-onto-step-theta-j-determined}
Let $R$ be a realization of a definite-side datum with a face $F$ that satisfies its characterization at
$F$ and condition (11) of Definition~\ref{is-primitive-rows}, and let $j$ be a lift, with $\Phi_j$ as in
Definition~\ref{is-ties-class-sums}(1). If $z\in L^{n_W(j)}$ and $\Phi_j(z)(v)=0$ for every tensor $v$ of
lowest $K$-type vectors of the model of $W_j$ at $\sigma_0$, then $z=0$.
\end{lemma}
\leanref{Development/AbelianVarietyTheory/Proofs/FoundationProofs/KernelSpanOverClasses.lean}{GlobalTheory.TraceClassRank.formsDetermined_of}
\leanref{Development/AbelianVarietyTheory/Proofs/FoundationProofs/KernelSpanOverClasses.lean}{GlobalTheory.TraceClassRank.FormsDetermined}
\begin{proof}
By the core conditions on the models, part of the characterization, the map that sends a vector of
the Hecke module $H_W(j,\sigma_0)$ to its automorphic form is injective, with image the multiplicity
space of the model of $W_j$ at $\sigma_0$. So $\Phi_j(z)$ lies in this multiplicity space, and
Lemma~\ref{is-onto-step-theta-j-low-types} gives $\Phi_j(z)=0$. Hence $j_{\sigma_0}(z)=0$. By (11),
$j_{\sigma_0}$ carries $L$-linearly independent families to $\mathbf C$-linearly independent ones; if
$z\ne0$, the family $(z)$ is independent and $(j_{\sigma_0}(z))=(0)$ is not. So $z=0$.
% Lean: formsDetermined_of with hres = ResInj.of_char LowTypeRestriction.resInj, hW = hC.models.HW (h1
% injectivity, h2 the range) and hu = LUnif, the second link of the weight rows.
\end{proof}

\begin{lemma}\label{is-onto-step-theta-j-fg}
Let $R$ be a realization of a definite-side datum with a face $F$, with subgroup $P_0$ and character
$\psi_0$, let $j$ be a lift and $v$ a twisted place, and suppose that the conclusion of
Lemma~\ref{is-onto-step-theta-j-determined} holds for $R$ and $j$. Suppose that the base theta lattice
$\Theta^{\mathrm b}_v(j)$ is the span of $\bigcup_c\Theta_\kappa(h_c)$ for every base theta kernel
$\kappa$ at $(j,v)$ and every family $(h_c)$ that meets the classes of the level $J_v$
(Definition~\ref{is-onto-step-kernels}). Then $\Theta^{\mathrm b}_v(j)$ is a finitely generated
$\mathcal O_L$-module.
\end{lemma}
\leanref{Development/AbelianVarietyTheory/Proofs/FoundationProofs/KernelSpanOverClasses.lean}{GlobalTheory.TraceClassRank.thetaBase_fg}
\leanref{Development/AbelianVarietyTheory/Proofs/FoundationProofs/KernelSpanOverClasses.lean}{GlobalTheory.TraceClassRank.seriesAt_fg}
\leanref{Development/AbelianVarietyTheory/Proofs/FoundationProofs/KernelSpanOverClasses.lean}{GlobalTheory.TraceClassRank.fg_of_graph}
\begin{proof}
If there is no base theta kernel at $(j,v)$, then $\Theta^{\mathrm b}_v(j)=0$. Otherwise let
$\kappa=(\omega,\varphi)$ be one, and use the notation of Definition~\ref{is-ties-lattice-systems} for the
tuple $t=\mu_v$: the representatives $Y\subset G_f$ of the presentation of the lift forms at $(j,t)$,
with the level $J$ of these forms, the rational coefficient space $\mathcal V$, the integral vectors
$\mathcal N_y$ and the map $\mathrm{coef}$. Every element of $G_f$ is $\gamma yk$ with $\gamma\in\Gamma$,
$y\in Y$ and $k\in J$, and $J\subset J_v$, since $J_v$ is the level of the data at the subset $Q_v$ of the
ramified set $Q_t$. So the family $(y)_{y\in Y}$ meets the classes of $J_v$, and
$\Theta^{\mathrm b}_v(j)$ is the span of the finite union $\bigcup_{y\in Y}\Theta_\kappa(y)$. It suffices
to show that the span of $\Theta_\kappa(y)$ is finitely generated for each $y\in Y$.

Fix $y$. If $w$ is integral for the lattice of $y$, then $w=\iota_{\lambda'}(P)$ for a family tie and
some $P\in\nabla_{\lambda'}(L_y^\vee)=e(\mathcal N_y)$, so $w=\mathrm{coef}(Q)$ with $Q\in\mathcal N_y$,
as $\mathrm{coef}=\iota_{\lambda'}\circ e$ (Definition~\ref{is-ties-lattice-systems}(4)). The module
$\mathcal N_y$ is finitely generated, as $e$ is injective and $\nabla_{\lambda'}(L_y^\vee)$ is a lattice in
$V_{\lambda'}(L)$. Write $\operatorname{tr}$ for the trace $\sum_{g\in P_0}\sigma_D(\psi_0(g))\,g$ and
$\vartheta(y,w)$ for the theta series of $\kappa$ at $y$ for the polynomial of $w$, and let $\Gamma_y$ be
the set of pairs $(Q,x)\in\mathcal V\times L^{n_W(j)}$ with $\Phi_j(x)(v)=\operatorname{tr}\vartheta(y,
\mathrm{coef}(Q))(v)$ for every tensor $v$ of lowest $K$-type vectors. As $\mathrm{coef}$ and $\Phi_j$ are
$\sigma_0$-semilinear and $\vartheta$ and $\operatorname{tr}$ are linear, $\Gamma_y$ is an $L$-subspace.
If $(0,x)\in\Gamma_y$, then $\Phi_j(x)(v)=0$ for every $v$, and $x=0$ by hypothesis. So the first
projection is injective on the $\mathcal O_L$-module $\Gamma_y\cap(\mathcal N_y\times L^{n_W(j)})$, which
is therefore isomorphic to a submodule of $\mathcal N_y$, finitely generated over the noetherian ring
$\mathcal O_L$. Its second projection is finitely generated and contains $\Theta_\kappa(y)$.
% Lean: thetaBase_fg (with RepsCompare.existsUnique_lift for the h_c, the reps of the presentation and
% levJ_le_dataLevel), seriesAt_fg (exists_coef_fg through coefTie_class and nabAtr_hLat_fg, the graph
% Gamma, and fg_of_graph). The Lean proves it for any (U, psi, phase); the book's case is (P0, psi0^-1, 1).
\end{proof}

\begin{lemma}\label{is-onto-step-theta-j-type}
Let $R$ be a realization of a definite-side datum with a face $F$, with group $G$, subgroup $P_0$ and
character $\psi_0$, let $j$ be a lift and $v$ a twisted place, with the notation of
Definition~\ref{is-ties-class-sums}. Suppose that $R$ satisfies the core conditions at $F$, condition
(11) of Definition~\ref{is-primitive-rows} and hypothesis (2) of Lemma~\ref{is-onto-step-kernel-span};
that the conclusion of Lemma~\ref{is-onto-step-theta-j-determined} holds for $R$ and $j$; that the values
of $\delta_v$ lie in $\sigma_D(L)$; and that $\chi_j$ is trivial on $P_0$. Then
$\varrho_j(u)x=\psi_0(u)^{-1}x$ for every $x\in\Theta^{\mathrm b}_v(j)$ and $u\in P_0$.
\end{lemma}
\leanref{Development/AbelianVarietyTheory/Proofs/DivisorProofs/BaseTupleThetaSpanHasTypeStatement.lean}{GlobalTheory.T.baseTupleThetaSpan_hasType}
\leanref{Development/AbelianVarietyTheory/Proofs/DivisorProofs/BaseTupleThetaSpanHasType.lean}{GlobalTheory.BaseKernelType.thetaBase_hasType}
\leanref{Development/AbelianVarietyTheory/Proofs/DivisorProofs/BaseTupleThetaSpanHasType.lean}{GlobalTheory.BaseKernelType.singleSeries_hasType}
\leanref{Development/AbelianVarietyTheory/Proofs/DivisorProofs/BaseTupleThetaSpanHasType.lean}{GlobalTheory.BaseKernelType.actL_levelTrace}
\leanref{Formalism/LocalToGlobal/IndefiniteReal/SubgroupCharExtensionOneApply.lean}{Formalism.IndefR.Lev.actG0}
\leanref{Formalism/LocalToGlobal/IndefiniteReal/SubgroupCharExtensionOneApply.lean}{Formalism.IndefR.Lev.level_action_spec}
\leanref{Development/AbelianVarietyTheory/Proofs/TwistedOperators/LevelGroupActionEqZeroOfNoLift.lean}{GlobalTheory.ProjComm.levelGroupAction_mul_of_lifts}
\leanref{Development/AbelianVarietyTheory/Proofs/TwistedOperators/LevelGroupActionEqZeroOfNoLift.lean}{GlobalTheory.ProjComm.levelGroupAction_eq_zero_of_no_lift}
\begin{proof}
Write $\Phi_j$ as in Definition~\ref{is-ties-class-sums}(1); it is semilinear for $\sigma_0=\sigma_D$. The
action $g\vartheta$ of $g\in G$ on functions on $U(W_j)(\mathbf A_{F_h})$ in the traces of
Definition~\ref{is-onto-step-kernels}(2) is the level action. For each $g$ that has a level lift
(Lemma~\ref{is-eigenvalue-core-group}) fix one, $k_g$, and put $(g\vartheta)(h)=\chi_K(k_g)^{-1}\vartheta(hk_g)$,
where $\chi_K$ is the level character of the trivial level datum of $W_j$ on its parahoric level $K_W$;
put $g\vartheta=0$ if $g$ has no level lift. The character $\chi_K$ is a product of characters of blocks of
the reductions at the places of $P_g$ and at the ramified places of the datum; the trivial level datum
has no ramified places and no blocks with characters, so $\chi_K=1$. Say that $\vartheta$ \emph{has type}
$1$ if $\vartheta(hk)=\vartheta(h)$ for every $h$ and every level lift $k$ of $1$.

\emph{Step 1: the level action on functions of type $1$.} Let $\vartheta$ have type $1$. If $k$ is any level
lift of $g$, then $k_g^{-1}k$ is a level lift of $1$, as the reductions of elements of $K_W$ multiply; so
$(g\vartheta)(h)=\vartheta(hk_g)=\vartheta(hk_g\cdot k_g^{-1}k)=\vartheta(hk)$. Let $g,g'$ have level lifts $k_g$ and
$k'$. Then $k_gk'$ is a level lift of $gg'$, and
\[
 \big(g(g'\vartheta)\big)(h)=(g'\vartheta)(hk_g)=\vartheta(hk_gk')=\big((gg')\vartheta\big)(h).
\]
If $g$ has a level lift $k$ and $g'$ has none, then $gg'$ has none, since otherwise $g'=g^{-1}(gg')$ would
have the level lift $k^{-1}k''$; so $g(g'\vartheta)=0=(gg')\vartheta$. Hence, if $u\in P_0$ has a level lift,
\[
 u\sum_{g\in P_0}\sigma_D(\psi_0(g))\,g\vartheta=\sum_{g\in P_0}\sigma_D(\psi_0(g))\,(ug)\vartheta
 =\sigma_D(\psi_0(u))^{-1}\sum_{g\in P_0}\sigma_D(\psi_0(g))\,g\vartheta,
\]
by the substitution $g\mapsto u^{-1}g$ in $P_0$.

\emph{Step 2: the theta series have type $1$.} Let $\kappa=(\omega,\varphi)$ be a base theta kernel at
$(j,v)$. By Lemma~\ref{is-onto-step-kernel-span-twisted-base}, which uses the core conditions and hypothesis (2)
of Lemma~\ref{is-onto-step-kernel-span}, there are $\chi$ and $(\xi_P)$ such that $(\chi,(\xi_P),\omega,\varphi)$
is a twisted kernel of $R$ at $j$ for $(\delta_v,Q_v)$. By Lemma~\ref{is-onto-step-kernel-span-level-vector}(2),
for this twisted kernel and with the values of $\delta_v$ in $\sigma_D(L)$, $\omega(1,k)\varphi=\varphi$ for every
level lift $k$ of $1$, read with archimedean component $1$. By
Lemma~\ref{is-eigenvalue-conjugate-theta-facts}(3) and (4), $\varphi$ and $\varphi\otimes P$ are theta
functions for every polynomial $P$, and $\omega(1,k)(\varphi\otimes P)=(\omega(1,k)\varphi)\otimes P=\varphi\otimes P$,
as the archimedean components of $(1,k)$ are $1$. So for the theta series $\vartheta$ of $\omega$ and $\varphi$ at
$h'\in U(V_j)(\mathbf A_{F_h})$ for the polynomial $P$ of a coefficient vector and a tensor of lowest
$K$-type vectors,
\[
 \vartheta(hk)=\big(\omega(h',h)\,\omega(1,k)(\varphi\otimes P)\big)(0)=\vartheta(h).
\]
So $\vartheta$ has type $1$.

\emph{Step 3: a generator.} Let $u\in P_0$, $\kappa$ a base theta kernel at $(j,v)$ and $y\in\Theta_\kappa$:
there are $h'$ and $w$ such that $\Phi_j(y)(t)=\sum_{g\in P_0}\sigma_D(\psi_0(g))\,g\vartheta_t$ for every tensor
$t$ of lowest $K$-type vectors, where $\vartheta_t$ is the theta series of $\kappa$ at $h'$ for the polynomial of
$w$ and $t$. By condition (11), $\Phi_j(\varrho_j(u)y)(t)=\sigma_0(\chi_j(u))\,u\big(\Phi_j(y)(t)\big)$, and
$\chi_j(u)=1$. If $u$ has a level lift, Steps 1 and 2 give
\[
 \Phi_j(\varrho_j(u)y)(t)=\sigma_D(\psi_0(u))^{-1}\Phi_j(y)(t)=\Phi_j\big(\psi_0(u)^{-1}y\big)(t)
\]
for every $t$, so $\Phi_j(\varrho_j(u)y-\psi_0(u)^{-1}y)$ vanishes on the tensors of lowest $K$-type vectors, and
$\varrho_j(u)y=\psi_0(u)^{-1}y$ by the conclusion of Lemma~\ref{is-onto-step-theta-j-determined}. If $u$ has
no level lift, $\Phi_j(\varrho_j(u)y)(t)=0$ for every $t$, so $\varrho_j(u)y=0$, and $y=\varrho_j(u^{-1})\varrho_j(u)y=0$;
the identity holds again.

\emph{Step 4: the span.} The base theta lattice is the $\mathcal O_L$-span of the sets $\Theta_\kappa$, and both
sides of $\varrho_j(u)x=\psi_0(u)^{-1}x$ are $\mathcal O_L$-linear in $x$.
% Lean: T.baseTupleThetaSpan_hasType = BaseKernelType.thetaBase_hasType. Step 4 is span_induction over
% thetaBase; Step 3 is hgen: hy from seriesSet (with phase 1 and levelTrace R i P0 psi0^-1, whose
% coefficients are embLL(psi0 g)), hrho from the third clause of LUnif ((hU i (sg0 D A)).2.2 u y) with
% hchi u, then adForm_sub, adForm_smul and FormsDetermined (hd) if u has a level lift, and
% levelGroupAction_eq_zero_of_no_lift otherwise. Step 2 is singleSeries_hasType: hfix is the second clause
% of T.thetaKernel_vector_moved_by_rootsOfUnity for the KerDat of T.baseTuple_thetaKernel_tameCharacter
% (Lemma is-onto-step-kernel-span-level-vector(2)), hphi from VecTheta.vecThetaSD, levChi_flag1 (rfl rfl:
% s = 0 and Q = empty for datWp i triv), and OrientedHeis.om_fin_pureTensorO. Step 1 is actL_levelTrace:
% levelGroupAction_mul_of_lifts (level_action_spec, isLift_mul, levChi_flagL_mul_of_mem_parLev),
% isLift_inv and isLift_mul for the case without a lift, and Fintype.sum_bijective for the substitution.
% The level action is actL (Wc D i) (R.betaW i) (R.datWp i R.triv) (R.eG g) = actG0, with chi_K =
% levChi for flagL, which is 1 by EqCut.levChi_flag (s = 0, Q = empty); the Lean's Step 1 carries the
% character through levChi_flagL_mul_of_mem_parLev. HasType for flag1L is the book's type 1.
\end{proof}

\begin{lemma}\label{is-onto-step-theta-j-idempotent}
Let $K$ be a field of characteristic $0$ with an algebraic closure $\bar K$, and $G$ a finite group.
Call a function $\chi\colon G\to K$ an \emph{absolutely irreducible character} if it is the character of
an irreducible representation of $G$ on a finite-dimensional $\bar K$-vector space. For such $\chi$ put
\[
 e_\chi=\frac{\chi(1)}{|G|}\sum_{g\in G}\chi(g^{-1})\,g\in K[G],
\]
and for a finite set $\mathcal X$ of them put $e_{\mathcal X}=\sum_{\chi\in\mathcal X}e_\chi$. Let $\chi$ and
$\chi'$ be absolutely irreducible characters. Then:
\begin{enumerate}
\item $\chi(sts^{-1})=\chi(t)$ for all $s,t\in G$, and $e_\chi$ lies in the centre of $K[G]$;
\item $e_\chi e_{\chi'}=e_{\chi'}$ if $\chi=\chi'$, and $e_\chi e_{\chi'}=0$ otherwise;
\item $e_{\mathcal X}$ is an idempotent in the centre of $K[G]$.
\end{enumerate}
\end{lemma}
\leanref{Formalism/LocalToGlobal/Divisibility/ProjLemmas.lean}{Formalism.L2.DIV.idempotent_central}
\leanref{Formalism/LocalToGlobal/Divisibility/GoodPartHermitianSpec.lean}{Formalism.L2.DIV.IsAbsIrrChar}
\leanref{Formalism/LocalToGlobal/Divisibility/ProjLemmas.lean}{Formalism.L2.DIV.IsAbsIrrChar.conj}
\leanref{Formalism/LocalToGlobal/Divisibility/ProjLemmas.lean}{Formalism.L2.DIV.chiIdem_comm_single}
\leanref{Formalism/LocalToGlobal/Divisibility/ProjLemmas.lean}{Formalism.L2.DIV.convA}
\leanref{Formalism/LocalToGlobal/Divisibility/ProjLemmas.lean}{Formalism.L2.DIV.conv_char}
\leanref{Formalism/LocalToGlobal/Divisibility/ProjLemmas.lean}{Formalism.L2.DIV.chiIdem_mul}
\begin{proof}
Let $\rho$ and $\rho'$ be irreducible representations of $G$ over $\bar K$ with characters $\chi$ and
$\chi'$, and $n'=\chi'(1)$ the dimension of $\rho'$.

(1) $\chi(sts^{-1})=\operatorname{tr}\big(\rho(s)\rho(t)\rho(s)^{-1}\big)=\operatorname{tr}\rho(t)=\chi(t)$. Let
$s,u\in G$. The coefficient of $u$ in $s\,e_\chi$ is that of $s^{-1}u$ in $e_\chi$, which is
$\frac{\chi(1)}{|G|}\chi(u^{-1}s)$; the coefficient of $u$ in $e_\chi s$ is that of $us^{-1}$ in $e_\chi$,
which is $\frac{\chi(1)}{|G|}\chi(su^{-1})$. As $u^{-1}s=s^{-1}(su^{-1})s$, the two agree. So $e_\chi$
commutes with every element of $G$, and so with $K[G]$.

(2) \emph{A convolution.} Let $\varphi\colon G\to\bar K$ satisfy $\varphi(gsg^{-1})=\varphi(s)$ for all
$g,s$. The operator $T=\sum_s\varphi(s)\rho'(s)$ commutes with every $\rho'(g)$, as the substitution
$s\mapsto gsg^{-1}$ shows, so by Schur's lemma \cite[\S2.2, Proposition~4]{Ser77} it is a scalar
$c\in\bar K$. The trace of $T\rho'(u)$ gives $\sum_s\varphi(s)\chi'(su)=c\,\chi'(u)$ for every $u$, and
$u=1$ gives $\sum_s\varphi(s)\chi'(s)=c\,n'$. Hence
\[
 n'\sum_s\varphi(s)\chi'(su)=\chi'(u)\sum_s\varphi(s)\chi'(s).
\]

\emph{The orthogonality.} By the orthogonality relations \cite[\S2.3, Theorem~3]{Ser77},
$\sum_g\chi'(g)\chi(g^{-1})$ is $|G|$ if $\rho$ and $\rho'$ are isomorphic and $0$ otherwise.
Isomorphic representations have the same character. Conversely, if $\chi=\chi'$, the relation for
$(\rho,\rho)$ gives $\sum_g\chi'(g)\chi(g^{-1})=\sum_g\chi(g)\chi(g^{-1})=|G|\ne0$, so $\rho$ and $\rho'$
are isomorphic. Hence $\sum_s\chi(s^{-1})\chi'(s)$ is $|G|$ if $\chi=\chi'$ and $0$ otherwise.

\emph{The product.} By (1), $\varphi(s)=\chi(s^{-1})$ satisfies the hypothesis of the convolution, so
\[
 \chi'(1)\sum_s\chi(s^{-1})\chi'(su)=\begin{cases}|G|\,\chi'(u)&\text{if }\chi=\chi',\\ 0&\text{otherwise.}\end{cases}
\]
The coefficient of $u$ in $e_\chi e_{\chi'}$ is
$\sum_s\frac{\chi(1)}{|G|}\chi(s^{-1})\frac{\chi'(1)}{|G|}\chi'(u^{-1}s)$. By (1),
$\chi'(u^{-1}s)=\chi'(su^{-1})$, so this coefficient is
$\frac{\chi(1)}{|G|^2}\,\chi'(1)\sum_s\chi(s^{-1})\chi'(su^{-1})$. By the display at $u^{-1}$, it is $0$ if
$\chi\ne\chi'$, and $\frac{\chi'(1)}{|G|}\chi'(u^{-1})$, the coefficient of $u$ in $e_{\chi'}$, if
$\chi=\chi'$.

(3) By (2), $e_{\mathcal X}e_{\mathcal X}=\sum_{\chi,\chi'\in\mathcal X}e_\chi e_{\chi'}=\sum_{\chi\in\mathcal X}e_\chi
=e_{\mathcal X}$, and $e_{\mathcal X}$ is central by (1).
% Lean: idempotent_central = (idemOf_mul_self, idemOf_central). (1) is IsAbsIrrChar.conj (trace_mul_comm)
% and chiIdem_comm_single; (2) is chiIdem_mul, through conv_char, which applies convA (schur_scalar,
% from Mathlib's IsIrreducible.algebraMap_intertwiningMap_bijective_of_isAlgClosed) with phi(s) =
% chi(s^-1), and Mathlib's char_orthonormal with card_inv_mul_sum_char_mul_char_eq_finrank and
% finrank_intertwiningMap_self for the case chi = chi'; (3) is idemOf_mul_self (Finset.sum_ite_eq) and
% idemOf_central.
\end{proof}

\begin{lemma}\label{is-onto-step-theta-j}
Let $p$ be a point of a family that satisfies the point conditions and the further conditions, $r$
a pair of runs at $p$ pinned to the own places, and $O$ a record object recorded at $(p,r)$, with the
notation of Definitions~\ref{is-record-object} and~\ref{is-onto-record-numbers}. Let $R$ be the
realization of $p$, with the scalars of its characterization, and $F$ its face, with subgroup $P_0$
and character $\psi_0$.
\begin{enumerate}
\item Every $e_j$ is an idempotent with image $V_0(j)$ that commutes with $\rho_j(P)$.
\item Let $v$ be a twisted place and $j$ a lift. Suppose that the base theta lattice
  $\Theta^{\mathrm b}_v(j)$ is the span of $\bigcup_c\Theta_\kappa(h_c)$ for every base theta kernel
  $\kappa$ at $(j,v)$ and every family $(h_c)$ meeting the classes of the level $J_v$
  (Definition~\ref{is-onto-step-kernels}), and that the comparison away from the own place of
  Lemma~\ref{is-onto-step-compare} holds at $v$ and $j$ for the run used at $v$. Then $\Theta_J(v,j)$
  is a finitely generated $\mathcal O_L$-submodule of $V_\psi(j)$, and
  $\Theta_0(v,j)_{(v)}\subset\Theta_J(v,j)_{(v)}$.
\end{enumerate}
\end{lemma}
\leanref{Development/AbelianVarietyTheory/Proofs/DivisorProofs/RankAtPlacesFromComparison.lean}{GlobalTheory.BigOn.tie_thetaJ}
\leanref{Development/AbelianVarietyTheory/Proofs/DivisorProofs/BaseTupleThetaSpanHasTypeStatement.lean}{GlobalTheory.T.baseTupleThetaSpan_hasType}
\leanref{Development/AbelianVarietyTheory/Proofs/DivisorProofs/NotOntoOfCharAlone.lean}{GlobalTheory.RealIndexBound.level_of_tied}
\leanref{Development/AbelianVarietyTheory/Proofs/DivisorProofs/RankAtPlacesFromComparison.lean}{GlobalTheory.BigOn.tie_e_comm_rho}
\begin{proof}
Let $\alpha_j$ and $\mathrm{inc}$ be as in Lemma~\ref{is-ties-level-tie}. By the links
(Lemma~\ref{is-onto-step-links}), $R$ satisfies its characterization at $F$ and condition (11) of
Definition~\ref{is-primitive-rows}, and its family of definite forms satisfies the family conditions.

(1) By Lemma~\ref{is-ties-level-tie}(3), $V_0(j)=e_j(L^{n_Z(j)})$ and
$e_j=\alpha_j^{-1}e_{\mathrm{good}}\alpha_j$. The map $g\mapsto\chi_j(g)^{-1}\varrho_j(g)$ is a
representation of $G$, and $e_{\mathrm{good}}$ is the image of $e_{\mathcal X}$ under its extension to
$L[G]$. By the family conditions the characters in $\mathcal X$ are absolutely irreducible, so
$e_{\mathcal X}$ is an idempotent in the centre of $L[G]$ by Lemma~\ref{is-onto-step-theta-j-idempotent}(3).
Hence $e_{\mathrm{good}}$ is an idempotent that commutes with every $\chi_j(g)^{-1}\varrho_j(g)$, and so
with every $\varrho_j(g)$. So $e_j$ is an idempotent, and for $g\in P$, by Lemma~\ref{is-ties-level-tie}(1),
\[
 \alpha_je_j\rho_j(g)=e_{\mathrm{good}}\varrho_j(\mathrm{inc}(g))\alpha_j
 =\varrho_j(\mathrm{inc}(g))e_{\mathrm{good}}\alpha_j=\alpha_j\rho_j(g)e_j ,
\]
so $e_j\rho_j(g)=\rho_j(g)e_j$.

(2) \emph{Finiteness.} Lemma~\ref{is-onto-step-theta-j-determined} applies to $R$ and $j$, so by the
hypothesis and Lemma~\ref{is-onto-step-theta-j-fg}, $\Theta^{\mathrm b}_v(j)$ is finitely generated. By
Lemma~\ref{is-onto-step-theta-j-ties}(1) and the bijectivity of $\alpha_j$, $\alpha_j$ maps
$\Theta_J(v,j)$ isomorphically onto $\Theta^{\mathrm b}_v(j)$, so $\Theta_J(v,j)$ is finitely generated.

\emph{The type.} The characterization gives hypothesis (2) of Lemma~\ref{is-onto-step-kernel-span}, and
its condition on the group makes $\chi_j$ trivial on $P_0$; the core ties put the values of $\delta_v$ in
$\sigma_D(L)$. So Lemma~\ref{is-onto-step-theta-j-type} applies. Let $x\in\Theta_J(v,j)$ and $u\in U$.
Then $\alpha_j(x)\in\Theta^{\mathrm b}_v(j)$, and $\mathrm{inc}(u)\in P_0$ by
Lemma~\ref{is-ties-level-tie}(2). By (1) and (2) of that lemma,
\[
 \alpha_j(\rho_j(u)x)=\varrho_j(\mathrm{inc}(u))\alpha_j(x)=\psi_0(\mathrm{inc}(u))^{-1}\alpha_j(x)
 =\alpha_j(\psi(u)x),
\]
so $\rho_j(u)x=\psi(u)x$, and $x\in V_\psi(j)$.

\emph{The local containment.} Let $s$ be the run used at $v$ and $q$ its own prime. The residue
characteristic of $v$ is not $q$: if $v$ lies over $p_{de}$, then $q=p'_{de}\ne p_{de}$, and otherwise
$q=p_{de}$. So the comparison of Lemma~\ref{is-onto-step-compare} holds at $\lambda=v$: with $T$ the
trace lattice of $s$, $T_{(v)}\subset\Theta^{\mathrm b}_v(j)_{(v)}$. Let $x\in\Theta_0(v,j)_{(v)}$, so
$ax\in\Theta_0(v,j)$ for some $a\in\mathcal O_L$ outside $v$. By
Lemma~\ref{is-onto-step-theta-j-ties}(2), $a\alpha_j(x)=\alpha_j(ax)\in T$, so $\alpha_j(x)\in
T_{(v)}\subset\Theta^{\mathrm b}_v(j)_{(v)}$, and $b\alpha_j(x)\in\Theta^{\mathrm b}_v(j)$ for some
$b\in\mathcal O_L$ outside $v$. By Lemma~\ref{is-onto-step-theta-j-ties}(1), $bx\in\Theta_J(v,j)$, and
$x\in\Theta_J(v,j)_{(v)}$.
% Lean: from rank_at_twisted_places_of_rec_span (rankAtPlaces_rec). (1) is RealIndexBound.level_of_tied
% (LeafProj: e_eOp, goodOp_idem with idempotent_central (Lemma is-onto-step-theta-j-idempotent) and
% hX = hOk.irr from T.second_link_of_weight_rows, V0_eq) and tie_e_comm_rho (goodOp_rho_comm, goodOp_comm with idempotent_central,
% rho_eq_of_actionPin at a place v_1 from exists_place3). (2) is tie_thetaJ: ThJ_eq with
% toForms_bijective and thetaBase_fg (hd = formsDetermined_of, hM = SpanOverClasses) for finiteness;
% hty = T.baseTupleThetaSpan_hasType (with hC.tup.mu_on, delAt_mem of the core ties, chi_eq_one from
% grp_chiW), rho_eq_of_actionPin and psi_eq for the type; trace_lattice_le_theta_lattice_at_place
% (Th0_eq_comap, the RunPin of the run, CompareOffOwn at lam = v, localization_comap_le) for the
% containment. Size: about 13,600 lines of Lean below these (text count), nearly all below Lemmas
% is-onto-step-theta-j-low-types and is-onto-step-theta-j-type.
\end{proof}

\subsection{Run models and the level hypothesis}\label{subsection-is-run-models}

\begin{definition}\label{is-onto-step-run-model}
Let $p$ be a point of a family, with definite-side datum of rank $r_0$, lifts $1,\dots,m$, host field
$E_h$, group $G$ and one-map data with weight families $\kappa_i$, and let $r$ be a pair of runs at
$p$ pinned to the own places (Definition~\ref{is-onto-step-runs}). Let $i$ be a lift, $C$ a finite set
with an action of $G$, $M$ two-level data of rank $r_0$ over $E_h$ for $G$ and $C$
(Definition~\ref{is-onto-share-two-level}), and $\kappa$ a weight for $M$. Then $(C,M,\kappa)$ is a
\emph{run model} of the lift $i$ if there are a finite place $v$ of $L$ and a bijection $e$ from $C$ to
the set of components of the unitary data of the run $s$ used at $v$ at level $i+1$ and $\sigma_0$,
such that:
\begin{enumerate}
\item for every $c\in C$, the flag data $X_c$ of $M$ are the flag data of the component $e(c)$, and
  $\kappa$ is the weight of $s$ at that component for the weight family $\kappa_i$;
\item there are scalars $a(g,c)\in\mathbf C^\times$ such that for every cusp form $F$ of $s$ at
  level $i+1$ and weight $\kappa_i$, every $g\in G$, $c\in C$ and $x\in\mathcal D_M$, the reading of
  the transform of $F$ by $\ell_i(g^{-1})$ at $e(c)$ and $x$ is $a(g,c)$ times the reading of $F$ at
  $e(gc)$ and $t(g,c)x$;
\item for every $c$, every $\gamma\in\Gamma_c$ and every $w\in L'_c$, $\gamma w-w\in M_bL'_c$, where
  $M_b$ is the level ideal of the core;
\item there are isomorphisms $\phi_c$ from the level module of the core at level $i+1$ onto
  $L'_c/M_bL'_c$ such that for every $c$ and every $u\in G$ with $uc=c$ there is an integer $n$ with
  $t(u,c)w\equiv\phi_c\big(\ell_i(u)^n\phi_c^{-1}(w)\big)\pmod{M_bL'_c}$ for all $w\in L'_c$.
\end{enumerate}
\end{definition}
\leanref{Development/AbelianVarietyTheory/Proofs/DivisorProofs/RunModel.lean}{GlobalTheory.ShareOn.RunModel}

\begin{lemma}\label{is-onto-step-run-model-exists-components-lattice}
Let $E$ be a CM field with complex conjugation $a\mapsto\bar a$ and different $\mathfrak d_E$ over $\mathbf Q$,
let $r\geq0$ be an integer, and let $X$ be data (1) to (6) of flag data of rank $r$ over $E$
(Definition~\ref{is-onto-share-flag-data}), with form $h$, module $\Lambda$, group $\Gamma$ and element $\xi_0$.
Let $\mathfrak A$ be a fractional ideal of $E$, with complex conjugate $\bar{\mathfrak A}$, such that
$\xi_0\mathfrak A\bar{\mathfrak A}\mathfrak d_E=\mathcal O_E$; let $k\le1$ be a real number with
$|\tau(\xi_0)|=k$ for every embedding $\tau\colon E\to\mathbf C$; and let $g>0$ be an integer. Suppose that
$\operatorname{Tr}_{E/\mathbf Q}(-\xi_0h(x,y))\in\mathbf Z$ for $x,y\in\Lambda$, that every homomorphism
$\Lambda\to\mathbf Z$ is $y\mapsto\operatorname{Tr}_{E/\mathbf Q}(-\xi_0h(x,y))$ for some $x\in\Lambda$, and that
$\Gamma$ contains every element of $GL_r(E)$ that preserves $h$ and $\Lambda$ and satisfies $\gamma y-y\in
g\Lambda$ for $y\in\Lambda$. Then $X$ extends to flag data with $g_X=g$ whose self-dual module $L'$ has the
property that $\Lambda$ is the $\mathcal O_E$-span of the $au$ with $a\in\mathfrak A$ and $u\in L'$.
\end{lemma}
\leanref{Development/AbelianVarietyTheory/Proofs/CubicalLineBundles/ExistsPosNatMemFractionalIdeal.lean}{GlobalTheory.exists_dataX_of_perfect_traceForm}
\begin{proof}
For an $\mathcal O_E$-submodule $N$ of $E^r$ let $N^*$ be the set of the $x\in E^r$ with $h(x,y)$ integral for
every $y\in N$. As $h$ is linear in its first variable and conjugate linear in its second,
$(\mathfrak bN)^*=\bar{\mathfrak b}^{-1}N^*$ for every nonzero fractional ideal $\mathfrak b$. The fractional
ideal $\mathfrak A$ is not $0$, as $\xi_0\mathfrak A\bar{\mathfrak A}\mathfrak d_E\ne0$.

\emph{The dual of $\Lambda$.} For $x\in E^r$, the set $h(x,\Lambda)$ is an $\mathcal O_E$-module, so by the
definition of the different, $\operatorname{Tr}_{E/\mathbf Q}(-\xi_0h(x,y))\in\mathbf Z$ for every $y\in\Lambda$ if
and only if $\xi_0h(x,\Lambda)\subset\mathfrak d_E^{-1}$, that is, $x\in(\xi_0\mathfrak d_E)^{-1}\Lambda^*$. The
trace form is nondegenerate on $E^r$, as $\det A\ne0$ and $\xi_0\ne0$, so every homomorphism $\Lambda\to\mathbf Z$
is $y\mapsto\operatorname{Tr}_{E/\mathbf Q}(-\xi_0h(x,y))$ for one $x\in E^r$, as $\Lambda$ spans $E^r$. So the
two hypotheses on the trace form say that $(\xi_0\mathfrak d_E)^{-1}\Lambda^*$ contains $\Lambda$ and is contained
in it: $\Lambda^*=\xi_0\mathfrak d_E\Lambda$.

\emph{The module $L'$.} Put $L'=\mathfrak A^{-1}\Lambda$. Then
\[
 L'^*=\bar{\mathfrak A}\Lambda^*=\xi_0\mathfrak A\bar{\mathfrak A}\mathfrak d_E\cdot\mathfrak A^{-1}\Lambda=L' .
\]
So $h(x,y)$ is integral for $x,y\in L'$, and every $x$ with $h(x,L')$ integral lies in $L'$: $L'$ is self-dual
for $h$. It is finitely generated, as $\Lambda$ is, and $\Lambda=\mathfrak AL'$ is the $\mathcal O_E$-span of the
$au$ with $a\in\mathfrak A$ and $u\in L'$.

\emph{The integers.} The nonzero fractional ideals $\mathfrak A^{-1}$ and $\mathfrak A$ contain positive integers
$N_A$ and $N_C$, and then $N_A\Lambda\subset\mathfrak A^{-1}\Lambda=L'$ and $N_CL'\subset\mathfrak AL'=\Lambda$. Put
$k_0=2k$, so that $0<k_0\le2$ and $|\tau(\xi_0)|=k_0/2$. With $g_X=g$ and the hypothesis on $\Gamma$, the data
(7) and (8) of Definition~\ref{is-onto-share-flag-data} hold.
\end{proof}
% Lean: exists_dataX_of_perfect_traceForm (ExistsPosNatMemFractionalIdeal.lean): hbc (the conjugate ideal is
% conjFrac AA), hd = perfect_iff_hermDual_eq (hermDual Gm Lambda = (xi0 d_E) Lambda), hself through hermDual_smul,
% conjFrac_inv and hone, hermDual_eq_self_iff for L'_int and L'_dual, exists_pos_nat_mem_fractionalIdeal for NA
% and NC, k0 = 2k, and fractionalIdeal_smul_inv_smul for Lambda = AA L'. The Lean's L' is AA^-1 Lambda.

\begin{lemma}\label{is-onto-step-run-model-exists-components-flag}
Let $D$ be a definite-side datum of rank $r_0$, with host field $E_h$ of discriminant $d_{E_h}$, imaginary
quadratic subfield $K=\mathbf Q(\sqrt{-d})$, host element $\xi_0=-1/(\mathrm N\mathfrak c\cdot\sqrt{-d})$,
embedding $\varphi_0$, ideal $\mathfrak c$ of $\mathcal O_K$ and ideal $\mathfrak b$ of $\mathcal O_{E_h}$
(Definitions~\ref{th-definite-datum} and~\ref{hf-host}), and let $\mathfrak A=\mathfrak c\mathcal
O_{E_h}\cdot\mathfrak b^{-1}$ be the \emph{ideal of the host}. Let a core of $D$ satisfy the laws of the core,
with level $M$, let $l$ be a level of a lift and $s$ a good embedding of the coefficient field of the core,
and for each component $c$ of the uniformization of the core at level $l$ at $s$ let $X_c$ be its data (1)
to (6) of Definition~\ref{is-onto-share-flag-data}, with form $h_c$, module $\Lambda_c$ and group $\Gamma_c$.
Suppose that every $\Gamma_c$ contains every element of $GL_{r_0}(E_h)$ that preserves $h_c$ and $\Lambda_c$
and satisfies $\gamma u-u\in M\Lambda_c$ for $u\in\Lambda_c$. Then every $X_c$ extends to flag data with
$k_0=2|\varphi_0(\xi_0)|$, $g_X=M$, $N_A\le|d_{E_h}|$ and $N_C\le\mathrm N\mathfrak c$, whose self-dual module
$L'$ has the property that $\Lambda_c$ is the $\mathcal O_{E_h}$-span of the $au$ with $a\in\mathfrak A$ and
$u\in L'$.
\end{lemma}
\leanref{Development/AbelianVarietyTheory/Proofs/CubicalLineBundles/ExistsRankOneCuspLatticeDataOfPrincipal.lean}{GlobalTheory.LatticeData.exists_dataX_of_hasPrinc}
\leanref{Development/AbelianVarietyTheory/Proofs/CubicalLineBundles/TraceFormIntegralPerfectOfIsUnf.lean}{GlobalTheory.traceForm_integral_perfect_of_isUnf}
\leanref{Formalism/LocalToGlobal/IndefiniteReal/FinitePlaces/FactPrincipalDatum.lean}{Formalism.IndefR.Fin.factPrincipalDatum}
\begin{proof}
Let $c$ be a component, with element $\xi_{0,c}$, and put $\xi_R=|\varphi_0(\xi_0)|=1/(\mathrm N\mathfrak
c\sqrt d)$, so that $0<\xi_R\le1$.

\emph{The element.} By the laws of the core, $|\tau(\xi_{0,c})|=\xi_R$ for every embedding $\tau$ of $E_h$. Both
$\xi_{0,c}$ and $\xi_0$ are purely imaginary, by Definition~\ref{is-onto-share-flag-data}(6) and as $\sqrt{-d}$
is. So $\varphi_0(\xi_{0,c}^2)=-\xi_R^2=\varphi_0(\xi_0^2)$, and $\xi_{0,c}=\pm\xi_0$.

\emph{The ideal.} The different of $E_h$ is $\mathfrak d_{E_h/K}\mathfrak d_K$. Here $\mathfrak d_{E_h/K}=\mathfrak
b\bar{\mathfrak b}$ (Proposition~\ref{hf-host-field-properties}(f)), and $\mathfrak d_K=\sqrt{-d}\,\mathcal O_K$, as
$d\equiv3\pmod4$ (Definition~\ref{hf-framed-host-field}). The conjugate of $\mathfrak A$ is $\bar{\mathfrak
c}\mathcal O_{E_h}\cdot\bar{\mathfrak b}^{-1}$, and $\mathfrak c\bar{\mathfrak c}=\mathrm N\mathfrak c\,\mathcal O_K$.
So $\xi_0\mathfrak A\bar{\mathfrak A}\mathfrak d_{E_h}$ is generated by $\xi_0\mathrm N\mathfrak c\sqrt{-d}=-1$,
and $\xi_{0,c}\mathfrak A\bar{\mathfrak A}\mathfrak d_{E_h}=\mathcal O_{E_h}$.

\emph{The trace form.} By the unfolding law of the core, the uniformization at level $l$ at $s$ satisfies
the compatibility conditions of the moduli problem, as in
Situation~\ref{is-eigenvalue-conjugate-frame-uniformization}. By (i) there, $\varphi_c$ carries the alternating
form of the lattice $L_c$ of a PEL datum, which is integral and perfect, to
$(u,v)\mapsto\operatorname{Tr}_{E_h/\mathbf Q}(-\xi_{0,c}h_c(u,v))$ on $\Lambda_c$. So this form is integral on
$\Lambda_c$, and every homomorphism $\Lambda_c\to\mathbf Z$ is $v\mapsto\operatorname{Tr}_{E_h/\mathbf
Q}(-\xi_{0,c}h_c(u,v))$ for some $u\in\Lambda_c$.

\emph{The extension.} Lemma~\ref{is-onto-step-run-model-exists-components-lattice}, with $k=\xi_R$, $g=M$ and
the hypothesis on $\Gamma_c$, extends $X_c$ to flag data with $g_X=M$ and module $L'$ with $\Lambda_c=\mathfrak
AL'$; its proof gives $k_0=2\xi_R$. Replace $N_A$ and $N_C$ by $\mathrm N\mathfrak b$ and $\mathrm N\mathfrak c$.
These are allowed: $\mathrm N\mathfrak b\in\mathfrak b$, so $\mathrm N\mathfrak b\,\mathfrak A\subset\mathfrak
c\mathcal O_{E_h}\subset\mathcal O_{E_h}$ and $\mathrm N\mathfrak b\,\Lambda_c=\mathrm N\mathfrak b\,\mathfrak
AL'\subset L'$; and $\mathrm N\mathfrak c\in\mathfrak c\mathcal O_{E_h}\subset\mathfrak A$, so $\mathrm
N\mathfrak c\,L'\subset\Lambda_c$. Finally $\mathfrak b$ divides $\mathfrak d_{E_h}$, so $\mathrm N\mathfrak b\le\mathrm
N\mathfrak d_{E_h}=|d_{E_h}|$.
\end{proof}
% Lean: LatticeData.exists_dataX_of_hasPrinc: numerics_of_pins (xiR D <= 1, hx1), the law h.xiAbs (habs),
% traceForm_integral_perfect_of_isUnf (the field EF of the moduli datum through phi_k and its unimod, for the
% unfolding law isUnf_ofPj of h.unf), eq_or_eq_neg_of_purelyImaginary_of_norm_eq with flagDomainConj_hostXi_eq_neg
% (hxi), Fin.factPrincipalDatum (= factPrincipalDatum_of with diffKPrincipal, the different of K generated by
% sqrt(-d), and ccNorm, c cbar = (N c)), exists_dataX_of_perfect_traceForm (Lemma
% is-onto-step-run-model-exists-components-lattice) with AA = ccE * b^-1 and AAbar, then NA = absNorm b
% (absNorm_mul_mem_one, absNorm_halfDifferent_le) and NC = absNorm ccK (absNorm_ccK_mem).

\begin{lemma}\label{is-onto-step-run-model-exists-components-congruence-exponential}
Let a PEL datum over a CM field $E$ with lattice $\mathfrak L$ be given (Situation~\ref{cv-function-situation}),
with a level $(M,\mathfrak M,N)$ (Definition~\ref{cv-model-toroidal-semiabelian}(1)), and let a family
uniformization of its moduli scheme over an embedding, for a type $k\mapsto\sigma_k$ of length $g$, be given
(Theorem~\ref{cv-uniformization-family}), with lattices $L_i$, forms $\langle\ ,\ \rangle_i$, maps
$t_i\colon\mathfrak L\to L_i$ and points $y_i(P)$. For a period $P$ of $L_i$ of this type let
$\operatorname{ex}_{i,P}\colon\mathbf C^g\to A_{y_i(P)}(\mathbf C)$ be the quotient map followed by the isomorphism
$\mathbf C^g/P(L_i)\to A_{y_i(P)}(\mathbf C)$ of the family uniformization. Let $P$ and $P'$ be periods of $L_i$ and
$L_j$ of this type with $y_i(P)=y_j(P')$. Then there are a $\mathbf C$-linear automorphism $h$ of $\mathbf C^g$ and
an $\mathcal O_E$-linear isomorphism $\gamma\colon L_i\to L_j$ such that:
\begin{enumerate}
\item $\operatorname{ex}_{j,P'}(hz)=\operatorname{ex}_{i,P}(z)$ for every $z\in\mathbf C^g$;
\item $P'(\gamma a)=h(P(a))$ and $\langle\gamma a,\gamma a'\rangle_j=\langle a,a'\rangle_i$ for $a,a'\in L_i$;
\item $\gamma(t_ia)-t_ja\in ML_j$ for $a\in\bar{\mathfrak M}\mathfrak L$, and $\gamma(t_ia)-t_ja\in NL_j$ for
  $a\in\mathfrak L$.
\end{enumerate}
\end{lemma}
\leanref{Development/AbelianVarietyTheory/Proofs/UnitaryExistence/FamilyUnifExpUniqueStatement.lean}{GlobalTheory.T.familyUnif_exp_unique_alt}
\leanref{Development/AbelianVarietyTheory/Proofs/UnitaryExistence/FamilyUnifExpUnique.lean}{GlobalTheory.familyUnif_exp_unique}
\begin{proof}
\openstep{The comparison of the two exponentials at the common point, which lifts the isomorphism of complex
tori between them to a linear map $h$ and reads the Weil pairings and the level structures on the lattices,
is not written out.}
% Lean: T.familyUnif_exp_unique_alt = familyUnif_exp_unique: ExpUniqFamily.fibrewise for the exponentials U.ex i P
% and U.ex j P' with the laws ex_add, ex_ker, ex_lin, ex_weil, ex_diff of the family uniformization and
% ExpUniqFamily.isHolTo_ex, and its last clause h4 applied to the level structures ex_eta (at M) and ex_etaA (at
% the full level A = N). Size, following imports and stopping at statements the book cites: about 960 lines
% below familyUnif_exp_unique_alt (830 in 51 theorems), about 920 of them below no other open step of this
% section.
\end{proof}

\begin{lemma}\label{is-onto-step-run-model-exists-components-congruence-rigidity}
Let a PEL datum over a CM field $E$, a level $(M,\mathfrak M,N)$ of it and a family uniformization of its moduli
scheme, for a type of length $g$, be given as in
Lemma~\ref{is-onto-step-run-model-exists-components-congruence-exponential}, with lattices $L_i$, forms
$\langle\ ,\ \rangle_i$, maps $t_i$, points $y_i(P)$ and exponentials $\operatorname{ex}_{i,P}$. Let $P$ and $P'$ be
periods of $L_i$ of the type, $\gamma$ an $\mathcal O_E$-linear automorphism of $L_i$ that preserves
$\langle\ ,\ \rangle_i$, with $\gamma(a)-a\in ML_i$ for $a\in\bar{\mathfrak M}L_i$ and $\gamma(a)-a\in NL_i$ for
$a\in L_i$, and $h_0\in M_g(\mathbf C)$ with $P'(\gamma a)=h_0P(a)$ for every $a\in L_i$. Then $h_0$ is
invertible, $y_i(P)=y_i(P')$, and $\operatorname{ex}_{i,P'}(h_0z)=\operatorname{ex}_{i,P}(z)$ for every
$z\in\mathbf C^g$.
\end{lemma}
\leanref{Development/AbelianVarietyTheory/Proofs/UnitaryExistence/UniformizationPeriodLatticeChange.lean}{GlobalTheory.FamRigid.ex_eq_of_level_group}
\leanref{Development/AbelianVarietyTheory/Proofs/UnitaryExistence/UniformizationPeriodLatticeChange.lean}{GlobalTheory.FamRigid.finite_order}
\leanref{Development/AbelianVarietyTheory/Proofs/ThetaShiftProofs/ExistsInvariantNatLengthLattice.lean}{GlobalTheory.latMinkowski}
\begin{proof}
\emph{The point.} The image of the $\mathbf C$-linear map $h_0$ contains $P'(\gamma a)$ for every $a$, so it
contains $P'(L_i)$, which spans $\mathbf C^g$ over $\mathbf R$ (Definition~\ref{cv-cm-point-periods}(3)). So $h_0$
is surjective, hence invertible. Then $P'=h_0\circ P\circ\gamma^{-1}$, and $\gamma^{-1}$ preserves the form and
satisfies the same congruences as $\gamma$; by the description of the fibres of the family uniformization
(Theorem~\ref{cv-uniformization-family}(2), with the congruence modulo $N$ for the full level),
$y_i(P)=y_i(P')$.

\emph{A second isomorphism.} By Lemma~\ref{is-onto-step-run-model-exists-components-congruence-exponential} there
are a $\mathbf C$-linear automorphism $h_1$ and an automorphism $\gamma_1$ of $L_i$ that preserves the form, with
$\operatorname{ex}_{i,P'}(h_1z)=\operatorname{ex}_{i,P}(z)$, $P'(\gamma_1a)=h_1P(a)$ and
$\gamma_1(t_ia)-t_ia\in ML_i$ for $a\in\bar{\mathfrak M}\mathfrak L$. The map $t_i$ induces an isomorphism
$\mathfrak L/M'\mathfrak L\to L_i/M'L_i$ for an integer $M'$ divisible by $M$
(Theorem~\ref{cv-uniformization-family}). So every $b\in L_i$ is $t_i(c')+M'e$ with $c'\in\mathfrak L$ and $e\in L_i$,
and for $v\in\bar{\mathfrak M}$, $\gamma_1(vb)-vb=\big(\gamma_1(t_i(vc'))-t_i(vc')\big)+M'(\gamma_1(ve)-ve)\in ML_i$. So
$\gamma_1(a)-a\in ML_i$ for $a\in\bar{\mathfrak M}L_i$, and by Lemma~\ref{cv-volume-key-models-half}, applied to
$\gamma$ and to $\gamma_1$, $\gamma(a)-a$ and $\gamma_1(a)-a$ lie in $ML_i$ for every $a\in L_i$.

\emph{Finite order.} Put $\delta=\gamma\circ\gamma_1^{-1}$, an automorphism of $L_i$ that preserves the form, and
$k=h_0h_1^{-1}$. Then $P'(\delta a)=k(P'(a))$ for every $a$. Let $H$ be a positive definite Hermitian form on
$\mathbf C^g$ with $\operatorname{Im}H(P'(a),P'(b))=\langle a,b\rangle_i$ (Definition~\ref{cv-cm-point-periods}(3)).
As $P'(L_i)$ spans $\mathbf C^g$ over $\mathbf R$, the $\mathbf C$-linear map $k$ preserves $\operatorname{Im}H$,
hence also $\operatorname{Re}H(z,z)=\operatorname{Im}H(iz,z)$. So for $a\in L_i$ all the points $P'(\delta^na)$,
$n\ge0$, lie in the bounded set $\{z:\operatorname{Re}H(z,z)\le\operatorname{Re}H(P'(a),P'(a))\}$ and in the lattice
$P'(L_i)$, which has finitely many points there. As $P'$ is injective, $\delta^na=a$ for some $n\ge1$, and as
$L_i$ is finitely generated, $\delta$ has finite order.

\emph{Minkowski's lemma.} For $a\in L_i$, $\delta(\gamma_1a)-\gamma_1a=(\gamma a-a)-(\gamma_1a-a)\in ML_i$, so
$\delta(b)-b\in ML_i$ for every $b$, and $M\ge3$. Then $\delta=1$. Indeed, $L_i$ is a free $\mathbf Z$-module of
finite rank; let $|y|$ be the sum of the absolute values of the coordinates of $y$ in a basis, $n\ge1$ with
$\delta^n=1$, and $p(y)=\sum_{j<n}|\delta^jy|$. Then $p(y-y')\le p(y)+p(y')$, $p(my)=mp(y)$ for integers $m\ge0$,
$p(y)=0$ only for $y=0$, and $p(\delta y)=p(y)$. We show $\delta y=y$ by induction on $p(y)$. If $p(y)=0$, then
$y=0$. Otherwise $\delta y=y+Mz$ with $z\in L_i$, and $Mp(z)=p(\delta y-y)\le2p(y)$, so $p(z)<p(y)$ as $M\ge3$.
By induction $\delta z=z$, so $\delta^jy=y+jMz$ for every $j$, and $j=n$ gives $nMz=0$, so $z=0$ and $\delta y=y$.

\emph{The exponentials.} So $\gamma=\gamma_1$, and $h_0P(a)=P'(\gamma a)=h_1P(a)$ for every $a$. As $P(L_i)$ spans
$\mathbf C^g$ over $\mathbf R$, $h_0=h_1$, and $\operatorname{ex}_{i,P'}(h_0z)=\operatorname{ex}_{i,P}(z)$.
\end{proof}
% Lean: FamRigid.ex_eq_of_level_group (with the hypothesis latMinkowski), followed: hsurj and hu for h0; hpt from
% the field fibre of FamilyUnif; T.familyUnif_exp_unique_alt (Lemma
% is-onto-step-run-model-exists-components-congruence-exponential) for h1, gam1, with hM1' through F.t_surj and
% l.M | l.M'; f0, f1 from LevelHalf.trivial_mod_of_trivial_on_half (Lemma cv-volume-key-models-half); finite_order
% (hB, hQ: the real part of Hm through re_eq_imForm, and Metric.finite_isBounded_inter_isClosed on the discrete
% range); latMinkowski = end_eq_one_of_pow_eq_one_of_congr (exists_invariant_nat_length_lattice, the strong
% induction key); hg and hlin.

\begin{lemma}\label{is-onto-step-run-model-exists-components-congruence-orbit}
In Situation~\ref{is-eigenvalue-conjugate-frame-uniformization}, suppose that the integer $N$ of the level
$(M,\mathfrak M,N)$ is $1$. Let $c$ be a component, with form $h_c$, and let $\gamma\in GL_r(E)$ preserve $h_c$ and
$\Lambda_c$ and satisfy $\gamma u-u\in M\Lambda_c$ for $u\in\Lambda_c$. Then for every $x\in\mathcal D_c$ there are
$\gamma'\in\Gamma_c$ and $b\in\mathcal B^{\Psi_c}$ with $\gamma x=\gamma'xb$.
\end{lemma}
\leanref{Development/AbelianVarietyTheory/Proofs/CubicalLineBundles/CongruenceActionMemFrameOrbit.lean}{GlobalTheory.congr_lact_mem_orbit}
\begin{proof}
As $\gamma$ preserves $h_c$, the Gram matrices of the columns of $\gamma x$ and of $x$ are equal, so $\gamma x\in
\mathcal D_c$. By (ii) of the situation there are periods $P$ and $P'$ of $L_c$ at $x$ and at $\gamma x$: the $j$-th
coordinate of $P(a)$ is the period of $\varphi_c(a)$ at $x$ at the letter $\ell_j$, and that of $P'(a)$ the period
at $\gamma x$; the points of $y(c,x)$ and $y(c,\gamma x)$ are $y_c(P)$ and $y_c(P')$; and the frame vectors of
$y(c,x)$ and $y(c,\gamma x)$ are the derivatives at $0$ of $\operatorname{ex}_{c,P}$ and of
$\operatorname{ex}_{c,P'}$ in the directions of the standard basis.

Let $\tilde\gamma=\varphi_c^{-1}\circ\gamma\circ\varphi_c$, an $\mathcal O_E$-linear automorphism of $L_c$. By (i),
it preserves $\langle\ ,\ \rangle_c$, as $\gamma$ preserves $h_c$; and $\tilde\gamma(a)-a\in ML_c$ for every $a$, as
$\gamma u-u\in M\Lambda_c$; the congruence modulo $N=1$ holds trivially. As $\gamma$ preserves $h_c$, the periods of
$\gamma v$ at $\gamma x$ are those of $v$ at $x$ (Definition~\ref{is-eigenvalue-conjugate-frame-fibres}(1)):
$h_\tau(\gamma^\tau x_\tau e_a,\tau(\gamma v))=h_\tau(x_\tau e_a,\tau(v))$ and $(\gamma^\tau x_\tau)^{-1}\tau(\gamma
v)=x_\tau^{-1}\tau(v)$. So $P'(\tilde\gamma a)=P(a)$ for every $a\in L_c$. By
Lemma~\ref{is-onto-step-run-model-exists-components-congruence-rigidity} with $h_0=1$, $y(c,x)$ and $y(c,\gamma x)$
have the same point and $\operatorname{ex}_{c,P'}=\operatorname{ex}_{c,P}$; so they have the same frame vectors.
Two framed points with the same point have sections that generate the same line, so they differ by a nonzero
scalar; with the same frame vectors, the two framed points are triangularly equivalent, with the identity
matrices as the triangular change. By the compatibility conditions of the uniformization, $y(c,x)$ and
$y(c,\gamma x)$ are triangularly equivalent only if $\gamma x=\gamma'xb$ for some $\gamma'\in\Gamma_c$ and
$b\in\mathcal B^{\Psi_c}$.
\end{proof}
% Lean: congr_lact_mem_orbit, followed: dom_lact, the field pt of the witness at x and at gamma x (hpt k), gam =
% phi^-1 o latAut gamma o phi, FamRigid.ex_eq_of_level_group with h0 = 1 (Lemma
% is-onto-step-run-model-exists-components-congruence-rigidity; the form through hEF and the first clause of
% IsUStab, the congruence from hM, the periods from perFr_lact), then the triangular law of IsUnf (hfib) at
% CongrOrbit.frBor_of_eq (Model.exists_frChg for the scalar of the sections) with VerticalDeriv.isDerivAt0_unique
% for the frame vectors.

\begin{lemma}\label{is-onto-step-run-model-exists-components-congruence-roots}
In Situation~\ref{is-eigenvalue-conjugate-frame-uniformization}, suppose that the integer $N$ of the level
$(M,\mathfrak M,N)$ is $1$. Let $c$ be a component, with form $h_c$, and let $\gamma\in GL_r(E)$ preserve $h_c$ and
$\Lambda_c$ and satisfy $\gamma u-u\in M\Lambda_c$ for $u\in\Lambda_c$. Then there are $\gamma'\in\Gamma_c$, $u\in
E$ and an integer $n\ge1$ with $u^n=1$ and $\gamma'=u\gamma$.
\end{lemma}
\leanref{Development/AbelianVarietyTheory/Proofs/CubicalLineBundles/UnfGroupContainsCongruenceUptoRootsOfUnity.lean}{GlobalTheory.unf_group_contains_congruence_upto_rootsOfUnity}
\begin{proof}
By Lemma~\ref{is-onto-step-run-model-exists-components-congruence-orbit}, $\gamma x\in\Gamma_cx\mathcal B^{\Psi_c}$
for every $x\in\mathcal D_c$. By Lemma~\ref{cv-uniformization-scalar-orbits}, applied to the flag data $X_c$,
whose flag domain $\mathcal D_c$ is the set of frames, there are $\gamma'\in\Gamma_c$ and $u\in E$ with
$\gamma'=u\gamma$; and $u\ne0$, as $\gamma'$ is invertible. It remains to show that $u$ is a root of unity.

By Definition~\ref{is-onto-share-flag-data}(4), $\Gamma_c$ has finite index in the group of the elements of
$GL_r(E)$ that preserve $h_c$ and $\Lambda_c$: there is a finite set $S$ of such elements such that every such
element $g$ has $s^{-1}g\in\Gamma_c$ for some $s\in S$. For every $i\ge0$ the element $\gamma^i$ preserves $h_c$ and
$\Lambda_c$; choose $s_i\in S$ with $s_i^{-1}\gamma^i\in\Gamma_c$. As the scalar $u$ commutes with $\gamma$,
$\gamma'^i(s_i^{-1}\gamma^i)^{-1}=u^is_i$ lies in $\Gamma_c$. As $S$ is finite, there are $a<b$ with $s_a=s_b$, and
then $u^{b-a}=(u^bs_b)(u^as_a)^{-1}\in\Gamma_c$. The flag domain $\mathcal D_c$ is not empty: by
Definition~\ref{is-onto-share-flag-data}(2), a basis with Gram matrix $J$ at every $\tau\in\Psi_c$ is a point of it.
For $x\in\mathcal D_c$, $u^{b-a}x=x\beta$, where $\beta_\tau$ is the scalar matrix $\tau(u^{b-a})$, which lies in
$\mathcal B$. By Definition~\ref{is-onto-share-flag-data}(4), $u^{b-a}=1$, and $n=b-a$ works.
\end{proof}
% Lean: unf_group_contains_congruence_upto_rootsOfUnity, followed: congr_lact_mem_orbit (Lemma
% is-onto-step-run-model-exists-components-congruence-orbit), OrbitScalar.exists_smul_of_orbits (Lemma
% cv-uniformization-scalar-orbits, stated in the Lean for any flag data), hu0, Data.Gamma_index with isUStab_pow',
% hmem, Finite.exists_ne_map_eq_of_infinite for a < b, OrbitScalar.exists_mem_dom and Data.Gamma_free.

\begin{lemma}\label{is-onto-step-run-model-exists-components-congruence-lattice}
Let $F\subset E$ be number fields, $n\ge0$, $L\subset E^n$ a finitely generated $\mathcal O_E$-submodule that spans
$E^n$, with closure $\hat L$ in $(\mathbf A_{F,f}\otimes_FE)^n$, and $g\in GL_n(\mathbf A_{F,f}\otimes_FE)$. Then
$\Lambda=E^n\cap g\hat L$ is a finitely generated $\mathcal O_E$-submodule that spans $E^n$, and its closure is
$g\hat L$. Moreover, let $M\ge1$ be an integer, $\kappa\in M_n(\mathbf A_{F,f}\otimes_FE)$ with $\kappa z-z\in M\hat L$
for $z\in\hat L$, and $\delta\in M_n(E)$ with $\delta g=g\kappa$. Then $\delta v-v\in M\Lambda$ for every
$v\in\Lambda$.
\end{lemma}
\leanref{Development/AbelianVarietyTheory/Proofs/UnitaryExistence/ExistsLatticeRationalVectorsInAdelicTranslateStatement.lean}{GlobalTheory.T.exists_lattice_rationalVectorsInAdelicTranslate}
\leanref{Development/AbelianVarietyTheory/Proofs/UnitaryExistence/CoreArithmeticGroupContainsPrincipalCongruence.lean}{GlobalTheory.PrincCore.exists_lattice_congr}
\begin{proof}
\openstep{The first assertion, that the rational vectors of an adelic translate of a lattice form a lattice
with that closure, is not written out.}
% Lean: T.exists_lattice_rationalVectorsInAdelicTranslate = AdLat.exists_lattice_latOf. Size, following imports
% and stopping at statements the book cites: about 750 lines (710 in 47 theorems), none of them below the other
% open step of this section.

For the second, let $v\in\Lambda$, so $v=gz$ with $z\in\hat L$, and $\kappa z-z=My$ with $y\in\hat L$. Then
$\delta v-v=g(\kappa z-z)=M\,gy$. As $M$ is invertible in $E$, $w=M^{-1}(\delta v-v)$ lies in $E^n$, and $w=gy\in
g\hat L$; so $w\in\Lambda$ and $\delta v-v=Mw$.
\end{proof}
% Lean: PrincCore.exists_lattice_congr, followed: exists_lattice_rationalVectorsInAdelicTranslate for Lambda (h1,
% h2, h3), then for x in Lambda the element (M:E)^-1 (delta x - x), with e1, e2, e3 for the computation in the
% adeles.

\begin{lemma}\label{is-onto-step-run-model-exists-components-congruence-root}
Let $D$ be a definite-side datum with host field $E/F$, let $V$ be a Hermitian space of $D$ of rank $r\ge1$
(Definition~\ref{th-spaces}) with Gram matrix $G_V$, lattice $L_0\subset E^r$ in the coordinates of its basis and
closure $\hat L_0$ of $L_0$, and let a dictionary for $V$ be given (Definition~\ref{is-eigenvalue-conjugate-dictionary}(3)),
with components $c$, flag data $X_c$ with sets $\Psi_c$, forms $h_c$, modules $\Lambda_c$ and groups $\Gamma_c$, map
$g\mapsto(c(g),x(g))$ and multiplier $\mu$. Let $M\ge3$, and let $K$ be the group of the $k\in U(V)(\mathbf A)$ with
archimedean component $1$ whose finite part satisfies $kz-z\in M\hat L_0$ for $z\in\hat L_0$. Suppose that the
dictionary satisfies the uniformization condition of Definition~\ref{is-eigenvalue-conjugate-dictionary}(4)(c) for
$K$, with homomorphisms $\theta_{c,\tau}$, and that $\mu$ is multiplicative with $|\mu|=1$, trivial on
$U(V)(F)$, and $|d_i(x(g),\tau)|=1$ for all $g$, $\tau$ and $i$; only the last condition on $\mu$ is used. Let $c$ be a
component, $\gamma\in GL_r(E)$ with $\gamma v-v\in M\Lambda_c$ for $v\in\Lambda_c$, $\gamma'\in\Gamma_c$, and $u\in
E$ with $u^n=1$ for some $n\ge1$ and $\gamma'=u\gamma$. Then $u=1$.
\end{lemma}
\leanref{Development/AbelianVarietyTheory/Proofs/UnitaryExistence/CoreArithmeticGroupContainsPrincipalCongruence.lean}{GlobalTheory.PrincCore.root_eq_one}
\leanref{Development/AbelianVarietyTheory/Proofs/UnitaryExistence/CoreArithmeticGroupContainsPrincipalCongruence.lean}{GlobalTheory.PrincCore.unitary_of_first_minor}
\leanref{Development/AbelianVarietyTheory/Proofs/UnitaryExistence/AdelicUnitaryExistsWithPrescribedPartsStatement.lean}{GlobalTheory.T.adelicUnitary_exists_with_prescribed_parts}
\begin{proof}
We use three facts on the adelic unitary group. Every $\gamma_0\in U(V)(F)$ has a matrix $\delta\in GL_r(E)$, which is
its finite part and satisfies $\theta_{c,\tau}(\gamma_0)=\tau(\delta)$ for every $c$ and $\tau\in\Psi_c$. As
$\Psi_c$ contains exactly one of $\tau$ and $\bar\tau$ for every $\tau$, the $\theta_{c,\tau}$, $\tau\in\Psi_c$,
identify $(\mathbf A_{F,\infty}\otimes_FE)$ with $\prod_{\tau\in\Psi_c}\mathbf C$; so for every family $(w_\tau)$ of
matrices with $\bar w_\tau^{\mathsf T}\tau(G_V)w_\tau=\tau(G_V)$ there is $h\in U(V)(\mathbf A)$ with finite component
$1$ and $\theta_{c,\tau}(h)=w_\tau$ for every $\tau\in\Psi_c$. And $\theta_{c,\tau}$ is multiplicative.

\emph{The base point.} By the uniformization condition there are $x_c\in\mathcal D_c$ and $g$ with $(c(g),x(g))=(c,x_c)$,
and for every $h$ with finite component $1$ there is $\gamma''\in\Gamma_c$ with
$(c(gh),x(gh))=(c,\gamma''x')$, where $x'_\tau=\theta_{c,\tau}(h)x_{c,\tau}$.

\emph{Unitarity.} Fix $\tau\in\Psi_c$. Let $A_\tau(z,w)=z^{\mathsf T}\tau(A_c)\bar w$ be the form $h_{c,\tau}$ on
$\mathbf C^r$, where $A_c$ is the matrix of $X_c$, and $B_\tau(z,w)=z^{\mathsf T}\tau(G_V)^{\mathsf T}\bar w=\bar
w^{\mathsf T}\tau(G_V)z$, the form of $V$ at $\tau$ with its arguments exchanged; both are Hermitian, and $B_\tau$ is
nondegenerate. A matrix $w$ preserves $B_\tau$ if and only if $\bar w^{\mathsf T}\tau(G_V)w=\tau(G_V)$. We show that
$\tau(\gamma')$ preserves $B_\tau$. Let $v$ be the first column of $x_{c,\tau}$, let $w$ preserve $B_\tau$, and
choose $h$ with finite component $1$, $\theta_{c,\tau}(h)=w$ and $\theta_{c,\tau'}(h)=1$ for $\tau'\ne\tau$. Then
$x(gh)=\gamma''x'$ with $x'_\tau=wx_{c,\tau}$, and $\gamma''$ preserves $h_c$, so
$|A_\tau(wv,wv)|=|d_1(x(gh),\tau)|=1$. By Lemma~\ref{cv-uniformization-unitary-orbit}, $A_\tau=tB_\tau$ for a real
$t\ne0$. As $\gamma'\in\Gamma_c$ preserves $h_c$, $\tau(\gamma')$ preserves $A_\tau$, hence $B_\tau$.

\emph{The rational element.} By the second fact there is $h$ with finite component $1$ and $\theta_{c,\tau}(h)=\tau(\gamma')$
for every $\tau\in\Psi_c$. The base point gives $\gamma''\in\Gamma_c$ with $(c(gh),x(gh))=(c,\gamma''\gamma'x_c)$, and
$\gamma''\gamma'\in\Gamma_c$. By the uniformization condition, $gh=\gamma_0gk$ with $\gamma_0\in U(V)(F)$ and
$k\in K$. Let $\delta$ be the matrix of $\gamma_0$. The finite parts give $\delta g_f=g_fk_f^{-1}$, as $h_f=1$, and
$k_f^{-1}z-z\in M\hat L_0$ for $z\in\hat L_0$. The archimedean parts give $\theta_{c,\tau}(\gamma_0)\theta_{c,\tau}(g)=
\theta_{c,\tau}(g)\theta_{c,\tau}(h)$, as $k$ has archimedean component $1$.

\emph{The conclusion.} By Lemma~\ref{is-onto-step-run-model-exists-components-congruence-lattice}, applied to $L_0$,
$g_f$, $\kappa=k_f^{-1}$ and $\delta$, the module $\Lambda'=E^r\cap g_f\hat L_0$ is finitely generated, spans $E^r$, and
$\delta v-v\in M\Lambda'$ for $v\in\Lambda'$. Choose $\tau\in\Psi_c$ and put $Q=\theta_{c,\tau}(g)^{-1}$. Then
\[
 \tau(u\gamma)\,Q=\tau(\gamma')\,Q=\theta_{c,\tau}(h)\theta_{c,\tau}(g)^{-1}=\theta_{c,\tau}(g)^{-1}\theta_{c,\tau}(\gamma_0)
 =Q\,\tau(\delta).
\]
By Lemma~\ref{cv-uniformization-root-congruence}, for $\Lambda_c$, $\Lambda'$, $\gamma$, $\delta$ and $M\ge3$, $u=1$.
\end{proof}
% Lean: PrincCore.root_eq_one, followed: the fields x0, archMat and fibre of DictIsUnif (hx0, key, hfib);
% T.adelicUnitary_exists_with_prescribed_parts (= UnitaryGlue.adelic_glue: hrat for delta, hsur for h, the
% multiplicativity hm of archM); hun from unitary_of_first_minor (UnitaryOrbitForm.eq_smul_of_norm_eq_one_on_orbit,
% Lemma cv-uniformization-unitary-orbit, for G.map tau and J = V.gram.map tau transposed, through the last field of
% MuUnit, dmin_lact and hsur) and PairUnfoldPf.gramE; IsLevelK for kappa^-1; hfin and hinf; exists_lattice_congr
% (Lemma is-onto-step-run-model-exists-components-congruence-lattice) with latHat V = hat (L0); a tau in Psi from
% cmType; root_of_unity_eq_one_of_conj_congruence (Lemma cv-uniformization-root-congruence) with the archimedean
% identity e.

\begin{lemma}\label{is-onto-step-run-model-exists-components-congruence}
Let a core of a definite-side datum $D$ of rank $r_0$ with host field $E_h$ satisfy the laws of the core, with
level $M$, let $l$ be a level of a lift and $s$ a good embedding of the coefficient field of the core, and for
each component $c$ of the uniformization of the core at level $l$ at $s$ let $X_c$ be its data (1) to (6) of
Definition~\ref{is-onto-share-flag-data}, with form $h_c$, module $\Lambda_c$ and group $\Gamma_c$. Then every
element $\gamma$ of $GL_{r_0}(E_h)$ that preserves $h_c$ and $\Lambda_c$ and satisfies $\gamma u-u\in M\Lambda_c$
for $u\in\Lambda_c$ lies in $\Gamma_c$.
\end{lemma}
\leanref{Development/AbelianVarietyTheory/Proofs/UnitaryExistence/CoreArithmeticGroupContainsPrincipalCongruence.lean}{GlobalTheory.PrincCore.hasPrinc_core}
\begin{proof}
By the pins of the core, $M\ge3$, the integer $N$ of the level of the core is $1$, and $r_0\ge1$. By the
unfolding law, the uniformization of the core at level $l$ at $s$ satisfies the compatibility conditions of the
moduli problem, as in Situation~\ref{is-eigenvalue-conjugate-frame-uniformization} with $r=r_0$ and the level
$(M,\mathfrak M,1)$ of the core. Let $c$ be a component and $\gamma$ as in the statement.

\emph{A root of unity.} By Lemma~\ref{is-onto-step-run-model-exists-components-congruence-roots} there are
$\gamma'\in\Gamma_c$, $u\in E_h$ and $n\ge1$ with $u^n=1$ and $\gamma'=u\gamma$.

\emph{The root is $1$.} By the dictionary law of the core at level $l$ and $s$, the dictionary of the core there
is a dictionary for the space $V$ of the core there, of rank $r_0$, with the components $c$ and the flag data
$X_c$ of the uniformization, that satisfies the uniformization condition of
Definition~\ref{is-eigenvalue-conjugate-dictionary}(4)(c) for the level group $K_f$ of the core there, with
homomorphisms $\theta_{c,\tau}$ that extend $\tau$. By the level law and the pins, $K_f$ is the group of the
elements with archimedean component $1$ whose finite part $k$ satisfies $kz-z\in M\hat L_0$ for $z$ in the
closure $\hat L_0$ of the lattice of $V$. By the multiplier law, the multiplier $\mu$ is multiplicative with
$|\mu|=1$, trivial on the rational points, and $|d_i(x(g),\tau)|=1$ for all $g$, $\tau$ and $i$. So
Lemma~\ref{is-onto-step-run-model-exists-components-congruence-root} gives $u=1$, and $\gamma=\gamma'\in\Gamma_c$.
\end{proof}
% Lean: PrincCore.hasPrinc_core, followed: numerics_of_pins (M >= 3 from GoodAt 2, the full level l.A = 1,
% 1 <= r), T.unf_group_contains_congruence_upto_rootsOfUnity (Lemma
% is-onto-step-run-model-exists-components-congruence-roots) for the unfolding law isUnf_ofPj of h.unf, the
% dictionary law h.dictU (th, hrk, hth, hD = DictIsUnif), h.levelK with the pin h.pins.2.1.1 for the level,
% h.muUnit, and PrincCore.root_eq_one (Lemma is-onto-step-run-model-exists-components-congruence-root). The same
% argument appears as Steps A to D of the proof of Lemma cv-cover-uniformization. Two narrower open steps remain:
% Lemma is-onto-step-run-model-exists-components-congruence-exponential (about 960 lines) and the first assertion
% of Lemma is-onto-step-run-model-exists-components-congruence-lattice (about 750 lines).

\begin{lemma}\label{is-onto-step-run-model-exists-components-back-periods}
Let $E$ be a CM field, and let $X$ and $X'$ be flag data of rank $r$ over $E$
(Definition~\ref{is-onto-share-flag-data}) with the same set $\Psi$, integer $p_0$ and element $\xi_0$, forms $h$ and
$h'$, modules $\Lambda$ and $\Lambda'$ and flag domains $\mathcal D$ and $\mathcal D'$. Let $a\in GL_r(E)$ with
$h'(au,av)=h(u,v)$ for $u,v\in E^r$, and $v\in\Lambda$ if and only if $av\in\Lambda'$. For an embedding $\tau$ of $E$
put $m_\tau=r-p_0$ if $\tau\in\Psi$ and $m_\tau=p_0$ otherwise; for $1\le i\le m_\tau$, a family $x$ of invertible
matrices and $v\in E^r$, the \emph{period} of $v$ at $x$ at $(\tau,i)$ for $X$ is $\pi_{\tau,i}(x,v)=
(x_\tau^{-1}\tau(v))_{p_0+i}$ if $\tau\in\Psi$ and $\pi_{\tau,i}(x,v)=h_{\bar\tau}(x_{\bar\tau}e_{p_0+1-i},\bar\tau(v))$
if $\bar\tau\in\Psi$, and $\pi'_{\tau,i}$ is the same for $X'$. Let $x\in\mathcal D$, $x'\in\mathcal D'$ and complex
numbers $A_{\tau,ij}$ with $A_{\tau,ij}=0$ for $j<i$, such that
\[
 \pi'_{\tau,i}(x',av)=\sum_{j=1}^{m_\tau}A_{\tau,ij}\,\pi_{\tau,j}(x,v)
\]
for every $\tau$, every $1\le i\le m_\tau$ and every $v\in\Lambda$. Then there is $b\in\mathcal B^\Psi$ with
$x'=(ax)b$. If moreover $|d_i(x,\tau)|=1$ and $|d_i(x',\tau)|=1$ for all $\tau\in\Psi$ and $i$, then
$|b_{\tau,ii}|=1$ for all $\tau$ and $i$, and $|A_{\tau,ii}|=1$ for all $\tau$ and $1\le i\le m_\tau$.
\end{lemma}
\leanref{Development/AbelianVarietyTheory/Proofs/JetHarmonics/FramesDifferByIsometryAndTriangularOfPeriodsStatement.lean}{GlobalTheory.T.frames_differ_by_isometry_and_triangular_of_periods}
\begin{proof}
\openstep{The comparison, at each embedding in $\Psi$, of the two frames through the periods, which shows
that $(ax)_\tau^{-1}x'_\tau$ is upper triangular, is not written out.}
% Lean: T.frames_differ_by_isometry_and_triangular_of_periods = PeriodBorel.borel_of_periods (53 lines): hz
% (IsMatch.mem_dom), hlat, PeriodBorel.loc_borel at each tau in Psi for the periods at tau (perFr_mem) and at
% the conjugate embedding (perFr_conj), b = (a x)^-1 x', and FrameDegreeData.ratio_rmul for the unit frames.
% Size, following imports and stopping at statements the book cites: about 530 lines (460 in 28 theorems),
% none of them below the other open steps of this section.
\end{proof}

\begin{lemma}\label{is-onto-step-run-model-exists-components-back}
In Situation~\ref{is-eigenvalue-conjugate-frame-uniformization}, suppose that the integer $N$ of the level
$(M,\mathfrak M,N)$ is $1$. Let $c$ be a component, with form $h_c$, let $\gamma\in\Gamma_c$ and
$x\in\mathcal D_c$. Then there are an element $\eta\in GL_r(E)$ that preserves $h_c$ and $\Lambda_c$ and
satisfies $\eta w-w\in M\Lambda_c$ for $w\in\Lambda_c$, and $b\in\mathcal B^{\Psi_c}$, such that
$x=(\eta\gamma x)b$.
\end{lemma}
\leanref{Development/AbelianVarietyTheory/Proofs/DivisorProofs/ExistsPrincipalCongruenceBackTriangularFrameStatement.lean}{GlobalTheory.T.exists_principal_congruence_back_triangular_frame}
\begin{proof}
Write $(p,q)$ for the signature of the PEL datum and $m_{\tau'}$ for its multiplicities. By the compatibility
conditions, $r=p+q$, the integer $p_0$ of $X_c$ is $q$, and $\Psi_c$ is the set of the embeddings that send
$\sqrt{-d}$ to the square root of $-d$ fixed by the datum. The multiplicity $m_{\tau'}$ is $p$ if $\tau'$ sends
$\sqrt{-d}$ to that root and $q$ otherwise, and as $s$ is good, $\tau'\mapsto s\circ\varepsilon_{\tau'}$ is a
bijection onto the embeddings of $E$ that preserves this property. So for a letter $(\tau',i)$ with
$\tau=s\circ\varepsilon_{\tau'}$, $m_{\tau'}=p=r-p_0$ if $\tau\in\Psi_c$ and $m_{\tau'}=q=p_0$ otherwise, and
$(\tau',i)\mapsto(\tau,i)$ is a bijection from the letters onto the pairs of Lemma~\ref{is-onto-step-run-model-exists-components-back-periods}
for $X_c$, under which the period of $v$ at $x$ at the letter $(\tau',i)$
(Definition~\ref{is-eigenvalue-conjugate-frame-fibres}(5)) is $\pi_{\tau,i}(x,v)$.

\emph{The two framed points.} As $\gamma$ preserves $h_c$, $\gamma x\in\mathcal D_c$. As $\gamma x=\gamma x1$ with
$\gamma\in\Gamma_c$, the compatibility conditions make $y(c,x)$ and $y(c,\gamma x)$ triangularly equivalent: they
have the same point, and there are complex numbers $c_{\tau',ij}$, with $c_{\tau',ij}=0$ for $j>i$, such that the
frame vector $u'_{\tau',i}$ of $y(c,\gamma x)$ is $\sum_jc_{\tau',ij}u_{\tau',j}$, where the $u_{\tau',j}$ are the
frame vectors of $y(c,x)$. By (ii) there are periods $P$ and $P'$ of $L_c$ at $x$ and at $\gamma x$, with the
periods of $\varphi_c(a)$ at the letters as coordinates, such that $y(c,x)$ and $y(c,\gamma x)$ have the points
$y_c(P)$ and $y_c(P')$ and their frame vectors are the derivatives at $0$ of $\operatorname{ex}_{c,P}$ and
$\operatorname{ex}_{c,P'}$ along the standard basis vectors $\varepsilon_j$; write $\varepsilon_{\tau',i}$ for
$\varepsilon_j$ with $\ell_j=(\tau',i)$. As the points are equal, $y_c(P')=y_c(P)$.

\emph{An automorphism of $L_c$.} By Lemma~\ref{is-onto-step-run-model-exists-components-congruence-exponential},
applied to $P'$ and $P$, there are a $\mathbf C$-linear automorphism $h$ of $\mathbf C^g$ and an automorphism
$\tilde\gamma$ of $L_c$ that preserves $\langle\ ,\ \rangle_c$, with $\operatorname{ex}_{c,P}(hz)=\operatorname{ex}_{c,P'}(z)$,
$P(\tilde\gamma a)=h(P'(a))$, and $\tilde\gamma(t_ca)-t_ca\in ML_c$ for $a\in\bar{\mathfrak M}\mathfrak L$. As in the proof
of Lemma~\ref{is-onto-step-run-model-exists-components-congruence-rigidity}, $t_c$ is onto modulo $M'L_c$ for an
integer $M'$ divisible by $M$, so $\tilde\gamma(a)-a\in ML_c$ for $a\in\bar{\mathfrak M}L_c$, and by
Lemma~\ref{cv-volume-key-models-half} for every $a\in L_c$.

\emph{The columns of $h$.} The frame vector $u'_{\tau',i}$ is the derivative at $0$ of
$\operatorname{ex}_{c,P'}=\operatorname{ex}_{c,P}\circ h$ along $\varepsilon_{\tau',i}$, that is, the derivative of
$\operatorname{ex}_{c,P}$ along $h(\varepsilon_{\tau',i})$; and $\sum_jc_{\tau',ij}u_{\tau',j}$ is the derivative of
$\operatorname{ex}_{c,P}$ along $\sum_jc_{\tau',ij}\varepsilon_{\tau',j}$. The derivative at $0$ of
$\operatorname{ex}_{c,P}$ is an injective $\mathbf C$-linear map (Lemma~\ref{cv-volume-key-normalized-tangent}(1)), so
$h(\varepsilon_{\tau',i})=\sum_jc_{\tau',ij}\varepsilon_{\tau',j}$. Put $A_{\tau,ij}=c_{\tau',ji}$ for
$\tau=s\circ\varepsilon_{\tau'}$; then $A_{\tau,ij}=0$ for $j<i$, and for every $a\in L_c$ the coordinate of
$P(\tilde\gamma a)=h(P'(a))$ at the letter $(\tau',i)$ is $\sum_jA_{\tau,ij}$ times the coordinate of $P'(a)$ at
$(\tau',j)$.

\emph{The isometry $\eta$.} Let $g_m=\varphi_c\circ\tilde\gamma\circ\varphi_c^{-1}$, an $\mathcal O_E$-linear
automorphism of $\Lambda_c$. As $\Lambda_c$ spans $E^r$, $g_m$ is the restriction of an $E$-linear map $\eta$ of $E^r$,
which is invertible as $g_m$ is. By (i), $\operatorname{Tr}_{E/\mathbf Q}(-\xi_0h_c(\eta u,\eta v))=\operatorname{Tr}_{E/\mathbf
Q}(-\xi_0h_c(u,v))$ for $u,v\in\Lambda_c$; applied to $\alpha u$ for $\alpha\in\mathcal O_E$, this gives
$\operatorname{Tr}_{E/\mathbf Q}\big(\alpha\xi_0(h_c(\eta u,\eta v)-h_c(u,v))\big)=0$ for every $\alpha$, so
$h_c(\eta u,\eta v)=h_c(u,v)$, as the trace form is nondegenerate and $\xi_0\ne0$; and this extends to $E^r$, as
$\Lambda_c$ spans $E^r$. Also $v\in\Lambda_c$ if and only if $\eta v\in\Lambda_c$. For $w=\varphi_c(a)\in\Lambda_c$,
$\eta w-w=\varphi_c(\tilde\gamma a-a)\in M\Lambda_c$. By the columns of $h$, for every $v\in\Lambda_c$, every
embedding $\tau$ and every $1\le i\le m_\tau$,
\[
 \pi_{\tau,i}(x,\eta v)=\sum_{j=1}^{m_\tau}A_{\tau,ij}\,\pi_{\tau,j}(\gamma x,v).
\]

\emph{The triangular element.} Lemma~\ref{is-onto-step-run-model-exists-components-back-periods}, with $X=X'=X_c$,
$a=\eta$, the point $\gamma x$ in place of $x$ and the point $x$ in place of $x'$, gives $b\in\mathcal B^{\Psi_c}$ with
$x=(\eta\gamma x)b$.
\end{proof}
% Lean: CongrOrbit.exists_congr_back, followed: dom_lact; the triangular law hU.2.2.2.1 of IsUnf at b1 = 1 for
% FrBor (with the coefficients cf, cf i = t, cf j = 0 for j > i, and IsCombAt for the frame vectors); hpt k at x
% and at gamma x; T.familyUnif_exp_unique_alt (Lemma is-onto-step-run-model-exists-components-congruence-exponential)
% for P' and P; hM' through F.t_surj and l.M | l.M', and LevelHalf.trivial_mod_of_trivial_on_half (Lemma
% cv-volume-key-models-half); hcol from DerLin.isCombAt_iff_isDerivAt0, ExpChain.isDerivAt0_comp_linear and
% isDerivAt0_ex_isTanAt (cited in Lemma cv-volume-key-normalized-tangent, stated in the Lean for any family
% uniformization); A tau i j = cf tau j i and hAr (apply_of_cols); gm and hgiso; T.exists_rational_isometry_of_
% trace_isometry (= TraceIsometry.exists_gl_of_trace_isometry, 90 lines, followed above for X = X': its epsilon is
% 1 as xi0 is the same, the matrix A from FrameFibre.exists_matrix and FrameFibre.eq_zero_of_trace); hcong;
% Relab.relab_good with the type of Psi and p0 = q (m1, m2); and T.frames_differ_by_isometry_and_triangular_of_periods
% (Lemma is-onto-step-run-model-exists-components-back-periods) with isMatch_of_isUStab. Two narrower open steps
% remain: Lemmas is-onto-step-run-model-exists-components-congruence-exponential (about 960 lines, shared with the
% congruence) and is-onto-step-run-model-exists-components-back-periods (about 530 lines).

\begin{lemma}\label{is-onto-step-run-model-exists-components}
Let $p$ be a point of a family whose further conditions hold, with definite-side datum of rank $r_0$,
host field $E_h$ and ideal $\mathfrak A$ of the host
(Lemma~\ref{is-onto-step-run-model-exists-components-flag}), and let $r$ be a pair of runs at $p$ pinned to
the own places. Let $i$ be a lift, $v$ a finite place of $L$, $s$ the run used at $v$ and $C_s$ the set of
components of the unitary data of $s$ at level $i+1$ and $\sigma_0$. For $c\in C_s$ let $X_c$ be the
data (1) to (6) of Definition~\ref{is-onto-share-flag-data} that the component $c$ carries, with set
$\Psi_c$, matrix $A_c$, integer $p_{0,c}$, module $\Lambda_c$, group $\Gamma_c$, characters
$\chi_{c,\kappa}$ and element $\xi_{0,c}$, and put
$L'_c=\{x\in E_h^{r_0}:ax\in\Lambda_c\text{ for all }a\in\mathfrak A\}$. Then:
\begin{enumerate}
\item $C_s$ is not empty, and $\Psi_c$, $A_c$, $p_{0,c}$ and $\xi_{0,c}$ do not depend on $c$;
\item every $\chi_{c,\kappa}$ is trivial on $\Gamma_c$;
\item every $X_c$ extends to flag data (Definition~\ref{is-onto-share-flag-data}) with self-dual module
  $L'_c$;
\item $\gamma w-w\in M_bL'_c$ for every $c\in C_s$, $\gamma\in\Gamma_c$ and $w\in L'_c$, where $M_b$ is the
  level ideal of the core.
\end{enumerate}
\end{lemma}
\leanref{Development/AbelianVarietyTheory/Proofs/DivisorProofs/UsedRunUniformisationComponentsShareCmDataStatement.lean}{GlobalTheory.T.uniformisation_components_share_cmData}
\leanref{Development/AbelianVarietyTheory/Proofs/DivisorProofs/UsedRunUniformisationCharacterTrivialArithmeticGStatement.lean}{GlobalTheory.T.uniformisation_character_trivial_arithmeticGroup}
\leanref{Development/AbelianVarietyTheory/Proofs/DivisorProofs/MemLatPIffMulVec.lean}{GlobalTheory.ShareOn.exists_dataX_latP}
\leanref{Development/AbelianVarietyTheory/Proofs/DivisorProofs/LatticePreservedByComponentStabilizerStatement.lean}{GlobalTheory.T.lattice_preserved_by_component_stabilizer}
\begin{proof}
As $r$ is pinned to the own places, $s$ satisfies the laws. (1) The dictionary of the core gives a
component at level $i+1$ and $\sigma_0$, so $C_s$ is not empty. The uniformization law of $s$ gives
every component the same $\Psi$ and $p_0$, and its law on the models gives every component the same
$A$ and $\xi_0$. (2) This is the law of $s$ on the characters of its unitary data.

(3) The core of $s$ satisfies the laws of the core, and $\sigma_0$ is a good embedding of its coefficient
field. By Lemma~\ref{is-onto-step-run-model-exists-components-congruence}, at the level $i+1$ and
$\sigma_0$, every $\Gamma_c$ contains the principal congruence subgroup of level $M$, the elements that
preserve $h_c$ and $\Lambda_c$ and satisfy $\gamma u-u\in M\Lambda_c$ for $u\in\Lambda_c$. So
Lemma~\ref{is-onto-step-run-model-exists-components-flag} extends $X_c$ to flag data whose module
$\Lambda_c$ is the $\mathcal O_{E_h}$-span of the $au$ with $a\in\mathfrak A$ and $u$ in its self-dual
module $L'$. For such $L'$, $L'=L'_c$: if
$x\in L'$, then $ax\in\Lambda_c$ for $a\in\mathfrak A$; and if $ax\in\Lambda_c=\mathfrak AL'$ for every
$a\in\mathfrak A$, then $x\in\mathfrak A^{-1}\mathfrak AL'=L'$, as the fractional ideal $\mathfrak A$
is invertible.

(4) Let $c\in C_s$ and $\gamma\in\Gamma_c$, and let $x_0$ be the point of the flag domain of $X_c$ that the
dictionary law of $s$ gives at $c$. By the unfolding law of $s$, the uniformization of the core at level
$i+1$ and $\sigma_0$ is as in Situation~\ref{is-eigenvalue-conjugate-frame-uniformization}, with
$r=r_0\geq1$, and by the pins of the core the integer $N$ of its level is $1$. So
Lemma~\ref{is-onto-step-run-model-exists-components-back} gives an element $\eta$ that preserves $h_c$ and
$\Lambda_c$ and satisfies $\eta u-u\in M\Lambda_c$ for every $u\in\Lambda_c$, and an invertible upper
triangular $b\in\mathcal B^\Psi$ with $x_0=(\eta\gamma x_0)b$; and $\eta\in\Gamma_c$ by
Lemma~\ref{is-onto-step-run-model-exists-components-congruence}.
% Lean: (3) is LatticeData.exists_dataX_of_hasPrinc (Lemma is-onto-step-run-model-exists-components-flag,
% inside ModelLatticeCheck.holds_by_value) with PrincCore.hasPrinc_core (Lemma
% is-onto-step-run-model-exists-components-congruence); (4) is T.exists_principal_congruence_back_triangular_frame
% (Lemma is-onto-step-run-model-exists-components-back) with the unfolding law isUnf_ofPj, the full level
% l.A = 1 and 1 <= r from numerics_of_pins and the point x0 of the dictionary law dictU, and hasPrinc_core
% for eta in Gamma_c. Three narrower open steps remain below these, in Lemmas
% is-onto-step-run-model-exists-components-congruence-exponential (about 960 lines),
% -congruence-lattice (about 750 lines) and -back-periods (about 530 lines).
Then $(\eta\gamma)^{-1}\in\Gamma_c$ and $(\eta\gamma)^{-1}x_0=x_0b$, so $(\eta\gamma)^{-1}=1$ by (4) of
Definition~\ref{is-onto-share-flag-data}, and $\gamma=\eta^{-1}$. For $u\in\Lambda_c$, $\gamma u\in\Lambda_c$
and $\eta(\gamma u)=u$, so $\gamma u-u=-(\eta(\gamma u)-\gamma u)\in M\Lambda_c$. Let $w\in L'_c$ and
$z=M^{-1}(\gamma w-w)$. For $a\in\mathfrak A$, $aw\in\Lambda_c$ and $az=M^{-1}(\gamma(aw)-aw)\in\Lambda_c$;
so $z\in L'_c$. As $M_b\bar M_b=M\mathcal O_{E_h}$, $M\in M_b$, and $\gamma w-w=Mz\in M_bL'_c$.
% Lean: (1) is ModelLatticeCheck.comp_nonempty (the component of (T.dict j sigma 1)) and
% ShareOn.comps_share (the law unf for p0 and Psi, modelW for Gm and xi0); (2) is
% ShareOn.uAt_chi_eq_one (the law chiTriv, through OneEmbeddingProbe.chi_eq_one); (3) is
% exists_dataX_latP: ModelLatticeCheck.holds_by_value (LatticeData.exists_dataX_of_hasPrinc with
% hasPrinc_core; Fits: toData, and Lambda = span of AA L') and ModelLatExamples.eq_latP, with
% lattice2 D Lambda = {v : a v in Lambda for a in AA}; about 4,300 lines, 2,500 in theorems; (4) is
% ShareOn.sub_mem_half_smul_latP (Gamma_free for (eta gamma)^-1, hcong, hzP, natCast_mem_half).
\end{proof}

\begin{lemma}\label{is-onto-step-run-model-exists-transport}
Let $p$, $r$, $i$, $v$, $s$, $C_s$, $X_c$ and $L'_c$ be as in
Lemma~\ref{is-onto-step-run-model-exists-components}, with the group $G$ of $p$ and the maps $\ell_i$ of
$r$. Then there are an action of $G$ on $C_s$ and elements $t(g,c)\in GL_{r_0}(E_h)$, for $g\in G$ and
$c\in C_s$, such that:
\begin{enumerate}
\item $t(g,c)$ carries the Hermitian form of $A_c$ to that of $A_{gc}$, and
  $t(g,c)\Lambda_c=\Lambda_{gc}$;
\item $\gamma\in\Gamma_c$ if and only if $t(g,c)\gamma t(g,c)^{-1}\in\Gamma_{gc}$;
\item $t(gg',c)\big(t(g,g'c)t(g',c)\big)^{-1}\in\Gamma_{gg'c}$;
\item $g=1$ whenever $gc=c$ and $t(g,c)\in\Gamma_c$;
\item there are scalars $a(g,c)\in\mathbf C^\times$ such that for every cusp form $\phi$ of $s$ at level
  $i+1$ and weight $\kappa_i$, every $g\in G$, $c\in C_s$ and $x$ in the flag domain of $X_c$, the
  reading of the transform of $\phi$ by $\ell_i(g^{-1})$ at $c$ and $x$ is $a(g,c)$ times the reading of
  $\phi$ at $gc$ and $t(g,c)x$;
\item there are isomorphisms $\phi_c$ from the level module of the core at level $i+1$ onto
  $L'_c/M_bL'_c$ such that for every $c$ and every $u\in G$ with $uc=c$ there is an integer $n$ with
  $t(u,c)L'_c\subset L'_c$ and $t(u,c)w\equiv\phi_c\big(\ell_i(u)^n\phi_c^{-1}(w)\big)\pmod{M_bL'_c}$ for
  all $w\in L'_c$.
\end{enumerate}
\end{lemma}
\leanref{Development/AbelianVarietyTheory/Proofs/DivisorProofs/ExistsLevelGroupTransportUniformisationComponentStatement.lean}{GlobalTheory.T.exists_levelGroup_transport_uniformisation_components}
\begin{proof}
\openstep{The construction of the transports between the components through the level changes of
the group, their uniqueness up to the groups $\Gamma_c$, and the reading of cusp forms along them are
not written out.}
% Lean: T.exists_levelGroup_transport_uniformisation_components = ShareOn.exists_transport, which
% assembles T.exists_levelChange_transporters_all_components, T.unique_up_to_target_group,
% T.identity_transports_component_to_itself, T.isLevelComponentTransport_comp,
% T.isTransport_identity_of_mem_componentGroup, T.levelAction_mul_and_faithful_finiteGroup and
% T.levelComponentTransport_reads_cuspForms_const_mul: the action and t(g, c) are a chosen transporter
% of r.lev i g^-1, (1) is IsMatch (iso, lat), and (6) holds with n = 1. Size, following imports and
% stopping at statements the book cites: about 27,600 lines below exists_transport, 21,700 of them in
% theorems.
\end{proof}

\begin{lemma}\label{is-onto-step-run-model-exists}
Let $p$ be a point of a family whose further conditions hold, and $r$ a pair of runs at $p$ pinned to
the own places. Then every lift $i$ has a run model (Definition~\ref{is-onto-step-run-model}).
\end{lemma}
\leanref{Development/AbelianVarietyTheory/Proofs/DivisorProofs/ExistsRunModel.lean}{GlobalTheory.ShareOn.exists_runModel}
\begin{proof}
The ring $\mathcal O_L$ is not a field, so it has a nonzero prime ideal; let $v$ be this finite place of
$L$. Let $s$ be the run used at $v$ and $C=C_s$, and use the notation of
Lemma~\ref{is-onto-step-run-model-exists-components}. Take the action of $G$ on $C$, the elements
$t(g,c)$, the scalars $a(g,c)$ and the isomorphisms $\phi_c$ of
Lemma~\ref{is-onto-step-run-model-exists-transport}. By
Lemma~\ref{is-onto-step-run-model-exists-components}(3), extend each $X_c$ to flag data $X'_c$ with
self-dual module $L'_c$, and fix $c_0\in C$, by part (1) of that lemma.

\emph{The two-level data.} Let $M$ consist of the $\Psi$, $A$, $p_0$ and $\xi_0$ of $X_{c_0}$, the flag
data $X'_c$ for $c\in C$, and the elements $t(g,c)$. By
Lemma~\ref{is-onto-step-run-model-exists-components}(1), every $X'_c$ has these $\Psi$, $A$, $p_0$ and
$\xi_0$, so all of them have the same flag domain $\mathcal D_M$; by part (2), every $\chi_{c,\kappa}$ is
trivial on $\Gamma_c$. For $\beta\in GL_{r_0}(E_h)$ and $c,c'\in C$ with $\beta\Lambda_c=\Lambda_{c'}$, also
$\beta L'_c=L'_{c'}$, since $a\beta x=\beta(ax)$ for $a\in\mathfrak A$. The group $\Gamma_c$ preserves
$\Lambda_c$ ((4) of Definition~\ref{is-onto-share-flag-data}), so it preserves $L'_c$; and by
Lemma~\ref{is-onto-step-run-model-exists-transport}(1), $t(g,c)$ preserves the Hermitian form of $A$ and
$t(g,c)L'_c=L'_{gc}$. The remaining conditions of Definition~\ref{is-onto-share-two-level} are (2), (3)
and (4) of Lemma~\ref{is-onto-step-run-model-exists-transport}.

\emph{The weight.} Let $\kappa$ be the weight of $s$ at $c_0$ for the weight family $\kappa_i$. The weight
of $s$ at a component depends on the component only through its $\Psi$ and $p_0$, which are the same
for all components; so $\kappa$ is the weight of $s$ at every $c\in C$.

\emph{The run model.} Take $v$ and the identity of $C$ for $e$ in
Definition~\ref{is-onto-step-run-model}. Its condition (1) holds by the construction of $M$ and the
choice of $\kappa$; (2) is (5) of Lemma~\ref{is-onto-step-run-model-exists-transport}, as the flag domain
of every $X_c$ is $\mathcal D_M$; (3) is Lemma~\ref{is-onto-step-run-model-exists-components}(4); and
(4) is (6) of Lemma~\ref{is-onto-step-run-model-exists-transport}.
% Lean: exists_runModel: v from nonempty_place, hne = comp_nonempty, hcom, hchi, hcong and the
% transport of Lemmas is-onto-step-run-model-exists-components and -transport, Xx from
% exists_dataX_latP; M has comp = Xx, Psi, Gm, p0, xi0 of the component k0, chi_one from
% chi_one_of_toData, Gamma_L' and tr_L' from mem_localLattice_iff_mulVec (with Gamma_stab and
% IsMatch.lat), tr_herm from IsMatch.iso and the equal Gm; ka = laR at k0, equal at every component as
% laU depends on the component through p0 and Psi; e = Equiv.refl. Size, following imports and
% stopping at statements the book cites: about 28,600 lines below exists_runModel, 22,700 of them in
% theorems, nearly all in the transport.
\end{proof}

\begin{lemma}\label{is-onto-step-model-reading}
Let $p$ and the pair of runs be as in Definition~\ref{is-onto-step-run-model}, and suppose that the
number of bad places of the definite-side datum is positive and that $8fr+4\le n_V$, where $r$ is its
rank. Let $(C,M,\kappa)$ be a run
model of a lift $i$. Then the set $\Psi$ of embeddings of $M$ has $f$ elements, and there are a
representation $\rho$ of $G$ on the cusp forms $S_\kappa(M)$, pulled back from the action of $G$ on
$C$, an embedding $\sigma$ of $L$ in $\mathbf C$ and a $\sigma$-semilinear map $f$ from the
coordinates $L^{n_W(i)}$ of $FW_i$ to $S_\kappa(M)$ that is compatible with the actions of $G$ and
carries the standard basis to a basis of $S_\kappa(M)$.
\end{lemma}
\leanref{Development/AbelianVarietyTheory/Proofs/DivisorProofs/ReadingOfRunModel.lean}{GlobalTheory.ShareOn.Reading.reading_of_runModel}
\begin{proof}
\openstep{The reading of the rational structure of a lift as the cusp forms of a run model is not
written out.}
% Lean: ShareOn.Reading.reading_of_runModel. Size: about 146,000 lines of Lean in 8,600 theorems below
% it (text count, stopping at statements the book cites), about 500 of them below no other open step
% of Lemma is-onto-estimates.
\end{proof}

\begin{lemma}\label{is-onto-step-model-level-places}
Let $D$ be a definite-side datum with host field $E_h$, level primes $g_0,g_1$ and ranks with
$f=2n$, and let $k\in\{0,1\}$. Then every prime ideal of $\mathcal O_{E_h}$ that contains $g_k$ is one
of the $f$ places of $E_h$ over $g_k$ that $D$ gives, or the complex conjugate of one of them.
\end{lemma}
\leanref{Development/AbelianVarietyTheory/Proofs/DivisorProofs/PrimeOverTypePrimeFixedPlaceOrConjStatement.lean}{GlobalTheory.T.prime_over_typePrime_fixedPlace_or_conj}
\leanref{Development/AbelianVarietyTheory/Proofs/InertiaEigenvalues/EqPlGwOrIsConjPl.lean}{GlobalTheory.TypePlacesExhaust.eq_plGw_or_isConjPl}
\begin{proof}
A prime ideal containing $g_k$ is nonzero, so it is a place over $g_k$. The datum gives $f$ distinct
places over $g_k$, none of them the complex conjugate of another or of itself, and $g_k$ has
$[E_h:\mathbf Q]=2f$ places over it. The $f$ places and their $f$ conjugates are $2f$ distinct places
over $g_k$, so they are all of them.
\end{proof}

\begin{lemma}\label{is-onto-step-model-level-half}
Let $E$ be a number field, $r\geq0$ an integer, $L'$ an $\mathcal O_E$-submodule of $E^r$ and
$M_b\subset P$ ideals of $\mathcal O_E$. Let $A$ be an $\mathcal O_E$-module with an isomorphism
$\phi\colon A\to L'/M_bL'$, $\ell$ an automorphism of $A$ with $(\ell-1)^rA\subset PA$, and $n$ an
integer. Let $\gamma,t\in M_r(E)$ satisfy $\gamma L'\subset L'$, $(\gamma-1)L'\subset M_bL'$,
$tL'\subset L'$ and $tw\equiv\phi\big(\ell^n(\phi^{-1}(w))\big)\pmod{M_bL'}$ for $w\in L'$. Then
$(\gamma t-1)^rL'\subset PL'$.
\end{lemma}
\leanref{Development/AbelianVarietyTheory/Proofs/DivisorProofs/UnipotentModIdealOfTrivialModHalfStatement.lean}{GlobalTheory.T.unipotent_mod_ideal_of_trivial_mod_half}
\begin{proof}
Let $U$ be the endomorphism $w\mapsto\gamma tw-w$ of $L'$. Modulo $M_bL'$, $\gamma$ acts as the
identity, so $U$ induces $\phi(\ell^n-1)\phi^{-1}$ on $L'/M_bL'$, and $U^r$ induces
$\phi(\ell^n-1)^r\phi^{-1}$. Now $\ell^n-1=(\ell-1)X$ for an endomorphism $X$ of $A$ that commutes with
$\ell$: $X=1+\ell+\dots+\ell^{n-1}$ if $n\geq0$, and $X=-\ell^n(1+\ell+\dots+\ell^{-n-1})$ if $n<0$. So
$(\ell^n-1)^r=(\ell-1)^rX^r$ maps $A$ into $PA$. Hence $U^rw\in PL'+M_bL'=PL'$ for $w\in L'$.
\end{proof}

\begin{lemma}\label{is-onto-step-model-level-conj}
Let $X$ be flag data of rank $r$ over a CM field $E$ (Definition~\ref{is-onto-share-flag-data}), with
Hermitian form $h$ and self-dual module $L'$. Let $\beta\in GL_r(E)$ preserve $h$, with $\beta L'=L'$,
and let $P$ be an ideal of $\mathcal O_E$ with $(\beta-1)^rL'\subset PL'$. Then
$(\beta-1)^rL'\subset\bar PL'$.
\end{lemma}
\leanref{Development/AbelianVarietyTheory/Proofs/DivisorProofs/UnitaryUnipotentModIdealModConjIdealStatement.lean}{GlobalTheory.T.unitary_unipotent_mod_ideal_mod_conjIdeal}
\begin{proof}
We may assume $P\ne0$. Since $\beta$ preserves $h$, $h((\beta-1)x,z)=h(x,(\beta^{-1}-1)z)$, so
$h((\beta-1)^rx,z)=h(x,(\beta^{-1}-1)^rz)$; and $(\beta^{-1}-1)^r=(-\beta^{-1})^r(\beta-1)^r$ maps $L'$
into $PL'$, as $\beta^{-1}L'=L'$. Choose $c\ne0$ in $P$ and an ideal $P'$ with $PP'=c\mathcal O_E$. For
$p'\in P'$ and $x,z\in L'$, $c^{-1}p'(\beta^{-1}-1)^rz\in L'$, so
\[
 h\big(\bar c^{-1}\bar p'(\beta-1)^rx,z\big)=h\big(x,c^{-1}p'(\beta^{-1}-1)^rz\big)
\]
is an algebraic integer. By self-duality, $\bar c^{-1}\bar p'(\beta-1)^rx\in L'$ for every $p'\in P'$.
Since $\bar P\bar P'=\bar c\mathcal O_E$, this gives $(\beta-1)^rx\in\bar PL'$.
\end{proof}

\begin{lemma}\label{is-onto-step-model-level-coordinates}
Let $D$ be a definite-side datum with host field $E/F$ and the ideal $\mathfrak A$ of the host, and write
$\hat{\mathcal O}_E$ for the closure of $\mathcal O_E$ in $E\otimes_F\mathbf A_{F,f}$. Let $W$ be a
Hermitian space of rank $n$ over $E$ (Definition~\ref{th-spaces}), let $L_0\subset E^n$ be its lattice
in the coordinates of its basis, $\hat L_0$ the closure of $L_0$ in $(E\otimes_F\mathbf A_{F,f})^n$, and
$\hat L_W$ the subgroup of $(E\otimes_F\mathbf A_{F,f})^n$ generated by the $az$ with $a\in\mathfrak A$
and $z\in\hat L_0$. Let $\beta\in GL_n(E\otimes_F\mathbf A_{F,f})$ be a \emph{self-dual frame} of $W$:
$L_0$ is self-dual for the form of $W$, and a vector $z$ lies in $\hat L_0$ if and only if
$\beta^{-1}z\in\hat{\mathcal O}_E^n$. Let $\Lambda$ be an $\mathcal O_E$-module, $\iota\colon\Lambda\to E^n$
an injective $\mathcal O_E$-linear map with image $\mathfrak AL_0$, and $P$ a nonzero prime ideal of
$\mathcal O_E$ with residue field $k_P$. Then there is a surjective $\mathcal O_E$-linear map
$\phi\colon\Lambda\to k_P^n$ with the following property. Let $k\in U(W)(\mathbf A_F)$ have finite part
$k_f$ with $y=\beta^{-1}k_f\beta\in GL_n(\hat{\mathcal O}_E)$, let $M$ be an integer that lies in $P$,
and let $a,b\in\Lambda$ satisfy $\iota(b)-k_f\iota(a)\in M\hat L_W$. Then $\phi(b)=\bar y\,\phi(a)$, where
$\bar y\in GL_n(k_P)$ is the reduction of $y$ modulo $P$.
\end{lemma}
\leanref{Development/AbelianVarietyTheory/Proofs/DivisorProofs/LatticeCoordinatesModPrimeSelfDualStabilizerStatement.lean}{GlobalTheory.T.lattice_coordinates_mod_prime_selfDual_stabilizer}
\leanref{Development/AbelianVarietyTheory/Proofs/DivisorProofs/LatticeCoordinatesModPrimeSelfDualStabilizer.lean}{GlobalTheory.LatCoordRed.coord_red}
\begin{proof}
Write $\hat{\mathcal O}_F$ for the closure of $\mathcal O_F$ in $\mathbf A_{F,f}$, and identify
$\hat{\mathcal O}_E/P\hat{\mathcal O}_E$ with $k_P$.

\emph{A scaling.} The fractional ideal $\mathfrak A$ is nonzero, so $\mathfrak A^{-1}\mathfrak A=\mathcal
O_E$ is not contained in $P$, and there are $c\in\mathfrak A^{-1}$ and $a_1\in\mathfrak A$ with $ca_1\notin
P$. As $P$ is maximal, there is $c'\in\mathcal O_E$ with $c'ca_1\equiv1\pmod P$. Put $a_0=c'a_1\in\mathfrak
A$ and $e_0=ca_0$. Then $c\mathfrak A\subset\mathcal O_E$, $e_0\in\mathcal O_E$ and $e_0\equiv1\pmod P$.

\emph{The map.} For $z\in\hat L_W$ put $\rho(z)=c\beta^{-1}z$. If $z=az'$ with $a\in\mathfrak A$ and
$z'\in\hat L_0$, then $\rho(z)=(ca)\beta^{-1}z'$ lies in $\hat{\mathcal O}_E^n$, since $ca\in\mathcal O_E$;
so $\rho(\hat L_W)\subset\hat{\mathcal O}_E^n$. Since $\iota(\Lambda)=\mathfrak AL_0\subset\hat L_W$, we
may put $\phi(a)=\rho(\iota(a))\bmod P$. This map is $\mathcal O_E$-linear.

\emph{Surjectivity.} Let $v\in k_P^n$, and lift it to $v'\in\mathcal O_E^n$. The vector $z=a_0\beta v'$
lies in $\hat L_W$, as $\beta v'\in\hat L_0$, and $\rho(z)=e_0v'$. Let $N$ be the absolute norm of $P$, a
positive integer in $P$. Every element of $\hat{\mathcal O}_F$ is $s+Ny$ with $s\in\mathcal O_F$ and
$y\in\hat{\mathcal O}_F$, since $\mathcal O_F/N\mathcal O_F=\hat{\mathcal O}_F/N\hat{\mathcal O}_F$. The
group $\hat L_W$ is generated by the vectors $a(b\otimes x)$ with $a\in\mathfrak A$,
$b\in\hat{\mathcal O}_F$ and $x\in L_0$, and writing $b=s+Ny$ gives $a(b\otimes x)=asx+N\,a(y\otimes x)$
with $asx\in\mathfrak AL_0$. Hence $z=x+Nw$ with $x\in\mathfrak AL_0$ and $w\in\hat L_W$. Write
$x=\iota(a)$. Then $\rho(\iota(a))=e_0v'-N\rho(w)$ with $\rho(w)\in\hat{\mathcal O}_E^n$, and reducing
modulo $P$ gives $\phi(a)=v$, as $e_0\equiv1$ and $N\in P$.

\emph{The congruence.} Let $k,y,M,a,b$ be as in the statement, and write $\iota(b)=k_f\iota(a)+Mw$ with
$w\in\hat L_W$. As $\beta^{-1}k_f=y\beta^{-1}$,
\[
 \rho(\iota(b))=y\,\rho(\iota(a))+M\rho(w).
\]
The three vectors $\rho(\iota(b))$, $\rho(\iota(a))$ and $\rho(w)$ lie in $\hat{\mathcal O}_E^n$ and
$M\in P$, so reducing modulo $P$ gives $\phi(b)=\bar y\,\phi(a)$.
% Lean: LatCoordRed.coord_red. The scaling is exists_scale, rho is crd (c beta^{-1}, after the swap
% A_f (x) E = E (x) A_f) with crd_mem, phi is resS of crd; surjectivity is exists_of_coords (z = a0 beta
% v') and exists_approx (Formalism.UG.P.integral_adele_approximation_holds with N = absNorm P); the
% congruence is crd_mulVec with LevelActionPinned.map_finHom and resS_combo. The self-duality of L_0
% and the unitarity of k are not used; beta enters through mem_latHat_iff.
\end{proof}

\begin{definition}\label{is-onto-step-model-level-condition}
Let $p$ be a point of a family, with definite-side datum $D$ of host field $E_h/F_h$, rank $r$, ideal
$\mathfrak A$ of the host and lifts $1,\dots,m$, and with realization $R$, group $G$, identified with
$\prod_{P\in P_g}GL_{r}(k_P)$, and, for each lift $i$, frame
$\beta_i\in GL_{r}(E_h\otimes_{F_h}\mathbf A_{F_h,f})$ of $W_i$. The core of $D$ has a \emph{level}:
an integer $M\geq3$ and an ideal $M_b$ of $\mathcal O_{E_h}$, the \emph{level ideal}, with
$M_b\bar M_b=M\mathcal O_{E_h}$ and $M_b+\bar M_b=\mathcal O_{E_h}$. At each level $j$ it has a
\emph{lattice} $\Lambda_j$, a finitely generated projective $\mathcal O_{E_h}$-module, with an
$\mathcal O_{E_h}$-linear map $\iota_j\colon\Lambda_j\to E_h^{r}$; the \emph{level module} at level
$j$ is $\Lambda_j/M_b\Lambda_j$. For a lift $i$, let $\hat L_W$ be the subgroup of
$(E_h\otimes_{F_h}\mathbf A_{F_h,f})^{r}$ generated by the $az$ with $a\in\mathfrak A$ and $z$ in the
closure $\hat L_0$ of the lattice of $W_i$ in the coordinates of its basis. A map $\ell$ from $G$ to the
$\mathcal O_{E_h}$-linear automorphisms of $\Lambda_{i+1}/M_b\Lambda_{i+1}$ \emph{gives the level} of
$R$ at $i$ if for every $g\in G$ there is a level lift $k$ of $g$
(Lemma~\ref{is-eigenvalue-core-group}) with
\[
 \iota_{i+1}(s_g(a))-k\,\iota_{i+1}(a)\in M\hat L_W\qquad\text{for every }a\in\bar M_b\Lambda_{i+1},
\]
where $s_g\colon\Lambda_{i+1}\to\Lambda_{i+1}$ is an $\mathcal O_{E_h}$-linear map inducing $\ell(g)$,
fixed for each $g$. The level condition (c) of Definition~\ref{is-onto-step-runs}(5) at the lift $i$
says that $\ell_i$ gives the level of $R$ at $i$. Here a level lift $k$ of $g$ lies in the parahoric
level of $W_i$, so $y=\beta_i^{-1}k\beta_i$ lies in $GL_{r}(\hat{\mathcal O}_{E_h})$; and for $P\in P_g$
the upper left $r\times r$ block of the reduction of $y$ modulo $P$, which is the whole reduction
since $W_i$ has rank $r$, is the component $g_P$ of $g$.
\end{definition}
\leanref{Formalism/LocalToGlobal/IndefiniteReal/ThPinsR.lean}{Formalism.IndefR.FWT.LevIs}
\leanref{Formalism/LocalToGlobal/IndefiniteReal/SubgroupCharExtensionOneApply.lean}{Formalism.IndefR.Lev.IsLevelLift}
\leanref{Formalism/LocalToGlobal/ProjectiveLattices/LiftLaSpec.lean}{Formalism.L2.Pj.liftLa}

\begin{lemma}\label{is-onto-step-model-level-action}
Let $p$ be a point of a family, with definite-side datum of rank $r$ and group $G$, let a pair of runs
at $p$ pinned to the own places be given, with maps $\ell_i$ (Definition~\ref{is-onto-step-runs}), and
let $i$ be a lift, with the level module $A$ of the core at level $i+1$. Let $P\in P_g$ and let $u\in G$
be upper unitriangular at $P$. Then $(\ell_i(u)-1)^{r}A\subset PA$.
\end{lemma}
\leanref{Development/AbelianVarietyTheory/Proofs/DivisorProofs/LevelChangeUnitriangularUnipotentModPlaceStatement.lean}{GlobalTheory.T.levelChange_unitriangular_unipotent_modPlace}
\begin{proof}
Write $\mathcal O=\mathcal O_{E_h}$, and let $M$, $M_b$, $\Lambda=\Lambda_{i+1}$, $\iota=\iota_{i+1}$,
$\beta=\beta_i$ and $\hat L_W$ be as in Definition~\ref{is-onto-step-model-level-condition}, so that
$A=\Lambda/M_b\Lambda$. Choose a finite place $v$ of $L$. By (a) of Definition~\ref{is-onto-step-runs}(5)
the run used at $v$ satisfies the laws. They make $\iota$ injective with image $\mathfrak AL_0$, for the
lattice $L_0\subset E_h^{r}$ of $W_i$ in the coordinates of its basis, and they give $\Lambda$ the rank
$r$. The conditions of the characterization on the level characters, which are part of the point, make
$\beta$ a self-dual frame of $W_i$. So Lemma~\ref{is-onto-step-model-level-coordinates} gives a surjective
$\mathcal O$-linear map $\phi\colon\Lambda\to k_P^{r}$. It vanishes on $P\Lambda$, as $k_P^{r}$ is
killed by $P$.

By (c) of Definition~\ref{is-onto-step-runs}(5), $\ell_i$ gives the level of $R$ at $i$: there are a level
lift $k$ of $u$ and an $\mathcal O$-linear map $s\colon\Lambda\to\Lambda$ inducing $\ell_i(u)$ with
$\iota(s(a))-k\iota(a)\in M\hat L_W$ for every $a\in\bar M_b\Lambda$. Let $y=\beta^{-1}k\beta\in
GL_{r}(\hat{\mathcal O})$ and let $\bar y$ be its reduction modulo $P$. By (e) of
Definition~\ref{is-onto-step-runs}(5), $M_b=\prod_{Q\in P_g}Q$, which lies in $P$; and $M\in M_b$, as
$M_b\bar M_b=M\mathcal O$. So $M\in P$.

\emph{Step 1: $\phi(s(a))=\bar y\,\phi(a)$ for every $a\in\Lambda$.} As $M_b+\bar M_b=\mathcal O$, write
$1=c+w$ with $c\in M_b$ and $w\in\bar M_b$. Then $wa\in\bar M_b\Lambda$, and
Lemma~\ref{is-onto-step-model-level-coordinates}, applied to $k$, $M$, $wa$ and $s(wa)$, gives
$\phi(s(wa))=\bar y\,\phi(wa)$. Since $c\in P$, $\phi(wa)=\phi(a)-c\phi(a)=\phi(a)$ and
$\phi(s(a))=\phi(s(wa))+c\phi(s(a))=\phi(s(wa))$.

\emph{Step 2: $(\bar y-1)^{r}=0$.} Since $k$ is a level lift of $u$, $\bar y$ is the component $u_P$ of
$u$. The point also describes $G$ as the product of the groups $GL_{r}(\mathbf Z/g_k)$ over the places
over $g_0$ and $g_1$ that $D$ gives, and identifies these places with $P_g$; at the place $P$ over $g_k$,
$u_P$ is the image of the component of $u$ in $GL_{r}(\mathbf Z/g_k)$ under a ring isomorphism
$\mathbf Z/g_k\cong k_P$. That $u$ is upper unitriangular at $P$ means that this component is, so $u_P$ is
upper unitriangular. Hence $\bar y-1$ is strictly upper triangular, and $(\bar y-1)^{r}=0$.

\emph{Step 3: the kernel of $\phi$ is $P\Lambda$.} The map $\phi$ induces a surjection
$\Lambda/P\Lambda\to k_P^{r}$. As $\Lambda$ is finitely generated and projective of rank $r$,
$\Lambda/P\Lambda$ is a $k_P$-vector space of dimension $r$, so both sides have $|k_P|^{r}$ elements
and the surjection is a bijection.

\emph{Step 4.} By Step 1 and induction on $n$, $\phi\big((s-1)^na\big)=(\bar y-1)^n\phi(a)$ for $a\in\Lambda$
and $n\geq0$. For $n=r$ this is $0$ by Step 2, so $(s-1)^{r}a\in P\Lambda$ by Step 3. As $s$ induces
$\ell_i(u)$ on $A$, $(\ell_i(u)-1)^{r}$ maps the class of $a$ to the class of $(s-1)^{r}a$, which lies
in $PA$. Every element of $A$ is such a class.
% Lean: T.levelChange_unitriangular_unipotent_modPlace = ShareOn.UnipLevel.levUnip_mod_place: LevIs
% (hr.level v' i) gives kappa with Lev.IsLevelLift and the congruence modulo M latGHat on Mfr La;
% T.lattice_coordinates_mod_prime_selfDual_stabilizer (LatCoordRed.coord_red) gives surjective
% coordinates phi : La -> kP^r with phi(b) = red(kappa) phi(a) under that congruence; with
% Mb + Mfr = 1 and Mb <= P (hr.half), phi(f a) = N phi(a) for N = red(kappa), whose block is the image
% of u at P (p.grpRecPin, upper_block_val_submatrix); (N - 1)^r = 0 (pow_apply_eq_zero_of_strict), and
% ker phi = P La by counting (natCard_quot_smul_top_of_projective). Size: about 470 lines of Lean in
% the two files LevelChangeUnitriangularUnipotentModPlace.lean and
% LatticeCoordinatesModPrimeSelfDualStabilizer.lean, with their helpers. The place v is
% ShareOn.nonempty_place; the laws hr.laws v give LatEmb through TraceLattice.latEmb_iotaW and the rank
% r = p + q through hL.j.unf; the self-dual frame is p.charW.base.lev.betaW i. Step 1 is hfphi (1 = c1 + w
% from p.T.Q.l.coprime, hkill), Step 2 is h5, h6, h7, hsub, hstrict and hnil, Step 3 is hker, Step 4 is
% hpow, hmk and LevDesc.map_mem_smul_top.
\end{proof}

\begin{lemma}\label{is-onto-step-model-level}
\begin{enumerate}
\item Let $\mathcal H$ be the class of standard families. Eventually along $\mathcal H$, every family
  has a scale beyond which every point satisfying the point conditions and the further conditions,
  with definite-side datum of rank $r$ and host primes $g_0,g_1$, has $3r<g_0$ and $3r<g_1$.
\item Let $p$ and the pair of runs be as in Definition~\ref{is-onto-step-run-model}, with the primes
  $g_0,g_1$ of the host, and suppose that $3r<g_0$ and $3r<g_1$, where $r$ is the rank of the
  definite-side datum. Let $(C,M,\kappa)$ be a run model of a lift
  $i$. Then the level hypothesis of Definition~\ref{is-onto-share-two-level} holds for $M$ and every
  half-unipotent subgroup of $G$: every subgroup whose elements are upper unitriangular at all the
  places of $P_g$ over $g_0$, or at all those over $g_1$.
\end{enumerate}
\end{lemma}
\leanref{Development/AutomorphicArithmetic/ThetaCorrespondence/ThetaPairLifts/ThirdPlaceSeries/TypePrimesGtThreeR.lean}{GlobalTheory.BigOn.typePrimes_gt_three_r_alt}
\leanref{Development/AbelianVarietyTheory/Proofs/DivisorProofs/SplitsCompletelyTypePrime.lean}{GlobalTheory.ShareOn.levelCondition_restrict_of_halfUnipotent}
\begin{proof}
(1) The bound holds for every family, at every point above the threshold, with any largeness
threshold. Write $k,m$ for the ranks of the member, so that $r=k+m-1\le2k$ as $m+3\le k$. The
parameter choices of the member give a number $Y$ with $(kb)^8\le Y\le g_j$ for $j=0,1$. Since $k\geq1$
and $b=e^{x^\varepsilon}\geq2$, $(kb)^8\geq2k\cdot2^7>6k\geq3r$, so $3r<g_0$ and $3r<g_1$.

(2) Let $H$ be half-unipotent, and $k\in\{0,1\}$ such that every element of $H$ is upper unitriangular
at all the places of $P_g$ over $g_k$. We show the level hypothesis for $M|_H$ with the prime $g_k$.
The datum makes $g_k$ a prime with $[E_h:\mathbf Q]$ places of $E_h$ over it, so $g_k$ splits
completely in $E_h$, and $3r<g_k$ by hypothesis. Let $c\in C$, $u\in H$ with $uc=c$, $\gamma\in\Gamma_c$
and $P$ a prime ideal of $\mathcal O_{E_h}$ containing $g_k$, and put $\beta=\gamma t(u,c)$. We show
$(\beta-1)^{r}L'_c\subset PL'_c$.

The realization identifies $P_g$ with the places over $g_0$ and $g_1$ that the datum gives, so by
Lemma~\ref{is-onto-step-model-level-places}, $P$ is a place $P'\in P_g$ over $g_k$ or its complex
conjugate.

\emph{The case $P=P'$.} By (e) of Definition~\ref{is-onto-step-runs}(5), the level ideal $M_b$ of the
core is $\prod_{Q\in P_g}Q$, which lies in $P'$. Let $\phi_c$ and $n$ be as in (4) of
Definition~\ref{is-onto-step-run-model}, for $c$ and $u$. Since $u$ is upper unitriangular at $P'$,
Lemma~\ref{is-onto-step-model-level-action} gives $(\ell_i(u)-1)^{r}A\subset P'A$ for the level module
$A$ of the core at level $i+1$. The group $\Gamma_c$ preserves $L'_c$
(Definition~\ref{is-onto-share-two-level}), $(\gamma-1)L'_c\subset M_bL'_c$ by (3) of
Definition~\ref{is-onto-step-run-model}, $t(u,c)L'_c=L'_{uc}=L'_c$, and (4) is the congruence for
$t(u,c)$. So Lemma~\ref{is-onto-step-model-level-half}, with $A$, $\phi_c$, $\ell_i(u)$, $n$, $\gamma$
and $t(u,c)$, gives $(\beta-1)^{r}L'_c\subset P'L'_c$.

\emph{The case $P=\bar P'$.} The element $\beta$ preserves $h$, as $\gamma$ and $t(u,c)$ do, and
$\beta L'_c=L'_c$. By the first case $(\beta-1)^{r}L'_c\subset P'L'_c$, so
Lemma~\ref{is-onto-step-model-level-conj} gives $(\beta-1)^{r}L'_c\subset\bar P'L'_c=PL'_c$.
% Lean: (1) is typePrimes_gt_three_r_alt, through BigOn.three_r_lt_g_of_params (the field g_Y of
% ParameterChoice'). (2) is ShareOn.levelCondition_restrict_of_halfUnipotent, with
% splitsCompletely_typePrime (D.nd.g_split), T.prime_over_typePrime_fixedPlace_or_conj, half_le_plGw
% (h.half and p.grpRecPin), T.unipotent_mod_ideal_of_trivial_mod_half,
% T.levelChange_unitriangular_unipotent_modPlace and T.unitary_unipotent_mod_ideal_mod_conjIdeal.
\end{proof}

\begin{lemma}\label{is-onto-step-model-arch}
Let $\mathcal H$ be the class of standard families. Eventually along $\mathcal H$, every family has a
scale beyond which the following holds for every point $p$ that satisfies the point conditions and the
further conditions, at a member with integer $T_O$ of its parameter choices, every pair of runs at
$p$ pinned to the own places, every lift $i$ and every run model $(C,M,\kappa)$ of $i$: either
$p_0\in\{0,r\}$, or $0<p_0<r$ and $M$ satisfies the archimedean condition at $T_O$ for the weight
$\kappa$ (Definition~\ref{is-onto-share-two-level}), where $r$ is the rank of the definite-side
datum.
\end{lemma}
\leanref{Development/AutomorphicArithmetic/ThetaCorrespondence/ThetaPairLifts/ThirdPlaceSeries/RunModelArch.lean}{GlobalTheory.ShareOn.p0_extreme_or_archOn_of_runModel}
\begin{proof}
\openstep{The archimedean condition for the run models with $0<p_0<r$ is not written out.}
% Lean: ShareOn.p0_extreme_or_archOn_of_runModel. Size: about 683,000 lines of Lean in 38,000 theorems
% below it (text count, stopping at statements the book cites), nearly all shared with
% T.exists_constants_indexBound_hypotheses_liftWeights (Lemma is-onto-step-index-conditions); about
% 400 lines lie below no other open step of Lemma is-onto-estimates.
\end{proof}

\begin{lemma}\label{is-onto-step-models}
Let $\mathcal H$ be the class of standard families. Eventually along $\mathcal H$, every family has a
scale beyond which the following holds for every point $p$ that satisfies the point conditions and the
further conditions, at a member with integer $T_O$ of its parameter choices, every pair of runs at
$p$ pinned to the own places and every lift $i$. There are a finite set $C$ with an action of the group
$G$ of the realization, two-level data $M$ of rank $r$ over the host field for $G$ and $C$, where
$r$ is the rank of the definite-side datum, whose set $\Psi$ of embeddings has $f$ elements, a
weight $\kappa$, a representation $\rho$ of $G$ on $S_\kappa(M)$ pulled back from the action on $C$,
an embedding $\sigma$ of $L$ in $\mathbf C$ and a $\sigma$-semilinear map from $L^{n_W(i)}$ to
$S_\kappa(M)$ compatible with the actions of $G$ that carries the standard basis to a basis, such that
the level hypothesis of $M$ holds for every half-unipotent subgroup of $G$, and either
$p_0\in\{0,r\}$, or $0<p_0<r$ and $M$ satisfies the archimedean condition at $T_O$ for $\kappa$.
\end{lemma}
\leanref{Development/AutomorphicArithmetic/ThetaCorrespondence/ThetaPairLifts/ThirdPlaceSeries/ModelSt.lean}{GlobalTheory.ShareOn.modelSt_alt}
\leanref{Development/AbelianVarietyTheory/Definitions/Divisors/ExistsArchOne.lean}{GlobalTheory.ShareOn.ModelAt}
\begin{proof}
Take the scale beyond which Lemma~\ref{is-onto-step-model-arch} and part (1) of
Lemma~\ref{is-onto-step-model-level} hold. Let $p$, the pair of runs and $i$ be as in the statement,
and let $(C,M,\kappa)$ be a run model of $i$ (Lemma~\ref{is-onto-step-run-model-exists}), at a place
$v$. The core ties of $p$ make the field at $v$ good, so the number of bad places of the definite-side
datum is positive, and the parameter choices give $8fr+4\le n_V$. Lemma~\ref{is-onto-step-model-reading}
gives $|\Psi|=f$, $\rho$, $\sigma$ and the semilinear map. Since $3r<g_0$ and $3r<g_1$,
Lemma~\ref{is-onto-step-model-level}(2) gives the level hypothesis for the half-unipotent subgroups,
and Lemma~\ref{is-onto-step-model-arch} gives the alternative on $p_0$.
\end{proof}

\subsection{The bad-place bound}\label{subsection-is-bad-place-bound}

\begin{lemma}\label{is-onto-step-bad-weight-whittaker}
Let $p$ be a point of a family that satisfies the point conditions, with realization of group $G$,
face with subgroup $P_0$ and character $\psi_0$, rank $r$ of the definite-side datum, and the
embedding $\sigma_0$ of the base place of the coefficient field $L$. Identify $G$ with
$\prod_{k=0,1}\prod_{v\mid g_k}GL_r(\mathbf F_{g_k})$, the product over the level places, through
the isomorphism of $G$ with $\prod_{P\in P_g}GL_r(k_P)$ (Definition~\ref{is-realisation}(1)) and
identifications of the residue fields. Then there are nontrivial additive characters
$\psi_{k,v}\colon\mathbf F_{g_k}\to\mathbf C^\times$ such that $P_0$ is the subgroup of the elements
that are upper unitriangular at every level place, and
\[
 \sigma_0\big(\psi_0(u)\big)^{-1}=\prod_{k,v}\psi_{k,v}\Big(\sum_{i<r}u_{v,i,i+1}\Big)\qquad(u\in P_0).
\]
In particular $P_0$ is half-unipotent.
\end{lemma}
\leanref{Development/AutomorphicArithmetic/ThetaCorrespondence/ThetaPairLifts/ThirdPlaceSeries/BaseLevelIsWhitStatement.lean}{GlobalTheory.T.exists_whittaker_lam_psi_half_unip}
\leanref{Development/AutomorphicArithmetic/ThetaCorrespondence/ThetaPairLifts/ThirdPlaceSeries/BaseLevelIsWhit.lean}{GlobalTheory.ShareOn.isWhit_and_eq_embLL_inv_and_halfUnip}
\leanref{Development/AbelianVarietyTheory/Definitions/Divisors/BadElt.lean}{GlobalTheory.BadCount.IsWhit}
\begin{proof}
Among the point conditions is a level tie: there are nontrivial additive characters $\psi'_P\colon
k_P\to L^\times$, for $P\in P_g$, such that, under the isomorphism of $G$ with
$\prod_{P\in P_g}GL_r(k_P)$, $P_0$ is the subgroup of the elements that are upper unitriangular at every
$P$, and $\psi_0(u)=\prod_{P\in P_g}\psi'_P\big(\sum_{i<r}u_{P,i,i+1}\big)$ for $u\in P_0$; the places
of $P_g$ are the level places, and the identifications of their residue fields with the
$\mathbf F_{g_k}$ are compatible with the two isomorphisms of $G$. For the level place $v$ at $P$ put
$\psi_{k,v}(a)=\sigma_0\big(\psi'_P(-a)\big)$, transported to $\mathbf F_{g_k}$. It is a character, and
it is nontrivial because $\psi'_P$ is and $\sigma_0$ is injective. Then
$\sigma_0(\psi_0(u))^{-1}=\prod_P\sigma_0\big(\psi'_P(-\sum_iu_{P,i,i+1})\big)$, which is the formula.
An element upper unitriangular at every level place is so at those over $g_0$, so $P_0$ is
half-unipotent.
% Lean: isWhit_and_eq_embLL_inv_and_halfUnip from the field psi of the point conditions (LvlTie: gen,
% p0, chi0, comp; mem_iff_of_lvlTie and chi_of_lvlTie with edgeSum supEdge, the superdiagonal sum) and
% p.eGtie; lam is embLL o psi0^-1, and HalfUnip holds with k = 0.
\end{proof}

\begin{lemma}\label{is-onto-step-bad-weight-regular-determinant}
Let $K$ be a field, $\Gamma$ a finite group acting on a commutative $K$-algebra $R$ by $K$-algebra
automorphisms $g\mapsto\alpha_g$, $H\subset\Gamma$ a subgroup, and $u\in R$ such that the determinant of
the matrix $\big(\alpha_{hg}(u)\big)_{h,g\in H}$ is a nonzerodivisor of $R$. Let $V,V_1,V_2\subset R$ be
$K$-subspaces with $V_2$ finite-dimensional, such that $V$ is stable under every $\alpha_g$, $V_1$ is
stable under every $\alpha_h$ with $h\in H$, and $x\,\alpha_h(u)\in V_1$ for $x\in V$ and
$x\,\alpha_h(u)\in V_2$ for $x\in V_1$, for every $h\in H$. Then there is an integer $n\geq0$ with
$n|H|\le\dim V_2$ such that $V$, with the action of $H$ by the $\alpha_h$, embeds $H$-equivariantly into
the sum of $n$ copies of the left regular representation $K[H]$.
\end{lemma}
\leanref{Development/AbelianVarietyTheory/Proofs/DivisorProofs/EmbedsInRegularRepresentationCopiesOfGroupDetermStatement.lean}{GlobalTheory.T.embedsInRegularRepresentationCopies_of_groupDeterminant}
\leanref{Development/AbelianVarietyTheory/Proofs/DivisorProofs/EmbedsInRegularRepresentationCopiesOfGroupDeterm.lean}{GlobalTheory.BadCount.embReg_of_groupDet}
\begin{proof}
A square matrix over $R$ whose determinant is a nonzerodivisor kills no nonzero vector, as the
adjugate shows. Let $I$ be the space of the $H$-invariant elements of $V_1$, and
$\operatorname{tr}(x)=\sum_{g\in H}\alpha_g(x)$, which maps $V_1$ to $I$ and satisfies
$\operatorname{tr}\circ\alpha_h=\operatorname{tr}$ for $h\in H$.

\emph{The dimension of $I$.} On the space of functions $F\colon H\to I$, the map
$\Psi\colon F\mapsto\sum_{h\in H}F(h)\,\alpha_h(u)$ to $V_2$ is injective: if $\sum_hF(h)\alpha_h(u)=0$, applying $\alpha_g$ for $g\in H$ and using $\alpha_g(F(h))=F(h)$
gives $\sum_h\alpha_{gh}(u)F(h)=0$ for every $g$, so the vector $(F(h))_h$ is killed by the matrix
$(\alpha_{gh}(u))_{g,h}$, and $F=0$. Hence $|H|\dim I\le\dim V_2$; put $n=\dim I$, and fix a basis of
$I$.

\emph{The embedding.} For $h\in H$ let $T_h\colon V\to I$ be $v\mapsto\operatorname{tr}\big(v\,\alpha_h(u)\big)$,
and let $j(v)\in K[H]^n$ have $i$-th coordinate $\sum_{h\in H}b_i(T_hv)\,h$, where $b_i$ is the $i$-th
coordinate in the basis of $I$. If $j(v)=0$, then $T_hv=0$ for every $h$, that is,
$\sum_g\alpha_g(v)\alpha_{gh}(u)=0$ for every $h\in H$; so the vector $(\alpha_g(v))_g$ is killed by the
transpose of $(\alpha_{gh}(u))_{g,h}$, and $v=\alpha_1(v)=0$. For $h,k\in H$,
$T_k(\alpha_hv)=\operatorname{tr}\big(\alpha_h(v\,\alpha_{h^{-1}k}(u))\big)=T_{h^{-1}k}(v)$, so the
coefficient of $k$ in $j(\alpha_hv)$ is the coefficient of $h^{-1}k$ in $j(v)$, which is that of $k$ in
$h\,j(v)$. So $j$ is an $H$-equivariant injection.
% Lean: BadCount.embReg_of_groupDet with RegEmbed.invIn (I), trL (tr), trL_mem, eq_zero_of_sum (the
% injectivity of Psi), eq_zero_of_traces (the injectivity of j), sum_act_act and act_sum_act, and
% GlobalTheory.mulVec_eq_zero_of_det_nonZeroDivisor.
\end{proof}

\begin{lemma}\label{is-onto-step-bad-weight-regular-cocycle}
Let $M$ be two-level data over $E$ for a group $U$ acting on a set $C$
(Definition~\ref{is-onto-share-two-level}), $\kappa$ a weight and $\rho$ a representation of $U$ on
$S_\kappa(M)$ pulled back from the action on $C$. Suppose that every $f\in S_\kappa(M)$ satisfies
$f(c,\gamma x)=f(c,x)$ for all $c\in C$, $\gamma\in\Gamma_c$ and $x\in\mathcal D_M$. Then there are
$a(u,c)\in\mathbf C^\times$ with $a(hg,c)=a(g,c)\,a(h,gc)$ for all $g,h\in U$ and $c\in C$, such that
$\big(\rho(h^{-1})f\big)(c,x)=a(h,c)\,f\big(hc,t(h,c)x\big)$ for all $h\in U$, $f\in S_\kappa(M)$, $c\in C$
and $x\in\mathcal D_M$.
\end{lemma}
\leanref{Development/AbelianVarietyTheory/Proofs/DivisorProofs/ExistsCocyclePullbackRepresentationStatement.lean}{GlobalTheory.T.exists_cocycle_pullbackRepresentation}
\leanref{Development/AbelianVarietyTheory/Proofs/DivisorProofs/ExistsCocyclePullbackRepresentation.lean}{GlobalTheory.PullCocycle.exists_cocycle}
\begin{proof}
Let $a(h,c)$ be the scalars for which $\rho$ is pulled back. Both sides of the formula are linear in $f$,
so it holds for every $f\in S_\kappa(M)$, the span of the restrictions of the families. For
$g,h\in U$, $c\in C$, $x\in\mathcal D_M$ and $f\in S_\kappa(M)$, the element
$t(hg,c)\big(t(h,gc)t(g,c)\big)^{-1}$ lies in $\Gamma_{hgc}$, so the invariance of $f$ gives
$f\big(hgc,t(hg,c)x\big)=f\big(hgc,t(h,gc)t(g,c)x\big)$. Comparing
$\rho((hg)^{-1})=\rho(g^{-1})\rho(h^{-1})$ on $f$ at $(c,x)$ then gives
\[
 a(hg,c)\,f\big(hgc,t(hg,c)x\big)=a(g,c)\,a(h,gc)\,f\big(hgc,t(hg,c)x\big).
\]
Let $C_0$ be the set of $c$ such that some $f\in S_\kappa(M)$ is nonzero at some point of
$\{c\}\times\mathcal D_M$. By the formula and the invertibility of $\rho$, $C_0$ is stable under $U$. For
$c\in C_0$ choose $f'$ and $x$ with $f'(c,x)\ne0$; then $f=\rho(hg)f'$ satisfies
$a(hg,c)f\big(hgc,t(hg,c)x\big)=f'(c,x)\ne0$, and the displayed identity gives
$a(hg,c)=a(g,c)a(h,gc)$. Replace $a(h,c)$ by $1$ for $c\notin C_0$. The cocycle identity then holds for
every $c$, as $C\smallsetminus C_0$ is stable under $U$, and the formula still holds at $c\notin C_0$,
where both sides vanish since $hc\notin C_0$.
% Lean: PullCocycle.exists_cocycle: formula_span (the span), pull_mul_eq (the invariance under
% t(hg,c)(t(h,gc)t(g,c))^-1, from TwoLevels.tr_mul), C0, hstab and hstab', hcoc, and the scalars
% if C0 c then a h c else 1.
\end{proof}

\begin{lemma}\label{is-onto-step-bad-weight-regular-one-lift}
Let $M$ be two-level data of rank $r$ over $E$ for a finite group $U$ acting on a finite set $C$, for
which the level hypothesis holds, and let $s$ be an even integer with $s\geq2r+50r^2\dim X(M)+1$. Let
$c\in C$ and $x\in\mathcal D_M$ be such that there is no contraction at $x$ for $X_c$
(Definition~\ref{is-onto-share-flag-data}). Then there is a family $F$ of weight $\delta_s$ with
$F_c(x)\ne0$, $(h^*F)_c(x)=0$ for every $h\ne1$ in $U$, and $F_{c'}=0$ for every $c'$ outside the orbit
$Uc$.
\end{lemma}
\leanref{Development/AbelianVarietyTheory/Proofs/DivisorProofs/ExistsCuspFormNonzeroAtOneLiftStatement.lean}{GlobalTheory.T.exists_cuspForm_nonzero_at_one_lift}
\begin{proof}
Let $\Gamma^\circ_c$ be the set of products $\gamma\,t(v,c)$ with $v\in U$, $vc=c$ and $\gamma\in\Gamma_c$.
Lemma~\ref{is-onto-share-interpolation} gives a cusp form $P\in S_{\delta_s}(X_c)$ with $P(x)\ne0$ and
$P(\beta x)=0$ for every $\beta\in\Gamma^\circ_c$ outside $\Gamma_c$. Let $F$ be the family with $F_c=P$
and $F_{c'}=0$ for $c'\ne c$; it is a family of weight $\delta_s$, as $0$ is a cusp form of every weight
(Lemma~\ref{is-onto-share-sections}(1)).

Then $F_c(x)=P(x)\ne0$. A point $c'$ outside the orbit $Uc$ is not $c=1\cdot c$, so $F_{c'}=0$. Let
$h\in U$ with $h\ne1$. If $hc\ne c$, then $(h^*F)_c(x)=F_{hc}(t(h,c)x)=0$. If $hc=c$, then
$t(h,c)=1\cdot t(h,c)$ lies in $\Gamma^\circ_c$. It does not lie in $\Gamma_c$: otherwise $h=1$, by the
last condition on the $t(u,c)$ in Definition~\ref{is-onto-share-two-level}. Hence
$(h^*F)_c(x)=P(t(h,c)x)=0$.
% Lean: T.exists_cuspForm_nonzero_at_one_lift = LiftInterp.exists_form_one_lift. The group
% Gamma°_c is M.upper c (T.exists_upperLevel_subgroup_component), and the prime and the hypotheses of
% comp_interp come from T.exists_totallySplitPrime_upperLevel_unipotent and
% T.upperLevel_isUnitaryStabiliser_both_lattices, as in the proof of Lemma is-onto-share-interpolation;
% P is comp_interp (M.comp c) (M.chi_one c) G g. The family is u c' y = if c' = c then P y else 0, with
% zero_mem_SS off c; the vanishing of the translates is hPγ with the element <h, hc, 1, one_mem> of
% M.upper c and M.tr_faithful. The open steps of the indefinite case of Lemma is-onto-share-interpolation, below
% Lemmas is-onto-step-orbit-interpolation-large and is-onto-step-orbit-interpolation-small, are now written out.
\end{proof}

\begin{lemma}\label{is-onto-step-bad-weight-regular-translates}
Let $M$ be two-level data of rank $r$ over $E$ for a finite group $U$ acting on a finite set $C$, for
which the level hypothesis holds, and let $s$ be an even integer with $s\geq2r+50r^2\dim X(M)+1$. Then
there is a family $F$ of weight $\delta_s$ such that for every $c\in C$ there is $x\in\mathcal D_M$ with
\[
 \det\big(((hh')^*F)_c(x)\big)_{h,h'\in U}\ne0.
\]
\end{lemma}
\leanref{Development/AbelianVarietyTheory/Proofs/DivisorProofs/ExistsCuspFormIndependentTranslatesStatement.lean}{GlobalTheory.T.exists_cuspForm_independent_translates}
\leanref{Development/AbelianVarietyTheory/BlockTheory/PoincareCuspForms/TwoLevelLogarithm.lean}{GlobalTheory.IndefiniteRowSideFullHecke.FrameWeightGapBounds.IndependentTranslates}
\begin{proof}
The values of a family $G$ of weight $\lambda$ satisfy $G_c(\gamma x)=G_c(x)$ for $\gamma\in\Gamma_c$, as
$\chi_{c,\lambda}$ is trivial on $\Gamma_c$. With $u=v=1$ in the conditions on the $t(u,c)$,
$t(1,c)^{-1}\in\Gamma_c$, so $(1^*G)_c=G_c$; and $t(gk,c)\big(t(g,kc)t(k,c)\big)^{-1}\in\Gamma_{gkc}$
gives $(g^*G)_{kc}(t(k,c)x)=((gk)^*G)_c(x)$.

For each $c$, Lemma~\ref{is-onto-share-balanced} gives $x_c\in\mathcal D_M$ with no contraction for $X_c$,
and Lemma~\ref{is-onto-step-bad-weight-regular-one-lift} gives a family $F^{(c)}$ of weight $\delta_s$
with $F^{(c)}_c(x_c)\ne0$, $(h^*F^{(c)})_c(x_c)=0$ for $h\ne1$, and $F^{(c)}_{c'}=0$ off the orbit of
$c$. Choose a representative $c_q$ in each orbit $q$ and put $F=\sum_qF^{(c_q)}$, a family of weight
$\delta_s$, which agrees with $F^{(c_q)}$ on the orbit of $c_q$.

Let $c=kc_q$ with $k\in U$, and $x=t(k,c_q)x_{c_q}$. By the first paragraph, the entry at $(h,h')$ of the
matrix is $((hh'k)^*F^{(c_q)})_{c_q}(x_{c_q})$. It vanishes unless $hh'k=1$, and for $hh'k=1$ it is
$F^{(c_q)}_{c_q}(x_{c_q})\ne0$. So the matrix is this value times a permutation matrix, the permutation
being $h\mapsto h^{-1}k^{-1}$, and its determinant is $\pm F^{(c_q)}_{c_q}(x_{c_q})^{|U|}\ne0$.
% Lean: TranslatesDet.independentTranslates: T.existsBalancedFrameCMUnitary_proof (Lemma
% is-onto-share-balanced) for xb, T.exists_cuspForm_nonzero_at_one_lift for uc, the sum over the orbit
% representatives (ssSum, hsum, hrep), and detNeAt (ssLact, the entries, the permutation sigma and
% Matrix.det_permute).
\end{proof}

\begin{lemma}\label{is-onto-step-bad-weight-regular}
Let $M$ be two-level data of rank $r$ over $E$ for a finite group $U$ acting on a finite set $C$
(Definition~\ref{is-onto-share-two-level}), $\kappa$ a weight and $\rho$ a representation of $U$ on
$S_\kappa(M)$ pulled back from the action on $C$, and suppose that the cusp forms of $M$ have ranks
$(\varepsilon,c_0)$ at $\kappa$ and $s=s(M)$ for some real numbers $\varepsilon,c_0$. Let $H\subset U$
be a subgroup for which the level hypothesis holds. Then there is an integer $n\geq0$ with
$n|H|\le\dim S_{\kappa+2\delta_s}(M)$ such that $S_\kappa(M)$ embeds $H$-equivariantly into the sum of
$n$ copies of the left regular representation of $H$.
\end{lemma}
\leanref{Development/AbelianVarietyTheory/Proofs/DivisorProofs/TwoLevelCuspFormsEmbedRegularRepresentationCopieStatement.lean}{GlobalTheory.T.twoLevelCuspForms_embed_regularRepresentation_copies}
\leanref{Development/AbelianVarietyTheory/Definitions/Divisors/BadElt.lean}{GlobalTheory.BadCount.EmbReg}
\begin{proof}
\emph{An algebra of functions.} Let $A$ be the $\mathbf C$-algebra of the functions $f$ on
$C\times\mathcal D_M$ such that, for every $c$, $x\mapsto f(c,x)$ is holomorphic and $f(c,\gamma x)=f(c,x)$
for $\gamma\in\Gamma_c$. The restriction $(c,x)\mapsto G_c(x)$ of a family $G$ of any weight lies in $A$,
as $\chi_{c,\lambda}$ is trivial on $\Gamma_c$; so $S_\kappa(M)$ and $S_{\kappa+2\delta_s}(M)$ lie in $A$.
For $u\in U$ put $(\mathrm{pb}_uf)(c,x)=f\big(uc,t(u,c)x\big)$. This is an algebra endomorphism of $A$,
since $t(u,c)$ acts holomorphically and $t(u,c)\Gamma_ct(u,c)^{-1}=\Gamma_{uc}$; and
$\mathrm{pb}_1=1$ and $\mathrm{pb}_v\mathrm{pb}_u=\mathrm{pb}_{uv}$, by the invariance and the
conditions $t(1,c)^{-1}\in\Gamma_c$ and $t(uv,c)\big(t(u,vc)t(v,c)\big)^{-1}\in\Gamma_{uvc}$.

\emph{An action.} The elements of $S_\kappa(M)$ are invariant, so
Lemma~\ref{is-onto-step-bad-weight-regular-cocycle} gives scalars $a(u,c)$ with the cocycle identity and
$\big(\rho(g)f\big)(c,x)=a(g^{-1},c)\,(\mathrm{pb}_{g^{-1}}f)(c,x)$. The identity with $g=h=1$ gives
$a(1,c)=1$. On the polynomial ring $A[X]$ let $\tau_u$ be the $\mathbf C$-algebra endomorphism that is
$\mathrm{pb}_u$ on $A$ and sends $X$ to $a_uX$, where $a_u\in A$ is $(c,x)\mapsto a(u,c)$. Then
$\tau_1=1$, and $\tau_v\tau_u=\tau_{uv}$: on $A$ this is the law for $\mathrm{pb}$, and
$\tau_v\tau_u(X)=\mathrm{pb}_v(a_u)\,a_vX$ with $a(u,vc)a(v,c)=a(uv,c)$. So $\alpha_g=\tau_{g^{-1}}$ is an
action of $U$ on $A[X]$ by algebra automorphisms, and $\alpha_g(fX)=(\rho(g)f)X$ for $f\in S_\kappa(M)$.

\emph{The element $u$.} The number $s=s(M)=2r+50r^2\dim X(M)+2$ is even, and the level hypothesis holds
for $M|_H$, with $\dim X(M|_H)=\dim X(M)$. Lemma~\ref{is-onto-step-bad-weight-regular-translates} for
$M|_H$ gives a family $F$ of weight $\delta_s$ such that for every $c$ there is $x$ with
$\det\big(((hh')^*F)_c(x)\big)_{h,h'\in H}\ne0$. Let $w\in A$ be its restriction and $u=w\in A\subset A[X]$.
Then $\alpha_{hg}(u)=\mathrm{pb}_{g^{-1}h^{-1}}w$, whose value at $(c,x)$ is
$((g^{-1}h^{-1})^*F)_c(x)$; replacing $h,g$ by $h^{-1},g^{-1}$ and transposing does not change the
determinant, so $D=\det\big(\alpha_{hg}(u)\big)_{h,g\in H}\in A$ is not identically zero on
$\{c\}\times\mathcal D_M$, for every $c$. It is a nonzerodivisor of $A[X]$: if $pD=0$, each coefficient
$p_i$ of $p$ has $p_i(c,\cdot)D(c,\cdot)=0$ on $\mathcal D_M$, which is open and connected
(Lemma~\ref{is-onto-share-sections}(2)), so $p_i(c,\cdot)=0$ by the identity theorem.

\emph{The spaces.} Let $V=\{fX:f\in S_\kappa(M)\}$, $V_1$ the span of the $v\,\alpha_h(u)$ with $v\in V$
and $h\in H$, and $V_2=\{fX:f\in S_{\kappa+2\delta_s}(M)\}$, finite-dimensional since the cusp forms have
ranks $(\varepsilon,c_0)$. The space $V$ is stable under the $\alpha_g$, and so $V_1$ under the
$\alpha_h$, $h\in H$, as $\alpha_h(v\,\alpha_k(u))=\alpha_h(v)\,\alpha_{hk}(u)$. For $f\in S_\kappa(M)$ and
$h,k\in H$, $fX\,\alpha_k(u)\,\alpha_h(u)=f\,\mathrm{pb}_{k^{-1}}(w)\,\mathrm{pb}_{h^{-1}}(w)X$. If $f$ is
the restriction of a family $G$ of weight $\kappa$, this is the restriction of the family with components
$G_c\,((k^{-1})^*F)_c\,((h^{-1})^*F)_c$, which lie in $S_{\kappa+2\delta_s}(X_c)$ by
Lemma~\ref{is-onto-share-pullback} and Lemma~\ref{is-onto-share-sections}(1) applied twice; by linearity,
$V_1\alpha_h(u)\subset V_2$.

Lemma~\ref{is-onto-step-bad-weight-regular-determinant}, with $K=\mathbf C$, $\Gamma=U$ and $R=A[X]$,
gives $n$ with $n|H|\le\dim V_2=\dim S_{\kappa+2\delta_s}(M)$ and an $H$-equivariant embedding of $V$
into $\mathbf C[H]^n$. As $f\mapsto fX$ is an isomorphism from $S_\kappa(M)$ to $V$ that carries $\rho$ to
the $\alpha_g$, this is the claim.
% Lean: CuspEmb.cusp_embReg: HA (the algebra A, with IsHol and IsInv), pb and pbA, pbA_one and pbA_comp;
% T.twoLevelCuspForms_restrict_holInvariant (res_mem_HA); T.exists_cocycle_pullbackRepresentation; tw and
% T.twistedPullbackPolynomial_rightAction_of_cocycle (tw_law) for the action act k = tw a(k^-1) k^-1;
% T.exists_cuspForm_independent_translates for sub M H with sOf M (even_sOf, le_sOf); the determinant
% hnz with T.polynomialC_holInvariant_mem_nonZeroDivisors (C_mem_nonZeroDivisors, through
% flag_domain_frames_connected and the identity theorem eq_zero_of_mul_eq_zero); deg1 for V and V2, V1,
% hV1s, hV2 with mul_pb_mem_cusp (T.twoLevelCuspForms_mul_mem_add_weight and
% T.twoLevelCuspForms_pull_mem); T.embedsInRegularRepresentationCopies_of_groupDeterminant; hR.2.1 for
% the finite dimension. The remaining open step is Lemma is-onto-step-bad-weight-regular-one-lift.
\end{proof}

\begin{lemma}\label{is-onto-step-bad-weight-model}
Let $M$ be two-level data of rank $r$ over $E$ for a finite group $U$ acting on a finite set $C$,
$\kappa$ a weight, $\rho$ a representation of $U$ on $S_\kappa(M)$ pulled back from the action on $C$,
and $T$ an integer. Suppose that either $p_0\in\{0,r\}$, or $0<p_0<r$, $M$ satisfies the archimedean
condition at $T$ for $\kappa$, $s(M)\le T$ and the level hypothesis holds for the trivial subgroup. Let
$H,P\subset U$ be subgroups for which the level hypothesis holds and $\lambda$ a character of $P$.
Then there is an integer $n\geq0$ such that $S_\kappa(M)$ embeds $H$-equivariantly into the sum of $n$
copies of the left regular representation of $H$, and $n|H|\le2|P|\dim S_\kappa(M)^{P,\lambda}$.
\end{lemma}
\leanref{Development/AbelianVarietyTheory/Proofs/DivisorProofs/BadWStOfRegular.lean}{GlobalTheory.BadWOn.model_bound}
\begin{proof}
If $p_0\in\{0,r\}$, Lemma~\ref{is-onto-share-definite} for $H$ gives $n$ with $n|H|=\dim S_\kappa(M)$
and the embedding, and for $P$ it gives $|P|\dim S_\kappa(M)^{P,\lambda}=\dim S_\kappa(M)$. So
$n|H|=|P|\dim S_\kappa(M)^{P,\lambda}\le2|P|\dim S_\kappa(M)^{P,\lambda}$.

Otherwise put $s=s(M)$. By Lemma~\ref{is-onto-share-ranks} the cusp forms of $M$ have ranks
$(\varepsilon,2)$ at $\kappa$ and $s$, with $\varepsilon=\dim Y(M)(\dim Y(M)+6)/T$; so
$\dim S_{\kappa+2\delta_s}(M)\le2\dim S_{\kappa-\delta_s}(M)$. Lemma~\ref{is-onto-step-bad-weight-regular}
gives $n$ with the embedding and $n|H|\le\dim S_{\kappa+2\delta_s}(M)$, and
Lemma~\ref{is-onto-share-indefinite} for $P$ and $\lambda$ gives
$\dim S_{\kappa-\delta_s}(M)\le|P|\dim S_\kappa(M)^{P,\lambda}$. Together,
$n|H|\le2|P|\dim S_\kappa(M)^{P,\lambda}$.
% Lean: BadWOn.model_bound: T.definite_cuspForms_finite_regularRepresentation in the definite case;
% T.twoLevelFlagDomain_cuspForm_ranks_shifted_weights (hR, whose clause hR.2.2.2.1 is the bound with
% c0 = 2), T.twoLevelCuspForms_embed_regularRepresentation_copies and the first clause of
% T.twoLevelFlagDomain_eigenspace_share otherwise.
\end{proof}

\begin{lemma}\label{is-onto-step-bad-weight-transfer}
Let $K$ be a number field, $\Gamma$ a finite group with a representation $\rho_K$ on $K^N$, and
$\sigma,\sigma'\colon K\to\mathbf C$ embeddings. Let $V$ and $V'$ be complex representations of
$\Gamma$, and $f\colon K^N\to V$ and $f'\colon K^N\to V'$ maps that are additive,
$\sigma$-semilinear and $\sigma'$-semilinear respectively, carry the standard basis to a basis, and
satisfy $gf(x)=f(\rho_K(g)x)$ and $gf'(x)=f'(\rho_K(g)x)$. Let $H\subset\Gamma$ be a subgroup and $n\geq0$.
If $V$ embeds $H$-equivariantly into the sum of $n$ copies of the left regular representation of $H$,
then so does $V'$.
\end{lemma}
\leanref{Development/AbelianVarietyTheory/Proofs/DivisorProofs/EmbedsInRegularRepresentationCopiesComplexEmbeddStatement.lean}{GlobalTheory.T.embedsInRegularRepresentationCopies_complexEmbedding_indep}
\leanref{Development/AbelianVarietyTheory/Proofs/DivisorProofs/EmbedsInRegularRepresentationCopiesComplexEmbedd.lean}{GlobalTheory.BadCount.embReg_of_baseChange}
\begin{proof}
As $K$ is countable, there is a field automorphism $\tau$ of $\mathbf C$ with $\tau\circ\sigma'=\sigma$:
$\mathbf C$ has the same uncountable transcendence degree over the countable fields $\sigma'(K)$ and
$\sigma(K)$ and is algebraically closed, so the isomorphism $\sigma\circ\sigma'^{-1}$ between them
extends to $\mathbf C$. Let $b'_k=f'(e_k)$, the basis of $V'$, and define $T\colon V'\to V$ by
$T(\sum_kz_kb'_k)=\sum_k\tau(z_k)f(e_k)$. It is additive, $\tau$-semilinear and bijective, as the
$f(e_k)$ form a basis of $V$, and $T(f'(x))=\sum_k\tau(\sigma'(x_k))f(e_k)=f(x)$. So $T$ carries the
action of $g$ on $f'(x)$ to its action on $f(x)$, and, as both sides are $\tau$-semilinear in the
vector and the $f'(x)$ span $V'$ over $\mathbf C$, $T(gy)=gT(y)$ for all $y\in V'$. Let
$j\colon V\to\mathbf C[H]^n$ be an injective $H$-equivariant linear map. Then
$y\mapsto\tau^{-1}\big(j(T(y))\big)$, with $\tau^{-1}$ applied to the coefficients in $\mathbf C[H]$, is
injective, $\mathbf C$-linear and $H$-equivariant.
% Lean: embReg_of_baseChange: CompRat.exists_ringEquiv_complex_comp_eq for tau, the map T built from
% the basis b', and embReg_of_semilinear.
\end{proof}

\begin{lemma}\label{is-onto-step-bad-weight-share}
Assume Green's description of the cuspidal characters (Theorem~\ref{lit-green-cuspidal}). Let $D$ be a
definite-side datum with ranks $f,m$, rank $r\geq100$, level primes $g_0,g_1\geq2^{12}$ at which its
types are admissible, and host primes $q,d$, whose series idempotents are tied to the labels of
Definition~\ref{cv-separated}, as the recursion ties of a point give; the admissibility of the types is
not used in the proof. Let $\psi_0$ and $\psi_1$ be nontrivial additive characters of
$\mathbf F_{g_0}$ and $\mathbf F_{g_1}$, let $k\in\{0,1\}$, and let $S$ be a set of tuples of labels at
$g_k$ none of which is separated. Then the share $s_k(S)$ (Definition~\ref{cv-separated}(4)) satisfies
\[
 s_k(S)\le\frac{3f}{g_k}+\frac{23f\log g_k}{\sqrt{g_k}}
 +2fm\Big(\frac1{4\min(q,d)}+\frac{1+\log r}{g_k-1}\Big)
 +\frac r2\Big(\frac78\Big)^{f}+e^{-f/(5\log r)}.
\]
\end{lemma}
\leanref{Development/AbelianVarietyTheory/Proofs/DivisorProofs/WeightedIdempotentLocalShareBadTuplesLeStatement.lean}{GlobalTheory.T.weightedIdempotent_localShare_bad_tuples_le}
\leanref{Development/AbelianVarietyTheory/Proofs/InertiaEigenvalues/LemmaLSeriesWeightsOfGreenCuspidalParametrisatioStatement.lean}{GlobalTheory.T.admLq_of_green_of_eSerTie}
\begin{proof}
Write $g=g_k$, $w(h)=\operatorname{wt}_{\psi_k}(h)$ for a label $h$, $I$ for the set of the $f$ level
places over $g$, and $\mathcal S_r$ for the set of monic squarefree polynomials of degree $r$ over
$\mathbf F_g$ with nonzero constant term, the squarefree labels.

\emph{The weights.} By Step 1 of the proof of Lemma~\ref{cv-local-share}, which uses Green's description
and the tie of the series idempotents to the labels, $w(h)\geq0$ for every label $h$, and
$w(h)=\prod_{P\mid h}(g^{\deg P}-1)^{-1}$ for $h\in\mathcal S_r$, the product over the irreducible
factors; on $\mathcal S_r$, $w$ is the weight of Lemma~\ref{cv-squarefree-mass}. The series idempotents
of the labels sum to $1$ (the remark after Definition~\ref{cv-separated}), so
$\sum_h\big(e_U(\psi_k)e(h)\big)(1)=e_U(\psi_k)(1)=|U|^{-1}\ne0$, and the weights of the labels sum to
$1$. By Lemma~\ref{cv-squarefree-mass}, the labels in $\mathcal S_r$
have total weight at least $1-2/g$, so those outside $\mathcal S_r$ have total weight at most $2/g$; and
those in $\mathcal S_r$ have total weight at most $1$.

\emph{Splitting the tuples.} As the weights are nonnegative, for $t\in S$ the product
$\prod_vw(t_v)$ is at most the sum, over the $v\in I$ with $t_v\notin\mathcal S_r$, of $\prod_{v'}w(t_{v'})$,
plus $\prod_vw(t_v)$ if every $t_v$ lies in $\mathcal S_r$. For a fixed $v$, the sum of
$\prod_{v'}w(t_{v'})$ over all tuples of labels with $t_v\notin\mathcal S_r$ is at most $2/g$, as the
weights at the other places sum to $1$. So the first parts contribute at most $2f/g$ to $s_k(S)$.

\emph{The squarefree tuples.} It remains to bound the sum of $\prod_vw(t_v)$ over the $t\in S$ with every
$t_v\in\mathcal S_r$. Such a $t$ satisfies (a)
of Definition~\ref{cv-separated}(3) and is not separated, so it fails (b), (c), (d) or (e). The weight of
$\mathcal S_r$ is at most $1$, so the families failing a condition at one given place have weight at most
the weight of the polynomials failing it there, and the families failing a condition at every place have
weight at most the product of these weights. The datum gives the hypotheses of
Lemma~\ref{cv-local-estimates}(2) for $q_1=q$: as in Step 2 of the proof of Lemma~\ref{cv-local-share},
$q$ and $d$ are distinct odd primes greater than $4tr+1\geq16$, and $16qd$ divides $g-1$. With $g\geq2^{12}$
and $r\geq100$, Lemma~\ref{cv-local-estimates} gives:
\begin{itemize}
\item failures of (b) at some $v$: at most $f\cdot23\log g/\sqrt g$, by part (1);
\item failures of (c) at some $v$, for some $j\le m$ and for $\chi_{k,v,j}$ or its inverse: at most
  $2fm\big(\frac1{4\min(q,d)}+\frac{1+\log r}{g-1}\big)$, by part (2);
\item failures of (d), that is, for some $a$ with $1\le2a\le r$ every $t_v$ has a divisor of degree $a$:
  at most $\frac r2(\frac78)^f$, by part (3) with $0<a<r$;
\item failures of (e), that is, no $t_v$ has an irreducible factor of prime degree $p'$ with $r<2p'$ and
  $p'+3\le r$: at most $(1-\frac1{5\log r})^f\le e^{-f/(5\log r)}$, by part (4).
\end{itemize}

Adding up, $s_k(S)$ is at most
$\frac{2f}g+\frac{23f\log g}{\sqrt g}+2fm\big(\frac1{4\min(q,d)}+\frac{1+\log r}{g-1}\big)+\frac r2(\frac78)^f
+e^{-f/(5\log r)}$, which is at most the stated bound.
% Lean: T.admLq_of_green_of_eSerTie (= GLq.lemmaL_series_weights_of_greenCuspidalParametrisation, Step 1
% of the proof of Lemma cv-local-share) for hA0 and hA1, and T.weightedIdempotent_localShare_bad_tuples_le =
% PropWU.propW_eUk: htot from D.ntypes.eSer_sum and eU_coeff_one_ne, hoff (wtOf_off), hnsq from
% T.one_sub_two_div_le_squarefreeLabelMass, hsplit, hA (sum_pi_single), hB, and
% tupleMass_not_good_le with host_primes: the events B2 to B5 bounded by
% T.squarefreeLabelTorusMass_rootsNotDistinct_le, T.mass_not_rootsAvoid_le',
% T.squarefreeLabelTorusMass_hasDivisorOfDegree_le and mass_not_hasPrimeFactor_le_five, with
% T.squarefreeLabelTorusMass_divisor_bound (the mass of S_r is at most 1). The hypothesis TieAdmG is the
% unused argument _hT.
\end{proof}

\begin{remark}\label{is-onto-step-bad-weight-share-sharper}
The proof of Lemma~\ref{is-onto-step-bad-weight-share} gives the bound with $2f/g_k$ in place of
$3f/g_k$: the tuples with a component that is not squarefree contribute at most $2f/g_k$, and the other
four terms are unchanged.
\end{remark}

\begin{lemma}\label{is-onto-step-bad-weight}
Assume Green's description of the cuspidal characters of $GL_d(\mathbf F_g)$
(Theorem~\ref{lit-green-cuspidal}). Let $p$ be a point of a family that satisfies the point conditions
and whose recursion ties hold, at a member whose definite-side datum has rank $r\geq100$ and ranks
$f,m$, whose types are admissible at the primes $g_0$ and $g_1$ of the host, and whose host has the
primes $q,d$. Let $i$ be a lift and $T$ an integer with $2r+50r^2f\lfloor r^2/4\rfloor+2\le T$
such that $i$ has a model with parameter $T$ as in the conclusion of
Lemma~\ref{is-onto-step-models}. Let $\delta$ be a real number such that, for $k=0,1$, $g_k\geq2^{12}$
and
\[
 4\Big(\frac{3f}{g_k}+\frac{23f\log g_k}{\sqrt{g_k}}
 +2fm\Big(\frac1{4\min(q,d)}+\frac{1+\log r}{g_k-1}\Big)
 +\frac{r}2\Big(\frac78\Big)^{f}+e^{-f/(5\log r)}\Big)\le\delta.
\]
Then, with $FW_i^0$ the subspace of the coordinates of $FW_i$ on which $P_0$ acts through
$\chi_i\psi_0^{-1}$ and $e_{\mathrm{good}}$ the good operator (Definition~\ref{is-ties-class-sums}),
\[
 \dim\big(FW_i^{0}\cap\ker e_{\mathrm{good}}\big)\le\delta\dim FW_i^{0}.
\]
\end{lemma}
\leanref{Development/AutomorphicArithmetic/ThetaCorrespondence/ThetaPairLifts/ThirdPlaceSeries/BadWStOf.lean}{GlobalTheory.BadWOn.badW_of_del_bound_of_T_bound}
\leanref{Development/AbelianVarietyTheory/Definitions/Divisors/ExistsArchOne.lean}{GlobalTheory.ShareOn.BadW}
\begin{proof}
Let $(C,M,\kappa,\rho,\sigma,f)$ be the model of the lift $i$ with parameter $T$, so that $M$ has rank
$r$ and $|\Psi|=f$. Write $G$ for the group of the realization, $N=n_W(i)$, and $\chi_i$ and
$\varrho_i$ for the character and the action of $G$ on the coordinates $L^N$ of $FW_i$; put
$\rho_K(g)=\chi_i(g)^{-1}\varrho_i(g)$, a representation of $G$ on $L^N$.

\emph{Subgroups.} Identify $G$ with $\prod_{k,v}GL_{r}(\mathbf F_{g_k})$ as in
Lemma~\ref{is-onto-step-bad-weight-whittaker}, which gives characters $\psi_{k,v}$ for which $P_0$ is the
subgroup $P$ of Lemma~\ref{cv-bad-part} and $\sigma_0\circ\psi_0^{-1}$ is its character $\lambda$, and
shows that $P_0$ is half-unipotent. For $k=0,1$ let $H_k\subset G$ be the subgroup of the elements that
are upper unitriangular at the level places over $g_{1-k}$; it is half-unipotent, and so is the trivial
subgroup. By the definition of a model, the level hypothesis of $M$ holds for $P_0$, $H_0$, $H_1$ and
$\{1\}$.

\emph{The parameter.} For integers $0\le a\le b$, $4a(b-a)\le b^2$, so
$p_0(r-p_0)\le\lfloor r^2/4\rfloor$ and $\dim X(M)=fp_0(r-p_0)\le f\lfloor r^2/4\rfloor$. Hence
$s(M)=2r+50r^2\dim X(M)+2\le T$. So either $p_0\in\{0,r\}$, or $0<p_0<r$, $M$ satisfies the
archimedean condition at $T$ for $\kappa$, $s(M)\le T$ and the level hypothesis holds for $\{1\}$.

\emph{Regular embeddings.} Let $\rho'(g)=\sigma(\chi_i(g))^{-1}\rho(g)$ on $S_\kappa(M)$; it is again
pulled back from the action on $C$, with the scalars $a(u,c)$ multiplied by $\sigma(\chi_i(u))$, and
$\rho'(g)f(x)=f(\rho_K(g)x)$. Let $V=\mathbf C^N$ with the representation $\rho_V$ of $G$ for which the
map $x\mapsto\sigma_0(x)$ (coordinatewise) satisfies $\rho_V(g)\sigma_0(x)=\sigma_0(\rho_K(g)x)$. For
$k=0,1$, Lemma~\ref{is-onto-step-bad-weight-model}, with $H=H_k$, $P=P_0$ and the character
$\sigma\circ\psi_0^{-1}$, gives $n_k$ such that $\rho'|_{H_k}$ embeds into $n_k$ copies of the regular
representation of $H_k$ and $n_k|H_k|\le2|P_0|\dim S_\kappa(M)^{P_0,\sigma\circ\psi_0^{-1}}$.
Lemma~\ref{is-onto-step-bad-weight-transfer}, with $\rho_K$, $f$ and $x\mapsto\sigma_0(x)$, shows that
$\rho_V|_{H_k}$ embeds into $n_k$ copies of the regular representation of $H_k$.

\emph{Dimensions.} For $u\in P_0$ let $a_u=\rho_K(u)-\psi_0(u)^{-1}$, so that $FW_i^0=\bigcap_u\ker a_u$.
The maps $f$ and $x\mapsto\sigma_0(x)$ are semilinear, carry the standard basis to a basis, and carry
$a_u$ to $\rho'(u)-\sigma(\psi_0(u))^{-1}$ and to $\rho_V(u)-\lambda(u)$ respectively. As in the proof
of Lemma~\ref{cv-realization-linear}, the joint kernels of two such families have the same dimension, so
$\dim S_\kappa(M)^{P_0,\sigma\circ\psi_0^{-1}}=\dim FW_i^0=\dim V^\lambda$. By the ties of the point,
$\sigma_0$ carries the coefficient $c_{\mathcal X}(g)$ of $e_{\mathcal X}$ to the coefficient of
$e_{\mathrm{sep}}$ (Definition~\ref{cv-separated}) at the image of $g$. So $x\mapsto\sigma_0(x)$
carries $e_{\mathrm{good}}=\sum_gc_{\mathcal X}(g)\rho_K(g)$ to $\rho_V(e_{\mathrm{sep}})$, and adding
these to the families gives
$\dim(FW_i^0\cap\ker e_{\mathrm{good}})=\dim\big(V^\lambda\cap\ker\rho_V(e_{\mathrm{sep}})\big)$.

\emph{Bad part.} Lemma~\ref{cv-bad-part}, with $\Gamma=G$, $\pi$ and $\iota$ the identification and
its inverse, the representation $\rho_V$, $c=2$, the $n_k$ above, and $\psi_k(a)=e^{2\pi ia/g_k}$, gives
\[
 \dim\big(V^\lambda\cap\ker\rho_V(e_{\mathrm{sep}})\big)\le2\big(s_0(S_0)+s_1(S_1)\big)\dim V^\lambda,
\]
where $S_k$ is the set of tuples of labels at $g_k$ that are not separated. The recursion ties of $p$
tie the series idempotents of the datum to the labels, so Lemma~\ref{is-onto-step-bad-weight-share}
applies to each $S_k$, and with the hypothesis on $\delta$ it gives $4s_k(S_k)\le\delta$. Hence
\[
 \dim\big(FW_i^0\cap\ker e_{\mathrm{good}}\big)\le2\Big(\frac\delta4+\frac\delta4\Big)\dim FW_i^0
 =\delta\dim FW_i^0.
\]
% Lean: BadWOn.badW_of_del_bound_of_T_bound, with GreenCuspidal.green: the model hmod (ModelAt);
% T.exists_whittaker_lam_psi_half_unip (Lemma is-onto-step-bad-weight-whittaker) and
% T.exists_whittakerDatum_and_halfLevels (the subgroups Hh k with IsHalf, which are the H_k); hsOf;
% model_bound for rhoM = twist rho chiM (isPullRep_twist) and P' = P0;
% T.embedsInRegularRepresentationCopies_complexEmbedding_indep for rhoE = repE embLL rhoK;
% T.finrank_iInf_ker_semilinear_standardBasis for dM, dE and dBad, with hsp (spType as a joint kernel)
% and hE0 (the field hid of p.grpRecPin); T.bad_part_dim_le_sum_local_shares_whittaker (Lemma
% cv-bad-part) with psi0 k = ZMod.stdAddChar; T.admLq_of_green_of_eSerTie with hRec.eSer and
% T.weightedIdempotent_localShare_bad_tuples_le (Lemma is-onto-step-bad-weight-share).
\end{proof}

\begin{lemma}\label{is-onto-step-bad-numerics-bound}
For real numbers $g,f,m,N,s$ put
\[
 B(g,f,m,N,s)=\frac{3f}g+\frac{23f\log g}{\sqrt g}+2fm\Big(\frac1{4s}+\frac{1+\log N}{g-1}\Big)
 +\frac N2\Big(\frac78\Big)^{f}+e^{-f/(10\log N)}.
\]
Let $0<\varepsilon\le1$ and $C$ be real numbers. There is $x_0$ such that for every $x\geq x_0$ and all
real numbers $g,f,m,N,s$ with $e^{8x^{\varepsilon}}\le g$, $x^2\le f\le x^C$, $0\le m\le x^C$,
$x^{\varepsilon}\le N\le x^C$ and $e^{x^{\varepsilon/6}}\le s$, we have $N\geq10^5$,
$B(g,f,m,N,s)\geq0$ and $6B(g,f,m,N,s)\le e^{-x^{\varepsilon/7}}$.
\end{lemma}
\leanref{Development/AbelianVarietyTheory/Proofs/FoundationProofs/LogDivSqrtLeFourExp.lean}{GlobalTheory.BadShare.badBound_numbers}
\leanref{Development/AbelianVarietyTheory/Definitions/Foundations/SamePartIffOfNeZero.lean}{GlobalTheory.BadShare.badBound}
\begin{proof}
This is a computation. With $P=x^C$ and $E=e^{-x^{\varepsilon/6}}$, for $x$ large each term is at most
a multiple of $P^3E$: $1/g\le E$ and $1/(4s)\le E/4$; $\log g/\sqrt g\le4e^{-2x^\varepsilon}\le4E$, since
$\log g\le4g^{1/4}$; $(1+\log N)/(g-1)\le2NE$; $(7/8)^f\le e^{-f/8}\le e^{-x^2/8}\le E$; and
$e^{-f/(10\log N)}\le E$, since $10\log N\le20C\sqrt x$ and $f\geq x^2$. So $6B\le612P^3E$, which is at
most $e^{-x^{\varepsilon/7}}$ for $x$ large. The bound $N\geq10^5$ is $x^\varepsilon\geq10^5$.
\end{proof}

\begin{lemma}\label{is-onto-step-bad-numerics}
Let $\mathcal H$ be the class of standard families. Eventually along $\mathcal H$, every family has a
scale $x_0$ such that every member at a scale $x\geq x_0$, with exponent $\varepsilon_z$ of its scale
and integer $T_O$ of its parameter choices, has rank $r\geq100$ and types admissible at $g_0$ and
$g_1$, satisfies $\delta=\exp(-x^{\varepsilon_z/7})\le\frac14$ and, for $k=0,1$, $g_k\geq2^{12}$ and
the inequality of Lemma~\ref{is-onto-step-bad-weight} with this $\delta$; and has
$2r+50r^2f\lfloor r^2/4\rfloor+2\le T_O$.
\end{lemma}
\leanref{Development/AbelianVarietyTheory/Proofs/DivisorProofs/BadShareSizeBoundsAlongFamilyStatement.lean}{GlobalTheory.T.badShareSize_bounds_along_family}
\leanref{Development/AbelianVarietyTheory/Proofs/DivisorProofs/HodgeLineMultipleLeFamilyThresholdStatement.lean}{GlobalTheory.T.hodgeLineMultiple_le_familyThreshold}
\begin{proof}
The conclusion holds for every family of exponent $\varepsilon\le1$, with any largeness threshold. Let
$\mathcal F$ be such a family; its members share the exponent $\varepsilon_z$ of their scales, and
$0<\varepsilon_z\le\varepsilon$. Write $k,m,f$ for the ranks and $r=k+m-1$.

\emph{Size facts.} At a member at a scale $x\geq16$ above the threshold of $\mathcal F$, the size
conditions give, for $g\in\{g_0,g_1\}$,
\[
 e^{8x^{\varepsilon}}\le g,\quad x^2\le f\le x^5,\quad m\le x^5,\quad100\le r\le x^5,\quad
 e^{x^{\varepsilon_z/6}}\le\min(q,d).
\]
Indeed, $(kb)^8\le g$ and $(kb)^8\geq b^8=e^{8x^\varepsilon}$. Next,
$x^2\le x^{2+8\varepsilon-\varepsilon/100}\le n\le f$. As $k\le x^{9/2+10\varepsilon+\varepsilon/100}\le
x^{19/4}$, $f\le k$, $m\le k$ and $r\le2k$, each of $f,m,r$ is at most $2k\le x^{1/4}x^{19/4}=x^5$.
Since $f\geq x^2\geq256$, $f\le k$ and $m\geq2$, $r\geq k+1>100$. Finally $\log q\geq
x^{\varepsilon_z/6}$, and $\log d\geq x^{\varepsilon_z/5-\varepsilon_z/100}\geq x^{\varepsilon_z/6}$.

Let $x_b$ be the number of Lemma~\ref{is-onto-step-bad-numerics-bound} for the exponent
$\varepsilon_z$ and $C=5$, and let $x_0\geq\max(16,x_b)$ be such that $x^{\varepsilon_z/7}\geq4$ and
$x^{\varepsilon}\geq512$ for $x\geq x_0$. Let a member at a scale $x\geq x_0$ above the threshold be
given, with $\delta=e^{-x^{\varepsilon_z/7}}$, and let $g\in\{g_0,g_1\}$.

\emph{The first conditions.} $e^{x^{\varepsilon_z/7}}\geq1+x^{\varepsilon_z/7}\geq4$, so
$\delta\le\frac14$. The rank satisfies $r\geq100$ by the size facts. The types of every definite-side
datum are admissible at $g_0$ and $g_1$: the conditions of admissibility are among the data of the
datum and its host. And $g\geq e^{8x^{\varepsilon}}\geq1+8\cdot512>2^{12}$.

\emph{The inequality.} Let $W$ be the expression in parentheses in the inequality of
Lemma~\ref{is-onto-step-bad-weight}, at $g$, so that the inequality is $4W\le\delta$. Put
$N=\max(r,x^{\varepsilon_z})$ and $s=\min(q,d)$. By the size facts, $e^{8x^{\varepsilon_z}}\le
e^{8x^{\varepsilon}}\le g$, $x^2\le f\le x^5$, $0\le m\le x^5$, $x^{\varepsilon_z}\le N\le x^5$ and
$e^{x^{\varepsilon_z/6}}\le s$. So Lemma~\ref{is-onto-step-bad-numerics-bound} gives $B(g,f,m,N,s)\geq0$
and $6B(g,f,m,N,s)\le\delta$. Since $g\geq2$, $m\geq0$ and $100\le r\le N$, we have $\log r\le\log
N$, $\frac{r}2(\frac78)^f\le\frac N2(\frac78)^f$ and $5\log r\le10\log N$; comparing term by term,
$W\le B(g,f,m,N,s)$. Hence $4W\le4B\le6B\le\delta$.

\emph{The bound on $T_O$.} The five terms of $W$ are nonnegative and $4W\le\delta\le\frac14$, so
$\frac{r}2(\frac78)^f\le\frac1{16}$. As $r\geq100$, this forces $f\geq5$, since
$50(\frac78)^4>\frac1{16}$. The parameter choices give $11D_p^2\le T_O$ with $D_p=f\binom{r}2$. For
integers $r\geq100$ and $f\geq5$, a computation gives
$2r+50r^2f\lfloor r^2/4\rfloor+2\le11\big(f\binom{r}2\big)^2$, so the bound holds.
% Lean: T.badShareSize_bounds_along_family = BadNumOnto.badShare_numbers_onto (eps1 = 1, threshold
% <0, 0, 0, 0>), with family_eventually_size_bounds (sizes_record_at with Cs = 5, at x >= max 16 (x1 + 1)),
% BadShare.badBound_numbers at ez and C = 5, tieAdmG_of_record (g_prime, g_split, g_Mray, g_cond) and
% wg_le_badBound. T.hodgeLineMultiple_le_familyThreshold = HodgeMultNum.hodgeMult_le_TO, through
% five_le_of_pow, ParamsT.TO_11D2 with Dp_eq and nat_bound.
\end{proof}

\subsection{Weyl exponents and the exponent dictionary}\label{subsection-is-weyl-exponents}

\begin{lemma}\label{is-onto-step-exp-dominant}
Let $D$ be a definite-side datum with ranks $k,m,t$ and $r=k+m-1$ (Definition~\ref{th-ranks}), let
$\tau$ be one of its embeddings, with Harish-Chandra parameter $(u,v,b,e)$
(Definition~\ref{th-definite-datum}(2)), and let $1\le J\le m$. Put $p=k+J-1$, $q=m-J$ and $n_V=r+t$, the
rank of $V_{J}$. For $1\le i\le r$ let $\lambda_i=v_i$ if $i\le k-1$ and $\lambda_i=-b_{i-k+1}$ if $i\geq k$,
and for $1\le i\le n_V$ put
\[
 \lambda'_i=\begin{cases}
 \lambda_i+i-\frac{n_V}2-1,&1\le i\le p,\\[2pt]
 \frac{p-q-1}2-e_J,&p<i\le p+t,\\[2pt]
 -b_{J+c}+c+\frac{t+p-q}2-1,&i=p+t+c,\ 1\le c\le q.
 \end{cases}
\]
Then every $\lambda'_i$ is an integer, and $\lambda'_{i+1}+\rho_{i+1}<\lambda'_i+\rho_i$ for $1\le i<n_V$,
where $\rho_i=(n_V+1-2i)/2$; in particular $\lambda'_1\geq\lambda'_2\geq\dots\geq\lambda'_{n_V}$. So
$(\lambda'_1,\dots,\lambda'_{n_V})$ is a dominant weight with integral entries: the weight $\lambda'_\tau$
that the Harish-Chandra parameters give at level $J$ (Definition~\ref{is-ties-lattice-systems}(4)).
\end{lemma}
\leanref{Development/AbelianVarietyTheory/Proofs/DivisorProofs/LambdaPrimeAllPoints.lean}{GlobalTheory.LiftWeightForced.isLamPrime_lamLdef}
\leanref{Development/AbelianVarietyTheory/Proofs/DivisorProofs/LambdaPrimeAllPoints.lean}{GlobalTheory.LiftWeightForced.exists_isLamPrime}
\leanref{Development/AbelianVarietyTheory/Proofs/IwasawaDecomposition/LamGeoEqFormula.lean}{GlobalTheory.LowestKType.thetaSide}
\leanref{Formalism/LocalToGlobal/IndefiniteReal/BridgeTie.lean}{Formalism.IndefR.Bridge.IsLamPrime}
\begin{proof}
\emph{The sequence $\lambda$.} By Definition~\ref{th-hc-parameter}, $\lambda_{i+1}<\lambda_i$ for $1\le i<r$:
if $i+1\le k-1$, then $v_{i+1}<u_{i+1}<v_i$; if $i+1=k$, then $-b_1<u_k<v_{k-1}$; and if $i\geq k$, then
$b_{c+1}-b_c=(b_{c+1}-e_c)+(e_c-b_c)\geq t+3>0$ for $c=i-k+1$. Also $b_J<e_J$, and $e_J<b_{J+1}$ if $J<m$, by
(4) of that definition. The $\lambda_i$ are integers, and $\lambda_{p+c}=-b_{J+c}$ for $0\le c\le q$.

\emph{Integrality.} As $k$ and $m$ are even and $t$ is odd, $n_V=k+m-1+t$, $p-q-1=k-m+2J-2$ and
$t+p-q=t+k-m+2J-1$ are even. So the three expressions for $\lambda'_i$ are integers.

\emph{Monotonicity.} Put $T_i=\lambda'_i+\rho_i+\frac12$. A direct computation from the three cases gives
\[
 T_i=\lambda_i\ (1\le i\le p),\qquad T_i=-e_J+\tfrac{t+1}2-(i-p)\ (p<i\le p+t),\qquad
 T_{p+t+c}=-b_{J+c}\ (1\le c\le q).
\]
We show $T_{i+1}+1\le T_i$ for $1\le i<n_V$. For $i<p$ this is $\lambda_{i+1}+1\le\lambda_i$, as the
$\lambda_i$ are integers decreasing strictly. For $i=p$ it is $-e_J+\frac{t-1}2+1\le-b_J$, which follows from
$e_J-b_J\geq\frac{t+3}2$. For $p<i<p+t$ it is an equality. For $i=p+t$, so that $q\geq1$, it is
$-b_{J+1}+1\le-e_J-\frac{t-1}2$, which follows from $b_{J+1}-e_J\geq\frac{t+3}2$. For $i=p+t+c$ with
$c\geq1$ it is $\lambda_{p+c+1}+1\le\lambda_{p+c}$. Since $\rho_{i+1}=\rho_i-1$, $T_{i+1}+1\le T_i$ gives
$\lambda'_{i+1}+\rho_{i+1}<\lambda'_i+\rho_i$ and $\lambda'_{i+1}\le\lambda'_i$.
% Lean: LiftWeightForced.exists_isLamPrime builds the WeightData wd D j tau (k, m, t, J = j + 1, e = e_J,
% v (a + 2) = v_(a+1), b (a + 1) = b_(a+1)); wd_hyp and wd_gap check WeightData.Hyp and GapHyp (the first
% paragraph), and GlobalTheory.LowestKType.thetaSide gives ThetaSideConcl: lamPrime = thetaWt + dTau, which
% is lambda' above (aPlus, aMinus with aMinus_rev), integral by aPlus_int, aMinus_int and dTau_int
% (exists_half_r, weightData_t_odd), and lamPrime_param with target_step (uSideTarget is T). lamLdef is the
% chosen dominant weight with IsLamPrime, which exists by exists_isLamPrime, so isLamPrime_lamLdef holds;
% IsLamPrime fixes the entries, so lamLdef is lambda'.
\end{proof}

\begin{lemma}\label{is-onto-step-exp-dictionary}
Let $\mathcal H$ be the class of standard families. Eventually along $\mathcal H$, every family has a
scale beyond which, for every point at a member and every lift $j$, there are a parameter record $P$
(Definition~\ref{is-onto-step-exp-bound-weights}) and weights $W=(\lambda_\tau)_\tau$, indexed by the embeddings $\tau$ of the host field, in the
\emph{exponent dictionary} of the size data at $j$, and satisfying its numerical conditions. Here the
dictionary requires that $P$ have the ranks $k,m,t,f$ and the scale $x$, $\log b=x^\varepsilon$ and
$b$ of the member, that $f$ be the number of embeddings, that its rank $N_r$ be the rank of $V_j$, that
every $\lambda_\tau$ differ by a constant from the dominant weight $\lambda'_\tau$ that the
Harish-Chandra parameters give at level $j+1$ (Definition~\ref{is-ties-lattice-systems}(4)), that the
numbers $\lambda_{\tau,a}+N_r-1-a$, $a<N_r$, lie in $[0,\lfloor h\rfloor]$ for the number $h$ of $P$,
that $W$ have at most $|T_2|$ thick embeddings and as its special set a set of primes determined by
the member, and that $h<g_0$ and $h<g_1$. The numerical conditions are hypotheses on $P$, conditions
on $W$ relative to $P$, and two bounds on the rows of $P$.
\end{lemma}
\leanref{Development/AutomorphicArithmetic/ThetaCorrespondence/ThetaPairLifts/ThirdPlaceSeries/NumSt.lean}{GlobalTheory.ExpTie.numSt_alt}
\leanref{Development/AbelianVarietyTheory/Definitions/Divisors/SumExpSt.lean}{GlobalTheory.ExpTie.Dict}
\leanref{Development/AbelianVarietyTheory/Definitions/Divisors/SumExpSt.lean}{GlobalTheory.ExpTie.NumOk}
\begin{proof}
The conclusion holds for every family of exponent $\varepsilon\le1$, with any largeness threshold,
beyond a scale that depends on $\varepsilon$ and the constant $C_t$ of the family. Fix a member at a
scale $x$ above the threshold, with ranks $k,m,t,n$, $f=2n$ and $r=k+m-1$, and a lift $j$. Take for
$P$ the \emph{parameter record} of the member at $j$: the ranks $k,m,t,f,n$, the rank $N_r$ of $V_j$,
$L=\log b=x^\varepsilon$, the scale $x$, the quantity $b'$, and the size $\eta=\eta_{sz}$ and the cloud
size $S=S_c$ of the size data. Its derived numbers are $b=e^L$, $w''=b'/(xn)$, $N=k+m(t+1)$,
\[
 \ell^*=\frac{mb}{nt^2(N+t)L^2+kn^2(S+t+L)}\quad\text{and}\quad h=b+S+\eta+t+1.
\]
By construction $P$ has the ranks, the scale, $\log b$ and $b$ of the member, $N_r$ is the rank of
$V_j$, and $f$ is the number of embeddings of the host field. The rest comes from the following.
\begin{enumerate}
\item[(a)] (Hypotheses on $P$.) $x\geq1$, $b'>0$, $1\le m\le k$, $n\geq1$, $f\geq1$, $k+m+4\le t$,
  $3m<t$, $\eta\le b/(t^3fk^3)$, $N_r+1=k+m+t$, $t\le\eta$, $S\geq0$,
  $b'+S+\eta+b/\ell^*<\ell^*$ and $2\eta+t+8b/\ell^*+7<\ell^*$. The bounds on the ranks are conditions
  of the ranks of a definite-side datum ($m\geq2$, $m+3\le k$, $r+5\le t$, $3m<t$, $n\geq1$), and
  $N_r=r+t$ gives $N_r+1=k+m+t$. Every size data have $x>1$, $b'>0$ and $S>0$. The bounds on
  $\eta$ and the two bounds involving $\ell^*$ are conditions of the parameter choices.
\item[(b)] ($h<g_0$ and $h<g_1$.) The denominator of $\ell^*$ is at least $m$, so $\ell^*\le b$. By
  (a), $S+\eta<\ell^*\le b$ and $t\le\eta$, so $h<3b+1$. As $k\geq5$ and $b>1$, $3b+1\le kb\le(kb)^8$,
  and the parameter choices give a number $Y$ with $(kb)^8\le Y\le g_0,g_1$.
\item[(c)] (Dominant weights.) For every embedding $\tau$, the Harish-Chandra parameters give at level
  $j+1$ a dominant weight $\lambda'_\tau$ of length $N_r$ with integral entries. This is
  Lemma~\ref{is-onto-step-exp-dominant} with $J=j+1$, as $N_r=r+t$.
\item[(d)] (Rows.) Beyond a scale, with $X'=e^{x^{\varepsilon/3}}$ and $A=x^{2|C_t|+2}$: $b'=xb/X'$,
  $X'(S+\eta)\le b$, $kS\le w''$, $S\geq1$, $m\geq2$, $f=2n$, $f\le k$, $t\le A$ and $N\le Ax$; and
  $x\geq5$, $L\geq4$, $4(L+2)\le X'$, $5.2A^3x\le X'$, $40A^7x^2L\le X'$ and $40A^7x^2L^3\le e^L$.
\item[(e)] (Largeness.) $L\geq4$, $48k^2t^2<b$, $12ktb'<b$, $2(2b/\ell^*+t)<w''$ and
  $S+b'(1+1/x+1/x^2)+\eta+(t-1)/2\le2b'$.
\item[(f)] (Weights.) There are weights $W=(\lambda_\tau)_\tau$, with a set of thick embeddings and a
  special set, that satisfy the conditions of the numerical conditions relative to $P$, such that every
  $\lambda_\tau$ differs from $\lambda'_\tau$ by a constant, the numbers $\lambda_{\tau,a}+N_r-1-a$ for
  $a<N_r$ lie in $[0,\lfloor h\rfloor]$, there are at most $|T_2|$ thick embeddings, and the special set
  is the one the member determines.
\end{enumerate}

\emph{Proof of (d).} Put $C=2|C_t|+2$ and $A=x^C$. For every $\delta>0$, every fixed power of $x$ is at
most $e^{x^\delta}$ for large $x$. So there is $x_0$, depending on $\varepsilon$ and $C_t$, such that for
$x\geq x_0$: $x\geq5$, $L=x^\varepsilon\geq4$, the numbers $4(L+2)$, $5.2A^3x$ and $40A^7x^2L$ are at most
$X'=e^{x^{\varepsilon/3}}$, and $40A^7x^2L^3\le e^{L}$. Let $x\geq\max(x_0,2)$. By the definitions of the
size data (Definition~\ref{pd-size-data}(1)), $b'=xb/X'$ and $X'w''=b/n$. The parameter choices give
$kS\le w''$ and $m(t+1)\eta\le w''$; as $k\geq f=2n\geq2$ and $m(t+1)\geq2$, $S\le w''/2$ and
$\eta\le w''/2$, so $X'(S+\eta)\le X'w''=b/n\le b$. The cloud size is $S=4k\Delta_S$, and the parameter
choices give $\Delta_S\geq N(m)m^{10}\geq1$, so $S\geq1$. The bounds $m\geq2$, $f=2n$ and $f\le k$ are
conditions of the ranks (Definition~\ref{th-ranks}). The size data give $1\le t\le x^{C_t}$, so $t\le A$.
Finally $k\le t$ and $m\le t$ give $N=k+m(t+1)\le3t^2\le3x^{2C_t}$, and $3x^{2C_t}\le x^{2C_t+1}x\le Ax$,
as $x\geq2$ and $2C_t+1\le C$.

\emph{Proof of (e).} Put $a_k=\frac92+10\varepsilon+\frac\varepsilon{100}$ and
$a_n=2+8\varepsilon+\frac\varepsilon{100}$. Enlarging $x_0$, assume also that $4x^{a_k+a_n}\le e^{L/10}$,
$20x^{\varepsilon/3}\le L$, $10x^{\varepsilon/4}\le L$ and $5x^{2+2\varepsilon}\le e^{8L/35}$ for $x\geq x_0$.
Two further bounds come first.
\begin{enumerate}
\item[(e1)] $b^{4/5}S\le w''$. The size conditions of the member give $k\le x^{a_k}$ and $n\le x^{a_n}$, and
  the parameter choices give $\log\Delta_S\le x^{\varepsilon/4}/2$. So
  \[
   nS=4kn\Delta_S\le4x^{a_k+a_n}e^{x^{\varepsilon/4}/2}\le e^{L/10+x^{\varepsilon/4}/2}\le
   e^{L/5-x^{\varepsilon/3}},
  \]
  since $x^{\varepsilon/4}/2+x^{\varepsilon/3}\le L/20+L/20$. As $w''=e^{L-x^{\varepsilon/3}}/n$, this gives
  $e^{4L/5}S\le w''$.
\item[(e2)] $4x^2L^2b^{4/7}<b^{4/5}$, since $4x^2L^2=4x^{2+2\varepsilon}<e^{8L/35}$ and
  $\frac47+\frac8{35}=\frac45$.
\end{enumerate}
By (d), $A^7\le40A^7x^2L^3\le b$, so $A^4\le b^{4/7}$ and, by (e2), $4A^4x^2L^2<b^{4/5}$. By (a) and (d),
$n\le k\le t\le A$, $x\geq5$, $L\geq4$, $S\geq1$ and $t\le\eta$. Now the five conditions of (e).
$L\geq4$ is part of (d). Next, $48k^2t^2\le48A^4<40\cdot25\cdot64A^4\le40A^7x^2L^3\le b$. Next, $xL\geq20$,
so $12ktx\le12A^2x<40A^2x\cdot xL\le40A^7x^2L\le X'$, and $12ktb'=12ktxb/X'<b$. For the fourth, let $Q$ be the
denominator of $\ell^*$, so that $2b/\ell^*=2Q/m\le Q$ as $m\geq2$. As $n,k,t\le A$ and
$N+t\le Ax+A\le2Ax$, $Q\le2A^4xL^2+A^3(S+t+L)$, so $2b/\ell^*\le4A^4xL^2+2A^3(S+t+L)$. For $x\geq5$,
$L\geq4$, $A\geq1$, $S\geq0$ and $t\le A$ an elementary computation gives
$4A^4xL^2+2A^3(S+t+L)+t\le A^4x^2(L^2+S)$. As $L^2\geq1$ and $S\geq1$, $L^2+S\le2L^2S$, and therefore
\[
 2\Big(\frac{2b}{\ell^*}+t\Big)\le2A^4x^2(L^2+S)\le4A^4x^2L^2S<b^{4/5}S\le w''
\]
by (e1). For the last, $S+\eta\le b/X'=b'/x\le b'/5$ by (d), $(t-1)/2<\eta/2\le b'/10$, and
$b'/x+b'/x^2\le b'/5+b'/25$; the sum is less than $2b'$.
% Lean: (d) is rows_big_of_largeParameter (Scale.b', Scale.w'', ParameterChoice' kS_le, mteta,
% Del_ge.2.2.2, SizeData Sc_eq and t_win, Ranks f_le_k and two_le_m) with
% LargeParameterEstimates.eventually_large_parameter_bound at eps and 2|Ct| + 2 (eventually_power_le_exp_power
% and eventually_mul_pow_le_exp_rpow); the scale is max x0 2, and the members share eps (L_eps) and Ct
% (Ct_eq). (e) is large_of_rows_pow: (e1) and (e2) are parOf_further_alt (ev_large, w_over_S_real with
% SizeK, SizeN and ParameterChoice' logDel, pow_real), A^7 <= b from Big.b3, and
% LargeParameterEstimates.large_bound_of_constraints: forty_eight_k_sq_t_sq_lt_b, twelve_k_t_mul_lt_b,
% double_length_bound_lt_width (maximal_index_size_bound with polynomial_inequality_large_x) and
% top_level_sum_le_two_b, through large_parameter_facts_of_hyp.

\emph{Proof of (f).} Weights for $P$ consist of, for each embedding $\tau$, a sequence $\lambda_\tau$ of
natural numbers, natural numbers $j_\tau$ and $q_\tau$, integers $x^-_\tau$ and $x^+_\tau$ and real
numbers $c_\tau$ and $e_\tau$; a set $T_W$ of embeddings, the \emph{thick embeddings}; and a special set
of primes. Put $U_\tau(a)=\lambda_{\tau,a}+N_r-1-a$ and $p_\tau=k-1+j_\tau$. The conditions on $W$
relative to $P$ are the following, for every $\tau$.
\begin{enumerate}
\item[(W1)] $\lambda_{\tau,a'}\le\lambda_{\tau,a}$ for $a\le a'<N_r$, and $k-1+j_\tau+t+q_\tau=N_r$.
\item[(W2)] $\lambda_{\tau,a}=\lambda_{\tau,p_\tau}$ for $p_\tau\le a<p_\tau+t$, and
  $U_\tau(p_\tau+t-1)+\frac{t-1}2-c_\tau=-e_\tau$.
\item[(W3)] $c_\tau\le x^-_\tau\le x^+_\tau\le c_\tau+S$, and $x^-_\tau\le U_\tau(a)\le x^+_\tau$ for
  $a<k-1$.
\item[(W4)] $\frac{2b}{\ell^*}+\frac t2+3\le|U_\tau(a)-c_\tau+e_\tau|\le\eta$ for $k-1\le a<p_\tau$ and
  for $p_\tau+t\le a<N_r$.
\item[(W5)] If $\tau\in T_W$, then $b\le e_\tau\le b+\eta+1$ and $x^+_\tau\le h$.
\item[(W6)] If $\tau\notin T_W$, then, with the \emph{width} $U_\tau(0)-U_\tau(N_r-1)$, either there is
  $b_*$ with $b'\le b_*\le b'(1+x^{-2})$ and $b_*+b'/x\le e_\tau\le b_*+b'/x+\eta$, and the width is at
  most $S+b'(1+\frac1x+\frac1{x^2})+\eta+\frac{t-1}2$; or $w''\le e_\tau\le w''+\eta$, and the width is at
  most $S+w''+\eta+\frac{t-1}2$.
\end{enumerate}
Let $J=j+1$ and $p=k+J-1$, and for each $\tau$ let $\lambda'_\tau$ be the dominant weight of
Lemma~\ref{is-onto-step-exp-dominant} at $J$, with entries $\lambda'_{\tau,1}\geq\dots\geq\lambda'_{\tau,N_r}$.
Put $\lambda_{\tau,a}=\lambda'_{\tau,a+1}-\lambda'_{\tau,N_r}$ for $a<N_r$, a natural number,
$j_\tau=J$, $q_\tau=m-J$, $c_\tau=\frac{N_r}2-1-\lambda'_{\tau,N_r}$, $x^-_\tau=\lceil c_\tau\rceil$,
$x^+_\tau=\lfloor c_\tau+S\rfloor$ and $e_\tau=e_J(\tau)$; let $T_W$ be the set of embeddings that are
thick places of the host (Definition~\ref{hf-host}), and let the special set be the one the member
determines. Then $\lambda_\tau$ differs from $\lambda'_\tau$ by the constant $-\lambda'_{\tau,N_r}$, and
$|T_W|\le|T_2|$, as the thick places are the images of the elements of $T_2$ and distinct embeddings are
distinct elements of $G_\Phi$. Condition (W1) holds as $\lambda'_\tau$ is dominant and
$N_r=k+m-1+t$. By the formula of Lemma~\ref{is-onto-step-exp-dominant},
$U_\tau(a)-c_\tau=\lambda'_{\tau,a+1}+\frac{N_r}2-a$ equals
\[
 v_{a+1}\ (a<k-1),\quad -b_{a-k+2}\ (k-1\le a<p),\quad -e_J+\tfrac{t-1}2-(a-p)\ (p\le a<p+t),\quad
 -b_{J+a-p-t+1}\ (p+t\le a<N_r),\tag{$*$}
\]
for the Harish-Chandra parameter $(u,v,b,e)$ of $D$ at $\tau$.
\begin{itemize}
\item (W2). By $(*)$, $U_\tau(a)=U_\tau(p)-(a-p)$ for $p\le a<p+t$, so $\lambda_\tau$ is constant there, and
  $U_\tau(p+t-1)-c_\tau=-e_J-\frac{t-1}2$.
\item (W3). The size data give $0\le v_i\le S$. So $c_\tau\le U_\tau(a)\le c_\tau+S$ for $a<k-1$, by
  $(*)$; as $U_\tau(a)$ is an integer, $\lceil c_\tau\rceil\le U_\tau(a)\le\lfloor c_\tau+S\rfloor$. The case
  $a=0$ gives $x^-_\tau\le x^+_\tau$.
\item (W4). In these ranges $U_\tau(a)-c_\tau=-b_i$ for some $i\le m$, by $(*)$, so
  $|U_\tau(a)-c_\tau+e_\tau|=|b_i-e_J|$. The size data give $|b_i-e_J|\le\eta$, and the parameter choices give
  $|b_i-e_J|\geq2b/\ell^*+t/2+3$.
\item Two bounds. Since $\lambda_{\tau,N_r-1}=0$, $U_\tau(N_r-1)=0$. If $J<m$, $(*)$ at $a=N_r-1$ gives
  $c_\tau=b_m\le e_m-\frac{t+3}2$; if $J=m$, it gives $c_\tau=e_m+\frac{t-1}2$. In both cases
  $c_\tau\le e_m(\tau)+\frac{t-1}2$. The width is $U_\tau(0)=c_\tau+v_1\le c_\tau+S\le S+e_m(\tau)+\frac{t-1}2$.
\item (W5). If $\tau$ is a thick place, the windows of the size data give $b\le e_{J'}(\tau)\le b+\eta+1$
  for every $J'$. With $J'=J$ this is the first part, and $x^+_\tau\le c_\tau+S\le e_m(\tau)+\frac{t-1}2+S
  \le b+S+\eta+t+1=h$.
\item (W6). Let $\tau\notin T_W$. If $\tau=\tau_1(t')$ for a top class $t'$, the windows give
  $b_{t'}+b'/x\le e_{J'}(\tau)\le b_{t'}+b'/x+\eta$ for every $J'$, and $b_{t'}\geq b'$
  (Definition~\ref{pd-size-data}(2)). The parameter choices give $b_{t'}\le b'+N^*x^{\varepsilon}\Lambda_Gb'/k$
  with $N^*=2n\delta_1$, and $4N^*x^\varepsilon\Lambda_G/k+C_\rho\Lambda_G\log x/k\le x^{-2}$ for a constant
  $C_\rho>0$, where $\Lambda_G>0$ since $4S\le\Lambda_Gb'/k$; so $b_{t'}\le b'(1+x^{-2})$. The first alternative holds with
  $b_*=b_{t'}$, since the width is at most $S+e_m(\tau)+\frac{t-1}2\le S+b'(1+x^{-2})+b'/x+\eta+\frac{t-1}2$.
  Otherwise $\tau$ is neither a top place nor a thick place, and the windows give $w''\le
  e_{J'}(\tau)\le w''+\eta$ for every $J'$; the second alternative holds, as the width is at most
  $S+e_m(\tau)+\frac{t-1}2\le S+w''+\eta+\frac{t-1}2$.
\item The box. $U_\tau(a)\geq0$, as $\lambda_{\tau,a}\geq0$ and $a\le N_r-1$; and $U_\tau(a)\le U_\tau(0)$ by
  (W1). If $\tau$ is a thick place, $U_\tau(0)\le h$ as in (W5). Otherwise $e_m(\tau)\le
  b'(1+\frac1x+\frac1{x^2})+\eta$, by the windows as in (W6), and $w''\le b'$ as $xn\geq1$. So
  $U_\tau(0)\le S+b'(1+\frac1x+\frac1{x^2})+\eta+\frac{t-1}2\le2b'$ by (e), and $2b'\le12ktb'<b\le h$ by (e).
  As $U_\tau(a)$ is an integer, $0\le U_\tau(a)\le\lfloor h\rfloor$.
\end{itemize}
% Lean: (c) is LiftWeightForced.isLamPrime_lamLdef (T.isLamPrime_at_all_points), through
% exists_isLamPrime and GlobalTheory.LowestKType.thetaSide; (d) is T.exists_large_parameter_estimates_at_pts
% = ExpTie.rows_big_alt (Rows and Big of LargeParameterEstimates, from rows_big_of_largeParameter and
% eventually_large_parameter_bound); (e) is T.parameter_large_at_all_points = parOf_large_at_pts
% (large_of_rows_pow with parOf_further_alt); (f) is T.weight_lambda_shift_nonneg' = Lift.wts_ok_alt
% (the weights wts with W.Ok P: anti, what2, block and centre from liftWeightRho_block_constant_centred,
% window and cloud from liftWeightRho_window_contains_cloud, low from liftWeightRho_lows_distance_bounds,
% thick and other from liftWeightRho_thickPlace_centre_window and _nonThickPlace_centre_width; the box
% from U_nonneg and liftWeightRho_entries_le_floor_height; T2_card; the shift -detShift). In (f): W is
% Lift.wts (lam = lamT, the natW of lamL; T2 = the tau with eTau tau in thickPlaces; xm, xp, c = cT, e =
% wj D (j + 1)); (W1)-(W6) are the fields anti and what2, block and centre, window and cloud, low, thick,
% other of Wts.Ok (kc P = k - 1, bot = U at kc + j + t - 1, width = U 0 - U (Nr - 1)); (*) is ValsAt,
% from valsAt_of_isLamPrime with U_sub_c_eq_w; the two bounds are U_last and c_le; the size-data fields
% are cloud_v, be_close, win_thick, win_top, win_thin and bT_ge, and the parameter choices lows_far,
% bT_le, rho_le and S_le; the box is U_nonneg and wts_U_le_h0_of_large with Large.top_le and Large.bb'.
% The fields fr and cj of Wts (identity and false) do not enter Wts.Ok.

So $P$ and $W$ are in the exponent dictionary at $j$, by the construction, (b), (c) and (f), and they
satisfy its numerical conditions, by (a), (d), (e) and (f).
% Lean: ExpTie.numSt_alt, assembling T.parOf_hyp_at_pts (= parOf_hyp_alt, item (a): the ranks'
% two_le_m, m_add_three_le_k, rank_add_five_le, three_m_lt_t, D.sp.V_rank, and ParameterChoice'
% eta_le, t_le_eta, lstar1, lstar2), T.parOf_height_lt_genus_at_pts (= parOf_h_lt_g, item (b), with
% g_Y), (c), (d), (e) and (f) by atPts_and_alt; P is parOf (Fm.L x i) j.
\end{proof}

\begin{definition}\label{is-onto-step-exp-cut-exponent}
Let $N\geq0$ be an integer, $\nu=(\nu_1\geq\dots\geq\nu_N\geq0)$ integers and $l\geq0$ an integer, and let
$\nabla_{\mathbf Z}(\nu)$ and $\Delta_{\mathbf Z}(\nu)$ be the dual Weyl module and the Weyl module of
Definition~\ref{pd-weight-lattices}(1) over $\mathbf Z$, in the polynomials in the entries of an
$N\times N$ matrix $X$ of variables. The \emph{Weyl exponent} $e_l(\nu)$ is the least integer $e\geq0$
such that $ml^{e}\nabla_{\mathbf Z}(\nu)\subset\Delta_{\mathbf Z}(\nu)$ for some integer $m$ not
divisible by $l$, and $0$ if there is no such $e$. For a sequence $\lambda=(\lambda_0,\lambda_1,\dots)$ of natural
numbers whose first $N$ terms are nonincreasing, $e_l(N,\lambda)=e_l(\lambda_0,\dots,\lambda_{N-1})$.
\end{definition}
\leanref{Formalism/PieriBlockIdeals/StatementLA.lean}{Formalism.CohomologicalDescent.LA.InclZ}

\begin{lemma}\label{is-onto-step-exp-cut-attained}
Let $N\geq0$, $\nu\in\mathbf N^N$ nonincreasing and $l$ a prime. There is an integer $M\ne0$ with
$M\nabla_{\mathbf Z}(\nu)\subset\Delta_{\mathbf Z}(\nu)$. Hence there is an integer $m$ prime to $l$
with $ml^{e_l(\nu)}\nabla_{\mathbf Z}(\nu)\subset\Delta_{\mathbf Z}(\nu)$
(Definition~\ref{is-onto-step-exp-cut-exponent}).
\end{lemma}
\leanref{Formalism/PieriBlockIdeals/PieriColumnMinors/ProdInitialMinorsMemDetSpan.lean}{Formalism.CohomologicalDescent.LA.la_uniform}
\leanref{Formalism/PieriBlockIdeals/PieriColumnMinors/ProdInitialMinorsMemDetSpan.lean}{Formalism.CohomologicalDescent.LA.exact_z_specialization}
\begin{proof}
The module $\nabla_{\mathbf Z}(\nu)$ lies in the finitely generated $\mathbf Z$-module of homogeneous
polynomials of degree $\nu_1+\dots+\nu_N$, so it is finitely generated, say by $g_1,\dots,g_s$. Each
$g_i$ lies in $\nabla_{\mathbf Q}(\nu)$, which is the $\mathbf Q$-span of $\Delta_{\mathbf Z}(\nu)$
by Lemma~\ref{is-standard-basis}(2); so $c_ig_i\in\Delta_{\mathbf Z}(\nu)$ for some integer $c_i\ne0$.
Then $M=c_1\cdots c_s$ works. Write $|M|=l^vm$ with $l\nmid m$. The set of $e\geq0$ for which
$m'l^{e}\nabla_{\mathbf Z}(\nu)\subset\Delta_{\mathbf Z}(\nu)$ for some $m'$ prime to $l$ contains $v$,
so it is not empty, and its least element $e_l(\nu)$ belongs to it.
% Lean: la_uniform (la_dualWeyl_fg through homogeneousSubmodule_fg, and weylRational, the integer
% multiples of Lemma is-standard-basis(2)); exact_z_specialization (ExZSpec) with
% Nat.exists_eq_pow_mul_and_not_dvd and Nat.sInf_mem.
\end{proof}

\begin{lemma}\label{is-onto-step-exp-cut-beta}
Let $l$ be a prime and $w\geq0$ an integer, and write $\lfloor\log_lw\rfloor$ for the largest $v\geq0$
with $l^v\le w$ ($0$ if $w=0$). There is an integer $m$ prime to $l$ such that
$\Lambda=ml^{\lfloor\log_lw\rfloor}$ satisfies: $(j+k+1)!$ divides $\Lambda\,j!\,k!$ for all integers
$j,k\geq0$ with $j+k+1\le w$.
\end{lemma}
\leanref{Development/AbelianVarietyTheory/Proofs/DivisorProofs/BetaIntegralDenominatorBoundStatement.lean}{GlobalTheory.T.beta_integral_denominator_bound}
\leanref{Development/AbelianVarietyTheory/Proofs/DivisorProofs/NablaSchur.lean}{GlobalTheory.LeviCut.BetaDen}
\begin{proof}
For $1\le s\le w$ let $s'$ be the largest divisor of $s$ prime to $l$, and put $m=\prod_{s=1}^ws'$, which
is prime to $l$. Every $t$ with $1\le t\le w$ divides $\Lambda$: write $t=l^vt'$; then $t'$ divides $m$,
and $l^v\le t\le w$ gives $v\le\lfloor\log_lw\rfloor$.

The divisibility is proved by induction on $k$, for all $j$. For $k=0$, $(j+1)!=(j+1)\,j!$ and $j+1$
divides $\Lambda$. Let $j+k+2\le w$. Then
\[
 \Lambda\,j!\,(k+1)!+\Lambda\,(j+1)!\,k!=(j+k+2)\,\Lambda\,j!\,k!.
\]
By induction for $(j,k)$, $(j+k+2)!$ divides the right side, and by induction for $(j+1,k)$ it divides
$\Lambda\,(j+1)!\,k!$. So it divides $\Lambda\,j!\,(k+1)!$.
% Lean: LeviCut.betaDen, with m the product of the ordCompl[l] s for 1 <= s <= w,
% BetaDenProof.dvd_lam_of_le and BetaDenProof.beta_dvd.
\end{proof}

\begin{definition}\label{is-onto-step-exp-cut-grid}
Let $V$ be a $\mathbf Q$-vector space, $N\geq0$, and $e=(e_{ab})_{1\le a,b\le N}$ endomorphisms of $V$.
\begin{enumerate}
\item $e$ satisfies the \emph{$\mathfrak{gl}_N$ relations} if
  $e_{ab}e_{cd}-e_{cd}e_{ab}=\delta_{bc}e_{ad}-\delta_{da}e_{cb}$ for all $a,b,c,d$.
\item A \emph{weight vector} of weight $\mu\in\mathbf Z^N$ is an $x\in V$ with $e_{cc}x=\mu_cx$ for
  every $c$. The weights of $e$ are \emph{bounded} by an integer $K$ if every $x\in V$ with
  $e_{aa}x-e_{bb}x=nx$ for some $a,b$ and some integer $n>K$ is $0$. The family $e$ is \emph{dominant}
  for $\lambda\in\mathbf Z^N$ if every nonzero weight vector of weight $\mu$ has
  $\mu_1+\dots+\mu_k\le\lambda_1+\dots+\lambda_k$ for $k\le N$, with equality for $k=N$.
\item A $\mathbf Z$-submodule $L\subset V$ is \emph{stable under divided powers} if
  $\frac1{k!}e_{ab}^kL\subset L$ for all $a\ne b$ and $k\geq0$.
\item For a $\mathbf Q$-algebra $A$, integers $K,j\geq0$ and $X,Y\in A$ put
  \[
   P_{K,j}(X,Y)=\sum_{k=0}^{K}\frac{(-1)^k\,j!}{k!\,(k+j+1)!}\,Y^kX^k .
  \]
  For a family $D=(D_{ab})$ of elements of $A$ and integers $p,q\geq0$ with $p+q\le N$, the \emph{grid
  projector} is $G_K(D;p,q)=R_pR_{p-1}\cdots R_1$, where
  $R_i=P_{i,p+1}P_{i,p+2}\cdots P_{i,p+q}$ and $P_{ij}=P_{K,j-i-1}(D_{ij},D_{ji})$.
\item For an endomorphism $x$ of $V$, $\operatorname{ad}x$ is the endomorphism $T\mapsto xT-Tx$ of
  $\operatorname{End}(V)$, and $\operatorname{ad}e=(\operatorname{ad}e_{ab})$.
\item For $\nu\in\mathbf N^N$ nonincreasing, the \emph{row operators} on $\nabla_{\mathbf Q}(\nu)$
  (Definition~\ref{pd-weight-lattices}(1)) are $E_{ab}=\sum_cx_{ac}\,\partial/\partial x_{bc}$; they
  preserve $\nabla_{\mathbf Q}(\nu)$, as they act on the rows. The lattices $\nabla_{\mathbf Z}(\nu)$
  and $\Delta_{\mathbf Z}(\nu)$ are viewed inside $\nabla_{\mathbf Q}(\nu)$.
\end{enumerate}
\end{definition}
\leanref{Development/AbelianVarietyTheory/Proofs/DivisorProofs/NablaSchur.lean}{GlobalTheory.LeviCut.GlRel}
\leanref{Development/AbelianVarietyTheory/Proofs/DivisorProofs/NablaSchur.lean}{GlobalTheory.LeviCut.Bdd}
\leanref{Development/AbelianVarietyTheory/Proofs/DivisorProofs/NablaSchur.lean}{GlobalTheory.LeviCut.Dom}
\leanref{Development/AbelianVarietyTheory/Proofs/DivisorProofs/NablaSchur.lean}{GlobalTheory.LeviCut.DivStable}
\leanref{Development/AbelianVarietyTheory/Proofs/DivisorProofs/NablaSchur.lean}{GlobalTheory.LeviCut.gridP}
\leanref{Development/AbelianVarietyTheory/Proofs/DivisorProofs/NablaSchur.lean}{GlobalTheory.LeviCut.adj}
\leanref{Development/AbelianVarietyTheory/Proofs/DivisorProofs/NablaSchur.lean}{GlobalTheory.LeviCut.eV}

\begin{lemma}\label{is-onto-step-exp-cut-adjoint}
Let $V$ be a $\mathbf Q$-vector space and $e=(e_{ab})_{1\le a,b\le N}$ endomorphisms of $V$ that
satisfy the $\mathfrak{gl}_N$ relations, whose weights are bounded by $K$ and whose weight vectors span
$V$ (Definition~\ref{is-onto-step-exp-cut-grid}). Then $\operatorname{ad}e$ satisfies the
$\mathfrak{gl}_N$ relations on $\operatorname{End}(V)$, and its weights are bounded by $2K$.
\end{lemma}
\leanref{Development/AbelianVarietyTheory/Proofs/DivisorProofs/AdjointActionEndomorphismWeightBoundStatement.lean}{GlobalTheory.T.adjoint_action_endomorphism_weight_bound}
\leanref{Development/AbelianVarietyTheory/Proofs/DivisorProofs/NablaSchur.lean}{GlobalTheory.LeviCut.AdjBdd}
\begin{proof}
For endomorphisms $x,y$, $\operatorname{ad}x\operatorname{ad}y-\operatorname{ad}y\operatorname{ad}x=
\operatorname{ad}(xy-yx)$ and $\operatorname{ad}$ is linear, which gives the relations. Let
$T\in\operatorname{End}(V)$ with $\operatorname{ad}(e_{aa})T-\operatorname{ad}(e_{bb})T=nT$ and
$n>2K$, and let $x$ be a nonzero weight vector of weight $\mu$. Then $\mu_b-\mu_a\le K$, since
otherwise the bound applied to $(b,a)$ gives $x=0$. Evaluating at $x$ gives
$(e_{aa}-e_{bb})Tx=(n+\mu_a-\mu_b)Tx$, and $n+\mu_a-\mu_b>K$; so $Tx=0$. As the weight vectors span
$V$, $T=0$.
% Lean: adjBdd_holds, with adj_commutator_of_adjBdd and LinearMap.ext_on.
\end{proof}

\begin{lemma}\label{is-onto-step-exp-cut-invariance-rank-one}
Let $W$ be a $\mathbf Q$-vector space, $N\geq1$, and $D=(D_{cd})_{1\le c,d\le N}$ endomorphisms of $W$
that satisfy the $\mathfrak{gl}_N$ relations, with weights bounded by $2K$
(Definition~\ref{is-onto-step-exp-cut-grid}). Let $a\ne b$ and $T\in W$ with $D_{aa}T=D_{bb}T=0$. Then:
\begin{enumerate}
\item $D_{ab}\,P_{K,0}(D_{ab},D_{ba})\,T=0$;
\item if $D_{ab}T=0$, then $D_{ba}T=0$.
\end{enumerate}
\end{lemma}
\leanref{Development/AbelianVarietyTheory/Proofs/DivisorProofs/RankOneProjectorWeightZeroInvarianceStatement.lean}{GlobalTheory.T.rank_one_projector_weight_zero_invariance}
\leanref{Development/AbelianVarietyTheory/Proofs/DivisorProofs/NablaSchur.lean}{GlobalTheory.LeviCut.RankOne}
\begin{proof}
Put $X=D_{ab}$, $Y=D_{ba}$ and $H=D_{aa}-D_{bb}$. The relations give $XY-YX=H$, $HX-XH=2X$ and
$HY-YH=-2Y$, and $HT=0$. By induction, $HX^kw=X^kHw+2kX^kw$, $HY^kw=Y^kHw-2kY^kw$ and
\[
 XY^{k+1}w=Y^{k+1}Xw+(k+1)Y^kHw-(k+1)k\,Y^kw\qquad(w\in W,\ k\geq0).
\]
In particular $HX^{K+1}T=2(K+1)X^{K+1}T$ and $(D_{bb}-D_{aa})Y^{K+1}T=2(K+1)Y^{K+1}T$; as $2(K+1)>2K$,
the bound on the weights gives $X^{K+1}T=0$ and $Y^{K+1}T=0$.

(1) Let $c_k=(-1)^k/(k!\,(k+1)!)$, the coefficients of $P_{K,0}$; then $c_{k+1}(k+1)(k+2)=-c_k$. Since
$HX^{m+1}T=2(m+1)X^{m+1}T$, the formula gives
$X(Y^{m+1}X^{m+1}T)=Y^{m+1}X^{m+2}T+(m+1)(m+2)Y^mX^{m+1}T$. By induction on $n$,
$X\sum_{k=0}^nc_kY^kX^kT=c_nY^nX^{n+1}T$: adding the term $k=n+1$ adds
$c_{n+1}Y^{n+1}X^{n+2}T+c_{n+1}(n+1)(n+2)Y^nX^{n+1}T$, and the second summand cancels $c_nY^nX^{n+1}T$.
At $n=K$, $XP_{K,0}(X,Y)T=c_KY^KX^{K+1}T=0$.

(2) If $XT=0$, the formula gives $XY^{k+1}T=-(k+1)k\,Y^kT$. If $Y^{i+2}T=0$ for some $i\geq0$, then
$0=XY^{i+2}T=-(i+2)(i+1)Y^{i+1}T$, so $Y^{i+1}T=0$. Starting from $Y^{K+1}T=0$, this gives $YT=0$.
% Lean: LeviCut.rankOne: RankOnePf.rankOne_H_X_pow, rankOne_H_Y_pow, rankOne_X_Y_pow, rankOne_coef_succ;
% hterm and hsum for (1), with hXK from the bound Bdd; hYterm, hYK and the descent hdesc for (2).
\end{proof}

\begin{lemma}\label{is-onto-step-exp-cut-invariance-rank-two}
Let $W$, $N$, $K$ and $D$ be as in Lemma~\ref{is-onto-step-exp-cut-invariance-rank-one}. Let $a,b,c$ be
distinct, $j\geq0$, and $T\in W$ with $D_{aa}T=D_{bb}T=D_{cc}T=0$ and $D_{bc}T=0$. Then
\[
 D_{bc}\,P_{K,j}(D_{ab},D_{ba})\,P_{K,j+1}(D_{ac},D_{ca})\,T=0.
\]
\end{lemma}
\leanref{Development/AbelianVarietyTheory/Proofs/DivisorProofs/RankTwoConsecutiveProjectorInvarianceStatement.lean}{GlobalTheory.T.rank_two_consecutive_projector_invariance}
\leanref{Development/AbelianVarietyTheory/Proofs/DivisorProofs/NablaSchur.lean}{GlobalTheory.LeviCut.RankTwo}
\begin{proof}
Put $X=D_{ac}$, $Y=D_{ca}$, $X'=D_{ab}$, $Y'=D_{ba}$, $Z=D_{bc}$, $H=D_{aa}-D_{cc}$ and $H'=D_{aa}-D_{bb}$,
and let $c_{i,k}=(-1)^ki!/(k!\,(k+i+1)!)$, so that $P_{K,i}(A,B)=\sum_{k\le K}c_{i,k}B^kA^k$ and
$c_{i,k+1}(k+1)(k+i+2)=-c_{i,k}$. The relations give: $ZX'-X'Z=-X$, $ZY-YZ=Y'$, and $Z$ commutes with
$X$ and $Y'$; $X$ commutes with $X'$, and $Y$ with $Y'$; $XY-YX=H$, $HX-XH=2X$, $HY-YH=-2Y$;
$X'Y'-Y'X'=H'$, $H'X'-X'H'=2X'$; $H'X-XH'=X$ and $H'Y-YH'=-Y$. Also $HT=H'T=0$. We use three formulas,
each proved by induction on $k$:
\begin{enumerate}
\item[(i)] if $AB-BA=C$ and $C$ commutes with $B$, then $AB^{k+1}w=B^{k+1}Aw+(k+1)CB^kw$;
\item[(ii)] if $AB-BA=C$, $CB-BC=-2B$ and $Cw=\gamma w$, then
  $AB^{k+1}w=B^{k+1}Aw+(k+1)(\gamma-k)B^kw$;
\item[(iii)] if $AB-BA=C$, $CA-AC=2A$ and $Cw=\gamma w$, then
  $A^{k+1}Bw=BA^{k+1}w+(k+1)(\gamma+k)A^kw$.
\end{enumerate}

Since $Z$ commutes with $X$ and $ZT=0$, $ZX^kT=0$ for all $k$. Since $HT=0$, $HX^kT=2kX^kT$, and the bound on
the weights gives $X^{K+1}T=0$. Put
\[
 T'=P_{K,j+1}(X,Y)T=\sum_{k=0}^Kc_{j+1,k}Y^kX^kT,\qquad U=\sum_{k=0}^{K-1}(k+1)c_{j+1,k+1}Y^kX^{k+1}T.
\]

\emph{$ZT'=Y'U$.} By (i) for $Z$, $Y$ and $Y'$, and $ZX^{k+1}T=0$,
$ZY^{k+1}X^{k+1}T=(k+1)Y'Y^kX^{k+1}T$; and the term $k=0$ of $T'$ is killed by $Z$.

\emph{$XT'=-(j+1)U$.} By (ii) for $X$, $Y$, $H$ and $\gamma=2(k+1)$,
$XY^{k+1}X^{k+1}T=Y^{k+1}X^{k+2}T+(k+1)(k+2)Y^kX^{k+1}T$. The recursion of the $c_{j+1,k}$ gives
$c_{j+1,k}+c_{j+1,k+1}(k+1)(k+2)=-(j+1)(k+1)c_{j+1,k+1}$, and induction on $n$ gives
\[
 X\sum_{k=0}^nc_{j+1,k}Y^kX^kT=c_{j+1,n}Y^nX^{n+1}T-(j+1)\sum_{k=0}^{n-1}(k+1)c_{j+1,k+1}Y^kX^{k+1}T.
\]
At $n=K$ the first term vanishes.

\emph{$X'^KU=0$.} As $H'X^{k+1}T=(k+1)X^{k+1}T$ and $H'Y=YH'-Y$, $H'U=U$. Hence
$H'X'^KU=(2K+1)X'^KU$, and the bound on the weights gives $X'^KU=0$.

\emph{The projector.} For $l\geq0$, as $Z$ commutes with $Y'$, (i) for $Z$, $X'$ and $-X$ (which commutes
with $X'$) and the two identities above give
\[
 ZY'^{l+1}X'^{l+1}T'=Y'^{l+1}\big(X'^{l+1}Y'U+(l+1)(j+1)X'^lU\big),
\]
and (iii) for $X'$, $Y'$, $H'$ and $\gamma=1$ turns this into
$Y'^{l+2}X'^{l+1}U+(l+1)(l+j+2)Y'^{l+1}X'^lU$. Induction on $n$, starting from $ZT'=Y'U$ and using
$c_{j,n+1}(n+1)(n+j+2)=-c_{j,n}$, gives
$Z\sum_{l=0}^nc_{j,l}Y'^lX'^lT'=c_{j,n}Y'^{n+1}X'^nU$. At $n=K$ this is
$ZP_{K,j}(X',Y')T'=c_{j,K}Y'^{K+1}X'^KU=0$.
% Lean: RankTwoProof.rankTwo: the relations rZXp, rZXq, rZYp, rZYq, rXqXp, rYpYq, sXYp, sXYq and the
% weight relations; comm_pow_apply (i), sl2_X_Ypow (ii), sl2_Xpow_Y (iii), eig_pow, coef_succ_mul; hZv,
% hXqK, hZT', hXqT' (with key and hS), hHpU, hXpU, hf and hT5.
\end{proof}

\begin{lemma}\label{is-onto-step-exp-cut-invariance}
Let $W$ be a $\mathbf Q$-vector space and $D=(D_{ab})_{1\le a,b\le p+q}$ endomorphisms of $W$ that
satisfy the $\mathfrak{gl}_{p+q}$ relations, with weights bounded by $2K$. Let $T\in W$ satisfy
$D_{cd}T=0$ whenever $c,d\le p$ or $c,d>p$. Then $D_{cd}\big(G_K(D;p,q)T\big)=0$ for all $c,d$
(Definition~\ref{is-onto-step-exp-cut-grid}).
\end{lemma}
\leanref{Development/AbelianVarietyTheory/Proofs/DivisorProofs/GridInvarianceOfRankOneTwoStatement.lean}{GlobalTheory.T.grid_invariance_of_rank_one_two}
\leanref{Development/AbelianVarietyTheory/Proofs/DivisorProofs/NablaSchur.lean}{GlobalTheory.LeviCut.GridInv}
\begin{proof}
Write $N=p+q$, $P_{ij}=P_{K,j-i-1}(D_{ij},D_{ji})$ for $i\le p<j$, so that $G=G_K(D;p,q)=R_p\cdots R_1$
with $R_i=P_{i,p+1}\cdots P_{iN}$, and $S=GT$. By the relations,
$D_{ab}D_{cd}=D_{cd}D_{ab}$ when $b\ne c$ and $d\ne a$.

\emph{The diagonal.} Let $c\le N$. With $s=\delta_{ci}-\delta_{cj}$, the relations give
$D_{cc}D_{ij}-D_{ij}D_{cc}=sD_{ij}$ and $D_{cc}D_{ji}-D_{ji}D_{cc}=-sD_{ji}$, hence
$D_{cc}D_{ij}^k-D_{ij}^kD_{cc}=ksD_{ij}^k$ and $D_{cc}D_{ji}^k-D_{ji}^kD_{cc}=-ksD_{ji}^k$, so $D_{cc}$
commutes with $D_{ji}^kD_{ij}^k$ and with every $P_{ij}$. Hence $D_{cc}S=G(D_{cc}T)=0$.

\emph{The simple roots above the diagonal.} Let $1\le k<N$. Say that $x\in W$ has property $(Z_k)$ if
$D_{cc}x=0$ for every $c$ and $D_{k,k+1}x=0$, and that an endomorphism keeps $(Z_k)$ if it carries
vectors with $(Z_k)$ to vectors with $(Z_k)$. The vector $T$ has $(Z_k)$ when $k,k+1\le p$ or $k>p$. By
the diagonal step, every $P_{ij}$ commutes with the $D_{cc}$; and if $i,j\notin\{k,k+1\}$, $D_{ij}$ and
$D_{ji}$ commute with $D_{k,k+1}$, so $P_{ij}$ keeps $(Z_k)$. Three cases.
\begin{itemize}
\item $k+1\le p$. The rows $R_i$ with $i\ne k,k+1$ keep $(Z_k)$. For $j\ne j'$ in $\{p+1,\dots,N\}$ the
  four indices $k,k+1,j,j'$ are distinct, so $P_{k+1,j'}$ commutes with $P_{kj}$, and
  $R_{k+1}R_k=\prod_{j=p+1}^N\big(P_{k+1,j}P_{kj}\big)$. The family $D'_{cd}=-D_{dc}$ satisfies the
  $\mathfrak{gl}_N$ relations, with weights bounded by $2K$, and $P_{K,i}(-A,-B)=P_{K,i}(A,B)$.
  Lemma~\ref{is-onto-step-exp-cut-invariance-rank-two} for $D'$, the indices $(j,k+1,k)$ and $j-k-2$ gives
  $D_{k,k+1}P_{k+1,j}P_{kj}x=0$ for $x$ with $(Z_k)$; so each $P_{k+1,j}P_{kj}$ keeps $(Z_k)$, and so does
  $G$. Hence $S$ has $(Z_k)$.
\item $k=p$. Then $G=P_{p,p+1}x'$ with $x'=P_{p,p+2}\cdots P_{pN}R_{p-1}\cdots R_1$, and $D_{cc}x'T=0$ by
  the diagonal step. As $P_{p,p+1}=P_{K,0}(D_{p,p+1},D_{p+1,p})$,
  Lemma~\ref{is-onto-step-exp-cut-invariance-rank-one}(1) gives $D_{p,p+1}S=0$.
\item $k>p$. In each row $R_i$ the factors $P_{ij}$ with $j\ne k,k+1$ keep $(Z_k)$, and the adjacent pair
  $P_{ik}P_{i,k+1}$ keeps $(Z_k)$ by Lemma~\ref{is-onto-step-exp-cut-invariance-rank-two} for the indices
  $(i,k,k+1)$ and $k-i-1$. So $G$ keeps $(Z_k)$, and $S$ has $(Z_k)$.
\end{itemize}

\emph{The simple roots below the diagonal.} For $1\le k<N$, $D_{kk}S=D_{k+1,k+1}S=0$ and $D_{k,k+1}S=0$, so
Lemma~\ref{is-onto-step-exp-cut-invariance-rank-one}(2) gives $D_{k+1,k}S=0$.

\emph{All roots.} For $c<d<N$, $D_{c,d+1}=D_{cd}D_{d,d+1}-D_{d,d+1}D_{cd}$ and
$D_{d+1,c}=D_{d+1,d}D_{dc}-D_{dc}D_{d+1,d}$; induction on $d-c$ gives $D_{cd}S=0$ for all $c\ne d$, and
$D_{cc}S=0$ by the diagonal step.
% Lean: gridInv_of_rankOne_rankTwo with RankOne (T.rank_one_projector_weight_zero_invariance) and RankTwo
% (T.rank_two_consecutive_projector_invariance): GridProof.commute_diag_rootP (with
% CommutatorPowerFormula.commutator_pow_of_commutator_eq_smul and commute_pair_pow) for the diagonal; PZ
% for (Z_k), keeps_rootP_far; the three cases of hup: keeps_grid_pair with keeps_rows_pair and
% keeps_pair_twist (glRel_twist, bdd_twist, proj_neg), the rank one case for k + 1 = a, and keeps_grid
% with keeps_range'_pair and keeps_pair_row; hdown from RankOne (2); fill.
\end{proof}

\begin{lemma}\label{is-onto-step-exp-cut-scalar-root}
Let $V$ be a $\mathbf Q$-vector space, $E,F,H,T$ endomorphisms of $V$, $v\in V$, $\phi$ a linear form
on $V$ and $n,K,j\geq0$ integers. Suppose that $EF-FE=H$, $HF-FH=-2F$, $Ev=0$, $Hv=nv$, $n\le K$,
$\phi\circ F=0$, and $T(F^kv)=0$ for every $k\geq1$. Then, with $P_{K,j}$ as in
Definition~\ref{is-onto-step-exp-cut-grid}(4),
\[
 \phi\big(P_{K,j}(\operatorname{ad}E,\operatorname{ad}F)(T)\,v\big)\cdot(n+j+1)=\phi(Tv).
\]
\end{lemma}
\leanref{Development/AbelianVarietyTheory/Proofs/DivisorProofs/ProjectorValueHighestWeightVectorStatement.lean}{GlobalTheory.T.projector_value_highest_weight_vector}
\leanref{Development/AbelianVarietyTheory/Proofs/DivisorProofs/NablaSchur.lean}{GlobalTheory.LeviCut.StepScalar}
\begin{proof}
From $HF=FH-2F$ and $EF=FE+H$, induction on $i$ gives $HF^iv=(n-2i)F^iv$ and
$EF^{i+1}v=(i+1)(n-i)F^iv$.

We show by induction on $k$ that $(\operatorname{ad}E)^kT$ kills $F^iv$ for $i>k$ and that
\[
 (\operatorname{ad}E)^kT(F^kv)=(-1)^kk!\prod_{t<k}(n-t)\cdot Tv.
\]
For $k=0$ this is the hypothesis on $T$. For the induction step,
$(\operatorname{ad}E)^{k+1}T(y)=E\big((\operatorname{ad}E)^kT(y)\big)-(\operatorname{ad}E)^kT(Ey)$. If
$i>k+1$, $EF^iv$ is a multiple of $F^{i-1}v$ with $i-1>k$, so both terms vanish at $y=F^iv$. At
$y=F^{k+1}v$ the first term vanishes, and the second is $-(k+1)(n-k)(\operatorname{ad}E)^kT(F^kv)$, which
gives the formula.

Since $\phi\circ F=0$, $\phi\big((\operatorname{ad}F)(S)(y)\big)=\phi(FSy)-\phi(SFy)=-\phi(S(Fy))$ for every
endomorphism $S$, and so $\phi\big((\operatorname{ad}F)^k(S)(y)\big)=(-1)^k\phi\big(S(F^ky)\big)$. With the
formula above,
\[
 \phi\big((\operatorname{ad}F)^k(\operatorname{ad}E)^k(T)\,v\big)=(-1)^k\phi\big((\operatorname{ad}E)^kT(F^kv)\big)
 =k!\prod_{t<k}(n-t)\cdot\phi(Tv).
\]
By the definition of $P_{K,j}$,
\[
 \phi\big(P_{K,j}(\operatorname{ad}E,\operatorname{ad}F)(T)\,v\big)
 =\sum_{k=0}^K\frac{(-1)^kj!}{(k+j+1)!}\prod_{t<k}(n-t)\cdot\phi(Tv).
\]
Induction on $K$ gives
\[
 (n+j+1)\sum_{k=0}^K\frac{(-1)^kj!}{(k+j+1)!}\prod_{t<k}(n-t)
 =1+\frac{(-1)^Kj!}{(K+j+1)!}\prod_{t\le K}(n-t):
\]
for $K=0$ both sides are $1+\frac n{j+1}$, and adding the term $K+1$ changes both sides by
$(-1)^{K+1}j!\,(n+j+1)\prod_{t\le K}(n-t)/(K+j+2)!$. As $n\le K$, the product $\prod_{t\le K}(n-t)$ has
the factor $n-n=0$. Hence $(n+j+1)\,\phi\big(P_{K,j}(\operatorname{ad}E,\operatorname{ad}F)(T)v\big)=\phi(Tv)$.
% Lean: StepScalarProof.stepScalar: H_pow, E_pow, adjE_pow (the induction on k), phi_adjF_pow,
% partial_sum (the induction on K), and the vanishing of prod_{t <= K}(n - t) at t = n.
\end{proof}

\begin{lemma}\label{is-onto-step-exp-cut-scalar}
Let $V$ be a $\mathbf Q$-vector space and $e=(e_{ab})_{1\le a,b\le p+q}$ endomorphisms of $V$ that
satisfy the $\mathfrak{gl}_{p+q}$ relations and are dominant for $\lambda\in\mathbf Z^{p+q}$. Let $v$
be a weight vector of weight $\lambda$, $\phi$ a linear form on $V$ and $T\in\operatorname{End}(V)$,
and $K\geq0$ an integer. Suppose that for $i<j$: $e_{ij}v=0$, $\phi\circ e_{ji}=0$ and
$\lambda_j\le\lambda_i\le\lambda_j+K$; and that $Tx=0$ for every weight vector $x$ of weight $\mu$ with
$\mu_1+\dots+\mu_p<\lambda_1+\dots+\lambda_p$. Then
\[
 \phi\big(G_K(\operatorname{ad}e;p,q)(T)\,v\big)\cdot\prod_{i=1}^{p}\prod_{j=p+1}^{p+q}
 (\lambda_i-\lambda_j+j-i)=\phi(Tv).
\]
\end{lemma}
\leanref{Development/AbelianVarietyTheory/Proofs/DivisorProofs/ProjectorGridScalarOfProjectorScalarStatement.lean}{GlobalTheory.T.projectorGridScalar_of_projectorScalar}
\leanref{Development/AbelianVarietyTheory/Proofs/DivisorProofs/NablaSchur.lean}{GlobalTheory.LeviCut.GridScalar}
\begin{proof}
Write $N=p+q$, $\varepsilon_c$ for the $c$-th unit vector of $\mathbf Z^N$, and $P_{ij}$ for the factors
$P_{K,j-i-1}(\operatorname{ad}e_{ij},\operatorname{ad}e_{ji})$ of the grid projector. For a set
$S\subset\{1,\dots,N\}$ let $\mathcal V(S)$ be the space of the $U\in\operatorname{End}(V)$ with $Ux=0$
for every weight vector $x$ of a weight $\mu$ with $\sum_{c\in S}\mu_c<\sum_{c\in S}\lambda_c$. The
hypothesis on $T$ is $T\in\mathcal V(\{1,\dots,p\})$.

\emph{Weights.} The relation $e_{kk}e_{cd}-e_{cd}e_{kk}=(\delta_{kc}-\delta_{kd})e_{cd}$ shows that if $x$
has weight $\mu$, then $e_{cd}x$ has weight $\mu+\varepsilon_c-\varepsilon_d$, and $e_{cd}^kx$ has weight
$\mu+k(\varepsilon_c-\varepsilon_d)$. Hence:
\begin{enumerate}
\item[(a)] if $c,d\in S$ or $c,d\notin S$, then $\operatorname{ad}e_{cd}$ preserves $\mathcal V(S)$,
  since $(\operatorname{ad}e_{cd})(U)x=e_{cd}Ux-U(e_{cd}x)$ and $e_{cd}x$ has the same sum over $S$ as $x$;
  so $P_{ij}$ preserves $\mathcal V(S)$ if $i,j\in S$ or $i,j\notin S$;
\item[(b)] if $U\in\mathcal V(S)$, $i\in S$ and $j\notin S$, then $U(e_{ji}^kv)=0$ for $k\geq1$, since the
  sum over $S$ of the weight of $e_{ji}^kv$ is $\sum_{c\in S}\lambda_c-k$.
\end{enumerate}

\emph{Dominance.} For $1\le i\le p+1$ and $p<s\le N+1$,
$\mathcal V(\{1,\dots,p\})\subset\mathcal V(\{i,\dots,p\})\subset\mathcal V(\{i,\dots,p\}\cup\{s,\dots,N\})$.
Let $x\ne0$ be a weight vector of weight $\mu$. For the first inclusion, dominance gives
$\sum_{c<i}\mu_c\le\sum_{c<i}\lambda_c$, so if $\sum_{c=i}^p\mu_c<\sum_{c=i}^p\lambda_c$, then
$\sum_{c\le p}\mu_c<\sum_{c\le p}\lambda_c$. For the second, dominance at $s-1$ and the equality of the
total sums give $\sum_{c\geq s}\mu_c\geq\sum_{c\geq s}\lambda_c$, so a smaller sum over
$\{i,\dots,p\}\cup\{s,\dots,N\}$ forces a smaller sum over $\{i,\dots,p\}$.

\emph{One root.} Let $i<j$, and let $U$ satisfy $U(e_{ji}^kv)=0$ for $k\geq1$. Then
$\phi\big(P_{ij}(U)v\big)(\lambda_i-\lambda_j+j-i)=\phi(Uv)$. Indeed, Lemma~\ref{is-onto-step-exp-cut-scalar-root}
applies to $E=e_{ij}$, $F=e_{ji}$, $H=e_{ii}-e_{jj}$, $n=\lambda_i-\lambda_j$ and the index $j-i-1$: the
relations give $EF-FE=H$ and $HF-FH=-2F$, $Ev=0$, $Hv=nv$ with $0\le n\le K$, and $\phi\circ F=0$; and
$n+(j-i-1)+1=\lambda_i-\lambda_j+j-i$.

\emph{One row.} Let $i\le p$ and $U\in\mathcal V(\{i,\dots,p\})$. We show by descending induction on
$s$, for $p<s\le N+1$, that
\[
 \phi\big(P_{is}P_{i,s+1}\cdots P_{iN}(U)\,v\big)\prod_{j=s}^N(\lambda_i-\lambda_j+j-i)=\phi(Uv).
\]
For $s=N+1$ there is nothing to prove. Let $s\le N$, $S=\{i,\dots,p\}\cup\{s+1,\dots,N\}$ and
$W=P_{i,s+1}\cdots P_{iN}(U)$. By the dominance inclusions $U\in\mathcal V(S)$, and the factors $P_{ij}$
with $j>s$ have $i,j\in S$; so $W\in\mathcal V(S)$ by (a). As $i\in S$ and $s\notin S$, (b) gives
$W(e_{si}^kv)=0$ for $k\geq1$, and the case of one root gives
$\phi(P_{is}(W)v)(\lambda_i-\lambda_s+s-i)=\phi(Wv)$. With the induction hypothesis this is the claim at
$s$. At $s=p+1$ it reads $\phi(R_i(U)v)\prod_{j>p}(\lambda_i-\lambda_j+j-i)=\phi(Uv)$.

\emph{The grid.} For $0\le t\le p$, $R_t\cdots R_1(T)\in\mathcal V(\{i,\dots,p\})$ whenever $t<i\le p+1$: for
$t=0$ by the dominance inclusions, and the factors $P_{t+1,j}$ of $R_{t+1}$ have $t+1,j\notin\{i,\dots,p\}$
when $i>t+1$, so (a) gives the induction step. Hence, by induction on $t$, using the row formula for
$i=t+1$ and $U=R_t\cdots R_1(T)\in\mathcal V(\{t+1,\dots,p\})$,
\[
 \phi\big(R_t\cdots R_1(T)\,v\big)\prod_{i=1}^t\prod_{j=p+1}^N(\lambda_i-\lambda_j+j-i)=\phi(Tv).
\]
At $t=p$ this is the claim.
% Lean: gridScalar_of_stepScalar = GridScalarProof.grid_scalar: vanSub (the spaces V(S)), isWt_shift and
% isWt_pow (weights), adj_mem_vanSub, rootP_mem and rowP_mem ((a)), kill_pow ((b)), van_Ico_of_range and
% van_union (dominance), step (one root, from StepScalar with H = e_ii - e_jj and n = (lam i - lam j).toNat),
% row_scalar (one row), grid_van and grid_scalar (the grid).
\end{proof}

\begin{lemma}\label{is-onto-step-exp-cut-integrality}
Let $V$ be a $\mathbf Q$-vector space, $e=(e_{ab})_{1\le a,b\le p+q}$ endomorphisms of $V$, and $L,M$
$\mathbf Z$-submodules of $V$ stable under divided powers. Let $K,w,\Lambda\geq0$ be integers with
$K+p+q\le w+1$ such that $(j+k+1)!$ divides $\Lambda\,j!\,k!$ whenever $j+k+1\le w$. If $c\geq0$ and
$T\in\operatorname{End}(V)$ satisfy $cT(L)\subset M$, then
$c\Lambda^{pq}\,G_K(\operatorname{ad}e;p,q)(T)(L)\subset M$.
\end{lemma}
\leanref{Development/AbelianVarietyTheory/Proofs/DivisorProofs/ProjectorGridDenominatorBoundStatement.lean}{GlobalTheory.T.projector_grid_denominator_bound}
\leanref{Development/AbelianVarietyTheory/Proofs/DivisorProofs/NablaSchur.lean}{GlobalTheory.LeviCut.GridInt}
\begin{proof}
Call $U\in\operatorname{End}(V)$ \emph{$c$-integral} if $cU(L)\subset M$, and write
$x^{(k)}=\frac1{k!}x^k$ for an endomorphism $x$ of $V$.

\emph{One root.} Left and right multiplication by $x$ commute, so the binomial formula gives
$(\operatorname{ad}x)^kU=\sum_{t=0}^k(-1)^{k-t}\binom kt\,x^tUx^{k-t}$, that is,
\[
 \tfrac1{k!}(\operatorname{ad}x)^kU=\sum_{t=0}^k(-1)^{k-t}\,x^{(t)}\,U\,x^{(k-t)}.
\]
Hence, if $L$ and $M$ are stable under the $x^{(t)}$ and $U$ is $c$-integral, so is
$\frac1{k!}(\operatorname{ad}x)^kU$. Now let $X,Y$ be endomorphisms of $V$ under whose divided powers $L$
and $M$ are stable, and let $j_0\geq0$ be such that $(j_0+k+1)!$ divides $\Lambda\,j_0!\,k!$ for
$k\le K$, say $\Lambda\,j_0!\,k!=m_k(j_0+k+1)!$. The $k$-th term of
$P_{K,j_0}(\operatorname{ad}X,\operatorname{ad}Y)(U)$ is
\[
 \frac{(-1)^kj_0!\,k!}{(k+j_0+1)!}\cdot\tfrac1{k!}(\operatorname{ad}Y)^k\big(\tfrac1{k!}(\operatorname{ad}X)^kU\big),
\]
so $\Lambda$ times it is $(-1)^km_k$ times an endomorphism that is $c$-integral by the preceding
remark, applied to $X$ and then to $Y$. Hence, if $U$ is $c$-integral,
$P_{K,j_0}(\operatorname{ad}X,\operatorname{ad}Y)(U)$ is $c\Lambda$-integral.

\emph{The grid.} Let $i\le p<j\le p+q$. Then $i\ne j$, so $L$ and $M$ are stable under the divided
powers of $e_{ij}$ and of $e_{ji}$; and for $k\le K$, $(j-i-1)+k+1\le p+q-1+K\le w$, so the
divisibility holds for $j_0=j-i-1$. By the preceding paragraph the factor $P_{ij}$ of the grid
projector carries $c$-integral endomorphisms to $c\Lambda$-integral ones, for every $c$. The factor
$R_i$ is a product of $q$ such factors, so it carries them to $c\Lambda^q$-integral ones, and
$G_K(\operatorname{ad}e;p,q)=R_p\cdots R_1$ to $c\Lambda^{pq}$-integral ones. Applied to $T$ this is
the claim.
% Lean: gridInt_holds: GridIntProof.adj_pow_apply and intOn_divAdj (one root), intOn_proj (the
% projector, with m from the divisibility), hstep for rootP i j (j0 = j - i - 1, with the bound
% K + (a + b) <= w + 1), intOn_listProd for rowP (Lambda^b) and for gridP (Lambda^(a b)).
\end{proof}

\begin{lemma}\label{is-onto-step-exp-cut-module}
Let $N\geq1$ and $\nu\in\mathbf N^N$ nonincreasing, and let $E=(E_{ab})$ be the row operators on
$\nabla_{\mathbf Q}(\nu)$ (Definition~\ref{is-onto-step-exp-cut-grid}(6)).
\begin{enumerate}
\item $E$ satisfies the $\mathfrak{gl}_N$ relations, its weight vectors span $\nabla_{\mathbf Q}(\nu)$,
  it is dominant for $\nu$, and $\nabla_{\mathbf Z}(\nu)$ and $\Delta_{\mathbf Z}(\nu)$ are stable
  under divided powers.
\item If $\nu_i\le\nu_j+K$ for all $i,j$, the weights of $E$ are bounded by $K$.
\item If an endomorphism $S$ of $\nabla_{\mathbf Q}(\nu)$ commutes with every $E_{cd}$ and
  $[h_\nu,Sh_\nu]=[h_\nu,h_\nu]$ for the Fock pairing, then $S=1$.
\end{enumerate}
\end{lemma}
\leanref{Development/AbelianVarietyTheory/Proofs/DivisorProofs/DualWeylModuleStructureGlProofStatement.lean}{GlobalTheory.T.dualWeylModuleStructureGl_proof}
\leanref{Development/AbelianVarietyTheory/Proofs/DivisorProofs/DualWeylWeightSpreadBoundProofStatement.lean}{GlobalTheory.T.dualWeylWeightSpreadBound_proof}
\leanref{Development/AbelianVarietyTheory/Proofs/DivisorProofs/DualWeylSchurEndomorphismStatement.lean}{GlobalTheory.T.dualWeyl_schur_endomorphism}
\begin{proof}
(1) The relations hold for the derivations $E_{ab}$ of $\mathbf Q[X]$. A monomial is a weight vector,
of weight its vector of row degrees. The operators $E_{cc}$ commute and preserve
$\nabla_{\mathbf Q}(\nu)$, so the components of an element of $\nabla_{\mathbf Q}(\nu)$ of given row
degrees lie in $\nabla_{\mathbf Q}(\nu)$, and the weight vectors span. A nonzero weight vector of weight
$\mu$ has a monomial of row degrees $\mu$; by the proof of Lemma~\ref{pd-gain-holds-hecke-rows}(1) its
degree in the first $k$ rows is at most $\nu_1+\dots+\nu_k$, and its total degree is
$\nu_1+\dots+\nu_N$. For $a\ne b$, $\frac1{k!}E_{ab}^k$ is the operator $E^{(k)}_{ab}$ of
Definition~\ref{pd-weight-lattices}(1), by Taylor's formula; it preserves $\Delta_{\mathbf Z}(\nu)$ by
definition and $\nabla_{\mathbf Z}(\nu)$ because it acts on the rows with integer coefficients.

(2) Permuting the rows preserves $\nabla_{\mathbf Q}(\nu)$, so by the bound in (1) every row of a
monomial of an element of $\nabla_{\mathbf Q}(\nu)$ has degree at most $\nu_1$, and, as the other rows
have total degree at most $\nu_1+\dots+\nu_{N-1}$, at least $\nu_N$. If $x\ne0$ and
$E_{aa}x-E_{bb}x=nx$, comparing the coefficients of a monomial of $x$ shows that $n$ is its degree in
row $a$ minus its degree in row $b$, so $n\le\nu_1-\nu_N\le K$.

(3) The polynomial $h_\nu$ is killed by the $E_{ab}$ with $a<b$, since each leading principal minor
$\delta_s$ is (it has two equal rows if $b\le s$, and does not involve row $b$ if $b>s$). So $Sh_\nu$ is
killed by them. In characteristic $0$, $\nabla_{\mathbf Q}(\nu)$ is the irreducible representation of
highest weight $\nu$, whose vectors killed by the $E_{ab}$, $a<b$, form the line of $h_\nu$
\cite[II.2]{Jan03}; so
$Sh_\nu=ch_\nu$, and $c=1$ as $[h_\nu,h_\nu]\ne0$. The kernel of $S-1$ is stable under the $E_{ab}$, hence
under the $E^{(k)}_{ab}$, and contains $h_\nu$; so it contains $\Delta_{\mathbf Q}(\nu)$, which is
$\nabla_{\mathbf Q}(\nu)$ by Lemma~\ref{is-standard-basis}(2). Hence $S=1$.
% Lean: (1) is LeviCut.nablaGl (glRel_eV, wtSpan_eV, dom_eV with DualWeyl.rowDeg_dominated_of_mem_nabla,
% divStable_eV with rowDivStable and weylRowDivStable); (2) is LeviCut.nablaBdd (nablaBdd_rowDeg_bounds
% with dualWeylRowStable for the row permutations); (3) is LeviCut.nablaSchur, with der_hw_raise,
% GradedWeylDivisibility.hwLine (the highest weight line), LR.hw_norm_pos, LR.closed_of_der and
% GradedWeylDivisibility.closedEqNabla (a subspace of nabla_Q stable under the divided powers and
% containing h_nu is all of it).
\end{proof}

\begin{lemma}\label{is-onto-step-exp-cut-levi-duality}
Let $N\geq0$ and $\nu\in\mathbf N^N$ nonincreasing, and let $[\ ,\ ]$ be the Fock pairing
(Definition~\ref{pd-weight-lattices}(2)).
\begin{enumerate}
\item $\nabla_{\mathbf Z}(\nu)$ is the set of the $Q\in\nabla_{\mathbf Q}(\nu)$ with
  $[P,Q]\in\mathbf Z\cdot[h_\nu,h_\nu]$ for every $P\in\Delta_{\mathbf Z}(\nu)$.
\item $\Delta_{\mathbf Z}(\nu)$ is the set of the $P\in\nabla_{\mathbf Q}(\nu)$ with
  $[P,Q]\in\mathbf Z\cdot[h_\nu,h_\nu]$ for every $Q\in\nabla_{\mathbf Z}(\nu)$.
\end{enumerate}
\end{lemma}
\leanref{Development/AbelianVarietyTheory/Proofs/FoundationProofs/StandardWeylLatticesFockDual.lean}{GlobalTheory.StdDual.weylEqDualNabla}
\leanref{Development/AbelianVarietyTheory/Proofs/FoundationProofs/ApplyWordMemDualWeyl.lean}{GlobalTheory.WeylDuality.nablaEqDualWeylOf}
\leanref{Development/AbelianVarietyTheory/Proofs/FoundationProofs/ApplyWordMemDualWeyl.lean}{GlobalTheory.WeylDuality.weylEqDualNablaOf}
\begin{proof}
If $N=0$, the three modules are $\mathbf Q$, $\mathbf Z$ and $\mathbf Z$, and $h_\nu=1$; let $N\geq1$.
For a composite $u$ of operators $E^{(k)}_{ab}$ let $u^t$ be the composite of the corresponding
$E^{(k)}_{ba}$ in the reverse order; by Lemma~\ref{pd-kappa-gram-fock}(2), $[uf,f']=[f,u^tf']$. The
$E^{(k)}_{ab}$ preserve $\nabla_{\mathbf Z}(\nu)$ and $\nabla_{\mathbf Q}(\nu)$, and
$h_\nu\in\nabla_{\mathbf Z}(\nu)$, as in the proof of Lemma~\ref{pd-gain-weyl-comparison}; so
$\Delta_{\mathbf Z}(\nu)\subset\nabla_{\mathbf Z}(\nu)$, and $\Delta_{\mathbf Z}(\nu)$ is the
$\mathbf Z$-span of the $uh_\nu$.

(1) Let $Q\in\nabla_{\mathbf Z}(\nu)$. For $P=uh_\nu$, $[P,Q]=[h_\nu,u^tQ]$ with
$u^tQ\in\nabla_{\mathbf Z}(\nu)$, so Lemma~\ref{pd-kappa-gram-fock}(4) gives
$[P,Q]=(u^tQ)(1)\,[h_\nu,h_\nu]$, and $(u^tQ)(1)\in\mathbf Z$. By linearity in $P$, $Q$ lies in the set
of (1). Conversely, let $Q$ lie in that set, and let $Z^e$ be a monomial with $e_{as}=0$ for $a<s$. If
$\sum_ae_{as}\ne\nu_s$ for some $s$, the coefficient $Q_e$ of $Z^e$ is $0$, as $Q$ has degree $\nu_s$ in
the $s$-th column. Otherwise, with $w_e$ as in Lemma~\ref{pd-kappa-gram-fock}(5), parts (5), (4) and (2)
of that lemma give $Q_e=(w_eQ)(1)=[h_\nu,w_eQ]/[h_\nu,h_\nu]=[w_e^th_\nu,Q]/[h_\nu,h_\nu]$, which lies in
$\mathbf Z$ since $w_e^th_\nu\in\Delta_{\mathbf Z}(\nu)$. So $Q\in\nabla_{\mathbf Z}(\nu)$ by
Lemma~\ref{pd-kappa-gram-fock}(6) for $A=\mathbf Z$.

(2) Put $B(x,y)=[x,y]/[h_\nu,h_\nu]$ on $\nabla_{\mathbf Q}(\nu)$. It is symmetric, and nondegenerate,
since $[f,f]=\sum_\kappa\kappa!\,f_\kappa^2>0$ for $f\ne0$. For a $\mathbf Z$-submodule $M$ let $M^*$ be
the set of $y\in\nabla_{\mathbf Q}(\nu)$ with $B(x,y)\in\mathbf Z$ for all $x\in M$. Then (1) says
$\nabla_{\mathbf Z}(\nu)=\Delta_{\mathbf Z}(\nu)^*$, and (2) says
$\Delta_{\mathbf Z}(\nu)=\nabla_{\mathbf Z}(\nu)^*$. The module $\Delta_{\mathbf Z}(\nu)$ is a lattice in
$\nabla_{\mathbf Q}(\nu)$: it is finitely generated, as a submodule of $\nabla_{\mathbf Z}(\nu)$, which is
free of finite rank by Lemma~\ref{pd-kappa-defects-restriction-standard}(2); and it spans
$\nabla_{\mathbf Q}(\nu)$, which is the $\mathbf Q$-span of the $uh_\nu$ by
Lemma~\ref{is-standard-basis}(2). For a lattice $M$ with a basis $(b_i)$, $M^*$ is the lattice with the dual
basis for $B$, so $M^{**}=M$. Hence
$\nabla_{\mathbf Z}(\nu)^*=\Delta_{\mathbf Z}(\nu)^{**}=\Delta_{\mathbf Z}(\nu)$.
% Lean: (1) is WeylDuality.nablaEqDualWeylOf (NablaEqDualWeyl, bDual) with CoeffWord.coeffWord (part (5) of
% Lemma pd-kappa-gram-fock), FockHighestWeight.fockQ_hw (part (4)), FockAdj.fockQ_word_rowAct_adj (part
% (2)) and WeylInt.lowerTriIntegral (part (6)); (2) is weylEqDualNablaOf: inV_bDual, normB_nondeg,
% normB_symm, weyl_isLattice (dualWeyl_int_fg from sSBasis, weylLeDual, weylEqNablaQ) and
% dual_dual_of_isLattice.
\end{proof}

\begin{lemma}\label{is-onto-step-exp-cut-levi-substitution}
Let $R$ be a commutative ring, $0\le k\le N$ integers, $\tau\colon\{1,\dots,k\}\to\{1,\dots,N\}$ strictly
increasing, $X$ an $N\times N$ and $Z$ a $k\times k$ matrix of variables, and $\varphi\colon R[X]\to R[Z]$
the $R$-algebra homomorphism with $\varphi(X_{xy})=\psi(x,y)$, where $\psi(\tau a,\tau s)=Z_{as}$,
$\psi(x,\tau s)=0$ for $x\notin\tau(\{1,\dots,k\})$, and $\psi(x,y)\in R$ for
$y\notin\tau(\{1,\dots,k\})$. Let $\nu\in\mathbf Z_{\geq0}^N$, and let $\nabla_R(\nu)$ and
$\nabla_R(\nu\circ\tau)$ be as in Lemma~\ref{pd-kappa-defects-restriction-polynomials}. Then
$\varphi(\nabla_R(\nu))\subset\nabla_R(\nu\circ\tau)$.
\end{lemma}
\leanref{Development/AbelianVarietyTheory/Proofs/DivisorProofs/DualWeylLatticeSubstIncreasingMapStableStatement.lean}{GlobalTheory.T.dualWeylLattice_subst_increasing_map_stable}
\leanref{Development/AbelianVarietyTheory/Proofs/DivisorProofs/NablaSchur.lean}{GlobalTheory.LeviCut.SubstMem}
\begin{proof}
Let $f\in\nabla_R(\nu)$. If $y=\tau s$, $\psi(x,y)$ is $Z_{as}$ or $0$, homogeneous of degree $1$ in the
$s$-th column and $0$ in the others; if $y\notin\tau(\{1,\dots,k\})$, $\psi(x,y)$ is a constant. So a
monomial of column degrees $\nu$ goes to a polynomial of column degrees $\nu\circ\tau$, and so does $f$.
For $l<j\le k$ let $C_{lj}$ be the substitution $X_{x,\tau j}\mapsto X_{x,\tau j}+yX_{x,\tau l}$ on
$R[X,y]$, which is $f(X)\mapsto f(X(1+yE_{\tau l,\tau j}))$, and $C'_{lj}$ the substitution
$Z_{aj}\mapsto Z_{aj}+yZ_{al}$ on $R[Z,y]$. Then $\varphi C_{lj}=C'_{lj}\varphi$, with $\varphi(y)=y$, as one
checks on the variables: $X_{\tau a,\tau j}+yX_{\tau a,\tau l}\mapsto Z_{aj}+yZ_{al}$;
$X_{x,\tau j}+yX_{x,\tau l}\mapsto0$ for $x$ not in the image; and the other variables are fixed by
$C_{lj}$ and go to elements fixed by $C'_{lj}$. Since $\tau l<\tau j$, $C_{lj}f=f$, so
$C'_{lj}\varphi(f)=\varphi(f)$.
% Lean: LeviCut.substMem: colPull and aeval_isWeightedHomogeneous for the degrees; psiO, hA and hB (the
% two composites agree on the variables) for the invariance.
\end{proof}

\begin{lemma}\label{is-onto-step-exp-cut-levi-restriction}
Let $p,q\geq0$, $\nu\in\mathbf N^{p+q}$ nonincreasing, $\nu'=(\nu_1,\dots,\nu_p)$ and
$\nu''=(\nu_{p+1},\dots,\nu_{p+q})$. Let $X$, $X'$ and $X''$ be matrices of variables of sizes $p+q$, $p$
and $q$, and $\rho\colon\mathbf Q[X]\to\mathbf Q[X']\otimes\mathbf Q[X'']$ the ring homomorphism with
$\rho(X_{as})=X'_{as}\otimes1$ for $a,s\le p$, $\rho(X_{p+b,p+t})=1\otimes X''_{bt}$, and $\rho(X_{xy})=0$
for the other variables (the restriction to block diagonal matrices). Write $E$, $E'$ and $E''$ for the
row operators $E_{ab}=\sum_cx_{ac}\partial/\partial x_{bc}$ on $\mathbf Q[X]$, $\mathbf Q[X']$ and
$\mathbf Q[X'']$.
\begin{enumerate}
\item $\rho(\nabla_{\mathbf Z}(\nu))$ lies in the $\mathbf Z$-span of the $f_1\otimes f_2$ with
  $f_1\in\nabla_{\mathbf Z}(\nu')$ and $f_2\in\nabla_{\mathbf Z}(\nu'')$.
\item $\rho(h_\nu)=h_{\nu'}\otimes h_{\nu''}$.
\item $\rho\circ E_{cd}=(E'_{cd}\otimes1)\circ\rho$ for $c,d\le p$, and
  $\rho\circ E_{p+c,p+d}=(1\otimes E''_{cd})\circ\rho$ for $c,d\le q$.
\item If $f\in\mathbf Q[X]$ has degree $\nu_s$ in the $s$-th column for every $s$, and
  $\sum_{c\le p}E_{cc}f=tf$ with $t\ne\nu_1+\dots+\nu_p$, then $\rho(f)=0$.
\end{enumerate}
\end{lemma}
\leanref{Development/AbelianVarietyTheory/Proofs/DivisorProofs/NablaSchur.lean}{GlobalTheory.LeviCut.resT}
\leanref{Development/AbelianVarietyTheory/Proofs/DivisorProofs/DualWeylRestrictionLatticeTensorOfSubstitutionStStatement.lean}{GlobalTheory.T.dualWeylRestrictionLatticeTensor_of_substitutionStable}
\leanref{Development/AbelianVarietyTheory/Proofs/DivisorProofs/DualWeylHighestWeightRestrictionBlockDiagonalStatement.lean}{GlobalTheory.T.dualWeyl_highest_weight_restriction_blockDiagonal}
\leanref{Development/AbelianVarietyTheory/Proofs/DivisorProofs/DualWeylRestrictionRowOperatorsStatement.lean}{GlobalTheory.T.dualWeyl_restriction_rowOperators}
\leanref{Development/AbelianVarietyTheory/Proofs/DivisorProofs/NablaSchur.lean}{GlobalTheory.LeviCut.ResLat}
\leanref{Development/AbelianVarietyTheory/Proofs/DivisorProofs/NablaSchur.lean}{GlobalTheory.LeviCut.ResHw}
\leanref{Development/AbelianVarietyTheory/Proofs/DivisorProofs/NablaSchur.lean}{GlobalTheory.LeviCut.ResDer}
\begin{proof}
(1) Let $\rho_{\mathbf Z}\colon\mathbf Z[X]\to\mathbf Z[X']\otimes\mathbf Z[X'']$ be the same map over
$\mathbf Z$, so that $\rho$ on $\mathbf Z[X]$ is $\rho_{\mathbf Z}$ followed by the map to
$\mathbf Q[X']\otimes\mathbf Q[X'']$, and let $f\in\nabla_{\mathbf Z}(\nu)$. Put $A=\mathbf Z[X']$; then $f$
lies in $\nabla_A(\nu)$, since the defining conditions survive the extension of scalars. Apply
Lemma~\ref{is-onto-step-exp-cut-levi-substitution} over $A$ with $k=q$, $\tau(s)=p+s$ and $\psi$ sending
$X_{as}$ to the constant $X'_{as}$ for $a,s\le p$, $X_{p+b,p+t}$ to $Z_{bt}$, and the other variables to
$0$: it gives $\varphi(f)\in\nabla_A(\nu'')$. By
Lemma~\ref{pd-kappa-defects-restriction-standard}(1), $\nabla_A(\nu'')$ is the $A$-span of the
bideterminants of shape $\nu''$, which have integer coefficients and lie in $\nabla_{\mathbf Z}(\nu'')$.
Identifying $A[Z]$ with $\mathbf Z[X']\otimes\mathbf Z[X'']$, $\varphi(f)=\rho_{\mathbf Z}(f)$, which
therefore lies in the image of $\mathbf Z[X']\otimes\nabla_{\mathbf Z}(\nu'')$. In the same way, with the
roles of the blocks exchanged, it lies in the image of $\nabla_{\mathbf Z}(\nu')\otimes\mathbf Z[X'']$.
The quotient $\mathbf Z[X'']/\nabla_{\mathbf Z}(\nu'')$ is torsion-free, as the conditions defining
$\nabla_{\mathbf Z}(\nu'')$ are linear and $\mathbf Z[X'']$ is torsion-free, hence flat over $\mathbf Z$.
For submodules $N_1\subset M_1$ and $N_2\subset M_2$ with $M_2/N_2$ flat, an element of $M_1\otimes M_2$
that is the image of $w\in N_1\otimes M_2$ and also lies in the image of $M_1\otimes N_2$ is the image of
an element of $N_1\otimes N_2$: the image of $w$ in $N_1\otimes M_2/N_2$ maps to $0$ in
$M_1\otimes M_2/N_2$, injectively by flatness, so it is $0$, and $w$ comes from $N_1\otimes N_2$ by right
exactness. Hence $\rho_{\mathbf Z}(f)$ is a $\mathbf Z$-combination of $f_1\otimes f_2$ as in (1), and so
is $\rho(f)$.

(2) Write $\delta_s$, $\delta'_s$, $\delta''_s$ for the leading principal minors, so that
$h_\nu=\prod_{s=1}^{p+q}\delta_s^{\nu_s-\nu_{s+1}}$ with $\nu_{p+q+1}=0$. The map $\rho$ sends the leading
$s\times s$ submatrix of $X$ to a block diagonal matrix, so $\rho(\delta_s)=\delta'_s\otimes1$ for $s\le p$
and $\rho(\delta_{p+u})=\delta'_p\otimes\delta''_u$ for $1\le u\le q$. Hence $\rho(h_\nu)$ is
$\prod_{s<p}\delta'^{\,\nu_s-\nu_{s+1}}_s\cdot\delta'^{\,e}_p\otimes\prod_{u=1}^q
\delta''^{\,\nu_{p+u}-\nu_{p+u+1}}_u$ with $e=(\nu_p-\nu_{p+1})+\sum_{u=1}^q(\nu_{p+u}-\nu_{p+u+1})=\nu_p$,
which is $h_{\nu'}\otimes h_{\nu''}$.

(3) Both sides are derivations along $\rho$, so it suffices to compare them on the variables, where
$E_{cd}X_{xy}=\delta_{xd}X_{cy}$. Let $c,d\le p$. If $x,y\le p$ both sides give
$\delta_{xd}X'_{cy}\otimes1$; if $x\le p<y$ both give $0$; and if $x>p$, $\delta_{xd}=0$ and $\rho(X_{xy})$
is $0$ or of the form $1\otimes X''_{bt}$, which $E'_{cd}\otimes1$ kills. The second block is the same.

(4) Let $m$ be a monomial of $f$ with $\rho(m)\ne0$. Then $m$ involves no variable $X_{xy}$ with exactly
one of $x,y$ at most $p$, so its degree in the first $p$ rows equals its degree in the first $p$
columns, which is $\nu_1+\dots+\nu_p$. As $E_{cc}$ multiplies a monomial by its degree in the $c$-th row,
comparing the coefficients of $m$ in $\sum_{c\le p}E_{cc}f=tf$ gives $t=\nu_1+\dots+\nu_p$, a
contradiction. So $\rho$ kills every monomial of $f$.
% Lean: (1) resLat_of_substMem: resT_mem_lTensor and resT_mem_rTensor (psiR and psiL with SubstMem,
% standardMonomialSpanDomain, bidet_map', phiR_psiR), flat_quot_dualWeyl (dualWeylSat), mem_range_map_of_flat,
% toQ_resT; (2) resHw_holds (resT_lpMinor_le, resT_lpMinor_ge, hw_first); (3) dualWeyl_resDer with
% resT_der_left_X and resT_der_right_X; (4) resT_wrong_degree (resT_monomial_eq_zero, coeff_der_diag).
\end{proof}

\begin{lemma}\label{is-onto-step-exp-cut-levi}
Let $\nu\in\mathbf N^{p+q}$ be nonincreasing, $\nu'=(\nu_1,\dots,\nu_p)$ and
$\nu''=(\nu_{p+1},\dots,\nu_{p+q})$, and let $c_1,c_2\geq0$ be integers with
$c_1\nabla_{\mathbf Z}(\nu')\subset\Delta_{\mathbf Z}(\nu')$ and
$c_2\nabla_{\mathbf Z}(\nu'')\subset\Delta_{\mathbf Z}(\nu'')$. Then there is an endomorphism $T$ of
$\nabla_{\mathbf Q}(\nu)$ such that, with the row operators $E_{cd}$ and the Fock pairing:
\begin{enumerate}
\item $T$ commutes with $E_{cd}$ whenever $c,d\le p$ or $c,d>p$;
\item $c_1c_2T(\nabla_{\mathbf Z}(\nu))\subset\Delta_{\mathbf Z}(\nu)$;
\item $Tx=0$ for every weight vector $x$ of weight $\mu$ with $\mu_1+\dots+\mu_p<\nu_1+\dots+\nu_p$;
\item $[h_\nu,Th_\nu]=[h_\nu,h_\nu]$.
\end{enumerate}
\end{lemma}
\leanref{Development/AbelianVarietyTheory/Proofs/DivisorProofs/DualWeylLeviProjectionDenominatorsOfHypothesesStatement.lean}{GlobalTheory.T.dualWeylLeviProjectionDenominators_of_hypotheses}
\leanref{Development/AbelianVarietyTheory/Proofs/DivisorProofs/NablaSchur.lean}{GlobalTheory.LeviCut.LeviProj}
\begin{proof}
Let $\rho$, $E'$ and $E''$ be as in Lemma~\ref{is-onto-step-exp-cut-levi-restriction}, and give
$\mathbf Q[X']\otimes\mathbf Q[X'']$ the pairing $[a\otimes b,a'\otimes b']=[a,a'][b,b']$, the tensor
product of the Fock pairings. For $f\ne0$, $[f,f]=\sum_\kappa\kappa!\,f_\kappa^2>0$; so the Fock pairing
is nondegenerate on the finite-dimensional space $\nabla_{\mathbf Q}(\nu)$, and
$[h_{\nu'},h_{\nu'}]$ and $[h_{\nu''},h_{\nu''}]$ are nonzero. Put
\[
 \kappa=\frac{[h_\nu,h_\nu]}{[h_{\nu'},h_{\nu'}]\,[h_{\nu''},h_{\nu''}]},\qquad
 \beta(x,y)=\kappa\,[\rho x,\rho y]\quad(x,y\in\nabla_{\mathbf Q}(\nu)),
\]
and let $T$ be the endomorphism of $\nabla_{\mathbf Q}(\nu)$ with $[x,Ty]=\beta(x,y)$ for all $x,y$.

(4) By Lemma~\ref{is-onto-step-exp-cut-levi-restriction}(2),
$[h_\nu,Th_\nu]=\kappa[h_{\nu'},h_{\nu'}][h_{\nu''},h_{\nu''}]=[h_\nu,h_\nu]$.

(3) Let $x$ be a weight vector of weight $\mu$ with $\mu_1+\dots+\mu_p<\nu_1+\dots+\nu_p$. It has degree
$\nu_s$ in the $s$-th column, and $\sum_{c\le p}E_{cc}x=(\mu_1+\dots+\mu_p)x$; so $\rho x=0$ by
Lemma~\ref{is-onto-step-exp-cut-levi-restriction}(4). Then $[y,Tx]=0$ for all $y$, and $Tx=0$.

(1) Let $c,d\le p$, and $x,y\in\nabla_{\mathbf Q}(\nu)$. By Lemma~\ref{pd-kappa-gram-fock}(2) for $k=1$
($E^{(1)}_{ab}=E_{ab}$), $[x,E_{cd}z]=[E_{dc}x,z]$, and the same holds for $E'$; on the tensor product,
$[(E'_{dc}\otimes1)u,v]=[u,(E'_{cd}\otimes1)v]$, by bilinearity from the pure tensors. With
Lemma~\ref{is-onto-step-exp-cut-levi-restriction}(3),
\[
 [x,E_{cd}Ty]=[E_{dc}x,Ty]=\kappa[(E'_{dc}\otimes1)\rho x,\rho y]=\kappa[\rho x,(E'_{cd}\otimes1)\rho y]
 =\kappa[\rho x,\rho E_{cd}y]=[x,TE_{cd}y].
\]
So $E_{cd}T=TE_{cd}$. The case $c,d>p$ is the same with $E''$.

(2) Let $x\in\nabla_{\mathbf Z}(\nu)$. By Lemma~\ref{is-onto-step-exp-cut-levi-duality}(2) it suffices
that $[g,c_1c_2Tx]\in\mathbf Z\cdot[h_\nu,h_\nu]$ for every $g\in\nabla_{\mathbf Z}(\nu)$, and
$[g,c_1c_2Tx]=c_1c_2\kappa[\rho g,\rho x]$. By Lemma~\ref{is-onto-step-exp-cut-levi-restriction}(1),
$\rho g$ and $\rho x$ are $\mathbf Z$-combinations of tensors $f_1\otimes f_2$ with
$f_1\in\nabla_{\mathbf Z}(\nu')$ and $f_2\in\nabla_{\mathbf Z}(\nu'')$. For two such tensors,
$c_1c_2[f_1\otimes f_2,f'_1\otimes f'_2]=[f_1,c_1f'_1]\,[f_2,c_2f'_2]$; as $c_1f'_1\in\Delta_{\mathbf Z}(\nu')$
and $c_2f'_2\in\Delta_{\mathbf Z}(\nu'')$, Lemma~\ref{is-onto-step-exp-cut-levi-duality}(2) for $\nu'$
and $\nu''$ puts the two factors in $\mathbf Z\cdot[h_{\nu'},h_{\nu'}]$ and
$\mathbf Z\cdot[h_{\nu''},h_{\nu''}]$. Hence $c_1c_2\kappa[\rho g,\rho x]\in\mathbf Z\cdot[h_\nu,h_\nu]$.
% Lean: leviProj_of with lpT (the endomorphism, lpT_spec), lpBeta and lpK (beta and kappa), lpTen (the
% tensor pairing), lpF_nondeg; (4) lp_hw (ResHw); (3) lp_vanish (ResDer, third part); (1) lp_comm with
% lp_comm_fin, lpTen_rTensor, lpTen_lTensor and LR.fock_der (ResDer, first two parts); (2) lp_lat with
% lp_ten_int and lp_blk (ResLat and StdDual.weylEqDualNabla).
\end{proof}

\begin{lemma}\label{is-onto-step-exp-cut-blocks}
Let $\nu\in\mathbf N^{p+q}$ be nonincreasing, with $\nu'$ and $\nu''$ as in
Lemma~\ref{is-onto-step-exp-cut-levi}, and let $c_1,c_2,K,w,\Lambda\geq0$ be integers with
$c_1\nabla_{\mathbf Z}(\nu')\subset\Delta_{\mathbf Z}(\nu')$,
$c_2\nabla_{\mathbf Z}(\nu'')\subset\Delta_{\mathbf Z}(\nu'')$, $\nu_i\le\nu_j+K$ for all $i,j$,
$K+p+q\le w+1$, and $(j+k+1)!$ dividing $\Lambda\,j!\,k!$ whenever $j+k+1\le w$. Then
$c_1c_2\Lambda^{pq}\nabla_{\mathbf Z}(\nu)\subset\Delta_{\mathbf Z}(\nu)$.
\end{lemma}
\leanref{Development/AbelianVarietyTheory/Proofs/DivisorProofs/WeylLatticeExactExponentBlockCut.lean}{GlobalTheory.LeviCut.cut_incl}
\begin{proof}
If $p+q=0$, both lattices are $\mathbf Z$ and there is nothing to prove; let $p+q\geq1$. Let $T$ be the
endomorphism of Lemma~\ref{is-onto-step-exp-cut-levi}, $E$ the row operators on
$\nabla_{\mathbf Q}(\nu)$, and $S_0=G_K(\operatorname{ad}E;p,q)(T)$. By
Lemma~\ref{is-onto-step-exp-cut-module}(1) and~(2) and Lemma~\ref{is-onto-step-exp-cut-adjoint},
$\operatorname{ad}E$ satisfies the $\mathfrak{gl}_{p+q}$ relations with weights bounded by $2K$. By
Lemma~\ref{is-onto-step-exp-cut-levi}(1), $\operatorname{ad}(E_{cd})T=0$ when $c,d\le p$ or $c,d>p$, so
Lemma~\ref{is-onto-step-exp-cut-invariance} gives $\operatorname{ad}(E_{cd})S_0=0$: $S_0$ commutes
with every $E_{cd}$.

Apply Lemma~\ref{is-onto-step-exp-cut-scalar} with $v=h_\nu$, $\lambda=\nu$ and $\phi=[h_\nu,\cdot\,]$.
The polynomial $h_\nu$ has degree $\nu_c$ in the $c$-th row, so it is a weight vector of weight $\nu$;
for $i<j$, $E_{ij}h_\nu=0$ (proof of Lemma~\ref{is-onto-step-exp-cut-module}(3)),
$[h_\nu,E_{ji}x]=[E_{ij}h_\nu,x]=0$ by Lemma~\ref{pd-kappa-gram-fock}(2), and
$\nu_j\le\nu_i\le\nu_j+K$. The family $E$ is dominant for $\nu$ by
Lemma~\ref{is-onto-step-exp-cut-module}(1), and $T$ kills the weight vectors of
Lemma~\ref{is-onto-step-exp-cut-levi}(3). With $z=\prod_{i\le p<j}(\nu_i-\nu_j+j-i)$, an integer, this
and Lemma~\ref{is-onto-step-exp-cut-levi}(4) give $z[h_\nu,S_0h_\nu]=[h_\nu,Th_\nu]=[h_\nu,h_\nu]$. So
$zS_0$ commutes with every $E_{cd}$ and $[h_\nu,zS_0h_\nu]=[h_\nu,h_\nu]$, and
Lemma~\ref{is-onto-step-exp-cut-module}(3) gives $zS_0=1$.

By Lemma~\ref{is-onto-step-exp-cut-module}(1), $\nabla_{\mathbf Z}(\nu)$ and $\Delta_{\mathbf Z}(\nu)$
are stable under divided powers, and $c_1c_2T(\nabla_{\mathbf Z}(\nu))\subset\Delta_{\mathbf Z}(\nu)$.
So Lemma~\ref{is-onto-step-exp-cut-integrality} gives
$c_1c_2\Lambda^{pq}S_0(\nabla_{\mathbf Z}(\nu))\subset\Delta_{\mathbf Z}(\nu)$. For
$x\in\nabla_{\mathbf Z}(\nu)$, multiplying by the integer $z$ gives
$c_1c_2\Lambda^{pq}x=c_1c_2\Lambda^{pq}zS_0x\in\Delta_{\mathbf Z}(\nu)$.
% Lean: LeviCut.cut_incl, with the hypotheses hL (LeviProj), hN (NablaGl), hA and hB (AdjBdd, NablaBdd),
% hG (GridInv), hS (GridScalar, with LR.der_diag_iff and LR.hw_rowHom for the weight, LR.der_hw_raise
% and LR.fock_der), hSch (NablaSchur) and hI (GridInt).
\end{proof}

\begin{lemma}\label{is-onto-step-exp-cut}
For an integer $N\geq0$, a sequence $\lambda$ of natural numbers whose first $N$ terms are
nonincreasing, and a prime $l$, let $e_l(N,\lambda)$ be the Weyl exponent of
Definition~\ref{is-onto-step-exp-cut-exponent}: the least $e\geq0$ such that $ml^{e}$ times the
dual Weyl lattice of $\mathbf Z$-rank $N$ and weight $(\lambda_0,\dots,\lambda_{N-1})$ lies in the Weyl
lattice, for some $m$ prime to $l$. Let $u_a=\lambda_a+N-1-a$, and suppose that there are integers
$\mathrm{lo}$ and $w\geq0$ with $\mathrm{lo}\le u_a\le\mathrm{lo}+w$ for $a<N$. Then for every
$a\le N$,
\[
 e_l(N,\lambda)\le e_l(a,\lambda)+e_l\big(N-a,(\lambda_{a+i})_i\big)+a(N-a)\lfloor\log_lw\rfloor ,
\]
with $\lfloor\log_lw\rfloor$ as in Lemma~\ref{is-onto-step-exp-cut-beta}.
\end{lemma}
\leanref{Development/AbelianVarietyTheory/Proofs/DivisorProofs/WeylLatticeExactExponentBlockCutStatement.lean}{GlobalTheory.T.weylLattice_exact_exponent_block_cut}
\leanref{Formalism/PieriBlockIdeals/StatementE.lean}{Formalism.CohomologicalDescent.E.ExCut}
\leanref{Formalism/PieriBlockIdeals/StatementLA.lean}{Formalism.CohomologicalDescent.LA.exZ}
\begin{proof}
Let $a\le N$, put $b=N-a$, and let $\nu=(\lambda_0,\dots,\lambda_{N-1})$,
$\nu'=(\lambda_0,\dots,\lambda_{a-1})$ and $\nu''=(\lambda_a,\dots,\lambda_{N-1})$, all nonincreasing.
By Lemma~\ref{is-onto-step-exp-cut-attained} there are integers $m_1,m_2$ prime to $l$ with
$m_1l^{e_1}\nabla_{\mathbf Z}(\nu')\subset\Delta_{\mathbf Z}(\nu')$ and
$m_2l^{e_2}\nabla_{\mathbf Z}(\nu'')\subset\Delta_{\mathbf Z}(\nu'')$, where $e_1=e_l(a,\lambda)$ and
$e_2=e_l(b,(\lambda_{a+i})_i)$. Lemma~\ref{is-onto-step-exp-cut-beta} gives $m$ prime to $l$ such that
$\Lambda=ml^{\lfloor\log_lw\rfloor}$ has the divisibility property for $w$.

If $N=0$ there is nothing to prove, so let $N\geq1$. For $i\le j<N$, $u_i\le\mathrm{lo}+w\le u_j+w$
gives $\lambda_i+j\le\lambda_j+w+i$. With $i=0$ and $j=N-1$, $\lambda_0+N-1\le\lambda_{N-1}+w$; as
$\lambda_0\geq\lambda_{N-1}$, $N\le w+1$. Put $K=w+1-N\geq0$. Then $\lambda_i\le\lambda_0\le
\lambda_{N-1}+K\le\lambda_j+K$ for all $i,j<N$, and $K+N\le w+1$. So
Lemma~\ref{is-onto-step-exp-cut-blocks}, with $p=a$, $q=b$, $c_1=m_1l^{e_1}$ and $c_2=m_2l^{e_2}$,
gives
\[
 m_1m_2m^{ab}\,l^{e_1+e_2+ab\lfloor\log_lw\rfloor}\,\nabla_{\mathbf Z}(\nu)\subset\Delta_{\mathbf Z}(\nu).
\]
As $m_1m_2m^{ab}$ is prime to $l$, $e_l(N,\lambda)\le e_1+e_2+ab\lfloor\log_lw\rfloor$.
% Lean: T.weylLattice_exact_exponent_block_cut = LeviCut.exCut = ExCut₂ applied to the nine statements
% (GridInv, AdjBdd, GridScalar, GridInt, NablaGl, NablaBdd, NablaSchur, LeviProj, BetaDen); ExCut₂
% uses exact_z_specialization twice, BetaDen, the spread bound hsp, hK with K = w + 1 - (a + b), hw',
% LeviCut.cut_incl (Lemma is-onto-step-exp-cut-blocks) and Nat.sInf_le.
\end{proof}

\begin{lemma}\label{is-onto-step-exp-sum-shift}
Let $N\geq1$, $\nu\in\mathbf N^N$ nonincreasing, $c\geq0$ an integer and $l$ a prime. Then
$e_l(\nu_1+c,\dots,\nu_N+c)=e_l(\nu)$ (Definition~\ref{is-onto-step-exp-cut-exponent}).
\end{lemma}
\leanref{Formalism/PieriBlockIdeals/PieriColumnMinors/HighestWeightShiftEqDetMul.lean}{Formalism.CohomologicalDescent.LA.exact_z_twist_by_det}
\leanref{Formalism/PieriBlockIdeals/PieriColumnMinors/HighestWeightShiftEqDetMul.lean}{Formalism.CohomologicalDescent.LA.incl_z_shift_by_one_iff}
\begin{proof}
By induction on $c$ it suffices to treat $c=1$; write $\nu+1$ for $(\nu_1+1,\dots,\nu_N+1)$ and
$\delta=\det X$. Adding $y$ times the $a$-th row of $X$ to the $b$-th does not change $\delta$ if
$a\ne b$, and multiplies it by $1+y$ if $a=b$. So, for the operators $E^{(k)}_{ab}$ of
Definition~\ref{pd-weight-lattices}(1), $E^{(k)}_{ab}(\delta f)=\delta E^{(k)}_{ab}f$ if $a\ne b$, and
$E^{(k)}_{aa}(\delta f)=\delta\big(E^{(k)}_{aa}f+E^{(k-1)}_{aa}f\big)$, with $E^{(-1)}_{aa}=0$. Also
$h_{\nu+1}=\delta h_\nu$, as $\delta$ is the last leading principal minor. Hence $\delta\Delta_{\mathbf
Z}(\nu)$ contains $h_{\nu+1}$ and is stable under the operators, so $\Delta_{\mathbf
Z}(\nu+1)\subset\delta\Delta_{\mathbf Z}(\nu)$; and, by induction on the number of operators, $\delta$
times every polynomial obtained from $h_\nu$ by the operators lies in $\Delta_{\mathbf Z}(\nu+1)$. So
$\Delta_{\mathbf Z}(\nu+1)=\delta\Delta_{\mathbf Z}(\nu)$.

Next, $\delta\nabla_{\mathbf Z}(\nu)\subset\nabla_{\mathbf Z}(\nu+1)$, since $\delta$ has degree $1$ in
each column and is invariant under the column operations. Conversely let $f\in\nabla_{\mathbf
Z}(\nu+1)$. By Lemma~\ref{is-onto-step-exp-cut-attained} there is $M\ne0$ with
$Mf\in\Delta_{\mathbf Z}(\nu+1)=\delta\Delta_{\mathbf Z}(\nu)$, so $\delta$ divides $Mf$ in
$\mathbf Z[X]$. The coefficients of $\delta$ are $\pm1$, so by Gauss's lemma $f=\delta g$ with
$g\in\mathbf Z[X]$. Then $g$ has degree $\nu_b$ in the $b$-th column, and it is invariant under the
column operations because $\delta$ is and $\delta\ne0$; so $g\in\nabla_{\mathbf Z}(\nu)$. Hence
$\nabla_{\mathbf Z}(\nu+1)=\delta\nabla_{\mathbf Z}(\nu)$.

As multiplication by $\delta$ is injective, $ml^{e}\nabla_{\mathbf Z}(\nu+1)\subset\Delta_{\mathbf
Z}(\nu+1)$ if and only if $ml^{e}\nabla_{\mathbf Z}(\nu)\subset\Delta_{\mathbf Z}(\nu)$, for all $m$
and $e$. So $e_l(\nu+1)=e_l(\nu)$.
% Lean: exact_z_twist_by_det by induction with incl_z_shift_by_one_iff: weyl_shift_eq_map_det_mul
% (highest_weight_shift_eq_det_mul, divided_power_stable_det_mul with divided_power_succ_det_mul and
% divided_power_det_mul_of_ne, genericDet_mul_applyWord_mem), det_mul_mem_dual_weyl_shift,
% det_dvd_of_mem_dual_weyl_shift (la_uniform and genericDet_dvd_of_nsmul_eq_mul) and
% mem_dual_weyl_of_det_mul (weighted_homogeneous_cancel).
\end{proof}

\begin{lemma}\label{is-onto-step-exp-sum-localize}
Let $K$ be a number field, $\mathcal T$ a finite set, $n\geq1$, $\lambda=(\lambda_\tau)_{\tau\in\mathcal
T}$ dominant weights of length $n$, and $M=(M_\tau)$ a family of $\mathcal O_K$-lattices in $K^n$ that
span $K^n$, with the lattices $\nabla_\lambda$, $\Delta_\lambda$, $\nabla_\lambda(M)$,
$\Delta_\lambda(M)$ and the action $\rho_\lambda$ of Definition~\ref{pd-weight-lattices}. Let $w$ be a
finite place of $K$ and $B=(B_\tau)$ matrices in $GL_n(K)$ whose columns form a basis of the
localization of $M_\tau$ at $w$. Then the localizations at $w$ satisfy
\[
 \nabla_\lambda(M)_{(w)}=\rho_\lambda(B)\big((\nabla_\lambda)_{(w)}\big),\qquad
 \Delta_\lambda(M)_{(w)}=\rho_\lambda(B)\big((\Delta_\lambda)_{(w)}\big).
\]
\end{lemma}
\leanref{Formalism/LocalToGlobal/Divisibility/CoefModLemmas.lean}{Formalism.L2.DIV.NabDelLoc}
\leanref{Formalism/LocalToGlobal/Columns/CoefficientsModP/NablaIndependenceHolds.lean}{Formalism.L2.DIV.nablaDeltaLocal_holds}
\begin{proof}
Let $T$ be $\nabla_\lambda$ or $\Delta_\lambda$, and $T(M)$ the corresponding lattice; write
$\mathcal O_{(w)}$ for the local ring at $w$ and $\nu_\tau=\lambda_\tau^\natural$.

\emph{Stability.} For $u\in GL_n(\mathcal O_{(w)})^{\mathcal T}$, $\rho_\lambda(u)$ preserves
$T_{(w)}$. Indeed $T_{(w)}$ is the $\mathcal O_{(w)}$-span of the products $\otimes_\tau f_\tau$ with
$f_\tau$ in $\nabla_{\mathcal O_{(w)}}(\nu_\tau)$, respectively $\Delta_{\mathcal O_{(w)}}(\nu_\tau)$,
as these modules are spanned by the integral polynomials that span them over $\mathcal O_K$ (the
bideterminants of shape $\nu_\tau$ \cite{DEP80}, respectively the polynomials obtained from
$h_{\nu_\tau}$ by the operators $E^{(k)}_{ab}$). By Lemma~\ref{pd-gain-holds-hecke-rows}(2),
$f_\tau({}^tu_\tau X)$ lies in the same module, and $\det(u_\tau)$ is a unit of $\mathcal O_{(w)}$.

\emph{Local pieces.} If $B'$ is another family whose columns form bases of the localizations of the
$M_\tau$ at $w$, then $B_\tau'^{-1}B_\tau\in GL_n(\mathcal O_{(w)})$, so by stability
$\rho_\lambda(B)^{-1}x\in T_{(w)}$ if and only if $\rho_\lambda(B')^{-1}x\in T_{(w)}$. For a finite
place $w'$ put $A_{w'}=\rho_\lambda(B^{w'})(T_{(w')})$, for any family $B^{w'}$ of bases at $w'$. Then
$A_{w'}$ is its own localization at $w'$, and $T(M)=\bigcap_{w'}A_{w'}$ by
Definition~\ref{pd-weight-lattices}(5).

\emph{A common constant.} The lattice $\nabla_\lambda$ is finitely generated, it contains
$\Delta_\lambda$ (as in the proof of Lemma~\ref{pd-gain-holds-hecke-rows}, the Weyl module lies in the
dual Weyl module), and by Lemma~\ref{is-standard-basis}(2) the two span the same $K$-space; so
$c\nabla_\lambda\subset\Delta_\lambda$ for some $c\in\mathcal O_K\smallsetminus\{0\}$. Choose
$d\in\mathcal O_K\smallsetminus\{0\}$ with $dM_\tau\subset\mathcal O_K^n$ and $d\mathcal O_K^n\subset
M_\tau$ for every $\tau$, and put $N=\sum_\tau(|\nu_\tau|+n|\lambda_{\tau,n}|)$. If $g\in
GL_n(K)^{\mathcal T}$ and the entries of the $dg_\tau$ and $dg_\tau^{-1}$ lie in $\mathcal O_{(w')}$,
then $d^N\rho_\lambda(g)$ preserves $(\nabla_\lambda)_{(w')}$: substituting ${}^t(dg_\tau)X$ for $X$
preserves $\nabla_{\mathcal O_{(w')}}(\nu_\tau)$, as it acts on the rows and commutes with the column
operations, and $d^{n|\lambda_{\tau,n}|}\det(g_\tau)^{\lambda_{\tau,n}}$ is integral at $w'$, since
$\det(dg_\tau)$ and $\det(dg_\tau^{-1})$ are. The families $B^{w'}$ and $(B^{w'})^{-1}$ have this
property, so $e=cd^N$ satisfies $eA_{w'}\subset(\nabla_\lambda)_{(w')}$ and $e\nabla_\lambda\subset
A_{w'}$ for every $w'$.

\emph{Localization.} Clearly $T(M)_{(w)}\subset(A_w)_{(w)}=A_w$. Let $x\in A_w$. Choose $s\notin w$
with $se\,x\in\nabla_\lambda$, and $t\in\mathcal O_K\smallsetminus w$ with $t\in e^2\mathcal O_{(w')}$
for every $w'\ne w$ (by the Chinese remainder theorem, at the finitely many places dividing $e$). Let
$w'\ne w$. Then $e(se\,x)\in A_{w'}$, so $tsx=(t/e^2)\,e(sex)\in A_{w'}$. Also $tsx\in A_w$. So
$tsx\in T(M)$, and $x\in T(M)_{(w)}$. Hence $T(M)_{(w)}=A_w=\rho_\lambda(B)(T_{(w)})$.
% Lean: nablaDeltaLocal_holds = nb_loc for nabStd and delStd: nb_stab_nab and nb_stab_del (stability,
% with factSqStable and nb_mem_nabW), nb_indep, nb_at_closed, nb_at_comm (NbComm with nb_comm_nab and
% nb_comm_del, nb_sandwich_fam, nb_entries and nb_scale with the exponent nbN) and localized_iInf
% (with lc_primeTo for t). The Lean's nabP is a submodule of all polynomials; it lies in V_lambda(K)
% by nb_le_VV.
\end{proof}

\begin{lemma}\label{is-onto-step-exp-sum-local}
Let $E$ and $K$ be number fields, $\mathcal T$ a finite set with embeddings $\tau\colon E\to K$, $n\geq1$,
$\lambda=(\lambda_\tau)$ dominant weights of length $n$, and $L'\subset E^n$ an $\mathcal O_E$-lattice
that spans $E^n$, with the lattices $\nabla_\lambda(L'^\vee)$ and $\Delta_\lambda(L'^\vee)$ of
Definition~\ref{pd-weight-lattices}(6). Let $w$ be a finite place of $K$ with residue characteristic
$l$, and put $J=\sum_\tau e_l(\lambda_\tau^\natural)$ (Definition~\ref{is-onto-step-exp-cut-exponent}).
Then for every $x\in\nabla_\lambda(L'^\vee)$ there is $s\in\mathcal O_K\smallsetminus w$ with
$sl^Jx\in\Delta_\lambda(L'^\vee)$.
\end{lemma}
\leanref{Development/AbelianVarietyTheory/Proofs/DivisorProofs/SumExactExponentsSmulMemLocalizedCoeffLatticeStatement.lean}{GlobalTheory.T.sum_exactExponents_smul_mem_localizedCoeffLattice}
\leanref{Development/AbelianVarietyTheory/Proofs/DivisorProofs/SumExactExponentsSmulMemLocalizedCoeffLattice.lean}{GlobalTheory.OntoExponentsLat.exZ_sum_smul_mem_locP_delAtr}
\begin{proof}
Let $M_\tau$ be the dual of the $\mathcal O_K$-span of $\tau(L')$. As $L'$ spans $E^n$, that span is a
lattice spanning $K^n$, and so is $M_\tau$. The local ring $\mathcal O_{(w)}$ is a discrete valuation
ring, and the localization of $M_\tau$ at $w$ is finitely generated and torsion free, hence free of
rank $n$; let $B_\tau\in GL_n(K)$ have as columns a basis of it. By
Lemma~\ref{is-onto-step-exp-sum-localize} and Definition~\ref{pd-weight-lattices}(6),
$x=\rho_\lambda(B)z$ with $z\in(\nabla_\lambda)_{(w)}$, and
$\Delta_\lambda(L'^\vee)_{(w)}=\rho_\lambda(B)\big((\Delta_\lambda)_{(w)}\big)$. So it suffices to
show $l^J(\nabla_\lambda)_{(w)}\subset(\Delta_\lambda)_{(w)}$.

Fix $\tau$, and put $\nu=\lambda_\tau^\natural$ and $e=e_l(\nu)$. By
Lemma~\ref{is-onto-step-exp-cut-attained} there is an integer $m$ prime to $l$ with
$ml^e\nabla_{\mathbf Z}(\nu)\subset\Delta_{\mathbf Z}(\nu)$. The module $\nabla_{\mathcal O_K}(\nu)$
is spanned over $\mathcal O_K$ by the bideterminants of shape $\nu$ \cite{DEP80}, which lie in
$\nabla_{\mathbf Z}(\nu)$, and $\Delta_{\mathbf Z}(\nu)\subset\Delta_{\mathcal O_K}(\nu)$, as the
operators $E^{(k)}_{ab}$ have integer coefficients. So $ml^e\nabla_{\mathcal O_K}(\nu)\subset
\Delta_{\mathcal O_K}(\nu)$, and as $m\notin w$, $l^e\nabla_{\mathcal O_K}(\nu)$ lies in the
localization of $\Delta_{\mathcal O_K}(\nu)$ at $w$.

Now let $f_\tau\in\nabla_{\mathcal O_K}(\lambda_\tau^\natural)$ for every $\tau$, with exponents
$e_\tau=e_l(\lambda_\tau^\natural)$. Choose $t_\tau\in\mathcal O_K\smallsetminus w$ with
$t_\tau l^{e_\tau}f_\tau\in\Delta_{\mathcal O_K}(\lambda_\tau^\natural)$. Then
$\big(\prod_\tau t_\tau\big)l^J\otimes_\tau f_\tau=\otimes_\tau(t_\tau l^{e_\tau}f_\tau)\in
\Delta_\lambda$. Such products span $\nabla_\lambda$, so $l^J\nabla_\lambda\subset(\Delta_\lambda)_{(w)}$,
and hence $l^J(\nabla_\lambda)_{(w)}\subset(\Delta_\lambda)_{(w)}$.
% Lean: exZ_sum_smul_mem_locP_delAtr: du_dualFr_lattice and tauLatLattice (M_tau is a lattice),
% LatBasisAt.basisAtEx (the basis at w), nabDelLoc (Lemma is-onto-step-exp-sum-localize), factor_mem
% (StdDual.dualWeyl_eq and image_bidSet for the bideterminants, WeylRowAll.map_mem_weyl for the Weyl
% module, JetOpWeylLat.smul_mem_localizationAt_of_valuation for 1/m) with exact_z_specialization,
% tensor_mem, and mem_localization_delAtr.
\end{proof}

\begin{definition}\label{is-onto-step-exp-sum-lattices}
Let $R$ be a realization of a definite-side datum with a face $F$, $j$ a lift and $t$ a tuple, with the
notation of Definitions~\ref{is-ties-class-sums} and~\ref{is-ties-lattice-systems}, and let
$c_{0,j,t}\colon L^{d_Z(j,t)}\to\{\text{lift forms at }(j,t)\}$ be the coordinate map of the face at
level $0$ (Definition~\ref{is-complex-data}).
\begin{enumerate}
\item Besides the integral vectors $\mathcal N_y$, the presentation of the lift forms at $(j,t)$ gives
  $\mathcal O_L$-submodules $\mathcal D_y\subset\mathcal V$ for $y\in Y$, the \emph{Weyl vectors} at
  $y$. The isomorphism $e$ of a family tie of $R$ at $(j,t)$ also satisfies
  $e(\mathcal D_y)=\Delta_{\lambda'}(L_y^\vee)$ for every $y\in Y$.
\item The \emph{cut lattice} at $(j,t)$ is the set of $z\in L^{d_Z(j,t)}$ with
  $c_{0,j,t}(z)(y)\in\mathcal N_y$ for every $y\in Y$, and the \emph{Weyl cut lattice} is the set of
  $z$ with $c_{0,j,t}(z)(y)\in\mathcal D_y$ for every $y\in Y$.
\end{enumerate}
\end{definition}
\leanref{Formalism/LocalToGlobal/Divisibility/RowLat.lean}{Formalism.L2.DIV.BlockDat.cutZ}
\leanref{Formalism/LocalToGlobal/Divisibility/RowLat.lean}{Formalism.L2.DIV.BlockDat.delZ}
\leanref{Formalism/LocalToGlobal/Divisibility/ClassSum.lean}{Formalism.L2.DIV.CSD.delLat}

\begin{lemma}\label{is-onto-step-exp-sum}
Let $\mathcal H$ be the class of standard families. Eventually along $\mathcal H$, every family has a
scale beyond which the following holds for every point $p$ that satisfies the point conditions and the
further conditions, every lift $j$, and every parameter record $P$ and weights $W$ in the exponent dictionary
at $j$ satisfying its numerical conditions (Lemma~\ref{is-onto-step-exp-dictionary}). For every tuple
$t$ and finite place $w$ of $L$, with $J_A=\sum_\tau e_{p_w}(N_r,\lambda_\tau)$
(Lemma~\ref{is-onto-step-exp-cut}), $p_w^{J_A}$ times the cut lattice at $(j,t)$ of the realization of
$p$ and its face (Definition~\ref{is-onto-step-exp-sum-lattices}) lies in the localization at $w$ of
the Weyl cut lattice at $(j,t)$.
\end{lemma}
\leanref{Development/AutomorphicArithmetic/ThetaCorrespondence/ThetaPairLifts/ThirdPlaceSeries/SumExpStStatement.lean}{GlobalTheory.T.sum_expansion'}
\leanref{Development/AbelianVarietyTheory/Definitions/Divisors/SumExpSt.lean}{GlobalTheory.ExpTie.LIncl}
\begin{proof}
The conclusion holds at every scale. Let $p$, $j$, $P$, $W$, $t$ and $w$ be given, let $R$ be the
realization of $p$, put $l=p_w$, and let $n=N_r$ be the rank of $V_j$. The weak characterization of
$R$, a datum of the point, gives a family tie of $R$; let $e\colon\mathcal V\to V_{\lambda'}(L)$ be its
isomorphism at $(j,t)$, so that $e(\mathcal N_y)=\nabla_{\lambda'}(L_y^\vee)$ and
$e(\mathcal D_y)=\Delta_{\lambda'}(L_y^\vee)$ for $y\in Y$
(Definitions~\ref{is-ties-lattice-systems}(4) and~\ref{is-onto-step-exp-sum-lattices}). Each $L_y$
is an $\mathcal O_{E_h}$-lattice that spans $E_h^n$.

\emph{The exponent.} The dictionary gives, for every $\tau$, an integer $c_\tau$ with
$\lambda_{\tau,a}=\lambda'_{\tau,a+1}+c_\tau$ for $a<n$, where $\lambda'_\tau$ is the dominant weight
that the Harish-Chandra parameters give at level $j+1$. As $\lambda_{\tau,n-1}\geq0$, this says
$\lambda_{\tau,a}=\lambda'^{\natural}_{\tau,a+1}+\lambda_{\tau,n-1}$ for $a<n$, with the natural weight
$\lambda'^\natural_\tau$ (Definition~\ref{is-polynomial-model}(1)). By
Lemma~\ref{is-onto-step-exp-sum-shift}, $e_l(n,\lambda_\tau)=e_l(\lambda'^\natural_\tau)$, so
$J_A=\sum_\tau e_l(\lambda'^\natural_\tau)$.

\emph{One class.} Let $y\in Y$ and $x\in\mathcal N_y$. Then $e(x)\in\nabla_{\lambda'}(L_y^\vee)$, and
Lemma~\ref{is-onto-step-exp-sum-local}, with the embeddings $\tau_L$ and $L'=L_y$, gives
$s\in\mathcal O_L\smallsetminus w$ with $sl^{J_A}e(x)\in\Delta_{\lambda'}(L_y^\vee)=e(\mathcal D_y)$. As
$e$ is injective and $\mathcal O_L$-linear, $sl^{J_A}x\in\mathcal D_y$.

\emph{All classes.} Let $z$ be in the cut lattice. For each $y\in Y$ the previous step, applied to
$x=c_{0,j,t}(z)(y)$, gives $s_y\in\mathcal O_L\smallsetminus w$ with
$s_yl^{J_A}c_{0,j,t}(z)(y)\in\mathcal D_y$. Put $s=\prod_ys_y$, which is not in $w$. Then
$c_{0,j,t}(sl^{J_A}z)(y)=\big(\prod_{y'\ne y}s_{y'}\big)s_yl^{J_A}c_{0,j,t}(z)(y)\in\mathcal D_y$ for
every $y$. So $sl^{J_A}z$ lies in the Weyl cut lattice, and $l^{J_A}z$ lies in its localization at $w$.
% Lean: T.sum_expansion' = ExpTie.sumExpSt_alt, which holds at every point (EventuallyAtPoints.of_forall)
% by lIncl_sumExp: the tie p.charW.base.fam.tie and its coefL j t (e with hNab and hDel); hJ, from
% exZ_res_congr (Dict.Nr), Dict.lam and exZ_res_eq_natW (exact_z_twist_by_det, Lemma
% is-onto-step-exp-sum-shift); cutLat_locIncl (all classes), locIncl_of_equiv (one class) and
% T.sum_exactExponents_smul_mem_localizedCoeffLattice (Lemma is-onto-step-exp-sum-local) with
% RepsCompare.isLattice_hLat. Only the fields Nr and lam of the dictionary and none of its numerical
% conditions enter.
\end{proof}

\begin{lemma}\label{is-onto-step-exp-model}
Assume the published results of Theorem~\ref{lit-bundle-unitary} and the local Langlands laws of
Chapter~\ref{chapter-galois-representations}; of these, base change from definite unitary groups
(Theorem~\ref{gal-base-change-definite}), the Galois representations of
Definition~\ref{gal-galois-representations}, Theorem~\ref{lit-clozel-cohomological}, the Chebotarev
theorem for primes of degree one (Theorem~\ref{lit-chebotarev-degree-one}), the $\ell$-adic
characters attached to algebraic Hecke characters, the local Langlands laws and the crystalline
lattices of Bhatt and Scholze (Theorem~\ref{lit-bhatt-scholze}) enter. Let $\mathcal H$ be the
class of standard families. Eventually along $\mathcal H$, every family has a scale beyond which the
following holds for every point $p$ that satisfies the point conditions and the further conditions,
every lift $j$, every parameter record $P$ and weights $W$ in the exponent dictionary at $j$ satisfying its
numerical conditions, and every twisted place $w$ whose residue characteristic $p_w$ is in the range
of the weights ($l_J<p_w\le h$, $p_w$ not excluded by $P$ and not in the special set of $W$). There is
an \emph{exponent model} at $w$: a family of Hida settings, in each of which $p_w$ is not a unit, that
covers the block data of $p$ at $(j,\mu_w)$ at $w$ and is linked to the weights $W$, twisted by
Frobenius data, at $p_w$.
\end{lemma}
\leanref{Development/AutomorphicArithmetic/ThetaCorrespondence/ThetaPairLifts/ThirdPlaceSeries/ModelStOfFramesStatement.lean}{GlobalTheory.T.model_statement_of_crystalline_chebotarev_frame_gal}
\leanref{Development/AbelianVarietyTheory/Definitions/Divisors/SumExpSt.lean}{GlobalTheory.ExpTie.ModelAt}
\begin{proof}
\openstep{The construction of the exponent models from the Galois representations, the Chebotarev
theorem and the crystalline lattices is not written out.}
% Lean: T.model_statement_of_crystalline_chebotarev_frame_gal applied to T.small_estimate,
% T.frame_statement_of_small_statement and T.galois_frame_compatibility_alt (which takes the published
% inputs ClozelCohomological, HeckeLAdicCharacter, ChebotarevDegreeOne, CuspUAlt, BaseChangeDefinite,
% BhattScholze, HeckeLAdicCharacterAbove and RecTExists). The small, frame and Galois statements are
% natural separate pieces. Size: about 198,000 lines of Lean in 12,800 theorems below them (text count,
% stopping at statements the book cites and at the published statements), about 148,000 of them below
% no other open step of Lemma is-onto-estimates.
\end{proof}

\subsection{The recursive exponent and its total bound}\label{subsection-is-recursive-exponent}

\begin{definition}\label{is-onto-step-exp-bound-lift-exponent}
Let $R$ be a realization of a definite-side datum with a face, $j$ a lift and $t$ a tuple, and let
$C_{j,t}$ and $D_{j,t}$ be the cut lattice and the Weyl cut lattice at $(j,t)$, which are
$\mathcal O_L$-submodules of $L^{d_Z(j,t)}$ (Definition~\ref{is-onto-step-exp-sum-lattices}). Let $w$ be
a finite place of $L$, with residue characteristic $p_w$ and ramification index $e_w$, and write
$M_{(w)}$ for the localization at $w$ of an $\mathcal O_L$-module $M$. The \emph{lift exponent} at
$(j,t,w)$ is $J_0(t,w,j)=J_w/e_w$, where $J_w$ is the least integer $J\geq0$ such that
$az\in(D_{j,t})_{(w)}$ for every $a\in w^{J}$ and every $z\in C_{j,t}$, and $J_w=0$ if there is no
such $J$. For a point $p$ (Definition~\ref{is-point}), the function $J_0$ of its scale data $Z(p)$, and
so of the numerical bounds $N_p(C,C_H)$, is the lift exponent of its realization and face.
\end{definition}
\leanref{Formalism/LocalToGlobal/Divisibility/RowLat.lean}{Formalism.L2.DIV.BlockDat.JexpU}
\leanref{Formalism/LocalToGlobal/Divisibility/RowLat.lean}{Formalism.L2.DIV.expW}
\leanref{Development/AbelianVarietyTheory/Proofs/DivisorProofs/ExpansionOfNumericsModel.lean}{GlobalTheory.ExpTie.Asm.splitPartNumbers_exponent_eq_liftExponent}

\begin{lemma}\label{is-onto-step-exp-bound-lift}
Let $R$, $j$, $t$ and $w$ be as in Definition~\ref{is-onto-step-exp-bound-lift-exponent}, and let
$J\geq0$ be an integer with $p_w^{J}C_{j,t}\subset(D_{j,t})_{(w)}$. Then $0\le J_0(t,w,j)\le J$.
\end{lemma}
\leanref{Development/AbelianVarietyTheory/Proofs/DivisorProofs/ExponentDatumLeOfPrimePowerLatticeInclusionStatement.lean}{GlobalTheory.T.exponent_datum_le_of_primePowerLatticeInclusion}
\leanref{Development/AbelianVarietyTheory/Proofs/DivisorProofs/ExponentDatumLeOfPrimePowerLatticeInclusion.lean}{GlobalTheory.ExpTie.LIncl.jexpU_bounds}
\begin{proof}
As $J_w\geq0$ and $e_w>0$, $J_0(t,w,j)\geq0$. For the upper bound it suffices to show that $Je_w$ has
the property that defines $J_w$, for then $J_w\le Je_w$. Let $a\in w^{Je_w}$ and $z\in C_{j,t}$. In the
local ring $\mathcal O_{(w)}$ of $L$ at $w$, a discrete valuation ring, $w^{e_w}\mathcal
O_{(w)}=p_w\mathcal O_{(w)}$, so $a\in p_w^{J}\mathcal O_{(w)}$: there are $b\in\mathcal O_L$ and
$s'\in\mathcal O_L\smallsetminus w$ with $s'a=p_w^{J}b$. By hypothesis there is $s\in\mathcal
O_L\smallsetminus w$ with $sp_w^{J}z\in D_{j,t}$. Then $(ss')(az)=b\,(sp_w^{J}z)\in D_{j,t}$ and
$ss'\notin w$, so $az\in(D_{j,t})_{(w)}$.
% Lean: LIncl.jexpU_bounds, with exists_mul_eq_prW_pow (algebraMap_mem_maximalIdeal_pow_iff and
% span_prW_ow: the maximal ideal of O_(w) to the power e_w is p_w O_(w)), plc w = prW w
% (eq_prW_of_mem) and Nat.sInf_le for J_w <= J e_w.
\end{proof}

\begin{definition}\label{is-onto-step-exp-bound-weights}
A \emph{parameter record} $P$ consists of natural numbers $k,m,t,f,n,N_r$ and real numbers
$L,x,b',\eta,S$. Its derived numbers are $b=e^{L}$, $w''=b'/(xn)$, $N=k+m(t+1)$,
\[
 \ell^*=\frac{mb}{nt^2(N+t)L^2+kn^2(S+t+L)},\qquad h=b+S+\eta+t+1,\qquad
 l_J=\frac{mbL}{k(m+t)\log(3h)},
\]
$\kappa_{\max}=\lfloor2b/\ell^*\rfloor$, $h_0=\lfloor h\rfloor$ and $\kappa(l)=\lfloor h/(l-1)\rfloor$
for natural numbers $l$, where $\lfloor y\rfloor=0$ for $y<0$ and $y/0=0$.
\begin{enumerate}
\item $P$ satisfies the \emph{hypotheses} if $x\geq1$, $b'>0$, $1\le m\le k$, $n\geq1$, $f\geq1$,
  $k+m+4\le t$, $3m<t$, $\eta\le b/(t^3fk^3)$, $N_r+1=k+m+t$, $t\le\eta$, $S\geq0$,
  $b'+S+\eta+b/\ell^*<\ell^*$ and $2\eta+t+8b/\ell^*+7<\ell^*$. It is \emph{large} if $L\geq4$,
  $48k^2t^2<b$, $12ktb'<b$, $2(2b/\ell^*+t)<w''$ and $S+b'(1+x^{-1}+x^{-2})+\eta+(t-1)/2\le2b'$.
\item For real numbers $X'$ and $A$, $P$ has \emph{rows} $(X',A)$ if $b'=xb/X'$, $X'(S+\eta)\le b$,
  $kS\le w''$, $S\geq1$, $m\geq2$, $f=2n$, $f\le k$, $t\le A$ and $N\le Ax$. The pair $(X',A)$ is
  \emph{big} for $P$ if $x\geq5$, $L\geq4$, $4(L+2)\le X'$, $5.2A^3x\le X'$, $40A^7x^2L\le X'$ and
  $40A^7x^2L^3\le e^{L}$.
\item \emph{Weights} $W$ for $P$, indexed by a finite set $\mathcal T$, consist of: for each
  $\tau\in\mathcal T$, a sequence $\lambda_\tau=(\lambda_{\tau,0},\lambda_{\tau,1},\dots)$ of natural
  numbers, natural numbers $j_\tau$ and $q_\tau$, integers $x^-_\tau$ and $x^+_\tau$, and real numbers
  $c_\tau$ and $e_\tau$; a subset $T_W\subset\mathcal T$, whose elements are \emph{thick}; a finite set
  $\mathrm{Sp}$ of natural numbers, the \emph{special set}; and \emph{Frobenius data}, maps
  $\phi_l\colon\mathcal T\to\mathcal T$ and $s_l\colon\mathcal T\to\{0,1\}$ for $l\in\mathbf N$. Put
  $U_\tau(a)=\lambda_{\tau,a}+N_r-1-a$ for $a\geq0$, $p_\tau=k-1+j_\tau$,
  $\mathrm{bot}_\tau=U_\tau(p_\tau+t-1)$, $\mathrm{top}_\tau=U_\tau(p_\tau)$ and
  $\mathcal U_\tau=\{U_\tau(a):a<N_r\}$; the \emph{width} at $\tau$ is $U_\tau(0)-U_\tau(N_r-1)$.
\item The weights satisfy the \emph{weight conditions} if for every $\tau$:
  \begin{enumerate}
  \item[(W1)] $\lambda_{\tau,a'}\le\lambda_{\tau,a}$ for $a\le a'<N_r$, and $p_\tau+t+q_\tau=N_r$;
  \item[(W2)] $\lambda_{\tau,a}=\lambda_{\tau,p_\tau}$ for $p_\tau\le a<p_\tau+t$, and
    $\mathrm{bot}_\tau+(t-1)/2-c_\tau=-e_\tau$;
  \item[(W3)] $c_\tau\le x^-_\tau\le x^+_\tau\le c_\tau+S$, and $x^-_\tau\le U_\tau(a)\le x^+_\tau$ for
    $a<k-1$;
  \item[(W4)] $2b/\ell^*+t/2+3\le|U_\tau(a)-c_\tau+e_\tau|\le\eta$ for $k-1\le a<p_\tau$ and for
    $p_\tau+t\le a<N_r$;
  \item[(W5)] if $\tau\in T_W$, then $b\le e_\tau\le b+\eta+1$ and $x^+_\tau\le h$;
  \item[(W6)] if $\tau\notin T_W$, then either there is $b_*$ with $b'\le b_*\le b'(1+x^{-2})$ and
    $b_*+b'/x\le e_\tau\le b_*+b'/x+\eta$, and the width at $\tau$ is at most
    $S+b'(1+x^{-1}+x^{-2})+\eta+(t-1)/2$; or $w''\le e_\tau\le w''+\eta$, and the width at $\tau$ is
    at most $S+w''+\eta+(t-1)/2$.
  \end{enumerate}
  These conditions do not involve the special set or the Frobenius data.
\item $P$ and $W$ satisfy the \emph{numerical conditions} if $P$ satisfies the hypotheses and is
  large, $W$ satisfies the weight conditions, and $P$ has rows $(X',A)$ for some pair $(X',A)$ that is
  big for $P$.
\item A natural number $l$ is \emph{excluded} by $P$ if it divides an integer $z$ with
  $b-(t-1)/2-\kappa_{\max}\le z\le b+S+\eta+1+(t-1)/2+\kappa_{\max}$. It is in the \emph{range} of $W$
  if $l_J<l\le h$, $l$ is not excluded by $P$ and $l\notin\mathrm{Sp}$.
\end{enumerate}
The parameter records and weights of Lemma~\ref{is-onto-step-exp-dictionary} are of this kind, and
its numerical conditions are those of (5).
\end{definition}
\leanref{Formalism/PieriBlockIdeals/MertensReciprocalElementary.lean}{Formalism.CohomologicalDescent.LogSmallConstantBounds.Par}
\leanref{Formalism/PieriBlockIdeals/MertensReciprocalElementary.lean}{Formalism.CohomologicalDescent.LogSmallConstantBounds.Par.Hyp}
\leanref{Formalism/PieriBlockIdeals/MertensReciprocalElementary.lean}{Formalism.CohomologicalDescent.LogSmallConstantBounds.Par.Large}
\leanref{Formalism/PieriBlockIdeals/MertensReciprocalElementary.lean}{Formalism.CohomologicalDescent.LogSmallConstantBounds.Par.Excl}
\leanref{Formalism/PieriBlockIdeals/TailInequalityEventually.lean}{Formalism.CohomologicalDescent.LargeParameterEstimates.Rows}
\leanref{Formalism/PieriBlockIdeals/TailInequalityEventually.lean}{Formalism.CohomologicalDescent.LargeParameterEstimates.Big}
\leanref{Formalism/PieriBlockIdeals/LatticeInclusionRealSettingAlt.lean}{Formalism.CohomologicalDescent.ShortWeightBounds.Wts}
\leanref{Formalism/PieriBlockIdeals/LatticeInclusionRealSettingAlt.lean}{Formalism.CohomologicalDescent.ShortWeightBounds.Wts.Ok}
\leanref{Formalism/PieriBlockIdeals/LatticeInclusionRealSettingAlt.lean}{Formalism.CohomologicalDescent.ShortWeightBounds.Wts.WRange}
\leanref{Development/AbelianVarietyTheory/Definitions/Divisors/SumExpSt.lean}{GlobalTheory.ExpTie.NumOk}

\begin{definition}\label{is-onto-step-exp-bound-recursion}
Let $P$ be a parameter record and $W$ weights for $P$ indexed by $\mathcal T$
(Definition~\ref{is-onto-step-exp-bound-weights}), and let $l$ be a natural number. Below, $y/0=0$
for every real $y$, a congruence modulo $0$ is an equality, and the least element of an empty set of
natural numbers is $0$; for a prime $l$ these conventions do not enter.
\begin{enumerate}
\item The \emph{Weyl exponents} at $\tau\in\mathcal T$ are $e_\tau(l)=e_l(N_r,\lambda_\tau)$ and
  $e^{\mathrm c}_\tau(l)=e_l(k-1,\lambda_\tau)$, with $e_l$ defined as in
  Definition~\ref{is-onto-step-exp-cut-exponent} for every sequence of natural numbers, and the
  \emph{Weyl exponent sum} is $J_A(l)=\sum_{\tau}e_\tau(l)$.
\item For an integer $\beta$ and a finite set $\mathcal U$ of integers, let $n_l(\beta,\mathcal U)$ be
  the number of $u\in\mathcal U$ outside $[\beta,\beta+t-1]$ that are congruent modulo $l$ to an
  integer in $[\beta-\kappa(l),\beta+t-1+\kappa(l)]$. The \emph{window counts} at $\tau$ are
  $n_\tau(l)=n_l(\mathrm{bot}_\tau,\mathcal U_\tau)$ and $n^{\mathrm F}_\tau(l)$, which is
  $n_l(\mathrm{bot}_{\tau'},\mathcal U_{\tau'})$ if $s_l(\tau)=0$ and
  $n_l(h_0-\mathrm{top}_{\tau'},\{h_0-u:u\in\mathcal U_{\tau'}\})$ if $s_l(\tau)=1$, where
  $\tau'=\phi_l(\tau)$. The \emph{depth} at $\tau$ is
  \[
   i_\tau(l)=\big\lfloor q_\tau h+\big(l\,n^{\mathrm F}_\tau(l)+n_\tau(l)\big)\big(t+2\kappa(l)\big)\big\rfloor.
  \]
\item Let $d_\tau(l)$ be the least integer $d\geq1$ with $l\mid x^-_\tau-\mathrm{top}_\tau-d$; let
  $\delta_\tau(l)$ be $1$ if some integer in $[x^-_\tau,x^+_\tau]$ is congruent to $\mathrm{top}_\tau$
  modulo $l$, and $0$ otherwise; let $\gamma_\tau(l)=\min\big(k-1,t,i_\tau(l)/d_\tau(l)\big)$ if
  $d_\tau(l)<x^-_\tau-\mathrm{top}_\tau$, and $\gamma_\tau(l)=0$ otherwise; and write
  $\log^+y=\max(0,\log y)$ for $y>0$ and $\log^+0=0$. The \emph{short bound} and the \emph{coarse bound} at $\tau$ are
  \[
   C^{\mathrm s}_\tau(l)=(q_\tau+2j_\tau)\frac{h_0+N_r}{l-1}+e^{\mathrm c}_\tau(l)
   +(k-1)\Big(\frac{x^+_\tau-x^-_\tau}{l}+\delta_\tau(l)\Big)+\gamma_\tau(l)
   +\frac{i_\tau(l)}{l}\Big(1+\log^+\frac{h_0}{l}\Big)
  \]
  and $C_\tau(l)=C^{\mathrm s}_\tau(l)+6(q_\tau+2j_\tau)$.
\item The \emph{short weight bound} and the \emph{coarse weight bound} are
  \[
   J^{\mathrm s}_W(l)=\sum_{\tau\in T_W}\min\big(e_\tau(l),\lfloor C^{\mathrm s}_\tau(l)\rfloor\big)
   +\sum_{\tau\notin T_W}e_\tau(l),\qquad
   J_W(l)=\sum_{\tau\in T_W}\min\big(e_\tau(l),\lfloor C_\tau(l)\rfloor\big)+\sum_{\tau\notin T_W}e_\tau(l).
  \]
\item The \emph{recursive exponent} of $P$ and $W$ at $l$ is $J'(l)=J_A(l)$ if $l\le l_J$, $l$ is
  excluded by $P$ or $l\in\mathrm{Sp}$, and $J'(l)=J_W(l)$ otherwise.
\end{enumerate}
\end{definition}
\leanref{Formalism/PieriBlockIdeals/LatticeInclusionRealSettingAlt.lean}{Formalism.CohomologicalDescent.ShortWeightBounds.Wts.JA}
\leanref{Formalism/PieriBlockIdeals/Statement.lean}{Formalism.CohomologicalDescent.mwCount}
\leanref{Formalism/PieriBlockIdeals/LatticeInclusionRealSettingAlt.lean}{Formalism.CohomologicalDescent.ShortWeightBounds.Wts.mwF}
\leanref{Formalism/PieriBlockIdeals/LatticeInclusionRealSettingAlt.lean}{Formalism.CohomologicalDescent.ShortWeightBounds.Wts.iW}
\leanref{Formalism/PieriBlockIdeals/FactorialValuationSum.lean}{Formalism.CohomologicalDescent.UnitNormBounds.cCoarse}
\leanref{Formalism/PieriBlockIdeals/LatticeInclusionRealSettingAlt.lean}{Formalism.CohomologicalDescent.ShortWeightBounds.cShort6}
\leanref{Formalism/PieriBlockIdeals/LatticeInclusionRealSettingAlt.lean}{Formalism.CohomologicalDescent.ShortWeightBounds.Wts.JW}
\leanref{Formalism/PieriBlockIdeals/LatticeInclusionRealSettingAlt.lean}{Formalism.CohomologicalDescent.ShortWeightBounds.Wts.coarseWeightBound}
\leanref{Formalism/PieriBlockIdeals/LatticeInclusionRealSettingAlt.lean}{Formalism.CohomologicalDescent.ShortWeightBounds.Wts.Jrec7}

\begin{lemma}\label{is-onto-step-exp-bound-numbers}
Let $P$ be a parameter record that satisfies the hypotheses and is large
(Definition~\ref{is-onto-step-exp-bound-weights}). Then $\ell^*>7$, $2\ell^*\le b$, $2\eta<b$, $h\le2b$,
$0<\log(3h)\le2L$, $b/(4kt)\le l_J$, $L\ell^*\le l_J$, $\ell^*+1<l_J$, $3b'<l_J$, $w''\le b'$ and
$\kappa_{\max}+t<w''/2$. For every integer $l>l_J$, moreover, $\ell^*\le l-1$, $\kappa(l)\le2b/\ell^*$,
$\kappa(l)\le\kappa_{\max}$ and $\eta+t+\kappa(l)<l$.
\end{lemma}
\leanref{Formalism/PieriBlockIdeals/PieriColumnMinors/ParameterBounds.lean}{Formalism.CohomologicalDescent.LogSmallConstantBounds.paramsBasic_holds}
\leanref{Formalism/PieriBlockIdeals/PieriColumnMinors/WeightBounds.lean}{Formalism.CohomologicalDescent.ShortWeightBounds.two_eta_lt_b}
\leanref{Formalism/PieriBlockIdeals/PieriColumnMinors/ParameterBounds.lean}{Formalism.CohomologicalDescent.LogSmallConstantBounds.paramsLStar_holds}
\leanref{Formalism/PieriBlockIdeals/PieriColumnMinors/ParameterBounds.lean}{Formalism.CohomologicalDescent.LogSmallConstantBounds.paramsTop_holds}
\leanref{Formalism/PieriBlockIdeals/PieriColumnMinors/ParameterBounds.lean}{Formalism.CohomologicalDescent.LogSmallConstantBounds.paramsThin_holds}
\leanref{Formalism/PieriBlockIdeals/PieriColumnMinors/ParameterBounds.lean}{Formalism.CohomologicalDescent.LogSmallConstantBounds.paramsLargePrime_holds}
\begin{proof}
The hypotheses give $t\geq k+m+4\geq6$, $m<t$, $n\geq1$ and $\eta\geq t\geq0$, and largeness gives
$L\geq4$; also $b=e^L>1$. Let $D_*$ be the denominator of $\ell^*$. As $n\geq1$ and $N\geq k$,
$D_*\geq t^2(N+t)L^2\geq t^2kL^2$ and $D_*\geq t^3L^2\geq t>2m$. Hence $\ell^*>0$ and
$2\ell^*\le b$. The last hypothesis gives $\ell^*>2\eta+t+8b/\ell^*+7>7$ and $t+1<\ell^*$, and the one
before it gives $S+\eta<\ell^*$, as $b'>0$, $S\geq0$ and $b/\ell^*>0$. So $2\eta<2\ell^*\le b$, and
$h=b+(S+\eta)+(t+1)<b+2\ell^*\le2b$. Then $3h>3b>1$ and $\log(3h)\le\log(6b)=L+\log6\le2L$.

As $m<t$, $k(m+t)\log(3h)\le4ktL$, so $l_J\geq mbL/(4ktL)\geq b/(4kt)$. As $tL\geq4$,
$D_*\geq t^2kL^2\geq4ktL\geq k(m+t)\log(3h)$, so $l_J\geq mbL/D_*=L\ell^*\geq4\ell^*>\ell^*+1$. As
$12ktb'<b$, $3b'<b/(4kt)\le l_J$. As $x\geq1$ and $n\geq1$, $w''=b'/(xn)\le b'$. As
$\kappa_{\max}\le2b/\ell^*$, largeness gives $\kappa_{\max}+t<w''/2$.

Let $l>l_J$ be an integer. Then $l-1>\ell^*>0$, so $h/(l-1)\le2b/\ell^*$; hence $\kappa(l)\le2b/\ell^*$
and, the floor being monotone, $\kappa(l)\le\kappa_{\max}$. Finally
$\eta+t+\kappa(l)\le\eta+t+2b/\ell^*<\ell^*<l$ by the last hypothesis.
% Lean: paramsBasic_holds (Elem0) and lstar_bounds_of_large; two_mul_lstar_le_b (le_lstarDenom,
% lstarDenom_pos), seven_lt_lstar, two_eta_lt_b; log_three_mul_h_pos_le; b_div_le_lJ (lJDenom_le);
% paramsLStar_holds (lJDenom_le_lstarDenom); paramsTop_holds (Elem_e); paramsThin_holds (Elem_f, with
% kmax_le); paramsLargePrime_holds (Elem_d, through lstar_add_one_lt_lJ).
\end{proof}

\begin{lemma}\label{is-onto-step-exp-bound-total-rank-two}
Let $N\le2$, $\nu=(\nu_1\geq\dots\geq\nu_N\geq0)$ integers and $l$ a prime, with $\nabla_{\mathbf Z}(\nu)$
and $\Delta_{\mathbf Z}(\nu)$ as in Definition~\ref{is-onto-step-exp-cut-exponent}, and let $v_l$ be the
$l$-adic valuation. For $1\le a<b\le N$ put $n_{ab}=\nu_a-\nu_b+b-a-1$, and let
\[
 E=\sum_{1\le a<b\le N}\ \max_{0\le i\le n_{ab}}v_l\binom{n_{ab}}i.
\]
Then there is an integer $m$ prime to $l$ with $ml^{E}\nabla_{\mathbf Z}(\nu)\subset\Delta_{\mathbf
Z}(\nu)$. If $N\le1$, then $E=0$, and so $e_l(\nu)=0$.
\end{lemma}
\leanref{Formalism/PieriBlockIdeals/PieriColumnMinors/ProdInitialMinorsMemDetSpan.lean}{Formalism.CohomologicalDescent.LA.lattice_z_rank_le_two}
\leanref{Formalism/PieriBlockIdeals/PieriColumnMinors/ProdInitialMinorsMemDetSpan.lean}{Formalism.CohomologicalDescent.LA.lattice_z_zero_columns}
\leanref{Formalism/PieriBlockIdeals/PieriColumnMinors/ProdInitialMinorsMemDetSpan.lean}{Formalism.CohomologicalDescent.LA.lattice_z_rank_one}
\leanref{Formalism/PieriBlockIdeals/PieriColumnMinors/ProdInitialMinorsMemDetSpan.lean}{Formalism.CohomologicalDescent.LA.lattice_z_rank_two_two_columns}
\begin{proof}
If $N\le1$ there are no pairs $a<b$, so $E=0$, and $m=1$ qualifies. Indeed, for $N=0$ the polynomials
in no variables are $\mathbf Z$, $h_\nu=1$ lies in $\Delta_{\mathbf Z}(\nu)$, and every element of
$\nabla_{\mathbf Z}(\nu)$ is an integer multiple of $1$. For $N=1$, with the single variable $x$, an
element of $\nabla_{\mathbf Z}(\nu)$ is homogeneous of degree $\nu_1$, so it is an integer multiple of
$x^{\nu_1}=h_\nu\in\Delta_{\mathbf Z}(\nu)$. Then $e=0$ has the property that defines $e_l(\nu)$, so
$e_l(\nu)=0$.

Let $N=2$. Write $x$ and $y$ for the entries of the first column of $X$, $D=\det X$ and
$d=\nu_1-\nu_2$, so that $E=\max_{0\le i\le d}v_l\binom di$. Let $M$ be the $\mathbf Z$-span of the
$x^py^q$ and $B$ the $\mathbf Z$-span of the $\binom dqx^py^q$, for $p+q=d$.
\begin{enumerate}
\item[(a)] $\Delta_{\mathbf Z}(\nu)=D^{\nu_2}B$. Here $h_\nu=x^{d}D^{\nu_2}\in D^{\nu_2}B$, and
  $D^{\nu_2}B$ is stable under the operators $E^{(k)}_{ab}$ of Definition~\ref{pd-weight-lattices}(1):
  adding a multiple of one row to another does not change $D$, and multiplying a row by $1+s$
  multiplies $D$ by $1+s$, so it suffices that $B$ is stable. On $x^py^q$, the operators that replace
  the first row give $\binom pk x^{p-k}y^q$ times $x^k$ or $y^k$, and those that replace the second
  row give $\binom qk x^py^{q-k}$ times $x^k$ or $y^k$; the coefficients stay in $B$ since
  $\binom{d}{k+q}$ divides $\binom dq\binom{p}{k}$ and $\binom{d}{q-k}$ divides $\binom dq\binom qk$.
  So $\Delta_{\mathbf Z}(\nu)\subset D^{\nu_2}B$. Conversely the operator that adds $s$ times the second
  row to the first sends $h_\nu$ to the coefficient of $s^q$ in $(x+sy)^dD^{\nu_2}$, which is
  $\binom dqx^{d-q}y^qD^{\nu_2}$.
\item[(b)] $\nabla_{\mathbf Z}(\nu)=D^{\nu_2}M$. By
  Lemma~\ref{pd-kappa-defects-restriction-standard}(1), $\nabla_{\mathbf Z}(\nu)$ is spanned by the
  bideterminants of shape $\nu$: products of $\nu_2$ minors of size $2$, each equal to $D$, and $d$
  minors of size $1$ on the first column, each equal to $x$ or $y$. So $\nabla_{\mathbf
  Z}(\nu)\subset D^{\nu_2}M$. Conversely $\binom dqD^{\nu_2}x^py^q$ lies in $\Delta_{\mathbf Z}(\nu)$, which
  lies in $\nabla_{\mathbf Z}(\nu)$ (as in the proof of Lemma~\ref{pd-gain-holds-hecke-rows}); as the
  conditions that define $\nabla_{\mathbf Z}(\nu)$ in $\mathbf Z[X]$ are linear, a polynomial with a
  nonzero multiple in $\nabla_{\mathbf Z}(\nu)$ lies in it, so $D^{\nu_2}x^py^q\in\nabla_{\mathbf Z}(\nu)$.
\item[(c)] Let $m=\prod_{i=0}^d\binom di l^{-v_l\binom di}$, an integer prime to $l$. Every $\binom dq$
  divides $ml^{E}$, so $ml^{E}D^{\nu_2}x^py^q\in D^{\nu_2}B$ for $p+q=d$.
\end{enumerate}
By (a), (b) and (c), $ml^{E}\nabla_{\mathbf Z}(\nu)\subset\Delta_{\mathbf Z}(\nu)$.
% Lean: lattice_z_rank_le_two, by cases: lattice_z_zero_columns (hw = 1, x = coeff 0 x * 1) for N = 0,
% lattice_z_rank_one (highest_weight_one_column, the support of a weighted homogeneous x) for N = 1,
% and lattice_z_rank_two_two_columns = sharp_rank_two_two_columns for N = 2: la_weyl_r2n2 is (a), with
% la_binLat_stable (la_rowDiv_xy_zero and _one, choose_dvd_choose_mul_right and _left) and la_tw_stable
% (la_rowDiv_tw_ne and _eq); la_nabla_r2n2 is (b), with standardMonomialSpan,
% prod_initial_minors_mem_det_span, weylLeDual and dualWeylSat; la_core with la_arith is (c), and
% la_eExp_r2n2 identifies E. The exponent E is eExp (lamU N N nu) l, a sum over posPairs N of maxVal.
% The Lean states it for r x n matrices with n <= r <= 2; the book takes r = n = N.
\end{proof}

\begin{lemma}\label{is-onto-step-exp-bound-interval}
Assume the exact cut bound of Lemma~\ref{is-onto-step-exp-cut}. Let $N\geq0$, let $\lambda$ be a
sequence of natural numbers whose first $N$ terms are nonincreasing, $l$ a prime, and $\mathrm{lo}$ and
$w\geq0$ integers with $\mathrm{lo}\le\lambda_a+N-1-a\le\mathrm{lo}+w$ for $a<N$. Then
$e_l(N,\lambda)\le\binom N2\lfloor\log_lw\rfloor$ (Definition~\ref{is-onto-step-exp-cut-exponent},
with $\lfloor\log_lw\rfloor$ as in Lemma~\ref{is-onto-step-exp-cut-beta}). In particular
$e_l(N,\lambda)=0$ if $w<l$.
\end{lemma}
\leanref{Formalism/PieriBlockIdeals/PieriColumnMinors/RecursionInclusionExceptional.lean}{Formalism.CohomologicalDescent.E.exZ_le_of_exCut}
\leanref{Formalism/PieriBlockIdeals/PieriColumnMinors/RecursionInclusionExceptional.lean}{Formalism.CohomologicalDescent.E.exIntervalLt_of_exCut}
\leanref{Formalism/PieriBlockIdeals/PieriColumnMinors/RecursionInclusionExceptional.lean}{Formalism.CohomologicalDescent.E.exZ_rank_le_one}
\begin{proof}
First, $e_l(N,\lambda)=0$ if $N\le1$, by Lemma~\ref{is-onto-step-exp-bound-total-rank-two}: its
exponent $E$ is a sum over the pairs $a<b\le N$, of which there are none.

The bound is proved by induction on $N$, for all sequences; the case $N=0$ is done. Let $N\geq1$, and
put $\lambda'=(\lambda_{1+i})_{i\geq0}$. Lemma~\ref{is-onto-step-exp-cut} with $a=1$ gives
\[
 e_l(N,\lambda)\le e_l(1,\lambda)+e_l(N-1,\lambda')+(N-1)\lfloor\log_lw\rfloor,
\]
and $e_l(1,\lambda)=0$. The first $N-1$ terms of $\lambda'$ are nonincreasing, and for $a<N-1$,
$\lambda'_a+(N-1)-1-a=\lambda_{1+a}+N-1-(1+a)$ lies in $[\mathrm{lo},\mathrm{lo}+w]$. By induction
$e_l(N-1,\lambda')\le\binom{N-1}2\lfloor\log_lw\rfloor$, and $\binom{N-1}2+N-1=\binom N2$. If $w<l$, then
$\lfloor\log_lw\rfloor=0$.
% Lean: exZ_le_of_exCut, by induction on N with ExCut at a = 1, exZ_rank_le_one and res_anti_shift;
% exIntervalLt_of_exCut with Nat.log_of_lt. For N <= 1, exZ_rank_le_one applies LA.lattice_z_rank_le_two
% (Lemma is-onto-step-exp-bound-total-rank-two), whose exponent eExp is a sum over the empty posPairs
% N, and Nat.sInf_le.
\end{proof}

\begin{lemma}\label{is-onto-step-exp-bound-beyond}
Assume the exact cut bound of Lemma~\ref{is-onto-step-exp-cut}. Let $P$ be a parameter record that
satisfies the hypotheses and is large, $W$ weights for $P$ that satisfy the weight conditions, and $l$ a
prime with $l>l_J$ and $l>h$ (Definition~\ref{is-onto-step-exp-bound-weights}). Then $e_\tau(l)=0$ for
every $\tau$, and $J_A(l)=J^{\mathrm s}_W(l)=0$ (Definition~\ref{is-onto-step-exp-bound-recursion}).
\end{lemma}
\leanref{Formalism/PieriBlockIdeals/PieriColumnMinors/RecursionInclusionExceptional.lean}{Formalism.CohomologicalDescent.E.beyond_height_exceptional}
\leanref{Formalism/PieriBlockIdeals/PieriColumnMinors/RecursionInclusionExceptional.lean}{Formalism.CohomologicalDescent.E.thick_exceptional}
\leanref{Formalism/PieriBlockIdeals/PieriColumnMinors/RecursionInclusionExceptional.lean}{Formalism.CohomologicalDescent.E.thin_exceptional}
\begin{proof}
Let $\tau$ be given. As $N_r+1=k+m+t$, $N_r\geq1$, and by (W1), $U_\tau(N_r-1)\le U_\tau(a)\le U_\tau(0)$
for $a<N_r$.

Let $\tau\notin T_W$. In the first case of (W6) the width at $\tau$ is at most
$S+b'(1+x^{-1}+x^{-2})+\eta+(t-1)/2\le2b'$, by largeness. In the second case it is at most
$S+w''+\eta+(t-1)/2$, which is at most the same number, as $w''\le b'$
(Lemma~\ref{is-onto-step-exp-bound-numbers}) and $x>0$. So the width is at most $2b'<3b'<l_J<l$, as
$b'>0$. Lemma~\ref{is-onto-step-exp-bound-interval}, with $\mathrm{lo}=U_\tau(N_r-1)$ and $w$ the width,
gives $e_\tau(l)=0$.

Let $\tau\in T_W$. We show $0\le U_\tau(a)\le h_0$ for $a<N_r$. Then, as $h_0\le h<l$,
Lemma~\ref{is-onto-step-exp-bound-interval} with $\mathrm{lo}=0$ and $w=h_0$ gives $e_\tau(l)=0$. First,
$U_\tau(a)\geq0$, as $\lambda_{\tau,a}\geq0$ and $a\le N_r-1$. By (W5), $x^+_\tau\le h$, and $x^+_\tau$ is
an integer, so $x^+_\tau\le h_0$. If $a<k-1$, then $U_\tau(a)\le x^+_\tau$ by (W3). Let $k-1\le a<N_r$.
Then $|U_\tau(a)-c_\tau+e_\tau|\le\eta$: if $p_\tau\le a<p_\tau+t$, then by (W2)
$U_\tau(a)=\mathrm{bot}_\tau+(p_\tau+t-1-a)$, so $U_\tau(a)-c_\tau+e_\tau=(p_\tau+t-1-a)-(t-1)/2$
lies in $[-(t-1)/2,(t-1)/2]$, and $(t-1)/2\le t\le\eta$; otherwise this is (W4). By (W5)
$e_\tau\geq b$, by (W3) $c_\tau\le x^-_\tau$, and $2\eta<b$ by
Lemma~\ref{is-onto-step-exp-bound-numbers}. So $U_\tau(a)\le c_\tau-e_\tau+\eta\le x^-_\tau-b+\eta<
x^-_\tau\le x^+_\tau$.

Hence $e_\tau(l)=0$ for every $\tau$, so $J_A(l)=0$, and each term of $J^{\mathrm s}_W(l)$ is at most
$e_\tau(l)=0$.
% Lean: beyond_height_exceptional: thick_exceptional (first clause; shape_a_real_holds gives
% 0 <= U a <= h0 through weight_nonneg', weight_le_upper_translation_of_thick,
% weight_lt_lower_translation_of_thick, abs_weight_sub_near_translation_le, two_eta_lt_b and
% upper_translation_le_floor_height) and thin_exceptional (first clause; weight_in_width_range,
% width_nonneg) for tau not in T2, with paramsTop_holds, paramsThin_holds and Large.top_le; then JA and
% JW are sums of zeros. The interval bound enters as ExIntervalLt (exIntervalLt_of_exCut).
\end{proof}

\begin{lemma}\label{is-onto-step-exp-bound-compare}
Assume the exact cut bound of Lemma~\ref{is-onto-step-exp-cut}. Let $P$ be a parameter record that
satisfies the hypotheses and is large, $W$ weights for $P$ that satisfy the weight conditions, with
recursive exponent $J'$, and $l$ a prime (Definitions~\ref{is-onto-step-exp-bound-weights}
and~\ref{is-onto-step-exp-bound-recursion}). If $l$ is not in the range of $W$, then $J_A(l)\le J'(l)$.
If $l$ is in the range of $W$, then $J_W(l)=J'(l)$.
\end{lemma}
\leanref{Development/AbelianVarietyTheory/Proofs/DivisorProofs/ExponentOfRecordVsFormulas.lean}{GlobalTheory.exponent_of_record_vs_formulas}
\begin{proof}
If $l$ is in the range of $W$, then $l>l_J$, $l$ is not excluded by $P$ and $l\notin\mathrm{Sp}$, so
$J'(l)=J_W(l)$ by definition. Let $l$ not be in the range. If $l\le l_J$, $l$ is excluded or
$l\in\mathrm{Sp}$, then $J'(l)=J_A(l)$. Otherwise $l>l_J$, $l$ is not excluded and $l\notin\mathrm{Sp}$,
so $l>h$, as $l$ is not in the range; Lemma~\ref{is-onto-step-exp-bound-beyond} gives $J_A(l)=0\le
J'(l)$.
% Lean: exponent_of_record_vs_formulas, with seventh_recursion_eq_a_of_small_prime,
% seventh_recursion_eq_w_of_large_prime and E.beyond_height_exceptional (with exIntervalLt_of_exCut).
\end{proof}

\begin{lemma}\label{is-onto-step-exp-bound-tail}
Let $P$ be a parameter record that satisfies the hypotheses and has rows $(X',A)$ for a pair $(X',A)$
that is big for $P$ (Definition~\ref{is-onto-step-exp-bound-weights}), and let $T\le f$ be a natural
number. Put
\[
 \mathrm{Tot}(T)=T\big(30mbL^2+2k(S+1)(L+2)\big)+1.3fN_r^2b'+(S+\eta+t+2\kappa_{\max}+6)fN_r^2L.
\]
Then $\mathrm{Tot}(T)\le30TmbL^2+2mb$.
\end{lemma}
\leanref{Formalism/PieriBlockIdeals/TailInequalityEventually.lean}{Formalism.CohomologicalDescent.LargeParameterEstimates.TailInequalitySum}
\leanref{Formalism/PieriBlockIdeals/PieriColumnMinors/LargeParameterEstimates.lean}{Formalism.CohomologicalDescent.LargeParameterEstimates.tail_bound_holds}
\begin{proof}
This is a computation from the hypotheses, the rows and the bounds that make $(X',A)$ big: the terms of
$\mathrm{Tot}(T)$ other than $30TmbL^2$ have sum at most $4b$, and $4b\le2mb$ as $m\geq2$.
% Lean: tail_bound_holds, from sum_of_small_terms_bound (small_term_bound_first, _second and _third,
% with large_parameter_facts_of_hyp); the bound for the three terms is SmallSum.
\end{proof}

\begin{lemma}\label{is-onto-step-exp-bound-total-numbers}
Let $P$ be a parameter record that satisfies the hypotheses and is large
(Definition~\ref{is-onto-step-exp-bound-weights}). Then $N_r\le2t-5$, $b+7\le h$, $b'<b$,
$8<l_J\le b$, $\log l_J\le L$, $3b<l_J^2$, $h+N_r\le3b$, $0<\log h\le L+0.7$ and
$\log(h+N_r)\le L+1.1$. Every integer $z$ with
$b-(t-1)/2-\kappa_{\max}\le z\le b+S+\eta+1+(t-1)/2+\kappa_{\max}$ satisfies $0<z<l_J^2$, and every
integer $l>l_J$ satisfies $S+k<l$, $t+2\kappa(l)<l$, $N_r+1<l$ and $h+N_r<l^2$.
\end{lemma}
\leanref{Formalism/PieriBlockIdeals/PieriColumnMinors/ParameterBounds.lean}{Formalism.CohomologicalDescent.LogSmallConstantBounds.natBounds}
\leanref{Formalism/PieriBlockIdeals/PieriColumnMinors/ParameterBounds.lean}{Formalism.CohomologicalDescent.LogSmallConstantBounds.b_add_seven_le_h}
\leanref{Formalism/PieriBlockIdeals/PieriColumnMinors/ParameterBounds.lean}{Formalism.CohomologicalDescent.LogSmallConstantBounds.smaller_exponent_lt}
\leanref{Formalism/PieriBlockIdeals/PieriColumnMinors/ParameterBounds.lean}{Formalism.CohomologicalDescent.LogSmallConstantBounds.eight_lt_lJ}
\leanref{Formalism/PieriBlockIdeals/PieriColumnMinors/ParameterBounds.lean}{Formalism.CohomologicalDescent.LogSmallConstantBounds.paramsLog_holds}
\leanref{Formalism/PieriBlockIdeals/PieriColumnMinors/ParameterBounds.lean}{Formalism.CohomologicalDescent.LogSmallConstantBounds.log_lJ_bounds}
\leanref{Formalism/PieriBlockIdeals/PieriColumnMinors/ParameterBounds.lean}{Formalism.CohomologicalDescent.LogSmallConstantBounds.paramsLJ_holds}
\leanref{Formalism/PieriBlockIdeals/PieriColumnMinors/ParameterBounds.lean}{Formalism.CohomologicalDescent.LogSmallConstantBounds.log_h_bounds}
\leanref{Formalism/PieriBlockIdeals/PieriColumnMinors/ParameterBounds.lean}{Formalism.CohomologicalDescent.LogSmallConstantBounds.h_add_Nr_bounds}
\leanref{Formalism/PieriBlockIdeals/PieriColumnMinors/ParameterBounds.lean}{Formalism.CohomologicalDescent.LogSmallConstantBounds.exceptionalRange_holds}
\begin{proof}
The hypotheses give $k\geq m\geq1$, $t\geq k+m+4\geq6$, $\eta\geq t$, $S\geq0$ and $N_r+1=k+m+t$. So
$N_r\le2t-5$ and $h=b+S+\eta+t+1\geq b+7$. As $12kt\geq1$, largeness gives $b'\le12ktb'<b$.

By Lemma~\ref{is-onto-step-exp-bound-numbers}, $l_J>\ell^*+1>8$. As $3h\geq b$, $L=\log b\le\log(3h)$,
and $m\le k(m+t)$; so $mbL\le b\,k(m+t)\log(3h)$, that is, $l_J\le b$, and $\log l_J\le L$. By
Lemma~\ref{is-onto-step-exp-bound-numbers}, $b/(4kt)\le l_J$, and largeness gives $48k^2t^2<b$; so
$l_J^2\geq b^2/(16k^2t^2)>3b$.

The hypotheses give $S+\eta<\ell^*$ and $2\eta+t+8b/\ell^*+7<\ell^*$, so $3t<\ell^*$ and
$t+2b/\ell^*<\ell^*$, as $t\le\eta$; and $2\ell^*\le b$ by Lemma~\ref{is-onto-step-exp-bound-numbers}.
Hence $h+N_r\le b+(S+\eta)+3t-4<b+2\ell^*\le2b\le3b$. As $1<b+7\le h\le2b$, $0<\log h\le\log2+L\le L+0.7$, and
$\log(h+N_r)\le\log(3b)=\log3+L\le L+1.1$.

As $\kappa_{\max}\le2b/\ell^*$ and $(t-1)/2+1\le t$, the upper end of the interval for $z$ is at most
$b+(S+\eta)+t+2b/\ell^*<b+2\ell^*\le2b<l_J^2$, and its lower end is at least $b-t-2b/\ell^*>b-\ell^*\geq
b/2>0$. Finally, let $l>l_J$ be an integer. Then $k\le t\le\eta$ gives $S+k\le S+\eta<\ell^*<l_J<l$. By
Lemma~\ref{is-onto-step-exp-bound-numbers}, $\kappa(l)\le2b/\ell^*$ and $\ell^*\le l-1$; as
$t+4b/\ell^*<\ell^*$ by the last hypothesis, $t+2\kappa(l)<\ell^*<l$. As $m\le k\le t$,
$N_r+1=k+m+t\le3t<\ell^*<l$. As $l>l_J>0$, $l^2>l_J^2>3b\geq h+N_r$.
% Lean: natBounds (nr2), b_add_seven_le_h, smaller_exponent_lt, eight_lt_lJ (lstar_add_one_lt_lJ,
% seven_lt_lstar), paramsLog_holds (lJ <= b, through log_three_mul_h_bounds) and log_lJ_bounds,
% paramsLJ_holds (b_div_le_lJ with Large.b_gt for 3b < lJ^2, size_bounds_of_large for h + Nr <= 3b and
% for the upper end), log_h_bounds, h_add_Nr_bounds, exceptionalRange_holds, and the clauses S + k < l,
% t + 2 kap l < l, Nr + 1 < l and h + Nr < l^2 of paramsLargePrime_holds (linarith and
% mul_self_lt_mul_self). The Lean states h + Nr <= 3b and the upper end <= 3b.
\end{proof}

\begin{lemma}\label{is-onto-step-exp-bound-total-primes}
For a real number $y$ let $\vartheta(y)=\sum_{p\le y}\log p$, the sum over the primes $p\le y$.
\begin{enumerate}
\item $\vartheta(y)\le y\log4$ for $y\geq0$.
\item $\sum_{p\le y}\log p/p\le\log y+\log4$ for $y\geq1$.
\item $\sum_{i=1}^n1/i\le1+\log n$ for every integer $n\geq1$, and $\sum_{p\le y}1/p\le\log y$ for
  $y\geq1$.
\end{enumerate}
\end{lemma}
\leanref{Formalism/PieriBlockIdeals/PieriColumnMinors/ParameterBounds.lean}{Formalism.CohomologicalDescent.LogSmallConstantBounds.chebyshevThetaElementary_holds}
\leanref{Formalism/PieriBlockIdeals/PieriColumnMinors/ParameterBounds.lean}{Formalism.CohomologicalDescent.LogSmallConstantBounds.mertensLogElementary_holds}
\leanref{Formalism/PieriBlockIdeals/PieriColumnMinors/ParameterBounds.lean}{Formalism.CohomologicalDescent.LogSmallConstantBounds.mertensReciprocalElementary_holds}
\leanref{Formalism/PieriBlockIdeals/PieriColumnMinors/SharperSizeNeedsMleKFalse.lean}{Formalism.CohomologicalDescent.sums_harmonic_real}
\begin{proof}
(1) For an integer $n\geq1$, the product of the primes at most $n$ is less than $4^n$, that is,
$\vartheta(n)<2n\log2$ \cite[Theorem 415]{HW08}; for $n=0$ the product is $1=4^0$. Take
$n=\lfloor y\rfloor\le y$.

(2) Let $n=\lfloor y\rfloor\geq1$. As $v_p(n!)\geq\lfloor n/p\rfloor$,
$\sum_{p\le n}\lfloor n/p\rfloor\log p\le\log n!\le n\log n$. As $n/p\le\lfloor n/p\rfloor+1$, (1) gives
$n\sum_{p\le n}\log p/p\le n\log n+\vartheta(n)\le n(\log n+\log4)$, and $\log n\le\log y$. This is a
form of Mertens' first theorem \cite{Mer74}.

(3) For $i\geq2$, $1/i\le\int_{i-1}^i dx/x$, so $\sum_{i=1}^n1/i\le1+\log n$. With $n=\lfloor y\rfloor$,
the primes at most $y$ are among $2,\dots,n$, so $\sum_{p\le y}1/p\le\sum_{i=2}^n1/i\le\log n\le\log y$.
% Lean: chebyshevThetaElementary_holds (ConvolutionKernelIntegrable) from chebyshevThetaLe_holds, which is
% Mathlib's Chebyshev.theta_le_log4_mul_x; mertensLogElementary_holds (MertensLogElementary) through
% sum_log_div_le_log_add_log_four and sum_div_mul_log_le_log_factorial' (padicValNat_factorial);
% mertensReciprocalElementary_holds (MertensReciprocalElementary) with sums_harmonic_real, which is
% Mathlib's harmonic_le_one_add_log.
\end{proof}

\begin{lemma}\label{is-onto-step-exp-bound-total-excluded}
Let $P$ be a parameter record that satisfies the hypotheses and is large, and $Q$ a finite set of
primes $l>l_J$ that are excluded by $P$ (Definition~\ref{is-onto-step-exp-bound-weights}). Then
$|Q|\le S+\eta+t+2\kappa_{\max}+1$.
\end{lemma}
\leanref{Formalism/PieriBlockIdeals/MertensReciprocalElementary.lean}{Formalism.CohomologicalDescent.LogSmallConstantBounds.ExceptionalCount}
\leanref{Formalism/PieriBlockIdeals/PieriColumnMinors/ParameterBounds.lean}{Formalism.CohomologicalDescent.LogSmallConstantBounds.exceptionalCount_holds}
\leanref{Formalism/PieriBlockIdeals/PieriColumnMinors/LogarithmEquivalence.lean}{Formalism.CohomologicalDescent.largePrimeDivisorCount_holds}
\begin{proof}
Let $I$ be the set of integers $z$ with $b-(t-1)/2-\kappa_{\max}\le z\le
b+S+\eta+1+(t-1)/2+\kappa_{\max}$. This interval has length $S+\eta+t+2\kappa_{\max}$, so
$|I|\le S+\eta+t+2\kappa_{\max}+1$. Every $l\in Q$ divides some $z_l\in I$, and $0<z_l<l_J^2$ by
Lemma~\ref{is-onto-step-exp-bound-total-numbers}. If $z_l=z_{l'}$ with $l\ne l'$, the distinct primes
$l,l'>l_J$ both divide $z_l$, and $l_J^2<ll'\le z_l<l_J^2$, which is absurd. So $l\mapsto z_l$ is
injective, and $|Q|\le|I|$.
% Lean: exceptionalCount_holds: the z_l lie in Finset.Icc (ceil of the lower end) (floor of the upper
% end), exceptionalRange_holds gives 0 < z < lJ^2, and largePrimeDivisorCount_holds is the injectivity
% (eq_of_dvd_of_abs_lt_sq), with Int.card_Icc for the count.
\end{proof}

\begin{lemma}\label{is-onto-step-exp-bound-total-no-window}
Let $P$ be a parameter record that satisfies the hypotheses and is large, $W$ weights for $P$ that
satisfy the weight conditions (Definition~\ref{is-onto-step-exp-bound-weights}), and $l>l_J$ a prime
that is not excluded by $P$. Then the window counts $n_\tau(l)$ and $n^{\mathrm F}_\tau(l)$ vanish
for every $\tau$ (Definition~\ref{is-onto-step-exp-bound-recursion}(2)).
\end{lemma}
\leanref{Formalism/PieriBlockIdeals/MertensReciprocalElementary.lean}{Formalism.CohomologicalDescent.LogSmallConstantBounds.SizeData₂.NoWindow}
\leanref{Formalism/PieriBlockIdeals/PieriColumnMinors/ParameterBounds.lean}{Formalism.CohomologicalDescent.LogSmallConstantBounds.noWindowOf_holds}
\leanref{Formalism/PieriBlockIdeals/PieriColumnMinors/WeightBounds.lean}{Formalism.CohomologicalDescent.ShortWeightBounds.places_real_holds}
\begin{proof}
Let $\tau$ be given. The set $\mathcal U_\tau$, $\beta=\mathrm{bot}_\tau$, $c_\tau$, $e_\tau$ and
$\sigma=1$ satisfy the conditions of Lemma~\ref{is-onto-step-exp-bound-window} on $\mathcal U$. Indeed,
(W2) gives $\mathrm{bot}_\tau+(t-1)/2-c_\tau=-e_\tau$. By (W1) the $U_\tau(a)$, $a<N_r$, decrease
strictly, and by (W2) $U_\tau(a)=\mathrm{bot}_\tau+(p_\tau+t-1-a)$ for $p_\tau\le a<p_\tau+t$; so if
$U_\tau(a)$ lies outside $[\mathrm{bot}_\tau,\mathrm{bot}_\tau+t-1]$, then $a<p_\tau$ or $a\geq p_\tau+t$.
If $a<k-1$, (W3) gives $0\le U_\tau(a)-c_\tau\le S$; otherwise (W4) gives
$2b/\ell^*+t/2+3\le|U_\tau(a)-c_\tau+e_\tau|\le\eta$. By (W5) and (W6), $e_\tau$ satisfies one of the
three conditions on $e$ of that lemma, the first one together with the fact that $l$ is not excluded.
So $n_\tau(l)=0$.

Let $\tau'=\phi_l(\tau)$. If $s_l(\tau)=0$, then $n^{\mathrm F}_\tau(l)=n_{\tau'}(l)=0$. Let $s_l(\tau)=1$.
The set $\{h_0-u:u\in\mathcal U_{\tau'}\}$, $\beta'=h_0-\mathrm{top}_{\tau'}$, $c'=h_0-c_{\tau'}$,
$e_{\tau'}$ and $\sigma=-1$ satisfy the conditions of Lemma~\ref{is-onto-step-exp-bound-window}: as
$\mathrm{top}_{\tau'}=\mathrm{bot}_{\tau'}+t-1$ by (W2),
$-(\beta'+(t-1)/2-c')=\mathrm{bot}_{\tau'}+(t-1)/2-c_{\tau'}=-e_{\tau'}$; $h_0-u$ lies in
$[\beta',\beta'+t-1]$ if and only if $u$ lies in $[\mathrm{bot}_{\tau'},\mathrm{bot}_{\tau'}+t-1]$; and
$-((h_0-u)-c')=u-c_{\tau'}$. So $n^{\mathrm F}_\tau(l)=0$.
% Lean: the second part of places_real_holds, which is noWindowOf_holds (noWindowResidue_holds, Lemma
% is-onto-step-exp-bound-window, at the shapes of placesOfParameters). The same argument is the second
% paragraph of the proof of Lemma is-onto-step-exp-bound-frobenius, where l is in the range.
\end{proof}

\begin{lemma}\label{is-onto-step-exp-bound-total-exponents}
Assume the exact cut bound of Lemma~\ref{is-onto-step-exp-cut}. Let $P$ be a parameter record that
satisfies the hypotheses and is large, $W$ weights for $P$ that satisfy the weight conditions, $\tau$ an
index, $w_\tau$ the width at $\tau$, and $l$ a prime (Definitions~\ref{is-onto-step-exp-bound-weights}
and~\ref{is-onto-step-exp-bound-recursion}).
\begin{enumerate}
\item If $\tau\notin T_W$, then $w_\tau\geq0$; $e_\tau(l)=0$ if $w_\tau<l$, and
  $e_\tau(l)\log l\le\binom{N_r}2\log w_\tau$ if $l\le w_\tau$.
\item If $\tau\in T_W$, then $e_\tau(l)=0$ if $h<l$, and $e_\tau(l)\log l\le\binom{N_r}2\log h$; if
  moreover $2\eta+t\le l$, then $e_\tau(l)\log l\le k(m+t)\log(3h)+k^2\log(S+k)$.
\item If $\tau\in T_W$ and $S+k<l$, then $e^{\mathrm c}_\tau(l)=0$.
\item If $\tau\in T_W$, then $0\le\mathrm{top}_\tau<x^-_\tau\le x^+_\tau\le h$,
  $x^+_\tau-x^-_\tau\le S$ and $q_\tau+j_\tau\le m$.
\end{enumerate}
\end{lemma}
\leanref{Formalism/PieriBlockIdeals/MertensReciprocalElementary.lean}{Formalism.CohomologicalDescent.LogSmallConstantBounds.SizeData₂.ThinA}
\leanref{Formalism/PieriBlockIdeals/MertensReciprocalElementary.lean}{Formalism.CohomologicalDescent.LogSmallConstantBounds.SizeData₂.ThickA}
\leanref{Formalism/PieriBlockIdeals/MertensReciprocalElementary.lean}{Formalism.CohomologicalDescent.LogSmallConstantBounds.SizeData₂.CloudA}
\leanref{Formalism/PieriBlockIdeals/MertensReciprocalElementary.lean}{Formalism.CohomologicalDescent.LogSmallConstantBounds.SizeData₂.ShapeT}
\leanref{Formalism/PieriBlockIdeals/PieriColumnMinors/RecursionInclusionExceptional.lean}{Formalism.CohomologicalDescent.E.thin_exceptional}
\leanref{Formalism/PieriBlockIdeals/PieriColumnMinors/RecursionInclusionExceptional.lean}{Formalism.CohomologicalDescent.E.thick_exceptional}
\leanref{Formalism/PieriBlockIdeals/PieriColumnMinors/RecursionInclusionExceptional.lean}{Formalism.CohomologicalDescent.E.cloud_exceptional}
\leanref{Formalism/PieriBlockIdeals/PieriColumnMinors/WeightBounds.lean}{Formalism.CohomologicalDescent.ShortWeightBounds.shape_t_real_holds}
\begin{proof}
Write $U(a)=U_\tau(a)$. By (W1) the first $N_r$ terms of $\lambda_\tau$ are nonincreasing, so $U$
decreases strictly on $[0,N_r)$; as $N_r\geq1$, $U(N_r-1)\le U(a)\le U(0)$ for $a<N_r$. For an integer
$y\geq0$ put $\Lambda(y)=\lfloor\log_ly\rfloor\log l$, with $\lfloor\log_ly\rfloor$ as in
Lemma~\ref{is-onto-step-exp-cut-beta}. Then $\Lambda(0)=0$, and $\Lambda(y)\le\log y$ for $y\geq1$, as
$l^{\lfloor\log_ly\rfloor}\le y$.

(1) The width $w_\tau=U(0)-U(N_r-1)$ is at least $0$, and $U(N_r-1)\le U(a)\le U(N_r-1)+w_\tau$ for
$a<N_r$. Lemma~\ref{is-onto-step-exp-bound-interval}, with $\mathrm{lo}=U(N_r-1)$ and $w=w_\tau$, gives
$e_\tau(l)\le\binom{N_r}2\lfloor\log_lw_\tau\rfloor$, and $e_\tau(l)=0$ if $w_\tau<l$. If $l\le w_\tau$,
then $e_\tau(l)\log l\le\binom{N_r}2\Lambda(w_\tau)\le\binom{N_r}2\log w_\tau$.

(2) As in the proof of Lemma~\ref{is-onto-step-exp-bound-beyond}, $0\le U(a)\le h_0$ for $a<N_r$, and
$|U(a)-c_\tau+e_\tau|\le\eta$ for $k-1\le a<N_r$. Lemma~\ref{is-onto-step-exp-bound-interval}, with
$\mathrm{lo}=0$ and $w=h_0$, gives $e_\tau(l)\le\binom{N_r}2\lfloor\log_lh_0\rfloor$, and $e_\tau(l)=0$ if
$h_0<l$, in particular if $h<l$. As $1\le h_0\le h$ (Lemma~\ref{is-onto-step-exp-bound-total-numbers}),
$e_\tau(l)\log l\le\binom{N_r}2\Lambda(h_0)\le\binom{N_r}2\log h$.

Let now $2\eta+t\le l$, and put $K=k-1$ and $s=x^+_\tau-x^-_\tau$, which is at least $0$ by (W3). For
$a<K$, (W3) gives $x^-_\tau\le U(a)\le x^+_\tau$, so $\lambda_{\tau,a}+K-1-a=U(a)-(N_r-K)$ lies in
$[x^-_\tau-(N_r-K),x^-_\tau-(N_r-K)+s]$. For $K\le a<N_r$, $U(N_r-1)\le U(a)\le U(K)$, and the bound
$|U(a)-c_\tau+e_\tau|\le\eta$ at $a=K$ and $a=N_r-1$ gives $U(K)-U(N_r-1)\le2\eta<l$, as $t\geq1$.
Lemma~\ref{is-onto-step-exp-cut}, with $a=K$, $\mathrm{lo}=0$ and $w=h_0$, gives
\[
 e_\tau(l)\le e_l(K,\lambda_\tau)+e_l(N_r-K,\lambda'')+K(N_r-K)\lfloor\log_lh_0\rfloor,\qquad
 \lambda''=(\lambda_{\tau,K+i})_{i\geq0}.
\]
Lemma~\ref{is-onto-step-exp-bound-interval} gives $e_l(K,\lambda_\tau)\le\binom K2\lfloor\log_ls\rfloor$.
For $a<N_r-K$, $\lambda''_a+(N_r-K)-1-a=U(K+a)$ lies in $[U(N_r-1),U(K)]$, an interval of length less
than $l$; so the same lemma gives $e_l(N_r-K,\lambda'')=0$. Hence
\[
 e_\tau(l)\log l\le\binom K2\Lambda(s)+K(N_r-K)\Lambda(h_0).
\]
Here $N_r-K=m+t$ and $\Lambda(h_0)\le\log h_0\le\log(3h)$. Also $\binom K2\le K^2$, and
$\Lambda(s)\le\log(S+k)$: if $s\geq1$, then $\Lambda(s)\le\log s\le\log(S+k)$, as $s\le S$ by (W3); if
$s=0$, then $\Lambda(s)=0\le\log(S+k)$, as $S+k\geq1$. As $K\le k$, the bound follows.

(3) As in (2), for $a<K$ the numbers $\lambda_{\tau,a}+K-1-a$ lie in an interval of length $s\le S<l$,
and the first $K$ terms of $\lambda_\tau$ are nonincreasing. Lemma~\ref{is-onto-step-exp-bound-interval}
with $N=K$ gives $e^{\mathrm c}_\tau(l)=e_l(K,\lambda_\tau)=0$.

(4) By (W1), $p_\tau+t+q_\tau=N_r$, so $p_\tau\le N_r-1$ and $\mathrm{top}_\tau=U(p_\tau)\geq0$. As
$p_\tau\geq k-1$, the proof of Lemma~\ref{is-onto-step-exp-bound-beyond} gives $U(p_\tau)<x^-_\tau$.
(W3) gives $x^-_\tau\le x^+_\tau$ and $x^+_\tau-x^-_\tau\le S$, and (W5) gives $x^+_\tau\le h$. Finally
$p_\tau=k-1+j_\tau$ and $N_r+1=k+m+t$ give $q_\tau+j_\tau=m$, so $q_\tau+j_\tau\le m$.
% Lean: (1) is thin_exceptional (interval bound ExIntervalLt at lo = U (Nr - 1) and w = width, with
% weight_in_width_range, width_nonneg and exponent_mul_log_le_choose_mul_log); (2) is thick_exceptional:
% shape_a_real_holds gives the Shape and 0 <= U <= h0, the first two clauses are the interval bound at
% lo = 0, w = h0, and the third is ExCut at a = kc P = K with the interval bounds for the first K rows
% (cloud) and the rest (rest_near_holds, U K - U (Nr - 1) <= 2 eta < l), log_le_log_bound (Lambda(y) <=
% log y) and the second part of SL2BinomialLattices.size_large (K (Nr - K) log h0 + C(K,2) log s <=
% K (m + t) log (3 h0) + K^2 log (s + K)), then kc_cast_le_k; (3) is cloud_exceptional; (4) is
% shape_t_real_holds (weight_nonneg', weight_lt_lower_translation_of_thick, window, thick, and
% j_add_q_eq_m weakened to q + j <= m by omega). The interval bounds come from ExCut by Lemma
% is-onto-step-exp-bound-interval.
\end{proof}

\begin{definition}\label{is-onto-step-exp-bound-total-real}
Let $P$ be a parameter record, $W$ weights for $P$ indexed by $\mathcal T$, $\tau\in\mathcal T$ and
$l\geq2$ an integer (Definitions~\ref{is-onto-step-exp-bound-weights}
and~\ref{is-onto-step-exp-bound-recursion}). The \emph{real depth} at $\tau$ is
\[
 D_\tau(l)=q_\tau h+\big(l\,n^{\mathrm F}_\tau(l)+n_\tau(l)\big)\big(t+2\kappa(l)\big),
\]
so that $i_\tau(l)=\lfloor D_\tau(l)\rfloor$. Let $\bar\gamma_\tau(l)=\min\big(k,t,D_\tau(l)/d_\tau(l)\big)$
if $d_\tau(l)<x^-_\tau-\mathrm{top}_\tau$, and $\bar\gamma_\tau(l)=0$ otherwise. The \emph{real coarse
bound} at $\tau$ is
\[
 \bar C_\tau(l)=(q_\tau+2j_\tau)\Big(\frac{h+N_r}{l-1}+6\Big)+e^{\mathrm c}_\tau(l)
 +k\Big(\frac Sl+\delta_\tau(l)\Big)+\bar\gamma_\tau(l)+\frac{D_\tau(l)}l\Big(1+\log^+\frac hl\Big).
\]
\end{definition}
\leanref{Formalism/PieriBlockIdeals/MertensReciprocalElementary.lean}{Formalism.CohomologicalDescent.LogSmallConstantBounds.SizeData₂.depth}
\leanref{Formalism/PieriBlockIdeals/MertensReciprocalElementary.lean}{Formalism.CohomologicalDescent.LogSmallConstantBounds.SizeData₂.cW}
\leanref{Formalism/PieriBlockIdeals/PieriColumnMinors/WeightBounds.lean}{Formalism.CohomologicalDescent.ShortWeightBounds.correction_weight_eq_formula}

\begin{lemma}\label{is-onto-step-exp-bound-total-weight-bound}
Let $P$ be a parameter record that satisfies the hypotheses and is large, $W$ weights for $P$ that
satisfy the weight conditions, and $l\geq2$ an integer (Definition~\ref{is-onto-step-exp-bound-weights}).
Then $\lfloor C_\tau(l)\rfloor\le\bar C_\tau(l)$ for every $\tau\in T_W$
(Definitions~\ref{is-onto-step-exp-bound-recursion} and~\ref{is-onto-step-exp-bound-total-real}), and
\[
 J_W(l)\le\sum_{\tau\in T_W}\min\big(e_\tau(l),\bar C_\tau(l)\big)+\sum_{\tau\notin T_W}e_\tau(l).
\]
Let $J'$ be the recursive exponent of $P$ and $W$, and $J$ a function from the natural numbers to
themselves with $J(p)\le J'(p)$ for every prime $p$. If $l$ is a prime with $l>l_J$, $l$ not excluded by
$P$, $l\notin\mathrm{Sp}$, $h+N_r<l^2$, $N_r+1<l$ and $t+2\kappa(l)<l$, then $J(l)$ is at most the right
side. The last three of these conditions are not used in the proof.
\end{lemma}
\leanref{Formalism/PieriBlockIdeals/LatticeInclusionRealSettingAlt.lean}{Formalism.CohomologicalDescent.ShortWeightBounds.CoarseLeCW}
\leanref{Formalism/PieriBlockIdeals/PieriColumnMinors/WeightBounds.lean}{Formalism.CohomologicalDescent.ShortWeightBounds.coarse_le_c_w_holds}
\leanref{Formalism/PieriBlockIdeals/PieriColumnMinors/WeightBounds.lean}{Formalism.CohomologicalDescent.ShortWeightBounds.short_bound_nonneg_holds}
\leanref{Formalism/PieriBlockIdeals/PieriColumnMinors/WeightBounds.lean}{Formalism.CohomologicalDescent.ShortWeightBounds.j_w_seven_le_holds}
\leanref{Formalism/PieriBlockIdeals/PieriColumnMinors/WeightBounds.lean}{Formalism.CohomologicalDescent.ShortWeightBounds.sizeData_boundW_of_le_seventh_recursion}
\begin{proof}
Let $\tau\in T_W$. Every term of $C^{\mathrm s}_\tau(l)$ is at least $0$, as $l\geq2$, $x^-_\tau\le
x^+_\tau$ by (W3), and $\delta_\tau(l)$, $\gamma_\tau(l)$ and $i_\tau(l)$ are at least $0$. So
$C_\tau(l)\geq C^{\mathrm s}_\tau(l)\geq0$, and $\lfloor C_\tau(l)\rfloor\le C_\tau(l)$. Compare
$C_\tau(l)$ with $\bar C_\tau(l)$ term by term: $h_0\le h$; $k-1\le k$ and $x^+_\tau-x^-_\tau\le S$
by (W3); $D_\tau(l)\geq0$ and $i_\tau(l)\le D_\tau(l)$, so $\gamma_\tau(l)\le\bar\gamma_\tau(l)$; and
$\log^+(h_0/l)\le\log^+(h/l)$. The terms $e^{\mathrm c}_\tau(l)$ agree. So $C_\tau(l)\le\bar C_\tau(l)$,
and $\min(e_\tau(l),\lfloor C_\tau(l)\rfloor)\le\min(e_\tau(l),\bar C_\tau(l))$. Summing over $\tau$
gives the bound for $J_W(l)$. Let $l$ be a prime with $l>l_J$, not excluded by $P$ and not in
$\mathrm{Sp}$. Then $J'(l)=J_W(l)$ by Definition~\ref{is-onto-step-exp-bound-recursion}(5), so
$J(l)\le J'(l)=J_W(l)$, which is at most the right side.
% Lean: the second part of j_w_seven_le_holds (CoarseWeightBoundLe), from coarse_le_c_w_holds (with
% correction_weight_eq_formula, sizeData_depth_nonneg, floor_height_le_height, kc_cast_le_k and
% Real.posLog_le_posLog), short_bound_nonneg_holds and short_bound_le_coarse_holds; then
% sizeData_boundW_of_le_seventh_recursion (BoundW for the size data W.sizeData₂ P J) with
% seventh_recursion_eq_w_of_large_prime. BoundW's hypotheses h + Nr < l^2, Nr + 1 < l and
% t + 2 kap l < l are introduced as _ and never used.
\end{proof}

\begin{lemma}\label{is-onto-step-exp-bound-total-cost}
Assume the exact cut bound of Lemma~\ref{is-onto-step-exp-cut}. Let $P$ be a parameter record that
satisfies the hypotheses and is large, $W$ weights for $P$, indexed by a set of $f$ elements, that
satisfy the weight conditions (Definition~\ref{is-onto-step-exp-bound-weights}), and $l$ a prime. Then
$e_\tau(l)\log l\le N_r^2L$ for every $\tau$, and $\big(\sum_\tau e_\tau(l)\big)\log l\le fN_r^2L$
(Definition~\ref{is-onto-step-exp-bound-recursion}).
\end{lemma}
\leanref{Formalism/PieriBlockIdeals/MertensReciprocalElementary.lean}{Formalism.CohomologicalDescent.LogSmallConstantBounds.CostB}
\leanref{Formalism/PieriBlockIdeals/PieriColumnMinors/ParameterBounds.lean}{Formalism.CohomologicalDescent.LogSmallConstantBounds.costBound_holds}
\leanref{Formalism/PieriBlockIdeals/PieriColumnMinors/ParameterBounds.lean}{Formalism.CohomologicalDescent.LogSmallConstantBounds.ex_mul_log_le_of_thickA}
\leanref{Formalism/PieriBlockIdeals/PieriColumnMinors/ParameterBounds.lean}{Formalism.CohomologicalDescent.LogSmallConstantBounds.ex_mul_log_le_of_thinA}
\begin{proof}
Note $\binom{N_r}2\le N_r^2/2$ and $L\geq4$. Let $\tau\in T_W$. By
Lemmas~\ref{is-onto-step-exp-bound-total-exponents}(2) and~\ref{is-onto-step-exp-bound-total-numbers},
$e_\tau(l)\log l\le\binom{N_r}2\log h\le\frac{N_r^2}2(L+0.7)\le N_r^2L$. Let $\tau\notin T_W$. As in the
proof of Lemma~\ref{is-onto-step-exp-bound-beyond}, the width $w_\tau$ is at most $2b'$. If $w_\tau<l$,
then $e_\tau(l)=0$ by Lemma~\ref{is-onto-step-exp-bound-total-exponents}(1). Otherwise $2\le l\le
w_\tau\le2b'<2b$ (Lemma~\ref{is-onto-step-exp-bound-total-numbers}), so $\log w_\tau\le\log2+L\le2L$, and
Lemma~\ref{is-onto-step-exp-bound-total-exponents}(1) gives $e_\tau(l)\log l\le\frac{N_r^2}2\cdot2L$.
Summing over the $f$ indices gives the second bound.
% Lean: costBound_holds, through ex_mul_log_le_of_thickA (log_h_pos_and_le, choose_two_le_sq_div_two) and
% ex_mul_log_le_of_thinA (hw_le_two_mul_b', smaller_exponent_lt).
\end{proof}

\begin{lemma}\label{is-onto-step-exp-bound-total-thin}
Assume the exact cut bound of Lemma~\ref{is-onto-step-exp-cut}, and that the number of primes at most
$y$ is less than $1.26\,y/\log y$ for every real $y>1$. Let $P$ be a parameter record that satisfies
the hypotheses and is large, $W$ weights for $P$ that satisfy the weight conditions, and $\tau\notin T_W$
(Definition~\ref{is-onto-step-exp-bound-weights}). Then for every finite set $Q$ of primes,
\[
 \sum_{l\in Q}e_\tau(l)\log l\le1.3N_r^2b'.
\]
\end{lemma}
\leanref{Formalism/PieriBlockIdeals/MertensReciprocalElementary.lean}{Formalism.CohomologicalDescent.LogSmallConstantBounds.SumA}
\leanref{Formalism/PieriBlockIdeals/PieriColumnMinors/ParameterBounds.lean}{Formalism.CohomologicalDescent.LogSmallConstantBounds.thinSum_holds}
\begin{proof}
Let $w=w_\tau$ be the width at $\tau$. As in the proof of Lemma~\ref{is-onto-step-exp-bound-total-cost},
$0\le w\le2b'$, and by Lemma~\ref{is-onto-step-exp-bound-total-exponents}(1), $e_\tau(l)=0$ for $l>w$.
If $w\le1$, every term vanishes, as $l\geq2$. Let $w>1$. The remaining terms have $l\le w$, each is at
most $\binom{N_r}2\log w$, and there are fewer than $1.26\,w/\log w$ of them. So
\[
 \sum_{l\in Q}e_\tau(l)\log l\le1.26\,\frac w{\log w}\binom{N_r}2\log w=1.26\,w\binom{N_r}2
 \le1.26\cdot2b'\cdot\frac{N_r^2}2\le1.3N_r^2b'.
\]
% Lean: thinSum_holds (SumA), with hw_le_two_mul_b' and choose_two_le_sq_div_two; the primes of Q up to
% the width form a subset of pr (width) (subset_primesLe).
\end{proof}

\begin{lemma}\label{is-onto-step-exp-bound-total-small}
Assume the exact cut bound of Lemma~\ref{is-onto-step-exp-cut}. Let $P$ be a parameter record that
satisfies the hypotheses and is large, $W$ weights for $P$ that satisfy the weight conditions, and
$\tau\in T_W$ (Definition~\ref{is-onto-step-exp-bound-weights}).
\begin{enumerate}
\item For every prime $l\le l_J$, $e_\tau(l)\log l\le2mbL/l$.
\item For every finite set $Q$ of primes at most $l_J$, $\sum_{l\in Q}e_\tau(l)\log l\le2mbL^2$.
\end{enumerate}
\end{lemma}
\leanref{Formalism/PieriBlockIdeals/MertensReciprocalElementary.lean}{Formalism.CohomologicalDescent.LogSmallConstantBounds.TermC}
\leanref{Formalism/PieriBlockIdeals/PieriColumnMinors/ParameterBounds.lean}{Formalism.CohomologicalDescent.LogSmallConstantBounds.thickTerm_holds}
\leanref{Formalism/PieriBlockIdeals/MertensReciprocalElementary.lean}{Formalism.CohomologicalDescent.LogSmallConstantBounds.SumCE}
\leanref{Formalism/PieriBlockIdeals/PieriColumnMinors/ParameterBounds.lean}{Formalism.CohomologicalDescent.LogSmallConstantBounds.smallPrimeSumElementary_holds}
\begin{proof}
(1) Suppose first that $2\eta+t\le l$. By Lemma~\ref{is-onto-step-exp-bound-total-exponents}(2),
$e_\tau(l)\log l\le k(m+t)\log(3h)+k^2\log(S+k)$. Here $k^2\le k(m+t)$, as $k\le t$, and $1\le S+k\le
S+t\le3h$; so $e_\tau(l)\log l\le2k(m+t)\log(3h)$. By the definition of $l_J$,
$l_J\,k(m+t)\log(3h)=mbL$, so $l\le l_J$ gives $l\,e_\tau(l)\log l\le2mbL$.

Suppose now that $l<2\eta+t$, so that $l<3\eta$, as $t\le\eta$. By
Lemmas~\ref{is-onto-step-exp-bound-total-exponents}(2) and~\ref{is-onto-step-exp-bound-total-numbers},
$e_\tau(l)\log l\le\binom{N_r}2\log h\le\frac{N_r^2}2\cdot2L\le4t^2L$, as $N_r\le2t$ and $\log h\le
L+0.7\le2L$. So $l\,e_\tau(l)\log l\le12t^2L\eta$. The hypotheses give $\eta t^3fk^3\le b$, and
$mftk^3\geq6$, as $m,f,k\geq1$ and $t\geq6$. So $12t^2L\eta\le2t^2L\eta\cdot mftk^3=2mL\cdot\eta
t^3fk^3\le2mbL$.

(2) By Lemma~\ref{is-onto-step-exp-bound-total-numbers}, $l_J>8$ and $\log l_J\le L$. By (1) and
Lemma~\ref{is-onto-step-exp-bound-total-primes}(3),
\[
 \sum_{l\in Q}e_\tau(l)\log l\le2mbL\sum_{p\le l_J}\frac1p\le2mbL\log l_J\le2mbL^2.
\]
% Lean: thickTerm_holds (TermC), with the third clause of ThickA in the first case (indexDenominator, the
% denominator of lJ) and the second clause with H.eta_le in the other (natBounds); then
% smallPrimeSumElementary_holds (SumCE) with mertensReciprocalElementary_holds and log_lJ_bounds.
\end{proof}

\begin{lemma}\label{is-onto-step-exp-bound-total-large}
Assume the exact cut bound of Lemma~\ref{is-onto-step-exp-cut}. Let $P$ be a parameter record that
satisfies the hypotheses and is large, $W$ weights for $P$ that satisfy the weight conditions,
$\tau\in T_W$, and $Q$ a finite set of primes $l>l_J$ that are not excluded by $P$
(Definition~\ref{is-onto-step-exp-bound-weights}). Then
\[
 \sum_{l\in Q}\min\big(e_\tau(l),\bar C_\tau(l)\big)\log l\le18mbL^2+2k(S+1)(L+2),
\]
with $\bar C_\tau(l)$ as in Definition~\ref{is-onto-step-exp-bound-total-real}.
\end{lemma}
\leanref{Formalism/PieriBlockIdeals/MertensReciprocalElementary.lean}{Formalism.CohomologicalDescent.LogSmallConstantBounds.SumDE}
\leanref{Formalism/PieriBlockIdeals/PieriColumnMinors/ParameterBounds.lean}{Formalism.CohomologicalDescent.LogSmallConstantBounds.largePrimeSumElementary_holds}
\leanref{Formalism/PieriBlockIdeals/PieriColumnMinors/ParameterBounds.lean}{Formalism.CohomologicalDescent.LogSmallConstantBounds.sum_min_cW_le_split}
\leanref{Formalism/PieriBlockIdeals/PieriColumnMinors/LogarithmEquivalence.lean}{Formalism.CohomologicalDescent.residueIndicatorPrimeSum_holds}
\leanref{Formalism/PieriBlockIdeals/PieriColumnMinors/SharperSizeNeedsMleKFalse.lean}{Formalism.CohomologicalDescent.gatedPrimeSum_holds}
\begin{proof}
\emph{Splitting.} Let $R$ be the set of $l\in Q$ with $l\le h$. If $l\in Q\smallsetminus R$, then
$e_\tau(l)=0$ by Lemma~\ref{is-onto-step-exp-bound-total-exponents}(2), and the term at $l$ is
$\min(0,\bar C_\tau(l))\log l\le0$. Let $l\in R$. By Lemma~\ref{is-onto-step-exp-bound-total-numbers},
$S+k<l$, so $e^{\mathrm c}_\tau(l)=0$ by Lemma~\ref{is-onto-step-exp-bound-total-exponents}(3); and
$n_\tau(l)=n^{\mathrm F}_\tau(l)=0$ by Lemma~\ref{is-onto-step-exp-bound-total-no-window}, so
$D_\tau(l)=q_\tau h$. So the term at $l$ is at most $\bar C_\tau(l)\log l$, which is
\[
 (q_\tau+2j_\tau)(h+N_r)\frac{\log l}{l-1}+6(q_\tau+2j_\tau)\log l+kS\frac{\log l}l+k\,\delta_\tau(l)\log l
 +\bar\gamma_\tau(l)\log l+q_\tau h\Big(1+\log^+\frac hl\Big)\frac{\log l}l,
\]
where $\bar\gamma_\tau(l)=\min(k,t,q_\tau h/d_\tau(l))$ if $d_\tau(l)<x^-_\tau-\mathrm{top}_\tau$, and
$\bar\gamma_\tau(l)=0$ otherwise.

\emph{The sums over $R$.} Put $\Theta=\sum_{l\in R}\log l$ and $M=\sum_{l\in R}\log l/l$. By
Lemma~\ref{is-onto-step-exp-bound-total-numbers}, $l_J>8$, $h+N_r\le3b$, $\log(h+N_r)\le L+1.1$, $0<\log
h\le L+0.7$ and $h_0\le h<l_J^2$; and $\log4\le1.39$.
\begin{enumerate}
\item[(a)] Every $l\in R$ is at least $9$, so $l/(l-1)\le9/8$ and $\sum_{l\in R}\log l/(l-1)\le\frac98M$.
\item[(b)] The primes in $R$ are at most $h+N_r$, so Lemma~\ref{is-onto-step-exp-bound-total-primes}
  gives $\Theta\le(h+N_r)\log4\le1.39\cdot3b$ and $M\le\log(h+N_r)+\log4\le L+2.49$.
\item[(c)] $\sum_{l\in R}\delta_\tau(l)\log l\le(S+1)\log h$. Indeed, by
  Lemma~\ref{is-onto-step-exp-bound-total-exponents}(4), $0\le\mathrm{top}_\tau<x^-_\tau\le x^+_\tau\le
  h_0$, as $x^+_\tau$ is an integer at most $h$. If $\delta_\tau(l)=1$, there is $y_l$ in
  $[x^-_\tau,x^+_\tau]$ with $l\mid y_l-\mathrm{top}_\tau$; then $0<y_l-\mathrm{top}_\tau\le h_0$, so
  $\log l\le\log h_0\le\log h$. Two distinct primes $l,l'>l_J$ cannot both divide
  $y_l-\mathrm{top}_\tau<l_J^2$, so $l\mapsto y_l$ is injective, and there are at most
  $x^+_\tau-x^-_\tau+1\le S+1$ such $l$.
\item[(d)] $\sum_{l\in R}\bar\gamma_\tau(l)\log l\le q_\tau h(1+\log h)\log h$. Indeed, put
  $D=x^-_\tau-\mathrm{top}_\tau$, so that $1\le D\le h_0$ by (c), and $d_\tau(l)\le D$. If
  $d_\tau(l)<D$, then $l$ divides $D-d_\tau(l)>0$, so $l\le h_0$ and
  $\bar\gamma_\tau(l)\log l\le q_\tau h\log h_0/d_\tau(l)$. Two distinct primes $l,l'>l_J$ with
  $d_\tau(l)=d_\tau(l')<D$ would both divide $D-d_\tau(l)<l_J^2$; so $l\mapsto d_\tau(l)$ is injective on
  these $l$, with values in $[1,h_0]$. By Lemma~\ref{is-onto-step-exp-bound-total-primes}(3) the sum is
  at most $q_\tau h\log h_0\sum_{d=1}^{h_0}1/d\le q_\tau h(1+\log h_0)\log h_0$, and $h_0\le h$.
\item[(e)] $\sum_{l\in R}(1+\log^+(h/l))\log l/l\le(L+1.7)M$, as $\log(h/l)\le\log h\le L+0.7$.
\end{enumerate}

\emph{The total.} By Lemma~\ref{is-onto-step-exp-bound-total-exponents}(4), $q_\tau+j_\tau\le m$, so
$q_\tau+2j_\tau\le2m$; and $q_\tau h\le2mb$, as $h\le2b$ (Lemma~\ref{is-onto-step-exp-bound-numbers}). By
(a) to (e), the sum over $Q$ is at most
\[
 2m\cdot3b\cdot\tfrac98(L+2.49)+12m\cdot1.39\cdot3b+kS(L+2.49)+k(S+1)(L+0.7)
 +2mb(L+1.7)(L+0.7)+2mb(L+1.7)(L+2.49).
\]
The third and fourth terms add up to at most $2k(S+1)(L+2)$, as
$2(S+1)(L+2)-S(L+2.49)-(S+1)(L+0.7)=0.81S+L+3.3\geq0$. The other four are at most
$mb\big(12(L+2.5)+51+2(L+2)^2+2(L+2)(L+2.5)\big)=mb(4L^2+29L+99)$, and
$18L^2-(4L^2+29L+99)=14(L-4)^2+83(L-4)+9\geq0$, as $L\geq4$.
% Lean: largePrimeSumElementary_holds (SumDE): sum_min_cW_le_split (thickA for l > h; cloudA with the
% clause S + k < l of paramsLargePrime_holds, NoWindow and cW_eq_of_eq_zero for l <= h, with
% gatedCoefficientTerm the gamma-bar term); (a) is pred, (b) is chebyshevThetaElementary_holds and
% mertensLogElementary_holds over subset_primesLe_h_add_Nr, (c) is sum_residueIndicator_mul_log_le
% (residueIndicatorPrimeSum_holds, ISum, with eq_of_dvd_of_abs_lt_sq and floor_h_bounds), (d) is
% sum_gateTerm_mul_log_le (gatedPrimeSum_holds, GatedSum: sums_le_of_gate, sums_inj and
% sums_harmonic_real), (e) is sum_one_add_log_div_le; q_j_bounds, small_terms_le_of_log_four (the third
% and fourth terms) and the final nlinarith with 14 (L - 4)^2 + 83 (L - 4) + 9.
\end{proof}

\begin{lemma}\label{is-onto-step-exp-bound-total}
Assume the exact cut bound of Lemma~\ref{is-onto-step-exp-cut}, and that the number of primes at most
$y$ is less than $1.26\,y/\log y$ for every real $y>1$. Let $P$ be a parameter record that satisfies
the hypotheses and is large, and $W$ weights for $P$, indexed by a set of $f$ elements, that satisfy
the weight conditions and whose special set has at most $5$ elements
(Definition~\ref{is-onto-step-exp-bound-weights}). Let $J'$ be the recursive exponent of $P$ and $W$
(Definition~\ref{is-onto-step-exp-bound-recursion}), and $\mathrm{Tot}$ as in
Lemma~\ref{is-onto-step-exp-bound-tail}. Then for every finite set $Q$ of primes,
\[
 \sum_{l\in Q}J'(l)\log l\le\mathrm{Tot}(|T_W|).
\]
\end{lemma}
\leanref{Formalism/PieriBlockIdeals/PieriColumnMinors/RecursionInclusionExceptional.lean}{Formalism.CohomologicalDescent.E.RecSum7E}
\leanref{Formalism/PieriBlockIdeals/PieriColumnMinors/RecursionInclusionExceptional.lean}{Formalism.CohomologicalDescent.E.recursion_sum_exceptional}
\leanref{Formalism/PieriBlockIdeals/MertensReciprocalElementary.lean}{Formalism.CohomologicalDescent.LogSmallConstantBounds.totalBound}
\begin{proof}
Let $Q$ be a finite set of primes. Let $A$ be the set of $l\in Q$ with $l\le l_J$, $E$ the set of
$l\in Q$ with $l>l_J$ that are excluded by $P$ or lie in $\mathrm{Sp}$, and $V$ the rest of $Q$. By
Definition~\ref{is-onto-step-exp-bound-recursion}(5), $J'(l)=J_A(l)=\sum_\tau e_\tau(l)$ for $l\in A\cup
E$. Let $l\in V$. As $l>l_J$, Lemma~\ref{is-onto-step-exp-bound-total-numbers} gives $h+N_r<l^2$,
$N_r+1<l$ and $t+2\kappa(l)<l$, and Lemma~\ref{is-onto-step-exp-bound-total-weight-bound}, with $J=J'$,
gives $J'(l)\le\sum_\tau c_\tau(l)$, where $c_\tau(l)=\min(e_\tau(l),\bar C_\tau(l))$ if $\tau\in T_W$ and
$c_\tau(l)=e_\tau(l)$ otherwise (Definition~\ref{is-onto-step-exp-bound-total-real}).

\emph{The primes in $E$.} By Lemma~\ref{is-onto-step-exp-bound-total-cost}, $J'(l)\log l\le fN_r^2L$ for
$l\in E$. At most $S+\eta+t+2\kappa_{\max}+1$ primes of $E$ are excluded by $P$
(Lemma~\ref{is-onto-step-exp-bound-total-excluded}), and the others lie in $\mathrm{Sp}$, which has at
most $5$ elements. So
\[
 \sum_{l\in E}J'(l)\log l\le(S+\eta+t+2\kappa_{\max}+6)fN_r^2L.
\]

\emph{The primes in $A$ and $V$.} Exchanging the sums,
\[
 \sum_{l\in A\cup V}J'(l)\log l\le\sum_\tau s_\tau,\qquad
 s_\tau=\sum_{l\in A}e_\tau(l)\log l+\sum_{l\in V}c_\tau(l)\log l.
\]
Let $\tau\in T_W$. Lemma~\ref{is-onto-step-exp-bound-total-small}(2) bounds the first sum by $2mbL^2$.
The primes of $V$ are greater than $l_J$ and not excluded, so
Lemma~\ref{is-onto-step-exp-bound-total-large} bounds the second by $18mbL^2+2k(S+1)(L+2)$. Let
$\tau\notin T_W$. Then $s_\tau$ is the sum of $e_\tau(l)\log l$ over the disjoint union of $A$ and $V$,
which is at most $1.3N_r^2b'$ by Lemma~\ref{is-onto-step-exp-bound-total-thin}. At most $f$ indices lie
outside $T_W$. So
\[
 \sum_{l\in A\cup V}J'(l)\log l\le|T_W|\big(20mbL^2+2k(S+1)(L+2)\big)+1.3fN_r^2b'.
\]
Adding the two bounds, and as $20mbL^2\le30mbL^2$, $\sum_{l\in Q}J'(l)\log l\le\mathrm{Tot}(|T_W|)$.
% Lean: E.recursion_sum_exceptional (RecSum7E, which also takes ExIntervalLt, Lemma
% is-onto-step-exp-bound-interval), with D = W.sizeData₂ P (W.Jrec7 P): recursion_inputs_exceptional
% gives D.Inputs (boundA from seventh_recursion_eq_a_of_small_prime, boundW from
% sizeData_boundW_of_le_seventh_recursion = Lemma is-onto-step-exp-bound-total-weight-bound, thinA,
% thickA, cloudA and shapeT = Lemma is-onto-step-exp-bound-total-exponents, widths from (W6)), and the
% second part of places_real_holds gives D.NoWindow (Lemma is-onto-step-exp-bound-total-no-window). Then
% totalElementary_holds: hXC = smallPrimeSumElementary_holds (Lemma is-onto-step-exp-bound-total-small),
% hXD = largePrimeSumElementary_holds (Lemma is-onto-step-exp-bound-total-large), and
% sum_J_mul_log_le_totalBound, whose sets A, E and W are the book's A, E and V (costBound_holds,
% exceptionalCount_holds, thinSum_holds, I.card_Sp; I.boundW on W with clauses h8, h6 and h4 of
% paramsLargePrime_holds), with totalBound_mono for 2 + 18 <= 30. Size before
% this step was written out: about 4,000 lines of Lean in 270 theorems.
\end{proof}

\begin{lemma}\label{is-onto-step-exp-bound-sum}
Assume the exact cut bound of Lemma~\ref{is-onto-step-exp-cut}. Let $D$ be a definite-side datum, with
rank $f=2n$, whose host has the primes $q$, $d$, $p_{de}$, $p'_{de}$ and the integer $Q_S$
(Definition~\ref{hf-host}). Let $P$ and $W$ be a parameter record and weights for $P$ that satisfy the
numerical conditions (Definition~\ref{is-onto-step-exp-bound-weights}), such that the number $f$ of $P$
is $f$, $W$ is indexed by a set of $f$ elements, and the special set of $W$ is the set of prime divisors
of $Q_S$ together with $q$, $d$, $p_{de}$ and $p'_{de}$. Let $T\geq|T_W|$ be an integer. Then for every finite set
$Q$ of primes,
\[
 \sum_{l\in Q}J'(l)\log l\le30TmbL^2+2mb,
\]
where $J'$ is the recursive exponent of $P$ and $W$ (Definition~\ref{is-onto-step-exp-bound-recursion})
and $m$, $b$ and $L$ are those of $P$.
\end{lemma}
\leanref{Development/AutomorphicArithmetic/PrimeDiscriminants/ExcSet.lean}{UpperTheory.ExcSet.sum_exponent_log_le}
\leanref{Development/AbelianVarietyTheory/Proofs/DivisorProofs/SumExponentLogLeOfCut.lean}{GlobalTheory.sum_exponent_log_le_of_cut}
\leanref{Development/AutomorphicArithmetic/PrimeDiscriminants/ExcSet.lean}{UpperTheory.ExcSet.small_le_lJ}
\begin{proof}
Let $W^\circ$ be the weights $W$ with the special set $\{q,d,p_{de},p'_{de}\}$, which has at most $4$
elements. The weight conditions do not involve the special set, so $W^\circ$ satisfies them, and the
Weyl exponents and the coarse weight bound of $W^\circ$ are those of $W$.

The recursive exponents of $W$ and $W^\circ$ agree at every prime. Indeed, the special set of $W$ is
that of $W^\circ$ together with the prime divisors of $Q_S$, and every prime divisor $l$ of $Q_S$
satisfies $l\le l_J$: the host gives $Q_S\le n^2$ (Definition~\ref{hf-host-field-rows}(3)), the rows of
$P$ give $2n=f\le k$, and $k\le t$, so $l\le n^2\le kt$; largeness gives $4k^2t^2<b$, so $kt\le
b/(4kt)\le l_J$ by Lemma~\ref{is-onto-step-exp-bound-numbers}. So at every prime $l$ the condition that
$l\le l_J$, $l$ is excluded or $l$ lies in the special set holds for $W$ if and only if it holds for
$W^\circ$, and the two recursive exponents are given by the same case of
Definition~\ref{is-onto-step-exp-bound-recursion}(5).

The number of primes at most $y$ is less than $1.26\,y/\log y$ for every real $y>1$, a weaker form of
the bound of Rosser and Schoenfeld \cite[(3.6)]{RS62}. So Lemma~\ref{is-onto-step-exp-bound-total}, for $P$ and $W^\circ$, gives
$\sum_{l\in Q}J'(l)\log l\le\mathrm{Tot}(|T_W|)$. As $|T_W|\le f$,
Lemma~\ref{is-onto-step-exp-bound-tail} gives $\mathrm{Tot}(|T_W|)\le30|T_W|mbL^2+2mb$, and
$30|T_W|mbL^2\le30TmbL^2$ as $mbL^2\geq0$.
% Lean: sum_exponent_log_le: W' with Sp = {q, d, pde, pde'} (Finset.card_le_four), hJ from small_le_lJ
% (D.host.QS_le, Rows.f_le, paramsLJ_holds and Large.b_gt); then sum_exponent_log_le_of_cut:
% E.recursion_sum_exceptional (Lemma is-onto-step-exp-bound-total) with exIntervalLt_of_exCut and
% Top.prime_count_bound_rs = prime_count_bound (the Lean proves the bound with the constant 1.26 itself,
% through PI.PC.primeCountBound_of), tail_bound_holds with T = |T2| <= f, and |T2| <= T.
\end{proof}

\begin{lemma}\label{is-onto-step-exp-bound-window}
Let $P$ be a parameter record that satisfies the hypotheses and is large
(Definition~\ref{is-onto-step-exp-bound-weights}), and $l>l_J$ a prime. Let $\mathcal U$ be a finite set
of integers, $\beta$ an integer, $c$ and $e$ real numbers and $\sigma\in\{1,-1\}$, such that
$\sigma(\beta+(t-1)/2-c)=-e$, and every $u\in\mathcal U$ outside $[\beta,\beta+t-1]$ satisfies
$0\le\sigma(u-c)\le S$ or $2b/\ell^*+t/2+3\le|\sigma(u-c)+e|\le\eta$. Suppose that $b\le e\le b+\eta+1$
and $l$ is not excluded by $P$; or that there is $b_*$ with $b'\le b_*\le b'(1+x^{-2})$ and
$b_*+b'/x\le e\le b_*+b'/x+\eta$; or that $w''\le e\le w''+\eta$. Then $n_l(\beta,\mathcal U)=0$
(Definition~\ref{is-onto-step-exp-bound-recursion}(2)).
\end{lemma}
\leanref{Formalism/PieriBlockIdeals/MertensReciprocalElementary.lean}{Formalism.CohomologicalDescent.LogSmallConstantBounds.NoResidueInWindow}
\leanref{Formalism/PieriBlockIdeals/PieriColumnMinors/ParameterBounds.lean}{Formalism.CohomologicalDescent.LogSmallConstantBounds.noWindowResidue_holds}
\begin{proof}
Let $u\in\mathcal U$ lie outside $[\beta,\beta+t-1]$, and let $x$ be an integer in
$[\beta-\kappa(l),\beta+t-1+\kappa(l)]$. We show that $l$ does not divide $u-x$. Put
$d=\sigma(x-\beta-(t-1)/2)$, so that $|d|\le(t-1)/2+\kappa(l)$, and $y=\sigma(u-x)$, an integer. By the
condition on $\beta$,
\[
 y=\big(\sigma(u-c)+e\big)-d.
\]
By Lemma~\ref{is-onto-step-exp-bound-numbers}, $\kappa(l)\le\kappa_{\max}\le2b/\ell^*$,
$\eta+t+\kappa(l)<l$, $\kappa_{\max}+t<w''/2$, $w''\le b'$ and $3b'<l_J<l$.

Suppose first that $0\le\sigma(u-c)\le S$. If $b\le e\le b+\eta+1$, then
$b-(t-1)/2-\kappa_{\max}\le y\le b+S+\eta+1+(t-1)/2+\kappa_{\max}$; as $l$ is not excluded by $P$, $l$
does not divide $y$. If $b_*+b'/x\le e\le b_*+b'/x+\eta$ with $b'\le b_*\le b'(1+x^{-2})$, then
$y\geq b'-(t-1)/2-\kappa_{\max}>b'-b'/2>0$, as $\kappa_{\max}+t<w''/2\le b'/2$; and
$y\le S+b'(1+x^{-1}+x^{-2})+\eta+(t-1)/2+\kappa_{\max}\le2b'+\kappa_{\max}<3b'<l$, by largeness. If
$w''\le e\le w''+\eta$, then $y\geq w''-(t-1)/2-\kappa_{\max}>w''/2>0$, and
$y\le S+w''+\eta+(t-1)/2+\kappa_{\max}\le2b'+\kappa_{\max}<l$ in the same way, as $w''\le b'$. In the last
two cases $0<y<l$.

Suppose now that $2b/\ell^*+t/2+3\le|\sigma(u-c)+e|\le\eta$. Then
$|y|\geq2b/\ell^*+t/2+3-(t-1)/2-\kappa(l)\geq7/2$ and $|y|\le\eta+(t-1)/2+\kappa(l)<\eta+t+\kappa(l)<l$.
So $0<|y|<l$, and $l$ does not divide $y$.
% Lean: noWindowResidue_holds, through windowResidueCountEqZero_holds; the numbers are
% size_bounds_of_large, paramsLargePrime_holds, paramsTop_holds, paramsThin_holds, Large.w_gt and
% Large.top_le. In the first case the integer z of P.Excl is sigma (u - x).
\end{proof}

\begin{lemma}\label{is-onto-step-exp-bound-frobenius}
Let $P$ be a parameter record that satisfies the hypotheses and is large, $W$ weights for $P$ that
satisfy the weight conditions (Definition~\ref{is-onto-step-exp-bound-weights}), and $W'$ the weights
obtained from $W$ by replacing its Frobenius data by other Frobenius data. Then $W'$ satisfies the
weight conditions, and $J_{W'}(l)=J_W(l)$ for every prime $l$ in the range of $W$
(Definition~\ref{is-onto-step-exp-bound-recursion}).
\end{lemma}
\leanref{Development/AbelianVarietyTheory/Proofs/DivisorProofs/ExpoFrobIndep.lean}{GlobalTheory.expo_frob_indep}
\leanref{Formalism/PieriBlockIdeals/PieriColumnMinors/WeightBounds.lean}{Formalism.CohomologicalDescent.ShortWeightBounds.places_real_holds}
\leanref{Formalism/PieriBlockIdeals/PieriColumnMinors/WeightBounds.lean}{Formalism.CohomologicalDescent.ShortWeightBounds.weights_have_positive_shape}
\leanref{Formalism/PieriBlockIdeals/PieriColumnMinors/WeightBounds.lean}{Formalism.CohomologicalDescent.ShortWeightBounds.shape_reflect}
\begin{proof}
The weight conditions do not involve the Frobenius data, so $W'$ satisfies them. Let $l$ be a prime in
the range of $W$, so that $l>l_J$ and $l$ is not excluded by $P$. We show that the window counts
$n_\tau(l)$ and $n^{\mathrm F}_\tau(l)$ vanish, for $W$ and for $W'$, at every $\tau$. Then
$i_\tau(l)=\lfloor q_\tau h\rfloor$ for both, and the coarse weight bound depends on the Frobenius data
only through the depths $i_\tau(l)$, so $J_{W'}(l)=J_W(l)$.

Let $V$ be $W$ or $W'$, with Frobenius data $(\phi_l,s_l)$, and let $\tau$ be given. The set
$\mathcal U_\tau$, $\beta=\mathrm{bot}_\tau$, $c_\tau$, $e_\tau$ and $\sigma=1$ satisfy the conditions
of Lemma~\ref{is-onto-step-exp-bound-window} on $\mathcal U$. Indeed, (W2) gives
$\mathrm{bot}_\tau+(t-1)/2-c_\tau=-e_\tau$. By (W1) the $U_\tau(a)$, $a<N_r$, decrease strictly, and by
(W2) $U_\tau(a)=\mathrm{bot}_\tau+(p_\tau+t-1-a)$ for $p_\tau\le a<p_\tau+t$; so if $U_\tau(a)$ lies outside
$[\mathrm{bot}_\tau,\mathrm{bot}_\tau+t-1]$, then $a<p_\tau$ or $a\geq p_\tau+t$. If $a<k-1$, (W3) gives
$0\le U_\tau(a)-c_\tau\le S$; otherwise (W4) gives $2b/\ell^*+t/2+3\le|U_\tau(a)-c_\tau+e_\tau|\le\eta$.
By (W5) and (W6), $e_\tau$ satisfies one of the three conditions on $e$ of that lemma, the first one
together with the fact that $l$ is not excluded. So $n_\tau(l)=0$.

Let $\tau'=\phi_l(\tau)$. If $s_l(\tau)=0$, then $n^{\mathrm F}_\tau(l)=n_{\tau'}(l)=0$. Let
$s_l(\tau)=1$. By (W2), $\mathrm{top}_{\tau'}=\mathrm{bot}_{\tau'}+t-1$. The set
$\{h_0-u:u\in\mathcal U_{\tau'}\}$, $\beta'=h_0-\mathrm{top}_{\tau'}$, $c'=h_0-c_{\tau'}$, $e_{\tau'}$ and
$\sigma=-1$ satisfy the conditions of Lemma~\ref{is-onto-step-exp-bound-window}: first,
$-(\beta'+(t-1)/2-c')=\mathrm{bot}_{\tau'}+(t-1)/2-c_{\tau'}=-e_{\tau'}$; next, $h_0-u$ lies in
$[\beta',\beta'+t-1]$ if and only if $u$ lies in $[\mathrm{bot}_{\tau'},\mathrm{bot}_{\tau'}+t-1]$;
and $-\big((h_0-u)-c'\big)=u-c_{\tau'}$. The condition on $e_{\tau'}$ holds as for $\tau$. So
$n^{\mathrm F}_\tau(l)=0$.
% Lean: expo_frob_indep: the fields of Wts.Ok carried over, and places_real_holds for W and for the
% reframed weights, whose second part is noWindowOf_holds (noWindowResidue_holds) for the places
% placesOfParameters (weights_have_positive_shape; shape_reflect with top_eq_bot_add_t_sub_one for the
% conjugated case; thick_or_top_or_thin); then iW agree and coarseWeightBound unfolds to the same sum.
\end{proof}

\subsection{Linkage and the depth bound}\label{subsection-is-linkage-depth}

\begin{definition}\label{is-onto-step-exp-bound-link-depth-series}
Let $l$ be a prime and $O$ a commutative ring that is a $\mathbf Z_l$-algebra. Power series are in one
variable $X$, and a series $f$ is evaluated only at series $g$ with constant term $0$, giving $f(g)$.
\begin{enumerate}
\item For $c\in\mathbf Z_l$ put $\beta_c=(1+X)^c-1=\sum_{m\geq1}\binom cmX^m$ in $O[[X]]$, and
  $\mathrm{act}_c(f)=f(\beta_c)$ and $\varphi(f)=f\big((1+X)^l-1\big)$ for $f\in O[[X]]$.
\item Let $\mu_{l-1}$ be the set of roots of $x^{l-1}=1$ in $\mathbf Z_l$. Put
  $\pi_0=\prod_{d\in\mu_{l-1}}\beta_d$, $q_0=\pi_0-l$ and $q=\sum_{m\geq0}\binom l{m+1}X^m$, so that
  $Xq=(1+X)^l-1$; and let $\iota(g)=g(\pi_0)$ for $g\in O[[X]]$.
\item For $u\in\mathbf Z_l$, a series $u^\flat\in O[[X]]$ is a \emph{unit series} for $u$ if
  $\mathrm{act}_u(\pi_0)=\pi_0\big(1+\iota(u^\flat)q_0\big)$.
\item For a commutative ring $R$ and $v\in R[[X]]$ put $\gamma_v(y)=y\big(X(1+vX)\big)$ and
  $\mathrm{Fr}_l(y)=y(X^l)$ for $y\in R[[X]]$. Both are ring endomorphisms of $R[[X]]$, and they act on
  matrices entry by entry.
\item For $x\in R$, an integer $r\geq0$ and a matrix $P$ over $R$, $x$ \emph{divides the $r$-minors} of
  $P$ if it divides the determinant of every $r\times r$ submatrix of $P$ whose rows and columns are
  taken in increasing order.
\end{enumerate}
\end{definition}
\leanref{Formalism/PieriBlockIdeals/CapelliCentralEnveloping.lean}{Formalism.CohomologicalDescent.HeightDescentMinors.act}
\leanref{Formalism/PieriBlockIdeals/CapelliCentralEnveloping.lean}{Formalism.CohomologicalDescent.HeightDescentMinors.phi}
\leanref{Formalism/PieriBlockIdeals/CapelliCentralEnveloping.lean}{Formalism.CohomologicalDescent.HeightDescentMinors.mu}
\leanref{Formalism/PieriBlockIdeals/CapelliCentralEnveloping.lean}{Formalism.CohomologicalDescent.HeightDescentMinors.pi0}
\leanref{Formalism/PieriBlockIdeals/CapelliCentralEnveloping.lean}{Formalism.CohomologicalDescent.HeightDescentMinors.q0}
\leanref{Formalism/PieriBlockIdeals/CapelliCentralEnveloping.lean}{Formalism.CohomologicalDescent.HeightDescentMinors.q}
\leanref{Formalism/PieriBlockIdeals/CapelliCentralEnveloping.lean}{Formalism.CohomologicalDescent.HeightDescentMinors.inc}
\leanref{Formalism/PieriBlockIdeals/LemmaWFromLattice.lean}{Formalism.CohomologicalDescent.FrobeniusPeriodChange.IsUn}
\leanref{Formalism/PieriBlockIdeals/LemmaWFromLattice.lean}{Formalism.CohomologicalDescent.FrobeniusPeriodChange.gamSer}
\leanref{Formalism/PieriBlockIdeals/LemmaWFromLattice.lean}{Formalism.CohomologicalDescent.FrobeniusPeriodChange.gam0}
\leanref{Formalism/PieriBlockIdeals/LemmaWFromLattice.lean}{Formalism.CohomologicalDescent.FrobeniusPeriodChange.frob0}
\leanref{Formalism/PieriBlockIdeals/CapelliCentralEnveloping.lean}{Formalism.CohomologicalDescent.HeightDescentMinors.MinorsDvd}
% Lean: br l O c = binomialSeries O c - 1, act and phi by PowerSeries.subst, mu l = nthRootsFinset (l - 1) 1,
% inc g = g.subst (pi0 l O). The book's gamma_v is gam0 0 v, the substitution of gamSer 0 v = X (1 + v X);
% gam0 with a nonzero lambda does not enter the period change.

\begin{definition}\label{is-onto-step-exp-bound-link-depth-data}
Let $l$ be a prime, $O$ a commutative ring that is a $\mathbf Z_l$-algebra, $E$ a finite set,
$n,t,h\geq0$ integers, $u\in\mathbf Z_l^\times$ and $c\in\mathbf Z_l$; the notation is that of
Definition~\ref{is-onto-step-exp-bound-link-depth-series}. For an integer $y$ write $y^+=\max(0,y)$.
\begin{enumerate}
\item A \emph{lattice datum} $N$ of rank $n$ over $O$, indexed by $E$, consists of matrices
  $P_i,D_i\in M_n(O[[X]])$ and $\Theta_i\in M_n(O)$ for $i\in E$. Let $U'=(U'_i)_{i\in E}$ be integer
  sequences of length $n$. Then $N$ has \emph{height $h$ with weights $U'$} if for every $i$:
  \begin{enumerate}
  \item[(a)] $P_iM=(X-l)^h$ for some $M\in M_n(O[[X]])$;
  \item[(b)] for all integers $r,S\geq0$: if $S\le\sum_{j\in J}U'_i(j)$ for every set $J$ of $r$
    indices, then $(X-l)^S$ divides the $r$-minors of $P_i$;
  \item[(c)] $\det P_i=v(X-l)^{(\sum_jU'_i(j))^+}$ for a unit $v$ of $O[[X]]$;
  \item[(d)] $X-l$ divides $D_i(a,b)-\Theta_i(a,b)$ in $O[[X]]$ for all $a$ and $b$;
  \item[(e)] there are $r_j\in\mathbf Z_l$ with $u^{-U'_i(j)}=1+lr_j$ for every $j$ and
    $\det(Y-\Theta_i)=\prod_j(Y-r_j)$ in $O[Y]$.
  \end{enumerate}
\item Let $L$ be a field, $\mathrm{em}\colon O\to L$ a ring homomorphism and $x\in O$. A ring
  homomorphism $\jmath\colon O[[X]]\to L[[X]]$ is a \emph{local exponential} at $x$ for $u$ if it maps
  every constant $a$ to $\mathrm{em}(a)$, $\jmath(X-x)$ is $X$ times a unit, and
  $\jmath(\mathrm{act}_uf)(X)=\jmath(f)\big(\mathrm{em}(u)X\big)$ for every $f\in O[[X]]$.
\item Given such $L$ and em, and ring homomorphisms $\jmath_b\colon O[[X]]\to L[[X]]$ for $b\geq0$, the
  lattice datum $N$ has \emph{local height $h$ with weights $U'$} if for every $i$: $\iota(P_i)Q=q^h$
  for some $Q\in M_n(O[[X]])$; (d) holds; and for every $0<b<l$ there are $G,H\in GL_n(L[[X]])$ with
  \[
   \jmath_b\big(\iota(P_i)\big)=G\Delta_iH\quad\text{and}\quad
   \jmath_b\big(\iota(1+XD_i)\big)\cdot(G\Delta_i)\big(\mathrm{em}(u)X\big)=G\Delta_i,
  \]
  where $\Delta_i$ is the diagonal matrix with entries $X^{U'_i(a)^+}$, $a<n$.
\item Let $F$ be a field. \emph{Witt lattice data} $\mathcal W$ for $(E,t,n,O,F)$ consist of a ring
  homomorphism $\mathrm{red}\colon O\to F$; a series $u^\flat\in O[[X]]$; a lattice datum $N$ of rank
  $n$ and lattice data $N^\psi_s$ of rank $1$ for $s<t$, all indexed by $E$; matrices
  $g_i\in GL_n(F[[X]])$ for $i\in E$; maps $b\colon\{0,\dots,n-1\}\to\mathbf N$ and
  $\mathrm{pos}\colon\{0,\dots,t-1\}\to\{0,\dots,n-1\}$; and units $w_i(s)$ of $F[[X]]$ for $i\in E$ and
  $s<t$. A bar denotes reduction of coefficients by red, and $\bar u=\overline{u^\flat}$. Let $P_i,D_i$
  be the matrices of $N$, and $P^\psi_{s,i},D^\psi_{s,i}$ the single entries of those of $N^\psi_s$. For
  a permutation $\sigma$ of $E$ put
  \[
   \Gamma_i=g_i^{-1}(1+X\bar D_i)\gamma_{\bar u}(g_i)\qquad\text{and}\qquad
   \Phi_i=g_i^{-1}\bar P_i\,\mathrm{Fr}_l\big(g_{\sigma(i)}\big).
  \]
  A matrix $M$ indexed by $\{0,\dots,n-1\}$ is \emph{block upper triangular} for $b$ if $M(a,a')=0$
  whenever $b(a')<b(a)$.
\item Let $U'_i$ be integer sequences of length $n$, and $\mathrm{wt}_i(s)$ and $\nu_i(s)$ integers, for
  $i\in E$ and $s<t$. Then $\mathcal W$ is a \emph{frame} for
  $(\sigma,l,h,u,c,U',\mathrm{wt},\nu)$ if:
  \begin{enumerate}
  \item[(a)] $u=1+lc$, $c$ is a unit of $\mathbf Z_l$, and $u^\flat$ is a unit series for $u$;
  \item[(b)] $0\le U'_i(j)\le h$ and $0\le\mathrm{wt}_i(s)\le h$ for all $i$, $j$ and $s$;
  \item[(c)] pos is injective, and $b(j)=b(\mathrm{pos}(s))$ only for $j=\mathrm{pos}(s)$;
  \item[(d)] every $\Phi_i$ and every $\Gamma_i$ is block upper triangular for $b$;
  \item[(e)] in the Laurent series $F((X))$, for all $i$ and $s$, with $p=\mathrm{pos}(s)$,
    \[
     \mathrm{Fr}_l\big(w_{\sigma(i)}(s)\big)X^{l\nu_{\sigma(i)}(s)}\bar P^\psi_{s,i}
     =\Phi_i(p,p)\,w_i(s)X^{\nu_i(s)};
    \]
  \item[(f)] in $F((X))$, for all $i$ and $s$, with $p=\mathrm{pos}(s)$,
    \[
     \gamma_{\bar u}\big(w_i(s)\big)\big(X(1+\bar uX)\big)^{\nu_i(s)}\big(1+X\bar D^\psi_{s,i}\big)
     =\Gamma_i(p,p)\,w_i(s)X^{\nu_i(s)}.
    \]
  \end{enumerate}
  It has the \emph{weight inputs} for $(\sigma,l,h,u,c,U',\mathrm{wt},\nu)$ if moreover $N$ has height
  $h$ with weights $U'$ and every $N^\psi_s$ has height $h$ with weights $(\mathrm{wt}_i(s))_{i\in E}$.
\item Let $I$ be an ideal of $O$, $\zeta\in O$, $L$ a field, $\mathrm{em}\colon O\to L$ a ring
  homomorphism and $\jmath_b\colon O[[X]]\to L[[X]]$ ring homomorphisms for $b\geq0$. Then $\mathcal W$
  satisfies the \emph{inputs of the Frobenius period change} for
  $(\sigma,l,h,u,c,U',\mathrm{wt},\nu,I,\zeta,\mathrm{em},\jmath)$ if it is a frame for
  $(\sigma,l,h,u,c,U',\mathrm{wt},\nu)$; $\zeta$ is a primitive $l$-th root of unity, $\zeta-1\in I$, em
  is injective, and $\jmath_b$ is a local exponential at $\zeta^b-1$ for $u$ for every $0<b<l$; $N$ has
  local height $h$ with weights $U'$; and every $N^\psi_s$ has local height $h$ with weights
  $(\mathrm{wt}_i(s))_{i\in E}$.
\end{enumerate}
\end{definition}
\leanref{Formalism/PieriBlockIdeals/LemmaWFromLattice.lean}{Formalism.CohomologicalDescent.FrobeniusPeriodChange.Lat}
\leanref{Formalism/PieriBlockIdeals/LemmaWFromLattice.lean}{Formalism.CohomologicalDescent.FrobeniusPeriodChange.LatW}
\leanref{Formalism/PieriBlockIdeals/LemmaWFromLattice.lean}{Formalism.CohomologicalDescent.FrobeniusPeriodChange.IsLocExp}
\leanref{Formalism/PieriBlockIdeals/LemmaWFromLattice.lean}{Formalism.CohomologicalDescent.FrobeniusPeriodChange.LatN}
\leanref{Formalism/PieriBlockIdeals/LemmaWFromLattice.lean}{Formalism.CohomologicalDescent.FrobeniusPeriodChange.WLat}
\leanref{Formalism/PieriBlockIdeals/LemmaWFromLattice.lean}{Formalism.CohomologicalDescent.FrobeniusPeriodChange.WLat.Gam}
\leanref{Formalism/PieriBlockIdeals/LemmaWFromLattice.lean}{Formalism.CohomologicalDescent.FrobeniusPeriodChange.WLat.Phi}
\leanref{Formalism/PieriBlockIdeals/LemmaWFromLattice.lean}{Formalism.CohomologicalDescent.FrobeniusPeriodChange.WFrame}
\leanref{Formalism/PieriBlockIdeals/LemmaWFromLattice.lean}{Formalism.CohomologicalDescent.FrobeniusPeriodChange.WInputs}
\leanref{Formalism/PieriBlockIdeals/LemmaWFromLattice.lean}{Formalism.CohomologicalDescent.FrobeniusPeriodChange.NInputs}
% Lean: (1) is Lat with LatW (the exponent (sum U')^+ is Int.toNat); (2) is IsLocExp, with rescale for
% f(em(u) X) and Associated for "X times a unit"; (3) is LatN; (4) is WLat, with Gam = chgGam unbar (g i)
% (red D) and Phi = chgPhi l (g i) (g (sigma i)) (red P), and Matrix.BlockTriangular for "block upper
% triangular"; (5) is WFrame (fields u = 1 + l c, IsUnit c, IsUn; U' and wt in [0, h]; pos; the two
% BlockTriangular; the two identities in LaurentSeries F) and WInputs; (6) is NInputs. The book's nu is
% the Lean's kap.

\begin{definition}\label{is-onto-step-exp-bound-link-depth-reflected}
Let $E$ be a finite set with a permutation $\sigma$, $n,t,h\geq0$ integers and $F$ a field.
\begin{enumerate}
\item \emph{Labels} of length $n$ on $E$ consist of a map $i\mapsto\bar i$ of $E$ to itself, a subset
  $\Phi\subset E$, and for every $i\in E$ integer sequences $U_i,U'_i$ of length $n$ and an integer
  $\mathrm{bot}_i$. They are \emph{reflected} about $h$ if for all $i\in E$ and $a<n$:
  $\bar{\bar i}=i$, $\sigma(\bar i)=\overline{\sigma(i)}$, $i\in\Phi$ or $\bar i\in\Phi$,
  $U_{\bar i}(a)=h-U_i(n-1-a)$, $U'_{\bar i}(a)=h-U'_i(n-1-a)$ and
  $\mathrm{bot}_{\bar i}=h-(\mathrm{bot}_i+t-1)$.
\item For integers $z_i$ ($i\in E$), the labels are \emph{linked} over $F$ with the constants $z_i$ if
  for every $i\in\Phi$: $U_i$ and $U'_i$ decrease strictly, $U_i$ takes the values $\mathrm{bot}_i+s$
  for $0\le s<t$, and $\prod_{a<n}\big(Y+U_i(a)+z_i\big)=\prod_{a<n}\big(Y+U'_i(a)+z_i\big)$ in $F[Y]$.
\item Integer sequences $J_i$ of length $n$ and integers $\beta_i$ ($i\in E$) are \emph{carried} by the
  labels through a map $\epsilon\colon E\to E$ if $\epsilon$ is the identity or $i\mapsto\bar i$, and
  $J_{\epsilon(i)}=U'_i$ and $\beta_{\epsilon(i)}=\mathrm{bot}_i$ for every $i$.
\item For integers $l\geq1$, $\kappa\geq0$ and $\beta$ and a finite set $\mathcal U$ of integers, let
  $n_{l,\kappa}(\beta,\mathcal U)$ be the number of $u\in\mathcal U$ outside $[\beta,\beta+t-1]$ that
  are congruent modulo $l$ to an integer in $[\beta-\kappa,\beta+t-1+\kappa]$. Thus the number
  $n_l(\beta,\mathcal U)$ of Definition~\ref{is-onto-step-exp-bound-recursion}(2) is
  $n_{l,\kappa(l)}(\beta,\mathcal U)$.
\end{enumerate}
\end{definition}
\leanref{Formalism/PieriBlockIdeals/MultiplicityZeroThickEmbedding.lean}{Formalism.CohomologicalDescent.WorkedExample.Labels}
\leanref{Formalism/PieriBlockIdeals/MultiplicityZeroThickEmbedding.lean}{Formalism.CohomologicalDescent.WorkedExample.Labels.Refl}
\leanref{Formalism/PieriBlockIdeals/MultiplicityZeroThickEmbedding.lean}{Formalism.CohomologicalDescent.WorkedExample.Labels.Linked}
\leanref{Formalism/PieriBlockIdeals/MultiplicityZeroThickEmbedding.lean}{Formalism.CohomologicalDescent.WorkedExample.Labels.Carried}
\leanref{Formalism/PieriBlockIdeals/Statement.lean}{Formalism.CohomologicalDescent.mwCount}
% Lean: Labels (cc, Phi, U, U', bot) with Refl sigma h t, Linked F t cst (the z_i are cst), Carried ep J
% botJ, and mwCount l b t kappa U for n_{l,kappa}(b, U); Labels.mw is mwCount at the label's bot and the
% image of its U.

\begin{definition}\label{is-onto-step-exp-bound-hida}
Let $N\geq0$, let $\mathcal T$ be a finite set and $\lambda=(\lambda_\tau)_{\tau\in\mathcal T}$ sequences
of natural numbers whose first $N$ terms are nonincreasing, and put
$\nu_\tau=(\lambda_{\tau,0},\dots,\lambda_{\tau,N-1})$. For a commutative ring $O$ let
$\nabla_O(\lambda)=\bigotimes_{\tau}\nabla_O(\nu_\tau)$, the tensor product over $O$ of the dual Weyl
modules of Definition~\ref{pd-weight-lattices}(1) in the entries of an $N\times N$ matrix, and for a
subset $s\subset\mathcal T$ let $\Delta^s_O(\lambda)$ be the image in $\nabla_O(\lambda)$ of the tensor
product of the modules $\nabla_O(\nu_\tau)$ for $\tau\in s$ and $\Delta_O(\nu_\tau)$ for $\tau\notin s$;
put $\Delta_O(\lambda)=\Delta^{\emptyset}_O(\lambda)$.
\begin{enumerate}
\item A \emph{Hida setting} of rank $N$ and weights $\lambda$ consists of a Noetherian local domain $O$
  of characteristic $0$; a ring $A$ with an $O$-algebra structure; a group $G$ with subgroups $\Gamma$
  and $K$ and a homomorphism of monoids $j\colon K\to A$; an $A$-module structure on
  $\nabla_O(\lambda)$ that extends its $O$-module structure; a commutative $O$-algebra $T$ that acts
  $O$-linearly on $F(\nabla_O(\lambda))$; and a submonoid $S\subset T$. Here, for an $A$-module $V$,
  $F(V)$ is the $O$-module of \emph{forms} with values in $V$: the maps $f\colon G\to V$ with
  $f(\gamma g\kappa)=j(\kappa^{-1})f(g)$ for $\gamma\in\Gamma$, $g\in G$ and $\kappa\in K$.
\item For integers $l,J\geq0$, the setting has the \emph{inclusion with exponent $J$} at $l$ if for
  every $y$ in the localization $S^{-1}F(\nabla_O(\lambda))$ there are $s\in S$ and a form $f$ with
  values in $\Delta_O(\lambda)$ with $l^{J}y=f/s$. Its \emph{least exponent} at $l$ is the least such
  $J$, and $0$ if there is none.
\item Let $R$, $j$, $t$ and $w$ be as in Definition~\ref{is-onto-step-exp-bound-lift-exponent}. A
  family of Hida settings of rank $N$ and weights $\lambda$ \emph{covers} the block data at $(j,t)$ at
  $w$ if $p_w^{J}C_{j,t}\subset(D_{j,t})_{(w)}$ for every $J\geq0$ such that every member of the family
  has the inclusion with exponent $J$ at $p_w$.
\item Let $P$ be a parameter record and $W$ weights for $P$ with sequences $\lambda$
  (Definition~\ref{is-onto-step-exp-bound-weights}), $l$ a prime, and $X$ a Hida setting of rank $N_r$
  and weights $\lambda$. \emph{Constituent data} for $X$ at $l$ consist of $\varpi\in O$ and
  $n_0\geq0$ such that $O/lO$ is finite and $l$ divides $\varpi^{n_0}$; elements $e_\tau(\omega)\in A$
  for $\tau\in\mathcal T$ and $\omega\in\mathbf N^{N_r}$; $A$-modules $L_\mu$, for $\mu$ in a set
  $\Lambda$, with compatible $O$-module structures and actions of $T$ on $F(L_\mu)$; and weights
  $\omega_{\mu,\tau}\in\mathbf N^{N_r}$; such that:
  \begin{enumerate}
  \item for every $s\subset\mathcal T$, $A$ preserves $\Delta^s_O(\lambda)$; and $e_\tau(\omega)$ maps a
    pure tensor $\bigotimes_{\tau'}x_{\tau'}$ to the pure tensor in which $x_\tau$ is replaced by its
    component of row degree $\omega$, of degree $\omega_a$ in the entries of the $a$-th row for every
    $a$;
  \item $T$ preserves the forms with values in every $A$-submodule $U$ of $\nabla_O(\lambda)$, and
    commutes with the maps of forms induced by $A$-linear maps from such $U$ to the $L_\mu$;
  \item every simple $A$-module that is a subquotient of $\nabla_O(\lambda)/\varpi\nabla_O(\lambda)$
    is isomorphic to some $L_\mu$, and $e_\tau(\omega_{\mu,\tau})$ acts on $L_\mu$ by a nonzero map for
    all $\mu$ and $\tau$.
  \end{enumerate}
  An index $\mu$ is \emph{relevant} if $L_\mu$ is isomorphic to a subquotient of
  $\nabla_O(\lambda)/\varpi\nabla_O(\lambda)$ and $S^{-1}F(L_\mu)\ne0$. The setting has the
  \emph{weight inputs} for $W$ at $l$ if it has constituent data at $l$ such that, for every relevant
  $\mu$ and every $\tau\in T_W$,
  \[
   \sum_{a<k-1+j_\tau}\big(\lambda_{\tau,a}-\omega_{\mu,\tau,a}\big)\le i_\tau(l)
  \]
  (Definition~\ref{is-onto-step-exp-bound-recursion}(2)).
\item Let $n_\tau=k-1+j_\tau+t+q_\tau$. A weight $\omega\in\mathbf N^{N_r}$ is \emph{linked} to $W$ at
  $l$ at $\tau$ if there are:
  \begin{enumerate}
  \item[(i)] a field $F$ of characteristic $l$, a finite set $E$ with a permutation $\sigma$, and labels
    $(\bar{\ },\Phi,U,U',\mathrm{bot})$ of length $n_\tau$ on $E$ that are reflected about $h_0$ and
    linked over $F$ with some constants $z_i$; and sequences $J_i$ and integers $\beta_i$ carried by the
    labels through some map $\epsilon$ (Definition~\ref{is-onto-step-exp-bound-link-depth-reflected});
  \item[(ii)] a domain $O'$ that is a $\mathbf Z_l$-algebra and complete for the $I$-adic topology of an
    ideal $I$, a field $L'$, $u\in\mathbf Z_l^\times$, $c\in\mathbf Z_l$, $\zeta\in O'$, ring
    homomorphisms $\mathrm{em}\colon O'\to L'$ and $\jmath_b\colon O'[[X]]\to L'[[X]]$ for $b\geq0$,
    integers $\nu_i(s)$, and Witt lattice data for $(E,t,n_\tau,O',F)$ that satisfy the inputs of the
    Frobenius period change for $(\sigma,l,h_0,u,c,J,\mathrm{wt},\nu,I,\zeta,\mathrm{em},\jmath)$, where
    $\mathrm{wt}_i(s)=\beta_i+s$ (Definition~\ref{is-onto-step-exp-bound-link-depth-data});
  \item[(iii)] a map $\iota\colon\mathcal T\to E$ with $U_{\iota(\tau')}(a)=U_{\tau'}(a)$ for $a<n_\tau$
    and $\mathrm{bot}_{\iota(\tau')}=\mathrm{bot}_{\tau'}$ for every $\tau'$,
    $\sigma(\iota(\tau))=\iota(\phi_l(\tau))$ if $s_l(\tau)=0$ and
    $\sigma(\iota(\tau))=\overline{\iota(\phi_l(\tau))}$ if $s_l(\tau)=1$,
    $U'_{\iota(\tau)}(a)=\omega_a+N_r-1-a$ for $a<\min(n_\tau,N_r)$, and
    $\sum_{a<n_\tau}U_{\iota(\tau)}(a)=\sum_{a<n_\tau}U'_{\iota(\tau)}(a)$.
  \end{enumerate}
  The setting $X$ is \emph{linked} to $W$ at $l$ if
  it has constituent data at $l$ such that $\omega_{\mu,\tau}$ is linked to $W$ at $l$ at $\tau$ for every
  relevant $\mu$ and every $\tau\in T_W$.
\end{enumerate}
\end{definition}
\leanref{Formalism/PieriBlockIdeals/LatticeInclusionRealSettingAlt.lean}{Formalism.CohomologicalDescent.ShortWeightBounds.HidaSetting}
\leanref{Formalism/PieriBlockIdeals/LatticeInclusionRealSettingAlt.lean}{Formalism.CohomologicalDescent.ShortWeightBounds.HidaSetting.Incl}
\leanref{Formalism/PieriBlockIdeals/BlockIdealBound5Derived.lean}{Formalism.CohomologicalDescent.DepthNonnegative.TopIncl}
\leanref{Development/AbelianVarietyTheory/Definitions/Divisors/SumExpSt.lean}{GlobalTheory.ExpTie.Covers}
\leanref{Formalism/PieriBlockIdeals/LatticeInclusionRealSettingAlt.lean}{Formalism.CohomologicalDescent.ShortWeightBounds.HidaSetting.WInputs}
\leanref{Formalism/PieriBlockIdeals/LatticeInclusionRealSettingAlt.lean}{Formalism.CohomologicalDescent.ShortWeightBounds.WInputsLink}
\leanref{Formalism/PieriBlockIdeals/LatticeInclusionRealSettingAlt.lean}{Formalism.CohomologicalDescent.ShortWeightBounds.LatticeLink}
% Part (5) is LatticeLink: (i) is WorkedExample.Labels D with D.Refl sigma h0 t, D.Linked F t cst and
% D.Carried ep J botJ (Definition is-onto-step-exp-bound-link-depth-reflected); (ii) is a
% FrobeniusPeriodChange.WLat over O' (CommRing, IsDomain, Algebra Z_l, IsAdicComplete I) with NInputs
% (Definition is-onto-step-exp-bound-link-depth-data); (iii) is the last five clauses (hU, hbot, hsigma,
% hU' over the a : Fin n_tau and a' : Fin Nr with equal values, hsum). The labels have length
% kc P + W.j tau + P.t + W.q tau, which is N_r under (W1).

\begin{lemma}\label{is-onto-step-exp-bound-link-depth-content}
Let $l$ be a prime, $K$ a field of characteristic $l$, $n\geq0$, $U,U'\in\mathbf Z^n$ and
$z\in\mathbf Z$, with $\prod_{a<n}(Y+U(a)+z)=\prod_{a<n}(Y+U'(a)+z)$ in $K[Y]$. Then the multisets of the
residues modulo $l$ of $U(0),\dots,U(n-1)$ and of $U'(0),\dots,U'(n-1)$ agree.
\end{lemma}
\leanref{Formalism/PieriBlockIdeals/Statement.lean}{Formalism.CohomologicalDescent.Linkage}
\leanref{Formalism/PieriBlockIdeals/PieriColumnMinors/ContentLinkageHolds.lean}{Formalism.CohomologicalDescent.contentLinkage_holds}
\begin{proof}
The roots in $K$ of the two sides, with multiplicities, are the images under the ring homomorphism
$\mathbf Z/l\to K$ of the multisets of the residues of the $-(U(a)+z)$ and of the $-(U'(a)+z)$. That
homomorphism is injective, so these two multisets of residues agree. Applying $y\mapsto-y-z$ gives the
claim.
% Lean: contentLinkage_holds, through Polynomial.roots of both products (comb_roots_prod), ZMod.castHom
% and Multiset.map_injective.
\end{proof}

\begin{lemma}\label{is-onto-step-exp-bound-link-depth-shift}
Let $l\geq1$, $b\in\mathbf Z$ and $t,\kappa\geq0$ be integers with $t+2\kappa<l$. Let $\mathcal U$ be a
finite set of integers that contains $b+s$ for $0\le s<t$, and $\nu_0,\dots,\nu_{t-1}$ integers with
$|\nu_s|\le\kappa$, such that the multiset of the residues modulo $l$ of the $b+s-\nu_s$, $s<t$, is
contained in the multiset of the residues of the elements of $\mathcal U$. Then
\[
 \Big|\sum_{s<t}\nu_s\Big|\le n_{l,\kappa}(b,\mathcal U)\,(t+2\kappa)
\]
(Definition~\ref{is-onto-step-exp-bound-link-depth-reflected}(4), with this $t$).
\end{lemma}
\leanref{Formalism/PieriBlockIdeals/Statement.lean}{Formalism.CohomologicalDescent.KapSum}
\leanref{Formalism/PieriBlockIdeals/PieriColumnMinors/WeightExchangeBoundHolds.lean}{Formalism.CohomologicalDescent.shiftSumBound_holds}
\leanref{Formalism/PieriBlockIdeals/PieriColumnMinors/WeightExchangeBoundHolds.lean}{Formalism.CohomologicalDescent.we_msum}
\begin{proof}
Put $\mathrm{lo}=b-\kappa$ and $\mathrm{hi}=b+t-1+\kappa$, so that $\mathrm{hi}<\mathrm{lo}+l$. For an
integer $u$ let $g(u)$ be the integer in $[\mathrm{lo},\mathrm{lo}+l-1]$ congruent to $u$ modulo $l$.
It depends only on the residue of $u$, and $g(v)=v$ for $v\in[\mathrm{lo},\mathrm{hi}]$. Let $Y$ be the
multiset of the $y_s=b+s-\nu_s$, $s<t$, which lie in $[\mathrm{lo},\mathrm{hi}]$. Applying to the
hypothesis the map that sends a residue to its representative in $[\mathrm{lo},\mathrm{lo}+l-1]$, $Y$ is
contained in the multiset of the $g(u)$, $u\in\mathcal U$, and so in its part $A$ of elements at most
$\mathrm{hi}$.

An element $u\in\mathcal U$ has $g(u)\le\mathrm{hi}$ if and only if it is congruent to an integer in
$[\mathrm{lo},\mathrm{hi}]$. The elements of $\mathcal U$ in $[b,b+t-1]$ are the $b+s$, which $g$ fixes.
So $A=B+Z$, where $B$ is the multiset of the $b+s$, $s<t$, and $Z$ the multiset of the $g(u)$ for the
$u\in\mathcal U$ outside $[b,b+t-1]$ with $g(u)\le\mathrm{hi}$. Thus $Z$ has $n_{l,\kappa}(b,\mathcal U)$
elements, and all elements of $A$ lie in $[\mathrm{lo},\mathrm{hi}]$; if $\mathrm{hi}<\mathrm{lo}$, then
$A$ is empty.

Now $Y$ and $B$ both have $t$ elements and $Y\le B+Z$. So the multisets $Y-B$ and $B-Y$ have the same
number $m$ of elements, and $Y-B\le Z$ gives $m\le|Z|$. As $\sum Y-\sum B=\sum(Y-B)-\sum(B-Y)$, with both
sums over $m$ elements of $[\mathrm{lo},\mathrm{hi}]$, $|\sum B-\sum Y|\le m(\mathrm{hi}-\mathrm{lo})\le
|Z|(\mathrm{hi}-\mathrm{lo})$. Finally $\sum B-\sum Y=\sum_s\nu_s$, and
$\mathrm{hi}-\mathrm{lo}=t-1+2\kappa\le t+2\kappa$.
% Lean: shiftSumBound_holds (KapSum), with g u = lo + (u - lo) % l (g_ge, g_fix, g_mod, g_cong) and the
% map rho on ZMod l; hY and hYA (Y <= A), hA with hfirst (A = B + Z), hR and hcard (|Z| = mwCount); the
% last paragraph is we_msum (Multiset.sub_add_inter, sum_le_card_nsmul, card_nsmul_le_sum).
\end{proof}

\begin{lemma}\label{is-onto-step-exp-bound-link-depth-arith}
Let $k,j,t,q,l,h,\kappa,n,n'\geq0$ be integers, $N=k+j+t+q$, $U,U'\in\mathbf Z^N$ strictly decreasing,
and $K_c,K_p$ integers. Suppose that $U(a)\geq0$ and $U'(a)\le h$ for every $a<N$,
$\sum_aU(a)=\sum_aU'(a)$, $|K_c|\le n(t+2\kappa)$, $|K_p|\le n'(t+2\kappa)$ and
\[
 \sum_{k+j+q\le a<N}U'(a)\le\sum_{k+j\le a<k+j+t}U(a)+lK_p-K_c.
\]
Then
\[
 \sum_{a<k+j}\big(U(a)-U'(a)\big)\le qh+(ln'+n)(t+2\kappa).
\]
That $U$ and $U'$ decrease strictly is not used in the proof.
\end{lemma}
\leanref{Formalism/PieriBlockIdeals/Statement.lean}{Formalism.CohomologicalDescent.WEArith}
\leanref{Formalism/PieriBlockIdeals/PieriColumnMinors/WeightExchangeBoundHolds.lean}{Formalism.CohomologicalDescent.weightExchangeBound_holds}
\begin{proof}
For a condition $C$ on $a$ let $[C]$ be $1$ if $C$ holds and $0$ otherwise. For every $a<N$,
\[
 [a<k+j]\big(U(a)-U'(a)\big)+[k+j\le a<k+j+t]\,U(a)\le[k+j\le a<k+j+q]\,h+[a\geq k+j+q]\,U'(a)
 +U(a)-U'(a):
\]
for $a<k+j$ both sides are $U(a)-U'(a)$, and for $a\geq k+j$ the inequality follows from $U'(a)\le h$
and $U(a)\geq0$. Sum over $a$. At most $q$ indices satisfy $k+j\le a<k+j+q$, and
$\sum_a(U(a)-U'(a))=0$, so
\[
 \sum_{a<k+j}\big(U(a)-U'(a)\big)+\sum_{k+j\le a<k+j+t}U(a)\le qh+\sum_{k+j+q\le a<N}U'(a).
\]
With the hypothesis, $\sum_{a<k+j}(U(a)-U'(a))\le qh+lK_p-K_c\le qh+(ln'+n)(t+2\kappa)$.
% Lean: weightExchangeBound_holds (WEArith): hq (the count of the indices in [k + j, k + j + q)), hpt (by
% split_ifs and omega), the sum hsumle and nlinarith. StrictAnti U and StrictAnti U' are introduced as _.
\end{proof}

\begin{lemma}\label{is-onto-step-exp-bound-link-depth-chain}
Let $k,j,t,q,h,\kappa\geq0$ be integers and $l$ a prime with $t+2\kappa<l$, let $N=k+j+t+q$, and let $K$
be a field of characteristic $l$. Let $U,U',V,V'\in\mathbf Z^N$ be strictly decreasing,
$b,b',z,z'\in\mathbf Z$, and $\nu^c,\nu^p\in\mathbf Z^t$. Suppose that:
\begin{enumerate}
\item $U(a)\geq0$ and $U'(a)\le h$ for every $a$, $\sum_aU(a)=\sum_aU'(a)$, and $U(a)=b+(k+j+t-1-a)$ for
  $k+j\le a<k+j+t$;
\item $V$ takes the values $b'+s$ for $0\le s<t$;
\item $\prod_a(Y+U(a)+z)=\prod_a(Y+U'(a)+z)$ and $\prod_a(Y+V(a)+z')=\prod_a(Y+V'(a)+z')$ in $K[Y]$;
\item $|\nu^c_s|\le\kappa$ and $|\nu^p_s|\le\kappa$ for every $s$;
\item the multiset of the residues modulo $l$ of the $b+s-\nu^c_s$, $s<t$, is contained in that of the
  $U'(a)$, $a<N$, and the multiset of the residues of the $b'+s-\nu^p_s$ in that of the $V'(a)$;
\item $\sum_{k+j+q\le a<N}U'(a)\le\sum_{s<t}\big(b+s+l\nu^p_s-\nu^c_s\big)$.
\end{enumerate}
Let $\mathcal U$ and $\mathcal V$ be the sets of values of $U$ and $V$. Then
\[
 \sum_{a<k+j}\big(U(a)-U'(a)\big)\le qh+\big(l\,n_{l,\kappa}(b',\mathcal V)+n_{l,\kappa}(b,\mathcal U)\big)(t+2\kappa)
\]
(Definition~\ref{is-onto-step-exp-bound-link-depth-reflected}(4)). That $U'$ and $V'$ decrease
strictly is not used in the proof.
\end{lemma}
\leanref{Formalism/PieriBlockIdeals/WEChain.lean}{Formalism.CohomologicalDescent.WEChain}
\leanref{Formalism/PieriBlockIdeals/PieriColumnMinors/WeightExchangeChainHolds.lean}{Formalism.CohomologicalDescent.weightExchangeChain_holds}
\begin{proof}
By (3) and Lemma~\ref{is-onto-step-exp-bound-link-depth-content}, the residues modulo $l$ of the $U(a)$
and of the $U'(a)$ form the same multiset, and so do those of the $V(a)$ and of the $V'(a)$. As $U$ and
$V$ are injective, these are the multisets of the residues of the elements of $\mathcal U$ and
$\mathcal V$. By (1), $\mathcal U$ contains $b+s=U(k+j+t-1-s)$ for $s<t$, and by (2), $\mathcal V$
contains $b'+s$. So (4), (5) and Lemma~\ref{is-onto-step-exp-bound-link-depth-shift} give
$|\sum_s\nu^c_s|\le n_{l,\kappa}(b,\mathcal U)(t+2\kappa)$ and
$|\sum_s\nu^p_s|\le n_{l,\kappa}(b',\mathcal V)(t+2\kappa)$. By (1),
$\sum_{k+j\le a<k+j+t}U(a)=\sum_{s<t}(b+s)$, so (6) reads
\[
 \sum_{k+j+q\le a<N}U'(a)\le\sum_{k+j\le a<k+j+t}U(a)+l\sum_s\nu^p_s-\sum_s\nu^c_s.
\]
Lemma~\ref{is-onto-step-exp-bound-link-depth-arith}, with $K_c=\sum_s\nu^c_s$ and $K_p=\sum_s\nu^p_s$,
gives the bound.
% Lean: weightExchangeChain_holds (WEChain): contentLinkage_holds for (U, U') and (V, V'), hmemU and
% hmemV, shiftSumBound_holds twice (map_cast_image_val_eq), sum_filter_block_eq for hfilt', and
% weightExchangeBound_holds. The hypothesis StrictAnti V' (hV') is not used, and StrictAnti U' only
% reaches WEArith, which does not use it.
\end{proof}

\begin{definition}\label{is-onto-step-exp-bound-link-depth-datum}
Let $F$ be a field, $E$ a finite set and $t,n\geq0$ integers. A \emph{weight datum} $\mathcal D$ for
$(E,t,n,F)$ consists of $\mathrm a\in F$, matrices $\Theta_i\in M_n(F)$ and $\Phi_i\in M_n(F[[X]])$ for
$i\in E$, and maps $b\colon\{0,\dots,n-1\}\to\mathbf N$ and
$\mathrm{pos}\colon\{0,\dots,t-1\}\to\{0,\dots,n-1\}$. Let $\sigma$ be a permutation of $E$ and $l\geq0$
an integer, and for $i\in E$ and $s<t$ let $U'_i\in\mathbf Z^n$, $\mathrm{bot}_i,\nu_i(s)\in\mathbf Z$ and
$m_i\in\mathbf N$; put $\alpha_i(s)=l\nu_{\sigma(i)}(s)+\mathrm{bot}_i+s-\nu_i(s)$. Then $\mathcal D$ has
the \emph{weight properties} for $(\sigma,l,U',\mathrm{bot},\nu,m)$ if, for all $i$ and $s$:
\begin{enumerate}
\item $\mathrm a\ne0$; pos is injective, and $b(j)=b(\mathrm{pos}(s))$ only for $j=\mathrm{pos}(s)$;
\item $\Theta_i$ is block upper triangular for $b$
  (Definition~\ref{is-onto-step-exp-bound-link-depth-data}(4)),
  $\det(Y-\Theta_i)=\prod_{j<n}\big(Y+\mathrm aU'_i(j)\big)$, and
  $\Theta_i(\mathrm{pos}(s),\mathrm{pos}(s))=-\mathrm a\big(\mathrm{bot}_i+s-\nu_i(s)\big)$;
\item $\Phi_i$ is block upper triangular for $b$, $\Phi_i(\mathrm{pos}(s),\mathrm{pos}(s))$ is a unit of
  $F[[X]]$ times $X^{\alpha_i(s)^+}$, and $X^{m_i}$ divides the determinant of the $t\times t$ matrix
  $\big(\Phi_i(\mathrm{pos}(s),\mathrm{pos}(s'))\big)_{s,s'<t}$.
\end{enumerate}
\end{definition}
\leanref{Formalism/PieriBlockIdeals/LemmaWAssembled.lean}{Formalism.CohomologicalDescent.WData}
\leanref{Formalism/PieriBlockIdeals/LemmaWAssembled.lean}{Formalism.CohomologicalDescent.WProps}

\begin{lemma}\label{is-onto-step-exp-bound-link-depth-assembled}
Let $l$ be a prime, $t,n,h\geq0$ integers, $F$ a field of characteristic $l$, $E$ a finite set with a
permutation $\sigma$, and for $i\in E$ and $s<t$ let $U'_i\in\mathbf Z^n$, $\mathrm{bot}_i,\nu_i(s)\in\mathbf
Z$ and $m_i\in\mathbf N$. Let $\mathcal D$ be a weight datum for $(E,t,n,F)$ with the weight properties
for $(\sigma,l,U',\mathrm{bot},\nu,m)$, and $\alpha_i(s)$ as there
(Definition~\ref{is-onto-step-exp-bound-link-depth-datum}). Suppose that $0\le\mathrm{bot}_i+s\le h$ and
$0\le\alpha_i(s)\le h$ for all $i\in E$ and $s<t$. Then for every $i$:
\begin{enumerate}
\item $(l-1)|\nu_i(s)|\le h$ for every $s<t$;
\item the multiset of the residues modulo $l$ of the $\mathrm{bot}_i+s-\nu_i(s)$, $s<t$, is contained in
  that of the $U'_i(j)$, $j<n$;
\item $m_i\le\sum_{s<t}\big(\mathrm{bot}_i+s+l\nu_{\sigma(i)}(s)-\nu_i(s)\big)$.
\end{enumerate}
\end{lemma}
\leanref{Formalism/PieriBlockIdeals/LemmaWAssembled.lean}{Formalism.CohomologicalDescent.LemmaWAssembled}
\leanref{Formalism/PieriBlockIdeals/PieriColumnMinors/RegularCuspidalTrace.lean}{Formalism.CohomologicalDescent.twisted_block_bounds}
\begin{proof}
(1) Fix $s$, choose $i_0\in E$ with $|\nu_{i_0}(s)|$ largest, and put $i_1=\sigma^{-1}(i_0)$. Then
$0\le l\nu_{i_0}(s)+\mathrm{bot}_{i_1}+s-\nu_{i_1}(s)\le h$ and $0\le\mathrm{bot}_{i_1}+s\le h$. If
$\nu_{i_0}(s)\geq0$, the upper bound gives $l\nu_{i_0}(s)\le h+\nu_{i_1}(s)$; otherwise the lower bound
gives $-l\nu_{i_0}(s)\le h-\nu_{i_1}(s)$. In both cases
$l|\nu_{i_0}(s)|\le h+|\nu_{i_1}(s)|\le h+|\nu_{i_0}(s)|$, so
$(l-1)|\nu_i(s)|\le(l-1)|\nu_{i_0}(s)|\le h$.

(2) As $\Theta_i$ is block upper triangular, $\det(Y-\Theta_i)$ is the product of the determinants of
the diagonal blocks of $Y-\Theta_i$. The block of $\mathrm{pos}(s)$ is $1\times1$, with entry
$Y-\Theta_i(\mathrm{pos}(s),\mathrm{pos}(s))$, and these blocks are distinct, as $b\circ\mathrm{pos}$ is
injective. So $\prod_s\big(Y+\mathrm a(\mathrm{bot}_i+s-\nu_i(s))\big)$ divides
$\prod_j\big(Y+\mathrm aU'_i(j)\big)$, and the multiset of the $-\mathrm a(\mathrm{bot}_i+s-\nu_i(s))$ is
contained in that of the $-\mathrm aU'_i(j)$. As $\mathrm a\ne0$ and $\mathbf Z/l\to F$ is injective,
$y\mapsto-\mathrm ay$ is injective on the image of $\mathbf Z/l$, which gives (2).

(3) The matrix $\big(\Phi_i(\mathrm{pos}(s),\mathrm{pos}(s'))\big)$ is block upper triangular for the
injective map $b\circ\mathrm{pos}$, so its determinant is the product of its diagonal entries, a unit
times $X^{\sum_s\alpha_i(s)}$, as $\alpha_i(s)\geq0$. As $X$ is prime in $F[[X]]$ and $X^{m_i}$ divides
this determinant, $m_i\le\sum_s\alpha_i(s)$, which is the right side of (3).
% Lean: twisted_block_bounds (LemmaWAssembled): (1) is sub_one_mul_abs_kap_le (Finite.exists_max); (2)
% goes through Matrix.BlockTriangular.charmatrix, lw_det_block, Finset.prod_dvd_prod_of_subset,
% Polynomial.roots.le_of_dvd and map_cast_le_of_map_neg_mul_le; (3) is le_sum_of_dvd_det_submatrix
% (lw_det_inj, pow_dvd_pow_iff).
\end{proof}

\begin{lemma}\label{is-onto-step-exp-bound-link-depth-exchange}
Let $k,j,t,q,h\geq0$ be integers and $l$ a prime with $t+2\kappa<l$, where $\kappa=\lfloor h/(l-1)\rfloor$;
let $N=k+j+t+q$, $F$ a field of characteristic $l$, and $E$ a finite set with a permutation $\sigma$. For
$i\in E$ and $s<t$ let $U_i,U'_i\in\mathbf Z^N$, $\mathrm{bot}_i,z_i,\nu_i(s)\in\mathbf Z$ and
$m_i\in\mathbf N$, and let $\tau\in E$. Let $\mathcal D$ be a weight datum for $(E,t,N,F)$ with the weight
properties for $(\sigma,l,U',\mathrm{bot},\nu,m)$, and $\alpha_i(s)$ as there
(Definition~\ref{is-onto-step-exp-bound-link-depth-datum}). Suppose that for all $i$ and
$s<t$:
\begin{enumerate}
\item $0\le\mathrm{bot}_i+s\le h$ and $0\le\alpha_i(s)\le h$;
\item $U_i$ and $U'_i$ decrease strictly, $U_i$ takes the value $\mathrm{bot}_i+s$, and
  $\prod_a(Y+U_i(a)+z_i)=\prod_a(Y+U'_i(a)+z_i)$ in $F[Y]$;
\item $U_\tau(a)\geq0$ and $0\le U'_\tau(a)\le h$ for every $a$, $\sum_aU_\tau(a)=\sum_aU'_\tau(a)$,
  $U_\tau(a)=\mathrm{bot}_\tau+(k+j+t-1-a)$ for $k+j\le a<k+j+t$, and
  $m_\tau=\sum_{k+j+q\le a<N}U'_\tau(a)$.
\end{enumerate}
Let $n_i=n_{l,\kappa}(\mathrm{bot}_i,\{U_i(a):a<N\})$
(Definition~\ref{is-onto-step-exp-bound-link-depth-reflected}(4)). Then
\[
 \sum_{a<k+j}\big(U_\tau(a)-U'_\tau(a)\big)\le qh+\big(l\,n_{\sigma(\tau)}+n_\tau\big)(t+2\kappa).
\]
The bound $U'_\tau(a)\geq0$ is not used in the proof.
\end{lemma}
\leanref{Formalism/PieriBlockIdeals/LemmaWAssembled.lean}{Formalism.CohomologicalDescent.WEFull}
\leanref{Formalism/PieriBlockIdeals/PieriColumnMinors/RegularCuspidalTrace.lean}{Formalism.CohomologicalDescent.weightExchangeFull_holds}
\begin{proof}
By (1) and Lemma~\ref{is-onto-step-exp-bound-link-depth-assembled}, for every $i$:
$(l-1)|\nu_i(s)|\le h$, so $|\nu_i(s)|\le\kappa$ as $l\geq2$; the residues of the
$\mathrm{bot}_i+s-\nu_i(s)$ are contained in those of the $U'_i(j)$; and
$m_i\le\sum_s(\mathrm{bot}_i+s+l\nu_{\sigma(i)}(s)-\nu_i(s))$. Apply
Lemma~\ref{is-onto-step-exp-bound-link-depth-chain} to $U=U_\tau$, $U'=U'_\tau$, $V=U_{\sigma(\tau)}$,
$V'=U'_{\sigma(\tau)}$, $b=\mathrm{bot}_\tau$, $b'=\mathrm{bot}_{\sigma(\tau)}$, $z=z_\tau$,
$z'=z_{\sigma(\tau)}$, $\nu^c=\nu_\tau$ and $\nu^p=\nu_{\sigma(\tau)}$. Its condition (1) is (3), its (2)
and (3) come from (2), its (4) and (5) from the bounds just proved at $\tau$ and $\sigma(\tau)$, and its
(6) is the last of them at $\tau$, as $m_\tau=\sum_{k+j+q\le a<N}U'_\tau(a)$. Its conclusion is the
claim.
% Lean: weightExchangeFull_holds (WEFull): twisted_block_bounds, then hkappa (Int.le_ediv_iff_mul_le) and
% weightExchangeChain_holds at tau and sigma tau. Only the upper bound of hU'h is used.
\end{proof}

\begin{lemma}\label{is-onto-step-exp-bound-link-depth-matrices}
Let $R$ be a domain, $\pi\in R$ nonzero, $h,n\geq0$ integers, $b\colon\{0,\dots,n-1\}\to\mathbf N$, and
$\Phi,\Phi'\in M_n(R)$ with $\Phi\Phi'=\pi^h$ and $\Phi$ block upper triangular for $b$. Then $\Phi'$ is
block upper triangular for $b$, and $\Phi(s,s)\Phi'(s,s)=\pi^h$ for every $s$ such that $b(j)\ne b(s)$
for $j\ne s$.
\end{lemma}
\leanref{Formalism/PieriBlockIdeals/PieriColumnMinors/HeckeLFunctionIdealEulerProduct.lean}{Formalism.CohomologicalDescent.FrobeniusPeriodChange.blockTriangular_of_mul_eq_pow_smul_one}
\leanref{Formalism/PieriBlockIdeals/LemmaWFromLattice.lean}{Formalism.CohomologicalDescent.FrobeniusPeriodChange.BlockHeightLine}
\leanref{Formalism/PieriBlockIdeals/PieriColumnMinors/HeckeLFunctionIdealEulerProduct.lean}{Formalism.CohomologicalDescent.FrobeniusPeriodChange.block_height_on_line_holds}
\begin{proof}
Over the field of fractions of $R$, $\Phi$ is invertible with inverse $\pi^{-h}\Phi'$, and the inverse of
an invertible block upper triangular matrix is block upper triangular; so $\Phi'$ is. The $(s,s)$ entry
of $\Phi\Phi'$ is $\sum_j\Phi(s,j)\Phi'(j,s)$. Here $\Phi(s,j)=0$ if $b(j)<b(s)$, $\Phi'(j,s)=0$ if
$b(s)<b(j)$, and $b(j)=b(s)$ only for $j=s$; so this entry is $\Phi(s,s)\Phi'(s,s)$.
% Lean: blockTriangular_of_mul_eq_pow_smul_one (FractionRing, Matrix.blockTriangular_inv_of_blockTriangular)
% and block_height_on_line_holds (BlockHeightLine, Finset.sum_eq_single).
\end{proof}

\begin{lemma}\label{is-onto-step-exp-bound-link-depth-reduction}
Let $l$ be a prime, $O$ a commutative ring that is a $\mathbf Z_l$-algebra, $F$ a field of characteristic
$l$ and $\mathrm{red}\colon O\to F$ a ring homomorphism.
\begin{enumerate}
\item Let $u\in\mathbf Z_l^\times$ and $c\in\mathbf Z_l$ with $u=1+lc$, $m\in\mathbf Z$, and
  $r\in\mathbf Z_l$ with $u^{-m}=1+lr$. Then $\mathrm{red}(r)=-m\,\mathrm{red}(c)$.
\item Let $\lambda,a\in O$ and $g\in O[[X]]$, with $X-\lambda$ dividing $g-a$. Then $g(0)=a+\lambda e$ for
  some $e\in O$; if moreover $\mathrm{red}(\lambda)=0$, then $\mathrm{red}(g(0))=\mathrm{red}(a)$.
\item Reduction of coefficients by red maps $(X-l)^m$ to $X^m$ for every $m\geq0$.
\end{enumerate}
\end{lemma}
\leanref{Formalism/PieriBlockIdeals/PieriColumnMinors/HeckeLFunctionIdealEulerProduct.lean}{Formalism.CohomologicalDescent.FrobeniusPeriodChange.reduction_log_unit_pow_eq}
\leanref{Formalism/PieriBlockIdeals/LemmaWFromLattice.lean}{Formalism.CohomologicalDescent.FrobeniusPeriodChange.EvalCong}
\leanref{Formalism/PieriBlockIdeals/PieriColumnMinors/HeckeLFunctionIdealEulerProduct.lean}{Formalism.CohomologicalDescent.FrobeniusPeriodChange.evaluation_congruence_holds}
\leanref{Formalism/PieriBlockIdeals/PieriColumnMinors/HeckeLFunctionIdealEulerProduct.lean}{Formalism.CohomologicalDescent.FrobeniusPeriodChange.map_X_sub_C_prime_pow_eq_X_pow}
\begin{proof}
(1) By induction on $k\geq0$, $u^k=1+lkc+l^2s_k$ with $s_k\in\mathbf Z_l$. If $m=k\geq0$, then
$1=u^{-k}u^k=(1+lr)(1+lkc+l^2s_k)$, so $l\big(r+kc+l(rkc+s_k+lrs_k)\big)=0$; as $\mathbf Z_l$ is a domain
and $l\ne0$, $r+kc$ is divisible by $l$. If $m=-k$, then $1+lr=u^k$, so $r=kc+ls_k$. As $\mathrm{red}(l)=0$, (1)
follows. (2) Write $g-a=(X-\lambda)p$; then $g(0)-a=-\lambda p(0)$. (3) $\mathrm{red}(l)=0$.
% Lean: reduction_log_unit_pow_eq (the induction key and Int.eq_nat_or_neg), evaluation_congruence_holds
% (EvalCong) and map_X_sub_C_prime_pow_eq_X_pow; red(r) stands for red (algebraMap Z_l O r).
\end{proof}

\begin{lemma}\label{is-onto-step-exp-bound-link-depth-rank-one}
Let $l$ be a prime, $O$ a commutative ring that is a $\mathbf Z_l$-algebra, $F$ a field of characteristic
$l$, $\mathrm{red}\colon O\to F$ a ring homomorphism, $u\in\mathbf Z_l^\times$ and $c\in\mathbf Z_l$ with
$u=1+lc$, $h\geq0$, and $\mathrm{wt}_i$ integers for $i$ in a set $E$. Let $N$ be a lattice datum
of rank $1$ over $O$, indexed by $E$, that has height $h$ with weights $(\mathrm{wt}_i)_{i\in E}$
(Definition~\ref{is-onto-step-exp-bound-link-depth-data}(1)), and write $P_i$ and $D_i$ for the single
entries of its matrices. Then for every $i$, $P_i=v(X-l)^{\mathrm{wt}_i^+}$ for a unit $v$ of $O[[X]]$,
and $\mathrm{red}(D_i(0))=-\mathrm{wt}_i\,\mathrm{red}(c)$.
\end{lemma}
\leanref{Formalism/PieriBlockIdeals/LemmaWFromLattice.lean}{Formalism.CohomologicalDescent.FrobeniusPeriodChange.RankOneModel}
\leanref{Formalism/PieriBlockIdeals/PieriColumnMinors/HeckeLFunctionIdealEulerProduct.lean}{Formalism.CohomologicalDescent.FrobeniusPeriodChange.rank_one_model_holds}
\begin{proof}
The first claim is condition (c) for $n=1$. By (e), $u^{-\mathrm{wt}_i}=1+lr$ for some $r\in\mathbf Z_l$
with $Y-\Theta_i=Y-r$, so $\Theta_i=r$. By (d) and
Lemma~\ref{is-onto-step-exp-bound-link-depth-reduction}(2) with $\lambda=l$,
$\mathrm{red}(D_i(0))=\mathrm{red}(\Theta_i)=\mathrm{red}(r)$, which is $-\mathrm{wt}_i\,\mathrm{red}(c)$
by Lemma~\ref{is-onto-step-exp-bound-link-depth-reduction}(1).
% Lean: rank_one_model_holds (RankOneModel), with Matrix.det_fin_one, the constant coefficient of the
% charpoly, evaluation_congruence_holds and reduction_log_unit_pow_eq.
\end{proof}

\begin{lemma}\label{is-onto-step-exp-bound-link-depth-frobenius-entry}
Let $F$ be a field, $l\geq1$ and $s\geq0$ integers, $p\in F[[X]]$, $v,w,w'$ units of $F[[X]]$ and
$\nu,\nu'\in\mathbf Z$, with $\mathrm{Fr}_l(w')X^{l\nu'}vX^s=p\,wX^{\nu}$ in $F((X))$
(Definition~\ref{is-onto-step-exp-bound-link-depth-series}(4)). Then $e=l\nu'+s-\nu\geq0$, and $p$ is a
unit of $F[[X]]$ times $X^e$.
\end{lemma}
\leanref{Formalism/PieriBlockIdeals/LemmaWFromLattice.lean}{Formalism.CohomologicalDescent.FrobeniusPeriodChange.PhiDiag}
\leanref{Formalism/PieriBlockIdeals/PieriColumnMinors/HeckeLFunctionIdealEulerProduct.lean}{Formalism.CohomologicalDescent.FrobeniusPeriodChange.frobenius_diagonal_holds}
\begin{proof}
As $\mathrm{Fr}_l$ is a ring endomorphism, $\mathrm{Fr}_l(w')$ is a unit, so
$A=\mathrm{Fr}_l(w')vw^{-1}$ is a unit of $F[[X]]$, and $p=AX^e$ in $F((X))$. If $e<0$, then
$A=pX^{-e}$ would have constant term $0$. So $e\geq0$ and $p=AX^e$.
% Lean: frobenius_diagonal_holds (PhiDiag), with frobenius_laws_hold and HahnSeries.ofPowerSeries.
\end{proof}

\begin{lemma}\label{is-onto-step-exp-bound-link-depth-base-change}
Let $k$ be a commutative ring, $n,h\geq0$ and $l\geq1$ integers, $\bar u\in k[[X]]$,
$g,g'\in GL_n(k[[X]])$ and $D,P\in M_n(k[[X]])$, and put $\Gamma=g^{-1}(1+XD)\gamma_{\bar u}(g)$ and
$\Phi=g^{-1}P\,\mathrm{Fr}_l(g')$ (Definition~\ref{is-onto-step-exp-bound-link-depth-series}(4)).
\begin{enumerate}
\item $\Gamma(0)=1$; the coefficient $\Theta$ of $X$ in $\Gamma$ is $g(0)^{-1}D(0)g(0)$; and
  $\det(Y-\Theta)=\det(Y-D(0))$.
\item If $PM=X^h$ for some $M\in M_n(k[[X]])$, then $\Phi M'=X^h$ for some $M'$. If $x\in k[[X]]$
  divides the $r$-minors of $P$, then it divides the $r$-minors of $\Phi$.
\end{enumerate}
\end{lemma}
\leanref{Formalism/PieriBlockIdeals/LemmaWFromLattice.lean}{Formalism.CohomologicalDescent.FrobeniusPeriodChange.ThetaBaseChange}
\leanref{Formalism/PieriBlockIdeals/PieriColumnMinors/HeckeLFunctionIdealEulerProduct.lean}{Formalism.CohomologicalDescent.FrobeniusPeriodChange.theta_base_change_holds}
\leanref{Formalism/PieriBlockIdeals/LemmaWFromLattice.lean}{Formalism.CohomologicalDescent.FrobeniusPeriodChange.PhiBaseChange}
\leanref{Formalism/PieriBlockIdeals/PieriColumnMinors/HeckeLFunctionIdealEulerProduct.lean}{Formalism.CohomologicalDescent.FrobeniusPeriodChange.frobenius_base_change_holds}
\begin{proof}
(1) As $\gamma_{\bar u}(X)=X+\bar uX^2$, $\gamma_{\bar u}$ keeps the constant term and the coefficient of
$X$ of a series. For matrices $A,B$ over $k[[X]]$, with coefficients $A_0,A_1,\dots$, the coefficient of
$X$ in $AB$ is $A_1B_0+A_0B_1$. Let $A=g^{-1}$; then $A_0g_0=1$ and $A_1g_0+A_0g_1=0$. As
$(1+XD)_0=1$ and $(1+XD)_1=D_0$, $\Gamma_0=A_0g_0=1$ and
$\Gamma_1=A_1g_0+A_0D_0g_0+A_0g_1=A_0D_0g_0$. The characteristic polynomial is invariant under
conjugation.

(2) Put $M'=\mathrm{Fr}_l(g')^{-1}Mg$; then $\Phi M'=g^{-1}PMg=X^h$. For matrices $G$ and $H$, every
$r\times r$ minor of $GPH$ is, by multilinearity of the determinant in the rows and in the columns, a
linear combination of determinants of $r\times r$ submatrices of $P$, each of which is a minor of $P$ up
to sign or is $0$.
% Lean: theta_base_change_holds (ThetaBaseChange: gamma_laws_hold, hcoef, Matrix.charpoly_mul_comm) and
% frobenius_base_change_holds (PhiBaseChange), whose minor part is minors_dvd_mul_holds
% (det_submatrix_mul_eq_sum and dvd_det_submatrix_of_minorsDvd).
\end{proof}

\begin{lemma}\label{is-onto-step-exp-bound-link-depth-theta-entry}
Let $k$ be a field, $\bar u,g,y\in k[[X]]$, $w$ a unit of $k[[X]]$ and $\nu\in\mathbf Z$, with
$\gamma_{\bar u}(w)\big(X(1+\bar uX)\big)^{\nu}(1+Xg)=y\,wX^{\nu}$ in $k((X))$. Then $y(0)=1$, and the
coefficient of $X$ in $y$ is $\nu\bar u(0)+g(0)$.
\end{lemma}
\leanref{Formalism/PieriBlockIdeals/LemmaWFromLattice.lean}{Formalism.CohomologicalDescent.FrobeniusPeriodChange.ThetaDiag}
\leanref{Formalism/PieriBlockIdeals/PieriColumnMinors/HeckeLFunctionIdealEulerProduct.lean}{Formalism.CohomologicalDescent.FrobeniusPeriodChange.theta_diagonal_holds}
\begin{proof}
Put $V=1+\bar uX$, a unit with constant term $1$ and coefficient $\bar u(0)$ of $X$. For every
$\nu\in\mathbf Z$, $V^\nu$ has constant term $1$ and coefficient $\nu\bar u(0)$ of $X$, by the rule for
the first two coefficients of a product. Cancelling $X^\nu$, $y=\gamma_{\bar u}(w)V^\nu(1+Xg)w^{-1}$ in
$k[[X]]$. As $\gamma_{\bar u}$ keeps the first two coefficients, the constant term of $y$ is
$w(0)w^{-1}(0)=1$, and its coefficient of $X$ is
$w_1w^{-1}_0+w_0w^{-1}_1+\nu\bar u(0)+g(0)=\nu\bar u(0)+g(0)$, as $ww^{-1}=1$.
% Lean: theta_diagonal_holds (ThetaDiag): gamma_laws_hold, hc1 (the coefficient of X of a product), hV1
% for the powers of V, the cancellation hc in LaurentSeries, and hww0, hww1.
\end{proof}

\begin{lemma}\label{is-onto-step-exp-bound-link-depth-unit-series-binomial}
Let $l$ be a prime and $O$ a commutative ring that is a $\mathbf Z_l$-algebra, with the notation of
Definition~\ref{is-onto-step-exp-bound-link-depth-series}, and write $B_c=1+\beta_c=(1+X)^c$ for
$c\in\mathbf Z_l$.
\begin{enumerate}
\item For $c,d\in\mathbf Z_l$, $B_{c+d}=B_cB_d$ and $\beta_d(\beta_c)=\beta_{dc}$. For an integer $n\geq0$,
  $\beta_n=(1+X)^n-1$; in particular $\beta_1=X$.
\item For $c\in\mathbf Z_l$, $\mathrm{act}_c$ is a ring endomorphism of $O[[X]]$ that fixes the
  constants. For $c,d\in\mathbf Z_l$, $\mathrm{act}_c\circ\mathrm{act}_d=\mathrm{act}_{dc}$;
  $\mathrm{act}_1$ is the identity; $\varphi=\mathrm{act}_l$; and
  $\varphi\circ\mathrm{act}_c=\mathrm{act}_c\circ\varphi$.
\item $Xq=(1+X)^l-1=\varphi(X)$.
\item For $c\in\mathbf Z_l$, $\beta_c=X\beta'_c$ for a series $\beta'_c\in O[[X]]$ with constant term $c$.
\item For $d,e\in\mathbf Z_l$, $\beta_{d+le}=\beta_d+Xqs$ for some $s\in O[[X]]$.
\end{enumerate}
\end{lemma}
\leanref{Formalism/PieriBlockIdeals/CapelliCentralEnveloping.lean}{Formalism.CohomologicalDescent.HeightDescentMinors.ActLaws}
\leanref{Formalism/PieriBlockIdeals/PieriColumnMinors/HeightDescentMinors.lean}{Formalism.CohomologicalDescent.HeightDescentMinors.action_laws_holds}
\leanref{Formalism/PieriBlockIdeals/PieriColumnMinors/HeightDescentMinors.lean}{Formalism.CohomologicalDescent.HeightDescentMinors.binomialSeries_subst_comp}
\leanref{Formalism/PieriBlockIdeals/CapelliCentralEnveloping.lean}{Formalism.CohomologicalDescent.HeightDescentMinors.QMul}
\leanref{Formalism/PieriBlockIdeals/PieriColumnMinors/HeightDescentMinors.lean}{Formalism.CohomologicalDescent.HeightDescentMinors.q_multiplicative_holds}
\leanref{Formalism/PieriBlockIdeals/PieriColumnMinors/HeightDescentMinors.lean}{Formalism.CohomologicalDescent.HeightDescentMinors.br_eq_X_mul_const}
\leanref{Formalism/PieriBlockIdeals/PieriColumnMinors/HeightDescentMinors.lean}{Formalism.CohomologicalDescent.HeightDescentMinors.binomialSeriesSubOne_add_prime_mul_cong}
\begin{proof}
(1) The binomial coefficients $\binom cm$ of $c\in\mathbf Z_l$ lie in $\mathbf Z_l$. So it suffices to
prove the identities in $\mathbf Z_l[[X]]$ and map the coefficients to $O$, which commutes with
substitution. The first identity is Vandermonde's identity
$\binom{c+d}m=\sum_{i\le m}\binom ci\binom d{m-i}$, and the third is the binomial theorem. For the second,
the identity $(m+1)\binom c{m+1}=(c-m)\binom cm$ says that $(1+X)B_c'=cB_c$, where the prime is the
derivative. A series $H\in\mathbf Z_l[[X]]$ with $(1+X)H'=aH$ for some $a\in\mathbf Z_l$ and $H(0)=0$ is
$0$: the coefficient of $X^m$ gives $(m+1)H_{m+1}=(a-m)H_m$, and $\mathbf Z_l$ has no torsion. So two
solutions of $(1+X)H'=aH$ with the same constant term are equal. Put $F=B_d(\beta_c)$. By the chain rule,
and as $1+\beta_c=B_c$,
\[
 (1+X)F'=B_d'(\beta_c)\,(1+X)B_c'=c\,B_c\,B_d'(\beta_c)=c\,\big((1+X)B_d'\big)(\beta_c)=cd\,F,
\]
and $F(0)=1$. So $F=B_{dc}$, and subtracting $1$ gives $\beta_d(\beta_c)=\beta_{dc}$.

(2) Substitution of a series with constant term $0$ is a ring endomorphism that fixes the constants, and
$f(a)(b)=f\big(a(b)\big)$. So $\mathrm{act}_c(\mathrm{act}_df)=f\big(\beta_d(\beta_c)\big)=\mathrm{act}_{dc}f$
by (1). Also $\mathrm{act}_1f=f(X)=f$, and $\varphi f=f(\beta_l)=\mathrm{act}_lf$ by (1). Finally
$\varphi\circ\mathrm{act}_c=\mathrm{act}_{cl}=\mathrm{act}_c\circ\varphi$.

(3) By the binomial theorem, $(1+X)^l-1=\sum_{m\geq0}\binom l{m+1}X^{m+1}=Xq$; and
$\varphi(X)=(1+X)^l-1$.

(4) The series $\beta_c$ has constant term $0$ and coefficient $\binom c1=c$ of $X$.

(5) By (1) and (2), $\beta_{el}=\mathrm{act}_l(\beta_e)=\varphi(\beta_e)$. By (4) and (3),
$\varphi(\beta_e)=\varphi(X)\varphi(\beta'_e)=Xq\,\varphi(\beta'_e)$. So $B_{le}=1+Xqs'$ with
$s'=\varphi(\beta'_e)$, and by (1), $B_{d+le}=B_dB_{le}=B_d+Xq\,B_ds'$. Subtracting $1$ gives the claim
with $s=B_ds'$.
% Lean: action_laws_holds (ActLaws): binomialAction_binomialSeriesSubOne from binomialSeries_subst_comp
% (eq_of_binomial_ode_of_constantCoeff_eq, binomialSeries_one_add_X_mul_derivative, derivative_subst) and
% binomialSeries_subst_sub_one_eq_mul (map_binomialSeries_alt); binomialAction_comp_eq_mul, phi_eq_act_l;
% q_multiplicative_holds (X_mul_q_eq_one_add_X_pow_sub_one); br_eq_X_mul_const (coeff_one_br);
% binomialSeriesSubOne_add_prime_mul_cong (binomialSeries_add, dvd_subst_of_constantCoeff_zero).
\end{proof}

\begin{lemma}\label{is-onto-step-exp-bound-link-depth-unit-series-roots}
Let $l$ be a prime and $\mu_{l-1}$ the set of roots of $x^{l-1}=1$ in $\mathbf Z_l$.
\begin{enumerate}
\item Reduction modulo $l$ maps $\mu_{l-1}$ bijectively onto $\mathbf F_l^\times$. So $\mu_{l-1}$ has
  $l-1$ elements.
\item There is $d\in\mu_{l-1}$ such that $d^j-1$ is a unit of $\mathbf Z_l$ for every $0<j<l-1$.
\item If $l\ne2$, then $\prod_{d\in\mu_{l-1}}d=-1$.
\end{enumerate}
\end{lemma}
\leanref{Formalism/PieriBlockIdeals/CapelliCentralEnveloping.lean}{Formalism.CohomologicalDescent.HeightDescentMinors.MuFacts}
\leanref{Formalism/PieriBlockIdeals/PieriColumnMinors/HeightDescentMinors.lean}{Formalism.CohomologicalDescent.HeightDescentMinors.mu_facts_hold}
\leanref{Formalism/PieriBlockIdeals/PieriColumnMinors/HeightDescentMinors.lean}{Formalism.CohomologicalDescent.HeightDescentMinors.exists_mu_generator_pow_sub_one_isUnit}
\begin{proof}
(1) An element $d$ of $\mu_{l-1}$ has nonzero reduction, as $d^{l-1}=1$. Let $x\in\mathbf F_l^\times$ and
let $a\in\mathbf Z_l$ reduce to $x$. Then $a^{l-1}-1$ reduces to $x^{l-1}-1=0$, and the derivative
$(l-1)a^{l-2}$ of $Y^{l-1}-1$ at $a$ is a unit. By Hensel's lemma there is a unique root $d$ of
$Y^{l-1}-1$ in $\mathbf Z_l$ with $d\equiv a$ modulo $l$. These are the Teichmüller representatives
\cite[Chapter~II, \S4]{Ser79}.

(2) The group $\mathbf F_l^\times$ is cyclic of order $l-1$. By (1), some $d\in\mu_{l-1}$ reduces to a
generator $g$. For $0<j<l-1$, $g^j\ne1$, so $d^j-1$ has nonzero reduction and is a unit.

(3) The product $p$ of the elements of $\mu_{l-1}$ lies in $\mu_{l-1}$, and so does $-1$, as $l-1$ is
even. By (1), the reduction of $p$ is the product of the elements of $\mathbf F_l^\times$, which is $-1$
by Wilson's theorem \cite[Theorem~80]{HW08}. As reduction is injective on $\mu_{l-1}$, $p=-1$.
% Lean: mu_facts_hold (MuFacts): existsUnique_teichmuller_lift (hensels_lemma for X^(l-1) - 1),
% card_mu_eq_sub_one, prod_roots_unity_eq_neg_one (wilson_zmod_prod_units, from ZMod.prod_Ico_one_prime);
% exists_mu_generator_pow_sub_one_isUnit (IsCyclic.exists_generator for (ZMod l)^x).
\end{proof}

\begin{lemma}\label{is-onto-step-exp-bound-link-depth-unit-series-pi}
Let $l$ be a prime and $O$ a commutative ring that is a $\mathbf Z_l$-algebra, with the notation of
Definition~\ref{is-onto-step-exp-bound-link-depth-series}.
\begin{enumerate}
\item The series $\pi_0$ has constant term $0$. The map $\iota$ is a ring endomorphism of $O[[X]]$ that
  fixes the constants, with $\iota(X)=\pi_0$ and $\iota(X-l)=q_0$; $\iota(f)$ has the constant term of
  $f$; and $\iota\big(f(p)\big)=f\big(\iota(p)\big)$ for $f,p\in O[[X]]$ with $p(0)=0$.
\item For $d\in\mu_{l-1}$, $\mathrm{act}_d(\pi_0)=\pi_0$.
\item If $l\ne2$, then $\pi_0=wX^{l-1}$ for a unit $w$ of $O[[X]]$ with constant term $-1$.
\end{enumerate}
\end{lemma}
\leanref{Formalism/PieriBlockIdeals/LemmaWFromLattice.lean}{Formalism.CohomologicalDescent.FrobeniusPeriodChange.IncLaws}
\leanref{Formalism/PieriBlockIdeals/PieriColumnMinors/HeckeLFunctionIdealEulerProduct.lean}{Formalism.CohomologicalDescent.FrobeniusPeriodChange.pi0Subst_laws_hold}
\leanref{Formalism/PieriBlockIdeals/CapelliCentralEnveloping.lean}{Formalism.CohomologicalDescent.HeightDescentMinors.DivisibilityInvariant}
\leanref{Formalism/PieriBlockIdeals/PieriColumnMinors/HeightDescentMinors.lean}{Formalism.CohomologicalDescent.HeightDescentMinors.r2_inverse_holds}
\leanref{Formalism/PieriBlockIdeals/CapelliCentralEnveloping.lean}{Formalism.CohomologicalDescent.HeightDescentMinors.R1}
\leanref{Formalism/PieriBlockIdeals/PieriColumnMinors/HeightDescentMinors.lean}{Formalism.CohomologicalDescent.HeightDescentMinors.r1_holds}
\begin{proof}
(1) As $1\in\mu_{l-1}$ and $\beta_1=X$ has constant term $0$, so has $\pi_0$. The rest are properties of
the substitution of a series with constant term $0$, and $\iota(X-l)=\pi_0-l=q_0$.

(2) By Lemma~\ref{is-onto-step-exp-bound-link-depth-unit-series-binomial}(1) and~(2),
$\mathrm{act}_d(\pi_0)=\prod_{e\in\mu_{l-1}}\mathrm{act}_d(\beta_e)=\prod_{e\in\mu_{l-1}}\beta_{ed}$. As
$d^{l-1}=1$, $d\ne0$, and $e\mapsto ed$ maps $\mu_{l-1}$ injectively to itself, with $ed^{l-2}$ a preimage
of $e$. So the product is $\pi_0$.

(3) By Lemma~\ref{is-onto-step-exp-bound-link-depth-unit-series-binomial}(4), $\beta_d=X\beta'_d$ with
$\beta'_d(0)=d$. As $\mu_{l-1}$ has $l-1$ elements
(Lemma~\ref{is-onto-step-exp-bound-link-depth-unit-series-roots}(1)), $\pi_0=wX^{l-1}$ with
$w=\prod_d\beta'_d$, whose constant term is $\prod_dd=-1$ by
Lemma~\ref{is-onto-step-exp-bound-link-depth-unit-series-roots}(3). A series whose constant term is a
unit is a unit.
% Lean: pi0Subst_laws_hold (IncLaws, whose injectivity clause is Lemma
% is-onto-step-exp-bound-link-depth-unit-series-invariants(1)): constantCoeff_pi0 (one_mem_mu),
% inc_X_sub_l_eq_q0, subst_comp_subst_apply; r2_inverse_holds (act_finset_prod, Finset.prod_nbij for
% e -> e * d); r1_holds (br_eq_X_mul_shift, coeff_one_br, prod_roots_unity_eq_neg_one, card_mu_eq_sub_one).
\end{proof}

\begin{lemma}\label{is-onto-step-exp-bound-link-depth-unit-series-division}
Let $R$ be a commutative ring, $n\geq1$ an integer, $w$ a unit of $R[[X]]$ and $P=wX^n$.
\begin{enumerate}
\item Every $f\in R[[X]]$ can be written uniquely as $f=\sum_{j<n}g_j(P)X^j$ with $g_j\in R[[X]]$.
\item Let $s\in R[[X]]$ with $P(Xs)=P$, such that $s(0)^j-1$ is a unit of $R$ for every $0<j<n$. If
  $f\in R[[X]]$ and $f(Xs)=f$, then $f=g(P)$ for some $g\in R[[X]]$.
\end{enumerate}
\end{lemma}
\leanref{Formalism/PieriBlockIdeals/CapelliCentralEnveloping.lean}{Formalism.CohomologicalDescent.HeightDescentMinors.BasisGeneral}
\leanref{Formalism/PieriBlockIdeals/PieriColumnMinors/HeightDescentMinors.lean}{Formalism.CohomologicalDescent.HeightDescentMinors.general_basis_holds}
\leanref{Formalism/PieriBlockIdeals/PieriColumnMinors/HeightDescentMinors.lean}{Formalism.CohomologicalDescent.HeightDescentMinors.exists_subst_of_invariant}
\begin{proof}
Write $f_i$ for the coefficient of $X^i$ in $f$, and put $\delta(f)=w^{-1}\sum_{i\geq0}f_{i+n}X^i$. Then
$f=\sum_{j<n}f_jX^j+P\delta(f)$, and by induction on $K$,
\[
 f=\sum_{k<K}P^k\sum_{j<n}\big(\delta^kf\big)_jX^j+P^K\delta^K(f).
\]

(1) Existence. Put $g_j=\sum_{k\geq0}(\delta^kf)_jX^k$. For a series $g=\sum_kg_kX^k$ and every $K$,
$g(P)-\sum_{k<K}g_kP^k$ is divisible by $P^K$. So $f-\sum_{j<n}g_j(P)X^j$ is divisible by $P^K$, hence by
$X^{nK}$, and its coefficients of $X^m$ for $m<K$ vanish. As $K$ is arbitrary, it is $0$.

Uniqueness. It suffices that $\sum_{j<n}d_j(P)X^j=0$ implies $d_j=0$ for all $j$. We show by induction on
$k$ that the coefficient of $X^k$ in every $d_j$ vanishes, for all such families at once. As
$d_j(P)\equiv d_j(0)$ modulo $P$, hence modulo $X^n$, the identity gives
$\sum_{j<n}d_j(0)X^j\equiv0$ modulo $X^n$, so $d_j(0)=0$ for every $j$. So $d_j=Xd'_j$ and
$d_j(P)=Pd'_j(P)$; then $P\sum_jd'_j(P)X^j=0$, and as $w$ is a unit and $X^n$ is not a zero divisor,
$\sum_jd'_j(P)X^j=0$. The coefficient of $X^{k+1}$ in $d_j$ is that of $X^k$ in $d'_j$, which vanishes by
induction.

(2) First, $f(Xs)=f$ implies $f_j=0$ for $0<j<n$. Indeed, suppose $f_i=0$ for $0<i<j$. Modulo $X^{j+1}$,
$f(Xs)\equiv f_0+f_j(Xs)^j\equiv f_0+f_js(0)^jX^j$, so the coefficients of $X^j$ in $f(Xs)=f$ give
$f_j\big(s(0)^j-1\big)=0$, and $f_j=0$ as $s(0)^j-1$ is a unit. Hence $f=f_0+P\delta(f)$. Substituting
$Xs$, $f=f(Xs)=f_0+P(Xs)\,\delta(f)(Xs)=f_0+P\,\delta(f)(Xs)$, and cancelling $P$ gives
$\delta(f)(Xs)=\delta(f)$. By induction on $k$, every $\delta^k(f)$ is fixed by the substitution, so
$(\delta^kf)_j=0$ for $0<j<n$ and every $k$. In the decomposition of (1), $g_j=0$ for $0<j<n$, and
$f=g_0(P)$.
% Lean: general_basis_holds (BasisGeneral): divQuotient is delta, eq_sum_coeff_add_division,
% eq_sum_iterated_division, eq_sum_subst_unit_mul_X_pow (subst_eq_sum_range_add_pow_mul, X_pow_dvd_iff),
% and coeff_eq_zero_of_subst_sum_eq_zero (constantCoeff_eq_zero_of_sum_subst, cancel_unit_mul_X_pow);
% exists_subst_of_invariant: coeff_eq_zero_of_subst_invariant (coeff_subst_X_mul_of_low_zero) and
% averagedSeries_subst_invariant.
\end{proof}

\begin{lemma}\label{is-onto-step-exp-bound-link-depth-unit-series-invariants}
Let $l\ne2$ be a prime and $O$ a commutative ring that is a $\mathbf Z_l$-algebra, with the notation of
Definition~\ref{is-onto-step-exp-bound-link-depth-series}.
\begin{enumerate}
\item Every $f\in O[[X]]$ can be written uniquely as $f=\sum_{j<l-1}\iota(g_j)X^j$ with
  $g_j\in O[[X]]$. In particular $\iota$ is injective.
\item A series $f\in O[[X]]$ satisfies $\mathrm{act}_d(f)=f$ for every $d\in\mu_{l-1}$ if and only if
  $f=\iota(g)$ for some $g\in O[[X]]$.
\end{enumerate}
\end{lemma}
\leanref{Formalism/PieriBlockIdeals/CapelliCentralEnveloping.lean}{Formalism.CohomologicalDescent.HeightDescentMinors.PowerSeriesFreeBasis}
\leanref{Formalism/PieriBlockIdeals/PieriColumnMinors/HeightDescentMinors.lean}{Formalism.CohomologicalDescent.HeightDescentMinors.r2_basis_holds}
\leanref{Formalism/PieriBlockIdeals/CapelliCentralEnveloping.lean}{Formalism.CohomologicalDescent.HeightDescentMinors.RootsOfUnityInvariants}
\leanref{Formalism/PieriBlockIdeals/PieriColumnMinors/HeightDescentMinors.lean}{Formalism.CohomologicalDescent.HeightDescentMinors.fixed_relation_holds}
\begin{proof}
(1) By Lemma~\ref{is-onto-step-exp-bound-link-depth-unit-series-pi}(3), $\iota(g)=g(\pi_0)=g(wX^{l-1})$
with $w$ a unit, so this is Lemma~\ref{is-onto-step-exp-bound-link-depth-unit-series-division}(1) with
$n=l-1$. If $\iota(f)=\iota(g)$, the two decompositions $\iota(f)X^0+\sum_{0<j<l-1}\iota(0)X^j$ and
$\iota(g)X^0+\sum_{0<j<l-1}\iota(0)X^j$ of the same series agree, so $f=g$.

(2) Suppose that $f$ is fixed. Take $d$ as in
Lemma~\ref{is-onto-step-exp-bound-link-depth-unit-series-roots}(2), and write $\beta_d=Xs$ with
$s(0)=d$ (Lemma~\ref{is-onto-step-exp-bound-link-depth-unit-series-binomial}(4)). Then
$f(Xs)=\mathrm{act}_d(f)=f$, and $\pi_0(Xs)=\mathrm{act}_d(\pi_0)=\pi_0$ by
Lemma~\ref{is-onto-step-exp-bound-link-depth-unit-series-pi}(2). For $0<j<l-1$, $s(0)^j-1$ is the image
of the unit $d^j-1$ of $\mathbf Z_l$, so it is a unit of $O$.
Lemma~\ref{is-onto-step-exp-bound-link-depth-unit-series-division}(2) with $P=\pi_0=wX^{l-1}$ gives
$f=g(\pi_0)=\iota(g)$. Conversely, $\mathrm{act}_d\big(\iota(g)\big)=g\big(\pi_0(\beta_d)\big)
=g\big(\mathrm{act}_d(\pi_0)\big)=\iota(g)$ by
Lemma~\ref{is-onto-step-exp-bound-link-depth-unit-series-pi}(2).
% Lean: r2_basis_holds = r2_basis_of_r1 r1_holds general_basis_holds; the injectivity is the last clause
% of pi0Subst_laws_hold (r2_basis_holds at the index 0). fixed_relation_holds =
% fixed_relation_of_mu_and_inverse (exists_mu_generator_pow_sub_one_isUnit, br_eq_X_mul_const,
% exists_subst_of_invariant; subst_comp_subst_apply and r2_inverse_holds for the converse).
\end{proof}

\begin{lemma}\label{is-onto-step-exp-bound-link-depth-unit-series-q}
Let $l\ne2$ be a prime and $O$ a commutative ring that is a $\mathbf Z_l$-algebra, with the notation of
Definition~\ref{is-onto-step-exp-bound-link-depth-series}.
\begin{enumerate}
\item $q_0=gq$ for a unit $g$ of $O[[X]]$ with constant term $-1$.
\item Let $S\geq0$ be an integer and $m,y\in O[[X]]$ with $\iota(m)=q^Sy$. Then $m=(X-l)^Sm'$ for some
  $m'\in O[[X]]$.
\end{enumerate}
\end{lemma}
\leanref{Formalism/PieriBlockIdeals/CapelliCentralEnveloping.lean}{Formalism.CohomologicalDescent.HeightDescentMinors.UnitFactorQ}
\leanref{Formalism/PieriBlockIdeals/PieriColumnMinors/HeightDescentMinors.lean}{Formalism.CohomologicalDescent.HeightDescentMinors.r3_unit_holds}
\leanref{Formalism/PieriBlockIdeals/CapelliCentralEnveloping.lean}{Formalism.CohomologicalDescent.HeightDescentMinors.R4}
\leanref{Formalism/PieriBlockIdeals/PieriColumnMinors/HeightDescentMinors.lean}{Formalism.CohomologicalDescent.HeightDescentMinors.fourth_relation_holds}
\begin{proof}
(1) It suffices to prove this for $O=\mathbf Z_l$ and map the coefficients to $O$, as $\pi_0$, $q_0$ and
$q$ are the images of the same series over $\mathbf Z_l$. For $d\in\mu_{l-1}$ let $c_d\in\{1,\dots,l-1\}$
be the integer with $c_d\equiv d$ modulo $l$, and write $d=c_d+le_d$ with $e_d\in\mathbf Z_l$. By
Lemma~\ref{is-onto-step-exp-bound-link-depth-unit-series-binomial}(1) and~(5), $q$ divides
$\beta_d-\big((1+X)^{c_d}-1\big)$. As $d\mapsto c_d$ is a bijection of $\mu_{l-1}$ onto
$\{1,\dots,l-1\}$ (Lemma~\ref{is-onto-step-exp-bound-link-depth-unit-series-roots}(1)), and a difference
of two products whose factors are congruent modulo $q$ is divisible by $q$, $q$ divides
$\pi_0-\prod_{0<c<l}\big((1+X)^c-1\big)$.

Next, let $\Phi_l=\sum_{i<l}Y^i\in\mathbf Z[Y]$ be the $l$-th cyclotomic polynomial and
$\zeta\in\mathbf C$ a primitive $l$-th root of unity. Then $\Phi_l$ is the minimal polynomial of $\zeta$
over $\mathbf Q$ \cite[Chapter~2]{Was97}. As $l-1$ is even,
$\prod_{0<c<l}(\zeta^c-1)=\prod_{0<c<l}(1-\zeta^c)=\Phi_l(1)=l$. So $\zeta$ is a root of
$\prod_{0<c<l}(Y^c-1)-l$, and as $\Phi_l$ is monic with integer coefficients, it divides this polynomial
in $\mathbf Z[Y]$. Substitute $1+X$ for $Y$. As $X\Phi_l(1+X)=(1+X)^l-1=Xq$
(Lemma~\ref{is-onto-step-exp-bound-link-depth-unit-series-binomial}(3)), $\Phi_l(1+X)=q$, so $q$ divides
$\prod_{0<c<l}\big((1+X)^c-1\big)-l$. Adding, $q$ divides $\pi_0-l=q_0$. Write $q_0=gq$. The constant
terms are $q_0(0)=-l$ and $q(0)=\binom l1=l$, so $g(0)=-1$, and $g$ is a unit.

(2) By (1), $q=g^{-1}q_0$. By Lemma~\ref{is-onto-step-exp-bound-link-depth-unit-series-invariants}(1),
$g^{-S}y=\sum_{j<l-1}\iota(a_j)X^j$ for some $a_j$. As $\iota\big((X-l)^S\big)=q_0^S$
(Lemma~\ref{is-onto-step-exp-bound-link-depth-unit-series-pi}(1)),
\[
 \iota(m)=q_0^Sg^{-S}y=\sum_{j<l-1}\iota\big((X-l)^Sa_j\big)X^j.
\]
Also $\iota(m)=\iota(m)X^0+\sum_{0<j<l-1}\iota(0)X^j$. By the uniqueness in
Lemma~\ref{is-onto-step-exp-bound-link-depth-unit-series-invariants}(1), $m=(X-l)^Sa_0$.
% Lean: r3_unit_holds = r3_unit_of_laws: q_dvd_q_zero_of_ne_two (q_dvd_binomialSeriesSubOne_sub_pow,
% padicInt_eq_digit_add_mul, prod_mu_toZMod_val_eq_prod_Ioo, dvd_prod_sub_prod_of_dvd_sub, and
% q_dvd_prod_sub_l from cyclotomic_dvd_prod_X_pow_sub_one_sub: minpoly.isIntegrallyClosed_dvd and
% IsPrimitiveRoot.prod_pow_sub_one_eq_order), then the constant terms and Units.map to O.
% fourth_relation_holds = fourth_relation_of_unit_and_basis r3_unit_holds r2_basis_holds.
\end{proof}

\begin{lemma}\label{is-onto-step-exp-bound-link-depth-unit-series}
Let $l\ne2$ be a prime, $O$ a domain that is a $\mathbf Z_l$-algebra, and $u,c\in\mathbf Z_l$ with
$u=1+lc$. Then there is a unit series for $u$ in $O[[X]]$
(Definition~\ref{is-onto-step-exp-bound-link-depth-series}(3)); it is unique; and its constant term is
$c+le$ for some $e\in O$.
\end{lemma}
\leanref{Formalism/PieriBlockIdeals/LemmaWFromLattice.lean}{Formalism.CohomologicalDescent.FrobeniusPeriodChange.UnConst}
\leanref{Formalism/PieriBlockIdeals/PieriColumnMinors/HeckeLFunctionIdealEulerProduct.lean}{Formalism.CohomologicalDescent.FrobeniusPeriodChange.unit_series_constant_holds}
\begin{proof}
Existence. First let $O=\mathbf Z_l$. For $d\in\mu_{l-1}$ write $\beta_d=X\beta'_d$ with $\beta'_d(0)=d$
(Lemma~\ref{is-onto-step-exp-bound-link-depth-unit-series-binomial}(4)), and put
$B=\prod_{d\in\mu_{l-1}}\beta'_d$ and $B'=\prod_{d\in\mu_{l-1}}\beta'_{du}$. As $\mu_{l-1}$ has $l-1$
elements, $\pi_0=X^{l-1}B$; and by
Lemma~\ref{is-onto-step-exp-bound-link-depth-unit-series-binomial}(1) and~(2),
$\mathrm{act}_u(\pi_0)=\prod_d\mathrm{act}_u(\beta_d)=\prod_d\beta_{du}=X^{l-1}B'$. The constant terms are
$B(0)=\prod_dd=-1$ and $B'(0)=\prod_ddu=-u^{l-1}$
(Lemma~\ref{is-onto-step-exp-bound-link-depth-unit-series-roots}(3)). So $B$ is a unit; put
$r=B'B^{-1}$. Then $\mathrm{act}_u(\pi_0)=\pi_0r$.

The series $q$ divides $r-1$. Indeed, for $d\in\mu_{l-1}$, $du=d+l(dc)$, so
$\beta_{du}=\beta_d+Xqs$ for some $s$ by
Lemma~\ref{is-onto-step-exp-bound-link-depth-unit-series-binomial}(5); cancelling $X$,
$\beta'_{du}-\beta'_d=qs$. So $q$ divides $B'-B$, and $r-1=(B'-B)B^{-1}$.

The series $r$ is fixed by $\mathrm{act}_d$ for every $d\in\mu_{l-1}$. Indeed, by
Lemma~\ref{is-onto-step-exp-bound-link-depth-unit-series-binomial}(2) and
Lemma~\ref{is-onto-step-exp-bound-link-depth-unit-series-pi}(2),
$\mathrm{act}_d\big(\mathrm{act}_u(\pi_0)\big)=\mathrm{act}_{ud}(\pi_0)
=\mathrm{act}_u\big(\mathrm{act}_d(\pi_0)\big)=\mathrm{act}_u(\pi_0)$, while
$\mathrm{act}_d(\pi_0r)=\pi_0\,\mathrm{act}_d(r)$. As $\pi_0=X^{l-1}B\ne0$ and $\mathbf Z_l[[X]]$ is a
domain, $\mathrm{act}_d(r)=r$.

By Lemma~\ref{is-onto-step-exp-bound-link-depth-unit-series-invariants}(2), $r=\iota(\rho)$ for some
$\rho$. Then $\iota(\rho-1)=r-1=qy$ for some $y$, and
Lemma~\ref{is-onto-step-exp-bound-link-depth-unit-series-q}(2) with $S=1$ gives
$\rho-1=(X-l)u^\flat$ for some $u^\flat$. So
\[
 r=\iota\big(1+(X-l)u^\flat\big)=1+\iota(u^\flat)q_0
\]
by Lemma~\ref{is-onto-step-exp-bound-link-depth-unit-series-pi}(1), and
$\mathrm{act}_u(\pi_0)=\pi_0\big(1+\iota(u^\flat)q_0\big)$: $u^\flat$ is a unit series for $u$.

Its constant term. The constant term of $r$ is $B'(0)B(0)^{-1}=u^{l-1}$, and that of
$1+\iota(u^\flat)q_0$ is $1-l\,u^\flat(0)$, as $q_0(0)=-l$. So $u^{l-1}=1-l\,u^\flat(0)$. As
$(1+lc)^{l-1}=1+(l-1)lc+(lc)^2e'$ for some $e'\in\mathbf Z_l$,
$l\,u^\flat(0)=l\big(c+l(-c-c^2e')\big)$, and as $\mathbf Z_l$ is a domain,
$u^\flat(0)=c+l(-c-c^2e')$.

For a general $O$, map the coefficients of $u^\flat$ to $O$. The series $\beta_u$, $\pi_0$ and $q_0$ over
$O$ are the images of those over $\mathbf Z_l$, and the map commutes with substitution, so the image of
$u^\flat$ is a unit series for $u$ in $O[[X]]$, with constant term $c+le$, where $e$ is the image of
$-c-c^2e'$.

Uniqueness. Let $u^\flat_1$ and $u^\flat_2$ be unit series for $u$ in $O[[X]]$. Then
$\pi_0\big(1+\iota(u^\flat_1)q_0\big)=\pi_0\big(1+\iota(u^\flat_2)q_0\big)$. Here $\pi_0=\iota(X)$ and
$q_0=\iota(X-l)$ are nonzero, as $\iota$ is injective
(Lemma~\ref{is-onto-step-exp-bound-link-depth-unit-series-invariants}(1)), and $O[[X]]$ is a domain, as
$O$ is. So $\iota(u^\flat_1)=\iota(u^\flat_2)$, and $u^\flat_1=u^\flat_2$. Hence every unit series for $u$
is the one constructed above, and its constant term is $c+le$.
% Lean: unit_series_constant_holds (UnConst): r6_holds = sixth_relation_of_action_and_mu_laws, whose
% first part act_unit_pi0_eq_mul_padic is the case O = Z_l (B, B', r, hQ, hinv, the fixed relation, R4
% with S = 1, and one_add_pow_eq_linear_add_sq_mul for the constant term), mapped to O by map_act_padic,
% map_algebraMap_pi0Subst, map_q0_padic_alt and map_pi0_padicInt_algebraMap; uniqueness from
% pi0Subst_laws_hold (inc X = pi0, inc (X - l) = q0, inc injective), mul_left_cancel and mul_right_cancel.
% The period change uses only the constant term, in theta_diagonal_entry_eq.
\end{proof}

\begin{lemma}\label{is-onto-step-exp-bound-link-depth-alpha}
Let $l\ne2$ be a prime, $O$ a domain that is a $\mathbf Z_l$-algebra, $F$ a field of characteristic $l$,
$E$ a finite set with a permutation $\sigma$, $t,n,h\geq0$ integers, $u\in\mathbf Z_l^\times$,
$c\in\mathbf Z_l$, integer sequences $U'_i$ of length $n$ and integers $\mathrm{wt}_i(s)$ and $\nu_i(s)$
for $i\in E$ and $s<t$. Let $\mathcal W$ be Witt lattice data for $(E,t,n,O,F)$ with the weight inputs for
$(\sigma,l,h,u,c,U',\mathrm{wt},\nu)$ (Definition~\ref{is-onto-step-exp-bound-link-depth-data}(5)). For
$i\in E$ and $s<t$ put
$e=l\nu_{\sigma(i)}(s)+\mathrm{wt}_i(s)-\nu_i(s)$. Then $0\le e\le h$, and
$\Phi_i(\mathrm{pos}(s),\mathrm{pos}(s))$ is a unit of $F[[X]]$ times $X^e$. That $l\ne2$ and that $O$ is
a domain are not used in the proof.
\end{lemma}
\leanref{Formalism/PieriBlockIdeals/LemmaWFromLattice.lean}{Formalism.CohomologicalDescent.FrobeniusPeriodChange.AlphaBounds}
\leanref{Formalism/PieriBlockIdeals/PieriColumnMinors/HeckeLFunctionIdealEulerProduct.lean}{Formalism.CohomologicalDescent.FrobeniusPeriodChange.alpha_bounds_hold}
\leanref{Formalism/PieriBlockIdeals/PieriColumnMinors/HeckeLFunctionIdealEulerProduct.lean}{Formalism.CohomologicalDescent.FrobeniusPeriodChange.frobenius_diagonal_entry_eq_unit_mul}
\begin{proof}
Let $p=\mathrm{pos}(s)$. As $N^\psi_s$ has height $h$ with weights $(\mathrm{wt}_{i'}(s))_{i'}$, the frame
condition (a) and Lemma~\ref{is-onto-step-exp-bound-link-depth-rank-one} give
$P^\psi_{s,i}=v(X-l)^{\mathrm{wt}_i(s)}$ for a unit $v$, as $\mathrm{wt}_i(s)\geq0$ by (b). By
Lemma~\ref{is-onto-step-exp-bound-link-depth-reduction}(3), $\bar P^\psi_{s,i}=\bar vX^{\mathrm{wt}_i(s)}$
with $\bar v$ a unit. The frame condition (e) and
Lemma~\ref{is-onto-step-exp-bound-link-depth-frobenius-entry} then give $e\geq0$ and
$\Phi_i(p,p)=u_0X^e$ with $u_0$ a unit.

By (a) of the height of $N$, $P_iM=(X-l)^h$ for some $M$. Reducing, $\bar P_i\bar M=X^h$, and
Lemma~\ref{is-onto-step-exp-bound-link-depth-base-change}(2) gives $M'$ with $\Phi_iM'=X^h$. As
$\Phi_i$ is block upper triangular for $b$ by (d), and $p$ is alone in its block by (c),
Lemma~\ref{is-onto-step-exp-bound-link-depth-matrices} gives $\Phi_i(p,p)M'(p,p)=X^h$. So $X^e$ divides
$X^h$, and $e\le h$.
% Lean: alpha_bounds_hold (AlphaBounds): frobenius_diagonal_entry_eq_unit_mul (rank_one_model_holds,
% map_X_sub_C_prime_pow_eq_X_pow, frobenius_diagonal_holds), frobenius_mul_eq_X_pow_smul_one (with
% frobenius_base_change_holds) and block_height_on_line_holds; hl2 is not used.
\end{proof}

\begin{lemma}\label{is-onto-step-exp-bound-link-depth-properties}
Let $l\ne2$ be a prime, $O$ a domain that is a $\mathbf Z_l$-algebra, $F$ a field of characteristic $l$,
$E$ a finite set with a permutation $\sigma$, $t,n,h\geq0$ integers, $u\in\mathbf Z_l^\times$,
$c\in\mathbf Z_l$, and for $i\in E$ and $s<t$ integer sequences $U'_i$ of length $n$ and integers
$\mathrm{bot}_i$ and $\nu_i(s)$. Let $\mathcal W$ be Witt lattice data for $(E,t,n,O,F)$ with the weight
inputs for $(\sigma,l,h,u,c,U',\mathrm{wt},\nu)$, where $\mathrm{wt}_i(s)=\mathrm{bot}_i+s$
(Definition~\ref{is-onto-step-exp-bound-link-depth-data}(5)). Let $m_i\in\mathbf N$ with
$m_i\le\sum_{j\in J}U'_i(j)$ for every $i$ and every set $J$ of $t$ indices, and let $\Theta_i$ be the
coefficient of $X$ in $\Gamma_i$. Then the weight datum $(\mathrm{red}(c),\Theta,\Phi,b,\mathrm{pos})$
for $(E,t,n,F)$ has the weight properties for $(\sigma,l,U',\mathrm{bot},\nu,m)$
(Definition~\ref{is-onto-step-exp-bound-link-depth-datum}).
\end{lemma}
\leanref{Formalism/PieriBlockIdeals/LemmaWFromLattice.lean}{Formalism.CohomologicalDescent.FrobeniusPeriodChange.WPropsOfInputs}
\leanref{Formalism/PieriBlockIdeals/PieriColumnMinors/HeckeLFunctionIdealEulerProduct.lean}{Formalism.CohomologicalDescent.FrobeniusPeriodChange.weight_properties_of_inputs}
\begin{proof}
Write $\mathrm a=\mathrm{red}(c)$, which is nonzero as $c$ is a unit by the frame condition (a). The
conditions on pos and $b$ are (c). As $\Gamma_i$ is block upper triangular for $b$ by (d), so is
$\Theta_i$.

By Lemma~\ref{is-onto-step-exp-bound-link-depth-base-change}(1),
$\det(Y-\Theta_i)=\det(Y-\bar D_i(0))$. By (d) of the height of $N$ and
Lemma~\ref{is-onto-step-exp-bound-link-depth-reduction}(2), $\bar D_i(0)$ is the reduction of the
matrix $\Theta^N_i\in M_n(O)$ of the lattice datum $N$; by (e), $\det(Y-\Theta^N_i)=\prod_j(Y-r_j)$ with
$u^{-U'_i(j)}=1+lr_j$, and $\mathrm{red}(r_j)=-\mathrm aU'_i(j)$ by
Lemma~\ref{is-onto-step-exp-bound-link-depth-reduction}(1). So
$\det(Y-\Theta_i)=\prod_j(Y+\mathrm aU'_i(j))$.

Let $s<t$ and $p=\mathrm{pos}(s)$. The frame condition (f) and
Lemma~\ref{is-onto-step-exp-bound-link-depth-theta-entry} give
$\Theta_i(p,p)=\nu_i(s)\bar u(0)+\bar D^\psi_{s,i}(0)$. Here
$\bar D^\psi_{s,i}(0)=-\mathrm a\,\mathrm{wt}_i(s)$ by
Lemma~\ref{is-onto-step-exp-bound-link-depth-rank-one}, and $\bar u(0)=\mathrm a$, as $u^\flat$ is a unit
series for $u=1+lc$ and its constant term is $c+le$ by
Lemma~\ref{is-onto-step-exp-bound-link-depth-unit-series}. So
$\Theta_i(p,p)=-\mathrm a(\mathrm{bot}_i+s-\nu_i(s))$. The matrix $\Phi_i$ is block upper triangular by
(d), and $\Phi_i(p,p)$ is a unit times $X^{\alpha_i(s)}$ with $\alpha_i(s)\geq0$ by
Lemma~\ref{is-onto-step-exp-bound-link-depth-alpha}.

Finally, by (b) of the height of $N$, with $r=t$ and $S=m_i$, $(X-l)^{m_i}$ divides the $t$-minors of
$P_i$. Reduction of coefficients maps minors to minors, so $X^{m_i}$ divides the $t$-minors of $\bar P_i$
by Lemma~\ref{is-onto-step-exp-bound-link-depth-reduction}(3), and those of $\Phi_i$ by
Lemma~\ref{is-onto-step-exp-bound-link-depth-base-change}(2). The determinant of
$(\Phi_i(\mathrm{pos}(s),\mathrm{pos}(s')))$ is $0$ or, up to sign, a $t$-minor of $\Phi_i$, so $X^{m_i}$
divides it.
% Lean: weight_properties_of_inputs (WPropsOfInputs): reduction_algebraMap_unit_ne_zero; the Theta block
% triangularity from hGamma; charpoly_theta_eq_prod (theta_base_change_holds, evaluation_congruence_holds,
% reduction_log_unit_pow_eq); theta_diagonal_entry_eq (rank_one_model_holds, theta_diagonal_holds and the
% constant term from unit_series_constant_holds); frobenius_diagonal_entry_eq_unit_mul; and
% minorsDvd_frobenius_of_sum_bound (minors_divide_all_holds, minor_reduction_holds,
% frobenius_base_change_holds).
\end{proof}

\begin{lemma}\label{is-onto-step-exp-bound-link-depth-inputs}
Let $k,j,t,q,h\geq0$ be integers and $l\ne2$ a prime with $t+2\kappa<l$, where
$\kappa=\lfloor h/(l-1)\rfloor$, and let $N=k+j+t+q$. Let $O$ be a domain that is a $\mathbf Z_l$-algebra,
$F$ a field of characteristic $l$, $E$ a finite set with a permutation $\sigma$, $u\in\mathbf Z_l^\times$,
$c\in\mathbf Z_l$, and $\mathrm{bot}_i,z_i,\nu_i(s)\in\mathbf Z$ and $U_i,U'_i\in\mathbf Z^N$ for $i\in E$
and $s<t$. Let $\mathcal W$ be Witt lattice data for $(E,t,N,O,F)$ with the weight inputs for
$(\sigma,l,h,u,c,U',\mathrm{wt},\nu)$, where $\mathrm{wt}_i(s)=\mathrm{bot}_i+s$
(Definition~\ref{is-onto-step-exp-bound-link-depth-data}(5)). Suppose
that for every $i$, $U_i$ and $U'_i$ decrease strictly, $U_i$ takes the values $\mathrm{bot}_i+s$ for
$0\le s<t$, and $\prod_a(Y+U_i(a)+z_i)=\prod_a(Y+U'_i(a)+z_i)$ in $F[Y]$; and let $\tau\in E$ with
$U_\tau(a)\geq0$ for every $a$, $\sum_aU_\tau(a)=\sum_aU'_\tau(a)$, and
$U_\tau(a)=\mathrm{bot}_\tau+(k+j+t-1-a)$ for $k+j\le a<k+j+t$. Let
$n_i=n_{l,\kappa}(\mathrm{bot}_i,\{U_i(a):a<N\})$
(Definition~\ref{is-onto-step-exp-bound-link-depth-reflected}(4)). Then
\[
 \sum_{a<k+j}\big(U_\tau(a)-U'_\tau(a)\big)\le qh+\big(l\,n_{\sigma(\tau)}+n_\tau\big)(t+2\kappa).
\]
\end{lemma}
\leanref{Formalism/PieriBlockIdeals/LemmaWFromLattice.lean}{Formalism.CohomologicalDescent.FrobeniusPeriodChange.WEFromInputs}
\leanref{Formalism/PieriBlockIdeals/PieriColumnMinors/HeckeLFunctionIdealEulerProduct.lean}{Formalism.CohomologicalDescent.FrobeniusPeriodChange.weight_exponents_from_inputs}
\begin{proof}
By Lemma~\ref{is-onto-step-exp-bound-link-depth-alpha}, $0\le\alpha_i(s)\le h$, with
$\alpha_i(s)=l\nu_{\sigma(i)}(s)+\mathrm{bot}_i+s-\nu_i(s)$; the frame condition (b) gives
$0\le U'_i(j)\le h$ and $0\le\mathrm{bot}_i+s\le h$. Put $m_\tau=\sum_{k+j+q\le a<N}U'_\tau(a)$, which is
at least $0$, and $m_i=0$ for $i\ne\tau$. Then $m_i\le\sum_{j'\in J}U'_i(j')$ for every $i$ and every set
$J$ of $t$ indices. For $i\ne\tau$ this holds as $U'_i\geq0$. For $i=\tau$, the set $T$ of the indices
$a\geq k+j+q$ has $t$ elements, and $U'_\tau$ is decreasing: the elements of $T$ not in $J$ have values
at most $U'_\tau(k+j+q)$, the elements of $J$ not in $T$ have values at least $U'_\tau(k+j+q)$, and there
are as many of the one as of the other; so the sum over $T$ is at most the sum over $J$.

By Lemma~\ref{is-onto-step-exp-bound-link-depth-properties}, the weight datum
$(\mathrm{red}(c),\Theta,\Phi,b,\mathrm{pos})$ has the weight properties for
$(\sigma,l,U',\mathrm{bot},\nu,m)$. All the conditions of
Lemma~\ref{is-onto-step-exp-bound-link-depth-exchange} now hold, and it gives the bound.
% Lean: weight_exponents_from_inputs (WEFromInputs): alpha_bounds_hold, the fields hW.1.2.1 and
% hW.1.2.2.1 of the frame, the m_i as cc (card_filter_le_val_eq_sub and sum_tail_le_sum_of_antitone),
% weight_properties_of_inputs, and weightExchangeFull_holds for W.toWData l c sigma.
\end{proof}

\begin{lemma}\label{is-onto-step-exp-bound-link-depth-norm-division}
Let $O$ be a domain that is complete for the $I$-adic topology of an ideal $I$, and $c\in I$.
\begin{enumerate}
\item Every $f\in O[[X]]$ can be written uniquely as $f=(X-c)f_1+a$ with $f_1\in O[[X]]$ and $a\in O$.
\item If $X-c$ divides a constant $b\in O$ in $O[[X]]$, then $b=0$. The series $X-c$ is a prime element of
  $O[[X]]$.
\item Let $s$ be a finite set, $c_b\in I$ for $b\in s$ pairwise distinct, $S\geq0$ an integer and
  $m\in O[[X]]$ with $(X-c_b)^S\mid m$ for every $b\in s$. Then $\big(\prod_{b\in s}(X-c_b)\big)^S\mid m$.
\end{enumerate}
\end{lemma}
\leanref{Formalism/PieriBlockIdeals/CapelliCentralEnveloping.lean}{Formalism.CohomologicalDescent.HeightDescentMinors.LinDivision}
\leanref{Formalism/PieriBlockIdeals/PieriColumnMinors/HeightDescentMinors.lean}{Formalism.CohomologicalDescent.HeightDescentMinors.linear_division_holds}
\leanref{Formalism/PieriBlockIdeals/CapelliCentralEnveloping.lean}{Formalism.CohomologicalDescent.HeightDescentMinors.LinPrime}
\leanref{Formalism/PieriBlockIdeals/PieriColumnMinors/HeightDescentMinors.lean}{Formalism.CohomologicalDescent.HeightDescentMinors.linear_prime_holds}
\leanref{Formalism/PieriBlockIdeals/CapelliCentralEnveloping.lean}{Formalism.CohomologicalDescent.HeightDescentMinors.RootsProduct}
\leanref{Formalism/PieriBlockIdeals/PieriColumnMinors/HeightDescentMinors.lean}{Formalism.CohomologicalDescent.HeightDescentMinors.roots_product_holds}
\begin{proof}
(1) As $c\in I$, the image of $X-c$ in $(O/I)[[X]]$ is $X$, of order $1$, so $X-c$ is a distinguished
polynomial of degree $1$, and this is Weierstrass division by it \cite[Proposition~7.2]{Was97}.

(2) If $(X-c)y=b$, then $(X-c)y+0=(X-c)0+b$, and uniqueness in (1) gives $b=0$. So $X-c$ does not divide
$1$ and is not a unit; and $X-c\ne0$, as its coefficient of $X$ is $1$. Let $X-c$ divide $fg$. Write
$f=(X-c)f_1+a$ and $g=(X-c)g_1+b$ by (1). Then $ab=fg-(X-c)(f_1g+ag_1)$ is divisible by $X-c$, so
$ab=0$. As $O$ is a domain, $a=0$ or $b=0$, and $X-c$ divides $f$ or $g$.

(3) By induction on $s$. Suppose $m=\big(\prod_{b\in s}(X-c_b)\big)^Sm'$, and let $a\notin s$ with
$c_a\in I$ distinct from the $c_b$, and $(X-c_a)^S\mid m$. The prime $X-c_a$ does not divide
$\big(\prod_{b\in s}(X-c_b)\big)^S$: otherwise it divides some $X-c_b$, hence the constant
$c_a-c_b=(X-c_b)-(X-c_a)$, and $c_a=c_b$ by (2). So $(X-c_a)^S$ divides $m'$, and the claim holds for
$s\cup\{a\}$.
% Lean: linear_division_holds (LinDivision) is Mathlib's Weierstrass division
% (IsWeierstrassDivisorAt.isWeierstrassDivisionAt_div_mod and eq_of_mul_add_eq_mul_add) for
% X_sub_C_isDistinguished; eq_zero_of_X_sub_C_dvd_C and linear_prime_holds (LinPrime); roots_product_holds
% (RootsProduct), with Prime.pow_dvd_of_dvd_mul_left.
\end{proof}

\begin{lemma}\label{is-onto-step-exp-bound-link-depth-norm-descent}
Let $l\ne2$ be a prime and $O$ a commutative ring that is a $\mathbf Z_l$-algebra, with the notation of
Definition~\ref{is-onto-step-exp-bound-link-depth-series}.
\begin{enumerate}
\item Let $P,Q\in M_n(O[[X]])$ and $h\geq0$ with $\iota(P)Q=q^h$. Then $PQ'=(X-l)^h$ for some
  $Q'\in M_n(O[[X]])$.
\item Suppose that $O$ is a domain and $\zeta\in O$ a primitive $l$-th root of unity. Then
  $q=\prod_{0<b<l}\big(X-(\zeta^b-1)\big)$.
\item Suppose moreover that $O$ is complete for the $I$-adic topology of an ideal $I$ with $\zeta-1\in I$,
  and let $S\geq0$ and $m\in O[[X]]$ with $\big(X-(\zeta^b-1)\big)^S\mid\iota(m)$ for every $0<b<l$.
  Then $(X-l)^S\mid m$.
\item Let $O$ be a domain that is complete for the $I$-adic topology of an ideal $I$, $x\in I$, $L$ a
  field, $\mathrm{em}\colon O\to L$ an injective ring homomorphism and $\jmath\colon O[[X]]\to L[[X]]$ a
  ring homomorphism that maps every constant $a$ to $\mathrm{em}(a)$, with $\jmath(X-x)$ equal to $X$
  times a unit. If $S\geq0$, $m\in O[[X]]$ and $X^S\mid\jmath(m)$, then $(X-x)^S\mid m$.
\end{enumerate}
\end{lemma}
\leanref{Formalism/PieriBlockIdeals/CapelliCentralEnveloping.lean}{Formalism.CohomologicalDescent.HeightDescentMinors.HeightDescent}
\leanref{Formalism/PieriBlockIdeals/PieriColumnMinors/HeightDescentMinors.lean}{Formalism.CohomologicalDescent.HeightDescentMinors.height_descent_holds}
\leanref{Formalism/PieriBlockIdeals/CapelliCentralEnveloping.lean}{Formalism.CohomologicalDescent.HeightDescentMinors.QRoots}
\leanref{Formalism/PieriBlockIdeals/PieriColumnMinors/HeightDescentMinors.lean}{Formalism.CohomologicalDescent.HeightDescentMinors.q_roots_holds}
\leanref{Formalism/PieriBlockIdeals/CapelliCentralEnveloping.lean}{Formalism.CohomologicalDescent.HeightDescentMinors.RootsDescent}
\leanref{Formalism/PieriBlockIdeals/PieriColumnMinors/HeightDescentMinors.lean}{Formalism.CohomologicalDescent.HeightDescentMinors.roots_descent_holds}
\leanref{Formalism/PieriBlockIdeals/LemmaWFromLattice.lean}{Formalism.CohomologicalDescent.FrobeniusPeriodChange.OrderDescent}
\leanref{Formalism/PieriBlockIdeals/PieriColumnMinors/HeckeLFunctionIdealEulerProduct.lean}{Formalism.CohomologicalDescent.FrobeniusPeriodChange.order_descent_holds}
\begin{proof}
(1) By Lemma~\ref{is-onto-step-exp-bound-link-depth-unit-series-q}(1), $q_0=gq$ with $g$ a unit. By
Lemma~\ref{is-onto-step-exp-bound-link-depth-unit-series-invariants}(1), write
$g^hQ(a,b)=\sum_{j<l-1}\iota\big(G_{ab}(j)\big)X^j$ for every entry. Then for all $a,b$,
\[
 q_0^h\,\delta_{ab}=g^h\big(\iota(P)Q\big)(a,b)=\sum_c\iota\big(P(a,c)\big)\sum_{j<l-1}\iota\big(G_{cb}(j)\big)X^j
 =\sum_{j<l-1}\iota\Big(\sum_cP(a,c)G_{cb}(j)\Big)X^j,
\]
while $q_0^h\delta_{ab}=\iota\big((X-l)^h\delta_{ab}\big)X^0$ by
Lemma~\ref{is-onto-step-exp-bound-link-depth-unit-series-pi}(1). By the uniqueness in
Lemma~\ref{is-onto-step-exp-bound-link-depth-unit-series-invariants}(1),
$\sum_cP(a,c)G_{cb}(0)=(X-l)^h\delta_{ab}$. So $Q'=\big(G_{ab}(0)\big)_{a,b}$ works.

(2) As $O$ is a domain and $\zeta$ is a primitive $l$-th root of unity, $Y^l-1=\prod_{b<l}(Y-\zeta^b)$ in
$O[Y]$. Substituting $1+X$ for $Y$ gives $(1+X)^l-1=X\prod_{0<b<l}\big(X-(\zeta^b-1)\big)$. As
$(1+X)^l-1=Xq$ (Lemma~\ref{is-onto-step-exp-bound-link-depth-unit-series-binomial}(3)), cancelling $X$
gives the claim.

(3) For $0<b<l$, $\zeta^b-1=(\zeta-1)(1+\zeta+\dots+\zeta^{b-1})$ lies in $I$, and these elements are
pairwise distinct as $\zeta$ is primitive. By
Lemma~\ref{is-onto-step-exp-bound-link-depth-norm-division}(3) and (2), $q^S$ divides $\iota(m)$, and
Lemma~\ref{is-onto-step-exp-bound-link-depth-unit-series-q}(2) gives $(X-l)^S\mid m$.

(4) By induction on $S$; the case $S=0$ is clear. Let $X^{S+1}\mid\jmath(m)$, and write $\jmath(X-x)=Xv$
with $v$ a unit. By Lemma~\ref{is-onto-step-exp-bound-link-depth-norm-division}(1),
$m=(X-x)f_1+a$. Then $X$ divides $\jmath(m)=Xv\jmath(f_1)+\mathrm{em}(a)$, hence divides $\mathrm{em}(a)$,
so $\mathrm{em}(a)=0$ and $a=0$, as em is injective. Now $X^{S+1}$ divides $Xv\jmath(f_1)$; as $L[[X]]$ is a
domain and $v$ a unit, $X^S\mid\jmath(f_1)$. By induction $(X-x)^S\mid f_1$, so $(X-x)^{S+1}\mid m$.
% Lean: height_descent_holds = height_descent_of_r3_r2 r3_unit_holds r2_basis_holds (key, with the
% uniqueness of r2_basis_holds at q0^h * 1 a b); q_roots_holds (QRoots: X_pow_sub_one_eq_prod,
% X_mul_q_eq_one_add_X_pow_sub_one); roots_descent_holds = rootsDescent_of_roots_and_fourth_relation
% (roots_product_holds, q_roots_holds, fourth_relation_holds with S); order_descent_holds (OrderDescent),
% by induction on S with linear_division_holds.
\end{proof}

\begin{lemma}\label{is-onto-step-exp-bound-link-depth-norm-local}
Let $l\ne2$ be a prime, $O$ a domain that is a $\mathbf Z_l$-algebra, $I$ an ideal of $O$, $\zeta\in O$ a
primitive $l$-th root of unity with $\zeta-1\in I$, $L$ a field, $\mathrm{em}\colon O\to L$ an injective
ring homomorphism, $u\in\mathbf Z_l$, $0<b<l$ an integer, and $\jmath\colon O[[X]]\to L[[X]]$ a local
exponential at $\zeta^b-1$ for $u$ (Definition~\ref{is-onto-step-exp-bound-link-depth-data}(2)). Then
$l\in I$, $l\ne0$ in $O$, and $\jmath(q_0)$ is $X$ times a unit. If $y\in O[[X]]$ and $a\in O$ with
$X-l\mid y-a$, then $\jmath\big(\iota(y)\big)$ has constant term $\mathrm{em}(a)$.
\end{lemma}
\leanref{Formalism/PieriBlockIdeals/LemmaWFromLattice.lean}{Formalism.CohomologicalDescent.FrobeniusPeriodChange.LocFacts}
\leanref{Formalism/PieriBlockIdeals/PieriColumnMinors/HeckeLFunctionIdealEulerProduct.lean}{Formalism.CohomologicalDescent.FrobeniusPeriodChange.localization_facts_hold}
\leanref{Formalism/PieriBlockIdeals/CapelliCentralEnveloping.lean}{Formalism.CohomologicalDescent.HeightDescentMinors.RootOfUnityRemainder}
\leanref{Formalism/PieriBlockIdeals/PieriColumnMinors/HeightDescentMinors.lean}{Formalism.CohomologicalDescent.HeightDescentMinors.r3_divisibility_holds}
\begin{proof}
For $0<b'<l$, $\zeta^{b'}-1=(\zeta-1)(1+\zeta+\dots+\zeta^{b'-1})\in I$. By
Lemma~\ref{is-onto-step-exp-bound-link-depth-unit-series-q}(1) and
Lemma~\ref{is-onto-step-exp-bound-link-depth-norm-descent}(2),
$q_0=g\prod_{0<b'<l}\big(X-(\zeta^{b'}-1)\big)$ with $g$ a unit. So $q_0=\big(X-(\zeta^b-1)\big)s$ for some
$s$, and comparing constant terms, $-l=-(\zeta^b-1)s(0)$; so $l\in I$. As $O$ is a domain with a primitive
$l$-th root of unity, $l\ne0$ in $O$: in characteristic $l$, $Y^l-1=(Y-1)^l$, and $1$ would be the only
$l$-th root of unity.

Write $\jmath(X-(\zeta^b-1))=Xv$ with $v$ a unit. For $b'\ne b$, $0<b'<l$,
\[
 \jmath\big(X-(\zeta^{b'}-1)\big)=Xv+\mathrm{em}(\zeta^b-\zeta^{b'})
\]
has constant term $\mathrm{em}(\zeta^b-\zeta^{b'})\ne0$, as $\zeta^b\ne\zeta^{b'}$ and em is injective; so
it is a unit of $L[[X]]$. Hence $\jmath(q_0)=X\cdot v\,\jmath(g)\prod_{b'\ne b}\jmath\big(X-(\zeta^{b'}-1)\big)$
is $X$ times a unit.

Finally, write $y=(X-l)s'+a$. Then $\iota(y)=q_0\iota(s')+a$ by
Lemma~\ref{is-onto-step-exp-bound-link-depth-unit-series-pi}(1), and
$\jmath\big(\iota(y)\big)=\jmath(q_0)\jmath\big(\iota(s')\big)+\mathrm{em}(a)$, whose constant term is
$\mathrm{em}(a)$ as $X$ divides $\jmath(q_0)$.
% Lean: localization_facts_hold (LocFacts): hbI (mul_geom_sum); hlI from r3_divisibility_holds
% (RootOfUnityRemainder = r3_divisibility_of_unit_roots r3_unit_holds q_roots_holds); hl0 from
% IsPrimitiveRoot.neZero'; hq0 (r3_unit_holds, q_roots_holds, Finset.mul_prod_erase); and the last
% clause through inc_X_sub_l_eq_q0 and pi0Subst_C.
\end{proof}

\begin{lemma}\label{is-onto-step-exp-bound-link-depth-norm-matrices}
\begin{enumerate}
\item Let $R$ be a commutative ring, $x\in R$, $n,t,S\geq0$ integers, $a_j\geq0$ integers for $j<n$, and
  $G,H\in M_n(R)$. If $S\le\sum_{j\in J}a_j$ for every set $J$ of $t$ indices, then $x^S$ divides the
  $t$-minors of $G\,\mathrm{diag}(x^{a_j})\,H$.
\item Let $L$ be a field, $\Gamma,G\in M_n(L[[X]])$, $\upsilon\in L$ and $a_j\geq0$ integers for $j<n$, and
  write $\Delta=\mathrm{diag}(X^{a_j})$. If $\Gamma\cdot(G\Delta)(\upsilon X)=G\Delta$, then
  $\Gamma(0)G(0)\,\mathrm{diag}(\upsilon^{a_j})=G(0)$.
\item Let $O$ be a commutative ring, $L$ a field, $\mathrm{em}\colon O\to L$ an injective ring
  homomorphism, $\Theta\in M_n(O)$, $\lambda\in O$ nonzero, $r_j\in O$ for $j<n$ and $G\in GL_n(L)$ with
  $\mathrm{em}(1+\lambda\Theta)\,G=G\,\mathrm{diag}\big(\mathrm{em}(1+\lambda r_j)\big)$. Then
  $\det(Y-\Theta)=\prod_j(Y-r_j)$.
\item Let $l$ be a prime, $u\in\mathbf Z_l^\times$ and $c\in\mathbf Z_l$ with $u=1+lc$, and $m\in\mathbf Z$.
  Then $u^{-m}=1+l(-mc+le)$ for some $e\in\mathbf Z_l$.
\end{enumerate}
\end{lemma}
\leanref{Formalism/PieriBlockIdeals/CapelliCentralEnveloping.lean}{Formalism.CohomologicalDescent.HeightDescentMinors.DiagMinors}
\leanref{Formalism/PieriBlockIdeals/PieriColumnMinors/HeightDescentMinors.lean}{Formalism.CohomologicalDescent.HeightDescentMinors.diagonal_minors_holds}
\leanref{Formalism/PieriBlockIdeals/PieriColumnMinors/HeckeLFunctionIdealEulerProduct.lean}{Formalism.CohomologicalDescent.FrobeniusPeriodChange.constantCoeff_conj_diagonal_rescale}
\leanref{Formalism/PieriBlockIdeals/CapelliCentralEnveloping.lean}{Formalism.CohomologicalDescent.HeightDescentMinors.CharpolyDiag}
\leanref{Formalism/PieriBlockIdeals/PieriColumnMinors/HeightDescentMinors.lean}{Formalism.CohomologicalDescent.HeightDescentMinors.charpoly_diagonal_holds}
\leanref{Formalism/PieriBlockIdeals/CapelliCentralEnveloping.lean}{Formalism.CohomologicalDescent.HeightDescentMinors.ThetaCongPadic}
\leanref{Formalism/PieriBlockIdeals/PieriColumnMinors/HeightDescentMinors.lean}{Formalism.CohomologicalDescent.HeightDescentMinors.theta_congruence_padic}
\begin{proof}
(1) Let the rows $f(0)<\dots<f(t-1)$ and the columns $g(0)<\dots<g(t-1)$ give a minor of
$\mathrm{diag}(x^{a_j})$. Its $i$-th row is $x^{a_{f(i)}}$ times the $i$-th row of the corresponding
submatrix of the identity, so the minor is $x^{\sum_ia_{f(i)}}$ times a minor of the identity, and
$\sum_ia_{f(i)}\geq S$ for the set $J=\{f(0),\dots,f(t-1)\}$. Every $t\times t$ minor of a product
$GMH$ is a linear combination of $t\times t$ minors of $M$, as in the proof of
Lemma~\ref{is-onto-step-exp-bound-link-depth-base-change}(2).

(2) As $(X^{a_j})(\upsilon X)=\upsilon^{a_j}X^{a_j}$, the hypothesis says
$\big(\Gamma\cdot G(\upsilon X)\big)(i,j)\,\upsilon^{a_j}X^{a_j}=G(i,j)X^{a_j}$ for all $i,j$. Cancelling
$X^{a_j}$ in the domain $L[[X]]$ and taking constant terms, which the substitution $X\mapsto\upsilon X$
keeps, gives $\big(\Gamma(0)G(0)\big)(i,j)\,\upsilon^{a_j}=G(0)(i,j)$.

(3) Here $\mathrm{em}(1+\lambda\Theta)=1+\mathrm{em}(\lambda)\mathrm{em}(\Theta)$ and
$\mathrm{diag}\big(\mathrm{em}(1+\lambda r_j)\big)=1+\mathrm{em}(\lambda)\mathrm{diag}\big(\mathrm{em}(r_j)\big)$,
so the hypothesis gives
$\mathrm{em}(\lambda)\,\mathrm{em}(\Theta)G=\mathrm{em}(\lambda)\,G\,\mathrm{diag}\big(\mathrm{em}(r_j)\big)$.
As $\mathrm{em}(\lambda)\ne0$, $\mathrm{em}(\Theta)=G\,\mathrm{diag}\big(\mathrm{em}(r_j)\big)G^{-1}$, whose
characteristic polynomial is $\prod_j(Y-\mathrm{em}(r_j))$. This is the image under em of
$\det(Y-\Theta)$ and of $\prod_j(Y-r_j)$, and em is injective.

(4) Let $U$ be a unit of a commutative ring with $U=1+xv$. For $n\geq0$,
$(1+xv)^n=1+nxv+x^2s$ for some $s$, by the binomial theorem. Then
$U^n\big(U^{-n}-(1-nxv)\big)=1-(1+nxv+x^2s)(1-nxv)=x^2\big(n^2v^2-s(1-nxv)\big)$, so
$U^{-n}\equiv1-nxv$ modulo $x^2$ as $U^n$ is a unit. So $U^\kappa\equiv1+\kappa xv$ modulo $x^2$ for every
$\kappa\in\mathbf Z$. With $U=u$, $x=l$, $v=c$ and $\kappa=-m$, $u^{-m}=1-mlc+l^2e=1+l(-mc+le)$.
% Lean: diagonal_minors_holds (DiagMinors: Matrix.det_mul_column, minors_dvd_mul_holds);
% constantCoeff_conj_diagonal_rescale (rescale_X, mul_right_cancel, coeff_rescale at 0);
% charpoly_diagonal_holds (CharpolyDiag: smul_right_injective, Matrix.charpoly_units_conj,
% Matrix.charpoly_diagonal, Polynomial.map_injective); theta_congruence_padic (ThetaCongPadic) from
% zpow_congruence_general (Int.eq_nat_or_neg, sq_dvd_one_add_pow_sub_linear).
\end{proof}

\begin{lemma}\label{is-onto-step-exp-bound-link-depth-norm}
Let $l\ne2$ be a prime, $O$ a domain that is a $\mathbf Z_l$-algebra and complete for the $I$-adic
topology of an ideal $I$, $L$ a field, $\zeta\in O$ a primitive $l$-th root of unity with $\zeta-1\in I$,
$\mathrm{em}\colon O\to L$ an injective ring homomorphism, $u\in\mathbf Z_l^\times$ and $c\in\mathbf Z_l$
with $u=1+lc$, and $\jmath_b\colon O[[X]]\to L[[X]]$ ring homomorphisms for $b\geq0$, with $\jmath_b$ a
local exponential at $\zeta^b-1$ for $u$ for every $0<b<l$. Let $N$ be a lattice datum of rank $n$ over
$O$, indexed by a set $E$, $U'_i$ integer sequences of length $n$ with $U'_i(a)\geq0$ for all $i$ and
$a$, and $h\geq0$ an integer. If $N$ has local height $h$ with weights $U'$, then it has height $h$ with
weights $U'$
(Definition~\ref{is-onto-step-exp-bound-link-depth-data}).
\end{lemma}
\leanref{Formalism/PieriBlockIdeals/LemmaWFromLattice.lean}{Formalism.CohomologicalDescent.FrobeniusPeriodChange.LatWOfN}
\leanref{Formalism/PieriBlockIdeals/PieriColumnMinors/HeckeLFunctionIdealEulerProduct.lean}{Formalism.CohomologicalDescent.FrobeniusPeriodChange.weight_lattice_of_norm}
\begin{proof}
Let $i\in E$. By the local height, $\iota(P_i)Q=q^h$ for some $Q$, condition (d) holds, and for every
$0<b<l$ there are $G_b,H_b\in GL_n(L[[X]])$ with $\jmath_b\big(\iota(P_i)\big)=G_b\Delta_iH_b$ and
$\jmath_b\big(\iota(1+XD_i)\big)\cdot(G_b\Delta_i)\big(\mathrm{em}(u)X\big)=G_b\Delta_i$, where
$\Delta_i=\mathrm{diag}(X^{U'_i(a)})$ as $U'_i\geq0$. By
Lemma~\ref{is-onto-step-exp-bound-link-depth-norm-local} for $b=1$, $l\in I$, $l\ne0$ in $O$ and
$\jmath_1(q_0)$ is $X$ times a unit. We check (a) to (e) of the height.

(a) Lemma~\ref{is-onto-step-exp-bound-link-depth-norm-descent}(1) gives $M$ with $P_iM=(X-l)^h$.

(b) Let $r,S\geq0$ with $S\le\sum_{j\in J}U'_i(j)$ for every set $J$ of $r$ indices, and let $m$ be an
$r$-minor of $P_i$. Let $0<b<l$. Then $\zeta^b-1\in I$, and $\jmath_b\big(\iota(m)\big)$ is the
corresponding minor of $\jmath_b\big(\iota(P_i)\big)=G_b\Delta_iH_b$, which is divisible by $X^S$ by
Lemma~\ref{is-onto-step-exp-bound-link-depth-norm-matrices}(1). As $\jmath_b$ is a local exponential at
$\zeta^b-1$, Lemma~\ref{is-onto-step-exp-bound-link-depth-norm-descent}(4) gives
$\big(X-(\zeta^b-1)\big)^S\mid\iota(m)$. By Lemma~\ref{is-onto-step-exp-bound-link-depth-norm-descent}(3),
$(X-l)^S\mid m$.

(c) By (a), $\det P_i\det M=(X-l)^{hn}$. As $l\in I$, $X-l$ is a prime element of $O[[X]]$
(Lemma~\ref{is-onto-step-exp-bound-link-depth-norm-division}(2)), so $\det P_i=v(X-l)^e$ for a unit $v$
and some $e\geq0$. Apply $\jmath_1\circ\iota$: on one side
$\jmath_1\big(\iota(\det P_i)\big)=\det G_1\cdot X^{\sum_jU'_i(j)}\cdot\det H_1$, which is
$X^{\sum_jU'_i(j)}$ times a unit; on the other,
$\jmath_1\big(\iota((X-l)^e)\big)=\jmath_1(q_0)^e$ is $X^e$ times a unit. So
$e=\sum_jU'_i(j)=\big(\sum_jU'_i(j)\big)^+$.

(d) This is part of the local height.

(e) For every $j$, Lemma~\ref{is-onto-step-exp-bound-link-depth-norm-matrices}(4) gives
$u^{-U'_i(j)}=1+lr_j$ with $r_j=-U'_i(j)c+le_j$ for some $e_j\in\mathbf Z_l$. Put
$\Gamma=\jmath_1\big(\iota(1+XD_i)\big)$ and $G=G_1$. By
Lemma~\ref{is-onto-step-exp-bound-link-depth-norm-matrices}(2) with $\upsilon=\mathrm{em}(u)$,
$\Gamma(0)G(0)\,\mathrm{diag}\big(\mathrm{em}(u)^{U'_i(j)}\big)=G(0)$. By (d),
$X-l$ divides $(1+XD_i)(a,b)-(1+l\Theta_i)(a,b)=X\big(D_i(a,b)-\Theta_i(a,b)\big)+(X-l)\Theta_i(a,b)$, so
$\Gamma(0)=\mathrm{em}(1+l\Theta_i)$ by Lemma~\ref{is-onto-step-exp-bound-link-depth-norm-local}. As
$u^{U'_i(j)}(1+lr_j)=1$, $\mathrm{diag}\big(\mathrm{em}(u)^{U'_i(j)}\big)$ is the inverse of
$\mathrm{diag}\big(\mathrm{em}(1+lr_j)\big)$, so
$\mathrm{em}(1+l\Theta_i)G(0)=G(0)\,\mathrm{diag}\big(\mathrm{em}(1+lr_j)\big)$. The matrix $G(0)$ is
invertible, as $G$ is. As $l\ne0$ in $O$,
Lemma~\ref{is-onto-step-exp-bound-link-depth-norm-matrices}(3) gives $\det(Y-\Theta_i)=\prod_j(Y-r_j)$.
% Lean: weight_lattice_of_norm (LatWOfN): hW1 from height_descent_holds; localization_facts_hold at b = 1;
% the minors through roots_descent_holds, order_descent_holds, pi0Subst_det_map and diagonal_minors_holds;
% the determinant through linear_prime_holds, dvd_prime_pow and Associated.map by jj 1 o inc
% (pow_dvd_pow_iff for X); the characteristic polynomial through theta_congruence_padic,
% constantCoeff_conj_diagonal_rescale, one_add_X_smul_sub_C_dvd (gamma_mod_q_holds), hunit and
% charpoly_diagonal_holds with lam = l.
\end{proof}

\begin{lemma}\label{is-onto-step-exp-bound-link-depth-period}
Let $k,j,t,q,h\geq0$ be integers and $l\ne2$ a prime with $t+2\kappa<l$, where
$\kappa=\lfloor h/(l-1)\rfloor$, and let $N=k+j+t+q$. Let $O$ be a domain that is a $\mathbf Z_l$-algebra
and complete for the $I$-adic topology of an ideal $I$, $F$ and $L$ fields with $F$ of characteristic
$l$, $E$ a finite set with a permutation $\sigma$, $u\in\mathbf Z_l^\times$, $c\in\mathbf Z_l$,
$\zeta\in O$, ring homomorphisms $\mathrm{em}\colon O\to L$ and $\jmath_b\colon O[[X]]\to L[[X]]$ for
$b\geq0$, and $\mathrm{bot}_i,z_i,\nu_i(s)\in\mathbf Z$ and $U_i,U'_i\in\mathbf Z^N$ for $i\in E$ and
$s<t$. Let $\mathcal W$ be Witt lattice data for $(E,t,N,O,F)$ that satisfy the inputs of the Frobenius
period change for $(\sigma,l,h,u,c,U',\mathrm{wt},\nu,I,\zeta,\mathrm{em},\jmath)$, where
$\mathrm{wt}_i(s)=\mathrm{bot}_i+s$ (Definition~\ref{is-onto-step-exp-bound-link-depth-data}(6)). Suppose that for every $i$, $U_i$ and
$U'_i$ decrease strictly, $U_i$ takes the values $\mathrm{bot}_i+s$ for $0\le s<t$, and
$\prod_a(Y+U_i(a)+z_i)=\prod_a(Y+U'_i(a)+z_i)$ in $F[Y]$; and let $\tau\in E$ with $U_\tau(a)\geq0$ for
every $a$, $\sum_aU_\tau(a)=\sum_aU'_\tau(a)$, and $U_\tau(a)=\mathrm{bot}_\tau+(k+j+t-1-a)$ for
$k+j\le a<k+j+t$. Let $n_i=n_{l,\kappa}(\mathrm{bot}_i,\{U_i(a):a<N\})$. Then
\[
 \sum_{a<k+j}\big(U_\tau(a)-U'_\tau(a)\big)\le qh+\big(l\,n_{\sigma(\tau)}+n_\tau\big)(t+2\kappa).
\]
\end{lemma}
\leanref{Formalism/PieriBlockIdeals/LemmaWFromLattice.lean}{Formalism.CohomologicalDescent.FrobeniusPeriodChange.WEFromLattice}
\leanref{Formalism/PieriBlockIdeals/PieriColumnMinors/HeckeLFunctionIdealEulerProduct.lean}{Formalism.CohomologicalDescent.FrobeniusPeriodChange.weight_exponents_from_lattice}
\leanref{Formalism/PieriBlockIdeals/LemmaWFromLattice.lean}{Formalism.CohomologicalDescent.FrobeniusPeriodChange.WInputsOfN}
\leanref{Formalism/PieriBlockIdeals/PieriColumnMinors/HeckeLFunctionIdealEulerProduct.lean}{Formalism.CohomologicalDescent.FrobeniusPeriodChange.weight_inputs_of_norm}
\begin{proof}
The frame condition (a) gives $u=1+lc$, and (b) gives $U'_i(j)\geq0$ and $\mathrm{bot}_i+s\geq0$.
Lemma~\ref{is-onto-step-exp-bound-link-depth-norm}, for $N$ with the weights $U'$ and for every
$N^\psi_s$ with the weights $(\mathrm{bot}_i+s)_{i\in E}$, shows that $\mathcal W$ has the weight inputs
for $(\sigma,l,h,u,c,U',\mathrm{wt},\nu)$. Lemma~\ref{is-onto-step-exp-bound-link-depth-inputs} gives the
bound.
% Lean: weight_exponents_from_lattice = weight_exponents_from_inputs applied to weight_inputs_of_norm
% (WInputsOfN), which is weight_lattice_of_norm for W.N and for each W.Npsi s.
\end{proof}

\begin{lemma}\label{is-onto-step-exp-bound-link-depth-orient}
Let $k,j,t,q,h\geq0$ be integers and $l\ne2$ a prime with $t+2\kappa<l$, where
$\kappa=\lfloor h/(l-1)\rfloor$, and let $N=k+j+t+q$. Let $F$ be a field of characteristic $l$, $E$ a
finite set with a permutation $\sigma$, and $(\bar{\ },\Phi,U,U',\mathrm{bot})$ labels of length $N$ on
$E$ that are reflected about $h$ and linked over $F$ with constants $z_i$; and let sequences $J_i$ and
integers $\beta_i$ be carried by the labels through a map $\epsilon$
(Definition~\ref{is-onto-step-exp-bound-link-depth-reflected}). Let $O$ be a domain that is a
$\mathbf Z_l$-algebra and complete for the $I$-adic topology of an ideal $I$, $L$ a field,
$u\in\mathbf Z_l^\times$, $c\in\mathbf Z_l$, $\zeta\in O$, ring homomorphisms $\mathrm{em}\colon O\to L$
and $\jmath_b\colon O[[X]]\to L[[X]]$ for $b\geq0$, integers $\nu_i(s)$ for $i\in E$ and $s<t$, and
$\mathcal W$ Witt lattice data for $(E,t,N,O,F)$ that satisfy the inputs of the Frobenius period change
for $(\sigma,l,h,u,c,J,\mathrm{wt},\nu,I,\zeta,\mathrm{em},\jmath)$, where $\mathrm{wt}_i(s)=\beta_i+s$
(Definition~\ref{is-onto-step-exp-bound-link-depth-data}(6)). Let $\tau\in E$ with $U_\tau(a)\geq0$ for
every $a$, $\sum_aU_\tau(a)=\sum_aU'_\tau(a)$, and $U_\tau(a)=\mathrm{bot}_\tau+(k+j+t-1-a)$ for
$k+j\le a<k+j+t$. Let $n_i=n_{l,\kappa}(\mathrm{bot}_i,\{U_i(a):a<N\})$. Then
\[
 \sum_{a<k+j}\big(U_\tau(a)-U'_\tau(a)\big)\le qh+\big(l\,n_{\sigma(\tau)}+n_\tau\big)(t+2\kappa).
\]
\end{lemma}
\leanref{Formalism/PieriBlockIdeals/MultiplicityZeroThickEmbedding.lean}{Formalism.CohomologicalDescent.WorkedExample.WEOrientLattice}
\leanref{Formalism/PieriBlockIdeals/PieriColumnMinors/WorkedExample.lean}{Formalism.CohomologicalDescent.WorkedExample.orientation_lattice_holds}
\begin{proof}
First, $\epsilon(\epsilon(i))=i$ and $\sigma(\epsilon(i))=\epsilon(\sigma(i))$ for every $i$: for
$\epsilon$ the identity this is clear, and for $\epsilon$ the map $i\mapsto\bar i$ it is part of the
reflection. So $J_i=U'_{\epsilon(i)}$ and $\beta_i=\mathrm{bot}_{\epsilon(i)}$ for every $i$.

Next, every $i\in E$ satisfies the linked conditions with the constant $0$: $U_i$ and $U'_i$ decrease
strictly, $U_i$ takes the values $\mathrm{bot}_i+s$ for $s<t$, and
$\prod_a(Y+U_i(a))=\prod_a(Y+U'_i(a))$ in $F[Y]$. For $i\in\Phi$, substitute $Y-z_i$ for $Y$ in the
identity of the linkage. Otherwise $i'=\bar i$ lies in $\Phi$ and $i=\bar{i'}$. Then
$U_i(a)=h-U_{i'}(N-1-a)$ and $U'_i(a)=h-U'_{i'}(N-1-a)$ decrease strictly; $U_i$ takes the value
$h-(\mathrm{bot}_{i'}+s')=\mathrm{bot}_i+(t-1-s')$ for every $s'<t$; and substituting $-Y-h-z_{i'}$ for
$Y$ in the identity for $i'$ and multiplying by $(-1)^N$ gives
$\prod_a(Y+h-U_{i'}(a))=\prod_a(Y+h-U'_{i'}(a))$, which after the change of index $a\mapsto N-1-a$ is the
identity for $i$.

Apply Lemma~\ref{is-onto-step-exp-bound-link-depth-period} to the sequences $V_i=U_{\epsilon(i)}$ and
$J_i$, the integers $\beta_i$ and the constants $0$, at the index $\epsilon(\tau)$. Its conditions hold:
$V_i$ and $J_i=U'_{\epsilon(i)}$ satisfy the linked conditions with the constant $0$, and $V_i$ takes the
values $\beta_i+s$; and $V_{\epsilon(\tau)}=U_\tau$, $J_{\epsilon(\tau)}=U'_\tau$ and
$\beta_{\epsilon(\tau)}=\mathrm{bot}_\tau$, so the conditions at $\epsilon(\tau)$ are those assumed at
$\tau$. As $\sigma(\epsilon(\tau))=\epsilon(\sigma(\tau))$, $V_{\sigma(\epsilon(\tau))}=U_{\sigma(\tau)}$
and $\beta_{\sigma(\epsilon(\tau))}=\mathrm{bot}_{\sigma(\tau)}$, so its conclusion is the bound.
% Lean: orientation_lattice_holds (WEOrientLattice): reflection_carried_facts, linked_all_holds
% (prod_X_add_C_shift_eq for i in Phi; linkage_reflection_holds and Fin.revPerm otherwise), and
% weight_exponents_from_lattice at ep tau with U := D.U o ep and cst := 0, simplified by e1 to e6.
\end{proof}

\begin{lemma}\label{is-onto-step-exp-bound-link-depth}
Let $P$ be a parameter record that satisfies the hypotheses and is large, $W$ weights for $P$ that
satisfy the weight conditions (Definition~\ref{is-onto-step-exp-bound-weights}), $l>l_J$ a prime,
$\tau$ an index and $\omega\in\mathbf N^{N_r}$ a weight linked to $W$ at $l$ at $\tau$
(Definition~\ref{is-onto-step-exp-bound-hida}(5)). Then
$\sum_{a<k-1+j_\tau}(\lambda_{\tau,a}-\omega_a)\le i_\tau(l)$.
\end{lemma}
\leanref{Formalism/PieriBlockIdeals/LatticeInclusionRealSettingAlt.lean}{Formalism.CohomologicalDescent.ShortWeightBounds.DepthOfLink}
\leanref{Formalism/PieriBlockIdeals/PieriColumnMinors/WeightBounds.lean}{Formalism.CohomologicalDescent.ShortWeightBounds.depth_of_linkage_holds}
\begin{proof}
Take the data of the linkage, as in Definition~\ref{is-onto-step-exp-bound-hida}(5), and put
$\kappa=\kappa(l)$. By (W1), $n_\tau=k-1+j_\tau+t+q_\tau=N_r$. By
Lemma~\ref{is-onto-step-exp-bound-total-numbers}, $N_r+1<l$ and $t+2\kappa<l$; as $N_r\geq1$ by the
hypotheses, $l\ne2$. As $l-1$ is a positive integer, $\kappa=\lfloor h/(l-1)\rfloor=\lfloor
h_0/(l-1)\rfloor$.

The labels satisfy the conditions of Lemma~\ref{is-onto-step-exp-bound-link-depth-orient} at
$\iota(\tau)$, with $(k-1,j_\tau,t,q_\tau,h_0)$ in place of $(k,j,t,q,h)$. Indeed, for $a<N_r$,
$U_{\iota(\tau)}(a)=U_\tau(a)=\lambda_{\tau,a}+N_r-1-a\geq0$; the equality of the sums is part of the
linkage; and for $p_\tau\le a<p_\tau+t$, (W2) gives $\lambda_{\tau,a}=\lambda_{\tau,p_\tau+t-1}$, so
$U_\tau(a)=\mathrm{bot}_\tau+(p_\tau+t-1-a)$, with $\mathrm{bot}_\tau=\mathrm{bot}_{\iota(\tau)}$ and
$p_\tau=k-1+j_\tau$. With $n_i=n_{l,\kappa}(\mathrm{bot}_i,\{U_i(a):a<N_r\})$ for $i\in E$, that lemma
gives
\[
 \sum_{a<k-1+j_\tau}\big(U_{\iota(\tau)}(a)-U'_{\iota(\tau)}(a)\big)
 \le q_\tau h_0+\big(l\,n_{\sigma(\iota(\tau))}+n_{\iota(\tau)}\big)(t+2\kappa).
\]

For every $\tau'$, $n_{\iota(\tau')}=n_{\tau'}(l)$, as the values of $U_{\iota(\tau')}$ form the set
$\mathcal U_{\tau'}$ and $\mathrm{bot}_{\iota(\tau')}=\mathrm{bot}_{\tau'}$. The map $u\mapsto h_0-u$
carries $[\beta,\beta+t-1]$ onto $[h_0-(\beta+t-1),h_0-\beta]$ and $[\beta-\kappa,\beta+t-1+\kappa]$ onto
$[h_0-(\beta+t-1)-\kappa,h_0-\beta+\kappa]$, and it preserves congruence modulo $l$; so
\[
 n_{l,\kappa}\big(h_0-(\beta+t-1),\{h_0-u:u\in\mathcal U\}\big)=n_{l,\kappa}(\beta,\mathcal U)
\]
for every integer $\beta$ and finite set $\mathcal U$. As the values of $U_{\bar i}$ are $h_0$ minus
those of $U_i$ and $\mathrm{bot}_{\bar i}=h_0-(\mathrm{bot}_i+t-1)$, this gives $n_{\bar i}=n_i$ for
every $i\in E$. Let $\tau'=\phi_l(\tau)$. If $s_l(\tau)=0$, then $\sigma(\iota(\tau))=\iota(\tau')$, and
$n_{\sigma(\iota(\tau))}=n_{\tau'}(l)=n^{\mathrm F}_\tau(l)$. If $s_l(\tau)=1$, then
$\sigma(\iota(\tau))=\overline{\iota(\tau')}$ and $n_{\sigma(\iota(\tau))}=n_{\tau'}(l)$; and as
$\mathrm{top}_{\tau'}=\mathrm{bot}_{\tau'}+t-1$ by (W2), the identity above gives
$n^{\mathrm F}_\tau(l)=n_l(h_0-\mathrm{top}_{\tau'},\{h_0-u:u\in\mathcal U_{\tau'}\})=n_{\tau'}(l)$. In
both cases $n_{\sigma(\iota(\tau))}=n^{\mathrm F}_\tau(l)$.

As $U'_{\iota(\tau)}(a)=\omega_a+N_r-1-a$ for $a<N_r$, the left side is
$\sum_{a<k-1+j_\tau}(\lambda_{\tau,a}-\omega_a)$. As $h_0\le h$, the right side is at most
$q_\tau h+\big(l\,n^{\mathrm F}_\tau(l)+n_\tau(l)\big)\big(t+2\kappa(l)\big)$; being an integer, it is at
most the floor $i_\tau(l)$ of that number (Definition~\ref{is-onto-step-exp-bound-recursion}(2)).
% Lean: depth_of_linkage_holds (DepthOfLink): paramsLargePrime_holds (Nr + 1 < l and t + 2 kap l < l,
% Lemma is-onto-step-exp-bound-total-numbers) with rank_pos_of_hyp for l /= 2; kappa_floor_holds
% (Nat.floor_div_natCast); hshape from weight_nonneg', hsum and weight_on_block_eq_bot_add;
% WorkedExample.orientation_lattice_holds (Lemma is-onto-step-exp-bound-link-depth-orient); hmw
% (image_univ_eq_image_range) and hmwF (weight_letters_holds and weight_reflection_holds, which is the
% displayed identity, with top_eq_bot_add_t_sub_one); linkage_depth_eq_sum_difference
% (depth_as_weight_sum_holds) and depth_nat_le_holds (Nat.le_floor, h0 <= h). Size before this step was
% written out: about 5,200 lines of Lean in 390 theorems.
\end{proof}

\begin{lemma}\label{is-onto-step-exp-bound-link}
Let $P$ be a parameter record that satisfies the hypotheses and is large, $W$ weights for $P$ that
satisfy the weight conditions, $l>l_J$ a prime, and $X$ a Hida setting of rank $N_r$ and the weights
$\lambda$ of $W$ that is linked to $W$ at $l$ (Definition~\ref{is-onto-step-exp-bound-hida}). Then $X$
has the weight inputs for $W$ at $l$.
\end{lemma}
\leanref{Formalism/PieriBlockIdeals/LatticeInclusionRealSettingAlt.lean}{Formalism.CohomologicalDescent.ShortWeightBounds.WInputsOfLink}
\leanref{Formalism/PieriBlockIdeals/PieriColumnMinors/WeightBounds.lean}{Formalism.CohomologicalDescent.ShortWeightBounds.weight_inputs_of_linkage_holds}
\begin{proof}
Take the constituent data at $l$ of the linkage. For every relevant $\mu$ and every $\tau\in T_W$, the
weight $\omega_{\mu,\tau}$ is linked to $W$ at $l$ at $\tau$, so
Lemma~\ref{is-onto-step-exp-bound-link-depth} gives
$\sum_{a<k-1+j_\tau}(\lambda_{\tau,a}-\omega_{\mu,\tau,a})\le i_\tau(l)$.
% Lean: weight_inputs_of_linkage_holds keeps the data and applies first_input_of_link, which is
% depth_of_linkage_holds at every relevant mu and tau.
\end{proof}

\subsection{Block modules and Jucys--Murphy operators}\label{subsection-is-jucys-murphy}

\begin{definition}\label{is-onto-step-exp-bound-short-branching}
Let $l$ be a prime. For a nonincreasing $\nu\in\mathbf N^r$, $\nabla_{\mathbf Z}(\nu)$ and
$\Delta_{\mathbf Z}(\nu)$ are the dual Weyl module and the Weyl module of
Definition~\ref{pd-weight-lattices}(1) over $\mathbf Z$, in the polynomials in the entries of an
$r\times r$ matrix. A polynomial has \emph{row degree} $\mu\in\mathbf N^r$ if it is of degree $\mu_a$ in
the entries of the $a$-th row for every $a$.
\begin{enumerate}
\item For $\nu,\mu\in\mathbf N^r$ with $\nu$ nonincreasing and an integer $e\geq0$, the pair
  $(\nu,\mu)$ has \emph{exponent at most $e$} at $l$ if for every $x\in\nabla_{\mathbf Z}(\nu)$ of row
  degree $\mu$ there is an integer $m$ not divisible by $l$ with $ml^ex\in\Delta_{\mathbf Z}(\nu)$.
\item For $\xi\in\mathbf Z^r$ and $\lambda,\mu\in\mathbf N^r$ put
  $\mathrm{dp}_\xi(\lambda,\mu)=\sum_{a<r}\xi_a(\lambda_a-\mu_a)$. For integers $p,q,w\geq0$ and
  $\kappa\in\mathbf N^q$, let $w^p\kappa\in\mathbf N^{p+q}$ be the sequence
  $(w,\dots,w,\kappa_0,\dots,\kappa_{q-1})$ and $\kappa w^p\in\mathbf N^{q+p}$ the sequence
  $(\kappa_0,\dots,\kappa_{q-1},w,\dots,w)$, each with $p$ entries $w$.
\item For $\lambda_2,\tau\in\mathbf N^q$, indexed by $0\le a<q$, $c\in\mathbf Z$ and $y\geq0$, let
  $b(\lambda_2,\tau,c)$ be the number of $a$ with $\lambda_{2,a}\le c+a<\tau_a$, and
  $\mathrm{col}(\lambda_2,\tau,y)$ the number of $a$ with $\lambda_{2,a}\le y<\tau_a$. For integers
  $p,w,i\geq0$, let $\mathrm{Sh}(p,w,i,\lambda_2)$ be the set of nonincreasing $\tau\in\mathbf N^q$ with
  $\lambda_{2,a}\le\tau_a\le\lambda_{2,a}+i$ for every $a$, $\sum_a\tau_a=\sum_a\lambda_{2,a}+i$,
  $\mathrm{col}(\lambda_2,\tau,y)\le p$ for every $a$ and every $y<\tau_a$, and $b(\lambda_2,\tau,w)=0$.
  Let $\mu(c)$ be the largest $b(\lambda_2,\tau,c)$ for $\tau\in\mathrm{Sh}(p,w,i,\lambda_2)$, and $0$ if
  that set is empty, and put
  \[
   D_n(p,w,i,\lambda_2)=\prod_{0\le y<q+w+i}(w+q-y)^{\mu(y-q)}.
  \]
  Let $v_l$ be the $l$-adic valuation on $\mathbf Z$, with $v_l(0)=0$.
\item The \emph{block exponent bound holds in graded form} if the following holds for every prime $l$, all integers
  $p,q,w\geq0$, every nonincreasing $\lambda_2\in\mathbf N^q$ with $\lambda_{2,b}\le w$ for every $b$,
  every nonincreasing $\xi\in\mathbf Z^{p+q}$, every $i'\in\mathbf Z$ and every integer $e_L\geq0$. If
  $(\lambda_2,\mu_2)$ has exponent at most $e_L$ at $l$ for every $\mu_2\in\mathbf N^q$ with
  $\mathrm{dp}_\xi(w^p\lambda_2,w^p\mu_2)\le i'$, then $(w^p\lambda_2,\mu)$ has exponent at most
  $e_L+v_l\big(D_n(p,w,i,\lambda_2)\big)$ at $l$ for every integer $i\geq0$ and every
  $\mu\in\mathbf N^{p+q}$ with $\sum_{a<p}\mu_a+i=pw$ and $\mathrm{dp}_\xi(w^p\lambda_2,\mu)\le i'$.
\item The \emph{block exponent bound holds in last graded form} if the following holds for every prime $l$, all integers
  $p,q,w,M\geq0$, every nonincreasing $\lambda_2\in\mathbf N^q$ with $w\le\lambda_{2,b}\le M$ for every
  $b$, every nonincreasing $\xi\in\mathbf Z^{q+p}$, every $i'\in\mathbf Z$ and every integer $e_L\geq0$.
  If $(\lambda_2,\mu_2)$ has exponent at most $e_L$ at $l$ for every $\mu_2\in\mathbf N^q$ with
  $\mathrm{dp}_\xi(\lambda_2w^p,\mu_2w^p)\le i'$, then $(\lambda_2w^p,\mu)$ has exponent at most
  $e_L+v_l\big(D_n(p,(M-w)^+,i,\lambda_2^\vee)\big)$ at $l$ for every integer $i\geq0$ and every
  $\mu\in\mathbf N^{q+p}$ with $\sum_{a<p}\mu_{q+a}=pw+i$ and $\mathrm{dp}_\xi(\lambda_2w^p,\mu)\le i'$;
  here $\lambda^\vee_{2,b}=M-\lambda_{2,q-1-b}$ and $(M-w)^+=\max(0,M-w)$.
\end{enumerate}
\end{definition}
\leanref{Formalism/PieriBlockIdeals/StatementLD.lean}{Formalism.CohomologicalDescent.LD.nablaWt}
\leanref{Formalism/PieriBlockIdeals/StatementLD.lean}{Formalism.CohomologicalDescent.LD.KilledAt}
\leanref{Formalism/PieriBlockIdeals/StatementLD.lean}{Formalism.CohomologicalDescent.LD.ExpLe}
\leanref{Formalism/PieriBlockIdeals/StatementLD.lean}{Formalism.CohomologicalDescent.LD.depth2}
\leanref{Formalism/PieriBlockIdeals/StatementLD.lean}{Formalism.CohomologicalDescent.LD.blockWt}
\leanref{Formalism/PieriBlockIdeals/StatementLD.lean}{Formalism.CohomologicalDescent.LD.blockWtLast}
\leanref{Formalism/PieriBlockIdeals/StatementLD.lean}{Formalism.CohomologicalDescent.LD.boxesOfContent}
\leanref{Formalism/PieriBlockIdeals/StatementLD.lean}{Formalism.CohomologicalDescent.LD.colLen}
\leanref{Formalism/PieriBlockIdeals/StatementLD.lean}{Formalism.CohomologicalDescent.LD.shapes}
\leanref{Formalism/PieriBlockIdeals/StatementLD.lean}{Formalism.CohomologicalDescent.LD.mu}
\leanref{Formalism/PieriBlockIdeals/StatementLD.lean}{Formalism.CohomologicalDescent.LD.Dn}
\leanref{Formalism/PieriBlockIdeals/StatementLD.lean}{Formalism.CohomologicalDescent.LD.IsDepth}
\leanref{Formalism/PieriBlockIdeals/StatementLD.lean}{Formalism.CohomologicalDescent.LD.IsDepthLast}
\leanref{Formalism/PieriBlockIdeals/StatementLD.lean}{Formalism.CohomologicalDescent.LD.LemmaDGradedAt}
\leanref{Formalism/PieriBlockIdeals/StatementLD.lean}{Formalism.CohomologicalDescent.LD.LemmaDGraded}
\leanref{Formalism/PieriBlockIdeals/BranchingBoundLastGradedOneRow.lean}{Formalism.CohomologicalDescent.LD.LemmaDLastGradedAt}
\leanref{Formalism/PieriBlockIdeals/BranchingBoundLastGradedOneRow.lean}{Formalism.CohomologicalDescent.LD.LemmaDLastGraded}
% Lean: ExpLe l r nu mu e = KilledAt l r nu mu (l^e), over nablaWt r nu mu = dualWeyl Z r r nu meet the
% weighted homogeneous polynomials of rowW-degree mu; blockWt p q w kappa = w^p kappa and blockWtLast p q w
% kappa = kappa w^p; IsDepth and IsDepthLast are the conditions on the sums of mu; padicValInt for v_l.

\begin{definition}\label{is-onto-step-exp-bound-short-graded-block}
Let $K$ be a field, $I_1$ a finite set, $q\geq0$ and $I_2=\{0,\dots,q-1\}$. Write $\alpha,\beta$ for
elements of $I_1$, $b,b'$ for elements of $I_2$, and $\epsilon_c$ ($c\in I_1\sqcup I_2$) for the basis
vectors of $\mathbf Z^{I_1\sqcup I_2}$.
\begin{enumerate}
\item A \emph{block module} over $K$ for $I_1$ and $I_2$ is a $K$-vector space $V$ with endomorphisms
  $E_{cd}$ for $c,d\in I_1\sqcup I_2$, a subspace $V^+$ and an element $w\in K$, such that
  $E_{cd}E_{c'd'}-E_{c'd'}E_{cd}=\delta_{dc'}E_{cd'}-\delta_{d'c}E_{c'd}$ for all $c,d,c',d'$, and
  $E_{\alpha b}u=0$, $E_{\alpha\beta}u=\delta_{\alpha\beta}wu$ and $E_{bb'}u\in V^+$ for every $u\in V^+$.
  For $\kappa\in\mathbf Z^{I_1\sqcup I_2}$, $V_\kappa$ is the set of $v\in V$ with $E_{cc}v=\kappa_cv$ for
  every $c$.
\item For $v=(b,\alpha)\in I_2\times I_1$ put $f_v=E_{b\alpha}$; these commute. For a multiset $m$ of $n$
  elements of $I_2\times I_1$, $F_m$ is the product of the $f_v$ over the elements $v$ of $m$, counted
  with multiplicity, and $m!$ is the product of the factorials of the multiplicities. Let $O$ be a
  commutative ring with a homomorphism $O\to K$, acting on $V$ through it, and $\Delta_L\subset V^+$ an
  $O$-submodule. Then $\Delta_n(\Delta_L)$ is the $O$-span of the vectors $(m!)^{-1}F_mu$ with
  $u\in\Delta_L$ and $m$ of $n$ elements, and $\Delta(\Delta_L)=\sum_{n\geq0}\Delta_n(\Delta_L)$. For a
  bilinear form $B$ on $V$ and an $O$-submodule $N\subset V$, $N^B$ is the set of $v\in V$ with $B(v,x)$
  in the image of $O$ for every $x\in N$; and $\nabla_L$ is the set of $u\in V^+$ with $B(u,x)$ in the
  image of $O$ for every $x\in\Delta_L$.
\item For $n\geq0$, $T_n$ is the space of functions $z$ on $(I_2\times I_1)^n$ with values in $V^+$; an
  endomorphism of $V^+$ acts on $T_n$ through the values. For $s<n$, $e^{(s)}_{bb'}$ is the endomorphism of
  $T_n$ with $(e^{(s)}_{bb'}z)(g)=z(g')$ if $g_s=(b,\alpha)$ for some $\alpha$, where $g'$ is $g$ with
  $g_s$ replaced by $(b',\alpha)$, and $(e^{(s)}_{bb'}z)(g)=0$ otherwise. For $j\le n$ put
  $D_j(b,b')=E_{bb'}+\sum_{s<j}e^{(s)}_{bb'}$ and $C_j=\sum_{b,b'}D_j(b,b')D_j(b',b)$. The
  \emph{symmetrization} of $z\in T_n$ is the function $z^{\mathrm{sym}}$ on the multisets $m$ of $n$
  elements of $I_2\times I_1$ with $z^{\mathrm{sym}}(m)=\sum_gz(g)$, the sum over the $g$ whose entries
  form $m$.
\item For $\kappa\in\mathbf N^q$ put $\mathrm{cas}(\kappa)=\sum_{b<q}\kappa_b(\kappa_b+q-1-2b)$. For
  $\lambda_2,\tau\in\mathbf N^q$, a \emph{box path} of length $n$ from $\lambda_2$ to $\tau$ is a sequence
  $\kappa_0=\lambda_2,\kappa_1,\dots,\kappa_n=\tau$ of nonincreasing elements of $\mathbf N^q$, each
  $\kappa_{j+1}$ obtained from $\kappa_j$ by raising one entry by $1$. The \emph{box path space}
  $T_n(\lambda_2,\tau)$ is the span of the $z\in T_n$ for which some box path of length $n$ from
  $\lambda_2$ to $\tau$ has $C_jz=\mathrm{cas}(\kappa_j)z$ for every $j\le n$.
\end{enumerate}
\end{definition}
\leanref{Formalism/PieriBlockIdeals/IntegralQuotient.lean}{Formalism.CohomologicalDescent.Symmetrization.HidaSetting}
\leanref{Formalism/PieriBlockIdeals/IntegralQuotient.lean}{Formalism.CohomologicalDescent.Symmetrization.HidaSetting.wt}
\leanref{Formalism/PieriBlockIdeals/IntegralQuotient.lean}{Formalism.CohomologicalDescent.Symmetrization.HidaSetting.f}
\leanref{Formalism/PieriBlockIdeals/IntegralQuotient.lean}{Formalism.CohomologicalDescent.Symmetrization.HidaSetting.Fs}
\leanref{Formalism/PieriBlockIdeals/IntegralQuotient.lean}{Formalism.CohomologicalDescent.Symmetrization.mfact}
\leanref{Formalism/PieriBlockIdeals/IntegralQuotient.lean}{Formalism.CohomologicalDescent.Symmetrization.HidaSetting.DeltaN}
\leanref{Formalism/PieriBlockIdeals/LemmaDSkeletonUnit.lean}{Formalism.CohomologicalDescent.BinaryFormsExample.DeltaAll}
\leanref{Formalism/PieriBlockIdeals/LemmaDSkeletonUnit.lean}{Formalism.CohomologicalDescent.BinaryFormsExample.NablaL}
\leanref{Formalism/PieriBlockIdeals/IntegralQuotient.lean}{Formalism.CohomologicalDescent.Symmetrization.onM}
\leanref{Formalism/PieriBlockIdeals/IntegralQuotient.lean}{Formalism.CohomologicalDescent.Symmetrization.unitY}
\leanref{Formalism/PieriBlockIdeals/LemmaDSkeletonUnit.lean}{Formalism.CohomologicalDescent.BinaryFormsExample.Dpart}
\leanref{Formalism/PieriBlockIdeals/LemmaDSkeletonUnit.lean}{Formalism.CohomologicalDescent.BinaryFormsExample.Cas}
\leanref{Formalism/PieriBlockIdeals/IntegralQuotient.lean}{Formalism.CohomologicalDescent.Symmetrization.toSym}
\leanref{Formalism/PieriBlockIdeals/LemmaDSkeletonUnit.lean}{Formalism.CohomologicalDescent.BinaryFormsExample.casVal}
\leanref{Formalism/PieriBlockIdeals/LemmaDSkeletonUnit.lean}{Formalism.CohomologicalDescent.BinaryFormsExample.IsPath}
\leanref{Formalism/PieriBlockIdeals/LemmaDSkeletonUnit.lean}{Formalism.CohomologicalDescent.BinaryFormsExample.PathSpace}
% Lean: HidaSetting K V P R (fields E, comm, W, w, kill, scal, stab), with P = I_1 and R = Fin q = I_2;
% the Lean's E is indexed by P + R. f v = E (inr v.1) (inl v.2), Fs m the product over m.toList, mfact m,
% DeltaN n DelL and DeltaAll; B.dualSubmodule is N^B, and NablaL S B DelL = (B.restrict W).dualSubmodule
% DelL. Tens R P W n = (Fin n -> R x P) -> W, onM, unitY j a b (the book's e^(j)_ab), Dpart j b b', Cas j,
% toSym (with cont g the multiset of entries), casVal, IsPath (start, step, stop, dom) and PathSpace.

\begin{definition}\label{is-onto-step-exp-bound-short-graded-skeleton-jucys}
Let $K$ be a field, $I_1$ a finite set, $q,n\geq0$ integers, $I_2=\{0,\dots,q-1\}$, and $M$ a $K$-vector
space with endomorphisms $E_{bb'}$ for $b,b'\in I_2$. Define $T_n$, the operators $e^{(s)}_{bb'}$, $D_j(b,b')$
and $C_j$, the symmetrization $z\mapsto z^{\mathrm{sym}}$, and for $\lambda_2,\tau\in\mathbf N^q$ the box paths
from $\lambda_2$ to $\tau$ and the box path space $T_n(\lambda_2,\tau)$, as in
Definition~\ref{is-onto-step-exp-bound-short-graded-block}(3) and~(4), with $M$ in place of $V^+$ and these
$E_{bb'}$. Write $g\in(I_2\times I_1)^n$ as $g_s=(b_s,\alpha_s)$ for $s<n$, and let $\sigma$ be a permutation of
$\{0,\dots,n-1\}$.
\begin{enumerate}
\item $(\pi_\sigma z)(g)=z(g\circ\sigma)$; $(\pi^{\mathrm Y}_\sigma z)(g)=z(g')$ with $g'_s=(b_{\sigma(s)},\alpha_s)$; and
  $(\pi^{\mathrm X}_\sigma z)(g)=z(g'')$ with $g''_s=(b_s,\alpha_{\sigma(s)})$. An endomorphism $T$ of $M$ acts on
  $T_n$ through the values, $(Tz)(g)=T(z(g))$.
\item For $j<n$, the \emph{Jucys--Murphy operator} is
  $x_j=\sum_{b,b'\in I_2}e^{(j)}_{bb'}E_{b'b}+\sum_{j'<j}\pi^{\mathrm Y}_{(j'\,j)}$, where $(j'\,j)$ is a
  transposition. For $t\in K$ put $A(t)=(t-x_0)(t-x_1)\cdots(t-x_{n-1})$, and let
  $A(X)=(X-x_0)\cdots(X-x_{n-1})$ be the same product in the polynomials in $X$ with coefficients in
  $\operatorname{End}(T_n)$.
\item For $c\in\mathbf Z^n$ and $\lambda_2\in\mathbf N^q$, an element $z\in T_n$ has \emph{weight} $c$ for
  $\lambda_2$ if $C_0z=\mathrm{cas}(\lambda_2)z$ and $x_jz=c_jz$ for every $j<n$.
\item For a box path $\kappa_0,\dots,\kappa_n$, let $\rho_j$ be the index raised at step $j$, so that
  $\kappa_{j+1}=\kappa_j+\epsilon_{\rho_j}$, with $\epsilon_r$ the $r$-th basis vector. The box added at step
  $j$ has \emph{column} $y_j=\kappa_{j,\rho_j}$ and \emph{content} $c_j=\kappa_{j,\rho_j}-\rho_j$. A
  \emph{content sequence} of length $n$ for $\lambda_2$ is the sequence $(c_j)_{j<n}$ of a box path of length $n$
  from $\lambda_2$.
\item Write $[g]$ for the multiset of the entries of $g$, and $m!$ for the product of the factorials of the
  multiplicities of a multiset $m$. For a function $y$ on the multisets of $n$ elements of $I_2\times I_1$ with
  values in $M$, $Ny\in T_n$ is $(Ny)(g)=[g]!\,y([g])$.
\end{enumerate}
\end{definition}
\leanref{Formalism/PieriBlockIdeals/IntegralQuotient.lean}{Formalism.CohomologicalDescent.Symmetrization.perm}
\leanref{Formalism/PieriBlockIdeals/IntegralQuotient.lean}{Formalism.CohomologicalDescent.Symmetrization.permY}
\leanref{Formalism/PieriBlockIdeals/IntegralQuotient.lean}{Formalism.CohomologicalDescent.Symmetrization.permX}
\leanref{Formalism/PieriBlockIdeals/IntegralQuotient.lean}{Formalism.CohomologicalDescent.Symmetrization.Om}
\leanref{Formalism/PieriBlockIdeals/IntegralQuotient.lean}{Formalism.CohomologicalDescent.Symmetrization.xJM}
\leanref{Formalism/PieriBlockIdeals/IntegralQuotient.lean}{Formalism.CohomologicalDescent.Symmetrization.AA}
\leanref{Formalism/PieriBlockIdeals/IntegralQuotient.lean}{Formalism.CohomologicalDescent.Symmetrization.AApoly}
\leanref{Formalism/PieriBlockIdeals/StatementLC.lean}{Formalism.CohomologicalDescent.LC.IsWt}
\leanref{Formalism/PieriBlockIdeals/StatementLC.lean}{Formalism.CohomologicalDescent.LC.colSeq}
\leanref{Formalism/PieriBlockIdeals/StatementLC.lean}{Formalism.CohomologicalDescent.LC.cseq}
\leanref{Formalism/PieriBlockIdeals/StatementLC.lean}{Formalism.CohomologicalDescent.LC.IsContentSeq}
\leanref{Formalism/PieriBlockIdeals/IntegralQuotient.lean}{Formalism.CohomologicalDescent.Symmetrization.NmS}
% Lean: perm, permY, permX and onM on Tens R P M n; Om k EM j = sum over b, b' of unitY j b b' * onM (EM b' b)
% and xJM = Om + the permY (swap j' j) for j' < j; AA k EM n t and AApoly k EM n; IsWt K EM lam2 c z; colSeq,
% cseq and IsContentSeq (a box path given by kappa and the rows rho).

\begin{lemma}\label{is-onto-step-exp-bound-short-graded-skeleton-jucys-relations}
Let $K$, $I_1$, $q$, $n$, $M$ and $E_{bb'}$ be as in
Definition~\ref{is-onto-step-exp-bound-short-graded-skeleton-jucys}, with the operators defined there, and let
$\sigma,\sigma'$ be permutations of $\{0,\dots,n-1\}$.
\begin{enumerate}
\item $\pi^{\mathrm Y}_{\sigma\sigma'}=\pi^{\mathrm Y}_\sigma\pi^{\mathrm Y}_{\sigma'}$, $\pi^{\mathrm Y}_1=1$ and
  $\pi_\sigma=\pi^{\mathrm Y}_\sigma\pi^{\mathrm X}_\sigma$. The operators $\pi^{\mathrm Y}_\sigma$ and
  $\pi^{\mathrm X}_\sigma$ commute with every endomorphism of $M$ acting through the values,
  $\pi^{\mathrm X}_\sigma$ commutes with every $e^{(s)}_{bb'}$, and
  $\pi^{\mathrm Y}_\sigma e^{(s)}_{bb'}=e^{(\sigma(s))}_{bb'}\pi^{\mathrm Y}_\sigma$.
\item $\sum_{b,b'}e^{(s)}_{bb'}e^{(s')}_{b'b}=\pi^{\mathrm Y}_{(s\,s')}$ for $s\ne s'$, and
  $\sum_{b,b'}e^{(s)}_{bb'}e^{(s)}_{b'b}=q$.
\item $C_{j+1}-C_j=2x_j+q$ for $j<n$.
\item Let $j+1<n$ and $s=\pi^{\mathrm Y}_{(j\,j+1)}$. Then $s^2=1$, $sx_j=x_{j+1}s-1$, and $s$ commutes with
  $x_l$ for $l\ne j,j+1$.
\item $\pi^{\mathrm X}_\sigma$ commutes with $C_0$ and with every $x_j$, and $\pi^{\mathrm Y}_\sigma$ commutes
  with $C_0$.
\item $(\pi_\sigma z)^{\mathrm{sym}}=z^{\mathrm{sym}}$, $(\pi^{\mathrm Y}_\sigma z)^{\mathrm{sym}}=
  (\pi^{\mathrm X}_{\sigma^{-1}}z)^{\mathrm{sym}}$ and $N(z^{\mathrm{sym}})=\sum_\sigma\pi_\sigma z$, the sum over
  all permutations, for every $z\in T_n$; and every function on the multisets of $n$ elements with values in $M$ is
  $z^{\mathrm{sym}}$ for some $z\in T_n$.
\item $\mathrm{cas}(\kappa+\epsilon_r)-\mathrm{cas}(\kappa)-q=2(\kappa_r-r)$ for $\kappa\in\mathbf N^q$ and $r<q$.
\end{enumerate}
\end{lemma}
\leanref{Formalism/PieriBlockIdeals/PieriColumnMinors/Symmetrization.lean}{Formalism.CohomologicalDescent.Symmetrization.permOnTensor_mul}
\leanref{Formalism/PieriBlockIdeals/PieriColumnMinors/Symmetrization.lean}{Formalism.CohomologicalDescent.Symmetrization.permTensor_eq_permSecond_mul_permFirst}
\leanref{Formalism/PieriBlockIdeals/PieriColumnMinors/Symmetrization.lean}{Formalism.CohomologicalDescent.Symmetrization.permFirst_commute_unitSecond}
\leanref{Formalism/PieriBlockIdeals/PieriColumnMinors/Symmetrization.lean}{Formalism.CohomologicalDescent.Symmetrization.permSecond_mul_unitSecond}
\leanref{Formalism/PieriBlockIdeals/PieriColumnMinors/Symmetrization.lean}{Formalism.CohomologicalDescent.Symmetrization.sum_unitSecond_products_eq_swap}
\leanref{Formalism/PieriBlockIdeals/IntegralQuotient.lean}{Formalism.CohomologicalDescent.Symmetrization.CasimirY}
\leanref{Formalism/PieriBlockIdeals/PieriColumnMinors/Symmetrization.lean}{Formalism.CohomologicalDescent.Symmetrization.casimir_on_tensor_holds}
\leanref{Formalism/PieriBlockIdeals/PieriColumnMinors/Symmetrization.lean}{Formalism.CohomologicalDescent.Symmetrization.jucysMurphy_eq_sum_position_products}
\leanref{Formalism/PieriBlockIdeals/IntegralQuotient.lean}{Formalism.CohomologicalDescent.Symmetrization.CasimirStep}
\leanref{Formalism/PieriBlockIdeals/PieriColumnMinors/Symmetrization.lean}{Formalism.CohomologicalDescent.Symmetrization.casimir_recursion_holds}
\leanref{Formalism/PieriBlockIdeals/LemmaDSkeletonUnit.lean}{Formalism.CohomologicalDescent.BinaryFormsExample.CasimirJM}
\leanref{Formalism/PieriBlockIdeals/PieriColumnMinors/BinaryFormsExample.lean}{Formalism.CohomologicalDescent.BinaryFormsExample.casimir_jucys_murphy_holds}
\leanref{Formalism/PieriBlockIdeals/IntegralQuotient.lean}{Formalism.CohomologicalDescent.Symmetrization.HeckeRel}
\leanref{Formalism/PieriBlockIdeals/PieriColumnMinors/Symmetrization.lean}{Formalism.CohomologicalDescent.Symmetrization.hecke_relations_hold}
\leanref{Formalism/PieriBlockIdeals/PieriColumnMinors/Symmetrization.lean}{Formalism.CohomologicalDescent.Symmetrization.commute_jucysMurphy_of_commute}
\leanref{Formalism/PieriBlockIdeals/PieriColumnMinors/Symmetrization.lean}{Formalism.CohomologicalDescent.Symmetrization.toSym_perm_eq}
\leanref{Formalism/PieriBlockIdeals/StatementLC.lean}{Formalism.CohomologicalDescent.LC.ToSymYX}
\leanref{Formalism/PieriBlockIdeals/PieriColumnMinors/PermutationWeightInvariant.lean}{Formalism.CohomologicalDescent.LC.to_symmetric_y_x}
\leanref{Formalism/PieriBlockIdeals/PieriColumnMinors/Symmetrization.lean}{Formalism.CohomologicalDescent.Symmetrization.norm_toSym_eq_sum_perm}
\leanref{Formalism/PieriBlockIdeals/PieriColumnMinors/Symmetrization.lean}{Formalism.CohomologicalDescent.Symmetrization.toSym_surjective}
\leanref{Formalism/PieriBlockIdeals/IntegralQuotient.lean}{Formalism.CohomologicalDescent.Symmetrization.CasimirContent}
\leanref{Formalism/PieriBlockIdeals/PieriColumnMinors/Symmetrization.lean}{Formalism.CohomologicalDescent.Symmetrization.casimir_content_holds}
\begin{proof}
(1) These follow from the definitions: $e^{(s)}_{bb'}$ changes only the entry of $I_2$ at the position $s$,
$\pi^{\mathrm X}_\sigma$ moves only the entries of $I_1$, and $\pi^{\mathrm Y}_\sigma$ moves the entries of $I_2$
by $\sigma$, so that the entry at $s$ before $\pi^{\mathrm Y}_\sigma$ is the entry at $\sigma(s)$ after it.

(2) For $s\ne s'$ and $g\in(I_2\times I_1)^n$,
$\big(\sum_{b,b'}e^{(s)}_{bb'}e^{(s')}_{b'b}z\big)(g)$ is the value of $z$ at $g$ with the entries of $I_2$ at $s$
and $s'$ exchanged, which is $(\pi^{\mathrm Y}_{(s\,s')}z)(g)$. For $s=s'$ the term for $(b,b')$ is $z(g)$ if
$b_s=b$ and $0$ otherwise; summing over $b$ and $b'$ gives $qz(g)$.

(3) Put $\varepsilon_0(b,b')=E_{bb'}$, acting through the values, and $\varepsilon_{s+1}(b,b')=e^{(s)}_{bb'}$ for
$s<n$. For $i\ne i'$, every $\varepsilon_i(b,b')$ commutes with every $\varepsilon_{i'}(c,d)$, by (1) for $i=0$
and as $e^{(s)}$ and $e^{(s')}$ change different positions otherwise. Then
$D_j(b,b')=\sum_{i\le j}\varepsilon_i(b,b')$, and
\[
 x_j=\sum_{a\le j}\sum_{b,b'}\varepsilon_a(b,b')\,\varepsilon_{j+1}(b',b):
\]
the term $a=0$ is $\sum_{b,b'}E_{bb'}e^{(j)}_{b'b}=\sum_{b,b'}e^{(j)}_{bb'}E_{b'b}$, and the term $a=j'+1$ is
$\pi^{\mathrm Y}_{(j'\,j)}$ by (2). Now $D_{j+1}(b,b')=D_j(b,b')+\varepsilon_{j+1}(b,b')$, so
$D_{j+1}(b,b')D_{j+1}(b',b)-D_j(b,b')D_j(b',b)$ is
\[
 D_j(b,b')\varepsilon_{j+1}(b',b)+\varepsilon_{j+1}(b,b')D_j(b',b)+\varepsilon_{j+1}(b,b')\varepsilon_{j+1}(b',b).
\]
As $\varepsilon_{j+1}(b,b')D_j(b',b)=D_j(b',b)\varepsilon_{j+1}(b,b')$, the first two terms have the same sum over
$b,b'$, which is $x_j$; and the sum of the last is $q$ by (2).

(4) By (1), $s^2=\pi^{\mathrm Y}_1=1$, and for every permutation $\sigma$ and $l<n$,
\[
 \pi^{\mathrm Y}_\sigma x_l=\Big(\sum_{b,b'}e^{(\sigma(l))}_{bb'}E_{b'b}
 +\sum_{j'<l}\pi^{\mathrm Y}_{(\sigma(j')\,\sigma(l))}\Big)\pi^{\mathrm Y}_\sigma.
\]
For $\sigma=(j\,j+1)$ and $l=j$ the bracket is $x_{j+1}-\pi^{\mathrm Y}_{(j\,j+1)}$, as the $j'<j$ are fixed by
$\sigma$; so $sx_j=x_{j+1}s-s^2=x_{j+1}s-1$. For $l\ne j,j+1$, $\sigma(l)=l$, and $\sigma$ permutes the
$j'<l$, as $j$ and $j+1$ are both below or both above $l$; so the bracket is $x_l$.

(5) By (1), $\pi^{\mathrm X}_\sigma$ commutes with every $\varepsilon_i(b,b')$ of the proof of (3), hence with
$C_0=\sum_{b,b'}E_{bb'}E_{b'b}$ and with $x_j$ by the formula there; and $\pi^{\mathrm Y}_\sigma$ commutes with
$C_0$ by (1).

(6) The map $g\mapsto g\circ\sigma$ permutes the $g$ whose entries form a given multiset, which gives the first
claim. For the second, $(\pi_\sigma\pi^{\mathrm X}_{\sigma^{-1}}z)(g)$ is the value of $z$ at the $h$ with
$h_s=(b_{\sigma(s)},\alpha_s)$, so $\pi_\sigma\pi^{\mathrm X}_{\sigma^{-1}}=\pi^{\mathrm Y}_\sigma$, and the first claim
applies to $\pi^{\mathrm X}_{\sigma^{-1}}z$. Next, $(N(z^{\mathrm{sym}}))(g)=[g]!\sum_{[g']=[g]}z(g')$, and for each
$g'$ with $[g']=[g]$ the number of permutations $\sigma$ with $g\circ\sigma=g'$ is $[g]!$; so this is
$\sum_\sigma z(g\circ\sigma)$. Finally, a function $y$ on the multisets is $z^{\mathrm{sym}}$ for the $z$ that takes
the value $y(m)$ at one chosen $g$ with $[g]=m$, for each $m$, and $0$ elsewhere.

(7) Only the term $b=r$ of $\mathrm{cas}$ changes, and
$(\kappa_r+1)(\kappa_r+q-2r)-\kappa_r(\kappa_r+q-1-2r)=2\kappa_r+q-2r$.
% Lean: (1) permOnTensor_mul, permSecond_one, permTensor_eq_permSecond_mul_permFirst,
% permSecond_commute_onModule, permFirst_commute_onModule, permFirst_commute_unitSecond,
% permSecond_mul_unitSecond; (2) sum_unitSecond_products_eq_swap and casimir_on_tensor_holds; (3)
% casimir_jucys_murphy_holds: factorGenerator (the book's varepsilon_i), factorGenerator_commute_of_ne,
% depthPart_eq_sum_positionGenerators, jucysMurphy_eq_sum_position_products and casimir_recursion_holds; (4)
% hecke_relations_hold (swap_mul_jucysMurphy_eq_sub_one, swap_adjacent_commute_jucysMurphy_far, through
% permSecond_mul_jucysMurphy); (5) commute_jucysMurphy_of_commute and commute_casimir_zero; (6) toSym_perm_eq,
% to_symmetric_y_x, norm_toSym_eq_sum_perm (card_perms_matching_content_eq_mfact) and toSym_surjective; (7)
% casimir_content_holds. None of these uses the gl relations of the E_bb'.
\end{proof}

\begin{lemma}\label{is-onto-step-exp-bound-short-graded-skeleton-path-eigen}
Let $K$ be a field of characteristic $0$, and $I_1$, $q$, $n$, $M$ and $E_{bb'}$ as in
Definition~\ref{is-onto-step-exp-bound-short-graded-skeleton-jucys}. Let $\kappa_0,\dots,\kappa_n\in\mathbf N^q$ and
$\rho_0,\dots,\rho_{n-1}\in I_2$ with $\kappa_{j+1}=\kappa_j+\epsilon_{\rho_j}$, put
$c_j=\kappa_{j,\rho_j}-\rho_j$, and let $z\in T_n$ with $C_jz=\mathrm{cas}(\kappa_j)z$ for every $j\le n$. Then
$x_jz=c_jz$ for every $j<n$, and for every $a\geq0$ the coefficient of $X^a$ in $A(X)$ maps $z$ to the
coefficient of $X^a$ in $\prod_{j<n}(X-c_j)$ times $z$.
\end{lemma}
\leanref{Formalism/PieriBlockIdeals/LemmaDSkeletonUnit.lean}{Formalism.CohomologicalDescent.BinaryFormsExample.pathPoly}
\leanref{Formalism/PieriBlockIdeals/LemmaDSkeletonUnit.lean}{Formalism.CohomologicalDescent.BinaryFormsExample.PathEigen}
\leanref{Formalism/PieriBlockIdeals/PieriColumnMinors/BinaryFormsExample.lean}{Formalism.CohomologicalDescent.BinaryFormsExample.path_eigen_holds}
\leanref{Formalism/PieriBlockIdeals/PieriColumnMinors/BinaryFormsExample.lean}{Formalism.CohomologicalDescent.BinaryFormsExample.coeff_list_prod_X_sub_C_apply_eigen}
\begin{proof}
By Lemma~\ref{is-onto-step-exp-bound-short-graded-skeleton-jucys-relations}(3) and~(7),
\[
 2x_jz=(C_{j+1}-C_j-q)z=\big(\mathrm{cas}(\kappa_j+\epsilon_{\rho_j})-\mathrm{cas}(\kappa_j)-q\big)z=2c_jz,
\]
and $2$ is invertible in $K$. For the coefficients, induct on $n$. Write $A(X)=(X-x_0)Q(X)$ with
$Q(X)=(X-x_1)\cdots(X-x_{n-1})$, and $\prod_j(X-c_j)=(X-c_0)Q'(X)$. The coefficient of $X^a$ in $A(X)$ is
$Q_{a-1}-x_0Q_a$, with $Q_{-1}=0$. By induction $Q_{a-1}z=Q'_{a-1}z$ and $Q_az=Q'_az$, so it maps $z$ to
$Q'_{a-1}z-Q'_ac_0z$, the coefficient of $X^a$ in $(X-c_0)Q'(X)$ times $z$.
% Lean: path_eigen_holds: casimir_jucys_murphy_holds and casimir_content_holds for x_j, then
% coeff_list_prod_X_sub_C_apply_eigen (induction on n with coeff_zero_X_sub_C_mul_alt and
% coeff_succ_X_sub_C_mul_alt) for the coefficients of AApoly against pathPoly.
\end{proof}

\begin{lemma}\label{is-onto-step-exp-bound-short-graded-skeleton-commute}
Let $K$, $I_1$, $q$, $n$, $M$ and $E_{bb'}$ be as in
Definition~\ref{is-onto-step-exp-bound-short-graded-skeleton-jucys}, and suppose that
$E_{ab}E_{cd}-E_{cd}E_{ab}=\delta_{bc}E_{ad}-\delta_{da}E_{cb}$ for all $a,b,c,d\in I_2$. Write $A_a$ for the
coefficient of $X^a$ in $A(X)$.
\begin{enumerate}
\item The $x_j$ commute with each other, and so do the $A_a$.
\item $A(t)=t^n+\sum_{a<n}t^aA_a$ for $t\in K$. For every permutation $\sigma$ of $\{0,\dots,n-1\}$, the operators
  $\pi^{\mathrm Y}_\sigma$, $\pi^{\mathrm X}_\sigma$ and $\pi_\sigma$ commute with every $A_a$ and with every $A(t)$.
\item Suppose that $K$ has characteristic $0$. For $a<n$, let $c_a$ be the endomorphism
  $y\mapsto(n!)^{-1}(A_aNy)^{\mathrm{sym}}$ of the space of functions on the multisets of $n$ elements of
  $I_2\times I_1$ with values in $M$. Then $c_a(z^{\mathrm{sym}})=(A_az)^{\mathrm{sym}}$ for $z\in T_n$, the $c_a$
  commute, and $(A(t)z)^{\mathrm{sym}}=\big(t^n+\sum_{a<n}t^ac_a\big)(z^{\mathrm{sym}})$ for $t\in K$.
\item In (3), let $O$ be a commutative ring with a homomorphism $O\to K$, acting on $M$ through it, and
  $L(m)\subset M$ an $O$-submodule for each multiset $m$ of $n$ elements. If every $x_j$ maps the set of the
  $z\in T_n$ with $z(g)\in L([g])$ for all $g$ into itself, then every $c_a$ maps the set of the functions $y$
  with $y(m)\in L(m)$ for all $m$ into itself.
\end{enumerate}
\end{lemma}
\leanref{Formalism/PieriBlockIdeals/IntegralQuotient.lean}{Formalism.CohomologicalDescent.Symmetrization.JMCommAbstract}
\leanref{Formalism/PieriBlockIdeals/PieriColumnMinors/Symmetrization.lean}{Formalism.CohomologicalDescent.Symmetrization.jucys_murphy_commute_holds}
\leanref{Formalism/PieriBlockIdeals/PieriColumnMinors/Symmetrization.lean}{Formalism.CohomologicalDescent.Symmetrization.casimir_pair_commute_sum}
\leanref{Formalism/PieriBlockIdeals/IntegralQuotient.lean}{Formalism.CohomologicalDescent.Symmetrization.XComm}
\leanref{Formalism/PieriBlockIdeals/PieriColumnMinors/Symmetrization.lean}{Formalism.CohomologicalDescent.Symmetrization.XComm_of}
\leanref{Formalism/PieriBlockIdeals/IntegralQuotient.lean}{Formalism.CohomologicalDescent.Symmetrization.SymCommAbstract}
\leanref{Formalism/PieriBlockIdeals/PieriColumnMinors/Symmetrization.lean}{Formalism.CohomologicalDescent.Symmetrization.symmetric_commute_abstract_holds}
\leanref{Formalism/PieriBlockIdeals/IntegralQuotient.lean}{Formalism.CohomologicalDescent.Symmetrization.AAComm}
\leanref{Formalism/PieriBlockIdeals/PieriColumnMinors/Symmetrization.lean}{Formalism.CohomologicalDescent.Symmetrization.AAComm_of}
\leanref{Formalism/PieriBlockIdeals/PieriColumnMinors/Symmetrization.lean}{Formalism.CohomologicalDescent.Symmetrization.jucysMurphyProduct_expand_coeff}
\leanref{Formalism/PieriBlockIdeals/IntegralQuotient.lean}{Formalism.CohomologicalDescent.Symmetrization.CoeffOps}
\leanref{Formalism/PieriBlockIdeals/PieriColumnMinors/Symmetrization.lean}{Formalism.CohomologicalDescent.Symmetrization.coefficient_operators_hold}
\begin{proof}
(1) Let $\varepsilon_i$ be as in the proof of
Lemma~\ref{is-onto-step-exp-bound-short-graded-skeleton-jucys-relations}(3). Each family
$(\varepsilon_i(b,b'))_{b,b'}$ satisfies the relations of the statement: $\varepsilon_0$ by assumption, and
$\varepsilon_{s+1}$ as $e^{(s)}_{ab}e^{(s)}_{cd}=\delta_{bc}e^{(s)}_{ad}$. For $i\ne i'$ the two families commute.
Put $P(i,j)=\sum_{b,b'}\varepsilon_i(b,b')\varepsilon_j(b',b)$, so that $x_j=\sum_{a\le j}P(a,j+1)$. For $i\ne j$, a
direct computation from these relations shows that $P(i,j)$ commutes with $\varepsilon_i(c,d)+\varepsilon_j(c,d)$
for all $c,d$: of the four sums that arise, two cancel against the other two. Hence, for $l\ne i,j$, $P(i,j)$
commutes with $P(i,l)+P(j,l)=\sum_{c,d}\big(\varepsilon_i(c,d)+\varepsilon_j(c,d)\big)\varepsilon_l(d,c)$, and it
commutes with $P(a,b)$ when $a,b\notin\{i,j\}$. Now let $i<j$ and $a<i$. Then $P(a,i)$ commutes with
$P(a,j)+P(i,j)$ and with $P(b,j)$ for $b<j$, $b\ne a,i$, hence with $\sum_{b<j}P(b,j)$. So $\sum_{a<i}P(a,i)$ commutes
with $\sum_{b<j}P(b,j)$, and for $i+1$ and $j+1$ in place of $i$ and $j$ this says that $x_i$ and $x_j$ commute.
Hence $x_j$ commutes with $A(X)$; taking coefficients, every $A_a$ commutes with every $x_j$, so $A_a$ commutes
with $A(X)$, and $A_a$ with $A_b$.

(2) By induction on $n$, $A(X)$ is monic of degree $n$, and $A(t)$ is obtained from it by substituting the
central element $t$; this gives the formula for $A(t)$, and it suffices to show that the three operators commute
with $A(X)$. By Lemma~\ref{is-onto-step-exp-bound-short-graded-skeleton-jucys-relations}(5), $\pi^{\mathrm X}_\sigma$
commutes with every $x_j$, hence with $A(X)$. As $\pi^{\mathrm Y}$ is multiplicative, it suffices to treat
$s=\pi^{\mathrm Y}_{(j\,j+1)}$. By Lemma~\ref{is-onto-step-exp-bound-short-graded-skeleton-jucys-relations}(4), $s$
commutes with $X-x_l$ for $l\ne j,j+1$; and $sx_j=x_{j+1}s-1$ and $s^2=1$ give $sx_{j+1}=x_js+1$. So
$s(X-x_j)=(X-x_{j+1})s+1$ and $s(X-x_{j+1})=(X-x_j)s-1$, and as $x_j$ and $x_{j+1}$ commute,
\[
 s(X-x_j)(X-x_{j+1})=(X-x_{j+1})\big((X-x_j)s-1\big)+(X-x_{j+1})=(X-x_j)(X-x_{j+1})s.
\]
So $s$ commutes with $A(X)$. Finally $\pi_\sigma=\pi^{\mathrm Y}_\sigma\pi^{\mathrm X}_\sigma$.

(3) By Lemma~\ref{is-onto-step-exp-bound-short-graded-skeleton-jucys-relations}(6),
$N(z^{\mathrm{sym}})=\sum_\sigma\pi_\sigma z$, so by (2) and that lemma
$c_a(z^{\mathrm{sym}})=(n!)^{-1}\sum_\sigma(\pi_\sigma A_az)^{\mathrm{sym}}=(A_az)^{\mathrm{sym}}$. Every function on
the multisets is some $z^{\mathrm{sym}}$, so $c_ac_b(z^{\mathrm{sym}})=(A_aA_bz)^{\mathrm{sym}}$ and (1) shows that the
$c_a$ commute. The last formula follows from (2).

(4) The endomorphisms of $T_n$ that map the set $Z_L$ of the $z$ with $z(g)\in L([g])$ for all $g$ into itself
form a subring, which contains every $x_j$. By induction on $n$, the coefficients of $(X-x_0)\cdots(X-x_{n-1})$
are sums of products of the $-x_j$ and $1$, so every $A_a$ lies in it. Let $y(m)\in L(m)$ for every $m$, and let $z$
take the value $y(m)$ at one chosen $g$ with $[g]=m$, for each $m$, and $0$ elsewhere. Then $z\in Z_L$ and
$z^{\mathrm{sym}}=y$, so $c_a(y)=(A_az)^{\mathrm{sym}}$, whose value at $m$ is a sum of values of $A_az$ at $g$ with
$[g]=m$; these lie in $L(m)$.
% Lean: (1) XComm_of jucys_murphy_commute_holds (glRel_factorGenerator, factorGenerator_commute_of_ne,
% casimir_pair_commute_sum, transposition_commute_sum_shared, commute_partial_jucysMurphy_sums) and, for the
% coefficients, commute_C_jucysMurphyPoly and commute_coeff_of_commute_C; (2) jucysMurphyProduct_expand_coeff
% (prod_X_sub_C_coeff_top), AAComm_of symmetric_commute_abstract_holds jucys_murphy_commute_holds
% (commute_C_permSecond_jucysMurphyPoly by Submonoid.dense_induction over the adjacent swaps,
% commute_list_prod_of_commute_adjacent, evalAtCentral); (3) and (4) coefficient_operators_hold
% (descendToCoinvariants, descendToCoinvariants_toSym, coeff_prod_X_sub_C_mem_subring with stabilizerSubring).
\end{proof}

\begin{definition}\label{is-onto-step-exp-bound-short-graded-skeleton-maps}
Let $K$ be a field, $I_1$ a finite set, $q\geq0$, $I_2=\{0,\dots,q-1\}$, and $(V,E,V^+,w)$ a block module over
$K$ for $I_1$ and $I_2$, with $f_v$, $F_m$, $m!$ and $T_n$ as in
Definition~\ref{is-onto-step-exp-bound-short-graded-block}. The operators of
Definition~\ref{is-onto-step-exp-bound-short-graded-skeleton-jucys} on $T_n$ are those for $M=V^+$ and the
restrictions to $V^+$ of the $E_{bb'}$, $b,b'\in I_2$. Let $n\geq0$.
\begin{enumerate}
\item For $v=(b,\alpha)\in I_2\times I_1$ put $e_v=E_{\alpha b}$. For $g\in(I_2\times I_1)^n$ put
  $F_g=f_{g_0}f_{g_1}\cdots f_{g_{n-1}}$ and $E_g=e_{g_0}e_{g_1}\cdots e_{g_{n-1}}$. For a multiset $m$ of $n$
  elements of $I_2\times I_1$, $E_m$ is the product of the $e_v$ over the elements $v$ of $m$, counted with
  multiplicity.
\item $\mu_n\colon T_n\to V$ is the map $\mu_n(z)=\sum_gF_g\,z(g)$. For $v\in V$, $\iota_n(v)$ is the function on
  the multisets $m$ of $n$ elements with $\iota_n(v)(m)=(m!)^{-1}E_mv$. If $K$ has characteristic $0$,
  $\mathrm{act}_n(y)=(n!)^{-1}\mu_n(Ny)$ for every function $y$ on these multisets with values in $V^+$.
\item For $j\le n$ and $v\in I_2\times I_1$, $\mathrm{ct}_{j,v}\colon T_{n+1}\to T_n$ is the map with
  $(\mathrm{ct}_{j,v}z)(g)=z(g')$, where $g'$ is $g$ with $v$ inserted at the position $j$.
\item For $\alpha,\beta\in I_1$ and $s<n$, $u^{(s)}_{\alpha\beta}$ is the endomorphism of $T_n$ with
  $(u^{(s)}_{\alpha\beta}z)(g)=z(g'')$ if $g_s=(b,\alpha)$ for some $b$, where $g''$ is $g$ with $g_s$ replaced by
  $(b,\beta)$, and $(u^{(s)}_{\alpha\beta}z)(g)=0$ otherwise.
\item For a multiset $m$ of elements of $I_2\times I_1$, $\mathrm{wt}(m)\in\mathbf Z^{I_1\sqcup I_2}$ has entry at
  $\alpha\in I_1$ the number of elements of $m$ with second entry $\alpha$, and entry at $b\in I_2$ minus the number
  of elements of $m$ with first entry $b$.
\end{enumerate}
\end{definition}
\leanref{Formalism/PieriBlockIdeals/IntegralQuotient.lean}{Formalism.CohomologicalDescent.Symmetrization.HidaSetting.e}
\leanref{Formalism/PieriBlockIdeals/IntegralQuotient.lean}{Formalism.CohomologicalDescent.Symmetrization.HidaSetting.Fw}
\leanref{Formalism/PieriBlockIdeals/IntegralQuotient.lean}{Formalism.CohomologicalDescent.Symmetrization.HidaSetting.Ew}
\leanref{Formalism/PieriBlockIdeals/IntegralQuotient.lean}{Formalism.CohomologicalDescent.Symmetrization.HidaSetting.Es}
\leanref{Formalism/PieriBlockIdeals/IntegralQuotient.lean}{Formalism.CohomologicalDescent.Symmetrization.HidaSetting.EW}
\leanref{Formalism/PieriBlockIdeals/IntegralQuotient.lean}{Formalism.CohomologicalDescent.Symmetrization.HidaSetting.mu}
\leanref{Formalism/PieriBlockIdeals/IntegralQuotient.lean}{Formalism.CohomologicalDescent.Symmetrization.HidaSetting.iota}
\leanref{Formalism/PieriBlockIdeals/IntegralQuotient.lean}{Formalism.CohomologicalDescent.Symmetrization.HidaSetting.act}
\leanref{Formalism/PieriBlockIdeals/IntegralQuotient.lean}{Formalism.CohomologicalDescent.Symmetrization.contr}
\leanref{Formalism/PieriBlockIdeals/IntegralQuotient.lean}{Formalism.CohomologicalDescent.Symmetrization.unitX}
\leanref{Formalism/PieriBlockIdeals/IntegralQuotient.lean}{Formalism.CohomologicalDescent.Symmetrization.beta}
% Lean: HidaSetting.e, Fw, Ew, Fs, Es, EW (the E_bb' restricted to W), mu, iota and act; contr k j v; unitX k j
% alpha beta (the book's u^(j)); beta m (the book's wt(m)); diagY and diagX are the operators of the lemma below.

\begin{lemma}\label{is-onto-step-exp-bound-short-graded-skeleton-mu}
Let $K$, $I_1$, $q$ and $(V,E,V^+,w)$ be as in Definition~\ref{is-onto-step-exp-bound-short-graded-skeleton-maps},
$p=\#I_1$, and $n\geq0$.
\begin{enumerate}
\item $f_vf_{v'}=f_{v'}f_v$ and $e_ve_{v'}=e_{v'}e_v$ for all $v,v'$. So $F_g=F_{[g]}$ and $E_g=E_{[g]}$, and
  $\mu_n(\pi_\sigma z)=\mu_n(z)$ for every permutation $\sigma$ and every $z\in T_n$.
\item $\mu_0(T_0)=V^+$, and $e_v\mu_0(z)=0$ for $z\in T_0$. For $z\in T_{n+1}$ and $j\le n$,
  $\mu_{n+1}(z)=\sum_vf_v\,\mu_n(\mathrm{ct}_{j,v}z)$.
\item For $z\in T_n$, $c,d\in I_2$ and $\alpha,\beta\in I_1$,
  \[
   E_{cd}\mu_n(z)=\mu_n\Big(E_{cd}z+\sum_{s<n}e^{(s)}_{cd}z\Big),\qquad
   E_{\alpha\beta}\mu_n(z)=\mu_n\Big(\delta_{\alpha\beta}wz-\sum_{s<n}u^{(s)}_{\beta\alpha}z\Big),
  \]
  where $E_{cd}$ acts on $T_n$ through the values. In particular
  $\sum_{\alpha\in I_1}E_{\alpha\alpha}\mu_n(z)=(pw-n)\mu_n(z)$.
\item For $z\in T_{n+1}$ and $v\in I_2\times I_1$, $e_v\mu_{n+1}(z)=\mu_n(\rho_v(z))$, where
  $\Omega_j=\sum_{b,b'}e^{(j)}_{bb'}E_{b'b}$ and
  \[
   \rho_v(z)=\sum_{j\le n}\mathrm{ct}_{j,v}(wz-\Omega_jz)
   -\sum_{j<l\le n}\big(\mathrm{ct}_{j,v}+\mathrm{ct}_{l,v}\big)(\pi^{\mathrm Y}_{(j\,l)}z).
  \]
  If $\pi_\sigma z=z$ for every $\sigma$, then $\rho_v(z)=(n+1)\,\mathrm{ct}_{n,v}(wz-x_nz)$.
\item Suppose that every subspace of $V$ that contains $V^+$ and is stable under every $E_{cd}$ is $V$. Then
  $V=\sum_{n\geq0}\mu_n(T_n)$. If moreover $K$ has characteristic $0$ and $w$ is an integer, then every $v\in V_\nu$
  with $\sum_{\alpha\in I_1}\nu_\alpha+n=pw$ lies in $\mu_n(T_n)$.
\end{enumerate}
\end{lemma}
\leanref{Formalism/PieriBlockIdeals/PieriColumnMinors/Symmetrization.lean}{Formalism.CohomologicalDescent.Symmetrization.word_products_of_content}
\leanref{Formalism/PieriBlockIdeals/PieriColumnMinors/Symmetrization.lean}{Formalism.CohomologicalDescent.Symmetrization.symMap_succ_eq_sum_contract}
\leanref{Formalism/PieriBlockIdeals/IntegralQuotient.lean}{Formalism.CohomologicalDescent.Symmetrization.diagY}
\leanref{Formalism/PieriBlockIdeals/IntegralQuotient.lean}{Formalism.CohomologicalDescent.Symmetrization.diagX}
\leanref{Formalism/PieriBlockIdeals/IntegralQuotient.lean}{Formalism.CohomologicalDescent.Symmetrization.MuEquivariant}
\leanref{Formalism/PieriBlockIdeals/PieriColumnMinors/Symmetrization.lean}{Formalism.CohomologicalDescent.Symmetrization.symMap_equivariant_holds}
\leanref{Formalism/PieriBlockIdeals/LemmaDSkeletonUnit.lean}{Formalism.CohomologicalDescent.BinaryFormsExample.DepthEigen}
\leanref{Formalism/PieriBlockIdeals/PieriColumnMinors/BinaryFormsExample.lean}{Formalism.CohomologicalDescent.BinaryFormsExample.depth_eigen_holds}
\leanref{Formalism/PieriBlockIdeals/IntegralQuotient.lean}{Formalism.CohomologicalDescent.Symmetrization.OneStep}
\leanref{Formalism/PieriBlockIdeals/PieriColumnMinors/Symmetrization.lean}{Formalism.CohomologicalDescent.Symmetrization.raising_symMap_eq_one_level_down}
\leanref{Formalism/PieriBlockIdeals/IntegralQuotient.lean}{Formalism.CohomologicalDescent.Symmetrization.OneStepInv}
\leanref{Formalism/PieriBlockIdeals/PieriColumnMinors/Symmetrization.lean}{Formalism.CohomologicalDescent.Symmetrization.one_step_invariance_holds}
\leanref{Formalism/PieriBlockIdeals/LemmaDSkeletonUnit.lean}{Formalism.CohomologicalDescent.BinaryFormsExample.DepthSpan}
\leanref{Formalism/PieriBlockIdeals/PieriColumnMinors/BinaryFormsExample.lean}{Formalism.CohomologicalDescent.BinaryFormsExample.depth_span_holds}
\leanref{Formalism/PieriBlockIdeals/LemmaDSkeletonUnit.lean}{Formalism.CohomologicalDescent.BinaryFormsExample.DepthOfWeight}
\leanref{Formalism/PieriBlockIdeals/PieriColumnMinors/BinaryFormsExample.lean}{Formalism.CohomologicalDescent.BinaryFormsExample.depth_of_weight_holds}
\begin{proof}
(1) For $v=(b,\alpha)$ and $v'=(b',\alpha')$, the relations give
$E_{b\alpha}E_{b'\alpha'}-E_{b'\alpha'}E_{b\alpha}=\delta_{\alpha b'}E_{b\alpha'}-\delta_{\alpha'b}E_{b'\alpha}=0$, as $I_1$
and $I_2$ are disjoint; likewise for the $e_v$. So $F_g$ and $E_g$ depend only on $[g]$, and
$\mu_n(\pi_\sigma z)=\sum_gF_gz(g\circ\sigma)=\sum_gF_{g\circ\sigma^{-1}}z(g)=\mu_n(z)$.

(2) An element of $T_0$ is one value $z(\varnothing)\in V^+$, and $\mu_0(z)=z(\varnothing)$; and $e_v=E_{\alpha b}$
kills $V^+$. For the last formula, $g\mapsto(g_j,g^-)$, with $g^-$ the sequence $g$ without its entry at $j$, is a
bijection, and $F_g=f_{g_j}F_{g^-}$ by (1).

(3) Induct on $n$, using (2) with $j=0$. For $n=0$ the formulas say that $E_{cd}u=E_{cd}u$ and
$E_{\alpha\beta}u=\delta_{\alpha\beta}wu$ for $u\in V^+$. For the step, the relations give
$E_{cd}f_{(b,\gamma)}=f_{(b,\gamma)}E_{cd}+\delta_{db}f_{(c,\gamma)}$ and
$E_{\alpha\beta}f_{(b,\gamma)}=f_{(b,\gamma)}E_{\alpha\beta}-\delta_{\gamma\alpha}f_{(b,\beta)}$. Apply these to
$\mu_{n+1}(z)=\sum_vf_v\mu_n(\mathrm{ct}_{0,v}z)$, use the induction hypothesis, and compare with
$\mathrm{ct}_{0,v}$ of the operators on $T_{n+1}$: for $v=(b,\gamma)$,
$\mathrm{ct}_{0,v}\big(E_{cd}z+\sum_{s\le n}e^{(s)}_{cd}z\big)=E_{cd}\mathrm{ct}_{0,v}z+\sum_{s<n}e^{(s)}_{cd}\mathrm{ct}_{0,v}z
+\delta_{bc}\mathrm{ct}_{0,(d,\gamma)}z$, and similarly for the $u^{(s)}$. Reindexing the extra sums over $v$ gives
the claim. Finally $\sum_\alpha u^{(s)}_{\alpha\alpha}=1$ for every $s$.

(4) Write $v=(b_0,\alpha_0)$ and induct on $n$. By (2) with $j=0$ and the relation
$e_vf_{(b,\gamma)}=f_{(b,\gamma)}e_v+\delta_{b_0b}E_{\alpha_0\gamma}-\delta_{\gamma\alpha_0}E_{bb_0}$,
\[
 e_v\mu_{n+1}(z)=\sum_{v'}f_{v'}e_v\mu_n(\mathrm{ct}_{0,v'}z)+\mu_n(\Lambda),
\]
where, by (3),
$\Lambda=\sum_\gamma\big(\delta_{\alpha_0\gamma}w-\sum_su^{(s)}_{\gamma\alpha_0}\big)\mathrm{ct}_{0,(b_0,\gamma)}z
-\sum_b\big(E_{bb_0}+\sum_se^{(s)}_{bb_0}\big)\mathrm{ct}_{0,(b,\alpha_0)}z$. Split $\rho_v(z)$ into the terms with
$j=0$, which form $\rho'_v(z)$, and the others, which form $\rho''_v(z)$. Evaluating at each $g$ gives
$\Lambda=\mathrm{ct}_{0,v}(wz-\Omega_0z)-\sum_{t<n}\varphi_t-\sum_{t<n}\mathrm{ct}_{0,v}(\pi^{\mathrm Y}_{(0\,t+1)}z)$, where
$\varphi_t(g)$ is the value of $z$ at the sequence with first entry $(b_0,\alpha_t)$ followed by $g$ with $g_t$
replaced by $(b_t,\alpha_0)$. Expanding $\mu_n$ by (2) at the position $t$ for $\varphi_t$, and at the position $0$
for $\mathrm{ct}_{t+1,v}(\pi^{\mathrm Y}_{(0\,t+1)}z)$, the two expansions agree term by term; so
$\mu_n(\varphi_t)=\mu_n(\mathrm{ct}_{t+1,v}(\pi^{\mathrm Y}_{(0\,t+1)}z))$, and $\mu_n(\Lambda)=\mu_n(\rho'_v(z))$. If
$n=0$, then $\rho''_v(z)=0$ and $e_v\mu_0(\mathrm{ct}_{0,v'}z)=0$ by (2). If $n\geq1$, then
$\mathrm{ct}_{0,v'}\rho''_v(z)=\rho_v(\mathrm{ct}_{0,v'}z)$, as $\mathrm{ct}_{0,v'}$ commutes with the contractions,
the $\Omega$ and the $\pi^{\mathrm Y}$ after the shift of positions by one; by induction and (2),
$\sum_{v'}f_{v'}e_v\mu_n(\mathrm{ct}_{0,v'}z)=\sum_{v'}f_{v'}\mu_{n-1}(\mathrm{ct}_{0,v'}\rho''_v(z))=\mu_n(\rho''_v(z))$.
In both cases $e_v\mu_{n+1}(z)=\mu_n(\rho'_v(z)+\rho''_v(z))$.

If $z$ is invariant under every $\pi_\sigma$, the value of $z$ at a sequence with $v$ inserted at the position $j$
does not depend on $j$. So $\mathrm{ct}_{j,v}(wz-\Omega_jz)=\mathrm{ct}_{n,v}(wz-\Omega_nz)$ for every $j$, and each
of the $n+1$ positions $s$ contributes $\sum_{t<n}\mathrm{ct}_{n,v}(\pi^{\mathrm Y}_{(t\,n)}z)$ to the second sum.
As $x_n=\Omega_n+\sum_{t<n}\pi^{\mathrm Y}_{(t\,n)}$, $\rho_v(z)=(n+1)\,\mathrm{ct}_{n,v}(wz-x_nz)$.

(5) The subspace $U=\sum_n\mu_n(T_n)$ contains $V^+=\mu_0(T_0)$. It is stable under the $E_{cd}$: under
$E_{b\alpha}=f_{(b,\alpha)}$ as $f_v\mu_n(z)=\mu_{n+1}(z')$ with $z'(g)=z(g^-)$ if $g_0=v$ and $0$ otherwise, by (2);
under $E_{\alpha b}=e_{(b,\alpha)}$ by (4) and (2); and under $E_{cd}$ and $E_{\alpha\beta}$ by (3). So $U=V$. By
(3), $\mu_n(T_n)$ lies in the eigenspace of $\sum_\alpha E_{\alpha\alpha}$ for the eigenvalue $pw-n$, and these
eigenvalues are distinct in characteristic $0$. The vector $v$ lies in the eigenspace for
$\sum_\alpha\nu_\alpha=pw-n$. Write $v=\sum_{n'}v_{n'}$ with $v_{n'}\in\mu_{n'}(T_{n'})$; then $v-v_n$ lies in that
eigenspace and in the sum of the eigenspaces for the other eigenvalues, so it is $0$, and $v=v_n$.
% Lean: (1) lowering_operators_commute_pair, raising_operators_commute, word_products_of_content,
% symMap_perm_invariant; (2) symMap_zero_eq_default, symMap_succ_eq_sum_contract (wordProduct_insertNth) and
% depth_span_holds; (3) symMap_diagSecond_intertwines and symMap_diagFirst_intertwines (with
% diagSecond_lowering_commutator, diagFirst_lowering_commutator, contract_zero_diagSecond,
% contract_zero_diagFirst), then depth_eigen_holds (sum_diagFirst_eq_smul); (4) raising_symMap_eq_one_level_down:
% raising_symMap_succ_eq (leviRemainder), raisingRemainder_eq_first_add_tail, contraction_remainder_eq
% (flippedLetterTerm), symMap_swap_contraction_eq, contract_zero_tailRemainder, and for invariant z
% raisingRemainder_of_perm_invariant (apply_insertNth_eq_of_perm_invariant, sum_ordered_pairs_eq_sum_succAbove,
% jucysMurphy_last_eq), which gives one_step_invariance_holds; (5) depth_span_holds (lowering_symMap_eq_symMap_succ)
% and depth_of_weight_holds (mem_eigen_piece_of_mem_iSup, Module.End.eigenspaces_iSupIndep).
\end{proof}

\begin{lemma}\label{is-onto-step-exp-bound-short-graded-skeleton-iota}
Let $K$ be a field of characteristic $0$, $I_1$, $q$ and $(V,E,V^+,w)$ as in
Definition~\ref{is-onto-step-exp-bound-short-graded-skeleton-maps}, and $n\geq0$; $A(t)$ is the operator of
Definition~\ref{is-onto-step-exp-bound-short-graded-skeleton-jucys}(2) for $M=V^+$.
\begin{enumerate}
\item If $z\in T_n$ and $\pi_\sigma z=z$ for every $\sigma$, then $E_h\mu_n(z)=n!\,(A(w)z)(h)$ for every
  $h\in(I_2\times I_1)^n$.
\item $\iota_n(\mu_n(z))(m)=(A(w)z)^{\mathrm{sym}}(m)$ for every $z\in T_n$ and every multiset $m$ of $n$ elements.
\item $\mathrm{act}_n(z^{\mathrm{sym}})=\mu_n(z)$ for $z\in T_n$, and $\mathrm{act}_n(y)=F_mu$ if $y$ takes the
  value $u$ at $m$ and $0$ elsewhere.
\item Let $O$ be a commutative ring with a homomorphism $O\to K$, acting on $V$ through it, and $\Delta_L\subset V^+$
  an $O$-submodule. If $y(m)\in\Delta_L$ for every $m$, then $\mathrm{act}_n(y)\in\Delta_n(\Delta_L)$.
\end{enumerate}
\end{lemma}
\leanref{Formalism/PieriBlockIdeals/PieriColumnMinors/Symmetrization.lean}{Formalism.CohomologicalDescent.Symmetrization.generators_satisfy_gl_relations}
\leanref{Formalism/PieriBlockIdeals/IntegralQuotient.lean}{Formalism.CohomologicalDescent.Symmetrization.Iteration}
\leanref{Formalism/PieriBlockIdeals/PieriColumnMinors/Symmetrization.lean}{Formalism.CohomologicalDescent.Symmetrization.Iteration_of}
\leanref{Formalism/PieriBlockIdeals/IntegralQuotient.lean}{Formalism.CohomologicalDescent.Symmetrization.IotaMu}
\leanref{Formalism/PieriBlockIdeals/PieriColumnMinors/Symmetrization.lean}{Formalism.CohomologicalDescent.Symmetrization.IotaMu_of}
\leanref{Formalism/PieriBlockIdeals/IntegralQuotient.lean}{Formalism.CohomologicalDescent.Symmetrization.ActSpec}
\leanref{Formalism/PieriBlockIdeals/PieriColumnMinors/Symmetrization.lean}{Formalism.CohomologicalDescent.Symmetrization.ActSpec_of}
\leanref{Formalism/PieriBlockIdeals/IntegralQuotient.lean}{Formalism.CohomologicalDescent.Symmetrization.ActLattice}
\leanref{Formalism/PieriBlockIdeals/PieriColumnMinors/Symmetrization.lean}{Formalism.CohomologicalDescent.Symmetrization.ActLattice_of}
\begin{proof}
The relations of the block module show that the restrictions of the $E_{bb'}$ to $V^+$ satisfy the relations of
Lemma~\ref{is-onto-step-exp-bound-short-graded-skeleton-commute}, which therefore applies to $M=V^+$.

(1) Induct on $n$; for $n=0$ both sides are $z(\varnothing)$. Let $z\in T_{n+1}$ be invariant,
$h\in(I_2\times I_1)^{n+1}$, $v=h_n$ and $h'=(h_0,\dots,h_{n-1})$, so that $E_h=E_{h'}e_v$. Put $X=A(w)z$ and
$R=\mathrm{ct}_{n,v}(wz-x_nz)$. By Lemma~\ref{is-onto-step-exp-bound-short-graded-skeleton-mu}(4),
$e_v\mu_{n+1}(z)=(n+1)\mu_n(R)$. By Lemma~\ref{is-onto-step-exp-bound-short-graded-skeleton-commute}(2), $X$ is
invariant, hence so is $\mathrm{ct}_{n,v}X$. The contraction $\mathrm{ct}_{n,v}$ carries $x_j$ on $T_{n+1}$ to $x_j$
on $T_n$ for $j<n$, and $A(w)$ on $T_{n+1}$ is the product for $j<n$ followed by $w-x_n$; so $A(w)R=\mathrm{ct}_{n,v}X$
on $T_n$. Put $Y=\sum_\sigma\pi_\sigma R$, which is invariant. By
Lemma~\ref{is-onto-step-exp-bound-short-graded-skeleton-mu}(1), $\mu_n(Y)=n!\,\mu_n(R)$, and by
Lemma~\ref{is-onto-step-exp-bound-short-graded-skeleton-commute}(2),
$A(w)Y=\sum_\sigma\pi_\sigma\mathrm{ct}_{n,v}X=n!\,\mathrm{ct}_{n,v}X$. By induction,
$n!\,E_{h'}\mu_n(R)=E_{h'}\mu_n(Y)=n!\,(A(w)Y)(h')=(n!)^2X(h)$. So
$E_h\mu_{n+1}(z)=(n+1)E_{h'}\mu_n(R)=(n+1)!\,X(h)$.

(2) Let $h$ have $[h]=m$, so that $E_m=E_h$ by
Lemma~\ref{is-onto-step-exp-bound-short-graded-skeleton-mu}(1). Put $Z=N(z^{\mathrm{sym}})=\sum_\sigma\pi_\sigma z$
(Lemma~\ref{is-onto-step-exp-bound-short-graded-skeleton-jucys-relations}(6)), which is invariant, with
$\mu_n(Z)=n!\,\mu_n(z)$. By (1) and Lemma~\ref{is-onto-step-exp-bound-short-graded-skeleton-commute}(2),
$n!\,E_h\mu_n(z)=n!\,(A(w)Z)(h)$, and $A(w)Z=\sum_\sigma\pi_\sigma A(w)z=N\big((A(w)z)^{\mathrm{sym}}\big)$, whose value
at $h$ is $m!\,(A(w)z)^{\mathrm{sym}}(m)$. Dividing by $n!\,m!$ gives the claim.

(3) By Lemma~\ref{is-onto-step-exp-bound-short-graded-skeleton-jucys-relations}(6) and
Lemma~\ref{is-onto-step-exp-bound-short-graded-skeleton-mu}(1),
$\mu_n(N(z^{\mathrm{sym}}))=\sum_\sigma\mu_n(\pi_\sigma z)=n!\,\mu_n(z)$. For the second claim, $y=z^{\mathrm{sym}}$ for the
$z$ with value $u$ at one $g$ with $[g]=m$ and $0$ elsewhere, so $\mathrm{act}_n(y)=\mu_n(z)=F_gu=F_mu$.

(4) By (3), $\mathrm{act}_n(y)=\sum_mF_my(m)=\sum_mm!\,\big((m!)^{-1}F_my(m)\big)$, a sum of elements of
$\Delta_n(\Delta_L)$.
% Lean: generators_satisfy_gl_relations; (1) Iteration_of one_step_invariance_holds lowering_operators_commute
% symMap_equivariant_holds (jucysMurphyProduct_contract_last through contract_last_jucysMurphy_castSucc,
% contract_last_perm_invariant, perm_sum_over_perms_eq); (2) IotaMu_of (symMap_norm_toSym_eq_factorial_smul,
% jucysMurphyProduct_norm_toSym_comm, norm_apply_eq_mfact_smul); (3) ActSpec_of
% (symmetric_action_eq_symMap, symMap_single_eq_wordProduct); (4) ActLattice_of.
\end{proof}

\begin{lemma}\label{is-onto-step-exp-bound-short-graded-skeleton-quotient}
Let $O\to K$ be a homomorphism of commutative rings, $M$ a $K$-module, on which $O$ acts through $K$, $n\geq0$,
and $c_0,\dots,c_{n-1}$ commuting $O$-linear endomorphisms of $M$. Let $(M_\tau)_{\tau\in J}$ be $K$-submodules of
$M$ and $\chi_\tau\in K^n$ such that $c_am=\chi_{\tau,a}m$ for every $a$ and $m\in M_\tau$, and let $D\in O[X]$ be
such that $X^n+\sum_{a<n}\chi_{\tau,a}X^a$ divides the image of $D$ in $K[X]$ for every $\tau$. Let $w\in O$. Then
there is a $Q$ in the $O$-subalgebra of $\operatorname{End}_O(M)$ generated by the $c_a$ such that
$D(w)m=\big(w^n+\sum_{a<n}w^ac_a\big)(Qm)$ for every $m\in\sum_\tau M_\tau$.
\end{lemma}
\leanref{Formalism/PieriBlockIdeals/IntegralQuotient.lean}{Formalism.CohomologicalDescent.Symmetrization.IntegralQuotient}
\leanref{Formalism/PieriBlockIdeals/PieriColumnMinors/Symmetrization.lean}{Formalism.CohomologicalDescent.Symmetrization.integral_quotient_holds}
\begin{proof}
Let $R=O[Y_0,\dots,Y_{n-1}]$ and $\varphi\colon R\to\operatorname{End}_O(M)$ the homomorphism with $\varphi(Y_a)=c_a$,
which exists as the $c_a$ commute; its image is the subalgebra generated by the $c_a$. For $\tau\in J$ let
$\psi_\tau\colon R\to K$ be the homomorphism with $\psi_\tau(Y_a)=\chi_{\tau,a}$; by induction on the polynomial,
$\varphi(r)m=\psi_\tau(r)m$ for $r\in R$ and $m\in M_\tau$. The polynomial $\mathcal A=X^n+\sum_aY_aX^a\in R[X]$ is
monic; divide: $D=D_0+\mathcal AD_1$ in $R[X]$, with $D_0$ of degree $<n$. Applying $\psi_\tau$ to the
coefficients, the image of $\mathcal A$ is $X^n+\sum_a\chi_{\tau,a}X^a$, which divides the image of $D$; as division
by a monic polynomial commutes with $\psi_\tau$, the image of $D_0$ is $0$. Evaluate at $X=w$:
$D(w)=D_0(w)+\mathcal A(w)D_1(w)$ in $R$, with $\varphi(\mathcal A(w))=w^n+\sum_aw^ac_a$. Put $Q=\varphi(D_1(w))$. For
$m\in M_\tau$, $\varphi(D_0(w))m=\psi_\tau(D_0(w))m=0$, so $D(w)m=(w^n+\sum_aw^ac_a)(Qm)$. Both sides are additive in
$m$.
% Lean: integral_quotient_holds: exists_mvPolynomial_algHom_commuting (phi with values in Algebra.adjoin),
% AAp monic, modByMonic_add_div, map_modByMonic and modByMonic_eq_zero_iff_dvd (hRem), then
% Submodule.iSup_induction.
\end{proof}

\begin{lemma}\label{is-onto-step-exp-bound-short-graded-skeleton-weights}
Let $K$, $I_1$, $q$ and $(V,E,V^+,w)$ be as in Definition~\ref{is-onto-step-exp-bound-short-graded-skeleton-maps},
$n\geq0$ and $j<n$.
\begin{enumerate}
\item Let $O$ be a commutative ring with a homomorphism $O\to K$, acting on $V$ through it, and $L\subset V^+$ an
  $O$-submodule with $E_{bb'}L\subset L$ for all $b,b'$. If $z\in T_n$ has $z(g)\in L$ for every $g$, then
  $(x_jz)(g)\in L$ for every $g$.
\item Let $\nu\in\mathbf Z^{I_1\sqcup I_2}$. If $z(g)\in V_{\nu+\mathrm{wt}([g])}$ for every $g$, then
  $(x_jz)(g)\in V_{\nu+\mathrm{wt}([g])}$ for every $g$.
\item If $v\in V_\nu$, then $\iota_n(v)(m)\in V_{\nu+\mathrm{wt}(m)}$ for every multiset $m$ of $n$ elements.
\item If $\xi\colon\mathbf Z^{I_1\sqcup I_2}\to\mathbf Z$ is a homomorphism with $\xi(\epsilon_\alpha-\epsilon_b)\geq0$
  for all $\alpha$ and $b$, then $\xi(\mathrm{wt}(m))\geq0$ for every multiset $m$.
\end{enumerate}
\end{lemma}
\leanref{Formalism/PieriBlockIdeals/IntegralQuotient.lean}{Formalism.CohomologicalDescent.Symmetrization.XLattice}
\leanref{Formalism/PieriBlockIdeals/PieriColumnMinors/Symmetrization.lean}{Formalism.CohomologicalDescent.Symmetrization.x_lattice_holds}
\leanref{Formalism/PieriBlockIdeals/IntegralQuotient.lean}{Formalism.CohomologicalDescent.Symmetrization.XWeight}
\leanref{Formalism/PieriBlockIdeals/PieriColumnMinors/Symmetrization.lean}{Formalism.CohomologicalDescent.Symmetrization.x_weight_holds}
\leanref{Formalism/PieriBlockIdeals/IntegralQuotient.lean}{Formalism.CohomologicalDescent.Symmetrization.IotaWeight}
\leanref{Formalism/PieriBlockIdeals/PieriColumnMinors/Symmetrization.lean}{Formalism.CohomologicalDescent.Symmetrization.iota_weight_holds}
\leanref{Formalism/PieriBlockIdeals/IntegralQuotient.lean}{Formalism.CohomologicalDescent.Symmetrization.BetaDominant}
\leanref{Formalism/PieriBlockIdeals/PieriColumnMinors/Symmetrization.lean}{Formalism.CohomologicalDescent.Symmetrization.beta_dominant_holds}
\begin{proof}
Write $g_s=(b_s,\alpha_s)$. By the definitions,
$(x_jz)(g)=\sum_{b'}E_{b'b_j}z(g^{(b')})+\sum_{j'<j}z(g_{(j'j)})$, where $g^{(b')}$ is $g$ with $b_j$ replaced by $b'$
and $g_{(j'j)}$ is $g$ with $b_{j'}$ and $b_j$ exchanged. This gives (1). For (2), the relations give
$E_{cd}V_\kappa\subset V_{\kappa+\epsilon_c-\epsilon_d}$; moreover
$\mathrm{wt}([g^{(b')}])+\epsilon_{b'}-\epsilon_{b_j}=\mathrm{wt}([g])$ and $\mathrm{wt}([g_{(j'j)}])=\mathrm{wt}([g])$. For
(3) and (4), $\mathrm{wt}(m)=\sum_{(b,\alpha)}(\epsilon_\alpha-\epsilon_b)$, the sum over the elements of $m$, and
$e_{(b,\alpha)}=E_{\alpha b}$ maps $V_\kappa$ into $V_{\kappa+\epsilon_\alpha-\epsilon_b}$.
% Lean: x_lattice_holds and x_weight_holds (jucysMurphy_apply, gl_generator_shifts_weight, weight_content_update,
% weight_content_swap_invariant); iota_weight_holds (list_raising_mem_weightSpace); beta_dominant_holds
% (weightOfMultiset_eq_sum_roots).
\end{proof}

\begin{lemma}\label{is-onto-step-exp-bound-short-graded-skeleton-assembly}
Let $K$ be a field of characteristic $0$, $O$ a commutative ring with a homomorphism $O\to K$, $I_1$, $q$ and
$(V,E,V^+,w)$ as in Definition~\ref{is-onto-step-exp-bound-short-graded-skeleton-maps}, with $O$ acting on $V$
through $K$, $n\geq0$, and $w_O\in O$ with image $w$ in $K$. Let $\nabla',\Delta_L\subset V^+$ be $O$-submodules,
$(M_\tau)_{\tau\in J}$ subspaces of the space of functions on the multisets of $n$ elements with values in $V^+$,
$\chi_\tau\in K^n$, $D\in O[X]$, $\xi\colon\mathbf Z^{I_1\sqcup I_2}\to\mathbf Z$ a homomorphism,
$\lambda\in\mathbf Z^{I_1\sqcup I_2}$, $l\in O$, and $e\geq0$ and $i'$ integers. Write $A_a$ for the coefficient of
$X^a$ in $A(X)$ (Definition~\ref{is-onto-step-exp-bound-short-graded-skeleton-jucys}). Suppose that:
\begin{enumerate}
\item[(a)] $E_{bb'}\nabla'\subset\nabla'$ for all $b,b'\in I_2$;
\item[(b)] $\mu_n(z)=0$ for every $z\in T_n$ with $\iota_n(\mu_n(z))=0$;
\item[(c)] $(A_az)^{\mathrm{sym}}=\chi_{\tau,a}z^{\mathrm{sym}}$ for every $\tau$, $a<n$ and $z\in T_n$ with
  $z^{\mathrm{sym}}\in M_\tau$;
\item[(d)] $X^n+\sum_{a<n}\chi_{\tau,a}X^a$ divides the image of $D$ in $K[X]$ for every $\tau$;
\item[(e)] $(A(w)z)^{\mathrm{sym}}\in\sum_\tau M_\tau$ for every $z\in T_n$;
\item[(f)] $\xi(\epsilon_\alpha-\epsilon_b)\geq0$ for all $\alpha\in I_1$ and $b\in I_2$;
\item[(g)] $l^eu\in\Delta_L$ for every $\kappa$ with $\xi(\lambda-\kappa)\le i'$ and every $u\in\nabla'\cap V_\kappa$.
\end{enumerate}
Let $\nu\in\mathbf Z^{I_1\sqcup I_2}$ with $\xi(\lambda-\nu)\le i'$, and $z\in T_n$ with $\mu_n(z)\in V_\nu$ and
$\iota_n(\mu_n(z))(m)\in\nabla'$ for every $m$. Then $l^eD(w_O)\mu_n(z)\in\Delta_n(\Delta_L)$.
\end{lemma}
\leanref{Formalism/PieriBlockIdeals/IntegralQuotient.lean}{Formalism.CohomologicalDescent.Symmetrization.LemmaDSkeleton}
\leanref{Formalism/PieriBlockIdeals/PieriColumnMinors/Symmetrization.lean}{Formalism.CohomologicalDescent.Symmetrization.LemmaDSkeleton_of}
\leanref{Formalism/PieriBlockIdeals/PieriColumnMinors/Symmetrization.lean}{Formalism.CohomologicalDescent.Symmetrization.lemma_d_skeleton_holds}
\begin{proof}
Let $c_a$ be as in Lemma~\ref{is-onto-step-exp-bound-short-graded-skeleton-commute}(3) for $M=V^+$ (see the proof
of Lemma~\ref{is-onto-step-exp-bound-short-graded-skeleton-iota}). They are $O$-linear and commute, and
$c_ay=\chi_{\tau,a}y$ for $y\in M_\tau$, by (c), as $y=z'^{\mathrm{sym}}$ for some $z'$. By (d) and
Lemma~\ref{is-onto-step-exp-bound-short-graded-skeleton-quotient}, there is a $Q$ in the $O$-algebra generated by
the $c_a$ with $D(w_O)y=(w^n+\sum_aw^ac_a)(Qy)$ for every $y\in\sum_\tau M_\tau$.

Put $y_0=(A(w)z)^{\mathrm{sym}}$; by Lemma~\ref{is-onto-step-exp-bound-short-graded-skeleton-iota}(2),
$y_0(m)=\iota_n(\mu_n(z))(m)$. Let $L(m)=\nabla'\cap V_{\nu+\mathrm{wt}(m)}$. Then $y_0(m)\in L(m)$: it lies in
$\nabla'$ by hypothesis and in $V_{\nu+\mathrm{wt}(m)}$ by
Lemma~\ref{is-onto-step-exp-bound-short-graded-skeleton-weights}(3). By (a) and parts (1) and (2) of that lemma,
every $x_j$ maps the set of the $z'$ with $z'(g)\in L([g])$ for all $g$ into itself. So by
Lemma~\ref{is-onto-step-exp-bound-short-graded-skeleton-commute}(4) every $c_a$, and hence $Q$, maps the set of
the $y$ with $y(m)\in L(m)$ for all $m$ into itself. So $y_1=Qy_0$ has $y_1(m)\in L(m)$ for every $m$. By
Lemma~\ref{is-onto-step-exp-bound-short-graded-skeleton-weights}(4) and (f),
$\xi(\lambda-\nu-\mathrm{wt}(m))=\xi(\lambda-\nu)-\xi(\mathrm{wt}(m))\le i'$, so (g) gives $l^ey_1(m)\in\Delta_L$.

Choose $z_1\in T_n$ with $z_1^{\mathrm{sym}}=l^ey_1$
(Lemma~\ref{is-onto-step-exp-bound-short-graded-skeleton-jucys-relations}(6)). By
Lemma~\ref{is-onto-step-exp-bound-short-graded-skeleton-iota}(3) and~(4),
$\mu_n(z_1)=\mathrm{act}_n(l^ey_1)\in\Delta_n(\Delta_L)$. Put $r=l^eD(w_O)$. By (e) and
Lemma~\ref{is-onto-step-exp-bound-short-graded-skeleton-commute}(3),
$(A(w)z_1)^{\mathrm{sym}}=(w^n+\sum_aw^ac_a)(l^eQy_0)=l^eD(w_O)y_0=ry_0$. So, by
Lemma~\ref{is-onto-step-exp-bound-short-graded-skeleton-iota}(2),
$\iota_n(\mu_n(z_1-rz))(m)=(A(w)z_1)^{\mathrm{sym}}(m)-ry_0(m)=0$ for every $m$, and (b) gives $\mu_n(z_1)=r\mu_n(z)$.
Hence $r\mu_n(z)\in\Delta_n(\Delta_L)$.
% Lean: LemmaDSkeleton_of coefficient_operators_hold x_lattice_holds x_weight_holds iota_weight_holds
% beta_dominant_holds action_lattice_holds integral_quotient_holds iota_symMap_holds action_on_symmetric_holds:
% c' (restrictScalars O), hchi, Q from hIQ, y0, the lattices L m = nabL meet wt (nu + beta m), hQL by
% Algebra.adjoin_induction, hy2, z1 from toSym_surjective, hmem, hAy, hzero with iota_sub_smul.
\end{proof}

\begin{lemma}\label{is-onto-step-exp-bound-short-graded-skeleton-form}
Let $K$ be a field of characteristic $0$, $I_1$, $q$ and $(V,E,V^+,w)$ as in
Definition~\ref{is-onto-step-exp-bound-short-graded-skeleton-maps}, $p=\#I_1$, $n\geq0$, and $B$ a bilinear form on
$V$ with $B(E_{cd}x,y)=B(x,E_{dc}y)$ for all $c,d$ and $x,y\in V$.
\begin{enumerate}
\item $B(F_mx,y)=B(x,E_my)$ and $B(x,F_my)=B(E_mx,y)$ for every multiset $m$ and $x,y\in V$. So
  $B\big((m!)^{-1}F_mu,v\big)=B(u,\iota_n(v)(m))$ and $B\big(v,(m!)^{-1}F_mu\big)=B(\iota_n(v)(m),u)$ for $m$ of $n$
  elements.
\item If $n'\ne n$, then $B(\mu_n(z),\mu_{n'}(z'))=0$ for $z\in T_n$ and $z'\in T_{n'}$.
\item Suppose that every subspace of $V$ that contains $V^+$ and is stable under every $E_{cd}$ is $V$, and that
  $v=0$ if $B(x,v)=0$ for every $x\in V$. Then $\mu_n(z)=0$ for every $z\in T_n$ with $\iota_n(\mu_n(z))=0$.
\item Let $O$, $\Delta_L$, $\Delta(\Delta_L)$, $N^B$ and $\nabla_L$ be as in
  Definition~\ref{is-onto-step-exp-bound-short-graded-block}(2). For $z\in T_n$, $\mu_n(z)\in\Delta(\Delta_L)^B$ if
  and only if $\iota_n(\mu_n(z))(m)\in\nabla_L$ for every $m$.
\item If $E_{bb'}\Delta_L\subset\Delta_L$ for all $b,b'$, then $E_{bb'}\nabla_L\subset\nabla_L$ for all $b,b'$.
\end{enumerate}
\end{lemma}
\leanref{Formalism/PieriBlockIdeals/LemmaDSkeletonUnit.lean}{Formalism.CohomologicalDescent.BinaryFormsExample.Contra}
\leanref{Formalism/PieriBlockIdeals/LemmaDSkeletonUnit.lean}{Formalism.CohomologicalDescent.BinaryFormsExample.ContraWord}
\leanref{Formalism/PieriBlockIdeals/PieriColumnMinors/BinaryFormsExample.lean}{Formalism.CohomologicalDescent.BinaryFormsExample.contravariant_word_holds}
\leanref{Formalism/PieriBlockIdeals/LemmaDSkeletonUnit.lean}{Formalism.CohomologicalDescent.BinaryFormsExample.ContraIota}
\leanref{Formalism/PieriBlockIdeals/PieriColumnMinors/BinaryFormsExample.lean}{Formalism.CohomologicalDescent.BinaryFormsExample.contravariant_iota_holds}
\leanref{Formalism/PieriBlockIdeals/LemmaDSkeletonUnit.lean}{Formalism.CohomologicalDescent.BinaryFormsExample.DepthOrthogonal}
\leanref{Formalism/PieriBlockIdeals/PieriColumnMinors/BinaryFormsExample.lean}{Formalism.CohomologicalDescent.BinaryFormsExample.depth_orthogonal_holds}
\leanref{Formalism/PieriBlockIdeals/LemmaDSkeletonUnit.lean}{Formalism.CohomologicalDescent.BinaryFormsExample.IotaKernel}
\leanref{Formalism/PieriBlockIdeals/PieriColumnMinors/BinaryFormsExample.lean}{Formalism.CohomologicalDescent.BinaryFormsExample.iota_kernel_holds}
\leanref{Formalism/PieriBlockIdeals/LemmaDSkeletonUnit.lean}{Formalism.CohomologicalDescent.BinaryFormsExample.IotaInjective}
\leanref{Formalism/PieriBlockIdeals/PieriColumnMinors/BinaryFormsExample.lean}{Formalism.CohomologicalDescent.BinaryFormsExample.iota_injective_holds}
\leanref{Formalism/PieriBlockIdeals/LemmaDSkeletonUnit.lean}{Formalism.CohomologicalDescent.BinaryFormsExample.NablaCriterion}
\leanref{Formalism/PieriBlockIdeals/PieriColumnMinors/BinaryFormsExample.lean}{Formalism.CohomologicalDescent.BinaryFormsExample.nabla_criterion_holds}
\leanref{Formalism/PieriBlockIdeals/LemmaDSkeletonUnit.lean}{Formalism.CohomologicalDescent.BinaryFormsExample.DualStable}
\leanref{Formalism/PieriBlockIdeals/PieriColumnMinors/BinaryFormsExample.lean}{Formalism.CohomologicalDescent.BinaryFormsExample.dual_stable_holds}
\begin{proof}
(1) The hypothesis gives $B(f_vx,y)=B(x,e_vy)$. Iterating over the elements of $m$ gives $B(F_mx,y)=B(x,E'y)$ with
$E'$ the product of the $e_v$ in the reverse order, which is $E_m$ as the $e_v$ commute
(Lemma~\ref{is-onto-step-exp-bound-short-graded-skeleton-mu}(1)). The form $(x,y)\mapsto B(y,x)$ satisfies the same
hypothesis, which gives the second identity.

(2) The operator $\sum_\alpha E_{\alpha\alpha}$ is symmetric for $B$, so by
Lemma~\ref{is-onto-step-exp-bound-short-graded-skeleton-mu}(3),
$(pw-n)B(\mu_n(z),\mu_{n'}(z'))=(pw-n')B(\mu_n(z),\mu_{n'}(z'))$, and $n\ne n'$ in characteristic $0$.

(3) Let $v=\mu_n(z)$. For $z'\in T_n$, (1) and Lemma~\ref{is-onto-step-exp-bound-short-graded-skeleton-mu}(1) give
$B(\mu_n(z'),v)=\sum_gB(z'(g),E_gv)=\sum_g[g]!\,B\big(z'(g),\iota_n(v)([g])\big)=0$. For $n'\ne n$ and $z'\in T_{n'}$,
$B(\mu_{n'}(z'),v)=0$ by (2). By Lemma~\ref{is-onto-step-exp-bound-short-graded-skeleton-mu}(5),
$V=\sum_{n'}\mu_{n'}(T_{n'})$, so $B(x,v)=0$ for every $x\in V$, and $v=0$.

(4) Suppose $\mu_n(z)\in\Delta(\Delta_L)^B$, and let $m$ be a multiset and $u\in\Delta_L$. By
Lemma~\ref{is-onto-step-exp-bound-short-graded-skeleton-iota}(2), $\iota_n(\mu_n(z))(m)\in V^+$, and by (1),
$B(\iota_n(\mu_n(z))(m),u)=B(\mu_n(z),(m!)^{-1}F_mu)$, which lies in the image of $O$ as $(m!)^{-1}F_mu\in\Delta_n(\Delta_L)$.
Conversely, $\Delta(\Delta_L)$ is spanned over $O$ by the $(m!)^{-1}F_mu$ with $u\in\Delta_L$ and $m$ of any size
$n'$. For $n'=n$, $B(\mu_n(z),(m!)^{-1}F_mu)=B(\iota_n(\mu_n(z))(m),u)$ by (1). For $n'\ne n$, $F_mu=\mu_{n'}(z')$ for
the $z'$ with value $u$ at one $g$ with $[g]=m$ and $0$ elsewhere, and the pairing is $0$ by (2).

(5) For $u\in\nabla_L$ and $x\in\Delta_L$, $B(E_{bb'}u,x)=B(u,E_{b'b}x)$, which lies in the image of $O$ as
$E_{b'b}x\in\Delta_L$.
% Lean: (1) contravariant_word_holds (pairing_list_reverse_adjoint, pairing_list_adjoint, contravariant_flip) and
% contravariant_iota_holds; (2) depth_orthogonal_holds (depth_eigen_holds); (3) iota_kernel_holds and
% iota_injective_holds (depth_span_holds, Submodule.iSup_induction); (4) nabla_criterion_holds
% (iota_symMap_holds, symMap_single_eq_wordProduct, depth_orthogonal_holds); (5) dual_stable_holds.
\end{proof}

\begin{lemma}\label{is-onto-step-exp-bound-short-graded-skeleton-shapes}
Let $p,w,i,q\geq0$ be integers and $\lambda_2\in\mathbf N^q$, with $b(\lambda_2,\tau,c)$, $\mathrm{col}$, $\mu(c)$,
$\mathrm{Sh}(p,w,i,\lambda_2)$ and $D_n(p,w,i,\lambda_2)$ as in
Definition~\ref{is-onto-step-exp-bound-short-branching}(3). For $\tau\in\mathbf N^q$ put
$s_{\lambda_2,\tau}(X)=\prod_{a<q}\prod_{\lambda_{2,a}\le y<\tau_a}\big(X-(y-a)\big)\in\mathbf Z[X]$, and let
$D(X)=\prod_{0\le y<q+w+i}\big(X-(y-q)\big)^{\mu(y-q)}$. Let $\mathrm{Sh}^*(p,i,\lambda_2)$ be the set of the
nonincreasing $\tau\in\mathbf N^q$ with $\lambda_{2,a}\le\tau_a\le\lambda_{2,a}+i$ for every $a$,
$\sum_a\tau_a=\sum_a\lambda_{2,a}+i$, and $\mathrm{col}(\lambda_2,\tau,y)\le p$ for every $a$ and every $y<\tau_a$.
\begin{enumerate}
\item For $c\in\mathbf Z$, $s_{\lambda_2,\tau}(c)=0$ if and only if $b(\lambda_2,\tau,c)\ne0$. For
  $\tau\in\mathrm{Sh}^*(p,i,\lambda_2)$, $s_{\lambda_2,\tau}$ is monic of degree $i$. The set
  $\mathrm{Sh}(p,w,i,\lambda_2)$ consists of the $\tau\in\mathrm{Sh}^*(p,i,\lambda_2)$ with $b(\lambda_2,\tau,w)=0$.
\item $D(w)=D_n(p,w,i,\lambda_2)$. If $\lambda_{2,a}\le w$ for every $a$ and $\tau\in\mathrm{Sh}(p,w,i,\lambda_2)$,
  then $s_{\lambda_2,\tau}$ divides $D$.
\item Let $\kappa$ be a box path of length $n$ from $\lambda_2$ to $\tau$, with contents $c_j$
  (Definition~\ref{is-onto-step-exp-bound-short-graded-skeleton-jucys}(4)). Then
  $\prod_{j<n}(X-c_j)=s_{\lambda_2,\tau}$, $\lambda_{2,a}\le\tau_a\le\lambda_{2,a}+n$ for every $a$, and
  $\sum_a\tau_a=\sum_a\lambda_{2,a}+n$.
\item Let $K$, $I_1$, $n$, $M$ and $E_{bb'}$ be as in
  Definition~\ref{is-onto-step-exp-bound-short-graded-skeleton-jucys}, with $K$ of characteristic $0$ and the
  $E_{bb'}$ satisfying the relations of Lemma~\ref{is-onto-step-exp-bound-short-graded-skeleton-commute}. If
  $z\in T_n$ and $z^{\mathrm{sym}}=z'^{\mathrm{sym}}$ for some $z'\in T_n(\lambda_2,\tau)$, then
  $(A_az)^{\mathrm{sym}}=s_az^{\mathrm{sym}}$ for every $a$, where $A_a$ and $s_a$ are the coefficients of $X^a$ in
  $A(X)$ and in $s_{\lambda_2,\tau}$.
\item In (4), without the relations, let $\#I_1=p$, and suppose that $T_n$ is the sum of the
  $T_n(\lambda_2,\tau)$, $\tau\in\mathbf N^q$, and that $z^{\mathrm{sym}}=0$ for every $z\in T_n(\lambda_2,\tau)$
  whenever $\mathrm{col}(\lambda_2,\tau,y)>p$ for some $a$ and some $y<\tau_a$. Then every function on the
  multisets of $n$ elements with values in $M$ is a sum of elements $z^{\mathrm{sym}}$ with $z\in T_n(\lambda_2,\tau)$
  and $\tau\in\mathrm{Sh}^*(p,n,\lambda_2)$.
\end{enumerate}
\end{lemma}
\leanref{Formalism/PieriBlockIdeals/IntegralQuotient.lean}{Formalism.CohomologicalDescent.Symmetrization.shapePoly}
\leanref{Formalism/PieriBlockIdeals/IntegralQuotient.lean}{Formalism.CohomologicalDescent.Symmetrization.DnPoly}
\leanref{Formalism/PieriBlockIdeals/LemmaDSkeletonUnit.lean}{Formalism.CohomologicalDescent.BinaryFormsExample.shapesAll}
\leanref{Formalism/PieriBlockIdeals/PieriColumnMinors/BinaryFormsExample.lean}{Formalism.CohomologicalDescent.BinaryFormsExample.shape_eval_zero_holds}
\leanref{Formalism/PieriBlockIdeals/PieriColumnMinors/BinaryFormsExample.lean}{Formalism.CohomologicalDescent.BinaryFormsExample.shape_coefficient_holds}
\leanref{Formalism/PieriBlockIdeals/PieriColumnMinors/BinaryFormsExample.lean}{Formalism.CohomologicalDescent.BinaryFormsExample.mem_shapes_holds}
\leanref{Formalism/PieriBlockIdeals/PieriColumnMinors/Symmetrization.lean}{Formalism.CohomologicalDescent.Symmetrization.polynomial_evaluation_holds}
\leanref{Formalism/PieriBlockIdeals/PieriColumnMinors/Symmetrization.lean}{Formalism.CohomologicalDescent.Symmetrization.shape_divisibility_holds}
\leanref{Formalism/PieriBlockIdeals/PieriColumnMinors/BinaryFormsExample.lean}{Formalism.CohomologicalDescent.BinaryFormsExample.shape_divisibility_field_holds}
\leanref{Formalism/PieriBlockIdeals/LemmaDSkeletonUnit.lean}{Formalism.CohomologicalDescent.BinaryFormsExample.PathShape}
\leanref{Formalism/PieriBlockIdeals/PieriColumnMinors/BinaryFormsExample.lean}{Formalism.CohomologicalDescent.BinaryFormsExample.path_shape_holds}
\leanref{Formalism/PieriBlockIdeals/LemmaDSkeletonUnit.lean}{Formalism.CohomologicalDescent.BinaryFormsExample.PathSpaceEigen}
\leanref{Formalism/PieriBlockIdeals/PieriColumnMinors/BinaryFormsExample.lean}{Formalism.CohomologicalDescent.BinaryFormsExample.path_space_eigen_holds}
\leanref{Formalism/PieriBlockIdeals/LemmaDSkeletonUnit.lean}{Formalism.CohomologicalDescent.BinaryFormsExample.PathDecomp}
\leanref{Formalism/PieriBlockIdeals/PieriColumnMinors/BinaryFormsExample.lean}{Formalism.CohomologicalDescent.BinaryFormsExample.path_decomposition_holds}
\begin{proof}
(1) $s_{\lambda_2,\tau}(c)=0$ if and only if $y-a=c$ for some $a$ and $y$ with $\lambda_{2,a}\le y<\tau_a$, that is,
$\lambda_{2,a}\le c+a<\tau_a$. The polynomial is a product of $\sum_a(\tau_a-\lambda_{2,a})=i$ monic linear factors. The
last claim compares the definitions.

(2) The first claim is the definition of $D_n$. Every factor $X-(y-a)$ of $s_{\lambda_2,\tau}$ has
$0\le y+q-a<q+w+i$, as $y<\tau_a\le\lambda_{2,a}+i\le w+i$ and $a<q$. Group the factors by $k=y+q-a$: the factor
$X-(k-q)$ occurs once for each $a$ with $\lambda_{2,a}\le k-q+a<\tau_a$, that is, $b(\lambda_2,\tau,k-q)$ times, and
$b(\lambda_2,\tau,k-q)\le\mu(k-q)$ as $\tau\in\mathrm{Sh}(p,w,i,\lambda_2)$.

(3) Induct on the length: if $\lambda_{2,r}\le\kappa_r$, then
$s_{\lambda_2,\kappa+\epsilon_r}=\big(X-(\kappa_r-r)\big)s_{\lambda_2,\kappa}$, and the bounds and the sum change by
one box.

(4) The space $T_n(\lambda_2,\tau)$ is spanned by the $z''$ for which a box path from $\lambda_2$ to $\tau$ has
$C_jz''=\mathrm{cas}(\kappa_j)z''$ for every $j\le n$; for these, $A_az''=s_az''$ by
Lemma~\ref{is-onto-step-exp-bound-short-graded-skeleton-path-eigen} and (3). So $A_az'=s_az'$. With $c_a$ as in
Lemma~\ref{is-onto-step-exp-bound-short-graded-skeleton-commute}(3),
$(A_az)^{\mathrm{sym}}=c_a(z^{\mathrm{sym}})=c_a(z'^{\mathrm{sym}})=(A_az')^{\mathrm{sym}}=s_az^{\mathrm{sym}}$.

(5) By Lemma~\ref{is-onto-step-exp-bound-short-graded-skeleton-jucys-relations}(6) every function is some
$z^{\mathrm{sym}}$, and $z=\sum_\tau z_\tau$ with $z_\tau\in T_n(\lambda_2,\tau)$. If $\tau$ has a column of length
$\mathrm{col}(\lambda_2,\tau,y)>p$ with $y<\tau_a$, then $z_\tau^{\mathrm{sym}}=0$. If there is no box path of length $n$
from $\lambda_2$ to $\tau$, then $T_n(\lambda_2,\tau)=0$. Otherwise $\tau=\kappa_n$ is nonincreasing, and (3) gives the
bounds and the sum, so $\tau\in\mathrm{Sh}^*(p,n,\lambda_2)$.
% Lean: (1) shape_eval_zero_holds, shapePoly_monic, natDegree_shapePoly_eq, shape_coefficient_holds and
% mem_shapes_holds; (2) polynomial_evaluation_holds and shape_divisibility_holds (Finset.prod_fiberwise_of_maps_to,
% then the injection into boxesOfContent and Finset.le_sup), assembled with the field maps in
% shape_divisibility_field_holds; (3) path_shape_holds (path_bounds_and_shapePoly, shapePoly_add_single); (4)
% path_space_eigen_holds (coefficient_operators_hold, path_eigen_holds); (5) path_decomposition_holds
% (toSym_surjective).
\end{proof}

\begin{lemma}\label{is-onto-step-exp-bound-short-graded-skeleton}
Let $K$ be a field of characteristic $0$, $I_1$ a finite set with $p$ elements, $q,n,w\geq0$ integers,
$I_2=\{0,\dots,q-1\}$, and $(V,E,V^+,w)$ a block module over $K$ for $I_1$ and $I_2$ whose element of $K$
is the integer $w$ (Definition~\ref{is-onto-step-exp-bound-short-graded-block}). Let $l$ be a prime and $O$
a commutative ring with a homomorphism $O\to K$, acting on $V$ through it, in which every integer not
divisible by $l$ is a unit. Let $B$ be a bilinear form on $V$, $\lambda_2\in\mathbf N^q$, $\Delta_L\subset V^+$
an $O$-submodule, $\xi\colon\mathbf Z^{I_1\sqcup I_2}\to\mathbf Z$ a homomorphism,
$\lambda\in\mathbf Z^{I_1\sqcup I_2}$, and $e\geq0$ and $i'$ integers. Suppose that:
\begin{enumerate}
\item $\lambda_{2,b}\le w$ for every $b$;
\item $B(E_{cd}x,y)=B(x,E_{dc}y)$ for all $c,d$ and $x,y\in V$; every subspace of $V$ that contains $V^+$
  and is stable under every $E_{cd}$ is $V$; and $v=0$ if $B(x,v)=0$ for every $x\in V$;
\item $E_{bb'}\Delta_L\subset\Delta_L$ for all $b,b'$;
\item $T_n$ is the sum of the box path spaces $T_n(\lambda_2,\tau)$, $\tau\in\mathbf N^q$;
\item if $\tau\in\mathbf N^q$ and $\mathrm{col}(\lambda_2,\tau,y)>p$ for some $b$ and some $y<\tau_b$
  (Definition~\ref{is-onto-step-exp-bound-short-branching}(3)), then $z^{\mathrm{sym}}=0$ for every
  $z\in T_n(\lambda_2,\tau)$;
\item $\xi(\epsilon_\alpha-\epsilon_b)\geq0$ for all $\alpha$ and $b$;
\item $l^eu\in\Delta_L$ for every $\kappa$ with $\xi(\lambda-\kappa)\le i'$ and every $u\in\nabla_L$ that lies in
  $V_\kappa$.
\end{enumerate}
Let $\nu\in\mathbf Z^{I_1\sqcup I_2}$ with $\sum_\alpha\nu_\alpha+n=pw$ and $\xi(\lambda-\nu)\le i'$, and let
$v\in V_\nu$ lie in $\Delta(\Delta_L)^B$. Then $l^{e+v_l(D_n(p,w,n,\lambda_2))}v$ lies in $\Delta(\Delta_L)$ and
in $V_\nu$.
\end{lemma}
\leanref{Formalism/PieriBlockIdeals/LemmaDSkeletonUnit.lean}{Formalism.CohomologicalDescent.BinaryFormsExample.LemmaDInputs3}
\leanref{Formalism/PieriBlockIdeals/LemmaDSkeletonUnit.lean}{Formalism.CohomologicalDescent.BinaryFormsExample.Contra}
\leanref{Formalism/PieriBlockIdeals/LemmaDSkeletonUnit.lean}{Formalism.CohomologicalDescent.BinaryFormsExample.Generated}
\leanref{Formalism/PieriBlockIdeals/LemmaDSkeletonUnit.lean}{Formalism.CohomologicalDescent.BinaryFormsExample.LemmaDSkeleton3}
\leanref{Formalism/PieriBlockIdeals/PieriColumnMinors/BinaryFormsExample.lean}{Formalism.CohomologicalDescent.BinaryFormsExample.depth_skeleton_third_form_holds}
\leanref{Formalism/PieriBlockIdeals/LemmaDSkeletonUnit.lean}{Formalism.CohomologicalDescent.BinaryFormsExample.InputsFromPieri}
\leanref{Formalism/PieriBlockIdeals/PieriColumnMinors/BinaryFormsExample.lean}{Formalism.CohomologicalDescent.BinaryFormsExample.inputs_from_pieri_holds}
\leanref{Formalism/PieriBlockIdeals/LemmaDSkeletonUnit.lean}{Formalism.CohomologicalDescent.BinaryFormsExample.LemmaDSkeleton2Lit}
\leanref{Formalism/PieriBlockIdeals/PieriColumnMinors/BinaryFormsExample.lean}{Formalism.CohomologicalDescent.BinaryFormsExample.depth_skeleton_literal_holds}
\leanref{Formalism/PieriBlockIdeals/LemmaDSkeletonUnit.lean}{Formalism.CohomologicalDescent.BinaryFormsExample.LemmaDSkeleton2Val}
\leanref{Formalism/PieriBlockIdeals/PieriColumnMinors/BinaryFormsExample.lean}{Formalism.CohomologicalDescent.BinaryFormsExample.depth_skeleton_valuation_holds}
\leanref{Formalism/PieriBlockIdeals/LemmaDSkeletonUnit.lean}{Formalism.CohomologicalDescent.BinaryFormsExample.LemmaDSkeleton2}
\leanref{Formalism/PieriBlockIdeals/PieriColumnMinors/BinaryFormsExample.lean}{Formalism.CohomologicalDescent.BinaryFormsExample.depth_skeleton_second_form_holds}
\leanref{Formalism/PieriBlockIdeals/LemmaDSkeletonUnit.lean}{Formalism.CohomologicalDescent.BinaryFormsExample.GoodPart}
\leanref{Formalism/PieriBlockIdeals/PieriColumnMinors/BinaryFormsExample.lean}{Formalism.CohomologicalDescent.BinaryFormsExample.good_part_holds}
\leanref{Formalism/PieriBlockIdeals/LemmaDSkeletonUnit.lean}{Formalism.CohomologicalDescent.BinaryFormsExample.ValuationUnit}
\leanref{Formalism/PieriBlockIdeals/PieriColumnMinors/BinaryFormsExample.lean}{Formalism.CohomologicalDescent.BinaryFormsExample.valuation_unit_holds}
\begin{proof}
By (2), as $w$ is an integer, Lemma~\ref{is-onto-step-exp-bound-short-graded-skeleton-mu}(5) gives $z\in T_n$ with
$v=\mu_n(z)$. The restrictions of the $E_{bb'}$ to $V^+$ satisfy the relations of
Lemma~\ref{is-onto-step-exp-bound-short-graded-skeleton-commute}. For
$\tau\in\mathrm{Sh}^*=\mathrm{Sh}^*(p,n,\lambda_2)$
(Lemma~\ref{is-onto-step-exp-bound-short-graded-skeleton-shapes}), let $M_\tau$ be the space of the
$z'^{\mathrm{sym}}$ with $z'\in T_n(\lambda_2,\tau)$, and $\chi_{\tau,a}$ the coefficient of $X^a$ in
$s_{\lambda_2,\tau}$; as $s_{\lambda_2,\tau}$ is monic of degree $n$, $X^n+\sum_{a<n}\chi_{\tau,a}X^a=s_{\lambda_2,\tau}$.
By Lemma~\ref{is-onto-step-exp-bound-short-graded-skeleton-shapes}(4),
$(A_az')^{\mathrm{sym}}=\chi_{\tau,a}z'^{\mathrm{sym}}$ whenever $z'^{\mathrm{sym}}\in M_\tau$; and by (4), (5) and
part (5) of that lemma, the $M_\tau$ with $\tau\in\mathrm{Sh}^*$ span the space of the functions on the multisets
of $n$ elements with values in $V^+$.

Let $c_a$ be as in Lemma~\ref{is-onto-step-exp-bound-short-graded-skeleton-commute}(3) for $M=V^+$. It acts on
$M_\tau$ by $\chi_{\tau,a}$, so $w^n+\sum_aw^ac_a$ acts on $M_\tau$ by $s_{\lambda_2,\tau}(w)$, which is $0$ for
$\tau\in\mathrm{Sh}^*\setminus\mathrm{Sh}(p,w,n,\lambda_2)$ by
Lemma~\ref{is-onto-step-exp-bound-short-graded-skeleton-shapes}(1). Writing $z'^{\mathrm{sym}}=\sum_\tau y_\tau$
with $y_\tau\in M_\tau$, Lemma~\ref{is-onto-step-exp-bound-short-graded-skeleton-commute}(3) gives
$(A(w)z')^{\mathrm{sym}}=\sum_\tau s_{\lambda_2,\tau}(w)y_\tau$, which lies in the sum of the $M_\tau$ with
$\tau\in\mathrm{Sh}(p,w,n,\lambda_2)$, for every $z'\in T_n$.

Apply Lemma~\ref{is-onto-step-exp-bound-short-graded-skeleton-assembly} with $J=\mathrm{Sh}(p,w,n,\lambda_2)$,
these $M_\tau$ and $\chi_\tau$, $\nabla'=\nabla_L$, $D$ the image in $O[X]$ of the polynomial $D(X)$ of
Lemma~\ref{is-onto-step-exp-bound-short-graded-skeleton-shapes} for $(p,w,n,\lambda_2)$, $w_O$ the image of $w$,
and the given $\xi$, $\lambda$, $l$, $e$ and $i'$. Its hypotheses hold: (a) by (3) and
Lemma~\ref{is-onto-step-exp-bound-short-graded-skeleton-form}(5); (b) by (2) and
Lemma~\ref{is-onto-step-exp-bound-short-graded-skeleton-form}(3); (c) and (e) as just shown; (d) by
Lemma~\ref{is-onto-step-exp-bound-short-graded-skeleton-shapes}(2), as $\lambda_{2,b}\le w$ by (1); (f) and (g) are
(6) and (7). As $v\in\Delta(\Delta_L)^B$, Lemma~\ref{is-onto-step-exp-bound-short-graded-skeleton-form}(4) gives
$\iota_n(v)(m)\in\nabla_L$ for every $m$. So $l^eD(w)v\in\Delta_n(\Delta_L)$, and $D(w)=D_n(p,w,n,\lambda_2)$ by
Lemma~\ref{is-onto-step-exp-bound-short-graded-skeleton-shapes}(2).

Write $D_n=D_n(p,w,n,\lambda_2)$. It is a nonzero integer: in its defining product the only base that can vanish is
$w-(y-q)$ for $y-q=w$, whose exponent is $\mu(w)=0$
(Lemma~\ref{is-onto-step-exp-bound-short-top-window-products}(1)). So $D_n=\eta\,l^{v_l(D_n)}$ with $\eta$ an integer
not divisible by $l$, a unit of $O$. Hence $l^{e+v_l(D_n)}v=\eta^{-1}l^eD_nv$ lies in
$\Delta_n(\Delta_L)\subset\Delta(\Delta_L)$, and in $V_\nu$ as $v$ does.
% Lean: depth_skeleton_third_form_holds = depth_skeleton_literal_holds applied to inputs_from_pieri_holds
% (path_space_eigen_holds and path_decomposition_holds, which give the eigen and decomp fields of LemmaDInputs
% with Mt tau the image of PathSpace under toSym); depth_skeleton_literal_holds uses depth_of_weight_holds and
% depth_skeleton_valuation_holds = valuation_unit_holds (LD.d_n_ne_zero) with depth_skeleton_second_form_holds:
% shape_divisibility_field_holds, good_part_holds (good_part_abstract_holds) for the shapes, and
% lemma_d_skeleton_holds with dual_stable_holds, iota_injective_holds and nabla_criterion_holds, then
% unit_clause_holds. Size before this step was written out: about 3,400 lines of Lean in 210 theorems.
\end{proof}

\subsection{The rational block module and Pieri's rule}\label{subsection-is-pieri-rule}

\begin{definition}\label{is-onto-step-exp-bound-short-graded-rational}
Let $p,q,w\geq0$ be integers and $\lambda_2\in\mathbf N^q$ nonincreasing with $\lambda_{2,b}\le w$ for every
$b$, and put $\nu=w^p\lambda_2\in\mathbf N^{p+q}$ (Definition~\ref{is-onto-step-exp-bound-short-branching}(2)),
which is nonincreasing. Let $I_1=\{0,\dots,p-1\}$ and $I_2=\{0,\dots,q-1\}$, and identify
$I_1\sqcup I_2$ with $\{0,\dots,p+q-1\}$ by $\alpha\mapsto\alpha$ and $b\mapsto p+b$. The \emph{rational
block module} of $(p,q,w,\lambda_2)$ consists of the following data.
\begin{enumerate}
\item The space $V=\nabla_{\mathbf Q}(\nu)$ with the row operators $E_{cd}=\sum_sx_{cs}\,\partial/\partial
  x_{ds}$ (Definition~\ref{is-onto-step-exp-cut-grid}(6), with the rows indexed from $0$), the subspace
  $W_0$ of the $u\in V$ with $E_{\alpha b}u=0$ and $E_{\alpha\beta}u=\delta_{\alpha\beta}wu$ for all
  $\alpha,\beta\in I_1$ and $b\in I_2$, and the element $w$.
\item The bilinear form $B(x,y)=[x,y]/[h_\nu,h_\nu]$ on $V$, where $[\ ,\ ]$ is the Fock pairing
  (Definition~\ref{pd-weight-lattices}(2)).
\item For a prime $l$, the ring $\mathbf Z_{(l)}$ of the rational numbers whose denominator is prime to $l$,
  and the $\mathbf Z_{(l)}$-submodule $\Delta_L\subset W_0$ of the $u$ with $mu\in\Lambda$ for some integer $m$
  prime to $l$, where $\Lambda$ is the $\mathbf Z$-span of the polynomials obtained from $h_\nu$ by
  repeatedly applying the operators $E^{(k)}_{cd}$ with $p\le c,d<p+q$ and $k\geq0$
  (Definition~\ref{pd-weight-lattices}(1)).
\end{enumerate}
\end{definition}
\leanref{Formalism/PieriBlockIdeals/StatementLR.lean}{Formalism.CohomologicalDescent.LR.VQ}
\leanref{Formalism/PieriBlockIdeals/StatementLR.lean}{Formalism.CohomologicalDescent.LR.EV}
\leanref{Formalism/PieriBlockIdeals/PieriColumnMinors/RowGenDerivative.lean}{Formalism.CohomologicalDescent.LR.der}
\leanref{Formalism/PieriBlockIdeals/StatementLR.lean}{Formalism.CohomologicalDescent.LR.W0}
\leanref{Formalism/PieriBlockIdeals/StatementLR.lean}{Formalism.CohomologicalDescent.LR.rationalBlockModuleSetting}
\leanref{Formalism/PieriBlockIdeals/StatementLR.lean}{Formalism.CohomologicalDescent.LR.formB}
\leanref{Formalism/PieriBlockIdeals/PieriColumnMinors/LrcFockZeroRight.lean}{Formalism.CohomologicalDescent.LR.Zl}
\leanref{Formalism/PieriBlockIdeals/PieriColumnMinors/LrcFockZeroRight.lean}{Formalism.CohomologicalDescent.LR.locSpan}
\leanref{Formalism/PieriBlockIdeals/PieriColumnMinors/XiOfDepth2.lean}{Formalism.CohomologicalDescent.LR.LevZ}
\leanref{Formalism/PieriBlockIdeals/PieriColumnMinors/XiOfDepth2.lean}{Formalism.CohomologicalDescent.LR.MW}
\leanref{Formalism/PieriBlockIdeals/PieriColumnMinors/XiOfDepth2.lean}{Formalism.CohomologicalDescent.LR.DelLR}
% Lean: VQ p q w lam2 = dualWeyl Q (p + q) (p + q) (blockWt p q w lam2), EV a b = der (idx a) (idx b)
% restricted to it (idx = Sum.elim (Fin.castAdd q) (Fin.natAdd p)), W0, and rationalBlockModuleSetting;
% formB is fockQ divided by fockQ (hw) (hw). Lambda is LevZ (applyWord of liftWord p omega to hw), MW is
% W0 meet LevZ, and Delta_L is DelLR l = locSpan l (MW), over Zl l.

\begin{lemma}\label{is-onto-step-exp-bound-short-graded-rational-module-closed}
Let $N\geq1$, $\nu\in\mathbf N^N$ nonincreasing, and $M\subset\nabla_{\mathbf Q}(\nu)$ a $\mathbf Q$-subspace
stable under every operator $E^{(k)}_{ab}$ (Definition~\ref{pd-weight-lattices}(1)).
\begin{enumerate}
\item If $M\ne0$, then $h_\nu\in M$.
\item If $h_\nu\in M$, then $M=\nabla_{\mathbf Q}(\nu)$.
\end{enumerate}
\end{lemma}
\leanref{Formalism/GradedWeylDivisibility/StatementB.lean}{Formalism.GradedWeylDivisibility.HWInClosed}
\leanref{Formalism/GradedWeylDivisibility/FockSpaceTableaux/HWInClosed.lean}{Formalism.GradedWeylDivisibility.hwInClosed_of}
\leanref{Formalism/GradedWeylDivisibility/FockSpaceTableaux/LHwinLemmas.lean}{Formalism.GradedWeylDivisibility.lhwin_exists_killed}
\leanref{Formalism/GradedWeylDivisibility/FockSpaceTableaux/RowDivClosedGraded.lean}{Formalism.GradedWeylDivisibility.row_divided_closed_graded}
\leanref{Formalism/GradedWeylDivisibility/StatementB.lean}{Formalism.GradedWeylDivisibility.ClosedEqNabla}
\leanref{Formalism/GradedWeylDivisibility/FockSpaceTableaux/ClosedEqNabla.lean}{Formalism.GradedWeylDivisibility.closedEqNabla_of}
\begin{proof}
(1) First, $M$ contains the component of each of its elements of each row degree. Indeed, $E^{(1)}_{aa}$
multiplies a monomial by its degree in row $a$. Let $f\in M$, $t\geq0$, and $S$ a finite set of integers
containing the degrees in row $a$ of the monomials of $f$. Then
$\prod_{s\in S,s\ne t}(t-s)^{-1}(E^{(1)}_{aa}-s)f$ lies in $M$, and it is the sum of the terms of $f$ of
degree $t$ in row $a$. Doing this for one row after another gives the claim.

Call $\sum_aa\,\mu_a$ the height of a monomial of degree $\mu_a$ in row $a$ for every $a$; it is $\geq0$. As
$M\ne0$, it contains a nonzero polynomial of some row degree, which is homogeneous for the height. Let
$v\in M$ be nonzero and homogeneous for the height, of the least height $H$ among such elements of $M$. Let
$a<b$ and $k\geq1$. The operator $E^{(k)}_{ab}$ moves $k$ units of degree from row $b$ to row $a$, so
$E^{(k)}_{ab}v$ is homogeneous of height $H-k(b-a)<H$. It lies in $M$, so it is $0$ by the choice of $v$.
Thus $v\in\nabla_{\mathbf Q}(\nu)$ is killed by every $E^{(k)}_{ab}$ with $a<b$ and $k\geq1$. In
characteristic $0$, $\nabla_{\mathbf Q}(\nu)$ is the irreducible representation of highest weight $\nu$,
whose vectors killed by these operators form the line of $h_\nu$ \cite[II.2]{Jan03}. So $v=ch_\nu$ with
$c\ne0$, and $h_\nu=c^{-1}v\in M$.

(2) Let $M'$ be the set of $x\in\nabla_{\mathbf Q}(\nu)$ with $[x,g]=0$ for every $g\in M$, for the Fock
pairing. The operators $E^{(k)}_{ab}$ preserve $\nabla_{\mathbf Q}(\nu)$, as in the proof of
Lemma~\ref{pd-gain-weyl-comparison}, and $[E^{(k)}_{ab}x,g]=[x,E^{(k)}_{ba}g]$ by
Lemma~\ref{pd-kappa-gram-fock}(2); so $M'$ is stable under every $E^{(k)}_{ab}$. The pairing satisfies
$[f,f]=\sum_\kappa\kappa!\,f_\kappa^2>0$ for $f\ne0$. If $M'\ne0$, then (1) gives $h_\nu\in M'$, so
$[h_\nu,h_\nu]=0$ as $h_\nu\in M$, which is false as $h_\nu\ne0$. So $M'=0$. The pairing restricted to the
finite-dimensional space $M$ is nondegenerate. So for $f\in\nabla_{\mathbf Q}(\nu)$ there is $p\in M$ with
$[p,g]=[f,g]$ for every $g\in M$. Then $f-p\in M'=0$, and $f=p\in M$.
% Lean: (1) is hwInClosed = hwInClosed_of hwLine rowDivClosedGraded: lhwin_exists_killed
% (row_divided_closed_graded through lgraded_multi and lgraded_sep, the Lagrange interpolation in the
% Euler operator rowDiv a a 1; the height lhwinHt; Nat.find for the least height; lhwin_rowDiv_shift and
% lhwin_nonneg), then hwLine, cited here. (2) is closedEqNabla = closedEqNabla_of hwInClosed fockQPos
% fockQBilin fockQRowDivAdj: the orthogonal N with rowDivStable, hwNe, and BilinForm.toDual on M.
\end{proof}

\begin{lemma}\label{is-onto-step-exp-bound-short-graded-rational-module}
Let $p,q,w$, $\lambda_2$ and $\nu$ be as in Definition~\ref{is-onto-step-exp-bound-short-graded-rational},
with the rational block module $(V,E,W_0,w)$, the form $B$ and, for a prime $l$, the module $\Delta_L$.
\begin{enumerate}
\item $(V,E,W_0,w)$ is a block module over $\mathbf Q$ for $I_1$ and $I_2$
  (Definition~\ref{is-onto-step-exp-bound-short-graded-block}).
\item $B(E_{cd}x,y)=B(x,E_{dc}y)$ for all $c,d$ and $x,y\in V$, and $v=0$ if $B(x,v)=0$ for every $x\in V$.
\item $E_{bb'}\Delta_L\subset\Delta_L$ for all $b,b'\in I_2$.
\item For $\mu\in\mathbf N^{p+q}$, an $x\in V$ has row degree $\mu$ if and only if $x\in V_\mu$.
\item $h_\nu\in W_0$, and every subspace of $V$ that contains $W_0$ and is stable under every $E_{cd}$ is $V$.
\end{enumerate}
\end{lemma}
\leanref{Formalism/PieriBlockIdeals/PieriColumnMinors/RowGenDerivative.lean}{Formalism.CohomologicalDescent.LR.der_comm}
\leanref{Formalism/PieriBlockIdeals/PieriColumnMinors/RowGenDerivative.lean}{Formalism.CohomologicalDescent.LR.der_mem}
\leanref{Formalism/PieriBlockIdeals/StatementLR.lean}{Formalism.CohomologicalDescent.LR.w_zero_stable_right_block}
\leanref{Formalism/PieriBlockIdeals/StatementLR.lean}{Formalism.CohomologicalDescent.LR.RealContra}
\leanref{Formalism/PieriBlockIdeals/PieriColumnMinors/HighestWeightMemWZero.lean}{Formalism.CohomologicalDescent.LR.real_contragredient}
\leanref{Formalism/PieriBlockIdeals/StatementLR.lean}{Formalism.CohomologicalDescent.LR.RealNondeg}
\leanref{Formalism/PieriBlockIdeals/PieriColumnMinors/HighestWeightMemWZero.lean}{Formalism.CohomologicalDescent.LR.real_nondegenerate}
\leanref{Formalism/PieriBlockIdeals/PieriColumnMinors/RealOfSkeleton.lean}{Formalism.CohomologicalDescent.LR.delL_stable}
\leanref{Formalism/PieriBlockIdeals/StatementLR.lean}{Formalism.CohomologicalDescent.LR.RealWeight}
\leanref{Formalism/PieriBlockIdeals/PieriColumnMinors/HighestWeightMemWZero.lean}{Formalism.CohomologicalDescent.LR.real_weight}
\leanref{Formalism/PieriBlockIdeals/PieriColumnMinors/HighestWeightMemWZero.lean}{Formalism.CohomologicalDescent.LR.hwV_mem}
\leanref{Formalism/PieriBlockIdeals/PieriColumnMinors/HighestWeightMemWZero.lean}{Formalism.CohomologicalDescent.LR.der_hw_lower}
\leanref{Formalism/PieriBlockIdeals/StatementLR.lean}{Formalism.CohomologicalDescent.LR.RealGenerated}
\leanref{Formalism/PieriBlockIdeals/PieriColumnMinors/HighestWeightMemWZero.lean}{Formalism.CohomologicalDescent.LR.real_generated}
\leanref{Formalism/PieriBlockIdeals/PieriColumnMinors/HighestWeightMemWZero.lean}{Formalism.CohomologicalDescent.LR.top_of_stable}
\begin{proof}
(1) The row operators preserve $V$ and satisfy the $\mathfrak{gl}_{p+q}$ relations
(Definition~\ref{is-onto-step-exp-cut-grid}(6), Lemma~\ref{is-onto-step-exp-cut-module}(1)). The first two
conditions on $W_0$ hold by its definition. Let $u\in W_0$, $b,b'\in I_2$ and $u'=E_{bb'}u$. For $\alpha\in I_1$
and every $c$, the relations give $E_{\alpha c}u'=E_{bb'}E_{\alpha c}u+\delta_{cb}E_{\alpha b'}u-\delta_{b'\alpha}
E_{bc}u$, and $\delta_{b'\alpha}=0$. For $c\in I_2$ all terms vanish; for $c=\beta\in I_1$, $\delta_{\beta b}=0$ and
$E_{\alpha\beta}u'=E_{bb'}(\delta_{\alpha\beta}wu)=\delta_{\alpha\beta}wu'$. So $u'\in W_0$.

(2) By Lemma~\ref{pd-kappa-gram-fock}(2) with $k=1$, $[E_{cd}x,y]=[x,E_{dc}y]$; and $[h_\nu,h_\nu]$ is a positive
integer by Lemma~\ref{pd-kappa-gram-fock}(3). If $B(x,v)=0$ for every $x$, then $[v,v]=0$; as $[v,v]$ is the
sum of the squares of the coefficients of $v$ weighted by positive integers, $v=0$.

(3) Let $u\in\Delta_L$, with $mu\in\Lambda$ and $m$ prime to $l$. Then $E_{bb'}u\in W_0$ by (1), and
$mE_{bb'}u=E^{(1)}_{p+b,p+b'}(mu)\in\Lambda$, as $\Lambda$ is stable under these operators.

(4) $E_{cc}$ multiplies each monomial by its degree in row $c$. So $x$ has row degree $\mu$ if and only if
$E_{cc}x=\mu_cx$ for every $c$.

(5) The polynomial $h_\nu$ has degree $\nu_c$ in row $c$, so $E_{\alpha\alpha}h_\nu=wh_\nu$. It is killed by
the $E_{cd}$ with $c<d$, as each leading principal minor is: such a minor has two equal rows after the
substitution if row $d$ is among its rows, and does not involve row $d$ otherwise. This gives
$E_{\alpha b}h_\nu=0$, and $E_{\alpha\beta}h_\nu=0$ for $\alpha<\beta$. Let $\alpha>\beta$ and $y=E_{\alpha\beta}h_\nu$. By
(2), the relations and $E_{\beta\alpha}h_\nu=0$,
\[
 [y,y]=[h_\nu,E_{\beta\alpha}E_{\alpha\beta}h_\nu]=[h_\nu,E_{\alpha\beta}E_{\beta\alpha}h_\nu+(E_{\beta\beta}-E_{\alpha\alpha})h_\nu]
 =(\nu_\beta-\nu_\alpha)[h_\nu,h_\nu]=0,
\]
so $y=0$. Hence $h_\nu\in W_0$. Now let $U\supset W_0$ be stable under every $E_{cd}$. By induction on $k$,
using $(k+1)E^{(k+1)}_{cd}=E^{(k)}_{cd}E_{cd}-\delta_{cd}kE^{(k)}_{cd}$, $U$ is stable under every
$E^{(k)}_{cd}$. As $h_\nu\in W_0\subset U$ and $h_\nu\ne0$, $U\ne0$. So
Lemma~\ref{is-onto-step-exp-bound-short-graded-rational-module-closed}(1) gives $h_\nu\in U$, and part (2) of
that lemma gives $U=V$.
% Lean: (1) rationalBlockModuleSetting (EV_comm from der_comm, der_mem, mem_orthogonal_frame_iff,
% w_zero_stable_right_block); (2) real_contragredient (fock_der, which is fockQRowDivAdj at k = 1) and
% real_nondegenerate (fockQPos, hw_norm_pos); (3) delL_stable (levZ_rowDiv); (4) real_weight (der_diag_iff,
% coeff_der_diag); (5) hwV_mem (der_hw_raise from hwRaise, hw_rowHom, der_hw_lower) and real_generated
% (U /= 0 by hwV_mem and hwV_ne) = top_of_stable (closed_of_stable through closed_of_der with rowDiv_succ,
% then hwInClosed and closedEqNabla, Lemma is-onto-step-exp-bound-short-graded-rational-module-closed).
\end{proof}

\begin{lemma}\label{is-onto-step-exp-bound-short-graded-rational-lattice}
Let $p,q,w$, $\lambda_2$ and $\nu$ be as in Definition~\ref{is-onto-step-exp-bound-short-graded-rational},
with the rational block module $(V,E,W_0,w)$, the form $B$, a prime $l$ and the module $\Delta_L$, and let
$\Delta=\Delta(\Delta_L)\subset V$ for $O=\mathbf Z_{(l)}$
(Definition~\ref{is-onto-step-exp-bound-short-graded-block}(2)).
\begin{enumerate}
\item Every element of $\Delta$ is $y/m$ for some $y\in\Delta_{\mathbf Z}(\nu)$ and some integer $m$ prime to
  $l$.
\item $\nabla_{\mathbf Z}(\nu)\subset\Delta^B$.
\item If $x\in\nabla_{\mathbf Z}(\nu)$ and $c\geq0$ with $l^cx\in\Delta$, then $ml^cx\in\Delta_{\mathbf Z}(\nu)$ for
  some integer $m$ prime to $l$.
\end{enumerate}
\end{lemma}
\leanref{Formalism/PieriBlockIdeals/PieriColumnMinors/LriEVPowCoe.lean}{Formalism.CohomologicalDescent.LR.deltaAll_le}
\leanref{Formalism/PieriBlockIdeals/PieriColumnMinors/FsDivMem.lean}{Formalism.CohomologicalDescent.LR.Fs_div_mem}
\leanref{Formalism/PieriBlockIdeals/PieriColumnMinors/LriEVPowCoe.lean}{Formalism.CohomologicalDescent.LR.mem_dual}
\leanref{Formalism/PieriBlockIdeals/PieriColumnMinors/LriEVPowCoe.lean}{Formalism.CohomologicalDescent.LR.concl}
\begin{proof}
(1) The set $L'$ of the $y/m$ is a $\mathbf Z_{(l)}$-submodule of $V$. It contains $\Lambda$, as
$\Lambda\subset\Delta_{\mathbf Z}(\nu)$, and hence $\Delta_L$. For $v=(b,\alpha)$ and $k\geq0$,
$(k!)^{-1}f_v^k$ is the operator $E^{(k)}_{p+b,\alpha}$, as $p+b\ne\alpha$
(Lemma~\ref{is-onto-step-exp-cut-module}(1)); it preserves $\Delta_{\mathbf Z}(\nu)$, hence $L'$. As the
$f_v$ commute, $(m!)^{-1}F_m$ is the product over $v$ of the $(k_v!)^{-1}f_v^{k_v}$, with $k_v$ the
multiplicity of $v$ in $m$, so it preserves $L'$. Hence $\Delta\subset L'$.

(2) Let $x\in\nabla_{\mathbf Z}(\nu)$ and $z\in\Delta$, and write $z=y/m$ by (1). By
Lemma~\ref{is-onto-step-exp-cut-levi-duality}(1), $[y,x]=a[h_\nu,h_\nu]$ for some integer $a$, so
$B(x,z)=a/m\in\mathbf Z_{(l)}$.

(3) By (1), $l^cx=y/m$ with $y\in\Delta_{\mathbf Z}(\nu)$ and $m$ prime to $l$, so $ml^cx=y$.
% Lean: (1) deltaAll_le (MW_le_MV with lri_LevZ_le, lri_hL with der_pow, Fs_div_mem through lrd_core);
% (2) mem_dual, which uses fock_int (fock_word: the pairing with a word applied to hw, through
% fock_space_word_adjoint and fock_space_highest_weight) in place of the duality lemma; (3) concl, with
% lri_intQ_inj for the injectivity of Z[x] -> Q[x].
\end{proof}

\begin{lemma}\label{is-onto-step-exp-bound-short-graded-rational-pieri-irreducible}
Let $p,q,w$, $\lambda_2$ and $\nu$ be as in Definition~\ref{is-onto-step-exp-bound-short-graded-rational}, with the
rational block module $(V,E,W_0,w)$. Let $Z\subset W_0$ be a nonzero subspace with $E_{bb'}Z\subset Z$ for all
$b,b'\in I_2$. Then $h_\nu\in Z$.
\end{lemma}
\leanref{Formalism/PieriBlockIdeals/StatementLR.lean}{Formalism.CohomologicalDescent.LR.RealIrred}
\leanref{Formalism/PieriBlockIdeals/PieriColumnMinors/HighestWeightMemWZero.lean}{Formalism.CohomologicalDescent.LR.real_irreducible}
\begin{proof}
The data $(V,E,Z,w)$ form a block module over $\mathbf Q$ for $I_1$ and $I_2$: the conditions on $Z$ hold as
$Z\subset W_0$ and $Z$ is stable under the $E_{bb'}$. It satisfies the generation condition of
Lemma~\ref{is-onto-step-exp-bound-short-graded-skeleton-mu}(5). Indeed, let $U\supset Z$ be a subspace stable under
every $E_{cd}$. As in the proof of Lemma~\ref{is-onto-step-exp-bound-short-graded-rational-module}(5), $U$ is
stable under every $E^{(k)}_{cd}$, and $U\ne0$ as $Z\ne0$; so
Lemma~\ref{is-onto-step-exp-bound-short-graded-rational-module-closed} gives $U=V$. By
Lemma~\ref{is-onto-step-exp-bound-short-graded-rational-module}(4), $h_\nu\in V_\nu$, and
$\sum_{\alpha<p}\nu_\alpha+0=pw$. So Lemma~\ref{is-onto-step-exp-bound-short-graded-skeleton-mu}(5) with $n=0$ and part
(2) of that lemma give $h_\nu\in\mu_0(T_0)=Z$.
% Lean: real_irreducible: the HidaSetting SZ with W = Z, hgen through top_of_stable, hwt from real_weight and
% hw_rowHom, then depth_of_weight_holds at n = 0 and the first part of depth_span_holds.
\end{proof}

\begin{lemma}\label{is-onto-step-exp-bound-short-graded-rational-pieri-highest}
Let $K$ be an ordered field, $N$ a $K$-vector space, $q\geq0$, $T_{ab}$ ($a,b<q$) endomorphisms of $N$ with
$T_{ab}T_{cd}-T_{cd}T_{ab}=\delta_{bc}T_{ad}-\delta_{da}T_{cb}$, and $B$ a symmetric bilinear form on $N$. Suppose that
$B(x,x)>0$ for $x\ne0$, that $B(T_{ab}x,y)=B(x,T_{ba}y)$ for all $a,b,x,y$, and that every $T_{aa}$ is
diagonalizable: $N$ is the sum of its eigenspaces. A \emph{highest weight vector} of weight $\mu\in K^q$ is a
$v\in N$ with $T_{ab}v=0$ for $a<b$ and $T_{aa}v=\mu_av$ for every $a$; a subspace is \emph{stable} if it is stable
under every $T_{ab}$; and $\mathrm{Cas}(T)=\sum_{b,b'}T_{bb'}T_{b'b}$.
\begin{enumerate}
\item If $v\ne0$ is a highest weight vector of weight $\mu$, then $\mu_b\le\mu_a$ for $a<b$.
\item If $N$ is finite-dimensional, every nonzero stable subspace contains a nonzero highest weight vector.
\item If $N$ is finite-dimensional, $S$ and $S'$ are stable, and every highest weight vector in $S$ lies in $S'$,
  then $S\subset S'$.
\item Let $\kappa\in K^q$. If every stable subspace that contains all the highest weight vectors of weight
  $\kappa$ is $N$, then every nonzero highest weight vector has weight $\kappa$.
\item $\mathrm{Cas}(T)$ commutes with every $T_{cd}$. If $v\ne0$ is a highest weight vector of weight
  $\kappa\in\mathbf N^q$, then $\kappa$ is nonincreasing and $\mathrm{Cas}(T)v=\mathrm{cas}(\kappa)v$
  (Definition~\ref{is-onto-step-exp-bound-short-graded-block}(4)).
\item Suppose that $N$ is finite-dimensional, and that $h\ne0$ lies in every nonzero stable subspace. Then every
  stable subspace that contains $h$ is $N$.
\end{enumerate}
\end{lemma}
\leanref{Formalism/PieriBlockIdeals/StatementLP.lean}{Formalism.CohomologicalDescent.LP.IsHW}
\leanref{Formalism/PieriBlockIdeals/StatementLP.lean}{Formalism.CohomologicalDescent.LP.FormMod}
\leanref{Formalism/PieriBlockIdeals/StatementLP.lean}{Formalism.CohomologicalDescent.LP.casOp}
\leanref{Formalism/PieriBlockIdeals/StatementLP.lean}{Formalism.CohomologicalDescent.LP.HWDominant}
\leanref{Formalism/PieriBlockIdeals/PieriColumnMinors/HighestWeightDominantHolds.lean}{Formalism.CohomologicalDescent.LP.highestWeightDominant_holds}
\leanref{Formalism/PieriBlockIdeals/StatementLP.lean}{Formalism.CohomologicalDescent.LP.ExistsHW}
\leanref{Formalism/PieriBlockIdeals/PieriColumnMinors/HighestWeightDominantHolds.lean}{Formalism.CohomologicalDescent.LP.existsHighestWeight_holds}
\leanref{Formalism/PieriBlockIdeals/StatementLP.lean}{Formalism.CohomologicalDescent.LP.CompleteRed}
\leanref{Formalism/PieriBlockIdeals/PieriColumnMinors/HighestWeightDominantHolds.lean}{Formalism.CohomologicalDescent.LP.completeReducibility_holds}
\leanref{Formalism/PieriBlockIdeals/StatementLP.lean}{Formalism.CohomologicalDescent.LP.GeneratedHW}
\leanref{Formalism/PieriBlockIdeals/PieriColumnMinors/GeneratedHighestWeightHolds.lean}{Formalism.CohomologicalDescent.LP.generatedHighestWeight_holds}
\leanref{Formalism/PieriBlockIdeals/StatementLP.lean}{Formalism.CohomologicalDescent.LP.CasCentral}
\leanref{Formalism/PieriBlockIdeals/PieriColumnMinors/HighestWeightSuccessive.lean}{Formalism.CohomologicalDescent.LP.casimir_central'}
\leanref{Formalism/PieriBlockIdeals/PieriColumnMinors/HighestWeightSuccessive.lean}{Formalism.CohomologicalDescent.LP.lp_hw_cas}
\leanref{Formalism/PieriBlockIdeals/StatementLPR.lean}{Formalism.CohomologicalDescent.LPR.IrredGen}
\leanref{Formalism/PieriBlockIdeals/PieriColumnMinors/IrreducibleGeneration.lean}{Formalism.CohomologicalDescent.LPR.irreducible_generation}
\begin{proof}
(1) The relations give $(T_{ab}T_{ba}-T_{ba}T_{ab})v=(T_{aa}-T_{bb})v$, and $T_{ab}v=0$, so
$T_{ab}T_{ba}v=(\mu_a-\mu_b)v$. Hence $0\le B(T_{ba}v,T_{ba}v)=B(v,T_{ab}T_{ba}v)=(\mu_a-\mu_b)B(v,v)$, with
$B(v,v)>0$.

(2) The $T_{aa}$ commute, by the relations, and are diagonalizable; a nonzero stable subspace $C$ is stable under
them, so, intersecting with their eigenspaces one at a time, $C$ contains a nonzero common eigenvector. The set
$W$ of the weights of the nonzero common eigenvectors of the $T_{aa}$ in $C$ is finite, as each $T_{aa}$ has
finitely many eigenvalues, and nonempty. Choose $\mu\in W$ with the largest value of $\sum_c-c\,\mu_c$, and $v\in C$
a nonzero common eigenvector of weight $\mu$. If $a<b$ and $T_{ab}v\ne0$, then by the relations $T_{ab}v\in C$ is a
common eigenvector of weight $\mu+\epsilon_a-\epsilon_b$, for which the value is larger by $b-a>0$; this is
impossible. So $v$ is a highest weight vector.

(3) The subspace $S\cap S'$ is stable, and so is $U=S\cap(S\cap S')^\perp$, as $B(T_{ab}x,y)=B(x,T_{ba}y)$. If
$U\ne0$, (2) gives a nonzero highest weight vector $v\in U$; then $v\in S$, so $v\in S'$ and $v\in S\cap S'$, and
$B(v,v)=0$, which is false. So $U=0$. As $B$ is positive definite, $N=(S\cap S')\oplus(S\cap S')^\perp$. Let $x\in S$
and write $x=y+z$ accordingly; then $z=x-y\in U=0$, and $x=y\in S'$.

(4) Let $L$ be the span of the vectors obtained from highest weight vectors of weight $\kappa$ by applying operators
$T_{ab}$ with $b<a$ repeatedly. By induction on the number of operators applied, using the relations, $L$ is stable;
it contains every highest weight vector of weight $\kappa$, so $L=N$. Let $v\ne0$ be a highest weight vector of
weight $\mu$, and suppose $\mu_a\ne\kappa_a$. For a highest weight vector $u$ of weight $\kappa$,
$\mu_aB(v,u)=B(T_{aa}v,u)=B(v,T_{aa}u)=\kappa_aB(v,u)$, so $B(v,u)=0$; and for $d<c$ and $y\in N$,
$B(v,T_{cd}y)=B(T_{dc}v,y)=0$. So $B(v,L)=0$, and $B(v,v)=0$, which is false.

(5) The first claim is a direct computation from the relations: the four sums that arise in
$\mathrm{Cas}(T)T_{cd}-T_{cd}\mathrm{Cas}(T)$ cancel in pairs. Let $v$ be as stated; $\kappa$ is nonincreasing by
(1). For $b<b'$, $T_{bb'}T_{b'b}v=T_{b'b}T_{bb'}v+(T_{bb}-T_{b'b'})v=(\kappa_b-\kappa_{b'})v$; for $b=b'$ the term is
$\kappa_b^2v$; and for $b>b'$, $T_{b'b}v=0$. So
$\mathrm{Cas}(T)v=\big(\sum_b\kappa_b^2+\sum_{b<b'}(\kappa_b-\kappa_{b'})\big)v
=\sum_b\kappa_b(\kappa_b+q-1-2b)\,v=\mathrm{cas}(\kappa)v$.

(6) Let $S$ be stable with $h\in S$. Then $S^\perp$ is stable. If $S^\perp\ne0$, then $h\in S^\perp$ and
$B(h,h)=0$, which is false. So $S^\perp=0$, and $N=S\oplus S^\perp=S$ as $B$ is positive definite.
% Lean: (1) highestWeightDominant_holds; (2) existsHighestWeight_holds (lpa_joint with
% Submodule.inf_iSup_genEigenspace, Set.exists_max_image for the weighted sum, weighted_sum_add_single_sub_single);
% (3) completeReducibility_holds; (4) generatedHighestWeight_holds (LowW, lowSpan, lowW_all, lowSpan_stab); (5)
% casimir_central' and lp_hw_cas; (6) irreducible_generation (B.orthogonal S,
% isCompl_orthogonal_of_restrict_nondegenerate).
\end{proof}

\begin{lemma}\label{is-onto-step-exp-bound-short-graded-rational-pieri-tower}
Let $K$ be an ordered field and $N$ a $K$-vector space.
\begin{enumerate}
\item Let $q\geq0$, and $T_{ab}$ and $\varepsilon_{ab}$ ($a,b<q$) endomorphisms of $N$ with
  $\varepsilon_{ab}\varepsilon_{cd}=\delta_{bc}\varepsilon_{ad}$, $\sum_a\varepsilon_{aa}=1$, and every $T_{ab}$
  commuting with every $\varepsilon_{cd}$. If $z\ne0$ is a highest weight vector of weight $\mu$ for the operators
  $T_{ab}+\varepsilon_{ab}$, then there is $b$ with $\varepsilon_{bb}z\ne0$, and $\varepsilon_{bb}z$ is a highest weight
  vector of weight $\mu-\epsilon_b$ for the $T_{ab}$.
\item Let $N$ be finite-dimensional, $n\geq0$, and for $j\le n$ and $j'<n$ let $D_j=(D_j(a,b))$ and
  $\varepsilon_{j'}=(\varepsilon_{j'}(a,b))$ be families of endomorphisms of $N$, and $B$ a bilinear form on $N$, such
  that: each $(N,D_j,B)$ satisfies the hypotheses of
  Lemma~\ref{is-onto-step-exp-bound-short-graded-rational-pieri-highest}; $D_{j+1}(a,b)=D_j(a,b)+\varepsilon_j(a,b)$;
  each $\varepsilon_j$ satisfies the conditions of (1); and $D_i(a,b)$ commutes with $\varepsilon_j(c,d)$ for $i\le j$.
  Let $\lambda_2\in\mathbf N^q$, and suppose that every nonzero highest weight vector for $D_0$ has weight
  $\lambda_2$. Then $N$ is spanned by the $z$ for which some box path $\kappa$ of length $n$ from $\lambda_2$ has
  $\mathrm{Cas}(D_j)z=\mathrm{cas}(\kappa_j)z$ for every $j\le n$.
\end{enumerate}
\end{lemma}
\leanref{Formalism/PieriBlockIdeals/StatementLP.lean}{Formalism.CohomologicalDescent.LP.StepHW}
\leanref{Formalism/PieriBlockIdeals/PieriColumnMinors/HighestWeightSuccessive.lean}{Formalism.CohomologicalDescent.LP.highest_weight_successive}
\leanref{Formalism/PieriBlockIdeals/StatementLP.lean}{Formalism.CohomologicalDescent.LP.Tower}
\leanref{Formalism/PieriBlockIdeals/StatementLP.lean}{Formalism.CohomologicalDescent.LP.TowerPieri}
\leanref{Formalism/PieriBlockIdeals/PieriColumnMinors/HighestWeightSuccessive.lean}{Formalism.CohomologicalDescent.LP.TowerPieri_of}
\leanref{Formalism/PieriBlockIdeals/PieriColumnMinors/HighestWeightSuccessive.lean}{Formalism.CohomologicalDescent.LP.lp_claim}
\leanref{Formalism/PieriBlockIdeals/PieriColumnMinors/HighestWeightSuccessive.lean}{Formalism.CohomologicalDescent.LP.lp_key}
\begin{proof}
(1) As $z=\sum_a\varepsilon_{aa}z\ne0$, some $\varepsilon_{aa}z$ is nonzero; let $b$ be the largest such $a$. For
$a<c$, $T_{ac}z=-\varepsilon_{ac}z$, so $T_{ac}\varepsilon_{bb}z=\varepsilon_{bb}T_{ac}z=-\varepsilon_{bb}\varepsilon_{ac}z$,
which is $0$ unless $a=b$, and then it is $-\varepsilon_{bc}z=-\varepsilon_{bc}\varepsilon_{cc}z=0$ as $c>b$. And
$T_{aa}z=\mu_az-\varepsilon_{aa}z$, so $T_{aa}\varepsilon_{bb}z=\mu_a\varepsilon_{bb}z-\delta_{ab}\varepsilon_{bb}z$.

(2) For a sequence $\kappa_0,\dots,\kappa_j$ let $\mathrm{Eig}_j(\kappa)$ be the set of the $z$ with
$\mathrm{Cas}(D_i)z=\mathrm{cas}(\kappa_i)z$ for $i\le j$, and call $\kappa$ admissible if it is a box path from
$\lambda_2$ of length $j$. By Lemma~\ref{is-onto-step-exp-bound-short-graded-rational-pieri-highest}(5) and induction,
$\mathrm{Cas}(D_i)$ commutes with $D_{j'}(c,d)$ for $i\le j'\le n$ and with $\varepsilon_{j'}(c,d)$ for $i\le j'<n$; so
$\mathrm{Eig}_j(\kappa)$ is stable under $D_{j'}$ for $j\le j'\le n$ and under $\varepsilon_j(b,b)$. We show by induction
on $j\le n$: (i) $N$ is the sum of the $\mathrm{Eig}_j(\kappa)$, $\kappa$ admissible; and (ii) for admissible $\kappa$,
every nonzero $z\in\mathrm{Eig}_j(\kappa)$ that is a highest weight vector for $D_j$ has weight $\kappa_j$. For $j=n$,
(i) is the claim.

First, a consequence of (ii) at $j<n$. Let $\kappa$ be admissible of length $j$, and $z\in\mathrm{Eig}_j(\kappa)$ nonzero
and a highest weight vector for $D_{j+1}$ of weight $\mu$. By (1) applied to $D_j$ and $\varepsilon_j$, some
$\varepsilon_j(b,b)z\ne0$ is a highest weight vector for $D_j$ of weight $\mu-\epsilon_b$; it lies in
$\mathrm{Eig}_j(\kappa)$, so $\mu=\kappa_j+\epsilon_b$ by (ii). By
Lemma~\ref{is-onto-step-exp-bound-short-graded-rational-pieri-highest}(5) for $D_{j+1}$, $\kappa_j+\epsilon_b$ is
nonincreasing and $\mathrm{Cas}(D_{j+1})z=\mathrm{cas}(\kappa_j+\epsilon_b)z$.

For $j=0$, (ii) is the hypothesis. If $\lambda_2$ is nonincreasing, apply
Lemma~\ref{is-onto-step-exp-bound-short-graded-rational-pieri-highest}(3) to $S=N$ and $S'=\mathrm{Eig}_0(\lambda_2)$:
a nonzero highest weight vector for $D_0$ has weight $\lambda_2$, so lies in $S'$ by part (5) of that lemma; this gives
(i). Otherwise, part (5) shows that $D_0$ has no nonzero highest weight vector, and part (3) with $S'=0$ gives $N=0$.

Let $j<n$. For (i), let $\kappa$ be admissible of length $j$, and apply
Lemma~\ref{is-onto-step-exp-bound-short-graded-rational-pieri-highest}(3) for $D_{j+1}$ to $S=\mathrm{Eig}_j(\kappa)$
and $S'$ the sum of the $\mathrm{Eig}_{j+1}(\kappa')$, $\kappa'$ admissible of length $j+1$. A nonzero highest weight
vector $v\in S$ for $D_{j+1}$ has, by the consequence above, $\mathrm{Cas}(D_{j+1})v=\mathrm{cas}(\kappa_j+\epsilon_b)v$
with $\kappa_j+\epsilon_b$ nonincreasing; so $v\in\mathrm{Eig}_{j+1}(\kappa')$ for $\kappa'$ the extension of $\kappa$ by
$\kappa_j+\epsilon_b$, which is admissible. So $S\subset S'$, and (i) for $j+1$ follows from (i) for $j$. For (ii), let
$\kappa$ be admissible of length $j+1$, with $\kappa_{j+1}=\kappa_j+\epsilon_\rho$, and $z\in\mathrm{Eig}_{j+1}(\kappa)$
a nonzero highest weight vector for $D_{j+1}$ of weight $\mu$. By the consequence above, $\mu=\kappa_j+\epsilon_b$ and
$\mathrm{cas}(\kappa_j+\epsilon_b)z=\mathrm{Cas}(D_{j+1})z=\mathrm{cas}(\kappa_j+\epsilon_\rho)z$. By
Lemma~\ref{is-onto-step-exp-bound-short-graded-skeleton-jucys-relations}(7),
$\kappa_{j,b}-b=\kappa_{j,\rho}-\rho$, and as $\kappa_j$ is nonincreasing, $b=\rho$. So $\mu=\kappa_{j+1}$.
% Lean: (1) highest_weight_successive; (2) TowerPieri_of completeReducibility_holds highestWeightDominant_holds,
% through lp_claim (by induction on j: PathTo, Eig, lp_Eig_stab with lp_cas_D and lp_cas_eps, lp_key, lp_hw_cas,
% lp_casVal_step and lp_content_inj).
\end{proof}

\begin{lemma}\label{is-onto-step-exp-bound-short-graded-rational-pieri-general}
Let $K$ be an ordered field, $I_1$ a finite set, $q\geq0$, $I_2=\{0,\dots,q-1\}$, $M$ a finite-dimensional $K$-vector
space with endomorphisms $E_{bb'}$ ($b,b'\in I_2$), and $B$ a bilinear form on $M$, such that $(M,E,B)$ satisfies the
hypotheses of Lemma~\ref{is-onto-step-exp-bound-short-graded-rational-pieri-highest}. Let $\lambda_2\in\mathbf N^q$,
and suppose that every $E$-stable subspace of $M$ that contains all the highest weight vectors of weight $\lambda_2$
is $M$. Then for every $n\geq0$, $T_n$ is the sum of the box path spaces $T_n(\lambda_2,\tau)$, $\tau\in\mathbf N^q$
(Definition~\ref{is-onto-step-exp-bound-short-graded-skeleton-jucys}).
\end{lemma}
\leanref{Formalism/PieriBlockIdeals/StatementLP.lean}{Formalism.CohomologicalDescent.LP.PieriGenerated}
\leanref{Formalism/PieriBlockIdeals/PieriColumnMinors/PieriGenerated.lean}{Formalism.CohomologicalDescent.LP.pieri_generated}
\leanref{Formalism/PieriBlockIdeals/StatementLP.lean}{Formalism.CohomologicalDescent.LP.Pieri}
\leanref{Formalism/PieriBlockIdeals/PieriColumnMinors/PieriGenerated.lean}{Formalism.CohomologicalDescent.LP.pieri_rule}
\leanref{Formalism/PieriBlockIdeals/StatementLP.lean}{Formalism.CohomologicalDescent.LP.tensForm}
\leanref{Formalism/PieriBlockIdeals/StatementLP.lean}{Formalism.CohomologicalDescent.LP.TensTower}
\leanref{Formalism/PieriBlockIdeals/PieriColumnMinors/LpcDpartCommEps.lean}{Formalism.CohomologicalDescent.LP.tensor_tower}
\leanref{Formalism/PieriBlockIdeals/StatementLP.lean}{Formalism.CohomologicalDescent.LP.TensBase}
\leanref{Formalism/PieriBlockIdeals/PieriColumnMinors/LpcDpartCommEps.lean}{Formalism.CohomologicalDescent.LP.tensor_base_case}
\begin{proof}
By Lemma~\ref{is-onto-step-exp-bound-short-graded-rational-pieri-highest}(4), every nonzero highest weight vector of
$M$ has weight $\lambda_2$. On $T_n$, which is finite-dimensional, put $B_n(x,y)=\sum_gB(x(g),y(g))$, and for $j<n$ let
$\varepsilon_j(a,b)=e^{(j)}_{ab}$. Then $(T_n,D_j,B_n)$ satisfies the hypotheses of
Lemma~\ref{is-onto-step-exp-bound-short-graded-rational-pieri-highest} for every $j\le n$. The relations: $D_0$ is $E$
acting through the values, and $D_{j+1}=D_j+\varepsilon_j$, where $\varepsilon_j$ satisfies the relations as
$e^{(j)}_{ab}e^{(j)}_{cd}=\delta_{bc}e^{(j)}_{ad}$ and commutes with $D_j$; and the sum of two commuting families that
satisfy the relations satisfies them. $B_n$ is symmetric and positive definite as $B$ is. Contravariance: $B_n(E_{ab}x,y)
=B_n(x,E_{ba}y)$ termwise, and $B_n(e^{(s)}_{ab}x,y)=B_n(x,e^{(s)}_{ba}y)$ after the change of variable that replaces the
entry $a$ at the position $s$ by $b$. Diagonalizability: $(D_j(a,a)z)(g)=E_{aa}z(g)+c\,z(g)$, with $c$ the number of
$s<j$ whose entry of $I_2$ is $a$; so on the functions supported at one $g$, $D_j(a,a)$ is $E_{aa}$ plus a constant.
Moreover $\varepsilon_j(a,b)\varepsilon_j(c,d)=\delta_{bc}\varepsilon_j(a,d)$, $\sum_a\varepsilon_j(a,a)=1$, and $D_i$ commutes
with $\varepsilon_j$ for $i\le j$, as they act on the values and on different positions. A nonzero highest weight
vector $z$ for $D_0$ has a value $z(g)\ne0$, which is a highest weight vector of $M$ of the same weight; so its weight
is $\lambda_2$. Lemma~\ref{is-onto-step-exp-bound-short-graded-rational-pieri-tower}(2) now shows that $T_n$ is
spanned by the $z$ for which a box path $\kappa$ of length $n$ from $\lambda_2$ has $C_jz=\mathrm{cas}(\kappa_j)z$ for
every $j\le n$, as $C_j=\mathrm{Cas}(D_j)$; such a $z$ lies in $T_n(\lambda_2,\kappa_n)$.
% Lean: pieri_generated = pieri_rule with generatedHighestWeight_holds; pieri_rule: tensor_tower (lpc_formMod:
% lpc_Dpart_glrel with lpc_glrel_add and lpc_eps_glrel, lpc_form_apply, lpc_form_unitY, lpc_Dpart_diag; lpc_Dpart_succ,
% lpc_eps_mul, lpc_eps_one, lpc_Dpart_comm_eps), tensor_base_case, and tower_pieri = TowerPieri_of
% completeReducibility_holds highestWeightDominant_holds.
\end{proof}

\begin{lemma}\label{is-onto-step-exp-bound-short-graded-rational-pieri-real}
Let $p,q,w$, $\lambda_2$ and $\nu$ be as in Definition~\ref{is-onto-step-exp-bound-short-graded-rational}, with the
rational block module $(V,E,W_0,w)$ and the form $B$.
\begin{enumerate}
\item $V$ is finite-dimensional.
\item $W_0$ with the operators $E_{bb'}$, $b,b'\in I_2$, and the restriction of $B$ satisfies the hypotheses of
  Lemma~\ref{is-onto-step-exp-bound-short-graded-rational-pieri-highest}.
\item $h_\nu$ is a nonzero highest weight vector of $W_0$ of weight $\lambda_2$.
\item Every subspace of $W_0$ that is stable under the $E_{bb'}$ and contains all the highest weight vectors of weight
  $\lambda_2$ is $W_0$.
\end{enumerate}
\end{lemma}
\leanref{Formalism/PieriBlockIdeals/StatementLPR.lean}{Formalism.CohomologicalDescent.LPR.RealFinite}
\leanref{Formalism/PieriBlockIdeals/PieriColumnMinors/IrreducibleGeneration.lean}{Formalism.CohomologicalDescent.LPR.real_form_finite}
\leanref{Formalism/PieriBlockIdeals/StatementLPR.lean}{Formalism.CohomologicalDescent.LPR.RealFormMod}
\leanref{Formalism/PieriBlockIdeals/PieriColumnMinors/IrreducibleGeneration.lean}{Formalism.CohomologicalDescent.LPR.real_form_modular}
\leanref{Formalism/PieriBlockIdeals/StatementLPR.lean}{Formalism.CohomologicalDescent.LPR.DiagTransfer}
\leanref{Formalism/PieriBlockIdeals/PieriColumnMinors/IrreducibleGeneration.lean}{Formalism.CohomologicalDescent.LPR.diagonal_transfer}
\leanref{Formalism/PieriBlockIdeals/StatementLPR.lean}{Formalism.CohomologicalDescent.LPR.EulerDiag}
\leanref{Formalism/PieriBlockIdeals/PieriColumnMinors/IrreducibleGeneration.lean}{Formalism.CohomologicalDescent.LPR.euler_diagonal}
\leanref{Formalism/PieriBlockIdeals/StatementLPR.lean}{Formalism.CohomologicalDescent.LPR.RealHW}
\leanref{Formalism/PieriBlockIdeals/PieriColumnMinors/IrreducibleGeneration.lean}{Formalism.CohomologicalDescent.LPR.real_highest_weight}
\leanref{Formalism/PieriBlockIdeals/StatementLPR.lean}{Formalism.CohomologicalDescent.LPR.RealGen}
\leanref{Formalism/PieriBlockIdeals/PieriColumnMinors/IrreducibleGeneration.lean}{Formalism.CohomologicalDescent.LPR.real_form_generation}
\begin{proof}
(1) $V$ lies in the space of polynomials of degree $\nu_s$ in column $s$ for every $s$, which has a finite basis of
monomials.

(2) The relations hold by Lemma~\ref{is-onto-step-exp-bound-short-graded-rational-module}(1). $B$ is symmetric, and
$B(x,x)=[x,x]/[h_\nu,h_\nu]>0$ for $x\ne0$, as in the proof of part (2) of that lemma; contravariance is part (2). For
diagonalizability, $E_{p+b,p+b}$ multiplies each monomial by its degree in row $p+b$, so the polynomial ring is the sum
of its eigenspaces. Let $x\in W_0$, and write $x=\sum_{\mu\in F}x_\mu$ with $F$ finite and $E_{p+b,p+b}x_\mu=\mu x_\mu$.
These $x_\mu$ lie in $W_0$: by induction on $\#F$, applying $E_{p+b,p+b}-\mu_0$ for one $\mu_0\in F$, which preserves
$W_0$, gives $\sum_{\mu\ne\mu_0}(\mu-\mu_0)x_\mu\in W_0$, so the $x_\mu$ with $\mu\ne\mu_0$ lie in $W_0$, and then so does
$x_{\mu_0}$.

(3) $h_\nu\ne0$, $h_\nu\in W_0$ by Lemma~\ref{is-onto-step-exp-bound-short-graded-rational-module}(5),
$E_{p+a,p+b}h_\nu=0$ for $a<b$ as in its proof, and $E_{p+a,p+a}h_\nu=\nu_{p+a}h_\nu=\lambda_{2,a}h_\nu$.

(4) Let $S$ be such a subspace; $h_\nu\in S$ by (3). By
Lemma~\ref{is-onto-step-exp-bound-short-graded-rational-pieri-irreducible}, $h_\nu$ lies in every nonzero stable
subspace of $W_0$. So Lemma~\ref{is-onto-step-exp-bound-short-graded-rational-pieri-highest}(6), with (1) and (2), gives
$S=W_0$.
% Lean: (1) real_form_finite (dualWeyl inside the weightedHomogeneousSubmodule for colW, spanned by finitely many
% monomials); (2) real_form_modular (generators_satisfy_gl_relations, fock_symm, fockQPos with hw_norm_pos,
% real_contragredient, diagonal_transfer with euler_diagonal through lpr_inv); (3) real_highest_weight (hwNe,
% der_hw_raise, der_diag_iff); (4) real_form_generation (irreducible_generation with real_irreducible).
\end{proof}

\begin{lemma}\label{is-onto-step-exp-bound-short-graded-rational-pieri}
Let $p,q,w$ and $\lambda_2$ be as in Definition~\ref{is-onto-step-exp-bound-short-graded-rational}, and
$n\geq0$. For the rational block module of $(p,q,w,\lambda_2)$, $T_n$ is the sum of the box path spaces
$T_n(\lambda_2,\tau)$, $\tau\in\mathbf N^q$ (Definition~\ref{is-onto-step-exp-bound-short-graded-block}).
\end{lemma}
\leanref{Formalism/PieriBlockIdeals/StatementLR.lean}{Formalism.CohomologicalDescent.LR.PieriReal}
\leanref{Formalism/PieriBlockIdeals/PieriColumnMinors/IrreducibleGeneration.lean}{Formalism.CohomologicalDescent.LPR.pieri_real_form}
\begin{proof}
By Lemma~\ref{is-onto-step-exp-bound-short-graded-rational-pieri-real}(1), (2) and~(4),
Lemma~\ref{is-onto-step-exp-bound-short-graded-rational-pieri-general} applies to $K=\mathbf Q$, $M=W_0$ with the
operators $E_{bb'}$, and the restriction of $B$. Its $T_n$ and box path spaces are those of the rational block
module, which gives the claim.
% Lean: pieri_real_form = pieri_generated Q (Fin p) W0 q EW (formB restricted to W0) lam2 real_form_modular
% real_form_generation n, with real_form_finite. Size before this step was written out: about 1,400 lines of Lean in
% 60 theorems beyond those shared with the skeleton and with real_generated.
\end{proof}

\begin{lemma}\label{is-onto-step-exp-bound-short-graded-rational-columns-paths}
Let $q,n\geq0$ be integers, $\lambda_2,\tau\in\mathbf N^q$, and $\kappa$ a box path of length $n$ from
$\lambda_2$ to $\tau$, with rows $\rho_j$, columns $y_j$ and contents $c_j$
(Definition~\ref{is-onto-step-exp-bound-short-graded-skeleton-jucys}(4)).
\begin{enumerate}
\item Let $\kappa'$ be a box path of length $n$ from $\lambda_2$, with contents $c'_j$, and $m\le n$ with
  $c_j=c'_j$ for $j<m$. Then $\kappa_l=\kappa'_l$ for $l\le m$.
\item Let $j+1<n$ with $y_{j+1}\le y_j$. Then $\rho_j<\rho_{j+1}$ and $y_{j+1}=\kappa_{j,\rho_{j+1}}$.
\item Let $j+1<n$ with $y_{j+1}<y_j$. Then $c_j-c_{j+1}\geq2$, and there is a box path of length $n$ from
  $\lambda_2$ to $\tau$ whose rows are those of $\kappa$ with the entries $j$ and $j+1$ exchanged; its columns
  and contents are those of $\kappa$ with the entries $j$ and $j+1$ exchanged.
\item Let $j+1<n$ with $y_{j+1}=y_j$. Then $c_{j+1}=c_j-1$, and $c$ with the entries $j$ and $j+1$ exchanged
  is not a content sequence of length $n$ for $\lambda_2$.
\item Suppose that $y_j\le y_{j+1}$ whenever $j+1<n$, and let $p,y\geq0$ be integers with
  $\mathrm{col}(\lambda_2,\tau,y)>p$ (Definition~\ref{is-onto-step-exp-bound-short-branching}(3)). Then there is
  an $i$ with $i+p<n$ and $y_l=y$ for $i\le l\le i+p$.
\item For $u\in\mathbf N^n$ put $\mathrm{ms}(u)=\sum_{l<n}(n-l)u_l$. If $j+1<n$ and $u_{j+1}<u_j$, then
  $\mathrm{ms}(u')<\mathrm{ms}(u)$, where $u'$ is $u$ with the entries $j$ and $j+1$ exchanged.
\end{enumerate}
\end{lemma}
\leanref{Formalism/PieriBlockIdeals/StatementLC.lean}{Formalism.CohomologicalDescent.LC.meas}
\leanref{Formalism/PieriBlockIdeals/PieriColumnMinors/PathEqOfContentSequence.lean}{Formalism.CohomologicalDescent.LC.path_eq_of_content_sequence}
\leanref{Formalism/PieriBlockIdeals/PieriColumnMinors/PathEqOfContentSequence.lean}{Formalism.CohomologicalDescent.LC.lca_rows_lt}
\leanref{Formalism/PieriBlockIdeals/StatementLC.lean}{Formalism.CohomologicalDescent.LC.SwapPath}
\leanref{Formalism/PieriBlockIdeals/PieriColumnMinors/PathEqOfContentSequence.lean}{Formalism.CohomologicalDescent.LC.swap_path_exists}
\leanref{Formalism/PieriBlockIdeals/StatementLC.lean}{Formalism.CohomologicalDescent.LC.SameCol}
\leanref{Formalism/PieriBlockIdeals/PieriColumnMinors/PathEqOfContentSequence.lean}{Formalism.CohomologicalDescent.LC.same_column}
\leanref{Formalism/PieriBlockIdeals/StatementLC.lean}{Formalism.CohomologicalDescent.LC.MonoRun}
\leanref{Formalism/PieriBlockIdeals/PieriColumnMinors/AntitoneShapes.lean}{Formalism.CohomologicalDescent.LC.monotone_run}
\leanref{Formalism/PieriBlockIdeals/StatementLC.lean}{Formalism.CohomologicalDescent.LC.MeasSwap}
\leanref{Formalism/PieriBlockIdeals/PieriColumnMinors/MeasureSwapInvariant.lean}{Formalism.CohomologicalDescent.LC.measure_swap_invariant}
\begin{proof}
On a nonincreasing $\kappa_l$, the map $r\mapsto\kappa_{l,r}-r$ is strictly decreasing, hence injective.

(1) Induct on $l$; $\kappa_0=\lambda_2=\kappa'_0$. If $\kappa_l=\kappa'_l$ and $l<m$, then $c_l=c'_l$ says
$\kappa_{l,\rho_l}-\rho_l=\kappa_{l,\rho'_l}-\rho'_l$, so $\rho_l=\rho'_l$ and $\kappa_{l+1}=\kappa'_{l+1}$.

(2) If $\rho_{j+1}\le\rho_j$, then, as $\kappa_{j+1}$ is nonincreasing,
$y_{j+1}=\kappa_{j+1,\rho_{j+1}}\geq\kappa_{j+1,\rho_j}=y_j+1$, contrary to the hypothesis. So
$\rho_j<\rho_{j+1}$, and $y_{j+1}=\kappa_{j+1,\rho_{j+1}}=\kappa_{j,\rho_{j+1}}$.

(3) By (2), $\rho_j<\rho_{j+1}$ and $\kappa_{j,\rho_{j+1}}=y_{j+1}<y_j=\kappa_{j,\rho_j}$, so
$c_j-c_{j+1}=(\kappa_{j,\rho_j}-\kappa_{j,\rho_{j+1}})+(\rho_{j+1}-\rho_j)\geq2$. Let $\kappa'$ agree with $\kappa$
except that $\kappa'_{j+1}=\kappa_j+\epsilon_{\rho_{j+1}}$. Then $\kappa'_{j+1}+\epsilon_{\rho_j}=\kappa_{j+2}$, so
$\kappa'$ raises $\rho_{j+1}$ at step $j$, $\rho_j$ at step $j+1$, and $\rho_l$ at the other steps. The sequence
$\kappa'_{j+1}$ is nonincreasing: only its entry at $\rho_{j+1}$ exceeds that of $\kappa_j$, and for
$a<\rho_{j+1}$, either $a\le\rho_j$ and $\kappa_{j,a}\geq\kappa_{j,\rho_j}>\kappa_{j,\rho_{j+1}}$, or
$\rho_j<a<\rho_{j+1}$ and $\kappa_{j,a}=\kappa_{j+2,a}\geq\kappa_{j+2,\rho_{j+1}}=\kappa_{j,\rho_{j+1}}+1$. The box
added by $\kappa'$ at step $j$ has column $\kappa_{j,\rho_{j+1}}=y_{j+1}$, and the one at step $j+1$ has column
$\kappa'_{j+1,\rho_j}=\kappa_{j,\rho_j}=y_j$. So the columns and the contents at $j$ and $j+1$ are exchanged, and
the others are unchanged.

(4) By (2), $\rho_j<\rho_{j+1}$ and $\kappa_{j,\rho_{j+1}}=\kappa_{j,\rho_j}$. If some $a$ had
$\rho_j<a<\rho_{j+1}$, then $\kappa_{j,a}=\kappa_{j+2,a}\geq\kappa_{j+2,\rho_{j+1}}=\kappa_{j,\rho_j}+1$, while
$\kappa_{j,a}\le\kappa_{j,\rho_j}$. So $\rho_{j+1}=\rho_j+1$, and $c_{j+1}=\kappa_{j,\rho_j}-\rho_j-1=c_j-1$. Suppose
that a box path $\kappa'$ of length $n$ from $\lambda_2$ has the contents $c$ with the entries $j$ and $j+1$
exchanged. Its contents agree with $c$ before $j$, so $\kappa'_j=\kappa_j$ by (1). Its row $\rho'_j$ has
$\kappa_{j,\rho'_j}-\rho'_j=c_{j+1}=\kappa_{j,\rho_{j+1}}-\rho_{j+1}$, so $\rho'_j=\rho_{j+1}$. Then
$\kappa'_{j+1}=\kappa_j+\epsilon_{\rho_{j+1}}$ has the entry $\kappa_{j,\rho_j}$ at $\rho_j$ and $\kappa_{j,\rho_j}+1$
at $\rho_{j+1}>\rho_j$, so it is not nonincreasing, a contradiction.

(5) Let $S$ be the set of the $l<n$ with $y_l=y$. For each $a$ with $\lambda_{2,a}\le y<\tau_a$, the entry $a$
of $\kappa_l$ rises by steps of $1$ from $\lambda_{2,a}$ to $\tau_a$, so some step $l$ raises the entry $a$ when
$\kappa_{l,a}=y$; then $l\in S$ and $\rho_l=a$. So $\rho$ maps $S$ onto a set of at least
$\mathrm{col}(\lambda_2,\tau,y)>p$ elements, and $\#S>p$. Let $i$ and $i'$ be the least and the largest elements
of $S$. Then $i'-i+1\geq\#S\geq p+1$, so $i+p\le i'<n$; and for $i\le l\le i+p$, $y=y_i\le y_l\le y_{i'}=y$, as
the columns are nondecreasing.

(6) Only the terms $l=j$ and $l=j+1$ change, and
$(n-j)u_{j+1}+(n-j-1)u_j-(n-j)u_j-(n-j-1)u_{j+1}=u_{j+1}-u_j<0$.
% Lean: (1) path_eq_of_content_sequence with lca_content_inj; (2) lca_rows_lt; (3) swap_path_exists = lca_swap
% (the antitone check is anti, with d0 and d2); (4) same_column = lca_same (hrow, then path_eq_of_content_sequence,
% lca_content_inj and H'.dom); (5) monotone_run (lcb_ivt, Finset.card_image_le, min' and max' of S, lcb_mono);
% (6) measure_swap_invariant.
\end{proof}

\begin{lemma}\label{is-onto-step-exp-bound-short-graded-rational-columns-weights}
Let $K$ be a field of characteristic $0$, $I_1$ a finite set with $p$ elements, $q,n\geq0$ integers,
$I_2=\{0,\dots,q-1\}$, $M$ a $K$-vector space with endomorphisms $E_{bb'}$ for $b,b'\in I_2$, and
$\lambda_2\in\mathbf N^q$, with $T_n$ and the operators of
Definition~\ref{is-onto-step-exp-bound-short-graded-skeleton-jucys}. Let $c\in\mathbf Z^n$. For $j+1<n$ write
$s_j=\pi^{\mathrm Y}_{(j\,j+1)}$, and $c^{(j)}$ for $c$ with the entries $j$ and $j+1$ exchanged. Weights are
for $\lambda_2$.
\begin{enumerate}
\item If $z\in T_n$ has weight $c$, then $C_lz=\big(\mathrm{cas}(\lambda_2)+\sum_{i<l}(2c_i+q)\big)z$ for $l\le n$,
  and $\pi^{\mathrm X}_\sigma z$ has weight $c$ for every permutation $\sigma$ of $\{0,\dots,n-1\}$.
\item Let $\kappa_0=\lambda_2,\kappa_1,\dots,\kappa_n\in\mathbf N^q$ and $\rho_j\in I_2$ with
  $\kappa_{j+1}=\kappa_j+\epsilon_{\rho_j}$ for $j<n$, and let $c_j=\kappa_{j,\rho_j}-\rho_j$. Then an element
  $z\in T_n$ has $C_jz=\mathrm{cas}(\kappa_j)z$ for every $j\le n$ if and only if it has weight $c$.
\item If $z$ has weight $c$, $j+1<n$ and $c_j\ne c_{j+1}$, then $s_jz+(c_j-c_{j+1})^{-1}z$ has weight $c^{(j)}$.
\item Let $j+1<n$ with $c_j-c_{j+1}\geq2$. If $z'^{\mathrm{sym}}=0$ for every $z'$ of weight $c^{(j)}$, then
  $z^{\mathrm{sym}}=0$ for every $z$ of weight $c$.
\item Suppose that $T_n$ is the sum of the box path spaces $T_n(\lambda_2,\tau)$, $\tau\in\mathbf N^q$. If $z$ has
  weight $c$ and $c$ is not a content sequence of length $n$ for $\lambda_2$, then $z=0$. If $z$ has weight $c$,
  $j+1<n$, $c_{j+1}=c_j-1$ and $c^{(j)}$ is not a content sequence of length $n$ for $\lambda_2$, then
  $s_jz=-z$.
\item Let $u\in T_n$ and $i\geq0$ with $i+p<n$ and $s_ju=-u$ for $i\le j<i+p$. Then $u^{\mathrm{sym}}=0$.
\end{enumerate}
\end{lemma}
\leanref{Formalism/PieriBlockIdeals/StatementLC.lean}{Formalism.CohomologicalDescent.LC.CasFromX}
\leanref{Formalism/PieriBlockIdeals/PieriColumnMinors/PermutationWeightInvariant.lean}{Formalism.CohomologicalDescent.LC.casimir_from_x}
\leanref{Formalism/PieriBlockIdeals/StatementLC.lean}{Formalism.CohomologicalDescent.LC.PermXWt}
\leanref{Formalism/PieriBlockIdeals/PieriColumnMinors/PermutationWeightInvariant.lean}{Formalism.CohomologicalDescent.LC.permutation_weight_invariant}
\leanref{Formalism/PieriBlockIdeals/StatementLC.lean}{Formalism.CohomologicalDescent.LC.WtIff}
\leanref{Formalism/PieriBlockIdeals/PieriColumnMinors/PermutationWeightInvariant.lean}{Formalism.CohomologicalDescent.LC.weight_iff_condition}
\leanref{Formalism/PieriBlockIdeals/PieriColumnMinors/PermutationWeightInvariant.lean}{Formalism.CohomologicalDescent.LC.lc_c_casVal_path}
\leanref{Formalism/PieriBlockIdeals/StatementLC.lean}{Formalism.CohomologicalDescent.LC.Exchange}
\leanref{Formalism/PieriBlockIdeals/PieriColumnMinors/PermutationWeightInvariant.lean}{Formalism.CohomologicalDescent.LC.exchange'}
\leanref{Formalism/PieriBlockIdeals/StatementLC.lean}{Formalism.CohomologicalDescent.LC.GoodSwap}
\leanref{Formalism/PieriBlockIdeals/PieriColumnMinors/PermutationWeightInvariant.lean}{Formalism.CohomologicalDescent.LC.good_swap_exists}
\leanref{Formalism/PieriBlockIdeals/StatementLC.lean}{Formalism.CohomologicalDescent.LC.NoTuple}
\leanref{Formalism/PieriBlockIdeals/PieriColumnMinors/MeasureSwapInvariant.lean}{Formalism.CohomologicalDescent.LC.no_bad_tuple}
\leanref{Formalism/PieriBlockIdeals/StatementLC.lean}{Formalism.CohomologicalDescent.LC.NoWt}
\leanref{Formalism/PieriBlockIdeals/PieriColumnMinors/PermutationWeightInvariant.lean}{Formalism.CohomologicalDescent.LC.NoWt_of}
\leanref{Formalism/PieriBlockIdeals/StatementLC.lean}{Formalism.CohomologicalDescent.LC.AntiStep}
\leanref{Formalism/PieriBlockIdeals/PieriColumnMinors/LemmaDSkeletonFour.lean}{Formalism.CohomologicalDescent.LC.antitone_increment}
\leanref{Formalism/PieriBlockIdeals/StatementLC.lean}{Formalism.CohomologicalDescent.LC.Anti}
\leanref{Formalism/PieriBlockIdeals/PieriColumnMinors/AntitoneShapes.lean}{Formalism.CohomologicalDescent.LC.antitone_shapes}
\begin{proof}
(1) Induct on $l$. For $l=0$ this is the condition on $C_0$. For $l<n$,
Lemma~\ref{is-onto-step-exp-bound-short-graded-skeleton-jucys-relations}(3) gives
$C_{l+1}z=C_lz+2x_lz+qz$, and $x_lz=c_lz$. By part (5) of that lemma, $\pi^{\mathrm X}_\sigma$ commutes with $C_0$
and with every $x_j$, so it preserves the conditions defining the weight.

(2) If $C_jz=\mathrm{cas}(\kappa_j)z$ for every $j\le n$, then $C_0z=\mathrm{cas}(\lambda_2)z$, and
Lemma~\ref{is-onto-step-exp-bound-short-graded-skeleton-path-eigen} gives $x_jz=c_jz$. Conversely, by
Lemma~\ref{is-onto-step-exp-bound-short-graded-skeleton-jucys-relations}(7) and induction on $j$,
$\mathrm{cas}(\kappa_j)=\mathrm{cas}(\lambda_2)+\sum_{i<j}(2c_i+q)$, so (1) gives $C_jz=\mathrm{cas}(\kappa_j)z$.

(3) Put $s=s_j$, $\beta=(c_j-c_{j+1})^{-1}$ and $z'=sz+\beta z$. By
Lemma~\ref{is-onto-step-exp-bound-short-graded-skeleton-jucys-relations}(4), $sx_j=x_{j+1}s-1$ and $s^2=1$;
multiplying by $s$ on both sides gives $x_js=sx_{j+1}-1$. Hence $x_{j+1}(sz)=c_jsz+z$ and
$x_j(sz)=c_{j+1}sz-z$. As $\beta(c_j-c_{j+1})=1$,
\[
 x_jz'=c_{j+1}sz+(\beta c_j-1)z=c_{j+1}z',\qquad x_{j+1}z'=c_jsz+(1+\beta c_{j+1})z=c_jz'.
\]
For $l\ne j,j+1$, $s$ commutes with $x_l$, so $x_lz'=c_lz'$; and $s$ commutes with $C_0$ by part (5) of that
lemma, so $C_0z'=\mathrm{cas}(\lambda_2)z'$.

(4) Let $z$ have weight $c$, and let $s$, $\beta$ and $z'$ be as in (3). Then $z'$ has weight $c^{(j)}$, so
$z'^{\mathrm{sym}}=0$. By (1), $\pi^{\mathrm X}_{(j\,j+1)}z'$ also has weight $c^{(j)}$, and
$(sz')^{\mathrm{sym}}=(\pi^{\mathrm X}_{(j\,j+1)}z')^{\mathrm{sym}}=0$ by
Lemma~\ref{is-onto-step-exp-bound-short-graded-skeleton-jucys-relations}(6). As $s^2=1$,
$sz'-\beta z'=(1-\beta^2)z$, so $(1-\beta^2)z^{\mathrm{sym}}=0$. As $|c_j-c_{j+1}|\geq2$, $\beta^2\ne1$, and
$z^{\mathrm{sym}}=0$.

(5) First a fact of linear algebra. Let $(T_i)_{i\in J}$ be endomorphisms of a $K$-vector space and $E$ a set of
common eigenvectors. If $v$ lies in the span of $E$ and is a common eigenvector whose eigenvalues $(t_i)$ are not
those of any element of $E$, then $v=0$. Indeed, write $v=\sum_{e\in F}f_ee$ with $F\subset E$ finite, and induct
on $\#F$. For $F$ empty, $v=0$. Otherwise pick $e_0\in F$, with eigenvalues $(t_i(e_0))$, and $i$ with
$t_i\ne t_i(e_0)$. Then
$(t_i-t_i(e_0))v=(T_i-t_i(e_0))v=\sum_{e\in F\setminus\{e_0\}}f_e\big(t_i(e)-t_i(e_0)\big)e$ is a common
eigenvector with the eigenvalues $(t_i)$, so it is $0$ by induction, and $v=0$.

Apply this to the $C_j$, $j\le n$, and the set $E$ of the elements $z''$ for which some box path $\kappa$ of
length $n$ from $\lambda_2$ has $C_jz''=\mathrm{cas}(\kappa_j)z''$ for every $j\le n$; by assumption $E$ spans
$T_n$. By (1), $z$ has the eigenvalues $t_j=\mathrm{cas}(\lambda_2)+\sum_{i<j}(2c_i+q)$. If these were the
$\mathrm{cas}(\kappa_j)$ of a box path $\kappa$ from $\lambda_2$, then by
Lemma~\ref{is-onto-step-exp-bound-short-graded-skeleton-jucys-relations}(7),
$2c_i+q=\mathrm{cas}(\kappa_{i+1})-\mathrm{cas}(\kappa_i)=2c'_i+q$ with $c'$ the contents of $\kappa$, and $c$ would
be a content sequence. So $z=0$. For the second claim, $c_j-c_{j+1}=1$, so by (3) $s_jz+z$ has weight $c^{(j)}$;
by the first claim it is $0$.

(6) First, $\pi^{\mathrm Y}_{(a\,b)}u=-u$ for $i\le a<b\le i+p$, by induction on $b-a$: for $b-a\geq2$, put
$d=b-1$; then $(a\,b)=(d\,b)(a\,d)(d\,b)$, and
Lemma~\ref{is-onto-step-exp-bound-short-graded-skeleton-jucys-relations}(1) gives
$\pi^{\mathrm Y}_{(a\,b)}u=\pi^{\mathrm Y}_{(d\,b)}\pi^{\mathrm Y}_{(a\,d)}\pi^{\mathrm Y}_{(d\,b)}u=-u$. Let
$g\in(I_2\times I_1)^n$, with $g_s=(b_s,\alpha_s)$. The $p+1$ positions $i,\dots,i+p$ carry entries $\alpha_s$ in
$I_1$, which has $p$ elements; so two of them, $a<b$, have $\alpha_a=\alpha_b$. Choose such a pair $(a(g),b(g))$
depending only on the entries $\alpha_s$ of $g$, and put $\iota(g)=g\circ(a(g)\,b(g))$. Then $\iota(g)$ has the
same entries $\alpha_s$ as $g$, so $\iota(\iota(g))=g$, and its entries form the same multiset. As
$\alpha_{a(g)}=\alpha_{b(g)}$, $u(\iota(g))=(\pi^{\mathrm Y}_{(a(g)\,b(g))}u)(g)=-u(g)$. In the sum
$u^{\mathrm{sym}}(m)=\sum_gu(g)$, the terms for $g$ and $\iota(g)\ne g$ cancel, and $u(g)=0$ if $\iota(g)=g$, as
$2u(g)=0$. So $u^{\mathrm{sym}}=0$.
% Lean: (1) casimir_from_x (casimir_jucys_murphy_holds) and permutation_weight_invariant (lc_c_eig with
% commute_casimir_zero and commute_jucysMurphy_of_commute); (2) weight_iff_condition (path_eigen_holds;
% casimir_from_x with lc_c_casVal_path, which iterates lc_c_casVal_step); (3) exchange' (hecke_relations_hold,
% lc_c_ss); (4) good_swap_exists (exchange', to_symmetric_y_x, permutation_weight_invariant); (5) NoWt_of
% no_bad_tuple, and antitone_increment = AntiStep_of no_bad_tuple; (6) antitone_shapes (lcb_transp,
% Fintype.exists_ne_map_eq_of_card_lt, toSym_apply_sum_fiber, Finset.sum_involution).
\end{proof}

\begin{lemma}\label{is-onto-step-exp-bound-short-graded-rational-columns-general}
Let $K$ be a field of characteristic $0$, $I_1$ a finite set with $p$ elements, $q,n\geq0$ integers,
$I_2=\{0,\dots,q-1\}$, $M$ a $K$-vector space with endomorphisms $E_{bb'}$ for $b,b'\in I_2$, and
$\lambda_2\in\mathbf N^q$, with $T_n$ and the box path spaces $T_n(\lambda_2,\tau)$ of
Definition~\ref{is-onto-step-exp-bound-short-graded-skeleton-jucys}. Suppose that $T_n$ is the sum of the
$T_n(\lambda_2,\tau)$, $\tau\in\mathbf N^q$. Let $\tau\in\mathbf N^q$ with $\mathrm{col}(\lambda_2,\tau,y)>p$ for some
$b$ and some integer $0\le y<\tau_b$ (Definition~\ref{is-onto-step-exp-bound-short-branching}(3)). Then
$z^{\mathrm{sym}}=0$ for every $z\in T_n(\lambda_2,\tau)$.
\end{lemma}
\leanref{Formalism/PieriBlockIdeals/StatementLC.lean}{Formalism.CohomologicalDescent.LC.ColumnsFromPieri}
\leanref{Formalism/PieriBlockIdeals/PieriColumnMinors/MeasureSwapInvariant.lean}{Formalism.CohomologicalDescent.LC.ColumnsFromPieri_of}
\leanref{Formalism/PieriBlockIdeals/PieriColumnMinors/LemmaDSkeletonFour.lean}{Formalism.CohomologicalDescent.LC.columns_from_pieri}
\begin{proof}
For a box path $\kappa$ of length $n$ from $\lambda_2$ to $\tau$, with columns $y(\kappa)$ and contents
$c(\kappa)$, we show that $z^{\mathrm{sym}}=0$ for every $z$ of weight $c(\kappa)$ for $\lambda_2$, by induction on
$\mathrm{ms}(y(\kappa))$ (Lemma~\ref{is-onto-step-exp-bound-short-graded-rational-columns-paths}(6)). Write
$c=c(\kappa)$ and $y_j$ for the columns.

Suppose first that $y_{j+1}<y_j$ for some $j+1<n$. By
Lemma~\ref{is-onto-step-exp-bound-short-graded-rational-columns-paths}(3), $c_j-c_{j+1}\geq2$, and there is a box
path $\kappa'$ from $\lambda_2$ to $\tau$ with contents $c^{(j)}$ and columns $y$ with the entries $j$ and $j+1$
exchanged; by part (6) of that lemma, $\mathrm{ms}(y(\kappa'))<\mathrm{ms}(y(\kappa))$. By induction,
$z'^{\mathrm{sym}}=0$ for every $z'$ of weight $c^{(j)}$, and
Lemma~\ref{is-onto-step-exp-bound-short-graded-rational-columns-weights}(4) gives the claim for $\kappa$.

Otherwise $y_j\le y_{j+1}$ whenever $j+1<n$, and
Lemma~\ref{is-onto-step-exp-bound-short-graded-rational-columns-paths}(5) gives an $i$ with $i+p<n$ and $y_l=y$ for
$i\le l\le i+p$. For $i\le j<i+p$, $y_{j+1}=y_j$, so by part (4) of that lemma $c_{j+1}=c_j-1$ and $c^{(j)}$ is not
a content sequence of length $n$ for $\lambda_2$. Let $z$ have weight $c$. By
Lemma~\ref{is-onto-step-exp-bound-short-graded-rational-columns-weights}(5), $s_jz=-z$ for $i\le j<i+p$, and part
(6) of that lemma gives $z^{\mathrm{sym}}=0$.

Finally, $T_n(\lambda_2,\tau)$ is spanned by the $z$ for which some box path $\kappa$ of length $n$ from
$\lambda_2$ to $\tau$ has $C_jz=\mathrm{cas}(\kappa_j)z$ for every $j\le n$. Such a $z$ has weight $c(\kappa)$ by
Lemma~\ref{is-onto-step-exp-bound-short-graded-rational-columns-weights}(2), so $z^{\mathrm{sym}}=0$; and the
symmetrization is linear.
% Lean: columns_from_pieri = ColumnsFromPieri_of swap_path_exists same_column measure_swap_invariant
% monotone_run weight_iff_condition good_swap_exists antitone_increment antitone_shapes: key, by
% Nat.strong_induction_on the measure meas (colSeq kappa rho), splitting on a descent of colSeq (hSwap, hGood,
% hMeas) or none (hMono, hSame, hAntiStep, hAnti); then PathSpace <= ker toSym by Submodule.span_le with hWt.
\end{proof}

\begin{lemma}\label{is-onto-step-exp-bound-short-graded-rational-columns}
Let $p,q,w$ and $\lambda_2$ be as in Definition~\ref{is-onto-step-exp-bound-short-graded-rational}, and
$n\geq0$, and suppose that $T_n$ is the sum of the box path spaces $T_n(\lambda_2,\tau)$ for the rational
block module of $(p,q,w,\lambda_2)$. Let $\tau\in\mathbf N^q$ with $\mathrm{col}(\lambda_2,\tau,y)>p$ for some $b$
and some $y<\tau_b$ (Definition~\ref{is-onto-step-exp-bound-short-branching}(3)). Then $z^{\mathrm{sym}}=0$ for
every $z\in T_n(\lambda_2,\tau)$.
\end{lemma}
\leanref{Formalism/PieriBlockIdeals/StatementLR.lean}{Formalism.CohomologicalDescent.LR.ColumnsReal}
\leanref{Formalism/PieriBlockIdeals/LemmaDGradedOfPieriReal.lean}{Formalism.CohomologicalDescent.ColumnMinorsFromPieri.ColumnsRealOfPieriReal}
\leanref{Formalism/PieriBlockIdeals/PieriColumnMinors/ColumnsPieriSkeleton.lean}{Formalism.CohomologicalDescent.ColumnMinorsFromPieri.columns_real_of_pieri_real}
\begin{proof}
The space $W_0$ of the rational block module is a $\mathbf Q$-vector space with the endomorphisms $E_{bb'}$,
$b,b'\in I_2$, and $I_1=\{0,\dots,p-1\}$ has $p$ elements. Its $T_n$, its symmetrization and its box path spaces are
those of Definition~\ref{is-onto-step-exp-bound-short-graded-skeleton-jucys} for $K=\mathbf Q$, $M=W_0$ and these
$E_{bb'}$. So the claim is
Lemma~\ref{is-onto-step-exp-bound-short-graded-rational-columns-general} for them.
% Lean: columns_real_of_pieri_real applies LC.columns_from_pieri Q (Fin p) W0 q (EW) n lam2 to the Pieri
% hypothesis at the same p, q, w, lam2 and n (Lemma
% is-onto-step-exp-bound-short-graded-rational-columns-general). Size before this step was written out: about
% 800 lines of Lean in 40 theorems beyond those shared with the skeleton and the Pieri step.
\end{proof}

\begin{lemma}\label{is-onto-step-exp-bound-short-graded-rational-exponent-restriction}
Let $p,q,w$, $\lambda_2$ and $\nu$ be as in Definition~\ref{is-onto-step-exp-bound-short-graded-rational}, with the
rational block module $(V,E,W_0,w)$ and the lattice $\Lambda$. For a commutative ring $R$, let
$\mathrm{res}\colon R[x_{cd}:c,d<p+q]\to R[x_{bt}:b,t<q]$ be the homomorphism of $R$-algebras that substitutes
$\delta_{cd}$ for $x_{cd}$ if $c<p$ or $d<p$, and $x_{bt}$ for $x_{p+b,p+t}$. Indices of rows and columns start at
$0$.
\begin{enumerate}
\item $\mathrm{res}\circ E^{(k)}_{p+b,p+b'}=E^{(k)}_{bb'}\circ\mathrm{res}$ for $b,b'<q$ and $k\geq0$, and
  $\mathrm{res}(f)(1)=f(1)$, the values at the identity matrices.
\item If $f$ has row degree $\mu\in\mathbf N^{p+q}$, then $\mathrm{res}(f)$ has row degree $(\mu_{p+b})_{b<q}$.
\item $\mathrm{res}$ maps $\nabla_R(\nu)$ into $\nabla_R(\lambda_2)$, and $\mathrm{res}(h_\nu)=h_{\lambda_2}$.
\item $\Lambda\subset W_0$, and every element of $\Delta_{\mathbf Z}(\lambda_2)$ is $\mathrm{res}(d)$ for some
  $d\in\Lambda$.
\item If $u\in W_0$ and $\mathrm{res}(u)=0$, then $u=0$.
\end{enumerate}
\end{lemma}
\leanref{Formalism/PieriBlockIdeals/PieriColumnMinors/LrbPiSingleCast.lean}{Formalism.CohomologicalDescent.LR.res}
\leanref{Formalism/PieriBlockIdeals/PieriColumnMinors/LrbPiSingleCast.lean}{Formalism.CohomologicalDescent.LR.liftWord}
\leanref{Formalism/PieriBlockIdeals/PieriColumnMinors/LrbPiSingleCast.lean}{Formalism.CohomologicalDescent.LR.res_rowDiv}
\leanref{Formalism/PieriBlockIdeals/PieriColumnMinors/LrbPiSingleCast.lean}{Formalism.CohomologicalDescent.LR.res_applyWord}
\leanref{Formalism/PieriBlockIdeals/PieriColumnMinors/LrbPiSingleCast.lean}{Formalism.CohomologicalDescent.LR.res_eval_id}
\leanref{Formalism/PieriBlockIdeals/PieriColumnMinors/LrbPiSingleCast.lean}{Formalism.CohomologicalDescent.LR.res_rowHom}
\leanref{Formalism/PieriBlockIdeals/PieriColumnMinors/LrbPiSingleCast.lean}{Formalism.CohomologicalDescent.LR.res_mem}
\leanref{Formalism/PieriBlockIdeals/PieriColumnMinors/LrbPiSingleCast.lean}{Formalism.CohomologicalDescent.LR.res_hw}
\leanref{Formalism/PieriBlockIdeals/PieriColumnMinors/LrjTransposeLift.lean}{Formalism.CohomologicalDescent.LR.levZ_sub}
\leanref{Formalism/PieriBlockIdeals/PieriColumnMinors/LrjTransposeLift.lean}{Formalism.CohomologicalDescent.LR.liftWeyl}
\leanref{Formalism/PieriBlockIdeals/PieriColumnMinors/LrjTransposeLift.lean}{Formalism.CohomologicalDescent.LR.resInj}
\begin{proof}
(1) The substitution replaces $X$ by the block diagonal matrix with blocks $1_p$ and $X'$, where $X'$ is the
$q\times q$ matrix of the $x_{bt}$. Adding $y$ times row $p+b$ to row $p+b'$ of this matrix gives the block diagonal
matrix with $X'$ modified by adding $y$ times row $b$ to row $b'$. Comparing the coefficients of $y^k$ gives the
first claim. At $X'=1$ the block matrix is $1_{p+q}$.

(2) The variables of row $p+b$ go to variables of row $b$ or to constants, and those of the first $p$ rows to
constants.

(3) Likewise the variables of column $p+t$ go to variables of column $t$ or to $0$, so $\mathrm{res}$ carries a
polynomial of degree $\nu_{p+t}=\lambda_{2,t}$ in column $p+t$ to one of degree $\lambda_{2,t}$ in column $t$; and
the column operations of Definition~\ref{pd-weight-lattices}(1) on the last $q$ columns commute with the
substitution in the same way as the row operations in (1). So $\mathrm{res}(\nabla_R(\nu))\subset\nabla_R(\lambda_2)$.
For $b<b'$ and $k\geq1$, $E^{(k)}_{p+b,p+b'}h_\nu=0$, so $\mathrm{res}(h_\nu)\in\nabla_{\mathbf Q}(\lambda_2)$ is killed
by $E^{(k)}_{bb'}$ by (1). In characteristic $0$ the vectors of $\nabla_{\mathbf Q}(\lambda_2)$ killed by these
operators form the line of $h_{\lambda_2}$ \cite[II.2]{Jan03}, so $\mathrm{res}(h_\nu)=ch_{\lambda_2}$ for some
$c\in\mathbf Q$. Evaluating at $1$ by (1), $c=h_\nu(1)=1$.

(4) By Lemma~\ref{is-onto-step-exp-bound-short-graded-rational-module}(1), $W_0$ is stable under the $E_{bb'}$,
$b,b'\in I_2$, hence, as in the proof of part (5) of that lemma, under the $E^{(k)}_{p+b,p+b'}$. It contains $h_\nu$
by part (5) of that lemma, hence the generators of $\Lambda$, and $\Lambda$. The module $\Delta_{\mathbf Z}(\lambda_2)$
is the $\mathbf Z$-span of the $\omega h_{\lambda_2}$ for composites $\omega$ of operators $E^{(k)}_{bb'}$. Let
$\omega'$ be the same composite of the $E^{(k)}_{p+b,p+b'}$; then $\omega'h_\nu\in\Lambda$, and by (1) and (3),
$\mathrm{res}(\omega'h_\nu)=\omega\,\mathrm{res}(h_\nu)=\omega h_{\lambda_2}$.

(5) The subspace $Z$ of the $u\in W_0$ with $\mathrm{res}(u)=0$ is stable under the $E_{bb'}$, by
Lemma~\ref{is-onto-step-exp-bound-short-graded-rational-module}(1) and (1) with $k=1$. If $Z\ne0$, then
$h_\nu\in Z$ by Lemma~\ref{is-onto-step-exp-bound-short-graded-rational-pieri-irreducible}, so
$h_{\lambda_2}=\mathrm{res}(h_\nu)=0$ by (3), which is false.
% Lean: (1) lrb_rowGen with res_rowDiv and res_applyWord (liftWord), and res_eval_id; (2) res_rowHom (lrbPi,
% lrb_var_row); (3) res_mem (res_colHom, lrb_resO_colAdd with rowst_mem_dualWeyl_iff) and res_hw (hwLine,
% hwRaise, lehw_eval_idFrame_hw); (4) levZ_sub (closed_of_der, w_zero_stable_right_block, hwV_mem) and liftWeyl
% (span_induction over weyl Z, lclosedeq_applyWord_map, lclosedeq_hw_map); (5) resInj (real_irreducible).
\end{proof}

\begin{lemma}\label{is-onto-step-exp-bound-short-graded-rational-exponent-dual}
Let $l$ be a prime, $r\geq1$, $\nu\in\mathbf N^r$ nonincreasing, and $g\in\nabla_{\mathbf Q}(\nu)$
(Definition~\ref{pd-weight-lattices}(1)). Suppose that $(\omega g)(1)\in\mathbf Z_{(l)}$ for every composite $\omega$
of operators $E^{(k)}_{ab}$, where $\mathbf Z_{(l)}$ is the ring of the rational numbers whose denominator is prime to
$l$. Then $mg\in\nabla_{\mathbf Z}(\nu)$ for some integer $m$ prime to $l$.
\end{lemma}
\leanref{Formalism/PieriBlockIdeals/PieriColumnMinors/LrcFockZeroRight.lean}{Formalism.CohomologicalDescent.LR.dual_int_loc}
\begin{proof}
Let $Z^e$ be a monomial with $e_{as}=0$ for $a<s$. If the degree of $Z^e$ in some column $s$ differs from $\nu_s$,
its coefficient in $g$ is $0$. Otherwise it is $(w_eg)(1)\in\mathbf Z_{(l)}$, with $w_e$ as in
Lemma~\ref{pd-kappa-gram-fock}(5). Let $m$ be the product of the denominators of these finitely many coefficients; it
is prime to $l$, and the coefficients of $mg\in\nabla_{\mathbf Q}(\nu)$ at all these monomials are integers. So
$mg\in\nabla_{\mathbf Z}(\nu)$ by Lemma~\ref{pd-kappa-gram-fock}(6) for $A=\mathbf Z$.
% Lean: dual_int_loc (hC : CoeffWord, Lemma pd-kappa-gram-fock(5); hL : LowerTriIntegral, which is part (6) and is
% proved as lower_triangular_integral), with M the product of the denominators over g.support filtered by
% LowerTri, and dualWeylIntRat.
\end{proof}

\begin{lemma}\label{is-onto-step-exp-bound-short-graded-rational-exponent}
Let $p,q,w$, $\lambda_2$ and $\nu$ be as in Definition~\ref{is-onto-step-exp-bound-short-graded-rational},
with the rational block module, the form $B$, a prime $l$ and the module $\Delta_L$, and let $\nabla_L$ be
as in Definition~\ref{is-onto-step-exp-bound-short-graded-block}(2) for $O=\mathbf Z_{(l)}$. Let
$\xi\in\mathbf Z^{p+q}$, $i'\in\mathbf Z$ and $e_L\geq0$ be such that $(\lambda_2,\mu_2)$ has exponent at most
$e_L$ at $l$ for every $\mu_2\in\mathbf N^q$ with $\mathrm{dp}_\xi(w^p\lambda_2,w^p\mu_2)\le i'$
(Definition~\ref{is-onto-step-exp-bound-short-branching}). Then $l^{e_L}u\in\Delta_L$ for every
$\kappa\in\mathbf Z^{p+q}$ with $\sum_c\xi_c(\nu_c-\kappa_c)\le i'$ and every $u\in\nabla_L$ that lies in
$V_\kappa$.
\end{lemma}
\leanref{Formalism/PieriBlockIdeals/PieriColumnMinors/LrjTransposeLift.lean}{Formalism.CohomologicalDescent.LR.expL_real}
\leanref{Formalism/PieriBlockIdeals/PieriColumnMinors/LrjTransposeLift.lean}{Formalism.CohomologicalDescent.LR.wt_extract}
\leanref{Formalism/PieriBlockIdeals/PieriColumnMinors/LrcFockZeroRight.lean}{Formalism.CohomologicalDescent.LR.fock_word}
\leanref{Formalism/PieriBlockIdeals/PieriColumnMinors/XiOfDepth2.lean}{Formalism.CohomologicalDescent.LR.xiOf_depth2}
\begin{proof}
Let $\kappa$ and $u$ be as in the statement, and use the restriction map $\mathrm{res}$ of
Lemma~\ref{is-onto-step-exp-bound-short-graded-rational-exponent-restriction}. If $u=0$ there is nothing to prove;
let $u\ne0$. As $E_{cc}u=\kappa_cu$ and $E_{cc}$ multiplies each monomial by its degree in row $c$, every monomial of $u$
has degree $\kappa_c$ in row $c$; so $\kappa=\mu\in\mathbf N^{p+q}$ and $u$ has row degree $\mu$. As $u\in W_0$,
$E_{\alpha\alpha}u=wu$, so $\mu_\alpha=w$ for $\alpha<p$. Put $\mu_2=(\mu_{p+b})_{b<q}$. Then $\mu=w^p\mu_2$, and
$\mathrm{dp}_\xi(w^p\lambda_2,w^p\mu_2)=\sum_c\xi_c(\nu_c-\kappa_c)\le i'$.

Put $g=\mathrm{res}(u)$, which lies in $\nabla_{\mathbf Q}(\lambda_2)$ by
Lemma~\ref{is-onto-step-exp-bound-short-graded-rational-exponent-restriction}(3). Let $\omega$ be a composite of
operators $E^{(k)}_{bb'}$ with $b,b'<q$, let $\omega^{\mathsf t}$ be the composite of the $E^{(k)}_{b'b}$ in the
reverse order, and let $\omega'$ and $\omega'^{\mathsf t}$ be the same composites with $p+b$ and $p+b'$ in place of $b$
and $b'$. Then $y=\omega'^{\mathsf t}h_\nu\in\Lambda$, so $y\in\Delta_L$ by
Lemma~\ref{is-onto-step-exp-bound-short-graded-rational-exponent-restriction}(4), and $B(u,y)\in\mathbf Z_{(l)}$ as
$u\in\nabla_L$. By Lemma~\ref{pd-kappa-gram-fock}(2) and~(4), as $\omega'u\in\nabla_{\mathbf Q}(\nu)$,
$[u,\omega'^{\mathsf t}h_\nu]=[\omega'u,h_\nu]=(\omega'u)(1)[h_\nu,h_\nu]$, so $B(u,y)=(\omega'u)(1)$. By
Lemma~\ref{is-onto-step-exp-bound-short-graded-rational-exponent-restriction}(1),
$(\omega'u)(1)=\mathrm{res}(\omega'u)(1)=(\omega g)(1)$. So $(\omega g)(1)\in\mathbf Z_{(l)}$ for every $\omega$, and
Lemma~\ref{is-onto-step-exp-bound-short-graded-rational-exponent-dual} gives an integer $m$ prime to $l$ with
$g'=mg\in\nabla_{\mathbf Z}(\lambda_2)$. By
Lemma~\ref{is-onto-step-exp-bound-short-graded-rational-exponent-restriction}(2), $g'$ has row degree $\mu_2$.

By the hypothesis on $(\lambda_2,\mu_2)$, there is an integer $m'$ prime to $l$ with
$m'l^{e_L}g'\in\Delta_{\mathbf Z}(\lambda_2)$. By
Lemma~\ref{is-onto-step-exp-bound-short-graded-rational-exponent-restriction}(4) there is $d\in\Lambda\subset W_0$
with $\mathrm{res}(d)=m'l^{e_L}g'$. Put $c=m'l^{e_L}m$. Then $cu-d\in W_0$ and
$\mathrm{res}(cu-d)=m'l^{e_L}(mg-g')=0$, so $cu=d$ by part (5) of that lemma. Hence $(m'm)(l^{e_L}u)\in\Lambda\cap W_0$
with $m'm$ prime to $l$, that is, $l^{e_L}u\in\Delta_L$.
% Lean: expL_real: wt_extract (real_weight), hfirst (mem_orthogonal_frame_iff, mem_wt), hblock and hd
% (xiOf_depth2), res_mem, hω (levZ_sub with transposeWord, le_locSpan, formB_apply, fock_word,
% lrj_transpose_lift, res_eval_id, res_applyWord), dual_int_loc, lrj_hom_of_smul with res_rowHom, heL, liftWeyl,
% levZ_sub, resInj, mem_locSpan. Size before this step was written out: about 1,150 lines of Lean in 80 theorems
% beyond the skeleton and real_generated.
\end{proof}

\begin{lemma}\label{is-onto-step-exp-bound-short-graded}
The block exponent bound holds in graded form
(Definition~\ref{is-onto-step-exp-bound-short-branching}(4)).
\end{lemma}
\leanref{Formalism/PieriBlockIdeals/PieriColumnMinors/CorollaryOfParamsLinkage.lean}{Formalism.CohomologicalDescent.Down.hG}
\leanref{Formalism/PieriBlockIdeals/StatementLR.lean}{Formalism.CohomologicalDescent.LR.RealOfPieri}
\leanref{Formalism/PieriBlockIdeals/PieriColumnMinors/RealOfSkeleton.lean}{Formalism.CohomologicalDescent.LR.real_of_pieri}
\leanref{Formalism/PieriBlockIdeals/StatementLR.lean}{Formalism.CohomologicalDescent.LR.RealOfSkeleton}
\leanref{Formalism/PieriBlockIdeals/PieriColumnMinors/RealOfSkeleton.lean}{Formalism.CohomologicalDescent.LR.real_of_skeleton}
\begin{proof}
Let $l$, $p,q,w$, $\lambda_2$, $\xi$, $i'$ and $e_L$ be as in
Definition~\ref{is-onto-step-exp-bound-short-branching}(4), with the hypothesis on the pairs
$(\lambda_2,\mu_2)$, and put $\nu=w^p\lambda_2$. Let $i\geq0$ and $\mu\in\mathbf N^{p+q}$ with
$\sum_{a<p}\mu_a+i=pw$ and $\mathrm{dp}_\xi(\nu,\mu)\le i'$, let $x\in\nabla_{\mathbf Z}(\nu)$ have row degree
$\mu$, and put $c=e_L+v_l\big(D_n(p,w,i,\lambda_2)\big)$. It suffices to find an integer $m$ prime to $l$ with
$ml^cx\in\Delta_{\mathbf Z}(\nu)$.

Take the rational block module $(V,E,W_0,w)$ of $(p,q,w,\lambda_2)$, with the form $B$ and the module
$\Delta_L$ for $l$ (Definition~\ref{is-onto-step-exp-bound-short-graded-rational}), and $O=\mathbf Z_{(l)}$,
in which every integer prime to $l$ is a unit. It is a block module over $\mathbf Q$ whose element of
$\mathbf Q$ is the integer $w$, by Lemma~\ref{is-onto-step-exp-bound-short-graded-rational-module}(1). With
$I_1\sqcup I_2$ identified with $\{0,\dots,p+q-1\}$, let $\xi'(\kappa)=\sum_c\xi_c\kappa_c$ for
$\kappa\in\mathbf Z^{p+q}$. Apply Lemma~\ref{is-onto-step-exp-bound-short-graded-skeleton} with $n=i$, this
$w$, $\lambda_2$, $\Delta_L$, $\xi'$, $\lambda=\nu$, $e=e_L$ and $i'$. Its hypotheses hold: (1) by assumption;
(2) by Lemma~\ref{is-onto-step-exp-bound-short-graded-rational-module}(2) and~(5); (3) by
Lemma~\ref{is-onto-step-exp-bound-short-graded-rational-module}(3); (4) by
Lemma~\ref{is-onto-step-exp-bound-short-graded-rational-pieri}; (5) by
Lemma~\ref{is-onto-step-exp-bound-short-graded-rational-columns}, with (4); (6) as
$\xi'(\epsilon_\alpha-\epsilon_b)=\xi_\alpha-\xi_{p+b}\geq0$, $\xi$ being nonincreasing and $\alpha<p+b$; and (7) by
Lemma~\ref{is-onto-step-exp-bound-short-graded-rational-exponent}.

The vector $x$ lies in $V$, in $V_\mu$ by Lemma~\ref{is-onto-step-exp-bound-short-graded-rational-module}(4),
and in $\Delta(\Delta_L)^B$ by Lemma~\ref{is-onto-step-exp-bound-short-graded-rational-lattice}(2). Moreover
$\sum_{\alpha<p}\mu_\alpha+i=pw$ and $\xi'(\nu-\mu)=\mathrm{dp}_\xi(\nu,\mu)\le i'$. So
Lemma~\ref{is-onto-step-exp-bound-short-graded-skeleton} gives $l^cx\in\Delta(\Delta_L)$, and
Lemma~\ref{is-onto-step-exp-bound-short-graded-rational-lattice}(3) gives the integer $m$.
% Lean: Down.hG = real_of_pieri pieri_real_form (columns_real_of_pieri_real pieri_real_form), and
% real_of_pieri = real_of_skeleton coefficient_word lower_triangular_integral (Lemma pd-kappa-gram-fock(5)
% and (6)). real_of_skeleton: hx' (dualWeylMap to Q), the inputs (scalar, le_w, real_contragredient,
% real_generated, real_nondegenerate, delL_stable, pieri, columns, dominant_xi, expL_real), hdep, hxi
% (xiOf_depth2), depth_skeleton_third_form_holds with zl_unit, mem_dual and wt_of_hom, then concl. Size
% before this step was written out: about 7,600 lines of Lean in 450 theorems below it (text count). Its four
% inputs, Lemmas is-onto-step-exp-bound-short-graded-skeleton (about 3,400 lines), -rational-pieri (about 1,400
% beyond the shared parts), -rational-columns (about 800) and -rational-exponent (about 1,150), are now written
% out to the depth floor.
\end{proof}

\subsection{The top weight theorem and the exponent bound}\label{subsection-is-top-weight}

\begin{lemma}\label{is-onto-step-exp-bound-short-last-dual}
Let $r,M\geq0$ be integers. For $\kappa\in\mathbf N^r$ with $\kappa_a\le M$ for every $a$, let
$\kappa^\vee\in\mathbf N^r$ be the sequence $\kappa^\vee_a=M-\kappa_{r-1-a}$. The notation is that of
Definition~\ref{is-onto-step-exp-bound-short-branching}.
\begin{enumerate}
\item $(\kappa^\vee)^\vee=\kappa$ and $\kappa^\vee_a\le M$; if $\kappa$ is nonincreasing, so is $\kappa^\vee$.
\item Let $p,q,w\geq0$ with $w\le M$, and $\kappa\in\mathbf N^q$ with $\kappa_b\le M$. Then
  $(\kappa w^p)^\vee=(M-w)^p\kappa^\vee$. If $\kappa$ is nonincreasing and $w\le\kappa_b$ for every $b$,
  then $\kappa w^p$ is nonincreasing.
\item Let $p,q,w,i\geq0$ with $w\le M$, and $\mu\in\mathbf N^{q+p}$ with $\mu_a\le M$. Then
  $\sum_{a<p}\mu_{q+a}=pw+i$ if and only if $\sum_{a<p}\mu^\vee_a+i=p(M-w)$.
\item Let $\xi\in\mathbf Z^r$ and $\lambda,\mu\in\mathbf N^r$ with entries at most $M$, and put
  $\xi'_a=-\xi_{r-1-a}$. Then $\mathrm{dp}_{\xi'}(\lambda^\vee,\mu^\vee)=\mathrm{dp}_\xi(\lambda,\mu)$, and if
  $\xi$ is nonincreasing, so is $\xi'$.
\end{enumerate}
\end{lemma}
\leanref{Formalism/PieriBlockIdeals/LastGradedOfGradedAlone.lean}{Formalism.CohomologicalDescent.LD.dualWt}
\leanref{Formalism/PieriBlockIdeals/LastGradedOfGradedAlone.lean}{Formalism.CohomologicalDescent.LD.DualWtSpec}
\leanref{Formalism/PieriBlockIdeals/PieriColumnMinors/HighestWeightProdLeadingMinors.lean}{Formalism.CohomologicalDescent.LD.dualWtSpec}
\leanref{Formalism/PieriBlockIdeals/LastGradedOfGradedAlone.lean}{Formalism.CohomologicalDescent.LD.BlockDual}
\leanref{Formalism/PieriBlockIdeals/PieriColumnMinors/HighestWeightProdLeadingMinors.lean}{Formalism.CohomologicalDescent.LD.blockDual}
\leanref{Formalism/PieriBlockIdeals/PieriColumnMinors/LastGradedImpliesOneBlock.lean}{Formalism.CohomologicalDescent.LD.block_weight_last_antitone}
\leanref{Formalism/PieriBlockIdeals/LastGradedOfGradedAlone.lean}{Formalism.CohomologicalDescent.LD.DepthDual}
\leanref{Formalism/PieriBlockIdeals/PieriColumnMinors/HighestWeightProdLeadingMinors.lean}{Formalism.CohomologicalDescent.LD.depthDual}
\leanref{Formalism/PieriBlockIdeals/BranchingBoundLastGradedOneRow.lean}{Formalism.CohomologicalDescent.LD.Depth2Dual}
\leanref{Formalism/PieriBlockIdeals/PieriColumnMinors/LastGradedImpliesOneBlock.lean}{Formalism.CohomologicalDescent.LD.depth2Dual}
\begin{proof}
(1) $M-(M-\kappa_a)=\kappa_a$; and $a\le a'$ gives $r-1-a'\le r-1-a$.

(2) For $a<p$, $(\kappa w^p)^\vee_a=M-(\kappa w^p)_{q+p-1-a}=M-w$; for $b<q$,
$(\kappa w^p)^\vee_{p+b}=M-\kappa_{q-1-b}=\kappa^\vee_b$. The second claim is clear.

(3) As $a\mapsto p-1-a$ permutes $\{0,\dots,p-1\}$, $\sum_{a<p}\mu^\vee_a=\sum_{a<p}(M-\mu_{q+a})
=pM-\sum_{a<p}\mu_{q+a}$, and $pw\le pM$.

(4) After the change of index $a\mapsto r-1-a$,
$\mathrm{dp}_{\xi'}(\lambda^\vee,\mu^\vee)=\sum_a-\xi_a\big((M-\lambda_a)-(M-\mu_a)\big)
=\sum_a\xi_a(\lambda_a-\mu_a)$.
% Lean: dualWtSpec, blockDual (with Fin.cast between q + p and p + q), block_weight_last_antitone, depthDual
% (sum_complement_add_sum) and depth2Dual (DualWtSpec, BlockDual, DepthDual, Depth2Dual).
\end{proof}

\begin{lemma}\label{is-onto-step-exp-bound-short-last-weights}
Let $r,M\geq0$ be integers, $\nu\in\mathbf N^r$ nonincreasing with $\nu_a\le M$ for every $a$, and
$\mu\in\mathbf N^r$. If $\nabla_{\mathbf Z}(\nu)$ contains a nonzero polynomial of row degree $\mu$
(Definition~\ref{is-onto-step-exp-bound-short-branching}), then $\mu_a\le M$ for every $a$ and
$\sum_a\mu_a=\sum_a\nu_a$.
\end{lemma}
\leanref{Formalism/PieriBlockIdeals/LastGradedOfGradedAlone.lean}{Formalism.CohomologicalDescent.LD.WtBound}
\leanref{Formalism/PieriBlockIdeals/PieriColumnMinors/HighestWeightProdLeadingMinors.lean}{Formalism.CohomologicalDescent.LD.wtBound}
\leanref{Formalism/PieriBlockIdeals/PieriColumnMinors/LastGradedImpliesOneBlock.lean}{Formalism.CohomologicalDescent.LD.dominance_of_nabla_weight}
\begin{proof}
By Lemma~\ref{pd-kappa-defects-restriction-standard}(1), $\nabla_{\mathbf Z}(\nu)$ is spanned by the
bideterminants of shape $\nu$: products over $k<K$ of the minors on rows $I_k(0),\dots,I_k(s_k-1)$ and the
first $s_k$ columns, with $\#\{k:t<s_k\}=\nu_t$ for every $t$. Such a minor is $0$ if two of its rows
agree, and otherwise each of its terms takes one entry from each of its rows; so it has degree at most $1$
in every row, and degree $0$ if $s_k=0$. So every monomial of a bideterminant of shape $\nu$, and of every
element of $\nabla_{\mathbf Z}(\nu)$, has degree at most $\#\{k:0<s_k\}=\nu_0\le M$ in every row. A monomial
of the given polynomial has degree $\mu_a$ in row $a$, so $\mu_a\le M$. It has degree $\nu_t$ in column $t$
for every $t$, so its total degree is $\sum_t\nu_t$, and also $\sum_a\mu_a$.
% Lean: wtBound (WtBound): standardMonomialSpan, initial_minor_degree_le_one and dual_weyl_degree_bound;
% the sum is the second part of dominance_of_nabla_weight (dualWeylDominanceZ, from colSums).
\end{proof}

\begin{lemma}\label{is-onto-step-exp-bound-short-last-cofactor}
Let $r\geq0$, let $X=(x_{a,s})$ be the $r\times r$ matrix of variables over $\mathbf Z$, $\delta=\det X$,
$\mathrm{adj}(X)$ its adjugate, and $C$ the ring endomorphism of $\mathbf Z[X]$ with
$C(x_{a,s})=\mathrm{adj}(X)_{r-1-s,r-1-a}$. Write $\kappa^\vee$ as in
Lemma~\ref{is-onto-step-exp-bound-short-last-dual}.
\begin{enumerate}
\item Let $A$ be an $r\times r$ matrix over a commutative ring, $S$ and $T$ finite sets with $s$ and $t$
  elements, and $\rho,\gamma$ bijections of $S\sqcup T$ onto $\{0,\dots,r-1\}$. Then
  \[
   \det\big(A_{\rho(i),\gamma(j)}\big)_{i,j\in S\sqcup T}\det\big(\mathrm{adj}(A)_{\gamma(j),\rho(i)}\big)_{i,j\in S}
   =\det(A)^s\det\big(A_{\rho(i),\gamma(j)}\big)_{i,j\in T}.
  \]
\item $C(\delta)=\delta^{r-1}$, and $\delta^NC\big(C(x)\big)=\delta^{(r-1)N}x$ for every $x$ homogeneous of
  degree $N$.
\item Let $x$ have row degree $\mu$. Then $C(x)$ has row degree $(\sum_a\mu_a-\mu_{r-1-c})_{c<r}$. If
  moreover $M\geq0$, $\mu_a\le M$ for every $a$, and $\delta^{\sum_a\mu_a}g=\delta^MC(x)$, then $g$ has row
  degree $\mu^\vee$.
\end{enumerate}
\end{lemma}
\leanref{Formalism/PieriBlockIdeals/LastGradedOfGradedAlone.lean}{Formalism.CohomologicalDescent.LD.genDet}
\leanref{Formalism/PieriBlockIdeals/LastGradedOfGradedAlone.lean}{Formalism.CohomologicalDescent.LD.cofMap}
\leanref{Formalism/PieriBlockIdeals/PieriColumnMinors/HighestWeightProdLeadingMinors.lean}{Formalism.CohomologicalDescent.LD.jacobi_complementary_minor}
\leanref{Formalism/PieriBlockIdeals/LastGradedOfGradedAlone.lean}{Formalism.CohomologicalDescent.LD.CofCof}
\leanref{Formalism/PieriBlockIdeals/PieriColumnMinors/HighestWeightProdLeadingMinors.lean}{Formalism.CohomologicalDescent.LD.cofCof}
\leanref{Formalism/PieriBlockIdeals/LastGradedOfGradedAlone.lean}{Formalism.CohomologicalDescent.LD.DualWeight}
\leanref{Formalism/PieriBlockIdeals/PieriColumnMinors/HighestWeightProdLeadingMinors.lean}{Formalism.CohomologicalDescent.LD.dualWeight}
\begin{proof}
(1) Put $A'=(A_{\rho(i),\gamma(j)})_{i,j}$ and $B=(\mathrm{adj}(A)_{\gamma(i),\rho(j)})_{i,j}$, indexed by
$S\sqcup T$. Then $A'B=\det(A)\cdot1$. Write both in blocks for $S\sqcup T$. The identities
$A'_{SS}B_{SS}+A'_{ST}B_{TS}=\det(A)\cdot1$ and $A'_{TS}B_{SS}+A'_{TT}B_{TS}=0$ say that
\[
 A'\begin{pmatrix}B_{SS}&0\\B_{TS}&1\end{pmatrix}=\begin{pmatrix}\det(A)\cdot1&A'_{ST}\\0&A'_{TT}\end{pmatrix}.
\]
Taking determinants gives $\det A'\det B_{SS}=\det(A)^s\det A'_{TT}$, and $\det B_{SS}$ is the determinant
of its transpose $(\mathrm{adj}(A)_{\gamma(j),\rho(i)})_{i,j\in S}$.

(2) The matrix $C(X)=(C(x_{a,s}))$ is the transpose of $\mathrm{adj}(X)$ with its rows and columns
reversed, so $C(\delta)=\det\mathrm{adj}(X)=\delta^{r-1}$. Applying $C$ to the entries of
$\mathrm{adj}(X)$ gives the adjugate of $C(X)$, so $C(C(x_{a,s}))=\mathrm{adj}(\mathrm{adj}(X))_{a,s}$. From
$X\mathrm{adj}(X)=\delta$ and $\mathrm{adj}(X)\mathrm{adj}(\mathrm{adj}(X))=\det\mathrm{adj}(X)=\delta^{r-1}$,
$\delta\,\mathrm{adj}(\mathrm{adj}(X))=X\mathrm{adj}(X)\mathrm{adj}(\mathrm{adj}(X))=\delta^{r-1}X$. So
$\delta\,C(C(x_v))=\delta^{r-1}x_v$ for every variable $x_v$; multiplying $N$ such identities gives the
claim for a monomial of degree $N$, and summing gives it for $x$.

(3) The entry $\mathrm{adj}(X)_{i,j}$ is the determinant of $X$ with its $j$-th row replaced by the $i$-th
unit vector, so each of its terms has degree $0$ in row $j$ and degree $1$ in every other row. So
$C(x_{a,s})$ has degree $0$ in row $r-1-a$ and $1$ in every other row, and a monomial of row degree $\mu$
goes to a polynomial of row degree $(\sum_a\mu_a-\mu_{r-1-c})_c$. As $\delta$ has degree $1$ in every row,
$\delta^MC(x)$ has row degree $(M+\sum_a\mu_a-\mu_{r-1-c})_c=(\sum_a\mu_a+\mu^\vee_c)_c$. As
$\delta^{\sum_a\mu_a}$ is nonzero of row degree $(\sum_a\mu_a,\dots,\sum_a\mu_a)$ and $\mathbf Z[X]$ is a
domain, $g$ has row degree $\mu^\vee$.
% Lean: jacobi_complementary_minor (Matrix.fromBlocks); cofCof (CofCof): cofactor_det_eq_pow
% (cofactor_map_generic_matrix, Matrix.det_adjugate), cofactor_cofactor_variable (AlgHom.map_adjugate,
% adjugate_adjugate_smul) and pow_mul_eval_homogeneous_eq; dualWeight (DualWeight): cofactor_weighted_homog
% (adjugate_entry_weighted_homog, weight_omega_add_degree), det_pow_weighted_homog and
% weighted_homogeneous_cancel.
\end{proof}

\begin{lemma}\label{is-onto-step-exp-bound-short-last-nabla}
Let $r,M\geq0$ be integers and $\nu\in\mathbf N^r$ nonincreasing with $\nu_a\le M$ for every $a$, with
$\delta$ and $C$ as in Lemma~\ref{is-onto-step-exp-bound-short-last-cofactor}. For every
$x\in\nabla_{\mathbf Z}(\nu)$ there is $g\in\nabla_{\mathbf Z}(\nu^\vee)$ with
$\delta^{\sum_a\nu_a}g=\delta^MC(x)$. The same holds with the Weyl modules $\Delta_{\mathbf Z}$ in place of
$\nabla_{\mathbf Z}$.
\end{lemma}
\leanref{Formalism/PieriBlockIdeals/LastGradedOfGradedAlone.lean}{Formalism.CohomologicalDescent.LD.NablaDual}
\leanref{Formalism/PieriBlockIdeals/PieriColumnMinors/HighestWeightProdLeadingMinors.lean}{Formalism.CohomologicalDescent.LD.nablaDual}
\leanref{Formalism/PieriBlockIdeals/LastGradedOfGradedAlone.lean}{Formalism.CohomologicalDescent.LD.WeylDual}
\leanref{Formalism/PieriBlockIdeals/PieriColumnMinors/HighestWeightProdLeadingMinors.lean}{Formalism.CohomologicalDescent.LD.weylDual}
\begin{proof}
For $s\le r$ and rows $I(0),\dots,I(s-1)$ let $m_I$ be the minor of $X$ on these rows and the first $s$
columns; it lies in $\nabla_{\mathbf Z}(\omega_s)$, where $\omega_s=(1,\dots,1,0,\dots,0)$ has $s$ entries
$1$. For $\nu'\in\mathbf N^r$ and $c\geq0$ write $\nu'+c$ for $(\nu'_a+c)_a$. By the proof of
Lemma~\ref{is-onto-step-exp-sum-shift}, $\nabla_{\mathbf Z}(\nu'+c)=\delta^c\nabla_{\mathbf Z}(\nu')$ and
$\Delta_{\mathbf Z}(\nu'+c)=\delta^c\Delta_{\mathbf Z}(\nu')$ for nonincreasing $\nu'$, and
$E^{(k)}_{ab}(\delta f)=\delta E^{(k)}_{ab}f$ for $a\ne b$, with the operators of
Definition~\ref{pd-weight-lattices}(1).

The dual Weyl module. Let $t=r-s$. If two of the rows $I(p)$ agree, $m_I=0$. Otherwise extend
$p\mapsto r-1-I(p)$ to a bijection $\rho$ of $\{0,\dots,s-1\}\sqcup\{0,\dots,t-1\}$ onto
$\{0,\dots,r-1\}$, and let $\gamma$ map $q$ in the first part to $r-1-q$ and $q$ in the second to $q$. By
Lemma~\ref{is-onto-step-exp-bound-short-last-cofactor}(1) for $A=X$, as $C(m_I)$ is the determinant of
$(\mathrm{adj}(X)_{r-1-q,r-1-I(p)})_{p,q<s}$, $\pm\delta\,C(m_I)=\delta^sm'$, where $m'$ is the minor on the
rows $\rho(i)$, $i<t$, of the second part and the first $t$ columns, an element of
$\nabla_{\mathbf Z}(\omega_t)$. By Lemma~\ref{pd-kappa-defects-restriction-standard}(1) it suffices to treat
a bideterminant $x=\prod_{k<K}m_{I_k}$ of shape $\nu$, with sizes $s_k$. Multiplying,
$\delta^KC(x)=\delta^{\sum_ks_k}g$ with $g\in\nabla_{\mathbf Z}(\sum_k\omega_{r-s_k})$, as products of
elements of dual Weyl modules lie in the dual Weyl module of the sum. Here $\sum_ks_k=\sum_t\nu_t$, and for
every $t$,
$\big(\sum_k\omega_{r-s_k}\big)_t+\nu_{r-1-t}=\#\{k:t<r-s_k\}+\#\{k:r-1-t<s_k\}=K$. If $K\le M$, then
$\delta^{M-K}g\in\nabla_{\mathbf Z}(\nu^\vee)$ and $\delta^{\sum\nu}\delta^{M-K}g=\delta^MC(x)$. If
$K\geq M$, then $g\in\nabla_{\mathbf Z}(\nu^\vee+(K-M))$, so $g=\delta^{K-M}g'$ with
$g'\in\nabla_{\mathbf Z}(\nu^\vee)$, and cancelling $\delta^{K-M}$ gives $\delta^{\sum\nu}g'=\delta^MC(x)$.

The Weyl module. Let $L$ be the set of $x$ with $\delta^MC(x)\in\delta^{\sum\nu}\Delta_{\mathbf Z}(\nu^\vee)$, a
subgroup. First, $h_\nu\in L$. Put $\nu_r=0$, so that $h_\nu=\prod_{s<r}\delta_{s+1}^{\nu_s-\nu_{s+1}}$ with
$\delta_s$ the leading principal $s\times s$ minor. Lemma~\ref{is-onto-step-exp-bound-short-last-cofactor}(1)
with $\rho=\gamma$ as above gives $\delta\,C(\delta_s)=\delta^s\delta_{r-s}$. As
$\sum_s(\nu_s-\nu_{s+1})=\nu_0$ and $\sum_s(s+1)(\nu_s-\nu_{s+1})=\sum_s\nu_s$,
$\delta^{\nu_0}C(h_\nu)=\delta^{\sum\nu}Q$ with $Q=\prod_{s<r}\delta_{r-s-1}^{\nu_s-\nu_{s+1}}$. Reindexing,
$h_{\nu^\vee}=\delta^{M-\nu_0}Q$, as $\nu^\vee_{r-1}=M-\nu_0$ and
$\nu^\vee_t-\nu^\vee_{t+1}=\nu_{r-2-t}-\nu_{r-1-t}$ for $t<r-1$. So
$\delta^{\sum\nu}h_{\nu^\vee}=\delta^MC(h_\nu)$.

Next, $L$ is stable under $E^{(k)}_{ab}$ for $a\ne b$. The substitution that adds $y$ times row $a$ of $X$
to row $b$ multiplies $X$ on the left by the transvection $1+ye_{ba}$, whose adjugate is $1-ye_{ba}$; so it
multiplies $\mathrm{adj}(X)$ on the right by $1-ye_{ba}$. Comparing entries, it acts on $C(f)$ as $C$ of
the substitution that adds $-y$ times row $r-1-b$ to row $r-1-a$, acting on $f$; taking coefficients of
$y^k$, $E^{(k)}_{ab}C(f)=(-1)^kC(E^{(k)}_{r-1-b,r-1-a}f)$. So if $\delta^MC(f)=\delta^{\sum\nu}g$ with
$g\in\Delta_{\mathbf Z}(\nu^\vee)$, then
$\delta^MC(E^{(k)}_{ab}f)=(-1)^k\delta^ME^{(k)}_{r-1-b,r-1-a}C(f)=\delta^{\sum\nu}(-1)^kE^{(k)}_{r-1-b,r-1-a}g$,
and $\Delta_{\mathbf Z}(\nu^\vee)$ is stable under the operators.

Finally, $\Delta_{\mathbf Z}(\nu)\subset L$. Let $L'$ be the set of $x$ all of whose components of fixed row
degree lie in $L$; then $L'\subset L$. The vector $h_\nu$ has row degree $\nu$, so it lies in $L'$. An
operator $E^{(k)}_{ab}$ with $a\ne b$ maps a polynomial of row degree $\mu'$ to one of row degree
$\mu'+k(e_a-e_b)$, so $L'$ is stable under it; and $E^{(k)}_{aa}$ multiplies a polynomial of row degree
$\mu'$ by $\binom{\mu'_a}k$, so $L'$ is stable under it too. As $\Delta_{\mathbf Z}(\nu)$ is spanned by the
images of $h_\nu$ under products of the operators, $\Delta_{\mathbf Z}(\nu)\subset L'\subset L$.
% Lean: nablaDual (NablaDual): standardMonomialSpan, cofactor_bideterminant (cofactor_initial_minor with
% injective_extends_to_equiv, finSumEquivOfAddEq, jacobi_complementary_minor, det_submatrix_eq_or_neg,
% initMinor_mem; dualWeyl_prod), shape_total_degree, shape_complement_sum, det_pow_mul_mem_dual_weyl and
% dual_weyl_shift_eq_det_pow_mul. weylDual (WeylDual): weyl_le_of_stable_submodule (la_weyl_le_of_stable,
% hwRowWeight, divided_power_weight_vector, divided_power_diag_weight_vector), highest_weight_dual_cofactor
% (cofactor_leading_minor, cofactor_highest_weight, highest_weight_dual_weight), divided_power_cofactor
% (cofactor_intertwines_generator: generator_map_generic_matrix, adjugate_transvection', eval_neg_x_coeff),
% divided_power_det_pow_mul and weylRowDivStable.
\end{proof}

\begin{lemma}\label{is-onto-step-exp-bound-short-last-duality}
Let $l$ be a prime, $r,M\geq0$ and $e\geq0$ integers, $\nu\in\mathbf N^r$ nonincreasing and
$\mu\in\mathbf N^r$, with all entries at most $M$. Then $(\nu,\mu)$ has exponent at most $e$ at $l$ if and
only if $(\nu^\vee,\mu^\vee)$ does (Definition~\ref{is-onto-step-exp-bound-short-branching}(1),
Lemma~\ref{is-onto-step-exp-bound-short-last-dual}).
\end{lemma}
\leanref{Formalism/PieriBlockIdeals/LastGradedOfGradedAlone.lean}{Formalism.CohomologicalDescent.LD.ExpDual}
\leanref{Formalism/PieriBlockIdeals/PieriColumnMinors/HighestWeightProdLeadingMinors.lean}{Formalism.CohomologicalDescent.LD.expDual}
\leanref{Formalism/PieriBlockIdeals/PieriColumnMinors/HighestWeightProdLeadingMinors.lean}{Formalism.CohomologicalDescent.LD.exp_le_of_dual_exp_le}
\begin{proof}
As $(\nu^\vee)^\vee=\nu$, $(\mu^\vee)^\vee=\mu$ and $\nu^\vee$ is nonincreasing
(Lemma~\ref{is-onto-step-exp-bound-short-last-dual}(1)), it suffices to show that the condition for
$(\nu^\vee,\mu^\vee)$ implies it for $(\nu,\mu)$. Let $x\in\nabla_{\mathbf Z}(\nu)$ have row degree $\mu$; as
$0\in\Delta_{\mathbf Z}(\nu)$, we may assume $x\ne0$. Put $n=\sum_a\nu_a$; then $\sum_a\mu_a=n$ by
Lemma~\ref{is-onto-step-exp-bound-short-last-weights}, and $x$ is homogeneous of degree $n$. By
Lemma~\ref{is-onto-step-exp-bound-short-last-nabla}, $\delta^ng=\delta^MC(x)$ for some
$g\in\nabla_{\mathbf Z}(\nu^\vee)$, and $g$ has row degree $\mu^\vee$ by
Lemma~\ref{is-onto-step-exp-bound-short-last-cofactor}(3). By assumption $ml^eg\in\Delta_{\mathbf Z}(\nu^\vee)$
for some $m$ not divisible by $l$. Lemma~\ref{is-onto-step-exp-bound-short-last-nabla} for $\nu^\vee$, whose
dual is $\nu$, gives $y\in\Delta_{\mathbf Z}(\nu)$ with $\delta^{n^\vee}y=\delta^MC(ml^eg)$, where
$n^\vee=\sum_a\nu^\vee_a=rM-n$.

Apply $C$ to $\delta^ng=\delta^MC(x)$: by
Lemma~\ref{is-onto-step-exp-bound-short-last-cofactor}(2), $\delta^{(r-1)n}C(g)=\delta^{(r-1)M}C(C(x))$ and
$\delta^nC(C(x))=\delta^{(r-1)n}x$. Hence
\[
 \delta^{n^\vee+(r-1)n+n}y=ml^e\delta^{M+n}\delta^{(r-1)n}C(g)=ml^e\delta^{M+(r-1)M}\delta^nC(C(x))
 =ml^e\delta^{rM+(r-1)n}x.
\]
As $n^\vee+rn=rM+(r-1)n$ and $\mathbf Z[X]$ is a domain, $y=ml^ex$. So $ml^ex\in\Delta_{\mathbf Z}(\nu)$.
% Lean: expDual = expDualOfLattices nablaDual weylDual dualWeight cofCof; exp_le_of_dual_exp_le
% (dominance_of_nabla_weight for the sum, weighted_homog_is_homogeneous, dual_weight_total, the exponent
% identity hE and mul_left_cancel); the case r = 0 has delta = 1.
\end{proof}

\begin{lemma}\label{is-onto-step-exp-bound-short-last}
If the block exponent bound holds in graded form, then it holds in last graded form
(Definition~\ref{is-onto-step-exp-bound-short-branching}(4) and~(5)).
\end{lemma}
\leanref{Formalism/PieriBlockIdeals/LastGradedOfGradedAlone.lean}{Formalism.CohomologicalDescent.LD.LastGradedOfGradedAlone}
\leanref{Formalism/PieriBlockIdeals/PieriColumnMinors/HighestWeightProdLeadingMinors.lean}{Formalism.CohomologicalDescent.LD.last_graded_of_graded_alone}
\begin{proof}
Let $l$, $p,q,w,M$, $\lambda_2$, $\xi$, $i'$ and $e_L$ be as in
Definition~\ref{is-onto-step-exp-bound-short-branching}(5), with the hypothesis on the pairs
$(\lambda_2,\mu_2)$. Suppose first that $q>0$, so that $w\le\lambda_{2,0}\le M$. Write $\kappa^\vee$ as in
Lemma~\ref{is-onto-step-exp-bound-short-last-dual}, with this $M$, and put $\xi'_a=-\xi_{q+p-1-a}$, which is
nonincreasing. By Lemma~\ref{is-onto-step-exp-bound-short-last-dual}(1) and~(2), $\lambda_2^\vee$ is
nonincreasing with $\lambda^\vee_{2,b}\le M-w$, and $\lambda_2w^p$ is nonincreasing with entries at most
$M$.

The pairs $(\lambda_2^\vee,\mu'_2)$ have exponent at most $e_L$ for every $\mu'_2\in\mathbf N^q$ with
$\mathrm{dp}_{\xi'}\big((M-w)^p\lambda_2^\vee,(M-w)^p\mu'_2\big)\le i'$. Indeed, let
$y\in\nabla_{\mathbf Z}(\lambda_2^\vee)$ be nonzero of row degree $\mu'_2$; then $\mu'_{2,b}\le M$ by
Lemma~\ref{is-onto-step-exp-bound-short-last-weights}. Put $\mu_2=\mu_2'^\vee$, so that $\mu_2^\vee=\mu'_2$.
By Lemma~\ref{is-onto-step-exp-bound-short-last-dual}(2) and~(4),
$\mathrm{dp}_\xi(\lambda_2w^p,\mu_2w^p)=\mathrm{dp}_{\xi'}\big((M-w)^p\lambda_2^\vee,(M-w)^p\mu'_2\big)\le i'$,
so $(\lambda_2,\mu_2)$ has exponent at most $e_L$ by hypothesis, and so has $(\lambda_2^\vee,\mu'_2)$ by
Lemma~\ref{is-onto-step-exp-bound-short-last-duality}. This gives the condition for $y$.

Now let $i\geq0$ and $\mu\in\mathbf N^{q+p}$ with $\sum_{a<p}\mu_{q+a}=pw+i$ and
$\mathrm{dp}_\xi(\lambda_2w^p,\mu)\le i'$, and let $x\in\nabla_{\mathbf Z}(\lambda_2w^p)$ be nonzero of row
degree $\mu$. By Lemma~\ref{is-onto-step-exp-bound-short-last-weights}, $\mu_a\le M$ for every $a$. By
Lemma~\ref{is-onto-step-exp-bound-short-last-dual}(3) and~(4),
$\sum_{a<p}\mu^\vee_a+i=p(M-w)$ and
$\mathrm{dp}_{\xi'}\big((M-w)^p\lambda_2^\vee,\mu^\vee\big)=\mathrm{dp}_\xi(\lambda_2w^p,\mu)\le i'$. The
graded form for $(l,p,q,M-w,\lambda_2^\vee,\xi',i',e_L)$ now shows that
$\big((M-w)^p\lambda_2^\vee,\mu^\vee\big)=\big((\lambda_2w^p)^\vee,\mu^\vee\big)$ has exponent at most
$e_L+v_l\big(D_n(p,M-w,i,\lambda_2^\vee)\big)$. By
Lemma~\ref{is-onto-step-exp-bound-short-last-duality}, so has $(\lambda_2w^p,\mu)$, which gives the
condition for $x$; the condition holds for $x=0$ trivially.

If $q=0$, the argument above applies with $\max(M,w)$ in place of $M$, as it needs only $w\le M$ besides
the hypotheses. For $q=0$ every $D_n(p,\cdot,i,\cdot)$ is $1$, as $\mu(c)=0$ for every $c$, so the bound it
gives is $e_L$, the bound required.
% Lean: last_graded_of_graded_alone = lastGradedOfGradedAlone = lastGradedOfGraded expDual, with expDual =
% expDualOfLattices nablaDual weylDual dualWeight cofCof (Lemma is-onto-step-exp-bound-short-last-duality).
% lastGradedOfGraded splits q = 0 (max M w and d_n_empty_eq_one) from q > 0; last_graded_at_of_exp_dual uses
% block_weight_last_antitone, wtBound, exp_le_of_nonzero_weight_vectors, blockDual, depthDual,
% depth_two_of_dual_weights (depth2Dual) and dualWtSpec.
\end{proof}

\begin{definition}\label{is-onto-step-exp-bound-short-top-data}
Let $l$ be a prime and $N,k,t,h\geq0$ integers, $\mathcal T$ a finite set and $T_2\subset\mathcal T$. For
$\tau\in\mathcal T$ let $\lambda_\tau$ be a sequence of natural numbers, $j_\tau,q_\tau,e^{\mathrm
c}_\tau,e_\tau,i_\tau\geq0$ integers and $x^-_\tau,x^+_\tau$ integers. Put $U_\tau(a)=\lambda_{\tau,a}+N-1-a$
and $\mathrm{top}_\tau=U_\tau(k+j_\tau)$, and let $\lfloor y\rfloor=0$ for $y<0$.
\begin{enumerate}
\item An index $\tau$ has the \emph{shape} for these data if $k+j_\tau+t+q_\tau=N$;
  $\lambda_{\tau,a'}\le\lambda_{\tau,a}$ for $a\le a'<N$; $\lambda_{\tau,a}=\lambda_{\tau,k+j_\tau}$ for
  $k+j_\tau\le a<k+j_\tau+t$; $\lambda_{\tau,0}+N\le h+\lambda_{\tau,N-1}+1$;
  $x^-_\tau\le U_\tau(a)\le x^+_\tau$ for $a<k$; $U_\tau(a)<x^-_\tau$ for $k\le a<k+j_\tau$;
  $\mathrm{top}_\tau<x^-_\tau\le x^+_\tau$; and $x^+_\tau-\mathrm{top}_\tau\le h$.
\item Let $d_\tau$ be the least integer $d\geq1$ with $l\mid x^-_\tau-\mathrm{top}_\tau-d$; let
  $\delta_\tau$ be $1$ if some integer in $[x^-_\tau,x^+_\tau]$ is congruent to $\mathrm{top}_\tau$ modulo
  $l$, and $0$ otherwise; let $\gamma_\tau=\min(k,t,i_\tau/d_\tau)$ if $d_\tau<x^-_\tau-\mathrm{top}_\tau$,
  and $\gamma_\tau=0$ otherwise; and let $H_\tau=0$ if $h<l$ and $H_\tau=(i_\tau/l)\big(1+\log(h/l)\big)$
  if $h\geq l$. Put
  \[
   C'_\tau=(q_\tau+2j_\tau)\frac{h-1}{l-1}+e^{\mathrm c}_\tau+k\Big(\frac{x^+_\tau-x^-_\tau}l+\delta_\tau\Big)
   +\gamma_\tau+H_\tau,
  \]
  \[
   C^{\mathrm s}_\tau=(q_\tau+2j_\tau)\frac{h+N}{l-1}+e^{\mathrm c}_\tau
   +k\Big(\frac{x^+_\tau-x^-_\tau}l+\delta_\tau\Big)+\gamma_\tau+\frac{i_\tau}l\Big(1+\log^+\frac hl\Big),
  \]
  with $\log^+$ as in Definition~\ref{is-onto-step-exp-bound-recursion}(3), and
  $J'=\sum_{\tau\in T_2}\min(e_\tau,\lfloor C'_\tau\rfloor)+\sum_{\tau\notin T_2}e_\tau$; let
  $J^{\mathrm s}$ be the same sum with $C^{\mathrm s}_\tau$ in place of $C'_\tau$.
\item Let $X$ consist of the data of a Hida setting of rank $N$ and weights
  $\lambda=(\lambda_\tau)_\tau$ (Definition~\ref{is-onto-step-exp-bound-hida}(1)), except that its ring
  $O$ is only required to be a Noetherian domain. Then $X$ has the \emph{top inputs} for these data if:
  every integer not divisible by $l$ is a unit of $O$; $h<l^2$ and $t\geq1$; every $\tau\in T_2$ has the
  shape; for every $\tau\in T_2$ and every $\mu\in\mathbf N^k$, the pair
  $\big((\lambda_{\tau,a})_{a<k},\mu\big)$ has exponent at most $e^{\mathrm c}_\tau$ at $l$
  (Definition~\ref{is-onto-step-exp-bound-short-branching}(1)); for every $\tau$, with
  $\nu_\tau=(\lambda_{\tau,a})_{a<N}$, $ml^{e_\tau}\nabla_{\mathbf Z}(\nu_\tau)\subset\Delta_{\mathbf
  Z}(\nu_\tau)$ for some integer $m$ not divisible by $l$; and $X$ has constituent data at $l$
  (Definition~\ref{is-onto-step-exp-bound-hida}(4), with $N$ in place of $N_r$) such that
  $\sum_{a<k+j_\tau}(\lambda_{\tau,a}-\omega_{\mu,\tau,a})\le i_\tau$ for every relevant $\mu$ and every
  $\tau\in T_2$.
\end{enumerate}
\end{definition}
\leanref{Formalism/PieriBlockIdeals/BlockIdealBound5Derived.lean}{Formalism.CohomologicalDescent.DepthNonnegative.Shape}
\leanref{Formalism/PieriBlockIdeals/Statement.lean}{Formalism.CohomologicalDescent.d0}
\leanref{Formalism/PieriBlockIdeals/Statement.lean}{Formalism.CohomologicalDescent.IlInd}
\leanref{Formalism/PieriBlockIdeals/FactorialValuationSum.lean}{Formalism.CohomologicalDescent.UnitNormBounds.harm}
\leanref{Formalism/PieriBlockIdeals/FactorialValuationSum.lean}{Formalism.CohomologicalDescent.UnitNormBounds.labelSumBoundPos}
\leanref{Formalism/PieriBlockIdeals/FactorialValuationSum.lean}{Formalism.CohomologicalDescent.UnitNormBounds.cPrime}
\leanref{Formalism/PieriBlockIdeals/BlockIdealBound5Derived.lean}{Formalism.CohomologicalDescent.DepthNonnegative.cFloor}
\leanref{Formalism/PieriBlockIdeals/LatticeInclusionRealSettingAlt.lean}{Formalism.CohomologicalDescent.ShortWeightBounds.cShort6}
\leanref{Formalism/PieriBlockIdeals/LatticeInclusionRealSettingAlt.lean}{Formalism.CohomologicalDescent.ShortWeightBounds.cFloorShort6}
\leanref{Formalism/PieriBlockIdeals/BlockIdealBound5Derived.lean}{Formalism.CohomologicalDescent.DepthNonnegative.cTau}
\leanref{Formalism/PieriBlockIdeals/BlockIdealBound5Derived.lean}{Formalism.CohomologicalDescent.DepthNonnegative.TopInputs}
% Lean: Shape N k j t q h lam xm xp (fields what2, anti, block, spread, cloud, lows, above, le, top), with
% top_tau = UN N lam (k + j); C'_tau is cPrime (harm, labelSumBoundPos with IlInd for delta and d0 for d),
% floored by cFloor; C^s_tau is cShort6, floored by cFloorShort6; J' and J^s are the sums of cTau. In
% TopInputs, the exponent clauses are ExpLe l k (res k lam) mu ecl and LA.InclZ N N (res N lam) l ex, and
% the clauses on A, e, T, the L_mu and the relevant mu are those of Definition
% is-onto-step-exp-bound-hida(4) (FrameA, FrameT, SimpleFacts, FirstInputWeightElim).

\begin{lemma}\label{is-onto-step-exp-bound-short-top-window-dominance}
Let $r\geq0$, $\nu\in\mathbf N^r$ nonincreasing and $\mu\in\mathbf N^r$. If $\nabla_{\mathbf Z}(\nu)$
contains a nonzero polynomial of row degree $\mu$
(Definition~\ref{is-onto-step-exp-bound-short-branching}), then $\sum_{a<c}\mu_a\le\sum_{a<c}\nu_a$ for
every $c\geq0$, and $\sum_a\mu_a=\sum_a\nu_a$.
\end{lemma}
\leanref{Formalism/PieriBlockIdeals/PieriColumnMinors/LastGradedImpliesOneBlock.lean}{Formalism.CohomologicalDescent.LD.dominance_of_nabla_weight}
\leanref{Formalism/GradedWeylDivisibility/FinalC.lean}{Formalism.GradedWeylDivisibility.dualWeylDominanceZ}
\begin{proof}
The equality of the sums is Lemma~\ref{is-onto-step-exp-bound-short-last-weights}. Let $m$ be a monomial
of the polynomial; it has degree $\mu_a$ in row $a$. By
Lemma~\ref{pd-kappa-defects-restriction-standard}(1) the polynomial is an integer combination of
bideterminants of shape $\nu$, so $m$ is a monomial of one of them, $\prod_{k<K}m_{I_k}$, where $m_{I_k}$
is the minor on the rows $I_k(0),\dots,I_k(s_k-1)$ and the first $s_k$ columns, and
$\#\{k:t<s_k\}=\nu_t$ for every $t$. So $m=\prod_km_k$ with $m_k$ a monomial of $m_{I_k}$, and $m_k$ has
degree $1$ in each of $s_k$ distinct rows and degree $0$ in the others. Hence $m_k$ has degree at most
$\min(c,s_k)$ in the first $c$ rows, and
\[
 \sum_{a<c}\mu_a\le\sum_{k<K}\min(c,s_k)=\sum_{k<K}\#\{t<c:t<s_k\}=\sum_{t<c}\nu_t.
\]
% Lean: dominance_of_nabla_weight reads dualWeylDominanceZ at a monomial of the support:
% standardMonomialSpan with Submodule.span_induction, bidet_rowPart_lemma (initMinor_support) and
% sum_rowNu_lemma_2; the equality of the sums comes from colSums.
\end{proof}

\begin{lemma}\label{is-onto-step-exp-bound-short-top-window-products}
Let $p,w,i,m,q\geq0$ be integers and $\lambda_2\in\mathbf N^q$, and let $\mathrm{Sh}(p,w,i,\lambda_2)$,
$b(\lambda_2,\tau,c)$, $\mu(c)$ and $D_n=D_n(p,w,i,\lambda_2)$ be as in
Definition~\ref{is-onto-step-exp-bound-short-branching}(3).
\begin{enumerate}
\item $\mu(w)=0$, and $b(\lambda_2,\tau,c)\le p$ for every $\tau\in\mathrm{Sh}(p,w,i,\lambda_2)$ and every
  $c\in\mathbf Z$, so $\mu(c)\le p$.
\item Suppose that $m\le\lambda_{2,b}\le w$ for every $b$. If $\mu(c)\ne0$, then $m-q<c<w$. The integer
  $D_n$ is positive and divides $(n!)^p$, where $n=\max(0,w+q-m-1)$, and $(l-1)v_l(D_n)\le pn$ for every
  prime $l$.
\end{enumerate}
\end{lemma}
\leanref{Formalism/PieriBlockIdeals/PieriColumnMinors/DividedPowerZeroZeroVector.lean}{Formalism.CohomologicalDescent.LD.ld_mu_w}
\leanref{Formalism/PieriBlockIdeals/PieriColumnMinors/DividedPowerZeroZeroVector.lean}{Formalism.CohomologicalDescent.LD.ld_boxes_le}
\leanref{Formalism/PieriBlockIdeals/StatementLD.lean}{Formalism.CohomologicalDescent.LD.MuLe}
\leanref{Formalism/PieriBlockIdeals/PieriColumnMinors/DividedPowerZeroZeroVector.lean}{Formalism.CohomologicalDescent.LD.muLe}
\leanref{Formalism/PieriBlockIdeals/StatementLD.lean}{Formalism.CohomologicalDescent.LD.MuRange}
\leanref{Formalism/PieriBlockIdeals/PieriColumnMinors/DividedPowerZeroZeroVector.lean}{Formalism.CohomologicalDescent.LD.muRange}
\leanref{Formalism/PieriBlockIdeals/StatementLD.lean}{Formalism.CohomologicalDescent.LD.DnPos}
\leanref{Formalism/PieriBlockIdeals/PieriColumnMinors/DividedPowerZeroZeroVector.lean}{Formalism.CohomologicalDescent.LD.dnPos}
\leanref{Formalism/PieriBlockIdeals/StatementLD.lean}{Formalism.CohomologicalDescent.LD.DnDvd}
\leanref{Formalism/PieriBlockIdeals/PieriColumnMinors/DividedPowerZeroZeroVector.lean}{Formalism.CohomologicalDescent.LD.dnDvd}
\leanref{Formalism/PieriBlockIdeals/StatementLD.lean}{Formalism.CohomologicalDescent.LD.DnVal}
\leanref{Formalism/PieriBlockIdeals/PieriColumnMinors/DividedPowerZeroZeroVector.lean}{Formalism.CohomologicalDescent.LD.dnVal}
\begin{proof}
(1) Every $\tau\in\mathrm{Sh}(p,w,i,\lambda_2)$ has $b(\lambda_2,\tau,w)=0$, so $\mu(w)=0$. Let $\tau$ be in
$\mathrm{Sh}(p,w,i,\lambda_2)$, and let $A$ be the set of $a$ with $\lambda_{2,a}\le c+a<\tau_a$. Suppose that
$A$ is nonempty, let $a^*$ be its largest element, and put $y=c+a^*$, so that $0\le y<\tau_{a^*}$. For
$a\in A$, $a\le a^*$, so $\tau_{a^*}\le\tau_a$, and $\lambda_{2,a}\le c+a\le y<\tau_{a^*}\le\tau_a$. So
$b(\lambda_2,\tau,c)=\#A\le\mathrm{col}(\lambda_2,\tau,y)\le p$.

(2) If $\mu(c)\ne0$, there are $\tau\in\mathrm{Sh}(p,w,i,\lambda_2)$ and $a<q$ with
$\lambda_{2,a}\le c+a<\tau_a$. Then $c\geq\lambda_{2,a}-a>m-q$. As $b(\lambda_2,\tau,w)=0$ and
$\lambda_{2,a}\le w\le w+a$, $\tau_a\le w+a$, so $c<w$. In the product defining $D_n$, a factor
$(w+q-y)^{\mu(y-q)}$ with $\mu(y-q)\ne0$ has $m+1\le y<w+q$; so $w+q-y$ is a positive integer, and the factor
divides $(w+q-y)^p$ by (1). The other factors are $1$. Hence $D_n>0$, and $D_n$ divides
$\prod_{m+1\le y<w+q}(w+q-y)^p=(n!)^p$. So $v_l(D_n)\le p\,v_l(n!)$, and $(l-1)v_l(n!)=n-s_l(n)\le n$ by
Legendre's formula, with $s_l(n)$ the sum of the digits of $n$ in base $l$.
% Lean: ld_mu_w; ld_boxes_le (Finset.max' of the set A, with the column condition of shapes) and muLe;
% muRange; dnPos; dnDvd (Finset.prod_dvd_prod_of_dvd, Finset.prod_ite_mem, Nat.descFactorial_self);
% dnVal (padicValInt_dvd_iff and Mathlib's sub_one_mul_padicValNat_factorial, Legendre's formula).
\end{proof}

\begin{lemma}\label{is-onto-step-exp-bound-short-top-window-peel}
Suppose that the block exponent bound holds in last graded form
(Definition~\ref{is-onto-step-exp-bound-short-branching}(5)). Let $l$ be a prime, $n,s,h\geq0$ integers,
$\lambda$ a sequence of natural numbers with $\lambda_b\le\lambda_a$ for $a\le b<n+s$ and, if $n+s\geq1$,
$\lambda_0+n+s\le h+\lambda_{n+s-1}+1$, $\xi$ a nonincreasing sequence of integers, $I\in\mathbf Z$ and
$e\geq0$. For $r\geq0$ write $\lambda|_r=(\lambda_a)_{a<r}$ and $\xi|_r=(\xi_a)_{a<r}$. Suppose that
$(\lambda|_n,\mu)$ has exponent at most $e$ at $l$ for every $\mu\in\mathbf N^n$ with
$\mathrm{dp}_{\xi|_n}(\lambda|_n,\mu)\le I$. Then $(\lambda|_{n+s},\mu)$ has exponent at most
$e+s\lfloor\max(0,h-1)/(l-1)\rfloor$ at $l$ for every $\mu\in\mathbf N^{n+s}$ with
$\mathrm{dp}_{\xi|_{n+s}}(\lambda|_{n+s},\mu)\le I$.
\end{lemma}
\leanref{Formalism/PieriBlockIdeals/BranchingBoundLastGradedOneRow.lean}{Formalism.CohomologicalDescent.LD.ExpUpTo}
\leanref{Formalism/PieriBlockIdeals/BranchingBoundLastGradedOneRow.lean}{Formalism.CohomologicalDescent.LD.PeelRows}
\leanref{Formalism/PieriBlockIdeals/PieriColumnMinors/LastGradedImpliesOneBlock.lean}{Formalism.CohomologicalDescent.LD.peelRows}
\leanref{Formalism/PieriBlockIdeals/BranchingBoundLastGradedOneRow.lean}{Formalism.CohomologicalDescent.LD.PeelBoundLast}
\leanref{Formalism/PieriBlockIdeals/PieriColumnMinors/LastGradedImpliesOneBlock.lean}{Formalism.CohomologicalDescent.LD.peelBoundLast}
\leanref{Formalism/PieriBlockIdeals/PieriColumnMinors/LastGradedImpliesOneBlock.lean}{Formalism.CohomologicalDescent.LD.exponent_bound_peel_one_row}
\leanref{Formalism/PieriBlockIdeals/PieriColumnMinors/LastGradedImpliesOneBlock.lean}{Formalism.CohomologicalDescent.LD.depth_exists_last_block}
\begin{proof}
An exponent bound for a pair persists for every larger exponent, as $ml^{e'}x=l^{e'-e}(ml^ex)$, and it
holds for $x=0$. So it suffices to treat nonzero $x$.

First, one row. Let $q\geq0$ with $\lambda_b\le\lambda_a$ for $a\le b\le q$, and $e'\geq0$ such that
$(\lambda|_q,\mu_2)$ has exponent at most $e'$ for every $\mu_2\in\mathbf N^q$ with
$\mathrm{dp}_{\xi|_q}(\lambda|_q,\mu_2)\le I$. We show that $(\lambda|_{q+1},\mu)$ has exponent at most
$e'+\lfloor\max(0,\lambda_0+q-\lambda_q-1)/(l-1)\rfloor$ for every $\mu\in\mathbf N^{q+1}$ with
$\mathrm{dp}_{\xi|_{q+1}}(\lambda|_{q+1},\mu)\le I$. Put $w=\lambda_q$, $M=\lambda_0$ and $\lambda_2=\lambda|_q$,
which is nonincreasing with $w\le\lambda_{2,b}\le M$; then $\lambda|_{q+1}=\lambda_2w^1$ in the notation of
Definition~\ref{is-onto-step-exp-bound-short-branching}(2). For $\mu_2\in\mathbf N^q$,
$\mathrm{dp}_{\xi|_{q+1}}(\lambda_2w^1,\mu_2w^1)=\mathrm{dp}_{\xi|_q}(\lambda_2,\mu_2)$, so the hypothesis of the
last graded form holds for $(l,1,q,w,M,\lambda_2,\xi|_{q+1},I,e')$. Let $\mu$ be as above and
$x\in\nabla_{\mathbf Z}(\lambda|_{q+1})$ nonzero of row degree $\mu$. As $\lambda|_{q+1}$ is nonincreasing,
Lemma~\ref{is-onto-step-exp-bound-short-top-window-dominance} gives
$\sum_{a<q}\mu_a\le\sum_b\lambda_{2,b}$ and $\sum_{a\le q}\mu_a=\sum_b\lambda_{2,b}+w$. So
$i=\sum_b\lambda_{2,b}-\sum_{a<q}\mu_a\geq0$ and $\mu_q=w+i$. The last graded form shows that $(\lambda_2w^1,\mu)$
has exponent at most $e'+v_l\big(D_n(1,M-w,i,\lambda_2^\vee)\big)$, with
$\lambda^\vee_{2,b}=M-\lambda_{2,q-1-b}\in[0,M-w]$. By
Lemma~\ref{is-onto-step-exp-bound-short-top-window-products}(2) with $m=0$,
$(l-1)v_l\big(D_n(1,M-w,i,\lambda_2^\vee)\big)\le\max(0,M-w+q-1)$, which gives the bound for $x$.

Now induct on $s$; for $s=0$ there is nothing to prove. Suppose the hypotheses hold for $s+1$. They hold for
$s$, as $\lambda_{n+s}\le\lambda_{n+s-1}$ if $n+s\geq1$. By induction, the conclusion holds for $s$, with
exponent $e+s\kappa_h$, where $\kappa_h=\lfloor\max(0,h-1)/(l-1)\rfloor$. The step above with $q=n+s$
gives the exponent $e+s\kappa_h+\lfloor\max(0,\lambda_0+n+s-\lambda_{n+s}-1)/(l-1)\rfloor$ for $n+s+1$
rows. As $\lambda_0+n+s+1\le h+\lambda_{n+s}+1$, the last term is at most $\kappa_h$.
% Lean: peelRows by induction on s; exponent_bound_peel_one_row (restriction_eq_block_weight, depth2Levi)
% applies peelBoundLast through last_graded_implies_one_block'; peelBoundLast uses
% exp_le_of_nonzero_weight_vectors, depth_exists_last_block (dominance_of_nabla_weight, depthLastIff),
% dnVal with m = 0, expLeMono and Nat.le_div_iff_mul_le; the comparison is Nat.div_le_div_right.
\end{proof}

\begin{lemma}\label{is-onto-step-exp-bound-short-top-window-block}
Suppose that the block exponent bound holds in last graded form
(Definition~\ref{is-onto-step-exp-bound-short-branching}(5)). Let $l$ be a prime, $n\geq0$ and $t\geq1$
integers, and $\lambda$ a sequence of natural numbers with $\lambda_b\le\lambda_a$ for $a\le b<n+t$ and
$\lambda_a=\lambda_n$ for $n\le a<n+t$. Write $\lambda|_r=(\lambda_a)_{a<r}$, and let $e_L\geq0$ be such that
$(\lambda|_n,\mu_2)$ has exponent at most $e_L$ at $l$ for every $\mu_2\in\mathbf N^n$. For an integer
$I\geq0$ let $B(I)$ be the largest of the numbers $v_l\big(D_n(t,\lambda_0-\lambda_n,i,\lambda^\vee)\big)$ for
$0\le i\le I$, where $\lambda^\vee=(\lambda_0-\lambda_{n-1-b})_{b<n}$. Then for every integer $I\geq0$ and every
$\mu\in\mathbf N^{n+t}$ with $\sum_{a<n}(\lambda_a-\mu_a)\le I$, $(\lambda|_{n+t},\mu)$ has exponent at most
$e_L+B(I)$ at $l$.
\end{lemma}
\leanref{Formalism/PieriBlockIdeals/BranchingBoundLastGradedOneRow.lean}{Formalism.CohomologicalDescent.LD.blockCost}
\leanref{Formalism/PieriBlockIdeals/BranchingBoundLastGradedOneRow.lean}{Formalism.CohomologicalDescent.LD.BlockStep}
\leanref{Formalism/PieriBlockIdeals/PieriColumnMinors/LastGradedImpliesOneBlock.lean}{Formalism.CohomologicalDescent.LD.blockStep}
\leanref{Formalism/PieriBlockIdeals/StatementLD.lean}{Formalism.CohomologicalDescent.LD.LemmaDLast}
\leanref{Formalism/PieriBlockIdeals/BranchingBoundLastGradedOneRow.lean}{Formalism.CohomologicalDescent.LD.LastGradedImpliesLast}
\leanref{Formalism/PieriBlockIdeals/PieriColumnMinors/LastGradedImpliesOneBlock.lean}{Formalism.CohomologicalDescent.LD.lastGradedImpliesLast}
\begin{proof}
Put $w=\lambda_n$, $M=\lambda_0$ and $\lambda_2=\lambda|_n$, which is nonincreasing with
$w\le\lambda_{2,b}\le M$; then $\lambda|_{n+t}=\lambda_2w^t$. Apply the last graded form to
$(l,t,n,w,M,\lambda_2,0,0,e_L)$, with $\xi=0$ and $i'=0$. As $\mathrm{dp}_0=0$, its hypothesis holds by the
choice of $e_L$, and it shows that $(\lambda_2w^t,\mu)$ has exponent at most
$e_L+v_l\big(D_n(t,M-w,i,\lambda^\vee)\big)$ for every $i\geq0$ and every $\mu$ with
$\sum_{a<t}\mu_{n+a}=tw+i$. Let $\mu$ be as in the statement and $x\in\nabla_{\mathbf Z}(\lambda|_{n+t})$
nonzero of row degree $\mu$; an exponent bound holds for $x=0$ and persists for larger exponents. As
$\lambda|_{n+t}$ is nonincreasing, Lemma~\ref{is-onto-step-exp-bound-short-top-window-dominance} gives
$i=\sum_{a<n}(\lambda_a-\mu_a)\geq0$ and $\sum_{a<t}\mu_{n+a}=tw+i$. As $i\le I$, the exponent
$e_L+v_l\big(D_n(t,M-w,i,\lambda^\vee)\big)$ is at most $e_L+B(I)$.
% Lean: blockStep (hD = lastGradedImpliesLast, the last graded form at xi = 0 and i' = 0):
% restriction_eq_block_weight, restriction_antitone, exp_le_of_nonzero_weight_vectors,
% depth_exists_last_block, expLeMono and Finset.le_sup for blockCost. The hypothesis 1 <= t is not used.
\end{proof}

\begin{lemma}\label{is-onto-step-exp-bound-short-top-window-steps}
Suppose that the block exponent bound holds in last graded form
(Definition~\ref{is-onto-step-exp-bound-short-branching}(5)). Let $l$ be a prime, $k,j,t,q,h\geq0$
integers with $t\geq1$, $N=k+j+t+q$, and $\lambda$ a sequence of natural numbers with $\lambda_b\le\lambda_a$
for $a\le b<N$, $\lambda_a=\lambda_{k+j}$ for $k+j\le a<k+j+t$, and $\lambda_0+N\le h+\lambda_{N-1}+1$.
Write $\lambda|_r=(\lambda_a)_{a<r}$, let $e^{\mathrm c}\geq0$ be such that $(\lambda|_k,\mu')$ has exponent
at most $e^{\mathrm c}$ at $l$ for every $\mu'\in\mathbf N^k$, and let $B(I)$ be as in
Lemma~\ref{is-onto-step-exp-bound-short-top-window-block} for $n=k+j$. Then for every integer $I\geq0$ and
every $\mu\in\mathbf N^N$ with $\sum_{a<k+j}(\lambda_a-\mu_a)\le I$, $(\lambda|_N,\mu)$ has exponent at most
$e^{\mathrm c}+(j+q)\lfloor\max(0,h-1)/(l-1)\rfloor+B(I)$ at $l$.
\end{lemma}
\leanref{Formalism/PieriBlockIdeals/BranchingBoundLastGradedOneRow.lean}{Formalism.CohomologicalDescent.LD.ApplicationSteps}
\leanref{Formalism/PieriBlockIdeals/PieriColumnMinors/LastGradedImpliesOneBlock.lean}{Formalism.CohomologicalDescent.LD.applicationSteps}
\begin{proof}
Write $\kappa_h=\lfloor\max(0,h-1)/(l-1)\rfloor$. Apply
Lemma~\ref{is-onto-step-exp-bound-short-top-window-peel} with $n=k$, $s=j$, $\xi=0$, $I=0$ and
$e=e^{\mathrm c}$. Its hypotheses hold: $\mathrm{dp}_0=0$; and if $k+j\geq1$, then
$\lambda_0+k+j\le h+\lambda_{k+j-1}+1$, as $\lambda_{N-1}\le\lambda_{k+j-1}$ and $t\geq1$. So
$(\lambda|_{k+j},\mu_2)$ has exponent at most $e^{\mathrm c}+j\kappa_h$ for every $\mu_2\in\mathbf N^{k+j}$.
Lemma~\ref{is-onto-step-exp-bound-short-top-window-block} with $n=k+j$ and $e_L=e^{\mathrm c}+j\kappa_h$
shows that $(\lambda|_{k+j+t},\mu')$ has exponent at most $e_L+B(I)$ for every $\mu'$ with
$\sum_{a<k+j}(\lambda_a-\mu'_a)\le I$. Let $\varpi$ be the sequence with $\varpi_a=1$ for $a<k+j$ and
$\varpi_a=0$ otherwise; it is nonincreasing, and $\mathrm{dp}_{\varpi|_r}(\lambda|_r,\mu')=
\sum_{a<k+j}(\lambda_a-\mu'_a)$ for $r\geq k+j$ and $\mu'\in\mathbf N^r$. So
Lemma~\ref{is-onto-step-exp-bound-short-top-window-peel} with $n=k+j+t$, $s=q$, $\xi=\varpi$, this $I$ and
$e=e_L+B(I)$ gives the exponent $e^{\mathrm c}+j\kappa_h+B(I)+q\kappa_h$.
% Lean: applicationSteps: peelRows with xi = fun _ => 0 and I = 0 (s3, s3'), blockStep with
% lastGradedImpliesLast (s2), peelRows with varpi (k + j) and varpi_antitone (s1), and the reordering e.
\end{proof}

\begin{lemma}\label{is-onto-step-exp-bound-short-top-window-label-sum-base}
Let $l$ be a prime and $k,t,h\geq0$ integers with $h<l^2$ and $l\le2h$. For a real $I\geq0$ and integers
$\beta,x^-,x^+$ with $0\le\beta<x^-\le x^+\le h$, let $\delta$ be $1$ if some integer in $[x^-,x^+]$ is
congruent to $\beta$ modulo $l$ and $0$ otherwise, $d$ the least integer $d\geq1$ with $l\mid x^--\beta-d$,
and $\gamma(I)=\min(k,t,I/d)$ if $d<x^--\beta$ and $\gamma(I)=0$ otherwise. Then
\[
 \sum_{x^-\le y\le x^+}k\,v_l(y-\beta)+\sum_{\beta<y<x^-}\min\Big(k,t,\frac I{x^--y}\Big)v_l(y-\beta)
 \le k\Big(\frac{x^+-x^-}l+\delta\Big)+\gamma(I)+\frac Il\Big(1+\log\frac hl\Big).
\]
\end{lemma}
\leanref{Formalism/PieriBlockIdeals/Statement.lean}{Formalism.CohomologicalDescent.labelSumBound}
\leanref{Formalism/PieriBlockIdeals/Statement.lean}{Formalism.CohomologicalDescent.LabelSum}
\leanref{Formalism/PieriBlockIdeals/PieriColumnMinors/LabelSumNeedsLeFalse.lean}{Formalism.CohomologicalDescent.labelSumLe_holds}
\begin{proof}
Every $y$ in the two sums has $0<y-\beta\le x^+\le h<l^2$, as $\beta\geq0$. So $v_l(y-\beta)\le1$, and
$v_l(y-\beta)=0$ unless $l\mid y-\beta$.

The first sum is therefore at most $k$ times the number of $y\in[x^-,x^+]$ with $l\mid y-\beta$. If there
are no such $y$, this number is $0$. Otherwise $\delta=1$. The map $y\mapsto\lfloor(y-x^-)/l\rfloor$ is
injective on these $y$, as they are congruent modulo $l$, and takes values in
$\{0,\dots,\lfloor(x^+-x^-)/l\rfloor\}$; so they number at most $(x^+-x^-)/l+1$. In both cases the first sum is
at most $k\big((x^+-x^-)/l+\delta\big)$.

The second sum is at most the sum of $\min\big(k,t,I/(x^--y)\big)$ over the set $S$ of the $y$ with
$\beta<y<x^-$ and $l\mid y-\beta$. Let $y\in S$. Then $l$ divides $(x^--\beta-d)-(y-\beta)=(x^--y)-d$. Also
$x^--y\geq1$ and $l$ divides $(x^--\beta)-(x^--y)=y-\beta$, so $d\le x^--y$ by the choice of $d$. Hence
$x^--y=d+l\,g(y)$ with $g(y)\geq0$ an integer. If $g(y)=0$, then $x^--y=d$, so $d<x^--\beta$ and the term for
$y$ is $\gamma(I)$. If $g(y)\geq1$, the term is at most $I/(x^--y)\le I/(l\,g(y))$. The map $g$ is injective on
$S$, and $l\,g(y)\le x^--y<x^-\le h$, as $y>\beta\geq0$; so $g(y)\le\lfloor h/l\rfloor$. With
$H_n=\sum_{1\le g\le n}1/g$, the second sum is therefore at most
\[
 \gamma(I)+\sum_{g=1}^{\lfloor h/l\rfloor}\frac I{lg}=\gamma(I)+\frac Il\,H_{\lfloor h/l\rfloor}.
\]
If $\lfloor h/l\rfloor\geq1$, then $H_{\lfloor h/l\rfloor}\le1+\log\lfloor h/l\rfloor\le1+\log(h/l)$. If
$\lfloor h/l\rfloor=0$, then $H_0=0$, while $\log(h/l)\geq1-l/h\geq-1$ as $l\le2h$, so $1+\log(h/l)\geq0$.
Adding the bounds for the two sums gives the claim.
% Lean: labelSumLe_holds: label_sum_le_filter (with label_vl_le_ind) for both sums; for the first,
% label_card_le (the injective map (x - xm) / l into Icc 0 ((xp - xm) / l)) or the empty filter, split on
% IlInd; for the second, label_second (leastCongruentPos_mem and leastCongruentPos_le for d, the map
% g = (xm - ct - d0) / l, injective with values in range (h / l + 1), and label_harm_le, which uses
% Mathlib's harmonic_le_one_add_log, and Real.one_sub_inv_le_log_of_pos when h / l = 0).
\end{proof}

\begin{lemma}\label{is-onto-step-exp-bound-short-top-window-label-sum}
Let $l$ be a prime and $k,t,h\geq0$ integers with $h<l^2$. For a real $I\geq0$ and integers
$\beta<x^-\le x^+$, let $\delta$ be $1$ if some integer in $[x^-,x^+]$ is congruent to $\beta$ modulo $l$ and
$0$ otherwise, $d$ the least integer $d\geq1$ with $l\mid x^--\beta-d$, $\gamma(I)=\min(k,t,I/d)$ if
$d<x^--\beta$ and $\gamma(I)=0$ otherwise, $H(I)=0$ if $h<l$ and $H(I)=(I/l)\big(1+\log(h/l)\big)$ if
$h\geq l$, and
\[
 L(I;x^-,x^+,\beta)=k\Big(\frac{x^+-x^-}l+\delta\Big)+\gamma(I)+H(I).
\]
\begin{enumerate}
\item If $x^+-\beta\le h$, then
  \[
   \sum_{x^-\le y\le x^+}k\,v_l(y-\beta)+\sum_{\beta<y<x^-}\min\Big(k,t,\frac I{x^--y}\Big)v_l(y-\beta)
   \le L(I;x^-,x^+,\beta).
  \]
\item $L(I;x^-,x^+,\beta)$ is nondecreasing in $I$, and $L(I;x^-+c,x^++c,\beta+c)=L(I;x^-,x^+,\beta)$ for
  every integer $c$.
\end{enumerate}
\end{lemma}
\leanref{Formalism/PieriBlockIdeals/Statement.lean}{Formalism.CohomologicalDescent.labelLHS}
\leanref{Formalism/PieriBlockIdeals/FactorialValuationSum.lean}{Formalism.CohomologicalDescent.UnitNormBounds.labelSumBoundPos}
\leanref{Formalism/PieriBlockIdeals/FactorialValuationSum.lean}{Formalism.CohomologicalDescent.UnitNormBounds.LabelSumPlus}
\leanref{Formalism/PieriBlockIdeals/PieriColumnMinors/UnitNormBounds.lean}{Formalism.CohomologicalDescent.UnitNormBounds.label_sum_plus_holds}
\leanref{Formalism/PieriBlockIdeals/PieriColumnMinors/UnitNormBounds.lean}{Formalism.CohomologicalDescent.UnitNormBounds.labelLHS_shift_to_zero}
\leanref{Formalism/PieriBlockIdeals/PieriColumnMinors/UnitNormBounds.lean}{Formalism.CohomologicalDescent.UnitNormBounds.indexIntegral_shift}
\leanref{Formalism/PieriBlockIdeals/PieriColumnMinors/UnitNormBounds.lean}{Formalism.CohomologicalDescent.UnitNormBounds.labelRHSPlus_monotone}
\leanref{Formalism/PieriBlockIdeals/PieriColumnMinors/UnitNormBounds.lean}{Formalism.CohomologicalDescent.UnitNormBounds.labelRHSPlus_shift_invariant}
\begin{proof}
(1) Every $y$ in the two sums has $0<y-\beta\le x^+-\beta\le h$. If $h<l$, then $0<y-\beta<l$, so
$v_l(y-\beta)=0$ for all these $y$. The left side is then $0$, while every term of $L(I;x^-,x^+,\beta)$ is
nonnegative.

Let $h\geq l$, so that $H(I)=(I/l)\big(1+\log(h/l)\big)$. First move $\beta$ to $0$. The substitution
$y\mapsto y-\beta$ carries the two sums for $(x^-,x^+,\beta)$ term by term to the two sums for
$(x^--\beta,x^+-\beta,0)$. An integer $y\in[x^-,x^+]$ is congruent to $\beta$ modulo $l$ if and only if
$y-\beta\in[x^--\beta,x^+-\beta]$ is congruent to $0$, so $\delta$ is the same for the two triples; and $d$ and
$\gamma(I)$ depend only on $x^--\beta$. So both sides of (1) are unchanged when $(x^-,x^+,\beta)$ is replaced
by $(x^--\beta,x^+-\beta,0)$. For the new triple, $0<x^--\beta\le x^+-\beta\le h$; moreover $h<l^2$ and
$l\le h\le2h$. So Lemma~\ref{is-onto-step-exp-bound-short-top-window-label-sum-base}, applied with $\beta=0$
to $x^--\beta$ and $x^+-\beta$, gives (1).

(2) $\gamma(I)$ and $H(I)$ are nondecreasing in $I$, as $\log(h/l)\geq0$ if $h\geq l$. The numbers
$x^+-x^-$, $\delta$ and $d$ depend only on the differences of $x^-$, $x^+$ and $\beta$.
% Lean: label_sum_plus_holds: for h < l, label_vl_zero_small, indexIntegral_nonneg and gated_min_nonneg;
% for h >= l, labelLHS_shift_to_zero and indexIntegral_shift move beta to 0, and labelSumLe_holds
% (Lemma is-onto-step-exp-bound-short-top-window-label-sum-base) applies at (xm - mB, xp - mB, 0) with
% l <= 2h. (2) is labelRHSPlus_monotone (harmonic_monotone, gated_min_monotone) and
% labelRHSPlus_shift_invariant.
\end{proof}

\begin{lemma}\label{is-onto-step-exp-bound-short-top-window-labels}
Let $k,j,t,w,M\geq0$ be integers, $\lambda\in\mathbf N^{k+j}$, and $x^-$, $x^+$ integers. Put $\beta=w+t-1$,
$W=\max(0,M-w)$, $U(a)=\lambda_a+k+j+t-1-a$ for $a<k+j$, and $\lambda^\vee=(M-\lambda_{k+j-1-b})_{b<k+j}$.
Suppose that $w\le\lambda_a\le M$ for every $a$, that $x^-\le U(a)\le x^+$ for $a<k$ and $U(a)<x^-$ for
$k\le a<k+j$, and that $\beta<x^-\le x^+$. For an integer $i\geq0$ let $\mu(c)$ be as in
Definition~\ref{is-onto-step-exp-bound-short-branching}(3) for $(t,W,i,\lambda^\vee)$, and put $c(y)=W+\beta-y$
for $y\in\mathbf Z$.
\begin{enumerate}
\item For every integer $y$, $\mu(c(y))\le j+\min\big(k,t,i/(x^--y)\big)$ if $y<x^-$, and
  $\mu(c(y))\le\min(k,t)$ if $y\geq x^-$.
\item $v_l\big(D_n(t,W,i,\lambda^\vee)\big)=\sum_{\beta<y<x^+}\mu(c(y))\,v_l(y-\beta)$ for every prime $l$.
\item Let $l$ be a prime and $h\geq0$ an integer with $x^+-\beta\le h<l^2$, and let $L$ be as in
  Lemma~\ref{is-onto-step-exp-bound-short-top-window-label-sum} for $l$, $k$, $t$ and $h$. Then for every
  real $I\geq i$,
  \[
   v_l\big(D_n(t,W,i,\lambda^\vee)\big)\le j\,\frac{h-1}{l-1}+L(I;x^-,x^+,\beta).
  \]
\end{enumerate}
\end{lemma}
\leanref{Formalism/PieriBlockIdeals/FactorialValuationSum.lean}{Formalism.CohomologicalDescent.UnitNormBounds.Rows}
\leanref{Formalism/PieriBlockIdeals/FactorialValuationSum.lean}{Formalism.CohomologicalDescent.UnitNormBounds.cont}
\leanref{Formalism/PieriBlockIdeals/FactorialValuationSum.lean}{Formalism.CohomologicalDescent.UnitNormBounds.MuBound}
\leanref{Formalism/PieriBlockIdeals/PieriColumnMinors/UnitNormBounds.lean}{Formalism.CohomologicalDescent.UnitNormBounds.multiplicity_bound_holds}
\leanref{Formalism/PieriBlockIdeals/FactorialValuationSum.lean}{Formalism.CohomologicalDescent.UnitNormBounds.DnValSum}
\leanref{Formalism/PieriBlockIdeals/PieriColumnMinors/UnitNormBounds.lean}{Formalism.CohomologicalDescent.UnitNormBounds.depth_valuation_sum_holds}
\leanref{Formalism/PieriBlockIdeals/FactorialValuationSum.lean}{Formalism.CohomologicalDescent.UnitNormBounds.FactorialValuationSum}
\leanref{Formalism/PieriBlockIdeals/PieriColumnMinors/UnitNormBounds.lean}{Formalism.CohomologicalDescent.UnitNormBounds.first_bound_holds}
\leanref{Formalism/PieriBlockIdeals/FactorialValuationSum.lean}{Formalism.CohomologicalDescent.UnitNormBounds.Step2Val}
\leanref{Formalism/PieriBlockIdeals/PieriColumnMinors/UnitNormBounds.lean}{Formalism.CohomologicalDescent.UnitNormBounds.second_valuation_holds}
\leanref{Formalism/PieriBlockIdeals/FactorialValuationSum.lean}{Formalism.CohomologicalDescent.UnitNormBounds.Step2ValUpTo}
\leanref{Formalism/PieriBlockIdeals/PieriColumnMinors/UnitNormBounds.lean}{Formalism.CohomologicalDescent.UnitNormBounds.second_valuation_up_to_holds}
\begin{proof}
First, the labels of boxes. Let $\tau\in\mathrm{Sh}(t,W,i,\lambda^\vee)$, $b<k+j$ and $s\geq0$ with
$\lambda^\vee_b\le s<\tau_b$, and put $a=k+j-1-b$. Then $w\le\lambda_a\le M$, so $W=M-w$ and
$\lambda^\vee_b=M-\lambda_a$, and $U(a)=M-\lambda^\vee_b+t+b$. Put $y=M+t-1-(s-b)$, so that $c(y)=s-b$. Then
\[
 \beta<y\le U(a)-1. \tag{$*$}
\]
Indeed $y=U(a)-(s-\lambda^\vee_b+1)\le U(a)-1$. If $y\le\beta$, then $s-b\geq M-w$, so $W+b\le s<\tau_b$; as
also $\lambda^\vee_b\le W\le W+b$, the row $b$ is counted in $b(\lambda^\vee,\tau,W)$, which is $0$, a
contradiction.

(1) If $\mathrm{Sh}(t,W,i,\lambda^\vee)$ is empty, $\mu=0$. Let $\tau$ be in it, and let $R$ be the set of
$b$ with $\lambda^\vee_b\le c(y)+b<\tau_b$, so that $\#R=b(\lambda^\vee,\tau,c(y))$; by $(*)$ for $s=c(y)+b$,
every $b\in R$ has $\beta<y\le U(k+j-1-b)-1$. Let $R_1$ be the set of $b\in R$ with $k+j-1-b\geq k$, and
$R_0=R\setminus R_1$. Every $b\in R_1$ has $b<j$, so $\#R_1\le j$; and if $y\geq x^-$, then $R_1$ is empty,
as $U(k+j-1-b)<x^-$ for $b\in R_1$. The map $b\mapsto k+j-1-b$ sends $R_0$ into $\{0,\dots,k-1\}$, so
$\#R_0\le k$, and $\#R_0\le\#R\le t$ by Lemma~\ref{is-onto-step-exp-bound-short-top-window-products}(1). Let
$y<x^-$ and $b\in R_0$. Then $U(k+j-1-b)\geq x^-$, so $\lambda^\vee_b\le M+t+b-x^-$, while
$\tau_b\geq c(y)+b+1=M+t+b-y$; so $\tau_b-\lambda^\vee_b\geq x^--y$. As $\tau_b\geq\lambda^\vee_b$ for every
$b$ and $\sum_b(\tau_b-\lambda^\vee_b)=i$, $\#R_0\cdot(x^--y)\le i$. Hence $\#R\le j+\min\big(k,t,i/(x^--y)\big)$
if $y<x^-$, and $\#R=\#R_0\le\min(k,t)$ if $y\geq x^-$. Taking the largest over $\tau$ gives (1).

(2) Writing $c=y'-(k+j)$ for the index $y'$ of the product,
$D_n(t,W,i,\lambda^\vee)=\prod_{-(k+j)\le c<W+i}(W-c)^{\mu(c)}$. The factor with $c=W$ is $1$, as $\mu(W)=0$
by Lemma~\ref{is-onto-step-exp-bound-short-top-window-products}(1), and the other bases are nonzero integers.
So $v_l(D_n)=\sum_c\mu(c)v_l(W-c)$. Substitute $c=c(y)$, so that $W-c=y-\beta$. If $\mu(c(y))\ne0$, there are
$\tau$ and $b$ with $\lambda^\vee_b\le c(y)+b<\tau_b$, and $(*)$ gives $\beta<y\le U(a)-1$ with $a=k+j-1-b$;
here $U(a)\le x^+$ if $a<k$, and $U(a)<x^-\le x^+$ otherwise, so $y<x^+$. Conversely, if $\mu(c(y))\ne0$, then
$-(k+j)<c(y)<W$ by Lemma~\ref{is-onto-step-exp-bound-short-top-window-products}(2) with $m=0$, as
$0\le\lambda^\vee_b\le W$. So the nonzero terms of the two sums correspond, and (2) follows.

(3) Split the sum in (2) at $x^-$. By (1), the terms with $\beta<y<x^-$ are at most
$j\,v_l(y-\beta)+\min\big(k,t,i/(x^--y)\big)v_l(y-\beta)$, and those with $x^-\le y<x^+$ are at most
$k\,v_l(y-\beta)$. By Legendre's formula,
$\sum_{\beta<y<x^-}v_l(y-\beta)=v_l\big((x^--\beta-1)!\big)\le(x^--\beta-1)/(l-1)$. Adding the nonnegative
term $k\,v_l(x^+-\beta)$, Lemma~\ref{is-onto-step-exp-bound-short-top-window-label-sum}(1) with $i$ in place
of $I$ bounds the rest by $L(i;x^-,x^+,\beta)$. Finally $x^--\beta\le x^+-\beta\le h$, and $L(i)\le L(I)$ by
Lemma~\ref{is-onto-step-exp-bound-short-top-window-label-sum}(2).
% Lean: rowStart_rev_eq and content_label_eq (U(a) and c(y) for w <= M), good_labels_holds (the
% bound (*)), counts_rows_holds (R_1, with card_filter_not_rev_lt_le), counts_cloud_holds
% (card_filter_rev_lt_le and ld_boxes_le), counts_deep_holds and boxesOfLabel_le_bound, whose largest
% value over the shapes is multiplicity_bound_holds; depth_valuation_sum_holds (padicValInt_prod_eq_sum,
% mem_label_interval_of_mu_ne_zero, muSupport, Finset.sum_bij_ne_zero); second_valuation_holds
% (first_bound_holds, which is Legendre's formula through sum_padicVal_shift_eq_factorial, and
% label_sum_plus_holds) and second_valuation_up_to_holds (labelRHSPlus_monotone).
\end{proof}

\begin{lemma}\label{is-onto-step-exp-bound-short-top-window}
Suppose that the block exponent bound holds in last graded form
(Definition~\ref{is-onto-step-exp-bound-short-branching}(5)). Let $l$ be a prime and $N,k,t,h\geq0$
integers with $t\geq1$ and $h<l^2$, and let $\lambda_\tau$, $j_\tau$, $q_\tau$, $e^{\mathrm c}_\tau$,
$i_\tau$, $x^-_\tau$ and $x^+_\tau$ be data at an index $\tau$ as in
Definition~\ref{is-onto-step-exp-bound-short-top-data}, such that $\tau$ has the shape and
$\big((\lambda_{\tau,a})_{a<k},\mu'\big)$ has exponent at most $e^{\mathrm c}_\tau$ at $l$ for every
$\mu'\in\mathbf N^k$. Put $\nu_\tau=(\lambda_{\tau,a})_{a<N}$. Then for every $\mu\in\mathbf N^N$ with
$\sum_{a<k+j_\tau}(\lambda_{\tau,a}-\mu_a)\le i_\tau$, the pair $(\nu_\tau,\mu)$ has exponent at most
$\lfloor C'_\tau\rfloor$ at $l$.
\end{lemma}
\leanref{Formalism/PieriBlockIdeals/BlockIdealBound5Derived.lean}{Formalism.CohomologicalDescent.DepthNonnegative.ApplicationWindow}
\leanref{Formalism/PieriBlockIdeals/PieriColumnMinors/DepthNonnegative.lean}{Formalism.CohomologicalDescent.DepthNonnegative.application_window}
\begin{proof}
Write $\lambda=\lambda_\tau$, $j=j_\tau$, $q=q_\tau$, $i=i_\tau$, $x^\pm=x^\pm_\tau$, $e^{\mathrm c}=e^{\mathrm
c}_\tau$, $\mathrm{top}=\mathrm{top}_\tau$ and $\lambda|_r=(\lambda_a)_{a<r}$. As $\tau$ has the shape,
$N=k+j+t+q$, $\lambda_b\le\lambda_a$ for $a\le b<N$, $\lambda_a=\lambda_{k+j}$ for $k+j\le a<k+j+t$, and
$\lambda_0+N\le h+\lambda_{N-1}+1$. By Lemma~\ref{is-onto-step-exp-bound-short-top-window-steps} with
$I=i$, $(\nu_\tau,\mu)$ has exponent at most
\[
 E=e^{\mathrm c}+(j+q)\lfloor\max(0,h-1)/(l-1)\rfloor+B(i)
\]
for every $\mu$ in question, where $B(i)$ is the largest of the numbers
$v_l\big(D_n(t,\lambda_0-\lambda_{k+j},i',\lambda^\vee)\big)$ for $0\le i'\le i$, with
$\lambda^\vee=(\lambda_0-\lambda_{k+j-1-b})_{b<k+j}$. An exponent bound persists for larger exponents, and
$E$ is an integer, so it suffices to show that $E\le C'_\tau$.

By the shape, $\mathrm{top}<x^-\le x^+$ and $x^+-\mathrm{top}\le h$, so $h\geq1$ and
$\lfloor\max(0,h-1)/(l-1)\rfloor\le(h-1)/(l-1)$. Put $w=\lambda_{k+j}$, $M=\lambda_0$ and $\beta=w+t-1$. For
$a<k+j$, $\lambda_a+k+j+t-1-a=U_\tau(a)-q$, and $\mathrm{top}=U_\tau(k+j)=\beta+q$. So the shape gives the
hypotheses of Lemma~\ref{is-onto-step-exp-bound-short-top-window-labels} for $k$, $j$, $t$, $w$, $M$,
$\lambda|_{k+j}$ and the integers $x^--q$ and $x^+-q$: $w\le\lambda_a\le M$ for $a<k+j$, as $t\geq1$;
$x^--q\le U_\tau(a)-q\le x^+-q$ for $a<k$; $U_\tau(a)-q<x^--q$ for $k\le a<k+j$; and
$\beta<x^--q\le x^+-q$, as $\mathrm{top}<x^-$. Moreover $(x^+-q)-\beta=x^+-\mathrm{top}\le h<l^2$, and
$\max(0,M-w)=\lambda_0-\lambda_{k+j}$. Let $i'\le i$ attain $B(i)$. Part (3) of that lemma, at $i'$ with
$I=i$, and Lemma~\ref{is-onto-step-exp-bound-short-top-window-label-sum}(2) give
\[
 B(i)\le j\,\frac{h-1}{l-1}+L(i;x^--q,x^+-q,\beta)=j\,\frac{h-1}{l-1}+L(i;x^-,x^+,\mathrm{top}),
\]
with $L$ for $l$, $k$, $t$ and $h$. By Definition~\ref{is-onto-step-exp-bound-short-top-data}(2),
$L(i;x^-,x^+,\mathrm{top})=k\big((x^+-x^-)/l+\delta_\tau\big)+\gamma_\tau+H_\tau$. Hence
\[
 E\le e^{\mathrm c}+(q+2j)\frac{h-1}{l-1}+k\Big(\frac{x^+-x^-}l+\delta_\tau\Big)+\gamma_\tau+H_\tau=C'_\tau.
\]
% Lean: application_window (ApplicationWindow) = expLeMono with applicationSteps hD and Nat.le_floor: hb
% (the bounds w <= lam a <= M), eU and e3 (U_tau(a) = U k j t (res (k + j) lam) a + q and top = maxB + q),
% hR (Rows k j t w M (res (k + j) lam) (xm - q) (xp - q)), hxp, hh1 (1 <= h), hcost (blockCost attained by
% Finset.exists_mem_eq_sup and bounded by second_valuation_up_to_holds), hP (Nat.cast_div_le) and
% labelRHSPlus_shift_invariant, then cPrime unfolded and linarith. Size before this step was written out:
% about 3,800 lines of Lean in 190 theorems.
\end{proof}

\begin{lemma}\label{is-onto-step-exp-bound-short-top-base}
Let $O$ be a domain in which every integer not divisible by a prime $l$ is a unit, $r\geq0$, $e\geq0$,
and $\nu\in\mathbf N^r$ nonincreasing.
\begin{enumerate}
\item If $ml^e\nabla_{\mathbf Z}(\nu)\subset\Delta_{\mathbf Z}(\nu)$ for some integer $m$ not divisible by
  $l$, then $l^e\nabla_O(\nu)\subset\Delta_O(\nu)$.
\item If $\mu\in\mathbf N^r$ and $(\nu,\mu)$ has exponent at most $e$ at $l$
  (Definition~\ref{is-onto-step-exp-bound-short-branching}(1)), then $l^ex_\mu\in\Delta_O(\nu)$ for every
  $x\in\nabla_O(\nu)$, where $x_\mu$ is the component of $x$ of row degree $\mu$.
\end{enumerate}
\end{lemma}
\leanref{Formalism/PieriBlockIdeals/PieriColumnMinors/ProdInitialMinorsMemDetSpan.lean}{Formalism.CohomologicalDescent.LA.la_incl_R_of_Z}
\leanref{Formalism/PieriBlockIdeals/PieriColumnMinors/DepthNonnegative.lean}{Formalism.CohomologicalDescent.DepthNonnegative.pow_smul_weighted_component_mem_weyl}
\begin{proof}
By Lemma~\ref{pd-kappa-defects-restriction-standard}(1), $\nabla_O(\nu)$ is the $O$-span of the
bideterminants of shape $\nu$, which are the images of the bideterminants over $\mathbf Z$; these lie in
$\nabla_{\mathbf Z}(\nu)$ and have a row degree, as each minor on rows $I(0),\dots,I(s-1)$ has degree $1$ in
each of these rows. The map $\mathbf Z[X]\to O[X]$ carries $\Delta_{\mathbf Z}(\nu)$ into $\Delta_O(\nu)$,
as both are generated from $h_\nu$ by the same operators. As the claims are $O$-linear in $x$, it suffices
to take for $x$ the image of a bideterminant $b$ of shape $\nu$.

(1) Here $ml^eb\in\Delta_{\mathbf Z}(\nu)$, so $ml^ex\in\Delta_O(\nu)$, and $m$ is a unit of $O$.

(2) If $b$ has row degree $\mu$, then $x_\mu=x$, and by assumption $ml^eb\in\Delta_{\mathbf Z}(\nu)$ for some
$m$ not divisible by $l$; so $l^ex_\mu\in\Delta_O(\nu)$ as in (1). Otherwise $x_\mu=0$.
% Lean: la_incl_R_of_Z and pow_smul_weighted_component_mem_weyl, both through standardMonomialSpanDomain,
% bidetMem, ldbasis_bidet_rowHom, la_map_weyl, ldlat_map_bidet and the inverse of the unit m.
\end{proof}

\begin{lemma}\label{is-onto-step-exp-bound-short-top-forms}
Let $O$ be a commutative ring, $A$ an $O$-algebra, $G$ a group with subgroups $\Gamma$ and $K$, $j\colon
K\to A$ a homomorphism of monoids, $V$ an $A$-module, $T$ a commutative ring acting $O$-linearly on the
forms $F(V)$ (Definition~\ref{is-onto-step-exp-bound-hida}(1)) and $S\subset T$ a submonoid. Suppose that
$T$ preserves the forms with values in every $A$-submodule of $V$.
\begin{enumerate}
\item Let $\varpi\in O$, $n\geq0$, and $F\subset F'$ $A$-submodules of $V$ with $\varpi^nV\subset F$ and
  $F'/F$ simple. Then $F'/F$ is a quotient of an $A$-submodule of $V/\varpi V$.
\item Let $F\subset F'$ be $A$-submodules of $V$, $W$ an $A$-module on whose forms $T$ acts, and
  $\psi\colon F'\to W$ an $A$-linear map with kernel in $F$, such that $T$ commutes with the maps of forms
  induced by $\psi$ and $S^{-1}F(W)=0$. Then for every form $f$ with values in $F'$ there is $\sigma\in S$
  with $\sigma f$ valued in $F$.
\item Suppose that $V$ is a Noetherian $O$-module, and let $D_0$ be an $O$-submodule with $V/D_0$ Artinian,
  and $D\subset E$ $A$-submodules with $D_0\subset D$. Suppose that for all $A$-submodules $F\subset F'$
  with $D\subset F$, $F'\subset E$ and $F'/F$ simple, and every form $f$ with values in $F'$, some $\sigma f$
  with $\sigma\in S$ has values in $F$. Then for every form $f$ with values in $E$, some $\sigma f$ with
  $\sigma\in S$ has values in $D$.
\end{enumerate}
\end{lemma}
\leanref{Formalism/PieriBlockIdeals/BlockIdealBound5Derived.lean}{Formalism.CohomologicalDescent.DepthNonnegative.SimpleSubquotOfPower}
\leanref{Formalism/PieriBlockIdeals/PieriColumnMinors/DepthNonnegative.lean}{Formalism.CohomologicalDescent.DepthNonnegative.simple_subquotient_of_power_holds}
\leanref{Formalism/PieriBlockIdeals/BlockIdealBound5Derived.lean}{Formalism.CohomologicalDescent.DepthNonnegative.FormsStep}
\leanref{Formalism/PieriBlockIdeals/PieriColumnMinors/DepthNonnegative.lean}{Formalism.CohomologicalDescent.DepthNonnegative.forms_induction_holds}
\leanref{Formalism/PieriBlockIdeals/BlockIdealBound5Derived.lean}{Formalism.CohomologicalDescent.DepthNonnegative.FormsDevissage}
\leanref{Formalism/PieriBlockIdeals/PieriColumnMinors/DepthNonnegative.lean}{Formalism.CohomologicalDescent.DepthNonnegative.forms_devissage_holds}
\begin{proof}
(1) Put $F_a=F+(F'\cap\varpi^aV)$. Then $F_0=F'$ and $F_n=F$, and as $F'/F$ is simple, each $F_a$ is $F$
or $F'$. Let $a$ be the largest index with $F_a=F'$; then $a<n$ and $F_{a+1}=F$. Let
$E_0=\{v\in V:\varpi^av\in F'\}$ and $\psi_0\colon E_0\to F'/F$, $v\mapsto\varpi^av$ modulo $F$. It is onto,
as $F'=F+(F'\cap\varpi^aV)$. If $v\in E_0\cap\varpi V$, say $v=\varpi v'$, then
$\varpi^av=\varpi^{a+1}v'\in F'\cap\varpi^{a+1}V\subset F_{a+1}=F$, so $\psi_0(v)=0$. So $\psi_0$ factors
through the image of $E_0$ in $V/\varpi V$.

(2) The map $g\mapsto\psi(f(g))$ is a form with values in $W$, as $\psi$ is $A$-linear. As $S^{-1}F(W)=0$,
some $\sigma\in S$ kills it. As $\sigma f$ has values in $F'$ and $T$ commutes with $\psi$,
$\psi\big((\sigma f)(g)\big)=\big(\sigma(\psi\circ f)\big)(g)=0$ for every $g$, so $\sigma f$ has values in
$F$.

(3) Among the $A$-submodules $F$ with $D\subset F\subset E$ such that some $\sigma f$, $\sigma\in S$, has
values in $F$ (for instance $E$), choose one whose image in $V/D_0$ is minimal, which is possible as $V/D_0$
is Artinian. If $F\subset D$ we are done. Otherwise $D\subsetneq F$; as $V$ is Noetherian over $O$, hence
over $A$, there is a maximal $A$-submodule $F_0$ with $D\subset F_0\subsetneq F$, and $F/F_0$ is simple. By
hypothesis some $\sigma'\sigma f$ has values in $F_0$, and as $D_0\subset D\subset F_0$, the image of $F_0$ in
$V/D_0$ is strictly smaller than that of $F$, a contradiction.
% Lean: simple_subquotient_of_power_holds (SimpleSubquotOfPower; covBy_iff_quot_is_simple, Nat.find for a,
% mulOntoScaledSubmodule); forms_induction_holds (FormsStep, with HeckeNat for the commutation and
% LocalizedModule.mk_eq); forms_devissage_holds (FormsDevissage; IsArtinian.set_has_minimal,
% set_has_maximal_iff_noetherian, isNoetherian_of_tower).
\end{proof}

\begin{lemma}\label{is-onto-step-exp-bound-short-top-core}
Let $l$ be a prime, $N\geq0$ and $k\geq0$ integers, $\mathcal T$ a finite set, $T_2\subset\mathcal T$,
and for $\tau\in\mathcal T$ let $\lambda_\tau$ be a sequence of natural numbers whose first $N$ terms are
nonincreasing, $\nu_\tau=(\lambda_{\tau,a})_{a<N}$, and $j_\tau,i_\tau,e_\tau,c'_\tau\geq0$ integers. Let $X$
consist of the data of a Hida setting of rank $N$ and weights $\lambda$
(Definition~\ref{is-onto-step-exp-bound-hida}(1)), except that its ring $O$ is only required to be a
Noetherian domain, with constituent data at $l$ (Definition~\ref{is-onto-step-exp-bound-hida}(4), with $N$
in place of $N_r$). Suppose that:
\begin{enumerate}
\item $\nabla_O(\lambda)/\Delta_O(\lambda)$ is an Artinian $O$-module, and $l^{e_\tau}\nabla_O(\nu_\tau)\subset
  \Delta_O(\nu_\tau)$ for every $\tau$;
\item for every $\tau\in T_2$ and every relevant $\mu$,
  $\sum_{a<k+j_\tau}(\lambda_{\tau,a}-\omega_{\mu,\tau,a})\le i_\tau$;
\item for every $\tau\in T_2$, every $w\in\mathbf N^N$ with $\sum_{a<k+j_\tau}(\lambda_{\tau,a}-w_a)\le i_\tau$
  and every $x\in\nabla_O(\nu_\tau)$, $l^{c'_\tau}x_w\in\Delta_O(\nu_\tau)$, where $x_w$ is the component of
  $x$ of row degree $w$.
\end{enumerate}
Put $c_\tau=\min(e_\tau,c'_\tau)$ for $\tau\in T_2$ and $c_\tau=e_\tau$ otherwise. Then $X$ has the inclusion
with exponent $\sum_\tau c_\tau$ at $l$ (Definition~\ref{is-onto-step-exp-bound-hida}(2)).
\end{lemma}
\leanref{Formalism/PieriBlockIdeals/BlockIdealBound5Derived.lean}{Formalism.CohomologicalDescent.DepthNonnegative.TheoremWCore}
\leanref{Formalism/PieriBlockIdeals/PieriColumnMinors/DepthNonnegative.lean}{Formalism.CohomologicalDescent.DepthNonnegative.theorem_w_central_of_simple_quotients}
\leanref{Formalism/PieriBlockIdeals/BlockIdealBound5Derived.lean}{Formalism.CohomologicalDescent.DepthNonnegative.cTau}
\begin{proof}
Write $V=\nabla_O(\lambda)$, $N_\tau=\nabla_O(\nu_\tau)$, $B_\tau=\Delta_O(\nu_\tau)$ and $D_s=\Delta^s_O(\lambda)$
for $s\subset\mathcal T$; so $D_\emptyset=\Delta_O(\lambda)$, $D_{\mathcal T}=V$, $D_s\subset D_{s'}$ for
$s\subset s'$, and $A$ preserves every $D_s$. For $a\notin s$, an $O$-linear map $p$ of $N_a$ with
$l^cp(N_a)\subset B_a$ satisfies $l^cp_a(D_{s\cup\{a\}})\subset D_s$, where $p_a$ applies $p$ to the factor
$a$ of a pure tensor; and by (1), $l^{\sum_\tau e_\tau}V\subset D_\emptyset$. Each $N_\tau$ lies in the
polynomials of degree $\sum_a\nu_{\tau,a}$, a finitely generated $O$-module, so $V$ is a finitely generated
module over the Noetherian ring $O$, hence Noetherian.

Step. Let $a\notin s$ and let $f$ be a form with values in $D_{s\cup\{a\}}$. Then $\sigma l^{c_a}f$ has
values in $D_s$ for some $\sigma\in S$. If $c_a=e_a$, take $\sigma=1$ and $p$ the identity. Otherwise
$a\in T_2$ and $c_a=c'_a<e_a$. Let $E=D_s+l^{c'_a}D_{s\cup\{a\}}$, an $A$-submodule containing $D_s$, in
which $l^{c'_a}f$ takes its values. We apply Lemma~\ref{is-onto-step-exp-bound-short-top-forms}(3) with
$D_0=\Delta_O(\lambda)$, $D=D_s$ and this $E$. Let $D_s\subset F\subset F'\subset E$ with $F'/F$ simple. As
$l\mid\varpi^{n_0}$, $\varpi^{n_0\sum e_\tau}V\subset D_\emptyset\subset F$, so by
Lemma~\ref{is-onto-step-exp-bound-short-top-forms}(1) and constituent datum (c), $F'/F\cong L_\mu$ for some
$\mu$, and $L_\mu$ is a subquotient of $V/\varpi V$. Let $\psi\colon F'\to L_\mu$ be the quotient map
followed by this isomorphism. Then $\sum_{a'<k+j_a}(\lambda_{a,a'}-\omega_{\mu,a,a'})>i_a$. Indeed,
otherwise take $z\in L_\mu$ with $e_a(\omega_{\mu,a})z\ne0$, and $x\in F'$ with $\psi(x)=z$; write
$x=d+l^{c'_a}n$ with $d\in D_s$ and $n\in D_{s\cup\{a\}}$. Then $e_a(\omega_{\mu,a})d\in D_s$, and
$e_a(\omega_{\mu,a})l^{c'_a}n=l^{c'_a}p_a(n)\in D_s$ for $p$ the projection to the row degree
$\omega_{\mu,a}$, by (3). So $e_a(\omega_{\mu,a})x\in F$, and $e_a(\omega_{\mu,a})z=0$, a contradiction. By
(2), $\mu$ is not relevant, so $S^{-1}F(L_\mu)=0$, and
Lemma~\ref{is-onto-step-exp-bound-short-top-forms}(2) moves every form with values in $F'$ into $F$ by an
element of $S$. Lemma~\ref{is-onto-step-exp-bound-short-top-forms}(3) now gives the step.

Applying the step to the elements of $\mathcal T$ one at a time, from $s=\mathcal T$ down to $\emptyset$,
every form $f$ (with values in $D_{\mathcal T}=V$) has some $\sigma l^{\sum_\tau c_\tau}f$, $\sigma\in S$, with
values in $D_\emptyset=\Delta_O(\lambda)$. For $y=f/\sigma_0$ in $S^{-1}F(V)$, then
$l^{\sum c_\tau}y=\big(\sigma l^{\sum c_\tau}f\big)/(\sigma\sigma_0)$.
% Lean: theorem_w_central_of_simple_quotients (TheoremWCore, stated for an abstract depth function and
% tensor factors): DA s = tensSub (chainB B s), chain_ends_holds, chain_monotone_holds, hpow
% (tensor_reduction_holds), hstep (tensor_weight_kill_holds, forms_devissage_holds, simple_subquotient_of_power,
% hclass, subquotient_iso, hdeep with hhw and hIN2, hvan with hIN1, forms_induction_holds), hall by
% Finset.induction_on, and LocalizedModule.induction_on.
\end{proof}

\begin{lemma}\label{is-onto-step-exp-bound-short-top}
Suppose that the block exponent bound holds in last graded form
(Definition~\ref{is-onto-step-exp-bound-short-branching}(5)). Let $l$, $N$, $k$, $t$, $h$, $\mathcal T$,
$T_2$, $\lambda$, $j$, $q$, $e^{\mathrm c}$, $e$, $i$, $x^-$ and $x^+$ be data as in
Definition~\ref{is-onto-step-exp-bound-short-top-data}, and let $X$ consist of the data of a Hida setting
of rank $N$ and weights $\lambda$, except that its ring $O$ is only required to be a Noetherian domain,
with the top inputs for these data. Then $X$ has the inclusion with exponent $J'$ at $l$
(Definition~\ref{is-onto-step-exp-bound-hida}(2)), and its least exponent at $l$ is at most $J'$.
\end{lemma}
\leanref{Formalism/PieriBlockIdeals/BlockIdealBound5Derived.lean}{Formalism.CohomologicalDescent.DepthNonnegative.TheoremWTop}
\leanref{Formalism/PieriBlockIdeals/PieriColumnMinors/DepthNonnegative.lean}{Formalism.CohomologicalDescent.DepthNonnegative.weight_theorem_top}
\begin{proof}
By the top inputs, $l$ is a prime, every integer not divisible by $l$ is a unit of $O$, $O/lO$ is finite,
$h<l^2$ and $t\geq1$. For every $\tau$, $ml^{e_\tau}\nabla_{\mathbf Z}(\nu_\tau)\subset\Delta_{\mathbf Z}(\nu_\tau)$
for some $m$ not divisible by $l$, so $l^{e_\tau}\nabla_O(\nu_\tau)\subset\Delta_O(\nu_\tau)$ by
Lemma~\ref{is-onto-step-exp-bound-short-top-base}(1). Hence $l^{\sum_\tau e_\tau}$ kills
$\nabla_O(\lambda)/\Delta_O(\lambda)$, which is finitely generated over $O$ (as in the proof of
Lemma~\ref{is-onto-step-exp-bound-short-top-core}). As $O$ is Noetherian and $O/lO$ is finite,
$O/l^{\sum e_\tau}O$ is finite, so the quotient is finite, hence Artinian.

Let $\tau\in T_2$. Then $\tau$ has the shape and $\big((\lambda_{\tau,a})_{a<k},\mu'\big)$ has exponent at
most $e^{\mathrm c}_\tau$ for every $\mu'$, so by Lemma~\ref{is-onto-step-exp-bound-short-top-window},
$(\nu_\tau,w)$ has exponent at most $\lfloor C'_\tau\rfloor$ at $l$ for every $w\in\mathbf N^N$ with
$\sum_{a<k+j_\tau}(\lambda_{\tau,a}-w_a)\le i_\tau$. By
Lemma~\ref{is-onto-step-exp-bound-short-top-base}(2), $l^{\lfloor C'_\tau\rfloor}x_w\in\Delta_O(\nu_\tau)$ for
every $x\in\nabla_O(\nu_\tau)$.

The constituent data of the top inputs satisfy condition (2) of
Lemma~\ref{is-onto-step-exp-bound-short-top-core}. That lemma, with $e_\tau$ and
$c'_\tau=\lfloor C'_\tau\rfloor$, gives the inclusion with exponent
$\sum_{\tau\in T_2}\min(e_\tau,\lfloor C'_\tau\rfloor)+\sum_{\tau\notin T_2}e_\tau=J'$. So the least exponent
is at most $J'$.
% Lean: weight_theorem_top = theorem_w_top_of_simple_quotients simple_subquotient_of_power_holds
% subquotient_iso_holds: hexO (la_incl_R_of_Z), hArt (finite_quotient_of_pow_smul_mem,
% tensor_reduction_holds, isArtinian_of_finite), he' (factor_action_on_tensor_product_holds), hIN2
% (application_window and pow_smul_weighted_component_mem_weyl), theorem_w_central_of_simple_quotients and
% Nat.sInf_le.
\end{proof}

\begin{lemma}\label{is-onto-step-exp-bound-short-top-bound}
Suppose that the block exponent bound holds in last graded form
(Definition~\ref{is-onto-step-exp-bound-short-branching}(5)). Let $l$, $N$, $k$, $t$, $h$, $\mathcal T$,
$T_2$, $\lambda$, $j$, $q$, $e^{\mathrm c}$, $e$, $i$, $x^-$ and $x^+$ be data as in
Definition~\ref{is-onto-step-exp-bound-short-top-data}, and let $X$ consist of the data of a Hida setting
of rank $N$ and weights $\lambda$, except that its ring $O$ is only required to be a Noetherian domain,
with the top inputs for these data. Then $X$ has the inclusion with exponent $J^{\mathrm s}$ at $l$, and
its least exponent at $l$ is at most $J^{\mathrm s}$.
\end{lemma}
\leanref{Formalism/PieriBlockIdeals/LatticeInclusionRealSettingAlt.lean}{Formalism.CohomologicalDescent.ShortWeightBounds.TheoremWTopShort6}
\leanref{Formalism/PieriBlockIdeals/PieriColumnMinors/WeightBounds.lean}{Formalism.CohomologicalDescent.ShortWeightBounds.theorem_w_top_short_bound_holds}
\leanref{Formalism/PieriBlockIdeals/LatticeInclusionRealSettingAlt.lean}{Formalism.CohomologicalDescent.ShortWeightBounds.CPrimeLeShort6}
\leanref{Formalism/PieriBlockIdeals/PieriColumnMinors/WeightBounds.lean}{Formalism.CohomologicalDescent.ShortWeightBounds.c_prime_le_short_bound_holds}
\leanref{Formalism/PieriBlockIdeals/PieriColumnMinors/DepthNonnegative.lean}{Formalism.CohomologicalDescent.DepthNonnegative.top_inclusion_monotone_holds}
\begin{proof}
Let $\tau\in T_2$. Then $C'_\tau\le C^{\mathrm s}_\tau$: as $l\geq2$,
$(h-1)/(l-1)\le(h+N)/(l-1)$; the next three terms agree; and
$H_\tau\le(i_\tau/l)(1+\log^+(h/l))$, as $H_\tau=0$ if $h<l$, and $\log(h/l)\le\log^+(h/l)$ otherwise. So
$\min(e_\tau,\lfloor C'_\tau\rfloor)\le\min(e_\tau,\lfloor C^{\mathrm s}_\tau\rfloor)$, and
$J'\le J^{\mathrm s}$. By Lemma~\ref{is-onto-step-exp-bound-short-top}, $X$ has the inclusion with
exponent $J'$. It then has the inclusion with exponent $J^{\mathrm s}$: if $l^{J'}y=f/s$ with $s\in S$ and
$f$ a form with values in $\Delta_O(\lambda)$, then $l^{J^{\mathrm s}}y=l^{J^{\mathrm s}-J'}f/s$, and
$l^{J^{\mathrm s}-J'}f$ is again such a form. So the least exponent is at most $J^{\mathrm s}$.
% Lean: theorem_w_top_short_bound_holds (TheoremWTopShort6): top_inclusion_monotone_holds (TopInclMono)
% applied to the first part of weight_theorem_top, the sums compared through Nat.floor_le_floor and
% c_prime_le_short_bound_holds (CPrimeLeShort6, with harmonic_positive_log_holds), and Nat.sInf_le.
\end{proof}

\begin{lemma}\label{is-onto-step-exp-bound-short-shape}
Let $P$ be a parameter record that satisfies the hypotheses and is large, $W$ weights for $P$ that
satisfy the weight conditions (Definition~\ref{is-onto-step-exp-bound-weights}), and $\tau\in T_W$. Then
$\tau$ has the shape (Definition~\ref{is-onto-step-exp-bound-short-top-data}(1)) for $N=N_r$, $k-1$ in
place of $k$, $t$, and $h_0$ in place of $h$, with the sequence $\lambda_\tau$ and the integers $j_\tau$,
$q_\tau$, $x^-_\tau$ and $x^+_\tau$ of $W$.
\end{lemma}
\leanref{Formalism/PieriBlockIdeals/LatticeInclusionRealSettingAlt.lean}{Formalism.CohomologicalDescent.ShortWeightBounds.ShapeWReal}
\leanref{Formalism/PieriBlockIdeals/PieriColumnMinors/WeightBounds.lean}{Formalism.CohomologicalDescent.ShortWeightBounds.shape_w_real_holds}
\begin{proof}
Let $p_\tau=k-1+j_\tau$. By (W1), $p_\tau+t+q_\tau=N_r$ and the first $N_r$ terms of $\lambda_\tau$ are
nonincreasing; (W2) gives $\lambda_{\tau,a}=\lambda_{\tau,p_\tau}$ for $p_\tau\le a<p_\tau+t$; and (W3)
gives $x^-_\tau\le U_\tau(a)\le x^+_\tau$ for $a<k-1$, and $x^-_\tau\le x^+_\tau$. By (W5),
$x^+_\tau\le h$, so $x^+_\tau\le h_0$, as $x^+_\tau$ is an integer. As in the proof of
Lemma~\ref{is-onto-step-exp-bound-beyond}, $U_\tau(a)<x^-_\tau$ for $k-1\le a<N_r$; this gives the
condition for $k-1\le a<p_\tau$, and $\mathrm{top}_\tau=U_\tau(p_\tau)<x^-_\tau$, as $p_\tau<N_r$. Hence
$U_\tau(0)\le x^+_\tau\le h_0$, that is, $\lambda_{\tau,0}+N_r-1\le h_0$, so
$\lambda_{\tau,0}+N_r\le h_0+\lambda_{\tau,N_r-1}+1$. Finally $\mathrm{top}_\tau\geq0$, as
$\lambda_{\tau,p_\tau}\geq0$ and $p_\tau\le N_r-1$; so $x^+_\tau-\mathrm{top}_\tau\le h_0$.
% Lean: shape_w_real_holds (ShapeWReal): rank_fits_block_and_t_ge_six, upper_translation_le_floor_height
% (le_floor_h), hW.what2, hW.anti, hW.block, hW.cloud, (hW.window tau).2.1,
% weight_le_upper_translation_of_thick at a = 0, weight_lt_lower_translation_of_thick
% (abs_weight_sub_near_translation_le, two_eta_lt_b) and weight_nonneg'.
\end{proof}

\begin{lemma}\label{is-onto-step-exp-bound-short-inputs}
Let $P$ be a parameter record that satisfies the hypotheses and is large, $W$ weights for $P$ that
satisfy the weight conditions, indexed by $\mathcal T$, $l>l_J$ a prime, and $X$ a Hida setting of rank
$N_r$ and the weights $\lambda$ of $W$, with $l$ not a unit in its ring $O$, that has constituent data at
$l$ with $\sum_{a<k-1+j_\tau}(\lambda_{\tau,a}-\omega_{\mu,\tau,a})\le i_\tau(l)$ for every relevant $\mu$
and every $\tau\in T_W$ (Definitions~\ref{is-onto-step-exp-bound-weights},
\ref{is-onto-step-exp-bound-recursion} and~\ref{is-onto-step-exp-bound-hida}). Then $X$ has the top
inputs (Definition~\ref{is-onto-step-exp-bound-short-top-data}(3)), with the same constituent data, for
$l$, $N_r$, $k-1$, $t$, $h_0$, $\mathcal T$ and $T_W$, and for every $\tau$ the sequence $\lambda_\tau$ and
the integers $j_\tau$, $q_\tau$, $e^{\mathrm c}_\tau(l)$, $e_\tau(l)$, $i_\tau(l)$, $x^-_\tau$ and
$x^+_\tau$.
\end{lemma}
\leanref{Formalism/PieriBlockIdeals/LatticeInclusionRealSettingAlt.lean}{Formalism.CohomologicalDescent.ShortWeightBounds.TopInputsReal}
\leanref{Formalism/PieriBlockIdeals/PieriColumnMinors/WeightBounds.lean}{Formalism.CohomologicalDescent.ShortWeightBounds.top_inputs_of_shape_w}
\leanref{Formalism/PieriBlockIdeals/PieriColumnMinors/RealLocalDeltaNabla.lean}{Formalism.CohomologicalDescent.RealLocalDeltaNabla.units_of_local_holds}
\leanref{Formalism/PieriBlockIdeals/PieriColumnMinors/DepthNonnegative.lean}{Formalism.CohomologicalDescent.DepthNonnegative.exact_exponents_fit_holds}
\begin{proof}
Let $m$ be an integer not divisible by $l$. Then $la+mb=1$ for some integers $a$ and $b$. As $O$ is local
and $l$ is not a unit of $O$, $la$ is not a unit, so $mb$ is, and so is $m$. By
Lemma~\ref{is-onto-step-exp-bound-total-numbers}, $h+N_r<l^2$, so $h_0\le h<l^2$; and $t\geq1$ by the
hypotheses. Every $\tau\in T_W$ has the shape by Lemma~\ref{is-onto-step-exp-bound-short-shape}. By (W1),
the first $N_r$ terms of $\lambda_\tau$ are nonincreasing, and so are its first $k-1$ terms, as
$k-1\le N_r$. So Lemma~\ref{is-onto-step-exp-cut-attained}, for $e_\tau(l)=e_l(N_r,\lambda_\tau)$ and
$e^{\mathrm c}_\tau(l)=e_l(k-1,\lambda_\tau)$ (Definition~\ref{is-onto-step-exp-bound-recursion}(1)), gives
integers $m$ and $m'$ prime to $l$ with $ml^{e_\tau(l)}\nabla_{\mathbf Z}(\nu_\tau)\subset\Delta_{\mathbf
Z}(\nu_\tau)$ and $m'l^{e^{\mathrm c}_\tau(l)}\nabla_{\mathbf Z}(\nu'_\tau)\subset\Delta_{\mathbf Z}(\nu'_\tau)$,
where $\nu'_\tau=(\lambda_{\tau,a})_{a<k-1}$; the second inclusion holds in particular for the elements of
row degree $\mu$, for every $\mu$. The remaining conditions are the constituent data and the bound on the
depths.
% Lean: top_inputs_of_shape_w (TopInputsReal): RealLocalDeltaNabla.units_of_local_holds
% (Nat.gcd_eq_gcd_ab, IsLocalRing.isUnit_or_isUnit_of_isUnit_add); hsq from the clause h + Nr < l^2 of
% paramsLargePrime_holds and floor_height_le_height; shape_w_real_holds; and exact_exponents_fit_holds
% (LA.exact_z_specialization and exponent_le_of_integral_inclusion_holds) with restriction_lambda_antitone
% and restriction_lambda_antitone_kc.
\end{proof}

\begin{lemma}\label{is-onto-step-exp-bound-short}
Let $P$ be a parameter record that satisfies the hypotheses and is large, $W$ weights for $P$ that
satisfy the weight conditions, $l>l_J$ a prime, and $X$ a Hida setting of rank $N_r$ and the weights
$\lambda$ of $W$, with $l$ not a unit in its ring $O$, that has the weight inputs for $W$ at $l$
(Definition~\ref{is-onto-step-exp-bound-hida}). Then $X$ has the inclusion with exponent
$J^{\mathrm s}_W(l)$ at $l$ (Definition~\ref{is-onto-step-exp-bound-recursion}), and its least exponent at
$l$ is at most $J^{\mathrm s}_W(l)$.
\end{lemma}
\leanref{Formalism/PieriBlockIdeals/StatementDown.lean}{Formalism.CohomologicalDescent.Down.InclWReal}
\leanref{Formalism/PieriBlockIdeals/PieriColumnMinors/CorollaryOfParamsLinkage.lean}{Formalism.CohomologicalDescent.Down.inclusion_w_real_holds}
\begin{proof}
Take the constituent data at $l$ of the weight inputs. By
Lemma~\ref{is-onto-step-exp-bound-short-inputs}, $X$ has the top inputs for $l$, $N_r$, $k-1$, $t$,
$h_0$, the index set and $T_W$ of $W$, and the numbers $\lambda_\tau$, $j_\tau$, $q_\tau$,
$e^{\mathrm c}_\tau(l)$, $e_\tau(l)$, $i_\tau(l)$, $x^-_\tau$ and $x^+_\tau$. By
Lemmas~\ref{is-onto-step-exp-bound-short-graded} and~\ref{is-onto-step-exp-bound-short-last}, the
block exponent bound holds in last graded form. So Lemma~\ref{is-onto-step-exp-bound-short-top-bound} shows that $X$ has the
inclusion with exponent $J^{\mathrm s}$ at $l$, and that its least exponent at $l$ is at most
$J^{\mathrm s}$. For these data, $\mathrm{top}_\tau=U_\tau(k-1+j_\tau)$, and $d_\tau$, $\delta_\tau$ and
$\gamma_\tau$ are $d_\tau(l)$, $\delta_\tau(l)$ and $\gamma_\tau(l)$ of
Definition~\ref{is-onto-step-exp-bound-recursion}(3); so $C^{\mathrm s}_\tau$ is the short bound
$C^{\mathrm s}_\tau(l)$, and $J^{\mathrm s}=J^{\mathrm s}_W(l)$.
% Lean: Down.inclusion_w_real_holds = ShortWeightBounds.inclusion_w_real_holds hG, and
% inclusion_w_real_holds = inclusion_w_of_shape_w shape_w_real_holds: the data of WInputs,
% top_inputs_of_shape_w, and theorem_w_top_short_bound_holds (last_graded_of_graded_alone hD) with
% hD = hG; the sum of cTau T2 ex (cFloorShort6 ...) is W.JW P l by definition. Size before this step was
% written out: about 7,300 lines of Lean in 540 theorems.
\end{proof}

\begin{lemma}\label{is-onto-step-exp-bound-coarse}
Let $P$, $W$, $l$ and $X$ be as in Lemma~\ref{is-onto-step-exp-bound-short}. Then $X$ has the inclusion
with exponent $J_W(l)$ at $l$, and its least exponent at $l$ is at most $J_W(l)$.
\end{lemma}
\leanref{Formalism/PieriBlockIdeals/HidaLatticeInclusionAtCoarseWeightBound.lean}{Formalism.CohomologicalDescent.Down7.HidaLatticeInclusionAtCoarseWeightBound}
\leanref{Formalism/PieriBlockIdeals/PieriColumnMinors/CorollaryOfParamsSevenLinkage.lean}{Formalism.CohomologicalDescent.Down7.inclusion_w_seven_real_holds}
\begin{proof}
For every $\tau$, $C^{\mathrm s}_\tau(l)\le C_\tau(l)$, as $q_\tau+2j_\tau\geq0$; so
$\min(e_\tau(l),\lfloor C^{\mathrm s}_\tau(l)\rfloor)\le\min(e_\tau(l),\lfloor C_\tau(l)\rfloor)$, and
$J^{\mathrm s}_W(l)\le J_W(l)$. If $X$ has the inclusion with exponent $J$ at $l$ and $J\le J'$, it has
the inclusion with exponent $J'$: for $y\in S^{-1}F(\nabla_O(\lambda))$, apply the inclusion with
exponent $J$ to $l^{J'-J}y$. Lemma~\ref{is-onto-step-exp-bound-short} gives the inclusion with exponent
$J^{\mathrm s}_W(l)$ and least exponent at most $J^{\mathrm s}_W(l)$, so both statements follow.
% Lean: Down7.inclusion_w_seven_real_holds: Down.inclusion_w_real_holds, JW_le_coarseWeightBound
% (cShort6_le_cCoarse) and Incl_mono (pow_smul_mem_mono').
\end{proof}

\begin{lemma}\label{is-onto-step-exp-bound-models}
Let $R$, $j$, $t$ and $w$ be as in Definition~\ref{is-onto-step-exp-bound-lift-exponent}. Let $P$ and
$W$ be a parameter record and weights for $P$ that satisfy the numerical conditions
(Definition~\ref{is-onto-step-exp-bound-weights}), with $p_w$ in the range of $W$. Let $W'$ be the
weights $W$ with some other Frobenius data, and let $(X_i)_{i\in I}$ be a family of Hida settings of rank
$N_r$ and the weights $\lambda$ of $W$ that covers the block data at $(j,t)$ at $w$, such that $p_w$ is
not a unit in the ring of any $X_i$ and every $X_i$ is linked to $W'$ at $p_w$
(Definition~\ref{is-onto-step-exp-bound-hida}). Then $p_w^{J_W(p_w)}C_{j,t}\subset(D_{j,t})_{(w)}$.
\end{lemma}
\leanref{Development/AbelianVarietyTheory/Proofs/DivisorProofs/PrimePowerLatticeInclusionOfLocalModelsStatement.lean}{GlobalTheory.T.primePowerLatticeInclusion_of_localModels}
\leanref{Development/AbelianVarietyTheory/Proofs/DivisorProofs/PrimePowerLatticeInclusionOfLocalModels.lean}{GlobalTheory.ExpTie.lIncl_of_models}
\begin{proof}
By Lemma~\ref{is-onto-step-exp-bound-frobenius}, $W'$ satisfies the weight conditions and
$J_{W'}(p_w)=J_W(p_w)$; the sequences $\lambda$ of $W'$ are those of $W$. Let $i\in I$. As $p_w$ is in the
range of $W$, $p_w>l_J$, so Lemma~\ref{is-onto-step-exp-bound-link}, for $P$ and $W'$, shows that $X_i$
has the weight inputs for $W'$ at $p_w$. Lemma~\ref{is-onto-step-exp-bound-coarse}, for $P$ and $W'$,
then shows that $X_i$ has the inclusion with exponent $J_{W'}(p_w)=J_W(p_w)$ at $p_w$. As this holds for
every $i$ and the family covers the block data at $(j,t)$ at $w$,
$p_w^{J_W(p_w)}C_{j,t}\subset(D_{j,t})_{(w)}$.
% Lean: lIncl_of_models: expo_frob_indep, weight_inputs_of_linkage_holds and
% Top.inclusion_w_seven_real_holds (= Down7.inclusion_w_seven_real_holds) for reFr W fr cj, then the
% covering hC. Of the numerical conditions only the hypotheses, largeness and the weight conditions
% are used.
\end{proof}

\begin{lemma}\label{is-onto-step-exp-bound-place}
Assume the exact cut bound of Lemma~\ref{is-onto-step-exp-cut}. Let $R$, $j$, $t$ and $w$ be as in
Definition~\ref{is-onto-step-exp-bound-lift-exponent}, and let $P$ and $W$ be a parameter record and
weights for $P$ that satisfy the numerical conditions, with recursive exponent $J'$
(Definitions~\ref{is-onto-step-exp-bound-weights} and~\ref{is-onto-step-exp-bound-recursion}). Suppose
that:
\begin{enumerate}
\item if $p_w$ is not in the range of $W$, then $p_w^{J_A(p_w)}C_{j,t}\subset(D_{j,t})_{(w)}$;
\item if $p_w$ is in the range of $W$, then there are Frobenius data and a family of Hida settings of rank
  $N_r$ and the weights $\lambda$ of $W$ that covers the block data at $(j,t)$ at $w$, such that $p_w$
  is not a unit in any of their rings and each of them is linked at $p_w$ to the weights $W$ with
  these Frobenius data (Definition~\ref{is-onto-step-exp-bound-hida}).
\end{enumerate}
Then $J_0(t,w,j)\le J'(p_w)$.
\end{lemma}
\leanref{Development/AbelianVarietyTheory/Proofs/DivisorProofs/ExpansionOfNumericsModel.lean}{GlobalTheory.ExpTie.Asm.jexpU_le_jrec7}
\begin{proof}
If $p_w$ is in the range of $W$, Lemma~\ref{is-onto-step-exp-bound-models} and (2) give
$p_w^{J_W(p_w)}C_{j,t}\subset(D_{j,t})_{(w)}$, Lemma~\ref{is-onto-step-exp-bound-lift} gives
$J_0(t,w,j)\le J_W(p_w)$, and $J_W(p_w)=J'(p_w)$ by Lemma~\ref{is-onto-step-exp-bound-compare}. If not,
(1) and Lemma~\ref{is-onto-step-exp-bound-lift} give $J_0(t,w,j)\le J_A(p_w)$, and
$J_A(p_w)\le J'(p_w)$ by Lemma~\ref{is-onto-step-exp-bound-compare}.
% Lean: jexpU_le_jrec7: exponent_of_record_vs_formulas (Lemma is-onto-step-exp-bound-compare),
% T.primePowerLatticeInclusion_of_localModels (Lemma is-onto-step-exp-bound-models) and
% T.exponent_datum_le_of_primePowerLatticeInclusion (Lemma is-onto-step-exp-bound-lift). The Lean
% states it for any block data B, level l and subspace Z; the book states it for the cut lattices of a
% realization, the only case it uses.
\end{proof}

\begin{lemma}\label{is-onto-step-exp-bound}
Assume the exact cut bound of Lemma~\ref{is-onto-step-exp-cut}. Let $p$ be a point of a family
(Definition~\ref{is-point}), at a member with rank $m$, scale $x$ and exponent $\varepsilon$ whose host
has the set $T_2$. Let $O$ be a record object, with any constants, for the indefinite-side datum and the
group of the record objects recorded at $p$ (Definition~\ref{is-record-object}), whose numbers $m$, $x$,
$\varepsilon$ and $T_2$ are $m$, $x$, $\varepsilon$ and $|T_2|$; by Lemma~\ref{is-onto-recorded-numbers}
this holds if $O$ is recorded at $p$ and a pair of runs. For every lift $j$ let $P_j$ and $W_j$ be a
parameter record and weights in the exponent dictionary at $j$ satisfying its numerical conditions
(Lemma~\ref{is-onto-step-exp-dictionary}, Definition~\ref{is-onto-step-exp-bound-weights}). Suppose
that:
\begin{enumerate}
\item for every lift $j$, tuple $t$ and finite place $w$, the inclusion
  $p_w^{J_A(p_w)}C_{j,t}\subset(D_{j,t})_{(w)}$ of Lemma~\ref{is-onto-step-exp-sum} holds for the cut
  lattices of the realization of $p$ and its face, where $J_A$ is the Weyl exponent sum of $W_j$
  (Definitions~\ref{is-onto-step-exp-bound-lift-exponent} and~\ref{is-onto-step-exp-bound-recursion});
\item for every lift $j$ and every twisted place $w$ with $p_w$ in the range of $W_j$, there is an
  exponent model at $w$: Frobenius data and a family of Hida settings of rank $N_r$ and the weights
  $\lambda$ of $W_j$ that covers the block data at $(j,\mu_w)$ at $w$, such that $p_w$ is not a unit in
  any of their rings and each of them is linked at $p_w$ to the weights $W_j$ with these Frobenius data
  (Definition~\ref{is-onto-step-exp-bound-hida}).
\end{enumerate}
Then $O$ has the exponent bound for $N_p(C,C_H)$ (Definition~\ref{is-onto-local-estimates}(5)) for all
$C$ and $C_H$, with $J'(j,l)$ the recursive exponent of $P_j$ and $W_j$ at $l$
(Definition~\ref{is-onto-step-exp-bound-recursion}) for every natural number $l$.
\end{lemma}
\leanref{Development/AbelianVarietyTheory/Proofs/DivisorProofs/ExpansionOfNumericsModel.lean}{GlobalTheory.ExpTie.Asm.hasExpBound_at}
\begin{proof}
Let $C$ and $C_H$ be real numbers. The function $J_0$ of $N_p(C,C_H)$ is the lift exponent of the
realization of $p$ and its face (Definition~\ref{is-onto-step-exp-bound-lift-exponent}). The numbers
$J'(j,l)$ are natural numbers, so $J'(j,l)\geq0$. We check the remaining conditions of
Definition~\ref{is-onto-local-estimates}(5).

\emph{Nonnegativity.} Let $w$ be a twisted place and $j$ a lift. Hypothesis (1) at $(j,\mu_w,w)$ and
Lemma~\ref{is-onto-step-exp-bound-lift} give $J_0(\mu_w,w,j)\geq0$.

\emph{The sum.} Let $j$ be a lift. The dictionary gives: $W_j$ is indexed by the $f$ embeddings of the
host field, where $f$ is the number $f$ of $P_j$ and the rank $f=2n$ of the member; the special set of
$W_j$ is the set determined by the member, the set of prime divisors of $Q_S$ together with $q$, $d$,
$p_{de}$ and $p'_{de}$; $W_j$ has at most $|T_2|$ thick elements; and $P_j$ has the rank $m$ of the
member, $L=x^{\varepsilon}$ and $b=\exp(x^{\varepsilon})$. With the numerical conditions,
Lemma~\ref{is-onto-step-exp-bound-sum} applies with $T=|T_2|$: for every finite set $Q$ of primes,
\[
 \sum_{l\in Q}J'(j,l)\log l\le30\,|T_2|\,mb\,x^{2\varepsilon}+2mb.
\]
As the numbers $m$, $x$, $\varepsilon$ and $T_2$ of $O$ are $m$, $x$, $\varepsilon$ and $|T_2|$, the
right side is the number $E_J$ of $O$, with $b=\exp(x^{\varepsilon})$
(Definition~\ref{is-onto-record-numbers}(2)).

\emph{The bound at the twisted places.} Let $w$ be a twisted place and $j$ a lift. Then
Lemma~\ref{is-onto-step-exp-bound-place} applies to the realization of $p$ and its face, the lift $j$,
the tuple $\mu_w$ and the place $w$, with $P_j$ and $W_j$: its hypothesis (1) is hypothesis (1) at
$(j,\mu_w,w)$, and its hypothesis (2) is hypothesis (2) at $w$, which is twisted. So
$J_0(\mu_w,w,j)\le J'(j,p_w)$.
% Lean: ExpTie.Asm.hasExpBound_at. J0 is JexpU of the block data (splitPartNumbers_exponent_eq_liftExponent,
% by rfl); the nonnegativity is T.exponent_datum_le_of_primePowerLatticeInclusion with hS at
% ((S2 p).mu w); Jp j l = Jrec7 (P_j, W_j) l for every l; the sum is UpperTheory.ExcSet.sum_exponent_log_le
% with the fields card, f, special, thick, log, b and m of Dict and hT2, rewritten with hx, heps and hm;
% the bound is jexpU_le_jrec7 with hS and hM. The fields k, t, Nr, x, lam, forced, box and big of Dict
% are not used.
\end{proof}

\subsection{From the estimates to the surjectivity bound}\label{subsection-is-surjectivity-bound}

\begin{lemma}\label{is-onto-estimates}
Assume the published results on unitary groups of Theorem~\ref{lit-bundle-unitary} and the local
Langlands laws of Chapter~\ref{chapter-galois-representations}. Let $\mathcal H$ be the class of
standard families. Eventually along $\mathcal H$, every family $\mathcal F$ has the following
properties, where $p$ runs over the points at members at large scales that satisfy the point
conditions and the further conditions (Definition~\ref{is-point}).
\begin{enumerate}
\item (Runs and rationality.) Every such $p$ has a pair of runs pinned to the own places, and its
  realization is rational at the twisted places.
\item (Export.) At every twisted place $v$ and every lift $j$, the twisted lattices of the realization
  are exported to the local models at $v$, and they span over the classes of lattices of the own
  place of $v$.
\item There is $c_0$ such that, for all constants $c$ with $C\ge c_0$ as in
  Definition~\ref{is-record-object}, at the points $p$ beyond a scale that may depend on $c$, every
  record object $O$ recorded at $(p,r)$, for a pair of runs $r$ pinned to the own places, satisfies
  the following, where $C'$ and $C_H$ run over the real numbers and $N_p(C',C_H)$ are the numerical
  bounds of $p$.
  \begin{enumerate}
  \item (Page.) The page estimates (Definition~\ref{is-onto-page-estimates}) for the residue fields
    $k_P$, $P\in P_g$, of the realization of $p$ (Definition~\ref{is-realisation}(1)) and $f_N=f$;
    the number $b$ of $O$ is the quantity $b$ of the size data of the member; $E_1^{u}\geq0$ for the
    scale data $Z(p)$; and the function $E_1$ of $N_p(C',C_H)$ is $E_1(j)=C'E_1^{u}(O)$ for every
    lift $j$.
  \item (Geometry.) The geometric estimates (Definition~\ref{is-onto-geometric-estimates}) and the
    depth estimates for $N_p(C',C_H)$ (Definition~\ref{is-onto-depth-estimates}).
  \item (Rank.) The rank at the twisted places (Definition~\ref{is-onto-local-estimates}(1)).
  \item (Twist.) A twist bound (Definition~\ref{is-onto-local-estimates}(2)).
  \item (Share.) The share bound (Definition~\ref{is-onto-local-estimates}(3)).
  \item (Bad places.) The bad-place bound (Definition~\ref{is-onto-local-estimates}(4)).
  \item (Exponents.) The exponent bound for $N_p(C',C_H)$
    (Definition~\ref{is-onto-local-estimates}(5)).
  \end{enumerate}
\end{enumerate}
\end{lemma}
\leanref{Development/AutomorphicArithmetic/ThetaCorrespondence/ThetaPairLifts/ThirdPlaceSeries/ExportSt.lean}{GlobalTheory.BigOn.ForallRecordObjects}
\leanref{Development/AutomorphicArithmetic/ThetaCorrespondence/ThetaPairLifts/ThirdPlaceSeries/ExportSt.lean}{GlobalTheory.BigOn.ExportedStatement}
\leanref{Development/AutomorphicArithmetic/ThetaCorrespondence/ThetaPairLifts/ThirdPlaceSeries/ExportSt.lean}{GlobalTheory.BigOn.PageStatement}
\leanref{Development/AutomorphicArithmetic/ThetaCorrespondence/ThetaPairLifts/ThirdPlaceSeries/ExportSt.lean}{GlobalTheory.BigOn.GeometricStatement'}
\leanref{Development/AutomorphicArithmetic/ThetaCorrespondence/ThetaPairLifts/ThirdPlaceSeries/ExportSt.lean}{GlobalTheory.BigOn.TwistStatement'}
\leanref{Development/AutomorphicArithmetic/ThetaCorrespondence/ThetaPairLifts/ThirdPlaceSeries/ExportSt.lean}{GlobalTheory.BigOn.ShareStatement}
\leanref{Development/AutomorphicArithmetic/ThetaCorrespondence/ThetaPairLifts/ThirdPlaceSeries/ExportSt.lean}{GlobalTheory.BigOn.BadPlacesStatement}
\leanref{Development/AutomorphicArithmetic/ThetaCorrespondence/ThetaPairLifts/ThirdPlaceSeries/ExportSt.lean}{GlobalTheory.BigOn.ExponentStatement}
\leanref{Development/AutomorphicArithmetic/ThetaCorrespondence/ThetaPairLifts/ThirdPlaceSeries/OntoXOfPrints.lean}{GlobalTheory.BigOn.ontoX_of_hypotheses_alt_alt}
\begin{proof}
Each of (1), (2) and (3a) to (3g) is proved separately, each with its own scale and, for (3a) to
(3g), its own constant; $c_0$ is the largest of the constants, and the scales are the largest of the
scales. In Steps 3 to 9, $p$ is a point at a member at a scale $x$ that satisfies the point
conditions and the further conditions, $r$ is a pair of runs pinned to the own places, and $O$ is a
record object with constants $c$ recorded at $(p,r)$. Its indefinite-side datum $S$ is the second
presentation of the rows of $p$: the rows of the realization of $p$ at its face with the scale data
$Z(p)$, after its scalars $\kappa_v$ and $\lambda$ are replaced by those of the characterization of
the realization. It has the twisted places of $S_0(p)$, and its rank $D$ is
positive by the links of the further conditions. The group $P$ of $O$ is the product of the
mirabolic subgroups of the factors $GL_{r_0}(k_P)$ of $G$, where $r_0=k+m-1$.

\emph{Step 1: runs and rationality.} No published input enters, and (1) holds at every scale. Let $p$
be a point that satisfies the point conditions and the further conditions, with realization $R$,
taken with the scalars of its characterization, and face $F$. The last of the further conditions is
the hypothesis on the twisted places of Lemma~\ref{is-onto-step-links}, so $R$ and $F$ have the links.
The core ties hold by the further conditions. Lemma~\ref{is-onto-step-runs-at-prime} with $l=p_{de}$
and with $l=p'_{de}$ gives runs $r_1$ and $r_2$, and Lemma~\ref{is-onto-step-level-structure} gives
maps $\ell_i$. Put $r=(r_1,r_2,\ell)$, and let $v$ be a finite place. If $v$ lies over $p_{de}$, the
run used at $v$ is $r_2$ and the own prime of $v$ is $p'_{de}$; otherwise the run used at $v$ is $r_1$
and the own prime is $p_{de}$. In both cases the run used at $v$ satisfies the laws, is covered and
rational at every lift, has $p_d(\delta_l(n))$ equal to the own prime of $v$, and is pinned at $v$ and
every lift if $v$ is twisted. These are (a), (b) and (f) of Definition~\ref{is-onto-step-runs}(5). As
it satisfies the laws, it satisfies (c) and (d) with $\ell$, and (e) is a condition of the core ties.
So $r$ is pinned to the own places.

For the rationality, Lemma~\ref{is-onto-step-rational} applies to $R$ and $F$. The links give the
characterization, condition (11), the rank condition, the base links and the pins of the one-map data.
The point conditions give an extremal vector that pairs to the delta at every lift, and they give
every lift as one map with the canonical constants of the core and the phase of the point; the change
to the scalars of the characterization makes these phases $1$. The ramified sets of a point satisfy
the guard. The recursion ties of the
further conditions contain, for every tuple $t$, the ray class field of modulus $M_t$ in the form of
that lemma. So $R$ is rational at the twisted places.

\emph{Step 2: export and span.} Global Artin reciprocity, the existence of ray class fields, the rapid
decay of cusp forms and the covering of the discrete series are among the results of
Theorem~\ref{lit-bundle-unitary}, and the Casselman-Wallach theorem is
Theorem~\ref{lit-casselman-wallach}; so Lemma~\ref{is-onto-step-export} gives the export in (2), with
its scale. The span over the classes of lattices of the own place uses no published input. Let $p$ be
as in Step 1, $v$ a twisted place and $j$ a lift. By the links, $R$ satisfies its characterization at
$F$. The characterization contains the core conditions at $F$ and, among its conditions on the tuples,
condition (2) of Lemma~\ref{is-onto-step-kernel-span}; and the core ties put the values of $\delta_v$ in
$\sigma_D(L)$. Lemma~\ref{is-onto-step-span},
with Lemmas~\ref{is-onto-step-kernel-unique} and~\ref{is-onto-step-kernel-span}, gives that the base
theta lattice $\Theta^{\mathrm b}_v(j)$ is the span of $\bigcup_c\Theta_\kappa(h_c)$ for every base
theta kernel $\kappa$ at $(j,v)$ and every family $(h_c)$ that meets the classes of the level $J_v$.

\emph{Step 3: the page estimates (3a).} By Lemma~\ref{is-onto-recorded-numbers}, $b=\exp(x^\varepsilon)$
is the quantity $b$ of the size data. The scale data of a point have
$E_1^{u}=x^\varepsilon(fmb\,x^\varepsilon+N_c)$, which is $E_1^{u}(O)$ by the same lemma, and the
function $E_1$ of $N_p(C',C_H)$ is $C'E_1^{u}$ (Definition~\ref{is-numerical-bounds}). This number
is nonnegative because $x,b>0$ and $N_c=|T_2|kbx^{2\varepsilon}/n^2\geq0$.

Conditions (1) to (4) of the page estimates hold by Lemma~\ref{is-onto-page-numerics} once
$C\geq1$ and $x\geq3C$, so the constant of (3a) is at least $1$ and its scale at least $3C$.

For (5), the realization identifies $P_g$ with the disjoint union of two sets of $f$ places, the
residue fields at the places of the $i$-th set being $\mathbf F_{g_i}$; and $g_0\ne g_1$ are primes.
So every $k_P$ has $g_0$ or $g_1$ elements, and $f$ of them have $g_i$ elements for $i=0,1$.

For (6), the parameter choices of the member give a number $Y_g$ with $(kb)^8\le Y_g\le g_i$ for
$i=0,1$. Since $x>1$, $b=\exp(x^\varepsilon)\geq1+x^\varepsilon\geq2$, and $1\le f\le k$, so
\[
 \min(g_0,g_1)-1\geq256k^8-1\geq8k^7\geq8f^4k^3.
\]
This is $2f/(\min(g_0,g_1)-1)\le1/(4(fk)^3)=y$.

\emph{Step 4: the geometric estimates (3b).} No published input enters. Let $C_K$ and $x_a$ be the
constants of Lemma~\ref{is-onto-step-index-conditions}, $x_b$ the number that
Lemma~\ref{is-onto-step-leaves} gives for $C_K$, and $C_d$ the constant of
Lemma~\ref{is-onto-step-depth}. The constant of (3b) is the larger of the constant of (3a) and $C_d$,
and its scale is the largest of the scale of (3a) and the scales of Lemma~\ref{is-onto-step-numerics}
for $x_1=x_a$ and for $x_1=x_b$. Let $p$, $r$ and $O$ be as above, at such a scale. By
Lemma~\ref{is-onto-step-numerics}, the member satisfies the conditions of
Lemma~\ref{is-onto-step-index-conditions} that involve $x_a$ and those of
Lemma~\ref{is-onto-step-leaves} that involve $x_b$, and has $3\le x$, $kb'\le b$, $C_{nW}=10$ and
$3\le x^{\varepsilon_z/5}$. The core ties hold and $r$ is pinned to the own places, so the runs of $r$
satisfy the index-parameter conditions with $C_K$ by Lemma~\ref{is-onto-step-index-conditions}. By the
links (Step 1), the realization satisfies its characterization at its face, so
Lemma~\ref{is-onto-step-leaves} gives the supremum, span, containment and data estimates (1) to (4).
By (3a), the number $b$ of $O$ is the quantity $b$ of the size data, so
Lemma~\ref{is-onto-step-index-bound} gives the index bound (5).

For the depth estimates, let $C'$ and $C_H$ be real numbers and $N=N_p(C',C_H)$. The function $F_0$ of
$N$ is that of the scale data $Z(p)$, which by the construction of the scale data of a point is, at
the lift $j$, the local term of the definite-side datum at level $j+1$. By (3a), $O$ satisfies (1)
and (2) of the page estimates, and $C\geq C_d$. So Lemma~\ref{is-onto-step-depth} gives the depth
estimates for $N$.

\emph{Step 5: the rank (3c).} No published input enters. The constant of (3c) is that of (3b), and its
scale is the largest of the scales of (3b), of Step 2 and of Lemma~\ref{is-onto-step-compare}. Let
$p$, $r$ and $O$ be as above. By Lemma~\ref{is-onto-step-theta-j}(1), every $e_j$ is an idempotent
with image $V_0(j)$ that commutes with $\rho_j(P)$. Let $v$ be a twisted place or the conjugate of
one, which is twisted by the characterization of the realization, and $j$ a lift. Step 2 gives the span
of the base theta lattice over the classes at $v$ and $j$, and Lemma~\ref{is-onto-step-compare} the
comparison away from the own place for the run used at $v$. So Lemma~\ref{is-onto-step-theta-j}(2)
gives that $\Theta_J(v,j)$ is a finitely generated submodule of $V_\psi(j)$ with
$\Theta_0(v,j)_{(v)}\subset\Theta_J(v,j)_{(v)}$. The rest is short. By the span estimate,
$\Theta_K(v,j)$ is a lattice in $K^{n_Z(j)}$, so $\Theta_K(v,j)\cap V_\psi(j)$ spans $V_\psi(j)$; by
the data estimate, $\Theta_0(v,j)$ is finitely generated, lies in $V_\psi(j)$ and contains
$|U|(\Theta_K(v,j)\cap V_\psi(j))$, so it spans $V_\psi(j)$. Localization does not change the span,
so $\Theta_J(v,j)$ also spans $V_\psi(j)$. Since $e_j$ commutes with $\rho_j(U)$ and fixes $V_0(j)$,
$e_j(V_\psi(j))=Z_g(j)$. Hence $e_j\Theta_0(v,j)$ and $e_j\Theta_J(v,j)$ are lattices in $Z_g(j)$.

\emph{Step 6: the twist bound (3d).} No published input enters beyond the export of (2). Standard
families lie in $\operatorname{Des}(\mathcal C_N,3)$, hence in $(\mathcal C_N)_{\mathrm{tw}}$
(Definition~\ref{asm-designed}), so they satisfy the twist condition
(Definition~\ref{asm-twist-condition}): beyond a scale every member has numbers $N_0(i)$ and sets
$S(i,j,b)$ as in (2) of that definition. Fix a lift $j$ and a twisted place $v$ of $S$; it is
twisted for the member of $p$. Let $(i,b,Q,k',\psi)$ be the twist choice at $v$, so that
$\ell_i=p_v$; put $\varphi=\psi^{(b)}\circ\operatorname{art}$, let $\Sigma$ be the set of theta data
that the twist choices give at the member for $(i,j,b)$, written $\Sigma_D(i,j,b)$ in that
definition, and let $\epsilon_{j',b'}\geq0$ be the numbers of (2)(c) at $v$. The twist choice also
carries a discrete valuation ring $\mathcal O'$ containing $\mathcal O_{i,k'}$ and finite over it,
and an embedding
$\sigma_v$ of $L$ in the fraction field $F'$ of $\mathcal O'$ such that $v(x)\geq0$ if and only if
$\sigma_v(x)\in\mathcal O'$.

The tie of $O$ at $(v,j)$ is an $L$-linear bijection $t$ from $K^{n_Z(j)}$, where $K=L$ is the
field of $S$, onto the coordinates $L^{n_W(j)}$ of $FW_j$ with $t\circ e_j=e_{\mathrm{good}}\circ t$,
where $e_{\mathrm{good}}$ is the operator of Step 8 by which the good idempotent of the family acts
on $FW_j$. It carries $\Theta_J(v,j)$ onto the \emph{base theta lattice} $\Theta^{\mathrm b}_v$, the
$\mathcal O_L$-span of the coordinate vectors of the theta series, traced to the level
$(P_0,\psi_0^{-1})$, of the theta kernels of the realization for its character $\delta'$ if
$p_v=p_{de}$, and $\delta$ otherwise, with level at the places of $Q_{\mu_v}$ over the own
place of $v$; and $X_3(v,j)=e_jt^{-1}(\Theta_{\mu_v})$, where
$\Theta_{\mu_v}$ is the $\mathcal O_L$-span of the coordinate vectors of the theta series of the
Weil datum of the realization at the tuple $\mu_v$, traced to the same level. The export of (2) at
$(v,j)$ is a $\sigma_v$-semilinear map $\mathrm{rd}$ from $L^{n_W(j)}$ to
$F'\otimes_{\mathcal O_i}F_{i,j,b}$ that carries some basis to a linearly independent family, satisfies
$\mathrm{rd}\circ e_{\mathrm{good}}=e_{i,b}\circ\mathrm{rd}$, and carries $\Theta^{\mathrm b}_v$ and
$\Theta_{\mu_v}$ to sets that span $\mathcal O'\otimes\Theta(\Sigma)$ and
$\mathcal O'\Theta_{\varphi^{(b)},Q}(\Sigma)$ over $\mathcal O'$. So $R=\mathrm{rd}\circ t$ is
$\sigma_v$-semilinear, carries a basis of $K^{n_Z(j)}$ to a linearly independent family, and
\[
 \mathcal O'R\big(e_j\Theta_J(v,j)\big)=\mathcal O'\otimes e_{i,b}\Theta(\Sigma),\qquad
 \mathcal O'R\big(X_3(v,j)\big)=\mathcal O'\Theta',\qquad\Theta'=e_{i,b}\Theta_{\varphi^{(b)},Q}(\Sigma),
\]
in the notation of Lemma~\ref{is-onto-twist-transport}.

By (2)(a), $e_{i,b}\Theta(\Sigma)$ is a lattice in $V_0(i,j,b)$: the module $|G_0|e_{i,b}\Theta(\Sigma)$
is spanned by the finitely many vectors $|G_0|e_{i,b}\theta(a)$, $a\in S(i,j,b)$, and spans
$V_0(i,j,b)$. By (2)(c), $\Theta'$ is a lattice in $V_0(i,j,b)\otimes\mathcal O_{i,k'}$. So both
$\mathcal O'$-modules in the display are lattices in $V_0(i,j,b)\otimes\mathcal O'$, and the index
$[\mathcal O'\otimes e_{i,b}\Theta(\Sigma):\mathcal O'\Theta']$ of Lemma~\ref{is-onto-twist-transport},
which is that of Definition~\ref{asm-twist-condition} since $\ell_i=p_v$, is defined. The index bound
of (2)(c) with $M=\mathcal O'\otimes e_{i,b}\Theta(\Sigma)$ gives
$[\mathcal O'\otimes e_{i,b}\Theta(\Sigma):\mathcal O'\Theta']\le\epsilon_{j,b}$.

The supremum estimate of (3b) makes $M_s(j)$ a lattice in $K^{n_Z(j)}$, so $M_s(j)\cap Z_g(j)$ is a
lattice in $Z_g(j)$: it is finitely generated, and every vector of $Z_g(j)$ has a nonzero multiple
in $M_s(j)$. By (3c), $e_j\Theta_J(v,j)$ is a lattice in $Z_g(j)$. Lemma~\ref{is-onto-twist-transport},
with $Z=Z_g(j)$, $M=M_s(j)\cap Z_g(j)$, $\Theta=e_j\Theta_J(v,j)$, $X=X_3(v,j)$ and
$\varepsilon=\epsilon_{j,b}$, gives the inequality of the twist bound at $(v,j)$ with
$\varepsilon'(v,j)=\epsilon_{j,b}\geq0$.

For the sum, (2)(c) gives $\sum_j\varepsilon'(v,j)\le\sum_{j'=1}^m\sum_{b'}\epsilon_{j',b'}\le
\ell_i^{-2}=p_v^{-2}$. The conjugate $\bar w$ of a twisted place $w$ is twisted, with $p_{\bar w}=p_w$,
so $\sum_j\big(\varepsilon'(w,j)+\varepsilon'(\bar w,j)\big)\le2/p_w^{2}$.

\emph{Step 7: the share bound (3e).} None of the published results assumed in the statement enters;
the Borel-Weil theorem enters through Lemma~\ref{is-onto-share-definite-finite}, the Witt decomposition
and the Hermite normal form through Lemma~\ref{is-onto-share-semistable}, and Fourier-Jacobi expansions
through Lemma~\ref{is-onto-share-sections}. Call a subgroup of $G$
\emph{half-unipotent} if its elements are upper unitriangular at all the places of $P_g$ over
$g_0$, or at all those over $g_1$. A \emph{model} of the $j$-th lift with parameter $T$ consists of
a finite set $C$ with an action of $G$; two-level data $M$ of rank $r_0$ over the host field for $G$
and $C$ (Definition~\ref{is-onto-share-two-level}), whose set $\Psi$ of embeddings has $f$ elements; a
weight $\kappa$, with the space $S_\kappa=S_\kappa(M)$ of cusp forms and a representation $\rho$ of $G$
on $S_\kappa$ pulled back from the action of $G$ on $C$; an embedding $\sigma$ of $L$ in $\mathbf C$; and a $\sigma$-semilinear map $f$ from the coordinates
$L^{n_W(j)}$ of $FW_j$ to $S_\kappa$ that is compatible with the actions of $G$ and carries the
standard basis to a basis of $S_\kappa$. The level hypothesis of $M$ must hold for every
half-unipotent subgroup of $G$; and either $p_0\in\{0,r_0\}$, the \emph{definite} case, or
$0<p_0<r_0$ and $M$ satisfies the archimedean condition at $T$ for the weight $\kappa$. By
Lemma~\ref{is-onto-step-models}, beyond a scale the $j$-th lift of every such point has a model with
parameter $T$, the number $T$ of $O$, which is the integer $T_O$ of the parameter choices of the
member; there $r_0$ is the rank of the definite-side datum, which is the number $r_0$ of $O$ since $O$
is recorded at $(p,r)$.

Put $N_u=fr_0(r_0-1)/2$, which is $\dim Y(M)$, so that $s(M)=2r_0+50r_0^2fp_0(r_0-p_0)+2$. Two
dimension statements hold for the cusp forms of a model. Let $H$ be a subgroup of $G$ for which the
level hypothesis holds and $\lambda$ a character of $H$, and let $S_\kappa^\lambda=S_\kappa(M)^{H,\lambda}$
be the subspace of $S_\kappa$ on which $H$ acts through $\lambda$. In the definite case,
$|H|\dim S_\kappa^\lambda=\dim S_\kappa$ by Lemma~\ref{is-onto-share-definite}. If $0<p_0<r_0$, $M$
satisfies the archimedean condition at $T$ for $\kappa$, the level hypothesis holds for the trivial
subgroup and $s(M)\le T$, then the cusp forms of $M$ have ranks $(N_u(N_u+6)/T,2)$ at $\kappa$ and
$s(M)$ by Lemma~\ref{is-onto-share-ranks}, and Lemma~\ref{is-onto-share-indefinite} gives
$(1-N_u(N_u+6)/T)\dim S_\kappa\le|H|\dim S_\kappa^\lambda$. Beyond a further scale,
$2r_0+50r_0^2fs(r_0-s)+2\le T$ for $0\le s\le r_0$ and $N_u(N_u+6)/T\le y$, where $y=1/(4(fk)^3)$.

Take $H=P_0$, which is half-unipotent by the links of the point conditions, and
$\lambda=\sigma\circ(\chi_j\psi_0^{-1})$ on $P_0$. The $\mathbf C$-linear extension of $f$ is an
isomorphism $\mathbf C\otimes_{\sigma}L^{n_W(j)}\to S_\kappa$ compatible with the actions of $G$, and
it carries $\mathbf C\otimes_\sigma FW_j^{0}$ onto $S_\kappa^\lambda$, with $FW_j^{0}$ as in
Lemma~\ref{is-onto-recorded-numbers}(3). So $\dim S_\kappa=n_W(j)$ and
$\dim S_\kappa^\lambda=\dim FW_j^{0}$. In the definite case the first statement gives
$|P_0|\dim FW_j^{0}=n_W(j)$. Otherwise the trivial subgroup is half-unipotent, so the level
hypothesis holds for it, and $s(M)\le T$ by the first numerical bound with $s=p_0$; the second
statement and $N_u(N_u+6)/T\le y$ give $(1-y)n_W(j)\le|P_0|\dim FW_j^{0}$. Since $y\geq0$, in both
cases the \emph{shared weight bound}
\[
 \big(1-y\big)\,n_W(j)\le|P_0|\dim FW_j^{0}
\]
holds. Lemma~\ref{is-onto-recorded-numbers}(3) turns it into $(1-y)n_Z(j)\le|U|\,d_0(j)$, which is
the share bound with $\epsilon=y$.

\emph{Step 8: the bad-place bound (3f).} The good idempotent of the family acts on $FW_j$ as an
idempotent $e_{\mathrm{good}}$ commuting with $G$, so it preserves $FW_j^{0}$, and
\[
 \dim FW_j^{0}=\dim\big(FW_j^{0}\cap e_{\mathrm{good}}FW_j\big)+\dim\big(FW_j^{0}\cap\ker e_{\mathrm{good}}\big).
\]
The first space is the good part of the rows at the lift $j$, of dimension $d_W(j)$, which is the
rank $D$ of $S$ by the links of the further conditions; and $\dim FW_j^{0}=d_0(j)$ by
Lemma~\ref{is-onto-recorded-numbers}(3). Beyond the scales of Lemmas~\ref{is-onto-step-models}
and~\ref{is-onto-step-bad-numerics}, Lemma~\ref{is-onto-step-bad-weight} applies to $p$, the lift $j$,
$T=T_O$ and $\delta=\exp(-x^{\varepsilon_z/7})$: the point satisfies the point conditions, its
recursion ties hold by the further conditions, Lemma~\ref{is-onto-step-bad-numerics} gives the
conditions on $r_0$, on the types, on $\delta$ and on $T_O$, and Lemma~\ref{is-onto-step-models} gives
a model of the lift $j$ with parameter $T_O$. It gives the \emph{bad weight bound}
$\dim(FW_j^{0}\cap\ker e_{\mathrm{good}})\le\delta\dim FW_j^{0}$, and $\delta\le\frac14$ by
Lemma~\ref{is-onto-step-bad-numerics}. Hence $(1-\delta)d_0(j)\le D$ with $0\le\delta\le\frac14$.

\emph{Step 9: the exponent bound (3g).} The constant of (3g) is $0$, and its scale is the largest of
the scales of Lemmas~\ref{is-onto-step-exp-dictionary}, \ref{is-onto-step-exp-sum}
and~\ref{is-onto-step-exp-model}. For every lift $j$, choose a parameter record $P_j$ and weights $W_j$ in
the exponent dictionary at $j$ satisfying its numerical conditions
(Lemma~\ref{is-onto-step-exp-dictionary}). Lemma~\ref{is-onto-step-exp-sum} gives the inclusions of
the block lattices for all $j$, $t$ and $w$, and Lemma~\ref{is-onto-step-exp-model} the exponent models
at the twisted places in the range of $W_j$. With the exact cut bound of
Lemma~\ref{is-onto-step-exp-cut}, Lemma~\ref{is-onto-step-exp-bound} gives the exponent bound for
$N_p(C',C_H)$ for all $C'$ and $C_H$. The Galois input enters only through
Lemma~\ref{is-onto-step-exp-model}.
% The open parts of this proof are the open steps of the is-onto-step-* lemmas above, and those of
% is-onto-share-interpolation (the free ball and the orbit interpolation; used by
% is-onto-share-translates) and is-onto-share-counts, used in Step 7. That lemma is written out through
% is-onto-share-counts-model, -jets, -reduction, -maximum, -no-maximum, -upper, -section,
% -interpolation and -lower; only -reduction, -section and -interpolation still have open steps.
% The semistable majorant (is-onto-share-semistable, through the
% is-onto-semistable-* lemmas) and the Fourier-Jacobi facts (is-onto-sections-fourier, cited) are
% written out.
% Step 1 is RunsExist.exists_runs_at_alt (bigOn_runsSt), with ownPin_of for (f), and
% twistedRational_R2_F2 (bigOn_ratSt) with ioExt_R2, one_map_R2 (oneMapRowR_charOf, phEff_eq_one),
% splitPartRealisation_satisfies_qGuard and RecTies.rayQ. Step 2 is
% T.exportSt_of_reciprocity_and_casselman_wallach, after HeckeLAdicCharacter.globalArtin,
% GlobalArtinReciprocity.rayClassFieldSt and seriesRatS_and_serInt_of_globalArtin in
% ontoX_of_hypotheses_alt_alt, and maxSt_alt (spanOverClasses_of, with CoreTies.delAt_mem and FO_le).
% Step 3 is pageSt_alt: numericLeaves (Lemma is-onto-page-numerics), levelSq_yBound_atRec_alt
% (leafLevelSq from RealDat.grp_record, leafYBound from ParameterChoice'.g_Y) and
% twoRunDatum_boundUnit_at_indefinitePoint. Step 4 is geoSt_of_pageSt_alt: runsHold_of_ties,
% leaves_sup_span_data, hasIndexBound_of_runs and leafDepth_all_record, with
% T.runNumericalConditions_hold_along_family for x_a and x_b and
% indefinitePoint_splitPartContent_eq_forcedContent (numOf_Fj_eq_ofRecord) for F_0.
% Step 5 is rankLocSt_of_max_geo_runs = rank_at_twisted_places_of_rec_span (conjPl_mem_twPl,
% tie_thetaJ, maxSt_alt, compareOffOwn_atRuns_alt, and T.eigenlattice_local_containment_at_pairs for the
% short part). Step 6 (written out) is twistSt_of_alt =
% hasTwistBound_at: exists_leafTW1_of_twGoodAt with FamHyp's TwGood, TiedAt.lix_le (OnTie.ThJ_eq,
% X₃_eq, e_eq; ExportTie = Reads), Reads.lix_le (Lemma is-onto-twist-transport), sum_side_le,
% l_eq_plc, goodPartLattice_of_leafSup and leafTW1_of_places. Step 7 is modelSt_alt (ModelAt),
% shareWSt_of_modelSt_alt (shareW_of_modelAt, shareSize_bounds_along_family) and
% leafShareI_of_shareW; in shareW_of_modelAt the two dimension statements are
% T.definite_cuspForms_finite_regularRepresentation (Lemma is-onto-share-definite, written out down to
% Closed.semistable = Lemma is-onto-share-semistable) and T.twoLevelFlagDomain_cuspForm_ranks_shifted_weights
% (Lemma is-onto-share-ranks, with ShareRanks.compT = Lemma is-onto-share-counts, written out through
% UpperModel.upper and T.card_independent_boundedSections_shiftedWeight_le (is-onto-share-counts-upper),
% T.exists_nonzero_boundedSection_shiftedWeight (-section) and T.lowerRank_cuspForms_detWeight_oneComponent
% (-lower), with ArchOne = ArchInM, Shares.Margins, FlatM (-model); the open steps left there are in
% -reduction (FD.reductionTheory_holds), -section and -interpolation)
% with T.twoLevelFlagDomain_eigenspace_share (Lemma is-onto-share-indefinite, using
% T.exists_cuspForm_independent_translates (open below LiftInterp.comp_interp), siegelSpaceAlgebra_holds
% (open below moderateCoefficientPositivity_holds and coefficientConvolution_holds) and
% flag_domain_frames_connected (written out));
% the transport is Alg.finrank_eig_eq and Alg.finrank_eq_of_semilinear, and P_0 is half-unipotent by
% T.exists_whittaker_lam_psi_half_unip. The models are modelSt_alt (Lemma is-onto-step-models).
% Step 8 is badWSt_of_greenCuspidal_of_modelSt (badW_of_del_bound_of_T_bound,
% badShareSize_bounds_along_family, hodgeLineMultiple_le_familyThreshold) and badSt_of_badW_alt
% (bad_rec_alt). Step 9 is expSt_of_base_change_and_crystalline_lattices =
% T.expSt_of_num_sumExp_model = ExpTie.Asm.expSt_of_alt, with hasExpBound_at.
\end{proof}

\begin{definition}\label{is-places-error-terms}
Let $O$ be a record object for an indefinite-side datum over a number field $K$, a finite group $P$
and constants $c=(C,C_{b,1},C_{b,2})$, with the notation of Definitions~\ref{is-record-object}
and~\ref{is-onto-record-numbers}. For a finite place $v$ of $K$, with residue characteristic $p_v$
and ramification index $e_v$, put
\[
 v_U(v)=\frac{v_v(|U|)}{e_v},\qquad
 g_q(v)=\begin{cases}2fr_0^{2}&\text{if }p_v\in\{g_0,g_1\},\\0&\text{otherwise.}\end{cases}
\]
For a lift $j$ and an integer $l\geq1$, the \emph{error at $l$} is
\[
 \mathrm{err}_P(j,l)=\begin{cases}
 4T+4fr_0^{2}&\text{if }l\in\{g_0,g_1\},\\
 2s_C(j)+2T+2&\text{if }l\in S_m\setminus\{g_0,g_1\},\\
 2\big(s_B(j)+T+1\big)+4\beta(l)&\text{if }l\in\{q,d\}\setminus\big(S_m\cup\{g_0,g_1\}\big),\\
 2a_E(j,l)+2&\text{if $l$ is not a bad prime and }l\le Y,\\
 0&\text{if $l$ is not a bad prime and }l>Y,
 \end{cases}
\]
where $\beta(l)=fkb'$ if $l=d$ and $\beta(l)=0$ otherwise. Put $C_b=C_{b,1}+C_{b,2}+4$, so that
$4C^2+25C+C_b+44=4C^2+25C+C_{b,1}+C_{b,2}+48$, and
\[
 E_A(j)=2\Big(A_1(j)\log Y+A_w(j)\frac Ck+Cmb+\frac{2.04\,A_w(j)}{X_2}+Y\log4\Big)
 +\sum_{l\text{ bad}}\mathrm{err}_P(j,l)\log l,
\]
where $l$ runs over the bad primes.
\end{definition}
\leanref{Formalism/LocalToGlobal/IndefiniteForms/IndefiniteTree/ExampleSatisfiesConditionsAtFive.lean}{Formalism.L2.Indef.Tree.OnObj.vU0}
\leanref{Formalism/LocalToGlobal/IndefiniteForms/IndefiniteTree/ExampleSatisfiesConditionsAtFive.lean}{Formalism.L2.Indef.Tree.OnObj.gq}
\leanref{Formalism/LocalToGlobal/IndefiniteForms/IndefiniteTree/ExampleSatisfiesConditionsAtFive.lean}{Formalism.L2.Indef.Tree.OnObj.errP}
\leanref{Formalism/MirabolicMackeyTheory/StatementT.lean}{Formalism.EigenvalueLevelCeiling.errGood}
\leanref{Formalism/LocalToGlobal/IndefiniteForms/IndefiniteTree/ExampleSatisfiesConditionsAtFive.lean}{Formalism.L2.Indef.Tree.OnObj.Cb}
\leanref{Formalism/LocalToGlobal/IndefiniteForms/IndefiniteTree/ExampleSatisfiesConditionsAtFive.lean}{Formalism.L2.Indef.Tree.OnConst.CE}

\begin{lemma}\label{is-places-prime-sums}
Let $Q$ be a finite set of primes, and let $Y>1$ and $X>0$ be real numbers. Then
\[
 \sum_{l\in Q,\ l\le Y}\frac{\log l}{l}\le\log Y,\qquad
 \sum_{l\in Q,\ l\le Y}\log l\le Y\log4,\qquad
 \sum_{l\in Q,\ l>X}\frac{\log l}{l^{2}}\le\frac{2.04}{X}.
\]
\end{lemma}
\leanref{Formalism/MirabolicMackeyTheory/MirabolicMackeyProofs/AbsorbLocalConditionsOfExact.lean}{Formalism.EigenvalueLevelCeiling.primeBoundsExact}
\leanref{Formalism/PeriodIntegrals/OrbitalBoundsConclusion/PrimeLogOverP.lean}{sum_log_p_over_p}
\leanref{Formalism/MirabolicMackeyTheory/MirabolicMackeyProofs/LocalConditionsComposeGeneral.lean}{Formalism.EigenvalueLevelCeiling.chebyshev_theta_log_four}
\leanref{Formalism/PeriodIntegrals/OrbitalBoundsConclusion/PrimeLogOverPSqTail.lean}{prime_log_over_sq_tail}
\leanref{Formalism/PeriodIntegrals/UnitaryGroupMassFormula/MertensLogSum/PrimeLogOverPTop.lean}{PI.PrimeLogOverPWork.sum_log_div_prime_le}
\leanref{Formalism/PeriodIntegrals/UnitaryGroupMassFormula/PrimeLogarithmTail/SumLogDivSqLeDiv.lean}{PI.TailWork.sum_log_div_sq_le_div}
\leanref{Formalism/PeriodIntegrals/UnitaryGroupMassFormula/PrimeLogarithmTail/SumLogBlock.lean}{PI.TailWork.tail_odd}
\begin{proof}
Write $\theta(t)=\sum_{l\le t}\log l$ and $S(n)=\sum_{l\le n}\log l/l$, the sums over the primes $l$.
All three sums have nonnegative terms, so it suffices to bound the sums over all primes in the given
ranges.

\emph{The second bound.} It is $\theta(Y)\le Y\log4$, which follows from
$\prod_{l\le n}l\le4^n$ for $n=\lfloor Y\rfloor$. This holds by induction on $n$: for even $n>2$ the
product is that for $n-1$, and for $n=2m+1$ the primes $l$ with $m+1<l\le2m+1$ divide the binomial
coefficient $\binom{2m+1}{m}$, which is at most $4^m$.

\emph{The first bound.} Let $n=\lfloor Y\rfloor\geq1$. Since $\log n\le\log Y$, it suffices to show
$S(n)\le\log n$. For $n\le63$ this is a finite check: $S$ is nondecreasing, so it suffices that
$S(b)\le\log a$ for $(a,b)=(1,1),(2,2),(3,4),(5,10),(11,30),(31,63)$. Let $n\geq64$. Legendre's formula
gives $\log n!=\sum_{d\le n}\Lambda(d)\lfloor n/d\rfloor$, with $\Lambda$ the von Mangoldt function.
Keeping only the primes and $d\in\{4,8,16,9,27,25\}$, and using $\lfloor n/d\rfloor\geq n/d-1$,
\[
 \log n!\geq nS(n)-\theta(n)+\Big(\frac{7n}{16}-3\Big)\log2+\Big(\frac{4n}{27}-2\Big)\log3
 +\Big(\frac n{25}-1\Big)\log5 .
\]
On the other side $\log n!\le n\log n-n+1+\frac12\log n$, because the sequence
$n!\,e^n/n^{n+1/2}$ is nonincreasing, and $\theta(n)\le n\log4$ by the second bound. With
$\log3\geq\frac{19}{12}\log2$, $\log5\geq\frac{51}{22}\log2$, $0.6931471803<\log2<0.6931471808$ and
$\log n\le n/64-1+6\log2$, these give $nS(n)\le n\log n$.

\emph{The third bound.} Put $g(m)=1/(m(m-1))$ and $h(m)=1/(m-1)+1/m$ for $m\geq2$. Let first $k\geq3$
and suppose that every element of $Q$ is at least $2k+1$; its elements are then odd. For an odd prime
$l\geq2k+1$ a computation gives $g(\frac{l+1}2)-g(l)\geq(3-\frac1{2k+1})/l^2$. Write
$g(\frac{l+1}2)-g(l)$ as the sum of $g(m)-g(m+1)$ over $\frac{l+1}2\le m<l$ and exchange the sums.
The primes $l$ with $\frac{l+1}2\le m<l$ lie in $(m,2m]$, and these divide $\binom{2m}m\le4^m$, so
their logarithms sum to at most $m\log4$. Since $(g(m)-g(m+1))m=h(m)-h(m+1)$, this gives
\[
 \Big(3-\frac1{2k+1}\Big)\sum_{l\in Q}\frac{\log l}{l^2}\le\sum_{l\in Q}\log l\,
 \Big(g\Big(\frac{l+1}2\Big)-g(l)\Big)\le\log4\sum_{m>k}\big(h(m)-h(m+1)\big)
 \le\log4\Big(\frac1k+\frac1{k+1}\Big),
\]
and the resulting bound is at most $2.04/(2k+1)$ because $k\geq3$. If $X\geq5$, take
$k=\lfloor(\lfloor X\rfloor+1)/2\rfloor\geq3$: every prime $l>X$ is odd and at least $2k+1$, and
$X\le2k+1$, so the sum is at most $2.04/(2k+1)\le2.04/X$. If $X<5$, the primes among $2,3,5$ that
exceed $X$ contribute directly, and the others are at least $7$ and contribute at most $2.04/7$ by the
case $k=3$; with $\log2<0.6931471808$, $\log3\le1.11$ and $\log5\le1.62$ a direct check bounds the
total by $2.04/(\lfloor X\rfloor+1)\le2.04/X$.
\end{proof}

\begin{remark}\label{is-places-prime-sums-rs}
Lemma~\ref{is-places-prime-sums} also follows from the bounds of Rosser and Schoenfeld
\cite[\S3]{RS62}. They give the first bound, and $\theta(t)<1.01624\,t$ for $t>0$, with
$1.01624<\log4$, which gives the second. For the third, let $Z$ be the largest element of $Q$, which
we may assume is larger than $X$. By partial summation and $\theta(X)\geq0$,
\[
 \sum_{X<l\le Z}\frac{\log l}{l^{2}}=\frac{\theta(Z)}{Z^{2}}-\frac{\theta(X)}{X^{2}}
 +2\int_X^Z\frac{\theta(t)}{t^{3}}\,dt
 \le\frac{1.01624}{Z}+2.03248\Big(\frac1X-\frac1Z\Big)\le\frac{2.04}{X}.
\]
\end{remark}

\begin{lemma}\label{is-places-error-sum}
Let $O$ be a record object for an indefinite-side datum over a number field $K$, a finite group $P$
and constants $c=(C,C_{b,1},C_{b,2})$, with the notation of Definitions~\ref{is-record-object},
\ref{is-onto-record-numbers} and~\ref{is-places-error-terms}. Suppose that $O$ satisfies (1), (2)
and (4) of the page estimates (Definition~\ref{is-onto-page-estimates}), that $\vartheta(l)\geq0$ and
$\alpha(l)\geq0$ for every $l$, and that $\sum_{l\in Q}\vartheta(l)\log l\le C/k$ and
$T\sum_{l\in Q}\alpha(l)\log l\le Cmb$ for every finite set $Q$ of primes. Let $\mathcal S$ be a
finite set of finite places of $K$, with $\ell_w$, $e_w$ and $p_w$ as in
Definition~\ref{is-indefinite-datum}, and let $j$ be a lift. Then:
\begin{enumerate}
\item for every function $h\geq0$ on the primes,
  $\sum_{w\in\mathcal S}\ell_we_w\,h(p_w)\le\sum_{l}h(l)\log l$, where $l$ runs over the residue
  characteristics of the places in $\mathcal S$;
\item $\mathrm{err}_P(j,l)\geq0$ for every $l\geq1$;
\item $\sum_{w\in\mathcal S}\ell_we_w\,\mathrm{err}_P(j,p_w)\le E_A(j)$;
\item $\sum_{l\text{ bad}}\mathrm{err}_P(j,l)\log l\le C_b\,fmb\log b$.
\end{enumerate}
\end{lemma}
\leanref{Formalism/MirabolicMackeyTheory/MirabolicMackeyProofs/RamBasic.lean}{Formalism.EigenvalueLevelCeiling.place_sum}
\leanref{Formalism/LocalToGlobal/IndefiniteForms/IndefiniteTree/ExampleSatisfiesConditionsAtFive.lean}{Formalism.L2.Indef.Tree.errorTerm_nonneg}
\leanref{Formalism/MirabolicMackeyTheory/MirabolicMackeyProofs/AbsorbLocalConditionsOfExact.lean}{Formalism.EigenvalueLevelCeiling.local_sum_exact}
\leanref{Formalism/LocalToGlobal/IndefiniteForms/IndefiniteTree/ExampleSatisfiesConditionsAtFive.lean}{Formalism.L2.Indef.Tree.bad_sum}
\begin{proof}
(1) Let $l$ be a prime. The places $w$ above $l$ have $Nw=l^{f_w}$ and
$\sum_{w\mid l}e_wf_w=[K:\mathbf Q]$, so $\sum_{w\in\mathcal S,\,w\mid l}\ell_we_w\le\log l$.
Multiply by $h(l)\geq0$ and sum over $l$.

(2) The page estimates give $T\geq0$, $b'\geq0$, $S_e\geq0$, $\mathrm{wt}(j)\geq0$ and $k\geq1$, so
$s_B(j)$, $s_C(j)$, $\beta(l)$, $A_1(j)$ and $A_w(j)$ are nonnegative. With $\vartheta,\alpha\geq0$
this gives $a_E(j,l)\geq0$, and every case of $\mathrm{err}_P(j,l)$ is nonnegative.

(3) By (1) and (2), the sum is at most $\sum_l\mathrm{err}_P(j,l)\log l$ over the residue
characteristics $l$ of $\mathcal S$. The bad primes among them contribute at most the last term of
$E_A(j)$. A prime $l$ that is not bad contributes $0$ if $l>Y$, and if $l\le Y$ it contributes
\[
 \frac{2A_1(j)\log l}{l}+2A_w(j)\vartheta(l)\log l+2T\alpha(l)\log l+2\log l,
\]
plus $2A_w(j)\log l/l^{2}$ if $l>X_2$. Sum over these primes. Lemma~\ref{is-places-prime-sums}
applies, since $Y>1$ and $X_2\geq k\geq1$ by the sizes: the first terms sum to at most
$2A_1(j)\log Y$, the terms $2\log l$ to at most $2Y\log4$, and the terms for $l>X_2$ to at most
$2.04\cdot2A_w(j)/X_2$. The hypotheses on $\vartheta$ and $\alpha$
bound the second and third terms by $2A_w(j)C/k$ and $2Cmb$.

(4) By (1) of the page estimates, $g_0,g_1,q,d$ are distinct and not in $S_m$, so
\[
 \sum_{l\text{ bad}}\mathrm{err}_P(j,l)\log l=(4T+4fr_0^{2})\log(g_0g_1)
 +2(s_B(j)+T+1)\log(qd)+4fkb'\log d+(2s_C(j)+2T+2)\sum_{l\in S_m}\log l.
\]
As $\log q\geq0$ and $fkb'\geq0$, $4fkb'\log d\le4fkb'\log(qd)$. By (4) of the page estimates the
sum is therefore at most $4b+C_{b,1}fmb+C_{b,2}fmb\log b$. Since $f\geq8$, $m\geq1$ and
$\log b\geq1$, we have $fm\log b\geq1$, so $4b\le4fmb\log b$ and
$C_{b,1}fmb\le C_{b,1}fmb\log b$. The sum is at most $(C_{b,1}+C_{b,2}+4)fmb\log b=C_bfmb\log b$.
\end{proof}

\begin{lemma}\label{is-places-absorb}
Let $C,C_b,L,\xi,f,m,b,k,w_s,S_e,Y,X_2,B$ be real numbers with $C\geq1$, $C_b\geq0$,
$1\le L\le\xi$, $f\geq1$, $m\geq1$, $b>0$, $k\geq1$, $\log k\le L$, $0\le w_s\le2fb$, $S_e\geq0$,
$kS_eL\le mb$, $1<Y\le mb$, $\log Y\le2L$, $k\le X_2$ and $B\le C_bfmbL$. Then
\[
 2\Big(C\big(mw_s\log k+fkS_e\big)\log Y+C^{2}w_s+Cmb+\frac{2.04\,Ckw_s}{X_2}+Y\log4\Big)+B+12fmb
 \le\big(4C^{2}+25C+C_b+12\big)\,\xi\,fmbL.
\]
\end{lemma}
\leanref{Formalism/MirabolicMackeyTheory/MirabolicMackeyProofs/AbsorbLocalConditionsOfExact.lean}{Formalism.EigenvalueLevelCeiling.absorb_local_conditions_exact}
\leanref{Formalism/MirabolicMackeyTheory/ErrSumGoodExact.lean}{Formalism.EigenvalueLevelCeiling.AbsorbOnExact}
\begin{proof}
Write $Z=fmb$. First the terms in the bracket, doubled. Since $\log Y\le2L$, the first term gives at
most $4CLmw_s\log k$, and $2CfkS_e\log Y\le4Cf\cdot kS_eL\le4CZ$. Since $w_s\le2fb$ and $m\geq1$,
$2C^{2}w_s\le4C^{2}Z$. Since $f\geq1$, $2Cmb\le2CZ$. Since $k\le X_2$, $2\cdot2.04\,Ckw_s/X_2\le
4.08\,Cw_s\le8.16\,CZ$. Since $\log4<1.3863$, $Y\le mb$ and $Cf\geq1$, $2Y\log4\le2.7726\,CZ$. As
$4+2+8.16+2.7726\le17$, the doubled bracket is at most $4CLmw_s\log k+(4C^{2}+17C)Z$.

Now $\log k\le L$ and $w_s\le2fb$ give $4CLmw_s\log k\le8CZL^{2}\le8C\xi ZL$. Since $\xi L\geq1$,
$(4C^{2}+17C)Z\le(4C^{2}+17C)\xi ZL$ and $12Z\le12\xi ZL$, and since $\xi\geq1$,
$B\le C_bZL\le C_b\xi ZL$. Adding gives the bound.
\end{proof}

\begin{lemma}\label{is-places-index-subspace}
Let $K$ be a number field, $w$ a finite place of $K$ with valuation $v_w$, and $Z$ a
finite-dimensional $K$-vector space, with $\operatorname{idx}_w$, localizations $A_{(w)}$ and
lattices as in Definition~\ref{is-onto-record-numbers}. Let $A$ and $B$ be lattices in $Z$.
\begin{enumerate}
\item For every $c\in K^\times$, $A\cap cB$ is a lattice in $Z$.
\item Suppose that $B\subset A$. Let $a\in K^\times$ with $v_w(a)=E$ for an integer $E\geq0$, and
  let $r$ be a real number. Then $\operatorname{idx}_w(A\cap W,B\cap W)\le E\dim W+r$ for every
  subspace $W\subset Z$ if and only if $\operatorname{idx}_w(A,A\cap a^{-1}B)\le r$.
\end{enumerate}
\end{lemma}
\leanref{Formalism/MirabolicMackeyTheory/StatementL.lean}{Formalism.EigenvalueLevelCeiling.LatInter}
\leanref{Formalism/MirabolicMackeyTheory/MirabolicMackeyProofs/IdxRules.lean}{Formalism.EigenvalueLevelCeiling.lattice_intersection}
\leanref{Formalism/MirabolicMackeyTheory/StatementL.lean}{Formalism.EigenvalueLevelCeiling.ThreeForms}
\leanref{Formalism/MirabolicMackeyTheory/MirabolicMackeyProofs/ThreeForms.lean}{Formalism.EigenvalueLevelCeiling.three_forms}
\leanref{Formalism/MirabolicMackeyTheory/MirabolicMackeyProofs/ThreeForms.lean}{Formalism.EigenvalueLevelCeiling.threeForms_of}
\begin{proof}
(1) The module $cB$ is finitely generated and spans $Z$, since $B$ does. The module $A\cap cB$ is
finitely generated, as a submodule of the finitely generated module $A$ over the noetherian ring
$\mathcal O_K$. Let $z\in Z$. Since $A$ and $cB$ span $Z$, there are nonzero $d,d'\in\mathcal O_K$
with $dz\in A$ and $d'z\in cB$, and then $dd'z\in A\cap cB$. So $A\cap cB$ spans $Z$.

(2) Write $\mathcal O$ for the localization of $\mathcal O_K$ at $w$, a discrete valuation ring with
valuation $v_w$ and a uniformizer $\pi$, and put $A'=A_{(w)}$, $B'=B_{(w)}$ and $C'=a^{-1}B'$.
Lengths below are lengths of $\mathcal O$-modules. Localization at $w$ commutes with finite
intersections, with multiplication by $a^{-1}$ and with intersection with a subspace, and for
lattices $Y\subset X$ in a subspace, $\operatorname{idx}_w(X,Y)$ is the length of
$X_{(w)}/Y_{(w)}$. Hence, for every subspace $W\subset Z$,
\[
 \operatorname{idx}_w(A,A\cap a^{-1}B)=\operatorname{len}\big(A'/(A'\cap C')\big),\qquad
 \operatorname{idx}_w(A\cap W,B\cap W)=\operatorname{len}\big((A'\cap W)/(B'\cap W)\big).
\]
Since $v_w(a)=E\geq0$, $a\in\mathcal O$, so $B'\subset C'$. For a subspace $W$, the module
$B\cap W$ is a lattice in $W$: it is finitely generated, as a submodule of $B$, and for $z\in W$
some nonzero $d\in\mathcal O_K$ has $dz\in B\cap W$. So its localization $B'\cap W$ is free of rank
$\dim W$ over $\mathcal O$. As $C'\cap W=a^{-1}(B'\cap W)$, the quotient
$(C'\cap W)/(B'\cap W)$ is isomorphic to $(\mathcal O/a\mathcal O)^{\dim W}$, and
\[
 \operatorname{len}\big((C'\cap W)/(B'\cap W)\big)=E\dim W .
\]

\emph{The subspace bound gives the bound for $A\cap a^{-1}B$.} The $\mathcal O$-module $A'$ is free
of rank $n=\dim Z$, and $B'\subset A'$ has the same rank. By the elementary divisor theorem over
$\mathcal O$, there are a basis $x_1,\dots,x_n$ of $A'$ and integers $m_1,\dots,m_n\geq0$ such that
$\pi^{m_1}x_1,\dots,\pi^{m_n}x_n$ is a basis of $B'$. Let $W$ be the $K$-span of the $x_i$ with
$m_i>E$. Then:
\begin{itemize}
\item $A'\subset(A'\cap W)+(A'\cap C')$. Indeed, the $x_i$ with $m_i>E$ lie in $A'\cap W$, and if
  $m_i\le E$, then $ax_i\in\pi^{m_i}\mathcal Ox_i\subset B'$, so $x_i\in C'$.
\item $C'\cap W\subset A'$. Indeed, let $x=\sum_{m_i>E}c_ix_i$ with $c_i\in K$ and $ax\in B'$. Then
  $v_w(ac_i)\geq m_i$, so $v_w(c_i)\geq m_i-E>0$, and $x\in A'$.
\end{itemize}
By the first, $A'\cap W$ maps onto $A'/(A'\cap C')$; the kernel is $A'\cap W\cap C'$, which is
$C'\cap W$ by the second. So
$\operatorname{len}(A'/(A'\cap C'))\le\operatorname{len}((A'\cap W)/(C'\cap W))$. By the second,
$B'\cap W\subset C'\cap W\subset A'\cap W$, so
\[
 \operatorname{len}\big((A'\cap W)/(B'\cap W)\big)
 =E\dim W+\operatorname{len}\big((A'\cap W)/(C'\cap W)\big)
 \geq E\dim W+\operatorname{idx}_w(A,A\cap a^{-1}B).
\]
The left side is $\operatorname{idx}_w(A\cap W,B\cap W)\le E\dim W+r$, so
$\operatorname{idx}_w(A,A\cap a^{-1}B)\le r$.

\emph{The bound for $A\cap a^{-1}B$ gives the subspace bound.} Let $W\subset Z$ be a subspace.
Since $B'\cap W\subset A'\cap C'\cap W\subset A'\cap W$,
\[
 \operatorname{len}\big((A'\cap W)/(B'\cap W)\big)
 =\operatorname{len}\big((A'\cap C'\cap W)/(B'\cap W)\big)
 +\operatorname{len}\big((A'\cap W)/(A'\cap C'\cap W)\big).
\]
The first term is at most $\operatorname{len}((C'\cap W)/(B'\cap W))=E\dim W$. The second is at
most $\operatorname{len}(A'/(A'\cap C'))\le r$, since $(A'\cap W)/(A'\cap C'\cap W)$ embeds in
$A'/(A'\cap C')$. So $\operatorname{idx}_w(A\cap W,B\cap W)\le E\dim W+r$.
\end{proof}

\begin{lemma}\label{is-places-index-projector}
Let $K$ be a number field and $Z$ a finite-dimensional $K$-vector space. Let $P$ be a finite group
acting linearly on $Z$ by $\rho$, let $U\subset P$ be a subgroup and $\psi\colon P\to\mathcal O_K$ a
function, and suppose that $(U,\psi)$ is a \emph{Mackey pair}: $\psi(1)=1$,
$\psi(uu')=\psi(u)\psi(u')$ for $u,u'\in U$, and for every $h\in P\setminus U$ there is $u\in U$ with
$h^{-1}uh\in U$ and $\psi(h^{-1}uh)\ne\psi(u)$. Put
$V_\psi=\{z\in Z:\rho(u)z=\psi(u)z\text{ for all }u\in U\}$ and let $\epsilon$ be the endomorphism
\[
 \epsilon=\sum_{u\in U}\psi(u^{-1})\rho(u)
\]
of $Z$. Then:
\begin{enumerate}
\item $\epsilon(Z)\subset V_\psi$;
\item $\epsilon(z)=|U|z$ for every $z\in V_\psi$;
\item $\epsilon(L)\subset L$ for every $\mathcal O_K$-submodule $L\subset Z$ stable under $\rho(P)$;
\item $\epsilon(\rho(h)z)=0$ for every $h\in P\setminus U$ and $z\in V_\psi$.
\end{enumerate}
\end{lemma}
\leanref{Formalism/MirabolicMackeyTheory/MirabolicMackeyProofs/LemmaGG.lean}{Formalism.EigenvalueLevelCeiling.ggE}
\leanref{Formalism/MirabolicMackeyTheory/MirabolicMackeyProofs/LemmaGG.lean}{Formalism.EigenvalueLevelCeiling.ggE_mem_eig}
\leanref{Formalism/MirabolicMackeyTheory/MirabolicMackeyProofs/LemmaGG.lean}{Formalism.EigenvalueLevelCeiling.ggE_eig}
\leanref{Formalism/MirabolicMackeyTheory/MirabolicMackeyProofs/LemmaGG.lean}{Formalism.EigenvalueLevelCeiling.ggE_stable}
\leanref{Formalism/MirabolicMackeyTheory/StatementL.lean}{Formalism.EigenvalueLevelCeiling.OrthOfMackey}
\leanref{Formalism/MirabolicMackeyTheory/MirabolicMackeyProofs/OrthogonalityOfMackey.lean}{Formalism.EigenvalueLevelCeiling.orthogonality_of_mackey}
\begin{proof}
For $u\in U$, $\psi(u^{-1})\psi(u)=\psi(1)=1$.

(1) Let $u_0\in U$ and $x\in Z$. Then $\rho(u_0)\epsilon(x)=\sum_{u\in U}\psi(u^{-1})\rho(u_0u)x$.
Substitute $u=u_0^{-1}u'$: since $\psi(u'^{-1}u_0)=\psi(u'^{-1})\psi(u_0)$, the sum is
$\psi(u_0)\epsilon(x)$. So $\epsilon(x)\in V_\psi$.

(2) For $z\in V_\psi$ and $u\in U$, $\psi(u^{-1})\rho(u)z=\psi(u^{-1})\psi(u)z=z$. Summing over
$U$ gives $|U|z$.

(3) Each term $\psi(u^{-1})\rho(u)$ maps $L$ into $L$, since $\psi(u^{-1})\in\mathcal O_K$.

(4) Choose $u_0\in U$ with $h^{-1}u_0h\in U$ and $\psi(h^{-1}u_0h)\ne\psi(u_0)$, and put
$S=\epsilon(\rho(h)z)$. For $y\in Z$, substituting $u=u'u_0^{-1}$ in
$\epsilon(\rho(u_0)y)=\sum_{u\in U}\psi(u^{-1})\rho(uu_0)y$ and using
$\psi(u_0u'^{-1})=\psi(u_0)\psi(u'^{-1})$ gives $\epsilon(\rho(u_0)y)=\psi(u_0)\epsilon(y)$. With
$y=\rho(h)z$ this is $\epsilon(\rho(u_0)\rho(h)z)=\psi(u_0)S$. On the other hand
$\rho(u_0)\rho(h)z=\rho(h)\rho(h^{-1}u_0h)z=\psi(h^{-1}u_0h)\rho(h)z$, as $z\in V_\psi$, so
$\epsilon(\rho(u_0)\rho(h)z)=\psi(h^{-1}u_0h)S$. Hence $(\psi(u_0)-\psi(h^{-1}u_0h))S=0$, and
$S=0$.
\end{proof}

\begin{lemma}\label{is-places-index-eigen}
Let $K$ be a number field, $w$ a finite place of $K$ with valuation $v_w$, and $Z$ a
finite-dimensional $K$-vector space, with $\operatorname{idx}_w$ and lattices as in
Definition~\ref{is-onto-record-numbers}. Let $P$ be a finite group acting linearly on $Z$ by $\rho$,
let $U\subset P$ be a subgroup and $\psi\colon P\to\mathcal O_K$ a function such that $(U,\psi)$ is
a Mackey pair (Lemma~\ref{is-places-index-projector}), and put
$V_\psi=\{z\in Z:\rho(u)z=\psi(u)z\text{ for all }u\in U\}$. Let $B'\subset B$ be lattices in $Z$
stable under $\rho(P)$. Then
\[
 [P:U]\operatorname{idx}_w(B\cap V_\psi,B'\cap V_\psi)\le\operatorname{idx}_w(B,B')
 +\dim Z\cdot v_w(|U|).
\]
\end{lemma}
\leanref{Formalism/MirabolicMackeyTheory/StatementL.lean}{Formalism.EigenvalueLevelCeiling.LemmaGG}
\leanref{Formalism/MirabolicMackeyTheory/MirabolicMackeyProofs/LemmaGG.lean}{Formalism.EigenvalueLevelCeiling.good_index_GG}
\leanref{Formalism/MirabolicMackeyTheory/MirabolicMackeyProofs/LemmaGG.lean}{Formalism.EigenvalueLevelCeiling.lemmaGG_of}
\begin{proof}
Write $W=V_\psi$, $N=|U|$, $k=[P:U]$, and $\epsilon$ for the endomorphism of
Lemma~\ref{is-places-index-projector}. Choose representatives $h_1,\dots,h_k$ of the left cosets of
$U$ in $P$, so that $h_i^{-1}h_j\in U$ only if $i=j$.

\emph{Step 1.} For $z\in W$, $\epsilon(\rho(h_i)^{-1}\rho(h_j)z)$ is $Nz$ if $i=j$ and $0$
otherwise. Indeed, $\rho(h_i)^{-1}\rho(h_j)=\rho(h_i^{-1}h_j)$. For $i=j$ this is the identity, and
Lemma~\ref{is-places-index-projector}(2) applies. For $i\ne j$, $h_i^{-1}h_j\notin U$, and
Lemma~\ref{is-places-index-projector}(4) applies.

\emph{Step 2.} $k\dim W\le\dim Z$. The linear map $W^k\to Z$, $(z_j)\mapsto\sum_j\rho(h_j)z_j$, is
injective: applying $\epsilon\rho(h_i)^{-1}$ to the image of $(z_j)$ gives $Nz_i$ by Step 1, and
$N\ne0$ in $K$.

\emph{Step 3.} The modules $B\cap W$ and $B'\cap W$ are lattices in $W$: they are finitely
generated, as submodules of $B$, and for $z\in W$ some nonzero $d\in\mathcal O_K$ has $dz\in B'$,
so $dz\in B'\cap W\subset B\cap W$. Let $\mathcal O$ be the localization of $\mathcal O_K$ at $w$, and take lengths
of $\mathcal O$-modules. Put $A=(B\cap W)_{(w)}$ and $A'=(B'\cap W)_{(w)}$, which lie in $W$, and
\[
 Q_A=A/A',\qquad Q_B=B_{(w)}/B'_{(w)}.
\]
Since $B'\subset B$, $\operatorname{idx}_w(B,B')=\operatorname{len}Q_B$ and
$\operatorname{idx}_w(B\cap W,B'\cap W)=\operatorname{len}Q_A$, and both are finite. Let
$s_i=\rho(h_i)$ and $t_i=\epsilon\rho(h_i)^{-1}$. Since $B$ and $B'$ are stable under $\rho(P)$,
$s_i$ maps $A$ into $B_{(w)}$ and $A'$ into $B'_{(w)}$, and by
Lemma~\ref{is-places-index-projector}(1) and~(3), $t_i$ maps $B_{(w)}$ into $A$ and $B'_{(w)}$ into
$A'$. So they induce $\mathcal O$-linear maps $\bar s_i\colon Q_A\to Q_B$ and
$\bar t_i\colon Q_B\to Q_A$. By Step 1, $\bar t_i\bar s_j$ is multiplication by $N$ if $i=j$ and
$0$ otherwise.

\emph{Step 4.} $k\operatorname{len}(NQ_A)\le\operatorname{len}Q_B$. The image of the map
$\bar t=(\bar t_i)_i\colon Q_B\to Q_A^k$ contains $(NQ_A)^k$: for $y_1,\dots,y_k\in Q_A$,
$\bar t(\sum_j\bar s_jy_j)=(Ny_i)_i$ by Step 3. So
$k\operatorname{len}(NQ_A)=\operatorname{len}((NQ_A)^k)\le\operatorname{len}\bar t(Q_B)
\le\operatorname{len}Q_B$.

\emph{Step 5.} The quotient $Q_A/NQ_A=A/(A'+NA)$ is a quotient of $A/NA$. Since $A$ is the
localization of the lattice $B\cap W$ in $W$, it is free of rank $\dim W$ over $\mathcal O$, so
$\operatorname{len}(A/NA)=v_w(N)\dim W$. Hence
\[
 \operatorname{len}Q_A=\operatorname{len}(NQ_A)+\operatorname{len}(Q_A/NQ_A)
 \le\operatorname{len}(NQ_A)+v_w(N)\dim W .
\]
Multiplying by $k$ and using Step 4, then Step 2 and $v_w(N)\geq0$,
\[
 k\operatorname{idx}_w(B\cap W,B'\cap W)\le\operatorname{idx}_w(B,B')+k\dim W\cdot v_w(N)
 \le\operatorname{idx}_w(B,B')+\dim Z\cdot v_w(N).
\]
\end{proof}

\begin{lemma}\label{is-places-local-index}
Let $K$ be a number field, $w$ a finite place of $K$, and $Z$ a finite-dimensional $K$-vector space,
with $\operatorname{idx}_w$, localizations $A_{(w)}$ and lattices as in
Definition~\ref{is-onto-record-numbers}. Let $P$ be a finite group acting linearly on $Z$ by $\rho$,
let $U\subset P$ be a subgroup and $\psi\colon P\to\mathcal O_K$ a function, and suppose that
$(U,\psi)$ is a \emph{Mackey pair}: $\psi(1)=1$, $\psi(uu')=\psi(u)\psi(u')$ for $u,u'\in U$, and
for every $h\in P\setminus U$ there is $u\in U$ with $h^{-1}uh\in U$ and $\psi(h^{-1}uh)\ne\psi(u)$.
Put $V_\psi=\{z\in Z:\rho(u)z=\psi(u)z\text{ for all }u\in U\}$. Let $e$ be an idempotent
endomorphism of $Z$ with image $V_0$, and put $Z_g=V_\psi\cap V_0$. Let $M_s$, $\Theta$, $\Theta_K$,
$\Theta_1$, $\Theta_0$, $\Theta_J$ and $X_3$ be $\mathcal O_K$-submodules of $Z$, let
$\sigma,a\in K^\times$, let $E\geq0$ be an integer and let $r,\varepsilon$ be real numbers. Let $X$
be a finite-dimensional $K$-vector space with a linear map $\theta'\colon X\to Z$, endomorphisms
$e_X$ and $\mathrm{pr}$, and an $\mathcal O_K$-submodule $M_t$, and put $\mathrm{Im}=e\theta'(M_t)$.
Let $\varphi\colon K^{D'}\to Z$ be an injective linear map, and let $M$ and $T$ be
$\mathcal O_K$-submodules of $K^{D'}$. Suppose:
\begin{enumerate}
\item $M_s$, $\Theta$, $\Theta_K$ and $\Theta_1$ are lattices in $Z$, $M_s$ and $\Theta$ are stable
  under $\rho(P)$, $\Theta_1\subset\Theta\subset M_s$ and $\Theta\subset\sigma\Theta_K$;
\item $v_w(a)=E$, and $\operatorname{idx}_w(M_s\cap W,\Theta_1\cap W)\le E\dim W+r$ for every
  subspace $W\subset Z$;
\item $\Theta_0$ is finitely generated and contained in $V_\psi$,
  $|U|\cdot(\Theta_K\cap V_\psi)\subset\Theta_0\subset\Theta_J$, and $e(\Theta_0)$ and $e(\Theta_J)$
  are lattices in $Z_g$;
\item $\theta'e_X=e\theta'$, and $\theta'(x)=\theta'(\mathrm{pr}(x))$ for every $x$ in the image of
  $e_X$;
\item $X_3$ and $\mathrm{Im}$ are lattices in $Z_g$, $X_{3,(w)}\subset\mathrm{Im}_{(w)}$, and
  $\operatorname{idx}_w(M_s\cap Z_g,X_3)\le\operatorname{idx}_w(M_s\cap Z_g,e(\Theta_J))+\varepsilon$;
\item $\varphi(K^{D'})=Z_g$, $\varphi(M)=M_s\cap Z_g$ and $\varphi(T)=\theta'(\mathrm{pr}(e_X(M_t)))$.
\end{enumerate}
Then
\[
 \operatorname{idx}_w(M,T)\le\big(E-v_w(\sigma)+v_w(|U|)\big)D'
 +\frac{r+\dim Z\cdot v_w(|U|)}{[P:U]}+\varepsilon.
\]
\end{lemma}
\leanref{Formalism/LocalToGlobal/IndefiniteForms/IndefiniteTree/ExampleSatisfiesConditionsAtFive.lean}{Formalism.L2.Indef.Tree.local_bound}
\leanref{Formalism/MirabolicMackeyTheory/MirabolicMackeyProofs/LocalConditionsWeightedConditional.lean}{Formalism.EigenvalueLevelCeiling.local_conditions_weighted_at_place}
\leanref{Formalism/MirabolicMackeyTheory/MirabolicMackeyProofs/OnLocal.lean}{Formalism.EigenvalueLevelCeiling.onLocal_of}
\leanref{Formalism/MirabolicMackeyTheory/StatementL.lean}{Formalism.EigenvalueLevelCeiling.IsMackey}
\begin{proof}
For lattices $A,B,B'$ in one subspace, $\operatorname{idx}_w(A,B)$ depends only on $A_{(w)}$ and
$B_{(w)}$, $\operatorname{idx}_w(A,B')=\operatorname{idx}_w(A,B)+\operatorname{idx}_w(B,B')$,
$\operatorname{idx}_w(A,B')\le\operatorname{idx}_w(A,B)$ if $B\subset B'$, and
$\operatorname{idx}_w(B,cB)$ is $v_w(c)$ times the dimension of the subspace for $c\in K^\times$:
compare each with a common sublattice.
Write $\Lambda=M_s\cap Z_g$; by (6), $\dim Z_g=D'$.

(a) Since $\Theta_1\cap W\subset\Theta\cap W$, (2) also holds with $\Theta$ in place of $\Theta_1$.

(b) The module $N=M_s\cap a^{-1}\Theta$ is a lattice in $Z$ by
Lemma~\ref{is-places-index-subspace}(1). It is stable under $\rho(P)$, since $M_s$ and $\Theta$ are
and each $\rho(g)$ commutes with multiplication by $a^{-1}$. By (1), $\Theta\subset M_s$, and
$v_w(a)=E\geq0$ by (2). So (a) and Lemma~\ref{is-places-index-subspace}(2), applied to $M_s$,
$\Theta$ and $a$, give $\operatorname{idx}_w(M_s,N)\le r$.

(c) The lattices $N\subset M_s$ are stable under $\rho(P)$, so Lemma~\ref{is-places-index-eigen}
gives
$[P:U]\operatorname{idx}_w(M_s\cap V_\psi,N\cap V_\psi)\le\operatorname{idx}_w(M_s,N)+\dim Z\cdot
v_w(|U|)$.

(d) Since $N\cap Z_g\subset N\cap V_\psi$ and $N\subset M_s$, the quotient $\Lambda/(N\cap Z_g)$
embeds in $(M_s\cap V_\psi)/(N\cap V_\psi)$, so by (b) and (c)
$\operatorname{idx}_w(\Lambda,N\cap Z_g)\le(r+\dim Z\cdot v_w(|U|))/[P:U]$. The module
$\Theta_0\cap V_0$ is a lattice in $Z_g$: it is finitely generated, lies in $V_\psi\cap V_0$, and
contains $|U|(\Theta_K\cap Z_g)$, which spans $Z_g$. If $z\in N\cap Z_g$, then $az\in\Theta\subset
\sigma\Theta_K$, so $a\sigma^{-1}z\in\Theta_K\cap V_\psi$ and $|U|a\sigma^{-1}z\in\Theta_0\cap V_0$
by (3). Hence $|U|a\sigma^{-1}(N\cap Z_g)\subset\Theta_0\cap V_0$, and
\[
 \operatorname{idx}_w(\Lambda,\Theta_0\cap V_0)\le\operatorname{idx}_w\big(\Lambda,|U|a\sigma^{-1}
 (N\cap Z_g)\big)=\operatorname{idx}_w(\Lambda,N\cap Z_g)+\big(v_w(|U|)+E-v_w(\sigma)\big)D'.
\]

(e) Since $e$ is the identity on $V_0$, $\Theta_0\cap V_0\subset e(\Theta_0)\subset e(\Theta_J)$. So
$\operatorname{idx}_w(\Lambda,e(\Theta_J))\le\operatorname{idx}_w(\Lambda,\Theta_0\cap V_0)$, and by
(5), $\operatorname{idx}_w(\Lambda,X_3)\le\operatorname{idx}_w(\Lambda,\Theta_0\cap V_0)+\varepsilon$.

(f) By (5), $X_3\cap\mathrm{Im}$ is a lattice in $Z_g$ with the same localization at $w$ as $X_3$.
So $\operatorname{idx}_w(\Lambda,\mathrm{Im})\le\operatorname{idx}_w(\Lambda,X_3\cap\mathrm{Im})=
\operatorname{idx}_w(\Lambda,X_3)$.

(g) For $m\in M_t$, (4) gives $\theta'(\mathrm{pr}(e_Xm))=\theta'(e_Xm)=e\theta'(m)$, so
$\theta'(\mathrm{pr}(e_X(M_t)))=\mathrm{Im}$. By (6), the injective map $\varphi$ carries $M$ and $T$
onto $\Lambda$ and $\mathrm{Im}$, so $\operatorname{idx}_w(M,T)=\operatorname{idx}_w(\Lambda,
\mathrm{Im})$.

Combining (g), (f), (e) and (d) gives the bound.
\end{proof}

\begin{lemma}\label{is-places-combination}
Let $S$ be an indefinite-side datum over a number field $K$ with twisted places $\mathrm{Tw}$, tuple
map $w\mapsto\mu_w$, lattices $M_j$ and $T_{j,t}$, matrices $\theta_{j,t}$ and rank $D\geq1$
(Definition~\ref{is-indefinite-datum}), and suppose that $d_W(j)=D$ for every lift $j$. Let $N$ be
numerical bounds for $S$ with functions $F_0,J_0,E_1$ (Definition~\ref{is-numerical-bounds}). Let $O$
be a record object for $S$, a finite group $P$ and constants $c=(C,C_{b,1},C_{b,2})$, with the
notation of Definitions~\ref{is-record-object}, \ref{is-onto-record-numbers}
and~\ref{is-places-error-terms}, and let $C_E\geq4C^2+25C+C_{b,1}+C_{b,2}+48$ be a real number. In
(3) and (4), $v$ runs over the twisted places $w$ and their conjugates $\bar w$, $j$ over the lifts,
and $e_X$ and $\mathrm{pr}$ are the endomorphisms of $O$ at $(v,j)$. Suppose:
\begin{enumerate}
\item $(U,\psi)$ is a Mackey pair (Lemma~\ref{is-places-local-index}), and each $e_j$ is an idempotent
  with image $V_0(j)$;
\item $(1-y)|U|\le[P:U]$, $r_0\geq1$, and for every twisted place $w$,
  $v_w(|U|)=e_wfr_0(r_0-1)/2$ if $p_w\in\{g_0,g_1\}$ and $v_w(|U|)=0$ otherwise;
\item $\theta'_{v,j}e_X=e_j\theta'_{v,j}$ and $\theta'_{v,j}(x)=\theta'_{v,j}(\mathrm{pr}(x))$ for
  every $x$ in the image of $e_X$; and $X_3(v,j)$ and $\mathrm{Im}(v,j)=e_j\theta'_{v,j}(M_t(v,j))$ are
  lattices in $Z_g(j)$ with $X_3(v,j)_{(v)}\subset\mathrm{Im}(v,j)_{(v)}$;
\item $\varphi_j$ is injective with image $Z_g(j)$, $\varphi_j(M_j)=M_s(j)\cap Z_g(j)$, and
  $\varphi_j(\theta_{j,\mu_w}T_{j,\mu_w})=\theta'_{v,j}(\mathrm{pr}(e_X(M_t(v,j))))$ for
  $v\in\{w,\bar w\}$;
\item $O$ satisfies (1) to (4) of the page estimates (Definition~\ref{is-onto-page-estimates}), the
  geometric estimates (Definition~\ref{is-onto-geometric-estimates}) and the depth estimates for $N$
  (Definition~\ref{is-onto-depth-estimates}), and has the rank at the twisted places, a twist bound,
  the share bound, the bad-place bound and the exponent bound for $N$
  (Definition~\ref{is-onto-local-estimates});
\item $E_1(j)=C_EE_1^{u}(O)$ for every lift $j$.
\end{enumerate}
Then there are real numbers $\mathrm{err}(w,j)$ such that $S$ satisfies (1), (2) and (3) of
Definition~\ref{is-surjectivity-bound} for $N$ with $\mathrm{err}$ in place of $\mathrm{err}_0$.
\end{lemma}
\leanref{Development/AbelianVarietyTheory/Proofs/DivisorProofs/TwistedPlacesBoundChainTwoRunChain.lean}{GlobalTheory.OnAtPlacesBase.on_of_has_atPlaces}
\leanref{Development/AbelianVarietyTheory/Proofs/FoundationProofs/TraceLatticeLeThetaLatticeAtPlace.lean}{GlobalTheory.TraceClassRank.on_core_atPlaces}
\leanref{Formalism/LocalToGlobal/IndefiniteForms/IndefiniteTree/ExampleSatisfiesConditionsAtFive.lean}{Formalism.L2.Indef.Tree.on_core}
\leanref{Formalism/LocalToGlobal/IndefiniteForms/IndefiniteTree/ExampleSatisfiesConditionsAtFive.lean}{Formalism.L2.Indef.Tree.OnObj.errOn}
\leanref{Formalism/MirabolicMackeyTheory/MirabolicMackeyProofs/LocalConditionsGlobalForm.lean}{Formalism.EigenvalueLevelCeiling.local_conditions_global_form}
\leanref{Formalism/MirabolicMackeyTheory/MirabolicMackeyProofs/UnipotentNumeric.lean}{Formalism.EigenvalueLevelCeiling.share_numeric}
\leanref{Formalism/MirabolicMackeyTheory/MirabolicMackeyProofs/UnipotentNumeric.lean}{Formalism.EigenvalueLevelCeiling.depth_cases}
\leanref{Formalism/MirabolicMackeyTheory/MirabolicMackeyProofs/UnipotentNumeric.lean}{Formalism.EigenvalueLevelCeiling.budget_numeric}
\begin{proof}
The conjugation of $K$ carries $w$ to $\bar w$, so $e_{\bar w}=e_w$, $p_{\bar w}=p_w$,
$\ell_{\bar w}=\ell_w$, and $v_{\bar w}(n)=v_w(n)$ for every integer $n$.

\emph{Step 0: the data of the estimates.} Replace $\Theta_1$, $a$, $r_1$, $r_2$ and $\delta_2$ of $O$
by $\Theta_1'$, $a'$, $r_1'$, $r_2'$ and $\delta_2'$ of the index bound (geometric estimates (5)),
$\varepsilon_0$ by the numbers $\varepsilon'$ of the twist bound, $J_p$ by the numbers $J'$ of the
exponent bound, and $\Theta_J(v,j)$ by $\Theta_J(v,j)+\Theta_0(v,j)$. The hypotheses still hold: by
the rank at the twisted places, $e_j(\Theta_J+\Theta_0)=e_j\Theta_J+e_j\Theta_0$ is a lattice in
$Z_g(j)$, and as $\Theta_0(v,j)_{(v)}\subset\Theta_J(v,j)_{(v)}$ it has the same localization at $v$
as $e_j\Theta_J(v,j)$, so the twist bound is unchanged. Now $\Theta_0(v,j)\subset\Theta_J(v,j)$, the
index bound holds with the data of $O$, $\varepsilon_0(v,j)\geq0$ satisfies the twist bound, and
$J_0(\mu_w,w,j)\le J_p(j,p_w)$ with $J_p\geq0$ and $\sum_{l\in Q}J_p(j,l)\log l\le E_J$. Put
\[
 \mathrm{err}(w,j)=\mathrm{err}_P(j,p_w)+\frac{r(w,j)/e_w+r(\bar w,j)/e_{\bar w}}{[P:U]\,D}.
\]

\emph{Step 1: the local index bounds.} Let $w\in\mathrm{Tw}$, $v\in\{w,\bar w\}$ and $j$ a lift.
Lemma~\ref{is-places-local-index} applies at $v$ to $Z=K^{n_Z(j)}$ with $\rho_j$, $U$, $\psi$, $e_j$,
the modules $M_s(j)$, $\Theta(v,j)$, $\Theta_K(v,j)$, $\Theta_1(v,j)$, $\Theta_0(v,j)$,
$\Theta_J(v,j)$, $X_3(v,j)$, the numbers $\sigma(v,j)$, $a(v,j)$, $E(v,j)$, $r(v,j)$,
$\varepsilon=e_v\varepsilon_0(v,j)$, the maps $\theta'_{v,j}$, $e_X$, $\mathrm{pr}$, the module
$M_t(v,j)$, and $\varphi_j$ with $D'=D$, $M=M_j$ and $T=\theta_{j,\mu_w}T_{j,\mu_w}$. Hypothesis (1)
gives the Mackey pair and the idempotent. The hypotheses (1) and (2) of the lemma are (1), (2), (3)
and (5)(a) of the geometric estimates; its (3) is (4) of the geometric estimates, Step 0 and the rank
at the twisted places; its (4) and (5) are hypothesis (3) and the twist bound; and its (6) is
hypothesis (4). Since $e_ve_0(v,j)=E(v,j)-v_v(\sigma(v,j))$ and $e_vv_U(v)=v_v(|U|)$, it gives
\[
 I_v(j,\mu_w)\le e_v\Big(\big(e_0(v,j)+v_U(v)\big)D+\frac{r(v,j)/e_v+n_Z(j)v_U(v)}{[P:U]}
 +\varepsilon_0(v,j)\Big).\tag{$*$}
\]

\emph{Step 2: the share.} For every lift $j$, $n_Z(j)/[P:U]\le2D\le3D$. Indeed $fk\geq8$ by (1) of
the page estimates, so $y=1/(4(fk)^3)\le1/8$. The share bound gives $\epsilon\le y$ with
$(1-\epsilon)n_Z(j)\le|U|d_0(j)$, so $\frac78n_Z(j)\le|U|d_0(j)$; by hypothesis (2),
$\frac78|U|\le[P:U]$; and the bad-place bound gives $\frac34d_0(j)\le D$. Hence
$n_Z(j)\le(\frac87)^2\frac43[P:U]D\le2[P:U]D$.

\emph{Step 3: the order of $U$.} For $w\in\mathrm{Tw}$, $v_U(w)\geq0$ and $v_U(\bar w)=v_U(w)$, since
$|U|$ is an integer. By hypothesis (2), $v_U(w)$ is $fr_0(r_0-1)/2$ if $p_w\in\{g_0,g_1\}$ and $0$
otherwise, so $4v_U(w)\le g_q(w)$.

\emph{Step 4: the depth bound.} For $w\in\mathrm{Tw}$ and every lift $j$,
\[
 e_0(w,j)+e_0(\bar w,j)-2F_0(w,j)+2g_q(w)\le\mathrm{err}_P(j,p_w).\tag{$**$}
\]
If $p_w\in\{g_0,g_1\}$, (4) of the depth estimates gives $e_0(w,j)+e_0(\bar w,j)\le4T-\nu_w(j)
-\nu_{\bar w}(j)=4T+2F_0(w,j)$, and $2g_q(w)=4fr_0^2$. Otherwise $g_q(w)=0$. If $p_w\in S_m$, (3)
of the depth estimates gives $e_0(w,j)+e_0(\bar w,j)\le2s_C(j)+2T+2$ and $F_0(w,j)=0$. If
$p_w\in\{q,d\}$ and $p_w\notin S_m$, (2) of the depth estimates gives $e_0(w,j)+e_0(\bar w,j)\le
2(s_B(j)+T+1)$ and $-2F_0(w,j)\le4\beta(p_w)$, as $q\ne d$. If $p_w$ is not bad, (1) of the depth
estimates gives $e_0(w,j)+e_0(\bar w,j)-2F_0(w,j)\le2a_E(j,p_w)+2$ when $p_w\le Y$, and
$e_0(w,j)+e_0(\bar w,j)-2F_0(w,j)\le0$ when $p_w>Y$. In each case this is $(**)$.

\emph{Step 5: condition (1).} Let $w\in\mathrm{Tw}$ and $j$ a lift. Add $(*)$ for $v=w$ and
$v=\bar w$. By Steps 2 and 3, $2n_Z(j)v_U(w)/[P:U]\le6Dv_U(w)$, so the sum of the right sides is at
most
\[
 e_w\Big(\big(e_0(w,j)+e_0(\bar w,j)+8v_U(w)\big)D+\frac{r(w,j)/e_w+r(\bar w,j)/e_{\bar w}}{[P:U]}
 +\varepsilon_0(w,j)+\varepsilon_0(\bar w,j)\Big).
\]
By Step 3 and $(**)$, $e_0(w,j)+e_0(\bar w,j)+8v_U(w)\le2F_0(w,j)+\mathrm{err}_P(j,p_w)$. So
\[
 I_w(j,\mu_w)+I_{\bar w}(j,\mu_w)\le\big(2e_wF_0(w,j)+e_w\,\mathrm{err}(w,j)\big)D
 +e_w\big(\varepsilon_0(w,j)+\varepsilon_0(\bar w,j)\big).
\]
Sum over $j$. The twist bound gives $\sum_j(\varepsilon_0(w,j)+\varepsilon_0(\bar w,j))\le2/p_w^2$,
and $D\geq1$, so the total is at most
$\big(\sum_j(2e_wF_0(w,j)+e_w\,\mathrm{err}(w,j))+2e_w/p_w^2\big)D$.

\emph{Step 6: condition (2).} By Lemma~\ref{is-places-error-sum}(2), $\mathrm{err}_P\geq0$, and
$r_1,r_2\geq0$ by the index bound, so $\mathrm{err}(w,j)\geq0$.

\emph{Step 7: three sums.} Fix a lift $j$. The page and depth estimates give the hypotheses of
Lemma~\ref{is-places-error-sum}, so by its part (3) for $\mathcal S=\mathrm{Tw}$,
$\sum_{w\in\mathrm{Tw}}\ell_we_w\,\mathrm{err}_P(j,p_w)\le E_A(j)$. For every finite set
$\mathcal S'$ of finite places, $r\le r_1+r_2$ and the index bound with $\delta_2\le1$ give
$\sum_{u\in\mathcal S'}\ell_ur(u,j)\le2n_Z(j)fmb$. By Step 0 and part (1) of
Lemma~\ref{is-places-error-sum},
$\sum_{w\in\mathrm{Tw}}\ell_we_wJ_0(\mu_w,w,j)\le\sum_lJ_p(j,l)\log l\le E_J$.

\emph{Step 8: condition (3).} Fix a lift $j$. As $\ell_{\bar w}=\ell_w$ and $e_{\bar w}=e_w$,
\[
 \sum_{w\in\mathrm{Tw}}\ell_w\big(e_w\,\mathrm{err}(w,j)+e_wJ_0(\mu_w,w,j)\big)
 =\sum_w\ell_we_w\,\mathrm{err}_P(j,p_w)+\frac{\sum_w\ell_wr(w,j)+\sum_w\ell_{\bar w}r(\bar w,j)}
 {[P:U]\,D}+\sum_w\ell_we_wJ_0(\mu_w,w,j).
\]
The map $w\mapsto\bar w$ is injective, so Step 7 applies to $\mathrm{Tw}$ and to its image, and with
Step 2 the middle term is at most $4n_Z(j)fmb/([P:U]D)\le12fmb$. So the sum is at most
$E_A(j)+12fmb+E_J$. Lemma~\ref{is-places-absorb}, with $L=\log b$, $\xi=x^\varepsilon$,
$w_s=\mathrm{wt}(j)$ and $B=\sum_{l\text{ bad}}\mathrm{err}_P(j,l)\log l$, applies by (1), (2) and
(4) of the page estimates, $\log b=x^\varepsilon$ and Lemma~\ref{is-places-error-sum}(4); since
$A_1(j)=C(m\,\mathrm{wt}(j)\log k+fkS_e)$ and $A_w(j)C/k=C^2\mathrm{wt}(j)$, it gives
$E_A(j)+12fmb\le(4C^2+25C+C_b+12)x^\varepsilon fmb\log b$. Since $T_2\le f$ by (3) of the page
estimates, $x^{2\varepsilon}=x^\varepsilon\log b$, $f\geq1$ and $x^\varepsilon\log b\geq1$,
$E_J=30T_2mbx^{2\varepsilon}+2mb\le32x^\varepsilon fmb\log b$. Adding,
\[
 \sum_{w\in\mathrm{Tw}}\ell_w\big(e_w\,\mathrm{err}(w,j)+e_wJ_0(\mu_w,w,j)\big)
 \le\big(4C^2+25C+C_{b,1}+C_{b,2}+48\big)x^\varepsilon fmb\log b\le C_EE_1^{u}(O)=E_1(j),
\]
by the bound assumed for $C_E$, which also gives $C_E\geq0$, and $N_c\geq0$.
\end{proof}

\begin{lemma}\label{is-places-level-index}
Let $r\geq1$ be an integer and $(k_i)_{i\in I}$ a finite family of finite fields. For each $i$ let
$P_i\subset GL_r(k_i)$ be the mirabolic subgroup of matrices whose last row is $(0,\dots,0,1)$ and
$U_i\subset P_i$ the subgroup of upper triangular unipotent matrices, and put $P=\prod_iP_i$ and
$U=\prod_iU_i$. Then:
\begin{enumerate}
\item $|U|=\prod_i|k_i|^{r(r-1)/2}$ and $[P:U]=\prod_i\prod_{a=1}^{r-1}(|k_i|^a-1)$;
\item if $y$ is a real number with $\sum_i1/(|k_i|-1)\le y$, then $(1-y)|U|\le[P:U]$;
\item if $g_0\ne g_1$ are primes and $f\geq0$ an integer such that $f$ of the $k_i$ have $g_0$
  elements and the other $f$ have $g_1$ elements, then $\sum_i1/(|k_i|-1)\le2f/(\min(g_0,g_1)-1)$,
  and for every finite place $w$ of a number field, with residue characteristic $p_w$ and
  ramification index $e_w$, $v_w(|U|)=e_wfr(r-1)/2$ if $p_w\in\{g_0,g_1\}$ and $v_w(|U|)=0$
  otherwise.
\end{enumerate}
\end{lemma}
\leanref{Formalism/LocalToGlobal/IndefiniteForms/IndefiniteTree/ONLeaves.lean}{Formalism.L2.Indef.Tree.leafIndexU_of}
\leanref{Formalism/LocalToGlobal/IndefiniteForms/IndefiniteTree/ONLeaves.lean}{Formalism.L2.Indef.Tree.leafUnip_of}
\leanref{Formalism/PeriodIntegrals/OrbitalBoundsConclusion/IndexUBound.lean}{product_lowerBound}
\begin{proof}
(1) A matrix of $U_i$ has free entries above the diagonal, so $|U_i|=|k_i|^{r(r-1)/2}$. A matrix of
$P_i$ is determined by its upper left block in $GL_{r-1}(k_i)$ and the first $r-1$ entries of its
last column, so $|P_i|=|k_i|^{r-1}|GL_{r-1}(k_i)|=|k_i|^{r(r-1)/2}\prod_{a=1}^{r-1}(|k_i|^a-1)$.

(2) By (1), $[P:U]=|U|\prod_i\prod_{a=1}^{r-1}(1-|k_i|^{-a})$. For numbers $t_s\in[0,1]$,
$\prod_s(1-t_s)\geq1-\sum_st_s$, by induction on their number, and
$\sum_{a\geq1}|k_i|^{-a}=1/(|k_i|-1)$. So the product is at least $1-y$.

(3) There are $2f$ fields, each with at least $\min(g_0,g_1)$ elements. By (1),
$|U|=(g_0g_1)^{fr(r-1)/2}$, and $v_w(g_0)$ is $e_w$ if $p_w=g_0$ and $0$ otherwise, and likewise for
$g_1$.
\end{proof}

\begin{lemma}\label{is-ties-mirabolic-mackey}
Let $r\geq1$ be an integer, $(k_i)_{i\in I}$ a finite family of finite fields and $A$ a commutative
ring. For each $i$ let $P_i\subset GL_r(k_i)$ be the mirabolic subgroup of matrices whose last row is
$(0,\dots,0,1)$, $U_i\subset P_i$ the subgroup of upper triangular unipotent matrices, and
$\chi_i\colon k_i\to A$ a nontrivial additive character: $\chi_i(0)=1$,
$\chi_i(a+b)=\chi_i(a)\chi_i(b)$ and $\chi_i\ne1$. For a matrix $h$ put
$s(h)=\sum_{a=1}^{r-1}h_{a,a+1}$. Put $P=\prod_iP_i$, $U=\prod_iU_i$ and
$\psi(h)=\prod_i\chi_i(s(h_i))$ for $h=(h_i)\in P$. Then $(U,\psi)$ is a Mackey pair
(Lemma~\ref{is-places-local-index}).
\end{lemma}
\leanref{Formalism/MirabolicMackeyTheory/MirabolicMackeyProofs/Mirabolic.lean}{Formalism.EigenvalueLevelCeiling.mirabolic_is_mackey}
\leanref{Formalism/MirabolicMackeyTheory/MirabolicMackeyProofs/Mirabolic.lean}{Formalism.EigenvalueLevelCeiling.mirabolic_mackey_pair}
\leanref{Formalism/MirabolicMackeyTheory/MirabolicMackeyProofs/OnLocal.lean}{Formalism.EigenvalueLevelCeiling.mackey_product}
\leanref{Formalism/MirabolicMackeyTheory/MirabolicMackeyProofs/IsMackeyPiUpperUnipotentProdCharacter.lean}{Formalism.EigenvalueLevelCeiling.isMackey_pi_upper_unipotent_prod_character}
\begin{proof}
If $u$ and $u'$ are upper triangular unipotent, then $(uu')_{a,a+1}=u_{a,a+1}+u'_{a,a+1}$, so
$s(uu')=s(u)+s(u')$. Hence $\psi(1)=1$ and $\psi(uu')=\psi(u)\psi(u')$ for $u,u'\in U$. Every value
$\chi_i(z)$ is invertible, with inverse $\chi_i(-z)$.

First let $I$ have one element, so that $P=P_1$, and write $\chi=\chi_1$. Let $h\in P\setminus U$.
Call a row $a$ of $h$ \emph{unipotent} if $h_{aa}=1$ and $h_{ab}=0$ for $b<a$. The last row is
unipotent as $h\in P$, and since $h\notin U$ some row is not; let $q<r$ be the largest such row. If
$h_{qq}\ne1$ put $l=q$, and otherwise choose $l<q$ with $h_{ql}\ne0$. Then $l\le q$, every row below
$q$ is unipotent, and $c=\delta_{lq}-h_{ql}\ne0$. Let $e_1,\dots,e_r$ be the standard basis of $k^r$,
and for $t\in k$ put
\[
 u=1+t\,he_le_{q+1}^{\mathsf T},\qquad v=1+t\,e_le_{q+1}^{\mathsf T}h.
\]
Then $uh=h+t\,he_le_{q+1}^{\mathsf T}h=hv$, so $h^{-1}uh=v$. The entries of $u-1$ are the
$th_{al}$ in column $q+1$; for $a\geq q+1$ the row $a$ is unipotent and $l<a$, so they vanish, and
$u$ is upper triangular unipotent. The entries of $v-1$ are the $th_{q+1,b}$ in row $l$; for $b\le l$
they vanish because row $q+1$ is unipotent and $b<q+1$, so $v$ is upper triangular unipotent.
Moreover $s(u)=th_{ql}$, and $s(v)=th_{q+1,l+1}=t\delta_{lq}$ since row $q+1$ is unipotent. Choose
$z\in k$ with $\chi(z)\ne1$ and put $t=z/c$. Then $s(v)=s(u)+z$, so
$\psi(h^{-1}uh)=\psi(u)\chi(z)\ne\psi(u)$, as $\psi(u)$ is invertible.

In general let $h\in P\setminus U$, and choose $i$ with $h_i\notin U_i$. The first case gives
$u_i\in U_i$ with $h_i^{-1}u_ih_i\in U_i$ and $\chi_i(s(h_i^{-1}u_ih_i))\ne\chi_i(s(u_i))$. Let $u\in U$
have component $u_i$ at $i$ and $1$ elsewhere. Then $h^{-1}uh\in U$, and since $\chi_{i'}(s(1))=1$
for every $i'$, $\psi(h^{-1}uh)=\chi_i(s(h_i^{-1}u_ih_i))\ne\chi_i(s(u_i))=\psi(u)$.
\end{proof}

\begin{definition}\label{is-ties-class-sums}
Let $R$ be a realization of a definite-side datum $D$, with coefficient field $L$, embedding
$\sigma_0$ of its base place and group $G$, and let $F$ be a face, with subgroup $P_0$ and character
$\psi_0$ (Definition~\ref{is-realisation}). Write $\mathcal X$ for the finite set of characters of $G$
of the family of definite forms, and $e_{\mathcal X}=\sum_{\chi\in\mathcal X}\frac{\chi(1)}{|G|}
\sum_{g\in G}\chi(g^{-1})\,g$ for its idempotent in the group algebra $L[G]$, with coefficients
$c_{\mathcal X}(g)$. Let $j$ be a lift and $t$ a tuple.
\begin{enumerate}
\item The rational structure $FW_j$ of the $j$-th lift (Definition~\ref{is-realisation}(8)) has
  coordinates $L^{n_W(j)}$, an action $\varrho_j$ of $G$, a character $\chi_j$ of $G$ and
  $\sigma$-semilinear maps $j_\sigma\colon L^{n_W(j)}\to H_W(j,\sigma)$. Write $\Lambda_j$ for its
  integral lattice, the $\mathcal O_L$-span of its integral vectors, and $\Phi_j(z)$ for the
  automorphic form on $U(W_j)$ of the vector $j_{\sigma_0}(z)$ of the Hecke module $H_W(j,\sigma_0)$.
\item For $c\colon G\to L$ put $e_c=\sum_gc(g)\chi_j(g)^{-1}\varrho_j(g)$. The \emph{good operator}
  is $e_{\mathrm{good}}=e_{c_{\mathcal X}}$. With $FW_j^0$ the subspace of $L^{n_W(j)}$ on which
  $P_0$ acts through $\chi_j\psi_0^{-1}$, the \emph{good part} is
  $W^{g}_j=FW_j^0\cap e_{\mathrm{good}}(L^{n_W(j)})$. The rows of $R$ at $F$
  (Definition~\ref{is-complex-data}) have $d_W(j)=\dim W^{g}_j$, the coordinate map
  $c_j\colon L^{d_W(j)}\to W^{g}_j$ of a basis of $W^g_j$, and the lattice $M_j=c_j^{-1}(\Lambda_j)$.
\item The \emph{level forms} at $(j,t)$ are the lift forms $f$ at $(j,t)$ with $gf=\psi_0(g)^{-1}f$
  for $g\in P_0$. Write $f_{j,t}\colon L^{s(j,t)}\to\{\text{level forms}\}$ for the coordinate map of a
  basis, $N_{j,t}\subset L^{s(j,t)}$ for the coordinates of the integral level forms, $H_{j,t}$ for
  the matrix of the class-sum pairing in these coordinates, and $e_{\mathcal X,j,t}$ for the matrix of
  $e_{\mathcal X}$, which preserves the level forms because it is central. For the coordinate map
  $c_{0,j,t}\colon L^{d_Z(j,t)}\to\{\text{lift forms}\}$ of the face at level $0$
  (Definition~\ref{is-complex-data}), write $\pi_{j,t}\colon
  L^{s(j,t)}\to L^{d_Z(j,t)}$ for the map such that $c_{0,j,t}(\pi_{j,t}(x))$ is the orthogonal
  projection of $f_{j,t}(x)$, for the class-sum pairing, onto the image of $c_{0,j,t}$. The lattice of
  the rows is $T_{j,t}=\pi_{j,t}(N_{j,t})$.
\item A \emph{class sum} at $(j,t)$ is a linear map $\tau\colon L^{s(j,t)}\to L^{n_W(j)}$ with
  $\Phi_j(\tau(x))=\mathrm{th}_F(j,t)\big(j_V(f_{j,t}(x))\big)$ for every $x$, where $j_V$ is the
  $\sigma_0$-semilinear map from the lift forms at $(j,t)$ to $H_V(j,t)$ and $\mathrm{th}_F(j,t)$ is
  the theta map of Definition~\ref{is-realisation}(7). The realization is \emph{rational at the
  twisted places} if there is a class sum at $(j,\mu_w)$ for every lift $j$ and every twisted place
  $w$.
\item The \emph{theta lattice} $\Xi_{j,t}\subset L^{n_W(j)}$ of the tuple $t$ is the
  $\mathcal O_L$-span of the vectors $z$ for which $\Phi_j(z)$ is the trace
  $\sum_{g\in P_0}\sigma_D(\psi_0(g))\,g\vartheta$, with $\sigma_D$ the embedding of $L$ into $\mathbf C$
  of $D$ (Definition~\ref{is-realisation}), of a theta series $\vartheta$ of the Weil datum
  and the test vector of $(V_j,W_j)$ at $t$ (Definition~\ref{is-realisation}(6)), at an element $h$
  of the finite adelic group of $U(V_j)$ and a coefficient vector that is integral for the lattice of
  $h$.
\end{enumerate}
\end{definition}
\leanref{Formalism/LocalToGlobal/Divisibility/GoodPartHermitianSpec.lean}{Formalism.L2.DIV.idemOf}
\leanref{Formalism/LocalToGlobal/IndefiniteReal/GramCEqTranspose.lean}{Formalism.IndefR.Arith.FormsL.eOp}
\leanref{Formalism/LocalToGlobal/IndefiniteReal/GramCEqTranspose.lean}{Formalism.IndefR.Arith.FormsL.typeP}
\leanref{Formalism/LocalToGlobal/IndefiniteReal/GramCEqTranspose.lean}{Formalism.IndefR.Arith.FormsL.Wgood}
\leanref{Formalism/LocalToGlobal/IndefiniteReal/LinksDefs.lean}{Formalism.IndefR.Links.RealDat.WL}
\leanref{Development/AbelianVarietyTheory/Definitions/AutomorphicMembership/LiftWeightComplexPointEqTwist.lean}{GlobalTheory.OntoChain.goodOp}
\leanref{Development/AbelianVarietyTheory/Definitions/AutomorphicMembership/LiftWeightComplexPointEqTwist.lean}{GlobalTheory.OntoChain.coordMap}
\leanref{Development/AbelianVarietyTheory/Definitions/AutomorphicMembership/LiftWeightComplexPointEqTwist.lean}{GlobalTheory.OntoChain.supLat}
\leanref{Development/AbelianVarietyTheory/Definitions/Foundations/ThetaSeriesTies.lean}{GlobalTheory.OnTies.srcForm}
\leanref{Development/AbelianVarietyTheory/Definitions/AutomorphicMembership/LiftWeightComplexPointEqTwist.lean}{GlobalTheory.OntoChain.srcLat}
\leanref{Development/AbelianVarietyTheory/Definitions/AutomorphicMembership/LiftWeightComplexPointEqTwist.lean}{GlobalTheory.OntoChain.srcGood}
\leanref{Formalism/LocalToGlobal/Divisibility/RowLat.lean}{Formalism.L2.DIV.BlockDat.lPr}
\leanref{Formalism/LocalToGlobal/Divisibility/RowLat.lean}{Formalism.L2.DIV.BlockDat.Tl}
\leanref{Development/AbelianVarietyTheory/Definitions/Foundations/ThetaSeriesTies.lean}{GlobalTheory.OnTies.IsClassSum}
\leanref{Development/AbelianVarietyTheory/Definitions/Foundations/ThetaSeriesTies.lean}{GlobalTheory.OnTies.ClassSumRational}
\leanref{Development/AbelianVarietyTheory/Proofs/DivisorProofs/NotOntoOfCharAlone.lean}{GlobalTheory.RealIndexBound.TwistedRational}
\leanref{Development/AbelianVarietyTheory/Definitions/Foundations/ThetaSeriesTies.lean}{GlobalTheory.OnTies.thetaTuple}
\leanref{Development/AbelianVarietyTheory/Definitions/Foundations/ThetaSeriesTies.lean}{GlobalTheory.OnTies.seriesSet}
\leanref{Development/AbelianVarietyTheory/Definitions/Foundations/ThetaSeriesTies.lean}{GlobalTheory.OnTies.levelTrace}

\begin{definition}\label{is-ties-lattice-systems}
Let $R$ be a realization of a definite-side datum $D$ (Definition~\ref{is-realisation}), with host
field $E_h/F_h$, coefficient field $L$, embedding $\sigma_0$ of its base place and set $\mathcal T$
of embeddings of $E_h$ (Definition~\ref{th-definite-datum}(2)). Let $j$ be a lift and $t$ a tuple,
and put $n=\operatorname{rk}V_j$. Write $\Gamma=U(V_j)(F_h)$ and $G_f=U(V_j)(\mathbf A_{F_h,f})$;
$J\subset G_f$ for the level of the lift forms at $(j,t)$ and $\chi\colon J\to\mathbf C^\times$ for
its character; and $W_c$ for the coefficient space of $V_j$ (Definition~\ref{is-realisation}(3)),
with its action $\rho$ of $U(V_j)(\mathbf A_{F_h})$.
\begin{enumerate}
\item (Classes.) The realization of the definite family by algebraic modular forms presents the lift
  forms at $(j,t)$ through a finite set $Y\subset G_f$ of representatives: every $h\in G_f$ is
  $\gamma yk$ with $\gamma\in\Gamma$ and $k\in J$ for exactly one $y\in Y$. For each $y\in Y$ it gives
  a finite group $\Gamma_y$ with an isomorphism $g\mapsto\gamma_g$ onto the stabilizer
  $\{\gamma\in\Gamma:y^{-1}\gamma y\in J\}$; put $k_g=(y^{-1}\gamma_gy)^{-1}\in J$, so that
  $\gamma_gyk_g=y$. It gives an $L$-vector space $\mathcal V$, the \emph{rational coefficient space},
  with an action $\rho_y$ of each $\Gamma_y$ and $\mathcal O_L$-submodules $\mathcal N_y\subset\mathcal
  V$, the \emph{integral vectors} at $y$; and a $\sigma_0$-semilinear map $\mathrm{coef}\colon\mathcal
  V\to W_c$ that carries a basis of $\mathcal V$ to a basis of $W_c$ and satisfies
  $\mathrm{coef}(\rho_y(g)x)=\chi(k_g)\,\rho(\gamma_g)\,\mathrm{coef}(x)$. The lift forms at $(j,t)$
  are the families $f=(f(y))_{y\in Y}$ of vectors $f(y)\in\mathcal V$ fixed by $\Gamma_y$, and $f$ is
  \emph{integral} if $f(y)\in\mathcal N_y$ for every $y$. For $y\in Y$ and $x\in\mathcal V$, the
  \emph{averaged form} $a_y(x)$ is the lift form with value $\sum_{g\in\Gamma_y}\rho_y(g)x$ at $y$
  and $0$ at the other representatives.
\item (Theta series.) For $h\in G_f$ and $w\in W_c$, $\vartheta_t(h,w)$ is the theta series of the
  Weil datum and the test vector of $(V_j,W_j)$ at $t$ (Definition~\ref{is-realisation}(6)) at the
  element of $U(V_j)(\mathbf A_{F_h})$ with finite part $h$ and archimedean component $1$, for the
  polynomial $\mathrm{io}_P(j)(v\otimes w)$ (Definition~\ref{is-realisation}(7)). It is a function of
  the tensors $v$ of lowest $K$-type vectors and of $g\in U(W_j)(\mathbf A_{F_h})$, and it is linear
  in $w$.
\item (Lattice systems.) A \emph{lattice system} for the presentation consists of $L$-linear maps
  $r(\gamma)$ of $\mathcal V$ for $\gamma\in\Gamma$ and $\mathcal O_L$-submodules
  $\Lambda(h)\subset\mathcal V$ for $h\in G_f$ such that
  $\mathrm{coef}(r(\gamma)x)=\rho(\gamma)\,\mathrm{coef}(x)$ for all $\gamma$ and $x$,
  $r(\gamma)\Lambda(h)\subset\Lambda(\gamma hk)$ for all $\gamma\in\Gamma$, $h\in G_f$ and $k\in J$,
  and $\Lambda(y)=\mathcal N_y$ for every $y\in Y$.
\item (Family ties.) Let $\lambda'=(\lambda'_\tau)_{\tau\in\mathcal T}$ be the dominant weights of
  length $n$ that the Harish-Chandra parameters of $D$ give at level $j+1$, and let
  $\tau_L\colon E_h\to L$ be the embeddings with $\sigma_0\circ\tau_L=\tau$. With these embeddings,
  $V_{\lambda'}(L)$, $\rho_{\lambda'}$ and the lattices $\nabla_{\lambda'}(M^\vee)$ for
  $\mathcal O_{E_h}$-lattices $M\subset E_h^n$ are those of Definition~\ref{pd-weight-lattices}. For
  $\gamma\in\Gamma$ let $\gamma_E\in GL_n(E_h)$ be its matrix and
  $\check\gamma=(\tau_L(\gamma_E)^{-\mathsf T})_{\tau\in\mathcal T}$. The \emph{lattice} of $h\in G_f$
  is $L_h=E_h^n\cap h\hat L_0$, where $L_0$ is the standard lattice of $V_j$ and $\hat L_0$ its
  closure. Among the data and conditions of a \emph{family tie} of $R$ are a $\sigma_0$-semilinear map
  $\iota_{\lambda'}\colon V_{\lambda'}(L)\to W_c$ with
  $\iota_{\lambda'}(\rho_{\lambda'}(\check\gamma)P)=\rho(\gamma)\,\iota_{\lambda'}(P)$ for
  $\gamma\in\Gamma$, and, for each tuple, an $L$-linear isomorphism $e\colon\mathcal V\to
  V_{\lambda'}(L)$ with $\mathrm{coef}=\iota_{\lambda'}\circ e$ and
  $e(\mathcal N_y)=\nabla_{\lambda'}(L_y^\vee)$ for every $y\in Y$. The \emph{lattice system of the
  family tie} at $(j,t)$ is $r(\gamma)=e^{-1}\rho_{\lambda'}(\check\gamma)e$ and
  $\Lambda(h)=e^{-1}\big(\nabla_{\lambda'}(L_h^\vee)\big)$. A vector $w\in W_c$ is \emph{integral for
  the lattice of $h$} if $w=\iota_{\lambda'}(P)$ for the map $\iota_{\lambda'}$ of some family tie of
  $R$ and some $P\in\nabla_{\lambda'}(L_h^\vee)$.
\end{enumerate}
The theta series of Definition~\ref{is-ties-class-sums}(5) are the $\vartheta_t(h,w)$ with $h\in G_f$
and $w$ integral for the lattice of $h$.
\end{definition}
\leanref{Formalism/LocalToGlobal/Divisibility/ClassSum.lean}{Formalism.L2.DIV.CSD}
\leanref{Formalism/LocalToGlobal/IndefiniteReal/Bridge.lean}{Formalism.IndefR.Bridge.ClassPres}
\leanref{Development/AbelianVarietyTheory/Proofs/SatakeParameters/ClassPresBasic.lean}{GlobalTheory.PresBasic.avgForm}
\leanref{Development/AbelianVarietyTheory/Definitions/Foundations/ThetaSeriesTies.lean}{GlobalTheory.OnTies.singleSeries}
\leanref{Development/AbelianVarietyTheory/Proofs/FoundationProofs/GoodLiftOnto.lean}{GlobalTheory.TraceIncl.kern}
\leanref{Development/AbelianVarietyTheory/Proofs/FoundationProofs/InjAndRatOfLi.lean}{GlobalTheory.TraceIncl.LatSys}
\leanref{Formalism/LocalToGlobal/IndefiniteReal/BridgeTie.lean}{Formalism.IndefR.Bridge.FamTie}
\leanref{Formalism/LocalToGlobal/IndefiniteReal/BridgeTie.lean}{Formalism.IndefR.Bridge.CoefTie}
\leanref{Formalism/LocalToGlobal/IndefiniteReal/BridgeTie.lean}{Formalism.IndefR.Bridge.CoefMapTie}
\leanref{Formalism/LocalToGlobal/IndefiniteReal/LinksChar.lean}{Formalism.IndefR.Links.RealDat.fam_tie}
\leanref{Formalism/LocalToGlobal/IndefiniteReal/BridgeTie.lean}{Formalism.IndefR.Bridge.hLat}
\leanref{Formalism/LocalToGlobal/Divisibility/CoefMod.lean}{Formalism.L2.DIV.nabAtr}
\leanref{Development/AbelianVarietyTheory/Definitions/Foundations/ThetaSeriesTies.lean}{GlobalTheory.OnTies.InCoefLat}

\begin{lemma}\label{is-ties-lattice-facts}
Let $R$ be a realization of a definite-side datum with a face $F$, let $j$ be a lift and $t$ a tuple,
with the notation of Definitions~\ref{is-ties-class-sums} and~\ref{is-ties-lattice-systems}. Suppose
that $R$ satisfies its characterization at $F$ and that its family of definite forms satisfies the
family conditions.
\begin{enumerate}
\item Let $\phi$ be an injective $\sigma_0$-semilinear map from an $L$-vector space to a complex
  vector space such that $\phi(x)=c\,\phi(x')$ with $x'\ne0$ implies $c\in\sigma_0(L)$. If
  $\phi(x)=\zeta\phi(x')$ for a root of unity $\zeta$, then $x=ax'$ for some $a\in\mathcal O_L$. The
  map $\mathrm{coef}$ has these two properties, and so has every $\sigma_0$-semilinear map that
  carries $L$-linearly independent vectors to $\mathbf C$-linearly independent ones.
\item The values of $\chi$ and of $\psi_0$ are roots of unity, and $G$ carries integral lift forms
  to integral lift forms.
\item For $\gamma\in\Gamma$, $h\in G_f$, $k\in J$ and $w\in W_c$,
  $\vartheta_t(\gamma h,\rho(\gamma)w)=\vartheta_t(h,w)$ and
  $\vartheta_t(hk,w)=\chi(k)^{-1}\vartheta_t(h,w)$.
\item For every lift form $f$ at $(j,t)$,
  $\mathrm{th}_F(j,t)\big(j_V(f)\big)=\sum_{y\in Y}|\Gamma_y|^{-1}\vartheta_t\big(y,\mathrm{coef}(f(y))\big)$.
  In particular $\mathrm{th}_F(j,t)\big(j_V(a_y(x))\big)=\vartheta_t\big(y,\mathrm{coef}(x)\big)$
  for $y\in Y$ and $x\in\mathcal V$.
\item The lattice system of every family tie of $R$ at $(j,t)$ is a lattice system.
\item Let $(r,\Lambda)$ be a lattice system. If $y\in Y$ and $x\in\mathcal N_y$, then $a_y(x)$ is
  integral. If $h\in G_f$ and $x\in\Lambda(h)$, then there are an integral lift form $f$ and $k\in J$
  with $\vartheta_t\big(h,\mathrm{coef}(x)\big)=\chi(k)^{-1}\,\mathrm{th}_F(j,t)\big(j_V(f)\big)$.
\item For every integral lift form $f$ at $(j,t)$, $\sum_{g\in P_0}\psi_0(g)\,gf$ is an integral
  level form.
\end{enumerate}
\end{lemma}
\leanref{Development/AbelianVarietyTheory/Proofs/FoundationProofs/InjAndRatOfLi.lean}{GlobalTheory.TraceIncl.exists_int_smul}
\leanref{Development/AbelianVarietyTheory/Proofs/FoundationProofs/InjAndRatOfLi.lean}{GlobalTheory.TraceIncl.inj_and_rat_of_li}
\leanref{Development/AbelianVarietyTheory/Proofs/FoundationProofs/InjAndRatOfLi.lean}{GlobalTheory.TraceIncl.coef_rat}
\leanref{Development/AbelianVarietyTheory/Proofs/SatakeParameters/ClassPresBasic.lean}{GlobalTheory.PresBasic.coef_injective}
\leanref{Development/AbelianVarietyTheory/Proofs/FoundationProofs/ResidueCharIsOfFinOrder.lean}{GlobalTheory.LevelCharOrder.levEnv_chi_isOfFinOrder}
\leanref{Formalism/LocalToGlobal/Divisibility/GoodPartHermitianSpec.lean}{Formalism.L2.DIV.CSD.ActGood}
\leanref{Development/AbelianVarietyTheory/Proofs/FoundationProofs/ThetaKernelFiniteAdelicLaws.lean}{GlobalTheory.TraceIncl.kern_laws}
\leanref{Development/AbelianVarietyTheory/Proofs/FoundationProofs/ThetaKernelFiniteAdelicLaws.lean}{GlobalTheory.TraceIncl.KernLaws.kern_fin}
\leanref{Development/AbelianVarietyTheory/Proofs/FoundationProofs/MemRangeSrcForm.lean}{GlobalTheory.TraceIncl.thF_eq_clSum}
\leanref{Development/AbelianVarietyTheory/Proofs/FoundationProofs/InjAndRatOfLi.lean}{GlobalTheory.TraceIncl.clSum_avgForm}
\leanref{Development/AbelianVarietyTheory/Proofs/FoundationProofs/MemRangeSrcForm.lean}{GlobalTheory.TraceIncl.exists_latSys}
\leanref{Development/AbelianVarietyTheory/Proofs/FoundationProofs/MulVecMemHatSwap.lean}{GlobalTheory.HLatMove.hLat_rat_mul_stab}
\leanref{Formalism/LocalToGlobal/Columns/CoefColLemmas.lean}{Formalism.L2.Col.nabAtr_gmap}
\leanref{Development/AbelianVarietyTheory/Proofs/FoundationProofs/InjAndRatOfLi.lean}{GlobalTheory.TraceIncl.LatSys.avg_mem}
\leanref{Development/AbelianVarietyTheory/Proofs/FoundationProofs/InjAndRatOfLi.lean}{GlobalTheory.TraceIncl.LatSys.exists_trace}
\leanref{Development/AbelianVarietyTheory/Proofs/FoundationProofs/MemRangeSrcForm.lean}{GlobalTheory.TraceIncl.exists_src_trU}
\leanref{Development/AbelianVarietyTheory/Proofs/FoundationProofs/MemRangeSrcForm.lean}{GlobalTheory.TraceIncl.trU_eigen}
\begin{proof}
(1) If $x'=0$, then $\phi(x)=0$ and $x=0=0\cdot x'$. Otherwise $\zeta=\sigma_0(a)$ for some $a\in L$.
If $\zeta^m=1$, then $\sigma_0(a^m)=1$, so $a^m=1$ and $a\in\mathcal O_L$. Then
$\phi(x)=\sigma_0(a)\phi(x')=\phi(ax')$, so $x=ax'$. Let $(b_i)$ be a basis of $\mathcal V$ whose
image under $\mathrm{coef}$ is a basis of $W_c$, and write $x=\sum_ix_ib_i$, $x'=\sum_ix'_ib_i$. Then
$\mathrm{coef}(x)=c\,\mathrm{coef}(x')$ says $\sigma_0(x_i)=c\,\sigma_0(x'_i)$ for every $i$. For $c=1$
this gives $x=x'$, and if $x'_i\ne0$ it gives $c=\sigma_0(x_i/x'_i)$. Now let $\phi$ carry
independent vectors to independent ones. A nonzero vector is independent, so $\phi$ is injective. If
$\phi(x)=c\,\phi(x')$ with $x'\ne0$, then $\phi(x)$ and $\phi(x')$ are dependent, so $x=ax'$ for some
$a\in L$; then $\sigma_0(a)\phi(x')=c\,\phi(x')$ with $\phi(x')\ne0$, and $c=\sigma_0(a)$.

(2) The character of the level is a product of characters of residue fields, so its values are roots
of unity. The group $P_0$ is finite, so the values of $\psi_0$ are roots of unity. That $G$ preserves
the integral lift forms is part of the family conditions.

(3) For $h'\in U(V_j)(\mathbf A_{F_h})$ let $\Theta(h',w)$ be the theta series of the Weil datum and the
test vector at $t$ at $h'$, for the polynomial $\mathrm{io}_P(j)(v\otimes\rho(h')^{-1}w)$. By the
characterization of the coefficient spaces, $\rho$ is trivial on the elements with archimedean
component $1$. By the characterization of the maps $\mathrm{io}_P$, the Weil representation of an
archimedean element $u$ acts on the test function of the polynomial $\mathrm{io}_P(j)(v\otimes w)$ as
$\rho(u)$ acts on $w$; so $\Theta(h'u,w)=\Theta(h',w)$, and $\Theta(h',w)=\vartheta_t(h'_f,w)$, where
$h'_f$ is the finite part of $h'$. By the characterization of the Weil data, theta series are
invariant under $U(V_j)(F_h)$ on the left. So $\Theta(\gamma h',\rho(\gamma)w)=\Theta(h',w)$, which
gives the first identity. By the conditions of the characterization on the level characters and the
test vectors, the element with finite part $k\in J$ and archimedean component $1$ multiplies the test
vector at $t$ by $\chi(k)^{-1}$, and it acts on the finite component of the test function only. So
right translation by it multiplies $\Theta(h',w)$ by $\chi(k)^{-1}$, which gives the second identity.

(4) By the characterization of the theta maps, $\mathrm{th}_F(j,t)\big(j_V(f)\big)$ is, at $v$ and
$g$, the integral over the fundamental domain of $U(V_j)(F_h)$ of the theta series at $(h',g)$ for
the polynomial $\mathrm{io}_P(j)(v\otimes\rho(h')^{-1}\varphi(h'_f))$, where $\varphi$ is the function
of $j_V(f)$ on $G_f$. By the characterization of the models,
$\varphi(\gamma yk)=\chi(k)\rho(\gamma)\,\mathrm{coef}(f(y))$ for $\gamma\in\Gamma$, $y\in Y$ and
$k\in J$. By the proof of (3) the integrand is $\vartheta_t(h'_f,\varphi(h'_f))$, and by (3) and the
transformation of $\varphi$, $h\mapsto\vartheta_t(h,\varphi(h))$ is invariant under $\Gamma$ on the
left and under $J$ on the right. The measure gives the preimage of $J$ volume $1$, so the integral
over the fundamental domain is $\sum_{y\in Y}|\Gamma_y|^{-1}\vartheta_t(y,\varphi(y))$, and
$\varphi(y)=\mathrm{coef}(f(y))$. For $f=a_y(x)$ only the term at $y$ remains, and it is
$|\Gamma_y|^{-1}\sum_{g\in\Gamma_y}\vartheta_t\big(y,\mathrm{coef}(\rho_y(g)x)\big)$. Here
$\mathrm{coef}(\rho_y(g)x)=\chi(k_g)\rho(\gamma_g)\,\mathrm{coef}(x)$ and $y=\gamma_gyk_g$, so by (3)
$\vartheta_t\big(y,\mathrm{coef}(\rho_y(g)x)\big)=\chi(k_g)\vartheta_t\big(yk_g,\mathrm{coef}(x)\big)
=\vartheta_t\big(y,\mathrm{coef}(x)\big)$ for every $g$.

(5) Let a family tie be given, with $\iota_{\lambda'}$ and $e$. Let $\gamma\in\Gamma$. By the
conditions of the characterization on the level characters, $L_{\gamma hk}=\gamma_EL_h$ for $h\in
G_f$ and $k\in J$. Let $M\subset E_h^n$ be an $\mathcal O_{E_h}$-lattice and $w$ a finite place of
$L$. The columns of matrices $B_\tau$ form bases at $w$ of the duals of the $\mathcal O_L$-spans of
the $\tau_L(M)$ if and only if the columns of the $\tau_L(\gamma_E)^{-\mathsf T}B_\tau$ form such bases
for $\gamma_EM$; so by Definition~\ref{pd-weight-lattices}(5),
$\nabla_{\lambda'}\big((\gamma_EM)^\vee\big)=\rho_{\lambda'}(\check\gamma)\nabla_{\lambda'}(M^\vee)$.
Now $\mathrm{coef}(r(\gamma)x)=\iota_{\lambda'}\big(\rho_{\lambda'}(\check\gamma)e(x)\big)
=\rho(\gamma)\,\iota_{\lambda'}(e(x))=\rho(\gamma)\,\mathrm{coef}(x)$. If $x\in\Lambda(h)$, then
$e(r(\gamma)x)=\rho_{\lambda'}(\check\gamma)e(x)$ lies in
$\rho_{\lambda'}(\check\gamma)\nabla_{\lambda'}(L_h^\vee)=\nabla_{\lambda'}\big((\gamma_EL_h)^\vee\big)
=\nabla_{\lambda'}(L_{\gamma hk}^\vee)$, so $r(\gamma)x\in\Lambda(\gamma hk)$. Finally
$\Lambda(y)=e^{-1}\big(\nabla_{\lambda'}(L_y^\vee)\big)=\mathcal N_y$ for $y\in Y$.

(6) Let $x\in\mathcal N_y$ and $g\in\Gamma_y$. As $\gamma_gyk_g=y$, $r(\gamma_g)x\in\Lambda(y)=\mathcal
N_y$. Moreover $\mathrm{coef}(\rho_y(g)x)=\chi(k_g)\rho(\gamma_g)\,\mathrm{coef}(x)
=\chi(k_g)\,\mathrm{coef}(r(\gamma_g)x)$, and $\chi(k_g)$ is a root of unity by (2). By (1),
$\rho_y(g)x=a\,r(\gamma_g)x$ for some $a\in\mathcal O_L$, so $\rho_y(g)x\in\mathcal N_y$. Summing over
$g$, the value of $a_y(x)$ at $y$ lies in $\mathcal N_y$, and its other values are $0$.

Now let $x\in\Lambda(h)$, and write $h=\gamma yk$ with $\gamma\in\Gamma$, $y\in Y$ and $k\in J$. Then
$r(\gamma^{-1})x\in\Lambda(\gamma^{-1}hk^{-1})=\Lambda(y)=\mathcal N_y$, so $f=a_y(r(\gamma^{-1})x)$ is
integral by the first part. By (4),
$\mathrm{th}_F(j,t)\big(j_V(f)\big)=\vartheta_t\big(y,\rho(\gamma)^{-1}\mathrm{coef}(x)\big)$, and by
(3) this is $\vartheta_t\big(\gamma y,\mathrm{coef}(x)\big)=\chi(k)\,\vartheta_t\big(h,\mathrm{coef}(x)\big)$.

(7) Put $f'=\sum_{g\in P_0}\psi_0(g)\,gf$. For $u\in P_0$,
$uf'=\sum_g\psi_0(g)\,(ug)f=\sum_g\psi_0(u^{-1}g)\,gf=\psi_0(u)^{-1}f'$, so $f'$ is a level form. Each
$gf$ is integral and each $\psi_0(g)$ lies in $\mathcal O_L$, by (2), so $f'$ is integral.
\end{proof}

\begin{lemma}\label{is-ties-class-sum-facts}
Let $R$ be a realization of a definite-side datum with a face $F$, with the notation of
Definition~\ref{is-ties-class-sums}. Suppose that $R$ satisfies its characterization at $F$, that its
family of definite forms satisfies the family conditions and its rows $S_0(R,F)$ have the base links
(Lemma~\ref{is-eigenvalue-base-inputs}(1)), and that its rational structures satisfy (11) of
Definition~\ref{is-primitive-rows}. Let $j$ be a lift such that the model of discrete series of
$U(W_j)$ at $\sigma_0$ has an extremal vector (Definition~\ref{is-bounded-rows}(1)), let $t$ be a
tuple and $\tau$ a class sum at $(j,t)$. Then:
\begin{enumerate}
\item $\tau(e_{\mathcal X,j,t}x)=e_{\mathrm{good}}\tau(x)$ for every $x\in L^{s(j,t)}$;
\item $\tau(x)\in FW_j^0$ for every $x\in L^{s(j,t)}$;
\item $\Xi_{j,t}\subset\tau(N_{j,t})$;
\item $\tau(N_{j,t})$ lies in the $L$-span of $\Xi_{j,t}$;
\item $W^g_j\subset e_{\mathrm{good}}\tau(L^{s(j,t)})$;
\item if $e_{\mathrm{good}}\tau(x)\in W^g_j$, then $e_{\mathrm{good}}\tau(x)=c_j(\theta_{j,t}\pi_{j,t}(x))$,
  where $\theta_{j,t}$ is the matrix of the rows (Definition~\ref{is-complex-data}(c)).
\end{enumerate}
\end{lemma}
\leanref{Development/AbelianVarietyTheory/Proofs/FoundationProofs/MemRangeSrcForm.lean}{GlobalTheory.TraceIncl.HidaSetting}
\leanref{Development/AbelianVarietyTheory/Proofs/DivisorProofs/SubOrthProjPerpMem.lean}{GlobalTheory.ClassSumIdentity.classSum_eOp}
\leanref{Development/AbelianVarietyTheory/Proofs/DivisorProofs/SubOrthProjPerpMem.lean}{GlobalTheory.ClassSumIdentity.lvlE_lvlOp}
\leanref{Development/AbelianVarietyTheory/Proofs/FoundationProofs/MemRangeSrcForm.lean}{GlobalTheory.TraceIncl.HidaSetting.theta_mem_typeP}
\leanref{Development/AbelianVarietyTheory/Proofs/FoundationProofs/MemRangeSrcForm.lean}{GlobalTheory.TraceIncl.HidaSetting.thetaTuple_le}
\leanref{Development/AbelianVarietyTheory/Proofs/FoundationProofs/MemRangeSrcForm.lean}{GlobalTheory.TraceIncl.HidaSetting.theta_mem_span}
\leanref{Development/AbelianVarietyTheory/Proofs/DivisorProofs/GoodLiftOntoOfBase.lean}{GlobalTheory.BaseTwistIncl.goodLiftOnto_of_base}
\leanref{Development/AbelianVarietyTheory/Proofs/DivisorProofs/ExistsTieCoordBase.lean}{GlobalTheory.BigOnBase.goodOp_classSum_base}
\begin{proof}
The characterization, the uniformity (11) and the extremal vector give the following. A vector $z$ of
$L^{n_W(j)}$ is determined by $\Phi_j(z)$, and a vector of $H_W(j,\sigma_0)$ by its automorphic
form; $\Phi_j(\varrho_j(g)z)=\sigma_0(\chi_j(g))\,g\Phi_j(z)$ for $g\in G$; and
$\mathrm{th}_F(j,t)\circ j_V$ commutes with the actions of $G$. Since $e_{\mathcal X}$ is central,
$f_{j,t}(e_{\mathcal X,j,t}x)=\sum_gc_{\mathcal X}(g)\,gf_{j,t}(x)$. By the family conditions, the
image of $c_{0,j,t}$ consists of level forms.

(1) As $j_V$ is $\sigma_0$-semilinear,
\[
 \Phi_j\big(e_{\mathrm{good}}\tau(x)\big)=\sum_g\sigma_0(c_{\mathcal X}(g))\,g\Phi_j(\tau(x))
 =\mathrm{th}_F(j,t)\,j_V\Big(\sum_gc_{\mathcal X}(g)\,gf_{j,t}(x)\Big)
 =\Phi_j\big(\tau(e_{\mathcal X,j,t}x)\big),
\]
so $e_{\mathrm{good}}\tau(x)=\tau(e_{\mathcal X,j,t}x)$.

(2) For $u\in P_0$, $uf_{j,t}(x)=\psi_0(u)^{-1}f_{j,t}(x)$, so
$\Phi_j(\varrho_j(u)\tau(x))=\sigma_0(\chi_j(u))\,u\,\mathrm{th}_F(j,t)j_V(f_{j,t}(x))
=\sigma_0(\chi_j(u)\psi_0(u)^{-1})\Phi_j(\tau(x))$. Hence
$\varrho_j(u)\tau(x)=\chi_j(u)\psi_0(u)^{-1}\tau(x)$.

(3) Let $z$ be one of the vectors that span $\Xi_{j,t}$. With the notation of
Definition~\ref{is-ties-lattice-systems},
$\Phi_j(z)=\sum_{g\in P_0}\sigma_D(\psi_0(g))\,g\,\vartheta_t(h,w)$ for some $h\in G_f$ and
$w=\iota_{\lambda'}(P)$, where $\iota_{\lambda'}$ is the map of a family tie of $R$ and
$P\in\nabla_{\lambda'}(L_h^\vee)$. Let $(r,\Lambda)$ be the lattice system of this family tie at
$(j,t)$, with its isomorphism $e$; it is a lattice system by Lemma~\ref{is-ties-lattice-facts}(5).
The vector $q=e^{-1}(P)$ lies in $\Lambda(h)$, and $\mathrm{coef}(q)=w$. By part (6) of that lemma
there are an integral lift form $f$ and $k\in J$ with
$\vartheta_t(h,w)=\chi(k)^{-1}\,\mathrm{th}_F(j,t)\big(j_V(f)\big)$, and by its part (7)
$\sum_{g\in P_0}\psi_0(g)\,gf=f_{j,t}(x)$ for some $x\in N_{j,t}$. As $j_V$ is $\sigma_0$-semilinear,
$\sigma_0=\sigma_D$ (Definition~\ref{is-realisation}) and $\mathrm{th}_F(j,t)\circ j_V$ commutes with the
actions of $G$,
\[
 \Phi_j(\tau(x))=\mathrm{th}_F(j,t)\,j_V\Big(\sum_{g\in P_0}\psi_0(g)\,gf\Big)
 =\sum_{g\in P_0}\sigma_D(\psi_0(g))\,g\,\mathrm{th}_F(j,t)\big(j_V(f)\big)=\chi(k)\,\Phi_j(z).
\]
So $j_{\sigma_0}(z)=\chi(k)^{-1}j_{\sigma_0}(\tau(x))$. By (11) of Definition~\ref{is-primitive-rows},
$j_{\sigma_0}$ carries $L$-linearly independent vectors to $\mathbf C$-linearly independent ones, and
$\chi(k)^{-1}$ is a root of unity, so parts (1) and (2) of Lemma~\ref{is-ties-lattice-facts} give
$a\in\mathcal O_L$ with $z=a\tau(x)=\tau(ax)$. As $ax\in N_{j,t}$ and $\tau(N_{j,t})$ is an
$\mathcal O_L$-module, $\Xi_{j,t}\subset\tau(N_{j,t})$.

(4) Let $x\in N_{j,t}$, so that $F=f_{j,t}(x)$ is an integral level form. By the characterization,
$R$ has a family tie; let $(r,\Lambda)$ be its lattice system at $(j,t)$, with its map
$\iota_{\lambda'}$ and its isomorphism $e$. It is a lattice system by
Lemma~\ref{is-ties-lattice-facts}(5). Let $y\in Y$. As $F(y)\in\mathcal N_y$, the averaged form
$a_y(F(y))$ is integral by part (6) of that lemma, and by its part (7) there is $z_y\in N_{j,t}$ with
$f_{j,t}(z_y)=\sum_{g\in P_0}\psi_0(g)\,g\,a_y(F(y))$. By its part (4) and the $G$-equivariance, as in
(3),
\[
 \Phi_j(\tau(z_y))=\sum_{g\in P_0}\sigma_D(\psi_0(g))\,g\,\vartheta_t\big(y,\mathrm{coef}(F(y))\big),
\]
and $\mathrm{coef}(F(y))=\iota_{\lambda'}\big(e(F(y))\big)$ with
$e(F(y))\in e(\mathcal N_y)=\nabla_{\lambda'}(L_y^\vee)$. So $\tau(z_y)\in\Xi_{j,t}$. Next,
$|P_0|\,x=\sum_{y\in Y}|\Gamma_y|^{-1}z_y$. Indeed, $F$ is a level form, so
$\sum_{g\in P_0}\psi_0(g)\,gF=|P_0|\,F$; and $F=\sum_y|\Gamma_y|^{-1}a_y(F(y))$, because $F(y)$ is
fixed by $\Gamma_y$, so that $a_y(F(y))$ has value $|\Gamma_y|\,F(y)$ at $y$ and $0$ at the other
representatives. Hence $f_{j,t}(|P_0|\,x)=\sum_y|\Gamma_y|^{-1}f_{j,t}(z_y)$, and $f_{j,t}$ is
injective. So $\tau(x)=|P_0|^{-1}\sum_y|\Gamma_y|^{-1}\tau(z_y)$ lies in the $L$-span of $\Xi_{j,t}$.
% (3) is HidaSetting.thetaTuple_le: exists_latSys for the tie of the InCoefLat witness,
% LatSys.exists_trace, exists_src_trU, kern_fin and adForm_theta_trU (clSum_trU), eq_of_low, and
% exists_int_smul with j_inj_rat. (4) is HidaSetting.theta_mem_span: exists_latSys for R.fam_tie,
% LatSys.avg_mem and exists_src_trU for each class, exists_fin and clSum_avgForm for seriesSet, and
% srcForm_injective with trU_of_eigen and sum_avgForm. The trace levelTrace has the weights
% embLL (psi0 g) = sigma_D (psi0 g), as does clSum_trU for the sigma_D-semilinear class sum (presV); the
% book's sigma_0 is sg0 D A, equal to embLL by TraceIncl.sign_at_base_place_eq_embedding (RealDat.hv0).

For (5) and (6), first $j_{\sigma_0}(e_{\mathrm{good}}\tau(x))=\mathrm{th}_G(j,t)\,j_V(f_{j,t}(x))$
for every $x$. Indeed, by the characterization of the theta maps the automorphic form of
$\mathrm{th}_G(j,t)(\varphi)$ is $\sum_g\sigma_0(c_{\mathcal X}(g))\,g\,\mathrm{th}_F(j,t)(\varphi)$,
which for $\varphi=j_V(f_{j,t}(x))$ is $\Phi_j(e_{\mathrm{good}}\tau(x))$ by the first equality in (1).

(5) Let $z\in W^g_j$. Up to the normalization of the rows, $j_{\sigma_0}\circ c_j$ is the map $w_0(j)$
of Definition~\ref{is-complex-data}(c). By the base links $t_0(j,t)$ has the same image as $w_0(j)$, so
$j_{\sigma_0}(z)$ is a $\mathbf C$-linear combination of the vectors
$\mathrm{th}_G(j,t)\,j_V(c_{0,j,t}(y))$. Each $c_{0,j,t}(y)$ is a level form $f_{j,t}(x)$, so
$j_{\sigma_0}(z)$ lies in the $\mathbf C$-span of the vectors $j_{\sigma_0}(e_{\mathrm{good}}\tau(x))$.
As $j_{\sigma_0}$ carries $L$-linearly independent vectors to $\mathbf C$-linearly independent ones,
$z$ lies in the $L$-span of $e_{\mathrm{good}}\tau(L^{s(j,t)})$, which is that subspace.

(6) Let $z=e_{\mathrm{good}}\tau(x)\in W^g_j$. By the base links $\theta_{j,t}$ is invertible, so
$z=c_j(\theta_{j,t}y)$ for some $y$, and by the definition of $\theta_{j,t}$,
$j_{\sigma_0}(z)=\mathrm{th}_G(j,t)\,j_V(c_{0,j,t}(y))$. Write $c_{0,j,t}(y)=f_{j,t}(x')$, so that
$\pi_{j,t}(x')=y$. Then $\mathrm{th}_G(j,t)\,j_V$ vanishes on $f_{j,t}(x-x')$. The image of
$c_{0,j,t}$ lies in the definite subspace, whose image under $j_V$ lies in the image of the adjoint of
$\mathrm{th}_G(j,t)$ (Definition~\ref{is-complex-data}), and $j_V$ carries the class-sum pairing to
the inner product by the characterization of the models. So $f_{j,t}(x-x')$ is orthogonal to the
image of $c_{0,j,t}$, and $\pi_{j,t}(x-x')=0$. Hence $\pi_{j,t}(x)=y$ and
$z=c_j(\theta_{j,t}\pi_{j,t}(x))$.
\end{proof}

\begin{lemma}\label{is-ties-level-tie}
Let $p$ be a point at a member of a family (Definition~\ref{is-point}) whose definite-side datum has
ranks $k$ and $m$. Let $R$ be the realization of $p$, with the scalars of its characterization, and
$F$ its face, with the notation of Definition~\ref{is-ties-class-sums}. Let a pair of runs at $p$ be
given, and let $O$ be a record object recorded at $p$ and these runs
(Definition~\ref{is-record-object}), with the notation there and in
Definition~\ref{is-onto-record-numbers}. Then $r_0=r$, where $r=k+m-1$ is the rank of the
definite-side datum, the group $P$ of $O$ is the product over $P'\in P_g$ of the mirabolic subgroups
of $GL_{r}(k_{P'})$, and the isomorphism $G\cong\prod_{P'\in P_g}GL_{r}(k_{P'})$ of the realization
(Definition~\ref{is-realisation}(1))
gives an injective homomorphism $\mathrm{inc}\colon P\to G$. Moreover there are $L$-linear
isomorphisms $\alpha_j\colon L^{n_Z(j)}\to L^{n_W(j)}$, for the lifts $j$, and
$\beta_{v,j}\colon L^{n_X(v,j)}\to L^{s(j,\mu_v)}$, for the finite places $v$ and the lifts $j$, such
that for every $v$ and $j$, with $t=\mu_v$:
\begin{enumerate}
\item $\alpha_j(\rho_j(g)x)=\varrho_j(\mathrm{inc}(g))\,\alpha_j(x)$ for $g\in P$;
\item $U=\mathrm{inc}^{-1}(P_0)$, and $\psi(g)=\psi_0(\mathrm{inc}(g))^{-1}$ for $g\in U$;
\item $\alpha_je_j=e_{\mathrm{good}}\alpha_j$, and $V_0(j)=e_j(L^{n_Z(j)})$;
\item $M_s(j)=\alpha_j^{-1}(\Lambda_j)$;
\item $X_3(v,j)=e_j\alpha_j^{-1}(\Xi_{j,t})$ and $M_t(v,j)=\beta_{v,j}^{-1}(N_{j,t})$;
\item $\beta_{v,j}e_X=e_{\mathcal X,j,t}\beta_{v,j}$ and $\beta_{v,j}\,\mathrm{pr}=\varpi\beta_{v,j}$,
  where $e_X$ and $\mathrm{pr}$ are the endomorphisms of $O$ at $(v,j)$ and $\varpi$ is the orthogonal
  projection, for $H_{j,t}$, onto the orthogonal complement of $\beta_{v,j}(\ker\theta'_{v,j})$;
\item if there is a class sum at $(j,t)$, then $\alpha_j\theta'_{v,j}\beta_{v,j}^{-1}$ is one.
\end{enumerate}
\end{lemma}
\leanref{Development/AbelianVarietyTheory/Definitions/Divisors/RealisedChainTorusEq.lean}{GlobalTheory.BigOn.IsRec}
\leanref{Development/AbelianVarietyTheory/Proofs/DivisorProofs/NotOntoOfCharAlone.lean}{GlobalTheory.RealIndexBound.OnObjLevelCompatibility}
\leanref{Development/AbelianVarietyTheory/Definitions/Foundations/ThetaSeriesTies.lean}{GlobalTheory.OnTies.OnTie}
\leanref{Development/AbelianVarietyTheory/Definitions/Foundations/ThetaSeriesTies.lean}{GlobalTheory.OnTies.toForms_bijective}
\leanref{Development/AbelianVarietyTheory/Definitions/Foundations/ThetaSeriesTies.lean}{GlobalTheory.OnTies.OnTie.rho_eq_of_actionPin}
\leanref{Development/AbelianVarietyTheory/Proofs/DivisorProofs/NotOntoOfCharAlone.lean}{GlobalTheory.RealIndexBound.mem_Ms_iff}
\leanref{Development/AbelianVarietyTheory/Definitions/AutomorphicMembership/LiftWeightComplexPointEqTwist.lean}{GlobalTheory.OntoChain.Mir}
\leanref{Development/AbelianVarietyTheory/Definitions/AutomorphicMembership/LiftWeightComplexPointEqTwist.lean}{GlobalTheory.OntoChain.mirIncl}
\begin{proof}
The value of $r_0$ is that of Lemma~\ref{is-onto-recorded-numbers}(1). The group of a record object at
a point is the product of the mirabolic subgroups, and $\mathrm{inc}$ is injective because the
isomorphism of the realization is. The rest are the conditions of the level tie of $O$ at $v$ and
$j$, which is part of being recorded. There $\alpha_j$ is the composite of the isomorphism of the tie
from $L^{n_Z(j)}$ to the cusp forms at the level of the $j$-th lift, which is the same at every
place, with the map from these cusp forms to $L^{n_W(j)}$ given by the weight tie of the
realization, which is bijective; (1) uses the action pin of the tie, and (4) the weight tie.
\end{proof}

\begin{lemma}\label{is-twisted-places-chain}
Let $p$ be a point at a member of a family (Definition~\ref{is-point}) whose realization has the
links of the further conditions, let a pair of runs pinned to the own places be given, and let $O$ be a
record object with constants $c$ recorded at $p$ and this pair (Definition~\ref{is-record-object}).
Suppose that the realization is rational at the twisted places (Definition~\ref{is-ties-class-sums}) and that
$O$ satisfies (3a) to (3g) of Lemma~\ref{is-onto-estimates}. Let $C_E=4C^2+25C+C_{b,1}+C_{b,2}+48$.
Then there are real numbers
$\mathrm{err}(w,j)$ such that the rows $S_0(p)$ satisfy, for the numerical bounds $N_p(C_E,C_H)$:
\begin{enumerate}
\item condition (1) of Definition~\ref{is-surjectivity-bound} with $\mathrm{err}$ in place of
  $\mathrm{err}_0$;
\item $\mathrm{err}(w,j)\geq0$ for every twisted place $w$ and lift $j$;
\item $\sum_{w\in\mathrm{Tw}}\ell_w\big(e_w\,\mathrm{err}(w,j)+e_wJ_0(\mu_w,w,j)\big)\le C_EE_1^{u}$
  for every lift $j$.
\end{enumerate}
\end{lemma}
\leanref{Development/AbelianVarietyTheory/Proofs/DivisorProofs/TwistedPlacesBoundChainTwoRunChainStatement.lean}{GlobalTheory.T.twistedPlacesBoundChain_twoRunChain}
\leanref{Development/AbelianVarietyTheory/Proofs/DivisorProofs/OnClausesWithErrorOfSplitPartStatement.lean}{GlobalTheory.T.onClauses_with_error_of_splitPart}
\leanref{Development/AbelianVarietyTheory/Proofs/DivisorProofs/SubOrthProjPerpMem.lean}{GlobalTheory.ClassSumIdentity.leafIdentity_of_isRec}
\leanref{Development/AbelianVarietyTheory/Proofs/DivisorProofs/GoodLiftOntoOfBase.lean}{GlobalTheory.BaseTwistIncl.leafTW2_of_isRec}
\begin{proof}
Let $P$ be the group of $O$, and let $S$ be the second presentation of the rows of the realization at
its face, the indefinite-side datum of which $O$ is a record object.

\emph{Step 1: the ties of the record.} Write $R$ for the realization of $p$, with the scalars of its
characterization, and $F$ for its face, with subgroup $P_0$ and character $\psi_0$, so that $S$ is the
rows of $R$ at $F$, with the notation of Definition~\ref{is-ties-class-sums} and
Lemma~\ref{is-ties-level-tie}. The links of the further conditions include the characterization of
$R$ at $F$, the family conditions, the base links of the rows, the condition (11) of
Definition~\ref{is-primitive-rows} on the rational structures, and the rank condition in (d) below;
the weight tie of $p$ gives extremal vectors of the models of the $W_j$ at $\sigma_0$. So
Lemma~\ref{is-ties-class-sum-facts} applies to $R$, $F$ and every lift. Since $O$ is recorded at
$p$ and the pair of runs, and the realization is rational at the twisted places and has the links of
the further conditions, the following hold.
\begin{enumerate}
\item[(a)] $P$ is the product over $P'\in P_g$ of the mirabolic subgroups of $GL_{r}(k_{P'})$ of
  Lemma~\ref{is-places-level-index}, and $U$ is the product of their upper triangular unipotent
  subgroups. Moreover $(U,\psi)$ is a Mackey pair (Lemma~\ref{is-places-local-index}), and each
  $e_j$ is an idempotent with image $V_0(j)$.
\item[(b)] Hypothesis (3) of Lemma~\ref{is-places-combination} holds: the identity of the class sums
  and the inclusion of the base twists.
\item[(c)] After a change of the maps $\varphi_j$ of $O$, which changes none of its other data,
  hypothesis (4) of Lemma~\ref{is-places-combination} holds.
\item[(d)] The rank $D$ of $S$ is at least $1$, and $d_W(j)=D$ for every lift $j$.
\end{enumerate}

(a) By Lemma~\ref{is-ties-level-tie}, $r_0=r=k+m-1\geq2$, where $r$ is the rank of the definite-side
datum, $P$ is the product of the mirabolic subgroups and $U=\mathrm{inc}^{-1}(P_0)$. The
characterization of $R$ on its group says that, under the isomorphism
$G\cong\prod_{P'}GL_{r}(k_{P'})$, $P_0$ is the product of the upper triangular
unipotent subgroups, and that there are nontrivial complex characters $\psi'_{P'}$ of the additive
groups of the $k_{P'}$ with $\sigma_0(\psi_0(g))=\prod_{P'}\psi'_{P'}(s(g_{P'}))$ for $g\in P_0$,
where $s$ is as in Lemma~\ref{is-ties-mirabolic-mackey}. Upper triangular unipotent matrices are
mirabolic and $\mathrm{inc}$ is injective, so $P_0\subset\mathrm{inc}(P)$ and $U$ is the product of
the upper triangular unipotent subgroups. The values of $\psi_0$ are roots of unity, so lie in
$\mathcal O_L^\times$. For $a\in k_{P'}$ let $u_a\in U$ have component $1+aE_{12}$ at $P'$, with
$E_{12}$ the matrix unit, and $1$ elsewhere, and put $\chi_{P'}(a)=\psi_0(\mathrm{inc}(u_a))^{-1}$.
Then $\sigma_0\circ\chi_{P'}=(\psi'_{P'})^{-1}$, so $\chi_{P'}$ is a nontrivial additive character
of $k_{P'}$ with values in $\mathcal O_L$, and by Lemma~\ref{is-ties-level-tie}(2)
$\psi(u)=\psi_0(\mathrm{inc}(u))^{-1}=\prod_{P'}\chi_{P'}(s(u_{P'}))$ for $u\in U$. By
Lemma~\ref{is-ties-mirabolic-mackey} with $A=\mathcal O_L$, $U$ and the function
$h\mapsto\prod_{P'}\chi_{P'}(s(h_{P'}))$ form a Mackey pair. The conditions for a Mackey pair
evaluate the function only at elements of $U$, where it agrees with $\psi$, so $(U,\psi)$ is a Mackey
pair.

By Lemma~\ref{is-ties-level-tie}(3), $V_0(j)=e_j(L^{n_Z(j)})$ and
$e_j=\alpha_j^{-1}e_{\mathrm{good}}\alpha_j$. The map $g\mapsto\chi_j(g)^{-1}\varrho_j(g)$ is a
representation of $G$, and $e_{\mathrm{good}}$ is the image of $e_{\mathcal X}$ under its extension
to $L[G]$. By the family conditions the characters in $\mathcal X$ are absolutely irreducible, so
$e_{\mathcal X}$ is a central idempotent of $L[G]$ \cite[\S2.6, Theorem~8]{Ser77}. Hence
$e_{\mathrm{good}}$ is an idempotent that commutes with $\varrho_j(G)$, and $e_j$ is an idempotent.

Next, for every lift $j$,
\[
 Z_g(j)=\alpha_j^{-1}(W^g_j).\tag{$\dagger$}
\]
By the characterization of $R$ on its group, $\chi_j$ is a product of characters of the determinants
of the components, so it is trivial on $P_0$. Let $x\in L^{n_Z(j)}$. By (1) and (2) of
Lemma~\ref{is-ties-level-tie}, $x\in V_\psi(j)$ if and only if
$\varrho_j(\mathrm{inc}(u))\alpha_j(x)=\psi_0(\mathrm{inc}(u))^{-1}\alpha_j(x)$ for every $u\in U$.
As $P_0\subset\mathrm{inc}(P)$ and $U=\mathrm{inc}^{-1}(P_0)$, this says that $P_0$ acts on
$\alpha_j(x)$ through $\chi_j\psi_0^{-1}$, that is, $\alpha_j(x)\in FW_j^0$. By (3) of that lemma,
$x\in V_0(j)$ if and only if $\alpha_j(x)\in e_{\mathrm{good}}(L^{n_W(j)})$. This is $(\dagger)$.

(b) Let $w$ be a twisted place, $v\in\{w,\bar w\}$ and $j$ a lift, and put $t=\mu_w$. The
characterization gives $\mu_{\bar w}=\mu_w$, so $t=\mu_v$. By the rationality there is a class sum at
$(j,t)$, so $\tau=\alpha_j\theta'_{v,j}\beta^{-1}$ is one by Lemma~\ref{is-ties-level-tie}(7), where
$\beta=\beta_{v,j}$. Write $\theta'=\theta'_{v,j}$, and $e_X$, $\mathrm{pr}$ for the endomorphisms of
$O$ at $(v,j)$. For every $x$, (6) and (3) of Lemma~\ref{is-ties-level-tie} and
Lemma~\ref{is-ties-class-sum-facts}(1) give
\[
 \alpha_j\theta'(e_Xx)=\tau(e_{\mathcal X,j,t}\beta x)=e_{\mathrm{good}}\tau(\beta x)
 =e_{\mathrm{good}}\alpha_j\theta'(x)=\alpha_je_j\theta'(x),
\]
so $\theta'e_X=e_j\theta'$. By the family conditions $H_{j,t}$ is positive definite at $\sigma_0$.
For such a form and a subspace $N$, a vector differs from its orthogonal projection onto $N^\perp$ by
an element of $N^{\perp\perp}=N$. With $N=\beta(\ker\theta')$, Lemma~\ref{is-ties-level-tie}(6)
gives $\beta(x-\mathrm{pr}(x))\in\beta(\ker\theta')$, so $\theta'(\mathrm{pr}(x))=\theta'(x)$ for
every $x$, in particular for $x$ in the image of $e_X$. This is the identity of the class sums.

For the inclusion of the base twists put $\mathrm{Im}=e_j\theta'(M_t(v,j))$. By (5) and (3) of
Lemma~\ref{is-ties-level-tie}, $\alpha_j(\mathrm{Im})=e_{\mathrm{good}}\tau(N_{j,t})$ and
$\alpha_j(X_3(v,j))=e_{\mathrm{good}}(\Xi_{j,t})$, so $X_3(v,j)\subset\mathrm{Im}$ by
Lemma~\ref{is-ties-class-sum-facts}(3). By the family conditions $N_{j,t}$ is a lattice, so
$\mathrm{Im}$ is finitely generated, and so is its submodule $X_3(v,j)$. For every $x$,
$\tau(x)\in FW_j^0$ by Lemma~\ref{is-ties-class-sum-facts}(2), and $e_{\mathrm{good}}$ commutes with
$\varrho_j(G)$, so $e_{\mathrm{good}}\tau(x)\in W^g_j$; by $(\dagger)$, $\mathrm{Im}\subset Z_g(j)$.
Conversely let $z\in Z_g(j)$. By $(\dagger)$ and Lemma~\ref{is-ties-class-sum-facts}(5),
$\alpha_j(z)=e_{\mathrm{good}}\tau(x)$ for some $x$. As $N_{j,t}$ spans $L^{s(j,t)}$, part (4) of that
lemma puts $\tau(x)$ in the $L$-span of $\Xi_{j,t}$, so $z=e_j\alpha_j^{-1}\tau(x)$ lies in the
$L$-span of $e_j\alpha_j^{-1}(\Xi_{j,t})=X_3(v,j)$. Hence $X_3(v,j)\subset\mathrm{Im}\subset Z_g(j)$
are finitely generated and span $Z_g(j)$: both are lattices in $Z_g(j)$, and
$X_3(v,j)_{(v)}\subset\mathrm{Im}_{(v)}$.

(c) Replace $\varphi_j$ by $\varphi'_j=\alpha_j^{-1}c_j$; no other datum of $O$ changes. The map
$\varphi'_j$ is injective, and by $(\dagger)$ its image is $\alpha_j^{-1}(W^g_j)=Z_g(j)$. As
$M_j=c_j^{-1}(\Lambda_j)$ and $M_s(j)=\alpha_j^{-1}(\Lambda_j)$ (Lemma~\ref{is-ties-level-tie}(4)),
$\varphi'_j(M_j)=\alpha_j^{-1}(\Lambda_j\cap W^g_j)=M_s(j)\cap Z_g(j)$. Now let $w$, $v$, $j$, $t$,
$\tau$ and $\beta$ be as in the proof of (b), and let $x\in M_t(v,j)$. By (b),
$e_j\theta'(x)\in\mathrm{Im}\subset Z_g(j)$, so $e_{\mathrm{good}}\tau(\beta x)=\alpha_je_j\theta'(x)$
lies in $W^g_j$, and Lemma~\ref{is-ties-class-sum-facts}(6) gives
$e_j\theta'(x)=\varphi'_j(\theta_{j,t}\pi_{j,t}\beta x)$. By (b) also
$\theta'(\mathrm{pr}(e_Xx))=\theta'(e_Xx)=e_j\theta'(x)$. Since $\beta(M_t(v,j))=N_{j,t}$ and
$T_{j,t}=\pi_{j,t}(N_{j,t})$, both $\varphi'_j(\theta_{j,t}T_{j,t})$ and
$\theta'(\mathrm{pr}(e_X(M_t(v,j))))$ are the set of the $e_j\theta'(x)$ with $x\in M_t(v,j)$.

(d) This is the rank condition among the links.
% (a) is GlobalTheory.RealIndexBound.level_of_tied: U_eq with mem_base_iff (grp_levels), psi_eq with
% levelChar_spec (exists_isLevelChar, exists_psiU_eq), leafMackey = leafMackey_of
% (isMackey_pi_upper_unipotent_prod_character), and e_eOp, V0_eq and goodOp_idem (idempotent_central);
% IsRec.rr_eq and IsRec.depth give TieLevelNum (f = D.R.f, rr = rankPred + 1). (dagger) is mem_Zg_iff
% (chi_eq_one from grp_chiW). (b) is GlobalTheory.ClassSumIdentity.leafIdentity_of_isRec =
% identity_of_tie (classSum_eOp, lvlE_lvlOp, sub_orthProj_perp_mem with lf_lvlH_pos) and
% GlobalTheory.BaseTwistIncl.leafTW2_of_isRec = TraceIncl.lattices_of_tie, which
% GlobalTheory.BigOn.OntoPath.onto_path derives and twistedPlacesBoundChain_twoRunChain takes as
% hypotheses. (c) is GlobalTheory.BigOnBase.exists_tieCoord_base (withCoord, coord_lattice,
% coord_tuple, goodOp_classSum_base). (d) is RankRows in GlobalTheory.BigOn.Links2, from
% GlobalTheory.T.second_link_of_weight_rows. Steps 1 to 3 are on_base_atPlaces; Step 4 is
% onClauses_with_error_of_splitPart (onWith_of_split: iw_split_eq, d_split_eq).

\emph{Step 2: the group.} By Lemma~\ref{is-onto-recorded-numbers}(1), $r_0=k+m-1\geq1$ and $f$ is
the number $f$ of the member. By (a), the levels estimate and the $y$-bound of the page estimates in
(3a), and Lemma~\ref{is-places-level-index} with $r=r_0$, hypothesis (2) of
Lemma~\ref{is-places-combination} holds: $(1-y)|U|\le[P:U]$, $r_0\geq1$, and for every twisted place
$w$, $v_w(|U|)=e_wfr_0(r_0-1)/2$ if $p_w\in\{g_0,g_1\}$ and $v_w(|U|)=0$ otherwise.

\emph{Step 3: the combination.} Put $C_E=4C^2+25C+C_{b,1}+C_{b,2}+48$ and $N=N_p(C_E,C_H)$. By (3a)
with $C'=C_E$, $O$ satisfies (1) to (4) of the page estimates and $E_1(j)=C_EE_1^{u}(O)$ for every
lift $j$; (3b) gives the geometric estimates and the depth estimates for $N$; and (3c) to (3g) give
the rank at the twisted places, a twist bound, the share bound, the bad-place bound and the exponent
bound for $N$. With (a) to (d) and Step 2, Lemma~\ref{is-places-combination} gives real numbers
$\mathrm{err}(w,j)$ for which $S$ satisfies (1), (2) and (3) of
Definition~\ref{is-surjectivity-bound} for $N$, with $\mathrm{err}$ in place of $\mathrm{err}_0$.

\emph{Step 4: the rows.} The datum $S$ and the rows $S_0(p)$ have the same twisted places, tuple map
and rank $D$, and the same indices $I_w(j,\mu_w)$ and $I_{\bar w}(j,\mu_w)$; the functions $F_0$,
$J_0$ and $E_1$ of $N$ are the same for both. So $S_0(p)$ satisfies the three conditions with the
same numbers $\mathrm{err}(w,j)$.
\end{proof}

\subsection{Eigenvalues of theta lifts and the base links}\label{subsection-is-eigenvalues-base-links}

\begin{definition}\label{is-eigenvalue-lift-data}
Let $R$ be a realization of a definite-side datum $D$ with a face $F$
(Definition~\ref{is-realisation}). Write $E_h/F_h$ for the host field and $K$ for the subfield of
the host; $L$ for the coefficient field and $\sigma_0$ for the embedding of its base place; $f$, $r$
and $t$ for the ranks (Definition~\ref{th-ranks}) and $\mathcal T$ for the set of $f$ embeddings of
$E_h$ (Definition~\ref{th-definite-datum}(2)); $V_i$ and $W_i$ for the Hermitian spaces of the
$i$-th lift, $n_V=\operatorname{rk}V_i$, and $(p_\rho,q_\rho)$ for the signature of $W_i$ at
$\rho\in\mathcal T$; $\mu^0_i$, $\eta$ and $\xi_\mu$ for the Hecke characters and $Q_\mu$ for the
ramified sets of the tuples $\mu$; $P_g$ for the finite set of places of $E_h$ and
$G=\prod_{P\in P_g}GL_r(k_P)$; and $H_W(i,\sigma)$ for the Hecke modules. An element $U$ of
$H_W(i,\sigma)$ gives, for each tensor $v=\otimes_\rho v_\rho$ of vectors of the lowest $K$-types
$\kappa_\rho(\sigma)$ of the models of discrete series of $U(W_i)$ at $\sigma$, a function $U(v)$ on
$U(W_i)(\mathbf A)$; write $\langle\ ,\ \rangle_W$ for the Petersson pairing of such functions over
the fundamental domain of the realization, $\pi_\rho$ for the action of $U(p_\rho,q_\rho)$ on the
model at $\rho$, and $K_W\subset U(W_i)(\mathbf A_{F_h,f})$ for the level of $W_i$. For $\rho\in\mathcal
T$ let $a^+_\rho(\sigma)$ and $a^-_\rho(\sigma)$ be the positive and negative parts of the weights of
the $i$-th lift at the label of $\rho$ under $\sigma$, of lengths $p_\rho$ and $q_\rho$, and for a
decreasing sequence $a$ of integers let $\mathrm{hook}(a)$ be the product of the hook lengths of the
Young diagram with rows $a$.

The \emph{core conditions} at $F$ are conditions (1) to (6) of the characterization of $R$ at $F$
(Definition~\ref{is-realisation-characterization}). A tuple $\mu$ is \emph{prime to the host} if no place in $P_g\cup Q_\mu$
contains the integer $Q_Sqd\,\mathrm N(\mathfrak c_K)$ of the host; then $S_\mu$ is the set of places
of $F_h$ below $P_g\cup Q_\mu$. Suppose that $R$ satisfies the core conditions at $F$. \emph{Lift
data} for $R$ at $F$ consist of the following.
\begin{enumerate}
\item A lift $i$, a tuple $\mu$ prime to the host, an embedding $\sigma\colon L\to\mathbf C$ that
  agrees with $\sigma_0$ on $K$, and an automorphism $s$ of $\mathbf C$ with $\sigma=s\circ\sigma_0$.
\item A \emph{local level} $\mathcal P$: a subgroup $H_{\mathcal P}$ of $G$ that is a product of
  subgroups of the factors $GL_r(k_P)$, with a character.
\item Hecke characters $\chi'$ and $\eta'$ of $E_h$ whose values at the finite places, normalized
  with the weights $0$ and $\frac12$, are the images under $s$ of those of $\mu^0_i\xi_\mu$ and of
  $\eta$, also viewed as unitary characters of $\mathbf A_{E_h}^\times/E_h^\times$; and for each
  $P\in P_g\cup Q_\mu$ a character $m_P$ of $k_P^\times$, the restriction of $\chi'$ to the units at
  $P$.
\item A Weil datum $\omega$ for $(V_i,W_i)$ with characters $(\chi',\chi'\eta'^{-1})$
  (Definition~\ref{th-weil-datum}).
\item For each $\rho\in\mathcal T$ a basis of $W_i$ at $\rho$, orthonormal for its form of signature
  $(p_\rho,q_\rho)$ and adapted to the Cartan involution of the model at $\sigma$.
\item A \emph{level vector} $\phi$: the theta function obtained from a Fock vacuum vector by the level
  function of $\mathcal P$ and of the characters $m_P$ at the places of $P_g\cup Q_\mu$. For every
  polynomial $Q$ in the creation operators, the vector $\phi_Q=\phi\otimes Q$ of the Fock model is a
  theta function.
\item A \emph{coefficient datum}: a nonzero finite-dimensional Hilbert space $W_c$ with an irreducible
  continuous unitary representation of $U(V_i)(\mathbf A)$ that is trivial on the finite adelic
  points; and a linear map $\mathrm{io}$ from $\bigotimes_\rho\kappa_\rho(\sigma)\otimes W_c$ to the
  polynomials such that $\phi_{\mathrm{io}(\cdot)}$ intertwines $\omega$ with the actions of
  $U(V_i)(F_h\otimes\mathbf R)$ on $W_c$ and of the maximal compact subgroups of
  $U(W_i)(F_h\otimes\mathbf R)$ on the lowest $K$-types, and $\mathrm{io}(v\otimes w)$ is the delta
  polynomial $\Delta_i(\sigma)$ of the sequences $a^\pm_\rho(\sigma)$ for some unit vectors
  $v=\otimes_\rho v_\rho$ and $w$.
\end{enumerate}
The lift data determine an open subgroup $K_V$ of $U(V_i)(\mathbf A_{F_h,f})$, their \emph{level},
with a character $\chi_V$. A \emph{measure} for the lift data is a Haar measure $\nu$ on
$U(V_i)(\mathbf A)$ that gives the preimage of $K_V$ volume $1$, with a fundamental domain $\mathcal F$
for $U(V_i)(F_h)$. A \emph{form space} is a finite-dimensional Hilbert space $H_V$ identified with the
space of functions $\varphi$ on $U(V_i)(\mathbf A_{F_h,f})$ with values in $W_c$ that transform under
$U(V_i)(F_h)$ through $W_c$ and under $K_V$ through $\chi_V$, with the inner product
$\sum_y|\Gamma_y|^{-1}\langle\varphi(y),\varphi'(y)\rangle$ over representatives $y$ of the double
cosets, $\Gamma_y$ the stabilizer of $y$. For $h\in U(V_i)(\mathbf A)$ with finite part $h_f$ put
$\tilde\varphi(h)=h^{-1}\varphi(h_f)$. A \emph{theta lift} is a linear map $\theta\colon H_V\to
H_W(i,\sigma)$ such that for every $\varphi\in H_V$, every tensor $v$ of lowest $K$-type vectors and
every $g\in U(W_i)(\mathbf A)$ the function $h\mapsto\theta_\omega\big(h,g;\phi_{\mathrm{io}(v\otimes
\tilde\varphi(h))}\big)$ is integrable on $\mathcal F$ for $\nu$, with integral $\theta(\varphi)(v)(g)$.
For a system $c$ of Hecke eigenvalues at the finite places of $E_h$ not above $S_\mu$, the \emph{type
part} $H_W(i,\sigma)_c$ is the subspace of the $U\in H_W(i,\sigma)$ whose functions $U(v)$ have the
level of $\mathcal P$ and on which every Hecke operator at those places acts by its eigenvalue in $c$.
\end{definition}
\leanref{Development/AbelianVarietyTheory/Definitions/TwistedInnerProducts/RealLowestKTypeCoeffSquareIntegrable.lean}{GlobalTheory.LiftRows.IsLiftDataC}
\leanref{Development/AbelianVarietyTheory/Definitions/GaloisGeometry/NormalFam.lean}{GlobalTheory.Hk3.Core}
\leanref{Development/AbelianVarietyTheory/BlockTheory/DivisorBase/TwistedLevelValues.lean}{GlobalTheory.IndefiniteRowSideFullHecke.TwistedLevelValues.PrimeToHost}
\leanref{Development/AbelianVarietyTheory/BlockTheory/DivisorBase/TwistedLevelValues.lean}{GlobalTheory.IndefiniteRowSideFullHecke.TwistedLevelValues.IsLift}
\leanref{Development/AbelianVarietyTheory/BlockTheory/DivisorBase/TwistedLevelValues.lean}{GlobalTheory.IndefiniteRowSideFullHecke.TwistedLevelValues.typePart}

\begin{lemma}\label{is-eigenvalue-zeta-inputs}
Let $R$ be a realization of a definite-side datum $D$ with a face $F$ that satisfies the core
conditions, and let lift data for $R$ at $F$ be given, with a measure $(\nu,\mathcal F)$ and the
notation of Definition~\ref{is-eigenvalue-lift-data}. Identify $\prod_\rho U(p_\rho,q_\rho)$ with
$U(W_i)(F_h\otimes\mathbf R)$ through a matrix at each $\rho$ that carries the form of $W_i$ at $\rho$
to the standard form of signature $(p_\rho,q_\rho)$, let $m=\prod_\rho m_\rho$ be the
product of the standard Haar measures $m_\rho$ on the $U(p_\rho,q_\rho)$, let $x$ run over a set of
representatives of $U(W_i)(\mathbf A_{F_h,f})/K_W$, and let $\langle\ ,\ \rangle$ and $\|\ \|$ be the
$L^2$-pairing and norm of theta functions on a fundamental domain of covolume $1$ for the rational
points.
\begin{enumerate}
\item (Cusp forms.) Assume rapid decay of cusp forms. For every function $\varphi$ on
  $U(V_i)(\mathbf A_{F_h,f})$ with values in $W_c$ that transforms under $U(V_i)(F_h)$ through $W_c$
  and under $K_V$ through $\chi_V$, there is $U\in H_W(i,\sigma)$ such that $U(v)(g)$ is the integral
  over $\mathcal F$ of $h\mapsto\theta_\omega\big(h,g;\phi_{\mathrm{io}(v\otimes\tilde\varphi(h))}\big)$
  for every tensor $v$ of lowest $K$-type vectors and every $g\in U(W_i)(\mathbf A)$.
\item (Inner product formula.) Assume rapid decay of cusp forms, Theorem~\ref{lit-weil-siegel} and the
  continuation of Siegel Eisenstein series of Theorem~\ref{lit-doubling}(1), and suppose that $D$ has
  parameter choices. Let $U_1,U_2\in H_W(i,\sigma)$, let $v_1,v_2$ be tensors of lowest $K$-type
  vectors and $Q_1,Q_2$ polynomials, put $F_k=U_k(v_k)$ and
  $\Theta_k(h)=\langle\theta_\omega(h,\cdot\,;\phi_{Q_k}),F_k\rangle_W$ for $h\in U(V_i)(\mathbf A)$.
  Then the integrals below converge absolutely, and
  \[
   \int_{\mathcal F}\Theta_1\overline{\Theta_2}\,d\nu=\nu(\mathcal F)\int\sum_x
   \langle\omega(1,xb)\phi_{Q_1},\phi_{Q_2}\rangle\,\overline{\langle F_1(\cdot\,xb),F_2\rangle_W}\,dm(b).
  \]
\item (Finite zeta integrals.) Assume rapid decay of cusp forms, Bernstein-Zelevinsky finiteness and
  the unramified computation of the doubling integral for $GL$. For every system $c$ of Hecke
  eigenvalues with $H_W(i,\sigma)_c\ne0$ there are Satake parameters $\beta_{w,a}$ of $c$, for the
  finite places $w$ of $E_h$ not above $S_\mu$ and $1\le a\le r$, with $|\beta_{w,a}|\le\mathrm
  Nw^{(r-1)/2}$, such that, with $\alpha_{w,a}=\beta_{w,a}\chi'(\varpi_w)^{-1}\mathrm Nw^{-(1+t)/2}$,
  \[
   L^S(\alpha)=\prod_{w,a}(1-\alpha_{w,a})^{-1},\qquad
   d^S=\prod_{v\notin S_\mu}\prod_{j=t+1}^{n_V}(1-\varepsilon_v^j\mathrm Nv^{-j})^{-1},
  \]
  where $\varepsilon_v$ is $1$ if $v$ splits in $E_h$ and $-1$ otherwise, the following holds for
  every $U_1\in H_W(i,\sigma)_c$, every $U_2\in H_W(i,\sigma)$ and all tensors $v_1,v_2$:
  \[
   \sum_x\langle\omega(1,x)\phi,\phi\rangle\,\overline{\langle U_1(v_1)(\cdot\,x),U_2(v_2)\rangle_W}
   =|H_{\mathcal P}|\,\|\phi\|^2\,\frac{L^S(\alpha)}{d^S}\,\overline{\langle U_1(v_1),U_2(v_2)\rangle_W},
  \]
  the sum converging. In particular $|\alpha_{w,a}|\le\mathrm Nw^{(r-1)/2-(1+t)/2}$.
\item (Formal degrees.) Assume the formal degree of the discrete series in terms of the compact dual.
  For every $\rho$ and every unit vector $e\in\kappa_\rho(\sigma)$, the function
  $b\mapsto|\langle e,\pi_\rho(b)e\rangle|^2$ is integrable for $m_\rho$, its integral $J_\rho$ is
  positive, and
  \[
   \log J_\rho^{-1}\le2+r^2\log\big(2\sqrt2\,r^2n_V(\deg_\rho+2)\big),
  \]
  where $\deg_\rho$ is the sum of the entries of $a^+_\rho(\sigma)$ and of $a^-_\rho(\sigma)$.
\end{enumerate}
\end{lemma}
\leanref{Development/AbelianVarietyTheory/Definitions/Divisors/FormalDegreeUpperC.lean}{GlobalTheory.LowRow.LiftIsCuspC}
\leanref{Development/AbelianVarietyTheory/Proofs/DivisorProofs/AdelicLiftMemCuspFormMultiplicitySpaceStatement.lean}{GlobalTheory.T.adelicLift_mem_cuspForm_multiplicitySpace}
\leanref{Development/AbelianVarietyTheory/Definitions/TwistedInnerProducts/RealLowestKTypeCoeffSquareIntegrable.lean}{GlobalTheory.LiftRows.InnerProductFormulaC}
\leanref{Development/AbelianVarietyTheory/Proofs/DecoratedLatticeProofs/TwistedThetaLiftInnerProductFormulaOfPublishedStatement.lean}{GlobalTheory.T.twistedThetaLift_innerProductFormula_of_hypotheses}
\leanref{Development/AbelianVarietyTheory/Definitions/Divisors/FormalDegreeUpperC.lean}{GlobalTheory.LowRow.FiniteZetaC}
\leanref{Development/AbelianVarietyTheory/Proofs/TwistedOperators/FiniteZetaIntegralThetaLiftEqPartialLValueStatement.lean}{GlobalTheory.T.finiteZetaIntegral_thetaLift_eq_partialLValue}
\leanref{Development/AbelianVarietyTheory/Proofs/DivisorProofs/FiniteZetaIntegralEulerProductOfPartialLStatement.lean}{GlobalTheory.T.finiteZetaIntegral_eulerProduct_of_partialL}
\leanref{Development/AbelianVarietyTheory/Definitions/Divisors/FormalDegreeUpperC.lean}{GlobalTheory.LowRow.FormalDegreeUpperC}
\leanref{Development/AbelianVarietyTheory/Proofs/DivisorProofs/FormalDegreeRealPlaceUpperBoundOfHarishChandraStatement.lean}{GlobalTheory.T.formalDegreeRealPlaceUpperBound_of_harishChandra}
\begin{proof}
(1) The theta integral of $\tilde\varphi$ is continuous, invariant under the level $K_W$, killed by
the lowering operators at the infinite places because the lowering part kills the products of minors
of the Fock vectors, and cuspidal because the average of the level Fock vector over the Heisenberg
group of an isotropic subspace vanishes. By rapid decay of cusp forms \cite{HC68,MW95} it lies in the
multiplicity space of the model of $W_i$ at $\sigma$, which is the image of $H_W(i,\sigma)$ by the
core conditions on the models.

(2) This is the Rallis inner product formula \cite{Ral84,Li92} in the doubling form. By
Theorem~\ref{lit-weil-siegel}, the integral over $U(V_i)(F_h)\backslash U(V_i)(\mathbf A)$ of the theta
function of the doubled space is a Siegel Eisenstein series, continued by
Theorem~\ref{lit-doubling}(1), and the unfolding of the doubling integral turns its pairing with
$F_1\otimes\overline{F_2}$ into the integral over $U(W_i)(\mathbf A)$ on the right. Rapid decay of cusp
forms gives the absolute convergence.

(3) This is the unramified computation of the doubling zeta integral \cite{GPSR87,LR05}, with
Bernstein-Zelevinsky finiteness for the Satake parameters and their bound: the sum is
$|H_{\mathcal P}|\,\|\phi\|^2$ times the partial $L$-function of $c$ twisted by $\chi'^{-1}$ at
$(1+t)/2$, divided by $d^S$, times $\overline{\langle U_1(v_1),U_2(v_2)\rangle_W}$. Since $t\ge r+5$,
the partial $L$-function converges at $(1+t)/2$, and its Euler product is $L^S(\alpha)$.

(4) By Harish-Chandra's formula for the formal degree in terms of the compact dual
\cite{HC66,HII08} and Schur orthogonality \cite[Chapter~14]{Dix77}, $J_\rho^{-1}$ is the formal degree
of the discrete series at $\rho$, $\prod_{a<b}|\lambda_a-\lambda_b|/(b-a)$ for its Harish-Chandra
parameter $\lambda$, divided by the volume of the compact dual $U(r)$ for $m_\rho$. The parameter
$\lambda$ is strictly decreasing with $|\lambda_a-\lambda_b|\le\deg_\rho+n_V+3$, and the bound is an
elementary estimate of the product and the volume.
\end{proof}

\begin{lemma}\label{is-eigenvalue-local-factors}
Let $R$ be a realization of a definite-side datum $D$ with a face $F$ that satisfies the core
conditions, and let lift data for $R$ at $F$ be given, with the notation of
Definition~\ref{is-eigenvalue-lift-data} and Lemma~\ref{is-eigenvalue-zeta-inputs}. Let
$K_{\mathrm{lat}}\subset U(V_i)(\mathbf A_{F_h,f})$ be the stabilizer of the standard lattice of $V_i$,
and let $\operatorname{mass}(V_i)=\sum_y|\Gamma_y|^{-1}$, the sum over representatives $y$ of
$U(V_i)(F_h)\backslash U(V_i)(\mathbf A_{F_h,f})/K_{\mathrm{lat}}$ and $\Gamma_y$ the stabilizer of
$y$. For a finite place $v$ of $F_h$ put $d_v=\prod_{j=t+1}^{n_V}(1-\varepsilon_v^j\mathrm
Nv^{-j})^{-1}$.
\begin{enumerate}
\item (Level law.) $\omega(\kappa,1)\phi_Q=\chi_V(\kappa)^{-1}\phi_Q$ for every $\kappa\in K_V$, viewed
  in $U(V_i)(\mathbf A)$ with archimedean component $1$, and every polynomial $Q$.
\item (Fock coefficients.) For $x\in U(W_i)(\mathbf A_{F_h,f})$, $g=(g_\rho)\in\prod_\rho
  U(p_\rho,q_\rho)$, tensors $e,e'$ of lowest $K$-type vectors and $w,w'\in W_c$,
  \[
   \langle\omega(1,xg)\phi_{\mathrm{io}(e\otimes w)},\phi_{\mathrm{io}(e'\otimes w')}\rangle
   =\|\Delta_i(\sigma)\|^2\prod_\rho\langle e'_\rho,\pi_\rho(g_\rho)e_\rho\rangle\,\langle w',w\rangle\,
   \langle\omega(1,x)\phi,\phi\rangle,
  \]
  where $\|\Delta_i(\sigma)\|$ is the Fock norm of the polynomial $\Delta_i(\sigma)$.
\item (Delta norm.) $\|\Delta_i(\sigma)\|^2=\prod_\rho\mathrm{hook}(a^+_\rho(\sigma))\,
  \mathrm{hook}(a^-_\rho(\sigma))$.
\item (Level index.) $K_V\subset K_{\mathrm{lat}}$, and
  $[K_{\mathrm{lat}}:K_V]\,\|\phi\|^2=\prod_{v\in S_\mu}d_v^{-1}$.
\item (Volume.) If $[K_{\mathrm{lat}}:K_V]$ is finite and nonzero and $(\nu,\mathcal F)$ is a measure for
  the lift data, then $\nu(\mathcal F)=[K_{\mathrm{lat}}:K_V]\operatorname{mass}(V_i)$.
\item (Mass.) If $V_i$ is positive definite and has a basis orthonormal for its form that spans its
  standard lattice, then $\operatorname{mass}(V_i)\ge(|\mu(E_h)|^{n_V}\,n_V!)^{-1}$, where $\mu(E_h)$ is
  the group of roots of unity of $E_h$.
\end{enumerate}
\end{lemma}
\leanref{Development/AbelianVarietyTheory/Definitions/TwistedInnerProducts/RealLowestKTypeCoeffSquareIntegrable.lean}{GlobalTheory.LiftRows.SourceLevelLawC}
\leanref{Development/AbelianVarietyTheory/Proofs/TwistedOperators/FiniteDatumSourceLevelCharacterLawStatement.lean}{GlobalTheory.T.finiteDatum_sourceLevelCharacter_law}
\leanref{Development/AbelianVarietyTheory/Definitions/TwistedInnerProducts/RealLowestKTypeCoeffSquareIntegrable.lean}{GlobalTheory.LiftRows.FockCoefficientC}
\leanref{Development/AbelianVarietyTheory/Proofs/TwistedOperators/FockVectorMatrixCoefficientFormulaStatement.lean}{GlobalTheory.T.fockVector_matrixCoefficient_formula}
\leanref{Development/AbelianVarietyTheory/Proofs/DivisorProofs/InnerProductFormula.lean}{GlobalTheory.AdelicScalar.DeltaNorm}
\leanref{Development/AbelianVarietyTheory/Proofs/DivisorProofs/FinitenessInputsRealData.lean}{GlobalTheory.LowRowAdel.deltaNorm_all}
\leanref{Development/AbelianVarietyTheory/Definitions/TwistedInnerProducts/RealLowestKTypeCoeffSquareIntegrable.lean}{GlobalTheory.LiftRows.LevelIndexNormC}
\leanref{Development/AbelianVarietyTheory/Proofs/TwistedOperators/ThetaLevelIndexMulSchwartzNormFormulaStatement.lean}{GlobalTheory.T.thetaLevelIndex_mul_schwartzNorm_formula}
\leanref{Development/AbelianVarietyTheory/Proofs/DivisorProofs/VolEqIndexMulMass.lean}{GlobalTheory.VolIndex.vol_eq_index_mul_mass}
\leanref{Development/AbelianVarietyTheory/Proofs/DivisorProofs/InnerProductFormula.lean}{GlobalTheory.AdelicScalar.MassLower}
\leanref{Development/AbelianVarietyTheory/Proofs/DivisorProofs/UnitaryMassLowerBoundOfOrthonormalLatticeFrames.lean}{GlobalTheory.LowRow.StdMass.massLower_of_stdV}
\begin{proof}
Items (1) to (4) are computations in the Fock model at the infinite places and with the level
functions at the places of $S_\mu$; (2) uses the harmonic decomposition of the Fock space under the
compact groups, and (3) is the Fock norm of a product of minors. For (5), the preimage of $K_V$ has
volume $1$, so the preimage of $K_{\mathrm{lat}}$ has volume $[K_{\mathrm{lat}}:K_V]$, and a fundamental
domain for $U(V_i)(F_h)$ has the volume of that preimage times $\sum_y|\Gamma_y|^{-1}$. For (6), the
mass is at least the term $|\Gamma_1|^{-1}$ of the class of $1$, where $\Gamma_1$ is the group of
matrices $g$ with entries in $\mathcal O_{E_h}$ and $\bar g^{\mathsf T}g=1$, since the standard lattice
is $\mathcal O_{E_h}^{n_V}$ with the form $\sum_a\bar x_ax_a$. Each column $x$ of such a $g$ has
$\sum_a\bar x_ax_a=1$; each nonzero $\bar x_ax_a$ is a totally positive integer of $F_h$ and so has
trace at least $[F_h:\mathbf Q]$, so exactly one $x_a$ is nonzero, and it is a root of unity. Hence $g$
is a monomial matrix with entries in $\mu(E_h)$, and $|\Gamma_1|\le|\mu(E_h)|^{n_V}n_V!$.
\end{proof}

\begin{lemma}\label{is-eigenvalue-lift-computations}
Let $R$ be a realization of a definite-side datum $D$ with a face $F$ that satisfies the core
conditions, and let lift data for $R$ at $F$ be given, with a measure $(\nu,\mathcal F)$, a form space
$H_V$ and the notation of Definition~\ref{is-eigenvalue-lift-data} and
Lemma~\ref{is-eigenvalue-zeta-inputs}. Put $N=\prod_\rho\dim\kappa_\rho(\sigma)$.
\begin{enumerate}
\item The inner product of $H_W(i,\sigma)$ is $N\langle U,U'\rangle_{\mathrm{Pet}}$, where
  $\langle\ ,\ \rangle_{\mathrm{Pet}}$ is the Petersson inner product of the families of functions
  $U(v)$.
\item (Adjoints.) Suppose that the level law of Lemma~\ref{is-eigenvalue-local-factors}(1) holds and
  that the integrals of Lemma~\ref{is-eigenvalue-zeta-inputs}(2) converge. Let $\theta$ be a theta lift,
  $v$ a tensor of unit lowest $K$-type vectors and $(w_k)$ an orthonormal basis of $W_c$, and for $U\in
  H_W(i,\sigma)$ put $\Theta_k(U)(h)=\langle\theta_\omega(h,\cdot\,;\phi_{\mathrm{io}(v\otimes
  w_k)}),U(v)\rangle_W$. Then for $U_1,U_2\in H_W(i,\sigma)$
  \[
   \langle\theta^*U_1,\theta^*U_2\rangle=N^2\sum_k\int_{\mathcal
   F}\Theta_k(U_1)\overline{\Theta_k(U_2)}\,d\nu.
  \]
\item (Archimedean factorization.) Let $\phi'$ be a theta function, $v$ a tensor of lowest $K$-type
  vectors, $c_0,z\in\mathbf C$ and $U_1,U_2\in H_W(i,\sigma)$. Suppose that
  $\langle\omega(1,xg)\phi',\phi'\rangle=c_0\prod_\rho\langle v_\rho,\pi_\rho(g_\rho)v_\rho\rangle\,
  \langle\omega(1,x)\phi,\phi\rangle$ for all $x\in U(W_i)(\mathbf A_{F_h,f})$ and
  $g\in\prod_\rho U(p_\rho,q_\rho)$, and that for all tensors $v_1,v_2$ the sum
  $\sum_x\langle\omega(1,x)\phi,\phi\rangle\,\overline{\langle U_1(v_1)(\cdot\,x),U_2(v_2)\rangle_W}$
  converges to $z\,\overline{\langle U_1(v_1),U_2(v_2)\rangle_W}$. Then
  \[
   \int\sum_x\langle\omega(1,xb)\phi',\phi'\rangle\,\overline{\langle
   U_1(v)(\cdot\,xb),U_2(v)\rangle_W}\,dm(b)
   =c_0\,z\,\overline{\langle U_2,U_1\rangle_{\mathrm{Pet}}}\prod_\rho\int|\langle
   v_\rho,\pi_\rho(b)v_\rho\rangle|^2\,dm_\rho(b).
  \]
\item (Scalars.) Let $\theta\colon X\to Y$ be a linear map of finite-dimensional Hilbert spaces,
  $Y_0\subset Y$ a nonzero subspace, $n\ge1$ an integer, $a\in\mathbf R$ and $E\in\mathbf C$ with
  $a|E|>0$. If $\langle\theta^*y,\theta^*y'\rangle=naE\langle y,y'\rangle$ for all $y\in Y_0$ and
  $y'\in Y$, then $n^{-1}\theta\theta^*y=a|E|\,y$ for all $y\in Y_0$.
\end{enumerate}
\end{lemma}
\leanref{Development/AbelianVarietyTheory/Proofs/DivisorProofs/AdelCOfInnerProductFormula.lean}{GlobalTheory.LowRow.inner_HW_lowC}
\leanref{Development/AbelianVarietyTheory/Proofs/DivisorProofs/AdelCOfInnerProductFormula.lean}{GlobalTheory.LowRow.form_key_lowC}
\leanref{Development/AbelianVarietyTheory/Proofs/DivisorProofs/FormOfClassSums.lean}{GlobalTheory.AdelicScalar.form_eq_sum}
\leanref{Development/AbelianVarietyTheory/Proofs/DivisorProofs/ExistsLinearLift.lean}{GlobalTheory.AdelicScalar.zeta_eval}
\leanref{Development/AbelianVarietyTheory/Proofs/DivisorProofs/ExistsLinearLift.lean}{GlobalTheory.AdelicScalar.scalar_adjLev}
\begin{proof}
(1) is the normalization of the inner product in the core conditions on the models. (2) is a
computation from the definition of a theta lift: by the level law the integrand of $\theta(\varphi)$
transforms under the preimage of $K_V$ through $\chi_V$, so the integral over $\mathcal F$ is a sum over
the double cosets weighted by $|\Gamma_y|^{-1}$, and the adjoint of $\theta$ for the inner product of
$H_V$ is computed coordinate by coordinate in the basis $(w_k)$. (3) is a computation: the right
translate of $U_1(v)$ by $b$ is $U_1(\pi(b)v)$, the sum over $x$ then gives
$z\,\overline{\langle U_1(\pi(b)v),U_2(v)\rangle_W}$, and by the structure of the multiplicity space
$\langle U_1(\pi(b)v),U_2(v)\rangle_W=\langle U_2,U_1\rangle_{\mathrm{Pet}}\prod_\rho\langle
v_\rho,\pi_\rho(b)v_\rho\rangle$, so the integrand is $c_0z\,\overline{\langle
U_2,U_1\rangle_{\mathrm{Pet}}}\prod_\rho|\langle v_\rho,\pi_\rho(b)v_\rho\rangle|^2$. For (4),
$\theta\theta^*y$ is determined by
its inner products with all $y'$, so $\theta\theta^*y=naEy$ for $y\in Y_0$. Taking $y\ne0$ in $Y_0$,
$naE\|y\|^2=\|\theta^*y\|^2\ge0$, so $aE$ is real and nonnegative; as $a|E|>0$, $aE=a|E|$.
\end{proof}

\begin{lemma}\label{is-eigenvalue-size-bounds}
Let $D$ be a definite-side datum with host field $E_h/F_h$ and ranks $f$, $r$ and $t$, put
$n_V=\operatorname{rk}V_i$, and let $\mu(E_h)$ be the group of roots of unity of $E_h$.
\begin{enumerate}
\item $|\mu(E_h)|\le8f^2$.
\item Let $S$ be a finite set of finite places of $F_h$, and let $\alpha_{w,a}$, for the finite places
  $w$ of $E_h$ not above $S$ and $1\le a\le r$, be complex numbers with
  $|\alpha_{w,a}|\le\mathrm Nw^{(r-1)/2-(1+t)/2}$. Then $|\prod_{w,a}(1-\alpha_{w,a})^{-1}|\ge
  e^{-2rf}$.
\item For every finite set $S$ of finite places of $F_h$,
  $0<\prod_{v\notin S}d_v\le e^{4rf}$ and $0<\prod_{v\in S}d_v\le e^{4rf}$, with $d_v$ as in
  Lemma~\ref{is-eigenvalue-local-factors}.
\item Let $b$ be real and $i$ a lift at which the lifts have the sizes of the size data with $b$, and
  let $L_p$ be the logarithmic size of the size data (Lemma~\ref{is-hook-lower-bound}). Then
  $L_p\ge5.5$, $\log n_V\le L_p$, $b\ge0$ and $\log(3rb+2)\le2L_p$, and for every embedding $\sigma$
  of the coefficient field and every $\rho$ the number $\deg_\rho$ of
  Lemma~\ref{is-eigenvalue-zeta-inputs}(4) is at most $3rb$.
\end{enumerate}
\end{lemma}
\leanref{Development/AutomorphicArithmetic/CyclotomicBounds/RootsOfUnity/HostTorsionOrderLe.lean}{UpperTheory.Mu.host_torsionOrder_le}
\leanref{Development/AbelianVarietyTheory/Proofs/DivisorProofs/IntegralEqClassSumOfHom.lean}{GlobalTheory.AdelicScalar.eulerOff_lower}
\leanref{Development/AbelianVarietyTheory/Proofs/DivisorProofs/IntegralEqClassSumOfHom.lean}{GlobalTheory.AdelicScalar.dOff_bound}
\leanref{Development/AbelianVarietyTheory/Proofs/DivisorProofs/IntegralEqClassSumOfHom.lean}{GlobalTheory.AdelicScalar.dLoc_bound}
\leanref{Development/AbelianVarietyTheory/Proofs/DivisorProofs/ToNatFloorLe.lean}{GlobalTheory.AdelicScalar.SizeBounds.sizes}

\begin{lemma}\label{is-eigenvalue-numeric}
Let $f,n,r,w$ be integers with $f,r,w\ge1$, $n\ge8$, $2r\le n$ and $w\le8f^2$, and let $\mathcal T$ be a
set of $f$ elements. Let $L'$, $d_{\max}$, $M$, $\delta$, $I$, $N_\phi$, $E$, $d_O$ and $d_S$ be real
numbers, and for $\rho\in\mathcal T$ let $H_\rho$, $k_\rho$, $J_\rho$ and $\deg_\rho$ be real numbers,
such that $L'\ge5.5$, $\log n\le L'$, $d_{\max}\ge0$ and $\log(d_{\max}+2)\le2L'$;
$M\ge(w^nn!)^{-1}$; $\delta\ge1$; $IN_\phi=d_S^{-1}$; $E\ge e^{-2rf}$; $0<d_O\le e^{4rf}$ and
$0<d_S\le e^{4rf}$; and, for every $\rho$, $H_\rho>0$, $k_\rho\ge1$, $J_\rho>0$, $0\le\deg_\rho\le
d_{\max}$ and $\log J_\rho^{-1}\le2+r^2\log\big(2\sqrt2\,r^2n(\deg_\rho+2)\big)$. Put
\[
 l=\Big(\prod_\rho k_\rho\Big)\,\delta\,IM\,\Big(\prod_\rho H_\rho\Big)\,\frac{N_\phi}{d_O}\,
 \Big(\prod_\rho J_\rho\Big)\,E.
\]
Then $l>0$ and $\log l\ge\sum_\rho\log H_\rho-8fn^2L'$.
\end{lemma}
\leanref{Development/AutomorphicArithmetic/CyclotomicBounds/RootsOfUnity/MassNumeric.lean}{UpperTheory.Mu.numeric_bound}
\leanref{Development/AbelianVarietyTheory/Proofs/DivisorProofs/LamLowNumericalEnd.lean}{GlobalTheory.lamLow_numerical_end}
\begin{proof}
Since $IN_\phi=d_S^{-1}$, $l=P\cdot E/(d_Od_S)$ with $P=M\delta\prod_\rho(k_\rho H_\rho J_\rho)$, a
product of positive numbers.

\emph{Step 1: the main factor.} Put $M'=(w/2)^nM\delta$. Since $M\geq(w^nn!)^{-1}$ and $\delta\geq1$,
$M'\geq(2^nn!)^{-1}$, so $\log M'\geq-n\log2-n\log n$. For each $\rho$, $2\sqrt2\le3$ and
$\deg_\rho\le d_{\max}$ give $\log J_\rho^{-1}\le2+r^2\log\big(3r^2n(d_{\max}+2)\big)$. As
$3r^2\le n^2$, the argument of this logarithm is at most $n^3(d_{\max}+2)$, whose logarithm is at most
$5L'$; and $r^2\le n^2/4$. So $\log J_\rho^{-1}\le2+\frac54n^2L'$, and with $k_\rho\geq1$,
$\log(k_\rho H_\rho J_\rho)\geq\log H_\rho-2-\frac54n^2L'$. Summing over the $f$ elements $\rho$,
\[
 \log\big((w/2)^nP\big)\geq\sum_\rho\log H_\rho-n\log2-n\log n-f\Big(2+\frac54n^2L'\Big).
\]
Since $\log2\le1$, $\log n\le L'$, $fL'\geq1$, $n\geq8$ and $L'\geq5.5$, the terms $n\log2$,
$n\log n$ and $2f$ are each at most $fnL'$, and $8fnL'\le fn^2L'$; so the right side is at least
$\sum_\rho\log H_\rho-2fn^2L'$.

\emph{Step 2: the other factors.} Since $w/2\le(2f)^2$ and $\log(2f)\le2f-1$,
$\log\big((w/2)^n\big)\le4nf$. Moreover $\log E\geq-2rf$, $\log d_O\le4rf$ and $\log d_S\le4rf$.
Hence
\[
 \log l=\log\big((w/2)^nP\big)-\log\big((w/2)^n\big)+\log E-\log d_O-\log d_S
 \geq\sum_\rho\log H_\rho-2fn^2L'-4nf-10rf.
\]
Since $r\le n\le n^2$, $f\geq1$ and $L'\geq5.5$, $10rf+4nf\le14fn^2\le6fn^2L'$, and
$\log l\geq\sum_\rho\log H_\rho-8fn^2L'$.
\end{proof}

\begin{remark}\label{is-eigenvalue-numeric-direct}
The bound of Lemma~\ref{is-eigenvalue-numeric} can also be read off factor by factor. Taking
logarithms, $\log l-\sum_\rho\log H_\rho$ is at least
\[
 -n\log w-\log n!-2rf-8rf-\sum_\rho\Big(2+r^2\log\big(2\sqrt2\,r^2n(d_{\max}+2)\big)\Big).
\]
Each term is bounded with $w\le8f^2$, $\log n!\le n\log n\le nL'$, $2r\le n$, $\log r\le\log n\le
L'$ and $\log(d_{\max}+2)\le2L'$, and the sum of the bounds is at most $8fn^2L'$ because $n\ge8$ and
$L'\ge5.5$.
\end{remark}

\begin{lemma}\label{is-eigenvalue-adelic}
Assume rapid decay of cusp forms, the Weil-Siegel formula and the Siegel-Weil formula in the convergent
range (an assumption the proof does not use), the doubling method with its unramified computation for
$GL$, Bernstein-Zelevinsky finiteness and the formal degree of the compact dual
(Theorem~\ref{lit-bundle-unitary}). Let $D$ be a
definite-side datum that has parameter choices and standard spaces $V_i$, and let $R$ be a realization
of $D$ with a face $F$ that satisfies the core conditions. Let $b$ be real, and let lift data for $R$ at
$F$ be given (Definition~\ref{is-eigenvalue-lift-data}) whose lift $i$ has the sizes of the size data
with $b$, with a measure $(\nu,\mathcal F)$ and a form space $H_V$. Let $L_p$ be the logarithmic size of
the size data (Lemma~\ref{is-hook-lower-bound}). Then:
\begin{enumerate}
\item there is a theta lift $\theta\colon H_V\to H_W(i,\sigma)$;
\item for every theta lift $\theta$ and every system $c$ of Hecke eigenvalues with
  $H_W(i,\sigma)_c\ne0$ there is $l>0$ such that $|H_{\mathcal P}|^{-1}\theta\theta^*$ is
  multiplication by $l$ on $H_W(i,\sigma)_c$, and
  \[
   \log l\ge\sum_{\rho\in\mathcal
   T}\log\big(\mathrm{hook}(a^+_\rho(\sigma))\,\mathrm{hook}(a^-_\rho(\sigma))\big)
   -8fn_V^2L_p.
  \]
\end{enumerate}
\end{lemma}
\leanref{Development/AbelianVarietyTheory/Definitions/Divisors/FormalDegreeUpperC.lean}{GlobalTheory.LowRow.AdelC}
\leanref{Development/AutomorphicArithmetic/ThetaCorrespondence/ThetaPairLifts/LocalTypeDictionary/FinitenessInputsRealData.lean}{GlobalTheory.T.adelC_of_parameter_choice_stdV}
\leanref{Development/AbelianVarietyTheory/Proofs/DivisorProofs/AdelCOfInnerProductFormula.lean}{GlobalTheory.LowRow.adelC_of_innerProductFormula}
\begin{proof}
Lemmas~\ref{is-eigenvalue-zeta-inputs}, \ref{is-eigenvalue-local-factors}
and~\ref{is-eigenvalue-lift-computations} apply to the lift data; the mass bound applies because $D$
has standard spaces and $V_i$ is positive definite by the core conditions on the models. Use their
notation, and put $N=\prod_\rho\dim\kappa_\rho(\sigma)$.

(1) For $\varphi\in H_V$, a tensor $v$ and $g\in U(W_i)(\mathbf A)$, the function
$h\mapsto\theta_\omega\big(h,g;\phi_{\mathrm{io}(v\otimes\tilde\varphi(h))}\big)$ is integrable on
$\mathcal F$, because $V_i$ is positive definite, $\phi_{\mathrm{io}(\cdot)}$ consists of theta
functions and $\tilde\varphi$ is continuous. By Lemma~\ref{is-eigenvalue-zeta-inputs}(1) its integral
is $U(v)(g)$ for some $U\in H_W(i,\sigma)$. The integral is linear in $\varphi$, so a linear choice of
$U$ on a basis of $H_V$ is a theta lift.

(2) Let $\theta$ be a theta lift and $c$ a system with $H_W(i,\sigma)_c\ne0$. By the lift data there
are unit vectors $v=\otimes_\rho v_\rho$ and $w$ with $\mathrm{io}(v\otimes w)=\Delta_i(\sigma)$. Let
$(w_k)$ be an orthonormal basis of $W_c$, put $\phi_k=\phi_{\mathrm{io}(v\otimes w_k)}$, fix at each
$\rho$ a matrix as in Lemma~\ref{is-eigenvalue-zeta-inputs}, which exists by the core conditions on the
models, and let $J_\rho=\int|\langle v_\rho,\pi_\rho(b)v_\rho\rangle|^2\,dm_\rho(b)$. Let
$U_1\in H_W(i,\sigma)_c$ and $U_2\in H_W(i,\sigma)$.

By Lemma~\ref{is-eigenvalue-lift-computations}(2), whose hypotheses are
Lemma~\ref{is-eigenvalue-local-factors}(1) and the convergence in
Lemma~\ref{is-eigenvalue-zeta-inputs}(2),
$\langle\theta^*U_1,\theta^*U_2\rangle=N^2\sum_k\int_{\mathcal
F}\Theta_k(U_1)\overline{\Theta_k(U_2)}\,d\nu$.
By the inner product formula (Lemma~\ref{is-eigenvalue-zeta-inputs}(2)) the $k$-th integral is
$\nu(\mathcal F)$ times
$\int\sum_x\langle\omega(1,xb)\phi_k,\phi_k\rangle\,\overline{\langle
U_1(v)(\cdot\,xb),U_2(v)\rangle_W}\,dm(b)$.
Lemma~\ref{is-eigenvalue-local-factors}(2), with $e=e'=v$ and $w=w'=w_k$, a unit vector, gives the
first hypothesis of Lemma~\ref{is-eigenvalue-lift-computations}(3) for $\phi'=\phi_k$ with
$c_0=\|\Delta_i(\sigma)\|^2$. Lemma~\ref{is-eigenvalue-zeta-inputs}(3) gives the second, with
$z=|H_{\mathcal P}|\,\|\phi\|^2L^S(\alpha)/d^S$, because $U_1$ lies in the type part. So the $k$-th
integral is
\[
 \nu(\mathcal F)\,\|\Delta_i(\sigma)\|^2\,|H_{\mathcal P}|\,\|\phi\|^2\,\frac{L^S(\alpha)}{d^S}\,
 \overline{\langle U_2,U_1\rangle_{\mathrm{Pet}}}\,\prod_\rho J_\rho,
\]
independent of $k$. Put $H=\prod_\rho\mathrm{hook}(a^+_\rho(\sigma))\,\mathrm{hook}(a^-_\rho(\sigma))$,
which is $\|\Delta_i(\sigma)\|^2$ by Lemma~\ref{is-eigenvalue-local-factors}(3). Summing over the
$\dim W_c$ values of $k$ and using $\langle U_1,U_2\rangle=N\,\overline{\langle
U_2,U_1\rangle_{\mathrm{Pet}}}$ (Lemma~\ref{is-eigenvalue-lift-computations}(1)),
\[
 \langle\theta^*U_1,\theta^*U_2\rangle=|H_{\mathcal P}|\,a\,L^S(\alpha)\,\langle U_1,U_2\rangle,\qquad
 a=N\dim W_c\,\nu(\mathcal F)\,H\,\frac{\|\phi\|^2}{d^S}\prod_\rho J_\rho.
\]

Put $l=a|L^S(\alpha)|$. By Lemma~\ref{is-eigenvalue-local-factors}(4),
$I=[K_{\mathrm{lat}}:K_V]$ satisfies $I\|\phi\|^2=\prod_{v\in S_\mu}d_v^{-1}>0$, so $I$ is finite and
nonzero, and $\nu(\mathcal F)=I\operatorname{mass}(V_i)$ by
Lemma~\ref{is-eigenvalue-local-factors}(5). Apply Lemma~\ref{is-eigenvalue-numeric} with $n=n_V$,
$w=|\mu(E_h)|$, $L'=L_p$, $d_{\max}=3rb$, $M=\operatorname{mass}(V_i)$, $\delta=\dim W_c$, $I$,
$N_\phi=\|\phi\|^2$, $E=|L^S(\alpha)|$, $d_O=d^S$, $d_S=\prod_{v\in S_\mu}d_v$, and, for each $\rho$,
$H_\rho=\mathrm{hook}(a^+_\rho(\sigma))\,\mathrm{hook}(a^-_\rho(\sigma))$, which is a positive integer,
$k_\rho=\dim\kappa_\rho(\sigma)$, $J_\rho$ and $\deg_\rho$. Its hypotheses hold: $n_V\ge8$ and
$2r\le n_V$ by the inequalities of the rank datum; $|\mu(E_h)|\le8f^2$, the bounds for $L_p$, $3rb$
and $\deg_\rho$, the lower bound for $|L^S(\alpha)|$, which uses the bound for $\alpha$ in
Lemma~\ref{is-eigenvalue-zeta-inputs}(3), and the bounds for $d^S$ and $\prod_{v\in S_\mu}d_v$ are
Lemma~\ref{is-eigenvalue-size-bounds}; the mass bound is Lemma~\ref{is-eigenvalue-local-factors}(6);
and the bounds for $J_\rho$ are Lemma~\ref{is-eigenvalue-zeta-inputs}(4). It gives $l>0$ and $\log
l\ge\sum_\rho\log H_\rho-8fn_V^2L_p$. Finally Lemma~\ref{is-eigenvalue-lift-computations}(4), with
$Y_0=H_W(i,\sigma)_c$, $n=|H_{\mathcal P}|$ and $E=L^S(\alpha)$, shows that $|H_{\mathcal
P}|^{-1}\theta\theta^*$ is multiplication by $l$ on $H_W(i,\sigma)_c$.
\end{proof}

\begin{lemma}\label{is-eigenvalue-type-parts}
Let $R$ be a realization of a definite-side datum $D$ with a face $F$ that satisfies the core
conditions, except possibly the one on its chain of theta maps, with the notation of
Definition~\ref{is-eigenvalue-lift-data}. Let $i$ be a lift, $\sigma$ an embedding of $L$ and $S$ the
set of places of $F_h$ below $P_g$. For a system $c$ of Hecke eigenvalues at the finite places of $E_h$
not above $S$, let $H_W(i,\sigma)_c$
be the subspace of $H_W(i,\sigma)$ on which every Hecke operator of level $K_W$ at those places acts by
its eigenvalue in $c$, and call $c$ \emph{normalized} if it vanishes at the operators that are not
defined at that level.
\begin{enumerate}
\item The spaces $H_W(i,\sigma)_c$, for the normalized systems $c$, span $H_W(i,\sigma)$.
\item Each $H_W(i,\sigma)_c$ is stable under the action of $G$.
\item Let $\sigma$ agree with $\sigma_0$ on $K$, and let $s$ be an automorphism of $\mathbf C$ with
  $\sigma=s\circ\sigma_0$. Let $\chi'$ be a $\sigma$-conjugate of $\mu^0_i$: a Hecke character, trivial
  on $\mathbf A_{F_h}^\times$, with archimedean exponents $-e_i$ at the labels of $\sigma$, whose values
  at the finite ideles are the images under $\sigma$ of the rational values of $\mu^0_i$, and which is
  $\mu^0_i$ if $\sigma$ agrees with $\sigma_0$ on the field $K_s$ of the realization. Then $\chi'$ is
  the unitary character of a
  Hecke character whose values at the finite places, normalized with the weight $0$, are the images
  under $s$ of those of $\mu^0_i\xi_{\mu_0}$.
\item With $\sigma$ and $s$ as in (3), every $\sigma$-conjugate $\eta'$ of $\eta$, a Hecke character
  unramified everywhere whose values at uniformizers, divided by the square roots of the norms, are the
  images under $\sigma$ of those of $\eta$, is the unitary character of a Hecke character whose values,
  normalized with the weight $\frac12$, are the images under $s$ of those of $\eta$.
\end{enumerate}
Moreover, if the core conditions hold, the type parts of lift data with the lift $i$, the embedding
$\sigma$, the base tuple $\mu_0$ and the trivial level are the spaces $H_W(i,\sigma)_c$.
\end{lemma}
\leanref{Development/AbelianVarietyTheory/Proofs/DivisorProofs/PrincipalLevelEigenpartsAndConjugateCharactersStatement.lean}{GlobalTheory.T.principalLevel_eigenparts_and_conjugate_characters}
\leanref{Development/AbelianVarietyTheory/Proofs/DivisorProofs/TypePartBase.lean}{GlobalTheory.LiftLower.typePart_base}
\begin{proof}
(1) By the core conditions on the models, the functions of the elements of $H_W(i,\sigma)$ form the
multiplicity space of the model of level $K_W$, on which the Petersson pairing is positive definite by
the Schur multiplicity property. The Hecke operators of level $K_W$ at the places not above $S$
preserve this space and commute, and the adjoint of each for the Petersson pairing is the operator of
the dual label. A commuting family of operators on a finite-dimensional Hilbert space that is closed
under adjoints has a basis of joint eigenvectors, so these spaces span. (2) The action of $G$
commutes with these Hecke
operators. (3) and (4) compare the local components of the characters at each finite place, where
$\sigma=s\circ\sigma_0$ carries the values of $\mu^0_i$ and $\eta$ to those of $\chi'$ and $\eta'$;
for (3), $\xi_{\mu_0}$ is trivial by the core conditions on the level characters. For the last
assertion, $Q_{\mu_0}$ is empty by the core conditions on the level
characters, so $S_{\mu_0}$ is the set of places below $P_g$, and the trivial level has the level $K_W$.
\end{proof}

\begin{lemma}\label{is-eigenvalue-projection}
Let $T\colon X\to Y$ be a linear map of finite-dimensional complex Hilbert spaces, let $m>0$, and let
$(Y_k)$ be a family of subspaces that spans $Y$ such that, for each $k$, $TT^*$ is multiplication by a
scalar $l_k\ge m$ on $Y_k$. Let $e$ be an orthogonal projection of $Y$, that is $e^2=e=e^*$. Then:
\begin{enumerate}
\item $m\|u\|^2\le\|T^*u\|^2$ for every $u\in Y$;
\item if $x\in Y$ is nonzero and $l>0$ with $(eT)(eT)^*x=lx$, then $m\le l$;
\item $m\|x\|^2\le\|eTx\|^2$ for every $x\in X$ orthogonal to the kernel of $eT$.
\end{enumerate}
\end{lemma}
\leanref{Development/AbelianVarietyTheory/Proofs/DivisorProofs/MulNormSqLeNormProjApplySqOfSpanningParts.lean}{GlobalTheory.mul_norm_sq_le_norm_adjoint_sq_of_spanning_parts}
\leanref{Development/AbelianVarietyTheory/Proofs/DivisorProofs/HookPlusLeMainTerm.lean}{GlobalTheory.LowTopBound.le_eigenvalue}
\leanref{Development/AbelianVarietyTheory/Proofs/DivisorProofs/MulNormSqLeNormProjApplySqOfSpanningParts.lean}{GlobalTheory.mul_norm_sq_le_norm_proj_apply_sq_of_spanning_parts}
\leanref{Development/AbelianVarietyTheory/Proofs/DivisorProofs/ExistsFullLift.lean}{GlobalTheory.LowForm.lower_of_parts}
\begin{proof}
(1) The operator $TT^*$ is self-adjoint, and $Y$ is spanned by eigenvectors of $TT^*$ with eigenvalues
at least $m$, so every eigenvalue of $TT^*$ is at least $m$. Hence
$\|T^*u\|^2=\langle TT^*u,u\rangle\ge m\|u\|^2$.

(2) Since $(eT)^*=T^*e$, $lx=eTT^*ex$ lies in the image of $e$, so $ex=x$. Then
$l\|x\|^2=\langle eTT^*x,x\rangle=\|T^*x\|^2\ge m\|x\|^2$ by (1), and $m\le l$ as $x\ne0$.

(3) The orthogonal complement of the kernel of $eT$ is the image of $(eT)^*=T^*e$, so $x=T^*u$ with
$u=ey$ for some $y$. Then $\|x\|^2=\langle TT^*u,u\rangle=\langle eTx,u\rangle\le\|eTx\|\,\|u\|$, and
$m\|u\|^2\le\|T^*u\|^2=\|x\|^2$ by (1). Hence $\|x\|^2\le\|eTx\|\,\|x\|/\sqrt m$, which is (3).
\end{proof}

\begin{definition}\label{is-eigenvalue-conjugate-datum}
Let $R$ be a realization of a definite-side datum $D$ with a face $F$, with the notation of
Definition~\ref{is-eigenvalue-lift-data}, let $i$ be a lift and $\sigma$ an embedding of $L$. Write
$j_V$ for the $\sigma_0$-semilinear map from the lift forms at $(i,\mu_0)$ to $H_V(i,\mu_0)$
(Definition~\ref{is-realisation}(4)); $\langle x,u\rangle_\sigma$ for the image under $\sigma$ of the
class-sum pairing of lift forms $x,u$ at $(i,\mu_0)$; and $e$ for the operator $\sum_ge_C(g)g$ of
$H_W(i,\sigma)$, where $e_C$ is the central idempotent of the group algebra of $G$ of the realization.
Let $c_0$ be the coefficient map of a family tie of $R$, which identifies the values of the forms in
$H_V(i,\mu_0)$ with rational coefficient vectors. A \emph{conjugate datum} at $\sigma$ for the lift
$i$ and the tie consists of:
\begin{enumerate}
\item a $\sigma$-conjugate $\chi'$ of $\mu^0_i$ and a $\sigma$-conjugate $\eta'$ of $\eta$, as in
  Lemma~\ref{is-eigenvalue-type-parts}(3) and~(4);
\item a coefficient datum $W_c$ (Definition~\ref{is-eigenvalue-lift-data}(7)) with a
  $\sigma$-semilinear isomorphism from the rational coefficient space of $V_i$ onto $W_c$ that
  intertwines the action of $U(V_i)(F_h)$ and carries the rational pairing, followed by $\sigma$, to
  the inner product;
\item a form space $H$, as in Definition~\ref{is-eigenvalue-lift-data}, for the level of $V_i$ at
  the base tuple with the character whose values are the images under $\sigma$ of those of the
  character of the realization;
\item a $\sigma$-semilinear map $j$ from the lift forms at $(i,\mu_0)$ to $H$ such that, for every lift
  form $x$ and every $h$, the values of $j_V(x)$ and $j(x)$ at $h$ are the images of one rational
  coefficient vector under $c_0$ and under the isomorphism of (2);
\item a linear map $\mathrm{th}\colon H\to H_W(i,\sigma)$ such that, for every $\varphi\in H$ and every
  tensor $v$ of lowest $K$-type vectors, $\mathrm{th}(\varphi)(v)$ is $e$ applied to the theta integral
  $g\mapsto\int\theta_\omega\big(h,g;\phi_{\mathrm{io}(v\otimes\tilde\varphi(h))}\big)\,d\nu(h)$ over
  the fundamental domain of the realization at $(i,\mu_0)$, for a Weil datum $\omega$ with characters
  $(\chi',\chi'\eta'^{-1})$, a level vector $\phi$ of the trivial level, a basis as in
  Definition~\ref{is-eigenvalue-lift-data}(5) and a map $\mathrm{io}$ as in
  Definition~\ref{is-eigenvalue-lift-data}(7).
\end{enumerate}
\end{definition}
\leanref{Development/AbelianVarietyTheory/BlockTheory/DivisorBase/WeilGramSignature.lean}{GlobalTheory.IndefiniteRowSideFullHecke.WeilGramSignature.IsConjDatum}
\leanref{Development/AbelianVarietyTheory/BlockTheory/DivisorBase/WeilGramSignature.lean}{GlobalTheory.IndefiniteRowSideFullHecke.WeilGramSignature.IsConjForms}
\leanref{Development/AbelianVarietyTheory/BlockTheory/DivisorBase/WeilGramSignature.lean}{GlobalTheory.IndefiniteRowSideFullHecke.WeilGramSignature.IsConjLift}

\begin{lemma}\label{is-baselinks-rational-tests}
Let $K$ be a field with an involution $a\mapsto a^*$, and let $\sigma\colon K\to\mathbf C$ be an
embedding.
\begin{enumerate}
\item Let $E$ be a complex vector space with linear functionals $\psi_1,\dots,\psi_N$ that separate
  its points, and let $f\colon\mathbf C^d\to E$ be an injective linear map with
  $\psi_i(f(e_l))\in\sigma(K)$ for all $i$ and $l$. If $c\in\mathbf C^d$ and
  $\psi_i(f(c))\in\sigma(K)$ for every $i$, then $c\in\sigma(K)^d$.
\item Let $F$ be a finite-dimensional $K$-vector space with a form $h$ that is semilinear for the
  involution in the first variable and linear in the second, let $V$ and $E$ be finite-dimensional
  complex Hilbert spaces, and let $j\colon F\to V$ be a $\sigma$-semilinear map that carries a basis
  of $F$ to a basis of $V$ and satisfies $\langle j(x),j(y)\rangle=\sigma(h(x,y))$ for all
  $x,y\in F$. Let $\Theta\colon V\to E$ be a linear map, and suppose that $E$ has linear functionals
  $\psi_1,\dots,\psi_N$ that separate its points with $\psi_i(\Theta(j(x)))\in\sigma(K)$ for every
  $i$ and every $x\in F$. Then $\Theta^*(E)$ is spanned by $j(Z)$, where
  $Z=\{x\in F:j(x)\in\Theta^*(E)\}$.
\end{enumerate}
\end{lemma}
\leanref{Formalism/LocalToGlobal/IndefiniteReal/RationalAdjoints/CoordRat.lean}{Formalism.IndefR.Rat.coord_rat}
\leanref{Formalism/LocalToGlobal/IndefiniteReal/RationalAdjoints/RangeAdjointRat.lean}{Formalism.IndefR.Rat.range_adjoint_rat}
\begin{proof}
(1) Let $A$ be the $N\times d$ matrix over $K$ with $\sigma(A_{il})=\psi_i(f(e_l))$. If $Ax=0$ for
some $x\in K^d$, then $\psi_i(f(\sigma(x)))=\sigma((Ax)_i)=0$ for every $i$, so $f(\sigma(x))=0$ and
$x=0$. Hence $A$ has a left inverse $B$ over $K$. The vector $\sigma(A)c=(\psi_i(f(c)))_i$ is
$\sigma(b)$ for some $b\in K^N$, so $c=\sigma(B)\sigma(A)c=\sigma(Bb)$.

(2) Let $b_1,\dots,b_d$ be the basis; we may assume $d\ge1$. Comparing
$\langle j(ab_1),j(b_1)\rangle=\overline{\sigma(a)}\langle j(b_1),j(b_1)\rangle$ with
$\sigma(h(ab_1,b_1))=\sigma(a^*)\langle j(b_1),j(b_1)\rangle$ gives
$\sigma(a^*)=\overline{\sigma(a)}$ for every $a\in K$. Let $w_i\in E$ represent $\psi_i$, that is
$\langle w_i,u\rangle=\psi_i(u)$ for $u\in E$; the $w_i$ span $E$ because the $\psi_i$ separate its
points. Fix $i$ and write $\Theta^*w_i=\sum_kc_kj(b_k)$. The functionals
$\langle j(b_l),\cdot\,\rangle$ separate the points of $V$, take the values $\sigma(h(b_l,b_k))$ at
the $j(b_k)$, and take at $\Theta^*w_i$ the value
$\langle\Theta(j(b_l)),w_i\rangle=\overline{\psi_i(\Theta(j(b_l)))}$, which lies in $\sigma(K)$
because $\overline{\sigma(a)}=\sigma(a^*)$. By (1), $c_k=\sigma(a_k)$ with $a_k\in K$, so
$\Theta^*w_i=j(\sum_ka_kb_k)$ and $\sum_ka_kb_k\in Z$. Hence $\Theta^*(E)$, which the $\Theta^*w_i$
span, lies in the span of $j(Z)$. The other inclusion holds by the definition of $Z$.
\end{proof}

\begin{lemma}\label{is-baselinks-group-algebra}
Let $G$ be a finite group, $K$ a field of characteristic $0$ with an embedding
$\sigma\colon K\to\mathbf C$, and $F$ a $K$-vector space with a representation $r$ of $G$. Let $V$
and $E$ be finite-dimensional complex Hilbert spaces, let $j\colon F\to V$ be a $\sigma$-semilinear
map whose image spans $V$, and let $\rho(g)$, for $g\in G$, be linear operators on $V$ with
$j(r(g)x)=\rho(g)j(x)$ and $\langle j(r(g)x),j(r(g)y)\rangle=\langle j(x),j(y)\rangle$ for all
$x,y\in F$. Let $G$ act unitarily on $E$, and let $\Theta\colon V\to E$ be a linear map with
$\Theta(j(r(g)x))=g\,\Theta(j(x))$ for all $g$ and $x$. For $a=\sum_ga_gg$ in the group algebra
$K[G]$ put $a_F=\sum_ga_gr(g)$, $a_V=\sum_g\sigma(a_g)\rho(g)$ and $a_E=\sum_g\sigma(a_g)g$.
\begin{enumerate}
\item $\rho$ is a unitary representation of $G$, $\Theta\rho(g)=g\Theta$ for every $g\in G$,
  $a_V\Theta^*=\Theta^*a_E$ for every $a\in K[G]$, and $a_V(V)$ is spanned by $j(a_F(F))$.
\item Let $e\in K[G]$ with $e_F^2=e_F$. For every $u\in e_V(V)$ there is
  $u'\in\Theta^*(E)\cap e_V(V)$ with $\Theta(u')=\Theta(u)$.
\item Let $e\in K[G]$ with $e_F^2=e_F$, and suppose that $\Theta^*(E)$ is spanned by $j(Z)$, where
  $Z=\{x\in F:j(x)\in\Theta^*(E)\}$. Then $\Theta^*(E)\cap e_V(V)$ is spanned by
  $j(Z\cap e_F(F))$.
\item Let $X$ be a finite set of functions $G\to K$, and put
  $e_X=\sum_{\chi\in X}\frac{\chi(1)}{|G|}\sum_{g\in G}\chi(g^{-1})\,g$. Let $P\subset G$ be a subgroup
  with a character $\psi\colon P\to K^\times$, and put $e_{P,\psi}=|P|^{-1}\sum_{h\in P}\psi(h)^{-1}h$.
  If $v\in E$ satisfies $(e_X)_Ev=v$ and $hv=\sigma(\psi(h))v$ for every $h\in P$, then
  $\Theta^*v\in(e_Xe_{P,\psi})_V(V)$.
\item Let $h\in G$, $c\in K$ and $S\subset F$ with $r(h)x=cx$ for every $x\in S$. Then
  $h\,\Theta(u)=\sigma(c)\Theta(u)$ for every $u$ in the span of $j(S)$.
\end{enumerate}
\end{lemma}
\leanref{Formalism/LocalToGlobal/IndefiniteReal/LiftGoodPartComplexEqRange.lean}{Formalism.IndefR.CxPass.LevelOperators.opV_adj}
\leanref{Formalism/LocalToGlobal/IndefiniteReal/LiftGoodPartComplexEqRange.lean}{Formalism.IndefR.CxPass.LevelOperators.Hyp.range_opV}
\leanref{Formalism/LocalToGlobal/IndefiniteReal/LiftGoodPartComplexEqRange.lean}{Formalism.IndefR.CxPass.LevelOperators.Hyp.dn}
\leanref{Formalism/LocalToGlobal/IndefiniteReal/LiftGoodPartComplexEqRange.lean}{Formalism.IndefR.CxPass.LevelOperators.Hyp.ratLev}
\leanref{Formalism/LocalToGlobal/IndefiniteReal/LiftGoodPartComplexEqRange.lean}{Formalism.IndefR.CxPass.LevelOperators.Hyp.up}
\leanref{Formalism/LocalToGlobal/IndefiniteReal/LiftGoodPartComplexEqRange.lean}{Formalism.IndefR.CxPass.LevelOperators.Hyp.type}
\begin{proof}
(1) For $x,y\in F$ and $g,g'\in G$, $\langle\rho(g)j(x),\rho(g)j(y)\rangle=\langle j(x),j(y)\rangle$
and $\rho(gg')j(x)=j(r(g)r(g')x)=\rho(g)\rho(g')j(x)$. Since $j(F)$ spans $V$, $\rho$ is a unitary
representation. In the same way $\Theta\rho(g)=g\Theta$. Taking adjoints,
$\rho(g)^{-1}\Theta^*=\Theta^*g^{-1}$ for every $g$, so $a_V\Theta^*=\Theta^*a_E$. Finally
$a_Vj(x)=j(a_Fx)$ because $j$ is $\sigma$-semilinear, so $a_V(V)=a_V(\operatorname{span}j(F))$ is
spanned by $j(a_F(F))$.

(2) Since $e_Vj(x)=j(e_Fx)$, $e_V$ is idempotent on $j(F)$, hence on $V$, and $e_Vu=u$. Write
$u=p+k$ with $p\in\Theta^*(E)$ and $k$ orthogonal to $\Theta^*(E)$. Then $\Theta k=0$, since
$\langle\Theta k,\Theta k\rangle=\langle k,\Theta^*\Theta k\rangle=0$. Put $u'=e_Vp$, which lies in
$\Theta^*(E)$ by (1). As $\Theta e_Vk=e_E\Theta k=0$, $\Theta(u')=\Theta(e_Vu)=\Theta(u)$.

(3) If $x\in Z\cap e_F(F)$, say $x=e_Fy$, then $j(x)\in\Theta^*(E)$ and $j(x)=e_Vj(y)\in e_V(V)$.
Conversely, let $v\in\Theta^*(E)\cap e_V(V)$. Then $v=e_Vv$, and $v$ is a combination of vectors
$j(z)$ with $z\in Z$, so $v$ is a combination of the vectors $e_Vj(z)=j(e_Fz)$. Each $e_Fz$ lies in
$e_F(F)$, and in $Z$ because $j(e_Fz)=e_Vj(z)$ lies in $e_V\Theta^*(E)\subset\Theta^*(E)$ by (1).

(4) Since $hv=\sigma(\psi(h))v$ for $h\in P$, $(e_{P,\psi})_Ev=v$. Moreover $(ab)_V=a_Vb_V$ for
$a,b\in K[G]$, because both sides send $j(x)$ to $j((ab)_Fx)$ and $j(F)$ spans $V$. By (1), applied to
$e_{P,\psi}$ and to $e_X$,
$(e_Xe_{P,\psi})_V\Theta^*v=(e_X)_V\Theta^*(e_{P,\psi})_Ev=(e_X)_V\Theta^*v=\Theta^*(e_X)_Ev=\Theta^*v$.

(5) For $x\in S$, $h\,\Theta(j(x))=\Theta(j(r(h)x))=\Theta(j(cx))=\sigma(c)\Theta(j(x))$; both sides
are linear in $j(x)$.
\end{proof}

\begin{lemma}\label{is-baselinks-lift-cut}
Let $\theta\colon F\to W$ be a linear map of finite-dimensional complex Hilbert spaces, let
$u_1,\dots,u_d\in F$ be linearly independent, let $\Gamma\subset F$ and $W_0\subset W$ be subspaces,
and let $(W_p)_{p\in\Pi}$ be a finite family of subspaces of $W$ with positive numbers $\lambda_p$.
Suppose that:
\begin{enumerate}
\item the $u_k$ span $\theta^*(W)\cap\Gamma$;
\item $\theta(\Gamma)\subset W_0$ and $\theta^*(W_0)\subset\Gamma$;
\item $\theta\theta^*w=\lambda_pw$ for every $p$ and every $w\in W_p\cap W_0$, and the spaces
  $W_p\cap W_0$ span $W_0$.
\end{enumerate}
Then the map $t\colon\mathbf C^d\to W$, $c\mapsto\theta(\sum_kc_ku_k)$, is injective with image
$W_0$.
\end{lemma}
\leanref{Formalism/LocalToGlobal/IndefiniteReal/EqZeroOfRangeAdjoint.lean}{Formalism.IndefR.Cx.CxData.cutLifts}
\begin{proof}
If $t(c)=0$, then $\sum_kc_ku_k$ lies in $\theta^*(W)$ by (1) and in $\ker\theta$, the orthogonal
complement of $\theta^*(W)$; so $\sum_kc_ku_k=0$ and $c=0$. By (1) and (2),
$t(c)\in\theta(\Gamma)\subset W_0$. Conversely, by (3) it suffices to show that each
$w\in W_p\cap W_0$ lies in the image of $t$. By (2), $\theta^*w\in\theta^*(W)\cap\Gamma$, so
$\theta^*w=\sum_kc_ku_k$ for some $c$ by (1), and
$w=\lambda_p^{-1}\theta\theta^*w=t(\lambda_p^{-1}c)$ by (3).
\end{proof}

\begin{lemma}\label{is-baselinks-lifted-parts}
Let $\theta\colon A\to B$ be a linear map of finite-dimensional complex Hilbert spaces, let
$(T_i)_{i\in I}$ be commuting operators on $A$ such that the adjoint of each $T_i$ is some $T_{i'}$,
and let $(T'_i)_{i\in I}$ be operators on $B$ with $\theta T_i=T'_i\theta$. Let $\Gamma\subset A$ be
a subspace stable under every $T_i$, and $\Gamma'\subset B$ a subspace with
$\theta(\Gamma)\subset\Gamma'$. For $c\in\mathbf C^I$ let $A_c$ and $B_c$ be the joint eigenspaces
of the $T_i$ and of the $T'_i$ for the eigenvalues $c_i$, and call $c$ \emph{lifted} if
$A_c\cap\Gamma\ne0$ and $B_c\cap\Gamma'\ne0$. If $A_c\cap\Gamma\subset\theta^*(B)$ for every lifted
$c$, then
\[
 (\ker\theta)^\perp\cap\Gamma=\sum_{c\text{ lifted}}A_c\cap\Gamma.
\]
\end{lemma}
\leanref{Formalism/LocalToGlobal/IndefiniteReal/NormalDecomposition/Decomp.lean}{Formalism.IndefR.U2.decompOfNormal}
\leanref{Formalism/LocalToGlobal/IndefiniteReal/NormalDecomposition/DecompOfNormalStatement.lean}{Formalism.IndefR.U2.liftsFull_of}
\begin{proof}
Each $T_i$ commutes with its adjoint, and $\Gamma$ is stable under the $T_i$ and their adjoints. A
commuting family of normal operators on the finite-dimensional Hilbert space $\Gamma$ has an
orthonormal basis of joint eigenvectors, so $\Gamma$ is the orthogonal sum of the nonzero spaces
$A_c\cap\Gamma$. Let $\Gamma_+$ be the sum of those with $c$ lifted and $\Gamma_-$ the sum of the
others. If $A_c\cap\Gamma\ne0$ and $c$ is not lifted, then $\theta(A_c\cap\Gamma)\subset
B_c\cap\Gamma'=0$; so $\Gamma_-\subset\ker\theta$. Let $x\in(\ker\theta)^\perp\cap\Gamma$ and write
$x=l+m$ with $l\in\Gamma_+$ and $m\in\Gamma_-$. Then $0=\langle m,x\rangle=\langle m,m\rangle$, so
$x=l\in\Gamma_+$. Conversely, for lifted $c$,
$A_c\cap\Gamma\subset\theta^*(B)=(\ker\theta)^\perp$.
\end{proof}

\begin{lemma}\label{is-baselinks-definite-coordinates}
Let $R$ be a realization of a definite-side datum $D$ for a family of definite forms, with a face $F$
(Definition~\ref{is-realisation}), and let $K_0\subset L$ be the subfield on which the involution is
complex conjugation at every embedding. Let $B$ be a space of forms of the family, the row forms at a
tuple or the lift forms at a lift and a tuple, with its rational basis, its class-sum pairing, its action
of $G$ and its Hecke and inert Hecke operators. Let $X$ be a finite set of functions $G\to L$, let
$P\subset G$ be a subgroup with a character $\psi$, let $e_X$ be as in
Lemma~\ref{is-baselinks-group-algebra}(4), put $e_{P,\psi^{-1}}=|P|^{-1}\sum_{h\in P}\psi(h)h$, and let
$\Gamma$ be the image of $e_Xe_{P,\psi^{-1}}$ on $B$. For a subspace $Z\subset B$, the \emph{definite
side} of $B$ at $(X,P,\psi)$ for $Z$ is built from a basis of $\Gamma$, chosen with coordinates in $K_0$
in the rational basis whenever $\Gamma$ has such a basis, the matrix $H_\Gamma$ of the
class-sum pairing divided by $|P|$ in that basis, and a basis matrix $E$ of the coordinate space
$Z_0$ of $Z\cap\Gamma$ in that basis, chosen with entries in $K_0$ whenever $Z_0$ has such a basis: its
coordinate map $c$ sends $y$ to the form in $\Gamma$ with coordinates $Ey$, its Hermitian matrix is
$H=E^*H_\Gamma E$, and its metric $X_Z$ is the quotient metric of $H_\Gamma$ along the orthogonal
projection onto $Z_0$. At a level $b$, with the subgroup $P_b$ and the character $\psi_b$ of the face,
the complex data of $R$ and $F$ (Definition~\ref{is-complex-data}) take the definite side at
$(\mathcal X,P_b,\psi_b)$, for the set $\mathcal X$ of the characters of $G$ in the family, of the lift
forms at $(j,t)$ for $Z_{j,t}$, with $c=c_{b,j,t}$, $H=H_T(j,t)$ and $X_Z=X_T(j,t)$, and of the row forms
at $t$ for $Z_t$; the image $\Gamma_b(j,t)$ of $e_{\mathcal X}e_{P_b,\psi_b^{-1}}$ on the lift forms at
$(j,t)$ is their \emph{level-$b$ part}.
\begin{enumerate}
\item The image of $c$ is $Z\cap\Gamma$.
\item Suppose that the involution of $L$ is complex conjugation at $v_0$, that the functions in $X$ are
  the characters of absolutely irreducible representations of $G$, and that $B$ satisfies the conditions
  that the family conditions impose on its forms. Then $X_{Z,v_0}$ is invertible, and for all $x,y$ the
  image under $\sigma_0$ of the class-sum pairing of $c(x)$ and $c(y)$ is
  $|P|\,(X_{Z,v_0}^{-1}\sigma_0(x))^*(X_{Z,v_0}^{-1}\sigma_0(y))$.
\item Suppose the hypotheses of (2), and let $(T_i)$ be a family of operators on $B$ whose matrices in
  the rational basis have entries in $K_0$, that contains the adjoint of each of its members for the
  class-sum pairing, whose members commute with each other and with $G$, and each of whose members has
  generalized eigenspaces that span $B$. If $Z\cap\Gamma$ is the intersection of $\Gamma$ with the sum of
  the joint generalized eigenspaces of the $T_i$ for the systems of eigenvalues in some set, then the
  entries of $H$ lie in $K_0$ and $|y|^2_{X_Z,v}=\sigma_v(y)^*\sigma_v(H)\sigma_v(y)$ for every infinite
  place $v$ of $L$ and every $y$.
\end{enumerate}
\end{lemma}
\leanref{Formalism/LocalToGlobal/IndefiniteReal/LiftGoodPartComplexEqRange.lean}{Formalism.IndefR.CxPass.LevelOperators.range_toF}
\leanref{Formalism/LocalToGlobal/IndefiniteReal/LiftGoodPartComplexEqRange.lean}{Formalism.IndefR.CxPass.LevelOperators.goodR_eq}
\leanref{Formalism/LocalToGlobal/IndefiniteReal/LiftCoordinatesInGoodPartOfDefinite.lean}{Formalism.IndefR.Cx.BlockGood.csum_toF_at}
\leanref{Formalism/LocalToGlobal/IndefiniteReal/EquivFunOM.lean}{Formalism.IndefR.Cx.GroupActionMatrices.eig_ratG}
\leanref{Formalism/LocalToGlobal/Divisibility/LiftLatLemmas.lean}{Formalism.L2.DIV.CyclotomicLifting.xq_rat}
\begin{proof}
(1) The image of $c$ is the set of forms in $\Gamma$ whose coordinates lie in $Z_0$, which is
$Z\cap\Gamma$.

(2) The class-sum pairing of $c(x)$ and $c(y)$ is $|P|\,x^*Hy$. Since the involution is complex
conjugation at $v_0$, its image under $\sigma_0$ is $|P|\,\sigma_0(x)^*\sigma_0(H)\sigma_0(y)$; and since
$H_\Gamma$ is positive definite at $v_0$ by the conditions on $B$, the quotient metric at $v_0$ is
invertible with Hermitian form $\sigma_0(H)$.

(3) Let $n$ be the dimension of $B$. Under the hypotheses of (2), the matrix of $e_Xe_{P,\psi^{-1}}$ in
the rational basis has entries in $K_0$ and is self-adjoint for the class-sum pairing, whose matrix has
entries in $K_0$ and is positive definite at every infinite place; so $\Gamma$, its image, has a basis
with coordinates in $K_0$, the basis of $\Gamma$ in the construction is one, and $H_\Gamma$ has entries
in $K_0$. The coordinates of $Z\cap\Gamma$ in the rational basis of $B$ form a subspace of $L^n$ that is
stable under every matrix commuting with the matrices of the $T_i$ and of $e_Xe_{P,\psi^{-1}}$, since it
is the intersection of the image of the latter with a sum of joint generalized eigenspaces of the
former. These matrices have entries in $K_0$, commute, and contain the adjoint of each of them. A subspace with this property has a basis in $K_0^n$, by descent along the involution
of $L$; so $Z_0$ has a basis with entries in $K_0$, and $E$ has entries in $K_0$. Then $H=E^*H_\Gamma E$
has entries in $K_0$, and the quotient metric $X_Z$ is the restriction of $H_\Gamma$ to the columns of
$E$, whose Hermitian form at $v$ is $\sigma_v(E)^*\sigma_v(H_\Gamma)\sigma_v(E)=\sigma_v(H)$, because the
involution is complex conjugation at every embedding on $K_0$.
\end{proof}

\subsection{Comparison of conjugate data}\label{subsection-is-conjugate-data}

\begin{lemma}\label{is-eigenvalue-core-group}
Let $R$ be a realization of a definite-side datum $D$ with a face $F$ that satisfies the core
conditions at $F$, except possibly the one on its chain, with the notation of
Definition~\ref{is-eigenvalue-lift-data}. Let $i$ be a lift and $\sigma$ an embedding of $L$, write
$(\nu_W,\mathcal F_W)$ for the Haar measure and fundamental domain of the realization on
$U(W_i)(\mathbf A)$ (Definition~\ref{is-realisation}(5)), and put $N_\sigma=\prod_{\rho\in\mathcal
T}\dim\kappa_\rho(\sigma)$. Call $k\in U(W_i)(\mathbf A_{F_h,f})$ a \emph{level lift} of $g\in G$ if
$k$ lies in the parahoric level $K_W$ of $W_i$ at $P_g$ and, for every $P\in P_g$, the upper left
$r\times r$ block of the reduction of $k$ at $P$ is $g_P$.
\begin{enumerate}
\item Every $g\in G$ has a level lift. If $k$ is a level lift of $g$, then $(gU)(v)(h)=U(v)(hk)$ for
  every $U\in H_W(i,\sigma)$, every tensor $v$ of vectors of the models of discrete series of $U(W_i)$
  at $\sigma$ and every $h\in U(W_i)(\mathbf A)$.
\item Every element of $G$ acts unitarily on $H_W(i,\sigma)$.
\item Let $v=\otimes_\rho v_\rho$ be a tensor of unit vectors of these models. Then
  \[
   \langle U,U'\rangle=N_\sigma\int_{\mathcal F_W}\overline{U(v)(h)}\,U'(v)(h)\,d\nu_W(h)
  \]
  for all $U,U'\in H_W(i,\sigma)$. In particular, if $U(v)=0$, then $U=0$.
\end{enumerate}
\end{lemma}
\leanref{Development/AbelianVarietyTheory/Proofs/DivisorProofs/ConjugateLiftUniqueAndBaseDatum.lean}{GlobalTheory.GalLift.act_eq_translation_five}
\leanref{Development/AbelianVarietyTheory/Proofs/TwistedOperators/ParahoricLevelLiftsExist.lean}{GlobalTheory.ParahoricLift.liftsExist}
\leanref{Development/AbelianVarietyTheory/Proofs/DivisorProofs/LiftEquivarianceWeilModel.lean}{GlobalTheory.GalLift.Transl.transl}
\leanref{Development/AbelianVarietyTheory/Proofs/TwistedOperators/NormLiftFormsFiniteGroupAction.lean}{GlobalTheory.gW_norm_eq}
\leanref{Development/AbelianVarietyTheory/Proofs/DivisorProofs/EigenvalueBarCompatibility.lean}{GlobalTheory.ConjInner.inner_eq_integral}
\leanref{Development/AbelianVarietyTheory/Proofs/TwistedOperators/IpMSelfEqZero.lean}{GlobalTheory.EvalAtUnitVector.evInj}
\begin{proof}
By the core conditions on the models, $U\mapsto(U(v))_v$ identifies $H_W(i,\sigma)$ with the
multiplicity space of the model of $U(W_i)$ at $\sigma$ at the level $K_W$, and $\langle
U,U'\rangle=N_\sigma c(U,U')$. Here $c(U,U')$ is the constant with
$\int_{\mathcal F_W}\overline{U(v)}\,U'(v')\,d\nu_W=c(U,U')\prod_\rho\langle v_\rho,v'_\rho\rangle$
for all tensors $v,v'$ if there is one, and $0$ otherwise.

(1) By the core conditions on the level characters, the frame of $W_i$ is a base frame, no two places
of $P_g$ are conjugate, and the trivial level datum of $W_i$ has no ramified places. Let $g\in G$ and
$P\in P_g$. As $P$ is not conjugate to itself, $E_h\otimes_{F_h}\mathbf A_{F_h}$ has an idempotent
supported at $P$. Lift the block diagonal matrix with blocks $g_P$ and $1$ in $GL_n(k_P)$ to a matrix
over $\mathcal O_{E_h}$ whose determinant is a unit at $P$. Through the idempotent this matrix gives an
element $k_P$ of $U(W_i)(\mathbf A_{F_h,f})$ that stabilizes the lattice of the frame, reduces at $P$
to that block diagonal matrix and reduces to $1$ at the other places of $P_g$. The product of the
$k_P$ over $P\in P_g$ lies in $K_W$ and is a level lift of $g$. By the core conditions on the models,
$g$ acts on the functions of $H_W(i,\sigma)$ by the level action: for a level lift $k$,
$(gU)(v)(h)=\chi_K(k)^{-1}U(v)(hk)$, where $\chi_K$ is the character of the trivial level datum on
$K_W$. This does not depend on $k$, because the functions $U(v)$ transform under $K_W$ through
$\chi_K$; and $\chi_K$ is trivial.

(2) Let $g\in G$, and let $k$ be a level lift of $g$, which exists by (1). The group $K_W$ is compact
and contains $k$, so right translation by $k$ preserves $\nu_W$: the image of $\nu_W$ is a Haar
measure, and it gives the same finite positive volume as $\nu_W$ to $CK_W$ for a compact neighbourhood
$C$ of $1$, because $CK_Wk=CK_W$. By (1), the functions of $gU$ are the right translates by $k$ of
those of $U$. These functions are invariant on the left under $U(W_i)(F_h)$, and the integral over a
fundamental domain of such a function does not change under a right translation that preserves the
measure. So $c(gU,gU)=c(U,U)$, and $\|gU\|=\|U\|$. A linear isometry preserves inner products, so $g$
acts unitarily.

(3) If $U\ne0$, the constant $c(U,U)$ exists, since otherwise $\langle U,U\rangle=0$; for $U=0$ it is
$0$. Applying this to $U$, $U'$, $U+U'$ and $U+iU'$, polarization gives the constant $c(U,U')$.
Taking $v'=v$, for which $\prod_\rho\langle v_\rho,v_\rho\rangle=1$, gives the formula. If $U(v)=0$,
the formula with $U'=U$ gives $\langle U,U\rangle=0$.
\end{proof}

\begin{lemma}\label{is-eigenvalue-conjugate-forms}
Let $E_1,\dots,E_n$ and $W$ be complex inner product spaces. For $v,v'\in\prod_\tau E_\tau$ and
$w,w'\in W$ put $P(v,w;v',w')=\prod_\tau\langle v_\tau,v'_\tau\rangle\,\langle w,w'\rangle$, and for
$a\in E_\tau$ let $v[\tau\leftarrow a]$ be $v$ with its component at $\tau$ replaced by $a$.
\begin{enumerate}
\item Let $E$ be a finite-dimensional complex Hilbert space and $(g_k)_{k\in I}$ a family of linear
  maps $E\to E$ that preserve the inner product, such that $0$ and $E$ are the only subspaces of $E$
  stable under every $g_k$. Let $b$ be a form on $E$, conjugate-linear in the first variable and
  linear in the second, with $b(g_kx,g_ky)=b(x,y)$ for all $k$, $x$ and $y$. Then there is
  $c\in\mathbf C$ with $b(x,y)=c\langle x,y\rangle$ for all $x,y\in E$.
\item Suppose that $W$ has a unit vector, and let $B(v,w;v',w')\in\mathbf C$ be given for
  $v,v'\in\prod_\tau E_\tau$ and $w,w'\in W$. Suppose that for all $\tau$, $v,v'$ and $w,w'$ there is
  $c\in\mathbf C$ with $B(v[\tau\leftarrow a],w;v'[\tau\leftarrow a'],w')=c\langle a,a'\rangle$ for
  all $a,a'\in E_\tau$, and that for all $v,v'$ there is $c\in\mathbf C$ with
  $B(v,w;v',w')=c\langle w,w'\rangle$ for all $w,w'\in W$. Then there is $c\in\mathbf C$ with $B=cP$.
\item Let $Y$ be a complex vector space with a form $b$, conjugate-linear in the first variable and
  linear in the second, such that $b(y,x)=\overline{b(x,y)}$ and $\operatorname{Re}b(x,x)\ge0$ for
  all $x,y$, and $b(x,x)=0$ only for $x=0$. Let $F,F'\colon\prod_\tau E_\tau\times W\to Y$ and
  $c_1,c_2,c_3\in\mathbf C$ satisfy
  \[
   b(F(v,w),F(v',w'))=c_1P,\qquad b(F'(v,w),F'(v',w'))=c_2P,\qquad b(F(v,w),F'(v',w'))=c_3P
  \]
  with $P=P(v,w;v',w')$, for all $v,w,v',w'$. Suppose that $F(v_0,w_0)=l\,F'(v'_0,w'_0)$ for some
  $v_0,v'_0$ whose components are unit vectors, unit vectors $w_0,w'_0$ and $l\in\mathbf C$ with
  $|l|=1$. Then there is $\kappa\in\mathbf C$ with $|\kappa|=1$ and $F'(v,w)=\kappa F(v,w)$ for all
  $v,w$.
\end{enumerate}
\end{lemma}
\leanref{Development/AbelianVarietyTheory/Proofs/DivisorProofs/FormEqMulInner.lean}{GlobalTheory.ConjLiftUnique.form_eq_mul_inner}
\leanref{Development/AbelianVarietyTheory/Proofs/DivisorProofs/FormEqMulInner.lean}{GlobalTheory.ConjLiftUnique.exists_const}
\leanref{Development/AbelianVarietyTheory/Proofs/DivisorProofs/FormEqMulInner.lean}{GlobalTheory.ConjLiftUnique.exists_phase}
\begin{proof}
(1) There is a linear map $A$ of $E$ with $b(x,y)=\langle Ax,y\rangle$ for all $x,y$. Each $g_k$
preserves the inner product, so it is injective, hence bijective. For all $x,z$,
$\langle Ag_kx,g_kz\rangle=b(g_kx,g_kz)=b(x,z)=\langle Ax,z\rangle=\langle g_kAx,g_kz\rangle$; since
$g_k$ is surjective, $Ag_k=g_kA$. If $E=0$, take $c=0$. Otherwise $A$ has an eigenvalue $\lambda$.
Its eigenspace is nonzero and stable under every $g_k$, so it is $E$, and
$b(x,y)=\langle\lambda x,y\rangle=\bar\lambda\langle x,y\rangle$.

(2) If some $E_\tau$ is $0$, the first hypothesis with $a=v_\tau$ and $a'=v'_\tau$ gives
$B(v,w;v',w')=c\langle v_\tau,v'_\tau\rangle=0$, and $c=0$ works. Otherwise choose unit vectors
$e_\tau\in E_\tau$, put $e=(e_\tau)$, and let $w_0$ be a unit vector of $W$. By the second
hypothesis, $B(v,w;v',w')=B(v,w_0;v',w_0)\langle w,w'\rangle$. For a set $s$ of indices and $v,v'$
with $v_\tau=v'_\tau=e_\tau$ for $\tau\notin s$,
$B(v,w_0;v',w_0)=B(e,w_0;e,w_0)\prod_{\tau\in s}\langle v_\tau,v'_\tau\rangle$, by induction on $s$:
for $\tau\in s$ the first hypothesis gives $c$ with $B(v,w_0;v',w_0)=c\langle v_\tau,v'_\tau\rangle$
and $B(v[\tau\leftarrow e_\tau],w_0;v'[\tau\leftarrow e_\tau],w_0)=c$, and the induction hypothesis
for $s\smallsetminus\{\tau\}$ applies to the latter. For the set of all indices this gives $B=cP$ with
$c=B(e,w_0;e,w_0)$.

(3) The value of $P$ at $(v_0,w_0;v_0,w_0)$ and at $(v'_0,w'_0;v'_0,w'_0)$ is $1$. So
$c_1=b(F(v_0,w_0),F(v_0,w_0))=|l|^2c_2=c_2$. If $c_1=0$, then $b(F(v,w),F(v,w))$ and
$b(F'(v,w),F'(v,w))$ vanish for all $v,w$, so $F=F'=0$ and $\kappa=1$ works. Let $c_1\ne0$, and put
$p=P(v_0,w_0;v'_0,w'_0)$ and $q=P(v'_0,w'_0;v_0,w_0)=\bar p$. The third formula at
$(v_0,w_0;v'_0,w'_0)$ gives $c_3p=b(lF'(v'_0,w'_0),F'(v'_0,w'_0))=\bar lc_1$, and at
$(v_0,w_0;v_0,w_0)$ it gives $c_3=b(lF'(v'_0,w'_0),F'(v_0,w_0))=\bar lc_1q$. Hence
$\bar lc_1(qp-1)=0$, so $qp=1$, and $|q|=1$ because $|p|=|q|$. Put $\kappa=\bar lq$, so that
$|\kappa|=1$ and $c_3=\kappa c_1$. Fix $v,w$ and put $x=F(v,w)$ and $y=F'(v,w)$. The three formulas
give $b(y,y)=b(x,x)$ and $b(x,y)=\kappa b(x,x)$, hence $b(y,x)=\bar\kappa b(x,x)$, as $b(x,x)$ is
real. Then $b(y-\kappa x,y-\kappa x)=b(y,y)-\kappa b(y,x)-\bar\kappa b(x,y)+|\kappa|^2b(x,x)=0$, so
$y=\kappa x$.
\end{proof}

\begin{lemma}\label{is-eigenvalue-conjugate-invariant-form}
Let $\omega$ be a Weil datum for Hermitian spaces $(S,T)$ of a definite-side datum
(Definition~\ref{th-weil-datum}(3)), and write $p$ for elements of $U(S)(\mathbf A)\times
U(T)(\mathbf A)$.
\begin{enumerate}
\item There is a form $b$ on the space of theta functions for $(S,T)$, conjugate-linear in the first
  variable and linear in the second, such that $b(y,x)=\overline{b(x,y)}$,
  $\operatorname{Re}b(x,x)\ge0$, $b(x,x)=0$ only for $x=0$, and $b(\omega(p)x,\omega(p)y)=b(x,y)$ for
  all theta functions $x,y$ and all $p$.
\item If $f$ is a theta function and $\omega(p)f=cf$ for some $p$ and $c\in\mathbf C$, then
  $\omega(p)f=c'f$ for some $c'\in\mathbf C$ with $|c'|=1$.
\end{enumerate}
\end{lemma}
\leanref{Development/AbelianVarietyTheory/Proofs/DivisorProofs/ThetaLiftUniqueUpToUnit.lean}{GlobalTheory.ConjLiftUnique.exists_invForm}
\leanref{Development/AbelianVarietyTheory/Proofs/DivisorProofs/ThetaLiftUniqueUpToUnit.lean}{GlobalTheory.ConjLiftUnique.exists_unit_eigen}
\begin{proof}
(1) The group $\Gamma=M_{\operatorname{rk}T\times\operatorname{rk}S}(E)$ of rational points of
$\mathbb W$ is discrete, and some compact set $C\subset\mathbb W$ meets every coset of $\Gamma$,
because $E$ is discrete and cocompact in $\mathbf A_E$. So there are a Haar measure $\mu$ on
$\mathbb W$ and a fundamental domain $\mathcal F$ for $\Gamma$ with $\mu(\mathcal F)=1$. A theta
function $f$ is continuous, and $|f(\gamma+w)|=|f(w)|$ for $\gamma\in\Gamma$ because $\psi$ is
unitary; so $|f|$ is bounded by its maximum on $C$. Put $b(x,y)=\int_{\mathcal F}y\,\bar x\,d\mu$. It
is conjugate-linear in $x$ and linear in $y$, $b(y,x)=\overline{b(x,y)}$, and
$b(x,x)=\int_{\mathcal F}|x|^2d\mu\ge0$. If $b(x,x)=0$, then the integral of $|x|^2$ over
$\mathbb W$, the sum of its integrals over the translates $\gamma+\mathcal F$, is $0$, since $|x|$ is
invariant under $\Gamma$; so the continuous function $x$ vanishes almost everywhere, hence everywhere,
as a Haar measure charges every nonempty open set. By Definition~\ref{th-weil-datum}(3), $\omega(p)$
preserves theta functions and their $L^2$-norm on $\mathcal F$; by polarization it preserves $b$.

(2) If $f=0$, take $c'=1$. Otherwise $b(f,f)\ne0$, and $b(f,f)=b(\omega(p)f,\omega(p)f)=|c|^2b(f,f)$
by (1), so $|c|=1$.
\end{proof}

\begin{lemma}\label{is-eigenvalue-conjugate-coefficients}
Let $V$ be a Hermitian space of a definite-side datum $D$ with host field $E_h/F_h$, and let
$(W_c,\rho)$ and $(W'_c,\rho')$ be coefficient data for $V$ (Definition~\ref{is-eigenvalue-lift-data}(7)):
nonzero finite-dimensional Hilbert spaces with irreducible continuous unitary representations of
$U(V)(\mathbf A)$ that are trivial on $U(V)(\mathbf A_{F_h,f})$. Let $X$ be a finite-dimensional vector
space over a field $K$ with a basis $(b_a)$, $\sigma\colon K\to\mathbf C$ an embedding, and
$r_\gamma\colon X\to X$ a map for each $\gamma\in U(V)(F_h)$. Let $c_1\colon X\to W_c$ and
$c'_1\colon X\to W'_c$ be $\sigma$-semilinear maps that carry $(b_a)$ to bases, with
$\langle c'_1(b_a),c'_1(b_{a'})\rangle=\langle c_1(b_a),c_1(b_{a'})\rangle$ for all $a,a'$, and with
$c_1(r_\gamma x)=\rho(\gamma)c_1(x)$ and $c'_1(r_\gamma x)=\rho'(\gamma)c'_1(x)$ for all $\gamma$ and
$x$. Then there is a linear isomorphism $\Phi\colon W_c\to W'_c$ that preserves norms, with
$\Phi\circ c_1=c'_1$ and $\Phi\rho(h)=\rho'(h)\Phi$ for every $h\in U(V)(\mathbf A)$.
\end{lemma}
\leanref{Development/AbelianVarietyTheory/Proofs/DivisorProofs/ThetaLiftUniqueUpToUnit.lean}{GlobalTheory.ConjLiftUnique.exists_isometry}
\begin{proof}
Let $\Phi$ be the linear isomorphism with $\Phi(c_1(b_a))=c'_1(b_a)$ for all $a$. The maps
$\Phi\circ c_1$ and $c'_1$ are $\sigma$-semilinear and agree on the basis, so they are equal. Writing
$x\in W_c$ in the basis $(c_1(b_a))$, the hypothesis on inner products gives
$\langle\Phi x,\Phi x\rangle=\langle x,x\rangle$. For $\gamma\in U(V)(F_h)$ the maps $\Phi\rho(\gamma)$
and $\rho'(\gamma)\Phi$ agree on the basis, since
$\Phi\rho(\gamma)c_1(b_a)=\Phi c_1(r_\gamma b_a)=c'_1(r_\gamma b_a)=\rho'(\gamma)c'_1(b_a)$; so they
are equal. Let $h\in U(V)(\mathbf A)$ and $x\in W_c$, and write $h=uk$ with $u\in
U(V)(F_h\otimes\mathbf R)$ and $k\in U(V)(\mathbf A_{F_h,f})$. Both representations are trivial on
$k$, so it suffices to show that $\Phi\rho(u)x=\rho'(u)\Phi x$. If not, the set of $h'\in
U(V)(\mathbf A)$ with $\Phi\rho(h')x\ne\rho'(h')\Phi x$ is open, by continuity, and contains $u$. By
weak approximation for unitary groups \cite[Chapter~7]{PR94} it contains an element $\gamma k'$ with
$\gamma\in U(V)(F_h)$ and $k'\in U(V)(\mathbf A_{F_h,f})$. But
$\Phi\rho(\gamma k')x=\Phi\rho(\gamma)x=\rho'(\gamma)\Phi x=\rho'(\gamma k')\Phi x$.
\end{proof}

\begin{definition}\label{is-eigenvalue-conjugate-k-types}
Let $V$ be a Hermitian space of rank $n$ of a definite-side datum, and let $M$ be a model of discrete
series for $U(V)$ (Definition~\ref{asm-dependence-rows-models}(1)) with Hilbert spaces $H_\tau$,
representations $\sigma_\tau$, parameters $(\lambda_\tau,\epsilon_\tau)$ and signature matrices
$J_\tau$. Let $K_\tau$ be the group of unitary matrices that commute with $J_\tau$.
\begin{enumerate}
\item A vector $v\in H_\tau$ has \emph{weight} $\mu\in\mathbf R^n$ if
  $\sigma_\tau(\operatorname{diag}(e^{it_1},\dots,e^{it_n}))v=e^{i\sum_a\mu_at_a}v$ for every
  $t\in\mathbf R^n$.
\item For $\Lambda\in\mathbf R^n$, a \emph{$K$-type of weight $\Lambda$} of $\sigma_\tau$ is a
  finite-dimensional subspace $\kappa\subset H_\tau$, stable under $\sigma_\tau(K_\tau)$, whose only
  subspaces stable under $\sigma_\tau(K_\tau)$ are $0$ and $\kappa$, that contains a nonzero vector of
  weight $\Lambda$, and in which, for every nonzero vector of a weight $\mu$, $\Lambda-\mu$ is a
  combination with coefficients in $\mathbf Z_{\ge0}$ of the vectors $e_a-e_b$ with $a<b$ and
  $\epsilon_{\tau,a}=\epsilon_{\tau,b}$.
\item Let $\Lambda_\tau=\lambda_\tau+\rho_n-\rho_c$, where $2\rho_{n,a}$ is the number of $b>a$ with
  $\epsilon_{\tau,b}\ne\epsilon_{\tau,a}$ minus the number of $b<a$ with
  $\epsilon_{\tau,b}\ne\epsilon_{\tau,a}$, and $2\rho_{c,a}$ is the same difference for the $b$ with
  $\epsilon_{\tau,b}=\epsilon_{\tau,a}$. The \emph{lowest $K$-type} $\kappa_\tau$ of $M$ at $\tau$ is
  the sum of the $K$-types of $\sigma_\tau$ of weight $\Lambda_\tau$. The model has \emph{smooth
  irreducible lowest $K$-types} if at every $\tau$ the space $\kappa_\tau$ consists of smooth vectors
  and is itself a $K$-type of $\sigma_\tau$ of weight $\Lambda_\tau$.
\end{enumerate}
\end{definition}
\leanref{Formalism/LocalToGlobal/IndefiniteReal/FockTheta/DiscreteSeriesUnitary.lean}{Formalism.IndefR.Fock.SquareIntegrableRepresentations.IsWtOn}
\leanref{Formalism/LocalToGlobal/IndefiniteReal/FockTheta/DiscreteSeriesUnitary.lean}{Formalism.IndefR.Fock.SquareIntegrableRepresentations.IsKTypeOf}
\leanref{Formalism/LocalToGlobal/IndefiniteReal/FockTheta/DiscreteSeriesUnitary.lean}{Formalism.IndefR.Fock.SquareIntegrableRepresentations.lowKSub}
\leanref{Formalism/FormalDegreeBessel/StatementDS.lean}{Formalism.QuantumDeterminants.lowK}
\leanref{Formalism/LocalToGlobal/IndefiniteReal/ExistsChosenExtremalVector.lean}{Formalism.IndefR.FWT.LowSm}
\leanref{Development/AbelianVarietyTheory/BlockTheory/TwistedDiscreteSeries/HolomorphicLowestType.lean}{GlobalTheory.IndefiniteRowSideFullHecke.HolomorphicLowestType.LowIrr}

\begin{lemma}\label{is-eigenvalue-conjugate-theta-closed}
Let $D$ be a definite-side datum with host field $E_h/F_h$, let $\psi$ be a standard character of
$\mathbf A_{F_h}$, let $\theta$ be an archimedean family, and let $S$ and $T$ be Hermitian spaces of
$D$, with $\mathbb W$, $\langle\ ,\ \rangle$ and $\rho$ as in Definition~\ref{th-weil-datum}. A
\emph{real direction} is a continuous additive map $L\colon\mathbf R\to\mathbb W$ with finite
component $0$. For a function $f$ on $\mathbb W$ let $\partial_Lf(w)$ be the derivative at $t=0$ of
$t\mapsto(\rho(L(t))f)(w)$. For $\tau\in\mathcal T$ and a complex matrix $z$ of size $\operatorname{rk}
T\times\operatorname{rk}S$ let $L_{\tau,z}$ be the real direction whose component at $\tau$ under
$\theta$ is $t\mapsto tz$ and whose components at the other embeddings are $0$, and put
$\partial_\tau(z)=\partial_{L_{\tau,z}}$. Let $f$ be a theta function for $(S,T)$. Then:
\begin{enumerate}
\item $\rho(u)f$ is a theta function for every $u\in\mathbb W$;
\item for every real direction $L$ and every $w$, the function $t\mapsto(\rho(L(t))f)(w)$ is
  differentiable with derivative $(\rho(L(t))\partial_Lf)(w)$ at $t$, and $\partial_Lf$ is a theta
  function;
\item $f$ is bounded, and for every real direction $L$ there is $C$ with
  $|t^{-1}(\rho(L(t))f-f)-\partial_Lf|\le Ct$ on $\mathbb W$ for every $t>0$.
\end{enumerate}
\end{lemma}
\leanref{Development/AbelianVarietyTheory/Proofs/TwistedOperators/InfPartOneSubEInfSmul.lean}{GlobalTheory.ThetaClosed.thetaClosed}
\leanref{Development/AbelianVarietyTheory/Proofs/TwistedOperators/InfPartOneSubEInfSmul.lean}{GlobalTheory.ThetaClosed.isThetaFunO_heisO_all}
\leanref{Development/AbelianVarietyTheory/Proofs/AdelicLatticeProofs/HasDerivAtZeroAndIsTheta.lean}{GlobalTheory.WeilVacMove.isTheta_heisO_real}
\leanref{Development/AbelianVarietyTheory/Proofs/AdelicLatticeProofs/SmulEqZeroOfInf.lean}{GlobalTheory.FinTranslate.isThetaFunO_heisO}
\leanref{Development/AbelianVarietyTheory/Proofs/AdelicLatticeProofs/ComplexEmbeddingHermitianQuadrupleSum.lean}{GlobalTheory.OrientedHeis.theta_dLine}
\leanref{Development/AbelianVarietyTheory/Proofs/AdelicLatticeProofs/ComplexEmbeddingHermitianQuadrupleSum.lean}{GlobalTheory.OrientedHeis.hasDerivAt_heisO}
\leanref{Development/AbelianVarietyTheory/Proofs/AdelicLatticeProofs/EqZeroOfLintegralEqZero.lean}{GlobalTheory.OrientedHeis.theta_bounded}
\leanref{Development/AbelianVarietyTheory/Proofs/AdelicLatticeProofs/EqZeroOfLintegralEqZero.lean}{GlobalTheory.OrientedHeis.quotient_bound}
\begin{proof}
Recall that $\rho(w)\rho(u)=\psi(\frac12\langle w,u\rangle)\rho(w+u)$, and that $\langle w,u\rangle=0$
when $w$ has finite component $0$ and $u$ archimedean component $0$, the pairing being computed place by
place.

(1) Write $u=u_f+u_\infty$ with $u_f$ of archimedean component $0$ and $u_\infty$ of finite
component $0$; then $\rho(u)=\rho(u_f)\rho(u_\infty)$, and it suffices to treat $u_f$ and $u_\infty$.
For any $u$ and rational $\gamma$,
$(\rho(u)f)(\gamma+w)=\psi(\frac12\langle\gamma+w,u\rangle)\psi(-\frac12\langle\gamma,w+u\rangle)f(w+u)
=\psi(-\frac12\langle\gamma,w\rangle)(\rho(u)f)(w)$, as the terms $\pm\frac12\langle\gamma,u\rangle$
cancel.
Let $U$ be an open subgroup with $\rho(x)f=f$ for the $x\in U$ of archimedean component $0$. For such
$x$, $\rho(x)$ commutes with $\rho(u_\infty)$, so $\rho(u_\infty)f$ is invariant under them. The set
of $x$ with $\psi(\frac12\langle x,u_f\rangle)=1$ is an open subgroup, as $\psi$ is continuous and the
pairing with $u_f$ has archimedean component $0$; on its intersection with $U$, $\rho(x)$ commutes with
$\rho(u_f)$, so $\rho(u_f)f$ is invariant there. For a continuous additive map $L$ from some
$\mathbf R^a$ to $\mathbb W$ with finite component $0$,
$(\rho(u_f)f)(w+L(t))=\psi(\frac12\langle w,u_f\rangle)f(w+u_f+L(t))$ is smooth in $t$, and
$(\rho(u_\infty)f)(w+L(t))$ is $f(w+u_\infty+L(t))$ times the exponential of an affine function of
$t$, because $\psi$ is standard; both are smooth, and continuous on $\mathbb W$. A function that is
smooth along the real directions and invariant under $\rho(x)$ for the $x$ of archimedean component $0$
in an open subgroup is continuous, since $g(w+a+x)=\psi(-\frac12\langle w+a,x\rangle)g(w+a)$ for such
$x$ and real $a$.

(2) Since $\langle L(s),L(t)\rangle=st\langle L(1),L(1)\rangle=0$,
$\rho(L(s+t))=\rho(L(s))\rho(L(t))$, so $(\rho(L(s+t))f)(w)=\psi(\frac12\langle w,L(s)\rangle)
(\rho(L(t))f)(w+L(s))$. At $s=0$ the function $t\mapsto(\rho(L(t))f)(w)$ is $f(w+L(t))$ times the
exponential of a linear function of $t$, hence differentiable; the identity gives the derivative at
every $s$. By (1), the $\rho(L(t))f$ are theta functions, and differentiating their quasi-periodicity in
$t$ gives that of $\partial_Lf$. For $x$ of archimedean component $0$ in the open subgroup of $f$,
$\rho(x)\rho(L(t))f=\rho(L(t))\rho(x)f=\rho(L(t))f$, and $\rho(x)$ commutes with $\partial/\partial t$;
so $\partial_Lf$ is invariant under these $\rho(x)$. It is smooth along the real directions, as $f$ is,
and so continuous, as in (1).

(3) As $|\psi|=1$, $|f|$ is invariant under the rational points, and a compact subset of $\mathbb W$
meets every coset of them, because $E_h$ is discrete and cocompact in $\mathbf A_{E_h}$; so the
continuous function $f$ is bounded. By (2), applied twice, $\varphi(t)=(\rho(L(t))f)(w)$ has derivatives
$\varphi'(t)=(\rho(L(t))\partial_Lf)(w)$ and $\varphi''(t)=(\rho(L(t))\partial_L\partial_Lf)(w)$, and
$|\varphi''|\le C$ for a bound $C$ of the theta function $\partial_L\partial_Lf$. Then
$|\varphi'(s)-\varphi'(0)|\le Cs$, so
$|\varphi(t)-\varphi(0)-t\varphi'(0)|\le Ct^2$ for $t>0$, which is the claim.
\end{proof}

\begin{lemma}\label{is-eigenvalue-conjugate-theta-schur}
Let $D$ be a definite-side datum with host field $E_h/F_h$, let $\psi$ be a standard character of
$\mathbf A_{F_h}$, and let $S$ and $T$ be Hermitian spaces of $D$, with $\mathbb W$, $\langle\ ,\
\rangle$ and $\rho$ as in Definition~\ref{th-weil-datum}. Let $M$ be a linear operator on the functions
on $\mathbb W$ that carries theta functions to theta functions, satisfies $M\rho(w)f=\rho(w)Mf$ for
every $w\in\mathbb W$ and every theta function $f$, and preserves the $L^2$-norm of theta functions on
every fundamental domain of every Haar measure of $\mathbb W$ for its rational points
$\Gamma=M_{\operatorname{rk}T\times\operatorname{rk}S}(E_h)$. Then there is $c\in\mathbf C$ with $|c|=1$
and $Mf=cf$ for every theta function $f$.
\end{lemma}
\leanref{Development/AbelianVarietyTheory/Proofs/DivisorProofs/SchurLemmaThetaFunctionsOrientedPairStatement.lean}{GlobalTheory.T.schurLemma_thetaFunctions_orientedPair}
\leanref{Development/AbelianVarietyTheory/Proofs/DivisorProofs/SchurLemmaThetaFunctionsOrientedPair.lean}{GlobalTheory.ThetaSchur.schur}
\leanref{Development/AbelianVarietyTheory/Proofs/DivisorProofs/SchurLemmaThetaFunctionsOrientedPair.lean}{GlobalTheory.ThetaSchur.prod_eq}
\leanref{Development/AbelianVarietyTheory/Proofs/IwasawaDecomposition/RationalSymplecticPointsSelfAnnihilator.lean}{GlobalTheory.RatAnnih.ratL_annihilator}
\leanref{Development/AbelianVarietyTheory/Proofs/IwasawaDecomposition/EqZeroOfCoefficientsSeparatingCharactersEqZero.lean}{GlobalTheory.SeparatingCharacters.eq_zero_of_coefficients_eq_zero}
\leanref{Development/AbelianVarietyTheory/Proofs/IwasawaDecomposition/ThetaFunctionScalarProductPreservedIsometry.lean}{GlobalTheory.Mem.LiftPolar.scalar_eq}
\begin{proof}
\emph{Step 1: two identities.} Let $f,f'$ be theta functions. For $\gamma\in\Gamma$,
$f(\gamma+x)\overline{f'(\gamma+x)}=f(x)\overline{f'(x)}$, because the factors
$\psi(\mp\frac12\langle\gamma,x\rangle)$ of the quasi-periodicity have absolute value $1$. And
$(\rho(\gamma)f)(x)=\psi(\frac12\langle x,\gamma\rangle)f(\gamma+x)=\psi(\langle x,\gamma\rangle)f(x)$,
as $\langle\ ,\ \rangle$ is alternating.

\emph{Step 2: the pairing.} A compact subset of $\mathbb W$ meets every coset of $\Gamma$, because $E_h$
is discrete and cocompact in $\mathbf A_{E_h}$; so there are a Haar measure $\mu$ and a fundamental
domain $\mathcal F$ for $\Gamma$ of finite volume, and theta functions are bounded by Step 1. For theta
functions $f,f'$ the integral $(f,f')=\int_{\mathcal F}f\,\overline{f'}\,d\mu$ is a Hermitian form whose
quadratic form is the square of the $L^2$-norm. As $M$ preserves that norm on theta functions, which form
a vector space, polarization gives $(Mf,Mf')=(f,f')$.

\emph{Step 3: products.} Let $f,f'$ be theta functions and
$\Phi(x)=Mf(x)\overline{Mf'(x)}-f(x)\overline{f'(x)}$. It is continuous and, by Step 1 for $f,f'$ and
for $Mf,Mf'$, invariant under $\Gamma$. For $\gamma\in\Gamma$ put $e_\gamma(x)=\psi(\langle
x,\gamma\rangle)$, a continuous character of $\mathbb W$ trivial on $\Gamma$. By Step 1,
$e_\gamma f=\rho(\gamma)f$ and $e_\gamma Mf=\rho(\gamma)Mf=M\rho(\gamma)f$, so by Step 2
\[
 \int_{\mathcal F}e_\gamma\Phi\,d\mu=(M\rho(\gamma)f,Mf')-(\rho(\gamma)f,f')=0,
\]
$\rho(\gamma)f$ being a theta function (Lemma~\ref{is-eigenvalue-conjugate-theta-closed}(1)). The
characters $e_\gamma$ separate the points of the compact group $\mathbb W/\Gamma$. Indeed, if
$\psi(\langle x,\gamma\rangle)=1$ for every $\gamma\in\Gamma$, then for every entry $x_{ij}$ of $x$ and
every $a\in E_h$ there is $\gamma\in\Gamma$ with $\langle x,\gamma\rangle=\operatorname{Tr}_{\mathbf
A_{E_h}/\mathbf A_{F_h}}(ax_{ij})$, so $\psi\circ\operatorname{Tr}$, a nontrivial character of
$\mathbf A_{E_h}$ trivial on $E_h$, takes the value $1$ at $ax_{ij}$ for every $a\in E_h$, and
$x_{ij}\in E_h$ by Theorem~\ref{asm-theta-agree-invariance-tate}(4); so $x\in\Gamma$. The linear span
of the $e_\gamma$ is an algebra of continuous functions on $\mathbb W/\Gamma$ that contains the
constants, is closed under complex conjugation and separates points, so it is dense in the continuous
functions by the Stone-Weierstrass theorem. Hence $\int_{\mathcal F}\bar\Phi\Phi\,d\mu=0$, so $\Phi$
vanishes almost everywhere on $\mathcal F$, hence on $\mathbb W$ as it is invariant under $\Gamma$; as
$\Phi$ is continuous and a Haar measure charges every nonempty open set, $\Phi=0$.

\emph{Step 4: two points.} Let $f,f'$ be theta functions and $x,y\in\mathbb W$. Apply Step 3 to $f$ and
the theta function $\rho(y-x)f'$ at the point $x$. Since $(\rho(y-x)g)(x)=\psi(\frac12\langle
x,y\rangle)g(y)$ for every $g$, and $M$ commutes with $\rho(y-x)$, this gives
$Mf(x)\overline{Mf'(y)}=f(x)\overline{f'(y)}$ after cancelling the factor of absolute value $1$.

\emph{Step 5: the scalar.} If every theta function vanishes, $c=1$ works. Otherwise let $f_0$ be a theta
function and $y$ a point with $f_0(y)\ne0$. Step 4 with $f=f'=f_0$ and $x=y$ gives
$|Mf_0(y)|=|f_0(y)|\ne0$. Put $c=\overline{f_0(y)}/\overline{Mf_0(y)}$, so that $|c|=1$. For every theta
function $f$ and every $x$, Step 4 with $f'=f_0$ gives $Mf(x)\overline{Mf_0(y)}=f(x)\overline{f_0(y)}$,
that is, $Mf(x)=cf(x)$.
\end{proof}

\begin{lemma}\label{is-eigenvalue-conjugate-vacuum-exists}
Let $D$ be a definite-side datum with host field $E_h/F_h$, let $\psi$ be a standard character of
$\mathbf A_{F_h}$, let $\theta$ be an archimedean family whose member at each $\tau\in\mathcal T$ extends
$\tau$, and let $S$ and $T$ be Hermitian spaces of $D$ with Cartan involutions $s_S=(s_{S,\tau})$ and
$s_T=(s_{T,\tau})$ at the $\tau\in\mathcal T$ and self-dual lattices $\Lambda_S$ and $\Lambda_T$. Let
$\Lambda$ be the lattice of the pair given by $\Lambda_S$ and $\Lambda_T$.
\begin{enumerate}
\item There is a nonzero function $f_0$ on $\mathbb W$ such that, for every $c\in\mathbf C$, $cf_0$ is a
  theta function for $(S,T)$, a lattice vacuum for $\Lambda$, and annihilated by the annihilation
  operators for $s_S$ and $s_T$ at every $\tau$.
\item If $S$ and $T$ have correct signatures, there is a Fock vacuum vector for $s_S$, $s_T$ and
  $\Lambda$.
\end{enumerate}
\end{lemma}
\leanref{Development/AbelianVarietyTheory/Proofs/AdelicLatticeProofs/MapFinNeg.lean}{GlobalTheory.VacuumO.exists_ne_zero}
\leanref{Development/AbelianVarietyTheory/Proofs/AdelicLatticeProofs/MapFinNeg.lean}{GlobalTheory.VacuumO.vacExists}
\leanref{Development/AbelianVarietyTheory/Proofs/AdelicLatticeProofs/ContinuousBicharacterWeilSetupStatement.lean}{GlobalTheory.T.continuous_bicharacter_weil_setup}
\leanref{Development/AbelianVarietyTheory/Proofs/AdelicLatticeProofs/ExistsNormaliseMeanSquareOfCocompactStatement.lean}{GlobalTheory.T.exists_normalise_meanSquare_of_cocompact'}
\begin{proof}
Part (1) is the theta series of the Gaussian \cite{Wei64}: at the archimedean places, the Gaussian of
the complex structures given by $s_S$ and $s_T$ is killed by the annihilation operators, and at the
finite places the characteristic function of the closure of $\Lambda$, twisted by the character of
$\Lambda$ that the lattice vacuum conditions prescribe, is fixed by $\rho(2u)$ for $u\in\Lambda$ and by
$\rho(u)$ for $u\in\Lambda$ in a split half; the sum over the rational points of the translates of
their product converges to a nonzero theta function with these properties, and so do its multiples.

For (2), let $f_0$ be as in (1). A compact set meets every coset of the rational points $\Gamma$, so
there are Haar measures and fundamental domains of covolume $1$, and any two such give the same integral
of a $\Gamma$-invariant function. The function $|f_0|$ is $\Gamma$-invariant, by the quasi-periodicity,
and continuous and nonzero, so its $L^2$-norm $a$ on a fundamental domain of covolume $1$ is positive and
does not depend on the domain. Then $a^{-1}f_0$ is a Fock vacuum vector.
\end{proof}

\begin{lemma}\label{is-eigenvalue-conjugate-split-form}
Let $D$ be a definite-side datum with host field $E_h/F_h$, let $\psi$ be a standard character of
$\mathbf A_{F_h}$, and let $S$ and $T$ be Hermitian spaces of $D$. Call $e\in\mathbf A_{E_h}$ a
\emph{split idempotent} if $e^2=e$, $e\bar e=0$ and $e+\bar e$ is the idempotent of $\mathbf A_{F_h}$ of
a finite place $v$ (which then splits in $E_h$); then $\bar e$ is a split idempotent too. The
\emph{half} of $\mathbb W$ attached to $e$ is the set of $u\in\mathbb W$ with $eu=u$. For every split
idempotent $e$ there is a linear form $\ell$ on the functions on $\mathbb W$ such that
$\ell(\rho(u)f)=\ell(f)$ for every theta function $f$ and every $u$ in the half attached to $\bar e$,
and $\ell(f)\ne0$ for some theta function $f$.
\end{lemma}
\leanref{Development/AbelianVarietyTheory/Proofs/DivisorProofs/ExistsSplitPlaceHalfHeisenbergInvariantFormStatement.lean}{GlobalTheory.T.exists_split_place_half_heisenberg_invariant_form}
\leanref{Development/AbelianVarietyTheory/Proofs/DivisorProofs/ExistsSplitPlaceHalfHeisenbergInvariantForm.lean}{GlobalTheory.HalfCoinv.exists_invariant_form}
\leanref{Development/AbelianVarietyTheory/Proofs/DivisorProofs/ExistsSplitPlaceHalfHeisenbergInvariantForm.lean}{GlobalTheory.HalfCoinv.exists_form}
\leanref{Development/AbelianVarietyTheory/Proofs/DivisorProofs/ExistsSplitPlaceHalfHeisenbergInvariantForm.lean}{GlobalTheory.HalfCoinv.exists_sum_eq_zero}
\leanref{Development/AbelianVarietyTheory/Proofs/DivisorProofs/ExistsSplitPlaceHalfHeisenbergInvariantForm.lean}{GlobalTheory.HalfCoinv.exists_separating}
\leanref{Development/AbelianVarietyTheory/Proofs/AdelicLatticeProofs/IsSelfDualCLatFrame.lean}{GlobalTheory.PlaceModel.mem_of_dual}
\begin{proof}
Choose an archimedean family $\theta$ whose member at each $\tau$ extends $\tau$, Cartan involutions of
$S$ and $T$ at every $\tau$, which exist because the Gram matrices are Hermitian and invertible, and let
$\Lambda$ be the lattice of the pair given by the self-dual lattices of the frames of $S$ and $T$. Let
$f_0\ne0$ be as in Lemma~\ref{is-eigenvalue-conjugate-vacuum-exists}(1); it is a theta function and a
lattice vacuum for $\Lambda$. Let $Y$ be the half attached to $\bar e$, and write
$B(u,y)=\psi(\langle u,y\rangle)$. Then:
\begin{enumerate}
\item[(a)] $\rho(u)\rho(y)=B(u,y)\rho(y)\rho(u)$ for all $u,y\in\mathbb W$, and $\rho(x)\rho(x')=\rho(x+x')$
  for $x,x'\in Y$, because $\langle x,x'\rangle=\langle\bar ex,\bar ex'\rangle=0$ as $e\bar e=0$.
\item[(b)] For $u\in Y$ and a theta function $g$ there is $k>0$ with $\rho(ku)g=g$. Indeed $u$ has
  archimedean component $0$, $g$ is invariant under $\rho(x)$ for the $x$ of archimedean component $0$
  in an open subgroup, and every element of archimedean component $0$ has a positive multiple in that
  subgroup.
\item[(c)] For $y\in Y$, either $\rho(y)f_0=f_0$, or there is $u$ with $\rho(u)f_0=f_0$ and $B(u,y)\ne1$.
  Indeed, let $\Lambda_v$ be the part of $\Lambda$ at the place $v$ of $e$. If $B(u,y)=1$ for every
  $u\in\Lambda_v$ in the half attached to $e$, then for every $u'\in\Lambda$, $eu'$ is such an element,
  and $\psi(\langle y,u'\rangle)=\psi(\langle y,eu'\rangle)=1$, as $y=\bar ey$. By
  Lemma~\ref{th-lattice-signs}(1), $y\in\Lambda$; so $\rho(y)f_0=f_0$, as $f_0$ is a lattice vacuum and
  $y$ lies in the half attached to the split idempotent $\bar e$. Otherwise some $u\in\Lambda$ in the
  half attached to $e$ has $B(u,y)\ne1$, and $\rho(u)f_0=f_0$ for the same reason.
\end{enumerate}
Let $N$ be the span of the functions $\rho(u)g-g$ with $u\in Y$ and $g$ a theta function. We show that
$f_0\notin N$; then a linear form $\ell$ that vanishes on $N$ with $\ell(f_0)\ne0$ has the required
properties.

First, for every $m\in N$ there are finitely many $y_1,\dots,y_s\in Y$, one of them $0$, with
$\sum_i\rho(y_i)m=0$. For $m=\rho(u)g-g$, take $y_i=(i-1)u$ for $1\le i\le k$, with $k$ as in (b):
by (a) the sum telescopes to $\rho(ku)g-g=0$. If $m$ and $m'$ have families $(y_i)$ and $(z_j)$, the
family of the $y_i+z_j$ works for $m+m'$, since by (a)
$\sum_{i,j}\rho(y_i+z_j)m=\sum_j\rho(z_j)\sum_i\rho(y_i)m=0$ and likewise for $m'$; and the family of $m$
works for its multiples.

Second, for every finite family $y_1,\dots,y_s$ in $Y$ there is a linear operator $P$ with $Pf_0\ne0$
such that, for every $i$, $\rho(y_i)f_0=f_0$ or $P\rho(y_i)f_0=0$. By induction on $s$: for $s=0$ take
$P=1$. Given $P$ for $y_1,\dots,y_{s-1}$, if $\rho(y_s)f_0=f_0$ keep $P$. Otherwise take $u$ as in (c)
and replace $P$ by $P'=P\circ(\rho(u)-B(u,y_s))$. By (a) and $\rho(u)f_0=f_0$,
$(\rho(u)-B(u,y_s))\rho(y_i)f_0=(B(u,y_i)-B(u,y_s))\rho(y_i)f_0$, which vanishes for $i=s$ and is a
multiple of $\rho(y_i)f_0$ for $i<s$; and $P'f_0=(1-B(u,y_s))Pf_0\ne0$.

Now suppose $f_0\in N$, with a family $y_1,\dots,y_s$ as in the first part, and let $P$ be as in the
second. Applying $P$ to $\sum_i\rho(y_i)f_0=0$ gives $nPf_0=0$, where $n\geq1$ is the number of $i$ with
$\rho(y_i)f_0=f_0$, which includes an $i$ with $y_i=0$. This contradicts $Pf_0\ne0$.
\end{proof}

\begin{lemma}\label{is-eigenvalue-conjugate-split-ideals}
Let $E\supset F$ be number fields.
\begin{enumerate}
\item For every nonzero ideal $I$ of $\mathcal O_E$ there is a nonzero $s\in I$ such that every prime
  ideal of $\mathcal O_E$ that contains $s$ contains $I$.
\item For every nonzero $s\in\mathcal O_E$ there is a nonzero $m\in\mathcal O_F$, divisible by $s$ in
  $\mathcal O_E$, such that every prime ideal of $\mathcal O_F$ that contains $m$ lies under a prime ideal
  of $\mathcal O_E$ that contains $s$.
\item Let $Q$ be a set of prime ideals of $\mathcal O_E$, and let $x_1,\dots,x_r\in E$ be integral at
  every nonzero prime ideal of $\mathcal O_E$ that is not in $Q$. Then there is a nonzero
  $s\in\mathcal O_E$ with $sx_1,\dots,sx_r\in\mathcal O_E$ such that every prime ideal of $\mathcal O_E$
  that contains $s$ lies in $Q$.
\end{enumerate}
\end{lemma}
\leanref{Development/AbelianVarietyTheory/Proofs/AdelicLatticeProofs/ExistsRationalUnitarySplitDenominator.lean}{GlobalTheory.UnRep.LatForm.exists_mem_primes}
\leanref{Development/AbelianVarietyTheory/Proofs/AdelicLatticeProofs/ExistsUnitaryAdelicDecompositionSplitSupport.lean}{GlobalTheory.UnRep.SplitRep.exists_multiple}
\leanref{Development/AbelianVarietyTheory/Proofs/AdelicLatticeProofs/ExistsDenominatorPrimesPrescribed.lean}{GlobalTheory.UnRep.exists_denominator_primes_prescribed}
\begin{proof}
The class groups of $E$ and $F$ are finite; let $h_E$ and $h_F$ be their orders.

(1) The ideal $I^{h_E}$ is principal, say $I^{h_E}=s\mathcal O_E$. Then $s\in I$ and $s\ne0$. A prime
ideal that contains $s$ contains $I^{h_E}$, hence $I$.

(2) The ideal $J=s\mathcal O_E\cap\mathcal O_F$ is nonzero, as $s$ is nonzero and integral over
$\mathcal O_F$. Write $J^{h_F}=m\mathcal O_F$. Then $m\in J$, so $s$ divides $m$ in $\mathcal O_E$. A
prime ideal $\mathfrak p$ of $\mathcal O_F$ that contains $m$ contains $J^{h_F}$, hence $J$. The ring
$\mathcal O_E/s\mathcal O_E$ is integral over its subring $\mathcal O_F/J$, so there is a prime ideal of
it over $\mathfrak p/J$ \cite[Theorem~5.10]{AM69}, and its inverse image in $\mathcal O_E$ lies over
$\mathfrak p$ and contains $s$.

(3) The set $A$ of the $a\in\mathcal O_E$ with $ax_i\in\mathcal O_E$ for every $i$ is an ideal, nonzero
because the $x_i$ have a common denominator. Let $s$ be as in (1) for $A$. Suppose that a prime ideal
$P$ that contains $s$ is not in $Q$. Then $P$ contains $A$. Since each $x_i$ is integral at $P$, there is
$u_i\in\mathcal O_E\smallsetminus P$ with $u_ix_i\in\mathcal O_E$. The product of the $u_i$ lies in $A$
and not in $P$, a contradiction.
\end{proof}

\begin{lemma}\label{is-eigenvalue-conjugate-split-norm-one}
Let $E/F$ be a quadratic extension of number fields with nontrivial automorphism $x\mapsto\bar x$, let
$u\in E$ with $u\bar u=1$, and let $N,D>0$ be integers. There is $u'\in E$ with $u'\bar u'=1$ such that:
\begin{enumerate}
\item $t(u'/u-1)\in D\mathcal O_E$ for some natural number $t$ prime to $N$;
\item there is a nonzero $s\in\mathcal O_E$ with $s(u'-1)\in D\mathcal O_E$ and
  $s(u'^{-1}-1)\in D\mathcal O_E$, such that every prime ideal $P$ of $\mathcal O_E$ that contains $s$
  satisfies $\bar P\ne P$ or $N\in P$.
\end{enumerate}
\end{lemma}
\leanref{Development/AbelianVarietyTheory/Proofs/AdelicLatticeProofs/NormOneApproximationAtLevelPrimesStatement.lean}{GlobalTheory.T.normOne_approximation_at_level_primes}
\leanref{Development/AbelianVarietyTheory/Proofs/AdelicLatticeProofs/NormOneApproximationAtLevelPrimes.lean}{GlobalTheory.UnRep.normOne_approx}
\leanref{Development/AbelianVarietyTheory/Proofs/AdelicLatticeProofs/NormOneApproximationAtLevelPrimes.lean}{GlobalTheory.UnRep.NormOne.hilbert90_matrix_subalgebra_norm_one}
\leanref{Development/AbelianVarietyTheory/Proofs/AdelicLatticeProofs/ValuationEqOneOfNormOneOfNotMoved.lean}{GlobalTheory.UnRep.valuation_eq_one_of_norm_one_of_not_moved}
\begin{proof}
If $N=1$, then $u'=1$, $t=0$ and $s=1$ will do. Let $N>1$. By Hilbert's Theorem 90 there is a nonzero
$\alpha\in\mathcal O_E$ with $u\bar\alpha=\alpha$; put $\beta=\bar\alpha$, so that $u=\alpha/\beta$. Let
$n_1=\alpha\gamma_1$ and $n_2=\beta\gamma_2$, with $\gamma_1,\gamma_2\in\mathcal O_E$, be positive
integers, for example the absolute norms of $\alpha\mathcal O_E$ and $\beta\mathcal O_E$. Write
$Dn_1n_2=D_1D_2$ with $D_1$ dividing a power of $N$ and $D_2$ prime to $N$. Then $D_2$ is prime to
$D_1N$, and $pD_2+qD_1N=1$ for some integers $p,q$. Put
\[
 b=qD_1N,\qquad\alpha'=pD_2\alpha+b,\qquad\beta'=pD_2\beta+b=\bar\alpha',\qquad
 z=-2qD_1+qD_1b+pqD\alpha\gamma_1\gamma_2.
\]
Since $pD_2=1-b$ and $D\alpha\gamma_1\beta\gamma_2=D_1D_2$,
\[
 pD_2\beta'=(1-b)^2\beta+pD_2b=\beta(1+Nz).\tag{$*$}
\]
The element $1+Nz$ is nonzero, as otherwise $N$ would be a unit of $\mathcal O_E$. So $\beta'\ne0$ by
$(*)$, and $u'=\alpha'/\beta'$ satisfies $u'\bar u'=1$.

(1) Modulo the nonzero ideal $(1+Nz)$ the integer $N$ is a unit, with inverse $-z$, and the quotient is
finite; so $N^k\equiv1$ modulo $(1+Nz)$ for some $k\geq1$. Then $t=N^k-1$ is prime to $N$, and
$t=(1+Nz)c$ with $c\in\mathcal O_E$. As $\alpha'\beta-\beta'\alpha=b(\beta-\alpha)$, the identity $(*)$
and $D_1D_2=D\alpha\gamma_1\beta\gamma_2$ give
\[
 t\Big(\frac{u'}u-1\Big)=\frac{(1+Nz)c\,b(\beta-\alpha)}{\beta'\alpha}
 =D\cdot cpqN\gamma_1\gamma_2(\beta-\alpha)\in D\mathcal O_E.
\]

(2) Let $P$ be a prime ideal of $\mathcal O_E$ with $\bar P=P$ and $N\notin P$. The valuations of $u'$
and $\bar u'$ at $P$ agree, as $\bar P=P$, and add up to $0$; so $u'$ is a unit at $P$. If $D\notin P$,
then $(u'-1)/D$ and $(u'^{-1}-1)/D$ are integral at $P$. If $D\in P$, then $D_2\in P$, since $D_1$
divides a power of $N$. Then $\beta'\equiv b=1-pD_2\equiv1$ modulo $P$, and likewise $\alpha'\notin P$;
and
\[
 \frac{u'-1}D=\frac{pD_2(\alpha-\beta)}{D\beta'}=\frac{pn_1n_2(\alpha-\beta)}{D_1\beta'},\qquad
 \frac{u'^{-1}-1}D=\frac{pn_1n_2(\beta-\alpha)}{D_1\alpha'}
\]
are integral at $P$, as $D_1\notin P$. Lemma~\ref{is-eigenvalue-conjugate-split-ideals}(3), applied to
the set of the prime ideals $P$ with $\bar P\ne P$ or $N\in P$ and to the elements $(u'-1)/D$ and
$(u'^{-1}-1)/D$, gives $s$.
\end{proof}

\begin{lemma}\label{is-eigenvalue-conjugate-split-correction}
Let $E/F$ be a quadratic extension of number fields with nontrivial automorphism $x\mapsto\bar x$, let
$G\in GL_n(E)$ with $\bar G^{\mathsf T}=G$, and let $L\subset E^n$ be a finitely generated
$\mathcal O_E$-submodule that spans $E^n$. Write $U(G)(E)$ for the group of the $\gamma\in GL_n(E)$ with
$\bar\gamma^{\mathsf T}G\gamma=G$. Let $\gamma_0\in U(G)(E)$, and let $N,D>0$ be integers. There are
$\gamma\in U(G)(E)$, a natural number $t$ prime to $N$ and a nonzero $s\in\mathcal O_E$ such that
$t(\gamma\gamma_0^{-1}-1)L\subset DL$, $s\gamma L\subset L$ and $s\gamma^{-1}L\subset L$, and every
prime ideal $P$ of $\mathcal O_E$ that contains $s$ satisfies $\bar P\ne P$ or $N\in P$.
\end{lemma}
\leanref{Development/AbelianVarietyTheory/Proofs/AdelicLatticeProofs/ExistsUnitaryCorrectionCongruentOneAtLevelStatement.lean}{GlobalTheory.T.exists_unitary_correction_congruent_one_at_level}
\leanref{Development/AbelianVarietyTheory/Proofs/AdelicLatticeProofs/ExistsUnitaryCorrectionCongruentOneAtLevel.lean}{GlobalTheory.UnRep.Correct.exists_unitary_correction}
\leanref{Development/AbelianVarietyTheory/Proofs/AdelicLatticeProofs/ExistsUnitaryCorrectionCongruentOneAtLevel.lean}{GlobalTheory.UnRep.Correct.congAt_mul}
\leanref{Development/AbelianVarietyTheory/Proofs/AdelicLatticeProofs/ExistsUnitaryCorrectionCongruentOneAtLevel.lean}{GlobalTheory.UnRep.Correct.congAt_conj}
\leanref{Development/AbelianVarietyTheory/Proofs/AdelicLatticeProofs/ExistsNormOneReflectionCorrectionAtLevelStatement.lean}{GlobalTheory.T.exists_norm_one_reflection_correction_at_level}
\leanref{Development/AbelianVarietyTheory/Proofs/AdelicLatticeProofs/ExistsNormOneReflectionCorrectionAtLevel.lean}{GlobalTheory.UnRep.exists_refl_correction}
\leanref{Development/AbelianVarietyTheory/Proofs/AdelicLatticeProofs/UnitaryQuasiReflectionLawsStatement.lean}{GlobalTheory.T.unitaryQuasiReflection_laws}
\leanref{Development/AbelianVarietyTheory/Proofs/AdelicLatticeProofs/UnitaryMatrixEqProdQuasiReflectionsStatement.lean}{GlobalTheory.T.unitaryMatrix_eq_prod_quasiReflections}
\begin{proof}
Call $\delta\in M_n(E)$ \emph{congruent to $1$ modulo $D$ at $N$} if $t(\delta-1)L\subset DL$ for some
natural number $t$ prime to $N$. Two facts about this notion:
\begin{enumerate}
\item[(a)] If $A$ and $B$ are congruent to $1$ modulo $D$ at $N$, with $t_1$ and $t_2$, then so is $AB$,
  with $t_1t_2$. Indeed, for $v\in L$, $t_2Bv=t_2v+t_2(B-1)v$ lies in $L$, and
  $t_1t_2(AB-1)v=t_1(A-1)(t_2Bv)+t_1\,t_2(B-1)v\in DL$.
\item[(b)] Let $R\in GL_n(E)$ and $d,e>0$ with $dRL\subset L$ and $eR^{-1}L\subset L$. If $A$ is
  congruent to $1$ modulo $Dde$ at $N$, with $t$, then so is $RAR^{-1}$ modulo $D$, with $t$. Indeed,
  for $v\in L$, $t(A-1)(eR^{-1}v)=Ddew$ with $w\in L$; applying $R$ and dividing by $e$ gives
  $t(RAR^{-1}-1)v=D\,dRw\in DL$.
\end{enumerate}
For an anisotropic $x\in E^n$, that is $\bar x^{\mathsf T}Gx\ne0$, let $P_x$ be the orthogonal projection
onto $Ex$, $P_xv=(\bar x^{\mathsf T}Gv/\bar x^{\mathsf T}Gx)\,x$, and for $a\in E^\times$ let
$r_{x,a}=1+(a-1)P_x$. Then $r_{x,1}=1$, $r_{x,a}r_{x,a'}=r_{x,aa'}$, and $r_{x,a}\in U(G)(E)$ if
$a\bar a=1$. The group $U(G)(E)$ is generated by these quasi-reflections \cite{Die71}: $\gamma_0$ is a
product $r_{x_1,u_1}\cdots r_{x_k,u_k}$ with $u_i\bar u_i=1$. We argue by induction on $k$. If $k=0$,
then $\gamma_0=1$, and $\gamma=1$, $t=1$ and $s=1$ will do.

Let $k\geq1$, put $x=x_1$, $R=r_{x,u_1}$ and $\gamma'_0=R^{-1}\gamma_0$, a product of $k-1$
quasi-reflections, and choose integers $d,e>0$ with $dRL\subset L$ and $eR^{-1}L\subset L$, which exist
as $L$ is finitely generated and spans $E^n$. The induction hypothesis, applied to $\gamma'_0$ with $Dde$
in place of $D$, gives $\gamma_1\in U(G)(E)$ with $\gamma_1{\gamma'_0}^{-1}$ congruent to $1$ modulo
$Dde$ at $N$ and an element $s_1$ with the properties of the statement for $\gamma_1$.

Next, choose an integer $d'>0$ with $d'P_xL\subset L$, and apply
Lemma~\ref{is-eigenvalue-conjugate-split-norm-one} to $u_1$, $N$ and $Dd'$; it gives $u'$, $t'$ and
$s'$. Since $r_{x,a}-1=(a-1)P_x$, for $v\in L$
\[
 t'(r_{x,u'/u_1}-1)v=t'\Big(\frac{u'}{u_1}-1\Big)P_xv\in D\mathcal O_E\,d'P_xv\subset DL,\qquad
 s'r_{x,u'}^{\pm1}v=s'v+s'(u'^{\pm1}-1)P_xv\in L,
\]
so $r_{x,u'/u_1}$ is congruent to $1$ modulo $D$ at $N$, and $s'r_{x,u'}L$ and $s'r_{x,u'}^{-1}L$ lie
in $L$.

Put $\gamma=r_{x,u'}\gamma_1$ and $s=s's_1$. Then $\gamma\in U(G)(E)$; $s\gamma L\subset L$ and
$s\gamma^{-1}L\subset L$; a prime ideal that contains $s$ contains $s'$ or $s_1$; and
\[
 \gamma\gamma_0^{-1}=r_{x,u'}\gamma_1{\gamma'_0}^{-1}R^{-1}=r_{x,u'/u_1}\cdot R(\gamma_1{\gamma'_0}^{-1})R^{-1},
\]
a product of two matrices congruent to $1$ modulo $D$ at $N$, by (b) for the second factor. By (a),
$\gamma\gamma_0^{-1}$ is congruent to $1$ modulo $D$ at $N$.
\end{proof}

\begin{theorem}[Weak approximation {\cite[Chapter~7]{PR94}}]\label{is-eigenvalue-conjugate-split-approximation}
Let $E/F$ be a quadratic extension of number fields with nontrivial automorphism $x\mapsto\bar x$, let
$G\in GL_n(E)$ with $\bar G^{\mathsf T}=G$, and let $L\subset E^n$ be a finitely generated
$\mathcal O_E$-submodule that spans $E^n$. Put $\mathbf A_{E,f}=\mathbf A_{F,f}\otimes_FE$, and let
$\hat\Lambda\subset\mathbf A_{E,f}^n$ be the subgroup generated by the vectors $ax$ with $a$ an integral
finite adele of $F$ and $x\in L$. Let $g\in GL_n(\mathbf A_{E,f})$ with $\bar g^{\mathsf T}Gg=G$, and let
$N,D>0$ be integers. Then there are $\gamma\in GL_n(E)$ with $\bar\gamma^{\mathsf T}G\gamma=G$ and a
natural number $t$ prime to $N$ with $t(\gamma g-1)\hat\Lambda\subset D\hat\Lambda$.
\end{theorem}
\leanref{Development/AbelianVarietyTheory/Proofs/AdelicLatticeProofs/UnitaryGroupWeakApproximationAtPrimesStatement.lean}{GlobalTheory.T.unitaryGroup_weak_approximation_at_primes}
\leanref{Development/AbelianVarietyTheory/Proofs/AdelicLatticeProofs/UnitaryGroupWeakApproximationAtPrimes.lean}{GlobalTheory.UnRep.WeakApprox.weak_approximation}

\begin{lemma}\label{is-eigenvalue-conjugate-split-finite}
Let $E/F$ be a quadratic extension of number fields with nontrivial automorphism $x\mapsto\bar x$, let
$G\in GL_n(E)$ with $\bar G^{\mathsf T}=G$, and let $L\subset E^n$ be a finitely generated
$\mathcal O_E$-submodule that spans $E^n$. Put $\mathbf A_{E,f}=\mathbf A_{F,f}\otimes_FE$, let
$\hat\Lambda\subset\mathbf A_{E,f}^n$ be the subgroup generated by the vectors $ax$ with $a$ an integral
finite adele of $F$ and $x\in L$, and write $U(G)(E)$ and $U(G)(\mathbf A_f)$ for the groups of the
$\gamma\in GL_n(E)$ and of the $g\in GL_n(\mathbf A_{E,f})$ with $\bar\gamma^{\mathsf T}G\gamma=G$ and
$\bar g^{\mathsf T}Gg=G$. Let $g\in U(G)(\mathbf A_f)$, and let $N>0$ be an integer. Then $g=\gamma hk$
with $\gamma\in U(G)(E)$ and $h,k\in U(G)(\mathbf A_f)$ such that:
\begin{enumerate}
\item $h-1=\varepsilon(h-1)$, where $\varepsilon\in\mathbf A_{F,f}$ is the idempotent of a finite set of
  places of $F$ that split in $E$ and do not divide $N$;
\item $(k-1)\hat\Lambda\subset N\hat\Lambda$ and $k^{-1}\hat\Lambda\subset\hat\Lambda$.
\end{enumerate}
\end{lemma}
\leanref{Development/AbelianVarietyTheory/Proofs/AdelicLatticeProofs/ExistsUnitaryAdelicDecompositionSplitSupportStatement.lean}{GlobalTheory.T.exists_unitary_adelic_decomposition_split_support}
\leanref{Development/AbelianVarietyTheory/Proofs/AdelicLatticeProofs/ExistsUnitaryAdelicDecompositionSplitSupport.lean}{GlobalTheory.UnRep.exists_rep_suppSplit}
\leanref{Development/AbelianVarietyTheory/Proofs/AdelicLatticeProofs/ExistsRationalUnitarySplitDenominator.lean}{GlobalTheory.UnRep.LatForm.exists_rep_split_denominator}
\leanref{Development/AbelianVarietyTheory/Proofs/AdelicLatticeProofs/ExistsRationalUnitarySplitDenominator.lean}{GlobalTheory.UnRep.LatForm.inv_integral}
\leanref{Development/AbelianVarietyTheory/Proofs/AdelicLatticeProofs/IdempotentAdelesAndSmallAdelicMatricesStatement.lean}{GlobalTheory.T.idempotent_adeles_and_small_adelic_matrices}
\leanref{Development/AbelianVarietyTheory/Proofs/AdelicLatticeProofs/ExistsUnitaryAdelicDecompositionSplitSupport.lean}{GlobalTheory.UnRep.SplitRep.isUnA_mix}
\begin{proof}
Three properties of $\hat\Lambda$ are used. (i) For every $M\in M_n(\mathbf A_{E,f})$ there is an integer
$d>0$ with $dM\hat\Lambda\subset\hat\Lambda$, as the entries of $M$ have bounded denominators. (ii)
$c\hat\Lambda\subset\hat\Lambda$ for every integral finite adele $c$ of $F$ and every $c\in\mathcal O_E$.
(iii) If $\delta\in M_n(E)$, $s\in\mathcal O_E$ and $s\delta L\subset L$, then
$s\delta\hat\Lambda\subset\hat\Lambda$; and if $t(\delta-1)L\subset DL$, then
$t(\delta-1)\hat\Lambda\subset D\hat\Lambda$; for $\delta$ commutes with the multiplication by finite
adeles of $F$. As in (a) of the proof of Lemma~\ref{is-eigenvalue-conjugate-split-correction}, if
$t_1(A-1)\hat\Lambda\subset D\hat\Lambda$ and $t_2(B-1)\hat\Lambda\subset D\hat\Lambda$, then
$t_1t_2(AB-1)\hat\Lambda\subset D\hat\Lambda$.

\emph{Step 1: a rational correction.} We show that there are $\gamma_1\in U(G)(E)$, a natural number $t$
prime to $N$ and a nonzero $s\in\mathcal O_E$ such that $x=\gamma_1g$ satisfies $t(x-1)\hat\Lambda\subset
N\hat\Lambda$, $sx\hat\Lambda\subset\hat\Lambda$ and $sx^{-1}\hat\Lambda\subset\hat\Lambda$, and every
prime ideal $P$ of $\mathcal O_E$ that contains $s$ satisfies $\bar P\ne P$ and $N\notin P$.

By (i) there are integers $d_1,d_2>0$ with $d_1g\hat\Lambda\subset\hat\Lambda$ and
$d_2g^{-1}\hat\Lambda\subset\hat\Lambda$; put $M=Nd_1d_2$.
Theorem~\ref{is-eigenvalue-conjugate-split-approximation}, with $M$ in place of $N$ and of $D$, gives
$\gamma_0\in U(G)(E)$ and $t_0$ prime to $M$ with $t_0(\gamma_0g-1)\hat\Lambda\subset M\hat\Lambda$.
Lemma~\ref{is-eigenvalue-conjugate-split-correction}, applied to $\gamma_0$ with $M$ in place of $N$ and
of $D$, gives $\gamma_1$, $t_1$ prime to $M$ and $s_1$. By (iii) and the product rule, $x=\gamma_1g=
(\gamma_1\gamma_0^{-1})(\gamma_0g)$ satisfies $t(x-1)\hat\Lambda\subset M\hat\Lambda\subset N\hat\Lambda$
with $t=t_1t_0$, which is prime to $M$, hence to $N$.

Next, there is a natural number $t'$ prime to $M$ with $t'x^{-1}\hat\Lambda\subset\hat\Lambda$. Indeed,
let $m>0$ with $mx^{-1}\hat\Lambda\subset\hat\Lambda$, by (i). For $w\in\hat\Lambda$ write
$t(x-1)w=Mw'$ with $w'\in\hat\Lambda$; applying $x^{-1}$ gives $tx^{-1}w=tw-Mx^{-1}w'$. By induction on
$j$, for every $z\in\hat\Lambda$ there are $y,w\in\hat\Lambda$ with $t^jx^{-1}z=y+M^jx^{-1}w$. Choose $j$
and $u$ prime to $M$ such that $m$ divides $uM^j$. Then $ut^jx^{-1}z=uy+(uM^j/m)\,mx^{-1}w$ lies in
$\hat\Lambda$, and $t'=ut^j$ will do.

Let $I$ be the ideal of the $a\in\mathcal O_E$ with $ax\hat\Lambda\subset\hat\Lambda$ and
$ax^{-1}\hat\Lambda\subset\hat\Lambda$. It contains $tt'$, as $txz=tz+t(x-1)z$ and $t'x^{-1}z$ lie in
$\hat\Lambda$ for $z\in\hat\Lambda$. It contains $d_1d_2s_1$, as $s_1\gamma_1(d_1gz)\in\hat\Lambda$ and
$d_2g^{-1}(s_1\gamma_1^{-1}z)\in\hat\Lambda$ by (iii). So $I\ne0$; let $s$ be as in
Lemma~\ref{is-eigenvalue-conjugate-split-ideals}(1) for $I$. Then $sx\hat\Lambda\subset\hat\Lambda$ and
$sx^{-1}\hat\Lambda\subset\hat\Lambda$. A prime ideal $P$ that contains $s$ contains $tt'$ and
$d_1d_2s_1$. As $tt'$ is prime to $M$, the two generate the unit ideal, so $M\notin P$. Hence $N\notin P$
and $d_1d_2\notin P$, so $s_1\in P$, and $\bar P\ne P$ by the property of $s_1$.

\emph{Step 2: the splitting.} Let $\gamma_1$, $t$, $s$ and $x$ be as in Step 1. By
Lemma~\ref{is-eigenvalue-conjugate-split-ideals}(2) there is a nonzero $m\in\mathcal O_F$ with $m=sd$ for
some $d\in\mathcal O_E$, such that every prime ideal $\mathfrak p$ of $\mathcal O_F$ that contains $m$ lies
under a prime ideal $P$ that contains $s$. Such a $P$ satisfies $\bar P\ne P$ and $N\notin P$, so
$\mathfrak p$ splits in $E$ and does not divide $N$. Let $\varepsilon$ be the idempotent of the finite set
of the places $\mathfrak p$ that contain $m$. Then $a=(1-\varepsilon)/m$ is an integral finite adele, as
$m$ is a unit at the other places, and $1-\varepsilon=ads$. Put
\[
 h=\varepsilon x+(1-\varepsilon),\qquad k=(1-\varepsilon)x+\varepsilon.
\]
As $\varepsilon$ is an idempotent fixed by the conjugation, $h$ and $k$ lie in $U(G)(\mathbf A_f)$, with
inverses $\varepsilon x^{-1}+1-\varepsilon$ and $(1-\varepsilon)x^{-1}+\varepsilon$, and $hk=x$. So
$g=\gamma hk$ with $\gamma=\gamma_1^{-1}$.

(1) $h-1=\varepsilon(x-1)=\varepsilon(h-1)$.

(2) By (ii), $(1-\varepsilon)y=ad(sy)$ lies in $\hat\Lambda$ whenever $sy\in\hat\Lambda$. Let
$z\in\hat\Lambda$, and put $y=kz-z=(1-\varepsilon)(xz-z)$. It lies in $\hat\Lambda$, as $sxz\in\hat\Lambda$
and $(1-\varepsilon)z\in\hat\Lambda$. Write $t(xz-z)=Nw$ with $w\in\hat\Lambda$; then $ty=Nw''$ with
$w''=(1-\varepsilon)w\in\hat\Lambda$. With integers $\alpha,\beta$ such that $\alpha t+\beta N=1$,
$y=\alpha ty+\beta Ny=N(\alpha w''+\beta y)\in N\hat\Lambda$. Finally
$k^{-1}z=(1-\varepsilon)x^{-1}z+\varepsilon z$ lies in $\hat\Lambda$, as $sx^{-1}z\in\hat\Lambda$.
\end{proof}

\begin{lemma}\label{is-eigenvalue-conjugate-split-decomposition}
Let $D$ be a definite-side datum with host field $E_h/F_h$, let $V$ be a Hermitian space of $D$ with the
completion $\hat\Lambda$ of the lattice of its frame, and let $N>0$. Every $g\in U(V)(\mathbf A)$ is a
product $g=\gamma ahk$ with $\gamma\in U(V)(F_h)$, $a$ of finite component $1$, and $h$ and $k$ of
archimedean component $1$, such that $(k-1)\hat\Lambda\subset N\hat\Lambda$, and such that $h-1$ is
supported at finitely many places of $F_h$ that split in $E_h$ and do not divide $N$.
\end{lemma}
\leanref{Development/AbelianVarietyTheory/Proofs/AdelicLatticeProofs/ExistsRepSuppSplit.lean}{GlobalTheory.UnRep.Space.exists_rep_suppSplit}
\leanref{Development/AbelianVarietyTheory/Proofs/AdelicLatticeProofs/ExistsUnitaryAdelicDecompositionSplitSupportStatement.lean}{GlobalTheory.T.exists_unitary_adelic_decomposition_split_support}
\leanref{Development/AbelianVarietyTheory/Proofs/UnitaryExistence/AdelicUnitaryEqRealMulFiniteSplit.lean}{GlobalTheory.UnitarySplit.real_finite_split}
\leanref{Development/AbelianVarietyTheory/Proofs/UnitaryExistence/RationalPointsAdelicUnitaryEqUnitary.lean}{GlobalTheory.Member.RatElt.rational_elements}
\begin{proof}
Write $n=\operatorname{rk}V$ and $G=G_V$, and let $L\subset E_h^n$ be the lattice of the frame of $V$ in
its coordinates, a finitely generated $\mathcal O_{E_h}$-submodule that spans $E_h^n$; then $\hat\Lambda$ is
the subgroup of Theorem~\ref{is-eigenvalue-conjugate-split-approximation} for $L$. As
$\mathbf A_{E_h}$ is the product of its archimedean part and $\mathbf A_{E_h,f}$, an element of
$U(V)(\mathbf A)$ is a pair of an archimedean component and a finite component $g_f\in U(G)(\mathbf A_f)$,
and every such pair is an element; and $U(V)(F_h)=U(G)(E_h)$, embedded diagonally.

Let $g\in U(V)(\mathbf A)$. Lemma~\ref{is-eigenvalue-conjugate-split-finite}, applied to $g_f$ and $N$,
gives $g_f=\gamma h'k'$ with $\gamma\in U(V)(F_h)$ and $h',k'\in U(G)(\mathbf A_f)$. Let $h$ and $k$ be
the elements of $U(V)(\mathbf A)$ with finite components $h'$ and $k'$ and archimedean component $1$. Then
$a=\gamma^{-1}gk^{-1}h^{-1}$ has finite component $\gamma^{-1}g_fk'^{-1}h'^{-1}=1$, and $g=\gamma ahk$.
The conditions on $h$ and $k$ are (1) and (2) of Lemma~\ref{is-eigenvalue-conjugate-split-finite}.
\end{proof}

\begin{lemma}\label{is-eigenvalue-conjugate-split-character}
Let $D$ be a definite-side datum with host field $E_h/F_h$, let $V$ be a Hermitian space of $D$, and
let $\nu$ be a continuous homomorphism from $U(V)(\mathbf A)$ to the complex numbers of absolute value
$1$. Suppose that $\nu(\gamma)=1$ for every $\gamma\in U(V)(F_h)$, and that $\nu(g)=1$ for every $g$
\emph{supported at} a split idempotent $e$ (Lemma~\ref{is-eigenvalue-conjugate-split-form}), that is,
with $(e+\bar e)(g-1)=g-1$. Then $\nu=1$.
\end{lemma}
\leanref{Development/AbelianVarietyTheory/Proofs/DivisorProofs/UnitaryAdelicCharacterTrivialRatAndSplitStatement.lean}{GlobalTheory.T.unitary_adelic_character_trivial_rat_and_split}
\leanref{Development/AbelianVarietyTheory/Proofs/DivisorProofs/UnitaryAdelicCharacterTrivialRatAndSplit.lean}{GlobalTheory.SplitCharTrivial.eq_one_of_rat_of_split}
\leanref{Development/AbelianVarietyTheory/Proofs/DivisorProofs/UnitaryAdelicCharacterTrivialRatAndSplit.lean}{GlobalTheory.SplitCharTrivial.eq_one_of_finset}
\leanref{Development/AbelianVarietyTheory/Proofs/DivisorProofs/UnitaryAdelicCharacterTrivialRatAndSplit.lean}{GlobalTheory.SplitCharTrivial.eq_one_refl}
\leanref{Development/AbelianVarietyTheory/Proofs/DivisorProofs/UnitaryAdelicCharacterTrivialRatAndSplit.lean}{GlobalTheory.SplitCharTrivial.eq_one_inf}
\leanref{Development/AbelianVarietyTheory/Proofs/DivisorProofs/UnitaryAdelicCharacterTrivialRatAndSplit.lean}{GlobalTheory.SplitCharTrivial.exists_seq}
\begin{proof}
For an idempotent $\varepsilon$ of $\mathbf A_{E_h}$ with $\bar\varepsilon=\varepsilon$ and $g\in
U(V)(\mathbf A)$, the matrix $\varepsilon g+(1-\varepsilon)$ lies in $U(V)(\mathbf A)$, and
$g=(\varepsilon g+1-\varepsilon)((1-\varepsilon)g+\varepsilon)$.

\emph{Step 1: split support.} Let $\varepsilon$ be the idempotent of a finite set $\Sigma$ of places
of $F_h$ that split in $E_h$, and let $g$ satisfy $\varepsilon(g-1)=g-1$. Then $\nu(g)=1$. By induction
on $\Sigma$: for $v\in\Sigma$ with idempotent $\varepsilon_v=e+\bar e$, $e$ a split idempotent, write
$g=(\varepsilon_vg+1-\varepsilon_v)((1-\varepsilon_v)g+\varepsilon_v)$. The first factor is supported at
$e$, and the second satisfies the hypothesis for $\Sigma\smallsetminus\{v\}$.

\emph{Step 2: rational quasi-reflections.} Let $x\in E_h^{\operatorname{rk}V}$ be anisotropic for the
form $h_V$ of $V$, let $P$ be the orthogonal projection onto $E_hx$, and for $u\in E_h$ with $u\bar u=1$
let $r_u=1+(u-1)P\in U(V)(F_h)$. Let $r_u^\infty$ be the element with archimedean component $r_u$ and
finite component $1$. Then $\nu(r_u^\infty)=1$. By Hilbert's Theorem 90, $u=\alpha/\bar\alpha$ with
$\alpha\in\mathcal O_{E_h}$ nonzero. For $k\ge0$ put $N_k=(k+2)!$ and
$w_k=(1+N_k\alpha)/(1+N_k\bar\alpha)$, so that $w_k\bar w_k=1$ and $w_k\to u$ at every archimedean
place. The element $(w_k-1)/N_k$ of $E_h$ is integral at every prime $P$ with $\bar P=P$: if $N_k\in
P$, it is $(\alpha-\bar\alpha)/(1+N_k\bar\alpha)$ with $1+N_k\bar\alpha$ a unit at $P$; otherwise $w_k$,
of norm $1$, is a unit at $P$, and $N_k$ is one too. So its denominator involves only primes $P\ne\bar
P$, which lie over places of $F_h$ that split in $E_h$; let $\varepsilon_k$ be the idempotent of these
finitely many places. Then $(1-\varepsilon_k)(w_k-1)$ is $N_k$ times an integral adele, so that
$(1-\varepsilon_k)(w_k-1)\to0$ in the finite adeles. Write $r_{w_k}=(\varepsilon_kr_{w_k}+1-\varepsilon_k)((1-\varepsilon_k)r_{w_k}+\varepsilon_k)$,
with $\varepsilon_k$ placed at the finite component. The left side is rational and the first factor has
split support, so $\nu$ of the second factor is $1$ by Step 1. The second factor tends to
$r_u^\infty$, so $\nu(r_u^\infty)=1$ by continuity.

\emph{Step 3: archimedean elements.} Let $a$ have finite component $1$, and suppose $\nu(a)\ne1$. The
set $\{g:\nu(g)\ne1\}$ is open and contains $a$. By weak approximation for $U(V)$ at the archimedean
places \cite[Chapter~7]{PR94}, it contains an element $\gamma^\infty$ with archimedean component a
rational $\gamma$ and finite component $1$. The group $U(V)(F_h)$ is generated by the quasi-reflections
$r_u$ of Step 2 \cite{Die71}, so $\gamma^\infty$ is a product of elements $r_u^\infty$, and
$\nu(\gamma^\infty)=1$, a contradiction.

\emph{Step 4.} Let $g\in U(V)(\mathbf A)$ and $n\geq1$. By
Lemma~\ref{is-eigenvalue-conjugate-split-decomposition} with $N=n!$, $g=\gamma_na_nh_nk_n$, and by
the hypothesis, Step 3 and Step 1, $\nu(g)=\nu(k_n)$. As $n!\to0$ in the finite adeles, the conditions
on $k_n$ give $k_n\to1$, so $\nu(g)=\lim\nu(k_n)=1$.
% Step 2 is eq_one_refl: exists_seq (hilbert90_matrix_subalgebra_norm_one, exists_conj, exists_cut with
% exists_denominator_primes_prescribed, tendsto_ratio_inf) and cutA_mul, rfA_mem, tendsto_rfA. Step 3 is
% eq_one_inf: RealApprox.weak_approx_real and UnRep.isUnE_eq_prod_refl, with eq_one_rat by induction on
% the reflections. Step 4 is eq_one_of_rat_of_split with tendsto_one_of_finCong and eq_one_of_suppSplit.
\end{proof}

\begin{theorem}[Theta functions of a principal polarization \cite{Mum83,BL04}]\label{is-eigenvalue-conjugate-vacuum-classical}
Let $U$ be a finite-dimensional real vector space with an $\mathbf R$-linear map $J$ with $J^2=-1$ and an
alternating bilinear form $\omega$ with $\omega(Jx,Jy)=\omega(x,y)$ for all $x,y$ and $\omega(Jx,x)<0$
for $x\ne0$. Let $V$ be a finite-dimensional $\mathbf Q$-vector space, $p\colon V\to U$ an injective
additive map whose image spans $U$, and $L\subset V$ a subgroup such that every $v\in V$ has a positive
multiple in $L$, and $v\in L$ if and only if $\omega(p(v),p(y))\in\mathbf Z$ for every $y\in L$. Let
$c\colon V\to\mathbf C$ take values of absolute value $1$. Then any two differentiable functions
$F_1,F_2$ on $U$ such that, for $k=1,2$, all $w,z\in U$ and all $v\in L$,
\[
 \big(dF_k(w)z-\pi i\,\omega(w,z)F_k(w)\big)+i\big(dF_k(w)(Jz)-\pi i\,\omega(w,Jz)F_k(w)\big)=0,\qquad
 F_k(z+p(v))=c(v)\,e^{-\pi i\,\omega(z,p(v))}F_k(z),
\]
are linearly dependent.
\end{theorem}
\leanref{Development/AbelianVarietyTheory/Proofs/AdelicLatticeProofs/NotLinearIndependentStatement.lean}{GlobalTheory.T.not_linearIndependent_alt}

\begin{lemma}\label{is-eigenvalue-conjugate-vacuum-coordinates}
Let $D$ be a definite-side datum with host field $E_h/F_h$ and element $\xi_0\in E_h$, let $\psi$ be a
standard character of $\mathbf A_{F_h}$, let $\theta$ be an archimedean family whose member at each
$\tau\in\mathcal T$ extends $\tau$, and let $S$ and $T$ be Hermitian spaces of $D$. For $\tau\in\mathcal
T$ put $G_{S,\tau}=\tau(G_S)$ and $G_{T,\tau}=\tau(G_T)$, and let $s_{S,\tau}$ and $s_{T,\tau}$ be Cartan
involutions of them: $s^2=1$, $\bar s^{\mathsf T}Gs=G$ and $Gs$ is positive definite, for
$(G,s)=(G_{S,\tau},s_{S,\tau})$ and for $(G,s)=(G_{T,\tau},s_{T,\tau})$. Let
$R=\prod_{\tau\in\mathcal T}M_{\operatorname{rk}T\times\operatorname{rk}S}(\mathbf C)$, let
$\pi\colon\mathbb W\to R$ be the archimedean component read through $\theta$, and let $\iota\colon
R\to\mathbb W$ be the additive map with $\pi\circ\iota=1$ whose values have finite component $0$. Put
$J(z)_\tau=i\,s_{T,\tau}z_\tau s_{S,\tau}$, and let $\omega$ be the real bilinear form on $R$ with
\[
 \omega(x,y)=\sum_{\tau\in\mathcal T}2\operatorname{Re}\big(\tau(\xi_0)\,h_\tau(x_\tau,y_\tau)\big),\qquad
 h_\tau(x_\tau,y_\tau)=\operatorname{tr}\big(G_{S,\tau}^{-1}\bar y_\tau^{\mathsf T}G_{T,\tau}x_\tau\big).
\]
Then:
\begin{enumerate}
\item $\psi(\frac12\langle x,y\rangle)=e^{-\pi i\,\omega(\pi(x),\pi(y))}$ whenever $x$ or $y$ has finite
  component $0$;
\item $J^2=-1$, $\omega$ is alternating, $\omega(Jx,Jy)=\omega(x,y)$, and $\omega(Jx,x)<0$ for $x\ne0$;
\item for every theta function $f$ annihilated by the annihilation operators for $s_S$ and $s_T$ at
  every $\tau$, $F=f\circ\iota$ is differentiable and satisfies the first equation of
  Theorem~\ref{is-eigenvalue-conjugate-vacuum-classical};
\item the map $p=\pi|_{M_{\operatorname{rk}T\times\operatorname{rk}S}(E_h)}$ is injective and its image
  spans $R$ over $\mathbf R$.
\end{enumerate}
\end{lemma}
\leanref{Development/AbelianVarietyTheory/Proofs/AdelicLatticeProofs/ExistsInverseChart.lean}{GlobalTheory.RealLineAt.exists_inverse_chart}
\leanref{Development/AbelianVarietyTheory/Proofs/AdelicLatticeProofs/StdCharacterHalfSymplecticFormRealPointStatement.lean}{GlobalTheory.T.stdCharacter_halfSymplecticForm_realPoint}
\leanref{Development/AbelianVarietyTheory/Proofs/AdelicLatticeProofs/RealSymplecticFormAlternatingInvariantDefiniteStatement.lean}{GlobalTheory.T.realSymplecticForm_alternating_invariant_definite}
\leanref{Development/AbelianVarietyTheory/Proofs/AdelicLatticeProofs/FockVacuumAnnihilationEquationRealCoordinatesStatement.lean}{GlobalTheory.T.fockVacuum_annihilation_equation_real_coordinates}
\leanref{Development/AbelianVarietyTheory/Proofs/HeightProofs/SpanRealProjectionRationalPointsEqTop.lean}{GlobalTheory.FockCompare.span_piW_ratW_eq_top}
\leanref{Development/AbelianVarietyTheory/Proofs/AdelicLatticeProofs/ApplyAddInfRat.lean}{GlobalTheory.VacUnique.infRat_injective}
\begin{proof}
(1) Let $x$ or $y$ have finite component $0$. Then so does
$\operatorname{tr}(G_S^{-1}\bar y^{\mathsf T}G_Tx)$, and with it $a=\langle x,y\rangle\in\mathbf A_{F_h}$.
The member $\theta_\tau$ of $\theta$ extends $\tau$ and commutes with the conjugation, so
$\theta_\tau$ carries the archimedean component of $\operatorname{tr}(G_S^{-1}\bar y^{\mathsf T}G_Tx)$ to
$h_\tau(\pi(x)_\tau,\pi(y)_\tau)$, and that of $a$, a trace from $\mathbf A_{E_h}$, to
\[
 \tau(\xi_0)h_\tau(\pi(x)_\tau,\pi(y)_\tau)+\overline{\tau(\xi_0)h_\tau(\pi(x)_\tau,\pi(y)_\tau)}
 =2\operatorname{Re}\big(\tau(\xi_0)h_\tau(\pi(x)_\tau,\pi(y)_\tau)\big).
\]
As $\psi(b)=\exp(-2\pi i\sum_\tau\tau(b))$ for $b\in F_h\otimes\mathbf R$,
$\psi(\frac12a)=\exp\big(-\pi i\sum_\tau2\operatorname{Re}(\tau(\xi_0)h_\tau)\big)
=e^{-\pi i\,\omega(\pi(x),\pi(y))}$.

(2) Fix $\tau$ and write $G=G_{S,\tau}$, $s=s_{S,\tau}$, $H=G_{T,\tau}$, $t=s_{T,\tau}$ and
$h=h_\tau$. From $s^2=1$ and $\bar s^{\mathsf T}Gs=G$ follow $(Gs)^{-1}=sG^{-1}$ and
$sG^{-1}\bar s^{\mathsf T}=G^{-1}$, and likewise for $H$ and $t$. The following hold for complex matrices
$x,y$ of size $\operatorname{rk}T\times\operatorname{rk}S$.
\begin{itemize}
\item $J^2=-1$: $i\,t(i\,txs)s=-t^2xs^2=-x$.
\item $h(y,x)=\overline{h(x,y)}$, by taking the conjugate transpose inside the trace, as $G$ and $H$ are
  Hermitian. Since $\xi_0=-1/(\mathrm N\mathfrak c\sqrt{-d})$ and $\tau(\sqrt{-d})=i\sqrt d$ for
  $\tau\in\mathcal T$, $\tau(\xi_0)=i/(\mathrm N\mathfrak c\sqrt d)$ has real part $0$ and positive
  imaginary part; so $\operatorname{Re}(\tau(\xi_0)\overline{h(x,y)})=-\operatorname{Re}(\tau(\xi_0)h(x,y))$,
  and $\omega(y,x)=-\omega(x,y)$.
\item $h(i\,txs,y)=i\operatorname{tr}\big((Gs)^{-1}\bar y^{\mathsf T}(Ht)x\big)$, by
  $s G^{-1}=(Gs)^{-1}$ and the cyclicity of the trace; and
  $h(i\,txs,i\,tys)=\operatorname{tr}\big(sG^{-1}\bar s^{\mathsf T}\bar y^{\mathsf T}\bar t^{\mathsf T}Htx\big)
  =h(x,y)$. So $\omega(Jx,Jy)=\omega(x,y)$.
\item Let $x\ne0$. The matrices $P=(Gs)^{-1}$ and $Q=Ht$ are positive definite; writing
  $P=\bar B^{\mathsf T}B$ and $Q=\bar C^{\mathsf T}C$,
  $\operatorname{tr}(P\bar x^{\mathsf T}Qx)=\operatorname{tr}(\bar Y^{\mathsf T}Y)>0$ with $Y=Cx\bar
  B^{\mathsf T}\ne0$. Hence $2\operatorname{Re}\big(\tau(\xi_0)\,i\operatorname{tr}(P\bar x^{\mathsf
  T}Qx)\big)=-2\operatorname{Im}\tau(\xi_0)\operatorname{tr}(P\bar x^{\mathsf T}Qx)<0$.
\end{itemize}
Summing over $\tau$, the terms of $\omega(Jx,x)$ are $\le0$, and negative at a $\tau$ with
$x_\tau\ne0$; so $\omega(Jx,x)<0$ for $x\ne0$.

(3) Since $f$ is smooth along continuous additive maps from real vector spaces to $\mathbb W$ with finite
component $0$, such as $\iota$, $F$ is differentiable. Let $\tau\in\mathcal T$, let $z$ be a complex
matrix, and let $z^{(\tau)}\in R$ be $z$ at $\tau$ and $0$ at the other embeddings. The real direction
$L_{\tau,z}$ of Lemma~\ref{is-eigenvalue-conjugate-theta-closed} is $t\mapsto\iota(tz^{(\tau)})$, as both
have finite component $0$ and the same image under $\pi$. By (1), for $w\in R$,
\[
 (\rho(\iota(tz^{(\tau)}))f)(\iota(w))=\psi\big(\tfrac12\langle\iota(w),\iota(tz^{(\tau)})\rangle\big)
 f(\iota(w+tz^{(\tau)}))=e^{-\pi it\,\omega(w,z^{(\tau)})}F(w+tz^{(\tau)}),
\]
whose derivative at $t=0$ is $A(z^{(\tau)})$, where $A(y)=dF(w)y-\pi i\,\omega(w,y)F(w)$ for $y\in R$. As
$J(z^{(\tau)})=(J(z)_\tau)^{(\tau)}$, the annihilation operator $\partial_\tau(z)+i\,\partial_\tau(J_\tau
z)$ applied to $f$ at $\iota(w)$ is $A(z^{(\tau)})+iA(J(z^{(\tau)}))$, which is $0$. The map $A$ is
additive, and every $y\in R$ is the sum of the $y_\tau^{(\tau)}$; so $A(y)+iA(Jy)=0$ for every $y$,
which is the first equation of Theorem~\ref{is-eigenvalue-conjugate-vacuum-classical} for $F$.

(4) The set $\mathcal T$ is not empty, and each $\tau$ is injective; so $p(\gamma)_\tau$, the matrix
$\tau(\gamma)$, determines $\gamma$. For the span, let $\Sigma\subset\mathbf C^{\mathcal T}$ be the real
span of the vectors $(\tau(e))_\tau$ with $e\in E_h$, a closed subspace. For $x\in E_h$ the map
$b\mapsto(\theta_\tau(b\otimes x))_\tau=(\theta_\tau(b\otimes1)\tau(x))_\tau$ on $F_h\otimes\mathbf R$ is
continuous and carries $F_h$ into $\Sigma$, hence carries all of $F_h\otimes\mathbf R$ into $\Sigma$, as
$F_h$ is dense there. So $\Sigma$ contains the image of $E_h\otimes_{F_h}(F_h\otimes\mathbf R)$ under the
$\theta_\tau$, which is $\mathbf C^{\mathcal T}$, as it contains the components of the values of $\iota$.
Hence the matrices $p(eE_{ij})$, with $e\in E_h$ and $E_{ij}$ the matrix units, span $R$.
\end{proof}

\begin{lemma}\label{is-eigenvalue-conjugate-vacuum-unique}
Let $D$ be a definite-side datum with host field $E_h/F_h$, let $\psi$ be a standard character of
$\mathbf A_{F_h}$, let $\theta$ be an archimedean family whose member at each $\tau\in\mathcal T$ extends
$\tau$, and let $S$ and $T$ be Hermitian spaces of $D$ with correct signatures, Cartan involutions $s_S$
and $s_T$ at the $\tau$, and self-dual lattices $\Lambda_S$ and $\Lambda_T$. Two Fock vacuum vectors
$f,f'$ for $s_S$, $s_T$ and the lattice $\Lambda$ of the pair given by $\Lambda_S$ and $\Lambda_T$
satisfy $f'=af$ for some $a\in\mathbf C$ with $|a|=1$.
\end{lemma}
\leanref{Formalism/LocalToGlobal/IndefiniteReal/LocalNewformFunction.lean}{Formalism.IndefR.Datum.VacUniqueO}
\leanref{Development/AbelianVarietyTheory/Proofs/IwasawaDecomposition/VacuumUniqueAll.lean}{GlobalTheory.vacuum_unique_all}
\leanref{Development/AbelianVarietyTheory/Proofs/AdelicLatticeProofs/ApplyAddInfRat.lean}{GlobalTheory.VacUnique.vacUniqueO}
\leanref{Development/AbelianVarietyTheory/Proofs/AdelicLatticeProofs/ApplyAddInfRat.lean}{GlobalTheory.VacUnique.apply_add_infRat}
\leanref{Development/AbelianVarietyTheory/Proofs/AdelicLatticeProofs/RationalPointsFinitePartSelfDualLatticeIff.lean}{GlobalTheory.RatSelfDual.finRat_mem_iff}
\begin{proof}
Use the notation of Lemma~\ref{is-eigenvalue-conjugate-vacuum-coordinates}. For
$v\in M=M_{\operatorname{rk}T\times\operatorname{rk}S}(E_h)$ write $v_f$ for the element of $\mathbb W$
with finite component $v$ and archimedean component $0$, so that $v=\iota(p(v))+v_f$. By
Lemma~\ref{th-lattice-signs}(2), every $u\in\Lambda$ has a sign $\varepsilon(u)=\pm1$ with
$\rho(u)g=\varepsilon(u)g$ for every lattice vacuum $g$ for $\Lambda$.

\emph{Step 1: values on $\iota(R)$.} Every $w\in\mathbb W$ is $\gamma+\iota(z)+u$ with $\gamma\in M$,
$z\in R$ and $u\in\Lambda$: the element $w-\iota(\pi(w))$ has archimedean component $0$, and its finite
component is a rational point plus an element of the open subgroup $\Lambda$, by
Theorem~\ref{asm-theta-agree-invariance-tate}(1). For a theta function $g$ that is a lattice vacuum
for $\Lambda$, the quasi-periodicity and
$(\rho(u)g)(\iota(z))=\psi(\frac12\langle\iota(z),u\rangle)g(\iota(z)+u)=g(\iota(z)+u)$, where
$\langle\iota(z),u\rangle=0$ as the two have disjoint supports, give
\[
 g(\gamma+\iota(z)+u)=\psi(-\tfrac12\langle\gamma,\iota(z)+u\rangle)\,\varepsilon(u)\,g(\iota(z)).
\]
The factor does not depend on $g$. So if $g'$ is another such function with $g'\circ\iota=a\,g\circ\iota$,
then $g'=ag$.

\emph{Step 2: the law.} Let $L$ be the set of $v\in M$ with $v_f\in\Lambda$, and let $g$ be as in
Step 1. For $v\in L$ and $z\in R$, $\iota(z+p(v))=v+\iota(z)+(-v)_f$, so the quasi-periodicity, the sign
of $(-v)_f$ and Lemma~\ref{is-eigenvalue-conjugate-vacuum-coordinates}(1) give
\[
 g(\iota(z+p(v)))=c(v)\,e^{-\pi i\,\omega(z,p(v))}\,g(\iota(z)),\qquad
 c(v)=\psi(-\tfrac12\langle v,(-v)_f\rangle)\,\varepsilon((-v)_f),
\]
with $|c(v)|=1$ and $c$ independent of $g$.

\emph{Step 3: the lattice $L$.} Every $v\in M$ has a positive multiple in $L$, as $\Lambda$ is open in
the finite part. For $x,y\in M$, $\psi(\langle x,y\rangle)=1$ since $x,y$ are rational, so
$\psi(\langle x_f,y_f\rangle)=\psi(\langle\iota(p(x)),\iota(p(y))\rangle)^{-1}
=e^{2\pi i\,\omega(p(x),p(y))}$ by Lemma~\ref{is-eigenvalue-conjugate-vacuum-coordinates}(1); so
$\psi(\langle x_f,y_f\rangle)=1$ if and only if $\omega(p(x),p(y))\in\mathbf Z$. If $x\in L$, then
$\omega(p(x),p(y))\in\mathbf Z$ for $y\in L$ by Lemma~\ref{th-lattice-signs}(1). Conversely, suppose
$\omega(p(x),p(y))\in\mathbf Z$ for every $y\in L$, and let $N>0$ with $Nx_f\in\Lambda$. Every
$u'\in\Lambda$ is $N\gamma'$ with $\gamma'$ of archimedean component $0$, and $\gamma'=(y_0)_f+u$ with
$y_0\in M$ and $u\in\Lambda$ by Step 1; so $u'=(Ny_0)_f+Nu$ with $Ny_0\in L$, and
$\psi(\langle x_f,u'\rangle)=\psi(\langle x_f,(Ny_0)_f\rangle)\psi(\langle Nx_f,u\rangle)=1$. By
Lemma~\ref{th-lattice-signs}(1), $x_f\in\Lambda$, that is, $x\in L$.

\emph{Step 4.} By Lemma~\ref{is-eigenvalue-conjugate-vacuum-coordinates}(2)--(4) and Steps 2 and 3,
Theorem~\ref{is-eigenvalue-conjugate-vacuum-classical} applies to $U=R$, $V=M$, $p$, $L$, $c$ and the
functions $F=f\circ\iota$ and $F'=f'\circ\iota$, which are therefore linearly dependent. If $F=0$, then
$f=0$ by Step 1, which contradicts $\|f\|=1$. So $F'=aF$ for some $a$, and $f'=af$ by Step 1. Taking
norms on a fundamental domain of covolume $1$ gives $1=|a|^2$.
\end{proof}

\begin{lemma}\label{is-eigenvalue-conjugate-vacuum-span}
Let $D$, $\psi$, $\theta$, $S$, $T$, $s_S$, $s_T$, $\Lambda_S$, $\Lambda_T$ and $\Lambda$ be as in
Lemma~\ref{is-eigenvalue-conjugate-vacuum-unique}, and let $f_0$ be a Fock vacuum vector for $s_S$, $s_T$
and $\Lambda$. Every theta function $f$ annihilated by the annihilation operators for $s_S$ and $s_T$ at
every $\tau$ is a linear combination of functions $\rho(x)f_0$ with $x\in\mathbb W$ of archimedean
component $0$.
\end{lemma}
\leanref{Development/AbelianVarietyTheory/Proofs/DivisorProofs/MaximalCompactActsOnRealVacuaByDetPowers.lean}{GlobalTheory.VacuumChar.mem_span}
\leanref{Development/AbelianVarietyTheory/Proofs/DivisorProofs/MaximalCompactActsOnRealVacuaByDetPowers.lean}{GlobalTheory.VacuumChar.vac_line}
\leanref{Development/AbelianVarietyTheory/Proofs/DivisorProofs/MemSpanHeisenbergTranslatesOfVacuumStatement.lean}{GlobalTheory.T.mem_span_heisenberg_translates_of_vacuum}
\leanref{Development/AbelianVarietyTheory/Proofs/DivisorProofs/MemSpanHeisenbergTranslatesOfVacuum.lean}{GlobalTheory.HeisVacuum.mem_span_translates_vacuum}
\begin{proof}
Let $X$ be the group of elements of $\mathbb W$ of archimedean component $0$, which contains $\Lambda$,
and $\mathcal M$ the space of theta functions annihilated by the annihilation operators. For $x\in X$,
$\rho(x)$ preserves $\mathcal M$, by Lemma~\ref{is-eigenvalue-conjugate-theta-closed}(1) and because
$\rho(x)$ commutes with $\rho(\iota(z))$ for $z\in R$. Put $\beta(x,y)=\psi(\frac12\langle x,y\rangle)$, a
bicharacter with $\beta(x,x)=1$ and $\rho(x)\rho(y)=\beta(x,y)\rho(x+y)$, so that
$\rho(x)\rho(y)=\beta(x,y)^2\rho(y)\rho(x)$. Then:
\begin{enumerate}
\item[(a)] $x\in X$ lies in $\Lambda$ if and only if $\beta(x,a)^2=1$ for every $a\in\Lambda$, by
  Lemma~\ref{th-lattice-signs}(1); every element of $X$ is divisible by every $N>0$ in $X$; and
  $\Lambda/N\Lambda$ is finite for $N>0$.
\item[(b)] $\rho(u)f_0=\varepsilon(u)f_0$ with $\varepsilon(u)=\pm1$ for $u\in\Lambda$, and
  $\rho(2u)f_0=f_0$ (Lemma~\ref{th-lattice-signs}(2)).
\item[(c)] If $g\in\mathcal M$ satisfies $\rho(u)g=\varepsilon(u)g$ for every $u\in\Lambda$, then $g$ is
  a multiple of $f_0$. Indeed $g$ is then a lattice vacuum, as $f_0$ is; if $g\ne0$, then $g$ divided
  by its norm on a fundamental domain of covolume $1$ is a Fock vacuum vector, and
  Lemma~\ref{is-eigenvalue-conjugate-vacuum-unique} applies.
\item[(d)] $f$ is fixed by $\rho(Nu)$ for every $u\in\Lambda$, for some $N>0$: $f$ is fixed by $\rho(x)$
  for the $x\in X$ in an open subgroup, which contains $N\Lambda$ for some $N>0$.
\end{enumerate}
Put $K=2N$ and $A(u)=\varepsilon(u)^{-1}\rho(u)$ for $u\in\Lambda$. By (b),
$A(u)A(u')=A(u+u')$, $A(Ku)f_0=f_0$ and $A(Ku)f=f$, so $A$ is a representation of the finite group
$Q=\Lambda/K\Lambda$ on the functions fixed by the $\rho(Ku)$. For such a function $g$, the average
$\sum_{q\in Q}A(q)g$ is fixed by every $A(u)$, so it is a multiple of $f_0$ by (c). Let
$T_K=\{t\in X:Kt\in\Lambda\}$. For $t\in T_K$, the function $\rho(-t)f$ is fixed by the $\rho(Ku)$,
since $\beta(Ku,-t)^2=\beta(u,-Kt)^2=1$ by (a); so $\sum_qA(q)\rho(-t)f=c_tf_0$, and applying $\rho(t)$,
\[
 \sum_{q\in Q}\beta(t,q)^2A(q)f=c_t\,\rho(t)f_0,
\]
which lies in the span $\mathcal S$ of the functions $\rho(x)f_0$, $x\in X$. The characters
$t\mapsto\beta(t,q)^2$ of $T_K$ for $q\in Q$ are distinct: if two agree, the difference $q'-q=Ky$ with
$y\in X$ satisfies $\beta(Kt,y)^2=1$ for every $t\in T_K$, that is, $\beta(a,y)^2=1$ for every $a\in\Lambda$,
so $y\in\Lambda$ by (a) and $q=q'$ in $Q$. Distinct characters are linearly independent. For a linear
form $\phi$ vanishing on $\mathcal S$, $\sum_q\beta(t,q)^2\phi(A(q)f)=0$ for every $t\in T_K$, so
$\phi(A(q)f)=0$ for every $q$, and $\phi(f)=0$ for $q=0$. Hence $f\in\mathcal S$.
\end{proof}

\begin{lemma}\label{is-eigenvalue-conjugate-theta-facts}
Let $D$ be a definite-side datum with host field $E_h/F_h$ and coefficient field $L$, let $\psi$ be a
standard character of $\mathbf A_{F_h}$, let $\theta$ be an archimedean family whose member at each
$\tau\in\mathcal T$ extends $\tau$, and let $S$ and $T$ be Hermitian spaces of $D$ with correct
signatures (Definition~\ref{th-spaces}) and Cartan involutions $s_S=(s_{S,\tau})$ and
$s_T=(s_{T,\tau})$ at the $\tau\in\mathcal T$. For a theta function $f$ for $(S,T)$
(Definition~\ref{th-weil-datum}(2)), families $B_S=(B_{S,\tau})$ and $B_T=(B_{T,\tau})$ of bases of $S$
and $T$ at the $\tau$, and a polynomial $P$ in the creation variables, let $f\otimes_{s,B}P$ be the
function obtained from $f$ by applying the polynomial $P$ in the creation operators at the $\tau$
formed with $s_S$, $s_T$, $B_S$ and $B_T$. Let $\Lambda_S$ and $\Lambda_T$ be self-dual lattices of $S$
and $T$, and call a function a \emph{Fock vacuum vector} if it is one for $s_S$, $s_T$ and the lattice
of the pair given by $\Lambda_S$ and $\Lambda_T$ (Chapter~\ref{chapter-unitary-theta}, after
Definition~\ref{th-weil-datum}). For adelic matrices $\beta_S$ and $\beta_T$, a finite set $P_0$ of
places of $E_h$ and characters $\xi$ of their residue fields, a \emph{level vector} with vacuum $f_0$
is a function $f$ on $\mathbb W$ such that $f_0$ is a Fock vacuum vector and $f=T(\Phi)$ for every
place model $T$ of $f_0$ at the places of $P_0$, of which there is at least one. Here $\Phi$ is the
level function of $\beta_S$, $\beta_T$, $P_0$ and $\xi$, and a \emph{place model} of $f_0$ is a linear
map $T$ on functions on $\mathbb W$ that carries the translation by an element $x$ of the part of
$\mathbb W$ at these places, and the multiplication by $t\mapsto\psi(\langle t,y\rangle)$ for an element
$y$ of the part at the complex conjugate places, to the operators $\rho(x)$ and $\rho(y)$ of
Definition~\ref{th-weil-datum}, and that carries the indicator function of the part of the lattice of
the pair at these places to $f_0$.
\begin{enumerate}
\item (Vacua.) Two Fock vacuum vectors $f,f'$ satisfy $f'=cf$ for some $c\in\mathbf C$ with $|c|=1$.
\item (Weil data.) If $\omega$ and $\omega'$ are Weil data for $(S,T)$ with the same characters, then
  $\omega(p)f=\omega'(p)f$ for every $p\in U(S)(\mathbf A)\times U(T)(\mathbf A)$ and every theta
  function $f$.
\item (Level vectors.) A level vector is a theta function annihilated by the annihilation operators
  for $s_S$ and $s_T$ at every $\tau$.
\item (Pure tensors.) If $f$ is a theta function, then so is $f\otimes_{s,B}P$ for every $P$. Let
  $\omega$ be a Weil datum for $(S,T)$, let $(g,h)\in U(S)(\mathbf A)\times U(T)(\mathbf A)$ have
  components $g_\tau$ and $h_\tau$ at $\tau$, and let $s'_S,s'_T,B'_S,B'_T$ satisfy
  $s'_{S,\tau}g_\tau=g_\tau s_{S,\tau}$, $s'_{T,\tau}h_\tau=h_\tau s_{T,\tau}$,
  $B'_{S,\tau}=g_\tau B_{S,\tau}$ and $B'_{T,\tau}=h_\tau B_{T,\tau}$ for every $\tau$. Then
  $\omega(g,h)(f\otimes_{s,B}P)=(\omega(g,h)f)\otimes_{s',B'}P$ for every theta function $f$ and
  every $P$.
\item (Vacua under compact elements.) Let $\omega$ be a Weil datum for $(S,T)$ and $f$ a theta
  function annihilated by the annihilation operators for $s_S$ and $s_T$ at every $\tau$. If
  $g\in U(S)(\mathbf A)$ and $h\in U(T)(\mathbf A)$ have finite components $1$, and their components
  at every $\tau$ commute with $s_{S,\tau}$ and $s_{T,\tau}$, then $\omega(g,h)f=cf$ for some
  $c\in\mathbf C$.
\item (Lowest $K$-types.) Let $R$ be a realization of $D$ (Definition~\ref{is-realisation}), $i$ a
  lift and $\sigma$ an embedding of $L$ that agrees with $\sigma_0$ on $K$. Then the model of $W_i$ at
  $\sigma$ has smooth irreducible lowest $K$-types
  (Definition~\ref{is-eigenvalue-conjugate-k-types}(3)).
\end{enumerate}
\end{lemma}
\leanref{Formalism/LocalToGlobal/IndefiniteReal/LocalNewformFunction.lean}{Formalism.IndefR.Datum.VacUniqueO}
\leanref{Development/AbelianVarietyTheory/Proofs/IwasawaDecomposition/VacuumUniqueAll.lean}{GlobalTheory.vacuum_unique_all}
\leanref{Development/AbelianVarietyTheory/Proofs/DivisorProofs/WeilDatumUniqueOnThetaFunctions.lean}{GlobalTheory.WeilUniq.weilDatum_unique}
\leanref{Formalism/LocalToGlobal/IndefiniteReal/DatumVec.lean}{Formalism.IndefR.Datum.IsThPVecOf}
\leanref{Development/AbelianVarietyTheory/Proofs/DivisorProofs/ThetaPureTensorO.lean}{GlobalTheory.GalLift.Vec.theta_of_thetaAt}
\leanref{Development/AbelianVarietyTheory/Proofs/TwistedOperators/IsFockVacAtOfIsThetaAt.lean}{GlobalTheory.DatumFock.isFockVacAt_thPVec}
\leanref{Development/AbelianVarietyTheory/Proofs/DivisorProofs/ThetaPureTensorO.lean}{GlobalTheory.GalLift.Vec.theta_pureTensorO}
\leanref{Development/AbelianVarietyTheory/Proofs/AdelicLatticeProofs/PureTensorLIntertwine.lean}{GlobalTheory.OrientedHeis.om_pureTensorO}
\leanref{Development/AbelianVarietyTheory/Proofs/IwasawaDecomposition/DefiniteUnitaryRealPointsActFockSubstitution.lean}{GlobalTheory.DefLift.FockAct.om_vac}
\leanref{Development/AbelianVarietyTheory/Proofs/DivisorProofs/HolShapeOfOverK.lean}{GlobalTheory.GalLift.ProjLift.holShape_of_overK}
\leanref{Development/AbelianVarietyTheory/Proofs/IwasawaDecomposition/HolomorphicDiscreteSeriesLowestKTypeSmoothIrredu.lean}{GlobalTheory.HolLowOnce.lowSm_and_lowIrr}
\leanref{Development/AbelianVarietyTheory/Proofs/TwistedOperators/IsFockVacAtOfIsThetaAt.lean}{GlobalTheory.DatumFock.isFockVacAt_of_isThetaAt}
\leanref{Development/AbelianVarietyTheory/Proofs/AdelicLatticeProofs/EqZeroOfLintegralEqZero.lean}{GlobalTheory.OrientedHeis.om_dLine}
\leanref{Development/AbelianVarietyTheory/Proofs/AdelicLatticeProofs/EqZeroOfLintegralEqZero.lean}{GlobalTheory.OrientedHeis.om_dRho}
\leanref{Development/AbelianVarietyTheory/Proofs/AdelicLatticeProofs/SqNormIntegralFundamentalDomainEqL2NormSq.lean}{GlobalTheory.WeilVacMove.isVacO_om}
\begin{proof}
Write $\partial_\tau(z)$ as in Lemma~\ref{is-eigenvalue-conjugate-theta-closed}, and put
$J_\tau z=i\,s_{T,\tau}zs_{S,\tau}$. The annihilation operators at $\tau$ for $s_S$ and $s_T$ are the
operators $\partial_\tau(z)+i\,\partial_\tau(J_\tau z)$, for the complex matrices $z$. The creation
operator of the variable $(\tau,i,j)$ for $s_S,s_T,B_S,B_T$ is a positive multiple of
$\partial_\tau(u)-i\,\partial_\tau(J_\tau u)$ with $u=B_{T,\tau}E_{ij}B_{S,\tau}^{-1}$, $E_{ij}$ the
matrix unit, and $f\otimes_{s,B}P$ is the combination, with the coefficients of $P$, of the iterates of
these operators given by the monomials of $P$, applied to $f$.

(1) This is Lemma~\ref{is-eigenvalue-conjugate-vacuum-unique}.

(2) We may assume $f\ne0$. Choose $w$ with $f(w)\ne0$; then $f_0=\rho(w)f$ is a theta function by
Lemma~\ref{is-eigenvalue-conjugate-theta-closed}(1), and $f_0(0)=f(w)\ne0$. For every $p$,
$\omega'(p)f_0\ne0$, as $\omega'(p^{-1})\omega'(p)=1$. Fix $p=(g,h)$. The operator
$\omega'(p^{-1})\omega(p)$ preserves theta functions and their $L^2$-norms on fundamental domains, and
commutes with every $\rho(w')$ on theta functions, because $\omega(p)\rho(w')=\rho(hw'g^{-1})\omega(p)$
and $\omega'(p^{-1})\rho(hw'g^{-1})=\rho(w')\omega'(p^{-1})$ there. By
Lemma~\ref{is-eigenvalue-conjugate-theta-schur} it acts on them as a scalar $c(p)$ with $|c(p)|=1$; so
$\omega(p)=c(p)\omega'(p)$ on theta functions. Evaluating at $f_0$:
\begin{itemize}
\item $c$ is multiplicative, since
  $c(pq)\omega'(pq)f_0=\omega(p)\omega(q)f_0=c(q)\omega(p)\omega'(q)f_0=c(p)c(q)\omega'(pq)f_0$ and
  $\omega'(pq)f_0\ne0$;
\item $c$ is continuous, since near $p$ it is the quotient $(\omega(q)f_0)(x)/(\omega'(q)f_0)(x)$ for an
  $x$ with $(\omega'(p)f_0)(x)\ne0$, and both are continuous in $q$;
\item $c(\gamma,\gamma')=1$ for $\gamma\in U(S)(F)$ and $\gamma'\in U(T)(F)$, since
  $(\omega(\gamma,\gamma')f_0)(0)=f_0(0)=(\omega'(\gamma,\gamma')f_0)(0)$ and $f_0(0)\ne0$;
\item $c(g,h)=1$ if $g$ and $h$ are supported at a split idempotent $e$
  (Lemma~\ref{is-eigenvalue-conjugate-split-character}). Indeed, by
  Lemma~\ref{is-eigenvalue-conjugate-split-form} there are a linear form $\ell$ invariant under $\rho(u)$
  on theta functions for the $u$ in the half of $\mathbb W$ attached to $\bar e$, and a theta function
  $f'$ with $\ell(f')\ne0$. By the condition on
  Weil data at split places, $\ell(\omega(g,h)f')=r\chi\,\ell(f')$ and $\ell(\omega'(g,h)f')=r'\chi\,
  \ell(f')$ with $r,r'>0$ and the same nonzero value $\chi$ of the characters at the determinants. With
  $\omega(g,h)f'=c(g,h)\omega'(g,h)f'$ this gives $r=c(g,h)r'$; as $|c(g,h)|=1$, $r=r'$ and $c(g,h)=1$.
\end{itemize}
The restrictions of $c$ to $U(S)(\mathbf A)$ and $U(T)(\mathbf A)$ are continuous unitary characters
trivial on the rational points and on the elements supported at split idempotents, hence trivial by
Lemma~\ref{is-eigenvalue-conjugate-split-character}. So
$c(p)=c(g,1)c(1,h)=1$, and $\omega(p)f=\omega'(p)f$.

(3) Let $f$ be a level vector with vacuum $f_0$, and let $T$ be a place model of $f_0$. The elements
$x$ and $y$ in the definition of a place model have archimedean component $0$. So if $A$ is a linear
operator on the functions on $\mathbb W$ with $A\rho(u)=\rho(u)A$ for every $u\in\mathbb W$ of
archimedean component $0$, and $Af_0=0$, then $T+A\circ T$ is again a place model of $f_0$: it carries
the translation by $x$ and the multiplication attached to $y$ to $\rho(x)(T+A\circ T)$ and
$\rho(y)(T+A\circ T)$, and the indicator function to $f_0+Af_0=f_0$. Then $f=(T+A\circ T)(\Phi)=f+Af$,
so $Af=0$. We apply this twice.

First, let $\ell$ be a linear form on the functions on $\mathbb W$ that vanishes on the theta functions,
and put $(Ag)(w)=\ell(\rho(w)g)$. Since
$\rho(w)\rho(u)=\psi(\frac12\langle w,u\rangle)\rho(w+u)$, $A$ commutes with every $\rho(u)$. The Fock
vacuum vector $f_0$ is a theta function, so each $\rho(w)f_0$ is one by
Lemma~\ref{is-eigenvalue-conjugate-theta-closed}(1), and $Af_0=0$. Hence $Af=0$, and
$\ell(f)=(Af)(0)=0$. As $\ell$ was arbitrary, $f$ is a theta function.

Second, fix $\tau$ and $z$, and let $l$ be a linear form on the functions $\mathbf R\to\mathbf C$ that
is the derivative at $0$ on those differentiable at $0$. Put
\[
 (Ag)(w)=l\big(t\mapsto(\rho(L_{\tau,z}(t))g)(w)\big)+i\,l\big(t\mapsto(\rho(L_{\tau,J_\tau z}(t))g)(w)\big).
\]
For $u$ of archimedean component $0$, $\rho(L(t))\rho(u)=\rho(u)\rho(L(t))$ for every real direction
$L$, as $\langle L(t),u\rangle=0$; so $(A\rho(u)g)(w)=\psi(\frac12\langle w,u\rangle)(Ag)(w+u)$, that is
$A\rho(u)=\rho(u)A$. On a theta function $g$ the functions $t\mapsto(\rho(L(t))g)(w)$ are differentiable
by Lemma~\ref{is-eigenvalue-conjugate-theta-closed}(2), so $Ag$ is the annihilation operator
$\partial_\tau(z)g+i\,\partial_\tau(J_\tau z)g$. Hence $Af_0=0$, so $Af=0$, and $f$, a theta function by
the first application, is annihilated by this operator.

(4) By Lemma~\ref{is-eigenvalue-conjugate-theta-closed}(2) every $\partial_\tau(z)$ maps theta functions
to theta functions; so do the creation operators and their iterates, and $f\otimes_{s,B}P$ is a theta
function.

For the second assertion we first show, for a theta function $f$ and real directions $L$ and
$L'=hLg^{-1}$,
\[
 \omega(g,h)\,\partial_Lf=\partial_{L'}\,\omega(g,h)f.\tag{$*$}
\]
Put $f'=\omega(g,h)f$ and $\delta=\partial_{L'}f'-\omega(g,h)\partial_Lf$, theta functions by
Definition~\ref{th-weil-datum}(3) and Lemma~\ref{is-eigenvalue-conjugate-theta-closed}(2). For $t>0$
put $r_t=t^{-1}(\rho(L(t))f-f)-\partial_Lf$ and $r'_t=t^{-1}(\rho(L'(t))f'-f')-\partial_{L'}f'$; by
Lemma~\ref{is-eigenvalue-conjugate-theta-closed}(3) there are $C,C'$ with $|r_t|\le Ct$ and $|r'_t|\le
C't$ on $\mathbb W$. As $\omega(g,h)\rho(L(t))=\rho(L'(t))\omega(g,h)$ on theta functions,
$\delta=\omega(g,h)r_t-r'_t$. Fix a Haar measure on $\mathbb W$ and a fundamental domain $\mathcal F$ of
volume $1$ for its rational points, and let $\|\ \|$ be the $L^2$-norm on $\mathcal F$. Since
$\omega(g,h)$ preserves it on theta functions, and $r_t$ is one by
Lemma~\ref{is-eigenvalue-conjugate-theta-closed}(1) and~(2), $\|\delta\|\le\|r_t\|+\|r'_t\|\le(C+C')t$ for every
$t>0$. So $\|\delta\|=0$; as $\delta$ is continuous and $|\delta|$ is invariant under the rational
points, $\delta=0$, as in the proof of Lemma~\ref{is-eigenvalue-conjugate-invariant-form}(1).

The real direction $hL_{\tau,z}g^{-1}$ has component $h_\tau zg_\tau^{-1}$ at $\tau$ and $0$ at the
other embeddings, so it is $L_{\tau,h_\tau zg_\tau^{-1}}$, and $(*)$ gives
$\omega(g,h)\partial_\tau(z)f=\partial_\tau(h_\tau zg_\tau^{-1})\omega(g,h)f$. Now
$h_\tau B_{T,\tau}E_{ij}B_{S,\tau}^{-1}g_\tau^{-1}=B'_{T,\tau}E_{ij}{B'_{S,\tau}}^{-1}$, and
$h_\tau(i\,s_{T,\tau}us_{S,\tau})g_\tau^{-1}=i\,s'_{T,\tau}(h_\tau ug_\tau^{-1})s'_{S,\tau}$ because
$h_\tau s_{T,\tau}=s'_{T,\tau}h_\tau$ and $s_{S,\tau}g_\tau^{-1}=g_\tau^{-1}s'_{S,\tau}$. So
$\omega(g,h)$ carries each creation operator for $s,B$, applied to a theta function, to the creation
operator for $s',B'$ applied to its image. By induction on the monomials of $P$,
$\omega(g,h)(f\otimes_{s,B}P)=(\omega(g,h)f)\otimes_{s',B'}P$.

(5) Let $\Lambda^0_S$ and $\Lambda^0_T$ be the lattices of the frames of $S$ and $T$, which are
self-dual, and let $f_0$ be a Fock vacuum vector for $s_S$, $s_T$ and the lattice of the pair given by
$\Lambda^0_S$ and $\Lambda^0_T$, which exists by Lemma~\ref{is-eigenvalue-conjugate-vacuum-exists}(2).
Then $f_1=\omega(g,h)f_0$ is another one. It is a theta function by
Definition~\ref{th-weil-datum}(3). The elements $u$ of the lattice of the pair have archimedean component
$0$, so $hug^{-1}=u$, as $g$ and $h$ have finite components $1$; hence $\rho(2u)f_1=\omega(g,h)\rho(2u)f_0
=f_1$, and $\rho(u)f_1=f_1$ for the $u$ in the half of $\mathbb W$ attached to a split place. Its norm
is $1$, as $\omega(g,h)$ preserves the $L^2$-norm on fundamental domains. By $(*)$, and since $g_\tau$
and $h_\tau$ commute with $s_{S,\tau}$ and $s_{T,\tau}$, $\omega(g,h)$ carries
$(\partial_\tau(z)+i\,\partial_\tau(J_\tau z))f_0$ to
$(\partial_\tau(z')+i\,\partial_\tau(J_\tau z'))f_1$ with $z'=h_\tau zg_\tau^{-1}$; every $z'$ has this
form, so $f_1$ is annihilated by the annihilation operators. By (1), applied to $\Lambda^0_S$ and
$\Lambda^0_T$, $f_1=cf_0$ for some $c$. The theta function $f$, annihilated by the annihilation
operators, is a combination of functions $\rho(x)f_0$ with $x\in\mathbb W$ of archimedean component
$0$, by Lemma~\ref{is-eigenvalue-conjugate-vacuum-span}. For such $x$, $hxg^{-1}=x$, so
$\omega(g,h)\rho(x)f_0=\rho(x)\omega(g,h)f_0=c\,\rho(x)f_0$. Hence $\omega(g,h)f=cf$.

(6) The label of $\sigma$ at $\tau$ is the element $\tau'\in\mathcal T$ for which $\sigma$, composed with
the embedding of $E_h$ into $L$ attached to $\tau'$, is $\tau$ or its complex conjugate. As $\sigma$
agrees with $\sigma_0$ on $K$, it is $\tau$ at every $\tau$. The model of $W_i$ at $\sigma$ then has at
$\tau$ the holomorphic parameter of the Harish-Chandra parameter $\lambda(\tau')$ of $D$:
$\lambda_\tau=(v_1,\dots,v_{k-1},-b_1,\dots,-b_m)$, a strictly decreasing sequence of integers, and,
with $(p,q)$ the signature of $W_i$, $\epsilon_{\tau,a}$ is $1$ for the first $p$ indices $a$ and $-1$
for the others. The rank $r$ of $W_i$ is odd, so $\lambda_{\tau,a}-\frac{r-1}2$ is an integer for every
$a$, and $\lambda_\tau$ is a Harish-Chandra
parameter: $\sigma_\tau$ is a holomorphic discrete series representation, whose lowest $K$-type has the
weight $\Lambda_\tau$. This $K$-type occurs in $\sigma_\tau$ with multiplicity one \cite{Sch75}, so the
sum $\kappa_\tau$ of the $K$-types of weight $\Lambda_\tau$ is one such $K$-type; and its vectors are
$K$-finite, hence smooth \cite{HC53}.
% (1) is VacUnique.vacUniqueO, proved in Lemma is-eigenvalue-conjugate-vacuum-unique.
% (2) inputs: Lemmas is-eigenvalue-conjugate-theta-schur, is-eigenvalue-conjugate-split-form and
% is-eigenvalue-conjugate-split-character. (5) inputs: Lemmas is-eigenvalue-conjugate-vacuum-exists(2)
% and is-eigenvalue-conjugate-vacuum-span.
\end{proof}

\begin{lemma}\label{is-eigenvalue-conjugate-level-vectors}
In the setting of Lemma~\ref{is-eigenvalue-conjugate-theta-facts}, two level vectors $f$ and $f'$ for
the same $\beta_S$, $\beta_T$, $P_0$ and $\xi$ satisfy $f'=af$ for some $a\in\mathbf C$ with $|a|=1$.
\end{lemma}
\leanref{Development/AbelianVarietyTheory/Proofs/DivisorProofs/ThetaLiftUniqueUpToUnit.lean}{GlobalTheory.ConjLiftUnique.vec_unique}
\begin{proof}
Let $f_0$ and $f'_0$ be the vacua of $f$ and $f'$. The idempotent of $\mathbf A_{E_h}$ attached to the
places of $P_0$ is the sum of the idempotents attached to its places, each of which is unique; so the
place models for $f$ and $f'$ are taken at the same idempotent. By
Lemma~\ref{is-eigenvalue-conjugate-theta-facts}(1), $f'_0=af_0$ for some $a$ with $|a|=1$. Let $T$ be
a place model of $f_0$. The conditions on a place model are linear in $T$, except the last, which is
multiplied by $a$ when $T$ is; so $aT$ is a place model of $af_0=f'_0$. Hence $f=T(\Phi)$ and
$f'=(aT)(\Phi)=af$.
\end{proof}

\begin{lemma}\label{is-eigenvalue-conjugate-pure-tensors}
Let $D$, $\psi$ and $\theta$ be as in Lemma~\ref{is-eigenvalue-conjugate-theta-facts}, and let $V$
and $W$ be Hermitian spaces of $D$ with correct signatures such that $1$ is a Cartan involution of $V$
at every $\tau$. Let $M$ be a model of discrete series for $U(W)$
(Definition~\ref{asm-dependence-rows-models}(1)) with representations $\sigma_\tau$, frames $f_\tau$,
signature matrices $J_\tau$ and Cartan involutions $s_\tau=f_\tau^{-1}J_\tau f_\tau$; let $K_\tau$ be
the group of unitary matrices that commute with $J_\tau$, and for $k\in K_\tau$ let $k^\theta\in
U(W)(\mathbf A)$ be the image of $f_\tau^{-1}kf_\tau$ under $\theta$ at $\tau$. Suppose that $M$ has
smooth irreducible lowest $K$-types $\kappa_\tau$ (Definition~\ref{is-eigenvalue-conjugate-k-types}(3)):
at every $\tau$ the sum $\kappa_\tau$ of the $K$-types of $\sigma_\tau$ of weight $\Lambda_\tau$
consists of smooth vectors and is itself a $K$-type of weight $\Lambda_\tau$. Fix integers $p_\tau$, let
$\epsilon_\tau$ be the diagonal matrix with entries $1$ at the first $p_\tau$ places and $-1$ at the
others, and call a family $B=(B_\tau)$ of matrices \emph{adapted} if
$B_\tau^*G_W(\tau)B_\tau=\epsilon_\tau$ and $s_\tau B_\tau=B_\tau\epsilon_\tau$ for every $\tau$. Fix
a family $B_V$ of bases of $V$ at the $\tau$, and for a theta function $f$ for $(V,W)$ write
$f\otimes_BP$ for $f\otimes_{s,B}P$ formed with the Cartan involutions $1$ and $s_\tau$ and the bases
$B_V$ and $B$.

Let $\omega$ be a Weil datum for $(V,W)$, $f$ a theta function, $B$ an adapted family, $\Delta$ a
polynomial and $(W_c,\rho)$ a coefficient datum for $V$ (Lemma~\ref{is-eigenvalue-conjugate-coefficients}).
An \emph{io map} for $(\omega,f,B,\Delta)$ and $W_c$ is a linear map $\mathrm{io}$ from
$\bigotimes_\tau\kappa_\tau\otimes W_c$ to the polynomials such that:
\begin{enumerate}
\item[(i)] $\omega(u,1)(f\otimes_B\mathrm{io}(z\otimes w))=f\otimes_B\mathrm{io}(z\otimes\rho(u)w)$
  for every $u\in U(V)(\mathbf A)$ with finite component $1$, every $z\in\bigotimes_\tau\kappa_\tau$
  and every $w\in W_c$;
\item[(ii)] $\omega(1,k^\theta)(f\otimes_B\mathrm{io}(v\otimes w))=f\otimes_B\mathrm{io}(v'\otimes w)$
  for every $\tau$, every $k\in K_\tau$, every tensor $v=\otimes_\rho v_\rho$ with $v_\rho\in\kappa_\rho$
  and every $w$, where $v'$ is $v$ with $v_\tau$ replaced by $\sigma_\tau(k)v_\tau$;
\item[(iii)] $\mathrm{io}(v\otimes w)=\Delta$ for some tensor $v$ of unit vectors and some unit vector
  $w$.
\end{enumerate}
Let $\omega$ and $\omega'$ be Weil data for $(V,W)$ with the same characters, $f$ a theta function
annihilated by the annihilation operators for $1$ and $s_\tau$ at every $\tau$, $a\in\mathbf C$ with
$|a|=1$, $B$ and $B'$ adapted families and $\Delta$ a polynomial. Let $(W_c,\rho)$ and $(W'_c,\rho')$ be
coefficient data for $V$ with a linear isomorphism $\Phi\colon W_c\to W'_c$ that preserves norms and
satisfies $\Phi\rho(h)=\rho'(h)\Phi$ for every $h$. Let $\mathrm{io}$ be an io map for
$(\omega,f,B,\Delta)$ and $W_c$, and $\mathrm{io}'$ one for $(\omega',af,B',\Delta)$ and $W'_c$. Then
there is $\kappa\in\mathbf C$ with $|\kappa|=1$ and
\[
 (af)\otimes_{B'}\mathrm{io}'(z\otimes\Phi w)=\kappa\,f\otimes_B\mathrm{io}(z\otimes w)
\]
for every $z\in\bigotimes_\tau\kappa_\tau$ and every $w\in W_c$.
\end{lemma}
\leanref{Development/AbelianVarietyTheory/Proofs/DivisorProofs/ThetaLiftUniqueUpToUnit.lean}{GlobalTheory.ConjLiftUnique.io_unique}
\leanref{Development/AbelianVarietyTheory/Proofs/DivisorProofs/ThetaLiftUniqueUpToUnit.lean}{GlobalTheory.ConjLiftUnique.exists_move}
\leanref{Development/AbelianVarietyTheory/Proofs/DivisorProofs/ThetaLiftUniqueUpToUnit.lean}{GlobalTheory.ConjLiftUnique.exists_kset}
\leanref{Development/AbelianVarietyTheory/Proofs/DivisorProofs/ThetaLiftUniqueUpToUnit.lean}{GlobalTheory.ConjLiftUnique.exists_const_of_eqv}
\leanref{Development/AbelianVarietyTheory/Proofs/DivisorProofs/ThetaLiftUniqueUpToUnit.lean}{GlobalTheory.ConjLiftUnique.exists_kop}
\leanref{Development/AbelianVarietyTheory/Proofs/IwasawaDecomposition/EqTopOfTranslates.lean}{GlobalTheory.DefLift.eq_top_of_translates}
\begin{proof}
\emph{Step 0: the lowest $K$-types.} Fix $\tau$. As $\kappa_\tau$ is a $K$-type, it is
finite-dimensional and $\sigma_\tau(k)$ maps it into itself for every $k\in K_\tau$; its vectors and
their translates are smooth by hypothesis. Each $\sigma_\tau(k)$ preserves the inner product, because
$\sigma_\tau$ is a unitary representation of the unitary group of $J_\tau$, which contains $K_\tau$.
The space $\kappa_\tau$ is irreducible under these operators. Indeed, let $U\subset\kappa_\tau$ be a
subspace stable under them that contains a vector $v\ne0$, and let $P$ be the span of the vectors
$\sigma_\tau(k)v$ with $k\in K_\tau$. Then $P\subset U\subset\kappa_\tau$; $P$ is stable under
$\sigma_\tau(K_\tau)$, since $\sigma_\tau(k)\sigma_\tau(k')v=\sigma_\tau(kk')v$ and $K_\tau$ is a group;
and $P\ni\sigma_\tau(1)v=v\ne0$. As $\kappa_\tau$ is a $K$-type, its only stable subspaces are $0$ and
$\kappa_\tau$, so $P=\kappa_\tau$ and $U=\kappa_\tau$.

For $z\in\bigotimes_\tau\kappa_\tau$ and $w\in W_c$ put $F(z,w)=f\otimes_B\mathrm{io}(z\otimes w)$ and
$F'(z,w)=(af)\otimes_{B'}\mathrm{io}'(z\otimes\Phi w)$. These are theta functions by
Lemma~\ref{is-eigenvalue-conjugate-theta-facts}(4), and $\omega'$ agrees with $\omega$ on them by
Lemma~\ref{is-eigenvalue-conjugate-theta-facts}(2). So (i) and (ii) for $\mathrm{io}'$, with
$\Phi\rho(u)=\rho'(u)\Phi$, give $\omega(u,1)F'(z,w)=F'(z,\rho(u)w)$ and
$\omega(1,k^\theta)F'(v,w)=F'(v',w)$, and (i) and (ii) for $\mathrm{io}$ give the same identities for
$F$. Let $b$ be the form of Lemma~\ref{is-eigenvalue-conjugate-invariant-form}(1) for $\omega$, and
write $F(v,w)$ for $F(\otimes_\tau v_\tau,w)$ when $v\in\prod_\tau\kappa_\tau$.

\emph{Step 1: the pairings.} Let $G_1,G_2$ be two of $F,F'$, possibly equal. There is
$c\in\mathbf C$ with $b(G_1(v,w),G_2(v',w'))=cP(v,w;v',w')$ for all $v,v'\in\prod_\tau\kappa_\tau$ and
$w,w'\in W_c$, with $P$ as in Lemma~\ref{is-eigenvalue-conjugate-forms}. Indeed, fix $\tau$, $v,v'$ and
$w,w'$. The form $(a,a')\mapsto b(G_1(v[\tau\leftarrow a],w),G_2(v'[\tau\leftarrow a'],w'))$ on
$\kappa_\tau$ is invariant under the operators $\sigma_\tau(k)$, $k\in K_\tau$, by (ii) and the
invariance of $b$ under $\omega(1,k^\theta)$. By Step 0, $\kappa_\tau$ is finite-dimensional, these
operators preserve the inner product, and $\kappa_\tau$ is irreducible under them; so by
Lemma~\ref{is-eigenvalue-conjugate-forms}(1) the form is a multiple of $\langle a,a'\rangle$. Fix
$v,v'$. The form $(w,w')\mapsto b(G_1(v,w),G_2(v',w'))$ on $W_c$ is invariant under $\rho(u)$ for
$u\in U(V)(\mathbf A)$ with finite component $1$, by (i) and the invariance of $b$. These operators
preserve the inner product, and $W_c$ is irreducible under them, because every element of
$U(V)(\mathbf A)$ is such a $u$ times an element of $U(V)(\mathbf A_{F_h,f})$, on which $\rho$ is
trivial. So by Lemma~\ref{is-eigenvalue-conjugate-forms}(1) this form is a multiple of
$\langle w,w'\rangle$. As $W_c$ has a unit vector by (iii), Lemma~\ref{is-eigenvalue-conjugate-forms}(2)
gives $c$.

\emph{Step 2: moving the bases.} Let $v_0$ and $w_0$ be the unit vectors of (iii) for $\mathrm{io}$, so
that $F(v_0,w_0)=f\otimes_B\Delta$. For a set $s$ of embeddings let $B_s$ be the family equal to $B'$
at the $\tau\in s$ and to $B$ elsewhere. By induction on $s$ there are $v_s\in\prod_\tau\kappa_\tau$
with unit components and $l_s\in\mathbf C$ with $|l_s|=1$ and $F(v_s,w_0)=l_s\,f\otimes_{B_s}\Delta$:
start with $v_\emptyset=v_0$ and $l_\emptyset=1$, and let $\tau\notin s$. The matrix
$U=B'_\tau B_\tau^{-1}$ satisfies $U^*G_W(\tau)U=G_W(\tau)$ and $Us_\tau=s_\tau U$, because $B_\tau$
and $B'_\tau$ satisfy the same equations with the same $\epsilon_\tau$. As
$f_\tau^*J_\tau f_\tau=G_W(\tau)$, the matrix $k=f_\tau Uf_\tau^{-1}$ satisfies $k^*J_\tau k=J_\tau$
and $kJ_\tau=J_\tau k$, so $k^*k=k^*J_\tau kJ_\tau=1$ and $k\in K_\tau$. The element $h=k^\theta$ of
$U(W)(\mathbf A)$ has finite component $1$, component $U$ at $\tau$ and $1$ at the other embeddings,
and $U$ commutes with $s_\tau$. So Lemma~\ref{is-eigenvalue-conjugate-theta-facts}(5) gives
$\omega(1,h)f=a'f$ for some $a'$, and $|a'|=1$ by Lemma~\ref{is-eigenvalue-conjugate-invariant-form}(2).
Lemma~\ref{is-eigenvalue-conjugate-theta-facts}(4), with the same Cartan involutions and the bases
$hB_s=B_{s\cup\{\tau\}}$, gives
$\omega(1,h)(f\otimes_{B_s}\Delta)=(\omega(1,h)f)\otimes_{B_{s\cup\{\tau\}}}\Delta=a'f\otimes_{B_{s\cup\{\tau\}}}\Delta$.
By (ii), if $v_{s\cup\{\tau\}}$ is $v_s$ with its component at $\tau$ replaced by the unit vector
$\sigma_\tau(k)v_{s,\tau}$, then
$F(v_{s\cup\{\tau\}},w_0)=\omega(1,h)F(v_s,w_0)=l_sa'\,f\otimes_{B_{s\cup\{\tau\}}}\Delta$. For the set
of all embeddings this gives $v_1$ and $l$ with $|l|=1$ and $F(v_1,w_0)=l\,f\otimes_{B'}\Delta$.

\emph{Step 3: the phase.} Let $v'_0$ and $w'_0$ be the unit vectors of (iii) for $\mathrm{io}'$, and
let $w_1\in W_c$ with $\Phi w_1=w'_0$, a unit vector. Then
$F'(v'_0,w_1)=(af)\otimes_{B'}\Delta=a\,f\otimes_{B'}\Delta$, so $F(v_1,w_0)=la^{-1}F'(v'_0,w_1)$ with
$|la^{-1}|=1$. By Step 1, Lemma~\ref{is-eigenvalue-conjugate-forms}(3) applies to $b$, $F$ and $F'$,
and gives $\kappa$ with $|\kappa|=1$ and $F'(v,w)=\kappa F(v,w)$ for all $v$ and $w$. Both sides are
linear in the tensor $\otimes_\tau v_\tau$, so $F'(z,w)=\kappa F(z,w)$ for every $z$.
\end{proof}

\begin{lemma}\label{is-eigenvalue-conjugate-good}
Let $R$ be a realization of a definite-side datum $D$ (Definition~\ref{is-realisation}), with host field
$E_h$, imaginary quadratic subfield $K=\mathbf Q(\sqrt{-d})$ and coefficient field $L$, and with one-map
data that have a weight tie (Definition~\ref{is-one-map-data}). The core comes with a square root $r$
of $-d$ in $E_h$, which every $\tau\in\mathcal T$ carries to one complex number $\delta_0$, and its
coefficient field $K_s$ is a subfield of $\mathbf C$ that contains $\delta_0$ and maps to $L$. Call an
embedding $\sigma$ of $L$ \emph{good} if $\sigma(\delta_0)=\delta_0$, with $\delta_0$ read in $L$. For
$\tau\in\mathcal T$, the \emph{label} of $\sigma$ at $\tau$ is the element $\tau'\in\mathcal T$ for
which $\sigma\circ\iota_{\tau'}$ is $\tau$ or $\bar\tau$, where $\iota_{\tau'}\colon E_h\to L$ is the
embedding attached to $\tau'$, so that $\sigma_0\circ\iota_{\tau'}=\tau'$.
\begin{enumerate}
\item $\sigma_0$ is good, and for every $\sigma$ exactly one of $\sigma$ and $\bar\sigma$ is good.
\item Let $\tau\in\mathcal T$, with the label $\tau'$ of $\sigma$ at $\tau$. Then $\sigma$ is good if
  and only if $\sigma\circ\iota_{\tau'}=\tau$.
\item If $\sigma$ agrees with $\sigma_0$ on $K$, then $\sigma$ is good. If not, then $\bar\sigma$ agrees
  with $\sigma_0$ on $K$, and $\sigma$ is not good.
\end{enumerate}
\end{lemma}
\leanref{Formalism/LocalToGlobal/IndefiniteReal/FWExist.lean}{Formalism.IndefR.FWT.base_embedding_good}
\leanref{Formalism/LocalToGlobal/IndefiniteReal/FWExist.lean}{Formalism.IndefR.FWT.good_resC}
\leanref{Development/AbelianVarietyTheory/Counterexamples/TwistSurvey/ConjModelCast.lean}{GlobalTheory.ConjReading.not_good_conj}
\leanref{Development/AbelianVarietyTheory/Counterexamples/TwistSurvey/ConjModelCast.lean}{GlobalTheory.ConjReading.branch_of_pins}
\leanref{Development/AbelianVarietyTheory/Proofs/DivisorProofs/ExistsPsiOfOverK.lean}{GlobalTheory.GalLift.labels_overK}
\begin{proof}
(1) The weight tie makes the map $K_s\to L$ extend the map of rings of integers of the core, and the
pins of the core make its composite with $\sigma_0$ the inclusion of $K_s$ in $\mathbf C$; so
$\sigma_0(\delta_0)=\delta_0$. As $\delta_0^2=-d<0$, $\delta_0$ is not real, and $\bar\delta_0=-\delta_0$.
For every $\sigma$, $\sigma(\delta_0)$ is a square root of $-d$, so it is $\delta_0$ or
$-\delta_0=\bar\delta_0$, and $\bar\sigma(\delta_0)$ is then $\bar\delta_0\ne\delta_0$ or $\delta_0$.

(2) The element $\iota_{\tau'}(r)$ of $L$ is $\delta_0$: by (1), $\sigma_0(\delta_0)=\delta_0=\tau'(r)
=\sigma_0(\iota_{\tau'}(r))$, and $\sigma_0$ is injective. So
$\sigma(\delta_0)=(\sigma\circ\iota_{\tau'})(r)$. If $\sigma\circ\iota_{\tau'}=\tau$, this is
$\tau(r)=\delta_0$. Otherwise $\sigma\circ\iota_{\tau'}=\bar\tau$, and it is
$\overline{\tau(r)}=\bar\delta_0\ne\delta_0$.

(3) Let $\sigma$ agree with $\sigma_0$ on $K$, let $\tau\in\mathcal T$ and let $\tau'$ be the label of
$\sigma$ at $\tau$. Every embedding in $\mathcal T$ carries $\sqrt{-d}$ to $i\sqrt d$, and so does the CM
embedding $\varphi_0=\sigma_0\circ\iota$ of the host, where $\iota\colon E_h\to L$ is the embedding of
$D$. Hence $\sigma_0(\iota_{\tau'}(\sqrt{-d}))=\tau'(\sqrt{-d})=\sigma_0(\iota(\sqrt{-d}))$, so
$\iota_{\tau'}(\sqrt{-d})=\iota(\sqrt{-d})$, and
$\sigma(\iota_{\tau'}(\sqrt{-d}))=\sigma_0(\iota(\sqrt{-d}))=i\sqrt d$. As $\bar\tau(\sqrt{-d})=-i\sqrt d$,
$\sigma\circ\iota_{\tau'}=\tau$, and $\sigma$ is good by (2). If $\sigma$ does not agree with
$\sigma_0$ on $K$, then $\sigma(\iota(\sqrt{-d}))=-i\sqrt d$, as it is a square root of $-d$ different
from $\sigma_0(\iota(\sqrt{-d}))=i\sqrt d$. So $\bar\sigma(\iota(\sqrt{-d}))=i\sqrt d$, and $\bar\sigma$,
which then agrees with $\sigma_0$ on $K=\mathbf Q(\sqrt{-d})$, is good; by (1), $\sigma$ is not good.
\end{proof}

\begin{lemma}\label{is-eigenvalue-conjugate-models}
\begin{enumerate}
\item Let $R$ be a realization of a definite-side datum $D$ that satisfies the core conditions on its
  models, with one-map data that have a weight tie (Definition~\ref{is-one-map-data}). Let $i$ be a
  lift and $\sigma$ an embedding of $L$ that is not good
  (Lemma~\ref{is-eigenvalue-conjugate-good}), and write
  $\bar\sigma$ for its complex conjugate. For every $\tau$ there is a conjugate-linear bijective
  isometry $c_\tau$ from the space of the model of $W_i$ at $\sigma$ at $\tau$ onto that of the model
  at $\bar\sigma$ such that $c_\tau$ carries the action of each $g\in U(G_{W_i}(\tau))$ on the first
  to its action on the second.
\item Let $V$ be a Hermitian space of a definite-side datum, $\theta$ an archimedean family, and $M$,
  $M'$ models of discrete series for $U(V)$ (Definition~\ref{asm-dependence-rows-models}(1)) with
  conjugate-linear bijective isometries $c_\tau$ from the space of $M'$ at $\tau$ onto that of $M$
  that carry the action of each $g\in U(G_V(\tau))$ on the first to its action on the second. Then
  $c_\tau$ and $c_\tau^{-1}$ preserve smooth vectors, and for every $J\subset U(V)(\mathbf A)$, every
  function $\chi$ on $U(V)(\mathbf A)$ and every $T\in\operatorname{Mult}_\theta(M;J,\chi)$
  (Definition~\ref{asm-dependence-rows-models}(2)) there is
  $T'\in\operatorname{Mult}_\theta(M';J,\bar\chi)$ with $T'(v')(g)=\overline{T(v)(g)}$ for all $g$ and
  all $v'$ and $v$ with $v_\tau=c_\tau(v'_\tau)$ for every $\tau$.
\end{enumerate}
\end{lemma}
\leanref{Development/AbelianVarietyTheory/Proofs/DivisorProofs/ExistsContinuousRepresentativeSobolevBoundStatement.lean}{GlobalTheory.T.exists_antiunitary_intertwiner_conjugateEmbedding_models}
\leanref{Development/AbelianVarietyTheory/Proofs/TwistedOperators/ComplexConjugationMapsMultiplicitySpace.lean}{GlobalTheory.ConjMult.conj_mult}
\leanref{Development/AbelianVarietyTheory/Proofs/DivisorProofs/ExistsContinuousRepresentativeSobolevBound.lean}{GlobalTheory.ConjEq.exists_conj_equiv}
\leanref{Development/AbelianVarietyTheory/Proofs/DivisorProofs/ExistsConjModelExt.lean}{GlobalTheory.ConjExt.exists_conjModel_ext}
\leanref{Development/AbelianVarietyTheory/Proofs/DivisorProofs/HolomorphicDiscreteSeriesModelsUnitarilyEquivaleStatement.lean}{GlobalTheory.T.holomorphic_discrete_series_models_unitarily_equivalent}
\leanref{Development/AbelianVarietyTheory/Proofs/DivisorProofs/HolomorphicShapeIntegralWeightGoodEmbedding.lean}{GlobalTheory.HolDict.holShape_of_good}
\begin{proof}
(1) Let $M$ be the model of $W_i$ at $\sigma$. Since $\sigma$ is not good, $\bar\sigma$ is good, by
Lemma~\ref{is-eigenvalue-conjugate-good}(1). So at every $\tau$, with the label $\tau'$ of
$\bar\sigma$ at $\tau$, $\bar\sigma\circ\iota_{\tau'}=\tau$, by
Lemma~\ref{is-eigenvalue-conjugate-good}(2). As in the proof of
Lemma~\ref{is-eigenvalue-conjugate-theta-facts}(6), the model of $W_i$ at $\bar\sigma$ then has
\emph{holomorphic shape}: at every $\tau$ its $\lambda_\tau$ is a Harish-Chandra parameter and its signs
are $1$ at the first $p$ indices and $-1$ at the others, $(p,q)$ the signature of $W_i$; its extremal
weights, the entries of $\Lambda_\tau$ (Definition~\ref{is-eigenvalue-conjugate-k-types}(3)), are
integers, as those of $\lambda_\tau$ are and the rank of $W_i$ is odd. By the weight tie, $v_0(i,\sigma)$ is an extremal vector of $M$: at every $\tau$
its component is a unit vector of the lowest $K$-type of extremal weight.

The complex conjugates $\bar H_\tau$ of the spaces $H_\tau$ of $M$, on which $g\in U(G_{W_i}(\tau))$
acts by its action on $H_\tau$, form a model $M'$ of discrete series for $U(W_i)$, whose parameter at
$\tau$ is the dual $\{(-\lambda,\epsilon):(\lambda,\epsilon)\in S_\tau\}$ of the parameter $S_\tau$ of
$M$, and $M'$ has an extremal vector, the conjugate of a translate of $v_0(i,\sigma)$: the complex
conjugate of a discrete series representation is the discrete series representation with the opposite
Harish-Chandra parameter \cite{HC66}. Since the label of $\sigma$ at $\tau$ carries the embedding of
$E_h$ to $\tau$ exactly when the label of $\bar\sigma$ does not, the parameter of the model of $W_i$ at
$\bar\sigma$ is the dual of that of $M$, which is the parameter of $M'$. Two models of discrete series
for $U(W_i)$ with the same parameters, of holomorphic shape and with extremal vectors, are unitarily
equivalent, since a discrete series representation is determined by its Harish-Chandra parameter
\cite{HC66}: there are unitary maps $U_\tau$ from the spaces of $M'$ onto those of the model at
$\bar\sigma$ that carry the action of each $g\in U(G_{W_i}(\tau))$ on the first to its action on the
second. Let $c_\tau$ be the identity map $H_\tau\to\bar H_\tau$, which is conjugate-linear, isometric
and compatible with the actions, followed by $U_\tau$.

(2) A vector $y$ of the space of $M'$ at $\tau$ is smooth if $t\mapsto g(t)y$ is smooth for every
product $g(t)$ of exponentials $\exp(t_kX_k)$ of elements of the Lie algebra of $U(G_V(\tau))$, and
the smooth vectors are the span of these. As $c_\tau$ is a continuous real-linear bijection with
$c_\tau(gy)=gc_\tau(y)$, it carries smooth vectors of $M'$ to smooth vectors of $M$, and so does
$c_\tau^{-1}$ in the other direction; as $c_\tau$ is conjugate-linear, it preserves spans.

Let $T\in\operatorname{Mult}_\theta(M;J,\chi)$. The multiplicity maps of level $(J,\chi)$ form a
subspace, as their conditions are linear, so $T$ is one. For $v'\in\prod_\tau H'^\infty_\tau$ let
$c(v')=(c_\tau(v'_\tau))_\tau$ and put $T'(v')=\overline{T(c(v'))}$. It is linear in each variable,
since $c_\tau$ and complex conjugation are conjugate-linear. We check the conditions of a multiplicity
map of level $(J,\bar\chi)$ for $T'$.
\begin{itemize}
\item Each $T'(v')$ is a cuspidal automorphic form of the required kind. The invariance under
  $U(V)(F)$ on the left, the invariance under an open subgroup at the finite places, continuity,
  smoothness at the archimedean places and moderate growth pass to the complex conjugate, and the
  constant terms of $\overline{T(c(v'))}$ are the conjugates of those of $T(c(v'))$, which vanish. The
  right translates of $T(c(v'))$ by the archimedean points are again values of $T$, by the
  equivariance of $T$. Complex conjugation commutes with the right derivatives along the real Lie
  algebras, so it carries the action of each generator of the centre of an enveloping algebra on the
  span of these translates to an element of the algebra these generators generate; hence a
  finite-dimensional space containing $T(c(v'))$ and stable under the centres is carried to one
  containing $T'(v')$.
\item For $k\in J$, $T'(v')(gk)=\overline{\chi(k)T(c(v'))(g)}=\bar\chi(k)T'(v')(g)$.
\item If $v'_\tau$ is replaced by its translate under $g\in U(G_V(\tau))$, then $c(v')_\tau$ is
  replaced by its translate under $g$, so $T(c(v'))$, and with it $T'(v')$, by its right translate by
  the image of $g$ under $\theta$.
\item Let $x$ be a point and $\ell$ a functional on the space of $M$ at $\tau$, continuous for a
  Sobolev norm, with $\ell(w)=T(c(v')[\tau\leftarrow w])(x)$ for smooth $w$. Then
  $\ell'(w')=\overline{\ell(c_\tau w')}$ is linear, satisfies
  $\ell'(w')=T'(v'[\tau\leftarrow w'])(x)$, and is continuous for the Sobolev norm of $M'$ given by
  the same elements of the Lie algebra, since $c_\tau$ is isometric and commutes with their derived
  actions.
\end{itemize}
So $T'\in\operatorname{Mult}_\theta(M';J,\bar\chi)$, and $T'(v')(g)=\overline{T(v)(g)}$ when
$v_\tau=c_\tau(v'_\tau)$ for every $\tau$.
\end{proof}

\begin{lemma}\label{is-eigenvalue-conjugate-lowest-types}
Let $D$ be a definite-side datum, $\theta$ an archimedean family and $W$ a Hermitian space of $D$, and
let $M$ be a model of discrete series for $U(W)$ (Definition~\ref{asm-dependence-rows-models}(1)) with
representations $\sigma_\tau$, frames $f_\tau$ and signature matrices $J_\tau$ that has smooth
irreducible lowest $K$-types $\kappa_\tau$ (Definition~\ref{is-eigenvalue-conjugate-k-types}(3)). Let
$K_\tau$ be the group of unitary matrices that commute with $J_\tau$, and for $k\in K_\tau$ let
$k^\theta\in U(W)(\mathbf A)$ be the image of $f_\tau^{-1}kf_\tau$ under $\theta$ at $\tau$. Call a map
$U$ from $\prod_\tau\kappa_\tau$ to functions on $U(W)(\mathbf A)$, linear in each variable,
\emph{$K$-equivariant} if $U(v')(g)=U(v)(gk^\theta)$ for every $\tau$, every $k\in K_\tau$, every
$v\in\prod_\tau\kappa_\tau$ and every $g\in U(W)(\mathbf A)$, where $v'$ is $v$ with $v_\tau$ replaced
by $\sigma_\tau(k)v_\tau$.
\begin{enumerate}
\item For every $J\subset U(W)(\mathbf A)$ and every function $\chi$ on $U(W)(\mathbf A)$, every element
  of $\operatorname{Mult}_\theta(M;J,\chi)$ (Definition~\ref{asm-dependence-rows-models}(2)) is
  $K$-equivariant on $\prod_\tau\kappa_\tau$.
\item If $U$ and $U'$ are $K$-equivariant and $U(e)=U'(e)$ for some $e\in\prod_\tau\kappa_\tau$ with
  $e_\tau\ne0$ for every $\tau$, then $U=U'$.
\end{enumerate}
\end{lemma}
\leanref{Development/AbelianVarietyTheory/Proofs/IwasawaDecomposition/EqTopOfTranslates.lean}{GlobalTheory.DefLift.KEquiv}
\leanref{Development/AbelianVarietyTheory/Proofs/IwasawaDecomposition/EqTopOfTranslates.lean}{GlobalTheory.DefLift.eq_top_of_translates}
\leanref{Development/AbelianVarietyTheory/Proofs/IwasawaDecomposition/EqTopOfTranslates.lean}{GlobalTheory.DefLift.mem_of_translates}
\leanref{Development/AbelianVarietyTheory/Proofs/DivisorProofs/HolShapeOfOverK.lean}{GlobalTheory.GalLift.ProjLift.eq_on_low}
\begin{proof}
(1) Let $k\in K_\tau$ and put $g=f_\tau^{-1}kf_\tau$. As $f_\tau^*J_\tau f_\tau=G_W(\tau)$ and
$k^*J_\tau k=J_\tau$, $g^*G_W(\tau)g=f_\tau^*k^*J_\tau kf_\tau=G_W(\tau)$, so $g\in U(G_W(\tau))$, and
$g$ acts on the space of $M$ at $\tau$ by $\sigma_\tau(f_\tau gf_\tau^{-1})=\sigma_\tau(k)$. The vectors
of $\kappa_\tau$ are smooth, and $\sigma_\tau(k)$ preserves $\kappa_\tau$. So, for a multiplicity map $T$,
replacing $v_\tau$ by $\sigma_\tau(k)v_\tau$ replaces $T(v)$ by its right translate by the image of $g$
under $\theta$, which is $k^\theta$. This property is preserved by sums and scalar multiples, so it
holds for every element of the span $\operatorname{Mult}_\theta(M;J,\chi)$.

(2) Put $B=U-U'$. It is $K$-equivariant, and $B(e)=0$. First let $\tau$ be given and $w\in\prod_\rho
\kappa_\rho$ with $B(w)=0$ and $w_\tau\ne0$, and let $N$ be the set of $u\in\kappa_\tau$ for which $B$
vanishes at $w$ with $w_\tau$ replaced by $u$. It is a subspace, as $B$ is linear in the $\tau$-th
variable, and it contains $\sigma_\tau(k)w_\tau$ for every $k\in K_\tau$, by the $K$-equivariance of $B$.
The span of these vectors is stable under $\sigma_\tau(K_\tau)$, lies in $\kappa_\tau$ and contains
$w_\tau=\sigma_\tau(1)w_\tau\ne0$. As $\kappa_\tau$ is a $K$-type, its only subspaces stable under
$\sigma_\tau(K_\tau)$ are $0$ and $\kappa_\tau$; so the span is $\kappa_\tau$, and $N=\kappa_\tau$.

Now let $v\in\prod_\tau\kappa_\tau$, number the $\tau$ as $\tau_1,\dots,\tau_f$, and let $w^{(j)}$ have
the components of $v$ at $\tau_1,\dots,\tau_j$ and those of $e$ at the others. Then
$B(w^{(0)})=B(e)=0$. If $B(w^{(j-1)})=0$, the first step at $\tau_j$, where $w^{(j-1)}$ has the
component $e_{\tau_j}\ne0$, gives $B(w^{(j)})=0$. Hence $B(v)=B(w^{(f)})=0$.
\end{proof}

\begin{lemma}\label{is-eigenvalue-conjugate-class-sums}
Let $R$ be a realization of a definite-side datum $D$ (Definition~\ref{is-realisation}), with
coefficient field $L$ and base tuple $\mu_0$, and let one-map data for $R$ be given
(Definition~\ref{is-one-map-data}), with a core whose coefficient field is $K_s\subset L$, weight
families $\kappa_j$ and a weight tie. Let $i$ be a lift at which the one-map data are pinned and give
the $i$-th lift as one map with some constants and some phase, and write $j_V$ for the
$\sigma_0$-semilinear map from the lift forms at $(i,\mu_0)$ to $H_V(i,\mu_0)$. Write $L'\supset K_s$
for the coefficient field of the one-map data at $i$ and $\iota\colon L\to L'$ for the isomorphism
that the pins give, and call the cusp forms of weight $\kappa_i$ of the core at level $i+1$ the
\emph{core forms}.
\begin{enumerate}
\item The pins give representatives $h_1,\dots,h_N\in U(V_i)(\mathbf A)$ of the classes of the
  one-map data at $i$ such that, for every lift form $x$ at $(i,\mu_0)$, the values of $j_V(x)$ at the
  finite parts of the $h_k$ are the images under the coefficient map of the one-map data of vectors $x'_1,\dots,x'_N$ of
  the rational coefficient space of $V_i$ over $L'$; and the theta class data of the one-map data are
  the families $x'=(x'_k)$ that arise in this way.
\item Let $x$ be a lift form at $(i,\mu_0)$ and $x'$ as in (1). Then there are a section $f'$ of
  weight $\kappa_i$ of the core at level $i+1$ over $L'$ that is built from $x'$, core forms
  $f_1,\dots,f_n$ and $a_1,\dots,a_n\in L$ such that $f'=\sum_k\iota(a_k)f_k$ after extension of
  scalars from $K_s$ to $L'$. Here $f'$ is \emph{built from} $x'$ if, at every point of the core over an
  algebraically closed field, its value is the value there of the delta section of the one-map data
  times $\sum_k\mathrm{vol}\,\gamma_k^{-1}\varepsilon_k\vartheta_k(x'_k)$, where $\mathrm{vol}$ is the
  volume, $\gamma_k$ the stabilizer order and $\varepsilon_k$ the local constant of the $k$-th class, and
  $\vartheta_k(x'_k)$ is the theta value of the $k$-th lattice for the polynomial that the polynomial
  realization of the coefficients attaches to $x'_k$.
\end{enumerate}
\end{lemma}
\leanref{Development/AbelianVarietyTheory/Proofs/DivisorProofs/GalLiftOfOneMap.lean}{GlobalTheory.GalLift.Assembly.exists_reps}
\leanref{Development/AbelianVarietyTheory/Proofs/DivisorProofs/ReadingsExistRow.lean}{GlobalTheory.GalLift.RepsAt}
\leanref{Development/AbelianVarietyTheory/Proofs/DivisorProofs/ReadingsExistRow.lean}{GlobalTheory.GalLift.IsClassSum}
\leanref{Formalism/LocalToGlobal/ProjectiveLattices/ThClass.lean}{Formalism.L2.Pj.IsThClsD}
\leanref{Formalism/LocalToGlobal/ProjectiveLattices/ThClassPins.lean}{Formalism.L2.Pj.IsBC}
\leanref{Formalism/LocalToGlobal/ProjectiveLattices/ThClassPins.lean}{Formalism.L2.Pj.CuspBC}
\leanref{Development/AbelianVarietyTheory/Proofs/DivisorProofs/OscillatorModelsEvaluateSecond.lean}{GlobalTheory.GalLift.BaseId.tie_baseChange}
\leanref{Development/AbelianVarietyTheory/Proofs/DivisorProofs/ConjugateLiftUniqueAndBaseDatum.lean}{GlobalTheory.GalLift.BaseId.base_identity_five}
\begin{proof}
(1) This is the condition on representatives in the pins of the one-map data at $i$, at the base tuple.

(2) By (1), $x'$ is a theta class datum. The one-map data give the $i$-th lift as one map, and the
second condition of this, the one-map tie, gives two facts. First, for every theta class datum $x'$
there is a cusp form $f'$ of weight $\kappa_i$ of the core at level $i+1$ over $L'$ that is built from
$x'$; it is in particular a section of that weight. Second, every such cusp form over $L'$ comes from
the core forms by extension of scalars: there are core forms $f_1,\dots,f_n$ and $b_1,\dots,b_n\in L'$
with $f'=\sum_kb_kf_k$ after extension of scalars from $K_s$ to $L'$. The pins identify $L'$ with the
image of $L$, through $\iota$; so $b_k=\iota(a_k)$ for some $a_k\in L$. The one-map tie states these
facts for forms whose length is the rank of $V_i$, which the weight tie makes equal to the length of
$\kappa_i$; so the $f_k$ are core forms.
\end{proof}

\subsection{Dictionaries, multiplicity maps and phases}\label{subsection-is-dictionaries-phases}

\begin{definition}\label{is-eigenvalue-conjugate-dictionary}
Let $D$ be a definite-side datum with host field $E_h/F_h$, element $\xi_0$ of the host and set
$\mathcal T$ of embeddings of $E_h$ (Definition~\ref{th-definite-datum}). Let $V$ be a Hermitian space
of $D$ of rank $n$ with Gram matrix $G_V$, let $\theta$ be an archimedean family whose member
$\theta_\rho$ extends $\rho$ for every $\rho\in\mathcal T$, and let $M$ be a model of discrete series
for $U(V)$ (Definition~\ref{asm-dependence-rows-models}(1)) with frames $f_\rho$, signs $\epsilon_\rho$
and representations $\pi_\rho$; let $q_\rho$ be the number of signs $-1$ in $\epsilon_\rho$. For
$g\in U(V)(\mathbf A)$ write $g_\rho\in U(G_V(\rho))$ for the image of $g$ under $\theta_\rho$, put
$U_\infty=\prod_\rho U(G_V(\rho))$, and let $\mathcal B\subset GL_n(\mathbf C)$ be the group of
invertible upper triangular matrices.
\begin{enumerate}
\item (Frames.) The \emph{frame} $x_{0,\rho}$ of $M$ at $\rho$ is the matrix whose $a$-th column is
  the $(n+1-a)$-th column of $f_\rho^{-1}$ if $a\le q_\rho$, and its $(a-q_\rho)$-th column if
  $a>q_\rho$. The \emph{frame domain} $\mathcal D_M$ is the set of $x\in M_n(\mathbf C)^{\mathcal T}$
  with $\det x_\rho\ne0$ and $(-1)^{\max(0,i-q_\rho)}\operatorname{Re}d_i(x,\rho)>0$ for all $\rho$ and
  $0\le i\le n$, where $d_i(x,\rho)$ is the leading principal $i\times i$ minor of
  $-x_\rho^*G_V(\rho)x_\rho$.
\item (Holomorphic forms.) Let $w=(w_{\rho,a})$ be integers. A function $f$ on $U_\infty$ is
  \emph{holomorphic of weight $w$} if there is a function $P$ on $M_n(\mathbf C)^{\mathcal T}$,
  holomorphic on $\mathcal D_M$, with $P(xb)=\prod_{\rho,a}(b_\rho)_{aa}^{w_{\rho,a}}P(x)$ for
  $x\in\mathcal D_M$ and $b\in\mathcal B^{\mathcal T}$, such that $f(u)=P\big((u_\rho
  x_{0,\rho})_\rho\big)$ for $u\in U_\infty$. A function $\phi$ on $U(V)(\mathbf A)$ is
  \emph{holomorphic of weight $w$} if for every $g$ some function $f$ on $U_\infty$, holomorphic of
  weight $w$, satisfies $\phi(gu)=f((u_\rho)_\rho)$ for every $u\in U(V)(\mathbf A)$ with finite
  component $1$. For a subset $J$ and a function $\chi$ of $U(V)(\mathbf A)$, a \emph{bounded
  holomorphic automorphic form} of weight $w$ and level $(J,\chi)$ is a bounded continuous function
  $\phi$ on $U(V)(\mathbf A)$, holomorphic of weight $w$, invariant on the left under $U(V)(F_h)$ and
  on the right under the elements with archimedean component $1$ of some open subgroup, such that
  $\phi(gk)=\chi(k)\phi(g)$ for every $g$ and every $k\in J$.
\item (Constant terms.) For $\xi\in E_h^n$ with $\xi\ne0$ and $\xi^*G_V\xi=0$ let $Z_\xi$ be the
  subgroup of $U(V)(\mathbf A)$ generated by the elements $g$ with $g-1=a\,\xi\xi^*G_V$ for some
  $a\in\mathbf A_{F_h}\otimes_{F_h}E_h$. The \emph{constant terms} of a function $\phi$ on
  $U(V)(\mathbf A)$ \emph{vanish} if $\int_{\mathcal F}\phi(zg)\,d\mu(z)=0$ for every such $\xi$, every
  $g$, every Haar measure $\mu$ on $Z_\xi$ and every fundamental domain $\mathcal F$ of
  $Z_\xi\cap U(V)(F_h)$ in $Z_\xi$.
\item (Dictionaries.) A \emph{dictionary} for $V$ consists of a set $C$; for each $c\in C$ flag data
  $X_c$ of rank $n$ over $E_h$ (Definition~\ref{is-onto-share-flag-data}), with set $\Psi_c$, group
  $\Gamma_c$, characters $\chi_{c,\kappa}$ and flag domain $\mathcal D_c$, a weight $\kappa_c$ of $X_c$
  and a constant $c_c\in\mathbf C$; a map $g\mapsto(c(g),x(g))$ from $U(V)(\mathbf A)$ to the pairs of
  some $c\in C$ and some $x\in M_n(\mathbf C)^{\Psi_c}$; and a function $\mu$ on $U(V)(\mathbf A)$, the
  \emph{multiplier}. The \emph{adelization} of a family $P=(P_c)$ of functions $P_c$ on
  $M_n(\mathbf C)^{\Psi_c}$ is the function
  $\operatorname{ad}(P)(g)=c_{c(g)}\,\mu(g)\,P_{c(g)}(x(g))$ on $U(V)(\mathbf A)$.
\item (Good dictionaries.) Let $w=(w_{\rho,a})$ be integers and $K$ a subgroup of $U(V)(\mathbf A)$.
  A dictionary for $V$ is \emph{good} for $(M,w,K)$ if it satisfies the following conditions.
  \begin{enumerate}
  \item (Shape.) $C$ is finite; for every $c$, $\Psi_c=\mathcal T$, the integer $p_0$ of $X_c$ is
    $q_\rho$ for every $\rho$, the matrix of $X_c$ is $-G_V^{\mathsf T}$, its element $\xi_0$ is
    $-\xi_0$, $\chi_{c,\kappa_c}$ is trivial on $\Gamma_c$ and $c_c\ne0$; and the numbers $|c_c|$ are
    bounded.
  \item (Base frame.) For every $h\in U(V)(\mathbf A)$ there is $\gamma\in U(V)(F_h)$ with
    $x(h)_\rho=(\gamma h)_\rho x_{0,\rho}$ for every $\rho$ and such that, for every $u$ with finite
    component $1$, $(c(hu),x(hu))=(c(h),\gamma'x)$ for some $\gamma'\in\Gamma_{c(h)}$, where
    $x_\rho=(\gamma hu)_\rho x_{0,\rho}$.
  \item (Uniformization.) For every $g$, every $\gamma\in U(V)(F_h)$ and every $k\in K$, the pairs
    $(c(\gamma g),x(\gamma g))$ and $(c(gk),x(gk))$ are $(c(g),\gamma'x(g))$ for some
    $\gamma'\in\Gamma_{c(g)}$; $x(g)\in\mathcal D_{c(g)}$; if $(c(g'),x(g'))=(c(g),\gamma'x(g))$ with
    $\gamma'\in\Gamma_{c(g)}$, then $g'\in U(V)(F_h)gK$; every $c$ and every $x\in\mathcal D_c$ have some
    $g$, $\gamma'\in\Gamma_c$ and $b\in\mathcal B^{\Psi_c}$ with $(c(g),x(g))=(c,\gamma'xb)$; and there
    are points $x_c=x(g_c)\in\mathcal D_c$ with $c(g_c)=c$ and, for every $c$ and $\tau\in\Psi_c$, a
    ring homomorphism $\theta_{c,\tau}$ from $\mathbf A_{F_h,\infty}\otimes_{F_h}E_h$ to $\mathbf C$
    extending $\tau$, such that $(c(gh),x(gh))=(c(g),\gamma'x(g)b)$ for some $\gamma'\in\Gamma_{c(g)}$
    whenever $h$ has finite component $1$, where $b_\tau=x_{c(g),\tau}^{-1}h_{c(g),\tau}x_{c(g),\tau}$
    and $h_{c,\tau}$ is the image of $h$ under $\theta_{c,\tau}$.
  \item (Multiplier.) $\mu$ is a continuous homomorphism with $|\mu|=1$ that is trivial on $U(V)(F_h)$
    and on $K$; and $|d_i(x(g),\tau)|=1$ for all $g$, $\tau$ and $i$, with the minors $d_i$ of
    Definition~\ref{is-onto-share-flag-data} for $X_{c(g)}$.
  \item (Level.) The elements of $K$ have archimedean component $1$, and $K$ contains the elements
    with archimedean component $1$ of some open subgroup of $U(V)(\mathbf A)$.
  \item (Weight.) There are integers $e_\rho$ with $\mu(u)=\prod_\rho\det(u_\rho)^{e_\rho}$ for every
    $u$ with finite component $1$, and $w_{\rho,a}=\kappa_{c,\rho,a}+e_\rho$ for all $c$, $\rho$ and
    $a$.
  \end{enumerate}
\end{enumerate}
\end{definition}
\leanref{Formalism/LocalToGlobal/IndefiniteReal/ExistsChosenExtremalVector.lean}{Formalism.IndefR.FWT.x0Mat}
\leanref{Development/AbelianVarietyTheory/Definitions/Divisors/IsSmoothVector.lean}{GlobalTheory.HolDict.frDom}
\leanref{Development/AbelianVarietyTheory/Definitions/Divisors/IsSmoothVector.lean}{GlobalTheory.HolDict.HolFn}
\leanref{Development/AbelianVarietyTheory/Definitions/Divisors/IsSmoothVector.lean}{GlobalTheory.HolDict.HolAd}
\leanref{Development/AbelianVarietyTheory/Definitions/Divisors/IsSmoothVector.lean}{GlobalTheory.HolDict.HolForm}
\leanref{Development/AbelianVarietyTheory/Definitions/Divisors/IsSmoothVector.lean}{GlobalTheory.HolDict.ZI}
\leanref{Development/AbelianVarietyTheory/Definitions/Divisors/IsSmoothVector.lean}{GlobalTheory.HolDict.ZCusp}
\leanref{Formalism/LocalToGlobal/ProjectiveLattices/EquivariantFiniteGroup.lean}{Formalism.L2.Pj.adel}
\leanref{Formalism/LocalToGlobal/ProjectiveLattices/FamilyUniformisationBasePointOrthogonalBis.lean}{Formalism.L2.Pj.muPin}
\leanref{Development/AbelianVarietyTheory/Definitions/Divisors/IsSmoothVector.lean}{GlobalTheory.HolDict.GoodDict}
\leanref{Formalism/LocalToGlobal/IndefiniteReal/ExistsChosenExtremalVector.lean}{Formalism.IndefR.FWT.BaseFr}
\leanref{Formalism/LocalToGlobal/ProjectiveLattices/EquivariantFiniteGroup.lean}{Formalism.L2.Pj.DictIsUnif}
\leanref{Formalism/LocalToGlobal/ProjectiveLattices/EquivariantFiniteGroup.lean}{Formalism.L2.Pj.MuUnit}
\leanref{Formalism/LocalToGlobal/ProjectiveLattices/FamilyUniformisationBasePointOrthogonalBis.lean}{Formalism.L2.Pj.ModelW}

\begin{lemma}\label{is-eigenvalue-conjugate-dictionary-facts}
Let $D$, $V$, $G_V$, $\theta$, $M$, $q_\rho$, $U_\infty$ and $\mathcal B$ be as in
Definition~\ref{is-eigenvalue-conjugate-dictionary}.
\begin{enumerate}
\item (Archimedean lift.) There is a continuous homomorphism $L\colon U_\infty\to U(V)(\mathbf A)$ such
  that $L(u)$ has finite component $1$ and $L(u)_\rho=u_\rho$ for every $\rho$, and $L((g_\rho)_\rho)=g$
  for every $g$ with finite component $1$. Every $g\in U(V)(\mathbf A)$ is $au$ with $a$ of
  archimedean component $1$ and $u$ of finite component $1$, and such $a$ and $u$ commute.
\item (Frames.) Suppose that $M$ has holomorphic shape (Lemma~\ref{is-eigenvalue-conjugate-models}(1)).
  Then $x_{0,\rho}^*G_V(\rho)x_{0,\rho}$ is the diagonal matrix with entries $-1$ at the first $q_\rho$
  indices and $1$ at the others, and $\mathcal D_M$ is open, contains $(x_{0,\rho})_\rho$ and contains
  $(u_\rho x_\rho b_\rho)_\rho$ for all $x\in\mathcal D_M$, $u\in U_\infty$ and $b\in\mathcal
  B^{\mathcal T}$. For every component $X_c$ of a dictionary for $V$ that is good for some $(M,w,K)$,
  $\mathcal D_c=\mathcal D_M$, and $\rho(\gamma)\in U(G_V(\rho))$ for every $\gamma\in\Gamma_c$ and every
  $\rho$.
\item (Transvections.) Let $\xi$ and $Z_\xi$ be as in
  Definition~\ref{is-eigenvalue-conjugate-dictionary}(3). Then $Z_\xi$ is closed, there is an
  isomorphism of topological groups from $\mathbf A_{F_h}$ onto $Z_\xi$ that carries $F_h$ onto
  $Z_\xi\cap U(V)(F_h)$, and every element of $Z_\xi$ fixes $\xi$. For $z\in Z_\xi$ and $\rho$,
  $z_\rho=1+is\,\zeta_\rho\zeta_\rho^*G_V(\rho)$ for some $s\in\mathbf R$, where $\zeta_\rho=\rho(\xi)$,
  and every family of matrices of this form is $(z_\rho)_\rho$ for some $z\in Z_\xi$ with finite
  component $1$. For every Haar measure on $Z_\xi$, every fundamental domain of $Z_\xi\cap U(V)(F_h)$ in
  $Z_\xi$ has finite positive measure. For every open subgroup $U$ of $U(V)(\mathbf A)$, every
  $z\in Z_\xi$ is $\gamma uk$ with $\gamma,u,k\in Z_\xi$, $\gamma$ rational, $u$ of finite component
  $1$ and $k\in U$ of archimedean component $1$.
\item (Flows.) Let $X$ be flag data with matrix $A$ (Definition~\ref{is-onto-share-flag-data}), let
  $\xi\in\Lambda$ be isotropic, and let $\phi_\xi(s)$ act by the matrices $f_\tau(s)$. For
  $s,s'\in\mathbf R^\Psi$, $\phi_\xi(s)\phi_\xi(s')=\phi_\xi(s+s')$, $\det f_\tau(s)=1$, and
  $\phi_\xi(s)$ carries $\mathcal D_X$ into itself. If $A=-G^{\mathsf T}$ and $\tau(\xi)=N\zeta$ with
  $N\in\mathbf N$, then $f_\tau(s)=1+ic_\tau N^2s\,\zeta\zeta^*G^\tau$. For every weight $\kappa$, the
  group $\Pi_\kappa$ has a measurable fundamental domain of finite positive volume.
\end{enumerate}
\end{lemma}
\leanref{Development/AbelianVarietyTheory/Proofs/DivisorProofs/ExistsArchimedeanLiftUnitaryGroupProduct.lean}{GlobalTheory.RealLift.exists_real_lift}
\leanref{Development/AbelianVarietyTheory/Proofs/DivisorProofs/DiscreteSeriesModelUnitaryIrreducibleOrthonormalStatement.lean}{GlobalTheory.T.discreteSeriesModel_unitary_irreducible_orthonormal}
\leanref{Development/AbelianVarietyTheory/Proofs/DivisorProofs/NegativeFrameDomainProductOpenPreconnectedStable.lean}{GlobalTheory.HolDict.frDoms_open_preconnected_stable}
\leanref{Development/AbelianVarietyTheory/Proofs/DivisorProofs/FlagDomainFramesEqGramFramesUnitaryStatement.lean}{GlobalTheory.T.flagDomain_frames_eq_gramFrames_unitary}
\leanref{Development/AbelianVarietyTheory/Proofs/DivisorProofs/AdelicTransvectionSubgroupEqAdelesClosed.lean}{GlobalTheory.HolDict.TransvGroup.transvection_group}
\leanref{Development/AbelianVarietyTheory/Proofs/DivisorProofs/TransvectionSubgroupCocompactHaarFundamentalDomaStatement.lean}{GlobalTheory.T.transvectionSubgroup_cocompact_haar_fundamentalDomain}
\leanref{ReciprocityTheory/ReductiveGroups/ArithmeticGeometry/FundamentalDomains/EmbeddingHermitianPairs.lean}{ReciprocityTheory.ChartRestriction.G.FD.EmbeddingHermitianPairs.flow_flow}
\leanref{ReciprocityTheory/ReductiveGroups/ArithmeticGeometry/FundamentalDomains/EmbeddingHermitianPairs.lean}{ReciprocityTheory.ChartRestriction.G.FD.EmbeddingHermitianPairs.det_flowMat}
\leanref{ReciprocityTheory/ReductiveGroups/ArithmeticGeometry/FundamentalDomains/TransvectionPeriodicity.lean}{ReciprocityTheory.ChartRestriction.G.FD.TransvectionPeriodicity.flow_real_mem_dom}
\leanref{Development/AbelianVarietyTheory/Proofs/DivisorProofs/AdelicCuspFormTransvectionConstantTermsZero.lean}{GlobalTheory.HolDict.TransvCusp.flowMat_eq_transv}
\leanref{ReciprocityTheory/ReductiveGroups/ArithmeticGeometry/FundamentalDomains/PeriodLatticeCountable.lean}{ReciprocityTheory.ChartRestriction.G.FD.PeriodLatticeCountable.fd_exists}

\begin{lemma}\label{is-eigenvalue-conjugate-spectral-facts}
Let $D$, $V$, $\theta$, $M$, $\pi_\rho$, $q_\rho$, $U_\infty$ and $x_{0,\rho}$ be as in
Definition~\ref{is-eigenvalue-conjugate-dictionary}, and let $H_\rho$ be the space of $\pi_\rho$.
\begin{enumerate}
\item (Finite volume \cite{BHC62}.) There are a right-invariant Haar measure $\nu$ on $U(V)(\mathbf A)$
  and a fundamental domain $\mathcal F$ of $U(V)(F_h)$ with $\nu(\mathcal F)<\infty$.
\item (The regular representation.) For such $\nu$ and $\mathcal F$, let $\mathcal H$ be the space of
  square-integrable functions on $U(V)(F_h)\backslash U(V)(\mathbf A)$ for the measure that $\nu$ induces
  through $\mathcal F$, and say that a function $\psi$ on $U(V)(\mathbf A)$ \emph{represents}
  $F\in\mathcal H$ if they agree almost everywhere. There is a strongly continuous representation $R$ of
  $U(V)(\mathbf A)$ on $\mathcal H$ by isometries with $(R(g)F)(x)=F(xg)$ almost everywhere. Every
  bounded continuous function invariant on the left under $U(V)(F_h)$ represents an element of
  $\mathcal H$; two continuous functions that represent the same element are equal; a function
  represents at most one element; if $\psi$ represents $F$, then $x\mapsto\psi(xg)$ represents $R(g)F$;
  representing is linear; and if $\psi,\psi'$ represent $F,F'$, then $\bar\psi\psi'$ is integrable on
  $\mathcal F$ with integral $\langle F,F'\rangle$.
\item (Holomorphic functions on $U_\infty$.) Suppose that $M$ has holomorphic shape, let $w$ be
  integers, and for $s\in\mathbf R^{\mathcal T\times\{1,\dots,n\}}$ put
  $t(s)_\rho=x_{0,\rho}\operatorname{diag}(e^{is_{\rho,a}})_ax_{0,\rho}^{-1}$ and
  $\chi_w(s)=\prod_{\rho,a}e^{is_{\rho,a}w_{\rho,a}}$.
  \begin{enumerate}
  \item Every function $f$ on $U_\infty$ that is holomorphic of weight $w$ is smooth, and
    $f(ut(s))=\chi_w(s)f(u)$ for all $u$ and $s$.
  \item If $f$ is holomorphic of weight $w$, $f(t(s)u)=\chi_w(s)f(u)$ for all $u$ and $s$, and
    $f(1)=0$, then $f=0$.
  \item Let $(f_o)_{o\in\Omega}$ be functions on $U_\infty$, holomorphic of weight $w$, indexed by a
    finite measure space, with $|f_o(u)|\le C$ for all $o$ and $u$ and $o\mapsto f_o(u)$ measurable
    for every $u$. Then $u\mapsto\int f_o(u)\,do$ is holomorphic of weight $w$.
  \item Let $v$ be an extremal vector of $M$, a tensor of unit vectors $v_\rho$ of the lowest $K$-types
    of extremal weight, and let $w_{\rho,a}$ be the extremal weights of $M$ in the order of the columns
    of $x_{0,\rho}$. Then $u\mapsto\prod_\rho\langle v_\rho,\pi_\rho(u_\rho)v_\rho\rangle$ is
    holomorphic of weight $w$, $\pi_\rho(t(s)_\rho)v_\rho=\prod_ae^{is_{\rho,a}w_{\rho,a}}v_\rho$, each
    $v_\rho$ is smooth, and for $a\le c$ the complexified derived action of
    $x_{0,\rho}E_{ac}x_{0,\rho}^{-1}$ kills $v_\rho$ if $a<c$ and multiplies it by $w_{\rho,a}$ if
    $a=c$; here $E_{ac}$ is the matrix unit.
  \end{enumerate}
\item (Representations.)
  \begin{enumerate}
  \item Let $\sigma_\rho$ be irreducible unitary representations of $U(G_V(\rho))$ on Hilbert spaces
    $H'_\rho$, $v_\rho\in H'_\rho$ unit vectors, $R'$ a representation of $U_\infty$ by isometries on a
    Hilbert space $\mathcal K$, and $k\in\mathcal K$ with $\langle k,R'(u)k\rangle=\|k\|^2\prod_\rho
    \langle v_\rho,\sigma_\rho(u_\rho)v_\rho\rangle$ for every $u$. Then there is a continuous
    multilinear map $\Phi\colon\prod_\rho H'_\rho\to\mathcal K$ with $\Phi(v)=k$, with values in the
    closed span of the $R'(u)k$, that carries the replacement of $y_\rho$ by $\sigma_\rho(g)y_\rho$ to
    $R'$ of the element of $U_\infty$ with component $g$ at $\rho$ and $1$ elsewhere.
  \item For $\nu$, $\mathcal F$, $\mathcal H$ and $R$ as in (2), if a continuous multilinear map
    $\Phi\colon\prod_\rho H_\rho\to\mathcal H$ carries the replacement of $y_\rho$ by $\pi_\rho(g)y_\rho$
    to $R$ of the image of $g$ under $\theta$ at $\rho$, then $\Phi$ carries tensors of smooth vectors to
    smooth vectors for the archimedean points, and the derivative of $R$ along an element $X$ of the Lie
    algebra at $\rho$ applied to $\Phi(y)$ is $\Phi$ of $y$ with $y_\rho$ replaced by $d\pi_\rho(X)y_\rho$.
  \item A smooth vector $F\in\mathcal H$ fixed by the elements with archimedean component $1$ of an
    open subgroup $U$ is represented by a continuous function $\psi$, invariant on the left under
    $U(V)(F_h)$ and on the right under those elements, smooth along products of one-parameter subgroups
    at the archimedean places, whose iterated right derivatives along Lie algebra elements are continuous
    and represent the iterated derivatives of $R$ applied to $F$.
  \item (Sobolev bound.) For an open subgroup $U$ there are $N$ and finitely many words $l$ in Lie
    algebra elements such that for every $x$ there is $C$ with
    $|\psi(xu)|\le C\|u\|^N\sum_l\|l\cdot F\|$ for every $u$ with finite component $1$, every smooth
    vector $F$ fixed as in (c) and every continuous $\psi$ representing $F$; here $\|u\|$ is a fixed
    norm of the archimedean components and $l\cdot F$ the iterated derivative.
  \item For every $\rho$, the smooth vectors of $\pi_\rho$ form a subspace stable under $\pi_\rho$ and
    under the derived actions of the Lie algebra and of its complexification, and the derivatives of
    $t\mapsto\ell(\pi_\rho(\exp tX)y)$ at $0$ are $\ell(d\pi_\rho(X)y)$ for smooth $y$ and every linear
    functional $\ell$ continuous for a Sobolev norm.
  \item (Infinitesimal characters \cite{HC53}.) Let $\sigma$ be an irreducible unitary representation
    of $U(G_V(\rho))$ with a nonzero smooth vector on which the complexified derived actions of the
    $x_{0,\rho}E_{ac}x_{0,\rho}^{-1}$, $a\le c$, act as in (3)(d) for some integers $w_a$. Then every
    Gelfand invariant of the Lie algebra acts on the smooth vectors of $\sigma$ by a scalar.
  \item A continuous function on $U(V)(\mathbf A)$, invariant on the left under $U(V)(F_h)$ and on the
    right under the elements with archimedean component $1$ of an open subgroup, whose constant terms
    vanish, is cuspidal: its constant terms along the unipotent radicals of all parabolic subgroups
    vanish.
  \item Let $\phi$ be a function as in (g) that represents $F_0\in\mathcal H$. Every continuous function,
    invariant on the left under $U(V)(F_h)$ and on the right under the elements with archimedean
    component $1$ of an open subgroup, that represents an element of the closed span of the $R(u)F_0$,
    $u\in U(V)(\mathbf A)$, has vanishing constant terms.
  \end{enumerate}
\end{enumerate}
\end{lemma}
\leanref{Development/AbelianVarietyTheory/Proofs/DivisorProofs/ExistsFinPair.lean}{GlobalTheory.HolDict.exists_finPair}
\leanref{Development/AbelianVarietyTheory/Proofs/IwasawaDecomposition/FiniteVolumeAdelicUnitaryQuotient.lean}{GlobalTheory.Mem.VolFin.finVol}
\leanref{Development/AbelianVarietyTheory/Proofs/DivisorProofs/RegularRepresentationOfFinitePairStatement.lean}{GlobalTheory.T.regular_representation_of_finite_pair}
\leanref{Development/AbelianVarietyTheory/Proofs/DivisorProofs/BoundedAutomorphicHasRepresentativeStatement.lean}{GlobalTheory.T.bounded_automorphic_has_representative}
\leanref{Development/AbelianVarietyTheory/Proofs/DivisorProofs/SmFnAndKilledFnOfHolFn.lean}{GlobalTheory.HolDict.smFn_and_killedFn_of_holFn}
\leanref{Development/AbelianVarietyTheory/Proofs/DivisorProofs/TorusLawOfKilled.lean}{GlobalTheory.HolDict.torus_law_of_killed}
\leanref{Development/AbelianVarietyTheory/Proofs/DivisorProofs/HolomorphicTypeEqZeroOfTorusLawAtOne.lean}{GlobalTheory.HolDict.Uniq.eq_zero_of_torus_law}
\leanref{Development/AbelianVarietyTheory/Proofs/DivisorProofs/IntegralHolomorphicTypeFamilyHolomorphicType.lean}{GlobalTheory.HolDict.holFn_integral}
\leanref{Development/AbelianVarietyTheory/Proofs/DivisorProofs/ExtremalVectorMatrixCoeffHolomorphicType.lean}{GlobalTheory.HolDict.holFn_coef_extVec}
\leanref{Development/AbelianVarietyTheory/Proofs/DivisorProofs/BaseFrameBorelActsExtremalVectorByWeightStatement.lean}{GlobalTheory.T.baseFrameBorel_acts_extremalVector_by_weight}
\leanref{Development/AbelianVarietyTheory/Proofs/DivisorProofs/ExistsEquivariantMultilinearOfMatrixCoeffEqUnitaStatement.lean}{GlobalTheory.T.exists_equivariant_multilinear_of_matrixCoeff_eq_unitaryProduct}
\leanref{Development/AbelianVarietyTheory/Proofs/DivisorProofs/EquivariantMultilinearMapsSmoothVectorsIntertwinStatement.lean}{GlobalTheory.T.equivariantMultilinear_maps_smoothVectors_intertwines}
\leanref{Development/AbelianVarietyTheory/Proofs/DivisorProofs/SmoothVectorAutomorphicDerivativesStatement.lean}{GlobalTheory.T.smooth_vector_automorphic_derivatives}
\leanref{Development/AbelianVarietyTheory/Proofs/DivisorProofs/SobolevBoundForSmoothVectorsStatement.lean}{GlobalTheory.T.sobolev_bound_for_smooth_vectors}
\leanref{Development/AbelianVarietyTheory/Proofs/DivisorProofs/SmoothVectorsSubspaceStableSobolevFormDerivStatement.lean}{GlobalTheory.T.smoothVectors_subspace_stable_sobolevForm_deriv}
\leanref{Development/AbelianVarietyTheory/Proofs/DivisorProofs/GelfandInvariantsActByScalarsHighestWeightStatement.lean}{GlobalTheory.T.gelfand_invariants_act_by_scalars_highest_weight}
\leanref{Development/AbelianVarietyTheory/Proofs/DivisorProofs/IsCuspFunOfTransvectionConstantTermsZeroStatement.lean}{GlobalTheory.T.isCuspFun_of_transvection_constantTerms_zero}
\leanref{Development/AbelianVarietyTheory/Proofs/DivisorProofs/CuspidalityPreservedInClosureStatement.lean}{GlobalTheory.T.cuspidality_preserved_in_closure}

\begin{lemma}\label{is-eigenvalue-conjugate-dictionary-good}
Let $R$ be a realization of a definite-side datum $D$ that satisfies the core conditions on its level
characters and on its models, with one-map data that have a weight tie and are pinned at some lift
(Definition~\ref{is-one-map-data}), and with the notation of Definition~\ref{is-eigenvalue-lift-data}.
Let $i$ be a lift, $\sigma$ a good embedding of $L$ (Lemma~\ref{is-eigenvalue-conjugate-good}) with
restriction $\sigma_s$ to $K_s$, and $M$ the model of $W_i$ at $\sigma$; let $w$ be integers such that
the $w_{\rho,a}$ are the extremal weights of $M$ in the order of the columns of the frames of $M$. The
weight tie identifies the space of the core at level $i+1$ and $\sigma_s$ with $W_i$; through this
identification the dictionary of the core there, with the weights of $\kappa_i$ on its components and
its constants, is a dictionary for $W_i$ (Definition~\ref{is-eigenvalue-conjugate-dictionary}(4)). Let
$K_f$ be the level group of the core there, and $K_W$ and $\chi_K$ the level and the level character of
the trivial level datum of $W_i$.
\begin{enumerate}
\item $K_W=K_f$; $\chi_K$ is trivial on $K_W$; the elements of $K_W$ have archimedean component $1$;
  and $K_W$ contains every element with archimedean component $1$ of the open subgroup of the $k$ with
  $(k-1)\hat\Lambda\subset g_0g_1\hat\Lambda$, where $\hat\Lambda$ is the standard lattice of $W_i$
  over the finite adeles.
\item The multiplier of the dictionary is $1$ on $K_W$, and there are integers $e_\rho$ with
  $\mu(u)=\prod_\rho\det(u_\rho)^{e_\rho}$ for every $u$ with finite component $1$ and
  $w_{\rho,a}=\kappa_{c,\rho,a}+e_\rho$ for every component $c$ and all $\rho$ and $a$.
\item $\Psi_c=\mathcal T$ and the integer $p_0$ of $X_c$ is $q_\rho$ for every component $c$ and every
  $\rho$, and the dictionary satisfies the base frame condition of
  Definition~\ref{is-eigenvalue-conjugate-dictionary}(5)(b) for $M$.
\item The dictionary is good for $(M,w,K_W)$.
\end{enumerate}
\end{lemma}
\leanref{Development/AbelianVarietyTheory/Proofs/DivisorProofs/CoreDictionaryGoodAtGoodEmbeddingStatement.lean}{GlobalTheory.T.coreDictionary_good_at_good_embedding}
\leanref{Development/AbelianVarietyTheory/Proofs/DivisorProofs/RealisationPrincipalLevelEqCoreLevelOpenStatement.lean}{GlobalTheory.T.realisation_principalLevel_eq_coreLevel_open}
\leanref{Development/AbelianVarietyTheory/Proofs/DivisorProofs/LevelCharacterEmbeddingCompatStatement.lean}{GlobalTheory.T.level_character_embedding_compat}
\leanref{Development/AbelianVarietyTheory/Proofs/DivisorProofs/DictionaryCharacterTrivialOnPrincipalLevelStatement.lean}{GlobalTheory.T.dictionary_character_trivial_on_principal_level}
\leanref{Development/AbelianVarietyTheory/Proofs/DivisorProofs/DictionaryCharEqOneOnPrincipalLevelStatement.lean}{GlobalTheory.T.dictionaryChar_eq_one_on_principalLevel}
\leanref{Development/AbelianVarietyTheory/Proofs/DivisorProofs/ExtremalWeightEqCharacterExponentAddStatement.lean}{GlobalTheory.T.extremalWeight_eq_character_exponent_add}
\leanref{Development/AbelianVarietyTheory/Proofs/DivisorProofs/BaseFrameOfRealizationModelsStatement.lean}{GlobalTheory.T.base_frame_of_realization_models}
\begin{proof}
The pins of the one-map data make the primes of the core those of the realization. With this, (1)
to (3) are computations from the core conditions on the level characters and on the models and from
the laws of the weight tie, which this proof uses at the level $i+1$ and the good embedding $\sigma_s$.

(1) An element of $U(W_i)(\mathbf A)$ lies in $K_W$ exactly when it is a level lift of $1$ for the
trivial level datum, and comparing the lattice and the reductions at the places of $P_g$ of the two
descriptions shows that this holds exactly when it lies in $K_f$. The character of the trivial level
datum is the extension by $1$ of a trivial character, so $\chi_K$ is trivial, and the elements of
$K_W$ have archimedean component $1$ by the definition of the level. The weight tie requires $K_f$ to
consist of the elements with archimedean component $1$ and $(k-1)\hat\Lambda\subset
g_0g_1\hat\Lambda$; the set of $k$ with the second property is an open subgroup, as $g_0g_1>0$, and its
elements with archimedean component $1$ lie in $K_f=K_W$.

(2) The weight tie requires the multiplier to be $1$ on $K_f$, which is $K_W$ by (1). It also
requires the multiplier to be the value of a character of the ideles at a root of the determinant; at
an element with finite component $1$ this is $\prod_\rho\det(u_\rho)^{e_\rho}$ with $-e_\rho$ the
integer $e$ of the Harish-Chandra parameter at the label of $\sigma$ at $\rho$, and the weights of
$\kappa_i$ on the components are the extremal weights shifted by $-e_\rho$.

(3) The base frame comes from the uniformization law of the dictionary, the pins, the law that
the space of the core is $W_i$ with its matrix and element $\xi_0$, the law that the base points of the
uniformization are orthonormal, and the law that adapts the base classes to the level.

(4) The member of $\theta$ at each $\rho$ extends $\rho$ by the core conditions on the models. We check
the conditions of Definition~\ref{is-eigenvalue-conjugate-dictionary}(5).
\begin{itemize}
\item (a) The core has finitely many components at each level. By (3), $\Psi_c=\mathcal T$ and $p_0$ is
  $q_\rho$; the weight tie requires the matrix $-G_{W_i}^{\mathsf T}$ and the element $-\xi_0$, the
  triviality of $\chi_{c,\kappa_c}$ on $\Gamma_c$, and $|c_c|=1$, so that $c_c\ne0$ and the $|c_c|$ are
  bounded by $1$.
\item (b) is (3).
\item (c) The weight tie requires the dictionary to be a uniformization for $K_f$, which is $K_W$.
\item (d) The weight tie requires $\mu$ to be continuous, multiplicative, of absolute value $1$, trivial
  on $U(W_i)(F_h)$, and the minors $d_i(x(g),\tau)$ to have absolute value $1$; $\mu$ is trivial on
  $K_W$ by (2).
\item (e) is (1), and (f) is (2).\qedhere
\end{itemize}
\end{proof}
% (1) is HolDict.PrincLev.principal_level: key from mem_levG_triv_iff and T.level_character_embedding_compat
% (LiftLevel.isLift_one_iff_mem_level); levChi_triv_eq_one; hk.1 for GfC; congGroup (Member.FinCong V (g 0 *
% g 1)), isOpen_congGroup, and the law ht.laws.levelK (IsLevelK) through mem_of_isLevelK. The primes:
% ConjReading.pb_eq_of_omPins. (2) is PrincLevChar.char_eq_one (T.dictionary_character_trivial_on_principal_level)
% and HolDict.extWt_eq_laU_add (HolRead.mu_real_eq_prod_det with the law muVal and (ht.chiR i sigma).2.2.1;
% HolRead.SecWt.laU_eq_wtCol and extWt_arr_eq_wtCol). (3) is HolDict.BaseFrame.baseFrT_of_tieIn (baseFr_of with
% the laws dictU, adPin, modelW, unf, pins, familyUniformisationBasePointOrthogonal and ht.x0). (4) is
% HolDict.goodDict_core, field by field: inf = hM.inf, base, unif/model/muUnit/muCont through castDict_laws from
% the laws dictU, modelW, muUnit, muCont; chi = laws.chiTriv; cNe and cBdd from laws.cUnit; fin; levFin and
% levOpen from (1); muK from (2); wt from (2).

\begin{lemma}\label{is-eigenvalue-conjugate-adelization}
Let $D$, $V$, $\theta$ and $M$ be as in Definition~\ref{is-eigenvalue-conjugate-dictionary}, with $M$
of holomorphic shape (Lemma~\ref{is-eigenvalue-conjugate-models}(1)), let $w$ be integers and $K$ a
subgroup of $U(V)(\mathbf A)$, and let a dictionary for $V$ be good for $(M,w,K)$, with flag data
$X_c$ and weights $\kappa_c$. Let $P=(P_c)$ with $P_c$ a bounded section of weight $\kappa_c$ of $X_c$
(Definition~\ref{is-onto-share-flag-data}) for every $c$.
\begin{enumerate}
\item $\operatorname{ad}(P)$ is a bounded holomorphic automorphic form of weight $w$ and level $(K,1)$.
\item If $P_c\in S_{\kappa_c}(X_c)$ for every $c$, the constant terms of $\operatorname{ad}(P)$ vanish.
\end{enumerate}
\end{lemma}
\leanref{Development/AbelianVarietyTheory/Proofs/DivisorProofs/HolomorphicAdelicFormOfBoundedSectionStatement.lean}{GlobalTheory.T.holomorphicAdelicForm_of_boundedSection}
\leanref{Development/AbelianVarietyTheory/Proofs/DivisorProofs/AdelicCuspFormTransvectionConstantTermsZeroStatement.lean}{GlobalTheory.T.adelicCuspForm_transvection_constantTerms_zero}
\begin{proof}
Write $\phi=\operatorname{ad}(P)$, $\Gamma_c$, $\mathcal D_c$ and $\Lambda_c$ for the group, the flag
domain and the module of $X_c$, and use the conditions (a)--(f) of
Definition~\ref{is-eigenvalue-conjugate-dictionary}(5). As $\chi_{c,\kappa_c}$ is trivial on
$\Gamma_c$, $P_c(\gamma'x)=P_c(x)$ for $\gamma'\in\Gamma_c$ and $x\in\mathcal D_c$.

(1) \emph{Level.} Let $k\in K$. By (c), $(c(gk),x(gk))=(c(g),\gamma'x(g))$ with $\gamma'\in\Gamma_{c(g)}$,
and $x(g)\in\mathcal D_{c(g)}$; by (d), $\mu(gk)=\mu(g)\mu(k)=\mu(g)$. So $\phi(gk)=\phi(g)$.
\emph{Invariance.} In the same way $\phi(\gamma g)=\phi(g)$ for $\gamma\in U(V)(F_h)$, as $\mu(\gamma)=1$.

\emph{Holomorphy.} Let $h\in U(V)(\mathbf A)$, take $\gamma$ as in (b), and put $c=c(h)$ and
$A_\rho=(\gamma h)_\rho$, which lies in $U(G_V(\rho))$. By
Lemma~\ref{is-eigenvalue-conjugate-dictionary-facts}(2), $\mathcal D_c=\mathcal D_M$, and
$(A_\rho x_\rho)_\rho\in\mathcal D_M$ for $x\in\mathcal D_M$. Let $e_\rho$ be the integers of (f) and put
\[
 Q(x)=c_c\,\mu(h)\prod_\rho\det(x_{0,\rho})^{-e_\rho}\prod_\rho\det(x_\rho)^{e_\rho}\,P_c\big((A_\rho
 x_\rho)_\rho\big).
\]
It is holomorphic on $\mathcal D_M$, where $\det x_\rho\ne0$. For $b\in\mathcal B^{\mathcal T}$,
$\det(x_\rho b_\rho)=\det x_\rho\prod_a(b_\rho)_{aa}$ and
$P_c(yb)=\prod_{\rho,a}(b_\rho)_{aa}^{\kappa_{c,\rho,a}}P_c(y)$, so
$Q(xb)=\prod_{\rho,a}(b_\rho)_{aa}^{w_{\rho,a}}Q(x)$, since $w_{\rho,a}=\kappa_{c,\rho,a}+e_\rho$. Let $u$
have finite component $1$. By (b), $(c(hu),x(hu))=(c,\gamma'x)$ with $\gamma'\in\Gamma_c$ and
$x_\rho=A_\rho u_\rho x_{0,\rho}$, which lies in $\mathcal D_M$ because $(x_{0,\rho})_\rho$ does; and
$\mu(hu)=\mu(h)\prod_\rho\det(u_\rho)^{e_\rho}$ by (d) and (f). Hence
$\phi(hu)=c_c\mu(h)\prod_\rho\det(u_\rho)^{e_\rho}P_c(x)=Q\big((u_\rho x_{0,\rho})_\rho\big)$, and
$u\mapsto Q((u_\rho x_{0,\rho})_\rho)$ is holomorphic of weight $w$ on $U_\infty$.

\emph{Continuity.} Let $U$ be an open subgroup whose elements with archimedean component $1$ lie in $K$,
by (e), and let $L$ be the lift of Lemma~\ref{is-eigenvalue-conjugate-dictionary-facts}(1). Fix $h$, and
for $g$ put $A(g)=((h^{-1}g)_\rho)_\rho$ and $a(g)=L(A(g))^{-1}h^{-1}g$. Both depend continuously on $g$,
$a(g)$ has archimedean component $1$, and $a(h)=1$; so $a(g)\in U$ for $g$ near $h$. For such $g$,
$g=hL(A(g))a(g)$ and $\phi(g)=\phi(hL(A(g)))=Q\big((A(g)_\rho x_{0,\rho})_\rho\big)$ with the $Q$ of $h$,
which is continuous in $g$ because $Q$ is holomorphic on $\mathcal D_M$.

\emph{Boundedness.} By (d),
$|\phi(g)|=|c_{c(g)}|\,|P_{c(g)}(x(g))|=|c_{c(g)}|\,\|P_{c(g)}\|_{\kappa_{c(g)}}(x(g))$, since all
$|d_i(x(g),\tau)|=1$. The norms $\|P_c\|_{\kappa_c}$ are bounded on $\mathcal D_c$, $C$ is finite and the
$|c_c|$ are bounded, by (a); so $\phi$ is bounded.

(2) Let $\xi$ be isotropic, $\mu_Z$ a Haar measure on $Z_\xi$, $\mathcal F$ a fundamental domain and $g$
an element. By (1), $\phi$ is continuous, bounded, invariant on the left under $U(V)(F_h)$, and invariant
on the right under the elements with archimedean component $1$ of an open subgroup $U$. Put $c=c(g)$ and
take $\gamma$ as in (b) for $h=g$, with rational matrix $m$. As $\Lambda_c$ spans $E_h^n$,
$\xi'=Nm\xi\in\Lambda_c$ for some integer $N>0$, and $\xi'$ is isotropic for the form of $X_c$, whose
matrix is $-G_V^{\mathsf T}$, because $m$ preserves $G_V$. Let $B$ be the fundamental domain of
$\Pi_{\kappa_c}$ that defines the constant terms, of finite positive volume
(Lemma~\ref{is-eigenvalue-conjugate-dictionary-facts}(4)). As $P_c$ is a cusp form,
$\int_BP_c(\phi_{\xi'}(s)x)\,ds=0$ for $x\in\mathcal D_c$.

\emph{Step 1: means.} The group $Z_\xi$ is commutative, by
Lemma~\ref{is-eigenvalue-conjugate-dictionary-facts}(3). For $z'\in Z_\xi$, $z'\mathcal F$ is again a
fundamental domain, the function $z\mapsto\phi(zg)$ is invariant under $Z_\xi\cap U(V)(F_h)$, and $\mu_Z$
is invariant; so the integral of $z\mapsto\phi(z'zg)$ over $\mathcal F$ is that of $z\mapsto\phi(zg)$.
For a continuous map $s\mapsto u(s)$ from $\mathbf R^{\mathcal T}$ to $Z_\xi$, integrating over $B$ and
exchanging the integrals, which converge absolutely as $\phi$ is bounded and $\mu_Z(\mathcal F)$ is
finite, gives
\[
 \operatorname{vol}(B)\int_{\mathcal F}\phi(zg)\,d\mu_Z(z)=\int_{\mathcal F}\int_B\phi(u(s)zg)\,ds\,d\mu_Z(z).
\]
It suffices to find $u$ with $\int_B\phi(u(s)zg)\,ds=0$ for every $z\in Z_\xi$.

\emph{Step 2: the flow.} For $\rho\in\mathcal T$ and $t\in\mathbf R$ put $\zeta_\rho=\rho(\xi)$,
$T_\rho(t)=1+it\,\zeta_\rho\zeta_\rho^*G_V(\rho)$ and $k_\rho=c_\rho N^2$, with $c_\rho=2|\rho(\xi_0)|\ne0$
the constant of $X_c$. Since $\gamma_\rho$ preserves $G_V(\rho)$ and $\rho(\xi')=N\gamma_\rho\zeta_\rho$,
Lemma~\ref{is-eigenvalue-conjugate-dictionary-facts}(4) gives $\gamma_\rho T_\rho(k_\rho
s)=f_\rho(s)\gamma_\rho$, where $f_\rho(s)$ are the matrices of the flow $\phi_{\xi'}(s)$ of $X_c$. By
Lemma~\ref{is-eigenvalue-conjugate-dictionary-facts}(1) and~(3), $u(s)=L\big((T_\rho(k_\rho
s_\rho))_\rho\big)$ lies in $Z_\xi$ and depends continuously on $s$.

\emph{Step 3: values.} Let $v$ have finite component $1$ and $\gamma_\rho v_\rho=f_\rho(s)\gamma_\rho$ for
every $\rho$. Then $\phi(vg)=c_c\mu(g)P_c(\phi_{\xi'}(s)x(g))$. Indeed $g^{-1}vg$ has finite component $1$,
so by (b), $(c(vg),x(vg))=(c,\gamma'x)$ with $\gamma'\in\Gamma_c$ and
$x_\rho=(\gamma vg)_\rho x_{0,\rho}=f_\rho(s)(\gamma g)_\rho x_{0,\rho}=f_\rho(s)x(g)_\rho$, that is
$x=\phi_{\xi'}(s)x(g)$, which lies in $\mathcal D_c$ by
Lemma~\ref{is-eigenvalue-conjugate-dictionary-facts}(4). By (f), $\mu(v)=\prod_\rho\det(v_\rho)^{e_\rho}=1$,
as $\det f_\rho(s)=1$. So $\phi(vg)=c_c\mu(g)\mu(v)P_c(\gamma'x)=c_c\mu(g)P_c(x)$.

\emph{Step 4.} Let $z\in Z_\xi$. By Lemma~\ref{is-eigenvalue-conjugate-dictionary-facts}(3), applied to
the open subgroup $gUg^{-1}$, $z=\gamma_1vk$ with $\gamma_1,v,k\in Z_\xi$, $\gamma_1$ rational, $v$ of
finite component $1$, and $k$ of archimedean component $1$ with $g^{-1}kg\in U$. Write
$v_\rho=T_\rho(t_\rho)$ and put $s_0=(t_\rho/k_\rho)_\rho$. As $Z_\xi$ is commutative,
$u(s)zg=\gamma_1\,(u(s)vg)\,(g^{-1}kg)$, so $\phi(u(s)zg)=\phi(u(s)vg)$. The element $u(s)v$ has finite
component $1$ and $\gamma_\rho u(s)_\rho v_\rho=f_\rho(s)f_\rho(s_0)\gamma_\rho=f_\rho(s+s_0)\gamma_\rho$ by
Step 2 and Lemma~\ref{is-eigenvalue-conjugate-dictionary-facts}(4). By Step 3,
$\phi(u(s)zg)=c_c\mu(g)P_c\big(\phi_{\xi'}(s)\phi_{\xi'}(s_0)x(g)\big)$. As
$\phi_{\xi'}(s_0)x(g)\in\mathcal D_c$, the integral over $s\in B$ vanishes.
\end{proof}
% (1) is HolDict.BackHol.holForm_adel: hinv from IsSec and G.chi; hlev and the left invariance from G.unif
% (hl, hk, hdom) and G.muUnit; holomorphy through G.base, FdDom (T.flagDomain_frames_eq_gramFrames_unitary),
% frDoms_open_preconnected_stable and holSec_of_isSec with the constant c mu(h) prod det(x0)^(-m); continuity
% is continuous_of_holAd; boundedness through pnorm_eq_of_unit and G.cBdd. (2) is HolDict.TransvCusp.zCusp_adel:
% Step 1 is setIntegral_eq_zero_of_means; xi' from exists_rat and exists_isIso; B from exists_mean_zero
% (PeriodLatticeCountable.fd_exists); Step 2 is hAT (mul_transv, flowMat_eq_transv) and htv; Step 3 is hval;
% Step 4 is the decomposition hdec of T.transvectionSubgroup_cocompact_haar_fundamentalDomain and key.

\begin{lemma}\label{is-eigenvalue-conjugate-matrix-coefficient}
Let $D$, $V$, $\theta$ and $M$ be as in Definition~\ref{is-eigenvalue-conjugate-dictionary}, with $M$ of
holomorphic shape, let $v$ be an extremal vector of $M$ and $w$ the extremal weights of $M$ in the order
of the columns of its frames, as in Lemma~\ref{is-eigenvalue-conjugate-spectral-facts}(3)(d). Let $\nu$,
$\mathcal F$, $\mathcal H$ and $R$ be as in Lemma~\ref{is-eigenvalue-conjugate-spectral-facts}(2). If a
bounded holomorphic automorphic form $\phi$ of weight $w$ and some level represents $F\in\mathcal H$, then
\[
 \langle F,R(u)F\rangle=\|F\|^2\prod_\rho\langle v_\rho,\pi_\rho(u_\rho)v_\rho\rangle
\]
for every $u\in U(V)(\mathbf A)$ with finite component $1$.
\end{lemma}
\leanref{Development/AbelianVarietyTheory/Proofs/DivisorProofs/L2MatrixCoeffBoundedHolomorphicFormEqExtremalVecStatement.lean}{GlobalTheory.T.inner_rho_eq_normSq_mul_prod_of_holForm}
\begin{proof}
Let $L$ be the lift of Lemma~\ref{is-eigenvalue-conjugate-dictionary-facts}(1), and $t(s)$ and $\chi_w$
as in Lemma~\ref{is-eigenvalue-conjugate-spectral-facts}(3); the $t(s)$ lie in $U_\infty$, as the frames
$x_{0,\rho}$ are orthonormal in the sense of Lemma~\ref{is-eigenvalue-conjugate-dictionary-facts}(2). For
every $x$, the function $u\mapsto\phi(xL(u))$ on $U_\infty$ is holomorphic of weight $w$, by the
holomorphy of $\phi$ at $x$ and $L(u)_\rho=u_\rho$. By
Lemma~\ref{is-eigenvalue-conjugate-spectral-facts}(3)(a), $\phi(xL(t(s)))=\chi_w(s)\phi(x)$; so the right
translate of $\phi$ by $L(t(s))$ is $\chi_w(s)\phi$, and by
Lemma~\ref{is-eigenvalue-conjugate-spectral-facts}(2), $R(L(t(s)))F=\chi_w(s)F$.

Put $A(u)=\langle F,R(L(u))F\rangle$ and $B(u)=\prod_\rho\langle v_\rho,\pi_\rho(u_\rho)v_\rho\rangle$. By
Lemma~\ref{is-eigenvalue-conjugate-spectral-facts}(2),
$A(u)=\int_{\mathcal F}\overline{\phi(x)}\,\phi(xL(u))\,d\nu(x)$, an integral over a finite measure of
functions of $u$ that are holomorphic of weight $w$ and bounded by $C^2$ for a bound $C$ of $\phi$; so
$A$ is holomorphic of weight $w$, by Lemma~\ref{is-eigenvalue-conjugate-spectral-facts}(3)(c). If $U'$
is an isometry and $U'F=cF$ with $|c|=1$, then $\langle F,U'G\rangle=c\langle F,G\rangle$ for every $G$.
Applied to $U'=R(L(t(s)))$, with $L(t(s)u)=L(t(s))L(u)$, this gives $A(t(s)u)=\chi_w(s)A(u)$. By
Lemma~\ref{is-eigenvalue-conjugate-spectral-facts}(3)(d), $B$ is holomorphic of weight $w$, and the same
argument in each $H_\rho$ gives $B(t(s)u)=\chi_w(s)B(u)$. Moreover $A(1)=\|F\|^2$ and $B(1)=1$, as the
$v_\rho$ are unit vectors. By Lemma~\ref{is-eigenvalue-conjugate-spectral-facts}(3)(b), applied to
$A-\|F\|^2B$, $A=\|F\|^2B$ on $U_\infty$. For $u$ with finite component $1$, $u=L((u_\rho)_\rho)$ by
Lemma~\ref{is-eigenvalue-conjugate-dictionary-facts}(1), which gives the formula.
\end{proof}
% This is HolDict.inner_reg_eq_coef_extremal: horb, htor (smFn_and_killedFn_of_holFn, torus_law_of_killed),
% the eigenvector law hrho-t (repr_properties), hA (holFn_integral), hB and hBt (holFn_coef_extVec),
% inner_iso_eigen, and eq_const_mul (Uniq.eq_zero_of_torus_law); htm (TorusLaw.mul_torEl_mem) for t(s).

\begin{lemma}\label{is-eigenvalue-conjugate-multiplicity}
Let $D$, $V$, $\theta$ and $M$ be as in Definition~\ref{is-eigenvalue-conjugate-dictionary}, with $M$ of
holomorphic shape, and let $v$ and $w$ be as in Lemma~\ref{is-eigenvalue-conjugate-matrix-coefficient}.
Suppose that $U(V)(F_h)$ has a fundamental domain of finite volume for some right-invariant Haar measure
on $U(V)(\mathbf A)$. Let $J$ be a set of elements of $U(V)(\mathbf A)$ with archimedean component $1$,
$\chi$ a function on $U(V)(\mathbf A)$, and $\phi$ a bounded holomorphic automorphic form of weight $w$
and level $(J,\chi)$ whose constant terms vanish. Then there is a multiplicity map $T$ of $M$ of level
$(J,\chi)$ (Definition~\ref{asm-dependence-rows-models}(2)) with $T(v)=\phi$.
\end{lemma}
\leanref{Development/AbelianVarietyTheory/Proofs/DivisorProofs/BoundedHolomorphicFormEqMultiplicitySpaceValueStatement.lean}{GlobalTheory.T.boundedHolomorphicForm_eq_multiplicitySpace_value}
\begin{proof}
Let $\nu$ and $\mathcal F$ be the measure and the fundamental domain, $\mathcal H$ and $R$ as in
Lemma~\ref{is-eigenvalue-conjugate-spectral-facts}(2), $L$ the lift of
Lemma~\ref{is-eigenvalue-conjugate-dictionary-facts}(1), and $F_0\in\mathcal H$ the element that $\phi$
represents. For $\tau$ and $g\in U(G_V(\tau))$ write $g^{(\tau)}$ for the image of $g$ under $\theta$
at $\tau$, and $y[\tau\leftarrow y']$ for $y$ with its component at $\tau$ replaced by $y'$.

\emph{Step 1: an equivariant map.} By Lemma~\ref{is-eigenvalue-conjugate-matrix-coefficient} and
$L(u)_\rho=u_\rho$,
$\langle F_0,R(L(u))F_0\rangle=\|F_0\|^2\prod_\rho\langle v_\rho,\pi_\rho(u_\rho)v_\rho\rangle$ for
$u\in U_\infty$. The representations $\pi_\rho$ are unitary and irreducible, and $R\circ L$ is a
representation of $U_\infty$ by isometries. By Lemma~\ref{is-eigenvalue-conjugate-spectral-facts}(4)(a)
there is a continuous multilinear map $\Phi\colon\prod_\rho H_\rho\to\mathcal H$ with $\Phi(v)=F_0$, with
values in the closed span of the $R(L(u))F_0$, $u\in U_\infty$, and with
$\Phi(y[\tau\leftarrow\pi_\tau(g)y_\tau])=R(g^{(\tau)})\Phi(y)$. If $k$ has archimedean component $1$
and $\phi(gk)=c\phi(g)$ for all $g$, then $R(k)F_0=cF_0$, by
Lemma~\ref{is-eigenvalue-conjugate-spectral-facts}(2); as $k$ commutes with the $L(u)$,
$R(k)R(L(u))F_0=cR(L(u))F_0$, and by continuity $R(k)\Phi(y)=c\Phi(y)$ for all $y$.

\emph{Step 2: representatives.} Let $U$ be an open subgroup under whose elements of archimedean component
$1$ the function $\phi$ is invariant; by Step 1 they fix every $\Phi(y)$. For a tensor $y$ of smooth
vectors, $\Phi(y)$ is smooth, by Lemma~\ref{is-eigenvalue-conjugate-spectral-facts}(4)(b), and so is
represented by a function $T(y)$ as in Lemma~\ref{is-eigenvalue-conjugate-spectral-facts}(4)(c). As
representing is linear and continuous representatives are unique, $T$ is multilinear, and $T(v)=\phi$.
For an element $X$ of the Lie algebra at $\tau$, the right derivative of $T(y)$ along $X$ represents the
derivative of $R$ applied to $\Phi(y)$, which is $\Phi(y[\tau\leftarrow d\pi_\tau(X)y_\tau])$; the vector
$d\pi_\tau(X)y_\tau$ is smooth by Lemma~\ref{is-eigenvalue-conjugate-spectral-facts}(4)(e). By
uniqueness, this right derivative is $T(y[\tau\leftarrow d\pi_\tau(X)y_\tau])$. Iterating, every iterated
derivative of $T(y)$ along Lie algebra elements is a value of $T$, and each Gelfand invariant acts on
$T(y[\tau\leftarrow y'])$ as $T(y[\tau\leftarrow\,\cdot\,])$ of its action on $y'$.

\emph{Step 3: the conditions of a multiplicity map.}
\begin{itemize}
\item Each $T(y)$ is continuous, invariant on the left under $U(V)(F_h)$ and on the right under the
  elements of archimedean component $1$ of $U$, and smooth at the archimedean places, by Step 2.
\item Finiteness under the centres: by Lemma~\ref{is-eigenvalue-conjugate-spectral-facts}(3)(d),
  $v_\tau$ is a nonzero smooth vector on which the matrix units act as there, so each Gelfand invariant
  acts by a scalar on the smooth vectors of $\pi_\tau$, by
  Lemma~\ref{is-eigenvalue-conjugate-spectral-facts}(4)(f). By Step 2 it acts on $T(y)$ by the same
  scalar, so the line through $T(y)$ is stable under the centres.
\item Moderate growth: by Step 2 every iterated derivative of $T(y)$ along Lie algebra elements is a
  value $T(y')$, and by Lemma~\ref{is-eigenvalue-conjugate-spectral-facts}(4)(d), applied to $U$,
  $|T(y')(xu)|\le C\|u\|^N\sum_l\|l\cdot\Phi(y')\|$ for $u$ with finite component $1$.
\item Cuspidality: $T(y)$ represents $\Phi(y)$, which lies in the closed span of the $R(u)F_0$. As the
  constant terms of $\phi$ vanish, those of $T(y)$ vanish by
  Lemma~\ref{is-eigenvalue-conjugate-spectral-facts}(4)(h), and $T(y)$ is cuspidal by
  Lemma~\ref{is-eigenvalue-conjugate-spectral-facts}(4)(g).
\item Level: for $k\in J$, $R(k)\Phi(y)=\chi(k)\Phi(y)$ by Step 1. The right translate of $T(y)$ by $k$
  represents $R(k)\Phi(y)$, and $\chi(k)T(y)$ represents $\chi(k)\Phi(y)$; both are continuous, so
  $T(y)(gk)=\chi(k)T(y)(g)$.
\item Equivariance: $T(y[\tau\leftarrow\pi_\tau(g)y_\tau])$ and the right translate of $T(y)$ by
  $g^{(\tau)}$ both represent $R(g^{(\tau)})\Phi(y)$, so they are equal.
\item Continuity: let $x$ be a point and $\tau$ an embedding. By
  Lemma~\ref{is-eigenvalue-conjugate-spectral-facts}(4)(d) at $u=1$,
  $|T(y)(x)|\le C\sum_l\|l\cdot\Phi(y)\|$. By Lemma~\ref{is-eigenvalue-conjugate-spectral-facts}(4)(b),
  $l\cdot\Phi(y)$ is $\Phi$ of $y$ with each component replaced by the letters of $l$ at its place
  applied to it. So, as $\Phi$ is continuous, the functional $y'\mapsto T(y[\tau\leftarrow y'])(x)$ on the
  smooth vectors is bounded by a constant times $\sum_l\|l_\tau y'\|$, where $l_\tau$ is the word of the
  letters of $l$ at $\tau$. Any linear extension of it to the space of $\pi_\tau$ is continuous for this
  Sobolev norm.\qedhere
\end{itemize}
\end{proof}
% This is GlobalTheory.BackForm.exists_multMap: Step 1 is exists_map (exists_map_abs through
% T.exists_equivariant_multilinear_of_matrixCoeff_eq_unitaryProduct, eig_closure, eig_class); Step 2 is the
% choice psi of T.smooth_vector_automorphic_derivatives with uniq, MultilinearMap.mk' and hT; Step 3 is the
% refine of IsMultMap: zfin_value (gelf_value, T.gelfand_invariants_act_by_scalars_highest_weight with
% T.baseFrameBorel_acts_extremalVector_by_weight), the growth bound from T.sobolev_bound_for_smooth_vectors
% and word_value, cuspidality from T.cuspidality_preserved_in_closure and
% T.isCuspFun_of_transvection_constantTerms_zero, the level through e4, the equivariance through e2, and
% cont_value (word_map, sum_bound).

\begin{lemma}\label{is-eigenvalue-conjugate-characters}
Let $R$ be a realization of a definite-side datum $D$ with a face $F$ that satisfies the core
conditions at $F$, except possibly the one on its chain, with the notation of
Definition~\ref{is-eigenvalue-lift-data}. Let $i$ be a lift and $\sigma$ an embedding of $L$ that
agrees with $\sigma_0$ on $K$.
\begin{enumerate}
\item For every finite place $P$ of $E_h$, the value of $\eta$ at a uniformizer at $P$, divided by
  $|k_P|^{1/2}$, is $\sigma_0(y_P)$ for some $y_P\in L$.
\item There is a $\sigma$-conjugate $\eta'$ of $\eta$ (Lemma~\ref{is-eigenvalue-type-parts}(4)).
\item Every $\sigma$-conjugate $\eta'$ of $\eta$ restricts on $\mathbf A_{F_h}^\times$ to the quadratic
  character of $E_h/F_h$, and has archimedean exponent $\frac12$ at every $\rho\in\mathcal T$.
\item Let $\chi'$ be a $\sigma$-conjugate of $\mu^0_i$ (Lemma~\ref{is-eigenvalue-type-parts}(3)) and
  $\eta'$ a $\sigma$-conjugate of $\eta$. Then there is a
  Weil datum for $(V_i,W_i)$ with characters $(\chi',\chi'\eta'^{-1})$
  (Definition~\ref{th-weil-datum}(3)).
\end{enumerate}
\end{lemma}
\leanref{Development/AbelianVarietyTheory/Proofs/DivisorProofs/ConjugateCharactersEmbeddingsOfCharacterisationStatement.lean}{GlobalTheory.T.conjugateCharacters_embeddings_of_characterisation}
\leanref{Development/AbelianVarietyTheory/Proofs/DivisorProofs/ConjugateCharactersEmbeddingsOfCharacterisation.lean}{GlobalTheory.EtaValues.etaVals_five}
\leanref{Development/AbelianVarietyTheory/Proofs/DivisorProofs/ExistsIsEtaConj.lean}{GlobalTheory.EtaConjExist.exists_conj}
\leanref{Development/AbelianVarietyTheory/Proofs/DivisorProofs/CoefMapNormEquivariantBasisOfFive.lean}{GlobalTheory.EtaConjType.resEta_hasInfExp_of_isEtaConj_ofFive}
\leanref{Development/AbelianVarietyTheory/Proofs/DivisorProofs/ExistsConjWeil.lean}{GlobalTheory.GalLift.ConjWeil.exists_of_datum}
\leanref{Development/AbelianVarietyTheory/Proofs/DivisorProofs/ExistsWeilDatumTwistHeckeCharacters.lean}{GlobalTheory.WeilChange.exists_weilDatum_twist}
\begin{proof}
(1) This is the argument of the proof of Lemma~\ref{is-eigenvalue-base-inputs}(3)(iii), that $\eta$ is
a $\sigma_0$-conjugate of itself.

(2) By the core conditions on the level characters, $\eta$ is the character of the host, an unramified
unitary Hecke character of archimedean type $1$ on the embeddings $\varphi_g$ of the CM type of the
host. Let $y_P$ be as in (1) and put $c(P)=\sigma(y_P)/\sigma_0(y_P)$. For $k\in E_h^\times$, the product
formula for $\eta$ at $k$ gives $\prod_P\sigma_0(y_P)^{\operatorname{ord}_Pk}\prod_g\varphi_g(k)=1$. The
$\varphi_g$ are the composites of the CM embedding with the automorphisms of $E_h$ over $K$, so
$\prod_g\varphi_g(k)$ is $\sigma_0$ of the norm $n$ of $k$ to $K$. So
$\prod_Py_P^{\operatorname{ord}_Pk}=n^{-1}$ in $L$, and as $\sigma$ and $\sigma_0$ agree on $K$,
$\prod_Pc(P)^{\operatorname{ord}_Pk}=1$. Thus $P\mapsto c(P)$ is a character of the ideal class group of
$E_h$; let $\xi$ be the unramified Hecke character of finite order with $\xi(\varpi_P)=c(P)$ at
uniformizers $\varpi_P$. Then $\eta'=\eta\xi$ is an unramified Hecke character with
$\eta'(\varpi_P)|k_P|^{-1/2}=\sigma_0(y_P)c(P)=\sigma(y_P)$, a $\sigma$-conjugate of $\eta$.

(3) By the core conditions on the level characters, $\eta$ is unitary and unramified, restricts on
$\mathbf A_{F_h}^\times$ to the quadratic character of $E_h/F_h$, and has the archimedean type of (2).
As $\sigma$ agrees with $\sigma_0$ on $K$, the values of $\eta'$ at the uniformizers, which are those of
$\eta$ carried by $\sigma$, show that $\eta'$ has the same four properties; this is a computation with the
product formula, as in (2). A unitary character with this archimedean type has exponent $\frac12$ at
every $\rho$.

(4) The core conditions on the theta lifts give a Weil datum $\omega_P$ for $(V_i,W_i)$ with characters
$(\chi_S,\chi_T)$ whose restrictions to $\mathbf A_{F_h}^\times$ are the quadratic character to the powers
$n_V$ and $r$. The rank $n_V$ of $V_i$ is even and the rank $r$ of $W_i$ is odd. As $\chi'$ is trivial
on $\mathbf A_{F_h}^\times$, as a $\sigma$-conjugate of $\mu^0_i$, and $\eta'$ restricts to the quadratic character by (3), the Hecke characters
$\nu_S=\chi'\chi_S^{-1}$ and $\nu_T=\chi'\eta'^{-1}\chi_T^{-1}$ are trivial on $\mathbf A_{F_h}^\times$.
Then $\omega'(g,h)=\nu_T(x)^{-1}\nu_S(y)\,\omega_P(g,h)$, for ideles $x$ and $y$ that are roots of the
determinants of $g$ and of $h$, does not depend on the roots and is a Weil datum with characters
$(\chi_S\nu_S,\chi_T\nu_T)=(\chi',\chi'\eta'^{-1})$.
\end{proof}
% (1) is EtaValues.etaVals_five (hv_mem_host with the family tie's etaVal). (2) is
% EtaConjExist.exists_conj: exists_norm, finprod_alg (the product formula with IsEtaT: unramified with
% infinity type 1 on the CM type embPsi), finprod_quot, exists_classChar, CharTransport.exists_chrEq.
% (3) is resEta_hasInfExp_of_isEtaConj_ofFive (isEtaT_of_conj, Mem.EtaTilde.isEta_res, val_infinite).
% (4) is ConjWeil.exists_of_datum with the base Weil datum hR.thP.omegaP i mu0, hecke_mul_inv,
% resEta_parity, even_im_hermitian_sub (nV even, r odd) and WeilChange.exists_weilDatum_twist (the twist by
% IsDetRoot values). In the Lean, the hypotheses of (4) are carried by HypConj, which
% GalLift.hypConj_of_chars_ofFive builds from ChiAt, IsEtaConj and (3).

\begin{lemma}\label{is-eigenvalue-conjugate-coefficient-forms}
Let $R$ be a realization of a definite-side datum $D$ with a face $F$ that satisfies the core
conditions at $F$, except possibly the one on its chain, with the notation of
Definitions~\ref{is-eigenvalue-lift-data} and~\ref{is-eigenvalue-conjugate-datum}. Let $i$ be a lift
and $\sigma$ an embedding of $L$.
\begin{enumerate}
\item If $\sigma$ agrees with $\sigma_0$ on $K$, there are a coefficient datum $W_c$ and an isomorphism
  $c_1$ as in (2) of Definition~\ref{is-eigenvalue-conjugate-datum}.
\item For every family tie of $R$ and all $W_c$ and $c_1$ as in (2), there are a form space $H$ and a
  map $j$ as in (3) and (4) for that tie.
\item If $W_c$, $c_1$, $H$ and $j$ satisfy (2)--(4) for one family tie, then $W_c$, $c'_1$, $H$ and $j$
  satisfy them for any other family tie, for some $c'_1$.
\end{enumerate}
\end{lemma}
\leanref{Development/AbelianVarietyTheory/Proofs/DivisorProofs/ConjugateLiftOfCharacterisationClausesStatement.lean}{GlobalTheory.T.conjugateLift_of_characterisationClauses}
\leanref{Development/AbelianVarietyTheory/Proofs/DivisorProofs/ExistsCoefAtOfPosDef.lean}{GlobalTheory.CoefModel.exists_coefAt_of_posDef}
\leanref{Development/AbelianVarietyTheory/Proofs/DivisorProofs/ExistsBaseChange.lean}{GlobalTheory.ConjForms.exists_baseChange}
\leanref{Development/AbelianVarietyTheory/Proofs/DivisorProofs/ConjFormsRescale.lean}{GlobalTheory.GalLift.rescale_abstract}
\begin{proof}
(1) By the core conditions on the models, $V_i$ is positive definite, so each $G_{V_i}(\rho)$ is
positive definite. As $\sigma$ agrees with $\sigma_0$ on $K$, the embeddings $\sigma\circ\iota_\tau$ of
$E_h$, $\tau\in\mathcal T$, are the elements of $\mathcal T$, in the notation of
Lemma~\ref{is-eigenvalue-conjugate-good}, so that $p(\tau)=\sigma\circ\iota_\tau$ is a bijection of
$\mathcal T$. Let $W_c$ be the complex coefficient space of $V_i$ with the weights of $V_i$, with the
pairing given at $\tau$ by the inverse transpose of $G_{V_i}(p(\tau))$, and let $U(V_i)(\mathbf A)$ act
on it at $\tau$ through the inverse transpose of the archimedean component at $p(\tau)$, trivially at
the finite places. Let $c_1$ apply $\sigma$ to the
coefficients of a rational vector. A rational element of $U(V_i)(F_h)$ has archimedean components the
images of its rational matrix, so $c_1$ intertwines the actions. The rational basis goes to a basis, and
the pairing of $c_1(x)$ and $c_1(y)$ is $\sigma$ of the rational pairing of $x$ and $y$, because the
rational pairing has integral structure constants and the inverse transposes of the $G_{V_i}(p(\rho))$
are the images under $\sigma$ of the rational Gram matrices at the labels. This is a computation.

(2) The form space $H_V(i,\mu_0)$ has the class presentation of the tie: a form is determined by its
values at the representatives of the classes, and its values are images under the coefficient map
$c_0$ of the tie of rational vectors. The values of the character of the realization on the level of
$V_i$ lie in $\sigma_0(L)$, by the core conditions on the level characters and the group; let $\chi'$ be
the character whose values are their images under $\sigma$. Let $H$ be the space of functions on
$U(V_i)(\mathbf A_{F_h,f})$ with values in $W_c$ that transform under $U(V_i)(F_h)$ through $W_c$ and
under the level through $\chi'$, with the inner product of Definition~\ref{is-eigenvalue-lift-data}.
For a lift form $x$ and $h$, the value of $j_V(x)$ at $h$ is $c_0(v)$ for one rational vector $v$, as
$c_0$ is injective; put $j(x)(h)=c_1(v)$. The rational action and the level character transform $v$ as
they transform $c_0(v)$, so $j(x)\in H$, and $j$ is $\sigma$-semilinear.

(3) Let $c_0$ and $c'_0$ be the coefficient maps of the two ties. Both are given by identifications
$e$ and $e'$ of the coefficient space of the class presentation with the rational coefficient space
that carry its values to the same values and its pairing to the rational pairing, with coordinates in
the subfield $K_0$ of $L$ on which the involution is complex conjugation at every embedding. Put
$c'_1=c_1\circ e'\circ e^{-1}$. It carries the rational basis to a basis and, as the coordinates lie in
$K_0$, the rational pairing, followed by $\sigma$, to the inner product. It intertwines the rational
action: $c_0$ and $c'_0$ both intertwine it, $c'_0e'e^{-1}=c_0$, and $c'_0$ is injective, so $e'e^{-1}$
commutes with it. Every value $c'_0(v')$ is $c_0(v)$ with $v=e(e'^{-1}v')$, and then $c_1(v')=c'_1(v)$; so
$j$ satisfies (4) for the other tie.
\end{proof}
% (1) is CoefModel.exists_coefAt_of_posDef (exists_coef_round_at: coefOf lam with toHilb, realGL reindexed
% by the label bijection, invTrK, mapVV sigma, realGL_rat, pairK_mapVV_int), with hR.models.vpos and
% labels_overK. (2) is ConjForms.exists_conjForms_five through exists_baseChange (exists_conjChar,
% AdFormsModel.exists_hmod_adForms, PresBasic.coef_injective, HilbertModel.exists_lift_fun) and levChar_mem.
% (3) is GalLift.conjForms_rescale_five through rescale_abstract (rescale_mem with cmMax, rescale_pair,
% rescale_inj).

\begin{lemma}\label{is-eigenvalue-conjugate-level-data}
Let $R$ be a realization of a definite-side datum $D$ with a face $F$ that satisfies the core
conditions at $F$, except possibly the one on its chain, with the notation of
Definitions~\ref{is-eigenvalue-lift-data} and~\ref{is-eigenvalue-conjugate-datum}. Let $i$ be a lift
and $\sigma$ an embedding of $L$ that agrees with $\sigma_0$ on $K$.
\begin{enumerate}
\item There are a level vector $\phi$ of the trivial level (Definition~\ref{is-eigenvalue-lift-data}(6),
  for the trivial local level, no ramified places and the trivial characters) and bases $B$ as in
  Definition~\ref{is-eigenvalue-lift-data}(5), and $\phi_Q=\phi\otimes_BQ$ is a theta function for every
  polynomial $Q$.
\item Let $\chi'$ be a $\sigma$-conjugate of $\mu^0_i$, $\eta'$ a
  $\sigma$-conjugate of $\eta$, $\omega$ a Weil datum for $(V_i,W_i)$ with characters
  $(\chi',\chi'\eta'^{-1})$, $W_c$ and $c_1$ as in Lemma~\ref{is-eigenvalue-conjugate-coefficient-forms}(1),
  and $\phi$ and $B$ as in (1). Then there is a map $\mathrm{io}$ as in
  Definition~\ref{is-eigenvalue-lift-data}(7) for $\omega$, $\phi$, $B$ and the delta polynomial
  $\Delta_i(\sigma)$.
\item Let moreover $H$ and $j$ be as in Lemma~\ref{is-eigenvalue-conjugate-coefficient-forms}(2). For
  every $\varphi\in H$, every tensor $v$ of lowest $K$-type vectors and every $g\in U(W_i)(\mathbf A)$,
  the function $h\mapsto\theta_\omega\big(h,g;\phi_{\mathrm{io}(v\otimes\tilde\varphi(h))}\big)$ is
  integrable on the fundamental domain of the realization at $(i,\mu_0)$; so the theta integral
  $\Theta$ of these data is defined on $H$.
\end{enumerate}
\end{lemma}
\leanref{Development/AbelianVarietyTheory/Proofs/DivisorProofs/ConjugateLiftOfCharacterisationClauses.lean}{GlobalTheory.GalLift.Vec.exists_vecData_five}
\leanref{Development/AbelianVarietyTheory/Proofs/IwasawaDecomposition/GramMapMulInv.lean}{GlobalTheory.PlaceVec.thPVecExists}
\leanref{Development/AbelianVarietyTheory/Proofs/DivisorProofs/ExistsIoConjOfFive.lean}{GlobalTheory.MovedWeights.exists_io_conj_ofFive}
\leanref{Development/AbelianVarietyTheory/Proofs/DivisorProofs/ExistsIsThetaPForms.lean}{GlobalTheory.ConjThetaIntegral.IsThetaP₂}
\begin{proof}
(1) By the core conditions on the level characters, the lattices of the frames of $V_i$ and $W_i$ are
self-dual, and the base tuple has no ramified places. The Cartan involution of $V_i$ at every $\rho$
is $1$, as $V_i$ has an orthonormal basis there by the core conditions on the models, and those of $W_i$
are the Cartan involutions of the model at $\sigma$. By
Lemma~\ref{is-eigenvalue-conjugate-vacuum-exists}(2) there is a Fock vacuum vector $f_0$ for these, and
the place model of $f_0$ at the places of $P_g$ applied to the level function gives a level vector
$\phi$ with vacuum $f_0$. Since $W_i$ has signature $(p_\rho,q_\rho)$ at $\rho$ and the Cartan
involutions of the model are Cartan involutions of $W_i$, bases $B$ as in (5) exist. By
Lemma~\ref{is-eigenvalue-conjugate-theta-facts}(3) and~(4), $\phi$ and then every $\phi\otimes_BQ$ are
theta functions.

(2) This is a computation in the Fock model at the archimedean places, which constructs $\mathrm{io}$
from a polynomial model of $W_c$ given by $c_1$ and the orthonormal bases of $V_i$ of the realization.
It uses that the maps $\tau\mapsto\sigma\circ\iota_\tau$ of (1) of
Lemma~\ref{is-eigenvalue-conjugate-coefficient-forms} are a bijection onto $\mathcal T$, that $W_i$ has
the signature of the datum at every $\rho$, that the model at $\sigma$ has at every $\rho$ the
holomorphic parameter of the label of $\sigma$ at $\rho$ (Lemma~\ref{is-eigenvalue-conjugate-good}),
that the archimedean exponents of $\chi'$ are $-e_i$ at the labels and those of $\eta'$ are $\frac12$
(Lemma~\ref{is-eigenvalue-conjugate-characters}(3)), and that $\phi$ is a theta function whose vacuum
is a Fock vacuum vector.

(3) The function $\tilde\varphi$ is continuous and transforms under $U(V_i)(F_h)$ through $W_c$, as
$\varphi\in H$ and the character of $H$ is trivial where that of the realization is. As $V_i$ is
positive definite, the theta kernel of the theta functions $\phi_{\mathrm{io}(v\otimes w)}$ against such
a function is integrable on the fundamental domain for the measure of the realization.
\end{proof}
% (1) is GalLift.Vec.exists_vecData_five: SDFrameLattice.isSelfDualC_of_isSDFrame for betaV, betaW; Q empty
% from lev.Q_ram and xi_base; Chain.VacVec.cart_isCartanInv; Member.vacuum_oriented_exists (= VacuumO.vacExists,
% Lemma is-eigenvalue-conjugate-vacuum-exists(2)); PlaceVec.thPVecExists; exists_isONBasis; theta_pureTensorO
% with theta_of_thetaAt. (2) is MovedWeights.exists_io_conj_ofFive (exists_poly_model, weights_label,
% exists_isIoW with mu_inf and hasInfExp_mul/inv of eta_inf, DatumFock.isFockVacAt_thPVec). (3) is
% ConjThetaIntegral.exists_isThetaP_conj_five through IsThetaP2 (ext_continuous,
% DefiniteQuotient.integrableOn_theta, ThetaIntegralLinear.exists_isThetaP_of_integrable).

\begin{lemma}\label{is-eigenvalue-conjugate-phase}
Let $R$ be a realization of a definite-side datum $D$ that satisfies the core conditions on its models,
with the notation of Definition~\ref{is-eigenvalue-lift-data}; let $i$ be a lift, $\sigma$ an embedding
of $L$, and $K_W$ the level of the trivial level datum of $W_i$. Let $u$ be a function on
$U(W_i)(\mathbf A)$ with $|u|=1$, invariant on the left under $U(W_i)(F_h)$ and on the right under the
elements with finite component $1$ and under $K_W$. Then there is a unitary operator $U$ of
$H_W(i,\sigma)$ with $U(Y)(v)=u\,Y(v)$ for all $Y$ and all tensors $v$.
\end{lemma}
\leanref{Development/AbelianVarietyTheory/Proofs/DivisorProofs/CentralShiftPhaseConjReadingsFiveClauses.lean}{GlobalTheory.CompPhase.exists_isCompPhase_five}
\leanref{Development/AbelianVarietyTheory/Proofs/DivisorProofs/CompatiblePhaseExists.lean}{GlobalTheory.CompPhase.compFun_of_level}
\leanref{Development/AbelianVarietyTheory/Proofs/DivisorProofs/CompatiblePhaseExists.lean}{GlobalTheory.CompPhase.CompFun.multMap}
\begin{proof}
Every element is a product of one with finite component $1$ and one with archimedean component $1$;
so $u$ is invariant on the right under the open subgroup of the elements whose finite component lies in
the finite part of $K_W$, and $u$ is locally constant. The elements $n$ with $u(ng)=u(g)$ for all $g$
form a closed subgroup that contains $U(W_i)(F_h)$ and the elements with finite component $1$, since
$u(kg)=u(g\,g^{-1}kg)$; so it contains the unipotent radical of the parabolic subgroup of every isotropic
subspace, by approximation in that radical.

Let $T$ be a multiplicity map of the model at $\sigma$ of level $(K_W,\chi_K)$. Then $uT$ is one: $uT(v)$
is continuous, invariant on the left under $U(W_i)(F_h)$ and on the right under an open subgroup, and
smooth with $u$ constant along the archimedean directions, so its derivatives are $u$ times those of
$T(v)$, which gives finiteness under the centres and moderate growth; its constant terms are $u(g)$ times
those of $T(v)$, by the invariance under the unipotent radicals, and vanish; the level law and the
equivariance hold as $u$ is invariant under $K_W$ and the archimedean points; and the Sobolev bounds are
those of $T$ times $|u(x)|$. The Petersson pairing of $uT$ and $uT'$ is that of $T$ and $T'$, as $|u|=1$.
By the core conditions on the models, $H_W(i,\sigma)$ is the multiplicity space of this level with an
inner product given by the Petersson pairing; so multiplication by $u$ is a unitary operator $U$ of
$H_W(i,\sigma)$ with $U(Y)(v)=uY(v)$.
\end{proof}
% This is CompPhase.exists_isCompPhase_five: compFun_of_level (UnitarySplit.real_finite_split,
% ParahoricCompact.isOpen_comap_lev), CompFun.cont, CompFun.unip (NI_le), CompFun.multMap (CompFun.aut,
% CompFun.gelf, CompFun.cusp), ipM_phase and exists_unitary with hR.models.HW.

\begin{lemma}\label{is-eigenvalue-conjugate-extension}
Let $R$ be a realization of a definite-side datum $D$ over a field $K_s\subset L$
(Definition~\ref{is-realisation}), with one-map data that have a weight tie
(Definition~\ref{is-one-map-data}), and let $i$ be a lift at which the one-map data are pinned. Write
$L'\subset\mathbf C$ for the coefficient field of the one-map data at $i$, which contains $K_s$, and
$\iota\colon L\to L'$ for the map that the pins give. Let $\sigma$ be an embedding of $L$.
\begin{enumerate}
\item There is an embedding $\sigma'$ of $L'$ with $\sigma'\circ\iota=\sigma$.
\item If $\sigma'$ is an embedding of $L'$ with $\sigma'\circ\iota=\sigma$, then $\sigma'$ and $\sigma$
  agree on $K_s$.
\end{enumerate}
\end{lemma}
\leanref{Development/AbelianVarietyTheory/Proofs/DivisorProofs/OscillatorModelsWeightCompat.lean}{GlobalTheory.GalOne.exists_ext}
\leanref{Development/AbelianVarietyTheory/Proofs/DivisorProofs/CoefMapNormEquivariantBasisOfFive.lean}{GlobalTheory.GalOne.Facts.factsAt_ofFive}
\begin{proof}
(1) The pins make $L'$ the image $\sigma_0(L)$, and $\iota(x)=\sigma_0(x)$ for $x\in L$. So $\iota$ is a
bijection, and $\sigma'=\sigma\circ\iota^{-1}$ is an embedding of $L'$ with $\sigma'\circ\iota=\sigma$.

(2) By the weight tie and the pins of the core, $\sigma_0(x)=x$ for $x\in K_s$, as in the proof of
Lemma~\ref{is-eigenvalue-conjugate-good}(1). So $\iota(x)=x$, and
$\sigma'(x)=\sigma'(\iota(x))=\sigma(x)$.
\end{proof}
% (1) is GalOne.exists_ext, from the field pin O.c'.F = fieldRange of embLL and hP.lift.field
% (inC c' o jL = embLL). (2) is the first clause res T sigma = sOf O sigma' of FactsAt, proved at the
% start of Facts.factsAt_ofFive (hj from hP.lift.field and res_embLL with tie.alg_eq and tie.laws.pins).

\begin{lemma}\label{is-eigenvalue-conjugate-phase-function}
Let $R$ be a realization of a definite-side datum $D$ over a field $K_s$ that satisfies the core
conditions on its level characters and on its models, with one-map data that have a weight tie and
are pinned at some lift (Definition~\ref{is-one-map-data}), and with the notation of
Definition~\ref{is-eigenvalue-lift-data}. Let $i$ be a lift, $\sigma$ a good embedding of $L$
(Lemma~\ref{is-eigenvalue-conjugate-good}) with restriction $s$ to $K_s$, and $h\mapsto(c(h),x(h))$ the
dictionary of the core at level $i+1$ and $s$, a dictionary for $W_i$ by
Lemma~\ref{is-eigenvalue-conjugate-dictionary-good}. Let $a$ be a function on its set of components with
$|a|=1$, and put $u(h)=a(c(h))$ for $h\in U(W_i)(\mathbf A)$. Then $|u|=1$, and $u$ is invariant on the
left under $U(W_i)(F_h)$ and on the right under the elements with finite component $1$ and under
$K_W$.
\end{lemma}
\leanref{Development/AbelianVarietyTheory/Proofs/DivisorProofs/CoefMapNormEquivariantBasisOfFive.lean}{GlobalTheory.GalOne.comp_inv_ofFive}
\begin{proof}
By Lemma~\ref{is-eigenvalue-conjugate-dictionary-good}(4), the dictionary is good for $(M,w,K_W)$, where
$M$ is the model of $W_i$ at $\sigma$ and $w$ its extremal weights. By the uniformization condition of
Definition~\ref{is-eigenvalue-conjugate-dictionary}(5)(c), $c(\gamma h)=c(h)$ and $c(hk)=c(h)$ for
$\gamma\in U(W_i)(F_h)$ and $k\in K_W$, and $c(hg)=c(h)$ for $g$ with finite component $1$. So $u$ has
the three invariances, and $|u|=1$ as $|a|=1$.
\end{proof}
% This is GalOne.comp_inv_ofFive: (HolDict.goodDict_core ... hpb ...).unif gives h1 (rational elements),
% h2 (the level) and h3 (finite component 1), each keeping the component. The hypothesis hpb
% (T.pb = R.C.P.pb) comes from the pins, as in Lemma is-eigenvalue-conjugate-dictionary-good.

\begin{lemma}\label{is-eigenvalue-conjugate-class-values}
Let $R$ be a realization of a definite-side datum $D$ with a face $F$ that satisfies the core
conditions at $F$, except possibly the one on its chain, with one-map data that have a weight tie, and
let $i$ be a lift at which they are pinned, with $L'$, $\iota$ and the coefficient map $c'_0$ of the
one-map data at $i$ as in Lemma~\ref{is-eigenvalue-conjugate-class-sums}. For a rational coefficient
vector $v$ of $V_i$ over $L$, write $\iota(v)$ for the vector over $L'$ whose coefficients are those of
$v$ mapped by $\iota$. Let $c_0$ be the coefficient map of a family tie of $R$ with
$c'_0(\iota(v))=c_0(v)$ for every $v$, let $\sigma$ be an embedding of $L$, and let $W_c$, $c_1$, $H$
and $j$ satisfy (2)--(4) of Definition~\ref{is-eigenvalue-conjugate-datum} for this tie. Let $x$ be a
lift form at $(i,\mu_0)$, $y\in U(V_i)(\mathbf A_{F_h,f})$ and $x'$ a rational coefficient vector of
$V_i$ over $L'$ with $j_V(x)(y)=c'_0(x')$. Then there is a rational coefficient vector $x''$ of $V_i$
over $L$ with $\iota(x'')=x'$ and $j(x)(y)=c_1(x'')$.
\end{lemma}
\leanref{Development/AbelianVarietyTheory/Proofs/DivisorProofs/CoefMapNormEquivariantBasisOfFive.lean}{GlobalTheory.GalOne.exists_values_ofFive}
\begin{proof}
By (4) of Definition~\ref{is-eigenvalue-conjugate-datum} there is a rational coefficient vector $v$ over
$L$ with $j_V(x)(y)=c_0(v)$ and $j(x)(y)=c_1(v)$, and $c_0(v)=c'_0(\iota(v))$. The map $c'_0$ is
injective. Indeed, the pins of the one-map data at $i$ include a family tie of $R$ whose coefficient map
$c^{\mathrm p}_0$ satisfies $c'_0(\iota(v))=c^{\mathrm p}_0(v)$ for every $v$ and carries the rational
basis to linearly independent vectors. The images under $\iota$ of the rational basis form a basis over
$L'$, which $c'_0$ carries to these vectors, and $c'_0$ is semilinear for the inclusion of $L'$ in
$\mathbf C$. Hence $c'_0(x')=c'_0(\iota(v))$ gives $x'=\iota(v)$, and $x''=v$ has the two properties.
\end{proof}
% This is GalOne.exists_values_ofFive: hj.base gives v, hc the base change, and the injectivity is
% inj_of_indep with the last clause of AdTh.coefMap_norm_equivariant_basis_ofFive' (CoefExt.linearIndependent
% from the tie of the pins hP.lift.c0 and its mapL).

\begin{lemma}\label{is-eigenvalue-conjugate-class-action}
Let a core of a definite-side datum $D$ be given (Definition~\ref{is-one-map-data}), with coefficient
field $K_s$, a level $j$ and a subfield $L'\subset\mathbf C$ containing $K_s$. A section of the core at
level $j$ over $K_s$, or over $L'$, assigns a value to every framed point of the moduli of the core at
level $j$ over a $K_s$-algebra, or over an $L'$-algebra $S$. A section $f'$ over $L'$ is
$\sum_ka_kf_k$ \emph{after extension of scalars}, for sections $f_k$ over $K_s$ and $a_k\in L'$, if
$f'(y)=\sum_ka_kf_k(y|_{K_s})$ for every framed point $y$ over an $L'$-algebra $S$, where $y|_{K_s}$ is $y$
viewed over $S$ as a $K_s$-algebra. Let a group $G$ act on the framed points over $K_s$-algebras by
changes of their level structures through endomorphisms $\ell_g$ of the lattice of the core, with the
laws of an action, and on the sections over $K_s$ by $(gf)(y)=f(y_g)$, where $y_g$ is a framed point
obtained from $y$ by the change through $\ell_g$.
\begin{enumerate}
\item The changes through the $\ell_g$ satisfy the laws of an action on the framed points over
  $L'$-algebras, and so define an action $f'\mapsto gf'$ of $G$ on the sections over $L'$.
\item Let $f'$ be a section over $L'$ and $f_1,\dots,f_n$ sections over $K_s$, all of one weight, and let
  $a_1,\dots,a_n\in L'$ with $f'=\sum_ka_kf_k$ after extension of scalars. Then $gf'=\sum_ka_k(gf_k)$
  after extension of scalars for every $g\in G$.
\end{enumerate}
\end{lemma}
\leanref{Development/AbelianVarietyTheory/Proofs/CuspFormEvaluations/ActLevAndActBC.lean}{GlobalTheory.FieldBC.actLev_and_actBC}
\begin{proof}
(1) A framed point over an $L'$-algebra $S$ restricts to a framed point over $S$ viewed as a
$K_s$-algebra, and the change of level structure through $\ell_g$ commutes with this restriction. Each
law of an action for $L'$-algebras is the same law for $K_s$-algebras, applied to the restricted points.

(2) Let $y$ be a framed point over an $L'$-algebra $S$, and $y_g$ a framed point obtained from $y$ by
the change through $\ell_g$. Its restriction $y_g|_{K_s}$ is obtained from $y|_{K_s}$ by the same change.
The value of a section of a weight at a framed point depends only on its isomorphism class, and two
framed points obtained from one point by the change through $\ell_g$ are isomorphic. So
$(gf')(y)=f'(y_g)=\sum_ka_kf_k(y_g|_{K_s})=\sum_ka_k(gf_k)(y|_{K_s})$.
\end{proof}
% This is FieldBC.actLev_and_actBC: (1) is actLev_of_le (each law of ActLev for c' is the law for c at
% y.res); (2) is Pj.actBC_of, with evalAt_actL, levActR_res for the restriction of the moved point, and
% evalAt_congr_id with EvalWD from the chart cover of the core (Frames.chartCoverF_of_secBaseChange, from
% abelianSchemeLaws_all.secBaseChange).

\begin{lemma}\label{is-eigenvalue-conjugate-unfolding}
Let $V$ be a Hermitian space of rank $n$ of a definite-side datum $D$, with an adelic matrix $\beta$
and a level datum $L$. Write $K\subset GL_n(E_h\otimes_{F_h}\mathbf A_{F_h,f})$ for the open subgroup
that $\beta$ and $L$ cut out, and $J$ for the level of $V$ for $\beta$ and $L$: the intersection with
$U(V)(\mathbf A_{F_h,f})$ of the subgroup that $\beta$ and $L$ with its local level replaced by the
trivial one cut out. Suppose that these two subgroups are equal. Suppose also that $V$ is positive definite and that its signature is correct
(Definition~\ref{th-spaces}); these two hypotheses are not used in the proof. Let $\nu$ be a Haar
measure on $U(V)(\mathbf A)$ that gives the set of elements with finite part in $K$ and component in
$B_\rho$ at each archimedean place $\rho$ the measure $\prod_\rho m_\rho(B_\rho)$, for probability Haar
measures $m_\rho$ on the archimedean components, and let $\mathcal F$ be a fundamental domain for
$U(V)(F_h)$. Let $J'\subset J$ be a subgroup of finite index, let
$h_1,\dots,h_N\in U(V)(\mathbf A)$ have archimedean component $1$ and finite parts that represent the
double cosets $U(V)(F_h)\backslash U(V)(\mathbf A_{F_h,f})/J'$, each once, and let $\Gamma_k$ be the group
of $\gamma\in U(V)(F_h)$ with $h_k^{-1}\gamma h_k\in J'$. Let $\Phi$ be a function on $U(V)(\mathbf A)$
that is invariant on the left under $U(V)(F_h)$ and on the right under every element whose finite part
lies in $J$. Then
\[
 \int_{\mathcal F}\Phi\,d\nu=\sum_k[J:J']^{-1}|\Gamma_k|^{-1}\Phi(h_k).
\]
\end{lemma}
\leanref{Development/AbelianVarietyTheory/Proofs/DivisorProofs/IntegralEqClassSumOfLe.lean}{GlobalTheory.ClsInt.integral_eq_classSum_of_le}
\leanref{Development/AbelianVarietyTheory/Proofs/DivisorProofs/IntegralEqClassSumOfLe.lean}{GlobalTheory.ClsInt.integral_eq_sum_of_le}
\begin{proof}
As $\beta$ and $L$ with the trivial local level cut out $K$, $J=K\cap U(V)(\mathbf A_{F_h,f})$, and the
elements with finite part in $J$ are those with finite part in $K$. They form an open set, of volume $\prod_\rho m_\rho(U(V)(F_{h,\rho}))=1$
for $\nu$. By the invariances, $\Phi$ is constant on each set of the elements whose finite part lies in one double
coset $U(V)(F_h)yJ$. Group the $h_k$ by these double cosets, and choose one $y$ in each of them. Unfolding
the integral over $\mathcal F$ along the double cosets of $J$, for the measure that gives the elements
with finite part in $J$ volume $1$, gives $\int_{\mathcal F}\Phi\,d\nu=\sum_y|\Gamma^J_y|^{-1}\Phi(y)$,
where $\Gamma^J_y$ is the group of $\gamma\in U(V)(F_h)$ with $y^{-1}\gamma y\in J$, which is finite. For
each $y$, counting the double cosets of $J'$ in $U(V)(F_h)yJ$ gives
$\sum_k[J:J']^{-1}|\Gamma_k|^{-1}=|\Gamma^J_y|^{-1}$, the sum over the $h_k$ in the double coset of $y$.
As $\Phi(h_k)=\Phi(y)$ for these $h_k$, the formula follows.
\end{proof}
% The general statement is ClsInt.integral_eq_sum_of_le: the grouping is the setoid r and hreps, the
% unfolding along the double cosets of J is ClassUnfold.integral_eq_sum_of_surjective with the finiteness
% hfin (InnerClassSum.finite_stab_of_measure), and the count is sum_inv_card_stab. In
% integral_eq_classSum_of_le the level is (levEnv V beta L).J, the level of flag1L L, and the measure
% satisfies IsMeasDef V (lev beta L) (IsMeasLevDef and IsFundamentalDomain). The flag1L condition hL gives
% hK, which moves the volume measure_comap_finE from lev beta L to J; the openness is
% ThetaIntegrand.eventually_level. The hypotheses IsPosDef and SigOK are unnamed and not used.

\subsection{Frames and Zak theta functions}\label{subsection-is-zak-theta}

\begin{lemma}\label{is-eigenvalue-conjugate-frame-vacuum}
Let $D$ be a definite-side datum with host field $E_h/F_h$, let $S$ and $T$ be Hermitian spaces of $D$
with Gram matrices $G_S$ and $G_T$, and let $(g,h)\in U(S)(\mathbf A)\times U(T)(\mathbf A)$ act on
$\mathbb W$ by $x\mapsto hxg^{-1}$ (Definition~\ref{th-weil-datum}). Write $\hat{\mathcal O}$ for the
profinite completion of $\mathcal O_{E_h}$ and $\mathbf A_{E_h,f}=E_h\otimes_{F_h}\mathbf A_{F_h,f}$. For
$\beta_S\in GL_{\operatorname{rk}S}(\mathbf A_{E_h,f})$, an element $g\in U(S)(\mathbf A)$ whose finite
part $g_f$ satisfies $\beta_S^{-1}g_f\beta_S\in GL_{\operatorname{rk}S}(\hat{\mathcal O})$, and a finite
place $P$ of $E_h$, write $\bar g_P\in GL_{\operatorname{rk}S}(k_P)$ for the reduction of
$\beta_S^{-1}g_f\beta_S$ at $P$, and similarly for $T$. For an idempotent $e$ of $\mathbf A_{E_h}$ and
$g\in U(S)(\mathbf A)$ put $\det_eg=e\det g+1-e$.
\begin{enumerate}
\item (Vacua \cite{Kud94}.) Let $\psi$ be a character of $\mathbf A_{F_h}$, $\Lambda_S$ and $\Lambda_T$
  self-dual lattices of $S$ and $T$, $\Lambda\subset\mathbb W$ the lattice of the pair that they give,
  and $\omega$ a Weil datum for $(S,T)$ with characters $(\chi_S,\chi_T)$
  (Definition~\ref{th-weil-datum}). Let $\beta_S\in GL_{\operatorname{rk}S}(\mathbf A_{E_h,f})$ and
  $\beta_T\in GL_{\operatorname{rk}T}(\mathbf A_{E_h,f})$ satisfy
  $\beta_S\hat{\mathcal O}^{\operatorname{rk}S}=\hat{\mathcal O}\Lambda_S$ and
  $\beta_T\hat{\mathcal O}^{\operatorname{rk}T}=\hat{\mathcal O}\Lambda_T$, and call $(g,h)$ a
  \emph{stabilizer pair} if $g$ and $h$ have archimedean component $1$,
  $\beta_S^{-1}g_f\beta_S\in GL_{\operatorname{rk}S}(\hat{\mathcal O})$ and
  $\beta_T^{-1}h_f\beta_T\in GL_{\operatorname{rk}T}(\hat{\mathcal O})$. Suppose that the quadratic
  character $\eta$ of $E_h/F_h$ is the restriction to $\mathbf A_{F_h}^\times$ of a Hecke character of
  $E_h$ that is unramified at every finite place and has archimedean exponents $\frac12$. Let $P$ be a
  finite set of places of $E_h$ and $e$ the sum of the idempotents of $\mathbf A_{E_h}$ at the places of
  $P$, with $e\bar e=0$, and suppose that $\chi_S$ and $\chi_T$ are trivial on every idele that has
  archimedean component $1$, is a unit at every finite place and has component $1$ at the places of $P$
  and at their complex conjugates. Then
  \[
   \omega(g,h)f=\chi_T(\det_eg)^{-1}\chi_S(\det_eh)\,f
  \]
  for every stabilizer pair $(g,h)$ and every lattice vacuum $f$ for $\Lambda$.
\item (Level functions.) Let $\beta_S\in GL_{\operatorname{rk}S}(\mathbf A_{E_h,f})$ and
  $\beta_T\in GL_{\operatorname{rk}T}(\mathbf A_{E_h,f})$, let $P_g$ and $Q$ be finite sets of places of
  $E_h$, $\xi_P$ a character of $k_P^\times$ for each place $P$, $a\in\mathbf A_{E_h}$, and
  $r\le\operatorname{rk}S$ with $\operatorname{rk}T=r$. For $x\in\mathbb W$ let
  $y(x)=\beta_S^{-1}(G_S^{-1}\bar x^{\mathsf T}G_T)_f\beta_T$, an $\operatorname{rk}S\times r$ matrix over
  $\mathbf A_{E_h,f}$, and, when its entries lie in $\hat{\mathcal O}$, let $y_P(x)$ be its reduction at
  $P$. The \emph{level function} $\Phi$ of $\beta_S$, $\beta_T$, $P_g$, $Q$ and $\xi$ on $a\mathbb W$ is
  $0$ at $x$ unless $ax=x$ and the entries of $y(x)$ lie in $\hat{\mathcal O}$; then $\Phi(x)$ is the
  product of two factors. The first is $1$ if at every $P\in P_g$ the rows of $y_P(x)$ below the $r$-th
  vanish and its top $r\times r$ block is $1$, and $0$ otherwise. The second is
  $\prod_{P\in Q}\xi_P(\det t_P)^{-1}$ if at every $P\in Q$ the rows of $y_P(x)$ below the $r$-th vanish
  and its top $r\times r$ block $t_P$ is invertible, and $0$ otherwise. Let $g\in U(S)(\mathbf A)$ and
  $h\in U(T)(\mathbf A)$ satisfy $\beta_S^{-1}g_f\beta_S\in GL_{\operatorname{rk}S}(\hat{\mathcal O})$ and
  $\beta_T^{-1}h_f\beta_T\in GL_{\operatorname{rk}T}(\hat{\mathcal O})$, such that $\bar g_P$ is block
  upper triangular for the blocks $(r,\operatorname{rk}S-r)$ for every $P\in P_g\cup Q$, and its upper
  left $r\times r$ block $\bar g^{\mathrm{ul}}_P$ is $\bar h_P$ for every $P\in P_g$. Then for every
  $x\in\mathbb W$
  \[
   \Phi(hxg^{-1})=\prod_{P\in Q}\xi_P(\det\bar g^{\mathrm{ul}}_P)^{-1}\xi_P(\det\bar h_P)\,\Phi(x).
  \]
\item (Moving place models.) Let $\psi$ be a character of $\mathbf A_{F_h}$, $\Lambda_S$ and
  $\Lambda_T$ self-dual lattices of $S$ and $T$, $\Lambda\subset\mathbb W$ the lattice of the pair that
  they give, and $\omega$ a Weil datum for $(S,T)$ with characters $(\chi_S,\chi_T)$. Let $a$ be an
  idempotent of $\mathbf A_{E_h}$ with archimedean component $0$, $f_0$ a theta function, $\Phi$ a
  function on $\mathbb W$, and $f$ a function with $f=T(\Phi)$ for every place model $T$ of $f_0$ at
  $a$, of which there is at least one; a place model at $a$ is defined as in
  Lemma~\ref{is-eigenvalue-conjugate-theta-facts}, with the parts $a\mathbb W$ and $\bar a\mathbb W$
  for the parts of $\mathbb W$ at the places and at their conjugates. Let
  $(g,h)\in U(S)(\mathbf A)\times U(T)(\mathbf A)$ satisfy $hug^{-1}\in\Lambda$ if and only if
  $u\in\Lambda$, for every $u\in\mathbb W$, and let $c,d\in\mathbf C$ with $c\ne0$,
  $\omega(g,h)f_0=cf_0$ and $\Phi(hxg^{-1})=d\,\Phi(x)$ for every $x\in\mathbb W$. Then
  $d\,\omega(g,h)f=cf$.
\end{enumerate}
\end{lemma}
\leanref{Development/AbelianVarietyTheory/Proofs/TwistedOperators/WeilLatticeVacuumStabilizerSplittingCharactersStatement.lean}{GlobalTheory.T.weil_latticeVacuum_stabilizer_splitting_characters}
\leanref{Development/AbelianVarietyTheory/Proofs/TwistedOperators/AdjointHalfDatumStabilizerPairMoveStatement.lean}{GlobalTheory.T.adjointHalfDatum_stabilizerPair_move}
\leanref{Development/AbelianVarietyTheory/Proofs/TwistedOperators/WeilActionSchrodingerModelFiniteIdempotentStatement.lean}{GlobalTheory.T.weilAction_schrodingerModel_finiteIdempotent}
\begin{proof}
(2) Let $y(x)$ and $y_P(x)$ be as in the statement. As $g$ and $h$ are unitary,
$G_S^{-1}\overline{g^{-1}}^{\mathsf T}=gG_S^{-1}$ and $\bar h^{\mathsf T}G_T=G_Th^{-1}$, so
$y(hxg^{-1})=(\beta_S^{-1}g_f\beta_S)\,y(x)\,(\beta_T^{-1}h_f\beta_T)^{-1}$. As $a$ is a scalar,
$a\,hxg^{-1}=h(ax)g^{-1}$, so $hxg^{-1}$ satisfies $a\,hxg^{-1}=hxg^{-1}$ exactly when $ax=x$. As
$\beta_S^{-1}g_f\beta_S$ and $\beta_T^{-1}h_f\beta_T$ lie in the groups $GL(\hat{\mathcal O})$, the
entries of $y(hxg^{-1})$ lie in $\hat{\mathcal O}$ exactly when those of $y(x)$ do, and then
$y_P(hxg^{-1})=\bar g_Py_P(x)\bar h_P^{-1}$. At a place $P\in
P_g\cup Q$, $\bar g_P$ is block upper triangular, so the lower rows of $\bar g_Py_P(x)\bar h_P^{-1}$
vanish exactly when those of $y_P(x)$ do, and then its top block is $\bar g^{\mathrm{ul}}_Pt\bar h_P^{-1}$,
where $t$ is the top block of $y_P(x)$. For $P\in P_g$ this is $\bar h_Pt\bar h_P^{-1}$, which is $1$
exactly when $t=1$; so the first factor does not change. For $P\in Q$ its determinant is
$\det\bar g^{\mathrm{ul}}_P\det t\,(\det\bar h_P)^{-1}$, so the second factor is multiplied by
$\prod_{P\in Q}\xi_P(\det\bar g^{\mathrm{ul}}_P)^{-1}\xi_P(\det\bar h_P)$.

(3) Let $T$ be a place model of $f_0$ at $a$. Its values on the Schwartz-Bruhat functions on $a\mathbb
W$ are theta functions. Indeed, $T$ carries the indicator function of $\Lambda\cap a\mathbb W$ to the
theta function $f_0$; these Schwartz-Bruhat functions are finite linear combinations of functions
obtained from that indicator function by translations by elements of $a\mathbb W$ and by
multiplications by $t\mapsto\psi(\langle t,y\rangle)$ with $y\in\bar a\mathbb W$; $T$ carries these
operations to the operators $\rho(x)$ and $\rho(y)$; and these preserve theta functions
(Lemma~\ref{is-eigenvalue-conjugate-theta-closed}(1)).

For a function $\Psi$ on $\mathbb W$ put $\Psi^{g,h}(t)=\Psi(htg^{-1})$, and let
$T'(\Psi)=c^{-1}\omega(g,h)T(\Psi^{g,h})$. The map $t\mapsto htg^{-1}$ preserves $a\mathbb W$, $\bar a\mathbb
W$, the pairing and $\Lambda$, so $\Psi^{g,h}$ is a Schwartz-Bruhat function on $a\mathbb W$ if $\Psi$ is.
For $x\in a\mathbb W$ the translate of $\Psi$ by $x$ goes to the translate of $\Psi^{g,h}$ by
$x'=h^{-1}xg$, and on theta functions $\omega(g,h)\rho(x')=\rho(x)\omega(g,h)$ (Definition~\ref{th-weil-datum}(3));
so $T'$ carries the translation by $x$ to $\rho(x)$. In the same way it carries the multiplication by
$t\mapsto\psi(\langle t,y\rangle)$, $y\in\bar a\mathbb W$, to $\rho(y)$. The indicator function of
$\Lambda\cap a\mathbb W$ is invariant under $(g,h)$, so $T'$ carries it to $c^{-1}\omega(g,h)f_0=f_0$.
Hence $T'$ is a place model of $f_0$ at $a$, and
\[
 f=T'(\Phi)=c^{-1}\omega(g,h)T(\Phi^{g,h})=c^{-1}\omega(g,h)T(d\,\Phi)=c^{-1}d\,\omega(g,h)f.
\]
\end{proof}
% (1) is VecInv.Vac.stab_vacuum: the pairs (g, h) with om (g, h) f = splitFacO chiS chiT (detAt e g) (detAt e h) f
% form a closed subgroup (eigSub, closed by om_cont and the continuity of the characters), and
% T.mem_closedSubgroup_of_localStabilizerPairs (StabDense.mem_of_local) reduces to stabilizer pairs supported
% at one place v. At a place of P (or of its conjugates) and at a split place outside it this is
% T.splitPlace_weilRep_latticeVacuum_character (VecInv.om_latVac_split), with detAt e = detAt e' by
% idempotentInterpolation; at a nonsplit place outside P both determinants are 1 and it is
% T.weilOperator_lattice_stabiliser_fixes_vacuum_nonsplit (VecInvInert.om_fixes_vacuum), the analogue of
% Lemma th-nonsplit-vacua, which uses the Hecke character eta-tilde (IsEtaTilde). The hypotheses are IsFrame
% beta Lambda, IsIdemSum e P, UnrOff (e + cA e) chiS and chiT, Stab S betaS g, Stab T betaT h, IsThetaFunO and
% IsLatVac. (2) is AdjDatumMove.adjDatum_actO (adjDatum: frMapAdj, omegaK at Pg, omegaQ at Q with exponent -1;
% red_actO, omegaK_mul, omegaQ_mul). Its hypotheses are only betaS, betaT, Pg, Q, xi, an arbitrary a : AE D
% (IsHalf a x is a * x = x), r <= rk S, r <= rk T with rk T = r, finE g in St betaS, finE h in St betaT and the
% block conditions: no psi, lattices, IsFrame or Weil datum. (3) is SchrMove.smul_om_eq, with
% SchrTheta.isTheta_placeModel (proved by placeModel_induction with LatFinMod.latO_finite_mod) and
% isPlaceModel_move; its hypotheses are IsWeilDatumO, IsSelfDualC for both lattices, a * a = a, infPart a = 0,
% the lattice condition, IsThetaFunO f0, IsThetaAt at a with c = 1, c /= 0 and the two moves: no adelic
% matrices and no eta-tilde.

\begin{lemma}\label{is-eigenvalue-conjugate-frame-level-vector}
Let $D$, $\psi$, $S$, $T$, $\Lambda_S$, $\Lambda_T$, $\Lambda$, $\beta_S$, $\beta_T$, $\omega$, $\chi_S$
and $\chi_T$ be as in Lemma~\ref{is-eigenvalue-conjugate-frame-vacuum}(1), and suppose that $\eta$
extends to a Hecke character as in (1) there. Let $\theta$ be an archimedean family, $s_S$ and $s_T$
Cartan involutions of $S$ and $T$, $P_g$ and $Q$ finite sets of places of $E_h$, $\xi_P$ characters of
the $k_P^\times$, and $f$ a level vector with vacuum $f_0$ for these data
(Lemma~\ref{is-eigenvalue-conjugate-theta-facts}). Then there is a sum $e$ of the idempotents of
$\mathbf A_{E_h}$ at the places of $P_g\cup Q$, with $e\bar e=0$, with the following property. Suppose
that $\chi_S$ and $\chi_T$ are trivial on the ideles of
Lemma~\ref{is-eigenvalue-conjugate-frame-vacuum}(1) for $P=P_g\cup Q$, let $r\le\operatorname{rk}S$ with
$\operatorname{rk}T=r$, and let $(g,h)$ be a stabilizer pair that satisfies the conditions of
Lemma~\ref{is-eigenvalue-conjugate-frame-vacuum}(2). Then
\[
 \prod_{P\in Q}\xi_P(\det\bar g^{\mathrm{ul}}_P)^{-1}\xi_P(\det\bar h_P)\cdot\omega(g,h)f
 =\chi_T(\det_eg)^{-1}\chi_S(\det_eh)\,f.
\]
\end{lemma}
\leanref{Development/AbelianVarietyTheory/Proofs/TwistedOperators/WeilActionLiftPairStepVectorInvLevelChar.lean}{GlobalTheory.VecInv.om_vec}
\begin{proof}
By the definition of a level vector, $f_0$ is a Fock vacuum vector, hence a theta function and a lattice
vacuum for $\Lambda$, and there is a sum $e$ of the idempotents at the places of $P_g\cup Q$, with
$e\bar e=0$, such that $f=T(\Phi)$ for every place model $T$ of $f_0$ at $\bar e$, of which there is
one, where $\Phi$ is the level function on $\bar e\mathbb W$. The idempotent $\bar e$ has archimedean
component $0$. Let $(g,h)$ be as in the statement. As the finite parts of $g$ and $h$ stabilize
$\hat\Lambda_S$ and $\hat\Lambda_T$, $(g,h)$ preserves $\Lambda$. By
Lemma~\ref{is-eigenvalue-conjugate-frame-vacuum}(1), $\omega(g,h)f_0=cf_0$ with
$c=\chi_T(\det_eg)^{-1}\chi_S(\det_eh)$, and $c\ne0$ as $|c|=1$. By (2) of that lemma,
$\Phi(hxg^{-1})=d\,\Phi(x)$, where $d$ is the product on the left of the assertion. By (3),
$d\,\omega(g,h)f=cf$.
\end{proof}
% om_vec: hvac (IsVacO, with .theta and .lat) and e, he, hth from IsThPVecOf; the place model is at cA D e
% (idemSum_infPart, infPart_cA), the lattice condition is mem_latO_iff frS frT, h0 is
% weil_latticeVacuum_stabilizer_splitting_characters, and the last step is
% weilAction_schrodingerModel_finiteIdempotent with adjointHalfDatum_stabilizerPair_move.

\begin{lemma}\label{is-eigenvalue-conjugate-frame-level-character}
Let $R$ be a realization of a definite-side datum $D$ that satisfies the core conditions on its level
characters, with the notation of Definition~\ref{is-eigenvalue-lift-data}; let $i$ be a lift and $t$ a
tuple with $Q_t=\emptyset$, and let $e$ be the sum of the idempotents of $\mathbf A_{E_h}$ at the places
of $P_g$, with $e\bar e=0$. Let $k\in U(V_i)(\mathbf A)$ have archimedean component $1$ and finite part
$k_f$ in the level $J_t$ of $V_i$ at $t$, and let $\chi_{J_t}$ be the level character of $R$ on $J_t$.
For $P\in P_g$ let $d_P$ be the lower right block of size $\operatorname{rk}V_i-r$ of the reduction
$\bar k_P$ of $k_f$ at $P$, for the adelic matrix of $V_i$ (Lemma~\ref{is-eigenvalue-conjugate-frame-vacuum}).
Then
\[
 \mu^0_i\xi_t(\det_ek)=\chi_{J_t}(k_f)=\prod_{P\in P_g}\chi_{i+1,P}(\det d_P)^{-1}.
\]
\end{lemma}
\leanref{Development/AbelianVarietyTheory/Proofs/DivisorProofs/ConjugateWeilRepresentationSecond.lean}{GlobalTheory.GalOne.mu_detAt_eq_levChi}
\leanref{Formalism/LocalToGlobal/IndefiniteReal/SubgroupCharExtensionOneApply.lean}{Formalism.IndefR.Lev.levFn_flag1}
\begin{proof}
As $k_f\in J_t$, $k_f$ stabilizes the lattice, and for $P\in P_g$ the reduction $\bar k_P$ is block upper
triangular for the blocks $(r,\operatorname{rk}V_i-r)$ with upper left block $1$; so $\det\bar k_P=\det
d_P$. The determinant of $k_f$ is a unit at every finite place, with residue $\det\bar k_P$ at $P$, and
$\det_ek$ is the idele with this determinant at the places of $P_g$ and component $1$ elsewhere. By the
core conditions on the level characters, the restriction of $\mu^0_i$ to the units at $P\in P_g$ is the
character $\chi_{i+1,P}^{-1}$ of the residue, and $\xi_t$ is unramified at the places of $P_g$. So
$\mu^0_i\xi_t(\det_ek)=\prod_{P\in P_g}\chi_{i+1,P}(\det d_P)^{-1}$. The level character of $R$ at an
element of $J_t$ whose reductions at $P_g$ are block upper triangular is the product of this same
factor, of factors at the places of $Q_t$, which is empty, and of a factor of the local level at the
upper left blocks. This last factor is $1$ for the trivial local level, which defines $\chi_{J_t}$, and
for the full local level with the trivial character, for which the product is computed. So
$\chi_{J_t}(k_f)$ is the same product.
\end{proof}
% mu_detAt_eq_levChi: exists_det_int (aV and hrV), mu_val (from tameHeckeChar_splitSupported_eq_prod_residue for
% mu0 with hL.chi, for xi and for eta), levChi_flag_V and T.det_twoDiagBlocks_properties. The level character
% (levEnv V beta (datVt i t)).chi is levFn for flag1L; levFn_flag1 identifies it with levFn for flagL, the
% value of levChi in mu_detAt_eq_levChi (both ext1 factors are 1 by subgroupCharExtension_one_apply). The
% hypothesis on k is IsLevelLift (Vc D i) (betaV i) (datVt i t) 1 k.

\begin{definition}\label{is-eigenvalue-conjugate-frame-fibres}
Let $X$ be flag data of rank $r$ over a CM field $E$ (Definition~\ref{is-onto-share-flag-data}), with set
$\Psi$, matrix $A$, integer $p_0$, module $\Lambda$, element $\kappa_0$ and forms $h$ and $h_\tau$, and
put $d=|\Psi|r$. Let $x=(x_\tau)_{\tau\in\Psi}$ be a family of matrices in $M_r(\mathbf C)$, and write
$e_a$ for the standard basis of $\mathbf C^r$.
\begin{enumerate}
\item (Coordinates and periods.) An \emph{indexing} is a bijection $j\mapsto(\tau_j,a_j)$ from
  $\{1,\dots,d\}$ onto $\Psi\times\{1,\dots,r\}$; its \emph{row type} at $j$ is $\rho_j=\bar\tau_j$ if
  $a_j\le p_0$ and $\rho_j=\tau_j$ if $a_j>p_0$. For a family $u=(u_\tau)_{\tau\in\Psi}$ of vectors in
  $\mathbf C^r$, the \emph{coordinates} of $u$ at $x$ are the vector $z_x(u)\in\mathbf C^d$ with
  \[
   z_x(u)_j=h_{\tau_j}(x_{\tau_j}e_{a_j},u_{\tau_j})\quad\text{if }a_j\le p_0,\qquad
   z_x(u)_j=(x_{\tau_j}^{-1}u_{\tau_j})_{a_j}\quad\text{if }a_j>p_0.
  \]
  The \emph{periods} of $v\in E^r$ at $x$ are $p_x(v)=z_x\big((\tau(v))_{\tau\in\Psi}\big)$. They are
  additive in $v$, and $p_x(\alpha v)_j=\rho_j(\alpha)p_x(v)_j$ for $\alpha\in E$.
\item (Riemann forms and Gaussians.) A \emph{symplectic basis} of $\Lambda$ is a $\mathbf Z$-basis
  $b_1,\dots,b_d,b'_1,\dots,b'_d$ of $\Lambda$ such that the form
  $\langle u,v\rangle_X=\operatorname{Tr}_{E/\mathbf Q}\big(-\tfrac12\kappa_0h(u,v)\big)$ satisfies
  $\langle b_k,b'_l\rangle_X=-\delta_{kl}$, $\langle b'_l,b_k\rangle_X=\delta_{kl}$ and
  $\langle b_k,b_l\rangle_X=\langle b'_k,b'_l\rangle_X=0$. For such a basis and an indexing, let
  $\Omega_x,\Omega'_x\in M_d(\mathbf C)$ be the matrices whose $k$-th columns are $p_x(b_k)$ and
  $p_x(b'_k)$, put $\tau_x=\Omega_x^{-1}\Omega'_x$, and let
  \[
   H_x(z,w)=(\Omega_x^{-1}z)^{\mathsf T}\,(\operatorname{Im}\tau_x)^{-1}\,\overline{\Omega_x^{-1}w}
   \qquad(z,w\in\mathbf C^d),
  \]
  with the imaginary part taken entry by entry. The \emph{Gaussian} at $x$ is the function
  \[
   \mathcal G_x(z)=\exp\Big(-\frac\pi2\sum_{k=1}^nH_x(z_k,z_k)\Big)
  \]
  of $z=(z_1,\dots,z_n)\in(\mathbf C^d)^n$, for each $n\ge1$.
\item (Analytic fibres.) Let $y$ be a complex framed point of the moduli scheme of a PEL datum over $E$ at
  a level (Definition~\ref{cv-model-toroidal-semiabelian}(1)), such as the moduli of a core at some level
  (Definition~\ref{is-one-map-data}), for an embedding $s$ of its coefficient field $K_s$. It has a fibre
  $A_y$ of the universal abelian scheme, with an action of $\mathcal O_E$; the line bundle $\mathcal L$
  of the polarization, with trivializing opens $U_\beta$, transition functions $g_{\beta\beta'}$ and the
  rigidifying section at the origin; a frame $(u_\ell)$ of the tangent space of $A_y$ at the origin,
  indexed by the \emph{letters} $\ell=(\tau',i)$ of the datum, with $\tau'$ an embedding
  of $E$ and $1\le i\le m_{\tau'}$ for the multiplicities $m_{\tau'}$ of the datum, $u_{\tau',i}$ in the
  eigenspace of $\mathcal O_E$ for $s\circ\varepsilon_{\tau'}$, where $\varepsilon_{\tau'}$ is the
  embedding of $E$ in $K_s$ attached to $\tau'$; and a section $s_y$ of $\mathcal L$ on $A_y$ that
  generates its sections. Let a symplectic basis and an indexing be given, a bijection $j\mapsto\ell_j$
  from $\{1,\dots,d\}$ onto the letters, and an $\mathcal O_E$-submodule $\Lambda'\subset E^r$. An
  \emph{analytic fibre} at $y$ for these data and $x$ consists of:
  \begin{enumerate}
  \item[(a)] a surjective holomorphic group homomorphism $\operatorname{ex}\colon\mathbf C^d\to A_y(\mathbf C)$
    whose kernel is $\{p_x(v):v\in\Lambda'\}$, such that
    $\operatorname{ex}\big((\rho_j(\alpha)z_j)_j\big)=\alpha\cdot\operatorname{ex}(z)$ for
    $\alpha\in\mathcal O_E$, and whose derivative at $0$ in the direction of the $j$-th standard basis
    vector is $u_{\ell_j}$;
  \item[(b)] for each $\beta$, a holomorphic function $\operatorname{fr}_\beta$ on the set of $z$ with
    $\operatorname{ex}(z)\in U_\beta$, nowhere zero there, with
    $\operatorname{fr}_\beta(z)=g_{\beta\beta'}(\operatorname{ex}(z))\operatorname{fr}_{\beta'}(z)$ when
    $\operatorname{ex}(z)\in U_\beta\cap U_{\beta'}$, and with $\operatorname{fr}_\beta(0)$ the value in
    $U_\beta$ of the rigidifying section when $0\in\operatorname{ex}^{-1}(U_\beta)$;
  \item[(c)] a function $\chi$ on $E^r$ with $|\chi(v)|=1$ for $v\in\Lambda'$ such that
    \[
     \operatorname{fr}_\beta(z)=\chi(v)\exp\Big(\pi H_x(z,p_x(v))+\frac\pi2H_x(p_x(v),p_x(v))\Big)
     \operatorname{fr}_\beta(z+p_x(v))
    \]
    for every $v\in\Lambda'$, every $\beta$ and every $z$ with $\operatorname{ex}(z)\in U_\beta$.
  \end{enumerate}
\item (Theta readings.) A \emph{theta reading} of a section $F$ of $\mathcal L^k$ on $A_y$, for an
  analytic fibre at $y$, is a function $\vartheta$ on $\mathbf C^d$ such that the value of $F$ at
  $\operatorname{ex}(z)$ in the trivialization over $U_\beta$ is $\vartheta(z)\operatorname{fr}_\beta(z)^k$,
  for every $\beta$ and every $z$ with $\operatorname{ex}(z)\in U_\beta$.
\item (Matched indexings.) An indexing and a bijection $j\mapsto\ell_j$ onto the letters are
  \emph{matched} for $s$ if, for every $j$, with $\ell_j=(\tau',i)$ and $\tau=s\circ\varepsilon_{\tau'}$,
  the row type $\rho_j$ is $\tau$, and $p_x(v)_j$ is the \emph{period of $v$ at the letter} $\ell_j$ for
  every family $x$ of invertible matrices and every $v\in E^r$: $(x_\tau^{-1}\tau(v))_{p_0+i}$ if
  $\tau\in\Psi$, and $h_{\bar\tau}(x_{\bar\tau}e_{p_0+1-i},\bar\tau(v))$ if $\bar\tau\in\Psi$.
\item (Theta series.) A function $\varphi$ on $M_{n\times d}(\mathbf Q)$ is \emph{Schwartz-Bruhat} if
  there is an integer $N\ge1$ such that $\varphi(m)\ne0$ only if $Nm$ has integer entries, and
  $\varphi(m+Nm')=\varphi(m)$ for every $m'$ with integer entries. For such $\varphi$, a symplectic
  basis and an indexing, the \emph{theta series} of $\varphi$ at $x$ is
  \[
   \theta_\varphi(x;z)=\det(\Omega_x^{-1})^{\lfloor n/2\rfloor}
   \exp\Big(\frac\pi2\sum_{k=1}^n\zeta_k^{\mathsf T}(\operatorname{Im}\tau_x)^{-1}\zeta_k\Big)
   \sum_{m\in M_{n\times d}(\mathbf Q)}\varphi(m)\exp\Big(\pi i\sum_{k=1}^n\big(m_k^{\mathsf T}\tau_xm_k
   +2m_k^{\mathsf T}\zeta_k\big)\Big)
  \]
  for $z\in(\mathbf C^d)^n$, where $\zeta_k=\Omega_x^{-1}z_k$ and $m_k$ is the $k$-th row of $m$; and
  for families $u_1,\dots,u_n$ as in (1) its \emph{theta function} is
  $\Theta_\varphi(x;u)=\mathcal G_x(z)\,\theta_\varphi(x;z)$ with $z=(z_x(u_1),\dots,z_x(u_n))$.
\item (Mean one.) An $\mathcal O_E$-submodule $\Lambda'\subset E^r$ is \emph{self-dual} for $X$ if it
  is finitely generated and $v\in\Lambda'$ exactly when $\langle v,w\rangle_X\in\mathbf Z$ for every
  $w\in\Lambda'$. A function $\vartheta$ on $\mathbf C^d$ has \emph{mean one} for $x$ and $\Lambda'$, for a
  symplectic basis and an indexing, if there is a fundamental domain $\mathcal F$ of the group
  $\{p_x(v):v\in\Lambda'\}$ in $\mathbf C^d$, for Lebesgue measure, with
  \[
   \frac1{\operatorname{vol}(\mathcal F)}\int_{\mathcal F}|\vartheta(z)|^2e^{-\pi\operatorname{Re}H_x(z,z)}\,dz=1.
  \]
\end{enumerate}
\end{definition}
\leanref{Development/AbelianVarietyTheory/Definitions/ThetaValues/PolyDerivSmulLetter.lean}{GlobalTheory.ThFr.per}
\leanref{Development/AbelianVarietyTheory/Definitions/ThetaValues/PolyDerivSmulLetter.lean}{GlobalTheory.ThFr.rowType}
\leanref{Development/AbelianVarietyTheory/Definitions/SpanningFamilies/IsStabForm.lean}{GlobalTheory.Dict.pv}
\leanref{Development/AbelianVarietyTheory/Definitions/SpanningFamilies/Fischer.lean}{GlobalTheory.Ker.zR}
\leanref{Development/AbelianVarietyTheory/Definitions/ThetaValues/PolyDerivSmulLetter.lean}{GlobalTheory.ThFr.IsSympl}
\leanref{Development/AbelianVarietyTheory/Definitions/ThetaValues/PolyDerivSmulLetter.lean}{GlobalTheory.ThFr.tauMat}
\leanref{Development/AbelianVarietyTheory/Definitions/ThetaValues/PolyDerivSmulLetter.lean}{GlobalTheory.ThFr.imInv}
\leanref{Development/AbelianVarietyTheory/Definitions/SpanningFamilies/IsStabForm.lean}{GlobalTheory.Dict.hx}
\leanref{Development/AbelianVarietyTheory/Definitions/SpanningFamilies/Fischer.lean}{GlobalTheory.Ker.gaussR}
\leanref{Development/AbelianVarietyTheory/Definitions/SpanningFamilies/IsStabForm.lean}{GlobalTheory.Dict.AnFib}
\leanref{Development/AbelianVarietyTheory/Definitions/SpanningFamilies/IsStabForm.lean}{GlobalTheory.Dict.AnFib.Reads}
\leanref{Development/AbelianVarietyTheory/Proofs/IwasawaDecomposition/FibExtMarked.lean}{GlobalTheory.OneMapRef.Matched}
\leanref{Formalism/LocalToGlobal/ProjectiveLattices/PetersonNormSubspace.lean}{Formalism.L2.Pj.perFr}
\leanref{Formalism/LocalToGlobal/ProjectiveLattices/IsCombAt.lean}{Formalism.L2.Pj.FrPt}
\leanref{Development/AbelianVarietyTheory/Definitions/SpanningFamilies/IMemSiegel.lean}{GlobalTheory.SW.IsSB}
\leanref{Development/AbelianVarietyTheory/Definitions/SpanningFamilies/IMemSiegel.lean}{GlobalTheory.SW.HasLevel}
\leanref{Development/AbelianVarietyTheory/Definitions/SpanningFamilies/IMemSiegel.lean}{GlobalTheory.SW.thS}
\leanref{Development/AbelianVarietyTheory/Definitions/SpanningFamilies/IsSystem.lean}{GlobalTheory.SW.thA}
\leanref{Development/AbelianVarietyTheory/Definitions/SpanningFamilies/Fischer.lean}{GlobalTheory.Ker.thR}
\leanref{Development/AbelianVarietyTheory/Definitions/SpanningFamilies/IsSystem.lean}{GlobalTheory.SW.SelfD}
\leanref{Development/AbelianVarietyTheory/Definitions/SpanningFamilies/IsStabForm.lean}{GlobalTheory.Dict.MeanOne}
% (6) is SW.IsSB (HasLevel R with R > 0), SW.thA (det (hMat)^(n/2), natural division, the exponential of
% (pi/2) sum (hMat z_k) . imInv (hMat z_k), and SW.thS phi tauMat (hMat z_k), a tsum over Q-matrices) and
% Ker.thR = gaussR * thA at zR. (7) is SW.SelfD (FG, and EF-integrality characterizes membership) and
% Dict.MeanOne (a fundamental domain Fd of the period group for volume, with the normalized integral 1).
% (1) is ThFr.per (FD.hermC tau Gm (col (mat x tau) a) (emb tau v) for a < p0, ((mat x tau)^-1 v)_a otherwise),
% Dict.pv = per o idx, and Ker.zR, the same formula for complex vectors (v tau k); rowType is conj tau for
% a < p0 and tau otherwise, and per_add, per_smul give the last sentence. (2) is IsSympl (EF (bas a) (bas b)
% = -omegaStd a b, EF u v = Tr(-(kappa0/2) herm Gm u v)), perMat with Sum.inl and Sum.inr, hMat = (perMat
% inl)^-1, tauMat = hMat * perMat inr, imInv = ((tauMat).map im)^-1, Dict.hx z w = (hMat z) . imInv (star
% (hMat w)) and Ker.gaussR. Inverses are Lean's nonsingular inverses (0 for a singular matrix). (3) is
% Dict.AnFib X bas idx e x L y: ex (ex_add, ex_onto, ex_ker, ex_lin, hol, der with y.u), fr (fr_hol, fr_ne,
% fr_g, fr_0 with the rigidification Xo.univ.t.PP.rig.s), chi with chi_norm, and fr_per. The framed point
% is Formalism.L2.Pj.FrPt (x, u with tan, eig and basis, and s with gen); the letters are Lt1 Dw, pairs of an
% embedding and an index below Dw.mult. (4) is AnFib.Reads k F theta. (5) is OneMapRef.Matched c sigma X idx e:
% pv X idx x v j = perFr X x v (relab c sigma (e j).1) (e j).2 and rowType X (idx j) = relab c sigma (e j).1,
% with relab c sigma tau = sigma o embF c tau; perFr is 0 when neither tau nor its conjugate is in Psi or
% the index is out of range.

\begin{situation}\label{is-eigenvalue-conjugate-frame-situation}
Let $D$ be a definite-side datum with host field $E_h/F_h$, element $\xi_0$ and set $\mathcal T$ of
embeddings of $E_h$ (Definition~\ref{th-definite-datum}), let $\psi$ be a standard character of
$\mathbf A_{F_h}$ (Definition~\ref{th-weil-datum}(1)), and let $\theta$ be an archimedean family whose
member $\theta_\rho$ extends $\rho$ for every $\rho\in\mathcal T$. Let $S$ and $T$ be Hermitian spaces of
$D$, where $S$ has Gram matrix $1$ and even rank $n$; put $r=\operatorname{rk}T$ and
$G_T(\rho)=\rho(G_T)$. Let $X$ be flag data of rank $r$ over $E_h$
(Definition~\ref{is-onto-share-flag-data}) with set $\Psi=\mathcal T$, matrix $-G_T^{\mathsf T}$,
element $-\xi_0$, $\kappa_0=-2\xi_0$, integer $p_0$ and module $\Lambda$, together with a symplectic
basis of $\Lambda$ and an indexing (Definition~\ref{is-eigenvalue-conjugate-frame-fibres}(1)
and~(2)); put $d=|\mathcal T|r$. Coordinates $z_x$, Gaussians $\mathcal G_x$ and theta functions
$\Theta_\varphi(x;u)$ (Definition~\ref{is-eigenvalue-conjugate-frame-fibres}) are formed with this
basis and indexing.
\begin{enumerate}
\item (Signature families.) Let $N$ be the diagonal $r\times r$ matrix with entries $-1$ at the first
  $p_0$ places and $1$ at the others. A family $Z=(Z_\rho)_{\rho\in\mathcal T}$ of complex $r\times r$
  matrices is a \emph{signature family} if $Z_\rho^*G_T(\rho)Z_\rho=N$ for every $\rho$. For such $Z$
  put $s_{Z,\rho}=Z_\rho NZ_\rho^{-1}$, and write $x^Z$ for $Z$ read as a point of
  $M_r(\mathbf C)^\Psi$. For $g\in U(T)(\mathbf A)$ write $g_\rho=\theta_\rho(g)$ and
  $gZ=(g_\rho Z_\rho)_\rho$.
\item (Real points.) For $w\in\mathbb W$ with finite component $0$ (Definition~\ref{th-weil-datum})
  let $u(w)=(u_1,\dots,u_n)$, where $(u_k)_\rho$ is the $k$-th column of $\theta_\rho(w)$.
\end{enumerate}
\end{situation}
\leanref{Development/AbelianVarietyTheory/Proofs/IwasawaDecomposition/FibExtMarked.lean}{GlobalTheory.OneMapRef.IsSet}
\leanref{Development/AbelianVarietyTheory/Proofs/IwasawaDecomposition/FibExtMarked.lean}{GlobalTheory.OneMapRef.IsONeg}
\leanref{Development/AbelianVarietyTheory/Proofs/IwasawaDecomposition/FibExtMarked.lean}{GlobalTheory.OneMapRef.cartZ}
\leanref{Development/AbelianVarietyTheory/Proofs/IwasawaDecomposition/FibExtMarked.lean}{GlobalTheory.OneMapRef.mvZ}
\leanref{Development/AbelianVarietyTheory/Proofs/IwasawaDecomposition/FibExtMarked.lean}{GlobalTheory.OneMapRef.frOf}
\leanref{Development/AbelianVarietyTheory/Proofs/IwasawaDecomposition/FibExtMarked.lean}{GlobalTheory.OneMapRef.rePt}
\leanref{Development/AbelianVarietyTheory/Proofs/IwasawaDecomposition/FibExtMarked.lean}{GlobalTheory.OneMapRef.Zak}
% IsSet psi theta S T hr X eT bas: IsStdChar, IsInfFam, S.gram = 1, Even S.rank, ModelW T hr X (matrix
% -G_T^T, element -xi0), kappa0 = -(2 xi0), (eT rho).1 = tauE rho (so Psi is T, and eT is determined) and
% IsSympl X bas. IsONeg T q Z is Z^H gramAt Z = negM T q; cartZ is s_Z; mvZ theta h Z is gZ; frOf is x^Z;
% rePt is u(w). Zak ... x phi f (f a theta function with f(w) = Theta_phi(x; u(w)) at the real points) is
% written out in each statement.

\begin{lemma}\label{is-eigenvalue-conjugate-frame-glue}
Let $D$ be a definite-side datum with host field $E_h/F_h$ and set $\mathcal T$ of embeddings of $E_h$,
let $\theta$ be an archimedean family whose member $\theta_\rho$ extends $\rho$ for every $\rho\in\mathcal
T$, and let $V$ be a Hermitian space of $D$ with Gram matrix $G_V$. Put $G_V(\rho)=\rho(G_V)$ and, for
$g\in U(V)(\mathbf A)$, $g_\rho=\theta_\rho(g)$, the image of the archimedean component of $g$.
\begin{enumerate}
\item If $g\in U(V)(\mathbf A)$ satisfies $g_\rho=1$ for every $\rho\in\mathcal T$, then $g$ has
  archimedean component $1$.
\item For every family $(u_\rho)_{\rho\in\mathcal T}$ of complex matrices with
  $u_\rho^*G_V(\rho)u_\rho=G_V(\rho)$ for every $\rho$ there is $g\in U(V)(\mathbf A)$ with finite
  component $1$ and $g_\rho=u_\rho$ for every $\rho$.
\end{enumerate}
\end{lemma}
\leanref{Development/AbelianVarietyTheory/Proofs/HeightProofs/AdelicUnitaryRealPointsStructureRealPlaces.lean}{GlobalTheory.SmoothForms.RealPts.glue}
\leanref{Development/AbelianVarietyTheory/Proofs/UnitaryExistence/AdelicUnitaryExistsWithPrescribedParts.lean}{GlobalTheory.UnitaryGlue.adelic_glue}
\begin{proof}
The set $\mathcal T$ contains exactly one of every two complex conjugate embeddings of $E_h$, and each
$\theta_\rho$ extends $\rho$ and is continuous on $\mathbf A_{F_h,\infty}$. So the map
$z\mapsto(\theta_\rho(z))_{\rho\in\mathcal T}$ is a bijective ring homomorphism from
$\mathbf A_{F_h,\infty}\otimes_{F_h}E_h$ onto $\mathbf C^{\mathcal T}$, and
$\theta_\rho(\bar z)=\overline{\theta_\rho(z)}$.

(1) The archimedean component of $g$ and the matrix $1$ have the same image under every $\theta_\rho$,
entry by entry, so they are equal.

(2) Each $u_\rho$ is invertible, as $\det G_V(\rho)\ne0$. By the bijectivity there are matrices $M$ and
$N$ over $\mathbf A_{F_h,\infty}\otimes_{F_h}E_h$ with $\theta_\rho(M)=u_\rho$ and $\theta_\rho(N)=u_\rho^{-1}$
for every $\rho$. Then $MN=NM=1$ and $\bar M^{\mathsf T}G_VM=G_V$, since both sides of each identity have
the same image under every $\theta_\rho$. The element $g$ of $U(V)(\mathbf A)$ with archimedean component
$M$ and finite component $1$ has $g_\rho=u_\rho$.
\end{proof}
% glue: adelic_glue for the finite set Psi = tauE(T) (tauE_cm: it contains tau exactly when it does not contain
% conj tau), with the bijection of GlobalTheory.InfAdeleCM.infAdele_tensor_bijective (hb) and its conjugation
% clause (hconj); (1) is the first clause of adelic_glue's second part (Units.ext through hb.1); (2) is its last
% clause: gram_det_ne_zero for the invertibility, exists_glue for M (with N for the inverse and the unitarity
% checked through inj), and exists_inf_fin with finite component 1.

\begin{lemma}\label{is-eigenvalue-conjugate-frame-signature}
In Situation~\ref{is-eigenvalue-conjugate-frame-situation}, let $Z$ be a signature family.
\begin{enumerate}
\item The point $x^Z$ lies in the flag domain $\mathcal D_X$ and is orthonormal, and
  $|d_i(x^Z,\rho)|=1$ for all $\rho$ and $i$ (Definition~\ref{is-onto-share-flag-data}).
\item For every $\rho$, $s_{Z,\rho}$ is a Cartan involution of $G_T(\rho)$: $s_{Z,\rho}^2=1$,
  $s_{Z,\rho}^*G_T(\rho)s_{Z,\rho}=G_T(\rho)$ and $G_T(\rho)s_{Z,\rho}$ is positive definite.
\item Let $g,g'\in U(T)(\mathbf A)$. Then $gZ$ is a signature family, $s_{gZ,\rho}g_\rho=g_\rho
  s_{Z,\rho}$ for every $\rho$, and $(g'g)Z=g'(gZ)$; and $gZ=Z$ if and only if $g$ has archimedean
  component $1$.
\item For every signature family $Z'$ there is $a\in U(T)(\mathbf A)$ with finite component $1$ and
  $Z'=aZ$.
\end{enumerate}
\end{lemma}
\leanref{Development/AbelianVarietyTheory/Proofs/IwasawaDecomposition/OrthonormalNegFirstFamilyFrameCartanMove.lean}{GlobalTheory.OneMapRef.oneg_frames}
\begin{proof}
Each $g_\rho$ is unitary for $G_T(\rho)$, $g_\rho^*G_T(\rho)g_\rho=G_T(\rho)$, as $g$ is unitary for
$G_T$ and $\theta_\rho$ extends $\rho$ and commutes with the conjugation. The matrix $N$ satisfies
$N^2=1$ and $N^*=N$, and $Z_\rho$ is invertible, as $\det N\ne0$.

(1) As the matrix of $X$ is $-G_T^{\mathsf T}$, the form $h_\rho$ of $X$ is $h_\rho(u,v)=-v^*G_T(\rho)u$.
So the Gram matrix of the columns of $Z_\rho$ for $h_\rho$ is $-(Z_\rho^*G_T(\rho)Z_\rho)^{\mathsf
T}=-N$, the diagonal matrix with entries $1$ at the first $p_0$ places and $-1$ at the others. This is
the orthonormality of $x^Z$. The leading minors of this matrix are the $\epsilon_i=\pm1$, so
$|d_i(x^Z,\rho)|=1$ and $\epsilon_i\operatorname{Re}d_i(x^Z,\rho)=1>0$; as $\det Z_\rho\ne0$,
$x^Z\in\mathcal D_X$.

(2) Write $s=s_{Z,\rho}$, $G=G_T(\rho)$ and $Z=Z_\rho$; then $G=(Z^*)^{-1}NZ^{-1}$. So
$s^2=ZN^2Z^{-1}=1$, $s^*Gs=(Z^*)^{-1}N(Z^*GZ)NZ^{-1}=(Z^*)^{-1}N^3Z^{-1}=G$, and
$Gs=(Z^*)^{-1}N^2Z^{-1}=(ZZ^*)^{-1}$, which is positive definite.

(3) First, $(g_\rho Z_\rho)^*G_T(\rho)(g_\rho Z_\rho)=Z_\rho^*G_T(\rho)Z_\rho=N$. Next,
$s_{gZ,\rho}g_\rho=g_\rho Z_\rho NZ_\rho^{-1}g_\rho^{-1}g_\rho=g_\rho s_{Z,\rho}$, and
$(g'g)_\rho=g'_\rho g_\rho$ gives $(g'g)Z=g'(gZ)$. If $g$ has archimedean component $1$, then
$g_\rho=1$ for every $\rho$ and $gZ=Z$. Conversely, if $gZ=Z$, then
$g_\rho=g_\rho Z_\rho Z_\rho^{-1}=Z_\rho Z_\rho^{-1}=1$ for every $\rho$, as $Z_\rho$ is invertible; so
$g$ has archimedean component $1$ by Lemma~\ref{is-eigenvalue-conjugate-frame-glue}(1).

(4) Put $v_\rho=Z'_\rho Z_\rho^{-1}$. Then $v_\rho^*G_T(\rho)v_\rho=(Z_\rho^*)^{-1}(Z'^*_\rho G_T(\rho)Z'_\rho)
Z_\rho^{-1}=(Z_\rho^*)^{-1}(Z_\rho^*G_T(\rho)Z_\rho)Z_\rho^{-1}=G_T(\rho)$. By
Lemma~\ref{is-eigenvalue-conjugate-frame-glue}(2) there is $a\in U(T)(\mathbf A)$ with finite component $1$
and $a_\rho=v_\rho$ for every $\rho$. So $(aZ)_\rho=Z'_\rho Z_\rho^{-1}Z_\rho=Z'_\rho$.
\end{proof}
% oneg_frames: isOrtho_frOf (ModelW and (eT rho).1 = tauE rho), mem_dom and isUnitFr_of_isOrtho for (1);
% isCartanInv with negM_sq and negM_herm for (2); isONeg_mvZ (compAt_unitary'), cartZ_mvZ, mvZ_mul, mvZ_fin
% and fin_of_mvZ for (3); exists_move for (4). fin_of_mvZ and exists_move use SmoothForms.RealPts.glue
% (Lemma is-eigenvalue-conjugate-frame-glue).

\begin{lemma}\label{is-eigenvalue-conjugate-frame-det-roots}
Let $D$ be a definite-side datum with host field $E_h/F_h$, and $V$ a Hermitian space of $D$. A
\emph{determinant root} of $g\in U(V)(\mathbf A)$ is an idele $y\in\mathbf A_{E_h}^\times$ with
$y=\det(g)\,\bar y$.
\begin{enumerate}
\item If $y$ and $y'$ are determinant roots of $g$ and $g'$, then $yy'$ is a determinant root of $gg'$.
\item Let $\Lambda$ be a self-dual lattice of $V$, and $k\in U(V)(\mathbf A)$ an element with archimedean
  component $1$ and $k_f\hat\Lambda=\hat\Lambda$. Then $k$ has a determinant root with archimedean
  component $1$ that is a unit at every finite place. Of the self-duality of $\Lambda$ only its being a
  lattice is used.
\end{enumerate}
\end{lemma}
\leanref{Formalism/LocalToGlobal/BesselHermitianTheta/TiedWeil.lean}{Formalism.L2.Theta.IsDetRoot}
\leanref{Formalism/LocalToGlobal/ProjectiveLattices/FamilyUniformisationBasePointOrthogonalBis.lean}{Formalism.L2.Pj.FinUnit}
\leanref{Formalism/LocalToGlobal/ProjectiveLattices/FamilyUniformisationBasePointOrthogonal.lean}{Formalism.L2.Pj.detRootMul_holds}
\leanref{Development/AbelianVarietyTheory/Proofs/IwasawaDecomposition/ExistsFinUnitIsDetRoot.lean}{GlobalTheory.LatticeUnitRoot.exists_finUnit_isDetRoot}
\leanref{Development/AbelianVarietyTheory/Proofs/UnitaryExistence/ExistsIntegralUnitRootNormOneAdeleStatement.lean}{GlobalTheory.T.exists_integral_unit_root_norm_one_adele}
\begin{proof}
(1) $yy'=\det(g)\bar y\det(g')\bar y'=\det(gg')\,\overline{yy'}$.

(2) As $k_f$ maps $\hat\Lambda$ onto itself and $\hat\Lambda$ is free over the completion of
$\mathcal O_{E_h}$ at every finite place of $F_h$, $u=\det k_f$ is integral at every finite place; and
$u\bar u=1$, as $k$ is unitary. The extension $E_h/F_h$ is unramified at the finite places
(Chapter~\ref{chapter-unitary-theta}). By Hilbert's Theorem~90 for the units of an unramified extension
of local fields \cite{Ser79}, applied at every finite place of $F_h$ (at a split place it is immediate),
there is a finite idele $x$, a unit at every finite place, with $x=u\bar x$. Let $y$ be the idele with
archimedean component $1$ and finite component $x$. Then $y=\det(k)\bar y$ at the finite places, and at
the archimedean places both sides are $1$, as $k$ has archimedean component $1$.
\end{proof}
% exists_finUnit_isDetRoot: MarkDet.det_sub_one_of_cong with M = 1 (integrality of det k_f through a pseudo-basis
% adapted to Lambda), RecDict.det_finHom_norm, OddRoot.exists_root_odd with N = 1 (finite_root_odd, whose unit
% root is T.exists_integral_unit_root_norm_one_adele = AdLat.exists_integral_unit_root, under the condition
% D.sp.cm.cmUnram that t - c(t) is not in P for some integer t, at every prime P), DetRootCong.exists_idele, and
% the two components of IsDetRoot (infPart_sigR_one, finPart_sigR). (1) is detRootMul_holds.

\begin{lemma}[Vacua under finite stabilizers \cite{Kud94}]\label{is-eigenvalue-conjugate-frame-finite-vacuum}
Let $D$ be a definite-side datum with host field $E_h/F_h$, let $\psi$ be a character of
$\mathbf A_{F_h}$ and $\theta$ an archimedean family whose member at each $\tau\in\mathcal T$ extends
$\tau$, and let $S$ and $T$ be Hermitian spaces of $D$ with correct signatures, Cartan involutions $s_S$
and $s_T$ at the $\tau$, and self-dual lattices $\Lambda_S$ and $\Lambda_T$. Let $\omega$ be a Weil datum
for $(S,T)$ with characters $(\chi_S,\chi_T)$. Suppose that the quadratic character $\eta$ of
$E_h/F_h$ is the restriction of a Hecke character of $E_h$ that is unramified at every finite place
and has archimedean exponents $\frac12$, and that $\chi_S$ and $\chi_T$ are trivial on every idele that
is a unit at every finite place and differs from $1$ only at the places above one place of $F_h$ that
does not split in $E_h$. Let $f$ be a Fock vacuum vector for $s_S$, $s_T$ and the lattice of the pair
given by $\Lambda_S$ and $\Lambda_T$. Let $k\in U(S)(\mathbf A)$ and $k'\in U(T)(\mathbf A)$ have
archimedean components $1$, with $k_f\hat\Lambda_S=\hat\Lambda_S$ and $k'_f\hat\Lambda_T=\hat\Lambda_T$,
and let $x$ and $y$ be determinant roots of $k$ and $k'$
(Lemma~\ref{is-eigenvalue-conjugate-frame-det-roots}) with archimedean components $1$ that are units at
every finite place. Then
\[
 \omega(k,k')f=\chi_T(x)^{-1}\chi_S(y)\,f.
\]
\end{lemma}
\leanref{Development/AbelianVarietyTheory/Proofs/IwasawaDecomposition/WeilVacuumFiniteSelfDualLatticeSplittingStatement.lean}{GlobalTheory.T.weil_vacuum_finite_selfDualLattice_splitting}
\leanref{Development/AbelianVarietyTheory/Proofs/IwasawaDecomposition/WeilVacuumFiniteSelfDualLatticeSplitting.lean}{GlobalTheory.FinVac.om_vacuum_finite}

\begin{definition}\label{is-eigenvalue-conjugate-frame-heisenberg}
Let $X$ be data (1) to (6) of flag data of rank $r$ over a CM field $E$
(Definition~\ref{is-onto-share-flag-data}), with set $\Psi$ and module $\Lambda$, a symplectic basis
$b_1,\dots,b_d,b'_1,\dots,b'_d$ of $\Lambda$ and an indexing
(Definition~\ref{is-eigenvalue-conjugate-frame-fibres}(1) and~(2)); let $n\ge0$ and $w_0\in E$. The
symplectic basis is a basis of $E^r$ over $\mathbf Q$, so every $v\in E^r$ is
$\sum_l\alpha_lb_l+\sum_l\beta_lb'_l$ with $\alpha,\beta\in\mathbf Q^d$, its \emph{coordinates}.
\begin{enumerate}
\item (Heisenberg operators.) For $y\in M_{n\times r}(E)$ let $\alpha,\beta\in M_{n\times d}(\mathbf Q)$ have
  as $k$-th rows $\alpha_k$ and $\beta_k$ the coordinates of the $k$-th row $y_k$ of $y$. For a function
  $\varphi$ on $M_{n\times d}(\mathbf Q)$ put
  \[
   (\rho_X(y)\varphi)(m)=\varphi(m-\beta)\exp\Big(\pi i\sum_{k=1}^n\big(2m_k\cdot\alpha_k-\alpha_k\cdot\beta_k\big)\Big),
  \]
  where $m_k$ is the $k$-th row of $m$.
\item (Signs.) For $y\in M_{n\times r}(E)$ put
  $\varepsilon_{w_0}(y)=\exp\big(\pi i\sum_k\langle y_k,w_0y_k\rangle_X\big)$.
\item (Inner product.) For Schwartz-Bruhat functions $\varphi,\varphi'$ on $M_{n\times d}(\mathbf Q)$
  (Definition~\ref{is-eigenvalue-conjugate-frame-fibres}(6)) let $N\ge1$ be an integer such that
  $\varphi\bar\varphi'$ vanishes outside $M_{n\times d}(N^{-1}\mathbf Z)$ and is invariant under
  $M_{n\times d}(N\mathbf Z)$, and put
  \[
   \langle\varphi,\varphi'\rangle=N^{-nd}\sum_m\varphi(m)\overline{\varphi'(m)},
  \]
  the sum over a set of representatives of $M_{n\times d}(N^{-1}\mathbf Z)/M_{n\times d}(N\mathbf Z)$; this
  does not depend on $N$.
\item (Vacua.) For $\mathcal O_E$-submodules $\Lambda_S\subset E^n$ and $L\subset E^r$, the \emph{tensor
  lattice} $\Lambda_S\otimes L$ is the subgroup of $M_{n\times r}(E)$ generated by the matrices $vu^{\mathsf T}$
  with $v\in\Lambda_S$ and $u\in L$. A \emph{vacuum} of $\Lambda_S\otimes L$ for $w_0$ is a Schwartz-Bruhat
  function $\varphi$ on $M_{n\times d}(\mathbf Q)$ with $\langle\varphi,\varphi\rangle=1$ and
  $\varepsilon_{w_0}(y)\,\rho_X(y)\varphi=\varphi$ for every $y\in\Lambda_S\otimes L$.
\item (Level-one theta functions.) For $x\in M_r(\mathbf C)^\Psi$ and an $\mathcal O_E$-submodule
  $L\subset E^r$, a \emph{level-one theta function} for $w_0$, $x$ and $L$ is a holomorphic function
  $\vartheta$ on $\mathbf C^d$ with
  \[
   \epsilon(v)\,e^{-\pi H_x(z+\frac12p_x(v),\,p_x(v))}\,\vartheta(z+p_x(v))=\vartheta(z),\qquad
   \epsilon(v)=e^{\pi i\langle v,w_0v\rangle_X},
  \]
  for every $v\in L$ and every $z\in\mathbf C^d$.
\end{enumerate}
\end{definition}
\leanref{Development/AbelianVarietyTheory/Definitions/SpanningFamilies/IsSystem.lean}{GlobalTheory.SW.heisG}
\leanref{Development/AbelianVarietyTheory/Definitions/SpanningFamilies/IMemSiegel.lean}{GlobalTheory.SW.heisS}
\leanref{Development/AbelianVarietyTheory/Definitions/SpanningFamilies/IsSystem.lean}{GlobalTheory.SW.epsW}
\leanref{Development/AbelianVarietyTheory/Definitions/SpanningFamilies/IMemSiegel.lean}{GlobalTheory.SW.ip}
\leanref{Development/AbelianVarietyTheory/Definitions/SpanningFamilies/IsSystem.lean}{GlobalTheory.SW.tens}
\leanref{Development/AbelianVarietyTheory/Definitions/SpanningFamilies/IsSystem.lean}{GlobalTheory.SW.IsVac}
\leanref{Development/AbelianVarietyTheory/Definitions/SpanningFamilies/IsStabForm.lean}{GlobalTheory.Dict.IsTheta1}
% heisG y = heisS a b with a, b the matrices of crd (the inverse of vecQ, the rational combination of the basis
% vectors with Sum.inl and Sum.inr); heisS a b phi x = phi (x - b) * exp(pi i sum_k (2 x_k . a_k - a_k . b_k));
% epsW w y = exp(pi i sum_k EF (y k) (w y k)); ip phi psi = haar (phi * conj psi), with haar = haarAt R for a level
% R of the product (R^(-nd) times the finite sum over M(Q)/R M(Z)); tens is the closure of the vecMulVec v u;
% IsVac is IsSB, ip v v = 1 and the fixed-point condition. IsTheta1 is differentiability and the condition
% eps1 w v * fock x 1 (pv v) theta = theta for v in L.

\begin{situation}\label{is-eigenvalue-conjugate-frame-context}
Let $D$ be a definite-side datum with host field $E_h/F_h$, element $\xi_0$, ideal $\mathfrak A$ and set
$\mathcal T$ of embeddings of $E_h$, let $\psi$ be a standard character of $\mathbf A_{F_h}$, and let $\theta$
be an archimedean family whose member $\theta_\rho$ extends $\rho$ for every $\rho\in\mathcal T$. Let $S$ and
$T$ be Hermitian spaces of $D$, where $S$ has Gram matrix $1$ and even rank $n$; put $r=\operatorname{rk}T$
and $G_T(\rho)=\rho(G_T)$. Let $X'$ be flag data of rank $r$ over $E_h$
(Definition~\ref{is-onto-share-flag-data}) with set $\Psi'=\{\bar\rho:\rho\in\mathcal T\}$, matrix
$G_T^{\mathsf T}$, $\kappa'_0=2\xi_0$, integer $p'_0$ and module $\Lambda'$, together with a symplectic basis
of $\Lambda'$ and an indexing (Definition~\ref{is-eigenvalue-conjugate-frame-fibres}(1) and~(2)); put
$d=|\mathcal T|r$, and write $z'_{x'}$, $p'_{x'}$, $\mathcal G'_{x'}$, $\theta'_\varphi$ and $\Theta'_\varphi$
for the coordinates, periods, Gaussians, theta series and theta functions of $X'$ formed with them.
\begin{enumerate}
\item (Orthonormal families.) Let $J'$ be the diagonal $r\times r$ matrix with entries $1$ at the first
  $p'_0$ places and $-1$ at the others. A family $B=(B_\rho)_{\rho\in\mathcal T}$ of complex $r\times r$
  matrices is \emph{orthonormal for $X'$} if $B_\rho^*G_T(\rho)B_\rho=J'$ for every $\rho$; then
  $s'_{B,\rho}=B_\rho J'B_\rho^{-1}$, and $x'_B\in M_r(\mathbf C)^{\Psi'}$ has the component
  $\overline{B_\rho}$ at $\bar\rho$. For $g\in U(T)(\mathbf A)$ put $gB=(g_\rho B_\rho)_\rho$ with
  $g_\rho=\theta_\rho(g)$.
\item (Zak theta functions.) For $w\in\mathbb W$ with finite component $0$ let $u'(w)=(u'_1,\dots,u'_n)$,
  where $(u'_k)_{\bar\rho}$ is the complex conjugate of the $k$-th column of $\theta_\rho(w)$. For
  $x'\in M_r(\mathbf C)^{\Psi'}$ and a Schwartz-Bruhat function $\varphi$ on $M_{n\times d}(\mathbf Q)$, a
  \emph{Zak theta function} for $x'$ and $\varphi$ is a theta function $f$ for $(S,T)$ with
  $f(w)=\Theta'_\varphi(x';u'(w))$ for every $w\in\mathbb W$ with finite component $0$.
\end{enumerate}
\end{situation}
\leanref{Development/AbelianVarietyTheory/Definitions/SpanningFamilies/IsWeakOfDatum.lean}{GlobalTheory.Ker.Ctx}
\leanref{Development/AbelianVarietyTheory/Definitions/SpanningFamilies/IsWeakOfDatum.lean}{GlobalTheory.Ker.Match}
\leanref{Development/AbelianVarietyTheory/Definitions/SpanningFamilies/IsWeakOfDatum.lean}{GlobalTheory.Ker.Ctx.IsZak}
% Ker.Ctx S T: psi with IsStdChar, theta with IsInfFam, X : Data with Ker.Match S T X (S.gram = 1, X.Gm = T.gram^T,
% e : Tau = X.Psi with e tau = conj (tauE tau), X.kappa0 = 2 xi0), q0 with p0 + q0 = r, bas, idx, IsSympl and
% Even S.rank. IsB B is B^H gramAt B = sgnM X (J'), sB B is s'_B, xB B is x'_B, mvB g B is gB, realPt w is u'(w),
% IsZak x phi f is the Zak condition, and Th x phi is the Zak theta function of Lemma frame-zak(1).

\begin{lemma}\label{is-eigenvalue-conjugate-frame-conjugate-model}
In Situation~\ref{is-eigenvalue-conjugate-frame-situation}, there are flag data $X'$ of rank $r$ over $E_h$
with set $\Psi'=\{\bar\rho:\rho\in\mathcal T\}$, matrix $G_T^{\mathsf T}$, $\kappa'_0=2\xi_0$, integer
$p'_0=r-p_0$ and module $\Lambda$, a symplectic basis of $\Lambda$ for $X'$ and an indexing for $X'$, with
which $\psi$, $\theta$, $S$, $T$ and $X'$ are as in Situation~\ref{is-eigenvalue-conjugate-frame-context}, and
a permutation $\pi$ of $\{1,\dots,r\}$, with the following properties; we use the notation of that situation.
For a family $Z$ of complex $r\times r$ matrices let $Z^\pi$ be the family with $Z^\pi_\rho
e_a=Z_\rho e_{\pi(a)}$. For a polynomial $Q$ in the \emph{creation variables} $(\rho,a,k)$, $\rho\in\mathcal
T$, $1\le a\le r$, $1\le k\le n$, let $Q^\pi$ be $Q$ with each $(\rho,a,k)$ replaced by $(\rho,\pi^{-1}(a),k)$;
for a Schwartz-Bruhat function $\varphi$ on $M_{n\times d}(\mathbf Q)$ and $x\in M_r(\mathbf C)^{\mathcal T}$
let $D_x(Q)\varphi$ be the value at $z=0$ of $Q(\partial)\,\theta_\varphi(x;z)$, where $Q(\partial)$ replaces
each variable $(\rho,a,k)$ by $(2/f_\rho)\,\partial/\partial z_{k,j}$, with $f_\rho=2\sqrt{2\pi|\rho(\xi_0)|}$
and $j$ the index with $(\tau_j,a_j)=(\rho,a)$ in the indexing of $X$; and define $D'_{x'}(Q)\varphi$ for
$x'\in M_r(\mathbf C)^{\Psi'}$ in the same way with $\theta'_\varphi$ and the index $j$ with $(\tau'_j,a'_j)=
(\bar\rho,a)$ in the indexing of $X'$.
\begin{enumerate}
\item $\langle u,v\rangle_{X'}=\langle u,v\rangle_X$ for all $u,v\in E_h^r$, and the Heisenberg operators
  $\rho_{X'}(y)$ and $\rho_X(y)$ of Definition~\ref{is-eigenvalue-conjugate-frame-heisenberg}(1), formed with
  the symplectic bases of $X'$ and of $X$, are equal for every $y\in M_{n\times r}(E_h)$.
\item A family $Z$ is a signature family if and only if $Z^\pi$ is orthonormal for $X'$.
\item Let $Z$ be a signature family. Then $s'_{Z^\pi,\rho}=s_{Z,\rho}$ for every $\rho$; for every
  Schwartz-Bruhat function $\varphi$ on $M_{n\times d}(\mathbf Q)$ a function $f$ on $\mathbb W$ is a Zak
  theta function for $x'_{Z^\pi}$ and $\varphi$ if and only if $f$ is a theta function with
  $f(w)=\Theta_\varphi(x^Z;u(w))$ for every $w\in\mathbb W$ with finite component $0$; and
  $D'_{x'_{Z^\pi}}(Q^\pi)\varphi=D_{x^Z}(Q)\varphi$ for every such $\varphi$ and every polynomial $Q$.
\end{enumerate}
\end{lemma}
\leanref{Development/AbelianVarietyTheory/Proofs/IwasawaDecomposition/ExistsConjugateModelSameThetaFunctionsStatement.lean}{GlobalTheory.T.exists_conjugateModel_same_thetaFunctions}
\leanref{Development/AbelianVarietyTheory/Proofs/IwasawaDecomposition/ExistsConjugateModelSameThetaFunctions.lean}{GlobalTheory.OneMapRef.Flip.exists_flip_model}
\leanref{Development/AbelianVarietyTheory/Definitions/ThetaValues/PolyDerivSmulLetter.lean}{GlobalTheory.ThFr.letterDir}
\leanref{Development/AbelianVarietyTheory/Definitions/SpanningFamilies/IsSystem.lean}{GlobalTheory.SW.thP}
\begin{proof}
\emph{The data.} Write $A=-G_T^{\mathsf T}$ for the matrix of $X$ and $h,h'$ for the forms of $A$ and
$-A=G_T^{\mathsf T}$. The set $\Psi'$ contains $\sigma$ exactly when it does not contain $\bar\sigma$, as
$\mathcal T$ has this property; $-A$ is Hermitian with $\det(-A)\ne0$; and
$h'_{\bar\tau}(\bar u,\bar v)=-\overline{h_\tau(u,v)}$ for $u,v\in\mathbf C^r$. So if the columns of $P$
have Gram matrix $\operatorname{diag}(1,\dots,1,-1,\dots,-1)$, with $p_0$ entries $1$, for $h_\tau$,
then the columns of $\bar P$ in the reverse order have Gram matrix $J'$ for $h'_{\bar\tau}$. The module
$\Lambda$ and the element $-\xi_0$ are those of $X$, $\kappa'_0=-\kappa_0\in\{2(-\xi_0),-2(-\xi_0)\}$,
and $\operatorname{Im}\bar\tau(\kappa'_0)=\operatorname{Im}\tau(\kappa_0)<0$ for $\tau\in\mathcal T$.
Choose an integer $k\ge1$ such that $k\,\kappa_0h(x,y)$ is an algebraic integer for $x,y\in\Lambda$, and
put $m'=3k$. Let $\Gamma'$ be the group of elements of $GL_r(E_h)$ that preserve $h'$ and $\Lambda$ and
satisfy $\gamma v-v\in3\Lambda$ for $v\in\Lambda$, a subgroup of finite index, and $\chi'_\kappa=1$. A
transvection $v\mapsto v+m's\kappa'_0h'(v,\xi)\xi$ as in (6) lies in $\Gamma'$, as
$m's\kappa'_0h'(v,\xi)\xi=3\,s\,k\kappa_0h(v,\xi)\xi\in3\Lambda$. If $\gamma\in\Gamma'$ satisfies
$\gamma x=xb$ with $x\in\mathcal D_{X'}$ and $b$ upper triangular, then
$\gamma^\tau=x_\tau b_\tau x_\tau^{-1}$, and $b_\tau$ preserves the Gram matrix of the columns of $x_\tau$,
whose leading minors are not $0$; so $b_\tau$ is diagonal with entries of absolute value $1$, the powers
of $\gamma$ are bounded at every $\tau$, $\gamma$ has finite order as it preserves $\Lambda$, and
$\gamma=1$ as it is congruent to $1$ modulo $3$ \cite{Min87}. Give $X'$ the symplectic basis of $X$, which
is symplectic for $X'$ by (1) below, and the indexing $j\mapsto(\bar\tau_j,r+1-a_j)$, and let $\pi$ be the
reversal $a\mapsto r+1-a$. Then $\psi$, $\theta$, $S$, $T$ and $X'$ are as in
Situation~\ref{is-eigenvalue-conjugate-frame-context}, with $p'_0=r-p_0$: $S$ has Gram matrix $1$ and even
rank by Situation~\ref{is-eigenvalue-conjugate-frame-situation}.

(1) As $h'=-h$ and $\kappa'_0=-\kappa_0$,
$\operatorname{Tr}(-\frac12\kappa'_0h'(u,v))=\operatorname{Tr}(-\frac12\kappa_0h(u,v))$. The Heisenberg
operators depend only on the coordinates of the rows of $y$ in the symplectic basis, which is the same for
$X$ and $X'$.

(2) Let $P$ be the permutation matrix that reverses the order of the columns, so that $Z^\pi_\rho=Z_\rho P$,
$P^{-1}=P^{\mathsf T}=P$, and $PNP=J'$ for the matrix $N$ of
Situation~\ref{is-eigenvalue-conjugate-frame-situation}(1). Then
$(Z_\rho P)^*G_T(\rho)(Z_\rho P)=P(Z_\rho^*G_T(\rho)Z_\rho)P$, which is $J'$ exactly when
$Z_\rho^*G_T(\rho)Z_\rho=N$.

(3) First, $s'_{Z^\pi,\rho}=Z_\rho PJ'PZ_\rho^{-1}=Z_\rho NZ_\rho^{-1}=s_{Z,\rho}$. Put $x=x^Z$ and
$x'=x'_{Z^\pi}$, so that $x'_{\bar\rho}=\overline{x_\rho}P$ and $x'_{\bar\rho}e_{r+1-a}=\overline{x_\rho e_a}$.
As $x$ is orthonormal (Lemma~\ref{is-eigenvalue-conjugate-frame-signature}(1)),
$h_\rho(v,x_\rho e_a)=\operatorname{sgn}(a)(x_\rho^{-1}v)_a$ for every $v\in\mathbf C^r$. Let $u$ be a
family of vectors $u_\rho\in\mathbf C^r$ indexed by $\mathcal T$, and $u'$ the family with
$u'_{\bar\rho}=\overline{u_\rho}$. We show $z'_{x'}(u')=z_x(u)$. Let $j$ be an index with $(\tau_j,a_j)=
(\rho,a)$. If $a\le p_0$, then $r+1-a>p'_0$, and
\[
 z'_{x'}(u')_j=(x'^{-1}_{\bar\rho}\overline{u_\rho})_{r+1-a}=\overline{(x_\rho^{-1}u_\rho)_a}
 =\overline{h_\rho(u_\rho,x_\rho e_a)}=h_\rho(x_\rho e_a,u_\rho)=z_x(u)_j.
\]
If $a>p_0$, then $r+1-a\le p'_0$, and
\[
 z'_{x'}(u')_j=h'_{\bar\rho}(\overline{x_\rho e_a},\overline{u_\rho})=-\overline{h_\rho(x_\rho e_a,u_\rho)}
 =-h_\rho(u_\rho,x_\rho e_a)=(x_\rho^{-1}u_\rho)_a=z_x(u)_j.
\]
For $v\in E_h^r$ the family $(\bar\rho(v))_{\bar\rho}$ is $u'$ for $u=(\rho(v))_\rho$, so $p'_{x'}(v)=p_x(v)$.
Hence the period matrices $\Omega$, $\Omega'$ of the symplectic basis, the matrix $\tau_x$ and the form
$H_x$ are the same for $X'$ at $x'$ as for $X$ at $x$, and so are the Gaussians and the theta series. For
$w$ with finite component $0$, $u'(w)_k$ is the family $u'$ attached to $u=u(w)_k$, so
$\Theta'_\varphi(x';u'(w))=\Theta_\varphi(x;u(w))$, which gives the second assertion. For the third, the
theta series $\theta'_\varphi(x';z)$ and $\theta_\varphi(x;z)$ are the same function of $z$, and the variable
$(\rho,\pi^{-1}(a),k)=(\rho,r+1-a,k)$ of $Q^\pi$ gives the derivative in $z_{k,j}$ for the index $j$ with
$(\bar\tau_j,r+1-a_j)=(\bar\rho,r+1-a)$, that is, with $(\tau_j,a_j)=(\rho,a)$, with the factor $2/f_\rho$; so
$Q^\pi(\partial)$ for $X'$ is $Q(\partial)$ for $X$, and $D'_{x'}(Q^\pi)\varphi=D_x(Q)\varphi$.
\end{proof}
% exists_flip_model: exists_flip (FlagData.exists_data with M = 3, m = 3|k|; Gamma the congruence subgroup,
% Gamma_free by GamFree.gam_free with latMinkowski, chi = 1), the Match m (gS, gT, e = eT then conjugation with
% e_emb, ka), bas' with the same vectors (exists_bas), idx' = idx then (e, Fin.revPerm), pi = Fin.revPerm, hs
% from EF_flip, (2) from negM_rev, s_B = cartZ from negM_rev and inv_submatrix_equiv, and the Zak equivalence from
% zR_flip (herm_col for the orthonormal frame gram_frOf), hper and thA_congr; heisG X' bas' = heisG X bas is
% heisG_congr, and the transfer of thP (D_x(Q) phi is thP X hq bas idx phi (toP S hr eT hq Q) x, D' is thP X' with
% C.toH; f_rho is fockScale) is h1, h2 and polyDeriv_thA_rename with thA_congr. The Lean states X', m, bas', idx'
% and pi existentially; the proof's construction fixes them as above. The context is ctxOf of X' with these data.

\begin{lemma}\label{is-eigenvalue-conjugate-frame-zak-coordinates}
Let $X$, its symplectic basis and indexing be as in
Definition~\ref{is-eigenvalue-conjugate-frame-heisenberg}, with set $\Psi$, form $h$ and forms $h_\sigma$, let
$n\ge0$ be even (the parity is not used) and $x\in\mathcal D_X$. Let $U$ be the real vector space of the $n$-tuples
$u=(u_1,\dots,u_n)$ of families $u_k=(u_{k,\sigma})_{\sigma\in\Psi}$ of vectors of $\mathbf C^r$, and write
$z_x(u)=(z_x(u_1),\dots,z_x(u_n))\in(\mathbf C^d)^n$ for the coordinates
(Definition~\ref{is-eigenvalue-conjugate-frame-fibres}(1)). For $y\in M_{n\times r}(E)$ with rows $y_k$ let $\hat y\in
U$ have $\hat y_{k,\sigma}=\sigma(y_k)$, and for $u,u'\in U$ put
$E_X(u,u')=\sum_{\sigma\in\Psi}\sum_k2\operatorname{Re}\big(\sigma(-\tfrac12\kappa_0)\,h_\sigma(u_{k,\sigma},u'_{k,\sigma})\big)$.
\begin{enumerate}
\item The map $u\mapsto z_x(u)$ is an isomorphism of real vector spaces from $U$ onto $(\mathbf C^d)^n$.
\item $\operatorname{Im}\sum_kH_x(z_x(u_k),z_x(u'_k))=E_X(u,u')$ for all $u,u'\in U$.
\item $E_X(\hat y,\hat y')=\sum_k\langle y_k,y'_k\rangle_X$ for all $y,y'\in M_{n\times r}(E)$.
\item The map $y\mapsto\hat y$ has dense image in $U$.
\item Let $\gamma\in GL_r(E)$ preserve $h$. Then $z_{\gamma x}(\gamma u)=z_x(u)$ for every $u\in U$, where
  $(\gamma x)_\sigma=\sigma(\gamma)x_\sigma$ and $(\gamma u)_{k,\sigma}=\sigma(\gamma)u_{k,\sigma}$.
\end{enumerate}
\end{lemma}
\leanref{Development/AbelianVarietyTheory/Proofs/SiegelThetaTransformation/RealPeriodMapHomeomorphRiemannFormDensityStatement.lean}{GlobalTheory.T.realPeriodMap_homeomorph_riemannForm_density}
\leanref{Development/AbelianVarietyTheory/Proofs/SiegelThetaTransformation/RealPeriodMapHomeomorphRiemannFormDensity.lean}{GlobalTheory.Ker.zR_laws}
\leanref{Development/AbelianVarietyTheory/Definitions/SpanningFamilies/Fischer.lean}{GlobalTheory.Ker.EFR}
\begin{proof}
The map $u\mapsto z_x(u)$ is $\mathbf R$-linear, and $z_x(\hat y)_k=p_x(y_k)$ by the definition of the periods.

(3) As $h_\sigma(\sigma(u),\sigma(v))=\sigma(h(u,v))$ and $\Psi$ contains one of every two conjugate embeddings,
$E_X(\hat y,\hat y')=\sum_k\sum_{\sigma\in\Psi}2\operatorname{Re}\sigma\big(-\frac12\kappa_0h(y_k,y'_k)\big)=
\sum_k\operatorname{Tr}_{E/\mathbf Q}\big(-\frac12\kappa_0h(y_k,y'_k)\big)=\sum_k\langle y_k,y'_k\rangle_X$.

(1) and (4) By Lemma~\ref{is-eigenvalue-conjugate-frame-hermitian}(4), the periods $p_x(v)$, $v\in E^r$, span $\mathbf
C^d$ over $\mathbf R$, so the $\mathbf Q$-vector space $V=\{z_x(\hat y)\}=\{(p_x(y_k))_k\}$ spans $(\mathbf C^d)^n$ over
$\mathbf R$. A $\mathbf Q$-subspace that spans a finite-dimensional real vector space is dense in it: it contains a real
basis, and the rational combinations of a basis are dense. So the image of $z_x$ is a dense real subspace, hence all of
$(\mathbf C^d)^n$. As $\dim_{\mathbf R}U=2|\Psi|rn=2dn=\dim_{\mathbf R}(\mathbf C^d)^n$, $z_x$ is bijective, which is
(1); and the $\hat y$ are dense in $U$, as their images form the dense set $V$ and $z_x^{-1}$ is continuous, which is
(4).

(2) Both sides are continuous real bilinear functions of $(u,u')$. At $(\hat y,\hat y')$ the left side is
$\sum_k\operatorname{Im}H_x(p_x(y_k),p_x(y'_k))=\sum_k\langle y_k,y'_k\rangle_X$ by
Lemma~\ref{is-eigenvalue-conjugate-frame-hermitian}(2), which is the right side by (3). By (4) they agree everywhere.

(5) Both sides are continuous and $\mathbf R$-linear in $u$. At $u=\hat y$, $\gamma\hat y$ is $\hat y'$ for the matrix
$y'$ with rows $\gamma y_k$, and $z_{\gamma x}(\gamma\hat y)_k=p_{\gamma x}(\gamma y_k)=p_x(y_k)=z_x(\hat y)_k$ by
Lemma~\ref{is-eigenvalue-conjugate-frame-hermitian}(5). By (4) they agree everywhere.
\end{proof}
% Ker.zR_laws (RealPeriodMapHomeomorphRiemannFormDensity.lean): exists_equiv_zR (surjectivity from denseRange_zR_embR,
% denseRange_pv through Dict.span_range_pv and dense_of_span_real, and the dimension count), hB and hF by
% DenseRange.equalizer from Dict.hx_laws (Lemma is-eigenvalue-conjugate-frame-hermitian), EFR_embR and
% PerMove.emb_mulVec. The Lean's statement also takes q0 with p0 + q0 = r, which only names the number of indices
% above p0.

\begin{lemma}\label{is-eigenvalue-conjugate-frame-zak-series}
Let $X$, its symplectic basis and indexing be as in
Definition~\ref{is-eigenvalue-conjugate-frame-heisenberg}, let $n\ge0$, $x\in\mathcal D_X$, and $\varphi$ and $\psi$
Schwartz-Bruhat functions on $M_{n\times d}(\mathbf Q)$, with the theta series of
Definition~\ref{is-eigenvalue-conjugate-frame-fibres}(6).
\begin{enumerate}
\item The function $z\mapsto\theta_\varphi(x;z)$ is holomorphic on $(\mathbf C^d)^n$.
\item $\theta_{\varphi+\psi}(x;\cdot)=\theta_\varphi(x;\cdot)+\theta_\psi(x;\cdot)$, and $\theta_{a\varphi}(x;\cdot)=
  a\,\theta_\varphi(x;\cdot)$ for $a\in\mathbf C$.
\item If $\theta_\varphi(x;\cdot)=\theta_\psi(x;\cdot)$, then $\varphi=\psi$.
\end{enumerate}
\end{lemma}
\leanref{Development/AbelianVarietyTheory/Proofs/SiegelThetaTransformation/ThetaSeriesTestFunctionHolomorphicWeight.lean}{GlobalTheory.SW.ThetaHol.differentiable_thA}
\leanref{Development/AbelianVarietyTheory/Proofs/SiegelThetaTransformation/WeilActionPreservesThetaAndInnerProduct.lean}{GlobalTheory.SW.UnAct.thA_add}
\leanref{Development/AbelianVarietyTheory/Proofs/SiegelThetaTransformation/WeilActionPreservesThetaAndInnerProduct.lean}{GlobalTheory.SW.UnAct.thA_smul}
\leanref{Development/AbelianVarietyTheory/Proofs/SiegelThetaTransformation/WeilActionPreservesThetaAndInnerProduct.lean}{GlobalTheory.SW.UnAct.eq_of_thA_eq}
\begin{proof}
\openstep{The absolute and locally uniform convergence of the theta series in $\tau_x$, which lies in the Siegel upper
half space, the holomorphy that follows, and the recovery of a Schwartz-Bruhat function from its theta series as a
function of the variable, are not written out.}
% ThetaHol.differentiable_thA, UnAct.thA_add (SiegelSum.summable_norm_thetaTerm and tsum_add), UnAct.thA_smul (tsum_mul_left)
% and UnAct.eq_of_thA_eq (eq_of_thS_eq at tauMat, with hMat_mul_perMat and isUnit_det_hMat). Size, following imports and
% stopping at statements the book cites: about 900 lines for the four, mostly the convergence (SummableNormSiegelThetaSeriesTerms,
% DifferentiableOnOfTendstoLocallyUniformlyOn) and eq_of_thS_eq.
\end{proof}

\begin{lemma}\label{is-eigenvalue-conjugate-frame-zak-weil}
Let $X$, its symplectic basis and indexing be as in
Definition~\ref{is-eigenvalue-conjugate-frame-heisenberg}, with form $h$, and let $n\ge0$ be even. Let $\gamma\in
GL_r(E)$ preserve $h$, and write $(\gamma x)_\sigma=\sigma(\gamma)x_\sigma$. For a Schwartz-Bruhat function $\varphi$ on
$M_{n\times d}(\mathbf Q)$ let $w_\gamma\varphi$ be a Schwartz-Bruhat function $\psi$ with $\theta_\psi(\gamma
x;\cdot)=\theta_\varphi(x;\cdot)$ for every $x\in\mathcal D_X$, chosen once for all when there is one, and $0$ otherwise.
Let $\varphi$ and $\psi$ be Schwartz-Bruhat.
\begin{enumerate}
\item $w_\gamma\varphi$ is Schwartz-Bruhat.
\item $\theta_{w_\gamma\varphi}(\gamma x;\cdot)=\theta_\varphi(x;\cdot)$ for every $x\in\mathcal D_X$.
\item $\langle w_\gamma\varphi,w_\gamma\psi\rangle=\langle\varphi,\psi\rangle$, $w_\gamma(\varphi+\psi)=w_\gamma\varphi+w_\gamma\psi$, and
  for $x\in\mathcal D_X$, $\theta_\varphi(x;\cdot)=\theta_\psi(x;\cdot)$ implies $\varphi=\psi$.
\item $w_\gamma(a\varphi)=a\,w_\gamma\varphi$ for $a\in\mathbf C$.
\item $w_1\varphi=\varphi$.
\item $w_{\gamma'\gamma}\varphi=w_{\gamma'}(w_\gamma\varphi)$ for every $\gamma'\in GL_r(E)$ that preserves $h$.
\item $w_\gamma(\rho_X(y)\varphi)=\rho_X(y\gamma^{\mathsf T})(w_\gamma\varphi)$ for every $y\in M_{n\times r}(E)$.
\end{enumerate}
\end{lemma}
\leanref{Development/AbelianVarietyTheory/Proofs/SiegelThetaTransformation/WeilActionPreservesThetaAndInnerProduct.lean}{GlobalTheory.SW.wG_laws}
\leanref{Development/AbelianVarietyTheory/Definitions/SpanningFamilies/IsSystem.lean}{GlobalTheory.SW.wG}
\leanref{Development/AbelianVarietyTheory/Definitions/SpanningFamilies/IsSystem.lean}{GlobalTheory.SW.IsUn}
\begin{proof}
\emph{The Weil operator.} There is a map $\omega$ on the Schwartz-Bruhat functions such that $\omega\varphi$ is
Schwartz-Bruhat, $\theta_{\omega\varphi}(\gamma x;\cdot)=\theta_\varphi(x;\cdot)$ for every $x\in\mathcal D_X$, and
$\langle\omega\varphi,\omega\psi\rangle=\langle\varphi,\psi\rangle$, for all Schwartz-Bruhat $\varphi,\psi$.
\openstep{The construction of $\omega$, which reads $\gamma$ through its action on the periods of the symplectic
basis as a rational symplectic matrix, takes the Weil representation of the rational symplectic group on
Schwartz-Bruhat functions with the theta transformation formula, and uses that $n$ is even to remove the cocycle, is not
written out.}
% UnAct.exists_weil: periods_lact_symplectic (sigma symplectic with tauMat (lact gamma x) the image of tauMat x and hMat
% transformed by spDen), T.exists_weil_of_symplectic' (about 1,450 lines), T.canonicalThetaExponent_symplectic_invariant
% and T.symplecticSiegelAction_mem_upperHalfSpace_and_cocycle. Size, following imports and stopping at statements the
% book cites (with Lemma -zak-series): about 2,400 lines.

\emph{Uniqueness.} If $\psi'$ is Schwartz-Bruhat and $\theta_{\psi'}(\gamma x;\cdot)=\theta_\varphi(x;\cdot)$ for every
$x\in\mathcal D_X$, then $w_\gamma\varphi=\psi'$. Indeed, then a function as in the definition of $w_\gamma\varphi$ exists, so
$w_\gamma\varphi$ is one; the set $\mathcal D_X$ is not empty (it contains the orthonormal points given by data (2) of
Definition~\ref{is-onto-share-flag-data}), and for $x\in\mathcal D_X$, $\gamma x\in\mathcal D_X$
(Lemma~\ref{is-eigenvalue-conjugate-frame-hermitian}(5)) and $\theta_{w_\gamma\varphi}(\gamma x;\cdot)=\theta_{\psi'}(\gamma
x;\cdot)$, so $w_\gamma\varphi=\psi'$ by Lemma~\ref{is-eigenvalue-conjugate-frame-zak-series}(3).

(1) and (2) By uniqueness, $w_\gamma\varphi=\omega\varphi$.

(3) The first assertion is that of $\omega$. The function $w_\gamma\varphi+w_\gamma\psi$ is Schwartz-Bruhat
(Lemma~\ref{is-eigenvalue-conjugate-frame-vacua-mean}(1)), and its theta series at $\gamma x$ is
$\theta_\varphi(x;\cdot)+\theta_\psi(x;\cdot)=\theta_{\varphi+\psi}(x;\cdot)$ by (2) and
Lemma~\ref{is-eigenvalue-conjugate-frame-zak-series}(2); so it is $w_\gamma(\varphi+\psi)$ by uniqueness. The last
assertion is Lemma~\ref{is-eigenvalue-conjugate-frame-zak-series}(3).

(4) In the same way, $a\,w_\gamma\varphi$ is Schwartz-Bruhat with theta series $a\,\theta_\varphi(x;\cdot)=\theta_{a\varphi}
(x;\cdot)$ at $\gamma x$.

(5) The element $1$ preserves $h$, $1x=x$, and $\varphi$ has the property of $w_1\varphi$.

(6) The element $\gamma'\gamma$ preserves $h$, and $(\gamma'\gamma)x=\gamma'(\gamma x)$. The function $w_{\gamma'}(w_\gamma\varphi)$
is Schwartz-Bruhat, and for $x\in\mathcal D_X$, as $\gamma x\in\mathcal D_X$,
$\theta_{w_{\gamma'}(w_\gamma\varphi)}(\gamma'\gamma x;\cdot)=\theta_{w_\gamma\varphi}(\gamma x;\cdot)=\theta_\varphi(x;\cdot)$; so it is
$w_{\gamma'\gamma}\varphi$ by uniqueness.

(7) The rows of $y\gamma^{\mathsf T}$ are the $\gamma y_k$. Let $x\in\mathcal D_X$; by
Lemma~\ref{is-eigenvalue-conjugate-frame-hermitian}(5), $\gamma x\in\mathcal D_X$, $p_{\gamma x}(\gamma y_k)=p_x(y_k)$ and
$H_{\gamma x}=H_x$. By Lemma~\ref{is-eigenvalue-conjugate-frame-vacua-heisenberg}(5) at $\gamma x$ and at $x$, and (2),
\[
 \theta_{\rho_X(y\gamma^{\mathsf T})w_\gamma\varphi}(\gamma x;z)=e^{-\pi\sum_kH_x(z_k+\frac12p_x(y_k),p_x(y_k))}
 \theta_\varphi(x;z+p_x(y))=\theta_{\rho_X(y)\varphi}(x;z).
\]
The function $\rho_X(y\gamma^{\mathsf T})w_\gamma\varphi$ is Schwartz-Bruhat
(Lemma~\ref{is-eigenvalue-conjugate-frame-vacua-heisenberg}(2)), so it is $w_\gamma(\rho_X(y)\varphi)$ by uniqueness.
\end{proof}
% SW.wG_laws: key = exists_weil with wG_eq (exists_mem_dom, IsUn.lact_mem, eq_of_thA_eq), then
% T.rationalSchwartzBruhat_algebra_integral_linear and thA_add for the sum, thA_smul, lact_one', lact_mul' with
% isUn_mul, and heisG_laws with per_lact for the last clause.

\begin{lemma}\label{is-eigenvalue-conjugate-frame-zak-real}
Let $X$, its symplectic basis and indexing be as in
Definition~\ref{is-eigenvalue-conjugate-frame-heisenberg}, with set $\Psi$, form $h$ and forms $h_\sigma$, let
$n\ge0$ be even, $x\in\mathcal D_X$ and $\varphi$ a Schwartz-Bruhat function on $M_{n\times d}(\mathbf Q)$. Let
$U$ be the real vector space of the $n$-tuples $u=(u_1,\dots,u_n)$ of families $u_k=(u_{k,\sigma})_{\sigma\in\Psi}$ of
vectors of $\mathbf C^r$, on which the theta function $u\mapsto\Theta_\varphi(x;u)$ of
Definition~\ref{is-eigenvalue-conjugate-frame-fibres}(6) is defined. For $y\in M_{n\times r}(E)$ with rows $y_k$
let $\hat y\in U$ have $\hat y_{k,\sigma}=\sigma(y_k)$, and for $u,u'\in U$ put
$E_X(u,u')=\sum_{\sigma\in\Psi}\sum_k2\operatorname{Re}\big(\sigma(-\tfrac12\kappa_0)\,h_\sigma(u_{k,\sigma},u'_{k,\sigma})\big)$.
\begin{enumerate}
\item The function $u\mapsto\Theta_\varphi(x;u)$ is smooth on $U$.
\item For $y\in M_{n\times r}(E)$ and $u\in U$,
  $\Theta_{\rho_X(y)\varphi}(x;u)=e^{-\pi iE_X(u,\hat y)}\,\Theta_\varphi(x;u+\hat y)$.
\item Let $g\in GL_r(E)$ preserve $h$, and let $w_g\varphi$ be as in
  Lemma~\ref{is-eigenvalue-conjugate-frame-zak-weil}. Then $\Theta_{w_g\varphi}(gx;gu)=\Theta_\varphi(x;u)$ for every $u\in
  U$, where $(gx)_\sigma=\sigma(g)x_\sigma$ and $(gu)_{k,\sigma}=\sigma(g)u_{k,\sigma}$.
\item For Schwartz-Bruhat $\psi$, if $\Theta_\varphi(x;\cdot)=\Theta_\psi(x;\cdot)$ on $U$, then $\varphi=\psi$; and
  $\Theta_{\varphi+\psi}(x;\cdot)=\Theta_\varphi(x;\cdot)+\Theta_\psi(x;\cdot)$ on $U$.
\item $\Theta_{a\varphi}(x;\cdot)=a\,\Theta_\varphi(x;\cdot)$ on $U$, for $a\in\mathbf C$.
\end{enumerate}
\end{lemma}
\leanref{Development/AbelianVarietyTheory/Proofs/SiegelThetaTransformation/GaussRMulFactor.lean}{GlobalTheory.Ker.thR_laws}
\leanref{Development/AbelianVarietyTheory/Definitions/SpanningFamilies/Fischer.lean}{GlobalTheory.Ker.EFR}
\begin{proof}
By definition $\Theta_\varphi(x;u)=\mathcal G_x(z_x(u))\,\theta_\varphi(x;z_x(u))$, where $u\mapsto z_x(u)$ is an isomorphism of
real vector spaces (Lemma~\ref{is-eigenvalue-conjugate-frame-zak-coordinates}(1)) and $\mathcal G_x$ never vanishes.

(1) The Gaussian $\mathcal G_x$ is the exponential of a real polynomial in the real and imaginary parts of $z$, so it is
smooth; $\theta_\varphi(x;\cdot)$ is holomorphic (Lemma~\ref{is-eigenvalue-conjugate-frame-zak-series}(1)), so smooth;
and $z_x$ is linear.

(2) Put $z=z_x(u)$ and $a=z_x(\hat y)$, so that $a_k=p_x(y_k)$ and $z+a=z_x(u+\hat y)$. By
Lemma~\ref{is-eigenvalue-conjugate-frame-vacua-heisenberg}(5),
$\theta_{\rho_X(y)\varphi}(x;z)=e^{-\pi\sum_kH_x(z_k+\frac12a_k,a_k)}\theta_\varphi(x;z+a)$. Write $\Sigma=\sum_kH_x(z_k,a_k)$; as
$H_x(w,z)=\overline{H_x(z,w)}$ (Lemma~\ref{is-eigenvalue-conjugate-frame-hermitian}(1)),
\[
 \mathcal G_x(z)\,e^{-\pi\sum_kH_x(z_k+\frac12a_k,a_k)}=e^{-\frac\pi2\sum_k(H_x(z_k,z_k)+2H_x(z_k,a_k)+H_x(a_k,a_k))}
 =e^{-\frac\pi2(\Sigma-\bar\Sigma)}\mathcal G_x(z+a)=e^{-\pi i\operatorname{Im}\Sigma}\mathcal G_x(z+a).
\]
By Lemma~\ref{is-eigenvalue-conjugate-frame-zak-coordinates}(2), $\operatorname{Im}\Sigma=E_X(u,\hat y)$. So
$\Theta_{\rho_X(y)\varphi}(x;u)=e^{-\pi iE_X(u,\hat y)}\mathcal G_x(z+a)\theta_\varphi(x;z+a)=e^{-\pi iE_X(u,\hat
y)}\Theta_\varphi(x;u+\hat y)$.

(3) By Lemma~\ref{is-eigenvalue-conjugate-frame-hermitian}(5), $gx\in\mathcal D_X$ and $H_{gx}=H_x$, so
$\mathcal G_{gx}=\mathcal G_x$; by Lemma~\ref{is-eigenvalue-conjugate-frame-zak-coordinates}(5), $z_{gx}(gu)=z_x(u)$; and by
Lemma~\ref{is-eigenvalue-conjugate-frame-zak-weil}(2), $\theta_{w_g\varphi}(gx;\cdot)=\theta_\varphi(x;\cdot)$. Hence
$\Theta_{w_g\varphi}(gx;gu)=\mathcal G_x(z_x(u))\theta_\varphi(x;z_x(u))=\Theta_\varphi(x;u)$.

(4) and (5) If $\Theta_\varphi(x;\cdot)=\Theta_\psi(x;\cdot)$, then $\theta_\varphi(x;z)=\theta_\psi(x;z)$ for every $z$, as
$\mathcal G_x(z)\ne0$ and $z_x$ is onto; so $\varphi=\psi$ by Lemma~\ref{is-eigenvalue-conjugate-frame-zak-weil}(3) with
$\gamma=1$. The other identities follow from Lemma~\ref{is-eigenvalue-conjugate-frame-zak-series}(2).
\end{proof}
% Ker.thR_laws (GaussRMulFactor.lean), followed: T.realPeriodMap_homeomorph_riemannForm_density (Lemma -zak-coordinates:
% L, him, hinv), Dict.hx_laws (hsym, hga; Lemma is-eigenvalue-conjugate-frame-hermitian), ThetaHol.differentiable_thA
% and RealHeis.contDiff_gaussR for (1); heisG_laws (Lemma -vacua-heisenberg(5)) and RealHeis.gaussR_mul_factor for (2);
% wG_laws (2) for (3); wG_laws with gamma = 1 (its injectivity clause) for (4), and UnAct.thA_add and thA_smul.

\begin{lemma}\label{is-eigenvalue-conjugate-frame-zak-points}
In Situation~\ref{is-eigenvalue-conjugate-frame-context}, let $U'$ be the real vector space of the $n$-tuples
$u=(u_1,\dots,u_n)$ of families $u_k=(u_{k,\sigma})_{\sigma\in\Psi'}$ of vectors of $\mathbf C^r$, and for every
$w\in\mathbb W$ let $u'(w)\in U'$ have $u'(w)_{k,\bar\rho}$ equal to the complex conjugate of the $k$-th column of
$\theta_\rho(w)$, as in Situation~\ref{is-eigenvalue-conjugate-frame-context}(2). Then there is a continuous additive map
$R\colon U'\to\mathbb W$ such that:
\begin{enumerate}
\item $u'$ is continuous and additive;
\item for $v\in U'$, $R(v)$ has finite component $0$ and $u'(R(v))=v$; and $R(u'(w))=w$ for every $w$ with finite
  component $0$;
\item $w-R(u'(w))$ has archimedean component $0$ for every $w\in\mathbb W$, and $u'(w)=0$ if $w$ has archimedean
  component $0$;
\item for $\rho\in\mathcal T$, a complex $r\times n$ matrix $z$ and $t\in\mathbf R$, $L_{\rho,z}(t)=R(tv_z)$, where
  $L_{\rho,z}$ is the real direction of Lemma~\ref{is-eigenvalue-conjugate-theta-closed} and $v_z\in U'$ is concentrated
  at $\bar\rho$, with $(v_z)_{k,\bar\rho}$ the complex conjugate of the $k$-th column of $z$;
\item for $g\in U(S)(\mathbf A)$, $h\in U(T)(\mathbf A)$ and $w\in\mathbb W$, $hwg^{-1}$ has finite component $0$ if $w$
  has, and archimedean component $0$ if $w$ has, and $\theta_\rho(hwg^{-1})=\theta_\rho(h)\,\theta_\rho(w)\,
  \theta_\rho(g^{-1})$ for every $\rho\in\mathcal T$.
\end{enumerate}
\end{lemma}
\leanref{Development/AbelianVarietyTheory/Proofs/SiegelThetaTransformation/RealPointsClassicalModelRealVectorSpaceStatement.lean}{GlobalTheory.T.realPoints_classicalModel_realVectorSpace}
\leanref{Development/AbelianVarietyTheory/Proofs/SiegelThetaTransformation/RealPointsClassicalModelRealVectorSpace.lean}{GlobalTheory.KerReal.real_points}
\begin{proof}
As in the proof of Lemma~\ref{is-eigenvalue-conjugate-frame-glue}, $z\mapsto(\theta_\rho(z))_{\rho\in\mathcal T}$ is a
bijective ring homomorphism from $\mathbf A_{F_h,\infty}\otimes_{F_h}E_h$ onto $\mathbf C^{\mathcal T}$, and it is
$\mathbf R$-linear, so its inverse is continuous. For $v\in U'$ let $R(v)\in\mathbb W$ be the element with finite
component $0$ whose archimedean component has, under each $\theta_\rho$, the matrix whose $k$-th column is the complex
conjugate of $v_{k,\bar\rho}$. Then $R$ is continuous and additive, and $u'(R(v))=v$.

(1) The map $u'$ reads the archimedean component of $w$ through the continuous ring homomorphisms $\theta_\rho$ and
complex conjugation.

(2) The first part holds by construction. If $w$ has finite component $0$, then $R(u'(w))$ and $w$ have finite
components $0$ and the same images under every $\theta_\rho$, so they are equal by the bijectivity.

(3) The elements $w$ and $R(u'(w))$ have the same images under every $\theta_\rho$, so the archimedean component of
their difference is $0$ by the bijectivity. If $w$ has archimedean component $0$, every $\theta_\rho(w)$ is $0$.

(4) The map $t\mapsto R(tv_z)$ is continuous and additive, has finite component $0$, image $tz$ under $\theta_\rho$ and
$0$ under the other $\theta_{\rho'}$. A real direction is determined by these images, by the bijectivity, so it is
$L_{\rho,z}$.

(5) The finite and archimedean components of $hwg^{-1}$ are the products of those of $h$, $w$ and $g^{-1}$, and
$\theta_\rho$ is a ring homomorphism on the archimedean components.
\end{proof}
% KerReal.real_points: RealLineAt.exists_inverse_chart (sigma, with Chart.fin, Chart.comp, Chart.ext and
% Chart.cont), realOf C sigma for R, realPt_add, infPart_eq_zero for (3), isLineAt_apply_unique for (4), and
% map_actO for (5).

\begin{lemma}\label{is-eigenvalue-conjugate-frame-zak-forms}
In Situation~\ref{is-eigenvalue-conjugate-frame-context}, let $u'$ be as in
Lemma~\ref{is-eigenvalue-conjugate-frame-zak-points}, and $\hat y$ and $E_{X'}$ as in
Lemma~\ref{is-eigenvalue-conjugate-frame-zak-real} for $X'$; let $\langle\ ,\ \rangle$ be the pairing on $\mathbb W$
(Definition~\ref{th-weil-datum}). For $y\in M_{n\times r}(E_h)$ write $y^{\mathsf T}$ also for the rational point of
$\mathbb W$ with both components $y^{\mathsf T}$, and let $y_f\in\mathbb W$ have archimedean component $0$ and finite
component $y^{\mathsf T}$. Let $y,y'\in M_{n\times r}(E_h)$ and $w,w',w''\in\mathbb W$.
\begin{enumerate}
\item If $\langle y^{\mathsf T},y'^{\mathsf T}\rangle=t$ for some $t\in F_h$, then $\sum_k\langle y_k,y'_k\rangle_{X'}=
  -\operatorname{Tr}_{F_h/\mathbf Q}(t)$.
\item If $w$ and $w'$ have finite components $0$, then
  $\psi(\frac12\langle w,w'\rangle)=e^{\pi iE_{X'}(u'(w),u'(w'))}$.
\item $\psi(\frac12\langle y_f,y'_f\rangle)=e^{-\pi i\sum_k\langle y_k,y'_k\rangle_{X'}}$.
\item If $w$ has finite component $0$ and $u$ has archimedean component $0$, then $\langle w,u\rangle=\langle u,w\rangle=0$.
\item $\langle w,w\rangle=0$.
\item $\langle w+w',w''\rangle=\langle w,w''\rangle+\langle w',w''\rangle$ and $\langle w,w'\rangle=-\langle w',w\rangle$.
\item $\langle hwg^{-1},hw'g^{-1}\rangle=\langle w,w'\rangle$ for $g\in U(S)(\mathbf A)$ and $h\in U(T)(\mathbf A)$.
\item $u'(y^{\mathsf T})=\hat y$, $y_f$ has archimedean component $0$, $y^{\mathsf T}-y_f$ has finite component $0$, and
  $(y+y')_f=y_f+y'_f$.
\end{enumerate}
\end{lemma}
\leanref{Development/AbelianVarietyTheory/Proofs/SiegelThetaTransformation/SymplecticCharacterRealFinitePointsMatchStatement.lean}{GlobalTheory.T.symplecticCharacter_real_finite_points_match}
\leanref{Development/AbelianVarietyTheory/Proofs/SiegelThetaTransformation/SymplecticCharacterRealFinitePointsMatch.lean}{GlobalTheory.FormMatch.forms_match}
\begin{proof}
Write $h(x,x')=\operatorname{tr}(\bar x'^{\mathsf T}G_Tx)$ for $x,x'\in\mathbb W$, as $G_S=1$, so that
$\langle x,x'\rangle=\operatorname{Tr}_{\mathbf A_{E_h}/\mathbf A_{F_h}}(\xi_0h(x,x'))$, computed place by place; and
$\overline{h(x,x')}=h(x',x)$, $\bar\xi_0=-\xi_0$. The form of $X'$, of matrix $G_T^{\mathsf T}$, is
$h'(u,v)=\bar v^{\mathsf T}G_Tu$, so $\sum_kh'(y_k,y'_k)=h(y^{\mathsf T},y'^{\mathsf T})$, and
$\langle u,v\rangle_{X'}=\operatorname{Tr}_{E_h/\mathbf Q}(-\xi_0h'(u,v))$, as $\kappa'_0=2\xi_0$.

(1) $\langle y^{\mathsf T},y'^{\mathsf T}\rangle$ is the rational element $\operatorname{Tr}_{E_h/F_h}(\xi_0h(y^{\mathsf
T},y'^{\mathsf T}))$, so it is $t$, and $\sum_k\langle y_k,y'_k\rangle_{X'}=-\operatorname{Tr}_{F_h/\mathbf
Q}\operatorname{Tr}_{E_h/F_h}(\xi_0h(y^{\mathsf T},y'^{\mathsf T}))=-\operatorname{Tr}_{F_h/\mathbf Q}(t)$.

(2) The pairing $\langle w,w'\rangle$ has finite component $0$, and its image under $\rho\in\mathcal T$ is
$2\operatorname{Re}\big(\rho(\xi_0)\operatorname{tr}(W'^*_\rho G_T(\rho)W_\rho)\big)$, with $W_\rho=\theta_\rho(w)$ and
$W'_\rho=\theta_\rho(w')$, as $\theta_\rho$ extends $\rho$. As $\psi$ is standard,
$\psi(\frac12\langle w,w'\rangle)=\exp\big(-2\pi i\sum_\rho\operatorname{Re}(\rho(\xi_0)\operatorname{tr}(W'^*_\rho
G_T(\rho)W_\rho))\big)$. On the other side, with the $k$-th columns $W_{\rho,k}$ and $W'_{\rho,k}$,
$h'_{\bar\rho}(\overline{W_{\rho,k}},\overline{W'_{\rho,k}})=\overline{W'^*_{\rho,k}G_T(\rho)W_{\rho,k}}$ and
$\bar\rho(-\xi_0)=-\overline{\rho(\xi_0)}$, so $E_{X'}(u'(w),u'(w'))=-2\sum_\rho\operatorname{Re}\big(\rho(\xi_0)
\operatorname{tr}(W'^*_\rho G_T(\rho)W_\rho)\big)$.

(3) The pairing $\langle y_f,y'_f\rangle$ has archimedean component $0$ and finite component the rational element
$s=\operatorname{Tr}_{E_h/F_h}(\xi_0h(y^{\mathsf T},y'^{\mathsf T}))$. As $\psi$ is trivial on $F_h$ and standard,
$\psi$ of the element with archimedean component $0$ and finite component $x\in F_h$ equals $\psi$ of the element with
archimedean component $-x$ and finite component $0$, which is $e^{2\pi i\operatorname{Tr}_{F_h/\mathbf Q}(x)}$. So
$\psi(\frac12\langle y_f,y'_f\rangle)=e^{\pi i\operatorname{Tr}_{F_h/\mathbf Q}(s)}=e^{-\pi i\sum_k\langle
y_k,y'_k\rangle_{X'}}$, by the computation in (1).

(4) The pairing is computed place by place, and at every place one of the two arguments is $0$.

(5) and (6) The pairing is additive in each variable. As $\overline{\xi_0h(x,x')}=-\xi_0h(x',x)$ and the trace commutes
with conjugation, $\langle x,x'\rangle=-\langle x',x\rangle$; for $x=x'$ this gives $\langle x,x\rangle=0$.

(7) As $\bar g^{\mathsf T}g=1$ and $\bar h^{\mathsf T}G_Th=G_T$, $h(hxg^{-1},hx'g^{-1})=\operatorname{tr}(g\,\bar x'^{\mathsf
T}G_Txg^{-1})=h(x,x')$.

(8) The archimedean component of $y^{\mathsf T}$ has image $\rho(y)^{\mathsf T}$ under $\theta_\rho$, whose $k$-th
column has complex conjugate $\bar\rho(y_k)=\hat y_{k,\bar\rho}$. The other three assertions hold by the definitions.
\end{proof}
% FormMatch.forms_match: (1) symplO_rat and sum_EF (hRat_eq, C.m.ka); (2) OrientedChar.char_symplOH (through the inverse
% chart and real_adele_ext_and_char) and EFR_realPt (omO, the real form of the archimedean components); (3) symplOH_finW
% and StdRat.stdChar_fin_rat; (4) the last clause of OrientedFormLaws.symplOH_laws; (5) and (6)
% SymplecticOrthogonalSesquilinear.symplO_alt (symplO_skew); (7) hO_actO; (8) piW_ratW, finRat_inf and finRat_fin.

\begin{lemma}\label{is-eigenvalue-conjugate-frame-zak-space}
In Situation~\ref{is-eigenvalue-conjugate-frame-context}, let $f$ and $f'$ be theta functions for $(S,T)$
(Definition~\ref{th-weil-datum}(2)). For $y\in M_{n\times r}(E_h)$ let $y_f\in\mathbb W$ have archimedean
component $0$ and finite component $y^{\mathsf T}$, write $y^{\mathsf T}$ also for the rational point of $\mathbb W$, and
let $\hat y$ and $E_{X'}$ be as in Lemma~\ref{is-eigenvalue-conjugate-frame-zak-real} for $X'$, and $u'$ as in
Lemma~\ref{is-eigenvalue-conjugate-frame-zak-points}. Let $\hat\Lambda\subset\mathbb W$ be the subgroup generated by the
elements with archimedean component $0$ and finite component $\alpha y^{\mathsf T}$, with $\alpha$ in the closure of the
integers of $F_h$ in $\mathbf A_{F_h,f}$ and $y\in\mathcal O_{E_h}^n\otimes\Lambda'$
(Definition~\ref{is-eigenvalue-conjugate-frame-heisenberg}(4)).
\begin{enumerate}
\item $0$, $f+f'$ and $af$, for $a\in\mathbf C$, are theta functions; and if $f(w)=f'(w)$ for every $w\in\mathbb W$
  with finite component $0$, then $f=f'$.
\item For $w_0,w_1\in\mathbb W$, $\rho(w_0)f$ is a theta function, and $\rho(w_0)\rho(w_1)f=\psi(\frac12\langle
  w_0,w_1\rangle)\,\rho(w_0+w_1)f$.
\item For $y\in M_{n\times r}(E_h)$ and $w\in\mathbb W$ with finite component $0$,
  $(\rho(y_f)f)(w)=e^{\pi iE_{X'}(u'(w),\hat y)}f(w-y^{\mathsf T}+y_f)$.
\item For $\gamma\in U(S)(F_h)$ and $\gamma'\in U(T)(F_h)$, the function $w\mapsto f(\gamma'w\gamma^{-1})$ is a theta
  function.
\item There is an integer $N\ge1$ with $\rho(Nu)f=f$ for every $u\in\hat\Lambda$.
\end{enumerate}
\end{lemma}
\leanref{Development/AbelianVarietyTheory/Proofs/SiegelThetaTransformation/ThetaFunctionSpaceHeisenbergLawsStatement.lean}{GlobalTheory.T.thetaFunction_space_heisenberg_laws}
\begin{proof}
(1) This is Lemma~\ref{is-eigenvalue-conjugate-frame-theta-space}.

(2) The first assertion is Lemma~\ref{is-eigenvalue-conjugate-theta-closed}(1). By Definition~\ref{th-weil-datum},
$\rho(w_0)\rho(w_1)f(x)=\psi(\frac12\langle x,w_0\rangle+\frac12\langle x+w_0,w_1\rangle)f(x+w_0+w_1)$ and
$\rho(w_0+w_1)f(x)=\psi(\frac12\langle x,w_0+w_1\rangle)f(x+w_0+w_1)$; the arguments of $\psi$ differ by
$\frac12\langle w_0,w_1\rangle$.

(3) Put $z=y^{\mathsf T}-y_f$, which has finite component $0$, with $u'(z)=\hat y$, by
Lemma~\ref{is-eigenvalue-conjugate-frame-zak-forms}(8) and Lemma~\ref{is-eigenvalue-conjugate-frame-zak-points}(1)
and~(3). By Lemma~\ref{is-eigenvalue-conjugate-frame-zak-forms}(4) to~(6), $\langle w,y_f\rangle=0$, $\langle
z,z\rangle=0$ and $\langle y_f,w-z\rangle=0$, as $w-z$ has finite component $0$. Then
\[
 (\rho(y_f)f)(w)=f(w+y_f)=f\big(y^{\mathsf T}+(w-z)\big)=\psi\big(-\tfrac12\langle y^{\mathsf T},w-z\rangle\big)f(w-z),
\]
by the quasi-periodicity, and $-\langle y^{\mathsf T},w-z\rangle=-\langle y_f+z,w-z\rangle=-\langle z,w\rangle=\langle
w,z\rangle$. By Lemma~\ref{is-eigenvalue-conjugate-frame-zak-forms}(2), $\psi(\frac12\langle w,z\rangle)=
e^{\pi iE_{X'}(u'(w),\hat y)}$.

(4) Put $F(w)=f(\gamma'w\gamma^{-1})$ and $A(w)=\gamma'w\gamma^{-1}$, a continuous additive map, which preserves the pairing
(Lemma~\ref{is-eigenvalue-conjugate-frame-zak-forms}(7)), the rational points, and the points with finite or with
archimedean component $0$ (Lemma~\ref{is-eigenvalue-conjugate-frame-zak-points}(5)). For a rational $x$,
$F(x+w)=f(A(x)+A(w))=\psi(-\frac12\langle A(x),A(w)\rangle)f(A(w))=\psi(-\frac12\langle x,w\rangle)F(w)$. If $U$ is an open
subgroup for $f$ as in Definition~\ref{th-weil-datum}(2), then $A^{-1}(U)$ is one for $F$: for $u\in A^{-1}(U)$ with
archimedean component $0$, $(\rho(u)F)(w)=\psi(\frac12\langle A(w),A(u)\rangle)f(A(w)+A(u))=(\rho(A(u))f)(A(w))=F(w)$.
$F$ is continuous, and for a continuous additive $\ell\colon\mathbf R^a\to\mathbb W$ with finite component $0$,
$t\mapsto F(w+\ell(t))=f(A(w)+A(\ell(t)))$ is smooth, as $A\circ\ell$ has the same properties.

(5) Let $U$ be an open subgroup for $f$ as in Definition~\ref{th-weil-datum}(2). By
Lemma~\ref{is-eigenvalue-conjugate-frame-zak-exists-lattice}(4) for $\mathcal O_{E_h}^n$ and $\Lambda'$, which are finitely
generated and span, there is $N\ge1$ with $Nu\in U$ for every $u\in\hat\Lambda$; and $Nu$ has archimedean component $0$
(Lemma~\ref{is-eigenvalue-conjugate-frame-zak-exists-lattice}(3)), so $\rho(Nu)f=f$.
\end{proof}
% T.thetaFunction_space_heisenberg_laws = ThetaSpace.theta_space_laws (149 lines in its file), followed: isTheta_zero,
% isTheta_add, isTheta_smul, isTheta_sub and eq_zero_of_real for (1) (Lemma -theta-space); isTh_rho (Lemma
% theta-closed(1)) and heisO_heisO for (2); rho_finW for (3), from the clauses of
% T.symplecticCharacter_real_finite_points_match (Lemma -zak-forms) and T.realPoints_classicalModel_realVectorSpace
% (Lemma -zak-points); isTh_actO for (4) (exists_mt_eq for the rational point, the open subgroup U.comap A, hf.smooth);
% exists_level for (5) (h4 of T.completedTensorLattice_strongApprox_open_compact, Lemma -zak-exists-lattice(4)).

\begin{lemma}\label{is-eigenvalue-conjugate-frame-zak-exists-lattice}
In Situation~\ref{is-eigenvalue-conjugate-frame-context}, let $\Lambda\subset E_h^n$ and $L\subset E_h^r$ be finitely
generated $\mathcal O_{E_h}$-submodules that span $E_h^n$ and $E_h^r$ over $E_h$, put $A=\Lambda\otimes L$
(Definition~\ref{is-eigenvalue-conjugate-frame-heisenberg}(4)), and let $\hat\Lambda(A)\subset\mathbb W$ be the subgroup
generated by the elements with archimedean component $0$ and finite component $\alpha y^{\mathsf T}$, with $\alpha$ in the
closure $\hat{\mathcal O}$ of the integers of $F_h$ in $\mathbf A_{F_h,f}$ and $y\in A$. For $y\in M_{n\times r}(E_h)$ let
$y_f\in\mathbb W$ have archimedean component $0$ and finite component $y^{\mathsf T}$.
\begin{enumerate}
\item For every $u\in\mathbb W$ with archimedean component $0$ there is $y\in M_{n\times r}(E_h)$ with $u-y_f\in\hat\Lambda(A)$.
\item $y_f\in\hat\Lambda(A)$ if and only if $y\in A$.
\item Every $u\in\hat\Lambda(A)$ has archimedean component $0$, and for every integer $N\ge1$ there are $y\in A$ and
  $u'\in\hat\Lambda(A)$ with $u=y_f+Nu'$.
\item For every open subgroup $U$ of $\mathbb W$ there is an integer $N\ge1$ with $Nu\in U$ for every $u\in\hat\Lambda(A)$.
\item For every integer $N\ge1$, the set of the $w\in\mathbb W$ such that $w-Nu$ has finite component $0$ for some
  $u\in\hat\Lambda(A)$ is open.
\item $\hat\Lambda(A)$ is compact.
\end{enumerate}
\end{lemma}
\leanref{Development/AbelianVarietyTheory/Proofs/SiegelThetaTransformation/CompletedTensorLatticeStrongApproxOpenCompactStatement.lean}{GlobalTheory.T.completedTensorLattice_strongApprox_open_compact}
\begin{proof}
\openstep{The proofs, from strong approximation for the additive group of $F_h$, the density of $\mathcal O_{F_h}$ in
$\hat{\mathcal O}$ and its compactness, applied to the coordinates of $M_{n\times r}(E_h)$ in a basis contained in $A$, are not
written out.}
% T.completedTensorLattice_strongApprox_open_compact = Ker.HatLat.hatA_tens (CompletedTensorLatticeStrongApproxOpenCompact.lean).
% Clauses (1), (3) and (4) for stdL and C.X.Lambda are the facts (a), (b) and (c) of the proof of Lemma
% is-eigenvalue-conjugate-frame-theta-space, where they are derived from strong approximation. Size, following imports and
% stopping at statements the book cites: about 600 lines.
\end{proof}

\begin{lemma}\label{is-eigenvalue-conjugate-frame-zak-exists-extension}
Let $D$ be a definite-side datum, $\psi$ a standard character of $\mathbf A_{F_h}$, and $S$ and $T$ Hermitian spaces of
$D$, with $\mathbb W$ and $\langle\ ,\ \rangle$ as in Definition~\ref{th-weil-datum}; call the elements of $\mathbb W$
with both components in $M_{\operatorname{rk}T\times\operatorname{rk}S}(E_h)$ \emph{rational}. Let $V$ be a real normed
space, $p\colon\mathbb W\to V$ a continuous additive map, $F\colon V\to\mathbf C$ a smooth function, and $O\subset\mathbb W$
an open subgroup that contains every element with finite component $0$, such that:
\begin{enumerate}
\item[(a)] for every $w\in\mathbb W$ there is a rational $\gamma$ with $w-\gamma\in O$;
\item[(b)] $F(p(w+\delta))=\psi(-\frac12\langle\delta,w\rangle)F(p(w))$ for every rational $\delta\in O$ and every $w\in O$;
\item[(c)] for every $u\in O$ with archimedean component $0$, $p(u)=0$ and $\psi(\frac12\langle w,u\rangle)=1$ for every
  $w\in O$.
\end{enumerate}
Then there is a theta function $f$ for $(S,T)$ with $f(w)=F(p(w))$ for every $w$ with finite component $0$.
\end{lemma}
\leanref{Development/AbelianVarietyTheory/Proofs/SiegelThetaTransformation/AdelicZakThetaExistsUnique.lean}{GlobalTheory.Ker.ZakEx.exists_extension}
\begin{proof}
For rational $\gamma,\gamma'$, $\frac12\langle\gamma,\gamma'\rangle$ lies in $F_h$, so $\psi(\frac12\langle\gamma,\gamma'\rangle)=1$.
For each $w$ choose a rational $c(w)$ with $w-c(w)\in O$, and put
$f(w)=\psi(-\frac12\langle c(w),w\rangle)F(p(w-c(w)))$.

\emph{Independence.} If $\gamma$ is rational and $w-\gamma\in O$, then $f(w)=\psi(-\frac12\langle\gamma,w\rangle)F(p(w-\gamma))$.
Indeed, $\delta=\gamma-c(w)$ is rational and lies in $O$, so by (b) at $w-\gamma\in O$,
$F(p(w-c(w)))=\psi(-\frac12\langle\delta,w-\gamma\rangle)F(p(w-\gamma))$; and
$-\frac12\langle c(w),w\rangle-\frac12\langle\delta,w-\gamma\rangle=-\frac12\langle\gamma,w\rangle+\frac12\langle\gamma,\gamma\rangle-
\frac12\langle c(w),\gamma\rangle$, where the last two terms are pairings of rational elements.

\emph{A theta function.} For rational $\gamma'$, take $\gamma=\gamma'+c(w)$ for the point $\gamma'+w$: then
$f(\gamma'+w)=\psi(-\frac12\langle\gamma'+c(w),\gamma'+w\rangle)F(p(w-c(w)))=\psi(-\frac12\langle\gamma',w\rangle)f(w)$, after
removing the pairings of rational elements. For $u\in O$ with archimedean component $0$, $w+u-c(w)\in O$, so
$f(w+u)=\psi(-\frac12\langle c(w),w+u\rangle)F(p(w-c(w)))$ by (c); and $\psi(\frac12\langle w,u\rangle-\frac12\langle
c(w),u\rangle)=\psi(\frac12\langle w-c(w),u\rangle)=1$ by (c), so $(\rho(u)f)(w)=f(w)$. Near a point $w_0$, $w-c(w_0)\in O$, as
$O$ is open, so $f(w)=\psi(-\frac12\langle c(w_0),w\rangle)F(p(w-c(w_0)))$ there, and $f$ is continuous. For a continuous
additive map $\ell\colon\mathbf R^a\to\mathbb W$ with finite component $0$, $\ell(t)\in O$, so
$f(w+\ell(t))=\psi(-\frac12\langle c(w),w\rangle)\psi(-\frac12\langle c(w),\ell(t)\rangle)F\big(p(w-c(w))+p(\ell(t))\big)$; here
$t\mapsto\psi(-\frac12\langle c(w),\ell(t)\rangle)$ is a continuous unitary character of $\mathbf R^a$, so of the form
$e^{i\lambda(t)}$ with $\lambda$ linear, and $p\circ\ell$ is continuous and additive, hence linear. So $t\mapsto f(w+\ell(t))$ is
smooth, and $f$ is a theta function.

\emph{The values.} For $w$ with finite component $0$, $w\in O$, so $\gamma=0$ gives $f(w)=F(p(w))$.
\end{proof}
% Ker.ZakEx.exists_extension (AdelicZakThetaExistsUnique.lean), followed: c and hc, hcO by choice, key (with psi_rat, sym_sub_right
% and sym_sub_left), the quasi-periodicity, the open subgroup O, continuity through continuousAt, and smoothness with exists_log
% (the archimedean values of psi as the exponential of a continuous additive map) and contDiff_additive.

\begin{lemma}\label{is-eigenvalue-conjugate-frame-zak-exists}
In Situation~\ref{is-eigenvalue-conjugate-frame-context}, let $x'\in\mathcal D_{X'}$ and $\varphi$ a Schwartz-Bruhat
function on $M_{n\times d}(\mathbf Q)$. Then there is a Zak theta function for $x'$ and $\varphi$.
\end{lemma}
\leanref{Development/AbelianVarietyTheory/Proofs/SiegelThetaTransformation/AdelicZakThetaExistsUnique.lean}{GlobalTheory.Ker.ZakEx.exists_isZak}
\leanref{Development/AbelianVarietyTheory/Proofs/SiegelThetaTransformation/HeisGFixOfMemLattice.lean}{GlobalTheory.Ker.heisG_fix_of_mem_lattice}
\begin{proof}
Let $u'$, $R$ be as in Lemma~\ref{is-eigenvalue-conjugate-frame-zak-points}, $E_{X'}$ and $\hat y$ as in
Lemma~\ref{is-eigenvalue-conjugate-frame-zak-real} for $X'$, $H=\hat\Lambda(\mathcal O_{E_h}^n\otimes\Lambda')$ as in
Lemma~\ref{is-eigenvalue-conjugate-frame-zak-exists-lattice}, and call $w\in\mathbb W$ real if its finite component is $0$.

\emph{A level.} Let $R_0$ be a level of $\varphi$ and $N=2R_0$. If $y$ has its rows in $\Lambda'$, then $\rho_{X'}(jNy)\varphi=\varphi$
for every integer $j\ge0$. Indeed, the coordinates $(\alpha,\beta)$ of $y$ are integral, as $\Lambda'$ is spanned over
$\mathbf Z$ by the symplectic basis; those of $jNy$ are $jN\alpha$ and $jN\beta$; $\varphi(m-jN\beta)=\varphi(m)$; and where
$\varphi(m)\ne0$, $R_0m$ is integral and $2[m,jN\alpha]-[jN\alpha,jN\beta]=4j[R_0m,\alpha]-j^2N^2[\alpha,\beta]$ is an even
integer.

\emph{The pairing on $H$.} For $u,u'\in H$, $\psi(\langle u,u'\rangle)=1$. As $\psi(\langle\cdot,\cdot\rangle)$ is
multiplicative in each variable, it suffices to take $u=\alpha y_f$ and $u'=\beta y'_f$ with rows of $y,y'$ in $\Lambda'$ and
$\alpha,\beta\in\hat{\mathcal O}$. For $a\in\mathcal O_{F_h}$, $\psi(\langle(ay)_f,y'_f\rangle)=\psi(\frac12\langle(ay)_f,y'_f\rangle)^2
=e^{-2\pi i\sum_k\langle ay_k,y'_k\rangle_{X'}}=1$ (Lemma~\ref{is-eigenvalue-conjugate-frame-zak-forms}(3)), as the form of
$X'$ is integral on $\Lambda'$; by continuity of $\psi$ and density of $\mathcal O_{F_h}$ in $\hat{\mathcal O}$, as in the
proof of Lemma~\ref{is-eigenvalue-conjugate-frame-vacuum-line-signs}, $\psi(\langle u,u'\rangle)=1$.

\emph{The open subgroup.} Let $O$ be the set of the $w$ such that $w-2Nu$ is real for some $u\in H$. It is a subgroup, open
by Lemma~\ref{is-eigenvalue-conjugate-frame-zak-exists-lattice}(5), and it contains the real points. Every $u\in H$ has
archimedean component $0$ (Lemma~\ref{is-eigenvalue-conjugate-frame-zak-exists-lattice}(3)). For $w=r+2Nu$ and $w'=r'+2Nu'$
in $O$, with $r,r'$ real and $u,u'\in H$,
\[
 \psi(\tfrac12\langle w,w'\rangle)=\psi(\tfrac12\langle r,r'\rangle)\,\psi(\tfrac12\langle2Nu,2Nu'\rangle)
 =e^{\pi iE_{X'}(u'(r),u'(r'))}=e^{\pi iE_{X'}(u'(w),u'(w'))},
\]
by Lemma~\ref{is-eigenvalue-conjugate-frame-zak-forms}(4) and~(2), the pairing on $H$, and $u'(2Nu)=0$
(Lemma~\ref{is-eigenvalue-conjugate-frame-zak-points}(3)).

\emph{The extension.} Apply Lemma~\ref{is-eigenvalue-conjugate-frame-zak-exists-extension} to $V=U'$, $p=u'$, which is
continuous and additive (Lemma~\ref{is-eigenvalue-conjugate-frame-zak-points}(1)), the smooth function
$F=\Theta'_\varphi(x';\cdot)$ (Lemma~\ref{is-eigenvalue-conjugate-frame-zak-real}(1) for $X'$), and $O$.
\begin{enumerate}
\item[(a)] Let $w\in\mathbb W$. The element $w-R(u'(w))$ has archimedean component $0$
  (Lemma~\ref{is-eigenvalue-conjugate-frame-zak-points}(3)), so by
  Lemma~\ref{is-eigenvalue-conjugate-frame-zak-exists-lattice}(1) and~(3) it is $y_f+y'_f+2Nu$ with $y$, $y'$ matrices and $u\in
  H$. For the rational $\gamma=(y+y')^{\mathsf T}$, $w-\gamma-2Nu=R(u'(w))-\big((y+y')^{\mathsf T}-(y+y')_f\big)$ is real
  (Lemma~\ref{is-eigenvalue-conjugate-frame-zak-forms}(8)); so $w-\gamma\in O$.
\item[(b)] Let $\delta=y^{\mathsf T}$ be rational and in $O$, with $\delta-2Nu$ real, $u\in H$. Then $y_f-2Nu$ is real, as it is
  $(\delta-2Nu)-(y^{\mathsf T}-y_f)$, and has archimedean component $0$; so $y_f=2Nu$. Writing $y=2Ny_0$, $(y_0)_f=u\in H$, so
  $y_0\in\mathcal O_{E_h}^n\otimes\Lambda'$ (Lemma~\ref{is-eigenvalue-conjugate-frame-zak-exists-lattice}(2)) and its rows lie in
  $\Lambda'$ (Lemma~\ref{is-eigenvalue-conjugate-frame-vacua-tensor}(4)); so $\rho_{X'}(y)\varphi=\varphi$ by the level. For
  $w\in O$, $u'(w+\delta)=u'(w)+\hat y$ (Lemma~\ref{is-eigenvalue-conjugate-frame-zak-forms}(8) and
  Lemma~\ref{is-eigenvalue-conjugate-frame-zak-points}(1)), and by Lemma~\ref{is-eigenvalue-conjugate-frame-zak-real}(2),
  $\Theta'_\varphi(x';u'(w)+\hat y)=e^{\pi iE_{X'}(u'(w),\hat y)}\Theta'_{\rho_{X'}(y)\varphi}(x';u'(w))=e^{\pi iE_{X'}(u'(w),\hat
  y)}\Theta'_\varphi(x';u'(w))$. As $\delta\in O$, the pairing formula above gives
  $\psi(-\frac12\langle\delta,w\rangle)=\psi(\frac12\langle w,\delta\rangle)=e^{\pi iE_{X'}(u'(w),\hat y)}$.
\item[(c)] For $u\in O$ with archimedean component $0$, $u'(u)=0$
  (Lemma~\ref{is-eigenvalue-conjugate-frame-zak-points}(3)), and by the pairing formula $\psi(\frac12\langle
  w,u\rangle)=e^{\pi iE_{X'}(u'(w),0)}=1$ for $w\in O$.
\end{enumerate}
The lemma gives a theta function $f$ with $f(w)=\Theta'_\varphi(x';u'(w))$ for every real $w$: a Zak theta function for $x'$ and
$\varphi$.
\end{proof}
% Ker.ZakEx.exists_isZak (AdelicZakThetaExistsUnique.lean), followed: Ker.heisG_fix_of_mem_lattice (level 2R, through
% TensVac.heisG_eq), thR_laws (Lemma -zak-real: hsm, hth), T.completedTensorLattice_strongApprox_open_compact (h1, h2, h3, h5;
% Lemma -zak-exists-lattice), T.symplecticCharacter_real_finite_points_match (s2, s8; Lemma -zak-forms),
% T.realPoints_classicalModel_realVectorSpace (Lemma -zak-points), psi_hat (the pairing on H), hB, thick H (2N) for O, and
% ZakEx.exists_extension (Lemma -zak-exists-extension).

\begin{lemma}\label{is-eigenvalue-conjugate-frame-zak-fock}
In Situation~\ref{is-eigenvalue-conjugate-frame-context}, let $B$ be orthonormal for $X'$, and let $G$ be a
function on $\mathbb W$ and $\Phi$ a real-differentiable function on $(\mathbf C^d)^n$ with
$G(w)=\mathcal G'_{x'_B}(z)\,\Phi(z)$, $z=z'_{x'_B}(u'(w))$, for every $w$ with finite component $0$.
\begin{enumerate}
\item Let $\rho\in\mathcal T$ and $z_0$ a complex $r\times n$ matrix, put $Jz_0=i\,s'_{B,\rho}z_0$, and let
  $u_0\in(\mathbf C^d)^n$ be the coordinates $z'_{x'_B}(\bar z_0^{\mathsf T})$ of $z_0$ concentrated at $\rho$
  (Lemma~\ref{is-eigenvalue-conjugate-frame-zak}). Then for every $w$ with finite component $0$, with $z=z'_{x'_B}(u'(w))$
  and the derivatives $\partial_\rho$ of Lemma~\ref{is-eigenvalue-conjugate-theta-closed}, the annihilation operator
  $\partial_\rho(z_0)+i\,\partial_\rho(Jz_0)$ for the Cartan involutions $1$ of $S$ and $s'_{B,\rho}$ of $T$
  (Lemma~\ref{is-eigenvalue-conjugate-theta-facts}), applied to $G$, has the value
  $\mathcal G'_{x'_B}(z)\big(d\Phi_z(u_0)+i\,d\Phi_z(iu_0)\big)$ at $w$; and $\partial_\rho(z_0)-i\,\partial_\rho(Jz_0)$,
  applied to $G$, has the value
  \[
   \mathcal G'_{x'_B}(z)\cdot2\Big(-\pi\sum_kH'_{x'_B}(u_{0,k},z_k)\,\Phi(z)+\tfrac12\big(d\Phi_z(u_0)-i\,d\Phi_z(iu_0)\big)\Big)
  \]
  at $w$, where $H'$ is the form $H$ of $X'$.
\item If $\Phi$ is holomorphic, then for every polynomial $Q$ in the creation variables
  $(G\otimes_{s,B}Q)(0)$, with $s_S=1$, $s_T=s'_B$, $B_S=1$ and $B_T=B$, is the value at $0$ of the derivative of
  $\Phi$ by the polynomial in the variables $Y_{k,j}$ obtained from $Q$ by the substitution
  $(\rho,a,k)\mapsto(2/f_\rho)Y_{k,j}$ of Lemma~\ref{is-eigenvalue-conjugate-frame-zak}, $Y_{k,j}$ acting as
  $\partial/\partial z_{k,j}$.
\end{enumerate}
\end{lemma}
\leanref{Development/AbelianVarietyTheory/Proofs/SiegelThetaTransformation/AnnihilationCreationPeriodCoordinatesStatement.lean}{GlobalTheory.T.annihilation_creation_periodCoordinates}
\leanref{Development/AbelianVarietyTheory/Proofs/SiegelThetaTransformation/FockPolynomialVectorOriginEqPolyDerivStatement.lean}{GlobalTheory.T.fockPolynomialVector_origin_eq_polyDeriv}
\begin{proof}
(1) Write $x=x'_B$, $\mathcal G=\mathcal G'_x$, $H=H'_x$, and let $R$ be as in
Lemma~\ref{is-eigenvalue-conjugate-frame-zak-points}. \emph{A derivative.} Let $z_1$ be a complex $r\times n$ matrix, $v\in
U'$ the family of $z_1$ concentrated at $\bar\rho$ as in Lemma~\ref{is-eigenvalue-conjugate-frame-zak-points}(4), and
$u_1=z'_x(v)$. Let $w$ have finite component $0$ and put $\zeta=z'_x(u'(w))$. By
Lemma~\ref{is-eigenvalue-conjugate-frame-zak-points}(2) and~(4), $w+L_{\rho,z_1}(t)=R(u'(w)+tv)$, which has finite component
$0$, and $z'_x(u'(w)+tv)=\zeta+tu_1$. By Lemma~\ref{is-eigenvalue-conjugate-frame-zak-forms}(2) and
Lemma~\ref{is-eigenvalue-conjugate-frame-zak-coordinates}(2) for $X'$, $\psi(\frac12\langle w,L_{\rho,z_1}(t)\rangle)=e^{\pi
it\operatorname{Im}\Sigma}$ with $\Sigma=\sum_kH(\zeta_k,u_{1,k})$. So
\[
 \big(\rho(L_{\rho,z_1}(t))G\big)(w)=e^{\pi it\operatorname{Im}\Sigma}\,\mathcal G(\zeta+tu_1)\,\Phi(\zeta+tu_1).
\]
As $\mathcal G(\zeta+tu_1)=\mathcal G(\zeta)\exp\big(-\frac\pi2t\sum_k(H(\zeta_k,u_{1,k})+H(u_{1,k},\zeta_k))+O(t^2)\big)$ and
$H(u_{1,k},\zeta_k)=\overline{H(\zeta_k,u_{1,k})}$ (Lemma~\ref{is-eigenvalue-conjugate-frame-hermitian}(1)), the derivative at
$t=0$ is $\mathcal G(\zeta)\big(-\pi\bar\Sigma\,\Phi(\zeta)+d\Phi_\zeta(u_1)\big)$, with $\bar\Sigma=\sum_kH(u_{1,k},\zeta_k)$. So
$\partial_\rho(z_1)G(w)=\mathcal G(\zeta)\big(d\Phi_\zeta(u_1)-\pi\sum_kH(u_{1,k},\zeta_k)\Phi(\zeta)\big)$.

\emph{The two operators.} For $z_1=z_0$, $u_1=u_0$. For $z_1=Jz_0$, $u_1=iu_0$, by
Lemma~\ref{is-eigenvalue-conjugate-frame-zak}(3); and $H(iu_{0,k},\zeta_k)=iH(u_{0,k},\zeta_k)$. Write
$h=\sum_kH(u_{0,k},\zeta_k)\Phi(\zeta)$. Then
\[
 \big(\partial_\rho(z_0)\pm i\,\partial_\rho(Jz_0)\big)G(w)=\mathcal G(\zeta)\big(d\Phi_\zeta(u_0)-\pi h\pm i\,d\Phi_\zeta(iu_0)\pm\pi h\big),
\]
which is $\mathcal G(\zeta)(d\Phi_\zeta(u_0)+i\,d\Phi_\zeta(iu_0))$ for the sign $+$ and
$\mathcal G(\zeta)\big(-2\pi h+d\Phi_\zeta(u_0)-i\,d\Phi_\zeta(iu_0)\big)$ for the sign $-$.

(2) \openstep{The computation of the value at $0$ of the polynomial vector $G\otimes_{s,B}Q$, by iterating the creation
operators of (1) on $\mathcal G'_{x'_B}\cdot\Phi$ and evaluating at $0$, where the Gaussian factors and the terms with $H'$
drop out, is not written out.}
% T.fockPolynomialVector_origin_eq_polyDeriv = Ker.FockOrigin.fvec_zero_eq_polyDeriv: C.fvec B G Q 0 = polyDeriv (C.toH Q)
% letterDir Phi 0. Size, following imports and stopping at statements the book cites: about 975 lines.
\end{proof}
% (1) is T.annihilation_creation_periodCoordinates = Ker.annih_crea_coord, followed: HeisDeriv.dRho_eq (the derivative,
% through T.realPoints_classicalModel_realVectorSpace, Dict.hx_laws, zR_laws, the clause hpsi of
% T.symplecticCharacter_real_finite_points_match and hasDerivAt_translate) at z and at cplxO 1 (sB B tau) z, with the clause
% zR_cplxO of thetaKernel_orthonormalFrame_inDomain_periodLinear (Lemma -zak(3)). The holomorphy of the theta series, which
% this lemma stated before, is Lemma -zak-series(1).

\begin{lemma}\label{is-eigenvalue-conjugate-frame-zak}
In Situation~\ref{is-eigenvalue-conjugate-frame-context}, for a family $u$ of complex $r\times n$ matrices
$u_\rho$ indexed by $\mathcal T$ write $\bar u^{\mathsf T}$ for the family of $n$-tuples of vectors with
$k$-th vector at $\bar\rho$ the complex conjugate of the $k$-th column of $u_\rho$, and let
$z'_{x'}(\bar u^{\mathsf T})\in(\mathbf C^d)^n$ be its coordinates. For a polynomial $Q$ in the creation
variables $(\rho,a,k)$, $\rho\in\mathcal T$, $1\le a\le r$, $1\le k\le n$, a Schwartz-Bruhat function
$\varphi$ and $x'\in M_r(\mathbf C)^{\Psi'}$, let $D'_{x'}(Q)\varphi$ be the value at $z=0$ of
$Q(\partial)\theta'_\varphi(x';z)$, where $Q(\partial)$ replaces each $(\rho,a,k)$ by
$(2/f_\rho)\,\partial/\partial z_{k,j}$, with $f_\rho=2\sqrt{2\pi|\rho(\xi_0)|}$ and $j$ the index with
$(\tau'_j,a'_j)=(\bar\rho,a)$.
\begin{enumerate}
\item For every $x'\in\mathcal D_{X'}$ and every Schwartz-Bruhat function $\varphi$ on
  $M_{n\times d}(\mathbf Q)$ there is exactly one Zak theta function $\mathrm{Th}(x',\varphi)$ for $x'$ and
  $\varphi$.
\item Let $x'\in\mathcal D_{X'}$ and $\varphi$ be Schwartz-Bruhat. For every Schwartz-Bruhat $\varphi'$,
  $\mathrm{Th}(x',\varphi+\varphi')=\mathrm{Th}(x',\varphi)+\mathrm{Th}(x',\varphi')$, and
  $\mathrm{Th}(x',\varphi)=\mathrm{Th}(x',\varphi')$ only if $\varphi=\varphi'$. For every $a\in\mathbf C$,
  $\mathrm{Th}(x',a\varphi)=a\,\mathrm{Th}(x',\varphi)$. For every $y\in M_{n\times r}(E_h)$,
  $\rho(y_f)\mathrm{Th}(x',\varphi)=\mathrm{Th}(x',\rho_{X'}(-y)\varphi)$, where $y_f\in\mathbb W$ has
  archimedean component $0$ and finite component $y^{\mathsf T}$, and $\rho_{X'}$ is the Heisenberg operator
  of Definition~\ref{is-eigenvalue-conjugate-frame-heisenberg}(1). For every $\gamma\in U(T)(F_h)$,
  $\mathrm{Th}(\gamma x',w_\gamma\varphi)(w)=\mathrm{Th}(x',\varphi)(\gamma^{-1}w)$ for every $w\in\mathbb W$,
  where $(\gamma x')_{\sigma}=\sigma(\gamma)x'_\sigma$ and $w_\gamma\varphi$ is a Schwartz-Bruhat function
  $\psi$ with $\theta'_\psi(\gamma x'';\cdot)=\theta'_\varphi(x'';\cdot)$ for every $x''\in\mathcal D_{X'}$,
  chosen once for all when there is one, and $0$ otherwise.
\item There is a family orthonormal for $X'$. Let $B$ be orthonormal for $X'$. Then $x'_B$ is orthonormal
  for $X'$ and lies in $\mathcal D_{X'}$. For every $\rho$, $s'_{B,\rho}$ is a Cartan involution of
  $G_T(\rho)$, and $s'_{B,\rho}B_\rho=B_\rho J'$; for every complex $r\times n$ matrix $u$,
  $z'_{x'_B}\big(\overline{(is'_{B,\rho}u)}^{\mathsf T}\big)=i\,z'_{x'_B}(\bar u^{\mathsf T})$, the families being
  concentrated at $\rho$; and for $1\le a\le r$ and $1\le k\le n$, $z'_{x'_B}\big(\overline{(B_\rho
  E_{ak})}^{\mathsf T}\big)$, with the matrix unit $E_{ak}$ concentrated at $\rho$, is the unit vector at
  $(k,j)$ for the index $j$ with $(\tau'_j,a'_j)=(\bar\rho,a)$.
\item Let $B$ be orthonormal for $X'$ and $g,g'\in U(T)(\mathbf A)$. Then $gB$ is orthonormal for $X'$,
  $(g'g)B=g'(gB)$, $gB=B$ if $g$ has archimedean component $1$, $s'_{gB,\rho}g_\rho=g_\rho s'_{B,\rho}$ for
  every $\rho$, and $x'_{gB}$ has the component $\overline{g_\rho}\,x'_{B,\bar\rho}$ at $\bar\rho$; if $g$ is an
  element $\gamma$ of $U(T)(F_h)$, then $x'_{gB}=\gamma x'_B$.
\item Let $B$ be orthonormal for $X'$ and $\varphi$ Schwartz-Bruhat. Then $\mathrm{Th}(x'_B,\varphi)$ is
  annihilated by the annihilation operators for the Cartan involutions $1$ of $S$ and $s'_{B,\rho}$ of $T$
  at every $\rho$, and for every polynomial $Q$ in the creation variables
  \[
   \big(\mathrm{Th}(x'_B,\varphi)\otimes_{s,B}Q\big)(0)=D'_{x'_B}(Q)\varphi,
  \]
  with $s_S=1$, $s_T=s'_B$, $B_S=1$ and $B_T=B$ (Lemma~\ref{is-eigenvalue-conjugate-theta-facts}). The
  substitution $(\rho,a,k)\mapsto(2/f_\rho)Y_{k,j}$ of $Q(\partial)$ is a bijection from the polynomials in the
  creation variables onto the polynomials in variables $Y_{k,j}$.
\end{enumerate}
\end{lemma}
\leanref{Development/AbelianVarietyTheory/Proofs/SiegelThetaTransformation/AdelicZakThetaExistsUnique.lean}{GlobalTheory.Ker.ZakEx.isZak_Th}
\leanref{Development/AbelianVarietyTheory/Proofs/SiegelThetaTransformation/ThetaKernelTestFunctionMapLinearInjectiveEquivar.lean}{GlobalTheory.Ker.Ctx.th_laws}
\leanref{Development/AbelianVarietyTheory/Proofs/SiegelThetaTransformation/FockPolynomialVectorTestThetaZeroEqThetaSeries.lean}{GlobalTheory.Ker.fvec_Th_zero}
\leanref{Development/AbelianVarietyTheory/Proofs/SiegelThetaTransformation/ThetaKernelOrthonormalFrameInDomainPeriodLinear.lean}{GlobalTheory.Ker.onBasis_frame_period}
\begin{proof}
(4) Each $g_\rho$ is unitary for $G_T(\rho)$, so $(g_\rho B_\rho)^*G_T(\rho)(g_\rho B_\rho)=B_\rho^*G_T(\rho)
B_\rho=J'$; $(g'g)_\rho=g'_\rho g_\rho$; $g_\rho=1$ for every $\rho$ if $g$ has archimedean component $1$;
$s'_{gB,\rho}g_\rho=g_\rho B_\rho J'B_\rho^{-1}g_\rho^{-1}g_\rho=g_\rho s'_{B,\rho}$; and the component of
$x'_{gB}$ at $\bar\rho$ is $\overline{g_\rho B_\rho}=\overline{g_\rho}\,x'_{B,\bar\rho}$. If $g=\gamma$ is rational,
then $g_\rho=\rho(\gamma)$ and $\overline{\rho(\gamma)}=\bar\rho(\gamma)$, so $x'_{gB}=\gamma x'_B$.

(3) By (2) of Definition~\ref{is-onto-share-flag-data}, for every $\rho$ some matrix $P_\rho$ has columns with
Gram matrix $J'$ for the form $h'_{\bar\rho}$ of the matrix $G_T^{\mathsf T}$; this Gram matrix is the complex
conjugate of the transpose of $\overline{P_\rho}^*G_T(\rho)\overline{P_\rho}$, so $B=(\overline{P_\rho})_\rho$ is
orthonormal for $X'$. Now let $B$ be orthonormal. In the same way the Gram matrix of the columns of
$\overline{B_\rho}$ for $h'_{\bar\rho}$ is $J'$. So $x'_B$ is orthonormal for $X'$, its leading minors are
$\pm1$ with the signs of $X'$, and $x'_B\in\mathcal D_{X'}$. The computation of
Lemma~\ref{is-eigenvalue-conjugate-frame-signature}(2), with $J'$ for $N$, shows that $s'_{B,\rho}$ is a
Cartan involution of $G_T(\rho)$, and $s'_{B,\rho}B_\rho=B_\rho J'$.

For the coordinates, let $u$ be a complex $r\times n$ matrix concentrated at $\rho$, $1\le k\le n$, and $j$
an index with $\tau'_j=\bar\rho$ and $a=a'_j$; the coordinates at the other indices are $0$. Put
$w=B_\rho^{-1}u$. If $a\le p'_0$, then, as $\overline{G_T(\rho)}^{\mathsf T}=G_T(\rho)$ and
$B_\rho^*G_T(\rho)=J'B_\rho^{-1}$,
\[
 z'_{x'_B}(\bar u^{\mathsf T})_{k,j}=h'_{\bar\rho}\big(\overline{B_\rho}e_a,\overline{ue_k}\big)
 =e_a^{\mathsf T}B_\rho^*G_T(\rho)ue_k=(J'w)_{ak}=w_{ak};
\]
if $a>p'_0$, then $z'_{x'_B}(\bar u^{\mathsf T})_{k,j}=(\overline{B_\rho}^{-1}\overline{ue_k})_a=\overline{w_{ak}}$.
For $is'_{B,\rho}u$ in place of $u$, $w$ becomes $iJ'w$, so the coordinate becomes $iw_{ak}$ if $a\le p'_0$,
and $\overline{-iw_{ak}}=i\,\overline{w_{ak}}$ if $a>p'_0$. For $u=B_\rho E_{ak'}$, $w=E_{ak'}$, so the
coordinates form the unit vector at $(k',j)$ for the index $j$ with $(\tau'_j,a'_j)=(\bar\rho,a)$.

(1) By Lemma~\ref{is-eigenvalue-conjugate-frame-zak-exists} there is a Zak theta function for $x'$ and $\varphi$.
Two of them agree at every $w$ with finite component $0$, where both are $\Theta'_\varphi(x';u'(w))$, so they
are equal by Lemma~\ref{is-eigenvalue-conjugate-frame-zak-space}(1).

(2) Write $\mathrm{Th}(\varphi)$ for $\mathrm{Th}(x',\varphi)$, and call $w$ real if its finite component is $0$. The
functions $\varphi+\varphi'$, $a\varphi$ and $\rho_{X'}(-y)\varphi$ are Schwartz-Bruhat
(Lemma~\ref{is-eigenvalue-conjugate-frame-vacua-mean}(1) and Lemma~\ref{is-eigenvalue-conjugate-frame-vacua-heisenberg}(2)).
By (1), a theta function whose value at every real
$w$ is $\Theta'_\psi(x';u'(w))$ is $\mathrm{Th}(\psi)$. Lemma~\ref{is-eigenvalue-conjugate-frame-zak-real} is
applied to $X'$.
\begin{itemize}
\item $\mathrm{Th}(\varphi)+\mathrm{Th}(\varphi')$ and $a\,\mathrm{Th}(\varphi)$ are theta functions
  (Lemma~\ref{is-eigenvalue-conjugate-frame-zak-space}(1)), with the values $\Theta'_{\varphi+\varphi'}(x';u'(w))$ and
  $\Theta'_{a\varphi}(x';u'(w))$ at real $w$ by Lemma~\ref{is-eigenvalue-conjugate-frame-zak-real}(4) and~(5); so they are
  $\mathrm{Th}(\varphi+\varphi')$ and $\mathrm{Th}(a\varphi)$.
\item If $\mathrm{Th}(\varphi)=\mathrm{Th}(\varphi')$, then $\Theta'_\varphi(x';u')=\Theta'_{\varphi'}(x';u')$ for every
  $u'\in U'$, as every $u'$ is $u'(w)$ for a real $w$ (Lemma~\ref{is-eigenvalue-conjugate-frame-zak-points}(2)); so
  $\varphi=\varphi'$ by Lemma~\ref{is-eigenvalue-conjugate-frame-zak-real}(4).
\item $\rho(y_f)\mathrm{Th}(\varphi)$ is a theta function (Lemma~\ref{is-eigenvalue-conjugate-theta-closed}(1)). For
  real $w$, the element $w-y^{\mathsf T}+y_f$ is real with $u'(w-y^{\mathsf T}+y_f)=u'(w)-\hat y$, by
  Lemma~\ref{is-eigenvalue-conjugate-frame-zak-forms}(8) and Lemma~\ref{is-eigenvalue-conjugate-frame-zak-points}(1)
  and~(3). With $\widehat{-y}=-\hat y$ and $E_{X'}$ linear in its second argument,
  \[
   (\rho(y_f)\mathrm{Th}(\varphi))(w)=e^{\pi iE_{X'}(u'(w),\hat y)}\,\Theta'_\varphi(x';u'(w)-\hat y)
   =\Theta'_{\rho_{X'}(-y)\varphi}(x';u'(w)),
  \]
  by the translation formula of Lemma~\ref{is-eigenvalue-conjugate-frame-zak-space}(3) and the Heisenberg law of
  Lemma~\ref{is-eigenvalue-conjugate-frame-zak-real}(2) for $-y$. So it is $\mathrm{Th}(\rho_{X'}(-y)\varphi)$.
\item Let $\gamma\in U(T)(F_h)$ have matrix $g$. As $\bar g^{\mathsf T}G_Tg=G_T$ and
  $h'(u,v)=\bar v^{\mathsf T}G_Tu$, $g$ preserves $h'$; so $w_\gamma\varphi$ is Schwartz-Bruhat
  (Lemma~\ref{is-eigenvalue-conjugate-frame-zak-weil}(1)) and $\gamma x'\in\mathcal D_{X'}$
  (Lemma~\ref{is-eigenvalue-conjugate-frame-hermitian}(5)). The function
  $F(w)=\mathrm{Th}(\varphi)(\gamma^{-1}w)$ is a theta function (Lemma~\ref{is-eigenvalue-conjugate-frame-zak-space}(4)).
  For real $w$, $\gamma^{-1}w$ is real with $\theta_\rho(\gamma^{-1}w)=\rho(g^{-1})\theta_\rho(w)$
  (Lemma~\ref{is-eigenvalue-conjugate-frame-zak-points}(5)), so $u'(\gamma^{-1}w)=g^{-1}u'(w)$, as
  $\overline{\rho(g^{-1})}=\bar\rho(g^{-1})$; so
  $F(w)=\Theta'_\varphi(x';g^{-1}u'(w))=\Theta'_{w_\gamma\varphi}(\gamma x';u'(w))$ by
  Lemma~\ref{is-eigenvalue-conjugate-frame-zak-real}(3). So $F=\mathrm{Th}(\gamma x',w_\gamma\varphi)$ by (1) for
  $\gamma x'$.
\end{itemize}

(5) The substitution sends the creation variables bijectively to nonzero multiples of the variables $Y_{k,j}$:
$f_\rho>0$, as $\xi_0\ne0$, and $(\rho,a)\mapsto j$ with $(\tau'_j,a'_j)=(\bar\rho,a)$ is a bijection. The
substitution $Y_{k,j}\mapsto(f_\rho/2)(\rho,a,k)$ is its inverse. Let $f=\mathrm{Th}(x'_B,\varphi)$; by (3),
$x'_B\in\mathcal D_{X'}$. At real $w$, $f(w)=\mathcal G'_{x'_B}(z)\Phi(z)$ with $z=z'_{x'_B}(u'(w))$ and
$\Phi=\theta'_\varphi(x'_B;\cdot)$, which is holomorphic by Lemma~\ref{is-eigenvalue-conjugate-frame-zak-series}(1).
Let $\rho\in\mathcal T$ and $z_0$ be given, and let $Af$ be the image of $f$ under the annihilation operator at
$\rho$ and $z_0$. It is a theta function, by Lemma~\ref{is-eigenvalue-conjugate-theta-closed}(2) and
Lemma~\ref{is-eigenvalue-conjugate-frame-zak-space}(1). By Lemma~\ref{is-eigenvalue-conjugate-frame-zak-fock}(1),
$(Af)(w)=\mathcal G'_{x'_B}(z)\big(d\Phi_z(u_0)+i\,d\Phi_z(iu_0)\big)=0$ at every real $w$, as $d\Phi_z$ is complex
linear. So $Af=0$ by Lemma~\ref{is-eigenvalue-conjugate-frame-zak-space}(1). Finally, by
Lemma~\ref{is-eigenvalue-conjugate-frame-zak-fock}(2), $(f\otimes_{s,B}Q)(0)$ is the value at $0$ of $Q(\partial)$
applied to $\theta'_\varphi(x'_B;\cdot)$, which is $D'_{x'_B}(Q)\varphi$.
% (1) is Ker.ZakEx.isZak_Th (279 lines in its file, T.adelicZakTheta_exists_unique); (2) is Ctx.th_laws (126 lines in
% its file: additivity with injectivity, homogeneity, the Heisenberg clause and the clause for IsRatOf T ga gamma
% with wG ga phi, from T.realPoints_classicalModel_realVectorSpace, T.thetaFunction_space_heisenberg_laws, thR_laws,
% heisG_laws, symplecticCharacter_real_finite_points_match, rational_unitary_group_structure and wG_laws); Th 0 = 0
% and Th_ne_zero (ExistsRealCharacter.lean) follow from (2). (3) and (4) are onBasis_frame_period: exists_isB
% (Data.signature and gram_conj), isONFr_xB, mem_dom, isCartanInv, isONBasis, zR_single, zR_cplxO and zR_unitMap
% (the computation above), isB_mvB, mvB_mul, mvB_fin, sB_mvB, xB_mvB and xB_rat. (5) is fvec_Th_zero (83 lines in
% its file): toH_bijective, the IsFock clause, and fvec B (Th (xB B) phi) Q 0 = thP X hr bas idx phi (toH Q) (xB B).
% The proofs follow the Lean: (1) is exists_isZak (Lemma -zak-exists) and clause 2.1 of theta_space_laws (Lemma
% -zak-space(1)); (2) is th_laws with hth, hzak and hlaws (thR_laws, Lemma -zak-real), hsub and realPt_actO
% (Lemma -zak-points), symplecticCharacter_real_finite_points_match (Lemma -zak-forms(8)) and the formula rho_finW
% (Lemma -zak-space(3)), ZakLaws.embR_neg and
% EFR_neg_right; in the rational clause the book checks IsUn C.X ga from the unitarity of gamma, where the Lean
% takes it, with IsRatOf T ga^-1 gamma^-1 and gamma^-1 in GQ T, from T.rational_unitary_group_structure. (5) is
% toH_bijective (bijective_aeval_scale_rename, fockScale_pos), the IsFock clause from
% T.heisenbergDerivative_thetaFunction_laws (Lemma theta-closed(2)), T.annihilation_creation_periodCoordinates and
% the uniqueness (Lemma -zak-fock(1), -zak-space(1)), and T.fockPolynomialVector_origin_eq_polyDeriv (Lemma
% -zak-fock(2)); the holomorphy is the last clause of SW.theta_holomorphic_weight (zero weights, P = 0), which the
% book takes from Lemma -zak-series(1) (ThetaHol.differentiable_thA), the same assertion.
\end{proof}

\subsection{Archimedean characters and the overlap formula}\label{subsection-is-archimedean-characters}

\begin{lemma}\label{is-eigenvalue-conjugate-frame-root-values}
Let $D$ be a definite-side datum with host field $E_h/F_h$ and $V$ a Hermitian space of $D$.
\begin{enumerate}
\item Every $g\in U(V)(\mathbf A)$ has a determinant root
  (Lemma~\ref{is-eigenvalue-conjugate-frame-det-roots}).
\item If $y$ and $y'$ are determinant roots of $g$, then $y'=ya$ with $a\in\mathbf A_{F_h}^\times$.
\item Let $\chi$ be a character of $\mathbf A_{E_h}^\times$ whose restriction to $\mathbf A_{F_h}^\times$ is
  $\eta^n$ for an even integer $n$, $\eta$ the quadratic character of $E_h/F_h$. Then $\chi(y)=\chi(y')$
  for any two determinant roots $y,y'$ of an element $g$.
\item Every $\gamma\in U(V)(F_h)$ has a determinant root in $E_h^\times\subset\mathbf A_{E_h}^\times$.
\end{enumerate}
\end{lemma}
\leanref{Development/AbelianVarietyTheory/Proofs/UnitaryExistence/ExistsAlgebraMapOfSigRFixed.lean}{GlobalTheory.Member.isDetRoot_laws}
\leanref{Development/AbelianVarietyTheory/Proofs/UnitaryExistence/ExistsIsDetRoot.lean}{GlobalTheory.RecDict.exists_isDetRoot}
\leanref{Development/AbelianVarietyTheory/Proofs/IwasawaDecomposition/RealPointMovesAdelicThetaByDetCharacter.lean}{GlobalTheory.OneMapRef.RealMove.root_eq}
\begin{proof}
(1) As $g$ is unitary, $\det(g)\,\overline{\det(g)}=1$. By Hilbert's Theorem~90 for the quadratic extension
$\mathbf A_{E_h}$ of $\mathbf A_{F_h}$ \cite{Ser79}, there is $y\in\mathbf A_{E_h}^\times$ with
$y=\det(g)\,\bar y$.

(2) The idele $a=y^{-1}y'$ satisfies $\bar a=\overline{y}^{-1}\overline{y'}=(\det(g)^{-1}y)^{-1}\det(g)^{-1}y'=a$,
so it lies in $\mathbf A_{F_h}^\times$.

(3) By (2), $\chi(y')=\chi(y)\chi(a)=\chi(y)\eta(a)^n=\chi(y)$, as $\eta^2=1$ and $n$ is even.

(4) The determinant $u=\det\gamma$ lies in $E_h^\times$ and $u\bar u=1$. By Hilbert's Theorem~90 for $E_h/F_h$
\cite{Ser79} there is $x\in E_h^\times$ with $x=u\bar x$, and $x$ is a determinant root of $\gamma$.
\end{proof}
% exists_isDetRoot is AdLat.hilbert90_adeleRing_tensor_quadratic for det g with det_conj_mul; the second clause of
% isDetRoot_laws uses exists_unit_of_sigR_fixed; the third, (4), is exists_root_of_norm_one for the conjugation of
% E_h with det_conj_mul; the fourth is isDetRoot_mul (Lemma is-eigenvalue-conjugate-frame-det-roots(1)); root_eq is
% (3), with ResEta chi n and IsEta_mul_self.

\begin{lemma}\label{is-eigenvalue-conjugate-frame-compact-character}
Let $D$ be a definite-side datum with set $\mathcal T$ of embeddings, $\theta$ an archimedean family whose
member $\theta_\rho$ extends $\rho$ for every $\rho\in\mathcal T$, $V$ a Hermitian space of $D$ with Gram
matrix $G_V$, and $s=(s_\rho)_{\rho\in\mathcal T}$ a family of Cartan involutions of the $G_V(\rho)$:
$s_\rho^2=1$, $s_\rho^*G_V(\rho)s_\rho=G_V(\rho)$ and $G_V(\rho)s_\rho$ is positive definite. Let $c$ be a
function on $U(V)(\mathbf A)$ with $c(uu')=c(u)c(u')$ for all $u,u'$ with finite component $1$, and
$c(k)=1$ for every $k$ with finite component $1$ and $s_\rho k_\rho=k_\rho s_\rho$ for every $\rho$. Then
$c(u)=1$ for every $u$ with finite component $1$.
\end{lemma}
\leanref{Development/AbelianVarietyTheory/Proofs/DivisorProofs/MulFunEqOneOfEqOneOnMaxCompact.lean}{GlobalTheory.FrameEquiv.mulFun_eq_one_of_eq_one_on_maxCompact}
\leanref{Development/AbelianVarietyTheory/Proofs/DivisorProofs/CompAtMtMul.lean}{GlobalTheory.UnitaryChar.char_eq_one}
\leanref{Development/AbelianVarietyTheory/Proofs/DivisorProofs/CompAtMtMul.lean}{GlobalTheory.UnitaryChar.polar}
\begin{proof}
Let $u$ have finite component $1$, and fix $\rho$. Write $G=G_V(\rho)$ and $s=s_\rho$. By the polar
decomposition \cite{HJ13}, applied after conjugation by the positive square root $Q$ of $Gs$ (for which
$Q^{-1}GQ^{-1}$ is a Hermitian involution and $Q u_\rho Q^{-1}$ is unitary for it), $u_\rho=k_\rho q_\rho^2$
with $k_\rho$ and $q_\rho$ unitary for $G$, $sk_\rho=k_\rho s$ and $sq_\rho sq_\rho=1$. By
Lemma~\ref{is-eigenvalue-conjugate-frame-glue}(2) there are $k$, $q$ and $\sigma$ in $U(V)(\mathbf A)$ with
finite component $1$ and archimedean images $k_\rho$, $q_\rho$ and $s_\rho$ at every $\rho$; $s_\rho$ is
unitary for $G_V(\rho)$. By Lemma~\ref{is-eigenvalue-conjugate-frame-glue}(1), $u=kq^2$ and $\sigma q\sigma
q=1$, as both sides have the same archimedean images and finite component $1$. As $k$, $\sigma$ and $1$
commute with $s$ at every $\rho$, $c(k)=c(\sigma)=c(1)=1$. So
$1=c(\sigma q\sigma q)=c(q)^2$ and $c(u)=c(k)c(q)^2=1$.
\end{proof}
% char_eq_one: polar at each tau (B = CFC.sqrt (G s), polar_std for the involution S = C G C), glue for k', q' and
% sigma, ext_real (glue (1)) for e1 and e2, and the multiplicativity hm.

\begin{lemma}\label{is-eigenvalue-conjugate-frame-fock-character-lattice}
Let $D$ be a definite-side datum with host field $E_h/F_h$ and ideal $\mathfrak A$, $\psi$ a standard character of
$\mathbf A_{F_h}$, $S$ and $T$ Hermitian spaces of $D$ with self-dual lattices $\Lambda_S$ and $\Lambda_T$, and
$\chi_S,\chi_T$ continuous unitary characters of $\mathbf A_{E_h}^\times$, trivial on $E_h^\times$, whose restrictions
to $\mathbf A_{F_h}^\times$ are $\eta^{\operatorname{rk}S}$ and $\eta^{\operatorname{rk}T}$, $\eta$ the quadratic
character of $E_h/F_h$. Let $\hat{\mathcal O}$ be the closure of the integers of $F_h$ in $\mathbf A_{F_h,f}$, and for a
finite place $v$ of $F_h$ let $\Lambda_v\subset\mathbb W$, the \emph{part at $v$} of the lattice of the pair given by
$\Lambda_S$ and $\Lambda_T$, be the subgroup generated by the elements with archimedean component $0$ and finite
component $\alpha\,a\,t\,\bar s^{\mathsf T}G_S$, with $\alpha\in\hat{\mathcal O}$ whose components away from $v$ are $0$,
$a\in\mathfrak A$, $s\in\Lambda_S$ and $t\in\Lambda_T$. Then there is a homomorphism $\omega_f$ from
$U(S)(\mathbf A)\times U(T)(\mathbf A)$ to the linear endomorphisms of the space of functions on $\mathbb W$ such that
for every theta function $f$ (Definition~\ref{th-weil-datum}):
\begin{enumerate}
\item $\omega_f(p)f$ is a theta function, $p\mapsto(\omega_f(p)f)(w)$ is continuous for every $w\in\mathbb W$, and
  $\omega_f(p)$ preserves the $L^2$-norm of $f$ on every fundamental domain for
  $M_{\operatorname{rk}T\times\operatorname{rk}S}(E_h)$ in $\mathbb W$, for every Haar measure;
\item $\omega_f(g,h)f=f$ if $g$ and $h$ have finite components $1$;
\item $\omega_f(g,h)\rho(w_0)f=\rho(hw_0g^{-1})\,\omega_f(g,h)f$ for every $w_0\in\mathbb W$, if $g$ and $h$ have
  archimedean components $1$;
\item $\omega_f$ satisfies the condition of Definition~\ref{th-weil-datum}(3) at the places of $F_h$ split in $E_h$,
  for the characters $(\chi_S,\chi_T)$;
\item for every finite place $v$ of $F_h$ not split in $E_h$, every $g\in U(S)(\mathbf A)$ and $h\in
  U(T)(\mathbf A)$ with components $1$ away from $v$ such that $u\in\Lambda_v$ exactly when $hug^{-1}\in\Lambda_v$, and
  every theta function $f$ with $\rho(u)f=f$ for every $u\in\Lambda_v$, $\omega_f(g,h)f=f$.
\end{enumerate}
\end{lemma}
\leanref{Development/AbelianVarietyTheory/Proofs/DivisorProofs/ExistsIsFinWeil.lean}{GlobalTheory.FinWeilO.exists_isFinWeil}
\leanref{Development/AbelianVarietyTheory/Definitions/Divisors/IsUnitaryOnTheta.lean}{GlobalTheory.WeilParts.IsFinWeil}
\leanref{Development/AbelianVarietyTheory/Definitions/Divisors/IsUnitaryOnTheta.lean}{GlobalTheory.WeilParts.FixesVacAt}
\leanref{Development/AbelianVarietyTheory/Definitions/Divisors/IsUnitaryOnTheta.lean}{GlobalTheory.WeilParts.NonSplitAt}
\begin{proof}
\openstep{The construction of the lattice model, the action of the pairs on the theta functions through the
finite places twisted by characters of the unitary groups built from $\chi_S$ and $\chi_T$, and the proofs of its
five properties, are not written out.}
% FinWeilO.exists_isFinWeil: exists_goodNorm, the lattice representation exists_isLatRep (pres, cont, unit,
% real_triv, Heisenberg2, split and inert for the action Om on thetaSub), T.exists_splitPlaces_selecting_adelicIdempotent,
% the characters c and d of T.exists_unitaryGroup_character_of_heckeChar for (S, chi_T) and (T, chi_S) (c2, d2: trivial
% at archimedean elements; c3, d3: trivial at elements supported at a nonsplit place; c4, d4: the split values),
% z = c^-1 d and exists_extend for om = z . Om. Size, following imports and stopping at statements the book cites:
% about 3,940 lines, of which exists_isLatRep about 3,160 and the characters about 440.
\end{proof}

\begin{lemma}\label{is-eigenvalue-conjugate-frame-fock-character-oscillator}
Let $\iota$ be a finite set and $H$ an invertible Hermitian matrix indexed by $\iota$, and put
$\omega_H(x,y)=\operatorname{Im}(y^*Hx)$ for $x,y\in\mathbf C^\iota$. Let $B$ be a complex Banach space and $\rho$ a map
from $\mathbf C^\iota$ to the bounded operators on $B$ with $\rho(0)=1$,
$\rho(u)\rho(v)=e^{\pi i\omega_H(u,v)}\rho(u+v)$, $\|\rho(u)f\|=\|f\|$, and $u\mapsto\rho(u)f$ continuous for every
$f\in B$; call $f$ \emph{smooth} if $u\mapsto\rho(u)f$ is a smooth map of real vector spaces, and write $D_f$ for its
real derivative at $0$. Let $G$ be a group and $m\colon G\to GL_\iota(\mathbf C)$ a homomorphism with $m(g)^*Hm(g)=H$
for every $g$. Then there is a map $\Omega\colon G\times B\to B$, $(g,f)\mapsto\Omega(g)f$, such that:
\begin{enumerate}
\item on smooth vectors, $\Omega(g)$ is additive and complex homogeneous and preserves smoothness, $\Omega(1)f=f$,
  $\Omega(gg')f=\Omega(g)\Omega(g')f$, and $\Omega(g)\rho(v)f=\rho(m(g)v)\Omega(g)f$;
\item for a smooth $f$, $\Omega(g)Af=A\Omega(g)f$ for every bounded operator $A$ on $B$ that commutes with every
  $\rho(v)$, and $\|\Phi(\Omega(g)f)\|=\|\Phi(f)\|$ for every bounded linear map $\Phi$ from $B$ to a complex Hilbert
  space with $\|\Phi(\rho(v)f')\|=\|\Phi(f')\|$ for all $v$ and $f'$;
\item $\Omega(g)f=f$ for every smooth $f$ if $m(g)=1$;
\item if $G$ is a topological group and $m$ is continuous, $g\mapsto\Omega(g)f$ is continuous for every smooth $f$;
\item let $s$ be a matrix with $s^2=1$ and $Hs$ positive definite, $g\in G$ with $sm(g)=m(g)s$, and $f$ a smooth
  vector with $D_f(z)+i\,D_f(isz)=0$ for every $z\in\mathbf C^\iota$. Then
  $\Omega(g)f=\det\big(\tfrac12(1+s)m(g)+\tfrac12(1-s)\big)f$.
\end{enumerate}
\end{lemma}
\leanref{Development/AbelianVarietyTheory/Proofs/DivisorProofs/ExistsOscillatorActionUnitaryGroupHermitianStatement.lean}{GlobalTheory.T.exists_oscillatorAction_unitaryGroup_hermitian}
\leanref{Development/AbelianVarietyTheory/Proofs/DivisorProofs/ContinuousFourierGau.lean}{GlobalTheory.RealWeil.IsOsc}
\leanref{Development/AbelianVarietyTheory/Proofs/DivisorProofs/ContinuousFourierGau.lean}{GlobalTheory.RealWeil.IsHeis}
\leanref{Development/AbelianVarietyTheory/Proofs/DivisorProofs/ContinuousFourierGau.lean}{GlobalTheory.RealWeil.dP}
\begin{proof}
\openstep{The construction of the oscillator action of the unitary group of $H$ on the smooth vectors of a
Heisenberg representation, by Gaussian pairings, and the computation of its action on the vectors annihilated by the
annihilation operators of a Cartan involution $s$, are not written out.}
% T.exists_oscillatorAction_unitaryGroup_hermitian = RealWeil.UnitaryOsc.exists_osc (ExistsOscillatorActionUnitaryGroupHermitian.lean),
% through T.exists_oscillatorAction_of_gaussianPairings, GaussianPairingUnitaryGroupTwoThree,
% T.integral_complexGaussian_posDefHermitianPart, ProjEqSelfOfAnnihilated (the vacuum clause), CartanInvolutionConjOfPosDef
% and HermitianCongrInvolution. Size, following imports and stopping at statements the book cites: about 1,790 lines.
\end{proof}

\begin{lemma}\label{is-eigenvalue-conjugate-frame-fock-character-fock}
Let $D$ be a definite-side datum with host field $E_h/F_h$ and set $\mathcal T$ of embeddings, $\psi$ a standard
character of $\mathbf A_{F_h}$, $\theta$ an archimedean family whose member $\theta_\rho$ extends $\rho$ for every
$\rho\in\mathcal T$, and $S$ and $T$ Hermitian spaces of $D$. Then there is a homomorphism $\omega_\infty$ from
$U(S)(\mathbf A)\times U(T)(\mathbf A)$ to the linear endomorphisms of the space of functions on $\mathbb W$ such that
for every theta function $f$:
\begin{enumerate}
\item $\omega_\infty$ has property (1) of Lemma~\ref{is-eigenvalue-conjugate-frame-fock-character-lattice};
\item $\omega_\infty(g,h)f=f$ if $g$ and $h$ have archimedean components $1$;
\item $\omega_\infty(g,h)\rho(w_0)f=\rho(hw_0g^{-1})\,\omega_\infty(g,h)f$ for every $w_0\in\mathbb W$, if $g$ and $h$
  have finite components $1$;
\item let $s_S$ and $s_T$ be families of Cartan involutions of $S$ and $T$ at the $\rho\in\mathcal T$, let
  $g\in U(S)(\mathbf A)$ and $h\in U(T)(\mathbf A)$ have finite components $1$, with $s_{S,\rho}g_\rho=g_\rho
  s_{S,\rho}$ and $s_{T,\rho}h_\rho=h_\rho s_{T,\rho}$ for every $\rho$, and let $f$ be annihilated by the annihilation
  operators for $s_S$ and $s_T$ at every $\rho$ (Lemma~\ref{is-eigenvalue-conjugate-theta-facts}). Then
  $\omega_\infty(g,h)f=\prod_{\rho\in\mathcal T}\Delta_\rho\,f$, where $\Delta_\rho$ is the determinant of
  $\frac12(1+s_{T,\rho}\otimes s_{S,\rho}^{\mathsf T})(h_\rho\otimes(g_\rho^{-1})^{\mathsf T})+\frac12(1-s_{T,\rho}\otimes
  s_{S,\rho}^{\mathsf T})$.
\end{enumerate}
\end{lemma}
\leanref{Development/AbelianVarietyTheory/Proofs/DivisorProofs/LintegralEqOfNorms.lean}{GlobalTheory.RealWeil.exists_isInfWeil}
\leanref{Development/AbelianVarietyTheory/Definitions/Divisors/IsUnitaryOnTheta.lean}{GlobalTheory.WeilParts.IsInfWeil}
\leanref{Development/AbelianVarietyTheory/Definitions/Divisors/IsUnitaryOnTheta.lean}{GlobalTheory.WeilParts.VacCharPlus}
\begin{proof}
\emph{The model.} Let $n=\operatorname{rk}S$, $r=\operatorname{rk}T$ and $\iota=(\{1,\dots,r\}\times\{1,\dots,n\})
\times\mathcal T$, so that $\mathbf C^\iota$ is the space of families $v=(v_\rho)$ of complex $r\times n$ matrices. Let
$\tilde v\in\mathbb W$ be the element with finite component $0$ and $\theta_\rho(\tilde v)=v_\rho$ for every $\rho$. Let
$B$ be the Banach space, with the supremum norm, of the bounded continuous functions $f$ on $\mathbb W$ with
$f(\gamma+w)=\psi(-\frac12\langle\gamma,w\rangle)f(w)$ for every $\gamma\in M_{r\times n}(E_h)$, on which the operators
$\rho(w_0)$ of Definition~\ref{th-weil-datum} act, and put $\rho_B(v)=\rho(\tilde v)$. Let $H$ be the block diagonal
matrix with the blocks $2\operatorname{Im}\rho(\xi_0)\,G_T(\rho)\otimes(G_S(\rho)^{-1})^{\mathsf T}$, and $m$ the
homomorphism $(g,h)\mapsto$ the block diagonal matrix with the blocks $h_\rho\otimes(g_\rho^{-1})^{\mathsf T}$, so that
$m(g,h)$ acts by $v_\rho\mapsto h_\rho v_\rho g_\rho^{-1}$. Then $\rho_B$ satisfies the hypotheses of
Lemma~\ref{is-eigenvalue-conjugate-frame-fock-character-oscillator} for $H$, $m(g,h)^*Hm(g,h)=H$, and the theta functions
lie in $B$ and are smooth for $\rho_B$. Let $\Omega$ be the action that lemma gives for the group
$U(S)(\mathbf A)\times U(T)(\mathbf A)$ and $m$. It carries theta functions to theta functions, and $\omega_\infty$ is
$\Omega$ on the theta functions, extended by the identity on a complement of the space of theta functions.
\openstep{The realization of the theta functions as smooth vectors of $\rho_B$, the hypotheses of the oscillator
lemma for $\rho_B$, $H$ and $m$, the stability of the theta functions under $\Omega$, properties (1) to (3), and the
reading of the annihilation operators as the condition $D_f(z)+i\,D_f(isz)=0$ for the block matrix $s$ below, are not
written out.}
% RealWeil.exists_isInfWeil: RealLineAt.exists_inverse_chart (sigma), BB psi S T = QB (ratW) (the quasi-periodic bounded
% continuous functions), rhoV = rhoB o iotaV, HM S T and mM theta S T (blk b a = blockDiagonal (b tau (x) (a tau)^T)), HM_herm,
% HM_det, isHeis_rhoB, mM_unitary; sm_of_theta, theta_of_sm and the level argument hth; exists_hom; lintegral_eq_of_norms for
% the norms; mM_fin for (2); exists_parts, rhoB_comm, actO_fin, actO_iotaV and rhoB_add for (3); fderiv_vac (with
% VacUnique.Ann.lineAt_eq) for the annihilation condition. Size, following imports and stopping at the oscillator lemma
% and the three computations below: about 2,260 lines, of which fderiv_vac about 360.

\emph{Property (4).} Let $s_S$, $s_T$, $g$, $h$ and $f$ be as in (4), and let $s$ be the block diagonal matrix with the
blocks $s_{T,\rho}\otimes s_{S,\rho}^{\mathsf T}$. Then $s^2=1$, as $s_{S,\rho}^2=s_{T,\rho}^2=1$. The block of $Hs$ at
$\rho$ is $2\operatorname{Im}\rho(\xi_0)\,G_T(\rho)s_{T,\rho}\otimes\big(s_{S,\rho}G_S(\rho)^{-1}\big)^{\mathsf T}$. Here
$\operatorname{Im}\rho(\xi_0)>0$, $G_T(\rho)s_{T,\rho}$ is positive definite, and $s_{S,\rho}G_S(\rho)^{-1}=(G_S(\rho)
s_{S,\rho})^{-1}$ is positive definite, as $G_S(\rho)s_{S,\rho}=s_{S,\rho}^*G_S(\rho)$ by the Cartan conditions; a
Kronecker product of positive definite matrices is positive definite, so $Hs$ is positive definite. The blocks of
$s$ commute with those of $m(g,h)$: $s_{T,\rho}h_\rho=h_\rho s_{T,\rho}$, and $s_{S,\rho}g_\rho=g_\rho s_{S,\rho}$ gives
$s_{S,\rho}g_\rho^{-1}=g_\rho^{-1}s_{S,\rho}$, hence its transpose. By the model, $f$ is a smooth vector with
$D_f(z)+i\,D_f(isz)=0$ for every $z$, and $\omega_\infty(g,h)f=\Omega(g,h)f$. By
Lemma~\ref{is-eigenvalue-conjugate-frame-fock-character-oscillator}(5), $\Omega(g,h)f=\det\big(\frac12(1+s)m(g,h)+
\frac12(1-s)\big)f$, and the determinant of this block diagonal matrix is the product of the determinants of its
blocks, which are the $\Delta_\rho$.
\end{proof}
% Property (4): hvac of the oscillator lemma with HM_cartan (the two Cartan conditions; cartan_inv_mul, two_im_pos and
% blk_posDef), blk_comm_mM, fderiv_vac and detPlus_pairMatrix (detPlus_blockDiagonal).

\begin{lemma}\label{is-eigenvalue-conjugate-frame-fock-character-product}
Let $D$ be a definite-side datum with host field $E_h/F_h$ and set $\mathcal T$ of embeddings, $\psi$ a standard
character of $\mathbf A_{F_h}$, $\theta$ an archimedean family whose member $\theta_\rho$ extends $\rho$ for every
$\rho\in\mathcal T$, $S$ and $T$ Hermitian spaces of $D$ with self-dual lattices $\Lambda_S$ and $\Lambda_T$, and
$\chi_S,\chi_T$ characters as in Lemma~\ref{is-eigenvalue-conjugate-frame-fock-character-lattice}. Let $\omega_f$ be
a homomorphism with properties (1) to (5) of Lemma~\ref{is-eigenvalue-conjugate-frame-fock-character-lattice} for
these data, and $\omega_\infty$ one with properties (1) to (4) of
Lemma~\ref{is-eigenvalue-conjugate-frame-fock-character-fock}. Let $f$ be a theta function.
\begin{enumerate}
\item $\omega_f(p)\omega_\infty(q)f=\omega_\infty(q)\omega_f(p)f$ for all $p,q\in U(S)(\mathbf A)\times
  U(T)(\mathbf A)$.
\item $p\mapsto(\omega_f(p)\omega_\infty(p)f)(w)$ is continuous for every $w\in\mathbb W$.
\item Suppose that $(\chi_S,\chi_T)=(\tilde\eta^{\operatorname{rk}S},\tilde\eta^{\operatorname{rk}T})$ for a Hecke
  character $\tilde\eta$ of $E_h$ that is unramified at every finite place, restricts to $\eta$ on
  $\mathbf A_{F_h}^\times$ and has archimedean exponent $\frac12$ at every $\rho$. Then
  $(\omega_f(\gamma,1)\omega_\infty(\gamma,1)f)(0)=f(0)$ for $\gamma\in U(S)(F_h)$, and
  $(\omega_f(1,\gamma')\omega_\infty(1,\gamma')f)(0)=f(0)$ for $\gamma'\in U(T)(F_h)$.
\end{enumerate}
\end{lemma}
\leanref{Development/AbelianVarietyTheory/Proofs/DivisorProofs/ContinuousProduct.lean}{GlobalTheory.WeilProduct.fin_inf_comm}
\leanref{Development/AbelianVarietyTheory/Proofs/DivisorProofs/ContinuousProduct.lean}{GlobalTheory.WeilProduct.continuous_product}
\leanref{Development/AbelianVarietyTheory/Proofs/DivisorProofs/ThetaDistributionAutomorphySourceGroupStatement.lean}{GlobalTheory.T.thetaDistribution_automorphy_source_group}
\leanref{Development/AbelianVarietyTheory/Proofs/DivisorProofs/ThetaDistributionInvariantRationalTargetStatement.lean}{GlobalTheory.T.thetaDistribution_invariant_rationalTarget}
\begin{proof}
\openstep{The commutation of the lattice model with the Fock model, through the Schur lemma for theta functions and
the lattice representation, the continuity of their product, and the invariance of the theta distribution under the
rational points of the two unitary groups are not written out.}
% fin_inf_comm (exists_isLatRep, fin_eq_smul, exists_parts, the commutation h.comm of the lattice representation with
% operators commuting with the real Heisenberg operators, rhoF_real); continuous_product (eventually_smul);
% T.thetaDistribution_automorphy_source_group = DualSwap.automorphy_source and
% T.thetaDistribution_invariant_rationalTarget = RatAut.automorphy_target. Size, following imports and stopping at
% statements the book cites: about 19,090 lines for the four, 14,130 of them not shared with Lemmas
% -fock-character-lattice and -fock-character-fock.
\end{proof}

\begin{lemma}\label{is-eigenvalue-conjugate-frame-fock-character-model}
Let $D$ be a definite-side datum with host field $E_h/F_h$, $\psi$ a standard character of
$\mathbf A_{F_h}$ and $\theta$ an archimedean family whose member $\theta_\rho$ extends $\rho$ for every
$\rho\in\mathcal T$. Let $S$ and $T$ be Hermitian spaces of $D$ of ranks $n$ and $r$ with self-dual lattices
$\Lambda_S$ and $\Lambda_T$, and for a finite place $v$ of $F_h$ let $\Lambda_v$ be the part at $v$ of the lattice of
the pair given by them (Lemma~\ref{is-eigenvalue-conjugate-frame-fock-character-lattice}). Let $\tilde\eta$ be a Hecke
character of $E_h$ that is unramified at every finite place, restricts to the quadratic character $\eta$ of $E_h/F_h$
on $\mathbf A_{F_h}^\times$ and has archimedean exponent $\frac12$ at every $\rho$. Then there is a Weil datum
$\omega_0$ for $(S,T)$ with characters $(\tilde\eta^n,\tilde\eta^r)$ with the following two properties.
\begin{enumerate}
\item Let $v$ be a finite place of $F_h$ not split in $E_h$, let $g\in U(S)(\mathbf A)$ and $h\in U(T)(\mathbf A)$
  have components $1$ away from $v$, with $u\in\Lambda_v$ exactly when $hug^{-1}\in\Lambda_v$, and let $f$ be a theta
  function with $\rho(u)f=f$ for every $u\in\Lambda_v$. Then $\omega_0(g,h)f=f$.
\item Let $s_S$ and $s_T$ be families of Cartan involutions of $S$ and $T$ at the $\rho$, let
  $g\in U(S)(\mathbf A)$ and $h\in U(T)(\mathbf A)$ have finite components $1$, with
  $s_{S,\rho}g_\rho=g_\rho s_{S,\rho}$ and $s_{T,\rho}h_\rho=h_\rho s_{T,\rho}$ for every $\rho$, and let $f$ be a theta
  function annihilated by the annihilation operators for $s_S$ and $s_T$ at every $\rho$. Then
  \[
   \omega_0(g,h)f=\prod_{\rho\in\mathcal T}\Delta_\rho\,f,
  \]
  where $\Delta_\rho$ is the determinant of
  $\frac12(1+s_{T,\rho}\otimes s_{S,\rho}^{\mathsf T})(h_\rho\otimes(g_\rho^{-1})^{\mathsf T})+\frac12(1-s_{T,\rho}\otimes
  s_{S,\rho}^{\mathsf T})$.
\end{enumerate}
\end{lemma}
\leanref{Development/AbelianVarietyTheory/Proofs/DivisorProofs/ExistsWeilDatumPowersUnramifiedEtaTildeStatement.lean}{GlobalTheory.T.exists_weilDatum_powers_unramifiedEtaTilde}
\leanref{Development/AbelianVarietyTheory/Proofs/DivisorProofs/ExistsWeilDatumPowersUnramifiedEtaTilde.lean}{GlobalTheory.WeilProduct.exists_weilDatum_etaTilde}
\leanref{Development/AbelianVarietyTheory/Definitions/Divisors/IsUnitaryOnTheta.lean}{GlobalTheory.WeilParts.VacCharPlus}
\leanref{Development/AbelianVarietyTheory/Definitions/Divisors/IsUnitaryOnTheta.lean}{GlobalTheory.WeilParts.FixesVacAt}
\begin{proof}
The characters $\chi_S=\tilde\eta^n$ and $\chi_T=\tilde\eta^r$ are continuous, unitary and trivial on $E_h^\times$, and
restrict to $\eta^n$ and $\eta^r$ on $\mathbf A_{F_h}^\times$. Let $\omega_f$ be a homomorphism as in
Lemma~\ref{is-eigenvalue-conjugate-frame-fock-character-lattice} for $\Lambda_S$, $\Lambda_T$ and these characters, and
$\omega_\infty$ one as in Lemma~\ref{is-eigenvalue-conjugate-frame-fock-character-fock}. By
Lemma~\ref{is-eigenvalue-conjugate-frame-fock-character-product}(1), $p\mapsto\omega_f(p)\omega_\infty(p)$, restricted
to the space of theta functions, which both preserve, is a homomorphism. Let $\omega_0(p)$ be
$\omega_f(p)\omega_\infty(p)$ on the theta functions, extended by the identity on a complement of the space of theta
functions; it is a homomorphism. Write a pair $p=(g,h)$ as $(a,b)(k,l)$ with $a,b$ of finite components $1$ and $k,l$ of
archimedean components $1$, the two factors commuting. For a theta function $f$, if $g$ and $h$ have archimedean
components $1$, then $\omega_0(g,h)f=\omega_f(g,h)f$, as $\omega_\infty(g,h)f=f$
(Lemma~\ref{is-eigenvalue-conjugate-frame-fock-character-fock}(2)); if they have finite components $1$, then
$\omega_0(g,h)f=\omega_\infty(g,h)f$, by property (2) of Lemma~\ref{is-eigenvalue-conjugate-frame-fock-character-lattice}
applied to the theta function $\omega_\infty(g,h)f$.

The conditions of Definition~\ref{th-weil-datum}(3) hold: $\psi$ is standard and the characters are as required;
$\omega_0(p)$ preserves theta functions and their norms, by property (1) of the two lemmas; $\omega_0(p)=\omega_\infty(a,b)
\omega_f(k,l)$ on theta functions, and the Heisenberg condition for $p$ follows from property (3) of the two lemmas at
$(k,l)$ and $(a,b)$, with $hw_0g^{-1}=b(lw_0k^{-1})a^{-1}$; the continuity is
Lemma~\ref{is-eigenvalue-conjugate-frame-fock-character-product}(2); the invariance at
$(\gamma,\gamma')=(\gamma,1)(1,\gamma')$ is Lemma~\ref{is-eigenvalue-conjugate-frame-fock-character-product}(3), applied
to the theta functions $\omega_0(1,\gamma')f$ and $f$; and a pair supported at a split place has archimedean components
$1$, so the condition at the split places is property (4) of
Lemma~\ref{is-eigenvalue-conjugate-frame-fock-character-lattice}. So $\omega_0$ is a Weil datum with characters
$(\tilde\eta^n,\tilde\eta^r)$.

(1) The elements $g$ and $h$ have components $1$ away from the finite place $v$, so archimedean components $1$, and
$\omega_0(g,h)f=\omega_f(g,h)f=f$ by property (5) of
Lemma~\ref{is-eigenvalue-conjugate-frame-fock-character-lattice}.

(2) Here $g$ and $h$ have finite components $1$, so $\omega_0(g,h)f=\omega_\infty(g,h)f=\prod_\rho\Delta_\rho\,f$ by
Lemma~\ref{is-eigenvalue-conjugate-frame-fock-character-fock}(4).
\end{proof}
% exists_weilDatum_etaTilde: of from FinWeilO.exists_isFinWeil (Lemma -fock-character-lattice) for Lambda_S, Lambda_T and
% (eta^n, eta^r) (isHeckeChar_pow, resEta_pow), oi from RealWeil.exists_isInfWeil (Lemma -fock-character-fock),
% hcomm = fin_inf_comm, om from exists_extend on thetaSub, ho, hp, hfin (hi.triv) and hinf (hf.triv); the clauses
% of IsWeilDatumO: exists_parts with actO_mul, Heisenberg2 of both, thetaDistribution_automorphy_source_group and
% thetaDistribution_invariant_rationalTarget, continuous_product, the unit clauses, and hf.split with
% mem_GfC_of_suppAt; the inert clause is hf.inert with hfin (mem_GfC_of_supp, infPart_placeIdem), and the clause
% VacCharPlus is hi.vac with hinf. The same Lean statement is Lemma
% asm-dependence-test-assembly-first-layer-ratio-inputs-weil-powers of the assembly, stated there for the lattices of
% the frames and without the last two clauses.

\begin{lemma}\label{is-eigenvalue-conjugate-frame-fock-character}
Let $D$ be a definite-side datum with host field $E_h/F_h$, $\psi$ a standard character of
$\mathbf A_{F_h}$ and $\theta$ an archimedean family whose member $\theta_\rho$ extends $\rho$ for every
$\rho\in\mathcal T$. Let $\tilde\eta$ be a Hecke character of $E_h$ that is unramified at every finite
place, restricts to the quadratic character $\eta$ of $E_h/F_h$ on $\mathbf A_{F_h}^\times$ and has
archimedean exponent $\frac12$ at every $\rho$. Let $S$ and $T$ be Hermitian spaces of $D$ of ranks $n$
and $r$, $\omega$ a Weil datum for $(S,T)$ with characters $(\chi_S,\chi_T)$, and $s_S$ and $s_T$ families
of Cartan involutions of $S$ and $T$ at the $\rho$. Let $g\in U(S)(\mathbf A)$ and $h\in U(T)(\mathbf A)$
have finite components $1$, with $s_{S,\rho}g_\rho=g_\rho s_{S,\rho}$ and $s_{T,\rho}h_\rho=h_\rho s_{T,\rho}$
for every $\rho$, let $x$ and $y$ be determinant roots of $g$ and $h$, and let $f$ be a theta function
annihilated by the annihilation operators for $s_S$ and $s_T$ at every $\rho$. Then
\[
 \omega(g,h)f=(\chi_T\tilde\eta^{-r})(x)^{-1}(\chi_S\tilde\eta^{-n})(y)
 \prod_{\rho\in\mathcal T}\Delta_\rho\,f,
\]
where $\Delta_\rho$ is the determinant of
$\frac12(1+s_{T,\rho}\otimes s_{S,\rho}^{\mathsf T})(h_\rho\otimes(g_\rho^{-1})^{\mathsf T})+\frac12(1-s_{T,\rho}\otimes
s_{S,\rho}^{\mathsf T})$.
\end{lemma}
\leanref{Development/AbelianVarietyTheory/Proofs/DivisorProofs/WeilDatumRealPlacesFockVacuumCharacterStatement.lean}{GlobalTheory.T.weilDatum_realPlaces_fockVacuum_character}
\leanref{Development/AbelianVarietyTheory/Proofs/DivisorProofs/WeilDatumRealPlacesFockVacuumCharacter.lean}{GlobalTheory.WeilVacChar.vacuum_character}
\begin{proof}
Let $\omega_0$ be a Weil datum for $(S,T)$ with characters $(\tilde\eta^n,\tilde\eta^r)$ as in
Lemma~\ref{is-eigenvalue-conjugate-frame-fock-character-model}, for the self-dual lattices of the frames of $S$ and $T$,
so that $\omega_0(g,h)f=\prod_\rho\Delta_\rho\,f$ by part (2) of that lemma.
Put $\nu_S=\chi_S\tilde\eta^{-n}$ and $\nu_T=\chi_T\tilde\eta^{-r}$. These are Hecke characters, as $\chi_S$,
$\chi_T$ and $\tilde\eta$ are, and they are trivial on $\mathbf A_{F_h}^\times$: there $\chi_S$ and
$\tilde\eta^n$ are both $\eta^n$, and $\chi_T$ and $\tilde\eta^r$ are both $\eta^r$, by
Definition~\ref{th-weil-datum}(3). By Lemma~\ref{asm-dependence-test-assembly-first-layer-ratio-inputs-weil-twist}
there is a Weil datum $\omega'$ for $(S,T)$ with characters $(\tilde\eta^n\nu_S,\tilde\eta^r\nu_T)=(\chi_S,\chi_T)$
and $\omega'(g,h)=\nu_T(x)^{-1}\nu_S(y)\,\omega_0(g,h)$. By
Lemma~\ref{is-eigenvalue-conjugate-theta-facts}(2), $\omega(g,h)f=\omega'(g,h)f$, as $f$ is a theta function.
So
\[
 \omega(g,h)f=\nu_T(x)^{-1}\nu_S(y)\,\omega_0(g,h)f=(\chi_T\tilde\eta^{-r})(x)^{-1}(\chi_S\tilde\eta^{-n})(y)
 \prod_\rho\Delta_\rho\,f.\qedhere
\]
% vacuum_character (15 lines): T.exists_weilDatum_powers_unramifiedEtaTilde (Lemma -fock-character-model) with
% the lattices PlaceModel.isSelfDualC_latFrame, its clause hvac (VacCharPlus); T.exists_weilDatum_twist_heckeCharacters
% (WeilChange.exists_weilDatum_twist) with hecke_mul_inv and resEta_mul_inv from the fields hecS, resS, hecT, resT of
% the two Weil data, then mul_mul_inv_cancel for the characters, htw, and T.weilDatum_unique_on_thetaFunctions. The
% determinant is WeilParts.detPlus sS sT (g^-1) h.
\end{proof}

\begin{lemma}\label{is-eigenvalue-conjugate-frame-compact-vacuum}
In Situation~\ref{is-eigenvalue-conjugate-frame-context}, of which only $\psi$, $\theta$, the Gram matrix of
$S$ and the parity of $n$ are used, let $\omega$ be a Weil datum for $(S,T)$ with characters
$(\chi_S,\chi_T)$, and suppose that $\eta$ is the restriction of a Hecke character $\tilde\eta$
as in Lemma~\ref{is-eigenvalue-conjugate-frame-fock-character}. Let $s=(s_\rho)$ be a family of Cartan
involutions of $T$, $k\in U(T)(\mathbf A)$ an element with finite component $1$ and $k_\rho s_\rho=s_\rho
k_\rho$ for every $\rho$, $y$ a determinant root of $k$, and $f$ a theta function annihilated by the
annihilation operators for the Cartan involutions $1$ of $S$ and $s_\rho$ of $T$ at every $\rho$. Put
$d^\pm_\rho=\det\big(\frac12k_\rho(1\pm s_\rho)+\frac12(1\mp s_\rho)\big)$, the determinants of $k_\rho$ on
the $\pm1$-eigenspaces of $s_\rho$, and
\[
 v_s(k)=\prod_{\rho\in\mathcal T}(d^+_\rho)^{n/2}(d^-_\rho)^{-n/2}.
\]
Then $\omega(1,k)f=\chi_S(y)\,v_s(k)\,f$.
\end{lemma}
\leanref{Development/AbelianVarietyTheory/Proofs/IwasawaDecomposition/RealPointMovesAdelicThetaByDetCharacter.lean}{GlobalTheory.OneMapRef.RealMove.om_compact}
\leanref{Development/AbelianVarietyTheory/Proofs/IwasawaDecomposition/ExistsIdeleAt.lean}{GlobalTheory.OneMapRef.RealMove.exists_root}
\leanref{Development/AbelianVarietyTheory/Proofs/IwasawaDecomposition/ExistsIdeleAt.lean}{GlobalTheory.OneMapRef.RealMove.detPlus_one}
\begin{proof}
Each $k_\rho$ is unitary for $G_T(\rho)$, so $|\det k_\rho|=1$. Let $y_0$ be the idele with finite component
$1$ whose archimedean component has image under $\theta_\rho$ a square root $t_\rho$ of $\det k_\rho$ for every
$\rho$. Then $y_0=\det(k)\,\bar y_0$, as $t_\rho=\det(k_\rho)\bar t_\rho$, so $y_0$ is a determinant root
of $k$. For a character $\nu$ of $\mathbf A_{E_h}^\times$ with integral archimedean exponents $c_\rho$,
$\nu(y_0)=\prod_\rho(t_\rho/|t_\rho|)^{2c_\rho}=\prod_\rho\det(k_\rho)^{c_\rho}$.

The matrix $1$ is a Cartan involution of $S$, as $S$ has Gram matrix $1$, and $1\in U(S)(\mathbf A)$ has
the determinant root $1$. Lemma~\ref{is-eigenvalue-conjugate-frame-fock-character}, applied to $g=1$,
$x=1$, $h=k$ and $y_0$, gives $\omega(1,k)f=(\chi_S\tilde\eta^{-n})(y_0)\prod_\rho\Delta_\rho\,f$. Here
$s_{S,\rho}=1$, so $\Delta_\rho=\det\big(\frac12(1+s_\rho)k_\rho+\frac12(1-s_\rho)\big)^n=(d^+_\rho)^n$, as $k_\rho$
commutes with $s_\rho$; and $\det k_\rho=d^+_\rho d^-_\rho$. Write $n=2m$. The character
$\tilde\eta\cdot\tilde\eta^{-(n+1)}=\tilde\eta^{-n}$ has archimedean exponent $\frac12-\frac{n+1}2=-m$ at every
$\rho$, so $\tilde\eta^{-n}(y_0)=\prod_\rho\det(k_\rho)^{-m}$. Hence
\[
 (\chi_S\tilde\eta^{-n})(y_0)\prod_\rho\Delta_\rho
 =\chi_S(y_0)\prod_\rho(d^+_\rho d^-_\rho)^{-m}(d^+_\rho)^{2m}=\chi_S(y_0)\,v_s(k).
\]
Finally $\chi_S$ restricts to $\eta^n$ on $\mathbf A_{F_h}^\times$ and $n$ is even, so
$\chi_S(y_0)=\chi_S(y)$ by Lemma~\ref{is-eigenvalue-conjugate-frame-root-values}(3).
\end{proof}
% om_compact: exists_root (y0 and hv0 for nu with HasInfExp), isCartanInv one for S (C.m.gS), isDetRoot one,
% weilDatum_realPlaces_fockVacuum_character with g = 1, root_eq (hW.resS, C.even), hv0 for eta * (eta^(n+1))^-1
% with exponents -m (hasInfExp_twist), detPlus_one (d1 for detPlus 1 s 1 k, d2 for det k = detPos detNeg) and
% the definition of vacChar with a = 0 (detPos, detNeg).

\begin{lemma}\label{is-eigenvalue-conjugate-frame-theta-space}
In Situation~\ref{is-eigenvalue-conjugate-frame-context}, let $f$ and $f'$ be theta functions for $(S,T)$
(Definition~\ref{th-weil-datum}(2)).
\begin{enumerate}
\item The functions $0$, $f+f'$ and $af$, for $a\in\mathbf C$, are theta functions.
\item If $f(w)=f'(w)$ for every $w\in\mathbb W$ with finite component $0$, then $f=f'$.
\end{enumerate}
\end{lemma}
\leanref{Development/AbelianVarietyTheory/Proofs/SiegelThetaTransformation/ThetaFunctionSpaceHeisenbergLaws.lean}{GlobalTheory.ThetaSpace.theta_space_laws}
\leanref{Development/AbelianVarietyTheory/Proofs/SiegelThetaTransformation/ThetaFunctionSpaceHeisenbergLaws.lean}{GlobalTheory.ThetaSpace.eq_zero_of_real}
\leanref{Development/AbelianVarietyTheory/Proofs/SiegelThetaTransformation/ThetaFunctionSpaceHeisenbergLaws.lean}{GlobalTheory.ThetaSpace.exists_level}
\begin{proof}
(1) The quasi-periodicity, the continuity and the smoothness of Definition~\ref{th-weil-datum}(2) hold for $0$
and pass to sums and scalar multiples. For the open subgroup, take $\mathbb W$ for $0$, the subgroup of $f$ for $af$,
and the intersection of the subgroups of $f$ and $f'$ for $f+f'$.

(2) For $y\in M_{n\times r}(E_h)$ let $y_f\in\mathbb W$ have archimedean component $0$ and finite component
$y^{\mathsf T}$; then $y^{\mathsf T}-y_f$, with $y^{\mathsf T}$ a rational point of $\mathbb W$, has finite
component $0$. Let $L=\mathcal O_{E_h}^n\otimes\Lambda'$ be the tensor lattice
(Definition~\ref{is-eigenvalue-conjugate-frame-heisenberg}(4)), and $\hat L$ the subgroup of $\mathbb W$
generated by the elements with archimedean component $0$ and finite component $\alpha y^{\mathsf T}$, with $\alpha$
in the closure $\hat{\mathcal O}$ of the integers of $F_h$ in $\mathbf A_{F_h,f}$ and $y\in L$. As $L$ is finitely
generated and spans $M_{n\times r}(E_h)$ over $E_h$, strong approximation for the additive group of $F_h$,
$\mathbf A_{F_h,f}=F_h+\hat{\mathcal O}$ \cite{Wei67}, together with $\hat{\mathcal O}=\mathcal O_{F_h}+N\hat{\mathcal
O}$ for every integer $N\ge1$ and the compactness of $\hat{\mathcal O}$, gives:
\begin{enumerate}
\item[(a)] for every $u\in\mathbb W$ with archimedean component $0$ there is $y\in M_{n\times r}(E_h)$ with
  $u-y_f\in\hat L$;
\item[(b)] for every $u\in\hat L$ and every integer $N\ge1$ there are $y\in L$ and $u'\in\hat L$ with
  $u=y_f+Nu'$;
\item[(c)] for every open subgroup $U$ of $\mathbb W$ there is an integer $N\ge1$ with $Nu\in U$ for every
  $u\in\hat L$.
\end{enumerate}
Put $g=f-f'$, a theta function by (1), with $g(w)=0$ for every $w$ with finite component $0$. Let $U$ be an open
subgroup with $\rho(w_0)g=g$ for every $w_0\in U$ with archimedean component $0$, and $N$ as in (c) for $U$; then
$\rho(Nu')g=g$ for every $u'\in\hat L$. Let $w\in\mathbb W$, and write $w=u+v$, where $u$ has archimedean component
$0$ and $v$ has finite component $0$. By (a) and (b), $u=y_f+y'_f+Nu'$ with $y\in M_{n\times r}(E_h)$, $y'\in L$
and $u'\in\hat L$. Put $\gamma=(y+y')^{\mathsf T}$, a rational point, and $v'=v-(\gamma-(y+y')_f)$, which has
finite component $0$; then $w=\gamma+(v'+Nu')$. As $v'$ and $Nu'$ have no place at which both components are
nonzero, $\langle v',Nu'\rangle=0$, and
\[
 g(v'+Nu')=\psi\big(\tfrac12\langle v',Nu'\rangle\big)\,g(v'+Nu')=(\rho(Nu')g)(v')=g(v')=0.
\]
By the quasi-periodicity of $g$ under the rational points, $g(w)=\psi(-\frac12\langle\gamma,v'+Nu'\rangle)\,
g(v'+Nu')=0$. Hence $g=0$ and $f=f'$.
\end{proof}
% theta_space_laws (ThetaFunctionSpaceHeisenbergLaws.lean, 149 lines) also has the clauses on rho(w) f, the
% composition of the rho, rho_finW, the rational action and the level, not used here. (1) is isTheta_zero,
% isTheta_add (U inf U') and isTheta_smul (HasDerivAtZeroAndIsTheta.lean); (2) is isTheta_sub and eq_zero_of_real:
% exists_fin_add_real (u + v), T.completedTensorLattice_strongApprox_open_compact for stdL n and C.X.Lambda (h1 is
% (a), h3 is (b), h4 is (c), used in exists_level for N), symplecticCharacter_real_finite_points_match (the
% rational point ratW minus finW is real, finW additive), symplOH_fin_real (the pairing of a real and a finite point
% is 0) and hf.quasi.

\begin{lemma}\label{is-eigenvalue-conjugate-frame-compact-move}
In Situation~\ref{is-eigenvalue-conjugate-frame-context}, with the notation of
Lemma~\ref{is-eigenvalue-conjugate-frame-zak}, let $B$ be orthonormal for $X'$, $k\in
U(T)(\mathbf A)$ an element with finite component $1$ and $k_\rho s'_{B,\rho}=s'_{B,\rho}k_\rho$ for every
$\rho$, and $\varphi$ a Schwartz-Bruhat function. Then, with $v_{s'_B}(k)$ as in
Lemma~\ref{is-eigenvalue-conjugate-frame-compact-vacuum},
\[
 \mathrm{Th}(x'_{kB},\varphi)=v_{s'_B}(k)\,\mathrm{Th}(x'_B,\varphi).
\]
The hypothesis that $k$ has finite component $1$ is not used.
\end{lemma}
\leanref{Development/AbelianVarietyTheory/Proofs/DivisorProofs/DetBlockConj.lean}{GlobalTheory.FrameEquiv.CompactMove.Th_mvB}
\leanref{Development/AbelianVarietyTheory/Proofs/DivisorProofs/DetBlockConj.lean}{GlobalTheory.FrameEquiv.CompactMove.blocks}
\leanref{Development/AbelianVarietyTheory/Proofs/DivisorProofs/DetBlockConj.lean}{GlobalTheory.FrameEquiv.CompactMove.thR_block}
\leanref{Development/AbelianVarietyTheory/Proofs/DivisorProofs/DetBlockConj.lean}{GlobalTheory.FrameEquiv.CompactMove.thR_of_perC}
\begin{proof}
Let $P$ be the diagonal $r\times r$ matrix with entries $1$ at the first $p'_0$ places and $0$ at the others,
so that $J'=2P-1$.

\emph{The blocks.} Fix $\rho$ and put $t=k_\rho$, $b=B_\rho$, $s=s'_{B,\rho}=bJ'b^{-1}$ and $m=b^{-1}tb$, so that
$tb=bm$. As $t$ commutes with $s$, $m$ commutes with $J'$ and with $P$. As $t$ is unitary for $G_T(\rho)$ and
$b^*G_T(\rho)b=J'$, $m^*J'm=J'$, so $m^*m=m^*J'mJ'=J'^2=1$. The determinants of $t$ on the $\pm1$-eigenspaces of
$s$ are $d^+_\rho=\det\big(\frac12t(1+s)+\frac12(1-s)\big)=\det(mP+1-P)$ and
$d^-_\rho=\det\big(\frac12t(1-s)+\frac12(1+s)\big)=\det(m(1-P)+P)$, conjugating by $b$. As $m$ is unitary and
commutes with the real diagonal projection $P$, $(Pm+1-P)^*(Pm+1-P)=1$; so
$\det(P\bar m+1-P)=\overline{\det(Pm+1-P)}=(d^+_\rho)^{-1}$, and in the same way
$\det((1-P)\bar m+P)=(d^-_\rho)^{-1}$.

\emph{The change of frame.} Let $x'\in\mathcal D_{X'}$, and let $c=(c_\sigma)_{\sigma\in\Psi'}$ be invertible
matrices commuting with $P$; write $x'c$ for the family $(x'_\sigma c_\sigma)_\sigma$. For a family $u$ of vectors
and an index $j$ with $(\tau'_j,a'_j)=(\sigma,a)$, the definition of the coordinates
(Definition~\ref{is-eigenvalue-conjugate-frame-fibres}(1)) gives $z'_{x'c}(u)_j=\sum_{a'\le p'_0}(c_\sigma)_{a'a}\,
z'_{x'}(u)_{j'}$ if $a\le p'_0$, and $z'_{x'c}(u)_j=\sum_{a'>p'_0}(c_\sigma^{-1})_{aa'}\,z'_{x'}(u)_{j'}$ if $a>p'_0$,
where $j'$ is the index with $(\tau'_{j'},a'_{j'})=(\sigma,a')$. So $z'_{x'c}(u)=N\,z'_{x'}(u)$ for the matrix $N$
that is block diagonal, in the indexing, with the blocks $N_\sigma=Pc_\sigma^{\mathsf T}+(1-P)c_\sigma^{-1}$. The
period matrices are $\Omega_{x'c}=N\Omega_{x'}$ and $\Omega'_{x'c}=N\Omega'_{x'}$; so $\tau_{x'c}=\tau_{x'}$, and
$\Omega_{x'c}^{-1}z'_{x'c}(u)=\Omega_{x'}^{-1}z'_{x'}(u)$. Hence the form $H$ evaluated at the coordinates, the
Gaussian and the exponential sum in the theta series do not change, and the factor
$\det(\Omega^{-1})^{n/2}$ is multiplied by $\det(N)^{-n/2}$:
$\Theta'_\varphi(x'c;u)=\det(N)^{-n/2}\,\Theta'_\varphi(x';u)$ for every family $u=(u_1,\dots,u_n)$. As $c_\sigma$ commutes
with $P$, $\det N_\sigma=\det(Pc_\sigma+1-P)/\det((1-P)c_\sigma+P)$.

\emph{Conclusion.} The component of $x'_{kB}$ at $\bar\rho$ is $\overline{tb}=\overline{bm}=x'_{B,\bar\rho}\,\bar m$; so
$x'_{kB}=x'_Bc$ with $c_{\bar\rho}=\overline{m}$ for the $m$ of $\rho$, which is invertible and commutes with $P$.
By the blocks, $\det N_{\bar\rho}=(d^+_\rho)^{-1}/(d^-_\rho)^{-1}=d^-_\rho/d^+_\rho$, so $\det(N)^{-n/2}=\prod_\rho
(d^+_\rho)^{n/2}(d^-_\rho)^{-n/2}=v_{s'_B}(k)$. The points $x'_B$ and $x'_{kB}$ lie in $\mathcal D_{X'}$
(Lemma~\ref{is-eigenvalue-conjugate-frame-zak}(3) and~(4)), so $\mathrm{Th}(x'_B,\varphi)$ and
$\mathrm{Th}(x'_{kB},\varphi)$ are Zak theta functions for these points
(Lemma~\ref{is-eigenvalue-conjugate-frame-zak}(1)), and for $w$ with finite component $0$
\[
 \mathrm{Th}(x'_{kB},\varphi)(w)=\Theta'_\varphi(x'_{kB};u'(w))=v_{s'_B}(k)\,\Theta'_\varphi(x'_B;u'(w))
 =v_{s'_B}(k)\,\mathrm{Th}(x'_B,\varphi)(w).
\]
The function $v_{s'_B}(k)\,\mathrm{Th}(x'_B,\varphi)$ is a theta function by
Lemma~\ref{is-eigenvalue-conjugate-frame-theta-space}(1). So the two sides of the claim are theta functions that
agree at every $w$ with finite component $0$, and they are equal by
Lemma~\ref{is-eigenvalue-conjugate-frame-theta-space}(2).
\end{proof}
% Th_mvB (30 lines, in DetBlockConj.lean, 180 lines): isZak_Th for z1 and z2; ThetaSpace.theta_space_laws (l1: two
% theta functions equal at the real points are equal, Lemma frame-theta-space(2); l2 for the multiple, (1));
% blocks for each tau (m with m1 to m4:
% block_of_inv for the commutation with P = prFst and the determinants, det_block_conj for the conjugates); and
% thR_block (thR_of_perC with M the block diagonal matrix of P b^T + (1 - P) b^-1, perC_block, det_tr_inv_block)
% with b = conj m and mat_xB; vacChar with a = 0 is the product of detPos^(n/2) detNeg^(-n/2). The hypothesis
% _hk : k in GinfC T is not used.

\begin{definition}\label{is-eigenvalue-conjugate-frame-overlap-doubling}
Let $X$ be flag data of rank $r$ over a CM field $E$ (Definition~\ref{is-onto-share-flag-data}), with set $\Psi$,
matrix $A$, integer $p_0$, element $\kappa_0$ and forms $h_\tau$, together with a symplectic basis and an indexing
$j\mapsto(\tau_j,a_j)$ (Definition~\ref{is-eigenvalue-conjugate-frame-fibres}(1) and~(2)). Put $d=|\Psi|r$, let $J$ be
the diagonal $r\times r$ matrix with the entries $\operatorname{sgn}(a)$, write $j(\tau,a)$ for the index $j$ with
$(\tau_j,a_j)=(\tau,a)$ and $u\cdot v=\sum_au_av_a$ for $u,v\in\mathbf C^r$, and let $n\ge0$. The inverse of $0$ and of a
singular matrix is read as $0$.
\begin{enumerate}
\item For $x_1,x_2\in M_r(\mathbf C)^\Psi$ and $\tau\in\Psi$ let $A_+=A_+(x_1,x_2)_\tau$ and $A_-=A_-(x_1,x_2)_\tau$ be the
  $r\times r$ matrices with the entries
  \[
   (A_+)_{ac}=\sum_b\tau(A_{cb})\overline{(x_{2,\tau})_{ba}},\qquad(A_-)_{ac}=\sum_b(x_{1,\tau})_{ba}\tau(A_{bc})
   \qquad(a\le p_0),
  \]
  \[
   (A_+)_{ac}=(x_{1,\tau}^{-1})_{ac},\qquad(A_-)_{ac}=\overline{(x_{2,\tau}^{-1})_{ac}}\qquad(a>p_0).
  \]
  So the $a$-th entry of $A_+v$ is $h_\tau(v,x_{2,\tau}e_a)$ if $a\le p_0$ and $(x_{1,\tau}^{-1}v)_a$ if $a>p_0$.
\item Put $Z(x_1,x_2)_\tau=-\frac14\tau(\kappa_0)\big(A_+^{-1}JA_+\big)^{\mathsf T}A^\tau$. For $z,w\in(\mathbf C^d)^n$ and
  $1\le k\le n$ let $v_k,v'_k\in\mathbf C^r$ have the $a$-th entries $z_{k,j(\tau,a)}$ and $w_{k,j(\tau,a)}$ if
  $a\le p_0$, and $w_{k,j(\tau,a)}$ and $z_{k,j(\tau,a)}$ if $a>p_0$, and put
  \[
   \zeta(x_1,x_2;z,w)_{\tau,k}=-\tfrac12\tau(\kappa_0)\,A^\tau A_-^{-1}v_k,\qquad
   \zeta'(x_1,x_2;z,w)_{\tau,k}=\tfrac12\tau(\kappa_0)\,(A^\tau)^{\mathsf T}A_+^{-1}v'_k.
  \]
\item For a square complex matrix $M$ put $\operatorname{Im}^{\mathrm h}M=(2i)^{-1}(M-M^*)$. The \emph{overlap kernel} is
  \[
   K(x_1,x_2;z,w)=\prod_{\tau\in\Psi}\big(\det A_+\det A_-\big)^{-\lfloor n/2\rfloor}
   \exp\Big(\pi\sum_{k=1}^n\zeta'_{\tau,k}\cdot\big(\operatorname{Im}^{\mathrm h}Z_\tau\big)^{-1}\zeta_{\tau,k}\Big),
  \]
  with $A_\pm=A_\pm(x_1,x_2)_\tau$, $Z=Z(x_1,x_2)$, $\zeta=\zeta(x_1,x_2;z,w)$ and $\zeta'=\zeta'(x_1,x_2;z,w)$. Put
  $K(x_1,x_2)=K(x_1,x_2;0,0)$.
\item A function $\psi$ on $M_{n\times r}(E)$ is \emph{Schwartz-Bruhat} if there is an integer $N\ge1$ such that
  $\psi(y)\ne0$ only if $Ny$ has entries in $\mathcal O_E$, and $\psi(y+Ny')=\psi(y)$ for every $y'$ with entries in
  $\mathcal O_E$. For $Z\in M_r(\mathbf C)^\Psi$ and families $\zeta=(\zeta_{\tau,k})$ and $\zeta'=(\zeta'_{\tau,k})$ of
  vectors of $\mathbf C^r$ indexed by $\Psi\times\{1,\dots,n\}$, its \emph{Hermitian theta series} is
  \[
   \theta^{\mathrm h}_\psi(Z;\zeta,\zeta')=\sum_{y\in M_{n\times r}(E)}\psi(y)\exp\Big(2\pi i\sum_{\tau\in\Psi}
   \sum_{k=1}^n\big(\tau(y_k)\cdot Z_\tau\overline{\tau(y_k)}+\tau(y_k)\cdot\zeta_{\tau,k}
   +\zeta'_{\tau,k}\cdot\overline{\tau(y_k)}\big)\Big),
  \]
  where $y_k$ is the $k$-th row of $y$, and $0$ when the sum does not converge absolutely.
\item Let $\varphi_1,\varphi_2$ be Schwartz-Bruhat functions on $M_{n\times d}(\mathbf Q)$. A \emph{doubling} of
  $(\varphi_1,\varphi_2)$ is a Schwartz-Bruhat function $\psi$ on $M_{n\times r}(E)$ such that
  \[
   \theta_{\varphi_1}(x_1;z)\,\overline{\theta_{\varphi_2}(x_2;\bar w)}=K(x_1,x_2;z,w)\,
   \theta^{\mathrm h}_\psi\big(Z(x_1,x_2);\zeta(x_1,x_2;z,w),\zeta'(x_1,x_2;z,w)\big)
  \]
  for all $x_1,x_2\in\mathcal D_X$ and $z,w\in(\mathbf C^d)^n$, with the theta series of
  Definition~\ref{is-eigenvalue-conjugate-frame-fibres}(6). Let $\delta(\varphi_1,\varphi_2)$ be a doubling of
  $(\varphi_1,\varphi_2)$, chosen once for all, when there is one, and $0$ otherwise.
\end{enumerate}
\end{definition}
\leanref{Development/AbelianVarietyTheory/Definitions/SpanningFamilies/IsSystem.lean}{GlobalTheory.SW.aLin}
\leanref{Development/AbelianVarietyTheory/Definitions/SpanningFamilies/IsSystem.lean}{GlobalTheory.SW.aAnti}
\leanref{Development/AbelianVarietyTheory/Definitions/SpanningFamilies/IsSystem.lean}{GlobalTheory.SW.dblZ}
\leanref{Development/AbelianVarietyTheory/Definitions/SpanningFamilies/IsSystem.lean}{GlobalTheory.SW.dblz}
\leanref{Development/AbelianVarietyTheory/Definitions/SpanningFamilies/IsSystem.lean}{GlobalTheory.SW.dblz'}
\leanref{Development/AbelianVarietyTheory/Definitions/SpanningFamilies/IsSystem.lean}{GlobalTheory.SW.kZero}
\leanref{Development/AbelianVarietyTheory/Definitions/SpanningFamilies/FirstCoordinateAdmissible.lean}{GlobalTheory.SW.imH}
\leanref{Development/AbelianVarietyTheory/Definitions/SpanningFamilies/FirstCoordinateAdmissible.lean}{GlobalTheory.SW.expo}
\leanref{Development/AbelianVarietyTheory/Definitions/SpanningFamilies/FirstCoordinateAdmissible.lean}{GlobalTheory.SW.thH}
\leanref{Development/AbelianVarietyTheory/Definitions/SpanningFamilies/IsSystem.lean}{GlobalTheory.SW.IsDbl}
\leanref{Development/AbelianVarietyTheory/Definitions/SpanningFamilies/IsSystem.lean}{GlobalTheory.SW.dbl}
% aLin x1 x2 tau (a, c) = sum_b tau(Gm c b) conj(x2 tau b a) for a < p0 (0-based), (mat x1 tau)^-1 a c otherwise; aAnti
% = sum_b x1 tau b a tau(Gm b c), conj((mat x2 tau)^-1 a c). dblZ = (tau kappa0 / 4) (aLin^-1 diag(-sgn) aLin)^T Gm^tau;
% dblz row k = -(tau kappa0 / 2) Gm^tau (aAnti^-1 v_k), dblz' row k = (tau kappa0 / 2) (aLin^-1 v'_k) vecMul Gm^tau,
% with v_k, v'_k read from z1 and w through idx as in (2).
% kZero = prod_tau ((det aLin det aAnti)^(n/2))^-1 exp(pi sum_k dblz' . (imH dblZ)^-1 dblz). thH is a tsum (0 when
% not summable; for complex families summability is absolute summability). IsDbl phi1 phi2 psi is the displayed
% identity for x1, x2 in dom and all z1, w; dbl is the choice of an IsSB psi with IsDbl, and 0 if there is none. IsSB on
% Matrix (Fin n) (Fin r) E is HasLevel R (IntM: entries integral over Z).

\begin{lemma}\label{is-eigenvalue-conjugate-frame-overlap-determinants}
Let $X$ be flag data of rank $r$ over a CM field $E$ (Definition~\ref{is-onto-share-flag-data}), with set $\Psi$, let
$x_1,x_2\in\mathcal D_X$ and $\tau\in\Psi$, and let $A_\pm$, $Z$ and $\operatorname{Im}^{\mathrm h}$ be as in
Definition~\ref{is-eigenvalue-conjugate-frame-overlap-doubling}. Then $\det A_+(x_1,x_2)_\tau\ne0$,
$\det A_-(x_1,x_2)_\tau\ne0$, $A_+(x_2,x_1)_\tau=\overline{A_-(x_1,x_2)_\tau}$, and
$\operatorname{Im}^{\mathrm h}Z(x_1,x_2)_\tau$ is positive definite.
\end{lemma}
\leanref{Development/AbelianVarietyTheory/Proofs/SiegelThetaTransformation/DblZMemHn.lean}{GlobalTheory.SW.dblZ_mem_Hn}
\leanref{Development/AbelianVarietyTheory/Proofs/SiegelThetaTransformation/DblZMemHn.lean}{GlobalTheory.SW.DblPoint.det_aLin_ne_zero}
\leanref{Development/AbelianVarietyTheory/Proofs/SiegelThetaTransformation/DblZMemHn.lean}{GlobalTheory.SW.DblPoint.det_aAnti_ne_zero}
\leanref{Development/AbelianVarietyTheory/Proofs/SiegelThetaTransformation/DblZMemHn.lean}{GlobalTheory.SW.DblPoint.dblZ_mem}
\begin{proof}
Write $h=h_\tau$, $A_\pm=A_\pm(x_1,x_2)_\tau$, and $x_{i,1},\dots,x_{i,r}$ for the columns of $x_{i,\tau}$. As
$\bar A_{ab}=A_{ba}$, $h(v,v)$ is real.

\emph{Two signs.} (a) If $v\ne0$ and $(A_+v)_a=0$ for every $a>p_0$, then $h(v,v)>0$. Indeed,
$(x_{1,\tau}^{-1}v)_a=0$ for $a>p_0$, so $v$ is a combination of $x_{1,1},\dots,x_{1,p_0}$ with coefficients not all $0$,
and the Gram matrix of these columns for $h$ is positive definite, by
Lemma~\ref{is-onto-share-counts-section-kronecker-solve-chart-domain}. (b) If $v\ne0$ and $(A_+v)_a=0$ for every
$a\le p_0$, then $h(v,v)<0$. Indeed, $h(v,x_{2,a})=0$ for $a\le p_0$, and the first step of the proof of
Lemma~\ref{is-onto-share-counts-section-kronecker-solve-chart-domain}, which uses only that the Gram matrix of
$x_{2,1},\dots,x_{2,p_0}$ is positive definite, gives $h(v,v)<0$.

\emph{The determinants.} If $A_+v=0$ with $v\ne0$, then (a) and (b) both apply, which is absurd; so $\det A_+\ne0$.
Entry by entry, as $\overline{\tau(A_{bc})}=\tau(A_{cb})$, $A_+(x_2,x_1)_\tau=\overline{A_-}$. So
$\det A_-=\overline{\det A_+(x_2,x_1)_\tau}\ne0$, by the first assertion for $(x_2,x_1)$.

\emph{The imaginary part.} Put $S=A_+^{-1}(-J)A_+$, so that $Z_\tau=Z(x_1,x_2)_\tau=\frac14\tau(\kappa_0)S^{\mathsf T}A^\tau$.
For $v\in\mathbf C^r$, $v^*S^{\mathsf T}A^\tau v=\sum_{l,j}(S\bar v)_l\tau(A_{lj})v_j=h(S\bar v,\bar v)$. The matrix
$\operatorname{Im}^{\mathrm h}Z_\tau$ is Hermitian, with $v^*(\operatorname{Im}^{\mathrm h}Z_\tau)v=\operatorname{Im}(v^*Z_\tau v)$.
As $\bar\kappa_0=-\kappa_0$, $\tau(\kappa_0)$ is purely imaginary, and $\operatorname{Im}\tau(\kappa_0)<0$ by (6) of
Definition~\ref{is-onto-share-flag-data}; so $\operatorname{Im}(v^*Z_\tau v)=\frac14\operatorname{Im}\tau(\kappa_0)\,
\operatorname{Re}h(S\bar v,\bar v)$, and it suffices to show $\operatorname{Re}h(Sd,d)<0$ for $d\ne0$. Write $A_+d=y_++y_-$, with
$y_+$ supported on the indices $a\le p_0$ and $y_-$ on the others, and put $d_\pm=A_+^{-1}y_\pm$. Then $d=d_++d_-$ and
$Sd=d_--d_+$. As $\operatorname{Re}h(d_+,d_-)=\operatorname{Re}h(d_-,d_+)$,
\[
 \operatorname{Re}h(Sd,d)=h(d_-,d_-)-h(d_+,d_+).
\]
By (a) for $d_+$ and (b) for $d_-$, $h(d_+,d_+)>0$ unless $d_+=0$ and $h(d_-,d_-)<0$ unless $d_-=0$; as $d\ne0$, one of
them is not $0$, and $\operatorname{Re}h(Sd,d)<0$.
\end{proof}
% dblZ_mem_Hn (DblZMemHn.lean, 220 lines): its first clause IsType X.Psi is (1) of the flag data. Followed:
% pos_of_upper (a, through Data.gramPos_posDef), neg_of_lower (b, Data.neg_of_perp), det_aLin_ne_zero, aLin_swap,
% det_aAnti_ne_zero, re_hermC_invol_neg (yp, ym, d1, d2, hS, key), posDef_imH and dblZ_mem (re_kappa, kappa0_sign).

\begin{lemma}\label{is-eigenvalue-conjugate-frame-overlap-series}
Let $E$ be a CM field, $\Psi$ a set of embeddings $E\to\mathbf C$ that contains $\tau$ if and only if it does not
contain $\bar\tau$, and $n,N\ge0$; define Schwartz-Bruhat functions on $M_{n\times N}(E)$ and their Hermitian theta
series as in Definition~\ref{is-eigenvalue-conjugate-frame-overlap-doubling}(4), with $N$ in place of $r$. Let $\psi$
be a Schwartz-Bruhat function on $M_{n\times N}(E)$ and $Z\in M_N(\mathbf C)^\Psi$ with
$\operatorname{Im}^{\mathrm h}Z_\tau$ positive definite for every $\tau\in\Psi$. Then the series of
$\theta^{\mathrm h}_\psi(Z;\zeta,\zeta')$ converges absolutely for all $\zeta,\zeta'$; and if $\psi'$ is a
Schwartz-Bruhat function on $M_{n\times N}(E)$ with $\theta^{\mathrm h}_{\psi'}(Z;\zeta,\zeta')=\theta^{\mathrm
h}_\psi(Z;\zeta,\zeta')$ for all $\zeta,\zeta'$, then $\psi'=\psi$.
\end{lemma}
\leanref{Development/AbelianVarietyTheory/Proofs/SiegelThetaTransformation/HermitianThetaSummableAndDeterminesTestFnStatement.lean}{GlobalTheory.T.hermitianTheta_summable_and_determines_testFn}
\leanref{Development/AbelianVarietyTheory/Proofs/SiegelThetaTransformation/HermitianThetaSummableAndDeterminesTestFn.lean}{GlobalTheory.SW.thH_summable_and_inj}
\begin{proof}
\openstep{The proof, which writes the Hermitian theta series as a Siegel theta series in rational coordinates of
$M_{n\times N}(E)$ for a $\mathbf Q$-basis of $E$, with a period matrix in the Siegel upper half space, and uses the
absolute convergence of Siegel theta series and the fact that a Siegel theta series determines its Schwartz-Bruhat
function, is not written out.}
% thH_summable_and_inj (HermitianThetaSummableAndDeterminesTestFn.lean): thH_eq_thS_dictionary (ofQ, tauOf, zetaOf;
% hup: tauOf Z in Siegel for Z in Hn; hSB), SiegelSum.summable_norm_thetaTerm and eq_of_thS_eq. Size: 535 lines
% (lean_size --union of thH_summable_and_inj, Dict.ofQEquiv, Dict.tauOf_mem, Dict.zetaOf_surj, Dict.thH_eq,
% Dict.isSB_iff, summable_norm_thetaTerm and eq_of_thS_eq; lean_size of the theorem alone misses the _root_ names).
\end{proof}

\begin{lemma}\label{is-eigenvalue-conjugate-frame-overlap-doubling-exists}
Let $X$, a symplectic basis and an indexing be as in
Definition~\ref{is-eigenvalue-conjugate-frame-overlap-doubling}, let $n\ge0$ be even, and let $\varphi_1,\varphi_2$ be
Schwartz-Bruhat functions on $M_{n\times d}(\mathbf Q)$. Then $\delta(\varphi_1,\varphi_2)$ is a doubling of
$(\varphi_1,\varphi_2)$, and every doubling of $(\varphi_1,\varphi_2)$ is $\delta(\varphi_1,\varphi_2)$.
\end{lemma}
\leanref{Development/AbelianVarietyTheory/Proofs/SiegelThetaTransformation/DoubledHermitianTestFunctionSpec.lean}{GlobalTheory.SW.DblEx.dbl_spec}
\leanref{Development/AbelianVarietyTheory/Proofs/SiegelThetaTransformation/DoubledHermitianTestFunctionSpec.lean}{GlobalTheory.SW.DblEx.exists_isDbl}
\leanref{Development/AbelianVarietyTheory/Proofs/SiegelThetaTransformation/DoubledHermitianTestFunctionSpec.lean}{GlobalTheory.SW.DblEx.dblz_surj}
\begin{proof}
\emph{Existence.} There is a doubling of $(\varphi_1,\varphi_2)$.
\openstep{The proof of the existence of a doubling, which writes the product of the theta series at two points as a
Siegel theta series of the doubled symplectic space, passes to symplectic coordinates adapted to the doubled Hermitian
space, and transforms that theta series by the resulting rational symplectic matrix, is not written out.}
% exists_isDbl (DoubledHermitianTestFunctionSpec.lean, about 90 lines): T.exists_symplectic_doubled_hermitian_coordinates
% (F), T.siegelTheta_mul_conj_eq_blockTheta, can_trans (the transformation law of Siegel theta series),
% T.det_doubledCoordinates_to_rationalCoords (c, L), DblPer.dbl_per and zetaOf_period; the doubling is
% c^(n/2) psi o ofQ^-1. Size: 3897 lines (lean_size exists_isDbl).
So $\delta(\varphi_1,\varphi_2)$ is a doubling, by its definition.

\emph{Uniqueness.} Let $\psi$ be a doubling, and put $\psi_0=\delta(\varphi_1,\varphi_2)$. The flag domain $\mathcal D_X$
is not empty: by (2) of Definition~\ref{is-onto-share-flag-data} there is $x$ whose components have columns with Gram
matrix $J$ for $h_\tau$, and then $\det x_\tau\ne0$ and $d_i(x,\tau)=\epsilon_i$. Fix such an $x$. By
Lemma~\ref{is-eigenvalue-conjugate-frame-overlap-determinants}, the matrices $A_\pm(x,x)_\tau$ are invertible and the
$\operatorname{Im}^{\mathrm h}Z(x,x)_\tau$ are positive definite. Let $\zeta,\zeta'$ be families of vectors of $\mathbf
C^r$ indexed by $\Psi\times\{1,\dots,n\}$. As $\tau(\kappa_0)\ne0$ and $\det A^\tau\ne0$, the vectors
\[
 v_k=-\tfrac2{\tau(\kappa_0)}A_-(x,x)_\tau(A^\tau)^{-1}\zeta_{\tau,k},\qquad
 v'_k=\tfrac2{\tau(\kappa_0)}A_+(x,x)_\tau\big((A^\tau)^{\mathsf T}\big)^{-1}\zeta'_{\tau,k}
\]
are those of Definition~\ref{is-eigenvalue-conjugate-frame-overlap-doubling}(2) for the $z,w\in(\mathbf C^d)^n$ with
$z_{k,j(\tau,a)}=(v_k)_a$, $w_{k,j(\tau,a)}=(v'_k)_a$ if $a\le p_0$ and $z_{k,j(\tau,a)}=(v'_k)_a$,
$w_{k,j(\tau,a)}=(v_k)_a$ if $a>p_0$; so $\zeta(x,x;z,w)=\zeta$ and $\zeta'(x,x;z,w)=\zeta'$. The factor $K(x,x;z,w)$ is
not $0$, as its determinants are not. The identities of a doubling for $\psi$ and $\psi_0$ at $(x,x,z,w)$ have the same
left side; so $\theta^{\mathrm h}_\psi(Z(x,x);\zeta,\zeta')=\theta^{\mathrm h}_{\psi_0}(Z(x,x);\zeta,\zeta')$ for all
$\zeta,\zeta'$, and $\psi=\psi_0$ by Lemma~\ref{is-eigenvalue-conjugate-frame-overlap-series}.
\end{proof}
% dbl_spec (DoubledHermitianTestFunctionSpec.lean), followed: exists_isDbl and dif_pos for the first two clauses;
% FiniteNet.dom_nonempty, dblz_surj (the u and v above), DblPoint.dblZ_mem, the determinants, and
% T.hermitianTheta_summable_and_determines_testFn for the uniqueness.

\begin{lemma}\label{is-eigenvalue-conjugate-frame-overlap-kernel-shift}
Let $X$, a symplectic basis and an indexing be as in
Definition~\ref{is-eigenvalue-conjugate-frame-overlap-doubling}, let $n\ge0$ be even and $x_1,x_2\in\mathcal D_X$, and
let $U$ and $E_X$ be as in Lemma~\ref{is-eigenvalue-conjugate-frame-zak-real}; for $u\in U$ write
$z_i(u)=(z_{x_i}(u_1),\dots,z_{x_i}(u_n))\in(\mathbf C^d)^n$. Then for all $a,b,w\in U$
\[
 e^{\pi iE_X(w,a-b)}\,\mathcal G_{x_1}(z_1(w+a))\,\overline{\mathcal G_{x_2}(z_2(w+b))}\,
 K\big(x_1,x_2;z_1(w+a),\overline{z_2(w+b)}\big)
 =\mathcal G_{x_1}(z_1(a))\,\overline{\mathcal G_{x_2}(z_2(b))}\,K\big(x_1,x_2;z_1(a),\overline{z_2(b)}\big).
\]
\end{lemma}
\leanref{Development/AbelianVarietyTheory/Proofs/SiegelThetaTransformation/DoubledGaussianKernelDiagonalTranslationInvarianStatement.lean}{GlobalTheory.T.doubledGaussianKernel_diagonal_translation_invariant}
\leanref{Development/AbelianVarietyTheory/Proofs/SiegelThetaTransformation/DoubledGaussianKernelDiagonalTranslationInvarian.lean}{GlobalTheory.Ker.Diag.kernel_diag}
\begin{proof}
\openstep{The proof, which compares the exponents of the two sides as continuous functions of $w$ and checks that they
agree at the dense set of the $\hat y$, $y\in M_{n\times r}(E)$, through the periods of the rows of $y$ and the
transformation of the forms $H_{x_1}$ and $H_{x_2}$ under periods, is not written out.}
% kernel_diag (DoubledGaussianKernelDiagonalTranslationInvarian.lean): T.realPeriodMap_homeomorph_riemannForm_density
% (zR additive, Im sum hx = EFR, density of embR), Dict.fock_hx_add_left, _add_right, _smul_left, fock_hx_conj,
% kRow_per (the exponent of K at periods) and continuity. Size: 1906 lines (lean_size kernel_diag).
\end{proof}

\begin{lemma}\label{is-eigenvalue-conjugate-frame-overlap-mean}
Let $X$, a symplectic basis and an indexing be as in
Definition~\ref{is-eigenvalue-conjugate-frame-overlap-doubling}, with module $\Lambda$, let $n\ge0$ be even,
$x_1,x_2\in\mathcal D_X$, and $\varphi_1,\varphi_2$ Schwartz-Bruhat functions on $M_{n\times d}(\mathbf Q)$. Let $U$,
$\hat y$ and $E_X$ be as in Lemma~\ref{is-eigenvalue-conjugate-frame-zak-real}, with Lebesgue measure on $U$, write
$z_i(u)=(z_{x_i}(u_1),\dots,z_{x_i}(u_n))$ for $u\in U$, and for an integer $N>0$ let $\Gamma_N\subset U$ be the subgroup
generated by the $N\hat y$ with $y\in M_{n\times r}(E)$ whose rows lie in $\Lambda$. Then there is an integer $N_0>0$
such that for every positive multiple $N$ of $N_0$, every fundamental domain $F$ of $\Gamma_N$ in $U$ and all $a,b\in U$,
\[
 \frac1{\operatorname{vol}(F)}\int_Fe^{\pi iE_X(w,a-b)}\,\Theta_{\varphi_1}(x_1;w+a)\,
 \overline{\Theta_{\varphi_2}(x_2;w+b)}\,dw
 =\delta(\varphi_1,\varphi_2)(0)\,\mathcal G_{x_1}(z_1(a))\,\overline{\mathcal G_{x_2}(z_2(b))}\,
 K\big(x_1,x_2;z_1(a),\overline{z_2(b)}\big),
\]
with $\delta$ and $K$ as in Definition~\ref{is-eigenvalue-conjugate-frame-overlap-doubling}.
\end{lemma}
\leanref{Development/AbelianVarietyTheory/Proofs/SiegelThetaTransformation/MeanRealTorusNormalisedThetaMulConjStatement.lean}{GlobalTheory.T.mean_realTorus_normalisedTheta_mul_conj}
\leanref{Development/AbelianVarietyTheory/Proofs/SiegelThetaTransformation/MeanRealTorusNormalisedThetaMulConj.lean}{GlobalTheory.Ker.Mean.mean_thR_mul_conj}
\leanref{Development/AbelianVarietyTheory/Proofs/SiegelThetaTransformation/MeanRealTorusNormalisedThetaMulConj.lean}{GlobalTheory.Ker.Mean.expo_shift}
\leanref{Development/AbelianVarietyTheory/Proofs/SiegelThetaTransformation/MeanRealTorusNormalisedThetaMulConj.lean}{GlobalTheory.Ker.Mean.mean_tsum}
\begin{proof}
Write $\psi=\delta(\varphi_1,\varphi_2)$, $Z=Z(x_1,x_2)$, $A_\pm=A_\pm(x_1,x_2)_\tau$, and for $a,b\in U$ write
$\zeta_{a,b}=\zeta(x_1,x_2;z_1(a),\overline{z_2(b)})$, $\zeta'_{a,b}=\zeta'(x_1,x_2;z_1(a),\overline{z_2(b)})$ and
$K_{a,b}=\mathcal G_{x_1}(z_1(a))\overline{\mathcal G_{x_2}(z_2(b))}K(x_1,x_2;z_1(a),\overline{z_2(b)})$. The coordinates
$z_x(u)$ are additive in $u$. We use two facts on $E_X$: $E_X(\hat y,\hat y')=\sum_k\langle y_k,y'_k\rangle_X$ for
$y,y'\in M_{n\times r}(E)$, as $\Psi$ contains one embedding of each pair of conjugate embeddings, so that $\sum_{\sigma\in\Psi}
2\operatorname{Re}\sigma(\alpha)=\operatorname{Tr}_{E/\mathbf Q}\alpha$; and the $\hat y$ are dense in $U$, as $\alpha\mapsto(\sigma(\alpha))_{\sigma\in\Psi}$
identifies $E\otimes_{\mathbf Q}\mathbf R$ with $\mathbf C^\Psi$.

\emph{The integer $N_0$.} By Lemma~\ref{is-eigenvalue-conjugate-frame-overlap-doubling-exists}, $\psi$ is a doubling of
$(\varphi_1,\varphi_2)$; let $R\ge1$ be an integer for the Schwartz-Bruhat function $\psi$ as in
Definition~\ref{is-eigenvalue-conjugate-frame-overlap-doubling}(4). As $\Lambda$ spans $E^r$, there is an integer $M\ge1$
with $Mv\in\Lambda$ for every $v\in E^r$ with entries in $\mathcal O_E$. Put $N_0=MR$, and let $N=tN_0$, $F$, $a$ and $b$
be as in the statement.

\emph{The integrand as a series.} For $y\in M_{n\times r}(E)$ and $w\in U$ put $\chi_y(w)=e^{2\pi iE_X(\hat y,w)}$ and
\[
 c(y)=\psi(y)\exp\Big(2\pi i\sum_{\tau\in\Psi}\sum_{k=1}^n\big(\tau(y_k)\cdot Z_\tau\overline{\tau(y_k)}
 +\tau(y_k)\cdot(\zeta_{a,b})_{\tau,k}+(\zeta'_{a,b})_{\tau,k}\cdot\overline{\tau(y_k)}\big)\Big).
\]
We show that for every $w\in U$
\[
 e^{\pi iE_X(w,a-b)}\,\Theta_{\varphi_1}(x_1;w+a)\,\overline{\Theta_{\varphi_2}(x_2;w+b)}=K_{a,b}\sum_yc(y)\chi_y(w).
 \tag{$*$}
\]
By Definition~\ref{is-eigenvalue-conjugate-frame-fibres}(6),
$\Theta_{\varphi_i}(x_i;u)=\mathcal G_{x_i}(z_i(u))\,\theta_{\varphi_i}(x_i;z_i(u))$. As $\psi$ is a doubling, at
$z=z_1(w+a)$ and $\overline{z_2(w+b)}$,
\[
 \theta_{\varphi_1}(x_1;z_1(w+a))\,\overline{\theta_{\varphi_2}(x_2;z_2(w+b))}
 =K\big(x_1,x_2;z_1(w+a),\overline{z_2(w+b)}\big)\,\theta^{\mathrm h}_\psi(Z;\zeta_{w+a,w+b},\zeta'_{w+a,w+b}),
\]
and by Lemma~\ref{is-eigenvalue-conjugate-frame-overlap-kernel-shift} the left side of $(*)$ is $K_{a,b}\,\theta^{\mathrm
h}_\psi(Z;\zeta_{w+a,w+b},\zeta'_{w+a,w+b})$. To compare the series, let $u\in\mathbf C^r$ be the vector $w_{k,\tau}$. For
$a'\le p_0$ the entry of $z_1(w)_k$ at $j(\tau,a')$ is $h_\tau(x_{1,\tau}e_{a'},u)=(A_-\bar u)_{a'}$, and for $a'>p_0$ the
entry of $\overline{z_2(w)_k}$ there is $\overline{(x_{2,\tau}^{-1}u)_{a'}}=(A_-\bar u)_{a'}$; for $a'\le p_0$ the entry
of $\overline{z_2(w)_k}$ is $\overline{h_\tau(x_{2,\tau}e_{a'},u)}=h_\tau(u,x_{2,\tau}e_{a'})=(A_+u)_{a'}$, and for $a'>p_0$
the entry of $z_1(w)_k$ is $(x_{1,\tau}^{-1}u)_{a'}=(A_+u)_{a'}$. As $A_\pm$ are invertible
(Lemma~\ref{is-eigenvalue-conjugate-frame-overlap-determinants}),
\[
 (\zeta_{w+a,w+b})_{\tau,k}=(\zeta_{a,b})_{\tau,k}-\tfrac12\tau(\kappa_0)A^\tau\bar u,\qquad
 (\zeta'_{w+a,w+b})_{\tau,k}=(\zeta'_{a,b})_{\tau,k}+\tfrac12\tau(\kappa_0)(A^\tau)^{\mathsf T}u.
\]
As $\overline{\tau(\kappa_0)}=-\tau(\kappa_0)$, the two added terms contribute
$\tau(y_k)\cdot(-\frac12\tau(\kappa_0)A^\tau\bar u)+\frac12\tau(\kappa_0)(A^\tau)^{\mathsf T}u\cdot\overline{\tau(y_k)}
=2\operatorname{Re}\big(\tau(-\tfrac12\kappa_0)h_\tau(\tau(y_k),w_{k,\tau})\big)$ to the exponent of the term of $y$; summed over
$\tau$ and $k$ this is $E_X(\hat y,w)$. Both Hermitian theta series converge absolutely, by
Lemmas~\ref{is-eigenvalue-conjugate-frame-overlap-determinants} and~\ref{is-eigenvalue-conjugate-frame-overlap-series};
so $\theta^{\mathrm h}_\psi(Z;\zeta_{w+a,w+b},\zeta'_{w+a,w+b})=\sum_yc(y)\chi_y(w)$, which proves $(*)$.

\emph{The characters.} The sum $\sum_y|c(y)|$ is finite by
Lemma~\ref{is-eigenvalue-conjugate-frame-overlap-series}. Each $\chi_y$ is continuous, with $|\chi_y|=1$ and
$\chi_y(v+w)=\chi_y(v)\chi_y(w)$, as $E_X$ is real and additive in its second argument; and $\chi_0=1$.
\begin{itemize}
\item If $c(y)\ne0$, then $\chi_y=1$ on $\Gamma_N$. Indeed $\psi(y)\ne0$, so $Ry$ has entries in $\mathcal O_E$ and
  $MRy_k\in\Lambda$ for every $k$. For $y'$ with rows in $\Lambda$,
  $E_X(\hat y,N\hat y')=t\,E_X(\widehat{MRy},\hat y')=t\sum_k\langle MRy_k,y'_k\rangle_X$, which is an integer, as
  $\langle\ ,\ \rangle_X$ takes the values $0$ and $\pm1$ on the symplectic basis of $\Lambda$. So $\chi_y(N\hat y')=1$, and
  $\chi_y=1$ on the group $\Gamma_N$ that these elements generate.
\item If $y\ne0$, then $\chi_y(t')\ne1$ for some $t'\in U$. Let $q\in U$ have $q_{k,\tau}=\tau(-\frac12\kappa_0)(A^\tau)^{\mathsf
  T}\tau(y_k)$. Then $\tau(-\frac12\kappa_0)h_\tau(\tau(y_k),q_{k,\tau})=|q_{k,\tau}|^2$, so
  $E_X(\hat y,q)=2\sum_{k,\tau}|q_{k,\tau}|^2$, which is not $0$, as $\kappa_0\ne0$, $\det A\ne0$, $\Psi$ is not empty and
  some $y_k\ne0$. For $t'=(2E_X(\hat y,q))^{-1}q$, $\chi_y(t')=e^{\pi i}=-1$.
\item The countable group $\Gamma_N$ spans $U$ over $\mathbf R$: for $y\in M_{n\times r}(E)$ some integer $M'\ge1$ has
  $M'y_k\in\Lambda$ for every $k$, so $\hat y=(NM')^{-1}N\widehat{M'y}$ lies in the span, and the $\hat y$ are dense.
\end{itemize}

\emph{The mean.} As $\Gamma_N$ spans $U$, the fundamental domain $F$ has finite positive volume. By $(*)$, as $\sum_y|c(y)|$ is
finite and $|\chi_y|=1$, the left side of the claim is $K_{a,b}\sum_yc(y)\operatorname{vol}(F)^{-1}\int_F\chi_y(w)\,dw$. The
term of $y=0$ is $K_{a,b}c(0)=K_{a,b}\psi(0)$. Let $y\ne0$ with $c(y)\ne0$, and $t'$ as above. As $\chi_y$ is invariant under
$\Gamma_N$ and $t'+F$ is also a fundamental domain of $\Gamma_N$,
$\int_F\chi_y=\int_{t'+F}\chi_y=\chi_y(t')\int_F\chi_y$; so $\int_F\chi_y=0$. Hence the left side is
$K_{a,b}\psi(0)$, which is the claim.
\end{proof}
% mean_thR_mul_conj (MeanRealTorusNormalisedThetaMulConj.lean, 380 lines), followed: DblEx.dbl_spec (Lemma
% -overlap-doubling-exists), SW.dblZ_mem_Hn (Lemma -overlap-determinants), the clauses hEF and hdense of
% T.realPeriodMap_homeomorph_riemannForm_density (EFR_embR and the density of embR, which the Lean derives from the
% density of the periods, denseRange_zR_embR; the book gives the two short reasons above), SW.thH_summable_and_inj (Lemma
% -overlap-series), exists_int_mem (M, from exists_common), N0 = M R; hint with e1 (the doubling), e2
% (T.doubledGaussianKernel_diagonal_translation_invariant, Lemma -overlap-kernel-shift) and e3 (expo_shift through
% dblz_shift and dblz'_shift); mean_tsum (measure_lt_top_of_span, integral_char_eq_zero) with span_latR, the character
% conditions (IsSympl for the integrality of EF on Lambda, EFR_nsmul_left) and exists_EFR_ne.

\begin{lemma}\label{is-eigenvalue-conjugate-frame-real-character-overlap}
In Situation~\ref{is-eigenvalue-conjugate-frame-context}, with the notation of
Lemma~\ref{is-eigenvalue-conjugate-frame-zak}, let $\mu$ be a Haar measure on $\mathbb W$ and $\mathcal F$ a fundamental
domain of $M_{r\times n}(E_h)$ in $\mathbb W$ with $\mu(\mathcal F)=1$, and for theta functions $f,f'$ put
$\langle f,f'\rangle=\int_{\mathcal F}f\bar f'\,d\mu$. Let $\omega$ be a homomorphism from $U(S)(\mathbf A)\times U(T)(\mathbf A)$
to the linear endomorphisms of the space of functions on $\mathbb W$ with the four properties of
Lemma~\ref{is-eigenvalue-conjugate-frame-real-character}; of these only the preservation of theta functions and of the
$L^2$-norms of theta functions on fundamental domains is used.
\begin{enumerate}
\item Such $\mu$ and $\mathcal F$ exist. For theta functions $f,f'$ and $a\in\mathbf C$, $\langle af,f'\rangle=a\langle
  f,f'\rangle$, $\langle f,af'\rangle=\bar a\langle f,f'\rangle$, $\langle f,f\rangle=0$ only if $f=0$, and
  $\langle\omega(g,h)f,\omega(g,h)f'\rangle=\langle f,f'\rangle$ for all $(g,h)$.
\item Let $\delta$, $A_\pm$ and $K$ be as in Definition~\ref{is-eigenvalue-conjugate-frame-overlap-doubling} for $X'$,
  its symplectic basis and indexing, and $n$. For $x'_1,x'_2\in\mathcal D_{X'}$ and Schwartz-Bruhat functions
  $\varphi_1,\varphi_2$ on $M_{n\times d}(\mathbf Q)$,
  \[
   \langle\mathrm{Th}(x'_1,\varphi_1),\mathrm{Th}(x'_2,\varphi_2)\rangle=\delta(\varphi_1,\varphi_2)(0)\,K(x'_1,x'_2),
  \]
  and $K(x'_1,x'_2)=\prod_{\sigma\in\Psi'}\big(\det A_+(x'_1,x'_2)_\sigma\det A_-(x'_1,x'_2)_\sigma\big)^{-n/2}$.
\item Let $B$ and $B'$ be families of complex $r\times r$ matrices indexed by $\mathcal T$, with $B$ orthonormal for
  $X'$, and $u\in U(T)(\mathbf A)$; here $x'_{B'}$ has the component $\overline{B'_\rho}$ at $\bar\rho$ and
  $uB'=(u_\rho B'_\rho)_\rho$. Then $K(x'_{uB},x'_{uB'})=K(x'_B,x'_{B'})$.
\item Let $\varphi\ne0$ be a Schwartz-Bruhat function, $B$ and $B'$ orthonormal for $X'$, and $u\in U(T)(\mathbf A)$. Then
  $\langle\mathrm{Th}(x'_B,\varphi),\mathrm{Th}(x'_{B'},\varphi)\rangle\ne0$ and
  $\langle\mathrm{Th}(x'_{uB},\varphi),\mathrm{Th}(x'_{uB'},\varphi)\rangle=\langle\mathrm{Th}(x'_B,\varphi),
  \mathrm{Th}(x'_{B'},\varphi)\rangle$.
\end{enumerate}
\end{lemma}
\leanref{Development/AbelianVarietyTheory/Proofs/SiegelThetaTransformation/ThetaFunctionFundamentalDomainInnerProductBasic.lean}{GlobalTheory.Ker.ScalarProduct.ip_basic}
\leanref{Development/AbelianVarietyTheory/Proofs/DivisorProofs/InnerThetaFunctionsReadingTestFunctionFrames.lean}{GlobalTheory.AdTh.Overlap.ipO_Th}
\leanref{Development/AbelianVarietyTheory/Proofs/DivisorProofs/InnerThetaFunctionsReadingTestFunctionFrames.lean}{GlobalTheory.AdTh.Overlap.kZero_zero}
\leanref{Development/AbelianVarietyTheory/Proofs/DivisorProofs/InnerThetaFunctionsReadingTestFunctionFrames.lean}{GlobalTheory.AdTh.Overlap.kZero_mvB}
\leanref{Development/AbelianVarietyTheory/Proofs/SiegelThetaTransformation/InnerProductTranslatedVacuaEqGaussianKernel.lean}{GlobalTheory.Ker.GaussKer.ipO_rho_Th}
\begin{proof}
(1) The group $M_{r\times n}(E_h)$ is discrete and cocompact in $\mathbb W$ \cite{Wei67}, so it has a fundamental domain
of finite positive measure, and $\mu$ can be scaled. The two scalar rules are clear. If $\langle f,f\rangle=0$, then $f=0$
almost everywhere on $\mathcal F$, so almost everywhere on $\mathbb W$, as $|f|$ is invariant under the rational points
by quasi-periodicity; so $f=0$ by continuity. The invariance under $\omega(g,h)$ follows from the preservation of the
norms by polarization, as $\omega(g,h)$ is linear and $4\langle f,f'\rangle=\sum_{k=0}^3i^k\langle f+i^kf',f+i^kf'\rangle$.

(2) Let $x'_1,x'_2,\varphi_1,\varphi_2$ be given, and put $\mathrm{Th}_i=\mathrm{Th}(x'_i,\varphi_i)$. By
Lemma~\ref{is-eigenvalue-conjugate-frame-zak}(1), $\mathrm{Th}_i(w)=\Theta'_{\varphi_i}(x'_i;u'(w))$ for every
$w\in\mathbb W$ with finite component $0$. Let $N_0$ be as in Lemma~\ref{is-eigenvalue-conjugate-frame-torus-mean} for
$f=\mathrm{Th}_1$ and $f'=\mathrm{Th}_2$, and $N_1$ as in Lemma~\ref{is-eigenvalue-conjugate-frame-overlap-mean} for
$X'$, $x'_1$, $x'_2$, $\varphi_1$ and $\varphi_2$, as $n$ is even; put $N=N_0N_1$. The group $\Gamma'_N$ of
Lemma~\ref{is-eigenvalue-conjugate-frame-torus-mean} is the group $\Gamma_N$ of
Lemma~\ref{is-eigenvalue-conjugate-frame-overlap-mean} for $X'$, both being generated by the $N\hat y$ with the rows of
$y$ in $\Lambda'$, and it has a fundamental domain $F_N$. Lemma~\ref{is-eigenvalue-conjugate-frame-torus-mean}, with
$\Phi(u')=\Theta'_{\varphi_1}(x'_1;u')\overline{\Theta'_{\varphi_2}(x'_2;u')}$, and
Lemma~\ref{is-eigenvalue-conjugate-frame-overlap-mean}, with $a=b=0$, give
\[
 \langle\mathrm{Th}_1,\mathrm{Th}_2\rangle=\frac1{\operatorname{vol}(F_N)}\int_{F_N}\Phi(u')\,du'
 =\delta(\varphi_1,\varphi_2)(0)\,\mathcal G'_{x'_1}(0)\,\overline{\mathcal G'_{x'_2}(0)}\,K(x'_1,x'_2;0,0),
\]
as $E_{X'}(w,0)=0$ and the coordinates of $0$ are $0$. Here $\mathcal G'_{x'_i}(0)=1$, and $K(x'_1,x'_2;0,0)=K(x'_1,x'_2)$.
At $z=w=0$ the vectors $v_k$ and $v'_k$ of Definition~\ref{is-eigenvalue-conjugate-frame-overlap-doubling}(2) are $0$,
so $\zeta=\zeta'=0$, the exponential factor of the kernel is $1$, and $K(x'_1,x'_2)$ is the product of the determinant
factors, with $\lfloor n/2\rfloor=n/2$.

(3) Let $\rho\in\mathcal T$, $\sigma=\bar\rho$, $U_\rho=\theta_\rho(u)$ and $M=\overline{U_\rho}$. As in the proof of
Lemma~\ref{is-eigenvalue-conjugate-frame-zak}(4), the components of $x'_{uB}$ and $x'_{uB'}$ at $\sigma$ are
$Mx'_{B,\sigma}$ and $Mx'_{B',\sigma}$, and $U_\rho^*G_T(\rho)U_\rho=G_T(\rho)$. For $v,w\in\mathbf C^r$,
$h'_\sigma(v,w)=w^*G_T(\rho)^{\mathsf T}v$, as $G_T(\rho)$ is Hermitian; so
$h'_\sigma(Mv,Mw)=w^*\big(U_\rho^*G_T(\rho)U_\rho\big)^{\mathsf T}v=h'_\sigma(v,w)$. As $B_\rho^*G_T(\rho)B_\rho=J'$, $\det
G_T(\rho)\ne0$, and the determinant of $U_\rho^*G_T(\rho)U_\rho=G_T(\rho)$ gives $|\det M|=1$.

Let $x_1,x_2\in M_r(\mathbf C)^{\Psi'}$ and $y_1,y_2$ with $y_{i,\sigma}=Mx_{i,\sigma}$. Then
$A_+(y_1,y_2)_\sigma M=A_+(x_1,x_2)_\sigma$: for $v\in\mathbf C^r$ and $a\le p'_0$, by
Definition~\ref{is-eigenvalue-conjugate-frame-overlap-doubling}(1),
$(A_+(y_1,y_2)_\sigma Mv)_a=h'_\sigma(Mv,Mx_{2,\sigma}e_a)=h'_\sigma(v,x_{2,\sigma}e_a)=(A_+(x_1,x_2)_\sigma v)_a$, and for
$a>p'_0$, $(A_+(y_1,y_2)_\sigma Mv)_a=(x_{1,\sigma}^{-1}M^{-1}Mv)_a=(A_+(x_1,x_2)_\sigma v)_a$. Entry by entry, as in the proof
of Lemma~\ref{is-eigenvalue-conjugate-frame-overlap-determinants}, $A_-(x_1,x_2)_\sigma=\overline{A_+(x_2,x_1)_\sigma}$ for all
$x_1,x_2$. Hence
\[
 \det A_+(y_1,y_2)_\sigma\det A_-(y_1,y_2)_\sigma=\frac{\det A_+(x_1,x_2)_\sigma\,\overline{\det A_+(x_2,x_1)_\sigma}}{|\det M|^2}
 =\det A_+(x_1,x_2)_\sigma\det A_-(x_1,x_2)_\sigma.
\]
For $x_1=x'_B$ and $x_2=x'_{B'}$ this shows that the factors at $\sigma$ of $K(x'_{uB},x'_{uB'})$ and $K(x'_B,x'_{B'})$ in (2)
agree.

(4) By Lemma~\ref{is-eigenvalue-conjugate-frame-zak}(3) and~(4), the families $uB$ and $uB'$ are orthonormal and
$x'_B,x'_{B'},x'_{uB},x'_{uB'}\in\mathcal D_{X'}$. By (2) and
Lemma~\ref{is-eigenvalue-conjugate-frame-overlap-determinants}, $K(x'_B,x'_{B'})\ne0$. By (2),
$\langle\mathrm{Th}(x'_B,\varphi),\mathrm{Th}(x'_B,\varphi)\rangle=\delta(\varphi,\varphi)(0)K(x'_B,x'_B)$, and the left side is not
$0$ by (1), as $\mathrm{Th}(x'_B,\varphi)\ne0$ (Lemma~\ref{is-eigenvalue-conjugate-frame-zak}(2)); so
$\delta(\varphi,\varphi)(0)\ne0$. Then (2) gives
$\langle\mathrm{Th}(x'_B,\varphi),\mathrm{Th}(x'_{B'},\varphi)\rangle=\delta(\varphi,\varphi)(0)K(x'_B,x'_{B'})\ne0$, and (2) and (3)
give
$\langle\mathrm{Th}(x'_{uB},\varphi),\mathrm{Th}(x'_{uB'},\varphi)\rangle=\delta(\varphi,\varphi)(0)K(x'_{uB},x'_{uB'})
=\delta(\varphi,\varphi)(0)K(x'_B,x'_{B'})=\langle\mathrm{Th}(x'_B,\varphi),\mathrm{Th}(x'_{B'},\varphi)\rangle$.
\end{proof}
% ip_basic (ThetaFunctionFundamentalDomainInnerProductBasic.lean) for (1), with AdelicMatrixLattice.lattice and
% om_unit of Ker.IsWeak. (2) is ipO_Th (InnerThetaFunctionsReadingTestFunctionFrames.lean): delta is SW.dbl, K is SW.kZero
% at z = w = 0 (kZero_zero); its proof is GaussKer.ipO_rho_Th at w1 = w2 = 0 (rho 0 = id, zOf x 0 = 0, gaussR x 0 = 1),
% whose proof is T.thetaInnerProduct_eq_realTorusMean (Lemma -torus-mean) and T.mean_realTorus_normalisedTheta_mul_conj
% (Lemma -overlap-mean) for N0 N1. (3) is kZero_mvB, for B with C.IsB B (used only for det gramAt ne 0) and any B': mat_xB,
% ONBasis.compAt_unitary, hermC_conj, aLin_move, det_move (aLin_swap). (4) is the part of
% FrameEquiv.RealChar.scalar_eq (ExistsRealCharacter.lean) before its last lines: hd (delta(phi, phi)(0) ne 0 through
% Th_ne_zero and ip_basic) and hk (kZero_zero, DblPoint.det_aLin_ne_zero, det_aAnti_ne_zero).

\begin{lemma}\label{is-eigenvalue-conjugate-frame-real-character-scalar-fock}
In Situation~\ref{is-eigenvalue-conjugate-frame-context}, with the notation of
Lemma~\ref{is-eigenvalue-conjugate-frame-zak} and $\partial_\rho(z)$ as in
Lemma~\ref{is-eigenvalue-conjugate-theta-closed}, let $\omega$ be a homomorphism from $U(S)(\mathbf A)\times U(T)(\mathbf A)$ to
the linear endomorphisms of the space of functions on $\mathbb W$ with the four properties of
Lemma~\ref{is-eigenvalue-conjugate-frame-real-character}. Let $g\in U(S)(\mathbf A)$ and $u\in U(T)(\mathbf A)$, with
components $g_\rho=\theta_\rho(g)$ and $u_\rho=\theta_\rho(u)$, and let $f$ be a theta function for $(S,T)$.
\begin{enumerate}
\item $\omega(g,u)\,\partial_\rho(z)f=\partial_\rho(u_\rho zg_\rho^{-1})\,\omega(g,u)f$ for every $\rho\in\mathcal T$ and every
  complex $r\times n$ matrix $z$.
\item Let $B$ be orthonormal for $X'$. If $f$ is annihilated by the annihilation operators for the Cartan involutions $1$ of
  $S$ and $s'_{B,\rho}$ of $T$ at every $\rho$, then $\omega(g,u)f$ is annihilated by those for $1$ and $s'_{uB,\rho}$.
\item Let $B$ be orthonormal for $X'$. For every polynomial $Q$ in the creation variables,
  $\omega(g,u)(f\otimes_{s,B}Q)=(\omega(g,u)f)\otimes_{s',B'}Q$ (Lemma~\ref{is-eigenvalue-conjugate-theta-facts}), where
  $s_S=1$, $s_T=s'_B$, $B_S=1$, $B_T=B$, and $s'_S=1$, $s'_T=s'_{uB}$, $B'_S=(g_\rho)_\rho$, $B'_T=uB$.
\end{enumerate}
\end{lemma}
\leanref{Development/AbelianVarietyTheory/Proofs/SiegelThetaTransformation/WeakWeilRepCarriesDerivativesVacuaFockVectorsStatement.lean}{GlobalTheory.T.weakWeilRep_carries_derivatives_vacua_fockVectors}
\begin{proof}
\openstep{The proof, which differentiates the identity $\omega(g,u)\rho(w)=\rho(uwg^{-1})\omega(g,u)$ on theta functions along
the real directions and deduces the transport of the annihilation and creation operators, is not written out.}
% T.weakWeilRep_carries_derivatives_vacua_fockVectors = Ker.WeakMove.om_move. Size: 1542 lines (lean_size of the
% statement).
\end{proof}

\begin{lemma}\label{is-eigenvalue-conjugate-frame-real-character-scalar-schur}
Let $X$, its symplectic basis and indexing be as in Definition~\ref{is-eigenvalue-conjugate-frame-heisenberg}, let $n\ge0$ be
even (the parity is not used), and let $A$ be a map on the functions on $M_{n\times d}(\mathbf Q)$ that carries
Schwartz-Bruhat functions to Schwartz-Bruhat functions, with $A(\varphi+\psi)=A(\varphi)+A(\psi)$, $A(a\varphi)=aA(\varphi)$ and
$A(\rho_X(y)\varphi)=\rho_X(y)A(\varphi)$ for all Schwartz-Bruhat $\varphi,\psi$, $a\in\mathbf C$ and $y\in M_{n\times r}(E)$. Then
there is $c\in\mathbf C$ with $A(\varphi)=c\varphi$ for every Schwartz-Bruhat function $\varphi$.
\end{lemma}
\leanref{Development/AbelianVarietyTheory/Proofs/SiegelThetaTransformation/SchurHeisG.lean}{GlobalTheory.Ker.schur_heisG}
\begin{proof}
\openstep{The proof, which writes $\langle A\varphi,\psi\rangle$ as a sesquilinear form invariant under the Heisenberg operators and
shows that such a form on Schwartz-Bruhat functions is a multiple of the inner product, is not written out.}
% Ker.schur_heisG (SchurHeisG.lean), through exists_eq_mul_ip. Size: 923 lines (lean_size).
\end{proof}

\begin{lemma}\label{is-eigenvalue-conjugate-frame-real-character-scalar-norm}
In Situation~\ref{is-eigenvalue-conjugate-frame-context}, with the notation of
Lemma~\ref{is-eigenvalue-conjugate-frame-zak}, let $\mu$, $\mathcal F$ and $\langle\ ,\ \rangle$ be as in
Lemma~\ref{is-eigenvalue-conjugate-frame-real-character-overlap}. Then there is $a\ne0$ with
$\langle\mathrm{Th}(x'_B,\varphi),\mathrm{Th}(x'_B,\varphi)\rangle=a\langle\varphi,\varphi\rangle$ for every $B$ orthonormal for $X'$ and
every Schwartz-Bruhat function $\varphi$, the inner product on the right that of
Definition~\ref{is-eigenvalue-conjugate-frame-heisenberg}(3).
\end{lemma}
\leanref{Development/AbelianVarietyTheory/Proofs/SiegelThetaTransformation/WeilActionTranslatesThetaByUnitScalar.lean}{GlobalTheory.Ker.RealScalar.exists_ipO_Th}
\leanref{Development/AbelianVarietyTheory/Proofs/SiegelThetaTransformation/KZeroOnFrConstEntire.lean}{GlobalTheory.Ker.kZero_onFr_const_entire}
\leanref{Development/AbelianVarietyTheory/Proofs/SiegelThetaTransformation/DoubledHermitianTestFunctionZeroAndNorm.lean}{GlobalTheory.SW.dbl_zero_and_norm}
\begin{proof}
Let $\delta$, $A_\pm$ and $K$ be as in Lemma~\ref{is-eigenvalue-conjugate-frame-real-character-overlap}(2). For $B$ orthonormal,
$x'_B\in\mathcal D_{X'}$ (Lemma~\ref{is-eigenvalue-conjugate-frame-zak}(3)), and (2) of that lemma gives
$\langle\mathrm{Th}(x'_B,\varphi),\mathrm{Th}(x'_B,\varphi)\rangle=\delta(\varphi,\varphi)(0)K(x'_B,x'_B)$.

\emph{The kernel at orthonormal points.} Let $x$ be orthonormal for $X'$ and $\sigma\in\Psi'$, so that
$x_\sigma^{\mathsf T}(G_T^{\mathsf T})^\sigma\bar x_\sigma=J'$. For $w\in\mathbf C^r$ and $a\le p'_0$,
$(A_+(x,x)_\sigma x_\sigma w)_a=h'_\sigma(x_\sigma w,x_\sigma e_a)=(J'w)_a=w_a$, and for $a>p'_0$, $(A_+(x,x)_\sigma x_\sigma w)_a=w_a$; so
$A_+(x,x)_\sigma=x_\sigma^{-1}$. By the entrywise identity of Lemma~\ref{is-eigenvalue-conjugate-frame-overlap-determinants},
$A_-(x,x)_\sigma=\overline{A_+(x,x)_\sigma}$. So $\det A_+(x,x)_\sigma\det A_-(x,x)_\sigma=|\det x_\sigma|^{-2}=|\det G_T^\sigma|$, as
$|\det x_\sigma|^2|\det G_T^\sigma|=|\det J'|=1$. Hence $K(x,x)=\kappa$ with $\kappa=\prod_{\sigma\in\Psi'}|\det G_T^\sigma|^{-n/2}>0$,
for every orthonormal $x$, in particular for $x'_B$.

\emph{The doubling at $0$.} There is $a'\ne0$ with $\delta(\varphi_1,\varphi_2)(0)=a'\langle\varphi_1,\varphi_2\rangle$ for all
Schwartz-Bruhat $\varphi_1,\varphi_2$.
\openstep{The proof that the value at $0$ of the doubling is a fixed nonzero multiple of the inner product, which uses that both
are sesquilinear forms invariant under the Heisenberg operators, is not written out.}
% The first clause of SW.dbl_zero_and_norm (DoubledHermitianTestFunctionZeroAndNorm.lean), which also gives b > 0 with
% ip (dbl phi1 phi2) (dbl phi1 phi2) = b ip phi1 phi1 ip phi2 phi2; through T.doubledHermitianTestFunction_spec,
% T.doubledThetaTestFunction_sesquilinear_heisenberg and heisG_laws. Size: 3107 lines (lean_size).

Hence $a=a'\kappa$ works.
\end{proof}
% Ker.RealScalar.exists_ipO_Th (WeilActionTranslatesThetaByUnitScalar.lean), followed: dbl_zero_and_norm (a), the first
% clause of kZero_onFr_const_entire (kappa = prod_tau norm (det Gm^tau)^-1 ^ (n/2) on orthonormal frames, whose proof the
% book gives), T.innerProduct_translated_vacua_eq_gaussianKernel at the points 0 (Lemma -real-character-overlap(2)).

\begin{lemma}\label{is-eigenvalue-conjugate-frame-real-character-scalar}
In Situation~\ref{is-eigenvalue-conjugate-frame-context}, with the notation of
Lemma~\ref{is-eigenvalue-conjugate-frame-zak}, let $\omega$ be a homomorphism from $U(S)(\mathbf A)\times U(T)(\mathbf A)$ to the
linear endomorphisms of the space of functions on $\mathbb W$ with the four properties of
Lemma~\ref{is-eigenvalue-conjugate-frame-real-character}. Let $g\in U(S)(\mathbf A)$ and $u\in U(T)(\mathbf A)$ have finite
components $1$, and let $B$ be orthonormal for $X'$. Then there is $c\in\mathbf C$ with $|c|=1$ and
$\omega(g,u)\,\mathrm{Th}(x'_B,\varphi)=c\,\mathrm{Th}(x'_{uB},\varphi)$ for every Schwartz-Bruhat function $\varphi$.
\end{lemma}
\leanref{Development/AbelianVarietyTheory/Proofs/SiegelThetaTransformation/WeilActionTranslatesThetaByUnitScalar.lean}{GlobalTheory.Ker.RealScalar.om_real_Th}
\begin{proof}
Write $\omega=\omega(g,u)$. By Lemma~\ref{is-eigenvalue-conjugate-frame-zak}(3) and~(4), $uB$ is orthonormal and $x'_B$,
$x'_{uB}$ lie in $\mathcal D_{X'}$.

\emph{The map $A$.} Let $\varphi$ be Schwartz-Bruhat. By Lemma~\ref{is-eigenvalue-conjugate-frame-zak}(5), $\mathrm{Th}(x'_B,\varphi)$ is
a theta function annihilated by the annihilation operators for $1$ and $s'_B$. So $\omega\,\mathrm{Th}(x'_B,\varphi)$ is a theta
function, by the first property of $\omega$, annihilated by those for $1$ and $s'_{uB}$, by
Lemma~\ref{is-eigenvalue-conjugate-frame-real-character-scalar-fock}(2). By
Lemma~\ref{is-eigenvalue-conjugate-frame-vacuum-line-fock}, applied to $uB$, it is $\mathrm{Th}(x'_{uB},\varphi')$ for a
Schwartz-Bruhat function $\varphi'$; put $A(\varphi)=\varphi'$, and $A(\varphi)=\varphi$ for $\varphi$ not Schwartz-Bruhat.

\emph{Its laws.} By Lemma~\ref{is-eigenvalue-conjugate-frame-zak}(2), $\mathrm{Th}(x',\cdot)$ is additive, homogeneous and
injective on Schwartz-Bruhat functions, for $x'=x'_B$ and $x'=x'_{uB}$. For Schwartz-Bruhat $\varphi,\psi$, as $\omega$ is linear,
\[
 \mathrm{Th}(x'_{uB},A(\varphi+\psi))=\omega\,\mathrm{Th}(x'_B,\varphi)+\omega\,\mathrm{Th}(x'_B,\psi)
 =\mathrm{Th}(x'_{uB},A(\varphi)+A(\psi)),
\]
so $A(\varphi+\psi)=A(\varphi)+A(\psi)$; in the same way $A(a\varphi)=aA(\varphi)$. Let $y\in M_{n\times r}(E_h)$; then
$\rho_{X'}(y)\varphi$ is Schwartz-Bruhat (Lemma~\ref{is-eigenvalue-conjugate-frame-vacua-heisenberg}(2)). Let $(-y)_f\in\mathbb W$
have archimedean component $0$ and finite component $-y^{\mathsf T}$. As $g$ and $u$ have finite components $1$,
$u(-y)_fg^{-1}=(-y)_f$. By Lemma~\ref{is-eigenvalue-conjugate-frame-zak}(2) for $-y$ and the second property of $\omega$,
\[
 \mathrm{Th}(x'_{uB},A(\rho_{X'}(y)\varphi))=\omega\rho((-y)_f)\mathrm{Th}(x'_B,\varphi)=\rho((-y)_f)\,\omega\,\mathrm{Th}(x'_B,\varphi)
 =\mathrm{Th}(x'_{uB},\rho_{X'}(y)A(\varphi)),
\]
so $A(\rho_{X'}(y)\varphi)=\rho_{X'}(y)A(\varphi)$.

\emph{The scalar.} By Lemma~\ref{is-eigenvalue-conjugate-frame-real-character-scalar-schur} for $X'$, there is $c$ with $A(\varphi)=c\varphi$
for every Schwartz-Bruhat $\varphi$; so $\omega\,\mathrm{Th}(x'_B,\varphi)=\mathrm{Th}(x'_{uB},c\varphi)=c\,\mathrm{Th}(x'_{uB},\varphi)$. If
every Schwartz-Bruhat function is $0$, then $c=1$ also satisfies this, as $\mathrm{Th}(x',0)=0$. Otherwise let $\varphi\ne0$ be
Schwartz-Bruhat, $\mu$, $\mathcal F$ and $\langle\ ,\ \rangle$ as in Lemma~\ref{is-eigenvalue-conjugate-frame-real-character-overlap}, and
$a$ as in Lemma~\ref{is-eigenvalue-conjugate-frame-real-character-scalar-norm}. By that lemma for $B$ and for $uB$, and by the
invariance and the scalar rules of Lemma~\ref{is-eigenvalue-conjugate-frame-real-character-overlap}(1),
\[
 a\langle\varphi,\varphi\rangle=\langle\mathrm{Th}(x'_B,\varphi),\mathrm{Th}(x'_B,\varphi)\rangle
 =\langle\omega\,\mathrm{Th}(x'_B,\varphi),\omega\,\mathrm{Th}(x'_B,\varphi)\rangle=|c|^2\,a\langle\varphi,\varphi\rangle.
\]
As $a\ne0$ and $\langle\varphi,\varphi\rangle>0$ for $\varphi\ne0$ (Definition~\ref{is-eigenvalue-conjugate-frame-heisenberg}(3)), $|c|=1$.
\end{proof}
% Ker.RealScalar.om_real_Th (WeilActionTranslatesThetaByUnitScalar.lean), followed: onBasis_frame_period (Lemma -zak(3),
% (4)), T.adelicZakTheta_exists_unique, T.exists_testFunction_of_fockVacuum_thetaFunction (Lemma -vacuum-line-fock) with
% T.weakWeilRep_carries_derivatives_vacua_fockVectors (Lemma -real-character-scalar-fock) and
% T.fockPolynomialVector_testTheta_zero_eq_thetaSeries (Lemma -zak(5)); T.thetaKernel_testFunction_map_linear_injective_equivariant
% (Lemma -zak(2)), isSB_algebra_haar_linear, heisG_laws, om_heis and actO_finW; Ker.schur_heisG (Lemma
% -real-character-scalar-schur); for |c| = 1, T.thetaFunction_fundamentalDomain_innerProduct_basic (Lemma
% -real-character-overlap(1)), exists_ipO_Th (Lemma -real-character-scalar-norm) and SW.DblZero.ip_self_ne_zero.

\begin{lemma}\label{is-eigenvalue-conjugate-frame-real-character}
In Situation~\ref{is-eigenvalue-conjugate-frame-context}, with the notation of
Lemma~\ref{is-eigenvalue-conjugate-frame-zak}, let $\omega$ be a homomorphism from $U(S)(\mathbf A)\times
U(T)(\mathbf A)$ to the linear endomorphisms of the space of functions on $\mathbb W$ such that: $\omega(g,h)$
preserves theta functions; on theta functions $\omega(g,h)\rho(w_0)=\rho(hw_0g^{-1})\omega(g,h)$;
$(\omega(\gamma,\gamma')f)(0)=f(0)$ for $\gamma\in U(S)(F_h)$, $\gamma'\in U(T)(F_h)$ and theta functions $f$;
and $\omega(g,h)$ preserves the $L^2$-norm of theta functions on every fundamental domain of
$M_{r\times n}(E_h)$ in $\mathbb W$ for every Haar measure. A Weil datum for $(S,T)$ is such a homomorphism
(Definition~\ref{th-weil-datum}(3)). There is a function $\chi$ on $U(T)(\mathbf A)$ with $|\chi(u)|=1$ and $\chi(uu')=\chi(u)\chi(u')$ for all
$u,u'$ with finite component $1$, such that
\[
 \omega(1,u)\,\mathrm{Th}(x'_B,\varphi)=\chi(u)\,\mathrm{Th}(x'_{uB},\varphi)
\]
for every $u$ with finite component $1$, every family $B$ orthonormal for $X'$ and every Schwartz-Bruhat
function $\varphi$.
\end{lemma}
\leanref{Development/AbelianVarietyTheory/Proofs/DivisorProofs/ExistsRealCharacter.lean}{GlobalTheory.FrameEquiv.exists_real_character}
\leanref{Development/AbelianVarietyTheory/Definitions/SpanningFamilies/IsWeakOfDatum.lean}{GlobalTheory.Ker.IsWeak}
\begin{proof}
If every Schwartz-Bruhat function is $0$, $\chi=1$ works, as $\mathrm{Th}(x',0)=0$
(Lemma~\ref{is-eigenvalue-conjugate-frame-zak}(2)). Otherwise fix a Schwartz-Bruhat function $\varphi_0\ne0$ and a family $B_0$
orthonormal for $X'$ (Lemma~\ref{is-eigenvalue-conjugate-frame-zak}(3)). For $u$ with finite component $1$ let $\chi(u)$ be a scalar
$c$ of Lemma~\ref{is-eigenvalue-conjugate-frame-real-character-scalar} for $g=1$, $u$ and $B_0$; so $|\chi(u)|=1$.

\emph{Independence of the family.} Let $u$ have finite component $1$ and $B$ be orthonormal for $X'$, and let $c$ be as in
Lemma~\ref{is-eigenvalue-conjugate-frame-real-character-scalar} for $1$, $u$ and $B$. We show $c=\chi(u)$. Let $\mu$, $\mathcal F$
and $\langle\ ,\ \rangle$ be as in Lemma~\ref{is-eigenvalue-conjugate-frame-real-character-overlap}; the families $uB$ and $uB_0$ are
orthonormal (Lemma~\ref{is-eigenvalue-conjugate-frame-zak}(4)). Write $\mathrm{Th}_A=\mathrm{Th}(x'_A,\varphi_0)$. By (1) of that
lemma,
\[
 \langle\mathrm{Th}_B,\mathrm{Th}_{B_0}\rangle=\langle\omega(1,u)\mathrm{Th}_B,\omega(1,u)\mathrm{Th}_{B_0}\rangle
 =c\,\overline{\chi(u)}\,\langle\mathrm{Th}_{uB},\mathrm{Th}_{uB_0}\rangle,
\]
and by its (4) the inner products on the two sides are equal and not $0$. So $c\overline{\chi(u)}=1$, and as $|\chi(u)|=1$,
$c=\chi(u)$. Hence $\omega(1,u)\mathrm{Th}(x'_B,\varphi)=\chi(u)\mathrm{Th}(x'_{uB},\varphi)$ for every Schwartz-Bruhat $\varphi$.

\emph{Multiplicativity.} Let $u,u'$ have finite component $1$. As $\omega$ is a homomorphism, $(uu')B_0=u(u'B_0)$
(Lemma~\ref{is-eigenvalue-conjugate-frame-zak}(4)) and $u'B_0$ is orthonormal,
\[
 \chi(uu')\mathrm{Th}_{uu'B_0}=\omega(1,u)\omega(1,u')\mathrm{Th}_{B_0}=\chi(u')\,\omega(1,u)\mathrm{Th}_{u'B_0}
 =\chi(u')\chi(u)\,\mathrm{Th}_{uu'B_0}.
\]
As $\mathrm{Th}_{uu'B_0}\ne0$ (Lemma~\ref{is-eigenvalue-conjugate-frame-zak}(2)), $\chi(uu')=\chi(u)\chi(u')$.
\end{proof}
% exists_real_character (29 lines; ExistsRealCharacter.lean, 104 lines), followed: the case without nonzero phi0 (Th_zero),
% onBasis_frame_period for B0, om_real_Th at h = 1 (Lemma -real-character-scalar), scalar_eq (Lemma -real-character-overlap,
% ip_basic, ipO_Th, kZero_mvB), and the multiplicativity through mvB_mul and Th_ne_zero.
% Ker.RealScalar.om_real_Th (71 lines; its
% file 183 lines) for the scalar at one frame B0, scalar_eq for its independence of B through the overlap integrals
% AdTh.Overlap.ipO_Th (190 lines in its file), kZero and kZero_mvB, and the multiplicativity from mvB_mul. The
% hypothesis is Ker.IsWeak C.psi S T om, whose fields om_pres, om_heis, om_aut and om_unit are the four conditions.

\begin{lemma}\label{is-eigenvalue-conjugate-frame-real-move}
In Situation~\ref{is-eigenvalue-conjugate-frame-situation}, let $\omega$ be a Weil datum for $(S,T)$ with
characters $(\chi_S,\chi_T)$, and suppose that the quadratic character $\eta$ of $E_h/F_h$ is the
restriction of a Hecke character of $E_h$ that is unramified at every finite place and has archimedean
exponents $\frac12$. Let $Z$ be a signature family, $a\in U(T)(\mathbf A)$ an element with finite
component $1$, $y$ a determinant root of $a$ (Lemma~\ref{is-eigenvalue-conjugate-frame-det-roots}),
$\varphi$ a Schwartz-Bruhat function on $M_{n\times d}(\mathbf Q)$
(Definition~\ref{is-eigenvalue-conjugate-frame-fibres}(6)), and $f$ a theta function for $(S,T)$ with
$f(w)=\Theta_\varphi(x^Z;u(w))$ for every $w\in\mathbb W$ with finite component $0$. Then
$\omega(1,a)f$ is a theta function, and
\[
 (\omega(1,a)f)(w)=\chi_S(y)\,\Theta_\varphi(x^{aZ};u(w))
\]
for every $w\in\mathbb W$ with finite component $0$.
\end{lemma}
\leanref{Development/AbelianVarietyTheory/Proofs/IwasawaDecomposition/RealPointMovesAdelicThetaByDetCharacterStatement.lean}{GlobalTheory.T.real_point_moves_adelic_theta_by_det_character}
\leanref{Development/AbelianVarietyTheory/Proofs/IwasawaDecomposition/RealPointMovesAdelicThetaByDetCharacter.lean}{GlobalTheory.OneMapRef.RealMove.zak_om_real}
\begin{proof}
Let $X'$, its symplectic basis and indexing, and $\pi$ be as in
Lemma~\ref{is-eigenvalue-conjugate-frame-conjugate-model}; with them the data are as in
Situation~\ref{is-eigenvalue-conjugate-frame-context}, whose notation and that of
Lemma~\ref{is-eigenvalue-conjugate-frame-zak} we use. Put $B=Z^\pi$. By (2) and (3) of
Lemma~\ref{is-eigenvalue-conjugate-frame-conjugate-model}, $B$ is orthonormal for $X'$, $s'_B=s_Z$, and $f$ is a
Zak theta function for $x'_B$ and $\varphi$. By Lemma~\ref{is-eigenvalue-conjugate-frame-zak}(3) and~(1),
$x'_B\in\mathcal D_{X'}$ and $f=\mathrm{Th}(x'_B,\varphi)$. The family $aZ$ is a signature family
(Lemma~\ref{is-eigenvalue-conjugate-frame-signature}(3)), and $aB=(aZ)^\pi$; so $aB$ is orthonormal for $X'$
and $x'_{aB}\in\mathcal D_{X'}$.

We show $\omega(1,a)\mathrm{Th}(x'_B,\varphi)=\chi_S(y)\,\mathrm{Th}(x'_{aB},\varphi)$. If $\varphi=0$, both
sides are $0$, as $\mathrm{Th}(x',0)=0$ by Lemma~\ref{is-eigenvalue-conjugate-frame-zak}(2). Let
$\varphi\ne0$, and let $\chi$ be as in Lemma~\ref{is-eigenvalue-conjugate-frame-real-character}. By
Lemma~\ref{is-eigenvalue-conjugate-frame-root-values}(1) choose a determinant root $y(g)$ of every $g\in
U(T)(\mathbf A)$. The character $\chi_S$ restricts to $\eta^n$ on $\mathbf A_{F_h}^\times$, as $\omega$ is a
Weil datum, and $n$ is even; so $\chi_S$ takes the value $\chi_S(y(g))$ at every determinant root of $g$
(Lemma~\ref{is-eigenvalue-conjugate-frame-root-values}(3)). As $y(u)y(u')$ is a determinant root of $uu'$
(Lemma~\ref{is-eigenvalue-conjugate-frame-det-roots}(1)), the function $c(u)=\chi(u)/\chi_S(y(u))$ satisfies
$c(uu')=c(u)c(u')$ for $u,u'$ with finite component $1$. Let $k$ have finite component $1$, with $k_\rho
s'_{B,\rho}=s'_{B,\rho}k_\rho$ for every $\rho$. By Lemma~\ref{is-eigenvalue-conjugate-frame-zak}(5) and~(3),
$\mathrm{Th}(x'_B,\varphi)$ is a theta function annihilated by the annihilation operators for $1$ and
$s'_B$, and the $s'_{B,\rho}$ are Cartan involutions. A Weil datum satisfies the conditions of
Lemma~\ref{is-eigenvalue-conjugate-frame-real-character}. So
Lemmas~\ref{is-eigenvalue-conjugate-frame-compact-vacuum}, \ref{is-eigenvalue-conjugate-frame-real-character}
and~\ref{is-eigenvalue-conjugate-frame-compact-move} give
\[
 \chi(k)\,\mathrm{Th}(x'_{kB},\varphi)=\omega(1,k)\,\mathrm{Th}(x'_B,\varphi)
 =\chi_S(y(k))\,v_{s'_B}(k)\,\mathrm{Th}(x'_B,\varphi)=\chi_S(y(k))\,\mathrm{Th}(x'_{kB},\varphi).
\]
As $kB$ is orthonormal for $X'$ and $\varphi\ne0$, $\mathrm{Th}(x'_{kB},\varphi)\ne0$ by
Lemma~\ref{is-eigenvalue-conjugate-frame-zak}(2): otherwise
$\mathrm{Th}(x'_{kB},\varphi)=\mathrm{Th}(x'_{kB},0\cdot\varphi)$ would give $\varphi=0$. So $c(k)=1$. By
Lemma~\ref{is-eigenvalue-conjugate-frame-compact-character}, applied to the Cartan involutions $s'_B$,
$c(u)=1$ for every $u$ with finite component $1$; in particular $\chi(a)=\chi_S(y(a))=\chi_S(y)$, and
Lemma~\ref{is-eigenvalue-conjugate-frame-real-character} gives the claim.

By Lemma~\ref{is-eigenvalue-conjugate-frame-zak}(2),
$\omega(1,a)f=\chi_S(y)\,\mathrm{Th}(x'_{aB},\varphi)=\mathrm{Th}(x'_{aB},\chi_S(y)\varphi)$, a Zak theta function
for $x'_{(aZ)^\pi}$ and the Schwartz-Bruhat function $\chi_S(y)\varphi$. By
Lemma~\ref{is-eigenvalue-conjugate-frame-conjugate-model}(3), applied to the signature family $aZ$,
$\omega(1,a)f$ is a theta function with $(\omega(1,a)f)(w)=\Theta_{\chi_S(y)\varphi}(x^{aZ};u(w))=\chi_S(y)\,
\Theta_\varphi(x^{aZ};u(w))$ for every $w$ with finite component $0$.
\end{proof}
% zak_om_real: exists_conjugateModel_same_thetaFunctions with q0 = r - p0 (X', m, pi), ctxOf for C, hB from
% (H Z).1.1, em (aB = (aZ)^pi), hB' from onBasis_frame_period, hZ' from (H _).1.2, hd and hd' (xB_mem_dom), e1
% from isZak_Th (uniqueness) and (H Z).2, om_real, th_laws (Th (a phi) = a Th phi) and isZak_Th with (H _).2 for
% aZ. om_real: exists_real_character (chi, chimul, hchi); the case phi = 0 (Th_zero); isDetRoot_laws (rt, hmul);
% mulFun_eq_one_of_eq_one_on_maxCompact for chi / chiS o rt, multiplicative by root_eq and hmul, and 1 on the
% compact elements by om_compact (with the Fock clause of fvec_Th_zero), hchi, CompactMove.Th_mvB and Th_ne_zero;
% root_eq for chiS(rt a) = chiS(y).

\subsection{Moved vacua and level-one theta functions}\label{subsection-is-moved-vacua}

\begin{lemma}\label{is-eigenvalue-conjugate-frame-moved-vacuum}
Let $D$ be a definite-side datum with host field $E_h/F_h$, $\psi$ a character of $\mathbf A_{F_h}$ and
$\theta$ an archimedean family whose member at each $\tau\in\mathcal T$ extends $\tau$, and let $S$ and $T$
be Hermitian spaces of $D$ and $\omega$ a Weil datum for $(S,T)$ with characters $(\chi_S,\chi_T)$. Let
$s_S,s'_S$ and $s_T,s'_T$ be families of complex matrices of sizes $\operatorname{rk}S$ and
$\operatorname{rk}T$ indexed by $\mathcal T$, and $\Lambda_S,\Lambda'_S\subset E_h^{\operatorname{rk}S}$
and $\Lambda_T,\Lambda'_T\subset E_h^{\operatorname{rk}T}$ subsets. Let $(g,h)\in U(S)(\mathbf A)\times
U(T)(\mathbf A)$, with components $g_\tau$ and $h_\tau$ under $\theta_\tau$, satisfy
$s'_{S,\tau}g_\tau=g_\tau s_{S,\tau}$ and $s'_{T,\tau}h_\tau=h_\tau s_{T,\tau}$ for every $\tau$, and
suppose that the lattice of the pair given by $\Lambda'_S$ and $\Lambda'_T$ is the image under
$u\mapsto hug^{-1}$ of the lattice of the pair given by $\Lambda_S$ and $\Lambda_T$. If $f$ is a Fock vacuum
vector for $s_S$, $s_T$ and the lattice of the pair given by $\Lambda_S$ and $\Lambda_T$, then
$\omega(g,h)f$ is a Fock vacuum vector for $s'_S$, $s'_T$ and the lattice of the pair given by
$\Lambda'_S$ and $\Lambda'_T$. Here Fock vacuum vectors are defined as for Cartan involutions
(Chapter~\ref{chapter-unitary-theta}, after Definition~\ref{th-weil-datum}).
\end{lemma}
\leanref{Development/AbelianVarietyTheory/Proofs/AdelicLatticeProofs/SqNormIntegralFundamentalDomainEqL2NormSq.lean}{GlobalTheory.WeilVacMove.isVacO_om}
\begin{proof}
Put $f'=\omega(g,h)f$; it is a theta function, as $\omega(g,h)$ preserves theta functions. Let $u'$ lie
in the lattice of the pair given by $\Lambda'_S$ and $\Lambda'_T$, and write $u'=hug^{-1}$ with $u$ in the
lattice of the pair given by $\Lambda_S$ and $\Lambda_T$. On theta functions
$\omega(g,h)\rho(v)=\rho(hvg^{-1})\omega(g,h)$ (Definition~\ref{th-weil-datum}(3)); so
$\rho(2u')f'=\omega(g,h)\rho(2u)f=f'$. If $u'$ lies in the half of $\mathbb W$ attached to a split place,
so does $u=h^{-1}u'g$, as the half is stable under the action; then $\rho(u')f'=\omega(g,h)\rho(u)f=f'$. So
$f'$ is a lattice vacuum for the lattice of the pair given by $\Lambda'_S$ and $\Lambda'_T$. Every complex
matrix $z$ is $h_\tau z_0g_\tau^{-1}$ for $z_0=h_\tau^{-1}zg_\tau$. By $(*)$ in the proof of
Lemma~\ref{is-eigenvalue-conjugate-theta-facts}(4), $\omega(g,h)$ carries the annihilation operator
$\partial_\tau(z_0)+i\,\partial_\tau(i\,s_{T,\tau}z_0s_{S,\tau})$ applied to $f$ to
$\partial_\tau(z)+i\,\partial_\tau(h_\tau(i\,s_{T,\tau}z_0s_{S,\tau})g_\tau^{-1})$ applied to $f'$, and
$h_\tau s_{T,\tau}z_0s_{S,\tau}g_\tau^{-1}=s'_{T,\tau}zs'_{S,\tau}$. So $f'$ is annihilated by the
annihilation operators for $s'_S$ and $s'_T$. Finally $\omega(g,h)$ preserves the $L^2$-norm of theta
functions on fundamental domains, so $f'$ has norm $1$ on a fundamental domain of covolume $1$.
\end{proof}
% isVacO_om: hW.om_pres; actO_two_nsmul and hW.om_heis for the lattice; isHalf_actO for the halves; om_dRho
% and cplxO_move for the annihilation operators; sqNormIntegralFundamentalDomain_eq_l2Norm_sq and hW.om_unit for
% the norm.

\begin{lemma}\label{is-eigenvalue-conjugate-frame-move}
In Situation~\ref{is-eigenvalue-conjugate-frame-situation}, suppose that $S$ and $T$ have correct
signatures, let $\omega$ be a Weil datum for $(S,T)$ with characters $(\chi_S,\chi_T)$, and suppose that
the quadratic character $\eta$ of $E_h/F_h$ is the restriction of a Hecke character of $E_h$ that is
unramified at every finite place and has archimedean exponents $\frac12$, and that $\chi_S$ and $\chi_T$
are trivial on every idele that is a unit at every finite place and differs from $1$ only at the places
above one place of $F_h$ that does not split in $E_h$. Let $Z$ be a signature family, $\Lambda_S$ and
$\Lambda_T$ self-dual lattices of $S$ and $T$, and $f_0$ a Fock vacuum vector for the Cartan involutions
$1$ of $S$ and $s_Z$ of $T$ and the lattice of the pair given by $\Lambda_S$ and $\Lambda_T$. Let $\mu$ be a
function on $U(T)(\mathbf A)$ such that every element $g$ has a determinant root and $\mu(g)=\chi_S(y)$
for every determinant root $y$ of $g$ (Lemma~\ref{is-eigenvalue-conjugate-frame-det-roots}). Let
$g_0,g_1\in U(T)(\mathbf A)$ and a lattice $L$ of $T$ satisfy
$(g_0)_f\hat\Lambda_T=\hat L=(g_1)_f\hat\Lambda_T$, and let $\varphi$ be a Schwartz-Bruhat function on
$M_{n\times d}(\mathbf Q)$ such that $\omega(1,g_0)f_0$ is a theta function with
$(\omega(1,g_0)f_0)(w)=\Theta_\varphi(x^{g_0Z};u(w))$ for every $w\in\mathbb W$ with finite component
$0$. Then
\[
 (\omega(1,g_1)f_0)(w)=\mu(g_1)\mu(g_0)^{-1}\,\Theta_\varphi(x^{g_1Z};u(w))
\]
for every $w\in\mathbb W$ with finite component $0$.
\end{lemma}
\leanref{Development/AbelianVarietyTheory/Proofs/IwasawaDecomposition/OmegaFormScalarExtension.lean}{GlobalTheory.OneMapRef.FrameConst.om_vac_move}
\leanref{Development/AbelianVarietyTheory/Proofs/IwasawaDecomposition/OmegaFormScalarExtension.lean}{GlobalTheory.OneMapRef.FrameConst.isMov_quot}
\leanref{Formalism/LocalToGlobal/ProjectiveLattices/FamilyUniformisationBasePointOrthogonalBis.lean}{Formalism.L2.Pj.MuVal}
\begin{proof}
The value $\mu(g)$ is $\chi_S(y)$ for a determinant root $y$ of $g$, so $|\mu(g)|=1$. By
Lemma~\ref{is-eigenvalue-conjugate-frame-signature}(3), $g_0Z$ and $g_1Z$ are signature families, and by
part (4) of that lemma there is $a\in U(T)(\mathbf A)$ with finite component $1$ and $g_1Z=a(g_0Z)$. Put
$k=(ag_0)^{-1}g_1$, so that $g_1=ag_0k$. By Lemma~\ref{is-eigenvalue-conjugate-frame-signature}(3),
\[
 kZ=(ag_0)^{-1}(g_1Z)=(ag_0)^{-1}\big((ag_0)Z\big)=Z,
\]
so $k$ has archimedean component $1$. As $a$ has finite component $1$, $k_f=(g_0)_f^{-1}(g_1)_f$, so
$k_f\hat\Lambda_T=(g_0)_f^{-1}\hat L=\hat\Lambda_T$. By
Lemma~\ref{is-eigenvalue-conjugate-frame-det-roots}(2), $k$ has a determinant root $y$ with archimedean
component $1$ that is a unit at every finite place; $1$ is such a determinant root of $1\in U(S)(\mathbf
A)$. The matrix $1$ is a Cartan involution of $S$, as $S$ has Gram matrix $1$, and $s_Z$ is one of $T$
(Lemma~\ref{is-eigenvalue-conjugate-frame-signature}(2)). So
Lemma~\ref{is-eigenvalue-conjugate-frame-finite-vacuum}, applied to $(1,k)$, gives
$\omega(1,k)f_0=\chi_T(1)^{-1}\chi_S(y)f_0=\chi_S(y)f_0$.

Let $y_a$ and $y_0$ be determinant roots of $a$ and $g_0$. By
Lemma~\ref{is-eigenvalue-conjugate-frame-real-move}, applied to the signature family $g_0Z$, to $a$, $y_a$,
$\varphi$ and $f=\omega(1,g_0)f_0$, and as $a(g_0Z)=g_1Z$,
\[
 \big(\omega(1,a)\omega(1,g_0)f_0\big)(w)=\chi_S(y_a)\,\Theta_\varphi(x^{g_1Z};u(w))
\]
for every $w$ with finite component $0$. As $\omega$ is a homomorphism,
$\omega(1,g_1)f_0=\omega(1,a)\omega(1,g_0)\omega(1,k)f_0=\chi_S(y)\,\omega(1,a)\omega(1,g_0)f_0$. By
Lemma~\ref{is-eigenvalue-conjugate-frame-det-roots}(1), $y_ay_0y$ is a determinant root of $ag_0k=g_1$,
so $\mu(g_1)=\chi_S(y_a)\mu(g_0)\chi_S(y)$. Hence
\[
 (\omega(1,g_1)f_0)(w)=\chi_S(y)\chi_S(y_a)\,\Theta_\varphi(x^{g_1Z};u(w))
 =\mu(g_1)\mu(g_0)^{-1}\,\Theta_\varphi(x^{g_1Z};u(w)).\qedhere
\]
\end{proof}
% om_vac_move: oneg_frames (hc, hF), the last clause of oneg_frames at mvZ g0 Z0 for a and hae; hk from
% mvZ_mul, mvZ_fin and fin_of_mvZ; isMov_quot for hmov; LatticeUnitRoot.exists_finUnit_isDetRoot for y;
% T.weil_vacuum_finite_selfDualLattice_splitting with isCartanInv_one, finUnit_one and isDetRoot_one for h1;
% hmu.1 for ya and y0; T.real_point_moves_adelic_theta_by_det_character for h2; h3 from map_mul; h4 from
% detRootMul_holds; thR_smul.

\begin{lemma}\label{is-eigenvalue-conjugate-frame-vacua-mean}
Let $K$ be a number field, $a,b\ge0$ integers, $\mathcal M=M_{a\times b}(K)$ and $\mathcal M_0=M_{a\times b}(\mathcal
O_K)$. For an integer $R\ge1$, a function $\varphi\colon\mathcal M\to\mathbf C$ has \emph{level} $R$ if $\varphi(x)\ne0$
only when $Rx\in\mathcal M_0$, and $\varphi(x+Ry)=\varphi(x)$ for all $x\in\mathcal M$ and $y\in\mathcal M_0$; it is
\emph{Schwartz-Bruhat} if it has a level. The \emph{mean} of a Schwartz-Bruhat function $\varphi$ of level $R$ is
\[
 m(\varphi)=R^{-[K:\mathbf Q]ab}\sum_{q\in\mathcal M/R\mathcal M_0}\varphi(q),
\]
a sum with finitely many nonzero terms whose value does not depend on the level $R$. For $K=\mathbf Q$, $a=n$ and
$b=d$ these are the Schwartz-Bruhat functions of Definition~\ref{is-eigenvalue-conjugate-frame-fibres}(6), and
$\langle\varphi,\psi\rangle=m(\varphi\bar\psi)$ in Definition~\ref{is-eigenvalue-conjugate-frame-heisenberg}(3). Let
$\varphi$ and $\psi$ be Schwartz-Bruhat functions on $\mathcal M$.
\begin{enumerate}
\item The functions $\varphi+\psi$, $c\varphi$ for $c\in\mathbf C$, $\varphi\psi$ and $\bar\varphi$ are
  Schwartz-Bruhat.
\item $m(\varphi+\psi)=m(\varphi)+m(\psi)$, $m(c\varphi)=c\,m(\varphi)$ for $c\in\mathbf C$, and
  $m(\bar\varphi)=\overline{m(\varphi)}$.
\item For every $\beta\in\mathcal M$ the function $x\mapsto\varphi(x-\beta)$ is Schwartz-Bruhat, with mean $m(\varphi)$.
\end{enumerate}
\end{lemma}
\leanref{Development/AbelianVarietyTheory/Proofs/SiegelThetaTransformation/RationalSchwartzBruhatAlgebraIntegralLinearStatement.lean}{GlobalTheory.T.rationalSchwartzBruhat_algebra_integral_linear}
\leanref{Development/AbelianVarietyTheory/Definitions/SpanningFamilies/IMemSiegel.lean}{GlobalTheory.SW.haar}
\begin{proof}
If $\varphi$ and $\psi$ have levels $R_1$ and $R_2$, then both have the level $R=R_1R_2$: if $\varphi(x)\ne0$, then
$Rx=R_2(R_1x)\in\mathcal M_0$, and $\varphi(x+Ry)=\varphi(x+R_1(R_2y))=\varphi(x)$ for $y\in\mathcal M_0$.

(1) and (2) With the common level $R$, the four functions of (1) have level $R$, as the two conditions are kept by
sums, multiples, products and complex conjugates. The means of $\varphi+\psi$, $c\varphi$ and $\bar\varphi$, computed
with the level $R$, are sums over the same finite set of classes, which gives (2).

(3) Let $N\ge1$ be an integer with $N\beta\in\mathcal M_0$, and let $R$ be a level of $\varphi$. The translate
$\varphi_\beta(x)=\varphi(x-\beta)$ has level $RN$: if $\varphi_\beta(x)\ne0$, then $R(x-\beta)\in\mathcal M_0$ and
$RNx=N\cdot R(x-\beta)+R\cdot N\beta\in\mathcal M_0$; and $\varphi_\beta(x+RNy)=\varphi(x-\beta+R(Ny))=\varphi_\beta(x)$
for $y\in\mathcal M_0$. The function $\varphi$ also has level $RN$. Computed with the level $RN$, the mean of
$\varphi_\beta$ is $(RN)^{-[K:\mathbf Q]ab}\sum_q\varphi(q-\beta)$, and $q\mapsto q-\beta$ permutes the classes of
$\mathcal M/RN\mathcal M_0$; so it is $m(\varphi)$.
\end{proof}
% T.rationalSchwartzBruhat_algebra_integral_linear = SW.isSB_algebra_haar_linear, for any number field K and finite
% index types: exists_common_level for R, hasLevel_map2 for the four functions, haar_eq_haarAt (the mean computed
% at any level, through the finiteness of the support modulo R) for the means, and for (3)
% exists_natCast_smul_intM (N), hasLevel_translate (level R N), hasLevel_mul and finsum_translate. SW.haar is the
% mean (haarAt at a chosen level, and 0 for a function that is not Schwartz-Bruhat); SW.ip phi psi = haar (phi *
% conj psi).

\begin{lemma}\label{is-eigenvalue-conjugate-frame-vacua-heisenberg}
Let $X$, its symplectic basis and indexing be as in
Definition~\ref{is-eigenvalue-conjugate-frame-heisenberg}, let $n\ge0$ be even (the parity is not used), and let
$w_0\in E$. For
$y,y'\in M_{n\times r}(E)$ with rows $y_k$ and $y'_k$ put $\omega_X(y,y')=\sum_{k=1}^n\langle y_k,y'_k\rangle_X$, and
for $a,b,a',b'\in M_{n\times d}(\mathbf Q)$ put $[a,b]=\sum_ka_k\cdot b_k$, with rows $a_k$ and $b_k$.
\begin{enumerate}
\item For $v\in E^r$ with coordinates $(\alpha(v),\beta(v))$ (Definition~\ref{is-eigenvalue-conjugate-frame-heisenberg}),
  $v=\sum_l\alpha_l(v)b_l+\sum_l\beta_l(v)b'_l$; the coordinates are additive in $v$; and
  $\langle u,v\rangle_X=-\sum_l\big(\alpha_l(u)\beta_l(v)-\beta_l(u)\alpha_l(v)\big)$ for $u,v\in E^r$. The map
  $y\mapsto(\alpha,\beta)$ sending $y$ to its coordinates (Definition~\ref{is-eigenvalue-conjugate-frame-heisenberg}(1))
  is an additive bijection of $M_{n\times r}(E)$ onto $M_{n\times d}(\mathbf Q)^2$. If $(\alpha,\beta)$ and
  $(\alpha',\beta')$ are the coordinates of $y$ and $y'$, then $\omega_X(y,y')=-([\alpha,\beta']-[\alpha',\beta])$, and
  $\omega_X(y',y)=-\omega_X(y,y')$.
\item Let $\varphi$ be Schwartz-Bruhat. Then $\rho_X(y)\varphi$ is Schwartz-Bruhat, $\rho_X(0)\varphi=\varphi$,
  and $\rho_X(y)\rho_X(y')\varphi=e^{-\pi i\omega_X(y,y')}\rho_X(y+y')\varphi$.
\item For Schwartz-Bruhat $\varphi,\psi$, $\langle\rho_X(y)\varphi,\rho_X(y)\psi\rangle=\langle\varphi,\psi\rangle$.
\item $\varepsilon_{w_0}(0)=1$ and $|\varepsilon_{w_0}(y)|=1$.
\item Let $\varphi$ be Schwartz-Bruhat and $x\in\mathcal D_X$. For every $z=(z_1,\dots,z_n)\in(\mathbf C^d)^n$,
  \[
   \theta_{\rho_X(y)\varphi}(x;z)=\exp\Big(-\pi\sum_{k=1}^nH_x\big(z_k+\tfrac12p_x(y_k),p_x(y_k)\big)\Big)\,
   \theta_\varphi\big(x;z_1+p_x(y_1),\dots,z_n+p_x(y_n)\big),
  \]
  with the theta series and the form $H_x$ of Definition~\ref{is-eigenvalue-conjugate-frame-fibres}.
\end{enumerate}
\end{lemma}
\leanref{Development/AbelianVarietyTheory/Proofs/SiegelThetaTransformation/HMatMulVecPerCrd.lean}{GlobalTheory.SW.heisG_laws}
\leanref{Development/AbelianVarietyTheory/Proofs/SiegelThetaTransformation/TensorLatticeVacuaExistUnique.lean}{GlobalTheory.SW.TensVac.crdHom_bijective}
\leanref{Development/AbelianVarietyTheory/Proofs/SiegelThetaTransformation/TensorLatticeVacuaExistUnique.lean}{GlobalTheory.SW.TensVac.om_crdHom}
\leanref{Development/AbelianVarietyTheory/Proofs/SiegelThetaTransformation/TensorLatticeVacuaExistUnique.lean}{GlobalTheory.SW.TensVac.sumEF_swap}
\leanref{Development/AbelianVarietyTheory/Proofs/SiegelThetaTransformation/TensorLatticeVacuaExistUnique.lean}{GlobalTheory.SW.TensVac.epsW_zero}
\leanref{Development/AbelianVarietyTheory/Proofs/SiegelThetaTransformation/TensorLatticeVacuaExistUnique.lean}{GlobalTheory.SW.TensVac.epsW_mul_conj}
\begin{proof}
(1) The symplectic basis is a basis of $E^r$ over $\mathbf Q$, so the coordinates are additive, $v$ is the
combination of the basis with its coordinates, and every pair $(\alpha,\beta)$ comes from exactly one $y$. By the
values of $\langle\cdot,\cdot\rangle_X$ on the basis and the bilinearity over $\mathbf Q$, $\langle
u,v\rangle_X=-\sum_l\big(\alpha_l(u)\beta_l(v)-\beta_l(u)\alpha_l(v)\big)$; summing over the rows gives the formula
for $\omega_X$, and it is antisymmetric.

(2) Let $N$ be a level of $\varphi$, and let $N_2,N_1\ge1$ be integers with $N_2\beta$ and $N_1\alpha$ integral. By
Lemma~\ref{is-eigenvalue-conjugate-frame-vacua-mean}(3), $m\mapsto\varphi(m-\beta)$ has level $NN_2$. Its product with
$c(m)=e^{\pi i(2[m,\alpha]-[\alpha,\beta])}$, which is $\rho_X(y)\varphi$, has level $NN_2N_1$: it vanishes unless
$NN_2m$ is integral, and for an integral $m'$, translating $m$ by $NN_2N_1m'$ leaves $\varphi(m-\beta)$ unchanged and
multiplies $c(m)$ by $e^{2\pi iNN_2[m',N_1\alpha]}=1$. Clearly $\rho_X(0)\varphi=\varphi$. With $(\alpha',\beta')$ the
coordinates of $y'$,
\[
 \rho_X(y)\rho_X(y')\varphi(m)=\varphi(m-\beta-\beta')\,e^{\pi i(2[m,\alpha+\alpha']-[\alpha,\beta]-[\alpha',\beta']-2[\beta,\alpha'])},
\]
while $\rho_X(y+y')\varphi(m)=\varphi(m-\beta-\beta')\,e^{\pi i(2[m,\alpha+\alpha']-[\alpha+\alpha',\beta+\beta'])}$. The
exponents differ by $[\alpha,\beta']-[\alpha',\beta]=-\omega_X(y,y')$, by (1).

(3) The two exponentials in $\rho_X(y)\varphi(m)\overline{\rho_X(y)\psi(m)}$ cancel, as they have absolute value $1$,
so this product is $\varphi(m-\beta)\overline{\psi(m-\beta)}$. As $\langle\varphi,\psi\rangle$ is the mean of
$\varphi\bar\psi$, a Schwartz-Bruhat function, Lemma~\ref{is-eigenvalue-conjugate-frame-vacua-mean}(1) and~(3) give
$\langle\rho_X(y)\varphi,\rho_X(y)\psi\rangle=\langle\varphi,\psi\rangle$.

(4) $\varepsilon_{w_0}(y)$ is $e^{\pi it}$ with $t$ rational, and $t=0$ for $y=0$.

(5) Write $\Omega=\Omega_x$, $\Omega'=\Omega'_x$, $\tau=\tau_x$, $Y=\operatorname{Im}\tau$ and $\zeta_k=\Omega^{-1}z_k$. By
Lemma~\ref{is-eigenvalue-conjugate-frame-hermitian-riemann}(4), $\Omega$ is invertible and $\tau$ is symmetric; and
$H_x(z,w)=(\Omega^{-1}z)^{\mathsf T}Y^{-1}\overline{\Omega^{-1}w}$. The periods are additive and $p_x(av)=ap_x(v)$ for
$a\in\mathbf Q$ (Definition~\ref{is-eigenvalue-conjugate-frame-fibres}(1)), so by (1)
$p_x(y_k)=\Omega\alpha_k+\Omega'\beta_k$ and $l_k=\Omega^{-1}p_x(y_k)=\alpha_k+\tau\beta_k$, where $\alpha_k$ and $\beta_k$
are the rows of the coordinates of $y$. In the series of $\rho_X(y)\varphi$, substitute $m=m'+\beta$; as $\tau$ is
symmetric,
\[
 \sum_m\varphi(m-\beta)e^{\pi i\sum_k(2m_k\cdot\alpha_k-\alpha_k\cdot\beta_k+m_k^{\mathsf T}\tau m_k+2m_k^{\mathsf T}\zeta_k)}
 =e^{\pi i\sum_k(\beta_k^{\mathsf T}\tau\beta_k+2\beta_k^{\mathsf T}\zeta_k+\alpha_k\cdot\beta_k)}
 \sum_{m'}\varphi(m')e^{\pi i\sum_k(m_k'^{\mathsf T}\tau m'_k+2m_k'^{\mathsf T}(\zeta_k+l_k))}.
\]
The last sum is the sum in $\theta_\varphi(x;z+p_x(y))$, whose normalized coordinates are the $\zeta_k+l_k$. The
determinant factors of the two theta series are equal, so
$\theta_{\rho_X(y)\varphi}(x;z)=e^{\mathcal E}\,\theta_\varphi(x;z+p_x(y))$ with
\[
 \mathcal E=\pi i\sum_k(\beta_k^{\mathsf T}\tau\beta_k+2\beta_k^{\mathsf T}\zeta_k+\alpha_k\cdot\beta_k)
 +\frac\pi2\sum_k\big(\zeta_k^{\mathsf T}Y^{-1}\zeta_k-(\zeta_k+l_k)^{\mathsf T}Y^{-1}(\zeta_k+l_k)\big).
\]
As $l_k-\bar l_k=2iY\beta_k$, $Y^{-1}l_k=Y^{-1}\bar l_k+2i\beta_k$, and the second sum is
$-\pi\sum_k\big((\zeta_k+\frac12l_k)^{\mathsf T}Y^{-1}\bar l_k+2i\zeta_k^{\mathsf T}\beta_k+il_k^{\mathsf T}\beta_k\big)$, with
$l_k^{\mathsf T}\beta_k=\alpha_k\cdot\beta_k+\beta_k^{\mathsf T}\tau\beta_k$. So
$\mathcal E=-\pi\sum_k(\zeta_k+\frac12l_k)^{\mathsf T}Y^{-1}\bar l_k=-\pi\sum_kH_x(z_k+\frac12p_x(y_k),p_x(y_k))$.
\end{proof}
% heisG_laws: vecQ_crd, crd_add and EF_crd for (1), and for SB phi, with heisG y = heisS (crd of the rows,
% Sum.inl) (crd, Sum.inr), T.heisenbergOperatorRationalPoint_laws = SW.heisS_laws: h1 for (2) (isSB_mul_char, the
% level R D after the translate of level R N from Lemma -vacua-mean(3), as followed above), h3 and h4 for (2), h5 for
% (3), and h2 (the series computation, Equiv.addRight b) with hMat_mulVec_per_crd and heis_exponent for (5).
% crdHom_bijective (vecQ_crd, crd_vecQ), om_crdHom and sumEF_swap complete (1); epsW_zero and epsW_mul_conj give (4).
% The Lean's heisG_laws takes Even n, which its proof does not use. Clause (5) is also stated in Chapter 9 as Lemma
% cv-uniformization-theta-rule-translate, with the same leanref.

\begin{lemma}\label{is-eigenvalue-conjugate-frame-vacua-tensor}
Let $X$, its symplectic basis and indexing be as in
Definition~\ref{is-eigenvalue-conjugate-frame-heisenberg} (the basis, its being symplectic and the indexing
are not used), let $n\ge0$ be even (the parity is not used), let $w_0\in E$ be an algebraic integer with
$w_0+\bar w_0=1$, and let $L\subset E^r$ be an $\mathcal O_E$-submodule that is self-dual for $X$
(Definition~\ref{is-eigenvalue-conjugate-frame-fibres}(7)). For an integer $M\ge0$ let $\mathcal V_M$ be the set
of the finitely generated $\mathcal O_E$-submodules $\Lambda\subset E^n$ such that $x\in\Lambda$ exactly when
$\sum_ax_a\bar y_a$ is an algebraic integer for every $y\in\Lambda$, and such that some integer $t>0$ prime to
$M$ has $t\Lambda\subset\mathcal O_E^n$ and $t\mathcal O_E^n\subset\Lambda$ (this last condition is not used).
Let $M\ge0$ and $\Lambda_S\in\mathcal V_M$. Put $A=\Lambda_S\otimes L$ and
$\omega_X(y,y')=\sum_k\langle y_k,y'_k\rangle_X$. Then:
\begin{enumerate}
\item $A$ is finitely generated;
\item $y\in M_{n\times r}(E)$ lies in $A$ exactly when $\omega_X(y,y')\in\mathbf Z$ for every $y'\in A$;
\item for $y,y'\in A$, $\varepsilon_{w_0}(y)^2=1$ and
  $\varepsilon_{w_0}(y+y')=\varepsilon_{w_0}(y)\,\varepsilon_{w_0}(y')\,e^{-\pi i\omega_X(y,y')}$;
\item $y\in M_{n\times r}(E)$ lies in $\mathcal O_E^n\otimes L$ exactly when its rows $y_1,\dots,y_n$ lie in $L$;
\item for every $\mathcal O_E$-submodule $L'\subset E^r$ that is self-dual for $X$, the index of
  $A\cap(\Lambda_S\otimes L')$ in $A$ is $[L:L\cap L']^n$, and $A\cap(\Lambda_S\otimes L')=\Lambda_S\otimes(L\cap L')$;
\item for every $\Lambda'_S\in\mathcal V_M$, the index of $A\cap(\Lambda'_S\otimes L)$ in $A$ is
  $[\Lambda_S:\Lambda_S\cap\Lambda'_S]^r$.
\end{enumerate}
\end{lemma}
\leanref{Development/AbelianVarietyTheory/Proofs/SiegelThetaTransformation/TensorLatticeSelfDualSemicharacterIndexStatement.lean}{GlobalTheory.T.tensorLattice_selfDual_semicharacter_index}
\leanref{Development/AbelianVarietyTheory/Proofs/SiegelThetaTransformation/TensorLatticeSelfDualSemicharacterIndex.lean}{GlobalTheory.SW.TensLat.mem_tens_iff}
\leanref{Development/AbelianVarietyTheory/Proofs/SiegelThetaTransformation/TensorLatticeSelfDualSemicharacterIndex.lean}{GlobalTheory.SW.TensLat.tens_fg}
\leanref{Development/AbelianVarietyTheory/Proofs/SiegelThetaTransformation/TensorLatticeSelfDualSemicharacterIndex.lean}{GlobalTheory.SW.TensLat.tens_selfDual}
\leanref{Development/AbelianVarietyTheory/Proofs/SiegelThetaTransformation/TensorLatticeSelfDualSemicharacterIndex.lean}{GlobalTheory.SW.TensLat.epsW_semichar}
\leanref{Development/AbelianVarietyTheory/Proofs/SiegelThetaTransformation/TensorLatticeSelfDualSemicharacterIndex.lean}{GlobalTheory.SW.TensLat.mem_tens_stdL}
\leanref{Development/AbelianVarietyTheory/Proofs/SiegelThetaTransformation/TensorLatticeSelfDualSemicharacterIndex.lean}{GlobalTheory.SW.TensLat.relIndex_tens_right}
\leanref{Development/AbelianVarietyTheory/Proofs/SiegelThetaTransformation/TensorLatticeSelfDualSemicharacterIndex.lean}{GlobalTheory.SW.TensLat.relIndex_tens_left}
\leanref{Development/AbelianVarietyTheory/Proofs/CuspFormEvaluations/IsLocalizedModuleSubtype.lean}{GlobalTheory.exists_dual_family}
\begin{proof}
Write $(v,w)_1=\sum_av_a\bar w_a$ for $v,w\in E^n$, and for $y\in M_{n\times r}(E)$ and $w\in E^n$ put
$c_w(y)=\sum_k\bar w_ky_k\in E^r$, so that $c_w(vu^{\mathsf T})=(v,w)_1\,u$ for $v\in E^n$ and $u\in E^r$. As
$h(au,v)=h(u,\bar av)$ for $a\in E$, $\langle au,v\rangle_X=\langle u,\bar av\rangle_X$, and so
$\omega_X(y,vu^{\mathsf T})=\langle c_v(y),u\rangle_X$.

\emph{A dual family.} The module $\Lambda_S$ is finitely generated and torsion free over the Dedekind ring
$\mathcal O_E$, hence projective. It spans $E^n$: otherwise, as $(\cdot,\cdot)_1$ is nondegenerate, some
$x\ne0$ has $(x,y)_1=0$ for every $y\in\Lambda_S$, and then $Ex\subset\Lambda_S$ by self-duality, which is
impossible for a finitely generated $\mathbf Z$-module. By the dual basis
lemma there are $v_1,\dots,v_m\in\Lambda_S$ and $\mathcal O_E$-linear forms $\phi_i$ on $\Lambda_S$ with
$x=\sum_i\phi_i(x)v_i$ for $x\in\Lambda_S$; extending the $\phi_i$ to $E$-linear forms on $E^n$ and writing them
as $x\mapsto(x,w_i)_1$, the $w_i$ lie in $\Lambda_S$ by self-duality, as $(y,w_i)_1=\phi_i(y)$ is integral for
$y\in\Lambda_S$. So $x=\sum_i(x,w_i)_1v_i$ for every $x\in E^n$, and hence $y=\sum_iv_i\,c_{w_i}(y)^{\mathsf T}$
for every $y\in M_{n\times r}(E)$.

\emph{Membership.} We show that $y\in A$ exactly when $c_w(y)\in L$ for every $w\in\Lambda_S$. The map $c_w$ is
additive, and for $v\in\Lambda_S$ and $u\in L$, $c_w(vu^{\mathsf T})=(v,w)_1u\in L$, as $(v,w)_1$ is an
algebraic integer; so $c_w(A)\subset L$. Conversely, if every $c_{w_i}(y)$ lies in $L$, then
$y=\sum_iv_ic_{w_i}(y)^{\mathsf T}\in A$.

(1) By the membership criterion and the dual family, $A$ is the image of $L^m$ under
$(u_i)\mapsto\sum_iv_iu_i^{\mathsf T}$; $L$ is finitely generated over $\mathcal O_E$, hence over $\mathbf Z$.

(2) Let $y\in A$. For $y'=vu^{\mathsf T}$ with $v\in\Lambda_S$ and $u\in L$,
$\omega_X(y,y')=\langle c_v(y),u\rangle_X\in\mathbf Z$, as $c_v(y)\in L$ and $L$ is self-dual; and
$\omega_X(y,\cdot)$ is additive. Conversely, if $\omega_X(y,y')\in\mathbf Z$ for every $y'\in A$, then for
$w\in\Lambda_S$ and $u\in L$, $\langle c_w(y),u\rangle_X=\omega_X(y,wu^{\mathsf T})\in\mathbf Z$, so $c_w(y)\in L$
by self-duality of $L$, and $y\in A$.

(3) For $y'\in A$, $w_0y'\in A$, as $c_v(w_0y')=w_0c_v(y')\in L$; so $t=\omega_X(y,w_0y')\in\mathbf Z$ by (2).
By definition $\varepsilon_{w_0}(y)=e^{\pi i\omega_X(y,w_0y)}$, so $\varepsilon_{w_0}(y)^2=e^{2\pi i\omega_X(y,w_0y)}=1$.
By the antisymmetry of $\langle\cdot,\cdot\rangle_X$ and $\langle au,v\rangle_X=\langle u,\bar av\rangle_X$,
$\omega_X(y',w_0y)=-\omega_X(y,\bar w_0y')$, and $\omega_X(y,w_0y')+\omega_X(y,\bar w_0y')=\omega_X(y,y')$ as
$w_0+\bar w_0=1$. So
\[
 \omega_X(y+y',w_0(y+y'))=\omega_X(y,w_0y)+\omega_X(y',w_0y')-\omega_X(y,y')+2t,
\]
and $e^{2\pi it}=1$ gives the second identity.

The membership criterion holds, with the same proof, for every $\mathcal O_E$-submodule $B\subset E^r$ in place of
$L$: $y\in\Lambda_S\otimes B$ exactly when $c_w(y)\in B$ for every $w\in\Lambda_S$.

(4) The module $\mathcal O_E^n$ lies in $\mathcal V_M$, with $t=1$, so the criterion applies to it. For the
standard basis vector $e_k$, $c_{e_k}(y)=y_k$; so if $y\in\mathcal O_E^n\otimes L$, every row lies in $L$.
Conversely, if every row lies in $L$, then $c_w(y)=\sum_k\bar w_ky_k\in L$ for $w\in\mathcal O_E^n$, as the $\bar w_k$
are algebraic integers.

(5) and (6) \emph{An index formula.} Let $\Lambda\subset E^n$ be a finitely generated $\mathcal O_E$-submodule that
spans $E^n$, $B\subset E^r$ a finitely generated $\mathcal O_E$-submodule, and $B'\subset E^r$ an
$\mathcal O_E$-submodule that spans $E^r$. We show that the index of $(\Lambda\otimes B)\cap(\Lambda\otimes B')$ in
$\Lambda\otimes B$ is $[B:B\cap B']^n$. The module $\Lambda$ has a pseudo-basis: a basis $e_1,\dots,e_n$ of $E^n$ and
nonzero fractional ideals $\mathfrak c_a$ of $E$ with $\Lambda=\{\sum_ax_ae_a:x_a\in\mathfrak c_a\}$. For a fractional
ideal $\mathfrak c$ and an $\mathcal O_E$-submodule $B_0\subset E^r$ put $\mathfrak cB_0=\{u\in E^r:bu\in B_0\text{
for every }b\in\mathfrak c^{-1}\}$, and for $y\in M_{n\times r}(E)$ let $y^a\in E^r$ be the vector of the $a$-th
coordinates of its columns, so that $y=\sum_ae_a(y^a)^{\mathsf T}$. Then $y\in\Lambda\otimes B_0$ exactly when
$y^a\in\mathfrak c_aB_0$ for every $a$. Indeed, for $v=\sum_ax_ae_a\in\Lambda$ and $u\in B_0$, $(vu^{\mathsf
T})^a=x_au$, and $bx_au\in B_0$ for $b\in\mathfrak c_a^{-1}$, as $bx_a\in\mathcal O_E$; conversely, if
$y^a\in\mathfrak c_aB_0$, write $1=\sum_ix_ib_i$ with $x_i\in\mathfrak c_a$ and $b_i\in\mathfrak c_a^{-1}$; then
$e_a(y^a)^{\mathsf T}=\sum_i(x_ie_a)(b_iy^a)^{\mathsf T}$ with $x_ie_a\in\Lambda$ and $b_iy^a\in B_0$. So $y\mapsto(y^a)_a$
identifies $\Lambda\otimes B$ with $\prod_a\mathfrak c_aB$ and its intersection with $\Lambda\otimes B'$ with
$\prod_a(\mathfrak c_aB\cap\mathfrak c_aB')$, and the index is $\prod_a[\mathfrak c_aB:\mathfrak c_aB\cap\mathfrak
c_aB']$. It remains to see $[\mathfrak cB:\mathfrak cB\cap\mathfrak cB']=[B:B\cap B']$ for a nonzero fractional ideal
$\mathfrak c$. As $B$ is finitely generated and $B'$ spans $E^r$, there is a nonzero $N\in\mathcal O_E$ with
$NB\subset B'$. As $\mathcal O_E$ is a Dedekind ring there is $x\in\mathfrak c$ with $N\mathcal O_E+x\mathfrak
c^{-1}=\mathcal O_E$; write $1=sN+xb$ with $s\in\mathcal O_E$ and $b\in\mathfrak c^{-1}$. The map $v\mapsto xv$ carries
$B$ into $\mathfrak cB$, as $b'x\in\mathcal O_E$ for $b'\in\mathfrak c^{-1}$, and it induces a bijection
$B/(B\cap B')\to\mathfrak cB/(\mathfrak cB\cap\mathfrak cB')$:
\begin{itemize}
\item for $v\in B$, $v\in B'$ exactly when $xv\in\mathfrak cB'$: if $xv\in\mathfrak cB'$, then $v=sNv+b(xv)$ with
  $sNv\in B'$ and $b(xv)\in B'$;
\item for $u\in\mathfrak cB$, $v=bu$ lies in $B$, and $u-xv=sNu$ lies in $\mathfrak cB'$, as $b'sNu=sN(b'u)\in NB\subset
  B'$ for $b'\in\mathfrak c^{-1}$.
\end{itemize}

(5) The module $\Lambda_S$ spans $E^n$ (\emph{A dual family}), $L$ is finitely generated, and $L'$ spans $E^r$, by
the argument given there for $\Lambda_S$, with the nondegenerate form $\langle\cdot,\cdot\rangle_X$. So the index
formula gives the first assertion. By the membership criterion, $y$ lies in $A\cap(\Lambda_S\otimes L')$ exactly when
$c_w(y)\in L\cap L'$ for every $w\in\Lambda_S$, that is, when $y\in\Lambda_S\otimes(L\cap L')$.

(6) Transposition $y\mapsto y^{\mathsf T}$ carries $\Lambda\otimes B$ onto $B\otimes\Lambda$, the group generated by the
$uv^{\mathsf T}$. So the index of $A\cap(\Lambda'_S\otimes L)$ in $A$ is the index of $(L\otimes\Lambda_S)\cap
(L\otimes\Lambda'_S)$ in $L\otimes\Lambda_S$, which is $[\Lambda_S:\Lambda_S\cap\Lambda'_S]^r$ by the index formula
with the roles of $(n,E^n)$ and $(r,E^r)$ exchanged: $L$ is finitely generated and spans $E^r$, $\Lambda_S$ is finitely
generated, and $\Lambda'_S$ spans $E^n$.
\end{proof}
% T.tensorLattice_selfDual_semicharacter_index = TensLat.tens_selfDual_semichar_index: (1) to (3) are tens_fg,
% tens_selfDual and epsW_semichar; (4) is mem_tens_stdL (mem_tens_iff with Base.stdL_mem_Vset n 1); (5) is
% relIndex_tens_right (with span_of_selfD for L') and the intersection through mem_tens_iff; (6) is relIndex_tens_left
% (tens_eq_comap, relIndex_tens_right) with DualFamily.span_eq_top for Lambda'. The index formula is
% relIndex_tens_right: exists_pseudoBasis_adapted for the pseudo-basis (the book takes its existence as known, as in
% Step 1 of the proof of Lemma is-onto-step-leaves-det-congruence), tens_eq_of_pseudoBasis (fracMul is c B_0, colCoord the
% y^a), TensorLat.relIndex_pi_univ, exists_mul_le for N and relIndex_fracMul (with
% IsDedekindDomain.exists_add_spanSingleton_mul_eq for x). exists_dual_family (DualFamily.span_eq_top, projective
% lifting, exists_linearForm_extend_of_span_eq_univ, DualFamily.exists_herm_eq) gives v, w; eq_sum_ctr, ctr_pure,
% mem_tens_iff, sum_EF_pure (EF_smul), smul_mem_tens; in epsW_semichar, h2 (EF_swap, EF_smul) and h3 (w + cc w = 1).

\begin{lemma}\label{is-eigenvalue-conjugate-frame-vacua-lattice}
Let $n,d\ge0$, write $\mathcal M=M_{n\times d}(\mathbf Q)$ and $[a,b]=\sum_ka_k\cdot b_k$ for $a,b\in\mathcal M$
with rows $a_k$ and $b_k$, and put $\Omega((a,b),(a',b'))=[a,b']-[a',b]$ on $\mathcal M^2$. For $p=(a,b)\in\mathcal
M^2$ and a function $\varphi$ on $\mathcal M$ let $(\rho(p)\varphi)(x)=\varphi(x-b)\,e^{\pi i(2[x,a]-[a,b])}$. Let
$A\subset\mathcal M^2$ be a finitely generated subgroup such that $p\in A$ exactly when $\Omega(p,p')\in\mathbf Z$
for every $p'\in A$, and let $\varepsilon$ be a function on $\mathcal M^2$ with $|\varepsilon(p)|=1$ and
$\varepsilon(p+p')=\varepsilon(p)\varepsilon(p')e^{\pi i\Omega(p,p')}$ for $p,p'\in A$, and $\varepsilon(mp)=1$ for
every $p\in A$, for some integer $m>0$. Then there is a Schwartz-Bruhat function $v$ on $\mathcal M$ with
$\langle v,v\rangle=1$ and $\varepsilon(p)\rho(p)v=v$ for every $p\in A$, and every Schwartz-Bruhat $w$ with
$\varepsilon(p)\rho(p)w=w$ for every $p\in A$ is a multiple of $v$.
\end{lemma}
\leanref{Development/AbelianVarietyTheory/Proofs/SiegelThetaTransformation/ExistsExOfPow.lean}{GlobalTheory.SW.Vac.vacuum}
\leanref{Development/AbelianVarietyTheory/Proofs/SiegelThetaTransformation/ExistsExOfPow.lean}{GlobalTheory.SW.Vac.exists_graph}
\leanref{Development/AbelianVarietyTheory/Proofs/SiegelThetaTransformation/ExistsExOfPow.lean}{GlobalTheory.SW.Vac.exists_base}
\leanref{Development/AbelianVarietyTheory/Proofs/SiegelThetaTransformation/ExistsExOfPow.lean}{GlobalTheory.SW.Vac.exists_lift}
\begin{proof}
Write $e(t)=e^{\pi it}$ for $t\in\mathbf Q$, and for $p=(a,b)$ put $f_p(x)=\varepsilon(p)\,e(2[x,a]-[a,b])$. Then
$\varepsilon(p)\rho(p)w=w$ exactly when $w(x)=f_p(x)\,w(x-b)$ for every $x$, and the law of $\varepsilon$ gives
$f_{p+p'}(x)=f_p(x)\,f_{p'}(x-b)$ for $p=(a,b)$ and $p'$ in $A$, by expanding the exponents.

\emph{A graph lemma.} Let $\Gamma\subset\mathcal M\times\mathbf Q$ be a subgroup whose first projection lies in
a finitely generated subgroup of $\mathcal M$. Then there is $x\in\mathcal M$ with $(c,[x,c])\in\Gamma$ for every
$(c,s)\in\Gamma$. Indeed, the first projection $C$ of $\Gamma$ is finitely generated and torsion free, hence
free, so the projection $\Gamma\to C$ has a section $c\mapsto(c,\lambda(c))$ with $\lambda$ additive; as
$\mathbf Q$ is divisible, $\lambda$ extends to an additive map $\mathcal M\to\mathbf Q$, which is $\mathbf
Q$-linear and so of the form $c\mapsto[x,c]$.

\emph{A base point.} The pairs $(a,0)\in A$ form a group $A_0$, on which $a\mapsto\varepsilon(a,0)$ is a
character, as $\Omega$ vanishes there; its values are $m$-th roots of unity, so each is $e(2s)$ for a rational
$s$. Apply the graph lemma to $\Gamma=\{(a,s):(a,0)\in A,\ \varepsilon(a,0)e(2s)=1\}$, a subgroup that contains
such a pair for each $a$: there is $x_0$ with $\varepsilon(a,0)\,e(2[x_0,a])=1$ for every $(a,0)\in A$.

\emph{The vacuum.} Let $B$ be the set of the $b$ with $(a,b)\in A$ for some $a$. Put $v_0(x_0+b)=f_p(x_0+b)$ for
$p=(a,b)\in A$, and $v_0(x)=0$ for $x\notin x_0+B$. This does not depend on $a$: if $p'=(a',b)\in A$ too, then
$(a'-a,0)\in A$ and $f_{p'}(x_0+b)=f_p(x_0+b)f_{(a'-a,0)}(x_0)=f_p(x_0+b)$ by the choice of $x_0$. For $p=(a,b)\in A$
and $x=x_0+b'$ with $p'=(a',b')\in A$, $v_0(x)=f_{p'}(x)=f_p(x)f_{p'-p}(x-b)=f_p(x)v_0(x-b)$; and for
$x\notin x_0+B$ also $x-b\notin x_0+B$, so both sides vanish. So $v_0$ is fixed by every
$\varepsilon(p)\rho(p)$, and $v_0(x_0)=f_0(x_0)=\varepsilon(0)=1$.

\emph{Uniqueness.} Let $w$ be fixed by every $\varepsilon(p)\rho(p)$. Then $w(x_0+b)=f_p(x_0+b)w(x_0)$ for
$p=(a,b)\in A$, so $w=w(x_0)v_0$ on $x_0+B$. Let $x\notin x_0+B$ and suppose $w(x)\ne0$. For $(a,0)\in A$,
$w(x)=f_{(a,0)}(x)w(x)$ gives $\varepsilon(a,0)e(2[x,a])=1$, and with the choice of $x_0$, $e(2[x-x_0,a])=1$,
that is, $[x-x_0,a]\in\mathbf Z$. Put $c=x-x_0$ and apply the graph lemma to the group $\Gamma$ of the pairs
$(b',s)$ for which some $(a',b')\in A$ has $s-[a',c]\in\mathbf Z$. It gives $x'$ such that for every
$(a_1,b_1)\in A$ some $(a',b_1)\in A$ has $[x',b_1]-[a',c]\in\mathbf Z$; as $(a'-a_1,0)\in A$,
$[a'-a_1,c]\in\mathbf Z$, so $\Omega((x',c),(a_1,b_1))=[x',b_1]-[a_1,c]\in\mathbf Z$. By the self-duality of $A$,
$(x',c)\in A$, so $x\in x_0+B$, a contradiction. So $w=w(x_0)v_0$.

\emph{Level and norm.} Let $N_1$, $N_2$ and $D$ be positive integers with $N_1a$, $N_2b$ and $Dx_0$ integral
for every $(a,b)\in A$, and put $R=mN_1N_2D$. Then $v_0(x)\ne0$ only for $x\in x_0+B$, where $Rx$ is integral.
For an integral $z$, the pair $(0,N_1N_2Dz)$ lies in $A$ by self-duality, as $\Omega((0,N_1N_2Dz),(a,b))=-[a,N_1N_2Dz]$
is an integer; $m$ times it is $(0,Rz)$, with $\varepsilon(0,Rz)=1$ and $f_{(0,Rz)}=1$, so
$v_0(x+Rz)=v_0(x)$. So $v_0$ is Schwartz-Bruhat of level $R$, and $\langle v_0,v_0\rangle$ is a positive real
number, as $v_0(x_0)=1$. Put $v=\langle v_0,v_0\rangle^{-1/2}v_0$.
\end{proof}
% Vac.vacuum: fixed_iff and fac_add for f_p; exists_graph (Module.free_of_finite_type_torsion_free',
% projective_lifting_property, Module.Baer.of_divisible, exists_pr); exists_base for x0 (exists_ex_of_pow for the
% roots of unity); v with hv, hv0, hvx and hfix; huniq with exists_lift; the level m N1 N2 D (exists_den_fg,
% exists_den); T.schwartzBruhatInnerProduct_self_nonneg_bounds for 0 < ip v v, and the normalization.

\begin{lemma}\label{is-eigenvalue-conjugate-frame-vacua}
Let $X$, its symplectic basis and indexing be as in
Definition~\ref{is-eigenvalue-conjugate-frame-heisenberg}, let $n\ge0$ be even (the parity is not used), let
$w_0\in E$ be an algebraic integer with $w_0+\bar w_0=1$, and let $L\subset E^r$ be an $\mathcal O_E$-submodule
that is self-dual for $X$ (Definition~\ref{is-eigenvalue-conjugate-frame-fibres}(7)). For an integer $M\ge0$ let
$\mathcal V_M$ be the set of the finitely generated $\mathcal O_E$-submodules $\Lambda_S\subset E^n$ such
that $x\in\Lambda_S$ exactly when $\sum_ax_a\bar y_a$ is an algebraic integer for every $y\in\Lambda_S$, and
such that some integer $t>0$ prime to $M$ has $t\Lambda_S\subset\mathcal O_E^n$ and $t\mathcal
O_E^n\subset\Lambda_S$; it contains $\mathcal O_E^n$. Let $M\ge0$ and $\Lambda_S\in\mathcal V_M$.
\begin{enumerate}
\item The tensor lattice $\Lambda_S\otimes L$ has a vacuum for $w_0$.
\item If $v$ is such a vacuum and $\varphi$ is a Schwartz-Bruhat function with
  $\varepsilon_{w_0}(y)\,\rho_X(y)\varphi=\varphi$ for every $y\in\Lambda_S\otimes L$, then
  $\varphi=\langle\varphi,v\rangle\,v$.
\item Let $v$ be such a vacuum and $\varphi$ a Schwartz-Bruhat function, and let $O_\varphi$ be the set of
  the functions $\varepsilon_{w_0}(y)\,\rho_X(y)\varphi$, $y\in\Lambda_S\otimes L$. Put
  $\pi(\varphi)=|O_\varphi|^{-1}\sum_{\psi\in O_\varphi}\psi$ if $O_\varphi$ is finite, and $\pi(\varphi)=0$
  otherwise. Then $\pi(\varphi)=\langle\varphi,v\rangle\,v$.
\item Let $L'\subset E^r$ be a second $\mathcal O_E$-submodule self-dual for $X$, $\Lambda'_S\in\mathcal V_M$,
  and $v$ and $v'$ vacua of $\Lambda_S\otimes L$ and $\Lambda'_S\otimes L'$ for $w_0$. Then the index $j$ of
  $(\Lambda_S\otimes L)\cap(\Lambda'_S\otimes L')$ in $\Lambda_S\otimes L$ is finite, and
  $j\,|\langle v,v'\rangle|^2=1$.
\end{enumerate}
\end{lemma}
\leanref{Development/AbelianVarietyTheory/Proofs/SiegelThetaTransformation/TensorLatticeVacuaExistUnique.lean}{GlobalTheory.SW.TensVac.vacua}
\leanref{Development/AbelianVarietyTheory/Definitions/SpanningFamilies/FirstCoordinateAdmissible.lean}{GlobalTheory.SW.Vset}
\leanref{Development/AbelianVarietyTheory/Definitions/SpanningFamilies/IsSystem.lean}{GlobalTheory.SW.projA}
\leanref{Development/AbelianVarietyTheory/Definitions/SpanningFamilies/FirstCoordinateAdmissible.lean}{GlobalTheory.SW.Base.stdL_mem_Vset}
\begin{proof}
Write $\varepsilon=\varepsilon_{w_0}$, $\omega_X$ as in Lemma~\ref{is-eigenvalue-conjugate-frame-vacua-heisenberg},
and $\pi_y=\varepsilon(y)\rho_X(y)$. Call a subgroup $A\subset M_{n\times r}(E)$ \emph{good} if it has the three
properties of Lemma~\ref{is-eigenvalue-conjugate-frame-vacua-tensor}: it is finitely generated, $y\in A$ exactly
when $\omega_X(y,y')\in\mathbf Z$ for every $y'\in A$, and $\varepsilon(y)^2=1$ and
$\varepsilon(y+y')=\varepsilon(y)\varepsilon(y')e^{-\pi i\omega_X(y,y')}$ for $y,y'\in A$. By that lemma, $\Lambda_S\otimes L$
and $\Lambda'_S\otimes L'$ are good. By Lemma~\ref{is-eigenvalue-conjugate-frame-vacua-mean}, sums and complex
multiples of Schwartz-Bruhat functions are Schwartz-Bruhat, and $\langle\varphi,\psi\rangle$, the mean of
$\varphi\bar\psi$, is linear in $\varphi$ and satisfies $\langle\psi,\varphi\rangle=\overline{\langle\varphi,\psi\rangle}$.
Let $A$ be good. By
Lemma~\ref{is-eigenvalue-conjugate-frame-vacua-heisenberg}, for $y,y'\in A$ and a Schwartz-Bruhat $\varphi$:
$\pi_y\varphi$ is Schwartz-Bruhat; $\pi_0\varphi=\varphi$; $\pi_y\pi_{y'}\varphi=\pi_{y+y'}\varphi$, from (2) and the
law of $\varepsilon$; and $\langle\pi_y\varphi,\pi_y\psi\rangle=\langle\varphi,\psi\rangle$ for Schwartz-Bruhat $\psi$,
from (3) and (4). Also $\varepsilon(2y)=\varepsilon(y)^2e^{-\pi i\omega_X(y,y)}=1$, as $\omega_X(y,y)=0$.

\emph{Existence and uniqueness.} By Lemma~\ref{is-eigenvalue-conjugate-frame-vacua-heisenberg}(1), the
coordinates carry $A$ to a finitely generated subgroup $A^c$ of $M_{n\times d}(\mathbf Q)^2$ and $\omega_X$ to
$-\Omega$, in the notation of Lemma~\ref{is-eigenvalue-conjugate-frame-vacua-lattice}; and $\rho_X(y)$ is the
operator $\rho$ of the coordinates of $y$. So $A^c$ is self-dual for $\Omega$, and $\varepsilon$, read through the
coordinates, satisfies the hypotheses of that lemma with $m=2$. It gives a vacuum $v_0$ of $A$, and every
Schwartz-Bruhat $\varphi$ fixed by the $\pi_y$, $y\in A$, is $c\,v_0$ for some $c\in\mathbf C$. This is (1) for
$A=\Lambda_S\otimes L$. For (2), let $v$ be a vacuum and $\varphi$ fixed; then $v=c'v_0$ with
$|c'|^2=\langle v,v\rangle=1$ and $\varphi=cv_0=(c/c')v$, and $\langle\varphi,v\rangle=c\bar c'=c/c'$.

\emph{Finiteness of the orbits.} Let $\varphi$ be Schwartz-Bruhat of level $R$, and let $D\ge1$ be an integer
such that the coordinates of all elements of $A$ become integral when multiplied by $D$, which exists as $A$ is
finitely generated. Put $N=2RD$. For $b\in A$ the coordinates of $Nb$ are $2R$ times integral matrices
$(\alpha,\beta)$, so $\rho_X(Nb)\varphi(x)=\varphi(x-2R\beta)e^{\pi i(4R[x,\alpha]-4R^2[\alpha,\beta])}=\varphi(x)$: the
first factor by the level, and the exponential is $1$ where $\varphi(x)\ne0$, as $Rx$ is then integral. And
$\varepsilon(Nb)=1$, as $Nb$ is twice an element of $A$. So $\pi_{a+Nb}\varphi=\pi_a\pi_{Nb}\varphi=\pi_a\varphi$ for
$a,b\in A$, and the orbit $O_\varphi=\{\pi_y\varphi:y\in A\}$ is the image of the finite group $A/NA$.

(3) Let $v$ be a vacuum of $A=\Lambda_S\otimes L$. The orbit $O_\varphi$ is finite and contains $\varphi=\pi_0\varphi$.
Let $\Sigma$ be the sum of its elements, which is Schwartz-Bruhat. For $z\in A$, $\psi\mapsto\pi_z\psi$ is a
bijection of $O_\varphi$ onto itself, with inverse $\psi\mapsto\pi_{-z}\psi$, as $\pi_z\pi_y\varphi=\pi_{z+y}\varphi$; so
$\pi_z\Sigma=\Sigma$, and $\Sigma=\langle\Sigma,v\rangle v$ by (2). For $\psi=\pi_y\varphi$,
$\langle\psi,v\rangle=\langle\pi_y\varphi,\pi_yv\rangle=\langle\varphi,v\rangle$. So
$\langle\Sigma,v\rangle=|O_\varphi|\,\langle\varphi,v\rangle$, and $\pi(\varphi)=|O_\varphi|^{-1}\Sigma=\langle\varphi,v\rangle v$.

(4) Put $A=\Lambda_S\otimes L$ and $A'=\Lambda'_S\otimes L'$, both good, with vacua $v$ and $v'$; write $\pi_y$ for
$\varepsilon(y)\rho_X(y)$ for every $y$.

First, $\langle\pi_yv',v'\rangle=0$ for $y\in A\setminus A'$. As $y\notin A'$, some $y'\in A'$ has
$t=\omega_X(y,y')\notin\mathbf Z$. By Lemma~\ref{is-eigenvalue-conjugate-frame-vacua-heisenberg}(1) and (2),
$\pi_{y'}\pi_yv'=\varepsilon(y')\varepsilon(y)e^{-\pi i\omega_X(y',y)}\rho_X(y+y')v'$, while
$\pi_yv'=\pi_y\pi_{y'}v'=\varepsilon(y)\varepsilon(y')e^{-\pi i\omega_X(y,y')}\rho_X(y+y')v'$; so
$\pi_{y'}\pi_yv'=e^{2\pi it}\pi_yv'$. As $\pi_{y'}$ preserves inner products and fixes $v'$,
$\langle\pi_yv',v'\rangle=\langle\pi_{y'}\pi_yv',v'\rangle=e^{2\pi it}\langle\pi_yv',v'\rangle$, and $e^{2\pi it}\ne1$.

So for $z\in A$, $\pi_zv'=v'$ exactly when $z\in A'$: if $z\notin A'$, then
$\langle\pi_zv',v'\rangle=0\ne1=\langle v',v'\rangle$. Hence for $a,b\in A$, $\pi_av'=\pi_bv'$ exactly when
$\pi_{-a+b}v'=v'$, that is, when $-a+b\in A\cap A'$. So the orbit $T=\{\pi_yv':y\in A\}$, which is finite, has
$j=|T|$ elements, the index of $A\cap A'$ in $A$. The elements of $T$ are orthonormal:
$\langle\pi_bv',\pi_av'\rangle=\langle\pi_{-a+b}v',v'\rangle$, which is $1$ if $-a+b\in A'$ and $0$ otherwise, by the
first step. So the sum $\Sigma$ of the elements of $T$ has $\langle\Sigma,\Sigma\rangle=j$. By (3) for $\varphi=v'$,
$j^{-1}\Sigma=\langle v',v\rangle v$. Taking the inner product of each side with itself gives
$j^{-1}=|\langle v',v\rangle|^2=|\langle v,v'\rangle|^2$.

Finally, $\mathcal O_E^n\in\mathcal V_M$, with $t=1$: $x\in\mathcal O_E^n$ exactly when $\sum_ax_a\bar y_a$ is
integral for every $y\in\mathcal O_E^n$.
% TensVac.vacua (10 lines): Good X w (tens Lambda L) from T.tensorLattice_selfDual_semicharacter_index (Lemma
% -vacua-tensor); Good.vac (crdHom_bijective, om_crdHom, Vac.vacuum with m = 2 by Good.eps_two), Good.uniq,
% Good.finite (Vac.exists_den_fg, heisS_fix, finite_image_of_fg), Good.proj (sum_sb_ip, rho_ip, Good.rho_rho,
% Good.uniq), Good.ip_eq_zero and Good.overlap (ncard_image_eq for the index, the orthonormality horth and
% Good.proj at v'). The index in (4) is AddSubgroup.relIndex (tens Lambda' L') (tens Lambda L), cast to a real
% number. stdL_mem_Vset: O_E^n lies in Vset E n M with t = 1. The laws come from SW.heisG_laws (Lemma
% -vacua-heisenberg). The hypothesis Even n is passed on and not used.
\end{proof}

\begin{lemma}\label{is-eigenvalue-conjugate-frame-theta-product-vacuum}
Let $X$, $n$, $w_0$, $L$ be as in Lemma~\ref{is-eigenvalue-conjugate-frame-vacua}, put $d=|\Psi|r$, and let $v$ be a
vacuum of $\mathcal O_E^n\otimes L$ for $w_0$. Let $b_1,\dots,b_d,b'_1,\dots,b'_d$ be a $\mathbf Z$-basis of $L$ with
$\langle b_k,b'_l\rangle_X=-\delta_{kl}$, $\langle b'_l,b_k\rangle_X=\delta_{kl}$ and $\langle b_k,b_l\rangle_X=\langle
b'_k,b'_l\rangle_X=0$, and write $\epsilon(u)=e^{\pi i\langle u,w_0u\rangle_X}$. Then there are $c_1,c_2\in\{0,\frac12\}^d$
and $C\in\mathbf C^\times$ such that for every $x\in\mathcal D_X$, with $h=M_1^{-1}$ and $\tau=hM_2$, where $M_1$ and $M_2$
are the matrices with columns $p_x(b_k)$ and $p_x(b'_k)$:
\begin{enumerate}
\item $h$ is invertible, $\tau$ is symmetric with positive definite imaginary part $Y$, and
  $H_x(z,w)=(hz)^{\mathsf T}Y^{-1}\overline{hw}$ for all $z,w\in\mathbf C^d$;
\item $\epsilon(b_k)=e^{2\pi ic_{1,k}}$ and $\epsilon(b'_k)=e^{2\pi ic_{2,k}}$ for every $k$;
\item for every $z\in(\mathbf C^d)^n$, with the theta functions $\vartheta[c](\tau,\cdot)$ of
  Theorem~\ref{cv-uniformization-theta-transformation},
  \[
   \theta_v(x;z)=C\det(h)^{n/2}\prod_{k=1}^n\exp\Big(\frac\pi2(hz_k)^{\mathsf T}Y^{-1}(hz_k)\Big)\,
   \vartheta[c_1,c_2](\tau,hz_k).
  \]
\end{enumerate}
\end{lemma}
\leanref{Development/AbelianVarietyTheory/Proofs/SiegelThetaTransformation/PeriodCoordinateThetaOfSignedHeisenbergVacuumStatement.lean}{GlobalTheory.T.periodCoordinateTheta_of_signedHeisenbergVacuum}
\begin{proof}
\openstep{The proof, which writes the vacuum of the tensor lattice in the coordinates of the symplectic basis of $L$ as
a product of signed point masses and computes its theta series as a product of theta functions with characteristics,
is not written out.}
% GlobalTheory.T.periodCoordinateTheta_of_signedHeisenbergVacuum = Dict.VacForm.thA_vacuum (73 lines;
% PeriodCoordinateThetaOfSignedHeisenbergVacuum.lean, 522 lines). The Lean states the basis condition as EF X (b i)
% (b j) = -omegaStd (e_i) (e_j) and L as the Z-span of b; tau lies in SW.Siegel. Size (tools/lean_size.py thA_vacuum):
% 3338 lines. Top level: sigma = spB X bas b, the rational symplectic matrix of b in the basis bas (spB_mem, 96 lines);
% a transform omega of Schwartz-Bruhat functions for sigma with the transformation law of the theta series
% (SW.exists_weil_of_symplectic, 1448 lines, and DblEx.can_trans, 1821); eps(b_k) = +-1 (LatticeBoundTests.exists_half,
% 25); the periods of b (periods_spB, 890); the twisted Heisenberg operators of the basis vectors fix omega v
% (HeisFix.heis_fix, 2665); so omega v is C times the signed point mass (eq_smul_chi, 127), C ne 0 by the injectivity of
% the theta series (UnAct.eq_of_thA_eq, 787) and ip v v = 1; and the theta series of the signed point mass is the
% product of theta functions with characteristics (CharLink.thS_chi, 189). These overlap.
\end{proof}

\begin{lemma}\label{is-eigenvalue-conjugate-frame-theta-product-unique-fourier}
Let $g\ge0$ and $n\ge1$ be integers, $\tau\in M_g(\mathbf C)$ and $b_1,\dots,b_g\in\mathbf C^\times$, and write
$e_1,\dots,e_g$ for the standard basis of $\mathbf C^g$. Let $f_1,\dots,f_N$ be linearly independent holomorphic
functions on $\mathbf C^g$ with
\[
 f_k(z+e_i)=f_k(z),\qquad f_k(z+\tau e_j)=b_j^n\,e^{-2\pi inz_j}f_k(z)
\]
for all $k$, $i$, $j$ and $z\in\mathbf C^g$. Then $N\le n^g$.
\end{lemma}
\leanref{Development/AbelianVarietyTheory/Proofs/ToricGeometry/CardLePowOfLinearIndependent.lean}{GlobalTheory.ClassicalTheta.card_le_pow_of_linearIndependent}
\leanref{Development/AbelianVarietyTheory/Proofs/ToricGeometry/FourierCoeffUnitCubeTranslateOfEntirePeriodic.lean}{GlobalTheory.fourierCoeff_unitCube_translate_of_entire_periodic}
\leanref{Development/AbelianVarietyTheory/Proofs/ToricGeometry/EntirePeriodicEqZeroOfFourierCoeffEqZero.lean}{GlobalTheory.entire_periodic_eq_zero_of_fourierCoeff_eq_zero}
\begin{proof}
For a holomorphic function $F$ on $\mathbf C^g$ with $F(z+e_i)=F(z)$ for every $i$, and $m\in\mathbf Z^g$, put
$c_m(F)=\int_{[0,1]^g}F(x)e^{-2\pi i\,m\cdot x}\,dx$.

(a) For $w\in\mathbf C^g$, $\int_{[0,1]^g}F(x+w)e^{-2\pi i\,m\cdot x}\,dx=e^{2\pi i\,m\cdot w}c_m(F)$. Indeed,
$F_m(z)=F(z)e^{-2\pi i\,m\cdot z}$ is holomorphic and invariant under the $e_i$, and $G(w)=\int_{[0,1]^g}F_m(x+w)\,dx$ is
holomorphic in $w$, by differentiation under the integral over the compact cube. For real $t$, $G(w+t)=G(w)$, as the
integral of the periodic function $x\mapsto F_m(x+w)$ over $t+[0,1]^g$ is its integral over $[0,1]^g$. So the derivatives
of $G$ in the real directions vanish, hence, by complex linearity, all its derivatives, and $G$ is constant. Thus
$G(w)=G(0)=c_m(F)$, which is the claim.

(b) If $c_m(F)=0$ for every $m$, then $F=0$. Indeed, for fixed $w$ the function $x\mapsto F(x+w)$ on $\mathbf R^g$ is
continuous and invariant under $\mathbf Z^g$, and by (a) all its Fourier coefficients vanish; so it is $0$, as a
continuous function on $\mathbf R^g/\mathbf Z^g$ is determined by its Fourier coefficients. In particular $F(w)=0$.

(c) Let $F$ be a linear combination of the $f_k$; it satisfies the two laws of the statement. For $m\in\mathbf Z^g$ and
$j$, (a) with $w=\tau e_j$ and the second law give
\[
 e^{2\pi i\,m\cdot\tau e_j}c_m(F)=\int_{[0,1]^g}F(x+\tau e_j)e^{-2\pi i\,m\cdot x}\,dx
 =b_j^n\int_{[0,1]^g}F(x)e^{-2\pi i(m+ne_j)\cdot x}\,dx=b_j^nc_{m+ne_j}(F).
\]
As $b_j\ne0$, $c_m(F)=0$ if and only if $c_{m+ne_j}(F)=0$. Every $m\in\mathbf Z^g$ is $s+n\sum_jk_je_j$ with
$s\in\{0,\dots,n-1\}^g$ and $k\in\mathbf Z^g$. So if $c_s(F)=0$ for the $n^g$ vectors $s\in\{0,\dots,n-1\}^g$, then
$c_m(F)=0$ for every $m$, and $F=0$ by (b).

Hence the linear map $F\mapsto(c_s(F))_s$ from the span of the $f_k$ to $\mathbf C^{n^g}$ is injective, and $N\le n^g$.
\end{proof}
% ClassicalTheta.card_le_pow_of_linearIndependent (CardLePowOfLinearIndependent.lean), whose two hypotheses are
% fourierCoeff_unitCube_translate_of_entire_periodic ((a): setIntegral_unitCube_add_of_entire_periodic through
% differentiableOn_setIntegral_of_holomorphic_of_continuousOn and SCV.eq_of_invariant_real) and
% entire_periodic_eq_zero_of_fourierCoeff_eq_zero ((b): eq_zero_of_cube_fourierCoeff_eq_zero through
% UnitAddTorus.hasSum_mFourier_series_apply_of_summable); followed: key (forall_of_box_of_translate) and
% fintype_card_le_finrank.

\begin{lemma}\label{is-eigenvalue-conjugate-frame-theta-product-unique-normal-form}
Let $g\ge0$, let $\Lambda\subset\mathbf C^g$ be a discrete subgroup that spans $\mathbf C^g$ over $\mathbf R$, $H\in
M_g(\mathbf C)$ a Hermitian matrix with $\operatorname{Re}(z^{\mathsf T}H\bar z)>0$ for $z\ne0$, $\chi$ a function on
$\Lambda$ with $|\chi(l)|=1$, and $E$ a $\mathbf Z$-bilinear form on $\Lambda$ with $E(l,m)=\operatorname{Im}(l^{\mathsf
T}H\bar m)$ such that $l\mapsto E(l,\cdot)$ is a bijection from $\Lambda$ onto $\operatorname{Hom}(\Lambda,\mathbf Z)$. Put
$a(l,z)=\chi(l)\exp\big(\pi z^{\mathsf T}H\bar l+\frac\pi2l^{\mathsf T}H\bar l\big)$ for $l\in\Lambda$ and $z\in\mathbf
C^g$. Then there are an invertible $P\in M_g(\mathbf C)$, a symmetric $\tau\in M_g(\mathbf C)$ with positive definite
imaginary part, a holomorphic function $Q$ on $\mathbf C^g$ with $Q(0)=0$, and $b_1,\dots,b_g\in\mathbf C^\times$ such
that:
\begin{enumerate}
\item $Pz\in\Lambda$ if and only if $z=m+\tau m'$ with $m,m'\in\mathbf Z^g$;
\item $a(Pe_i,Pz)=e^{Q(z+e_i)-Q(z)}$ for all $i$ and $z$;
\item $a(P\tau e_j,Pz)=b_j\,e^{-2\pi iz_j}e^{Q(z+\tau e_j)-Q(z)}$ for all $j$ and $z$.
\end{enumerate}
\end{lemma}
\leanref{Development/AbelianVarietyTheory/Proofs/ToricGeometry/IsUnitOfLagrangian.lean}{GlobalTheory.SiegelNormalForm.exists_normal_form}
\leanref{Development/AbelianVarietyTheory/Proofs/ToricGeometry/IsUnitOfLagrangian.lean}{GlobalTheory.SiegelNormalForm.isUnit_of_lagrangian}
\leanref{Development/AbelianVarietyTheory/Proofs/ToricGeometry/IsUnitOfLagrangian.lean}{GlobalTheory.SiegelNormalForm.siegel_of_relations}
\begin{proof}
As $H$ is Hermitian, $l^{\mathsf T}H\bar l$ is real, so $E(l,l)=0$ and $E$ is alternating. As $\Lambda$ is discrete and
spans $\mathbf C^g$, it is free of rank $2g$. By the symplectic basis theorem for the unimodular alternating form $-E$
(proof of Lemma~\ref{cv-uniformization-adapted-basis}), $\Lambda$ has a basis $\lambda_1,\dots,\lambda_g,
\mu_1,\dots,\mu_g$ with $E(\lambda_i,\lambda_j)=E(\mu_i,\mu_j)=0$ and $E(\lambda_i,\mu_j)=-\delta_{ij}$. Let $P$ and $P'$
be the matrices with columns $\lambda_i$ and $\mu_j$, and $Y=P^{\mathsf T}H\bar P$, with entries $\lambda_i^{\mathsf
T}H\bar\lambda_j$. Their imaginary parts are $E(\lambda_i,\lambda_j)=0$, so $Y$ is real; it is Hermitian, so symmetric;
and $x^{\mathsf T}Yx=(Px)^{\mathsf T}H\overline{Px}$ for real $x$.

\emph{$P$ is invertible.} Otherwise $\det Y=0$, and there is a real $x\ne0$ with $Yx=0$. Then $(Px)^{\mathsf
T}H\overline{Px}=x^{\mathsf T}Yx=0$, so $Px=0$ by positivity, and
$0=\operatorname{Im}\big((Px)^{\mathsf T}H\bar\mu_j\big)=\sum_ix_iE(\lambda_i,\mu_j)=-x_j$ for every $j$, which is absurd.
So $Y$ is positive definite, and $(Pu)^{\mathsf T}H\overline{Pv}=u^{\mathsf T}Y\bar v$ for $u,v\in\mathbf C^g$.

\emph{The matrix $\tau$.} Put $\tau=P^{-1}P'$, so that $P\tau e_j=\mu_j$, and write $\tau=A+iB$ with $A,B$ real. Then
$-\delta_{ij}=E(\lambda_i,\mu_j)=\operatorname{Im}\big(e_i^{\mathsf T}Y\overline{\tau e_j}\big)=-(YB)_{ij}$, so $YB=1$, and
$B=Y^{-1}$ is symmetric and positive definite. Also $0=E(\mu_i,\mu_j)=\operatorname{Im}\big(\tau^{\mathsf T}Y\bar\tau\big)_{ij}$,
and $\operatorname{Im}(\tau^{\mathsf T}Y\bar\tau)=BYA-A^{\mathsf T}YB=A-A^{\mathsf T}$; so $A$, and with it $\tau$, is
symmetric. As $P$ is invertible, $Pz=\sum_im_i\lambda_i+\sum_jm'_j\mu_j=P(m+\tau m')$ describes the elements of $\Lambda$,
which is (1).

\emph{The function $Q$.} Choose $c_i$ with $e^{c_i}=\chi(\lambda_i)$, and put $Q(z)=\frac\pi2z^{\mathsf T}Yz+\sum_ic_iz_i$.
It is holomorphic with $Q(0)=0$, and $Q(z+w)-Q(z)=\pi z^{\mathsf T}Yw+Q(w)$, as $Y$ is symmetric. Then
$a(\lambda_i,Pz)=\chi(\lambda_i)\exp\big(\pi z^{\mathsf T}Ye_i+\frac\pi2Y_{ii}\big)=e^{Q(z+e_i)-Q(z)}$, which is (2). As
$Y\overline{\tau e_j}=Y\tau e_j-2iYBe_j=Y\tau e_j-2ie_j$, $\pi(Pz)^{\mathsf T}H\bar\mu_j=\pi z^{\mathsf T}Y\tau e_j-2\pi iz_j$,
and with $b_j=\chi(\mu_j)\exp\big(\frac\pi2\mu_j^{\mathsf T}H\bar\mu_j-Q(\tau e_j)\big)\ne0$,
\[
 a(\mu_j,Pz)=\chi(\mu_j)e^{\frac\pi2\mu_j^{\mathsf T}H\bar\mu_j}e^{-2\pi iz_j}e^{\pi z^{\mathsf T}Y\tau e_j}
 =b_je^{-2\pi iz_j}e^{Q(z+\tau e_j)-Q(z)},
\]
which is (3).
\end{proof}
% SiegelNormalForm.exists_normal_form (IsUnitOfLagrangian.lean), followed: hrank (ZLattice.rank), halt,
% exists_isSymplecticBasis_of_finrank (the book cites the symplectic basis theorem), hMY (Y real), isUnit_of_lagrangian,
% hrel, hYpos, siegel_of_relations (hYB, hBt, hsecond, hre), quad with c i = log chi (b (inl i)), quad_add_sub,
% hconj_col, and b j = chi (b (inr j)) exp(pi/2 H(b, b) - Q(tau e_j)).

\begin{lemma}\label{is-eigenvalue-conjugate-frame-theta-product-unique-count}
Let $g$, $\Lambda$, $H$, $\chi$, $E$ and $a(l,z)$ be as in
Lemma~\ref{is-eigenvalue-conjugate-frame-theta-product-unique-normal-form}, and let $n\ge1$. Let $f_1,\dots,f_N$ be
linearly independent holomorphic functions on $\mathbf C^g$ with $f_k(z+l)=a(l,z)^nf_k(z)$ for all $k$, $l\in\Lambda$ and
$z\in\mathbf C^g$. Then $N\le n^g$.
\end{lemma}
\leanref{Development/AbelianVarietyTheory/Proofs/ToricGeometry/CardLePowOfLinearIndependentThetaPrincipal.lean}{GlobalTheory.card_le_pow_of_linearIndependent_theta_principal}
\begin{proof}
Let $P$, $\tau$, $Q$ and $b$ be as in Lemma~\ref{is-eigenvalue-conjugate-frame-theta-product-unique-normal-form}, and put
$F_k(z)=f_k(Pz)e^{-nQ(z)}$. The $F_k$ are holomorphic, and linearly independent, as $z\mapsto Pz$ is bijective and
$e^{-nQ}$ has no zeros. By (1) of that lemma $Pe_i$ and $P\tau e_j$ lie in $\Lambda$, and (2) and (3) give
\[
 F_k(z+e_i)=a(Pe_i,Pz)^nf_k(Pz)e^{-nQ(z+e_i)}=F_k(z),\qquad
 F_k(z+\tau e_j)=a(P\tau e_j,Pz)^nf_k(Pz)e^{-nQ(z+\tau e_j)}=b_j^ne^{-2\pi inz_j}F_k(z).
\]
So $N\le n^g$ by Lemma~\ref{is-eigenvalue-conjugate-frame-theta-product-unique-fourier}.
\end{proof}
% card_le_pow_of_linearIndependent_theta_principal (CardLePowOfLinearIndependentThetaPrincipal.lean), stated with the
% normal form and the Fourier bound as hypotheses H1 and H2, which proportional_of_selfDual supplies
% (SiegelNormalForm.exists_normal_form, ClassicalTheta.card_le_pow_of_linearIndependent); followed: T, hT, the two laws.

\begin{lemma}\label{is-eigenvalue-conjugate-frame-theta-product-unique-proportional}
Let $V$ be a finite-dimensional $\mathbf Q$-vector space with a $\mathbf Q$-bilinear form $B$, $g\ge0$, $p\colon V\to\mathbf
C^g$ an injective additive map, and $L\subset V$ a finitely generated subgroup such that $x\in V$ lies in $L$ if and only
if $B(x,y)\in\mathbf Z$ for every $y\in L$, and such that $p(L)$ spans $\mathbf C^g$ over $\mathbf R$. Let $H\in
M_g(\mathbf C)$ be a Hermitian matrix with $\operatorname{Re}(z^{\mathsf T}H\bar z)>0$ for $z\ne0$ and
$\operatorname{Im}\big(p(u)^{\mathsf T}H\overline{p(v)}\big)=B(u,v)$ for $u,v\in V$, and $c$ a function on $V$ with
$|c(v)|=1$. Then any two holomorphic functions $f_1,f_2$ on $\mathbf C^g$ with
\[
 f_k(z+p(v))=c(v)\exp\Big(\pi z^{\mathsf T}H\overline{p(v)}+\frac\pi2p(v)^{\mathsf T}H\overline{p(v)}\Big)f_k(z)
\]
for all $k$, $v\in L$ and $z\in\mathbf C^g$ are linearly dependent.
\end{lemma}
\leanref{Development/AbelianVarietyTheory/Proofs/SiegelThetaTransformation/LevelOneThetaEqSmulOfSelfDualLattice.lean}{GlobalTheory.Dict.ThetaOne.proportional_of_selfDual}
\begin{proof}
By the hypothesis on $L$, $B(x,y)\in\mathbf Z$ for $x,y\in L$.

(i) If $x\in V$ and $B(x,y)=0$ for every $y\in L$, then $x=0$: then $qx\in L$ for every $q\in\mathbf Q$; as $L$ is a
finitely generated torsion-free group, it is free, and the coordinates of $x$ in a basis are divisible by every integer,
so they are $0$.

(ii) Every homomorphism $\psi\colon L\to\mathbf Z$ is $B(x,\cdot)$ for some $x\in L$. Let $b_1,\dots,b_s$ be a basis of $L$
and $G=(B(b_i,b_j))$, an integer matrix. If $vG=0$ for $v\in\mathbf Z^s$, then $x=\sum_iv_ib_i$ has $B(x,\cdot)=0$ on $L$,
so $x=0$ by (i) and $v=0$; so $\det G\ne0$. Put $c'=(\psi(b_j))_jG^{-1}\in\mathbf Q^s$ and $x=\sum_ic'_ib_i$. Then
$B(x,b_j)=\psi(b_j)$ for every $j$, so $B(x,y)=\psi(y)\in\mathbf Z$ for $y\in L$, and $x\in L$ by the hypothesis on $L$.

Let $\Lambda=p(L)$, a subgroup of $\mathbf C^g$ that spans it over $\mathbf R$, and, as $p$ is injective, let $E$ be the
$\mathbf Z$-bilinear form on $\Lambda$ with $E(p(u),p(v))=B(u,v)=\operatorname{Im}\big(p(u)^{\mathsf T}H\overline{p(v)}\big)$. By
(i) and (ii), $l\mapsto E(l,\cdot)$ is a bijection from $\Lambda$ onto $\operatorname{Hom}(\Lambda,\mathbf Z)$.

(iii) $\Lambda$ is discrete. Let $T$ be a finite set of generators of $L$, and $U$ the open neighbourhood of $0$ of the
$z$ with $|\operatorname{Im}(z^{\mathsf T}H\overline{p(s)})|<1$ for $s\in T$. If $x\in L$ and $p(x)\in U$, then $B(x,s)$ is
an integer of absolute value less than $1$, hence $0$, for $s\in T$; so $B(x,\cdot)=0$ on $L$ and $x=0$ by (i). So
$U\cap\Lambda=\{0\}$.

With $\chi(p(v))=c(v)$, the functions $f_1,f_2$ satisfy the law of
Lemma~\ref{is-eigenvalue-conjugate-frame-theta-product-unique-count} for $n=1$. If they were linearly independent, that
lemma would give $2\le1^g$. So they are linearly dependent.
\end{proof}
% proportional_of_selfDual (LevelOneThetaEqSmulOfSelfDualLattice.lean), followed: floor_bilinForm_eq_of_mem_selfDual,
% eq_zero_of_bilinForm_eq_zero_on_selfDual (i), exists_mem_bilinForm_eq_of_selfDual (ii), Ef and hbij,
% discreteTopology_map_of_integral_bilinForm (iii), and card_le_pow_of_linearIndependent_theta_principal with n = 1.

\begin{lemma}\label{is-eigenvalue-conjugate-frame-theta-product-unique}
Let $X$, $n$, $w_0$, $L$ be as in Lemma~\ref{is-eigenvalue-conjugate-frame-vacua}, of which $n$, the integrality of $w_0$
and $w_0+\bar w_0=1$ are not used, and $x\in\mathcal D_X$. If $\vartheta$ and $\vartheta'$ are level-one theta functions
for $w_0$, $x$ and $L$ (Definition~\ref{is-eigenvalue-conjugate-frame-heisenberg}(5)) and $\vartheta'\ne0$, then
$\vartheta=a\vartheta'$ for some $a\in\mathbf C$.
\end{lemma}
\leanref{Development/AbelianVarietyTheory/Proofs/SiegelThetaTransformation/LevelOneThetaEqSmulOfSelfDualLatticeStatement.lean}{GlobalTheory.T.levelOne_theta_eq_smul_of_selfDual_lattice}
\leanref{Development/AbelianVarietyTheory/Proofs/SiegelThetaTransformation/LevelOneThetaEqSmulOfSelfDualLattice.lean}{GlobalTheory.Dict.ThetaOne.exists_eq_smul_of_isTheta1}
\begin{proof}
Write $\langle\ ,\ \rangle=\langle\ ,\ \rangle_X$, and let $H$ be the Hermitian matrix with $H_x(z,w)=z^{\mathsf T}H\bar w$,
namely $H=(\Omega_x^{-1})^{\mathsf T}(\operatorname{Im}\tau_x)^{-1}\overline{\Omega_x^{-1}}$. By
Lemma~\ref{is-eigenvalue-conjugate-frame-hermitian}(1), (2) and~(4), $H$ is Hermitian with
$\operatorname{Re}(z^{\mathsf T}H\bar z)>0$ for $z\ne0$, $\operatorname{Im}\big(p_x(u)^{\mathsf T}H\overline{p_x(v)}\big)=\langle
u,v\rangle$ for $u,v\in E^r$, and $p_x$ is injective; it is additive. The module $L$ spans $E^r$ over $E$: otherwise, as
$\langle\ ,\ \rangle$ is nondegenerate on $E^r$, some $v\ne0$ has $\langle v,y\rangle=0$ for every $y\in L$, so $qv\in L$ for
every $q\in\mathbf Q$ by self-duality, which is impossible in the finitely generated module $L$. So, by
Lemma~\ref{is-eigenvalue-conjugate-frame-hermitian}(4), $p_x(L)$ spans $\mathbf C^d$ over $\mathbf R$.

Put $c(v)=e^{-\pi i\langle v,w_0v\rangle}$, of absolute value $1$ as $\langle v,w_0v\rangle$ is rational. Let $\vartheta_1$ be a
level-one theta function for $w_0$, $x$ and $L$, $v\in L$ and $p=p_x(v)$. As $H_x(z+\frac12p,p)=H_x(z,p)+\frac12H_x(p,p)$,
the identity of Definition~\ref{is-eigenvalue-conjugate-frame-heisenberg}(5) reads
\[
 \vartheta_1(z+p)=c(v)\exp\Big(\pi z^{\mathsf T}H\bar p+\frac\pi2p^{\mathsf T}H\bar p\Big)\vartheta_1(z).
\]
So $\vartheta$ and $\vartheta'$ satisfy the law of Lemma~\ref{is-eigenvalue-conjugate-frame-theta-product-unique-proportional}
for $V=E^r$, $B=\langle\ ,\ \rangle$, $g=d$, $p=p_x$, the group $L$, which is finitely generated and has the property
required there by self-duality, $H$ and $c$; and they are holomorphic. By that lemma there are $(s,t)\ne(0,0)$ with
$s\vartheta+t\vartheta'=0$. If $s=0$, then $t\ne0$ and $\vartheta'=0$, which is excluded. So $\vartheta=a\vartheta'$ with
$a=-t/s$.
\end{proof}
% GlobalTheory.T.levelOne_theta_eq_smul_of_selfDual_lattice = Dict.ThetaOne.exists_eq_smul_of_isTheta1
% (LevelOneThetaEqSmulOfSelfDualLattice.lean), followed: hx_laws (hsymm, hposx, him, hinj), TensLat.span_of_selfD,
% hx_eq_hermForm (Hm = ThFr.Hm), TensLat.EF_eq_bilin (B), hadd, span_image_pv, c, hlaw, proportional_of_selfDual and
% LinearIndependent.pair_iff. The hypotheses w integral and w + cc w = 1 (hw, hw1) are not used, and Even n only
% through hx_laws, which does not use it.

\begin{lemma}\label{is-eigenvalue-conjugate-frame-theta-product}
Let $X$, $n$, $w_0$, $L$ be as in Lemma~\ref{is-eigenvalue-conjugate-frame-vacua}, $v$ a vacuum of
$\mathcal O_E^n\otimes L$ for $w_0$, $x\in\mathcal D_X$, and $\vartheta$ a level-one theta function for
$w_0$, $x$ and $L$. Then there is $\kappa\in\mathbf C$ with
\[
 \prod_{k=1}^n\vartheta(z_k)=\kappa\,\theta_v(x;z)\qquad\text{for every }z=(z_1,\dots,z_n)\in(\mathbf C^d)^n,
\]
with the theta series of Definition~\ref{is-eigenvalue-conjugate-frame-fibres}(6).
\end{lemma}
\leanref{Development/AbelianVarietyTheory/Proofs/SiegelThetaTransformation/ProdLevelOneThetaEqVacuumTheta.lean}{GlobalTheory.Dict.VacProd.prod_eq_mul_thA}
\begin{proof}
Write $\langle\ ,\ \rangle=\langle\ ,\ \rangle_X$ and $\epsilon(u)=e^{\pi i\langle u,w_0u\rangle}$, and put $d=|\Psi|r$.

\emph{A symplectic basis.} The form $\langle\ ,\ \rangle$ is alternating, as $h(v,u)=\overline{h(u,v)}$ and
$\bar\kappa_0=-\kappa_0$. As $L$ is self-dual, it is a free $\mathbf Z$-module of rank $2d$ on which $\langle\ ,\
\rangle$ is integral and perfect; so it has a basis $b_1,\dots,b_d,b'_1,\dots,b'_d$ with $\langle
b_k,b'_l\rangle=-\delta_{kl}$, $\langle b'_l,b_k\rangle=\delta_{kl}$ and $\langle b_k,b_l\rangle=\langle
b'_k,b'_l\rangle=0$, by the symplectic basis theorem (proof of Lemma~\ref{cv-uniformization-adapted-basis}, with
$b'_l$ replaced by $-b'_l$). Let $c_1$, $c_2$ and $C$ be as in
Lemma~\ref{is-eigenvalue-conjugate-frame-theta-product-vacuum} for this basis, and $h$, $\tau$ and $Y$ as there for $x$.
For $m=(m',m'')\in\mathbf Z^d\times\mathbf Z^d$ put $u_m=\sum_km'_kb_k+\sum_km''_kb'_k$; then $hp_x(u_m)=m'+\tau m''$.

\emph{The signs.} For $u,u'\in L$, $\epsilon(u+u')=\epsilon(u)\epsilon(u')e^{-\pi i\langle u,u'\rangle}$. Indeed,
$h(u,w_0u')=\bar w_0h(u,u')$ and $\langle u',w_0u\rangle=-\langle w_0u,u'\rangle$, so, as $\bar w_0-w_0=1-2w_0$,
$\langle u,w_0u'\rangle+\langle u',w_0u\rangle=\langle u,u'\rangle-2\langle w_0u,u'\rangle$; here $w_0u\in L$ and
$\langle w_0u,u'\rangle\in\mathbf Z$ by self-duality, and $e^{\pi i\langle u,u'\rangle}=e^{-\pi i\langle u,u'\rangle}$ as
$\langle u,u'\rangle\in\mathbf Z$. Put
\[
 \mu(m)=\exp\big(\pi i\,m'\cdot m''+2\pi i(c_1\cdot m'-c_2\cdot m'')\big),\qquad P(m)=\epsilon(u_m)\mu(m).
\]
Then $\mu(m+n)=\mu(m)\mu(n)e^{\pi i(m'\cdot n''+n'\cdot m'')}$ and
$\langle u_m,u_n\rangle=-(m'\cdot n''-m''\cdot n')$, so $P(m+n)=P(m)P(n)e^{2\pi i\,m'\cdot n''}=P(m)P(n)$. At the basis
vectors, $P=\epsilon(b_k)e^{2\pi ic_{1,k}}=e^{4\pi ic_{1,k}}=1$, as $c_{1,k}\in\{0,\frac12\}$, and
$P=\epsilon(b'_k)e^{-2\pi ic_{2,k}}=1$, by Lemma~\ref{is-eigenvalue-conjugate-frame-theta-product-vacuum}(2). Hence
$\epsilon(u_m)\mu(m)=1$ for every $m$.

\emph{A level-one theta function.} Put $F(z)=\exp\big(\frac\pi2(hz)^{\mathsf T}Y^{-1}(hz)\big)\vartheta[c_1,c_2](\tau,hz)$, a
holomorphic function on $\mathbf C^d$. For $\zeta\in\mathbf C^d$ and $l=m'+\tau m''$, the quasi-periodicity of the theta
function with characteristics gives
\[
 F_0(\zeta+l)=\mu(m)\exp\big(\pi\,\zeta^{\mathsf T}Y^{-1}\bar l+\tfrac\pi2\,l^{\mathsf T}Y^{-1}\bar l\big)F_0(\zeta),\qquad
 F_0(\zeta)=\exp\big(\tfrac\pi2\zeta^{\mathsf T}Y^{-1}\zeta\big)\vartheta[c_1,c_2](\tau,\zeta),
\]
as $\tau$ is symmetric. Let $v'=u_m\in L$ and $p=p_x(v')$, so $hp=l$. By
Lemma~\ref{is-eigenvalue-conjugate-frame-theta-product-vacuum}(1),
$H_x(z+\frac12p,p)=(hz)^{\mathsf T}Y^{-1}\bar l+\frac12l^{\mathsf T}Y^{-1}\bar l$. So
$F(z+p)=\mu(m)e^{\pi H_x(z+\frac12p,p)}F(z)$, and
$\epsilon(v')e^{-\pi H_x(z+\frac12p,p)}F(z+p)=\epsilon(u_m)\mu(m)F(z)=F(z)$. So $F$ is a level-one theta function for
$w_0$, $x$ and $L$. It is not $0$, as $h$ is invertible and $\vartheta[c_1,c_2](\tau,\cdot)$ is not identically $0$.

\emph{Conclusion.} By Lemma~\ref{is-eigenvalue-conjugate-frame-theta-product-unique}, $\vartheta=aF$ for some $a\in\mathbf C$.
By Lemma~\ref{is-eigenvalue-conjugate-frame-theta-product-vacuum}(3), $\theta_v(x;z)=C\det(h)^{n/2}\prod_kF(z_k)$, with
$C\det(h)^{n/2}\ne0$. So $\prod_k\vartheta(z_k)=a^n\prod_kF(z_k)=\kappa\,\theta_v(x;z)$ with
$\kappa=a^n\big(C\det(h)^{n/2}\big)^{-1}$.
\end{proof}
% VacProd.prod_eq_mul_thA (92 lines; ProdLevelOneThetaEqVacuumTheta.lean, 282 lines), followed: VacNorm.exists_sympl_of_selfD
% for the basis b (the book cites the symplectic basis theorem); periodCoordinateTheta_of_signedHeisenbergVacuum (Lemma
% -theta-product-vacuum); F = CanTh.canTh tau a1 a2 (h z), with CanTh.canTh_add_period (canMult = mu) and
% thetaSign_mul_canonicalTheta_multiplier (the signs: thetaSignExpEF_add_selfDual_lattice, EF_sum, canMult_add_add, the
% basis values); CanTh.canTh_ne_zero; levelOne_theta_eq_smul_of_selfDual_lattice (Lemma -theta-product-unique).

\begin{lemma}\label{is-eigenvalue-conjugate-frame-vacuum-line-idempotents}
Let $D$ be a definite-side datum with host field $E_h/F_h$. Let $\mathcal O\subset\mathbf A_{E_h}$ be the subgroup of the
elements with archimedean component $0$ whose finite component is integral at every finite place, and $1_f\in\mathcal O$
the element with finite component $1$. A \emph{split idempotent} is an $e\in\mathbf A_{E_h}$ with $e^2=e$, $e\bar e=0$,
and $e+\bar e$ equal to the element with component $1$ at one finite place $v$ of $F_h$ and $0$ at all other places.
Let $w_0\in E_h$ be an algebraic integer with $w_0+\bar w_0=1$ and $w_0\bar w_0=2b$ for an algebraic integer $b$, and let
$\omega\in\mathcal O$ have finite component $w_0$. Then there is a finite set $s$ of split idempotents such that:
\begin{enumerate}
\item $ee'=0$ and $e\bar e'=0$ for distinct $e,e'\in s$;
\item $\omega-\sum_{e\in s}e\in2\mathcal O$ and $1_f-\sum_{e\in s}(e+\bar e)\in2\mathcal O$;
\item every split idempotent $e$ lies in $\mathcal O$, and $e\in s$, $\bar e\in s$ or $e\in2\mathcal O$.
\end{enumerate}
\end{lemma}
\leanref{Development/AbelianVarietyTheory/Proofs/SiegelThetaTransformation/ExistsSplitIdempotentsOfRelTraceOneModTwoStatement.lean}{GlobalTheory.T.exists_splitIdempotents_of_relTraceOneModTwo}
\begin{proof}
A split idempotent $e$ for $v$ is the idempotent of a place of $E_h$ above $v$, and $v$ splits in $E_h$: at the places
other than $v$ the components of $e$ and $\bar e$ are idempotents with sum $0$, hence $0$; and at $v$ the component of
$e$ is an idempotent of $E_h\otimes_{F_h}F_{h,v}$ whose conjugate is complementary to it, which is impossible if
$E_h\otimes_{F_h}F_{h,v}$ is a field and leaves the two idempotents of the factors if $v$ splits.

Let $v$ be a place of $F_h$ above $2$ and $P$ a place of $E_h$ above $v$. As $w_0\bar w_0\in2\mathcal O_{E_h}\subset P$,
one of $w_0$ and $\bar w_0$ lies in $P$, and not both, as $w_0+\bar w_0=1$. If $\bar P=P$, then $w_0\in P$ exactly when
$\bar w_0\in P$; so $\bar P\ne P$, and $v$ splits. Let $P_v$ be the place above $v$ at which $w_0$ is a unit; then
$\bar w_0$ is a unit at $\bar P_v$. Let $s$ be the set of the idempotents $e_v$ of the places $P_v$, $v\mid2$. They are
split idempotents at distinct places, which gives (1).

(2) At a finite place $Q$ of $E_h$ not above $2$, $2$ is a unit, so every integral element is $2$ times an integral
element. At $P_v$, $\omega-\sum_ee$ has the component $w_0-1=-\bar w_0$, and $v_{P_v}(\bar w_0)=v_{P_v}(w_0\bar w_0)\ge
v_{P_v}(2)$, as $w_0$ is a unit there. At $\bar P_v$ it has the component $w_0$, and $v_{\bar P_v}(w_0)=v_{\bar
P_v}(w_0\bar w_0)\ge v_{\bar P_v}(2)$. So $\omega-\sum_ee\in2\mathcal O$. The element $1_f-\sum_e(e+\bar e)$ has the
components $0$ above $2$ and $1$ at the other finite places, so it lies in $2\mathcal O$.

(3) The components of a split idempotent are $0$ or $1$, so it lies in $\mathcal O$. If it is the idempotent of a place
above $v\mid2$, it is $e_v$ or $\bar e_v$; otherwise its nonzero component is at a place where $2$ is a unit, and it lies
in $2\mathcal O$.
\end{proof}
% T.exists_splitIdempotents_of_relTraceOneModTwo = IdemModTwo.exists_split_idempotents (IdemTwo w: w * cc w = 2 b with b
% integral). intAE D is the book's O. Size of the Lean proof: about 550 lines.

\begin{lemma}\label{is-eigenvalue-conjugate-frame-vacuum-line-signs}
In Situation~\ref{is-eigenvalue-conjugate-frame-context}, let $w_0\in E_h$ be an algebraic integer with
$w_0+\bar w_0=1$ and $w_0\bar w_0=2b$ for an algebraic integer $b$, let $M\ge0$ and $\Lambda_S\in\mathcal V_M$
(Lemma~\ref{is-eigenvalue-conjugate-frame-vacua}), and let $L\subset E_h^r$ be an $\mathcal O_{E_h}$-submodule
self-dual for $X'$. Let $\hat\Lambda(\Lambda_S\otimes L)\subset\mathbb W$ be the subgroup generated by the
elements with archimedean component $0$ whose finite component is $\alpha y^{\mathsf T}$, with $\alpha$ in the
closure of the integers of $F_h$ in $\mathbf A_{F_h,f}$ and $y\in\Lambda_S\otimes L$; for $y\in M_{n\times
r}(E_h)$ let $y_f\in\mathbb W$ have archimedean component $0$ and finite component $y^{\mathsf T}$. Then a
theta function $f$ for $(S,T)$ is a lattice vacuum for $\hat\Lambda(\Lambda_S\otimes L)$ if and only if
$\varepsilon_{w_0}(y)\,\rho(y_f)f=f$ for every $y\in\Lambda_S\otimes L$, with the sign $\varepsilon_{w_0}$ of
$X'$ (Definition~\ref{is-eigenvalue-conjugate-frame-heisenberg}(2)).
\end{lemma}
\leanref{Development/AbelianVarietyTheory/Proofs/SiegelThetaTransformation/IsLatticeVacuumIffSignedRationalFixedStatement.lean}{GlobalTheory.T.isLatticeVacuum_iff_signedRational_fixed}
\leanref{Development/AbelianVarietyTheory/Proofs/SiegelThetaTransformation/IsLatticeVacuumIffSignedRationalFixed.lean}{GlobalTheory.Ker.SignVac.isLatVac_iff}
\begin{proof}
Write $A=\Lambda_S\otimes L$, $H=\hat\Lambda(A)$, $\varepsilon=\varepsilon_{w_0}$, and $\mathcal O$, $1_f$ and the split
idempotents as in Lemma~\ref{is-eigenvalue-conjugate-frame-vacuum-line-idempotents}; the lattice vacuum conditions ask
$\rho(2u)f=f$ for $u\in H$, and $\rho(u)f=f$ for $u\in H$ with $eu=u$ for a split idempotent $e$. For $a\in\mathbf
A_{E_h}$ and $x,x'\in\mathbb W$, $\langle ax,x'\rangle=\langle x,\bar ax'\rangle$, by Definition~\ref{th-weil-datum}. The
group $A$ is an $\mathcal O_{E_h}$-module, so $\mathcal OH\subset H$; and $1_fu=u$ for $u\in H$.

\emph{The pairing on $H$.} $\psi(\langle u,u'\rangle)=1$ for $u,u'\in H$. Both sides are multiplicative in each variable,
so it suffices to take $u=\alpha y_f$ and $u'=\beta y'_f$ with $y,y'\in A$ and $\alpha,\beta$ integral finite ideles of
$F_h$. Then $\langle u,u'\rangle$ is the element with archimedean component $0$ and finite component $\alpha\beta t$, where
$t\in F_h$ is the finite component of $\langle y_f,y'_f\rangle$. For $a\in\mathcal O_{F_h}$, $ay\in A$, and
$\psi(\langle(ay)_f,y'_f\rangle)=\psi(\frac12\langle(ay)_f,y'_f\rangle)^2=e^{-2\pi i\sum_k\langle ay_k,y'_k\rangle_{X'}}=1$, by
Lemma~\ref{is-eigenvalue-conjugate-frame-zak-forms}(3) and Lemma~\ref{is-eigenvalue-conjugate-frame-vacua-tensor}(2). So
$\psi$ is $1$ at the elements with finite component $at$, $a\in\mathcal O_{F_h}$, and by continuity at those with finite
component $\gamma t$ for $\gamma$ in the closure of $\mathcal O_{F_h}$, which contains $\alpha\beta$.

\emph{The element $\omega$.} Let $\omega\in\mathcal O$ have finite component $w_0$. Then $\omega H\subset H$, and
$\omega u+\bar\omega u=1_fu=u$ for $u\in H$. For $u\in H$ put $\operatorname{sg}(u)=\psi(\frac12\langle u,\omega u\rangle)$.
\begin{enumerate}
\item[(a)] For $y\in A$, $\operatorname{sg}(y_f)=\varepsilon(y)$: as $\omega y_f=(w_0y)_f$,
  Lemma~\ref{is-eigenvalue-conjugate-frame-zak-forms}(3) gives $\operatorname{sg}(y_f)=e^{-\pi i\sum_k\langle
  y_k,w_0y_k\rangle_{X'}}=\varepsilon(y)^{-1}$, which is $\varepsilon(y)$ as $\varepsilon(y)^2=1$
  (Lemma~\ref{is-eigenvalue-conjugate-frame-vacua-tensor}(3)).
\item[(b)] For $u,u'\in H$, $\operatorname{sg}(u+u')=\operatorname{sg}(u)\operatorname{sg}(u')\psi(\frac12\langle
  u,u'\rangle)$. Indeed, the cross terms are $\frac12\langle u,\omega u'\rangle+\frac12\langle u',\omega u\rangle$, and
  $\langle u',\omega u\rangle=-\langle\omega u,u'\rangle=-\langle u,\bar\omega u'\rangle$; with $\bar\omega u'=u'-\omega u'$ the
  cross terms are $\frac12\langle u,u'\rangle-\langle u,\bar\omega u'\rangle$, and $\psi(\langle u,\bar\omega u'\rangle)=1$ as
  $\bar\omega u'\in H$.
\item[(c)] As $\frac12\langle2u,2\omega u\rangle=2\langle u,\omega u\rangle$, $\operatorname{sg}(2u)=\psi(\langle
  u,\omega u\rangle)^2=1$ for $u\in H$. If $\omega u=u+2v$ or $\omega u=2v$ with $v\in H$,
  then $\operatorname{sg}(u)=\psi(\langle u,v\rangle)=1$, as $\langle u,u\rangle=0$.
\end{enumerate}

\emph{A group.} Let $G$ be the set of the $u\in H$ with $\operatorname{sg}(u)\rho(u)f=f$. It is a subgroup: $0\in G$; for
$u,u'\in G$, $\rho(u)\rho(u')f=\psi(\frac12\langle u,u'\rangle)\rho(u+u')f$
(Lemma~\ref{is-eigenvalue-conjugate-frame-zak-space}(2)), so by (b)
$\operatorname{sg}(u+u')\rho(u+u')f=\operatorname{sg}(u)\operatorname{sg}(u')\rho(u)\rho(u')f=f$; and (b) for $-u$ and $u$
gives $\operatorname{sg}(-u)\operatorname{sg}(u)\psi(-\frac12\langle u,u\rangle)=\operatorname{sg}(0)=1$, from which
$\operatorname{sg}(-u)\rho(-u)f=f$ follows in the same way.

\emph{The halves.} Let $s$ be as in Lemma~\ref{is-eigenvalue-conjugate-frame-vacuum-line-idempotents} for $w_0$, with
$\omega=\sum_{e'\in s}e'+2a$ and $1_f=\sum_{e'\in s}(e'+\bar e')+2c$, $a,c\in\mathcal O$. Let $e$ be a split idempotent and
$u\in H$ with $eu=u$. Then $\omega u=(\sum_{e'\in s}e'e)u+2au$. If $e\in s$, then $\sum_{e'}e'e=e$ and $\omega u=u+2au$; if
$\bar e\in s$, then every $e'e$ is $0$ and $\omega u=2au$; and if $e=2c'$ with $c'\in\mathcal O$, then $u=2(c'u)$. In each
case $\operatorname{sg}(u)=1$ by (c), as $au,c'u\in H$.

\emph{Only if.} Let $f$ be a lattice vacuum for $H$, and $y\in A$. For a split idempotent $e$, $ey_f\in H$, $e(ey_f)=ey_f$,
so $\operatorname{sg}(ey_f)=1$ and $\rho(ey_f)f=f$: $ey_f\in G$. Also $cy_f\in H$, $\operatorname{sg}(2cy_f)=1$ and
$\rho(2cy_f)f=f$: $2cy_f\in G$. As $y_f=1_fy_f=\sum_{e\in s}(ey_f+\bar ey_f)+2cy_f$ and $\bar e$ is a split idempotent, $y_f\in
G$, which is $\varepsilon(y)\rho(y_f)f=f$ by (a).

\emph{If.} Suppose $\varepsilon(y)\rho(y_f)f=f$ for every $y\in A$, so that $y_f\in G$ by (a). The module $A$ is finitely
generated and spans $M_{n\times r}(E_h)$, as $\Lambda_S$ and $L$ span $E_h^n$ and $E_h^r$ (proof of
Lemma~\ref{is-eigenvalue-conjugate-frame-vacua-tensor}). Let $U$ be an open subgroup for $f$ as in
Definition~\ref{th-weil-datum}(2), and $N$ as in Lemma~\ref{is-eigenvalue-conjugate-frame-zak-exists-lattice}(4) for
$\Lambda_S$, $L$ and $U$. Every $u\in H$ is $y_f+2Nu'$ with $y\in A$ and $u'\in H$, by
Lemma~\ref{is-eigenvalue-conjugate-frame-zak-exists-lattice}(3). Here $2Nu'=N(2u')\in U$ has archimedean component $0$, so
$\rho(2Nu')f=f$, and $\operatorname{sg}(2Nu')=1$ by (c); so $2Nu'\in G$, and $u\in G$. Hence
$G=H$. For $u\in H$, $2u\in G$ and $\operatorname{sg}(2u)=1$ give $\rho(2u)f=f$; and for $u\in H$ with $eu=u$ for a split
idempotent $e$, $\operatorname{sg}(u)=1$ and $u\in G$ give $\rho(u)f=f$. So $f$ is a lattice vacuum for $H$.
\end{proof}
% Ker.SignVac.isLatVac_iff (IsLatticeVacuumIffSignedRationalFixed.lean), followed: TensLat.smul_mem_tens (hA),
% char_sympl_hatA with TensLat.tens_selfDual and IntDense.addChar_eq_one_of_integral (the pairing on H), omega = pairA 0 1
% tensor w (hw, hw'), sg with heps, sg_add, sg_two, sg_of_cong, fixGroup, T.exists_splitIdempotents_of_relTraceOneModTwo
% (Lemma -vacuum-line-idempotents) and hhalf; the converse uses h3 and h4 of
% T.completedTensorLattice_strongApprox_open_compact for tens Lambda L (Lemma -zak-exists-lattice(3) and (4)).

\begin{lemma}\label{is-eigenvalue-conjugate-frame-vacuum-line-series}
Let $X$, its symplectic basis and indexing be as in Definition~\ref{is-eigenvalue-conjugate-frame-heisenberg}, with
module $\Lambda$, let $n\ge0$ be even (the parity is not used), and let $U$, $\hat y$ and $E_X$ be as in
Lemma~\ref{is-eigenvalue-conjugate-frame-zak-coordinates}. Let $x\in\mathcal D_X$, $N\ge1$ an integer, and $\Phi$ a
holomorphic function on $(\mathbf C^d)^n$ such that
\[
 e^{-\pi iE_X(v,N\hat y)}\,\mathcal G_x\big(z_x(v+N\hat y)\big)\,\Phi\big(z_x(v+N\hat y)\big)=\mathcal G_x\big(z_x(v)\big)\,
 \Phi\big(z_x(v)\big)
\]
for every $v\in U$ and every $y\in M_{n\times r}(E)$ whose rows lie in $\Lambda$. Then there is a Schwartz-Bruhat function
$\varphi$ on $M_{n\times d}(\mathbf Q)$ with $\theta_\varphi(x;\cdot)=\Phi$.
\end{lemma}
\leanref{Development/AbelianVarietyTheory/Proofs/SiegelThetaTransformation/ExistsThetaEqOfTranslationInvariantMultipleStatement.lean}{GlobalTheory.T.exists_theta_eq_of_translation_invariant_multiple}
\leanref{Development/AbelianVarietyTheory/Proofs/SiegelThetaTransformation/ExistsThetaEqOfTranslationInvariantMultiple.lean}{GlobalTheory.Ker.Onto.exists_thA_eq}
\begin{proof}
\openstep{The proof, which turns the hypothesis into the quasi-periodicity of $\Phi$ under $N$ times the periods of
$\Lambda$ in the normalized coordinates of the symplectic basis and expands such a holomorphic function as a finite
combination of theta series with rational characteristics, is not written out.}
% Ker.Onto.exists_thA_eq: zR_laws (Lemma -zak-coordinates), Dict.hx_laws, key (the quasi-periodicity in z, with
% exponent_eq), and exists_canonical at hMat, tauMat and imInv (the Fourier expansion). Size, following imports and
% stopping at statements the book cites: about 1,460 lines.
\end{proof}

\begin{lemma}\label{is-eigenvalue-conjugate-frame-vacuum-line-fock}
In Situation~\ref{is-eigenvalue-conjugate-frame-context}, with the notation of
Lemma~\ref{is-eigenvalue-conjugate-frame-zak}, let $B$ be orthonormal for $X'$, and let $f$ be a theta function
for $(S,T)$ that is annihilated by the annihilation operators for the Cartan involutions $1$ of $S$ and
$s'_{B,\rho}$ of $T$ at every $\rho$. Then $f=\mathrm{Th}(x'_B,\varphi)$ for a Schwartz-Bruhat function $\varphi$.
\end{lemma}
\leanref{Development/AbelianVarietyTheory/Proofs/SiegelThetaTransformation/ExistsTestFunctionOfFockVacuumThetaFunctionStatement.lean}{GlobalTheory.T.exists_testFunction_of_fockVacuum_thetaFunction}
\leanref{Development/AbelianVarietyTheory/Proofs/SiegelThetaTransformation/ExistsTestFunctionOfFockVacuumThetaFunction.lean}{GlobalTheory.Ker.exists_test_of_fock}
\begin{proof}
Put $x=x'_B$, which lies in $\mathcal D_{X'}$ (Lemma~\ref{is-eigenvalue-conjugate-frame-zak}(3)), and write $z=z'_x$, $\mathcal
G=\mathcal G'_x$. Let $R$ and $u'$ be as in Lemma~\ref{is-eigenvalue-conjugate-frame-zak-points}; by
Lemma~\ref{is-eigenvalue-conjugate-frame-zak-coordinates}(1) for $X'$, $z$ is an isomorphism of real vector spaces from
$U'$ onto $(\mathbf C^d)^n$.

\emph{The function $\Phi$.} Put $\Phi(\zeta)=\mathcal G(\zeta)^{-1}f\big(R(z^{-1}(\zeta))\big)$. As $R\circ z^{-1}$ is
continuous and additive with finite component $0$, the theta function $f$ is smooth along it
(Definition~\ref{th-weil-datum}(2)), and $\mathcal G$ is smooth and nowhere $0$; so $\Phi$ is differentiable as a function
of real variables. For $v\in U'$, $f(R(v))=\mathcal G(z(v))\Phi(z(v))$; and for $w$ with finite component $0$,
$R(u'(w))=w$ (Lemma~\ref{is-eigenvalue-conjugate-frame-zak-points}(2)), so $f(w)=\mathcal G(z(u'(w)))\Phi(z(u'(w)))$.

\emph{Holomorphy.} Let $\zeta\in(\mathbf C^d)^n$, $\rho\in\mathcal T$, $z_0$ a complex $r\times n$ matrix, and $u_0$ the
coordinates of $z_0$ concentrated at $\rho$ (Lemma~\ref{is-eigenvalue-conjugate-frame-zak}). Take $w=R(z^{-1}(\zeta))$, which
has finite component $0$ and $z(u'(w))=\zeta$. By Lemma~\ref{is-eigenvalue-conjugate-frame-zak-fock}(1), the annihilation
operator at $\rho$ and $z_0$, applied to $f$, has the value $\mathcal G(\zeta)\big(d\Phi_\zeta(u_0)+i\,d\Phi_\zeta(iu_0)\big)$
at $w$, and this is $0$ by the hypothesis on $f$. So $d\Phi_\zeta(iu_0)=i\,d\Phi_\zeta(u_0)$. Every $u\in(\mathbf C^d)^n$ is
$z(v)$ for some $v\in U'$, and $v$ is the sum of its parts concentrated at the $\bar\rho$, each of which is the family
attached to some $z_0$; so $u$ is a sum of such $u_0$, and $d\Phi_\zeta$ is complex linear. Hence $\Phi$ is holomorphic.

\emph{Translations.} By Lemma~\ref{is-eigenvalue-conjugate-frame-zak-space}(5) there is $N\ge1$ with $\rho(Nu)f=f$ for every
$u\in\hat\Lambda$, the group of that lemma. Let $y\in M_{n\times r}(E_h)$ have its rows in $\Lambda'$, so that $y\in\mathcal
O_{E_h}^n\otimes\Lambda'$ (Lemma~\ref{is-eigenvalue-conjugate-frame-vacua-tensor}(4)) and $(-Ny)_f=N(-y)_f\in N\hat\Lambda$.
Let $v\in U'$. By Lemma~\ref{is-eigenvalue-conjugate-frame-zak-space}(3) for $-Ny$ at $w=R(v)$,
\[
 f(R(v))=\big(\rho((-Ny)_f)f\big)(R(v))=e^{-\pi iE_{X'}(v,N\hat y)}f\big(R(v)+(Ny)^{\mathsf T}-(Ny)_f\big),
\]
and $R(v)+(Ny)^{\mathsf T}-(Ny)_f=R(v+N\hat y)$, as both have finite component $0$ and the same image under $u'$
(Lemma~\ref{is-eigenvalue-conjugate-frame-zak-forms}(8) and Lemma~\ref{is-eigenvalue-conjugate-frame-zak-points}). In terms of
$\Phi$ this is the hypothesis of Lemma~\ref{is-eigenvalue-conjugate-frame-vacuum-line-series} for $X'$, $x$ and $N$. So
there is a Schwartz-Bruhat $\varphi$ with $\theta'_\varphi(x;\cdot)=\Phi$.

\emph{Conclusion.} For $w$ with finite component $0$, $f(w)=\mathcal G(z(u'(w)))\theta'_\varphi(x;z(u'(w)))=\Theta'_\varphi(x;u'(w))$.
So $f$ is a Zak theta function for $x$ and $\varphi$, and $f=\mathrm{Th}(x,\varphi)$ by
Lemma~\ref{is-eigenvalue-conjugate-frame-zak}(1).
\end{proof}
% Ker.exists_test_of_fock (ExistsTestFunctionOfFockVacuumThetaFunction.lean), followed: T.realPoints_classicalModel_realVectorSpace
% (R, hR, hR'; Lemma -zak-points), thetaKernel_orthonormalFrame_inDomain_periodLinear (Lemma -zak(3)), Ker.zR_laws (L; Lemma
% -zak-coordinates), Phi = gaussR^-1 * f o R o L^-1 (differentiable_comp), hw, hone from the first clause of
% T.annihilation_creation_periodCoordinates (Lemma -zak-fock(1)), differentiable_of_fderiv_I, the level and rho_finW of
% T.thetaFunction_space_heisenberg_laws (Lemma -zak-space), T.exists_theta_eq_of_translation_invariant_multiple (Lemma
% -vacuum-line-series) and T.adelicZakTheta_exists_unique (Lemma -zak(1)).

\begin{lemma}\label{is-eigenvalue-conjugate-frame-vacuum-line}
In Situation~\ref{is-eigenvalue-conjugate-frame-context}, with the notation of
Lemma~\ref{is-eigenvalue-conjugate-frame-zak}, let $w_0\in E_h$ be an algebraic integer with
$w_0+\bar w_0=1$ and $w_0\bar w_0=2b$ for an algebraic integer $b$, let $M\ge0$ and $\Lambda_S\in\mathcal V_M$
(Lemma~\ref{is-eigenvalue-conjugate-frame-vacua}), and let $L\subset E_h^r$ be an $\mathcal
O_{E_h}$-submodule self-dual for $X'$. Let $\hat\Lambda(\Lambda_S\otimes L)$ be the subgroup of $\mathbb W$
generated by the elements with archimedean component $0$ whose finite component is $\alpha y^{\mathsf T}$,
with $\alpha$ in the closure of the integers of $F_h$ in $\mathbf A_{F_h,f}$ and $y\in\Lambda_S\otimes L$.
\begin{enumerate}
\item For $x'\in\mathcal D_{X'}$ and a Schwartz-Bruhat function $\varphi$, $\mathrm{Th}(x',\varphi)$ is a
  lattice vacuum for $\hat\Lambda(\Lambda_S\otimes L)$ if and only if
  $\varepsilon_{w_0}(y)\,\rho_{X'}(y)\varphi=\varphi$ for every $y\in\Lambda_S\otimes L$.
\item Let $B$ be orthonormal for $X'$, $v$ a vacuum of $\Lambda_S\otimes L$ for $X'$ and $w_0$, and $f$ a
  theta function for $(S,T)$ that is annihilated by the annihilation operators for the Cartan involutions $1$
  of $S$ and $s'_{B,\rho}$ of $T$ at every $\rho$ and is a lattice vacuum for $\hat\Lambda(\Lambda_S\otimes
  L)$. Then $f=t\,\mathrm{Th}(x'_B,v)$ for some $t\in\mathbf C$.
\end{enumerate}
\end{lemma}
\leanref{Development/AbelianVarietyTheory/Proofs/SiegelThetaTransformation/LatticeVacuaTensorCorrespondLine.lean}{GlobalTheory.Ker.latVac_th_line}
\leanref{Development/AbelianVarietyTheory/Definitions/SpanningFamilies/IsWeakOfDatum.lean}{GlobalTheory.Ker.hatA}
\begin{proof}
Write $A=\Lambda_S\otimes L$ and $\varepsilon=\varepsilon_{w_0}$. For $y\in M_{n\times r}(E_h)$,
$\varepsilon(-y)=\varepsilon(y)$, as $\langle-y_k,-w_0y_k\rangle_{X'}=\langle y_k,w_0y_k\rangle_{X'}$.

(1) By Lemma~\ref{is-eigenvalue-conjugate-frame-zak}(1), $f=\mathrm{Th}(x',\varphi)$ is a theta function, so by
Lemma~\ref{is-eigenvalue-conjugate-frame-vacuum-line-signs} it is a lattice vacuum for $\hat\Lambda(A)$ exactly
when $\varepsilon(y)\rho(y_f)f=f$ for every $y\in A$. By Lemma~\ref{is-eigenvalue-conjugate-frame-zak}(2),
$\rho((-y)_f)f=\mathrm{Th}(x',\rho_{X'}(y)\varphi)$ and $\varepsilon(y)\,\mathrm{Th}(x',\psi)=\mathrm{Th}(x',\varepsilon(y)\psi)$;
the functions $\rho_{X'}(y)\varphi$ are Schwartz-Bruhat by
Lemma~\ref{is-eigenvalue-conjugate-frame-vacua-heisenberg}(2). Suppose that $f$ is a lattice vacuum, and let
$y\in A$. Applied to $-y\in A$, the condition gives
$f=\varepsilon(-y)\rho((-y)_f)f=\mathrm{Th}(x',\varepsilon(y)\rho_{X'}(y)\varphi)$, and as $\mathrm{Th}(x',\cdot)$ is
injective, $\varepsilon(y)\rho_{X'}(y)\varphi=\varphi$. Conversely, if $\varepsilon(y)\rho_{X'}(y)\varphi=\varphi$ for every
$y\in A$, then for $y\in A$, as $-y\in A$,
$\varepsilon(y)\rho(y_f)f=\mathrm{Th}(x',\varepsilon(-y)\rho_{X'}(-y)\varphi)=\mathrm{Th}(x',\varphi)=f$.

(2) By Lemma~\ref{is-eigenvalue-conjugate-frame-vacuum-line-fock}, $f=\mathrm{Th}(x'_B,\varphi)$ for a
Schwartz-Bruhat function $\varphi$. The point $x'_B$ lies in
$\mathcal D_{X'}$ (Lemma~\ref{is-eigenvalue-conjugate-frame-zak}(3)), so
$\varepsilon_{w_0}(y)\rho_{X'}(y)\varphi=\varphi$ for every $y\in\Lambda_S\otimes L$ by (1), and
$\varphi=\langle\varphi,v\rangle v$ by Lemma~\ref{is-eigenvalue-conjugate-frame-vacua}(2). By
Lemma~\ref{is-eigenvalue-conjugate-frame-zak}(2), $f=\langle\varphi,v\rangle\,\mathrm{Th}(x'_B,v)$.
\end{proof}
% latVac_th_line: the first part is latVac_th_iff; the second is exists_testFunction_of_fockVacuum_thetaFunction for
% phi, thetaKernel_orthonormalFrame for the domain, latVac_th_iff, TensVac.vacua (second clause) and the linearity
% of thetaKernel_testFunction_map.

\begin{lemma}\label{is-eigenvalue-conjugate-frame-torus-mean-unfold}
Let $G$ and $V$ be locally compact Hausdorff abelian groups, $V$ second countable, with Haar measures $\mu$ and
$\nu$. Let $p\colon G\to V$ and $R\colon V\to G$ be continuous homomorphisms with $p\circ R=1$, and $U\subset G$ a
compact subgroup such that the \emph{tube} $T_U=\{w\in G:w-R(p(w))\in U\}$ is open. Let $\Gamma\subset G$ be a
countable subgroup with $\Gamma\cap U=\{0\}$, such that every $w\in G$ has some $\gamma\in\Gamma$ with
$\gamma+w\in T_U$, and put $L=p(\Gamma\cap T_U)$. Let $\mathcal F$ be a fundamental domain of $\Gamma$ in $G$ with
$\mu(\mathcal F)=1$, $F$ a fundamental domain of $L$ in $V$, and $h$ a continuous function on $G$ with
$h(\gamma+w)=h(w)$ and $h(w+u)=h(w)$ for all $\gamma\in\Gamma$, $u\in U$ and $w\in G$. Then
\[
 \int_{\mathcal F}h\,d\mu=\nu(F)^{-1}\int_Fh(R(v))\,d\nu(v).
\]
\end{lemma}
\leanref{Development/AbelianVarietyTheory/Proofs/SiegelThetaTransformation/ThetaInnerProductEqRealTorusMean.lean}{GlobalTheory.Ker.TorusMean.integral_eq_mean}
\leanref{Development/AbelianVarietyTheory/Proofs/SiegelThetaTransformation/ThetaInnerProductEqRealTorusMean.lean}{GlobalTheory.Ker.TorusMean.exists_map_eq_smul}
\begin{proof}
Write $T=T_U$. It is a subgroup, and $T+R(V)=T$.

\emph{The image measure.} Let $\nu'$ be the image under $p$ of the restriction of $\mu$ to $T$. For a compact $K\subset V$,
$p^{-1}(K)\cap T\subset U+R(K)$, which is compact; so $\nu'$ is finite on compact sets. For $v\in V$, translation by
$R(v)$ preserves $T$ and $\mu$ and commutes with $p$ up to translation by $v$; so $\nu'$ is translation invariant. Hence
$\nu'=c\,\nu$ for some $c\ge0$, by the uniqueness of Haar measure \cite{Fol95}. In particular $\mu(T\cap p^{-1}(N))=0$ for every $\nu$-null set $N$.

\emph{A second fundamental domain.} Put $\Phi=T\cap p^{-1}(F)$. It is a fundamental domain of $\Gamma$, up to null sets.
Indeed, let $w\in G$, with $\gamma+w\in T$ for some $\gamma\in\Gamma$; if some $l\in L$ has $l+p(\gamma+w)\in F$, write
$l=p(\gamma')$ with $\gamma'\in\Gamma\cap T$; then $\gamma'+\gamma+w\in T\cap p^{-1}(F)$. The $w$ for which no such $l$
exists lie in the countable union of the translates by $\Gamma$ of $T\cap p^{-1}(N)$, $N$ the null set of the $v$ with no
$l$ with $l+v\in F$; so they form a null set. If $\gamma_1\ne\gamma_2$ in $\Gamma$ and $\gamma_1+y_1=\gamma_2+y_2$ with
$y_1,y_2\in\Phi$, then $\delta=\gamma_1-\gamma_2=y_2-y_1\in\Gamma\cap T$, as $T$ is a subgroup, and $p(\delta)\ne0$,
as otherwise $\delta=\delta-R(p(\delta))\in\Gamma\cap U=\{0\}$. Then $p(y_2)=p(\delta)+p(y_1)$ lies in
$F\cap(p(\delta)+F)$, a $\nu$-null set as $0\ne p(\delta)\in L$. So the common points of $\gamma_1+\Phi$ and
$\gamma_2+\Phi$ lie in $\gamma_2+\big(T\cap p^{-1}(F\cap(p(\delta)+F))\big)$, a $\mu$-null set.

\emph{The integral.} For $w\in T$, $w=R(p(w))+u$ with $u\in U$, so $h(w)=h(R(p(w)))$. As $h$ is invariant under
$\Gamma$ and $\mathcal F$ and $\Phi$ are fundamental domains,
\[
 \int_{\mathcal F}h\,d\mu=\int_\Phi h\,d\mu=\int_{T\cap p^{-1}(F)}h(R(p(w)))\,d\mu(w)=\int_Fh(R(v))\,d\nu'(v)
 =c\int_Fh(R(v))\,d\nu(v).
\]
For $h=1$ this gives $1=\mu(\mathcal F)=\mu(\Phi)=c\,\nu(F)$, so $c=\nu(F)^{-1}$.
\end{proof}
% integral_eq_mean (135 lines) and exists_map_eq_smul (Measure.isAddLeftInvariant_eq_smul for the image measure, with
% exists_compact for its finiteness on compacts), followed: hfac, hGammaT, the fundamental domain hPhi (ae_covers through
% hcov and hL, aedisjoint through hinj), hint1, hint2 and hmass.

\begin{lemma}\label{is-eigenvalue-conjugate-frame-torus-mean}
In Situation~\ref{is-eigenvalue-conjugate-frame-context}, let $U'$ be the real
vector space of the families $u'=(u'_1,\dots,u'_n)$ in which each $u'_k$ is a family of vectors of
$\mathbf C^r$ indexed by $\Psi'$, with its Lebesgue measure, and for an integer $N>0$ let
$\Gamma'_N\subset U'$ be the subgroup generated by the families $\big(N\sigma(y_k)\big)_{k,\sigma}$ with
$y_1,\dots,y_n\in\Lambda'$. Let $\mu$ be a Haar measure on $\mathbb W$ and $\mathcal F$ a fundamental domain
of $M_{r\times n}(E_h)$ in $\mathbb W$ with $\mu(\mathcal F)=1$, and let $f,f'$ be theta functions for
$(S,T)$. Then there is $N_0>0$ such that for every positive multiple $N$ of $N_0$ the group $\Gamma'_N$ has
a fundamental domain in $U'$ for Lebesgue measure, and for every such fundamental domain $F_N$ and every
function $\Phi$ on $U'$ with $f(w)\overline{f'(w)}=\Phi(u'(w))$ for every $w\in\mathbb W$ with finite component
$0$,
\[
 \int_{\mathcal F}f\,\bar f'\,d\mu=\operatorname{vol}(F_N)^{-1}\int_{F_N}\Phi(u')\,du'.
\]
\end{lemma}
\leanref{Development/AbelianVarietyTheory/Proofs/SiegelThetaTransformation/ThetaInnerProductEqRealTorusMean.lean}{GlobalTheory.Ker.TorusMean.ipO_eq_mean}
\begin{proof}
For $y\in M_{n\times r}(E_h)$ let $y_f\in\mathbb W$ have archimedean component $0$ and finite component $y^{\mathsf T}$, and
let $\hat A=\hat\Lambda(\mathcal O_{E_h}^n\otimes\Lambda')$ be the subgroup of $\mathbb W$ of
Lemma~\ref{is-eigenvalue-conjugate-frame-vacuum-line}; here $\mathcal O_{E_h}^n\otimes\Lambda'$ is the set of the matrices
whose rows lie in $\Lambda'$. As $\mathcal O_{E_h}^n$ and $\Lambda'$ are finitely generated and span, strong approximation
for the additive group of $F_h$ \cite{Wei67} gives:
\begin{enumerate}
\item[(s1)] every $u\in\mathbb W$ with archimedean component $0$ is $y_f+a$ with $y\in M_{n\times r}(E_h)$ and $a\in\hat A$;
\item[(s2)] $y_f\in\hat A$ exactly when $y\in\mathcal O_{E_h}^n\otimes\Lambda'$;
\item[(s3)] the elements of $\hat A$ have archimedean component $0$, and for every integer $N>0$ every $a\in\hat A$ is
  $y_f+Na'$ with $y\in\mathcal O_{E_h}^n\otimes\Lambda'$ and $a'\in\hat A$;
\item[(s4)] for every open subgroup $O$ of $\mathbb W$ there is an integer $N>0$ with $N\hat A\subset O$;
\item[(s5)] for every integer $N>0$ the set of the $w$ such that $w-Na$ has finite component $0$ for some $a\in\hat A$ is
  open, and $\hat A$ is compact.
\end{enumerate}
As $z\mapsto(\theta_\rho(z))_\rho$ is a bijective ring homomorphism from $\mathbf A_{F_h,\infty}\otimes_{F_h}E_h$ onto
$\mathbf C^{\mathcal T}$ (proof of Lemma~\ref{is-eigenvalue-conjugate-frame-glue}), $w\mapsto u'(w)$ is a continuous additive
map $p\colon\mathbb W\to U'$ that vanishes on the elements with archimedean component $0$ and restricts to a real linear
homeomorphism from the elements with finite component $0$ onto $U'$; let $R$ be its inverse. Then $p\circ R=1$, $R(p(w))=w$
when $w$ has finite component $0$, and $w-R(p(w))$ has archimedean component $0$ for every $w$. For a rational point
$y^{\mathsf T}$, $y^{\mathsf T}-y_f$ has finite component $0$, so $R(p(y^{\mathsf T}))=y^{\mathsf T}-y_f$; and $p(y^{\mathsf
T})$ is the family $(\sigma(y_k))_{k,\sigma}$ of the rows $y_k$ of $y$.

\emph{The integer $N_0$.} By Definition~\ref{th-weil-datum}(2) and (s4) there are integers $N_f,N_{f'}>0$ with
$\rho(N_fa)f=f$ and $\rho(N_{f'}a)f'=f'$ for every $a\in\hat A$. Put $N_0=N_fN_{f'}$, let $N$ be a positive multiple of
$N_0$, and put $h=f\bar f'$, a continuous function. Then $h(w+Na)=h(w)$ for $a\in\hat A$, as $\rho(Na)$ fixes $f$ and
$f'$ and the factors $\psi(\frac12\langle w,Na\rangle)$, of absolute value $1$, cancel; and $h(\gamma+w)=h(w)$ for rational
$\gamma$, by the quasi-periodicity of $f$ and $f'$.

\emph{The data of Lemma~\ref{is-eigenvalue-conjugate-frame-torus-mean-unfold}.} Take $G=\mathbb W$ with $\mu$, $V=U'$ with
Lebesgue measure, $p$, $R$, $U=N\hat A$, which is compact by (s5), and $\Gamma$ the group of rational points
$M_{r\times n}(E_h)$, which is countable.
\begin{itemize}
\item The tube $T_U$ is the set of the $w$ such that $w-Na$ has finite component $0$ for some $a\in\hat A$: if $w-Na$ has
  finite component $0$, then $R(p(w))=R(p(w-Na))=w-Na$. So $T_U$ is open by (s5).
\item $\Gamma\cap U=\{0\}$: a rational point in $N\hat A$ has archimedean component $0$ by (s3), so it is $0$.
\item Let $w\in\mathbb W$. By (s1) and (s3), $w-R(p(w))=y_f+y'_f+Na'$ with $y\in M_{n\times r}(E_h)$,
  $y'\in\mathcal O_{E_h}^n\otimes\Lambda'$ and $a'\in\hat A$. For $\gamma=-(y+y')^{\mathsf T}$,
  $R(p(\gamma))=\gamma+(y+y')_f$, so $(\gamma+w)-R(p(\gamma+w))=w-R(p(w))-y_f-y'_f=Na'$, and $\gamma+w\in T_U$.
\item $p(\Gamma\cap T_U)=\Gamma'_N$. A rational point $y^{\mathsf T}$ lies in $T_U$ exactly when
  $y_f=y^{\mathsf T}-R(p(y^{\mathsf T}))\in N\hat A$, that is, by (s2), when $N^{-1}y\in\mathcal O_{E_h}^n\otimes\Lambda'$;
  and then $p(y^{\mathsf T})=(N\sigma(N^{-1}y_k))_{k,\sigma}$. These families form $\Gamma'_N$.
\end{itemize}
The group $\Gamma'_N$ is discrete: for a compact $K\subset U'$, $p^{-1}(K)\cap T_U$ lies in the compact set $U+R(K)$, which
meets the discrete group $\Gamma$ in a finite set, so $\Gamma'_N\cap K$ is finite. Hence $\Gamma'_N$ has a fundamental domain in
$U'$ for Lebesgue measure. Let $F_N$ be one and $\Phi$ as in the statement. For $v\in U'$, $R(v)$ has finite component $0$
and $p(R(v))=v$, so $h(R(v))=\Phi(v)$. Lemma~\ref{is-eigenvalue-conjugate-frame-torus-mean-unfold} gives
\[
 \int_{\mathcal F}f\bar f'\,d\mu=\operatorname{vol}(F_N)^{-1}\int_{F_N}h(R(v))\,dv=\operatorname{vol}(F_N)^{-1}\int_{F_N}\Phi(v)\,dv.
\]
\end{proof}
% Ker.TorusMean.ipO_eq_mean (166 lines; ThetaInnerProductEqRealTorusMean.lean, 544 lines), followed:
% T.realPoints_classicalModel_realVectorSpace (R, its continuity and the clauses used above),
% T.completedTensorLattice_strongApprox_open_compact for stdL n and C.X.Lambda ((s1) to (s5): hA1, hA2, hA3, hA5, hA6),
% the levels Nf, Nf' from T.thetaFunction_space_heisenberg_laws, hstep and hGamma (mul_conj_phase), U = N hatA, hmemT, hO,
% hcov', hinj, hL (mem_latR_iff), hUc, AdelicMatrixLattice.lattice for the topology, discrete_image,
% exists_isAddFundamentalDomain_addSubgroup, and integral_eq_mean (Lemma -torus-mean-unfold).

\begin{lemma}\label{is-eigenvalue-conjugate-frame-tensor-lattice}
In Situation~\ref{is-eigenvalue-conjugate-frame-context}, let $\Lambda_S,\Lambda\subset E_h^n$ be $\mathcal
O_{E_h}$-submodules with $v\in\Lambda$ exactly when $\bar v\in\Lambda_S$, and $L_T,L\subset E_h^r$ $\mathcal
O_{E_h}$-submodules such that $L=\mathfrak AL_T$, the group generated by the products $\alpha t$ with
$\alpha\in\mathfrak A$ and $t\in L_T$. Then the lattice of the pair given by $\Lambda_S$ and $L_T$ is
$\hat\Lambda(\Lambda\otimes L)$ (Lemma~\ref{is-eigenvalue-conjugate-frame-vacuum-line}).
\end{lemma}
\leanref{Development/AbelianVarietyTheory/Proofs/SiegelThetaTransformation/LatOEqHatA.lean}{GlobalTheory.Ker.LatTens.latO_eq_hatA}
\leanref{Development/AbelianVarietyTheory/Definitions/SpanningFamilies/IsWeakOfDatum.lean}{GlobalTheory.Ker.IsAA}
\begin{proof}
Let $\hat{\mathcal O}$ be the closure of the integers of $F_h$ in $\mathbf A_{F_h,f}$. The lattice of the pair given by
$\Lambda_S$ and $L_T$ is the subgroup of $\mathbb W$ generated by the elements with archimedean component $0$ and finite
component $\alpha\,a\,t\,\bar s^{\mathsf T}G_S$, for $\alpha\in\hat{\mathcal O}$, $a\in\mathfrak A$, $s\in\Lambda_S$ and
$t\in L_T$. As $S$ has Gram matrix $G_S=1$, this finite component is $\alpha\,y^{\mathsf T}$ with $y=\bar s(at)^{\mathsf
T}$; write $[\alpha y]$ for the element of $\mathbb W$ with archimedean component $0$ and finite component $\alpha
y^{\mathsf T}$, which is additive in $y$.

Each generator of the lattice of the pair lies in $\hat\Lambda(\Lambda\otimes L)$: here $\bar s\in\Lambda$, as
$s\in\Lambda_S$, and $at\in\mathfrak AL_T=L$, so $y=\bar s(at)^{\mathsf T}\in\Lambda\otimes L$. Conversely, fix
$\alpha\in\hat{\mathcal O}$ and $v\in\Lambda$. The $u\in E_h^r$ for which $[\alpha\,vu^{\mathsf T}]$ lies in the lattice
of the pair form a group, which contains $at$ for $a\in\mathfrak A$ and $t\in L_T$, as $\bar v\in\Lambda_S$; so it
contains $L$. Hence the $y\in M_{n\times r}(E_h)$ for which $[\alpha y]$ lies in the lattice of the pair form a group
that contains the generators $vu^{\mathsf T}$, $v\in\Lambda$ and $u\in L$, of $\Lambda\otimes L$, and so contains
$\Lambda\otimes L$. So every generator of $\hat\Lambda(\Lambda\otimes L)$ lies in the lattice of the pair, and the two
groups are equal.
\end{proof}
% Ker.LatTens.latO_eq_hatA (78 lines in LatOEqHatA.lean), followed, for a Ker.Ctx (it uses S.gram = 1 from C.m.gS),
% with (forall v, v in Lambda iff star v in LS) and IsAA L LT: the lattice of the pair is
% Formalism.IndefR.Datum.latO S T LS LT (fun _ => True), whose generators gen_eq writes as gen S T alpha
% (vecMulVec (star s) (a . t)); le_antisymm, with hmem and the two closure inductions hT and hL'.

\begin{lemma}\label{is-eigenvalue-conjugate-frame-unit}
In Situation~\ref{is-eigenvalue-conjugate-frame-situation}, let $w_0\in E_h$ be an algebraic integer
with $w_0+\bar w_0=1$ such that $w_0\bar w_0=2b$ for an algebraic integer $b$. Let $Z$ be a signature
family and $x=x^Z$; let $L\subset E_h^r$ be an $\mathcal O_{E_h}$-submodule that is self-dual for $X$
(Definition~\ref{is-eigenvalue-conjugate-frame-fibres}(7)), and $L_T$ an $\mathcal O_{E_h}$-submodule of
$E_h^r$ with $L=\mathfrak AL_T$, the group generated by the products $\alpha v$ with $\alpha$ in the ideal
$\mathfrak A$ fixed by the host and $v\in L_T$. Let $f$ be a Fock vacuum vector for the Cartan involutions
$1$ of $S$ and $s_Z$ of $T$ and the lattice of the pair given by $\mathcal O_{E_h}^n$ and $L_T$. Let
$\vartheta$ be a level-one theta function for $w_0$, $x$ and $L$
(Definition~\ref{is-eigenvalue-conjugate-frame-heisenberg}(5)), formed with the symplectic basis and
indexing of $X$, and suppose that $\vartheta$ has mean one for $x$ and $L$. Then there is $a\in\mathbf C$ with $|a|=1$ such that
\[
 f(w)=a\,\mathcal G_x(z)\prod_{k=1}^n\vartheta(z_k),\qquad z=z_x(u(w)),
\]
for every $w\in\mathbb W$ with finite component $0$.
\end{lemma}
\leanref{Development/AbelianVarietyTheory/Proofs/IwasawaDecomposition/NormSqGaussProd.lean}{GlobalTheory.OneMapRef.VacUnit.vac_eq_unit_mul}
\leanref{Development/AbelianVarietyTheory/Definitions/SpanningFamilies/IsStabForm.lean}{GlobalTheory.Dict.IsTheta1}
\leanref{Development/AbelianVarietyTheory/Proofs/IwasawaDecomposition/NormSqGaussProd.lean}{GlobalTheory.OneMapRef.VacUnit.norm_periodic}
\leanref{Development/AbelianVarietyTheory/Proofs/IwasawaDecomposition/NormSqGaussProd.lean}{GlobalTheory.OneMapRef.VacUnit.norm_sq_gauss_prod}
\leanref{Development/AbelianVarietyTheory/Proofs/IwasawaDecomposition/NormSqGaussProd.lean}{GlobalTheory.OneMapRef.VacUnit.integral_prod_eq}
\begin{proof}
By Lemma~\ref{is-eigenvalue-conjugate-frame-signature}(1), $x\in\mathcal D_X$. Let $X'$, its symplectic
basis and indexing, and $\pi$ be as in Lemma~\ref{is-eigenvalue-conjugate-frame-conjugate-model}; with them the
data are as in Situation~\ref{is-eigenvalue-conjugate-frame-context}, whose notation and that of
Lemma~\ref{is-eigenvalue-conjugate-frame-zak} we use. Put $B=Z^\pi$. By (2) and (3) of
Lemma~\ref{is-eigenvalue-conjugate-frame-conjugate-model}, $B$ is orthonormal for $X'$ and $s'_B=s_Z$. The
flag data $X$ and $X'$ have the same module and, by (1) of that lemma, the same form
$\langle\ ,\ \rangle$ and the same Heisenberg operators; so $L$ is self-dual for $X'$, and the Heisenberg
operators, the signs and the inner product of Definition~\ref{is-eigenvalue-conjugate-frame-heisenberg} are
the same for $X$ and $X'$.

\emph{The vacuum.} By Lemma~\ref{is-eigenvalue-conjugate-frame-vacua}(1), applied with
$\Lambda_S=\mathcal O_{E_h}^n\in\mathcal V_1$, $\mathcal O_{E_h}^n\otimes L$ has a vacuum $v$ for $X$ and
$w_0$, which is also a vacuum for $X'$. The function $f$ is a theta function annihilated by the
annihilation operators for $1$ and $s_Z=s'_B$, and a lattice vacuum for the lattice of the pair given by
$\mathcal O_{E_h}^n$ and $L_T$, which is $\hat\Lambda(\mathcal O_{E_h}^n\otimes L)$ by
Lemma~\ref{is-eigenvalue-conjugate-frame-tensor-lattice}, as $\mathcal O_{E_h}^n$ is stable under complex
conjugation. So Lemma~\ref{is-eigenvalue-conjugate-frame-vacuum-line}(2), with $\Lambda_S=\mathcal
O_{E_h}^n$ and $M=1$, gives $t\in\mathbf C$ with $f=t\,\mathrm{Th}(x'_B,v)$. The function $\mathrm{Th}(x'_B,v)$ is a Zak theta function for
$x'_B$ and $v$, so by Lemma~\ref{is-eigenvalue-conjugate-frame-conjugate-model}(3) its value at $w$ with
finite component $0$ is $\Theta_v(x;u(w))=\mathcal G_x(z)\theta_v(x;z)$, $z=z_x(u(w))$
(Definition~\ref{is-eigenvalue-conjugate-frame-fibres}(6)). Hence $f(w)=t\,\mathcal G_x(z)\theta_v(x;z)$.

\emph{The constant.} By Lemma~\ref{is-eigenvalue-conjugate-frame-theta-product} there is $\kappa$ with
$\prod_k\vartheta(z_k)=\kappa\,\theta_v(x;z)$ for every $z\in(\mathbf C^d)^n$. As $\vartheta$ has mean one, it
is not identically $0$; if $\vartheta(z_0)\ne0$, then $\prod_k\vartheta(z_0)\ne0$, so $\kappa\ne0$. Put
$a=t/\kappa$. Then $f(w)=a\,\mathcal G_x(z)\prod_k\vartheta(z_k)$ for every $w$ with finite component $0$.

It remains to show $|a|=1$. Put $\phi(z)=|\vartheta(z)|^2e^{-\pi\operatorname{Re}H_x(z,z)}$, a continuous
function $\ge0$ on $\mathbf C^d$. Since $\mathcal G_x(z)=\exp(-\frac\pi2\sum_kH_x(z_k,z_k))$,
\[
 \Big|\mathcal G_x(z)\prod_k\vartheta(z_k)\Big|^2=\prod_k\phi(z_k).
\]
The function $\phi$ is invariant under the group $A=\{p_x(v):v\in L\}$. Indeed, let $v\in L$ and
$p=p_x(v)$. As $x\in\mathcal D_X$, $H_x$ is Hermitian, so $\operatorname{Re}H_x(z+p,z+p)=\operatorname{Re}H_x(z,z)
+2\operatorname{Re}H_x(z,p)+\operatorname{Re}H_x(p,p)$, and $|\epsilon(v)|=1$; taking absolute values in
the law of $\vartheta$ gives $\phi(z+p)=\phi(z)$.

As $f$ is a Fock vacuum vector, there are a Haar measure $\mu$ on $\mathbb W$ and a fundamental domain
$\mathcal F$ of $M_{r\times n}(E_h)$ with $\mu(\mathcal F)=1$ and $\int_{\mathcal F}|f|^2\,d\mu=1$. Let $U'$,
$\Gamma'_N$ and $N_0$ be as in Lemma~\ref{is-eigenvalue-conjugate-frame-torus-mean} for $f'=f$. The module
$\Lambda$ is finitely generated and $L$ spans $E_h^r$, so there is $N_1>0$ with $N_1\Lambda\subset L$. Let
$N=N_0N_1$ and $F_N$ a fundamental domain of $\Gamma'_N$ in $U'$, by that lemma. For $u'\in U'$ let
$J(u')$ be the family whose $k$-th member has the component $\overline{(u'_k)_{\bar\rho}}$ at $\rho\in\mathcal
T$, and put $e(u')=\big(z_x(J(u')_1),\dots,z_x(J(u')_n)\big)$. As $x\in\mathcal D_X$, $z_x$ is a real linear
isomorphism, so $e$ is a real linear isomorphism of $U'$ onto $(\mathbf C^d)^n$. By the definitions,
$J(u'(w))=u(w)$, so $e(u'(w))=z_x(u(w))$ for $w$ with finite component $0$. And $e$ carries $\Gamma'_N$ into
$A^n$: $J$ carries $(N\sigma(y_k))_{k,\sigma}$ to $(N\rho(y_k))_{k,\rho}$, whose image under $e$ is
$(p_x(Ny_k))_k$, and $Ny_k=N_0(N_1y_k)\in L$. Let $F\subset\mathbf C^d$ be a fundamental domain of $A$ with
$\operatorname{vol}(F)^{-1}\int_F\phi=1$, which exists as $\vartheta$ has mean one. Then
\[
 \int_{F_N}\prod_k\phi(e(u')_k)\,du'=\operatorname{vol}(F_N).
\]
Indeed, $F^n$ is a fundamental domain of $A^n$, with $\int_{F^n}\prod_k\phi(z_k)\,dz=\operatorname{vol}(F)^n
=\operatorname{vol}(F^n)$. The group $C=e^{-1}(A^n)$ contains $\Gamma'_N$, is countable, and has the
fundamental domain $e^{-1}(F^n)$, on which the integral of the $C$-invariant function
$\prod_k\phi(e(u')_k)$ is its volume, as $e$ carries Lebesgue measure to a multiple of Lebesgue measure.
Cutting $F_N$ along the translates of $e^{-1}(F^n)$ by $C$, both the integral over $F_N$ and the volume of
$F_N$ are the sums over $C/\Gamma'_N$ of the integral and the volume over $e^{-1}(F^n)$; so they are equal.

The function $\Phi(u')=|a|^2\prod_k\phi(e(u')_k)$ satisfies $f(w)\overline{f(w)}=\Phi(u'(w))$ for $w$ with
finite component $0$, by the formula for $f$ and the identity for $|\mathcal G_x(z)\prod_k\vartheta(z_k)|^2$.
So Lemma~\ref{is-eigenvalue-conjugate-frame-torus-mean} and the identity above give
$1=\int_{\mathcal F}|f|^2\,d\mu=\operatorname{vol}(F_N)^{-1}\int_{F_N}\Phi(u')\,du'=|a|^2$, where the last
step uses $\operatorname{vol}(F_N)\ne0$, which holds as the middle term is $1$. So $|a|=1$.
\end{proof}
% vac_eq_unit_mul: Flip.exists_flip_model (X', m, bas', idx', pi, hLambda, hEF, hhe, H), ctxOf for the context C
% and B = Z o pi; TensVac.vacua (with stdL_mem_Vset) for v and hvac'; Ker.latVac_th_line for t with
% LatTens.latO_eq_hatA (L = AA L_T) and the Fock condition; Ker.ZakEx.isZak_Th with hZak for hz;
% VacProd.prod_eq_mul_thA for kappa; hkappa0 from MeanOne; hfp. The norm: hv.norm, Ker.TorusMean.ipO_eq_mean for
% N0, exists_nsmul_mem for N1, Ker.zR_laws for e0, conjPt for J, hJ and hJ', norm_periodic (hx_laws for the
% Hermitian symmetry), integral_prod_eq (lintegral_sub, isFD_pi), norm_sq_gauss_prod and the final computation.

\subsection{Frame constants of the core}\label{subsection-is-frame-constants}

\begin{situation}\label{is-eigenvalue-conjugate-frame-core}
Let $R$ be a realization of a definite-side datum $D$ over a field $K_s\subset L$ with a face $F$ that
satisfies the core conditions at $F$, except possibly the one on its chain, with the notation of
Definitions~\ref{is-eigenvalue-lift-data} and~\ref{is-eigenvalue-conjugate-datum}. Let one-map data for
$R$ be given (Definition~\ref{is-one-map-data}), with a core whose coefficient field is $K_s$, weight
families $\kappa_j$, vectors $v_0(j,\sigma)$ and a weight tie, pinned at a lift $i$, with $L'$ and $\iota$
as in Lemma~\ref{is-eigenvalue-conjugate-class-sums}. Let $\sigma$ be an embedding of $L$ that agrees with
$\sigma_0$ on $K$, $\sigma'$ an embedding of $L'$ with $\sigma'\circ\iota=\sigma$, and $s$ the restriction
of $\sigma$ to $K_s$. Write $h\mapsto(c(h),x(h))$ for the dictionary of the core at level $i+1$ and $s$
and $\mu$ for its multiplier (Lemma~\ref{is-eigenvalue-conjugate-dictionary-good}); for a component $c$,
$X_c$ for its flag data, with set $\Psi_c$, module $\Lambda_c$ and flag domain $\mathcal D_c$, and
$y(c,x)$ for the complex framed point of the moduli of the core at level $i+1$ that the uniformization
of the core at $s$ attaches to $c$ and $x\in\mathcal D_c$, with its section $s_{y(c,x)}$
(Definition~\ref{is-eigenvalue-conjugate-frame-fibres}(3)). Let $\theta$ be the archimedean family of
$R$, $Z_0=(x_{0,\rho})_\rho$ the frames of the model of $W_i$ at $\sigma$, $n=\operatorname{rk}V_i$,
$r=\operatorname{rk}W_i$, $q$ the integer $q$ of the datum of the core at level $i+1$, $\mathfrak A$ the
ideal fixed by the host, $w_0=(1+\sqrt{-d})/2$ for the prime $d$ of the host, $\Lambda_V$ and $\Lambda_W$
the lattices of $V_i$ and $W_i$ of the realization, and $\hat L_0$ the completion of the lattice $L_0$ of
$W_i$. For $g\in U(W_i)(\mathbf A)$ let $L(g)$ be the $\mathcal O_{E_h}$-span of the products $\alpha v$
with $\alpha\in\mathfrak A$ and $v\in E_h^r$, $v\in g_f\hat L_0$.
\end{situation}
\leanref{Development/AbelianVarietyTheory/Proofs/DivisorProofs/FrAtEmb.lean}{GlobalTheory.GalOne.dictS}
\leanref{Development/AbelianVarietyTheory/Proofs/DivisorProofs/FrAtEmb.lean}{GlobalTheory.GalOne.zAt}
\leanref{Formalism/LocalToGlobal/ProjectiveLattices/FamilyUniformisationBasePointOrthogonalBis.lean}{Formalism.L2.Pj.latG}
% The hypotheses are those of GalOne.framed_point_at_ofFive: Core5 R F, FWTie T R kappa v0, O : OM T R i with
% Pinned, sigma'.comp O.jL = sigma, OverK sigma and hV. The dictionary is dictS hV (castDict of T.dict (i+1)
% (sOf O sigma')), the multiplier muS hV, Z0 = zAt R i sigma = x0Mat (R.MW i sigma), w0 = H.w0, L(g) = latG.

\begin{definition}\label{is-eigenvalue-conjugate-frame-framing}
In Situation~\ref{is-eigenvalue-conjugate-frame-core}, a \emph{framing} at $h\in U(W_i)(\mathbf A)$
consists of $\gamma\in U(W_i)(F_h)$, a symplectic basis of $\Lambda_{c(h)}$, an indexing and a bijection
onto the letters, the last two matched for $s$ (Definition~\ref{is-eigenvalue-conjugate-frame-fibres}),
such that: $\psi$, $\theta$, $S=V_i$, $T=W_i$, $X=X_{c(h)}$, the basis and the indexing satisfy the
conditions of Situation~\ref{is-eigenvalue-conjugate-frame-situation}; the integer $p_0$ of $X_{c(h)}$ is
$q$; $\Lambda_{c(h)}$ is self-dual for $X_{c(h)}$ (Definition~\ref{is-eigenvalue-conjugate-frame-fibres}(7))
and $\Lambda_{c(h)}=L(\gamma h)$; $x(h)=x^{\gamma hZ_0}$ and $x(h)\in\mathcal D_{c(h)}$; and there are an
analytic fibre at $y(c(h),x(h))$ for these data, $x(h)$ and $\Lambda_{c(h)}$, and a theta reading of
$s_{y(c(h),x(h))}$ for it that has mean one for $x(h)$ and $\Lambda_{c(h)}$. The element $\gamma$ is the
\emph{element} of the framing. Coordinates, Gaussians and theta functions at a framing are formed with its
symplectic basis and indexing.
\end{definition}
\leanref{Development/AbelianVarietyTheory/Proofs/IwasawaDecomposition/OmegaFormScalarExtension.lean}{GlobalTheory.OneMapRef.FrameConst.IsPt}
\leanref{Development/AbelianVarietyTheory/Proofs/DivisorProofs/FrAtEmb.lean}{GlobalTheory.GalOne.PointAt}

\begin{definition}\label{is-eigenvalue-conjugate-frame-marking}
In Situation~\ref{is-eigenvalue-conjugate-frame-core}, let $\mathfrak L$ be the lattice of the datum of the
core at level $i+1$, an $\mathcal O_{E_h}$-module, let $M$ and $\mathfrak M$ be the integer and the ideal of
the level of the core, and let $\iota\colon\mathfrak L\to E_h^r$ be the embedding of $\mathfrak L$ into $W_i$
given by the core. For $g\in U(W_i)(\mathbf A)$ let $\hat L(g)\subset\mathbf A_{E_h,f}^r$ be the subgroup
generated by the products $\alpha z$ with $\alpha\in\mathfrak A$ and $z\in g_f\hat L_0$.
\begin{enumerate}
\item Let $\ell$ be a level of the datum, with integer $M_\ell$ and ideal $\mathfrak M_\ell$
  (Definition~\ref{cv-model-toroidal-semiabelian}(1)). A \emph{marking} at $g$ for $\ell$ is an $\mathcal
  O_{E_h}$-linear map $m\colon\mathfrak L\to E_h^r$ with $m(a)-g_f\iota(a)\in M_\ell\hat L(g)$ for every
  $a\in\bar{\mathfrak M}_\ell\mathfrak L$. A \emph{marking} at $g$ is a marking at $g$ for the level of the
  core.
\item Let $y$ be a complex framed point of the moduli of the core at level $i+1$, whose level structure
  sends $a\in\mathfrak L$ to a point $\lambda_y(a)$ of $A_y$. An analytic fibre at $y$ for some data, a point
  $x$ and a module $L$ (Definition~\ref{is-eigenvalue-conjugate-frame-fibres}(3)) is \emph{marked} by an
  $\mathcal O_{E_h}$-linear map $m\colon\mathfrak L\to E_h^r$ if
  $\lambda_y(a)=\operatorname{ex}\big(p_x(M^{-1}m(a))\big)$ for every $a\in\bar{\mathfrak M}\mathfrak L$.
\end{enumerate}
\end{definition}
\leanref{Development/AbelianVarietyTheory/Proofs/IwasawaDecomposition/FibExtMarked.lean}{GlobalTheory.OneMapRef.MarkAd}
\leanref{Development/AbelianVarietyTheory/Definitions/SpanningFamilies/IsStabForm.lean}{GlobalTheory.Dict.AnFib.Marked}
\leanref{Formalism/LocalToGlobal/ProjectiveLattices/FamilyUniformisationBasePointOrthogonalBis.lean}{Formalism.L2.Pj.latGHat}
\leanref{Formalism/LocalToGlobal/ProjectiveLattices/ThShort.lean}{Formalism.L2.Pj.Mfr}
% MarkAd hr l iota h mk: for a in Mfr l * La (Mfr l = H.bar (l.Mb), the conjugate of the level ideal), finVec (mk a)
% - finHom h (finVec (iota a)) = l.M w with w in latGHat T h, the closure of the products (1 tensor a) z, a in AA,
% z in gLat T h. AnFib.Marked A mk: levPt Xo y.x a = A.ex (pv X idx x (M^-1 mk a)) for a in Mfr l * La. iota is
% iotaW T i, with LatEmb.

\begin{situation}\label{is-eigenvalue-conjugate-frame-uniformization}
Let a PEL datum of signature $(p,q)$ over a CM field $E$ be given (Situation~\ref{cv-function-situation}),
with prime $d$, lattice $\mathfrak L$ and multiplicities $m_{\tau'}$; its \emph{letters} are the pairs
$(\tau',i)$ of an embedding $\tau'$ of $E$ and $1\le i\le m_{\tau'}$, and $w_0=(1+\sqrt{-d})/2$. Let a level
$(M,\mathfrak M,N)$ of the datum (Definition~\ref{cv-model-toroidal-semiabelian}(1)) and its moduli
scheme be given, a coefficient field $K_c$ of the moduli scheme with a good embedding $s$
(Situation~\ref{cv-cusp-setting}), an integer $r\ge1$, and a uniformization of the moduli scheme at $s$ of
rank $r$ that satisfies the compatibility conditions of the moduli problem. It has finitely many
components $c$, each with flag data $X_c$ of rank $r$ over $E$ (Definition~\ref{is-onto-share-flag-data}),
with set $\Psi_c$, element $\xi_0$, numbers $c_\tau=2|\tau(\xi_0)|$, module $\Lambda_c$ and flag domain
$\mathcal D_c$, and a complex framed point $y(c,x)$ of the moduli scheme over $s$ at each $x\in\mathcal
D_c$, with its section $s_{y(c,x)}$ (Definition~\ref{is-eigenvalue-conjugate-frame-fibres}(3)). Among the
compatibility conditions are a bijection $j\mapsto\ell_j$ from $\{1,\dots,g\}$ onto the letters and a
family uniformization (Theorem~\ref{cv-uniformization-family}) that attaches to each component $c$ a
lattice $L_c$ with a map $t_c\colon\mathfrak L\to L_c$, an alternating form $\langle\ ,\ \rangle_c$ and a
signature $(p_c,q_c)$, and an $\mathcal O_E$-linear isomorphism $\varphi_c$ of $L_c$ onto $\Lambda_c$, with
the following properties, where $h_c$ is the form of $X_c$ and $H(z,w)=z^{\mathsf T}H\bar w$ for a matrix
$H\in M_g(\mathbf C)$.
\begin{enumerate}
\item[(i)] $\langle a,b\rangle_c=\operatorname{Tr}_{E/\mathbf Q}\big(-\xi_0h_c(\varphi_ca,\varphi_cb)\big)$ for
  all $a,b\in L_c$.
\item[(ii)] For every $x\in\mathcal D_c$ there is a period $P_x$ of $L_c$ (Definition~\ref{cv-cm-point-periods}(3))
  for the type $j\mapsto s\circ\varepsilon_{\tau'_j}$, where $\ell_j=(\tau'_j,i_j)$, whose $j$-th coordinate
  at $a\in L_c$ is the period of $\varphi_c(a)$ at the letter $\ell_j$
  (Definition~\ref{is-eigenvalue-conjugate-frame-fibres}(5)), such that $y(c,x)$ has the point that the
  family uniformization attaches to $P_x$, and such that the derivative at $0$ of the exponential
  $\operatorname{ex}_x\colon\mathbf C^g\to A_{y(c,x)}(\mathbf C)$ of this point in the direction of the $j$-th
  standard basis vector is the frame vector $u_{\ell_j}$ of $y(c,x)$.
\item[(iii)] If moreover $|d_i(x,\tau)|=1$ for all $\tau$ and $i$, there are a Hermitian matrix $H$ with
  $\operatorname{Re}H(z,z)>0$ for $z\ne0$ and $\operatorname{Im}H(P_x(a),P_x(b))=\langle a,b\rangle_c$ for
  $a,b\in L_c$; functions $\operatorname{fr}_\beta$ and $\chi$ with the properties (b) and (c) of
  Definition~\ref{is-eigenvalue-conjugate-frame-fibres}(3) for $\operatorname{ex}_x$, with $H$ in place of
  $H_x$ and the $P_x(a)$, $a\in L_c$, in place of the $p_x(v)$; a theta reading $\vartheta$ of $s_{y(c,x)}$
  for $\operatorname{ex}_x$ and these $\operatorname{fr}_\beta$; and a fundamental domain $\mathcal F$ of
  $P_x(L_c)$ in $\mathbf C^g$ with
  $\operatorname{vol}(\mathcal F)^{-1}\int_{\mathcal F}|\vartheta(z)|^2e^{-\pi\operatorname{Re}H(z,z)}\,dz=1$.
\end{enumerate}
A choice of these data is a \emph{uniformization witness}, and $m_c=\varphi_c\circ t_c\colon\mathfrak L\to
E^r$ is the \emph{marking} of $c$. Markings of analytic fibres at the framed points are read as in
Definition~\ref{is-eigenvalue-conjugate-frame-marking}(2), with the level structure of the moduli scheme.
\end{situation}
\leanref{Formalism/LocalToGlobal/ProjectiveLattices/UnFree.lean}{Formalism.L2.PjFree.Unf}
\leanref{Formalism/LocalToGlobal/ProjectiveLattices/UnFree.lean}{Formalism.L2.PjFree.IsUnf}
\leanref{Formalism/LocalToGlobal/ProjectiveLattices/UnFree.lean}{Formalism.L2.PjFree.GoodEmb}
\leanref{Development/AbelianVarietyTheory/Definitions/SpanningFamilies/IsBaseClassLevelAdapted.lean}{GlobalTheory.Dict.UW}
\leanref{Development/AbelianVarietyTheory/Definitions/SpanningFamilies/IsBaseClassLevelAdapted.lean}{GlobalTheory.Dict.UW.mark}
% Unf Xo c sigma r: components U.Comp (a Fintype), flag data U.X k and framed points U.un k x; IsUnf U is the list
% of compatibility conditions (r = Dw.p + Dw.q, p0 = Dw.q and the type of Psi, the triangular equivalences, the
% readings of sections, and the existential family uniformization); GoodEmb c sigma sends the base root of the
% level to H.rt0; 1 <= r. A UW structure W (Dict.nonempty_uw gives one) is the last clause of IsUnf: W.e, W.F,
% W.ik, W.phi, W.isom and W.pt; W.emb k = phi_k into E^r and W.mark k = W.emb k o F.t (ik k). The full level N is
% l.A. (i) is W.isom; (ii) and (iii) are W.pt k x: Pm with the periods perFr at the letters W.e j, IsPeriod for the
% type relab c sigma o W.e, (U.un k x).x = F.u (ik k) Pm and the derivatives IsDerivAt0 of F.ex, and, under
% |dmin| = 1, Hm, chil, fr, theta, Fd with their laws (Hm Hermitian with positive real part and Im = EF on the
% periods, fr_hol, fr_ne, the cocycle, fr at 0, the quasi-periodicity with |chil (Pm a)| = 1, the reading of
% (U.un k x).s, Fd a fundamental domain of Pm.range and the normalized integral 1). The exponential is F.ex
% (ik k) Pm, with the laws ex_add, ex_onto, ex_ker, ex_lin and ex_eta of the family uniformization.

\begin{lemma}\label{is-eigenvalue-conjugate-frame-dictionary-point}
In Situation~\ref{is-eigenvalue-conjugate-frame-core}, let the lattices $L_c$, the maps $t_c$ and the
isomorphisms $\varphi_c$ be those that the laws of the weight tie attach to the components $c$ of the
uniformization of the core at level $i+1$ at $s$, as in
Situation~\ref{is-eigenvalue-conjugate-frame-uniformization}, and $m_c=\varphi_c\circ t_c$. For every $h\in
U(W_i)(\mathbf A)$ there is $\gamma\in U(W_i)(F_h)$ with the following properties:
\begin{enumerate}
\item $\Lambda_{c(h)}=L(\gamma h)$, and $m_{c(h)}$ is a marking at $\gamma h$
  (Definition~\ref{is-eigenvalue-conjugate-frame-marking});
\item $x(h)=x^{\gamma hZ_0}$ (Situation~\ref{is-eigenvalue-conjugate-frame-situation}(1));
\item if $x(h)\in\mathcal D_{c(h)}$, every analytic fibre at $y(c(h),x(h))$ for a symplectic basis of
  $\Lambda_{c(h)}$, an indexing and a bijection onto the letters matched for $s$, the point $x(h)$ and the
  module $\Lambda_{c(h)}$ is marked by every marking at $\gamma h$ with values in $\Lambda_{c(h)}$.
\end{enumerate}
\end{lemma}
\leanref{Development/AbelianVarietyTheory/Proofs/IwasawaDecomposition/DictionaryRationalElementAnalyticPointStatement.lean}{GlobalTheory.T.dictionary_rational_element_analytic_point}
\begin{proof}
The laws of the weight tie at the level $i+1$ and $s$ contain the lattice law of the dictionary of the core:
the data $L_c$, $t_c$ and $\varphi_c$, with a bijection $j\mapsto\ell_j$ onto the letters, satisfy (i) and
(ii) of Situation~\ref{is-eigenvalue-conjugate-frame-uniformization}, and for every $h\in U(W_i)(\mathbf A)$
there is $\gamma\in U(W_i)(F_h)$ such that
\begin{enumerate}
\item[($\alpha$)] $\Lambda_{c(h)}=L(\gamma h)$;
\item[($\beta$)] $m_{c(h)}(a)-(\gamma h)_f\iota(a)\in M\hat L(\gamma h)$ for every $a\in\bar{\mathfrak
  M}\mathfrak L$, that is, $m_{c(h)}$ is a marking at $\gamma h$
  (Definition~\ref{is-eigenvalue-conjugate-frame-marking});
\item[($\gamma$)] $x(h)_\tau=\theta_{c(h),\tau}(\gamma h)\,x_{0,\tau}$ for every $\tau\in\Psi_{c(h)}=\mathcal T$, where
  $\theta_{c,\tau}$ are the ring homomorphisms of the dictionary
  (Definition~\ref{is-eigenvalue-conjugate-dictionary}(5)(c)), which extend $\tau$ and are continuous, and
  $x_{0,\tau}$ are the frames of the model of $W_i$ at $\sigma$, so that $Z_0=(x_{0,\rho})_\rho$.
\end{enumerate}
Let $h\in U(W_i)(\mathbf A)$, and fix such $\gamma$. Then (1) is ($\alpha$) and ($\beta$).

(2) The maps $\theta_{c(h),\tau}$ and $\theta_\tau$ are continuous ring homomorphisms from $\mathbf
A_{F_h,\infty}\otimes_{F_h}E_h$ to $\mathbf C$ that extend $\tau$, so they are equal, as $F_h$ is dense in
$\mathbf A_{F_h,\infty}$. So $x(h)_\tau=\theta_\tau(\gamma h)Z_{0,\tau}=(\gamma hZ_0)_\tau$ for every $\tau$.

(3) Put $c=c(h)$, $x=x(h)\in\mathcal D_c$ and $y=y(c,x)$, with level structure $\lambda_y$. Let a symplectic
basis, an indexing and a bijection $j\mapsto\ell'_j$ onto the letters be given, the last two matched for $s$, let
$\operatorname{ex}_A$ be the exponential of an analytic fibre $A$ at $y$ for these data, $x$ and $\Lambda_c$, and
let $m$ be a marking at $\gamma h$ with values in $\Lambda_c$. Let $a\in\bar{\mathfrak M}\mathfrak L$.

\emph{The exponentials.} Let $P=P_x$ and $\operatorname{ex}_x$ be as in (ii). Let $\pi$ be the permutation with
$\ell'_j=\ell_{\pi(j)}$, and $\varrho(z)=(z_{\pi(j)})_j$ for $z\in\mathbf C^g$. Then $\operatorname{ex}_A\circ\varrho$ and
$\operatorname{ex}_x$ are holomorphic group homomorphisms $\mathbf C^g\to A_y(\mathbf C)$, and their derivatives at
$0$ in the direction of the $m$-th standard basis vector $e_m$ are both $u_{\ell_m}$: for the first, because
$\varrho(e_m)=e_{\pi^{-1}(m)}$ and $\operatorname{ex}_A$ has the derivative $u_{\ell'_{\pi^{-1}(m)}}=u_{\ell_m}$ in
that direction. A holomorphic group homomorphism $e$ from $\mathbf C^g$ to a complex Lie group is determined by its
derivative at $0$, as $t\mapsto e(tz)$ is the one-parameter subgroup with tangent vector $de_0(z)$. So
$\operatorname{ex}_x=\operatorname{ex}_A\circ\varrho$.

\emph{The level structure.} By Theorem~\ref{cv-uniformization-family}, $\lambda_y(a)=\operatorname{ex}_x(M^{-1}P(t_ca))
=\operatorname{ex}_A(\varrho(M^{-1}P(t_ca)))$. The $j$-th entry of $\varrho(P(b))$ is $P(b)_{\pi(j)}$, the period of
$\varphi_c(b)$ at the letter $\ell_{\pi(j)}=\ell'_j$, which is $p_x(\varphi_cb)_j$ as the indexing is matched. As
$p_x(M^{-1}v)=M^{-1}p_x(v)$ (Definition~\ref{is-eigenvalue-conjugate-frame-fibres}(1)),
$\lambda_y(a)=\operatorname{ex}_A\big(p_x(M^{-1}m_c(a))\big)$.

\emph{The two markings.} By ($\beta$) and as $m$ is a marking at $\gamma h$, there are $w,w'\in\hat L(\gamma h)$ with
$m(a)-(\gamma h)_f\iota(a)=Mw$ and $m_c(a)-(\gamma h)_f\iota(a)=Mw'$. So the vector $v=M^{-1}(m(a)-m_c(a))$ of $E_h^r$
has the image $w-w'$ in $\hat L(\gamma h)$. As a lattice of $E_h^r$ is the set of the vectors of $E_h^r$ that lie
in its completion \cite{Wei67}, $v\in L(\gamma h)=\Lambda_c$. Then $p_x(M^{-1}m(a))=p_x(M^{-1}m_c(a))+p_x(v)$ with
$\operatorname{ex}_A(p_x(v))=1$, and
\[
 \lambda_y(a)=\operatorname{ex}_A\big(p_x(M^{-1}m_c(a))\big)=\operatorname{ex}_A\big(p_x(M^{-1}m(a))\big).
\]
So $A$ is marked by $m$.
\end{proof}
% T.dictionary_rational_element_analytic_point (OneMapRef.DictPt.dict_point, 128 lines in its file), applied
% with the laws dictU and adPin (DictLat) and the base frame tie.x0 (hx0 of factsAt_ofFive), as in
% framed_point_at_ofFive; MarkBy is the last clause of the statement. The Lean statement takes DictLat U W hr
% dict iota e0 F ik phi th x0 as a hypothesis: its first clause is (i), its second (ii) without (iii), and its
% third, for each h, gives gamma in GQ with (alpha), (beta) and (gamma) (archM (th k tau) for theta_{c,tau}, with
% IsInfEmb tau (th k tau)); x0 (tauE rho) = Z0 rho. dict_point follows the proof above: MarkBy from hD.2; (2)
% through SmoothForms.RealPts.infEmb_unique; (3) through DictPt.levPt_eq (AnUniq.exp_eq for the uniqueness of the
% exponential, with isDerivAt0_comp, isHolTo_ex and F.ex_eta), DictPt.exists_sub_eq (MarkCoset.latG_spec) and
% ex_add_ker, and per_smul for the division by M.

\begin{lemma}\label{is-eigenvalue-conjugate-frame-component-lattice-positive}
Let $X$ be flag data of rank $r$ over a CM field $E$ (Definition~\ref{is-onto-share-flag-data}), with set
$\Psi$, put $d=|\Psi|r$, and fix an indexing (Definition~\ref{is-eigenvalue-conjugate-frame-fibres}(1)). For
every $x\in\mathcal D_X$ there is a matrix $B\in M_d(\mathbf C)$ with $\operatorname{Re}(z^{\mathsf T}B\bar z)>0$
for $z\ne0$ and $\operatorname{Im}\big(p_x(u)^{\mathsf T}B\,\overline{p_x(v)}\big)=\langle u,v\rangle_X$ for all
$u,v\in E^r$.
\end{lemma}
\leanref{Development/AbelianVarietyTheory/Proofs/SiegelThetaTransformation/UniformisationComponentLatticesMarkings.lean}{GlobalTheory.Dict.UWLat.exists_pos_form}
\leanref{Development/AbelianVarietyTheory/Proofs/UnitaryExistence/HermFormConjDiag.lean}{GlobalTheory.FramePeriod.exists_loc}
\begin{proof}
Let $h$ and $h_\tau$ be the forms of $X$, $\kappa_0$ its element, and put $\xi=-\kappa_0/2$. Then $\bar\xi=-\xi$,
so $\tau(\xi)$ is purely imaginary, and $\operatorname{Im}\tau(\xi)>0$ for $\tau\in\Psi$, as
$\operatorname{Im}\tau(\kappa_0)<0$.

Fix $\tau\in\Psi$ and let $l_\tau(u)\in\mathbf C^r$, for $u\in\mathbf C^r$, have the entries
$h_\tau(x_\tau e_a,u)$ for $a\le p_0$ and $(x_\tau^{-1}u)_a$ for $a>p_0$, so that $p_x(v)$ has the entry
$l_\tau(\tau(v))_a$ at the index $j$ with $(\tau_j,a_j)=(\tau,a)$. The Gram matrix $\Gamma$ of the columns of
$x_\tau$ for $h_\tau$ is Hermitian, and its leading minors are real and nonzero with the signs $\epsilon_i$. So
there are an upper triangular matrix $U$ with diagonal entries $1$ and real numbers $\delta_a$ with
$U^*\Gamma U=\operatorname{diag}(\delta_a)$; the sign of $\delta_a$ is $\operatorname{sgn}(a)$. Put
$c(u)=U^{-1}x_\tau^{-1}u$, the coordinates of $u$ in the columns of $x_\tau U$. Then
$h_\tau(u,w)=\sum_a\delta_a\,c(u)_a\overline{c(w)_a}$. Let $S_\tau$ be the matrix whose $a$-th row is
$(U_{ba}/\delta_a)_b$ for $a\le p_0$ and the $a$-th row of $U^{-1}$ for $a>p_0$. As $U$ and $U^{-1}$ are upper
triangular, $(S_\tau l_\tau(u))_a=\delta_a^{-1}h_\tau(x_\tau Ue_a,u)=\overline{c(u)_a}$ for $a\le p_0$, and
$(S_\tau l_\tau(u))_a=c(u)_a$ for $a>p_0$. So $S_\tau l_\tau(u)=0$ only for $u=0$; hence $l_\tau$ is injective,
so it is a real linear bijection, and $S_\tau$ is invertible. Put
$H_\tau=S_\tau^{\mathsf T}\operatorname{diag}(\operatorname{sgn}(a)\delta_a)\overline{S_\tau}$, a Hermitian matrix.
Then $\operatorname{Re}(w^{\mathsf T}H_\tau\bar w)=\sum_a|\delta_a||(S_\tau w)_a|^2>0$ for $w\ne0$. For $a\le p_0$
the entry $a$ of $l_\tau(u)^{\mathsf T}H_\tau\overline{l_\tau(w)}$ is $\delta_a\overline{c(u)_a}c(w)_a$, and for
$a>p_0$ it is $-\delta_ac(u)_a\overline{c(w)_a}$; so
\[
 \operatorname{Im}\big(l_\tau(u)^{\mathsf T}H_\tau\overline{l_\tau(w)}\big)
 =-\sum_a\delta_a\operatorname{Im}\big(c(u)_a\overline{c(w)_a}\big)=-\operatorname{Im}h_\tau(u,w).
\]

Let $B$ be the block diagonal matrix, in the indexing, with the blocks $2\operatorname{Im}\tau(\xi)\,H_\tau$. Then
$\operatorname{Re}(z^{\mathsf T}B\bar z)>0$ for $z\ne0$, and for $u,v\in E^r$
\[
 \operatorname{Im}\big(p_x(u)^{\mathsf T}B\,\overline{p_x(v)}\big)=-\sum_{\tau\in\Psi}2\operatorname{Im}\tau(\xi)
 \operatorname{Im}h_\tau(\tau(u),\tau(v))=\sum_{\tau\in\Psi}2\operatorname{Re}\big(\tau(\xi)h_\tau(\tau(u),\tau(v))\big),
\]
as $\operatorname{Re}\tau(\xi)=0$. As $\Psi$ contains one of each pair of conjugate embeddings, this is
$\operatorname{Tr}_{E/\mathbf Q}(\xi h(u,v))=\langle u,v\rangle_X$.
\end{proof}
% exists_pos_form: exists_loc at each tau (the bijectivity of locPer, the Hermitian Hl = S^T diag(sgn d) conj S
% from exists_upper_diag, with Im Hl(locPer u, locPer w) = -Im hermC u w), the block matrix of the
% 2 Im tau(xi) Hl (hermForm_block), im_neg_half_pos, and EF_def with trace_re.

\begin{lemma}\label{is-eigenvalue-conjugate-frame-component-lattice}
In Situation~\ref{is-eigenvalue-conjugate-frame-uniformization}, fix a uniformization witness. Then for
every component $c$:
\begin{enumerate}
\item the element $\kappa_0$ of $X_c$ is $2\xi_0$, and $\Lambda_c$ is self-dual for $X_c$
  (Definition~\ref{is-eigenvalue-conjugate-frame-fibres}(7));
\item $m_c(\mathfrak L)\subset\Lambda_c$; if $a\in\bar{\mathfrak M}\mathfrak L$ and $m_c(a)\in M\Lambda_c$,
  then $a\in M\mathfrak L$; and every $v\in\bar{\mathfrak M}\Lambda_c$ is $m_c(a)+Mu$ with
  $a\in\bar{\mathfrak M}\mathfrak L$ and $u\in\Lambda_c$;
\item $r=p_c+q_c$; $\varphi_c$ is injective as a map into $E^r$ and its image spans $E^r$ over $E$; and the
  alternating form of $L_c$ is $(a,b)\mapsto\langle\varphi_c(a),\varphi_c(b)\rangle_{X_c}$
  (Definition~\ref{is-eigenvalue-conjugate-frame-fibres}(2)).
\end{enumerate}
\end{lemma}
\leanref{Development/AbelianVarietyTheory/Proofs/SiegelThetaTransformation/UniformisationComponentLatticesMarkings.lean}{GlobalTheory.Dict.uw_lattice}
\leanref{Development/AbelianVarietyTheory/Definitions/SpanningFamilies/IsBaseClassLevelAdapted.lean}{GlobalTheory.Dict.nonempty_uw}
\begin{proof}
Write $X=X_c$, $\Lambda=\Lambda_c$, $\Psi=\Psi_c$, $h$ and $h_\tau$ for its forms and $\langle\ ,\ \rangle_X$ for
its form of Definition~\ref{is-eigenvalue-conjugate-frame-fibres}(2). The form $\langle\ ,\ \rangle_c$ of $L_c$
is perfect, as $L_c$ is the lattice of a PEL datum (Definition~\ref{cv-cm-point-periods}); in particular
$na=nb$ for an integer $n\ne0$ implies $a=b$ in $\mathfrak L$ and in $L_c$.

\emph{The element $\kappa_0$.} By Definition~\ref{is-onto-share-flag-data}(6), $\kappa_0=\pm2\xi_0$. Suppose
that $\kappa_0=-2\xi_0$. Then $\langle u,v\rangle_X=\operatorname{Tr}_{E/\mathbf Q}(\xi_0h(u,v))$, so by (i) of
the situation $\langle\varphi_ca,\varphi_cb\rangle_X=-\langle a,b\rangle_c$ for $a,b\in L_c$. The flag domain
$\mathcal D_c$ is not empty: by Definition~\ref{is-onto-share-flag-data}(2) there are matrices $x_\tau$ whose
columns have Gram matrix $J$ for $h_\tau$, and then $d_i(x,\tau)=\epsilon_i$. Fix $x\in\mathcal D_c$ and let
$P=P_x$ be as in (ii) of the situation. As $P$ is a period of $L_c$, its image spans $\mathbf C^g$ over
$\mathbf R$, and there is a matrix $A$ with $\operatorname{Re}A(z,z)>0$ for $z\ne0$ and
$\operatorname{Im}A(P(a),P(b))=\langle a,b\rangle_c$ (Definition~\ref{cv-cm-point-periods}(3)).

By the compatibility conditions, $r=p+q$, the integer $p_0$ of $X$ is $q$, and $\Psi$ is the set of the $\tau$
with $\tau(\sqrt{-d})=\theta$. As $s$ is good, $\tau'\mapsto s\circ\varepsilon_{\tau'}$ is a bijection of the
embeddings of $E$, and $m_{\tau'}=p$ if $s\circ\varepsilon_{\tau'}\in\Psi$ and $m_{\tau'}=q$ otherwise. So the map
that sends the letter $(\tau',i)$ to $(\tau,q+i)$ if $\tau=s\circ\varepsilon_{\tau'}\in\Psi$, and to
$(\bar\tau,q+1-i)$ if $\tau\notin\Psi$, is a bijection of the letters onto $\Psi\times\{1,\dots,r\}$. Composed
with $j\mapsto\ell_j$ it is an indexing for which $p_x(v)_j$ is the period of $v$ at the letter $\ell_j$, by
Definition~\ref{is-eigenvalue-conjugate-frame-fibres}(1) and~(5). With this indexing $P(a)=p_x(\varphi_ca)$ for
$a\in L_c$, by (ii). Let $B$ be as in Lemma~\ref{is-eigenvalue-conjugate-frame-component-lattice-positive} for
$x$ and this indexing. Then for $a,b\in L_c$
\[
 \operatorname{Im}A(P(a),P(b))=\langle a,b\rangle_c=-\langle\varphi_ca,\varphi_cb\rangle_X
 =-\operatorname{Im}B(P(a),P(b)).
\]
So the function $F(z,w)=(A+B)(z,w)$, linear in $z$ and conjugate linear in $w$, has
$\operatorname{Im}F(z,w)=0$ for $z,w\in P(L_c)$. As $\operatorname{Im}F$ is real bilinear and $P(L_c)$ spans
$\mathbf C^g$ over $\mathbf R$, $\operatorname{Im}F=0$. Here $g=|\Psi|r\ge1$, as $r\ge1$ and $\Psi$ contains
one of every two conjugate embeddings. Let $\mathbf 1\in\mathbf C^g$ have all entries $1$. Then
$0=\operatorname{Im}F(i\mathbf 1,\mathbf 1)=\operatorname{Re}A(\mathbf 1,\mathbf 1)+\operatorname{Re}B(\mathbf
1,\mathbf 1)>0$, a contradiction. Hence $\kappa_0=2\xi_0$, and by (i)
\[
 \langle a,b\rangle_c=\operatorname{Tr}_{E/\mathbf Q}\big(-\xi_0h(\varphi_ca,\varphi_cb)\big)
 =\langle\varphi_ca,\varphi_cb\rangle_X\qquad(a,b\in L_c).\tag{$*$}
\]

\emph{Self-duality.} The module $\Lambda$ is finitely generated
(Definition~\ref{is-onto-share-flag-data}(3)). For $v=\varphi_ca$ and $w=\varphi_cb$ in $\Lambda$,
$\langle v,w\rangle_X=\langle a,b\rangle_c\in\mathbf Z$ by $(*)$. Conversely, let $v\in E^r$ with
$\langle v,w\rangle_X\in\mathbf Z$ for every $w\in\Lambda$. Then $b\mapsto\langle v,\varphi_cb\rangle_X$ is a
homomorphism $L_c\to\mathbf Z$, and as $\langle\ ,\ \rangle_c$ is perfect there is $a\in L_c$ with
$\langle a,b\rangle_c=\langle v,\varphi_cb\rangle_X$ for every $b$. So $\langle v-\varphi_ca,\varphi_cb\rangle_X=0$
for every $b$, by $(*)$. As $\Lambda$ is finitely generated and spans $E^r$, there is an integer $N>0$ with
$N(v-\varphi_ca)\in\Lambda$, say $N(v-\varphi_ca)=\varphi_ca'$. Then $\langle a',b\rangle_c=N\langle
v-\varphi_ca,\varphi_cb\rangle_X=0$ for every $b$, so $a'=0$ as the form is perfect, and $v=\varphi_ca\in\Lambda$.
This proves (1).

(2) Let $M'$ be the least common multiple of $M$ and the full level $N$, and $m=M'/M$, an integer $\ne0$. The map
$t_c$ induces an isomorphism $\mathfrak L/M'\mathfrak L\to L_c/M'L_c$ (Theorem~\ref{cv-uniformization-family}):
$t_c(a)\in M'L_c$ exactly when $a\in M'\mathfrak L$, and every $b\in L_c$ is $t_c(a)+M'e$ with $a\in\mathfrak L$
and $e\in L_c$. First, $m_c(a)=\varphi_c(t_ca)\in\Lambda$. Next, let $a\in\mathfrak L$ with $m_c(a)=Mu$, $u\in\Lambda$;
the hypothesis $a\in\bar{\mathfrak M}\mathfrak L$ is not needed. Write $u=\varphi_cb$. As $\varphi_c$ is injective,
$t_c(a)=Mb$, so $t_c(ma)=M'b\in M'L_c$, and $ma=M'b'=m\,Mb'$ for some $b'\in\mathfrak L$; so $a=Mb'$. Finally, the
elements $m_c(a)+Mu$ with $a\in\bar{\mathfrak M}\mathfrak L$ and $u\in\Lambda$ form a group, and $\bar{\mathfrak
M}\Lambda$ is generated by the products $sw$ with $s\in\bar{\mathfrak M}$ and $w\in\Lambda$. Write $w=\varphi_cb$
and $b=t_c(a)+M'e$. Then $sw=\varphi_c(t_c(sa))+M\varphi_c(sme)=m_c(sa)+Mu$ with $sa\in\bar{\mathfrak M}\mathfrak L$
and $u=\varphi_c(sme)\in\Lambda$.

(3) By the compatibility conditions $r=p+q$, and the lattices of the family uniformization have the signature
$(p,q)$ of the datum, so $r=p_c+q_c$. The map $\varphi_c$ is an isomorphism onto $\Lambda\subset E^r$, which
spans $E^r$ over $E$ (Definition~\ref{is-onto-share-flag-data}(3)); and the form is $(*)$.
\end{proof}
% Dict.uw_lattice (65 lines; UniformisationComponentLatticesMarkings.lean, 254 lines), under GoodEmb, IsUnf and
% 1 <= r, for any UW structure W: hkappa (kappa0_xi; frames_isConnected for a point x, where the book reads the
% nonemptiness from Data.signature; W.pt for P with IsPeriod (hspan, A); FrameUnif.exists_letter_equiv with
% relab_good (hbij, hrt) and the multiplicities Dw.mult for the indexing idx = W.e.trans f and hper; exists_pos_form
% for B; not_opposite, h1 = W.isom); hE (*); selfD_of_unimod (D.unimod, LatticeSpanMembership.exists_nat_smul_mem,
% EF_sub_left, EF_add_right, EF_nat_left); isMark_of with W.F.t_inj and W.F.t_surj (M' = Level.M' = lcm M A,
% nsmul_cancel); W.F.sig for r = p + q, and Lambda_span.

\begin{definition}\label{is-eigenvalue-conjugate-frame-letter-dictionary}
Let a PEL datum over a CM field $E$ with multiplicities $m_{\tau'}$, a level and an embedding $s$ of a
coefficient field $K_c$ of its moduli scheme be given, with letters as in
Situation~\ref{is-eigenvalue-conjugate-frame-uniformization}; for an embedding $\tau'$ of $E$ write
$s\tau'=s\circ\varepsilon_{\tau'}$ (Definition~\ref{is-eigenvalue-conjugate-frame-fibres}(3)). Let $X$ be flag
data of rank $r$ over $E$ with set $\Psi$, integer $p_0$ and module $\Lambda$, put $q_0=r-p_0$, let $n\ge0$,
and let $a$ assign an integer $a(\tau',i)\ge0$ to each embedding $\tau'$ of $E$ and each $i\ge1$. The
\emph{variables} of $X$ are the $Y_{\sigma,b,k}$ with $\sigma\in\Psi$, $1\le b\le r$ and $1\le k\le n$; the
\emph{letter variables} are the pairs $(\ell,k)$ of a letter and $1\le k\le n$. A \emph{letter dictionary}
for $X$, $n$ and $a$ consists of a symplectic basis of $\Lambda$, an indexing $j\mapsto(\tau_j,a_j)$ and a
bijection $j\mapsto\ell_j$ onto the letters, the last two matched for $s$
(Definition~\ref{is-eigenvalue-conjugate-frame-fibres}), and integers $k_{\sigma,b}\ge0$ ($\sigma\in\Psi$, $1\le
b\le q_0$) and $k'_{\sigma,b}\ge0$ ($1\le b\le p_0$), with the following properties. Its \emph{letter map}
sends the letter variable $(\ell_j,k)$ to $Y_{\tau_j,a_j,k}$.
\begin{enumerate}
\item (Weights.) For $\sigma\in\Psi$, $\sum_{b'\ge b}k_{\sigma,b'}=a(\tau',b)$ for $1\le b\le q_0$, where
  $s\tau'=\sigma$, and $\sum_{b'\ge b}k'_{\sigma,b'}=a(\tau'',b)$ for $1\le b\le p_0$, where
  $s\tau''=\bar\sigma$.
\item (Polynomials.) The polynomials in the letter variables over $K_c$ that are harmonic, invariant under
  the row unipotents and of the row degrees given by $a$ are carried, by $s$ on the coefficients and the
  letter map on the variables, into the space of polynomials in the variables of $X$ that are harmonic for
  the pairing of the rows $b>p_0$ with the rows $b\le p_0$ at each $\sigma$, homogeneous in each row of the
  degrees given by the partial sums of $k$ and $k'$, and invariant under the row transvections, and they
  span this space over $\mathbf C$.
\end{enumerate}
\end{definition}
\leanref{Development/AbelianVarietyTheory/Definitions/SpanningFamilies/IsStabForm.lean}{GlobalTheory.Dict.Dic}
\leanref{Formalism/LocalToGlobal/ProjectiveLattices/RowDeg.lean}{Formalism.L2.Pj.HHa}
\leanref{Development/AbelianVarietyTheory/Definitions/ThetaValues/ExistsMemRowHighestTotalDegreeEqTwo.lean}{GlobalTheory.Harm.rowHighest}
% Dic Dw c sigma X g q0 n a: hr (p0 + q0 = r), bas, idx, sympl, e, the letter variables eL : LtN Dw n = Harm.Var
% X.Psi q0 p0 (Fin n) with dir_eq (letterDir (eL l) is the unit vector at (k, e^-1 l), so eL is the letter map of
% the book once the harmonic variable (sigma, b, k) is read as Y at the row rowIdx b), k, k', per_eq and ty_eq (the
% matching), wt_eq (wtOf X hr k k' n = wtX c sigma X a n, which is (1) through partOf), mem and span (HHa against
% Harm.rowHighest k k'), (2).

\begin{lemma}\label{is-eigenvalue-conjugate-frame-letters}
In Situation~\ref{is-eigenvalue-conjugate-frame-uniformization}, suppose that $r=p+q$, and let $X$ be flag
data of rank $r$ over $E$ whose integer $p_0$ is $q$, whose set $\Psi$ is the CM type $\Phi$ of the datum
(the embeddings $\tau$ with $\tau(\sqrt{-d})=\theta$, Situation~\ref{cv-function-situation}), and whose
module $\Lambda$ is self-dual for $X$. Let $j\mapsto\ell_j$ be a bijection from $\{1,\dots,g\}$ onto the
letters, $n\ge0$ an even integer (the parity of $n$ is not used) and $a$ as in
Definition~\ref{is-eigenvalue-conjugate-frame-letter-dictionary} with $a(\tau',i)\ge a(\tau',i+1)$ for all
$\tau'$ and $i$. Then there is a letter dictionary for $X$, $n$ and $a$ with the bijection $j\mapsto\ell_j$; if
moreover $p>0$ and $a(\tau',m_{\tau'})>0$ for every $\tau'\in\Phi$ with $m_{\tau'}>0$, it can be chosen with
$k_{\sigma,p}\ge1$ for some $\sigma$. Only the datum, the level, $K_c$, the good embedding $s$ and $r$ of the
situation are used.
\end{lemma}
\leanref{Development/AbelianVarietyTheory/Proofs/SiegelThetaTransformation/ExistsLetterRowWeightDictionary.lean}{GlobalTheory.Dict.exists_dic}
\leanref{Development/AbelianVarietyTheory/Proofs/UnitaryExistence/RelabelEmbeddingsGoodBijectiveConjCompat.lean}{GlobalTheory.Relab.relab_good}
\leanref{Development/AbelianVarietyTheory/Proofs/SiegelThetaTransformation/ExistsLetterRowWeightDictionary.lean}{GlobalTheory.Dict.Letters.exists_isSympl}
\leanref{Development/AbelianVarietyTheory/Proofs/CuspFormEvaluations/ThetaLetterEquivConjugate.lean}{GlobalTheory.Harm.mem_HHa_iff_rename}
\leanref{Development/AbelianVarietyTheory/Proofs/CuspFormEvaluations/PushLettersIntCastOfIsHomogeneous.lean}{GlobalTheory.Harm.span_map_rowHighest}
\begin{proof}
Write $\bar\tau'$ for the complex conjugate of an embedding $\tau'$ of $E$; as $E$ is a CM field,
$\bar\tau'(x)=\tau'(\bar x)$ for $x\in E$. Every embedding of $E$ is $t$ or $\bar t$ for exactly one $t\in\Phi$,
and $m_t=p$ and $m_{\bar t}=q$ for $t\in\Phi$, as the datum has signature $(p,q)$. Put $q_0=r-p_0=p$.

\emph{Relabelling.} The map $\tau'\mapsto s\tau'$ is injective, as $s$ is, so it is a bijection of the
finite set of embeddings of $E$. It preserves $\Phi$: $\tau'(\sqrt{-d})$ is a square root of $-d$ in $K_c$, so
it is $\pm\theta$, and $s\tau'(\sqrt{-d})=s(\pm\theta)=\pm\theta$ with the same sign, as the good embedding $s$
fixes $\theta\ne0$. And $s\bar\tau'=\overline{s\tau'}$, since $s\bar\tau'(x)=s\tau'(\bar x)$ for $x\in E$. So
$t\mapsto st$ is a bijection of $\Phi$ onto $\Psi=\Phi$.

\emph{The indexing.} For $1\le j\le g$ put $(\tau_j,a_j)=(st,p_0+i)$ if $\ell_j=(t,i)$ with $t\in\Phi$ and
$1\le i\le p$, and $(\tau_j,a_j)=(st,p_0+1-i)$ if $\ell_j=(\bar t,i)$ with $t\in\Phi$ and $1\le i\le q=p_0$. As
$t\mapsto st$ is a bijection of $\Phi$ onto $\Psi$ and $p+q=r$, the map $j\mapsto(\tau_j,a_j)$ is a bijection
onto $\Psi\times\{1,\dots,r\}$, an indexing; in particular $g=|\Psi|r$.

\emph{The matching.} Let $x$ be a family of invertible matrices and $v\in E^r$. If $\ell_j=(t,i)$ with
$t\in\Phi$, then $a_j>p_0$, so the row type at $j$ is $\tau_j=st$ and
$p_x(v)_j=(x_{st}^{-1}st(v))_{p_0+i}$, the period of $v$ at the letter $(t,i)$, as $st\in\Psi$. If
$\ell_j=(\bar t,i)$ with $t\in\Phi$, then $a_j\le p_0$, so the row type at $j$ is $\overline{st}=s\bar t$ and
$p_x(v)_j=h_{st}(x_{st}e_{p_0+1-i},st(v))$, the period of $v$ at the letter $(\bar t,i)$, as the conjugate $st$ of
$s\bar t$ lies in $\Psi$. So the indexing and $j\mapsto\ell_j$ are matched for $s$.

\emph{The symplectic basis.} As a finitely generated $\mathcal O_E$-submodule that spans $E^r$, $\Lambda$ is a
free $\mathbf Z$-module of rank $[E:\mathbf Q]\,r=2|\Psi|r=2g$. As $\Lambda$ is self-dual,
$\langle u,v\rangle_X\in\mathbf Z$ for $u,v\in\Lambda$. The form is alternating: $h(v,u)=\overline{h(u,v)}$ and
$\bar\kappa_0=-\kappa_0$, so $\langle v,u\rangle_X=-\langle u,v\rangle_X$. It is perfect on $\Lambda$. The form
$(u,w)\mapsto\operatorname{Tr}_{E/\mathbf Q}(-\frac12\kappa_0h(u,w))$ on $E^r$ is nondegenerate, as $\kappa_0\ne0$
and $\det A\ne0$, and $\Lambda$ spans $E^r$ over $\mathbf Q$. So $u\in\Lambda$ is determined by the values
$\langle u,w\rangle_X$, $w\in\Lambda$; and a homomorphism $\Lambda\to\mathbf Z$ extends to a $\mathbf Q$-linear form
on $E^r$, which is $\langle u,\cdot\rangle_X$ for some $u\in E^r$, and $u\in\Lambda$ by self-duality.
Lemma~\ref{cv-uniformization-adapted-basis}, applied to $B=-\langle\cdot,\cdot\rangle_X$ with $N=1$ and
$P=Q=\Lambda$, gives a basis $(e_1,\dots,e_g,f_1,\dots,f_g)$ of $\Lambda$ with $B(e_a,e_b)=B(f_a,f_b)=0$ and
$B(e_a,f_b)=\delta_{ab}$. Then $b_k=e_k$ and $b'_k=f_k$ form a symplectic basis of $\Lambda$.

\emph{The weights.} For $t\in\Phi$ put $k_{st,b}=a(t,b)-a(t,b+1)$ for $1\le b<p$ and $k_{st,p}=a(t,p)$, and
$k'_{st,b}=a(\bar t,b)-a(\bar t,b+1)$ for $1\le b<q$ and $k'_{st,q}=a(\bar t,q)$. These integers are $\ge0$, as
$a(t,\cdot)$ and $a(\bar t,\cdot)$ decrease, and the sums telescope:
$\sum_{b'\ge b}k_{st,b'}=a(t,b)$ and $\sum_{b'\ge b}k'_{st,b'}=a(\bar t,b)$. As $s\bar t=\overline{st}$, this is
(1).

\emph{The polynomials.} The letter map sends the letter variables at $t$ and at $\bar t$, $t\in\Phi$, to the
variables of the rows $p_0+i$ and $p_0+1-i$ of $X$ at $st$. This renaming of the variables carries each of
the three conditions on a polynomial in the letter variables over $K_c$ (harmonic, invariant under the row
unipotents, of the row degrees given by $a$) to the corresponding condition of (2) over $K_c$ (harmonic for
the pairing of the rows, invariant under the row transvections, homogeneous in each row of the degree given by
the partial sums of $k$ and $k'$, which are the values of $a$ by the last paragraph). So the letter map is a
bijection from the polynomials over $K_c$ with the three properties onto the polynomials over $K_c$ in the
space of (2). The conditions defining that space are linear equations with integer coefficients on the
coefficients of a polynomial: the Laplacian of the pairing, the homogeneity, and the invariance under the
transvections by integers. So applying $s$ to the coefficients carries the polynomials over $K_c$ in the space
into the space over $\mathbf C$. Conversely, let $P$ lie in the space over $\mathbf C$, choose a basis $(c_\alpha)$
of $\mathbf C$ over $s(K_c)$, and write $P=\sum_\alpha c_\alpha P_\alpha$ with polynomials $P_\alpha$ with
coefficients in $s(K_c)$. Each linear condition, applied to $P$, gives $\sum_\alpha c_\alpha L(P_\alpha)=0$ with
$L(P_\alpha)$ over $s(K_c)$, so $L(P_\alpha)=0$ for every $\alpha$. So each $P_\alpha$ is $s$ applied to a
polynomial in the space over $K_c$, and $P$ lies in the span of the images. This is (2).

\emph{The last clause.} Suppose that $p>0$ and $a(\tau',m_{\tau'})>0$ for every $\tau'\in\Phi$ with
$m_{\tau'}>0$. The set $\Phi$ is not empty: for an embedding $\tau'$ of $E$, $\tau'$ or $\bar\tau'$ lies in
$\Phi$. For $t\in\Phi$, $m_t=p>0$, so $k_{st,p}=a(t,p)\ge1$.
\end{proof}
% Dict.exists_dic (100 lines): relab_good for GoodEmb (the bijection hbij, the type hrt, the conjugation hcj),
% f : {tau // InPsi H tau} = X.Psi through relab and hPsi, rho = halfRowEquiv (letters (t, i) with InPsi t to
% (t, inl i), letters (conj t, i) to (t, inr i), with mult_of_mem and mult_conjEmb_of_mem), idx = e then rho then
% rowOf (rowIdx: inl i to p0 + i, inr i to p0 - 1 - i, counted from 0); per_eq and ty_eq by per_of_le, per_of_lt
% and perFr, rowType_rowOf_inl and _inr; bas from Letters.exists_isSympl (ThetaHol.finrank_lattice, the integral
% values (hL.2), TensLat.EF_swap, eq_of_EFf_eq and exists_traceFormDual_eq, and
% exists_symplectic_coordinates_adapted with I = 0 and eps = -1; the book cites Lemma cv-uniformization-adapted-basis
% with N = 1 for the same basis); k0 = diffOf of a(t, .), k0' = diffOf of a(conj t, .), with partOf_diffOf for wt_eq
% (wtOf_inl, wtOf_inr, wtX); eL = typeLetterEquiv then f, dir_eq from letterDir; hmem from mem_HHa_iff_rename and
% Letters.rename_place_mem_iff; mem by map_mem_rowHighest, span by span_map_rowHighest (basis components) and
% rename_surjective; the last clause from LastPos at t with mult t = p. The second half of Dom a n (n even) is
% not used.

\begin{lemma}\label{is-eigenvalue-conjugate-frame-analytic-fibre}
In Situation~\ref{is-eigenvalue-conjugate-frame-uniformization}, fix a uniformization witness, with its
bijection $j\mapsto\ell_j$ onto the letters. Let $c$ be a component, fix a symplectic basis of $\Lambda_c$
and an indexing matched with $j\mapsto\ell_j$ for $s$, and let $x\in\mathcal D_c$ with $|d_i(x,\tau)|=1$
for all $\tau$ and $i$ (Definition~\ref{is-onto-share-flag-data}). Then there are an analytic fibre at
$y(c,x)$ for these data, $x$ and $\Lambda_c$ that is marked by the marking $m_c$ of $c$, and a theta reading
of $s_{y(c,x)}$ for it that has mean one for $x$ and $\Lambda_c$.
\end{lemma}
\leanref{Development/AbelianVarietyTheory/Proofs/SiegelThetaTransformation/AnalyticFibreUnitFrameUniformisation.lean}{GlobalTheory.Dict.uw_anFib}
\begin{proof}
Write $X=X_c$, $\Lambda=\Lambda_c$ and $y=y(c,x)$. Let $P=P_x$ and $\operatorname{ex}=\operatorname{ex}_x$ be as
in (ii) of the situation, and, as $|d_i(x,\tau)|=1$ for all $\tau$ and $i$, let $H$, $\operatorname{fr}_\beta$,
$\chi_0$ (the function $\chi$ there), $\vartheta$ and $\mathcal F$ be as in (iii).

\emph{The periods.} As the indexing is matched with $j\mapsto\ell_j$, the $j$-th coordinate of $p_x(\varphi_ca)$
is the period of $\varphi_c(a)$ at the letter $\ell_j$, so $P(a)=p_x(\varphi_ca)$ for $a\in L_c$. As $\varphi_c$
is onto $\Lambda$, $P(L_c)=\{p_x(v):v\in\Lambda\}$.

\emph{The form.} By Lemma~\ref{is-eigenvalue-conjugate-frame-component-lattice}(3),
$\langle a,b\rangle_c=\langle\varphi_ca,\varphi_cb\rangle_X$ for $a,b\in L_c$, and $\Lambda=\varphi_c(L_c)$ spans
$E^r$ over $E$. So $\operatorname{Im}H(p_x(u),p_x(v))=\langle u,v\rangle_X$ for $u,v\in\Lambda$, and
$H(z,w)=H_x(z,w)$ for all $z,w\in\mathbf C^g$ by Lemma~\ref{is-eigenvalue-conjugate-frame-hermitian}(4).

\emph{The analytic fibre.} By Theorem~\ref{cv-uniformization-family}, $\operatorname{ex}$ is a surjective
holomorphic group homomorphism $\mathbf C^g\to A_y(\mathbf C)$ with kernel $P(L_c)=p_x(\Lambda)$, and
$\operatorname{ex}((\sigma_j(\alpha)z_j)_j)=\alpha\cdot\operatorname{ex}(z)$ for $\alpha\in\mathcal O_E$, where
$\sigma_j=s\circ\varepsilon_{\tau'_j}$ is the type of $P$; as the indexing is matched, $\sigma_j$ is the row type
$\rho_j$. Its derivative at $0$ in the direction of the $j$-th standard basis vector is $u_{\ell_j}$, by (ii).
This is (a) of Definition~\ref{is-eigenvalue-conjugate-frame-fibres}(3). The functions $\operatorname{fr}_\beta$
satisfy (b), by (iii). For $v\in E^r$ put $\chi(v)=\chi_0(p_x(v))$ if $|\chi_0(p_x(v))|=1$, and $\chi(v)=1$
otherwise. For $v\in\Lambda$, $v=\varphi_ca$ with $p_x(v)=P(a)$, so $|\chi_0(p_x(v))|=1$ and $\chi(v)=\chi_0(P(a))$;
and the law (c) at $v$ is the law of (iii) at $a$, as $H=H_x$. So $\operatorname{ex}$, the $\operatorname{fr}_\beta$
and $\chi$ form an analytic fibre at $y$ for the data, $x$ and $\Lambda$.

\emph{The marking.} By Theorem~\ref{cv-uniformization-family}, the level structure of $y$ sends
$a\in\bar{\mathfrak M}\mathfrak L$ to $\operatorname{ex}(M^{-1}P(t_ca))$. As $p_x(\alpha v)_j=\rho_j(\alpha)p_x(v)_j$
for $\alpha\in E$ (Lemma~\ref{is-eigenvalue-conjugate-frame-hermitian}(3)) and $\rho_j(M^{-1})=M^{-1}$,
$p_x(M^{-1}m_c(a))=M^{-1}p_x(\varphi_c(t_ca))=M^{-1}P(t_ca)$. So the analytic fibre is marked by $m_c$.

\emph{The theta reading.} By (iii), $\vartheta$ is a theta reading of $s_y$ for this analytic fibre. The set
$\mathcal F$ is a fundamental domain of $P(L_c)=p_x(\Lambda)$, and
$\operatorname{Re}H(z,z)=\operatorname{Re}H_x(z,z)$; so the normalized integral of (iii) is that of
Definition~\ref{is-eigenvalue-conjugate-frame-fibres}(7), and $\vartheta$ has mean one for $x$ and $\Lambda$.
\end{proof}
% Dict.uw_anFib (90 lines; AnalyticFibreUnitFrameUniformisation.lean, 125 lines), under GoodEmb, IsUnf, 1 <= r, a
% UW structure W, IsSympl and the two clauses of Matched for W.e, at a point with IsUnitFr; the fibre is
% A.Marked (W.mark k). It follows the proof above: W.pt k x and its last clause hfr (IsUnitFr); exists_ex_transport
% for ex (F.ex with ex_add, ex_onto, ex_ker, ex_lin and ex_eta); hPv from hpv; uniformisation_component_lattices_
% markings (hspan, hEF); hx_laws (hsm for the smul clause, hL for H = H_x); famHol for the holomorphy, hder, the
% chi with the if, hM for the marking, hread for the reading and MeanOne with Pm.range.

\begin{lemma}\label{is-eigenvalue-conjugate-frame-points}
In Situation~\ref{is-eigenvalue-conjugate-frame-core}:
\begin{enumerate}
\item every $h\in U(W_i)(\mathbf A)$ has a framing whose symplectic basis, indexing and bijection onto the
  letters are those of a letter dictionary for $X_{c(h)}$, $n$ and the weights of $\kappa_i$
  (Definition~\ref{is-eigenvalue-conjugate-frame-letter-dictionary}), and whose element $\gamma$ has the
  following property: for every symplectic basis of $\Lambda_{c(h)}$, indexing and bijection onto the
  letters, the last two matched for $s$, every analytic fibre at $y(c(h),x(h))$ for these data, $x(h)$ and
  $\Lambda_{c(h)}$ is marked by every marking at $\gamma h$ with values in $\Lambda_{c(h)}$
  (Definition~\ref{is-eigenvalue-conjugate-frame-marking});
\item $Z_0$ is a signature family for $T=W_i$ and $p_0=q$
  (Situation~\ref{is-eigenvalue-conjugate-frame-situation}): $x_{0,\rho}^*G_{W_i}(\rho)x_{0,\rho}$ is the
  diagonal matrix with entries $-1$ at the first $q$ places and $1$ at the others; and the Cartan
  involution of the model of $W_i$ at $\sigma$ at every $\rho$ is $s_{Z_0,\rho}$;
\item for every $g\in U(W_i)(\mathbf A)$ there is a self-dual lattice $L'_g$ of $W_i$ with
  $g_f\hat\Lambda_W=\hat L'_g$ and $L(g)=\mathfrak AL'_g$, such that the lattice of the pair given by
  $\Lambda_V$ and $L'_g$ is the image under $u\mapsto gu$ of the lattice of the pair given by $\Lambda_V$
  and $\Lambda_W$;
\item $\mu$ is multiplicative and trivial on $U(W_i)(F_h)$; for every $\sigma$-conjugate $\chi'$ of
  $\mu^0_i$ (Lemma~\ref{is-eigenvalue-type-parts}(3)), every $g\in U(W_i)(\mathbf A)$ has a determinant
  root and $\mu(g)=\chi'(y)$ for every determinant root $y$ of $g$
  (Lemma~\ref{is-eigenvalue-conjugate-frame-det-roots}); and for every $\sigma$-conjugate $\eta'$ of
  $\eta$, $\chi'$ and $\chi'\eta'^{-1}$ are trivial on every idele that is a unit at every finite place
  and differs from $1$ only at the places above one place of $F_h$ that does not split in $E_h$;
\item $V_i$ has Gram matrix $1$ and even rank, $\Lambda_V=\mathcal O_{E_h}^n$, $\Lambda_V$ and
  $\Lambda_W$ are self-dual, and $r\ge1$.
\end{enumerate}
\end{lemma}
\leanref{Development/AbelianVarietyTheory/Proofs/DivisorProofs/CoefMapNormEquivariantBasisOfFive.lean}{GlobalTheory.GalOne.framed_point_at_ofFive}
\leanref{Development/AbelianVarietyTheory/Proofs/DivisorProofs/FrAtEmb.lean}{GlobalTheory.GalOne.PointAt}
\leanref{Development/AbelianVarietyTheory/Proofs/DivisorProofs/CoefMapNormEquivariantBasisOfFive.lean}{GlobalTheory.OneMapRef.FactsLat.facts_lat_ofFive}
\leanref{Development/AbelianVarietyTheory/Proofs/IwasawaDecomposition/OmegaFormIntegrality.lean}{GlobalTheory.OneMapRef.FactsLat.isONeg_x0Mat}
\leanref{Development/AbelianVarietyTheory/Proofs/IwasawaDecomposition/OmegaFormScalarExtension.lean}{GlobalTheory.OneMapRef.FrameConst.mu_rat_cast}
\begin{proof}
The embedding $s$ is good, by Lemma~\ref{is-eigenvalue-conjugate-extension}(2) and
Lemma~\ref{is-eigenvalue-conjugate-good}(3), as $\sigma$ agrees with $\sigma_0$ on $K$.

(2) By the core conditions on the models and the pins of the core, the signs of the model of $W_i$ at
$\sigma$ at $\rho$ are $1$ at the first $p$ indices and $-1$ at the others, where $p+q=r$. The columns of
$f_\rho^{-1}$, for the frame $f_\rho$ of the model, have Gram matrix the diagonal matrix $J_\rho$ of these
signs for $G_{W_i}(\rho)$, and $x_{0,\rho}$ puts the last $q$ of them first, in the reverse order
(Definition~\ref{is-eigenvalue-conjugate-dictionary}(1)). So $x_{0,\rho}^*G_{W_i}(\rho)x_{0,\rho}=N$ with
$N$ as in Situation~\ref{is-eigenvalue-conjugate-frame-situation} for $p_0=q$, and the Cartan involution
$f_\rho^{-1}J_\rho f_\rho$ of the model is $x_{0,\rho}Nx_{0,\rho}^{-1}=s_{Z_0,\rho}$.

(3) The lattices $\Lambda_V$ and $\Lambda_W$ are the lattices of the frames of $V_i$ and $W_i$, by the
core conditions on the level characters. Let $L'_g$ be the $\mathcal O_{E_h}$-span of the $v\in E_h^r$
with $g_f^{-1}v\in\hat\Lambda_W$. It is a self-dual lattice of $W_i$ with completion $g_f\hat\Lambda_W$,
$L'_1=\Lambda_W$, and $L(g)$ is the span of the products $\alpha v$ with $\alpha\in\mathfrak A$ and $v\in
L'_g$. An element $u$ of $\mathbb W$ lies in the lattice of the pair given by $\Lambda_V$ and $L'_g$
exactly when $g^{-1}u$ lies in the lattice of the pair given by $\Lambda_V$ and $\Lambda_W$, as the
finite part of $g^{-1}$ carries $\hat L'_g$ onto $\hat\Lambda_W$.

(4) The weight tie requires $\mu$ to be multiplicative and trivial on $U(W_i)(F_h)$, and every element
to have a determinant root at which the character of the core at the level $i+1$ and $s$ takes the value
of $\mu$; this character is a $\sigma$-conjugate of $\mu^0_i$, and $\sigma$-conjugates are unique, so it is
$\chi'$. Let $x$ be an idele as in the last assertion. By the core conditions on the level characters,
$\mu^0_i\xi_{\mu_0}$ is trivial on $x$ and $\xi_{\mu_0}=1$; as $\chi'(x)=\sigma(y)$ for some $y\in L$
with $\mu^0_i(x)=\sigma_0(y)$, $y=1$ and $\chi'(x)=1$. And $\eta'(x)=1$, as $\eta'$ is unramified everywhere
(Lemma~\ref{is-eigenvalue-type-parts}(4)).

(5) The pins give $V_i$ the Gram matrix $1$ and the standard lattice $\mathcal O_{E_h}^n$; $n$ is even,
by the parities of the ranks of $D$; the lattices of the adelic matrices of $V_i$ and $W_i$ are self-dual,
by the core conditions on the level characters; and $r\ge1$, by the ranks of $D$.

(1) By the laws of the weight tie at the level $i+1$ and the good embedding $s$, every component has
matrix $-G_{W_i}^{\mathsf T}$ and element $-\xi_0$, every $x(h)$ lies in $\mathcal D_{c(h)}$ with
$|d_i(x(h),\tau)|=1$ for all $\tau$ and $i$, and the uniformization of the core at $s$ satisfies the
compatibility conditions of the moduli problem. With the datum of the core at level $i+1$, its level, the
coefficient field $K_s$, the good embedding $s$ and $r\ge1$ (by (5)), this is
Situation~\ref{is-eigenvalue-conjugate-frame-uniformization}. Fix a uniformization witness, with its
bijection $j\mapsto\ell_j$ onto the letters. Let $h\in U(W_i)(\mathbf A)$, $c=c(h)$ and $x=x(h)$. The
following facts hold.
\begin{enumerate}
\item[(a)] There is $\gamma\in U(W_i)(F_h)$ with $\Lambda_c=L(\gamma h)$ and $x=x^{\gamma hZ_0}$ that has the
  property of (1), by Lemma~\ref{is-eigenvalue-conjugate-frame-dictionary-point}, as $x\in\mathcal D_c$.
\item[(b)] The element $\kappa_0$ of $X_c$ is twice its element $-\xi_0$, so $\kappa_0=-2\xi_0$, and
  $\Lambda_c$ is self-dual for $X_c$, by Lemma~\ref{is-eigenvalue-conjugate-frame-component-lattice}(1).
\item[(c)] $r=p+q$, $\Psi_c=\mathcal T$, and the integer $p_0$ of $X_c$ is $q$. Indeed, by the
  compatibility conditions of the uniformization, $r=p+q$ and $p_0=q$ for the signature $(p,q)$ of the datum
  of the core at level $i+1$, and $\Psi_c$ is its CM type, the set of embeddings $\tau$ of $E_h$ with
  $\tau(\sqrt{-d})=\theta$ for the square root $\theta$ of $-d$ given by the host of the core; by the pins of
  the core on its host, this set is $\mathcal T$.
\item[(d)] There is a letter dictionary for $X_c$, $n$ and the weights of $\kappa_i$ with the bijection
  $j\mapsto\ell_j$, by (b), (c) and Lemma~\ref{is-eigenvalue-conjugate-frame-letters}, applied to $X_c$, as
  $n$ is even (5) and the weights of $\kappa_i$ decrease at each embedding.
\item[(e)] For these data there are an analytic fibre at $y(c,x)$ and a theta reading of $s_{y(c,x)}$ for
  it that has mean one for $x$ and $\Lambda_c$, by Lemma~\ref{is-eigenvalue-conjugate-frame-analytic-fibre},
  as $|d_i(x,\tau)|=1$ for all $\tau$ and $i$.
\end{enumerate}
% (a) is T.dictionary_rational_element_analytic_point (with its marking clause hmk); (b) is Dict.uw_lattice for a
% UW structure W from nonempty_uw; (c) is the fields hU.1 and hU.2.1 of IsUnf (r = Dw.p + Dw.q, p0 = Dw.q, and tau
% in Psi iff tau(H.rt) = H.rt0) with OneMapRef.exists_type_equiv (the host pin hp.2.2.2 and
% Pf.exists_tau_of_inPsi); (d) is Dict.exists_dic with e = W.e; (e) is Dict.uw_anFib for W, with IsUnitFr from
% laws_cast. The laws are OneMapRef.laws_cast (ModelW, dom, IsUnitFr and the uniformization law unf).

With (a) to (e), the data satisfy the conditions of a framing at $h$ with element $\gamma$: $\psi$ is
standard, as the Weil datum of $R$ requires; $\theta$ is an archimedean family whose members extend the
$\rho\in\mathcal T$, by the core conditions on the models; $V_i$ has Gram matrix $1$ and even rank, by (5);
$X_c$ has matrix $-G_{W_i}^{\mathsf T}$, element $-\xi_0$, $\kappa_0=-2\xi_0$ and $\Psi_c=\mathcal T$, by the
laws and (b), (c); and the rest is (a) to (e). The property of $\gamma$ is that of (a). This proves (1).
\end{proof}
% framed_point_at_ofFive (PointAt with its marking clause; PointAt also packs the basis, indexing and bijection
% of the framing into a letter dictionary Dic for n = rk V_i and the weight wtA (kappa i), of which the book keeps
% the basis, the indexing, the bijection and the matching): facts_lat_ofFive (hg1, hev, hn, hr1),
% factsAt_ofFive (hsigma, hx0), laws_cast, T.dictionary_rational_element_analytic_point, nonempty_uw,
% uw_lattice, the field hU.2.1, exists_type_equiv, exists_dic and uw_anFib, assembled into IsSet (with
% (hC.thP.omegaP i A.mu0).std and hC.models.inf), Matched, p0 = q, SelfD, latW = latG, the point and the
% analytic fibre. (2) is factsAt_ofFive (FactsLat.isONeg_x0Mat and cart_eq_cartZ from signs_of_overK and
% the pin on q). (3) is the last clause of facts_lat_ofFive (MovedLat.moved_lattices: L'_g = hLat g, with
% isAA_span and isMov_hLat). (4) is the muUnit law (mu_rat_cast) and the last clause of factsAt_ofFive
% (FactsLat.muVal_cast with tie.laws.muVal, tie.chiR and ConjCharUnique.conjChar_unique_ofFive;
% unramAt_conj_ofFive). (5) is facts_lat_ofFive (hP.lift.std, the parities k_even, m_even, t_odd,
% SDFrameLattice.isSelfDualC_of_isSDFrame for betaV and betaW).

\begin{lemma}\label{is-eigenvalue-conjugate-frame-ortho-constant-norm-theta-mean}
Let $g\ge0$, $\tau\in M_g(\mathbf C)$ symmetric with $Y=\operatorname{Im}\tau$ positive definite, $h\in GL_g(\mathbf C)$,
$c_1,c_2\in\{0,\frac12\}^g$, and $A$ the group of the $\lambda\in\mathbf C^g$ with $h\lambda\in\mathbf Z^g+\tau\mathbf Z^g$. Then $A$
has a fundamental domain $F$ in $\mathbf C^g$ for Lebesgue measure with
\[
 \frac1{\operatorname{vol}(F)}\int_F\big|\vartheta[c_1,c_2](\tau,hz)\big|^2e^{-2\pi\,(\operatorname{Im}hz)^{\mathsf
 T}Y^{-1}(\operatorname{Im}hz)}\,dz=(2^g\det Y)^{-1/2},
\]
with the theta function of Theorem~\ref{cv-uniformization-theta-transformation}.
\end{lemma}
\leanref{Development/AbelianVarietyTheory/Proofs/UnitaryExistence/ExistsFundamentalDomainMeanNormSqThetaChar.lean}{GlobalTheory.ThetaMean.exists_fundamentalDomain_mean_norm_sq_thetaChar}
\begin{proof}
\openstep{The proof, which changes variables by $h$ to the lattice $\mathbf Z^g+\tau\mathbf Z^g$, integrates $|\vartheta|^2$ over the
real directions by orthogonality of the characters $e^{2\pi i\,m\cdot x}$, and evaluates the remaining Gaussian integral over the
imaginary directions, is not written out.}
% ThetaMean.exists_fundamentalDomain_mean_norm_sq_thetaChar (ExistsFundamentalDomainMeanNormSqThetaChar.lean). Size: 800
% lines (lean_size).
\end{proof}

\begin{lemma}\label{is-eigenvalue-conjugate-frame-ortho-constant-norm-mean}
Let $g$, $\tau$, $Y$, $h$, $c_1$, $c_2$ and $A$ be as in
Lemma~\ref{is-eigenvalue-conjugate-frame-ortho-constant-norm-theta-mean}, put $H(z,w)=(hz)^{\mathsf T}Y^{-1}\overline{hw}$
for $z,w\in\mathbf C^g$, and let $\vartheta$ be a function on $\mathbf C^g$ and $a\in\mathbf C$ with
\[
 \vartheta(z)=a\exp\Big(\frac\pi2(hz)^{\mathsf T}Y^{-1}(hz)\Big)\vartheta[c_1,c_2](\tau,hz)
\]
for every $z$. Suppose that for every $p\in A$ there is $\varepsilon\in\mathbf C$ with $|\varepsilon|=1$ and
$\varepsilon e^{-\pi H(z+\frac12p,p)}\vartheta(z+p)=\vartheta(z)$ for every $z$, and that some fundamental domain $F$ of $A$
has $\operatorname{vol}(F)^{-1}\int_F|\vartheta(z)|^2e^{-\pi\operatorname{Re}H(z,z)}\,dz=1$. Then $|a|^2=(2^g\det Y)^{1/2}$.
\end{lemma}
\leanref{Development/AbelianVarietyTheory/Proofs/SiegelThetaTransformation/ExistsSymplOfSelfDual.lean}{GlobalTheory.Dict.VacNorm.norm_sq_of_mean}
\begin{proof}
Put $f(z)=|\vartheta(z)|^2e^{-\pi\operatorname{Re}H(z,z)}$. As $Y^{-1}$ is real and symmetric, $\operatorname{Re}H(p,z)=
\operatorname{Re}H(z,p)$, so $\operatorname{Re}H(z+p,z+p)=\operatorname{Re}H(z,z)+2\operatorname{Re}H(z,p)+\operatorname{Re}H(p,p)$; and for
$p\in A$ the law gives $|\vartheta(z+p)|^2=e^{2\pi\operatorname{Re}H(z,p)+\pi\operatorname{Re}H(p,p)}|\vartheta(z)|^2$. So $f$ is
invariant under $A$. Let $F'$ be the fundamental domain of
Lemma~\ref{is-eigenvalue-conjugate-frame-ortho-constant-norm-theta-mean}; then $\int_Ff=\int_{F'}f$ and
$\operatorname{vol}(F)=\operatorname{vol}(F')$. For $\zeta\in\mathbf C^g$,
$\operatorname{Re}(\zeta^{\mathsf T}Y^{-1}\zeta)-\operatorname{Re}(\zeta^{\mathsf T}Y^{-1}\bar\zeta)=-2(\operatorname{Im}\zeta)^{\mathsf
T}Y^{-1}(\operatorname{Im}\zeta)$; with $\zeta=hz$ this gives
\[
 f(z)=|a|^2\,\big|\vartheta[c_1,c_2](\tau,hz)\big|^2e^{-2\pi(\operatorname{Im}hz)^{\mathsf T}Y^{-1}(\operatorname{Im}hz)}.
\]
So $1=\operatorname{vol}(F')^{-1}\int_{F'}f=|a|^2(2^g\det Y)^{-1/2}$, by
Lemma~\ref{is-eigenvalue-conjugate-frame-ortho-constant-norm-theta-mean}.
\end{proof}
% VacNorm.norm_sq_of_mean (ExistsSymplOfSelfDual.lean), followed: exists_fundamentalDomain_mean_norm_sq_thetaChar (Fd',
% hm'), hinv (e1 through re_dot_star_comm, e2), IsAddFundamentalDomain.setIntegral_eq and measure_eq, hint (re_dot,
% re_dot_star).

\begin{lemma}\label{is-eigenvalue-conjugate-frame-ortho-constant-norm-det}
Let $X$, a symplectic basis and an indexing be as in Definition~\ref{is-eigenvalue-conjugate-frame-fibres}(1) and~(2), with
set $\Psi$ and element $\xi_0$, and $c_\tau=2|\tau(\xi_0)|$; let $n$ be an even integer, which is not used. Let
$x\in\mathcal D_X$ with $|d_i(x,\tau)|=1$ for all $\tau$ and $i$, and $H\in M_d(\mathbf C)$ with $H_x(z,w)=z^{\mathsf T}H\bar w$
for all $z,w\in\mathbf C^d$. Then $\det H=\prod_{\tau\in\Psi}c_\tau^r$.
\end{lemma}
\leanref{Development/AbelianVarietyTheory/Proofs/SiegelThetaTransformation/DetHermitianFormAtFramePeriodCoordsOfUnitFrameStatement.lean}{GlobalTheory.T.det_hermitianFormAtFramePeriodCoords_of_unitFrame}
\leanref{Development/AbelianVarietyTheory/Proofs/SiegelThetaTransformation/DetHermitianFormAtFramePeriodCoordsOfUnitFrame.lean}{GlobalTheory.Dict.det_hx_of_isUnitFr}
\begin{proof}
By Lemma~\ref{is-onto-share-majorant}(3), $x=x'b$ with $x'$ orthonormal and $b\in\mathcal B^\Psi$. The Gram matrix of the columns
of $x_\tau$ for $h_\tau$ is $b_\tau^{\mathsf T}J\bar b_\tau$, whose leading minors are
$d_i(x,\tau)=\epsilon_i\prod_{a\le i}|(b_\tau)_{aa}|^2$; as these have absolute value $1$, $|(b_\tau)_{aa}|=1$ for all $\tau$ and $a$.
Let $N$ be the matrix $N_{b^{-1}}$ of the proof of Lemma~\ref{is-eigenvalue-conjugate-frame-hermitian-riemann}(2), with
$z_{x'}(u)=Nz_x(u)$; in a suitable order of the indices it is triangular with the diagonal entries $(b_\tau)_{aa}^{-1}$ and
$(b_\tau)_{aa}$, so $|\det N|=1$. As $xb^{-1}=x'$ is orthonormal, (4) of that lemma gives $H_x(z,w)=(Nz)^{\mathsf T}G\,\overline{Nw}$,
with $G$ diagonal with the entries $c_{\tau_j}$. So $H=N^{\mathsf T}G\bar N$, and $\det H=|\det N|^2\det G=\prod_\tau
c_\tau^r$.
\end{proof}
% det_hx_of_isUnitFr (DetHermitianFormAtFramePeriodCoordsOfUnitFrame.lean): UnitFrames.unit_frame_eq_orthonormal_mul_triangular
% (x = x0 b with x0 orthonormal and |b_ii| = 1; the book reads it from Lemma is-onto-share-majorant(3) and the minors),
% herm_mulVec_mulVec', ThFr.hermForm_Hm_rmul (Hm x0 = N^T Hm x conj N, N = borMat b), ThFr.det_Hm_of_isOrtho (det Hm
% x0 = prod c^r) and det_borMat (|det N| = 1); the book takes the last three from Lemma -hermitian-riemann(2) and (4).

\begin{lemma}\label{is-eigenvalue-conjugate-frame-ortho-constant-norm}
Let $X$, a symplectic basis and an indexing be as in Definition~\ref{is-eigenvalue-conjugate-frame-heisenberg}, let
$n\ge0$ be even, $w_0\in E$ an algebraic integer with $w_0+\bar w_0=1$, $L\subset E^r$ an $\mathcal O_E$-submodule
self-dual for $X$, and $v$ a vacuum of $\mathcal O_E^n\otimes L$ for $w_0$. Then there is a real number $\rho>0$ such that
$|\kappa|=\rho$ whenever $x\in\mathcal D_X$ has $|d_i(x,\tau)|=1$ for all $\tau$ and $i$, $\vartheta$ is a level-one theta
function for $w_0$, $x$ and $L$ that has mean one for $x$ and $L$, and $\kappa\in\mathbf C$ satisfies
$\prod_{k=1}^n\vartheta(z_k)=\kappa\,\theta_v(x;z)$ for every $z\in(\mathbf C^d)^n$.
\end{lemma}
\leanref{Development/AbelianVarietyTheory/Proofs/SiegelThetaTransformation/ExistsSymplOfSelfDual.lean}{GlobalTheory.Dict.VacNorm.norm_factor_eq}
\begin{proof}
Write $\langle\ ,\ \rangle=\langle\ ,\ \rangle_X$, $d=|\Psi|r$ and $c_\tau=2|\tau(\xi_0)|>0$. As $L$ is self-dual, it is a free
$\mathbf Z$-module of rank $2d$ on which $\langle\ ,\ \rangle$ is integral and perfect, so, by the symplectic basis theorem as in
the proof of Lemma~\ref{is-eigenvalue-conjugate-frame-theta-product}, it has a basis $b_1,\dots,b_d,b'_1,\dots,b'_d$ with
$\langle b_k,b'_l\rangle=-\delta_{kl}$, $\langle b'_l,b_k\rangle=\delta_{kl}$ and $\langle b_k,b_l\rangle=\langle b'_k,b'_l\rangle=0$. Let
$c_1$, $c_2$ and $C\ne0$ be as in Lemma~\ref{is-eigenvalue-conjugate-frame-theta-product-vacuum} for this basis, and put
\[
 \rho=\Big(2^d\Big/\prod_\tau c_\tau^r\Big)^{n/4}\Big/|C|.
\]
Let $x$, $\vartheta$ and $\kappa$ be as in the statement, and $h$, $\tau$ and $Y$ as in
Lemma~\ref{is-eigenvalue-conjugate-frame-theta-product-vacuum} for $x$. Put
$F_0(z)=\exp\big(\frac\pi2(hz)^{\mathsf T}Y^{-1}(hz)\big)\vartheta[c_1,c_2](\tau,hz)$. By (3) of that lemma,
$\prod_k\vartheta(z_k)=K\prod_kF_0(z_k)$ for every $z$, with $K=\kappa C\det(h)^{n/2}$.

If $n=0$, this says $1=\kappa C$, so $|\kappa|=1/|C|=\rho$. Let $n\ge1$. As $\vartheta$ has mean one, $\vartheta(z_0)\ne0$ for some
$z_0$. With all $z_k=z_0$, $\vartheta(z_0)^n=KF_0(z_0)^n$, so $F_0(z_0)\ne0$; with $z_1=z$ and the other $z_k=z_0$,
$\vartheta(z)\vartheta(z_0)^{n-1}=KF_0(z)F_0(z_0)^{n-1}$. So $\vartheta=aF_0$ with $a=KF_0(z_0)^{n-1}/\vartheta(z_0)^{n-1}$, and
$a^nF_0(z_0)^n=\vartheta(z_0)^n=KF_0(z_0)^n$ gives $a^n=K$.

For $m',m''\in\mathbf Z^d$, $hp_x\big(\sum_km'_kb_k+\sum_km''_kb'_k\big)=m'+\tau m''$, as $h=\Omega_x^{-1}$ and $\tau=h\Omega'_x$. So the
group $A=\{p_x(u):u\in L\}$ is the group of the $\lambda$ with $h\lambda\in\mathbf Z^d+\tau\mathbf Z^d$. By
Lemma~\ref{is-eigenvalue-conjugate-frame-theta-product-vacuum}(1), $H_x(z,w)=(hz)^{\mathsf T}Y^{-1}\overline{hw}$. The identity of
Definition~\ref{is-eigenvalue-conjugate-frame-heisenberg}(5), with $\epsilon(u)$ of absolute value $1$, and the mean one of
$\vartheta$ are the hypotheses of Lemma~\ref{is-eigenvalue-conjugate-frame-ortho-constant-norm-mean}, which gives
$|a|^2=(2^d\det Y)^{1/2}$. The matrix $h^{\mathsf T}Y^{-1}\bar h$ of $H_x$ has determinant $|\det h|^2/\det Y$, and by
Lemma~\ref{is-eigenvalue-conjugate-frame-ortho-constant-norm-det} it is $\prod_\tau c_\tau^r$. Hence
$|a|^2=\big(2^d/\prod_\tau c_\tau^r\big)^{1/2}|\det h|$.

Write $n=2k$. Then $|a|^n=|K|=|\kappa||C||\det h|^k$, and $|a|^{2k}=\big(2^d/\prod_\tau c_\tau^r\big)^{k/2}|\det h|^k$. As $\det h\ne0$,
$|\kappa|=\big(2^d/\prod_\tau c_\tau^r\big)^{n/4}/|C|=\rho$.
\end{proof}
% VacNorm.norm_factor_eq (ExistsSymplOfSelfDual.lean), followed: exists_sympl_of_selfD (the book cites the symplectic
% basis theorem), periodCoordinateTheta_of_signedHeisenbergVacuum (Lemma -theta-product-vacuum: hh, htau, hH, hth),
% rho = sqrt(2^g / prod_tau c_tau^r)^(n/2) / |C|, hK, the case n = 0, exists_factor, hpv and hA', norm_sq_of_mean (Lemma
% -ortho-constant-norm-mean), T.det_hermitianFormAtFramePeriodCoords_of_unitFrame (Lemma -ortho-constant-norm-det), e3
% and e4. The hypotheses w integral and w + cc w = 1 are passed to Lemma -theta-product-vacuum.

\begin{lemma}\label{is-eigenvalue-conjugate-frame-ortho-constant-unit-factor}
Let $O$ be an open subset of a complex normed space that is connected or empty, $I$ a set, $R_i$ and $T_i$ holomorphic
functions on $O$ for $i\in I$, and $\rho>0$. Suppose that for every $x\in O$ there is $k\in\mathbf C$ with $|k|=\rho$ and
$R_i(x)=kT_i(x)$ for every $i\in I$. Then there is $C\ne0$ with $R_i=CT_i$ on $O$ for every $i\in I$.
\end{lemma}
\leanref{Development/AbelianVarietyTheory/Proofs/SiegelThetaTransformation/ConstOfNormFactorUnitFr.lean}{GlobalTheory.Dict.const_of_norm_factor}
\begin{proof}
If every $T_i$ vanishes on $O$, then so does every $R_i$, and $C=\rho$ works. Otherwise let $T_j(x_0)\ne0$ with $x_0\in O$,
and $B\subset O$ an open ball around $x_0$ on which $T_j$ has no zeros. On $B$ the function $q=R_j/T_j$ is holomorphic, and
$q(y)=k$ for the $k$ of $y$, so $|q|=\rho$ on $B$. As $|q|$ attains its maximum at $x_0$, $q$ is constant on the ball $B$ by
the maximum modulus principle, equal to $C=q(x_0)$; and $C\ne0$, as $|C|=\rho$. For $y\in B$ the $k$ of $y$ is $q(y)=C$, so
$R_i=CT_i$ on $B$ for every $i$.

Let $S$ be the set of the $x\in O$ near which $R_i=CT_i$ for every $i$; it is open and contains $x_0$. It is closed in $O$.
Indeed, let $y\in O$ be in its closure, $\delta>0$ with the ball of radius $\delta$ around $y$ in $O$, $x_1\in S$ with
$|x_1-y|<\delta/4$, and $|z-y|<\delta/4$. For $|t|<3/2$ the point $x_1+t(z-x_1)$ lies in that ball, and the holomorphic
function $t\mapsto R_i(x_1+t(z-x_1))-CT_i(x_1+t(z-x_1))$ on this disc vanishes near $t=0$; by the identity theorem it
vanishes at $t=1$, so $R_i(z)=CT_i(z)$. As $O$ is connected, $S=O$.
\end{proof}
% Dict.const_of_norm_factor (ConstOfNormFactorUnitFr.lean), followed: the case hz, x0 and the ball, hqd, hqn,
% Complex.eqOn_of_isPreconnected_of_isMaxOn_norm, hloc, the set S, hcl (the identity theorem on the disc of radius 3/2,
% AnalyticOnNhd.eqOn_zero_of_preconnected_of_eventuallyEq_zero) and IsPreconnected.subset_of_closure_inter_subset.

\begin{lemma}\label{is-eigenvalue-conjugate-frame-ortho-constant-unit}
Let $X$ be flag data of rank $r$ over a CM field $E$ (Definition~\ref{is-onto-share-flag-data}), with set $\Psi$, let
$\lambda\in\mathbf Z^{\Psi\times\{1,\dots,r\}}$ be a weight, $\rho>0$, and $\kappa$ a function on the set of the $x\in\mathcal D_X$
with $|d_i(x,\tau)|=1$ for all $\tau$ and $i$, with $|\kappa(x)|=\rho$. Then there is $C\in\mathbf C$ with $|C|=\rho$ such that
$P_1(x)=C\,P_2(x)$ for every $x\in\mathcal D_X$, whenever $P_1$ and $P_2$ are functions on $M_r(\mathbf C)^\Psi$, holomorphic
on $\mathcal D_X$, with $P_j(xb)=\prod_{\tau,i}(b_\tau)_{ii}^{\lambda_{\tau,i}}P_j(x)$ for $x\in\mathcal D_X$ and
$b\in\mathcal B^\Psi$, and $P_1(x)=\kappa(x)P_2(x)$ at every such $x$ with $|d_i(x,\tau)|=1$.
\end{lemma}
\leanref{Development/AbelianVarietyTheory/Proofs/SiegelThetaTransformation/RdUEqConstMulClsSum.lean}{GlobalTheory.ClsVal.Const.exists_const}
\leanref{Development/AbelianVarietyTheory/Proofs/SiegelThetaTransformation/ConstOfNormFactorUnitFr.lean}{GlobalTheory.Dict.const_of_norm_factor_unitFr}
\leanref{Development/AbelianVarietyTheory/Proofs/SiegelThetaTransformation/ConstOfNormFactorUnitFr.lean}{GlobalTheory.Dict.exists_unitFr_rmul}
\begin{proof}
Call $x\in\mathcal D_X$ a \emph{unit point} if $|d_i(x,\tau)|=1$ for all $\tau$ and $i$, write
$\chi(b)=\prod_{\tau,i}(b_\tau)_{ii}^{\lambda_{\tau,i}}$, which is not $0$ for $b\in\mathcal B^\Psi$, and let $\mathcal P$ be the
set of the pairs $(P_1,P_2)$ as in the statement.

\emph{Unit points.} For every $x\in\mathcal D_X$ there is $b\in\mathcal B^\Psi$ such that $xb$ is a unit point. Indeed, by
Lemma~\ref{is-onto-share-majorant}(3), $x=x'b'$ with $x'$ orthonormal and $b'\in\mathcal B^\Psi$; for $b=b'^{-1}$, $xb=x'$
lies in $\mathcal D_X$ by the same lemma, and its Gram matrices are $J$, so $d_i(xb,\tau)=\epsilon_i=\pm1$.

\emph{A common factor.} Let $x\in\mathcal D_X$ and $b$ as above. For $(P_1,P_2)\in\mathcal P$,
$\chi(b)P_1(x)=P_1(xb)=\kappa(xb)P_2(xb)=\kappa(xb)\chi(b)P_2(x)$, so $P_1(x)=\kappa(xb)P_2(x)$, with $|\kappa(xb)|=\rho$. As
$\mathcal D_X$ is open and connected (Lemma~\ref{is-onto-share-sections}(2)),
Lemma~\ref{is-eigenvalue-conjugate-frame-ortho-constant-unit-factor}, for the family $\mathcal P$, gives $C'\ne0$ with
$P_1=C'P_2$ on $\mathcal D_X$ for every $(P_1,P_2)\in\mathcal P$.

\emph{The constant.} If $|C'|=\rho$, take $C=C'$. Otherwise take $C=\rho$. Let $(P_1,P_2)\in\mathcal P$, $x\in\mathcal D_X$ and
$b$ as above. Then $C'P_2(xb)=P_1(xb)=\kappa(xb)P_2(xb)$ with $|\kappa(xb)|=\rho\ne|C'|$, so $P_2(xb)=0$; hence
$P_2(x)=\chi(b)^{-1}P_2(xb)=0$, and $P_1(x)=C'P_2(x)=0=\rho P_2(x)$.
\end{proof}
% ClsVal.Const.exists_const (RdUEqConstMulClsSum.lean), followed: const_of_norm_factor_unitFr for the family of all
% pairs (the set of the q with the four conditions), with hopen, frames_isConnected, exists_unitFr_rmul (through
% UnitFrames.frame_mul_triangular_orthonormal; the book uses Lemma is-onto-share-majorant(3)) and bchar_ne_zero_of_mem_bor;
% then the case split on norm C = rho (h5, h6). bchar la b is the product of the (b_tau)_ii^la.

\begin{lemma}\label{is-eigenvalue-conjugate-frame-ortho-constant}
In Situation~\ref{is-eigenvalue-conjugate-frame-uniformization}, suppose that the level has full level
$N=1$; this hypothesis is not used. Let $c$ be a component whose numbers $c_\tau$, $\tau\in\Psi_c$, are all equal, let $m\ge0$ be an even
integer, fix a symplectic basis of $\Lambda_c$, an indexing and a bijection onto the letters, the last two
matched for $s$ (of the matching only the condition on periods is used), and let $v$ be a vacuum of
$\mathcal O_E^m\otimes\Lambda_c$ for $X_c$ and $w_0$
(Definition~\ref{is-eigenvalue-conjugate-frame-heisenberg}(4)). Then there is $C\in\mathbf C$ with
\[
 \prod_{k=1}^m\vartheta(z_k)=C\,\theta_v(x;z)
\]
for every orthonormal $x\in\mathcal D_c$, every theta reading $\vartheta$ of $s_{y(c,x)}$ for an analytic
fibre at $y(c,x)$ for these data, $x$ and $\Lambda_c$, and every $z\in(\mathbf C^d)^m$, with the theta
series formed with these data.
\end{lemma}
\leanref{Development/AbelianVarietyTheory/Proofs/IwasawaDecomposition/ExistsConstOrtho.lean}{GlobalTheory.OneMapRef.exists_const_ortho}
\leanref{Development/AbelianVarietyTheory/Proofs/IwasawaDecomposition/ExistsConstOrtho.lean}{GlobalTheory.OneMapRef.Kappa.exists_const}
\begin{proof}
Fix a uniformization witness, write $X=X_c$, $\Lambda=\Lambda_c$ and $j\mapsto\ell'_j$ for the given bijection onto the
letters. By Lemma~\ref{is-eigenvalue-conjugate-frame-component-lattice}(1) and~(2), $\Lambda$ is self-dual for $X$, and the
marking $m_c$ has values in $\Lambda$ and the two properties stated there. The element $w_0=(1+\sqrt{-d})/2$ is an algebraic
integer with $w_0+\bar w_0=1$. By the compatibility conditions and Lemma~\ref{is-eigenvalue-conjugate-frame-letters}, for the
weights $0$, there is a letter dictionary for $X$, $m$ and $0$ with the bijection $j\mapsto\ell'_j$. Its indexing is the given
one, as an indexing is determined by the periods at the letters; so the given indexing and bijection are matched for $s$,
with both conditions.

\emph{The constants at the unit frames.} Let $\rho$ be as in Lemma~\ref{is-eigenvalue-conjugate-frame-ortho-constant-norm}
for $X$, $m$, $w_0$, $\Lambda$ and $v$. Let $x\in\mathcal D_c$ with $|d_i(x,\tau)|=1$ for all $\tau$ and $i$. Let $\pi$ be the
permutation with $\ell'_j=\ell_{\pi(j)}$ for the bijection $j\mapsto\ell_j$ of the witness. With the indexing composed with
$\pi^{-1}$, which is matched with $j\mapsto\ell_j$, Lemma~\ref{is-eigenvalue-conjugate-frame-analytic-fibre} gives an analytic
fibre at $y(c,x)$, marked by $m_c$, with a theta reading of mean one. Renumbering the coordinates of $\mathbf C^d$ by $\pi$
turns it into an analytic fibre $A_x$ for the given indexing and bijection, still marked by $m_c$, with a theta reading
$\vartheta_x$ that still has mean one, as the renumbering preserves the periods, the forms and Lebesgue measure. By
Lemma~\ref{is-eigenvalue-conjugate-frame-level-one}, $\vartheta_x$ is a level-one theta function for $w_0$, $x$ and $\Lambda$,
and by Lemma~\ref{is-eigenvalue-conjugate-frame-theta-product} there is $\kappa(x)$ with
$\prod_k\vartheta_x(z_k)=\kappa(x)\theta_v(x;z)$. By Lemma~\ref{is-eigenvalue-conjugate-frame-ortho-constant-norm},
$|\kappa(x)|=\rho$.

\emph{The constant.} Let $C$ be as in Lemma~\ref{is-eigenvalue-conjugate-frame-ortho-constant-unit} for $X$, the weight
$\lambda$ of the sections of weight $0$ and exponent $m$ read on $c$, $\rho$ and $\kappa$.

\emph{The given point.} Let $x\in\mathcal D_c$ be orthonormal; then $|d_i(x,\tau)|=1$ for all $\tau$ and $i$. Let $A$ be an
analytic fibre at $y(c,x)$ for the given data, $x$ and $\Lambda$, and $\vartheta$ a theta reading of $s_{y(c,x)}$ for it. The
exponentials of $A$ and $A_x$ are holomorphic group homomorphisms with the same derivatives at $0$, so they are equal (as in
the proof of Lemma~\ref{is-eigenvalue-conjugate-frame-dictionary-point}); hence $A$ is marked by $m_c$. By
Lemma~\ref{is-eigenvalue-conjugate-frame-level-one} and Lemma~\ref{is-eigenvalue-conjugate-frame-theta-product} there is
$\kappa'$ with $\prod_k\vartheta(z_k)=\kappa'\theta_v(x;z)$, and it remains to show $\kappa'=C$.
\openstep{The remaining step, which builds from the letter dictionary above, the roots of the dual marking (Lemma~\ref{cv-uniformization-theta-marking}) and a level translate with a
nonvanishing jet, a section $f$ of weight $0$ and exponent $m$ and a function $P$ on $M_r(\mathbf C)^{\Psi_c}$, holomorphic on
$\mathcal D_c$ with the weight law of $\lambda$ and $P(x)\ne0$, such that the reading of $f$ at every $x'\in\mathcal D_c$ through
an analytic fibre marked by $m_c$ with $\prod_k\vartheta'(z_k)=\kappa''\theta_v(x';z)$ is $\kappa''P(x')$
(Lemma~\ref{cv-uniformization-theta-rule-fibre}), and then applies
Lemma~\ref{is-eigenvalue-conjugate-frame-ortho-constant-unit} to the reading of $f$ and $P$ to get $\kappa'P(x)=CP(x)$, is not
written out.}
% The remaining step: exists_dic (Lemma -frame-letters) with idx_unique (PowGeneratorEqConstMulVacuumTheta.lean), DualMark.dual_marking
% (hY), exists_level_translate_jet_ne_zero (ExistsLevelTranslateJetNeZero.lean, 91 lines, through
% T.exists_theta_sublattice_jet_ne_zero_alt and T.theta_ideal_lattice_eq_sum_level_translates), Kappa.exists_root,
% existsUnique_thRule and thetaRule_mem_sec (Lemma cv-uniformization-theta-rule), heisG_laws, theta_holomorphic_weight,
% the section clause hU.2.2.2.2.1 of IsUnf, Dc.wt_eq, and hval through thVal_eq_thP (Lemma cv-uniformization-theta-rule-fibre).
% The renumbering above is exists_anFib_trans and meanOne_trans; the equality of exponentials is Kappa.ex_eq
% (CuspFamilyCompare.exp_unique). Sizes (tools/lean_size.py): exists_level_translate_jet_ne_zero 3160 lines (the jet:
% y with M y_b in conj(Mfr) Lambda and kappa thP(heisG y' v, 1)(x) ne 0 for all y' = y mod Lambda, at an orthonormal x
% with all c_tau equal), SW.theta_holomorphic_weight 3086 lines (holomorphy and weight law of x -> thP(phi, P)(x)). The
% section clause of IsUnf (the reading of a section of weight 0 and exponent m is holomorphic on D_c with the weight law
% laU) is not stated in Situation -uniformization, so writing this step also needs that clause added there. The
% constant C is the one of Lemma -ortho-constant-unit for all pairs at once; the final equation is h4 with P(x) ne 0
% (h3) from the jet.
% exists_const_ortho = Kappa.exists_const (111 lines; ExistsConstOrtho.lean, 160 lines) for Unf.ofPj U, under
% GoodEmb c sigma, IsUnf U, 1 <= r, l.A = 1, the constant t, IsSympl and Matched (only hm.1 is passed on). Its
% proof takes W from nonempty_uw and hL from uw_lattice (the step above), then exists_dic, VacNorm.norm_factor_eq,
% uw_anFib, exists_anFib_trans, isTheta1_of_reads_generator, VacProd.prod_eq_mul_thA, Const.exists_const,
% DualMark.dual_marking and exists_level_translate_jet_ne_zero.
\end{proof}

\begin{lemma}\label{is-eigenvalue-conjugate-frame-hermitian-riemann}
Let $X$ be flag data of rank $r$ over a CM field $E$ (Definition~\ref{is-onto-share-flag-data}), with set $\Psi$,
integer $p_0$, element $\xi_0$ and forms $h_\tau$, with a symplectic basis and an indexing
(Definition~\ref{is-eigenvalue-conjugate-frame-fibres}(1) and~(2)). Put $d=|\Psi|r$ and $c_\tau=2|\tau(\xi_0)|$, and
let $G$ be the diagonal $d\times d$ matrix with entries $c_{\tau_j}$.
\begin{enumerate}
\item For every $x\in\mathcal D_X$ there is $b\in\mathcal B^\Psi$ such that $xb$ is orthonormal.
\item For every $b\in\mathcal B^\Psi$ there is an invertible matrix $N_b\in M_d(\mathbf C)$ such that
  $z_{xb}(u)=N_bz_x(u)$ for every family $x$ of invertible matrices and every family $u$ of vectors; in particular
  $p_{xb}(v)=N_bp_x(v)$ for $v\in E^r$.
\item If $x$ is orthonormal, then $\operatorname{Im}\big(p_x(u)^{\mathsf T}G\,\overline{p_x(v)}\big)=\langle
  u,v\rangle_X$ for all $u,v\in E^r$.
\item Let $x\in\mathcal D_X$ and $b$ be as in (1). Then $\Omega_x$ is invertible, $\tau_x$ is symmetric,
  $\operatorname{Im}\tau_x$ is positive definite, and $H_x(z,w)=(N_bz)^{\mathsf T}G\,\overline{N_bw}$ for all
  $z,w\in\mathbf C^d$.
\end{enumerate}
\end{lemma}
\leanref{ReciprocityTheory/ReductiveGroups/ArithmeticGeometry/FundamentalDomains/ProdCastLESucc.lean}{ReciprocityTheory.ChartRestriction.G.FD.FrameDegreeData.bOf}
\leanref{ReciprocityTheory/ReductiveGroups/ArithmeticGeometry/FundamentalDomains/ProdCastLESucc.lean}{ReciprocityTheory.ChartRestriction.G.FD.FrameDegreeData.ortho_bOf}
\leanref{Development/AbelianVarietyTheory/Proofs/CuspFormEvaluations/DifferentiableOnDetHMat.lean}{GlobalTheory.ThFr.per_rmul}
\leanref{Development/AbelianVarietyTheory/Proofs/CuspFormEvaluations/DifferentiableOnDetHMat.lean}{GlobalTheory.ThFr.im_hermForm_diagonal_c}
\leanref{Development/AbelianVarietyTheory/Proofs/UnitaryExistence/DifferentiableOnLeftInverse.lean}{GlobalTheory.RiemannRel.det_ne_zero}
\leanref{Development/AbelianVarietyTheory/Proofs/UnitaryExistence/DifferentiableOnLeftInverse.lean}{GlobalTheory.RiemannRel.symm_posDef}
\leanref{Development/AbelianVarietyTheory/Proofs/CuspFormEvaluations/DifferentiableOnDetHMat.lean}{GlobalTheory.ThFr.riemann_of_isOrtho}
\leanref{Development/AbelianVarietyTheory/Proofs/CuspFormEvaluations/DifferentiableOnDetHMat.lean}{GlobalTheory.ThFr.hermForm_Hm_eq_diagonal}
\begin{proof}
(1) Fix $\tau$, and let $M$ be the Gram matrix of the columns of $x_\tau$ for $h_\tau$, a Hermitian matrix
whose leading minors $d_i=d_i(x,\tau)$ are real and nonzero, with the sign $\epsilon_i$. By Gaussian
elimination there is an upper triangular matrix $u$ with diagonal entries $1$ such that $u^*Mu$ is diagonal; its
entries are $d_i/d_{i-1}$, of the sign $\epsilon_i\epsilon_{i-1}=\operatorname{sgn}(i)$. With $s$ the diagonal
matrix with entries $|d_i/d_{i-1}|^{-1/2}$, $b_\tau=us$ is upper triangular and invertible, and the Gram matrix of the
columns of $x_\tau b_\tau$ is $s\,u^*Mu\,s=J$.

(2) Fix $\tau$ and write $x=x_\tau$, $b=b_\tau$ and $y=u_\tau$. As $h_\tau$ is linear in its first variable and $b$
is upper triangular, for $a\le p_0$
\[
 h_\tau(xbe_a,y)=\sum_{c\le a}b_{ca}\,h_\tau(xe_c,y),
\]
where every $c\le a$ is $\le p_0$; and for $a>p_0$, $((xb)^{-1}y)_a=\sum_{c\ge a}(b^{-1})_{ac}(x^{-1}y)_c$, where every
$c\ge a$ is $>p_0$. So $z_{xb}(u)=N_bz_x(u)$ for the matrix $N_b$ whose entry at $(j,j')$ with $\tau_j=\tau_{j'}=\tau$
is $b_{a_{j'}a_j}$ if $a_j,a_{j'}\le p_0$ and $(b^{-1})_{a_ja_{j'}}$ if $a_j,a_{j'}>p_0$, and which is $0$ elsewhere. In
a suitable order of the indices it is triangular, with the nonzero diagonal entries $b_{aa}$ and $b_{aa}^{-1}$; so it is
invertible.

(3) Let $\tau\in\Psi$, and write $\tau(u)=x_\tau\alpha$ and $\tau(v)=x_\tau\beta$. As the columns of $x_\tau$ have Gram
matrix $J$, $p_x(u)_{\tau,a}=h_\tau(x_\tau e_a,x_\tau\alpha)=\bar\alpha_a$ for $a\le p_0$ and $p_x(u)_{\tau,a}=\alpha_a$
for $a>p_0$, and $h_\tau(\tau(u),\tau(v))=\sum_{a\le p_0}\alpha_a\bar\beta_a-\sum_{a>p_0}\alpha_a\bar\beta_a$. Hence
\[
 \operatorname{Im}\sum_ap_x(u)_{\tau,a}\overline{p_x(v)_{\tau,a}}=\operatorname{Im}\Big(\sum_{a\le p_0}\bar\alpha_a\beta_a
 +\sum_{a>p_0}\alpha_a\bar\beta_a\Big)=-\operatorname{Im}h_\tau(\tau(u),\tau(v)).
\]
As $\Psi$ contains one of each pair of conjugate embeddings, $\operatorname{Tr}_{E/\mathbf
Q}(t)=\sum_{\tau\in\Psi}2\operatorname{Re}\tau(t)$; and $\tau(-\frac12\kappa_0)=\frac i2c_\tau$, as
$\kappa_0=\pm2\xi_0$, $\tau(\xi_0)$ is purely imaginary and $\operatorname{Im}\tau(\kappa_0)<0$. So
$\langle u,v\rangle_X=\sum_\tau c_\tau\operatorname{Re}\big(i\,h_\tau(\tau(u),\tau(v))\big)=-\sum_\tau
c_\tau\operatorname{Im}h_\tau(\tau(u),\tau(v))$, which is the left side of (3) by the display.

(4) First let $x$ be orthonormal; then $x\in\mathcal D_X$. Write $A=\Omega_x$ and $B=\Omega'_x$, and
$G(z,w)=z^{\mathsf T}G\bar w$, a Hermitian form with $\operatorname{Re}G(z,z)>0$ for $z\ne0$. By (3) and the
symplectic basis, the matrices $\operatorname{Im}(A^{\mathsf T}G\bar A)$, $\operatorname{Im}(A^{\mathsf T}G\bar B)$
and $\operatorname{Im}(B^{\mathsf T}G\bar B)$ are $0$, $-1$ and $0$. These are the Riemann relations, and they give the
following \cite[Section~8.1]{BL04}.
\begin{itemize}
\item $A$ is invertible. Let $Av=0$ and write $v=s+it$ with $s,t$ real. Then $As=-iAt$, so
  $G(As,As)=i\,G(As,At)$, whose real part is $-\operatorname{Im}G(As,At)=0$, as $\operatorname{Im}(A^{\mathsf
  T}G\bar A)=0$ and $s,t$ are real. So $As=0$, and then $At=0$. As $\operatorname{Im}G(As,Be_k)=-s_k$ and
  $\operatorname{Im}G(At,Be_k)=-t_k$, $s=t=0$.
\item Put $S=A^{\mathsf T}G\bar A$, a Hermitian matrix with imaginary part $0$, so real symmetric, and positive
  definite: $s^{\mathsf T}Ss=G(As,As)>0$ for real $s\ne0$. With $\tau_x=X+iY$, $A^{\mathsf T}G\bar B=S\bar\tau_x$ has
  imaginary part $-SY=-1$, so $Y=S^{-1}$ is symmetric and positive definite. And $B^{\mathsf T}G\bar B=\tau_x^{\mathsf
  T}S\bar\tau_x$ has imaginary part $Y^{\mathsf T}SX-X^{\mathsf T}SY=X-X^{\mathsf T}=0$, so $\tau_x$ is symmetric.
\item $H_x(z,w)=(A^{-1}z)^{\mathsf T}S\,\overline{A^{-1}w}=z^{\mathsf T}G\bar w$.
\end{itemize}
Now let $x\in\mathcal D_X$ and $b$ be as in (1). By (2), $\Omega_{xb}=N_b\Omega_x$ and $\Omega'_{xb}=N_b\Omega'_x$, so
$\Omega_x$ is invertible, $\tau_x=\tau_{xb}$, and $H_x(z,w)=H_{xb}(N_bz,N_bw)$; the case just treated, for the
orthonormal point $xb$, gives (4).
\end{proof}
% (1) is FrameDegreeData.bOf = uOf * diagonal sOf (uMat of the Gram matrix, dOf its diagonal, sOf = |dOf|^(-1/2)),
% bOf_mem (upper triangular, invertible) and ortho_bOf. (2) is per_rmul with borMat (borMat_of_lt, borMat_of_le) and
% det_borMat_ne_zero (det_of_triangular); perMat_rmul, hMat_rmul, tauMat_rmul and hermForm_Hm_rmul transport the
% periods, tau and the form. (3) is im_hermForm_diagonal_c: EF_eq_sum (emb_neg_half, trace_eq_sum_re) and
% UnitFrame.im_sum_periods. (4) is riemann_of_isOrtho (hE from (3) and IsSympl; RiemannRel.det_ne_zero and
% RiemannRel.symm_posDef, which the proof above follows; Hm = G through imInv_eq and gram_mul_im), isUnit_det_perMat,
% tauMat_isSymm, tauMat_im_posDef and hermForm_Hm_eq_diagonal.

\begin{lemma}\label{is-eigenvalue-conjugate-frame-hermitian}
Let $X$ be flag data of rank $r$ over a CM field $E$ with a symplectic basis and an indexing
(Definition~\ref{is-eigenvalue-conjugate-frame-fibres}(1) and~(2)), let $n$ be an even integer, which is not
used, and let $x\in\mathcal D_X$.
\begin{enumerate}
\item $H_x(w,z)=\overline{H_x(z,w)}$ for all $z,w\in\mathbf C^d$, and $\operatorname{Re}H_x(z,z)>0$ for $z\ne0$.
\item $\operatorname{Im}H_x(p_x(u),p_x(v))=\langle u,v\rangle_X$ for all $u,v\in E^r$.
\item For $a\in E$, $H_x\big((\rho_j(a)z_j)_j,w\big)=H_x\big(z,(\rho_j(\bar a)w_j)_j\big)$ for all $z,w$, and
  $p_x(av)=(\rho_j(a)p_x(v)_j)_j$ for all $v\in E^r$.
\item $p_x$ is injective. For every $\mathcal O_E$-submodule $L\subset E^r$ that spans $E^r$ over $E$, the
  vectors $p_x(v)$, $v\in L$, span $\mathbf C^d$ over $\mathbf R$, and every Hermitian matrix $H'$ with
  $\operatorname{Im}\big(p_x(u)^{\mathsf T}H'\overline{p_x(v)}\big)=\langle u,v\rangle_X$ for all $u,v\in L$
  satisfies $z^{\mathsf T}H'\bar w=H_x(z,w)$ for all $z,w$.
\item Let $\gamma\in GL_r(E)$ with $h(\gamma u,\gamma v)=h(u,v)$ for the form $h$ of $X$. Then
  $\gamma x\in\mathcal D_X$, $p_{\gamma x}(\gamma v)=p_x(v)$ for every $v\in E^r$, $|d_i(\gamma x,\tau)|=1$ for all
  $\tau$ and $i$ if $|d_i(x,\tau)|=1$ for all $\tau$ and $i$, and $H_{\gamma x}=H_x$.
\end{enumerate}
\end{lemma}
\leanref{Development/AbelianVarietyTheory/Proofs/SiegelThetaTransformation/HermitianFormAtFrameLaws.lean}{GlobalTheory.Dict.hx_laws}
\begin{proof}
Let $b$, $N_b$ and $G$ be as in Lemma~\ref{is-eigenvalue-conjugate-frame-hermitian-riemann} for $x$, so that
$H_x(z,w)=(N_bz)^{\mathsf T}G\,\overline{N_bw}$. We use two facts.
\begin{enumerate}
\item[(a)] The vectors $p_x(v)$, $v\in E^r$, span $\mathbf C^d$ over $\mathbf R$. Indeed, let $z\in\mathbf C^d$ and
  $y=\Omega_x^{-1}z$. As $\operatorname{Im}\tau_x$ is invertible, there is a real vector $t$ with
  $\operatorname{Im}(\tau_xt)=\operatorname{Im}y$; then $s=y-\tau_xt$ is real, and $z=\Omega_xs+\Omega'_xt$ is a real
  combination of the periods of the symplectic basis.
\item[(b)] Let $F$ be a function on $\mathbf C^d\times\mathbf C^d$, linear in the first variable and conjugate linear
  in the second, with $\operatorname{Im}F(z,w)=0$ for all $z,w$ in a set that spans $\mathbf C^d$ over $\mathbf R$. Then
  $F=0$: $\operatorname{Im}F$ is real bilinear, so it vanishes everywhere, and $\operatorname{Re}F(z,w)=
  \operatorname{Im}F(iz,w)=0$.
\end{enumerate}

(1) and (2) By Lemma~\ref{is-eigenvalue-conjugate-frame-hermitian-riemann}(4), $H_x$ is the Hermitian form of the
positive definite diagonal matrix $G$ composed with the invertible matrix $N_b$, which gives (1). For (2),
$N_bp_x(v)=p_{xb}(v)$ by (2) of that lemma, so $\operatorname{Im}H_x(p_x(u),p_x(v))=\operatorname{Im}\big(p_{xb}(u)^{\mathsf
T}G\,\overline{p_{xb}(v)}\big)=\langle u,v\rangle_X$ by its (3), as $xb$ is orthonormal.

(3) The second assertion is Definition~\ref{is-eigenvalue-conjugate-frame-fibres}(1). For the first, both sides are
linear in $z$ and conjugate linear in $w$. At $z=p_x(u)$ and $w=p_x(v)$ their imaginary parts are $\langle
au,v\rangle_X$ and $\langle u,\bar av\rangle_X$, by the second assertion and (2); these are equal, as
$h(au,v)=ah(u,v)=h(u,\bar av)$. By (a) and (b) the two sides are equal.

(4) Let $p_x(u)=p_x(v)$. By (2), $\langle u,t\rangle_X=\langle v,t\rangle_X$ for every $t\in E^r$. The symplectic basis
is a basis of $E^r$ over $\mathbf Q$ on which $\langle\ ,\ \rangle_X$ has the matrix of the standard symplectic form, so
the pairings with its vectors determine the coordinates, and $u=v$. Let $L$ span $E^r$ over $E$. For $v\in E^r$ there
is an integer $k\ne0$ with $kv\in L$, and $p_x(v)=k^{-1}p_x(kv)$; so by (a) the $p_x(v)$, $v\in L$, span $\mathbf C^d$
over $\mathbf R$. If $H'$ is as in (4), then $(z,w)\mapsto z^{\mathsf T}H'\bar w-H_x(z,w)$ has imaginary part $0$ on
this spanning set by (2), and it vanishes by (b). That $H'$ is Hermitian is not used.

(5) For every $\tau$, $\gamma^\tau$ preserves $h_\tau$, so the Gram matrices of the columns of $\gamma^\tau x_\tau$ and of
$x_\tau$ are equal. Hence $\gamma x\in\mathcal D_X$, and $d_i(\gamma x,\tau)=d_i(x,\tau)$. For $a\le p_0$,
$h_\tau(\gamma^\tau x_\tau e_a,\gamma^\tau\tau(v))=h_\tau(x_\tau e_a,\tau(v))$, and for $a>p_0$,
$((\gamma^\tau x_\tau)^{-1}\gamma^\tau\tau(v))_a=(x_\tau^{-1}\tau(v))_a$; so $p_{\gamma x}(\gamma v)=p_x(v)$. Then by (2),
for $\gamma x$ and for $x$,
$\operatorname{Im}H_{\gamma x}(p_x(u),p_x(v))=\operatorname{Im}H_{\gamma x}(p_{\gamma x}(\gamma u),p_{\gamma x}(\gamma v))
=\langle\gamma u,\gamma v\rangle_X=\langle u,v\rangle_X=\operatorname{Im}H_x(p_x(u),p_x(v))$, and $H_{\gamma x}=H_x$ by (a)
and (b).
\end{proof}
% Dict.hx_laws (HermitianFormAtFrameLaws.lean, 116 lines before it), for X, bas, idx with IsSympl, n with Even n
% (unused) and x in X.dom: Hm_isHermitian and hermForm_Hm_pos for (1) (through riemann_of_isOrtho and
% hermForm_Hm_eq_diagonal, Lemma frame-hermitian-riemann), im_hermForm_Hm for (2), hermForm_eq_of_im with
% span_range_pv (fact (a)) and EF_smul_left for (3), form_eq_zero_of_im (fact (b)); hinj through vecQ_surjective
% and EF_vecQ, span_image_pv (LatFin.exists_int_smul_mem) and the uniqueness (the hypothesis Hm.IsHermitian is
% unused) for (4); IsUn.lact_mem, per_lact, the Gram matrices (PerMove.hermC_unitary) and hermForm_eq_of_im for (5).

\begin{lemma}\label{is-eigenvalue-conjugate-frame-level-one}
Let a PEL datum over a CM field $E$ with prime $d$ be given, and let $X$, a symplectic basis and an
indexing be as in Definition~\ref{is-eigenvalue-conjugate-frame-heisenberg}, over $E$. Let $y$, a bijection
$j\mapsto\ell_j$ onto the letters, $x\in\mathcal D_X$, an $\mathcal O_E$-submodule $L\subset E^r$ self-dual
for $X$, and an analytic fibre at $y$ for these data, $x$ and $L$ be as in
Definition~\ref{is-eigenvalue-conjugate-frame-fibres}(3). Then every theta reading of $s_y$ for this
analytic fibre is a level-one theta function for $w_0=(1+\sqrt{-d})/2$, $x$ and $L$
(Definition~\ref{is-eigenvalue-conjugate-frame-heisenberg}(5)).
\end{lemma}
\leanref{Development/AbelianVarietyTheory/Proofs/SiegelThetaTransformation/FramedPointGeneratorIsThetaLevelOne.lean}{GlobalTheory.Dict.isTheta1_of_reads_generator}
\leanref{Development/AbelianVarietyTheory/Proofs/SiegelThetaTransformation/ReadsExistsUniqueTheta.lean}{GlobalTheory.Dict.reads_existsUnique_theta}
\begin{proof}
Let $\vartheta$ be a theta reading of $s_y$, $\chi$ the function of (c) of the analytic fibre, $v\in L$ and
$p=p_x(v)$. Let $z\in\mathbf C^d$, and choose $\beta$ with $\operatorname{ex}(z)\in U_\beta$. As $p$ lies in
the kernel of $\operatorname{ex}$, $\operatorname{ex}(z+p)=\operatorname{ex}(z)$, and the section $s_y$ has one
value there: $\vartheta(z+p)\operatorname{fr}_\beta(z+p)=\vartheta(z)\operatorname{fr}_\beta(z)$. By (c),
$\operatorname{fr}_\beta(z)=\chi(v)\exp\big(\pi H_x(z,p)+\frac\pi2H_x(p,p)\big)\operatorname{fr}_\beta(z+p)$, and
$\operatorname{fr}_\beta(z+p)\ne0$. As $H_x$ is linear in its first variable, this gives
\[
 \chi(v)^{-1}e^{-\pi H_x(z+\frac12p,\,p)}\,\vartheta(z+p)=\vartheta(z).
\]
It remains to see that $\vartheta$ is holomorphic and that $\chi(v)^{-1}=e^{\pi i\langle v,w_0v\rangle_X}$.

\emph{Holomorphy.} Let $z\in\mathbf C^d$ and $\beta$ with $\operatorname{ex}(z)\in U_\beta$, and let $O_\beta$ be the
set of the $w$ with $\operatorname{ex}(w)\in U_\beta$, an open set as $\operatorname{ex}$ is holomorphic. The value
$s_\beta$ of $s_y$ in the trivialization over $U_\beta$ is a regular function, so $w\mapsto s_\beta(\operatorname{ex}(w))$
is holomorphic on $O_\beta$; and on $O_\beta$, $\vartheta(w)=s_\beta(\operatorname{ex}(w))\operatorname{fr}_\beta(w)^{-1}$
with $\operatorname{fr}_\beta$ holomorphic without zeros. So $\vartheta$ is holomorphic near $z$.

\emph{The sign.} Lemma~\ref{cv-uniformization-theta-half-sign-fibre}, applied to the PEL datum, $X$ with its
symplectic basis and indexing, $y$, the bijection $j\mapsto\ell_j$, $x$, $L$ and the analytic fibre, gives
$\chi(v)=e^{\pi i\langle v,w_0v\rangle_X}$. As $w_0\in\mathcal O_E$, $w_0v\in L$, so $\langle v,w_0v\rangle_X$ is an
integer by self-duality, and $e^{\pi i\langle v,w_0v\rangle_X}=\pm1$ is its own inverse. Hence $\chi(v)^{-1}=e^{\pi
i\langle v,w_0v\rangle_X}$.
\end{proof}
% isTheta1_of_reads_generator (43 lines in its file): reads_existsUnique_theta with k = 1 (Differentiable and
% (chi v)^-1 * fock x 1 (pv v) theta = theta; ReadsExistsUniqueTheta.lean), then
% analyticFibre_principalBundle_sign_eq (chi v in terms of the Riemann form) and
% thetaSignExpEF_inv_of_selfDual_lattice (that sign is eps1 w0 v for self-dual L). The first step above is the
% periodicity part of reads_existsUnique_theta; the holomorphy is its Differentiable clause (A.hol through
% isHolFam_fibrePoint and vtf_chart, times fr_hol to the power -1). analyticFibre_principalBundle_sign_eq, with the
% arguments IsSympl X bas, x in X.dom and SelfD X L passed on unchanged, is Lemma
% cv-uniformization-theta-half-sign-fibre. thetaSignExpEF_inv_of_selfDual_lattice: hwv (w0 v in L) and the
% integrality of EF on L.

\begin{lemma}\label{is-eigenvalue-conjugate-frame-renumbering}
Let a PEL datum over a CM field $E$, a level, its moduli scheme and an embedding $s$ of a coefficient
field $K_s$ of it be given. Let $X'$ and $X$ be flag data of rank $r$ over $E$ with the same set $\Psi$,
matrix, integer $p_0$, element $\xi_0$ and element $\kappa_0$, each with a symplectic basis of its module,
an indexing and a bijection onto the letters, the last two matched for $s$. Let $x\in\mathcal D_{X'}$,
$L\subset E^r$ an $\mathcal O_E$-submodule self-dual for $X'$ (the self-duality is not used), $\gamma\in
GL_r(E)$ with
$h(\gamma u,\gamma v)=h(u,v)$ for the form $h$ of $X$ and all $u,v\in E^r$, $y$ a complex framed point of
the moduli scheme over $s$, and $\vartheta'$ a theta reading of $s_y$ for an analytic fibre $A'$ at $y$ for
the data of $X'$, $x$ and $L$. Then there are an analytic fibre $A$ at $y$ for the data of $X$, the point
$\gamma x$ and the module $\gamma L$, and a theta reading $\vartheta$ of $s_y$ for $A$, such that $A$ is
marked by $a\mapsto\gamma m(a)$ whenever $A'$ is marked by $m$ (as in
Definition~\ref{is-eigenvalue-conjugate-frame-marking}(2), with the level structure of the moduli scheme),
and
\[
 \mathcal G_{\gamma x}(z)\prod_{k=1}^n\vartheta(z_k)=\mathcal G'_x(z')\prod_{k=1}^n\vartheta'(z'_k)
\]
for every $n\ge0$ and all families $u'=(u'_1,\dots,u'_n)$ and $u=(u_1,\dots,u_n)$ as in
Definition~\ref{is-eigenvalue-conjugate-frame-fibres}(1) with $(u_k)_\tau=\gamma^\tau(u'_k)_\tau$ for all $k$
and $\tau$, where $z=(z_{\gamma x}(u_k))_k$ and $\mathcal G$ are the coordinates and the Gaussian of $X$,
and $z'=(z'_x(u'_k))_k$ and $\mathcal G'$ those of $X'$.
\end{lemma}
\leanref{Development/AbelianVarietyTheory/Proofs/IwasawaDecomposition/SymplecticBasesMatched.lean}{GlobalTheory.OneMapRef.anFib_move_renum}
\begin{proof}
Write $p$, $H$, $z$ and $\ell_j$ for the periods, the forms, the coordinates and the bijection onto the letters
of $X$, and $p'$, $H'$, $z'$ and $\ell'_j$ for those of $X'$. Let $q$ be the bijection of the indices of $X'$ onto
those of $X$ with $\ell_{q(j)}=\ell'_j$, and for $w\in\mathbf C^d$, indexed as for $X$, let $w^q$ be the vector
with $(w^q)_j=w_{q(j)}$. Since $X$ and $X'$ have the same set, matrix and integer $p_0$, the period of a vector at
a letter is the same for both, and the matchings give $p_x(v)^q=p'_x(v)$ for every $v\in E^r$ and
$\rho_{q(j)}=\rho'_j$ for the row types.

\emph{The forms.} $H_x(w_1,w_2)=H'_x(w_1^q,w_2^q)$ for all $w_1,w_2$. Indeed, the difference of the two sides is
linear in $w_1$ and conjugate linear in $w_2$, and its imaginary part vanishes at $w_1=p_x(u)$, $w_2=p_x(v)$, as
both are $\langle u,v\rangle_X=\langle u,v\rangle_{X'}$ there (Lemma~\ref{is-eigenvalue-conjugate-frame-hermitian}(2),
as the two flag data have the same matrix and $\kappa_0$). As the $p_x(v)$ span $\mathbf C^d$ over $\mathbf R$
(Lemma~\ref{is-eigenvalue-conjugate-frame-hermitian}(4)) and a form of this kind with imaginary part $0$ on a
real spanning set is $0$ (its real part at $(w_1,w_2)$ is the imaginary part at $(iw_1,w_2)$), the difference
vanishes. By Lemma~\ref{is-eigenvalue-conjugate-frame-hermitian}(5), $\gamma x\in\mathcal D_X$,
$p_{\gamma x}(\gamma v)=p_x(v)$ and $H_{\gamma x}=H_x$. So $p_{\gamma x}(\gamma v)^q=p'_x(v)$ and
$H_{\gamma x}(w_1,w_2)=H'_x(w_1^q,w_2^q)$.

\emph{The fibre.} Let $\operatorname{ex}'$, $\operatorname{fr}'_\beta$ and $\chi'$ be the data of $A'$. Put
$\operatorname{ex}(w)=\operatorname{ex}'(w^q)$, $\operatorname{fr}_\beta(w)=\operatorname{fr}'_\beta(w^q)$ and
$\chi(v)=\chi'(\gamma^{-1}v)$, and $\vartheta(w)=\vartheta'(w^q)$. These satisfy (a) to (c) of
Definition~\ref{is-eigenvalue-conjugate-frame-fibres}(3) for the data of $X$, $\gamma x$ and $\gamma L$, and
$\vartheta$ is a theta reading of $s_y$ for them: $w\mapsto w^q$ is a linear bijection; $w$ lies in the kernel
of $\operatorname{ex}$ exactly when $w^q=p'_x(v)=p_{\gamma x}(\gamma v)^q$ for some $v\in L$, that is, $w\in
p_{\gamma x}(\gamma L)$; $\big((\rho_j(\alpha)w_j)_j\big)^q=(\rho'_j(\alpha)w^q_j)_j$; the unit vector at
$q(j)$ goes to the unit vector at $j$, and $\ell_{q(j)}=\ell'_j$; the trivializing functions are those of $A'$
read at $w^q$; and for $v=\gamma v'$ with $v'\in L$, $p_{\gamma x}(v)^q=p'_x(v')$ and the form $H_{\gamma x}$
at $w$ is $H'_x$ at $w^q$, so the law (c) for $A'$ at $v'$ is the law (c) for $A$ at $v$, with $\chi(v)=\chi'(v')$
of absolute value $1$. If $A'$ is marked by $m$, then for $a\in\bar{\mathfrak M}\mathfrak L$ the level structure
sends $a$ to $\operatorname{ex}'(p'_x(M^{-1}m(a)))=\operatorname{ex}'(p_{\gamma x}(M^{-1}\gamma m(a))^q)=
\operatorname{ex}(p_{\gamma x}(M^{-1}\gamma m(a)))$; so $A$ is marked by $a\mapsto\gamma m(a)$.

\emph{The identity.} Let $u'$ and $u$ be as in the statement. By the definition of the coordinates,
$z_{\gamma x}(u_k)=z_x(u'_k)$, as $\gamma^\tau$ is unitary for $h_\tau$; and $z_x(u'_k)^q=z'_x(u'_k)$: both sides
are real linear and continuous in $u'_k$, they agree on the families $(\tau(v))_\tau$, $v\in E^r$, where they
are periods, and these families are dense. So $z_k^q=z'_k$, $\mathcal G_{\gamma x}(z)=\exp(-\frac\pi2\sum_k
H'_x(z'_k,z'_k))=\mathcal G'_x(z')$ and $\prod_k\vartheta(z_k)=\prod_k\vartheta'(z'_k)$.
\end{proof}
% anFib_move_renum (73 lines; SymplecticBasesMatched.lean, 164 lines), for SameSp X' X, IsSympl and Matched for
% both data, x in X'.dom, SelfD X' L (hL, not used), IsUn X ga, a framed point y and A'.Reads 1 y.s th': q =
% e'.trans e.symm, hpv (perFr_castFr), hty, hx_laws for lact ga (Lemma frame-hermitian(5)), hx_renum (form_eq_zero_of_im
% with span_range_pv and the clause him of hx_laws), hlat for the kernel, the fibre field by field, the marking
% through hp, and the identity through zR_lact and zR_renum (denseRange_embR).

\begin{lemma}\label{is-eigenvalue-conjugate-frame-series}
In Situation~\ref{is-eigenvalue-conjugate-frame-uniformization}, suppose that the level has full level
$N=1$. Let $c$ be a component whose numbers $c_\tau$, $\tau\in\Psi_c$, are all equal and whose module
$\Lambda_c$ is self-dual for $X_c$, let $m\ge0$ be an even integer, and fix a symplectic basis of
$\Lambda_c$, an indexing and a bijection onto the letters, the last two matched for $s$. Then there is a
Schwartz-Bruhat function
$\varphi$ on $M_{m\times d}(\mathbf Q)$, $d=|\Psi_c|r$, with the following property. Let $x\in\mathcal D_c$
be orthonormal, let a second symplectic basis of $\Lambda_c$, indexing and bijection onto the letters be
given, the last two matched for $s$, and let $\vartheta$ be a theta reading of $s_{y(c,x)}$ for an
analytic fibre at $y(c,x)$ for the second data, $x$ and $\Lambda_c$. Then
\[
 \mathcal G'_x(z)\prod_{k=1}^m\vartheta(z_k)=\Theta_\varphi(x;u)
\]
for every family $u=(u_1,\dots,u_m)$ as in Definition~\ref{is-eigenvalue-conjugate-frame-fibres}(1),
where $z=(z'_x(u_1),\dots,z'_x(u_m))$, $\mathcal G'_x$ and $z'_x$ are formed with the second basis and
indexing, and $\Theta_\varphi$ with the first.
\end{lemma}
\leanref{Development/AbelianVarietyTheory/Proofs/IwasawaDecomposition/OmegaFormScalarExtension.lean}{GlobalTheory.OneMapRef.FrameConst.exists_const}
\begin{proof}
Coordinates and theta series are formed with the first data unless the second are named. The element
$w_0=(1+\sqrt{-d})/2$ is an algebraic integer with $w_0+\bar w_0=1$. By
Lemma~\ref{is-eigenvalue-conjugate-frame-vacua}(1), applied to $X_c$ with the first data, $n=m$,
$L=\Lambda_c$ and $\Lambda_S=\mathcal O_E^m$, which lies in $\mathcal V_1$, the tensor lattice $\mathcal
O_E^m\otimes\Lambda_c$ has a vacuum $v$ for $w_0$. Let $C$ be
as in Lemma~\ref{is-eigenvalue-conjugate-frame-ortho-constant} for $v$, and put $\varphi=Cv$, a
Schwartz-Bruhat function. Let $x$, the second data and $\vartheta$ be as in the statement. By
Lemma~\ref{is-eigenvalue-conjugate-frame-renumbering}, applied to $X_c$ with the second data as $X'$ and
with the first data as $X$, $\gamma=1$, $y=y(c,x)$ and $L=\Lambda_c$, for which $\gamma x=x$, $\gamma
L=\Lambda_c$ and $u=u'$, there are an analytic fibre at $y(c,x)$ for the first data, $x$ and $\Lambda_c$ and
a theta reading $\vartheta_1$ of $s_{y(c,x)}$ for it with $\mathcal G'_x(z')\prod_k\vartheta(z'_k)=\mathcal
G_x(z)\prod_k\vartheta_1(z_k)$ for every $u$, where $z'$ and $z$ are the coordinates of $u$ for the second
and the first data. For every $u$,
this and Lemma~\ref{is-eigenvalue-conjugate-frame-ortho-constant} give
\[
 \mathcal G'_x(z')\prod_k\vartheta(z'_k)=\mathcal G_x(z)\prod_k\vartheta_1(z_k)
 =\mathcal G_x(z)\,C\,\theta_v(x;z)=\Theta_\varphi(x;u),
\]
by Definition~\ref{is-eigenvalue-conjugate-frame-fibres}(6), as $\theta_{Cv}=C\theta_v$.
\end{proof}
% exists_const: TensVac.vacua for v, exists_const_ortho for C, phi = C v (TensVac.sb_smul), anFib_move_renum
% with e1 and e2, and thR_smul.

\begin{lemma}\label{is-eigenvalue-conjugate-frame-marking-laws}
In Situation~\ref{is-eigenvalue-conjugate-frame-core}, with the notation of
Definition~\ref{is-eigenvalue-conjugate-frame-marking}, let $\ell_g$ be the endomorphism of $\mathfrak L$
through which $g\in G$ changes the level structures of the framed points (the action of $G$ that the
one-map data at $i$ give, Lemma~\ref{is-eigenvalue-conjugate-class-action}).
\begin{enumerate}
\item Let $\ell$ be a level of the datum and $g\in U(W_i)(\mathbf A)$. There is a marking at $g$ for $\ell$
  with values in $L(g)$. Let $m$ be one. Then: if $a\in\bar{\mathfrak M}_\ell\mathfrak L$ and $m(a)\in
  M_\ell L(g)$, then $a\in M_\ell\mathfrak L$; every $v\in\bar{\mathfrak M}_\ell L(g)$ is $m(a)+M_\ell u$ with
  $a\in\bar{\mathfrak M}_\ell\mathfrak L$ and $u\in L(g)$; every $\mathcal O_{E_h}$-linear $m'$ with
  $m'(a)-m(a)\in M_\ell L(g)$ for all $a\in\bar{\mathfrak M}_\ell\mathfrak L$ is a marking at $g$ for $\ell$;
  for every level $\ell_2$ with $M_\ell\mid M_{\ell_2}$ and $\bar{\mathfrak M}_{\ell_2}\subset\bar{\mathfrak
  M}_\ell$ and every marking $m_2$ at $g$ for $\ell_2$ with values in $L(g)$, $m_2(a)-m(a)\in M_\ell L(g)$ for
  all $a\in\bar{\mathfrak M}_{\ell_2}\mathfrak L$; and for every $\delta\in U(W_i)(F_h)$, with matrix
  $\gamma\in GL_r(E_h)$, $L(\delta g)=\gamma L(g)$ and $\gamma\circ m$ is a marking at $\delta g$ for $\ell$.
  Of the situation only $W_i$, $L_0$, $\mathfrak A$ and $\iota$ are used, and of $\iota$ only that it is
  injective with image $L(1)$.
\item Let $g\in G$ and $k$ a level lift of $g$ (Lemma~\ref{is-eigenvalue-core-group}), read as an element
  of $U(W_i)(\mathbf A)$ with archimedean component $1$. Then $k_f\hat\Lambda_W=\hat\Lambda_W$;
  $\prod_{P\in P_g}\chi_{i+1,P}(\det g_P)$ times the multiplier of the one-map data at $i$ at $k$ is $1$;
  $L(hk)=L(h)$ for every $h\in U(W_i)(\mathbf A)$; $m\circ\ell_g$ is a marking at $hk$ for every marking $m$
  at $h$; $gf$ is a core form for every core form $f$; and there is $y\in K_s$ with
  $\sigma_0(y)=\prod_{P\in P_g}\chi_{i+1,P}(\det g_P)$ and $s(y)\mu(k)=1$.
\item For $\epsilon=\pm1$ there is a central element $z$ of $G$ with $\ell_z(a)-\epsilon a\in M\mathfrak L$ for
  every $a\in\bar{\mathfrak M}\mathfrak L$.
\end{enumerate}
\end{lemma}
\leanref{Development/AbelianVarietyTheory/Proofs/IwasawaDecomposition/LatticeMarkingAdmissible.lean}{GlobalTheory.OneMapRef.markAd_laws}
\leanref{Development/AbelianVarietyTheory/Proofs/DivisorProofs/CoefMapNormEquivariantBasisOfFive.lean}{GlobalTheory.OneMapRef.Lift.lift_laws_ofFive}
\leanref{Development/AbelianVarietyTheory/Proofs/DivisorProofs/ActLevMemCuspA.lean}{GlobalTheory.GeoLevel.mu_lift}
\leanref{Development/AbelianVarietyTheory/Proofs/DivisorProofs/ActLevMemCuspA.lean}{GlobalTheory.GeoLevel.actLev_mem_cuspA}
\begin{proof}
(1) Write $N=M_\ell$, view $E_h^r$ inside $\mathbf A_{E_h,f}^r$, and let $g_f$ act on $\mathbf A_{E_h,f}^r$. We use
the following facts, for $g'\in U(W_i)(\mathbf A)$: $\hat L(g')=g'_f\hat L(1)$, as $g'_f$ is $\mathcal
O_{E_h}$-linear; $\hat L(g')$ is an $\mathcal O_{E_h}$-module; $L(g')=E_h^r\cap\hat L(g')$, and every $z\in\hat L(g')$
is $v+Nw$ with $v\in L(g')$ and $w\in\hat L(g')$, as $L(g')$ is a lattice of $E_h^r$ with completion $\hat L(g')$
\cite{Wei67}; and $L(1)=\iota(\mathfrak L)$ with $\iota$ injective. In particular, if $v,v'\in E_h^r$ and
$v-v'\in N\hat L(g')$, then $N^{-1}(v-v')\in L(g')$.

\emph{Existence.} Let $P$ be the set of the pairs $(v,a)\in L(g)\times\mathfrak L$ with $v-g_f\iota(a)\in N\hat L(g)$,
an $\mathcal O_{E_h}$-submodule. Its projection to $\mathfrak L$ is onto: for $a\in\mathfrak L$,
$g_f\iota(a)\in g_f\hat L(1)=\hat L(g)$ is $v+Nw$ with $v\in L(g)$ and $w\in\hat L(g)$, and $(v,a)\in P$. As
$\mathfrak L$ is a projective $\mathcal O_{E_h}$-module, the projection has an $\mathcal O_{E_h}$-linear section
$a\mapsto(m_0(a),a)$. Then $m_0\colon\mathfrak L\to L(g)$ is $\mathcal O_{E_h}$-linear with $m_0(a)-g_f\iota(a)\in
N\hat L(g)$ for every $a$; in particular it is a marking at $g$ for $\ell$ with values in $L(g)$. It has two more
properties.
\begin{enumerate}
\item[(a)] If $m_0(a)=Nu$ with $u\in L(g)$, then $a\in N\mathfrak L$. Indeed, with $m_0(a)-g_f\iota(a)=Nw$,
  $g_f\iota(a)=N(u-w)$ and $u-w\in\hat L(g)=g_f\hat L(1)$; so $\iota(a)=Ny$ with $y\in\hat L(1)$. Then
  $y=N^{-1}\iota(a)\in E_h^r\cap\hat L(1)=L(1)$, so $y=\iota(b)$, and $a=Nb$ as $\iota$ is injective.
\item[(b)] Every $v\in L(g)$ is $m_0(a)+Nu$ with $a\in\mathfrak L$ and $u\in L(g)$. Indeed, $v=g_fy$ with $y\in\hat
  L(1)$, and $y=\iota(a)+Nw$ with $a\in\mathfrak L$ and $w\in\hat L(1)$. With $m_0(a)-g_f\iota(a)=Nw'$,
  $v-m_0(a)=N(g_fw-w')$, and $u=N^{-1}(v-m_0(a))$ lies in $E_h^r\cap\hat L(g)=L(g)$.
\end{enumerate}

\emph{Two markings.} Let $\ell_2$ be a level with $N\mid M_{\ell_2}$ and $\bar{\mathfrak M}_{\ell_2}\subset\bar{\mathfrak
M}_\ell$, $m$ a marking at $g$ for $\ell$ and $m_2$ one for $\ell_2$, both with values in $L(g)$, and $a\in\bar{\mathfrak
M}_{\ell_2}\mathfrak L\subset\bar{\mathfrak M}_\ell\mathfrak L$. Write $M_{\ell_2}=cN$. Then $m(a)-g_f\iota(a)=Nw$ and
$m_2(a)-g_f\iota(a)=cNw'$ with $w,w'\in\hat L(g)$, so $m_2(a)-m(a)=N(cw'-w)$ and $N^{-1}(m_2(a)-m(a))\in L(g)$. This is
the fourth assertion; with $\ell_2=\ell$ it gives $m(a)-m_0(a)\in NL(g)$ for every marking $m$ at $g$ for $\ell$ with
values in $L(g)$ and every $a\in\bar{\mathfrak M}_\ell\mathfrak L$.

\emph{The properties of a marking.} Let $m$ be a marking at $g$ for $\ell$ with values in $L(g)$. If $a\in\bar{\mathfrak
M}_\ell\mathfrak L$ and $m(a)=Nu$ with $u\in L(g)$, then $m_0(a)=N(u+u')$ with $u'\in L(g)$, and $a\in N\mathfrak L$ by
(a). The elements $m(a)+Nu$ with $a\in\bar{\mathfrak M}_\ell\mathfrak L$ and $u\in L(g)$ form a group; for $s\in\bar{\mathfrak
M}_\ell$ and $v\in L(g)$, (b) gives $sv=m_0(sa)+N(su)=m(sa)+N(su+u')$ with $u'\in L(g)$; so every element of
$\bar{\mathfrak M}_\ell L(g)$ is of this form. If $m'$ is $\mathcal O_{E_h}$-linear with $m'(a)=m(a)+Nu$, $u\in L(g)$, for
$a\in\bar{\mathfrak M}_\ell\mathfrak L$, then $m'(a)-g_f\iota(a)=N(w+u)$ with $w+u\in\hat L(g)$, so $m'$ is a marking at $g$
for $\ell$.

\emph{Rational elements.} Let $\delta\in U(W_i)(F_h)$ with matrix $\gamma$; then $\delta_f$ acts on $E_h^r$ as $\gamma$, and
$\hat L(\delta g)=\delta_f\hat L(g)$. For $v\in E_h^r$, $v\in L(\delta g)$ exactly when $\gamma^{-1}v=\delta_f^{-1}v\in\hat
L(g)$, that is, when $\gamma^{-1}v\in L(g)$; so $L(\delta g)=\gamma L(g)$. For $a\in\bar{\mathfrak M}_\ell\mathfrak L$,
$\gamma m(a)-(\delta g)_f\iota(a)=\delta_f\big(m(a)-g_f\iota(a)\big)\in N\delta_f\hat L(g)=N\hat L(\delta g)$; so $\gamma\circ m$
is a marking at $\delta g$ for $\ell$. Only $W_i$, $L_0$, $\mathfrak A$ and the injectivity and image of $\iota$ were used.

(2) and (3).
\openstep{The proofs of (2) and (3), the behaviour of the lattices, markings and multipliers under the level lifts
of the elements of $G$, the value $y$, the action on core forms, and the central elements acting by $\pm1$ on the level
structures, are not written out.}
% (2) and (3): the first clause of Lift.lift_laws_ofFive (77 lines in CoefMapNormEquivariantBasisOfFive.lean, through
% GeoLevel.exists_half_elements, exists_root_char, chiTie_of_chiR, mu_cast, Member.latG_move, LevelUnique.redFn_eq_of_lift,
% exists_cong, cong_on_half, mark_calc), GeoLevel.mu_lift (28 lines) and GeoLevel.actLev_mem_cuspA (18 lines); (3) is
% the second clause (redFn_sign, exists_cong_sign, sign_div, sign_calc). Size (tools/lean_size.py, following imports
% and stopping at what the book cites): lift_laws_ofFive 6482 lines, mu_lift 2200, actLev_mem_cuspA 13134, together
% 18906. The top level of the first clause: k in GfC and IsMov for the lattice of W_i are immediate from the parahoric
% level (isMov_of_mem_St, parLev_St); the multiplier identity is exists_half_elements and exists_root_char (a root x of
% the level character with value prod_P chi_P(det g_P)^-1), chiTie_of_chiR and mu_cast; latG (h k) = latG h is
% Member.latG_move with finCong_of_mem_St; the marking clause compares k with the lift k' of O.lev g (hO.lift.lev g)
% through redFn_eq_of_lift, exists_cong and cong_on_half, then mark_calc.
% (1) is markAd_laws (52 lines; LatticeMarkingAdmissible.lean, 129 lines, through AdMark.exists_marking), for any
% W, iota with LatEmb W iota, h, L with latW L = latG h and any level l: existence, IsMark (inj, surj), the
% change of marking mk', MarkOver for a second level l2 (l.M | l2.M, Mfr l2 <= Mfr l), and the rational element
% delta (mapL ga L, mvMap ga mk). (2) is the first clause of Lift.lift_laws_ofFive (77 lines) for g and k with
% IsLevelLift: k in GfC, IsMov for the lattice of W_i, the product of R.chi with O.mu k equal to 1,
% latG (h k) = latG h and MarkAd h mk -> MarkAd (h k) (mk o O.lev g); with GeoLevel.mu_lift (28 lines) for y and
% GeoLevel.actLev_mem_cuspA (18 lines) for g f. (3) is the second clause of lift_laws_ofFive (O.lev z a = eps a +
% M b for a in Mfr l * La).
\end{proof}

\begin{lemma}\label{is-eigenvalue-conjugate-frame-level-move}
Let a PEL datum over a CM field $E$ with lattice $\mathfrak L$, a level, its moduli scheme and an embedding $s$
of a coefficient field of it be given. Let $\ell$ be an $\mathcal O_E$-linear endomorphism of $\mathfrak L$, $y$
a complex framed point of the moduli scheme over $s$, and $y'$ the move of $y$ by a level automorphism
$(\rho,\nu)$ with matrix $\ell$ (Definition~\ref{cv-realization-level-move}) that also preserves the full level
structure; this last condition is not used. Let $X$ be flag data of rank $r$ over $E$ with a symplectic basis, an
indexing and a bijection onto the letters, $x\in\mathcal D_X$, $L\subset E^r$ an $\mathcal O_E$-submodule
self-dual for $X$, and $A$ an analytic fibre at $y$ for these data, $x$ and $L$; that the basis is symplectic,
that $x\in\mathcal D_X$ and that $L$ is self-dual are not used. Then there is an analytic
fibre $A'$ at $y'$ for the same data, $x$ and $L$, such that every theta reading of $s_y$ for $A$ is a theta
reading of $s_{y'}$ for $A'$, and $A'$ is marked by $m\circ\ell$ whenever $A$ is marked by $m$ (as in
Definition~\ref{is-eigenvalue-conjugate-frame-marking}(2), with the level structure of the moduli scheme).
\end{lemma}
\leanref{Development/AbelianVarietyTheory/Proofs/IwasawaDecomposition/AnalyticFibreDescriptionLevelMoveMarkingComp.lean}{GlobalTheory.Dict.FibreMove.anFib_levR_id}
\begin{proof}
Let $\operatorname{ex}$, $\operatorname{fr}_\beta$ and $\chi$ be the data of $A$, and $\lambda_y$, $\lambda_{y'}$ the
level structures of $y$ and $y'$. By Definition~\ref{cv-realization-level-move}, $\nu$ restricts to an
isomorphism $\nu_y\colon A_y\to A_{y'}$ of algebraic groups, compatible with the actions of $\mathcal O_E$, with
$\lambda_{y'}(a)=\nu_y(\lambda_y(\ell a))$ for $a\in\mathfrak L$. The move carries the frame vectors $u_\ell$ of $y$ to
the frame vectors $d\nu_y(u_\ell)$ of $y'$, and it comes with an isomorphism $w$ of the line bundle $\mathcal L$ of
the polarization with its pullback along $\nu$ that respects the rigidifications along the zero sections and
carries $s_y$ to $s_{y'}$. In the trivializations, $w$ is given by nowhere vanishing regular functions
$w_{\beta\gamma}$, defined at the points $p$ of $A_y$ with $p\in U_\gamma$ and $\nu_y(p)\in U_\beta$, such that
$w_{\beta\gamma}(p)g_{\gamma\gamma'}(p)=g_{\beta\beta'}(\nu_y(p))w_{\beta'\gamma'}(p)$; the value of $s_{y'}$ at
$\nu_y(p)$ over $U_\beta$ is $w_{\beta\gamma}(p)$ times the value of $s_y$ at $p$ over $U_\gamma$; and
$w_{\beta\gamma}(0)$ carries the value of the rigidification at $y$ over $U_\gamma$ to its value at $y'$ over
$U_\beta$.

Put $\operatorname{ex}'=\nu_y\circ\operatorname{ex}$, $\chi'=\chi$, and
$\operatorname{fr}'_\beta(z)=w_{\beta\gamma}(\operatorname{ex}(z))\operatorname{fr}_\gamma(z)$ for $z$ with
$\operatorname{ex}'(z)\in U_\beta$, where $\gamma$ is any index with $\operatorname{ex}(z)\in U_\gamma$; by the
relation for $w$ with $\beta'=\beta$ and the cocycle relation of the $\operatorname{fr}_\gamma$, the value does not
depend on $\gamma$.
\begin{itemize}
\item[(a)] $\operatorname{ex}'$ is a surjective holomorphic group homomorphism with the kernel
  $\{p_x(v):v\in L\}$ of $\operatorname{ex}$, as $\nu_y$ is an isomorphism; it is compatible with $\mathcal O_E$,
  as $\nu_y$ is; and its derivative at $0$ in the direction of the $j$-th basis vector is $d\nu_y(u_{\ell_j})$,
  the frame vector of $y'$.
\item[(b)] Near a point $z_0$, with $\gamma$ fixed, $\operatorname{fr}'_\beta$ is the product of the holomorphic
  functions $w_{\beta\gamma}\circ\operatorname{ex}$ and $\operatorname{fr}_\gamma$, so it is holomorphic, and it has
  no zeros. The relation for $w$ and the cocycle relation of the $\operatorname{fr}_\gamma$ give
  $\operatorname{fr}'_\beta(z)=g_{\beta\beta'}(\operatorname{ex}'(z))\operatorname{fr}'_{\beta'}(z)$. At $z=0$,
  $\operatorname{fr}_\gamma(0)$ is the value of the rigidification at $y$ over $U_\gamma$, so
  $\operatorname{fr}'_\beta(0)$ is its value at $y'$ over $U_\beta$.
\item[(c)] For $v\in L$, $\operatorname{ex}(z+p_x(v))=\operatorname{ex}(z)$, so the factor
  $w_{\beta\gamma}(\operatorname{ex}(z))$ is the same at $z$ and at $z+p_x(v)$, and the law (c) for $A$ gives the
  law (c) for $\operatorname{fr}'$ with the same $\chi$.
\end{itemize}
So $A'=(\operatorname{ex}',\operatorname{fr}',\chi')$ is an analytic fibre at $y'$ for the same data, $x$ and
$L$. If $\vartheta$ is a theta reading of $s_y$ for $A$, then the value of $s_{y'}$ at
$\operatorname{ex}'(z)=\nu_y(\operatorname{ex}(z))$ over $U_\beta$ is
$w_{\beta\gamma}(\operatorname{ex}(z))\vartheta(z)\operatorname{fr}_\gamma(z)=\vartheta(z)\operatorname{fr}'_\beta(z)$; so
$\vartheta$ is a theta reading of $s_{y'}$ for $A'$. Finally, let $A$ be marked by $m$, and
$a\in\bar{\mathfrak M}\mathfrak L$. Then $\ell a\in\bar{\mathfrak M}\mathfrak L$, as $\ell$ is $\mathcal
O_E$-linear, and
\[
 \lambda_{y'}(a)=\nu_y(\lambda_y(\ell a))=\nu_y\big(\operatorname{ex}(p_x(M^{-1}m(\ell a)))\big)
 =\operatorname{ex}'\big(p_x(M^{-1}(m\circ\ell)(a))\big).
\]
So $A'$ is marked by $m\circ\ell$.
\end{proof}
% Dict.FibreMove.anFib_levR_id = anFib_levR (102 lines; AnalyticFibreDescriptionLevelMoveMarkingComp.lean, 161
% lines), for any gM and levR (gM, id) y y', under IsSympl, x in dom and SelfD, which anFib_levR_id does not pass
% on (nor IsIso rho and the etaA clause of IsLevMove). The proof follows anFib_levR: MoveCalc.frMove_iff_carry
% (y'.x = y.x rho, y'.u = y.u nu, pm, the line bundle isomorphism w with CubIso and the Carry of the sections),
% exists_mulEquiv for nu_y (Phi), exists_frame_push for fr' (wv w and wv_change, val_g_self), val_rig for the
% value at 0, MoveValues.val_carry for the reading, and hnu.2.1 (eta (gM a) nu = rho eta a) for the marking.
% Equivar uses it with gM = O.lev g and the moved point mvPt of mvPt_spec.

\begin{lemma}\label{is-eigenvalue-conjugate-frame-compare-scalar}
Let a PEL datum over a CM field $E$ with lattice $\mathfrak L$, a level $(M,\mathfrak M,N)$ with full level $N=1$,
its moduli scheme and an embedding $s$ of a coefficient field $K_c$ of it be given. Let $X$ be flag data of rank
$r$ over $E$ with a symplectic basis, an indexing and a bijection onto the letters, $x\in\mathcal D_X$, and
$L\subset E^r$ an $\mathcal O_E$-submodule self-dual for $X$. Let $m\colon\mathfrak L\to E^r$ be an $\mathcal
O_E$-linear map with values in $L$ such that: if $a\in\bar{\mathfrak M}\mathfrak L$ and $m(a)\in ML$, then $a\in
M\mathfrak L$; and every $v\in\bar{\mathfrak M}L$ is $m(a)+Mu$ with $a\in\bar{\mathfrak M}\mathfrak L$ and $u\in L$.
Let $y$ and $y'$ be complex framed points of the moduli scheme over $s$, $A$ and $A'$ analytic fibres at $y$ and
$y'$ for these data, $x$ and $L$, both marked by $m$, and $\vartheta$ and $\vartheta'$ theta readings of $s_y$ and
$s_{y'}$ for $A$ and $A'$. Then there is $\lambda\in\mathbf C^\times$ with $\vartheta'=\lambda\vartheta$ and
$f(y')=\lambda^nf(y)$ for every section $f$ of the moduli scheme over $K_c$ whose weight has the exponent $n$ at
the line bundle of the polarization.
\end{lemma}
\leanref{Development/AbelianVarietyTheory/Proofs/IwasawaDecomposition/AnalyticFibreMatching.lean}{GlobalTheory.OneMapRef.exists_scalar_of_marked}
\leanref{Development/AbelianVarietyTheory/Proofs/DivisorProofs/ExistsObjMapPhi.lean}{GlobalTheory.Dict.anFib_ex_eq}
\leanref{Development/AbelianVarietyTheory/Proofs/DivisorProofs/AnFibFrEqOfExLeftEq.lean}{GlobalTheory.FrameEquiv.anFib_fr_eq_of_ex_left_eq}
\begin{proof}
Let $\operatorname{ex}$, $\operatorname{fr}_\beta$, $\chi$ and $\operatorname{ex}'$, $\operatorname{fr}'_\beta$, $\chi'$ be
the data of $A$ and $A'$.

\emph{Step 1: the exponentials.} The exponentials agree as maps to the universal abelian scheme:
$\operatorname{ex}'(z)=\operatorname{ex}(z)$ for every $z\in\mathbf C^d$.
\openstep{The proof of Step 1, which shows that the abelian varieties of $y$ and $y'$, described by the two analytic
fibres with the same periods, form and marking, are one object of the moduli problem, so that the two exponentials
agree, is not written out.}
% Step 1 is Dict.anFib_ex_eq (about 40 lines; ExistsObjMapPhi.lean, 82 lines), through Dict.anFib_point_eq,
% AnFib.exists_anPt, obj_pull, SameFibre.exists_holUnif, SameFibre.exists_sabMap_of_isAH (ExistsSabMapOfIsAH.lean,
% 178 lines) and SameExp.exists_objMap_phi, uses l.A = 1, IsMark and the markings; about 370 lines in its four files,
% with the descent of holomorphic maps of abelian schemes behind them. Size (tools/lean_size.py anFib_ex_eq): 9509
% lines. Top level: the two fibres give analytic points (exists_anPt); the universal object pulls back along y and y'
% (obj_pull) and is uniformized (exists_holUnif); the uniformizations have the same lattice, form and O_E-types, so
% exists_sabMap_of_isAH gives an isomorphism m of the two abelian varieties with m o ex = ex'; the markings and l.A = 1
% make it a map of objects of the moduli problem (exists_objMap_phi), and representability (moduli_rep_of_base,
% phi_eq_of_eq) gives y = y' with m the identity.

As $y$ and $y'$ are the images of $\operatorname{ex}(0)=\operatorname{ex}'(0)$, they have the same point, so
$A_y=A_{y'}$ with the same line bundle, transition functions and rigidification; and they have the same frame
vectors, the derivatives at $0$ of the common exponential.

\emph{Step 2: the trivializing functions.} For $z\in\mathbf C^d$ let $q(z)=\operatorname{fr}_\beta(z)/\operatorname{fr}'_\beta(z)$
for any $\beta$ with $\operatorname{ex}(z)\in U_\beta$; by the cocycle relations of (b), with the same transition
functions, this does not depend on $\beta$. The function $q$ is holomorphic and has no zeros. For $v\in L$,
$\operatorname{ex}(z+p_x(v))=\operatorname{ex}(z)$, and the laws (c) of $A$ and $A'$, with $|\chi(v)|=|\chi'(v)|=1$ and
the same exponential factor, give $|q(z+p_x(v))|=|q(z)|$. As $L$ is self-dual, it spans $E^r$, and the $p_x(v)$,
$v\in L$, span $\mathbf C^d$ over $\mathbf R$ (Lemma~\ref{is-eigenvalue-conjugate-frame-hermitian}(4)); so $|q|$ is
bounded, and $q$ is constant by Liouville's theorem. At $z=0$ both $\operatorname{fr}_\beta(0)$ and
$\operatorname{fr}'_\beta(0)$ are the value of the rigidification of the common point. So $q=1$ and
$\operatorname{fr}'_\beta=\operatorname{fr}_\beta$.

\emph{Step 3: the scalar.} The sections $s_y$ and $s_{y'}$ both generate the sections of the line bundle on the common
fibre, so $s_{y'}=\lambda s_y$ for some $\lambda\in\mathbf C^\times$. Then $\lambda\vartheta$ is a theta reading of
$s_{y'}$ for $A'$, by Steps 1 and 2; and a theta reading is determined by the section and the analytic fibre
(Definition~\ref{is-eigenvalue-conjugate-frame-fibres}(4), as $\operatorname{fr}_\beta$ has no zeros). So
$\vartheta'=\lambda\vartheta$. Finally $y'$ is $y$ with the same point and frame and the generator multiplied by
$\lambda$, so the transformation law of sections gives $f(y')=\lambda^nf(y)$ for every section $f$ whose weight has
the exponent $n$ at the line bundle of the polarization.
\end{proof}
% exists_scalar_of_marked (35 lines; AnalyticFibreMatching.lean, 57 lines), under IsSympl, x in dom, SelfD, l.A = 1,
% IsMark l L mk, A.Marked mk and A'.Marked mk, for readings th and th' with k = 1: anFib_ex_eq (Step 1), hyx and hu
% (isDerivAt0_unique), anFib_fr_eq_of_ex_left_eq (Step 2: q, key, hq, hper, norm_periodic_const_range with
% span_image_pv, h0 from fr_0), Frames.ChgOf.gen_unit' for lambda, reads_existsUnique_theta (uniqueness) with
% sec_values_complex_fibre for the reading, and LamT.evR_frChg_tri for the sections (the frame change with t = lambda
% at the generator and 1 elsewhere).

\begin{lemma}\label{is-eigenvalue-conjugate-frame-compare}
In Situation~\ref{is-eigenvalue-conjugate-frame-core}, with the notation of
Definition~\ref{is-eigenvalue-conjugate-frame-marking}, let $\delta\in U(W_i)(F_h)$ and $a'\in U(W_i)(\mathbf
A)$, and let $c,c'$ be components and $x\in\mathcal D_c$, $x'\in\mathcal D_{c'}$ with $\Lambda_c=L(\delta a')$
and $\Lambda_{c'}=L(a')$, both self-dual for their flag data, with the integer $p_0=q$ for $X_c$ and
$X_{c'}$, and with $x=x^{\delta a'Z_0}$ and $x'=x^{a'Z_0}$. Let symplectic bases of $\Lambda_c$ and
$\Lambda_{c'}$, indexings and bijections onto the letters be given, matched for $s$, such that $\psi$,
$\theta$, $S=V_i$, $T=W_i$, the two flag data, bases and indexings satisfy the conditions of
Situation~\ref{is-eigenvalue-conjugate-frame-situation}. Let $y$ and $y'$ be complex framed points of the
moduli of the core at level $i+1$, $A$ and $A'$ analytic fibres at $y$ and $y'$ for these data, $x$,
$\Lambda_c$ and $x'$, $\Lambda_{c'}$, $\vartheta$ and $\vartheta'$ theta readings of $s_y$ and $s_{y'}$ for
them, and $m$ and $m'$ markings at $\delta a'$ and at $a'$ with values in $\Lambda_c$ and $\Lambda_{c'}$ such
that $A$ is marked by $m$ and $A'$ by $m'$. Then there is $\lambda\in\mathbf C$ with the following
properties.
\begin{enumerate}
\item $f(y')=\lambda^nf(y)$ for every section $f$ of the core at level $i+1$ over $K_s$ whose weight has the
  exponent $n$ at the line bundle of the polarization.
\item For every $w\in\mathbb W$ with finite component $0$ there is $w'\in\mathbb W$ with finite component
  $0$ such that $\mathcal G'_{x'}(z')\prod_k\vartheta'(z'_k)=\lambda^n\mathcal G_x(z)\prod_k\vartheta(z_k)$,
  with $n=\operatorname{rk}V_i$, $z=z_x(u(w))$, $z'=z'_{x'}(u(w'))$ and the Gaussians of the two data, and
  $(\omega(1,\delta a')f)(w)=(\omega(1,a')f)(w')$ for every Weil datum $\omega$ for $(V_i,W_i)$ and every
  theta function $f$.
\end{enumerate}
\end{lemma}
\leanref{Development/AbelianVarietyTheory/Proofs/IwasawaDecomposition/OmegaFormHeckeEigen.lean}{GlobalTheory.OneMapRef.Equivar.compare}
\begin{proof}
By the laws of the weight tie, the level of the core has full level $N=1$, and $\iota$ is injective with image
$L(1)$. By Lemma~\ref{is-eigenvalue-conjugate-frame-points}(2), $Z_0$ is a signature family for $T=W_i$ and $p_0=q$.
Write $X=X_c$ and $X'=X_{c'}$, with forms $h$ and $h'$. By the conditions of
Situation~\ref{is-eigenvalue-conjugate-frame-situation}, both have the set $\mathcal T$, the matrix
$-G_{W_i}^{\mathsf T}$, the element $-\xi_0$ and $\kappa_0=-2\xi_0$, and both have $p_0=q$; so they have the same set,
matrix, integer $p_0$ and elements $\xi_0$ and $\kappa_0$.

\emph{The rational element.} Let $\gamma\in GL_r(E_h)$ be the matrix of $\delta$. As $\delta$ is unitary for $G_{W_i}$,
$h(\gamma u,\gamma v)=h(u,v)$ for $u,v\in E_h^r$. As $\theta_\rho$ extends $\rho$, $\delta_\rho=\rho(\gamma)$, so
$\gamma x^{Z}=x^{\delta Z}$ for every family $Z$; in particular, by
Lemma~\ref{is-eigenvalue-conjugate-frame-signature}(3), $\gamma x'=\gamma x^{a'Z_0}=x^{\delta a'Z_0}=x$. For
$w\in\mathbb W$ put $\delta w=(1,\delta)\cdot w$, the action of Definition~\ref{th-weil-datum}; then $w$ has finite
component $0$ exactly when $\delta^{-1}w$ has, and $u(\delta w)_k$ has the component $\rho(\gamma)u(w)_{k,\rho}$ at
$\rho$ (Situation~\ref{is-eigenvalue-conjugate-frame-situation}(2)). For a Weil datum $\omega$, a theta function $f$
and $w\in\mathbb W$, the function $\rho(w)f$ is a theta function, and by the axioms of
Definition~\ref{th-weil-datum}(3)
\[
 (\omega(1,\delta)f)(\delta w)=\big(\rho(\delta w)\,\omega(1,\delta)f\big)(0)=\big(\omega(1,\delta)\rho(w)f\big)(0)
 =(\rho(w)f)(0)=f(w).\tag{$*$}
\]

\emph{The markings.} By Lemma~\ref{is-eigenvalue-conjugate-frame-marking-laws}(1) for $g=a'$ and the marking $m'$,
$L(\delta a')=\gamma L(a')$, so $\Lambda_c=\gamma\Lambda_{c'}$, and $\gamma\circ m'$ is a marking at $\delta a'$ with
values in $\Lambda_c$.

\emph{Renumbering.} Lemma~\ref{is-eigenvalue-conjugate-frame-renumbering}, applied to $X'$ with its data as $X'$, $X$
with its data as $X$, the point $x'$, $L=\Lambda_{c'}$, $\gamma$, $y'$ and $\vartheta'$, gives an analytic fibre $A_1$
at $y'$ for the data of $X$, $\gamma x'=x$ and $\gamma\Lambda_{c'}=\Lambda_c$, marked by $\gamma\circ m'$, and a theta
reading $\vartheta_1$ of $s_{y'}$ for $A_1$ with
\[
 \mathcal G_x(z)\prod_k\vartheta_1(z_k)=\mathcal G'_{x'}(z')\prod_k\vartheta'(z'_k)\tag{$**$}
\]
whenever $(u_k)_\rho=\rho(\gamma)(u'_k)_\rho$ for all $k$ and $\rho$, with $z=(z_x(u_k))_k$ and $z'=(z'_{x'}(u'_k))_k$.
By Lemma~\ref{is-eigenvalue-conjugate-frame-marking-laws}(1), for the markings $m$ and $\gamma\circ m'$ at $\delta a'$
with values in $\Lambda_c$, $(\gamma\circ m')(a)-m(a)\in M\Lambda_c$ for $a\in\bar{\mathfrak M}\mathfrak L$; as the
exponential of $A_1$ is trivial on $p_x(\Lambda_c)$, $A_1$ is also marked by $m$. And $m$ has the two properties of
Lemma~\ref{is-eigenvalue-conjugate-frame-marking-laws}(1) for the level of the core.

\emph{The scalar.} Lemma~\ref{is-eigenvalue-conjugate-frame-compare-scalar}, applied to $X$ with its data, $x$,
$L=\Lambda_c$, $m$, the analytic fibres $A$ at $y$ and $A_1$ at $y'$ and the readings $\vartheta$ and $\vartheta_1$,
gives $\lambda\in\mathbf C^\times$ with $\vartheta_1=\lambda\vartheta$ and $f(y')=\lambda^nf(y)$ for every section $f$
of the core at level $i+1$ over $K_s$ whose weight has the exponent $n$. This is (1).

(2) Let $w\in\mathbb W$ have finite component $0$, and put $w'=\delta^{-1}w$, which has finite component $0$; then
$w=\delta w'$ and $u(w)_{k,\rho}=\rho(\gamma)u(w')_{k,\rho}$. By $(**)$ for $u=u(w)$ and $u'=u(w')$, and as
$\vartheta_1=\lambda\vartheta$ with $n=\operatorname{rk}V_i$ factors,
\[
 \mathcal G'_{x'}(z')\prod_k\vartheta'(z'_k)=\mathcal G_x(z)\prod_k\vartheta_1(z_k)=\lambda^n\mathcal
 G_x(z)\prod_k\vartheta(z_k).
\]
For a Weil datum $\omega$ and a theta function $f$, $\omega(1,a')f$ is a theta function, and $(*)$ gives
$(\omega(1,\delta a')f)(w)=(\omega(1,\delta)\omega(1,a')f)(\delta w')=(\omega(1,a')f)(w')$.
\end{proof}
% Equivar.compare (61 lines; OmegaFormHeckeEigen.lean, 172 lines), followed: sameSp_of_isSet; RatPt.rational_points
% (ga with IsUn, lact ga (frOf Z) = frOf (mvZ delta Z), the real points w = actO 1 delta w' with rePt (actO 1 delta w)
% = ga rePt w, and (*) through om_heis, om_aut, isTheta_heisO and heisO_at_zero); markAd_laws (rational clause) for
% mapL ga Lambda' = Lambda and MarkAd for mvMap ga mk'; oneg_frames (mvZ_mul) for eX; anFib_move_renum for A1, th1
% (r1, r2, r3); markAd_laws (MarkOver, IsMark) and marked_congr; exists_scalar_of_marked for lam (Lemma
% -compare-scalar). Under l.A = 1 and LatEmb iota (laws of the
% tie), IsSet for both data, Matched, p0 equal, SelfD for both lattices, delta in GQ, latW = latG for both, Z0
% with IsONeg, and the frames frOf (mvZ (delta a') Z0) and frOf (mvZ a' Z0); through RatPt.rational_points and
% the last clause of markAd_laws.

\subsection{Extremal vectors and the phase}\label{subsection-is-extremal-phase}

\begin{lemma}\label{is-eigenvalue-conjugate-frame-extremal-orthogonal}
Let $G\in M_n(\mathbf C)$ be a Hermitian matrix, $0\le q\le n$, and $x_0\in M_n(\mathbf C)$ such that $x_0^*Gx_0$ is
the diagonal matrix with entries $-1$ at the first $q$ indices and $1$ at the others. Let $U(G)$ be the group of the
$g$ with $g^*Gg=G$, $\mathfrak u$ its Lie algebra, the $X$ with $X^*G+GX=0$, so that every $Z\in M_n(\mathbf C)$ is
$A+iB$ with unique $A,B\in\mathfrak u$, and $\mathcal B$ the group of invertible upper triangular matrices. Let
$\mathcal D_G$ be the set of $x\in M_n(\mathbf C)$ with $\det x\ne0$ and $(-1)^{\max(0,i-q)}\operatorname{Re}d_i(x)>0$
for $0\le i\le n$, where $d_i(x)$ is the leading principal $i\times i$ minor of $-x^*Gx$. Let $w\in\mathbf Z^n$.
A function $f$ on $U(G)$ is \emph{holomorphic of weight $w$} if there is a function $P$ on $M_n(\mathbf C)$,
holomorphic on $\mathcal D_G$, with $P(xb)=\prod_ab_{aa}^{w_a}P(x)$ for $x\in\mathcal D_G$ and $b\in\mathcal B$, such
that $f(u)=P(ux_0)$ for $u\in U(G)$ (Definition~\ref{is-eigenvalue-conjugate-dictionary}(2), for one embedding). For
$s\in\mathbf R^n$ put $t(s)=x_0\operatorname{diag}(e^{is_1},\dots,e^{is_n})x_0^{-1}$ and $\chi_w(s)=e^{i\sum_as_aw_a}$.
For a function $f$ on $U(G)$, $X\in\mathfrak u$ and $Z=A+iB$ as above write $D_Xf(u)$ for the derivative at $t=0$ of
$t\mapsto f(u\exp(tX))$ and $D_Zf=D_Af+i\,D_Bf$; for a representation $\pi$ of $U(G)$ and a smooth vector $v$ write
$d\pi(Z)v=d\pi(A)v+i\,d\pi(B)v$. Write $E_{ac}$ for the matrix units.
\begin{enumerate}
\item Let $f$ be a function on $U(G)$ such that $t\mapsto f(u\exp(t_1X_1)\cdots\exp(t_kX_k))$ is smooth on
  $\mathbf R^k$ for all $u\in U(G)$ and $X_1,\dots,X_k\in\mathfrak u$, and such that for $a\le c$ the derivative
  $D_Zf$ along $Z=x_0E_{ac}x_0^{-1}$ is $0$ if $a<c$ and $w_af$ if $a=c$. Then $f$ is holomorphic of weight $w$.
\item If $f$ is holomorphic of weight $w$, $f(t(s)u)=\chi_w(s)f(u)$ for all $u\in U(G)$ and $s$, and $f(1)=0$, then
  $f=0$.
\end{enumerate}
Let $\pi$ be an irreducible unitary representation of $U(G)$ on a Hilbert space $H$, with the inner product
conjugate linear in the first variable, and $v\in H$ a nonzero smooth vector such that, for $a\le c$,
$d\pi(x_0E_{ac}x_0^{-1})v$ is $0$ if $a<c$ and $w_av$ if $a=c$.
\begin{enumerate}
\setcounter{enumi}{2}
\item For every $l\in H$ the function $u\mapsto\langle l,\pi(u)v\rangle$ on $U(G)$ is holomorphic of weight $w$.
\item If $l\in H$ satisfies $\pi(t(s))l=\chi_w(s)l$ for every $s$ and $\langle l,v\rangle=0$, then $l=0$.
\end{enumerate}
\end{lemma}
\leanref{Development/AbelianVarietyTheory/Proofs/DivisorProofs/HolomorphicTypeOfSmoothBorelKilled.lean}{GlobalTheory.HolDict.holFn_of_killed}
\leanref{Development/AbelianVarietyTheory/Proofs/DivisorProofs/HolomorphicTypeEqZeroOfTorusLawAtOne.lean}{GlobalTheory.HolDict.Uniq.eq_zero_of_torus_law}
\leanref{Development/AbelianVarietyTheory/Proofs/IwasawaDecomposition/HolomorphicDiscreteSeriesLowestKTypeSmoothIrredu.lean}{GlobalTheory.HolLowOnce.smFn_coef}
\leanref{Development/AbelianVarietyTheory/Proofs/IwasawaDecomposition/HolomorphicDiscreteSeriesLowestKTypeSmoothIrredu.lean}{GlobalTheory.HolLowOnce.derC_coef}
\leanref{Development/AbelianVarietyTheory/Proofs/IwasawaDecomposition/HolomorphicDiscreteSeriesLowestKTypeSmoothIrredu.lean}{GlobalTheory.HolLowOnce.eq_zero_of_orth}
% (1) is holFn_of_killed (87 lines; HolomorphicTypeOfSmoothBorelKilled.lean, 505 lines: the smooth Gram-Schmidt
% decomposition x = u x0 b on the frame domain, P(x) = chi_w(b) f(u), and its holomorphy from the killed
% derivatives), and (2) is Uniq.eq_zero_of_torus_law (176 lines; HolomorphicTypeEqZeroOfTorusLawAtOne.lean, 401
% lines, through frDoms_open_preconnected_stable and the identity theorem), both for one factor (iota = Unit); they
% are self-contained computations, and Lemma is-eigenvalue-conjugate-spectral-facts(3)(a) and (b) are their
% counterparts on the product of the places.
\begin{proof}
(3) Let $f(u)=\langle l,\pi(u)v\rangle$. As $v$ is smooth and $\pi$ is a representation by continuous operators,
$t\mapsto f(u\exp(t_1X_1)\cdots\exp(t_kX_k))=\langle l,\pi(u)\pi(\exp(t_1X_1)\cdots\exp(t_kX_k))v\rangle$ is smooth,
and $D_Xf(u)=\langle l,\pi(u)\,d\pi(X)v\rangle$ for $X\in\mathfrak u$, since $y\mapsto\langle l,\pi(u)y\rangle$ is
continuous and linear. For $Z=A+iB$ this gives $D_Zf(u)=\langle l,\pi(u)\,d\pi(Z)v\rangle$, as the inner product is
linear in the second variable. By the hypothesis on $v$, for $a\le c$ and $Z=x_0E_{ac}x_0^{-1}$ this is $0$ if $a<c$
and $w_af(u)$ if $a=c$. So $f$ is holomorphic of weight $w$ by (1).

(4) Let $f(u)=\langle l,\pi(u)v\rangle$, holomorphic of weight $w$ by (3). The $t(s)$ lie in $U(G)$, as $x_0^*Gx_0$ is
diagonal with entries $\pm1$. For $u\in U(G)$ and $s$, $\pi(t(s))$ is an isometry, so
\[
 f(u)=\langle\pi(t(s))l,\pi(t(s))\pi(u)v\rangle=\overline{\chi_w(s)}\,f(t(s)u),
\]
and $|\chi_w(s)|=1$; hence $f(t(s)u)=\chi_w(s)f(u)$. And $f(1)=\langle l,v\rangle=0$. By (2), $f=0$. Let $P$ be the
set of the $y\in H$ with $\langle l,\pi(g)y\rangle=0$ for every $g\in U(G)$. It is a closed subspace, stable under
$\pi$ as $U(G)$ is a group, and it contains $v$, as $f=0$. As $\pi$ is irreducible and $v\ne0$, $P=H$. So $l\in P$,
and $\langle l,l\rangle=0$ at $g=1$; hence $l=0$.
\end{proof}
% (3) is smFn_coef and derC_coef (through der_coef and the decomposition of Z by conjH) with holFn_of_killed; (4) is
% HolLowOnce.eq_zero_of_orth (54 lines), followed: hf, the torus law through inner_iso and hn, h1, hz from
% Uniq.eq_zero_of_torus_law, the closed stable subspace P and hirr.

\begin{lemma}\label{is-eigenvalue-conjugate-frame-extremal-fock}
\begin{enumerate}
\item Let $\mathcal M$ be a complex vector space, $I$ a set, $a_k$ and $b_k$ ($k\in I$) linear endomorphisms of
  $\mathcal M$ and $c_k\in\mathbf C$ nonzero, such that $a_ka_l=a_la_k$ and $b_ka_l=a_lb_k+\delta_{kl}c_k$ for all
  $k,l\in I$. Let $f\in\mathcal M$ be nonzero with $b_kf=0$ for every $k$. Then the linear map from the polynomial
  ring $\mathbf C[X_k:k\in I]$ to $\mathcal M$ that sends $P$ to $P(a)f$ is injective.
\item Let $D$, $\psi$, $\theta$, $S$ and $T$ be as in Lemma~\ref{is-eigenvalue-conjugate-theta-facts}, with
  $\partial_\tau(z)$ as in Lemma~\ref{is-eigenvalue-conjugate-theta-closed}. For $\tau,\tau'\in\mathcal T$, complex
  $\operatorname{rk}T\times\operatorname{rk}S$ matrices $z,z'$ and every theta function $f$ for $(S,T)$,
  \[
   \partial_\tau(z)\partial_{\tau'}(z')f-\partial_{\tau'}(z')\partial_\tau(z)f=-2\pi i\,\omega_{\tau\tau'}(z,z')f,
  \]
  where $\omega_{\tau\tau'}(z,z')=0$ if $\tau\ne\tau'$ and $\omega_{\tau\tau}(z,z')=2\operatorname{Re}\big(\tau(\xi_0)
  \operatorname{tr}(G_S(\tau)^{-1}z'^*G_T(\tau)z)\big)$, with $G_S(\tau)$ and $G_T(\tau)$ the images of the Gram matrices.
\item In the setting of Lemma~\ref{is-eigenvalue-conjugate-theta-facts}, let the bases $B_{S,\tau}$ and $B_{T,\tau}$ be
  orthonormal and adapted to $s_{S,\tau}$ and $s_{T,\tau}$: $B_{S,\tau}^*G_S(\tau)B_{S,\tau}$ is a diagonal matrix
  $\operatorname{diag}(\epsilon_\tau)$ of signs and $s_{S,\tau}B_{S,\tau}=B_{S,\tau}\operatorname{diag}(\epsilon_\tau)$,
  and the same for $T$ with signs $\epsilon'_\tau$. Let $f$ be a nonzero theta function annihilated by the annihilation
  operators for $s_S$ and $s_T$ at every $\tau$. Then $P\mapsto f\otimes_{s,B}P$ is injective on the polynomials in
  the creation variables.
\end{enumerate}
\end{lemma}
\leanref{Development/AbelianVarietyTheory/Proofs/DivisorProofs/PureTensorOInjective.lean}{GlobalTheory.FockVecInj.injective_of_canonicalCommutation}
\leanref{Development/AbelianVarietyTheory/Proofs/DivisorProofs/PureTensorOInjective.lean}{GlobalTheory.FockVecInj.injective_of_pderiv}
\leanref{Development/AbelianVarietyTheory/Proofs/DivisorProofs/PureTensorOInjective.lean}{GlobalTheory.FockVecInj.omO_single}
\leanref{Development/AbelianVarietyTheory/Proofs/DivisorProofs/PureTensorOInjective.lean}{GlobalTheory.FockVecInj.ccr_of_comm}
\leanref{Development/AbelianVarietyTheory/Proofs/DivisorProofs/PureTensorOInjective.lean}{GlobalTheory.FockVecInj.pureTensorO_injective}
\begin{proof}
(1) Write $E(P)=P(a)f$, a linear map; as the $a_k$ commute, $E(PX_l)=a_lE(P)$ and $E(1)=f$. First,
$b_kE(P)=c_kE(\partial_kP)$ for every $P$, where $\partial_k$ is the partial derivative in $X_k$: this holds for
constants, as $b_kf=0$; it is additive in $P$; and if it holds for $P$, then
$b_kE(PX_l)=b_ka_lE(P)=a_lb_kE(P)+\delta_{kl}c_kE(P)=c_k\big(a_lE(\partial_kP)+\delta_{kl}E(P)\big)=c_kE(\partial_k(PX_l))$.
Hence, as $c_k\ne0$, $E(P)=0$ implies $E(\partial_kP)=0$ for every $k$. We show by induction on $N$ that $E(Q)=0$ and
$Q$ of degree less than $N$ imply $Q=0$. For $N=0$, $Q=0$. Let $N>0$. Every $\partial_kQ$ has degree less than $N-1$
and $E(\partial_kQ)=0$, so $\partial_kQ=0$ by induction; so $Q$ is a constant $c$, and $0=E(Q)=cf$ gives $c=0$, as
$f\ne0$. Taking $N$ larger than the degree of $Q$ gives the claim.

(2) For a theta function $g$ and a real direction $L_{\tau,z}$, write $D_zg(w)$ for the derivative at $t=0$ of
$t\mapsto g(w+L_{\tau,z}(t))$. As $L_{\tau,z}(t)$ has finite component $0$ and $\psi$ is standard,
$\psi(\frac12\langle w,L_{\tau,z}(t)\rangle)=e^{-\pi it\,\omega_{\tau\tau}(w_\tau,z)}$, where $w_\tau$ is the component
of $w$ at $\tau$ under $\theta$; so
$\partial_\tau(z)g(w)=-\pi i\,\omega_{\tau\tau}(w_\tau,z)g(w)+D_zg(w)$. The form $\omega_{\tau\tau}$ is real bilinear and
alternating, as $\tau(\xi_0)$ is purely imaginary (Definition~\ref{th-definite-datum}) and
$\overline{\operatorname{tr}(G_S^{-1}z'^*G_Tz)}=\operatorname{tr}(G_S^{-1}z^*G_Tz')$. Applying the formula twice, using
that $D_z$ and $D_{z'}$ commute on $g$, which is smooth along the archimedean directions, and that $D_z$ of
$w\mapsto\omega(w_\tau,z')$ is $\omega(z,z')$, all terms but one are symmetric in $(z,\tau)$ and $(z',\tau')$, and
the commutator is $-\pi i\big(\omega_{\tau\tau'}(z,z')-\omega_{\tau'\tau}(z',z)\big)f=-2\pi i\,\omega_{\tau\tau'}(z,z')f$.
By Lemma~\ref{is-eigenvalue-conjugate-theta-closed}(2) all the functions involved are theta functions.

(3) Let $\mathcal M$ be the space of theta functions, which the operators $\partial_\tau(z)$ preserve
(Lemma~\ref{is-eigenvalue-conjugate-theta-closed}(2)). For a creation variable $k=(\tau,i,j)$ put
$u_k=B_{T,\tau}E_{ij}B_{S,\tau}^{-1}$ and $\eta_k=\epsilon'_{\tau,i}\epsilon_{\tau,j}$. As $s_{T,\tau}B_{T,\tau}=B_{T,\tau}
\operatorname{diag}(\epsilon'_\tau)$ and $B_{S,\tau}^{-1}s_{S,\tau}=\operatorname{diag}(\epsilon_\tau)B_{S,\tau}^{-1}$,
$J_\tau u_k=i\,s_{T,\tau}u_ks_{S,\tau}=i\eta_ku_k$. So, with the notation of the proof of
Lemma~\ref{is-eigenvalue-conjugate-theta-facts}, the creation operator of $k$ is
$a_k=r_\tau\big(\partial_\tau(u_k)-i\,\partial_\tau(i\eta_ku_k)\big)$ for some $r_\tau>0$, and
$b_k=\partial_\tau(u_k)+i\,\partial_\tau(i\eta_ku_k)$ is an annihilation operator, so $b_kf=0$. By the orthonormality of
the bases, $B_{T,\tau}^*G_T(\tau)B_{T,\tau}=\operatorname{diag}(\epsilon'_\tau)$ and $B_{S,\tau}^{-1}G_S(\tau)^{-1}
B_{S,\tau}^{-*}=\operatorname{diag}(\epsilon_\tau)$, so for $c,c'\in\mathbf C$
\[
 \omega_{\tau\tau'}(cu_k,c'u_l)=\begin{cases}2\operatorname{Re}\big(\tau(\xi_0)\,c\bar c'\eta_k\big)&k=l,\\0&k\ne l.\end{cases}
\]
By (2), with $\tau(\xi_0)$ purely imaginary, this gives $a_ka_l=a_la_k$, and $b_ka_l=a_lb_k+\delta_{kl}c_k$ with
$c_k=-8\pi r_\tau\operatorname{Im}\tau(\xi_0)$, which is nonzero as $\tau(\xi_0)\ne0$: for $k=l$ the terms of the
commutator are $\omega(u_k,u_k)=\omega(i\eta_ku_k,i\eta_ku_k)=0$ and $\omega(i\eta_ku_k,u_k)=-\omega(u_k,i\eta_ku_k)=
-2\operatorname{Im}\tau(\xi_0)$. As $f\otimes_{s,B}P=P(a)f$, (1) gives the claim.
\end{proof}
% (1) is injective_of_canonicalCommutation with injective_of_pderiv (PureTensorOInjective.lean, 274 lines). (2) is
% heisenberg_derivative_commutator with exists_heisenberg_operators, hD_eq_fderiv and omO_self, on the smooth functions
% Sm (theta functions are smooth by theta_sm); omega is omO, of which omO_single is the value on one place. (3) is
% pureTensorO_injective (100 lines): hJ (cplxO_unitMap), hom (omO_single, trace_unitMap), hcomm, ccr_of_comm, the
% constants fockScale^-1 (-8 pi Im tau(xi0)) (hsc, him from realPlace_imaginaryGenerator_re_eq_zero and
% embedding_xi_zero_ne_zero) and the vacuum condition hvac.

\begin{lemma}\label{is-eigenvalue-conjugate-frame-extremal-compact}
\begin{enumerate}
\item Let $D$, $\psi$, $\theta$, $S$ and $T$ be as in Lemma~\ref{is-eigenvalue-conjugate-theta-facts}, $\omega$ a Weil
  datum for $(S,T)$ with characters $(\chi_S,\chi_T)$, and $M$ a model of discrete series for $U(T)$ with frames $f_\rho$,
  signature matrices $J_\rho$ and Cartan involutions $s_\rho=f_\rho^{-1}J_\rho f_\rho$, whose signs at $\rho$ are $1$ at
  the first $p_\rho$ indices and $-1$ at the other $q_\rho$. Let $B_S$ be bases of $S$ with
  $B_{S,\rho}^*G_S(\rho)B_{S,\rho}=1$, and write $f\otimes_{f^{-1}}P$ for $f\otimes P$ formed with the Cartan involutions
  $1$ and $s_\rho$ and the bases $B_S$ and $f_\rho^{-1}$. Let $n=\operatorname{rk}S$, and let $\chi_S$ and $\chi_T$ have
  archimedean exponents $e_\rho$ and $e'_\rho$ at every $\rho$ with $x_{1,\rho}=e_\rho+n/2$,
  $x_{2,\rho}=e_\rho-n/2$ and $(p_\rho-q_\rho)/2+e'_\rho$ integers. Let $f$ be a theta function annihilated by the
  annihilation operators for $1$ and $s_\rho$ at every $\rho$. Let $\rho_0\in\mathcal T$, $a\in\mathbf
  R^{\operatorname{rk}T}$, $p=p_{\rho_0}$, and $h\in U(T)(\mathbf A)$ the element with finite component $1$, component
  $f_{\rho_0}^{-1}\operatorname{diag}(e^{ia_1},\dots,e^{ia_{\operatorname{rk}T}})f_{\rho_0}$ at $\rho_0$ and $1$ at the
  other embeddings. Then for every polynomial $P$ in the creation variables
  \[
   \omega(1,h)(f\otimes_{f^{-1}}P)=e^{i(x_{1,\rho_0}\sum_{j\le p}a_j+x_{2,\rho_0}\sum_{j>p}a_j)}\,f\otimes_{f^{-1}}P_a,
  \]
  where $P_a$ is $P$ with each variable $(\rho_0,j,k)$ replaced by $e^{ia_j}(\rho_0,j,k)$ if $j\le p$ and by
  $e^{-ia_j}(\rho_0,j,k)$ if $j>p$.
\item Let $R$, $D$, $i$ and $\sigma$ be as in Lemma~\ref{is-eigenvalue-conjugate-level-data}, $\rho\in\mathcal T$,
  $\tau'$ the label of $\sigma$ at $\rho$, $(p,q)$ the signature of $W_i$ at $\rho$, $n=\operatorname{rk}V_i$, $e_i$ the
  $i$-th number $e$ of the Harish-Chandra parameter $\lambda(\tau')$ of $D$ (Definition~\ref{th-hc-parameter}),
  $x_1=n/2-e_i$ and $x_2=-n/2-e_i$. Then $p+q=\operatorname{rk}W_i$, the signs of the model of $W_i$ at $\sigma$ at
  $\rho$ are $1$ at the first $p$ indices and $-1$ at the others, the sequences $a^\pm=a^\pm_\rho(\sigma)$
  (Definition~\ref{is-eigenvalue-lift-data}) decrease, and the weight $\Lambda_\rho$ of the lowest $K$-type
  (Definition~\ref{is-eigenvalue-conjugate-k-types}(3)) is
  $(a^+_1+x_1,\dots,a^+_p+x_1,-a^-_q+x_2,\dots,-a^-_1+x_2)$.
\end{enumerate}
\end{lemma}
\leanref{Development/AbelianVarietyTheory/Proofs/IwasawaDecomposition/WeilMaximalCompactFockSubstitutionCharacterStatement.lean}{GlobalTheory.T.weil_maximalCompact_fock_substitution_character}
\leanref{Development/AbelianVarietyTheory/Proofs/DivisorProofs/CoefficientConjugateThetaForm.lean}{GlobalTheory.MovedWeights.weights_label}
% Self-contained computations. (1) is the case g = 1, h diagonal at one place, of
% T.weil_maximalCompact_fock_substitution_character (FockCompactAction.om_pureTensorO_eq_fact): for compact (g, h) the
% Weil action on f (x) P is f (x) fact(...) P, the character fchar (detP^x1 detM^x2 det(u)^x3) times the substitution
% fsub; for the diagonal element k = diag(e^{i a}) and u = 1 these are the factor and the substitution P_a above. (2) is
% MovedWeights.weights_label (lowK eps lam = wtK p (aPl) (aMi) x1 x2 with 0-based indices, the parities and the
% monotonicity of aPl and aMi), for the model at sigma, whose parameter at rho is the holomorphic parameter of the
% label (hH.par, hH.sig of HypConj).

\begin{lemma}\label{is-eigenvalue-conjugate-frame-extremal-line}
Let $R$, $D$, $F$, $i$ and $\sigma$ be as in Lemma~\ref{is-eigenvalue-conjugate-level-data}, let $\chi'$, $\eta'$ and
$\omega$ be as in (2) there, and let $\phi$ be a level vector of the trivial level as in (1) there. Let $M$ be the
model of $W_i$ at $\sigma$, with representations $\sigma_\rho$, frames $f_\rho$, signature matrices $J_\rho$, Cartan
involutions $s_\rho=f_\rho^{-1}J_\rho f_\rho$, lowest $K$-types $\kappa_\rho$, groups $K_\rho$ of unitary matrices
commuting with $J_\rho$, and extremal weights $w_\rho$. Write $\phi\otimes_{B'}Q$ for $\phi\otimes Q$ formed with the
Cartan involutions $1$ and $s_\rho$, the bases of $V_i$ of the realization and bases $B'_\rho$ of $W_i$
(Lemma~\ref{is-eigenvalue-conjugate-theta-facts}). For $k\in K_\rho$ let $k^\theta\in U(W_i)(\mathbf A)$ be the element
with finite component $1$, component $f_\rho^{-1}kf_\rho$ at $\rho$ and $1$ at the other embeddings.
\begin{enumerate}
\item If $v$ and $v'$ are extremal vectors of $M$, then for every $\rho$ there is $c_\rho\in\mathbf C$ with
  $|c_\rho|=1$ and $v'_\rho=c_\rho v_\rho$.
\item The map $Q\mapsto\phi\otimes_{f^{-1}}Q$ on the polynomials in the creation variables is injective, and
  $\Delta_i(\sigma)\ne0$.
\item For every $\rho$ and $a\in\mathbf R^r$, with $t_a=\operatorname{diag}(e^{ia_1},\dots,e^{ia_r})\in K_\rho$,
  \[
   \omega(1,t_a^\theta)\big(\phi\otimes_{f^{-1}}\Delta_i(\sigma)\big)=e^{i\sum_jw_{\rho,j}a_j}\,
   \phi\otimes_{f^{-1}}\Delta_i(\sigma).
  \]
\item For every $\rho$, the only subspaces of $\kappa_\rho$ stable under $\sigma_\rho(K_\rho)$ are $0$ and
  $\kappa_\rho$.
\end{enumerate}
\end{lemma}
\leanref{Development/AbelianVarietyTheory/Proofs/DivisorProofs/ExtremalVectorsHolomorphicModelUniqueUpToPhaseStatement.lean}{GlobalTheory.T.extremal_vectors_holomorphic_model_unique_up_to_phase}
\leanref{Development/AbelianVarietyTheory/Proofs/DivisorProofs/CoefMapNormEquivariantBasisOfFive.lean}{GlobalTheory.GalOne.VecConj.vec_conj_ofFive}
\leanref{Development/AbelianVarietyTheory/Proofs/DivisorProofs/ConjugateWeilRepresentationFifth.lean}{GlobalTheory.GalOne.IoExtConj.deltaConj_ne}
\leanref{Development/AbelianVarietyTheory/Proofs/DivisorProofs/CoefMapNormEquivariantBasisOfFive.lean}{GlobalTheory.GalOne.MinorWeight.deltaConj_extremal_weight_ofFive}
\leanref{Development/AbelianVarietyTheory/Proofs/DivisorProofs/ExtremalVectorsHolomorphicModelUniqueUpToPhase.lean}{GlobalTheory.HolLowOnce.extLine_of_holShape}
\leanref{Development/AbelianVarietyTheory/Proofs/DivisorProofs/ExtensionOfLowSmoothLines.lean}{GlobalTheory.AdTh.IoExtPf.model_irred}
\leanref{Development/AbelianVarietyTheory/Proofs/DivisorProofs/ExtensionOfLowSmoothLines.lean}{GlobalTheory.AdTh.IoExtPf.lowK_irred}
\leanref{Development/AbelianVarietyTheory/Proofs/DivisorProofs/ExtensionOfLowSmoothLines.lean}{GlobalTheory.AdTh.IoExtPf.minor_ne}
\leanref{Development/AbelianVarietyTheory/Proofs/DivisorProofs/ThetaVectorNeZero.lean}{GlobalTheory.AdTh.VecNe.thPVec_ne_zero}
\leanref{Development/AbelianVarietyTheory/Proofs/DivisorProofs/ThetaVectorTorusAction.lean}{GlobalTheory.AdTh.MinorWeight.fact_diagonal}
\begin{proof}
As $\sigma$ agrees with $\sigma_0$ on $K$, $M$ has holomorphic shape (proof of
Lemma~\ref{is-eigenvalue-conjugate-theta-facts}(6)): at every $\rho$ its signs are $1$ at the first $p_\rho$ indices and
$-1$ at the others, $(p_\rho,q_\rho)$ the signature of $W_i$. Let $x_{0,\rho}$ be its frames and $t(s)$ the elements of
Lemma~\ref{is-eigenvalue-conjugate-spectral-facts}(3), and $\pi_\rho(g)=\sigma_\rho(f_\rho gf_\rho^{-1})$ the action of
$U(G_{W_i}(\rho))$ on the space $H_\rho$ of the model, an irreducible unitary representation
(Definition~\ref{asm-dependence-rows-models}(1)).

(1) Fix $\rho$, and let $w'_c$ be the extremal weight $w_\rho$ at the index of the $c$-th column of $x_{0,\rho}$. These are
integers: for $u\ne0$ of weight $w_\rho$, the weight condition at $t=2\pi e_a$ gives $u=e^{2\pi iw_{\rho,a}}u$. By
Lemma~\ref{is-eigenvalue-conjugate-spectral-facts}(3)(d), $v_\rho$ is smooth, for $a\le c$ the complexified derived action
of $x_{0,\rho}E_{ac}x_{0,\rho}^{-1}$ kills $v_\rho$ if $a<c$ and multiplies it by $w'_a$ if $a=c$, and
$\pi_\rho(t(s)_\rho)$ multiplies $v_\rho$, and likewise $v'_\rho$, by $\prod_ae^{is_aw'_a}$. By
Lemma~\ref{is-eigenvalue-conjugate-dictionary-facts}(2), $x_{0,\rho}^*G_{W_i}(\rho)x_{0,\rho}$ is the diagonal matrix with
entries $-1$ at the first $q_\rho$ indices and $1$ at the others. Put $c_\rho=\langle v_\rho,v'_\rho\rangle$ and
$l=v'_\rho-c_\rho v_\rho$. Then $\langle l,v_\rho\rangle=\langle v'_\rho,v_\rho\rangle-\bar c_\rho=0$, as $\|v_\rho\|=1$,
and $\pi_\rho(t(s)_\rho)l=\prod_ae^{is_aw'_a}\,l$. So
Lemma~\ref{is-eigenvalue-conjugate-frame-extremal-orthogonal}(4), applied to $G=G_{W_i}(\rho)$, $q=q_\rho$,
$x_0=x_{0,\rho}$, $\pi=\pi_\rho$, $v=v_\rho$ and the weights $w'$, gives $l=0$. Hence $v'_\rho=c_\rho v_\rho$, and
$|c_\rho|=\|v'_\rho\|=1$.

(4) Fix $\rho$, and let $N\subset\kappa_\rho$ be stable under $\sigma_\rho(K_\rho)$. By
Definition~\ref{asm-dependence-rows-models}(1), $\sigma_\rho$ has a $K$-type $S$ of weight $\Lambda_\rho$. Every $K$-type
$S'$ of weight $\Lambda_\rho$ contains a unit vector of weight $w_\rho$: if $u\in S'$ is nonzero of weight $\Lambda_\rho$,
the permutation matrix $\varpi$ that reverses the indices of the negative block lies in $K_\rho$, and
$\|u\|^{-1}\sigma_\rho(\varpi)u\in S'$ is a unit vector of weight $w_\rho$. So $M$ has an extremal vector $e$. Let
$u\in S$ and $u'\in S'$ be unit vectors of weight $w_\rho$. The families obtained from $e$ by replacing $e_\rho$ by $u$
and by $u'$ are extremal vectors, so $u'=cu$ for some $c$, by (1). Hence $S'\cap S\ne0$; as $S'\cap S$ is stable under
$\sigma_\rho(K_\rho)$ and $S'$ is a $K$-type, $S'\subset S$. So $\kappa_\rho$, the sum of the $K$-types of weight
$\Lambda_\rho$, is $S$, and $N$ is $0$ or $\kappa_\rho$, as $S$ is a $K$-type.

(2) \emph{The injectivity.} By Lemma~\ref{is-eigenvalue-conjugate-level-data}(1) and its proof, $\phi$ is a theta
function annihilated by the annihilation operators for $1$ and $s_\rho$ at every $\rho$
(Lemma~\ref{is-eigenvalue-conjugate-theta-facts}(3)), and the bases of $V_i$ of the realization are orthonormal for its
Gram matrix, with all signs $1$, so that they are adapted to the Cartan involution $1$. The bases $f_\rho^{-1}$ are
orthonormal and adapted to $s_\rho$, as $f_\rho^*J_\rho f_\rho=G_{W_i}(\rho)$ and $s_\rho f_\rho^{-1}=f_\rho^{-1}J_\rho$.
Moreover $\phi\ne0$. Indeed, $\phi=T(\Phi)$ for a place model $T$ of its vacuum $f_0$ at the places of $P_g$, where
$\Phi$ is the level function of the adelic matrices $\beta_V$ and $\beta_W$ of $R$, $P_g$, $Q=\emptyset$ and the trivial
characters (Lemmas~\ref{is-eigenvalue-conjugate-level-data}(1) and~\ref{is-eigenvalue-conjugate-frame-vacuum}(2)), with
the idempotent $a$ of $P_g$. The function $\Phi$ is a Schwartz-Bruhat function on $a\mathbb W$, and it is not $0$: as
$r=\operatorname{rk}W_i\le\operatorname{rk}V_i$, there is $x\in a\mathbb W$ for which $y(x)$ is the matrix with
components $a$ times the matrix whose top $r\times r$ block is $1$ and whose other rows are $0$, and then $\Phi(x)=1$.
The kernel of $T$ on the Schwartz-Bruhat functions on $a\mathbb W$ is stable under the translations by $a\mathbb W$ and
the multiplications by the characters $t\mapsto\psi(\langle t,y\rangle)$, $y\in\bar a\mathbb W$; by the Stone-von
Neumann theorem \cite[Chapitre~2]{MVW87} these functions form an irreducible representation of the Heisenberg group, so
the kernel is $0$, as $T$ carries the indicator function of the lattice to $f_0$, which has norm $1$. Hence
$\phi=T(\Phi)\ne0$. Lemma~\ref{is-eigenvalue-conjugate-frame-extremal-fock}(3), applied to $S=V_i$, $T=W_i$, the Cartan
involutions $1$ and $s_\rho$, the bases of the realization and $f_\rho^{-1}$, and $f=\phi$, gives the injectivity.

\emph{The delta polynomial.} By its definition, $\Delta_i(\sigma)$ is the product over $\rho\in\mathcal T$ of the powers
$(M^+_{\rho,k})^{a^+_k-a^+_{k+1}}$, $1\le k\le p_\rho$, and $(M^-_{\rho,k})^{a^-_k-a^-_{k+1}}$, $1\le k\le q_\rho$, where
$a^\pm=a^\pm_\rho(\sigma)$, $a^\pm_{k+1}=0$ beyond the last index, $M^+_{\rho,k}$ is the minor of the $r\times n_V$
matrix $(X_{\rho,a,j})$ of the creation variables on its first $k$ rows and its first $k$ columns, and $M^-_{\rho,k}$
the minor on the rows $p_\rho+1,\dots,p_\rho+k$ and the last $k$ columns. These minors exist, as $p_\rho+q_\rho=r\le n_V$.
Each is a nonzero polynomial: if its rows begin at $a_0$ and its columns at $j_0$, its value at the point where
$X_{\rho,a_0+\alpha,j_0+\alpha}=1$ for $0\le\alpha<k$ and the other variables are $0$ is the determinant of the identity
matrix. The polynomial ring is a domain, so $\Delta_i(\sigma)\ne0$.

(3) Let $\rho$ and $a$ be as in (3), and $x_1$, $x_2$ as in
Lemma~\ref{is-eigenvalue-conjugate-frame-extremal-compact}(2) at $\rho$; by
Lemma~\ref{is-eigenvalue-conjugate-characters}(3) the archimedean exponent of $\chi'$ at $\rho$ is $-e_i$, and that of
$\chi'\eta'^{-1}$ is $-e_i-\frac12$, so the integrality conditions of
Lemma~\ref{is-eigenvalue-conjugate-frame-extremal-compact}(1) hold, $n_V$ being even and $p_\rho-q_\rho$ odd. That lemma,
applied to $f=\phi$, $P=\Delta=\Delta_i(\sigma)$, $\rho_0=\rho$ and $h=t_a^\theta$, gives
$\omega(1,t_a^\theta)(\phi\otimes_{f^{-1}}\Delta)=e^{i(x_1\sum_{j\le p}a_j+x_2\sum_{j>p}a_j)}\,\phi\otimes_{f^{-1}}\Delta_a$,
with $p=p_\rho$. In $\Delta_a$ the minor $M^+_{\rho,k}$, whose rows are $\le p$, is multiplied by
$e^{i(a_1+\dots+a_k)}$, the minor $M^-_{\rho,k}$ by $e^{-i(a_{p+1}+\dots+a_{p+k})}$, and the factors at the other
places do not change. With the exponents of $\Delta$, $e^{ia_j}$ occurs with the exponent
$\sum_{k\ge j}(a^+_k-a^+_{k+1})=a^+_j$ for $j\le p$, and $e^{-ia_{p+l}}$ with the exponent $a^-_l$. So
$\Delta_a=e^{i(\sum_{j\le p}a^+_ja_j-\sum_{l\le q}a^-_la_{p+l})}\Delta$, and the total exponent of $e^{ia_j}$ is $a^+_j+x_1$
for $j\le p$ and $-a^-_{j-p}+x_2$ for $j>p$. By
Lemma~\ref{is-eigenvalue-conjugate-frame-extremal-compact}(2), $\Lambda_\rho$ is
$(a^+_1+x_1,\dots,a^+_p+x_1,-a^-_q+x_2,\dots,-a^-_1+x_2)$, and $w_\rho$, which takes the indices of the negative block
in the reverse order, has these exponents as its entries. This is (3).
\end{proof}
% (1) is T.extremal_vectors_holomorphic_model_unique_up_to_phase = HolLowOnce.extLine_of_holShape (48 lines), followed:
% wt_int for the integers w (WtIs M w), T.baseFrameBorel_acts_extremalVector_by_weight (Lemma
% is-eigenvalue-conjugate-spectral-facts(3)(d)), T.discreteSeriesModel_unitary_irreducible_orthonormal (unitary,
% irreducible, OrthFr x0Mat), torus_wt, c = <x, y>, horth, eq_zero_of_orth (Lemma -frame-extremal-orthogonal(4)); the
% holomorphic shape is holShape_of_overK. extWt M rho a = lowK (M.eps rho) (M.lam rho) (revNeg a). (4) is
% AdTh.IoExtPf.model_irred with lowK_irred and exists_ext_mem (pM (revNeg eps) in Kset, isWtOn_perm), with
% exists_isExtVec (under LowSm) for the family e; the Lean takes the K-type S from IsDS (M.ds rho). (2) is the
% injectivity clause of VecConj.vec_conj_ofFive: FockVecInj.pureTensorO_injective (Lemma -frame-extremal-fock(3)) with
% isCartanInv_one_of_onb, isONBasis_onb, isFockVacAt_thPVec and AdTh.VecNe.thPVec_ne_zero (54 lines: the vacuum norm,
% AdjDatumNe.exists_datum_ne_zero with exists_half_frMapAdj and rectOne, isSBHalf_of_red, and
% PlaceModelInj.placeModel_ne_zero, which the Lean proves its own way in Chain.StepVac.placeModel_span_ne, where the book
% cites Stone-von Neumann); and deltaConj_ne = IoExtPf.delta_ne with minor_ne (the evaluation at the shifted identity)
% and the ranks (W_sigval). deltaConj i sigma = delta (pW) (qW) (aPl at lab sigma) (aMi at lab sigma), with minorPlus
% and minorMinus (LocalNewformFunction.lean). (3) is MinorWeight.deltaConj_extremal_weight_ofFive (60 lines):
% T.weil_maximalCompact_fock_substitution_character (Lemma -frame-extremal-compact(1)), fact_diagonal (through
% DeltaEig.fsub_delta: the minors under the diagonal substitution), MovedWeights.weights_label (Lemma
% -frame-extremal-compact(2)) and revNeg_shape.

\begin{lemma}\label{is-eigenvalue-conjugate-frame-rebase-substitution}
Let $D$, $\psi$ and $\theta$ be as in Lemma~\ref{is-eigenvalue-conjugate-theta-facts}, and let $V$ and $W$ be Hermitian
spaces of $D$ such that $1$ is a Cartan involution of $V$ at every $\rho$. Let $M$ be a model of discrete series for
$U(W)$ with frames $f_\rho$, signature matrices $J_\rho$, Cartan involutions $s_\rho=f_\rho^{-1}J_\rho f_\rho$ and groups
$K_\rho$ of unitary matrices commuting with $J_\rho$; let $B_V$ be bases of $V$ at the $\rho$, and write $f\otimes_{B'}P$
for $f\otimes P$ formed with the Cartan involutions $1$ and $s_\rho$ and the bases $B_V$ and $B'$. Let $k_\rho\in K_\rho$
for every $\rho$, and let $s_k$ be the substitution of the creation variables
\[
 (\rho,a,j)\longmapsto\sum_bc_{\rho,ba}\,(\rho,b,j),\qquad c_{\rho,ba}=(k_\rho)_{ba}\text{ if }(J_\rho)_{aa}=1,\quad
 c_{\rho,ba}=\overline{(k_\rho)_{ba}}\text{ if }(J_\rho)_{aa}=-1.
\]
Then $f\otimes_{(f_\rho^{-1}k_\rho)_\rho}P=f\otimes_{f^{-1}}s_k(P)$ for every theta function $f$ for $(V,W)$ and every
polynomial $P$ in the creation variables.
\end{lemma}
\leanref{Development/AbelianVarietyTheory/Proofs/IwasawaDecomposition/FockPolynomialVectorMaximalCompactBasisSubstitut.lean}{GlobalTheory.FockSubst.pureTensorO_subst}
% A computation in the Fock model (pureTensorO_subst, FockPolynomialVectorMaximalCompactBasisSubstitut.lean): the
% matrix of the variable (rho, a, j) for the bases f^-1 k is the sum over b of (k_rho)_{ba} times that of (rho, b, j)
% for f^-1 (unitMap_mul); the creation operator is additive in it (aPlus_dir_sum), and on the matrix of (rho, b, j) for
% f^-1, an eigenvector of the complex structure with the sign of (J_rho)_{bb} (cplxO_unitMap), multiplication by a
% scalar becomes multiplication by it or by its conjugate (aPlus_dir_smul); k_rho commutes with J_rho, so
% (k_rho)_{ba} = 0 unless the signs at a and b agree; the creation operators commute (aPlus_comm), and induction on P.

\begin{lemma}\label{is-eigenvalue-conjugate-frame-rebase}
Let $R$, $D$, $F$, $i$ and $\sigma$ be as in Lemma~\ref{is-eigenvalue-conjugate-level-data}, let $\chi'$,
$\eta'$ and $\omega$ be as in (2) there, and let $\phi$ be a level vector of the trivial level with vacuum
$f_0$ as in (1) there. Let $s_\rho$ and $f_\rho$ be the Cartan involution and the frame of the model of
$W_i$ at $\sigma$ at $\rho\in\mathcal T$, and write $\phi\otimes_{B'}Q$ for $\phi\otimes Q$ formed with the
Cartan involutions $1$ and $s_\rho$, the bases of $V_i$ of the realization and the bases $B'_\rho$ of $W_i$
(Lemma~\ref{is-eigenvalue-conjugate-theta-facts}). Let $B$ be bases as in
Definition~\ref{is-eigenvalue-lift-data}(5), $W_c$ a coefficient datum and $\mathrm{io}$ a map as in
Definition~\ref{is-eigenvalue-lift-data}(7) for $\omega$, $\phi$, $B$ and $\Delta_i(\sigma)$. Then there is a
map $\mathrm{io}'$ as in Definition~\ref{is-eigenvalue-lift-data}(7) for $\omega$, $\phi$, the bases
$(f_\rho^{-1})_\rho$ and $\Delta_i(\sigma)$, with
$\phi\otimes_B\mathrm{io}(v\otimes w)=\phi\otimes_{f^{-1}}\mathrm{io}'(v\otimes w)$ for every tensor $v$ of
lowest $K$-type vectors and every $w\in W_c$.
\end{lemma}
\leanref{Development/AbelianVarietyTheory/Proofs/DivisorProofs/CoefMapNormEquivariantBasisOfFive.lean}{GlobalTheory.GalOne.Rebase.rebase_ofFive}
\begin{proof}
Use the notation of Lemma~\ref{is-eigenvalue-conjugate-frame-extremal-line}, and write $\Delta=\Delta_i(\sigma)$. A map
as in Definition~\ref{is-eigenvalue-lift-data}(7) for $\omega$, $\phi$, bases $B'$ and $\Delta$ is an io map for
$(\omega,\phi,B',\Delta)$ and $W_c$ in the sense of Lemma~\ref{is-eigenvalue-conjugate-pure-tensors}, with the
properties (i) to (iii) there.

\emph{The bases.} For every $\rho$, $B_\rho$ and $f_\rho^{-1}$ are adapted to $s_\rho$, orthonormal for
$G_{W_i}(\rho)$ with the same signs: $B_\rho$ by Definition~\ref{is-eigenvalue-lift-data}(5), and $f_\rho^{-1}$ as
$f_\rho^*J_\rho f_\rho=G_{W_i}(\rho)$ and $s_\rho f_\rho^{-1}=f_\rho^{-1}J_\rho$. So the computation of Step~2 of the
proof of Lemma~\ref{is-eigenvalue-conjugate-pure-tensors}, for the adapted families $f^{-1}$ and $B$, shows that
$k_\rho=f_\rho B_\rho$ lies in $K_\rho$; and $B_\rho=f_\rho^{-1}k_\rho$.

\emph{The substitution.} The level vector $\phi$ is a theta function, and $1$ is the Cartan involution of $V_i$ at
every $\rho$ (Lemma~\ref{is-eigenvalue-conjugate-level-data}(1) and its proof). So
Lemma~\ref{is-eigenvalue-conjugate-frame-rebase-substitution}, with $k=(k_\rho)$, gives $\phi\otimes_BP=\phi\otimes_{f^{-1}}
s_k(P)$ for every polynomial $P$. Put $\mathrm{io}'=s_k\circ\mathrm{io}$, a linear map. Then
$\phi\otimes_B\mathrm{io}(v\otimes w)=\phi\otimes_{f^{-1}}\mathrm{io}'(v\otimes w)$ for all $v$ and $w$, which is the
identity of the lemma; as the two sides are the same functions, the properties (i) and (ii) of $\mathrm{io}'$ for
$f^{-1}$ are those of $\mathrm{io}$ for $B$.

\emph{The unit vectors.} By (iii) for $\mathrm{io}$ there are a tensor $v^0$ of unit vectors $v^0_\rho\in\kappa_\rho$ and a
unit vector $w^0$ with $\mathrm{io}(v^0\otimes w^0)=\Delta$. Let $h$ be the product over $\rho$ of the elements
$(k_\rho^{-1})^\theta$; it has finite component $1$ and component $h_\rho=f_\rho^{-1}B_\rho^{-1}$ at $\rho$, with
$h_\rho B_\rho=f_\rho^{-1}$, and $h_\rho$ commutes with $s_\rho$, as $k_\rho$ commutes with $J_\rho$. As in Step~2 of the
proof of Lemma~\ref{is-eigenvalue-conjugate-pure-tensors}, Lemma~\ref{is-eigenvalue-conjugate-theta-facts}(5), for the level
vector $\phi$ (Lemma~\ref{is-eigenvalue-conjugate-theta-facts}(3)), gives $\omega(1,h)\phi=a\phi$, and $|a|=1$ by
Lemma~\ref{is-eigenvalue-conjugate-invariant-form}(2). By Lemma~\ref{is-eigenvalue-conjugate-theta-facts}(4),
$\omega(1,h)(\phi\otimes_B\Delta)=a\,\phi\otimes_{f^{-1}}\Delta$. By (ii) for $\mathrm{io}$, applied at each $\rho$ in
turn, $\omega(1,h)(\phi\otimes_B\mathrm{io}(v^0\otimes w^0))=\phi\otimes_B\mathrm{io}(v^1\otimes w^0)$, where
$v^1_\rho=\sigma_\rho(k_\rho^{-1})v^0_\rho$ is a unit vector of $\kappa_\rho$. Hence
\[
 \phi\otimes_{f^{-1}}\mathrm{io}'(v^1\otimes a^{-1}w^0)=a^{-1}\,\phi\otimes_B\mathrm{io}(v^1\otimes w^0)
 =\phi\otimes_{f^{-1}}\Delta,
\]
and $\mathrm{io}'(v^1\otimes a^{-1}w^0)=\Delta$ by Lemma~\ref{is-eigenvalue-conjugate-frame-extremal-line}(2). As
$|a|=1$, $a^{-1}w^0$ is a unit vector; this is (iii) for $\mathrm{io}'$.
\end{proof}
% Rebase.rebase_ofFive (31 lines with its statement): VecConj.vec_conj_ofFive, HolLowOnce.lowSm_and_lowIrr and
% rebase_gen, under ConjW, IsONBasis for B and IsIoW (ptAt B f) (deltaConj) io; io' is IsIoW for ptAt (onb M) f.
% rebase_gen (ConjugateWeilRepresentationSixth.lean, 45 lines) is followed: ConjLiftUnique.exists_kset for B = onb k,
% FockSubst.pureTensorO_subst for s, io' = io.compr2 s with e0, the clauses (i) and (ii) through e0, and (iii) by
% ConjLiftUnique.exists_move (the move of Step 2 of Lemma is-eigenvalue-conjugate-pure-tensors, with om_vac and
% exists_unit_eigen) for v1 and l, w = l^-1 w0 and the injectivity hinj (Lemma -frame-extremal-line(2)).

\begin{lemma}\label{is-eigenvalue-conjugate-frame-extremal-unit}
Let $R$, $D$, $F$, $i$, $\sigma$, $\chi'$, $\eta'$, $\omega$, $\phi$ and $f_\rho$ be as in
Lemma~\ref{is-eigenvalue-conjugate-frame-rebase}, $W_c$ a coefficient datum and $\mathrm{io}$ a map as in
Definition~\ref{is-eigenvalue-lift-data}(7) for $\omega$, $\phi$, the bases $(f_\rho^{-1})_\rho$ and
$\Delta_i(\sigma)$. Then for every extremal vector $v$ of the model of $W_i$ at $\sigma$ there is $w\in W_c$
with $\|w\|=1$ and $\mathrm{io}(v\otimes w)=\Delta_i(\sigma)$.
\end{lemma}
\leanref{Development/AbelianVarietyTheory/Proofs/DivisorProofs/CoefMapNormEquivariantBasisOfFive.lean}{GlobalTheory.GalOne.IoExtConj.ioExt_conj_ofFive}
\begin{proof}
Use the notation of Lemma~\ref{is-eigenvalue-conjugate-frame-extremal-line}, and write $\Delta=\Delta_i(\sigma)$. The
map $\mathrm{io}$ has the properties (i) to (iii) of an io map for $(\omega,\phi,f^{-1},\Delta)$ and $W_c$
(Lemma~\ref{is-eigenvalue-conjugate-pure-tensors}). By (iii) there are a tensor $v^0=\otimes_\rho v^0_\rho$ of unit
vectors $v^0_\rho\in\kappa_\rho$ and a unit vector $w^0\in W_c$ with $\mathrm{io}(v^0\otimes w^0)=\Delta$. For
$y\in\bigotimes_\rho\kappa_\rho$ put $F(y)=\phi\otimes_{f^{-1}}\mathrm{io}(y\otimes w^0)$, a linear map; for
$u\in\kappa_\rho$ write $v^0[\rho\leftarrow u]$ for $v^0$ with $v^0_\rho$ replaced by $u$. Then
$F(v^0)=\phi\otimes_{f^{-1}}\Delta\ne0$ by Lemma~\ref{is-eigenvalue-conjugate-frame-extremal-line}(2).

\emph{One factor.} Fix $\rho$. The map $u\mapsto F(v^0[\rho\leftarrow u])$ on $\kappa_\rho$ is injective. Its kernel $N$ is
stable under $\sigma_\rho(K_\rho)$: if $u\in N$ and $k\in K_\rho$, then $F(v^0[\rho\leftarrow\sigma_\rho(k)u])=
\omega(1,k^\theta)F(v^0[\rho\leftarrow u])=0$ by (ii). By
Lemma~\ref{is-eigenvalue-conjugate-frame-extremal-line}(4), $N$ is $0$ or $\kappa_\rho$, and $N\ne\kappa_\rho$ as
$v^0_\rho\notin N$.

\emph{The vector $v^0$ is extremal.} Let $\rho$ and $a\in\mathbf R^r$, and put $e=e^{i\sum_jw_{\rho,j}a_j}$ and
$u=\sigma_\rho(t_a)v^0_\rho\in\kappa_\rho$. By (ii) and
Lemma~\ref{is-eigenvalue-conjugate-frame-extremal-line}(3),
\[
 F(v^0[\rho\leftarrow u])=\omega(1,t_a^\theta)F(v^0)=e\,F(v^0)=F(v^0[\rho\leftarrow ev^0_\rho]).
\]
By the injectivity of one factor, $\sigma_\rho(t_a)v^0_\rho=e\,v^0_\rho$. So $v^0_\rho$ has the weight $w_\rho$
(Definition~\ref{is-eigenvalue-conjugate-k-types}(1)), and $v^0$, a tensor of unit vectors of the lowest $K$-types of
extremal weight, is an extremal vector of $M$.

\emph{Conclusion.} Let $v$ be an extremal vector of $M$. By
Lemma~\ref{is-eigenvalue-conjugate-frame-extremal-line}(1) there are $c_\rho$ with $|c_\rho|=1$ and $v_\rho=c_\rho
v^0_\rho$. Put $c=\prod_\rho c_\rho$ and $w=c^{-1}w^0$. Then $\|w\|=1$, and $v=c\,v^0$ as tensors, so
$\mathrm{io}(v\otimes w)=\mathrm{io}(v^0\otimes w^0)=\Delta$.
\end{proof}
% IoExtConj.ioExt_conj_ofFive (31 lines with its statement): holShape_of_overK, VecConj.vec_conj_ofFive,
% exists_unit_of_isIoW with lowSm_and_lowIrr, T.extremal_vectors_holomorphic_model_unique_up_to_phase,
% deltaConj_ne and MinorWeight.deltaConj_extremal_weight_ofFive. Extremal is IsExtVec. exists_unit_of_isIoW
% (ConjugateWeilRepresentationFifth.lean) is followed: f = pt o io.flip w, hv0, hv (the weight through hK, hwt and
% AdTh.IoExtPf.inj_factor, the one-factor injectivity, whose irreducibility the Lean takes from model_irred, that is
% lowK_irred with ExtLine, Lemma is-eigenvalue-conjugate-frame-extremal-line(4)), then ExtLine for
% the phases a and w = (prod a)^-1 w.

\begin{lemma}\label{is-eigenvalue-conjugate-frame-phase-proportional}
Let $\Gamma$ be a group, $\varrho$ a representation of $\Gamma$ on a nonzero finite-dimensional complex vector space
$\mathcal M$ whose only $\Gamma$-stable subspaces are $0$ and $\mathcal M$, and $\Sigma\subset\Gamma$ a set that generates
$\Gamma$ with $\operatorname{tr}\varrho(\gamma)\ne0$ for $\gamma\in\Sigma$.
\begin{enumerate}
\item Let $T$ be an injective linear endomorphism of $\mathcal M$ and $\chi$ a function on $\Gamma$ with
  $T\varrho(\gamma)=\chi(\gamma)\varrho(\gamma)T$ for every $\gamma$. Then $T$ is a scalar.
\item Let $\mathcal N$ be a complex vector space and $A,B\colon\mathcal M\to\mathcal N$ linear maps such that for every
  $\gamma$ there are a linear endomorphism $s_\gamma$ of $\mathcal N$ and $c_\gamma\in\mathbf C$ with
  $A\varrho(\gamma)=s_\gamma A$ and $B\varrho(\gamma)=c_\gamma s_\gamma B$. Let $w,z\in\mathcal M$ with $A(w)=B(z)\ne0$.
  Then there is $a\in\mathbf C$ with $A=aB$ and $aw=z$.
\end{enumerate}
\end{lemma}
\leanref{Development/AbelianVarietyTheory/Proofs/DivisorProofs/ScalarOfIntertwinesUpToScalar.lean}{GlobalTheory.AdTh.scalar_of_intertwines_up_to_scalar}
\leanref{Development/AbelianVarietyTheory/Proofs/DivisorProofs/OscillatorModelExtensionFirst.lean}{GlobalTheory.AdTh.PhaseMap.proportional}
\begin{proof}
(1) As $\mathcal M$ is finite-dimensional, $T$ is invertible, and $T\varrho(\gamma)T^{-1}=\chi(\gamma)\varrho(\gamma)$. For
$\gamma\in\Sigma$, taking traces gives $(\chi(\gamma)-1)\operatorname{tr}\varrho(\gamma)=0$, so $\chi(\gamma)=1$ and $T$
commutes with $\varrho(\gamma)$. The elements with which $T$ commutes form a subgroup, so $T$ commutes with
$\varrho(\Gamma)$. Let $a$ be an eigenvalue of $T$; its eigenspace is a nonzero $\Gamma$-stable subspace, so it is
$\mathcal M$, and $T=a$.

(2) The kernel of $A$ is $\Gamma$-stable, as $A\varrho(\gamma)=s_\gamma A$, and it does not contain $w$, as $A(w)\ne0$; so
$A$ is injective. Every $c_\gamma$ is nonzero: otherwise
$B(z)=B(\varrho(\gamma)\varrho(\gamma^{-1})z)=c_\gamma s_\gamma B(\varrho(\gamma^{-1})z)=0$. So the kernel of $B$ is $\Gamma$-stable, and $B$ is injective in the same
way, and $s_\gamma B(x)=B(c_\gamma^{-1}\varrho(\gamma)x)$. The subspace of the $x$ with $A(x)$ in the image of $B$ is
$\Gamma$-stable by this formula, and it contains $w$, which is not in the kernel of $A$; so it is $\mathcal M$. Let
$T=B^{-1}A$, an injective endomorphism of $\mathcal M$. Then $BT\varrho(\gamma)=A\varrho(\gamma)=s_\gamma A=s_\gamma
BT=B(c_\gamma^{-1}\varrho(\gamma)T)$, so $T\varrho(\gamma)=c_\gamma^{-1}\varrho(\gamma)T$, and $T=a$ by (1). Hence $A=aB$, and
$B(aw)=A(w)=B(z)$ gives $aw=z$.
\end{proof}
% scalar_of_intertwines_up_to_scalar (44 lines) and PhaseMap.proportional (36 lines), followed: hAi, hBi, hc, hs, hr,
% T = B^-1 A (LinearEquiv.ofInjective), and the scalar lemma with chi = c^-1.

\begin{lemma}\label{is-eigenvalue-conjugate-frame-phase-eta-extension}
Let $F$ be a totally real number field and $E$ a totally complex quadratic extension of $F$, with nontrivial
automorphism $c$ over $F$, and let $(\varphi_i)_{i\in I}$ be a finite family of embeddings $E\to\mathbf C$ such that
the map sending $i$ to the archimedean place of $\varphi_i$ is a bijection onto the archimedean places of $E$. For a
finite place $v$ of $F$ and $a\in F^\times$ put $\eta_v(a)=(-1)^{\operatorname{ord}_v(a)}$ if $E$ has exactly one
place over $v$, and $\eta_v(a)=1$ otherwise; for a real place $v$ of $F$ let $\eta_v(a)=\pm1$ be the sign of $a$ at
$v$. Let $n_i$, $i\in I$, be odd integers, and suppose:
\begin{enumerate}
\item[(a)] $\operatorname{ord}_P(a)=\operatorname{ord}_v(a)$ for every $a\in F^\times$ and every finite place $P$ of $E$,
  with $v$ the place of $F$ under $P$;
\item[(b)] $\prod_v\eta_v(a)=1$ for every $a\in F^\times$, the product over all places $v$ of $F$;
\item[(c)] for every root of unity $\zeta\in E$ there is $\delta\in E^\times$ with $\delta=\zeta\,c(\delta)$ and
  $\prod_{i\in I}\big(\varphi_i(\delta)/|\varphi_i(\delta)|\big)^{n_i}=1$ whose valuation $\operatorname{ord}_P(\delta)$ is
  even at every finite place $P$ of $E$ that is the only place of $E$ over its place of $F$.
\end{enumerate}
For a place $w$ of $E$ and $x\in E_w^\times$ write $x_{(w)}$ for the idele with component $x$ at $w$ and $1$
elsewhere, and for $k\in E^\times$ write $k_\infty$ for the idele with archimedean components $k$ and finite
components $1$. Then there is a continuous character $\chi$ of the ideles of $E$ with values in the unit circle,
trivial on $E^\times$, such that:
\begin{enumerate}
\item $\chi(x_{(P)})=1$ for every finite place $P$ of $E$ and every unit $x$ of $E_P$;
\item $\prod_w\chi(a_{(w)})=\eta_v(a)$ for every place $v$ of $F$ and every $a\in F^\times$, the product over the
  places $w$ of $E$ over $v$;
\item $\chi(k_\infty)=\prod_{i\in I}\big(\varphi_i(k)/|\varphi_i(k)|\big)^{n_i}$ for every $k\in E^\times$.
\end{enumerate}
\end{lemma}
\leanref{Development/AbelianVarietyTheory/Proofs/IwasawaDecomposition/ExistsUnramifiedResEtaStatement.lean}{GlobalTheory.T.exists_unramified_resEta_alt}
\leanref{Development/AbelianVarietyTheory/Proofs/IwasawaDecomposition/ExistsUnramifiedResEta.lean}{GlobalTheory.Mem.EtaExt.exists_unramified_resEta}
\leanref{Development/AbelianVarietyTheory/Proofs/IwasawaDecomposition/ExistsUnramifiedResEta.lean}{GlobalTheory.Mem.EtaExt.compat}
\leanref{Development/AbelianVarietyTheory/Proofs/IwasawaDecomposition/ExistsUnitaryCharOfOpenSubgroupAndHom.lean}{GlobalTheory.exists_unitary_char_of_open_subgroup_and_hom}
\leanref{Development/AbelianVarietyTheory/Proofs/AdelicLatticeProofs/ExistsContinuousCircleCharExtension.lean}{GlobalTheory.exists_continuous_circle_char_extension}
\begin{proof}
Write $\mathbf I_E$ for the ideles of $E$. For $i\in I$ let $w_i$ be the place of $\varphi_i$ and $e_i\colon
E_{w_i}\to\mathbf C$ the continuous extension of $\varphi_i$. The character
$\psi(x)=\prod_i\big(e_i(x_{w_i})/|e_i(x_{w_i})|\big)^{n_i}$ of $\mathbf I_E$ is continuous, has values in the unit
circle and is trivial on the ideles with archimedean components $1$; and
$\psi(k)=\psi(k_\infty)=\prod_i(\varphi_i(k)/|\varphi_i(k)|)^{n_i}$ for $k\in E^\times$. Let $U\subset\mathbf I_E$ be the
open subgroup of the ideles that are units at every finite place. For each finite place $v$ of $F$ fix
$\varpi_v\in F^\times$ with $\operatorname{ord}_v(\varpi_v)=1$, and let $\beta_v$ be the product of the ideles
$(\varpi_v)_{(P)}$ over the places $P$ of $E$ over $v$. Let $H$ be the group of the pairs $(k,m)$ of an element
$k\in E^\times$ and a finitely supported function $m$ from the finite places of $F$ to $\mathbf Z$, and put
$\alpha(k,m)=k\prod_v\beta_v^{m(v)}$ and $a(k,m)=\prod_v\eta_v(\varpi_v)^{m(v)}$, homomorphisms from $H$ to
$\mathbf I_E$ and to $\{\pm1\}$.

\emph{Step 1: if $\alpha(k,m)\in U$, then $a(k,m)=\psi(\alpha(k,m))$.} By (a), $\operatorname{ord}_P(\varpi_v)=1$ for
every place $P$ of $E$ over $v$. So $\alpha(k,m)\in U$ says that $\operatorname{ord}_P(k)=-m(v)$ for every finite place
$P$ of $E$, with $v$ the place under $P$. As $\psi$ is trivial on the $\beta_v$,
$\psi(\alpha(k,m))=\prod_i(\varphi_i(k)/|\varphi_i(k)|)^{n_i}$.
\begin{enumerate}
\item[(i)] $\zeta=k/c(k)$ is a root of unity. As $c^2=1$, $\zeta c(\zeta)=1$. For a finite place $P$ of $E$ let $P'$ be
  the conjugate place, with $\operatorname{ord}_{P'}(c(x))=\operatorname{ord}_P(x)$ for $x\in E^\times$; it lies over the
  same place $v$ of $F$, so $\operatorname{ord}_{P'}(\zeta)=\operatorname{ord}_{P'}(k)-\operatorname{ord}_P(k)=-m(v)+m(v)=0$.
  So $\zeta$ is a unit at every finite place, hence an algebraic integer. For every embedding $\varphi$ of $E$,
  $\varphi\circ c$ is the complex conjugate of $\varphi$, as both are embeddings different from $\varphi$ that agree
  with it on the totally real field $F$; so $|\varphi(\zeta)|^2=\varphi(\zeta c(\zeta))=1$. By Kronecker's theorem
  \cite{Kro57}, $\zeta$ is a root of unity.
\item[(ii)] Let $\delta$ be as in (c) for $\zeta$. Then $\delta c(k)=k\,c(\delta)$, so $\delta^{-1}k$ is fixed by $c$,
  and $k=\delta a$ with $a\in F^\times$.
\item[(iii)] $\eta_v(\varpi_v)^{m(v)}=\eta_v(a)^{-1}$ for every finite place $v$ of $F$. Let $P$ be a place of $E$ over
  $v$. By (a), $-m(v)=\operatorname{ord}_P(\delta a)=\operatorname{ord}_P(\delta)+\operatorname{ord}_v(a)$. If $P$ is the
  only place over $v$, then $\operatorname{ord}_P(\delta)$ is even and $\eta_v(\varpi_v)=-1$, so
  $\eta_v(\varpi_v)^{m(v)}=(-1)^{\operatorname{ord}_v(a)}=\eta_v(a)^{-1}$. Otherwise both sides are $1$.
\item[(iv)] Each real place of $F$ lies under exactly one archimedean place of $E$: as $E$ is totally complex of
  degree $2[F:\mathbf Q]$, both fields have $[F:\mathbf Q]$ archimedean places, and every place of $F$ lies under
  one of $E$. So the map sending $i$ to the place $v_i$ of $F$ under $w_i$ is a bijection onto the real places of
  $F$, and $\eta_{v_i}(a)=(\varphi_i(a)/|\varphi_i(a)|)^{n_i}$, as $\varphi_i(a)$ is real and $n_i$ is odd. By (iii),
  (b) and (c),
  \[
   a(k,m)=\prod_{v\text{ finite}}\eta_v(a)^{-1}=\prod_{v\text{ real}}\eta_v(a)=\prod_i\Big(\frac{\varphi_i(a)}{|\varphi_i(a)|}\Big)^{n_i}
   =\prod_i\Big(\frac{\varphi_i(\delta a)}{|\varphi_i(\delta a)|}\Big)^{n_i}=\psi(\alpha(k,m)).
  \]
\end{enumerate}

\emph{Step 2: the extension.} Give $H$ the discrete topology, and let $\Gamma\subset\mathbf I_E\times H$ be the
subgroup of the elements $(\alpha(y),y)$, $y\in H$. The character $(x,y)\mapsto\psi(x)a(y)^{-1}$ of the open subgroup
$U\times H$ is continuous with values in the unit circle, and by Step~1 it is trivial on $(U\times H)\cap\Gamma$. So
it induces a character of the subgroup $(U\times H)\Gamma/\Gamma$ of $(\mathbf I_E\times H)/\Gamma$. The unit circle is
a divisible abelian group, hence an injective one \cite[Chapter~XX, \S4]{Lan02}, so this character extends to
$(\mathbf I_E\times H)/\Gamma$. This gives a character $\tilde\chi$ of $\mathbf I_E\times H$ with values in the unit
circle, trivial on $\Gamma$, that is $\psi(x)a(y)^{-1}$ on $U\times H$; it is continuous, as it is continuous on an
open subgroup. Put $\chi(x)=\tilde\chi(x,1)$. Then $\chi$ is continuous with values in the unit circle, $\chi=\psi$ on
$U$, and $\chi(\alpha(y))=\tilde\chi(\alpha(y),y)\,\tilde\chi(1,y)^{-1}=a(y)$ for $y\in H$.

\emph{Step 3: the properties.} For $k\in E^\times$, $\chi(k)=\chi(\alpha(k,0))=a(k,0)=1$. For (1), $x_{(P)}$ lies in $U$
and has archimedean components $1$, so $\chi(x_{(P)})=\psi(x_{(P)})=1$. For (3), $k_\infty\in U$, so
$\chi(k_\infty)=\psi(k_\infty)$. For (2), let $v$ be finite and $a\in F^\times$. By (a), the idele
$\beta_v^{-\operatorname{ord}_v(a)}\prod_Pa_{(P)}$, the product over the places $P$ of $E$ over $v$, lies in $U$ and has
archimedean components $1$; so
\[
 \prod_P\chi(a_{(P)})=\chi(\beta_v)^{\operatorname{ord}_v(a)}=\eta_v(\varpi_v)^{\operatorname{ord}_v(a)}=\eta_v(a),
\]
as $\beta_v=\alpha(1,m_v)$ with $m_v$ the indicator function of $v$. For a real place $v$ of $F$, let $w_i$ be the
place over it, by (iv) of Step~1. Then $a_{(w_i)}\in U$, and $\chi(a_{(w_i)})=\psi(a_{(w_i)})=
(\varphi_i(a)/|\varphi_i(a)|)^{n_i}=\eta_v(a)$.
\end{proof}
% Mem.EtaExt.exists_unramified_resEta (about 330 lines in its file), followed: psi from ModelChar.exists_modelChar with
% S empty (hpsil, hpsiinf, hpsid); beta = freeHom of locF v (unif v), b = freeHom of etaFin v (unif v); the extension is
% exists_unitary_char_of_open_subgroup_and_hom for U = unitIdl (UnitOpen.unit_open_and_congr_open), proved with
% exists_continuous_circle_char_extension on G x A (circle_hom_extend_of_injective, the divisibility of the circle, and
% continuity on the open subgroup). Step 1 is compat: div_conj_isRoot (normOne_unit_isRoot, Kronecker through
% NumberField.Embeddings.pow_eq_one_of_norm_eq_one; existsUnique_conjPl and under_eq_iff_conjPl for P'), hroot,
% exists_upF for a, and cmHeckeChar_ideleCompat_finprod (e1 for (iii), hprod, prod_etaInf with etaInf_eq,
% sign_zpow_odd and comap_bijective for (iv)). Step 3: hchiU and hpsil' for (1); hchibeta, valAt_locF, valAt_free,
% freeHom_single and etaFin_unif_zpow for (2) at finite v; etaInf_eq and infIdele_mem at real v; hpsiinf for (3).

\begin{theorem}[Global reciprocity {\cite[Chapter~VI, \S5]{Neu99}}]\label{is-eigenvalue-conjugate-frame-phase-eta-product}
Let $D$ be a definite-side datum with host field $E_h/F_h$, and let $\eta_v$, for the places $v$ of $F_h$, be the
signs of Lemma~\ref{is-eigenvalue-conjugate-frame-phase-eta-extension} for $E_h/F_h$. Then $\prod_v\eta_v(a)=1$ for
every $a\in F_h^\times$.
\end{theorem}
\leanref{Development/AbelianVarietyTheory/Proofs/IwasawaDecomposition/QuadraticCharProductFormulaCMHostStatement.lean}{GlobalTheory.T.quadraticChar_product_formula_CM_host}
\leanref{Development/AbelianVarietyTheory/Proofs/IwasawaDecomposition/QuadraticCharProductFormulaCMHost.lean}{GlobalTheory.EtaProduct.eta_product_formula}

\begin{theorem}[Global norm index {\cite[Chapter~VI, \S5]{Neu99}}]\label{is-eigenvalue-conjugate-frame-phase-eta-unique}
Let $E/F$ be a quadratic extension of number fields, and put $\mathbf A_E=\mathbf A_F\otimes_FE$. There is exactly one
continuous character of $\mathbf A_F^\times$ with values in the unit circle that is trivial on $F^\times$ and on the norms
of the elements of $\mathbf A_E^\times$ and is not trivial: the quadratic character $\eta$ of $E/F$.
\end{theorem}
\leanref{Development/AbelianVarietyTheory/Proofs/IwasawaDecomposition/MemNormGroupIffEverywhereNorm.lean}{GlobalTheory.IdeleNormChar.globalNormIndexTwo}
\leanref{Development/AbelianVarietyTheory/Proofs/IwasawaDecomposition/MemNormGroupIffEverywhereNorm.lean}{GlobalTheory.IdeleNormChar.existsUnique_isEta}

\begin{lemma}\label{is-eigenvalue-conjugate-frame-phase-eta-roots}
Let $D$ be a definite-side datum with host field $E_h/F_h$, set $\mathcal T$ of embeddings of $E_h$ and
$N=Q_Sqd$. For every root of unity $\zeta\in E_h$ there is $\delta\in E_h^\times$ with $\delta=\zeta\bar\delta$ and
$\prod_{\tau\in\mathcal T}\tau(\delta)/|\tau(\delta)|=1$ that is a unit at every prime of $E_h$ not dividing $N$.
\end{lemma}
\leanref{Development/AutomorphicArithmetic/WeightCongruences/ExistsUnitEqRootOfUnityMulConjPhaseOne.lean}{UpperTheory.exists_unit_eq_rootOfUnity_mul_conj_phase_one}
\leanref{Development/AutomorphicArithmetic/WeightCongruences/ProdCmTypePhaseOneAddRootOfUnity.lean}{UpperTheory.prod_cmType_phase_one_add_rootOfUnity}
\leanref{Development/AutomorphicArithmetic/WeightCongruences/RootOfUnityGaloisActionCMDatum.lean}{UpperTheory.rootOfUnity_galois_action_CMDatum}
\begin{proof}
Write $E=E_h$, and let $\varphi_0$, $\gamma$, $G_\Phi$, $\Psi$, $\iota$ and $\iota_d$ be the CM embedding of the host,
its isomorphism $\operatorname{Gal}(E/\mathbf Q)\to G_\Phi$ and the elements and subgroup of $G_\Phi$ of
Definitions~\ref{hf-frame}, \ref{hf-framed-host-field} and~\ref{hf-host}. The elements of $\mathcal T$ are the
embeddings $\varphi_g=\varphi_0\circ\gamma^{-1}(g)$, $g\in\Psi$, each once. By
Definition~\ref{hf-framed-host-field}(2), $\gamma^{-1}(g)$ fixes $\sqrt{-d}$ for $g\in\Psi$, so
$\varphi_g(\sqrt{-d})=\varphi_0(\sqrt{-d})=i\sqrt d$. Complex conjugation $x\mapsto\bar x$ is $\gamma^{-1}(\iota)$ and
is induced by complex conjugation through every embedding of $E$
(Proposition~\ref{hf-host-field-properties}(a)); so $\zeta\bar\zeta=1$, as $|\varphi_0(\zeta)|=1$. The roots of unity of
$E$ have order dividing $N$: $E$ embeds in $\mathbf Q(\zeta_N)$ (Definition~\ref{hf-framed-host-field}(3)), whose roots
of unity are the $\pm\zeta_N^a$, and $N$ is even, as $2^{e+2}$ divides $Q_S$.

\emph{The case $\zeta=-1$.} Take $\delta=\sqrt{-d}$. Then $\bar\delta=-\delta$, as $\epsilon_2(\iota)=-1$, so
$\zeta\bar\delta=\delta$. As $\delta^2=-d$ and $d$ divides $N$, $\delta$ is a unit at every prime not dividing $N$. Each
factor $\varphi_g(\delta)/|\varphi_g(\delta)|$ is $i$, and $|\Psi|=2n$ is divisible by $4$: the surjection $h$ of
Definition~\ref{hf-frame} from $A$ onto $H_e$, a group of order $2^e$, shows that $2^e$ divides $n$, and $e\ge2$. So
the product is $i^{2n}=1$.

\emph{The case $\zeta\ne-1$.} Take $\delta=1+\zeta\ne0$. As $\bar\zeta=\zeta^{-1}$,
$\zeta\bar\delta=\zeta+1=\delta$. Let $P$ be a prime of $E$ not dividing $N$. The algebraic integer $\delta$ is not in
$P$: otherwise $\zeta\equiv-1$ modulo $P$, while $\zeta$ and $-1$ are distinct roots of $X^N-1$, which has distinct
roots modulo $P$ as $N\notin P$. So $\delta$ is a unit at $P$. For the product, put $s=\iota\iota_d$, an element of
$\Psi$ with $s^2=1$ and $s\ne1$. The automorphism $\sigma=\gamma^{-1}(\iota_d)$ fixes $\zeta$. Indeed, let $P_d$ be a prime
of $E$ over $d$; $\sigma$ lies in its inertia group, which has order $2$
(Proposition~\ref{hf-host-field-properties}(d)), so $\sigma(\zeta)\equiv\zeta$ modulo $P_d$. The order $m$ of $\zeta$
is prime to $d$. Otherwise $E$ contains a primitive $d$-th root of unity $\xi$. Then $d=\prod_{i=1}^{d-1}(1-\xi^i)$, the
value at $1$ of the $d$-th cyclotomic polynomial. As $d\in P_d$, some factor $1-\xi^j$ lies in $P_d$; every $\xi^i$ is a
power of $\xi^j$, so $1-\xi^j$ divides every factor, and $d\in P_d^{d-1}$. So the ramification index of $P_d$ over $d$
is at least $d-1\ge6$. But $E/\mathbf Q$ is Galois, so this index is the order $2$ of the inertia group. Hence $m\notin
P_d$. Then distinct $m$-th roots of unity $x,y$ are distinct modulo $P_d$: from $x^m=y^m$ and $x\ne y$,
$\sum_{i<m}x^iy^{m-1-i}=0$, and if $x\equiv y$ modulo $P_d$ this sum is $\equiv my^{m-1}$, which is not in $P_d$ as $y$
is a unit. So $\sigma(\zeta)=\zeta$. Hence $\gamma^{-1}(s)(\zeta)=\overline{\sigma(\zeta)}=\zeta^{-1}$. For $g\in\Psi$ put
$w=\varphi_g(\zeta)$, so that $|w|=1$; then $\varphi_{gs}(\delta)=\varphi_g(1+\zeta^{-1})=1+\bar w$ and
$\varphi_g(\delta)\varphi_{gs}(\delta)=|1+w|^2>0$. So the factors at $g$ and $gs$ have product $1$. As $g\mapsto gs$
is an involution of $\Psi$ without fixed points, the product over $\Psi$ is $1$.
\end{proof}
% exists_unit_eq_rootOfUnity_mul_conj_phase_one (40 lines), followed: for z = -1, delta = sqd (conj_sqd, val_sqrt with
% d_dvd_condN, emb_sqd, four_dvd from two_le_e and two_pow_dvd_ord); otherwise delta = 1 + z (rootOfUnity_galois_action_CMDatum
% for conj z = z^-1 through Der.kc_conj_mul_self, val_one_add with isMu_iff, w_dvd_condN, two_dvd_condN and
% Der.kc_eq_of_pow_eq_one), and prod_cmType_phase_one_add_rootOfUnity (prod_pair over g and g * sg with R.sg_mul_sg
% and R.sg_ne_one, phase_pair). The fixing by iota_d is Der.kc_iod_fix (C.inertia_d, kc_no_root_d through
% Mu.mu_le_ram and Mu.mu_card_le_two, kc_eq_of_pow_eq_one); sg = cc * iod is cc_mul_iod. The Lean states the product
% over an index type with a bijection onto Psi; it is used with D.eTau, and tauE D tau = H.emb (eTau tau).

\begin{lemma}\label{is-eigenvalue-conjugate-frame-phase-eta-split}
Let $D$ be a definite-side datum with host field $E_h/F_h$ and $N=Q_Sqd$. Then no prime of $E_h$ over a prime
dividing $N$ is fixed by complex conjugation: every place of $F_h$ over a prime dividing $N$ splits in $E_h$.
\end{lemma}
\leanref{Formalism/SharpRestriction/RamificationFrobenius/TotientGaloisClasses/InertiaGaloisSymmetry/CycGrp.lean}{Formalism.HC.Der.decompConj}
\leanref{Formalism/SharpRestriction/RamificationFrobenius/TotientGaloisClasses/InertiaGaloisSymmetry/KcBase.lean}{Formalism.HC.Der.eta}
\begin{proof}
Write $E=E_h$, and let $e\ge2$, $A$, $h$, $\iota$, $\iota_d$, $\gamma$, $j$ and $\operatorname{art}$ be the exponent, group,
homomorphism and elements of the Galois datum of the host field and the maps of the host field
(Definitions~\ref{hf-frame} and~\ref{hf-framed-host-field}). Complex conjugation is $c=\gamma^{-1}(\iota)$
(Proposition~\ref{hf-host-field-properties}(a)). Let $\ell$ be a prime dividing $N$ and $P$ a prime of $E$ over $\ell$,
and suppose that $c(P)=P$. Two facts on the host field are used.
\begin{enumerate}
\item[(i)] $16$ divides $Q_S$, as $2^{e+2}$ does; and $Q_Sq$ is prime to $d$, as $\gcd(Q_S,qd)=1$, $q\equiv3\pmod8$ and
  $d\equiv7\pmod{16}$.
\item[(ii)] There is an integer $s$ with $Q_Sq\mid s^2+d$. Indeed, $-d\equiv1\pmod8$ is a square modulo every power of $2$;
  for an odd prime $p\mid Q_Sq$, $\left(\frac{-d}{p}\right)=1$ (Definition~\ref{hf-framed-host-field}(1)), so $-d$ is a
  square modulo every power of $p$ by Hensel's lemma; and the Chinese remainder theorem combines these. As $s^2+d$ is
  even and $d$ is odd, $s=2a+1$ is odd; and $s$ is prime to $Q_Sq$, by (i).
\end{enumerate}

\emph{The case $\ell\mid Q_Sq$.} Every $\sigma\in\operatorname{Gal}(E/\mathbf Q)$ with $\sigma(P)=P$ fixes $\sqrt{-d}$.
Indeed, suppose that $\sigma(\sqrt{-d})=-\sqrt{-d}$. The number $\omega=(1+\sqrt{-d})/2$ is an algebraic integer, a root
of $X^2-X+(1+d)/4$, as $d\equiv3\pmod4$, and $\sigma(\omega)=1-\omega$. Then
\[
 (\omega-a-1)(\omega+a)=\omega^2-\omega-a^2-a=-\frac{s^2+d}4 .
\]
The integer $(s^2+d)/4$ is divisible by $\ell$, as $s^2+d$ is divisible by $16$ and by $\ell$. So the product lies in
$P$, and one of its factors does. As $\sigma$ fixes $P$ and the integers, $\sigma(\omega-a-1)=-(\omega+a)$ and
$\sigma(\omega+a)=-(\omega-a-1)$, both factors lie in $P$, and so does their difference $2a+1=s$. Then $\ell$ divides
$s$ and $s^2+d$, hence $d$, which is prime to $Q_Sq$ by (i): a contradiction. Now $c$ fixes $P$ and moves $\sqrt{-d}$,
as $\epsilon_2(\iota)=-1$ (Definition~\ref{hf-framed-host-field}(2)). So this case does not occur.

\emph{The case $\ell=d$.} Let $u$ be an integer with $u\equiv d\pmod{Q_Sq}$ and $u\equiv1\pmod d$; it is prime to $N$.
Then $\phi=\gamma^{-1}(\operatorname{art}(u))$ is a Frobenius element at $P$. Indeed, it corresponds through $j$ to the
restriction of $\sigma_u\colon\zeta_N\mapsto\zeta_N^u$ (Definition~\ref{hf-framed-host-field}(3)), and $\sigma_u$ acts as
$y\mapsto y^d$ on the residue field of every prime $\mathfrak Q$ of $\mathbf Q(\zeta_N)$ over $d$: the number
$\xi=\zeta_N^{u-d}$ is a $d$-th root of unity, so $(\xi-1)^d\equiv\xi^d-1=0$ modulo $d$, hence $\xi\in1+\mathfrak Q$ and
$\sigma_u(\zeta_N)\equiv\zeta_N^d$ modulo $\mathfrak Q$. The decomposition group of $P$ is generated by its inertia
group and $\phi$ \cite[Chapter~I, Proposition~9.4]{Neu99}. So $c=\rho\phi^k$ with $\rho$ in the inertia group of $P$
and $k\ge0$, and $\gamma(\rho)\in\{1,\iota_d\}$ by Proposition~\ref{hf-host-field-properties}(d). Hence
$\iota=\gamma(\rho)\operatorname{art}(u)^k$.

The first sign of $\operatorname{art}(u)$ is $-1$. Indeed, let $w$ be an integer prime to $N$ with $w\equiv s
\pmod{Q_Sq}$. Then $-w^2\equiv d\equiv u\pmod{Q_Sq}$, so $\operatorname{art}(u)=\operatorname{art}(-1)\operatorname{art}(w)^2g$
with $g\in\{1,\iota_d\}$ by Definition~\ref{hf-framed-host-field}(4). As $\operatorname{art}(-1)=\iota$ and the first
signs of $\iota_d$ and of a square are $1$, the first sign of $\operatorname{art}(u)$ is that of $\iota$, which is $-1$.

Comparing first signs in $\iota=\gamma(\rho)\operatorname{art}(u)^k$ gives $(-1)^k=-1$, so $k$ is odd. The
$A$-components of $\iota$ and $\iota_d$ are $0$, so $k$ times the $A$-component of $\operatorname{art}(u)$ is $0$. Its
image under $h$ is the class of $u$ in $H_e=(\mathbf Z/2^{e+2}\mathbf Z)^\times/\{\pm1\}$
(Definition~\ref{hf-framed-host-field}(4)), so $u^k\equiv\pm1$ modulo $2^{e+2}$, hence modulo $16$. But $u\equiv
d\equiv7\pmod{16}$ by (i), and $7^k\equiv7\pmod{16}$ for odd $k$, as $7^2\equiv1$: a contradiction.

So $c(P)\ne P$ in both cases. As $\operatorname{Gal}(E/F_h)=\{1,c\}$, the places of $E$ over the place of $F_h$ under $P$
are $P$ and $c(P)$, and that place splits in $E$.
\end{proof}
% Der.decompConj (about 75 lines in CycGrp.lean), followed. For l | Q_S q it is decomp_typeI (about 100 lines): sq_mod
% (sq_two_pow, sq_odd_pow and the coprime induction) with split_QSq for s, the parity of s, Decomp.omega_integral for
% omega, hprod, hA and hB through Decomp.smul_mem_iff_of_stab, hls, hld and cop_QSq_d; then cc_not_mem_Psi. For l = d:
% exists_unit_crt for uF, exists_over and tauN_frob (zO_sub_mem_iff) for the Frobenius, Decomp.stab_eq_inertia_mul,
% inertia_d, ps_d (sq_mod, ps_congr, art_neg_one) for the first sign, art_h for the class of d in H_e, and the powers
% of 7 modulo 16 (the even k by the first sign, the odd k by hpm). Der.eta (third clause) reads the lemma as etaF = 1
% through Eta.ncard_eq_two_iff; Mem.EtaT.card_above_of_prime uses that clause. About 300 lines with the helpers.

\begin{lemma}\label{is-eigenvalue-conjugate-frame-phase-eta-reading}
Let $D$ be a definite-side datum with host field $E_h/F_h$ and set $\mathcal T$ of embeddings of $E_h$, let $\theta$
be an archimedean family whose member $\theta_\rho$ extends $\rho$ for every $\rho\in\mathcal T$, and let
$n\colon\mathcal T\to\mathbf Z$. Read $\mathbf A_{E_h}^\times$ as the restricted product of the $E_{h,w}^\times$, with
the notation $x_{(w)}$ and $k_\infty$ of Lemma~\ref{is-eigenvalue-conjugate-frame-phase-eta-extension}. Let $\chi$ be
a continuous character of $\mathbf A_{E_h}^\times$ with values in the unit circle, trivial on $E_h^\times$, such that
$\chi(x_{(P)})=1$ for every finite place $P$ of $E_h$ and every unit $x$ of $E_{h,P}$, and
$\chi(k_\infty)=\prod_{\rho\in\mathcal T}(\rho(k)/|\rho(k)|)^{n(\rho)}$ for every $k\in E_h^\times$. For every finite
place $P$ fix $\varpi_P\in E_h^\times$ with $\operatorname{ord}_P(\varpi_P)=1$.
\begin{enumerate}
\item If $x\in\mathbf A_{E_h}^\times$ has archimedean component $1$, then
  $\chi(x)=\prod_P\chi\big((\varpi_P)_{(P)}\big)^{\operatorname{ord}_P(x_P)}$, a product in which all but finitely many
  factors are $1$.
\item If $x\in\mathbf A_{E_h}^\times$ has finite component $1$, then
  $\chi(x)=\prod_{\rho\in\mathcal T}\big(\theta_\rho(x)/|\theta_\rho(x)|\big)^{n(\rho)}$.
\end{enumerate}
\end{lemma}
\leanref{Development/AbelianVarietyTheory/Proofs/IwasawaDecomposition/ExistsUnramifiedHeckeCharHalfExponentRestrictQua.lean}{GlobalTheory.Mem.EtaTilde.val_finite}
\leanref{Development/AbelianVarietyTheory/Proofs/IwasawaDecomposition/ExistsUnramifiedHeckeCharHalfExponentRestrictQua.lean}{GlobalTheory.Mem.EtaTilde.val_infinite}
\leanref{Development/AbelianVarietyTheory/Proofs/FoundationProofs/HeckeCharEqProdLocalComponentsApproximating.lean}{GlobalTheory.CharLocProd.char_eq_prod_locK}
\begin{proof}
(1) As $\chi$ is continuous and a small neighbourhood of $1$ in the unit circle contains no subgroup other than $1$,
there is an integer $M\ge1$ such that $\chi(y)=1$ for every idele $y$ with archimedean component $1$ that at every
finite place $P$ is a unit with $y_P-1\in M\mathcal O_{E_h,P}$. Let $x$ have archimedean component $1$, and let $S$ be
the finite set of the finite places $P$ at which $x_P$ is not a unit or $M$ is not a unit. For $P\in S$ choose
$k_P\in E_h^\times$ with $x_P-k_P\in Mk_P\mathcal O_{E_h,P}$, as $E_h$ is dense in $E_{h,P}$; then
$\operatorname{ord}_P(k_P)=\operatorname{ord}_P(x_P)$. The idele $y=x\prod_{P\in S}(k_P)_{(P)}^{-1}$ has archimedean
component $1$ and is a unit congruent to $1$ modulo $M$ at every finite place, so $\chi(y)=1$. As
$k_P\varpi_P^{-\operatorname{ord}_P(k_P)}$ is a unit at $P$,
\[
 \chi(x)=\prod_{P\in S}\chi\big((k_P)_{(P)}\big)=\prod_{P\in S}\chi\big((\varpi_P)_{(P)}\big)^{\operatorname{ord}_P(x_P)},
\]
and $\operatorname{ord}_P(x_P)=0$ for $P\notin S$.

(2) As in the proof of Lemma~\ref{is-eigenvalue-conjugate-frame-glue}, $z\mapsto(\theta_\rho(z))_\rho$ is a bijective
ring homomorphism from $\mathbf A_{F_h,\infty}\otimes_{F_h}E_h$ onto $\mathbf C^{\mathcal T}$, and the elements of
$\mathcal T$ represent the archimedean places of $E_h$, each once. For $z\in(\mathbf C^\times)^{\mathcal T}$ let $T(z)$ be
the idele with finite component $1$ and $\theta_\rho(T(z))=z_\rho$ for every $\rho$. Then $T$ is a continuous
homomorphism onto the ideles with finite component $1$, and $T((\rho(k))_\rho)=k_\infty$ for $k\in E_h^\times$. The
continuous characters $z\mapsto\chi(T(z))$ and $z\mapsto\prod_\rho(z_\rho/|z_\rho|)^{n(\rho)}$ of
$(\mathbf C^\times)^{\mathcal T}$ agree at the points $(\rho(k))_\rho$, $k\in E_h^\times$, by the hypothesis on
$\chi(k_\infty)$. These points are dense: the map $E_h\otimes_{\mathbf Q}\mathbf R\to\mathbf C^{\mathcal T}$ induced by the
$\rho$ is an isomorphism, so the image of $E_h$ is dense, and $(\mathbf C^\times)^{\mathcal T}$ is open in $\mathbf
C^{\mathcal T}$. So the two characters are equal, which is (2).
\end{proof}
% (1) is val_finite (25 lines), followed: CharLocProd.char_eq_prod_locK (CharFin.char_finUnit_finite_order for M,
% finiteIdele_exists_local_approx for S and k, eq_prod_of_components) and HeckeProduct.locK_unr for the unramified
% factors. (2) is val_infinite (70 lines), followed: exists_real_chart and RealChart.exists_chart for T, its
% continuity Tc, hT, HeckeFinTrivial.mem_closure_range and denseRange_units for the density, and the induction on the
% dense range. In the Lean the character on the ideles of D is a transport c of a character chi on the restricted
% product (ChrEq records only the finite components), so main reads c(k_infinity) as chi.inf k through the product
% rule HeckeProduct.productRule with S empty and val_finite; the book takes the two as one character.

\begin{lemma}\label{is-eigenvalue-conjugate-frame-phase-eta}
Let $D$ be a definite-side datum with host field $E_h/F_h$ and set $\mathcal T$ of embeddings of $E_h$. Then there is
a Hecke character $\tilde\eta$ of $E_h$ that is unramified at every finite place, restricts to the quadratic character
$\eta$ of $E_h/F_h$ on $\mathbf A_{F_h}^\times$ and has archimedean exponent $\frac12$ at every $\rho\in\mathcal T$.
\end{lemma}
\leanref{Development/AbelianVarietyTheory/Proofs/IwasawaDecomposition/ExistsUnramifiedHeckeCharHalfExponentRestrictQuaStatement.lean}{GlobalTheory.T.exists_unramifiedHeckeChar_halfExponent_restrict_quadratic}
\leanref{Development/AbelianVarietyTheory/Proofs/IwasawaDecomposition/ExistsUnramifiedHeckeCharHalfExponentRestrictQua.lean}{GlobalTheory.Mem.EtaTilde.exists_etaTilde}
\leanref{Development/AbelianVarietyTheory/Proofs/IwasawaDecomposition/HostEtaTildeCharacterExists.lean}{GlobalTheory.Mem.EtaT.etaT_exists}
\begin{proof}
Write $E=E_h$ and $F=F_h$, and use the notation $x_{(w)}$ and $k_\infty$ of
Lemma~\ref{is-eigenvalue-conjugate-frame-phase-eta-extension}. By Proposition~\ref{hf-host-field-properties}(a)
and~(e), $E$ is a CM field, a totally complex quadratic extension of its maximal totally real subfield $F$ with
nontrivial automorphism $z\mapsto\bar z$, and no finite place of $F$ ramifies in $E$; so
$\operatorname{ord}_P(a)=\operatorname{ord}_v(a)$ for every $a\in F^\times$ and every finite place $P$ of $E$ over a
place $v$ of $F$. By the proof of Lemma~\ref{is-eigenvalue-conjugate-frame-glue}, every archimedean place of $E$ is the
place of exactly one $\tau\in\mathcal T$. Fix an archimedean family $\theta$ whose member $\theta_\rho$ extends $\rho$
for every $\rho\in\mathcal T$, and for every finite place $P$ of $E$ an element $\varpi_P\in E^\times$ with
$\operatorname{ord}_P(\varpi_P)=1$.

\emph{The character.} Lemma~\ref{is-eigenvalue-conjugate-frame-phase-eta-extension} applies to $E/F$, the family
$(\tau)_{\tau\in\mathcal T}$ and $n_\tau=1$. Its hypothesis (b) is
Theorem~\ref{is-eigenvalue-conjugate-frame-phase-eta-product}. For (c), let $\zeta$ be a root of unity and $\delta$ as
in Lemma~\ref{is-eigenvalue-conjugate-frame-phase-eta-roots}. If $P$ is a finite place of $E$ that is the only place
over its place of $F$, then $P$ is fixed by complex conjugation, so it does not divide $N$
(Lemma~\ref{is-eigenvalue-conjugate-frame-phase-eta-split}); hence $\delta$ is a unit at $P$ and
$\operatorname{ord}_P(\delta)=0$. Let $\tilde\eta$ be the character that the lemma gives, read as a character of
$\mathbf A_E^\times=(\mathbf A_F\otimes_FE)^\times$, the restricted product of the $E_w^\times$. It is continuous, with
values in the unit circle and trivial on $E^\times$: a Hecke character of $E$. By (1) and (3) of the lemma, with
$n_\tau=1$, $\tilde\eta$ satisfies the hypotheses of Lemma~\ref{is-eigenvalue-conjugate-frame-phase-eta-reading} for
$n=1$.

\emph{The exponents.} Let $\rho\in\mathcal T$, and let $x$ be an idele with finite component $1$ and
$\theta_{\rho'}(x)=1$ for $\rho'\ne\rho$. By Lemma~\ref{is-eigenvalue-conjugate-frame-phase-eta-reading}(2),
$\tilde\eta(x)=\theta_\rho(x)/|\theta_\rho(x)|$. So $\tilde\eta$ has archimedean exponent $\frac12$ at $\rho$
(Definition~\ref{th-first-lift-matching}(2)).

\emph{The finite places.} Let $x$ be an idele with archimedean component $1$ that is a unit at every finite place.
Then $\operatorname{ord}_P(x_P)=0$ for every $P$, and $\tilde\eta(x)=1$ by
Lemma~\ref{is-eigenvalue-conjugate-frame-phase-eta-reading}(1). So $\tilde\eta$ is unramified at every finite place.

\emph{The restriction.} Let $\nu$ be the restriction of $\tilde\eta$ to $\mathbf A_F^\times$. It is continuous with
values in the unit circle, and trivial on $F^\times\subset E^\times$.
\begin{enumerate}
\item[(i)] $\nu$ is not trivial. Let $\tau_0\in\mathcal T$, let $v_0$ be the real place of $F$ under it, and let
  $u\in\mathbf A_F^\times$ have component $-1$ at $v_0$ and $1$ at every other place. In $\mathbf A_E^\times$, $u$ has
  finite component $1$, $\theta_{\tau_0}(u)=-1$ and $\theta_\tau(u)=1$ for $\tau\ne\tau_0$, as $\tau_0$ is the only
  element of $\mathcal T$ over $v_0$. By Lemma~\ref{is-eigenvalue-conjugate-frame-phase-eta-reading}(2), $\nu(u)=-1$.
\item[(ii)] $\nu$ is trivial on the norms. The norm to $\mathbf A_F^\times$ of $y\in\mathbf A_E^\times$ is $y\bar y$, read
  in $\mathbf A_E^\times$. Let $x\in\mathbf A_E^\times$, and write $x=ab$ with $a$ of finite component $1$ and $b$ of
  archimedean component $1$. The idele $a\bar a$ has finite component $1$ and
  $\theta_\tau(a\bar a)=|\theta_\tau(a)|^2>0$ for every $\tau$, as $\theta_\tau(\bar a)=\overline{\theta_\tau(a)}$; so
  $\nu(a\bar a)=1$ by Lemma~\ref{is-eigenvalue-conjugate-frame-phase-eta-reading}(2). The idele $y=b\bar b$ of $F$ has
  archimedean component $1$. Let $S$ be a finite set of finite places of $F$ outside which $y$ is a unit, and for
  $v\in S$ choose $k_v\in F^\times$ with $\operatorname{ord}_v(k_v)=\operatorname{ord}_v(y_v)$. For a place $P$ of $E$
  over $v$, $\operatorname{ord}_P(y_P)$ is $\operatorname{ord}_P(k_v)$ if $v\in S$ and $0$ otherwise. By
  Lemma~\ref{is-eigenvalue-conjugate-frame-phase-eta-reading}(1), applied to $y$ and to the ideles $(k_v)_{(P)}$, and
  Lemma~\ref{is-eigenvalue-conjugate-frame-phase-eta-extension}(2),
  \[
   \nu(y)=\prod_{v\in S}\prod_{P\mid v}\tilde\eta\big((\varpi_P)_{(P)}\big)^{\operatorname{ord}_P(k_v)}
   =\prod_{v\in S}\prod_{P\mid v}\tilde\eta\big((k_v)_{(P)}\big)=\prod_{v\in S}\eta_v(k_v).
  \]
  If $E$ has two places over $v$, then $\eta_v(k_v)=1$. If $P$ is the only place of $E$ over $v$, then complex
  conjugation fixes $P$, so $\operatorname{ord}_P(\bar b_P)=\operatorname{ord}_P(b_P)$, and
  $\operatorname{ord}_v(k_v)=\operatorname{ord}_P(b_P\bar b_P)=2\operatorname{ord}_P(b_P)$ is even; so $\eta_v(k_v)=1$.
  Hence $\nu(x\bar x)=\nu(a\bar a)\,\nu(b\bar b)=1$.
\end{enumerate}
By Theorem~\ref{is-eigenvalue-conjugate-frame-phase-eta-unique}, $\nu$ is the quadratic character $\eta$ of $E/F$.
\end{proof}
% Mem.EtaTilde.exists_etaTilde (25 lines), followed: Mem.EtaT.etaT_exists (Mem.HostCM.host_cm for the CM data, hunr
% and the bijection of embPsi onto the places; T.exists_unramified_resEta_alt with
% T.quadraticChar_product_formula_CM_host, n = 1 and the root condition from
% exists_unit_eq_rootOfUnity_mul_conj_phase_one, with exists_prime_mem and card_above_of_prime, that is Lemma
% -phase-eta-split, for the places with one place above), CharTransport.exists_chrEq for the transport c of chi to
% the units of A_F tensor E (the book takes them as one character), hc.1 for IsHeckeChar, val_infinite for HasInfExp
% 1/2, IdeleNormChar.existsUnique_isEta with isEta_res for ResEta D c 1, and val_finite with val_mul_inv for
% UnramEverywhere. isEta_res (90 lines) is (i) and (ii): continuity, hc.1.2 on F^x, the parts a and b
% (IdeleParts.congr_open_and_parts), idF_norm (Published.algebraMap_norm_tensor: the norm read in A_E is x conj x),
% e1 through val_infinite and comp_cA, e2 through val_finite and finprod_norm (finiteIdele_exists_local_approx for S
% and k, etaFin, and val_conj for the place fixed by conjugation), and the idele w = 1 + rho(single tau0 (-2)) of
% exists_real_chart for (i). The assembly cites this lemma, as Lemma
% asm-dependence-test-assembly-first-layer-ratio-inputs-weil-eta, for the Lean's WeilParts.exists_isEtaTilde_unramified.

\begin{lemma}\label{is-eigenvalue-conjugate-frame-phase-vacuum}
Let $D$ be a definite-side datum with host field $E_h/F_h$ and set $\mathcal T$ of embeddings of $E_h$, let $\psi$
be a character of $\mathbf A_{F_h}$, and let $\theta$ be an archimedean family whose member $\theta_\rho$ extends
$\rho$ for every $\rho\in\mathcal T$. Let $S$ and $T$ be Hermitian spaces of $D$ with correct signatures, of ranks
$n$ and $r$, such that $1$ is a Cartan involution of the Gram matrix of $S$ at every $\rho$, and let $s=(s_\rho)$ be
a family of Cartan involutions of $T$ at the $\rho$. Let $\omega$ be a Weil datum for $(S,T)$ with characters
$(\chi_S,\chi_T)$ (Definition~\ref{th-weil-datum}(3)), let $a_\rho$ and $b_\rho$ be archimedean exponents of
$\chi_S$ and $\chi_T$ at $\rho$, and let $x_{1,\rho}$, $x_{2,\rho}$ and $x_{3,\rho}$ be integers with
\[
 x_{1,\rho}=a_\rho+\tfrac n2,\qquad x_{2,\rho}=a_\rho-\tfrac n2,\qquad
 x_{3,\rho}=-\big(\tfrac12\operatorname{tr}s_\rho+b_\rho\big)
\]
for every $\rho$. Let $f$ be a theta function annihilated by the annihilation operators for the Cartan involutions
$1$ of $S$ and $s_\rho$ of $T$ at every $\rho$ (Lemma~\ref{is-eigenvalue-conjugate-theta-facts}), and let $g\in
U(S)(\mathbf A)$ and $h\in U(T)(\mathbf A)$ have finite components $1$, with images $g_\rho=\theta_\rho(g)$ and
$h_\rho=\theta_\rho(h)$ and $s_\rho h_\rho=h_\rho s_\rho$ for every $\rho$. Put
$d^\pm_\rho=\det\big(\frac12(1\pm s_\rho)h_\rho+\frac12(1\mp s_\rho)\big)$. Then
\[
 \omega(g,h)f=\prod_{\rho\in\mathcal T}(d^+_\rho)^{x_{1,\rho}}(d^-_\rho)^{x_{2,\rho}}\det(g_\rho)^{x_{3,\rho}}\,f.
\]
\end{lemma}
\leanref{Development/AbelianVarietyTheory/Proofs/DivisorProofs/VacuumCharacterPositiveDefiniteSourceStatement.lean}{GlobalTheory.T.vacuum_character_positive_definite_source}
\leanref{Development/AbelianVarietyTheory/Proofs/DivisorProofs/VacuumCharacterPositiveDefiniteSource.lean}{GlobalTheory.VacuumCharacter.vacuum_character}
\begin{proof}
Let $\tilde\eta$ be a Hecke character of $E_h$ as in Lemma~\ref{is-eigenvalue-conjugate-frame-phase-eta}. The Gram
matrices $G_S(\rho)$ and $G_T(\rho)$ are invertible, as they have Cartan involutions: if $s'$ is a Cartan involution
of $G$, then $Gs'$ is positive definite.

\emph{Determinant roots.} Each $g_\rho$ is unitary for the invertible matrix $G_S(\rho)$, so $|\det g_\rho|=1$. Let
$x$ be the idele with finite component $1$ whose archimedean component has image under $\theta_\rho$ a square root
$t_\rho$ of $\det g_\rho$ with $|t_\rho|=1$, for every $\rho$. Then $x=\det(g)\,\bar x$, as $t_\rho=\det(g_\rho)\bar
t_\rho$ and $g$ has finite component $1$; so $x$ is a determinant root of $g$
(Lemma~\ref{is-eigenvalue-conjugate-frame-det-roots}). For a character $\nu$ of $\mathbf A_{E_h}^\times$ with
integral archimedean exponents $c_\rho$, $\nu(x)=\prod_\rho t_\rho^{2c_\rho}=\prod_\rho\det(g_\rho)^{c_\rho}$. In the
same way $h$ has a determinant root $y$ with $\nu(y)=\prod_\rho\det(h_\rho)^{c_\rho}$ for every such $\nu$.

\emph{The determinants.} Fix $\rho$, put $P=\frac12(1+s_\rho)$, an idempotent that commutes with $h_\rho$, and let
$p_\rho$ be its rank, so that $\operatorname{tr}s_\rho=2p_\rho-r$. For the Cartan involutions $1$ of $S$ and $s_\rho$
of $T$, the determinant $\Delta_\rho$ of Lemma~\ref{is-eigenvalue-conjugate-frame-fock-character} is that of
$(Ph_\rho)\otimes A+(1-P)\otimes1$ with $A=(g_\rho^{-1})^{\mathsf T}$. This matrix is the product of
$(Ph_\rho+1-P)\otimes1$ and $P\otimes A+(1-P)\otimes1$, so
$\Delta_\rho=(d^+_\rho)^n\det(g_\rho)^{-p_\rho}$. Likewise $h_\rho=(Ph_\rho+1-P)\big((1-P)h_\rho+P\big)$, so $\det
h_\rho=d^+_\rho d^-_\rho$; in particular $d^\pm_\rho\ne0$.

\emph{The scalar.} Lemma~\ref{is-eigenvalue-conjugate-frame-fock-character}, applied to $\tilde\eta$, the Cartan
involutions $1$ of $S$ and $s$ of $T$, the elements $g$ and $h$ with the determinant roots $x$ and $y$, and $f$, gives
\[
 \omega(g,h)f=(\chi_T\tilde\eta^{-r})(x)^{-1}(\chi_S\tilde\eta^{-n})(y)\prod_\rho\Delta_\rho\,f.
\]
The character $\chi_T\tilde\eta^{-r}$ has archimedean exponent $b_\rho-\frac r2=-x_{3,\rho}-p_\rho$ at $\rho$, and
$\chi_S\tilde\eta^{-n}$ has $a_\rho-\frac n2=x_{2,\rho}$; both are integers. So
$(\chi_T\tilde\eta^{-r})(x)=\prod_\rho\det(g_\rho)^{-x_{3,\rho}-p_\rho}$ and
$(\chi_S\tilde\eta^{-n})(y)=\prod_\rho(d^+_\rho d^-_\rho)^{x_{2,\rho}}$, and the scalar is
\[
 \prod_\rho\det(g_\rho)^{x_{3,\rho}+p_\rho}(d^+_\rho d^-_\rho)^{x_{2,\rho}}(d^+_\rho)^n\det(g_\rho)^{-p_\rho}
 =\prod_\rho(d^+_\rho)^{x_{1,\rho}}(d^-_\rho)^{x_{2,\rho}}\det(g_\rho)^{x_{3,\rho}},
\]
as $x_{1,\rho}=x_{2,\rho}+n$.
\end{proof}
% vacuum_character (49 lines), followed: eta from T.exists_unramifiedHeckeChar_halfExponent_restrict_quadratic (Lemma
% -frame-phase-eta), hdet (a Cartan involution gives det G /= 0), RealMove.exists_root for x and y (with hxv and
% hyv), RealMove.detPlus_one for p (tr s = 2p - r, d1: detPlus 1 s g^-1 h = (d+)^n det(g^-1)^p, through
% det_idem_kron, and d2: det h = d+ d-), T.weilDatum_realPlaces_fockVacuum_character (Lemma -frame-fock-character),
% hb' (-x3 - p = b - r/2), hasInfExp_twist for the exponents, and the algebra alg with e (x1 = x2 + n).

\begin{lemma}\label{is-eigenvalue-conjugate-frame-phase-connected}
Let $D$ be a definite-side datum with host field $E_h/F_h$, and let $V$ be a Hermitian space of $D$ of rank $n$ with
correct signature, whose signature is $(n,0)$ at every embedding of $E_h$. Put $R=\mathbf
A_{F_h,\infty}\otimes_{F_h}E_h$, with the conjugation $z\mapsto\bar z$ of $E_h$, and let $U(V)_\infty$ be the group
of the matrices $g\in GL_n(R)$ with $\bar g^{\mathsf T}G_Vg=G_V$, with the topology of $M_n(R)$. Then $U(V)_\infty$ is
connected.
\end{lemma}
\leanref{Development/AbelianVarietyTheory/Proofs/DivisorProofs/RationalPointsGenerate.lean}{GlobalTheory.RatGen.preconnected_inf}
\leanref{Development/AbelianVarietyTheory/Proofs/DivisorProofs/RationalPointsGenerate.lean}{GlobalTheory.RatGen.preconnected_unitary}
\leanref{Development/AbelianVarietyTheory/Proofs/DivisorProofs/RationalPointsGenerate.lean}{GlobalTheory.RatGen.det_one_add_ne_zero}
\leanref{Development/AbelianVarietyTheory/Proofs/AdelicLatticeProofs/CayleyTransformUnitarySkewInvolutive.lean}{GlobalTheory.CayleyInv.cayley_transform}
\leanref{Development/AbelianVarietyTheory/Proofs/FoundationProofs/ExistsUnitaryClassRepTrivialAtCm.lean}{GlobalTheory.NormSum.RepAt.exists_unit_det}
\leanref{Development/AbelianVarietyTheory/Proofs/FoundationProofs/ExistsUnitaryClassRepTrivialAtCm.lean}{GlobalTheory.NormSum.RepAt.normOne_infinite}
\begin{proof}
Write $G=G_V$, $U=U(V)_\infty$ and $m^*=\bar m^{\mathsf T}$ for a matrix $m$ over $R$, and call $S\in M_n(R)$
\emph{skew} if $S^*G+GS=0$. Let $\theta$ be an archimedean family whose member $\theta_\rho$ extends $\rho$ for
every $\rho\in\mathcal T$. As in the proof of Lemma~\ref{is-eigenvalue-conjugate-frame-glue}, $z\mapsto
(\theta_\rho(z))_\rho$ is a bijective ring homomorphism from $R$ onto $\mathbf C^{\mathcal T}$, with
$\theta_\rho(\bar z)=\overline{\theta_\rho(z)}$. Matrix inversion is continuous on $GL_n(R)$, and $\mathbf
A_{F_h,\infty}$, a product of copies of $\mathbf R$ and $\mathbf C$, is connected.

\emph{Step 1: if $S$ is skew, $1+S$ is invertible.} Fix $\rho$, and put $G_\rho=\theta_\rho(G)=\rho(G)$ and
$S_\rho=\theta_\rho(S)$. Then $S_\rho^*G_\rho+G_\rho S_\rho=0$, and $G_\rho$ is positive definite, as $V$ has correct
signature $(n,0)$ at $\rho$. If $(1+S_\rho)v=0$, then $S_\rho v=-v$ and
$0=v^*(S_\rho^*G_\rho+G_\rho S_\rho)v=-2v^*G_\rho v$, so $v=0$. Hence $\theta_\rho(\det(1+S))\ne0$ for every $\rho$,
and $\det(1+S)$ is a unit of $R$, by the bijectivity.

\emph{Step 2: the Cayley transform.} Let $k\in U$ with $1+k$ invertible, and put $S=(1-k)(1+k)^{-1}$. Then
$(1+k)^*(S^*G+GS)(1+k)=(1-k)^*G(1+k)+(1+k)^*G(1-k)=2(G-k^*Gk)=0$, so $S$ is skew; and $1+S=2(1+k)^{-1}$ and
$1-S=2k(1+k)^{-1}$, so $(1-S)(1+S)^{-1}=k$. Conversely, let $S$ be skew with $1+S$ invertible, and put
$k=(1-S)(1+S)^{-1}$. Then $(1+S)^*(k^*Gk-G)(1+S)=(1-S)^*G(1-S)-(1+S)^*G(1+S)=-2(S^*G+GS)=0$, so $k^*Gk=G$; as $G$ is
invertible, so is $k$, and $k\in U$.

\emph{Step 3: every $k\in U$ with $1+k$ invertible lies in the connected component $U^0$ of $1$.} Let $S$ be as in
Step~2. For $a\in\mathbf A_{F_h,\infty}$, $aS$ is skew, as $\bar a=a$, so $1+aS$ is invertible by Step~1, and
$c(a)=(1-aS)(1+aS)^{-1}$ lies in $U$ by Step~2. The map $a\mapsto c(a)$ is continuous, $c(0)=1$, and $c(1)=k$ by
Step~2. As $\mathbf A_{F_h,\infty}$ is connected, $k\in U^0$.

\emph{Step 4.} Let $g\in U$. As $E_h\ne F_h$, there is $z\in E_h$ with $\bar z\ne z$; put $d=\bar z-z$, so that $\bar
d=-d\ne0$. For $t\in\mathbf N$ the elements $1\pm td$ are conjugate with sum $2$, so both are nonzero, and
$\lambda_t=(1+td)/(1-td)$ satisfies $\lambda_t\bar\lambda_t=1$ and $1+\lambda_t=2/(1-td)\ne0$; these elements are
pairwise distinct. For every $\rho$, $\theta_\rho(\det(1+\lambda g))=\det(1+\rho(\lambda)g_\rho)$, with
$g_\rho=\theta_\rho(g)$, is the value at $\rho(\lambda)$ of the reversed characteristic polynomial of $-g_\rho$, which
has constant term $1$; so it vanishes for only finitely many $\lambda\in E_h$. Choose $t$ such that
$\lambda=\lambda_t$ avoids these finitely many values for every $\rho$; then $\det(1+\lambda g)$ is a unit of $R$. The
scalar matrix $\lambda\cdot1$ lies in $U$, and $1+\lambda\cdot1=(1+\lambda)\cdot1$ is invertible, so $\lambda\cdot1\in
U^0$ by Step~3. Likewise $\lambda g\in U$ and $1+\lambda g$ is invertible, so $\lambda g\in U^0$. Hence
$g=(\lambda\cdot1)^{-1}(\lambda g)\in U^0$, and $U$ is connected.
\end{proof}
% preconnected_inf (49 lines) applies preconnected_unitary (60 lines) to A = A_F,inf (preconnected_infAdele), the
% Artinian ring R (units_of_free for the continuity of inversion), p(a) = a . 1, J = gram and hdef from
% map_det_one_add_ne_zero and det_one_add_ne_zero (Step 1) through infiniteAdele_bijective (hunit). preconnected_unitary
% is followed: CayleyInv.cayley_transform (c1, c2, c3: Step 2), step (Step 3, the path a -> (1 - p(a) S)(1 + p(a) S)^-1
% with isPreconnected_range), and all (Step 4: NormSum.RepAt.exists_unit_det, through the maximal ideals of R and the
% reversed characteristic polynomial, with the infinite set of normOne_infinite, exists_conj_ne for z).

\begin{lemma}\label{is-eigenvalue-conjugate-frame-phase}
In Situation~\ref{is-eigenvalue-conjugate-frame-core}, let $\chi'$, $\eta'$, $\omega$, $\phi$ with vacuum
$f_0$, $s_\rho$ and $f_\rho$ be as in Lemma~\ref{is-eigenvalue-conjugate-frame-rebase}, and $W_c$ and $c_1$ as
in (2) of
Definition~\ref{is-eigenvalue-conjugate-datum}. Write $\phi\otimes_{1,f^{-1}}Q$ for $\phi\otimes Q$ formed with
the Cartan involutions $1$ and $s_\rho$ and the bases $1$ of $V_i$ and $f_\rho^{-1}$ of $W_i$. A \emph{letter
assignment} is a map $\mathrm{lt}$ from the letter variables $((\tau',i'),k)$ of the datum of the core at level
$i+1$, $1\le k\le n$, to the creation variables $(\rho,a,k)$ with the same $k$, such that either
$s\tau'=\rho$ and $a=i'$, or $s\tau'=\bar\rho$ and $a=p+i'$, for the signature $(p,q)$ of that datum; for a
polynomial $P$ in the letter variables, $\mathrm{lt}(P)$ is $P$ with each variable $\lambda$ replaced by
$\mathrm{lt}(\lambda)$. Let $\Lambda$ be a linear map from $W_c$ to the polynomials in the creation variables
such that $\omega(u,1)\big(\phi\otimes_{f^{-1}}\Lambda(w)\big)=\phi\otimes_{f^{-1}}\Lambda(uw)$ for every
$u\in U(V_i)(\mathbf A)$ with finite component $1$ and every $w$, and $\Lambda(w_0)=\Delta_i(\sigma)$ for some
$w_0\in W_c$ of norm $1$. Then there is $\epsilon\in\mathbf C$ with $|\epsilon|=1$ such that
\[
 \phi\otimes_{f^{-1}}\Lambda(c_1(x''))=\epsilon\cdot\phi\otimes_{1,f^{-1}}\mathrm{lt}\big(\sigma'(\mathrm{io}_O(\iota(x'')))\big)
\]
for every letter assignment $\mathrm{lt}$ and every rational coefficient vector $x''$ of $V_i$ over $L$, where
$\mathrm{io}_O(x')$ is the polynomial in the letter variables over $L'$ that the one-map data at $i$ attach to
a coefficient vector $x'$ of $V_i$ over $L'$, and $\sigma'$ acts on its coefficients.
\end{lemma}
\leanref{Development/AbelianVarietyTheory/Proofs/DivisorProofs/CoefMapNormEquivariantBasisOfFive.lean}{GlobalTheory.GalOne.PhaseAt.io_eq_phase_smul_at_ofFive}
\leanref{Development/AbelianVarietyTheory/Proofs/DivisorProofs/FrAtEmb.lean}{GlobalTheory.GalOne.IsLtAt}
\begin{proof}
Write $\Delta=\Delta_i(\sigma)$, $A_0(w)=\phi\otimes_{f^{-1}}\Lambda(w)$, a linear map on $W_c$, and
$\mathrm{PT}(P)=\phi\otimes_{1,f^{-1}}P$. We use the following facts.
\begin{enumerate}
\item[(P1)] $\mathrm{PT}$ is injective. Indeed, $V_i$ has Gram matrix $1$, so the bases $1$ are orthonormal and adapted
  to the Cartan involution $1$, and Lemma~\ref{is-eigenvalue-conjugate-frame-extremal-fock}(3) applies as in the proof
  of Lemma~\ref{is-eigenvalue-conjugate-frame-extremal-line}(2).
\item[(P2)] For every $g\in U(V_i)(\mathbf A)$ with finite component $1$ there is $\chi\in\mathbf C$ with $|\chi|=1$ such
  that $\omega(g,1)\,\mathrm{PT}(P)=\chi\cdot\phi\otimes P$ for every $P$, where $\phi\otimes P$ is formed with the
  Cartan involutions $1$ and $s_\rho$ and the bases $g_\rho$ of $V_i$ and $f_\rho^{-1}$ of $W_i$.
\item[(P3)] The rational points $\gamma\in U(V_i)(F_h)$ with $\operatorname{tr}\gamma|_{W_c}\ne0$ generate
  $U(V_i)(F_h)$, and the only subspaces of $W_c$ stable under $U(V_i)(F_h)$ are $0$ and $W_c$.
\item[(P4)] Let $v_1$ be the highest weight vector of the rational coefficient vectors of $V_i$ over $L$, the tensor of
  the extremal vectors. Then $\|c_1(v_1)\|=1$. For every letter assignment $\mathrm{lt}$ there is a linear map $B_0$ from
  $W_c$ to the polynomials with $B_0(c_1(x''))=\mathrm{lt}(\sigma'(\mathrm{io}_O(\iota(x''))))$ for every rational
  coefficient vector $x''$ over $L$, and $B_0(c_1(v_1))=\Delta$. The map $\mathrm{io}_O$ of the one-map data is injective.
\item[(P5)] For every letter assignment and every $\gamma\in U(V_i)(F_h)$, written $\gamma=ak$ with $a$ of finite
  component $1$ and $k$ of archimedean component $1$, there is $c_\gamma\in\mathbf C$ with
  $\mathrm{PT}(B_0(\gamma w))=c_\gamma\,\omega(a,1)\,\mathrm{PT}(B_0(w))$ for every $w\in W_c$.
\end{enumerate}
% (P2) is AdTh.PhaseMap.om_real (OscillatorModelExtensionFirst.lean); (P3) is RatGen.rational_points_generate (r1, r2;
% RationalPointsGenerate.lean); (P4) is n1 (inner_of_int, pairVr_extremal and the pairing pin hc1.pair),
% PhaseMap.exists_ext with n4 (CoefExt.linearIndependent), hv (the pin hwOK of O.io and map_extremal) and hpG (toG_pG
% with WeightExponents.weight_exponents, the pins d1, d2 and hT.kap); (P5) is the first argument of proportional in
% main: rat_parts, om_real for a (x), hc1.act, n2 (CoefExt.up_rhoVr), the pin i3 of O.io (equivariance through cga
% and gaStar), toG_gaStar and hsub (SourceBasis.pureTensorO_mul_source), with c_gamma = sigma'(cga gammaE) x^-1.

\emph{Proof of (P2).} Let $g$ have finite component $1$, and write $g_\rho$ for its image under $\theta_\rho$. The level
vector $\phi$ is a theta function annihilated by the annihilation operators for the Cartan involutions $1$ of $V_i$ and
$s_\rho$ of $W_i$ (Lemma~\ref{is-eigenvalue-conjugate-theta-facts}(3)), and $1$ is a Cartan involution of $V_i$, as $V_i$
has Gram matrix $1$. Recall $n=\operatorname{rk}V_i$ and $r=\operatorname{rk}W_i$; $n$ is even, $n=2k$, and $r$ is odd.
The spaces $V_i$ and $W_i$ have correct signatures. As in the proof of
Lemma~\ref{is-eigenvalue-conjugate-frame-extremal-line}(3), the archimedean exponent of $\chi'$ at $\rho$ is an integer
$a_\rho$, and that of $\chi'\eta'^{-1}$ is $a_\rho-\frac12=b_\rho+\frac r2$, with $b_\rho=a_\rho-\frac{r+1}2$ an
integer. As $s_\rho=f_\rho^{-1}J_\rho f_\rho$ for the signature matrix $J_\rho$ of the model of $W_i$ at $\sigma$
(Lemma~\ref{is-eigenvalue-conjugate-frame-extremal-line}), the number
$x_\rho=\frac12\operatorname{tr}s_\rho+\frac r2$ is the number of entries $1$ of $J_\rho$, an integer.
Lemma~\ref{is-eigenvalue-conjugate-frame-phase-vacuum}, applied to $S=V_i$, $T=W_i$, the Cartan involutions $s_\rho$,
$\omega$ with its characters $(\chi',\chi'\eta'^{-1})$ and their exponents $a_\rho$ and $b_\rho+\frac r2$, the integers
$x_{1,\rho}=a_\rho+k$, $x_{2,\rho}=a_\rho-k$ and $x_{3,\rho}=-(x_\rho+b_\rho)$, $f=\phi$, $g$ and $h=1$, gives
$\omega(g,1)\phi=\chi\phi$ with
\[
 \chi=\prod_{\rho\in\mathcal T}\det(g_\rho)^{-(x_\rho+b_\rho)},
\]
as $d^\pm_\rho=\det1=1$ for $h=1$. As $g_\rho$ is unitary for the Gram matrix $1$, $|\det g_\rho|=1$, so $|\chi|=1$.
Lemma~\ref{is-eigenvalue-conjugate-theta-facts}(4), applied to $(g,1)$ with $s'_{S,\rho}=g_\rho1g_\rho^{-1}=1$,
$s'_{T,\rho}=s_\rho$, $B'_{S,\rho}=g_\rho\cdot1$ and $B'_{T,\rho}=f_\rho^{-1}$, gives
$\omega(g,1)\,\mathrm{PT}(P)=(\omega(g,1)\phi)\otimes P=\chi\cdot\phi\otimes P$ with the bases $g_\rho$ of $V_i$, as
$\phi\otimes P$ is linear in $\phi$.
% om_real (32 lines), followed: a from hP.expV, b from hP.expW (IsPair from VecConj.vec_conj_ofFive: hexp from the
% integers e_i of the datum and mu_inf, b = a - s - 1 for r = 2s + 1), x from trace_parity (hB, the bases onb of the
% model), k from hP.even; T.vacuum_character_positive_definite_source (Lemma -frame-phase-vacuum) with h = 1 and
% x1 = a + k, x2 = a - k, x3 = -(x + b) (hvac); OrientedHeis.om_pureTensorO and pureTensorO_smul; the norm through
% e1, e2 (the determinants at h = 1 are det 1) and e3 (det(g_rho) of norm 1, compAt_unitary).

\emph{Proof of (P3).} Let $\varrho$ be the representation of $U(V_i)(\mathbf A)$ on $W_c$: it is continuous, irreducible
and trivial on the elements with archimedean component $1$ (Definition~\ref{is-eigenvalue-lift-data}(7)). Let $G_\infty$
be the group of the elements with finite component $1$. The group $U(V_i)(\mathbf A)$ is the product, as a topological
group, of the group $U(V_i)_\infty$ of its archimedean components
(Lemma~\ref{is-eigenvalue-conjugate-frame-phase-connected}) and the group of its finite components, and $G_\infty$ is
the first factor. As $V_i$ is positive definite, by the core conditions on the models, and has correct signature,
$G_\infty$ is connected by Lemma~\ref{is-eigenvalue-conjugate-frame-phase-connected}. Every $g$ is $g_\infty g_f$ with
$g_\infty\in G_\infty$ and $g_f$ of archimedean component $1$, and
$\varrho(g)=\varrho(g_\infty)$. By weak approximation for unitary groups \cite[Chapter~7]{PR94}, for $u\in G_\infty$ and a
neighbourhood $N$ of $u$ in $U(V_i)(\mathbf A)$ there are $\gamma\in U(V_i)(F_h)$ and $k$ of archimedean component $1$ with
$\gamma k\in N$; so the map $\gamma\mapsto\gamma_\infty$ from $U(V_i)(F_h)$ to $G_\infty$ has dense image.

Generation. The set $N=\{u\in G_\infty:\operatorname{tr}\varrho(u)\ne0\}$ is open, as $\varrho$ is continuous, and contains
$1$, as $\operatorname{tr}\varrho(1)=\dim W_c>0$. Let $H$ be the subgroup of $U(V_i)(F_h)$ generated by the $\gamma$ with
$\gamma_\infty\in N$, which are the $\gamma$ with $\operatorname{tr}\gamma|_{W_c}\ne0$, and let $C$ be the closure in
$G_\infty$ of the image of $H$, a subgroup. By density, every point of the open set $N$ is a limit of points
$\gamma_\infty\in N$, so $N\subset C$; so $C$ is an open subgroup, hence also closed, and $C=G_\infty$ by connectedness.
Let $\gamma\in U(V_i)(F_h)$. The set of the $u\in G_\infty$ with $u^{-1}\gamma_\infty\in N$ is an open neighbourhood of
$\gamma_\infty\in C$, so it contains $h_\infty$ for some $h\in H$. Then $h^{-1}\gamma$ has $(h^{-1}\gamma)_\infty\in N$, so
$h^{-1}\gamma\in H$, and $\gamma\in H$.

Irreducibility. Let $U\subset W_c$ be a subspace stable under $U(V_i)(F_h)$, $g\in U(V_i)(\mathbf A)$ and $x\in U$.
Suppose $\varrho(g)x=\varrho(g_\infty)x\notin U$. The set of the $g'$ with $\varrho(g')x\notin U$ is open, as $U$ is closed,
and contains $g_\infty$; by weak approximation it contains some $\gamma k$ with $\gamma$ rational and $k$ of archimedean
component $1$. But $\varrho(\gamma k)x=\varrho(\gamma)x\in U$. So $U$ is stable under $U(V_i)(\mathbf A)$, and $U=0$ or
$U=W_c$ by irreducibility.
% rational_points_generate (62 lines), followed: weak_approx_real (h3), preconnected_inf (Lemma
% -frame-phase-connected, the Cayley transform), Mem.VolSplit.exists_adelic_split for the product decomposition, closure_preimage_eq_top
% (the generation argument above, for phi = the archimedean projection of the rational points, dense by h3), and the
% irreducibility through real_finite_split and h3.

\emph{Proof of (P4).} The one-map data at $i$ have a polynomial realization of the coefficients of $V_i$ over $L'$,
with an injective linear map $\mathrm{io}_O$ to the polynomials in the letter variables, a highest weight vector $w_h$
and a function $c_O$ on the matrices over $E_h$, and their pins say: $w_h$ is the tensor of the extremal vectors over
$L'$; $\mathrm{io}_O$ is injective; for every $g\in M_n(E_h)$ with $\bar g^{\mathsf T}g=1$ and every coefficient vector
$y$, $\mathrm{io}_O(g\cdot y)=c_O(g)\,\bar g^\star\big(\mathrm{io}_O(y)\big)$, where $\bar g$ is $g$ with complex conjugation
applied to its entries and, for a matrix $m$, $m^\star$ is the substitution of the letter variables
$((\tau',i'),k)\mapsto\sum_j\varepsilon_{\tau'}(m_{jk})\,((\tau',i'),j)$, with $\varepsilon_{\tau'}$ the embedding of $E_h$ in
$L'$ attached to $\tau'$; and $\mathrm{io}_O(w_h)=p_G$, the product over the embeddings $\tau'$ of $E_h$ and
$1\le i'\le m_{\tau'}$ of powers of the minors of size $i'$ of the matrices of letter variables at $\tau'$ (leading for
$\tau'$ in the CM type, trailing otherwise), whose exponents are the steps of the weights of $\kappa_i$.

The norm. The vector $v_1$ has integer coefficients, and $\iota(v_1)=w_h$, as $\iota$ carries extremal vectors to
extremal vectors. The isomorphism $c_1$ carries the rational pairing, followed by $\sigma$, to the inner product
(Definition~\ref{is-eigenvalue-conjugate-datum}(2)), on a basis; the coordinates of a vector with integer coefficients
in that basis lie in the subfield of $L$ on which the involution is complex conjugation under $\sigma$, so
$\|c_1(v_1)\|^2=\sigma(\langle v_1,v_1\rangle)$. As $V_i$ has Gram matrix $1$, the pairing of the
tensor of extremal vectors with itself is $1$, and $\|c_1(v_1)\|=1$.

The map $B_0$. The map $x''\mapsto\mathrm{lt}(\sigma'(\mathrm{io}_O(\iota(x''))))$ is $\sigma$-semilinear, as
$\sigma'\circ\iota=\sigma$. The map $c_1$ is $\sigma$-semilinear and carries linearly independent families to linearly
independent families (Definition~\ref{is-eigenvalue-conjugate-datum}(2)); so it carries a basis of the rational
coefficient vectors to a basis of its span, and the linear map $B_0$ defined on these images and extended to $W_c$ has
$B_0(c_1(x''))=\mathrm{lt}(\sigma'(\mathrm{io}_O(\iota(x''))))$ for every $x''$. Then
$B_0(c_1(v_1))=\mathrm{lt}(\sigma'(\mathrm{io}_O(w_h)))=\mathrm{lt}(\sigma'(p_G))$. The letter assignment carries the letter
variables of the minors of $p_G$ at $\tau'$ with $s\tau'=\rho$ and at $\tau'$ with $s\tau'=\bar\rho$ to the creation
variables of the minors of the delta polynomial at $\rho$, as $\mathrm{lt}$ is injective and is the identity on the
index $k$; by the pins of the core, the datum has signature $(p,q)$ at every embedding with $p$ and $q$ those of $W_i$;
and by the weight tie the steps of the weights of $\kappa_i$ at $\tau'$ with $s\tau'=\rho$ and $s\tau'=\bar\rho$ are the
exponents $a^+_\rho(\sigma)$ and $a^-_\rho(\sigma)$ of the delta polynomial. So $\mathrm{lt}(\sigma'(p_G))=\Delta$. The
injectivity of $\mathrm{io}_O$ is a pin.
% n1: inner_of_int (the inner product of the images of vectors with integer coefficients is sigma of the pairing,
% from hc1.pair on the basis bVV and fk_descent) and CoefExt.pairVr_extremal; exists_ext with n4; hv from hO.io.1
% (hwOK) and map_extremal; hpG: the pin i4 (io hw = pG) and toG_pG, a renaming of the minors, with the bijection
% lab sigma, IsLt (mv (lab sigma) lt), the type hty from weight_exponents (w1, w2), the pins d1, d2 (p, q of W_i) and the
% exponents w3 with hT.kap. The pins of O.io are IoLFRec: hwOK and IoLF (injective io, range in the span of HHa,
% equivariance io (rhoM ga v) = cga ga * gaStar (conj ga) (io v) for ga^H ga = 1, io hw = pG).

\emph{Proof of (P5).} Let $\gamma\in U(V_i)(F_h)$ with matrix $\gamma_E\in GL_n(E_h)$. As $V_i$ has Gram matrix $1$,
$\bar\gamma_E^{\mathsf T}\gamma_E=1$; write $\gamma=ak$ with $a$ of finite component $1$ and $k$ of archimedean component
$1$, so that $U_\rho=\theta_\rho(a)=\rho(\gamma_E)$, and $(U_\rho^{-1})_{jk}=\overline{\rho((\gamma_E)_{kj})}$. Let $x$ be as in
(P2) for $a$. For a family $U=(U_\rho)$ of invertible matrices let $s_U$ be the substitution of the creation variables
\[
 (\rho,a',k)\mapsto\sum_bu_{a'}(k,b)\,(\rho,a',b),\qquad u_{a'}(k,b)=(U_\rho^{-1})_{kb}\text{ if }a'\le p,\quad
 u_{a'}(k,b)=\overline{(U_\rho^{-1})_{kb}}\text{ if }a'>p,
\]
for the signature $(p,q)$ of $W_i$. Two computations are used. First, the polynomial vectors formed with the bases
$U_\rho$ of $V_i$ are those formed with the bases $1$ after the substitution: $\phi\otimes Q=\mathrm{PT}(s_U(Q))$ for every
$Q$. Second, for every polynomial $P$ in the letter variables,
$\mathrm{lt}\big(\sigma'(\bar\gamma_E^\star P)\big)=s_U\big(\mathrm{lt}(\sigma'(P))\big)$. Indeed, $\bar\gamma_E^\star$ sends the letter
variable $((\tau',i'),k)$ to $\sum_j\varepsilon_{\tau'}(\overline{(\gamma_E)_{jk}})((\tau',i'),j)$, and $\sigma'\circ\varepsilon_{\tau'}=s\tau'$,
as $\sigma'\circ\iota=\sigma$. If $\mathrm{lt}$ sends $((\tau',i'),k)$ to $(\rho,a',k)$ with $a'\le p$, then $s\tau'=\rho$ and the
coefficient becomes $\overline{\rho((\gamma_E)_{jk})}=(U_\rho^{-1})_{kj}$; if $a'>p$, then $s\tau'=\bar\rho$ and it becomes
$\rho((\gamma_E)_{jk})=\overline{(U_\rho^{-1})_{kj}}$. In both cases this is the coefficient of $s_U$.

Let $x''$ be a rational coefficient vector over $L$. By Definition~\ref{is-eigenvalue-conjugate-datum}(2),
$\gamma\,c_1(x'')=c_1(\gamma_E\cdot x'')$, and $\iota(\gamma_E\cdot x'')=\gamma_E\cdot\iota(x'')$, as $\iota$ commutes with the
action of the matrices over $E_h$. By the pin on $\mathrm{io}_O$, with $g=\gamma_E$,
$\mathrm{io}_O(\gamma_E\cdot\iota(x''))=c_O(\gamma_E)\,\bar\gamma_E^\star\big(\mathrm{io}_O(\iota(x''))\big)$. Applying $\mathrm{lt}\circ\sigma'$
and the two computations,
\[
 \mathrm{PT}\big(B_0(\gamma c_1(x''))\big)=\sigma'(c_O(\gamma_E))\,\mathrm{PT}\big(s_U(B_0(c_1(x'')))\big)
 =\sigma'(c_O(\gamma_E))\,\phi\otimes B_0(c_1(x''))=\sigma'(c_O(\gamma_E))\,x^{-1}\,\omega(a,1)\,\mathrm{PT}\big(B_0(c_1(x''))\big),
\]
with $\phi\otimes$ formed with the bases $U_\rho=\theta_\rho(a)$ of $V_i$, by (P2) for $a$, as $|x|=1$. The vectors $c_1(x'')$
span $W_c$, and both sides are linear in $w=c_1(x'')$; so (P5) holds with $c_\gamma=\sigma'(c_O(\gamma_E))x^{-1}$.
% (P5) is the first argument of proportional in main (the proof of h8, extended from the range of c1 by n3):
% rat_parts (gammaE, a, k, h1 to h6), om_real for a (x, hxa), hc1.act, hB0, n2 (CoefExt.up_rhoVr), the pin i3 of O.io,
% hsm, hsub (SourceBasis.pureTensorO_mul_source, the first computation) and toG_gaStar (the second, with
% U t = compAt (mt a), h6 and the relabelings e1, e2).

By Lemma~\ref{is-eigenvalue-conjugate-level-data}(1) the bases of $V_i$ of the realization are unitary, as $V_i$ has
Gram matrix $1$, so there is $g\in U(V_i)(\mathbf A)$ with finite component $1$ whose archimedean components are these
bases. Let $\chi$ be as in (P2) for $g$; then $\omega(g,1)\,\mathrm{PT}(P)=\chi\cdot\phi\otimes_{f^{-1}}P$ for every
$P$. By hypothesis $A_0(uw)=\omega(u,1)A_0(w)$ for $u$ with finite component $1$, and $\Lambda(w_0)=\Delta$. Put
$u_1=g^{-1}(\chi w_0)\in W_c$. Then
\[
 A_0(u_1)=\omega(g^{-1},1)\big(\chi\cdot\phi\otimes_{f^{-1}}\Delta\big)=\omega(g^{-1},1)\,\omega(g,1)\,\mathrm{PT}(\Delta)
 =\mathrm{PT}(\Delta),
\]
and $\|u_1\|=|\chi|\,\|w_0\|=1$, as $U(V_i)(\mathbf A)$ acts unitarily on $W_c$.

Let $\mathrm{lt}$ be a letter assignment and $B_0$ as in (P4). Apply
Lemma~\ref{is-eigenvalue-conjugate-frame-phase-proportional}(2) to $\Gamma=U(V_i)(F_h)$ acting on $\mathcal M=W_c$,
with $\Sigma$ the elements of nonzero trace, which is possible by (P3); to $A=A_0$ and $B=\mathrm{PT}\circ B_0$; and to
$w=u_1$ and $z=c_1(v_1)$. For $\gamma=ak$ as in (P5), $\gamma$ acts on $W_c$ as $a$ does, as the action is trivial on
the finite adelic points; so $A_0(\gamma w)=\omega(a,1)A_0(w)$, and (P5) gives the condition on $B$ with
$s_\gamma=\omega(a,1)$. And $A_0(u_1)=\mathrm{PT}(\Delta)=\mathrm{PT}(B_0(c_1(v_1)))$, which is not $0$: the maps
$\mathrm{PT}$ (P1), $P\mapsto\mathrm{lt}(\sigma'(P))$ (a renaming of the variables by the injective map $\mathrm{lt}$,
after the injective $\sigma'$ on the coefficients) and $\mathrm{io}_O$ are injective, and $v_1\ne0$, as $\|c_1(v_1)\|=1$. So there is $a_{\mathrm{lt}}$ with
$A_0=a_{\mathrm{lt}}\,\mathrm{PT}\circ B_0$ and $a_{\mathrm{lt}}u_1=c_1(v_1)$.

If there is no letter assignment, $\epsilon=1$ will do. Otherwise fix one, $\mathrm{lt}_0$, and put
$\epsilon=a_{\mathrm{lt}_0}$. Then $|\epsilon|=|\epsilon|\,\|u_1\|=\|c_1(v_1)\|=1$ by (P4). For every letter assignment
$\mathrm{lt}$, $a_{\mathrm{lt}}u_1=c_1(v_1)=\epsilon u_1$ and $u_1\ne0$, so $a_{\mathrm{lt}}=\epsilon$; and for every
rational coefficient vector $x''$ over $L$,
\[
 \phi\otimes_{f^{-1}}\Lambda(c_1(x''))=A_0(c_1(x''))=\epsilon\,\mathrm{PT}(B_0(c_1(x'')))
 =\epsilon\cdot\phi\otimes_{1,f^{-1}}\mathrm{lt}\big(\sigma'(\mathrm{io}_O(\iota(x'')))\big).
\]
% PhaseAt.io_eq_phase_smul_at_ofFive (207 lines with its statement), followed: VecConj.vec_conj_ofFive, factsAt_ofFive,
% vec_fock_ofFive', the highest weight vector v1 (hv, n1 = |c1 v1| = 1), n2 (rhoVr against rhoM), n3, n4,
% RatGen.rational_points_generate (r1, r2), WeightExponents.weight_exponents, the pins, hBV, hon, hsub
% (pureTensorO_mul_source), the injectivity hinj (FockVecInj.pureTensorO_injective), the element g with
% archimedean images BV (RealPts.glue), om_real for chi, hA0, hw1, hn1, and main: PhaseMap.exists_ext for B0, hpG
% (toG_pG, the delta polynomial) and AdTh.PhaseMap.proportional with rat_parts, om_real and toG_gaStar; the phase is a,
% of norm 1 by hn1, and does not depend on lt (the case split hex).
\end{proof}

\subsection{Sections of weight zero and conjugate theta lifts}\label{subsection-is-weight-zero-sections}

\begin{lemma}\label{is-eigenvalue-conjugate-frame-central}
In Situation~\ref{is-eigenvalue-conjugate-frame-core}, with the notation of
Definition~\ref{is-eigenvalue-conjugate-frame-marking} and Lemma~\ref{is-eigenvalue-conjugate-frame-marking-laws},
let $m$ be an integer prime to $M$. Then there is a central element $z$ of $G$ with $\ell_z(a)-ma\in M\mathfrak L$
for every $a\in\bar{\mathfrak M}\mathfrak L$.
\end{lemma}
\leanref{Development/AbelianVarietyTheory/Proofs/DivisorProofs/CoefMapNormEquivariantBasisOfFive.lean}{GlobalTheory.GalCentral.central_lev_ofFive}
\leanref{Formalism/LocalToGlobal/IndefiniteReal/ThPinsR.lean}{Formalism.IndefR.FWT.LevIs}
\begin{proof}
By the pins of the one-map data at $i$, the ideal $\mathfrak M$ is the product of the places of $P_g$, and
$\mathfrak M\bar{\mathfrak M}=M\mathcal O_{E_h}$; so $m$ is a unit in every residue field $k_P$, $P\in P_g$. Let
$z\in G=\prod_{P\in P_g}GL_r(k_P)$ have the scalar $m$ at every $P$; it is central. By the pins, $z$ has a level
lift $k$ (Lemma~\ref{is-eigenvalue-core-group}) with $\iota(\ell_z(a))-k_f\iota(a)\in M\hat L(1)$ for every
$a\in\bar{\mathfrak M}\mathfrak L$. The reduction of $k_f$ at $P\in P_g$, for the adelic matrix of $W_i$, has upper
left block $z_P$, the scalar $m$, and this block is the whole matrix. So for $\beta\in\bar{\mathfrak M}$ and
$u\in\hat L(1)$, $\beta(k_fu-mu)\in M\hat L(1)$, as $\mathfrak M\bar{\mathfrak M}=M\mathcal O_{E_h}$; and as $\bar{\mathfrak
M}\mathfrak L$ is spanned by such products $\beta a$, $k_f\iota(a)-m\iota(a)\in M\hat L(1)$ for every
$a\in\bar{\mathfrak M}\mathfrak L$. Hence $\iota(\ell_z(a)-ma)\in M\hat L(1)$; as $\iota$ is injective with image
$L(1)$, $\ell_z(a)-ma\in M\mathfrak L$.
\end{proof}
% central_lev_ofFive (51 lines with its statement): hMb from the tuple pin (hO.tuple ()).half.1, T.Q.l.prod, hu
% (isUnit_of_coprime, cast_res), z from R.eG.surjective and scalGL, central by scalGL_comm; then hO.lift.lev z
% (LevIs: a level lift k with the congruence hf), redFn_int (the reductions of k are m), exists_cong_int,
% Lift.cong_on_half and Lift.sign_div (through exists_div and LatEmb iotaW).

\begin{lemma}\label{is-eigenvalue-conjugate-frame-cover-sections}
Let a PEL datum over a CM field $E$, a level of it and its moduli scheme be given, with a coefficient field $F$ of the
moduli scheme and an embedding $\sigma'$ of $F$ (Situation~\ref{cv-cusp-setting}). A \emph{test ring} is a reduced
finitely generated $F$-algebra. A section $f$ of the moduli scheme over $F$ assigns to every framed point $y$ over a
test ring $S$ a value $f(y)\in S$; its value at a framed point over any $F$-algebra is read through a framed point over
a test ring of which it is the base change. A complex framed point over $\sigma'$ is a framed point over $\mathbf C$,
an $F$-algebra through $\sigma'$.
\begin{enumerate}
\item Every complex framed point over $\sigma'$ is the base change of a framed point over a test ring. A homomorphism
  of $F$-algebras $\varphi\colon S\to S'$ carries every framed point $y$ over $S$ to a framed point $\varphi_*y$ over
  $S'$, its base change, and if $S$ is a test ring, then $f(\varphi_*y)=\varphi(f(y))$ for every section $f$ of some
  weight.
\item Two sections $f$ and $f'$ of the same weight with $f(y)=f'(y)$ for every complex framed point $y$ over $\sigma'$
  are equal.
\item If $f$ is a section of weight zero and $y$, $y'$ are complex framed points over $\sigma'$ with the same point of the
  moduli scheme, then $f(y)=f(y')$.
\end{enumerate}
\end{lemma}
\leanref{Development/AbelianVarietyTheory/Proofs/UnitaryExistence/FramedPointReducedFromTestRing.lean}{GlobalTheory.Frames.allTest_of_secBaseChange}
\leanref{Development/AbelianVarietyTheory/Proofs/UnitaryExistence/ExistsFramedPointBaseChangeAlgebraMap.lean}{GlobalTheory.Frames.FramedBaseChange.exists_frMap}
\leanref{Development/AbelianVarietyTheory/Proofs/UnitaryExistence/ShortSpanUnitOfExactValues.lean}{GlobalTheory.ShortSpanCompare.evR_of_frMap}
\leanref{Development/AbelianVarietyTheory/Proofs/FoundationProofs/ExistsCovPinNotIsXiDet.lean}{GlobalTheory.CoverPinCompare.eq_of_complex}
\leanref{Development/AbelianVarietyTheory/Proofs/UnitaryExistence/WeightZeroSectionValueConjPoint.lean}{GlobalTheory.WeightZeroValue.evR_eq_of_same_point}
% (1) is the framework of framed points and sections: allTest_of_secBaseChange (descent of a framed point to a test
% ring, TestDescent.exists_fg_descent), FramedBaseChange.exists_frMap (base change along an algebra map) and
% ShortSpanCompare.evR_of_frMap (evR at a base change); a section is a Rul, defined on the framed points over test
% rings. (3) is WeightZeroValue.evR_eq_of_same_point, a self-contained computation (WeightZeroSectionValueConjPoint.lean,
% 190 lines): two framed points over C with the same point differ by a unique frame change (g, s)
% (frChg_exists_unique), and a section of weight zero is invariant under the generator change s and under every change
% g of the bases (evR_bor, keeps_all through keeps_only and keeps_lower).
\begin{proof}
(2) Let $y$ be a framed point over a test ring $S$; we show $f(y)=f'(y)$. The ring $S$ is reduced, and it is a Jacobson
ring, being finitely generated over a field; so the intersection of its maximal ideals is $0$. Let $\mathfrak m$ be a
maximal ideal of $S$. The field $S/\mathfrak m$ is a finite extension of $F$, so it has an $F$-embedding into $\mathbf C$,
an $F$-algebra through $\sigma'$; let $\psi\colon S\to\mathbf C$ be the composite, a homomorphism of $F$-algebras with
kernel $\mathfrak m$. By (1), $\psi_*y$ is a complex framed point over $\sigma'$, and
$\psi(f(y)-f'(y))=f(\psi_*y)-f'(\psi_*y)=0$. So $f(y)-f'(y)\in\mathfrak m$ for every $\mathfrak m$, and $f(y)=f'(y)$.
\end{proof}
% (2) is CoverPinCompare.eq_of_complex, followed: IsJacobsonRing (isJacobsonRing_of_finiteType), the reducedness of
% p.R, exists_algHom_ker, exists_frMap and evR_of_frMap.

\begin{situation}\label{is-eigenvalue-conjugate-frame-cover-situation}
Let $D$ be a definite-side datum with host field $E_h/F_h$, element $\xi_0$, ideal $\mathfrak A$ and set $\mathcal T$
of embeddings of $E_h$, and let $\theta$ be an archimedean family whose member $\theta_\tau$ extends $\tau$ for every
$\tau\in\mathcal T$. Let $W$ be a Hermitian space of $D$ of rank $r\ge1$ with Gram matrix $G_W$ and lattice $L_0$, let
$\hat L_0$ be the completion of $L_0$, let $h$ be the form of the matrix $-G_W^{\mathsf T}$ and $h_\tau$ its forms at the
embeddings $\tau$ (Definition~\ref{is-onto-share-flag-data}), and put $\langle u,v\rangle=\operatorname{Tr}_{E_h/\mathbf
Q}(\xi_0h(u,v))$ for $u,v\in E_h^r$. For $g\in U(W)(\mathbf A)$ let $\hat L(g)\subset\mathbf A_{E_h,f}^r$ be the subgroup
generated by the products $\alpha z$ with $\alpha\in\mathfrak A$ and $z\in g_f\hat L_0$, and $L(g)$ the $\mathcal
O_{E_h}$-span of the products $\alpha v$ with $\alpha\in\mathfrak A$ and $v\in E_h^r\cap g_f\hat L_0$.

Let a PEL datum of signature $(p,q)$ over $E_h$ with $p+q\ge1$ be given (Situation~\ref{cv-function-situation}), with
lattice $\mathfrak L$ and form $\langle\ ,\ \rangle_{\mathfrak L}$, and an injective $\mathcal O_{E_h}$-linear map
$\iota\colon\mathfrak L\to E_h^r$ with image $L(1)$ and $\langle a,b\rangle_{\mathfrak L}=\langle\iota a,\iota b\rangle$
for $a,b\in\mathfrak L$. Let $\ell$ be a level of the datum with integer $M$, ideal $\mathfrak M$ and full level $1$
(Definition~\ref{cv-model-toroidal-semiabelian}(1)). A \emph{marking} at $g\in U(W)(\mathbf A)$ is an $\mathcal
O_{E_h}$-linear map $m\colon\mathfrak L\to E_h^r$ with $m(a)-g_f\iota(a)\in M\hat L(g)$ for every $a\in\bar{\mathfrak
M}\mathfrak L$.

Let the moduli scheme of the datum at $\ell$, a coefficient field $K_c$ of it with a good embedding $s$, and a
uniformization at $s$ of rank $r$ with a uniformization witness be given, as in
Situation~\ref{is-eigenvalue-conjugate-frame-uniformization}, with components $c$, flag data $X_c$ with groups
$\Gamma_c$, modules $\Lambda_c$ and flag domains $\mathcal D_c$, framed points $y(c,x)$, the bijection $j\mapsto\ell_j$
onto the letters, lattices $L_c$ with forms $\langle\ ,\ \rangle_c$, maps $t_c$, isomorphisms $\varphi_c$ and markings
$m_c=\varphi_c\circ t_c$. Suppose that for every component $c$ the flag data $X_c$ have matrix $-G_W^{\mathsf T}$,
element $-\xi_0$ and set $\mathcal T$, and that $\langle a,b\rangle_c=\langle\varphi_ca,\varphi_cb\rangle$ for
$a,b\in L_c$. Among the compatibility conditions of the uniformization are:
\begin{enumerate}
\item $r=p+q$, and the integer $p_0$ of every $X_c$ is $q$;
\item every complex framed point of the moduli scheme over $s$ is triangularly equivalent to some $y(c,x)$: the two
  have the same point, and their tangent frames and sections differ by a triangular change and a scalar;
\item $y(c,x)$ and $y(c',x')$ are triangularly equivalent if and only if $c'=c$ and $x'=\gamma xb$ for some
  $\gamma\in\Gamma_c$ and $b\in\mathcal B^{\mathcal T}$;
\item for every $c$ and $x\in\mathcal D_c$, the point of $y(c,x)$ is the point $y_c(P_x)$ that the family
  uniformization (Theorem~\ref{cv-uniformization-family}) attaches to the period $P_x$ of $L_c$ whose $j$-th
  coordinate at $a\in L_c$ is the period of $\varphi_c(a)$ at $x$ at the letter $\ell_j$, for $s$
  (Definition~\ref{is-eigenvalue-conjugate-frame-fibres}(5)).
\end{enumerate}
\end{situation}
\leanref{Formalism/LocalToGlobal/ProjectiveLattices/UnFree.lean}{Formalism.L2.PjFree.IsUnf}
\leanref{Development/AbelianVarietyTheory/Definitions/ExistenceOfLattices/CoeSelfDualLatticeInFrameCoordinatesEq.lean}{GlobalTheory.Member.IotaOK}
\leanref{Formalism/LocalToGlobal/ProjectiveLattices/FamilyUniformisationBasePointOrthogonalBis.lean}{Formalism.L2.Pj.ModelW}
\leanref{Development/AbelianVarietyTheory/Proofs/JetHarmonics/FormPlus.lean}{GlobalTheory.ShortMove.Marked}
% The data of TwistTie.exists_dict_over: D : D7, V : HSp D (W), H, Dw (the PEL datum), l with l.A = 1, Xo, c, sigma
% with GoodEmb c sigma, U : Unf Xo c sigma V.rank with IsUnf U (the clauses r = p + q and p0 = q, the type of Psi,
% the surjectivity hU.2.2.1 and the triangular law hU.2.2.2.1 are (1) to (3)), ModelW V rfl (U.X k) for every k
% (X.Gm = -V.gram^T, X.xi0 = -xi0), the witness e, F, ik, phi of DictLat, whose first clause is the form law and whose
% period clause is (4) (Pm a j = perFr (U.X k) x (phi k a) (relab c sigma (e j).1) (e j).2 and (U.un k x).x =
% F.u (ik k) Pm), and IotaOK Dw V io (rt, rt0: Psi is the set T; emb: LatEmb, io injective with range latG V 1;
% isom: the form law of iota). L(g) is latG V g, Lhat(g) latGHat V g, and a marking at g with Lambda = L(g) is
% Marked V rfl io l Lambda m g.

\begin{lemma}\label{is-eigenvalue-conjugate-frame-cover-isometry}
In Situation~\ref{is-eigenvalue-conjugate-frame-cover-situation}, of which only $D$, $W$, $\iota$ and $\ell$ are used,
let $X$ and $X'$ be flag data of rank $r$ over $E_h$ with matrix $-G_W^{\mathsf T}$ and element $-\xi_0$, with modules
$\Lambda$ and $\Lambda'$. Let $L$ and $L'$ be lattices of PEL data over $E_h$ with forms $\langle\ ,\ \rangle_L$ and
$\langle\ ,\ \rangle_{L'}$, let $\varphi\colon L\to\Lambda$ and $\varphi'\colon L'\to\Lambda'$ be $\mathcal
O_{E_h}$-linear isomorphisms with $\langle a,b\rangle_L=\langle\varphi a,\varphi b\rangle$ and $\langle
a,b\rangle_{L'}=\langle\varphi'a,\varphi'b\rangle$, and let $t\colon\mathfrak L\to L$, $t'\colon\mathfrak L\to L'$ and
$s'\colon\mathfrak L\to\mathfrak L$ be $\mathcal O_{E_h}$-linear maps. Let $g,\gamma,\kappa\in U(W)(\mathbf A)$ and
$\delta\in GL_r(E_h)$ satisfy: $h(\delta u,\delta v)=h(u,v)$ for $u,v\in E_h^r$; $\gamma_fv=\delta v$ for $v\in E_h^r$;
$\hat L(\kappa)=\hat L(1)$; and $\kappa_f\iota(a)-\iota(s'(a))\in M\hat L(1)$ for every $a\in\bar{\mathfrak M}\mathfrak
L$. Suppose that $\Lambda=L(g)$ and $\varphi\circ t$ is a marking at $g$, and that $\Lambda'=L(\gamma g\kappa)$ and
$\varphi'\circ t'$ is a marking at $\gamma g\kappa$. Then there is an $\mathcal O_{E_h}$-linear isomorphism
$\beta\colon L\to L'$ with $\langle\beta a,\beta b\rangle_{L'}=\langle a,b\rangle_L$ for $a,b\in L$,
$\beta(t(s'(a)))-t'(a)\in ML'$ for $a\in\bar{\mathfrak M}\mathfrak L$, and $\varphi'(\beta a)=\delta\varphi(a)$ for
$a\in L$.
\end{lemma}
\leanref{Development/AbelianVarietyTheory/Proofs/HeightProofs/MarkedIsometryHalf.lean}{GlobalTheory.TwistTie.marked_isometry_half}
\leanref{Development/AbelianVarietyTheory/Proofs/JetHarmonics/MarkedLatticeIsometryDoubleCosetMove.lean}{GlobalTheory.ShortMoveIso.transport}
\leanref{Development/AbelianVarietyTheory/Proofs/UnitaryExistence/FrameUniformEmbeddingSecond.lean}{GlobalTheory.MarkCoset.latG_spec}
\begin{proof}
For $g'\in U(W)(\mathbf A)$, $\hat L(g')=g'_f\hat L(1)$, as $g'_f$ is $\mathcal O_{E_h}$-linear; and $L(g')=E_h^r\cap\hat
L(g')$, as a lattice of $E_h^r$ is the set of the vectors of $E_h^r$ that lie in its completion \cite{Wei67}. Hence
\[
 \hat L(\gamma g\kappa)=\gamma_fg_f\kappa_f\hat L(1)=\gamma_fg_f\hat L(\kappa)=\gamma_fg_f\hat L(1)=\gamma_f\hat L(g),
\]
and for $v\in E_h^r$, $v\in\Lambda$ exactly when $\delta v=\gamma_fv\in\Lambda'$. So $v\mapsto\delta v$ is an $\mathcal
O_{E_h}$-linear isomorphism of $\Lambda$ onto $\Lambda'$, with inverse $v\mapsto\delta^{-1}v$, and
$\beta=\varphi'^{-1}\circ\delta\circ\varphi$ is an isomorphism of $L$ onto $L'$ with $\varphi'(\beta a)=\delta\varphi(a)$.
As $\delta$ preserves $h$, $\langle\beta a,\beta b\rangle_{L'}=\langle\delta\varphi a,\delta\varphi b\rangle=\langle\varphi
a,\varphi b\rangle=\langle a,b\rangle_L$.

Let $a\in\bar{\mathfrak M}\mathfrak L$; then $s'(a)\in\bar{\mathfrak M}\mathfrak L$, as $s'$ is linear. The marking at
$g$ at $s'(a)$, the hypothesis on $\kappa$ and the marking at $\gamma g\kappa$ give $z_1\in\hat L(g)$, $w\in\hat L(1)$
and $z_3\in\hat L(\gamma g\kappa)$ with $\varphi(t(s'a))-g_f\iota(s'a)=Mz_1$, $\kappa_f\iota(a)-\iota(s'a)=Mw$ and
$\varphi'(t'a)-(\gamma g\kappa)_f\iota(a)=Mz_3$. As $\varphi(t(s'a))\in E_h^r$, $\gamma_f\varphi(t(s'a))=\delta\varphi(t(s'a))$, and
\[
 \delta\varphi(t(s'a))-\varphi'(t'a)=\gamma_f\big(\varphi(t(s'a))-g_f\iota(s'a)\big)-\gamma_fg_f\big(\kappa_f\iota(a)-\iota(s'a)\big)
 -\big(\varphi'(t'a)-(\gamma g\kappa)_f\iota(a)\big)=M(\gamma_fz_1-\gamma_fg_fw-z_3).
\]
Here $\gamma_fz_1\in\gamma_f\hat L(g)$ and $\gamma_fg_fw\in\gamma_fg_f\hat L(1)$, both equal to $\hat L(\gamma g\kappa)$. So
$u=M^{-1}\big(\delta\varphi(t(s'a))-\varphi'(t'a)\big)$ lies in $E_h^r\cap\hat L(\gamma g\kappa)=\Lambda'$. With
$x=\varphi'^{-1}(u)$, $\varphi'(\beta(t(s'a)))=\delta\varphi(t(s'a))=\varphi'(t'a+Mx)$, so $\beta(t(s'a))=t'(a)+Mx$.
\end{proof}
% marked_isometry_half (64 lines), followed: latG_spec (L(g) = E^r cap Lhat(g) and Lhat(g) = g_f Lhat(1)), transport
% (ht1 for the lattices, ht2 for the three terms), psi = delta on X.Lambda, gam = phi'^-1 psi phi, the forms through
% ModelW and hdelta_u, and MarkCoset.finVec_div for the division by M.

\begin{lemma}\label{is-eigenvalue-conjugate-frame-cover-compare}
Let $\Lambda$ and $\Lambda'$ be the lattices of two PEL data over a CM field $E$
(Situation~\ref{cv-function-situation}), with forms $\langle\ ,\ \rangle$ and $\langle\ ,\ \rangle'$. Let $G$ be an
$\mathcal O_E$-semi-abelian scheme over a scheme $S$, $x\colon\operatorname{Spec}k\to S$ a point with $k$ an
algebraically closed field, and $((\zeta_m),(\iota_m))$ and $((\zeta'_m),(\iota'_m))$ Tate systems of $G$ at $x$ for
$\Lambda$ and for $\Lambda'$ (Definition~\ref{cv-toroidal-moduli-object}(2), with $\Lambda'$ in place of $\Lambda$ for the
second). Then there are $\mathcal O_E$-linear maps $f_m\colon\Lambda\to\Lambda'$ and integers $\nu_m$, $m\ge1$, such
that for all $m,m'\ge1$:
\begin{enumerate}
\item $\iota_m(a)=\iota'_m(f_m(a))$ for $a\in\Lambda$, and $\zeta_m=(\zeta'_m)^{\nu_m}$;
\item $f_{mm'}(a)-f_m(a)\in m\Lambda'$ for $a\in\Lambda$;
\item $f_m(a)\in m\Lambda'$ exactly when $a\in m\Lambda$, and $\Lambda'=f_m(\Lambda)+m\Lambda'$;
\item $\nu_m$ is prime to $m$, and $\nu_{mm'}\equiv\nu_m$ modulo $m$;
\item $\langle f_m(a),f_m(b)\rangle'\equiv\nu_m\langle a,b\rangle$ modulo $m$ for $a,b\in\Lambda$.
\end{enumerate}
\end{lemma}
\leanref{Development/AbelianVarietyTheory/Proofs/UnitaryExistence/TateIsoExistsLiftOfTwoData.lean}{GlobalTheory.tateIso_exists_lift_of_two_data}
\leanref{Development/AbelianVarietyTheory/Proofs/UnitaryExistence/TateIsoExistsLiftOfTwoData.lean}{GlobalTheory.tateIso_similar_of_two_data}
\begin{proof}
(1) Fix $m$. For $a\in\Lambda$, $m\iota_m(a)=\iota_m(ma)=0$, so $\iota_m(a)$ is $m$-torsion and equals $\iota'_m(a')$ for
some $a'\in\Lambda'$. The pairs $(a,a')\in\Lambda\times\Lambda'$ with $\iota_m(a)=\iota'_m(a')$ form an $\mathcal
O_E$-submodule, which the first projection maps onto $\Lambda$. As $\Lambda$ is projective, this projection has an
$\mathcal O_E$-linear section; its composite with the second projection is a map $f_m$ with
$\iota_m(a)=\iota'_m(f_m(a))$. As $\zeta_m$ and $\zeta'_m$ are primitive $m$-th roots of unity in $k$,
$\zeta_m=(\zeta'_m)^{\nu_m}$ for an integer $\nu_m\ge0$.

(2) $\iota'_m(f_{mm'}(a))=m'\iota'_{mm'}(f_{mm'}(a))=m'\iota_{mm'}(a)=\iota_m(a)=\iota'_m(f_m(a))$, and the kernel of
$\iota'_m$ is $m\Lambda'$.

(3) $f_m(a)\in m\Lambda'$ exactly when $\iota'_m(f_m(a))=\iota_m(a)$ is $0$, that is, when $a\in m\Lambda$. For
$a'\in\Lambda'$, $\iota'_m(a')$ is $m$-torsion, so it is $\iota_m(a)=\iota'_m(f_m(a))$ for some $a\in\Lambda$, and
$a'-f_m(a)\in m\Lambda'$.

(4) $(\zeta'_m)^{\nu_m}=\zeta_m$ is primitive of order $m$, so $\nu_m$ is prime to $m$. And
$(\zeta'_m)^{\nu_{mm'}}=(\zeta'_{mm'})^{m'\nu_{mm'}}=\zeta_{mm'}^{m'}=\zeta_m=(\zeta'_m)^{\nu_m}$, so $m$ divides
$\nu_{mm'}-\nu_m$.

(5) For the commutator pairing $e_m$,
$(\zeta'_m)^{\nu_m\langle a,b\rangle}=\zeta_m^{\langle a,b\rangle}=e_m(\iota_m(a),\iota_m(b))=e_m(\iota'_m(f_ma),\iota'_m(f_mb))
=(\zeta'_m)^{\langle f_m(a),f_m(b)\rangle'}$, and $\zeta'_m$ has order $m$.
\end{proof}
% tateIso_exists_lift_of_two_data (the submodule R, Module.projective_lifting_property for the section) and
% tateIso_similar_of_two_data (60 lines), followed: f n and nu n by choice, the clauses through tower, ker, onto,
% zeta_pow and zeta_prim (pow_iff_coprime, zpow_eq_one_iff_dvd), and isom with the uniqueness he of the Weil value.

\begin{lemma}\label{is-eigenvalue-conjugate-frame-cover-norm}
Let a PEL datum over a CM field $E$ be given (Situation~\ref{cv-function-situation}), with conjugation $b\mapsto\bar b$,
and let $(M,\mathfrak M,N)$ be a level of it (Definition~\ref{cv-model-toroidal-semiabelian}(1)).
\begin{enumerate}
\item For every $n\ge1$ and every integer $\nu$ prime to $n$ there is $t\in\mathcal O_E$ with $t\bar t-\nu\in
  n\mathcal O_E$ and, if $M$ divides $n$, $t-1\in\mathfrak M$.
\item Let $\nu_n$, $n\ge1$, be integers with $\nu_n$ prime to $n$ and $\nu_{nm}\equiv\nu_n$ modulo $n$ for all
  $n,m\ge1$. Then there are $t_n\in\mathcal O_E$, $n\ge1$, with $t_n\bar t_n-\nu_n\in n\mathcal O_E$ and
  $t_{nm}-t_n\in n\mathcal O_E$ for all $n,m\ge1$, and $t_M-1\in\mathfrak M$.
\end{enumerate}
\end{lemma}
\leanref{Development/AbelianVarietyTheory/Proofs/UnitaryExistence/ExistsHostNormModCoprimeLevelHalfStatement.lean}{GlobalTheory.T.exists_hostNorm_mod_coprime_level_half}
\leanref{Development/AbelianVarietyTheory/Proofs/UnitaryExistence/ExistsHostNormModCoprimeLevelHalf.lean}{GlobalTheory.NormMod.exists_norm_mod}
\leanref{Development/AbelianVarietyTheory/Proofs/UnitaryExistence/ExistsCompatibleNormTowerProfiniteStatement.lean}{GlobalTheory.T.exists_compatible_norm_tower_profinite}
\leanref{Development/AbelianVarietyTheory/Proofs/UnitaryExistence/ExistsCompatibleNormTowerProfinite.lean}{GlobalTheory.NormTower.exists_norm_tower}
\begin{proof}
(1) This is a self-contained computation in $\mathcal O_E$, in four steps. \emph{A prime.} For $n=2$, $\nu$ is odd and
$t=1$ will do. For an odd prime $p$: as no finite prime of the totally real subfield of $E$ ramifies in $E$, the
elements $b-\bar b$, $b\in\mathcal O_E$, generate the unit ideal, so the finite ring $\mathcal O_E/p\mathcal O_E$ has a
unit $\delta$ with $\bar\delta=-\delta$; the quadratic form $x^2-\delta^2y^2$ on the subring fixed by the conjugation
represents every unit of it, and $t=x+y\delta$ for $x^2-\delta^2y^2=\nu$ has $t\bar t=\nu$ modulo $p$. \emph{From $m$ to
$mp$.} Let $t\bar t-\nu=my$. If $p$ does not divide $m$, and $s$ is a solution modulo $p$, then $u=bpt+ams$ with
$am+bp=1$ is congruent to $t$ modulo $m$ and to $s$ modulo $p$, so $u\bar u\equiv\nu$ modulo $mp$. If $p$ divides $m$,
then $\bar y=y$, and with $a\nu+bp=1$ and $w_0=(1+\sqrt{-d})/2$, for which $w_0+\bar w_0=1$, the element
$t'=t(1-myaw_0)$ satisfies $t'\bar t'\equiv\nu+my(1-a\nu)=\nu+mybp\equiv\nu$ modulo $mp$; in both cases the new
solution is congruent to $t$ modulo $m$. \emph{Induction} on the prime factors of $k$ lifts a solution modulo $m$ to a
solution modulo $mk$, congruent to it modulo $m$, for every $k$ prime to $\nu$. \emph{The level.} Write $1=a+b$ with
$a\in\mathfrak M$ and $b\in\bar{\mathfrak M}$. Then $t=b+\nu a$ has $t-1=(\nu-1)a\in\mathfrak M$ and
$t\bar t-\nu=(1-\nu)b\bar b+(\nu^2-\nu)a\bar a\in\mathfrak M\bar{\mathfrak M}=M\mathcal O_E$. If $M$ divides $n$, lifting
this $t$ from $M$ to $n$ keeps $t-1\in\mathfrak M$, as $M\in\mathfrak M$; otherwise lift $t=1$ from $1$ to $n$.

(2) For $j\ge1$, (1) gives $s_j\in\mathcal O_E$ with $s_j\bar s_j-\nu_{j!}\in j!\,\mathcal O_E$, and $s_j-1\in\mathfrak M$
if $M$ divides $j!$. Fix a nonprincipal ultrafilter $\mathcal U$ on the positive integers; it contains every cofinite
set, in particular the set of the $j\ge n$, for which $n$ divides $j!$. For every $n$ the ring $\mathcal O_E/n\mathcal
O_E$ is finite, so there is $t_n\in\mathcal O_E$ such that the set of the $j$ with $s_j-t_n\in n\mathcal O_E$ lies in
$\mathcal U$. For $n\ge1$ choose $j$ in this set with $n\mid j!$, say $j!=nm$; then
$t_n\bar t_n\equiv s_j\bar s_j\equiv\nu_{nm}\equiv\nu_n$ modulo $n$. For $n,m\ge1$ choose $j$ in the sets of $t_n$ and
of $t_{nm}$; then $t_{nm}\equiv s_j\equiv t_n$ modulo $n$. Finally choose $j$ in the set of $t_M$ with $M\mid j!$; then
$t_M-1=(s_j-1)-(s_j-t_M)\in\mathfrak M$, as $M\in\mathfrak M$.
\end{proof}
% (1) is NormMod.exists_norm_mod (ExistsHostNormModCoprimeLevelHalf.lean, 330 lines): exists_norm_mod_prime (with
% span_sub_bar from H.unram, exists_isUnit_neg, exists_norm_eq and exists_sq_sub_mul_sq), exists_norm_mod_mul_prime
% (the two cases, with cmHost_bar_halfIntegerGenerator for w0 + bar w0 = 1), exists_norm_mod_mul (induction_on_primes),
% exists_norm_mod_level (t = b + nu a) and exists_norm_mod. (2) is NormTower.exists_norm_tower, followed: s j from
% T.exists_hostNorm_mod_coprime_level_half at j!, the hyperfilter, exists_eventually_dvd_sub (Ultrafilter.eq_pure_of_finite
% on the finite quotient) and the three clauses.

\begin{lemma}\label{is-eigenvalue-conjugate-frame-cover-tower}
In Situation~\ref{cv-cusp-setting}, let $L$, with form $\langle\ ,\ \rangle$, the level $(M',\mathfrak M')$, the moduli
scheme $X'$ and the type $k\mapsto\sigma_k$ be as in Theorem~\ref{cv-uniformization-family}, and let $L_i$, with form
$\langle\ ,\ \rangle_i$, and $t_i\colon L\to L_i$ be one of the lattices and maps that it gives. Let $P$ be a period of
$L_i$ of this type. Then there are $\mathcal O_E$-linear maps $f_N\colon L\to L_i$, $N\ge1$, such that for all
$N,N'\ge1$:
\begin{enumerate}
\item[(i)] $f_{NN'}(a)-f_N(a)\in NL_i$ for $a\in L$;
\item[(ii)] $f_N(a)\in NL_i$ exactly when $a\in NL$, and $L_i=f_N(L)+NL_i$;
\item[(iii)] $\langle f_Na,f_Na'\rangle_i-\langle a,a'\rangle\in N\mathbf Z$ for $a,a'\in L$;
\item[(iv)] $f_{M'}(a)-t_i(a)\in M'L_i$ for $a\in\bar{\mathfrak M}'L$.
\end{enumerate}
\end{lemma}
\leanref{Development/AbelianVarietyTheory/Proofs/UnitaryExistence/ExistsInverseUnitTower.lean}{GlobalTheory.familyUnif_lattice_genus}
\leanref{Development/AbelianVarietyTheory/Proofs/UnitaryExistence/TateIsoExistsLiftOfTwoData.lean}{GlobalTheory.familyUnif_lattice_similar}
\leanref{Development/AbelianVarietyTheory/Proofs/UnitaryExistence/ExistsInverseUnitTower.lean}{GlobalTheory.exists_inverse_unit_tower}
\begin{proof}
Let $y=y_i(P)$ be the complex point of $X'$ of Theorem~\ref{cv-uniformization-family}, $A$ its abelian variety and
$\eta$ its level structure, and write $[z]\in A$ for the image of $z\in\mathbf C^g$ under the isomorphism
$\mathbf C^g/P(L_i)\to A$ of that theorem.

\emph{Two Tate systems.} The universal object of $X'$ is an object at the level $(M',\mathfrak M',1)$
(Definition~\ref{cv-toroidal-moduli-object}(4) and~(5)), so $A$ has a Tate system $((\zeta_m),(\iota_m))$ at $y$ for
$L$ with $\iota_{M'}(a)=\eta(a)$ for $a\in\bar{\mathfrak M}'L$. Put $\zeta'_m=e^{-2\pi i/m}$ and
$\iota'_m(b)=[P(b)/m]$ for $b\in L_i$. This is a Tate system of $A$ at $y$ for $L_i$: $\zeta'_m$ is primitive of order
$m$ and $(\zeta'_{mm'})^{m'}=\zeta'_m$; $\iota'_m$ is additive, and $\mathcal O_E$-linear as the isomorphism is
compatible with the actions of $\mathcal O_E$ and $P$ has the type $k\mapsto\sigma_k$; $\iota'_m(b)=0$ exactly when
$P(b)/m\in P(L_i)$, that is, as $P$ is injective, when $b\in mL_i$; a point $[z]$ is $m$-torsion exactly when $mz\in
P(L_i)$, that is, when $[z]=\iota'_m(b)$ for some $b$; $m'\iota'_{mm'}(b)=[P(b)/m]=\iota'_m(b)$; and the Weil pairing
of $\iota'_m(a)$ and $\iota'_m(b)$ is $e^{-2\pi i\langle a,b\rangle_i/m}=(\zeta'_m)^{\langle a,b\rangle_i}$, as the
form of $L_i$ gives the Weil pairings. The level structure of $y$ is $\eta(a)=[P(t_ia)/M']=\iota'_{M'}(t_i(a))$ for
$a\in\bar{\mathfrak M}'L$.

\emph{The similitudes.} Lemma~\ref{is-eigenvalue-conjugate-frame-cover-compare} gives $\mathcal O_E$-linear maps
$g_N\colon L\to L_i$ and integers $\nu_N$ with its properties (1) to (5). For $a\in\bar{\mathfrak M}'L$,
$\iota'_{M'}(g_{M'}(a))=\iota_{M'}(a)=\eta(a)=\iota'_{M'}(t_i(a))$, so $g_{M'}(a)-t_i(a)\in M'L_i$.

\emph{The correction.} As $\nu_N$ is prime to $N$, there is an integer $\kappa_N$ with $\kappa_N\nu_N\equiv1$
modulo $N$; then $\kappa_N$ is prime to $N$, and $\kappa_{NN'}-\kappa_N\equiv\kappa_N\kappa_{NN'}(\nu_N-\nu_{NN'})\equiv0$
modulo $N$. Lemma~\ref{is-eigenvalue-conjugate-frame-cover-norm}(2), for the $\kappa_N$ and the level
$(M',\mathfrak M',1)$, gives $s_N\in\mathcal O_E$ with $s_N\bar s_N-\kappa_N\in N\mathcal O_E$, $s_{NN'}-s_N\in
N\mathcal O_E$ and $s_{M'}-1\in\mathfrak M'$. Put $r_N=\nu_N\bar s_N$, so that $r_Ns_N-1\in N\mathcal O_E$, and
$f_N=s_Ng_N$. These maps have the four properties.

(i) $f_{NN'}(a)-f_N(a)=s_{NN'}\big(g_{NN'}(a)-g_N(a)\big)+(s_{NN'}-s_N)g_N(a)\in NL_i$.

(ii) If $f_N(a)\in NL_i$, then $g_N(a)=r_Nf_N(a)-(r_Ns_N-1)g_N(a)\in NL_i$, so $a\in NL$; and if $a\in NL$, then
$f_N(a)\in NL_i$. For $a'\in L_i$ write $a'=g_N(a)+Ne$; then $g_N(a)=f_N(r_Na)-(r_Ns_N-1)g_N(a)$, so $a'\in
f_N(L)+NL_i$.

(iii) As $\langle bx,x'\rangle=\langle x,\bar bx'\rangle$ for $b\in\mathcal O_E$,
$\langle s_Na,s_Na'\rangle=\langle a,\bar s_Ns_Na'\rangle\equiv\kappa_N\langle a,a'\rangle$ modulo $N$. So
$\langle f_Na,f_Na'\rangle_i=\langle g_N(s_Na),g_N(s_Na')\rangle_i\equiv\nu_N\langle s_Na,s_Na'\rangle
\equiv\nu_N\kappa_N\langle a,a'\rangle\equiv\langle a,a'\rangle$ modulo $N$.

(iv) Let $a\in\bar{\mathfrak M}'L$. As $s_{M'}-1\in\mathfrak M'$ and $\mathfrak M'\bar{\mathfrak M}'=M'\mathcal O_E$,
$s_{M'}a=a+M'e'$ with $e'\in L$. So $f_{M'}(a)=g_{M'}(s_{M'}a)=g_{M'}(a)+M'g_{M'}(e')\in t_i(a)+M'L_i$.
\end{proof}
% familyUnif_lattice_genus (55 lines), followed: familyUnif_lattice_similar (Xo.univ.geom at U.u i P for the Tate system
% I with hM, familyUnifTateIso for the analytic one with expNegUnit, U.ex_add, ex_lin, ex_ker, ex_onto and ex_weil,
% tateIso_similar_of_two_data with weilIs_unique_fieldPoint, and U.ex_eta with tateIso_ev_eq_iff_add for the clause
% at M), exists_inverse_unit_tower (kappa), T.exists_compatible_norm_tower_profinite (s, Lemma -frame-cover-norm(2)),
% hinv (r = nu bar s), the maps s n . f n and the four clauses, with D.adj for (iii) and level_smul_mem_half for (iv).
% The Lean's statement holds for every level l; its clause at the full level A is not used here.

\begin{lemma}\label{is-eigenvalue-conjugate-frame-cover-marked}
In Situation~\ref{is-eigenvalue-conjugate-frame-cover-situation}, for every component $c$ there is $g\in U(W)(\mathbf A)$
with $\Lambda_c=L(g)$ such that $m_c$ is a marking at $g$.
\end{lemma}
\leanref{Development/AbelianVarietyTheory/Proofs/JetHarmonics/ExistsAdelicElementOfMarkedComponentLattice.lean}{GlobalTheory.ShortMove.exists_marked_of_component}
\leanref{Development/AbelianVarietyTheory/Proofs/UnitaryExistence/ComponentAttainedAtBaseFrame.lean}{GlobalTheory.RecDict.comp_adelic_of_tower}
\leanref{Development/AbelianVarietyTheory/Proofs/UnitaryExistence/AdelicCompletionFractionalIdealDivisibleEqZero.lean}{GlobalTheory.AdLat.exists_adelic_matrix_of_tower}
\begin{proof}
Let $P$ be a period of $L_c$ of the type of the uniformization. By
Lemma~\ref{is-eigenvalue-conjugate-frame-cover-tower}, applied to the family uniformization of condition (4) of the
situation, its lattice $L_c$ with the map $t_c$, and $P$, there are $\mathcal O_{E_h}$-linear maps
$f_N\colon\mathfrak L\to L_c$, $N\ge1$, with the following properties:
\begin{enumerate}
\item[(i)] $f_{NN'}(a)-f_N(a)\in NL_c$ for $a\in\mathfrak L$ and $N,N'\ge1$;
\item[(ii)] $f_N(a)\in NL_c$ exactly when $a\in N\mathfrak L$, and $L_c=f_N(\mathfrak L)+NL_c$;
\item[(iii)] $\langle f_Na,f_Na'\rangle_c-\langle a,a'\rangle_{\mathfrak L}\in N\mathbf Z$ for $a,a'\in\mathfrak L$;
\item[(iv)] $f_M(a)-t_c(a)\in ML_c$ for $a\in\bar{\mathfrak M}\mathfrak L$.
\end{enumerate}
% familyUnif_lattice_genus (Lemma -frame-cover-tower). The period is PeriodExists.typeFn_period_exist with the type of
% the bijection onto the letters.

The maps $\varphi_c\circ f_N\circ\iota^{-1}$ from $L(1)=\iota(\mathfrak L)$ to $\Lambda_c$ satisfy (i) and (ii). Their limit
is an invertible matrix $g_f$ over $\mathbf A_{E_h,f}$ that carries the completion $\hat L(1)$ of $L(1)$ onto the
completion $\hat\Lambda_c$ of $\Lambda_c$, with $g_f\iota(a)-\varphi_c(f_N(a))\in N\hat\Lambda_c$ for every $N\ge1$ and
$a\in\mathfrak L$: by (i) the vectors $\varphi_c(f_N(a))$ converge in $\hat\Lambda_c$, and by (ii) the limit map is
bijective. It preserves the form $h$. Indeed, let $J$ be the $\mathcal O_{E_h}$-module of the $z\in E_h$ with
$\operatorname{Tr}_{E_h/\mathbf Q}(\xi_0bz)\in\mathbf Z$ for every $b\in\mathcal O_{E_h}$, a fractional ideal. As
$\langle\ ,\ \rangle_c$ takes integer values and $\Lambda_c$ is an $\mathcal O_{E_h}$-module, $h(u,v)\in J$ for
$u,v\in\Lambda_c$; and by (iii), applied to $ba$ and $a'$, $h(\varphi_cf_Na,\varphi_cf_Na')-h(\iota a,\iota a')\in NJ$ for
$a,a'\in\mathfrak L$. Letting $N$ grow, $h(g_f\iota a,g_f\iota a')=h(\iota a,\iota a')$, and by linearity and continuity
$g_f$ preserves $h$ on $\mathbf A_{E_h,f}^r$.

So there is $g\in U(W)(\mathbf A)$ with archimedean component $1$ and finite component $g_f$. Then $\hat L(g)=g_f\hat
L(1)=\hat\Lambda_c$, and $L(g)=E_h^r\cap\hat L(g)=E_h^r\cap\hat\Lambda_c=\Lambda_c$, as in the proof of
Lemma~\ref{is-eigenvalue-conjugate-frame-cover-isometry}. Let $a\in\bar{\mathfrak M}\mathfrak L$. By (iv),
$f_M(a)=t_c(a)+Me$ with $e\in L_c$, and $g_f\iota(a)=\varphi_c(f_M(a))+Mw$ with $w\in\hat\Lambda_c$. So
$m_c(a)-g_f\iota(a)=\varphi_c(t_c(a))-\varphi_c(t_c(a))-M\varphi_c(e)-Mw=-M(\varphi_c(e)+w)\in M\hat L(g)$, and $m_c$ is
a marking at $g$.
\end{proof}
% exists_marked_of_component (49 lines), followed: typeFn_period_exist for P, familyUnif_lattice_genus for f (h1 to
% h4 are (i) to (iv)), comp_latG, comp_adelic_of_tower (through AdLat.exists_adelic_matrix_of_tower for the limit, and
% the fractional ideal J = xi^-1 times the inverse different for the form), the first clause of
% T.adelicUnitary_exists_with_prescribed_parts for u, AdLat.hat_rational_points_and_approx for L(u) = Lambda, and the
% computation with he and hgw. The hypotheses LatEmb and FormPlus of the Lean are the conditions on iota.

\begin{lemma}\label{is-eigenvalue-conjugate-frame-cover-over}
In Situation~\ref{is-eigenvalue-conjugate-frame-cover-situation}, let $J$ be the diagonal $r\times r$ matrix with
entries $1$ at the first $q$ places and $-1$ at the others, and $x_0=(x_{0,\tau})_{\tau\in\mathcal T}$ a family of
invertible complex $r\times r$ matrices whose columns have Gram matrix $J$ for $h_\tau$. Let $g\mapsto(c(g),x(g))$ be
a map from $U(W)(\mathbf A)$ to the pairs of a component $c$ and a point of $\mathcal D_c$ such that for every $g$
there is $\gamma\in U(W)(F_h)$ with $\Lambda_{c(g)}=L(\gamma g)$, $m_{c(g)}$ a marking at $\gamma g$, and
$x(g)_\tau=\theta_\tau(\gamma g)\,x_{0,\tau}$ for every $\tau\in\mathcal T$. Then every complex framed point of the
moduli scheme over $s$ has the point of $y(c(g),x(g))$ for some $g\in U(W)(\mathbf A)$.
\end{lemma}
\leanref{Development/AbelianVarietyTheory/Proofs/HeightProofs/MarkedIsometryHalf.lean}{GlobalTheory.TwistTie.exists_dict_over}
\leanref{Development/AbelianVarietyTheory/Proofs/HeightProofs/MarkedIsometryHalf.lean}{GlobalTheory.TwistTie.perFr_of_unitary}
\leanref{Development/AbelianVarietyTheory/Proofs/HeightProofs/MarkedIsometryHalf.lean}{GlobalTheory.TwistTie.point_eq_of_isometry}
\leanref{Development/AbelianVarietyTheory/Proofs/UnitaryExistence/FrameMulTriangularOrthonormal.lean}{GlobalTheory.UnitFrames.frame_mul_triangular_orthonormal}
\leanref{Development/AbelianVarietyTheory/Proofs/UnitaryExistence/FrameUniformEmbeddingSecond.lean}{GlobalTheory.MarkCoset.rational_matrix}
\begin{proof}
Let $y$ be a complex framed point over $s$. By condition (2) of the situation, $y$ has the point of $y(c,x)$ for some
component $c$ and some $x\in\mathcal D_c$.

\emph{An orthonormal point.} The integer $p_0$ of $X_c$ is $q$, by condition (1). By the Gram--Schmidt process for the
forms $h_\tau$, there is $b\in\mathcal B^{\mathcal T}$ such that the columns of every $x_{1,\tau}$, for $x_1=xb$, have
Gram matrix $J$ for $h_\tau$; as the leading minors of $J$ are the signs $\epsilon_i$, $x_1\in\mathcal D_c$. By
condition (3), with $\gamma=1$, $y(c,x)$ and $y(c,x_1)$ are triangularly equivalent, so they have the same point. For
$\tau\in\mathcal T$ put $u_\tau=x_{1,\tau}x_{0,\tau}^{-1}$. The columns of $x_{1,\tau}=u_\tau x_{0,\tau}$ and of
$x_{0,\tau}$ have the same Gram matrix, so $u_\tau$ preserves $h_\tau$; equivalently, $u_\tau^*G_W(\tau)u_\tau=G_W(\tau)$
for $G_W(\tau)=\tau(G_W)$. By Lemma~\ref{is-eigenvalue-conjugate-frame-glue}(2) there is $a\in U(W)(\mathbf A)$ with
finite component $1$ and $\theta_\tau(a)=u_\tau$ for every $\tau\in\mathcal T$.

\emph{A marked element.} By Lemma~\ref{is-eigenvalue-conjugate-frame-cover-marked} there is $g_0$ with
$\Lambda_c=L(g_0)$ such that $m_c$ is a marking at $g_0$. Let $f\in U(W)(\mathbf A)$ have archimedean component $1$
and finite component $(g_0)_f$, and put $g=af$. Then $g_f=(g_0)_f$. The modules $L(\cdot)$ and $\hat L(\cdot)$ and the
markings at an element depend only on its finite component, so $\Lambda_c=L(g)$ and $m_c$ is a marking at $g$.

\emph{The isometry.} Let $\gamma$ be as in the hypothesis for $g$. The rational element $\gamma$ is a matrix
$\delta\in GL_r(E_h)$ that preserves the form of $W$, so $h(\delta u,\delta v)=h(u,v)$ for $u,v\in E_h^r$,
$\gamma_fv=\delta v$ for $v\in E_h^r$, and $\theta_\tau(\gamma)=\tau(\delta)$ for $\tau\in\mathcal T$, as $\theta_\tau$
extends $\tau$. Lemma~\ref{is-eigenvalue-conjugate-frame-cover-isometry}, applied to $X=X_c$, $X'=X_{c(g)}$,
$L=L_c$, $L'=L_{c(g)}$, $\varphi=\varphi_c$, $\varphi'=\varphi_{c(g)}$, $t=t_c$, $t'=t_{c(g)}$, $s'$ the identity, $g$,
$\gamma$, $\kappa=1$ and $\delta$, gives an $\mathcal O_{E_h}$-linear isomorphism $\beta\colon L_c\to L_{c(g)}$ that
preserves the forms, with $\beta(t_c(a))-t_{c(g)}(a)\in ML_{c(g)}$ for $a\in\bar{\mathfrak M}\mathfrak L$ and
$\varphi_{c(g)}(\beta a)=\delta\varphi_c(a)$ for $a\in L_c$.

\emph{The periods.} As $\theta_\tau$ is multiplicative and $f$ has archimedean component $1$,
$x(g)_\tau=\theta_\tau(\gamma)\theta_\tau(a)\theta_\tau(f)x_{0,\tau}=\tau(\delta)u_\tau x_{0,\tau}=\tau(\delta)x_{1,\tau}$ for
$\tau\in\mathcal T$. So for $v\in E_h^r$ the period of $\delta v$ at $x(g)$ at a letter $(\tau',i)$, with
$\tau=s\circ\varepsilon_{\tau'}$, is the period of $v$ at $x_1$ there: if $\tau\in\mathcal T$,
$(x(g)_\tau^{-1}\tau(\delta v))_{q+i}=(x_{1,\tau}^{-1}\tau(\delta)^{-1}\tau(\delta)\tau(v))_{q+i}=(x_{1,\tau}^{-1}\tau(v))_{q+i}$;
if $\bar\tau\in\mathcal T$, then, as $\bar\tau(\delta)$ preserves $h_{\bar\tau}$,
$h_{\bar\tau}(x(g)_{\bar\tau}e_{q+1-i},\bar\tau(\delta v))=h_{\bar\tau}(\bar\tau(\delta)x_{1,\bar\tau}e_{q+1-i},\bar\tau(\delta)\bar\tau(v))
=h_{\bar\tau}(x_{1,\bar\tau}e_{q+1-i},\bar\tau(v))$. Hence the periods $P=P_{x_1}$ of $L_c$ and $P'=P_{x(g)}$ of $L_{c(g)}$
of condition (4) satisfy $P'(\beta a)=P(a)$ for $a\in L_c$.

\emph{The point.} By condition (4), $y(c,x_1)$ has the point $y_c(P)$ and $y(c(g),x(g))$ the point $y_{c(g)}(P')$. The
level $\ell$ has full level $1$ and $s$ is good, so Lemma~\ref{cv-components-points}(1), applied to $\beta$ and the
identity of $\mathbf C^{g'}$, $g'=|\mathcal T|r$, shows that these points are equal. So $y$ has the point of
$y(c(g),x(g))$.
\end{proof}
% exists_dict_over (60 lines), followed: hU.2.2.1 for (k, x, t); frame_mul_triangular_orthonormal (h1: x1 and b with
% x1 = x b orthonormal; h4: two orthonormal frames differ by u with u^T Gm u-bar = Gm) and
% unit_frame_eq_orthonormal_mul_triangular for x1 in dom; the triangular law hU.2.2.2.1 with gamma = 1 for ht'; the
% archimedean lift hL of T.archimedeanMatrix_unitary_continuous_multiplicative for a (Lemma frame-glue(2), with the
% isInfEmb_unique comparison of the two archimedean families); ShortMove.exists_marked_of_component for h0 (Lemma
% frame-cover-marked); UnitarySplit.real_finite_split (h0 = a0 * f) and the second clause of
% T.adelicUnitaryLattice_rationalTranslate_invariance for a * f; hD.2 at a * f; MarkCoset.rational_matrix for delta;
% marked_isometry_half with s = id and kappa = 1 (Lemma frame-cover-isometry); hD.1.2 for the periods at x1 and at the
% dictionary point (condition (4)); point_eq_of_isometry (UnifPoint.point_eq_iff_isometry_closed, Lemma
% cv-components-points(1), with relab_good for the type and l.A = 1) and perFr_of_unitary for the periods.

\begin{lemma}\label{is-eigenvalue-conjugate-frame-cover-constant}
In Situation~\ref{is-eigenvalue-conjugate-frame-core}, there is $c_0\in\mathbf C$ with $|c_0|=1$ such that, as
$\sigma$ agrees with $\sigma_0$ on $K$, the delta section of the one-map data at $i$ takes the value $c_0$ at
every complex framed point of its moduli scheme over $\sigma'$.
\end{lemma}
\leanref{Development/AbelianVarietyTheory/Proofs/DivisorProofs/CoefMapNormEquivariantBasisOfFive.lean}{GlobalTheory.GalOne.cover_const_at_ofFive}
\leanref{Development/AbelianVarietyTheory/Proofs/IwasawaDecomposition/OneMapTypingsEquiv.lean}{GlobalTheory.OneMapRefCompare.cover_const}
\leanref{Development/AbelianVarietyTheory/Proofs/DivisorProofs/FrAtEmb.lean}{GlobalTheory.GalOne.CovAt}
\begin{proof}
We first show that there is $u\in L'$ with $u^m=1$ for some $m>0$ such that the delta section is $u$ times the
unit section. The delta data of the one-map data at $i$ come with a moduli scheme and a coefficient field
identified with $L'$; the embedding $\sigma'$ gives an embedding of the latter, and the pins of the one-map data
include a uniformization of that moduli scheme over it, with a dictionary $h\mapsto(c_2(h),x_2(h))$ on
$U(W_i)(\mathbf A)$, and the following clause on the delta section: there are $c\in\mathbf C$ and $m>0$ with
$c^m=1$ such that, for every $h$ and every $\gamma\in U(W_i)(F_h)$ that marks $h$ for the dictionary, the reading of
the delta section at $(c_2(h),x_2(h))$ is $c\,\xi(\det(\gamma h))$, where $\xi$ is the character of the delta data
at the base tuple. By the pins at the base tuple, $\xi(\det(\gamma h))=\xi_{\mu_0}(y)$ for a determinant root $y$ of
$\gamma h$ (Lemma~\ref{is-eigenvalue-conjugate-frame-root-values}(1)), and $\xi_{\mu_0}=1$ by the core conditions on
the level characters. As every $h$ is marked by some $\gamma$, the reading of the delta section at every dictionary
point is $c$. The reading at a point outside the flag domain of its component is $0$, and $c\ne0$; so the
dictionary point of $1$ lies in the flag domain, and gives a complex framed point of the moduli scheme of the delta
data over $\sigma'$. Two facts finish the argument:
\begin{enumerate}
\item[(A)] a section of weight zero whose readings at all the dictionary points are $c$ takes the value $c$ at every
  complex framed point over $\sigma'$;
\item[(B)] a section of weight zero that takes the value $c$ at every complex framed point over $\sigma'$, of which
  there is one, is $u\cdot1$ for some $u$ in the coefficient field with $\sigma'(u)=c$.
\end{enumerate}
By (A) and (B) the delta section is $u\cdot1$ with $\sigma'(u)=c$; then $\sigma'(u^m)=c^m=1$, and $u^m=1$ as $\sigma'$
is injective.

Put $c_0=\sigma'(u)=c$; then $c_0^m=1$, so $|c_0|=1$, and the value of the delta section at a complex framed point
over $\sigma'$ is $\sigma'(u)\cdot1=c_0$.

\emph{Proof of (A).} Let $f$ be a section of weight zero whose readings at all the dictionary points are $c$, and let $y$
be a complex framed point over $\sigma'$. Every dictionary point $(c_2(h),x_2(h))$ lies in the flag domain of its
component: by the lattice law of the dictionary, $x_2(h)_\tau$ is the image of some $\gamma h$, $\gamma\in U(W_i)(F_h)$, at
$\tau$ times the frame $x_{0,\tau}$, whose Gram matrix for $-G_{W_i}^{\mathsf T}$ is the diagonal matrix of the signs; as
the image of $\gamma h$ is unitary, $x_2(h)$ is orthonormal, and orthonormal frames lie in the flag domain. The data
of the delta data are as in Situation~\ref{is-eigenvalue-conjugate-frame-cover-situation}: with $D$, $\theta$, $W=W_i$,
the embedding $\iota$ of the lattice of the datum of the core at level $i+1$, which has image $L(1)$ and carries the
form of the lattice to $\langle\ ,\ \rangle$, the level of the delta data, which has full level $1$, its moduli scheme,
the coefficient field with the good embedding given by $\sigma'$, and the uniformization of the pins with its
witness, whose flag data have matrix $-G_{W_i}^{\mathsf T}$, element $-\xi_0$ and set $\mathcal T$ and whose lattices
$L_c$ have the forms $\langle\varphi_ca,\varphi_cb\rangle$, by the first law of the dictionary. With the frames
$x_{0,\tau}$ and the lattice law of the dictionary, Lemma~\ref{is-eigenvalue-conjugate-frame-cover-over} gives $h$
such that $y$ and the complex framed point $y(c_2(h),x_2(h))$ of the uniformization have the same point of the
moduli scheme. By Lemma~\ref{is-eigenvalue-conjugate-frame-cover-sections}(3), $f(y)=f(y(c_2(h),x_2(h)))$, and this
is the reading of $f$ at $(c_2(h),x_2(h))$, a point of the flag domain, which is $c$.
% TwistTie.exists_dict_over (Lemma is-eigenvalue-conjugate-frame-cover-over), with the hypotheses from the pins as in
% cover_const: TupleCover.iotaOK_and_frame (IotaOK for iotaW T i and the frames x0), hpos (1 <= p + q),
% (hP.cov ()).2.1 (l.A = 1), and from hP.cxi the IsUnf U2, ModelW, the witness and DictLat dict2.

\emph{Proof of (B).} Let $f$ be a section of weight zero over $F=L'$ with the value $c$ at every complex framed point over
$\sigma'$, and let $y$ be one. By Lemma~\ref{is-eigenvalue-conjugate-frame-cover-sections}(1), $y=\varphi_*y_1$ for a
framed point $y_1$ over a test ring $S$ and a homomorphism $\varphi\colon S\to\mathbf C$ of $F$-algebras. Let
$\mathfrak m$ be a maximal ideal of $S$ containing the kernel of $\varphi$, $K'=S/\mathfrak m$, a finite extension of $F$
as $S$ is finitely generated, and $a\in K'$ the image of $f(y_1)$. For every $F$-embedding $j$ of $K'$ into $\mathbf C$,
the composite $\psi\colon S\to K'\to\mathbf C$ carries $y_1$ to a complex framed point $\psi_*y_1$ over $\sigma'$, and
$j(a)=\psi(f(y_1))=f(\psi_*y_1)=c$ by (1) of that lemma. Summing over the $[K':F]$ embeddings,
$\sigma'(\operatorname{Tr}_{K'/F}a)=[K':F]\,c$, so $u=\operatorname{Tr}_{K'/F}(a)/[K':F]\in F$ satisfies $\sigma'(u)=c$. The
sections $f$ and $u\cdot1$ have weight zero, and $u\cdot1$ has the value $\sigma'(u)=c$ at every complex framed point over
$\sigma'$. By Lemma~\ref{is-eigenvalue-conjugate-frame-cover-sections}(2), $f=u\cdot1$.
\end{proof}
% cover_const_at_ofFive (39 lines with its statement): OneMapRefCompare.cover_const (hT.laws, hR.lev, the pins;
% OneMapTypingsEquiv.lean) for u, Complex.norm_eq_one_of_pow_eq_one, evR_smul and evR_one. cover_const (30 lines) is
% followed above: hP.cxi (CXiVal: c0, m, IsUnf U2, ModelW, DictLat dict2 and the readings rdU of cdelta at the
% dictionary points equal to c0 * xi (det (gamma h)) for MarkBy h gamma), the tuple pin (hP.tuple ()).xi with
% isDetRoot_laws, TwistFactorCompare.eq_one_of_chrEq_one with hL.xi_base, hv, hd (dif_neg outside the domain) and
% the framed point (O.U2 ()).un; (A) is Member.const_of_const_on_dictLat (ConstOfConstOnDictLat.lean, 66 lines,
% followed: hdom through hfr of DictLat, archimedeanMatrix_unitary_continuous_multiplicative and
% UnitFrames.unit_frame_eq_orthonormal_mul_triangular; TwistTie.exists_dict_over; WeightZeroValue.evR_eq_of_same_point;
% dif_pos for the reading), with iotaOK_and_frame for its hypotheses; (B) is OneMapRefCompare.eq_const_of_complex (20
% lines, followed: allTest_of_secBaseChange, Ideal.exists_le_maximal, finite_of_finite_type_of_isJacobsonRing,
% trace_value through trace_eq_sum_embeddings, exists_frMap, evR_of_frMap and CoverPinCompare.eq_of_complex).

\begin{lemma}\label{is-eigenvalue-conjugate-frame-inputs-roots}
Let a PEL datum over a CM field $E$ be given (Situation~\ref{cv-function-situation}), with lattice $\mathfrak L$, and let
$\ell$ be a level of it with integer $M_\ell$ and ideal $\mathfrak M_\ell$
(Definition~\ref{cv-model-toroidal-semiabelian}(1)); put $Q=\mathfrak L/\mathfrak M_\ell\mathfrak L$. Let $r,n\ge1$, let
$L\subset E^r$ be an $\mathcal O_E$-submodule, $m\colon\mathfrak L\to E^r$ an $\mathcal O_E$-linear map and $N\ge1$. A
\emph{family of roots} of $m$ for $\ell$, $L$ and $N$ is a map $Y$ from $Q^n$ to $M_{n\times r}(E)$ such that for every
$q\in Q^n$ and $1\le b\le n$ the $b$-th row $y$ of $Y(q)$ satisfies: $y-M_\ell^{-1}m(a)\in L$ for some
$a\in\bar{\mathfrak M}_\ell\mathfrak L$ in the class $q_b$; and there is $j\ge0$ with $sy\in2NL$ for every
$s\in\mathfrak M_\ell^j$.
\begin{enumerate}
\item If $m$ has values in $L$, there is a family of roots of $m$ for $\ell$, $L$ and $N$.
\item Let $\mu$ be an integer prime to $M_\ell$. Suppose that $m$ has values in $L$, that $a\in M_\ell\mathfrak L$
  whenever $a\in\bar{\mathfrak M}_\ell\mathfrak L$ and $m(a)\in M_\ell L$, and that every $v\in\bar{\mathfrak M}_\ell L$ is
  $m(a)+M_\ell u$ with $a\in\bar{\mathfrak M}_\ell\mathfrak L$ and $u\in L$. Then $\mu m$ has the same three properties.
\item Let $\mu$ be an integer prime to $M_\ell$, suppose that $m$ has values in $L$, and let $Y$ be a family of roots
  of $m$ for $\ell$, $L$ and $N$. Then $q\mapsto Y(\mu q)$ is a family of roots of $\mu m$ for $\ell$, $L$ and $N$,
  and $\sum_{q\in Q^n}F(\mu q)=\sum_{q\in Q^n}F(q)$ for every function $F$ on $Q^n$ with values in a complex vector
  space.
\end{enumerate}
\end{lemma}
\leanref{Development/AbelianVarietyTheory/Definitions/SpanningFamilies/IsStabForm.lean}{GlobalTheory.Dict.IsRoots}
\leanref{Development/AbelianVarietyTheory/Definitions/SpanningFamilies/IsStabForm.lean}{GlobalTheory.Dict.IsMark}
\leanref{Development/AbelianVarietyTheory/Proofs/SiegelThetaTransformation/DualHalfLevelMarkingExistsAndMatch.lean}{GlobalTheory.Dict.DualMark.exists_roots}
\leanref{Development/AbelianVarietyTheory/Proofs/SiegelThetaTransformation/DualHalfLevelMarkingExistsAndMatch.lean}{GlobalTheory.Dict.DualMark.exists_pow_le}
\leanref{Development/AbelianVarietyTheory/Proofs/DivisorProofs/MarkingRootsDilationCoprimeLevelStatement.lean}{GlobalTheory.T.marking_roots_dilation_coprime_level}
\leanref{Development/AbelianVarietyTheory/Proofs/DivisorProofs/MarkingRootsDilationCoprimeLevel.lean}{GlobalTheory.DilateLevel.dilate}
\begin{proof}
Write $M=M_\ell$ and $\mathfrak M=\mathfrak M_\ell$, so that $\mathfrak M\bar{\mathfrak M}=M\mathcal O_E$ and
$\mathfrak M+\bar{\mathfrak M}=\mathcal O_E$. Fix $c\in\bar{\mathfrak M}$ with $c-1\in\mathfrak M$, and a representative
$\tilde q\in\mathfrak L$ of every class $q\in Q$. Then $c\tilde q\in\bar{\mathfrak M}\mathfrak L$ lies in the class $q$,
as $(c-1)\tilde q\in\mathfrak M\mathfrak L$. Moreover, if $x\in\bar{\mathfrak M}\mathfrak L$ lies in the class $q$, then
$x-c\tilde q=c(x-\tilde q)-(c-1)x\in\bar{\mathfrak M}\mathfrak M\mathfrak L=M\mathfrak L$. If $m$ has values in $L$ and
$a\in\bar{\mathfrak M}\mathfrak L$, then $sa\in M\mathfrak L$ for $s\in\mathfrak M$, so $sM^{-1}m(a)\in L$.

(1) The ring $\mathcal O_E/2N\mathcal O_E$ is finite, so the decreasing chain of its submodules
$\mathfrak M^j(\mathcal O_E/2N\mathcal O_E)$ is stationary from some $j$ on; for that $j$ the finitely generated module
$P=\mathfrak M^j(\mathcal O_E/2N\mathcal O_E)$ satisfies $\mathfrak MP=P$, and by Nakayama's lemma
\cite[Corollary~2.5]{AM69} there is $t\in\mathcal O_E$ with $t-1\in\mathfrak M$ and $tP=0$, that is,
$t\mathfrak M^j\subset2N\mathcal O_E$. Let $Y(q)$ have the rows $tM^{-1}m(c\tilde q_b)$. Then
$Y(q)_b-M^{-1}m(c\tilde q_b)=(t-1)M^{-1}m(c\tilde q_b)\in L$, and $c\tilde q_b$ lies in the class $q_b$. For $s=xy$ with
$x\in\mathfrak M^j$ and $y\in\mathfrak M$, $sY(q)_b=(tx)\cdot yM^{-1}m(c\tilde q_b)\in2NL$; such products generate
$\mathfrak M^{j+1}$.

(2) Choose integers $\alpha,\beta$ with $\alpha\mu+\beta M=1$, so that $x=\alpha\mu x+\beta Mx$ for every $x$. The map
$\mu m$ has values in $L$. If $a\in\bar{\mathfrak M}\mathfrak L$ and $\mu m(a)=Mu$ with $u\in L$, then
$m(a)=\alpha Mu+\beta Mm(a)=M(\alpha u+\beta m(a))$, so $a\in M\mathfrak L$. If $v\in\bar{\mathfrak M}L$, then $\alpha
v=m(a_1)+Mu_1$ with $a_1\in\bar{\mathfrak M}\mathfrak L$ and $u_1\in L$, and $v=\mu\alpha v+\beta Mv=\mu
m(a_1)+M(\mu u_1+\beta v)$, with $\beta v\in L$.

(3) Let $q\in Q^n$ and $1\le b\le n$. As $Y$ is a family of roots, there is $a_0\in\bar{\mathfrak M}\mathfrak L$ in the
class $\mu q_b$ with $Y(\mu q)_b-M^{-1}m(a_0)\in L$, and the condition on $\mathfrak M^j$ holds for the row $Y(\mu q)_b$.
Put $a'=c\tilde q_b$, in the class $q_b$. The elements $a_0$ and $\mu a'$ of $\bar{\mathfrak M}\mathfrak L$ both lie in
the class $\mu q_b$, so each differs from $c\widetilde{\mu q_b}$ by an element of $M\mathfrak L$, and $a_0-\mu a'=Mz$
with $z\in\mathfrak L$. Then $Y(\mu q)_b-M^{-1}\mu m(a')=\big(Y(\mu q)_b-M^{-1}m(a_0)\big)+m(z)\in L$. Finally, $Mq=0$
for $q\in Q^n$, as $M\in\mathfrak M$; so $q\mapsto\mu q$ is a bijection of $Q^n$ with inverse $q\mapsto\alpha q$, and the
two sums agree.
\end{proof}
% exists_roots (30 lines), followed: exists_pow_le (IsArtinian.monotone_stabilizes and
% Submodule.exists_sub_one_mem_and_smul_eq_zero_of_fg_of_le_smul) for t and j, cSep_spec for c, key (smul_eq_nsmul with
% prod_le), and Y q b = t M^-1 mk (cSep levRep (q b)). marking_roots_dilation_coprime_level is DilateLevel.dilate (55
% lines), followed: hbez (Bezout), the IsMark clauses mem, inj and surj, the roots through sub_sepRep and
% cSep_smul_levQ, and for gphi the bijection e of the sum (finsum_comp_equiv). The Lean states the last clause for the
% class functions gphi X bas w Lambda L v c Y, the sums over q of c(q) heisG (Y q) (projA ...); the book states the
% reindexing it rests on.

\begin{lemma}\label{is-eigenvalue-conjugate-frame-inputs}
Let $R$ be a realization of a definite-side datum $D$ over a field $K_s\subset L$ with a face $F$ that
satisfies the core conditions at $F$, except possibly the one on its chain, with the notation of
Definitions~\ref{is-eigenvalue-lift-data} and~\ref{is-eigenvalue-conjugate-datum}. Let $i$ be a lift and
$\sigma$ an embedding of $L$ that agrees with $\sigma_0$ on $K$. Let $\chi'$ be a $\sigma$-conjugate of
$\mu^0_i$ and $\eta'$ a $\sigma$-conjugate of $\eta$ (Lemma~\ref{is-eigenvalue-type-parts}(3) and~(4)),
$\omega$ a Weil datum for $(V_i,W_i)$ with characters $(\chi',\chi'\eta'^{-1})$, and $\phi$ and $B$ a
level vector of the trivial level and bases as in Lemma~\ref{is-eigenvalue-conjugate-level-data}(1), with
$\phi_Q=\phi\otimes_BQ$.
\begin{enumerate}
\item Let $J$ be the level of $V_i$ at $\mu_0$, and $\chi'_J$ a character of $J$ whose values are the
  images under $\sigma$ of those of the level character of $R$ on $J$. Then
  $\omega(k,1)\phi_Q=\chi'_J(k_f)^{-1}\phi_Q$ for every polynomial $Q$ and every $k\in U(V_i)(\mathbf A)$
  with archimedean component $1$ whose finite part $k_f$ lies in $J$.
\item Let moreover one-map data for $R$ be given, with a core whose coefficient field is $K_s$, weight
  families $\kappa_j$, vectors $v_0(j,\sigma)$ and a weight tie, pinned at $i$, with $L'$, $\iota$ and
  the representatives of the classes of Lemma~\ref{is-eigenvalue-conjugate-class-sums}, the action
  $f\mapsto gf$ of $G$ on the core at level $i+1$ that the one-map data at $i$ give, and the core forms,
  the cusp forms of weight $\kappa_i$ of the core at level $i+1$. Let $\sigma'$ be an embedding of $L'$ with
  $\sigma'\circ\iota=\sigma$, and $s$ the restriction of $\sigma$ to $K_s$. Write $h\mapsto(c(h),x(h))$
  for the dictionary of the core at level $i+1$ and $s$ and $\mu$ for its multiplier
  (Lemma~\ref{is-eigenvalue-conjugate-dictionary-good}), $\mathrm{rd}(f)$ for the reading at $s$ of a
  section $f$ of the core at level $i+1$ over $K_s$, and, for a function $a$ on the components,
  $\mathrm{ad}_a(P)(h)=a(c(h))\,\mu(h)\,P_{c(h)}(x(h))$; let $\ell_i=(2\pi|\rho(\xi_0)|)^{-N/2}$, where $N$
  is the sum of the weights of $\kappa_i$ over the letters of the datum of the core at level $i+1$ and
  $|\rho(\xi_0)|$ does not depend on $\rho$, so that
  $\mathrm{ad}_\sigma(f)=\ell_i\,\mathrm{ad}_{c_E}(\mathrm{rd}(f))$ for every section $f$, where $c_E$
  are the canonical constants of the core. Let $W_c$ and $c_1$ be as in (2) of
  Definition~\ref{is-eigenvalue-conjugate-datum}, and $\mathrm{io}$ a map as in
  Definition~\ref{is-eigenvalue-lift-data}(7) for $\omega$, $\phi$, $B$ and $\Delta_i(\sigma)$. Then there
  are a function $a$ on the components with $|a|=1$, $\kappa\in\mathbf C$ with $|\kappa|=1$ and a central
  element $z$ of $G$ with the following properties.
  \begin{enumerate}
  \item[(a)] For every $g\in G$ with a level lift $k$ (Lemma~\ref{is-eigenvalue-core-group}) and every
    core form $f$, $gf$ is a core form, and there is $b\in K_s$ with
    $\sigma_0(b)=\prod_{P\in P_g}\chi_{i+1,P}(\det g_P)$, where $\chi_{i+1,P}$ is the character of
    $k_P^\times$ of the realization at the level $i+1$ and $P$, such that
    $\mathrm{ad}_a(\mathrm{rd}(gf))(h)=s(b)\,\mathrm{ad}_a(\mathrm{rd}(f))(hk)$ for every $h\in
    U(W_i)(\mathbf A)$.
  \item[(b)] Let $h_1,\dots,h_N$ be representatives of the classes with the properties of the pins
    (Lemma~\ref{is-eigenvalue-conjugate-class-sums}(1)), with the volume $\mathrm{vol}\in L'$ and the
    stabilizer orders $\gamma_j$ of the classes. Let $x_1,\dots,x_N$ be rational coefficient vectors of
    $V_i$ over $L'$ and $x''_1,\dots,x''_N$ over $L$ with $\iota(x''_j)=x_j$, $f'$ a section of weight
    $\kappa_i$ of the core at level $i+1$ over $L'$ built from $(x_j)$
    (Lemma~\ref{is-eigenvalue-conjugate-class-sums}(2)), and $f_1,\dots,f_n$ sections over $K_s$ and
    $a_1,\dots,a_n\in L'$ with $zf'=\sum_ka_kf_k$ after extension of scalars
    (Lemma~\ref{is-eigenvalue-conjugate-class-action}). Then for every $g\in U(W_i)(\mathbf A)$
    \[
     \kappa\sum_j\sigma'(\mathrm{vol}\,\gamma_j^{-1})\,\theta_\omega\big(h_j,g;
     \phi_{\mathrm{io}(v_0(i,\sigma)\otimes c_1(x''_j))}\big)
     =\ell_i\sum_k\sigma'(a_k)\,\mathrm{ad}_a(\mathrm{rd}(f_k))(g).
    \]
  \end{enumerate}
\end{enumerate}
\end{lemma}
\leanref{Development/AbelianVarietyTheory/Proofs/DivisorProofs/CoefMapNormEquivariantBasisOfFive.lean}{GlobalTheory.GalOne.omega_level_conjVec_ofFive}
\leanref{Development/AbelianVarietyTheory/Proofs/TwistedOperators/OmPureTensorOOfFinite.lean}{GlobalTheory.FiniteFock.om_pureTensorO_of_finite}
\leanref{Development/AbelianVarietyTheory/Proofs/DivisorProofs/CoefMapNormEquivariantBasisOfFive.lean}{GlobalTheory.GalOne.FrameConst.exists_frAtEmb_ofFive}
\leanref{Development/AbelianVarietyTheory/Proofs/DivisorProofs/CoefMapNormEquivariantBasisOfFive.lean}{GlobalTheory.GalOne.Equivar.eqAt_of_frAtEmb_ofFive}
\leanref{Development/AbelianVarietyTheory/Proofs/DivisorProofs/CoefMapNormEquivariantBasisOfFive.lean}{GlobalTheory.GalOne.IoCmp.ioAt_of_conj_ofFive}
\leanref{Development/AbelianVarietyTheory/Proofs/DivisorProofs/CoefMapNormEquivariantBasisOfFive.lean}{GlobalTheory.GalOne.identity_of_frConst_ofFive}
\begin{proof}
(1) First let $Q=1$, so that $\phi_Q=\phi$, a level vector with a vacuum $f_0$ for the trivial local
level, the set $P_g$, no ramified places and the trivial characters. By the core conditions on the level
characters, $\xi_{\mu_0}=1$ and $Q_{\mu_0}=\emptyset$; the lattices of the adelic matrices of $V_i$ and
$W_i$ are self-dual; and the Hecke character $\eta$ of $R$ is unramified at every finite place,
restricts to the quadratic character of $E_h/F_h$ on $\mathbf A_{F_h}^\times$ and has archimedean
exponents $\frac12$. So Lemma~\ref{is-eigenvalue-conjugate-frame-level-vector} applies to $S=V_i$,
$T=W_i$, the adelic matrices of $R$, the Weil datum $\omega$ with characters $(\chi',\chi'\eta'^{-1})$
and the level vector $\phi$, with $P_g$ and $Q=\emptyset$. Let $e$ be its idempotent, the sum of the
idempotents at the places of $P_g$.

The characters $\chi'$ and $\chi'\eta'^{-1}$ are trivial on every idele $x$ that has archimedean
component $1$, is a unit at every finite place and has component $1$ at the places of $P_g$ and their
conjugates. Indeed, by the core conditions on the level characters, $\mu^0_i=\mu^0_i\xi_{\mu_0}$ is
trivial on $x$, as it is unramified outside these places. As $\chi'$ is a $\sigma$-conjugate of
$\mu^0_i$, there is $y\in K_s$ with $\mu^0_i(x)=\sigma_0(y)$ and $\chi'(x)=\sigma(y)$; so $y=1$ and
$\chi'(x)=1$. And $\eta'$ is unramified everywhere, so $\eta'(x)=1$.

Let $k$ be as in (1) and $r=\operatorname{rk}W_i$. As $k_f\in J$, $k$ has archimedean component $1$ and
stabilizes the lattice of $V_i$, and for $P\in P_g$ the reduction $\bar k_P$ is block upper triangular for
the blocks $(r,\operatorname{rk}V_i-r)$ with upper left block $1$. So $(k,1)$ is a stabilizer pair that
satisfies the conditions of Lemma~\ref{is-eigenvalue-conjugate-frame-vacuum}(2), with $\bar
k^{\mathrm{ul}}_P=1=\bar 1_P$. As $Q$ is empty, Lemma~\ref{is-eigenvalue-conjugate-frame-level-vector}
gives
\[
 \omega(k,1)\phi=(\chi'\eta'^{-1})(\det_ek)^{-1}\chi'(\det_e1)\,\phi=\chi'(\det_ek)^{-1}\eta'(\det_ek)\,\phi.
\]
The idele $\det_ek$ has archimedean component $1$ and is a unit at every finite place, as $k_f$
stabilizes a lattice; so $\eta'(\det_ek)=1$. By
Lemma~\ref{is-eigenvalue-conjugate-frame-level-character}, with $t=\mu_0$,
$\mu^0_i(\det_ek)=\chi_J(k_f)$, the value of the level character of $R$. Let $y\in K_s$ with
$\mu^0_i(\det_ek)=\sigma_0(y)$ and $\chi'(\det_ek)=\sigma(y)$, and $y'\in L$ with
$\chi_J(k_f)=\sigma_0(y')$ and $\chi'_J(k_f)=\sigma(y')$. Then $\sigma_0(y)=\sigma_0(y')$, so $y=y'$ and
$\chi'(\det_ek)=\chi'_J(k_f)$. Hence $\omega(k,1)\phi=\chi'_J(k_f)^{-1}\phi$.

Now let $Q$ be any polynomial. The components of $k$ and of $1$ at the archimedean places are $1$, and
$\phi$ is a theta function (Lemma~\ref{is-eigenvalue-conjugate-level-data}(1)). So
Lemma~\ref{is-eigenvalue-conjugate-theta-facts}(4), applied to $(g,h)=(k,1)$ with the same Cartan
involutions and bases, gives $\omega(k,1)\phi_Q=(\omega(k,1)\phi)\otimes_BQ=\chi'_J(k_f)^{-1}\phi_Q$.

(2) The data of (2) are those of Situation~\ref{is-eigenvalue-conjugate-frame-core}, whose notation we
use; let $f_0$ be a vacuum of $\phi$. For $g\in U(W_i)(\mathbf A)$ write $x^g=x^{gZ_0}$, the family
$(g_\rho x_{0,\rho})_\rho$ indexed by $\mathcal T$, with $g_\rho$ the image of $g$ under $\theta_\rho$. For
$w\in\mathbb W$ with finite component $0$ and a family $x=(x_\rho)_{\rho\in\mathcal T}$ of matrices in
$M_r(\mathbf C)$ put $z_x(w)=(z_x(u_1),\dots,z_x(u_n))$, with $u(w)=(u_1,\dots,u_n)$ as in
Situation~\ref{is-eigenvalue-conjugate-frame-situation}(2). A
\emph{family of frame constants} is a function $a$ on the components with $|a|=1$ with the following
property. Let $h\in U(W_i)(\mathbf A)$ and $\gamma\in U(W_i)(F_h)$ with $x(h)=x^{\gamma h}$
(Definition~\ref{is-eigenvalue-conjugate-dictionary}(5)(b)) and $x(h)\in\mathcal D_{c(h)}$, and let a
symplectic basis of $\Lambda_{c(h)}$, an indexing and a bijection onto the letters be given, the last two
matched for $s$ (Definition~\ref{is-eigenvalue-conjugate-frame-fibres}), such that $\psi$, $\theta$,
$S=V_i$, $T=W_i$, $X=X_{c(h)}$, the basis and the indexing satisfy the conditions of
Situation~\ref{is-eigenvalue-conjugate-frame-situation}. Then for every
analytic fibre at $y(c(h),x(h))$ for these data, $x(h)$ and $\Lambda_{c(h)}$, every theta reading
$\vartheta$ of the section $s_{y(c(h),x(h))}$ for it, and every $w\in\mathbb W$ with finite component $0$,
\[
 \big(\omega(1,\gamma h)f_0\big)(w)=a(c(h))\,\mu(h)\,\mathcal G_{x(h)}(z)\prod_{k=1}^n\vartheta(z_k),
 \qquad z=z_{x(h)}(w),
\]
with the Gaussian $\mathcal G_{x(h)}$ for the given symplectic basis and indexing. The proof of (2)
combines four facts.
\begin{enumerate}
\item[(i)] Families of frame constants exist.
\item[(ii)] Every family of frame constants satisfies (a).
\item[(iii)] There is $\kappa_1\in\mathbf C$ with $|\kappa_1|=1$ with the following property. Let $x''$ be a
  rational coefficient vector of $V_i$ over $L$, and let $P=\sigma'(\mathrm{io}_O(\iota(x'')))$ be the image
  under $\sigma'$ of the polynomial in the letter variables that the one-map data at $i$ attach to
  $\iota(x'')$. Then there is a polynomial $Q$ in the creation variables $(\rho,a,k)$ such that:
  $\phi_{\mathrm{io}(v_0(i,\sigma)\otimes c_1(x''))}=\phi\otimes Q$, formed with the Cartan involutions $1$ of
  $V_i$ and $s_{Z_0,\rho}$ of $W_i$ and the bases $1$ and $Z_{0,\rho}$
  (Lemma~\ref{is-eigenvalue-conjugate-theta-facts}); and for every flag data $X$ of rank $r$ over $E_h$ with
  set $\mathcal T$ and every letter dictionary for $X$, $n$ and the weights of $\kappa_i$
  (Definition~\ref{is-eigenvalue-conjugate-frame-letter-dictionary}), with variables $Y_{\rho,a,k}$ and letter
  map $\lambda\mapsto Y(\lambda)$,
  \[
   \kappa_1\,Q\big((\rho,a,k)\mapsto(2/f_\rho)\,Y_{\rho,a,k}\big)=\ell_i\,P\big(\lambda\mapsto Y(\lambda)\big),
  \]
  where $f_\rho=2\sqrt{2\pi|\rho(\xi_0)|}$ and the arguments are the substitutions of the variables.
\item[(iv)] For every family $a$ of frame constants there are $c_0\in\mathbf C$ with $|c_0|=1$ and a
  central element $z$ of $G$ such that, for every $\kappa_1$ as in (iii) and all data as in (b), the
  identity of (b) holds with $\kappa_1c_0$ for $\kappa$.
\end{enumerate}
Let $a$ be a family of frame constants, by (i); it satisfies (a), by (ii). Let $c_0$ and $z$ be as in
(iv) for $a$, let $\kappa_1$ be as in (iii), and put $\kappa=\kappa_1c_0$. Then $|\kappa|=1$, and (iv)
gives (b).

\emph{Proof of (i).} We use the following facts.
\begin{itemize}
\item By Lemma~\ref{is-eigenvalue-conjugate-frame-points}(2) and~(5), $Z_0$ is a signature family for
  $T=W_i$ and $p_0=q$, the Cartan involutions of the model of $W_i$ at $\sigma$ are the $s_{Z_0,\rho}$,
  $\Lambda_V=\mathcal O_{E_h}^n$, and $\Lambda_V$ and $\Lambda_W$ are self-dual. So $f_0$, a Fock vacuum
  vector for the Cartan involutions $1$ of $V_i$ and those of the model and the lattice of the pair given
  by $\Lambda_V$ and $\Lambda_W$ (Lemma~\ref{is-eigenvalue-conjugate-level-data}(1)), is one for the
  Cartan involutions $1$ and $s_{Z_0}$ and the lattice of the pair given by $\mathcal O_{E_h}^n$ and
  $\Lambda_W$.
\item As in the proof of (1), $\eta$ is the restriction of a Hecke character of $E_h$ that is unramified
  at every finite place and has archimedean exponents $\frac12$. The spaces $V_i$ and $W_i$ of $D$ have
  correct signatures, and $\omega$ is a Weil datum for $(V_i,W_i)$ with characters
  $(\chi',\chi'\eta'^{-1})$. By Lemma~\ref{is-eigenvalue-conjugate-frame-points}(4), every $g$ has a
  determinant root and $\mu(g)=\chi'(y)$ for every determinant root $y$ of $g$, so that $|\mu|=1$;
  $\mu(\gamma g)=\mu(g)$ for $\gamma\in U(W_i)(F_h)$; and $\chi'$ and $\chi'\eta'^{-1}$ are trivial on the
  ideles at the places that do not split as in Lemma~\ref{is-eigenvalue-conjugate-frame-move}.
\end{itemize}

Let $c$ be a component. If $c\ne c(h)$ for every $h$, put $a(c)=1$; the condition on frame constants
concerns only the components $c(h)$. Otherwise choose $h_0$ with $c(h_0)=c$, and by
Lemma~\ref{is-eigenvalue-conjugate-frame-points}(1) a framing at $h_0$ with element $\gamma_0$, an
analytic fibre and a theta reading $\vartheta_0$ of mean one; put $g_0=\gamma_0h_0$, so that
$x(h_0)=x^{g_0}$. The data of the framing satisfy the conditions of
Situation~\ref{is-eigenvalue-conjugate-frame-situation} with $S=V_i$, $T=W_i$, $X=X_c$ and $p_0=q$, and
theta functions are formed at this framing. By Lemma~\ref{is-eigenvalue-conjugate-frame-signature}(3),
$g_0Z_0$ is a signature family and $s_{g_0Z_0,\rho}(g_0)_\rho=(g_0)_\rho s_{Z_0,\rho}$ for every $\rho$.
Let $L'_{g_0}$ be as in Lemma~\ref{is-eigenvalue-conjugate-frame-points}(3). By
Lemma~\ref{is-eigenvalue-conjugate-frame-moved-vacuum}, applied to $(1,g_0)$, $\omega(1,g_0)f_0$ is a Fock
vacuum vector for the Cartan involutions $1$ and $s_{g_0Z_0}$ and the lattice of the pair given by
$\mathcal O_{E_h}^n$ and $L'_{g_0}$. The module $\Lambda_c=L(g_0)=\mathfrak AL'_{g_0}$ is self-dual for
$X_c$, and $\vartheta_0$ is a level-one theta function for $w_0$, $x(h_0)$ and $\Lambda_c$
(Lemma~\ref{is-eigenvalue-conjugate-frame-level-one}, applied to the datum of the core at level $i+1$ and
$X=X_c$) of mean one. The element $w_0$ is an algebraic
integer with $w_0+\bar w_0=1$ and $w_0\bar w_0=(1+d)/4=2b$ with $b=(d+1)/8$, an integer as $d$ is $7$
modulo $16$. So Lemma~\ref{is-eigenvalue-conjugate-frame-unit}, applied to $g_0Z_0$, $L=\Lambda_c$,
$L_T=L'_{g_0}$, $f=\omega(1,g_0)f_0$ and $\vartheta_0$, gives $\lambda$ with $|\lambda|=1$ and
\[
 (\omega(1,g_0)f_0)(w)=\lambda\,\mathcal G_{x(h_0)}(z)\prod_k\vartheta_0(z_k),\qquad z=z_{x(h_0)}(w),
\]
for every $w$ with finite component $0$. By the laws of the weight tie at the level $i+1$ and the good
embedding $s$, the uniformization of the core at $s$ is as in
Situation~\ref{is-eigenvalue-conjugate-frame-uniformization} (proof of
Lemma~\ref{is-eigenvalue-conjugate-frame-points}(1)), the level of the core has full level $N=1$, and the
numbers $c_\tau=2|\tau(\xi_0)|$ of the flag data of every component do not depend on $\tau$. Let $\varphi$
be as in Lemma~\ref{is-eigenvalue-conjugate-frame-series} for $c$, $m=n$, which is even, and the data of the
framing at $h_0$. The point $x(h_0)=x^{g_0Z_0}$ lies in $\mathcal D_c$ and is orthonormal
(Lemma~\ref{is-eigenvalue-conjugate-frame-signature}(1)); so that lemma, with the data of the framing as
second data, gives
\[
 (\omega(1,g_0)f_0)(w)=\Theta_{\lambda\varphi}(x^{g_0};u(w))
\]
for every $w$ with finite component $0$; and $\omega(1,g_0)f_0$ is a theta function. Put
$a(c)=\lambda\,\mu(h_0)^{-1}$; then $|a(c)|=1$.

Now let $h$, $\gamma$ and the data be as in the definition of frame constants, with $c(h)=c$, an
analytic fibre and a theta reading $\vartheta$. By Lemma~\ref{is-eigenvalue-conjugate-frame-points}(1)
$h$ has a framing, with some element $\gamma_1$. Then $x^{\gamma h}=x(h)=x^{\gamma_1h}$, so
$(\gamma h)Z_0=(\gamma_1h)Z_0$, and $k=(\gamma h)^{-1}\gamma_1h$ satisfies
$kZ_0=(\gamma h)^{-1}\big((\gamma h)Z_0\big)=Z_0$ by
Lemma~\ref{is-eigenvalue-conjugate-frame-signature}(3); by the same lemma $k$ has archimedean component
$1$. Hence so does its conjugate
$\gamma^{-1}\gamma_1=hkh^{-1}$. This is a rational element; its matrix $\delta$ over $E_h$ satisfies
$\rho(\delta)=1$ for any $\rho$, so $\delta=1$ and $\gamma=\gamma_1$. By the framings at $h$ and
$h_0$, $L(\gamma h)=\Lambda_c=L(g_0)$. Let $L'_{\gamma h}$ and $L'_{g_0}$ be as in
Lemma~\ref{is-eigenvalue-conjugate-frame-points}(3); their products with $\mathfrak A$ span the same
lattice. As $\mathfrak A$ is invertible, they are equal: for $u'\in L'_{\gamma h}$,
$\beta\in\mathfrak A^{-1}$ and $\alpha\in\mathfrak A$, $\beta\alpha u'\in\mathfrak A^{-1}\mathfrak
AL'_{g_0}=L'_{g_0}$, and $1$ is a sum of such products $\beta\alpha$; by symmetry the two lattices are
equal. So $(\gamma h)_f\hat\Lambda_W=\hat L'_{g_0}=(g_0)_f\hat\Lambda_W$, and
Lemma~\ref{is-eigenvalue-conjugate-frame-move}, applied with the data of the framing at $h_0$, $Z=Z_0$,
$\Lambda_S=\mathcal O_{E_h}^n$, $\Lambda_T=\Lambda_W$, $f_0$, $\chi_S=\chi'$, $\mu$, $g_0$, $g_1=\gamma h$,
$L=L'_{g_0}$ and $\lambda\varphi$, gives
\[
 (\omega(1,\gamma h)f_0)(w)=\mu(\gamma h)\mu(g_0)^{-1}\lambda\,\Theta_\varphi(x(h);u(w)).
\]
The point $x(h)=x^{\gamma hZ_0}$ lies in $\mathcal D_c$ and is orthonormal
(Lemma~\ref{is-eigenvalue-conjugate-frame-signature}(1)), so
Lemma~\ref{is-eigenvalue-conjugate-frame-series}, with the given data as second data, gives
$\Theta_\varphi(x(h);u(w))=\mathcal G_{x(h)}(z)\prod_k\vartheta(z_k)$, $z=z_{x(h)}(w)$, with the Gaussian
and coordinates of the given basis and indexing. Finally $\mu(\gamma h)=\mu(h)$ and $\mu(g_0)=\mu(h_0)$, so
$(\omega(1,\gamma h)f_0)(w)=a(c)\,\mu(h)\,\mathcal G_{x(h)}(z)\prod_k\vartheta(z_k)$. Hence $a$ is a family
of frame constants.

\emph{Proof of (ii).} Let $a$ be a family of frame constants, $g\in G$ with a level lift $k$, read as an
element of $U(W_i)(\mathbf A)$ with archimedean component $1$, $f$ a core form and $h\in U(W_i)(\mathbf A)$.
For a component $c$ and $x\in\mathcal D_c$, the reading $\mathrm{rd}(f)_c(x)$ is the value of $f$ at the
framed point $y(c,x)$.

\emph{Reduction.} By Lemma~\ref{is-eigenvalue-conjugate-frame-marking-laws}(2), $gf$ is a core form and
there is $b\in K_s$ with $\sigma_0(b)=\prod_{P\in P_g}\chi_{i+1,P}(\det g_P)$ and $s(b)\mu(k)=1$. As $\mu$ is
multiplicative (Lemma~\ref{is-eigenvalue-conjugate-frame-points}(4)),
$s(b)\,\mathrm{ad}_a(\mathrm{rd}(f))(hk)=s(b)\mu(k)\,a(c(hk))\,\mu(h)\,\mathrm{rd}(f)_{c(hk)}(x(hk))
=a(c(hk))\,\mu(h)\,\mathrm{rd}(f)_{c(hk)}(x(hk))$, and $\mu(h)\ne0$. So (a) for $g$, $f$ and $h$ follows from
\[
 (*)\qquad a(c(h))\,\mathrm{rd}(gf)_{c(h)}(x(h))=a(c(hk))\,\mathrm{rd}(f)_{c(hk)}(x(hk)).
\]

\emph{The framed points.} Put $c=c(h)$, $x=x(h)$, $c'=c(hk)$ and $x'=x(hk)$. By
Lemma~\ref{is-eigenvalue-conjugate-frame-points}(1), $h$ and $hk$ have framings with elements $\gamma$ and
$\gamma'$ that have the marking property there; let $A_0$ and $A'$ be their analytic fibres at $y(c,x)$ and
$y(c',x')$, and $\vartheta$ and $\vartheta'$ their theta readings of mean one. By
Lemma~\ref{is-eigenvalue-conjugate-frame-marking-laws}(1), for the level of the core, there are markings $m$
at $\gamma h$ with values in $L(\gamma h)=\Lambda_c$ and $m'$ at $\gamma'hk$ with values in $\Lambda_{c'}$; so
$A_0$ is marked by $m$ and $A'$ by $m'$.

\emph{The moved point.} Let $y_g$ be the framed point obtained from $y(c,x)$ by the change of level
structure through $\ell_g$. By the definition of the action of $G$
(Lemma~\ref{is-eigenvalue-conjugate-class-action}), $\mathrm{rd}(gf)_c(x)=(gf)(y(c,x))=f(y_g)$. By
Lemma~\ref{is-eigenvalue-conjugate-frame-level-move}, applied to $\ell=\ell_g$, $y=y(c,x)$ and $y'=y_g$, there
is an analytic fibre $A_1$ at $y_g$ for the data of the framing at $h$, $x$ and $\Lambda_c$, for which
$\vartheta$ is a theta reading of $s_{y_g}$, marked by $m\circ\ell_g$.

\emph{The comparison.} Put $\delta=\gamma\gamma'^{-1}$, a rational element, and $a'=\gamma'hk$, so that
$\delta a'=\gamma hk$. By Lemma~\ref{is-eigenvalue-conjugate-frame-marking-laws}(2), $L(\gamma hk)=L(\gamma
h)=\Lambda_c$ and $m\circ\ell_g$ is a marking at $\gamma hk$. As $k$ has archimedean component $1$,
$x^{\gamma hkZ_0}=x^{\gamma hZ_0}=x$ (Lemma~\ref{is-eigenvalue-conjugate-frame-signature}(3)). The integers
$p_0$ of $X_c$ and $X_{c'}$ are $q$, by the framings. So
Lemma~\ref{is-eigenvalue-conjugate-frame-compare} applies to $\delta$, $a'$, $c$, $c'$, $x$, $x'$, the data
of the two framings, the points $y_g$ and $y(c',x')$, the fibres $A_1$ and $A'$, the readings $\vartheta$
and $\vartheta'$, and the markings $m\circ\ell_g$ and $m'$; let $\lambda$ be as there, and $n=\operatorname{rk}V_i$.

\emph{The constants.} By Lemma~\ref{is-eigenvalue-conjugate-frame-det-roots}(2), $k$ has a determinant root
$u$ with archimedean component $1$ that is a unit at every finite place. The facts listed at the
beginning of the proof of (i) let Lemma~\ref{is-eigenvalue-conjugate-frame-finite-vacuum} apply to $(1,k)$,
with the determinant roots $1$ and $u$, and give
$\omega(1,k)f_0=(\chi'\eta'^{-1})(1)^{-1}\chi'(u)f_0=\mu(k)f_0$, as $\mu(k)=\chi'(u)$
(Lemma~\ref{is-eigenvalue-conjugate-frame-points}(4)). As $\vartheta$ has mean one, it is not identically
$0$, and $u\mapsto z_x(u)$ is onto; so there is $w$ with finite component $0$ and
$\mathcal G_x(z)\prod_k\vartheta(z_k)\ne0$, $z=z_x(u(w))$. Let $w'$ be as in
Lemma~\ref{is-eigenvalue-conjugate-frame-compare}(2), and $z'=z'_{x'}(u(w'))$. The frame constants at $h$,
with $\gamma$, $A_0$ and $\vartheta$, and at $hk$, with $\gamma'$, $A'$ and $\vartheta'$, give
\[
 (\omega(1,\gamma h)f_0)(w)=a(c)\mu(h)\,\mathcal G_x(z)\prod_k\vartheta(z_k),\qquad
 (\omega(1,\gamma'hk)f_0)(w')=a(c')\mu(hk)\,\mathcal G'_{x'}(z')\prod_k\vartheta'(z'_k).
\]
By Lemma~\ref{is-eigenvalue-conjugate-frame-compare}(2), applied to the theta function $f_0$,
$(\omega(1,\gamma hk)f_0)(w)=(\omega(1,\gamma'hk)f_0)(w')$. The left side is
$(\omega(1,\gamma h)\omega(1,k)f_0)(w)=\mu(k)\,a(c)\mu(h)\,\mathcal G_x(z)\prod_k\vartheta(z_k)$, and the right
side is $a(c')\mu(hk)\lambda^n\,\mathcal G_x(z)\prod_k\vartheta(z_k)$. As $\mu(hk)=\mu(h)\mu(k)\ne0$, this gives
$a(c)=a(c')\lambda^n$.

\emph{Conclusion.} The weight $\kappa_i$ of the core forms has the exponent $n=\operatorname{rk}V_i$ at the
line bundle of the polarization, so Lemma~\ref{is-eigenvalue-conjugate-frame-compare}(1) gives
$f(y(c',x'))=\lambda^nf(y_g)$. Hence
$a(c)\,\mathrm{rd}(gf)_c(x)=a(c')\lambda^nf(y_g)=a(c')f(y(c',x'))=a(c')\,\mathrm{rd}(f)_{c'}(x')$, which is
$(*)$.
% (ii) is GalOne.Equivar.eqAt_of_frAtEmb_ofFive (100 lines with its statement): GeoLevel.eqAt_iff (mu_lift,
% actLev_mem_cuspA and castDict_mu_mul) for the reduction; framed_point_at_ofFive at h and h * k (with its
% marking clause p9, q9) and markAd_laws for mk and mk'; mvPt_spec and evalAt_actL for the moved point (hev);
% Dict.FibreMove.anFib_levR_id for A1; compare with delta = gamma gamma'^-1 and a' = gamma' (h k), using
% Lift.lift_laws_ofFive (l1, l2, l4) and oneg_frames for the frame; exists_finUnit_isDetRoot and
% weil_vacuum_finite_selfDualLattice_splitting with hmu (f11.2) for omega (1, k) f0; exists_real_ne_zero for w;
% the frame condition hcc at h and at h * k (e1, e2); z2 for w' (e3); castDict_mu_mul (e5, e7); z1 with f4
% ((kappa i).n = rank V_i) for e6; and rdU_eq for the readings.
\emph{Proof of (iii).} Let $v=v_0(i,\sigma)$, and let $s_\rho$ and $f_\rho$ be the Cartan involution and the
frame of the model of $W_i$ at $\sigma$, with the notation $\phi\otimes_{B'}$ and $\phi\otimes_{1,f^{-1}}$ of
Lemmas~\ref{is-eigenvalue-conjugate-frame-rebase} and~\ref{is-eigenvalue-conjugate-frame-phase}; so
$\phi_{\mathrm{io}(\cdot)}=\phi\otimes_B\mathrm{io}(\cdot)$. Let $(p,q)$ be the signature of the datum of the core
at level $i+1$.

\emph{The phase.} By Lemma~\ref{is-eigenvalue-conjugate-frame-rebase} there is a map $\mathrm{io}'$ as in
Definition~\ref{is-eigenvalue-lift-data}(7) for $\omega$, $\phi$, $(f_\rho^{-1})_\rho$ and $\Delta_i(\sigma)$ with
$\phi\otimes_B\mathrm{io}(v\otimes w)=\phi\otimes_{f^{-1}}\mathrm{io}'(v\otimes w)$ for every $w\in W_c$. By the
weight tie, $v$ is an extremal vector of the model, so
Lemma~\ref{is-eigenvalue-conjugate-frame-extremal-unit} gives $w_0\in W_c$ of norm $1$ with
$\mathrm{io}'(v\otimes w_0)=\Delta_i(\sigma)$. The linear map $\Lambda(w)=\mathrm{io}'(v\otimes w)$ satisfies
$\omega(u,1)(\phi\otimes_{f^{-1}}\Lambda(w))=\phi\otimes_{f^{-1}}\Lambda(uw)$ for $u$ with finite component $1$, by the
intertwining property of $\mathrm{io}'$. So Lemma~\ref{is-eigenvalue-conjugate-frame-phase} gives
$\epsilon$ with $|\epsilon|=1$ and
\[
 \phi\otimes_B\mathrm{io}(v\otimes c_1(x''))=\phi\otimes_{f^{-1}}\Lambda(c_1(x''))
 =\epsilon\cdot\phi\otimes_{1,f^{-1}}\mathrm{lt}\big(\sigma'(\mathrm{io}_O(\iota(x'')))\big)
\]
for every letter assignment $\mathrm{lt}$ and every $x''$. Put $\kappa_1=\epsilon^{-1}$.

\emph{A letter assignment.} Let $\lambda=((\tau',i'),k)$ be a letter variable. The multiplicity $m_{\tau'}$ is
$p$ if $\tau'$ lies in the CM type of the host of the core and $q$ otherwise, and as $s$ is good, $s\tau'$
lies in the CM type exactly when $\tau'$ does; by the pins of the core on its host, the CM type is $\mathcal T$,
and an embedding outside it is $\bar\rho$ for some $\rho\in\mathcal T$. Put $\rho(\lambda)=s\tau'$ and
$c(\lambda)=q+i'$ if $s\tau'\in\mathcal T$, where $i'\le p$, and $\rho(\lambda)=\overline{s\tau'}$ and
$c(\lambda)=q+1-i'$ otherwise, where $i'\le q$; so $1\le c(\lambda)\le r$, as $p+q=r$ by the pins. For $1\le c\le
r$ let $\pi(c)=r+1-c$ if $c\le q$ and $\pi(c)=c-q$ if $c>q$. Then
$\mathrm{lt}(\lambda)=(\rho(\lambda),\pi(c(\lambda)),k)$ is a letter assignment: $\pi(q+i')=i'$ and
$\pi(q+1-i')=p+i'$. Put $\mathrm{lt}'(\lambda)=(\rho(\lambda),c(\lambda),k)$.

\emph{The frames.} The frame $x_{0,\rho}=Z_{0,\rho}$ puts the last $q$ columns of $f_\rho^{-1}$ first, in the
reverse order (proof of Lemma~\ref{is-eigenvalue-conjugate-frame-points}(2)), so $Z_{0,\rho}e_c=f_\rho^{-1}e_{\pi(c)}$
for every $c$; and $s_{Z_0,\rho}=s_\rho$ (Lemma~\ref{is-eigenvalue-conjugate-frame-points}(2)). The creation
operator of the variable $(\rho,c,k)$ for the Cartan involutions $1$ and $s_\rho$ and the bases $1$ and $Z_0$ is
built from the matrix $u=Z_{0,\rho}E_{ck}$ (proof of Lemma~\ref{is-eigenvalue-conjugate-theta-facts}), which is
the matrix $f_\rho^{-1}E_{\pi(c)k}$ of the variable $(\rho,\pi(c),k)$ for the bases $1$ and $f^{-1}$; so the two
operators are equal, and $\phi\otimes_{1,f^{-1}}\mathrm{lt}(P')=\phi\otimes Q'$ for every polynomial $P'$ in the
letter variables, with
$Q'=\mathrm{lt}'(P')$ and $\phi\otimes Q'$ formed as in (iii). With
$Q=\epsilon\,\mathrm{lt}'(P)$ this gives the first assertion of (iii).

\emph{The letter map.} Let $X$ and a letter dictionary be as in (iii), with indexing $j\mapsto(\tau_j,a_j)$,
bijection $j\mapsto\ell_j$ and integer $p_0$; then $p_0+p=r$, so $p_0=q$. Let $\lambda=(\ell_j,k)$ with
$\ell_j=(\tau',i')$. By the matching, the row type of $j$ is $s\tau'$, and $p_x(v)_j$ is the period of $v$ at
$\ell_j$ for all $x$ and $v$. If $s\tau'\in\mathcal T$, the row type $\bar\tau_j$ of an index with $a_j\le p_0$
is not in $\mathcal T$; so $a_j>p_0$, $\tau_j=s\tau'$, and $(x_{\tau_j}^{-1}\tau_j(v))_{a_j}=(x_{\tau_j}^{-1}
\tau_j(v))_{p_0+i'}$ for all $x$ and $v$ gives $a_j=p_0+i'=c(\lambda)$. Otherwise $a_j\le p_0$, $\tau_j=
\overline{s\tau'}$, and $h_{\tau_j}(x_{\tau_j}e_{a_j},\tau_j(v))=h_{\tau_j}(x_{\tau_j}e_{p_0+1-i'},\tau_j(v))$ for
all $x$ and $v$ gives $a_j=p_0+1-i'=c(\lambda)$. So the letter map sends $\lambda$ to
$Y_{\rho(\lambda),c(\lambda),k}$.

\emph{The constant.} By the pins of the one-map data, $\mathrm{io}_O(\iota(x''))$ lies in the span of the
highest weight polynomials of the weights of $\kappa_i$, so $P$ is homogeneous of degree $N$. The number
$|\rho(\xi_0)|$ does not depend on $\rho$, so $2/f_\rho=(2\pi|\rho(\xi_0)|)^{-1/2}$ is the same for every $\rho$.
Hence, by the first assertion and the letter map,
\[
 \kappa_1\,Q\big((\rho,a,k)\mapsto(2/f_\rho)Y_{\rho,a,k}\big)
 =P\big(\lambda\mapsto(2/f_{\rho(\lambda)})Y_{\rho(\lambda),c(\lambda),k}\big)
 =(2\pi|\rho(\xi_0)|)^{-N/2}\,P\big(\lambda\mapsto Y(\lambda)\big)=\ell_i\,P\big(\lambda\mapsto Y(\lambda)\big).
\]
% (iii) is GalOne.IoCmp.ioAt_of_conj_ofFive (135 lines with its statement): IoAt kappa O sigma sigma' c1 f Lambda ph
% for Lambda w = ptAt R i sigma B f (io (tprod v) w) with v = v0 i sigma; Q is ka * rename lt' (map sigma' (O.io x'))
% and ph = kappa_1 = ka^-1. Its proof: factsAt_ofFive (hgood, hord, hcart), Rebase.rebase_ofFive for io',
% hext from (t.hv0 i sigma).1, IoExtConj.ioExt_conj_ofFive (hw), PhaseAt.io_eq_phase_smul_at_ofFive (ka);
% WeightExponents.weight_exponents (w1, w2), Pj.relab_rt_iff (hrl), the pins d1, PinNum.numerics_of_pins (hx, hpq),
% hsig and hneg (q negative signs), the choice hex of rho and c (0-based: c = q + i or q - 1 - i), lt' and
% lt = (rho, arr q c, k) with IsLtAt; the first clause is hio'.2, hio and IoCmp.pureTensorO_rename (zAt = x0Mat
% is onb with the columns arr q); the second is letter_relab (heL, from dir_eq, ty_eq, per_eq,
% coord_eq_of_per_mem and coord_eq_of_per_conj), isHomogeneous_of_span with the field i3 of hO.io (the range of
% O.io is the span of HHa), two_div_fockScale, aeval_scale and lieFac_eq_pow (ell_i = lieFac with xiR D =
% |phi0 xi0|). The book's letter map is Dc.eL read on rows through rowIdx.
\emph{Proof of (iv).} Let $a$ be a family of frame constants. The delta data of the one-map data at $i$ include
a second level $\ell_2$ of the datum of the core at level $i+1$, whose integer $M_2$ is a multiple of $M$, and
for each class $k$ a function $w_k$ with values in $L'$ on the $n$-tuples of classes of the lattice of the datum
modulo $\ell_2$. We use the following facts.
\begin{enumerate}
\item[($\alpha$)] There is $\epsilon=\pm1$ such that, for every $g'\in U(W_i)(\mathbf A)$, the functions
  $\omega(h_k,g')\phi$ are the theta functions, at the frames $x^{g'Z_0}$, of Schwartz-Bruhat functions built
  from a vacuum, the values $w_k$ and roots of $\epsilon$ times a marking at $g'$ for $\ell_2$, whenever
  $\omega(1,g')f_0$ is the theta function there of a multiple of the vacuum.
\item[($\beta$)] There is $c_0$ as in Lemma~\ref{is-eigenvalue-conjugate-frame-cover-constant}.
\item[($\gamma$)] There is an integer $m$ prime to $M_2$ with $\sigma'(w_k(t))=w_k(mt)$ for every $k$ and every
  tuple $t$.
\item[($\delta$)] For an integer prime to $M$ there is a central element of $G$ as in
  Lemma~\ref{is-eigenvalue-conjugate-frame-central}.
\end{enumerate}
Let $\epsilon$, $c_0$ and $m$ be as in ($\alpha$) to ($\gamma$), and $z$ as in ($\delta$) for $m\epsilon$, which is
prime to $M$ as $M$ divides $M_2$. Let $\kappa_1$ be as in (iii), the data as in (b), and $g\in U(W_i)(\mathbf A)$. By
Lemma~\ref{is-eigenvalue-conjugate-frame-points}(1), $g$ has a framing with element $\gamma$, whose basis,
indexing and bijection are those of a letter dictionary, with an analytic fibre $A_0$ at $y=y(c(g),x(g))$ and a
theta reading $\vartheta$ of $s_y$ of mean one; put $X=X_{c(g)}$ and $\Lambda=\Lambda_{c(g)}=L(\gamma g)$. As $a$ is
a family of frame constants, $(\omega(1,\gamma g)f_0)(w)=a(c(g))\mu(g)\,\mathcal G_{x(g)}(\zeta)\prod_k
\vartheta(\zeta_k)$, $\zeta=z_{x(g)}(u(w))$, for $w$ with finite component $0$. By
Lemma~\ref{is-eigenvalue-conjugate-frame-vacua}(1), with $\Lambda_S=\mathcal O_{E_h}^n\in\mathcal V_1$, the
tensor lattice $\mathcal O_{E_h}^n\otimes\Lambda$ has a vacuum $v$ for $w_0$. By
Lemma~\ref{is-eigenvalue-conjugate-frame-level-one}, $\vartheta$ is a level-one theta function for $w_0$,
$x(g)$ and $\Lambda$, so by Lemma~\ref{is-eigenvalue-conjugate-frame-theta-product} there is $\kappa'$ with
$\prod_k\vartheta(\zeta_k)=\kappa'\theta_v(x(g);\zeta)$ for every $\zeta$. Hence $\omega(1,\gamma g)f_0$, a theta
function, has the value
$\Theta_{a(c(g))\mu(g)\kappa'v}(x(g);u(w))$ at every $w$ with finite component $0$. By
Lemma~\ref{is-eigenvalue-conjugate-frame-marking-laws}(1) there are markings $m_1$ at $\gamma g$ for the level of
the core and $m_2$ at $\gamma g$ for $\ell_2$, with values in $\Lambda$; $A_0$ is marked by $m_1$
(Lemma~\ref{is-eigenvalue-conjugate-frame-points}(1)). For every class $k$, with the integer $N_k\ge1$ that the delta
data attach to it, Lemma~\ref{is-eigenvalue-conjugate-frame-inputs-roots}(1) gives a family $Y_k$ of roots of $\epsilon
m_2$ for $\ell_2$, $\Lambda$ and $N_k$, as $\epsilon m_2$ has values in $\Lambda$; as $m$ is prime to $M_2$, by (3) of
that lemma $q\mapsto Y_k(mq)$ is a family of roots of $m\epsilon m_2$ for $\ell_2$, $\Lambda$ and $N_k$.

With these data the two sides of (b) at $g$ are multiples of one class sum $C$, formed at $x(g)$ from the
polynomials $\mathrm{io}_O(x'_j)$, the vacuum $v$, the values $w_k$ and the roots $Y_k(mq)$ of $m\epsilon m_2$:
\[
 \kappa_1\sum_j\sigma'(\mathrm{vol}\,\gamma_j^{-1})\,\theta_\omega\big(h_j,g;\phi_{\mathrm{io}(v_0(i,\sigma)\otimes
 c_1(x''_j))}\big)=a(c(g))\mu(g)\kappa'\,\ell_i\,C,\qquad (zf')(y)=c_0\kappa'C.
\]
As $zf'=\sum_ka_kf_k$ after extension of scalars, $\sum_k\sigma'(a_k)\,\mathrm{rd}(f_k)_{c(g)}(x(g))=(zf')(y)=
c_0\kappa'C$. Hence
\[
 \kappa_1c_0\sum_j\sigma'(\mathrm{vol}\,\gamma_j^{-1})\,\theta_\omega(\cdots)=a(c(g))\mu(g)\,\ell_i\,c_0\kappa'C
 =\ell_i\sum_k\sigma'(a_k)\,a(c(g))\,\mu(g)\,\mathrm{rd}(f_k)_{c(g)}(x(g))=\ell_i\sum_k\sigma'(a_k)\,
 \mathrm{ad}_a(\mathrm{rd}(f_k))(g),
\]
which is the identity of (b) with $\kappa_1c_0$ for $\kappa$.

\emph{Proof of ($\gamma$).} The coefficient field $L'$ of the one-map data is a subfield of $\mathbf C$; write $\mathfrak M_2$
for the ideal of $\ell_2$ and $\langle\ ,\ \rangle_{\mathfrak L}$ for the form of the lattice $\mathfrak L$ of the datum. By the
pins of the one-map data at $i$, the delta data have a function $\varphi_\delta$ with values in $L'$ on the $n$-tuples of
classes of $\mathfrak L$ modulo $\bar{\mathfrak M}_2\mathfrak L$, and $w_k$ is its finite Fourier transform for every class
$k$:
\[
 w_k(t)=\sum_x\varphi_\delta(x)\,\zeta^{\langle x,t\rangle},\qquad \langle x,t\rangle=\sum_{b=1}^n\langle x_b,e\,t_b\rangle_{\mathfrak L},
\]
where $\zeta\in L'$ is a primitive $M_2$-th root of unity, $e\in\bar{\mathfrak M}_2$ is an element with $e-1\in\mathfrak M_2$, and
the exponent is formed with representatives of the classes; it is well defined modulo $M_2$, as $e\mathfrak M_2\subset
\bar{\mathfrak M}_2\mathfrak M_2=M_2\mathcal O_{E_h}$ and $\langle\ ,\ \rangle_{\mathfrak L}$ takes integer values. By the pins of the one-map
data at the base tuple, $\varphi_\delta(x)$ is a product of a factor that is $0$ or $1$, given by the reductions of a matrix
at the places of $P_g$, and of a product of character values over the ramified places of the base tuple; there are no
such places, by the core conditions on the level characters. So $\varphi_\delta$ takes the values $0$ and $1$. The
complex numbers $\zeta$ and $\sigma'(\zeta)$ are primitive $M_2$-th roots of unity, so $\sigma'(\zeta)=\zeta^m$ for an integer
$m$ prime to $M_2$. As $mt_b$ represents the class of $mt_b$, $\langle x,mt\rangle=m\langle x,t\rangle$ modulo $M_2$, and
\[
 \sigma'(w_k(t))=\sum_x\varphi_\delta(x)\,\zeta^{m\langle x,t\rangle}=\sum_x\varphi_\delta(x)\,\zeta^{\langle x,mt\rangle}=w_k(mt)
\]
for every $k$ and $t$.
% conj_coeff_eq_dilate_ofFive (56 lines), followed: hQ (R.Q mu0 = empty, chars_of_lift_ofFive), hq and hk (omegaQ over
% the empty set is 1, omegaK is 0 or 1), hphi from the tuple pin phi (PhiIs), GalOne.map_phiHat
% (CoprimeScalarForEmbedding.lean: m from IsPrimitiveRoot.eq_pow_of_pow_eq_one for zetaM, and pairL_levRep_zsmul) for
% the embedding sigma' o eF^-1 of the field c2 of the delta data, and the pin w (WPin: eF (w k q) = phiHat (eF o phi) q).
% phiHat c phi q is the finsum over x of phi x * zetaM c ^ pairL (qRep x) (levRep q), with pairL x0 x = sum_b EF (x0 b)
% (cSep l * x b), cSep in Mfr l with cSep - 1 in l.Mb.
\openstep{The class term identity ($\alpha$), which writes each function $\omega(h_k,g')\phi$ as the theta function of a
Schwartz-Bruhat function built from a vacuum, the values $w_k$ and roots of $\epsilon$ times a marking, is not written
out.}
% Size of (alpha), Term.zak_class_term_conj_ofFive, following imports and stopping at statements the book cites:
% about 21,000 lines (17,700 in 1,290 theorems), about 16,400 of them below neither of the two values.
\openstep{The first value, which computes the conjugated theta sum at $g$ as $a(c(g))\mu(g)\kappa'\ell_i$ times the class
sum $C$ from ($\alpha$), Fact~(iii) and the dilation by $m$ of the roots, is not written out.}
% Size of the first value, adelic_value_ofFive (59 lines; its proof uses OneMapRef.kernel_value, ShortRead.theta_smul2,
% the clause of (alpha) at h k, and Lemma is-eigenvalue-conjugate-frame-inputs-roots(3) for the dilation): about 4,600
% lines (2,400 in 200 theorems), about 240 of them below neither (alpha) nor the second value.
\openstep{The second value, which computes the moved section $zf'$ at the framed point $y$ as $c_0\kappa'$ times the class
sum $C$ through the level move by $z$ and a marked lift to the second level, is not written out.}
% Size of the second value, geometric_value_ofFive (61 lines; its proof uses evR_eq_clsSum, MarkLift.exists_lift_marked,
% T.analyticFibre_description_levelMove_marking_comp and the constant c0 of (beta)): about 33,200 lines (28,700 in
% 1,920 theorems), about 31,200 of them below neither (alpha) nor the first value.
% (iv) is GalOne.identity_of_frConst_ofFive (83 lines with its statement), followed above: (alpha) is
% Term.zak_class_term_conj_ofFive (97 lines; TermAt O sigma sigma' omega f0 f eps, for RepsR, IsSet, p0 = q, SelfD,
% latW L = latG g', IsVac for tens stdL L, Zak at frOf (mvZ g' zAt) of a v, MarkAd at l2 for mk2 with values in L
% and IsRoots l2 L (eps mk2) (Nc k) Y: Zak of a sigma'(ep k (mIdx)^-1 (card)^-1 zCard) gphi X bas w0 (La k) L v (w k) Y
% for omega (h k, g') f); (beta) is cover_const_at_ofFive (Lemma -frame-cover-constant); (gamma) is
% conj_coeff_eq_dilate_ofFive (56 lines; sigma' (O.w k q) = O.w k (m q) with m coprime to l2.M); (delta) is
% GalCentral.central_lev_ofFive (Lemma -frame-central) for m * eps, coprime to T.Q.l.M as l.M | l2.M
% ((hP.cov ()).1.1.1). At k0 = g: framed_point_at_ofFive (gamma, the Dic Dc, A0, th, MeanOne, the marking
% clause k5), the frame condition hFr, TensVac.vacua with stdL_mem_Vset, isTheta1_of_reads_generator,
% VacProd.prod_eq_mul_thA (kappa'), hzak; markAd_laws at T.Q.l (mk) and at l2 (mk2); the roots Y from the last
% clause of DualMark.dual_marking, which is DualMark.exists_roots (Lemma -frame-inputs-roots(1); the Lean passes
% isMark_sign of IsMark from markAd_laws and a primitive root of unity of order M2, which that clause does not use),
% and hY' from T.marking_roots_dilation_coprime_level (Lemma -frame-inputs-roots(3)). The two values: adelic_value_ofFive (ha: ph * sum = a * lieFac * clsSum
% with the roots Y j (m q)) and geometric_value_ofFive (hg: evR of actLev z f' at the point = c0 * kappa' * clsSum),
% joined by the base change hbc (IsBC) at the framed point (e1); the last step is the algebra alg.
% What the open step needs, from the Lean: (alpha) rests on OneMapRef.FinDatum.thetaVector_eq_sum_translates_of_core_five
% (78 lines), GalOne.vec_eq_mean_translates_ofFive (79), ClassConj.class_scalar_conj_ofFive (93),
% T.splitFinitePoint_moves_vacuum_to_projection, OneMapRef.Term.heis_root_zak (50), Term.zak_sum (30) and
% Const.isSB_gphi. Dict.gphi X bas w La L v W Y is the sum over q of W q times heisG (Y q) applied to projA (tens La L) v,
% the projection of Lemma -frame-vacua(3). adelic_value_ofFive (59 lines) uses, besides (alpha) and (iii),
% OneMapRef.kernel_value (theta kernel of fv f Q at (h, k0) = thP of the Zak function at the moved frame, through
% om_pureTensorO, T.adelicTheta_of_testFunction_exists_laws_fockVectors and om_aut), ShortRead.theta_smul2 and
% Lemma -frame-inputs-roots(3) for the dilation by m. ClsVal.clsSum weights thP of the class functions gphi by
% sigma'(vol gam^-1 ep mIdx^-1 card^-1 zCard), where GenLat.mIdx and zCard are the r-th powers of the indices of
% N times the standard lattice in Ln and in La of the class lattice. geometric_value_ofFive (61 lines) uses
% evR_eq_clsSum, MarkLift.exists_lift_marked and T.analyticFibre_description_levelMove_marking_comp (compare Lemma
% -frame-level-move).
% (1) is GalOne.omega_level_conjVec_ofFive for Q = 1, extended to every Q by
% FiniteFock.om_pureTensorO_of_finite (as in theta_eq_classSum_of_conjugate, hfin). In
% omega_level_conjVec_ofFive: chars_of_lift_ofFive gives IsEtaTilde R.eta, R.xi mu0 = 1 and R.Q mu0 = empty;
% VecInv.om_vec (Lemma is-eigenvalue-conjugate-frame-level-vector) with the self-dual lattices and IsFrame
% from hR.lev.betaV and betaW; v1 (UnrOff of mu') from VecInv.unrOff_mu and the values clause of ChiAt; v2
% from eta' unramified; the block conditions from IsLevelLift ... 1 k (parLev_g, blockTri_blkOf_le, rTup =
% 1); hinf and Vac.finPart_detAt for n1; m0 is mu_detAt_eq_levChi with levFn_flag1
% (Lemma is-eigenvalue-conjugate-frame-level-character); m1 compares the two ConjVal witnesses through the
% injectivity of embLL; smul_of_splitFac unfolds splitFacO mu' (mu' eta'^-1) x 1 = (mu' eta'^-1)(x)^-1 mu'(1).
% om_pureTensorO_of_finite is OrientedHeis.om_pureTensorO (Lemma is-eigenvalue-conjugate-theta-facts(4))
% with compAt_mt_of_finite for k in GfC.
% (2): the frame condition is FrAtEmb O sigma sigma' hV omega f0 cst = OneMapRef.FrConst (om (1, gamma * h) f0 w
% = cst (dict h).1 * mu h * (Ker.gaussR ... * prod_b th (Ker.zR ...)) at Ker.IsRealPt w, for every IsSet,
% Matched, dict h .2 = frOf (mvZ theta (gamma * h) (zAt R i sigma)), AnFib A and A.Reads 1 _ th). (i) is
% FrameConst.exists_frAtEmb_ofFive (OneMapRef.FrameConst.exists_frConst, 74 lines, with
% GalOne.framed_point_at_ofFive, factsAt_ofFive and facts_lat_ofFive); the proof of (i) above follows
% exists_frConst: the case split on IsPt at k, the base point (x0, h0, gamma0, A0, th0, MeanOne), hz, cc k =
% a / mu h0, then he (eT unique), frOf_injective, hfin and eq_one_of_rat for gamma = gamma1, le_of_isAA both
% ways for L0 = L1, om_vac_move, thR_smul, hall and hQ. A framing is OneMapRef.FrameConst.IsPt (with
% latW = latG), Definition is-eigenvalue-conjugate-frame-framing. The inputs of exists_frConst, assembled in
% exists_frAtEmb_ofFive, are Lemmas is-eigenvalue-conjugate-frame-points (framed_point_at_ofFive, factsAt_ofFive,
% facts_lat_ofFive, the muUnit law and isTheta1_of_reads_generator), -frame-moved-vacuum (isVacO_om),
% -frame-unit (vac_eq_unit_mul, with H.w0 through hostHalfRoot_add_conj_eq_one and
% isRelTraceOneIdempotentModTwo_hostElement), -frame-series (exists_const), -frame-move (om_vac_move) and
% -frame-signature (oneg_frames: the orthonormality hortho and the gamma = gamma1 step through mvZ_mul and
% fin_of_mvZ). The sizes of (ii) to (iv) are the line counts of the declarations named beside the open steps.
% The assembly with kappa = ph * c0 is that of oneMap_of_ext_ofFive.
% (ii) gives EqAt T R kappa O (sOf O sigma') hV cst, with phiAt = ad_a of rd and the product of R.chi (i+1)
% at the places of R.Pg; (iii) is used with Lambda w = ptAt (io (tprod v1) w). The hypotheses: ConjW R i sigma mu' eta' omega f0 f
% (OverK, ChiAt, IsEtaConj, IsWeilDatumO, IsThPVecOf f0 f, theta f), Pinned and FWTie; for (2)(b),
% CoefAt c c1 (the first five fields of IsConjForms), IsONBasis for B and IsIoW for io.
\end{proof}

\begin{lemma}\label{is-eigenvalue-conjugate-theta-sum}
Let $R$, $D$, $F$, $i$, $\sigma$, $\chi'$, $\eta'$, $\omega$, $\phi$, $B$, $J$ and $\chi'_J$ be as in
Lemma~\ref{is-eigenvalue-conjugate-frame-inputs}(1). Let $W_c$, $c_1$, $H$ and $j$ satisfy (2)--(4) of
Definition~\ref{is-eigenvalue-conjugate-datum} for a family tie of $R$, where $\chi'_J$ is the character
of the form space $H$, let $\mathrm{io}$ be a map as in Definition~\ref{is-eigenvalue-lift-data}(7) for
$\omega$, $\phi$, $B$ and $\Delta_i(\sigma)$, and let $\Theta$ be the theta integral of these data on
$H$ (Lemma~\ref{is-eigenvalue-conjugate-level-data}(3)), over the fundamental domain $\mathcal F$ of the
realization at $(i,\mu_0)$ for its measure $\nu$. Let $K'$ be a field with a ring homomorphism
$\sigma'\colon K'\to\mathbf C$, and let the following be given.
\begin{enumerate}
\item[(a)] A subgroup $J'\subset J$ of finite index, and $\mathrm{vol}\in K'$, the inverse of the
  index $[J:J']$.
\item[(b)] Elements $h_1,\dots,h_N\in U(V_i)(\mathbf A)$ with archimedean component $1$ whose finite
  parts represent the double cosets $U(V_i)(F_h)\backslash U(V_i)(\mathbf A_{F_h,f})/J'$, each once,
  and the orders $\gamma_k$ of the groups $\Gamma_k$ of $\gamma\in U(V_i)(F_h)$ with
  $h_k^{-1}\gamma h_k\in J'$.
\item[(c)] An integer $b\ne0$ that lies in every place of $P_g\cup Q_{\mu_0}$, such that each $h_k$
  has component $1$ outside finitely many places of $F_h$ that split in $E_h$ into places not
  containing $b$; and the lattices $\Lambda_k=V_i\cap h_k\hat L_0$, where $L_0$ is the lattice of
  $V_i$.
\end{enumerate}
The data and conditions of (c) are not used in the proof. Then for every $y\in H$, every tensor $v$ of
lowest $K$-type vectors and every $g\in U(W_i)(\mathbf A)$,
\[
 \Theta(y)(v)(g)=\sum_k\sigma'(\mathrm{vol}\,\gamma_k^{-1})\,\theta_\omega\big(h_k,g;\phi_{\mathrm{io}(v\otimes
 y(h_k))}\big),
\]
where $y(h_k)$ is the value of $y$ at the finite part of $h_k$, and
$\sigma'(\mathrm{vol}\,\gamma_k^{-1})=[J:J']^{-1}|\Gamma_k|^{-1}$.
\end{lemma}
\leanref{Development/AbelianVarietyTheory/Proofs/DivisorProofs/CoefMapNormEquivariantBasisOfFive.lean}{GlobalTheory.GalOne.theta_eq_classSum_of_conjugate}
\leanref{Development/AbelianVarietyTheory/Proofs/DivisorProofs/ConjugateWeilRepresentationThird.lean}{GlobalTheory.GalOne.integral_theta_chi}
\begin{proof}
The functions $y\in H$ satisfy $y(\gamma z)=\gamma\,y(z)$ for $\gamma\in U(V_i)(F_h)$ and
$y(zk)=\chi'_J(k)\,y(z)$ for $k\in J$. Put $\Lambda(w)=\phi_{\mathrm{io}(v\otimes w)}$ for $w\in W_c$, and
$\Phi(z)=\theta_\omega(z,g;\Lambda(\tilde y(z)))$ for $z\in U(V_i)(\mathbf A)$, with $\tilde
y(z)=z^{-1}y(z_f)$, so that $\Theta(y)(v)(g)=\int_{\mathcal F}\Phi\,d\nu$.

\emph{Step 1.} $\Phi$ is invariant on the left under $U(V_i)(F_h)$. For $\gamma\in U(V_i)(F_h)$, $\tilde
y(\gamma z)=z^{-1}\gamma^{-1}\gamma\,y(z_f)=\tilde y(z)$. Each $\Lambda(w)$ is a theta function, by
Lemma~\ref{is-eigenvalue-conjugate-level-data}(1), and $\omega$ is a Weil datum, so the theta series of
$\omega(\gamma z,g)\Lambda(w)$ is that of $\omega(z,g)\Lambda(w)$.

\emph{Step 2.} $\Phi$ is invariant on the right under every $u$ whose finite part lies in $J$. Write
$u=ka$ with $k$ of archimedean component $1$ and $k_f=u_f$, and $a$ of finite component $1$. The
representation of $U(V_i)(\mathbf A)$ on $W_c$ is trivial on $k$, so $\tilde
y(zu)=a^{-1}k^{-1}z^{-1}y(z_fk_f)=\chi'_J(k_f)\,a^{-1}\tilde y(z)$. By
Lemma~\ref{is-eigenvalue-conjugate-frame-inputs}(1), $\omega(k,1)\Lambda(w)=\chi'_J(k_f)^{-1}\Lambda(w)$,
and since $\phi_{\mathrm{io}(\cdot)}$ intertwines $\omega$ with the action of the archimedean points on
$W_c$, $\omega(a,1)\Lambda(w)=\Lambda(aw)$. Hence
\[
 \omega(zu,g)\Lambda(\tilde y(zu))=\omega(z,g)\,\omega(k,1)\,\omega(a,1)\Lambda\big(\chi'_J(k_f)\,a^{-1}\tilde
 y(z)\big)=\omega(z,g)\Lambda(\tilde y(z)),
\]
and $\Phi(zu)=\Phi(z)$.

\emph{Step 3.} By the core conditions on the models, $V_i$ is positive definite and $\nu$ is normalized
as in Lemma~\ref{is-eigenvalue-conjugate-unfolding} for the frame and the level datum of $V_i$ at
$\mu_0$; this level datum has the trivial local level, so the hypothesis of that lemma on the two
subgroups holds; and $V_i$ has correct signature (Definition~\ref{th-spaces}(7)).
Lemma~\ref{is-eigenvalue-conjugate-unfolding}, with (a) and (b), gives
$\int_{\mathcal F}\Phi\,d\nu=\sum_k[J:J']^{-1}|\Gamma_k|^{-1}\Phi(h_k)$. As $h_k$ has archimedean component
$1$, it acts trivially on $W_c$, so $\tilde y(h_k)=y(h_k)$, and
$\Phi(h_k)=\theta_\omega(h_k,g;\phi_{\mathrm{io}(v\otimes y(h_k))})$. Finally $\sigma'$ is a ring
homomorphism into $\mathbf C$, so $\sigma'(\mathrm{vol}\,\gamma_k^{-1})=\sigma'([J:J'])^{-1}\sigma'(\gamma_k)^{-1}
=[J:J']^{-1}|\Gamma_k|^{-1}$.
\end{proof}
% This is GalOne.theta_eq_classSum_of_conjugate, which takes RepsR R i mu0 h bad J' Lh gam vol over K'
% with sigma'. (a) is hle, hN and hvol; (b) is hy, (hk k).1.1 (h k in GfC, the first clause of
% SplitOffBad) and (hk k).2.2 (gam); (c) is bad /= 0, bad in the places of R.Pg and R.Q mu0, the rest of
% SplitOffBad (an idempotent sum e over places not containing bad with SuppAt e) and Lh k = hLat (fin (h k)),
% none of which the proof uses. The last equality is h3. The proof follows integral_theta_chi: Step 1 is ThetaIntegrand.ext_invariant and
% hom.om_aut; Step 2 is LiftUnfold.exists_fin_mul_real and LiftUnfold.integrand_right (form_right for y,
% hc.1 for the trivial action of k, hLk from hfin, hLa from the first clause of IsIoW); Step 3 is
% ClsInt.integral_eq_classSum_of_le with D.sp.V_sig, hR.models.vpos, the flag1L condition by Iff.rfl (datVt
% has LevelP.triv) and hR.models.nuV, and ext_fin. The equality of Theta with the
% integral is hL.2.2.2.eq (Lift.IsThetaP).

\begin{lemma}\label{is-eigenvalue-conjugate-inputs}
Let $R$ be a realization of a definite-side datum $D$ with a face $F$, with the notation of
Definitions~\ref{is-eigenvalue-lift-data} and~\ref{is-eigenvalue-conjugate-datum} and level lifts as
in Lemma~\ref{is-eigenvalue-core-group}. Let one-map data for $R$ be given
(Definition~\ref{is-one-map-data}), with a core whose coefficient field is $K_s\subset L$, weight
families $\kappa_j$ and vectors $v_0(j,\sigma)$; let $i$ be a lift, and fix a family tie of $R$.
Suppose that $R$ satisfies the core conditions at $F$, except possibly the one on its chain, that the
one-map data have a weight tie with a conjugation tie and are pinned at every lift, that the rational
structures $FW_j$ satisfy (11) of Definition~\ref{is-primitive-rows}, and that the one-map data at $i$
give the $i$-th lift as one map with the canonical constants of the core and some phase.

The \emph{core forms} are the cusp forms of weight $\kappa_i$ of the core at level $i+1$, a
$K_s$-vector space; the one-map data at $i$ let $G$ act on the core at level $i+1$, $f\mapsto gf$. The
weight tie gives a $K_s$-linear injection $f\mapsto y(f)$ of the core forms into $L^{n_W(i)}$ whose
image has a $K_s$-basis that is an $L$-basis of $L^{n_W(i)}$. Write $\chi_W$ for the character of $G$
of $FW_i$. For a good embedding $\sigma$ of $L$ (Lemma~\ref{is-eigenvalue-conjugate-good}), which is
the case if $\sigma$ agrees with $\sigma_0$ on $K$, the laws of the core in the weight tie make the
space of the core at level $i+1$ and the restriction of $\sigma$ to $K_s$ equal to $W_i$, and
$\mathrm{ad}_\sigma(f)$ is the function on $U(W_i)(\mathbf A)$ that the core attaches there to $f$.
The pins of the one-map data at $i$ give representatives of the classes of their lattices in $V_i$. A
lift form $x$ at $(i,\mu_0)$ is a \emph{class sum} $\sum_ka_kf_k$, with core forms $f_k$ and $a_k\in
L$, if the section of the core that the one-map data build from the values of $j_V(x)$ at these
representatives is $\sum_ka_kf_k$ after extension of scalars to $L$. The \emph{theta integral} of a
conjugate datum is the map $\Theta$ with
$\Theta(\varphi)(v)(g)=\int\theta_\omega\big(h,g;\phi_{\mathrm{io}(v\otimes\tilde\varphi(h))}\big)\,
d\nu(h)$, in the notation of Definition~\ref{is-eigenvalue-conjugate-datum}(5), so that
$\mathrm{th}(\varphi)(v)=e\,\Theta(\varphi)(v)$. Then:
\begin{enumerate}
\item[(A)] (Evaluation.) If $\sigma$ is good, there is a map $J$ from the core
  forms to $H_W(i,\sigma)$, semilinear for the restriction of $\sigma$ to $K_s$, such that the
  function of $J(f)$ at $v_0(i,\sigma)$ is $\mathrm{ad}_\sigma(f)$ for every core form $f$. If $\sigma$
  is not good, there are a tensor $v$ of unit vectors of the models of
  discrete series of $U(W_i)$ at $\sigma$ and such a map $J$ with the function of $J(f)$ at $v$ equal
  to $\overline{\mathrm{ad}_{\bar\sigma}(f)}$ for every $f$.
\item[(B)] (Uniqueness.) Let $\sigma$ agree with $\sigma_0$ on $K$. If two conjugate data at $\sigma$
  for $i$ and the tie have the same characters $\chi',\eta'$, maps $j,j'$ and theta integrals
  $\Theta,\Theta'$, then there is $\kappa\in\mathbf C$ with $|\kappa|=1$ and
  $\Theta'(j'(x))=\kappa\,\Theta(j(x))$ for every lift form $x$ at $(i,\mu_0)$.
\item[(C)] (Conjugate lifts.) Let $\sigma\ne\sigma_0$ agree with $\sigma_0$ on $K$. There are a
  conjugate datum at $\sigma$ for $i$ and the tie, with map $j$ and theta integral $\Theta$, a
  function $u$ on $U(W_i)(\mathbf A)$ of absolute value $1$ with a unitary operator $U$ of
  $H_W(i,\sigma)$ such that $U(Y)(v)=u\,Y(v)$ for all $Y$ and $v$, a central element $z_2$ of $G$ and
  $\kappa\in\mathbf C$ with $|\kappa|=1$, such that:
  \begin{enumerate}
  \item[(a)] $u(h)\,\mathrm{ad}_\sigma(gf)(h)=\sigma(\chi_W(g))\,u(hk)\,\mathrm{ad}_\sigma(f)(hk)$ for
    every $g\in G$ with a level lift $k$, every core form $f$ and every $h\in U(W_i)(\mathbf A)$;
  \item[(b)] $\kappa\,\Theta(j(x))(v_0(i,\sigma))=\sum_k\sigma(a_k)\,u\,\mathrm{ad}_\sigma(z_2f_k)$ for
    every lift form $x$ at $(i,\mu_0)$ that is a class sum $\sum_ka_kf_k$.
  \end{enumerate}
\end{enumerate}
\end{lemma}
\leanref{Development/AbelianVarietyTheory/Proofs/TwistedOperators/JEvExistsAndJEvBar.lean}{GlobalTheory.ConjReading.jEvExists_and_jEvBar}
\leanref{Development/AbelianVarietyTheory/Proofs/IwasawaDecomposition/JEvExistsAndJEvBar.lean}{GlobalTheory.CuspValues.jEvExists}
\leanref{Development/AbelianVarietyTheory/Proofs/DivisorProofs/ConjugateLiftUniqueAndBaseDatum.lean}{GlobalTheory.ConjLiftUnique.lift_unique_five}
\leanref{Development/AbelianVarietyTheory/Proofs/DivisorProofs/ConjugatesOfRationalThetaLiftFromCharacterisatio.lean}{GlobalTheory.ConjRatLiftFive.oneMapConj5}
\leanref{Development/AbelianVarietyTheory/Proofs/DivisorProofs/CoefMapNormEquivariantBasisOfFive.lean}{GlobalTheory.GalOne.oneMap_of_ext_ofFive}
\leanref{Development/AbelianVarietyTheory/Proofs/DivisorProofs/ConjugateLiftOfCharacterisationClausesStatement.lean}{GlobalTheory.T.conjugateLift_of_characterisationClauses}
\begin{proof}
\emph{Proof of the first assertion of (A).} Let $\sigma$ be good, write $\sigma_s$ for its restriction
to $K_s$, $M$ for the model of $W_i$ at $\sigma$ and $v_0=v_0(i,\sigma)$. By
Lemma~\ref{is-eigenvalue-conjugate-good}(2), as in the proof of
Lemma~\ref{is-eigenvalue-conjugate-models}(1), $M$ has holomorphic shape and integral extremal weights,
and $v_0$ is an extremal vector of $M$ by the weight tie. By the definition of the core,
$\mathrm{ad}_\sigma(f)$ is a constant multiple, independent of $f$, of the adelization through the
dictionary of the core at level $i+1$ and $\sigma_s$ of the reading of $f$ at $\sigma_s$, which is
$\sigma_s$-semilinear in $f$; and at a good embedding the reading of a core form lies in the span of the
bounded holomorphic cusp sections of weight $\kappa_i$, the families $P=(P_c)$ of cusp forms
$P_c\in S_{\kappa_c}(X_c)$ on the components, $\kappa_c$ the weight of $\kappa_i$ on $X_c$. Let $P$ be
such a family, and let $w$ be the extremal weights of $M$ in the order of the columns of its frames. By
Lemma~\ref{is-eigenvalue-conjugate-dictionary-good}, the dictionary of the core there is good for
$(M,w,K_W)$, the elements of $K_W$ have archimedean component $1$, and the level character $\chi_K$ of
the trivial level datum is trivial on $K_W$. So the adelization of $P$ is a bounded holomorphic
automorphic form of weight $w$ and level $(K_W,1)=(K_W,\chi_K)$ whose constant terms vanish, by
Lemma~\ref{is-eigenvalue-conjugate-adelization}. As $U(W_i)(F_h)$ has a fundamental domain of finite
volume (Lemma~\ref{is-eigenvalue-conjugate-spectral-facts}(1)), it is the value at the extremal vector
$v_0$ of a multiplicity map of $M$ of level $(K_W,\chi_K)$, by
Lemma~\ref{is-eigenvalue-conjugate-multiplicity}. By the core conditions on the
models, every such multiplicity map is the family of functions of an element of $H_W(i,\sigma)$. Hence,
by linearity, for every core form $f$ there is an element $J(f)$ of $H_W(i,\sigma)$ whose function at
$v_0$ is $\mathrm{ad}_\sigma(f)$, and it is unique by Lemma~\ref{is-eigenvalue-core-group}(3). The
uniqueness makes $J$ additive and $\sigma_s$-semilinear, as $f\mapsto\mathrm{ad}_\sigma(f)$ is.

\emph{Proof of the second assertion of (A), from the first.} Let $\sigma$ not be good. Then
$\bar\sigma$ is good, by Lemma~\ref{is-eigenvalue-conjugate-good}(1), and the
restriction of $\bar\sigma$ to $K_s$ is the complex conjugate of that of $\sigma$. The first assertion at
$\bar\sigma$ gives a map $\bar J$ from the core forms to $H_W(i,\bar\sigma)$, semilinear for the
restriction of $\bar\sigma$ to $K_s$, such that the function of $\bar J(f)$ at $v_0(i,\bar\sigma)$ is
$\mathrm{ad}_{\bar\sigma}(f)$. Let $c_\tau$ be the maps of Lemma~\ref{is-eigenvalue-conjugate-models}(1),
and put $v_\tau=c_\tau^{-1}(v_0(i,\bar\sigma)_\tau)$ and $v=\otimes_\tau v_\tau$. By
Lemma~\ref{is-eigenvalue-conjugate-models}(2), applied to the $c_\tau$ from the model at $\sigma$ to that
at $\bar\sigma$, each $v_\tau$ is smooth; it is a unit vector, because $c_\tau$ is isometric and the
weight tie makes $v_0(i,\bar\sigma)$ a tensor of unit vectors. Let $Y\in H_W(i,\bar\sigma)$. By the core
conditions on the models, the functions of $Y$ form an element of the multiplicity space of the model
at $\bar\sigma$ of level $(K_W,\chi_K)$, with the level character $\chi_K$ of the trivial level datum,
which is trivial. Lemma~\ref{is-eigenvalue-conjugate-models}(2) gives an element of the multiplicity
space of the model at $\sigma$ of level $(K_W,\bar\chi_K)=(K_W,\chi_K)$ whose value at $v$ is the
complex conjugate of the function of $Y$ at $v_0(i,\bar\sigma)$. By the core conditions on the models
again, it is the family of functions of some $X(Y)\in H_W(i,\sigma)$. Put $J(f)=X(\bar J(f))$, so that
the function of $J(f)$ at $v$ is $\overline{\mathrm{ad}_{\bar\sigma}(f)}$. An element of $H_W(i,\sigma)$
is determined by its function at $v$, by Lemma~\ref{is-eigenvalue-core-group}(3). So $J$ is additive,
and $J(af)=\sigma(a)J(f)$ for $a\in K_s$, since the function of $J(af)$ at $v$ is the conjugate of
$\bar\sigma(a)$ times that of $\bar J(f)$ at $v_0(i,\bar\sigma)$.
% This is CuspValues.jEvExists_and_jEvBar, second half: T.hilbertModel_conjugation_antilinear_bijection =
% ConjHilb.conj_hilbert (key: T.exists_antiunitary_intertwiner_conjugateEmbedding_models, ConjMult.conj_mult
% for smoothness, IsUnitV from (ht.hv0 i _).norm, exists_conj_value with ConjEval.levChi_principal and
% models_HW; then ConjEval.evEx_conj with evAt_inj and JEvExists at conj o sigma, through good_resC).

\emph{Proof of (B).} Let two conjugate data at $\sigma$ for $i$ and the tie be given, with the same
characters $\chi',\eta'$. Write $(W_c,\rho)$ and $(W'_c,\rho')$ for their coefficient data, $c_1$ and
$c'_1$ for the isomorphisms of (2) of Definition~\ref{is-eigenvalue-conjugate-datum}, and
$\omega,\omega'$, $\phi,\phi'$, $B,B'$ and $\mathrm{io},\mathrm{io}'$ for the Weil data, level vectors,
bases and maps of (5). With the bases $B_V$ of $V_i$ of the realization and the notation of
Lemma~\ref{is-eigenvalue-conjugate-pure-tensors}, $\phi_Q=\phi\otimes_BQ$, so that
\[
 \Theta(y)(z)(g)=\int_{\mathcal F}\theta_\omega\big(h,g;\phi\otimes_B\mathrm{io}(z\otimes\tilde
 y(h))\big)\,d\nu(h)
\]
for $y$ in the form space of the first datum, $z\in\bigotimes_\rho\kappa_\rho(\sigma)$ and $g\in
U(W_i)(\mathbf A)$, and likewise for $\Theta'$ with $\omega',\phi',B',\mathrm{io}'$. By the core
conditions on the level characters, the lattices of the frames of $V_i$ and $W_i$ are self-dual. By the
core conditions on the models, the member of $\theta$ at each $\tau$ extends $\tau$, and $V_i$ has an
orthonormal basis at every $\tau$, so that $1$ is a Cartan involution of $V_i$. The Cartan involutions
of the model of $W_i$ at $\sigma$ are Cartan involutions of $W_i$.

\emph{Step 1.} The level vectors $\phi$ and $\phi'$ are formed from the same frames, the set $P_g$ and
the trivial characters, so $\phi'=a\phi$ for some $a$ with $|a|=1$, by
Lemma~\ref{is-eigenvalue-conjugate-level-vectors}.

\emph{Step 2.} Lemma~\ref{is-eigenvalue-conjugate-coefficients} applies to $c_1$ and $c'_1$, with the
rational basis of the coefficient space of $V_i$ and, for $\gamma\in U(V_i)(F_h)$, the rational action
$r_\gamma$ of the matrix of $\gamma$: both maps carry the rational basis to bases, their inner products
on it are the images under $\sigma$ of the rational pairing, and both intertwine the rational action.
It gives an isomorphism $\Phi\colon W_c\to W'_c$ that preserves norms, with $\Phi\circ c_1=c'_1$ and
$\Phi\rho(h)=\rho'(h)\Phi$.

\emph{Step 3.} By Lemma~\ref{is-eigenvalue-conjugate-theta-facts}(3), $\phi$ is a theta function
annihilated by the annihilation operators for $1$ and the Cartan involutions of the model at every
$\tau$. Since $\sigma$ agrees with $\sigma_0$ on $K$, Lemma~\ref{is-eigenvalue-conjugate-theta-facts}(6)
gives the hypothesis on the lowest $K$-types of Lemma~\ref{is-eigenvalue-conjugate-pure-tensors}. The
maps $\mathrm{io}$ and $\mathrm{io}'$ are io maps for $(\omega,\phi,B,\Delta_i(\sigma))$ and
$(\omega',a\phi,B',\Delta_i(\sigma))$, by Definition~\ref{is-eigenvalue-lift-data}(7). So
Lemma~\ref{is-eigenvalue-conjugate-pure-tensors} gives $\kappa$ with $|\kappa|=1$ and
$\phi'\otimes_{B'}\mathrm{io}'(z\otimes\Phi w)=\kappa\,\phi\otimes_B\mathrm{io}(z\otimes w)$ for all $z$
and $w$.

\emph{Step 4.} Let $x$ be a lift form at $(i,\mu_0)$ and $h\in U(V_i)(\mathbf A_{F_h,f})$. By (4) of
Definition~\ref{is-eigenvalue-conjugate-datum} there are rational coefficient vectors $u,u'$ with
$j_V(x)(h)=c_0(u)=c_0(u')$, $j(x)(h)=c_1(u)$ and $j'(x)(h)=c'_1(u')$. The family tie makes $c_0$ carry
the rational basis to linearly independent vectors, so $c_0$ is injective and $u=u'$. Hence
$j'(x)(h)=c'_1(u)=\Phi(j(x)(h))$, and $\widetilde{j'(x)}(h)=\rho'(h^{-1})\Phi(j(x)(h_f))
=\Phi(\widetilde{j(x)}(h))$ for every $h\in U(V_i)(\mathbf A)$.

\emph{Step 5.} The functions $\phi'\otimes_{B'}Q$ are theta functions, on which $\omega'$ agrees with
$\omega$ (Lemma~\ref{is-eigenvalue-conjugate-theta-facts}(4) and~(2)). By Steps 3 and 4, for every
lift form $x$, $z$, $g$ and $h$,
\[
 \theta_{\omega'}\big(h,g;\phi'\otimes_{B'}\mathrm{io}'(z\otimes\widetilde{j'(x)}(h))\big)
 =\theta_\omega\big(h,g;\phi'\otimes_{B'}\mathrm{io}'(z\otimes\Phi(\widetilde{j(x)}(h)))\big)
 =\kappa\,\theta_\omega\big(h,g;\phi\otimes_B\mathrm{io}(z\otimes\widetilde{j(x)}(h))\big).
\]
Integrating over $\mathcal F$ gives $\Theta'(j'(x))=\kappa\,\Theta(j(x))$.
% (B) is ConjLiftUnique.lift_unique_five. The data: hC.models.inf (theta), SDFrameLattice.isSelfDualC_of_isSDFrame
% for hC.lev.betaV i and betaW i, Chain.VacVec.cart_isCartanInv, GalLift.Vec.isCartanInv_one (models.BV),
% HolLowOnce.lowSm_and_lowIrr with GalLift.ProjLift.holShape_of_overK. Step 1 is vec_unique; Step 2 is
% exists_isometry with hj.coef, hj.indep, hj.span, hj.pair and hj.act; Step 3 is io_unique with
% Vec.theta_of_thetaAt and DatumFock.isFockVacAt_thPVec; Step 4 is hb (hj.base, hj'.base, rescale_inj with
% (hT.mapL i).2.1); Step 5 is thetaP_eq with WeilUniq.weilDatum_unique and p3.

\emph{Proof of (C).} Let $\sigma\ne\sigma_0$ agree with $\sigma_0$ on $K$; then $\sigma$ is good, by
Lemma~\ref{is-eigenvalue-conjugate-good}(3). The proof assembles the following inputs, which come from
the characterization of the realization and from the one-map data at $i$.
\begin{enumerate}
\item[(C1)] (Characters.) The weight tie gives a $\sigma$-conjugate $\chi'$ of $\mu^0_i$, the
  characterization gives a $\sigma$-conjugate $\eta'$ of $\eta$, and there is a Weil datum $\omega$ for
  $(V_i,W_i)$ with characters $(\chi',\chi'\eta'^{-1})$.
\item[(C2)] (Coefficients and forms.) For every family tie of $R$ there are a coefficient datum $W_c$ with
  a map $c_1$ as in (2) of Definition~\ref{is-eigenvalue-conjugate-datum}, and a form space $H$ with a
  map $j$ as in (3) and~(4) for that tie. If $(W_c,c_1,H,j)$ satisfy (2)--(4) for one family tie, then
  $(W_c,c'_1,H,j)$ satisfy them for any other, for some $c'_1$.
\item[(C3)] (Vectors and the theta integral.) There are a level vector $\phi$ of the trivial level and
  bases $B$ as in Definition~\ref{is-eigenvalue-lift-data}(5); for $\omega$ as in (C1) and these, a map
  $\mathrm{io}$ as in Definition~\ref{is-eigenvalue-lift-data}(7); and the integrals that define the
  theta integral $\Theta$ of these data converge on $H$.
\item[(C4)] (Comparison with the one-map data.) The pins of the one-map data at $i$ provide a family tie
  of $R$. For that tie, and data as in (C1)--(C3) for it, there are a function $u$ on
  $U(W_i)(\mathbf A)$ of absolute value $1$, invariant on the left under $U(W_i)(F_h)$ and on the right
  under the elements of finite component $1$ and under $K_W$, a central element $z_2$ of $G$ and
  $\kappa\in\mathbf C$ with $|\kappa|=1$, such that (a) holds, and (b) holds for $\Theta$ and $j$. For such a $u$ there is a unitary
  operator $U$ of $H_W(i,\sigma)$ with $U(Y)(v)=u\,Y(v)$ for all $Y$ and $v$.
\item[(C5)] (The theta map.) For data as in (C1)--(C3), if $\Theta(j(x))(v_0(i,\sigma))$ is the function
  at $v_0(i,\sigma)$ of an element of $H_W(i,\sigma)$ for every lift form $x$ at $(i,\mu_0)$, then there
  is a linear map $\mathrm{th}\colon H\to H_W(i,\sigma)$ with $\mathrm{th}(\varphi)(v)=e\,\Theta(\varphi)(v)$
  for every $\varphi\in H$ and every tensor $v$ of lowest $K$-type vectors.
\end{enumerate}
Take $\chi'$, $\eta'$ and $\omega$ by (C1), $W_c$, $c_1$, $H$ and $j$ by (C2) for the tie of the pins,
$\phi$, $B$, $\mathrm{io}$ and $\Theta$ by (C3), and $u$, $U$, $z_2$ and $\kappa$ by (C4). Let $x$ be a lift
form at $(i,\mu_0)$. By Lemma~\ref{is-eigenvalue-conjugate-class-sums}, applied to the pinned one-map
data at $i$, $x$ is a class sum $\sum_ka_kf_k$. By (A),
$\mathrm{ad}_\sigma(z_2f_k)$ is the function at $v_0(i,\sigma)$ of $J(z_2f_k)$, so
$u\,\mathrm{ad}_\sigma(z_2f_k)$ is that of $U(J(z_2f_k))$, and by (b),
\[
 \Theta(j(x))(v_0(i,\sigma))=\kappa^{-1}\sum_k\sigma(a_k)\,u\,\mathrm{ad}_\sigma(z_2f_k)
\]
is the function at $v_0(i,\sigma)$ of $\kappa^{-1}\sum_k\sigma(a_k)\,U(J(z_2f_k))$. So (C5) gives
$\mathrm{th}$. Finally (C2) replaces $c_1$ by a map $c'_1$ for the given tie, with the same $H$ and $j$.
Then $\chi'$, $\eta'$, $W_c$, $c'_1$, $H$, $j$ and $\mathrm{th}$ form a conjugate datum at $\sigma$ for
$i$ and the tie, with theta integral $\Theta$, and $u$, $U$, $z_2$ and $\kappa$ satisfy (a) and (b).

\emph{Proof of (C1)--(C3) and of the last assertion of (C4).} (C1) The weight tie gives at every
good embedding a $\sigma$-conjugate $\chi'$ of $\mu^0_i$. By
Lemma~\ref{is-eigenvalue-conjugate-characters}(2) there is a $\sigma$-conjugate $\eta'$ of $\eta$, and by
Lemma~\ref{is-eigenvalue-conjugate-characters}(4) a Weil datum $\omega$ with characters
$(\chi',\chi'\eta'^{-1})$. (C2) is Lemma~\ref{is-eigenvalue-conjugate-coefficient-forms}. (C3) is
Lemma~\ref{is-eigenvalue-conjugate-level-data}. The last assertion of (C4) is
Lemma~\ref{is-eigenvalue-conjugate-phase}.

\emph{Proof of (C5).} Let data as in (C1)--(C3) be given, with the form space $H$, the map $j$, the
Weil datum $\omega$, the level vector $\phi$, the bases $B$, the map $\mathrm{io}$ and the theta
integral $\Theta$, and for every lift form $x$ let $Y_x\in H_W(i,\sigma)$ have the function
$\Theta(j(x))(v_0)$ at $v_0=v_0(i,\sigma)$. Let $M$ be the model of $W_i$ at $\sigma$, with the
notation of Lemma~\ref{is-eigenvalue-conjugate-lowest-types}. As $\sigma$ agrees with $\sigma_0$ on
$K$, $M$ has smooth irreducible lowest $K$-types $\kappa_\tau$, by
Lemma~\ref{is-eigenvalue-conjugate-theta-facts}(6), and the components of $v_0$ are unit vectors of the
$\kappa_\tau$, by the weight tie.

\emph{Step 1.} For every $\varphi\in H$, $\Theta(\varphi)$ is $K$-equivariant on
$\prod_\tau\kappa_\tau$. Indeed, let $k\in K_\tau$, and let $v'$ be $v$ with $v_\tau$ replaced by
$\sigma_\tau(k)v_\tau$. For $h\in U(V_i)(\mathbf A)$ and $g\in U(W_i)(\mathbf A)$, the theta series of
$\omega(h,gk^\theta)\phi_Q$ is that of $\omega(h,g)\big(\omega(1,k^\theta)\phi_Q\big)$, and
$\omega(1,k^\theta)\phi_{\mathrm{io}(v\otimes w)}=\phi_{\mathrm{io}(v'\otimes w)}$ for every $w\in W_c$,
because $\mathrm{io}$ is an io map (Lemma~\ref{is-eigenvalue-conjugate-pure-tensors}(ii)), as in Step 3
of the proof of (B). So the integrands that define $\Theta(\varphi)(v')(g)$ and
$\Theta(\varphi)(v)(gk^\theta)$ agree, and so do the integrals.

\emph{Step 2.} By the core conditions on the models, the functions $Y(v)$ of each $Y\in H_W(i,\sigma)$
form an element of the multiplicity space of $M$ of level $(K_W,\chi_K)$, which is $K$-equivariant by
Lemma~\ref{is-eigenvalue-conjugate-lowest-types}(1). For a lift form $x$, the functions of $Y_x$ and
$\Theta(j(x))$ agree at $v_0$, whose components are nonzero, so they agree at every tensor of vectors of
the $\kappa_\tau$, by Step 1 and Lemma~\ref{is-eigenvalue-conjugate-lowest-types}(2).

\emph{Step 3.} Let $X$ be the space of maps from $\prod_\tau\kappa_\tau$ to functions on
$U(W_i)(\mathbf A)$, and let $L_1\colon H_W(i,\sigma)\to X$ and $L_2\colon H\to X$ be the linear maps
$Y\mapsto(v\mapsto Y(v))$ and $\varphi\mapsto(v\mapsto\Theta(\varphi)(v))$. By Step 2,
$L_2(j(x))=L_1(Y_x)$ lies in the image of $L_1$ for every lift form $x$. The image of $j$ spans $H$, by
the proof of Lemma~\ref{is-eigenvalue-base-inputs}(4), which uses only (2)--(4) of
Definition~\ref{is-eigenvalue-conjugate-datum}. So $L_2(H)$ lies in the image of $L_1$. Choosing a
basis of $H$ and a preimage under $L_1$ of the image under $L_2$ of each of its elements gives a linear
map $\lambda\colon H\to H_W(i,\sigma)$ with $L_1\circ\lambda=L_2$. Put $\mathrm{th}=e\circ\lambda$. By
the core conditions on the models, the functions of $eY$ are those of $Y$ with $e$ applied, so
$\mathrm{th}(\varphi)(v)=e\,\Theta(\varphi)(v)$ for every $\varphi\in H$ and every tensor $v$ of vectors
of the $\kappa_\tau$.
% (C5) is GalLift.ProjLift.exists_isProjOf_five: hspan = ConjSpan.span_conjForms; hlow =
% HolLowOnce.lowSm_and_lowIrr with holShape_of_overK; Step 1 is kEquiv_thetaP (hio.2.1 and thetaKO of a
% product); Step 2 is eq_on_low inside exists_lift, with models_HW for the multiplicity maps; Step 3 is
% exists_factor (Module.projective_lifting_property) and proj_apply with models_gW. The division by
% kappa in the assembly is Assembly.exists_ev.

\emph{Proof of the rest of (C4).} Write $L'$, $\iota$ and $c'_0$ for the coefficient field, the map
$L\to L'$ and the coefficient map of the one-map data at $i$, and $\iota(v)$ for a rational coefficient
vector $v$ of $V_i$ over $L$ with its coefficients mapped by $\iota$, as in
Lemma~\ref{is-eigenvalue-conjugate-class-values}. The pins of the one-map data at $i$ include a family
tie of $R$, with coefficient map $c_0$, and a base change of it to the one-map data:
$c'_0(\iota(v))=c_0(v)$ for every $v$. Let $\chi'$, $\eta'$, $\omega$, $W_c$, $c_1$, $H$, $j$, $\phi$,
$B$, $\mathrm{io}$ and $\Theta$ be data as in (C1)--(C3) for this tie. By
Lemma~\ref{is-eigenvalue-conjugate-extension} there is an embedding $\sigma'$ of $L'$ with
$\sigma'\circ\iota=\sigma$, and $\sigma'$ and $\sigma$ have the same restriction $s$ to $K_s$. As $\sigma$
agrees with $\sigma_0$ on $K$, it is good (Lemma~\ref{is-eigenvalue-conjugate-good}(3)). Use the notation of Lemma~\ref{is-eigenvalue-conjugate-frame-inputs}(2), and let $a$,
$\kappa$ and $z_2=z$ be as there. The canonical constants have $|c_E|=1$, by the weight tie. Put
$u(h)=a(c(h))/c_E(c(h))$. By Lemma~\ref{is-eigenvalue-conjugate-phase-function}, applied to the function
$a/c_E$ on the components, $|u|=1$ and $u$ is invariant on the left under $U(W_i)(F_h)$ and on the right
under the elements with finite component $1$ and under $K_W$. As
$\mathrm{ad}_\sigma(f)(h)=\ell_i\,c_E(c(h))\,\mu(h)\,\mathrm{rd}(f)_{c(h)}(x(h))$,
\[
 u(h)\,\mathrm{ad}_\sigma(f)(h)=\ell_i\,\mathrm{ad}_a(\mathrm{rd}(f))(h)\tag{$*$}
\]
for every section $f$ over $K_s$ and every $h\in U(W_i)(\mathbf A)$.

(a) Let $g\in G$ with a level lift $k$, and $f$ a core form. Let $b\in K_s$ be as in
Lemma~\ref{is-eigenvalue-conjugate-frame-inputs}(2)(a). By the core conditions on the group,
$\sigma_0(\chi_W(g))=\prod_{P\in P_g}\chi_{i+1,P}(\det g_P)=\sigma_0(b)$; so $\chi_W(g)=b$ in $L$, and
$\sigma(\chi_W(g))=s(b)$. By $(*)$ and that lemma, for every $h$,
\[
 u(h)\,\mathrm{ad}_\sigma(gf)(h)=\ell_i\,s(b)\,\mathrm{ad}_a(\mathrm{rd}(f))(hk)
 =\sigma(\chi_W(g))\,u(hk)\,\mathrm{ad}_\sigma(f)(hk).
\]

(b) Let $x$ be a lift form at $(i,\mu_0)$ that is a class sum $\sum_ka_kf_k$, for the representatives
$h_1,\dots,h_N$ that the pins give (Lemma~\ref{is-eigenvalue-conjugate-class-sums}(1)). So there are
rational coefficient vectors $x_1,\dots,x_N$ of $V_i$ over $L'$ with $j_V(x)(h_j)=c'_0(x_j)$, and a section
$f'$ of weight $\kappa_i$ over $L'$ built from $(x_j)$ with $f'=\sum_k\iota(a_k)f_k$ after extension of
scalars. By Lemma~\ref{is-eigenvalue-conjugate-class-values}, $x_j=\iota(x''_j)$ for rational coefficient
vectors $x''_j$ over $L$ with $j(x)(h_j)=c_1(x''_j)$. By
Lemma~\ref{is-eigenvalue-conjugate-class-action}(2), $zf'=\sum_k\iota(a_k)(zf_k)$ after extension of
scalars. Let $g\in U(W_i)(\mathbf A)$. Lemma~\ref{is-eigenvalue-conjugate-frame-inputs}(2)(b), with the
$x''_j$, the sections $zf_k$ and the coefficients $\iota(a_k)$, gives the following identity, in which
$c_1(x''_j)$ is written $j(x)(h_j)$ and $\sigma'(\iota(a_k))$ is written $\sigma(a_k)$:
\[
 \kappa\sum_j\sigma'(\mathrm{vol}\,\gamma_j^{-1})\,\theta_\omega\big(h_j,g;
 \phi_{\mathrm{io}(v_0(i,\sigma)\otimes j(x)(h_j))}\big)
 =\ell_i\sum_k\sigma(a_k)\,\mathrm{ad}_a(\mathrm{rd}(zf_k))(g).
\]
The representatives have archimedean component $1$ and represent the double cosets of a subgroup
$J'$ of finite index in the level $J$ of $V_i$ at $\mu_0$, with $\mathrm{vol}=[J:J']^{-1}$ and $\gamma_j$
the order of the group $\Gamma_j$ of Lemma~\ref{is-eigenvalue-conjugate-unfolding}; so
$\sigma'(\mathrm{vol}\,\gamma_j^{-1})=[J:J']^{-1}|\Gamma_j|^{-1}$. By
Lemma~\ref{is-eigenvalue-conjugate-theta-sum}, applied to $y=j(x)$ and $v=v_0(i,\sigma)$, whose
components are lowest $K$-type vectors by the weight tie, the left side is
$\kappa\,\Theta(j(x))(v_0(i,\sigma))(g)$. By $(*)$ the right side is
$\sum_k\sigma(a_k)\,u(g)\,\mathrm{ad}_\sigma(zf_k)(g)$. This is (b).
% (C4) is GalOne.oneMap_of_ext_ofFive, through ConjRatLiftFive.oneMapConj5, applied with the tie of
% hP.2.lift.c0 (ThPinsR.c0: FamTie and CoefficientBaseChange). oneMapConj5 takes sigma' from GalOne.exists_ext
% (Lemma is-eigenvalue-conjugate-extension(1)) and the theta property of f from hL.2.1.theta 1. In
% oneMap_of_ext_ofFive: hF is Facts.factsAt_ofFive, of which the proof uses res T sigma = sOf O sigma'
% (Lemma is-eigenvalue-conjugate-extension(2)) and GoodEmb (overK_iff_good_ofFive); hpb is
% CoverHalf.primes_eq_of_halfIs; hcE is tie.laws.cUnit; hkn ((kappa i).n = rank Vc) is from
% FactsLat.facts_lat_ofFive. The constants cst, c0, z and ph are those of
% Lemma is-eigenvalue-conjugate-frame-inputs(2), with kappa = ph * c0; u = cst/cE is PhaseFun by
% comp_inv_ofFive. (a) is eqAt_of_frAtEmb_ofFive with e1 (hR.grp.chiW and the injectivity of embLL),
% e2 (hF.1) and div_mul_phiC, which is (*). (b) is exists_values_ofFive (h6), FieldBC.actLev_and_actBC
% (hbc), identity_of_frConst_ofFive (h7) and theta_eq_classSum_of_conjugate (h8), with div_mul_phiC.
% (C1)-(C3) and U are written in Lemmas is-eigenvalue-conjugate-characters, -coefficient-forms,
% -level-data and -phase.
% (A) is GlobalTheory.ConjReading.jEvExists_and_jEvBar = CuspValues.jEvExists_and_jEvBar. Its first
% assertion JEvExists = CuspValues.jEvExists is proved above, following its proof: for P in holSub,
% ad P is a HolForm (Lemma is-eigenvalue-conjugate-adelization(1), with Lemma
% is-eigenvalue-conjugate-dictionary-good for GoodDict and the level) with ZCusp (Lemma
% is-eigenvalue-conjugate-adelization(2)); Lemma is-eigenvalue-conjugate-multiplicity (BackForm.exists_multMap,
% with HolDict.exists_finPair) gives the multiplicity map, which is (HW i sigma).e x by models_HW; then
% EvalAtUnitVector.evInj and CuspValues.evEx_of_range give J with ev (J f) = lieFac * ad (rdU f) = phiC
% (phiC_eq), using mem_cusp for f. The pins give T.pb = R.C.P.pb (ConjReading.pb_eq_of_omPins). The second
% assertion (JEvBar) is proved above from the first.
% (C) is the construction in ConjRatLiftFive.conjugateRationalLift_at_nonBaseEmbedding with conj_lift5,
% whose assembly is written above. The inputs: (C1) is fw.chiR (ChiAt) with
% GalLift.overK_iff_good_ofFive, clauses 2 and 3 of T.conjugateCharacters_embeddings_of_characterisation
% (EtaVals, IsEtaConj), hypConj5 (HypConj) and its seventh clause (IsWeilDatumO); (C2) is clauses 1, 2
% and 3 of T.conjugateLift_of_characterisationClauses (CoefAt = CoefModel.exists_coefAt_of_posDef,
% IsConjForms = ConjForms.exists_conjForms_five through exists_baseChange, conjForms_rescale_five);
% (C3) is clause 4 (VecData = Vec.exists_vecData_five), MovedWeights.exists_io_conj_ofFive (IsIoW) and
% clause 5 (Lift.IsThetaP = ConjThetaIntegral.exists_isThetaP_conj_five); (C4) is
% ConjRatLiftFive.oneMapConj5 = GalOne.oneMap_of_ext_ofFive (PhaseFun u, z central, |kappa| = 1,
% EquivBy for (a), the class-sum identity for (b)), applied with the tie of hP.2.lift.c0 and its
% CoefficientBaseChange, and the second clause of T.centralShift_phase_conjReadings_fiveClauses
% (CompPhase.exists_isCompPhase_five, IsCompPhase) for U; (C5) is clause 6 (IsProjOf =
% ProjLift.exists_isProjOf_five), proved above, with Assembly.exists_ev for the division by kappa. The
% class-sum presentation is base_lift5 (the last two clauses of T.conjugateLift_unique_and_baseDatum, with
% exists_reps for the representatives), written as Lemma is-eigenvalue-conjugate-class-sums, and the
% evaluation of U(J(z2 f_k)) is hJ in conj_lift5 (fw.Evaluates from JEvExists, and IsCompPhase).
% For the next writer on (C1)-(C4): EtaVals (EtaValues.etaVals_five) is the argument already written in
% the proof of Lemma is-eigenvalue-base-inputs(iii); IsEtaConj comes from EtaConjExist.exists_conj; the
% Weil datum from GalLift.ConjWeil.exists_of_datum (the base Weil datum hR.thP.omegaP twisted, with the
% parities resEta_parity of hH.mu_res and hH.eta_res); HypConj from GalLift.hypConj_of_chars_ofFive and
% EtaConjType.resEta_hasInfExp_of_isEtaConj_ofFive. CoefAt is CoefModel.exists_coef_round_at (coefOf
% with toHilb and the Fock pairing pairK); VecData uses Member.vacuum_oriented_exists,
% PlaceVec.thPVecExists and exists_isONBasis; the theta integral is IsThetaP2 in ExistsIsThetaPForms.
\end{proof}

\subsection{The lower norm bound and the surjectivity bound}\label{subsection-is-lift-bounds}

\begin{lemma}\label{is-eigenvalue-base-inputs}
Let $D$ be a definite-side datum and $R$ a realization of $D$ with a face $F$
(Definition~\ref{is-realisation}), with the notation of Definitions~\ref{is-eigenvalue-lift-data}
and~\ref{is-eigenvalue-conjugate-datum}.
\begin{enumerate}
\item (Base links.) Suppose that the family of definite forms satisfies the \emph{family conditions}
  (the involution of $L$ is complex conjugation at $v_0$, there are finitely many tuples, the
  characters of $G$ in the family are absolutely irreducible, and its row and lift data satisfy their
  conditions), with the conditions on its inert indices, and has lift-distinct types, and that the
  realization satisfies the weak
  characterization at $F$ and the existence conditions for its theta maps, has primitive rows at $F$
  (Definition~\ref{is-primitive-rows}) and has all its phases equal to $1$. Then $R$ satisfies the core
  conditions at $F$, and the rows $S_0(R,F)$ have the \emph{base links}: for every lift $j$ and tuple
  $t$ the map $t_0(j,t)$ of Definition~\ref{is-complex-data}(c) is injective with the same image as
  $w_0(j)$, and the matrix that expresses it through $w_0(j)$ is the image under $\sigma_0$ of a matrix
  over $L$; and for every lift $j$ the entries of $H_T(j,\mu_0)$ lie in the subfield $K_0$ of the base
  data and $|y|^2_{X_T(j,\mu_0),v}=\sigma_v(y)^*\sigma_v(H_T(j,\mu_0))\sigma_v(y)$ for every infinite
  place $v$ of $L$ and every $y$.
\item (Prime to the host.) If $R$ satisfies the core conditions at $F$, except possibly the one on its
  chain, then the base tuple $\mu_0$ is prime to the host.
\item (Rational lifts.) Let one-map data for $R$ be given (Definition~\ref{is-one-map-data}), with
  weight families $\kappa_i$ and vectors $v_0(i,\sigma)$, let $i$ be a lift and $\phi\in\mathbf C$.
  Suppose that $R$ satisfies the core conditions at $F$, except possibly the one on its chain, that
  $\psi$ is the standard character with sign $-1$, that the one-map data have a weight tie with a
  conjugation tie and are pinned at every lift, that (1) of Definition~\ref{is-bounded-rows} holds,
  that the rational structures $FW_i$ satisfy (11) of Definition~\ref{is-primitive-rows}, and that the
  one-map data at $i$ give the $i$-th lift as one map with the canonical constants of the core and
  phase $\phi$. Call $z\in L^{n_W(i)}$ a \emph{rational lift} of a lift form $x$ at $(i,\mu_0)$ if
  $FW_i(\sigma_0)(z)=\phi\,\mathrm{th}_G(i,\mu_0)(j_V(x))$. Then:
  \begin{enumerate}
  \item every lift form at $(i,\mu_0)$ has exactly one rational lift;
  \item for every family tie of $R$ and every embedding $\sigma$ of $L$ that agrees with $\sigma_0$ on
    $K$, the sums
    $\sum_\rho\log\big(\mathrm{hook}(a^+_\rho(\sigma))\,\mathrm{hook}(a^-_\rho(\sigma))\big)$ at
    $\sigma$ and at $\sigma_0$ are equal, there is a conjugate datum at $\sigma$ for the lift $i$ and the
    tie, and
    there are a unitary operator $U$ of $H_W(i,\sigma)$ that multiplies the functions $U'(v)$ by a
    function of absolute value $1$, equal to $1$ if $\sigma=\sigma_0$, and an element $z_0$ of the
    center of $G$ that acts unitarily on $H_W(i,\sigma)$, such that for every conjugate datum at
    $\sigma$ for $i$ and the tie there is $\kappa\in\mathbf C$ with $|\kappa|=1$ and
    $U\big(FW_i(\sigma)(z)\big)=\kappa\,z_0\,\mathrm{th}(j(x))$ for every lift form $x$ with rational
    lift $z$;
  \item for every embedding $\sigma$ of $L$ that does not agree with $\sigma_0$ on $K$, the embedding
    $\bar\sigma$ agrees with $\sigma_0$ on $K$, and $\langle FW_i(\sigma)z,FW_i(\sigma)z'\rangle=
    \overline{\langle FW_i(\bar\sigma)z,FW_i(\bar\sigma)z'\rangle}$ for all $z,z'\in L^{n_W(i)}$.
  \end{enumerate}
\item (Conjugate forms.) If $R$ satisfies the core conditions at $F$, then for every family tie of $R$,
  every lift $i$, every embedding $\sigma$ of $L$ and every conjugate datum at $\sigma$ for $i$ and the
  tie, the image of $j$ spans $H$,
  and $\langle j(x),j(u)\rangle=\langle x,u\rangle_\sigma$ for all lift forms $x,u$ at $(i,\mu_0)$.
\end{enumerate}
\end{lemma}
\leanref{Development/AbelianVarietyTheory/Proofs/DivisorProofs/SplitPartBaseLinksOfPrimitiveRowsStatement.lean}{GlobalTheory.T.splitPart_baseLinks_of_primitiveRows}
\leanref{Development/AbelianVarietyTheory/Definitions/Divisors/RealisedChainTorusEq.lean}{GlobalTheory.BigOn.BaseLinks}
\leanref{Development/AbelianVarietyTheory/Proofs/DivisorProofs/BadPlacesPrimeToHostBaseTupleStatement.lean}{GlobalTheory.T.badPlacesPrimeToHost_baseTuple}
\leanref{Development/AbelianVarietyTheory/Proofs/DivisorProofs/ConjugatesOfRationalThetaLiftFromCharacterisatioStatement.lean}{GlobalTheory.T.conjugatesOfRationalThetaLift_from_characterisation}
\leanref{Development/AbelianVarietyTheory/Definitions/Divisors/FormalDegreeUpperC.lean}{GlobalTheory.LowRow.GalOut}
\leanref{Development/AbelianVarietyTheory/BlockTheory/DivisorBase/WeilGramSignature.lean}{GlobalTheory.IndefiniteRowSideFullHecke.WeilGramSignature.ConjAt}
\leanref{Development/AbelianVarietyTheory/Proofs/DivisorProofs/InnerConjFormsEqCsumAt.lean}{GlobalTheory.LiftLower.inner_conjForms_eq_csumAt}
\leanref{Development/AbelianVarietyTheory/Proofs/DivisorProofs/SpanEqTopOfTie.lean}{GlobalTheory.ConjSpan.span_conjForms}
\begin{proof}
(1) The weak characterization and the existence conditions for the theta maps give the
characterization of $R$ at $F$, whose conditions on the level characters, the group, the family, the
models, the chain and the theta lifts are the core conditions. The family conditions of the family
with its inert indices contain those of the underlying family. For the base links, fix a lift $j$
and a tuple $t$, write $\theta=\theta'_{j,t}=\mathrm{th}_G(j,t)\colon F_{j,t}\to W_j$ and $W_0$ for
the image of $w_0(j)$, let $\Gamma_0(j,t)$ be the level-$0$ part of the lift forms at $(j,t)$ and
$e=e_Xe_{P_0,\psi_0^{-1}}$ (Lemma~\ref{is-baselinks-definite-coordinates}), and let
$\Gamma\subset F_{j,t}$ be the span of $j_{j,t}(\Gamma_0(j,t))$.

\emph{Step 1: the group.} By the core conditions on the models, $j_{j,t}$ carries a basis of the
lift forms at $(j,t)$ to a basis of $F_{j,t}$, and $\langle j_{j,t}(x),j_{j,t}(y)\rangle$ is the
image under $\sigma_0$ of the class-sum pairing of $x$ and $y$. So each element of $G$ and each Hecke
operator on the lift forms induces an operator on $F_{j,t}$ that $j_{j,t}$ intertwines with it. By
the conditions of the lift data, $G$ preserves the class-sum pairing; by (3) of
Definition~\ref{is-primitive-rows}, $G$ acts unitarily on $W_j$ and $\theta$ is $G$-equivariant. So
Lemma~\ref{is-baselinks-group-algebra} applies to the lift forms, $j_{j,t}$, $F_{j,t}$, $W_j$ and
$\theta$, and likewise at every tuple. The characters of $G$ in the family are absolutely
irreducible, so $e_X$ is a central idempotent, and $e_{P_0,\psi_0^{-1}}$ is an idempotent; hence $e$
acts on the lift forms as an idempotent, and $\Gamma=e_V(F_{j,t})$ by
Lemma~\ref{is-baselinks-group-algebra}(1).

\emph{Step 2: rationality.} By the models and (5) of Definition~\ref{is-primitive-rows},
Lemma~\ref{is-baselinks-rational-tests}(2) applies to $j_{j,t}$ and $\theta$, with the class-sum
pairing as $h$. So $\theta^*(W_j)$ is spanned by $j_{j,t}(Z_{j,t})$.

\emph{Step 3: the first base link.} Apply Lemma~\ref{is-baselinks-lift-cut} to $\theta$, the vectors
$u_k=j_{j,t}(c_{0,j,t}(e_k))$, $\Gamma$, $W_0$, the parts $W_p(j)$ and
$\lambda_p=\lambda^{(3)}_{j,t,p,v_0}$. By the models and
Lemma~\ref{is-baselinks-definite-coordinates}(2), the Gram matrix of the $u_k$ is
$|P_0|\,(X^{-1})^*X^{-1}$ for the invertible matrix $X=X_T(j,t)_{v_0}$, which is positive definite;
so the $u_k$ are linearly independent. Hypothesis (1) of that lemma holds by
Lemma~\ref{is-baselinks-definite-coordinates}(1), Lemma~\ref{is-baselinks-group-algebra}(3) and
Step 2. For the first half of hypothesis (2), let $u\in\Gamma$. By
Lemma~\ref{is-baselinks-group-algebra}(2) there is $u'\in\theta^*(W_j)\cap\Gamma$ with
$\theta(u')=\theta(u)$, and $u'$ is a combination of vectors $j_{j,t}(x)$ with
$x\in Z_{j,t}\cap\Gamma_0(j,t)$. The image of $e$ lies in that of $e_X$, which is the level-$0$ part
for the face with the trivial subgroup and character; so by
Lemma~\ref{is-baselinks-definite-coordinates}(1) for that face each such $x$ is in the image of its
coordinate map, and $\theta(u)$ lies in the image of $t^{\mathrm{pr}}(j,t)$. Every $h\in P_0$ acts on
the image of $e$ by $\psi_0(h)^{-1}$, because $e_X$ is central and
$he_{P_0,\psi_0^{-1}}=\psi_0(h)^{-1}e_{P_0,\psi_0^{-1}}$; so $h$ acts on $\theta(u)$ by
$\sigma_0(\psi_0(h))^{-1}$, by Lemma~\ref{is-baselinks-group-algebra}(5). By (8) of
Definition~\ref{is-primitive-rows}, $\theta(u)\in W_0$. For the second half, let $w\in W_0$. By (8),
$P_0$ acts on $w$ through $\sigma_0\circ\psi_0^{-1}$, and $w$ is a linear combination of vectors
$\theta'_{j,t'}(j_{j,t'}(x))$ with $x$ in the image of $e_X$ on the lift forms at $(j,t')$. By
Lemma~\ref{is-baselinks-group-algebra}(1) at $t'$, $e_X$ fixes each of these vectors, so it fixes
$w$, and Lemma~\ref{is-baselinks-group-algebra}(4) with $e_X$ and $(P_0,\psi_0^{-1})$ gives
$\theta^*w\in e_V(F_{j,t})=\Gamma$. Hypothesis (3) holds by (9), (7) and (10) of
Definition~\ref{is-primitive-rows}. So $t_0(j,t)$, which is the map $t$ of
Lemma~\ref{is-baselinks-lift-cut}, is injective with image $W_0$.

\emph{Step 4: the second base link.} The map $e_{0,j}$ is $\sqrt{|P_0|}$ times an inclusion, and the
vectors $w_{0,j,v_0}(k)$ are $|P_0|^{-1/2}$ times the images under $FW_j(\sigma_0)$ of an $L$-basis
of a subspace of $L^{n_W(j)}$, so they are linearly independent by (11) of
Definition~\ref{is-primitive-rows}; hence $w_0(j)$ is injective. By Step 3, $t_0(j,t)(e_k)=w_0(j)(c)$
for some $c\in\mathbf C^{d_W(0,j)}$. By (5) of Definition~\ref{is-primitive-rows} there are
functionals separating the points of $W_j$ whose values at the vectors $w_0(j)(e_l)$ and at
$t_0(j,t)(e_k)=\theta(j_{j,t}(c_{0,j,t}(e_k)))$ lie in $\sigma_0(L)$. By
Lemma~\ref{is-baselinks-rational-tests}(1), $c\in\sigma_0(L)^{d_W(0,j)}$. So the matrix that
expresses $t_0(j,t)$ through $w_0(j)$ is the image under $\sigma_0$ of a matrix over $L$.

\emph{Step 5: the third base link.} Let $t=\mu_0$, and let $T_i$ be the operators on $F_{j,\mu_0}$
induced by the Hecke operators and the inert Hecke operators on the lift forms at $(j,\mu_0)$
(Step 1). They commute and preserve $\Gamma$, because the operators on the lift forms commute with
each other and with $G$ (family conditions), and they contain the adjoint of each of them, by (1) of
Definition~\ref{is-primitive-rows} and the isometry of $j_{j,\mu_0}$. By (2) of
Definition~\ref{is-primitive-rows} there are operators $T'_i$ on $W_j$ with $\theta T_i=T'_i\theta$,
and $\theta(\Gamma)\subset W_0$ by Step 3. Call a system $c$ of eigenvalues lifted if its joint
eigenspaces in $\Gamma$ and in $W_0$ are nonzero, as in Lemma~\ref{is-baselinks-lifted-parts} with
$\Gamma'=W_0$. By (6) of Definition~\ref{is-primitive-rows}, the joint eigenspace of every lifted
$c$ in $\Gamma$ lies in $\theta^*(W_j)$. So Lemma~\ref{is-baselinks-lifted-parts} gives
$(\ker\theta)^\perp\cap\Gamma=\sum_{c\text{ lifted}}A_c\cap\Gamma$.

Let $C$ be the set of systems $y$ of eigenvalues in $L$ of the operators on the lift forms at
$(j,\mu_0)$ for which $\sigma_0(y)$ is lifted. These operators split over $L$ and commute with $G$
(family conditions), so every $f\in\Gamma_0(j,\mu_0)$ is the sum of its components in their joint
eigenspaces, these components lie in $\Gamma_0(j,\mu_0)$, and $j_{j,\mu_0}$ carries the component for
$y$ into $A_{\sigma_0(y)}\cap\Gamma$. Now $f\in Z_{j,\mu_0}$ if and only if
$j_{j,\mu_0}(f)\in\theta^*(W_j)=(\ker\theta)^\perp$. Since joint eigenspaces of distinct systems are
independent and $j_{j,\mu_0}$ is injective, this holds if and only if the components of $f$ for
$y\notin C$ vanish. So $Z_{j,\mu_0}\cap\Gamma_0(j,\mu_0)$ is the intersection of $\Gamma_0(j,\mu_0)$
with the sum of the joint eigenspaces of the $y\in C$, and
Lemma~\ref{is-baselinks-definite-coordinates}(3), with (1) of Definition~\ref{is-primitive-rows},
gives the third base link.
% (1) is GlobalTheory.LowRow.baseLinks_of_prims: RealDat3.characterisation_of_weak gives Char3, and
% complex_data_conditions_base_level_of_primitive_links builds ComplexDataConditionsBaseLevel from
% PrimRows3; lifts0 (CxData.cutLifts), ratTh (Rat.ratThOfLiftTest) and proviso0
% (proviso_of_lift_eigenvalues) give H5, RATth and Proviso, after CharRealise.toCx3_eq_one for the
% phases 1. Step 1 is CxPass.lift_hypotheses_of_unitary (kFormL_of, jIsoL_of from the models);
% Step 2 is Rat.ratZOfTests; Step 3 is cutLifts with lift_isometry_of_norms,
% lift_rationality_of_equivariance, lift_descent_at_lower_type, upper_level_lift_at_lower_type,
% thmAop_of and sumW; Step 4 is ratThOfLiftTest with RealDat.ewNorm and RealDat.kForm; Step 5 is
% U2.eigCOf (liftsFull_adj), U2.eig0Of and proviso_of_lift_eigenvalues (eig_ratG, xq_rat).

(2) By the core conditions on the level characters, $Q_{\mu_0}$ is empty and every place of $P_g$ lies
over one of the primes $g_0,g_1$ of the host. These satisfy $Q_Sqd\equiv1\pmod{g_k}$ and are prime to
$\mathrm N(\mathfrak c_K)$, so no place of $P_g$ contains $Q_Sqd\,\mathrm N(\mathfrak c_K)$.
% (2) is LowRow.primeToHost_base with LiftScalarCompare.primeToHost and LiftLower.Q_base.

(3) Fix the lift $i$ and a family tie of $R$, the given one in (3b). The hypotheses of (3) contain
those of Lemma~\ref{is-eigenvalue-conjugate-inputs}; use its notation: the core forms $f$ with the
action $f\mapsto gf$ of $G$, the vectors $y(f)\in L^{n_W(i)}$, the character $\chi_W$ of $FW_i$, the
functions $\mathrm{ad}_\sigma(f)$, class sums and theta integrals. Write $e_C(g)\in L$ for the
coefficients of the central idempotent $e_C$, which is $e_X$, so that $e$ is $\sum_g\sigma(e_C(g))g$
on $H_W(i,\sigma)$. The weight tie makes each vector $v_0(i,\sigma)$ extremal, so a tensor of unit
vectors. The argument uses the following facts.
\begin{enumerate}
\item[(0)] If $\sigma$ agrees with $\sigma_0$ on $K$, the function of $FW_i(\sigma)(y(f))$ at
  $v_0(i,\sigma)$ is $\mathrm{ad}_\sigma(f)$ for every core form $f$. If $\sigma$ does not, there is a
  tensor $v$ of unit vectors of the models at $\sigma$ such that the function of $FW_i(\sigma)(y(f))$
  at $v$ is $\overline{\mathrm{ad}_{\bar\sigma}(f)}$ for every core form $f$.
\item[(i)] $G$ acts on the core forms, and
  $FW_i(\sigma_0)(y(gf))=\sigma_0(\chi_W(g))\,g\,FW_i(\sigma_0)(y(f))$.
\item[(ii)] There are a central element $z_1$ of $G$ and a root of unity $\zeta\in L$ such that every
  lift form $x$ at $(i,\mu_0)$ is a class sum $\sum_ka_kf_k$ with
  $\phi\,\mathrm{th}_G(i,\mu_0)(j_V(x))=\sum_g\sigma_0(e_C(g))\,g\sum_k\sigma_0(\zeta a_k)\,
  FW_i(\sigma_0)(y(z_1f_k))$.
\item[(iii)] Two conjugate data at an embedding $\sigma$ that agrees with $\sigma_0$ on $K$, for $i$
  and the tie, with maps $j,\mathrm{th}$ and $j',\mathrm{th}'$, satisfy
  $\mathrm{th}'(j'(x))=\kappa'\,\mathrm{th}(j(x))$ for all $x$, for some $\kappa'$ with
  $|\kappa'|=1$. The characters $\mu^0_i,\eta$, the coefficients of $R$, the form space
  $H_V(i,\mu_0)$, the map $j_V$ and $\mathrm{th}_G(i,\mu_0)$ are a conjugate datum at $\sigma_0$.
\item[(iv)] Let $\sigma\ne\sigma_0$ agree with $\sigma_0$ on $K$. There are a conjugate datum at
  $\sigma$ for $i$ and the tie, with maps $j$ and $\mathrm{th}$, a function $u$ of absolute value $1$
  with a unitary operator $U$ of $H_W(i,\sigma)$ that multiplies the functions by $u$, a central
  element $z_2$ of $G$ and $\kappa$ with $|\kappa|=1$, such that the identity
  $U(FW_i(\sigma)(y(gf)))=\sigma(\chi_W(g))\,g\,U(FW_i(\sigma)(y(f)))$ holds for all $g$ and $f$,
  and for every lift form $x$ that is a class sum $\sum_ka_kf_k$, $\kappa$ times the theta integral
  of $j(x)$ at the vector $v_0(i,\sigma)$ is the function of
  $\sum_k\sigma(a_k)\,U(FW_i(\sigma)(y(z_2f_k)))$ at $v_0(i,\sigma)$. An element of
  $H_W(i,\sigma)$ is determined by its function at $v_0(i,\sigma)$, and $e$ commutes with taking
  functions at $v_0(i,\sigma)$.
\item[(v)] If $\sigma$ does not agree with $\sigma_0$ on $K$, then
  $\langle FW_i(\sigma)z,FW_i(\sigma)z'\rangle=
  \overline{\langle FW_i(\bar\sigma)z,FW_i(\bar\sigma)z'\rangle}$ for all $z,z'\in L^{n_W(i)}$.
\item[(vi)] Every element of $G$ acts unitarily on $H_W(i,\sigma)$, for every $\sigma$.
\end{enumerate}

\emph{Proof of (0).} Let $\sigma$ agree with $\sigma_0$ on $K$. Then $\sigma$ is good, by
Lemma~\ref{is-eigenvalue-conjugate-good}(3); let $J$ be the map of
Lemma~\ref{is-eigenvalue-conjugate-inputs}(A). A $K_s$-basis of the image of $y$ is an $L$-basis of
$L^{n_W(i)}$, so the $\sigma$-semilinear map $J'$ on $L^{n_W(i)}$ that takes the values of
$J\circ y^{-1}$ on this basis agrees with $J\circ y^{-1}$ on its $K_s$-span, and $J'(y(f))=J(f)$ for
every core form $f$. The weight tie requires that the function of $FW_i(\sigma)(y(f))$ at
$v_0(i,\sigma)$ be $\mathrm{ad}_\sigma(f)$ for every $f$ as soon as there is a $\sigma$-semilinear map
$J'$ with this property in place of $FW_i(\sigma)$, at every good $\sigma$. If $\sigma$ does not agree
with $\sigma_0$ on $K$, then $\sigma$ is not good, by Lemma~\ref{is-eigenvalue-conjugate-good}(3), and
the conjugation tie requires the second assertion, at every embedding that is not good, as soon as
there are $v$ and $J$ as in the second part of Lemma~\ref{is-eigenvalue-conjugate-inputs}(A).

\emph{Proof of (i).} Let $g\in G$, let $f$ be a core form, write $v_0=v_0(i,\sigma_0)$, and let $k$
be a level lift of $g$ (Lemma~\ref{is-eigenvalue-core-group}(1)). The one-map data at $i$ give the
lift as one map, so they are equivariant: $gf$ is a core form, and
$\mathrm{ad}_{\sigma_0}(gf)(h)=\chi(g)\,\mathrm{ad}_{\sigma_0}(f)(hk)$ for every $h$, where
$\chi(g)=\prod_{P\in P_g}\chi_{i,P}(\det g_P)$ for the characters $\chi_{i,P}$ of the residue fields
of the realization. The one-map data state this for their own functions of the core forms, which are
the $\mathrm{ad}_{\sigma_0}(f)$ divided by a constant of the core. By the core conditions on the
group, $\chi(g)=\sigma_0(\chi_W(g))$. By (0) and Lemma~\ref{is-eigenvalue-core-group}(1), for every
$h$,
\[
 FW_i(\sigma_0)(y(gf))(v_0)(h)=\sigma_0(\chi_W(g))\,\mathrm{ad}_{\sigma_0}(f)(hk)
 =\sigma_0(\chi_W(g))\,\big(g\,FW_i(\sigma_0)(y(f))\big)(v_0)(h).
\]
So the two sides of (i) have the same function at $v_0$, and they are equal by
Lemma~\ref{is-eigenvalue-core-group}(3).

\emph{Proof of (ii).} Write $v_0=v_0(i,\sigma_0)$. The one-map data at $i$ give the lift as one map
with the canonical constants of the core and the phase $\phi$. Their one-map tie provides a central
element $z_1$ of $G$ and a root of unity in $\sigma_0(L)$, which is $\sigma_0(\zeta)$ for a root of
unity $\zeta\in L$, with the following property. Let $x'$ be a theta class datum of the one-map data:
a family of values at the class representatives that comes from a lift form. Then there are core forms
$f_k$ and coefficients $a_k$, which lie in $L$ because the pins identify the coefficient field of the
one-map data with $\sigma_0(L)$, such that the section of the core built from $x'$ is
$\sum_ka_kf_k$ after extension of scalars, and
\[
 \phi\,\vartheta(x')(q)=\sum_k\sigma_0(\zeta a_k)\,\mathrm{ad}_{\sigma_0}(z_1f_k)(q)
\]
at a representative $q$ of every point of $U(W_i)(F_h)\backslash U(W_i)(\mathbf A)$, where
$\vartheta(x')$ is the theta function of $x'$ that the one-map data give. The $z_1f_k$ are core forms
by (i). Let $x$ be a lift form at $(i,\mu_0)$. By the pins, the values of $j_V(x)$ at the class
representatives form a theta class datum $x'$, and the function of $\mathrm{th}_F(i,\mu_0)(j_V(x))$
at $v_0$ is $\vartheta(x')$; so $x$ is the class sum $\sum_ka_kf_k$. By the core conditions on the
theta lifts, $\mathrm{th}_F(i,\mu_0)(j_V(x))$ is the family of functions of some $Y\in
H_W(i,\sigma_0)$, and the functions of $\mathrm{th}_G(i,\mu_0)(j_V(x))$ are those of $Y$ with $e$
applied. Put $Y'=\sum_k\sigma_0(\zeta a_k)\,FW_i(\sigma_0)(y(z_1f_k))$. By (0), the functions of
$\phi Y$ and $Y'$ at $v_0$ agree at a representative of every point of
$U(W_i)(F_h)\backslash U(W_i)(\mathbf A)$. Both are invariant on the left under $U(W_i)(F_h)$, so they
are equal, and $\phi Y=Y'$ by Lemma~\ref{is-eigenvalue-core-group}(3). By the core conditions on the
models, $G$ acts on the functions of $H_W(i,\sigma_0)$ as it acts on $H_W(i,\sigma_0)$, and an
element is determined by its functions; so
$\phi\,\mathrm{th}_G(i,\mu_0)(j_V(x))=\sum_g\sigma_0(e_C(g))\,g\,Y'$, which is (ii).

\emph{Proof of (iii).} First the conjugate datum at $\sigma_0$. At every embedding $\sigma$ that agrees
with $\sigma_0$ on $K$ the weight tie gives a $\sigma$-conjugate of $\mu^0_i$, which it requires to be
$\mu^0_i$ when $\sigma$ agrees with $\sigma_0$ on $K_s$; so $\mu^0_i$ is a $\sigma_0$-conjugate of
itself. By the core conditions on the level characters, $\eta$ is the character $\eta$ of the host,
an unramified Hecke character. Its values at uniformizers, divided by the square roots of the norms,
lie in $\sigma_0(L)$: at the split places outside a finite set by the core conditions on the family; at
an inert place, where the value is $-1$ divided by the norm of the place of $F_h$ below; and at each
remaining place $P$ by the product formula for $\eta$ on the principal ideal of an element of
$\mathcal O_{E_h}$ of valuation $1$ at $P$ that is a unit at the other remaining places, since the
archimedean component of $\eta$ at that element is the inverse of a product of its conjugates, which
lie in $\sigma_0(L)$. So $\eta$ is a $\sigma_0$-conjugate of itself. By the core conditions on the models, the
coefficients of $R$ form a
coefficient datum, and the coefficient map $c_0$ of the tie is a $\sigma_0$-semilinear isomorphism
onto it that intertwines the rational action of $U(V_i)(F_h)$ and, by the class presentation of the
tie, carries the rational pairing, followed by $\sigma_0$, to the inner product. The space
$H_V(i,\mu_0)$ is a form space for the level of $V_i$ at $\mu_0$ with the level character of the
realization, whose values lie in $\sigma_0(L)$. The class presentation writes the value of $j_V(x)$
at $h$ as a value of this character times the image under $c_0$ of a rational vector moved by an
element of $U(V_i)(F_h)$, which is again the image under $c_0$ of a rational vector; so (4) of
Definition~\ref{is-eigenvalue-conjugate-datum} holds with $j=j_V$. Finally, the core conditions on the
theta lifts give the Weil datum of $R$ at $(i,\mu_0)$, with characters $(\mu^0_i,\mu^0_i\eta^{-1})$
since $\xi_{\mu_0}$ is trivial, its level vector, of the trivial level since $Q_{\mu_0}$ is empty, the
bases of the models, the map $\mathrm{io}_P(i)$, whose delta polynomial is $\Delta_i(\sigma_0)$, and the
theta integral of $R$, and they give $\mathrm{th}_G(i,\mu_0)$ as $e$ applied to it. This is (5) of
Definition~\ref{is-eigenvalue-conjugate-datum} with $\mathrm{th}=\mathrm{th}_G(i,\mu_0)$.

Now let two conjugate data at $\sigma$ be given, with characters $\chi',\eta'$ and $\chi'',\eta''$ and
theta integrals $\Theta,\Theta'$. At a finite idele where $\mu^0_i$ takes the value $\sigma_0(a)$,
$a\in L$, both $\chi'$ and $\chi''$ take the value $\sigma(a)$, and $a$ is unique; so
$\chi'\chi''^{-1}$ is a Hecke character trivial on the finite ideles. Such a character is trivial,
because $E_h$ is dense in $E_h\otimes\mathbf R$; so $\chi'=\chi''$. In the same way $\eta'\eta''^{-1}$
is an unramified Hecke character that is trivial at every uniformizer, hence trivial, and
$\eta'=\eta''$. By Lemma~\ref{is-eigenvalue-conjugate-inputs}(B) there is $\kappa'$ with $|\kappa'|=1$
and $\Theta'(j'(x))=\kappa'\,\Theta(j(x))$ for every lift form $x$. Applying $e$, the functions of
$\mathrm{th}'(j'(x))$ and $\kappa'\,\mathrm{th}(j(x))$ at $v_0(i,\sigma)$ agree, so they are equal by
Lemma~\ref{is-eigenvalue-core-group}(3).

\emph{Proof of (iv).} Take the conjugate datum, $u$, $U$, $z_2$ and $\kappa$ of
Lemma~\ref{is-eigenvalue-conjugate-inputs}(C), and write $v_0=v_0(i,\sigma)$. By (0), the function of
$U(FW_i(\sigma)(y(f)))$ at $v_0$ is $u\,\mathrm{ad}_\sigma(f)$ for every core form $f$. Let $g\in G$,
with a level lift $k$. By (C)(a) and Lemma~\ref{is-eigenvalue-core-group}(1), for every $h$,
\[
 U(FW_i(\sigma)(y(gf)))(v_0)(h)=u(h)\,\mathrm{ad}_\sigma(gf)(h)
 =\sigma(\chi_W(g))\,u(hk)\,\mathrm{ad}_\sigma(f)(hk)
 =\sigma(\chi_W(g))\,\big(g\,U(FW_i(\sigma)(y(f)))\big)(v_0)(h),
\]
and the identity of (iv) follows by Lemma~\ref{is-eigenvalue-core-group}(3). For a lift form $x$ that
is a class sum $\sum_ka_kf_k$, (C)(b) gives
$\kappa\,\Theta(j(x))(v_0)=\sum_k\sigma(a_k)\,u\,\mathrm{ad}_\sigma(z_2f_k)$, which is the function of
$\sum_k\sigma(a_k)\,U(FW_i(\sigma)(y(z_2f_k)))$ at $v_0$. An element of $H_W(i,\sigma)$ is determined by
its function at $v_0$ by Lemma~\ref{is-eigenvalue-core-group}(3). By the core conditions on the models,
$G$ acts on the functions at $v_0$ as it acts on $H_W(i,\sigma)$, so $e$, a combination of elements
of $G$, commutes with taking functions at $v_0$.

\emph{Proof of (v).} Let $\sigma$ not agree with $\sigma_0$ on $K$. Then $\bar\sigma$ does, because
$K$ is imaginary quadratic. By (0) at $\bar\sigma$ and at $\sigma$, there is a tensor $v$ of unit
vectors at $\sigma$ such that, for every core form $f$, the function of $FW_i(\sigma)(y(f))$ at $v$ is
the complex conjugate of the function of $FW_i(\bar\sigma)(y(f))$ at $v_0(i,\bar\sigma)$. The
embeddings $\sigma$ and $\bar\sigma$ have the same labels, so $N_\sigma=N_{\bar\sigma}$, a positive
integer. By Lemma~\ref{is-eigenvalue-core-group}(3), at $\sigma$ with $v$ and at $\bar\sigma$ with
$v_0(i,\bar\sigma)$,
$\langle FW_i(\sigma)y(f),FW_i(\sigma)y(f')\rangle=
\overline{\langle FW_i(\bar\sigma)y(f),FW_i(\bar\sigma)y(f')\rangle}$ for all core forms $f,f'$.
Both sides of (v) are additive in $z$ and in $z'$, and they change in the same way when $z$ or $z'$ is
multiplied by $a\in L$, because $\sigma(a)=\overline{\bar\sigma(a)}$. The vectors $y(f)$ span
$L^{n_W(i)}$ over $L$, so (v) holds for all $z,z'$.

\emph{Proof of (vi).} This is Lemma~\ref{is-eigenvalue-core-group}(2).
% (3) is GlobalTheory.ConjRatLiftFive.conjRatLift5 (base_lift5, conjAt_base5,
% conjugateRationalLift_at_nonBaseEmbedding, conj_lift5 with GalLift.Assembly.read_sum and
% conj_abstract, conjugateRationalLift_at_embedding_of_datum, lift_unique5); ConjAt also contains
% HypConj, not stated here because the proofs below do not use it.
% (0) is FWTie.evaluates (exists_ext_of_isKsForm and the field ev of the weight tie) and the bar tie
% BarTie, from ConjReading.jEvExists_and_jEvBar (Lemma is-eigenvalue-conjugate-inputs(A)).
% (i) is GalLift.BaseEquiv.base_equiv_five = reading_actG_five: act_eq_translation_five, EquivAt
% (first clause of OneMapOf), ev_base (Evaluates with phi_cE), hc.grp.chiW, alg_step and EvInj.
% (ii) is GalLift.BaseId.base_identity_five: tie_baseChange (the one-map tie, second clause of
% OneMapOf), the pins (field, RepsAt from exists_reps), thF_mem and thG of the core theta lifts,
% phi_cE with Evaluates, eq_of_reps (EvInj and mult_leftInv) and proj_model.
% (iii) is GalLift.BaseDatum.isConjDatum_base_five (fw.chiR at sg0 with the last clause of ChiAt,
% isEtaConj_base with EtaVals = EtaValues.etaVals_five, clause 2 of
% conjugateCharacters_embeddings_of_characterisation: hv_mem_host with the family tie's etaVal,
% hv_mem_of_inert and hv_mem through prod_hv; isConjForms_base_five with pair_val, ext_val and
% ConjForms.levChar_mem, isConjLift_base_five) and ConjRatLiftFive.lift_unique5
% (ConjCharUnique.conjChar_unique_ofFive: eq_of_conj_on_finite, UnramCharUnique.isEtaConj_unique;
% lift_unique_five, Lemma is-eigenvalue-conjugate-inputs(B)).
% (iv) is ConjRatLiftFive.conj_lift5 for the datum of conjugateRationalLift_at_nonBaseEmbedding
% (Lemma is-eigenvalue-conjugate-inputs(C)): hJ, hJa and hJe for the identity, hS for the theta
% integral, EvInj and Transl.proj.
% (v) is ConjInner.inner_reading_conj_five (clause 3 of centralShift_phase_conjReadings_fiveClauses):
% good_resC, inner_conj_of_values (inner_eq_integral, dim_conj, conj_integral) and inner_conj_of_span
% with fw.ks. (vi) is gW_norm_eq with CentralShift.inner_eq_of_norm; the Lean uses it for central
% elements, through IsCentralShift.unitary.

(3a) Let $x=\sum_ka_kf_k$ be a lift form, as in (ii), and put
$z_x=\sum_ge_C(g)\chi_W(g)^{-1}\sum_k\zeta a_k\,y(gz_1f_k)$. By (i) and the semilinearity of
$FW_i(\sigma_0)$,
$FW_i(\sigma_0)(z_x)=\sum_g\sigma_0(e_C(g))\,g\sum_k\sigma_0(\zeta a_k)FW_i(\sigma_0)(y(z_1f_k))$,
which is $\phi\,\mathrm{th}_G(i,\mu_0)(j_V(x))$ by (ii). So $z_x$ is a rational lift of $x$. It is
the only one, because $FW_i(\sigma_0)$ is injective by (11) of Definition~\ref{is-primitive-rows}.

(3b) Let $\sigma$ agree with $\sigma_0$ on $K$. Then $\sigma$ permutes the labels of the embeddings
of $E_h$, and the sum of the logarithms of the hook products is invariant under this permutation; so
the sums at $\sigma$ and $\sigma_0$ are equal. If $\sigma=\sigma_0$, the base conjugate datum of
(iii) exists; take $U=1$ and $z_0=1$. For a conjugate datum with maps $j'$ and $\mathrm{th}'$, (iii)
gives $\kappa'$ with $\mathrm{th}'(j'(x))=\kappa'\,\mathrm{th}_G(i,\mu_0)(j_V(x))$, so
$FW_i(\sigma_0)(z)=\phi\kappa'^{-1}\mathrm{th}'(j'(x))$ for every lift form $x$ with rational lift
$z$, and $|\phi\kappa'^{-1}|=1$ because the one-map data at $i$ have phase of absolute value $1$.

Let $\sigma\ne\sigma_0$. Take the conjugate datum at $\sigma$, with maps $j$ and $\mathrm{th}$, and
$u$, $U$, $z_2$ and $\kappa$ of (iv). For $x=\sum_ka_kf_k$ with rational lift $z_x$ as in (3a), the
identity for $U$ in (iv) gives
\[
 U(FW_i(\sigma)(z_x))=\sum_g\sigma(e_C(g))\,g\sum_k\sigma(\zeta a_k)\,U(FW_i(\sigma)(y(z_1f_k)))
 =\sigma(\zeta)\,\sigma(\chi_W(z_1))\,\sigma(\chi_W(z_2))^{-1}\,z_1z_2^{-1}
 \sum_g\sigma(e_C(g))\,g\sum_k\sigma(a_k)\,U(FW_i(\sigma)(y(z_2f_k))),
\]
since $z_1$ and $z_2$ are central. The last sum over $g$ is $e$ applied to
$\sum_k\sigma(a_k)U(FW_i(\sigma)(y(z_2f_k)))$, whose function at $v_0(i,\sigma)$ is $\kappa$ times
the theta integral of $j(x)$ by (iv). Applying $e$, which gives $\mathrm{th}(j(x))$ from the theta
integral, the last sum and $\kappa\,\mathrm{th}(j(x))$ have the same function at $v_0(i,\sigma)$, so
they are equal by (iv). Hence $U(FW_i(\sigma)(z_x))=\kappa_1z_0\,\mathrm{th}(j(x))$ with
$z_0=z_1z_2^{-1}$, central and acting unitarily by (vi), and
$\kappa_1=\sigma(\zeta)\sigma(\chi_W(z_1))\sigma(\chi_W(z_2))^{-1}\kappa$, of absolute value $1$ as
$\zeta$ and the values of $\chi_W$ are roots of unity. For another conjugate datum at $\sigma$, (iii)
changes $\kappa_1$ by a factor of absolute value $1$.

(3c) The subfield $K=\mathbf Q(\sqrt{-d})$ of the host is imaginary quadratic, so if $\sigma$ does
not agree with $\sigma_0$ on $K$, then $\bar\sigma$ does. The identity is (v).

(4) Fix a family tie, a lift $i$, an embedding $\sigma$ and a conjugate datum at $\sigma$, with the
coefficient datum $W_c$, the $\sigma$-semilinear isomorphism $c_1$ of
Definition~\ref{is-eigenvalue-conjugate-datum}(2), the form space $H$ and the map $j$. The family tie
presents the lift forms at $(i,\mu_0)$ by their values at representatives $y$
of the double cosets of their level, with stabilizers $\Gamma_y$: a lift form $x$ has values $x(y)$
in the rational coefficient space of $V_i$, the class-sum pairing is
$\sum_y|\Gamma_y|^{-1}\langle x(y),x'(y)\rangle$ for the rational pairing, and the lift forms have a
basis $b_1,\dots,b_n$ whose values have coordinates in $K_0$ in the rational basis of the coefficient
space, because the stabilizers act on that basis by matrices over $K_0$ and the character of the
level takes roots of unity as values. The map $c_0$ of the tie is injective, and $j_V$ is the map of
the models of the core conditions, so (4) of Definition~\ref{is-eigenvalue-conjugate-datum} gives
$j(x)(y)=c_1(x(y))$.

\emph{Inner products.} By the inner product of $H$,
$\langle j(x),j(x')\rangle=\sum_y|\Gamma_y|^{-1}\langle c_1(x(y)),c_1(x'(y))\rangle$. For vectors
$a,a'$ of the coefficient space with coordinates in $K_0$, $\langle c_1(a),c_1(a')\rangle$ is the
image under $\sigma$ of their rational pairing, by (2) of
Definition~\ref{is-eigenvalue-conjugate-datum} and because $\sigma(\bar b)=\overline{\sigma(b)}$ for
$b\in K_0$. So $\langle j(b_k),j(b_l)\rangle=\sigma(H_{kl})$ for the matrix $H$ of the class-sum
pairing in the basis $(b_k)$. As $j$ is $\sigma$-semilinear,
$\langle j(x),j(x')\rangle=\sigma(\xi)^*\sigma(H)\sigma(\xi')$ for the coordinate vectors $\xi,\xi'$
of $x,x'$ in this basis, which is the image $\langle x,x'\rangle_\sigma$ under $\sigma$ of the
class-sum pairing.

\emph{Spanning.} The stabilizer $\Gamma_y$ acts on the rational coefficients through $U(V_i)(F_h)$
and the character of the level, and $c_1$ carries this action to the action on $W_c$ twisted by the
character of $H$, by (2) and (3) of Definition~\ref{is-eigenvalue-conjugate-datum}. For a
representative $y$ and a rational coefficient vector $a$, let $x_{y,a}$ be the lift form that
vanishes at the other representatives and has value $\sum_{\gamma\in\Gamma_y}\gamma a$ at $y$. Then
$j(x_{y,a})$ vanishes at the other representatives and has value
$\sum_{\gamma}\chi'(\gamma)\gamma\,c_1(a)$ at $y$, where $\chi'$ is the character of $H$. Since the
image of $c_1$ spans $W_c$, the span of the image of $j$ contains, for each $y$ and each $w\in W_c$,
an element $\varphi_{y,w}$ that vanishes at the other representatives and has value
$\sum_\gamma\chi'(\gamma)\gamma w$ at $y$. An element $\varphi$ of $H$ is determined by its values
at the representatives, and $\chi'(\gamma)\gamma\varphi(y)=\varphi(y)$ for $\gamma\in\Gamma_y$; so
$\varphi=\sum_y\varphi_{y,|\Gamma_y|^{-1}\varphi(y)}$. Hence the image of $j$ spans $H$.
% (4) is LiftLower.inner_conjForms_eq_csumAt (inner_eq_csumAt_of_pres, with the rational basis from
% HeckeRational.dcOp_matrix_mem_cmMax and the family tie Bridge.FamTie.coefL) and
% ConjSpan.span_conjForms (span_eq_top_of_tie, span_eq_top), from R.models_jV.
\end{proof}

\begin{lemma}\label{is-eigenvalue-base-place}
Let $D$ be a definite-side datum, let $R$ be a realization of $D$ with a face $F$ that satisfies the
core conditions, and let $b$ be real. Suppose that the family of definite forms satisfies the family
conditions of Lemma~\ref{is-eigenvalue-base-inputs}(1), that the conclusions of
Lemma~\ref{is-eigenvalue-adelic} hold for all lift data for $R$ at $F$ whose lift has the sizes of the
size data with some real $b'$, with all their measures and form spaces, that the base tuple $\mu_0$ is
prime to the host, and that every lift has the sizes of the size data with $b$. Let $i$ be a lift, put
$\theta_i=\mathrm{th}_G(i,\mu_0)\colon H_V(i,\mu_0)\to H_W(i,\sigma_0)$, and let $\mathrm{low}(i)$ be as
in Lemma~\ref{is-hook-lower-bound}. If $x\ne0$ and $l>0$ satisfy $\theta_i\theta_i^*x=lx$, then $\log
l\ge\mathrm{low}(i)$.
\end{lemma}
\leanref{Development/AbelianVarietyTheory/Proofs/DivisorProofs/LiftAdjointEigenvalueLowerBoundBaseEmbeddingStatement.lean}{GlobalTheory.T.liftAdjoint_eigenvalue_lower_bound_base_embedding}
\leanref{Development/AbelianVarietyTheory/Definitions/Divisors/RealisedChainTorusEq.lean}{GlobalTheory.BigOn.LowTop}
\begin{proof}
By the core conditions on the theta lifts, the full theta map $\mathrm{th}_F(i,\mu_0)$ takes values in
the functions of elements of $H_W(i,\sigma_0)$, so it is given by a linear map $\tilde\theta\colon
H_V(i,\mu_0)\to H_W(i,\sigma_0)$, and $\theta_i=e\tilde\theta$ with $e$ as in
Definition~\ref{is-eigenvalue-conjugate-datum}. The operator $e$ is an orthogonal projection: $e^2=e$
because $e_C$ is a central idempotent, and $e^*=e$ because $G$ acts unitarily (core conditions on the
group and the models) and every central idempotent of the group algebra is a sum of the self-adjoint
idempotents of the irreducible characters.

The core conditions also show that $\tilde\theta$ is a theta lift for the lift data
(Definition~\ref{is-eigenvalue-lift-data}) with the lift $i$, the base tuple $\mu_0$, the embedding
$\sigma_0$ and $s$ the identity, the trivial level, the characters $\mu^0_i\xi_{\mu_0}$ and $\eta$, the
residue characters at the places of $P_g$ given by the level characters of the realization, and the
Weil datum, level vector, bases, coefficients, map $\mathrm{io}$, measure and form space of the
realization at $(i,\mu_0)$; here $Q_{\mu_0}$ is empty, and the core conditions on the level characters
give the residues of $\mu^0_i\xi_{\mu_0}$ at $P_g$. These lift data are defined because $\mu_0$ is
prime to the host. By Lemma~\ref{is-eigenvalue-type-parts} their type parts are the
spaces $H_W(i,\sigma_0)_c$, and the level is trivial, so $|H_{\mathcal P}|=1$. By
Lemma~\ref{is-eigenvalue-adelic}(2), for every normalized $c$ with $H_W(i,\sigma_0)_c\ne0$,
$\tilde\theta\tilde\theta^*$ is multiplication by some $l_c>0$ on $H_W(i,\sigma_0)_c$ with
\[
 \log l_c\ge\sum_\rho\log\big(\mathrm{hook}(a^+_\rho(\sigma_0))\,\mathrm{hook}(a^-_\rho(\sigma_0))\big)
 -8fn_V^2L_p\ge\mathrm{low}(i),
\]
since $\mathrm{low}(i)=\sum_\rho\log\mathrm{hook}(a^+_\rho(\sigma_0))-8fn_V^2L_p$ and every hook product
is at least $1$. On the zero spaces $H_W(i,\sigma_0)_c$ any scalar $l_c\ge e^{\mathrm{low}(i)}$ will do.

By Lemma~\ref{is-eigenvalue-type-parts}(1) these spaces span $H_W(i,\sigma_0)$, so
Lemma~\ref{is-eigenvalue-projection}(1), for $T=\tilde\theta$ and $m=e^{\mathrm{low}(i)}$, gives
$e^{\mathrm{low}(i)}\|u\|^2\le\|\tilde\theta^*u\|^2$ for every $u$. If $x\ne0$ and $l>0$ satisfy
$\theta_i\theta_i^*x=lx$, Lemma~\ref{is-eigenvalue-projection}(2) gives $e^{\mathrm{low}(i)}\le l$, that
is, $\log l\ge\mathrm{low}(i)$.
\end{proof}

\begin{lemma}\label{is-eigenvalue-other-places}
Let $D$ be a definite-side datum, let $R$ be a realization of $D$ with a face $F$ that satisfies the
core conditions, and let $b$ be real. Suppose that the family of definite forms satisfies the family
conditions of Lemma~\ref{is-eigenvalue-base-inputs}(1); that the conclusions of
Lemma~\ref{is-eigenvalue-adelic} hold for all lift data for $R$ at $F$ whose lift has the sizes of the
size data with some real $b'$, with all their measures and form spaces; that the conclusions of
Lemma~\ref{is-eigenvalue-base-inputs}(3) hold with phase $1$ at every lift; that for every lift $i$ and
every embedding $\sigma$ of $L$ the model of $W_i$ at $\sigma$ has an extremal vector
(Definition~\ref{is-bounded-rows}(1)); that the rational structures $FW_i$ satisfy (11) of
Definition~\ref{is-primitive-rows}; that the rows $S_0(R,F)$ have the base links of
Lemma~\ref{is-eigenvalue-base-inputs}(1); that the base tuple $\mu_0$ is prime to the host; and that
every lift has the sizes of the size data with $b$. Let $i$ be a lift and $\mathrm{low}(i)$ as in
Lemma~\ref{is-hook-lower-bound}. Then for every infinite place $v\ne v_0$ of $L$ and every $y\in
L^{d_Z(i,\mu_0)}$,
\[
 e^{\mathrm{low}(i)}\,|y|^2_{X_T(i,\mu_0),v}\le|w_v(\theta_{i,\mu_0}y)|^2,
\]
where $\theta_{i,\mu_0}$ and $X_T(i,\mu_0)$ are those of the rows and $w_v$ is the map to
$W_{0,i,v}$ given by the vectors $w_{0,i,v}(k)$ of Definition~\ref{is-complex-data}(4).
\end{lemma}
\leanref{Development/AbelianVarietyTheory/Proofs/DivisorProofs/RationalLiftLowerBoundOffBasePlaceStatement.lean}{GlobalTheory.T.rationalLift_lowerBound_off_basePlace}
\leanref{Development/AbelianVarietyTheory/Definitions/Divisors/RealisedChainTorusEq.lean}{GlobalTheory.BigOn.LowForm}
\begin{proof}
Write $\langle\ ,\ \rangle_\sigma$ and $e$ as in Definition~\ref{is-eigenvalue-conjugate-datum}, fix a
family tie of $R$, which exists by the core conditions on the family, and let $n_P=|P_0|$ be the order
of the subgroup of the face at the rows.

\emph{Reduction.} Let $v\ne v_0$, let $\sigma_v$ be its embedding, let $y\in L^{d_Z(i,\mu_0)}$, let $x$
be the lift form $c_{0,i,\mu_0}(y)$ with coordinates $y$ in the face, and let $z\in L^{n_W(i)}$ be
the vector of
coordinates of $\theta_{i,\mu_0}y$ in the rational structure $FW_i$. By the first two base links,
$FW_i(\sigma_0)(z)=\mathrm{th}_G(i,\mu_0)(j_V(x))$, so $z$ is the rational lift of $x$ with phase $1$.
By the third base link, $|y|^2_{X_T(i,\mu_0),v}=n_P^{-1}\langle x,x\rangle_{\sigma_v}$, and
$\langle x,u\rangle_{\sigma'}=0$ for every embedding $\sigma'$ and every lift form $u$ with
$\mathrm{th}_G(i,\mu_0)(j_V(u))=0$, because $x$ lies in the definite subspace $Z_{i,\mu_0}$, the
preimage under $j_V$ of the image of $\mathrm{th}_G(i,\mu_0)^*$, which is orthogonal to such $u$ for the
class-sum pairing. By the definition of $w_v$,
$|w_v(\theta_{i,\mu_0}y)|^2=n_P^{-1}\|FW_i(\sigma_v)(z)\|^2$. It therefore suffices to show that
\[
 e^{\mathrm{low}(i)}\langle x,x\rangle_\sigma\le\|FW_i(\sigma)(z)\|^2\tag{$*$}
\]
for every embedding $\sigma$ of $L$ and every lift form $x$ at $(i,\mu_0)$ with rational lift $z$ such
that $\langle x,u\rangle_{\sigma'}=0$ for every embedding $\sigma'$ and every lift form $u$ with
$\mathrm{th}_G(i,\mu_0)(j_V(u))=0$.

\emph{Embeddings that agree with $\sigma_0$ on $K$.} Let $\sigma$ be one. By
Lemma~\ref{is-eigenvalue-base-inputs}(3b) there are a conjugate datum at $\sigma$ for $i$ and the tie,
with characters
$\chi',\eta'$, coefficients $W_c$, form space $H$ and maps $j$ and $\mathrm{th}$, and $U$, $z_0$ and
$\kappa$ with $U(FW_i(\sigma)(z))=\kappa\,z_0\,\mathrm{th}(j(x))$ for every lift form $x$ with rational
lift $z$. Choose an automorphism $s$ of $\mathbf C$ with $\sigma=s\circ\sigma_0$. By
Lemma~\ref{is-eigenvalue-type-parts}(3) and~(4), $\chi'$ and $\eta'$ come from the images under $s$ of
$\mu^0_i\xi_{\mu_0}$ and $\eta$. Hence the Weil datum, level vector, basis and map $\mathrm{io}$ of the
conjugate datum, with the residue characters the images under $s$ of those of the realization at the
places of $P_g$, are lift data for $R$ at $F$ with the lift $i$, the base tuple, the embedding
$\sigma$, the automorphism $s$ and the trivial level; the measure of the realization at $(i,\mu_0)$ is
a measure for them, and $H$ is a form space for them, its character being the one of these lift data.
By Lemma~\ref{is-eigenvalue-adelic}(1) they have a theta lift $T\colon H\to H_W(i,\sigma)$. The maps
$\mathrm{th}$ and $eT$ give the same function at every tensor of lowest $K$-type vectors, and an element
of $H_W(i,\sigma)$ whose function at an extremal vector vanishes is zero, so $\mathrm{th}=eT$.

By Lemma~\ref{is-eigenvalue-adelic}(2), for every normalized $c$ with $H_W(i,\sigma)_c\ne0$, $TT^*$ is
multiplication on $H_W(i,\sigma)_c$ by some $l_c>0$ with
$\log l_c\ge\sum_\rho\log\big(\mathrm{hook}(a^+_\rho(\sigma))\,\mathrm{hook}(a^-_\rho(\sigma))\big)
-8fn_V^2L_p$. By Lemma~\ref{is-eigenvalue-base-inputs}(3b) the sum equals the one at $\sigma_0$, so
$\log l_c\ge\mathrm{low}(i)$ as in the proof of Lemma~\ref{is-eigenvalue-base-place}. The spaces
$H_W(i,\sigma)_c$ span $H_W(i,\sigma)$ (Lemma~\ref{is-eigenvalue-type-parts}(1)), and $e$ is an
orthogonal projection, so Lemma~\ref{is-eigenvalue-projection}(3) gives
$e^{\mathrm{low}(i)}\|h\|^2\le\|\mathrm{th}(h)\|^2$ for every $h\in H$ orthogonal to the kernel of
$\mathrm{th}$.

Let $x$ and $z$ be as in $(*)$. By Lemma~\ref{is-eigenvalue-base-inputs}(3a) the rational lift defines an
$L$-linear map $\theta_L$ from the lift forms at $(i,\mu_0)$ to $L^{n_W(i)}$ with
$FW_i(\sigma_0)\theta_L=\mathrm{th}_G(i,\mu_0)j_V$, and $U\,FW_i(\sigma)\,\theta_L=\kappa\,z_0\,
\mathrm{th}\,j$. If $\theta_L(u)=0$, then $\mathrm{th}_G(i,\mu_0)(j_V(u))=0$.
Since the image of $j$ spans $H$ (Lemma~\ref{is-eigenvalue-base-inputs}(4)) and $FW_i(\sigma)$ carries
$L$-linearly independent vectors to $\mathbf C$-linearly independent ones
(Definition~\ref{is-primitive-rows}(11)),
the kernel of $\mathrm{th}$ lies in the span of $j(\ker\theta_L)$. By hypothesis and
Lemma~\ref{is-eigenvalue-base-inputs}(4), $\langle j(x),j(u)\rangle=\langle x,u\rangle_\sigma=0$ for
$u\in\ker\theta_L$, so $j(x)$ is orthogonal to the kernel of $\mathrm{th}$. Therefore
$e^{\mathrm{low}(i)}\|j(x)\|^2\le\|\mathrm{th}(j(x))\|^2$. Finally $\|j(x)\|^2=\langle x,x\rangle_\sigma$
by Lemma~\ref{is-eigenvalue-base-inputs}(4), and $\|\mathrm{th}(j(x))\|=\|FW_i(\sigma)(z)\|$ because $U$
and $z_0$ are unitary and $|\kappa|=1$. This is $(*)$.

\emph{The other embeddings.} If $\sigma$ does not agree with $\sigma_0$ on $K$, then $\bar\sigma$ does
(Lemma~\ref{is-eigenvalue-base-inputs}(3c)), and $(*)$ holds at $\bar\sigma$. Both sides of $(*)$ are
real, and those at $\sigma$ are the complex conjugates of those at $\bar\sigma$: for
$\langle x,x\rangle_\sigma$ because the class-sum pairing is taken to $\mathbf C$ by $\sigma$, and for
$\|FW_i(\sigma)(z)\|^2$ by Lemma~\ref{is-eigenvalue-base-inputs}(3c). So $(*)$ holds at $\sigma$.
\end{proof}

\begin{lemma}\label{is-lift-eigenvalue-bound}
Assume rapid decay of cusp forms, the Weil-Siegel formula and the Siegel-Weil formula in the
convergent range (an assumption the proof does not use), the doubling method with its unramified
computation for $GL$,
Bernstein-Zelevinsky finiteness and the formal degree of the compact dual
(Theorem~\ref{lit-bundle-unitary}). Let $D$ be a definite-side datum that has parameter choices and
standard spaces $V_i$, let $R$ be a realization of $D$ with a face $F$ and one-map data
(Definitions~\ref{is-realisation} and~\ref{is-one-map-data}), and let $C,C_H$ and $b$ be real.
Suppose that the family of definite forms satisfies the family conditions, with the conditions on its
inert indices (Lemma~\ref{is-eigenvalue-base-inputs}(1)), and has lift-distinct types
(Definition~\ref{is-realisation}),
the realization has primitive rows at $F$ (Definition~\ref{is-primitive-rows}), satisfies the weak
characterization at $F$ and the existence conditions for its theta maps, has all its
phases equal to $1$ and the standard character $\psi$ with sign $-1$, satisfies (1) of
Definition~\ref{is-bounded-rows}, the one-map data are pinned at every lift and give each lift as one
map for some constants and phase, and, if the rational structures $FW_i$ satisfy (11) of
Definition~\ref{is-primitive-rows}, they have the \emph{phase tie} at every lift $i$: if the one-map data
at $i$ give the $i$-th lift as one map, they give it with the canonical constants of the core and the
phase $\phi_i$ of the realization, replaced by $\bar\phi_i$ if $\sigma_0$ is not the embedding of the base
place $v_0$. Suppose finally that the lifts have the sizes of the size data with $b$, and that
$\ell_6(j)\le\ell_6^{m}(j)\le\mathrm{low}(j)$ for the numerical bounds $N(R,F;C,C_H)$ and every lift
$j$ (Lemma~\ref{is-hook-lower-bound}). Then the columns $S_1(R,F)$ satisfy the lower norm bound for
$N(R,F;C,C_H)$ (Definition~\ref{is-lower-norm-bound}).
\end{lemma}
\leanref{Development/AutomorphicArithmetic/ThetaCorrespondence/ThetaPairLifts/LocalTypeDictionary/FinitenessInputsRealDataThree.lean}{GlobalTheory.T.finiteness_inputs_real_data_three_alt}
\leanref{Development/AutomorphicArithmetic/ThetaCorrespondence/ThetaPairLifts/LocalTypeDictionary/FinitenessInputsRealDataThree.lean}{GlobalTheory.LowRow.lamLow_closing_alt}
\begin{proof}
Write $\theta_j=\mathrm{th}_G(j,\mu_0)$ and $\mathrm{low}(j)$ as in Lemma~\ref{is-hook-lower-bound}.

\emph{Step 1: inputs.} By Lemma~\ref{is-eigenvalue-base-inputs}(1), the family conditions, the
lift-distinct types, the weak characterization with the existence conditions for the theta maps, the
primitive rows and the phases give the core conditions at $F$ and the base links of the rows
$S_0(R,F)$; by Lemma~\ref{is-eigenvalue-base-inputs}(2) the base tuple is prime to the host. By
Lemma~\ref{is-eigenvalue-adelic}, the published results, the parameter choices and the standard spaces
give the conclusions of that lemma for all lift data for $R$ at $F$. Since all phases are $1$, the
phase tie of the one-map data, which applies by (11) of Definition~\ref{is-primitive-rows}, makes the
one-map data at each lift give that lift as one map with the canonical constants of the core and
phase $1$. So Lemma~\ref{is-eigenvalue-base-inputs}(3) applies with phase $1$ at every lift; its other
hypotheses are the standard character, the weight and conjugation ties, the pins, (1) of
Definition~\ref{is-bounded-rows} and (11) of Definition~\ref{is-primitive-rows}. The weight tie also
makes each vector $v_0(i,\sigma)$ extremal.

\emph{Step 2: the two eigenvalue bounds.} From these inputs and the sizes of the lifts:
\begin{enumerate}
\item[(a)] at the base place, every positive eigenvalue $l$ of $\theta_j\theta_j^*$ on
  $H_W(j,\sigma_0)$ satisfies $\log l\geq\mathrm{low}(j)$, by Lemma~\ref{is-eigenvalue-base-place};
\item[(b)] at every infinite place $v\ne v_0$, the matrix $\theta_{j,\mu_0}$ of the rows satisfies
  $e^{\mathrm{low}(j)}|y|^2_{X_T(j,\mu_0),v}\le|w_v(\theta_{j,\mu_0}y)|^2$ for every vector $y$ of
  coordinates, where $X_T(j,\mu_0)$ is the metric of the rows and $w_v$ the map to $W_{0,j,v}$ given
  by the vectors $w_{0,j,v}(k)$ of Definition~\ref{is-complex-data}(4), by
  Lemma~\ref{is-eigenvalue-other-places}.
\end{enumerate}
% Step 1 is lamLow_closing_alt: adelC_of_parameter_choice_stdV (Lemma is-eigenvalue-adelic),
% splitPart_baseLinks_of_primitiveRows and badPlacesPrimeToHost_baseTuple (Lemma is-eigenvalue-base-inputs
% (1), (2)); in lamLow_of_weak, conjugatesOfRationalThetaLift_from_characterisation (ConjRatLift5, (3) of
% that lemma) gives GalOut R i 1 through the phase tie htie i prim.unif (hrow0 i), and fw.hv0 gives the
% extremal vectors. Step 2(a) is liftAdjoint_eigenvalue_lower_bound_base_embedding (LowTop, Lemma
% is-eigenvalue-base-place) and 2(b) is rationalLift_lowerBound_off_basePlace (LowForm, Lemma
% is-eigenvalue-other-places). Steps 3 to 5 are GlobalTheory.LowRow.lamLow_of_weak with
% l6m_le_log_lam3_base_place and l6m_le_log_lam3_other_places.

\emph{Step 3: the base place.} The weak characterization contains the conventions for the scalars:
if $W_p(j)=0$ then $\lambda_{j,\mu_0,p,v_0}=e^{\ell_6^m(j)}$, and otherwise
$\lambda_{j,\mu_0,p,v_0}$ is the eigenvalue of $\theta_j\theta_j^*$ on $W_p(j)$, which exists by the
scalars of the primitive rows (Definition~\ref{is-primitive-rows}(9)) and is positive by their
positivity. In the first case $\log\lambda_{j,\mu_0,p,v_0}=\ell_6^m(j)$. In the second, choose
$x\ne0$ in $W_p(j)$; then (a) gives $\log\lambda_{j,\mu_0,p,v_0}\geq\mathrm{low}(j)\geq\ell_6^m(j)$.

\emph{Step 4: the other places.} Let $v\ne v_0$. If no $y$ with $t_0(j,\mu_0)\sigma_0(y)\in W_p(j)$
has $|y|_{X_T(j,\mu_0),v}\ne0$, the weak characterization sets $\lambda_{j,\mu_0,p,v}=
e^{\ell_6^m(j)}$. Otherwise take such a $y$. As all phases are $1$, the conjugation condition of the
primitive rows (Definition~\ref{is-primitive-rows}(12)) gives $|w_v(\theta_{j,\mu_0}y)|^2=
\lambda_{j,\mu_0,p,v}|y|^2_{X_T(j,\mu_0),v}$, and with (b),
$e^{\mathrm{low}(j)}\le\lambda_{j,\mu_0,p,v}$. In both cases $\log\lambda_{j,\mu_0,p,v}\geq
\ell_6^m(j)$, using $\ell_6^m(j)\le\mathrm{low}(j)$.

\emph{Step 5.} Since $\ell_6(j)\le\ell_6^m(j)$, Steps 3 and 4 give
$\log\lambda_{j,\mu_0,p,v}\geq\ell_6(j)$ for every lift $j$, part $p$ and infinite place $v$, which
is the lower norm bound for the columns.
\end{proof}


\begin{theorem}\label{is-surjectivity-and-lower-norm}
Assume the published results on unitary groups of Theorem~\ref{lit-bundle-unitary}. Let
$\mathcal H$ be the class of standard families. Then the following hold.
\begin{enumerate}
\item Eventually along $\mathcal H$, every family $\mathcal F$ that satisfies the division bound at
  twisted places (Definition~\ref{is-twisted-division-bound}; a hypothesis that the proof does not use)
  has a constant $C_0\ge0$ and a scale
  $x_0$ with the following property. For every scale $x\ge x_0$, every member $i$ at $x$ and every
  point $p$ at $(x,i)$ that satisfies the point conditions and the further conditions
  (Definition~\ref{is-point}), the rows $S_0(p)$ satisfy the surjectivity bound
  (Definition~\ref{is-surjectivity-bound}) for $N_p(C,C_H)$ for every $C\ge C_0$.
\item Eventually along $\mathcal H$, every family $\mathcal F$ has a scale $x_0$ with the following
  property. For every scale $x\ge x_0$, every member $i$ at $x$ and every point $p$ at $(x,i)$ that
  satisfies the point conditions, the columns $S_1(p)$ satisfy the lower norm bound
  (Definition~\ref{is-lower-norm-bound}) for $N_p(C,C_H)$ for every real $C$.
\end{enumerate}
\end{theorem}
\leanref{Development/AutomorphicArithmetic/ThetaCorrespondence/ThetaPairLifts/ThirdPlaceSeries/SurjectiveAndLowWeightOfUnramifiedExponents.lean}{GlobalTheory.Fc.S0.ontoX_and_lamLow_of_unitaryGroupHypotheses'}
\leanref{Development/AutomorphicArithmetic/ThetaCorrespondence/ThetaPairLifts/ThirdPlaceSeries/ExportSt.lean}{GlobalTheory.BigOn.SurjectiveOntoX}
\leanref{Development/AutomorphicArithmetic/ThetaCorrespondence/ThetaPairLifts/ThirdPlaceSeries/Realised.lean}{GlobalTheory.BigTop.LambdaLowerBound₂}
\begin{proof}
(2) Eventually along $\mathcal H$ the size constant $C_{sz}$ of the family is at least $8$, and
eventually along each family the lifts have the sizes required by the size data. Fix a family, a
scale $x$ beyond that threshold, a member $i$, a point $p$ at $(x,i)$ that satisfies the point
conditions, and a real number $C$. The point conditions give the hypotheses of
Lemma~\ref{is-lift-eigenvalue-bound} on the realization of $p$ and its face: lift-distinct types,
primitive rows, the weak characterization, phases $1$, the standard character, condition (1) of
Definition~\ref{is-bounded-rows}, and pinned one-map data giving each lift as one map. The family
supplies parameter choices for $D$. The scale data of $p$ have $\ell_6^{u}=0$ and
$\ell_6^{m}(j)=\sum_\rho\log\mathrm{hook}_j(\rho)-C_{sz}fn_V^2L_p$, so
$\ell_6(j)=\ell_6^{m}(j)-C\cdot0=\ell_6^{m}(j)$, and $\ell_6^{m}(j)\le\mathrm{low}(j)$ by
Lemma~\ref{is-hook-lower-bound}. Lemma~\ref{is-lift-eigenvalue-bound} gives the lower norm bound of
$S_1(p)$ for $N_p(C,C_H)$.

(1) Lemma~\ref{is-onto-estimates} gives, eventually along $\mathcal H$, for each family $\mathcal F$,
scales beyond which (1) and (2) of that lemma hold, and constants $c_G,c_K,c_T,c_S,c_B,c_E,c_P$ for the
seven estimates (3b), (3c), (3d), (3e), (3f), (3g) and (3a). Let $c_0$ be the largest of these
constants and $1$, and take $c=(c_0,14,66c_0+4)$; then $C=c_0$ is at least each constant, and
$14\le C_{b,1}$, $66C+4\le C_{b,2}$. Let $C_E$ be the number of Lemma~\ref{is-twisted-places-chain}
for $c$, put $C_0=2\max(C_E,0)$, and let $x_0$ be the largest of the scales of the estimates.

The division bound at twisted places is not used. Let $x\ge x_0$, let $i$ be a member at $x$, let $p$
be a point at $(x,i)$ that satisfies the point conditions and the further conditions, and let
$C'\ge C_0$. By (1) of Lemma~\ref{is-onto-estimates}
there is a pair of runs $r$ pinned to the own places, and the realization of $p$ is rational at the
twisted places. The further conditions tie the core of $p$ to its realization, so there is a record
object $O$ with constants $c$ recorded at $(p,r)$, and $O$ satisfies (3a) to (3g). The further
conditions also give the links of the realization, among them $d_Q(\mu_0)>0$ for the rows, so $D>0$.
Lemma~\ref{is-twisted-places-chain} gives error terms $\mathrm{err}$ for which $S_0(p)$ satisfies its
three conditions with the bound $C_EE_1^{u}$.

Apply Lemma~\ref{is-canonical-error} to the twisted places of $S_0(p)$ and its lifts, with
$I(w,j)=I_w(j,\mu_w)+I_{\bar w}(j,\mu_w)$, $F=F_0$, $J_0(w,j)=J_0(\mu_w,w,j)$, which is nonnegative by
(3g), $e_w$, $c_w=2e_w/p_w^2$, $\ell_w$, $E=C_EE_1^{u}$ and $E'=C'E_1^{u}$. Since $E_1^{u}\geq0$ by
(3a) and $2C_E\le C_0\le C'$, $2E\le E'$. The point is compatible with the canonical error terms: the
error terms $\mathrm{err}_0$ of $S_0(p)$ are the numbers $\mathrm{err}_0$ of the lemma for $F_0$. So
$S_0(p)$ satisfies (1), (2) and (3) of Definition~\ref{is-surjectivity-bound} with
$E_1(j)=C'E_1^{u}$, which is the surjectivity bound for $N_p(C',C_H)$.
\end{proof}

\section{The definite rows}\label{section-is-definite-rows}

\begin{definition}\label{is-cauchy-binet}
Let $K$ be a number field, let $c_0,\dots,c_{m'}$ be integers, let $X_i$ be metrics on $K^{c_i}$
invertible at every infinite place, and let $K_i\in M_{c_i\times c_{i-1}}(K)$ for $1\leq i\leq m'$.
\emph{Cauchy-Binet data} for the chain $(X_i,K_i)$ consist of the following. For each level $i$:
a finite set $A_i$, finite sets $\mathrm{Cl}_i(a)$ of \emph{classes} for $a\in A_i$, each with at most
$h_i$ elements, weights $0<\gamma_i(p)\leq1$ for $p$ in the disjoint union $\mathrm{Cl}_i$ of the
$\mathrm{Cl}_i(a)$, and at each infinite place $v$ a normed complex space $W_{i,v}$ and a linear map
$\Phi_{i,v}\colon\mathbf C^{c_i}\to W_{i,v}^{\mathrm{Cl}_i}$ with
$\sum_p\gamma_i(p)\|\Phi_{i,v}(x)_p\|^2=|X_{i,v}^{-1}x|^2$. For each step $i\geq1$ and place $v$: a
real number $t_i(v)\geq0$, and for $p\in\mathrm{Cl}_{i-1}$ and $q\in\mathrm{Cl}_i$ a complex number
$\epsilon_{i,v}(p,q)$ of absolute value $1$ and a linear map $T_{i,v}(p,q)\colon W_{i-1,v}\to W_{i,v}$
with $\|T_{i,v}(p,q)\|^2\leq t_i(v)$, such that
\[
 \Phi_{i,v}\big(\sigma_v(K_i)x\big)_q=\sum_{p\in\mathrm{Cl}_{i-1}}\gamma_{i-1}(p)\,\epsilon_{i,v}(p,q)\,
 T_{i,v}(p,q)\big(\Phi_{i-1,v}(x)_p\big).
\]
The \emph{square factor} at the step $i$ and $v$ is $|A_{i-1}|\,|A_i|\,h_{i-1}h_i\,t_i(v)$.
\end{definition}
\leanref{Formalism/LocalToGlobal/Divisibility/Inputs.lean}{Formalism.L2.DIV.CBData}
\leanref{Formalism/ThetaSupremumBounds/StatementCB.lean}{Formalism.DimensionVariation.CBHyp}
\leanref{Formalism/ThetaSupremumBounds/StatementCB.lean}{Formalism.DimensionVariation.cbSq}

\begin{lemma}\label{is-cauchy-binet-bound}
Let Cauchy-Binet data be given for a chain $(X_i,K_i)$ (Definition~\ref{is-cauchy-binet}). Then at
every step $i$ and infinite place $v$, the operator norm of $K_i$ from $X_{i-1}$ to $X_i$ satisfies
$\|K_i\|_v^2\leq|A_{i-1}|\,|A_i|\,h_{i-1}h_i\,t_i(v)$.
\end{lemma}
\leanref{Formalism/LocalToGlobal/Divisibility/DivisionHeightBoundStatement.lean}{Formalism.L2.DIV.nK_of_cb}
\leanref{Formalism/ThetaSupremumBounds/GaussianShiftBounds/PerEmbeddingCoupledBound.lean}{Formalism.DimensionVariation.chain_coupled_bound}
\begin{proof}
Fix $x\in\mathbf C^{c_{i-1}}$ and $q\in\mathrm{Cl}_i$. By the triangle inequality and the
Cauchy-Schwarz inequality,
\[
 \|\Phi_{i,v}(\sigma_v(K_i)x)_q\|\leq\sum_p\gamma_{i-1}(p)\,t_i(v)^{1/2}\|\Phi_{i-1,v}(x)_p\|
 \leq t_i(v)^{1/2}\Big(\sum_p\gamma_{i-1}(p)\Big)^{1/2}|X_{i-1,v}^{-1}x|.
\]
Since $\gamma_{i-1}\leq1$, $\sum_p\gamma_{i-1}(p)\leq|\mathrm{Cl}_{i-1}|\leq|A_{i-1}|h_{i-1}$. So
$|X_{i,v}^{-1}\sigma_v(K_i)x|^2=\sum_q\gamma_i(q)\|\Phi_{i,v}(\sigma_v(K_i)x)_q\|^2\leq
|A_i|h_i\,t_i(v)\,|A_{i-1}|h_{i-1}\,|X_{i-1,v}^{-1}x|^2$.
\end{proof}

\begin{definition}\label{is-definite-side-condition}
Let $D$ be a definite-side datum (Definition~\ref{th-definite-datum}) with coefficient field $L$,
words $u$ and division lattices $N_u$, let $b$ be a bundle at $D$, that is, a realization of $D$ by
indefinite-side data with square-integrable columns and a determined row side, and let $C,C_H$ be
real. The \emph{definite-side condition} holds for $b$ and $(C,C_H)$ if, whenever the realization
satisfies the primitivity conditions and the covolume, surjectivity and isometry conditions with
constants $(C,C_H)$, and the record numbers of every word hold with constants
$\lambda_0,C_0,c_e,c,S_w,C_1$, then for all division data tied to the realization and every word $u$
the following \emph{rows condition} holds.
\begin{enumerate}
\item There is a chain of lifts for $u$: integers $c_0,\dots,c_{m'}$, where $c_0$ is the dimension
  of the column of $u$, metrics $X_i$ on $L^{c_i}$ with $X_0$ the metric of the column, matrices
  $K_i\in M_{c_i\times c_{i-1}}(L)$ for $1\le i\le m'$, a map from lifts to steps, and a linear map
  $P$ from $L^{c_{m'}}$ to the dual of the top space of the base tuple, such that the period
  pairing of $u$ at the base tuple is $(\varphi,F)\mapsto P(K_{m'}\cdots K_1\varphi)(F)$.
\item The chain is tied to the theta kernels of the realization, and its first columns are those of
  the realization.
\item There is a metric $X^{0}$ for the Hermitian matrix $H_V(\mu_0)$ of the columns, that is,
  $(X^{0}_v)^{-1*}(X^{0}_v)^{-1}=\sigma_v(H_V(\mu_0))$ at every infinite place $v$.
\item There are height numbers for the record of $u$ (Definition~\ref{is-height-numbers}), with a
  period exponent $h_P$, such that
  $\|P\|_v^{2}\le e^{h_P}$ at every infinite place $v$, the norm taken from $X_{m'}$ to the dual
  metric of $X^{0}$.
\item There are Cauchy-Binet data for the chain (Definition~\ref{is-cauchy-binet}) such that, at every step $i$ and infinite place
  $v$, the logarithm of their square factor at $(i,v)$ is at most the sum of the logarithmic sizes
  $\mathrm{sz}(j,v)$ of the lifts $j$ assigned to the step $i$. Here
  \[
   \mathrm{sz}(j,v)=\sum_\tau\Big(\sum_i\phi(\nu_{j,\tau,i})-\Big(\sum_i\nu_{j,\tau,i}\Big)\log\Omega_\tau(v)\Big)
   +(\text{lower order terms}),\qquad \phi(y)=y\log y-y,
  \]
  where $\nu_{j,\tau,i}$ are the weights of the $j$-th lift at the type $\tau$ and $\Omega_\tau(v)$
  are the periods of the record.
\item The metric of the column of $u$ is invertible at every infinite place, $c_0>0$, and $N_u$,
  viewed in $L^{c_0}$, is a lattice.
\end{enumerate}
\end{definition}
\leanref{Development/AutomorphicArithmetic/ThetaCorrespondence/ThetaPairLifts/ThirdPlaceSeries/DefiniteSideCondition.lean}{GlobalTheory.Fc.DefiniteSideCondition}
\leanref{Formalism/LocalToGlobal/PtKOfReal.lean}{Formalism.L2.RowDataTransport.RowsU}
\leanref{Formalism/LocalToGlobal/Divisibility/Inputs.lean}{Formalism.L2.DIV.ColChain}
\leanref{Formalism/LocalToGlobal/Divisibility/Inputs.lean}{Formalism.L2.DIV.HtNums}
\leanref{Formalism/ThetaSupremumBounds/DivIneqExplicit6.lean}{Formalism.DimensionVariation.liftedMapLogSize}

The main term $\phi(a_t)-a_t\log\Omega$ of $\mathrm{sz}(j,v)$ at a top class $t$ is the cost of one
lift; summed over the $m$ lifts it gives the term $m(\phi(a_t)-a_t\log\Omega)$ of
$\mathrm{MAIN}(u)$ in the division link of Section~\ref{section-is-division}.
Section~\ref{section-is-kernel-data} proves the definite-side condition for the bundles
attached to linked parameters (Theorem~\ref{is-definite-side-condition-holds}).

\begin{definition}\label{is-slot-rows}
A \emph{slot} $Q$ is a condition on the cores of the stages of the construction, and \emph{link
data} $\Lambda$ are a condition on the bundles at base points. A \emph{selection} for a class
$\mathcal A$ of families chooses a bundle at every member of every family of $\mathcal A$ at every
scale beyond its threshold. A property $H$ of classes and selections holds \emph{along}
$(Q,\Lambda)$ if there is a choice $E$, for each parameter, of a condition on bundles, satisfied
eventually by every bundle that $\Lambda$ links to a base point of $Q$, such that $H(\mathcal A,s)$
holds for every class $\mathcal A$ and every selection $s$ for $\mathcal A$ whose bundles are built
from $(Q,\Lambda)$ and satisfy $E$.
\end{definition}
\leanref{Development/AbelianVarietyTheory/Definitions/GlobalTwisting/ValidForConstructedPointClassesEventuallyWorking.lean}{GlobalTheory.Fc.Row}
\leanref{Formalism/LocalToGlobal/ClassOfRecordOfChoice.lean}{Formalism.L2.Top3.D3}

\begin{theorem}\label{is-definite-rows}
Assume the covering of discrete series, the finiteness of Harish-Chandra modules, rapid decay of
cusp forms, the Casselman-Wallach theorem and the classification of discrete series
(Theorem~\ref{lit-bundle-unitary}), and assume parts (1) and (2) of
Theorem~\ref{is-surjectivity-and-lower-norm}. Let $Q$ be a slot and $\Lambda$ link data
(Definition~\ref{is-slot-rows}). Suppose that the following hold along $(Q,\Lambda)$, for every
class $\mathcal A$ and selection $s$:
\begin{enumerate}
\item at every member of every family of $\mathcal A$ at every scale, the realization of the
  selected bundle has primitive rows, and its family of definite forms has lift-distinct types
  (Definition~\ref{is-realisation});
\item at every such member, the division rows of the selected bundle hold;
\item eventually along $\mathcal A^{\mathrm{des}}$, eventually along each family, every row of the
  definite-side datum is good;
\item at every member, the selected bundle has a comparison part whose point has the normalized
  embedding of its setting and satisfies the class-sum pin, the compatibility conditions and the
  further conditions of Definition~\ref{is-point};
\item eventually along $\mathcal A^{\mathrm{des}}$, every family has a constant $C_0\ge0$ such that
  for every $C\ge C_0$, eventually along the family, the rows of the realization satisfy the
  covolume bound (Definition~\ref{is-numerical-bounds}) for $N(C,C_H)$;
\item the same holds for the functional bound (Definition~\ref{is-functional-bound}), for the
  families of $\mathcal A^{\mathrm{des}}$ that satisfy the division bound at twisted places;
\item eventually along $\mathcal A^{\mathrm{des}}$, eventually along each family, the multiplicity
  condition of the realization holds;
\item eventually along $\mathcal A^{\mathrm{des}}$, eventually along each family, the archimedean
  scalars $\kappa$ of the columns are those of the division data;
\item eventually along $\mathcal A^{\mathrm{des}}$, every family has a constant $C_0\ge0$ such that
  for every $C\ge C_0$, eventually along the family, the definite-side condition
  (Definition~\ref{is-definite-side-condition}) holds for the selected bundle and $(C,C_H)$.
\end{enumerate}
Then the following holds along $(Q,\Lambda)$. For every class $\mathcal A$ and selection $s$,
eventually along the designed families $\mathcal A^{\mathrm{des}}$, every family $\mathcal F$ has a
constant $C\ge0$ and a scale $x_0$ such that at every scale $x\ge x_0$ and every member, the
realization $R$ of the selected bundle satisfies:
\begin{enumerate}
\item every place $w$ at which an index $I_w(j,\mu_0)$ or $I_{\bar w}(j,\mu_0)$ of the rows, a
  relative index of $T_{j,\mu_0}$ and its dual, or a value $F_0(w,j)$ or $F_0(\bar w,j)$ of
  $N(C,C_H)$ is nonzero, is twisted;
\item the rows of $R$ satisfy the covolume bound for $N(C,C_H)$;
\item if $\mathcal F$ satisfies the division bound at twisted places, the rows satisfy the
  surjectivity bound and the functional bound for $N(C,C_H)$;
\item if every lift of the selected bundle has a single row map, the columns of $R$ satisfy the
  lower norm bound for $N(C,C_H)$;
\item the isometry of the row space, the multiplicity condition, the passage from rows to columns
  and the equality of the archimedean scalars hold;
\item the definite-side condition holds for the selected bundle and $(C,C_H)$.
\end{enumerate}
\end{theorem}
\leanref{Development/AutomorphicArithmetic/ThetaCorrespondence/ThetaPairLifts/ThirdPlaceSeries/TupAllUsedAtPart.lean}{GlobalTheory.Fc.StageZeroRow.row_rRows_stage0_alt}
\leanref{Development/AutomorphicArithmetic/ThetaCorrespondence/ThetaPairLifts/ThirdPlaceSeries/CountsIrreducibilityConstants.lean}{GlobalTheory.Fc.R1.SelectionDefiniteRows}
\begin{proof}
Conclusions (2), (6) and the multiplicity and scalar parts of (5) are hypotheses (5), (9), (7) and
(8). Hypothesis (6) gives the functional bound in (3).

Choose, at every member, the comparison part of hypothesis (4); this is a compatible selection.
Every tuple of characters used by the selected realization is either the base tuple or the tuple
$\mu_w$ of a place, because the realization attaches $\mu_w$ to $w$ and is tied to the tuples of
$D$. Hence the division rows of hypothesis (2) give the divisor conditions of the selected parts.
Every family of a selection built from $(Q,\Lambda)$ carries its seven duties. Rapid decay of cusp
forms, the finiteness of Harish-Chandra modules, the covering of discrete series and the
Casselman-Wallach theorem, with hypotheses (1), (2), (3) and the class-sum pin, show that
eventually the rows of the selected parts are tied to their divisors.

At such parts, the further conditions hold by hypothesis (4), so part (1) of
Theorem~\ref{is-surjectivity-and-lower-norm} gives the surjectivity bound in (3), and part (2) gives
the lower norm bound in (4), without using the single row maps. The classification of discrete
series shows that the theta lift attached to each selected part does not vanish at the indefinite scale, given the normalized embedding of
hypothesis (4); with the primitivity of hypothesis (1) this gives the smallness of the summands,
which consists of (1) and the isometry and passage parts of (5).

Each of these conclusions holds for $C$ above its own constant and for $x$ above its own scale.
Taking $C$ the largest of the constants and $x_0$ the largest of the scales gives all of them at
once.
\end{proof}

\section{The polynomial model and the branching norm}\label{section-is-polynomial-model}

\begin{definition}\label{is-polynomial-model}
Let $K$ be a field of characteristic $0$ and $N\geq0$ an integer.
\begin{enumerate}
\item A \emph{dominant weight} of length $N$ is a sequence of integers
  $\lambda=(\lambda_1\geq\cdots\geq\lambda_N)$. Its \emph{shift} is $\operatorname{sh}(\lambda)=\lambda_N$ ($0$ if $N=0$),
  and its \emph{natural weight} is $\lambda^{\natural}_i=\lambda_i-\operatorname{sh}(\lambda)$.
\item Let $\mathrm{Pol}_N=K[x_{ab}:1\le a,b\le N]$ be the polynomial ring on the entries of an
  $N\times N$ matrix. For a non-increasing $\nu\in\mathbf N^N$, $\nabla(\nu)\subset\mathrm{Pol}_N$ is
  the space of polynomials $f$ that are homogeneous of degree $\nu_b$ in the entries of the $b$-th
  column for every $b$, and satisfy $f(xu)=f(x)$ for every upper unitriangular $u$; equivalently,
  adding $y$ times the $l$-th column to the $j$-th column, for $l<j$ and an indeterminate $y$, does
  not change $f$. For a dominant weight $\lambda$, $\nabla(\lambda)=\nabla(\lambda^{\natural})$.
\item For $1\le s\le N$ let $\Delta_s$ be the leading principal $s\times s$ minor of $x$. The
  \emph{extremal vector} of $\lambda$ is
  $h_\lambda=\prod_{s=1}^{N}\Delta_s^{\lambda^{\natural}_s-\lambda^{\natural}_{s+1}}$, with
  $\lambda^{\natural}_{N+1}=0$.
\item The \emph{Fock pairing} on a polynomial ring is the bilinear form
  $\langle f,g\rangle_F=\sum_mf_mg_m\,m!$, where $f=\sum_mf_mx^m$, $g=\sum_mg_mx^m$ and
  $m!=\prod_km_k!$ for an exponent vector $m$. If $K$ has an involution $a\mapsto\bar a$, its
  sesquilinear form is $\{f,g\}_F=\sum_m\bar f_mg_m\,m!$, and $\bar f=\sum_m\bar f_mx^m$.
\item Let $\iota$ be a finite set and $\lambda=(\lambda_\tau)_{\tau\in\iota}$ dominant weights of
  length $N$. Let $\mathrm{Pol}_N^{\iota}=K[x_{\tau,a,b}]$, and for $f_\tau\in\mathrm{Pol}_N$ write
  $\otimes_\tau f_\tau$ for the product of the $f_\tau$ in the variables $x_\tau$. Then $V(\lambda)$
  is the span of the $\otimes_\tau f_\tau$ with $f_\tau\in\nabla(\lambda_\tau)$, and
  $h_\lambda=\otimes_\tau h_{\lambda_\tau}$. The \emph{normalized pairing} on $V(\lambda)$ is
  $\langle x,y\rangle_\lambda=\langle x,y\rangle_F/\prod_\tau\langle h_{\lambda_\tau},h_{\lambda_\tau}
  \rangle_F$.
\item The group $GL_N(K)^{\iota}$ acts on $V(\lambda)$ by
  $(\rho_\lambda(g)f)(x)=\prod_\tau\det(g_\tau)^{\operatorname{sh}(\lambda_\tau)}\,f\big((g_\tau^{\mathsf T}x_\tau)_\tau\big)$.
\item If $K$ has an involution, for matrices $G=(G_\tau)$ in $M_N(K)$ the \emph{Hermitian pairing}
  on $V(\lambda)$ is
  \[
   \langle x,y\rangle_{\lambda,G}=\frac{\prod_\tau\det(G_\tau)^{\operatorname{sh}(\lambda_\tau)}}
   {\prod_\tau\{h_{\lambda_\tau},h_{\lambda_\tau}\}_F}\,
   \big\{x,\ y\big((G_\tau^{\mathsf T}x_\tau)_\tau\big)\big\}_F .
  \]
  For an embedding $s\colon K\to\mathbf C$ the \emph{norm} of $f$ is
  $\|f\|^2_{s,G}=\operatorname{Re}s\langle f,f\rangle_{\lambda,G}$.
\end{enumerate}
\end{definition}
\leanref{Formalism/LocalRingCardinality/DataCoeff.lean}{Formalism.BasicNotions.DomWeight}
\leanref{Formalism/GradedWeylDivisibility/StatementA.lean}{Formalism.GradedWeylDivisibility.dualWeyl}
\leanref{Formalism/LocalRingCardinality/DataCoeff.lean}{Formalism.BasicNotions.nablaSq}
\leanref{Formalism/LocalRingCardinality/DataCoeff.lean}{Formalism.BasicNotions.extremal}
\leanref{Formalism/LocalToGlobal/Columns/CoefModB.lean}{Formalism.L2.DIV.fockB}
\leanref{Formalism/LocalToGlobal/Divisibility/CoefMod.lean}{Formalism.L2.DIV.VV}
\leanref{Formalism/LocalToGlobal/Columns/CoefModB.lean}{Formalism.L2.DIV.pairB}
\leanref{Formalism/LocalToGlobal/Divisibility/CoefMod.lean}{Formalism.L2.DIV.rhoK}
\leanref{Formalism/LocalToGlobal/Divisibility/CoefMod.lean}{Formalism.L2.DIV.pairP}
\leanref{Development/AbelianVarietyTheory/BlockTheory/DivisorBase/ScaledCoisometryNorms.lean}{GlobalTheory.IndefiniteRowSideFullHecke.ScaledCoisometryNorms.nrmSq}

The space $\nabla(\lambda)$ is the polynomial model of the irreducible representation of $GL_N$ of
highest weight $\lambda^{\natural}$, the dual Weyl module, and $\rho_\lambda$ twists it by
$\det^{\operatorname{sh}(\lambda)}$. The Fock pairing is the pairing for which differentiation in $x_{ab}$ is adjoint
to multiplication by $x_{ab}$.

\begin{lemma}\label{is-standard-basis}
Let $K$ be a field, $N\geq0$, and $\nu\in\mathbf N^N$ non-increasing, with $\nabla(\nu)\subset
\mathrm{Pol}_N$ as in Definition~\ref{is-polynomial-model}.
\begin{enumerate}
\item $\nabla(\nu)$ has a $K$-basis indexed by the semistandard tableaux of shape $\nu$ with entries
  at most $N$, consisting of bideterminants: products of minors of $x$, the $k$-th on rows
  $I_k(1)<\cdots<I_k(s_k)$ and on the first $s_k$ columns, where $\#\{k:s_k\geq t\}=\nu_t$ for every
  $t$. These polynomials have integer coefficients, and each is a weight vector for the diagonal
  matrices acting on the rows, of the weight given by the content of its tableau.
\item If $K$ has characteristic $0$, then $\nabla(\nu)$ is the Weyl module $\Delta_K(\nu)$ of
  Definition~\ref{pd-weight-lattices}: the $K$-span of the polynomials obtained from $h_\nu$ by
  repeatedly applying the operators $E^{(k)}_{ab}$, where $E^{(k)}_{ab}f$ is the coefficient of $y^k$
  in $f$ with the $b$-th row of $x$ replaced by its sum with $y$ times the $a$-th row.
\end{enumerate}
\end{lemma}
\leanref{Formalism/GradedWeylDivisibility/FinalD.lean}{Formalism.GradedWeylDivisibility.sSBasis}
\leanref{Formalism/LocalToGlobal/Columns/PolyFnLemmas.lean}{Formalism.L2.Col.dualWeyl_eq_weyl}

Part (1) is the standard basis theorem for bitableaux \cite{DEP80}. Part (2) holds because over the
rationals the Weyl module and the dual Weyl module have the same dimension, the number of
semistandard tableaux of shape $\nu$, and the first lies in the second \cite{CL74}.

\begin{lemma}\label{is-fock-facts}
Let $K$, $\iota$ and $\lambda$ be as in Definition~\ref{is-polynomial-model}.
\begin{enumerate}
\item $\langle\otimes_\tau f_\tau,\otimes_\tau g_\tau\rangle_F=\prod_\tau\langle f_\tau,g_\tau\rangle_F$,
  and $\langle h_{\lambda_\tau},h_{\lambda_\tau}\rangle_F\ne0$.
\item The normalized pairing is symmetric and nondegenerate on $V(\lambda)$.
\item $\langle\rho_\lambda(g)x,y\rangle_\lambda=\langle x,\rho_\lambda(g^{\mathsf T})y\rangle_\lambda$
  for $g\in GL_N(K)^{\iota}$, where $g^{\mathsf T}=(g_\tau^{\mathsf T})_\tau$.
\item $V(\lambda)$ is spanned by the vectors $\rho_\lambda(g)h_\lambda$, $g\in GL_N(K)^{\iota}$.
\item Suppose that $K$ has an involution and $s\colon K\to\mathbf C$ is an embedding with
  $s(\bar a)=\overline{s(a)}$. If the $G_\tau$ are invertible, then
  $\langle x,y\rangle_{\lambda,G}=\langle\bar x,\rho_\lambda(G)y\rangle_\lambda$. If the matrices
  $s(G_\tau)$ are positive definite, then $\overline{\langle x,y\rangle_{\lambda,G}}=
  \langle y,x\rangle_{\lambda,G}$ and $\|x\|^2_{s,G}\geq0$.
\item If $K=\mathbf R$ and $f\ne0$, then $\langle f,f\rangle_F>0$.
\end{enumerate}
\end{lemma}
\leanref{Development/AbelianVarietyTheory/Proofs/DivisorProofs/BranchingAdmissibleFacts.lean}{GlobalTheory.BrCompare.fockB_ten}
\leanref{Development/AbelianVarietyTheory/Proofs/DivisorProofs/LastMapNorm.lean}{GlobalTheory.BranchNorm.pairBK_adj}
\leanref{Development/AbelianVarietyTheory/Proofs/DivisorProofs/LastMapNorm.lean}{GlobalTheory.BranchNorm.pairK_eq}
\leanref{Development/AbelianVarietyTheory/Proofs/DivisorProofs/LastMapNorm.lean}{GlobalTheory.BranchNorm.pairK_pos}
\leanref{Development/AbelianVarietyTheory/Proofs/DivisorProofs/LastMapNorm.lean}{GlobalTheory.BranchNorm.fockB_pos}
\begin{proof}
(1) A monomial in the variables of different $\tau$ is a product of monomials, and $m!$ is
multiplicative. The extremal vector is a product of minors, a nonzero polynomial with integer
coefficients, so $\langle h,h\rangle_F=\sum_mh_m^2m!$ is a positive integer, which is nonzero in $K$
because $K$ has characteristic $0$.

(2) The formula for $\langle f,g\rangle_F$ is symmetric in $f$ and $g$. For each $\tau$ choose the
basis of $\nabla(\lambda_\tau)$ of Lemma~\ref{is-standard-basis}(1). Polynomials in disjoint sets of
variables multiply to linearly independent products, so the products $b_i=\otimes_\tau b_{i,\tau}$ of
basis vectors form a basis $(b_i)$ of $V(\lambda)$, and each $b_i$ has integer coefficients. Let
$Q=(\langle b_i,b_k\rangle_F)_{i,k}$, a symmetric matrix with rational entries. For a nonzero rational
vector $c$, the polynomial $p=\sum_ic_ib_i$ has rational coefficients and is nonzero, so
$c^{\mathsf T}Qc=\langle p,p\rangle_F=\sum_mp_m^2m!>0$. Hence $Q$ is positive definite over
$\mathbf Q$ and $\det Q\ne0$. As $K$ has characteristic $0$, $\det Q$ is nonzero in $K$, so the Fock
pairing is nondegenerate on $V(\lambda)$. The normalized pairing differs from it by the factor
$\prod_\tau\langle h_{\lambda_\tau},h_{\lambda_\tau}\rangle_F^{-1}$, which is defined and nonzero by
(1).

(3) Since $\langle f,g\rangle_F$ is the constant term of $f(\partial)g$, the Fock adjoint of the
substitution $f(x)\mapsto f(g^{\mathsf T}x)$ is $f(x)\mapsto f(gx)$, on all polynomials; the
determinant twists agree since $\det g^{\mathsf T}=\det g$.

(4) Let $S$ be the span of the vectors $\rho_\lambda(g)h_\lambda$. The substitutions
$f(x)\mapsto f(ux)$ act on the rows of $x$ and commute with the column operations that define
$\nabla(\nu)$, so they preserve each $\nabla(\nu)$; hence $V(\lambda)$ is stable under $\rho_\lambda$
and $S\subset V(\lambda)$.

For the converse fix $\tau$, write $\nu=\lambda_\tau^{\natural}$, and let $M\subset\mathrm{Pol}_N$ be
a subspace that is stable under the substitutions $f(x)\mapsto f(ux)$, $u\in GL_N(K)$. Let $f\in M$,
$1\le a,b\le N$ and $c\in\mathbf N$. The matrix $1+cE_{ba}$ is invertible: if $a\ne b$ its
determinant is $1$, and if $a=b$ it is $1+c\ne0$, since $K$ has characteristic $0$. The polynomial
$f((1+yE_{ba})x)=\sum_ky^kE^{(k)}_{ab}f$ in the indeterminate $y$ has the value
$f((1+cE_{ba})x)\in M$ at $y=c$. A polynomial in $y$ with coefficients in $\mathrm{Pol}_N$ whose values at
infinitely many points of $K$ lie in $M$ has all its coefficients in $M$: a linear functional that
vanishes on $M$, applied to the coefficients, gives a polynomial over $K$ with infinitely many
roots. So $E^{(k)}_{ab}f\in M$ for all $k$, $a$ and $b$. Take for $M$ the span of the polynomials
$h_\nu(u^{\mathsf T}x)$, $u\in GL_N(K)$, which is stable and contains $h_\nu$. Then $M$ contains the
Weyl module $\Delta_K(\nu)$, which is $\nabla(\nu)$ by Lemma~\ref{is-standard-basis}(2).

So every $f_\tau\in\nabla(\lambda_\tau)$ is a combination of polynomials
$h_{\lambda_\tau}(u_\tau^{\mathsf T}x_\tau)$, and $\otimes_\tau f_\tau$ is a combination of products
$\otimes_\tau h_{\lambda_\tau}(u_\tau^{\mathsf T}x_\tau)=\prod_\tau\det(u_\tau)^{-\operatorname{sh}(\lambda_\tau)}
\,\rho_\lambda(u)h_\lambda$. Hence $V(\lambda)\subset S$.

(5) Since $\{f,g\}_F=\langle\bar f,g\rangle_F$ and $h_\lambda$ has rational coefficients, so that
$\{h,h\}_F=\langle h,h\rangle_F$, the first identity is the definition of $\rho_\lambda(G)$.

Suppose that the matrices $s(G_\tau)$ are positive definite. Then so are the $s(G_\tau)^{\mathsf T}$;
let $B_\tau$ be the positive square root of $s(G_\tau)^{\mathsf T}$, so that
$s(G_\tau)^{\mathsf T}=B_\tau B_\tau^*$ with $B_\tau$ invertible. For a polynomial $x$ over $K$ write
$s(x)$ for the complex polynomial obtained by applying $s$ to its coefficients, and for complex
matrices $A=(A_\tau)$ let $S_A$ be the substitution $F\mapsto F((A_\tau x_\tau)_\tau)$, so that
$S_{B^*}S_B=S_{BB^*}$. On complex polynomials the Hermitian Fock form satisfies
$\{F,S_AF'\}_F=\{S_{A^*}F,F'\}_F$, by the argument of (3) with $\bar F$ in place of $f$. As
$s(\bar a)=\overline{s(a)}$, $s\{x,z\}_F=\{s(x),s(z)\}_F$. With
$c=\prod_\tau\det s(G_\tau)^{\operatorname{sh}(\lambda_\tau)}/\prod_\tau\langle h_{\lambda_\tau},
h_{\lambda_\tau}\rangle_F$, which is a positive real number, this gives
\[
 s\langle x,y\rangle_{\lambda,G}=c\,\{s(x),S_{B^*}S_Bs(y)\}_F=c\,\{S_Bs(x),S_Bs(y)\}_F .
\]
The right side is Hermitian in $(x,y)$, so $s\big(\overline{\langle x,y\rangle_{\lambda,G}}\big)=
\overline{s\langle x,y\rangle_{\lambda,G}}=s\langle y,x\rangle_{\lambda,G}$, and the Hermitian symmetry
follows since $s$ is injective. At $y=x$ the right side is $c$ times the Fock norm of $S_Bs(x)$, which
is a sum of nonnegative terms, so $\|x\|^2_{s,G}\geq0$.

(6) $\langle f,f\rangle_F=\sum_mf_m^2m!$ is a sum of nonnegative terms, positive at a monomial of
$f$.
\end{proof}

\begin{definition}\label{is-branching-data}
Let $K$ be a field of characteristic $0$, $n\geq0$, and let $\lambda$ and $\sigma$ be dominant
weights of lengths $n+1$ and $n$ (Definition~\ref{is-polynomial-model}).
\begin{enumerate}
\item For $1\le i\le n$ put $a_i=\sigma_i-\lambda_{i+1}$ and $b_i=\lambda_i-\sigma_i$. The weights
  \emph{interlace} if all $a_i,b_i\geq0$, and have the \emph{same shift} if $\operatorname{sh}(\sigma)=\operatorname{sh}(\lambda)$.
  Put $d_0=\sum_ia_i$, $n_{\lambda,\sigma}=\sum_i\lambda_i-\sum_i\sigma_i$, let $R_i$ be the set of
  integers $v$ with $\lambda_{i+1}\le v\le\min(\lambda_i,\lambda_{i+1}+d_0)$, and
  \[
   c_{\lambda,\sigma}=\prod_{i=1}^{n}\ \prod_{v\in R_i\smallsetminus\{\sigma_i\}}\ \prod_{j=1}^{n}
   \big((v-i)-(\sigma_j-j)\big)\in\mathbf Z.
  \]
\item On $(n+1)\times(n+1)$ matrices let $Z_i$ be the $i\times i$ minor on the rows $1,\dots,i-1,n+1$
  and the columns $1,\dots,i$. The \emph{branching vector} is
  $\xi_{\lambda,\sigma}=\prod_{i=1}^{n}\Delta_i^{a_i}Z_i^{b_i}$, and
  $\kappa_{\lambda,\sigma}=\langle\xi_{\lambda,\sigma},\xi_{\lambda,\sigma}\rangle_F/\langle h_\lambda,
  h_\lambda\rangle_F$.
\item For $\beta\in\mathbf N^n$, the \emph{coefficient map} $C_\beta\colon\mathrm{Pol}_{n+1}\to
  \mathrm{Pol}_n$ sends a monomial $x^m$ to the monomial $x^{m'}$ of its upper left $n\times n$ block
  if $m$ has degree $\beta_s$ in $x_{n+1,s}$ for $s\le n$ and degree $0$ in the last column, and to $0$
  otherwise.
\item For a finite set $\iota$ and weights $\lambda_\tau,\sigma_\tau$ that interlace, put
  $c=\prod_\tau c_{\lambda_\tau,\sigma_\tau}\kappa_{\lambda_\tau,\sigma_\tau}$, let $C$ be the
  coefficient map with exponents $b_i(\tau)$ in the variables of $\tau$, and $\mathrm{pr}_\sigma$ the
  Fock-orthogonal projection of $\mathrm{Pol}_n^{\iota}$ onto $V(\sigma)$. The \emph{branching map} is
  $\mathrm{br}=c\cdot\mathrm{pr}_\sigma\circ C\colon V(\lambda)\to V(\sigma)$, and for
  $t\in(K^\times)^{\iota}$, $\mathrm{br}_t=\prod_\tau t_\tau^{-n_{\lambda_\tau,\sigma_\tau}}\cdot\mathrm{br}$.
\end{enumerate}
\end{definition}
\leanref{Formalism/LocalToGlobal/Columns/FiniteSumNonzeroLargeDepth.lean}{Formalism.L2.Col.Br.Interlace}
\leanref{Formalism/LocalToGlobal/Columns/FiniteSumNonzeroLargeDepth.lean}{Formalism.L2.Col.Br.cBr}
\leanref{Formalism/LocalToGlobal/Columns/FiniteSumNonzeroLargeDepth.lean}{Formalism.L2.Col.Br.xiBr}
\leanref{Formalism/LocalToGlobal/Columns/FiniteSumNonzeroLargeDepth.lean}{Formalism.L2.Col.Br.kapBr}
\leanref{Formalism/LocalToGlobal/Columns/FiniteSumNonzeroLargeDepth.lean}{Formalism.L2.Col.Br.coefLast}
\leanref{Formalism/LocalToGlobal/Columns/FiniteSumNonzeroLargeDepth.lean}{Formalism.L2.Col.Br.brStd}
\leanref{Formalism/LocalToGlobal/Columns/FiniteSumNonzeroLargeDepth.lean}{Formalism.L2.Col.Br.brAt}

\begin{lemma}[Capelli identity]\label{is-capelli}
Let $K$ be a field of characteristic $0$ and $n\geq0$. On polynomials in the entries of a matrix $x$
with at least $n$ rows let $E_{ab}=\sum_cx_{a,c}\,\partial/\partial x_{b,c}$ for $a,b\le n$, and for
$u\in\mathbf Z$ let
\[
 C_n(u)=\sum_{\alpha\in S_n}\operatorname{sgn}(\alpha)\,F_{\alpha(1),1}F_{\alpha(2),2}\cdots F_{\alpha(n),n},
 \qquad F_{a,k}=(u+k-1)\delta_{ak}-E_{ak},
\]
the product taken in the order written. For $\mu\in\mathbf Z^n$ put
$\chi_\mu(u)=\prod_{j=1}^n(u+j-1-\mu_j)$.
\begin{enumerate}
\item $C_n(u)$ commutes with every $E_{ab}$, $a,b\le n$. In particular $C_n(u)C_n(u')=C_n(u')C_n(u)$
  for all $u,u'\in\mathbf Z$.
\item If $\mu\in\mathbf N^n$ is non-increasing, then $C_n(u)$ acts on
  $\nabla(\mu)\subset\mathrm{Pol}_n$ (Definition~\ref{is-polynomial-model}) as multiplication by
  $\chi_\mu(u)$, for every $u\in\mathbf Z$.
\end{enumerate}
\end{lemma}
\leanref{Formalism/BranchSandwichDivisors/BranchIsometryProofs/BranchMap.lean}{Formalism.BranchingDeterminantCap.capDet_central}
\leanref{Development/AbelianVarietyTheory/Proofs/DivisorProofs/CapelliDeterminantFirstRowsMemBmW.lean}{GlobalTheory.BrCompare.rCap_comm}
\leanref{Development/AbelianVarietyTheory/Proofs/DivisorProofs/BranchingAdmissibleFacts.lean}{GlobalTheory.BrCompare.capDet_dualWeyl}

The operator $C_n(u)$ is Capelli's column determinant \cite{Cap90}. It lies in the centre of the
algebra generated by the $E_{ab}$, which is (1), since $C_n(u')$ lies in that algebra. On an
irreducible module it acts by the scalar that it takes on a highest weight vector, which for $h_\mu$
is $\chi_\mu(u)$ \cite{HU91}.

\begin{lemma}\label{is-top-component}
Let $K$ be a field of characteristic $0$, $n\geq1$, and $\lambda\in\mathbf N^n$ non-increasing, with
$\lambda_{n+1}=0$. Let $\nabla(\lambda)$ be the space of polynomials $y$ in the entries of an
$(n+1)\times n$ matrix $x$ that are homogeneous of degree $\lambda_s$ in the $s$-th column for every
$s$ and invariant under adding $t$ times the $l$-th column to the $j$-th, for $l<j$ and an
indeterminate $t$. For $\beta\in
\mathbf Z^n$ let $y_\beta$ be the sum of the terms of $y$ of degree $\beta_s$ in $x_{n+1,s}$ for every
$s$, and put $\omega(\beta)=\sum_ss\beta_s$ and $z^\beta=\prod_sx_{n+1,s}^{\beta_s}$. Let
$y\in\nabla(\lambda)$, and let $\beta$ satisfy $y_\beta\ne0$ and $\omega(\beta')\le\omega(\beta)$ for
every $\beta'$ with $y_{\beta'}\ne0$. Then:
\begin{enumerate}
\item $0\le\beta_s\le\lambda_s-\lambda_{s+1}$ for $1\le s\le n$;
\item $y_\beta=z^\beta f$ for a polynomial $f\in\nabla(\lambda-\beta)$ in the entries of the upper
  $n\times n$ block of $x$;
\item $C_n(u)y_\beta=\chi_{\lambda-\beta}(u)\,y_\beta$ for every $u\in\mathbf Z$, with $C_n(u)$ and
  $\chi$ as in Lemma~\ref{is-capelli}.
\end{enumerate}
\end{lemma}
\leanref{Formalism/BranchSandwichDivisors/BranchIsometryProofs/BranchMap.lean}{Formalism.BranchingDeterminantCap.bm_eig}
\leanref{Development/AbelianVarietyTheory/Proofs/DivisorProofs/LastRowTopComponentDualWeylEqMonomialMulStatement.lean}{GlobalTheory.T.lastRow_top_component_dualWeyl_eq_monomial_mul}
\leanref{Development/AbelianVarietyTheory/Proofs/DivisorProofs/BranchingAdmissibleFacts.lean}{GlobalTheory.BrCompare.rCap_top}
\begin{proof}
For $l\ne j$ let $D_{lj}=\sum_{a\le n+1}x_{a,l}\,\partial/\partial x_{a,j}$. Since
$y(x(1+tE_{lj}))=\sum_kt^kD_{lj}^ky/k!$, a polynomial of column degrees $\mu$ lies in $\nabla(\mu)$ if
and only if $D_{lj}$ kills it for all $l<j$; this also holds for polynomials on the upper block, with
the sum over $a\le n$. Write $D_{lj}=D'_{lj}+Z_{lj}$ with $Z_{lj}=x_{n+1,l}\,\partial/\partial
x_{n+1,j}$. The operator $D'_{lj}$ does not change the degrees in the last row, and $Z_{lj}$ maps the
$\gamma$-part of a polynomial into its $(\gamma+e_l-e_j)$-part.

(2) Each monomial of $y_\beta$ has degree $\beta_s$ in $x_{n+1,s}$, so $y_\beta=z^\beta f$ with $f$ a
polynomial on the upper block, of column degrees $\lambda_s-\beta_s$. Let $l<j$. The $\beta$-part of
$D_{lj}y=0$ is $D'_{lj}y_\beta+Z_{lj}y_{\beta-e_l+e_j}$. As
$\omega(\beta-e_l+e_j)=\omega(\beta)+j-l>\omega(\beta)$, the second term is $0$. So
$z^\beta D'_{lj}f=D'_{lj}(z^\beta f)=0$, and $f$ is killed by the column operators of the upper block.
Once (1) shows that $\lambda-\beta$ is non-increasing, this gives $f\in\nabla(\lambda-\beta)$.

(1) Clearly $\beta_s\geq0$, and $\beta_n\le\lambda_n$ because $y$ has degree $\lambda_n$ in the $n$-th
column. Let $s<n$, $j=s+1$ and $k=\beta_s$, and suppose $k>\lambda_s-\lambda_j$. The operator
$D_{js}$ lowers the degree in the $s$-th column by one, so $D_{js}^{\lambda_s+1}y=0$. The operator
$D_{ss}-D_{jj}$ multiplies $D_{js}^iy$ by $\lambda_s-\lambda_j-2i$, and
$D_{sj}D_{js}-D_{js}D_{sj}=D_{ss}-D_{jj}$; as $D_{sj}y=0$, induction on $i$ gives
$D_{sj}D_{js}^{i+1}y=(i+1)(\lambda_s-\lambda_j-i)D_{js}^iy$. For $i>\lambda_s-\lambda_j$ the factor
$(i+1)(\lambda_s-\lambda_j-i)$ is nonzero in $K$, so descending induction from $i=\lambda_s$ gives
$D_{js}^iy=0$ for every $i>\lambda_s-\lambda_j$; in particular $D_{js}^ky=0$. Since $D'_{js}$ and
$Z_{js}$ commute,
$D_{js}^k=\sum_i\binom kiD'^{\,k-i}_{js}Z^i_{js}$, and the term with index $i$ maps the
$\gamma$-part of $y$ into the $(\beta+k(e_j-e_s))$-part only for $\gamma=\beta+(k-i)(e_j-e_s)$, which
has $\omega(\gamma)=\omega(\beta)+(k-i)(j-s)>\omega(\beta)$ unless $i=k$. So the
$(\beta+k(e_j-e_s))$-part of $D_{js}^ky$ is $Z_{js}^k(z^\beta f)=\beta_s!\,z^{\beta+k(e_j-e_s)}f\ne0$,
a contradiction. Hence $\beta_s\le\lambda_s-\lambda_{s+1}$. In particular $\lambda_s-\beta_s\geq
\lambda_{s+1}\geq\lambda_{s+1}-\beta_{s+1}$, so $\lambda-\beta$ is non-increasing, which completes (2).

(3) The operators $E_{ab}$ with $a,b\le n$ act on the upper block and commute with multiplication by
$z^\beta$. So $C_n(u)y_\beta=z^\beta C_n(u)f=\chi_{\lambda-\beta}(u)z^\beta f$ by (2) and
Lemma~\ref{is-capelli}(2).
\end{proof}

\begin{lemma}\label{is-capelli-separation}
Let $n\geq1$, let $\lambda\in\mathbf N^n$ be non-increasing, put $\lambda_{n+1}=0$, and let
$\sigma\in\mathbf N^n$ with $\lambda_{i+1}\le\sigma_i\le\lambda_i$ for $1\le i\le n$. Put
$b_i=\lambda_i-\sigma_i$ and $d=\sum_ib_i$, and let $\chi_\mu$ be as in Lemma~\ref{is-capelli}. Call
$\beta\in\mathbf Z^n$ \emph{admissible} if $\sum_s\beta_s=d$ and
$0\le\beta_s\le\lambda_s-\lambda_{s+1}$ for $1\le s\le n$; thus $b$ is admissible. Let $R_i$ and
$c_{\lambda,\sigma}$ be those of Definition~\ref{is-branching-data} for the weights
$(\lambda_1,\dots,\lambda_n,0)$ and $\sigma$, and let $\ell$ be the list of the integers $v-i+1$, for
$1\le i\le n$ and $v\in R_i\smallsetminus\{\sigma_i\}$.
\begin{enumerate}
\item If $\beta\ne\beta'$ are admissible, then there is $u\in\mathbf Z$ with
  $\chi_{\lambda-\beta'}(u)=0\ne\chi_{\lambda-\beta}(u)$.
\item If $\beta\ne b$ is admissible, then $\chi_{\lambda-\beta}(u)=0$ for some entry $u$ of $\ell$.
\item The product of the $\chi_\sigma(u)$ over the entries $u$ of $\ell$ is $c_{\lambda,\sigma}$.
\end{enumerate}
\end{lemma}
\leanref{Formalism/BranchSandwichDivisors/BranchIsometryProofs/BranchMap.lean}{Formalism.BranchingDeterminantCap.bm_sep}
\leanref{Formalism/BranchSandwichDivisors/BranchIsometryProofs/BranchMap.lean}{Formalism.BranchingDeterminantCap.bm_list_kills}
\leanref{Formalism/BranchSandwichDivisors/BranchIsometryProofs/BranchMap.lean}{Formalism.BranchingDeterminantCap.bm_list_prod}
\begin{proof}
(1) Choose $i$ with $\beta_i\ne\beta'_i$, and put $u=\lambda_i-\beta'_i-i+1$. The factor $j=i$ of
$\chi_{\lambda-\beta'}(u)=\prod_j(u+j-1-\lambda_j+\beta'_j)$ is $0$. Put $\mu=\lambda-\beta$ and
$\mu'=\lambda-\beta'$; by admissibility $\lambda_{s+1}\le\mu_s\le\lambda_s$ and
$\lambda_{s+1}\le\mu'_s\le\lambda_s$ for every $s$. The $j$-th factor of $\chi_{\lambda-\beta}(u)$ is
$(\mu'_i-i)-(\mu_j-j)$. For $j=i$ it is $\beta_i-\beta'_i\ne0$. For $j<i$,
$\mu'_i\le\lambda_i\le\lambda_{j+1}\le\mu_j$, so it is at most $j-i<0$. For $j>i$,
$\mu'_i\geq\lambda_{i+1}\geq\lambda_j\geq\mu_j$, so it is at least $j-i>0$. So
$\chi_{\lambda-\beta}(u)\ne0$.

(2) Choose $i$ with $\beta_i\ne b_i$, and put $v=\lambda_i-\beta_i$ and $u=v-i+1$. The factor $j=i$ of
$\chi_{\lambda-\beta}(u)$ is $u+i-1-v=0$. It remains to see that $v\in R_i\smallsetminus\{\sigma_i\}$.
First $v\ne\lambda_i-b_i=\sigma_i$. Next $\lambda_{i+1}\le v\le\lambda_i$ by admissibility. Finally,
every term $\lambda_s-\beta_s-\lambda_{s+1}$ is nonnegative, so
\[
 v-\lambda_{i+1}\le\sum_s(\lambda_s-\beta_s-\lambda_{s+1})=\sum_s(\lambda_s-\lambda_{s+1})-d
 =\sum_s(\sigma_s-\lambda_{s+1})=d_0,
\]
the integer $d_0$ of Definition~\ref{is-branching-data}. So $v\le\lambda_{i+1}+d_0$.

(3) For $u=v-i+1$, $\chi_\sigma(u)=\prod_j(v-i+j-\sigma_j)=\prod_j\big((v-i)-(\sigma_j-j)\big)$.
The product of these over $1\le i\le n$ and $v\in R_i\smallsetminus\{\sigma_i\}$ is the definition of
$c_{\lambda,\sigma}$.
\end{proof}

\begin{lemma}\label{is-capelli-kill}
Let $K$ be a field of characteristic $0$, and let $n$, $\lambda$, $\sigma$, $b$, $d$, the admissible
vectors and the list $\ell$ be as in Lemma~\ref{is-capelli-separation}. On polynomials in the entries
of an $(n+1)\times n$ matrix, let $\nabla(\lambda)$, the components $y_\gamma$ and $\omega$ be as in
Lemma~\ref{is-top-component}, and let $C_n(u)$ be the operators of Lemma~\ref{is-capelli}. Let $W$ be
the space of the $y\in\nabla(\lambda)$ that are homogeneous of degree $d$ in the entries of the last
row.
\begin{enumerate}
\item Each $C_n(u)$ maps $W$ into $W$, and $C_n(u)(y_\gamma)=(C_n(u)y)_\gamma$ for every polynomial
  $y$ and every $\gamma\in\mathbf Z^n$.
\item Let $u_1,\dots,u_r$ be the entries of $\ell$. If $y\in W$ and $y_b=0$, then
  $C_n(u_1)\cdots C_n(u_r)y=0$.
\end{enumerate}
\end{lemma}
\leanref{Development/AbelianVarietyTheory/Proofs/DivisorProofs/CapelliDeterminantFirstRowsMemBmW.lean}{GlobalTheory.BrCompare.capelliDeterminantFirstRows_mem_bmW}
\leanref{Development/AbelianVarietyTheory/Proofs/DivisorProofs/CapelliDeterminantFirstRowsMemBmW.lean}{GlobalTheory.BrCompare.zcomp_comm_rCap}
\leanref{Development/AbelianVarietyTheory/Proofs/DivisorProofs/BranchingAdmissibleFacts.lean}{GlobalTheory.BrCompare.rCap_kill}
\leanref{Formalism/BranchSandwichDivisors/BranchIsometryProofs/BranchMap.lean}{Formalism.BranchingDeterminantCap.bm_kill}
\leanref{Formalism/BranchSandwichDivisors/BranchIsometryProofs/BranchMap.lean}{Formalism.BranchingDeterminantCap.bmE_sep_op}
\begin{proof}
(1) Let $a,b\le n$. The operator $E_{ab}$ does not involve the last row, so it preserves the degree
in each entry $x_{n+1,s}$ and commutes with $y\mapsto y_\gamma$. It preserves the degree in each
column, and a direct computation shows that it commutes with the column operators $D_{lj}$ of the
proof of Lemma~\ref{is-top-component}. As $\nabla(\lambda)$ is the space of polynomials of column
degrees $\lambda$ killed by the $D_{lj}$ with $l<j$, $E_{ab}$ maps $\nabla(\lambda)$ and $W$ into
themselves. So does $C_n(u)$, a polynomial in these operators.

(2) Write $C(u)=C_n(u)$ and $\chi_\beta=\chi_{\lambda-\beta}$. For a finite set $S\subset\mathbf Z^n$
let $W_S$ be the space of $y\in W$ with $y_\gamma=0$ for $\gamma\notin S$; by (1) it is stable under
every $C(u)$. Two facts are used.
\begin{enumerate}
\item[(a)] If $y\in W$, $y_\beta\ne0$ and $\omega(\beta)\geq\omega(\gamma)$ for every $\gamma$ with
  $y_\gamma\ne0$, then $\beta$ is admissible and $C(u)y_\beta=\chi_\beta(u)y_\beta$ for every $u$. This
  is Lemma~\ref{is-top-component}(1) and (3), with $\sum_s\beta_s=d$ because $y$ is homogeneous of
  degree $d$ in the last row.
\item[(b)] The $C(u)$ commute with each other, by Lemma~\ref{is-capelli}(1).
\end{enumerate}

\emph{Step 1: separating operators.} Let $\beta_0$ be admissible and $S$ a finite set with
$\beta_0\notin S$. Then there are an integer $q\ne0$ and an operator $Q$, a product of operators
$C(u)-c$ with integers $u$ and $c$, such that $Qy=0$ for every $y\in W_S$, and $Qf=qf$ for every
polynomial $f$ with $C(u)f=\chi_{\beta_0}(u)f$ for all $u$. The proof is by induction on $|S|$. If $S$
is empty, $W_S=0$; take $Q=1$ and $q=1$. Otherwise choose $\beta_1\in S$ with $\omega(\beta_1)$
maximal, and let $Q'$ and $q'$ be given by induction for $\beta_0$ and $S\smallsetminus\{\beta_1\}$.
If $\beta_1$ is not admissible, then by (a) no $y\in W_S$ has $y_{\beta_1}\ne0$, so
$W_S=W_{S\smallsetminus\{\beta_1\}}$; take $Q=Q'$ and $q=q'$. If $\beta_1$ is admissible, then
$\beta_1\ne\beta_0$, and Lemma~\ref{is-capelli-separation}(1) gives $u$ with
$\chi_{\beta_0}(u)\ne\chi_{\beta_1}(u)$. Put $T=C(u)-\chi_{\beta_1}(u)$, $Q=Q'T$ and
$q=q'\big(\chi_{\beta_0}(u)-\chi_{\beta_1}(u)\big)\ne0$. For $f$ as above,
$Tf=(\chi_{\beta_0}(u)-\chi_{\beta_1}(u))f$, so $Qf=qf$. For $y\in W_S$, $Ty\in W_S$ and
$(Ty)_{\beta_1}=T(y_{\beta_1})=0$: either $y_{\beta_1}=0$, or $\beta_1$ has the largest $\omega$ among
the $\gamma$ with $y_\gamma\ne0$, and (a) gives $C(u)y_{\beta_1}=\chi_{\beta_1}(u)y_{\beta_1}$. So
$Ty\in W_{S\smallsetminus\{\beta_1\}}$, and $Qy=Q'Ty=0$.

\emph{Step 2: the product kills.} Put $P=C(u_1)\cdots C(u_r)$. By (1) and (b), $P$ and the operators
$Q$ of Step 1 commute with each other, with every $C(u)$ and with $y\mapsto y_\gamma$, and they map
$W$ into $W$. Let $y\in W$ with $y_b=0$; we show $Py=0$ by induction on the number of $\gamma$ with
$y_\gamma\ne0$. If $y=0$ there is nothing to show. Otherwise choose $\beta_0$ with $y_{\beta_0}\ne0$ and
$\omega(\beta_0)$ maximal. By (a), $\beta_0$ is admissible and
$C(u)y_{\beta_0}=\chi_{\beta_0}(u)y_{\beta_0}$ for every $u$; and $\beta_0\ne b$ because $y_b=0$. Let
$S$ be the set of $\gamma\ne\beta_0$ with $y_\gamma\ne0$, and let $Q$ and $q$ be as in Step 1 for
$\beta_0$ and $S$.

For every $u$, $C(u)y-\chi_{\beta_0}(u)y$ lies in $W_S$: its $\beta_0$-component is
$C(u)y_{\beta_0}-\chi_{\beta_0}(u)y_{\beta_0}=0$, and its other components vanish where those of $y$
do. So $Q$ kills it, and $C(u)Qy=\chi_{\beta_0}(u)Qy$. By
Lemma~\ref{is-capelli-separation}(2), $\chi_{\beta_0}(u_k)=0$ for some $k$, so
$PQy=\chi_{\beta_0}(u_1)\cdots\chi_{\beta_0}(u_r)\,Qy=0$.

The polynomial $qy-Qy$ lies in $W$. Its $\beta_0$-component is $qy_{\beta_0}-Qy_{\beta_0}=0$, by
Step 1 applied to $f=y_{\beta_0}$; its $b$-component is $qy_b-Qy_b=0$; and its other nonzero
components are among those of $y$. By induction $P(qy-Qy)=0$. Hence $qPy=PQy+P(qy-Qy)=0$, and $Py=0$
because $q\ne0$ in $K$.
\end{proof}

\begin{lemma}\label{is-branching-vector}
Let $\lambda,\sigma$ be dominant weights of lengths $n+1$ and $n$ that interlace and have the same
shift (Definition~\ref{is-branching-data}), and $b=(b_1,\dots,b_n)$.
\begin{enumerate}
\item $\xi_{\lambda,\sigma}\in\nabla(\lambda)$, $C_b(\xi_{\lambda,\sigma})=h_\sigma$,
  $\xi_{\lambda,\sigma}\ne0$, and $c_{\lambda,\sigma}\ne0$. Over $\mathbf R$, $\kappa_{\lambda,\sigma}>0$.
\item For every $x\in\nabla(\lambda)$,
  $\langle x,\xi_{\lambda,\sigma}\rangle_F\,\langle h_\sigma,h_\sigma\rangle_F
  =\langle\xi_{\lambda,\sigma},\xi_{\lambda,\sigma}\rangle_F\,\langle C_b(x),h_\sigma\rangle_F$.
\item For a finite set $\iota$ and such pairs $(\lambda_\tau,\sigma_\tau)$, let $g\in GL_n(K)^{\iota}$,
  $t\in(K^\times)^\iota$ and $g'\in GL_{n+1}(K)^{\iota}$ with $g'_\tau$ block diagonal with blocks
  $g_\tau$ and $t_\tau$. Then $\mathrm{br}(\rho_\lambda(g')x)=\prod_\tau t_\tau^{n_{\lambda_\tau,
  \sigma_\tau}}\rho_\sigma(g)\,\mathrm{br}(x)$ for $x\in V(\lambda)$.
\end{enumerate}
\end{lemma}
\leanref{Development/AbelianVarietyTheory/Proofs/DivisorProofs/LastMapNorm.lean}{GlobalTheory.BranchNorm.xi_mem}
\leanref{Development/AbelianVarietyTheory/Proofs/DivisorProofs/LastMapNorm.lean}{GlobalTheory.BranchNorm.coefOne_xi}
\leanref{Development/AbelianVarietyTheory/Proofs/DivisorProofs/LastMapNorm.lean}{GlobalTheory.BranchNorm.xi_ne}
\leanref{Development/AbelianVarietyTheory/Proofs/DivisorProofs/LastMapNorm.lean}{GlobalTheory.BranchNorm.cBr_ne}
\leanref{Development/AbelianVarietyTheory/Proofs/DivisorProofs/LastMapNorm.lean}{GlobalTheory.BranchNorm.kapBr_pos}
\leanref{Development/AbelianVarietyTheory/Proofs/DivisorProofs/LastMapNorm.lean}{GlobalTheory.BranchNorm.fock_xi_one}
\leanref{Development/AbelianVarietyTheory/Proofs/DivisorProofs/BranchingAdmissibleFacts.lean}{GlobalTheory.BrCompare.fock_xi}
\leanref{Development/AbelianVarietyTheory/Proofs/GlobalInterpolation/StatementWithTransposeOfBranching.lean}{GlobalTheory.BranchFormula.brStd_rhoK}
\begin{proof}
Since $\operatorname{sh}(\sigma)=\operatorname{sh}(\lambda)=\lambda_{n+1}$, subtracting $\lambda_{n+1}$
from all entries of $\lambda$ and $\sigma$ changes neither the $a_i$, $b_i$, $c_{\lambda,\sigma}$ and
$\xi_{\lambda,\sigma}$ nor the spaces $\nabla(\lambda)$, $\nabla(\sigma)$ and the vector $h_\sigma$.
So assume $\lambda_{n+1}=\operatorname{sh}(\sigma)=0$; then $\lambda=\lambda^{\natural}$,
$\sigma=\sigma^{\natural}$ and $\lambda_{i+1}\le\sigma_i\le\lambda_i$.

(1) The minors $\Delta_i$ and $Z_i$, $i\le n$, lie on the first $i$ columns, so
$\xi=\xi_{\lambda,\sigma}$ does not involve the last column, and neither does any element of
$\nabla(\lambda)$, which has degree $\lambda_{n+1}=0$ there. So work with polynomials on
$(n+1)\times n$ matrices, with the notation of Lemma~\ref{is-top-component}. For non-increasing
$\mu\in\mathbf N^n$ and $\beta\in\mathbf N^n$, call a polynomial $G$ a \emph{lift} of a polynomial $f$
on the upper $n\times n$ block, of weight $\mu$ and exponent $\beta$, if $G\in\nabla(\mu)$, $G$ is
homogeneous of degree $\sum_s\beta_s$ in the last row, and every monomial of $G-z^\beta f$ has degree
$\gamma$ in the last row with $\omega(\gamma)<\omega(\beta)$; here $z^\beta=\prod_sx_{n+1,s}^{\beta_s}$.

Lifts multiply. Let $G$ be a lift of $f$ of weight $\mu$ and exponent $\beta$, and $G'$ a lift of
$f'$ of weight $\mu'$ and exponent $\beta'$. Then $GG'$ has column degrees $\mu+\mu'$ and is invariant
under the column additions, so $GG'\in\nabla(\mu+\mu')$; it is homogeneous of degree
$\sum_s(\beta_s+\beta'_s)$ in the last row; and in
$GG'-z^{\beta+\beta'}ff'=(G-z^\beta f)G'+z^\beta f(G'-z^{\beta'}f')$ every monomial has a degree
$\gamma$ in the last row with $\omega(\gamma)<\omega(\beta)+\omega(\beta')$. Indeed $\omega$ is
additive, every monomial of $G'$ has a last-row degree $\gamma'$ with $\omega(\gamma')\le\omega(\beta')$,
and every monomial of $z^\beta f$ has last-row degree $\beta$. So $GG'$ is a lift of $ff'$ of weight
$\mu+\mu'$ and exponent $\beta+\beta'$.

Write $1^i\in\mathbf N^n$ for the vector with $i$ entries $1$ followed by zeros, and $e_i$ for the
$i$-th unit vector. A minor on the first $i$ columns lies in $\nabla(1^i)$: adding $y$ times the
$l$-th column to the $j$-th, for $l<j$, leaves it unchanged, because if $j\le i$ both columns are
among its columns, and if $j>i$ it does not involve the $j$-th column. So $\Delta_i$ is a lift of
itself of weight $1^i$ and exponent $0$, as it does not involve the last row. Expanding $Z_i$ along
its last row gives $Z_i=\sum_{s\le i}\pm x_{n+1,s}M_s$ with minors $M_s$ of the upper block, where
the term $s=i$ is $x_{n+1,i}\Delta_{i-1}$ ($\Delta_0=1$) and the term $s$ has last-row degree $e_s$
with $\omega(e_s)=s$. So $Z_i$ is a lift of $\Delta_{i-1}$ of weight $1^i$ and exponent $e_i$.
Multiplying, $\xi=\prod_i\Delta_i^{a_i}Z_i^{b_i}$ is a lift of
$\prod_i\Delta_i^{a_i}\Delta_{i-1}^{b_i}$ of exponent $\sum_ib_ie_i=b$ and weight
$\sum_i(a_i+b_i)1^i$, whose $s$-th entry is $\sum_{i\geq s}(\lambda_i-\lambda_{i+1})=\lambda_s$.
Moreover $\prod_i\Delta_i^{a_i}\Delta_{i-1}^{b_i}=h_\sigma$, since the exponent of $\Delta_i$ is
$a_i+b_{i+1}=\sigma_i-\sigma_{i+1}$ for $i<n$ and $a_n=\sigma_n$.

Hence $\xi\in\nabla(\lambda)$, $\xi$ is homogeneous of degree $d=\sum_sb_s$ in the last row, its
component of last-row degree $b$ is $\xi_b=z^bh_\sigma$, and every $\gamma\ne b$ with $\xi_\gamma\ne0$
has $\omega(\gamma)<\omega(b)$. In particular $C_b(\xi)=h_\sigma$, and $\xi\ne0$ because
$h_\sigma\ne0$.

Each factor $(v-i)-(\sigma_j-j)$ of $c_{\lambda,\sigma}$ is nonzero. For $j=i$ it is $v-\sigma_i\ne0$.
For $j<i$, $\sigma_j\geq\lambda_{j+1}\geq\lambda_i\geq v$, so $\sigma_j-j>v-i$. For $j>i$,
$\sigma_j\le\lambda_j\le\lambda_{i+1}\le v$, so $\sigma_j-j<v-i$. Over $\mathbf R$,
$\kappa_{\lambda,\sigma}$ is a quotient of two positive Fock squares (Lemma~\ref{is-fock-facts}(6)).

(2) Work on $(n+1)\times n$ matrices, as in (1). Write $b!=\prod_sb_s!$ and
$\rho=(\sigma_1,\dots,\sigma_n,d)$, and put
\[
 \psi(x)=\langle h_\sigma,h_\sigma\rangle_F\,b!\,\langle x,\xi\rangle_F-\langle\xi,\xi\rangle_F\,
 \langle x,z^bh_\sigma\rangle_F .
\]
For any polynomial $g$ on the upper $n\times n$ block, the monomials of $z^bg$ are the products of
$z^b$ with those of $g$, and $(z^bx^m)!=b!\,m!$; so
$\langle x,z^bg\rangle_F=b!\,\langle C_b(x),g\rangle_F$. Since differentiation in $x_{a,c}$ is
Fock adjoint to multiplication by $x_{a,c}$, the Fock adjoint of $E_{jl}$ is $E_{lj}$.

\emph{Three properties of $\psi$.} (i) $\psi(E_{jl}x)=0$ for $l<j\le n$. Indeed $E_{lj}$ kills $\xi$
and $z^bh_\sigma$. It is a derivation, and it kills every minor whose rows are $1,\dots,i$ or
$1,\dots,i-1,n+1$: if the row $j$ of such a minor is replaced by the row $l$, then either $j$ is not
one of its rows, or $l<j$ is one as well and the result has two equal rows. Moreover $E_{lj}$ does not
involve the last row. (ii) $\psi(E_{aa}x)=\rho_a\psi(x)$ for $a\le n+1$, since $E_{aa}$ multiplies a
polynomial homogeneous of degree $e$ in the $a$-th row by $e$, and $\xi$ and $z^bh_\sigma$ have degree
$\rho_a$ in the $a$-th row. For $\xi$ and $a\le n$ this degree is
$\sum_{i\geq a}a_i+\sum_{i>a}b_i=\sigma_a-\lambda_{n+1}=\sigma_a$, and for $a=n+1$ it is
$\sum_ib_i=d$. (iii) $\psi(\xi)=0$, because $\langle\xi,z^bh_\sigma\rangle_F=b!\langle C_b(\xi),h_\sigma
\rangle_F=b!\langle h_\sigma,h_\sigma\rangle_F$ by (1).

\emph{The Capelli operators.} With $\chi_\sigma(u)=\prod_{j=1}^n(u+j-1-\sigma_j)$ as in
Lemma~\ref{is-capelli}, $\psi(C_n(u)x)=\chi_\sigma(u)\psi(x)$ for all $x$ and $u$. Expand $C_n(u)$ as the sum over
$\alpha\in S_n$ of $\operatorname{sgn}(\alpha)F_{\alpha(1),1}\cdots F_{\alpha(n),n}$. If $\alpha$ is
the identity, each factor $F_{k,k}=(u+k-1)-E_{kk}$ multiplies $\psi$ by $u+k-1-\sigma_k$, by (ii).
Otherwise let $k$ be the least index with $\alpha(k)\ne k$. The factors before the $k$-th are
$F_{j,j}$, which multiply $\psi$ by scalars, and $\alpha(k)>k$ because $\alpha$ fixes $1,\dots,k-1$;
so the $k$-th factor is $-E_{\alpha(k),k}$ and contributes $0$ by (i).

\emph{The induction.} Each row-weight component of an $x\in\nabla(\lambda)$ lies in $\nabla(\lambda)$,
because the $E_{aa}$ commute with the column operators. If a component $x'$ has weight $\rho'\ne\rho$,
choose $a$ with $\rho'_a\ne\rho_a$; then $\rho'_a\psi(x')=\psi(E_{aa}x')=\rho_a\psi(x')$ by (ii), so
$\psi(x')=0$. It remains to show $\psi(y)=0$ for $y\in\nabla(\lambda)$ of row weight $\rho$. Such a
$y$ is homogeneous of degree $\rho_{n+1}=d$ in the last row, so it lies in the space $W$ of
Lemma~\ref{is-capelli-kill}. Argue by induction on the number $t(y)$ of $\gamma$ with
$y_\gamma\ne0$. If $y=0$ there is nothing to prove. Otherwise choose $\beta$ with $y_\beta\ne0$ and
$\omega(\beta)\geq\omega(\gamma)$ for every $\gamma$ with $y_\gamma\ne0$. By
Lemma~\ref{is-top-component}(1), and since $\sum_s\beta_s=d$, $\beta$ is admissible in the sense of
Lemma~\ref{is-capelli-separation}, and so is $b$.

Suppose first that $\beta\ne b$. By Lemma~\ref{is-capelli-separation}(1) there is $u\in\mathbf Z$ with
$\chi_{\lambda-\beta}(u)\ne\chi_{\lambda-b}(u)=\chi_\sigma(u)$. Let
$y'=C_n(u)y-\chi_{\lambda-\beta}(u)y$. It lies in $W$ by Lemma~\ref{is-capelli-kill}(1), and it has
row weight $\rho$, because $C_n(u)$ commutes with the $E_{aa}$, $a\le n$, by
Lemma~\ref{is-capelli}(1), and does not change the degree in the last row. Its components are
$y'_\gamma=C_n(u)y_\gamma-\chi_{\lambda-\beta}(u)y_\gamma$, by Lemma~\ref{is-capelli-kill}(1). They
vanish where $y_\gamma$ does, and $y'_\beta=0$ by Lemma~\ref{is-top-component}(3). So $t(y')<t(y)$,
and $\psi(y')=0$ by induction. By the Capelli operators,
$\psi(y')=\big(\chi_\sigma(u)-\chi_{\lambda-\beta}(u)\big)\psi(y)$, so $\psi(y)=0$.

Suppose now that $\beta=b$. By Lemma~\ref{is-top-component}(2), $y_b=z^bf$ with
$f\in\nabla(\lambda-b)=\nabla(\sigma)$ on the upper block, and $f$ has row weight $\sigma$ because
$y$ has row weight $\rho$. In the basis of Lemma~\ref{is-standard-basis}(1), the vectors of row weight
$\sigma$ are those whose tableau has content $\sigma$; the only semistandard tableau of shape and
content $\sigma$ has all entries of its $i$-th row equal to $i$, and its bideterminant is $h_\sigma$.
So $f=ah_\sigma$ for some $a\in K$. By (1), $\xi$ lies in $W$ and $\xi_b=z^bh_\sigma$. Hence
$y-a\xi\in W$ and $(y-a\xi)_b=0$. Let $u_1,\dots,u_r$ be the entries of the list $\ell$ of
Lemma~\ref{is-capelli-separation}. Lemma~\ref{is-capelli-kill}(2) gives
$C_n(u_1)\cdots C_n(u_r)(y-a\xi)=0$. Applying $\psi$ and the Capelli operators $r$ times, and then
Lemma~\ref{is-capelli-separation}(3),
\[
 0=\chi_\sigma(u_1)\cdots\chi_\sigma(u_r)\,\psi(y-a\xi)=c_{\lambda,\sigma}\,\psi(y-a\xi).
\]
As $c_{\lambda,\sigma}\ne0$ by (1), $\psi(y-a\xi)=0$, and $\psi(y)=a\psi(\xi)=0$ by (iii).

So $\psi=0$ on $\nabla(\lambda)$. Substituting $\langle x,z^bh_\sigma\rangle_F=b!\langle C_b(x),
h_\sigma\rangle_F$ and dividing by $b!$ gives (2).

(3) Since $g'_\tau$ is block diagonal, the substitution $x\mapsto g'^{\mathsf T}x$ multiplies the last
row by $t_\tau$ and keeps the last column in the last column, so it commutes with $C$ up to
$\prod_\tau t_\tau^{\sum_sb_s(\tau)}$ and the action of $g$ on the upper left block. The projection
$\mathrm{pr}_\sigma$ commutes with the substitution by $g^{\mathsf T}$: $V(\sigma)$ is stable under the
substitutions by $g^{\mathsf T}$ and by $g$, and by the adjointness of Lemma~\ref{is-fock-facts}(3)
the substitution by $g^{\mathsf T}$ preserves the Fock-orthogonal complement of $V(\sigma)$. The determinant
twists differ by $\prod_\tau t_\tau^{\operatorname{sh}(\lambda_\tau)}$, and $\sum_sb_s(\tau)+\operatorname{sh}(\lambda_\tau)=
n_{\lambda_\tau,\sigma_\tau}$ because the shifts agree.
\end{proof}

\begin{remark}\label{is-branching-vector-order}
Part (2) of Lemma~\ref{is-branching-vector} has a shorter proof that does not use
$c_{\lambda,\sigma}\ne0$ or Lemma~\ref{is-capelli-kill}: the induction can be run in another order.
For $y\ne0$ of row weight $\rho$, let $\omega(y)$ be the largest $\omega(\gamma)$ with $y_\gamma\ne0$,
and $t'(y)$ the number of $\gamma$ with $y_\gamma\ne0$ and $\omega(\gamma)=\omega(y)$; argue by
induction on $(\omega(y),t'(y))$ in the lexicographic order, a well order on these pairs since
$\omega(y)\geq0$. When the top index $\beta$ differs from $b$, the polynomial $y'$ of the proof is $0$
or smaller than $y$, since $y'_\beta=0$ and its nonzero components are among those of $y$. When
$\beta=b$, every $\gamma\ne b$ with $\xi_\gamma\ne0$ has $\omega(\gamma)<\omega(b)$, by the proof of
(1). So $(y-a\xi)_\gamma=y_\gamma$ for every $\gamma\ne b$ with $\omega(\gamma)\geq\omega(b)$, and
$y-a\xi$ is $0$ or smaller than $y$. The induction hypothesis gives $\psi(y-a\xi)=0$ directly, and
$\psi(y)=a\psi(\xi)=0$.
\end{remark}

\begin{lemma}\label{is-abstract-coisometry}
Let $K$ be a field with an involution and $s\colon K\to\mathbf C$ an embedding with
$s(\bar a)=\overline{s(a)}$. Let $B$ and $B'$ be sesquilinear forms on $K$-vector spaces $V$ and
$V'$ with $\overline{B(x,y)}=B(y,x)$, $\overline{B'(x,y)}=B'(y,x)$ and
$\operatorname{Re}s\,B(x,x)\geq0$. Let $P\colon V\to V'$ and $A\colon V'\to V$ be linear maps with
$B(Ay,x)=B'(y,Px)$, and $k\in K$ with $s(k)=\kappa>0$ and $PA=k$. Then
$\operatorname{Re}s\,B'(Pf,Pf)\le\kappa\operatorname{Re}s\,B(f,f)$ for all $f\in V$.
\end{lemma}
\leanref{Development/AbelianVarietyTheory/Proofs/DivisorProofs/LastMapNorm.lean}{GlobalTheory.BranchNorm.norm_le}
\begin{proof}
Expand $0\le\operatorname{Re}s\,B(kf-APf,kf-APf)$. Here
$B(f,APf)=\overline{B(APf,f)}=\overline{B'(Pf,Pf)}=B'(Pf,Pf)$, and $B(APf,APf)=B'(Pf,PAPf)=
kB'(Pf,Pf)$. Since $s(k)=s(\bar k)=\kappa$, the expansion is
$\kappa^2\operatorname{Re}sB(f,f)-2\kappa\operatorname{Re}sB'(Pf,Pf)+\kappa\operatorname{Re}sB'(Pf,Pf)
\geq0$, which gives the claim after dividing by $\kappa$.
\end{proof}

\begin{proposition}\label{is-branching-norm}
Let $K$ be a field with an involution and $s\colon K\to\mathbf C$ an embedding with
$s(\bar a)=\overline{s(a)}$. Let $\iota$ be a finite set, and for $\tau\in\iota$ let $\lambda_\tau$
and $\sigma_\tau$ be dominant weights of lengths $n+1$ and $n$ that interlace and have the same shift
(Definition~\ref{is-branching-data}). Let $G'_\tau\in M_n(K)$ with $s(G'_\tau)$ positive definite,
and let $g_\tau,t_\tau\in K$ and $T>0$ with $|s(t_\tau)|^2s(g_\tau)=T$. Let $G_\tau$ be the block
diagonal matrix with blocks $G'_\tau$ and $g_\tau$, and put
\[
 \beta=\prod_\tau T^{-n_{\lambda_\tau,\sigma_\tau}}\,c_{\lambda_\tau,\sigma_\tau}^2\,\kappa_{\lambda_\tau,\sigma_\tau},
\]
with $\kappa$ computed over $\mathbf R$. Then $\beta>0$, and:
\begin{enumerate}
\item there is a linear map $A\colon V(\sigma)\to V(\lambda)$ with $\langle Ay,x\rangle_{\lambda,G}=
  \langle y,\mathrm{br}_tx\rangle_{\sigma,G'}$ and $s\langle Ay,Ay'\rangle_{\lambda,G}=\beta\,
  s\langle y,y'\rangle_{\sigma,G'}$ for all $x,y,y'$, so $\mathrm{br}_t$ is a scaled coisometry;
\item $\|\mathrm{br}_tf\|^2_{s,G'}\le\beta\|f\|^2_{s,G}$ for every $f\in V(\lambda)$;
\item equality holds in (2) for some $f\ne0$.
\end{enumerate}
\end{proposition}
\leanref{Development/AbelianVarietyTheory/BlockTheory/DivisorBase/ScaledCoisometryNorms.lean}{GlobalTheory.IndefiniteRowSideFullHecke.ScaledCoisometryNorms.BrNorm}
\leanref{Development/AbelianVarietyTheory/Proofs/DivisorProofs/LastMapNorm.lean}{GlobalTheory.BranchNorm.brNorm}
\leanref{Development/AbelianVarietyTheory/Proofs/DivisorProofs/LastMapNorm.lean}{GlobalTheory.BranchNorm.lastMapNorm}
\begin{proof}
Since $s$ is injective, $K$ has characteristic $0$. Each $|s(t_\tau)|^2$ is positive, since
$T>0$, so $s(g_\tau)=T/|s(t_\tau)|^2$ is real and positive and $s(G_\tau)$ is positive definite; in
particular $G_\tau$ and $G'_\tau$ are invertible. Each factor of $\beta$ is positive by
Lemma~\ref{is-branching-vector}(1).

\emph{The transpose.} By Lemma~\ref{is-fock-facts}(2) there is a unique linear map
$\mathrm{tr}\colon V(\sigma)\to V(\lambda)$ with $\langle\mathrm{tr}\,y,x\rangle_\lambda=\langle y,
\mathrm{br}\,x\rangle_\sigma$. Put $\xi=\otimes_\tau\xi_{\lambda_\tau,\sigma_\tau}\in V(\lambda)$ and
$c_0=\prod_\tau c_{\lambda_\tau,\sigma_\tau}$.

\emph{Step 1: $\mathrm{br}\,\xi=c\,h_\sigma$.} The coefficient map is multiplicative over $\tau$, so
$C(\xi)=\otimes_\tau C_{b(\tau)}(\xi_{\lambda_\tau,\sigma_\tau})=h_\sigma$ by
Lemma~\ref{is-branching-vector}(1), and $\mathrm{pr}_\sigma h_\sigma=h_\sigma$.

\emph{Step 2: $\mathrm{tr}\,h_\sigma=c_0\,\xi$.} For $F=\otimes_\tau f_\tau$ with $f_\tau\in
\nabla(\lambda_\tau)$, Lemma~\ref{is-branching-vector}(2) for each $\tau$, multiplied over $\tau$ with
Lemma~\ref{is-fock-facts}(1), gives $c_0\langle F,\xi\rangle_\lambda=\langle\mathrm{br}\,F,h_\sigma
\rangle_\sigma$, using that $\mathrm{pr}_\sigma$ is orthogonal and $h_\sigma\in V(\sigma)$. By linearity
this holds for all $F\in V(\lambda)$. With the symmetry of the pairings it says
$\langle\mathrm{tr}\,h_\sigma,F\rangle_\lambda=\langle c_0\xi,F\rangle_\lambda$ for all $F$, and
nondegeneracy gives the claim.

\emph{Step 3: $\mathrm{br}\circ\mathrm{tr}=c_0c$.} By Lemma~\ref{is-fock-facts}(4) it suffices to
check this on $\rho_\sigma(g)h_\sigma$. Let $g'_\tau$ be block diagonal with blocks $g_\tau$ and $1$.
For all $y$ and $z$, by Lemma~\ref{is-fock-facts}(3) and Lemma~\ref{is-branching-vector}(3) with
$t=1$ applied to $g'^{\mathsf T}$,
$\langle\rho_\lambda(g')\mathrm{tr}\,y,z\rangle_\lambda=\langle y,\mathrm{br}\,\rho_\lambda(g'^{\mathsf T})
z\rangle_\sigma=\langle y,\rho_\sigma(g^{\mathsf T})\mathrm{br}\,z\rangle_\sigma=\langle\mathrm{tr}\,
\rho_\sigma(g)y,z\rangle_\lambda$, so $\mathrm{tr}\circ\rho_\sigma(g)=\rho_\lambda(g')\circ\mathrm{tr}$. So
$\mathrm{br}\,\mathrm{tr}\,\rho_\sigma(g)h_\sigma=\mathrm{br}\,\rho_\lambda(g')\,c_0\xi=c_0\rho_\sigma(g)
\,\mathrm{br}\,\xi=c_0c\,\rho_\sigma(g)h_\sigma$, by Steps 1 and 2 and
Lemma~\ref{is-branching-vector}(3).

\emph{Step 4: adjointness for $G$.} Put
$A=\overline{\prod_\tau(t_\tau g_\tau)^{-n_\tau}}\cdot\mathrm{tr}$, with
$n_\tau=n_{\lambda_\tau,\sigma_\tau}$. Write $\bar y$ for the coefficientwise conjugate. Conjugation
commutes with $\mathrm{br}$, hence with $\mathrm{tr}$. By Lemma~\ref{is-fock-facts}(5) and (3), and
Lemma~\ref{is-branching-vector}(3) with the blocks $G'_\tau$ and $g_\tau$,
\[
 \langle Ay,x\rangle_{\lambda,G}
 =\prod_\tau(t_\tau g_\tau)^{-n_\tau}\langle\mathrm{tr}\,\bar y,\rho_\lambda(G)x\rangle_\lambda
 =\prod_\tau(t_\tau g_\tau)^{-n_\tau}g_\tau^{n_\tau}\langle\bar y,\rho_\sigma(G')\mathrm{br}\,x
 \rangle_\sigma=\langle y,\mathrm{br}_tx\rangle_{\sigma,G'}.
\]

\emph{Step 5: the norm.} By Step 3, $\mathrm{br}_tA=k$ with
$k=\prod_\tau t_\tau^{-n_\tau}\cdot\overline{\prod_\tau(t_\tau g_\tau)^{-n_\tau}}\cdot c_0c$. Since
$s(t_\tau)\overline{s(t_\tau)s(g_\tau)}=|s(t_\tau)|^2s(g_\tau)=T$ and $s(\kappa)=\kappa$ over
$\mathbf R$, $s(k)=\prod_\tau T^{-n_\tau}c_\tau^2\kappa_\tau=\beta$. By Step 4,
$s\langle Ay,Ay'\rangle_{\lambda,G}=s\langle y,\mathrm{br}_tAy'\rangle_{\sigma,G'}=\beta\,s\langle
y,y'\rangle_{\sigma,G'}$, which is (1). Lemma~\ref{is-abstract-coisometry}, with $B$ and $B'$ the two
Hermitian pairings, which are Hermitian and positive by Lemma~\ref{is-fock-facts}(5), gives (2).

For (3) take $f=Ah_\sigma$. It is nonzero: otherwise $kh_\sigma=\mathrm{br}_tAh_\sigma=0$, while
$k\ne0$ and $h_\sigma\ne0$. And $\|\mathrm{br}_tf\|^2_{s,G'}=|s(k)|^2\|h_\sigma\|^2_{s,G'}=\beta\cdot\beta
\|h_\sigma\|^2_{s,G'}=\beta\|f\|^2_{s,G}$ by (1).
\end{proof}

\section{Chains of lifts and height numbers}\label{section-is-definite-side-proof}

\begin{definition}\label{is-height-numbers}
Let $D$ be a definite-side datum (Definition~\ref{th-definite-datum}) with ranks $k,m$, $f=2n$ and
$N=N(m)$ (Definition~\ref{th-ranks}), set $\mathcal T$ of $f$ places, Harish-Chandra parameter
$\lambda(\tau)$ with entries $u_i(\tau)$ for $\tau\in\mathcal T$, and coefficient field $L$ with
$r_\infty$ infinite places. Let size data at scale $x$ with exponent $\varepsilon$ be given
(Definition~\ref{pd-size-data}), with the quantity $a$, the quantity $b'$, the width $\eta$, the
cloud size $S_c$, and the set $T$ of item (2) viewed as a subset of $\mathcal T$. Let $u$ be a word
with auxiliary prime $T_u$ and record constants $\lambda_0,C_0,c_e,c,S_w,C_1$
(Definition~\ref{is-definite-side-condition}). Write $\nu_{j,\tau,i}$ for the weights of the $j$-th
lift at $\tau$, $|\nu_{j,\tau}|=\sum_i\nu_{j,\tau,i}$, $a_\tau=u_1(\tau)-\frac{k-1}2$ for
$\tau\in T$, $L_a=\log(a\,r_E)$ with $r_E$ the root discriminant of the host field,
$\phi(y)=y\log y-y$, and $\log\Omega$ as in Definition~\ref{is-division-link}. The \emph{record} of
$u$ consists of these numbers, $\log T_u$ and the constant $c$.
\begin{enumerate}
\item \emph{Kernel-side numbers} with constant $C$ for the record are real numbers $g_0(j)$,
  $M(j,v)$, $\log\Omega_\tau(v)$, $\log\theta_0(v)$, $\log h_j$, $\log h'_j$, $v_0$ and $\ell_g$,
  for lifts $1\le j\le m$, places $\tau\in\mathcal T$ and infinite places $v$ of $L$, such that for
  all $j$, $\tau$ and $v$:
  \begin{enumerate}
  \item $\frac1{r_\infty}\sum_{v}\log\Omega_\tau(v)=\log\Omega$;
  \item $M(j,v)\le Cfm^2N^2(\ell_g+\log T_u)+CfN^2\ell_g$;
  \item $\frac1{r_\infty}\sum_{v}\log\theta_0(v)\le Cx^{\varepsilon}$;
  \item $\log h_j\le CfN^2(x^{\varepsilon}+\log T_u+\ell_g)$ and
    $\log h'_j\le CfN^2(x^{\varepsilon}+\log T_u+\ell_g)$;
  \item $0\le g_0(j)\le N^2$;
  \item $\log(8+32fg_0(j)v_0)\le CL_a$ and $\log(2|\nu_{j,\tau}|+2)\le CL_a$;
  \item $\ell_g\le Cx^{\varepsilon}$.
  \end{enumerate}
  Their \emph{sizes} are
  \begin{multline*}
   \mathrm{sz}(j,v)=\sum_{\tau\in\mathcal T}\Big(\sum_i\phi(\nu_{j,\tau,i})-|\nu_{j,\tau}|\log\Omega_\tau(v)\Big)
   +fg_0(j)\big(\log\theta_0(v)+\log(8+32fg_0(j)v_0)\big)\\
   +2M(j,v)+\sum_{\tau\in\mathcal T}\big(1+2g_0(j)\log(2|\nu_{j,\tau}|+2)\big)+\log h_j+\log h'_j.
  \end{multline*}
\item \emph{Last-side numbers} with constants $C_T$ and $A$ for the record are real numbers $n_\tau$,
  $\lambda_1(\tau)$, $D_0(\tau)$, $D_1(\tau)$, $\log c_\tau$, $\log\xi_\tau$ and $\beta_\tau$, for
  $\tau\in\mathcal T$, such that for all $\tau$:
  \begin{enumerate}
  \item $A\le-\sum_{\tau}n_\tau\log T_u+\sum_{\tau}(2\log c_\tau+\log\xi_\tau)$;
  \item $\log c_\tau\le(N-1)D_1(\tau)\log(2a)$ and $\log\xi_\tau\le\frac{N(N-1)}2\log(2a)$;
  \item $n_\tau=\lambda_1(\tau)-D_0(\tau)$;
  \item $\lambda_1(\tau)=a_\tau-\frac{N-k}2$ if $\tau\in T$, and
    $\lambda_1(\tau)\ge S_c-\frac{N-1}2$ if $\tau\notin T$;
  \item $D_0(\tau)\le S_c+m\eta+\beta_\tau$ and $\beta_\tau\le2b'$ if $\tau\in T$,
    $D_0(\tau)\le S_c+m\eta$ if $\tau\notin T$, and $D_1(\tau)\le3D_0(\tau)+S_c+\eta$;
  \item $0\le\log T_u\le C_Tmx^{\varepsilon}$.
  \end{enumerate}
\item \emph{Height numbers} with period exponent $h_P$ for the record consist of kernel-side numbers
  in which the constant $C$ is replaced by nonnegative constants, $C_2$ and $C_e$ in the two terms
  of (1b), $C_\theta$ in (1c), $C_4$ in (1d), $c_5$ and $c_6$ in the two parts of (1f) and $C_8$ in
  (1g); real numbers $e_v$ and $C_{ev}\ge0$ with $e_v\le C_{ev}fN^2L_a$; and last-side numbers with
  constants $C_T\ge0$ and $h_P-e_v$; such that
  \[
   2(16+C_T+C_{ev})\le c,\qquad C_\theta+c_5+2c_6+1+2C_4+2C_eC_8\le c,\qquad
   2(C_2+C_4)(C_T+C_8)\le c.
  \]
  Their sizes are those of the kernel-side numbers.
\end{enumerate}
\end{definition}
\leanref{Formalism/LocalToGlobal/Divisibility/Inputs.lean}{Formalism.L2.DIV.RecNums}
\leanref{Formalism/LocalToGlobal/Divisibility/Inputs.lean}{Formalism.L2.DIV.HtNums}
\leanref{Development/AbelianVarietyTheory/Definitions/Divisors/FactorBounds.lean}{GlobalTheory.ColRows.KerSide}
\leanref{Development/AbelianVarietyTheory/Definitions/Divisors/FactorBounds.lean}{GlobalTheory.ColRows.LastSide}
\leanref{Formalism/ThetaSupremumBounds/DivIneqExplicit6.lean}{Formalism.DimensionVariation.liftedMapLogSize}

\begin{lemma}\label{is-height-numbers-from-sides}
Let a record with constant $c$ be given (Definition~\ref{is-height-numbers}), and let $C\ge0$,
$C_T\ge0$ and $A$ be real numbers with
\[
 2(16+C_T)\le c,\qquad 6C+1+2C^2\le c,\qquad 4C(C_T+C)\le c.
\]
If the record has kernel-side numbers with constant $C$ and last-side numbers with constants $C_T$
and $A$, then it has height numbers with period exponent $A$ whose sizes are those of the
kernel-side numbers.
\end{lemma}
\leanref{Development/AbelianVarietyTheory/Proofs/DivisorProofs/ExistsHtNumsOfKerSideLastSide.lean}{GlobalTheory.ColRows.exists_htNums_of_kerSide_lastSide}
\begin{proof}
Take $C_2=C_e=C_\theta=C_4=c_5=c_6=C_8=C$, $e_v=0$ and $C_{ev}=0$. The kernel-side and last-side
numbers then satisfy the conditions of Definition~\ref{is-height-numbers}(3) other than the three
inequalities for $c$. These read $2(16+C_T)\le c$, $6C+1+2C^2\le c$ and $2\cdot2C\cdot(C_T+C)\le c$,
which are the hypotheses.
\end{proof}

\begin{lemma}\label{is-chain-from-kernels}
Let $L$ be a totally complex number field with an involution, and write $\sigma_v$ for the
embedding attached to an infinite place $v$. Let $H$ be a Hermitian matrix of size $n_V$ over $L$
with a metric $X^0$, that is, $(X^0_v)^{-1*}(X^0_v)^{-1}=\sigma_v(H)$ at every $v$; let $c_0\ge0$,
let $X_c$ be a metric on $L^{c_0}$ that is invertible at every $v$, and let $\mathrm{Pt}$ be a
bilinear pairing on $L^{c_0}\times L^{n_V}$. Let $F_0,\dots,F_n$ and $S$ be finite-dimensional
$L$-vector spaces with Hermitian forms $\langle\ ,\ \rangle_v$ at the infinite places, each
\emph{metrized}: there are an isomorphism with some $L^{d}$ and a metric $Y$ on $L^d$, invertible
at every $v$, with $\langle y,y\rangle_v=|Y_v^{-1}\sigma_v(y)|^2$. Let $s_i\colon F_i\to F_{i+1}$ for
$0\le i<n$ and $s_n\colon F_n\to S$ be linear maps, put $\operatorname{col}=s_n\circ\cdots\circ s_0$,
and let $f_0\colon L^{c_0}\to F_0$ and $g\colon S\to\operatorname{Hom}_L(L^{n_V},L)$ be linear maps
with
\[
 \langle f_0y,f_0y\rangle_v=|X_{c,v}^{-1}\sigma_v(y)|^2,\qquad
 \mathrm{Pt}(\varphi,F)=g(\operatorname{col}f_0\varphi)(F).
\]
Let a record (Definition~\ref{is-height-numbers}) with set of lifts $J$ have height numbers with
period exponent $h_P$ and sizes $\mathrm{sz}(j,v)$, let $o\colon J\to\mathbf N$ be a map, and put
$t_i(v)=\sum_j\mathrm{sz}(j,v)$, the sum over the lifts $j$ with $\min(o(j),n)=i$. Suppose that at
every infinite place $v$:
\begin{enumerate}
\item $\langle s_iy,s_iy\rangle_v\le e^{t_i(v)}\langle y,y\rangle_v$ for $0\le i\le n$ and
  $y\in F_i$;
\item $|\sigma_v(g(\psi)(F))|^2\le e^{h_P}\langle\psi,\psi\rangle_v\,|X^0_v{}^{-1}\sigma_v(F)|^2$
  for $\psi\in S$ and $F\in L^{n_V}$.
\end{enumerate}
Then there are a chain of lifts with $m'=n+1$ steps, metric $X_0=X_c$, period map $P=g$ and the lift
$j$ assigned to the step $\min(o(j),n)+1$, which factors $\mathrm{Pt}$ as in
Definition~\ref{is-definite-side-condition}(1) and is tied to the maps $s_i$ with first column
$f_0$, and Cauchy-Binet data for this chain, such that at every infinite place $v$ the norm of $P$
from $X_{m'}$ to the dual metric of $X^0$ satisfies $\|P\|_v^2\le e^{h_P}$, and the logarithm of
the square factor at the step $i+1$ and $v$ is at most $t_i(v)$. In particular conditions (1), (3),
(4) and (5) of the rows condition hold for these data.
\end{lemma}
\leanref{Development/AbelianVarietyTheory/Proofs/DivisorProofs/KerChainMatricesWithNormBoundsStatement.lean}{GlobalTheory.T.kerChain_matrices_with_norm_bounds}
\leanref{Development/AbelianVarietyTheory/Proofs/DivisorProofs/KerChainMatricesWithNormBounds.lean}{GlobalTheory.DefRowsChain.rowsBody_of_framed}
\leanref{Development/AbelianVarietyTheory/Definitions/Divisors/FactorBounds.lean}{GlobalTheory.ColRows.RowsBody}
\begin{proof}
Choose the isomorphisms $F_i\cong L^{d_i}$ and $S\cong L^{d_S}$ and the metrics $Y_i$ and
$Y_S$ of the metrizations. Put $c_i=d_i$ and $X_i=Y_i$ for $1\le i\le n$, $c_{n+1}=d_S$ and
$X_{n+1}=Y_S$. Let $K_1$ be the matrix of $s_0\circ f_0$, let $K_i$ be the matrix of $s_{i-1}$ for
$2\le i\le n+1$, and let $P$ be $g$ read in the coordinates of $S$. Then $K_{n+1}\cdots K_1$ is the
matrix of $\operatorname{col}\circ f_0$, so $\mathrm{Pt}(\varphi,F)=P(K_{n+1}\cdots K_1\varphi)(F)$.

The coordinates identify the levels of the chain with $L^{c_0},F_1,\dots,F_n,S$, compatibly with the
maps $s_i$, and they carry the norm of $X_i$ to $\langle\ ,\ \rangle_v$: for $i\ge1$ by the
metrizations, and for $i=0$ through $f_0$ by the hypothesis on $f_0$. This is the tie, with first
column $f_0$.

By hypothesis (1), applied to $f_0y$ when $i=0$, every $y$ satisfies
$|X_{i+1,v}^{-1}\sigma_v(K_{i+1}y)|^2\le e^{t_i(v)}|X_{i,v}^{-1}\sigma_v(y)|^2$. So the operator
$X_{i+1,v}^{-1}\sigma_v(K_{i+1})X_{i,v}$ has norm at most $e^{t_i(v)/2}$. Take Cauchy-Binet data (Definition~\ref{is-cauchy-binet})
with a single class at each level, unit weights and these operators. Their square factor at the
step $i+1$ and $v$ is $e^{t_i(v)}$.

For $y\in L^{c_{n+1}}$ with image $\psi\in S$, hypothesis (2) gives
$|\sigma_v(P(y)(F))|^2\le e^{h_P}|X_{n+1,v}^{-1}\sigma_v(y)|^2\,|X^0_v{}^{-1}\sigma_v(F)|^2$ for all
$F$. A pairing bounded in this way has norm at most $e^{h_P/2}$ from the metric $X_{n+1}$ to the dual
metric of $X^0$, so $\|P\|_v^2\le e^{h_P}$.
\end{proof}

\begin{definition}\label{is-kernel-ties}
Let $D$ be a definite-side datum (Definition~\ref{th-definite-datum}) with ranks $k,m,t$, $f=2n$
and $N(j)$, set $\mathcal T$ of places, Harish-Chandra parameters $\lambda_u(\tau)$ with entries
$u_i,v_i,b_i,e_i$, coefficient field $L$, host field with root discriminant $r_E$, auxiliary primes
$q,d$ and level primes $g_0,g_1$, and words $u$ with auxiliary primes $T_u$. Let size data at scale
$x$ with exponent $\varepsilon$ be given (Definition~\ref{pd-size-data}), with the quantities
$a,b,b',X',w''$, the width $\eta$, the cloud size $S_c$, a size $\eta_{sz}$ and the set $T$ of item
(2), viewed as a subset of $\mathcal T$. \emph{Kernel data} for $D$ consist of number fields and
primes attached to the host, weights $W$ for the forms of every word at every level, Hecke
characters, and for every tuple $\mu$ of characters local constants and, for every word $u$, the
\emph{kernel chain} of $u$ at $\mu$. Its spaces are $F_0,\dots,F_{m-2}$, where $F_j$ consists of
the forms of weight $\lambda_{u,j}$ (Definition~\ref{th-definite-datum}(7)) on the lattice classes
of level $j$, with the inner products
$\langle\varphi,\psi\rangle_v=\sum_c|\Gamma_c|^{-1}\langle\varphi_c,\psi_c\rangle_v$ at the infinite
places $v$ of $L$; its maps are the lifts $s_j\colon F_j\to F_{j+1}$ for $j<m-2$, given by the theta
kernels as in Definition~\ref{th-iterated-lift}, and the coupling $s_{m-2}\colon F_{m-2}\to S_u$,
which composes the last two lifts through a coupling set at level $m-1$ and evaluates at the top,
with values in a space $S_u$ of forms of weight $\lambda_{u,m}$ with its inner products.
\begin{enumerate}
\item The kernel data satisfy the \emph{kernel ties} with weights $W$ if: their ranks, auxiliary
  primes, level primes and host are those of $D$, and their fields and primes satisfy the standing
  conditions of the host; for some identification of $\mathcal T$ with the embeddings of their
  fields, the weights $W$ are well formed and give, for every word $u$, the parameters
  $\lambda_u(\tau)$, with the interior entries read from the letters of $u$; their characters are
  tied to the $e$-parts of the parameters and to the Weil data; their discriminant term, character
  and degeneracy constants are those of the host; the tuple of characters at the base tuple
  $\mu_0$ is trivial and its local constants are sharp; the kernel chains have $m-2$ lifts before
  the coupling; and for every word $u$ and every $i<m$ the theta jets of the $i$-th lift at $\mu_0$
  have coefficients in their field of definition.
\item The \emph{size conditions} with constant $C_Y$ hold for a word $u$ if: $\log r_E$, $\log q$,
  $\log d$, $\log(4x)$ and $|\operatorname{Re}\frac{L'}{L}(1,\chi_d)|$ are at most $x^{\varepsilon}$;
  $(kb)^8\le g_i$ and $\log g_i\le C_Yx^{\varepsilon}+1$ for $i=0,1$; $fN(m)\le b$; the weights of
  the lifts for $u$ lie in $[0,3a]$; and for the parameter $\lambda_u(\tau)$ at every $\tau$,
  $|b_i(\tau)-e_j(\tau)|\le\eta$ for all $i,j$, the entries $v_i(\tau)$ and the interior entries
  $u_i(\tau)$, $1<i<k$, lie in $[0,S_c]$, $b'\le-u_k(\tau)-\frac{k-1}2\le2b'$ if $\tau\in T$, and
  $u_k(\tau)=v_{k-1}(\tau)-\frac12$, $S_c\le u_1(\tau)-v_1(\tau)$ and $u_1(\tau)\le3S_c$ if
  $\tau\notin T$.
\item The kernel data satisfy the \emph{kernel sizes} with constants $(C_Y,C_T)$ if
  $g_i-1>2fk$ for $i=0,1$, $2fN(m)+1<4a$, the weights of every level $j\le m$ have spread at most
  $2a$, and for every word $u$ we have $4a<T_u$, $\log T_u\le C_Tmx^{\varepsilon}$ and the size
  conditions with constant $C_Y$.
\item The \emph{height conditions} with constant $C_Y$ hold for the size data if parameter choices
  exist; $\log r_E$, $|\operatorname{Re}\frac{L'}{L}(1,\chi_d)|$ and $\log(4x)$ are at most
  $x^{\varepsilon}$; $\log g_i\le C_Yx^{\varepsilon}+1$ for $i=0,1$; $4N(m)^3\le\eta_{sz}$,
  $t\le\eta_{sz}$, $m(t+1)\eta_{sz}\le w''$ and $kS_c\le w''$; $-u_k(\tau)-\frac{k-1}2\le2b'$ for
  $\tau\in T$; the number $\Delta_S$ of Definition~\ref{th-definite-datum}(8) is an integer; and the
  weights of the lifts at the thick places of the host are at most $b+\eta_{sz}+1$.
\item A bundle at $D$ (Definition~\ref{is-definite-side-condition}) satisfies the \emph{evaluation
  conditions} if its theta realization is tied at the real place; its columns have symmetric forms
  whose characters are trivial at step zero; its column characters are compatible with its rows;
  the kernel data of its realization satisfy the kernel ties with the weights $W$ of its columns;
  and for all real $C_Y$ and $C_T$ these kernel data satisfy the kernel sizes with $(C_Y,C_T)$
  whenever the height conditions with constant $C_Y$ hold, $X'\ge8$, $x\ge2$, and $4a<T_u$ and
  $\log T_u\le C_Tmx^{\varepsilon}$ for every word $u$.
\end{enumerate}
\end{definition}
\leanref{Development/AbelianVarietyTheory/Definitions/Divisors/KerSizes.lean}{GlobalTheory.ColRows.KerTies}
\leanref{Development/AbelianVarietyTheory/Definitions/Divisors/KerSizes.lean}{GlobalTheory.ColRows.KerSizes}
\leanref{Development/AbelianVarietyTheory/Definitions/LatticeFrames/ColHeightSharpWith.lean}{GlobalTheory.Core.Ht2.SizeRowsS}
\leanref{Development/AbelianVarietyTheory/Definitions/LatticeFrames/RecHeightSWith.lean}{GlobalTheory.CoreHt.HtRowsS}
\leanref{Formalism/LocalToGlobal/Divisibility/DivisionHeightBoundStatement.lean}{Formalism.L2.DIV.KerChain}
\leanref{Development/AutomorphicArithmetic/ThetaCorrespondence/ThetaPairLifts/ThirdPlaceSeries/PointEvaluationSide.lean}{GlobalTheory.Fc.RL.DefSide.PointEvaluationSide}

\section{Kernel data and the definite-side condition}\label{section-is-kernel-data}

\subsection{Jet operators and the block bound}\label{subsection-is-jet-operators}

\begin{definition}\label{is-jet-operator}
Let $\iota$ be a finite set and $a\le b$ integers, and write $\varphi(y)=y\log y-y$.
\begin{enumerate}
\item The \emph{Fock space} of the pair is $V_{a,b}=M_{b\times a}(\mathbf C)^{\iota}$, with coordinates
  $z_{\tau,\beta,\alpha}$. For $r\in\mathbf R_{>0}^{\iota}$ put
  $|z|_r^2=\sum_{\tau,\beta,\alpha}|r_\tau z_{\tau,\beta,\alpha}|^2$.
\item For an entire function $G$ on $V_{a,b}$, $c\in\mathbf C^{\iota}$ and an exponent vector $m$ on
  the coordinates, the \emph{jet} of $G$ is $j_c(G)(m)=\prod_\tau c_\tau^{|m_\tau|}\,\partial^mG(0)$,
  where $|m_\tau|$ is the degree of $m$ in the coordinates of $\tau$.
\item Let $K$ be a field. For a function $J$ from exponent vectors to $K$ and $n\in\mathbf N^{\iota}$,
  the \emph{jet map} $J_{J,n}\colon\mathrm{Pol}_a^{\iota}\to\mathrm{Pol}_b^{\iota}$ multiplies $P$ by
  $\prod_\tau\det(x_\tau)^{n_\tau}$, substitutes $x_{\tau,\alpha,s}=\sum_{\beta\le b}y_{\tau,\beta,s}
  u_{\tau,\beta,\alpha}$ with new variables $y$ (an $b\times b$ matrix for each $\tau$) and $u$, and
  replaces each monomial $u^m$ by $J(m)$. For dominant weights $\lambda$ of length $a$ and $\lambda'$ of
  length $b$ (Definition~\ref{is-polynomial-model}), the \emph{jet operator}
  $\mathrm{Jet}(J,n)\colon V(\lambda)\to V(\lambda')$ is $J_{J,n}$ followed by a projection of
  $\mathrm{Pol}_b^{\iota}$ onto $V(\lambda')$.
\item Dominant weights $\lambda,\nu$ of length $a$ and $\lambda'$ of length $b$ are \emph{step
  weights} if $\nu_i\geq0$, $\nu-\lambda$ is constant, and $\lambda'$ plus a constant is $\nu$
  followed by $b-a$ zeros. The \emph{last entry} of $\nu$ is $\ell(\nu)=\max(\nu_a,0)$ and its size is
  $|\nu|=\sum_i\max(\nu_i,0)$.
\item For $S\geq0$, $\Omega\in\mathbf R_{>0}^{\iota}$ and weights $\nu_\tau$ of length $a$,
  \[
   t(b,\nu,S,\Omega)=S^2\prod_\tau\Big(\Omega_\tau^{-|\nu_\tau|}\,e\,(2|\nu_\tau|+2)^{2ab}
   \exp\Big(\sum_i\varphi(\nu_{\tau,i})\Big)\Big).
  \]
\end{enumerate}
\end{definition}
\leanref{Formalism/LocalToGlobal/KernelLattices/ThDistC.lean}{Formalism.L2.Ker.jetC}
\leanref{Formalism/LocalToGlobal/KernelLattices/ThDistC.lean}{Formalism.L2.Ker.jetMap}
\leanref{Formalism/LocalToGlobal/KernelLattices/ThDistC.lean}{Formalism.L2.Ker.jetOp}
\leanref{Development/AbelianVarietyTheory/BlockTheory/DivisorBase/PaddedLastEntry.lean}{GlobalTheory.IndefiniteRowSideFullHecke.PaddedLastEntry.StepWeights}
\leanref{Development/AbelianVarietyTheory/BlockTheory/DivisorBase/PaddedLastEntry.lean}{GlobalTheory.IndefiniteRowSideFullHecke.PaddedLastEntry.unitSq}
\leanref{Development/AbelianVarietyTheory/BlockTheory/DivisorBase/PaddedLastEntry.lean}{GlobalTheory.IndefiniteRowSideFullHecke.PaddedLastEntry.tmax}

\begin{lemma}\label{is-fock-substitution}
Let $a\le b$ and $A\in M_{a\times b}(\mathbf C)$. There are $t_1,\dots,t_a\geq0$ with
$\sum_it_i^2=\sum_{\alpha,\beta}|A_{\alpha\beta}|^2$ such that for every non-increasing
$\nu\in\mathbf N^a$ and every $f\in\nabla(\nu)$ over $\mathbf C$, the polynomial
$f_A(x)=f(Ax')$ on $b\times b$ matrices $x$, where $x'$ is the $b\times a$ matrix of the first $a$
columns of $x$, satisfies $\{f_A,f_A\}_F\le\prod_it_i^{2\nu_i}\{f,f\}_F$.
\end{lemma}
\leanref{Development/AbelianVarietyTheory/Proofs/DivisorProofs/FockSubstLe.lean}{GlobalTheory.JetBlock.fock_subst_le}
\begin{proof}
For a complex matrix $M$ with $a$ rows let $S_M$ be the substitution $F(x)\mapsto F(Mx)$; for an
$a\times a$ matrix it acts on polynomials in the entries of an $a\times a$ matrix, and
$f_A=S_A(f)$, with $f$ read as a polynomial in the first $a$ columns of a $b\times b$ matrix. The
Hermitian Fock form satisfies $\{S_MF,S_MF\}_F=\{F,S_{MM^*}F\}_F$: the adjoint of $S_M$ is $S_{M^*}$,
as in the proof of Lemma~\ref{is-fock-facts}(5), and $S_{M^*}S_M=S_{MM^*}$.

\emph{Step 1: a square root of $AA^*$.} The matrix $AA^*$ is positive semidefinite. By the spectral
theorem there are a unitary $U\in M_a(\mathbf C)$ and reals $t_1\geq\cdots\geq t_a\geq0$, the square
roots of its eigenvalues, with $AA^*=MM^*$ for $M=U\operatorname{diag}(t)$. Then
$\sum_it_i^2=\operatorname{tr}(AA^*)=\sum_{\alpha,\beta}|A_{\alpha\beta}|^2$.

\emph{Step 2: reduction to a diagonal matrix.} By the identity above,
$\{f_A,f_A\}_F=\{f,S_{AA^*}f\}_F=\{f,S_{MM^*}f\}_F=\{S_Mf,S_Mf\}_F$. Put $g=S_Uf$. The substitution by
an invertible matrix acts on the rows and preserves $\nabla(\nu)$, so $g\in\nabla(\nu)$, and
$\{g,g\}_F=\{f,S_{UU^*}f\}_F=\{f,f\}_F$. Moreover $S_Mf=S_{\operatorname{diag}(t)}g$, the polynomial
$g(\operatorname{diag}(t)x)$.

\emph{Step 3: dominance.} Let $r_\alpha(m)$ be the degree of a monomial $x^m$ in the $\alpha$-th row.
Every monomial of $g$ satisfies $\sum_{\alpha\le k}r_\alpha(m)\le\sum_{t\le k}\nu_t$ for all $k$, with
equality for $k=a$. Indeed, by Lemma~\ref{is-standard-basis}(1), $g$ is a combination of products of
minors on the first $s_1,s_2,\dots$ columns with $\#\{l:s_l\geq t\}=\nu_t$; a minor of size $s_l$ has
degree at most $\min(s_l,k)$ in the first $k$ rows, and $\sum_l\min(s_l,k)=\sum_{t\le k}\nu_t$; for
$k=a$ the total degree is $\sum_t\nu_t$.

\emph{Step 4: the diagonal bound.} The substitution by $\operatorname{diag}(t)$ multiplies $x^m$ by
$\prod_\alpha t_\alpha^{r_\alpha(m)}$. For $m$ as in Step 3, $\prod_\alpha t_\alpha^{r_\alpha(m)}\le
\prod_\alpha t_\alpha^{\nu_\alpha}$. If all $t_\alpha>0$ this follows by summation by parts from the
dominance, since $\log t_\alpha$ is non-increasing. Otherwise let $\alpha_0$ be the number of $\alpha$ with $t_\alpha>0$. If
$\nu_\alpha>0$ for some $\alpha>\alpha_0$, then $\sum_{\alpha\le\alpha_0}r_\alpha(m)\le
\sum_{\alpha\le\alpha_0}\nu_\alpha<\sum_\alpha r_\alpha(m)$, so $r_\alpha(m)>0$ for some
$\alpha>\alpha_0$ and both sides are $0$. Otherwise $\nu_\alpha=0$ for $\alpha>\alpha_0$, the same inequality gives
$r_\alpha(m)=0$ for $\alpha>\alpha_0$, and the first case applies to the first $\alpha_0$ rows. The
Hermitian Fock form is diagonal in the monomials, so
\[
 \{S_Mf,S_Mf\}_F=\sum_m|g_m|^2m!\prod_\alpha t_\alpha^{2r_\alpha(m)}\le\prod_\alpha t_\alpha^{2\nu_\alpha}
 \sum_m|g_m|^2m!=\prod_\alpha t_\alpha^{2\nu_\alpha}\{f,f\}_F .
\]
With Step 2 this is the claim. The numbers $t_i$ depend only on $A$.
\end{proof}

\begin{lemma}\label{is-gaussian-moments}
Let $\kappa$ be a finite set with $q$ elements, write $|u|^2=\sum_{x\in\kappa}|u_x|^2$ for
$u\in\mathbf C^\kappa$, and let $du$ be Lebesgue measure on $\mathbf C^\kappa=\mathbf R^{2q}$.
Write $\varphi(y)=y\log y-y$, with $\varphi(0)=0$.
\begin{enumerate}
\item Let $F=\sum_ef_eu^e$ be an entire function on $\mathbf C^\kappa$ with $|F(u)|\le Se^{|u|^2/2}$
  for all $u$. Then for every exponent vector $e$ the function $F(u)\bar u^ee^{-|u|^2}$ is integrable
  and $\int F(u)\bar u^ee^{-|u|^2}du=\pi^q\,e!\,f_e$.
\item Let $\nu_1,\dots,\nu_r\geq0$ be real, $|\nu|=\sum_i\nu_i$, and let $s_1,\dots,s_r$ be
  nonnegative measurable functions on $\mathbf C^\kappa$ with $\sum_is_i(u)^2=|u|^2$. Then
  \[
   \frac1{\pi^q}\int\prod_is_i(u)^{\nu_i}e^{-|u|^2/2}du\le e^{1/2}(2|\nu|+2)^q
   \exp\Big(\frac12\sum_i\varphi(\nu_i)\Big).
  \]
\end{enumerate}
\end{lemma}
\leanref{Formalism/ThetaSupremumBounds/GaussianShiftBounds/FockJet.lean}{Formalism.DimensionVariation.gaussianCoeffIntegral}
\leanref{Formalism/ThetaSupremumBounds/GaussianShiftBounds/GaussCoord.lean}{Formalism.DimensionVariation.blockIntegralBound}
\begin{proof}
(1) The integrand is bounded by $S|u^e|e^{-|u|^2/2}$, which is integrable. Let $D_R$ be the polydisc
$|u_x|\le R$. On $D_R$ the Taylor series of $F$ converges uniformly, and in polar coordinates in each
variable $\int_{D_R}u^{e'}\bar u^ee^{-|u|^2}du=0$ for $e'\ne e$. So
$\int_{D_R}F(u)\bar u^ee^{-|u|^2}du=f_e\int_{D_R}|u^e|^2e^{-|u|^2}du$. As $R\to\infty$ the left side
tends to the integral over $\mathbf C^\kappa$ by dominated convergence, and the right side to
$f_e\prod_x\int_{\mathbf C}|z|^{2e_x}e^{-|z|^2}dz=f_e\prod_x\pi\,e_x!$. This is the orthogonality of
the monomials in the Fock space of Bargmann \cite{Bar61}.

(2) Put $\epsilon=1/(|\nu|+1)$. For $y\geq0$ and $s\geq0$,
$s^ye^{-(1-\epsilon)s^2/2}\le(1-\epsilon)^{-y/2}e^{\varphi(y)/2}$: for $y>0$ the left side is largest
at $s^2=y/(1-\epsilon)$, where it equals the right side, and for $y=0$ both sides are at most $1$.
Multiplying over $i$ and using $\sum_is_i^2=|u|^2$,
\[
 \prod_is_i(u)^{\nu_i}e^{-|u|^2/2}\le(1-\epsilon)^{-|\nu|/2}e^{\frac12\sum_i\varphi(\nu_i)}
 e^{-\epsilon|u|^2/2}\le e^{1/2}e^{\frac12\sum_i\varphi(\nu_i)}e^{-\epsilon|u|^2/2},
\]
since $(1-\epsilon)^{-|\nu|/2}$ is $1$ if $|\nu|=0$ and $(1+1/|\nu|)^{|\nu|/2}\le e^{1/2}$
otherwise. Finally
$\int e^{-\epsilon|u|^2/2}du=(2\pi/\epsilon)^q=\pi^q(2|\nu|+2)^q$.
\end{proof}

\begin{lemma}\label{is-fock-cauchy}
Let $\iota$ be a finite set, $a\le b$, and for each $\tau$ let $\lambda_\tau,\nu_\tau,\lambda'_\tau$ be
step weights (Definition~\ref{is-jet-operator}); let $P\in V(\lambda)$ over $\mathbf C$. Let
$r,\Omega\in\mathbf R_{>0}^{\iota}$ and $c\in\mathbf C^{\iota}$ with $|c_\tau r_\tau|^2=1/\Omega_\tau$, and
let $G$ be an entire function on $V_{a,b}$ with $|G(z)|\le S\exp(|z|_r^2/2)$. Let $\phi$ be a linear
functional on $\mathrm{Pol}_b^{\iota}$ such that for every $A\in M_{a\times b}(\mathbf C)^{\iota}$,
\[
 \Big|\phi\Big(\big(P\textstyle\prod_\tau\det(x_\tau)^{\ell(\nu_\tau)}\big)_A\Big)\Big|
 \le\prod_{\tau,i}t_i(A_\tau)^{\nu_{\tau,i}},
\]
with the substitution and the $t_i$ of Lemma~\ref{is-fock-substitution}, taken in each $\tau$. Then
$|\phi(J_{j_c(G),\ell(\nu)}P)|^2\le t(b,\nu,S,\Omega)$.
\end{lemma}
\leanref{Formalism/ThetaSupremumBounds/LiftOperatorBoundOfFockJet.lean}{Formalism.DimensionVariation.blockBound}
\leanref{Development/AbelianVarietyTheory/Proofs/DivisorProofs/NormFockLSqLe.lean}{GlobalTheory.JetBlock.scalar_block}
\begin{proof}
Write $n_\tau=\ell(\nu_\tau)$ and $Q=P\prod_\tau\det(x_\tau)^{n_\tau}$, and let $q=|\iota|ab$ be the
number of coordinates $u_{\tau,\beta,\alpha}$ ($\beta\le b$, $\alpha\le a$). The jet map substitutes
$x_\tau=u_\tau^{\mathsf T}y'_\tau$ in $Q$, where $u_\tau$ is the $b\times a$ matrix of the
$u_{\tau,\beta,\alpha}$ and $y'_\tau$ the first $a$ columns of $y_\tau$. Write
$Q((u_\tau^{\mathsf T}y'_\tau)_\tau)=\sum_eu^eA_e$ with $A_e\in\mathrm{Pol}_b^{\iota}$; then
$J_{J,n}P=\sum_eJ(e)A_e$ for every $J$.

\emph{Degrees.} Since $P\in V(\lambda)$, $Q$ has degree $\sum_i(\lambda^{\natural}_{\tau,i}+n_\tau)$ in
$x_\tau$. As $\nu_\tau-\lambda_\tau$ is constant and $\nu_{\tau,i}\geq0$,
$\lambda^{\natural}_{\tau,i}+n_\tau=\lambda_{\tau,i}-\lambda_{\tau,a}+\nu_{\tau,a}=\nu_{\tau,i}$, so
this degree is $|\nu_\tau|$. Hence $\sum_{\beta,\alpha}e_{\tau,\beta,\alpha}=|\nu_\tau|$ whenever
$A_e\ne0$.

\emph{The jets as Gaussian integrals.} Let $F(u)=G((u_\tau/r_\tau)_\tau)$. It is entire, and
$|F(u)|\le Se^{|u|^2/2}$ with $|u|^2=\sum|u_{\tau,\beta,\alpha}|^2$, since $|z|_r^2=\sum_\tau|r_\tau
z_\tau|^2$. If $G=\sum_eg_ez^e$ and $F=\sum_ef_eu^e$, then $f_e=g_e\prod_\tau r_\tau^{-|e_\tau|}$, so
$j_c(G)(e)=\prod_\tau c_\tau^{|e_\tau|}\,e!\,g_e=e!\,f_e\prod_\tau(c_\tau r_\tau)^{|e_\tau|}$. For $e$
with $A_e\ne0$ this is $\gamma\,e!\,f_e$ with $\gamma=\prod_\tau(c_\tau r_\tau)^{|\nu_\tau|}$, and
$|\gamma|^2=\prod_\tau\Omega_\tau^{-|\nu_\tau|}$. By Lemma~\ref{is-gaussian-moments}(1),
\[
 \phi(J_{j_c(G),n}P)=\gamma\sum_ee!\,f_e\,\phi(A_e)=\frac{\gamma}{\pi^q}\int F(u)e^{-|u|^2}
 \sum_e\bar u^e\phi(A_e)\,du=\frac{\gamma}{\pi^q}\int F(u)e^{-|u|^2}\phi(Q_{A(u)})\,du,
\]
where $A(u)=(\bar u_\tau^{\mathsf T})_\tau\in M_{a\times b}(\mathbf C)^{\iota}$, because
$\sum_e\bar u^eA_e=Q((\bar u_\tau^{\mathsf T}y'_\tau)_\tau)=Q_{A(u)}$.

\emph{The bound.} By hypothesis $|\phi(Q_{A(u)})|\le\prod_{\tau,i}t_i(\bar u_\tau^{\mathsf T})^{\nu_{\tau,i}}$,
and $\sum_it_i(\bar u_\tau^{\mathsf T})^2=|u_\tau|^2$ by Lemma~\ref{is-fock-substitution}. The $t_i$ are
the square roots of the eigenvalues of $A(u)_\tau A(u)_\tau^*$, so they are continuous in
$u$. With $|F(u)|e^{-|u|^2}\le S\prod_\tau e^{-|u_\tau|^2/2}$ and Fubini,
\[
 |\phi(J_{j_c(G),n}P)|\le|\gamma|\,S\prod_\tau\frac1{\pi^{ab}}\int\prod_it_i(\bar u_\tau^{\mathsf T})^{\nu_{\tau,i}}
 e^{-|u_\tau|^2/2}\,du_\tau\le|\gamma|\,S\prod_\tau e^{1/2}(2|\nu_\tau|+2)^{ab}
 e^{\frac12\sum_i\varphi(\nu_{\tau,i})},
\]
by Lemma~\ref{is-gaussian-moments}(2) for each $\tau$, with $s_i(u_\tau)=t_i(\bar u_\tau^{\mathsf T})$.
Squaring gives $|\phi(J_{j_c(G),n}P)|^2\le S^2\prod_\tau\Omega_\tau^{-|\nu_\tau|}e(2|\nu_\tau|+2)^{2ab}
\exp(\sum_i\varphi(\nu_{\tau,i}))=t(b,\nu,S,\Omega)$.
\end{proof}

\begin{proposition}\label{is-block-bound}
Let $\iota$ be a finite set, $a\le b$, and for each $\tau$ let $\lambda_\tau,\nu_\tau,\lambda'_\tau$
be step weights (Definition~\ref{is-jet-operator}), with $\ell(\nu_\tau)=0$ if $a=b$. Let
$r,\Omega\in\mathbf R_{>0}^{\iota}$ and $c\in\mathbf C^{\iota}$ with $|c_\tau r_\tau|^2=1/\Omega_\tau$, and
let $G$ be an entire function on $V_{a,b}$ with $|G(z)|\le S\exp(|z|_r^2/2)$. Then for every
$P\in V(\lambda)$ over $\mathbf C$,
\[
 \|\mathrm{Jet}(j_c(G),\ell(\nu))P\|^2\le t(b,\nu,S,\Omega)\,\|P\|^2,
\]
the norms being those of Definition~\ref{is-polynomial-model}(7) with $G_\tau=1$.
\end{proposition}
\leanref{Development/AbelianVarietyTheory/BlockTheory/DivisorBase/PaddedLastEntry.lean}{GlobalTheory.IndefiniteRowSideFullHecke.PaddedLastEntry.BlockBound}
\leanref{Development/AbelianVarietyTheory/Proofs/DivisorProofs/NormFockLSqLe.lean}{GlobalTheory.JetBlock.blockBound}
\begin{proof}
Write $n=\ell(\nu)$ and, for $A\in M_{a\times b}(\mathbf C)^{\iota}$, let $T_A$ multiply a polynomial by
$\prod_\tau\det(x_\tau)^{n_\tau}$ and substitute $A_\tau$ as in Lemma~\ref{is-fock-substitution}.

\emph{Step 1.} The jet map sends $V(\lambda)$ into $V(\lambda')$: for each $A$, $T_A$ sends
$\otimes_\tau f_\tau$ with $f_\tau\in\nabla(\lambda_\tau)$ to $\otimes_\tau(f_\tau\det^{n_\tau})_{A_\tau}$,
and $f\det^{n}$ lies in $\nabla$ of the natural weight of $\lambda$ plus $n$, whose padding by zeros is
the natural weight of $\lambda'$ because the weights are step weights and $a<b$ or $n=0$; the jet map
is a combination of the $T_A$. So the projection in the jet operator does nothing, and
$w=J_{j_c(G),n}P$ lies in $V(\lambda')$.

\emph{Step 2.} For $f$ in $V(\mu)$, $\|f\|^2=\{f,f\}_F/\prod_\tau e(\mu_\tau)$ with
$e(\mu_\tau)=\{h_{\mu_\tau},h_{\mu_\tau}\}_F>0$. Suppose $w\ne0$; otherwise there is nothing to prove.
Then $P\ne0$, and $N_w=\{w,w\}_F$ and $N_P=\{P,P\}_F$ are positive.

\emph{Step 3.} By Lemma~\ref{is-fock-substitution} and the multiplication by $\det^{n}$, which multiplies
$\{f,f\}_F$ by $e(\lambda')/e(\lambda)$ on $\nabla(\lambda)$, each tensor factor satisfies
$\{(f\det^n)_A,(f\det^n)_A\}_F\le\prod_it_i(A)^{2\nu_i}\frac{e(\lambda')}{e(\lambda)}\{f,f\}_F$. Since
the Fock Gram matrices of the factors are positive semidefinite, their Schur products are too, and
the bound passes to $V(\lambda)$:
$\{T_AP,T_AP\}_F\le\prod_{\tau,i}t_i(A_\tau)^{2\nu_{\tau,i}}\prod_\tau\frac{e(\lambda'_\tau)}{e(\lambda_\tau)}N_P$.

\emph{Step 4.} Let $\delta=N_w\prod_\tau\frac{e(\lambda'_\tau)}{e(\lambda_\tau)}N_P$ and
$\phi=\delta^{-1/2}\{w,\cdot\}_F$. By the Cauchy-Schwarz inequality and Step 3,
$|\phi(T_AP)|\le\prod_{\tau,i}t_i(A_\tau)^{\nu_{\tau,i}}$ for every $A$. Lemma~\ref{is-fock-cauchy}
gives $|\phi(w)|^2=N_w^2/\delta\le t(b,\nu,S,\Omega)$, that is,
$N_w/\prod_\tau e(\lambda'_\tau)\le t(b,\nu,S,\Omega)\,N_P/\prod_\tau e(\lambda_\tau)$, which is the
claim by Step 2.
\end{proof}

\subsection{Theta kernels and the growth of theta functions}\label{subsection-is-theta-growth}

\begin{definition}\label{is-vacuum-theta-vectors}
Let $d\equiv7\pmod8$ be a prime, $K=\mathbf Q(\sqrt{-d})$ and $\omega_K=(1+i\sqrt d)/2$. Let $E$ be a
CM field of degree $2f$ containing $K$, with complex conjugation $x\mapsto\bar x$, let $\mathcal T$ be
a set of $f$ embeddings of $E$ into $\mathbf C$ that agree on $K$, and let $\mathfrak A_0$ be a fractional
ideal of $E$. Let $\Lambda\subset\mathbf C$ be a lattice with complex multiplication by $\mathcal O_K$
(Definition~\ref{pd-cm-lattice}), of covolume $A$, and let $a,b\geq0$ be integers.
\begin{enumerate}
\item On $V_{a,b}=M_{b\times a}(\mathbf C)^{\mathcal T}$ (Definition~\ref{is-jet-operator}) put
  $\langle z,w\rangle=\sum_{\tau,\beta,\alpha}z_{\tau,\beta,\alpha}\overline{w_{\tau,\beta,\alpha}}$ and
  $H=A^{-1}\langle\ ,\ \rangle$, and let $U_x$ and $\alpha_{\sqrt d/2}$ be the operators and the
  normalization of Definition~\ref{pd-theta-function} for $H$. For $B\in M_{b\times a}(E)$ put
  $\iota(B)=(\tau(B))_{\tau\in\mathcal T}\in V_{a,b}$.
\item For a subgroup $\Gamma\subset V_{a,b}$, a \emph{theta function} for $\Gamma$ is a theta function
  for $(H,\Gamma)$ with normalization $\sqrt d/2$ (Definition~\ref{pd-theta-function}), and the
  \emph{norm} $\|G\|^2_\Gamma$ of a function $G$ on $V_{a,b}$ is the mean of $|G(z)|^2e^{-\pi H(z,z)}$
  over a fundamental domain of $\Gamma$. For subgroups $\Gamma,\Gamma'$ with $\Gamma'/(\Gamma\cap\Gamma')$
  finite, the \emph{averaging operator} is
  \[
   I_{\Gamma\to\Gamma'}G=[\Gamma':\Gamma\cap\Gamma']^{-1/2}\sum_x\alpha_{\sqrt d/2}(x)\,U_xG,
  \]
  the sum over representatives $x$ of $\Gamma'/(\Gamma\cap\Gamma')$.
\item For $\mathcal O_E$-lattices $L\subset E^a$ and $L'\subset E^b$, $\Lambda_{L,L'}\subset V_{a,b}$ is the
  subgroup generated by the $w\,\iota(sB)$ with $w\in\Lambda$, $s\in\mathfrak A_0$ and
  $B\in M_{b\times a}(E)$ with $BL\subset L'$. The pair $(L,L')$ is \emph{admitted} if $a$ and $b$ are
  even and $L$ and $L'$ are self-dual for the form $\sum_i\bar x_iy_i$.
\item An \emph{orthonormal basis} is a basis $v_1,\dots,v_r$ of $M_{b\times a}(E)$ over $K$, $r=fab$,
  with $\operatorname{Tr}_{E/K}\sum_{\beta,\alpha}(v_k)_{\beta\alpha}\overline{(v_l)_{\beta\alpha}}=\delta_{kl}$.
  Its lattice is $\Lambda_v=\sum_k\Lambda\,\iota(v_k)$, and the \emph{tensor power} along $v$ of a
  function $g$ on $\mathbf C$ is $g_v(z)=\prod_kg(\langle z,\iota(v_k)\rangle)$.
\item For $t\in\mathbf C\smallsetminus\Lambda$, $n(\Lambda,t)$ is the mean of
  $|\vartheta_{\Lambda,t}(z)|^2e^{-\pi|z|^2/A}$ over a fundamental domain of $\Lambda$, where
  $\vartheta_{\Lambda,t}$ is the normalized theta function of $\Lambda$ translated by $t$
  (Definition~\ref{pd-theta-function}). The point $t$ is an \emph{unkilled two-torsion point} of
  $\Lambda$ if $2t\in\Lambda$, $\omega_Kt\notin\Lambda$ and $\bar\omega_Kt\notin\Lambda$.
\item Let $t$ be an unkilled two-torsion point of $\Lambda$, and $L\subset E^a$ and $L'\subset E^b$
  $\mathcal O_E$-lattices. A function $G$ on $V_{a,b}$ is a \emph{theta vector up to units} for
  $(t,L,L')$ if
  $G=\kappa\,I_{\Lambda_v\to\Lambda_{L,L'}}\big((\vartheta_{\Lambda,t})_v\big)$ for an orthonormal basis
  $v$ and an algebraic number $\kappa$ all of whose conjugates have absolute value $1$.
\end{enumerate}
For the host of a definite-side datum, $d\equiv7\pmod{16}$ is the prime of its host field
(Definition~\ref{hf-framed-host-field}), $E$ is the host field with its embedding $\varphi_0$,
$\mathcal T$ is its set of types, which agree with $\varphi_0$ on $K$, and $\mathfrak A_0=\mathfrak b^{-1}$
for the ideal $\mathfrak b$ of Definition~\ref{hf-host}(c). A \emph{CM lattice} of the host is
$\Lambda=\delta\varphi_0(\mathfrak c)$ for a nonzero ideal $\mathfrak c$ of $\mathcal O_K$ of norm prime
to $2d$ and a $\delta\in\mathbf C$ with $\delta^{12}=\Delta(\varphi_0(\mathfrak c))$, where
$\Delta=g_2^3-27g_3^2$; then $\Delta(\Lambda)=1$. Let $M\subset\mathbf C$ be a number field containing
$\tau(E)$ for $\tau\in\mathcal T$, $g_2(\Lambda)$ and $g_3(\Lambda)$, and let $\rho$ be an embedding of
$M$. The \emph{conjugate} $\Lambda^\rho$ is the lattice with $g_k(\Lambda^\rho)=\rho(g_k(\Lambda))$ for
$k=2,3$. The objects of (1) to (6) formed with $\Lambda^\rho$ and the types $\rho\circ\tau$ in place of
$\Lambda$ and $\tau$ carry a superscript $\rho$. For a function $G$ on $V_{a,b}$ whose Taylor
coefficients at $0$ lie in $M$, $G^\rho$ is the function given by the power series whose coefficients
are their images under $\rho$.

The kernel data of the datum give a CM lattice $\Lambda$ of the host, an unkilled two-torsion point
$\omega_0/2$ of $\Lambda$, a $K$-linear isometry $\gamma_0$ from $E^4$ with the form
$\operatorname{Tr}_{E/K}\sum_ax_a\bar y_a$ to $K^{4\times\mathcal T}$ with its standard form, a number
$\kappa_0$, and a function $\chi$ on the fractional ideals of $E$ that is multiplicative on the ideals
$\mathfrak a$ with $\mathfrak a\bar{\mathfrak a}=\mathcal O_E$ and satisfies
$\chi(z\mathcal O_E)=\prod_\tau\tau(z)^{-1}$ whenever $z\bar z=1$. The \emph{vacuum vector of the host}
is a chosen theta function $\theta_0$ for $\Lambda_{\mathcal O_E,\mathcal O_E}\subset V_{1,1}=\mathbf
C^{\mathcal T}$ for which $\theta_0(z_1)\theta_0(z_2)\theta_0(z_3)\theta_0(z_4)$ is given by identity (4)
of Definition~\ref{pd-cm-datum}, read with $n=\mathcal T$, the form $H$ of (1), the subgroup
$\Lambda_{\mathcal O_E,\mathcal O_E}$ and the lattice $\Lambda$ in place of $\Lambda$ and $L$ there,
$\omega_0/2$ in place of $\lambda a_0/2$, $\varepsilon=\kappa_0$, and $g$ the complex-linear map with
$g(\iota(x))=\varphi_0(\gamma_0(x))$ for $x\in E^4$. For an admitted pair $(L,L')$ of ranks $a,b$, let
$\mathfrak d_L$ be the ideal generated by the determinants of the $B\in M_a(E)$ with
$B\mathcal O_E^a\subset L$. The \emph{vacuum vector of the pair} is
\[
 \theta_{L,L'}=\chi(\mathfrak d_L)^{b/2}\chi(\mathfrak d_{L'})^{-a/2}\,
 I_{\Lambda_{L,\mathcal O_E^b}\to\Lambda_{L,L'}}I_{\Lambda_{\mathcal O_E^a,\mathcal O_E^b}\to\Lambda_{L,\mathcal O_E^b}}
 \theta_0^{\otimes ab},\qquad
 \theta_0^{\otimes ab}(z)=\prod_{\beta,\alpha}\theta_0\big((z_{\tau,\beta,\alpha})_\tau\big).
\]
\end{definition}
\leanref{Formalism/SharpRestriction/GlobalSchurBounds/GeneralLinearActions/IsDiscriminantTwelfthRootOfIdealLattice.lean}{Formalism.MainNotions.TaylorLogMetric.La}
\leanref{Formalism/SharpRestriction/GlobalSchurBounds/GeneralLinearActions/IsDiscriminantTwelfthRootOfIdealLattice.lean}{Formalism.MainNotions.TaylorLogMetric.normSq}
\leanref{Formalism/SharpRestriction/GlobalSchurBounds/GeneralLinearActions/IsDiscriminantTwelfthRootOfIdealLattice.lean}{Formalism.MainNotions.TaylorLogMetric.Iop}
\leanref{Formalism/SharpRestriction/GlobalSchurBounds/GeneralLinearActions/IsDiscriminantTwelfthRootOfIdealLattice.lean}{Formalism.MainNotions.TaylorLogMetric.lamConj}
\leanref{Formalism/SharpRestriction/GlobalSchurBounds/GeneralLinearActions/IsDiscriminantTwelfthRootOfIdealLattice.lean}{Formalism.MainNotions.TaylorLogMetric.fnConj}
\leanref{Formalism/SharpRestriction/GlobalSchurBounds/StandardLatticesGL3.lean}{Formalism.MainNotions.StandardLatticesGL3.pullB}
\leanref{Formalism/SharpRestriction/GlobalSchurBounds/StandardLatticesGL3.lean}{Formalism.MainNotions.StandardLatticesGL3.nrm1Sq}
\leanref{Formalism/SharpRestriction/GlobalSchurBounds/StandardLatticesGL3.lean}{Formalism.MainNotions.StandardLatticesGL3.ThetaVectorUpToUnits}
\leanref{Formalism/SharpRestriction/GlobalSchurBounds/GL3Theta.lean}{Formalism.MainNotions.GL3Theta.IsTwoTorsionPointUnkilled}
\leanref{Formalism/SharpRestriction/GlobalSchurBounds/GL3Theta.lean}{Formalism.MainNotions.GL3Theta.th0}
\leanref{Formalism/LocalToGlobal/KernelLattices/LateralTheory.lean}{Formalism.L2.Ker.thM}
\leanref{Formalism/LocalToGlobal/KernelLattices/LateralTheory.lean}{Formalism.L2.Ker.Base.eNat}

\begin{definition}\label{is-theta-kernel}
Let $D$ be a definite-side datum with kernel data (Definition~\ref{is-kernel-ties}), host field $E$
with set $\mathcal T$ of types, coefficient field $L$ with the embedding $\sigma_0$, level primes
$g_0,g_1$ and words $u$ with auxiliary primes $T_u$. Let $a=N(i)\le b=N(i+1)$ be the ranks of two
consecutive levels.
\begin{enumerate}
\item A \emph{lattice of level $i$} is a self-dual $\mathcal O_E$-lattice $N\subset E^{a}$, together
  with a \emph{frame} at each place of the set $S_B$ of bad places, an isomorphism of $E^{a}$ carrying
  $N$ to the standard local lattice there. For a finite set $S$ of finite places of $E$, lattices
  $N\subset E^{a}$ and $N'\subset E^{b}$, and an integer $c\geq1$, let $H_S(N,N')$ be the
  $\mathcal O_E$-module of $x\in\operatorname{Hom}_E(E^{a},E^{b})$ with $xN\subset N'$ at every place of
  $S$, and $H_S(N,N')/c$ and $cH_S(N,N')$ its rescalings.
\item The \emph{coset theta function} of $x\in H_S(N,N')/c$ is the entire function
  $\Theta_{N,N',c,x}$ on $V_{a,b}$, the theta function of the coset of $x$ modulo $cH_S(N,N')$ built
  from the vacuum vector of the pair of levels and the CM lattice of the host
  (Definition~\ref{is-vacuum-theta-vectors}), normalized by the square root of the covolume. The \emph{theta coefficients} of $x$ are
  $J_x(m)=\sigma_0^{-1}\big(\chi_F(N,N')^{-1}j_c(\Theta_{N,N',c,x})(m)\big)$, with the periods $c_\tau$ of
  the kernel data and a constant $\chi_F(N,N')$.
\item The \emph{theta distribution} at weights $\lambda$, $\lambda'$ and $n\in\mathbf N^{\mathcal T}$
  sends $(N,N',c,x)$, for $c\mid D_0$ and $x\in H_S(N,N')/D_0$, to the operator
  $\mu(N)\mu(N')^{-1}\,\mathrm{Jet}(J_x,n)\colon V(\lambda)\to V(\lambda')$, and to $0$ otherwise; here
  $D_0$ is a fixed modulus and $\mu(N)\in L$ is the product of the characters $\mu_U$ and $\mu_I$ of the
  kernel data at the determinants of the frames of $N$.
\item For a function $\Xi$ on $\operatorname{Hom}_E(E^{a},E^{b})$ with values in $L$ and integers
  $c_A,c_B$, the \emph{kernel entry} between $N$ and $N'$ is
  \[
   K_\Xi(N,N')=\sum_{x\in(H_S(N,N')/c_A)/c_BH_S(N,N')}\Xi(x)\,\theta(N,N',c_B,x),
  \]
  with $\theta$ the theta distribution. The kernel of the $(i+1)$-th lift of the kernel chain
  (Definition~\ref{is-kernel-ties}) between classes of lattices $N$ and $N'$ is $K_\Xi(N,N')$ for the
  coefficient function $\Xi$ of the chain, with $S$ the places dividing the level, $c_A=(g_0g_1)^{fA}$
  and $c_B=T_u(g_0g_1)^{fB}$ for the local constants $A,B\geq1$.
\item The \emph{coupling set} of a lattice $N\subset E^{n}$ is the set of cosets
  $y+\Lambda_0$, where $y$ runs over the lattice that agrees with $N$ scaled by the inverse uniformizer
  at the places of the level set and by the $M$-th power of the inverse uniformizer at the other places
  of $S$, and $\Lambda_0$ is
  the lattice that agrees with $N$ at the places of the level set and with $N$ scaled by the $M$-th
  power of the uniformizer elsewhere in $S$; $M$ is a local constant of the kernel data.
\end{enumerate}
\end{definition}
\leanref{Formalism/LocalToGlobal/KernelLattices/ThDistC.lean}{Formalism.L2.Ker.Amb.thInd}
\leanref{Formalism/LocalToGlobal/KernelLattices/ThDistC.lean}{Formalism.L2.Ker.Amb.jetK}
\leanref{Formalism/LocalToGlobal/KernelLattices/ThDistC.lean}{Formalism.L2.Ker.thDistC}
\leanref{Formalism/LocalToGlobal/KernelLattices/ThDistC.lean}{Formalism.L2.Ker.muInvDet}
\leanref{Formalism/LocalToGlobal/Columns/FiniteSumNonzeroLargeDepth.lean}{Formalism.L2.Col.kerOf}
\leanref{Formalism/LocalToGlobal/Columns/FiniteSumNonzeroLargeDepth.lean}{Formalism.L2.Col.ybSet}

\begin{definition}\label{is-growth-theta-functions}
Let $D$ be a definite-side datum and $\Lambda$ a CM lattice of its host, with the notation of
Definition~\ref{is-vacuum-theta-vectors}: the space $V_{a,b}$ with the forms $\langle\ ,\ \rangle$ and
$H$, the vectors $\iota(x)$ for $x\in M_{b\times a}(E)$, the operators $U_x$, the normalization
$\alpha_{\sqrt d/2}$, the subgroups $\Lambda_{L,L'}$, the theta functions and norms for a subgroup,
the theta vectors up to units, the numbers $n(\Lambda,t)$, and the conjugates with their
superscript $\rho$.
\begin{enumerate}
\item A \emph{torsion point} of a subgroup $\Gamma\subset V_{a,b}$ is a $y\in V_{a,b}$ with
  $my\in\Gamma$ for some integer $m\geq1$.
\item When $g_2(\Lambda)$ and $g_3(\Lambda)$ are algebraic, the conjugates are formed in the same way
  for an embedding $\rho$ of $\overline{\mathbf Q}$ into $\mathbf C$. For a function $G$ on $V_{a,b}$
  whose Taylor coefficients $g_m$ at $0$ are algebraic, $G^\rho(z)=\sum_m\rho(g_m)z^m$ where the
  series converges.
\end{enumerate}
\end{definition}
\leanref{Formalism/SharpRestriction/GlobalSchurBounds/StandardLatticesGL3.lean}{Formalism.MainNotions.StandardLatticesGL3.LaQ}
\leanref{Formalism/SharpRestriction/GlobalSchurBounds/GeneralLinearActions/IsDiscriminantTwelfthRootOfIdealLattice.lean}{Formalism.MainNotions.TaylorLogMetric.lamConj}
\leanref{Formalism/SharpRestriction/GlobalSchurBounds/GeneralLinearActions/IsDiscriminantTwelfthRootOfIdealLattice.lean}{Formalism.MainNotions.TaylorLogMetric.fnConj}

\begin{lemma}\label{is-growth-period-lattice}
Let $D$ be a definite-side datum, $\Lambda=\delta\varphi_0(\mathfrak c)$ a CM lattice of its host, and
$L\subset E^a$, $L'\subset E^b$ lattices that are self-dual for $\sum_i\bar x_iy_i$, with the
notation of Definition~\ref{is-growth-theta-functions}.
\begin{enumerate}
\item $\operatorname{Im}H(x,y)\in\mathbf Z$ for all $x,y\in\Lambda_{L,L'}$.
\item If $x\in V_{a,b}$ and $\operatorname{Im}H(x,y)\in\mathbf Z$ for every $y\in\Lambda_{L,L'}$,
  then $x\in\Lambda_{L,L'}$.
\item $\Lambda_{L,L'}$ is finitely generated and spans $V_{a,b}$ over $\mathbf R$.
\item $\frac{\sqrt d}2H(x,x)\in\mathbf Z$ for every $x\in\Lambda_{L,L'}$.
\item For every $x\in M_{b\times a}(E)$ the vector $\delta\iota(x)$ is a torsion point of
  $\Lambda_{L,L'}$; and for torsion points $y,y'$ of $\Lambda_{L,L'}$ the numbers
  $\operatorname{Im}H(y,y')$ and $\frac{\sqrt d}2H(y,y)$ are rational.
\end{enumerate}
\end{lemma}
\leanref{Development/AbelianVarietyTheory/Proofs/FoundationProofs/ThetaSpaceSelfDualPairAtMostLine.lean}{GlobalTheory.ThetaLine.theta_line}
\leanref{Development/AbelianVarietyTheory/Proofs/DivisorProofs/RationalPoints.lean}{GlobalTheory.RationalPoints.rational_points}
\begin{proof}
Let $Y$ be the $\mathcal O_E$-span of the matrices $sB$ with $s\in\mathfrak b^{-1}$ and $BL\subset L'$,
so that $\Lambda_{L,L'}$ is generated by the $w\iota(x)$ with $w\in\Lambda$ and $x\in Y$, and recall
$H=A(\Lambda)^{-1}\langle\ ,\ \rangle$. Put
$h(x,y)=\operatorname{Tr}_{E/K}\sum_{\beta,\alpha}x_{\beta\alpha}\overline{y_{\beta\alpha}}$ on
$M_{b\times a}(E)$. The proof uses four facts.

(a) $\Lambda=\mathbf Z\omega_1+\mathbf Z\omega_2$ with
$A(\Lambda)=|\operatorname{Im}(\overline{\omega_1}\omega_2)|>0$, and $z\in\mathbf C$ lies in $\Lambda$
if and only if $\operatorname{Im}(z\overline w)/A(\Lambda)\in\mathbf Z$ for every $w\in\Lambda$: for
$z=s\omega_1+t\omega_2$ with $s,t$ real, the values at $w=\omega_2$ and $w=\omega_1$ are $\pm s$ and
$\pm t$.

(b) Every type restricts to $\varphi_0$ on $K$ and commutes with complex conjugation, and
$\sum_\tau\tau(t)=\varphi_0(\operatorname{Tr}_{E/K}t)$ for $t\in E$. So
$\langle\iota(x),\iota(y)\rangle=\varphi_0(h(x,y))$ for $x,y\in M_{b\times a}(E)$, and
$\iota(kx)=\varphi_0(k)\iota(x)$ for $k\in K$.

(c) $h(x,y)\in\mathcal O_K$ for $x,y\in Y$, and there are $y_1,\dots,y_r$ and $y'_1,\dots,y'_r$ in
$Y$ with $x=\sum_ih(x,y'_i)\,y_i$ for every $x\in M_{b\times a}(E)$. For the first, let $x=sB$ and
$y=s'B'$ with $s,s'\in\mathfrak b^{-1}$ and $BL,B'L\subset L'$. The adjoint $B'^*$ of $B'$ for
$\sum_i\bar x_iy_i$ maps $L'$ into $L$, by self-duality, so
$\sum_{\beta,\alpha}B_{\beta\alpha}\overline{B'_{\beta\alpha}}=\operatorname{tr}(B'^*B)$ lies in
$\mathcal O_E$. Since $\mathfrak D_{E/K}=\mathfrak b\overline{\mathfrak b}$
(Proposition~\ref{hf-host-field-properties}), $\mathfrak b^{-1}$ is self-dual for
$\operatorname{Tr}_{E/K}(x\overline y)$, so
$h(x,y)=\operatorname{Tr}_{E/K}(s\overline{s'}\operatorname{tr}(B'^*B))\in\mathcal O_K$; the general
case follows by additivity and $h(kx,y)=h(x,\overline ky)$. For the second, multiply dual families
of $\mathfrak b^{-1}$ for $\operatorname{Tr}_{E/K}(x\overline y)$ and of $L'$ and $L$ for
$\sum_i\bar x_iy_i$, which exist because these lattices are self-dual.

(d) The vectors $\iota(x)$, $x\in M_{b\times a}(E)$, span $V_{a,b}$ over $\mathbf C$, because
distinct embeddings of $E$ are linearly independent over $\mathbf C$.

\emph{Step 1.} If $z\in\Lambda_{L,L'}$ and $x\in Y$, then $\langle z,\iota(x)\rangle\in\Lambda$. For
$z=w\iota(x')$ with $w\in\Lambda$ and $x'\in Y$, $\langle z,\iota(x)\rangle=w\varphi_0(h(x',x))$ by
(b), which lies in $\Lambda$ by (c), since $\Lambda=\delta\varphi_0(\mathfrak c)$ is stable under
multiplication by $\varphi_0(\mathcal O_K)$. The general case follows by additivity.

\emph{Step 2: proof of (1).} For $y=w\iota(x)$ with $w\in\Lambda$ and $x\in Y$,
$\operatorname{Im}H(z,y)=\operatorname{Im}(\langle z,\iota(x)\rangle\overline w)/A(\Lambda)$, an
integer by Step 1 and (a). The general case follows by additivity in $y$.

\emph{Step 3: an expansion.} Every $z\in V_{a,b}$ satisfies
$z=\sum_i\langle z,\iota(y'_i)\rangle\,\iota(y_i)$. Both sides are $\mathbf C$-linear in $z$. For
$z=\iota(x)$ the right side is
$\sum_i\varphi_0(h(x,y'_i))\iota(y_i)=\iota\big(\sum_ih(x,y'_i)y_i\big)=\iota(x)$ by (b) and (c), and
(d) gives every $z$. So $z\in\Lambda_{L,L'}$ as soon as $\langle z,\iota(y'_i)\rangle\in\Lambda$ for
every $i$.

\emph{Step 4: proof of (2).} The hypothesis for $y=w\iota(y'_i)$, $w\in\Lambda$, gives
$\operatorname{Im}(\langle x,\iota(y'_i)\rangle\overline w)/A(\Lambda)\in\mathbf Z$ for every
$w\in\Lambda$, so $\langle x,\iota(y'_i)\rangle\in\Lambda$ by (a), and $x\in\Lambda_{L,L'}$ by
Step 3.

\emph{Step 5: proof of (3).} By Steps 1 and 3, $\Lambda_{L,L'}$ is generated by the $2r$ vectors
$\omega_k\iota(y_i)$, and by Step 3 these span $V_{a,b}$ over $\mathbf R$, since
$\mathbf C=\mathbf R\omega_1+\mathbf R\omega_2$.

\emph{Step 6: proof of (4).} We have $\omega_K=\varphi_0(\omega)$ for some $\omega\in\mathcal O_K$.
For $x\in\Lambda_{L,L'}$, $\langle\omega_Kx,\iota(y'_i)\rangle=\omega_K\langle x,\iota(y'_i)\rangle$
lies in $\Lambda$ by Step 1, so $\omega_Kx\in\Lambda_{L,L'}$ by Step 3. By (1),
$\operatorname{Im}H(\omega_Kx,x)=\operatorname{Im}(\omega_K)H(x,x)=\frac{\sqrt d}2H(x,x)$ is an
integer.

\emph{Step 7: proof of (5).} Let $x\in M_{b\times a}(E)$. There is an integer $m\geq1$ with
$mx\in Y$, and $n=\mathrm N\mathfrak c$ lies in $\mathfrak c$, so
$nm\,\delta\iota(x)=\delta\varphi_0(n)\,\iota(mx)\in\Lambda_{L,L'}$. If
$my,m'y'\in\Lambda_{L,L'}$ with $m,m'\geq1$, then $mm'\operatorname{Im}H(y,y')\in\mathbf Z$ by (1)
and $m^2\frac{\sqrt d}2H(y,y)\in\mathbf Z$ by (4).
\end{proof}

\begin{lemma}\label{is-growth-sup-mean-square}
Let $\kappa$ be a finite set with $g$ elements, give $\mathbf C^\kappa$ the norm
$|u|^2=\sum_{x\in\kappa}|u_x|^2$ and Lebesgue measure, and let $v\geq1$. Let
$\Gamma\subset\mathbf C^\kappa$ be a subgroup with $|x|^2\geq1/v$ for every nonzero $x\in\Gamma$, and
$\mathcal F$ a measurable fundamental domain for $\Gamma$ of volume $1$. Let $G$ be an entire
function on $\mathbf C^\kappa$ such that $N(u)=|G(u)|^2e^{-\pi|u|^2}$ satisfies $N(u+x)=N(u)$ for all
$x\in\Gamma$ and all $u$, and is integrable on $\mathcal F$. Then for every $z\in\mathbf C^\kappa$,
\[
 N(z)\le3(2gv+1)^g\int_{\mathcal F}N(u)\,du .
\]
\end{lemma}
\leanref{Development/AbelianVarietyTheory/Proofs/FoundationProofs/IntegralLeOfDisjointTranslates.lean}{GlobalTheory.ThetaSupBall.sup_le_mean_square}
\begin{proof}
Let $r>0$ with $r^2=1/(4v)$, and let $B$ be the open ball of radius $r$ about $z$.

\emph{Step 1: the ball and the fundamental domain.} The translates $y+B$, $y\in\Gamma$, are pairwise
disjoint: if $y_1+p=y_2+q$ with $p,q\in B$, then $|y_1-y_2|=|q-p|<2r$, so $|y_1-y_2|^2<1/v$ and
$y_1=y_2$. In particular $\Gamma$ is countable. As $N\geq0$ is invariant under $\Gamma$,
\[
 \int_BN=\sum_{y\in\Gamma}\int_{(y+B)\cap\mathcal F}N\le\int_{\mathcal F}N ,
\]
since the sets $(y+B)\cap\mathcal F$ are disjoint subsets of $\mathcal F$.

\emph{Step 2: a holomorphic function on the ball.} Put
$\phi(u)=G(u)\exp\big(-\pi\sum_x\overline{z_x}(u_x-z_x)\big)$, an entire function with
$\phi(z)=G(z)$. From $|u|^2=|z|^2+2\operatorname{Re}\sum_x\overline{z_x}(u_x-z_x)+|u-z|^2$,
\[
 |\phi(u)|^2e^{-\pi|z|^2}=e^{\pi|u-z|^2}N(u)\le e^{\pi/4}N(u)\le3N(u)\qquad(u\in B).
\]

\emph{Step 3: the mean value inequality.} The ball $B$ has volume $\pi^gr^{2g}/g!$, and
$|\phi(z)|^2\le\frac{g!}{\pi^gr^{2g}}\int_B|\phi(u)|^2du$. Indeed, expand
$\phi(u)=\sum_mc_m(u-z)^m$, which converges uniformly on $B$. The monomials $(u-z)^m$ are orthogonal
on $B$, because $B$ is invariant under $u_x-z_x\mapsto e^{i\theta_x}(u_x-z_x)$; so
$\int_B|\phi|^2=\sum_m|c_m|^2\int_B|(u-z)^m|^2\geq|c_0|^2\operatorname{vol}(B)$, and $c_0=\phi(z)$.

\emph{Step 4.} By Steps 3, 2 and 1,
\[
 N(z)=|\phi(z)|^2e^{-\pi|z|^2}\le\frac{g!}{\pi^gr^{2g}}\int_B|\phi|^2e^{-\pi|z|^2}
 \le\frac{3\,g!}{\pi^gr^{2g}}\int_{\mathcal F}N .
\]
Finally $g!/(\pi r^2)^g=g!\,(4v/\pi)^g\le(4gv/\pi)^g\le(2gv+1)^g$, since $g!\le g^g$ and $\pi\geq2$.
\end{proof}

\begin{lemma}\label{is-growth-translate-bound}
Let $D$ be a definite-side datum whose host field has root discriminant $r_E$, let $\Lambda$ be a CM
lattice of the host, $L\subset E^a$ and $L'\subset E^b$ lattices that are self-dual for
$\sum_i\bar x_iy_i$, and $G$ a theta function for $\Lambda_{L,L'}$
(Definition~\ref{is-growth-theta-functions}). Then for all $y,z\in V_{a,b}$,
\[
 |U_yG(z)|^2\le s_0(a,b)\,\|G\|^2_{\Lambda_{L,L'}}\,e^{\pi H(z,z)},\qquad
 s_0(a,b)=3(2fab\,r_E+1)^{fab}.
\]
\end{lemma}
\leanref{Development/AbelianVarietyTheory/Proofs/FoundationProofs/SupLeEuclidean.lean}{GlobalTheory.ThetaSupPair.sup_translate_le}
\begin{proof}
Write $\Gamma_0=\Lambda_{L,L'}$. Let $\kappa=\mathcal T\times\{1,\dots,b\}\times\{1,\dots,a\}$, with
$g=fab$ elements, and let $\phi(z)=A(\Lambda)^{-1/2}z$, a $\mathbf C$-linear isomorphism of $V_{a,b}$
onto $\mathbf C^\kappa$ with $\langle\phi(z),\phi(w)\rangle=H(z,w)$, where on $\mathbf C^\kappa$ we
write $\langle u,u'\rangle=\sum_xu_x\overline{u'_x}$. Put $\Gamma=\phi(\Gamma_0)$.

\emph{Step 1: the minimum.} For nonzero $x\in\Gamma_0$ the number $\frac{\sqrt d}2H(x,x)$ is a
positive integer by Lemma~\ref{is-growth-period-lattice}(4), so
$|\phi(x)|^2=H(x,x)\geq2/\sqrt d\geq1/r_E$. Here $\sqrt d\le r_E$, because the root discriminant of
$K\subset E$ is at least $\sqrt d$ and at most $r_E$; and $r_E\geq1$.

\emph{Step 2: covolume one.} By Lemma~\ref{is-growth-period-lattice}(1) to (3) and Step 1, $\Gamma$
is a lattice in $\mathbf C^\kappa$, and
$\Gamma=\{x:\operatorname{Im}\langle x,y\rangle\in\mathbf Z\text{ for all }y\in\Gamma\}$. As
$\operatorname{Im}\langle x,y\rangle=\operatorname{Re}\langle x,iy\rangle$, $\Gamma$ is the dual
lattice of $i\Gamma$ for the real inner product $\operatorname{Re}\langle\,,\rangle$, so the product
of the covolumes of $\Gamma$ and $i\Gamma$ is $1$. Multiplication by $i$ preserves volume, so
$\Gamma$ has covolume $1$, and each of its fundamental domains has volume $1$.

\emph{Step 3: periodicity.} For $x\in\Gamma_0$ we have $|\alpha_{\sqrt d/2}(x)|=1$ and
$|U_xG(z)|^2e^{-\pi H(z,z)}=|G(z+x)|^2e^{-\pi H(z+x,z+x)}$, because
$H(z+x,z+x)=H(z,z)+2\operatorname{Re}H(z,x)+H(x,x)$. Since $\alpha_{\sqrt d/2}(x)U_xG=G$, the continuous
function $N(u)=|G(\phi^{-1}u)|^2e^{-\pi|u|^2}$ is invariant under $\Gamma$. It is integrable on a
bounded fundamental domain, hence on every fundamental domain.

\emph{Step 4.} Lemma~\ref{is-growth-sup-mean-square}, with $v=r_E$, gives
$N(u)\le s_0(a,b)\int_{\mathcal F}N$ for every fundamental domain $\mathcal F$ of $\Gamma$. As $\phi$
multiplies Lebesgue measure by a constant and $\mathcal F$ has volume $1$, $\int_{\mathcal F}N$ is
the mean of $|G(z)|^2e^{-\pi H(z,z)}$ over the fundamental domain $\phi^{-1}(\mathcal F)$ of
$\Gamma_0$, that is, $\|G\|^2_{\Lambda_{L,L'}}$. Finally, by the computation of Step 3 with $y$ in
place of $x$, $|U_yG(z)|^2=N(\phi(z+y))\,e^{\pi H(z,z)}$.
\end{proof}

\begin{lemma}\label{is-growth-theta-genus-one}
Let $D$ be a definite-side datum, $\Lambda$ a CM lattice of its host
(Definition~\ref{is-vacuum-theta-vectors}), $\wp_\Lambda$ its Weierstrass function, and
$\sigma^*_\Lambda$, $H_\Lambda$, $U^\Lambda_t$ and $\vartheta_{\Lambda,t}$ as in
Definition~\ref{pd-theta-function}; put $q_\Lambda(x)=\frac{\sqrt d}2H_\Lambda(x,x)$. A \emph{torsion
point} of $\Lambda$ is a $t\in\mathbf C$ with $mt\in\Lambda$ for some integer $m\geq1$.
\begin{enumerate}
\item The numbers $g_2(\Lambda)$ and $g_3(\Lambda)$ are algebraic; so are $\wp_\Lambda(t)$ and
  $\wp'_\Lambda(t)$ for every torsion point $t\notin\Lambda$, and the Taylor coefficients at $0$ of
  $U^\Lambda_t\sigma^*_\Lambda$ for every torsion point $t$.
\item Let $\rho$ be an automorphism of the field $\mathbf C$ that is the identity on $\varphi_0(K)$, and
  $\Lambda^\rho$ the lattice with $g_k(\Lambda^\rho)=\rho(g_k(\Lambda))$ for $k=2,3$. Torsion points $t$ of
  $\Lambda$ and $t'$ of $\Lambda^\rho$ \emph{correspond} if $t\in\Lambda$ and $t'\in\Lambda^\rho$, or if
  $t\notin\Lambda$, $t'\notin\Lambda^\rho$, $\wp_{\Lambda^\rho}(t')=\rho(\wp_\Lambda(t))$ and
  $\wp'_{\Lambda^\rho}(t')=\rho(\wp'_\Lambda(t))$.
  \begin{enumerate}
  \item For every integer $n\geq1$, $\Lambda^\rho=\delta'\varphi_0(\mathfrak c')$ for a nonzero ideal
    $\mathfrak c'$ of $\mathcal O_K$ of norm prime to $2dn$ and a $\delta'$ with
    $\delta'^{12}=\Delta(\varphi_0(\mathfrak c'))$. In particular $\Lambda^\rho$ is a CM lattice of the host.
  \item Every torsion point of $\Lambda$ corresponds to a torsion point of $\Lambda^\rho$, and every
    torsion point of $\Lambda^\rho$ to one of $\Lambda$. If $t$ corresponds to $t_1$, then it corresponds
    to $t_2$ if and only if $t_1-t_2\in\Lambda^\rho$. If $t$ corresponds to $t'$ and $u$ to $u'$, then $t+u$
    corresponds to $t'+u'$, $t+\omega$ to $t'$ for $\omega\in\Lambda$, and $\varphi_0(a)t$ to
    $\varphi_0(a)t'$ for $a\in\mathcal O_K$; and $nt\in\Lambda$ if and only if $nt'\in\Lambda^\rho$, for
    every integer $n$.
  \item Let $m\geq1$, and let $x,x'\in\mathbf C$ with $mx\in2\Lambda$ and $x/2m$ corresponding to
    $x'/2m$. Then $\rho$ carries the Taylor coefficients at $0$ of $U^\Lambda_x\sigma^*_\Lambda$ to those
    of $U^{\Lambda^\rho}_{x'}\sigma^*_{\Lambda^\rho}$. For unkilled two-torsion points $t$ of $\Lambda$ and
    $t'$ of $\Lambda^\rho$, it carries those of $\vartheta_{\Lambda,t}$ to those of
    $\vartheta_{\Lambda^\rho,t'}$.
  \item Let $m\geq1$. If $ma,mb\in\Lambda$, $a/m$ corresponds to $a'/m$ and $b/m$ to $b'/m$, then
    $\rho\big(e^{2\pi i\operatorname{Im}H_\Lambda(a,b)}\big)=e^{2\pi i\operatorname{Im}H_{\Lambda^\rho}(a',b')}$.
    If $mx\in\Lambda$ and $x/2m$ corresponds to $x'/2m$, then
    $\rho\big(e^{\pi iq_\Lambda(x)}\big)=e^{\pi iq_{\Lambda^\rho}(x')}$.
  \end{enumerate}
\item Let $c$ be complex conjugation and $\bar\Lambda=c(\Lambda)$. Then
  $g_k(\bar\Lambda)=c(g_k(\Lambda))$, $A(\bar\Lambda)=A(\Lambda)$ and
  $\sigma^*_{\bar\Lambda}(z)=c\big(\sigma^*_\Lambda(c(z))\big)$.
\end{enumerate}
\end{lemma}
\leanref{Development/AbelianVarietyTheory/Proofs/FoundationProofs/IntegralJAndAlgebraicData.lean}{GlobalTheory.CurveAlgebraic.integral_j_and_algebraic_data}
\leanref{Development/AbelianVarietyTheory/Proofs/FoundationProofs/CmThetaTaylorCoeffAlgebraicAndConjugate.lean}{GlobalTheory.ThetaConj.cmTheta_taylorCoeff_algebraic_and_conjugate}
\leanref{Development/AbelianVarietyTheory/Proofs/FoundationProofs/ConjugateCurveDatumLatticeIsCurveLattice.lean}{GlobalTheory.ConjCurve.galCurve_holds}
\leanref{Development/AbelianVarietyTheory/Proofs/FoundationProofs/GalStarComplexMultiplicationPoints.lean}{GlobalTheory.ConjPoint.galStar_holds}
\leanref{Development/AbelianVarietyTheory/Proofs/FoundationProofs/GaloisConjTaylorThetaTranslate.lean}{GlobalTheory.JetConj.galois_translate_sgStar}
\leanref{Development/AbelianVarietyTheory/Proofs/FoundationProofs/WeilPairingGaloisEquivariantCmTorsionStatement.lean}{GlobalTheory.T.weilPairing_galois_equivariant_cm_torsion}
\leanref{Development/AbelianVarietyTheory/Proofs/FoundationProofs/CmQuadraticPhaseEquivariantEmbeddingsStatement.lean}{GlobalTheory.T.cmQuadraticPhase_equivariant_embeddings}
\leanref{Development/AbelianVarietyTheory/Proofs/FoundationProofs/CmThetaComplexConjugationCompatible.lean}{GlobalTheory.ThetaConj.useBar_holds}

These are facts on the elliptic curve $\mathcal E_\Lambda:y^2=4x^3-g_2(\Lambda)x-g_3(\Lambda)$, which is
parametrized by $t\mapsto(\wp_\Lambda(t),\wp'_\Lambda(t))$, has complex multiplication by
$\mathcal O_K$ through $t\mapsto\varphi_0(a)t$, and has discriminant $\Delta(\Lambda)=1$. In (1),
$j(\Lambda)=1728g_2(\Lambda)^3$ is an algebraic integer by Theorem~\ref{pd-j-invariant-integral}, and
$27g_3(\Lambda)^2=g_2(\Lambda)^3-1$. The points of $\mathcal E_\Lambda$ of a given finite order are
finitely many and are permuted by the automorphisms of $\mathbf C$ over $\overline{\mathbf Q}$, so their
coordinates are algebraic. These automorphisms fix $\Lambda$ and every torsion point corresponds to
itself under them, so by (2c) they fix the Taylor coefficients of $U^\Lambda_t\sigma^*_\Lambda$, which
are therefore algebraic. In (2), $\rho$ carries $\mathcal E_\Lambda$ to $\mathcal E_{\Lambda^\rho}$, its
group law and its points of finite order to those of $\mathcal E_{\Lambda^\rho}$, and, since it fixes
$\varphi_0(K)$, the endomorphism $\varphi_0(a)$ to $\varphi_0(a)$ \cite[Theorem~II.2.2]{Sil94}. This
gives (2b). It also gives $\varphi_0(\mathcal O_K)\Lambda^\rho\subset\Lambda^\rho$, so $\Lambda^\rho$ is
$\delta'$ times the image of an ideal of $\mathcal O_K$, which may be chosen in its class prime to $2dn$;
with $\Delta(\Lambda^\rho)=\rho(\Delta(\Lambda))=1$ this is (2a). Statement (2c) is a case of Shimura's
reciprocity law for theta functions with complex multiplication \cite{Shi76}; it contains the Galois
behaviour of the number $s_2(\Lambda)$ of Definition~\ref{pd-theta-function} for lattices with complex
multiplication \cite{Mas75}. The first identity of (2d) is the Galois invariance of the Weil pairing on
the points of order dividing $m^2$ \cite[Proposition~III.8.1]{Sil09}, which on $\mathbf C/\Lambda$ is
$(a/m,b/m)\mapsto e^{2\pi i\operatorname{Im}H_\Lambda(a,b)}$ up to a convention of sign common to
$\Lambda$ and $\Lambda^\rho$. The second is the same statement for the quadratic refinement
$e^{\pi iq_\Lambda}$ of this pairing, as $q_\Lambda(x)=\operatorname{Im}H_\Lambda(\omega_Kx,x)$.
Statement (3) follows from the definitions, as in the proof of Lemma~\ref{is-vacuum-genus-one}(4).

By (1) the numbers in (2) are algebraic, and every embedding of $\overline{\mathbf Q}$ into $\mathbf C$
extends to an automorphism of $\mathbf C$. So (2) holds for every embedding $\rho$ of
$\overline{\mathbf Q}$ into $\mathbf C$ that is the identity on $\varphi_0(K)$, with the lattice
$\Lambda^\rho$ of Definition~\ref{is-growth-theta-functions}.

\begin{lemma}\label{is-growth-theta-algebraic}
Let $D$ be a definite-side datum, $\Lambda$ a CM lattice of its host, and $v=(v_1,\dots,v_r)$ an
orthonormal basis of $M_{b\times a}(E)$, with the notation of
Definition~\ref{is-growth-theta-functions}; for a function $g$ on $\mathbf C$ let $g_v$ be its tensor
power along $v$. Let $\mathcal G_v$ be the set of entire functions $G$ on $V_{a,b}$ such that $U_xG$
has algebraic Taylor coefficients at $0$ for every torsion point $x$ of $\Lambda_v$.
\begin{enumerate}
\item Let $x=\sum_jc_j\iota(v_j)$ and let $g$ be a function on $\mathbf C$. Then
  $(U_xg_v)(z)=\prod_j(U^\Lambda_{c_j}g)(\langle z,\iota(v_j)\rangle)$, and $x$ is a torsion point of
  $\Lambda_v$ if and only if every $c_j$ is a torsion point of $\Lambda$.
\item The function $(\sigma^*_\Lambda)_v$ lies in $\mathcal G_v$.
\item If $G,G'\in\mathcal G_v$, $\kappa\in\overline{\mathbf Q}$ and $t$ is a torsion point of
  $\Lambda_v$, then $\kappa G$, $G+G'$ and $U_tG$ lie in $\mathcal G_v$.
\item For every standard subgroup $\Gamma'\subset V_{a,b}$ (Lemma~\ref{is-vacuum-averaging}) there is an
  integer $m\geq1$ with $m\Gamma'\subset\Lambda_v$. If $G\in\mathcal G_v$ and $\Gamma\subset V_{a,b}$ is a
  subgroup with $[\Gamma':\Gamma\cap\Gamma']$ finite, then $I_{\Gamma\to\Gamma'}G\in\mathcal G_v$.
\item Let $t$ be an unkilled two-torsion point of $\Lambda$, $L\subset E^a$ and $L'\subset E^b$ lattices
  that are self-dual for $\sum_i\bar x_iy_i$, and $\kappa$ an algebraic number. Then
  $F=\kappa I_{\Lambda_v\to\Lambda_{L,L'}}\big((\vartheta_{\Lambda,t})_v\big)$ lies in $\mathcal G_v$. In
  particular $F$ and $U_yF$, for every torsion point $y$ of $\Lambda_{L,L'}$, have algebraic Taylor
  coefficients at $0$.
\end{enumerate}
\end{lemma}
\leanref{Development/AbelianVarietyTheory/Proofs/FoundationProofs/ThetaTaylorCoeffsAlgebraicAtComplexMultiplicatio.lean}{GlobalTheory.ThetaAlgebraic.U_pullB}
\leanref{Development/AbelianVarietyTheory/Proofs/FoundationProofs/ThetaTaylorCoeffsAlgebraicAtComplexMultiplicatio.lean}{GlobalTheory.ThetaAlgebraic.good_sigma}
\leanref{Development/AbelianVarietyTheory/Proofs/FoundationProofs/ThetaTaylorCoeffsAlgebraicAtComplexMultiplicatio.lean}{GlobalTheory.ThetaAlgebraic.Good.transl}
\leanref{Development/AbelianVarietyTheory/Proofs/FoundationProofs/ThetaTaylorCoeffsAlgebraicAtComplexMultiplicatio.lean}{GlobalTheory.ThetaAlgebraic.la_le}
\leanref{Development/AbelianVarietyTheory/Proofs/FoundationProofs/ThetaTaylorCoeffsAlgebraicAtComplexMultiplicatio.lean}{GlobalTheory.ThetaAlgebraic.Good.iop}
\leanref{Development/AbelianVarietyTheory/Proofs/FoundationProofs/ThetaTaylorCoeffsAlgebraicAtComplexMultiplicatio.lean}{GlobalTheory.ThetaAlgebraic.good_theta}
\leanref{Development/AbelianVarietyTheory/Proofs/FoundationProofs/ThetaTaylorCoeffsAlgebraicAtComplexMultiplicatio.lean}{GlobalTheory.ThetaAlgebraic.taylor_algebraic}
\begin{proof}
Write $A=A(\Lambda)$, so that $H=A^{-1}\langle\ ,\ \rangle$ and $H_\Lambda(u,c)=u\bar c/A$.

(1) By Lemma~\ref{is-vacuum-tensor-norm}(1), $c_j=\langle x,\iota(v_j)\rangle$,
$\langle z,x\rangle=\sum_j\langle z,\iota(v_j)\rangle\overline{c_j}$ and $\langle x,x\rangle=\sum_j|c_j|^2$.
Since $\langle z+x,\iota(v_j)\rangle=\langle z,\iota(v_j)\rangle+c_j$, the values of $g$ in $(U_xg_v)(z)$
are those on the right side, and the factor $\exp(-\pi H(z,x)-\pi H(x,x)/2)$ is the product over $j$ of
$\exp\big(-\pi H_\Lambda(\langle z,\iota(v_j)\rangle,c_j)-\pi H_\Lambda(c_j,c_j)/2\big)$. The coordinates
$x\mapsto(\langle x,\iota(v_j)\rangle)_j$ carry $\Lambda_v=\sum_j\Lambda\iota(v_j)$ onto $\Lambda^r$, which
gives the second assertion.

(2) Let $x=\sum_jc_j\iota(v_j)$ be a torsion point of $\Lambda_v$. By (1), $U_x(\sigma^*_\Lambda)_v$ is the
product over $j$ of the functions $U^\Lambda_{c_j}\sigma^*_\Lambda$, whose Taylor coefficients at $0$ are
algebraic by Lemma~\ref{is-growth-theta-genus-one}(1), composed with the linear forms
$z\mapsto\langle z,\iota(v_j)\rangle=\sum_{\tau,\beta,\alpha}z_{\tau,\beta,\alpha}\overline{\tau((v_j)_{\beta\alpha})}$,
whose coefficients are algebraic. Composition with such linear forms, products and sums of entire
functions preserve the algebraicity of the Taylor coefficients at $0$.

(3) The operators $U_x$ are linear, which gives the first two assertions. For $x,t\in V_{a,b}$, expanding
the definition gives $U_xU_tG=e^{-\pi i\operatorname{Im}H(x,t)}U_{x+t}G$, as in
Lemma~\ref{pd-analytic-translation-system}(4). If $x$ and $t$ are torsion points of $\Lambda_v$, so is
$x+t$, and $\operatorname{Im}H(x,t)$ is rational: if $mx=\sum_jc_j\iota(v_j)$ and
$mt=\sum_jc'_j\iota(v_j)$ lie in $\Lambda_v$, then $m^2\operatorname{Im}H(x,t)=\sum_j\operatorname{Im}(c_j
\overline{c'_j})/A$ by (1), and $\operatorname{Im}(c\bar c')/A$ is an integer for $c,c'\in\Lambda$ (Step (a)
of the proof of Lemma~\ref{is-growth-period-lattice}). So $e^{-\pi i\operatorname{Im}H(x,t)}$ is a root of
unity, and $U_x(U_tG)$ has algebraic Taylor coefficients.

(4) Let $\Gamma'=\Lambda_{L,L'}$. For $s\in\mathfrak A_0$ and $B$ with $BL\subset L'$, $sB=\sum_jh_jv_j$ with
$h_j=\operatorname{Tr}_{E/K}\sum_{\beta,\alpha}(sB)_{\beta\alpha}\overline{(v_j)_{\beta\alpha}}\in K$, since $v$
is orthonormal over $K$, and $\iota(sB)=\sum_j\varphi_0(h_j)\iota(v_j)$ (Step (b) of the proof of
Lemma~\ref{is-growth-period-lattice}). As $\Lambda$ is an $\mathcal O_K$-module, $w\iota(sB)\in
m^{-1}\Lambda_v$ for $w\in\Lambda$ if $m$ is a common denominator of the $h_j$; and $\Lambda_{L,L'}$ is
finitely generated (Lemma~\ref{is-growth-period-lattice}(3)), so one $m$ serves for all of it. For
$\Gamma'=\Lambda_{v'}$, with another orthonormal basis $v'$, the same argument applies to the $v'_k$.
Now let $x_1,\dots,x_n$ be representatives of $\Gamma'/(\Gamma\cap\Gamma')$, so that
$I_{\Gamma\to\Gamma'}G=n^{-1/2}\sum_i\alpha_{\sqrt d/2}(x_i)U_{x_i}G$. Each $x_i$ is a torsion point of
$\Lambda_v$. The number $\frac{\sqrt d}2H(x_i,x_i)$ is an integer, by
Lemma~\ref{is-growth-period-lattice}(4) for $\Lambda_{L,L'}$ and by (1) and
Lemma~\ref{is-vacuum-genus-one}(2) for $\Lambda_{v'}$, so $\alpha_{\sqrt d/2}(x_i)=\pm1$; and $n^{-1/2}$ is
algebraic. So (3) applies.

(5) Since $2t\in\Lambda$ and $t\notin\Lambda$, the number $u=(U^\Lambda_t\sigma^*_\Lambda)(0)$ is algebraic by
Lemma~\ref{is-growth-theta-genus-one}(1), and nonzero because $\sigma_\Lambda$ vanishes only on $\Lambda$;
and $\vartheta_{\Lambda,t}=u^{-1}U^\Lambda_t\sigma^*_\Lambda$. With $x_t=\sum_jt\iota(v_j)$, a torsion point of
$\Lambda_v$, (1) gives $(\vartheta_{\Lambda,t})_v=u^{-r}U_{x_t}(\sigma^*_\Lambda)_v$, which lies in
$\mathcal G_v$ by (2) and (3). By (4) and (3), $F\in\mathcal G_v$. Finally $0$ and, by (4), every torsion
point of $\Lambda_{L,L'}$ are torsion points of $\Lambda_v$.
\end{proof}

\begin{lemma}\label{is-growth-theta-galois}
Let $D$ be a definite-side datum, $\Lambda$ a CM lattice of its host, $v=(v_1,\dots,v_r)$ an
orthonormal basis of $M_{b\times a}(E)$, and $\rho$ an embedding of $\overline{\mathbf Q}$ into
$\mathbf C$ that is the identity on $\varphi_0(K)$, with the notation of
Definition~\ref{is-growth-theta-functions}, and with $\Lambda^\rho$ and the correspondence of torsion
points of Lemma~\ref{is-growth-theta-genus-one}(2). Write $\iota^\rho(B)=(\rho(\tau(B)))_{\tau\in\mathcal T}$,
$g^\rho_v(z)=\prod_jg(\langle z,\iota^\rho(v_j)\rangle)$ for a function $g$ on $\mathbf C$, and, for a set
$Y\subset M_{b\times a}(E)$, $\Lambda_Y$ and $\Lambda^\rho_Y$ for the subgroups of $V_{a,b}$ generated by
the $w\iota(B)$ with $w\in\Lambda$ and $B\in Y$, and by the $w'\iota^\rho(B)$ with $w'\in\Lambda^\rho$ and
$B\in Y$. For $n\geq1$, vectors $x=\sum_jc_j\iota(v_j)$ and $x'=\sum_jc'_j\iota^\rho(v_j)$
\emph{correspond at level $n$} if $c_j/n$ corresponds to $c'_j/n$ for every $j$.
\begin{enumerate}
\item Let $n\geq1$. If $x$ and $y$ correspond at level $n$ to $x'$ and $y'$, then $x-y$ corresponds to
  $x'-y'$. Every torsion point of $\Lambda_v$ corresponds at level $n$ to some vector. Let
  $Y\subset M_{b\times a}(E)$ and $k\geq1$ with $k\mid n$, such that $kB$ lies in the $\mathcal O_K$-span of
  $v$ for every $B\in Y$, and $kv_j$ in the subgroup generated by $Y$ for every $j$. Then every element
  of $\Lambda_Y$ corresponds at level $n$ to an element of $\Lambda^\rho_Y$, every element of
  $\Lambda^\rho_Y$ to one of $\Lambda_Y$, and if $x$ corresponds to $x'$ at level $n$, then $x\in\Lambda_Y$
  if and only if $x'\in\Lambda^\rho_Y$. Such a $k$ exists for $Y=\{v_1,\dots,v_r\}$, with
  $\Lambda_Y=\Lambda_v$, and, for lattices $L\subset E^a$ and $L'\subset E^b$ that are self-dual for
  $\sum_i\bar x_iy_i$, for the set $Y$ of the $sB$ with $s\in\mathfrak A_0$ and $BL\subset L'$, with
  $\Lambda_Y=\Lambda_{L,L'}$.
\item Let $m\geq1$, let $G=\sum_i\kappa_iU_{t_i}(\sigma^*_\Lambda)_v$ with $\kappa_i\in\overline{\mathbf Q}$ and
  $mt_i\in\Lambda_v$, and let $t_i$ correspond to $t'_i$ at level $4m$. Then
  $G^\rho=\sum_i\rho(\kappa_i)U^\rho_{t'_i}(\sigma^*_{\Lambda^\rho})^\rho_v$. If moreover $mx\in\Lambda_v$ and
  $x$ corresponds to $x'$ at level $2m$, then $(U_xG)^\rho=U^\rho_{x'}G^\rho$ and
  $\rho\big(e^{\pi iq(x)}\big)=e^{\pi iq^\rho(x')}$, where $q(x)=\frac{\sqrt d}2H(x,x)$ and
  $q^\rho(x')=\frac{\sqrt d}2H^\rho(x',x')$.
\end{enumerate}
\end{lemma}
\leanref{Development/AbelianVarietyTheory/Proofs/FoundationProofs/TorsionMapCarriesTorusLattices.lean}{GlobalTheory.ThetaConj.adpB_lattices}
\leanref{Development/AbelianVarietyTheory/Proofs/FoundationProofs/ExistsHomLatticeCommensurableOrthonormalStatement.lean}{GlobalTheory.T.exists_homLattice_commensurable_orthonormal}
\leanref{Development/AbelianVarietyTheory/Proofs/FoundationProofs/GaloisConjTorsionTranslatesOddThetaProduct.lean}{GlobalTheory.ProdThetaConj.exists_conjPair}
\leanref{Development/AbelianVarietyTheory/Proofs/FoundationProofs/GaloisConjTorsionTranslatesOddThetaProduct.lean}{GlobalTheory.ProdThetaConj.galB_all}
\begin{proof}
(1) The types $\rho\circ\tau$ agree with $\varphi_0$ on $K$, because $\rho$ fixes $\varphi_0(K)$. For
$B\in Y$ write $kB=\sum_ja_j(B)v_j$ with $a_j(B)\in\mathcal O_K$; then
$kw\iota(B)=\sum_j\varphi_0(a_j(B))w\iota(v_j)$ and $kw'\iota^\rho(B)=\sum_j\varphi_0(a_j(B))w'\iota^\rho(v_j)$.
So the coordinates of the elements of $k\Lambda_Y$ are the $\varphi_0(\mathcal O_K)$-combinations of
elements of $\Lambda$ with the coefficients $a_j(B)$, for $\Lambda$ and for $\Lambda^\rho$ alike, and
$k\Lambda_v\subset\Lambda_Y$. Correspondence of torsion points is compatible with sums, with
multiplication by $\varphi_0(\mathcal O_K)$, and with membership in $\Lambda$ and $\Lambda^\rho$, and every
torsion point has a correspondent, determined modulo $\Lambda^\rho$
(Lemma~\ref{is-growth-theta-genus-one}(2b)). Applied to the coordinates, divided by $n$ with $k\mid n$,
this gives the assertions. The set $Y=\{v_j\}$ satisfies the hypothesis with $k=1$. For the set $Y$
attached to $(L,L')$, the subgroup generated by $Y$ is the $\mathcal O_E$-lattice $\mathfrak
A\operatorname{Hom}(L,L')$ of $M_{b\times a}(E)$; so the coordinates along $v$ of its elements have a
common denominator, and a positive integer multiple of each $v_j$ lies in it.

(2) Write $t_i=\sum_jw_{ij}\iota(v_j)$, $t'_i=\sum_jw'_{ij}\iota^\rho(v_j)$, $x=\sum_jc_j\iota(v_j)$ and
$x'=\sum_jc'_j\iota^\rho(v_j)$. The vectors $\iota^\rho(v_j)$ are orthonormal by
Lemma~\ref{is-vacuum-tensor-norm}(1), applied to the types $\rho\circ\tau$. By
Lemma~\ref{is-growth-theta-algebraic}(1), for $\Lambda$ and for $\Lambda^\rho$,
\[
 \big(U_xU_{t_i}(\sigma^*_\Lambda)_v\big)(z)=\prod_j\big(U^\Lambda_{c_j}U^\Lambda_{w_{ij}}\sigma^*_\Lambda\big)
 (\langle z,\iota(v_j)\rangle),
\]
and likewise with $\Lambda^\rho$, $c'_j$, $w'_{ij}$ and $\iota^\rho$. The factors compare as follows.

(a) $\rho$ carries the coefficients of the linear form $z\mapsto\langle z,\iota(v_j)\rangle$ to those of
$z\mapsto\langle z,\iota^\rho(v_j)\rangle$, because $E$ is a CM field, so that
$\rho(\overline{\tau(e)})=\rho(\tau(\bar e))=\overline{\rho(\tau(e))}$ for $e\in E$.

(b) Since $mw_{ij}\in\Lambda$ and $w_{ij}/4m$ corresponds to $w'_{ij}/4m$,
Lemma~\ref{is-growth-theta-genus-one}(2c) with $2m$ in place of $m$ shows that $\rho$ carries the Taylor
coefficients of $U^\Lambda_{w_{ij}}\sigma^*_\Lambda$ to those of
$U^{\Lambda^\rho}_{w'_{ij}}\sigma^*_{\Lambda^\rho}$.

(c) $\rho$ carries the Taylor coefficients of $U^\Lambda_{c_j}U^\Lambda_{w_{ij}}\sigma^*_\Lambda$ to those of
$U^{\Lambda^\rho}_{c'_j}U^{\Lambda^\rho}_{w'_{ij}}\sigma^*_{\Lambda^\rho}$. Indeed, choose $y_j$ with $c_j/4m$
corresponding to $y_j/4m$. Then $c_j/2m$ corresponds to $y_j/2m$ and to $c'_j/2m$, so
$c'_j-y_j\in2m\Lambda^\rho$ (Lemma~\ref{is-growth-theta-genus-one}(2b)). Since
$U^{\Lambda^\rho}_\omega\sigma^*_{\Lambda^\rho}=\sigma^*_{\Lambda^\rho}$ for $\omega\in2\Lambda^\rho$ (proof of
Lemma~\ref{pd-analytic-translation-system}), and $\operatorname{Im}H_{\Lambda^\rho}$ is integral on
$\Lambda^\rho$ while $my_j,mw'_{ij}\in\Lambda^\rho$, the phases in
$U^{\Lambda^\rho}_{c'_j}U^{\Lambda^\rho}_{w'_{ij}}=e^{-\pi i\operatorname{Im}H_{\Lambda^\rho}(c'_j,w'_{ij})}
U^{\Lambda^\rho}_{c'_j+w'_{ij}}$ show that replacing $c'_j$ by $y_j$ does not change this function. Now
$U^\Lambda_{c_j}U^\Lambda_{w_{ij}}=e^{-\pi i\operatorname{Im}H_\Lambda(c_j,w_{ij})}U^\Lambda_{c_j+w_{ij}}$, likewise
for $y_j,w'_{ij}$, and $(c_j+w_{ij})/4m$ corresponds to $(y_j+w'_{ij})/4m$; so (b) applies to the
translates, and Lemma~\ref{is-growth-theta-genus-one}(2d), applied with $2m$ to $a=w_{ij}$ and
$b=c_j/2$, gives $\rho\big(e^{-\pi i\operatorname{Im}H_\Lambda(c_j,w_{ij})}\big)=e^{-\pi i\operatorname{Im}
H_{\Lambda^\rho}(y_j,w'_{ij})}$, as $-\operatorname{Im}H(c,w)=2\operatorname{Im}H(w,c/2)$.

Conjugation of Taylor coefficients by $\rho$ commutes with sums, products, multiplication by algebraic
constants, and composition with linear forms with algebraic coefficients, the conjugate form being
taken. With (a) and (b) this gives the formula for $G^\rho$, and with (a) and (c) it gives
$(U_xG)^\rho=\sum_i\rho(\kappa_i)U^\rho_{x'}U^\rho_{t'_i}(\sigma^*_{\Lambda^\rho})^\rho_v=U^\rho_{x'}G^\rho$.
Finally $q(x)=\sum_jq_\Lambda(c_j)$ and $q^\rho(x')=\sum_jq_{\Lambda^\rho}(c'_j)$ by orthonormality, and
Lemma~\ref{is-growth-theta-genus-one}(2d) gives $\rho\big(e^{\pi iq_\Lambda(c_j)}\big)=
e^{\pi iq_{\Lambda^\rho}(c'_j)}$ for each $j$, since $mc_j\in\Lambda$.
\end{proof}

\begin{lemma}\label{is-growth-theta-odd-norm}
Let $\Lambda\subset\mathbf C$ be a lattice of covolume $A$ with $\Delta(\Lambda)=g_2(\Lambda)^3-27g_3(\Lambda)^2=1$,
with $\sigma^*_\Lambda$ and $U^\Lambda_t$ as in Definition~\ref{pd-theta-function}, and for a function $g$ on
$\mathbf C$ let $n_\Lambda(g)$ be the mean of $|g(z)|^2e^{-\pi|z|^2/A}$ over a fundamental domain of
$\Lambda$. Then $n_\Lambda(\sigma^*_\Lambda)^2=2\pi^2/A$, and
$n(\Lambda,t)=n_\Lambda(\sigma^*_\Lambda)/|(U^\Lambda_t\sigma^*_\Lambda)(0)|^2$ for every $t\in\mathbf C$ with
$(U^\Lambda_t\sigma^*_\Lambda)(0)\ne0$.
\end{lemma}
\leanref{Development/AbelianVarietyTheory/Proofs/FoundationProofs/OddThetaNormUnimodularLattice.lean}{GlobalTheory.OddThetaNorm.odd_theta_norm_sigma_star}
\begin{proof}
The first identity is a computation: Alpöge writes $\Lambda=a(\mathbf Z+\mathbf Z\tau)$ with
$\operatorname{Im}\tau>0$, identifies $\sigma^*_\Lambda$ as a multiple of the odd Jacobi theta function in
$z/a$, by the uniqueness of an entire function with the translation law of $\sigma^*_\Lambda$ and derivative
$1$ at $0$, computes the mean for that function, and evaluates the multiple with Jacobi's formula for the
eighth power of the derivative at $0$ of the odd theta function and the condition $\Delta(\Lambda)=1$. For
the second, $\vartheta_{\Lambda,t}=(U^\Lambda_t\sigma^*_\Lambda)(0)^{-1}U^\Lambda_t\sigma^*_\Lambda$, and
$|(U^\Lambda_t\sigma^*_\Lambda)(z)|^2e^{-\pi|z|^2/A}=|\sigma^*_\Lambda(z+t)|^2e^{-\pi|z+t|^2/A}$, whose mean over a
fundamental domain is that of $|\sigma^*_\Lambda(z)|^2e^{-\pi|z|^2/A}$, as this function is periodic for
$\Lambda$.
\end{proof}

\begin{lemma}\label{is-growth-padded-basis}
Let $D$ be a definite-side datum with host field $E$, and let $a,b\geq0$ be integers. Then
$M_{b\times4a}(E)$ has an orthonormal basis (Definition~\ref{is-vacuum-theta-vectors}(4)).
\end{lemma}
\leanref{Development/AbelianVarietyTheory/Proofs/FoundationProofs/ImaginaryQuadraticFieldStructureStatement.lean}{GlobalTheory.T.exists_orthonormal_frame_var}
\leanref{Development/AbelianVarietyTheory/Proofs/FoundationProofs/ImaginaryQuadraticFieldStructure.lean}{GlobalTheory.FourCopies.exists_onB_four}
\begin{proof}
For matrices $x,y$ of the same size over $E$ put
$h(x,y)=\operatorname{Tr}_{E/K}\sum_{\beta,\alpha}x_{\beta\alpha}\overline{y_{\beta\alpha}}\in K$. It is $K$-linear
in $x$, and $h(y,x)=\overline{h(x,y)}$. If $x\ne0$, then $h(x,x)$ is a positive rational number:
$\varphi_0(h(x,x))=\sum_\tau\sum_{\beta,\alpha}|\tau(x_{\beta\alpha})|^2$ by fact (b) in the proof of
Lemma~\ref{is-growth-period-lattice}, which is real and positive, and an element of $K$ with real image under
$\varphi_0$ is rational.

By the Gram--Schmidt process for $h$ there are nonzero $v_1,\dots,v_n\in M_{b\times a}(E)$, pairwise
orthogonal for $h$, that span $M_{b\times a}(E)$ over $K$; they are linearly independent, so $n=fab$. Put
$\mu_i=h(v_i,v_i)\in\mathbf Q_{>0}$. By Lagrange's four-square theorem, the product of the numerator and the
denominator of $\mu_i$ is a sum of four squares of integers; dividing these integers by the numerator gives
rationals $p,q,r,s$ (depending on $i$) with $\mu_i(p^2+q^2+r^2+s^2)=1$. The matrix
\[
 A_i=\begin{pmatrix}p&q&r&s\\-q&p&-s&r\\-r&s&p&-q\\-s&-r&q&p\end{pmatrix}
\]
has orthogonal rows of squared length $p^2+q^2+r^2+s^2$, so $\mu_iA_iA_i^{\mathsf T}=\mu_iA_i^{\mathsf T}A_i=1$.

For $c\in\mathbf Q^4$ and $v\in M_{b\times a}(E)$ let $[c]v\in M_{b\times4a}(E)$ be the matrix whose four
blocks of $a$ columns are $c_1v,\dots,c_4v$; then $h([c]v,[c']w)=(\sum_kc_kc'_k)\,h(v,w)$. Let
$u_{ij}=[A_{i,j}]v_i$ for $i\le n$ and $j\le4$, where $A_{i,j}$ is the $j$-th row of $A_i$. Then
$h(u_{ij},u_{i'j'})=0$ for $i\ne i'$, and $h(u_{ij},u_{ij'})=\mu_i(A_iA_i^{\mathsf T})_{jj'}=\delta_{jj'}$. For the
standard basis vector $e_k$ of $\mathbf Q^4$, $\sum_j\mu_i(A_i)_{jk}u_{ij}=[e_k]v_i$, because
$\mu_iA_i^{\mathsf T}A_i=1$. As the $v_i$ span $M_{b\times a}(E)$ and $[e_k]v$ is $K$-linear in $v$, every
$[e_k]x$ lies in the $K$-span of the $u_{ij}$, and every matrix in $M_{b\times4a}(E)$ is a sum
$\sum_k[e_k]x_k$. So the $4n=f\cdot4a\cdot b$ matrices $u_{ij}$ are orthonormal and span
$M_{b\times4a}(E)$; they form an orthonormal basis.
\end{proof}

\begin{lemma}\label{is-growth-padded-average}
Let $D$ be a definite-side datum, $\Lambda$ a CM lattice of its host and $a,b\geq0$ integers, with the
notation of Definition~\ref{is-growth-theta-functions}; write $\alpha=\alpha_{\sqrt d/2}$ and
$q(x)=\frac{\sqrt d}2H(x,x)$. Let $\Gamma\subset V_{a,b}$ be a subgroup with $\operatorname{Im}H(x,y)\in\mathbf Z$
and $q(x)\in\mathbf Z$ for all $x,y\in\Gamma$.
\begin{enumerate}
\item If $\omega_K\Gamma\subset\Gamma$, then $\alpha(x+y)=\alpha(x)\alpha(y)e^{-\pi i\operatorname{Im}H(x,y)}$ for all
  $x,y\in\Gamma$.
\item Let $F$ be a theta function for $\Gamma$, $m_0\geq1$ an integer, $x\in V_{a,b}$ with $m_0x\in\Gamma$,
  and $l\in\Gamma$. Then $U_{x+2m_0l}F=U_xF$.
\item Suppose that $\Gamma$ is finitely generated, that $\omega_K\Gamma\subset\Gamma$, and that every
  $x\in V_{a,b}$ with $\operatorname{Im}H(x,y)\in\mathbf Z$ for all $y\in\Gamma$ lies in $\Gamma$. Let $f\ne0$ be
  an entire function on $V_{a,b}$ and $n\geq1$ an integer with $U_{nx}f=f$ for every $x\in\Gamma$. Then
  there is a nonzero theta function for $\Gamma$ of the form $\sum_i\kappa_iU_{s_i}f$, a finite sum with
  $\kappa_i\in\mathbf C$ and $s_i$ torsion points of $\Gamma$.
\end{enumerate}
\end{lemma}
\leanref{Development/AbelianVarietyTheory/Proofs/FoundationProofs/SignedTranslationAddSelfDualLatticeStatement.lean}{GlobalTheory.T.signedTranslation_add_selfDual_lattice}
\leanref{Development/AbelianVarietyTheory/Proofs/FoundationProofs/ThetaTranslationCombPeriodicTwiceLevelStatement.lean}{GlobalTheory.T.thetaTranslation_comb_periodic_twice_level}
\leanref{Development/AbelianVarietyTheory/Proofs/FoundationProofs/ExistsNonzeroThetaOfFiniteLevelAverageStatement.lean}{GlobalTheory.T.exists_nonzero_theta_of_finite_level_average}
\begin{proof}
For $x,y\in V_{a,b}$, $U_xU_y=e^{-\pi i\operatorname{Im}H(x,y)}U_{x+y}$ (proof of
Lemma~\ref{is-growth-theta-algebraic}(3)), and hence $U_xU_y=e^{-2\pi i\operatorname{Im}H(x,y)}U_yU_x$. For
$x\in\Gamma$, $\alpha(x)=e^{\pi iq(x)}=(-1)^{q(x)}$.

(1) Let $x,y\in\Gamma$. Then $q(x+y)=q(x)+q(y)+\sqrt d\operatorname{Re}H(x,y)$. The numbers
$k'=\operatorname{Im}H(x,y)$ and $k=\operatorname{Im}H(\omega_Kx,y)=\frac12\operatorname{Im}H(x,y)
+\frac{\sqrt d}2\operatorname{Re}H(x,y)$ are integers, since $\omega_Kx\in\Gamma$. So
$\sqrt d\operatorname{Re}H(x,y)=2k-k'$, and $\alpha(x+y)=(-1)^{q(x)+q(y)+2k-k'}=\alpha(x)\alpha(y)e^{-\pi ik'}$.

(2) As $q(2m_0l)=4m_0^2q(l)$ is even, $\alpha(2m_0l)=1$ and $U_{2m_0l}F=F$. So
$U_xF=U_xU_{2m_0l}F=e^{-\pi i\operatorname{Im}H(x,2m_0l)}U_{x+2m_0l}F$, and
$\operatorname{Im}H(x,2m_0l)=2\operatorname{Im}H(m_0x,l)$ is an even integer.

(3) Put $N=2n$ and $W_x=\alpha(x)U_x$ for $x\in\Gamma$, and choose a set $Q$ of representatives of
$\Gamma/N\Gamma$, which is finite because $\Gamma$ is finitely generated.

\emph{Step 1.} $W_xW_y=W_{x+y}$ for $x,y\in\Gamma$, by (1) and the product formula for $U_xU_y$.

\emph{Step 2.} $W_{Nx}U_{y/N}f=U_{y/N}f$ for $x,y\in\Gamma$. Indeed $q(Nx)=N^2q(x)$ is even, so
$\alpha(Nx)=1$; $U_{Nx}f=f$, as $Nx=n\cdot2x$; and $U_{Nx}U_{y/N}=e^{-2\pi i\operatorname{Im}H(x,y)}U_{y/N}U_{Nx}$
with $\operatorname{Im}H(x,y)\in\mathbf Z$.

\emph{Step 3: the averages.} For $y\in\Gamma$ put $P_y=\sum_{a\in Q}W_aU_{y/N}f$. By Steps 1 and 2,
$W_{a+Nc}U_{y/N}f=W_aW_{Nc}U_{y/N}f=W_aU_{y/N}f$ for $c\in\Gamma$, so each summand depends only on the class of
$a$ modulo $N\Gamma$. For $x\in\Gamma$, $W_xP_y=\sum_{a\in Q}W_{x+a}U_{y/N}f=P_y$ by Step 1, because the $x+a$
also run over the classes. So $P_y$ is a theta function for $\Gamma$. It has the required form:
$W_aU_{y/N}f=\alpha(a)e^{-\pi i\operatorname{Im}H(a,y/N)}U_{a+y/N}f$, and $N(a+y/N)\in\Gamma$.

\emph{Step 4.} Suppose that $P_y=0$ for every $y\in\Gamma$. For $a,y\in\Gamma$ put
$\chi_a(y)=e^{-2\pi i\operatorname{Im}H(a,y)/N}$. Then $U_{y/N}W_af=\chi_a(y)^{-1}W_aU_{y/N}f$, so
$U_{y/N}\big(\sum_{a\in Q}\chi_a(y)W_af\big)=P_y=0$, and $\sum_{a\in Q}\chi_a(y)W_af=0$, since $U_{y/N}G=0$ only
for $G=0$.

\emph{Step 5.} Each $\chi_a$ is a character of $\Gamma$, trivial on $N\Gamma$, and $\chi_a=1$ for $a\in N\Gamma$,
because $\operatorname{Im}H$ is integral on $\Gamma$. If $a\in Q$ and $\sum_{y\in Q}\chi_a(y)\ne0$, then the
character $\chi_a$ of the finite group $\Gamma/N\Gamma$ is trivial, so $\operatorname{Im}H(a/N,y)\in\mathbf Z$
for every $y\in\Gamma$; then $a/N\in\Gamma$ by hypothesis, and $a\in N\Gamma$. So
$\sum_{y\in Q}\chi_a(y)=0$ for $a\in Q$ outside $N\Gamma$.

\emph{Step 6.} Let $a_0$ be the element of $Q$ in $N\Gamma$. Summing Step 4 over $y\in Q$ and using Step 5,
$0=\sum_{a\in Q}\big(\sum_{y\in Q}\chi_a(y)\big)W_af=|Q|\,W_{a_0}f$. By Step 2 with $y=0$, $W_{a_0}f=f$, so
$f=0$, which is excluded. Hence $P_y\ne0$ for some $y\in\Gamma$.
\end{proof}

\begin{lemma}\label{is-growth-padded-restriction}
Let $D$ be a definite-side datum, $\Lambda$ a CM lattice of its host, $L\subset E^a$ and $L'\subset E^b$
lattices that are self-dual for $\sum_i\bar x_iy_i$, $a'\geq a$ an integer and $u=(u_1,\dots,u_r)$ an
orthonormal basis of $M_{b\times a'}(E)$, with the notation of Definition~\ref{is-growth-theta-functions} on
$V_{a,b}$ and on $V_{a',b}$. For $z\in V_{a,b}$ let $\tilde z\in V_{a',b}$ be obtained by appending $a'-a$ zero
columns to each matrix $z_\tau$, and put $S=(\sigma^*_\Lambda)_u$.
\begin{enumerate}
\item There is an integer $R\geq1$ with $R\tilde x\in\Lambda_u$ for every $x\in\Lambda_{L,L'}$.
\item There is a $t\in V_{a',b}$ with $2t\in\Lambda_u$, $(U_tS)(0)\ne0$ and $U_{4w}U_tS=U_tS$ for every
  $w\in\Lambda_u$.
\item For every theta function $F$ for $\Lambda_{L,L'}$ there is a finite sum $G_1=\sum_i\kappa_iU_{t_i}S$,
  with $\kappa_i\in\mathbf C$ and $t_i$ torsion points of $\Lambda_u$, such that $F(z)=G_1(\tilde z)$ for every
  $z\in V_{a,b}$.
\end{enumerate}
\end{lemma}
\leanref{Development/AbelianVarietyTheory/Proofs/FoundationProofs/ExistsNsmulPaddedLatticeMemBasisProductStatement.lean}{GlobalTheory.T.exists_nsmul_padded_lattice_mem_basis_product}
\leanref{Development/AbelianVarietyTheory/Proofs/FoundationProofs/TranslateLevelFourStatement.lean}{GlobalTheory.T.U_pullB_ne_zero_and_invariant_four_smul}
\leanref{Development/AbelianVarietyTheory/Proofs/FoundationProofs/ThetaTranslateGaloisRationalityCmPoints.lean}{GlobalTheory.GalTheta.exists_restriction}
\begin{proof}
Write $B\mapsto\tilde B$ for the same operation on $M_{b\times a}(E)$; since $\tau(0)=0$,
$\iota(\tilde B)=\widetilde{\iota(B)}$. Since $\widetilde{z+x}=\tilde z+\tilde x$ and $H(\tilde z,\tilde x)=H(z,x)$:

$(\ast)$ if $G$ is a function on $V_{a',b}$ and $G_0(z)=G(\tilde z)$, then $(U_xG_0)(z)=(U_{\tilde x}G)(\tilde z)$
for $x,z\in V_{a,b}$.

(1) The $x\in V_{a,b}$ with $R_x\tilde x\in\Lambda_u$ for some integer $R_x\geq1$ form a subgroup. It contains
the generators $w\iota(B)$ of $\Lambda_{L,L'}$, $w\in\Lambda$, $B=sB'$ with $s\in\mathfrak A_0$ and
$B'L\subset L'$. Indeed $\widetilde{w\iota(B)}=w\iota(\tilde B)$; as $u$ is a basis over $K$,
$\tilde B=\sum_k\kappa_ku_k$ with $\kappa_k\in K$, and $w\iota(\tilde B)=\sum_k\varphi_0(\kappa_k)w\,\iota(u_k)$,
because the types agree with $\varphi_0$ on $K$; if $n\geq1$ and $n\kappa_k\in\mathcal O_K$ for every $k$,
then $nw\iota(\tilde B)\in\Lambda_u$, since $\Lambda$ is stable under multiplication by
$\varphi_0(\mathcal O_K)$. So the subgroup contains $\Lambda_{L,L'}$, which is finitely generated
(Lemma~\ref{is-growth-period-lattice}(3)); the product of the $R_x$ over a finite set of generators serves as
$R$.

(2) Write $\Lambda=\mathbf Z\omega_1+\mathbf Z\omega_2$; then $\omega_1/2\notin\Lambda$. By
Lemma~\ref{is-vacuum-tensor-norm}(1) the coordinates $z\mapsto(\langle z,\iota(u_k)\rangle)_k$ identify
$V_{a',b}$ with $\mathbf C^r$ and $H(z,z')$ with $A(\Lambda)^{-1}\sum_kz_k\overline{z'_k}$, and they carry
$\Lambda_u$ onto $\Lambda^r$. So $\operatorname{Im}H(x,x')=\sum_k\operatorname{Im}(x_k\overline{x'_k})/A(\Lambda)$
is an integer for $x,x'\in\Lambda_u$, by fact (a) in the proof of Lemma~\ref{is-growth-period-lattice}. Let $t$
be the vector all of whose coordinates are $\omega_1/2$; then $2t\in\Lambda_u$, and by
Lemma~\ref{is-growth-theta-algebraic}(1),
$(U_tS)(0)=\big((U^\Lambda_{\omega_1/2}\sigma^*_\Lambda)(0)\big)^r=\big(\sigma^*_\Lambda(\omega_1/2)
e^{-\pi H_\Lambda(\omega_1/2,\omega_1/2)/2}\big)^r$, which is nonzero because $\sigma_\Lambda$ vanishes only on
$\Lambda$. For $c\in\Lambda$, $U^\Lambda_c\sigma^*_\Lambda=\pm\sigma^*_\Lambda$ (proof of
Lemma~\ref{pd-analytic-translation-system}(2)) and $U^\Lambda_cU^\Lambda_c=U^\Lambda_{2c}$, so
$U^\Lambda_{2c}\sigma^*_\Lambda=\sigma^*_\Lambda$. Let $w\in\Lambda_u$, with coordinates $w_k\in\Lambda$. By
Lemma~\ref{is-growth-theta-algebraic}(1), $(U_{4w}S)(z)=\prod_k(U^\Lambda_{4w_k}\sigma^*_\Lambda)
(\langle z,\iota(u_k)\rangle)=S(z)$. Finally $\operatorname{Im}H(4w,t)=2\operatorname{Im}H(w,2t)$ is an even
integer, so the phases in $U_{4w}U_t=e^{-\pi i\operatorname{Im}H(4w,t)}U_{4w+t}$ and
$U_tU_{4w}=e^{-\pi i\operatorname{Im}H(t,4w)}U_{t+4w}$ are $1$, and $U_{4w}U_tS=U_tU_{4w}S=U_tS$.

(3) Let $R$ and $t$ be as in (1) and (2), and $f(z)=(U_tS)(\tilde z)$, an entire function with
$f(0)=(U_tS)(0)\ne0$. For $x\in\Lambda_{L,L'}$, $(U_{4Rx}f)(z)=(U_{4R\tilde x}U_tS)(\tilde z)=f(z)$, by $(\ast)$,
(1) and (2). The subgroup $\Gamma=\Lambda_{L,L'}$ satisfies the hypotheses of
Lemma~\ref{is-growth-padded-average}(3): it is finitely generated, $\operatorname{Im}H$ and
$\frac{\sqrt d}2H(x,x)$ are integral on it, and it contains every $x$ with
$\operatorname{Im}H(x,y)\in\mathbf Z$ for all $y\in\Gamma$, by Lemma~\ref{is-growth-period-lattice}(1) to (4);
and $\omega_K\Gamma\subset\Gamma$ (Step 6 of the proof of Lemma~\ref{is-growth-period-lattice}). So that lemma,
with $n=4R$, gives a nonzero theta function $P=\sum_i\kappa'_iU_{s_i}f$ for $\Lambda_{L,L'}$, with $s_i$
torsion points of $\Lambda_{L,L'}$. By Lemma~\ref{is-vacuum-averaging}(1), $F=\kappa P$ for some
$\kappa\in\mathbf C$. By $(\ast)$ and the product formula for the operators $U$,
\[
 F(z)=\sum_i\kappa\kappa'_i\,(U_{\tilde s_i}U_tS)(\tilde z)
 =\sum_i\kappa\kappa'_ie^{-\pi i\operatorname{Im}H(\tilde s_i,t)}\,(U_{\tilde s_i+t}S)(\tilde z).
\]
If $ms_i\in\Lambda_{L,L'}$ with $m\geq1$, then $Rm\tilde s_i\in\Lambda_u$ by (1), and
$2Rm(\tilde s_i+t)\in\Lambda_u$. So the $\tilde s_i+t$ are torsion points of $\Lambda_u$, and $G_1=\sum_i
\kappa\kappa'_ie^{-\pi i\operatorname{Im}H(\tilde s_i,t)}U_{\tilde s_i+t}S$ has the required form.
\end{proof}

\begin{lemma}\label{is-growth-padded-galois}
Let $D$ be a definite-side datum, $\Lambda$ a CM lattice of its host, $a',b\geq0$ integers,
$u=(u_1,\dots,u_r)$ an orthonormal basis of $M_{b\times a'}(E)$, and $\rho$ an automorphism of the field
$\mathbf C$ that is the identity on $\varphi_0(K)$, with the notation of
Definition~\ref{is-growth-theta-functions} on $V_{a',b}$. Form $\Lambda^\rho$ and the correspondence of torsion
points as in Lemma~\ref{is-growth-theta-genus-one}(2), and $\iota^\rho$, $H^\rho$, $U^\rho$ and the
correspondence of vectors at level $n$ along $u$ as in Lemma~\ref{is-growth-theta-galois}. For an entire
function $G$ on $V_{a',b}$ with Taylor coefficients $g_m$ at $0$, let $G^\rho(z)=\sum_m\rho(g_m)z^m$ where the
series converges.
\begin{enumerate}
\item Let $b_1,\dots,b_s\in M_{b\times a'}(E)$. There is an integer $R\geq1$ such that for every integer
  $n\geq1$, all torsion points $c_1,\dots,c_s$ of $\Lambda$ and all $c'_1,\dots,c'_s\in\mathbf C$ with
  $c_j/nR$ corresponding to $c'_j/nR$ for every $j$, the vectors $\sum_jc_j\iota(b_j)$ and
  $\sum_jc'_j\iota^\rho(b_j)$ correspond at level $n$.
\item Let $n_0,n_1\geq1$ be integers, $c$ a torsion point of $\Lambda$ and $c'\in\mathbf C$ with $c/n_0$
  corresponding to $c'/n_0$. Then there is an $\omega\in\Lambda$ with $(c+n_0\omega)/n_0n_1$ corresponding to
  $c'/n_0n_1$.
\item Let $m\geq1$ be an integer, $G=\sum_i\kappa_iU_{t_i}(\sigma^*_\Lambda)_u$ a finite sum with
  $\kappa_i\in\mathbf C$ and $mt_i\in\Lambda_u$, and $x,x'\in V_{a',b}$ with $mx\in\Lambda_u$ and $x$
  corresponding to $x'$ at level $2m$. Then $G^\rho$ is entire, and $(U_xG)^\rho=U^\rho_{x'}G^\rho$.
\item Let $a\le a'$, and for $z\in V_{a,b}$ let $\tilde z\in V_{a',b}$ be obtained by appending $a'-a$ zero
  columns to each matrix $z_\tau$. Let $G$ be an entire function on $V_{a',b}$ and $G_0(z)=G(\tilde z)$. Then
  for every $z\in V_{a,b}$ the series defining $G^\rho(\tilde z)$ is the series defining $G_0^\rho(z)$
  together with terms equal to $0$; in particular $G_0^\rho(z)=G^\rho(\tilde z)$ if $G^\rho$ is entire.
\end{enumerate}
\end{lemma}
\leanref{Development/AbelianVarietyTheory/Proofs/FoundationProofs/AdaptedPointChangeOfRationalBasisStatement.lean}{GlobalTheory.T.adapted_point_change_of_rational_basis}
\leanref{Development/AbelianVarietyTheory/Proofs/FoundationProofs/AdaptedTorsionLiftDeeperLevelStatement.lean}{GlobalTheory.T.adapted_torsion_lift_deeper_level}
\leanref{Development/AbelianVarietyTheory/Proofs/FoundationProofs/GaloisConjTorsionTranslatesOddThetaProductStatement.lean}{GlobalTheory.T.galois_conj_torsionTranslates_oddThetaProduct}
\leanref{Development/AbelianVarietyTheory/Proofs/FoundationProofs/TaylorConjugateTangentSpacePadZeroColumnsStatement.lean}{GlobalTheory.T.taylorConjugate_tangentSpacePadZeroColumns}
\begin{proof}
The restriction $\rho_0$ of $\rho$ to $\overline{\mathbf Q}$ is the identity on $\varphi_0(K)$. As
$g_2(\Lambda)$, $g_3(\Lambda)$, the $\tau(e)$ for $e\in E$, and the values of $\wp_\Lambda$ and $\wp'_\Lambda$ at
torsion points are algebraic (Lemma~\ref{is-growth-theta-genus-one}(1)), the objects formed with $\rho$ are
those formed with $\rho_0$. So Lemma~\ref{is-growth-theta-galois} applies to $\rho$.

(1) As $u$ spans $M_{b\times a'}(E)$ over $K$, $b_j=\sum_k\kappa_{jk}u_k$ with $\kappa_{jk}\in K$; let $R\geq1$
be an integer with $a_{jk}=R\kappa_{jk}\in\mathcal O_K$ for all $j,k$. The types $\tau$ and $\rho\circ\tau$
agree with $\varphi_0$ on $K$, since $\rho$ fixes $\varphi_0(K)$; so
$\iota(b_j)=\sum_k\varphi_0(\kappa_{jk})\iota(u_k)$ and $\iota^\rho(b_j)=\sum_k\varphi_0(\kappa_{jk})
\iota^\rho(u_k)$. Hence $\sum_jc_j\iota(b_j)=\sum_kC_k\iota(u_k)$ and
$\sum_jc'_j\iota^\rho(b_j)=\sum_kC'_k\iota^\rho(u_k)$, with
$C_k/n=\sum_j\varphi_0(a_{jk})\,(c_j/nR)$ and $C'_k/n=\sum_j\varphi_0(a_{jk})\,(c'_j/nR)$. These correspond,
because correspondence is compatible with sums and with multiplication by $\varphi_0(\mathcal O_K)$
(Lemma~\ref{is-growth-theta-genus-one}(2b)).

(2) Let $l\geq1$ with $lc\in\Lambda$. All correspondences below are those of
Lemma~\ref{is-growth-theta-genus-one}(2b). As $ln_0\cdot c/n_0\in\Lambda$, also
$ln_0n_1\cdot c'/n_0n_1=ln_0\cdot c'/n_0\in\Lambda^\rho$, so $c'/n_0n_1$ is a torsion point of $\Lambda^\rho$ and
corresponds to a torsion point $v$ of $\Lambda$. Then $n_1v$ corresponds to $n_1\cdot c'/n_0n_1=c'/n_0$, and
$-c/n_0$ to $-c'/n_0$, so $\omega=n_1v-c/n_0$ corresponds to $0$ and lies in $\Lambda$. Finally
$(c+n_0\omega)/n_0n_1=v$.

(3) Choose $t'_i$ corresponding to $t_i$ at level $4m$ (Lemma~\ref{is-growth-theta-galois}(1)). The proof of
Lemma~\ref{is-growth-theta-galois}(2) applies to the automorphism $\rho$ of $\mathbf C$. It uses the statements
of Lemma~\ref{is-growth-theta-genus-one}(2), which are stated for such $\rho$, and the fact that conjugation of
Taylor coefficients commutes with sums, products, multiplication by constants and composition with linear forms
with algebraic coefficients; as $\rho$ is an automorphism of $\mathbf C$, this holds for all complex constants,
the conjugate of $\kappa G$ being $\rho(\kappa)G^\rho$. It gives
$G^\rho=\sum_i\rho(\kappa_i)U^\rho_{t'_i}(\sigma^*_{\Lambda^\rho})^\rho_u$, which is entire, and
$(U_xG)^\rho=U^\rho_{x'}G^\rho$.

(4) A monomial in the coordinates of $V_{a',b}$ that involves a coordinate of an appended column vanishes at
$\tilde z$; the other monomials are the monomials in the coordinates of $V_{a,b}$, with the same values at
$\tilde z$ and at $z$. By the chain rule for the linear map $z\mapsto\tilde z$, the Taylor coefficient of $G_0$
at such a monomial is that of $G$.
\end{proof}

\begin{lemma}\label{is-growth-theta-conjugate}
Let $D$ be a definite-side datum, $\Lambda$ a CM lattice of its host with an unkilled two-torsion point
$t$, $L\subset E^a$ and $L'\subset E^b$ lattices that are self-dual for $\sum_i\bar x_iy_i$, with $ab$
even, $F$ a theta vector up to units for $(t,L,L')$, and $\rho$ an embedding of $\overline{\mathbf Q}$
into $\mathbf C$, with the notation of Definition~\ref{is-growth-theta-functions}. Let $c$ be complex
conjugation.
\begin{enumerate}
\item If $\rho$ is the identity on $\varphi_0(K)$, then $\Lambda^\rho$ is a CM lattice of the host, there is
  a permutation $\pi$ of $\mathcal T$ with $\rho\circ\tau=\pi(\tau)$ for every $\tau\in\mathcal T$, and
  there is an automorphism $s$ of $E$ over $K$ with $\rho\circ\tau=\tau\circ s$ for every
  $\tau\in\mathcal T$. Otherwise $\bar\rho=c\circ\rho$ is the identity on $\varphi_0(K)$,
  $\Lambda^\rho=c(\Lambda^{\bar\rho})$, and $G^\rho(z)=c\big(G^{\bar\rho}(c(z))\big)$ for every function $G$
  on $V_{a,b}$ with algebraic Taylor coefficients at $0$. In both cases $\Lambda^\rho$ has an unkilled
  two-torsion point.
\item The series $F^\rho$ converges to an entire function. It is a theta function for
  $\Lambda^\rho_{L,L'}$, and for every unkilled two-torsion point $t'$ of $\Lambda^\rho$ it is a theta vector
  up to units for $(t',L,L')$ formed with $\Lambda^\rho$ and the types $\rho\circ\tau$.
\item For every unkilled two-torsion point $t'$ of $\Lambda^\rho$ and every torsion point $y$ of
  $\Lambda_{L,L'}$, $\|(U_yF)^\rho\|^2_{\Lambda^\rho_{L,L'}}=n(\Lambda^\rho,t')^{fab}$; in particular
  $\|F^\rho\|^2_{\Lambda^\rho_{L,L'}}=n(\Lambda^\rho,t')^{fab}$. This number is positive.
\item If $\rho$ is the identity on $\varphi_0(K)$ and $y$ is a torsion point of $\Lambda_{L,L'}$, there is a
  torsion point $y'$ of $\Lambda^\rho_{L,L'}$ with $(U_yF)^\rho=U^\rho_{y'}F^\rho$.
\end{enumerate}
\end{lemma}
\leanref{Development/AbelianVarietyTheory/Proofs/FoundationProofs/ExistsTwoTorsionCMPointAllConjugates.lean}{GlobalTheory.TwoTorsionPt.exists_point}
\leanref{Development/AbelianVarietyTheory/Proofs/FoundationProofs/InducedTypeThetaGenusOnePermutes.lean}{GlobalTheory.InducedType.permutes}
\leanref{Development/AbelianVarietyTheory/Proofs/FoundationProofs/VacuumNormConjCmPointsAllEmbeddings.lean}{GlobalTheory.ConjNorm.exists_aut}
\leanref{Development/AbelianVarietyTheory/Proofs/FoundationProofs/ThetaVacuaConjugatesCmPoints.lean}{GlobalTheory.ThetaConj.pull_cjF}
\leanref{Development/AbelianVarietyTheory/Proofs/FoundationProofs/ThetaVacuaConjugatesCmPoints.lean}{GlobalTheory.ThetaConj.conj_iop}
\leanref{Development/AbelianVarietyTheory/Proofs/FoundationProofs/ThetaVacuaConjugatesCmPoints.lean}{GlobalTheory.ThetaConj.useConj_all}
\leanref{Development/AbelianVarietyTheory/Proofs/FoundationProofs/VacuumNormConjCmPointsAllEmbeddings.lean}{GlobalTheory.ConjNorm.norm_core}
\leanref{Development/AbelianVarietyTheory/Proofs/FoundationProofs/VacuumNormConjCmPointsAllEmbeddings.lean}{GlobalTheory.ConjNorm.useNorm_all}
\leanref{Development/AbelianVarietyTheory/Proofs/FoundationProofs/VacuumTranslateConjugateJetBound.lean}{GlobalTheory.ConjJet.thetaZeroNormSq_ne_zero_cmLattice}
\leanref{Development/AbelianVarietyTheory/Proofs/FoundationProofs/GaloisDataExists.lean}{GlobalTheory.GalExist.galEx_all}
\leanref{Development/AbelianVarietyTheory/Proofs/FoundationProofs/ThetaTranslateGaloisRationalityCmPoints.lean}{GlobalTheory.GalTheta.gal_all}
\begin{proof}
Write $F=\kappa I_{\Lambda_v\to\Lambda_{L,L'}}G_0$ with $G_0=(\vartheta_{\Lambda,t})_v$, for an orthonormal
basis $v$ and an algebraic $\kappa$ all of whose conjugates have absolute value $1$
(Definition~\ref{is-vacuum-theta-vectors}(6)). By Lemmas~\ref{is-vacuum-genus-one}(2),
\ref{is-vacuum-tensor-norm}(3) and~\ref{is-vacuum-averaging}(3), $G_0$ is a theta function for $\Lambda_v$
and $F$ one for $\Lambda_{L,L'}$. By the proof of Lemma~\ref{is-growth-theta-algebraic}(5),
$G_0=u^{-r}U_{x_t}(\sigma^*_\Lambda)_v$ with $u$ algebraic and $x_t=\sum_jt\iota(v_j)$. Let
$x_1,\dots,x_N$ be representatives of $\Lambda_{L,L'}/(\Lambda_v\cap\Lambda_{L,L'})$. Expanding $I$ and using
$U_{x_i}U_{x_t}=e^{-\pi i\operatorname{Im}H(x_i,x_t)}U_{x_i+x_t}$,
\[
 F=\sum_i\kappa_iU_{x_i+x_t}(\sigma^*_\Lambda)_v,\qquad
 \kappa_i=\kappa N^{-1/2}\alpha_{\sqrt d/2}(x_i)\,u^{-r}e^{-\pi i\operatorname{Im}H(x_i,x_t)},
\]
with $\kappa_i$ algebraic (proof of Lemma~\ref{is-growth-theta-algebraic}(3) and (4)). Fix $m\geq1$ with
$m\Lambda_{L,L'}\subset\Lambda_v$ and $mx_t\in\Lambda_v$ (Lemma~\ref{is-growth-theta-algebraic}(4)), and
$k$ as in Lemma~\ref{is-growth-theta-galois}(1) for the set $Y$ attached to $(L,L')$.

(1) Suppose that $\rho$ is the identity on $\varphi_0(K)$. Then $\Lambda^\rho$ is a CM lattice of the host by
Lemma~\ref{is-growth-theta-genus-one}(2a). For $\tau\in\mathcal T$, $\rho\circ\tau$ is an embedding of $E$
that agrees with $\varphi_0$ on $K$; the types are the $f=[E:K]$ such embeddings, so $\rho\circ\tau$ is a
type, and $\tau\mapsto\rho\circ\tau$ is injective. Fix a type $\tau_0$; every type is $\tau_0\circ g$ for an
automorphism $g$ of $E$ over $K$, and $\rho\circ\tau_0=\tau_0\circ s$ for one such $s$. The host field is
abelian over $\mathbf Q$ (Chapter~\ref{chapter-host-fields}), so
$\rho\circ\tau_0\circ g=\tau_0\circ s\circ g=\tau_0\circ g\circ s$. Otherwise $\rho$ is complex conjugation
on $\varphi_0(K)$, and $\bar\rho$ is the identity there. Then $g_k(\Lambda^\rho)=c(g_k(\Lambda^{\bar\rho}))
=g_k(c(\Lambda^{\bar\rho}))$ by Lemma~\ref{is-growth-theta-genus-one}(3), and a lattice is determined by its
invariants $g_2,g_3$, so $\Lambda^\rho=c(\Lambda^{\bar\rho})$; and $\rho=c\circ\bar\rho$ on the Taylor
coefficients gives $G^\rho(z)=c(G^{\bar\rho}(c(z)))$. A CM lattice of the host has an unkilled two-torsion
point (Lemma~\ref{is-vacuum-genus-one}(1)), and $c$ carries one of $\Lambda^{\bar\rho}$ to one of
$c(\Lambda^{\bar\rho})$, because $c(\omega_K)=\bar\omega_K$.

(2) Suppose first that $\rho$ is the identity on $\varphi_0(K)$, and let $t'$ be an unkilled two-torsion
point of $\Lambda^\rho$. Put $n=4mk$, and choose $x'_i\in\Lambda^\rho_{L,L'}$ corresponding to $x_i$ at
level $n$ (Lemma~\ref{is-growth-theta-galois}(1)).

\emph{Step 1: the tensor power.} $G_0$ is the product of the functions $\vartheta_{\Lambda,t}(\langle
z,\iota(v_j)\rangle)$. By Lemma~\ref{is-growth-theta-genus-one}(2c) and fact (a) in the proof of
Lemma~\ref{is-growth-theta-galois}(2), $G_0^\rho=(\vartheta_{\Lambda^\rho,t'})^\rho_v$.

\emph{Step 2: the representatives.} By Lemma~\ref{is-growth-theta-galois}(1), for the sets $Y$ attached
to $(L,L')$ and to $v$, the differences $x'_i-x'_l$ with $i\ne l$ correspond to $x_i-x_l\notin
\Lambda_v\cap\Lambda_{L,L'}$, so they do not lie in $\Lambda^\rho_v\cap\Lambda^\rho_{L,L'}$. Every
$x'\in\Lambda^\rho_{L,L'}$ corresponds to some $x\in\Lambda_{L,L'}$; then $x-x_i\in\Lambda_v\cap\Lambda_{L,L'}$
for some $i$, and $x'-x'_i\in\Lambda^\rho_v\cap\Lambda^\rho_{L,L'}$. So the $x'_i$ are representatives of
$\Lambda^\rho_{L,L'}/(\Lambda^\rho_v\cap\Lambda^\rho_{L,L'})$.

\emph{Step 3: the conjugate.} By definition $\alpha_{\sqrt d/2}(x)=e^{\pi iq(x)}$ and
$\alpha^\rho_{\sqrt d/2}(x')=e^{\pi iq^\rho(x')}$, with $q$ and $q^\rho$ as in
Lemma~\ref{is-growth-theta-galois}(2). The vectors $x_i$ and $x'_i$ correspond at level $2m$,
because they do at level $n$ and correspondence is compatible with sums. Lemma~\ref{is-growth-theta-galois}(2),
applied to $G_0$, gives $(U_{x_i}G_0)^\rho=U^\rho_{x'_i}G_0^\rho$ and
$\rho(\alpha_{\sqrt d/2}(x_i))=\alpha^\rho_{\sqrt d/2}(x'_i)$. Since $\rho(\sqrt N)^2=N$, $\rho(\sqrt
N)=\varepsilon\sqrt N$ with $\varepsilon=\pm1$. With Steps 1 and 2,
\[
 F^\rho=\rho(\kappa)\varepsilon N^{-1/2}\sum_i\alpha^\rho_{\sqrt d/2}(x'_i)\,U^\rho_{x'_i}
 (\vartheta_{\Lambda^\rho,t'})^\rho_v=\rho(\kappa)\varepsilon\,I^\rho_{\Lambda^\rho_v\to\Lambda^\rho_{L,L'}}
 (\vartheta_{\Lambda^\rho,t'})^\rho_v .
\]
The right side is entire, and all conjugates of $\rho(\kappa)\varepsilon$ have absolute value $1$, since
those of $\kappa$ do. So $F^\rho$ is a theta vector up to units for $(t',L,L')$ formed with $\Lambda^\rho$
and the types $\rho\circ\tau$.

\emph{Step 4: the theta function.} Let $x'\in\Lambda^\rho_{L,L'}$, and $x\in\Lambda_{L,L'}$ corresponding to
$x'$ at level $n$. As in Step 3, now with $F=\sum_i\kappa_iU_{x_i+x_t}(\sigma^*_\Lambda)_v$,
Lemma~\ref{is-growth-theta-galois}(2) gives $(\alpha_{\sqrt d/2}(x)U_xF)^\rho=\alpha^\rho_{\sqrt
d/2}(x')U^\rho_{x'}F^\rho$. As $\alpha_{\sqrt d/2}(x)U_xF=F$, this is $F^\rho$.

Suppose now that $\rho$ is complex conjugation on $\varphi_0(K)$. By (1),
$F^\rho(z)=c(F^{\bar\rho}(c(z)))$. The substitution $z\mapsto c(z)$ carries the objects formed with
$\Lambda^{\bar\rho}$ and the types $\bar\rho\circ\tau$ to those formed with $\Lambda^\rho$ and the types
$\rho\circ\tau=c\circ\bar\rho\circ\tau$: $\iota^\rho(B)=c(\iota^{\bar\rho}(B))$, $\Lambda^\rho_{L,L'}=
c(\Lambda^{\bar\rho}_{L,L'})$, $\Lambda^\rho_v=c(\Lambda^{\bar\rho}_v)$,
$H^\rho(c(z),c(w))=c(H^{\bar\rho}(z,w))$, and, for $G^c(z)=c(G(c(z)))$, $U^\rho_{c(x)}G^c=(U^{\bar\rho}_xG)^c$
and $\alpha^\rho_{\sqrt d/2}(c(x))=c(\alpha^{\bar\rho}_{\sqrt d/2}(x))$. By
Lemma~\ref{is-growth-theta-genus-one}(3), $\vartheta_{c(\Lambda'),c(t'')}=(\vartheta_{\Lambda',t''})^c$ for a
lattice $\Lambda'$. So the case of $\bar\rho$, with the unkilled point $c(t')$ of $\Lambda^{\bar\rho}$,
gives the case of $\rho$.

(3) Let $t'$ be an unkilled two-torsion point of $\Lambda^\rho$ and $y$ a torsion point of
$\Lambda_{L,L'}$. The case $y=0$, where $U_yF=F$, gives the assertion on $\|F^\rho\|^2$.

\emph{Step 1: the norms of the translates of a theta vector up to units.} Let $\Lambda'$ be a CM lattice
of the host with an unkilled two-torsion point $t''$, let $s$ be an automorphism of $E$ over $K$, and let
$G$ be a theta vector up to units for $(t'',L,L')$ formed with $\Lambda'$ and the types $\tau\circ s$. We
show that $G$ is a theta function for the subgroup $\Lambda'_{L,L'}$ formed with these types, and that
$\|U_xG\|^2=n(\Lambda',t'')^{fab}$ for every $x\in V_{a,b}$. The subgroups formed with $\Lambda'$ and the
types $\tau\circ s$ are those formed with $\Lambda'$ and the types $\tau$ after applying $s$ to the
matrices: there are lattices $L_1\subset E^a$ and $L'_1\subset E^b$, self-dual for $\sum_i\bar x_iy_i$,
with $s\big(\mathfrak A_0\operatorname{Hom}(L,L')\big)=\mathfrak A_0\operatorname{Hom}(L_1,L'_1)$, so that the
subgroup formed for $(L,L')$ with the types $\tau\circ s$ is $\Lambda'_{L_1,L'_1}$; and for an orthonormal
basis $w$, $s(w)=(s(w_j))_j$ is an orthonormal basis, as $s$ fixes $K$ and commutes with complex
conjugation, and the lattice and the tensor powers formed along $w$ with the types $\tau\circ s$ are
$\Lambda'_{s(w)}$ and those along $s(w)$. So
$G=\kappa I_{\Lambda'_{s(w)}\to\Lambda'_{L_1,L'_1}}(\vartheta_{\Lambda',t''})_{s(w)}$ for an orthonormal basis $w$
and a $\kappa$ with $|\kappa|=1$. By Lemma~\ref{is-vacuum-genus-one}(2),
Lemma~\ref{is-vacuum-tensor-norm}(2) and (3) apply to $\vartheta_{\Lambda',t''}$: the tensor power
$(\vartheta_{\Lambda',t''})_{s(w)}$ is a theta function for $\Lambda'_{s(w)}$, of squared norm
$n(\Lambda',t'')^{fab}$. By Lemma~\ref{is-vacuum-averaging}(3) its image under the averaging operator is a
theta function for $\Lambda'_{L_1,L'_1}$ with the same norm, and by Lemma~\ref{is-vacuum-averaging}(2)
so is every translate $U_x$ of it. Multiplying by $\kappa$ does not change these norms.

\emph{Step 2: $\rho$ the identity on $\varphi_0(K)$.} By (1), $\Lambda^\rho$ is a CM lattice of the host,
and there is an automorphism $s$ of $E$ over $K$ with $\rho\circ\tau=\tau\circ s$ for every type $\tau$.
By (2), $F^\rho$ is a theta vector up to units for $(t',L,L')$ formed with $\Lambda^\rho$ and the types
$\rho\circ\tau=\tau\circ s$. By (4) there is a torsion point $y'$ of $\Lambda^\rho_{L,L'}$ with
$(U_yF)^\rho=U^\rho_{y'}F^\rho$. Step 1, with $\Lambda'=\Lambda^\rho$, $t''=t'$, $G=F^\rho$ and $x=y'$, gives
$\|(U_yF)^\rho\|^2_{\Lambda^\rho_{L,L'}}=n(\Lambda^\rho,t')^{fab}$.

\emph{Step 3: $\rho$ complex conjugation on $\varphi_0(K)$.} By (1), $\bar\rho=c\circ\rho$ is the identity
on $\varphi_0(K)$, $\Lambda^\rho=c(\Lambda^{\bar\rho})$ and $(U_yF)^\rho(z)=c\big((U_yF)^{\bar\rho}(c(z))\big)$;
and $c(t')$ is an unkilled two-torsion point of $\Lambda^{\bar\rho}$, because $c(\omega_K)=\bar\omega_K$.
The substitution $z\mapsto c(z)$ of the proof of (2) preserves Lebesgue measure, carries
$\Lambda^{\bar\rho}_{L,L'}$ onto $\Lambda^\rho_{L,L'}$ and $e^{-\pi H^{\bar\rho}(z,z)}$ to $e^{-\pi H^\rho(z,z)}$. So
$\|(U_yF)^\rho\|^2_{\Lambda^\rho_{L,L'}}=\|(U_yF)^{\bar\rho}\|^2_{\Lambda^{\bar\rho}_{L,L'}}$, which is
$n(\Lambda^{\bar\rho},c(t'))^{fab}$ by Step 2, and this is $n(\Lambda^\rho,t')^{fab}$ by
Lemma~\ref{is-vacuum-genus-one}(4).

\emph{Step 4: positivity.} The number $n(\Lambda^\rho,t')$ is the mean of a nonnegative function, so it
suffices to show that it is not $0$. Put $\Lambda''=\Lambda^\rho$ and $t''=t'$ if $\rho$ is the identity on
$\varphi_0(K)$, and $\Lambda''=\Lambda^{\bar\rho}$ and $t''=c(t')$ otherwise; then
$n(\Lambda^\rho,t')=n(\Lambda'',t'')$, by Lemma~\ref{is-vacuum-genus-one}(4) in the second case, and
$\Lambda''$ is a CM lattice of the host by (1). So $\Delta(\Lambda'')=1$ and
$\varphi_0(\mathcal O_K)\Lambda''\subset\Lambda''$. If $t''$ lay in $\Lambda''$, so would $\omega_Kt''$, which
is excluded; so $t''\notin\Lambda''$. The Weierstrass function $\sigma_{\Lambda''}$ vanishes on
$\Lambda''$ and nowhere else, so $(U^{\Lambda''}_{t''}\sigma^*_{\Lambda''})(0)=\sigma^*_{\Lambda''}(t'')e^{-\pi H_{\Lambda''}(t'',t'')/2}\ne0$.
By Lemma~\ref{is-growth-theta-odd-norm}, $n(\Lambda'',t'')$ is the quotient of $n_{\Lambda''}(\sigma^*_{\Lambda''})$,
whose square is $2\pi^2/A(\Lambda'')>0$, by $|(U^{\Lambda''}_{t''}\sigma^*_{\Lambda''})(0)|^2$; so it is not $0$.
% Step 1 is GlobalTheory.ConjNorm.norm_core (exists_selfDual_pair_galoisConjugate_homLattice,
% orthonormalBasis_galois_map, halfPeriodTheta_differentiable_and_signedTranslate_invariant,
% orthonormalProductTheta_mem_thetaSpace, meanSquare_orthoProductTheta_eq_prod, unitary_all); Step 2 is
% norm_fix (conjugate_curveDatum_lattice_isCurveLattice, exists_aut, thetaVacua_conjugates_cmPoints,
% galois_data_exists, thetaTranslate_galoisRationality_cmPoints); Step 3 is norm_big
% (cmTheta_complexConjugation_compatible, meanSquare_evenTheta_conjugateLattice); the case y = 0 is
% useNorm_all with U_zero. Step 4 is ConjJet.thetaZeroNormSq_ne_zero_cmLattice.

(4) Let $m_0\geq1$ with $m_0y\in\Lambda_{L,L'}$, and put $n=2m_0$. By the definition of $\Lambda_{L,L'}$
(Definition~\ref{is-vacuum-theta-vectors}(3)), $m_0y=\sum_jw_j\iota(B_j)$ with $w_j\in\Lambda$ and $B_j$ in the
set $Y$ of the $sB$ with $s\in\mathfrak A_0$ and $BL\subset L'$; so $y=\sum_jc_j\iota(B_j)$ with $c_j=w_j/m_0$. By
Lemma~\ref{is-growth-theta-genus-one}(2b), applied to an automorphism of $\mathbf C$ extending $\rho$, choose
torsion points $t'_j$ of $\Lambda^\rho$ corresponding to $c_j/n$, and put
$y'=\sum_jnt'_j\iota^\rho(B_j)$. Then $c_j/n$ corresponds to $c'_j/n$ with $c'_j=nt'_j$, for every $j$:
$y$ and $y'$ are \emph{adapted at level $n$ along the $B_j$}. Since $m_0n(c_j/n)=w_j\in\Lambda$, also
$m_0nt'_j\in\Lambda^\rho$, so $m_0y'=\sum_j(m_0nt'_j)\iota^\rho(B_j)\in\Lambda^\rho_{L,L'}$, and $y'$ is a
torsion point of $\Lambda^\rho_{L,L'}$.

It remains to show $(U_yF)^\rho=U^\rho_{y'}F^\rho$. Alpöge proves this identity for every theta function
for $\Lambda_{L,L'}$, by passing to the space $V_{4a,b}$.

\emph{Step 1: an automorphism of $\mathbf C$.} Let $\hat\rho$ be an automorphism of $\mathbf C$ extending
$\rho$; it exists because every embedding of $\overline{\mathbf Q}$ into $\mathbf C$ extends to an automorphism
of $\mathbf C$, and it is the identity on $\varphi_0(K)$. Use the notation of
Lemma~\ref{is-growth-padded-galois} for $\hat\rho$ on $V_{a,b}$ and on $V_{4a,b}$. The numbers $g_2(\Lambda)$,
$g_3(\Lambda)$ and $\tau(e)$ for $e\in E$, the values of $\wp_\Lambda$ and $\wp'_\Lambda$ at torsion points
(Lemma~\ref{is-growth-theta-genus-one}(1)), and the Taylor coefficients at
$0$ of $F$ and of $U_yF$ (Lemma~\ref{is-growth-theta-algebraic}(5)) are algebraic. So
$\Lambda^{\hat\rho}=\Lambda^\rho$, $\hat\rho\circ\tau=\rho\circ\tau$ for every type $\tau$, the correspondences
formed with $\hat\rho$ are those formed with $\rho$, $F^{\hat\rho}=F^\rho$ and
$(U_yF)^{\hat\rho}=(U_yF)^\rho$; it suffices to show $(U_yF)^{\hat\rho}=U^{\hat\rho}_{y'}F^{\hat\rho}$. For
$z\in V_{a,b}$ and $B\in M_{b\times a}(E)$ let $\tilde z\in V_{4a,b}$ and $\tilde B\in M_{b\times4a}(E)$ be
obtained by appending $3a$ zero columns; then $\iota(\tilde B)=\widetilde{\iota(B)}$ and
$\iota^{\hat\rho}(\tilde B)=\widetilde{\iota^{\hat\rho}(B)}$.

\emph{Step 2: the padded space.} Let $\nu$ be an orthonormal basis of $M_{b\times4a}(E)$
(Lemma~\ref{is-growth-padded-basis}) and $S=(\sigma^*_\Lambda)_\nu$. By
Lemma~\ref{is-growth-padded-restriction}(3), with $a'=4a$, $F(z)=G_1(\tilde z)$ for a finite sum
$G_1=\sum_i\lambda_iU_{p_i}S$ with $\lambda_i\in\mathbf C$; fix $m_1\geq1$ with $m_1p_i\in\Lambda_\nu$ for every
$i$. Let $R$ be the integer of Lemma~\ref{is-growth-padded-restriction}(1), and $R'$ that of
Lemma~\ref{is-growth-padded-galois}(1) for $\nu$ and the matrices $\tilde B_j$.

\emph{Step 3: a deeper level.} Each $c_j$ is a torsion point of $\Lambda$, as $m_0c_j=w_j$. By
Lemma~\ref{is-growth-padded-galois}(2), with $n_0=n$ and $n_1=Rm_1R'$, there are $\omega_j\in\Lambda$ with
$(c_j+n\omega_j)/nRm_1R'$ corresponding to $c'_j/nRm_1R'$. Put $l=\sum_j\omega_j\iota(B_j)$, an element of
$\Lambda_{L,L'}$, and $x=y+nl=\sum_j(c_j+n\omega_j)\iota(B_j)$. By Lemma~\ref{is-growth-period-lattice}(1)
and (4), Lemma~\ref{is-growth-padded-average}(2) applies to $\Gamma=\Lambda_{L,L'}$; as $n=2m_0$, it gives
$U_xF=U_yF$.

\emph{Step 4: correspondence along $\nu$.} Put $m_2=m_0Rm_1$, so that $nRm_1=2m_2$. Now
$\tilde x=\sum_j(c_j+n\omega_j)\iota(\tilde B_j)$ and $\tilde y'=\sum_jc'_j\iota^{\hat\rho}(\tilde B_j)$, and the
$c_j+n\omega_j$ are torsion points of $\Lambda$. By the choice of $R'$ and Step 3, $\tilde x$ corresponds to
$\tilde y'$ at level $2m_2$. Moreover $m_0x=m_0y+m_0nl\in\Lambda_{L,L'}$, so
$m_2\tilde x=m_1\cdot R\,\widetilde{m_0x}\in\Lambda_\nu$; and $m_2p_i=m_0R\cdot m_1p_i\in\Lambda_\nu$.

\emph{Step 5: conclusion.} By Lemma~\ref{is-growth-padded-galois}(3) with $m=m_2$, $G_1^{\hat\rho}$ is entire and
$(U_{\tilde x}G_1)^{\hat\rho}=U^{\hat\rho}_{\tilde y'}G_1^{\hat\rho}$. As in $(\ast)$ in the proof of
Lemma~\ref{is-growth-padded-restriction}, $(U_xF)(z)=(U_{\tilde x}G_1)(\tilde z)$; and
$(U^{\hat\rho}_{y'}\Phi)(z)=(U^{\hat\rho}_{\tilde y'}\Phi_1)(\tilde z)$ whenever $\Phi(z)=\Phi_1(\tilde z)$, because
$H^{\hat\rho}(\tilde z,\tilde w)=H^{\hat\rho}(z,w)$. By Lemma~\ref{is-growth-padded-galois}(4), applied to $G_1$
and to $U_{\tilde x}G_1$, $F^{\hat\rho}(z)=G_1^{\hat\rho}(\tilde z)$ and
$(U_xF)^{\hat\rho}(z)=(U_{\tilde x}G_1)^{\hat\rho}(\tilde z)$. Hence
\[
 (U_yF)^{\hat\rho}(z)=(U_xF)^{\hat\rho}(z)=\big(U^{\hat\rho}_{\tilde y'}G_1^{\hat\rho}\big)(\tilde z)
 =\big(U^{\hat\rho}_{y'}F^{\hat\rho}\big)(z).
\]
% The choice of y' is GlobalTheory.GalExist.galEx_all (galEx_of_galStar: exists_sum_of_mem_La, the first and
% sixth clauses of galStar), as used in ConjTranslate.conj_fix and ConjNorm.norm_fix. The identity is
% GlobalTheory.GalTheta.gal_all with one term, coefficient 1 and translate 0. Step 1 is gal_all (an
% automorphism agreeing with rho on the data, T.exists_complexEndomorphism_eq_on_countable); Steps 2 to 5 are
% gal_top: Step 2 is exists_orthonormal_frame_var, exists_restriction,
% exists_nsmul_padded_lattice_mem_basis_product and adapted_point_change_of_rational_basis; Step 3 is
% adapted_torsion_lift_deeper_level and thetaTranslation_comb_periodic_twice_level (hp1); Step 4 is hy, e1, e2
% and hAd; Step 5 is galois_conj_torsionTranslates_oddThetaProduct (galB on the padded space), U_pad and
% taylorConjugate_tangentSpacePadZeroColumns.
\end{proof}

\begin{remark}\label{is-growth-theta-conjugate-direct}
Three steps of the proof of Lemma~\ref{is-growth-theta-conjugate} have shorter arguments that Alpöge
does not use. With the notation of that proof:
\begin{enumerate}
\item Statement (4) follows directly along the orthonormal basis $v$ of $F$, with no larger space. By
  Lemma~\ref{is-growth-theta-algebraic}(4), $y$ is a torsion point of $\Lambda_v$; enlarging $m$, we may
  assume $my\in\Lambda_v\cap\Lambda_{L,L'}$. Choose $y'$ corresponding to $y$ at level $2mk$
  (Lemma~\ref{is-growth-theta-galois}(1)). Then $y$ corresponds to $y'$ at level $2m$, and
  Lemma~\ref{is-growth-theta-galois}(2), applied to $F=\sum_i\kappa_iU_{x_i+x_t}(\sigma^*_\Lambda)_v$, gives
  $(U_yF)^\rho=U^\rho_{y'}F^\rho$. Moreover $my$ corresponds to $my'$ at level $2mk$, so
  $my'\in\Lambda^\rho_{L,L'}$ by Lemma~\ref{is-growth-theta-galois}(1), and $y'$ is a torsion point of
  $\Lambda^\rho_{L,L'}$.
\item For $y=0$, Step 2 of the proof of (3) does not need (4): $F^\rho$ itself is a theta vector up to
  units formed with the types $\tau\circ s$, and Step 1 with $x=0$ gives its norm.
\item The positivity of $n(\Lambda^\rho,t')$ in Step 4 also follows without
  Lemma~\ref{is-growth-theta-odd-norm}: the function $|\vartheta_{\Lambda^\rho,t'}(z)|^2e^{-\pi|z|^2/A(\Lambda^\rho)}$
  is continuous, nonnegative and equal to $1$ at $0$, and it is periodic for $\Lambda^\rho$, because
  $|U^{\Lambda^\rho}_\omega\vartheta_{\Lambda^\rho,t'}|=|\vartheta_{\Lambda^\rho,t'}|$ for $\omega\in\Lambda^\rho$
  (proof of Lemma~\ref{pd-analytic-translation-system}(2) and (4)); so its mean over a fundamental domain
  is positive.
\end{enumerate}
\end{remark}

\begin{definition}\label{is-growth-coset-split}
Let $W$ be a finite-dimensional complex vector space with a Hermitian form $H$, form the operators $U_x$
for $x\in W$ with $H$ as in Definition~\ref{pd-theta-function}(2), and let $\alpha$ be a function on $W$.
Self-duality for $H$ is that of Definition~\ref{pd-theta-function}(1), and semicharacters are those of
Lemma~\ref{is-vacuum-units-intertwiners}.
\begin{enumerate}
\item Subgroups $A$, $B$ and $\Gamma$ of $W$ form a \emph{split triple} for $(H,\alpha)$ if $A$ and $B$ are
  self-dual for $H$; $\alpha$ is a semicharacter on $A$, on $B$ and on $\Gamma$;
  $\Gamma\subset(\Gamma\cap A)+(\Gamma\cap B)$; $a\in\Gamma$ whenever $a\in A$, $b\in B$ and $a+b\in\Gamma$;
  every $a\in A$ with $\operatorname{Im}H(a,b)\in\mathbf Z$ for all $b\in B\cap\Gamma$ lies in $\Gamma$; and
  every $b\in B$ with $\operatorname{Im}H(a,b)\in\mathbf Z$ for all $a\in A\cap\Gamma$ lies in $\Gamma$.
\item Let $A$, $B$ and $\Gamma$ be subgroups of $W$, and let $S_X\subset A$, $S_A\subset\Gamma\cap A$ and
  $S_B\subset B$ be finite sets of representatives of $A/(A\cap\Gamma)$, of
  $(\Gamma\cap A)/(\Gamma\cap A\cap B)$ and of $B/(B\cap\Gamma)$. Let $v>0$ be a real number and $\theta$ a
  function on $W$. For a function $\Phi$ on $W$ and $s,y\in W$ put
  \[
   \hat\Phi(y)=v^{-1/2}\,\frac v{|S_A|}\sum_{t\in S_A}\Phi(t)\,e^{2\pi i\operatorname{Im}H(t,y)},\qquad
   \xi(\Phi)=\sum_{y\in S_B}\hat\Phi(y)\,\alpha(y)\,U_y\theta,
  \]
  \[
   \xi_s=\alpha(s)\,U_{-s}\,\xi(\mathbf 1_B),\qquad
   \Pi(\Phi)=\sum_{x\in S_X}\alpha(x)\,U_x\,\xi\big(\Phi(\cdot-x)\big),
  \]
  where $\mathbf 1_B$ is the indicator function of $B$ and $\Phi(\cdot-x)$ is the function
  $t\mapsto\Phi(t-x)$.
\end{enumerate}
\end{definition}
\leanref{Formalism/LocalToGlobal/KernelLattices/XiIndCClass.lean}{Formalism.L2.Ker.SplitHyp}
\leanref{Formalism/LocalToGlobal/PairingIndexBounds/ModelLatticeMoves/GoodLat.lean}{Formalism.L2.PairingIndexBounds.Mv.GoodLat}
\leanref{Formalism/LocalToGlobal/TorsionIntegrality/QuadraticSemicharacters/IsSemiChar.lean}{Formalism.L2.Th1.H.IsReps}
\leanref{Formalism/LocalToGlobal/KernelLattices/XiIndCClass.lean}{Formalism.L2.Ker.hatS}
\leanref{Formalism/LocalToGlobal/KernelLattices/XiIndCClass.lean}{Formalism.L2.Ker.xiS}
\leanref{Formalism/LocalToGlobal/KernelLattices/XiIndCClass.lean}{Formalism.L2.Ker.xiInd}
\leanref{Formalism/LocalToGlobal/KernelLattices/XiIndCClass.lean}{Formalism.L2.Ker.xiL}

\begin{lemma}\label{is-growth-coset-identities}
Let $W$ be a finite-dimensional complex vector space with a Hermitian form $H$, $\alpha$ a function on $W$,
and $(A,B,\Gamma)$ a split triple for $(H,\alpha)$ (Definition~\ref{is-growth-coset-split}). Let $S_X$,
$S_A$, $S_B$, $v>0$ and $\theta$ be as in Definition~\ref{is-growth-coset-split}(2), with
$\alpha(y)U_y\theta=\theta$ for every $y\in\Gamma$, and form $\hat\Phi$, $\xi(\Phi)$, $\xi_s$ and $\Pi(\Phi)$
there.
\begin{enumerate}
\item For $y\in B$, $\hat{\mathbf 1}_B(y)=v^{-1/2}v/|S_A|$. Hence
  $\xi(\mathbf 1_B)=v^{-1/2}\frac v{|S_A|}\sum_{y\in S_B}\alpha(y)U_y\theta$.
\item Let $A'\subset W$ be a subgroup on which $\alpha$ is a semicharacter, and let $s,m\in A'$ with
  $m\in B$. Then $\xi_{s+m}=\xi_s$.
\item For $x,s\in A$, $\alpha(x)U_x\xi_s=\xi_{s-x}$.
\item For every function $\Phi$ on $W$, $\xi(\Phi)=\sum_{s\in S_A}\Phi(s)\,\xi_s$.
\item Let $\Phi$ be a function on $W$ with $\Phi(t+m)=\Phi(t)$ for $t\in A$ and $m\in A\cap B$, and let
  $T\subset A$ be a finite set of representatives of $A/(A\cap B)$. Then
  $\Pi(\Phi)=\sum_{t\in T}\Phi(t)\,\xi_t$.
\item Let $(A_1,B_1,\Gamma)$ be a second split triple for $(H,\alpha)$, with finite sets of representatives
  $S^1_A\subset\Gamma\cap A_1$ of $(\Gamma\cap A_1)/(\Gamma\cap A_1\cap B_1)$ and $S^1_B\subset B_1$ of
  $B_1/(B_1\cap\Gamma)$, and form $\xi^1(\Phi)$ and $\xi^1_s$ with them, $v$ and $\theta$. Suppose that
  $A\cap B_1\subset\Gamma$, $B_1\subset(A\cap B_1)+(B\cap B_1)$ and $\Gamma\cap A\subset A_1$, and that for
  every $r\in\Gamma\cap A_1$ there is $u\in\Gamma\cap A$ with $u-r\in B_1$. Then, for $T$ as in (5) and
  every $s\in A$,
  \[
   \xi^1_s=\sum_{t\in T}\mathbf 1_{B_1}(t-s)\,\xi_t .
  \]
\end{enumerate}
\end{lemma}
\leanref{Formalism/LocalToGlobal/KernelLattices/KernelLocalFactors/AlternatingFormStandardHermitian.lean}{Formalism.L2.Ker.twist_induced_vanish}
\leanref{Formalism/LocalToGlobal/KernelLattices/KernelLocalFactors/ThDistC.lean}{Formalism.L2.Ker.lf_class}
\leanref{Formalism/LocalToGlobal/KernelLattices/KernelLocalFactors/AlternatingFormStandardHermitian.lean}{Formalism.L2.Ker.xiInd_shift}
\leanref{Formalism/LocalToGlobal/KernelLattices/KernelLocalFactors/AlternatingFormStandardHermitian.lean}{Formalism.L2.Ker.twist_expand_series}
\leanref{Formalism/LocalToGlobal/KernelLattices/KernelLocalFactors/AlternatingFormStandardHermitian.lean}{Formalism.L2.Ker.twist_expand_lattice}
\leanref{Formalism/LocalToGlobal/KernelLattices/KernelLocalFactors/ThDistC.lean}{Formalism.L2.Ker.rf_abs}
\begin{proof}
Write $E(x,y)=\operatorname{Im}H(x,y)$, so that $E(y,x)=-E(x,y)$. By
Lemma~\ref{is-vacuum-units-intertwiners}(1), $U_xU_y=e^{-\pi iE(x,y)}U_{x+y}$, so
$U_xU_y=e^{-2\pi iE(x,y)}U_yU_x$. Let $\alpha$ be a semicharacter on a subgroup $\Gamma'$. By
Lemma~\ref{is-vacuum-units-intertwiners}(2) and its proof, $E$ is integral on $\Gamma'\times\Gamma'$,
$\alpha(x)U_x(\alpha(y)U_yF)=\alpha(x+y)U_{x+y}F$ for $x,y\in\Gamma'$ and every function $F$, and if
$\alpha(y)U_yF=F$ for $y$ in a subgroup $\Gamma''$, then $\alpha(x)U_xF$, for $x\in\Gamma'$, depends only
on the class of $x$ modulo $\Gamma'\cap\Gamma''$. Moreover $\alpha(0)=1$ and $\alpha(-x)=\alpha(x)$ for
$x\in\Gamma'$, because $\alpha(0)=\alpha(0)^2$, $\alpha(x)\alpha(-x)=\alpha(0)$ and the values are $\pm1$.
The sum of a function over a set of representatives does not depend on the set when the function is
invariant under the subgroup. Put $\kappa=v^{-1/2}v/|S_A|$.

(1) Let $y\in B$. The only element of $S_A$ in $B$ is the representative $t_0$ of the class of $0$, and
$E(t_0,y)\in\mathbf Z$ because $\alpha$ is a semicharacter on $B$. So the sum that defines
$\hat{\mathbf 1}_B(y)$ has the single term $1$. The formula for $\xi(\mathbf 1_B)$ follows, as
$S_B\subset B$.

(2) By (1), $\xi(\mathbf 1_B)=\kappa G$ with $G=\sum_{y\in S_B}\alpha(y)U_y\theta$. Let $m'\in B$. Then
$\alpha(m')U_{m'}(\alpha(y)U_y\theta)=\alpha(m'+y)U_{m'+y}\theta$, which depends only on the class of
$m'+y$ modulo $B\cap\Gamma$, and $m'+S_B$ is again a set of representatives of $B/(B\cap\Gamma)$; so
$\alpha(m')U_{m'}G=G$. With $m'=-m$, and the semicharacter $\alpha$ on $A'$,
\[
 \xi_{s+m}=\alpha(-s-m)U_{-s-m}\,\xi(\mathbf 1_B)=\alpha(-s)U_{-s}\big(\alpha(-m)U_{-m}\,\xi(\mathbf 1_B)\big)
 =\alpha(-s)U_{-s}\,\xi(\mathbf 1_B)=\xi_s .
\]

(3) For $x,s\in A$, $\alpha(x)U_x(\alpha(-s)U_{-s}F)=\alpha(x-s)U_{x-s}F$ for every $F$. As
$\alpha(-s)=\alpha(s)$ and $\alpha(x-s)=\alpha(s-x)$, this is $\alpha(x)U_x\xi_s=\xi_{s-x}$ for
$F=\xi(\mathbf 1_B)$.

(4) Let $x\in\Gamma$. Since $\alpha(x)U_x\theta=\theta$, the commutation rule gives
$\alpha(x)U_x(\alpha(y)U_y\theta)=e^{-2\pi iE(x,y)}\alpha(y)U_y\theta$ for every $y\in W$. For
$s\in S_A\subset\Gamma$ this, with $x=-s$, $\alpha(s)=\alpha(-s)$ and (1), gives
\[
 \xi_s=\alpha(-s)U_{-s}\,\xi(\mathbf 1_B)=\kappa\sum_{y\in S_B}e^{2\pi iE(s,y)}\alpha(y)U_y\theta .
\]
Multiplying by $\Phi(s)$ and summing over $s\in S_A$ gives $\sum_{y\in S_B}\hat\Phi(y)\alpha(y)U_y\theta$,
which is $\xi(\Phi)$.

(5) First, $A\cap B\subset\Gamma$: an $m\in A\cap B$ has $E(m,b)\in\mathbf Z$ for every $b\in B$, as $\alpha$
is a semicharacter on $B$, so $m\in\Gamma$ by the definition of a split triple. Put $g(s)=\Phi(s)\,\xi_s$
for $s\in A$. By (2) with $A'=A$, and the hypothesis on $\Phi$, $g(t+m)=g(t)$ for $t\in A$ and
$m\in A\cap B$. For $x\in S_X$, (4) for the function $\Phi(\cdot-x)$ and (3) give
\[
 \alpha(x)U_x\,\xi\big(\Phi(\cdot-x)\big)=\sum_{t\in S_A}\Phi(t-x)\,\alpha(x)U_x\xi_t=\sum_{t\in S_A}g(t-x),
\]
so $\Pi(\Phi)$ is the sum of $g$ over the differences $t-x$ with $x\in S_X$ and $t\in S_A$. These
differences represent the classes of $A/(A\cap B)$, each once. If $(t-x)-(t'-x')\in B$, this element
lies in $A\cap B\subset\Gamma$, so $x-x'\in\Gamma$ and $x=x'$, and then $t-t'\in B$ gives $t=t'$. Every
$y\in A$ is congruent modulo $B$ to one of them: choose $x\in S_X$ with $y+x\in\Gamma$, and then
$t\in S_A$ with $y+x-t\in B$. Hence $\Pi(\Phi)=\sum_{t\in T}g(t)$.

(6) (a) If $y\in B$ and $E(t,y)\in\mathbf Z$ for every $t\in A\cap B_1$, then $y\in B_1$. Indeed $B_1$ is
self-dual, and every $w\in B_1$ is $w_1+w_2$ with $w_1\in A\cap B_1$ and $w_2\in B\cap B_1$, where
$E(y,w_1)=-E(w_1,y)\in\mathbf Z$, and $E(y,w_2)\in\mathbf Z$ because $\alpha$ is a semicharacter on $B$. As
$E$ is integral on $A$ and on $\Gamma$, and $A\cap B_1\subset\Gamma$, it follows that $A\cap B\subset B_1$
and $\Gamma\cap B\subset B_1$.

(b) Let $n$ be the number of elements of $S_A$ in $B_1$. These elements represent
$(\Gamma\cap A\cap B_1)/(\Gamma\cap A\cap B_1\cap B)$, because $\Gamma\cap A\cap B\subset B_1$ by (a). For
$y\in B$,
\[
 \sum_{t\in S_A}\mathbf 1_{B_1}(t)\,e^{2\pi iE(t,y)}=n\,\mathbf 1_{B_1}(y).
\]
If $y\in B_1$, every term with $t\in B_1$ is $1$, as $E$ is integral on $B_1$. If $y\notin B_1$, then by
(a) there is $t_0\in A\cap B_1\subset\Gamma\cap A\cap B_1$ with $E(t_0,y)\notin\mathbf Z$. The character
$t\mapsto e^{2\pi iE(t,y)}$ of $\Gamma\cap A\cap B_1$ is trivial on its intersection with $B$, as $E$ is
integral on $B$, and not at $t_0$, so its sum over a set of representatives is $0$.

(c) $|S_A|=n\,|S^1_A|$. For each $r\in S^1_A$ choose $u_r\in\Gamma\cap A$ with $u_r-r\in B_1$. Each
$t\in S_A$ lies in $\Gamma\cap A\subset\Gamma\cap A_1$, so it is congruent modulo $B_1$ to exactly one
$r\in S^1_A$, and then $t-u_r\in B_1$. For a fixed $r$ the number of $t\in S_A$ with $t-u_r\in B_1$ is
$n$, because translation by $-u_r\in\Gamma\cap A$ permutes the classes of
$(\Gamma\cap A)/(\Gamma\cap A\cap B)$ and $\mathbf 1_{B_1}$ is invariant under
$\Gamma\cap A\cap B\subset B_1$. Counting the pairs $(t,r)$ with $t-u_r\in B_1$ in two ways gives the
claim.

(d) The elements of $S_B$ in $B_1$ represent $B_1/(B_1\cap\Gamma)$. They represent
$(B\cap B_1)/(B\cap B_1\cap\Gamma)$, because $\Gamma\cap B\subset B_1$; and every $w\in B_1$ is $w_1+w_2$
with $w_1\in A\cap B_1\subset\Gamma$ and $w_2\in B\cap B_1$.

(e) By (b) and (c), then (d), and then (1) for the triple $(A_1,B_1,\Gamma)$,
\[
 \xi(\mathbf 1_{B_1})=\sum_{y\in S_B}v^{-1/2}\frac{v\,n}{|S_A|}\mathbf 1_{B_1}(y)\,\alpha(y)U_y\theta
 =v^{-1/2}\frac v{|S^1_A|}\sum_{y\in S_B\cap B_1}\alpha(y)U_y\theta
 =v^{-1/2}\frac v{|S^1_A|}\sum_{y\in S^1_B}\alpha(y)U_y\theta=\xi^1(\mathbf 1_{B_1}),
\]
where the third step holds because $\alpha(y)U_y\theta$, for $y\in B_1$, depends only on the class of $y$
modulo $B_1\cap\Gamma$.

(f) By (4) for $\Phi=\mathbf 1_{B_1}$, $\xi(\mathbf 1_{B_1})$ is the sum of the $\xi_s$ over the $s\in S_A$
in $B_1$. These represent $(A\cap B_1)/(A\cap B_1\cap B)$, by (b) and $\Gamma\cap A\cap B_1=A\cap B_1$; so
do the $t\in T$ in $B_1$, because $A\cap B\subset B_1$; and $\xi_t$ depends only on the class of $t$
modulo $A\cap B$, by (2). With (e), $\sum_{t\in T}\mathbf 1_{B_1}(t)\,\xi_t=\xi^1(\mathbf 1_{B_1})$.

(g) As $A\cap B\subset B_1$, the function $t\mapsto\mathbf 1_{B_1}(t-s)\,\xi_t$ on $A$ is invariant under
$A\cap B$, and $s+T$ represents $A/(A\cap B)$. So, by (3) with $x=-s$ and by (f),
\[
 \sum_{t\in T}\mathbf 1_{B_1}(t-s)\,\xi_t=\sum_{t\in T}\mathbf 1_{B_1}(t)\,\xi_{s+t}
 =\alpha(-s)U_{-s}\sum_{t\in T}\mathbf 1_{B_1}(t)\,\xi_t=\alpha(s)U_{-s}\,\xi^1(\mathbf 1_{B_1})=\xi^1_s .
\]
\end{proof}

\begin{definition}\label{is-growth-coset-lift}
Let $D$ be a definite-side datum with kernel data (Definition~\ref{is-kernel-ties}), with host field $E$,
the ideal $\mathfrak b$ of the host and the CM lattice $\Lambda=\delta\varphi_0(\mathfrak c)$ of the host
given by the kernel data, and form $V_{a,b}$, $H$, $U_x$, $\alpha=\alpha_{\sqrt d/2}$ and $\iota$ as in
Definition~\ref{is-vacuum-theta-vectors}. Let $N$ and $N'$ be lattices of levels $i$ and $i+1$, of ranks
$a$ and $b$, with the vacuum vector $\theta_{N,N'}$ of the pair, and let $S$ be a finite set of finite
places of $E$ that contains no two complex-conjugate places; write $\bar S$ for the set of the
conjugates of the places of $S$ (Definitions~\ref{is-vacuum-theta-vectors} and~\ref{is-theta-kernel}).
\begin{enumerate}
\item For an $\mathcal O_E$-lattice $M$ in a finite-dimensional $E$-vector space and an integer $g\geq1$,
  $M_A(g)$ is the set of $x$ with $x\in g^{-1}M_Q$ for $Q\in S$, $x\in gM_Q$ for $Q\in\bar S$ and $x\in M_Q$
  for the other finite places $Q$, where $M_Q$ is the localization of $M$ at $Q$; $M_B(g)$ is defined in
  the same way with $g$ and $g^{-1}$ exchanged.
\item For $\mathcal O_E$-lattices $L\subset E^a$ and $L'\subset E^b$, $P(L,L')$ is the $\mathcal O_E$-span of
  the $sB$ with $s\in\mathfrak c\mathfrak b^{-1}$ and $B\in\operatorname{Hom}_E(E^a,E^b)$ with $BL\subset L'$.
  Put $P=P(N,N')$, and let $\Gamma$, $A_g$ and $B_g$ be the images of $P$, $P_A(g)$ and $P_B(g)$ under the
  injective map $x\mapsto\delta\iota(x)$. Then $\Gamma=\Lambda_{N,N'}$, because every type restricts to
  $\varphi_0$ on $K$, so that $\delta\varphi_0(k)\iota(x)=\delta\iota(kx)$ for $k\in\mathfrak c$.
\item For $g\geq1$, let $S_X(g)$, $S_A(g)$ and $S_B(g)$ be chosen finite sets of representatives for
  $(A_g,B_g,\Gamma)$ as in Definition~\ref{is-growth-coset-split}(2) (the indices are finite by
  Lemma~\ref{is-growth-coset-triples}(2)), let $v>0$ be the covolume of
  Definition~\ref{is-theta-kernel}(2), and form $\xi^{(g)}(\Phi)$, $\xi^{(g)}_p$ and $\Pi^{(g)}(\Phi)$ with
  them, $v$ and $\theta=\theta_{N,N'}$.
\item For a function $\Phi$ on $\operatorname{Hom}_E(E^a,E^b)$ let $\tilde\Phi$ be the function on $V_{a,b}$
  with $\tilde\Phi(\delta\iota(t))=\Phi(t)$ that vanishes off the image of $x\mapsto\delta\iota(x)$. The
  \emph{Heisenberg lift} of $\Phi$ at the modulus $g$ is $\Pi_g(\Phi)=\Pi^{(g)}(\tilde\Phi)$; it is linear in
  $\Phi$.
\item Let $c\geq1$ be an integer and $x\in\operatorname{Hom}_E(E^a,E^b)$. A \emph{split representative} of
  $x$ at $c$ is an $s\in\operatorname{Hom}_E(E^a,E^b)$ with $s\in P_A(c_1)$ for some integer $c_1\geq1$ and
  $s-x\in cH_S(N,N')$. By its construction, the coset theta function of
  Definition~\ref{is-theta-kernel}(2) is $\Theta_{N,N',c,x}=v^{-1/2}\,\xi^{(c)}_{\delta\iota(s)}$ for a split
  representative $s$ of $x$ at $c$ chosen with $c$ and $x$, when $x$ has one.
\end{enumerate}
\end{definition}
\leanref{Formalism/LocalToGlobal/KernelLattices/XiIndCClass.lean}{Formalism.L2.Ker.twist₂}
\leanref{Formalism/LocalToGlobal/KernelLattices/XiIndCClass.lean}{Formalism.L2.Ker.halfAO}
\leanref{Formalism/LocalToGlobal/KernelLattices/XiIndCClass.lean}{Formalism.L2.Ker.halfBO}
\leanref{Formalism/LocalToGlobal/KernelLattices/XiIndCClass.lean}{Formalism.L2.Ker.perO}
\leanref{Formalism/LocalToGlobal/KernelLattices/XiIndCClass.lean}{Formalism.L2.Ker.LamC}
\leanref{Formalism/LocalToGlobal/KernelLattices/XiIndCClass.lean}{Formalism.L2.Ker.AC}
\leanref{Formalism/LocalToGlobal/KernelLattices/XiIndCClass.lean}{Formalism.L2.Ker.BC}
\leanref{Formalism/LocalToGlobal/KernelLattices/XiIndCClass.lean}{Formalism.L2.Ker.onV}
\leanref{Formalism/LocalToGlobal/KernelLattices/XiIndCClass.lean}{Formalism.L2.Ker.xiV}
\leanref{Formalism/LocalToGlobal/KernelLattices/XiIndCClass.lean}{Formalism.L2.Ker.xiC}
\leanref{Formalism/LocalToGlobal/KernelLattices/ThDistC.lean}{Formalism.L2.Ker.Amb.IsLift}
\leanref{Formalism/LocalToGlobal/KernelLattices/ThDistC.lean}{Formalism.L2.Ker.Amb.liftW}
\leanref{Formalism/LocalToGlobal/KernelLattices/ThDistC.lean}{Formalism.L2.Ker.cosInd}

\begin{lemma}\label{is-growth-coset-lattices}
Let $E$ be a CM field, $S$ a finite set of finite places of $E$ that contains no two complex-conjugate
places, and $g,c,D_1\geq1$ integers with $c\mid D_1$; form $M_A(g)$ and $M_B(g)$ for $\mathcal O_E$-lattices
$M$ as in Definition~\ref{is-growth-coset-lift}(1). For (4) to (7), let moreover $D$, $\mathfrak c$,
$\mathfrak b$, $N$, $N'$ and $P=P(N,N')$ be as in Definition~\ref{is-growth-coset-lift} for this $S$,
suppose that $\mathfrak c\mathfrak b^{-1}$ is a unit at every place of $S$, and write $H_S=H_S(N,N')$.
\begin{enumerate}
\item $gM\subset M_A(g)\cap M_B(g)$, and $gM_A(g)$ and $gM_B(g)$ lie in $M$.
\item There is $\varepsilon\in\mathcal O_E$ with $1-\varepsilon\in g\mathcal O_{E,Q}$ for $Q\in S$ and
  $\varepsilon\in g\mathcal O_{E,Q}$ for $Q\in\bar S$. Moreover $M\subset(M\cap M_A(g))+(M\cap M_B(g))$, and
  $a\in M$ whenever $a\in M_A(g)$, $b\in M_B(g)$ and $a+b\in M$.
\item $M_A(D_1)\cap M_B(c)\subset M$; $M_B(c)\subset(M_A(D_1)\cap M_B(c))+(M_B(D_1)\cap M_B(c))$;
  $M\cap M_A(D_1)\subset M_A(c)$; and for every $r\in M$ there is $u\in M\cap M_A(D_1)$ with
  $u-r\in M_B(c)$.
\item $N'_A(g)$ and $N'_B(g)$ are self-dual for $\sum_i\bar x_iy_i$, $P_A(g)=P(N,N'_A(g))$ and
  $P_B(g)=P(N,N'_B(g))$.
\item For $Q\in S$, an $x\in\operatorname{Hom}_E(E^a,E^b)$ lies in $P_Q$ if and only if $xN_Q\subset N'_Q$.
  Hence, for $w\in P_A(g)$, $w\in P_B(c)$ if and only if $w\in cH_S$.
\item If $D_1x\in H_S$, there is $s\in P_A(D_1)$ with $s-x\in D_1H_S$.
\item If $s\in P_A(c_1)$ and $s'\in P_A(c_2)$ for integers $c_1,c_2\geq1$, and $s'-s\in cH_S$, then $s$ and
  $s'$ lie in $P_A(g_0)$ for $g_0=\gcd(c_1,c_2)$, and $s'-s\in P_A(g_0)\cap P_B(c)$.
\end{enumerate}
\end{lemma}
\leanref{Formalism/LocalToGlobal/KernelLattices/KernelLocalFactors/ThDistC.lean}{Formalism.L2.Ker.hh_tw_of}
\leanref{Formalism/LocalToGlobal/KernelLattices/KernelLocalFactors/AlternatingFormStandardHermitian.lean}{Formalism.L2.Ker.exists_idem}
\leanref{Formalism/LocalToGlobal/KernelLattices/KernelLocalFactors/AlternatingFormStandardHermitian.lean}{Formalism.L2.Ker.tw_split}
\leanref{Formalism/LocalToGlobal/KernelLattices/KernelLocalFactors/ThDistC.lean}{Formalism.L2.Ker.rf_facts}
\leanref{Formalism/LocalToGlobal/KernelLattices/KernelLocalFactors/AlternatingFormStandardHermitian.lean}{Formalism.L2.Ker.twist_selfDual}
\leanref{Formalism/LocalToGlobal/KernelLattices/KernelLocalFactors/AlternatingFormStandardHermitian.lean}{Formalism.L2.Ker.twist_homomorphism}
\leanref{Formalism/LocalToGlobal/KernelLattices/KernelLocalFactors/ThDistC.lean}{Formalism.L2.Ker.lf_halfB_iff}
\leanref{Formalism/LocalToGlobal/KernelLattices/KernelLocalFactors/ThDistC.lean}{Formalism.L2.Ker.ti_exists_lift}
\leanref{Formalism/LocalToGlobal/KernelLattices/KernelLocalFactors/ThDistC.lean}{Formalism.L2.Ker.lf_gcd}
\begin{proof}
An element lies in an $\mathcal O_E$-lattice $M$ if and only if it lies in $M_Q$ for every finite place
$Q$, so each assertion can be checked one place at a time. At every place, $M_A(g)$ and $M_B(g)$ lie
between $gM_Q$ and $g^{-1}M_Q$; this is (1).

(2) The element $\varepsilon$ exists by the Chinese remainder theorem, as $S$ and $\bar S$ are disjoint.
For $x\in M$, $\varepsilon x\in M_A(g)$ and $(1-\varepsilon)x\in M_B(g)$: at $Q\in S$,
$\varepsilon x\in M_Q$ and $(1-\varepsilon)x\in gM_Q$; at $Q\in\bar S$, $\varepsilon x\in gM_Q$ and
$(1-\varepsilon)x\in M_Q$. If $a\in M_A(g)$, $b\in M_B(g)$ and $a+b\in M$, then at $Q\in S$ we have
$b\in gM_Q\subset M_Q$, so $a\in M_Q$; at $Q\in\bar S$, $a\in gM_Q$; and elsewhere $a\in M_Q$.

(3) An $x\in M_A(D_1)\cap M_B(c)$ lies in $cM_Q\subset M_Q$ for $Q\in S$, in $D_1M_Q\subset M_Q$ for
$Q\in\bar S$ and in $M_Q$ elsewhere, so in $M$. For $y\in M_B(c)$ take $\varepsilon$ as in (2) for $cD_1$:
at $Q\in S$, $y\in cM_Q$ gives $\varepsilon y\in M_Q$ and $(1-\varepsilon)y\in D_1M_Q$; at $Q\in\bar S$,
$y\in c^{-1}M_Q$ gives $\varepsilon y\in D_1M_Q$ and $(1-\varepsilon)y\in D_1^{-1}M_Q$, as $c\mid D_1$. So
$y=\varepsilon y+(1-\varepsilon)y$ with $\varepsilon y\in M_A(D_1)$ and $(1-\varepsilon)y\in M_B(D_1)$, and
both lie in the $\mathcal O_E$-module $M_B(c)$. An
$x\in M\cap M_A(D_1)$ lies in $M_Q\subset c^{-1}M_Q$ for $Q\in S$ and in $D_1M_Q\subset cM_Q$ for
$Q\in\bar S$. Finally, for $r\in M$ take $\varepsilon$ as in (2) for $D_1$ and $u=\varepsilon r$; then
$u\in M\cap M_A(D_1)$ by (2), and $u-r=-(1-\varepsilon)r$ lies in $D_1M_Q\subset cM_Q$ for $Q\in S$ and in
$M_Q$ elsewhere, so in $M_B(c)$.

(4) Let $e(Q)$ be $1$, $-1$ and $0$ for $Q$ in $S$, in $\bar S$ and elsewhere, so that
$M_A(g)_Q=g^{-e(Q)}M_Q$, $M_B(g)_Q=g^{e(Q)}M_Q$ and $e(\bar Q)=-e(Q)$. The localization at $Q$ of the dual
of a lattice $L$ for $\sum_i\bar x_iy_i$ is determined by $L_{\bar Q}$, and multiplying $L_{\bar Q}$ by
the rational number $g^{k}$ multiplies it by $g^{-k}$. As $N'$ is self-dual and $e(\bar Q)=-e(Q)$, the
duals of $N'_A(g)$ and $N'_B(g)$ are $N'_A(g)$ and $N'_B(g)$. For $P_A(g)=P(N,N'_A(g))$, compare
localizations at a place $Q$. If $BN\subset N'_A(g)$, then $g^{e(Q)}BN_Q\subset N'_Q$, and as $N$ is
finitely generated there is $t\in\mathcal O_E\smallsetminus Q$ with $tg^{e(Q)}BN\subset N'$; so every
generator $sB$ of $P(N,N'_A(g))$ lies in $g^{-e(Q)}P_Q$. Conversely, the Chinese remainder theorem gives
$t\in\mathcal O_E\smallsetminus Q$ with $tg^{e(Q')-e(Q)}$ integral at every place $Q'$; then
$tg^{-e(Q)}BN\subset N'_A(g)$ whenever $BN\subset N'$, so $tg^{-e(Q)}P\subset P(N,N'_A(g))$, and every
$z\in P_A(g)$, which has $g^{e(Q)}z\in P_Q$, lies in $P(N,N'_A(g))_Q$. The case of $P_B(g)$ is the same with
$-e$ in place of $e$.

(5) Let $Q\in S$. As $\mathfrak c\mathfrak b^{-1}$ is a unit at $Q$, $P_Q$ is the
$\mathcal O_{E,Q}$-span of the $B$ with $BN\subset N'$; as $N$ is finitely generated, this is the set of
$x$ with $xN_Q\subset N'_Q$. Let $w\in P_A(g)$. If $w\in P_B(c)$, then $c^{-1}w\in P_Q$ for $Q\in S$, so
$c^{-1}w\in H_S$ and $w\in cH_S$. If $w\in cH_S$, then $w\in cP_Q$ for $Q\in S$, $w\in gP_Q\subset c^{-1}P_Q$ for
$Q\in\bar S$, and $w\in P_Q$ elsewhere; so $w\in P_B(c)$.

(6) Choose $n\geq1$ with $nx\in P$; it exists as $\mathfrak c\mathfrak b^{-1}\neq0$, $N$ is finitely
generated and $N'$ spans $E^b$. By the Chinese remainder theorem choose $\varepsilon\in\mathcal O_E$ with
$1-\varepsilon\in D_1^2\mathcal O_{E,Q}$ for $Q\in S$ and $\varepsilon\in D_1n\mathcal O_{E,Q}$ for the
places $Q\notin S$ that divide $D_1n$, and put $s=\varepsilon x$. At $Q\in S$, $D_1x\in P_Q$ by (5), so
$D_1s\in P_Q$ and $D_1^{-1}(s-x)=-D_1^{-2}(1-\varepsilon)\,D_1x\in P_Q$. At $Q\notin S$,
$D_1^{-1}s=\varepsilon(D_1n)^{-1}\,nx\in P_Q$. So $s\in P_A(D_1)$, and $D_1^{-1}(s-x)\in H_S$ by (5).

(7) At $Q\in S$, $c_1s$ and $c_2s'$ lie in $P_Q$, and $s'-s\in cP_Q\subset P_Q$ by (5). Writing
$g_0=k_1c_1+k_2c_2$ with integers $k_1,k_2$, $g_0s=k_1c_1s+k_2c_2s'-k_2c_2(s'-s)\in P_Q$, and in the same
way $g_0s'\in P_Q$. At $Q\in\bar S$, $s\in c_1P_Q\subset g_0P_Q$ and $s'\in c_2P_Q\subset g_0P_Q$; elsewhere
$s,s'\in P_Q$. So $s,s'\in P_A(g_0)$, hence $s'-s\in P_A(g_0)$, and $s'-s\in P_B(c)$ by (5).
\end{proof}

\begin{lemma}\label{is-growth-coset-hermitian}
Let $D$ be a definite-side datum, $d$ the prime of its host field, and $\Lambda=\delta\varphi_0(\mathfrak c)$
a CM lattice of its host, with the notation of Definition~\ref{is-growth-theta-functions}; write
$\alpha=\alpha_{\sqrt d/2}$.
\begin{enumerate}
\item Let $L\subset E^a$ and $L'\subset E^b$ be lattices that are self-dual for $\sum_i\bar x_iy_i$. For all
  $x,y\in\Lambda_{L,L'}$ there are integers $r$ and $s$ with $r-s$ even and $H(x,y)=r/\sqrt d+si$.
\item Let $\Gamma\subset V_{a,b}$ be a subgroup such that for all $x,y\in\Gamma$ there are integers $r$ and
  $s$ with $r-s$ even and $H(x,y)=r/\sqrt d+si$. Then $\alpha$ is a semicharacter on $\Gamma$
  (Lemma~\ref{is-vacuum-units-intertwiners}).
\end{enumerate}
\end{lemma}
\leanref{Formalism/LocalToGlobal/GlobalTheory/GL3/PThIndep.lean}{Formalism.L2.GlobalTheoryTheory.IntegralLatticeDensity.hermitian_form_integral_values}
\leanref{Formalism/LocalToGlobal/GlobalTheory/GL3/PIopIso.lean}{Formalism.L2.GlobalTheoryTheory.IntegralLatticeDensity.IO.HfInt}
\leanref{Formalism/LocalToGlobal/GlobalTheory/GL3/PIopIso.lean}{Formalism.L2.GlobalTheoryTheory.IntegralLatticeDensity.IO.semiChar_of_HfInt}
\leanref{Formalism/LocalToGlobal/GlobalTheory/GL3/PHfGen.lean}{Formalism.L2.GlobalTheoryTheory.IntegralLatticeDensity.HfIntRec_of_gen}
\leanref{Formalism/LocalToGlobal/GlobalTheory/GL3/PHfArith.lean}{Formalism.L2.GlobalTheoryTheory.IntegralLatticeDensity.HfArithRec_of_int}
\leanref{Formalism/LocalToGlobal/GlobalTheory/GL3/PHfArith.lean}{Formalism.L2.GlobalTheoryTheory.IntegralLatticeDensity.AR.div_okc}
\begin{proof}
Let $G\subset\mathbf C$ be the set of the numbers $r/\sqrt d+si$ with integers $r,s$ and $r-s$ even; it is
an additive subgroup of $\mathbf C$.

(1) \emph{Step 1: the integers of $K$.} As $d\equiv3\pmod4$, $\mathcal O_K=\mathbf Z+\mathbf Z\omega_K$, and
$\varphi_0(\sqrt{-d})=i\sqrt d$. So $\varphi_0(\mathcal O_K)$ is the set of the numbers $(r+si\sqrt d)/2$
with integers $r,s$ and $r-s$ even. It is a ring, stable under complex conjugation, and $2u/\sqrt d\in G$
for every $u\in\varphi_0(\mathcal O_K)$.

\emph{Step 2: the covolume.} Put $e(w,w')=\operatorname{Im}(\bar ww')$ for $w,w'\in\mathbf C$, and
$\kappa=|\delta|^2\sqrt d/2$. Every element of $\Lambda$ is $\delta\varphi_0(c)$ with $c\in\mathfrak c$, and
$e(\delta\varphi_0(c),\delta\varphi_0(c'))=|\delta|^2\operatorname{Im}\big(\overline{\varphi_0(c)}\varphi_0(c')\big)$
lies in $\kappa\mathbf Z$ by Step 1. Let $\omega_1,\omega_2$ be a basis of $\Lambda$, so that
$A(\Lambda)=|e(\omega_1,\omega_2)|$ (Definition~\ref{pd-theta-function}). As $e$ is biadditive and
$e(w,w)=0$, $e(w,w')=(kl'-lk')e(\omega_1,\omega_2)$ for $w=k\omega_1+l\omega_2$ and
$w'=k'\omega_1+l'\omega_2$. So every value of $e$ on $\Lambda\times\Lambda$ lies in $A(\Lambda)\mathbf Z$, and
$A(\Lambda)=\kappa m$ for an integer $m\geq1$.

\emph{Step 3: the quotients.} Let $c,c'\in\mathfrak c$, and write
$z=\overline{\varphi_0(c')}\varphi_0(c)=(r+si\sqrt d)/2$ as in Step 1. Then
$e(\delta\varphi_0(c'),\delta\varphi_0(c))=\kappa s$. As $c\omega_K\in\mathfrak c$ and
$z\varphi_0(\omega_K)=\big((r-ds)+(r+s)i\sqrt d\big)/4$,
$e(\delta\varphi_0(c'),\delta\varphi_0(c\omega_K))=\kappa(r+s)/2$. By Step 2 both lie in $\kappa m\mathbf Z$,
so $s=mn_1$ and $r+s=2mn_2$ with integers $n_1,n_2$. Hence $z/m=\big((2n_2-n_1)+n_1i\sqrt d\big)/2$, which
lies in $\varphi_0(\mathcal O_K)$, as $(2n_2-n_1)-n_1$ is even.

\emph{Step 4: the values.} Let $Y$ and $h$ be as in the proof of Lemma~\ref{is-growth-period-lattice}. By
(b) and (c) there, $H(\iota(Z),\iota(Z'))=\varphi_0(h(Z,Z'))/A(\Lambda)$ and $h(Z,Z')\in\mathcal O_K$ for
$Z,Z'\in Y$. For $c,c'\in\mathfrak c$, Step 2 gives
\[
 H\big(\delta\varphi_0(c)\iota(Z),\delta\varphi_0(c')\iota(Z')\big)
 =\frac{|\delta|^2\varphi_0(c)\overline{\varphi_0(c')}\varphi_0(h(Z,Z'))}{\kappa m}=\frac{2u}{\sqrt d},
 \qquad u=\frac{\overline{\varphi_0(c')}\varphi_0(c)}m\,\varphi_0(h(Z,Z')),
\]
and $u\in\varphi_0(\mathcal O_K)$ by Steps 1 and 3; so this value lies in $G$. The group $\Lambda_{L,L'}$ is
generated by the $w\iota(Z)$ with $w\in\Lambda$ and $Z\in Y$ (proof of
Lemma~\ref{is-growth-period-lattice}), that is, by the $\delta\varphi_0(c)\iota(Z)$ with $c\in\mathfrak c$ and
$Z\in Y$. As $H$ is additive in each variable and $G$ is a group, $H(x,y)\in G$ for all
$x,y\in\Lambda_{L,L'}$.

(2) Write $q(x)=\frac{\sqrt d}2\operatorname{Re}H(x,x)$, so that $\alpha(x)=e^{\pi iq(x)}$. Let $x\in\Gamma$ and
$H(x,x)=r/\sqrt d+si$ as in the hypothesis. As $H(x,x)$ is real, $s=0$, so $r$ is even and $q(x)=r/2$ is an
integer; hence $\alpha(x)=\pm1$. Now let $x,y\in\Gamma$ and $H(x,y)=r/\sqrt d+si$. Since
$\operatorname{Re}H(y,x)=\operatorname{Re}H(x,y)$,
$\operatorname{Re}H(x+y,x+y)=\operatorname{Re}H(x,x)+\operatorname{Re}H(y,y)+2\operatorname{Re}H(x,y)$, so
$q(x+y)=q(x)+q(y)+r$. As $\operatorname{Im}H(x,y)=s$ and $r-s$ is even,
\[
 \alpha(x+y)=(-1)^{q(x)+q(y)+r}=(-1)^{q(x)+q(y)}(-1)^s=\alpha(x)\alpha(y)e^{\pi i\operatorname{Im}H(x,y)} .
\]
\end{proof}
% (1) is hermitian_form_integral_values = HfIntRec_of_gen (HfGenRec_of_arith arithmetic_hermitian_form):
% Step 1 is AR.okc_of, okc_mul and okc_conj, with good_of_okc; Steps 2 and 3 are AR.div_okc (Bf_lat,
% okc_pair, area_spec, ph_om); Step 4 is HfArithRec_of_int with FL.Hf_embV and Hf_smul, the integrality
% h(Z,Z') in O_K being hermitian_pairing_integral (TI.trOK_YY), and the closure is La_le_gen with
% HfInt_closure. (2) is IO.semiChar_of_HfInt (self_int, cf_int, Hf_re_swap).

\begin{lemma}\label{is-growth-coset-triples}
Let $D$, $\Lambda$, $N$, $N'$, $S$, $\Gamma=\Lambda_{N,N'}$, $A_g$ and $B_g$ be as in
Definition~\ref{is-growth-coset-lift}.
\begin{enumerate}
\item $\theta_{N,N'}$ is a theta function for $\Gamma$ (Definition~\ref{is-vacuum-theta-vectors}(2)); in
  particular $\alpha(y)U_y\theta_{N,N'}=\theta_{N,N'}$ for $y\in\Gamma$.
\item For every integer $g\geq1$, $(A_g,B_g,\Gamma)$ is a split triple for $(H,\alpha)$
  (Definition~\ref{is-growth-coset-split}), and the groups $A_g/(A_g\cap\Gamma)$,
  $(\Gamma\cap A_g)/(\Gamma\cap A_g\cap B_g)$ and $B_g/(B_g\cap\Gamma)$ are finite.
\end{enumerate}
\end{lemma}
\leanref{Formalism/LocalToGlobal/KernelLattices/KernelLocalFactors/ThDistC.lean}{Formalism.L2.Ker.heisHyp_of}
\leanref{Formalism/LocalToGlobal/KernelLattices/KernelLocalFactors/AlternatingFormStandardHermitian.lean}{Formalism.L2.Ker.half_periods_of_lattice}
\leanref{Formalism/LocalToGlobal/KernelLattices/KernelLocalFactors/AlternatingFormStandardHermitian.lean}{Formalism.L2.Ker.half_periods_complex}
\leanref{Development/AbelianVarietyTheory/Proofs/DivisorProofs/NormalisedFamilyPeriodLatticePropertiesHold.lean}{GlobalTheory.ColNorms.periodLattices}
\leanref{Formalism/LocalToGlobal/KernelLattices/KernelLocalFactors/ExistsIntegralMulOrderElement.lean}{Formalism.L2.Ker.laFacts_ambOf}
\leanref{Formalism/LocalToGlobal/KernelLattices/KernelLocalFactors/ExistsIntegralMulOrderElement.lean}{Formalism.L2.Ker.e_fixed_ambOf}
\leanref{Formalism/LocalToGlobal/KernelLattices/KernelLocalFactors/IsAdmissibleFlagPairOfUnitriangularLeftAutomorph.lean}{Formalism.L2.Ker.Base.latAll}
\begin{proof}
(a) Let $L\subset E^a$ and $L'\subset E^b$ be lattices that are self-dual for $\sum_i\bar x_iy_i$. Then
$\Lambda_{L,L'}$ is self-dual for $H$: $\operatorname{Im}H$ takes integer values on it by
Lemma~\ref{is-growth-coset-hermitian}(1), and Lemma~\ref{is-growth-period-lattice}(2) gives the converse
inclusion. By Lemma~\ref{is-growth-coset-hermitian}, $\alpha$ is a semicharacter on $\Lambda_{L,L'}$.

(1) The ranks $a$ and $b$ are even (Definition~\ref{th-ranks}), so by Lemma~\ref{is-vacuum-units-host}
$\theta_0^{\otimes ab}=\kappa\,I_{\Lambda_u\to\Gamma_0}(\vartheta_{\Lambda,\omega_0/2})_u$ for an orthonormal
basis $u$ and a number $\kappa$, where $\Gamma_0=\Lambda_{\mathcal O_E^a,\mathcal O_E^b}$. By
Lemma~\ref{is-vacuum-genus-one}(2), $\vartheta_{\Lambda,\omega_0/2}$ satisfies the hypothesis of
Lemma~\ref{is-vacuum-tensor-norm}(3), so $(\vartheta_{\Lambda,\omega_0/2})_u$ is a theta function for
$\Lambda_u$. By Lemma~\ref{is-vacuum-averaging}(3), applied three times, $\theta_0^{\otimes ab}$ is a theta
function for $\Gamma_0$ and
$I_{\Lambda_{N,\mathcal O_E^b}\to\Gamma}I_{\Gamma_0\to\Lambda_{N,\mathcal O_E^b}}\theta_0^{\otimes ab}$ is one
for $\Gamma$. By Definition~\ref{is-vacuum-theta-vectors}, $\theta_{N,N'}$ is a constant multiple of the
latter.

(2) By Lemma~\ref{is-growth-coset-lattices}(4), $N'_A(g)$ and $N'_B(g)$ are self-dual,
$P_A(g)=P(N,N'_A(g))$ and $P_B(g)=P(N,N'_B(g))$; so, as for $\Gamma$, $A_g=\Lambda_{N,N'_A(g)}$ and
$B_g=\Lambda_{N,N'_B(g)}$. By (a), $A_g$, $B_g$ and $\Gamma$ are self-dual for $H$, and $\alpha$ is a
semicharacter on each of them. As $x\mapsto\delta\iota(x)$ is additive and injective,
Lemma~\ref{is-growth-coset-lattices}(2) for $M=P$ gives $\Gamma\subset(\Gamma\cap A_g)+(\Gamma\cap B_g)$, and
$a\in\Gamma$ whenever $a\in A_g$, $b\in B_g$ and $a+b\in\Gamma$. Let $a\in A_g$ with
$\operatorname{Im}H(a,b)\in\mathbf Z$ for $b\in B_g\cap\Gamma$. Every $l\in\Gamma$ is $l_1+l_2$ with
$l_1\in\Gamma\cap A_g$ and $l_2\in\Gamma\cap B_g$; $\operatorname{Im}H(a,l_1)\in\mathbf Z$ because $\alpha$ is a
semicharacter on $A_g$ (proof of Lemma~\ref{is-vacuum-units-intertwiners}(2)), and
$\operatorname{Im}H(a,l_2)\in\mathbf Z$ by assumption; so $a\in\Gamma$, as $\Gamma$ is self-dual. The same
argument applies to $b\in B_g$. Finally, $P$ is finitely generated, and so are $P_A(g)$ and $P_B(g)$, which
lie in $g^{-1}P$ by Lemma~\ref{is-growth-coset-lattices}(1). By the same lemma $gA_g\subset A_g\cap\Gamma$,
$g(\Gamma\cap A_g)\subset\Gamma\cap A_g\cap B_g$ and $gB_g\subset B_g\cap\Gamma$; so each of the three groups
is a quotient of $M/gM$ for a finitely generated free abelian group $M$, and is finite.
\end{proof}

\begin{lemma}\label{is-growth-coset-indicator}
Let $D$ be a definite-side datum with kernel data (Definition~\ref{is-kernel-ties}), $N$ and $N'$
lattices of levels $i$ and $i+1$ (Definition~\ref{is-theta-kernel}), and $S$ a finite set of finite places
of $E$ that contains no two complex-conjugate places and at every place of which $\mathfrak c\mathfrak b^{-1}$
is a unit, with the notation of Definition~\ref{is-growth-coset-lift}; write $H_S=H_S(N,N')$. Let $c\geq1$
be an integer and $x\in\operatorname{Hom}_E(E^a,E^b)$.
\begin{enumerate}
\item If $s$ and $s'$ are split representatives of $x$ at $c$, then
  $\xi^{(c)}_{\delta\iota(s)}=\xi^{(c)}_{\delta\iota(s')}$. Hence
  $\Theta_{N,N',c,x}=v^{-1/2}\,\xi^{(c)}_{\delta\iota(s)}$ for every split representative $s$ of $x$ at $c$.
\item Let $D_0\geq1$ be an integer with $c\mid D_0$, and suppose that $x\in H_S/D_0$. Then $x$ has a split
  representative at $c$, and
  \[
   \Theta_{N,N',c,x}=v^{-1/2}\,\Pi_{D_0}\big(\mathbf 1_{x+cH_S}\big),
  \]
  where $\mathbf 1_{x+cH_S}$ is the indicator function of $x+cH_S$ on $\operatorname{Hom}_E(E^a,E^b)$.
\end{enumerate}
\end{lemma}
\leanref{Formalism/LocalToGlobal/KernelLattices/ThDistC.lean}{Formalism.L2.Ker.Amb.LiftIndep}
\leanref{Formalism/LocalToGlobal/KernelLattices/ThDistC.lean}{Formalism.L2.Ker.Amb.IndTieD}
\leanref{Formalism/LocalToGlobal/KernelLattices/KernelLocalFactors/ThDistC.lean}{Formalism.L2.Ker.liftIndep_of}
\leanref{Formalism/LocalToGlobal/KernelLattices/KernelLocalFactors/ThDistC.lean}{Formalism.L2.Ker.indTieD_of}
\leanref{Formalism/LocalToGlobal/KernelLattices/KernelLocalFactors/ThDistC.lean}{Formalism.L2.Ker.induced_kernel_degree_compat}
\leanref{Formalism/LocalToGlobal/KernelLattices/KernelLocalFactors/ThDistC.lean}{Formalism.L2.Ker.indTieD_of_la}
\begin{proof}
The map $y\mapsto\delta\iota(y)$ is additive and injective. By Lemma~\ref{is-growth-coset-triples},
$\alpha(y)U_y\theta_{N,N'}=\theta_{N,N'}$ for $y\in\Gamma$, and $(A_g,B_g,\Gamma)$ is a split triple for every
$g\geq1$; so Lemma~\ref{is-growth-coset-identities} applies to the operators of
Definition~\ref{is-growth-coset-lift}(3).

(1) Let $s\in P_A(c_1)$ and $s'\in P_A(c_2)$ be split representatives of $x$ at $c$, and put $p=\delta\iota(s)$ and
$m=\delta\iota(s'-s)$. Then $s'-s\in cH_S$, so by Lemma~\ref{is-growth-coset-lattices}(7), for
$g_0=\gcd(c_1,c_2)$, $p$ and $m$ lie in $A_{g_0}$, and $m\in B_c$. As $\alpha$ is a semicharacter on
$A_{g_0}$, Lemma~\ref{is-growth-coset-identities}(2), for the triple $(A_c,B_c,\Gamma)$ and $A'=A_{g_0}$,
gives $\xi^{(c)}_{\delta\iota(s')}=\xi^{(c)}_{p+m}=\xi^{(c)}_p$. The second assertion follows by
Definition~\ref{is-growth-coset-lift}(5).

(2) As $D_0x\in H_S$, Lemma~\ref{is-growth-coset-lattices}(6) gives $s\in P_A(D_0)$ with
$s-x\in D_0H_S\subset cH_S$. So $s$ is a split representative of $x$ at $c$, and, with $p=\delta\iota(s)$,
$\Theta_{N,N',c,x}=v^{-1/2}\xi^{(c)}_p$ by (1). The group $A_{D_0}/(A_{D_0}\cap B_{D_0})$ is finite (proof of
Lemma~\ref{is-growth-coset-identities}(5)); let $T\subset A_{D_0}$ be a set of representatives of it.

\emph{The coset theta function.} By Lemma~\ref{is-growth-coset-lattices}(3) for $M=P$ and $D_1=D_0$, the
triples $(A_{D_0},B_{D_0},\Gamma)$ and $(A_c,B_c,\Gamma)$ satisfy the hypotheses of
Lemma~\ref{is-growth-coset-identities}(6): $A_{D_0}\cap B_c\subset\Gamma$,
$B_c\subset(A_{D_0}\cap B_c)+(B_{D_0}\cap B_c)$, $\Gamma\cap A_{D_0}\subset A_c$, and for every
$r\in\Gamma\cap A_c$ there is $u\in\Gamma\cap A_{D_0}$ with $u-r\in B_c$. As $p\in A_{D_0}$, that lemma gives
\[
 \xi^{(c)}_p=\sum_{t\in T}\mathbf 1_{B_c}(t-p)\,\xi^{(D_0)}_t .
\]

\emph{The Heisenberg lift.} Let $\Phi=\tilde{\mathbf 1}_{x+cH_S}$ (Definition~\ref{is-growth-coset-lift}(4)).
Let $t=\delta\iota(t_0)\in A_{D_0}$ and $m=\delta\iota(m_0)\in A_{D_0}\cap B_{D_0}$. Then
$m_0\in P_A(D_0)\cap P_B(D_0)$, so $m_0\in D_0H_S\subset cH_S$ by Lemma~\ref{is-growth-coset-lattices}(5),
and $\Phi(t+m)=\mathbf 1_{x+cH_S}(t_0+m_0)=\Phi(t)$. So Lemma~\ref{is-growth-coset-identities}(5) gives
\[
 \Pi_{D_0}\big(\mathbf 1_{x+cH_S}\big)=\Pi^{(D_0)}(\Phi)=\sum_{t\in T}\Phi(t)\,\xi^{(D_0)}_t .
\]

\emph{Comparison.} Let $t=\delta\iota(t_0)\in T$, with $t_0\in P_A(D_0)$. Then $\Phi(t)=1$ if and only if
$t_0-x\in cH_S$. And $\mathbf 1_{B_c}(t-p)=1$ if and only if $t_0-s\in P_B(c)$; as $t_0-s\in P_A(D_0)$, this
holds if and only if $t_0-s\in cH_S$ (Lemma~\ref{is-growth-coset-lattices}(5)), that is, as
$s-x\in cH_S$, if and only if $t_0-x\in cH_S$. So the two sums agree term by term, and
$\Theta_{N,N',c,x}=v^{-1/2}\xi^{(c)}_p=v^{-1/2}\,\Pi_{D_0}(\mathbf 1_{x+cH_S})$.
\end{proof}

\begin{lemma}\label{is-growth-theta-facts}
Let $D$ be a definite-side datum, $\Lambda$ a CM lattice of its host with an unkilled two-torsion
point $t_0$, $L\subset E^a$ and $L'\subset E^b$ lattices that are self-dual for $\sum_i\bar x_iy_i$,
with $ab$ even, $F$ a theta vector up to units for $(t_0,L,L')$, and $y$ a torsion point of
$\Lambda_{L,L'}$ (Definition~\ref{is-growth-theta-functions}).
\begin{enumerate}
\item The numbers $g_2(\Lambda)$ and $g_3(\Lambda)$ are algebraic; $F$ is a theta function for
  $\Lambda_{L,L'}$; and $F$ and $U_yF$ have algebraic Taylor coefficients at $0$.
\item Let $\rho$ be an embedding of $\overline{\mathbf Q}$ into $\mathbf C$. Then $\Lambda^\rho$ has an
  unkilled two-torsion point, and $\|F^\rho\|^2_{\Lambda^\rho_{L,L'}}=n(\Lambda^\rho,t')^{fab}$ for every
  unkilled two-torsion point $t'$ of $\Lambda^\rho$.
\item Let $\rho$ be an embedding of $\overline{\mathbf Q}$ into $\mathbf C$ that is the identity on
  $\varphi_0(K)$. Then $\Lambda^\rho$ is a CM lattice of the host; there is a permutation $\pi$ of
  $\mathcal T$ with $\rho\circ\tau=\pi(\tau)$ for every $\tau\in\mathcal T$; $F^\rho$ is a nonzero
  theta function for $\Lambda^\rho_{L,L'}$; and there is a torsion point $y'$ of $\Lambda^\rho_{L,L'}$
  with $(U_yF)^\rho=U^\rho_{y'}F^\rho$.
\item Suppose that $D$ has kernel data (Definition~\ref{is-kernel-ties}) and that
  $\Lambda$ is the CM lattice of the host given by them, and let $N$ and $N'$ be lattices of levels
  $i$ and $i+1$ (Definition~\ref{is-theta-kernel}). Then $\chi_F(N,N')^{-1}\theta_{N,N'}$, for the
  vacuum vector $\theta_{N,N'}$ of the pair, is a theta vector up to units for $(t_0,N,N')$, and
  $\theta_0^{\otimes ab}$ is one for $(t_0,\mathcal O_E^a,\mathcal O_E^b)$
  (Definition~\ref{is-vacuum-theta-vectors}). Moreover, let $S$ be a finite set of finite places of $E$
  that contains no two complex-conjugate places and at each of which $\mathfrak c\mathfrak b^{-1}$ is a
  unit, for the ideal $\mathfrak c$ with $\Lambda=\delta\varphi_0(\mathfrak c)$ and the ideal $\mathfrak b$
  of the host, and form $H_S(N,N')$ and the coset theta functions with $S$. Then for integers
  $c,c'\geq1$ with $cc'$ dividing the modulus $D_0$ of the theta distribution and for
  $x\in H_S(N,N')/D_0$,
  \[
   \Theta_{N,N',c,x}=\sum_{y\in cH_S(N,N')/cc'H_S(N,N')}\Theta_{N,N',cc',x+y} .
  \]
\end{enumerate}
\end{lemma}
\leanref{Development/AbelianVarietyTheory/Proofs/FoundationProofs/IntegralJAndAlgebraicData.lean}{GlobalTheory.CurveAlgebraic.integral_j_and_algebraic_data}
\leanref{Development/AbelianVarietyTheory/Proofs/DivisorProofs/ConjTorsionTranslateVacuumEntireGaussianGrowth.lean}{GlobalTheory.ConjTranslate.mem_Th}
\leanref{Development/AbelianVarietyTheory/Proofs/FoundationProofs/ThetaTaylorCoeffsAlgebraicAtComplexMultiplicatio.lean}{GlobalTheory.ThetaAlgebraic.taylor_algebraic}
\leanref{Development/AbelianVarietyTheory/Proofs/FoundationProofs/ConjugateCurveDatumLatticeIsCurveLattice.lean}{GlobalTheory.ConjCurve.galCurve_holds}
\leanref{Development/AbelianVarietyTheory/Proofs/FoundationProofs/InducedTypeThetaGenusOnePermutes.lean}{GlobalTheory.InducedType.induced_type}
\leanref{Development/AbelianVarietyTheory/Proofs/FoundationProofs/ThetaVacuaConjugatesCmPoints.lean}{GlobalTheory.ThetaConj.useConj_all}
\leanref{Development/AbelianVarietyTheory/Proofs/FoundationProofs/VacuumNormConjCmPointsAllEmbeddings.lean}{GlobalTheory.ConjNorm.useNorm_all}
\leanref{Development/AbelianVarietyTheory/Proofs/FoundationProofs/ExistsTwoTorsionCMPointAllConjugates.lean}{GlobalTheory.TwoTorsionPt.exists_point}
\leanref{Development/AbelianVarietyTheory/Proofs/FoundationProofs/GaloisDataExists.lean}{GlobalTheory.GalExist.galEx_all}
\leanref{Development/AbelianVarietyTheory/Proofs/FoundationProofs/ThetaTranslateGaloisRationalityCmPoints.lean}{GlobalTheory.GalTheta.gal_all}
\leanref{Development/AbelianVarietyTheory/Proofs/DivisorProofs/NormalisedFamilyStandardLatticesEqBaseVectorStatement.lean}{GlobalTheory.T.normalisedFamily_standardLattices_eq_baseVector}
\leanref{Development/AbelianVarietyTheory/Proofs/DivisorProofs/IsVacuumMovedThetaVacuumStatement.lean}{GlobalTheory.T.thetaVectorUpToUnits_thMoved_of_thZero}
\leanref{Development/AbelianVarietyTheory/Proofs/DivisorProofs/CosetIndicatorThetaRefinementProofStatement.lean}{GlobalTheory.T.cosetIndicatorThetaRefinement_proof}
\begin{proof}
Write $F=\kappa I_{\Lambda_v\to\Lambda_{L,L'}}(\vartheta_{\Lambda,t_0})_v$ for an orthonormal basis $v$ and an
algebraic $\kappa$ all of whose conjugates have absolute value $1$
(Definition~\ref{is-vacuum-theta-vectors}(6)).

(1) The numbers $g_2(\Lambda)$ and $g_3(\Lambda)$ are algebraic by
Lemma~\ref{is-growth-theta-genus-one}(1). By Lemma~\ref{is-vacuum-genus-one}(2), $\vartheta_{\Lambda,t_0}$
satisfies the hypothesis of Lemma~\ref{is-vacuum-tensor-norm}(3), so $(\vartheta_{\Lambda,t_0})_v$ is a
theta function for $\Lambda_v$. By Lemma~\ref{is-vacuum-averaging}(3), its image under
$I_{\Lambda_v\to\Lambda_{L,L'}}$ is a theta function for $\Lambda_{L,L'}$, and so is $F$. The Taylor
coefficients of $F$ and $U_yF$ are algebraic by Lemma~\ref{is-growth-theta-algebraic}(5).

(2) By Lemma~\ref{is-growth-theta-conjugate}(1), $\Lambda^\rho$ has an unkilled two-torsion point, and the
norm of $F^\rho$ is given by Lemma~\ref{is-growth-theta-conjugate}(3).

(3) By Lemma~\ref{is-growth-theta-conjugate}(1), $\Lambda^\rho$ is a CM lattice of the host and the
permutation $\pi$ exists. By Lemma~\ref{is-growth-theta-conjugate}(2), $F^\rho$ is a theta function for
$\Lambda^\rho_{L,L'}$. It is nonzero, because its norm is positive by
Lemma~\ref{is-growth-theta-conjugate}(3), applied with an unkilled two-torsion point of $\Lambda^\rho$. The
torsion point $y'$ is given by Lemma~\ref{is-growth-theta-conjugate}(4).

(4) The ranks $a=N(i)$ and $b=N(i+1)$ are even (Definition~\ref{th-ranks}), and the lattices $N$, $N'$,
$\mathcal O_E^a$ and $\mathcal O_E^b$ are self-dual for $\sum_i\bar x_iy_i$; so the pairs $(N,N')$ and
$(\mathcal O_E^a,\mathcal O_E^b)$ are admitted. By Lemma~\ref{is-vacuum-genus-one}(1),
$t_0-\omega_0/2\in\Lambda$ for the unkilled two-torsion point $\omega_0/2$ given by the kernel data, and
$\vartheta_{\Lambda,t_0}=\vartheta_{\Lambda,\omega_0/2}$. So a theta vector up to units for
$(\omega_0/2,L_1,L'_1)$ is one for $(t_0,L_1,L'_1)$, for all lattices $L_1,L'_1$.

\emph{The tensor power.} By Lemma~\ref{is-vacuum-units-host}, $\theta_0^{\otimes ab}$ is a theta vector up
to units for $(\omega_0/2,\mathcal O_E^a,\mathcal O_E^b)$; by the previous paragraph it is one for
$(t_0,\mathcal O_E^a,\mathcal O_E^b)$.

\emph{The vacuum vector of the pair.} The constant $\chi_F(N,N')$ of Definition~\ref{is-theta-kernel}(2)
is the factor $\chi(\mathfrak d_N)^{b/2}\chi(\mathfrak d_{N'})^{-a/2}$ in the definition of $\theta_{N,N'}$
(Definition~\ref{is-vacuum-theta-vectors}), so
\[
 \chi_F(N,N')^{-1}\theta_{N,N'}=I_{\Lambda_{N,\mathcal O_E^b}\to\Lambda_{N,N'}}
 I_{\Lambda_{\mathcal O_E^a,\mathcal O_E^b}\to\Lambda_{N,\mathcal O_E^b}}\theta_0^{\otimes ab}.
\]
Let $G=\kappa'I_{\Lambda_w\to\Gamma}(\vartheta_{\Lambda,t_0})_w$ be a theta vector up to units for
$(t_0,L_1,L'_1)$, with $\Gamma=\Lambda_{L_1,L'_1}$, and let $\Gamma'=\Lambda_{L_2,L'_2}$ for self-dual lattices
$L_2,L'_2$. As in (1), $(\vartheta_{\Lambda,t_0})_w$ is a theta function for $\Lambda_w$, so
Lemma~\ref{is-vacuum-averaging}(4) gives $I_{\Gamma\to\Gamma'}G=\kappa'\mu
I_{\Lambda_w\to\Gamma'}(\vartheta_{\Lambda,t_0})_w$ with $\mu^8=1$; and all conjugates of $\kappa'\mu$ have
absolute value $1$. So $I_{\Gamma\to\Gamma'}G$ is a theta vector up to units for $(t_0,L_2,L'_2)$. Applied
to $\theta_0^{\otimes ab}$, first with $\Gamma'=\Lambda_{N,\mathcal O_E^b}$ and then with
$\Gamma'=\Lambda_{N,N'}$, this gives the first assertion.

\emph{The refinement.} Write $H=H_S(N,N')$. The group $cH/cc'H$ is finite, and the cosets $x+y+cc'H$, for
$y$ running over representatives of $cH/cc'H$, partition $x+cH$. So the indicator functions on
$\operatorname{Hom}_E(E^a,E^b)$ satisfy
\[
 \mathbf 1_{x+cH}=\sum_{y\in cH/cc'H}\mathbf 1_{x+y+cc'H}.
\]
By the hypotheses on $S$, Lemma~\ref{is-growth-coset-indicator}(2) applies, and the map
$\mathcal L=v^{-1/2}\Pi_{D_0}$ of Definition~\ref{is-growth-coset-lift}(4), linear from functions on
$\operatorname{Hom}_E(E^a,E^b)$ to functions on $V_{a,b}$, satisfies $\Theta_{N,N',c,x}=\mathcal L(\mathbf
1_{x+cH})$ and $\Theta_{N,N',cc',x+y}=\mathcal L(\mathbf 1_{x+y+cc'H})$, since $c$ and $cc'$ divide $D_0$ and
$x+y\in H/D_0$. Applying $\mathcal L$ to the identity of indicator functions gives the refinement.
% In the Lean the tensor power reaches (4) through FamilyTie (FamilyTieProof.familyTie, natTie, from
% T.base_theta_vector_is_vacuum_even_ranks) and T.normalisedFamily_standardLattices_eq_baseVector (eNat at
% the standard lattices is Base.thZero). The refinement is Amb.IndAdd (indAdd_of, ColNorms.coset_refinement:
% ad_xiL_sum and ad_onV_sum for the linearity, ad_cosInd_sum for the indicators,
% lattice_classes_finite_refinement for the finiteness) from Amb.IndTieD at the modulus D_0, which is
% Lemma is-growth-coset-indicator(2) (indTieD_of_la, from Amb.LaHyp through heisHyp_of; LaHyp is the
% hypothesis PeriodLattices, proved as T.normalisedFamily_periodLattice_properties_hold). The two
% conditions on S in (4) are the hypotheses hconj and AAUnit that PeriodLattices takes. For the places
% dividing the level, S = Sall sel mu u at mu = 1, kernel_le_sharp discharges them by
% Core.Ht2.BadPrimesS.bad_primes_sharp and T.periodLatticeIdeal_trivial_at_badPrimes
% (Lemma is-growth-coset-places(4), used in Proposition is-kernel-entry-bound).
\end{proof}

\begin{lemma}\label{is-growth-conjugate-translate}
Let $D$ be a definite-side datum, $\Lambda$ a CM lattice of its host with an unkilled two-torsion
point $t_0$, $L\subset E^a$ and $L'\subset E^b$ lattices that are self-dual for $\sum_i\bar x_iy_i$,
with $ab$ even, $F$ a theta vector up to units for $(t_0,L,L')$, $y$ a torsion point of
$\Lambda_{L,L'}$, and $\rho$ an embedding of $\overline{\mathbf Q}$ into $\mathbf C$
(Definition~\ref{is-growth-theta-functions}). Then $U_yF$ is entire with algebraic Taylor
coefficients at $0$; the series $(U_yF)^\rho$ converges to an entire function, whose Taylor
coefficients are the images under $\rho$ of those of $U_yF$; and for every $z\in V_{a,b}$
\[
 |(U_yF)^\rho(z)|^2\le s_0(a,b)\,\|F^\rho\|^2_{\Lambda^\rho_{L,L'}}\,
 \exp\big(\pi\langle z,z\rangle/A(\Lambda^\rho)\big),
\]
with $s_0(a,b)=3(2fab\,r_E+1)^{fab}$ as in Lemma~\ref{is-growth-translate-bound}.
\end{lemma}
\leanref{Development/AbelianVarietyTheory/Proofs/DivisorProofs/ConjTorsionTranslateVacuumEntireGaussianGrowthStatement.lean}{GlobalTheory.T.conj_torsion_translate_vacuum_entire_gaussian_growth}
\begin{proof}
By Lemma~\ref{is-growth-theta-facts}(1), $F$ is entire, so $U_yF$ is entire, and its Taylor
coefficients are algebraic. Note that $\langle z,z\rangle/A(\Lambda^\rho)=H^\rho(z,z)$.

\emph{Case 1: $\rho$ is the identity on $\varphi_0(K)$.} By Lemma~\ref{is-growth-theta-facts}(3),
$\Lambda^\rho$ is a CM lattice of the host, there is a permutation $\pi$ of $\mathcal T$ with
$\rho\circ\tau=\pi(\tau)$, $F^\rho$ is a theta function for $\Lambda^\rho_{L,L'}$, and
$(U_yF)^\rho=U^\rho_{y'}F^\rho$ for a torsion point $y'$ of $\Lambda^\rho_{L,L'}$. The right side is
entire, so the series $(U_yF)^\rho$ converges to an entire function with the stated Taylor
coefficients. Let $P$ be the permutation of coordinates $(Pz)_\tau=z_{\pi(\tau)}$. It preserves
$H^\rho$ and Lebesgue measure, carries the subgroup $\Lambda'_{L,L'}$ formed with
$\Lambda'=\Lambda^\rho$ and the types $\tau$ onto $\Lambda^\rho_{L,L'}$, which is formed with the
embeddings $\pi(\tau)$, and satisfies $(U^\rho_{Pw}G)(Pz)=U'_w(G\circ P)(z)$ for the operators $U'$
of $\Lambda'$ and the types. So $G=F^\rho\circ P$ is a theta function for $\Lambda'_{L,L'}$ with
$\|G\|^2_{\Lambda'_{L,L'}}=\|F^\rho\|^2_{\Lambda^\rho_{L,L'}}$, and
Lemma~\ref{is-growth-translate-bound} for the CM lattice $\Lambda'$ gives
\[
 |U^\rho_{y'}F^\rho(z)|^2=|U'_{P^{-1}y'}G(P^{-1}z)|^2\le s_0(a,b)\,
 \|F^\rho\|^2_{\Lambda^\rho_{L,L'}}\,e^{\pi H^\rho(z,z)} .
\]

\emph{Case 2: otherwise.} An embedding of $K$ sends $\sqrt{-d}$ to $\pm i\sqrt d$, so
$\bar\rho=\overline{\rho(\cdot)}$ is the identity on $\varphi_0(K)$. For a function $G$ with algebraic
Taylor coefficients, $G^\rho(z)=\overline{G^{\bar\rho}(\bar z)}$ where either side converges. The
lattice $\overline{\Lambda^\rho}$ has invariants $\overline{\rho(g_k(\Lambda))}=\bar\rho(g_k(\Lambda))$,
so $\Lambda^{\bar\rho}=\overline{\Lambda^\rho}$, and $A(\Lambda^{\bar\rho})=A(\Lambda^\rho)$. Since
$\bar\rho\circ\tau$ is the complex conjugate of $\rho\circ\tau$, the map $z\mapsto\bar z$ carries
$\Lambda^\rho_{L,L'}$ onto $\Lambda^{\bar\rho}_{L,L'}$, preserves Lebesgue measure and satisfies
$H^{\bar\rho}(\bar z,\bar w)=\overline{H^\rho(z,w)}$; hence
$\|F^{\bar\rho}\|^2_{\Lambda^{\bar\rho}_{L,L'}}=\|F^\rho\|^2_{\Lambda^\rho_{L,L'}}$. Case 1 for
$\bar\rho$ at $\bar z$ gives the assertions for $(U_yF)^\rho(z)=\overline{(U_yF)^{\bar\rho}(\bar z)}$.
\end{proof}

\begin{lemma}\label{is-growth-coset-translates}
Let $D$ be a definite-side datum with kernel data (Definition~\ref{is-kernel-ties}), $\Lambda$ the CM
lattice of the host given by them, $N$ and $N'$ lattices of levels $i$ and $i+1$
(Definition~\ref{is-theta-kernel}), $c\geq1$ an integer and $x\in\operatorname{Hom}_E(E^a,E^b)$. Let
$\theta_{N,N'}$ be the vacuum vector of the pair (Definition~\ref{is-vacuum-theta-vectors}) and
$F_0=\chi_F(N,N')^{-1}\theta_{N,N'}$. Then there are torsion points $y_1,\dots,y_k$ of
$\Lambda_{N,N'}$ (Definition~\ref{is-growth-theta-functions}), with $k\le c^{[E:\mathbf Q]ab}$, and
algebraic numbers $r_1,\dots,r_k$ with $|\rho(r_j)|\le1$ for every embedding $\rho$ of
$\overline{\mathbf Q}$ into $\mathbf C$, such that
\[
 \chi_F(N,N')^{-1}\,\Theta_{N,N',c,x}=\sum_{j\le k}r_j\,U_{y_j}F_0 .
\]
\end{lemma}
\leanref{Development/AbelianVarietyTheory/Proofs/DivisorProofs/CosetIndicatorThetaTorsionTranslatesStatement.lean}{GlobalTheory.T.cosetIndicatorTheta_torsionTranslates}
\begin{proof}
Write $\alpha=\alpha_{\sqrt d/2}$ and $\Lambda=\delta\varphi_0(\mathfrak c)$, and let $P$ be the
$\mathcal O_E$-span of the matrices $sB$ with $s\in\mathfrak c\mathfrak b^{-1}$ and $BN\subset N'$. Let
$P_A$ and $P_B$ be the lattices that agree with $P$ away from $S$ and its complex conjugate $\bar S$,
with $c^{-1}P$ and $cP$ respectively at the places of $S$, and with $cP$ and $c^{-1}P$ at those of
$\bar S$. Let $\Lambda_c$, $A_c$ and $B_c$ be the images of $P$, $P_A$ and $P_B$ under
$x'\mapsto\delta\iota(x')$; then $\Lambda_c=\Lambda_{N,N'}$, because every type restricts to
$\varphi_0$ on $K$, so that $\delta\varphi_0(k)\iota(x')=\delta\iota(kx')$ for $k\in\mathfrak c$. Let
$S_A\subset\Lambda_c\cap A_c$ be a set of representatives modulo $B_c$, $S_B\subset B_c$ a set of
representatives modulo $\Lambda_c$, $v$ the covolume of Definition~\ref{is-theta-kernel}(2), and
$p=\delta\iota(s)$ for the element $s$ with $s-x\in cH_S(N,N')$ chosen in the construction. By its
construction the coset theta function is
\[
 \Theta_{N,N',c,x}=v^{-1/2}\alpha(p)\,U_{-p}\Big(\sum_{y\in S_B}\hat h(y)\alpha(y)\,U_y\theta_{N,N'}\Big),
 \qquad
 \hat h(y)=v^{-1/2}\frac v{|S_A|}\sum_{t\in S_A}\mathbf 1_{B_c}(t)\,e^{2\pi i\operatorname{Im}H(t,y)} .
\]
A direct computation from the definition gives
$U_{-p}U_y=e^{\pi i\operatorname{Im}H(p,y)}U_{y-p}$. Hence
$\chi_F(N,N')^{-1}\Theta_{N,N',c,x}=\sum_{y\in S_B}r_yU_{y-p}F_0$ with
\[
 r_y=\big(\chi_F^{-1}\chi_F\big)\big(v^{-1}v\big)\Big(\frac1{|S_A|}\sum_{t\in S_A}\mathbf 1_{B_c}(t)\,
 e^{2\pi i\operatorname{Im}H(t,y)}\Big)\,\alpha(p)\alpha(y)\,e^{\pi i\operatorname{Im}H(p,y)} .
\]
The points $t$, $y$ and $p$ are of the form $\delta\iota(x')$, so they are torsion points of
$\Lambda_{N,N'}$ by Lemma~\ref{is-growth-period-lattice}(5), and so is $y-p$. By the same lemma
$\operatorname{Im}H(t,y)$, $\operatorname{Im}H(p,y)$, $\frac{\sqrt d}2H(y,y)$ and
$\frac{\sqrt d}2H(p,p)$ are rational, so the exponentials and $\alpha(y)$, $\alpha(p)$ are roots of
unity. Roots of unity, and averages and products of algebraic numbers whose conjugates have
absolute value at most $1$, are algebraic numbers whose conjugates have absolute value at most $1$;
so every $r_y$ is one. Finally $P_B\subset c^{-1}P$, so $B_c/(B_c\cap\Lambda_c)$ embeds in
$c^{-1}\Lambda_c/\Lambda_c$, which has $c^{[E:\mathbf Q]ab}$ elements because $P$ is free of rank
$[E:\mathbf Q]ab$ over $\mathbf Z$. So $|S_B|\le c^{[E:\mathbf Q]ab}$.
\end{proof}

\begin{lemma}\label{is-growth-coset-refinement}
Let $D$ be a definite-side datum with kernel data (Definition~\ref{is-kernel-ties}), $N$
and $N'$ lattices of levels $i$ and $i+1$, $S$ a finite set of finite places of $E$ that contains no two
complex-conjugate places and at each of which $\mathfrak c\mathfrak b^{-1}$ is a unit, for the CM lattice
$\delta\varphi_0(\mathfrak c)$ of the host given by the kernel data and the ideal $\mathfrak b$ of the host,
$H_S=H_S(N,N')$ with the coset theta functions formed with $S$ (Definition~\ref{is-theta-kernel}), and
$c_A,c,c'\geq1$ integers with $c_Acc'$ equal to the modulus $D_0$ of the theta distribution. Let
$\Xi$ be a complex function on $\operatorname{Hom}_E(E^a,E^b)$ that vanishes off $H_S$ and satisfies
$\Xi(x+y)=\Xi(x)$ for $y\in cH_S$. Then there is a set $R\subset H_S$ with at most
$c^{[E:\mathbf Q]ab}$ elements such that
\[
 \sum_{x\in(H_S/c_A)/cc'H_S}\Xi(x)\,\Theta_{N,N',cc',x}=\sum_{x\in R}\Xi(x)\,\Theta_{N,N',c,x} .
\]
\end{lemma}
\leanref{Development/AbelianVarietyTheory/Proofs/DivisorProofs/LatSumThIndRegroupStatement.lean}{GlobalTheory.T.exists_finset_latSum_thInd_eq_sum}
\begin{proof}
The function $\Theta_{N,N',cc',x}$ depends only on the class of $x$ modulo $cc'H_S$, and only
classes in $H_S$ contribute, since $\Xi$ vanishes off $H_S$. Let $R\subset H_S$ be a set of
representatives of $H_S/cH_S$. The group $H_S$ is torsion free of rank at most $[E:\mathbf Q]ab$, so
every finitely generated subgroup $F$ of $H_S$ has
$|F/(F\cap cH_S)|\le|F/cF|\le c^{[E:\mathbf Q]ab}$, and therefore $|R|\le c^{[E:\mathbf Q]ab}$. Every
class of $H_S$ modulo $cc'H_S$ is $x+y+cc'H_S$ for a unique $x\in R$ and a unique
$y\in cH_S/cc'H_S$, and $\Xi(x+y)=\Xi(x)$. So the left side is
$\sum_{x\in R}\Xi(x)\sum_{y\in cH_S/cc'H_S}\Theta_{N,N',cc',x+y}$, and the inner sum is
$\Theta_{N,N',c,x}$ by Lemma~\ref{is-growth-theta-facts}(4), because $x\in H_S\subset H_S/D_0$ and
$cc'$ divides $D_0$.
\end{proof}

\begin{lemma}\label{is-growth-coset-places}
Let $E$ be the host field of a definite-side datum, with complex conjugation $x\mapsto\bar x$, the
integers $Q_S$, $q$ and $d$ of the host, and the ideal $\mathfrak b$ of the host
(Definitions~\ref{hf-framed-host-field} and~\ref{hf-host}). Let $T$, $g_0$ and $g_1$ be primes with
$g_0\neq g_1$ and $T\notin\{g_0,g_1\}$. Let $\pi,\pi_0,\pi_1\in\mathcal O_E$ satisfy $\pi\bar\pi=T$,
$\pi\mathcal O_E+\bar\pi\mathcal O_E=\mathcal O_E$ and $\pi-1\in Q_Sqd(g_0g_1)^2\mathcal O_E$, and, for
$k=0,1$, $\pi_k\bar\pi_k=g_k$, $\pi_k\mathcal O_E+\bar\pi_k\mathcal O_E=\mathcal O_E$ and
$g_k\equiv1\pmod{Q_Sqd}$. Let $S$ be the set of finite places of $E$ that divide $\pi$, $\bar\pi_0$ or
$\bar\pi_1$.
\begin{enumerate}
\item $S$ contains no two complex-conjugate places.
\item Every place of $S$ lies over a prime $\ell\in\{T,g_0,g_1\}$ that does not divide $Q_Sqd$.
\item Let $Q\in S$, over the prime $\ell$ of (2), and let $\mathfrak c$ be a nonzero ideal of $\mathcal O_K$
  whose norm is prime to $\ell$. Then $\mathfrak c\mathfrak b^{-1}$ is a unit at $Q$: its localization at
  $Q$ is $\mathcal O_{E,Q}$.
\item Let $D$ be a definite-side datum with kernel data that satisfy the kernel ties
  (Definition~\ref{is-kernel-ties}), $u$ a word, $S$ the set of places dividing the level of $u$ at the
  base tuple (Definition~\ref{is-theta-kernel}(4)), and $\Lambda=\delta\varphi_0(\mathfrak c)$ the CM lattice
  of the host given by the kernel data. Then $S$ contains no two complex-conjugate places, and
  $\mathfrak c\mathfrak b^{-1}$ is a unit at each of its places.
\end{enumerate}
\end{lemma}
\leanref{Development/AbelianVarietyTheory/Proofs/DivisorProofs/CharacterSumLocalConstantSecond.lean}{GlobalTheory.ColNorms.BadPrimes.conjP_notMem}
\leanref{Development/AbelianVarietyTheory/Proofs/DivisorProofs/CharacterSumLocalConstantThird.lean}{GlobalTheory.ColNormsUnit.exists_prime}
\leanref{Development/AbelianVarietyTheory/Proofs/GlobalInterpolation/LocPAAAndPrimesNe.lean}{GlobalTheory.AAUnitAway.locP_AA}
\leanref{Development/AbelianVarietyTheory/Proofs/GlobalInterpolation/LocPAAAndPrimesNe.lean}{GlobalTheory.AAUnitAway.exists_den}
\leanref{Development/AbelianVarietyTheory/Proofs/DecoratedLatticeProofs/BadPrimesSharp.lean}{GlobalTheory.Core.Ht2.BadPrimesS.bad_primes_sharp}
\leanref{Development/AbelianVarietyTheory/Proofs/DecoratedLatticeProofs/SharpUnitAtBadPrimesStatement.lean}{GlobalTheory.T.periodLatticeIdeal_trivial_at_badPrimes}
\leanref{Development/AbelianVarietyTheory/Proofs/DecoratedLatticeProofs/SharpUnitAtBadPrimes.lean}{GlobalTheory.Core.Ht2.AAUnitS.aaUnit_bad_sharp}
\begin{proof}
Write $V_T$ for the set of places that divide $\pi$ and $V_k$ for the set of places that divide
$\bar\pi_k$, so that $S=V_T\cup V_0\cup V_1$, and $\bar Q$ for the complex conjugate of a place $Q$; then
$x\in\bar Q$ if and only if $\bar x\in Q$. Every place of $V_T$ contains $T=\pi\bar\pi$, and every place of
$V_k$ contains $g_k=\pi_k\bar\pi_k$. Two coprime integers never lie in one place, since $1$ is a
combination of them.

(1) Let $Q$ and $\bar Q$ both lie in $S$. A rational integer lies in $Q$ if and only if it lies in
$\bar Q$. As $T$, $g_0$ and $g_1$ are distinct primes, $Q$ and $\bar Q$ therefore lie in the same one of
$V_T$, $V_0$ and $V_1$. If both lie in $V_T$, then $\pi\in\bar Q$ gives $\bar\pi\in Q$, and $\pi\in Q$; this
is impossible, as $\pi$ and $\bar\pi$ generate $\mathcal O_E$. If both lie in $V_k$, then $\bar\pi_k\in\bar Q$
gives $\pi_k\in Q$, and $\bar\pi_k\in Q$; this is impossible in the same way.

(2) Let $Q\in V_k$, so that $g_k\in Q$. If $g_k$ divided $Q_Sqd$, then $g_k\equiv1\pmod{Q_Sqd}$ would give
$g_k\mid1$. Let $Q\in V_T$, so that $T\in Q$. If $T$ divided $Q_Sqd$, then $T$ would divide $\pi-1$ in
$\mathcal O_E$, so $\pi-1\in Q$, and $1=\pi-(\pi-1)\in Q$.

(3) The ideal $\mathfrak b$ is not contained in $Q$. Otherwise $Q$ divides
$\mathfrak b\bar{\mathfrak b}=\mathfrak D_{E/K}$, so $Q$ lies over a prime dividing $Q_Sq$ (proof of
Proposition~\ref{hf-host-field-properties}(f)) and $Q_Sq\in Q$; but $\ell\in Q$, and $\ell$ is prime to
$Q_Sq$ by (2). Choose $z\in\mathfrak b$ with $z\notin Q$. For $c\in\mathfrak c$ and $t\in\mathfrak b^{-1}$,
$zct\in\mathcal O_E$, because $zt\in\mathcal O_E$ and $c\in\mathcal O_K$; so
$\mathfrak c\mathfrak b^{-1}\subset z^{-1}\mathcal O_E\subset\mathcal O_{E,Q}$. Conversely, the norm $n$ of
$\mathfrak c$ lies in $\mathfrak c$, and $1\in\mathfrak b^{-1}$ as $\mathfrak b$ is integral, so
$n\in\mathfrak c\mathfrak b^{-1}$. As $n$ is prime to $\ell$ and $\ell\in Q$, $n\notin Q$; so
$\mathcal O_E\subset n^{-1}\mathfrak c\mathfrak b^{-1}$ lies in the localization of
$\mathfrak c\mathfrak b^{-1}$ at $Q$. Hence that localization is $\mathcal O_{E,Q}$.

(4) By the kernel ties, the primes of the kernel data satisfy the standing conditions of the host,
whose conductor is $Q_Sqd$. For the level primes $g_0\neq g_1$ of the host and the prime
$T_u\notin\{g_0,g_1\}$ of $u$, these conditions give elements $\pi_u,\pi_0,\pi_1\in\mathcal O_K$ with the
properties above, for $T=T_u$ and $\pi=\pi_u$. The kernel ties also take the ideal $\mathfrak c$ of the
host to be $\mathcal O_K$. At the base tuple the tuple of characters is trivial and adds no place to the
level, so $S$ is the set of places that divide $\pi_u$, $\bar\pi_0$ or $\bar\pi_1$. Now (1) gives the first
assertion, and (3), with $\mathfrak c=\mathcal O_K$ of norm $1$, gives the second.
\end{proof}
% (1) is ColNorms.BadPrimes.conjP_notMem (T_mem, g_mem, natCast_mem_of_conjP, eq_of_natCast_mem,
% false_of_isCoprime with piT_cop and piG_cop, g_injective); (2) is ColNormsUnit.exists_prime (piT_cong,
% g_cong); (3) is AAUnitAway.locP_AA with exists_den (D.host.𝔟_spec.1 and D.host.der.diffSupport,
% Ideal.absNorm_mem and one_mem_inverse_halfDifferent). (4) is the last clause of
% Core.Ht2.BadPrimesS.bad_primes_sharp and Core.Ht2.AAUnitS.aaUnit_bad_sharp
% (T.periodLatticeIdeal_trivial_at_badPrimes), from HypS: HostTie gives Pr.Ok, Pr.cond = D.host.condN and
% D.host.ccK = top, and mu = 1 gives rowSG = vGs (rowSG_one), so Sall sel mu u = vT u ∪ vG 0 ∪ vG 1.

\begin{proposition}\label{is-kernel-entry-bound}
For every integer $C\geq0$ there is $C_K\geq0$ with the following property. Let $D$ be a
definite-side datum with size data at scale $x$ and kernel data that satisfy the kernel ties with
weights $W$, the size conditions with a constant $C_Y$, and the kernel sizes
(Definition~\ref{is-kernel-ties}), for a word $u$, at the base tuple. Let $M\subset\mathbf C$ be a
number field containing the coefficients of the theta data and the image of $L$, and $\rho$ an
embedding of $M$. Let $0\le i<m$, and let $N$ and $N'$ be lattices of levels $i$ and $i+1$ with frames,
and $\Xi$ a function with values in $L$ that vanishes off $H_S(N,N')$, is invariant under
translation by $T_u g_0g_1H_S(N,N')$, and is bounded by $(g_0g_1)^{Cfm^2N_m^2}$ under every embedding.
Then the kernel entry $K_\Xi(N,N')$ of the $(i+1)$-th lift satisfies, for all $\varphi$,
\[
 \|\rho K_\Xi(N,N')\varphi\|^2\le e^{k_{i+1}(\rho)}\|\rho\varphi\|^2,
\]
the norms of Definition~\ref{is-polynomial-model}(7) after the embedding $\rho$, where, with
$U=fm^2N_m^2\log(g_0g_1T_u)$ and $\phi(y)=y\log y-y$,
\begin{multline*}
 k_j(\rho)=\sum_\tau\Big(\sum_i\phi(\nu_{j,\tau,i})-|\nu_{j,\tau}|\log\Omega_\tau(\rho)\Big)
 +N(j-1)N(j)\,\ell_\theta(\rho)+fN(j-1)N(j)\log\big(8+32fN(j-1)N(j)\tfrac{r_E}2\big)\\
 +2C_KU+\sum_\tau\big(1+2N(j-1)N(j)\log(2|\nu_{j,\tau}|+2)\big).
\end{multline*}
Here $\nu_{j,\tau}$ is the weight at $\tau$ of the $j$-th lift, $\Omega_\tau(\rho)$ is the period at
$\tau$ of the CM lattice of the host under $\rho$, $r_E$ is the root discriminant of the host field,
and $N(j-1)N(j)\,\ell_\theta(\rho)$ is the logarithm of the squared norm of the conjugate under $\rho$
of the vacuum vector $\theta_0^{\otimes N(j-1)N(j)}$ (Definition~\ref{is-vacuum-theta-vectors}).
\end{proposition}
\leanref{Development/AbelianVarietyTheory/Proofs/DecoratedLatticeProofs/ChainKernelBoundAtEmbeddingAnyLevelPrimeStatement.lean}{GlobalTheory.T.chainKernel_bound_at_embedding_anyLevelPrime}
\leanref{Development/AbelianVarietyTheory/Proofs/DecoratedLatticeProofs/ChainKernelBoundAtEmbeddingAnyLevelPrime.lean}{GlobalTheory.Core.Ht2.KernelS.kernel_le_sharp}
\leanref{Development/AbelianVarietyTheory/Proofs/DivisorProofs/JetCConstMul.lean}{GlobalTheory.ColNorms.block}
\leanref{Development/AbelianVarietyTheory/Proofs/DivisorProofs/JetBlockBoundLeExpKernelBoundStatement.lean}{GlobalTheory.T.exists_jetBlockBound_le_exp_kernelBound}
\begin{proof}
Let $a=N(i)$ and $b=N(i+1)$; then $a<b$, and $ab$ is even. The weights of the forms of $u$ at levels $i$
and $i+1$, with the weights $\nu$ of the $(i+1)$-th lift, are step weights, by the weight tie.

\emph{Step 1: the entry is one jet operator.} Choose representatives $x_q$ of the cosets $q$ in the
sum defining $K_\Xi(N,N')$. By the definition of the theta distribution,
$K_\Xi(N,N')=\mu(N)\mu(N')^{-1}\,\mathrm{Jet}(J,\ell(\nu))$ with
$J(m)=\sum_q\Xi(x_q)J_{x_q}(m)$, since the jet operator is linear in its coefficients. The characters
of the kernel data at the bad primes have values of absolute value $0$ or $1$, so
$|\rho(\mu(N)\mu(N')^{-1})|\le1$.

\emph{Step 2: the jet of one function.} Let $\Lambda$ be the CM lattice of the host given by the
kernel data, $t_0$ an unkilled two-torsion point of it, $\theta_{N,N'}$ the vacuum vector of the pair
and $F_0=\chi_F(N,N')^{-1}\theta_{N,N'}$, a theta vector up to units for $(t_0,N,N')$ by
Lemma~\ref{is-growth-theta-facts}(4) (Definitions~\ref{is-vacuum-theta-vectors}
and~\ref{is-growth-theta-functions}). Put
$c=T_ug_0g_1$ and $c'=(g_0g_1)^{fB-1}$, so that $c_B=cc'$, and $\beta=(g_0g_1)^{Cfm^2N_m^2}$. By
Lemma~\ref{is-growth-coset-translates} the coset theta functions are finite combinations of the
entire functions $U_yF_0$, so they are entire. As the theta jets of the lift have coefficients in
$L$ by the kernel ties, and the jet is linear,
$\sigma_0(J(m))=j_{(c_\tau)}(F)(m)$ for $F=\sum_q\sigma_0(\Xi(x_q))\chi_F^{-1}\Theta_{N,N',cc',x_q}$,
at the exponents $m$ of degree $|\nu_\tau|$ in each $\tau$.

(a) \emph{Regrouping.} The function $\sigma_0\circ\Xi$ vanishes off $H_S(N,N')$ and is invariant
under translation by $cH_S(N,N')$. By Lemma~\ref{is-growth-coset-places}(4), the set $S$ of places
dividing the level of $u$ contains no two complex-conjugate places, and $\mathfrak c\mathfrak b^{-1}$ is a
unit at each of its places, for $\Lambda=\delta\varphi_0(\mathfrak c)$. So
Lemma~\ref{is-growth-coset-refinement}, with the modulus
$D_0=c_Ac_B=c_Acc'$, gives a set $R\subset H_S(N,N')$ with at most $c^{[E:\mathbf Q]ab}$ elements and
$F=\sum_{x\in R}\sigma_0(\Xi(x))\chi_F^{-1}\Theta_{N,N',c,x}$.

(b) \emph{Torsion translates.} By Lemma~\ref{is-growth-coset-translates}, each
$\chi_F^{-1}\Theta_{N,N',c,x}$ is a combination of at most $c^{[E:\mathbf Q]ab}$ translates $U_yF_0$
by torsion points $y$ of $\Lambda_{N,N'}$, with algebraic coefficients whose conjugates have absolute
value at most $1$. So $F=\sum_ls_lU_{y_l}F_0$, a finite sum over torsion points $y_l$, where each
$s_l$ is $\sigma_0(\Xi(x))$ times such a coefficient for some $x\in R$. For an embedding $\rho_1$ of
$\overline{\mathbf Q}$, $\rho_1\circ\sigma_0$ is an embedding of $L$, so
$|\rho_1(\sigma_0(\Xi(x)))|\le\beta$, and
\[
 \sum_l|\rho_1(s_l)|\le|R|\,\beta\,c^{[E:\mathbf Q]ab}\le\beta c^{2[E:\mathbf Q]ab} .
\]

(c) \emph{Conjugation.} Extend $\rho$ to an embedding $\rho'$ of $\overline{\mathbf Q}$. As $M$
contains $\sigma_0(L)$, $g_2(\Lambda)$, $g_3(\Lambda)$ and the periods $c_\tau$, we have
$\rho'\circ\sigma_0=\rho\circ\sigma_0$ and $\Lambda^{\rho'}=\Lambda^\rho$. The lattices $N$ and $N'$
are self-dual and $ab$ is even, so by Lemma~\ref{is-growth-conjugate-translate} each $U_{y_l}F_0$ has
algebraic Taylor coefficients, $(U_{y_l}F_0)^{\rho'}$ is entire with the images under $\rho'$ as
Taylor coefficients, and, writing $|z|^2=\langle z,z\rangle$,
\[
 |(U_{y_l}F_0)^{\rho'}(z)|^2\le s_0(a,b)\,\|F_0^{\rho'}\|^2_{\Lambda^\rho_{N,N'}}\,
 e^{\pi|z|^2/A(\Lambda^\rho)},\qquad s_0(a,b)=3(2fab\,r_E+1)^{fab}.
\]
Put $F'=\sum_l\rho'(s_l)(U_{y_l}F_0)^{\rho'}$. It is entire, and its Taylor coefficients are the
images under $\rho'$ of those of $F$; so its jets at the conjugated periods $\rho(c_\tau)$ are the
images of the jets of $F$, and $\rho(\sigma_0(J(m)))=j_{(\rho(c_\tau))}(F')(m)$ at the exponents
above. By the triangle inequality and (b),
\[
 |F'(z)|^2\le\big(\beta c^{2[E:\mathbf Q]ab}\big)^2s_0(a,b)\,\|F_0^{\rho'}\|^2_{\Lambda^\rho_{N,N'}}\,
 e^{\pi|z|^2/A(\Lambda^\rho)} .
\]

(d) \emph{The vacuum norm.} By Lemma~\ref{is-growth-theta-facts}(4), $F_0$ and the vacuum theta
vector $\theta_0^{\otimes ab}$ of levels $i$ and $i+1$ are theta vectors up to units for $(t_0,N,N')$
and $(t_0,\mathcal O_E^a,\mathcal O_E^b)$. Let $t'$ be an unkilled two-torsion point of $\Lambda^\rho$.
By Lemma~\ref{is-growth-theta-facts}(2), applied to both,
$\|F_0^{\rho'}\|^2_{\Lambda^\rho_{N,N'}}=n(\Lambda^\rho,t')^{fab}
=\|(\theta_0^{\otimes ab})^{\rho'}\|^2_{\Lambda^\rho_{\mathcal O_E^a,\mathcal O_E^b}}$. The Taylor
coefficients of $\theta_0^{\otimes ab}$ lie in $M$, so its conjugate under $\rho'$ is its conjugate
under $\rho$, and the last norm is $e^{ab\,\ell_\theta(\rho)}$ by the definition of $\ell_\theta$. This
is the independence of the pair in Lemma~\ref{is-vacuum-norm-bound}.

(e) \emph{The growth bound.} With $r_\tau^2=\pi/A(\Lambda^\rho)$ we have
$|z|_r^2=\pi|z|^2/A(\Lambda^\rho)$, so (c) and (d) give $|F'(z)|\le S\exp(|z|_r^2/2)$ with
$S=\beta c^{2[E:\mathbf Q]ab}\sqrt{s_0(a,b)}\,e^{ab\,\ell_\theta(\rho)/2}$. The conjugated periods
satisfy $|\rho(c_\tau)r_\tau|^2=\pi|\rho(c_\tau)|^2/A(\Lambda^\rho)=1/\Omega_\tau(\rho)$, by the
definition of $\Omega_\tau(\rho)$.
% Step 2 follows GlobalTheory.ColNorms.block: hdiff and hJF (jet_finset_linear_combination,
% jetC_const_mul, JetsInField); (a) is T.exists_finset_latSum_thInd_eq_sum with IndAdd from
% CosetRefinement and PeriodLattices, whose conditions hconj and AAUnit on Sall sel mu u are
% Core.Ht2.BadPrimesS.bad_primes_sharp and T.periodLatticeIdeal_trivial_at_badPrimes; (b) is T.cosetIndicatorTheta_torsionTranslates and
% T.presentation_closed_under_bounded_combination (with n = |R| beta c^{[E:Q]ab}, then b1);
% (c) is T.conj_torsionTranslateCombination_entire_gaussianGrowth (ConjComb.conj_pres: sq_sum_le over
% T.conj_torsion_translate_vacuum_entire_gaussian_growth) with hJ' (derivAtZero_reindex_eq_taylor_coefficient);
% (d) is T.exists_algebraicNumbers_extension_conj_vacuum_normSq with
% T.thetaVectorUpToUnits_thMoved_of_thZero (hnorm', hlam); (e) is hblock, and Step 3 uses
% T.jetOperatorBound_at_embedding_of_complexOpNormSq.

\emph{Step 3: the block bound.} Proposition~\ref{is-block-bound} over $\mathbf C$, transported to the
embedding $\rho$ through $\sigma_0$, gives $\|\rho\,\mathrm{Jet}(J,\ell(\nu))\varphi\|^2\le
t(b,\nu,S,\Omega(\rho))\|\rho\varphi\|^2$, and Step 1 bounds the entry by the same quantity.

\emph{Step 4: the numerical bound.} Put $G=g_0g_1$, $A=fm^2N_m^2$ and $g=fab$, so that
$c=GT_u$, $U=A\log c$, and $1\le g\le A$ because $1\le a,b\le N_m$ and $m\geq1$. Then
$S^2=G^{2CA}c^{8g}s_0(a,b)\,e^{ab\,\ell_\theta(\rho)}$. Since $G\geq1$ and $T_u\geq1$,
$G^{2CA}\le c^{2CA}$ and $c^{8g}\le c^{8A}$; since $3\le8^g$,
$s_0(a,b)=3(2gr_E+1)^g\le(8+32g\,r_E/2)^g$. Hence
\[
 S^2\le\exp\big(ab\,\ell_\theta(\rho)+fab\log(8+32fab\,r_E/2)+2(C+4)U\big).
\]
For each $\tau$,
$\Omega_\tau^{-|\nu_\tau|}e\,(2|\nu_\tau|+2)^{2ab}\exp(\sum_i\varphi(\nu_{\tau,i}))=
\exp\big(\sum_i\varphi(\nu_{\tau,i})-|\nu_\tau|\log\Omega_\tau+1+2ab\log(2|\nu_\tau|+2)\big)$.
Multiplying these bounds, with $C_K=C+4$ and $ab=N(i)N(i+1)$, gives
$t(b,\nu,S,\Omega(\rho))\le e^{k_{i+1}(\rho)}$.
\end{proof}

\subsection{Vacuum norms and the definite-side condition}\label{subsection-is-vacuum-norms}

\begin{lemma}\label{is-transference}
Let $V$ be a $\mathbf Q$-vector space of dimension $M\geq1$, let $T$ be a positive definite
symmetric bilinear form on $V$, and let $\Lambda\subset V$ be the $\mathbf Z$-span of a basis of $V$.
Let $\lambda>0$, and suppose that $T(u,u)\geq\lambda^2$ for every nonzero $u\in V$ with
$T(u,y)\in\mathbf Z$ for all $y\in\Lambda$. Then $\Lambda$ contains $M$ linearly independent vectors
$w$ with $T(w,w)\le(M^2/\lambda)^2$.
\end{lemma}
\leanref{Development/AbelianVarietyTheory/Proofs/CuspFormEvaluations/TraceHermNondegenerate.lean}{GlobalTheory.ShortVectors.transference_rat}

With $T$ extended to $V\otimes\mathbf R$, $\Lambda$ is a lattice in a Euclidean space, and its dual
lattice consists of the $u\in V$ of the hypothesis, because $T$ is rational. Banaszczyk's
transference theorem $\lambda_i(\Lambda)\lambda_{M-i+1}(\Lambda^\#)\le M$ for the successive minima of
a lattice and of its dual \cite[Theorem 2.1]{Ban93}, with $i=M$, gives $M$ linearly independent
vectors $w\in\Lambda$ with $T(w,w)\le M^2/\lambda^2$, which is at most $(M^2/\lambda)^2$.

\begin{lemma}\label{is-selfdual-classes}
Let $E$ be a CM field whose maximal totally real subfield has degree $f$, let $r_E$ be the root
discriminant of $E$, and let $n\geq0$. Write $\langle x,y\rangle=\sum_i\bar x_iy_i$ on $E^n$. Up to
the action of the unitary group of $\langle\ ,\ \rangle$, there are at most
$(3(2f)^2n^4r_E^2)^{6fn^2}$ finitely generated $\mathcal O_E$-submodules $L\subset E^n$ that span $E^n$
and are self-dual, that is, $L=\{y:\langle x,y\rangle\in\mathcal O_E\text{ for all }x\in L\}$.
\end{lemma}
\leanref{Development/AbelianVarietyTheory/Proofs/FoundationProofs/ExistsFinsetSelfDualClassesRootDiscr.lean}{GlobalTheory.exists_finset_selfDual_classes_rootDiscr}
\begin{proof}
Put $B=4f^3n^4r_E^2$.

\emph{Step 1: a bounded basis.} Every such $L$ contains an $E$-basis $v_1,\dots,v_n$ of $E^n$ with
$|\sigma\langle v_i,v_k\rangle|\le B$ for every embedding $\sigma$ of $E$ and all $i,k$. For
$x,y\in E^n$ put $T(x,y)=\operatorname{Tr}_{E/\mathbf Q}\langle x,y\rangle$. Complex conjugation
preserves the trace, so $T$ is a symmetric $\mathbf Q$-bilinear form. Since
$\sigma(\bar a)=\overline{\sigma(a)}$ for every embedding $\sigma$ of the CM field $E$,
$\sigma\langle x,x\rangle=\sum_k|\sigma(x_k)|^2$ and $T(x,x)=\sum_\sigma\sum_k|\sigma(x_k)|^2$, the sum
over the $2f$ embeddings. So $T$ is positive definite, and $T(x,x)\geq2\sum_k|\sigma(x_k)|^2$ for each
$\sigma$, because $\sigma$ and $\bar\sigma\ne\sigma$ give the same inner sum. As a $\mathbf Z$-module
$L$ is finitely generated and torsion free, and it contains an $E$-basis of $E^n$; so it is the
$\mathbf Z$-span of a $\mathbf Q$-basis of $E^n$, of $M=2fn$ vectors.

Let $\mathfrak D$ be the different of $E$, so that
$\mathfrak D^{-1}=\{a\in E:\operatorname{Tr}_{E/\mathbf Q}(a\mathcal O_E)\subset\mathbf Z\}$ and
$N(\mathfrak D)=|d_E|$ for the discriminant $d_E$, and let $u\ne0$ in $E^n$ satisfy
$T(u,y)\in\mathbf Z$ for every $y\in L$. For $y\in L$ and $c\in\mathcal O_E$, $cy\in L$, so
$\operatorname{Tr}(c\langle u,y\rangle)\in\mathbf Z$; hence $\langle u,y\rangle\in\mathfrak D^{-1}$.
For $\delta\in\mathfrak D$ and $y\in L$, $\langle y,\delta u\rangle=\delta\,\overline{\langle u,y\rangle}$
lies in $\mathfrak D\,\overline{\mathfrak D}^{-1}=\mathcal O_E$, because complex conjugation preserves
$\mathfrak D$; so $\delta u\in L$ by self-duality. Then
$\operatorname{Tr}(c\delta\langle u,u\rangle)=T(u,c\delta u)\in\mathbf Z$ for $c\in\mathcal O_E$, so
$\langle u,u\rangle\mathfrak D\subset\mathfrak D^{-1}$, that is, $\langle u,u\rangle\in\mathfrak D^{-2}$.
Every $\sigma\langle u,u\rangle$ is positive, since $u\ne0$. A nonzero element of $\mathfrak D^{-2}$
has norm of absolute value at least $N(\mathfrak D)^{-2}$, so the inequality of arithmetic and
geometric means gives
\[
 T(u,u)=\sum_\sigma\sigma\langle u,u\rangle\geq2f\,\big|N_{E/\mathbf Q}\langle u,u\rangle\big|^{1/2f}
 \geq2f\,|d_E|^{-1/f}=2f/r_E^2.
\]
Lemma~\ref{is-transference}, with $\lambda^2=2f/r_E^2$ and $M=2fn$, gives $M$ linearly independent
vectors $w\in L$ with $T(w,w)\le M^4/\lambda^2=8f^3n^4r_E^2$, so that $\sum_k|\sigma(w_k)|^2\le B$ for
every $\sigma$. They span $E^n$ over $\mathbf Q$, hence over $E$, and $n$ of them form an $E$-basis
$v_1,\dots,v_n$. Finally
$|\sigma\langle v_i,v_k\rangle|\le\sum_l|\sigma(v_{i,l})|\,|\sigma(v_{k,l})|\le\frac12\sum_l
\big(|\sigma(v_{i,l})|^2+|\sigma(v_{k,l})|^2\big)\le B$.

\emph{Step 2: the type.} For such a basis let $A=(\langle v_i,v_k\rangle)_{i,k}$, which has entries
in $\mathcal O_E$ because $L$ is self-dual and so integral, and $\det A\ne0$; and let
$P=\{(\langle v_i,x\rangle)_i:x\in L\}\subset\mathcal O_E^n$, which contains the module $A\mathcal O_E^n$
spanned by the columns of $A$, the vectors $(\langle v_i,v_k\rangle)_i$. Call
$(A,P)$ the type of $L$. Two lattices $L,L'$ of the same type, with bases $v,v'$, are isometric: the
$E$-linear map $\gamma$ with $\gamma v_i=v'_i$ preserves $\langle\ ,\ \rangle$ because the Gram
matrices agree, so $\langle v'_i,\gamma x\rangle=\langle v_i,x\rangle$ for all $x$; as
$x\mapsto(\langle v_i,x\rangle)_i$ is injective, $\gamma L=\{y:(\langle v'_i,y\rangle)_i\in P\}=L'$.

\emph{Step 3: counting types.} The entries of $A$ are elements $a\in\mathcal O_E$ with
$|\sigma(a)|\le B$ for every $\sigma$, and there are at most $(1+2B)^{2f}$ of these. Indeed, choose
one embedding from each pair $\{\sigma,\bar\sigma\}$, which identifies $E\otimes\mathbf R$ with
$\mathbf C^f$. For distinct such $a,a'$, $|N_{E/\mathbf Q}(a-a')|\geq1$, so $|\sigma(a-a')|\geq1$ for
some chosen $\sigma$. Hence the polydiscs of radius $\frac12$ about the images of these $a$ are
disjoint; they lie in the polydisc of radius $B+\frac12$ about $0$, and comparing volumes bounds their
number by $\big((B+\frac12)/\frac12\big)^{2f}=(1+2B)^{2f}$. So there are at most $(1+2B)^{2fn^2}$
matrices $A$.

Fix such an $A$ with $\delta=\det A\ne0$. Then $|\sigma(\delta)|\le n!\,B^n\le(nB)^n$ for every
$\sigma$, so $|N_{E/\mathbf Q}(\delta)|\le(nB)^{2fn}$. The index of $A\mathcal O_E^n$ in
$\mathcal O_E^n$ is $|N_{E/\mathbf Q}(\delta)|$: the index of the image of a lattice under an
injective $\mathbf Q$-linear map is the absolute value of its determinant, and multiplication by $A$
on $E^n$ has determinant $N_{E/\mathbf Q}(\det A)$ over $\mathbf Q$. Let $P$ be an
$\mathcal O_E$-module with $A\mathcal O_E^n\subset P\subset\mathcal O_E^n$. For each prime
$\mathfrak p$ the localization $P_{\mathfrak p}$ is free of rank $n$ over the discrete valuation ring
$\mathcal O_{E,\mathfrak p}$, so $P_{\mathfrak p}/\delta P_{\mathfrak p}$ is generated by $n$ elements;
as $P/\delta P$ is the product of these localizations over the primes dividing $\delta$, the Chinese
remainder theorem shows that $P/\delta P$ is generated by $n$ elements. So
$P=\sum_j\mathcal O_Ex_j+\delta P$ for some $x_1,\dots,x_n\in P$. Since
$\delta\mathcal O_E^n=A\operatorname{adj}(A)\mathcal O_E^n\subset A\mathcal O_E^n$, we have
$\delta P\subset A\mathcal O_E^n$, and $P=\sum_j\mathcal O_Ex_j+A\mathcal O_E^n$ is determined by the
images of $x_1,\dots,x_n$ in $\mathcal O_E^n/A\mathcal O_E^n$. So there are at most
$|N_{E/\mathbf Q}(\delta)|^n\le(nB)^{2fn^2}$ modules $P$ for each $A$, and at most
$(1+2B)^{2fn^2}(nB)^{2fn^2}$ types.

\emph{Step 4: the bound.} Put $p=(2f)^2n^4r_E^2$, so that $B=fp$ and $f^2n\le p$. Then
$(1+2B)\,nB\le3nB^2=3nf^2p^2\le3p^3$, and
$(1+2B)^{2fn^2}(nB)^{2fn^2}\le(3p^3)^{2fn^2}\le(3p)^{6fn^2}$.
\end{proof}

\begin{lemma}\label{is-class-counts}
There is a constant $E\geq0$ with the following property. Let $D$ be a definite-side datum with size
data and kernel data as in Proposition~\ref{is-kernel-entry-bound}, at the base tuple, and put
$U=fm^2N_m^2\log(g_0g_1T_u)$ with $N_m=N(m)$. Then:
\begin{enumerate}
\item for every level $j\le m$ and every member of the chain at level $j$, the number of classes of
  lattices is at most $e^{EU}$, and so are the number of classes of the evaluation space and the
  number of top classes;
\item for every $j\le m-2$ the number of members of the chain at level $j$ is at most $e^{EU}$;
\item for every $n\le N_m$ and every lattice $N\subset E^{n}$, the coupling set of $N$ has at most
  $(g_0g_1)^{2fn}\le e^{2U}$ elements.
\end{enumerate}
\end{lemma}
\leanref{Development/AbelianVarietyTheory/Proofs/DecoratedLatticeProofs/CardOrbitsLeExp.lean}{GlobalTheory.ChainClassNumbersS.classNumbers_leS}
\leanref{Development/AbelianVarietyTheory/Proofs/DecoratedLatticeProofs/CardOrbitsLeExp.lean}{GlobalTheory.ChainClassNumbersS.card_orbits_le_exp}
\leanref{Development/AbelianVarietyTheory/Proofs/DecoratedLatticeProofs/CardMatrixStepIndicesLeExpStatement.lean}{GlobalTheory.T.card_matrixStep_indices_le_exp}
\leanref{Development/AbelianVarietyTheory/Proofs/DecoratedLatticeProofs/NcardYbSetLe.lean}{GlobalTheory.CouplingCosetsS.ncard_ybSet_le}
\leanref{Development/AbelianVarietyTheory/Proofs/DivisorProofs/CharacterSumLocalConstantTwelfth.lean}{GlobalTheory.CouplingCosets.finite_ncard_ybSet}
\begin{proof}
Write $d=[E:\mathbf Q]=2f$, and recall from the size data that $L=x^{\varepsilon}\geq1$,
$\log r_E\le L$, $fN_m\le e^L$ and $e^L\le g_0$.

(1) At the base tuple the set of twisted primes is empty, so the bad set is the set of places above
$g_0g_1$. Fix $n\le N_m$, with $n\geq1$, and call a \emph{decorated lattice} of rank $n$ a self-dual
$\mathcal O_E$-lattice $N\subset E^n$ for the standard Hermitian form, together with a nonzero linear
form $N/PN\to\mathcal O_E/P$ for each place $P$ of a set $S_T$ of places above $T_u$, and an isomorphism
$N/wN\cong(\mathcal O_E/w)^n$ for each place $w$ above $g_0g_1$. The unitary group of $E^n$ acts on
decorated lattices. Each set of classes in (1) is the set of orbits of a group $G'$ on a set $D'$ with
an injection $p$ of $D'$ into the decorated lattices of some rank $n\le N_m$, such that for every
unitary $g$ and $x\in D'$ there is $g'\in G'$ with $p(g'x)=g\,p(x)$. Then $p$ induces an injection
of the $G'$-orbits into the unitary orbits of decorated lattices, so the number of classes is at most
the number of the latter.

By Lemma~\ref{is-selfdual-classes} the underlying lattices fall into at most
$R=(3d^2n^4r_E^2)^{3dn^2}$ unitary classes. A fixed lattice $N$ carries at most
$\prod_{P\in S_T}NP^{\,n}\prod_{w\mid g_0g_1}Nw^{\,n^2}$ decorations, the number of the linear forms
and isomorphisms. The places above $T_u$ and $g_0g_1$ have norms whose product, with these
exponents, is at most $(g_0g_1T_u)^{fn^2}$. So the number of classes is at most
$h=R\,(g_0g_1T_u)^{fn^2}$.

Now $\log(3d^2n^4r_E^2)=\log12+2\log(fn)+2\log n+2\log r_E\le\log12+2L+2L+2L\le9L$, because
$fn\le fN_m\le e^L$, $n\le e^L$ and $\log12\le3\le3L$. Hence
$\log h\le6fn^2\cdot9L+fn^2\log(g_0g_1T_u)\le55fn^2\log(g_0g_1T_u)$, because
$L\le\log g_0\le\log(g_0g_1T_u)$. As $n\le N_m$ and $m\geq1$, this is at most $55U$.

(2) For $j=0$ there is one member. For $j\geq1$ the members of the chain at level $j$ are indexed by
level data whose matrices, of some size $n\le N_m$, have entries integral at the places $w$ above
$g_0g_1$. Two members whose matrices have the same reductions modulo every such $w$ coincide, by the
condition on their stabilizers. So the number of members is at most
$\prod_{w\mid g_0g_1}|M_n(\mathcal O_E/w)|=\prod_wNw^{\,n^2}=(g_0g_1)^{fn^2}\le(g_0g_1)^{fN_m^2}\le e^U$,
since $\prod_{w\mid g_0g_1}Nw=(g_0g_1)^f$.

(3) At the base tuple every place of $S$ lies in the level set, so the coupling set of $N$ is the set
of cosets of $N$ in its scaling by the inverse uniformizers at the level places. Their number is the
norm of the product of these places to the power $n$, which divides $(g_0g_1)^{[E:\mathbf Q]n}$, and
$[E:\mathbf Q]=2f$. Since $g_0g_1\le g_0g_1T_u$ and $n\le N_m\le m^2N_m^2$,
$(g_0g_1)^{2fn}\le(g_0g_1T_u)^{2fm^2N_m^2}=e^{2U}$.

Take $E=55+1+2$.
\end{proof}

\begin{lemma}\label{is-vacuum-averaging}
Let $E$, $\mathcal T$ and $\mathfrak A_0$ be those of the host of a definite-side datum, and $\Lambda$ a
CM lattice of the host (Definition~\ref{is-vacuum-theta-vectors}). Call a subgroup of $V_{a,b}$
\emph{standard} if it is $\Lambda_v$ for an orthonormal basis $v$, or $\Lambda_{L,L'}$ for
$\mathcal O_E$-lattices $L\subset E^a$ and $L'\subset E^b$ that are self-dual for $\sum_i\bar x_iy_i$.
Let $\Gamma_1,\Gamma_2,\Gamma_3\subset V_{a,b}$ be standard subgroups and $G$ a theta function for
$\Gamma_1$. Then:
\begin{enumerate}
\item the theta functions for $\Gamma_1$ form a space of dimension at most $1$;
\item $\|U_xG\|_{\Gamma_1}=\|G\|_{\Gamma_1}$ for every $x\in V_{a,b}$;
\item $I_{\Gamma_1\to\Gamma_2}G$ is a theta function for $\Gamma_2$, and
  $\|I_{\Gamma_1\to\Gamma_2}G\|_{\Gamma_2}=\|G\|_{\Gamma_1}$;
\item $I_{\Gamma_2\to\Gamma_3}I_{\Gamma_1\to\Gamma_2}G=\mu\,I_{\Gamma_1\to\Gamma_3}G$ for some $\mu$ with
  $\mu^8=1$.
\end{enumerate}
\end{lemma}
\leanref{Development/AbelianVarietyTheory/Proofs/FoundationProofs/ThetaSpaceSelfDualPairAtMostLineStatement.lean}{GlobalTheory.T.theta_space_selfDual_pair_at_most_line}
\leanref{Development/AbelianVarietyTheory/Proofs/FoundationProofs/IsHermSesqOf.lean}{GlobalTheory.ConjNorm.unitary_all}
\leanref{Development/AbelianVarietyTheory/Proofs/FoundationProofs/AdmissibleCharactersExist.lean}{GlobalTheory.TripleIntertwiner.triple_constant}

A standard subgroup is a lattice in $V_{a,b}$ that is self-dual for $\operatorname{Im}H$ and on which
$\frac{\sqrt d}2H(x,x)$ takes integer values. Statement (2) holds because
$|(U_xG)(z)|^2e^{-\pi H(z,z)}=|G(z+x)|^2e^{-\pi H(z+x,z+x)}$, and the mean of a function periodic for
$\Gamma_1$ over a fundamental domain does not change under translation. For
such lattices the theta functions form the lattice models of the Heisenberg representation attached
to $H$, and the averaging operators are the intertwining operators between them; (1), (3) and (4) are
the irreducibility of these models, the unitarity of the intertwiners, and the fact that a composite
of intertwiners differs from the direct one by an eighth root of unity (compare
\cite[Chapter~1]{LV80}).

\begin{lemma}\label{is-vacuum-tensor-norm}
Let $d\equiv7\pmod8$ be a prime and $K=\mathbf Q(\sqrt{-d})$. Let $E$ be a CM field of degree $2f$
containing $K$, $\mathcal T$ a set of $f$ embeddings of $E$ into $\mathbf C$ that agree on $K$,
$\Lambda\subset\mathbf C$ a lattice of covolume $A$, and $v_1,\dots,v_r$ an orthonormal basis of
$M_{b\times a}(E)$, $r=fab$ (Definition~\ref{is-vacuum-theta-vectors}).
\begin{enumerate}
\item The vectors $\iota(v_1),\dots,\iota(v_r)$ form an orthonormal basis of $V_{a,b}$ for
  $\langle\ ,\ \rangle$.
\item Let $g$ be a continuous function on $\mathbf C$ with $|(U^\Lambda_\omega g)(z)|=|g(z)|$ for all
  $\omega\in\Lambda$ and $z\in\mathbf C$, and let $n_g$ be the mean of $|g(z)|^2e^{-\pi|z|^2/A}$ over a
  fundamental domain of $\Lambda$. Then $\|g_v\|^2_{\Lambda_v}=n_g^{\,r}$.
\item Suppose that $m(\omega)=\frac{\sqrt d}2|\omega|^2/A$ is an integer for every $\omega\in\Lambda$, and
  let $g$ be an entire function with $(-1)^{m(\omega)}U^\Lambda_\omega g=g$ for every $\omega\in\Lambda$.
  Then $g_v$ is a theta function for $\Lambda_v$.
\end{enumerate}
\end{lemma}
\leanref{Development/AbelianVarietyTheory/Proofs/FoundationProofs/OrthonormalRealBasisComplexCoordinatesStatement.lean}{GlobalTheory.T.orthonormal_realBasis_complex_coordinates}
\leanref{Development/AbelianVarietyTheory/Proofs/FoundationProofs/MeanSquareOrthoProductThetaEqProdStatement.lean}{GlobalTheory.T.meanSquare_orthoProductTheta_eq_prod}
\leanref{Development/AbelianVarietyTheory/Proofs/FoundationProofs/OrthonormalProductThetaMemThetaSpaceStatement.lean}{GlobalTheory.T.orthonormalProductTheta_mem_thetaSpace}
\begin{proof}
(1) Let $\psi$ be the common restriction of the types to $K$. There are $[E:K]=f$ embeddings of $E$
extending $\psi$, so these are the types, and $\psi(\operatorname{Tr}_{E/K}u)=\sum_\tau\tau(u)$ for
$u\in E$. Since $E$ is a CM field, $\tau(\bar x)=\overline{\tau(x)}$. Hence
$\langle\iota(v_k),\iota(v_l)\rangle=\psi\big(\operatorname{Tr}_{E/K}\sum_{\beta,\alpha}(v_k)_{\beta\alpha}
\overline{(v_l)_{\beta\alpha}}\big)=\delta_{kl}$, and $r=\dim_{\mathbf C}V_{a,b}$.

(2) By (1), $c(z)=(\langle z,\iota(v_k)\rangle)_k$ is a unitary isomorphism of $V_{a,b}$ onto
$\mathbf C^r$. It carries $\Lambda_v$ onto $\Lambda^r$, $H(z,z)$ to $A^{-1}\sum_k|c_k|^2$, $g_v$ to
$\prod_kg(c_k)$ and Lebesgue measure to Lebesgue measure. A product of fundamental domains of $\Lambda$
is a fundamental domain of $\Lambda^r$, and on it $|g_v|^2e^{-\pi H}$ is the product of the functions
$|g(c_k)|^2e^{-\pi|c_k|^2/A}$. These are periodic for $\Lambda$ by the hypothesis on $g$, as in
Lemma~\ref{is-vacuum-averaging}(2), so Fubini's theorem gives the mean $n_g^{\,r}$.

(3) Let $x\in\Lambda_v$. Then $c_k(x)\in\Lambda$, $U_xg_v=\prod_k(U^\Lambda_{c_k(x)}g)(c_k)$ since $c$ is
unitary, and $\frac{\sqrt d}2H(x,x)=\sum_km(c_k(x))$ is an integer, so that
$\alpha_{\sqrt d/2}(x)=\prod_k(-1)^{m(c_k(x))}$. Hence $\alpha_{\sqrt d/2}(x)U_xg_v=g_v$, and $g_v$ is
entire.
\end{proof}

\begin{lemma}\label{is-vacuum-genus-one}
Let $d\equiv7\pmod8$ be a prime, $K=\mathbf Q(\sqrt{-d})$, and $\Lambda\subset\mathbf C$ a lattice with
complex multiplication by $\mathcal O_K$, of covolume $A$.
\begin{enumerate}
\item $\Lambda$ has an unkilled two-torsion point (Definition~\ref{is-vacuum-theta-vectors}), any two
  differ by an element of $\Lambda$, and $\vartheta_{\Lambda,t+\omega}=\vartheta_{\Lambda,t}$ for
  $t\notin\Lambda$ and $\omega\in\Lambda$.
\item Suppose that $\Lambda=\delta\mathfrak c$, read in $\mathbf C$ through $\sqrt{-d}\mapsto i\sqrt d$,
  for a nonzero ideal $\mathfrak c$ of $\mathcal O_K$ of norm prime to $2d$ and a $\delta$ with
  $\delta^{12}=\Delta(\mathfrak c)$, and let $t$ be an unkilled two-torsion point of $\Lambda$. Then for
  every $\omega\in\Lambda$ the number $m(\omega)=\frac{\sqrt d}2|\omega|^2/A$ is an integer and
  $(-1)^{m(\omega)}U^\Lambda_\omega\vartheta_{\Lambda,t}=\vartheta_{\Lambda,t}$.
\item Write $\Lambda=\omega_1(\mathbf Z+\mathbf Z\tau)$ with $\operatorname{Im}\tau>0$, let $2t\in\Lambda$
  with $t\notin\Lambda$, and for $c=(c_1,c_2)\in\{0,1\}^2$ let
  \[
   \theta_c(u,\tau)=\sum_{n\in\mathbf Z}\exp\Big(\pi i\big(n+\tfrac{c_1}2\big)^2\tau
   +2\pi i\big(n+\tfrac{c_1}2\big)\big(u+\tfrac{c_2}2\big)\Big).
  \]
  Then exactly one $c\in\{(0,0),(0,1),(1,0)\}$ has $\theta_c(t/\omega_1,\tau)=0$, and for it
  $n(\Lambda,t)^{-1}=\sqrt{2\operatorname{Im}\tau}\,|\theta_c(0,\tau)|^2$.
\item $n(\bar\Lambda,\bar t)=n(\Lambda,t)$ for $t\notin\Lambda$, where $\bar\Lambda$ is the complex
  conjugate of $\Lambda$.
\end{enumerate}
\end{lemma}
\leanref{Development/AbelianVarietyTheory/Definitions/GlobalTwisting/ExistsPowPrincipalFractionalIdeal.lean}{GlobalTheory.Glob.exists_isThetaZero}
\leanref{Development/AbelianVarietyTheory/Definitions/GlobalTwisting/ExistsPowPrincipalFractionalIdeal.lean}{GlobalTheory.Glob.isThetaZero_unique_lattice}
\leanref{Development/AbelianVarietyTheory/Proofs/FoundationProofs/ExistsTwoTorsionCMPointAllConjugates.lean}{GlobalTheory.TwoTorsionPt.isTwoTorsionCMPoint_of_odd_multiple}
\leanref{Development/AbelianVarietyTheory/Proofs/FoundationProofs/CmThetaTaylorCoeffAlgebraicAndConjugate.lean}{GlobalTheory.ThetaConj.two_torsion_unique}
\leanref{Development/AbelianVarietyTheory/Proofs/FoundationProofs/CmThetaTaylorCoeffAlgebraicAndConjugate.lean}{GlobalTheory.ThetaConj.evenTheta_add_lattice_mem}
\leanref{Development/AbelianVarietyTheory/Proofs/FoundationProofs/HalfPeriodThetaDifferentiableAndSignedTranslateIStatement.lean}{GlobalTheory.T.halfPeriodTheta_differentiable_and_signedTranslate_invariant}
\leanref{Development/AbelianVarietyTheory/BlockTheory/GlobalSignature/JacobiThetaHalfCharacteristic.lean}{GlobalTheory.IndefiniteRowSideFullHecke.JacobiThetaHalfCharacteristic.NormTheta0}
\leanref{Development/AbelianVarietyTheory/Proofs/GlobalInterpolation/MeanSquareNormalizedThetaTwoTorsion.lean}{GlobalTheory.ThetaNullNormSquares.normTheta0}
\leanref{Development/AbelianVarietyTheory/Proofs/FoundationProofs/MeanSquareEvenThetaConjugateLatticeStatement.lean}{GlobalTheory.T.meanSquare_evenTheta_conjugateLattice}
\begin{proof}
(1) As an $\mathcal O_K$-module $\Lambda$ is projective of rank one, so $\Lambda=\delta\mathfrak c$
for some $\delta\in\mathbf C^\times$ and a nonzero ideal $\mathfrak c$ of $\mathcal O_K$; since every
ideal class contains an ideal prime to $2$, the norm $N=N\mathfrak c$ may be taken odd. Every element
of $\Lambda$ is $\delta$ times an algebraic integer. Put $t=\delta N/2$. Then $2t=\delta N\in\Lambda$,
since $N\in\mathfrak c$. If $\omega_Kt\in\Lambda$, then $x=\omega_KN/2$ is an algebraic integer, and so
is $x+\bar x=N/2$, because $\omega_K+\bar\omega_K=1$; this is impossible, as $N$ is odd. In the same
way $\bar\omega_Kt\notin\Lambda$, so $t$ is an unkilled two-torsion point.

For the uniqueness, $\omega_K\bar\omega_K=(1+d)/4=2k$ with $k\in\mathbf Z$, because $d\equiv7\pmod8$,
and $\omega_K\Lambda\subset\Lambda$ and $\bar\omega_K\Lambda=(1-\omega_K)\Lambda\subset\Lambda$. Let $t$
and $t'$ be unkilled two-torsion points, and put $a=\omega_Kt$ and $b=\bar\omega_Kt$. Then $2a$ and
$2b$ lie in $\Lambda$, $a,b\notin\Lambda$, $a+b=t$, and $\omega_K\bar\omega_Kt=k\cdot2t\in\Lambda$. Also
$a-b\notin\Lambda$: otherwise $a=\omega_K(a-b)+2\omega_K\bar\omega_Kt$ would lie in $\Lambda$. So $0$,
$a$, $b$ and $a+b$ represent the four classes of $\frac12\Lambda/\Lambda$, and $t'$ lies in one of them.
It is not the class of $0$, since $\omega_Kt'\notin\Lambda$. If $t'-a\in\Lambda$, then
$\bar\omega_Kt'\equiv\bar\omega_K\omega_Kt\equiv0$ modulo $\Lambda$, which is excluded; if
$t'-b\in\Lambda$, then $\omega_Kt'\equiv0$ in the same way. Hence $t'-t\in\Lambda$. For
$\omega\in\Lambda$, $U^\Lambda_{t+\omega}\sigma^*_\Lambda$ is a constant
times $U^\Lambda_tU^\Lambda_\omega\sigma^*_\Lambda$, and $U^\Lambda_\omega\sigma^*_\Lambda=\pm\sigma^*_\Lambda$
by the quasi-periodicity of $\sigma_\Lambda$ (proof of Lemma~\ref{pd-analytic-translation-system}); the
constants disappear in the normalization at $0$.

(2) The covolume of $\delta\mathfrak c$ is $|\delta|^2N\mathfrak c\sqrt d/2$, so $m(\delta x)=N(x)/N\mathfrak c$
is an integer for $x\in\mathfrak c$. Since $U^\Lambda_\omega U^\Lambda_t=e^{-2\pi i\operatorname{Im}H_\Lambda(\omega,t)}
U^\Lambda_tU^\Lambda_\omega$, the operator $U^\Lambda_\omega$ multiplies $\vartheta_{\Lambda,t}$ by the
sign of $U^\Lambda_\omega\sigma^*_\Lambda=\pm\sigma^*_\Lambda$ times
$(-1)^{\operatorname{Im}H_\Lambda(\omega,2t)}$, where $H_\Lambda(z,w)=z\bar w/A$. The identity is then a
computation in
$\Lambda/2\Lambda\cong\mathcal O_K/2\mathcal O_K$: $m(\omega)$ plus the indicator of $\omega\notin2\Lambda$
is additive modulo $2$, and it agrees with $\operatorname{Im}H_\Lambda(\omega,2t)$ modulo $2$ when $t$ is
the unkilled point.

(3) The function $u\mapsto\vartheta_{\Lambda,t}(\omega_1u)$ and the function
$u\mapsto e^{\pi u^2/(2\operatorname{Im}\tau)}\theta_c(u,\tau)/\theta_c(0,\tau)$, for the $c$ with
$\theta_c(t/\omega_1,\tau)=0$, are entire, take the value $1$ at $0$, have the same simple zeros, at
$t/\omega_1+\mathbf Z+\mathbf Z\tau$, and the same factor of automorphy for $\mathbf Z+\mathbf Z\tau$ up
to signs; their quotient is entire with periodic absolute value, so it is $1$. The three even
characteristics have their zeros at $\frac{1+\tau}2$, $\frac\tau2$ and $\frac12$ modulo
$\mathbf Z+\mathbf Z\tau$, and $t/\omega_1$ is one of these. With $z=\omega_1u$ and
$A=|\omega_1|^2\operatorname{Im}\tau$, the mean defining $n(\Lambda,t)$ is the mean over
$0\le\operatorname{Re}u<1$, $0\le\operatorname{Im}u<\operatorname{Im}\tau$ of
$|\theta_c(u,\tau)|^2e^{-2\pi(\operatorname{Im}u)^2/\operatorname{Im}\tau}/|\theta_c(0,\tau)|^2$. Integrating
over $\operatorname{Re}u$ leaves the diagonal terms of the series, and the sum of the resulting
Gaussians over $n$ unfolds to $\int_{\mathbf R}e^{-2\pi y^2/\operatorname{Im}\tau}dy
=\sqrt{\operatorname{Im}\tau/2}$. Hence $n(\Lambda,t)=(2\operatorname{Im}\tau)^{-1/2}|\theta_c(0,\tau)|^{-2}$.

(4) The product formula for $\sigma_\Lambda$ and the definition of $s_2(\Lambda)$ give
$\sigma^*_{\bar\Lambda}(z)=\overline{\sigma^*_\Lambda(\bar z)}$, and $A(\bar\Lambda)=A$, so
$\vartheta_{\bar\Lambda,\bar t}(z)=\overline{\vartheta_{\Lambda,t}(\bar z)}$. The substitution
$z\mapsto\bar z$ carries a fundamental domain of $\Lambda$ to one of $\bar\Lambda$ and preserves $|z|$
and Lebesgue measure.
\end{proof}

\begin{remark}\label{is-vacuum-unkilled-splitting}
The existence and uniqueness of the unkilled two-torsion point in Lemma~\ref{is-vacuum-genus-one}(1)
also follow from the splitting of $2$ in $K$. As an $\mathcal O_K$-module $\Lambda$ is projective of
rank one, so $\frac12\Lambda/\Lambda\cong\mathcal O_K/2\mathcal O_K$. Since $-d\equiv1\pmod8$,
$2\mathcal O_K=\mathfrak p\bar{\mathfrak p}$ with $\mathfrak p\ne\bar{\mathfrak p}$, and
$\mathcal O_K/2\mathcal O_K\cong\mathcal O_K/\mathfrak p\times\mathcal O_K/\bar{\mathfrak p}$ is a product
of two fields with two elements. As $\omega_K\bar\omega_K=(1+d)/4$ is even and
$\omega_K+\bar\omega_K=1$, $\omega_K$ lies in one of $\mathfrak p,\bar{\mathfrak p}$ and $\bar\omega_K$ in
the other. So multiplication by $\omega_K$ kills one factor and is the identity on the other, and
$\bar\omega_K$ does the same with the factors exchanged. The unkilled two-torsion points are those
whose image has both components nonzero, which is one class modulo $\Lambda$.
\end{remark}

\begin{lemma}\label{is-vacuum-eta-weber}
Let $\Lambda=\omega_1(\mathbf Z+\mathbf Z\tau)$ be a lattice with $\operatorname{Im}\tau>0$ and
$\Delta(\Lambda)=1$, of covolume $A$. Let $2t\in\Lambda$ with $t\notin\Lambda$, and let $c$ be the
characteristic with $\theta_c(t/\omega_1,\tau)=0$ of Lemma~\ref{is-vacuum-genus-one}(3). Let $\eta$ be
Dedekind's eta function, $\wp_\Lambda$ the Weierstrass function of $\Lambda$, and
$Y=\frac14\big(12\wp_\Lambda(t)^2-g_2(\Lambda)\big)$. Then:
\begin{enumerate}
\item $A=4\pi^2\operatorname{Im}\tau\,|\eta(\tau)|^4$;
\item $Y\ne0$ and $Y^{-3}=\pm\big(\theta_c(0,\tau)/\eta(\tau)\big)^{12}$;
\item if $j(\Lambda)$ is an algebraic integer, then so is $Y^{-3}$.
\end{enumerate}
\end{lemma}
\leanref{Development/AbelianVarietyTheory/Proofs/GlobalInterpolation/IsPeriodLatticeIff.lean}{GlobalTheory.Theta0MeanProof.log_area}
\leanref{Development/AbelianVarietyTheory/Proofs/GlobalInterpolation/ExistsIntegralPeriodElement.lean}{GlobalTheory.IntegralX.modulus}
\leanref{Development/AbelianVarietyTheory/Proofs/GlobalInterpolation/ExistsIntegralPeriodElement.lean}{GlobalTheory.IntegralX.integral}
\leanref{Development/AbelianVarietyTheory/Proofs/GlobalInterpolation/ThetaWeberStatement.lean}{GlobalTheory.ThetaWeberObj.thetaWeberStatement}
\leanref{Development/AbelianVarietyTheory/Proofs/GlobalInterpolation/LatticePeriodEtaObj.lean}{GlobalTheory.LatticePeriodEtaObj.latticePeriodEta_obj}
\begin{proof}
(1) The discriminant is homogeneous of degree $-12$, and
$\Delta(\mathbf Z+\mathbf Z\tau)=(2\pi)^{12}\eta(\tau)^{24}$ \cite[Chapter~I]{Sil94}. So
$1=\Delta(\Lambda)=(2\pi/\omega_1)^{12}\eta(\tau)^{24}$, hence $|\omega_1|^2=4\pi^2|\eta(\tau)|^4$, and
$A=|\omega_1|^2\operatorname{Im}\tau$.

(2) Let $e=\wp_\Lambda(t)$, and let $e'$ and $e''$ be the values of $\wp_\Lambda$ at the other two
points of order $2$ of $\mathbf C/\Lambda$. Since $4x^3-g_2(\Lambda)x-g_3(\Lambda)=4(x-e)(x-e')(x-e'')$,
$Y=(e-e')(e-e'')$. Write $\wp_\tau$ for the Weierstrass function of $\mathbf Z+\mathbf Z\tau$, so that
$\wp_\Lambda(z)=\omega_1^{-2}\wp_\tau(z/\omega_1)$. The differences of the values of $\wp_\tau$ at
$\frac12$, $\frac\tau2$ and $\frac{1+\tau}2$ are $\pm\pi^2$ times the fourth powers of the theta
constants $\theta_{c'}(0,\tau)$, $c'\in\{(0,0),(0,1),(1,0)\}$ \cite[Chapter~XXI]{WW27}, and the two
differences in $Y$ involve the two constants other than $\theta_c(0,\tau)$. By Jacobi's identity
$\theta_{(0,0)}(0,\tau)\theta_{(0,1)}(0,\tau)\theta_{(1,0)}(0,\tau)=2\eta(\tau)^3$ \cite[Chapter~XXI]{WW27},
\[
 Y^3=\pm\frac{\pi^{12}}{\omega_1^{12}}\Big(\frac{2\eta(\tau)^3}{\theta_c(0,\tau)}\Big)^{12}
 =\pm\Delta(\Lambda)\Big(\frac{\eta(\tau)}{\theta_c(0,\tau)}\Big)^{12},
\]
with $\theta_c(0,\tau)\ne0$ because $\eta(\tau)\ne0$. Since $\Delta(\Lambda)=1$, this is (2).

(3) Let $\mathfrak f(\tau)=e^{-\pi i/24}\eta(\frac{\tau+1}2)/\eta(\tau)$,
$\mathfrak f_1(\tau)=\eta(\frac\tau2)/\eta(\tau)$ and $\mathfrak f_2(\tau)=\sqrt2\,\eta(2\tau)/\eta(\tau)$
be Weber's functions. Then $\theta_{(0,0)}(0,\tau)=\eta(\tau)\mathfrak f(\tau)^2$,
$\theta_{(0,1)}(0,\tau)=\eta(\tau)\mathfrak f_1(\tau)^2$ and
$\theta_{(1,0)}(0,\tau)=\eta(\tau)\mathfrak f_2(\tau)^2$ \cite{Web08}, so
$(\theta_c(0,\tau)/\eta(\tau))^{12}$ is one of $\mathfrak f(\tau)^{24}$, $\mathfrak f_1(\tau)^{24}$,
$\mathfrak f_2(\tau)^{24}$; and $\mathfrak f(\tau)^{24}$, $-\mathfrak f_1(\tau)^{24}$ and
$-\mathfrak f_2(\tau)^{24}$ are the roots of $(x-16)^3-j(\tau)x$ \cite[\S12]{Cox13}. Since
$j(\tau)=j(\Lambda)$, a root of this monic polynomial is an algebraic integer when $j(\Lambda)$ is, and
then so is $Y^{-3}$ by (2).
\end{proof}

\begin{lemma}\label{is-vacuum-eta-mean}
Let $\Lambda$ be a CM lattice of the host of a definite-side datum, with $d$ and $K$, let $\chi_d$ be
the Legendre symbol modulo $d$, and let $M\subset\mathbf C$ be a number field containing $\tau(E)$ for
the types $\tau$, $g_2(\Lambda)$ and $g_3(\Lambda)$ (Definition~\ref{is-vacuum-theta-vectors}). For each
embedding $\rho$ of $M$ let $t_\rho$ be an unkilled two-torsion point of $\Lambda^\rho$. Then:
\begin{enumerate}
\item for every $\rho$, $\Delta(\Lambda^\rho)=1$, and $\Lambda^\rho=\delta'\varphi_0(\mathfrak c')$ for a
  nonzero ideal $\mathfrak c'$ of $\mathcal O_K$ and some $\delta'$; in particular $\Lambda^\rho$ has
  complex multiplication by $\mathcal O_K$;
\item $\displaystyle\frac1{[M:\mathbf Q]}\sum_\rho\log A(\Lambda^\rho)=\log\pi+\gamma+\log\frac{2\pi}{\sqrt d}
  -\frac{L'}{L}(1,\chi_d)$;
\item $\displaystyle\frac1{[M:\mathbf Q]}\sum_\rho\log n(\Lambda^\rho,t_\rho)\le\frac12\Big(\log\sqrt d-\gamma
  +\frac{L'}{L}(1,\chi_d)\Big)$.
\end{enumerate}
\end{lemma}
\leanref{Development/AbelianVarietyTheory/Proofs/DecoratedLatticeProofs/ConjugateIdealLatticeDiscrOneHomotheticIdealStatement.lean}{GlobalTheory.T.conjugateIdealLattice_discr_one_homothetic_ideal}
\leanref{Development/AbelianVarietyTheory/Proofs/DecoratedLatticeProofs/MeanLogCovolumeConjugateLatticesEqStatement.lean}{GlobalTheory.T.mean_log_covolume_conjugateLattices_eq}
\leanref{Development/AbelianVarietyTheory/Proofs/GlobalInterpolation/HalfPeriodWeierstrassValueMemConjStatement.lean}{GlobalTheory.T.halfPeriod_weierstrassValue_mem_conj}
\leanref{Development/AbelianVarietyTheory/Proofs/DecoratedLatticeProofs/ExistsIntegralThetaEtaTwelfthPowerConjugatesStatement.lean}{GlobalTheory.T.exists_integral_thetaEta_twelfth_power_conjugates}
\leanref{Development/AbelianVarietyTheory/Proofs/DecoratedLatticeProofs/ThetaZeroMeanLogNormConjugatesOfWeber.lean}{GlobalTheory.Theta0MeanCore.core}
\begin{proof}
(1) $\Delta$ is a polynomial in $g_2,g_3$ with rational coefficients, so
$\Delta(\Lambda^\rho)=\rho(\Delta(\Lambda))=1$. The invariant $j(\Lambda^\rho)=\rho(j(\Lambda))$ is a
conjugate of $j(\varphi_0(\mathfrak c))$, hence equal to $j(\varphi_0(\mathfrak c'))$ for a nonzero ideal
$\mathfrak c'$ of $\mathcal O_K$ by the main theorem of complex multiplication
\cite[Chapter~II]{Sil94}, and lattices with the same $j$-invariant are homothetic.

(2) Write $\Lambda^\rho=\omega_\rho(\mathbf Z+\mathbf Z\tau_\rho)$ with $\operatorname{Im}\tau_\rho>0$. By (1)
and Lemma~\ref{is-vacuum-eta-weber}(1), $\log A(\Lambda^\rho)=\log\pi+
\log\big(4\pi\operatorname{Im}\tau_\rho|\eta(\tau_\rho)|^4\big)$, and the second term depends only on
the homothety class of $\Lambda^\rho$. By the main theorem of complex multiplication
\cite[Chapter~II]{Sil94}, $j(\varphi_0(\mathfrak c))\in M$ has degree $h_K$ over $\mathbf Q$ and its
conjugates are the $j(\varphi_0(\mathfrak a_C))$ for representatives $\mathfrak a_C$ of the $h_K$ ideal
classes $C$; so for each class $C$, $\Lambda^\rho$ is homothetic to $\varphi_0(\mathfrak a_C)$ for
$[M:\mathbf Q]/h_K$ embeddings $\rho$. The mean of the second term is therefore its mean over the ideal
classes, which is $\gamma+\log\frac{2\pi}{\sqrt d}-\frac{L'}{L}(1,\chi_d)$ by the Chowla--Selberg formula,
Theorem~\ref{cv-chowla-selberg}(1).

(3) Let $\iota_M$ be the inclusion of $M$ in $\mathbf C$, so that $\Lambda^{\iota_M}=\Lambda$, let
$e=\wp_\Lambda(t_{\iota_M})$ and $Y=\frac14(12e^2-g_2(\Lambda))$. The value of $\wp_{\Lambda^\rho}$ at an
unkilled two-torsion point is the root of $4x^3-g_2x-g_3$ at the point of order $2$ of the curve
$y^2=4x^3-g_2(\Lambda^\rho)x-g_3(\Lambda^\rho)$ that is killed by neither of the endomorphisms
$\omega_K,\bar\omega_K$. An automorphism of $\mathbf C$ that extends $\rho$ carries the curve of
$\Lambda$ to that of $\Lambda^\rho$ and the endomorphisms $\omega_K,\bar\omega_K$ to the same two
endomorphisms, in some order, so it carries $e$ to $\wp_{\Lambda^\rho}(t_\rho)$. For $\rho=\iota_M$ this
shows that $e$ is fixed by every automorphism of $\mathbf C$ over $M$, so $e\in M$ and $Y\in M$, and
then $\rho(Y)=\frac14(12\wp_{\Lambda^\rho}(t_\rho)^2-g_2(\Lambda^\rho))$ for every $\rho$.

Now write $\Lambda^\rho=\omega_\rho(\mathbf Z+\mathbf Z\tau_\rho)$ as in (2), and let $c_\rho$ be the
characteristic of Lemma~\ref{is-vacuum-genus-one}(3) for $t_\rho$. By
Lemma~\ref{is-vacuum-genus-one}(3) and Lemma~\ref{is-vacuum-eta-weber}(1),
\[
 \log n(\Lambda^\rho,t_\rho)+\tfrac12\log A(\Lambda^\rho)=\tfrac12\log2+\log\pi
 -2\log\big|\theta_{c_\rho}(0,\tau_\rho)/\eta(\tau_\rho)\big|.
\]
By Lemma~\ref{is-vacuum-eta-weber}(2), $X=Y^{-3}\in M$ satisfies
$|\rho(X)|=|\theta_{c_\rho}(0,\tau_\rho)/\eta(\tau_\rho)|^{12}$ for every $\rho$. Since $\Lambda$ has
complex multiplication by $\mathcal O_K$, $j(\Lambda)$ is an algebraic integer
(Theorem~\ref{pd-j-invariant-integral}), so $X$ is a nonzero algebraic integer by
Lemma~\ref{is-vacuum-eta-weber}(3), and $\sum_\rho\log|\rho(X)|=\log|N_{M/\mathbf Q}(X)|\geq0$. Hence
$\sum_\rho\log|\theta_{c_\rho}(0,\tau_\rho)/\eta(\tau_\rho)|\geq0$, and summing the display over $\rho$,
\[
 \sum_\rho\log n(\Lambda^\rho,t_\rho)\le[M:\mathbf Q]\big(\tfrac12\log2+\log\pi\big)
 -\tfrac12\sum_\rho\log A(\Lambda^\rho).
\]
By (2) the right side is $[M:\mathbf Q]\cdot\frac12\big(\log\sqrt d-\gamma+\frac{L'}{L}(1,\chi_d)\big)$.
\end{proof}

\begin{lemma}\label{is-vacuum-units-intertwiners}
Let $W$ be a finite-dimensional complex vector space with a Hermitian form $H$, and let $c\in\mathbf R$.
Form the operators $U_x$ for $x\in W$ and the function $\alpha_c$ with $H$ as in
Definition~\ref{pd-theta-function}. For a function $\alpha$ on $W$ and subgroups $\Gamma,\Gamma'\subset W$
with $\Gamma'/(\Gamma\cap\Gamma')$ finite, let $I^\alpha_{\Gamma\to\Gamma'}$ be the averaging operator of
Definition~\ref{is-vacuum-theta-vectors}(2) formed with $H$ and with $\alpha$ in place of
$\alpha_{\sqrt d/2}$, and put $I_{\Gamma\to\Gamma'}=I^{\alpha_c}_{\Gamma\to\Gamma'}$. Call $\alpha$ a
\emph{semicharacter} on a subgroup $\Gamma$ if $\alpha(x)=\pm1$ and
$\alpha(x+y)=\alpha(x)\alpha(y)e^{\pi i\operatorname{Im}H(x,y)}$ for all $x,y\in\Gamma$.
\begin{enumerate}
\item For $x,y\in W$, $U_xU_y=e^{-\pi i\operatorname{Im}H(x,y)}U_{x+y}$.
\item Let $d\geq1$ and $b_0$ be integers with $b_0$ odd, put $\omega=(b_0+i\sqrt d)/2$ and $c=\sqrt d/2$,
  and let $\Gamma\subset W$ be a subgroup with $\omega\Gamma\subset\Gamma$ on which
  $\operatorname{Im}H$ takes integer values. Then for $x,y\in\Gamma$ the number
  $q(x)=\frac{\sqrt d}2\operatorname{Re}H(x,x)$ is an integer and
  $q(x+y)-q(x)-q(y)-\operatorname{Im}H(x,y)$ is an even integer, and $\alpha_c$ is a semicharacter on
  $\Gamma$. Moreover, if $\alpha$ is a semicharacter on a subgroup $\Gamma_2\subset W$, then
  $\alpha(x)U_x\big(\alpha(y)U_yF\big)=\alpha(x+y)U_{x+y}F$ for all $x,y\in\Gamma_2$ and every function $F$
  on $W$. In particular, if $\Gamma_1\subset W$ is a subgroup with $\Gamma_2/(\Gamma_1\cap\Gamma_2)$
  finite and $\alpha(y)U_yF=F$ for $y\in\Gamma_1$, the sum that defines $I^\alpha_{\Gamma_1\to\Gamma_2}F$
  does not depend on the representatives.
\item Let $\Gamma_1,\Gamma_2\subset W$ be subgroups with $\Gamma_2/(\Gamma_1\cap\Gamma_2)$ finite, $\alpha$
  a semicharacter on $\Gamma_2$, and $F$ a function on $W$ with $\alpha(y)U_yF=F$ for $y\in\Gamma_1$. Let
  $W'$ be a complex vector space with a Hermitian form $H'$, and $\Phi\colon W'\to W$ a complex-linear
  isomorphism with $H(\Phi x,\Phi y)=H'(x,y)$ for all $x,y$. Then
  $I^{\alpha\circ\Phi}_{\Phi^{-1}\Gamma_1\to\Phi^{-1}\Gamma_2}(F\circ\Phi)=(I^\alpha_{\Gamma_1\to\Gamma_2}F)\circ\Phi$,
  where the left side is formed with $H'$.
\item Let $L\subset\mathbf C$ be a lattice of covolume $A$, let $t\in\mathbf C$, and let $\vartheta$ be
  $\vartheta_{L,t}$ (Definition~\ref{pd-theta-function}) if $(U^L_t\sigma^*_L)(0)\neq0$ and $0$ otherwise.
  For a finite set $S$, give $\mathbf C^S$ the form $H_S(z,w)=A^{-1}\sum_{s\in S}z_s\overline{w_s}$, let
  $L^S$ be the lattice of the vectors with all coordinates in $L$, put
  $\vartheta^S(y)=\prod_{s\in S}\vartheta(y_s)$, and form $U_x$, $\alpha_c$ and $I$ with $H_S$. Let $m$
  and $J$ be finite sets, write $y_k=(y_{k,p})_{p\in m}$ for $y\in\mathbf C^{J\times m}$, and for a
  subgroup $\Gamma\subset\mathbf C^m$ let $\Gamma^J$ be the set of $y\in\mathbf C^{J\times m}$ with
  $y_k\in\Gamma$ for every $k$. Let $\Gamma_2\subset\mathbf C^m$ be a subgroup on which $\alpha_c$ is a
  semicharacter, with $\Gamma_2/(L^m\cap\Gamma_2)$ finite, and suppose that
  $\alpha_c(y)U_y\vartheta^m=\vartheta^m$ for every $y\in L^m$. Then $\alpha_c$ is a semicharacter on
  $\Gamma_2^J$, $\alpha_c(y)U_y\vartheta^{J\times m}=\vartheta^{J\times m}$ for every $y\in L^{J\times m}$,
  $\Gamma_2^J/(L^{J\times m}\cap\Gamma_2^J)$ is finite, and
  \[
   \big(I_{L^{J\times m}\to\Gamma_2^J}\vartheta^{J\times m}\big)(y)
   =\prod_{k\in J}\big(I_{L^m\to\Gamma_2}\vartheta^m\big)(y_k).
  \]
\end{enumerate}
\end{lemma}
\leanref{Formalism/LocalToGlobal/TorsionIntegrality/QuadraticSemicharacters/CayleyGaussianIntegrals/MatrixFormFourfoldBlock.lean}{Formalism.L2.Th1.H.heisenbergMul}
\leanref{Formalism/LocalToGlobal/TorsionIntegrality/QuadraticSemicharacters/CayleyGaussianIntegrals/MatrixFormFourfoldBlockFirstBlockEmb.lean}{Formalism.L2.Th1.H.quadMultiplierSemicharacter}
\leanref{Formalism/LocalToGlobal/TorsionIntegrality/QuadraticSemicharacters/CayleyGaussianIntegrals/MatrixFormFourfoldBlock.lean}{Formalism.L2.Th1.H.semicharacterHom}
\leanref{Formalism/LocalToGlobal/TorsionIntegrality/QuadraticSemicharacters/QuadraticSemicharacterIdentityReformulated2.lean}{Formalism.L2.Th1.H.interReps}
\leanref{Development/AbelianVarietyTheory/Proofs/DivisorProofs/LatticeIntertwinerComapLinearIsometryStatement.lean}{GlobalTheory.T.latticeIntertwiner_comap_linear_isometry}
\leanref{Development/AbelianVarietyTheory/Proofs/DivisorProofs/LatticeIntertwinerProdOrthogonalBlocksStatement.lean}{GlobalTheory.T.latticeIntertwiner_prod_orthogonal_blocks}
\begin{proof}
(1) Both sides send $F$ to $z\mapsto F(z+x+y)$ times an exponential. Expanding $H(z+x,y)$ and
$H(x+y,x+y)$, the exponents differ by
$-\pi H(x,y)+\frac\pi2\big(H(x,y)+H(y,x)\big)=-\pi i\operatorname{Im}H(x,y)$.

(2) Write $E(x,y)=\operatorname{Im}H(x,y)$, so that $\alpha_c(x)=e^{\pi iq(x)}$. Let $x,y\in\Gamma$.
Since $H(\omega x,y)=\omega H(x,y)$, $E(\omega x,y)=\frac{b_0}2E(x,y)+\frac{\sqrt d}2\operatorname{Re}H(x,y)$,
and $E(x,x)=0$. So $q(x)=E(\omega x,x)$, an integer because $\omega x\in\Gamma$, and
\[
 q(x+y)-q(x)-q(y)-E(x,y)=\sqrt d\operatorname{Re}H(x,y)-E(x,y)=2E(\omega x,y)-(b_0+1)E(x,y),
\]
an even integer because $b_0$ is odd. Hence $\alpha_c(x)=\pm1$ and
$\alpha_c(x+y)=\alpha_c(x)\alpha_c(y)e^{\pi iE(x,y)}$. Now let $\alpha$ be a semicharacter on $\Gamma_2$ and
$x,y\in\Gamma_2$. Comparing $\alpha(x+y)$ with $\alpha(y+x)$ gives $e^{2\pi iE(x,y)}=1$, so $E(x,y)$ is an
integer and $e^{-\pi iE(x,y)}=e^{\pi iE(x,y)}$. By (1),
$\alpha(x)U_x\big(\alpha(y)U_yF\big)=\alpha(x)\alpha(y)e^{-\pi iE(x,y)}U_{x+y}F=\alpha(x+y)U_{x+y}F$. If
$\alpha(y)U_yF=F$ for $y\in\Gamma_1$, this gives $\alpha(x+y)U_{x+y}F=\alpha(x)U_xF$ for $x\in\Gamma_2$ and
$y\in\Gamma_1\cap\Gamma_2$, so the term of the sum at a representative depends only on its class.

(3) The function $\alpha\circ\Phi$ is a semicharacter on $\Phi^{-1}\Gamma_2$ for $H'$, because
$\operatorname{Im}H'(x,y)=\operatorname{Im}H(\Phi x,\Phi y)$. For the operators $U'$ of $H'$,
$U'_{\Phi^{-1}x}(F\circ\Phi)=(U_xF)\circ\Phi$, because $H'(z,\Phi^{-1}x)=H(\Phi z,x)$; so $F\circ\Phi$ is
fixed by $\alpha(\Phi y)U'_y$ for $y\in\Phi^{-1}\Gamma_1$. If $x_1,\dots,x_n$ represent
$\Gamma_2/(\Gamma_1\cap\Gamma_2)$, then the $\Phi^{-1}x_i$ represent
$\Phi^{-1}\Gamma_2/(\Phi^{-1}\Gamma_1\cap\Phi^{-1}\Gamma_2)$, and by (2) both averaging operators may be
computed with these representatives. The two sums then agree term by term.

(4) For $y,y'\in\mathbf C^{J\times m}$, $H_{J\times m}(y,y')=\sum_kH_m(y_k,y'_k)$. Hence
$\alpha_c(y)=\prod_k\alpha_c(y_k)$ and $\operatorname{Im}H_{J\times m}(y,y')=\sum_k\operatorname{Im}H_m(y_k,y'_k)$,
so the two identities of a semicharacter on $\Gamma_2$ give them on $\Gamma_2^J$. Moreover
$\vartheta^{J\times m}(z)=\prod_k\vartheta^m(z_k)$ and $(U_y\vartheta^{J\times m})(z)=\prod_k(U_{y_k}\vartheta^m)(z_k)$,
and $y\in L^{J\times m}$ if and only if $y_k\in L^m$ for every $k$; so $\vartheta^{J\times m}$ is fixed by
$\alpha_c(y)U_y$ for $y\in L^{J\times m}$. If $x_1,\dots,x_n$ represent $\Gamma_2/(L^m\cap\Gamma_2)$, the
$y$ whose blocks $y_k$ are among the $x_i$ represent $\Gamma_2^J/(L^{J\times m}\cap\Gamma_2^J)$, which
therefore has $n^{|J|}$ elements. By (2) both averaging operators may be computed with these
representatives, and expanding the product over $k$ of the sums that define
$(I_{L^m\to\Gamma_2}\vartheta^m)(y_k)$, with the constants $n^{-1/2}$, gives the sum that defines the
left side.
\end{proof}

\begin{lemma}\label{is-vacuum-units-host}
Let $D$ be a definite-side datum with kernel data, and let $\Lambda$, $\omega_0/2$, $\gamma_0$,
$\kappa_0$ and the vacuum vector $\theta_0$ of the host be those of
Definition~\ref{is-vacuum-theta-vectors}. Let $a=2n$ and $b=2n'$ be even. For $k=(i,j)$ with
$1\le i\le n'$ and $1\le j\le n$, the \emph{block} $k$ of a $b\times a$ matrix is its $2\times2$
submatrix in the rows $2i-1,2i$ and the columns $2j-1,2j$, read as a vector of length $4$ in a fixed
order of its entries. Let $e_w$, $w\in\{1,2,3,4\}\times\mathcal T$, be the standard basis of
$K^{4\times\mathcal T}$, and let $v_{k,w}\in M_{b\times a}(E)$ be the matrix whose block $k$ is
$\gamma_0^{-1}(e_w)$ and whose other entries are $0$. Then:
\begin{enumerate}
\item the $v_{k,w}$ form an orthonormal basis $v$ of $M_{b\times a}(E)$;
\item $\kappa_0$ is algebraic, and all its conjugates have absolute value $1$;
\item $\theta_0^{\otimes ab}=\kappa_0^{nn'}\,I_{\Lambda_v\to\Lambda_{\mathcal O_E^a,\mathcal O_E^b}}
  \big((\vartheta_{\Lambda,\omega_0/2})_v\big)$.
\end{enumerate}
In particular, $\theta_0^{\otimes ab}$ is a theta vector up to units for
$(\omega_0/2,\mathcal O_E^a,\mathcal O_E^b)$.
\end{lemma}
\leanref{Development/AbelianVarietyTheory/Proofs/DivisorProofs/BaseThetaVectorIsVacuumEvenRanksStatement.lean}{GlobalTheory.T.base_theta_vector_is_vacuum_even_ranks}
\leanref{Development/AbelianVarietyTheory/Proofs/DivisorProofs/OrthonormalBasisBlockVacuumLatticeBasisFamilyStatement.lean}{GlobalTheory.T.orthonormalBasis_blockVacuumLatticeBasisFamily}
\leanref{Development/AbelianVarietyTheory/Proofs/DivisorProofs/KappaProductPowAlgebraicModulusOneAllStatement.lean}{GlobalTheory.T.kappaProduct_pow_algebraicModulusOneAll}
\leanref{Development/AbelianVarietyTheory/Proofs/DivisorProofs/ThetaVacuumEqKappaPowMulBlockKernelsStatement.lean}{GlobalTheory.T.thetaVacuum_eq_kappa_pow_mul_block_kernels}
\begin{proof}
Write $\vartheta=\vartheta_{\Lambda,\omega_0/2}$, $\alpha=\alpha_{\sqrt d/2}$, $A$ for the covolume of
$\Lambda$, and $H_{\mathrm{std}}$ for the form $A^{-1}\sum_uy_u\overline{y'_u}$ on $\mathbf C^S$ for a finite
set $S$. For $x\in E^4$ put $\iota(x)=(\tau(x_c))_{c,\tau}\in\mathbf C^{4\times\mathcal T}$, and for
$z\in V_{a,b}$ let $z^{(k)}\in\mathbf C^{4\times\mathcal T}=(\mathbf C^{\mathcal T})^4$ be its block $k$,
the four entries $(z_{\tau,\beta,\alpha})_\tau$ with $(\beta,\alpha)$ in block $k$ in the fixed order. Let
$J$ be the set of blocks, which has $nn'$ elements.

(1) For $k\ne k'$ the matrices $v_{k,w}$ and $v_{k',w'}$ have disjoint supports, so
$\operatorname{Tr}_{E/K}\sum_{\beta,\alpha}(v_{k,w})_{\beta\alpha}\overline{(v_{k',w'})_{\beta\alpha}}=0$. For
$k=k'$ this number is $\operatorname{Tr}_{E/K}\sum_cx_c\bar y_c$ with $x=\gamma_0^{-1}(e_w)$ and
$y=\gamma_0^{-1}(e_{w'})$, which is the standard pairing of $e_w$ and $e_{w'}$, that is $\delta_{ww'}$,
because $\gamma_0$ is an isometry. So the $4fnn'=fab$ vectors $v_{k,w}$ are orthonormal, hence linearly
independent over $K$; as $\dim_KM_{b\times a}(E)=fab$, they form a basis.

(2) The kernel data construct $\kappa_0$ as $\prod_pr_p^{-y_p}$, the product over the primes $p$ that
are $2$ or split in $K$, where the $y_p$ are integers of which only finitely many are nonzero and
$r_p^{2h_K}=\varphi_0(\pi_p)/\varphi_0(\bar\pi_p)$ for the class number $h_K$ of $K$ and a generator
$\pi_p\in\mathcal O_K$ of $\mathfrak p^{h_K}$, $\mathfrak p$ a prime of $K$ over $p$. Each $r_p$ is
algebraic. Let $\sigma$ be an embedding of $\overline{\mathbf Q}$ into $\mathbf C$. Then
$\sigma\circ\varphi_0$ is an embedding of the CM field $E$, so
$\sigma\varphi_0(\bar\pi_p)=\overline{\sigma\varphi_0(\pi_p)}$ and $|\sigma(r_p)|^{2h_K}=1$. Hence
$|\sigma(\kappa_0)|=1$.

(3) \emph{Step 1: coordinates.} The vectors $\gamma_0^{-1}(e_w)$ form an orthonormal basis of
$E^4=M_{1\times4}(E)$, because $\gamma_0$ is an isometry. By Lemma~\ref{is-vacuum-tensor-norm}(1) their
images $\iota(\gamma_0^{-1}(e_w))$ form an orthonormal basis of $\mathbf C^{4\times\mathcal T}$ for
$\sum_uy_u\overline{y'_u}$, and the map $g$ of Definition~\ref{is-vacuum-theta-vectors} carries them to
$\varphi_0(e_w)=e_w$. So $g$ is unitary, and $H_{\mathrm{std}}(gy,gy')=H^{\oplus4}(y,y')$. Let
$\phi\colon V_{a,b}\to\mathbf C^{J\times4\times\mathcal T}$ be the map with $\phi(z)_k=g(z^{(k)})$.
Since the blocks partition the entries, $\phi$ is unitary, and $H_{\mathrm{std}}(\phi z,\phi z')=H(z,z')$.
Block $k$ of $\iota(v_{k,w})$ is $\iota(\gamma_0^{-1}(e_w))$ and its other blocks are $0$, so
$\phi(\iota(v_{k,w}))$ is the standard basis vector $e_{k,w}$. Hence:
\begin{enumerate}
\item[(a)] $\phi(z)_{k,w}=\langle\phi z,\phi\,\iota(v_{k,w})\rangle=\langle z,\iota(v_{k,w})\rangle$, so
  that $(\vartheta)_v=\vartheta^{\otimes}\circ\phi$, where $\vartheta^{\otimes}(y)=\prod_{k,w}\vartheta(y_{k,w})$;
\item[(b)] $\phi(\Lambda_v)$ is the lattice $\Lambda^{J\times4\times\mathcal T}$ of vectors with all
  coordinates in $\Lambda$.
\end{enumerate}

\emph{Step 2: the lattices.} Let $\Lambda'=g(\Lambda_{\mathcal O_E,\mathcal O_E}^4)\subset
\mathbf C^{4\times\mathcal T}$ and $F_0(y)=\prod_{c,\tau}\vartheta(y_{c,\tau})$, as in identity (4) of
Definition~\ref{pd-cm-datum} read for $\theta_0$. The subgroup $\Lambda_{\mathcal O_E^a,\mathcal O_E^b}$ is
the set of $z\in V_{a,b}$ whose entries $(z_{\tau,\beta,\alpha})_\tau$ all lie in
$\Lambda_{\mathcal O_E,\mathcal O_E}$: the entries of a generator $w\,\iota(sB)$ are the
$w\,\iota(sB_{\beta\alpha})$ with $sB_{\beta\alpha}\in\mathfrak A_0$, and conversely $w\,\iota(s)$ placed at the
entry $(\beta,\alpha)$ is $w\,\iota(sB)$ for the matrix unit $B$ at $(\beta,\alpha)$. So
$\phi(\Lambda_{\mathcal O_E^a,\mathcal O_E^b})=\Lambda'^{J}$ in the notation of
Lemma~\ref{is-vacuum-units-intertwiners}(4). In the same way $\Lambda_{\mathcal O_E,\mathcal O_E}^4$ is the
subgroup $\Lambda_{\mathcal O_E^2,\mathcal O_E^2}$ of $V_{2,2}$, read in the fixed order, which is standard
(Lemma~\ref{is-vacuum-averaging}); so $\operatorname{Im}H$ takes integer values on it, and
$\operatorname{Im}H_{\mathrm{std}}$ takes integer values on $\Lambda'$. Also
$\omega_K\Lambda'\subset\Lambda'$, because $\omega_K\Lambda\subset\Lambda$ and $g$ is complex-linear. For
$y\in\Lambda^{4\times\mathcal T}$, $\alpha(y)U_yF_0=F_0$ by Lemma~\ref{is-vacuum-genus-one}(2), as in the
proof of Lemma~\ref{is-vacuum-tensor-norm}(3). Finally, $\mathfrak A_0^4$ is finitely generated and
$\gamma_0$ is $K$-linear, so there is an integer $N\geq1$ with
$N\gamma_0(\mathfrak A_0^4)\subset\mathcal O_K^{4\times\mathcal T}$. The group
$\Lambda_{\mathcal O_E,\mathcal O_E}^4$ is generated by the $w\,\iota(x)$ with $w\in\Lambda$ and
$x\in\mathfrak A_0^4$, and $Ng(w\,\iota(x))=w\varphi_0(N\gamma_0(x))$ lies in $\Lambda^{4\times\mathcal T}$, as
$\varphi_0(\mathcal O_K)\Lambda\subset\Lambda$. So $N\Lambda'\subset\Lambda^{4\times\mathcal T}$, and
$\Lambda'/(\Lambda'\cap\Lambda^{4\times\mathcal T})$ is finite. By
Lemma~\ref{is-vacuum-units-intertwiners}(2), with $b_0=1$ and $\omega=\omega_K$, $\alpha$ formed with
$H_{\mathrm{std}}$ is a semicharacter on $\Lambda'$. So the hypotheses of
Lemma~\ref{is-vacuum-units-intertwiners}(4) hold for the lattice $\Lambda$, the point $t=\omega_0/2$, for
which $\vartheta^{4\times\mathcal T}=F_0$, the set $m=4\times\mathcal T$ and $\Gamma_2=\Lambda'$.

\emph{Step 3: the identity.} By Definition~\ref{is-vacuum-theta-vectors}, identity (4) of
Definition~\ref{pd-cm-datum} for $\theta_0$ reads
$\prod_{c=1}^4\theta_0(y_c)=\kappa_0\big(I_{\Lambda^{4\times\mathcal T}\to\Lambda'}F_0\big)(gy)$ for
$y=(y_1,\dots,y_4)\in(\mathbf C^{\mathcal T})^4$, with the averaging operator formed with
$H_{\mathrm{std}}$. Grouping the factors of $\theta_0^{\otimes ab}(z)$ by blocks, and using
Lemma~\ref{is-vacuum-units-intertwiners}(4) with $m=4\times\mathcal T$ and the set $J$ of blocks, for which
$\vartheta^{J\times m}=\vartheta^{\otimes}$,
\[
 \theta_0^{\otimes ab}(z)=\prod_{k\in J}\kappa_0\big(I_{\Lambda^{4\times\mathcal T}\to\Lambda'}F_0\big)(\phi(z)_k)
 =\kappa_0^{nn'}\big(I_{\Lambda^{J\times4\times\mathcal T}\to\Lambda'^{J}}\vartheta^{\otimes}\big)(\phi z).
\]
By Lemma~\ref{is-vacuum-units-intertwiners}(3) with $\Phi=\phi$, whose hypotheses on $\Lambda'^J$ and
$\vartheta^{\otimes}$ are conclusions of (4), and by Steps 1 and 2, the right side
is $\kappa_0^{nn'}\big(I_{\Lambda_v\to\Lambda_{\mathcal O_E^a,\mathcal O_E^b}}(\vartheta)_v\big)(z)$, the
averaging operator formed with $H$ and $\alpha$. This is (3). With (1) and (2), and since all
conjugates of $\kappa_0^{nn'}$ have absolute value $1$, it shows that $\theta_0^{\otimes ab}$ is a theta
vector up to units for $(\omega_0/2,\mathcal O_E^a,\mathcal O_E^b)$.
% The Lean proof is GlobalTheory.BaseVacuum.thetaVectorUpToUnits_thZero_of_even: the basis bas B n n'
% (blocks by posE and four, entries xv = gamma_0^{-1}(e_w)), the unit B.kappa^(n'n), Step 3 from
% thetaVacuum_eq_kappa_pow_mul_block_kernels, latticeIntertwiner_prod_orthogonal_blocks and
% latticeIntertwiner_comap_linear_isometry along phi B n n' (blockVacuumCoordinateEquiv_isometry),
% with blockVacuumLattice_eq_comap_standardLattice, la_tor0_stdL_eq_blkLat and
% pullback_basis_theta_eq_F0_form; Step 2 is Base.StdOK (IntegralLatticeDensity.SO.stdOK).
\end{proof}

\begin{lemma}\label{is-vacuum-norm-bound}
There is an absolute constant $C_t$ with the following property. Let $D$ be a definite-side datum with
kernel data as in Proposition~\ref{is-kernel-entry-bound}, and $M$ a number field containing the
coefficients of the theta data. Let $r_E$ be the root discriminant of the host field and $X\geq\log(r_E/2)+|\operatorname{Re}
\frac{L'}{L}(1,\chi_d)|$. Then for every admitted pair of lattices of ranks $a$ and $b$, the mean over
the embeddings $\rho$ of $M$ of the logarithm of the squared norm of the conjugate under $\rho$ of the
vacuum vector of the pair (Definition~\ref{is-vacuum-theta-vectors}) is at most $C_tfab\,X$; and this
norm does not depend on the admitted pair.
\end{lemma}
\leanref{Development/AbelianVarietyTheory/Proofs/DivisorProofs/NormalisedFamilyNormBoundAbsoluteConstantStatement.lean}{GlobalTheory.T.normalisedFamily_normBound_absoluteConstant}
\leanref{Development/AbelianVarietyTheory/Proofs/DivisorProofs/RelDetIdealSelfDualNormOneStatement.lean}{GlobalTheory.T.relDetIdeal_selfDual_normOne}
\leanref{Development/AbelianVarietyTheory/Proofs/DivisorProofs/NormOneCharacterModulusOneAllEmbeddingsStatement.lean}{GlobalTheory.T.normOneCharacter_modulusOne_all_embeddings}
\begin{proof}
Let $\Lambda$ be the CM lattice of the host given by the kernel data, $\theta_0$ the vacuum vector of
the host, and $\theta_{L,L'}$ the vacuum vector of an admitted pair $(L,L')$ of ranks $a,b$
(Definition~\ref{is-vacuum-theta-vectors}). Its squared norm is $\|\theta_{L,L'}\|^2_{\Lambda_{L,L'}}$,
and that of its conjugate under $\rho$ is $\|\theta^\rho_{L,L'}\|^2_{\Lambda^\rho_{L,L'}}$. For an
embedding $\rho$ of $M$, $\Lambda^\rho$ has complex multiplication by $\mathcal O_K$
(Lemma~\ref{is-vacuum-eta-mean}(1)), so it has an unkilled two-torsion point $t_\rho$, unique modulo
$\Lambda^\rho$, and $n(\rho)=n(\Lambda^\rho,t_\rho)$ does not depend on its choice
(Lemma~\ref{is-vacuum-genus-one}(1)).

\emph{Step 1: the tensor power.} Let $t$ be an unkilled two-torsion point of $\Lambda$, and call a
subgroup of $V_{a,b}$ standard as in Lemma~\ref{is-vacuum-averaging}. By
Lemma~\ref{is-vacuum-genus-one}(2), $\vartheta_{\Lambda,t}$ satisfies the hypotheses of
Lemma~\ref{is-vacuum-tensor-norm}(2) and (3). So for every orthonormal basis $v$, the tensor power
$(\vartheta_{\Lambda,t})_v$ is a theta function for $\Lambda_v$ with
$\|(\vartheta_{\Lambda,t})_v\|^2_{\Lambda_v}=n(\Lambda,t)^{fab}$.

(1a) \emph{$\theta_{L,L'}$ is a theta vector up to units for $(t,L,L')$.} Since $a$ and $b$ are even,
the tensor power $\theta_0^{\otimes ab}$ is a theta vector up to units for
$(\omega_0/2,\mathcal O_E^a,\mathcal O_E^b)$ by Lemma~\ref{is-vacuum-units-host}. Since
$\vartheta_{\Lambda,t}=\vartheta_{\Lambda,\omega_0/2}$ by Lemma~\ref{is-vacuum-genus-one}(1), it is also
one for $(t,\mathcal O_E^a,\mathcal O_E^b)$. If $G=\kappa I_{\Lambda_v\to\Gamma}(\vartheta_{\Lambda,t})_v$
with $\Gamma$ standard and $\Gamma'$ is standard, then
$I_{\Gamma\to\Gamma'}G=\kappa\mu\,I_{\Lambda_v\to\Gamma'}(\vartheta_{\Lambda,t})_v$ with $\mu^8=1$ by
Lemma~\ref{is-vacuum-averaging}(4), and all conjugates of $\kappa\mu$ have absolute value $1$. Applying
this with $\Gamma'=\Lambda_{L,\mathcal O_E^b}$ and then with $\Gamma'=\Lambda_{L,L'}$ shows that
$I_{\Lambda_{L,\mathcal O_E^b}\to\Lambda_{L,L'}}I_{\Lambda_{\mathcal O_E^a,\mathcal O_E^b}\to\Lambda_{L,\mathcal
O_E^b}}\theta_0^{\otimes ab}$ is a theta vector up to units for $(t,L,L')$.

Next, $\mathfrak d_L\bar{\mathfrak d}_L=\mathcal O_E$. For each prime $\mathfrak P$ of $E$ choose an
$E$-linear automorphism $f_{\mathfrak P}$ of $E^a$ that carries the localization of $\mathcal O_E^a$ at
$\mathfrak P$ onto that of $L$; then $\mathfrak d_L$ is generated at $\mathfrak P$ by
$\det f_{\mathfrak P}$. The form $\sum_i\bar x_iy_i$ pairs the localization of $L$ at $\mathfrak P$ with
that at $\bar{\mathfrak P}$, and $L$ is self-dual, so the localization of $L$ at $\bar{\mathfrak P}$ is
the image of that of $\mathcal O_E^a$ under $(f_{\mathfrak P}^*)^{-1}$, where $f^*$ is the conjugate
transpose. Hence $f_{\bar{\mathfrak P}}$ and $(f_{\mathfrak P}^*)^{-1}$ differ by an automorphism of the
standard lattice at $\bar{\mathfrak P}$, and $\overline{\det f_{\mathfrak P}}\det f_{\bar{\mathfrak P}}$ is
a unit at $\bar{\mathfrak P}$. So $\mathfrak d_L\bar{\mathfrak d}_L$ is generated by a unit at every
prime, and it is $\mathcal O_E$.

Let $h$ be the class number of $E$ and $\mathfrak d_L^h=y\mathcal O_E$. Then
$\bar{\mathfrak d}_L^h=\bar y\mathcal O_E$, so $y\bar y\mathcal O_E=\mathcal O_E$, and $z=y/\bar y$
satisfies $z\bar z=1$ and $z\mathcal O_E=y^2\mathcal O_E=\mathfrak d_L^{2h}$. The powers of
$\mathfrak d_L$ satisfy $\mathfrak a\bar{\mathfrak a}=\mathcal O_E$, so the multiplicativity of $\chi$ on
such ideals gives $\chi(\mathfrak d_L)^{2h}=\chi(z\mathcal O_E)=\prod_\tau\tau(z)^{-1}$. In particular
$\chi(\mathfrak d_L)$ is algebraic. For every embedding $\sigma$ of $E$ into $\mathbf C$,
$|\sigma(z)|^2=\sigma(z\bar z)=1$, because $E$ is a CM field; so all conjugates of
$\chi(\mathfrak d_L)^{2h}$, and hence of $\chi(\mathfrak d_L)$, have absolute value $1$, and likewise
for $\chi(\mathfrak d_{L'})$. Multiplying $\kappa$ by $\chi(\mathfrak d_L)^{b/2}\chi(\mathfrak d_{L'})^{-a/2}$
gives (1a).

(1b) \emph{For every embedding $\rho$ of $M$, $\|\theta^\rho_{L,L'}\|^2_{\Lambda^\rho_{L,L'}}=n(\rho)^{fab}$.}
By (1a), $\theta_{L,L'}$ is a theta vector up to units for $(t,L,L')$, and $ab$ is even. Extend $\rho$ to
an embedding of $\overline{\mathbf Q}$ into $\mathbf C$; the extension does not change the conjugates,
because $g_2(\Lambda)$, $g_3(\Lambda)$, the $\tau(E)$ and the Taylor coefficients of $\theta_{L,L'}$ lie in
$M$. Lemma~\ref{is-growth-theta-conjugate}(3), with $y=0$ and the unkilled two-torsion point $t_\rho$ of
$\Lambda^\rho$, gives $\|\theta^\rho_{L,L'}\|^2_{\Lambda^\rho_{L,L'}}=n(\Lambda^\rho,t_\rho)^{fab}=n(\rho)^{fab}$.
The right side does not involve the pair; this gives the last assertion.

% The fact from Lemma is-vacuum-units-host is GlobalTheory.T.base_theta_vector_is_vacuum_even_ranks.
% Step 1 is GlobalTheory.ColNorms.FamilyBound.conj_normSq and conj_indep: (1a) is
% GlobalTheory.FamilyTieProof.familyTie (genusOneVacuum_indep_twoTorsionPoint,
% GlobalTheory.MovedVacuum.isVacuum_movedThetaVacuum, chiFac_mod1All through
% relDetIdeal_selfDual_normOne and normOneCharacter_modulusOne_all_embeddings); (1b) is conj_normSq,
% which applies GlobalTheory.ConjNorm.useNorm_all (Lemma is-growth-theta-conjugate(3)) for the
% number field M directly; the point t_rho is GlobalTheory.Glob.exists_isThetaZero.

\emph{Step 2: the mean.} By Step 1 and Lemma~\ref{is-vacuum-eta-mean}(3),
\[
 \frac1{[M:\mathbf Q]}\sum_\rho\log\|\theta^\rho_{L,L'}\|^2_{\Lambda^\rho_{L,L'}}
 =\frac{fab}{[M:\mathbf Q]}\sum_\rho\log n(\rho)
 \le\frac{fab}2\Big(\log\sqrt d-\gamma+\frac{L'}{L}(1,\chi_d)\Big).
\]
The root discriminant of the host field satisfies $r_E\geq2^{e+1-2^{-e}}\sqrt d\geq2\sqrt d$, where
$e\geq2$ is the $2$-power exponent of its Galois datum, so $\log\sqrt d\le\log(r_E/2)$. Moreover
$\gamma>0$, and $\frac{L'}{L}(1,\chi_d)$ is real since $\chi_d$ is. Hence the mean is at most
$\frac{fab}2X$, and $C_t=\frac12$ has the stated property.
% Step 2 is GlobalTheory.ColNorms.FamilyBound.mean_le_of with GlobalTheory.Theta0MeanLeC.theta0MeanLeC
% (C = 1/2, from GlobalTheory.Theta0MeanCore.core and GlobalTheory.Theta0MeanProof.bound_le_half); the
% bound for r_E is GlobalTheory.Theta0MeanCompare.two_sqrt_le_rd, from the field rd_ge of
% Formalism.HC.HostDer.
\end{proof}

\begin{remark}\label{is-vacuum-determinant-dual}
The identity $\mathfrak d_L\bar{\mathfrak d}_L=\mathcal O_E$ in Step (1a) of the proof of
Lemma~\ref{is-vacuum-norm-bound} also follows without localizing: the dual of $L$ for
$\sum_i\bar x_iy_i$ has determinant ideal $\bar{\mathfrak d}_L^{-1}$, and $L$ is self-dual.
\end{remark}


\begin{proposition}\label{is-kernel-step-bound}
For every real $C_Y$ there is a constant $C\ge0$ with the following property. Let $D$ be a
definite-side datum with size data at scale $x$ and kernel data that satisfy the kernel ties with
weights $W$ (Definition~\ref{is-kernel-ties}). Suppose that $g_i-1>2fk$ for $i=0,1$,
$2fN(m)+1<4a$, the weights of every level $j\le m$ have spread at most $2a$, and $4a<T_{u'}$ for
every word $u'$. Let $u$ be a word for which the size conditions with constant $C_Y$ hold, and let
$s_0,\dots,s_{m-2}$ be the maps of the kernel chain of $u$ at the base tuple. Then the record of
$u$ (Definition~\ref{is-height-numbers}) has kernel-side numbers with constant $C$ whose sizes
satisfy, at every infinite place $v$ of $L$,
\[
 \langle s_i\varphi,s_i\varphi\rangle_v\le e^{\mathrm{sz}(i+1,v)}\langle\varphi,\varphi\rangle_v
 \quad(0\le i<m-2),\qquad
 \langle s_{m-2}\varphi,s_{m-2}\varphi\rangle_v\le e^{\mathrm{sz}(m-1,v)+\mathrm{sz}(m,v)}
 \langle\varphi,\varphi\rangle_v
\]
for all $\varphi$.
\end{proposition}
\leanref{Development/AbelianVarietyTheory/Proofs/DivisorProofs/KernelHalfChainNumbersAnyLevelPrimeStatement.lean}{GlobalTheory.T.kernelHalf_chainNumbers_any_levelPrime}
\leanref{Development/AbelianVarietyTheory/Proofs/DivisorProofs/KernelHalfChainNumbersAnyLevelPrime.lean}{GlobalTheory.KernelHalfS.kerSide_chainS}
\leanref{Development/AbelianVarietyTheory/Proofs/DecoratedLatticeProofs/NormsMatrixStepsCoupledFactorBoundedAnyLevelPrimStatement.lean}{GlobalTheory.T.norms_matrixSteps_coupledFactor_bounded_anyLevelPrime}
\begin{proof}
Put $U=fm^2N^2\log(g_0g_1T_u)$, where $N=N(m)$. The coefficients of the theta kernels lie in a
number field $M\subset\mathbf C$ that contains the image of $L$. Fix an infinite place $v$ of $L$.
There are weights $w_v(\rho)\ge0$ on the embeddings $\rho$ of $M$ that lie over $v$, with
$\sum_\rho w_v(\rho)=1$, and whose average over the places $v$ is $1/[M:\mathbf Q]$ for every
$\rho$.

First, a single kernel entry. The coefficient functions of the kernel chain at the base tuple vanish
off the modules $H_S(N,N')$ of their pair of lattices, are invariant under $T_ug_0g_1H_S(N,N')$, and
are bounded by $(g_0g_1)^{C_dfm^2N^2}$ under every embedding, for a constant $C_d$ that depends on
nothing; the same holds for the two kernels of the coupling. Let $C_K$ be the constant of
Proposition~\ref{is-kernel-entry-bound} for $C_d$. Then for lattices of levels $j-1$ and $j$ the
kernel entry of the $j$-th lift between them has squared operator norm at most $e^{k_j(\rho)}$ at the
embedding $\rho$, with the exponent $k_j(\rho)$ of that proposition; the norms of the forms are the
Hermitian norms of Definition~\ref{is-polynomial-model}(7) at $\rho$, which are nonnegative by
Lemma~\ref{is-fock-facts}(5).

Second, counting. Lemma~\ref{is-class-counts} gives a constant $E\ge0$ such that the number of
members and the number of lattice classes at each level $j\le m$, and the number of cosets in the
coupling set, are at most $e^{EU}$.

Third, Cauchy-Schwarz. The matrix of $s_i$ is a sum over pairs of lattice classes of kernel entries
weighted by the inverse orders of stabilizers. The Gram matrices of the forms are positive
semidefinite, so the Cauchy-Schwarz inequality over the classes bounds the squared norm of $s_i$ at
$\rho$ by the counts at the levels $i$ and $i+1$ times $e^{k_{i+1}(\rho)}$. The coupling is bounded
in the same way by the counts at the levels $m-2$, $m-1$ and $m$, the number of cosets, and
$e^{k_{m-1}(\rho)+k_m(\rho)}$. Since the inner product at $v$ does not depend on $\rho$ over $v$,
the bound at $v$ holds with the weighted average over $\rho$ of these exponents.

Fourth, averages. By the Chowla-Selberg formula (Theorem~\ref{cv-chowla-selberg}(1)), the average
of $\log\Omega_\tau(\rho)$ over all embeddings $\rho$ is $\log\Omega$. By the size conditions,
$\log(r_E/2)+|\operatorname{Re}\frac{L'}{L}(1,\chi_d)|\le2x^{\varepsilon}$, so
Lemma~\ref{is-vacuum-norm-bound} with $X=2x^{\varepsilon}$ and the ranks $N(0)$ and $N(1)$, divided by
$N(0)N(1)$, shows that the average of $\ell_\theta(\rho)$ is at most $2\max(C_t,0)fx^{\varepsilon}$.

Now take $g_0(j)=N(j-1)N(j)$, $\log\Omega_\tau(v)=\sum_\rho w_v(\rho)\log\Omega_\tau(\rho)$,
$\log\theta_0(v)=\frac1f\sum_\rho w_v(\rho)\ell_\theta(\rho)$, $\log h_j=\log h'_j=0$,
$v_0=r_E/2$, $\ell_g=\log(g_0g_1)$, and $M(j,v)$ equal to $C_KU$ plus half the logarithms of the
products of the counts at the levels $j-1$ and $j$. Then the averaged exponents are the sizes
$\mathrm{sz}(j,v)$. With $C=C_K+2E+2\max(C_t,0)+20+2\max(C_Y,0)$, condition (1a) is the
Chowla-Selberg average, (1b) follows from the counting bound, (1c) from the bound on the vacuum
vectors, and the remaining conditions from the size conditions with constant $C_Y$.
\end{proof}

\begin{proposition}\label{is-top-period-bound}
Let $D$ be a definite-side datum with size data at scale $x$ and kernel data that satisfy the kernel
ties with weights $W$ (Definition~\ref{is-kernel-ties}), and let $u$ be a word. For
$\tau\in\mathcal T$ let $\lambda(\tau)$ be the weight of the forms of $u$ at the top, of length
$N=N(m)$, and $\sigma(\tau)$ the weight of the evaluation space, of length $N-1$. Let
$n_\tau=\sum_i\lambda(\tau)_i-\sum_i\sigma(\tau)_i$; let $c_\tau\in\mathbf Z$ be the product of
the differences $(y-i)-(\sigma(\tau)_l-l)$ over the indices $i,l$ and the integers
$y\ne\sigma(\tau)_i$ of the branching interval of the $i$-th entry; let $\kappa_\tau>0$ be the ratio
of the Fock norms of the branching vector and of the extremal vector of weight $\lambda(\tau)$; and
put
\[
 \ell_u=-\Big(\sum_\tau n_\tau\Big)\log T_u+\log\prod_\tau c_\tau^2\kappa_\tau.
\]
\begin{enumerate}
\item For every $\tau$ the weights interlace,
  $\lambda(\tau)_{i+1}\le\sigma(\tau)_i\le\lambda(\tau)_i$, and $\lambda(\tau)$ and $\sigma(\tau)$
  have the same determinant shift.
\item At every infinite place $v$ of $L$, the period of a form $\psi$ of the evaluation space
  against a form $F$ of weight $\sigma$ at the base tuple satisfies
  \[
   |\sigma_v(\mathcal P(\psi,F))|^2\le e^{\ell_u}\langle\psi,\psi\rangle_v\langle F,F\rangle_v.
  \]
\item If $g_i-1>2fk$ for $i=0,1$, $2fN(m)+1<4a$, the weights of every level $j\le m$ have spread
  at most $2a$, $4a<T_{u'}$ for every word $u'$, the size conditions with constant $C_Y$ hold for
  $u$, and $\log T_u\le C_Tmx^{\varepsilon}$, then the record of $u$
  (Definition~\ref{is-height-numbers}) has last-side numbers with constants $C_T$ and $\ell_u$.
\end{enumerate}
\end{proposition}
\leanref{Development/AbelianVarietyTheory/Proofs/GlobalInterpolation/NonnegSubOfLt.lean}{GlobalTheory.BranchMembers.interlace}
\leanref{Development/AbelianVarietyTheory/BlockTheory/GlobalSignature/InterlacingEquivariance.lean}{GlobalTheory.IndefiniteRowSideFullHecke.InterlacingEquivariance.members_sameShift}
\leanref{Development/AbelianVarietyTheory/Proofs/DivisorProofs/IpAtNonnegOfDatFramed.lean}{GlobalTheory.ColRowsTop.period_le}
\leanref{Development/AbelianVarietyTheory/Proofs/DivisorProofs/LastMapNorm.lean}{GlobalTheory.BranchNorm.lastMapNorm}
\leanref{Development/AbelianVarietyTheory/Proofs/DivisorProofs/LastMapHalfRecNumsOfLogLevelPrimeLeStatement.lean}{GlobalTheory.T.lastMapHalf_recNums_of_log_levelPrime_le}
\leanref{Development/AbelianVarietyTheory/Definitions/Divisors/FactorBounds.lean}{GlobalTheory.ColRows.lastLog}
\begin{proof}
(1) The top weight of $u$ is given by the $u$-part of its parameter $\lambda_u(\tau)$ followed by
$m$ constant blocks determined by the $e$-part (Definition~\ref{th-definite-datum}(7)). The weight
$\sigma(\tau)$ is given by the entries $v_i(\tau)$ on its first $k-1$ entries and by the $b$- and
$e$-parts on the rest; these are the same for all words. On the first $k-1$ entries the interlacing
is the interlacing $u_1>v_1>u_2>\cdots$ of the parameter (Definition~\ref{th-hc-parameter}(2)). On
the rest it follows from $-b_1<u_k$ and the gaps between the $b$- and $e$-parts
(Definition~\ref{th-hc-parameter}(3) and (4)). The equality of the determinant shifts is a property
of well-formed weights.

(2) Let $\pi_u\in\mathcal O_E$ be an element of norm $T_u$, so that $\pi_u\bar\pi_u=T_u$, and put
$t_\tau=\tau(\bar\pi_u)$; then $|s(t_\tau)|^2=T_u$ for every embedding $s$. The forms $\psi$ and $F$
are functions on the classes $c$ of the top with values in $V(\lambda)$ and $V(\sigma)$
(Definition~\ref{is-polynomial-model}), the inner products are
$\langle\psi,\psi\rangle_v=\sum_cw_c^{-1}\|\psi_c\|^2_{\sigma_v,1}$ for the weights $w_c$ of the classes,
and the period is
\[
 \mathcal P(\psi,F)=\sum_cw_c^{-1}\langle\mathrm{br}_t\psi_c,F_c\rangle_\sigma
\]
with the branching map of Definition~\ref{is-branching-data}. By (1), Proposition~\ref{is-branching-norm}
applies with $G'_\tau=1$, $g_\tau=1$ and $T=T_u$, and gives
$\|\mathrm{br}_t\psi_c\|^2_{\sigma_v,1}\le\beta\|\psi_c\|^2_{\sigma_v,1}$ with
$\beta=\prod_\tau T_u^{-n_\tau}c_\tau^2\kappa_\tau=e^{\ell_u}$. The Hermitian pairings are positive
(Lemma~\ref{is-fock-facts}(5)), so the Cauchy-Schwarz inequality in $V(\sigma)$ and then over the
classes gives
\[
 |\sigma_v(\mathcal P(\psi,F))|^2\le\Big(\sum_cw_c^{-1}\|\mathrm{br}_t\psi_c\|\,\|F_c\|\Big)^2
 \le e^{\ell_u}\langle\psi,\psi\rangle_v\langle F,F\rangle_v.
\]

(3) Let $\lambda_1(\tau)$ be the first entry of $\lambda(\tau)$, let $D_0(\tau)$ and $D_1(\tau)$ be
the two depths of the branching pattern, so that $n_\tau=\lambda_1(\tau)-D_0(\tau)$, put
$\log c_\tau=(N-1)D_1(\tau)\log(2a)$, $\log\xi_\tau=\frac{N(N-1)}2\log(2a)$ and
$\beta_\tau=-u_k(\tau)-\frac{k-1}2$. Since the entries of $\lambda(\tau)$ have spread at most $2a$,
the branching constants satisfy $\log(c_\tau^2\kappa_\tau)\le2\log c_\tau+\log\xi_\tau$, which is
(2a). By the weight tie, $\lambda_1(\tau)=u_1(\tau)-\frac{k-1}2-\frac{m(t+1)}2$; since
$N-k=m(t+1)$ this is $a_\tau-\frac{N-k}2$ for $\tau\in T$, and for $\tau\notin T$ it is at least
$S_c-\frac{N-1}2$ because $u_1(\tau)\ge S_c$ by the size conditions. The bounds on $D_0(\tau)$ and
$D_1(\tau)$ follow from the windows $|b_i(\tau)-e_j(\tau)|\le\eta$ and $0\le v_i(\tau)\le S_c$, and
$\beta_\tau\le2b'$ for $\tau\in T$ is a size condition. Finally $\log T_u\ge0$ since $T_u$ is a
prime.
\end{proof}

\begin{theorem}\label{is-definite-side-condition-at-bundle}
Let $C_Y$ and $C_T\ge0$ be real numbers. There is a constant $C_0$ with the following property. Let
$D$ be a definite-side datum (Definition~\ref{th-definite-datum}) with coefficient field $L$,
division lattices $N_u$ and size data at scale $x$ with exponent $\varepsilon$
(Definition~\ref{pd-size-data}), with support, modules $N^\circ_u$ and numbers $s(u,w)$, and put
$I_u=\prod_w\mathfrak p_w^{-\lfloor e_ws(u,w)\rfloor}$, the product over the places $w$ of the
support. Let $c$ be real and let $b$ be a bundle at $D$ with constant $c$ that satisfies the
evaluation conditions (Definition~\ref{is-kernel-ties}). Suppose that:
\begin{enumerate}
\item the height conditions with constant $C_Y$ hold, $X'\ge8$, $x\ge2$, and $4a<T_u$ and
  $\log T_u\le C_Tmx^{\varepsilon}$ for every word $u$;
\item $N_u$ is the $\mathcal O_L$-module spanned by $I_uN^\circ_u$, for every word $u$;
\item every $N_u$ is finitely generated, spans its $L$-vector space, and has positive rank;
\item $C_0\le c$.
\end{enumerate}
Then the definite-side condition (Definition~\ref{is-definite-side-condition}) holds for $b$ and
$(C,C_H)$, for all real $C$ and $C_H$.
\end{theorem}
\leanref{Development/AutomorphicArithmetic/ThetaCorrespondence/ThetaPairLifts/ThirdPlaceSeries/PointEvaluationSide.lean}{GlobalTheory.Fc.RL.DefSide.definiteSideCondition_of_nprimeDef_of_nuFull}
\leanref{Development/AutomorphicArithmetic/ThetaCorrespondence/ThetaPairLifts/ThirdPlaceSeries/RowsUOfColPinsStatement.lean}{GlobalTheory.T.kernel_columns_from_ties_alt}
\leanref{Development/AbelianVarietyTheory/Proofs/DivisorProofs/MemberRowsOfKernelDatumSizesStatement.lean}{GlobalTheory.T.memberRows_of_kernelDatum_sizes}
\leanref{Development/AbelianVarietyTheory/Proofs/DivisorProofs/ColumnsMemberRowsOfLogLevelLeStatement.lean}{GlobalTheory.T.columnsMemberRows_of_logLevel_le}
\leanref{Development/AbelianVarietyTheory/Proofs/DivisorProofs/ColumnsMemberRowsOfLogLevelLe.lean}{GlobalTheory.ColRowsTopS.rowsBody_colFace}
\begin{proof}
Let $C$ be the constant of Proposition~\ref{is-kernel-step-bound} for $C_Y$, and put
$C_0=\max\big(2(16+C_T),\,6C+1+2C^2,\,4C(C_T+C)\big)$.

Assume the hypotheses of the definite-side condition for $(C,C_H)$: primitivity, the covolume,
surjectivity and isometry conditions, record numbers for every word with constants
$\lambda_0,C'_0,c_e,c,S_w,C_1$, and division data tied to the realization. They serve to form the
division package of the realization of $b$. The rows, columns and column face of this package are
those of $b$, and the division data attach to each word $u$ the column of $u$. So it suffices to
prove the rows condition for each word $u$, the columns of $b$ and the record of $u$.

By the conditions on the square-integrable columns of $b$, the Hermitian matrix $H_V(\mu_0)$ of the
columns is positive definite at every infinite place. By the Cholesky decomposition it has a metric
$X^0$.

Fix a word $u$. By the evaluation conditions and hypothesis (1), the kernel data of $b$ satisfy the
kernel ties with the weights of its columns and the kernel sizes with $(C_Y,C_T)$. In particular
the tuple of characters at $\mu_0$ is trivial. Let $s_0,\dots,s_{m-2}$ be the kernel chain of $u$
at $\mu_0$, with spaces $F_0,\dots,F_{m-2}$ and $S_u$.

The column of $u$ is $L^{c_0}$ with the metric $X_c$ of $M_{u,\mathrm{ram}}$, and the first column
map $f_0\colon L^{c_0}\to F_0$ is an isometry for it, with $X_c$ invertible at every infinite place.
The spaces $F_j$ and $S_u$ are metrized, because their Gram matrices are positive definite. The
period pairing of $u$ at $\mu_0$ is $\mathrm{Pt}(\varphi,F)=g(\operatorname{col}f_0\varphi)(F)$,
where $g(\psi)(F)=\mathcal P(\psi,\operatorname{ev}F)$ is the period of $\psi\in S_u$ against the
evaluation $\operatorname{ev}F$ of $F$ at the top.

Proposition~\ref{is-kernel-step-bound} gives kernel-side numbers with constant $C$ for the record
of $u$. Their sizes bound the maps $s_i$ as in hypothesis (1) of
Lemma~\ref{is-chain-from-kernels}, with $n=m-2$ and $o(j)=j-1$: the $j$-th lift belongs to the
step $j-1$ for $j<m-1$, and the last two lifts to the coupling.

Proposition~\ref{is-top-period-bound}(3) gives last-side numbers with constants $C_T$ and
$\ell_u$. Since $c\ge C_0$, Lemma~\ref{is-height-numbers-from-sides} gives height numbers with
period exponent $\ell_u$ whose sizes are those of the kernel-side numbers.

The column characters of $b$ are compatible with its rows, and this gives
$\langle\operatorname{ev}F,\operatorname{ev}F\rangle_v\le|X^0_v{}^{-1}\sigma_v(F)|^2$. With
Proposition~\ref{is-top-period-bound}(1) and (2),
\[
 |\sigma_v(g(\psi)(F))|^2\le e^{\ell_u}\langle\psi,\psi\rangle_v
 \langle\operatorname{ev}F,\operatorname{ev}F\rangle_v
 \le e^{\ell_u}\langle\psi,\psi\rangle_v|X^0_v{}^{-1}\sigma_v(F)|^2,
\]
which is hypothesis (2) of Lemma~\ref{is-chain-from-kernels} with $h_P=\ell_u$.

Lemma~\ref{is-chain-from-kernels} now gives a chain of $m-1$ lifts, tied to the kernel chain of the
realization with first columns $f_0$, and conditions (1) to (5) of the rows condition. For
condition (6): $X_c$ is invertible; $c_0$ is the rank of $N_u$, which is positive by hypothesis
(3); and $N_u$, viewed in $L^{c_0}$, is a lattice, because by hypotheses (2) and (3) it is a
finitely generated $\mathcal O_L$-module that spans.
\end{proof}

\begin{theorem}\label{is-definite-side-condition-holds}
Let $p$ be parameters for the construction of bundles, let $Q_p$ be the slot of the good base points
of $p$, and let $\Lambda_p$ be the link data that link a bundle to a base point when the bundle is
built from $p$ at that point (Definition~\ref{is-slot-rows}). Then the following holds along
$(Q_p,\Lambda_p)$. For every class $\mathcal A$ of families and every selection $s$, eventually along
the designed families $\mathcal A^{\mathrm{des}}$, every family $\mathcal F$ has the property that for
every $C\ge0$, eventually along $\mathcal F$, the definite-side condition
(Definition~\ref{is-definite-side-condition}) holds for the selected bundle and $(C,C_H)$. In
particular hypothesis (9) of Theorem~\ref{is-definite-rows} holds along $(Q_p,\Lambda_p)$, with
$C_0=0$.
\end{theorem}
\leanref{Development/AutomorphicArithmetic/ThetaCorrespondence/ThetaPairLifts/ThirdPlaceSeries/PointEvaluationSideStatement.lean}{GlobalTheory.T.weil_unitary_entries'}
\leanref{Development/AutomorphicArithmetic/ThetaCorrespondence/ThetaPairLifts/ThirdPlaceSeries/PointEvaluationSide.lean}{GlobalTheory.Fc.RL.linkedAt_prmR_df_of_cuspRapidDecay}
\leanref{Development/AutomorphicArithmetic/ThetaCorrespondence/ThetaPairLifts/ThirdPlaceSeries/PointEvaluationSide.lean}{GlobalTheory.Fc.RL.DefSide.linkedAt_df_of_parameters}
\leanref{Development/AutomorphicArithmetic/ThetaCorrespondence/ThetaPairLifts/ThirdPlaceSeries/RowGood.lean}{GlobalTheory.Fc.RL.RowDefectCondition}
\leanref{Development/AutomorphicArithmetic/ThetaCorrespondence/ThetaPairLifts/ThirdPlaceSeries/RowGood.lean}{GlobalTheory.Fc.RL.LinkedAt}
\begin{proof}
Take for the condition on bundles the evaluation conditions (Definition~\ref{is-kernel-ties}).
Every bundle that $\Lambda_p$ links to a base point satisfies them. Such a bundle comes with a
witness of its construction from $p$: its theta realization is tied at the real place, its columns
are the symmetric column forms, whose characters are trivial at step zero, its column characters
are compatible with its rows, and the kernel data of its realization are the standard kernel data
of its true columns. Standard kernel data satisfy the kernel ties and, under the hypotheses on the
size data in Definition~\ref{is-kernel-ties}(5), the kernel sizes.

For real $C_Y$ and $C_T$ let $c_0(C_Y,C_T)$ be the constant of
Theorem~\ref{is-definite-side-condition-at-bundle} for $C_Y$ and $\max(C_T,0)$. Let $\mathcal A$ be
a class and $s$ a selection whose bundles are built from $(Q_p,\Lambda_p)$ and satisfy the
evaluation conditions. The selected bundles have constant $c_{sz}$, the size constant of the family.
This is a constant of the middle list of the family, and the constants $C_{YL}$ and $C_{T_u}$ are in
its outer list (Definition~\ref{pd-family}). So the largeness threshold $\kappa_{\mathrm{mid}}=c_0$
shows that $c_{sz}\ge c_0(C_{YL},C_{T_u})$ eventually along $\mathcal A^{\mathrm{des}}$.

Fix such a family $\mathcal F$ in $\mathcal A^{\mathrm{des}}$ and $C\ge0$. The selection supplies
the second-layer conditions at the real place. With the conditions of a designed family, they show
that at all large scales every division lattice $N_u$ is finitely generated, spans, and has positive
rank. The size condition on $T_u$ of a designed family gives, at all large scales, the height
conditions with constant $C_{YL}$, $X'\ge8$, $x\ge2$, and $4a<T_u$ and
$\log T_u\le\max(C_{T_u},0)mx^{\varepsilon}$ for every word $u$. At every member the division
lattices are $N_u=I_uN^\circ_u$. Beyond the largest of these scales,
Theorem~\ref{is-definite-side-condition-at-bundle} applies at every member with $c=c_{sz}$, and
gives the definite-side condition for the selected bundle and $(C,C_H)$.
\end{proof}

\section{The division inequality}\label{section-is-division-inequality}

\begin{definition}\label{is-local-indices}
Let $K$ be a totally complex number field with an involution $a\mapsto\bar a$ and $r$ infinite
places, and let $d\geq1$. A \emph{lattice} in $K^d$ is a finitely generated $\mathcal O_K$-submodule
that spans $K^d$. Let $A,B\subset K^d$ be lattices and let $w$ be a finite place of $K$, with
conjugate place $\bar w$ and $\ell_w$ as in Definition~\ref{is-indefinite-datum}.
\begin{enumerate}
\item The \emph{index} of $B$ relative to $A$ at $w$ is
  $\operatorname{idx}_w(A,B)=\operatorname{len}_w\big(A/(A\cap B)\big)-\operatorname{len}_w\big(B/(A\cap B)\big)$.
  Thus $I_w(j,t)=\operatorname{idx}_w(M_j,\theta_{j,t}T_{j,t})$ in
  Definition~\ref{is-indefinite-datum}.
\item A \emph{local basis} of $A$ at $w$ is a basis $e=(e_1,\dots,e_d)$ of $K^d$ that is a basis
  of the localization of $A$ at the multiplicative set of elements of $\mathcal O_K$ lying in
  neither $w$ nor $\bar w$. For $G\in M_d(K)$, the \emph{Gram determinant} of $e$ for $G$ is
  $\det\big(\bar e_i^{\mathsf T}Ge_k\big)_{i,k}\in K$.
\item For an invertible Hermitian matrix $H\in M_d(K)$, let
  $A^\vee_H=\{y\in K^d:\bar x^{\mathsf T}Hy\in\mathcal O_K\text{ for all }x\in A\}$, and put
  $j_w(H,A)=\operatorname{idx}_w(A,A^\vee_H)$, the \emph{relative index} of $A$ and its dual at $w$.
\item For a metric $X$ on $K^d$ (Definition~\ref{is-indefinite-datum}) and an infinite place $v$,
  put $\langle y,z\rangle_{X,v}=(X_v^{-1}\sigma_v(y))^*X_v^{-1}\sigma_v(z)$. The \emph{Gram
  determinant} of $y_1,\dots,y_d$ in $X$ at $v$ is $\det(\langle y_i,y_k\rangle_{X,v})$. For a
  metric $X'$ on $K^{d'}$ and $\varphi\in M_{d'\times d}(K)$, $r_v(X,X',\varphi)$ is the Gram
  determinant of $\varphi e_1,\dots,\varphi e_d$ in $X'$ at $v$ divided by that of
  $e_1,\dots,e_d$ in $X$ at $v$, for the standard basis $(e_i)$ of $K^d$.
\item The \emph{mean} of a function $c$ on the infinite places is
  $\overline{\log c}=\frac1r\sum_v\log c(v)$.
\end{enumerate}
\end{definition}
\leanref{Formalism/QuotientMetricDivisors/Statement.lean}{Formalism.DivisionValues.idx}
\leanref{Formalism/QuotientMetricDivisors/Statement.lean}{Formalism.DivisionValues.IsLocBasis}
\leanref{Formalism/QuotientMetricDivisors/Statement.lean}{Formalism.DivisionValues.gramR}
\leanref{Formalism/QuotientMetricDivisors/Statement.lean}{Formalism.DivisionValues.jw}
\leanref{Formalism/QuotientMetricDivisors/Statement.lean}{Formalism.DivisionValues.detRatio}

\begin{lemma}\label{is-lattice-facts}
Let $K$, $r$ and $d$ be as in Definition~\ref{is-local-indices}, let $A,B,C\subset K^d$ be
lattices, let $w$ be a finite place and let $X$ be a metric on $K^d$.
\begin{enumerate}
\item $\operatorname{idx}_w(A,C)=\operatorname{idx}_w(A,B)+\operatorname{idx}_w(B,C)$, and
  $\operatorname{idx}_w(A,B)=0$ for all but finitely many $w$.
\item Let $e$ be a basis of $K^d$, let $\gamma\in GL_d(K)$, and suppose that $e\gamma$ is a local
  basis of $A$ at $w$. Then $\operatorname{idx}_w(\mathcal O_Ke,A)=v_w(\det\gamma)$ and
  $\operatorname{idx}_{\bar w}(\mathcal O_Ke,A)=v_{\bar w}(\det\gamma)$, where $\mathcal O_Ke$ is
  the lattice spanned by $e$. For every $G\in M_d(K)$, the Gram determinant of $e\gamma$ for $G$ is
  $\overline{\det\gamma}\,\det\gamma$ times that of $e$.
\item $v_w(\bar a)=v_{\bar w}(a)$ for $a\in K^\times$, and $\ell_{\bar w}=\ell_w$.
\item For a basis $e$ of $K^d$ with Gram determinants $g_v(e)$ in $X$,
  $\operatorname{lcv}_X(A)=\frac1r\sum_v\log g_v(e)+2\sum_w\ell_w\operatorname{idx}_w(\mathcal O_Ke,A)$.
\item For a metric $X'$ on $K^d$ and $\varphi\in GL_d(K)$,
  $\operatorname{lcv}_{X'}(\varphi A)=\operatorname{lcv}_X(A)+\frac1r\sum_v\log r_v(X,X',\varphi)$.
\item With the normalizations of Definition~\ref{hs-slopes}, the dual lattice $A^\vee$ in the dual
  metric has slope $\mu(A^\vee)=\operatorname{lcv}_X(A)/d$.
\item For $a\in K^\times$, $\frac1r\sum_v\log|\sigma_v(a)|=\sum_w\ell_wv_w(a)$.
\item Let $K_0\subset K$ be a subfield with $\sigma(\bar a)=\overline{\sigma(a)}$ for every complex
  embedding $\sigma$ and every $a\in K_0$. If $H\in M_d(K_0)$ is Hermitian and
  $(X_v^{-1})^*X_v^{-1}=\sigma_v(H)$ for every infinite place $v$, then
  $\operatorname{lcv}_X(A)=-\sum_w\ell_wj_w(H,A)$. For any invertible Hermitian $H$ and any local
  basis $f$ of $A$ at $w$, $j_w(H,A)=-v_w(\mathrm{Gram}_H(f))$, where $\mathrm{Gram}_H(f)$ is the
  Gram determinant of $f$ for $H$.
\end{enumerate}
\end{lemma}
\leanref{Formalism/QuotientMetricDivisors/ProdFormula.lean}{Formalism.DivisionValues.IdxChange}
\leanref{Formalism/QuotientMetricDivisors/HermitianTwistDivisors/IdxChange.lean}{Formalism.DivisionValues.index_change_formula}
\begin{proof}
(1) The lengths are additive along $A\supset A\cap B\cap C$, and $A$ and $B$ agree at almost all
places. (2) Both lattices are free over the semilocal ring at $\{w,\bar w\}$, and the index of two
free lattices with bases $e$ and $e\gamma$ is the valuation of $\det\gamma$ at each of the two
places; the Gram matrix of $e\gamma$ is $\bar\gamma^{\mathsf T}$ times that of $e$ times $\gamma$.
(3) The involution carries the valuation ring of $w$ onto that of $\bar w$ and preserves norms.
(4) For $A=\mathcal O_Ke$ the covolume is $\prod_vg_v(e)$ up to the powers of $2$ and of the
discriminant that $\operatorname{lcv}$ removes; a change of lattice multiplies the covolume by the
norm of the index, which is $\prod_w Nw^{\operatorname{idx}_w(\mathcal O_Ke,A)}$. (5) Apply (4) to
the bases $e$ of $A$ and $\varphi e$ of $\varphi A$, which have the same indices. (6) The dual basis
has the inverse Gram matrix. (7) This is the product formula with $[K:\mathbf Q]=2r$. (8) For the
standard basis the Gram determinant at $v$ is $\sigma_v(\det H)>0$, so (4) and (7) give
$\operatorname{lcv}_X(A)=\sum_w\ell_w\big(v_w(\det H)+2\operatorname{idx}_w(\mathcal O_K^d,A)\big)$.
By (2) and (3), $v_w(\mathrm{Gram}_H(f_w))=v_w(\det H)+\operatorname{idx}_w(\mathcal O_K^d,A)
+\operatorname{idx}_{\bar w}(\mathcal O_K^d,A)$ for local bases $f_w$, and summing over $w$ with
$\ell_{\bar w}=\ell_w$ gives $\operatorname{lcv}_X(A)=\sum_w\ell_wv_w(\mathrm{Gram}_H(f_w))$. The last
statement is the computation of the discriminant of a free Hermitian lattice: the dual of the
lattice with basis $f$ has the basis $f\,\mathrm{Gram}^{-1}$, whose index is
$-v_w(\mathrm{Gram}_H(f))$ by (2).
\end{proof}

\begin{lemma}\label{is-gram-covolume}
Let $K$ be a totally complex number field with an involution, let $K_0\subset K$ be a subfield
with $\sigma(\bar a)=\overline{\sigma(a)}$ for every complex embedding $\sigma$ and $a\in K_0$, and let
$v_0$ be an infinite place with $\sigma_{v_0}(\bar a)=\overline{\sigma_{v_0}(a)}$ for all $a\in K$. Let
$H\in M_n(K_0)$ be Hermitian and positive definite at every infinite place, let
$P\in M_{D\times n}(K)$ be surjective, let $X_Q$ be the quotient metric, with
$(X_{Q,v}^{-1})^*X_{Q,v}^{-1}=(\sigma_v(P)\sigma_v(H)^{-1}\sigma_v(P)^*)^{-1}$ at every $v$, and put
$Q=(PH^{-1}\bar P^{\mathsf T})^{-1}\in M_D(K)$. Let $M\subset K^n$ be a lattice and, for every
finite place $w$, let $e_w$ be a local basis of $\Lambda=PM$ at $w$. Then
$\operatorname{lcv}_{X_Q}(\Lambda)\leq\sum_w\ell_wv_w(\mathrm{Gram}_Q(e_w))$, where all but finitely
many terms vanish.
\end{lemma}
\leanref{Formalism/QuotientMetricDivisors/Statement.lean}{Formalism.DivisionValues.LemmaRb}
\leanref{Formalism/QuotientMetricDivisors/HermitianTwistDivisors/ReductionLemmaSecond.lean}{Formalism.DivisionValues.lemmaRb_of}
\begin{proof}
Let $e$ be the standard basis of $K^D$ and $G=\mathrm{Gram}_Q(e)=\det(PH^{-1}\bar P^{\mathsf
T})^{-1}$, which is nonzero because its image under $\sigma_{v_0}$ is the Gram determinant of $e$
in the positive definite form of $X_Q$ at $v_0$. Let $g_v$ be the Gram
determinant of $e$ in $X_Q$ at $v$, so $g_v=\det\big(\sigma_v(P)\sigma_v(H)^{-1}\sigma_v(P)^*\big)^{-1}$.
For an infinite place $v$ let $v'$ be the place of the embedding $a\mapsto\sigma_v(\bar a)$, so that
$\sigma_v(\bar a)=\overline{\sigma_{v'}(a)}$, and $\sigma_{v'}(H)=\sigma_v(H)$ because $H$ has
entries in $K_0$. Then
$\sigma_v(G)^{-1}=\det\big(\sigma_v(P)\sigma_v(H)^{-1}\sigma_{v'}(P)^*\big)$. The Cauchy-Schwarz
inequality for determinants of Gram matrices of the positive definite form $\sigma_v(H)^{-1}$ gives
$|\sigma_v(G)|^{-2}\leq g_v^{-1}g_{v'}^{-1}$, that is, $\log g_v+\log g_{v'}\leq2\log|\sigma_v(G)|$.
Since $v\mapsto v'$ permutes the infinite places, summing over $v$ and applying
Lemma~\ref{is-lattice-facts}(7) gives $\frac1r\sum_v\log g_v\leq\sum_w\ell_wv_w(G)$.

By Lemma~\ref{is-lattice-facts}(4),
$\operatorname{lcv}_{X_Q}(\Lambda)=\frac1r\sum_v\log g_v+2\sum_w\ell_w\operatorname{idx}_w(\mathcal
O_K^D,\Lambda)$. Write $e_w=e\gamma_w$. By Lemma~\ref{is-lattice-facts}(2) and (3),
$v_w(\mathrm{Gram}_Q(e_w))=v_w(G)+\operatorname{idx}_w(\mathcal O_K^D,\Lambda)
+\operatorname{idx}_{\bar w}(\mathcal O_K^D,\Lambda)$. Summing over $w$ with $\ell_{\bar
w}=\ell_w$ gives
$\sum_w\ell_wv_w(\mathrm{Gram}_Q(e_w))=\sum_w\ell_wv_w(G)+2\sum_w\ell_w\operatorname{idx}_w(\mathcal
O_K^D,\Lambda)$, and the two displays give the claim.
\end{proof}

\begin{lemma}\label{is-local-index-identity}
Let $K$ be a totally complex number field with an involution, let $v_0$ be an infinite place with
$\sigma_{v_0}(\bar a)=\overline{\sigma_{v_0}(a)}$ for all $a\in K$, and let $w$ be a finite place.
Let $X_E$ be a metric on $K^D$. For $i=0,1$ let $X_{Q,i}$ be a metric on $K^D$, $\Lambda_i\subset
K^D$ a lattice, $\iota_i\in GL_D(K)$, $e_i$ a local basis of $\Lambda_i$ at $w$, and $G_i\in K$ an
element whose image $\sigma_{v_0}(G_i)$ is the Gram determinant of $e_i$ in $X_{Q,i}$ at $v_0$. Let
$J$ be a finite set and, for $j\in J$, let $X_W(j)$ be a metric on $K^{d_j}$ and $M_j\subset
K^{d_j}$ a lattice, and for $i=0,1$ let $H_i(j)\in M_{d_j}(K)$ be invertible Hermitian, $X_{T,i}(j)$
a metric on $K^{d_j}$ with $(X_{T,i}(j)_{v_0}^{-1})^*X_{T,i}(j)_{v_0}^{-1}=\sigma_{v_0}(H_i(j))$,
$T_i(j)\subset K^{d_j}$ a lattice and $\theta_i(j)\in GL_{d_j}(K)$. Let $\kappa>0$, and let
$\rho\in K$ with $\sigma_{v_0}(\rho)=\rho_r$ real and positive and $v_w(\rho)=0$. Suppose that
\[
 r_{v_0}(X_{Q,0},X_E,\iota_0)=\kappa^D\prod_jr_{v_0}(X_{T,0}(j),X_W(j),\theta_0(j)),\qquad
 r_{v_0}(X_{Q,1},X_E,\iota_1)=(\kappa\rho_r)^D\prod_jr_{v_0}(X_{T,1}(j),X_W(j),\theta_1(j)).
\]
Write $I^i_y(j)=\operatorname{idx}_y(M_j,\theta_i(j)T_i(j))$ and
$\jmath^i(j)=j_w(H_i(j),T_i(j))$. Then
\begin{multline*}
 \operatorname{idx}_w(\iota_0\Lambda_0,\iota_1\Lambda_1)+\operatorname{idx}_{\bar w}(\iota_0\Lambda_0,\iota_1\Lambda_1)
 =v_w(G_1)-v_w(G_0)\\
 +\sum_{j\in J}\Big(\big(I^1_w(j)-I^0_w(j)\big)+\big(I^1_{\bar w}(j)-I^0_{\bar w}(j)\big)+\big(\jmath^1(j)-\jmath^0(j)\big)\Big).
\end{multline*}
\end{lemma}
\leanref{Formalism/QuotientMetricDivisors/HermitianTwistDivisors/ConjugateWeightLatticeOfGeneral.lean}{Formalism.DivisionValues.ConjugateWeightGram.local_idW}
\begin{proof}
For $i=0,1$ and $j\in J$ choose a local basis $f_i(j)$ of $T_i(j)$ at $w$. Write
$\iota_1e_1=(\iota_0e_0)\gamma$ and $\theta_1(j)f_1(j)=(\theta_0(j)f_0(j))\delta_j$ with $\gamma$ and
$\delta_j$ invertible matrices over $K$. By Lemma~\ref{is-lattice-facts}(1) and (2),
$\operatorname{idx}_y(\iota_0\Lambda_0,\iota_1\Lambda_1)=v_y(\det\gamma)$ and
$I^1_y(j)-I^0_y(j)=\operatorname{idx}_y(\theta_0(j)T_0(j),\theta_1(j)T_1(j))=v_y(\det\delta_j)$ for
$y\in\{w,\bar w\}$. By Lemma~\ref{is-lattice-facts}(8),
$\jmath^i(j)=-v_w(g_i(j))$, where $g_i(j)=\mathrm{Gram}_{H_i(j)}(f_i(j))$.

Now compute Gram determinants at $v_0$, where the involution is complex conjugation. By
Lemma~\ref{is-lattice-facts}(2), the Gram determinant of $\iota_1e_1$ in $X_E$ is
$|\sigma_{v_0}(\det\gamma)|^2$ times that of $\iota_0e_0$. By the definition of $r_{v_0}$, the Gram
determinant of $\iota_ie_i$ in $X_E$ is $r_{v_0}(X_{Q,i},X_E,\iota_i)\,\sigma_{v_0}(G_i)$. In the
same way the Gram determinant of $\theta_1(j)f_1(j)$ in $X_W(j)$ is $|\sigma_{v_0}(\det\delta_j)|^2$
times that of $\theta_0(j)f_0(j)$, and that of $\theta_i(j)f_i(j)$ is
$r_{v_0}(X_{T,i}(j),X_W(j),\theta_i(j))\,\sigma_{v_0}(g_i(j))$, because $X_{T,i}(j)$ is a metric for
$H_i(j)$ at $v_0$. Substituting the two hypotheses on $r_{v_0}$ and cancelling the nonzero factors
$\kappa^D$ and $r_{v_0}(X_{T,0}(j),X_W(j),\theta_0(j))$ gives
\[
 \sigma_{v_0}\Big(\rho^DG_1\prod_j\overline{\det\delta_j}\,\det\delta_j\,g_0(j)\Big)
 =\sigma_{v_0}\Big(\overline{\det\gamma}\,\det\gamma\,G_0\prod_jg_1(j)\Big).
\]
Since $\sigma_{v_0}$ is injective, the two elements of $K$ agree. All factors are nonzero, because
Gram determinants of bases in a metric are positive. Take $v_w$ of both sides, using $v_w(\rho)=0$ and
Lemma~\ref{is-lattice-facts}(3):
\[
 v_w(G_1)+\sum_j\big(v_{\bar w}(\det\delta_j)+v_w(\det\delta_j)+v_w(g_0(j))\big)
 =v_{\bar w}(\det\gamma)+v_w(\det\gamma)+v_w(G_0)+\sum_jv_w(g_1(j)).
\]
Substituting the indices of the first paragraph gives the identity.
\end{proof}

\begin{definition}\label{is-division-package}
Let $K$ be a totally complex number field with an involution, $K_0\subset K$ a subfield with
$\sigma(\bar a)=\overline{\sigma(a)}$ for every complex embedding $\sigma$ and $a\in K_0$, and $v_0$
an infinite place with $\sigma_{v_0}(\bar a)=\overline{\sigma_{v_0}(a)}$ for all $a\in K$. Let base
data (Definition~\ref{is-indefinite-datum}(1)) over $K$ with these $K_0$ and $v_0$ be given, with
tuples $\mathcal T$, base tuple $\mu_0$, parts $\Pi$ and lifts $J$, and let $N$ be numerical bounds
(Definition~\ref{is-numerical-bounds}). A definite side
$(H_V,M_t,P_t,X_Q,Q_p,H_T,X_T,T_{j,t},Z_p)$ is \emph{well formed} if for all $t$, $j$ and $v$:
$H_V(t)$ has entries in $K_0$ and is positive definite at every infinite place; $M_t$ and
$T_{j,t}$ are lattices; $P_t$ is surjective, $\Lambda_t=P_tM_t$ is a lattice and $X_Q(t)$ is the
quotient metric of $H_V(t)$ through $P_t$ (Lemma~\ref{is-gram-covolume}); the subspaces
$Q_p(\mu_0)$ are independent and span $K^{d_Q(\mu_0)}$, and the $Z_p(j,\mu_0)$ are independent;
$X_T(j,t)$ is invertible at every infinite place and is a metric for $H_T(j,t)$ at $v_0$. A
\emph{division package} over $K$ for $N$ consists of indefinite-side data $S_0$, the \emph{rows},
and $S_1$, the \emph{columns}, with these base data, on well-formed definite sides, and of
\emph{column data} for $S_1$, such that the following hold.
\begin{enumerate}
\item (Rows.) $d_Q(\mu_0)>0$; $\mu_{\bar w}=\mu_w$ for every $w$, and $w$ is twisted exactly when
  $\mu_w\neq\mu_0$; $X_E$ and the $X_W(j)$ are invertible and the $M_j$ are lattices;
  $\iota_{\mu_0}Q_p(\mu_0)=E_p$ and $\theta_{j,\mu_0}Z_p(j,\mu_0)=W_p(j)$; the matrices $\iota_t$ and
  $\theta_{j,t}$ are bijective, $\kappa_v>0$, $\lambda_{j,t,p,v}>0$, and the scalar conditions after
  Definition~\ref{is-indefinite-datum} hold; $\dim W_p(j)=\dim E_p$; the $Z_p(j,\mu_0)$ span;
  $H_T(j,\mu_0)$ has entries in $K_0$ and $X_T(j,\mu_0)$ is a metric for it at every infinite
  place. The rows satisfy the covolume bound, the surjectivity bound and the functional bound for
  $N$ (Definitions~\ref{is-numerical-bounds}, \ref{is-surjectivity-bound}
  and~\ref{is-functional-bound}), and they are \emph{untwisted}: for every finite place $w$ that is
  not twisted and every lift $j$, $F_0(w,j)=0$ and
  $I_w(j,\mu_w)+I_{\bar w}(j,\mu_w)+j_w(H_T(j,\mu_w),T_{j,\mu_w})=0$.
\item (Columns.) $d_Q(\mu_0)>0$; $\mu_{\bar w}=\mu_w$, and $w$ is twisted exactly when
  $\mu_w\neq\mu_0$; $X_E$ is invertible; $\iota_{\mu_0}Q_p(\mu_0)=E_p$; the $\iota_t$ are
  bijective, $\kappa_v>0$, $\lambda_{j,t,p,v}>0$, and $\iota_t$ multiplies the Hermitian form of
  $X_Q(t)$ at $v$ by $\kappa_v\prod_j\lambda_{j,t,p,v}$ on $\iota_t^{-1}(E_p)$ whenever $t=\mu_0$ or
  $v=v_0$; the columns satisfy the lower norm bound for $N$
  (Definition~\ref{is-lower-norm-bound}).
\item (Column data.) A set of words $u$, and for each word an integer $a(u)\geq0$, a metric $X_c(u)$
  on $K^{a(u)}$, bilinear \emph{period pairings}
  $\mathrm{Pt}_u(t)\colon K^{a(u)}\times K^{n_V(t)}\to K$ for the tuples $t$, a linear map
  $B_u\colon K^{a(u)}\to\operatorname{Hom}_K(K^D,K)$ and scalars $\zeta_u(t)\in K^\times$ with
  \[
   \mathrm{Pt}_u(t)(\varphi,F)=\zeta_u(t)\,B_u(\varphi)(\iota_tP_tF),\qquad \zeta_u(\mu_0)=1,\qquad
   v_w(\zeta_u(\mu_w))=0
  \]
  for all $t$, $\varphi$, $F$ and finite places $w$; here $D$, $\iota_t$ and $P_t$ are those of the
  columns.
\item (Passage.) The rows and the columns have the same map $w\mapsto\mu_w$, the same $\kappa_v$
  and the same $\lambda_{j,t,p,v}$. Write $RR_i$ for the glued lattice of $S_i$ in its metric $X_E$
  and $CC_i=RR_i^\vee$. If $0\leq\delta_e\leq\frac13$ and $\mu_{\max}(RR_0)\leq H$, then
  $\mu_{\max}(RR_1)\leq H+14\ell_a$ and
  \[
   \mu(CC_1)\leq\mu(CC_0)+3\delta_e\big(\mu(CC_0)+H+28\ell_a\big)+28\ell_a.
  \]
\end{enumerate}
\end{definition}
\leanref{Formalism/LocalToGlobal/IndefiniteForms/IndefSide.lean}{Formalism.L2.Indef.DefSide.WF}
\leanref{Formalism/LocalToGlobal/IndefiniteForms/IndefSide.lean}{Formalism.L2.Indef.IndefSide}
\leanref{Formalism/LocalToGlobal/IndefiniteForms/IndefSide.lean}{Formalism.L2.Indef.IndefData.Untwisted}
\leanref{Formalism/LocalToGlobal/IndefiniteForms/IndefSide.lean}{Formalism.L2.Indef.ColData.PairingFactorsThroughPeriods}
\leanref{Formalism/LocalToGlobal/IndefiniteForms/IndefSide.lean}{Formalism.L2.Indef.Passage}
\leanref{Formalism/LocalToGlobal/IndefiniteForms/PairingConditionOfInverseMatrixIdentity.lean}{Formalism.L2.Indef.IndefColsR}
\leanref{Formalism/LocalToGlobal/IndefiniteForms/PairingConditionOfInverseMatrixIdentity.lean}{Formalism.L2.Indef.IndefCoreR}

\begin{lemma}\label{is-glued-covolume}
Let a division package over $K$ for numerical bounds $N$ be given (Definition~\ref{is-division-package}),
and let $D$, $\mathrm{Tw}$, $X_E$, $RR_0$ and $CC_0=RR_0^\vee$ be those of its rows. For every tuple
$t$ and finite place $y$ let $e_{t,y}$ be a local basis of $\Lambda_t=P_tM_t$ at $y$, and let
$G_t(y)$ be its Gram determinant for $(P_tH_V(t)^{-1}\bar P_t^{\mathsf T})^{-1}$. Let $S_y$ be a
finite set of finite places containing $\mathrm{Tw}$ with $v_y(G_{\mu_y}(y))=0$ for $y\notin S_y$,
and let $L_G,E_3$ be real numbers, $E_2(j)$, $g_y$ and $\eta_y(j)$ real numbers for lifts $j$ and
finite places $y$, such that:
\begin{enumerate}
\item $v_y(G_{\mu_y}(y))\leq D\,g_y$ for $y\in S_y$, and $\sum_{y\in S_y}\ell_yg_y\leq L_G$;
\item $j_y(H_T(j,\mu_y),T_{j,\mu_y})\leq e_yJ_0(\mu_y,y,j)D+\eta_y(j)$ for $y\in\mathrm{Tw}$ and
  every $j$, and $\sum_{y\in\mathrm{Tw}}\ell_y\eta_y(j)/D\leq E_2(j)$ for every $j$;
\item $\sum_{y\in S_y}\ell_y\,2e_y/p_y^2\leq E_3$.
\end{enumerate}
Then
\[
 \mu(CC_0)\leq\overline{\log\kappa}+\sum_j\ell_3(j)
 +2\sum_j\sum_{y\in\mathrm{Tw}}\ell_ye_yF_0(y,j)+L_G+\sum_j\big(E_1(j)+E_2(j)\big)+E_3.
\]
\end{lemma}
\leanref{Formalism/LocalToGlobal/IndefiniteForms/IndefUse.lean}{Formalism.L2.Indef.IndefSide.boundsMu}
\leanref{Formalism/LocalToGlobal/IndefiniteForms/IndefUse.lean}{Formalism.L2.Indef.IndefRows.detId}
\leanref{Formalism/DivisorInequality/StatementObj.lean}{Formalism.DivisorInequality.BoundsMu}
\leanref{Formalism/DivisorInequality/DivisorInequalityProofs/BaseLevelBoundsWithCovolumeHolds.lean}{Formalism.DivisorInequality.slope_bounds}
\leanref{Formalism/QuotientMetricDivisors/HermitianTwistDivisors/ConjugateWeightLatticeOfGeneral.lean}{Formalism.DivisionValues.ConjugateWeightGram.stepBW_of}
\leanref{Formalism/QuotientMetricDivisors/HermitianTwistDivisors/BoundsSum.lean}{Formalism.DivisionValues.divisor_bounds_sum}
\begin{proof}
All objects are those of the rows. The dimensions $d_Q(t)$, $d_W(j)$ and $d_Z(j,t)$ all equal $D$:
$\iota_t$ and $\theta_{j,t}$ are bijective, and $d_W(j)=\sum_p\dim W_p(j)=\sum_p\dim E_p=D$ because the
parts are independent and span.
Write $A_y(j)=I_y(j,\mu_y)+I_{\bar y}(j,\mu_y)+j_y(H_T(j,\mu_y),T_{j,\mu_y})$ and
$\Lambda'=\iota_{\mu_0}\Lambda_{\mu_0}$.

\emph{Step 1: the covolume of $RR_0$.} At every finite place $y$, $RR_0$ agrees with
$\iota_{\mu_y}\Lambda_{\mu_y}$, so by Lemma~\ref{is-lattice-facts}(4) and (1),
$\operatorname{lcv}_{X_E}(RR_0)=\operatorname{lcv}_{X_E}(\Lambda')
+\sum_y\big(\ell_y\operatorname{idx}_y(\Lambda',RR_0)+\ell_{\bar y}\operatorname{idx}_{\bar y}(\Lambda',RR_0)\big)$,
where the indices are those of $\Lambda'$ and $\iota_{\mu_y}\Lambda_{\mu_y}$. The scalar conditions
of the rows say that $\iota_t$ multiplies the form of $X_Q(t)$ by $\kappa_v\prod_j\lambda_{j,t,p,v}$
on $\iota_t^{-1}(E_p)$, and $\theta_{j,t}$ multiplies the form of $X_T(j,t)$ by $\lambda_{j,t,p,v}$ on
$\theta_{j,t}^{-1}(W_p(j))$, when $t=\mu_0$ or $v=v_0$. The subspaces $\iota_t^{-1}(E_p)$ are
independent and span $K^D$, because $\iota_t$ is bijective and the $E_p=\iota_{\mu_0}Q_p(\mu_0)$ are
independent and span; the same holds for the $\theta_{j,t}^{-1}(W_p(j))$, and $\dim W_p(j)=\dim E_p$.
Computing both Gram determinants part by part therefore gives
\[
 r_v(X_Q(t),X_E,\iota_t)=\kappa_v^D\prod_jr_v(X_T(j,t),X_W(j),\theta_{j,t})
\]
whenever $t=\mu_0$ or $v=v_0$. Lemma~\ref{is-local-index-identity}, applied at each $y$ to the
tuples $\mu_0$ and $\mu_y$, with $\kappa=\kappa_{v_0}$, $\rho=1$ and $G_i$ the Gram determinants $G_{\mu_0}(y)$ and $G_{\mu_y}(y)$
(whose images at $v_0$ are the Gram determinants in $X_Q$, because $X_Q(t)$ is the quotient
metric and the involution is complex conjugation at $v_0$), gives
\[
 \operatorname{idx}_y(\Lambda',RR_0)+\operatorname{idx}_{\bar y}(\Lambda',RR_0)
 =v_y(G_{\mu_y}(y))-v_y(G_{\mu_0}(y))+\sum_j\big(A_y(j)-A^0_y(j)\big),
\]
where $A^0_y(j)=I_y(j,\mu_0)+I_{\bar y}(j,\mu_0)+j_y(H_T(j,\mu_0),T_{j,\mu_0})$.

By Lemma~\ref{is-lattice-facts}(5), $\operatorname{lcv}_{X_E}(\Lambda')
=\operatorname{lcv}_{X_Q(\mu_0)}(\Lambda_{\mu_0})+\frac1r\sum_v\log r_v(X_Q(\mu_0),X_E,\iota_{\mu_0})
=\operatorname{lcv}_{X_Q(\mu_0)}(\Lambda_{\mu_0})+D\,\overline{\log\kappa}
+\sum_j\frac1r\sum_v\log r_v(X_T(j,\mu_0),X_W(j),\theta_{j,\mu_0})$. For each $j$, by
Lemma~\ref{is-lattice-facts}(5), (4), (1) and (8) (the last applies because $H_T(j,\mu_0)$ is
rational and $X_T(j,\mu_0)$ is a metric for it at every place),
\[
 \frac1r\sum_v\log r_v(X_T(j,\mu_0),X_W(j),\theta_{j,\mu_0})
 =\operatorname{lcv}_{X_W(j)}(M_j)+2\sum_y\ell_yI_y(j,\mu_0)+\sum_y\ell_yj_y(H_T(j,\mu_0),T_{j,\mu_0}).
\]
Substituting, and using $\ell_{\bar y}=\ell_y$ to write $\sum_y\ell_yA^0_y(j)
=2\sum_y\ell_yI_y(j,\mu_0)+\sum_y\ell_yj_y(H_T(j,\mu_0),T_{j,\mu_0})$, the terms at $\mu_0$ cancel:
\[
 \operatorname{lcv}_{X_E}(RR_0)=\operatorname{lcv}_{X_Q(\mu_0)}(\Lambda_{\mu_0})-\sum_y\ell_yv_y(G_{\mu_0}(y))
 +D\,\overline{\log\kappa}+\sum_j\operatorname{lcv}_{X_W(j)}(M_j)
 +\sum_y\ell_y\Big(v_y(G_{\mu_y}(y))+\sum_jA_y(j)\Big).
\]
By Lemma~\ref{is-gram-covolume}, the first two terms on the right have a nonpositive sum.

\emph{Step 2: the local terms.} For $y\notin S_y$, $y$ is untwisted, so $A_y(j)=0$, and
$v_y(G_{\mu_y}(y))=0$; for $y\in S_y\setminus\mathrm{Tw}$, $A_y(j)=0$ and $F_0(y,j)=0$. So the last
sum runs over $S_y$, and by hypothesis (1) its $G$-part is at most $D\,L_G$. For $y\in\mathrm{Tw}$ the
surjectivity bound (Definition~\ref{is-surjectivity-bound}(1)) gives
$\sum_j\big(I_y(j,\mu_y)+I_{\bar y}(j,\mu_y)\big)\leq\big(\sum_j(2e_yF_0(y,j)+e_y\mathrm{err}_0(y,j))
+2e_y/p_y^2\big)D$, and hypothesis (2) bounds the relative indices. So
\[
 \sum_y\ell_y\sum_jA_y(j)\leq D\Big(2\sum_j\sum_{y\in\mathrm{Tw}}\ell_ye_yF_0(y,j)
 +\sum_j\sum_{y\in\mathrm{Tw}}\ell_ye_y\big(\mathrm{err}_0(y,j)+J_0(\mu_y,y,j)\big)
 +\sum_j\sum_{y\in\mathrm{Tw}}\frac{\ell_y\eta_y(j)}D+\sum_{y\in S_y}\frac{2\ell_ye_y}{p_y^2}\Big).
\]
The second sum is at most $\sum_jE_1(j)$ by Definition~\ref{is-surjectivity-bound}(3), the third at
most $\sum_jE_2(j)$ and the fourth at most $E_3$ by hypotheses (2) and (3).

\emph{Step 3.} By the covolume bound, $\operatorname{lcv}_{X_W(j)}(M_j)/D\leq\ell_3(j)$. Dividing
the identity of Step 1 by $D$ and inserting Step 2 gives the bound for
$\operatorname{lcv}_{X_E}(RR_0)/D$, which is $\mu(CC_0)$ by Lemma~\ref{is-lattice-facts}(6).
\end{proof}

\begin{lemma}\label{is-column-height}
Let $K$ be a totally complex number field. Let $H\in M_n(K)$, let $X^0$ be a metric on $K^n$ with
$(X^0_v{}^{-1})^*X^0_v{}^{-1}=\sigma_v(H)$ at every $v$, let $P\in M_{d\times n}(K)$ be surjective
with quotient metric $X_Q$ (Lemma~\ref{is-gram-covolume}), let $X_E$ be an invertible metric on
$K^D$ and $\iota\in M_{D\times d}(K)$ bijective. Let $Q_p\subset K^d$, for $p$ in a finite nonempty
set $\Pi$, be subspaces that span $K^d$, and let $\kappa_v>0$ and $\lambda_{j,p,v}>0$, for lifts $j$
in a finite set $J$, be such that $\iota$ multiplies the Hermitian form of $X_Q$ at $v$ by
$s_{p,v}=\kappa_v\prod_j\lambda_{j,p,v}$ on $Q_p$. Let $c_0,\dots,c_{m'}$ be integers, $X_i$ invertible
metrics on $K^{c_i}$, $K_i\in M_{c_i\times c_{i-1}}(K)$ for $1\leq i\leq m'$,
$\mathrm{Pt}\colon K^{c_{m'}}\to\operatorname{Hom}_K(K^n,K)$ linear, and
$B\colon K^{c_0}\to\operatorname{Hom}_K(K^D,K)$ a nonzero linear map with
$B(\varphi)(\iota PF)=\mathrm{Pt}(K_{m'}\cdots K_1\varphi)(F)$ for all $\varphi$ and $F$. Then at
every infinite place $v$, with operator norms from $X_0$ to the dual metric of $X_E$, from $X_{m'}$
to the dual metric of $X^0$, and from $X_{i-1}$ to $X_i$,
\[
 \log\|B\|_v^2\leq\log\|\mathrm{Pt}\|_v^2+\sum_{i=1}^{m'}\log\|K_i\|_v^2-\log\kappa_v
 -\min_{p\in\Pi}\sum_j\log\lambda_{j,p,v}.
\]
\end{lemma}
\leanref{Formalism/LocalToGlobal/Divisibility/DivisionHeightBoundStatement.lean}{Formalism.L2.DIV.complex_height_cast_at}
\leanref{Formalism/QuotientMetricDivisors/FinalW.lean}{Formalism.DivisionValues.height_constant_bound}
\begin{proof}
Fix $v$ and write all vectors after $\sigma_v$. Let $s=\min_ps_{p,v}$. The form
$\langle\iota x,\iota y\rangle_{X_E}$ equals $s_{p,v}\langle x,y\rangle_{X_Q}$ for $x\in Q_p$, so
the positive self-adjoint operator $\iota^*\iota$ (adjoint for $X_Q$ and $X_E$) acts on $Q_p$ as
$s_{p,v}$. Since the $Q_p$ span, its eigenvalues are among the $s_{p,v}$, so
$|\iota x|_{X_E}^2\geq s|x|_{X_Q}^2$ for all $x$.

Let $y\in K^D_v$ and put $x=\iota^{-1}y$. Since $X_Q$ is the quotient metric of $X^0$ through $P$,
there is $F$ with $PF=x$ and $|F|_{X^0}=|x|_{X_Q}$. Then
\[
 |B(\varphi)(y)|=|\mathrm{Pt}(K_{m'}\cdots K_1\varphi)(F)|
 \leq\|\mathrm{Pt}\|_v\prod_i\|K_i\|_v\,|\varphi|_{X_0}\,|F|_{X^0}
 \leq\|\mathrm{Pt}\|_v\prod_i\|K_i\|_v\,|\varphi|_{X_0}\,s^{-1/2}|y|_{X_E}.
\]
So $\|B\|_v^2\leq\|\mathrm{Pt}\|_v^2\prod_i\|K_i\|_v^2/s$, and
$\log s=\log\kappa_v+\min_p\sum_j\log\lambda_{j,p,v}$. Since $B\neq0$, the norm $\|B\|_v$ is
positive and the logarithm is defined.
\end{proof}

\begin{lemma}\label{is-shift-into-dual}
Let a division package over $K$ be given (Definition~\ref{is-division-package}), with columns of
rank $D$, glued lattice $RR=RR_1$ and $CC=RR^\vee$. Let $u$ be a word, $N\subset K^{a(u)}$ a lattice,
and $s(w)$ real numbers, one for each finite place $w$, with $\lfloor s(w)\rfloor=0$ for all but
finitely many $w$. Put $N'=\big(\prod_w\mathfrak p_w^{-\lfloor s(w)\rfloor}\big)N$. Suppose that for
every finite place $w$, every $\varphi\in N$ and every $F\in M_{\mu_w}$, the period
$\mathrm{Pt}_u(\mu_w)(\varphi,F)$ is zero or has $v_w\geq s(w)$. Then $B_u(N')\subset CC$, and
\[
 \mu(N')=\mu(N)+2\sum_w\ell_w\lfloor s(w)\rfloor,
\]
the slopes taken in the metric $X_c(u)$.
\end{lemma}
\leanref{Formalism/LocalToGlobal/IndefiniteForms/PairingConditionOfInverseMatrixIdentity.lean}{Formalism.L2.Indef.IndefColsR.intoDual}
\leanref{Formalism/QuotientMetricDivisors/DivIntoDualShift.lean}{Formalism.DivisionValues.DivIntoDualShift}
\leanref{Formalism/LocalToGlobal/Divisibility/Inputs.lean}{Formalism.L2.DIV.shiftLat}
\begin{proof}
Let $\varphi'\in N'$ and $y\in RR$; we show $B_u(\varphi')(y)\in\mathcal O_K$, that is,
$v_w(B_u(\varphi')(y))\geq0$ at every finite place $w$. Write $\varphi'=\sum_ka_k\varphi_k$ with
$\varphi_k\in N$ and $v_w(a_k)\geq-\lfloor s(w)\rfloor$. Locally at $w$, $RR$ agrees with
$\iota_{\mu_w}P_{\mu_w}M_{\mu_w}$, so $y=\sum_lb_l\iota_{\mu_w}P_{\mu_w}F_l$ with $F_l\in M_{\mu_w}$ and
$v_w(b_l)\geq0$. By the column data,
$B_u(\varphi_k)(\iota_{\mu_w}P_{\mu_w}F_l)=\zeta_u(\mu_w)^{-1}\mathrm{Pt}_u(\mu_w)(\varphi_k,F_l)$, and
$v_w(\zeta_u(\mu_w))=0$, so this value is zero or has valuation at least $s(w)$, hence at least
$\lceil s(w)\rceil\geq\lfloor s(w)\rfloor$, valuations being integers. So every term of
$B_u(\varphi')(y)=\sum_{k,l}a_kb_lB_u(\varphi_k)(\iota_{\mu_w}P_{\mu_w}F_l)$ has $v_w\geq0$.

Multiplying a lattice of rank $a(u)$ by the fractional ideal $\prod_w\mathfrak p_w^{-n_w}$ divides its
covolume by $\prod_wNw^{n_wa(u)}$, so $\operatorname{lcv}(N')=\operatorname{lcv}(N)-2a(u)\sum_w\ell_wn_w$
and $\mu(N')=\mu(N)+2\sum_w\ell_wn_w$ with $n_w=\lfloor s(w)\rfloor$.
\end{proof}

\begin{definition}\label{is-record-numbers}
\emph{Record numbers} consist of finite sets $J$ (the lifts) and $\mathcal P$ (the places), a subset
$T\subset\mathcal P$, an integer $r\geq0$ and real numbers
\[
 \log T_u,\ a_t\ (t\in T),\ \nu_{j,\rho,i},\ \log H_{j,\rho},\ x,\ \varepsilon,\ L,\ b,\ a,\ r_E,\
 L_a,\ B_0,\ N_c,\ m,\ f,\ k,\ N,\ S,\ t_0,\ \eta,\ w'',\ b',\ \lambda_0,\ C_0,\ c_e,\ c,\ n_V,\
 L_p,\ \log\Omega,\ q_j,\ w_{j,\rho},\ \epsilon_{j,\rho},\ \delta_{j,\rho}
\]
for $j\in J$, $\rho\in\mathcal P$ and $1\leq i\leq r$, such that, with $\phi(y)=y\log y-y$ and
$|\nu_{j,\rho}|=\sum_i\nu_{j,\rho,i}$:
\begin{enumerate}
\item $\sum_i\phi(\nu_{j,\rho,i})-\log H_{j,\rho}\leq[\rho\in T]\phi(a_\rho)+\epsilon_{j,\rho}$ and
  $|\nu_{j,\rho}|=[\rho\in T]a_\rho+k\,w_{j,\rho}+\delta_{j,\rho}$, where $[\rho\in T]$ is $1$ or $0$
  and the term is read as $0$ when $\rho\notin T$;
\item $x>1$, $\varepsilon>0$, $L=x^{\varepsilon}$, $b=e^{L}$, $a=xb$, $1\leq r_E$, $\log r_E\leq L$,
  $\log x\leq L$, $L_a=\log(ar_E)$, $B_0=fmbx^{\varepsilon}$ and $N_c\geq0$;
\item $|J|=m$, $|\mathcal P|\leq f$, $m,f,k\geq1$, $k+m\leq N$, and $S,t_0,\eta,b',\lambda_0,L_p\geq0$;
\item $0\leq w_{j,\rho}\leq2b$, $|q_j-k|\leq m$ and $|\log\Omega|\leq L_a$;
\item $m\lambda_0\leq b$, $N^2\leq b$, $m^3N^2L\leq b$, $m\eta NL\leq b$, $kN\leq b$, $kS\leq w''$,
  $m(t_0+1)\eta\leq w''$, $Nw''\leq b$, $b'N\leq a$, $mkb'\leq a$ and $0\leq n_V\leq N$;
\item $1\leq C_0\leq c_e$, $0\leq c\leq C_0$ and $cL_p\leq C_0L_a$;
\item $\epsilon_{j,\rho}\leq C_0\big(w_{j,\rho}+S+N^2+m(t_0+1)\eta\big)L_a$ and
  $|\delta_{j,\rho}|\leq C_0\big(w_{j,\rho}+k(N+S)+m(t_0+1)\eta+[\rho\in T]kb'\big)$.
\end{enumerate}
Their \emph{height}, \emph{right side} and \emph{error budget} are
\[
 h^{\mathrm{rec}}=h_A+\sum_{j\in J}\big(h_B(j)-h_X(j)\big),\qquad
 \mathrm{rhs}=\sum_{t\in T}\Big(m\big(\phi(a_t)-a_t\log\Omega\big)-a_t\log T_u\Big)
 +(50C_0+10)|T|aL_a+60c_emx^{\varepsilon}(B_0+N_c),
\]
and $\mathrm{err}_B=c_emx^{\varepsilon}(B_0+N_c)$, where
\begin{align*}
 h_A&=-\sum_{t\in T}a_t\log T_u+c\big(|T|b'N+f(m\eta+S)N+fN^2\big)L_a,\\
 h_B(j)&=\sum_{\rho\in\mathcal P}\big(\textstyle\sum_i\phi(\nu_{j,\rho,i})-|\nu_{j,\rho}|\log\Omega\big)+cfN^2(L_a+m^3L),\\
 h_X(j)&=\sum_{\rho\in\mathcal P}\log H_{j,\rho}-cfn_V^2L_p.
\end{align*}
The record of a word (Definition~\ref{is-height-numbers}) carries record numbers with
$\mathcal P=\mathcal T$, $T$ the top classes, $N=N(m)$, the weights of the lifts and the constant
$c$; Definition~\ref{is-height-numbers} uses $\log T_u$, $a_\tau$, $L=x^\varepsilon$, $L_a$ and
$\log\Omega$ from them.
\end{definition}
\leanref{Formalism/LocalToGlobal/Divisibility/Inputs.lean}{Formalism.L2.DIV.RecNums}
\leanref{Formalism/LocalToGlobal/Divisibility/Inputs.lean}{Formalism.L2.DIV.RecNums.hu}
\leanref{Formalism/LocalToGlobal/Divisibility/Inputs.lean}{Formalism.L2.DIV.RecNums.rhs}
\leanref{Formalism/LocalToGlobal/Divisibility/Inputs.lean}{Formalism.L2.DIV.RecNums.errB}

\begin{lemma}\label{is-record-height}
Let record numbers be given (Definition~\ref{is-record-numbers}), and height numbers for them with
period exponent $h_P$ and sizes $\mathrm{sz}(j,v)$ at the $r_\infty$ infinite places $v$ of a number
field (Definition~\ref{is-height-numbers}(3)). If $m_\lambda$ is a real number with
$\sum_jh_X(j)\leq m_\lambda$, then
\[
 h_P+\sum_j\frac1{r_\infty}\sum_v\mathrm{sz}(j,v)-m_\lambda\leq h^{\mathrm{rec}}.
\]
\end{lemma}
\leanref{Formalism/LocalToGlobal/Divisibility/DivisionHeightBoundStatement.lean}{Formalism.L2.DIV.HtNums.le_hu}
\leanref{Formalism/DivisorInequality/DivisorInequalityProofs/EvenDegreeWeightBounds.lean}{Formalism.DivisorInequality.height_bound_at_set}
\leanref{Formalism/DivisorInequality/DivisorInequalityProofs/EvenDegreeWeightBounds.lean}{Formalism.DivisorInequality.UniformCovolumeTail.stepDiSharp}
\begin{proof}
Write $C_2,C_e,C_\theta,C_4,c_5,c_6,C_8,C_T,C_{ev}$ for the constants of the height numbers and
$\ell_g$ for their number of that name. From Definition~\ref{is-record-numbers}(2),
$L_a=\log x+L+\log r_E$, so $1\leq L\leq L_a$, and $\log(2a)\leq2L_a$. Also $m\leq N$, so the last-side
condition (2f) gives $\log T_u\leq C_TmL\leq C_TNL_a$.

\emph{Kernel side: $\frac1{r_\infty}\sum_v\mathrm{sz}(j,v)\leq h_B(j)$ for every lift $j$.} By (1a) the
mean over $v$ of $\sum_\tau\big(\sum_i\phi(\nu_{j,\tau,i})-|\nu_{j,\tau}|\log\Omega_\tau(v)\big)$ is the
first sum of $h_B(j)$. For the remaining terms of $\mathrm{sz}(j,v)$ use $|\mathcal T|\leq f$ and
$0\leq g_0(j)\leq N^2$ (1e): the mean of $fg_0(j)\log\theta_0(v)$ is at most $fN^2C_\theta L$ by (1c);
$fg_0(j)\log(8+32fg_0(j)v_0)\leq fN^2c_5L_a$ by (1f); $2M(j,v)\leq2C_2fm^2N^2(C_8+C_Tm)L+2C_efN^2C_8L$
by (1b), (1g) and (2f); $\sum_\tau\big(1+2g_0(j)\log(2|\nu_{j,\tau}|+2)\big)\leq f(1+2N^2c_6L_a)\leq
fN^2(1+2c_6)L_a$ by (1f); and $\log h_j+\log h'_j\leq2C_4fN^2(1+C_Tm+C_8)L$ by (1d). With $L\leq L_a$ and
$1\leq m\leq m^2\leq m^3$, the sum of these is at most
\[
 fN^2\big((C_\theta+c_5+2c_6+1+2C_4+2C_eC_8)L_a+2(C_2+C_4)(C_T+C_8)m^3L\big)\leq cfN^2(L_a+m^3L),
\]
by the second and third conditions on $c$ in Definition~\ref{is-height-numbers}(3).

\emph{Last side: $h_P\leq h_A$.} By (2a), $h_P\leq-\sum_\tau n_\tau\log T_u+\sum_\tau(2\log c_\tau+\log
\xi_\tau)+e_v$. Fix $\tau$. By (2b), $\log\xi_\tau\leq\frac{N(N-1)}2\log(2a)\leq N^2L_a$ and $2\log
c_\tau\leq2(N-1)D_1(\tau)\log(2a)\leq4NL_aD_1(\tau)$. By (2e), $D_1(\tau)\leq4S+4m\eta$ if $\tau\notin T$
and $D_1(\tau)\leq4S+4m\eta+6b'$ if $\tau\in T$. By (2c) and (2d), if $\tau\in T$ then
$-n_\tau\log T_u=-a_\tau\log T_u+\big(\frac{N-k}2+D_0(\tau)\big)\log T_u$ with
$\frac{N-k}2+D_0(\tau)\leq\frac{N-1}2+S+m\eta+2b'$, and if $\tau\notin T$ then
$-n_\tau\leq\frac{N-1}2+m\eta$. With $\log T_u\leq C_TNL_a$ this gives
\[
 -n_\tau\log T_u+2\log c_\tau+\log\xi_\tau\leq\Big(\frac{N-1}2+S+m\eta\Big)C_TNL_a+16NL_a(S+m\eta)+N^2L_a
 +[\tau\in T]\big((2C_T+24)b'NL_a-a_\tau\log T_u\big).
\]
Summing over the at most $f$ places $\tau$ and adding $e_v\leq C_{ev}fN^2L_a$ gives
$h_P\leq-\sum_{t\in T}a_t\log T_u+(\frac{C_T}2+1+C_{ev})fN^2L_a+(C_T+16)f(m\eta+S)NL_a
+(2C_T+24)|T|b'NL_a$, which is at most $h_A$ because $2(16+C_T+C_{ev})\leq c$.

Adding the two bounds and $-m_\lambda\leq-\sum_jh_X(j)$ gives the claim.
\end{proof}

\begin{lemma}\label{is-record-division}
Let record numbers be given (Definition~\ref{is-record-numbers}), and let $\mu_C$, $\mu_0$, $F_o$,
$L_{\mathrm{pass}}$ and $E_b$ be real numbers with
\[
 \mu_C\leq\mu_0+L_{\mathrm{pass}},\qquad
 \mu_0\leq\sum_j\sum_\rho q_jw_{j,\rho}\log\Omega+2F_o+E_b,\qquad
 L_{\mathrm{pass}}\leq1+C_0m\lambda_0,\qquad E_b\leq\mathrm{err}_B.
\]
Then $\mu_C+h^{\mathrm{rec}}-2F_o\leq\mathrm{rhs}$.
\end{lemma}
\leanref{Formalism/LocalToGlobal/Divisibility/DivisionHeightBoundStatement.lean}{Formalism.L2.DIV.RecNums.divIneq}
\leanref{Formalism/QuotientMetricDivisors/HermitianTwistDivisors/DivAssembly.lean}{Formalism.DivisionValues.divisor_inequality_assembly}
\leanref{Formalism/QuotientMetricDivisors/HermitianTwistDivisors/AbsorbE.lean}{Formalism.DivisionValues.error_absorption}
\begin{proof}
\emph{Main terms.} Fix $j$ and $\rho$. By Definition~\ref{is-record-numbers}(1),
$q_jw_{j,\rho}\log\Omega-|\nu_{j,\rho}|\log\Omega=(q_j-k)w_{j,\rho}\log\Omega-[\rho\in T]a_\rho\log\Omega
-\delta_{j,\rho}\log\Omega$, so
\[
 q_jw_{j,\rho}\log\Omega+\sum_i\phi(\nu_{j,\rho,i})-\log H_{j,\rho}-|\nu_{j,\rho}|\log\Omega
 \leq[\rho\in T]\big(\phi(a_\rho)-a_\rho\log\Omega\big)+\epsilon_{j,\rho}
 +|q_j-k||w_{j,\rho}||\log\Omega|+|\delta_{j,\rho}||\log\Omega|.
\]
Sum over $\rho$ and over the $m$ lifts $j$, and add the constant parts of $h_A$, $h_B(j)$ and
$h_X(j)$. With the hypotheses on $\mu_C$ and $\mu_0$ this gives
$\mu_C+h^{\mathrm{rec}}-2F_o\leq\sum_{t\in T}\big(m(\phi(a_t)-a_t\log\Omega)-a_t\log T_u\big)+E$,
where $E$ is the sum of the eight error terms
\[
 L_{\mathrm{pass}},\quad E_b,\quad \sum_je_K,\quad e_6,\quad e_P,\quad \sum_{j,\rho}\epsilon_{j,\rho},\quad
 \sum_{j,\rho}|q_j-k||w_{j,\rho}||\log\Omega|,\quad \sum_{j,\rho}|\delta_{j,\rho}||\log\Omega|,
\]
with $e_K=cfN^2(L_a+m^3L)$, $e_6=cmfn_V^2L_p$ and $e_P=c(|T|b'N+f(m\eta+S)N+fN^2)L_a$.

\emph{Absorption.} Put $Y=\frac{c_e}{C_0}mx^{\varepsilon}(B_0+N_c)$. From
Definition~\ref{is-record-numbers}(2), $L\geq1$, $b\geq1$, $L_a\leq3L$ and $B_0=fmbL$; since
$m\geq1$ and $c_e/C_0\geq1$, $b\leq fmbL\leq mB_0\leq Y$. Now:
$L_{\mathrm{pass}}\leq1+C_0b\leq(1+C_0)Y$; $E_b\leq\mathrm{err}_B=C_0Y$; $e_K\leq C_0(3fbL+fb)$ by
$N^2\leq b$ and $m^3N^2L\leq b$, so $\sum_je_K\leq5C_0Y$; $e_6\leq C_0mfN^2L_a\leq3C_0Y$ by $cn_V^2L_p\leq
C_0N^2L_a$; $e_P\leq C_0|T|aL_a+9C_0Y$, by $b'N\leq a$, $m\eta NL_a\leq3b$, $SN\leq w''N\leq b$ and
$N^2\leq b$; each $\epsilon_{j,\rho}\leq15C_0bL$, because $w_{j,\rho}+S+N^2+m(t_0+1)\eta\leq5b$, so their
sum is at most $15C_0Y$; each $|q_j-k||w_{j,\rho}||\log\Omega|\leq6mbL$, so their sum is at most
$6Y$; and $|\delta_{j,\rho}||\log\Omega|\leq15C_0bL+[\rho\in T]C_0kb'L_a$, so their sum is at most
$15C_0Y+C_0|T|aL_a$ by $mkb'\leq a$. Hence
\[
 E\leq2C_0|T|aL_a+(49C_0+7)Y\leq(50C_0+10)|T|aL_a+(50C_0+10)\frac{c_e}{C_0}mx^{\varepsilon}(B_0+N_c),
\]
and $(50C_0+10)/C_0\leq60$ because $C_0\geq1$. So $E$ is at most the error part of $\mathrm{rhs}$.
\end{proof}

\begin{definition}\label{is-division-inputs}
Let a division package over $K$ for numerical bounds $N$ be given
(Definition~\ref{is-division-package}), with rows $S_0$ and columns $S_1$, let $u$ be a word, let
$N_u\subset K^{a(u)}$ be an $\mathcal O_K$-submodule, let $s(w)$ be real numbers indexed by the finite
places, and let record numbers be given (Definition~\ref{is-record-numbers}) with the lifts $J$ of
the package. \emph{Division inputs} for these consist of the following.
\begin{enumerate}
\item $N_u$ is a lattice, $\lfloor s(w)\rfloor=0$ for all but finitely many $w$, $X_c(u)$ is
  invertible at every place, and $a(u)>0$.
\item A metric $X^0$ for the Hermitian matrix $H_V(\mu_0)$ of the columns at every infinite place.
\item A chain of lifts for $u$: integers $c_0=a(u),c_1,\ldots,c_{m'}$, invertible metrics $X_i$ on
  $K^{c_i}$ with $X_0=X_c(u)$, matrices $K_i\in M_{c_i\times c_{i-1}}(K)$, a map $o$ from $J$ to the
  steps $\{1,\ldots,m'\}$ and a linear map $\mathrm{Pt}$ from $K^{c_{m'}}$ to the dual of $K^{n_V(\mu_0)}$,
  with $\mathrm{Pt}_u(\mu_0)(\varphi,F)=\mathrm{Pt}(K_{m'}\cdots K_1\varphi)(F)$.
\item Height numbers for the record numbers (Definition~\ref{is-height-numbers}(3)) over the
  infinite places of $K$, with period exponent $h_P$ and sizes $\mathrm{sz}(j,v)$, such that
  $\|\mathrm{Pt}\|_v^2\leq e^{h_P}$ from $X_{m'}$ to the dual metric of $X^0$, and Cauchy-Binet data
  for the chain (Definition~\ref{is-cauchy-binet}) whose square factor at the step $i$ and $v$ has
  logarithm at most $\sum_{o(j)=i}\mathrm{sz}(j,v)$.
\item $\sum_jh_X(j)\leq\sum_j\ell_6(j)$.
\item For the rows, local bases, a finite set $S_y$ and numbers $L_G,E_2(j),E_3,g_y,\eta_y(j)$ that
  satisfy the hypotheses of Lemma~\ref{is-glued-covolume}.
\item A real number $E_{\ell_3}$ with $\sum_j\ell_3(j)\leq\sum_j\sum_\rho q_jw_{j,\rho}\log\Omega+E_{\ell_3}$
  and $E_{\ell_3}+L_G+\sum_j(E_1(j)+E_2(j))+E_3\leq\mathrm{err}_B$.
\item A real number $Q\geq1$ with $0\leq\delta_e$, $3\delta_eQ\leq1$ and
  $B_{\mathrm{rows}}+h_N+28\ell_a\leq Q$, where $B_{\mathrm{rows}}=\sum_j\ell_3(j)+2F_o+L_G
  +\sum_j(E_1(j)+E_2(j))+E_3$ and $F_o=\sum_j\sum_{y\in\mathrm{Tw}}\ell_ye_yF_0(y,j)$; and
  $\ell_a\leq\lambda_0$, $28\leq C_0$.
\item (Valuations.) For every finite place $w$, every $\varphi\in N_u$ and every $F\in M_{\mu_w}$ of the
  columns, $\mathrm{Pt}_u(\mu_w)(\varphi,F)$ is zero or has $v_w\geq s(w)$.
\end{enumerate}
\end{definition}
\leanref{Formalism/LocalToGlobal/Divisibility/ConstantGeTwentyEight.lean}{Formalism.L2.DIV.DivInputs0R}
\leanref{Formalism/LocalToGlobal/Divisibility/ConstantGeTwentyEight.lean}{Formalism.L2.DIV.DivInputsR}
\leanref{Formalism/LocalToGlobal/Divisibility/ConstantGeTwentyEight.lean}{Formalism.L2.DIV.FoOfR}

\begin{theorem}\label{is-division-inequality}
Let a division package over $K$ be given (Definition~\ref{is-division-package}), with columns of
rank $D$, metric $X_E$, glued lattice $RR$, $CC=RR^\vee$ and scalars $\kappa_v$, and let $u$, $N_u$,
$s$ and record numbers carry division inputs (Definition~\ref{is-division-inputs}), with $F_o$ as
there. Put $N'_u=\big(\prod_w\mathfrak p_w^{-\lfloor s(w)\rfloor}\big)N_u$. Then:
\begin{enumerate}
\item $B_u(N'_u)\subset CC$;
\item there are $c_u(v)>0$, one for each infinite place $v$, with $\|B_u\|_v^2\leq c_u(v)$ from
  $X_c(u)$ to the dual metric of $X_E$ and $\overline{\log c_u}=h^{\mathrm{rec}}-\overline{\log\kappa}$;
  in particular $h(B_u)\leq h^{\mathrm{rec}}-\overline{\log\kappa}$ if $B_u\neq0$;
\item $\mu(CC)-\overline{\log\kappa}+h^{\mathrm{rec}}-2F_o\leq\mathrm{rhs}$;
\item $\mu(N'_u)=\mu(N_u)+2\sum_w\ell_w\lfloor s(w)\rfloor$.
\end{enumerate}
\end{theorem}
\leanref{Formalism/LocalToGlobal/Divisibility/ConstantGeTwentyEight.lean}{Formalism.L2.DIV.divClausesR_of}
\leanref{Formalism/LocalToGlobal/Divisibility/ConstantGeTwentyEight.lean}{Formalism.L2.DIV.DivClausesR}
\leanref{Formalism/LocalToGlobal/Divisibility/ConstantGeTwentyEight.lean}{Formalism.L2.DIV.height_chain_atR}
\leanref{Formalism/LocalToGlobal/Divisibility/DivisionHeightBoundStatement.lean}{Formalism.L2.DIV.exists_cu}
\begin{proof}
(1) and (4) are Lemma~\ref{is-shift-into-dual}, whose hypothesis is input (9).

(2) Suppose first $B_u=0$. Take $c_u(v)=\exp(h^{\mathrm{rec}}-\overline{\log\kappa})$ for every $v$.

Suppose $B_u\neq0$. By Lemma~\ref{is-cauchy-binet-bound} and input (4),
$\log\|K_i\|_v^2\leq\sum_{o(j)=i}\mathrm{sz}(j,v)$ for every step $i$. The columns are bijective with
scalars $\kappa_v\prod_j\lambda_{j,\mu_0,p,v}$ on the parts $Q_p(\mu_0)$, which span; their Hermitian
matrix $H_V(\mu_0)$ has the metric $X^0$ of input (2); and by the column data and input (3),
$B_u(\varphi)(\iota_{\mu_0}P_{\mu_0}F)=\mathrm{Pt}_u(\mu_0)(\varphi,F)=\mathrm{Pt}(K_{m'}\cdots
K_1\varphi)(F)$, since $\zeta_u(\mu_0)=1$. So Lemma~\ref{is-column-height} gives, at every $v$,
\[
 \log\|B_u\|_v^2\leq a(v):=h_P+\sum_j\mathrm{sz}(j,v)-\log\kappa_v-\min_p\sum_j\log\lambda_{j,\mu_0,p,v}.
\]
By the lower norm bound, $\log\lambda_{j,\mu_0,p,v}\geq\ell_6(j)$, so the mean of the $a(v)$ is at most
$h_P+\sum_j\frac1r\sum_v\mathrm{sz}(j,v)-\sum_j\ell_6(j)-\overline{\log\kappa}$, which is at most
$h^{\mathrm{rec}}-\overline{\log\kappa}$ by Lemma~\ref{is-record-height} with $m_\lambda=\sum_j\ell_6(j)$
and input (5). Put $c_u(v)=\exp\big(a(v)+h^{\mathrm{rec}}-\overline{\log\kappa}-\overline{a}\big)$, where
$\overline a$ is the mean of the $a(v)$. Then $c_u(v)\geq e^{a(v)}\geq\|B_u\|_v^2$ and
$\overline{\log c_u}=h^{\mathrm{rec}}-\overline{\log\kappa}$.

(3) Let $\mu_0=\mu(CC_0)-\overline{\log\kappa}$ for the rows. Lemma~\ref{is-glued-covolume}, with the
data of input (6), gives $\mu_0\leq B_{\mathrm{rows}}$. By input (8), $0\leq\delta_e$ and
$\delta_e\leq\frac1{3Q}\leq\frac13$. The functional bound of the rows says
$\mu_{\max}(RR_0)\leq h_N-\overline{\log\kappa}$, so the passage, with $H=h_N-\overline{\log\kappa}$,
and the equality of the scalars $\kappa_v$ of rows and columns give
\[
 \mu(CC)-\overline{\log\kappa}\leq\mu_0+L_{\mathrm{pass}},\qquad
 L_{\mathrm{pass}}=3\delta_e(\mu_0+h_N+28\ell_a)+28\ell_a.
\]
By input (8), $3\delta_e(\mu_0+h_N+28\ell_a)\leq3\delta_eQ\leq1$, and $28\ell_a\leq28\lambda_0\leq
C_0m\lambda_0$, because $28\leq C_0$ and $m\geq1$. So $L_{\mathrm{pass}}\leq1+C_0m\lambda_0$. By input
(7), $\mu_0\leq\sum_j\sum_\rho q_jw_{j,\rho}\log\Omega+2F_o+E_b$ with
$E_b=E_{\ell_3}+L_G+\sum_j(E_1(j)+E_2(j))+E_3\leq\mathrm{err}_B$. Lemma~\ref{is-record-division}, with
$\mu_C=\mu(CC)-\overline{\log\kappa}$, gives (3).
\end{proof}

\section{The division link}\label{section-is-division}

\begin{definition}\label{is-division-link}
Let $D$ be a definite-side datum (Definition~\ref{th-definite-datum}) with ranks $k,m,n$, $f=2n$,
coefficient field $L$, words $u$, auxiliary primes $T_u$ and division data
(Definition~\ref{th-definite-datum}(9)): Hermitian lattices $N_u$ over $L$
(Definition~\ref{hs-hermitian-lattice}) with real numbers $h_u$, the \emph{heights}, the glued
lattice $RR$ with dual lattice $CC=RR^{\vee}=\{\phi:\phi(RR)\subset\mathcal O_L\}$, and $L$-linear
maps $B_u$ from $LN_u$ to the dual of $L\,RR$. Let size data at scale $x$ with exponent
$\varepsilon$ be given (Definition~\ref{pd-size-data}), with top classes $T$ and exponents $a_t$, level primes $g_0,g_1$,
support, losses $F_j(w)$ and $E_{u,j}(w)$, and gains $\mathrm{gain}_j(w)$. Write
$\mathrm{sz}_w(y)=2\ell_we_wy$,
\[
 F_o=\sum_{j=1}^{m}\sum_{w\in\operatorname{supp}}\mathrm{sz}_w(F_j(w)),\qquad
 s=\frac{\sum_u\operatorname{rk}(N_u)\,(\mu(N_u)-h_u)}{\sum_u\operatorname{rk}(N_u)},
\]
let $h(B_u)$ be the height of $B_u$ (Definition~\ref{hs-hermitian-lattice}) for the metric of
$N_u$ and the dual metric of $RR$, let $\mu$ and $\mu_{\max}$ be the slopes of
Definition~\ref{hs-slopes}, and put
\[
 \mathrm{MAIN}(u)=\sum_{t\in T}\big(m(\phi(a_t)-a_t\log\Omega)-a_t\log T_u\big),\qquad
 \phi(y)=y\log y-y,
\]
with $\log\Omega=\gamma+\log(2\pi/r_E)-\operatorname{Re}\frac{L'}{L}(1,\chi_d)$, where $r_E$ is the
root discriminant of the host field and $\chi_d$ the quadratic character modulo its auxiliary prime
$d$, and $L_a=\log(a\,r_E)$, with $B_0$ and $N_c$ as in Definition~\ref{is-twisted-division-bound}.
Let $C,C_{abs},C_H$ be real. The \emph{division link} holds at $(D,C,C_{abs},C_H)$ if:
\begin{enumerate}
\item $B_u(N_u)\subset CC$ for every word $u$;
\item $s+\mu_{\max}(RR)\le x^{C_H}b$;
\item for every word $u$, every lift $j$ and every place $w$ whose residue characteristic exceeds
  $X''$, differs from $T_u$, $g_0$, $g_1$ and $d$, and does not divide the stabilizer denominator of
  $u$, we have $\mathrm{gain}_j(w)=0$ and $E_{u,j}(w)=0$;
\item the Bessel identity holds for $D$ (Definition~\ref{th-bessel-identity});
\item $RR$ is finitely generated and spans its $L$-vector space;
\item $h(B_u)\le h_u$ for every word $u$ with $B_u\ne0$;
\item for every word $u$,
\[
 \mu(CC)+h_u-F_o\le\mathrm{MAIN}(u)+C|T|\,a\,L_a+C\,m\,(B_0+N_c)\,x^{C_{abs}\varepsilon};
\]
\item for every word $u$ and every lift $1\le j\le m$,
\[
 \Big|\sum_{w\in\operatorname{supp}}\mathrm{sz}_w(E_{u,j}(w))-\Big(\sum_{t\in T}a_t\Big)\log(g_0g_1)\Big|\le CB_0;
\]
\item for all words $u,u'$,
\[
 \Big|\Big(h_u+\sum_{t\in T}a_t\log T_u\Big)-\Big(h_{u'}+\sum_{t\in T}a_t\log T_{u'}\Big)\Big|
 \le C\big(|T|\,a\,L_a+mB_0\big).
\]
\end{enumerate}
A family $\mathcal F$ satisfies the division link with constants $C_{abs},C_H$ if there are $C$ and
$x_0$ such that the division link holds at $(D_{x,i},C,C_{abs},C_H)$ for every scale $x\ge x_0$ and
every member $i$. It satisfies the \emph{conditional division link} if it satisfies the division
link whenever it satisfies the division bound at twisted places
(Definition~\ref{is-twisted-division-bound}) and the support condition: eventually, at every
member, every lift $j$ and every place $w$ with $\mu_w\ne\mu_0$, the local spaces at $w$ have total
rank equal to the rank of the top space of $\mu_w$ at level $j$.
\end{definition}
\leanref{Formalism/SharpRestriction/GlobalTheoremInputs.lean}{Formalism.MainNotions.LinkDIV_at}
\leanref{Formalism/SharpRestriction/GlobalTheoremInputs.lean}{Formalism.MainNotions.LinkDIV}
\leanref{Formalism/SharpRestriction/GlobalTheoremInputs.lean}{Formalism.MainNotions.LinkDIVt}
\leanref{Formalism/SharpRestriction/LevelLetterCocycles/LocsCoverTw.lean}{Formalism.MainNotions.N.LocsCoverTw}

\begin{definition}\label{is-division-kernel-bounds}
Let $\mathcal F$ be a family (Definition~\ref{pd-family}). \emph{Division kernel data} at a member
consist of the tied kernel data of the member, a row-side family whose rows satisfy the conditions
of the rows at the base tuple, a column record with its fields, characters and scale data, and the
compatibilities that identify the tuples, words, levels and metrics of the column record with those of
the row-side family and of the kernel data. Suppose that $\mathcal F$ has division kernel data at
every member at all large scales, and let $C_{abs},C_H$ be real. The family satisfies the
\emph{division kernel bounds} if:
\begin{enumerate}
\item if $\mathcal F$ satisfies the division bound at twisted places, there are $C$ and $x_0$ such
  that for every scale $x\ge x_0$ and member, the division kernel data carry a division package
  (Definition~\ref{is-division-package}) with constants $(C,C_{abs})$, whose tuples, evaluations and
  words are those of the data, together with the record numbers, the modules and the division
  inputs other than the valuations, and the three ties, used in the proof of
  Theorem~\ref{is-division-link-holds};
\item eventually the multiplicity bound of Chapter~\ref{chapter-period-divisibility}
  (Definition~\ref{pd-multiplicity-bound}) holds at every member;
\item eventually every used tuple of characters is cut;
\item there is $C$ such that eventually, at every member, the explicit heights satisfy the third
  complement: the quantities $h_u+(\sum_ta_t)\log T_u$ differ across words by at most
  $C(|T|aL_a+mB_0)$;
\item if $\mathcal F$ satisfies the division bound at twisted places, then eventually
  $s+\mu_{\max}(RR)\le x^{C_H}b$ for the explicit heights;
\item there is $C$ such that eventually, at every member, every word $u$ and every lift $j$,
  \[
   \Big|\sum_{w\in\operatorname{supp}}\mathrm{sz}_w(E_{u,j}(w))-\Big(\sum_ta_t\Big)\log(g_0g_1)\Big|\le CB_0
   \quad\text{and}\quad
   \sum_{w\in\operatorname{supp},\ p_w\notin\{g_0,g_1\}}|\mathrm{sz}_w(E_{u,j}(w))|\le CB_0;
  \]
\item eventually the Bessel pairings of every member agree with the standard Bessel periods: there is
  an isometry $V$ of the space of top forms such that for every word $u$ there are $a_u\ne0$ and an
  isometry $U_u$ with $B_u(\varphi,F)=a_u\mathcal B_u(U_u\varphi,VF)$ for all $\varphi,F$;
\item eventually the Bessel pairings of every member are rational;
\item eventually $RR$ is finitely generated and spans, at every member;
\item eventually $L_a\ge0$ at every member.
\end{enumerate}
\end{definition}
\leanref{Development/AutomorphicArithmetic/ThetaCorrespondence/ThetaPairLifts/ThirdPlaceSeries/DIVRecBridgeSel.lean}{Formalism.L2.DIV.DivisionKernelConstantBounds}
\leanref{Formalism/SharpRestriction/Divisibility/LocalQuadraticParts.lean}{Formalism.MainNotions.LocalQuadraticParts.adm11}
\leanref{Formalism/SharpRestriction/Divisibility/LocalQuadraticParts.lean}{Formalism.MainNotions.LocalQuadraticParts.adm13}
\leanref{Formalism/LocalToGlobal/Divisibility/RecBridge.lean}{Formalism.L2.DIV.Compl3}

\begin{theorem}\label{is-division-link-holds}
Let $\mathcal D$ be a class of families and $C_{abs}$ a real number, and suppose that every family
in $\mathcal D$ carries division kernel data at every member at all large scales. Suppose that,
eventually along $\mathcal D$, every family that satisfies the size conditions on $T_u$ and on the
stabilizer denominators (Definition~\ref{asm-designed}) satisfies the division kernel bounds
(Definition~\ref{is-division-kernel-bounds}) for $(C_{abs},C_H)$. Then, eventually along
$\mathcal D$, every family that satisfies these size conditions satisfies the conditional division
link (Definition~\ref{is-division-link}) with constants $C_{abs},C_H$.
\end{theorem}
\leanref{Development/AutomorphicArithmetic/ThetaCorrespondence/ThetaPairLifts/ThirdPlaceSeries/DIVRecBridgeSel.lean}{Formalism.L2.DIV.closed_linkDIVt_of_closed_bridge_kernel}
\leanref{Formalism/LocalToGlobal/Divisibility/RecBridge.lean}{Formalism.L2.DIV.linkDIVt_of_rowsF}
\leanref{Formalism/LocalToGlobal/Divisibility/RecBridge.lean}{Formalism.L2.DIV.linkDIV_at_words}
\leanref{Formalism/LocalToGlobal/Divisibility/RecBridge.lean}{Formalism.L2.DIV.DivisionPackageFace}
\leanref{Formalism/LocalToGlobal/Divisibility/RecBridge.lean}{Formalism.L2.DIV.DivisionPackageFace.wordFacts}
\leanref{Formalism/LocalToGlobal/Divisibility/RecBridge.lean}{Formalism.L2.DIV.wordFacts_of_face}
\leanref{Formalism/LocalToGlobal/Divisibility/LinkDIVL2R3Col.lean}{Formalism.L2.DIV.linkDIV_L2R_face}
\leanref{Formalism/LocalToGlobal/GroupLattices/ServJointF.lean}{Formalism.L2.G.ServJointF.val2AllAtRF}
\begin{proof}
The thresholds for $\mathcal D$ in the hypothesis serve for the conclusion, so it suffices to fix a
family $\mathcal F$ that satisfies the size conditions and the division kernel bounds, and the
division bound at twisted places, and to prove the division link for $\mathcal F$.

Let $x$ exceed the scales of all ten parts of Definition~\ref{is-division-kernel-bounds}, and let
$C$ be the largest of their constants and $0$. Fix a member $i$ at $x$, with coefficient field $L$.
Part (1) gives a division package over $L$ (Definition~\ref{is-division-package}) whose columns have
glued lattice $RR$, in the metric $X_E$, identified isometrically with the glued lattice of $D$, and,
for every word $u$, an isometric identification of the column $L^{a(u)}$ with the space of $N_u$
under which $B_u$ is the map of the division data. It also gives, for every word $u$: record numbers
(Definition~\ref{is-record-numbers}); an $\mathcal O_L$-module $N^{\flat}_u\subset L^{a(u)}$ with
$N_u=\big(\prod_w\mathfrak p_w^{-\lfloor e_ws(u,w)\rfloor}\big)N^{\flat}_u$, whose image lies in the
module $N^{\circ}_u$ of the size data because $N_u=I_uN^{\circ}_u$; the division inputs of
Definition~\ref{is-division-inputs} with $s(w)=\lfloor e_ws(u,w)\rfloor$, except the valuations (9);
and the three ties $h^{\mathrm{rec}}-\overline{\log\kappa}=h_u$, $2F^{\mathrm{pk}}_o=F_o$, where
$F^{\mathrm{pk}}_o$ is the number $F_o$ of Definition~\ref{is-division-inputs} and $F_o$ that of
Definition~\ref{is-division-link}, and
\[
 \mathrm{rhs}\leq\mathrm{MAIN}(u)+C|T|aL_a+Cm(B_0+N_c)x^{C_{abs}\varepsilon}.
\]

The valuations (9) come from the multiplicity bound of part (2). Let $w$ be a finite place,
$\varphi\in N^{\flat}_u$ and $F\in M_{\mu_w}$. The tuple $\mu_w$ is used by $D$, so it is cut by
part (3), and the columns of the package are tied to the periods of $D$; for a cut tuple these ties
identify $\mathrm{Pt}_u(\mu_w)(\varphi,F)$ with the period of $\varphi\in N^{\circ}_u$ against an
integral top form for $\mu_w$. By Definition~\ref{pd-multiplicity-bound} this period is zero or has
$v_w\geq\lfloor e_ws(u,w)\rfloor=s(w)$.

Theorem~\ref{is-division-inequality} now applies at every word $u$, with $N_u$ there equal to
$N^{\flat}_u$ and $N'_u$ there equal to $N_u$. Its parts (1) and (2), with the identifications, give
$B_u(N_u)\subset CC$ and $h(B_u)\leq h^{\mathrm{rec}}-\overline{\log\kappa}=h_u$ when $B_u\neq0$. Its
part (3), with the first two ties, gives $\mu(CC)+h_u-F_o\leq\mathrm{rhs}$, and the third tie gives
\[
 \mu(CC)+h_u-F_o\leq\mathrm{MAIN}(u)+C|T|aL_a+Cm(B_0+N_c)x^{C_{abs}\varepsilon}.
\]
The last inequality is the \emph{division inequality}.

These facts are conditions (1), (6) and (7) of the division link. The heights $h_u$ of the family
are the explicit heights, so part (5) gives condition (2) and part (4) gives condition (9). Part (6)
gives condition (8), since $\{g_0,g_1\}$ are the level primes. Parts (7) and (8) together are the
Bessel identity, condition (4). Part (9) is condition (5). Condition (3) holds at every member of
every family at every scale $x>x_1$: the parameter choices of a family give no gains and no losses
at places above $X''$ outside the listed primes.
\end{proof}

The division kernel bounds are themselves established in Chapter~\ref{chapter-assembly}. Part (7),
the agreement of the Bessel pairings with standard Bessel periods, follows from the multiplicity
condition of the realization, the finiteness of the spectral data and the fullness of the image of
the base vectors. The finiteness of the spectral data follows from the decay of the archimedean
Bessel functions $W_u$ (Definition~\ref{is-wu-decay}), the finite dimensionality of rows and the gap
hypothesis, which holds eventually along designed families. The twisted places of the rows $S_0$
are the places with $\mu_w\ne\mu_0$, and every such place lies in the support of the size data;
Chapter~\ref{chapter-assembly} uses this to build the division packages of part (1).

\section{Near-orthogonality of the columns}\label{section-is-orthogonality}

\begin{definition}\label{is-orthogonality-link}
Let $D$ be a definite-side datum (Definition~\ref{th-definite-datum}) with ranks $m,n$, $f=2n$,
words $u$ and normalizing dimension $D^{*}>0$, and with the Bessel-period data of
Definition~\ref{th-definite-datum}(8): a finite-dimensional Hilbert space $\mathcal H$, whose
dimension is the total number $D_{\mathrm{all}}$ of rows of $D$, and for each word $u$ a
finite-dimensional Hilbert space $\mathcal C_u$, isometric to $M_u$, with a linear map
$P_u:\mathcal C_u\to\mathcal H$ and positive volumes $\mathrm{vol}_N(u)$, $\mathrm{vol}_H(u)$. Let
$\dim\pi_u$ be the product over the types $\tau$ of the dimensions of the representations of
Harish-Chandra type attached to the parameter $\lambda_u(\tau)$. The \emph{Gram matrix} of the
columns of $u$ and $u'$ is
\[
 K_{uu'}=(c_uc_{u'})^{-1/2}P_u^{*}P_{u'}:\mathcal C_{u'}\to\mathcal C_u,\qquad
 c_u=\frac{\mathrm{vol}_N(u)\,\mathrm{vol}_H(u)}{\dim\pi_u}.
\]
Let size data at scale $x$ be given (Definition~\ref{pd-size-data}), with constant $C_\Gamma$, and put
$\Gamma=\exp(C_\Gamma m^{5}fx^{\varepsilon})$. For real $C_0,\eta_0,\eta_1,e_0$, the
\emph{orthogonality bounds} hold at $D$ if, whenever $C_0^{-1}\le(\dim M_u)\Gamma^{2}/D^{*}\le C_0$ for
every word $u$:
\begin{enumerate}
\item $\|K_{uu'}\|_{HS}^{2}\le(1+\eta_0)\dim M_u\dim M_{u'}/D^{*}$ for all words $u\ne u'$;
\item $\|K_{uu}-1\|_{HS}^{2}\le e_0\dim M_u$ for every word $u$;
\item $D_{\mathrm{all}}\le(1+\eta_1)D^{*}$.
\end{enumerate}
In the terms of Definition~\ref{hs-near-orthogonal}, these say that the family
$(P_u,c_u^{-1})_u$ is $(D^{*},\eta_0,\eta_1,e_0)$-nearly orthogonal in $\mathcal H$. A family
$\mathcal F$ of exponent $\varepsilon$ (Definition~\ref{pd-family}) satisfies the
\emph{orthogonality link} with constant $C_0$ if there is $x_0$ such that for every scale
$x\ge x_0$ and every member there are $\eta_0,\eta_1\le e^{-x^{\varepsilon/8}}$ and $e_0\ge0$ with
$\sqrt{e_0}\le e^{-x^{\varepsilon/8}}$ for which the orthogonality bounds hold.
\end{definition}
\leanref{Formalism/SharpRestriction/GlobalTheoremInputs.lean}{Formalism.MainNotions.LinkS}
\leanref{Formalism/SharpRestriction/GlobalTheoremInputs.lean}{Formalism.MainNotions.LinkS_at}
\leanref{Formalism/SharpRestriction/CuspSpanNormBounds/RelevantZeroGramTerms.lean}{Formalism.MainNotions.S.Kmat}
\leanref{Formalism/SharpRestriction/CuspSpanNormBounds/RelevantZeroGramTerms.lean}{Formalism.MainNotions.S.c0}

In the bundles of the proof, the space $\mathcal H$ is a space of square-integrable automorphic
forms on the indefinite unitary group, and $P_u$ sends a vector of $\mathcal C_u$ to a Poincaré
series. The summand of the Poincaré series is built from the archimedean Bessel function $W_u$ of the word $u$ and a character
of the finite adeles, and the series is summed over the rational points of the stabilizer of the
Bessel data.

\begin{definition}\label{is-wu-decay}
Let $D$ be a definite-side datum (Definition~\ref{th-definite-datum}) with $r=k+m-1$ and
Harish-Chandra parameters $b_1,\dots,b_m$ at the types $\tau$. The \emph{gap hypothesis} holds for
$D$ if $r\ge1801$ and $b_{j+1}(\tau)-b_j(\tau)\ge16r^{19}$ for every type $\tau$ and every $j<m$. For
a word $u$ of an automorphic setting attached to $D$, let $W_u:G_\infty\to V_{\pi_u}$ be the
archimedean Bessel function of $u$, attached to the matrix coefficient of the discrete series. The
\emph{decay of $W_u$} holds if $W_u$ is reproduced by convolution with a bounded measurable function
supported in a compact neighbourhood of $1$, and
\[
 \int|W_u(g)|\,\nu_u(g)^{N_1}\,d\nu_Q(g)<\infty,\qquad
 \nu_u(g)=\inf_{s\in S_\infty(u)}\prod_v\|s g\|_v,\qquad
 N_1=4n_I^{2}\sqrt{2\max_vq_v},
\]
where $S_\infty(u)$ is the archimedean stabilizer of the Bessel data, $\nu_Q$ the quotient measure,
$n_I$ the rank and $q_v$ the local invariants of the setting. The \emph{decay hypothesis for
$W_u$} says: for every definite-side datum satisfying the gap hypothesis and every automorphic
setting attached to it, every word satisfies the decay of $W_u$, and the archimedean matrix
coefficient is integrable.
\end{definition}
\leanref{Formalism/LocalToGlobal/SiegelLatticeTheory/FlagStabilizerTensors/RootLatticeThresholds/WuDecayAll.lean}{Formalism.L2.SG.Pf.WuDecayOfGap}
\leanref{Formalism/LocalToGlobal/SiegelLatticeTheory/HaarTransversals/WuDecay.lean}{Formalism.L2.SG.WuDecay.WuDecay}
\leanref{Formalism/LocalToGlobal/SiegelLatticeTheory/HaarTransversals/ConcreteSettingExists.lean}{Formalism.L2.SG.ConcreteSettingExists.GapHyp}

Alpöge proves the decay hypothesis for $W_u$ from the uniqueness of local Bessel maps under the gap
hypothesis.
\leanref{CompanionTheory/AutomorphicTheory/LocalBesselUniqueness/BesselModels/MultiplicityOneUnitaryBranchingOfGapHyp.lean}{CompanionTheory.Lime.BM.wuDecayOfGap_holds'}

\begin{definition}\label{is-gram-terms}
Let $D$ be a definite-side datum with the Bessel-period data of Definition~\ref{is-orthogonality-link}
and Gram matrices $K_{uu'}$. The \emph{unfolding} of the Bessel-period data consists, for all words
$u,u'$, of a set of \emph{terms} $t$, operators $a_t\colon\mathcal C_{u'}\to\mathcal C_u$, a set of
\emph{identity terms}, and elements $G_t$ of the totally real field $F$ of the host, such that
$\sum_t\|a_t\|$ converges and $\sum_ta_t=K_{uu'}$; it is the evaluation of the Gram matrices by
unfolding the Poincaré series, a term being a double coset of the stabilizer of the Bessel data. The
\emph{leading part} of $K_{uu'}$ is $K''_{uu'}=\sum_ta_t$ over the terms with $G_t=0$ that are not
identity terms. For size data with the quantity $b$ and the constant $C_\Gamma$ put
\[
 \eta_0=b^{-6},\qquad x_K=\exp(-kf\log k),\qquad \Gamma=\exp(C_\Gamma m^5fx^{\varepsilon}).
\]
\end{definition}
\leanref{Formalism/SharpRestriction/CuspSpanNormBounds/DataS.lean}{Formalism.MainNotions.S.Draft.SData}
\leanref{Formalism/SharpRestriction/CuspSpanNormBounds/RelevantZeroGramTerms.lean}{Formalism.MainNotions.S.KOf}
\leanref{Formalism/SharpRestriction/CuspSpanNormBounds/NonGenericAndKTwoTheorems.lean}{Formalism.MainNotions.NonGenericAndKTwoTheorems.Kpp}
\leanref{Formalism/SharpRestriction/CuspSpanNormBounds/AdmissibleBoundsFromSource.lean}{Formalism.MainNotions.AdmissibleBoundsFromSource.xK}
\leanref{Formalism/SharpRestriction/CuspSpanNormBounds/AdmissibleBoundsFromSource.lean}{Formalism.MainNotions.AdmissibleBoundsFromSource.eta0}

\begin{definition}\label{is-unfolding-term-data}
Let $I$ be a finite set and $H,H_1$ complex Hilbert spaces. \emph{Term data} in $(H,H_1)$ over $I$
consist of a set $T$ of \emph{terms}; for every term $t$ an element $a_t\in H$, a vector
$v_t\in H_1$, a \emph{row} $\rho(t)$ with a \emph{type} $\iota(\rho(t))\in I$, a real number $G_t$,
and the properties of being an \emph{identity term} and of being \emph{generic}; a complex matrix
$\Upsilon=(\Upsilon_{ij})_{i,j\in I}$ and a real number $\Upsilon_1$; and for every pair of terms an
integer $s_0(t,t')\geq0$, the \emph{corank}, and the property of being \emph{scalar}. The identity,
\emph{degenerate}, \emph{non-generic} and generic terms form the sets
\[
 T_{\mathrm{id}},\quad T_{\mathrm{deg}}=\{t\notin T_{\mathrm{id}}:G_t=0\},\quad
 T_{\mathrm{ng}}=\{t:G_t\neq0,\ t\text{ not generic}\},\quad
 T_{\mathrm{gen}}=\{t:G_t\neq0,\ t\text{ generic}\},
\]
and $T_{\neq0}=\{t:G_t\neq0\}$. For $S\subset T$ and $i\in I$ put
\[
 a(S)=\sum_{t\in S}a_t,\qquad v_i(S)=\sum_{t\in S,\ \iota(\rho(t))=i}v_t,\qquad
 \Phi(S)=\Upsilon_1\sum_{i\in I}\|v_i(S)\|^2,
\]
$K=a(T)$ and $K^0=a(T_{\mathrm{id}})+a(T_{\mathrm{gen}})+a(T_{\mathrm{ng}})$. A \emph{pair} of $S$
is a pair $(t,t')\in S\times S$ with $\rho(t)=\rho(t')$. Write $P(S)$ for the set of pairs of $S$,
$P_{\geq2}(S)$ and $P_1(S)$ for those with $s_0(t,t')\geq2$ and with $s_0(t,t')=1$, and
$P_1^{\mathrm{sc}}(S)$ and $P_1^{\mathrm{ns}}(S)$ for those in $P_1(S)$ that are scalar and that are
not. For a set $Q$ of pairs put
\[
 \Sigma(Q)=\Upsilon_1\sum_{(t,t')\in Q}\operatorname{Re}\langle v_t,v_{t'}\rangle,\qquad
 |\Sigma|(Q)=\Upsilon_1\sum_{(t,t')\in Q}|\langle v_t,v_{t'}\rangle|.
\]
The term data are
\begin{enumerate}
\item \emph{absolutely convergent} if $\sum_t\|a_t\|$ and $\sum_t\|v_t\|$ converge and $\Upsilon_1>0$;
  then all the sums above converge absolutely;
\item \emph{identity-degenerate} if $G_t=0$ for every identity term $t$;
\item \emph{row-orthogonal} if $\langle v_t,v_{t'}\rangle=0$ whenever
  $\iota(\rho(t))=\iota(\rho(t'))$ and $\rho(t)\neq\rho(t')$;
\item of \emph{positive corank} if $s_0(t,t')\geq1$ whenever $\rho(t)=\rho(t')$, $G_t\neq0$ and
  $G_{t'}\neq0$;
\item \emph{$\Upsilon$-bounded} by a real number $P$ if the matrix $P\Upsilon_1\cdot1-\Upsilon$ is
  positive semidefinite;
\item \emph{norm-expanding} on $S\subset T$ if
  $\|a(S)\|^2=\operatorname{Re}\sum_{i,j\in I}\Upsilon_{ij}\langle v_i(S),v_j(S)\rangle$.
\end{enumerate}
For real numbers $\varepsilon$, $A$ and $d$, the term data satisfy the \emph{large-row bound}
$(\varepsilon,d)$ if $|\Sigma|(P_{\geq2}(T_{\neq0}))\le\varepsilon d$, the \emph{non-scalar bound}
$(\varepsilon,d)$ if $|\Sigma|(P_1^{\mathrm{ns}}(T_{\mathrm{gen}}))\le\varepsilon d$, the
\emph{main-row bound} $(A,d)$ if $\Sigma(P_1^{\mathrm{sc}}(T_{\mathrm{gen}}))\le Ad$, and the
\emph{non-generic bound} $(\varepsilon,d)$ if $\Sigma(P_1(T_{\mathrm{ng}}))\le\varepsilon d$.
\end{definition}
\leanref{Formalism/HermitianCurvature/StatementPB.lean}{Formalism.FourierDecayAlt.PB.Terms}
\leanref{Formalism/HermitianCurvature/StatementPB.lean}{Formalism.FourierDecayAlt.PB.Terms.K0}
\leanref{Formalism/HermitianCurvature/StatementPB.lean}{Formalism.FourierDecayAlt.PB.Terms.GGBound}
\leanref{Formalism/HermitianCurvature/StatementPB.lean}{Formalism.FourierDecayAlt.PB.Terms.RowBig}
\leanref{Formalism/HermitianCurvature/StatementPB.lean}{Formalism.FourierDecayAlt.PB.Terms.RowMain}

\begin{lemma}\label{is-unfolding-identities}
Let term data (Definition~\ref{is-unfolding-term-data}) be absolutely convergent.
\begin{enumerate}
\item If they are identity-degenerate, then $T$ is the disjoint union of $T_{\mathrm{id}}$,
  $T_{\mathrm{deg}}$, $T_{\mathrm{gen}}$ and $T_{\mathrm{ng}}$, and $K=K^0+a(T_{\mathrm{deg}})$.
\item If they are norm-expanding on $S$ and $\Upsilon$-bounded by $P$, then
  $\|a(S)\|^2\le P\,\Phi(S)$.
\item If they are row-orthogonal and of positive corank, then $\Phi(S)=\Sigma(P(S))$ for every
  $S\subset T$, and
  \[
   \Phi(T_{\mathrm{gen}})=\Sigma(P_{\geq2}(T_{\mathrm{gen}}))+\Sigma(P_1^{\mathrm{ns}}(T_{\mathrm{gen}}))
   +\Sigma(P_1^{\mathrm{sc}}(T_{\mathrm{gen}})),\qquad
   \Phi(T_{\mathrm{ng}})=\Sigma(P_{\geq2}(T_{\mathrm{ng}}))+\Sigma(P_1(T_{\mathrm{ng}})).
  \]
\item If $Q\subset Q'$ are sets of pairs, then $|\Sigma|(Q)\le|\Sigma|(Q')$ and
  $|\Sigma(Q)|\le|\Sigma|(Q)$.
\end{enumerate}
\end{lemma}
\leanref{Formalism/LocalToGlobal/SiegelLatticeTheory/FlagStabilizerTensors/PackageTie.lean}{Formalism.L2.SG.PackageTie.K_eq_parts}
\leanref{Formalism/HermitianCurvature/HermitianCurvatureProofs/GrowthOfBasic.lean}{Formalism.FourierDecayAlt.P.price_of_index}
\leanref{Formalism/HermitianCurvature/HermitianCurvatureProofs/GrowthOfBasic.lean}{Formalism.FourierDecayAlt.P.equal_index}
\leanref{Formalism/HermitianCurvature/HermitianCurvatureProofs/GrowthOfBasic.lean}{Formalism.FourierDecayAlt.P.pair_split_alt}
\leanref{Formalism/HermitianCurvature/HermitianCurvatureProofs/GrowthOfBasic.lean}{Formalism.FourierDecayAlt.P.subfamilyAlt}
\begin{proof}
(1) An identity term has $G_t=0$. A term that is not an identity term is degenerate if $G_t=0$, and
generic or non-generic if $G_t\neq0$. So the four sets partition $T$, and summing the absolutely
convergent series $\sum_ta_t$ over each of them gives
$K=a(T_{\mathrm{id}})+a(T_{\mathrm{deg}})+a(T_{\mathrm{gen}})+a(T_{\mathrm{ng}})$.

(2) If $c\cdot1-\Upsilon$ is positive semidefinite, then
$\operatorname{Re}\sum_{i,j}\Upsilon_{ij}\langle w_i,w_j\rangle\le c\sum_i\|w_i\|^2$ for all vectors
$w_i\in H_1$. Take $c=P\Upsilon_1$ and $w_i=v_i(S)$; the norm expansion gives
$\|a(S)\|^2\le P\Upsilon_1\sum_i\|v_i(S)\|^2=P\,\Phi(S)$.

(3) Expanding, $\|v_i(S)\|^2$ is the sum of $\langle v_t,v_{t'}\rangle$ over the $t,t'\in S$ of
type $i$. By row orthogonality only the pairs with $\rho(t)=\rho(t')$ contribute, so
$\Phi(S)=\Upsilon_1\sum_{(t,t')\in P(S)}\langle v_t,v_{t'}\rangle$; this is real, so it equals
$\Sigma(P(S))$. Every pair of $T_{\mathrm{gen}}$ or of $T_{\mathrm{ng}}$ has $G_t\neq0$ and
$G_{t'}\neq0$, so $s_0(t,t')\geq1$ by positive corank. Hence these pairs split into those with
$s_0\geq2$ and those with $s_0=1$, and the latter into scalar and non-scalar pairs.

(4) The summands of $|\Sigma|$ are nonnegative, and $|\operatorname{Re}z|\le|z|$.
\end{proof}

\begin{lemma}\label{is-unfolding-pair-bound}
Let term data (Definition~\ref{is-unfolding-term-data}) be absolutely convergent, row-orthogonal,
of positive corank, norm-expanding on $T_{\mathrm{gen}}$ and on $T_{\mathrm{ng}}$, and
$\Upsilon$-bounded by $P\geq0$, and let $e\in H$ with $a(T_{\mathrm{id}})=e$. Let $A,\varepsilon,\varepsilon'$
be real numbers and $d\geq0$, and suppose that the term data satisfy the large-row and non-scalar
bounds $(\varepsilon,d)$, the main-row bound $(A,d)$ and the non-generic bound $(\varepsilon',d)$.
Then
\[
 \|K^0-e\|^2\le\Big(\sqrt{P(A+2\varepsilon)}+\sqrt{P(\varepsilon+\varepsilon')}\Big)^2d.
\]
\end{lemma}
\leanref{Formalism/LocalToGlobal/SiegelLatticeTheory/CokernelGLCounts/RowsOther.lean}{Formalism.L2.SG.RowsOther.diag_of_abstract}
\leanref{Formalism/LocalToGlobal/SiegelLatticeTheory/CokernelGLCounts/RowsOther.lean}{Formalism.L2.SG.RowsOther.nonIdentity_K0}
\leanref{Formalism/HermitianCurvature/HermitianCurvatureProofs/GrowthOfBasic.lean}{Formalism.FourierDecayAlt.P.pair_bound_abstract}
\leanref{Formalism/HermitianCurvature/HermitianCurvatureProofs/GrowthOfBasic.lean}{Formalism.FourierDecayAlt.P.GrowthPolynomialOperator.t_PairBoundAbstract}
\begin{proof}
By definition $K^0-e=a(T_{\mathrm{gen}})+a(T_{\mathrm{ng}})$. The sets $T_{\mathrm{gen}}$ and
$T_{\mathrm{ng}}$ lie in $T_{\neq0}$, so $P_{\geq2}(T_{\mathrm{gen}})$ and
$P_{\geq2}(T_{\mathrm{ng}})$ lie in $P_{\geq2}(T_{\neq0})$. Parts (3) and (4) of
Lemma~\ref{is-unfolding-identities} give
\[
 \Phi(T_{\mathrm{gen}})\le|\Sigma|(P_{\geq2}(T_{\neq0}))+|\Sigma|(P_1^{\mathrm{ns}}(T_{\mathrm{gen}}))
 +\Sigma(P_1^{\mathrm{sc}}(T_{\mathrm{gen}}))\le(A+2\varepsilon)d
\]
and $\Phi(T_{\mathrm{ng}})\le|\Sigma|(P_{\geq2}(T_{\neq0}))+\Sigma(P_1(T_{\mathrm{ng}}))
\le(\varepsilon+\varepsilon')d$. Since $P\geq0$, part (2) gives
$\|a(T_{\mathrm{gen}})\|^2\le P(A+2\varepsilon)d$ and
$\|a(T_{\mathrm{ng}})\|^2\le P(\varepsilon+\varepsilon')d$. The triangle inequality
$\|x+y\|\le\|x\|+\|y\|$ gives the bound.
\end{proof}

\begin{definition}\label{is-unfolding-evaluation}
Let $D$ be a definite-side datum (Definition~\ref{th-definite-datum}) with ranks $k,m,f$, level
primes $g_0,g_1$, and the Bessel-period data, Gram matrices $K_{uu'}$ and unfolding of
Definitions~\ref{is-orthogonality-link} and~\ref{is-gram-terms}. A \emph{concrete evaluation} of the
Gram matrices of $D$ consists of:
\begin{enumerate}
\item an automorphic setting attached to $D$, with the archimedean Bessel functions
  $W_u\colon G_\infty\to V_{\pi_u}$ of Definition~\ref{is-wu-decay}, the rows $\Lambda$ of every word
  $u$, column representatives $h_\varrho$ (each with a row), among them a base column representative
  $h_\Lambda$ with row $\Lambda$ for each row $\Lambda$, and the column space
  $\mathcal C^{\sharp}_u=\bigoplus_\varrho V_{\pi_u}$ of $u$, fundamental domains, a level basis
  $e^0_a,f^0_a,z^0,b_i,b'_i$ with its level subgroup $K_2$
  (Definition~\ref{asm-division-kernel-finite-neat-level}) and a character $\chi_2=\chi_N\chi_L$ of
  $K_2$, and reductions $\mathrm{res}_v$ at finitely many places $v$, with
  residue fields $\mathbf F_{q_v}$ of $q_v\in\{g_0,g_1\}$ elements and nontrivial additive
  characters $\psi_{0,v}$;
\item an isomorphism of the Bessel-period data of $D$ with those of the setting, and isometries
  $\sharp\colon\mathcal C_u\to\mathcal C^{\sharp}_u$;
\item the ties of the setting to $D$: its level subgroup is normalized and neat,
  $\dim\mathcal C^{\sharp}_u=\dim M_u$ for every word $u$, its normalizing dimension is $D^{*}$, the
  decay of $W_u$ holds and the archimedean matrix coefficient is integrable, its measures and
  characters are compatible with the local data, the local ties hold, and
  $1\le T_u\kappa_S^{2r_S}$ for every word $u$, where $r_S=m/2-1$.
\end{enumerate}
For words $u,u'$, the \emph{unfolded term data} $T_{uu'}$ are the following term data
(Definition~\ref{is-unfolding-term-data}) over $I=\prod_vGL_k(\mathbf F_{q_v})$, with $H$ the space of
operators $\mathcal C^{\sharp}_{u'}\to\mathcal C^{\sharp}_u$ with the Hilbert-Schmidt norm and
$H_1=\bigoplus_{(\Lambda,\Lambda')}\operatorname{HS}(V_{\pi_{u'}},V_{\pi_u})$, over the pairs of rows of
$u$ and of $u'$.
\begin{enumerate}
\item A term $t$ consists of rows $\Lambda$ of $u$ and $\Lambda'$ of $u'$ and a double coset
  $N_\Lambda\gamma N_{\Lambda'}$, under the groups $N_\Lambda$ and $N_{\Lambda'}$ of automorphisms of
  the rows that are unipotent along their flags, of a rational isometry $\gamma$ with $\gamma\Lambda'=\Lambda$ that is linked to an adelic point $x_t$ and
  satisfies the local conditions of the terms of the unfolding. It lies over the term of the unfolding
  of the setting given by $\Lambda$ and the flag $\gamma F'$, where $F'$ is the flag of $\Lambda'$;
  over each term of the unfolding lie finitely many terms $t$. The term $t$ is an identity term, is
  generic, or is \emph{relevant} when the term below it is; relevance means that the differentials
  of the characters of the Bessel data of the flag of $\Lambda$ and of $\gamma F'$ agree on the
  intersection of the Lie algebras of their unipotent radicals. Finally $G_t$ is $0$ or $1$
  according as the invariant $G$ of the term below $t$ vanishes or not.
\item For $\kappa\in G_{\mathrm f}$ let $\lambda(\kappa)\in I$ be the reduction at the places $v$ of
  the Levi block of $\kappa$: $\lambda_v(\kappa)$ is the $k\times k$ matrix
  $\big(\mathrm{res}_v\,h(b'_i,\kappa b_j)\big)_{ij}$ over $\mathbf F_{q_v}$ if it is invertible, and
  $1$ otherwise, where $\mathrm{res}_v$ is read as $0$ on non-integral elements. Put
  $k(t)=(h_\Lambda x_u)^{-1}x_{t,\mathrm f}(h_{\Lambda'}x_{u'})$, with the base points $x_u$ and
  $x_{u'}$ of the words and the base column representatives $h_\Lambda$ and $h_{\Lambda'}$. The row
  of $t$ is $\rho(t)=(\Lambda,\Lambda',\iota(t))$ with $\iota(t)=\lambda(k(t))$.
\item Let $\Omega_t\subset G_\infty$ be a fundamental domain for the discrete group $\Delta_t$ of
  archimedean components of the rational elements of $N_\Lambda\cap\gamma N_{\Lambda'}\gamma^{-1}$,
  and let
  \[
   O_t\colon V_{\pi_{u'}}\to V_{\pi_u},\qquad
   O_t\xi=\int_{\Omega_t}\big\langle W_{u'}(x_{t,\infty}^{-1}y),\xi\big\rangle\,W_u(y)\,d\mu(y)
  \]
  be the orbital integral. Put $c_{uu'}=\operatorname{vol}(K_2)(c_uc_{u'})^{-1/2}$ and
  $k_{\varrho\varrho'}(t)=(h_\varrho x_u)^{-1}x_{t,\mathrm f}(h_{\varrho'}x_{u'})$. The block of $a_t$ at column representatives $\varrho$ of $u$ and $\varrho'$ of $u'$
  is $c_{uu'}\chi_2(k_{\varrho\varrho'}(t))O_t$ if their rows are $\Lambda$ and $\Lambda'$ and
  $k_{\varrho\varrho'}(t)\in K_2$, and $0$ otherwise. The vector $v_t$ vanishes outside the
  coordinate $(\Lambda,\Lambda')$, where it is $\sqrt{w(\Lambda)w(\Lambda')}\,c_{uu'}
  \chi_N(k(t))O_t$, with the weights $w$ of the rows and $k(t)$ as in~(2).
\item $\Upsilon_{ij}=\overline{\upsilon(j^{-1}i)}$ and $\Upsilon_1=\operatorname{Re}\upsilon(1)$, where
  \[
   \upsilon(g)=\prod_v\frac{1}{|U_k(\mathbf F_{q_v})|}\sum_{y\in GL_k(\mathbf F_{q_v})}
   \psi_{N,v}(yg_vy^{-1})
  \]
  and $\psi_{N,v}$ is the function on $GL_k(\mathbf F_{q_v})$ equal to
  $\psi_{0,v}(\sum_in_{i,i+1})$ on upper unitriangular matrices $n$ and to $0$ elsewhere.
\item Let $t,t'$ be terms with the same rows $\Lambda,\Lambda'$, and let $\gamma$ and $\gamma'$ be the
  chosen representatives of their double cosets. The corank $s_0(t,t')$ and the scalarity of
  $(t,t')$ are invariants of $\Lambda$, its flag $F$, the two flags $\gamma F'$ and $\gamma'F'$ and the
  isometry $\gamma'\gamma^{-1}$: $s_0(t,t')=s_0(F,\gamma F',\gamma'F',\gamma'\gamma^{-1})$, with $k$ the
  rank of $D$, and $(t,t')$ is scalar if $s_0(t,t')=1$ and $(\gamma F',\gamma'F',\gamma'\gamma^{-1})$ is
  scalar for $F$, in the sense of
  Definition~\ref{asm-division-kernel-finite-pairs-invariants}(1) and~(2). For terms with
  different rows, $s_0(t,t')=0$ and $(t,t')$ is not scalar.
\end{enumerate}
\end{definition}
\leanref{Formalism/LocalToGlobal/SiegelLatticeTheory/FlagStabilizerTensors/ThmSOwn.lean}{Formalism.L2.SG.ThmSOwn.ConcK}
\leanref{Formalism/LocalToGlobal/SiegelLatticeTheory/CokernelGLCounts/RowsConc.lean}{Formalism.L2.SG.RowsConc.ConcTie}
\leanref{Formalism/LocalToGlobal/SiegelLatticeTheory/FlagStabilizerTensors/PoincareConstantEq.lean}{Formalism.L2.SG.Full.Full}
\leanref{Formalism/LocalToGlobal/SiegelLatticeTheory/FlagStabilizerTensors/ParametersTermRelations.lean}{Formalism.L2.SG.PackageTie2.FinX}
\leanref{Formalism/LocalToGlobal/SiegelLatticeTheory/ClosedLocalOrbits/TermOp.lean}{Formalism.L2.SG.TermOp.tOf}
\leanref{Formalism/LocalToGlobal/SiegelLatticeTheory/ClosedLocalOrbits/TermOp.lean}{Formalism.L2.SG.TermOp.GTerm}
\leanref{Formalism/LocalToGlobal/SiegelLatticeTheory/ClosedLocalOrbits/TermOp.lean}{Formalism.L2.SG.TermOp.OOt}
\leanref{Formalism/LocalToGlobal/SiegelLatticeTheory/ClosedLocalOrbits/TermOp.lean}{Formalism.L2.SG.TermOp.aTerm}
\leanref{Formalism/LocalToGlobal/SiegelLatticeTheory/ClosedLocalOrbits/TermOp.lean}{Formalism.L2.SG.TermOp.vec}
\leanref{Formalism/LocalToGlobal/SiegelLatticeTheory/ClosedLocalOrbits/TermOp.lean}{Formalism.L2.SG.TermOp.kOf}
\leanref{Formalism/LocalToGlobal/SiegelLatticeTheory/ClosedLocalOrbits/TermOp.lean}{Formalism.L2.SG.TermOp.iotaOf}
\leanref{Formalism/LocalToGlobal/SiegelLatticeTheory/ClosedLocalOrbits/GGMirabolic.lean}{Formalism.L2.SG.GGMirabolic.gam}

\begin{lemma}\label{is-unfolding-computations}
Let $D$ be a definite-side datum with a concrete evaluation of its Gram matrices
(Definition~\ref{is-unfolding-evaluation}), and let $u,u'$ be words. The unfolded term data $T_{uu'}$
are absolutely convergent, identity-degenerate, row-orthogonal, of positive corank, and
norm-expanding on $T_{\mathrm{deg}}$, on $T_{\mathrm{gen}}$ and on $T_{\mathrm{ng}}$. If $q_v\geq3$ for
every $v$, they are $\Upsilon$-bounded by $P_{\mathrm{ev}}=\prod_v(1+3/q_v)$. The sum $K$ of
$T_{uu'}$ represents the Gram matrix of the setting: $K=\sharp\circ K^{\mathrm{set}}_{uu'}\circ
\sharp^{-1}$, where $K^{\mathrm{set}}_{uu'}$ is the Gram matrix of the Bessel-period data of the
setting. For $u=u'$, $a(T_{\mathrm{id}})$ is the identity of $\mathcal C^{\sharp}_u$.
\end{lemma}
\leanref{Formalism/LocalToGlobal/SiegelLatticeTheory/FlagStabilizerTensors/LinkSFinal.lean}{Formalism.L2.SG.ThmSOwn.absConvTied_holds}
\leanref{Formalism/LocalToGlobal/SiegelLatticeTheory/ClosedVolumeStatements/FullExistsTiedHolds.lean}{Formalism.L2.SG.Closed.AbsConvC_outright}
\leanref{Formalism/LocalToGlobal/SiegelLatticeTheory/ClosedLocalOrbits/GaugeLatticeShifts/SatGram.lean}{Formalism.L2.SG.Pf.SatGram.idDegenerate_tO}
\leanref{Formalism/LocalToGlobal/SiegelLatticeTheory/ClosedLocalOrbits/TermOpStmt.lean}{Formalism.L2.SG.TermOp.rowOrth}
\leanref{Formalism/LocalToGlobal/SiegelLatticeTheory/CokernelGLCounts/RankPointCounts/KloostermanDiagonalOfConditions.lean}{Formalism.L2.SG.KappaOffDiagonal.corankPosC_holds}
\leanref{Formalism/LocalToGlobal/SiegelLatticeTheory/ClosedLocalOrbits/GaugeLatticeShifts/SatGram.lean}{Formalism.L2.SG.Pf.SatGram.corankPos_tO}
\leanref{Formalism/LocalToGlobal/SiegelLatticeTheory/CokernelGLCounts/RankPointCounts/MainEqHolds.lean}{Formalism.L2.SG.Pf.normExpDegC_holds}
\leanref{Formalism/LocalToGlobal/SiegelLatticeTheory/CokernelGLCounts/RankPointCounts/DiagCLC.lean}{Formalism.L2.SG.RowsConc.normExpSatTied_holds}
\leanref{Formalism/LocalToGlobal/SiegelLatticeTheory/ClosedLocalOrbits/TermOpStmt.lean}{Formalism.L2.SG.TermOp.ggBound}
\leanref{Formalism/LocalToGlobal/SiegelLatticeTheory/ClosedVolumeStatements/GenericRatioOfBoundHolds.lean}{Formalism.L2.SG.Closed.TermsGGBound_holds}
\leanref{Formalism/LocalToGlobal/SiegelLatticeTheory/FlagStabilizerTensors/LinkSFinal.lean}{Formalism.L2.SG.ThmSOwn.kExpTied_holds}
\leanref{Formalism/LocalToGlobal/SiegelLatticeTheory/ClosedVolumeStatements/UnipotentInLevelSubgroup.lean}{Formalism.L2.SG.Closed.KExpansionC_holds}
\leanref{Formalism/LocalToGlobal/SiegelLatticeTheory/FlagStabilizerTensors/RootLatticeThresholds/NormalizingConstantEq.lean}{Formalism.L2.SG.Pf.kmat_eq}
\leanref{Formalism/LocalToGlobal/SiegelLatticeTheory/CokernelGLCounts/RankPointCounts/DiagCLC.lean}{Formalism.L2.SG.RowsConc.idOneC_holds}
\begin{proof}
Row orthogonality: the vectors $v_t$ and $v_{t'}$ vanish outside the coordinates $(\Lambda,\Lambda')$
of their rows, and if $\rho(t)\neq\rho(t')$ while $\iota(t)=\iota(t')$ these coordinates differ. The
other statements are computations with the unfolded Poincaré series. The absolute convergence
rests on the decay of $W_u$ and the integrability of the archimedean matrix coefficient; the bound
for $\Upsilon$ rests on the nontriviality of the characters $\psi_{0,v}$ and on $q_v\geq3$; the
representation of the Gram matrix by $K$ is the unfolding of the inner products of the Poincaré
series.
\end{proof}

\begin{lemma}\label{is-unfolding-transport}
Let $D$ be a definite-side datum with a concrete evaluation of its Gram matrices
(Definition~\ref{is-unfolding-evaluation}), and let $u,u'$ be words. Under the isometries $\sharp$,
the Gram matrix $K_{uu'}$ becomes the sum $K$ of the unfolded term data $T_{uu'}$, its leading part
$K''_{uu'}$ (Definition~\ref{is-gram-terms}) becomes $a(T_{\mathrm{deg}})$, and for $u=u'$ the
identity of $\mathcal C_u$ becomes the identity of $\mathcal C^{\sharp}_u$. In particular
$\|K''_{uu'}\|_{HS}=\|a(T_{\mathrm{deg}})\|$, $\|K_{uu'}-K''_{uu'}\|_{HS}=\|K-a(T_{\mathrm{deg}})\|$, and
$\|K_{uu}-K''_{uu}-1\|_{HS}=\|K-a(T_{\mathrm{deg}})-a(T_{\mathrm{id}})\|$ for $T_{uu}$.
\end{lemma}
\leanref{Formalism/LocalToGlobal/SiegelLatticeTheory/FlagStabilizerTensors/ThmSOwn.lean}{Formalism.L2.SG.ThmSOwn.conjE_Kmat}
\leanref{Formalism/LocalToGlobal/SiegelLatticeTheory/FlagStabilizerTensors/ThmSOwn.lean}{Formalism.L2.SG.ThmSOwn.conjE_KOf}
\leanref{Formalism/LocalToGlobal/SiegelLatticeTheory/FlagStabilizerTensors/ThmSOwn.lean}{Formalism.L2.SG.ThmSOwn.kof_sOf}
\leanref{Formalism/LocalToGlobal/SiegelLatticeTheory/FlagStabilizerTensors/ThmSOwn.lean}{Formalism.L2.SG.ThmSOwn.conjE_Kpp}
\leanref{Formalism/LocalToGlobal/SiegelLatticeTheory/FlagStabilizerTensors/ThmSOwn.lean}{Formalism.L2.SG.ThmSOwn.hsNormSq_Kpp}
\leanref{Formalism/LocalToGlobal/SiegelLatticeTheory/FlagStabilizerTensors/ThmSOwn.lean}{Formalism.L2.SG.ThmSOwn.hsNormSq_off}
\leanref{Formalism/LocalToGlobal/SiegelLatticeTheory/FlagStabilizerTensors/ThmSOwn.lean}{Formalism.L2.SG.ThmSOwn.hsNormSq_diag}
\begin{proof}
The isomorphism of Bessel-period data carries the Gram matrix $K_{uu'}$ of $D$ to the Gram matrix
$K^{\mathrm{set}}_{uu'}$ of the setting, the identity to the identity, and the terms of the unfolding
of $D$ with $G=0$ that are not identity terms to those of the setting; it preserves
Hilbert-Schmidt norms. By Lemma~\ref{is-unfolding-computations}, $\sharp$ carries
$K^{\mathrm{set}}_{uu'}$ to $K$. The operator of a term of the unfolding of the setting is the sum of
the $a_t$ over the finitely many terms $t$ of $T_{uu'}$ above it, and $t$ is degenerate exactly when
the term below it has $G=0$ and is not an identity term. Since $\sum_t\|a_t\|$ converges
(Lemma~\ref{is-unfolding-computations}), the sum defining $K''$ for the setting can be regrouped as
$a(T_{\mathrm{deg}})$. The maps $\sharp$ are isometries, so conjugation by them preserves
Hilbert-Schmidt norms and carries the identity to the identity.
\end{proof}

\begin{lemma}\label{is-unfolding-schur}
Fix the constants of the standing hypotheses other than $C'$ and $C_l$, and let $C',C_l\geq0$. Let
$D$ be a definite-side datum with size data that satisfy the standing hypotheses with $C'$ and $C_l$,
let a concrete evaluation of the Gram matrices of $D$ be given
(Definition~\ref{is-unfolding-evaluation}), and let $u\neq u'$ be words. Then $a_t=0$ for every
degenerate term $t$ of $T_{uu'}$.
\end{lemma}
\leanref{Formalism/LocalToGlobal/SiegelLatticeTheory/CokernelGLCounts/RankPointCounts/FixedDecomposition.lean}{Formalism.L2.SG.FixedDecomposition.kppOffC_of_steps0}
\leanref{Formalism/LocalToGlobal/SiegelLatticeTheory/CokernelGLCounts/RankPointCounts/KappaOffDiagonal.lean}{Formalism.L2.SG.KappaOffDiagonal.words_differ}
\leanref{Formalism/LocalToGlobal/SiegelLatticeTheory/CokernelGLCounts/RankPointCounts/FixedDecomposition.lean}{Formalism.L2.SG.FixedDecomposition.OOt_eq_zero0}
\leanref{Formalism/LocalToGlobal/SiegelLatticeTheory/ClosedVolumeStatements/UnipotentInLevelSubgroup.lean}{Formalism.L2.SG.Closed.LemmaVC_holds}
\leanref{Formalism/LocalToGlobal/SiegelLatticeTheory/CokernelGLCounts/RankPointCounts/FixedDecomposition.lean}{Formalism.L2.SG.FixedDecomposition.covariantRepZero_of_lemmaV}
\leanref{Formalism/LocalToGlobal/SiegelLatticeTheory/ClosedLocalOrbits/GaugeLatticeShifts/Vanish.lean}{Formalism.L2.SG.TermOp.Vanish.AvgCore}
\leanref{Formalism/LocalToGlobal/SiegelLatticeTheory/CokernelGLCounts/RankPointCounts/KappaOffDiagonal.lean}{Formalism.L2.SG.KappaOffDiagonal.SchurNFC}
\leanref{Formalism/LocalToGlobal/SiegelLatticeTheory/CokernelGLCounts/RankPointCounts/SchurNFCOf.lean}{PI.SchurPlaces.schurNFC_of}
\leanref{Formalism/SharpRestriction/CuspSpanNormBounds/BruhatCellsKronecker/NormAdditiveCharacterTypePrimeEqOne.lean}{Formalism.MainNotions.BridgeLemma.schur_lemma_unitary_distinct_params}
\leanref{Formalism/LocalToGlobal/SiegelLatticeTheory/CokernelGLCounts/RankPointCounts/KappaOffDiagonal.lean}{Formalism.L2.SG.KappaOffDiagonal.aTerm_eq_zero_of_OOt}
\begin{proof}
The standing hypotheses contain the standing numerical data of $D$, and one of these says that for
distinct words the $u$-parts of the parameters $\lambda_u(\tau)$ and $\lambda_{u'}(\tau)$ differ for
every type $\tau$. Since $u\neq u'$, these $u$-parts differ for every $\tau$.

Let $t$ be a degenerate term, given by rows $\Lambda,\Lambda'$, an isometry $\gamma$ and a point
$x=x_t$; the invariant $G$ of the term below it vanishes. If $t$ is not relevant, then $a_t=0$: the
operators of the terms that are not relevant vanish, by a computation with the unfolded series.

Let $t$ be relevant. The ties give a fundamental domain $\Omega_t$ on which the integrand of $O_t$
is integrable. Let $R_t\subset G_\infty$ be the closed subgroup of the unipotent radical of the Bessel
data of $u$ that contains $\Delta_t$ cocompactly and whose conjugate $x_\infty^{-1}R_tx_\infty$ lies in
that of $u'$, and let $\psi_u,\psi_{u'}$ be the characters of the Bessel data. Then
$W_u(ry)=\psi_u(r)W_u(y)$ for $r$ in the radical of $u$, and similarly for $u'$, so the integrand of
$O_t$ transforms under $R_t$ by the character
$\chi_t(r)=\psi_u(r)\overline{\psi_{u'}(x_\infty^{-1}rx_\infty)}$, which is trivial on $\Delta_t$. Let
$H_t\subset G_\infty$ be the pointwise stabilizer of the subspace fixed by $t$. Then $O_t$ is
covariant:
\[
 \pi_u(h)\,O_t=O_t\,\pi_{u'}(x_\infty^{-1}hx_\infty)\qquad(h\in H_t).
\]
Indeed, if $\chi_t$ is nontrivial on $R_t$, then $O_t=0$: the integral over a fundamental domain for
$\Delta_t$ of a function that transforms under $R_t$ by a character trivial on $\Delta_t$ and
nontrivial on $R_t$ vanishes. If $\chi_t$ is trivial on $R_t$, let $h\in H_t$. Since $h$ and
$x_\infty^{-1}hx_\infty$ fix the vectors of the flags, the Bessel characters are trivial on them and
$W_u(hy)=\pi_u(h)W_u(y)$, $W_{u'}(x_\infty^{-1}hx_\infty y)=\pi_{u'}(x_\infty^{-1}hx_\infty)W_{u'}(y)$.
The function $y\mapsto\langle W_{u'}(x_\infty^{-1}y),\cdot\rangle W_u(hy)$ is invariant under $R_t$,
and the substitution $y\mapsto h^{-1}y$, which respects the Haar measures of $G_\infty$ and $R_t$,
gives the displayed identity.

Since $t$ is relevant and $G=0$, the group $H_t$ contains, at an archimedean place $v_0$, a compact
unitary group of rank $k$ that acts on $V_{\pi_u}$ and $V_{\pi_{u'}}$ through their tensor factors at
$v_0$, by irreducible representations of the Harish-Chandra types given by the $u$-parts of
$\lambda_u(\tau_0)$ and $\lambda_{u'}(\tau_0)$, where $\tau_0$ is the type of $v_0$. These differ by
the first paragraph, so by Schur's lemma an operator with the displayed covariance vanishes. Hence
$O_t=0$, and every block of $a_t$ is a multiple of $O_t$, so $a_t=0$.
\end{proof}

\begin{lemma}\label{is-unfolding-degenerate-pairs}
Assume the positivity of the row types of hypothesis (5) of
Theorem~\ref{is-orthogonality-link-holds}. Fix the constants of the standing hypotheses of size data
other than the constants $C'$ and $C_l$, which bound $k$ below by $C'm^5fx^{\varepsilon}$ and the gaps
of the Harish-Chandra parameters below by $C_lr^{19}f^5x^{\varepsilon}$. There are thresholds for
$C'$ and $C_l$ such that the following holds for every definite-side datum $D$ with the ties of a
member of a family and with size data that satisfy the standing hypotheses with $C'$ and $C_l$
beyond them, every concrete evaluation of the Gram matrices of $D$
(Definition~\ref{is-unfolding-evaluation}) and every word $u$: the sum
$\sum_{(t,t')\in P(T_{\mathrm{deg}})}\|v_t\|\|v_{t'}\|$ over the pairs of degenerate terms of
$T_{uu}$ converges, and
\[
 P_{\mathrm{ev}}\,\Upsilon_1\sum_{(t,t')\in P(T_{\mathrm{deg}})}\|v_t\|\,\|v_{t'}\|
 \le\dim M_u\cdot e^{-kf\log k}.
\]
\end{lemma}
\leanref{Formalism/LocalToGlobal/SiegelLatticeTheory/FlagStabilizerTensors/RowsPos.lean}{Formalism.L2.SG.RowsPos.PairsDegCutP}
\leanref{Formalism/LocalToGlobal/SiegelLatticeTheory/FlagStabilizerTensors/RowsPos.lean}{Formalism.L2.SG.RowsPos.pairs_degree_bound}
\leanref{Formalism/LocalToGlobal/SiegelLatticeTheory/CokernelGLCounts/RowsOther.lean}{Formalism.L2.SG.RowsOther.PairsDegAmp}

\begin{lemma}\label{is-unfolding-degenerate}
Assume the positivity of the row types of hypothesis (5) of
Theorem~\ref{is-orthogonality-link-holds}, and fix the constants of the standing hypotheses other
than $C'$ and $C_l$. There are thresholds for $C'$ and $C_l$ such that, for every definite-side
datum $D$ with the ties of a member of a family and with size data that satisfy the standing
hypotheses with $C'$ and $C_l$ beyond them, every concrete evaluation of the Gram matrices of $D$
(Definition~\ref{is-unfolding-evaluation}) and every word $u$, the sum
$\sum_{t\in T_{\mathrm{deg}}}\|a_t\|$ over the degenerate terms of $T_{uu}$ converges and
\[
 \|a(T_{\mathrm{deg}})\|^2\le\dim M_u\cdot e^{-kf\log k}.
\]
\end{lemma}
\leanref{Formalism/LocalToGlobal/SiegelLatticeTheory/FlagStabilizerTensors/RowsPos.lean}{Formalism.L2.SG.RowsPos.KppDiagP}
\leanref{Formalism/LocalToGlobal/SiegelLatticeTheory/FlagStabilizerTensors/RowsPos.lean}{Formalism.L2.SG.RowsPos.kernel_diagonal_positive}
\leanref{Formalism/LocalToGlobal/SiegelLatticeTheory/FlagStabilizerTensors/RowsPos.lean}{Formalism.L2.SG.RowsPos.diagonalBound_pos_of_degreeCutoff}
\leanref{Formalism/LocalToGlobal/SiegelLatticeTheory/CokernelGLCounts/RowsOther.lean}{Formalism.L2.SG.RowsOther.kppDiag_of_pairs}
\leanref{Formalism/LocalToGlobal/SiegelLatticeTheory/CokernelGLCounts/RowsOther.lean}{Formalism.L2.SG.RowsOther.pairAbs_le_ampSum}
\leanref{Formalism/LocalToGlobal/SiegelLatticeTheory/CokernelGLCounts/RowsConc.lean}{Formalism.L2.SG.RowsConc.card_ge_three}
\begin{proof}
Take the thresholds of Lemma~\ref{is-unfolding-degenerate-pairs}. By
Lemma~\ref{is-unfolding-computations}, $T_{uu}$ is absolutely convergent, which gives the convergence
of $\sum_{t\in T_{\mathrm{deg}}}\|a_t\|$; it is row-orthogonal, of positive corank and
norm-expanding on $T_{\mathrm{deg}}$. The standing hypotheses give $q_v\geq3$ for every $v$, so
$T_{uu}$ is $\Upsilon$-bounded by $P_{\mathrm{ev}}\geq0$. By parts (2), (3) and (4) of
Lemma~\ref{is-unfolding-identities},
\[
 \|a(T_{\mathrm{deg}})\|^2\le P_{\mathrm{ev}}\Phi(T_{\mathrm{deg}})
 =P_{\mathrm{ev}}\Sigma(P(T_{\mathrm{deg}}))\le P_{\mathrm{ev}}|\Sigma|(P(T_{\mathrm{deg}})).
\]
By the Cauchy-Schwarz inequality
$|\Sigma|(P(T_{\mathrm{deg}}))\le\Upsilon_1\sum_{(t,t')\in P(T_{\mathrm{deg}})}\|v_t\|\|v_{t'}\|$, and
Lemma~\ref{is-unfolding-degenerate-pairs} bounds $P_{\mathrm{ev}}$ times the right side by
$\dim M_u\cdot e^{-kf\log k}$.
\end{proof}

\begin{lemma}\label{is-unfolding-diagonal-numerics}
Let $C_B,C_4\geq0$, and fix the constants of the standing hypotheses other than $C'$ and $C_l$. There
are thresholds for $C'$ and $C_l$ such that the following holds for every definite-side datum $D$
with size data that satisfy the standing hypotheses with $C'$ and $C_l$ beyond them and with
$C_\Gamma\geq\max(C_B,C_4)+1$, every concrete evaluation of the Gram matrices of $D$
(Definition~\ref{is-unfolding-evaluation}) and every word $u$ such that the pair $(u,u)$ satisfies
the counting hypotheses. Put $\varepsilon_L=\exp(-\tfrac13kf\log k)$,
$\varepsilon_{NG}=2^{-(k-m-3)}\exp(C_4m^5f\log b)$, $A_B=\exp(C_Bm^5f\log b)$ and
$\Gamma=\exp(C_\Gamma m^5f\log b)$. Then
\[
 \Big(\sqrt{P_{\mathrm{ev}}(A_B+2\varepsilon_L)}+\sqrt{P_{\mathrm{ev}}(\varepsilon_L+\varepsilon_{NG})}\Big)^2
 \le\Gamma.
\]
\end{lemma}
\leanref{Formalism/LocalToGlobal/SiegelLatticeTheory/CokernelGLCounts/RowsConc.lean}{Formalism.L2.SG.RowsConc.DiagNumC}
\leanref{Formalism/LocalToGlobal/SiegelLatticeTheory/CokernelGLCounts/RankPointCounts/MainEqCut.lean}{Formalism.L2.SG.Pf.diagNumC_holds}

\begin{lemma}\label{is-unfolding-diagonal}
Let $C_B,C_4\geq0$, and fix the constants of the standing hypotheses other than $C'$ and $C_l$. For a
definite-side datum $D$ with size data, a concrete evaluation of its Gram matrices
(Definition~\ref{is-unfolding-evaluation}) and a word $u$, put $d=(\dim M_u)^2/D^{*}$ and let
$\varepsilon_L$, $\varepsilon_{NG}$, $A_B$ and $\Gamma$ be as in
Lemma~\ref{is-unfolding-diagonal-numerics}. Suppose that, whenever $D$ has the ties of a member of a
family, its size data satisfy the standing hypotheses with $C'$ and $C_l$ beyond some thresholds,
and $(u,u)$ satisfies the counting hypotheses, the term data $T_{uu}$ satisfy the main-row bound
$(A_B,d)$ and the non-generic bound $(\varepsilon_{NG},d)$. If $C_\Gamma\geq\max(C_B,C_4)+1$, there
are thresholds for $C'$ and $C_l$ such that, for every such $D$, evaluation and word $u$ for which
$T_{uu}$ also satisfies the large-row and non-scalar bounds $(\varepsilon_L,d)$,
\[
 \|K^0-a(T_{\mathrm{id}})\|^2\le\Gamma(\dim M_u)^2/D^{*}.
\]
\end{lemma}
\leanref{Formalism/LocalToGlobal/SiegelLatticeTheory/FlagStabilizerTensors/RowsPos.lean}{Formalism.L2.SG.RowsPos.DiagLP}
\leanref{Formalism/LocalToGlobal/SiegelLatticeTheory/FlagStabilizerTensors/RowsPos.lean}{Formalism.L2.SG.RowsPos.diagonal_positive_final}
\leanref{Formalism/LocalToGlobal/SiegelLatticeTheory/FlagStabilizerTensors/RowsPos.lean}{Formalism.L2.SG.RowsPos.diagLP_of_rows}
\begin{proof}
Take the larger of the thresholds in the hypothesis and those of
Lemma~\ref{is-unfolding-diagonal-numerics}. Fix $D$, the evaluation and $u$. By
Lemma~\ref{is-unfolding-computations}, $T_{uu}$ is absolutely convergent, row-orthogonal, of positive
corank and norm-expanding on $T_{\mathrm{gen}}$ and on $T_{\mathrm{ng}}$, and $a(T_{\mathrm{id}})$ is
the identity; the standing hypotheses give $q_v\geq3$ for every $v$, so $T_{uu}$ is $\Upsilon$-bounded by
$P_{\mathrm{ev}}\geq0$. The ties give $d\geq0$. By hypothesis $T_{uu}$ satisfies the main-row bound
$(A_B,d)$ and the non-generic bound $(\varepsilon_{NG},d)$. Lemma~\ref{is-unfolding-pair-bound}, with
$e=a(T_{\mathrm{id}})$, $P=P_{\mathrm{ev}}$, $A=A_B$, $\varepsilon=\varepsilon_L$ and
$\varepsilon'=\varepsilon_{NG}$, gives
\[
 \|K^0-a(T_{\mathrm{id}})\|^2\le\Big(\sqrt{P_{\mathrm{ev}}(A_B+2\varepsilon_L)}
 +\sqrt{P_{\mathrm{ev}}(\varepsilon_L+\varepsilon_{NG})}\Big)^2d,
\]
and Lemma~\ref{is-unfolding-diagonal-numerics} bounds the factor before $d$ by $\Gamma$.
\end{proof}

\begin{lemma}\label{is-gram-clauses}
Assume the hypotheses of Theorem~\ref{is-orthogonality-link-holds}. Eventually along the designed
families $\mathcal A^{\mathrm{des}}$, every family has a scale beyond which every member $D$ that
satisfies the standing hypotheses of its size data, with the constants of the auxiliary-prime
condition (3) and $C_0=3$, has an unfolding (Definition~\ref{is-gram-terms}) with:
\begin{enumerate}
\item $K''_{uu'}=0$ for all words $u\ne u'$;
\item $\|K''_{uu}\|_{HS}^2\le\dim M_u\cdot x_K$ for every word $u$;
\item $\|K_{uu'}-K''_{uu'}\|_{HS}^2\le(1+\eta_0)\dim M_u\dim M_{u'}/D^{*}$ for all words $u\ne u'$;
\item $\|K_{uu}-K''_{uu}-1\|_{HS}^2\le\Gamma(\dim M_u)^2/D^{*}$ for every word $u$.
\end{enumerate}
\end{lemma}
\leanref{Formalism/LocalToGlobal/SiegelLatticeTheory/FlagStabilizerTensors/ThmSOwn.lean}{Formalism.L2.SG.ThmSOwn.ClKppOff}
\leanref{Formalism/LocalToGlobal/SiegelLatticeTheory/FlagStabilizerTensors/ThmSOwn.lean}{Formalism.L2.SG.ThmSOwn.ClKppDiag}
\leanref{Formalism/LocalToGlobal/SiegelLatticeTheory/FlagStabilizerTensors/ThmSOwn.lean}{Formalism.L2.SG.ThmSOwn.ClOff}
\leanref{Formalism/LocalToGlobal/SiegelLatticeTheory/FlagStabilizerTensors/ThmSOwn.lean}{Formalism.L2.SG.ThmSOwn.ClDiag}
\leanref{Formalism/LocalToGlobal/SiegelLatticeTheory/FlagStabilizerTensors/ThmSOwn.lean}{Formalism.L2.SG.ThmSOwn.HypOwn}
\leanref{Formalism/LocalToGlobal/SiegelLatticeTheory/FlagStabilizerTensors/ThmSOwn.lean}{Formalism.L2.SG.ThmSOwn.hypOwnE_kppOff}
\leanref{Formalism/LocalToGlobal/SiegelLatticeTheory/FlagStabilizerTensors/LinkSPos.lean}{Formalism.L2.SG.ThmSOwn.hypOwnE_kppDiagP}
\leanref{Formalism/LocalToGlobal/SiegelLatticeTheory/FlagStabilizerTensors/ThmSOwn.lean}{Formalism.L2.SG.ThmSOwn.hypOwnS_off}
\leanref{Formalism/LocalToGlobal/SiegelLatticeTheory/FlagStabilizerTensors/LinkSPos.lean}{Formalism.L2.SG.ThmSOwn.hypOwnL_diagP}
\leanref{Formalism/LocalToGlobal/SiegelLatticeTheory/FlagStabilizerTensors/ThmSOwn.lean}{Formalism.L2.SG.ThmSOwn.kpp_zero}
\begin{proof}
Fix the constants of the standing hypotheses other than $C'$ and $C_l$: those of the auxiliary-prime
condition (3) and $C_0=3$. Eventually along $\mathcal A^{\mathrm{des}}$, the constants $C_\Gamma$,
$C'$ and $C_l$ of a family exceed any given thresholds that depend only on the fixed constants. Take
$C',C_l\geq29$, the constants $C_B,C_4\geq0$ that hypothesis (5) gives for the main-term identity and
the row bound for the non-generic terms, $C_\Gamma\geq\max(C_B,C_4)+1$, and the thresholds of
hypothesis (5) and of Lemmas~\ref{is-unfolding-degenerate} and~\ref{is-unfolding-diagonal}. With
$C',C_l\geq29$ the standing hypotheses give the gap hypothesis, so hypothesis (1) gives the decay of
$W_u$ and the integrability of the archimedean matrix coefficient. With hypothesis (2), this gives
the finiteness of the spectral data, and with hypothesis (3) and the local ties of the spectral
selection it shows that, eventually along $\mathcal A^{\mathrm{des}}$, every family has a scale beyond
which every member $D$ that satisfies the standing hypotheses has a concrete evaluation of its Gram
matrices (Definition~\ref{is-unfolding-evaluation}), given by the spectral selection. Take the scale
of a family beyond this one and beyond the scales of hypotheses (6) and (7), and fix a member $D$
beyond it that satisfies the standing hypotheses, with its unfolded term data $T_{uu'}$.

The standing hypotheses give: (a) $q_v\geq3$ for every place $v$ of the reductions; (b) for every
word $u$, the pair $(u,u)$ satisfies the counting hypotheses. By
Lemma~\ref{is-unfolding-computations} every $T_{uu'}$ is absolutely convergent and
identity-degenerate, so $K-a(T_{\mathrm{deg}})=K^0$ by Lemma~\ref{is-unfolding-identities}(1).
Lemma~\ref{is-unfolding-transport} then gives
\[
 \|K''_{uu'}\|_{HS}=\|a(T_{\mathrm{deg}})\|,\qquad\|K_{uu'}-K''_{uu'}\|_{HS}=\|K^0\|,\qquad
 \|K_{uu}-K''_{uu}-1\|_{HS}=\|K^0-a(T_{\mathrm{id}})\|,
\]
the last for $T_{uu}$.

(1) Let $u\ne u'$. By Lemma~\ref{is-unfolding-schur}, $a_t=0$ for every degenerate term $t$ of
$T_{uu'}$, so $a(T_{\mathrm{deg}})=0$ and $K''_{uu'}=0$.

(2) Hypothesis (5) includes the positivity of the row types, so Lemma~\ref{is-unfolding-degenerate}
gives $\|K''_{uu}\|_{HS}^2=\|a(T_{\mathrm{deg}})\|^2\le\dim M_u\cdot e^{-kf\log k}=\dim M_u\cdot x_K$.

(3) Let $u\ne u'$. The operator $K^0$ of $T_{uu'}$ is the leading kernel term $K^0_{uu'}$ of
hypothesis (7), so
$\|K_{uu'}-K''_{uu'}\|_{HS}^2=\|K^0_{uu'}\|^2\le(1+b^{-6})\dim M_u\dim M_{u'}/D^{*}$, and $b^{-6}=\eta_0$.

(4) Let $u$ be a word and $d=(\dim M_u)^2/D^{*}$. By (b) the pair $(u,u)$ satisfies the counting
hypotheses, and $T_{uu}$ is absolutely convergent, so hypothesis (6) says that $T_{uu}$ satisfies the
large-row and non-scalar bounds $(\varepsilon_L,d)$, with $\varepsilon_L=\exp(-\frac13kf\log k)$.
The main-term identity and the row bound for the non-generic terms of hypothesis (5) are the
main-row bound $(\exp(C_Bm^5f\log b),d)$ and the non-generic bound
$(2^{-(k-m-3)}\exp(C_4m^5f\log b),d)$ for $T_{uu}$. Since $C_\Gamma\geq\max(C_B,C_4)+1$,
Lemma~\ref{is-unfolding-diagonal} gives
$\|K_{uu}-K''_{uu}-1\|_{HS}^2=\|K^0-a(T_{\mathrm{id}})\|^2\le\Gamma(\dim M_u)^2/D^{*}$.
\end{proof}

\begin{lemma}\label{is-gram-combination}
Let $D$ be a definite-side datum with the Bessel-period data and an unfolding that satisfies (1)
to (4) of Lemma~\ref{is-gram-clauses}, with constants $\eta_0$, $x_K$ and $\Gamma>1$
(Definition~\ref{is-gram-terms}). Let $C_0\geq1$ and $c_7>0$, and suppose that
$D_{\mathrm{all}}\le(1+\eta_1)D^{*}$ with $\eta_1=\exp(-c_7rf\log r)$. If
$C_0^{-1}\le\dim M_u\,\Gamma^2/D^{*}\le C_0$ for every word $u$, then the orthogonality bounds hold at
$D$ (Definition~\ref{is-orthogonality-link}) with $\eta_0$, $\eta_1$ and
$e_0=2C_0/\Gamma+2x_K$. Suppose moreover that $5\le k\le r$, $f\geq1$, $8C_\Gamma m^5fx^{\varepsilon}\le
k$, and that $y\geq1$ satisfies $x^{\varepsilon}=y^8$, $(2C_0+2)e^{-C_\Gamma y^8}\le e^{-2y}$ and
$y\le8c_7C_\Gamma y^8$. Then $0\le e_0$ and $\eta_0,\eta_1,\sqrt{e_0}\le e^{-y}$.
\end{lemma}
\leanref{Formalism/SharpRestriction/CuspSpanNormBounds/StandingHypothesesHoldFamily.lean}{Formalism.MainNotions.S.linkS_concl}
\leanref{Formalism/SharpRestriction/CuspSpanNormBounds/StandingHypothesesHoldFamily.lean}{Formalism.MainNotions.S.diag_unproj}
\leanref{Formalism/SharpRestriction/CuspSpanNormBounds/StandingHypothesesHoldFamily.lean}{Formalism.MainNotions.S.eF_nonneg_and_le_exp_neg}
\begin{proof}
For $u\ne u'$, clause (1) gives $K_{uu'}=K_{uu'}-K''_{uu'}$, and clause (3) is the first orthogonality
bound. For the diagonal write $K_{uu}-1=(K_{uu}-K''_{uu}-1)+K''_{uu}$; since
$\|A+B\|_{HS}^2\le2\|A\|_{HS}^2+2\|B\|_{HS}^2$, clauses (4) and (2) give
$\|K_{uu}-1\|_{HS}^2\le2\Gamma(\dim M_u)^2/D^{*}+2x_K\dim M_u$. The window gives
$\dim M_u/D^{*}\le C_0/\Gamma^2$, so $\Gamma(\dim M_u)^2/D^{*}\le(C_0/\Gamma)\dim M_u$, and the second
bound holds with $e_0=2C_0/\Gamma+2x_K$. The third is the hypothesis on $D_{\mathrm{all}}$.

For the sizes put $\Lambda=C_\Gamma m^5fx^{\varepsilon}$, so $\Gamma=e^{\Lambda}$ and
$\Lambda\geq C_\Gamma y^8$ because $m,f\geq1$. Then $8\Lambda\le k$. Since $k\geq5$, $\log k\geq1$ and
$\log r\geq1$, so $kf\log k\geq k\geq8\Lambda\geq C_\Gamma y^8$ and $rf\log r\geq k\geq8C_\Gamma y^8$.
Hence $e_0=2C_0e^{-\Lambda}+2e^{-kf\log k}\le(2C_0+2)e^{-C_\Gamma y^8}\le e^{-2y}$, so
$\sqrt{e_0}\le e^{-y}$, and $e_0\geq0$. Next $\eta_0=b^{-6}=e^{-6x^{\varepsilon}}=e^{-6y^8}\le e^{-y}$
because $y\le y^8$. Finally $c_7rf\log r\geq8c_7C_\Gamma y^8\geq y$, so $\eta_1\le e^{-y}$.
\end{proof}

\begin{theorem}\label{is-orthogonality-link-holds}
Let $\mathcal A$ be a class of families with a twist selection, which chooses twisting data at every
member of every family of $\mathcal A$ beyond a threshold scale, and a spectral selection, which
realizes the columns at those members as Poincaré series in a space of automorphic forms. Assume:
\begin{enumerate}
\item the decay hypothesis for $W_u$ (Definition~\ref{is-wu-decay});
\item finite dimensionality of rows: for every automorphic setting, the image of a convolution
  operator by a locally constant compactly supported Hecke function, composed with integration
  against a square-integrable discrete-series matrix coefficient at the archimedean places, is
  finite-dimensional;
\item eventually along the designed families $\mathcal A^{\mathrm{des}}$, the auxiliary-prime condition holds
  with constants $(C_{YL}+1,\frac14,3,\frac12,C_{Tu},3)$: under the counting hypotheses, every
  word $u$ satisfies $1\le T_u\,\kappa_S^{2r_S}$, where $r_S=m/2-1$ and $\kappa_S>0$ is an invariant
  of the host field;
\item counting constants are chosen for every family, with $C\ge0$ and $A\ge2(C_{YL}+1)$;
\item the positivity of the row types, the row bound for the non-generic terms and the main-term
  identity for the diagonal rows;
\item for every choice of constants, eventually along $\mathcal A^{\mathrm{des}}$, for every pair of
  words $u,u'$ under the counting hypotheses whose terms converge absolutely, the row of the pair
  $(u,u')$ is large and not scalar;
\item the selected orthogonality bound: eventually along $\mathcal A^{\mathrm{des}}$, for all words
  $u\ne u'$, the leading kernel term of the Poincaré series satisfies
  \[
   \|K^{0}_{uu'}\|^{2}\le(1+b^{-6})\dim M_u\dim M_{u'}/D^{*};
  \]
\item eventually along $\mathcal A^{\mathrm{des}}$, the point counts hold with the chosen counting
  constants;
\item at every selected member, the volumes are independent of the signature and the class items
  exist for the rank of the first indefinite space.
\end{enumerate}
Then, eventually along $\mathcal A^{\mathrm{des}}$, every family that satisfies the size conditions on
$T_u$ satisfies the orthogonality link (Definition~\ref{is-orthogonality-link}) with constant $3$.
\end{theorem}
\leanref{Development/AutomorphicArithmetic/ThetaCorrespondence/ThetaPairLifts/ThirdPlaceSeries/LinkAltAltBis.lean}{GlobalTheory.Top.link_alt_alt_bis}
\leanref{Formalism/LocalToGlobal/SiegelLatticeTheory/FlagStabilizerTensors/LinkSPos.lean}{Formalism.L2.SG.ThmSOwn.theoremS_transport_link_zero}
\leanref{Formalism/LocalToGlobal/SiegelLatticeTheory/FlagStabilizerTensors/LinkSPos.lean}{Formalism.L2.SG.ThmSOwn.theoremS_link_concrete}
\begin{proof}
The proof takes $C_0=3$, which satisfies $1\le C_0$ and $\exp(1)\le C_0$.

The decay hypothesis for $W_u$ and finite dimensionality of rows give the finiteness of the spectral
data: at every selected member the Poincaré series converge, are square integrable, and span a
finite-dimensional space. With the local ties of the spectral selection and the auxiliary-prime condition (3),
this makes the selection admissible, so the Gram matrices $K_{uu'}$ have a concrete evaluation as
sums over the terms of the unfolded Poincaré series. The rows of every member are tied to an
automorphic realization of the row space, through the instance ties of the spectral selection, the
volume tie of the theta realization and hypothesis (9), and the invariant $q$ of the automorphic
setting is at least $1$.

Lemma~\ref{is-gram-clauses} gives, eventually along $\mathcal A^{\mathrm{des}}$, at every member
beyond a scale that satisfies the standing hypotheses, the four clauses for the leading parts
$K''_{uu'}$ of the Gram matrices. The residues of the evaluation are eventually small, and with the
rows tied to the automorphic realization they give the standing hypotheses of the members at large
scales. The point counts (8), with the counting constants (4), give
$D_{\mathrm{all}}\le(1+\exp(-c_7rf\log r))D^{*}$ for a constant $c_7>0$ of the counting constants.

Let $y=x^{\varepsilon/8}$. For large $x$, $y\geq1$, $(2C_0+2)e^{-C_\Gamma y^8}\le e^{-2y}$ and
$y\le8c_7C_\Gamma y^8$, and the standing hypotheses give $5\le k\le r$, $f\geq1$ and
$8C_\Gamma m^5fx^{\varepsilon}\le k$. At a member where $C_0^{-1}\le(\dim M_u)\Gamma^2/D^{*}\le C_0$ for
every word $u$, Lemma~\ref{is-gram-combination} gives the orthogonality bounds with
$\eta_0,\eta_1,\sqrt{e_0}\le e^{-y}=e^{-x^{\varepsilon/8}}$. At any other member the orthogonality
bounds hold with $\eta_0=\eta_1=e_0=0$, since their hypothesis fails. This is the orthogonality link
with constant $3$.
\end{proof}

\section{The row space}\label{section-is-row-space}

\begin{definition}\label{is-irreducibility-link}
Let $D$ be a definite-side datum (Definition~\ref{th-definite-datum}). Its row space has dimension
$D_{\mathrm{rows}}=(1-\delta)D_{\mathrm{all}}$, where $D_{\mathrm{all}}$ is the total number of rows
and $\delta\le1$ is the \emph{row deficit}. A family $\mathcal F$ of exponent $\varepsilon$ with
constant $J\ge1$ (Definition~\ref{pd-family}) satisfies the \emph{irreducibility link} if there is
$x_0$ such that for every scale $x\ge x_0$ and every member,
\[
 0\le\delta\le\exp\big(-x^{\varepsilon_z/7}\big),\qquad \varepsilon_z=\varepsilon/J.
\]
\end{definition}
\leanref{Formalism/SharpRestriction/GlobalTheoremInputs.lean}{Formalism.MainNotions.LinkIRR}
\leanref{Formalism/LocalRingCardinality/DataRows.lean}{Formalism.BasicNotions.Rows.Data}

\begin{definition}\label{is-row-labelling}
Let $D$ be a definite-side datum (Definition~\ref{th-definite-datum}) with ranks $f$ and
$r=k+m-1$, level primes $g_0,g_1$, and places $v$ of the host over $g_k$ for $k=0,1$. A
\emph{label} at $g_k$ is a monic polynomial of degree $r$ over $\mathbf F_{g_k}$ with nonzero
constant term. A \emph{labelling of all rows} of $D$ consists of: an automorphic realization of the
full row space; for every row $\tau\le D_{\mathrm{all}}$, every $k$ and every place $v$ over $g_k$, a
label $\mathrm{lab}_k(\tau,v)$; an injection of the $D_{\mathrm{rows}}$ good rows into all rows;
idempotents $e_{k,v}$ in the group algebra of $GL_r(\mathbf F_{g_k})$; and, for every set $S$ of
tuples of labels, a Hecke function $e_S$ with its convolution operator, such that $e_S$ is the
function of the realization when $S$ is everything, and for every $k$ and finite set $S$ of tuples
at $g_k$, $\operatorname{Re}e_S(1)=\operatorname{Re}e(1)\sum_{t\in S}\prod_v\mathrm{wt}(e_{k,v},t_v)$,
where $\mathrm{wt}(e,f)$ is the weight of the label $f$ in $e$. Let $\dim(S)$ be the dimension of the
image of the operator of $e_S$ on the realization, and $D^{*}(e)$ its expected dimension. The
labelling
\begin{enumerate}
\item \emph{counts labels and is injective on good labels} if $\dim(S)$ is the number of rows whose
  tuple of labels lies in $S$, for every $S$; the injection of good rows is injective; the labels of
  a good row are those of $D$; and a row has a good tuple of labels exactly when it is a good row;
\item \emph{joins all rows} if the rows of the realization are tied to those of $D$ and the levels
  are sharp;
\item has the \emph{row count} with constant $c>0$ if, for $\varepsilon_C=\exp(-c\,r\,f\log r)$, we
  have $0\le\varepsilon_C<1$, $D^{*}(e)>0$, and
  $|\dim(S)-D^{*}(e_S)|\le\varepsilon_CD^{*}(e_S)$ for every $k$ and every set $S$ of tuples that
  depends only on the labels at $g_k$.
\end{enumerate}
\end{definition}
\leanref{Formalism/SharpRestriction/LevelLetterCocycles/ArithmeticCylinders.lean}{Formalism.MainNotions.ArithmeticCylinders.AllRows}
\leanref{Formalism/SharpRestriction/LevelLetterCocycles/ArithmeticCylinders.lean}{Formalism.MainNotions.ArithmeticCylinders.LabelCountGoodInjective}
\leanref{Formalism/LocalToGlobal/HilbertRowConditions/AllJoin.lean}{Formalism.L2.GlobalTheoryTheory.AllJoin.AllJoinA}
\leanref{Formalism/LocalToGlobal/GlobalTheory/IrreducibleOfCountSizes.lean}{Formalism.L2.GlobalTheoryTheory.ConcN.CountAllRows}

\begin{theorem}\label{is-irreducibility-link-holds}
Assume Green's parametrization of the cuspidal characters of $GL_d(\mathbf F_g)$ by monic
irreducible polynomials of degree $d$ with nonzero constant term
(Corollary~\ref{lit-green-labelled}). Let $\mathcal A$ be a class of families with a twist selection,
and let $c>0$. Suppose that, eventually along the designed families $\mathcal A^{\mathrm{des}}$,
every family has a scale $x_0$ such that every member at every scale $x\ge x_0$ has a labelling of
all rows that counts labels and is injective on good labels, joins all rows, and has the row count
with constant $c$ (Definition~\ref{is-row-labelling}). Then, eventually along
$\mathcal A^{\mathrm{des}}$, every family satisfies the irreducibility link
(Definition~\ref{is-irreducibility-link}).
\end{theorem}
\leanref{Development/AutomorphicArithmetic/ThetaCorrespondence/ThetaPairLifts/ThirdPlaceSeries/LinkIrr.lean}{GlobalTheory.Top.Irr.closed_linkIRR_of_join_count}
\leanref{Development/AbelianVarietyTheory/Proofs/GlobalInterpolation/LinkIRROfJoinEx.lean}{GlobalTheory.Top.Irr.linkIRR_of_join_ex}
\leanref{Formalism/LocalToGlobal/GlobalTheory/ConcNP.lean}{Formalism.L2.GlobalTheoryTheory.ConcN.linkIRR_of_allRows}
\begin{proof}
Eventually along $\mathcal A^{\mathrm{des}}$ the exponent satisfies $\varepsilon\le1$. Fix a family
in $\mathcal A^{\mathrm{des}}$ and a scale $x$ beyond $x_0$, beyond the threshold of the twist
selection, beyond $x_1+1$ and beyond $4$, and a member with its labelling.

The twist selection at the member gives types for the realization. Green's parametrization shows
that the label weights at these types are admissible: for every place $v$ over $g_k$ the weights
$\mathrm{wt}(e_{k,v},f)$ are nonnegative, a squarefree label has weight within a factor
$1\pm O(r/g_k)$ of $(1-1/g_k)g_k^{-r}$, and the non-squarefree labels have total weight at most
$3/g_k$. The size recursion holds at the member: $g_k\ge\exp(8x^{\varepsilon_z})$,
$x^{2}\le f\le x^{6}$, $m\le x^{6}$, $100\le r\le x^{6}$ and $\min(q,d)\ge\exp(x^{\varepsilon_z/6})$,
where $q,d$ are the auxiliary primes of the host.

Since the labelling joins all rows and the label weights are admissible, the share of the bad tuples
at $g_k$ is at most
\[
 W_{g_k}=\frac{3f}{g_k}+\frac{23f\log g_k}{\sqrt{g_k}}
 +2fm\Big(\frac{1}{4\min(q,d)}+\frac{1+\log r}{g_k-1}\Big)+\frac r2\Big(\frac78\Big)^{f}
 +e^{-f/(5\log r)},
\]
and this weight bound holds at the member because $x^{\varepsilon_z}\ge3/2$. A row is bad exactly
when its tuple of labels is bad at $g_0$ or at $g_1$, so the counting of labels and the row count
bound the proportion of bad rows by $\frac{1+\varepsilon_C}{1-\varepsilon_C}(W_{g_0}+W_{g_1})$.
Under the size recursion each $W_{g_k}$ is at most $x^{21}\exp(-x^{\varepsilon_z/6})$ and
$\varepsilon_C\le1/2$, so for large $x$ the proportion is at most $\exp(-x^{\varepsilon_z/7})$.
The proportion of bad rows is the row deficit $\delta$, and $\delta\ge0$ because the good rows
inject into all rows.
\end{proof}
